\documentclass{amsbook}[openany]
\renewcommand\thesection{\thechapter.\arabic{section}}
\renewcommand\thepart{\arabic{part}}
\usepackage[T1]{fontenc}
\usepackage{amsmath,amssymb}
\usepackage{textcomp}
\usepackage{fancyvrb}
\usepackage{upquote}
%% ---- Kindle build, part 1 of 2: page geometry -----------------------------
%% manual-kindle.tex defines \VNBkindle and then \input's this file, so both
%% kindle blocks are inert in the ordinary build.  geometry goes here, BEFORE
%% hyperref, which is the supported order.  Part 2 (the code font) is below,
%% after the verbatim environments it overrides.
\ifdefined\VNBkindle
  \usepackage[paperwidth=105mm,paperheight=140mm,
              top=5mm,bottom=5mm,left=4mm,right=4mm]{geometry}
\fi
%% INDEX.  No \usepackage{makeidx}: amsbook already provides \see, \seealso,
%% the theindex environment and \printindex (amsbook.cls:1329,1190,1331), and
%% makeidx re-\newcommand's \see -- "Command \see already defined", a fatal
%% error.  \makeindex itself is in the LaTeX kernel, so this line is all that
%% is needed.  Entries come from the generators; see \printindex at the end.
\makeindex
\usepackage{hyperref}
\hypersetup{pdftitle={VNB Proof Checker: User Manual},pdfauthor={Claude (Anthropic), with design specifications, technical direction and review by F. J. Thayer}}
\usepackage[protrusion=true,expansion=false]{microtype}
\usepackage{array}
\usepackage{booktabs}
\usepackage{longtable}
\usepackage{enumitem}
%% List style, set ONCE here rather than per-list.  The manual had accumulated
%% six different labels for the same thing -- $\circ$ (36 lists, in three
%% spellings), $-$ (13), \textendash (4), $\ast$ (3), $\checkmark$ (1) -- each
%% carrying its own bracketed option block, plus two bare lists that fell back
%% to the amsbook bullet.  Levels are now distinguished by INDENTATION, with a
%% dash rather than a bullet at both: \textendash and not $-$, because $-$ is a
%% math minus, set in math italic at math spacing.
\setlist[itemize,1]{label=\textendash, leftmargin=0.85cm, itemsep=2mm}
\setlist[itemize,2]{label=\textendash, leftmargin=0.85cm, itemsep=1mm,
                    topsep=1mm}
\newtheorem{thm}{Theorem}[chapter]
\newtheorem{prop}[thm]{Proposition}
\newtheorem{cor}[thm]{Corollary}
\newtheorem{lem}[thm]{Lemma}
\newtheorem{claim}{Claim}

\theoremstyle{definition}
\newtheorem{dfn}{Definition}[chapter]
\newtheorem{example}{Example}[chapter]
\theoremstyle{remark}
\newtheorem{remark}{Remark}[chapter]
\theoremstyle{remark}
\pdfglyphtounicode{visiblespace}{0020}
\pdfgentounicode=1

%% Code blocks are set off by INDENTATION, not by a frame.  VNBcode used to be
%% Verbatim with frame=single (91 boxes) and VNBsnip the identical environment
%% without the frame (49 uses); nothing in the text distinguished them, so the
%% choice between box and no-box across 140 blocks was editorial noise.  They
%% are now one style.  The indent matches the lists' leftmargin (0.85cm), so
%% code and list items hang on the same left edge.
\DefineVerbatimEnvironment{VNBcode}{Verbatim}{xleftmargin=0.85cm,fontsize=\footnotesize}
\DefineVerbatimEnvironment{VNBsnip}{Verbatim}{xleftmargin=0.85cm,fontsize=\footnotesize}

%% ---- Kindle build, part 2 of 2: the code font -----------------------------
%% The widest verbatim line in the manual is 87 columns, and a code block is
%% precisely what a reflow tool (k2pdfopt and friends) mangles.  So the code
%% font is shrunk until 87 columns fit the text block instead; \VNBkindlecode
%% is set in manual-kindle.tex, where it can be retuned without touching this
%% file.  Prose keeps the ordinary body size.
\ifdefined\VNBkindle
  %% Same unframed style as the ordinary build, at the Kindle code size and a
  %% narrower indent -- 0.85cm of a 97 mm measure is a lot to give away.
  \DefineVerbatimEnvironment{VNBcode}{Verbatim}
    {xleftmargin=3mm,fontsize=\VNBkindlecode}
  \DefineVerbatimEnvironment{VNBsnip}{Verbatim}
    {xleftmargin=3mm,fontsize=\VNBkindlecode}
  \setlist[itemize,1]{label=\textendash, leftmargin=5mm, itemsep=1mm}
  \setlist[itemize,2]{label=\textendash, leftmargin=5mm, itemsep=0.5mm,
                      topsep=0.5mm}
  %% A 97 mm measure is narrower than any identifier this manual likes to name.
  %% `hyphenat[htt]' lets a long \texttt token break; the emergency stretch lets
  %% TeX buy a loose line rather than run into the margin; and the tables, which
  %% were laid out for a 345 pt measure, are set at \VNBkindletab.
  \usepackage[htt]{hyphenat}
  \usepackage{etoolbox}
  \setlength{\emergencystretch}{3em}
  \AtBeginEnvironment{tabular}{\VNBkindletab}
  %% BEFORE, not AtBegin, so the size covers the table's header rows too.
  \BeforeBeginEnvironment{longtable}{\VNBkindletab}
  \setlength{\tabcolsep}{2pt}
  %% Displayed formulas: the manual's are long enough that a dozen of them run
  %% clean off a 105 mm page -- not merely into the margin, off the paper, with
  %% the right-hand side of the equation gone.  So \[ ... \] is boxed, measured,
  %% and scaled down only when it does not fit.  A display that fits is centred
  %% and untouched, so nothing changes but the offenders.
  \usepackage{graphicx}
  \newsavebox{\VNBdispbox}
  \renewcommand{\[}{\par\smallskip\begin{lrbox}{\VNBdispbox}$\displaystyle}
  \renewcommand{\]}{$\end{lrbox}%
    \noindent
    \ifdim\wd\VNBdispbox>\linewidth
      \resizebox{\linewidth}{!}{\usebox{\VNBdispbox}}%
    \else
      \hbox to \linewidth{\hfil\usebox{\VNBdispbox}\hfil}%
    \fi
    \par\smallskip}
\fi

\title{\textbf{VNB Proof Checker}\\[4pt]
       \large User Manual}
\author{Claude (Anthropic)\\[6pt]
  {\normalfont\normalsize the author of the system and of this manual,\\
   with design specifications, technical direction and review\\
   by F.~J.~Thayer}}
\date{}

%% The list of kernel operations and its counts, generated from the registry the
%% startup check uses (gen-kernel-ops.py).  Defines \VNBkernelopcount, the four
%% per-kind counts and \VNBkerneloptable.
%% GENERATED by gen-kernel-ops.py from *kernel-rule-tags* (tactics-help.scm).  Do not edit.
\newcommand{\VNBkernelopcount}{65}
\newcommand{\VNBaxiomcount}{199}
\newcommand{\VNBkernelopcountlogic}{41}
\newcommand{\VNBkernelopcountschema}{13}
\newcommand{\VNBkernelopcountrewrite}{4}
\newcommand{\VNBkernelopcountoracle}{7}
\newcommand{\VNBkerneloptable}{%
\begin{longtable}{@{}p{0.41\textwidth}>{\raggedright\arraybackslash}p{0.45\textwidth}>{\raggedleft\arraybackslash}p{0.08\textwidth}@{}}
\textit{operation} & \textit{what it does} & \textit{figure}\\[2pt]
\endhead
\multicolumn{3}{@{}l}{\textbf{Logic, equality, induction and the structure eliminators} (41)}\\[2pt]
\hyperref[kop:and-intro]{\mbox{\texttt{and-intro}}}\index{and-intro@\texttt{and-intro} (kernel operation)} & goal AND: prove both (di) & p.~\pageref{kop:and-intro}\\
\hyperref[kop:or-intro-left]{\mbox{\texttt{or-intro-left}}}\index{or-intro-left@\texttt{or-intro-left} (kernel operation)} & goal OR: prove the left disjunct (oi-l) & p.~\pageref{kop:or-intro-left}\\
\hyperref[kop:or-intro-right]{\mbox{\texttt{or-intro-right}}}\index{or-intro-right@\texttt{or-intro-right} (kernel operation)} & goal OR: prove the right disjunct (oi-r) & p.~\pageref{kop:or-intro-right}\\
\hyperref[kop:implies-intro]{\mbox{\texttt{implies-intro}}}\index{implies-intro@\texttt{implies-intro} (kernel operation)} & goal IMPLIES: assume the antecedent (di) & p.~\pageref{kop:implies-intro}\\
\hyperref[kop:not-intro]{\mbox{\texttt{not-intro}}}\index{not-intro@\texttt{not-intro} (kernel operation)} & goal NOT: assume it, prove FALSITY (di) & p.~\pageref{kop:not-intro}\\
\hyperref[kop:iff-intro]{\mbox{\texttt{iff-intro}}}\index{iff-intro@\texttt{iff-intro} (kernel operation)} & goal IFF: prove both directions (di) & p.~\pageref{kop:iff-intro}\\
\hyperref[kop:forall-intro]{\mbox{\texttt{forall-intro}}}\index{forall-intro@\texttt{forall-intro} (kernel operation)} & goal FORALL: introduce an eigenvariable (di) & p.~\pageref{kop:forall-intro}\\
\hyperref[kop:forsome-intro]{\mbox{\texttt{forsome-intro}}}\index{forsome-intro@\texttt{forsome-intro} (kernel operation)} & goal FORSOME: supply a witness (ew) & p.~\pageref{kop:forsome-intro}\\
\hyperref[kop:truth-intro]{\mbox{\texttt{truth-intro}}}\index{truth-intro@\texttt{truth-intro} (kernel operation)} & goal TRUTH.  UNREACHABLE: no tactic calls pi-truth! & p.~\pageref{kop:truth-intro}\\
\hyperref[kop:and-elim]{\mbox{\texttt{and-elim}}}\index{and-elim@\texttt{and-elim} (kernel operation)} & split a conjunctive assumption (ai) & p.~\pageref{kop:and-elim}\\
\hyperref[kop:or-elim]{\mbox{\texttt{or-elim}}}\index{or-elim@\texttt{or-elim} (kernel operation)} & case-split on a disjunctive assumption (ai) & p.~\pageref{kop:or-elim}\\
\hyperref[kop:not-elim]{\mbox{\texttt{not-elim}}}\index{not-elim@\texttt{not-elim} (kernel operation)} & assumption NOT p with p in context closes the goal (ai) & p.~\pageref{kop:not-elim}\\
\hyperref[kop:iff-elim]{\mbox{\texttt{iff-elim}}}\index{iff-elim@\texttt{iff-elim} (kernel operation)} & an IFF assumption becomes its two implications (ai) & p.~\pageref{kop:iff-elim}\\
\hyperref[kop:forsome-elim]{\mbox{\texttt{forsome-elim}}}\index{forsome-elim@\texttt{forsome-elim} (kernel operation)} & open an existential assumption at a fresh eigenvariable (ai) & p.~\pageref{kop:forsome-elim}\\
\hyperref[kop:forall-elim]{\mbox{\texttt{forall-elim}}}\index{forall-elim@\texttt{forall-elim} (kernel operation)} & instantiate a universal (inst, inst+, fact) & p.~\pageref{kop:forall-elim}\\
\hyperref[kop:assumption]{\mbox{\texttt{assumption}}}\index{assumption@\texttt{assumption} (kernel operation)} & the goal is in the context (ass) & p.~\pageref{kop:assumption}\\
\hyperref[kop:theorem-assumption]{\mbox{\texttt{theorem-assumption}}}\index{theorem-assumption@\texttt{theorem-assumption} (kernel operation)} & cite an installed theorem (ta, fact) & p.~\pageref{kop:theorem-assumption}\\
\hyperref[kop:cut]{\mbox{\texttt{cut}}}\index{cut@\texttt{cut} (kernel operation)} & introduce a lemma (cut, have!) & p.~\pageref{kop:cut}\\
\hyperref[kop:weakening]{\mbox{\texttt{weakening}}}\index{weakening@\texttt{weakening} (kernel operation)} & drop an assumption (wk) & p.~\pageref{kop:weakening}\\
\hyperref[kop:detach]{\mbox{\texttt{detach}}}\index{detach@\texttt{detach} (kernel operation)} & modus ponens on a context implication (detach!) & p.~\pageref{kop:detach}\\
\hyperref[kop:backchain]{\mbox{\texttt{backchain}}}\index{backchain@\texttt{backchain} (kernel operation)} & reduce the goal to an implication's antecedent (bc, bc*) & p.~\pageref{kop:backchain}\\
\hyperref[kop:proof-by-contradiction]{\mbox{\texttt{proof-by-contradiction}}}\index{proof-by-contradiction@\texttt{proof-by-contradiction} (kernel operation)} & assume the negation, prove FALSITY (pbc) & p.~\pageref{kop:proof-by-contradiction}\\
\hyperref[kop:eq-subst]{\mbox{\texttt{eq-subst}}}\index{eq-subst@\texttt{eq-subst} (kernel operation)} & Leibniz rewriting by a context equation (subst) & p.~\pageref{kop:eq-subst}\\
\hyperref[kop:reflexivity]{\mbox{\texttt{reflexivity}}}\index{reflexivity@\texttt{reflexivity} (kernel operation)} & t = t, requiring t to be DEFINED (rfl) & p.~\pageref{kop:reflexivity}\\
\hyperref[kop:quasi-reflexivity]{\mbox{\texttt{quasi-reflexivity}}}\index{quasi-reflexivity@\texttt{quasi-reflexivity} (kernel operation)} & t == t, unconditionally (qrfl) & p.~\pageref{kop:quasi-reflexivity}\\
\hyperref[kop:if-true]{\mbox{\texttt{if-true}}}\index{if-true@\texttt{if-true} (kernel operation)} & IF-term, true branch & p.~\pageref{kop:if-true}\\
\hyperref[kop:if-false]{\mbox{\texttt{if-false}}}\index{if-false@\texttt{if-false} (kernel operation)} & IF-term, false branch & p.~\pageref{kop:if-false}\\
\hyperref[kop:cartesian-intro]{\mbox{\texttt{cartesian-intro}}}\index{cartesian-intro@\texttt{cartesian-intro} (kernel operation)} & membership in a CARTESIAN product (ci) & p.~\pageref{kop:cartesian-intro}\\
\hyperref[kop:cartesian-elim]{\mbox{\texttt{cartesian-elim}}}\index{cartesian-elim@\texttt{cartesian-elim} (kernel operation)} & consume one, componentwise (ce) & p.~\pageref{kop:cartesian-elim}\\
\hyperref[kop:tuples-intro]{\mbox{\texttt{tuples-intro}}}\index{tuples-intro@\texttt{tuples-intro} (kernel operation)} & membership in TUPLES (ti) & p.~\pageref{kop:tuples-intro}\\
\hyperref[kop:tuples-elim]{\mbox{\texttt{tuples-elim}}}\index{tuples-elim@\texttt{tuples-elim} (kernel operation)} & consume one (te) & p.~\pageref{kop:tuples-elim}\\
\hyperref[kop:union-intro]{\mbox{\texttt{union-intro}}}\index{union-intro@\texttt{union-intro} (kernel operation)} & membership in a UNION, at a chosen side (ui) & p.~\pageref{kop:union-intro}\\
\hyperref[kop:union-elim]{\mbox{\texttt{union-elim}}}\index{union-elim@\texttt{union-elim} (kernel operation)} & case-split a UNION membership (ue) & p.~\pageref{kop:union-elim}\\
\hyperref[kop:intersection-intro]{\mbox{\texttt{intersection-intro}}}\index{intersection-intro@\texttt{intersection-intro} (kernel operation)} & membership in an INTERSECTION (ii) & p.~\pageref{kop:intersection-intro}\\
\hyperref[kop:intersection-elim]{\mbox{\texttt{intersection-elim}}}\index{intersection-elim@\texttt{intersection-elim} (kernel operation)} & project one component (ie) & p.~\pageref{kop:intersection-elim}\\
\hyperref[kop:nth-reduce]{\mbox{\texttt{nth-reduce}}}\index{nth-reduce@\texttt{nth-reduce} (kernel operation)} & NTH of a literal LIST (nth-r) & p.~\pageref{kop:nth-reduce}\\
\hyperref[kop:length-reduce]{\mbox{\texttt{length-reduce}}}\index{length-reduce@\texttt{length-reduce} (kernel operation)} & LENGTH of a literal LIST (len-r) & p.~\pageref{kop:length-reduce}\\
\hyperref[kop:functoid-beta]{\mbox{\texttt{functoid-beta}}}\index{functoid-beta@\texttt{functoid-beta} (kernel operation)} & beta-reduce a functoid application (beta) & p.~\pageref{kop:functoid-beta}\\
\hyperref[kop:nn-induction]{\mbox{\texttt{nn-induction}}}\index{nn-induction@\texttt{nn-induction} (kernel operation)} & induction on NN (ni) & p.~\pageref{kop:nn-induction}\\
\hyperref[kop:transfinite-induction]{\mbox{\texttt{transfinite-induction}}}\index{transfinite-induction@\texttt{transfinite-induction} (kernel operation)} & STRONG transfinite induction on ORD (tfi) & p.~\pageref{kop:transfinite-induction}\\
\hyperref[kop:transfinite-induction-3cases]{\mbox{\texttt{transfinite-induction-3cases}}}\index{transfinite-induction-3cases@\texttt{transfinite-induction-3cases} (kernel operation)} & base/successor/limit transfinite induction (tfi3) & p.~\pageref{kop:transfinite-induction-3cases}\\
\\[-4pt]
\multicolumn{3}{@{}l}{\textbf{Set formation, descriptions and functions} (13)}\\[2pt]
\hyperref[kop:sep-sethood]{\mbox{\texttt{sep-sethood}}}\index{sep-sethood@\texttt{sep-sethood} (kernel operation)} & SEP over a set is a set & p.~\pageref{kop:sep-sethood}\\
\hyperref[kop:sep-mem-intro]{\mbox{\texttt{sep-mem-intro}}}\index{sep-mem-intro@\texttt{sep-mem-intro} (kernel operation)} & y in SEP(x,A,p) from y in A and p[x:=y] & p.~\pageref{kop:sep-mem-intro}\\
\hyperref[kop:sep-mem-elim]{\mbox{\texttt{sep-mem-elim}}}\index{sep-mem-elim@\texttt{sep-mem-elim} (kernel operation)} & consume y in SEP(x,A,p) into y in A and p[x:=y] & p.~\pageref{kop:sep-mem-elim}\\
\hyperref[kop:comp-mem-intro]{\mbox{\texttt{comp-mem-intro}}}\index{comp-mem-intro@\texttt{comp-mem-intro} (kernel operation)} & y in COMP(x,p) from y in SET and p[x:=y] & p.~\pageref{kop:comp-mem-intro}\\
\hyperref[kop:comp-mem-elim]{\mbox{\texttt{comp-mem-elim}}}\index{comp-mem-elim@\texttt{comp-mem-elim} (kernel operation)} & consume y in COMP(x,p) & p.~\pageref{kop:comp-mem-elim}\\
\hyperref[kop:big-union-sethood]{\mbox{\texttt{big-union-sethood}}}\index{big-union-sethood@\texttt{big-union-sethood} (kernel operation)} & set-indexed union of sets is a set & p.~\pageref{kop:big-union-sethood}\\
\hyperref[kop:big-union-mem-intro]{\mbox{\texttt{big-union-mem-intro}}}\index{big-union-mem-intro@\texttt{big-union-mem-intro} (kernel operation)} & membership in an indexed union, at a witness & p.~\pageref{kop:big-union-mem-intro}\\
\hyperref[kop:big-union-mem-elim]{\mbox{\texttt{big-union-mem-elim}}}\index{big-union-mem-elim@\texttt{big-union-mem-elim} (kernel operation)} & consume it, at a fresh eigenvariable & p.~\pageref{kop:big-union-mem-elim}\\
\hyperref[kop:iota-def]{\mbox{\texttt{iota-def}}}\index{iota-def@\texttt{iota-def} (kernel operation)} & the defining property of IOTA(x,p), under uniqueness & p.~\pageref{kop:iota-def}\\
\hyperref[kop:iota-in-elim]{\mbox{\texttt{iota-in-elim}}}\index{iota-in-elim@\texttt{iota-in-elim} (kernel operation)} & IOTA(x,p) DENOTES (per the context), so it satisfies p & p.~\pageref{kop:iota-in-elim}\\
\hyperref[kop:lambda-type]{\mbox{\texttt{lambda-type}}}\index{lambda-type@\texttt{lambda-type} (kernel operation)} & VNB-LAMBDA typing into FUN & p.~\pageref{kop:lambda-type}\\
\hyperref[kop:lambda-beta]{\mbox{\texttt{lambda-beta}}}\index{lambda-beta@\texttt{lambda-beta} (kernel operation)} & beta-reduce an applied VNB-LAMBDA in the GOAL & p.~\pageref{kop:lambda-beta}\\
\hyperref[kop:lambda-beta-hyp]{\mbox{\texttt{lambda-beta-hyp}}}\index{lambda-beta-hyp@\texttt{lambda-beta-hyp} (kernel operation)} & beta-reduce an applied VNB-LAMBDA in an ASSUMPTION & p.~\pageref{kop:lambda-beta-hyp}\\
\\[-4pt]
\multicolumn{3}{@{}l}{\textbf{Rewriting} (4)}\\[2pt]
\hyperref[kop:macete]{\mbox{\texttt{macete}}}\index{macete@\texttt{macete} (kernel operation)} & rewrite the GOAL by an installed macete (mac, macm) & p.~\pageref{kop:macete}\\
\hyperref[kop:macete-hyp]{\mbox{\texttt{macete-hyp}}}\index{macete-hyp@\texttt{macete-hyp} (kernel operation)} & rewrite an ASSUMPTION by one (mac-h, mac-h*) & p.~\pageref{kop:macete-hyp}\\
\hyperref[kop:cartesian-decompose]{\mbox{\texttt{cartesian-decompose}}}\index{cartesian-decompose@\texttt{cartesian-decompose} (kernel operation)} & IN x (CARTESIAN ...) to its existential chain (library.scm) & p.~\pageref{kop:cartesian-decompose}\\
\hyperref[kop:tuple-equality-decompose]{\mbox{\texttt{tuple-equality-decompose}}}\index{tuple-equality-decompose@\texttt{tuple-equality-decompose} (kernel operation)} & LIST = LIST to the conjunction of components (library.scm) & p.~\pageref{kop:tuple-equality-decompose}\\
\\[-4pt]
\multicolumn{3}{@{}l}{\textbf{Oracles} (7)}\\[2pt]
\hyperref[kop:arith-ground]{\mbox{\texttt{arith-ground}}}\index{arith-ground@\texttt{arith-ground} (kernel operation)} & decide a closed arithmetic sentence (arith) & p.~\pageref{kop:arith-ground}\\
\hyperref[kop:arith-forsome]{\mbox{\texttt{arith-forsome}}}\index{arith-forsome@\texttt{arith-forsome} (kernel operation)} & witness a ground existential (arith) & p.~\pageref{kop:arith-forsome}\\
\hyperref[kop:arith-simplify]{\mbox{\texttt{arith-simplify}}}\index{arith-simplify@\texttt{arith-simplify} (kernel operation)} & rewrite by ground context equations, then decide (arith) & p.~\pageref{kop:arith-simplify}\\
\hyperref[kop:ring-simplify]{\mbox{\texttt{ring-simplify}}}\index{ring-simplify@\texttt{ring-simplify} (kernel operation)} & normalise a ring identity (rs, simp) & p.~\pageref{kop:ring-simplify}\\
\hyperref[kop:comm-ring-simplify]{\mbox{\texttt{comm-ring-simplify}}}\index{comm-ring-simplify@\texttt{comm-ring-simplify} (kernel operation)} & normalise a commutative-ring identity (crs) & p.~\pageref{kop:comm-ring-simplify}\\
\hyperref[kop:ineq]{\mbox{\texttt{ineq}}}\index{ineq@\texttt{ineq} (kernel operation)} & the RR order calculus (ineq) & p.~\pageref{kop:ineq}\\
\hyperref[kop:sos]{\mbox{\texttt{sos}}}\index{sos@\texttt{sos} (kernel operation)} & sum-of-squares nonnegativity (sos) & p.~\pageref{kop:sos}\\
\\[-4pt]
\end{longtable}}


%% ---- Line breaking in code (2026-11-15) -----------------------------------
%% A \texttt token such as is-hom-uniformly-continuous used to run into the
%% margin: TeX may not break typewriter text at all.  Here it may break AFTER
%% a hyphen already in the name, and nowhere else: no hyphen is ever added,
%% since the typewriter text is set in a language with no hyphenation patterns.
%% The emergency stretch lets TeX accept a looser line rather than overrun.
\makeatletter
\ifdefined\l@nohyphenation\else\newlanguage\l@nohyphenation\fi
\newcommand{\VNBttbreaks}{\hyphenchar\font=`\-\relax\language\l@nohyphenation}
\AtBeginDocument{%
  \DeclareRobustCommand\ttfamily{%
    \not@math@alphabet\ttfamily\mathtt\fontfamily\ttdefault\selectfont
    \VNBttbreaks}}
\makeatother
\setlength{\emergencystretch}{2em}

%% Paragraph headings: amsbook's default is run-in italic at text size, which reads as
%% emphasis rather than as a heading (the user, 2026-11-17).  Bold, run-in, with space above.
\makeatletter
\def\paragraph{\@startsection{paragraph}{4}{\z@}{\medskipamount}{-0.6em}{\normalfont\bfseries}}
\makeatother

\begin{document}
\maketitle
\tableofcontents
\newpage

\chapter*{Preface}
\label{ch:preface}

\begin{sloppypar}The VNB Proof Checker is an interactive proof assistant implemented in
MIT~Scheme.  Its logical foundation is a first-order set theory in the
style of von~Neumann--Bernays (VNB): a single-sorted universe of
\emph{classes}, among which the \emph{sets} are distinguished as the
classes that are members of the class \textsc{set}.  Equality and
membership ($\in$) are the logical core; the primitive constants and
constructors described below (the number systems, the ordinals, function
spaces) extend it.\end{sloppypar}

The proof architecture follows the sequent-calculus approach of the
IMPS~2.0 system~\cite{IMPS}: proofs are \emph{deduction graphs} whose
nodes are sequents and whose edges are inference steps.  Every
inference step is validated by a small, auditable \emph{primitive
inference kernel}; higher-level tactics and a form of conditional
rewrite rules, called macetes in the IMPS system, are built on top of
the inference kernel and cannot introduce unsoundness.
The deduction-graph proof structure and the use of local contexts in
macete application are due to L.~G.~Monk~\cite{Monk1988}, whose work
was a direct influence on IMPS.

Other proof assistants in current use are built on type theories:
Lean~\cite{Lean} and Rocq~\cite{Rocq} on dependent type theory,
HOL~\cite{HOL} on higher-order logic.  VNB is built on set theory, the
formalism in which most mathematics is written.

This project is, among other things, a proof of concept: a rigorous
mathematical proof system built entirely by AI, with design
specifications and subsequent direction from F.~J.~Thayer
(Section~\ref{sec:vnb-lean}, item~1).


\part{The system}\label{part:system}

%% =========================================================
\chapter{Introduction}
\label{ch:intro}

This chapter compares VNB with Lean (Section~\ref{sec:vnb-lean}), lists its
primitives and the reductions of its notation to them
(Section~\ref{sec:primitives}), and outlines its architecture
(Section~\ref{sec:architecture}).

%% =========================================================
\section{VNB and comparable systems}
\label{sec:vnb-lean}

VNB is one of several proof checkers for mathematics.  Others are
HOL, Isabelle, Rocq, Mizar and Lean; of these, only Mizar is founded on
set theory, as VNB is.  The system most comparable to
VNB is Lean: both aim at a large library of ordinary mathematics whose
proofs are checked by a small kernel, and Lean's library is the usual
point of reference for such work.  The rest of this section compares VNB
with Lean.

Lean is based on dependent type theory (the calculus of inductive
constructions with universes); its mathematical library is Mathlib; it is
implemented in Lean~4 and C++; its proofs are checked by a small kernel,
and independent external checkers of that kernel exist.  Four differences
follow.

\begin{enumerate}
\item \textbf{Authorship and lineage.}  VNB is written entirely by a large
  language model.  It is the result of the VNB project, begun in 2026
  with design specifications and subsequent direction from
  F.~J.~Thayer.  It is partly inspired by IMPS~\cite{IMPS}: proofs are
  sequent-based deduction graphs, terms can be undefined, and rewriting
  is done by what IMPS called macetes.  The project is also indebted to
  the work of L.~G.~Monk (see the Preface), and to that
  of W.~M.~Farmer, among other things for the use of undefined terms and
  for the design of practical reasoning systems~\cite{Farmer2023}.

  VNB does not take over the formal type theory of IMPS (LUTINS, a typed logic of partial
  functions).  Where IMPS uses sorts, VNB uses set-theoretic membership
  guards: a function in \texttt{fun($A$, $B$)} is defined exactly on $A$,
  and a hypothesis $x \in A$ plays the part of a typing judgement.

  The little-theories approach of IMPS~\cite{LittleTheories} appears in a
  smaller role: the structures of the library and the relations between
  them (refinement, and viewing one structure as another,
  Section~\ref{sec:view-as}) are drawn as a graph, the structure graph,
  one of the reference pages that the workspace opens
  (Section~\ref{sec:emacs}).
\item \textbf{Naive set theory with a formal foundation.}  In practice VNB
  is a kind of naive set theory, although its foundation is formal
  (von~Neumann--Bernays--G\"odel set theory, NBG; consistency is relative to
  NBG, Section~\ref{sec:meta-relative}).  It is naive in that its structures follow
  ordinary mathematical practice: a group, a metric space or a normed field
  \emph{is} a tuple of sets (\texttt{declare-structure}, Chapter~\ref{ch:defs},
  ``Mathematical structures''), not a record of a dependent type and not an
  instance of a type class.
\item \textbf{Self-verification.}  VNB is written in Scheme.  This makes
  the verification of large parts of its own code less challenging: the
  kernel's data are S-expressions, which can be encoded as sets, and the
  kernel rules can then be stated and studied as relations on those sets.
  Part~\ref{part:meta} describes what has been built in this direction
  and what has not.  No claim is made here about the difficulty of the
  corresponding task for Lean.
\item \textbf{Provenance of a copy.}  A copy of VNB is one directory,
  the \emph{tree}: VNB's source files,
  its compiled files, the saved image of the loaded library, and the
  proof certificates.  A copy can be made to accept only what has been
  approved:
  \begin{enumerate}
  \item the copy is checked, most thoroughly by the exam, which reruns
    every proof through the kernel;
  \item the tree is then \emph{signed}: a list of a cryptographic hash
    of every file VNB loads, signed with a key that exists only on the
    computer that runs the copy;
  \item each later copy is compared with the signed lists before VNB is
    started from it.  If any file differs, the copy is refused, the
    differing files are named, and the previous copy stays in use.
  \end{enumerate}
  So nothing VNB loads (source, compiled code, certificates, axioms)
  changes without a new signature, and nothing on another machine can
  produce that signature.  The check is made at every start made with
  \texttt{vnb-start} (Section~\ref{sec:signed-trees}).  Chapter~\ref{ch:provenance}
  describes this together with the records of where each theorem comes
  from.
\end{enumerate}

%% =========================================================
\section{VNB: primitives and reductions}
\label{sec:primitives}

The theory originates with von~Neumann~\cite{vonNeumann1925}, was recast
by Bernays with classes and sets as primitives~\cite{Bernays1937}, and was
given its standard finite axiomatization by
G\"odel~\cite{Godel1940}.

Pure VNB set theory is a spare formalism.  The natural numbers, the
reals, functions, and ordinals can all be \emph{constructed} from the
axioms of set theory---but those constructions are long and bear little
resemblance to how mathematics is actually written or thought
about.  A proof that $\sin(\pi/2) = 1$ should not need to unfold the
Dedekind-cut definition of the reals.

This system therefore starts from a broad primitive base.  VNB set
theory supplies the logical foundation, but many mathematical objects
are taken as \emph{primitive constants} and \emph{primitive
constructors}, governed by axioms rather than by explicit
set-theoretic construction.  This means that, unlike the semantics of
pure von~Neumann--Bernays set theory, where each term has a determined
value in a set-theoretic universe, the semantics of VNB is
underdetermined.  In particular:

\begin{itemize}
\item \textbf{Number systems} \texttt{nn}, \texttt{zz}, \texttt{qq},
  \texttt{rr}, \texttt{cc} are primitive constants denoting sets.  Their
  elements and operations are characterised by ring, field, and order
  axioms, not by Dedekind cuts, Cauchy sequences, or ordinal arithmetic.
\item \textbf{Ordinals} \texttt{ord} form a primitive proper class with
  its own ordering predicate \texttt{ord-le}, independent of the
  set-membership relation.
\item \textbf{Function spaces} \texttt{fun($A$, $B$)} are primitive
  set constructors.  A member of \texttt{fun($A$, $B$)} is a set of
  ordered pairs, and the application \texttt{$f$($x$)} is the value
  paired with $x$ (axiom \texttt{app-graph}).
\item \textbf{Mathematical structures} such as groups, vector spaces,
  and Banach spaces are declared with \texttt{declare-structure}, which
  installs accessor operations and a typing predicate.
\end{itemize}

In principle every one of these primitives could be reduced to the bare
VNB axioms---ordinals by the von~Neumann construction, reals by Dedekind
cuts, functions by sets of ordered pairs, and so on.  This approach
simply takes these reductions as \emph{existence theorems} that
have been proved (or could be proved) elsewhere, and works with the
objects directly at arguably a more natural level of abstraction.
The two remarks fit together.  A reduction shows that the axioms
governing a primitive have at least one model inside pure VNB set
theory, and therefore that they are consistent relative to it; it does
not single that model out.  The standard constructions live side by
side in one set-theoretic universe, so the primitive base as a whole
has a model in which all the families hold together; assembling it is
a routine exercise and is not reproduced here.  The real numbers may be read as Dedekind
cuts or as classes of Cauchy sequences, and nothing provable in the
system distinguishes the two readings.  This is the sense in which the
semantics is underdetermined: a theorem of the system holds under every
reading of the primitives that satisfies their axioms.

A further departure from standard set theory is the treatment of
\emph{partial application}.  Any syntactic expression \texttt{f(x)}
is a well-formed term, but it may be \emph{undefined} if \texttt{f} is
not a function applicable to \texttt{x}.  The system provides both
strict equality \texttt{a = b} (false if either side is undefined) and
quasi-equality \texttt{a == b} (true when both sides are undefined,
or both are defined and equal), following the approach of IMPS~\cite{IMPS}.

A consequence of this broad primitive base is that membership in a
primitive constant has \emph{indeterminate truth value} in general.
The formula \texttt{x in 7} is syntactically well-formed --- VNB is
single-sorted and \texttt{in} accepts any two class expressions ---
but the system makes no claim about whether it holds. Please note that
this does not mean that the formula is undefined: whether it is true
or not depends on how $7$ is interpreted in VNB. This system
takes $7$ as a primitive number constant and makes no identification
of it with any set-theoretic object.  Membership statements of this
kind are grammatically meaningful but proof-theoretically inert.

%% =========================================================
\section{Architecture}
\label{sec:architecture}

\subsection{Proof architecture}

\begin{sloppypar}
The proof architecture follows the sequent-calculus approach of the
IMPS~2.0 system~\cite{IMPS}: proofs are \emph{deduction graphs} whose
nodes are sequents and whose edges are inference steps.  Every inference
step is validated by a small, auditable \emph{primitive inference kernel};
higher-level tactics and macetes are built on top of it and cannot
introduce unsoundness.
\end{sloppypar}

\subsection{Structures, not theories}
\label{sec:structures-not-theories}

VNB has no theories.  What a proof assistant built on the \emph{little
theories} methodology of IMPS~\cite{IMPS} would package as a theory---a
signature, its axioms, and the interpretations between theories---is
represented here inside one ambient axiom system (von~Neumann--Bernays set
theory with the number systems) by three mechanisms, each of which is an
ordinary definition of that system.

\begin{itemize}
\item \textbf{A theory is a structure declaration.}  \texttt{declare-structure}
  (Section~\ref{sec:declare-structure}) takes the signature as its slot
  list---carriers, operations, constants, substructures---and the axioms as its
  \texttt{property} and \texttt{law} clauses, and installs one predicate
  \texttt{is-\textit{name}}$(s)$ saying that the VNB list $s$ is a model.
  The models form the class \textit{name} $= \{s : \texttt{is-\textit{name}}(s)\}$,
  a proper class, and a theorem of the ``theory'' is a statement
  $\texttt{forall}([s \in \textit{name}], \ldots)$ of the ambient system.
  There is nothing to interpret between: a group, a ring and a metric space are
  three classes of the same universe.

\item \textbf{The models form a category.}  The declaration also generates the
  morphism predicate \texttt{is-hom-\textit{name}}$(a, b, f_1 \ldots f_k)$, with
  one map per independent carrier.  Each kind of slot contributes one rule
  (Section~\ref{sec:morphisms}):
  \begin{itemize}
  \item an operation or a constant is preserved;
  \item a substructure (the scalars of a module) is held equal;
  \item a \emph{family} slot \texttt{(family F C)}, a set of subsets of the
    carrier $C$ (a $\sigma$-algebra, a topology), is pulled back:
    $\forall u \in F(t).\ \mathtt{preimage}(s, f, u) \in F(s)$;
  \item a \emph{function on a family} \texttt{(family-fun M F R)} (a measure)
    is transported: $M(t)(u) = M(s)(\mathtt{preimage}(s, f, u))$ for
    $u \in F(t)$;
  \item a \emph{relation} slot \texttt{(relation R C)}, a set of pairs of points
    of $C$, is preserved:
    $\forall x_1, x_2 \in C(a).\ [x_1, x_2] \in R(a) \Rightarrow [f(x_1), f(x_2)] \in R(b)$.
  \end{itemize}

  The generated arrows are therefore the homomorphisms of an algebraic
  structure, the isometries of a metric space, the continuous maps of a
  topological space, the measurable maps of a measurable space, the
  measure-preserving measurable maps of a measure space, and the maps of a
  setoid that send related points to related points.  A species whose arrows
  are none of these replaces the definition with \texttt{declare-hom!}
  (\texttt{ringoid}, \texttt{metrizable-top-space},
  Section~\ref{sec:declare-hom}).

  For a structure with one or two independent carriers the declaration also
  installs the hom-set
  $\texttt{hom-\textit{name}}(a, b) = \{f \in \mathtt{fun}(c(a), c(b)) : \texttt{is-hom-\textit{name}}(a, b, f)\}$.
  With two carriers it is a set of pairs of maps.  The library proves, for
  each structure, that the identity is a morphism, that morphisms compose,
  and that the hom-set is a set (\texttt{theorem-library/hom-laws.scm} and
  \texttt{hom-kinds.scm}).  Three or more independent carriers are not yet
  covered: the tree has no separation over triples.

  \emph{Other categories.}  The generated category is only the default one;
  the same objects may carry others.  A metric space, for instance, has its
  isometries by default, but also its Lipschitz, uniformly continuous and
  continuous maps.  \texttt{(declare-category! 'CAT 'NAME '(a b f) BODY)}
  declares such a category on the models of \texttt{NAME}, with arrow
  predicate
  $\texttt{is-\textit{cat}-arrow}(a, b, f) \iff \texttt{is-\textit{name}}(a) \wedge \texttt{is-\textit{name}}(b) \wedge f \in \mathtt{fun}(c(a), c(b)) \wedge \mathtt{BODY}$
  and hom-set \texttt{hom-\textit{cat}}$(a, b)$.  The three category laws are
  then recorded as \emph{obligations}, not as theorems.  An obligation not yet
  proven stops an ordinary load and is listed by a keep-going one
  (Section~\ref{sec:meta-gates})
  (\texttt{category-obligations-audit}).  The library declares four such
  categories: \texttt{lipschitz}, \texttt{uniformly-continuous} and
  \texttt{continuous} on \texttt{metric-space}, and \texttt{bounded-linear} on
  \texttt{normed-vector-space}.  It proves their laws and the inclusions
  $\texttt{hom-metric-space} \subseteq \texttt{hom-lipschitz} \subseteq \texttt{hom-uniformly-continuous} \subseteq \texttt{hom-continuous}$
  and $\texttt{hom-normed-vector-space} \subseteq \texttt{hom-bounded-linear}$
  (\texttt{theorem-library/categories.scm}).

  The Hom functors $\mathrm{Hom}(a, -)$ and $\mathrm{Hom}(-, b)$ are post- and
  pre-composition by a morphism.  Their functor laws are proven once, over any
  set of functions (\texttt{hom-functors.scm}).

  \emph{Essentially small categories.}  The models of a structure form a
  proper class, so a category of them is not itself an object of VNB.  It
  may nevertheless be \emph{essentially small}: there is a set of models
  such that every model is isomorphic to a member of it.  Separable metric spaces under isometry
  are an example.  Every separable metric space embeds isometrically in the
  Urysohn space $\mathbb{U}$, so the subspaces of $\mathbb{U}$, which form a
  set, represent every isometry class.  Gromov's work on spaces of metric
  spaces rests on such sets of representatives: the Gromov--Hausdorff
  distance, for instance, is a metric on the isometry classes of compact
  metric spaces.

  Essential smallness is the reason VNB has no universe axiom.  Some systems
  extend set theory with Grothendieck universes, or adopt Tarski--Grothendieck
  set theory, so that a category of all groups or of all metric spaces can be
  a set.  For an essentially small category this is unnecessary.  A statement about
  all models up to isomorphism quantifies over the set of representatives,
  which the quantifiers reach; only a statement about the class itself is out
  of reach, and none of the mathematics formalized so far has needed one.

\item \textbf{Interpretations are views.}  Where IMPS translates one theory into
  another, VNB \emph{views} an instance of one structure as an instance of
  another: \texttt{def-functor} builds the view (a ring's additive group, a
  normed field as a metric space), and theorems are transported along it by
  specialisation (Section~\ref{sec:view-as}).  A view is a functor between the
  two categories: it sends objects to objects, and its action on arrows is the
  theorem that a morphism of the source is a morphism of the target
  (\texttt{structure-library/functoriality.scm}, \texttt{hom-views.scm}); the one
  view whose carrier changes, the multiplicative group of a field, still owes its
  arrow (the restriction to the non-zero elements).
\end{itemize}

No command takes a theory as an argument, and the record that holds the axioms,
definitions and theorems is called the \emph{library}
(\texttt{add-axiom!}, \texttt{add-definition!}, Chapter~\ref{ch:defs}).

\begin{remark}\label{remark:27137-61828-986360}
%
For the potential user interested in \emph{constructivist} mathematics,
this system is not appropriate. These potential users should consider
alternatives such as Rocq (formerly Coq)~\cite{Rocq}, which has a very well established library of
mathematics.
%
\end{remark}

\subsection{Design goals}

\begin{itemize}
\item \textbf{Expanded set-theoretic foundation.}  VNB set theory with
  primitive constants for standard mathematical objects; no type
  hierarchy; partial function application with undefinedness.
\item \textbf{Transparent kernel.}  A proof is built by
  \VNBkernelopcount{} \emph{kernel operations}, of which
  \VNBkernelopcountoracle{} are decision procedures.  They are listed and
  described in Section~\ref{sec:kernel-boundary}.  Every other command is a
  tactic, which acts on a proof only by calling kernel operations.  The
  code permitted to record an inference is a declared list of eight files
  (\texttt{*cert-kernel-caller-files*}, \texttt{certificates.scm}), whose
  contents are also hashed into every proof certificate; the list is
  re-checked at the end of every library load
  (Section~\ref{sec:kernel-checked}).
\item \textbf{Composable rewriting.}  A macete mechanism (borrowed from
  IMPS) allows theorems to be automatically converted into named rewrite
  strategies that can be sequenced and combined.
\item \textbf{N-ary constructors.}  Ordered tuples \texttt{[$e_1$,
  \ldots, $e_n$]} and Cartesian products \texttt{cartesian($A_1$, \ldots,
  $A_n$)} are first-class, supporting the definition of mathematical
  structures with multiple carriers and operations.
\end{itemize}

\subsection{Why MIT Scheme?}

MIT~Scheme~\cite{MITScheme} was chosen for several practical reasons: the language is
small and uniform, easy to learn and write code, free and readily
available on all major platforms, and has a long track record of
correctness and stability.  A further reason matters for auditing: the
system reads its own source with the same reader it uses for
mathematics, so the checks described in
Section~\ref{sec:kernel-boundary} are performed on the code as parsed,
not on its text.  Since all code is
open source, porting to another language with AI assistance should be
straightforward.

%% =========================================================

%% =========================================================

Part~\ref{part:system} describes the system: its classification of results, its expression language, its deductive system, the records of provenance, definitions and structures.  Part~\ref{part:meta} describes the self-verification of the kernel: the claim and what it is relative to, the admissibility of the kernel's rules over a core, the argument for functoids, and the limits.  Part~\ref{part:using} describes how to use it: the library, getting started, writing proofs, working with formulas from Scheme, and checking results and copies.

Chapter~\ref{ch:quickstart} (Getting started) gives the installation steps and walks through a first session in the workspace, including a complete proof and a macete example.


%% =========================================================
\chapter{Classification of results}
\label{ch:classification}

Every named result carries a \emph{classification}.
It answers two independent questions, and keeping them apart is the whole
point of this chapter:

\begin{itemize}
\item \textbf{What justifies the statement?} --- is it a primitive
  assumed as a base, a definition true by construction, a theorem with a
  machine proof, or a result trusted on a stated warrant without proof?
\item \textbf{May the proof engine activate it unbidden?} --- whether the
  prover is allowed to apply the result as a rewrite rule during a proof.
\end{itemize}

\noindent
These two questions are \emph{orthogonal}.  The first sorts results into
the four kinds below; the second is a separate matter, discussed in
Section~\ref{sec:firing}.  A common confusion --- ``is this proved?'' ---
conflates them, and most of what follows is aimed at pulling them apart.

%% =========================================================
\section{The four kinds}
\label{sec:four-kinds}

\begin{itemize}
\item \textbf{Axiom.} A primitive of VNB --- the base we build on.  The
  axiom set is small and essentially fixed; see
  Section~\ref{sec:primitives}.  The only reason to add one is the
  discovery that the existing primitives do not capture the full power of
  VNB.  An axiom is assumed, not proved.

\item \textbf{Definition.} True \emph{by construction}: a defining
  equation, or an unfolding of one.  A definition is \emph{not a theorem}
  --- there is nothing to prove, since it holds by the meaning of the
  symbol it introduces.  Consequently it is never submitted to the VNB
  test and is \emph{not} a member of the Proof Support Set.  It does,
  however, still \textbf{fire as a rewrite} (Section~\ref{sec:firing}):
  the defining equation is exactly the rule that unfolds the symbol.

\item \textbf{Theorem.} Carries a machine proof.  Its proof script runs
  end-to-end and is re-certified by the VNB test
  (Section~\ref{sec:vnb-test}).  This is the only kind whose truth rests
  on a checked derivation inside the system.

\item \textbf{Proof Support Set (PSS).} A theorem \emph{trusted on a
  warrant, without a machine proof} --- and so \emph{excused} from the VNB
  test.  Every PSS member carries credentials: a plain statement of what
  it says plus a short justification of why it is accepted
  (Section~\ref{sec:warrants}).  The PSS is the curated foundational
  engine of obvious-but-tedious facts (finite-sum identities, ordinal and
  set commonplaces) that are not in doubt and not worth the cost of
  mechanising.  A second motive is reuse: a proof typically calls on many
  small, well-established intermediate facts that recur across unrelated
  developments, and collecting them in the PSS makes each available
  everywhere at once rather than re-derived context after context.
\end{itemize}

\begin{remark}
The PSS exists for a practical reason.  Pushing into new
mathematical territory, a proof constantly bumps into basic --- and
genuinely important --- facts whose proofs are well known but tedious and
time-consuming to mechanise.  Stopping to grind each one out would stall
the advance.  The PSS is the mechanism for brushing them aside: assert the
fact on a warrant, keep moving.  That is exactly why it was introduced.
\end{remark}

\begin{remark}
``Asserted'' is not a fifth kind.  A result installed without a proof is
debt against the fixed axiom base: it is destined either to be discharged
(promoted to a \emph{theorem} with a proof) or recognised as a genuine
primitive gap (an \emph{axiom}), or --- if it is an obvious support fact
--- credentialed into the \emph{PSS}.  An asserted result carrying no
warrant at all is a backlog item, not a stable category.
\end{remark}

%% =========================================================
\section{Activation versus justification}
\label{sec:firing}

Whether the prover may apply a result as a rewrite rule is
\emph{independent} of the four-way classification above.  Axioms,
definitions, theorems, and PSS members can \emph{all} be rewrite rules,
and any of them can equally be inert (cited only by name).  Activation is
the business of the \emph{macete} mechanism, a separate registry described
in Section~\ref{sec:macetes}: installing a result as a macete with
\texttt{install-theorem!} is what lets the engine fire it, and that is a
decision about \emph{use}, not about what justifies the statement.

The clearest case is the definition: it is not a PSS member, yet its
defining equation is precisely the rewrite that unfolds the defined
symbol.  Conversely, a fully proved theorem may be left inert if it is too
specialised to be worth firing automatically.

\begin{center}
\begin{tabular}{@{}llll@{}}
\toprule
\textbf{Kind} & \textbf{Believed because} & \textbf{In VNB test?} & \textbf{May fire?}\\
\midrule
Axiom      & assumed (primitive)   & no  & yes\\
Definition & true by construction  & no  & yes\\
Theorem    & machine proof         & yes & yes\\
PSS        & a stated warrant      & no (excused) & yes\\
\bottomrule
\end{tabular}
\end{center}

\noindent
The last column is the same ``yes'' in every row: firing rights cut across
the classification entirely.

%% =========================================================
\section{Warrants}
\label{sec:warrants}

A definition needs no warrant: it is true by construction, so there is
nothing to justify.  A PSS member, by contrast, is an asserted theorem ---
something that \emph{could} be false and is not being proved here --- and
so it must carry credentials.  A \emph{warrant} is that justification: a
comprehensible statement of the fact together with at least a couple of
lines saying why we accept it without a machine proof.

A warrant is attached with

\begin{VNBsnip}
(warrant! 'name 'kind "free-text justification")
\end{VNBsnip}

\noindent
where \texttt{kind} is one of the recognised grades of justification (the
set \texttt{*warrant-kinds*}), weakest to strongest:

\begin{itemize}
\item \texttt{hand-wave} --- a heuristic or plausibility argument;
  convincing, not rigorous.
\item \texttt{well-known} --- a standard textbook fact, asserted without
  argument.
\item \texttt{reference} --- cites a specific source, named in the text.
\item \texttt{informal} --- a rigorous paper-proof exists (often sketched
  in the file comment or the text), just not mechanised in VNB.
\item \texttt{proof} --- a machine-checked VNB proof exists.
\end{itemize}

\begin{remark}
A warrant is a promissory note, not a permanent verdict: any weaker kind
may later be discharged into \texttt{proof}.  Nothing logical depends on
it --- it is metadata for review and triage.  A result with no warrant at
all ranks below even \texttt{hand-wave}.
\end{remark}

\begin{remark}
Warrants propagate.  When a biconditional theorem is installed, its
automatically generated \texttt{-rev} companion inherits the same warrant,
so the reverse-direction rewrite is credentialed identically to the
forward one.
\end{remark}

%% =========================================================
\section{Proof debt: bills and trust tiers}
\label{sec:proof-debt}

A machine proof is only as unconditional as the facts it cites.  When a
proof reaches \texttt{QED} the system computes its \emph{bill}: the set of
\emph{asserted} facts it ultimately rests on.  The bill is transitive.  A
cited \emph{proven} lemma contributes its own bill; a \emph{primitive}
axiom or a \emph{definition} contributes nothing (it is trusted base); an
\emph{asserted} fact --- a PSS member, or an as-yet-unwarranted assertion
--- contributes itself.  Each fact in the bill is a \emph{leaf}.

The bill is printed inline at \texttt{QED} and collected in the ledger
\texttt{reference/PROOF-DEBT.md}:

\begin{itemize}
\item \texttt{proven modulo 0} --- the \emph{empty bill}.  The proof rests
  only on the trusted base: the fixed primitive axioms, the kernel
  inference rules, and definitions.  Nothing asserted stands between the
  theorem and the kernel.  This is the strongest report a proof can make,
  and it carries \emph{no} trust tier.
\item \texttt{proven modulo \{a, b, \ldots\}} --- unconditional
  \emph{given} $a, b, \ldots$.  Discharging every leaf (proving it, or
  turning it into a definition) would collapse the bill to $0$.
\end{itemize}

\noindent
A non-empty bill is stamped with a \emph{trust tier}
\texttt{[trust:\ \textit{K}]}, and $K$ is the warrant of the bill's
\emph{weakest} leaf: a single shaky citation caps the whole proof, however
solid the rest.  The tiers are the warrant kinds of
Section~\ref{sec:warrants} with one addition below the weakest ---
\texttt{none}, an asserted leaf carrying no warrant at all --- so, weakest
to strongest:
\[
  \texttt{none} \;<\; \texttt{hand-wave} \;<\; \texttt{well-known}
  \;<\; \texttt{reference} \;<\; \texttt{informal} \;<\; \texttt{proof}.
\]
Thus \texttt{[trust:\ none]} is the \emph{weakest} thing a proof can
report --- some leaf justifies nothing --- and \texttt{modulo 0}, which
prints no tier at all, is the strongest.  The two are opposite ends of the
same scale, and are easily misread for one another; the ledger's header
carries this same legend for exactly that reason.

\begin{remark}
The tier is a moving target by design.  Discharging a \texttt{none} or
\texttt{hand-wave} leaf into a warranted or proven fact lifts the tier of
every proof that cites it.  The ledger's reverse index ranks the asserted
leaves by how many proofs each one caps, so the debts worth retiring
first are read off the top.
\end{remark}

%% =========================================================
\section{Circular dependency}
\label{sec:proof-cycle}

The bill of Section~\ref{sec:proof-debt} is computed from a \emph{citation
graph}: at each \texttt{QED} the system records which named results the finished
proof cited.  That graph carries a soundness obligation of its own, independent
of the trust tiers: \textbf{no proven theorem may depend --- transitively,
through its citations --- on itself.}

The failure it guards against is a circle that no single step betrays.  Assert
$X$; prove $Y$ citing $X$; then ``prove'' $X$ citing $Y$.  Each proof is locally
valid, yet together they rest on nothing --- $X$ is believed because of $Y$, and
$Y$ because of $X$.  A bill, a frozen snapshot taken at \texttt{QED}, cannot
reveal such a loop, so the check walks the \emph{live} citation graph instead.

\texttt{(proof-cycle-check)} returns every dependency cycle among the proven
theorems --- the empty list when there are none.  \texttt{load.scm} runs it as a
load-time gate and, on any surviving cycle, \textbf{aborts the load} rather than
merely warn: a genuine circular proof is a soundness bug and is treated as one.

Two kinds of citation edge are deliberately \emph{not} dependencies, and are
excluded when the graph is recorded.  Without the exclusion an ordinary
well-founded induction would read as a cycle:

\begin{itemize}
\item \textbf{A proof's citation of its own name.}  A proof by induction applies
  its induction hypothesis --- the theorem's own statement at a smaller
  instance --- as a rewrite; likewise a structure predicate is unfolded by a
  macete bearing the very name the finished theorem will carry.  Neither is a
  dependence of the theorem on itself.
\item \textbf{A citation that is a definition when it is recorded.}  Unfolding a
  definition adds no debt (Section~\ref{sec:proof-debt}) and no dependency; a
  definitional macete that happens to share a later theorem's name is excluded
  on the same ground.
\end{itemize}

\noindent
Excluding these leaves every genuine circular dependency intact --- such a loop
runs through \emph{proven} citations, which are always kept --- so the gate
neither cries wolf on sound induction nor lets a real circle through.

%% =========================================================
\section{The VNB test}
\label{sec:vnb-test}

The \emph{VNB test} is the standing examination that certifies the
\textbf{theorem} bucket.  It re-proves, in sequence, every result that has
a runnable proof script, and reports whether each still reaches
\texttt{QED}.  Definitions and PSS members are excused by construction:
the former have nothing to prove, the latter are accepted on a warrant.
Run it from the prover root:

\begin{VNBsnip}
./vnb-test
\end{VNBsnip}

\noindent
It harvests the proofs that run at load time and the standalone proof
scripts, runs each, and writes a \emph{certification stamp}
(\texttt{reference/certification.scm}) recording the date and which results
passed, together with a human-readable report
(\texttt{reference/VNB-TEST.md}).  The next time the prover loads, the
stamp is read back, so a result's status can report ``certified by the VNB
test on \textit{date}'' rather than merely ``has a proof script.''

\begin{remark}
The VNB test is \emph{not} the unit test suite.
\texttt{test-suite.scm} is a feature-and-regression suite
of checks, run during development; it checks that the
machinery behaves, not that the library's theorems still derive.  The VNB
test is the slower, periodic exam that re-runs the mathematics
itself --- re-deriving every proof in the library, as any serious proof
system must.  Keep both: they answer different questions.
\end{remark}

%% =========================================================
\section{Inspecting a result}
\label{sec:status-query}

Three commands expose the classification at the REPL.

\begin{itemize}
\item \texttt{(status 'name)} prints, on one screen, a result's
  provenance, whether it is a PSS member, any warrant it carries, and its
  certification state (``certified by the VNB test on \textit{date}'' or
  ``not certified since the last run'').

\item \texttt{(status-audit)} writes \texttt{reference/STATUS-AUDIT.md}: a
  census of the whole library by kind, flagging hygiene issues --- asserted
  results with no warrant, theorem-shaped names that have drifted out of
  the theorem bucket, and the like.

\item \texttt{(pss-letter 'name)} drafts a plain-English ``recommendation
  letter'' for a candidate PSS member: an English rendering of what it
  says, its current standing, and an assessment against the admission
  gates (true; a sound, terminating rewrite shape; general rather than
  bespoke; non-redundant).  It is an aid to deciding whether a result
  belongs in the curated support set, not an automatic admission.
\end{itemize}


%% =========================================================
\chapter{The expression language}
\label{sec:language}


The language has two syntactic layers.

The \emph{internal layer} represents every expression as a plain
Lisp S-expression.  This representation is time-tested and deliberately
simple: it requires no parser, supports straightforward recursive
processing, and is directly inspectable and constructible at the Scheme
REPL.  All prover machinery---inference rules, substitution, pattern
matching, macetes, arithmetic evaluation---operates on this layer.
Users are not required to engage with it directly; however, in the
tradition of early Lisp, they are free to do so.  Constructing or
manipulating S-expressions by hand, loading axioms as quoted lists,
or writing Scheme code that builds formulas programmatically are all
legitimate and supported uses.

The \emph{string syntax} is the surface notation accepted by
\texttt{make-wff-from-string} and emitted by \texttt{wff->string}
/ \texttt{pp}: binary operators are written infix, multi-argument
terms use functional notation, quantifiers use the binding-list form,
and set expressions use brace notation.
Section~\ref{sec:parser-printer} describes both layers in full.
The examples in this manual use string syntax for user-facing
formulas; internal S-expression forms appear only in library and
structure definitions.

MIT~Scheme reads symbols case-insensitively (folding to lower case),
so \texttt{FORALL}, \texttt{Forall}, and \texttt{forall} are the same
symbol once read.

As should be the case in any formal system, we distingish between
syntax and semantics. In what follows, we distinguish two distinct
types of syntactic categories, each with its own semantics: formulas
and terms. These syntactic categories can be viewed in either the
internal syntax or the string syntax. In the presentation below, we
adhere to the string syntax.


\section{Formulas}

Semantically, a formula is a syntactic expression that can have only
two values, truth and falsity.  The constant formulas \texttt{truth} and
\texttt{falsity} denote them.
%
\begin{remark}\label{remark:27137-52766-301830}
%
Note that ``UNDEFINED'' is not a possible value for a
formula. However, the value of any particular formula may not be
known. In fact, since we are silent about how elements in our system
may be interpreted in a pure von Neumann Bernays theory, different
interpretations may yield different values.
%
\end{remark}

Formulas can be produced from other ones by basic logical operators,
including quantification over variables that range over objects.
%
In the following table, the symbols $p$, $q$ are not variables that
range over formulas, but rather placeholders that can instantiated
with well-formed formulas.

\begin{itemize}
\item \texttt{truth} propositional truth $\top$
\item \texttt{falsity} propositional falsity $\bot$
\item \texttt{not($p$)} negation $\lnot p$
\item \texttt{$p$ and $q$} conjunction $p \land q$
\item \texttt{$p$ or $q$} disjunction $p \lor q$
\item \texttt{$p$ implies $q$} implication $p \Rightarrow q$
\item \texttt{$p$ iff $q$} biconditional $p \Leftrightarrow q$
\item \texttt{forall([$b_1$, \ldots, $b_n$], $p$)} universal quantification; see Section~\ref{sec:multiquant}
\item \texttt{forsome([$b_1$, \ldots, $b_n$], $p$)} existential quantification; see Section~\ref{sec:multiquant}
\item \texttt{$a$ = $b$} strict equality: both defined and equal
\item \texttt{$a$ == $b$} quasi-equality: both defined and equal, or both undefined
\item \texttt{$t$ := $s$} definition (sugar for \texttt{$t$ == $s$}): $t$ is defined iff $s$ is, and then equal
\item \texttt{$a$ in $b$} membership $a \in b$
\item \texttt{$A$ subset $B$} subclass $A \subseteq B$; see Section~\ref{sec:class-set-discipline}
\end{itemize}

\noindent
The quantifiers \texttt{forall} and \texttt{forsome} range over \emph{all
classes} (the system is single-sorted).  There is no restriction to sets.

\subsection{Quantification}
\label{sec:multiquant}

Both \texttt{forall} and \texttt{forsome} require a bracket-enclosed
\emph{binding list}, even for a single variable:

\begin{itemize}
\item \texttt{forall([binding, ...], body)}

\item \texttt{forsome([binding, ...], body)}
\end{itemize}

\noindent
\texttt{forall(x, x = x)} is a syntax error; the correct form is
\texttt{forall([x], x = x)}.

\noindent
Each \emph{binding} in the string syntax is one of:

\begin{itemize}
\item \texttt{$x$} unrestricted variable ranging over all classes, sets
  among them --- and over nothing else: a functoid is not in the range
  (Remark~\ref{rem:functoids-quantified}), and the instantiation rule
  refuses one as an instance
\item \texttt{$x$ in $A$} variable $x$ restricted to class $A$
\item \texttt{[$n_1$, \ldots, $n_k$] in $A$} tuple destructuring: fresh $t \in A$, components bound to \texttt{nth(1,~$t$)}, \ldots, \texttt{nth($k$,~$t$)}
\end{itemize}

\noindent
In the \emph{S-expression syntax} a binding can additionally group
several unrestricted variables into a single spec:
\[
  \texttt{(forall ((}x_1\ x_2\ \ldots\ x_n\texttt{) ...) body)}
\]
\begin{sloppypar}
The grouped binding \texttt{($x_1$ $x_2$ \ldots\ $x_n$)} is exactly
equivalent to the sequence of single-variable bindings
\texttt{($x_1$) ($x_2$) \ldots\ ($x_n$)}.  The expander
\texttt{expand-\allowbreak destructuring-\allowbreak quantifiers}
of \texttt{wff.scm} nests one
quantifier per variable of the spec, in place, so that the two surface
forms yield \emph{identical} internal expressions---equal, not merely
alpha-equivalent.  The same holds for a grouped spec occurring in a
longer binding list alongside restricted specs.
\end{sloppypar}

Three forms of the same formula are therefore in circulation:

\begin{enumerate}[label=(\arabic*),leftmargin=1.2cm]
\item the \emph{grouped} binding list,\\
      \texttt{(forall ((x y z)) body)};
\item the \emph{flat} binding list,\\
      \texttt{(forall ((x) (y) (z)) body)};
\item the \emph{fully nested} form,\\
      \texttt{(forall ((x)) (forall ((y)) (forall ((z)) body)))}.
\end{enumerate}

\noindent
The utility functions \texttt{vnb-unwind} and \texttt{vnb-wind} relate
form~(3) to the other two.  \texttt{vnb-unwind} rewrites any binding
list, grouped or not, into the fully nested form; \texttt{vnb-wind}
collapses a run of nested quantifiers \emph{of the same kind} into a
single binding list of single-variable specs.

\begin{VNBsnip}
(vnb-unwind '(FORALL ((x y z)) body))
;=> (FORALL ((x)) (FORALL ((y)) (FORALL ((z)) body)))

(vnb-unwind '(FORALL ((x) (y) (z)) body))
;=> (FORALL ((x)) (FORALL ((y)) (FORALL ((z)) body)))

(vnb-wind '(FORALL ((x)) (FORALL ((y)) (FORALL ((z)) body))))
;=> (FORALL ((x) (y) (z)) body)
\end{VNBsnip}

\noindent
The two are not mutually inverse.  \texttt{vnb-wind} never regroups a
spec, and never splits one, so it leaves a grouped binding list
untouched and the round trip normalizes form~(1) to form~(2):

\begin{VNBsnip}
(vnb-wind '(FORALL ((x y z)) body))
;=> (FORALL ((x y z)) body)

(vnb-wind (vnb-unwind '(FORALL ((x y z)) body)))
;=> (FORALL ((x) (y) (z)) body)
\end{VNBsnip}

\noindent
Two further properties are worth recording.  First, neither function
expands the binding sugar to primitive form---that is the office of
\texttt{make-wff}---so a restricted spec survives unwinding as a spec,
with no \texttt{implies} introduced.  Second, \texttt{vnb-wind} merges
only runs of the same quantifier; a \texttt{forall} over a
\texttt{forsome} is returned unchanged.

\begin{VNBsnip}
(vnb-unwind '(FORALL ((x) (IN z NN)) body))
;=> (FORALL ((x)) (FORALL ((IN z NN)) body))

(vnb-wind '(FORALL ((x)) (FORSOME ((y)) body)))
;=> (FORALL ((x)) (FORSOME ((y)) body))
\end{VNBsnip}

\noindent
Bindings are processed left to right, outermost first.  Thus in the
\emph{internal syntax} they expand to fully nested single-variable form
as follows:

\begin{VNBsnip}
forall([x, y in NN, z], y >= 0)
; => (FORALL x (FORALL y (IMPLIES (IN y NN) (FORALL z (>= y 0)))))

forall([[E, +, *] in METRIC], IS-METRIC(E, +, *))
; => (FORALL t_0 (IMPLIES (IN t_0 METRIC)
;                         (IS-METRIC (NTH 1 t_0) (NTH 2 t_0) (NTH 3 t_0))))
\end{VNBsnip}

\paragraph{Display.}
The printer reconstructs binding-list form from the internal nested
representation.  Consecutive same-type quantifiers are collapsed into
a single binding list; a change of quantifier type stops the collapsing.
For example, \texttt{forall([x in nn], forall([y], p))} prints as
\texttt{forall([x in nn, y], p)}.

\paragraph{Variable shadowing.}
Bound variable names may be reused within nested quantifiers.
The inner binding shadows the outer; \texttt{fresh-var} in
\texttt{subst-free} ensures substitution never captures a free variable.
Such formulas are valid but confusing; users should avoid them.

\subsection{Equality and quasi-equality}
\label{sec:equality}

Terms may be \emph{undefined} (e.g.\ an application \texttt{f(x)} when
\texttt{f} is not a function applicable to \texttt{x}).  The system
provides two equality predicates:

\paragraph{Strict equality \texttt{$a$ = $b$}.}
True iff both \texttt{a} and \texttt{b} are defined and equal.
False if either is undefined.  In particular:
\begin{itemize}
\item \texttt{george(s) = 75} is \emph{false} when \texttt{george}
  is not a function.
\item \texttt{george(s) = george(s)} is \emph{false} for the same
  reason.
\item A term $t$ is defined iff $t = t$:
  $\mathtt{DEF}(t) \;\equiv\; (t = t)$.  Section~\ref{sec:definedness}
  says more.
\end{itemize}


\paragraph{Quasi-equality \texttt{$a$ == $b$}.}
True iff both \texttt{a} and \texttt{b} are defined and equal, \emph{or}
both are undefined.  Key properties:
\begin{itemize}
\item \texttt{t == t} is \emph{always true}, for any term $t$,
  defined or not.
\item \texttt{george(s) == george(s)} is \emph{true}.
\item \texttt{$a$ = $b$} implies \texttt{$a$ == $b$}, but not vice versa.
\item \texttt{$a$ == $b$} and $\mathtt{DEF}(a)$ together imply
  \texttt{$a$ = $b$}.
\end{itemize}

\noindent
Undefinedness propagates through all operations: if \texttt{george(s)}
is undefined then \texttt{george(s) + 1} is also undefined,
and \texttt{george(s) + 1 = 2} is false.

The converse direction does not hold for the arithmetic symbols.
\texttt{+}, \texttt{*}, \texttt{-}, \texttt{recip}, \texttt{abs},
\texttt{conjugate}, \texttt{succ} are primitive operator symbols with
\emph{closure} axioms only ($a, b \in \mathtt{cc} \Rightarrow a + b \in
\mathtt{cc}$) and no domain: from the definedness of $a + b$ nothing follows
about $a$ or $b$, and for two objects outside every number system the term
$a + b$ is simply underdetermined --- not undefined, not defined, unspecified.
Every proof mechanism types the arguments before it concludes anything about
the sum, so this is invisible inside a proof; it is visible in a statement
that applies \texttt{+} to members of an unconstrained class.  The full
account is Section~\ref{sec:arith-symbols}.

\paragraph{Definitions: \texttt{$t$ := $s$}.}
A definition is a different speech act from an equation.  Writing
\texttt{$t$ := $s$} introduces $t$ \emph{as} $s$: $t$ is defined exactly when
$s$ is, and then equal --- $t$ inherits its definedness from $s$, having none
of its own.  That is precisely quasi-equality, so \texttt{:=} is surface
sugar for \texttt{==}: the parser reads \texttt{$t$ := $s$} as
\texttt{$t$ == $s$}.  Use \texttt{:=} for the defining equation of a partial
operation --- e.g.\ \texttt{(compose(f, g))(x) := f(g(x))} --- and reserve
\texttt{=} for a genuine \emph{identity}, asserted between two terms each
\emph{independently} known to be defined.  Writing a definition with
\texttt{=} over-claims: it asserts the defined term is defined at
\emph{every} argument, which is false off its domain (there both sides are
undefined, and \texttt{=}, being strict, is then false --- whereas
\texttt{==} is true).

\paragraph{Neither \texttt{=} nor \texttt{==} is a function.}
In VNB both are formula formers: $t$~\texttt{=}~$s$ and
$t$~\texttt{==}~$s$ are formulas, never terms.  Their values are truth
values, not objects of the universe, so nothing is ever applied to them,
and an undefined $t$ or $s$ raises no question of an application to an
undefined argument.  (In IMPS~\cite{IMPS}, where formulas are terms of a
sort of propositions, \texttt{==} is instead an abbreviation, a
\emph{quasi-constructor}, which the reader expands into
$(t{\downarrow} \lor s{\downarrow}) \supset t = s$.)

\section{Definedness}
\label{sec:definedness}
\index{definedness}

A formula always has a truth value (Remark~\ref{remark:27137-52766-301830}).
A term need not have a value: it may fail to denote, and is then
\emph{undefined}.  This section says which terms denote, how VNB states
and proves that a term denotes, how that can fail, and why there is no
separate predicate for it.

\subsection{Which terms denote}

Always defined:
\begin{itemize}
\item variables, and constants such as \texttt{ord}, \texttt{rr},
  \texttt{empty-set};
\item a total class constructor (\texttt{sep}, \texttt{fun},
  \texttt{interval}, \ldots, Section~\ref{sec:class-set-discipline})
  applied to defined arguments.
\end{itemize}
Defined only under a condition:
\begin{itemize}
\item an application $f(x)$: it is defined when $f \in \mathtt{fun}(A, B)$
  and $x \in A$ for some $A$, $B$ (Section~\ref{sec:funapp});
\item a description $\iota x.\, p$: it is defined when exactly one object
  satisfies $p$;
\item \texttt{choice($A$)}: it is defined when $A$ is nonempty;
\item a list $[e_1, \ldots, e_n]$: it is defined when every $e_i$ is a
  set.  So $[0, \mathtt{set}]$ denotes nothing, and neither does
  $\{\mathtt{set}\}$.
\end{itemize}
Undefinedness propagates: a term with an undefined subterm is undefined,
so if \texttt{george(s)} is undefined, so is \texttt{george(s) + 1}.

The arithmetic symbols are a different case.  On two objects outside
every number system, $a + b$ is not undefined but unspecified: no axiom
says what it is (Section~\ref{sec:arith-symbols}).

\subsection{$t = t$ says that $t$ is defined}

Equality is strict (Section~\ref{sec:equality}): $a = b$ is true exactly
when both terms are defined and equal.  So $t = t$ is true exactly when
$t$ is defined, and VNB uses it as the statement ``$t$ is defined'':
\[
  \mathtt{DEF}(t) \;\equiv\; (t = t).
\]

Every atomic formula is strict in the same way.  $t \in A$,
$a \le b$ and $P(t)$ are false when one of their terms is undefined.
So a true atomic formula shows that each of its terms denotes: $t \in A$
implies $t = t$ for any $A$.  The formula $t \in \mathtt{set}$ says more,
namely that $t$ denotes a \emph{set}.  It is the form most axioms ask
for; the image axiom, for instance, asks for the sethood of $f(x)$.

Two pitfalls:
\begin{itemize}
\item $t == t$ is true for every term, defined or not, so it says nothing
  about $t$;
\item a defining equation written with $=$ asserts that both sides are
  defined everywhere; use $:=$ (Section~\ref{sec:equality}).
\end{itemize}

\subsection{How definedness is proved}

A proof establishes $t = t$ in one of two ways.
\begin{itemize}
\item By the \emph{definedness certificate}
  (Section~\ref{sec:definedness-certificate}), a fixed test on $t$ and
  the context.  It accepts a variable or constant, a class term on
  defined arguments, a term typed or equated in a hypothesis, $f(a)$ with
  $f$ typed and $a$ in its domain, and a few other forms.  It never
  accepts a choice, a description or an untyped application.
\item By proving a typing fact $t \in A$, from which $t = t$ follows.
\end{itemize}
Two kernel operations need definedness and consult the certificate:
\begin{itemize}
\item \texttt{rfl} closes $t = t$ only when the certificate accepts $t$.
  \texttt{qrfl} closes $t == t$ with no condition.  So the strict
  reflexivity is not a triviality.
\item \texttt{forall-elim} concludes $p[x := t]$ from
  $\forall x.\, p$ only for a term $t$ that denotes.  When the
  certificate declines, it leaves $t = t$ as an open goal, to be closed
  by a typing fact.  Without that goal the rule would prove
  $\mathtt{recip}(0) = \mathtt{recip}(0)$.  The policy is stated in
  \texttt{docs/definedness-instantiation-2026-09-18.md}.
\end{itemize}

\subsection{How it can fail}

Within the logic, $t = t$ is false for an undefined $t$, and so is any
atomic formula that types or equates it.  A statement about an undefined
term is false, not meaningless: \texttt{george(s) = george(s)} is false
when \texttt{george} is not a function.

Within the checker, the certificate draws one inference from
strictness: a term that occurs in a true atomic hypothesis is defined.
That inference is sound only if every rule that concludes an atomic
formula has established that the formula's terms denote.  A rule that
concluded, for instance, $[\,] \in \mathtt{tuples}(A)$ for every $A$
would break it.  Take $i = \iota x.\,(x = x \wedge \textsc{falsity})$,
a description of nothing, for $A$.  The certificate would then accept
$i$, \texttt{rfl} would give $i = i$, and \texttt{iota-in-elim} would
give $i = i \wedge \textsc{falsity}$, hence \textsc{falsity}.  So each
kernel rule that concludes an atomic formula owes the definedness of
the terms in it, its class arguments included.

\subsection{Why there is no \texttt{is-defined($t$)}}

It would be $t = t$ under another name.  Every atomic formula is strict,
so a primitive \texttt{is-defined($t$)} would be an atomic formula true
exactly when $t$ denotes, which is the truth condition of $t = t$.  A
second primitive would need its own kernel rules and checkers, and the
definedness certificate would have to treat both forms alike: more code
to trust, and nothing new to say.

What to write instead:
\begin{itemize}
\item $t = t$, for ``$t$ is defined'';
\item $t \in \mathtt{set}$, for ``$t$ is a set'', the stronger and more
  common requirement;
\item $t \in A$, for ``$t$ is defined and lies in $A$''.
\end{itemize}
A surface spelling such as \texttt{defined($t$)}, read by the parser as
$t = t$ (as $:=$ is read as $==$), would change nothing in the kernel.
The system does not have one.

%% =========================================================
\section{Class constructors}
\label{sec:class-set-discipline}

Proper VNB has two kinds of object: \emph{sets} (elements of
the class $\mathtt{set}$) and \emph{proper classes} (collections that can have
members but are not themselves members of anything).  We retain that
distinction and add a third kind, \emph{functoids}.

A \textbf{functoid} is an operation specified by a \emph{signature declaration},
expressed as
\[
  A_1, \ A_2, \ \cdots, \ A_n \;\to\; B
\]
where $A_1, \ldots, A_n, B$ are classes; the arrow is pure metanotation
for the signature, not an object-language operation (no Cartesian
product is implied).  This declaration specifies that the functoid in
question takes $n$ arguments $a_1 \in A_1, a_2 \in A_2, \ldots, a_n
\in A_n$ producing a value $b \in B$.

%% sec-functoids.tex -- Sets, classes and functoids: what the quantifiers reach.
%% Written 2026-10-01 after the image-axiom inconsistency of 2026-09-30 and the decisions
%% of 2026-10-01 (functoids outside quantification; the beta value licence; the DOM
%% axiom; the tuple tripwire; fun-image-set).  \input from ch-expressions.tex, in
%% place of the remark that used to carry the label rem:functoids-quantified.

\subsection{Sets, classes and functoids: what the quantifiers reach}
\label{sec:functoids-doctrine}
\label{rem:functoids-quantified}

This subsection states the doctrine in one place: which things are objects of
the theory, which are not, what an unguarded quantifier ranges over, what an
application $f(x)$ means for each kind of $f$, what $\beta$-reduction of a
\texttt{vnb-lambda} is licensed to do, and where the axiom of replacement
enters.  Every statement here is enforced by a kernel rule or a gate, named
in the text; two remarks at the end show, on two terms the syntax admits,
what the rules are there to handle.

\subsubsection{Three kinds of thing}

\begin{enumerate}
\item A \textbf{set} is an element of the class $\mathtt{set}$.
\item A \textbf{proper class} has members but is a member of nothing.
  Sets and proper classes together are the \emph{objects} of the theory:
  every variable ranges over them, and every term that denotes denotes one
  of them.
\item A \textbf{functoid} is an operation on objects, given by a signature
  $A_1, \ldots, A_n \to B$ (metanotation: no Cartesian product is formed)
  and a rule.  It is \emph{not} an object.  Functoids enter in four ways:
  the primitive class constructors (\texttt{power}, \texttt{fun},
  \texttt{image}, \texttt{tuples}, \ldots); \texttt{def-functoid}
  (Section~\ref{sec:def-functoid}); the \texttt{lambdoid} binder; and
  \texttt{def-functor}, since every view between structures is the
  \texttt{def-functoid} that builds the target instance from the source's
  slots (Section~\ref{sec:view-as}).
\end{enumerate}

A functoid \emph{applied} to arguments, $\texttt{power}(X)$,
$\texttt{card}(A)$, $\mathit{view}(r)$, is a term, and the term denotes an
object or is undefined.  A functoid \emph{name standing alone} is not a
term of the object language.  The sentence ``for every functoid $F$'' is
metalanguage.

The arithmetic symbols \texttt{+}, \texttt{*}, \texttt{-}, \texttt{recip},
\texttt{abs}, \texttt{conjugate}, \texttt{succ} are none of the three.  They
are \emph{primitive operator symbols} (filed by the operator census under
``syntax: the head denotes nothing''): \texttt{(+ a b)} is a term, \texttt{+} alone is not, and the theory
characterises the \emph{applications} of the symbol by closure axioms and
guarded identities while saying nothing about the symbol --- in particular
nothing about its domain, so a sum of two objects outside every number system
is an underdetermined term.  A functoid has a signature and a rule; a
function is a set; \texttt{+} has neither.  When addition has to be named as
an object it is \texttt{binplus}, with \texttt{binplus(x, y) == x + y}, or the
\texttt{add} slot of a structure instance, which is a \texttt{vnb-lambda} over
that instance's carrier and hence a function with a domain.
Section~\ref{sec:arith-symbols} has the full account, including why the
function typing of \texttt{binplus} would prove FALSITY.

\subsubsection{What an unguarded quantifier reaches}

An unrestricted binding \texttt{forall([x], \ldots)} ranges over all
objects, sets and proper classes, and over nothing else.  The argument is
short: if functoids were in the range of a quantifier that ranges over
classes, they would be classes, and a class has a membership relation; no
axiom gives $x \in F$ a truth value for a functoid $F$; hence functoids are
outside the range.

The kernel enforces this at the one place where a quantifier meets a term:
\texttt{forall-elim} (\texttt{pi-instantiate!}, and its independent checker
in \texttt{rule-checkers-logic.scm}) refuses an instance term that is a
bare functoid, that is, a registered functoid name or a \texttt{lambdoid}
record.  An applied functoid is a class term and is as legal an instance as
any; so is a \texttt{vnb-lambda}, which is a class of pairs; so is a function
variable.  The library's functions on proper-class domains (the class
functions of \texttt{ord-no-injection-into-set} and of Zorn's lemma) are
classes of pairs, not functoids, and are quantified over as classes.

\begin{center}
\begin{tabular}{lll}
\hline
 & \texttt{vnb-lambda} & \texttt{lambdoid} / \texttt{def-functoid} \\
\hline
denotes & a set of ordered pairs & a functoid (meta-level record) \\
domains & sets & arbitrary classes, $\mathtt{ord}$ included \\
typing rule & \texttt{lam-t}: $\in \mathtt{fun}(A, B)$ & none \\
reduction & \texttt{lam-b}, under the licence below & \texttt{beta}, by unfolding \\
instance of a quantifier & yes & refused \\
in an \texttt{image}, a \texttt{dom}, an \texttt{app-graph} & yes & never formed \\
\hline
\end{tabular}
\end{center}

\subsubsection{Application and its graph}

For every object $f$ the axiom \texttt{app-graph} reads application off
membership of pairs:
\[
  f(x) \;==\; \iota v.\ [x, v] \in f .
\]
For $f \in \mathtt{fun}(A)$ the axiom \texttt{fun-no-junk} adds that every
member of $f$ is such a pair with $x \in A$; the two together say that a
function is exactly its graph.  Equality is partial: an atomic formula is
true only when its terms denote, and $t = t$ is the definedness predicate
of $t$.  The domain of an arbitrary $f$ is the class
\[
  x \in \mathtt{dom}(f) \iff x \in \mathtt{set} \;\wedge\; f(x) \in \mathtt{set}
\]
(\texttt{dom-membership}).  On a class of pairs the second conjunct
could equally be $f(x) = f(x)$, definedness, since $f(x)$ is defined
exactly when a unique set $v$ has $[x, v] \in f$, and that $v$ is a set.
The sethood form is the one stated because the image axiom (below) asks
for the sethood of $f(x)$, and the kernel has no rule from ``defined'' to
``a set''.  For $f \in \mathtt{fun}(A)$,
$x \in \mathtt{dom}(f) \iff x \in A$ (\texttt{dom-fun-membership}).

Since a functoid is refused as an instance of $f$, these axioms say nothing
about functoids.  The value of a functoid at an argument is given by its
defining rewrite alone, and that value may be a proper class
($\texttt{power}(\mathtt{ord})$).

\subsubsection{Pairs and tuples}

\texttt{list} is a primitive constructor; the lists it produces are
written $[e_1, \ldots, e_n]$ in the string syntax, and they are not reduced
to sets.  A list is defined exactly when its components are sets, and
then it is a set: $[e_1, \ldots, e_n]$ denotes only if every $e_i$ is a
set (the definedness certificate of the kernel decides it, so $\texttt{rfl}$
closes $[x, y] = [x, y]$ only once $x$ and $y$ are typed in context), and a
list of sets is a set (\texttt{pair-tuple-is-set}, \texttt{list-sethood}).
$[0, \mathtt{set}]$ denotes nothing, as $\{\mathtt{set}\}$ denotes nothing,
and an application whose value would be such a pair's second component is
undefined.  The axiom
\[
  L \in c \;\Longrightarrow\; \bigl(1 \le k \le \texttt{length}(L)
  \;\Longrightarrow\; \texttt{nth}(k, L) \in \mathtt{set}\bigr)
  \qquad (k \in \mathtt{nn})
\]
(\texttt{tuple-members-are-sets}, primitive) states this
sentence, which every reduction of pairs to sets supplies and the primitive
constructor does not.  Nothing cites it; it is a tripwire: should any
derivation produce a pair with a proper class as a member of
something, two steps (\texttt{nth-r}, then $\mathtt{set} \in \mathtt{set}$
against Burali-Forti) turn it into \texttt{falsity} instead of leaving it inert.

One consequence deserves its own sentence.  Because pairs are not sets,
``the union of the second components of the pairs of $f$'' is not a term
the axioms reach, and the ZF derivation of the sethood of an image from
union and separation (take $\bigcup\bigcup f$, then separate) is not
available here.  See \emph{Images and replacement} below.

\subsubsection{The $\beta$-licence of \texttt{vnb-lambda}}

The contraction
\[
  (\texttt{vnb-lambda}(x \in A, b))(u) \;==\; b[u/x]
\]
is a sound rewrite exactly when $b[u/x]$ is a set (then $[u, b[u/x]]$ is a member of the
lambda) or undefined (then both sides are undefined).  It fails in one
case: $b[u/x]$ a proper class, where the right side denotes and the left
does not.  The rule \texttt{lam-b} (\texttt{pi-lambda-beta!}; checker
\texttt{rcs--beta-licences-ok?}) therefore licenses a contraction on two
conditions, decided before either is posted:
\begin{itemize}
\item the \emph{argument} is in the domain, $u \in A$: owed as an
  obligation leaf $(\mathtt{in}\; u\; A)$ when not evident from the context;
\item the \emph{value} is not a proper class: owed as
  $(\mathtt{in}\; b[u/x]\; \mathtt{set})$ exactly when the value \emph{may}
  denote a proper class.  The test is negative: $\mathtt{set}$ and
  $\mathtt{ord}$; a bare variable with no typing in context; \texttt{power},
  \texttt{fun}, \texttt{tuples}, \texttt{injection}, \texttt{bijection} of a
  class that may be proper; a separation, a lambda, an image or a big union
  whose domain may be proper; a union, intersection, difference, relative
  complement or product with a part that may be proper; a complement
  (\texttt{comp}), always; a registered functoid applied, by its unfolded
  body; an \texttt{if} with such a branch.  A numeral, a
  typed variable, the set constants, arithmetic, a tuple, a choice or
  description, an application of a variable, and every other term owe
  nothing.  The test is negative because a lambda's value is typically
  an arithmetic term, an application or a choice: typed nowhere, and
  never a proper class.
\end{itemize}
The argument's membership, licensed or owed, counts as a typing of the
value, so
\[
  (\texttt{vnb-lambda}(z \in A, z))(c)
\]
owes $c \in A$ and nothing more.  A redex whose obligation would mention a variable bound where the
redex sits is left standing.

\paragraph{Example.}  The monomial is
\[
  \texttt{monomial}(R, c, n) := \texttt{vnb-lambda}(k \in \mathtt{nn},\;
  \texttt{if}(k = n, c, \mathtt{zero}(R))).
\]
The law
\[
  \texttt{monomial}(R, c, n)(n) = c
\]
is false when $c$ is unrestricted.  At a proper class $c$ the lambda has
no pair at $n$, the application is undefined, and the strict equation
fails.  Accordingly the contraction owes
\[
  \texttt{if}(n = n, c, \mathtt{zero}(R)) \in \mathtt{set},
\]
which nothing closes.  The law holds under the hypothesis
$c \in \mathtt{carr}(R)$.  This is the pitfall
Section~\ref{sec:class-set-discipline} describes: a strict equation whose
terms are untyped asserts their definedness, and here their sethood.

\subsubsection{Images and replacement}

The image of a class $S$ under an object $\phi$ is characterised by
\[
  w \in \mathtt{image}(\phi, S) \iff
  w \in \mathtt{set} \;\wedge\; \exists x \in S.\ \phi(x) = w
\]
(\texttt{image-membership-iff}, definitional).  The conjunct
$w \in \mathtt{set}$ is Bernays' definition made explicit: the image
collects the \emph{values that are sets}; a point where $\phi(x)$ is
undefined contributes nothing, and since $\phi$ is an object there is no
point where $\phi(x)$ is a proper class.  Two sethood facts stand beside
it.

\begin{description}
\item[\texttt{image-set} (replacement, primitive).]
  $S \in \mathtt{set} \Longrightarrow \mathtt{image}(\phi, S) \in
  \mathtt{set}$ for every object $\phi$.  This is the axiom of replacement
  and it is installed as such, on the primitive shelf.  It is load-bearing in the ordinal and cardinal
  layer: \texttt{ord-no-injection-into-set} (hence Zorn's lemma), the
  definition of $\mathtt{card}$, \texttt{card-image-finite},
  \texttt{fun-domain-in-set}.
\item[\texttt{fun-image-set} (proven).]
  $f \in \mathtt{fun}(A, B) \Longrightarrow B \in \mathtt{set}
  \Longrightarrow \mathtt{image}(f, S) \in \mathtt{set}$, with no
  hypothesis on $S$, from separation alone: every member of the image is
  $f(x)$ with $f(x)$ defined, hence $x \in A$ (\texttt{fun-domain-apply-def}),
  hence $f(x) \in B$, so the image is a subclass of the set $B$
  (\texttt{subclass-of-set-is-set}).  The form with $S \in \mathtt{set}$
  and no sethood of the codomain is not derivable without replacement, for
  the reason given under \emph{Pairs and tuples}.
\end{description}

Where a proof forms an image under a set function into a set, which is
every image in the analysis and algebra of the library, the second fact
suffices and the bill need not mention replacement.  Most of those proofs
still cite \texttt{image-set}, directly or through \texttt{card-image-finite},
whose map is an arbitrary object; routing them through
\texttt{fun-image-set} is recorded as open work.

A functoid is never the first argument of \texttt{image}: the axioms
quantify $\phi$ over objects, the instantiation rule refuses a functoid, and
the library forms images only of classes of pairs.  The image of a proper
class under a functoid, when it is wanted, is written as the class it is, by
separation or by a \texttt{vnb-lambda} over a set.

\begin{remark}
$\texttt{vnb-lambda}(x \in \mathtt{nn}, \mathtt{set})$ is a syntactically
valid term, and so is its application
$(\texttt{vnb-lambda}(x \in \mathtt{nn}, \mathtt{set}))(0)$.  $\beta$-reduction
formally returns the term $\mathtt{set}$.  But $\mathtt{set} \notin
\mathtt{set}$, so by the definition of $\mathtt{image}$ above
$\mathtt{set} \notin \mathtt{image}(\texttt{vnb-lambda}(x \in \mathtt{nn},
\mathtt{set}), \mathtt{nn})$, and the contradiction $\mathtt{set} \in
\mathtt{set}$ is avoided.  This does place a burden on the handling of the
\texttt{image} functoid, and it is carried in three places: the conjunct
$w \in \mathtt{set}$ of \texttt{image-membership-iff}; the refusal of the
instantiation rule to take a functoid as an instance, so that the axiom is
never read at $\phi(x) := \mathtt{set}$; and the value half of the
$\beta$ licence, under which the reduction above owes
$\mathtt{set} \in \mathtt{set}$ and \texttt{rfl} cannot close
$\mathtt{set} = \mathtt{set}$ from it.
\end{remark}

\begin{remark}
$\texttt{vnb-lambda}(x \in \mathtt{nn}, [x, \mathtt{set}])$ is also a valid
term.  $\beta$-reduction of
$(\texttt{vnb-lambda}(x \in \mathtt{nn}, [x, \mathtt{set}]))(0)$ formally
gives $[0, \mathtt{set}]$, which is undefined: a list denotes only when its
components are sets (\emph{Pairs and tuples} above).  In particular
\texttt{app-graph} proves nothing at such a term, since $[0, \mathtt{set}]
\in \phi$ would require $[0, \mathtt{set}]$ to denote, and
\texttt{forall-elim} posts that definedness as a side goal that nothing
closes.
\end{remark}

Open at the time of writing: the re-routing of the \texttt{image-set} citations
named above; and the rename of \texttt{def-functoid}
(\texttt{docs/functoids-and-functions.md}, \S8).


The \textbf{primitive class constant} \texttt{set} has the following
property: The formula \texttt{$a$ in set} expresses ``$a$ is a set'';
there is no separate predicate form \texttt{set($a$)}.  The key axiom
distinguishing VNB from typed systems: \[
  (a \in b) \;\Longrightarrow\; a \in \mathtt{set}
\]
Only sets can be members of anything: membership forces the \emph{left}
argument to be a set.  The right argument---the containing class---may be
a proper class.  ($3 \in \mathtt{ord}$ is well-formed even though the
class of all ordinals is a proper class.)

Function application is written $f(x_1, \ldots, x_n)$ in both cases:
for a function $f$ this abbreviates $\texttt{apply-function}(f, x_1, \ldots, x_n)$,
and for a functoid it abbreviates $\texttt{apply-functoid}(f, x_1, \ldots, x_n)$.

The subset relation
\[
  A \subseteq B \;\equiv\; \forall x.\; x \in A \Rightarrow x \in B
\]
is a primitive binary predicate written \texttt{A subset B}
(\texttt{subset-def}).  An immediate corollary of
\texttt{membership-implies-sethood} is that \emph{every class is a
subclass of $\mathtt{set}$} (axiom \texttt{subset-set}): only sets can
appear on the left of~$\in$.

The constructors below are \emph{total} functoids: each is defined for
\emph{any} class arguments and always produces a well-defined class.
Whether the result is a set depends on a separate \emph{sethood-closure}
condition on the arguments; those conditions are listed in
Section~\ref{sec:axioms} and Section~\ref{sec:funapp}.
Class comprehension $\{x \mid p\}$ is the sole exception: it is a
\emph{schema}, yielding one instance per syntactically correct formula
$p$, each taking $0$ arguments.

\begin{remark}
Totality is the VNB tradition: \texttt{fun($X$, $Y$)} denotes a
well-defined class for \emph{any} classes $X, Y$.  Three properties
must be kept separate:
\begin{itemize}
\item \emph{Definability} --- always.
\item \emph{Non-emptiness} --- only when $X \in \mathtt{set}$.
  Functions in VNB are sets of ordered pairs, so a proper-class
  domain forces $\mathtt{fun}(X) = \varnothing$; the existential in
  \texttt{is-fun-def} ($\exists A \in \mathtt{set}.\; f \in
  \mathtt{fun}(A)$) is load-bearing on this point.
\item \emph{Sethood} --- a set when $X, Y \in \mathtt{set}$
  (\texttt{fun-set-iff}).  The converse fails: $\mathtt{fun}(\emptyset, Y)$
  is a set for every class $Y$.
\end{itemize}
Functoids exist precisely for operations that genuinely act on
proper classes (e.g.\ $\mathtt{power}: \mathtt{CLASS} \to
\mathtt{CLASS}$): they cannot be encoded as elements of any
$\mathtt{fun}(\cdots)$.
\end{remark}

\begin{remark}
A related discipline: when we add axioms about a constructor we state
them at the abstract level---how \texttt{fun}, \texttt{pair},
\texttt{tuples}, \texttt{list}, etc.\ \emph{behave}---never in terms
of how they might be concretely encoded as sets in the underlying
set theory.  The user should never have to know whether ordered pairs
are Kuratowski or Wiener, or whether \texttt{nn} is built from von
Neumann or Zermelo numerals.

The axioms on lists fix a list by its length and its entries, and say
nothing about whether a list may coincide with an object of another
kind.  One such coincidence is excluded explicitly: the primitive
\texttt{nn-not-a-list} (\texttt{structure-library/list-recursion.scm})
states that no natural number is a list,
\[ \forall n.\; n \in \mathtt{nn} \Rightarrow
   \neg\,(n \in \mathtt{tuples}(\mathtt{set})). \]
It is stated in the same abstract terms, as a fact about $\mathtt{nn}$
and $\mathtt{tuples}$ and not about an encoding.  With it, a natural
number and a list are provably distinct; for example
$[2,3] \neq 23$.  It is on the primitive shelf by an explicit
foundational decision: it records that a natural number is a number and
not a finite sequence, which is part of what the two notions mean here.  The
strand-space theory, which uses naturals as atomic message tokens,
relies on it (\texttt{strand-tokens-nn.scm}).  It is consistent relative
to NBG with the rest of the base theory: with von~Neumann naturals and
lists encoded as Kuratowski pairs, every list axiom holds and no natural
is a list (the argument is in the comment above the axiom).
\end{remark}

\begin{remark}[Variadic functoids]
Several of the constructors that follow---\texttt{union},
\texttt{intersection}, \texttt{cartesian}, and the type constructor
\texttt{fun}---are \emph{variadic functoids}.  The number of arguments
varies between invocations but is fixed when the expression is written:
any \emph{finite} number is allowed.

This is distinct from variadic \emph{functions} of type
$\mathtt{fun}(\mathtt{tuples}(A), B)$, which take a single set-level
tuple as their one argument (see the remark on $\mathtt{tuples}$ in
Section~\ref{sec:constructor-examples}), and from fixed-arity functoids
such as ring addition
$\mathtt{add}(r) \in \mathtt{fun}(\mathtt{cartesian}(\mathtt{carr}(r), \mathtt{carr}(r)), \mathtt{carr}(r))$,
which have a single argument slot of fixed Cartesian type.
\end{remark}

\medskip\noindent The remaining constructors of this section are listed
below.  Each entry states the defining membership property and points to
the corresponding proof commands in
Section~\ref{sec:nary-proof-commands}.

\begin{itemize}

\item \texttt{union($A_1$,\ldots,$A_n$)}
  $x \in C$ iff $x \in A_i$ for some $i$.
  Axiom \texttt{union-membership} states this IFF for $n=2$.
  At any arity, the kernel rule \texttt{pi-union-intro!} reduces a goal
  $x \in \text{union}(A_1,\ldots,A_n)$ to the single subgoal $x \in A_k$
  for a chosen index $k$; \texttt{pi-union-elim!} consumes a hypothesis
  $x \in \text{union}(A_1,\ldots,A_n)$ and case-splits it into $n$
  branches, each with $x \in A_i$ assumed.  See
  Section~\ref{sec:nary-proof-commands}.
  The macete \texttt{union-decompose} rewrites
  $x \in \text{union}(A_1,\ldots,A_n)$ to the disjunction
  $(x \in A_1) \lor \cdots \lor (x \in A_n)$ in one step.
  Result is a set when all $A_i \in \mathtt{set}$
  (\texttt{union-set-closure}; n-ary by induction).

\item \texttt{intersection($A_1$,\ldots,$A_n$)}
  $x \in C$ iff $x \in A_i$ for all $i$.
  Axiom \texttt{intersection-membership} states this IFF for $n=2$.
  At any arity, the kernel rule \texttt{pi-intersection-intro!} reduces
  a goal $x \in \text{intersection}(A_1,\ldots,A_n)$ to $n$ subgoals
  $x \in A_i$; \texttt{pi-intersection-elim!} extracts $x \in A_k$ from
  a hypothesis $x \in \text{intersection}(A_1,\ldots,A_n)$ for a chosen
  index $k$.  See Section~\ref{sec:nary-proof-commands}.
  The macete \texttt{intersection-decompose} rewrites
  $x \in \text{intersection}(A_1,\ldots,A_n)$ to the conjunction
  $(x \in A_1) \land \cdots \land (x \in A_n)$ in one step.
  Result is a set when any $A_i \in \mathtt{set}$
  (\texttt{intersection-set-closure}).

\item \texttt{complement-in($A$, $B$)}
  Relative complement $A \setminus B$: $x \in C$ iff $x \in A$ and
  $x \notin B$ (\texttt{complement-in-membership}).
  $B$ need not be a set; result is a set when $A \in \mathtt{set}$
  (\texttt{complement-in-set-closure}).

\item \texttt{cartesian($A_1$,\ldots,$A_n$)}
  $x \in C$ iff $x = [a_1,\ldots,a_n]$ with $a_i \in A_i$.
  Result is a set when all $A_i \in \mathtt{set}$ (\texttt{cartesian-set-iff});
  for $n=2$ the converse holds when both factors are inhabited
  (\texttt{cartesian-set-converse-1}, \texttt{cartesian-set-converse-2}).

\item \texttt{tuples($A$)}
  $x \in C$ iff $x = [a_1,\ldots,a_n]$ with all $a_i \in A$,
  for some $n \ge 0$.
  Result is a set when $A \in \mathtt{set}$ (\texttt{tuples-sethood}).

\item \texttt{fun($A$)}, \texttt{fun($A$, $B$)}
  The unary $\mathtt{fun}(A)$ is the class of all total functions with
  domain $A$ (no codomain restriction); proper class in general.  The
  binary $\mathtt{fun}(A,B)$ is the restriction
  $\{f \in \mathtt{fun}(A) : \forall x \in A.\;f(x) \in B\}$.
  $\mathtt{fun}(A,B) \in \mathtt{set}$ when $A, B \in \mathtt{set}$
  (\texttt{fun-set-iff}); $\mathtt{fun}(A) \in \mathtt{set}$ never holds
  in general because the codomain is unconstrained.
  See Section~\ref{sec:funapp}.

  The axiom \texttt{fun-domain-apply-def} reads
  $f \in \mathtt{fun}(A) \Rightarrow ((f\,x) = (f\,x) \Leftrightarrow x \in A)$.
  Recall VNB equality is partial: $(t = t)$ is the definedness predicate
  for $t$, so the iff says \emph{$f(x)$ is defined exactly when $x \in A$}.
  Consequence: the witness $A$ is unique, justifying ``the'' domain of $f$.

\item \texttt{is-fun($f$)}
  The predicate ``$f$ is a function'': there exists a set $A$ with
  $f \in \mathtt{fun}(A)$.  Axiom \texttt{is-fun-def}:
  $\mathit{is\text{-}fun}(f) \Leftrightarrow \exists A \in \mathtt{set}.\;
   f \in \mathtt{fun}(A)$.

\item \texttt{dom($f$)}
  The \emph{domain} of $f$ as a class.  Characterised by membership
  (\texttt{dom-membership}):
  $x \in \mathtt{dom}(f) \Leftrightarrow x \in \mathtt{set}
   \land (f\,x) = (f\,x)$, i.e.\ $x$ is a set on which $f$ is defined.
  When $f \in \mathtt{fun}(A)$, $\mathtt{dom}(f)$ and $A$ agree on
  members (\texttt{dom-fun-membership}).  When additionally
  $A \in \mathtt{set}$, $\mathtt{dom}(f) \in \mathtt{set}$ and
  $\mathtt{dom}(f) = A$ as sets, by extensionality
  (\texttt{dom-of-fun}).

\item \texttt{res($f$, $B$)}
  The \emph{restriction} of $f$ to $B$.  Semantically
  $\mathtt{res}(f, B) = \lambda x \in B.\;f(x)$.  Total as a constructor
  (defined for all $f$ and $B$ -- bongo-style); typed by:
  \begin{itemize}
  \item \texttt{res-typing}:
    $f \in \mathtt{fun}(A) \land B \subseteq A
    \Rightarrow \mathtt{res}(f, B) \in \mathtt{fun}(B)$.
  \item \texttt{res-apply}:
    $f \in \mathtt{fun}(A) \land B \subseteq A \land x \in B
    \Rightarrow (\mathtt{res}(f, B))(x) = f(x)$.
  \item \texttt{res-codomain} (derived): codomain inherits from
    $f \in \mathtt{fun}(A, C)$.
  \end{itemize}
  Outside the case $B \subseteq A$ the term exists but carries no
  $\mathtt{fun}$-typing guarantees.

\item \texttt{partial-fun($A$)}, \texttt{partial-fun($A$, $C$)}
  Partial functions on $A$.  Members are total functions whose domain
  is some subset of $A$:
  $\mathit{partial\text{-}fun}(A) = \{f \mid \exists B \in \mathtt{power}(A).
   \; f \in \mathtt{fun}(B)\}$
  (\texttt{partial-fun-membership}, generally a proper class because
  the codomain is unconstrained).  The binary form fixes a codomain:
  $\mathit{partial\text{-}fun}(A, C) = \{f \mid \exists B \in \mathtt{power}(A).
   \; f \in \mathtt{fun}(B, C)\}$
  (\texttt{partial-fun-binary-membership}).  Sethood
  (\texttt{partial-fun-binary-sethood}):
  $A, C \in \mathtt{set} \Rightarrow \mathit{partial\text{-}fun}(A, C)
   \in \mathtt{set}$ (stated as axiom; provable once a class-union
  principle over a set-indexed family is in place).

\item \texttt{\{$x$ in $A$: $p$\}} / \texttt{\{$x$ in $A$ | $p$\}}
  Separation $\mathtt{sep}(x, A, p) = \{x \in A \mid p\}$: $x \in C$
  iff $x \in A$ and $p[x]$ holds.  Because $p$ is genuinely a schema
  (a formula, not a first-order term), the characterization is given
  by kernel rules rather than a single FOL axiom:
  \texttt{pi-sep-sethood!} ($\mathtt{sep}(x, A, p) \in \mathtt{set}$
  reduces to $A \in \mathtt{set}$);
  \texttt{pi-sep-mem-intro!} ($y \in \mathtt{sep}(x, A, p)$ splits
  into $y \in A$ and $p[x:=y]$);
  \texttt{pi-sep-mem-elim!} (assumption $y \in \mathtt{sep}(x, A, p)$
  yields both $y \in A$ and $p[x:=y]$ in context).
  Result is a set when $A \in \mathtt{set}$.

\item \texttt{\{$x$ | $p$\}}
  Class comprehension $\mathtt{comp}(x, p) = \{x \mid p\}$: $y \in C$
  iff $y \in \mathtt{set}$ and $p[x:=y]$ (members of any class are
  sets).  Same schema treatment as separation:
  \texttt{pi-comp-mem-intro!} and \texttt{pi-comp-mem-elim!}.
  Proper class in general unless provably bounded.

\item \texttt{big-union($z$, $A$, body)}
  Big-union over an index set:
  \[
    \mathit{big\text{-}union}(z, A, \mathit{body})
    \;=\; \bigcup_{z \in A} \mathit{body}
    \;=\; \{x \mid \exists z \in A.\; x \in \mathit{body}\}
  \]
  $z$ is bound in $\mathit{body}$; $A$ is the index class.
  As with separation, $\mathit{body}$ is a class schema in $z$, so the
  characterization is by kernel rules:
  \texttt{pi-big-union-sethood!}
  ($\mathit{big\text{-}union}(z, A, \mathit{body}) \in \mathtt{set}$
   reduces to $A \in \mathtt{set}$ and
   $\forall z \in A.\; \mathit{body} \in \mathtt{set}$);
  \texttt{pi-big-union-mem-intro!} (witness $w$:
   $x \in \mathit{big\text{-}union}(z, A, \mathit{body})$ splits into
   $w \in A$ and $x \in \mathit{body}[z := w]$);
  \texttt{pi-big-union-mem-elim!} (assumption gives a fresh
   eigenvariable $e$ with $e \in A$ and $x \in \mathit{body}[z := e]$).

  \begin{remark}
  Lists, numbers and functions are objects of the single-sorted
  universe: a list that is defined is a set (all its components are
  sets, and so is any component of a list belonging to a class,
  axiom \texttt{tuple-members-are-sets}), and the number classes
  \texttt{nn}, \texttt{zz}, \texttt{qq}, \texttt{rr}, \texttt{cc} are
  primitive constants (Section~\ref{sec:primitives}), whose elements are
  characterised by axioms and not by a set-theoretic construction.  The binder form $\mathit{big\text{-}union}(z, A,
  \mathit{body})$ makes the family structure explicit: $A$ is the
  index, $z \mapsto \mathit{body}$ is intended to be set-valued.
  The ZF-style $\bigcup S$ over a loose ``set of sets'' is the special
  case $\mathit{big\text{-}union}(s, S, s)$ and is not taken as
  primitive.
  \end{remark}

\item \texttt{power($A$)}, \texttt{power($A$, $B$)}
  Unary \texttt{power($A$)} is the power set: $x \in C$ iff
  $x \in \mathtt{set}$ and $x \subseteq A$ (\texttt{power-set}).
  Binary \texttt{power($A$, $B$)} is the function-space exponentiation
  $A^B$, axiomatized by \texttt{power-exp}:
  \[
     \mathtt{power}(A, B) \;=\; \mathtt{fun}(B, A)
  \]
  i.e.\ functions $B \to A$.  This convention matches the arithmetic
  evaluator: $\texttt{power}(2, 3) = 2^3 = 8$ counts functions from
  a 3-set to a 2-set.  Sethood and membership inherit from
  \texttt{fun-set-iff} / \texttt{fun-codomain-iff}.

\end{itemize}

\begin{remark}
\texttt{union($A_1$, \ldots, $A_n$)} takes a finite listed union.
For an index-set-driven union, use \texttt{big-union($z$, $A$, body)}
above: the binder form makes the family structure explicit.
\end{remark}

\section{Term constructors}
\label{sec:choice-iota-lambda}

\emph{Term constructors} build syntactic objects whose intended
representation are objects in the VNB universe.

\begin{itemize}
\item \texttt{[$e_1$, \ldots, $e_n$]} ordered $n$-tuple $(e_1, \ldots, e_n)$; always a set
\item \texttt{nth($k$, $e$)} $k$-th component of tuple $e$ (1-indexed); $k$ may be any term;
  \texttt{nth($k$, [$e_1$,\ldots,$e_n$])} reduces to $e_k$ via
  \texttt{cmd-nth-reduce} (Section~\ref{sec:commands})
\item \texttt{length($L$)} number of elements in tuple $L$; a natural number:
  \begin{itemize}
  \item \texttt{length-of-empty}: $\texttt{length}([\,]) = 0$
  \item \texttt{length-in-nn}: $L \in \mathtt{tuples}(\mathtt{set}) \Rightarrow
    \texttt{length}(L) \in \mathtt{nn}$
  \item \texttt{nth-in-range}: $L \in \mathtt{tuples}(A),\; i \in \mathtt{nn},\;
    1 \le i \le \texttt{length}(L) \Rightarrow \texttt{nth}(i,L) \in A$
  \end{itemize}
  The recursive characterisation is supplied by the constructor
  \texttt{cons}$(x,L)$ and the equation \texttt{length-cons},
  $\texttt{length}(\texttt{cons}(x,L)) = \texttt{succ}(\texttt{length}(L))$,
  in \texttt{structure-library/list-recursion.scm}; list induction over
  \texttt{tuples} is a proven theorem, \texttt{tuples-induction}, not a schema
  (Section~\ref{sec:list-recursion}).
\item \texttt{empty-set} the empty set $\emptyset$; always a set; $= \texttt{make-set}([\,])$
\item \texttt{\{$a$\}} singleton $\{a\}$; always a set; $= \texttt{make-set}([$a$])$
\item \texttt{make-set($L$)} unordered finite set from tuple $L \in \mathtt{tuples}(A)$,
  $A \in \mathtt{set}$; always a set:
  \begin{itemize}
  \item \texttt{make-set-membership}: $x \in \texttt{make-set}(L)
    \Longleftrightarrow \exists\,i \in \mathtt{nn}.\;
    1 \le i \le \texttt{length}(L) \land \texttt{nth}(i,L) = x$
    (the $1 \le i \le \texttt{length}(L)$ bounds matter:
    $\texttt{nth}$ out-of-range is unspecified)
  \item \texttt{make-set-sethood}: $L \in \mathtt{tuples}(A),\; A \in \mathtt{set}
    \Rightarrow \texttt{make-set}(L) \in \mathtt{set}$
  \item \texttt{make-set-empty}: $\texttt{make-set}([\,]) = \emptyset$
  \end{itemize}
  The notation $\{a,b,c,\ldots\}$ is sugar for $\texttt{make-set}([a,b,c,\ldots])$;
  repetitions are harmless.
\item \texttt{choice($A$)} primitive choice operator: an element of nonempty class~$A$
  (undefined when $A$ is empty).  Axiom \texttt{choice-axiom}:
  \[
    (\exists x.\; x \in A) \;\Longrightarrow\; \texttt{choice}(A) \in A
  \]
  This is the Bourbaki-style global choice: a single canonical operator
  defined uniformly on all nonempty classes, not just sets.
\item \texttt{iota($x$, $p$)} definite description: the unique object satisfying~$p$;
  a binder.  Defined only when $\exists! x.\; p(x)$ is provable.  The
  characterization is the kernel rule \texttt{pi-iota-def!}, invoked
  with the iota term as argument; it posts two subgoals:
  \begin{enumerate}[label=(\arabic*),leftmargin=1.2cm]
  \item the existence-and-uniqueness obligation
    $\exists x.\;p \land (\forall y.\; p[x:=y] \Rightarrow x = y)$, and
  \item the original goal with the defining property
    $p[x := \texttt{iota}(x, p)]$ added as an assumption.
  \end{enumerate}
  A planned \texttt{define-object} form will wrap this pattern:
\begin{VNBsnip}
(define-object sqtwo (IOTA x (= (* x x) 2)))
; requires a proved unique-existence theorem; in this language,
; spelt out as (FORSOME x (AND (= (* x x) 2)
;                              (FORALL y (IMPLIES (= (* y y) 2) (= x y)))))
; (a FORSOME-UNIQUE macro abbreviation is planned but not yet defined)
; installs sqtwo satisfying: sqtwo * sqtwo = 2
\end{VNBsnip}
\item \texttt{vnb-lambda([$x_1$ in $A_1$, \ldots, $x_n$ in $A_n$], $\mathit{body}$)}
  anonymous multi-argument \emph{set-function}; a binder.
  Each variable has its own independent domain, exactly as in quantification.
  All domains $A_i$ must belong to \texttt{set}; the result is an element
  of $\mathtt{fun}$.

  \begin{quote}
  \textbf{This is not the same construct as \texttt{lambdoid}, and the
  distinction is the one most worth getting right.}  A \texttt{vnb-lambda}
  denotes a \emph{set of ordered pairs} --- an element of some
  $\mathtt{fun}(A,B)$ --- and is an ordinary S-expression.  A
  \texttt{lambdoid} denotes a \emph{functoid}, a tagged record that need not
  belong to any \texttt{fun} class.  They carry different obligations and are
  reduced by different rules: \texttt{vnb-lambda} is typed by
  \texttt{pi-lambda-type!} (short form \texttt{(lam-t)}), which opens the
  $A \in \mathtt{set}$ subgoal a set of pairs requires, and reduced by
  \texttt{pi-lambda-beta!} (\texttt{(lam-b)}); a functoid is reduced by
  \texttt{cmd-functoid-beta} (\texttt{(beta)}) and has no such typing rule.
  Essentially every proof in the shipped library uses \texttt{vnb-lambda};
  none reduces a functoid.  An unadorned \texttt{lambda} is not part of the
  language: the parser rejects it with an error naming both binders.
  \end{quote}

  Its representation is the raw S-expression
    \texttt{(vnb-lambda $x$ $A$ $\mathit{body}$)} or
    \texttt{(vnb-lambda (list $x_1\,\ldots\,x_n$) $A$ $\mathit{body}$)},
    where $A$ is the \emph{domain} and lies outside the binder's scope
    (as in \texttt{sep} and \texttt{big-union}).  The domain is
    \emph{required}: a lambda that does not carry one does not determine
    a function, and without one the typing rule could certify one and the
    same domainless term into $\mathrm{fun}(A,B)$ for every $A$ at once.
    Used in axioms written directly as S-expressions
    (e.g.\ \texttt{ring-prod}, \texttt{cc-normed-field}, \texttt{sum}).
    At the surface, the binder list may carry its domains the way a
    quantifier's does --- \texttt{vnb-lambda([$x_1$ in $A_1$, \ldots,
    $x_n$ in $A_n$], $\mathit{body}$)} --- which is sugar for
    \texttt{vnb-lambda([$x_1,\ldots,x_n$], cartesian($A_1,\ldots,A_n$),
    $\mathit{body}$)} and parses to exactly that S-expression (to
    \texttt{(vnb-lambda $x_1$ $A_1$ $\mathit{body}$)} when $n = 1$).  A
    binder list that is only partly typed is rejected rather than
    completed by guesswork.
    Typing rule \texttt{pi-lambda-type!}: the goal
    $\mathit{vnb\text{-}lambda}(x, A, \mathit{body}) \in
    \mathtt{fun}(A, B)$ reduces to two subgoals,
    $\forall x.\; x \in A \Rightarrow \mathit{body} \in B$ and
    $A \in \mathtt{set}$.  The rule also accepts a
    binder \emph{list} when the domain is a \texttt{cartesian} of the
    same length, discharging the first subgoal componentwise as
    $\forall x_1 \in A_1 \ldots \forall x_n \in A_n.\;
    \mathit{body} \in B$; on a non-Cartesian domain, or a length
    mismatch, it declines and leaves the goal untouched.
    $\beta$-reduction \texttt{pi-lambda-beta!} (short form
    \texttt{(lam-b)}) rewrites
    $\bigl(\mathit{vnb\text{-}lambda}(\overline x, \mathit{body})\bigr)
    (\overline v) = \mathit{body}[\overline x := \overline v]$
    (parallel substitution).
  Cardinal-exponent typing schema (summarized; multi-binder derives from the
  single-binder rule plus uncurrying):
  \begin{multline*}
    \bigl(\forall x_i \in A_i.\;\mathit{body}(x_1,\ldots,x_n) \in B\bigr)
    \;\Longrightarrow\; \\
    \texttt{vnb-lambda}([x_1 \in A_1,\ldots,x_n \in A_n],\,\mathit{body})
      \in \mathtt{fun}(A_1 \times \cdots \times A_n,\,B)
  \end{multline*}

\item \texttt{lambdoid([$x_1$ in $A_1$, \ldots, $x_n$ in $A_n$], $\mathit{body}$)}
  the \emph{functoid} binder --- the class-domain sibling of
  \texttt{vnb-lambda}, and the only surface form that produces a functoid.
  The domains $A_i$ may be arbitrary classes (including proper classes such as
  \texttt{ord}).  The result is a tagged record at the meta-level; it does not
  necessarily belong to any \texttt{fun} class, has no \texttt{lam-t} typing
  rule, and $\beta$-reduces via \texttt{(beta)} rather than \texttt{(lam-b)}.
\end{itemize}

\begin{remark}[Functoids as meta-level objects]
\texttt{lambdoid} produces a \emph{functoid record} at the implementation
level --- a tagged Scheme structure distinct from plain S-expressions and from
VNB sets or classes.  (\texttt{vnb-lambda} does not: it is a plain
S-expression, and a set.)  This design follows the IMPS
principle that every structured entity carries its own recognisable tag; it
prevents the expression-traversal code from ever confusing a functoid with an
ordinary compound term.

Internally, the application
$\texttt{lambdoid}([x_1 \in A_1,\ldots], \mathit{body})(v_1,\ldots)$
is represented as
$(\texttt{apply-functoid}\;\langle\textit{record}\rangle\;v_1\;\cdots)$,
a mixed structure whose \texttt{car} is the symbol \texttt{apply-functoid}
and whose \texttt{cadr} is the record itself.  The surface syntax
$f(x)$ is used for \emph{all} application; the parser emits
\texttt{apply-functoid} only when $f$ is already tagged as a functoid.
\end{remark}

\begin{remark}
Tuples $[e_1,\ldots,e_n]$ are primitive here.  One could in principle reduce
them to iterated pairs (Kuratowski style), but we have no interest in doing so:
$[\ ]$ is perfectly well-defined as a term constructor of arity zero, and that
is all we need.
\end{remark}

\begin{remark}
The index $k$ in $\texttt{nth}(k, e)$ may be any term (not only a literal
integer).  The primitive inference \texttt{cmd-nth-reduce} does require a
literal integer to reduce $\texttt{nth}(k, [e_1,\ldots,e_n])$ to $e_k$, but
$\texttt{nth}(i, L)$ with a variable index $i$ is a perfectly valid term and
appears in \texttt{make-set-membership} and \texttt{nth-in-range}.
For example, $\texttt{nth}(3,\,[a,b,c,c,e])$ reduces to $c$.
\end{remark}

\begin{remark}[Indexing convention: lists are 1-indexed]
Tuples / lists in VNB are \emph{1-indexed}: \texttt{nth}$(i, L)$ is valid
for $1 \le i \le \texttt{length}(L)$, and $\texttt{nth}(1,\,[a, b, c]) = a$.
Lists and functions on initial segments of $\mathtt{nn}$ are genuinely
different things, and we do \emph{not} reduce one to the other.  When an
interval constructor is added, the canonical domain for an $n$-element list
is $\mathtt{interval}(1, n)$, \emph{not} $\mathtt{interval}(0, n-1)$.
The 1-indexing convention matches IMPS and is built into every axiom
that mentions \texttt{nth} or \texttt{length} (\texttt{make-set-membership},
\texttt{nth-in-range}, \texttt{length-of-empty}).
\end{remark}

\begin{remark}
\texttt{choice}, \texttt{iota}, \texttt{vnb-lambda}, and \texttt{lambdoid} all
participate correctly in \texttt{free-vars}, \texttt{subst-free},
\texttt{alpha-equiv?}, and macete pattern-matching.  \texttt{lambdoid} produces
a \emph{functoid record} --- a Scheme tagged structure separate from
S-expressions --- so it is never accidentally mistaken for an ordinary term by
the expression machinery.  \texttt{vnb-lambda} needs no such tag: it \emph{is}
an ordinary term, and each of those walkers has an explicit
\texttt{vnb-lambda} case treating it as the binder it is.
\end{remark}

\section{Function application}
\label{sec:funapp}

Function application is written as a compound term: \texttt{$f$($x$)} for
any object $f$ and any argument $x$.  Examples:
\begin{VNBsnip}
a + b               ; addition applied to a, b
NORM(E, v)          ; NORM of Banach space E applied to vector v
GEORGE(s)           ; syntactically valid, but may be undefined
\end{VNBsnip}

\begin{remark}[Curried vs.\ tupled application]
\label{sec:curried-tupled-apply}
The string-syntax form $f(a_1,\,\ldots,\,a_n)$ parses verbatim into the
\emph{curried} S-expression
\[
  (f\ a_1\ a_2\ \ldots\ a_n)
\]
— an $(n+1)$-element compound, head $f$, with one argument per
comma-separated term.  No implicit tupling at the parser level.

However, an operation declared via
\texttt{(op mul (cartesian A A) A)} (Section~\ref{sec:declare-structure})
is typed (Section~\ref{sec:declare-structure}) as the \emph{unary} function
\[
  (\textsc{mul}\ s)\ \in\ \textsc{fun}(\textsc{cartesian}(\mathtt{carr}(s),\,\mathtt{carr}(s)),\ \mathtt{carr}(s))
\]
whose single argument is an element of the Cartesian product.
Likewise the typing axiom \texttt{fun-apply-type} is unary
($f(x) \in B$ when $f \in \textsc{fun}(A,B)$ and $x \in A$).

The kernel's primitive application is unary, so it gives the
$n$-argument form no meaning by itself.  The library defines it, for
$n = 2$ and $n = 3$, as the unary apply to the tuple:
\begin{align*}
  \texttt{apply-tupling-2:}\quad & (f\ a\ b)\ \mathrel{==}\ (f\ (\textsc{list}\ a\ b))\\
  \texttt{apply-tupling-3:}\quad & (f\ a\ b\ c)\ \mathrel{==}\ (f\ (\textsc{list}\ a\ b\ c))
\end{align*}
These are definitions in \texttt{theorem-library/axioms.scm}, not
kernel rules.  No arity-4 definition exists yet; one is added by
declaration when needed.

There is none for $n = 1$, and there must not be one.  A
one-argument application \emph{is} the unary apply.  The equation
\[
  (f\ a)\ \mathrel{==}\ (f\ (\textsc{list}\ a)) \quad\text{for every } f
\]
would give $a = [a]$ for every set $a$: take $f$ the identity
function on $\{a\}$.  Then $f(a) = a$ is defined, so $f([a])$ is
defined, so $[a]$ lies in the domain $\{a\}$, and
\[
  a \;=\; f(a) \;=\; f([a]) \;=\; [a].
\]
Every set would then be a list of length one.  In particular
$[\,] = [[\,]]$ would have length $1$, against
$\mathtt{length}([\,]) = 0$.

Quasi-equality (not the partial $=$) is required because the
definition must hold even when $f$ is not a tuple-domain function:
both sides are then undefined and $\mathrel{==}$ holds vacuously,
whereas $=$ would fail since $t=t$ is the definedness predicate.
When $f$ \emph{is} so typed and the arguments lie in the components,
both sides are defined and \texttt{quasi-eq-def} upgrades
$\mathrel{==}$ to $=$ for substitution.

Together with
\texttt{fun-apply-type} and \texttt{quasi-eq-def}, these axioms
close $(f\ a_1\ \ldots\ a_n) \in B$ from the signature
$f \in \textsc{fun}(\textsc{cartesian}(A_1,\ldots,A_n),\,B)$ and the
memberships $a_i \in A_i$ — close the tupled apply by
\texttt{fun-apply-type}, lift $\mathrel{==}$ to $=$ via the
resulting definedness, then substitute — making the per-structure
``carrier-closed'' axioms (\texttt{monoid-carrier-closed-opr},
\texttt{ring-carrier-closed-add}, etc.) derivable rather than
primitive.
\end{remark}

\texttt{fun}($A$, $B$) (with $A, B \in \mathtt{set}$) is a primitive set constructor: its members are
the total functions from $A$ to $B$.  Since $A, B \in \mathtt{set}$,
\texttt{fun}($A$,$B$) is itself a set (\texttt{fun-set-iff}), and each
$f \in \mathtt{fun}(A,B)$ is a set (\texttt{membership-implies-sethood}).

\begin{remark}
The soundness of $\mathtt{fun}(A,B)$ is confirmed by the standard
set-theoretic construction: a function $f : A \to B$ is a set of pairs
$(a, b) \in A \times B$ that is functional (each $a$ has at most one $b$)
and total (every $a \in A$ appears).  Since $A \times B$ is a set, so
is any subset of it, making $f$ a set.
\end{remark}


\begin{remark}
Category theory motivates a natural extension of $\mathtt{fun}$.  A
functor between large categories---such as the category of all groups or
all metric spaces---maps a proper class of objects to a proper class of
objects.  Since $\mathtt{fun}$ requires a set-valued domain, functors
fall outside $\mathtt{fun}$ entirely.

They are \emph{functoids} (Remark~\ref{rem:functoids-quantified}), the
third kind of entity beside sets and classes, outside the range of the
quantifiers; a \texttt{def-functor} view (Section~\ref{sec:view-as}) is
one.  There is no class $\mathtt{oper}(A,B)$ of them---no such collection
can exist without violating the set/class dichotomy---and there is no
predicate that collects them.  A functoid applied to an argument is reduced by
substitution (\texttt{beta}); it has no typing rule of its own, and a
value that is a set is typed by the typing of the body.  Functoids are
not members of anything---they do not appear on the left of $\in$---so
they are consistent with \texttt{membership-implies-sethood}; the one
place this has to be enforced is in an axiom that turns a functoid's
value into a member, as the image axiom does, and such an axiom carries
the sethood of the value as a conjunct.
\end{remark}

%
The base library installs the following axioms.  The primary form of
\texttt{fun} is the unary $\mathtt{fun}(A)$ -- all total functions whose
domain is $A$, with no codomain restriction; the binary $\mathtt{fun}(A,B)$
is the derived form $\{f \in \mathtt{fun}(A) : \forall x \in A.\;f(x) \in B\}$.

\begin{itemize}
\item \texttt{fun-domain-apply-def}
  $f \in \mathtt{fun}(A) \land x \in A \;\Rightarrow\; f(x) = f(x)$
\item \texttt{fun-domain-extensionality}
  $f,g \in \mathtt{fun}(A) \land (\forall x \in A.\;f(x) = g(x)) \;\Rightarrow\; f = g$
\item \texttt{fun-codomain-iff}
  $f \in \mathtt{fun}(A,B) \;\Leftrightarrow\;
   (f \in \mathtt{fun}(A) \land \forall x \in A.\;f(x) \in B)$
\item \texttt{fun-apply-type}
  $f \in \mathtt{fun}(A,B) \land x \in A \;\Rightarrow\; f(x) \in B$
  (in \texttt{axioms.scm}; derivable from \texttt{fun-codomain-iff})
\item \texttt{fun-set-iff}
  $(A \in \mathtt{set}) \land (B \in \mathtt{set}) \;\Rightarrow\; \mathtt{fun}(A,B) \in \mathtt{set}$
\item \texttt{fun-elements-are-sets}
  $f \in \mathtt{fun}(A,B) \land A \in \mathtt{set} \;\Rightarrow\; f \in \mathtt{set}$
\item \texttt{fun-domain-elements-are-sets}
  $f \in \mathtt{fun}(A) \land A \in \mathtt{set} \;\Rightarrow\; f \in \mathtt{set}$
\item \texttt{cartesian-set-iff}
  $(A \in \mathtt{set}) \land (B \in \mathtt{set}) \;\Rightarrow\; A \times B \in \mathtt{set}$
\item \texttt{tuples-sethood}
  $A \in \mathtt{set} \;\Rightarrow\; \mathtt{tuples}(A) \in \mathtt{set}$
\end{itemize}

\begin{remark}
The constructor $\mathtt{tuples}(A)$ is the class of all finite sequences
$[a_1, \ldots, a_n]$ with each $a_i \in A$, for any $n \ge 0$ (including
the empty sequence $[\,]$).  It fills a gap that $\mathtt{cartesian}$ does
not cover: $\mathtt{cartesian}(A_1, \ldots, A_n)$ fixes the arity $n$ at
the time the expression is written, whereas $\mathtt{tuples}(A)$ lets $n$
range over all natural numbers.

This supports \emph{variadic} functions whose single argument is a tuple
of any length: an element of $\mathtt{fun}(\mathtt{tuples}(A), B)$ takes
one tuple and returns a value in $B$.  This is \emph{not} the same as a
fixed-arity $n$-ary function, which has type
$\mathtt{fun}(\mathtt{cartesian}(A_1, \ldots, A_n), B)$ and is bound to
exactly $n$ argument slots by its type.  The variadic flavour is the
right one when an operator must handle inputs of any length (e.g.\ a sum
over an indexed sequence); the fixed-arity flavour is the right one when
the arity is intrinsic to the operator (e.g.\ ring addition, which is
honestly binary).  Since $\mathtt{tuples}(A) \in \mathtt{set}$ when
$A \in \mathtt{set}$, the variadic case is an ordinary instance of
$\mathtt{fun}(A', B)$ with $A' = \mathtt{tuples}(A)$.
\end{remark}

\subsection{The constant-head registry}
\label{sec:constant-registry}

A compound term \texttt{$h$($a_1,\ldots,a_n$)} has a head $h$ of one
of three forms.
\begin{itemize}
\item A \emph{registered head}: a symbol in the constant-head registry
  below.  This covers the kernel operators (\texttt{cartesian},
  \texttt{union}, \texttt{in}, \ldots), structure accessors, defined
  functions, functoids and predicates.  It names something fixed and is
  never substituted for.
\item A \emph{compound term}, as in \texttt{mul(s)(a, b)}, whose head is
  \texttt{mul(s)}.  Its variables are collected and substituted like
  those of any other subterm.
\item Any other symbol: a \emph{variable} applied as a function.  It
  \emph{is} substituted for, and counted as free.
\end{itemize}
In this section ``constant'' means a registered head, not only a
0-ary constant.

Substitution (\texttt{subst-free}) and free-variable collection
(\texttt{free-vars}) must tell the first and third forms apart, and
nothing in the S-expression \texttt{($h$ \ldots)} marks which a symbol
$h$ is.

The \emph{constant-head registry} supplies the missing information.
Every constant head is registered with a kind tag; any head absent from
the registry is treated as a variable.  The kinds are:
\begin{itemize}
\item \texttt{operator} -- a kernel term- or formula-forming head
  (\texttt{+}, \texttt{in}, \texttt{forall}, \ldots), seeded at load
  time by \texttt{wff.scm};
\item \texttt{accessor} -- a structure field selector, registered by
  \texttt{declare-structure};
\item \texttt{defined-fn} -- a defined function or constant, registered
  by \texttt{add-definition!};
\item \texttt{functoid} -- registered by \texttt{def-functoid};
\item \texttt{predicate} -- registered by \texttt{def-predicate}.
\end{itemize}
The registry is therefore not fixed: it grows as a theory is built, so
declaring a structure or adding a definition extends the set of symbols
that count as constant heads.

\begin{remark}[The case-fold trap, and how the library forbids it]
\label{rem:case-fold}
The reader case-folds symbols.  A bound variable whose name case-folds
to a registered constant cannot be applied as a function --- the head
\texttt{($v$ \ldots)} reads as the constant, not the variable, scope-blind,
silently changing a formula's meaning.  Two safeguards keep the installed
library clean.  First, structure accessors are given \emph{distinctive}
names (\texttt{carr}, \texttt{pts}, \texttt{dist}, \texttt{iden}, \texttt{mul},
\texttt{vec}, \ldots) precisely so they never collide with the natural
element variables (\texttt{a}, \texttt{x}, \texttt{d}, \texttt{e}, \ldots).
Second, \texttt{constant-binder-audit} walks every installed formula and
\emph{hard-fails the load}, with a deliberately loud banner, if any binder
variable is named like a registered \texttt{accessor}, \texttt{operator},
\texttt{defined-fn}, \texttt{functoid}, or \texttt{predicate}.  Its companion
\texttt{case-fold-audit} catches the binder-over-binder case, and
\texttt{functoid-binder-audit} the same collision one level deeper, in a
\texttt{def-functoid} body --- which \texttt{case-fold-audit} cannot see, since
it walks the theorem table only.

The three do \emph{not} enforce alike, and the difference matters to anyone
reading a load log.  \texttt{constant-binder-audit} and
\texttt{functoid-binder-audit} hard-fail the load.  \texttt{case-fold-audit} is
\emph{warn-only}: it prints a \texttt{CASE-FOLD WARNING} banner naming each
offending formula and the load continues.  So the collision class is turned into
a build error where a binder collides with a registered constant, and into a
warning that must actually be read where one binder shadows another.
\end{remark}

\section{Constructor examples}
\label{sec:constructor-examples}

Quick reference for every class and term constructor.  Full descriptions
appear in Sections~\ref{sec:class-set-discipline}
and~\ref{sec:choice-iota-lambda}.

\paragraph{Functoids (class constructors).}  Each is defined for any class
arguments; the result is always a well-defined class.  Each row gives the
condition on $x$ for membership in the constructed class $C$, and the
sethood-closure condition on the arguments (when there is one).

\begin{center}
\footnotesize
\begin{tabular}{@{}p{4.5cm}p{4.8cm}p{3.0cm}@{}}
\toprule
Expression & $x \in C$ iff \ldots & $C \in \mathtt{set}$ when \\
\midrule
\texttt{union($A_1$,\ldots,$A_n$)} & $x \in A_i$ for some $i$ & all $A_i \in \mathtt{set}$ \\[2pt]
\texttt{intersection($A_1$,\ldots,$A_n$)} & $x \in A_i$ for all $i$ & any $A_i \in \mathtt{set}$ \\[2pt]
\texttt{complement-in($A$,$B$)} & $x \in A$ and $x \notin B$ & $A \in \mathtt{set}$ \\[2pt]
\texttt{cartesian($A_1$,\ldots,$A_n$)} & $x = [a_1,\ldots,a_n]$, $a_i \in A_i$ & all $A_i \in \mathtt{set}$ \\[2pt]
\texttt{tuples($A$)} & $x = [a_1,\ldots,a_n]$, $a_i \in A$, $n \ge 0$ & $A \in \mathtt{set}$ \\[2pt]
\texttt{fun($A$,$B$)} & $x$ is a total function from $A$ to $B$ & $A, B \in \mathtt{set}$ \\[2pt]
\texttt{\{$x$ in $A$: $p$\}} & $x \in A$ and $p(x)$ holds & $A \in \mathtt{set}$ \\[2pt]
\texttt{\{$x$ | $p$\}} & $p(x)$ holds & rarely \\[2pt]
\texttt{power($A$)} & $x \in \mathtt{set}$ and $x \subseteq A$ & $A \in \mathtt{set}$ \\
\bottomrule
\end{tabular}
\end{center}

\paragraph{Term constructors.}

\begin{center}
\footnotesize
\begin{tabular}{@{}p{5.2cm}p{7.1cm}@{}}
\toprule
Expression & Key property \\
\midrule
\texttt{empty-set} & empty set $\emptyset$; always a set; no members \\[2pt]
\texttt{\{$a$\}} & singleton $\{a\}$; always a set \\[2pt]
\texttt{\{$a$, $b$, \ldots\}} & finite set via \texttt{make-set}; repetitions harmless \\[2pt]
\texttt{[$e_1$, \ldots, $e_n$]} & ordered $n$-tuple; always a set \\[2pt]
\texttt{nth($k$, $e$)} & $k$-th component (1-indexed); reduces via \texttt{cmd-nth-reduce} \\[2pt]
\texttt{choice($A$)} & $(\exists x.\;x\in A) \Rightarrow \texttt{choice}(A) \in A$ \\[2pt]
\texttt{iota($x$, $p$)} & unique object satisfying $p$; requires $\exists!\,x.\;p(x)$ \\[2pt]
\texttt{vnb-lambda([$x_1$ in $A_1$,\ldots], $b$)} & multi-arg set-function, an element of \texttt{fun}; domains in \texttt{set}; typed by \texttt{(lam-t)}, $\beta$-reduces via \texttt{(lam-b)} \\[2pt]
\texttt{lambdoid([$x_1$ in $A_1$,\ldots], $b$)} & functoid; domains may be proper classes; no \texttt{(lam-t)}, $\beta$-reduces via \texttt{(beta)} \\
\bottomrule
\end{tabular}
\end{center}

%% =========================================================

\section{Destructuring quantification}

The syntax extends the bounded-variable binding form
\texttt{[$x$ in $A$]} to allow the bound variable to be a
\emph{tuple pattern}:
\[
  \texttt{[[$x_1$, \ldots, $x_k$] in $A$]}
\]
Here $A$ is \emph{any class} --- a Cartesian product, a number system,
a separation class, a declared structure, or any user-defined class.

This is notation.  \texttt{make-wff} expands it once, when the formula is
built, into a quantifier over a single new variable~$t_0$ guarded by
$t_0 \in A$, and replaces each name $x_i$ in the body by the $i$-th
component of~$t_0$:
\[
  x_i \;\longmapsto\; \texttt{nth}(i,\; t_0)
\]
When $A$ is a declared structure, the component is written with the
structure's $i$-th accessor instead, which is the form in which the
library states its theorems.  The scope of $x_1, \ldots, x_k$ is the
quantifier body.  The expansion is the whole mechanism: nothing is
attached to the formula, the names $x_i$ do not survive it, and a formula
is displayed in its expanded form.

\paragraph{Examples.}

Ranging over a Cartesian product (no structure involved):
\begin{VNBsnip}
forall([[x, y] in cartesian(RR, RR)], x + y = y + x)

; Expands to (internal form):
(FORALL t_0
  (IMPLIES (IN t_0 (CARTESIAN RR RR))
    (= (+ (NTH 1 t_0) (NTH 2 t_0))
       (+ (NTH 2 t_0) (NTH 1 t_0)))))
\end{VNBsnip}

Ranging over a declared structure.  The components are written with the
accessors of \texttt{metric-space}, \texttt{pts} and \texttt{dist}:
\begin{VNBsnip}
forall([[pp, dd] in metric-space], forall([u in pp], dd(u, u) = 0))

; Expands to (internal form):
(FORALL ms
  (IMPLIES (IN ms METRIC-SPACE)
    (FORALL u (IMPLIES (IN u (PTS ms))
      (= ((DIST ms) u u) 0)))))
\end{VNBsnip}

Destructuring bindings compose freely with plain and bounded bindings
(Section~\ref{sec:multiquant}):
\begin{VNBsnip}
forall([[pp, dd] in metric-space, u in pp], dd(u, u) = 0)
\end{VNBsnip}



%% =========================================================
\chapter{Proofs}
\label{ch:proofs}

\section{Sequents and deduction graphs}\label{section:sequents}

A \emph{sequent} is a pair $(L, B)$, where $L$ is a list of formulas
called \emph{assumptions} and $B$ is a single formula called the
\emph{assertion}.  A sequent is represented as
\[
  A_1,\, A_2,\, \ldots,\, A_n \;\Longrightarrow\; B
\]
where $A_1, \ldots, A_n$ are the assumptions (or hypotheses or the context) and $B$
is the assertion. In the course of a proof, many sequents may be created.

A \emph{deduction graph} is a directed acyclic graph of
\emph{sequent nodes} connected by \emph{inference nodes}.  An inference
node records which rule was applied, which sequent nodes are its
hypotheses, and which sequent node is its conclusion.

A sequent node is \emph{grounded} when it has an inference node all of
whose hypothesis nodes are themselves grounded.  Leaf inferences (with no
hypothesis nodes) are immediately grounded.  Grounding propagates upward.
A proof is \emph{complete} when the root sequent node is grounded.

\begin{remark}
A \emph{proof state} (\texttt{<proof-state>}) is not the same thing as a
deduction graph.  It is a thin wrapper that bundles three items together: the
deduction graph itself, the \emph{root} sequent node (the original goal,
fixed at \texttt{start-proof} time and checked by \texttt{qed}), and the
\emph{focus} sequent node (the subgoal currently being worked on).  The
deduction graph is the mathematical object that grows as rules are applied;
the proof state is the navigational view into it that the interactive session
and the \texttt{cmd-*} commands operate on.  In the interactive session,
\texttt{*ps*} always holds the current proof state.
\end{remark}

%% =========================================================
\section{Kernel operations}\label{sec:kernel-boundary}

A proof in VNB is carried out on a deduction graph, and the graph is
changed only by \emph{kernel operations}.  A kernel operation is applied
to one sequent node $N$.  It adds one inference node: the conclusion is
$N$, the rule recorded is the name of the operation, and the hypotheses
are sequent nodes that the operation builds from $N$ in a way fixed by
the operation.  Nothing else changes the graph.

There are \VNBkernelopcount{} kernel operations.  The list below is
produced from the list the system itself checks each time it starts
(Section~\ref{sec:kernel-checked}).

\VNBkerneloptable

The last \VNBkernelopcountoracle{} are \emph{oracles}.  An oracle runs a
decision procedure, written in Scheme, on the assertion of $N$ and some
of its assumptions.  If the procedure succeeds, the operation adds an
inference node, with no hypotheses except in the case of
\texttt{arith-simplify}, which posts one.  The graph records the name of
the oracle; for \texttt{ineq} and \texttt{sos} it also records the
certificate the procedure found.  The inference is written only if the
oracle's checker accepts it (Section~\ref{sec:inference-checked}).

The kernel operations, together with the axioms
(Section~\ref{sec:proof-debt}), are everything VNB assumes.  A result
reported \texttt{modulo 0} rests on these and on nothing else.

Every other proof command is a \emph{tactic}.  A tactic is a Scheme
procedure that calls kernel operations, possibly many, possibly after a
search.  It has no other access to the graph.  A tactic can fail to find
a proof.  It cannot add an inference that no kernel operation would add.
Section~\ref{sec:tactics-ref} describes the tactics, and
Section~\ref{sec:command-uses} lists, for each command, the kernel
operations it can reach.

\begin{remark}
One command may record different operations according to the form of the
assertion: \texttt{di} records \texttt{and-intro}, \texttt{implies-intro},
\texttt{forall-intro} and so on.  Two commands may record the same
operation.  The name recorded in the graph is that of the operation, not
of the command.
\end{remark}

\subsection{What the operations do}\label{sec:kernel-ops-detail}

%% Inference figures.  \VNBfig{operation}{hypotheses}{conclusion}, and
%% \VNBfigS with the hypotheses stacked one per row (separated by \\); the side
%% conditions follow in a VNBcond list.  Defined here because this is the only
%% section that uses them.
\providecommand{\VNBfig}[3]{\[\frac{\displaystyle #2}{\displaystyle #3}\;\;\mbox{\texttt{#1}}\]}
\providecommand{\VNBnohyp}{\phantom{\Gamma}}
\providecommand{\FV}{\mathrm{FV}}
\providecommand{\aeq}{=_{\alpha}}
\providecommand{\VNBfigS}[3]{\[\frac{\displaystyle\begin{array}{c}#2\end{array}}{\displaystyle #3}\;\;\mbox{\texttt{#1}}\]}
\newenvironment{VNBcond}{\sloppy\begin{description}[font=\normalfont\itshape,leftmargin=0.85cm,itemsep=0.5mm,topsep=0.5mm]}{\end{description}}

In this section $N$ is the node the operation is applied to, with
sequent $\Gamma \Longrightarrow G$.  ``Hypotheses'' are the sequents of
the new hypothesis nodes.  $B[x{:=}t]$ is $B$ with $t$
substituted for the free occurrences of $x$, bound variables being
renamed where needed to avoid capture.  $\FV(E)$ is the set of variables
free in $E$, and $\FV(\Gamma)$ the union over the formulas of $\Gamma$.
An operation that \emph{declines} leaves the graph unchanged.  Posting a
sequent that the graph already contains, up to renaming of bound
variables, returns the existing node, so a new hypothesis may already be
grounded.

An operation that takes an index or a pattern records it beside its
name: the graph holds \texttt{(cartesian-elim 2)}, \texttt{(macete
\textit{name} $L$ $R$)}, \texttt{(ineq \textit{certificate})}.  The name under which the inference is counted, and under which
it appears in the list above, is the leading symbol.  The arguments are
data about one application.

\emph{Definedness.}  VNB terms may fail to denote, and the sign of
denotation is $t = t$: strict equality holds only between defined terms
(Chapter~\ref{sec:language}).  Two operations, \texttt{forall-elim}
and \texttt{reflexivity}, therefore consult a \emph{definedness
certificate}, one test applied to a term and the assumptions of $N$.
$\Gamma$ certifies $t$ when $t$ is a variable or an atomic constant, a
numeral or a ground arithmetic term, a term that $\Gamma$ types
($t \in X$) or equates ($t = u$), a term occurring outside every binder
in an assumption of the form $\in$, $=$, $\leq$ or $<$, a total class
constructor applied to certified arguments, a separation, comprehension
or indexed union whose domain is certified, an accessor of a structure
$\Gamma$ carries, an applied structure operation whose arguments
$\Gamma$ types in its declared domain, or $f(a)$ with $f \in
\mathrm{fun}(D, \cdot)$ and $a \in D$ in $\Gamma$.  A term whose head is
\textsc{choice} or \textsc{iota}, and an application of an untyped
function, are never certified.  The policy is
\texttt{docs/definedness-instantiation-2026-09-18.md}.

\subsubsection*{Reading the figures}

Each operation is stated below as an inference figure.  The sequents above
the line are the hypotheses, the sequent below it is the sequent of $N$,
and the labelled conditions under the figure are the side conditions.  A
figure is the relation that the operation's \emph{checker} tests before
the inference is written (Section~\ref{sec:inference-checked}): an
inference of that form that meets the conditions is accepted, whichever
procedure produced it, and any other is refused.  The operation itself
produces one instance of its figure.  Where it has a choice to make
(which occurrences to rewrite, what to call a new variable), the prose
after the figure says which choice it makes.  A condition labelled
\emph{Not checked} is imposed by the operation but not tested by its
checker.

The figures share five conventions.
%% the conventions: header of rule-checkers-logic.scm, lines 14-27;
%% chk-set=? :82, chk-ctx+? :100, chk-ctx-+? :103, chk-arity? :133,
%% chk-match-under :167, chk-adds-nothing? :126
\begin{itemize}
\item \emph{Contexts are sets.}  $A \aeq B$ means that $A$ and $B$ are
  the same formula up to renaming of bound variables.  Two contexts are
  equal when every formula of each is $\aeq$ to a formula of the other.
  $\Gamma, A$ is $\Gamma \cup \{A\}$, so it is $\Gamma$ itself when
  $\Gamma$ already contains $A$; $\Gamma - F$ is $\Gamma$ with every
  formula $\aeq F$ removed.  The order of a context carries no meaning.
\item \emph{Arity.}  The number of hypotheses is the number drawn.  A
  figure with two sequents above the line refuses an inference with one
  or with three.  Where the number depends on the formula (one hypothesis
  for each component of a tuple), the figure says so.
\item \emph{Assumption.}  ``$F \in \Gamma$'' means that some member of
  $\Gamma$ is $\aeq F$.  $F$ is a member of the context, never a
  subformula of one.  The operation is told which member by its argument;
  the checker is not, and searches $\Gamma$ for one that fits.
\item \emph{Matching.}  A term that occurs in a hypothesis but not in the
  conclusion (the $t$ of \texttt{forall-elim}, the witness of
  \texttt{forsome-intro}, the eigenvariables of \texttt{forall-intro}
  and \texttt{forsome-elim}) is not recorded.  The checker recovers it by
  one-way matching: the conclusion's formula, with a hole in place of the
  variable, is matched against the hypothesis' formula.  A hole is never
  filled by a term that contains a variable bound inside the pattern.
\item \emph{Weakening case.}  \texttt{forall-elim}, \texttt{forsome-elim},
  \texttt{theorem-assumption}, \texttt{if-true} and \texttt{if-false}
  also accept a hypothesis whose assertion is $G$ and whose context is
  included in $\Gamma$.  That is what the operation produces when the
  formula it adds is already in $\Gamma$, and the inference is then an
  instance of \texttt{weakening}.
\end{itemize}

\subsection{Logic, equality, induction and the structure eliminators}

\subsubsection*{Direct inference (command \texttt{di})}

The command records one of five operations, according to the form of $G$
(\texttt{pi-direct-inference!}).  On any other form it declines.
%% pi-direct-inference!, primitive-inferences.scm:75

\begin{center}
\begin{tabular}{@{}ll@{}}
\toprule
$G$ & operation \\
\midrule
$A \wedge B$ & \texttt{and-intro} \\
$A \Rightarrow B$ & \texttt{implies-intro} \\
$A \Leftrightarrow B$ & \texttt{iff-intro} \\
$\neg A$ & \texttt{not-intro} \\
$\forall x.\, B$ & \texttt{forall-intro} \\
\bottomrule
\end{tabular}
\end{center}

%% rule-checkers-logic.scm:347 and-intro, :386 implies-intro, :412 iff-intro, :399 not-intro
\VNBfig{and-intro}{\Gamma \Longrightarrow A \qquad \Gamma \Longrightarrow B}{\Gamma \Longrightarrow A \wedge B}
\VNBfig{implies-intro}{\Gamma, A \Longrightarrow B}{\Gamma \Longrightarrow A \Rightarrow B}
\VNBfig{iff-intro}{\Gamma, A \Longrightarrow B \qquad \Gamma, B \Longrightarrow A}{\Gamma \Longrightarrow A \Leftrightarrow B}
\VNBfig{not-intro}{\Gamma, A \Longrightarrow \textsc{falsity}}{\Gamma \Longrightarrow \neg A}

To prove $\forall x.\, B$ one proves $B$ for an $x$ about which nothing
is assumed.  In a sequent that means a variable that does not occur free
in any assumption: if $\Gamma$ said something about $x$, a proof of $B$
from $\Gamma$ would be a proof about that particular $x$, not about every
$x$.  A variable introduced under this condition is called an
\emph{eigenvariable}\index{eigenvariable}, and the condition is the \emph{eigenvariable
condition}.  \texttt{forall-intro} removes a chain of leading universal
quantifiers in one step, and collects on the way the guards $y \in A$
that restrict them.  Write the assertion as $G_0 = G$ and, for
$1 \leq i \leq k$, $G_{i-1} = \forall x_i.\, C_i$.
%% chk-forall-strip :441, chk-forall-eigen-ok? :491, chk-forall-intro-at :514,
%% forall-intro checker :555, all rule-checkers-logic.scm
\VNBfig{forall-intro}{\Gamma, g_1, \ldots, g_m \Longrightarrow G_k}{\Gamma \Longrightarrow G_0}
\begin{VNBcond}
\item[Chain] $1 \leq k$, and for each $i$ the formula $C_i[x_i{:=}y_i]$
  is either $(y_i \in A_i) \Rightarrow G_i$, in which case $y_i \in A_i$
  is \emph{collected}, or else it is $G_i$.  $g_1, \ldots, g_m$ are the
  collected formulas.
\item[Eigenvariable] Each $y_i$ is a variable, and
  $y_i \notin \FV(\Gamma)$, $y_i \notin \FV(G_{i-1})$ and
  $y_i \notin V_{i-1}$, where $V_0 = \emptyset$ and $V_i$ is
  $V_{i-1} \cup \FV(y_i \in A_i)$ when step $i$ collected a guard and
  $V_{i-1} \cup \{y_i\}$ otherwise.
\item[Matching] The $y_i$ are read off the hypothesis.  The checker
  tries $k$ equal to the whole length of the chain first and then each
  shorter length, since a shorter chain is also a sound inference.
\end{VNBcond}
In words: an eigenvariable is free neither in the context, nor in the
formula still to be generalised at that step, nor in a guard collected
before it (domain included), and the eigenvariables are pairwise
distinct.

The operation (\texttt{peel-foralls-raw}) always takes the whole chain.
It keeps $y_i = x_i$ when $x_i$ is free neither in $\Gamma$ nor in a
guard already collected, and otherwise takes a new variable, spelled as
$x_i$ followed by an underscore and a number.
%% peel-foralls-raw, primitive-inferences.scm:41
A guard is collected only when it has the form $y_i \in A_i$ for the
variable just bound.  Any other antecedent stays in the assertion:
$\forall x.\, P(x) \Rightarrow C$ yields the hypothesis
$\Gamma \Longrightarrow P(x) \Rightarrow C$, and a second direct inference
gives $\Gamma, P(x) \Longrightarrow C$.

\subsubsection*{\texttt{or-intro-left}, \texttt{or-intro-right}
  (commands \texttt{oi-l}, \texttt{oi-r})}

%% rule-checkers-logic.scm:362, :374
\VNBfig{or-intro-left}{\Gamma \Longrightarrow A}{\Gamma \Longrightarrow A \vee B}
\VNBfig{or-intro-right}{\Gamma \Longrightarrow B}{\Gamma \Longrightarrow A \vee B}

\subsubsection*{\texttt{forsome-intro}: supplying a witness (command \texttt{ew})}

%% rule-checkers-logic.scm:574; pi-exists-witness!, primitive-inferences.scm:357
\VNBfig{forsome-intro}{\Gamma \Longrightarrow B[x{:=}t]}{\Gamma \Longrightarrow \exists x.\, B}
\begin{VNBcond}
\item[Matching] $t$ is read off the hypothesis.
\item[Not checked] Well-formedness of $B[x{:=}t]$.  The operation
  checks it and raises an error, recording nothing, on a malformed
  witness.
\end{VNBcond}
The argument of the operation is $t$.  No definedness obligation is
posted for $t$, so the operation is not the mirror of
\texttt{forall-elim} in that respect.

\subsubsection*{\texttt{truth-intro}}

%% rule-checkers-logic.scm:588
\VNBfig{truth-intro}{\VNBnohyp}{\Gamma \Longrightarrow \textsc{truth}}

This operation is unreachable.  No proof command calls the procedure
that records it, and a \textsc{truth} assertion is already closed by
\texttt{assumption} and by \texttt{reflexivity}, both of which accept it.
It is listed because the code is present and a future command could
reach it.

\subsubsection*{Antecedent inference (command \texttt{ai})}

The argument is an assumption $F$ of $N$.  The command records one of
five operations according to the form of $F$
(\texttt{pi-antecedent-inference!}).
%% pi-antecedent-inference!, primitive-inferences.scm:132
\begin{center}
\begin{tabular}{@{}ll@{}}
\toprule
$F$ & operation \\
\midrule
$A \wedge B$ & \texttt{and-elim} \\
$A \vee B$ & \texttt{or-elim} \\
$\neg A$, with $A$ in $\Gamma$ & \texttt{not-elim} \\
$\exists x.\, B$ & \texttt{forsome-elim} \\
$A \Leftrightarrow B$ & \texttt{iff-elim} \\
\bottomrule
\end{tabular}
\end{center}
The command declines when $F$ is not an assumption of $N$, when $F$ is
$\neg A$ and $A$ is not in $\Gamma$, and when $F$ has any other form: an
implication, a universal formula or an atomic formula.  An implication in
the context is used by \texttt{detach} and \texttt{backchain}, a universal
formula by \texttt{forall-elim}.

%% rule-checkers-logic.scm:597 and-elim, :615 or-elim, :633 not-elim, :645 iff-elim
\VNBfig{and-elim}{(\Gamma - F), A, B \Longrightarrow G}{\Gamma \Longrightarrow G}
\begin{VNBcond}
\item[Assumption] $F = A \wedge B$, $F \in \Gamma$.
\end{VNBcond}
\VNBfig{or-elim}{(\Gamma - F), A \Longrightarrow G \qquad (\Gamma - F), B \Longrightarrow G}{\Gamma \Longrightarrow G}
\begin{VNBcond}
\item[Assumption] $F = A \vee B$, $F \in \Gamma$.
\end{VNBcond}
\VNBfig{not-elim}{\VNBnohyp}{\Gamma \Longrightarrow G}
\begin{VNBcond}
\item[Assumption] $\neg A \in \Gamma$ and $A \in \Gamma$.
\end{VNBcond}
\VNBfig{iff-elim}{(\Gamma - F), A \Rightarrow B, B \Rightarrow A \Longrightarrow G}{\Gamma \Longrightarrow G}
\begin{VNBcond}
\item[Assumption] $F = A \Leftrightarrow B$, $F \in \Gamma$.
\end{VNBcond}

%% forsome-elim checker, rule-checkers-logic.scm:665
\VNBfig{forsome-elim}{(\Gamma - F), B[x{:=}y] \Longrightarrow G}{\Gamma \Longrightarrow G}
\begin{VNBcond}
\item[Assumption] $F = \exists x.\, B$, $F \in \Gamma$.
\item[Eigenvariable] $y$ is a variable, and $y \notin \FV(G)$,
  $y \notin \FV(\Gamma - F)$, $y \notin \FV(F)$.  When $x$ is not free
  in $B$ there is no $y$ to test.
\item[Matching] $y$ is read off the hypothesis.
\end{VNBcond}
Nothing else is required: in particular nothing about the other
assumptions of $N$.  The operation names $y$ by the counter
(\texttt{fresh-var}): $x$ followed by an underscore and $n$, with $n$ the
next value for which that variable is free in none of $B$, $G$ and
$\Gamma - F$.  It never reuses $x$ itself.  The weakening case of the
conventions arises here only when $x$ is not free in $B$ and $B$ is
already in $\Gamma - F$.
%% pi-antecedent-inference! FORSOME arm, primitive-inferences.scm:169;
%% fresh-var, expressions.scm:701

\subsubsection*{\texttt{forall-elim}: instantiation of a universal assumption
  (commands \texttt{inst}, \texttt{inst+}; used by \texttt{fact})}

%% forall-elim checker, rule-checkers-logic.scm:702; pi-instantiate!, primitive-inferences.scm:282
\VNBfig{forall-elim}{\Gamma, B[x{:=}t] \Longrightarrow G \qquad \bigl(\Gamma \Longrightarrow t = t\bigr)}{\Gamma \Longrightarrow G}
\begin{VNBcond}
\item[Assumption] $\forall x.\, B \in \Gamma$.  It stays in the context.
\item[Matching] $t$ is read off the first hypothesis.
\item[Definedness] The second hypothesis is present exactly when
  $\Gamma$ does not certify $t$ (the certificate at the head of this
  section).
\item[Not checked] The checker accepts the second hypothesis present or
  absent.  When it is present, the checker requires it to be
  $\Gamma \Longrightarrow u = u$ for some term $u$, and does not require
  $u$ to be $t$.  Definedness of $t$ is therefore enforced by the
  operation alone.
\end{VNBcond}
The arguments are the universal assumption and $t$.  The second
hypothesis is there because a term may be undefined ($f(a)$ for $a$
outside the domain of $f$, a description with no referent), and a
universal formula speaks only of defined values.  $B[x{:=}t]$ is checked
for well-formedness; a malformed term raises an error and leaves the
graph unchanged.  The same operation is recorded, once per term, by the
procedure behind \texttt{apply-thm}, which cites a theorem and then
instantiates it at a list of terms (\texttt{pi-spec!}).

\subsubsection*{\texttt{assumption}: the goal is already in the context
  (command \texttt{ass})}

%% rule-checkers-logic.scm:738
\VNBfig{assumption}{\VNBnohyp}{\Gamma \Longrightarrow G}
\begin{VNBcond}
\item[Assumption] $G \in \Gamma$, or $G$ is \textsc{truth}.
\end{VNBcond}

\subsubsection*{\texttt{theorem-assumption}: citing a theorem or an axiom
  (command \texttt{ta}; used by \texttt{fact})}

%% rule-checkers-logic.scm:750; chk-installed-theorem-name :329
\VNBfig{theorem-assumption}{\Gamma, S \Longrightarrow G}{\Gamma \Longrightarrow G}
\begin{VNBcond}
\item[Installed] $S \aeq$ the formula installed under some name in the
  theorem table.
\end{VNBcond}
The argument of the operation is a name, and $S$ is the formula installed
under it; the operation declines when there is none.  The name is not
recorded with the inference, so the checker looks the statement up in
the table.

This is the only way a formula enters a proof without being proved in
it.  $S$ is an axiom (there are \VNBaxiomcount{} of them, listed in the
catalog), a definition, a proven theorem, or an asserted support.  The
bill printed when a theorem is installed
(Section~\ref{sec:proof-debt}) lists the asserted supports reached
through this operation, directly or through the theorems cited.

\subsubsection*{\texttt{cut}: introducing a lemma (commands \texttt{cut},
  \texttt{have!})}

%% rule-checkers-logic.scm:767; pi-cut!, primitive-inferences.scm:223
\VNBfig{cut}{\Gamma \Longrightarrow A \qquad \Gamma, A \Longrightarrow G}{\Gamma \Longrightarrow G}

The argument is $A$.  It is checked for well-formedness first, and a
malformed lemma raises an error before anything is recorded.

\begin{remark}
The operation never declines.  When $\Gamma$ already contains $A$, the
second hypothesis is $\Gamma \Longrightarrow G$, which is the sequent of
$N$ itself, so the inference has $N$ among its own hypotheses and
grounds nothing.  The same happens when $A$ is $G$.  A command that cuts
a formula already present therefore appears to succeed and leaves the
proof where it was.
\end{remark}

\subsubsection*{\texttt{weakening}: dropping assumptions (command \texttt{wk})}

%% rule-checkers-logic.scm:781; pi-weaken!, primitive-inferences.scm:208
\VNBfig{weakening}{\Delta \Longrightarrow G}{\Gamma \Longrightarrow G}
\begin{VNBcond}
\item[Inclusion] $\Delta \subseteq \Gamma$.
\end{VNBcond}
The operation takes one assumption $F$ of $N$ and posts
$\Delta = \Gamma - F$; it declines when $F$ is not an assumption of $N$.

\subsubsection*{\texttt{detach}: forward modus ponens (command \texttt{detach!})}

%% rule-checkers-logic.scm:790; pi-detach!, primitive-inferences.scm:442
\VNBfig{detach}{\Gamma, B \Longrightarrow G}{\Gamma \Longrightarrow G}
\begin{VNBcond}
\item[Assumption] $A \Rightarrow B \in \Gamma$ and $A \in \Gamma$.
\end{VNBcond}
The argument is the implication.  Both it and its antecedent stay.

\subsubsection*{\texttt{backchain}: reducing the goal to an antecedent
  (commands \texttt{bc}, \texttt{bc*})}

%% rule-checkers-logic.scm:807; pi-backchain!, primitive-inferences.scm:419
\VNBfig{backchain}{\Gamma \Longrightarrow A}{\Gamma \Longrightarrow G}
\begin{VNBcond}
\item[Assumption] $A \Rightarrow G \in \Gamma$.
\end{VNBcond}

\subsubsection*{\texttt{proof-by-contradiction} (command \texttt{pbc})}

%% rule-checkers-logic.scm:824
\VNBfig{proof-by-contradiction}{\Gamma, \neg G \Longrightarrow \textsc{falsity}}{\Gamma \Longrightarrow G}

There are no arguments, and the operation never declines.

\subsubsection*{\texttt{contraposition}}

%% rule-checkers-logic.scm:835
\VNBfig{contraposition}{\Gamma \Longrightarrow \neg B}{\Gamma \Longrightarrow \neg A}
\begin{VNBcond}
\item[Assumption] $A \Rightarrow B \in \Gamma$.
\end{VNBcond}
This operation is unreachable: no proof command calls the procedure that
records it.

\subsubsection*{\texttt{eq-subst}: rewriting the assertion by a context
  equation (command \texttt{subst})}

%% rule-checkers-logic.scm:857; chk-replaced? :241
\VNBfig{eq-subst}{\Gamma \Longrightarrow G'}{\Gamma \Longrightarrow G}
\begin{VNBcond}
\item[Assumption] $\Gamma$ contains $s = t$, $t = s$, $s \equiv t$ or
  $t \equiv s$.
\item[Replacement] $G' \not\aeq G$, and $G'$ is $G$ with some
  occurrences of $s$ replaced by $t$, or with some occurrences of $t$
  replaced by $s$.  A binder whose bound variable the replacement renamed
  is matched with its renamed form.
\item[Not checked] That no replaced occurrence lies in the scope of a
  binder whose variable is free in $s$ or in $t$.  The checker compares
  $G$ and $G'$ position by position and does not test it.
\end{VNBcond}
The argument is the equation $s = t$ (or $s \equiv t$).  The operation
replaces \emph{every} occurrence of $s$ by $t$, in operator position as
well as in argument position, and declines when nothing changes.  It
leaves alone a subterm below a binder whose bound variable occurs free in
$s$ or in $t$, since rewriting there would capture; when $s$ is a
variable it replaces the free occurrences only.  $G'$ is checked for
well-formedness.
%% pi-eq-subst! :563, replace-term :496, primitive-inferences.scm

\begin{remark}
The direction of the rewrite is fixed by the argument, not by the
orientation in which $\Gamma$ happens to hold the equation.  To use a
context equation from right to left, name it the other way round.
\end{remark}

\subsubsection*{\texttt{reflexivity} and \texttt{quasi-reflexivity}
  (commands \texttt{rfl}, \texttt{qrfl})}

%% rule-checkers-logic.scm:879, :890
\VNBfig{reflexivity}{\VNBnohyp}{\Gamma \Longrightarrow t = t}
\begin{VNBcond}
\item[Definedness] $\Gamma$ certifies $t$.  The operation also accepts
  the assertion \textsc{truth}.
\item[Not checked] The definedness condition.  The checker accepts
  $\Gamma \Longrightarrow u = u$ for any term $u$.
\end{VNBcond}
\VNBfig{quasi-reflexivity}{\VNBnohyp}{\Gamma \Longrightarrow t \equiv t}

\texttt{quasi-reflexivity} asks nothing further: quasi-equality holds
between two undefined terms as well.

\subsubsection*{\texttt{if-true}, \texttt{if-false}
  (commands \texttt{if-true}, \texttt{if-false})}

%% chk-if-rule, rule-checkers-logic.scm:900-933
\VNBfig{if-true}{\Gamma \Longrightarrow p \qquad \Gamma, \mathrm{if}(p,a,b) = a \Longrightarrow G}{\Gamma \Longrightarrow G}
\VNBfig{if-false}{\Gamma \Longrightarrow \neg p \qquad \Gamma, \mathrm{if}(p,a,b) = b \Longrightarrow G}{\Gamma \Longrightarrow G}

The argument is the conditional term $\mathrm{if}(p, a, b)$.  It need not
occur in $G$.  The assertion is not rewritten: the branch equation
arrives as an assumption, and it takes an \texttt{eq-subst} to use it.
An argument that is not a conditional term raises an error rather than
declining.

\subsubsection*{\texttt{cartesian-intro}, \texttt{cartesian-elim}
  (commands \texttt{ci}, \texttt{ce})}

%% rule-checkers-logic.scm:939, :963
\VNBfig{cartesian-intro}{\Gamma \Longrightarrow a_1 \in A_1 \quad \cdots \quad \Gamma \Longrightarrow a_n \in A_n}{\Gamma \Longrightarrow [a_1, \ldots, a_n] \in A_1 \times \cdots \times A_n}
\begin{VNBcond}
\item[Arity] The literal tuple and the literal product have the same
  length $n$, and there is one hypothesis per component, in order.
\end{VNBcond}
\VNBfig{cartesian-elim $k$}{\Gamma, a_k \in A_k \Longrightarrow G}{\Gamma \Longrightarrow G}
\begin{VNBcond}
\item[Assumption] $[a_1, \ldots, a_n] \in A_1 \times \cdots \times A_n
  \in \Gamma$, with tuple and product of the same length.  It stays.
\item[Index] $k$ is recorded with the name, and $1 \leq k \leq n$.
\end{VNBcond}

\subsubsection*{\texttt{tuples-intro}, \texttt{tuples-elim}
  (commands \texttt{ti}, \texttt{te})}

$\mathrm{tuples}(A)$ is the class of finite tuples over $A$.
%% rule-checkers-logic.scm:991, :1012
\VNBfig{tuples-intro}{\Gamma \Longrightarrow a_1 \in A \quad \cdots \quad \Gamma \Longrightarrow a_n \in A}{\Gamma \Longrightarrow [a_1, \ldots, a_n] \in \mathrm{tuples}(A)}
\begin{VNBcond}
\item[Arity] One hypothesis per component; none for the empty tuple.
\end{VNBcond}
\VNBfig{tuples-elim $k$}{\Gamma, a_k \in A \Longrightarrow G}{\Gamma \Longrightarrow G}
\begin{VNBcond}
\item[Assumption] $[a_1, \ldots, a_n] \in \mathrm{tuples}(A) \in \Gamma$.
\item[Index] $k$ is recorded with the name, and $1 \leq k \leq n$.
\end{VNBcond}

\subsubsection*{\texttt{union-intro}, \texttt{union-elim}
  (commands \texttt{ui}, \texttt{ue})}

$\mathrm{union}$ takes $n$ arguments.
%% rule-checkers-logic.scm:1038, :1060
\VNBfig{union-intro $k$}{\Gamma \Longrightarrow x \in A_k}{\Gamma \Longrightarrow x \in A_1 \cup \cdots \cup A_n}
\begin{VNBcond}
\item[Index] $k$ is recorded with the name, and $1 \leq k \leq n$.
\end{VNBcond}
\VNBfig{union-elim}{(\Gamma - F), x \in A_1 \Longrightarrow G \quad \cdots \quad (\Gamma - F), x \in A_n \Longrightarrow G}{\Gamma \Longrightarrow G}
\begin{VNBcond}
\item[Assumption] $F = x \in A_1 \cup \cdots \cup A_n$, $F \in \Gamma$.
  It is consumed.
\item[Arity] One hypothesis per set, in order.
\end{VNBcond}

\subsubsection*{\texttt{intersection-intro}, \texttt{intersection-elim}
  (commands \texttt{ii}, \texttt{ie})}

%% rule-checkers-logic.scm:1079, :1098
\VNBfig{intersection-intro}{\Gamma \Longrightarrow x \in A_1 \quad \cdots \quad \Gamma \Longrightarrow x \in A_n}{\Gamma \Longrightarrow x \in A_1 \cap \cdots \cap A_n}
\VNBfig{intersection-elim $k$}{\Gamma, x \in A_k \Longrightarrow G}{\Gamma \Longrightarrow G}
\begin{VNBcond}
\item[Assumption] $x \in A_1 \cap \cdots \cap A_n \in \Gamma$.  It stays.
\item[Index] $k$ is recorded with the name, and $1 \leq k \leq n$.
\end{VNBcond}

\subsubsection*{\texttt{nth-reduce}, \texttt{length-reduce},
  \texttt{functoid-beta} (commands \texttt{nth-r}, \texttt{len-r},
  \texttt{beta})}

The three have the same figure and differ in the contraction.
%% chk-contracts? rule-checkers-logic.scm:284; checkers :1137, :1153, :1174
\VNBfig{}{\Gamma \Longrightarrow G'}{\Gamma \Longrightarrow G}
\begin{VNBcond}
\item[Contraction] $G' \not\aeq G$, and $G'$ is obtained from $G$ by
  contracting redexes, where a redex that a contraction creates may be
  contracted in its turn.  The redexes are, respectively:
  $\mathrm{nth}(k, [t_1, \ldots, t_n])$ with $1 \leq k \leq n$, which
  contracts to $t_k$; $\mathrm{length}([t_1, \ldots, t_n])$, which
  contracts to $n$; and a functoid applied to as many arguments as it
  binds, which contracts to its body with the arguments substituted in
  parallel for the bound variables.
\end{VNBcond}
None takes an argument.  Each operation contracts every redex of its
kind in $G$ and declines when $G'$ is $G$.  \texttt{nth-reduce} raises an
error on an index outside $1 \leq k \leq n$.  Nothing is asked of the
entries: a literal tuple has $n$ places whether or not the terms in them
denote.

\subsubsection*{\texttt{nn-induction} (command \texttt{ni})}

%% rule-checkers-logic.scm:1191
\VNBfig{nn-induction}{\Gamma \Longrightarrow B[n{:=}0] \qquad \Gamma \Longrightarrow \forall n.\, n \in \mathbb{N} \Rightarrow (B \Rightarrow B[n{:=}\mathrm{succ}\,n])}{\Gamma \Longrightarrow \forall n.\, n \in \mathbb{N} \Rightarrow B}
\begin{VNBcond}
\item[Shape] The guard is $n \in \mathbb{N}$ on the variable of the
  outermost quantifier.
\end{VNBcond}
The operation has no argument and declines on any other shape.  In
particular the induction variable must be the outermost one: a statement
that quantifies over something else first is not of this shape.

\subsubsection*{\texttt{transfinite-induction} (command \texttt{tfi})}

%% rule-checkers-logic.scm:1247; chk-tfi-beta-ok? :1241
\VNBfig{transfinite-induction}{\Gamma \Longrightarrow \forall \alpha.\, \bigl(\alpha \in \mathrm{ord} \wedge \forall \beta.\, \beta <_{\mathrm{ord}} \alpha \Rightarrow P[\alpha{:=}\beta]\bigr) \Rightarrow P}{\Gamma \Longrightarrow \forall \alpha.\, \alpha \in \mathrm{ord} \Rightarrow P}
\begin{VNBcond}
\item[Freshness] $\beta$ is a variable, $\beta \neq \alpha$,
  $\beta \notin \FV(P)$, $\beta \notin \FV(\Gamma)$.
\item[Matching] $\beta$ is read off the hypothesis.
\end{VNBcond}
This is strong induction: the hypothesis of the step is $P$ at every
smaller ordinal.

\begin{remark}
The operation \texttt{transfinite-induction} and the axiom of the same
name in \texttt{structure-library/ordinals.scm} are two different
objects.  One is code in the kernel, the other a formula in the catalog.
\end{remark}

\subsubsection*{\texttt{transfinite-induction-3cases} (command \texttt{tfi3})}

%% rule-checkers-logic.scm:1274
\[
\frac{\displaystyle
  \begin{array}{l}
  \Gamma \Longrightarrow P[\alpha{:=}0] \\[2pt]
  \Gamma \Longrightarrow \forall \alpha.\,
    (\alpha \in \mathrm{ord} \wedge P) \Rightarrow P[\alpha{:=}\mathrm{succ}_{\mathrm{ord}}\,\alpha] \\[2pt]
  \Gamma \Longrightarrow \forall \alpha.\,
    \bigl(\mathrm{limit\text{-}ord}(\alpha) \wedge
          \forall \beta.\, \beta <_{\mathrm{ord}} \alpha \Rightarrow P[\alpha{:=}\beta]\bigr)
    \Rightarrow P
  \end{array}}
 {\displaystyle \Gamma \Longrightarrow \forall \alpha.\, \alpha \in \mathrm{ord} \Rightarrow P}
\]
\begin{VNBcond}
\item[Freshness] As for \texttt{transfinite-induction}, with $\beta$ read
  off the third hypothesis.
\end{VNBcond}
The three hypotheses are the base, successor and limit cases, in that
order.

\subsection{Set formation, descriptions and functions}

The three binders of this group carry a formula or a class expression in
their body, with a variable of the binder free in it.  A first-order
axiom would need a variable ranging over formulas, which VNB does not
have, so each is characterised by operations instead.
$\mathrm{sep}(x, A, p)$ is $\{x \in A \mid p\}$,
$\mathrm{comp}(x, p)$ is $\{x \mid p\}$, and
$\mathrm{big\text{-}union}(z, A, b)$ is $\bigcup_{z \in A} b$.

\subsubsection*{\texttt{sep-sethood}, \texttt{sep-mem-intro},
  \texttt{sep-mem-elim}
  (commands \texttt{sep-set}, \texttt{sep-mi}, \texttt{sep-me})}

%% rcs-sep-sethood :107, rcs-sep-mem-intro :160, rcs-sep-mem-elim :167, rule-checkers-schema.scm
\VNBfig{sep-sethood}{\Gamma \Longrightarrow A \in \mathrm{set}}{\Gamma \Longrightarrow \mathrm{sep}(x,A,p) \in \mathrm{set}}
\VNBfig{sep-mem-intro}{\Gamma \Longrightarrow y \in A \qquad \Gamma \Longrightarrow p[x{:=}y]}{\Gamma \Longrightarrow y \in \mathrm{sep}(x,A,p)}
\VNBfig{sep-mem-elim}{(\Gamma - F), y \in A, p[x{:=}y] \Longrightarrow G}{\Gamma \Longrightarrow G}
\begin{VNBcond}
\item[Assumption] $F = y \in \mathrm{sep}(x,A,p)$, $F \in \Gamma$.  It
  is consumed.
\end{VNBcond}

\subsubsection*{\texttt{comp-mem-intro}, \texttt{comp-mem-elim}
  (commands \texttt{comp-mi}, \texttt{comp-me})}

A comprehension is in general a proper class, and its members are sets,
so the domain condition of separation is replaced by membership in
$\mathrm{set}$.
%% rcs-comp-mem-intro :177, rcs-comp-mem-elim :182, rule-checkers-schema.scm
\VNBfig{comp-mem-intro}{\Gamma \Longrightarrow y \in \mathrm{set} \qquad \Gamma \Longrightarrow p[x{:=}y]}{\Gamma \Longrightarrow y \in \mathrm{comp}(x,p)}
\VNBfig{comp-mem-elim}{(\Gamma - F), y \in \mathrm{set}, p[x{:=}y] \Longrightarrow G}{\Gamma \Longrightarrow G}
\begin{VNBcond}
\item[Assumption] $F = y \in \mathrm{comp}(x,p)$, $F \in \Gamma$.  It is
  consumed.
\end{VNBcond}
There is no sethood operation for comprehension.  A comprehension is a
set only when it is separately shown to be bounded by one.

\subsubsection*{\texttt{big-union-sethood}, \texttt{big-union-mem-intro},
  \texttt{big-union-mem-elim}
  (commands \texttt{bu-set}, \texttt{bu-mi}, \texttt{bu-me})}

Write $U$ for $\mathrm{big\text{-}union}(z, A, b)$.
%% rcs-big-union-sethood :194, rcs-big-union-mem-intro :234,
%% rcs-big-union-mem-elim :260, rule-checkers-schema.scm
\VNBfigS{big-union-sethood}{\Gamma \Longrightarrow A \in \mathrm{set} \\ \Gamma \Longrightarrow \forall v.\, v \in A \Rightarrow b[z{:=}v] \in \mathrm{set}}{\Gamma \Longrightarrow U \in \mathrm{set}}
\begin{VNBcond}
\item[Freshness] $v = z$, or $v \notin \FV(A)$.
\item[Not checked] That $v \notin \FV(b)$ when $v \neq z$.  Without it
  the quantifier of the second hypothesis captures a free $v$ of $b$.
  The operation takes $v$ new, free in none of $b$, $A$, $G$ and
  $\Gamma$ (\texttt{pi-big-union-sethood!}, primitive-inferences.scm).
\end{VNBcond}
%% pi-big-union-sethood!, primitive-inferences.scm:1651
\VNBfig{big-union-mem-intro}{\Gamma \Longrightarrow w \in A \qquad \Gamma \Longrightarrow x \in b[z{:=}w]}{\Gamma \Longrightarrow x \in U}
\begin{VNBcond}
\item[Matching] $w$ is read off the first hypothesis.
\end{VNBcond}
\VNBfig{big-union-mem-elim}{(\Gamma - F), e \in A, x \in b[z{:=}e] \Longrightarrow G}{\Gamma \Longrightarrow G}
\begin{VNBcond}
\item[Assumption] $F = x \in U$, $F \in \Gamma$.  It is consumed.
\item[Eigenvariable] $e$ is a variable free in none of $G$, $A$, $b$,
  $x$ and $\Gamma - F$.
\end{VNBcond}
\texttt{big-union-mem-intro} takes the witness $w$ as its argument, and
checks $x \in b[z{:=}w]$ for well-formedness.

\subsubsection*{\texttt{iota-def}, \texttt{iota-in-elim}
  (commands \texttt{iota-d}, \texttt{iota-e})}

$\mathrm{iota}(x, p)$ is the unique $x$ satisfying $p$, and denotes
nothing when there is no such $x$ or more than one.  Write $\iota$ for
$\mathrm{iota}(x, p)$.  Both operations take the term as their argument;
it need not occur in $G$.  Both grant the defining property
$p[x{:=}\iota]$ as an assumption, and they differ in what pays for it.
%% rcs-iota-def :308, rcs-iota-in-elim :372, rcs--context-defines? :350, rule-checkers-schema.scm
\VNBfigS{iota-def}{\Gamma \Longrightarrow \exists x.\, \bigl(p \wedge \forall y.\, p[x{:=}y] \Rightarrow x = y\bigr) \\ \Gamma, p[x{:=}\iota] \Longrightarrow G}{\Gamma \Longrightarrow G}
\begin{VNBcond}
\item[Freshness] $y$ is a variable, $y \neq x$, $y \notin \FV(p)$.
\end{VNBcond}
\VNBfig{iota-in-elim}{\Gamma, p[x{:=}\iota] \Longrightarrow G}{\Gamma \Longrightarrow G}
\begin{VNBcond}
\item[Denotation] $\Gamma$ contains $\iota \in X$ for some $X$, or an
  equation $\iota = u$ or $u = \iota$.
\end{VNBcond}
\texttt{iota-def} pays with a proof of existence and uniqueness.
\texttt{iota-in-elim} pays with the context: the two forms of the
Denotation condition are the first two clauses of the definedness
certificate, and the operation declines without one of them.

\subsubsection*{\texttt{lambda-type} (command \texttt{lam-t})}

$\mathrm{vnb\text{-}lambda}(\sigma, A, b)$ is the set of pairs
$\{(x, b) : x \in A\}$, with $\sigma$ the binder, a variable or a list of
variables.  The term carries its domain.
%% rcs-lambda-type :437, rcs--lambda-typing-ok? :414, rule-checkers-schema.scm
\VNBfig{lambda-type}{\Gamma \Longrightarrow \forall v.\, v \in A \Rightarrow b[x{:=}v] \in B \qquad \Gamma \Longrightarrow A \in \mathrm{set}}{\Gamma \Longrightarrow \mathrm{vnb\text{-}lambda}(x, A, b) \in \mathrm{fun}(A, B)}
\begin{VNBcond}
\item[Domain] The domain the term declares is $\aeq$ the domain of the
  $\mathrm{fun}$.
\item[Freshness] $v = x$, or $v \notin \FV(b) \cup \FV(B)$.
\item[Binder list] For $\sigma = [x_1, \ldots, x_n]$ the domain is
  $A_1 \times \cdots \times A_n$ with the same $n$, and the first
  hypothesis is the $n$-fold guarded universal
  $\forall v_1.\, v_1 \in A_1 \Rightarrow \cdots \Rightarrow
   \forall v_n.\, v_n \in A_n \Rightarrow b[x_1{:=}v_1, \ldots, x_n{:=}v_n] \in B$,
  each $v_i$ fresh as above and the $v_i$ pairwise distinct.
\end{VNBcond}
The operation has no argument, and declines unless the Domain and Binder
list conditions hold.  Neither condition is decoration.  Without the
first, one term types into $\mathrm{fun}(A, B)$ for every $A$.  Without
the sethood hypothesis, a function on a proper class, which is itself a
proper class, would be certified into a $\mathrm{fun}$.

\subsubsection*{\texttt{lambda-beta}, \texttt{lambda-beta-hyp}
  (commands \texttt{lam-b}, \texttt{lam-b-h})}

A \emph{redex} is a term
\[
  \mathrm{vnb\text{-}lambda}(\sigma, A, b)(u_1, \ldots, u_n),
\]
with $n = 1$ for a single binder and $n$ the length of $\sigma$ for a
binder list.  It contracts to $b$ with the $u_i$ substituted in parallel
for the variables of $\sigma$.  Write $o_1, \ldots, o_m$ for the
\emph{obligations}, $m \geq 0$.
%% rcs-lambda-beta :854, rcs-lambda-beta-hyp :881, rcs--beta-match/fuel :701,
%% rcs--beta-licences-ok? :816, rcs--beta-licensed? :775, rule-checkers-schema.scm
\VNBfigS{lambda-beta}{\Gamma \Longrightarrow G' \\ \Gamma \Longrightarrow o_1 \quad \cdots \quad \Gamma \Longrightarrow o_m}{\Gamma \Longrightarrow G}
\VNBfigS{lambda-beta-hyp}{(\Gamma - F), F' \Longrightarrow G \\ \Gamma \Longrightarrow o_1 \quad \cdots \quad \Gamma \Longrightarrow o_m}{\Gamma \Longrightarrow G}
\begin{VNBcond}
\item[Contraction] $G'$ (respectively $F'$, for some $F \in \Gamma$) is
  reached from $G$ (respectively $F$) by a sequence of contractions,
  each taken at a position of the formula reached so far.  A redex that
  a contraction creates may be contracted in its turn.
\item[Licence] Each redex contracted at a position $\pi$ is licensed
  there or owed.  It is \emph{licensed} when $u \in A$ (for $n = 1$) or
  $[u_1, \ldots, u_n] \in A$, or, when $A$ is a product
  $A_1 \times \cdots \times A_n$, every $u_i \in A_i$, is among the
  formulas known at $\pi$: the assumptions of $\Gamma$ in which no
  variable bound between the root and $\pi$ occurs free, together with
  the guards of the guarded universals that enclose $\pi$ and the domain
  memberships of the $\mathrm{vnb\text{-}lambda}$, $\mathrm{sep}$ and
  $\mathrm{big\text{-}union}$ binders that enclose it.  It is
  \emph{owed} when that membership is one of the $o_j$ and contains no
  variable bound between the root and $\pi$.  The argument may be taken
  as written or with the redexes inside it contracted; in the second
  case those inner redexes count as contracted and need their own
  licence.
\item[Obligations] Every $o_j$ is owed by some redex.
\item[Bound] The checker's search takes at most
  \texttt{rcs--beta-fuel-limit} contractions (\texttt{rule-checkers-schema.scm}),
  and refuses when that runs out.
\end{VNBcond}
\texttt{lambda-beta} takes no argument; \texttt{lambda-beta-hyp} takes
the assumption $F$.  Each reduces every redex it can license or owe, in
one bottom-up pass, and declines when the reduction changes nothing.
The obligations are posted after the reduced formula, one per redex
whose licence is not evident.  A member of $\mathrm{fun}(A, B)$ is
defined exactly on $A$, so reducing off the domain without the
obligation would make a term denote that the library says does not.

\begin{remark}
The obligation is posted in the context of $N$.  A redex under a binder
that has not yet been introduced therefore owes a membership in which the
bound variable is free, and that leaf cannot be closed.  Introducing the
binder and typing the argument first is not a preference but the
condition under which the obligation is provable.
\end{remark}

\subsection{Rewriting}

\subsubsection*{\texttt{macete}: rewriting the assertion by a theorem
  (commands \texttt{mac}, \texttt{macm})}

A \emph{macete} is a rewrite rule read off an installed theorem of the
form
\[
  \forall x_1 \ldots x_n.\; C_1 \Rightarrow (\cdots (C_k \Rightarrow L = R)),
\]
where $=$ may also be quasi-equality or $\Leftrightarrow$, and the
conditions $C_i$ may be absent.  A name may also denote a rule that is
not a theorem: a functoid $F$ with parameters $\vec p$ and body $e$
denotes $F(\vec p) \to e$, a structure accessor with slot $k$ denotes
$\mathrm{ACC}(s) \to \mathrm{nth}(k, s)$, and a functor projection
\texttt{F@ACC} denotes the corresponding component of $F$'s body.
%% rkw--macete-definition/compute, rule-checkers-rewrite.scm:593
The \emph{local context} of a position in a formula is $\Gamma$, with
conjunctions split into their conjuncts, extended on the way down:
entering $B$ in $A \Rightarrow B$ adds $A$; entering either conjunct of a
conjunction adds the other; entering either disjunct of $A \vee B$ adds
the negation of the other; entering the scope of a binder
($\forall$, $\exists$, $\mathrm{iota}$, $\mathrm{comp}$, the body of
$\mathrm{sep}$, $\mathrm{big\text{-}union}$ or $\mathrm{vnb\text{-}lambda}$,
a functoid) removes every formula in which the bound variable occurs
free (Monk~\cite{Monk1988}).
%% rkw--descend, rule-checkers-rewrite.scm:471
%% rkw--check-macete, rule-checkers-rewrite.scm:687; rkw--condition-discharged? :390
\VNBfig{macete $N$ $L$ $R$}{\Gamma \Longrightarrow G' \qquad \Gamma \Longrightarrow c_1 \quad \cdots \quad \Gamma \Longrightarrow c_m}{\Gamma \Longrightarrow G}
\begin{VNBcond}
\item[Licence] $N$ names an installed theorem, a functoid, an accessor
  or a functor projection, and denotes the rule
  $(C_1, \ldots, C_k;\ L_0 \to R_0)$.
\item[Tag] The recorded $L$ and $R$ are $\sigma L_0$ and $\sigma R_0$
  for one substitution $\sigma$.  (Not tested for the two variadic rules,
  \texttt{union-decompose} and \texttt{intersection-decompose}.)
\item[Rewrite] $G'$ is $G$ with some subterms of the form $\tau L_0$
  replaced by $\tau R_0$, $\tau$ a substitution of the rule's variables
  found by matching.  A pattern $\mathrm{succ}(P)$ matches a positive
  integer literal $m$ by matching $P$ against $m-1$.
\item[Conditions] At each rewritten position every $\tau C_i$ is
  \textsc{truth}, or is in the local context of the position, or is a
  closed arithmetic sentence that evaluates to true, or is one of the
  $c_j$.
\end{VNBcond}
The operation replaces each subterm that matches $L_0$ and whose
conditions hold in the local context; a match whose conditions do not
all hold is left alone, and the rewriter descends into its subterms
instead.  It declines when $G'$ is $G$.  Its local context is the one
above, except that inside a conjunction $A \wedge B$ it extends only the
context of $B$, by $A$.  Hypotheses $c_j$ appear only when the command
posts the conditions it could not discharge as subgoals.  The name of
the theorem and the two patterns are recorded with the operation, so
that the inference can be checked against the rule it claims to apply
(Section~\ref{sec:inference-checked}).

\subsubsection*{\texttt{macete-hyp}: rewriting an assumption
  (commands \texttt{mac-h}, \texttt{mac-h*})}

%% rkw--check-macete-hyp, rule-checkers-rewrite.scm:741
\VNBfig{macete-hyp $N$ $L$ $R$}{(\Gamma - H), H' \Longrightarrow G \qquad \Gamma \Longrightarrow c_1 \quad \cdots \quad \Gamma \Longrightarrow c_m}{\Gamma \Longrightarrow G}
\begin{VNBcond}
\item[Assumption] $H \in \Gamma$.  Exactly one formula of $\Gamma$ is
  missing from the first hypothesis, and at most one formula is new in
  it (none when $H'$ was already in $\Gamma$).
\item[Licence, Tag, Rewrite, Conditions] As for \texttt{macete}, with
  $H$ and $H'$ in place of $G$ and $G'$, and the local context started
  from all of $\Gamma$, $H$ included.
\item[Not checked] The two conditions below on which the operation
  declines.  Neither is needed for soundness: every rule a name denotes
  replaces a term by an equal or quasi-equal one, or a formula by an
  equivalent one, under its conditions.
\end{VNBcond}
The arguments are the name and $H$.  The rewrite is made whether or not
the conditions hold, and one hypothesis $\Gamma \Longrightarrow c_j$ is
posted for each instantiated condition that the local context does not
contain.  Besides declining when the rewrite changes nothing, the
operation declines on two conditions read off the named theorem.  The
theorem's core must be an equivalence, that is an $\Leftrightarrow$ or
one of the two equalities, since replacing an assumption by a
consequence of it would lose information.  And no variable of $R$ or of
a condition may be left undetermined by the match, since such a variable
would be conjured free into the context.

\subsubsection*{\texttt{cartesian-decompose} (command \texttt{mac})}

%% rkw--check-cartesian-decompose :870, rkw--cartesian-shape? :819, rule-checkers-rewrite.scm
\VNBfig{cartesian-decompose}{\Gamma \Longrightarrow G'}{\Gamma \Longrightarrow G}
\begin{VNBcond}
\item[Rewrite] $G'$ is $G$ with some subformulas
  $x \in A_1 \times \cdots \times A_n$, $n \geq 1$, replaced by
  \[
    \exists v_1.\, v_1 \in A_1 \wedge
    \bigl(\cdots \wedge
    \exists v_n.\, v_n \in A_n \wedge x = [v_1, \ldots, v_n] \bigr).
  \]
\item[Freshness] The $v_i$ are pairwise distinct, and none is free in
  $x$ or in any $A_j$.
\end{VNBcond}
No argument beyond the name.  The operation rewrites every such
subformula, with new $v_i$, and declines when $G'$ is $G$.  This is the
membership condition of the product, and it is not a formula in the
catalog: the fresh witnesses cannot be written as a fixed template, so it
is code.  It is on the primitive shelf, and contributes nothing to a
bill.

\subsubsection*{\texttt{tuple-equality-decompose} (command \texttt{mac})}

%% rkw--check-tuple-equality-decompose :902, rkw--tuple-conjunction :883, rule-checkers-rewrite.scm
\VNBfig{tuple-equality-decompose}{\Gamma \Longrightarrow G'}{\Gamma \Longrightarrow G}
\begin{VNBcond}
\item[Rewrite] $G'$ is $G$ with some subformulas
  $[a_1, \ldots, a_n] = [b_1, \ldots, b_n]$, literal tuples of the same
  length, replaced by $a_1 = b_1 \wedge (a_2 = b_2 \wedge \cdots)$; by
  $a_1 = b_1$ when $n = 1$, and by \textsc{truth} when $n = 0$.
\end{VNBcond}
No argument beyond the name.  The operation rewrites every such
subformula and declines when $G'$ is $G$.  Tuples of different lengths
are not rewritten.  That two of them cannot be equal is a separate fact
and not this operation's business.

\subsection{Oracles}\label{sec:oracles}

An oracle decides a class of assertions by computation.  Six of the
seven, on success, record an inference with no hypotheses, so $N$ is
grounded at once; \texttt{arith-simplify} posts one hypothesis.
\texttt{ineq} and \texttt{sos} record, with the name, the certificate the
computation produced, and print it; the others record the name alone.
Each inference is checked before it is written, as described in
Section~\ref{sec:inference-checked}.  The figures below state what the
checker of each oracle tests.
%% ineq tag: pi-ineq!, structure-library/ineq-oracle.scm:308;
%% sos tag: structure-library/sos-oracle.scm:120; arith-eval.scm:382-388

\subsubsection*{\texttt{arith-ground}, \texttt{arith-forsome},
  \texttt{arith-simplify} (command \texttt{arith})}

One procedure records these three, trying them in this order.
%% rco--check-arith-ground :242, rco--check-arith-forsome :265,
%% rco--check-arith-simplify :298, rule-checkers-oracle.scm
\VNBfig{arith-ground}{\VNBnohyp}{\Gamma \Longrightarrow G}
\begin{VNBcond}
\item[Evaluation] $G$ is closed, and exact evaluation gives true.
\end{VNBcond}
\VNBfig{arith-forsome}{\VNBnohyp}{\Gamma \Longrightarrow \exists x.\, x \in C \wedge e = x}
\begin{VNBcond}
\item[Evaluation] $e$ evaluates exactly to a number $v$, and the body
  evaluates to true at $x = v$.  The equation may also be written
  $x = e$.
\end{VNBcond}
\VNBfig{arith-simplify}{\Gamma \Longrightarrow G'}{\Gamma \Longrightarrow G}
\begin{VNBcond}
\item[Evaluation] $G'$ is $G$ with some ground subterms replaced by
  the numbers to which they evaluate exactly.
\end{VNBcond}

\texttt{arith-ground} fires when $G$ evaluates to true.  The evaluator
reads $=$ and $\leq$ between terms it can reduce to exact numbers,
membership in $\mathbb{N}$, $\mathbb{Z}$, $\mathbb{Q}$, $\mathbb{R}$ and
$\mathbb{C}$, and the connectives $\neg$, $\wedge$, $\vee$,
$\Rightarrow$.  Everything else is undecided: a free variable, a
quantifier, a biconditional, a strict $<$, or a value that is not exact.
Arithmetic is exact throughout; a computed floating-point value is
refused rather than rounded.

\texttt{arith-forsome} fires on the shape of its figure when $e$ reduces
to a number and that number lies in $C$.

\texttt{arith-simplify} is the fallback.  When replacing the ground
numeric subterms of $G$ by their values changes $G$, it posts the single
hypothesis $\Gamma \Longrightarrow G'$.

The command declines when none of the three applies.

\subsubsection*{\texttt{ring-simplify}: identities of rings
  (commands \texttt{rs}, \texttt{simp})}

%% rco--check-rs :939, rco--crs-number-ok? :881, rco--num-gens-ok? :846, rule-checkers-oracle.scm
\VNBfig{ring-simplify}{\VNBnohyp}{\Gamma \Longrightarrow \forall x_1.\, x_1 \in D_1 \Rightarrow \cdots \forall x_k.\, x_k \in D_k \Rightarrow e_1 = e_2}
\begin{VNBcond}
\item[Normal form] $e_1$ and $e_2$ have the same normal form in the
  free associative ring $\mathbb{Z}[\textit{generators}]$.
\item[Typing] Each $D_i$ is one of $\mathbb{N}$, $\mathbb{Z}$,
  $\mathbb{Q}$, $\mathbb{R}$, $\mathbb{C}$, and every generator of
  $e_1$ or $e_2$ is one of the $x_i$ or is typed in one of these domains
  by an assumption $g \in D$ of $\Gamma$.
\end{VNBcond}
No arguments.  A \emph{generator} is a maximal subterm that ground
evaluation cannot reduce to a number, and multiplication does not
commute.  The operation declines unless both conditions hold.  The
generators counted are those of the two sides as written, including any
that cancel.  Otherwise $x - x = 0$ would close for an $x$ of which
nothing is known, and in VNB a membership is also what says that a term
denotes.

\subsubsection*{\texttt{comm-ring-simplify}: identities of commutative rings
  (command \texttt{crs})}

%% rco--check-crs :929, rco--crs-structure-ok? :894, rule-checkers-oracle.scm
\VNBfig{comm-ring-simplify}{\VNBnohyp}{\Gamma \Longrightarrow \forall x_1.\, x_1 \in D_1 \Rightarrow \cdots \forall x_k.\, x_k \in D_k \Rightarrow e_1 = e_2}
\begin{VNBcond}
\item[Normal form] $e_1$ and $e_2$ have the same normal form in the
  free commutative ring $\mathbb{Z}[\textit{generators}]$, literal
  powers being expanded first.
\item[Typing] Either as for \texttt{ring-simplify}; or the operations
  are those of a structure $R$, $(\mathrm{add}\ R)$, $(\mathrm{mul}\ R)$,
  $(\mathrm{neg}\ R)$, $(\mathrm{zero}\ R)$, $(\mathrm{one}\ R)$,
  $\mathrm{is\text{-}commutative\text{-}ring}(R)$ is a peeled guard or is
  in $\Gamma$, and every generator is a variable bound by a guard
  $x_i \in \mathrm{carr}(R)$ or is typed in $\mathrm{carr}(R)$ by an
  assumption.
\end{VNBcond}
No arguments.  The operation tries the number reading first and the
structure reading second.

An existential assertion is outside this operation and every other
oracle except \texttt{arith-forsome}'s one shape: none of them supplies a
witness.  For $\exists a \in \mathbb{Z}.\, (x+y)^3 = x^3 + y \cdot a$ the
witness is given with \texttt{forsome-intro} (command \texttt{ew}), and
the equation that remains is an identity within reach of
\texttt{comm-ring-simplify}.

\subsubsection*{\texttt{ineq}: an oracle for linear arithmetic over the reals
  (command \texttt{ineq})}

%% rco--check-ineq :556, rco--farkas-ok? :515, rco--ineq-constraint :494, rule-checkers-oracle.scm
\VNBfig{ineq $c$}{\VNBnohyp}{\Gamma \Longrightarrow \forall x_1.\, x_1 \in \mathbb{R} \Rightarrow \cdots \forall x_k.\, x_k \in \mathbb{R} \Rightarrow G_0}
\begin{VNBcond}
\item[Certificate] $c$ is a list of pairs $(\mathit{id}\ .\ m)$,
  each $m$ an exact nonnegative rational, where $\mathit{id}$ is
  \texttt{goal}, the negation of $G_0$; or \texttt{h}$I$, the $I$-th
  assumption of $N$ counted from one; or \texttt{h}$I$\texttt{-rev}, the
  $I$-th assumption, an equation, read from right to left.  When $G_0$
  is an equation, $c$ is a pair of two such lists, one for each
  direction.
\item[Premises] Each assumption the certificate names is an equation or
  a $\leq$ or $<$ inequality.
\item[Absurdity] The sum of the named constraints, each written
  $\ell \leq 0$ or $\ell < 0$ and multiplied by its $m$, has every atom's
  coefficient zero, and its constant is $> 0$, or $\geq 0$ when a strict
  constraint enters with positive multiplier.
\item[Typing] Every atom of a constraint that enters with a positive
  multiplier is certified real: $\Gamma$ contains $a \in \mathbb{R}$, or $a$ is one of the
  $x_i$, or $a$ is $|t|$ with every atom of $t$ certified.
\item[Eigenvariable] No $x_i$ is free in a premise the certificate
  names.
\end{VNBcond}
The arguments of the operation are the positions in $\Gamma$ of the
assumptions to be used as premises, counted from one.  When the
procedure succeeds $N$ is grounded, and otherwise the operation
declines.

The assertion and each named premise are read as linear equations or
inequalities in their \emph{atoms}, the maximal subterms that are not
sums, differences, or products by a numeral.  A named premise that does
not read this way is left out, which can only make the procedure prove
less; an index outside the range of $\Gamma$ is an error and abandons the
call.  The operation requires the Typing condition of every atom, an
atom that cancels included, and declines when a variable bound by one of
the leading quantifiers occurs free in \emph{any} assumption of $N$.
Fourier--Motzkin elimination is run on the premises and the negation of
the assertion.  The procedure succeeds when the elimination reaches
$0 < 0$.  The nonnegative combination of premises that produces it, a
\emph{Farkas certificate}, is printed and is recorded in the graph with
the name of the operation.

\begin{remark}
The condition on bound variables is the eigenvariable condition of
\texttt{forall-intro}.  Until 2026-09-20 the oracle did not impose it,
and from $x \in \mathbb{R},\; x \leq 0$ it closed
$\forall x \in \mathbb{R}.\; x \leq 0$.  The defect was found when the
checker of Section~\ref{sec:inference-checked} was written for this
operation, and no proof in the library had used it.
\end{remark}

\subsubsection*{\texttt{sos}: nonnegativity by a sum of squares
  (command \texttt{sos})}

%% rco--check-sos :956, rule-checkers-oracle.scm; tag built at structure-library/sos-oracle.scm:120
\VNBfig{sos $c$}{\VNBnohyp}{\Gamma \Longrightarrow \forall x_1.\, x_1 \in \mathbb{R} \Rightarrow \cdots \forall x_k.\, x_k \in \mathbb{R} \Rightarrow a \leq b}
\begin{VNBcond}
\item[Certificate] $c$ is a list of pairs $(c_i\ .\ \lambda_i)$, each
  $\lambda_i$ an exact nonnegative rational.
\item[Identity] $b - a = \lambda_1 c_1^2 + \cdots + \lambda_n c_n^2$ as
  polynomials in the free commutative ring over the generators.
\item[Typing] Every generator of $a$, $b$ and the $c_i$ is certified
  real, as for \texttt{ineq}.
\end{VNBcond}
The assertion may also be written $b \geq a$.  The arguments of the
operation are the terms $c_1, \ldots, c_n$ to be squared.  An exact
linear program decides whether weights $\lambda_i \geq 0$ exist; on
success the pairs $(c_i\ .\ \lambda_i)$ are printed and recorded in the
graph with the name.

The oracle verifies a certificate, it does not search for one: the
squares are supplied by the caller.  A strict inequality is not accepted,
since a square may be zero.

\subsection{Which command reaches which operation}\label{sec:command-uses}

A tactic has no access to the graph except through the operations above,
so the operations one tactic can reach are the bound on what it can do.
The table below gives that bound for each proof command: the operations
recorded by any kernel procedure reachable from the command's own code,
whether or not the library happens to exercise it.  It is produced from
\texttt{reference/kernel-map-static.sexp}, the result of reading the
source of every loaded file with the Scheme reader.  One thing is left
out of every row: after any command that changes a proof, the system tries
to close the definedness obligations that instantiation has left
(Section~\ref{sec:kernel-ops-detail}, \texttt{forall-elim}), and that step
can reach kernel procedures of its own.  It is the same for every
command and is stated once in \texttt{reference/KERNEL-MAP.md}.

A command listed as \texttt{none} records nothing.  Those are the
navigation, search and session commands.

%% The table and its counts, generated by gen-tactic-uses.py from
%% reference/kernel-map-static.sexp and the command registry.
%% GENERATED by gen-tactic-uses.py from reference/kernel-map-static.sexp and
%% tactics-help.scm (*tactic-kind*).  Do not edit.
\newcommand{\VNBtacticusescount}{108}
\newcommand{\VNBtacticusesmax}{64}
\newcommand{\VNBtacticusestable}{%
\begin{longtable}{@{}p{0.17\textwidth}p{0.79\textwidth}@{}}
\toprule
command & operations it can reach \\
\midrule
\endfirsthead
\toprule
command & operations it can reach \\
\midrule
\endhead
\texttt{ai} & \texttt{and-elim} \texttt{forsome-elim} \texttt{iff-elim} \texttt{not-elim} \texttt{or-elim}\\
\texttt{apply-thm} & \texttt{forall-elim} \texttt{theorem-assumption}\\
\texttt{arith} & \texttt{arith-forsome} \texttt{arith-ground} \texttt{arith-simplify} \texttt{cartesian-decompose} \texttt{eq-subst} \texttt{macete} \texttt{tuple-equality-decompose}\\
\texttt{arith--decide} & \texttt{arith-forsome} \texttt{arith-ground} \texttt{arith-simplify}\\
\texttt{ass} & \texttt{assumption}\\
\texttt{ass-all} & \texttt{assumption}\\
\texttt{backup-one} & none\\
\texttt{bc} & \texttt{backchain}\\
\texttt{bc*} & \texttt{assumption} \texttt{backchain} \texttt{cut} \texttt{forall-elim} \texttt{theorem-assumption}\\
\texttt{bc*--attempt} & \texttt{assumption} \texttt{backchain} \texttt{cut} \texttt{forall-elim} \texttt{theorem-assumption}\\
\texttt{beta} & \texttt{functoid-beta}\\
\texttt{bu-me} & \texttt{big-union-mem-elim}\\
\texttt{bu-mi} & \texttt{big-union-mem-intro}\\
\texttt{bu-set} & \texttt{big-union-sethood}\\
\texttt{calc} & \emph{all 64 reachable operations}\\
\texttt{ce} & \texttt{cartesian-elim}\\
\texttt{choose} & \texttt{and-elim} \texttt{and-intro} \texttt{arith-forsome} \texttt{arith-ground} \texttt{arith-simplify} \texttt{assumption} \texttt{cartesian-decompose} \texttt{cut} \texttt{detach} \texttt{eq-subst} \texttt{forall-elim} \texttt{forall-intro} \texttt{forsome-elim} \texttt{forsome-intro} \texttt{iff-elim} \texttt{iff-intro} \texttt{implies-intro} \texttt{macete} \texttt{not-elim} \texttt{not-intro} \texttt{or-elim} \texttt{theorem-assumption} \texttt{tuple-equality-decompose}\\
\texttt{choose-pos} & \texttt{and-elim} \texttt{and-intro} \texttt{arith-forsome} \texttt{arith-ground} \texttt{arith-simplify} \texttt{assumption} \texttt{cartesian-decompose} \texttt{cut} \texttt{detach} \texttt{eq-subst} \texttt{forall-elim} \texttt{forall-intro} \texttt{forsome-elim} \texttt{forsome-intro} \texttt{iff-elim} \texttt{iff-intro} \texttt{implies-intro} \texttt{macete} \texttt{not-elim} \texttt{not-intro} \texttt{or-elim} \texttt{theorem-assumption} \texttt{tuple-equality-decompose}\\
\texttt{ci} & \texttt{cartesian-intro}\\
\texttt{comp-me} & \texttt{comp-mem-elim}\\
\texttt{comp-mi} & \texttt{comp-mem-intro}\\
\texttt{contra} & \texttt{and-elim} \texttt{and-intro} \texttt{arith-forsome} \texttt{arith-ground} \texttt{arith-simplify} \texttt{assumption} \texttt{cartesian-decompose} \texttt{cut} \texttt{detach} \texttt{eq-subst} \texttt{forall-elim} \texttt{forall-intro} \texttt{forsome-elim} \texttt{iff-elim} \texttt{iff-intro} \texttt{implies-intro} \texttt{ineq} \texttt{macete} \texttt{not-elim} \texttt{not-intro} \texttt{or-elim} \texttt{theorem-assumption} \texttt{tuple-equality-decompose}\\
\texttt{crs} & \texttt{comm-ring-simplify}\\
\texttt{cut} & \texttt{cut}\\
\texttt{detach!} & \texttt{detach}\\
\texttt{di} & \texttt{and-intro} \texttt{forall-intro} \texttt{iff-intro} \texttt{implies-intro} \texttt{not-intro}\\
\texttt{dial} & none\\
\texttt{dial-wff} & none\\
\texttt{dk-focus!} & none\\
\texttt{eps-part} & \texttt{and-elim} \texttt{and-intro} \texttt{arith-forsome} \texttt{arith-ground} \texttt{arith-simplify} \texttt{assumption} \texttt{cartesian-decompose} \texttt{cut} \texttt{detach} \texttt{eq-subst} \texttt{forall-elim} \texttt{forall-intro} \texttt{forsome-elim} \texttt{iff-elim} \texttt{iff-intro} \texttt{implies-intro} \texttt{macete} \texttt{not-elim} \texttt{not-intro} \texttt{or-elim} \texttt{theorem-assumption} \texttt{tuple-equality-decompose}\\
\texttt{ew} & \texttt{forsome-intro}\\
\texttt{ew-poly} & \texttt{and-intro} \texttt{comm-ring-simplify} \texttt{forsome-intro} \textsuperscript{r}\\
\texttt{expand} & \texttt{comm-ring-simplify} \texttt{cut} \texttt{eq-subst} \textsuperscript{r}\\
\texttt{fact} & \texttt{arith-forsome} \texttt{arith-ground} \texttt{arith-simplify} \texttt{cut} \texttt{detach} \texttt{forall-elim} \texttt{theorem-assumption}\\
\texttt{focus} & none\\
\texttt{focus-id} & none\\
\texttt{goal-status} & none\\
\texttt{grind} & \texttt{and-elim} \texttt{and-intro} \texttt{forall-intro} \texttt{forsome-elim} \texttt{iff-elim} \texttt{iff-intro} \texttt{implies-intro} \texttt{macete-hyp} \texttt{not-elim} \texttt{not-intro} \texttt{or-elim}\\
\texttt{grind-and-mp} & \texttt{and-elim} \texttt{and-intro} \texttt{arith-forsome} \texttt{arith-ground} \texttt{arith-simplify} \texttt{assumption} \texttt{cut} \texttt{detach} \texttt{forall-elim} \texttt{forall-intro} \texttt{forsome-elim} \texttt{iff-elim} \texttt{iff-intro} \texttt{implies-intro} \texttt{macete-hyp} \texttt{not-elim} \texttt{not-intro} \texttt{or-elim}\\
\texttt{have!} & \texttt{and-intro} \texttt{arith-forsome} \texttt{arith-ground} \texttt{arith-simplify} \texttt{assumption} \texttt{cartesian-decompose} \texttt{cut} \texttt{detach} \texttt{eq-subst} \texttt{forall-elim} \texttt{forall-intro} \texttt{iff-intro} \texttt{implies-intro} \texttt{macete} \texttt{not-intro} \texttt{theorem-assumption} \texttt{tuple-equality-decompose}\\
\texttt{ie} & \texttt{intersection-elim}\\
\texttt{if-false} & \texttt{if-false}\\
\texttt{if-true} & \texttt{if-true}\\
\texttt{ii} & \texttt{intersection-intro}\\
\texttt{in-rr} & \texttt{and-intro} \texttt{arith-forsome} \texttt{arith-ground} \texttt{arith-simplify} \texttt{assumption} \texttt{cartesian-decompose} \texttt{cartesian-intro} \texttt{cut} \texttt{detach} \texttt{eq-subst} \texttt{forall-elim} \texttt{forall-intro} \texttt{iff-intro} \texttt{implies-intro} \texttt{macete} \texttt{not-intro} \texttt{theorem-assumption} \texttt{tuple-equality-decompose}\\
\texttt{in-rr--refocus!} & none\\
\texttt{ineq} & \texttt{ineq}\\
\texttt{inst} & \texttt{forall-elim}\\
\texttt{inst+} & \texttt{arith-forsome} \texttt{arith-ground} \texttt{arith-simplify} \texttt{cut} \texttt{detach} \texttt{forall-elim}\\
\texttt{iota-d} & \texttt{iota-def}\\
\texttt{iota-e} & \texttt{iota-in-elim}\\
\texttt{lam-b} & \texttt{lambda-beta}\\
\texttt{lam-b-h} & \texttt{lambda-beta-hyp}\\
\texttt{lam-t} & \texttt{lambda-type}\\
\texttt{len-r} & \texttt{length-reduce}\\
\texttt{mac} & \texttt{cartesian-decompose} \texttt{macete} \texttt{tuple-equality-decompose}\\
\texttt{mac-h} & \texttt{macete-hyp}\\
\texttt{mac-h*} & \texttt{and-elim} \texttt{forsome-elim} \texttt{iff-elim} \texttt{macete-hyp} \texttt{not-elim} \texttt{or-elim}\\
\texttt{macm} & \texttt{cartesian-decompose} \texttt{macete} \texttt{tuple-equality-decompose}\\
\texttt{minimize!} & \texttt{and-elim} \texttt{and-intro} \texttt{arith-forsome} \texttt{arith-ground} \texttt{arith-simplify} \texttt{assumption} \texttt{cartesian-decompose} \texttt{cut} \texttt{detach} \texttt{eq-subst} \texttt{forall-elim} \texttt{forall-intro} \texttt{forsome-elim} \texttt{forsome-intro} \texttt{iff-elim} \texttt{iff-intro} \texttt{implies-intro} \texttt{macete} \texttt{not-elim} \texttt{not-intro} \texttt{or-elim} \texttt{reflexivity} \texttt{sep-mem-elim} \texttt{sep-mem-intro} \texttt{theorem-assumption} \texttt{tuple-equality-decompose}\\
\texttt{mp} & \texttt{arith-forsome} \texttt{arith-ground} \texttt{arith-simplify} \texttt{assumption} \texttt{cut} \texttt{detach} \texttt{forall-elim}\\
\texttt{ni} & \texttt{nn-induction}\\
\texttt{nth-r} & \texttt{nth-reduce}\\
\texttt{obtain} & \emph{all 64 reachable operations}\\
\texttt{obtain-at} & \texttt{and-elim} \texttt{and-intro} \texttt{arith-forsome} \texttt{arith-ground} \texttt{arith-simplify} \texttt{assumption} \texttt{cartesian-decompose} \texttt{cut} \texttt{detach} \texttt{eq-subst} \texttt{forall-elim} \texttt{forall-intro} \texttt{forsome-elim} \texttt{iff-elim} \texttt{iff-intro} \texttt{implies-intro} \texttt{macete} \texttt{not-elim} \texttt{not-intro} \texttt{or-elim} \texttt{theorem-assumption} \texttt{tuple-equality-decompose}\\
\texttt{oi-l} & \texttt{or-intro-left}\\
\texttt{oi-r} & \texttt{or-intro-right}\\
\texttt{orelse} & none\\
\texttt{pbc} & \texttt{proof-by-contradiction}\\
\texttt{prep} & none\\
\texttt{prop} & \texttt{and-elim} \texttt{and-intro} \texttt{arith-forsome} \texttt{arith-ground} \texttt{arith-simplify} \texttt{assumption} \texttt{cartesian-decompose} \texttt{cut} \texttt{detach} \texttt{eq-subst} \texttt{forall-elim} \texttt{forall-intro} \texttt{forsome-elim} \texttt{iff-elim} \texttt{iff-intro} \texttt{implies-intro} \texttt{macete} \texttt{not-elim} \texttt{not-intro} \texttt{or-elim} \texttt{or-intro-left} \texttt{or-intro-right} \texttt{proof-by-contradiction} \texttt{theorem-assumption} \texttt{tuple-equality-decompose}\\
\texttt{push-not-h} & \texttt{and-elim} \texttt{and-intro} \texttt{arith-forsome} \texttt{arith-ground} \texttt{arith-simplify} \texttt{assumption} \texttt{cartesian-decompose} \texttt{cut} \texttt{detach} \texttt{eq-subst} \texttt{forall-elim} \texttt{forall-intro} \texttt{forsome-elim} \texttt{forsome-intro} \texttt{iff-elim} \texttt{iff-intro} \texttt{implies-intro} \texttt{macete} \texttt{not-elim} \texttt{not-intro} \texttt{or-elim} \texttt{or-intro-left} \texttt{or-intro-right} \texttt{proof-by-contradiction} \texttt{theorem-assumption} \texttt{tuple-equality-decompose}\\
\texttt{qed} & none\\
\texttt{qrfl} & \texttt{quasi-reflexivity}\\
\texttt{quietly} & none\\
\texttt{repeat} & none\\
\texttt{replay-proof} & \emph{all 64 reachable operations}\\
\texttt{rfl} & \texttt{reflexivity}\\
\texttt{rs} & \texttt{ring-simplify}\\
\texttt{save-proof} & none\\
\texttt{scout} & \emph{all 64 reachable operations}\\
\texttt{scout-run} & none\\
\texttt{scout-show} & \emph{all 64 reachable operations}\\
\texttt{sep-me} & \texttt{sep-mem-elim}\\
\texttt{sep-mi} & \texttt{sep-mem-intro}\\
\texttt{sep-set} & \texttt{sep-sethood}\\
\texttt{show} & none\\
\texttt{simp} & \texttt{comm-ring-simplify} \texttt{cut} \texttt{eq-subst}\\
\texttt{slot} & \texttt{cartesian-decompose} \texttt{macete} \texttt{tuple-equality-decompose}\\
\texttt{slot-h} & \texttt{macete-hyp}\\
\texttt{sos} & \texttt{sos}\\
\texttt{sp} & none\\
\texttt{subst} & \texttt{eq-subst}\\
\texttt{supply} & \emph{all 64 reachable operations}\\
\texttt{ta} & \texttt{theorem-assumption}\\
\texttt{te} & \texttt{tuples-elim}\\
\texttt{tfi} & \texttt{transfinite-induction}\\
\texttt{tfi3} & \texttt{transfinite-induction-3cases}\\
\texttt{ti} & \texttt{tuples-intro}\\
\texttt{type-term} & \texttt{assumption} \texttt{comm-ring-simplify} \texttt{cut} \texttt{detach} \texttt{eq-subst} \texttt{forall-elim} \texttt{theorem-assumption} \textsuperscript{r}\\
\texttt{ue} & \texttt{union-elim}\\
\texttt{ui} & \texttt{union-intro}\\
\texttt{undo} & none\\
\texttt{use-at} & \texttt{and-intro} \texttt{arith-forsome} \texttt{arith-ground} \texttt{arith-simplify} \texttt{assumption} \texttt{cartesian-decompose} \texttt{cut} \texttt{detach} \texttt{eq-subst} \texttt{forall-elim} \texttt{forall-intro} \texttt{iff-intro} \texttt{implies-intro} \texttt{macete} \texttt{not-intro} \texttt{theorem-assumption} \texttt{tuple-equality-decompose}\\
\texttt{use-em} & \texttt{and-elim} \texttt{and-intro} \texttt{arith-forsome} \texttt{arith-ground} \texttt{arith-simplify} \texttt{assumption} \texttt{cartesian-decompose} \texttt{cut} \texttt{detach} \texttt{eq-subst} \texttt{forall-elim} \texttt{forall-intro} \texttt{forsome-elim} \texttt{iff-elim} \texttt{iff-intro} \texttt{implies-intro} \texttt{macete} \texttt{not-elim} \texttt{not-intro} \texttt{or-elim} \texttt{or-intro-left} \texttt{or-intro-right} \texttt{proof-by-contradiction} \texttt{theorem-assumption} \texttt{tuple-equality-decompose}\\
\texttt{vlet} & \texttt{and-elim} \texttt{cut} \texttt{forsome-elim} \texttt{iff-elim} \texttt{not-elim} \texttt{or-elim}\\
\texttt{wbc} & \texttt{and-elim} \texttt{arith-forsome} \texttt{arith-ground} \texttt{arith-simplify} \texttt{cut} \texttt{detach} \texttt{forall-elim} \texttt{forsome-elim} \texttt{iff-elim} \texttt{not-elim} \texttt{or-elim} \texttt{theorem-assumption}\\
\texttt{wk} & \texttt{weakening}\\
\bottomrule
\end{longtable}
\noindent\footnotesize \textsuperscript{r}: read from the command registry rather
than from the kernel map, which predates the command.\normalsize}

\VNBtacticusestable

\subsection{How an inference is recorded, and checked}\label{sec:inference-checked}

Every kernel operation ends by calling one procedure,
\texttt{dg-apply-rule!}\ (\texttt{deduction-graphs.scm}).  It is handed
the rule, the hypothesis sequents and the node $N$.  It checks the
inference (below); then it posts each hypothesis sequent as a node,
reusing a node the graph already has for the same sequent; then it
creates the inference node and writes the arrows.  No other procedure
writes to the graph.
%% dg-apply-rule!, deduction-graphs.scm:424; make-inference-node :59
The inference node records exactly three things: the rule, the list of
hypothesis nodes, and the conclusion node $N$.  The rule is the name of
the operation, or a list headed by that name when the operation carries
data about the application: an index (\texttt{(cartesian-elim 2)}), the
name and the patterns of a macete (\texttt{(macete \textit{name} $L$
$R$)}), or the certificate of an oracle (\texttt{(ineq $c$)},
\texttt{(sos $c$)}).

\subsubsection*{The construction and the check}

An operation has two parts.  The \emph{builder} is a procedure
\texttt{pi-\ldots!}\ (in \texttt{primitive-inferences.scm},
\texttt{macetes.scm} and the oracle files).  It reads $N$ and the
operation's argument and \emph{constructs} the hypothesis sequents: for
\texttt{and-intro}, \texttt{pi-direct-inference!} takes the two conjuncts
of $G$ and makes $\Gamma \Longrightarrow A$ and
$\Gamma \Longrightarrow B$; for \texttt{forall-intro} it also chooses the
eigenvariables; for an oracle it runs the decision procedure.  The
\emph{checker} is a separate procedure, registered for the operation's
name in one of four files (\texttt{rule-checkers-logic.scm},
\texttt{-schema.scm}, \texttt{-rewrite.scm}, \texttt{-oracle.scm}).
\texttt{dg-apply-rule!}\ hands it the finished result: the rule, the
hypothesis sequents and the sequent of $N$.  The checker does not call
the builder and does not compute hypotheses from $N$.  It decides whether
the sequents it is given stand in the relation of the operation's figure
(Section~\ref{sec:kernel-ops-detail}), and answers yes, or no with a
reason.

The checkers of the logic and schema operations use four things from the
expression layer: substitution, the free-variable computation, equality
of formulas up to renaming of bound variables ($\aeq$), and one-way
matching.  The rewrite checkers add two: a theorem is read as a rewrite
rule by the same procedure the rewriter uses
(\texttt{extract-rewrite-patterns}), because that reading is what the
name of a macete denotes; and a closed arithmetic condition is evaluated
by the arithmetic evaluator.  The oracle checkers carry their own exact
arithmetic, linear forms and polynomials, written separately from the
oracles.
%% headers of rule-checkers-logic.scm (1-13), rule-checkers-rewrite.scm (18-26),
%% rule-checkers-oracle.scm (1-14)

Two of the relations, written out as the checker tests them:
%% and-intro checker, rule-checkers-logic.scm:347; forall-elim checker :702
\begin{itemize}
\item \texttt{and-intro} accepts hypotheses $H_1, H_2$ and conclusion
  $\Gamma \Longrightarrow G$ exactly when $G$ is $A \wedge B$, there are
  two hypotheses, the context of each equals $\Gamma$ as a set, the
  assertion of $H_1$ is $\aeq A$ and that of $H_2$ is $\aeq B$.
\item \texttt{forall-elim} accepts exactly when there are one or two
  hypotheses, the assertion of $H_1$ is $\aeq G$, and either the context
  of $H_1$ is included in $\Gamma$, or $\Gamma$ contains a formula
  $\forall x.\, B$ and the context of $H_1$ is $\Gamma \cup \{D\}$ for a
  formula $D$ that matching shows to be $B[x{:=}t]$ for some term $t$;
  and, if there is a second hypothesis, it is
  $\Gamma \Longrightarrow u = u$ for some term $u$.
\end{itemize}

The eigenvariable condition of \texttt{forall-intro} is tested as
follows.  The checker counts the length of
the chain of leading universal quantifiers of $G$, passing through each
guard $x \in A$ on the variable just bound, and tries that length $k$
first and then each smaller one.  For a given $k$ it reads the
eigenvariables $y_1, \ldots, y_k$ off the hypothesis
$\Gamma' \Longrightarrow B'$ by matching: $G$ with its $k$ bound
variables replaced by holes is matched, guard by guard, against the
formulas of $\Gamma'$ that are not in $\Gamma$, and its body against
$B'$.  A hole that matching does not reach belongs to a variable absent
from the body, and is given the bound variable's own name.  It then
walks the chain forwards with a set $V$ of excluded variables, which
starts as $\FV(\Gamma)$.  At step $i$ the formula is
$\forall x_i.\, C_i$, and $y_i$ must be a variable, not in $V$, and not
free in $\forall x_i.\, C_i$.  If $C_i[x_i{:=}y_i]$ is
$(y_i \in A_i) \Rightarrow D$, the guard $y_i \in A_i$ is collected, its
free variables ($y_i$ and those of $A_i$) are added to $V$, and the walk
continues with $D$; otherwise $y_i$ is added to $V$ and the walk
continues with $C_i[x_i{:=}y_i]$.  The inference is accepted when the
formula reached after $k$ steps is $\aeq B'$ and $\Gamma'$ is $\Gamma$
together with the collected guards.
%% chk-forall-intro-at, rule-checkers-logic.scm:514; chk-forall-eigen-ok? :491;
%% chk-forall-fill :477; chk-vacuous-name :540; the k-loop in the checker at :555

If no checker is registered for the name, or the checker refuses,
\texttt{dg-apply-rule!}\ raises an error and the graph is unchanged.  A
refusal is printed even when output is suppressed, and the load ends by
reporting the number of inferences verified and the number refused.
%% dg-check-inference!, dg-check-report-refusal!, deduction-graphs.scm

An error in a builder goes unnoticed only if the checker's relation also
admits the inference the builder produced.  Where a figure of
Section~\ref{sec:kernel-ops-detail} carries a condition labelled
\emph{Not checked}, the relation is weaker than the operation, and at
that point the operation is the only safeguard.  There are four such
places: definedness in \texttt{forall-elim} and \texttt{reflexivity}, by
policy, since the certificate decides when an obligation may be omitted
and the checker accepts the inference either way; capture in
\texttt{eq-subst}; the variable of the second hypothesis of
\texttt{big-union-sethood}; and the two refusals of
\texttt{macete-hyp}, which are not needed for soundness.
%% policy: header of rule-checkers-logic.scm, lines 29-43

\subsubsection*{Oracles: a certificate check or a second decision}

An oracle has no proof to replay, so what its checker can establish
depends on what the oracle records.  There are two cases.
%% header of rule-checkers-oracle.scm, lines 22-62

\emph{Certificate check} (\texttt{ineq}, \texttt{sos}).  The oracle
records, with its name, the object that shows the assertion true, and
the checker verifies that object by arithmetic of its own, without
elimination, search or any code of the oracle.  For \texttt{ineq} the
object is a Farkas certificate: the checker reads the assertion and the
named premises as linear forms, multiplies each by its multiplier, adds,
and confirms that every atom cancels and that the constant left is an
absurd inequality.  For \texttt{sos} it is the list of squares and
weights: the checker expands $\lambda_1 c_1^2 + \cdots + \lambda_n c_n^2$
and $b - a$ as polynomials and compares them monomial by monomial.  Either
the arithmetic confirms the certificate, or the inference is refused.
%% rco--farkas-ok? :515, rco--check-sos :956, rule-checkers-oracle.scm

\begin{sloppypar}
\emph{Second decision} (\texttt{arith-ground}, \texttt{arith-forsome},
\texttt{arith-simplify}, \texttt{comm-ring-simplify},
\texttt{ring-simplify}).  These oracles record their name only, so the
checker decides the same question again with a second implementation of
the same mathematics: exact evaluation for \texttt{arith}, expansion of
both sides to a sum of monomials for the two ring oracles.  The two
implementations share no code beyond the expression layer.  A second
decision catches an error in how the oracle reads its input or reports
its answer: an atom misread, a typing it failed to require, a premise
dropped, a case it mishandles that the second implementation handles
correctly.  It does not catch an error that both implementations make
the same way, for instance a mistake in the mathematics of the normal
form itself.
\end{sloppypar}
%% rco--check-arith-ground :242, rco--check-crs :929, rco--check-rs :939

Neither kind of check adds proving power.  A checker only judges an
inference that an oracle has already produced; an assertion the oracle
declines is not proved by its checker.

\subsubsection*{Where a certificate is kept, and how it is shown}

The certificate of \texttt{ineq} is part of the rule recorded in the
inference node, beside the name of the operation.  The rule has the form
\texttt{(ineq $c$)}, where $c$ is a list of pairs
\texttt{($\mathit{id}$ .\ $m$)}: $m$ is an exact nonnegative rational,
and $\mathit{id}$ is \texttt{goal} (the negation of the assertion),
\texttt{h}$I$ (the $I$-th assumption of $N$, counted from one), or
\texttt{h}$I$\texttt{-rev} (the $I$-th assumption, an equation, read from
right to left).  For an assertion that is an equation, $c$ is a pair of
two such lists, one for each direction.  For example
\texttt{(ineq ((h1 .\ 1) (h3 .\ 2) (goal .\ 1)))} says that the first
assumption, twice the third, and the negated assertion add up to a
false inequality between constants.  The rule of \texttt{sos} is
\texttt{(sos (($c_1$ .\ $\lambda_1$) \ldots))}.
%% pi-ineq!, structure-library/ineq-oracle.scm:218 (tag built at :308);
%% ineq-hyp-constraints :167; sos-oracle.scm:120

The certificate therefore belongs to the deduction graph of the proof,
and lasts as long as that graph.  What \texttt{qed} keeps of a proof is
the statement, the script and the bill; the script records the command
as it was typed, for instance \texttt{(ineq 1 3)}, not the certificate.
Replaying the script, as the page audit does, runs the oracle again, and
the checker verifies the certificate recorded by that run.
%% proof-debt.scm:205 (*proof-oracles* holds oracle names, not certificates)

When it succeeds the oracle prints the certificate:
\begin{VNBcode}
;; ineq: closed by Farkas certificate ((h1 . 1) (h3 . 2) (goal . 1))
\end{VNBcode}
\texttt{sos} prints its decomposition in the same way.  \texttt{proof-tex}
prints, for each step, the command and the sequent it leaves, read from
the command trace recorded while the proof ran; since 2026-09-25 that trace
also keeps, with each step that closed a node by \texttt{ineq} or
\texttt{sos}, the certificate the oracle recorded together with the
formulas of the node it closed, and \texttt{proof-tex} prints it beside the
step.  A step of a composite tactic that runs an oracle with recording
suppressed (\texttt{zero-it} on a numeral difference) is recorded as the
composite alone: its certificate is on the node and in the printout, not on
the page.  The command \texttt{(proof-rules)} reads the graph of the
current proof and lists the rules its inferences recorded.
%% vnb--capture-step!, interactive.scm:189 (the sixth field, `certs');
%% proof-tex.scm:184; zi--run-step!, zero-it.scm (recording suppressed)
%% ineq-report, ineq-oracle.scm:215; sos-report, sos-oracle.scm:49;
%% proof-tex--steps-from-live, proof-tex.scm:94; proof-rules, suggest.scm:6006

\subsubsection*{Controls}

Each checker has two controls in the test suite: the proofs of the
library, which it must accept, and an inference that is wrong in a
stated way, which it must refuse.  A checker that accepted everything
would pass the first and fail the second.

\begin{remark}
Checking is on by default.  Setting the environment variable
\texttt{VNB\_NO\_RULE\_CHECK} turns it off.  Every inference of a load is
verified, and the load summary reports how many were verified and how many
refused; checking makes the load slower.
\end{remark}

\subsection{The boundary is checked, not merely documented}\label{sec:kernel-checked}

The restriction would be worth little if it were a convention.  The
files permitted to record inferences are named in a list in the source
(\texttt{*kernel-caller-files*}, \texttt{audit.scm}).  Every time the system starts, it reads
all of its own source with the Scheme reader and counts, in each file,
the calls to \texttt{dg-apply-rule!}\ and any use of that procedure as a
value --- the second being the way a file could otherwise hand the
capability to code outside the list.  If any file not on the list records
inferences, or if any file passes the procedure on, the system reports
the offending files and \emph{refuses to finish starting}.  A reader is
used rather than a text search so that a mention inside a comment, a
string or quoted data is not mistaken for code.

On a clean tree the check prints one line:

\begin{VNBcode}
;; kernel-callers-audit: ok (N call site(s) to dg-apply-rule!
;;   in M file(s), all allowed; 0 value use(s))
\end{VNBcode}

A companion check, \texttt{kernel-rules-audit}, compares the rules
actually recorded during the load against the documented list in both
directions, so a rule that is trusted but undocumented, or documented but
unreachable, is reported rather than left to be noticed.

Widening the trusted base is possible --- a new decision procedure is a
legitimate thing to add --- but it cannot happen quietly: the new file
must be added to the list by hand, which is a deliberate act that leaves
a record.  The complete inventory (every entry point, the rules it may
record, and which proof commands reach it) is
\texttt{reference/KERNEL-MAP.md}.

%% =========================================================
\section{Self-verification of the kernel (META)}\label{sec:meta}
\index{META}\index{self-verification}

This section collects what the tree does to check its own kernel.  The
operations and their checkers are described in
Sections~\ref{sec:kernel-boundary} and~\ref{sec:inference-checked}, the
closed list of files that may record an inference in
Section~\ref{sec:kernel-checked}, the proof certificate in
Section~\ref{sec:certificates}; they are not repeated here.  Consistency
of the whole is relative to NBG.

%% TODO(meta-ledger): report the K2 coverage ledger and the open findings here once the ledger is regenerated.
\subsection{Checking at the point of writing}

Every inference passes \texttt{dg-apply-rule!}, which asks the checker
registered for the operation's name to accept the pair (hypotheses,
conclusion) before anything is written to the graph
(Section~\ref{sec:inference-checked}).  The checkers are in
\texttt{rule-checkers-logic.scm}, \texttt{-schema.scm}, \texttt{-rewrite.scm}
and \texttt{-oracle.scm}.  A missing checker, or a refusal, raises an error;
the refusal is printed as a line beginning \texttt{;VNB RULE CHECKER REFUSED},
and the load summary reports the number of inferences verified and the number
refused.  In a strict load a nonzero number of refusals is an error.  The
builders (\texttt{pi-\ldots!}) and the tactics are therefore not part of what
must be trusted for a graph to be sound: a wrong builder can only fail to
prove.  What decides whether a false sequent can be grounded is the relation
each checker enforces, together with the expression layer the checkers use
(substitution, free variables, equality up to renaming, matching), the write
point and the propagation of grounding in \texttt{deduction-graphs.scm}, the
installation by \texttt{qed}, and the axioms.

\subsection{Predicativity: staying within NBG}

NBG allows class comprehension only over formulas whose quantifiers range over
sets.  The tree enforces this at four places.
\begin{itemize}
\item The kernel's membership rules for \texttt{comp} refuse a \texttt{comp}
  whose condition quantifies beyond sets, and the definedness certificate
  (\texttt{pi--body-shape-total?}, \texttt{primitive-inferences.scm}) treats
  a \texttt{comp} body as total only when its condition is predicative, and
  likewise a \texttt{sep}, \texttt{big-union} or \texttt{vnb-lambda} body.
\item \texttt{comp-predicativity-audit} lists every \texttt{comp} term named by
  a proof's commands, by an installed formula or by a functoid body whose
  condition quantifies beyond sets; \texttt{binder-predicativity-audit} does
  the same for \texttt{sep}, \texttt{big-union} and \texttt{vnb-lambda}
  (\texttt{proof-debt.scm}).  Both gates are fatal in a strict load and are
  reported under keep-going.  A certified load sees only the proofs it ran.
\item The \emph{NBG induction census}: an induction whose body has a binder not
  confined to sets.  The switch \texttt{*nbg-induction-fatal?*}
  (\texttt{rule-checkers-logic.scm}) is on, so the checkers of the induction
  operations refuse such a body, and the load prints one line,
  \texttt{;; NBG induction census: $n$ impredicative induction(s)}, with $n=0$
  in a load that completes.
\end{itemize}

\subsection{The core, written in VNB}

The library \texttt{meta} (\texttt{libraries.sexp}) encodes the kernel in the
kernel's own language.  \texttt{structure-library/meta-syntax.scm} encodes
expressions as sets (every node is a list whose first entry is a numeral tag);
\texttt{meta-core.scm} defines well-formedness, sequents, a set of core rule
relations and \texttt{DERIVABLE}$(T)$, the least set of sequents closed under
them, and transcribes the kernel's own rule relations;
\texttt{meta-cr.scm}, \texttt{meta-ind.scm} and \texttt{meta-def.scm} treat
contextual rewriting, the induction rules and the definedness clauses;
\texttt{meta-nbg.scm} transcribes the kernel forms of further rules, with the bridges that show each admissible over the core.  The laws proved of these
definitions are in the files of the same names in \texttt{theorem-library/}.
The aim, recorded in \texttt{docs/meta-k2-admissibility-2026-10-08.md}, is
that each non-oracle kernel rule is \emph{admissible} in the core: whenever
the hypotheses of an instance are derivable, so is its conclusion.  This is a
theorem about derivations regarded as data.  It does not give a semantics, and
it reduces what has to be trusted without removing it.

\subsection{Regression: the key and the exam}

The certificate of a theorem (Section~\ref{sec:certificates}) is valid only
for the kernel that issued it.  The kernel hash covers the files listed in
\texttt{*cert-kernel-files*} (\texttt{certificates.scm}): the files that may
record an inference, \texttt{deduction-graphs}, \texttt{expressions},
\texttt{sequents} and the four checker files; a change to any of them
invalidates every certificate.  The script \texttt{vnb-nightly-exam} runs the
exam each night from cron: every proof is run, the page of every proof is typed back, and the
certificates are rewritten (Section~\ref{sec:certificates}).

\subsection{What is not covered}

\begin{itemize}
\item The oracles are trusted code.  Their checkers verify a certificate
  (\texttt{ineq}, \texttt{sos}) or repeat the decision with a second
  implementation (Section~\ref{sec:inference-checked}); that adds no proving
  power and is not a proof of the oracle.
\item Tactics are untrusted; they act only through kernel operations.
\item That the Scheme code of a checker computes the relation written for it
  in the core is not a theorem: the transcription is made by hand, and the
  design note states that this faithfulness is not provable in VNB without a
  model of Scheme.
\item The expression layer, MIT Scheme, the certificate store and its hash,
  the undo journal and \texttt{qed} are assumed.
\item The axioms are assumed, and consistency is relative to NBG.  A
  statement or axiom stamped as definitional or primitive is invisible in the
  bill; the shelf of primitives is a list kept in the tree.
\end{itemize}

%% =========================================================
\section{Certificates}\label{sec:certificates}
\index{certificate}

VNB uses the word \emph{certificate} for several distinct objects.  What
they share is the division of labour: a certificate is evidence produced
by one part of the system and verified by another, and what a proof rests
on is the verification, never the producer.  This section lists every
kind, each with what it certifies, who issues it, its form, where it is
kept, who verifies it, and what it does not do.  The first five are
consumed by the inference check of Section~\ref{sec:inference-checked};
the sixth is the mechanism by which a load installs a theorem without
running its proof.

\subsection{The Farkas certificate}
\index{certificate!Farkas}\index{ineq@\texttt{ineq}!certificate}

\begin{description}
\item[Certifies] that an assertion $s \le t$, $s < t$ or $s = t$ over the
  reals, possibly under universals typed in $\mathbb{R}$, follows from the
  linear order facts among the assumptions of the node.
\item[Issued by] the oracle \texttt{ineq}.  Fourier--Motzkin elimination
  finds a nonnegative combination of the premises and the negated
  assertion whose atoms cancel and whose constant is absurd; that
  combination is the certificate.
\item[Form] a list of pairs $(\mathit{id}\ .\ m)$ with $m$ an exact
  nonnegative rational and $\mathit{id}$ one of \texttt{goal},
  \texttt{h}$I$, \texttt{h}$I$\texttt{-rev}; a pair of such lists for an
  equation (Section~\ref{sec:oracles}).
\item[Kept] as the rule tag \texttt{(ineq $c$)} on the inference node of
  the deduction graph, for the life of the graph; in the sixth field of the
  step's trace record, from which \texttt{proof-tex} prints it beside the
  step; and in the printout.  Not in the script, which records the command
  as typed, and not in the proof certificate below, whose bill names the
  oracle and not its combination.  Replay recomputes it.
\item[Verified by] the checker of the \texttt{ineq} operation: every
  multiplier nonnegative and exact, every $\mathit{id}$ naming an order
  formula read in the claimed direction, the combination cancelling to an
  absurd constant, every atom certified real (below), and no peeled
  variable free in a premise the certificate uses.
\item[Does not] search beyond linear arithmetic: a product of two
  non-constant terms is not an atom of the combination, and a premise
  containing one is dropped before the search.
\end{description}

\subsection{The sum-of-squares certificate}
\index{certificate!sum of squares}\index{sos@\texttt{sos}!certificate}

\begin{description}
\item[Certifies] $a \le b$ over the reals by exhibiting
  $b - a = \sum_i \lambda_i c_i^2$ with each $\lambda_i \ge 0$.
\item[Issued by] the user.  The squares $c_i$ and the weights are the
  arguments of the command; \texttt{sos} verifies a certificate, it does
  not search for one.  A certificate that does not account for $b - a$ is
  refused with the reason.
\item[Form] the list $((c_1\ .\ \lambda_1) \ldots)$.
\item[Kept and verified] as the Farkas certificate: the rule tag
  \texttt{(sos $c$)}, the trace's sixth field, the printout; the checker
  recomputes the identity $b - a = \sum \lambda_i c_i^2$ by the ring
  normal form and re-tests the typing of every generator.
\end{description}

\subsection{The atom certificate}
\index{certificate!atom (real)}

\begin{description}
\item[Certifies] that a term the arithmetic oracles treat as an atom
  denotes a real number, which is what licenses the order laws the
  certificate combines.
\item[Issued by] the context.  An atom $a$ is certified real when
  $\Gamma$ contains the standalone assumption $a \in \mathbb{R}$, when $a$
  is a variable the oracle peeled from a universal typed in $\mathbb{R}$,
  or when $a$ is $|t|$ with every atom of $t$ certified.  A typing in
  $\mathbb{N}$ or $\mathbb{Z}$ does not certify: the composites
  \texttt{contra}, \texttt{supply}, \texttt{prep} and the kit's
  \texttt{dk-ineq!} land the typing lemmas (the \emph{typing
  certificates}) before the oracle runs, by the inclusion theorems.
\item[Form] a decision, not an object: the test
  \texttt{ineq-atom-rr-ok?} on $\Gamma$ and $a$.
\item[Verified by] the checker, which re-decides it for every atom the
  Farkas certificate uses with a positive multiplier.
\item[Does not] complain: a premise with an uncertified atom is dropped
  silently, and the oracle then blames the goal.  Name the premises by
  formula after \texttt{keep} to see which one was dropped.
\end{description}

\subsection{The generator certification of the ring oracles}
\index{certificate!generator (ring)}

\begin{description}
\item[Certifies] that every generator of the terms \texttt{crs} or
  \texttt{rs} normalises, a variable or an opaque compound term, is an
  element of the ring's carrier.
\item[Issued by] the context: an assumption $v \in D$ for a number class
  $D$, or for $\mathrm{carr}(a)$ when the goal is over an abstract ring
  $a$ with $\mathrm{is\text{-}commutative\text{-}ring}(a)$ in context; the
  typing of a quantifier the oracle peels certifies its variable.
\item[Verified by] the checker of the rewrite, which is a second decision
  rather than a certificate check: it recomputes the normal forms of both
  sides under the same typing test and refuses when they differ.
\item[Does not] type a generator itself.  An uncertified generator leaves
  the goal untouched, with the generator named in the refusal; the
  command \texttt{type-term} lands the typing.
\end{description}

\subsection{The definedness certificate}
\index{certificate!definedness}\index{LUTINS}

\begin{description}
\item[Certifies] that a term $t$ denotes, in the context $\Gamma$, which
  is what \texttt{forall-elim} needs before instantiating a universal at
  $t$ and what \texttt{rfl} needs before closing $t = t$, since \texttt{=}
  is strict (the LUTINS rule, Section~\ref{sec:kernel-ops-detail}).
\item[Issued by] a fixed test on $\Gamma$ and $t$, \texttt{pi--defined?}:
  a variable or atomic constant, a ground number, a term typed or equated
  in context or occurring outside binders in a true membership, equation
  or order hypothesis, a class term on defined arguments, a lambda,
  separation or indexed union over a certified domain, arithmetic on
  number-typed arguments, an accessor of a structure whose predicate is in
  context, an applied structure operation on typed arguments, $f(a)$ with
  $f \in \mathrm{fun}(D, \_)$ and $a \in D$; and, for a predicate
  hypothesis whose defining iff is definitional, primitive or proven, what
  the conjuncts of that iff certify, for the parameters only.  Never a
  choice, a description, an untyped application, an entry without its
  range typing, a finite sum, or the description-bodied matrix
  constructors.
\item[Form] a decision, recomputed wherever it is needed; nothing is
  stored.  When the test declines, \texttt{forall-elim} posts the side
  sequent $\Gamma \Rightarrow t = t$ as an owed leaf, in the same context
  as its parent, and the driver discharges it by landing the typing
  (\texttt{dk-discharge-owed!}).
\item[Verified by] the checker of \texttt{forall-elim}, which re-decides
  whether the omitted side sequent was permitted, independently of the
  builder.
\item[Does not] reach a witness: the symbol a stripped existential
  introduces is fresh and certifies nothing until it is typed.  The policy
  and its cases are \texttt{docs/definedness-instantiation-2026-09-18.md}.
\end{description}

\subsection{The proof certificate}
\index{certificate!proof}\index{VNB\_CERTIFIED@\texttt{VNB\_CERTIFIED}}
\index{exam}

\begin{description}
\item[Certifies] that the proof of a named theorem ran to \texttt{qed} and
  passed the inference check, on a stated kernel, and what that proof
  depended on.  It is the mechanism by which the
  everyday load installs a theorem without running its proof.
\item[Issued by] a load that ran the proof: the \emph{exam},
  \texttt{VNB\_CERTIFIED=off ./prover --build-band}, which runs every proof
  and rewrites the store; or an ordinary load meeting a theorem with no
  valid record, which proves it and writes one.
\item[Form] one file per proof file, \texttt{certificates/}$\langle$load
  key$\rangle$\texttt{.cert}, one record per \texttt{qed} in installation
  order.  A record holds the theorem's name, the hash of its statement and
  the statement itself; \texttt{cites}, the names the ledger read off the
  script, each with the hash of its installed statement; \texttt{defs},
  the closure of the statement's defined heads through their definitions
  (the bodies of \texttt{def-functoid}s, the axioms of \texttt{def-predicate}s
  and \texttt{def-constant}s; proven theorems are not followed), with their hashes;
  \texttt{uses}, every macete the kernel saw fire in the proof (an accessor
  read by \texttt{slot} among them), and every \texttt{def-functoid} body and
  predicate defining iff a checker read while accepting an inference, with
  the hash of the rewrite rule or definition, the kernel's own list kept
  independent of the ledger's; the oracles the script called and the bill and oracles
  \texttt{qed} announced; and the kernel hash.  A record carries no date and no host.  A hash is
  the MD5 of the canonical text of a formula, with counter-minted names
  relabelled in order of first occurrence so that it does not depend on
  how much was proved before.  The \emph{kernel hash} is the MD5 of the
  source of the files in \texttt{*cert-kernel-files*}
  (\texttt{certificates.scm}): the files that may call
  \texttt{dg-apply-rule!}, the deduction-graph file, \texttt{expressions},
  \texttt{sequents} and the four checker files; tactics, the kit and
  the drivers are not in it, because every inference they produce is
  checked.
\item[Read by] \texttt{sp}, at load, under \texttt{VNB\_CERTIFIED=on}
  (the default) or \texttt{strict}.  A record is valid when the statement
  being started hashes to its statement hash, every cited, defined and
  used name hashes the same in this tree, and the kernel hash is this
  tree's.  On a valid record the theorem is installed with provenance
  \texttt{certified}, its bill recomputed from the recorded citations, and
  the driver's steps are skipped.  Under \texttt{on} a theorem without a
  valid record is proved and its record written; under \texttt{strict}
  nothing is proved, the uncovered theorems are listed and the load exits
  non-zero.  The bill line of such a theorem reads \texttt{certified modulo
  \{\ldots\}} (a record carries no date and no host, so an exam that changes
  nothing rewrites nothing); the load summary counts \texttt{proven} and
  \texttt{certified} separately, and the word \texttt{proven} alone means
  the proof ran in this image.  The page audit covers the proofs that ran.
\item[Invalidated by] a change to the statement, to the statement of a
  cited name, to a definition the statement mentions or a macete the proof
  fired, or to any kernel file, which invalidates every record at once.  A
  record without a \texttt{uses} field, written before 2026-09-26, is not
  valid.
\item[Does not] certify what the record cannot see: a functoid reduced by
  \texttt{beta} needs no entry (a \texttt{beta} contracts only a lambda record
  that reached the goal through the statement, a cited theorem or a named
  unfold, each of which is keyed), and the exam is the backstop for what
  the kernel hash does not cover: the tactics and the kit.  A certificate is not a proof:
  it says a proof ran and was checked.  The design is
  \texttt{docs/certificates-2026-09-24.md}.
\end{description}

\subsection{The bill, and \texttt{proof-rules}}
\index{bill}\index{proof-rules@\texttt{proof-rules}}

Two further things are certified about a completed proof, neither called
a certificate in the code.  The \emph{bill} that \texttt{qed} prints,
\texttt{proven modulo \{\ldots\} [trust: \ldots]}, is the kernel's
certification of the theorem relative to the asserted facts it names
(Section~\ref{sec:proof-debt}): the theorem holds if those do.  The
command \texttt{(proof-rules)} walks the grounded derivation of the
current proof from its root and tallies every kernel inference by its
rule tag; because every tag is one of the fixed set of kernel operations,
this shows that a proof driven by composite tactics rests on nothing but
those operations, independently of the command trace, which is lossy.

\subsection{The certification stamp of the VNB test}
\index{certificate!VNB test stamp}

The classification pages of Chapter~\ref{ch:classification} report a
result as ``certified by the VNB test on \textit{date}''.  That stamp is
written by \texttt{vnb-test} to \texttt{reference/certification.scm}, one
\texttt{(certify!\ '\textit{name} "\textit{date}")} per theorem the test
re-proved to \texttt{qed}, and read back at the next load.  It certifies a
run of the test, not a proof, and predates the proof certificate above,
which has replaced it as the everyday mechanism: the exam that writes the
store is the run that re-proves the library.

\subsection{What is not a certificate}

A \emph{refutation} is evidence that a goal is false: the countermodel
\texttt{prop} returns when the atoms do not entail the goal, the
counterexample the counterexample lane finds by evaluating the goal at a
value, the non-zero numeral \texttt{zero-it} reports when a difference
normalises to a constant.  Each is printed and none is recorded, because
no inference is made.  A \emph{warrant} (Section~\ref{sec:warrants})
records a claim about a support and its tier of trust; nobody verifies it,
which is why it is debt.  A provenance stamp, \texttt{primitive} or
\texttt{definitional}, is a decision of the user that a fact is not debt,
recorded in the file with its claim.

\section{Starting and inspecting a proof}

\begin{VNBcode}
; Start a new proof of a formula
(define ps (start-proof (make-wff-from-string
  "n in NN and m in NN implies n in NN")))

; Display the current state
(print-proof-state ps)

; Check if done
(proof-done? ps)   ; => #f or #t

; List open goals
(proof-open-goals ps)
\end{VNBcode}

The variable \texttt{ps} holds the \emph{proof state}: a deduction graph
plus the current \emph{focus} sequent node (the subgoal currently being
worked on).  Every proof command returns a new proof state with an
updated focus.

\texttt{start-proof} takes a \texttt{wff} (produced by \texttt{make-wff} or \texttt{make-wff-from-string}).
The root sequent has no assumptions: a proof begins from the bare goal,
and every hypothesis arrives through an inference.  All sequent assumptions and assertions are
\texttt{wff} objects throughout the proof; the library name propagates
automatically to every subformula produced by inference rules.

%% =========================================================
\section{Proof commands}
\label{sec:commands}

Commands are Scheme procedures of the form
\[
  \texttt{(cmd-}\mathit{name}\;\; \mathit{ps}\;\; \mathit{arg}_1 \ldots\texttt{)}
\]
Each command applies one primitive inference to the focus node, then
advances the focus to the first new ungrounded subgoal.

\subsection{Logical connectives}

\begin{itemize}

\item \texttt{(cmd-direct-inference ps)}
  Decompose the goal by its top-level connective:
  \begin{itemize}
  \item \texttt{p and q} $\to$ two subgoals $p$, $q$.
  \item \texttt{p implies q} $\to$ assume $p$, prove $q$.
  \item \texttt{p iff q} $\to$ two subgoals (each direction).
  \item \texttt{not(p)} $\to$ assume $p$, prove \texttt{falsity}.
  \item \texttt{forall([\ldots], p)} $\to$ peel the entire leading
    \texttt{forall} chain in one step, introducing a fresh variable for
    each bound variable.  If the substituted body is
    \texttt{$y$ in $A$ implies \ldots}, absorb \texttt{$y$ in $A$} into
    the assumptions and continue peeling.  Stop at the first
    non-\texttt{forall} head.
    Example: goal \texttt{forall([x in A, y], body)} yields one subgoal
    $x_0 \in A \;\Longrightarrow\; \texttt{body}[x/x_0,\,y/y_0]$.
  \end{itemize}

\item \texttt{(cmd-antecedent-inference ps $f$)}
  Decompose assumption $f$ by its top-level connective:
  \begin{itemize}
  \item \texttt{p and q} $\to$ replace with $p$ and $q$ separately.
  \item \texttt{p or q} $\to$ case split: branch on $p$, branch on $q$.
  \item \texttt{not(p)} $\to$ derive \texttt{falsity} if $p$ also in context.
  \item \texttt{forsome([\ldots], p)} $\to$ introduce a fresh witness for $x$.
  \end{itemize}

\item \texttt{(cmd-assumption ps)}
  Close goal immediately if it matches an assumption or is \texttt{truth}.

\item \texttt{(cmd-arith ps)}
  Close the current goal by ground arithmetic evaluation.  The goal must
  be a closed formula (no free variables) built from numeric literals,
  $+$, $*$, $-$, \texttt{recip}, \texttt{abs}, \texttt{succ},
  \texttt{conjugate}, $\leq$, $=$,
  $\in \mathtt{nn}/\mathtt{zz}/\mathtt{qq}/\mathtt{rr}/\mathtt{cc}$,
  and the propositional connectives.  Exact (rational) arithmetic only;
  no axiom chain is constructed.
  Related: \texttt{(vnb-do-arith \textit{f})} returns \texttt{\#t}/\texttt{\#f}
  for use outside a proof.

\item \texttt{(cmd-ring-simplify ps)}
  Close a goal by reducing both sides to their canonical form in the free
  associative $\mathbb{Z}$-algebra.  Addition is commutative (both sides
  are brought to the same sorted canonical order); multiplication is
  non-commutative.  Generators are detected automatically and no generator
  list need be supplied: \texttt{(rs)} treats any \emph{term} that
  \texttt{arith-eval} cannot reduce to a number as a generator of the free
  ring --- a bare symbol such as \texttt{x}, or a compound term such as
  \texttt{f(x)} whose head is not one of $+\,-\,*$.  A compound generator is
  a single opaque indeterminate: \texttt{f(x)} is atomic and its argument is
  not examined.

  The goal may be in either of two forms:
  \begin{itemize}
  \item \textbf{Bare equality} \texttt{$e_1$ = $e_2$}: generators are
    certified by \texttt{(in $v$ $D$)} assumptions already in the sequent
    context.  Typically reached after \texttt{(di)} peels the leading
    \texttt{forall} chain.
  \item \textbf{Universally quantified identity}
    \texttt{forall([$x_1$ in $D_1$, \ldots], $e_1$ = $e_2$)}: the
    quantifier typings themselves certify the generators.  \texttt{(di)}
    is not needed; \texttt{(rs)} closes the goal in one step.
  \end{itemize}

  In both cases every generator must be certified as an element of one of
  the five ring domains $\mathtt{nn}$, $\mathtt{zz}$, $\mathtt{qq}$,
  $\mathtt{rr}$, $\mathtt{cc}$ --- a symbol \texttt{x} by an
  \texttt{(in x $D$)} typing, a compound generator \texttt{f(x)} by an
  \texttt{(in f(x) $D$)} assumption.  This certification is also what keeps
  the rule sound on partial terms.  Since \texttt{=} is strict (it requires
  both sides defined), \texttt{3 + f(x) = 1 + f(x) + 2}
  holds only when \texttt{f(x)} is defined; membership entails definedness,
  so demanding \texttt{(in f(x) $D$)} supplies exactly that guarantee.  A
  generator that is not certified leaves the goal untouched --- the command
  never makes an unsound step.

  \begin{VNBcode}
;;; Directly on a universally quantified goal -- no (di) needed:
(sp (make-wff-from-string
  "forall([a in ZZ, b in ZZ, c in ZZ], a + b + c = c + b + a)"))
(rs)    ; closes in one step

;;; After (di), a compound term certified in context becomes a generator:
(sp (make-wff-from-string
  "forall([x in ZZ], (f(x) in ZZ) implies 3 + f(x) = 1 + f(x) + 2)"))
(di) (di) (rs)    ; di peels the forall then the hypothesis; rs then closes

;;; Coefficient arithmetic:
(sp (make-wff-from-string "forall([x in ZZ], 2*x + 3*x = 5*x)"))
(rs)
  \end{VNBcode}

  The primitive inference is \texttt{pi-ring-simplify!}, in
  \texttt{structure-library/ring-simplify.scm}; it normalizes each side to a
  canonical sum-of-monomials polynomial (\texttt{vnb->poly}) and compares.  The interactive short form is \texttt{(rs)}.

\item \texttt{(cmd-comm-ring-simplify ps)}
  The commutative counterpart of \texttt{(cmd-ring-simplify ps)}.  Where
  \texttt{(rs)} fixes the order of factors within a monomial,
  \texttt{(crs)} treats multiplication as commutative: a monomial is a
  \emph{multiset} of generators, canonicalised as a sorted generator
  list, so \texttt{x*y} and \texttt{y*x} share one normal form and
  commutativity lives in the representation rather than in any rewrite
  step.  Each side reduces to a sum of such monomials over \texttt{zz},
  and the goal closes iff the two normal forms coincide.  Use it for any
  identity whose proof needs reordering of factors, which \texttt{(rs)}
  cannot perform.

  The goal may take the same two shapes as for \texttt{(rs)} (a bare
  equality with generators certified by context typings, or a typed
  universally quantified identity), and additionally a \textbf{generic}
  form over an arbitrary ring: \texttt{forall($R$, is-commutative-ring($R$)
  implies \ldots)} built from \texttt{(add $R$)}, \texttt{(mul $R$)},
  \texttt{(neg $R$)}, \texttt{(zero $R$)}, \texttt{(one $R$)}, with the
  generators ranging over the carrier \texttt{(carr $R$)}.

  Generators are detected and certified exactly as for \texttt{(rs)}: a
  symbol or an opaque compound term, certified by an \texttt{(in $v$ $D$)}
  typing in the concrete form or an \texttt{(in $v$ (carr $R$))} assumption in
  the generic form.  The same partiality caveat applies --- a compound
  generator fires only once certified, which in VNB guarantees it is
  defined.

  \begin{VNBcode}
;;; Reordering of factors -- (rs) cannot close this, (crs) does:
(sp (make-wff-from-string
  "forall([x in ZZ, y in ZZ], x*y*x = x*x*y)"))
(crs)   ; closes in one step

;;; Commutative expansion:
(sp (make-wff-from-string
  "forall([x in ZZ, y in ZZ], (x+y)*(x+y) = x*x + 2*x*y + y*y)"))
(crs)
  \end{VNBcode}

  The procedure is a trusted oracle (it compares canonical normal forms
  rather than mechanising the ring properties), carrying the warrant
  \texttt{comm-ring-simplify} of kind \texttt{well-known}.  The primitive
  inference is \texttt{pi-comm-ring-simplify!} in
  \texttt{comm-ring-simplify.scm}.  The interactive short form is
  \texttt{(crs)}.

\item \texttt{(cmd-reflexivity ps)}
  Close goal \texttt{a = a} when \texttt{a} is demonstrably defined.
  Fails if \texttt{a} may be undefined.

\item \texttt{(cmd-quasi-reflexivity ps)}
  Close goal \texttt{a == a} unconditionally, for any term \texttt{a}.

\item \texttt{(cmd-cut ps $L$)}
  Introduce lemma $L$: subgoal~1 proves $L$; subgoal~2 has $L$ as assumption.

\item \texttt{(cmd-or-intro-left ps)}
  To prove \texttt{p or q}, prove $p$ instead.

\item \texttt{(cmd-or-intro-right ps)}
  To prove \texttt{p or q}, prove $q$ instead.

\item \texttt{(cmd-proof-by-contradiction ps)}
  To prove $p$, assume \texttt{not(p)} and derive \texttt{falsity}.

\item \texttt{(cmd-instantiate ps $F$\ $t$)}
  From \texttt{forall([$x$], $p$)} $= F$ in context, add $p[x/t]$.

\item \texttt{(cmd-exists-witness ps $t$)}
  To prove \texttt{forsome([$x$], $p$)}, prove $p[x/t]$.

\item \texttt{(cmd-weaken ps $f$)}
  Remove formula $f$ from the assumptions.

\item \texttt{(cmd-backchain ps $F$)}
  From \texttt{p implies goal} $= F$ in context, generate subgoal~$p$.

\item \texttt{(cmd-theorem-assumption ps $\mathit{name}$)}
  Add the named theorem (or axiom) to the assumptions.

\item \texttt{(cmd-apply-macete ps $\mathit{name}$)}
  Apply named macete: rewrite the goal using the associated theorem.

\item \begin{sloppypar}\texttt{(cmd-eq-subst ps $F$)}
  Equality substitution (Leibniz).  From an equality \texttt{$s$ = $t$} $= F$
  in the assumptions (either orientation), rewrite every occurrence of $s$
  in the goal to $t$.  Short form \texttt{(subst $F$)}.  This realises the
  schema \texttt{$s$ = $t$} $\wedge$ \texttt{$P[t]$} $\Rightarrow$
  \texttt{$P[s]$} for an arbitrary context $P$; sound for partial equality
  because \texttt{$s$ = $t$} already certifies both sides
  defined.\end{sloppypar}

\item \texttt{(cmd-qed ps $\mathit{name}$)}
  After a proof completes, install its root assertion as a named theorem.
  Errors if the proof is not complete.  This is the only command that
  records a result as \emph{proven}.

\end{itemize}

\subsection{N-ary constructors}
\label{sec:nary-proof-commands}

\begin{itemize}

\item \begin{sloppypar}\texttt{(cmd-cartesian-intro ps)}
  To prove \texttt{[$a_1$,\ldots,$a_n$] in cartesian($A_1$,\ldots,$A_n$)},
  generate $n$ subgoals \texttt{$a_i$ in $A_i$}.\end{sloppypar}

\item \begin{sloppypar}\texttt{(cmd-cartesian-elim ps $F$\ $k$)}
  From \texttt{[$a_1$,\ldots,$a_n$] in cartesian($A_1$,\ldots,$A_n$)} $= F$
  in context, add \texttt{$a_k$ in $A_k$}.\end{sloppypar}

\item \texttt{(cmd-tuples-intro ps)}
  To prove \texttt{[$a_1$,\ldots,$a_n$] in tuples($A$)},
  generate $n$ subgoals \texttt{$a_i$ in $A$}.
  For the empty sequence \texttt{[]} the goal closes immediately (0 subgoals).

\item \begin{sloppypar}\texttt{(cmd-tuples-elim ps $F$\ $k$)}
  From \texttt{[$a_1$,\ldots,$a_n$] in tuples($A$)} $= F$
  in context, add \texttt{$a_k$ in $A$}.\end{sloppypar}

\item \texttt{(cmd-nth-reduce ps)}
  Rewrite every occurrence of \texttt{nth($k$, [$\ldots$])} in the goal to
  the $k$-th element (1-indexed).

\item \begin{sloppypar}\texttt{(cmd-functoid-beta ps)}
  Rewrite every occurrence of
  \texttt{lambda([$x_1$ in $A_1$,\ldots], $\mathit{body}$)($v_1$,\ldots)}
  (or \texttt{lambdoid}) in the goal to
  $\mathit{body}[x_1 := v_1]\cdots$ (capture-avoiding).\end{sloppypar}

\item \texttt{(cmd-union-intro ps $k$)}
  To prove \texttt{$x$ in union($A_1$,\ldots,$A_n$)},
  prove \texttt{$x$ in $A_k$} instead.

\item \texttt{(cmd-union-elim ps $F$)}
  From \texttt{$x$ in union($A_1$,\ldots,$A_n$)} $= F$ in context,
  case-split: one branch per $A_i$ with \texttt{$x$ in $A_i$} assumed.

\item \begin{sloppypar}\texttt{(cmd-intersection-intro ps)}
  To prove \texttt{$x$ in intersection($A_1$,\ldots,$A_n$)},
  generate $n$ subgoals \texttt{$x$ in $A_i$}.\end{sloppypar}

\item \begin{sloppypar}\texttt{(cmd-intersection-elim ps $F$\ $k$)}
  From \texttt{$x$ in intersection($A_1$,\ldots,$A_n$)} $= F$ in context,
  add \texttt{$x$ in $A_k$}.\end{sloppypar}

\end{itemize}

%% =========================================================
\section{Two surfaces, one script}
\label{sec:two-surfaces}

Every proof step exists in two forms.

The \emph{builder} \texttt{(cmd-\textit{name}\ \textit{ps}\ \textit{arg}
\ldots)} (Section~\ref{sec:commands}) takes a proof state explicitly and
\emph{returns} a new one; it has no side effects and never prints.  It is
the form to call from Scheme code that manipulates proof states as values.

The \emph{interactive short form} \texttt{(\textit{name}\ \textit{arg}
\ldots)} (Section~\ref{sec:interactive}) is a thin wrapper around the
builder: it reads the global proof state \texttt{*ps*}, applies the
builder, stores the result back in \texttt{*ps*}, and calls \texttt{(show)}
so the new state is displayed.  This is the form typed in a live session.

\begin{center}
\begin{tabular}{@{}lll@{}}
\toprule
Builder & Short form & \\
\midrule
\texttt{(cmd-direct-inference ps)}        & \texttt{(di)}    & no argument \\
\texttt{(cmd-antecedent-inference ps $f$)}& \texttt{(ai $f$)}& formula argument \\
\texttt{(cmd-backchain ps $F$)}           & \texttt{(bc $F$)}& formula argument \\
\texttt{(cmd-instantiate ps $F$ $t$)}     & \texttt{(inst $F$ $t$)} & formula + term \\
\texttt{(cmd-apply-macete ps $n$)}        & \texttt{(mac $n$)} & theorem name \\
\bottomrule
\end{tabular}
\end{center}

\noindent
The wrapper does one more thing: it \emph{records itself}.  Each short form
appends an entry \texttt{(\textit{name}\ .\ \textit{args})} to the global
accumulator \texttt{*proof-script*}, which \texttt{(sp \ldots)} clears at
the start of a proof.  So there is no separate scripting language to learn:
\textbf{the script is the transcript} --- the very sequence of short forms
you typed.  \texttt{(qed \textit{name})} and \texttt{(save-proof
\textit{name})} store that sequence; \texttt{(replay-proof \textit{name})}
re-runs it (Section~\ref{sec:replay}).

\subsection{Supplying arguments}
\label{sec:arg-collection}

A short form takes one of a few kinds of argument, and each kind is
collected the same way wherever it appears.

\begin{itemize}
\item \textbf{Nothing.}  The structural steps --- \texttt{di}, \texttt{ass},
  \texttt{pbc}, \texttt{rs}, \texttt{rfl}, \texttt{ci}, \texttt{ni},
  \ldots\ --- read the focus goal and need no argument.

\item \textbf{A formula} (\texttt{ai}, \texttt{bc}, \texttt{wk},
  \texttt{cut}, \texttt{ue}, the assumption cited by \texttt{inst} and
  \texttt{mac-h}).  It may be supplied three interchangeable ways:
  \begin{itemize}
  \item a \textbf{surface string}, parsed for you ---
    \texttt{"n in nn and m in nn"};
  \item a \textbf{raw S-expression} ---
    \texttt{'(and (in n nn) (in m nn))};
  \item for a command that cites an existing \emph{assumption}, a
    \textbf{1-based assumption index}: a positive integer $k$ selects the
    $k$-th assumption of the focus sequent, exactly as numbered in the
    display.  \texttt{(ai 1)} stands in for retyping the first assumption.
  \end{itemize}

\item \textbf{A term} (\texttt{ew}, the witness slot of \texttt{inst},
  \texttt{if-true}/\texttt{if-false}).  A string or an S-expression; there
  is no index form, since a term is not an assumption to point at.

\item \textbf{A name} (\texttt{ta}, \texttt{mac}, \texttt{bc*},
  \texttt{qed}).  A symbol naming an installed theorem or macete.

\item \textbf{A small integer} (\texttt{ui $k$}, \texttt{ce $f$ $k$},
  \texttt{focus $n$}).  Selects a branch, a tuple component, or an open
  goal --- not an assumption index.
\end{itemize}

\begin{VNBcode}
(sp (make-wff-from-string "n in NN and m in NN implies n in NN"))
(di)        ; assume the conjunction (now assumption #1)
(ai 1)      ; crack assumption #1 -- the index, not the retyped formula
(ass)
(qed 'and-elim-demo)

(lookup-proof 'and-elim-demo)
; => ((di) (ai 1) (ass))     -- the index is recorded verbatim
\end{VNBcode}

\begin{remark}
The assumption-index form is accepted by the commands that \emph{cite} an
assumption --- \texttt{ai}, \texttt{bc}, \texttt{wk}, \texttt{inst},
\texttt{ce}, \texttt{te}, \texttt{ie}, \texttt{mac-h}.  It is recorded
verbatim in the script, so a replay re-resolves ``assumption~1'' against
the focus sequent it actually reaches --- robust to the formula's exact
text, but dependent on the assumption \emph{order} being the same on
replay.  When that is in doubt, give the formula itself.
\end{remark}

\subsection{Argument collection in a front end: \texttt{bc*}, \texttt{B}, \texttt{B+}}
\label{sec:bc-star-collection}

\texttt{bc*} (Section~\ref{sec:interactive}) matches a theorem's conclusion
against the goal and instantiates whatever that match determines.  A schema
variable the match leaves \emph{undetermined} must be supplied --- as a
binding list \texttt{((\textit{v}\ \textit{val})\ \ldots)} of Scheme
expressions.  The pure query \texttt{(bc*-undetermined '\textit{name})}
reports exactly which variables those are, without touching the proof.

A front end uses that query to collect the values \emph{by name} instead of
making the user write the quoted binding list: it asks for each undetermined
variable's value as an ordinary \emph{term string} (\texttt{"ran(f)"},
\texttt{"nn"}) and parses it.  In the Emacs Focus workspace this is the
\texttt{B} key --- the argument \emph{builder}.  Its companion key
\texttt{B+} is a saturating closer: it only commits a citation it
can first \emph{prove} will close the focus goal
(\texttt{bc*-can-close?} matches the conclusion \emph{and} checks every
hypothesis is already an assumption), so it never leaves the proof worse
than it found it.

%% =========================================================
\section{Interactive short-form commands}\label{sec:tactics-ref}
\label{sec:interactive}

\texttt{interactive.scm} provides the short-form commands --- one- or two-letter
wrappers around the \texttt{cmd-*} proof procedures --- that you type in a live
proof session, at the bare REPL or through the Emacs workspace.  Each mutates the
global proof state \texttt{*ps*} and calls \texttt{(show)} automatically; the
exceptions, which return a value rather than mutating state, are \texttt{sp},
\texttt{show}, \texttt{pp}, \texttt{fa}, \texttt{fs}, \texttt{vnb-unwind},
\texttt{vnb-wind}, \texttt{parse-string}, and \texttt{make-wff-from-string}.

The reference below is \emph{generated} from the live \texttt{*tactic-help*}
registry (through \texttt{reference/TACTICS.md}), so it lists the current commands
with their signatures and trust \emph{kind} and cannot drift from the prover.  At
the REPL, \texttt{(tactics)} prints the same menu and \texttt{(tactics 'name)} the
full per-command help.

%% AUTO-GENERATED by docs/gen-tactics-ref.py from reference/TACTICS.md -- do not edit.
%% Regenerate:  python3 docs/gen-tactics-ref.py   (after a prover load refreshes TACTICS.md)

\DeclareRobustCommand{\tactkind}[1]{{\footnotesize\textsc{#1}}}

Every tactic is tagged with a \emph{kind} grounded in the \texttt{dg-apply-rule!} tag it emits, so the trusted base is legible: \texttt{rule} is a single primitive kernel inference (the fixed trusted base); \texttt{oracle} is a trusted decision procedure run as a black box; \texttt{composite} only chains kernel rules, adding no new inference; \texttt{meta} performs no deduction (session, search, navigation).  A proof's trust surface is exactly its \texttt{rule} steps plus whichever \texttt{oracle}s and asserted premises it cites.

The same reference is available at the REPL with \texttt{(tactics 'name)} and in the browser file \texttt{reference/TACTICS.md}; all three are generated from one registry, so they never disagree.

An entry's \emph{Uses} line names every kernel operation the command's own code can cause to be recorded, those recorded by the procedures it calls included. It is the bound on what the command can do, whether or not any proof in the library exercises it, and it is measured by reading the source rather than declared alongside the command. One further set is common to every command and is stated once instead of in every entry: the hook that runs after any command which changed the proof, to close the definedness sequents an instantiation left owed. Section~\ref{sec:command-uses} tabulates the same figures for every command at once.

\subsection*{Starting \& finishing a proof}

\par\smallskip{\raggedright\noindent\texttt{(sp goal)}\quad\tactkind{meta}\index{sp@\texttt{sp} (tactic)}\par}\nobreak
\noindent Start a proof of `goal' (a ``string'', a raw S-expr, or a wff); clears the script.\par
\noindent\textit{Uses:} no kernel operation: it records no inference.\par
\noindent\textit{When useful:} starting a new proof\par

Begin a new proof: state what you want to prove.  The goal may be a ``string'' in surface syntax, a raw S-expr, or an already-built wff (e.g. (wff ``...'') or (make-wff '(...)));  sp coerces a string/S-expr for you.  This becomes your one open goal and clears any previous script.  (Technically: set the proof goal.)

\par\smallskip{\raggedright\noindent\texttt{(qed 'name)}\quad\tactkind{meta}\index{qed@\texttt{qed} (tactic)}\par}\nobreak
\noindent Install the finished proof as theorem `name' and save its replayable script.\par
\noindent\textit{Uses:} no kernel operation: it records no inference.\par
\noindent\textit{When useful:} no open goals remain -- record the result\par

Finish: once there are no open goals, this records the result as a named theorem you can cite later, and saves the replayable script.  (Technically: install-theorem! plus script capture; reports the asserted facts the proof still rests on, `proven modulo \{...\}'.)

\par\smallskip{\raggedright\noindent\texttt{(save-proof 'name)}\quad\tactkind{meta}\index{save-proof@\texttt{save-proof} (tactic)}\par}\nobreak
\noindent Save the current script under a name without finishing.\par
\noindent\textit{Uses:} no kernel operation: it records no inference.\par

Save the current (possibly unfinished) proof script under a name, to resume or replay later.  (Technically: snapshot the script without installing a theorem.)

\par\smallskip{\raggedright\noindent\texttt{(replay-proof 'name [subst])}\quad\tactkind{meta}\index{replay-proof@\texttt{replay-proof} (tactic)}\par}\nobreak
\noindent Re-run a saved script on the current goal, optionally renaming free vars.\par
\noindent\textit{Uses:} every one of the 64 kernel operations.\par

The optional subst is an alist ((old . new) ...) applied to every command argument before replay -- this is how one proof is reused at fresh eigenvariables.


\subsection*{Decomposing the goal}

\par\smallskip{\raggedright\noindent\texttt{(di)}\quad\tactkind{rule}\index{di@\texttt{di} (tactic)}\par}\nobreak
\noindent Direct inference: split an AND goal, move an IMPLIES antecedent into the assumptions, or introduce a fresh eigenvariable for a leading FORALL.\par
\noindent\textit{Uses:} \texttt{and-intro} \texttt{forall-intro} \texttt{iff-intro} \texttt{implies-intro} \texttt{not-intro}.\par
\noindent\textit{When useful:} the goal head is AND / IMPLIES / FORALL -- decompose before deciding\par

Break the goal into its pieces.  A conjunction `A and B' splits into two goals, A and B.  An implication `P implies Q' assumes P (P becomes a hypothesis you may use) and leaves you to prove Q.  A universal `for all x, ...' fixes an arbitrary x (a new constant standing for `any x') and asks for the statement at that x.  Apply it repeatedly until the goal is a single atomic statement.  (Technically: the goal-side introduction rules for AND / IMPLIES / FORALL; the constant introduced for a FORALL is an eigenvariable.  Pairs with mac, which unfolds a definition in the goal.)

\par\smallskip{\raggedright\noindent\texttt{(pbc)}\quad\tactkind{rule}\index{pbc@\texttt{pbc} (tactic)}\par}\nobreak
\noindent Proof by contradiction: assume the goal's negation, prove FALSITY.\par
\noindent\textit{Uses:} \texttt{proof-by-contradiction}.\par
\noindent\textit{When useful:} a direct argument stalls and assuming the negation gives something concrete to work with\par

Argue by contradiction: assume the goal is false and derive an absurdity.  (Technically: classical reductio -- adds the negation of the goal as a hypothesis and changes the goal to FALSITY.)

\par\smallskip{\raggedright\noindent\texttt{(oi-l)}\quad\tactkind{rule}\index{oi-l@\texttt{oi-l} (tactic)}\par}\nobreak
\noindent OR-intro left: reduce an (OR a b) goal to a.\par
\noindent\textit{Uses:} \texttt{or-intro-left}.\par
\noindent\textit{When useful:} the goal is (OR a b) and the LEFT alternative is provable\par

Prove a disjunction `A or B' by proving the LEFT alternative, A.  (Technically: or-introduction, left.)

\par\smallskip{\raggedright\noindent\texttt{(oi-r)}\quad\tactkind{rule}\index{oi-r@\texttt{oi-r} (tactic)}\par}\nobreak
\noindent OR-intro right: reduce an (OR a b) goal to b.\par
\noindent\textit{Uses:} \texttt{or-intro-right}.\par
\noindent\textit{When useful:} the goal is (OR a b) and the RIGHT alternative is provable\par

Prove a disjunction `A or B' by proving the RIGHT alternative, B.  (Technically: or-introduction, right.)

\par\smallskip{\raggedright\noindent\texttt{(ew term)}\quad\tactkind{rule}\index{ew@\texttt{ew} (tactic)}\par}\nobreak
\noindent Existential witness: discharge a FORSOME goal by supplying the witness term.\par
\noindent\textit{Uses:} \texttt{forsome-intro}.\par
\noindent\textit{When useful:} the goal is FORSOME and you have a witness term in mind\par

Prove `there exists an x with property P' by exhibiting a specific witness: you supply the term, and the goal becomes `P holds of that term'.  The everyday `take x = ...' step.  (Technically: existential introduction for a FORSOME goal.)

\par\smallskip{\raggedright\noindent\texttt{(ci)}\quad\tactkind{rule}\index{ci@\texttt{ci} (tactic)}\par}\nobreak
\noindent Cartesian intro: prove a CARTESIAN-product membership component-wise.\par
\noindent\textit{Uses:} \texttt{cartesian-intro}.\par
\noindent\textit{When useful:} the goal is membership in a CARTESIAN product\par

Prove that a pair (or tuple) belongs to a Cartesian product A x B by proving each coordinate lies in its factor.  (Technically: Cartesian-membership introduction, component-wise.)

\par\smallskip{\raggedright\noindent\texttt{(ti)}\quad\tactkind{rule}\index{ti@\texttt{ti} (tactic)}\par}\nobreak
\noindent Tuple intro: prove a tuple/LIST membership component-wise.\par
\noindent\textit{Uses:} \texttt{tuples-intro}.\par
\noindent\textit{When useful:} the goal is membership in a set of tuples / LIST\par

Prove that a list belongs to the set of tuples over A by proving each entry lies in A.  (Technically: tuple-membership introduction, component-wise.)

\par\smallskip{\raggedright\noindent\texttt{(ii)}\quad\tactkind{rule}\index{ii@\texttt{ii} (tactic)}\par}\nobreak
\noindent Intersection intro: prove (IN x (INTERSECTION ...)) for each branch.\par
\noindent\textit{Uses:} \texttt{intersection-intro}.\par
\noindent\textit{When useful:} the goal is membership in an INTERSECTION\par

Prove that x lies in an intersection of sets by proving it lies in each of them.  (Technically: intersection-membership introduction.)

\par\smallskip{\raggedright\noindent\texttt{(ui k)}\quad\tactkind{rule}\index{ui@\texttt{ui} (tactic)}\par}\nobreak
\noindent Union intro: reduce a goal (IN x (UNION ...)) to membership of x in the k-th set (1-based).\par
\noindent\textit{Uses:} \texttt{union-intro}.\par
\noindent\textit{When useful:} the goal is membership in a UNION and you know which branch\par

Prove that x lies in a union of sets by proving it lies in the k-th one (you choose which).  (Technically: union-membership introduction at branch k.)

\par\smallskip{\raggedright\noindent\texttt{(ni)}\quad\tactkind{rule}\index{ni@\texttt{ni} (tactic)}\par}\nobreak
\noindent Natural-number induction on the goal's leading FORALL over NN.\par
\noindent\textit{Uses:} \texttt{nn-induction}.\par
\noindent\textit{When useful:} the goal is `forall n in NN, P(n)'\par

Prove `for all natural numbers n, P(n)' by induction: show P(0), and that P(n) implies P(n+1).  (Technically: natural-number induction on the goal's leading FORALL over NN.)


\subsection*{Using a hypothesis}

\par\smallskip{\raggedright\noindent\texttt{(ai hyp)}\quad\tactkind{rule}\index{ai@\texttt{ai} (tactic)}\par}\nobreak
\noindent Antecedent inference: decompose a cited assumption -- AND-split, OR-into-cases, or FORSOME-elimination to a fresh eigenvariable.\par
\noindent\textit{Uses:} \texttt{and-elim} \texttt{forsome-elim} \texttt{iff-elim} \texttt{not-elim} \texttt{or-elim}.\par
\noindent\textit{When useful:} a hypothesis is an AND / OR / FORSOME to break apart\par

Break a hypothesis into its pieces -- the mirror image of di, but on the assumptions instead of the goal.  From a hypothesis `A and B' you get both A and B; from `A or B' you split into two cases (prove the goal in each); from `there exists x, P(x)' you get a fresh name for such an x together with P(x).  Cite the hypothesis by its number in the display, its formula, or a ``string''.  (Technically: the hypothesis-side elimination rules for AND / OR / FORSOME; the fresh name is an eigenvariable.)

\par\smallskip{\raggedright\noindent\texttt{(ass)}\quad\tactkind{rule}\index{ass@\texttt{ass} (tactic)}\par}\nobreak
\noindent Close the goal by an assumption alpha-equivalent to it.\par
\noindent\textit{Uses:} \texttt{assumption}.\par
\noindent\textit{When useful:} the goal already appears (up to bound-var renaming) among the hypotheses\par

Finish the goal because it is already one of your hypotheses -- the `that is exactly what we assumed' step.  Renaming of dummy/bound variables doesn't matter (`there exists x, P(x)' closes `there exists y, P(y)').  (Technically: the assumption rule, matching up to alpha-equivalence.)

\par\smallskip{\raggedright\noindent\texttt{(inst forall-hyp term)}\quad\tactkind{rule}\index{inst@\texttt{inst} (tactic)}\par}\nobreak
\noindent Instantiate a universally-quantified assumption at `term', adding the instance to context.\par
\noindent\textit{Uses:} \texttt{forall-elim}.\par
\noindent\textit{When useful:} you have a forall-hypothesis and a specific value to use it at\par

Use a `for all x, ...' hypothesis at a particular value: you supply the term, and the statement with x replaced by that term is added to your hypotheses.  (Technically: universal instantiation of an assumption.)

\par\smallskip{\raggedright\noindent\texttt{(inst+ forall-hyp term)}\quad\tactkind{composite}\index{inst+@\texttt{inst+} (tactic)}\par}\nobreak
\noindent Instantiate an in-context universal at `term', then forward-detach any guards whose antecedents are in context.\par
\noindent\textit{Uses:} \texttt{arith-forsome} \texttt{arith-ground} \texttt{arith-simplify} \texttt{cut} \texttt{detach} \texttt{forall-elim}.\par
\noindent\textit{When useful:} a GUARDED forall-hyp whose guards are dischargeable from context (e.g. the metric laws post-grind)\par

inst followed by detach: from `forall x. (x in S) => P(x)' and `t in S' already known, this lands `P(t)' directly (peeling nested guards level by level), instead of leaving the guarded implication for you to detach by hand.  The hypothesis-side analogue of `fact' (which assembles a THEOREM); this assembles an in-context UNIVERSAL.  scout's inst lane emits these to close witness-needing goals like the metric laws.  (Technically: pi-instantiate! then pi-detach! while the consequent stays a guard with an in-context antecedent.)

\par\smallskip{\raggedright\noindent\texttt{(mp)}\quad\tactkind{composite}\index{mp@\texttt{mp} (tactic)}\par}\nobreak
\noindent Close a goal that is an INSTANCE of a universal already in your context.  No arguments: the term is DERIVED by matching, not supplied.\par
\noindent\textit{Uses:} \texttt{arith-forsome} \texttt{arith-ground} \texttt{arith-simplify} \texttt{assumption} \texttt{cut} \texttt{detach} \texttt{forall-elim}.\par
\noindent\textit{When useful:} the goal IS an instance of a universal you already have -- no term to supply, it is derived by matching\par

The syllogism, with nothing to fill in.  From `forall([thing in human], thing in mortal)' together with `socrates in human', it closes `socrates in mortal'.  Where inst+ makes you name the term, mp works it out: it matches the universal's conclusion against the goal, which either determines the bound variable or fails outright, and then checks that the universal's guard at that term is something you already have.  Nothing is searched for and nothing is ranked -- if the goal is an instance there is exactly one term it can be.  It DECLINES, naming what it found, when no universal in the context has the goal as an instance, and when more than one does: a closer that silently picks between two universals is one you cannot predict.  One bound variable in this version; for `forall([a in A, b in B], ...)' use inst+.  (Technically: composite -- inst+ at the derived term, then ass.  It adds no kernel rule and no debt; the bill is what typing the two steps by hand would bill, which is nothing.)

\par\smallskip{\raggedright\noindent\texttt{(grind-and-mp)}\quad\tactkind{composite}\index{grind-and-mp@\texttt{grind-and-mp} (tactic)}\par}\nobreak
\noindent Tidy the goal with grind, then close it by modus ponens -- the one-button form of the syllogism, for a sentence typed exactly as it reads.\par
\noindent\textit{Uses:} \texttt{and-elim} \texttt{and-intro} \texttt{arith-forsome} \texttt{arith-ground} \texttt{arith-simplify} \texttt{assumption} \texttt{cut} \texttt{detach} \texttt{forall-elim} \texttt{forall-intro} \texttt{forsome-elim} \texttt{iff-elim} \texttt{iff-intro} \texttt{implies-intro} \texttt{macete-hyp} \texttt{not-elim} \texttt{not-intro} \texttt{or-elim}.\par
\noindent\textit{When useful:} a sentence typed as it reads -- `all men are mortal and socrates is a man implies socrates is mortal' -- closing in one move\par

Runs the two steps a syllogism actually takes on a sentence you have just typed.  `forall([thing in human], thing in mortal) and socrates in human implies socrates in mortal' is not yet a syllogism: it is an implication whose antecedent is a conjunction.  grind does the BOOKKEEPING -- split the AND, move the antecedents into the assumptions -- which leaves the goal `socrates in mortal' with the universal and `socrates in human' as hypotheses; mp then does the LOGIC (universal instantiation, then modus ponens).  They stay separate tactics because they are different kinds of step, and mp on its own is what to reach for once the goal is already tidy.  The recorded script keeps both steps, so the page reads as the two moves rather than as this wrapper.  (Technically: composite -- grind, quietly since on a tidy goal it legitimately declines, then mp.  Adds no kernel rule and no debt.)

\par\smallskip{\raggedright\noindent\texttt{(detach! impl)}\quad\tactkind{rule}\index{detach"!@\texttt{detach"!} (tactic)}\par}\nobreak
\noindent Forward modus ponens: from an in-context (IMPLIES A B) whose A is also in context, leave B in context.\par
\noindent\textit{Uses:} \texttt{detach}.\par
\noindent\textit{When useful:} you have both P and (P => Q) in context and want Q\par

If you have both `P' and `P implies Q' among your hypotheses, this adds `Q' to them.  It grows what you KNOW (the hypotheses) rather than changing the goal -- the forward counterpart of bc.  (Technically: forward modus ponens, the kernel rule pi-detach!; sound since P and P=>Q give Q.)

\par\smallskip{\raggedright\noindent\texttt{(fact 'thm term ...)}\quad\tactkind{composite}\index{fact@\texttt{fact} (tactic)}\par}\nobreak
\noindent Forward APPLICATION of a theorem: bring it in, instantiate its leading universals with the terms, and auto-detach every antecedent already in context, landing the consequent as a hypothesis.\par
\noindent\textit{Uses:} \texttt{arith-forsome} \texttt{arith-ground} \texttt{arith-simplify} \texttt{cut} \texttt{detach} \texttt{forall-elim} \texttt{theorem-assumption}.\par
\noindent\textit{When useful:} a library law `forall x. H(x) => P(x)' whose P you want landed as a hypothesis\par

Handles interleaved forall/implies (e.g. forall s. IS-X(s) => forall a. a in CARR(s) => P): consumes one term per FORALL, detaches each IMPLIES whose antecedent is in context.  The forward-assembly workhorse -- a law `forall x. H(x) => P(x)' becomes the usable fact P in one call, instead of ta + inst* + cut/backchain.  See theorem-library/module-zero-act.scm.

\par\smallskip{\raggedright\noindent\texttt{(ce hyp k)}\quad\tactkind{rule}\index{ce@\texttt{ce} (tactic)}\par}\nobreak
\noindent Cartesian elim: project the k-th component out of a CARTESIAN-membership assumption.\par
\noindent\textit{Uses:} \texttt{cartesian-elim}.\par
\noindent\textit{When useful:} a hypothesis is CARTESIAN-product membership\par

From a hypothesis that a tuple lies in a Cartesian product, extract that its k-th coordinate lies in the k-th factor.  (Technically: Cartesian-membership elimination, k-th projection.)

\par\smallskip{\raggedright\noindent\texttt{(te hyp k)}\quad\tactkind{rule}\index{te@\texttt{te} (tactic)}\par}\nobreak
\noindent Tuple elim: project the k-th component out of a tuple-membership assumption.\par
\noindent\textit{Uses:} \texttt{tuples-elim}.\par
\noindent\textit{When useful:} a hypothesis is tuple membership\par

From a hypothesis that something is a tuple over A, extract that its k-th entry lies in A.  (Technically: tuple-membership elimination, k-th projection.)

\par\smallskip{\raggedright\noindent\texttt{(ie hyp k)}\quad\tactkind{rule}\index{ie@\texttt{ie} (tactic)}\par}\nobreak
\noindent Intersection elim: extract the k-th branch of an INTERSECTION-membership assumption.\par
\noindent\textit{Uses:} \texttt{intersection-elim}.\par
\noindent\textit{When useful:} a hypothesis is INTERSECTION membership\par

From a hypothesis that x lies in an intersection, extract that x lies in the k-th set.  (Technically: intersection-membership elimination.)

\par\smallskip{\raggedright\noindent\texttt{(ue hyp)}\quad\tactkind{rule}\index{ue@\texttt{ue} (tactic)}\par}\nobreak
\noindent Union elim: split a cited (IN x (UNION ...)) membership assumption into one subgoal per set -- union-side case analysis, the dual of `ui'.\par
\noindent\textit{Uses:} \texttt{union-elim}.\par
\noindent\textit{When useful:} a hypothesis is UNION membership -- case-split on it\par

From a hypothesis that x lies in a union, split into cases -- one for each set x might belong to -- and prove the goal in each.  (Technically: union-membership elimination, case analysis; the dual of ui.)

\par\smallskip{\raggedright\noindent\texttt{(cut formula)}\quad\tactkind{rule}\index{cut@\texttt{cut} (tactic)}\par}\nobreak
\noindent Cut: prove `formula' as a side subgoal, then continue the main goal with `formula' added as an assumption (Gentzen cut).\par
\noindent\textit{Uses:} \texttt{cut}.\par
\noindent\textit{When useful:} you want to prove a lemma on the spot, then use it\par

Introduce a lemma you prove on the spot.  You state a formula; it becomes a side goal to prove, and on the main line you may then use it as a hypothesis.  The standard `we first show that ..., and now using it ...' move.  (Technically: Gentzen's cut -- nothing is left assumed-but-unproved, the side goal discharges it.)

\par\smallskip{\raggedright\noindent\texttt{(wk hyp)}\quad\tactkind{rule}\index{wk@\texttt{wk} (tactic)}\par}\nobreak
\noindent Weaken: drop a cited assumption from the context to tidy the hypothesis list.  `hyp' may be a formula, a ``string'', or a 1-based assumption index.\par
\noindent\textit{Uses:} \texttt{weakening}.\par
\noindent\textit{When useful:} the hypothesis list is cluttered with something no longer needed\par

Discard a hypothesis you no longer need, to keep the assumption list readable.  (Technically: weakening -- removing a hypothesis is always sound.)

\par\smallskip{\raggedright\noindent\texttt{(keep hyp ...)}\index{keep@\texttt{keep} (tactic)}\par}\nobreak
\noindent Keep only the cited assumptions and drop every other one, as ONE recorded step.  Each `hyp' may be a formula, a ``string'', or a 1-based assumption index; they are recorded as formulas.\par
\noindent\textit{When useful:} the context is deep and only a few hypotheses matter for the next step (prop, ineq)\par

Prune the context down to what the next step needs -- what a driver does before `prop' or `ineq' in a deep context.  (Technically: one weakening per dropped hypothesis, on the same leaf, recorded once; the driver kit's dk-only! is this command.)


\subsection*{Rewriting}

\par\smallskip{\raggedright\noindent\texttt{(mac 'name)}\quad\tactkind{rule}\index{mac@\texttt{mac} (tactic)}\par}\nobreak
\noindent Rewrite the GOAL with an equivalence macete (unfold a definition, apply an iff/=/== law).  Fires only where the macete's side-conditions already hold in context.\par
\noindent\textit{Uses:} \texttt{cartesian-decompose} \texttt{macete} \texttt{tuple-equality-decompose}.\par
\noindent\textit{When useful:} the GOAL has a defined predicate to unfold / an identity to apply, side-conditions already met\par

Rewrite the goal using a definition or a known equivalence/equality (a `macete').  For instance, replace a defined predicate by what it stands for, or apply an identity.  It fires only where the law's side-conditions already hold in your hypotheses; where they don't, it leaves that spot untouched and looks deeper inside.  (Technically: goal-side rewriting by an equivalence/equality macete; all-or-nothing -- it does NOT spawn unmet side-conditions as goals.  Its hypothesis-side cousin is mac-h; the minor-premise-spawning variant is macm.)

\par\smallskip{\raggedright\noindent\texttt{(slot 'acc)}\index{slot@\texttt{slot} (tactic)}\par}\nobreak
\noindent Reduce a structure ACCESSOR to its projection: (CARR s) becomes (NTH 1 s).  The one door for accessor reductions -- use it, never mac, on an accessor name.\par
\noindent\textit{Uses:} \texttt{cartesian-decompose} \texttt{macete} \texttt{tuple-equality-decompose}.\par

Replace an accessor by the tuple position it stands for: (CARR s) is slot 1, so it becomes (NTH 1 s).  Pair it with nth-r to compute on a concrete structure -- (MUL ZZ-RING) reduces to bintimes.  It refuses anything that is not an accessor.  (Technically: fires the accessor's macete, which def-structure installs at declaration.  It exists so that ALL accessor reductions go through ONE procedure: the reduction is currently global and unconditional -- index k for every argument -- and making it structure-relative later, guarded by IS-X(s), should change this door and not its callers.  accessor-callsite-audit fails the suite if any file fires an accessor macete by name.)

\par\smallskip{\raggedright\noindent\texttt{(slot-h 'acc hyp)}\index{slot-h@\texttt{slot-h} (tactic)}\par}\nobreak
\noindent Reduce a structure ACCESSOR to its projection inside a cited HYPOTHESIS -- `slot', hypothesis-side.\par
\noindent\textit{Uses:} \texttt{macete-hyp}.\par

The hypothesis-side door of slot: replace (PTS RR-MS) by the tuple position it stands for inside an assumption you cite, instead of in the goal.  Before it existed, a proof needing that reduction in a hypothesis reached for mac-h with the accessor macete's own name (rr-ms@pts) -- firing an accessor macete BY NAME, which is what the accessor pin (accessor-callsite-audit) exists to prevent, and which kept the suite red.  It refuses anything that is not an accessor, and it ERRORS rather than guess when the hypothesis mentions two different instances of the same accessor: after the first rewrite the assumption is a different formula, so a silent half-rewrite would be the failure mode.  (Technically: slot's projection-macete lookup run over the cited assumption, then cmd-apply-macete-to-assumption.)

\par\smallskip{\raggedright\noindent\texttt{(macm 'name)}\quad\tactkind{rule}\index{macm@\texttt{macm} (tactic)}\par}\nobreak
\noindent Like mac, but a CONDITIONAL macete fires even when its side-conditions are not yet in context: each unmet condition is left as a new subgoal.\par
\noindent\textit{Uses:} \texttt{cartesian-decompose} \texttt{macete} \texttt{tuple-equality-decompose}.\par
\noindent\textit{When useful:} the GOAL has a CONDITIONAL identity/definition to apply whose side-conditions you are willing to leave as subgoals\par

The minor-premises variant of mac.  mac only rewrites where the law's side-conditions ALREADY hold; macm rewrites regardless, spawning each unmet condition as a fresh goal to discharge -- the goal-side mirror of what mac-h already does on hypotheses.  This is what makes conditional finite-sum lemmas (finsum-*, sum-*, ...) usable as goal rewrites instead of only forward via fact.  (Technically: goal-side rewriting by a conditional macete; it emits the SAME single `macete' kernel rule as mac -- the unmet conditions become that one step's minor-premise subgoals.  It introduces NO new inference rule.)

\par\smallskip{\raggedright\noindent\texttt{(mac-h 'name hyp)}\quad\tactkind{rule}\index{mac-h@\texttt{mac-h} (tactic)}\par}\nobreak
\noindent Rewrite a cited ASSUMPTION in place with an equivalence macete; any side-condition not already in context is spawned as a new subgoal.\par
\noindent\textit{Uses:} \texttt{macete-hyp}.\par
\noindent\textit{When useful:} a HYPOTHESIS has a definition to unfold / an equivalence to apply\par

Like mac, but rewrites inside one of your HYPOTHESES instead of the goal -- replacing it by an equivalent statement (e.g. unfolding a definition in a hypothesis).  If the rewrite law has a side-condition you haven't established, that side-condition becomes a new goal to prove.  (Technically: hypothesis-side rewriting by Leibniz substitution of equivalents; sound because the macete is a genuine equivalence under its side-conditions, which are spawned as goals.)

\par\smallskip{\raggedright\noindent\texttt{(mac-h*)}\quad\tactkind{composite}\index{mac-h*@\texttt{mac-h*} (tactic)}\par}\nobreak
\noindent Saturating mac-h: repeatedly unfold every defined predicate in the hypotheses and split the conjunctions they expose, until nothing is left folded.  One step in the trace instead of a dozen.\par
\noindent\textit{Uses:} \texttt{and-elim} \texttt{forsome-elim} \texttt{iff-elim} \texttt{macete-hyp} \texttt{not-elim} \texttt{or-elim}.\par
\noindent\textit{When useful:} several hypotheses are folded definitions / conjunctions to open at once\par

The hands-free version of mac-h.  Instead of naming each definition and each conjunction by hand -- (mac-h 'IS-METRIC-SPACE A1), split, (mac-h 'is-metric A2), split, ... -- (mac-h*) keeps unfolding any assumption whose head is a defined predicate and splitting any AND assumption, restarting until a full pass changes nothing.  It records as a SINGLE step, so the proof (and its PDF) is far shorter.  It adds no kernel rule: every unfold is the same kernel-checked mac-h, every split the same ai -- just driven to a fixpoint.  Ask (what-now) to see which hypotheses it would touch first.

\par\smallskip{\raggedright\noindent\texttt{(grind)}\quad\tactkind{composite}\index{grind@\texttt{grind} (tactic)}\par}\nobreak
\noindent Saturate the no-choice moves at the focus: decompose the goal connective (di) and break open the hypotheses (mac-h*), repeatedly, until neither fires.  The whole introduce-everything prefix in one step.\par
\noindent\textit{Uses:} \texttt{and-elim} \texttt{and-intro} \texttt{forall-intro} \texttt{forsome-elim} \texttt{iff-elim} \texttt{iff-intro} \texttt{implies-intro} \texttt{macete-hyp} \texttt{not-elim} \texttt{not-intro} \texttt{or-elim}.\par
\noindent\textit{When useful:} right after sp, or any time the focus has connective/definitional structure -- normalize first\par

The deterministic normalizer.  It keeps applying di -- which strips a leading FORALL/IMPLIES/AND, introducing the bound variables and moving antecedents into your hypotheses -- and mac-h*, which unfolds defined predicates and splits conjunctions in the hypotheses, until the goal head is no longer a connective and no hypothesis is foldable.  None of these moves involves CHOOSING a lemma, so there is never anything to reconsider: grind just puts the focus into normal form (e.g. a metric-law goal becomes `(d s)(x,y) = (d s)(y,x)' with the metric's defining properties sitting unfolded in context).  It collapses the di...di mac-h* prefix that bloats proofs into a single step, and adds no kernel rule.

\par\smallskip{\raggedright\noindent\texttt{(scout [depth [branch [nodes]]])}\quad\tactkind{meta}\index{scout@\texttt{scout} (tactic)}\par}\nobreak
\noindent Speculatively search to a bounded depth across INDEPENDENT scratch branches; RETURNS a nested list (number-of-branches (d b) goal best-partials closing-branches) rather than printing.  Never touches the live proof.\par
\noindent\textit{Uses:} every one of the 64 kernel operations.\par
\noindent\textit{When useful:} you are stuck and want the machine to TRY move-sequences (returns data)\par

The copilot's deep lane.  Where (what-now)/(tt) suggest one move, scout actually TRIES sequences: it clones the focus into a fresh deduction graph per branch (so backtracking is free -- a dead branch is just discarded), and does a BEST-FIRST search (frontier ordered by open subgoals left, ties broken deeper-first so it dives toward a closure) whose alphabet is (grind), the closers (ass/rfl/crs/arith), the top rewrite-index (mac) rules, the parameterless backchain (bc*) lemmas, and the INST LANE -- (inst+ <universal hyp> <context-typed term>), instantiating a `forall' hypothesis at a term the context already types and detaching the guard.  It RETURNS the result as data -- (examined-count (depth branch) goal best-partials closing-branches), where best-partials is ((open-goals-left (form ...)) ...) most-reduced first and closing-branches is ((form ...) ...) shortest first -- so you can pick it apart programmatically; use (scout-show ...) for the readable REPL report.  Every move is the real, kernel-checked tactic, so a closing branch is a genuine proof when you adopt it with (scout-run k).  A CITATION GUARD suppresses circular closures: a branch that discharges the goal by citing a library theorem that is the goal -- DIRECTLY (statement alpha-equal to the goal, `P proved by P') or TRANSITIVELY (the cited theorem was itself proven using the goal, i.e. a name alpha-equal to the goal is in its proof-debt closure -- the `compact=>complete by compact=>bongo + bongo=>complete' trap) -- is dropped from closing-branches, the theorem name recorded in *scout-citations*, and scout-show notes `the goal is already theorem X'.  (Loops built purely from independent ASSERTED facts have no live dependency edge to detect; that is warrant/debt hygiene, not scout's to check.)  Bounds: depth (default 6), per-node fan-out branch (default 3), total nodes (default 600).  El cheapo: it does not reason about which move is wise (only the witness terms are guessed, from what the context already types -- a bad guess just dies on the clone); it brute-forces a small tree and you eyeball the survivors.  The inst lane closes the metric laws past grind; goals that need a witness scout can't type (a fresh existential, a constructed term) still surface under best-partials.

\par\smallskip{\raggedright\noindent\texttt{(scout-show [depth [branch [nodes]]])}\quad\tactkind{meta}\index{scout-show@\texttt{scout-show} (tactic)}\par}\nobreak
\noindent Like (scout), but PRINTS the human-readable report (examined count, goal, closing branches or best partials) to the REPL.  Returns the same nested list (scout) does.\par
\noindent\textit{Uses:} every one of the 64 kernel operations.\par
\noindent\textit{When useful:} same as scout but you want the readable report at the REPL\par

The eyeball version of scout: same search, same return value, but it also prints the numbered CLOSING branches (or, when none close, the best partials with how many goals each leaves open).  Use scout-show interactively, (scout) when you want to consume the result as data.

\par\smallskip{\raggedright\noindent\texttt{(scout-run k)}\quad\tactkind{composite}\index{scout-run@\texttt{scout-run} (tactic)}\par}\nobreak
\noindent Adopt closing branch k from the last (scout)/(scout-show) onto the live proof, running its tactics for real (they record and display normally).\par
\noindent\textit{Uses:} no kernel operation: it records no inference.\par
\noindent\textit{When useful:} scout found a closing branch you want to adopt onto the live proof\par

After scout finds CLOSING branches [1], [2], ..., (scout-run k) replays branch k's tactic forms through the real tactics on your live *ps*, from the same focus scout cloned -- so the proof advances and the steps are recorded for the script / PDF exactly as if you had typed them.  Works after either (scout) or (scout-show); both stash the closing branches.

\par\smallskip{\raggedright\noindent\texttt{(prep 'ineq)}\index{prep@\texttt{prep} (tactic)}\par}\nobreak
\noindent Why will a tactic not fire here, and what must be done first?  Runs the tactic's OWN preconditions one at a time, marks each ok / PREP / STOP, names the library lemma that repairs each unmet one, and prints a PLAN it has checked by running it on a scratch clone.  READ-ONLY: the live proof is never touched.  (prep 'ineq) is the only method so far.\par
\noindent\textit{Uses:} no kernel operation: it records no inference.\par

A tactic reports failure as a boolean, so every unmet precondition comes out as the same \#f and the same warning -- (ineq) says `goal not a linear-RR consequence of the named assumptions' whether your goal merely needs a di, or an atom needs a typing fact the library already proves, or the goal is simply false.  prep runs the same predicates separately and tells you WHICH failed.  For ineq the obligations are: the goal must BE an order relation (repair: di, counted by measuring, since one di consumes a typed FORALL and its guard together); every atom must be certified in RR by a LITERAL (IN t RR) scan -- (IN k NN) does not count, the oracle does no subtype reasoning (repair: a coercion lemma, e.g. nn-in-rr); some assumption must be order-shaped, since a fact like not(k=0) is invisible to Farkas (repair: a bridge lemma, e.g. nn-pos-of-nonzero); and the goal must actually follow, which only Fourier-Motzkin decides -- a STOP there means no prepping will help, which is the useful answer.  The repair search is by SHAPE over the whole library and is verified by SPECULATION: a candidate counts only if ineq then fires, not merely because it landed something order-shaped.  Every repair that works is reported, not just the first -- they come out alphabetical, which is no order of merit, and if the goal is itself a library theorem then citing IT is one of the closers.  Worked case: |- forall k in NN. \textasciitilde{}(k=0) => k < 2k.  prep returns (di) (di) (fact 'nn-in-rr 'k) (fact 'nn-pos-of-nonzero 'k) (ineq 4) -- six steps, where the hand-built cut/crs route through k = k+0 < k+k = 2k took thirteen and billed two extra transitivity lemmas.  A tactic joins the table with (prep-method! 'FUBA proc); crs and ass are the obvious next two.

\par\smallskip{\raggedright\noindent\texttt{(subst '(= s t))  |  (subst '(== s t))}\quad\tactkind{rule}\index{subst@\texttt{subst} (tactic)}\par}\nobreak
\noindent Rewrite s -> t throughout the goal, using an equation s = t -- or a quasi-equation s == t -- that is in context, in EITHER orientation (Leibniz substitution).  Reaches every position, operator (head) position included; a binder whose variable the equation mentions is not entered.\par
\noindent\textit{Uses:} \texttt{eq-subst}.\par
\noindent\textit{When useful:} you have an equation s = t (or a quasi-equation s == t) in context and want to rewrite the goal by it -- wherever the term occurs, operator position included\par

Use an equation among your hypotheses to replace s by t everywhere in the goal.  BOTH equalities license it: `=' is VNB's partial equality, and `s = t' already entails s = s and t = t, so t is defined wherever s was and the rewrite never replaces a defined term by an undefined one; `==' is quasi-equality (same definedness, equal where defined), which is a congruence and so substitutes exactly as `=' does -- which you need, since the partial-op recursion and bridge facts are stated with `=='.  The head you pass need not match the head in context, and neither need the orientation: all four of (= s t), (= t s), (== s t), (== t s) are searched for, and the goal is always rewritten s -> t.  LIMIT: the rewrite walk reaches argument positions only, so it is a silent no-op on a term in OPERATOR position -- (subst '(= (VADD md) ...)) will not touch the goal ((VADD md) x y).  Structure accessors are almost always in operator position; use the equation as a macete there (mac on the goal, mac-h on a hypothesis).  (Technically: Leibniz substitution from an in-context equality or quasi-equality.)

\par\smallskip{\raggedright\noindent\texttt{(beta)}\quad\tactkind{rule}\index{beta@\texttt{beta} (tactic)}\par}\nobreak
\noindent Beta-reduce a functoid application in the goal.\par
\noindent\textit{Uses:} \texttt{functoid-beta}.\par
\noindent\textit{When useful:} the goal has a functoid applied to an argument\par

Simplify a function-expression applied to an argument by substituting the argument into its body.  (Technically: beta-reduction of a functoid application in the goal.)

\par\smallskip{\raggedright\noindent\texttt{(nth-r)}\quad\tactkind{rule}\index{nth-r@\texttt{nth-r} (tactic)}\par}\nobreak
\noindent Reduce an NTH applied to a literal LIST in the goal.\par
\noindent\textit{Uses:} \texttt{nth-reduce}.\par
\noindent\textit{When useful:} the goal has NTH of a literal LIST\par

Simplify `the k-th entry of an explicit list [a1, a2, ...]' to that entry.  (Technically: reduce NTH applied to a literal LIST.)

\par\smallskip{\raggedright\noindent\texttt{(len-r)}\quad\tactkind{rule}\index{len-r@\texttt{len-r} (tactic)}\par}\nobreak
\noindent Reduce a LENGTH applied to a literal LIST in the goal.\par
\noindent\textit{Uses:} \texttt{length-reduce}.\par
\noindent\textit{When useful:} the goal has LENGTH of a literal LIST\par

Simplify `the length of an explicit list [a1, ..., an]' to the count n.  Structural (the dual of nth-r): the spine of a list literal is total, so its length is n regardless of whether the entries are defined.

\par\smallskip{\raggedright\noindent\texttt{(rfl)}\quad\tactkind{rule}\index{rfl@\texttt{rfl} (tactic)}\par}\nobreak
\noindent Close a reflexive equality goal (t = t).\par
\noindent\textit{Uses:} \texttt{reflexivity}.\par
\noindent\textit{When useful:} the goal is `t = t' with t manifestly defined\par

Close a goal `t = t' -- a thing equals itself.  (Technically: reflexivity of equality.)

\par\smallskip{\raggedright\noindent\texttt{(qrfl)}\quad\tactkind{rule}\index{qrfl@\texttt{qrfl} (tactic)}\par}\nobreak
\noindent Quasi-reflexivity: close t = t under the partial-equality definedness reading.\par
\noindent\textit{Uses:} \texttt{quasi-reflexivity}.\par
\noindent\textit{When useful:} the goal is `t = t' but t might be UNDEFINED (partial =)\par

Close `t = t' under the partial-equality reading, where asserting t = t also asserts that t is DEFINED.  Use this rather than rfl when t might be undefined.  (Technically: quasi-reflexivity; see the partial-equality convention, where `t = t' is the definedness predicate.)


\subsection*{Propositional logic}

\par\smallskip{\raggedright\noindent\texttt{(prop)}\index{prop@\texttt{prop} (tactic)}\par}\nobreak
\noindent Decide the goal by PROPOSITIONAL logic from the context, and close it if it follows.  Every non-connective formula -- `x in a', an equation, a whole `forall(...)' -- is one opaque atom; AND/OR/NOT/IMPLIES/IFF are read.  On a goal that does NOT follow it prints the countermodel (which atom must be true, which false) and leaves the proof untouched.  Adds no trust: it decides semantically, then discharges through di/ai/oi/ass/use-em/have!/detach!, so the bill is unchanged and the recorded script is the ordinary step-by-step proof.\par
\noindent\textit{Uses:} \texttt{and-elim} \texttt{and-intro} \texttt{arith-forsome} \texttt{arith-ground} \texttt{arith-simplify} \texttt{assumption} \texttt{cartesian-decompose} \texttt{cut} \texttt{detach} \texttt{eq-subst} \texttt{forall-elim} \texttt{forall-intro} \texttt{forsome-elim} \texttt{iff-elim} \texttt{iff-intro} \texttt{implies-intro} \texttt{macete} \texttt{not-elim} \texttt{not-intro} \texttt{or-elim} \texttt{or-intro-left} \texttt{or-intro-right} \texttt{proof-by-contradiction} \texttt{theorem-assumption} \texttt{tuple-equality-decompose}.\par
\noindent\textit{When useful:} the goal follows from the hypotheses by AND/OR/NOT/IMPLIES/IFF alone -- no quantifier or equality reasoning needed\par

Close a goal that follows from the hypotheses by pure propositional reasoning -- the leaves where you can SEE the answer and still have to pick between (oi-l)(ass), (oi-r) plus a conjunction split, and (ai) on a negation.  It treats each atomic statement as a black box, so it will not instantiate a quantifier or reason about equality: if the goal needs an instance, land the instance first (fact / inst+) and run it again.  When it declines it names a falsifying assignment, which usually tells you exactly which hypothesis is missing.  (Technically: three-valued evaluation over the atoms decides the entailment; the proof is then replayed through the kernel rules, case-splitting with excluded middle where a disjunctive goal needs it.)


\subsection*{Arithmetic \& ring oracles}

\par\smallskip{\raggedright\noindent\texttt{(arith)}\quad\tactkind{oracle}\index{arith@\texttt{arith} (tactic)}\par}\nobreak
\noindent Discharge a ground arithmetic goal by evaluation.\par
\noindent\textit{Uses:} \texttt{arith-forsome} \texttt{arith-ground} \texttt{arith-simplify} \texttt{cartesian-decompose} \texttt{eq-subst} \texttt{macete} \texttt{tuple-equality-decompose}.\par
\noindent\textit{When useful:} the goal is ground arithmetic -- no variables\par

Close a goal that is a concrete numerical fact with no variables -- e.g. 2 + 3 = 5, or 7 in NN -- by just computing it.  (Technically: decision by ground arithmetic evaluation.)

\par\smallskip{\raggedright\noindent\texttt{(rs)}\quad\tactkind{oracle}\index{rs@\texttt{rs} (tactic)}\par}\nobreak
\noindent Ring-simplify the goal (normal form over the ambient ring).\par
\noindent\textit{Uses:} \texttt{ring-simplify}.\par
\noindent\textit{When useful:} the goal is a (possibly non-commutative) ring-word identity\par

Simplify the goal to a normal form in the ambient ring (which need not be commutative).  The non-commutative companion of crs.  (Technically: ring-simplify to a canonical word form.)

\par\smallskip{\raggedright\noindent\texttt{(crs)}\quad\tactkind{oracle}\index{crs@\texttt{crs} (tactic)}\par}\nobreak
\noindent Commutative-ring decision procedure: prove a polynomial identity over ZZ[generators].  Expands literal powers, so (x+y)\textasciicircum{}2 = ... closes directly.\par
\noindent\textit{Uses:} \texttt{comm-ring-simplify}.\par
\noindent\textit{When useful:} the goal is a commutative-ring identity, e.g. (x+y)\textasciicircum{}2 = x\textasciicircum{}2+2xy+y\textasciicircum{}2\par

Prove a polynomial identity that holds in EVERY commutative ring -- e.g. (x+y)\textasciicircum{}2 = x\textasciicircum{}2 + 2xy + y\textasciicircum{}2 -- by reducing both sides to a canonical sum-of-monomials form and checking they agree.  A genuine decision procedure: a true commutative-ring identity closes, a non-identity is refused.  Literal powers are expanded for you.  (Technically: normal form over the free commutative ring ZZ[generators].)

\par\smallskip{\raggedright\noindent\texttt{(simp [target])}\quad\tactkind{oracle}\index{simp@\texttt{simp} (tactic)}\par}\nobreak
\noindent Rewrite a commutative-ring SUBTERM of the goal to canonical form, IN PLACE (e.g. (x+y)\textasciicircum{}2 inside a larger goal becomes x\textasciicircum{}2 + 2*x*y + y\textasciicircum{}2).  Works on BOTH surfaces: concrete number domains (+ * - \textasciicircum{} over NN/ZZ/QQ/RR/CC) and a generic ring s ((ADD s)/(MUL s)/(NEG s), carrier (CARR s)).  No arg = outermost ring subterm; ``term'' targets a specific one.  Sound by cut + crs + eq-subst (no new kernel rule); needs the subterm's generators typed in context (true post-di), else refuses and names them.\par
\noindent\textit{Uses:} \texttt{comm-ring-simplify} \texttt{cut} \texttt{eq-subst}.\par
\noindent\textit{When useful:} a ring SUBTERM of a larger goal should be put in normal form in place\par
\par\smallskip{\raggedright\noindent\texttt{(type-term)}\quad\tactkind{composite}\index{type-term@\texttt{type-term} (tactic)}\par}\nobreak
\noindent Close a typing goal (IN t C).  In a number class (NN ZZ QQ RR CC) t may be built from context-typed atoms (or atoms typed in a subclass: NN in ZZ, ZZ in RR, ...), numerals, + * -, and powers, typed by the class's closure laws, bottom-up.  In ANY class C, t may be an atom typed in a subclass, or an application f(a) with f in fun(A, C) in context and a typable in A, recursively: f(3 * x + 1) in B from f in fun(zz, B) and x in zz.  Declines -- with a warning, and nothing changed -- when there is no plan.\par
\noindent\textit{Uses:} \texttt{assumption} \texttt{comm-ring-simplify} \texttt{cut} \texttt{detach} \texttt{eq-subst} \texttt{forall-elim} \texttt{theorem-assumption}.\par
\noindent\textit{When useful:} the goal types a term: a polynomial in context-typed atoms in NN/ZZ/QQ/RR/CC, e.g. 3*x\textasciicircum{}2 + 3*x*y + y\textasciicircum{}2 in ZZ, or f(a) in any class C with f in fun(A, C) in context\par

The typing leaf that arithmetic leaves behind: `3 * x\textasciicircum{}2 + 3 * (x * y) + y\textasciicircum{}2 in ZZ' with x and y integers.  It plans the whole typing first -- which closure law types each subterm, or which subclass inclusion types an atom -- and only then runs it, so a term it cannot type leaves the proof untouched.  Each subterm's typing is proved on its own lane (cut) and lands in the context once.  A power x\textasciicircum{}k is typed by the power law in RR and CC; in ZZ, QQ and NN, where the library has no power law, by the identity x\textasciicircum{}k = x * ... * x (proved by crs) and the multiplication law, so the exponent must be a numeral there.  Subtraction in QQ goes through a - b = a + (-b) the same way; NN has no minus and declines.  Division and recip are not handled.  The kit twin (dk-type! t C) lands the typing in the context and returns it; in-rr hands its power goals to the same planner.  (Technically: composite -- cut, theorem-assumption, forall-elim, detach, assumption, eq-subst, and comm-ring-simplify for the power identities; no rule of its own.  Transactional: a run that does not close the goal is rolled back, script and undo stack included.)

\par\smallskip{\raggedright\noindent\texttt{(ew-poly)}\quad\tactkind{composite}\index{ew-poly@\texttt{ew-poly} (tactic)}\par}\nobreak
\noindent Close an existential goal forsome([a in C], L = R) (or unguarded) in which a occurs once, linearly: compute the witness by EXACT POLYNOMIAL DIVISION, supply it (ew), type it (type-term) and close the identity (crs).  One command for `(x+y)\textasciicircum{}3 = x\textasciicircum{}3 + y*a'.  Declines, with the reason and nothing changed, when the division leaves a remainder.\par
\noindent\textit{Uses:} \texttt{and-intro} \texttt{comm-ring-simplify} \texttt{forsome-intro}.\par
\noindent\textit{When useful:} the goal is forsome(a, L = R) with a occurring once, linearly, and the witness is a quotient of polynomials, e.g. (x+y)\textasciicircum{}3 = x\textasciicircum{}3 + y*a\par

The `find a such that ...' step when the answer is a quotient: for (x + y)\textasciicircum{}3 = x\textasciicircum{}3 + y * a the witness is ((x+y)\textasciicircum{}3 - x\textasciicircum{}3) / y = 3x\textasciicircum{}2 + 3xy + y\textasciicircum{}2.  Both sides are put in the sum-of-monomials normal form crs uses (integer coefficients over the generators, the witness variable among them); the equation must be of degree one in the witness variable, L - R = A*a + B, and the witness is -B/A by exact division of polynomials.  It declines, saying why, when a occurs other than exactly once, a side is not a polynomial, the equation is not linear in a, the coefficient A is 0, the division leaves a remainder, or the quotient has a fractional coefficient in ZZ or NN.  Then the recorded steps are ordinary ones: ew at the witness (written with + - * \textasciicircum{}, highest degree first), di to split the typing from the identity, type-term on the typing, crs on the identity.  what-now proposes it on such a goal, and scout's ew lane tries the same witness.  (Technically: composite -- forsome-intro, and-intro, the type-term steps, comm-ring-simplify; no rule of its own.  Everything is checked before anything runs, and the run is transactional besides.)

\par\smallskip{\raggedright\noindent\texttt{(expand)}\quad\tactkind{composite}\index{expand@\texttt{expand} (tactic)}\par}\nobreak
\noindent SHOW the expansion as a step: rewrite each polynomial side of an equation in the goal (also inside an existential or a conjunction) to its sum-of-monomials normal form, highest degree first.  (x + y)\textasciicircum{}3 = x\textasciicircum{}3 + y*a becomes x\textasciicircum{}3 + 3*x\textasciicircum{}2*y + 3*x*y\textasciicircum{}2 + y\textasciicircum{}3 = x\textasciicircum{}3 + y*a.  Declines, with the reason and nothing changed, when no side changes.\par
\noindent\textit{Uses:} \texttt{comm-ring-simplify} \texttt{cut} \texttt{eq-subst}.\par
\noindent\textit{When useful:} a side of an equation in the goal is an unexpanded polynomial and you want to SEE its expansion, e.g. (x+y)\textasciicircum{}3\par

The step a hand proof writes as `expanding, ...': the goal after it says what the equation says, term by term.  Each side t is proved equal to its normal form nf(t) by crs on a lane, and the goal is rewritten by subst, so the page shows the equation and the rewrite and both are checked.  A side is left alone when it mentions the variable an existential binds (that is not a term of the context -- the witness is what the goal asks for), when an atom of it is not typed in a number class (crs could not check it), or when it is already in normal form.  what-now offers it first on such a goal, and names the binomial theorem (binomial-theorem, n = k) when a power of a sum is expanded -- cited as what the expansion instantiates, not fired: the library states it over an abstract commutative ring as a finite sum.  (Technically: composite -- cut, comm-ring-simplify, eq-subst; no rule of its own; transactional.)

\par\smallskip{\raggedright\noindent\texttt{(zero-it)}\quad\tactkind{composite}\index{zero-it@\texttt{zero-it} (tactic)}\par}\nobreak
\noindent Move an equation goal P = Q (or P == Q, bare or under typed universals) to R = 0, where R is the normal form of P - Q (crs's calculator, highest degree first); prints `P - Q normalises to R' and returns R.  An inequality goal P <= Q or P < Q (the order of RR; atoms typed in NN, ZZ or QQ are lifted) is moved to 0 <= R or 0 < R, R the normal form of Q - P, the branch closed through rr-le-from-diff-nonneg / rr-lt-from-diff-pos; a numeral R decides it (0 <= -1: FALSE, nothing changed; 0 < 2: closed).  CC and an abstract ring have no order and are declined.  R = 0: closed by crs.  R a non-zero numeral: says the goal is FALSE and changes nothing.  Otherwise ONE recorded step leaves the goal R = 0, the atoms' typings landed or listed as owed side leaves.\par
\noindent\textit{Uses:} \texttt{comm-ring-simplify} \texttt{cut} \texttt{eq-subst}.\par
\noindent\textit{When useful:} an equation P = Q over a ring, or an inequality P <= Q / P < Q over the reals, and you want to see (and prove) what it reduces to: R = 0 (0 <= R), or a FALSE numeral that marks a dead path\par

The user's notes-42: `if in the end it means I have to prove 1 = 2, I can see I've been barking up the wrong tree'.  The ring is read off the typings of the atoms: the number surface (+ - * \textasciicircum{}, numerals; atoms typed in NN ZZ QQ RR CC, the largest class wins, NN lifting) or an abstract ring a with IS-COMMUTATIVE-RING(a) in context ((ADD a) (MUL a) (NEG a) (ZERO a) (ONE a)).  NN alone is refused for the rewrite (no subtraction).  Each atom of P and Q must be typed in the ring: an untyped one is typed by type-term's plan or, for recip(t) with not(t = 0) in context, by the recip-closure law; what cannot be typed is CUT as an owed side leaf `g in D', listed in the printout, and the new goal carries it as a hypothesis -- the notes' `[the eliminated terms all defined]'.  Then R = 0 is cut, and the branch R = 0 |- P = Q is closed by crs (P = Q + R), subst (P -> Q + R, then R -> 0) and crs; no sub-zero law and no typing of P or Q is needed.  Universals typed in the ring are peeled first (di, inside the same step).  The step records as (zero-it) with no argument: replay recomputes the normal form from the goal.  what-now shows the normal form on an equation goal and offers (zero-it).  (Technically: composite -- cut, comm-ring-simplify, eq-subst, plus di / theorem-assumption / forall-elim for the peeling and typings; the inner steps run with recording suppressed on the bc*-handler path, their citations kept for the bill; transactional.)

\par\smallskip{\raggedright\noindent\texttt{(supply)}\index{supply@\texttt{supply} (tactic)}\par}\nobreak
\noindent Close an inequality goal the way a hand proof would: beta-reduce any applied lambda, land the typing certificates and the standard bounds the linear oracle cannot see -- including the TRIANGLE inequality at the summands of an abs of a sum -- then call ineq over every premise.  Rehearses the whole sequence on a throwaway copy and does NOTHING unless it closes; on a miss it prints the sequence it tried.  Adds no trust: it drives lam-b, fact and ineq, each of which records itself.\par
\noindent\textit{Uses:} every one of the 64 kernel operations.\par

The committing form of what-now's SUPPLY THE ORACLE lane.  `ineq' reads abs(...), max(...) and (dist(s))(x,y) as opaque atoms and refuses any premise whose atoms are not certified real, so an inequality that is obviously true can fail for want of a typing nobody mentioned.  This lands them.  The one real idea is subadditivity: where the goal bounds abs(u + v), the useful intermediate term is abs(u) + abs(v), and the decomposition is in the term itself -- nothing is searched for.  Because a lambda reduction cannot be undone and can owe an unprovable leaf, the whole sequence is rehearsed before any of it is run.  (Technically: certificate closure to depth 2 over the forward-citation lane, plus the curated bound table, then Fourier-Motzkin.)

\par\smallskip{\raggedright\noindent\texttt{(ineq i1 i2 ...)}\quad\tactkind{oracle}\index{ineq@\texttt{ineq} (tactic)}\par}\nobreak
\noindent Close a linear-inequality goal over RR (<= < > >= = between RR terms) as a consequence of the named assumptions (1-based indices), via the Fourier-Motzkin/Farkas oracle.  Linearizes over + - * and the binplus/binneg/bintimes aliases; every MAXIMAL non-arithmetic subterm is an atom that must be certified in RR.  (Does NOT see through a generic ring's (ADD s)/(MUL s) -- those become opaque atoms.)\par
\noindent\textit{Uses:} \texttt{ineq}.\par
\noindent\textit{When useful:} a LINEAR <= / < over RR that follows from inequalities you can cite\par

Close a LINEAR inequality over the reals that follows from inequalities you cite (by their hypothesis numbers) -- by chaining them, adding them, and scaling by positive constants.  Anything that is not built from + - and multiplication-by-constants is treated as an opaque quantity, so it handles e.g. the triangle inequality where the distances are unknowns.  For NONLINEAR (polynomial) inequalities use sos instead.  (Technically: Fourier-Motzkin / Farkas over the ordered field RR; every atom must be certified real.)

\par\smallskip{\raggedright\noindent\texttt{(sos ``c1'' ``c2'' ...)}\quad\tactkind{oracle}\index{sos@\texttt{sos} (tactic)}\par}\nobreak
\noindent Sum-of-squares closer for a nonstrict polynomial inequality a <= b over RR (the nonlinear companion of (ineq)).  You supply the terms to be SQUARED; it finds the nonnegative coefficients.  (tactics 'sos) for the worked example.\par
\noindent\textit{Uses:} \texttt{sos}.\par
\noindent\textit{When useful:} a NONSTRICT POLYNOMIAL <= over RR (the nonlinear cousin of ineq)\par

A sum-of-squares closer.  To prove  a <= b  it is enough to exhibit b - a  as a sum of squares (each obviously >= 0), reducing the inequality to an algebraic identity.  You hand sos the terms to be SQUARED -- the c\_i, NOT the squares -- and it solves for nonnegative rationals lambda\_i with

\begin{VNBsnip}
      b - a  =  lambda_1 c_1^2 + ... + lambda_n c_n^2
\end{VNBsnip}

(matched coefficient-by-coefficient in crs's commutative-ring normal form), then PRINTS the lambda\_i it found.  You do not supply them.

Worked example.  Goal  forall([x in rr, y in rr], x*y <= x\textasciicircum{}2 + y\textasciicircum{}2).  Here b - a = x\textasciicircum{}2 - x*y + y\textasciicircum{}2 = 1/2(x-y)\textasciicircum{}2 + 1/2 x\textasciicircum{}2 + 1/2 y\textasciicircum{}2,  so the things squared are  x-y, x, y:

\begin{VNBsnip}
      (sos "x - y" "x" "y")
          ;; closes, printing  1/2(x-y)^2 + 1/2 x^2 + 1/2 y^2
\end{VNBsnip}

2*x*y <= x\textasciicircum{}2 + y\textasciicircum{}2 needs only one square:  (sos ``x - y'').  The three-variable a*b+b*c+c*a <= a\textasciicircum{}2+b\textasciicircum{}2+c\textasciicircum{}2  wants  (sos ``a - b'' ``b - c'' ``c - a'').

Notes. - Write the c\_i with the goal's own variable names.  di preserves them, so sos works before OR after di; skip di and sos peels the typed forall([... in rr]) wrapper itself. - Every generator (variable) must be certified in RR -- a goal binder x in rr  or an  (IN g RR)  assumption. - A wrong or insufficient certificate refuses cleanly: nothing is closed, you simply try other squares. - Strict (<) goals are refused -- a square can be 0, so squares alone never force a strict inequality.


\subsection*{Chained reasoning}

\par\smallskip{\raggedright\noindent\texttt{(calc L0 (rel1 L1 [just1]) (rel2 L2 [just2]) ...)}\quad\tactkind{composite}\index{calc@\texttt{calc} (tactic)}\par}\nobreak
\noindent Ground the focus goal (REL L0 Ln) by a CHAIN of intermediaries L0 rel1 L1 rel2 L2 ... reln Ln: proves each link and composes them into the endpoint.  A link no lane can close is LEFT OPEN as a leaf -- the refinement point, and the PSS candidate.  Relations: = == (folded by transitivity), < <= (folded through the co-*-trans lemmas, with = steps riding along), IFF (chained per direction).  RETURNS the links, the open ones, and their PSS candidates as an alist.\par
\noindent\textit{Uses:} every one of the 64 kernel operations.\par
\noindent\textit{When useful:} the goal is (REL L0 Ln) and you can WRITE the chain of intermediaries that gets there -- or you want the one link that will not close isolated as a PSS candidate\par

Write the argument the way you would on paper -- as a chain through intermediate terms -- and let the machine prove each link:

\begin{VNBsnip}
      Goal  k < 2*k   (with k in NN, ~(k=0))
      (calc 'k '(= (+ k 0)) '(< (+ k k)) '(= (* 2 k)))
          ;; i.e.  k = k+0 < k+k = 2*k
\end{VNBsnip}

Each step names a RELATION and the next LINE; calc cuts the link (rel L(i-1) Li), dispatches a lane at it, and on success composes the whole chain into the goal.  It is pure bookkeeping over the trusted tactics -- no kernel rule, no axiom of its own -- so a calc proof bills exactly the lemmas its lanes and composers cite.

The COMPOSER is forced by the relations you used, and there are three families: `cong' for = and ==, folded by eq-trans; `order' for < and <=, folded through the co-lt-trans / co-le-lt-trans / co-eq-lt-trans ... lemmas (a = step rewrites an endpoint of the running relation and rides along, which is why the chain above needs no separate arithmetic); and `iff', which di's the goal into its two directions and chains each with ai/detach!, because VNB cannot quantify over propositions and so has NO first-order iff-trans lemma to fold with.

The JUSTIFICATION of a link is optional and defaults to 'auto -- try the relation's lanes: crs then arith for = / ==; for < / <= the NN->RR bridge then ineq over the order-shaped premises; grind for IFF.  Otherwise pass 'scout (a small scout search), a macete/theorem NAME (mac it, then close), an explicit tactic form like '(fact 'nn-pos-of-nonzero 'k) eval'd at the link's focus, or 'open / \#f to leave the link open ON PURPOSE.

Sanity first: calc ERRORS before touching the proof if the goal's head is not the composed relation, or its LHS is not L0, or its RHS is not Ln.  A link already in the main context is taken as a GIVEN and skipped -- cutting it would be an alpha self-loop that opens no leaf.

What it is FOR is the open links.  A chain whose links are all closed is a proof; a chain with one open link has isolated your obstacle to a single formula, printed with its free variables and their context typing -- which is the PSS candidate to state as a lemma.  Refine by inserting more intermediaries until each link is discoverable.  (See also (prep 'ineq), which answers the other question: why a lane will not fire.)

\par\smallskip{\raggedright\noindent\texttt{(have! CLAIM [THUNK])}\quad\tactkind{composite}\index{have"!@\texttt{have"!} (tactic)}\par}\nobreak
\noindent Assert CLAIM as an intermediate step and carry on with it as a hypothesis: cut CLAIM, discharge the side goal from context (or by THUNK), and stay on the main branch.\par
\noindent\textit{Uses:} \texttt{and-intro} \texttt{arith-forsome} \texttt{arith-ground} \texttt{arith-simplify} \texttt{assumption} \texttt{cartesian-decompose} \texttt{cut} \texttt{detach} \texttt{eq-subst} \texttt{forall-elim} \texttt{forall-intro} \texttt{iff-intro} \texttt{implies-intro} \texttt{macete} \texttt{not-intro} \texttt{theorem-assumption} \texttt{tuple-equality-decompose}.\par
\noindent\textit{When useful:} you want to state an intermediate fact that follows immediately from context and continue with it -- the `we have X' step\par

The `we have X' step.  You state an intermediate fact CLAIM that follows immediately from what is already known; have! cuts it, proves the resulting side goal automatically, and returns you to the main branch with CLAIM now available as a hypothesis.  The automatic discharge (`from-context!') closes the everyday cases -- a conjunction, splitwise; an (IN (a*b) NN) by nn-mul-closed on the factors; a numeral membership by arith; otherwise the goal is already a context assumption up to alpha.  When the side goal needs more than that, pass a THUNK -- any tactic sequence -- as the second argument to discharge it your way.  It ERRORS, never silently no-ops, if CLAIM is already in context up to alpha: the cut would self-loop, opening one child and no main branch.  So a script reads as the argument does -- `we have q0*q0 = 3*(k*k); we have k =/= 0; ...'.  (Technically: composite -- cut then from-context! or THUNK; adds no kernel rule.)


\subsection*{Backchaining with a theorem}

\par\smallskip{\raggedright\noindent\texttt{(ta 'name)}\quad\tactkind{rule}\index{ta@\texttt{ta} (tactic)}\par}\nobreak
\noindent Theorem-assumption: bring the named installed theorem into context as an assumption.\par
\noindent\textit{Uses:} \texttt{theorem-assumption}.\par
\noindent\textit{When useful:} you want a library theorem available as a hypothesis\par

Bring an already-proved theorem into your current hypotheses so you can use it (instantiate it, detach from it, ...).  (Technically: adds the named installed theorem as an assumption.)

\par\smallskip{\raggedright\noindent\texttt{(wbc ['name])}\quad\tactkind{composite}\index{wbc@\texttt{wbc} (tactic)}\par}\nobreak
\noindent Witness-shape backchain: on an `exists v. ...' goal, cite a PSS lemma that MANUFACTURES a witness of that shape, so you can finish with inst+/grind/ew.\par
\noindent\textit{Uses:} \texttt{and-elim} \texttt{arith-forsome} \texttt{arith-ground} \texttt{arith-simplify} \texttt{cut} \texttt{detach} \texttt{forall-elim} \texttt{forsome-elim} \texttt{iff-elim} \texttt{not-elim} \texttt{or-elim} \texttt{theorem-assumption}.\par
\noindent\textit{When useful:} an EXISTENCE goal whose witness a PSS lemma manufactures (diagonalization / block-family)\par

The discovery move for existence proofs.  Faced with `there exists phi with ...', it looks up the produces-witness index for lemmas whose conclusion builds the RIGHT KIND of object -- the same existential-witness shape -- even when the rest of the statement differs (which is why plain bc* misses them: bc* demands the WHOLE conclusion unify).  E.g. on `exists phi. STRICTLY-MONO-NN(phi) and IS-CAUCHY-SEQ(SUBSEQ f phi)' it reaches for `diagonalization' (which makes a STRICTLY-MONO-NN phi from a nested family); on a nested-family goal it reaches for `block-family'.  With no argument it cites the top retrieved producer; (wbc 'name) cites the one you name.  It brings the lemma in (via ta); you then discharge its hypotheses (inst+), skolemize its existential (grind), and supply its witness to your goal (ew).  Circular `headline' lemmas that already conclude your exact goal are filtered out.  (See (witness-producers goal) for the raw retrieval; what-now lists the (wbc 'name) candidates on any existential goal.)

\par\smallskip{\raggedright\noindent\texttt{(bc impl)}\quad\tactkind{rule}\index{bc@\texttt{bc} (tactic)}\par}\nobreak
\noindent Backchain the goal through an (IMPLIES A B) already in context: if the goal matches B, the new goal is A.\par
\noindent\textit{Uses:} \texttt{backchain}.\par
\noindent\textit{When useful:} you have (P => Q) in context and the goal is Q\par

Work backwards through an implication you already have.  If `P implies Q' is among your hypotheses and your goal is Q, this reduces the goal to proving P.  The everyday `to get Q it's enough to show P'.  (Technically: backchaining the goal through an in-context implication whose conclusion matches.)

\par\smallskip{\raggedright\noindent\texttt{(bc* 'thm [((v val)...)] h1 h2 ...)}\quad\tactkind{composite}\index{bc*@\texttt{bc*} (tactic)}\par}\nobreak
\noindent Matching backchain: unify the theorem's conclusion against the goal, then leave its (instantiated) antecedents as subgoals -- optional handlers hk run on the k-th subgoal.\par
\noindent\textit{Uses:} \texttt{assumption} \texttt{backchain} \texttt{cut} \texttt{forall-elim} \texttt{theorem-assumption}.\par
\noindent\textit{When useful:} a library theorem's CONCLUSION matches your goal\par

Apply a known theorem to your goal.  If you have a theorem `if A and B then C' and your goal is its conclusion C -- matching the theorem's variables to yours, and the names of any dummy/bound variables don't matter (`there exists a net N' applies to a goal `there exists a net F') -- this replaces `prove the goal' with `prove A' and `prove B', the theorem's hypotheses.  The everyday `by Theorem X it suffices to show A and B'.  You may attach a handler to each hypothesis to dispatch it; values the match can't determine you supply as ((v val) ...).  (Technically: peels the theorem's leading universals and implications, matches the conclusion -- alpha-aware on bound variables -- and replays ta/inst/cut/bc automatically.)


\subsection*{Lambda, comprehension \& description}

\par\smallskip{\raggedright\noindent\texttt{(lam-t)}\quad\tactkind{rule}\index{lam-t@\texttt{lam-t} (tactic)}\par}\nobreak
\noindent VNB-LAMBDA typing: reduce (IN (VNB-LAMBDA ...) (FUN A B)) to TWO obligations -- the body, and that the domain is a set.\par
\noindent\textit{Uses:} \texttt{lambda-type}.\par
\noindent\textit{When useful:} the goal types a VNB-LAMBDA into FUN(A,B)\par

Show that a function defined by a formula (`x |-> ...') maps A into B -- reduces to showing that, for an arbitrary input in A, the value lies in B, AND that A is a set.  The second is not a formality: carrying the domain in the term stops one lambda from typing into FUN(A,B) for every A, but says nothing about A being a set, and a lambda over a proper class is not a function.  The lambda's declared domain must also be the FUN's, or the rule does not apply at all.  (Technically: VNB-LAMBDA typing into FUN A B.)

\par\smallskip{\raggedright\noindent\texttt{(lam-b)}\quad\tactkind{rule}\index{lam-b@\texttt{lam-b} (tactic)}\par}\nobreak
\noindent VNB-LAMBDA beta: reduce an applied lambda to its substituted body.\par
\noindent\textit{Uses:} \texttt{lambda-beta}.\par
\noindent\textit{When useful:} the goal has a VNB-LAMBDA applied to an argument\par

Simplify a function `x |-> e(x)' applied to an argument a to e(a) -- the body with the argument substituted in.  A lambda carries its DOMAIN, and it is only defined on that domain, so the reduction is licensed only where a is known to be in it: if the membership is neither in the context nor supplied by an enclosing binder, the step still fires but leaves (IN a A) as an extra subgoal.  Land the typing fact BEFORE the lam-b and no subgoal appears.  (Technically: VNB-LAMBDA beta-reduction.)

\par\smallskip{\raggedright\noindent\texttt{(lam-b-h hyp)}\quad\tactkind{rule}\index{lam-b-h@\texttt{lam-b-h} (tactic)}\par}\nobreak
\noindent VNB-LAMBDA beta in a cited ASSUMPTION -- what mac-h is to mac.\par
\noindent\textit{Uses:} \texttt{lambda-beta-hyp}.\par

Beta-reduce an applied lambda inside a hypothesis, in place.  Needed because a `fact' that instantiates a theorem's function variable at a lambda lands the APPLIED lambda in the CONTEXT, where the goal-side `lam-b' cannot reach it: union-of-opens-open at the identity family g := x |-> x lands is-open(md, big-union(i, fam, (x |-> x)(i))).  Without this the proof must detour through a cut beta-equation and a subst.  Cites nothing, so it adds no debt.

\par\smallskip{\raggedright\noindent\texttt{(sep-set)}\quad\tactkind{rule}\index{sep-set@\texttt{sep-set} (tactic)}\par}\nobreak
\noindent Separation sethood: the separation set \{x in A | p\} is a set.\par
\noindent\textit{Uses:} \texttt{sep-sethood}.\par
\noindent\textit{When useful:} the goal asserts a separation set \{ x in A | p \} is a set\par

Show that a set-builder set \{x in A | p(x)\} is genuinely a set.  (Technically: separation sethood -- a subclass of a set is a set.)

\par\smallskip{\raggedright\noindent\texttt{(sep-mi)}\quad\tactkind{rule}\index{sep-mi@\texttt{sep-mi} (tactic)}\par}\nobreak
\noindent Separation membership intro: prove (IN t \{x in A | p\}).\par
\noindent\textit{Uses:} \texttt{sep-mem-intro}.\par
\noindent\textit{When useful:} the goal is membership in a separation set \{ x in A | p \}\par

Prove that a term t belongs to \{x in A | p(x)\} by showing t lies in A and satisfies the condition p.  (Technically: separation-membership introduction.)

\par\smallskip{\raggedright\noindent\texttt{(sep-me hyp)}\quad\tactkind{rule}\index{sep-me@\texttt{sep-me} (tactic)}\par}\nobreak
\noindent Separation membership elim: split a separation-membership assumption into A-membership and the predicate.\par
\noindent\textit{Uses:} \texttt{sep-mem-elim}.\par
\noindent\textit{When useful:} a hypothesis is separation-set membership -- unpack it\par

From a hypothesis that t belongs to \{x in A | p(x)\}, extract the two facts it packs: t lies in A, and p(t) holds.  (Technically: separation-membership elimination.)

\par\smallskip{\raggedright\noindent\texttt{(comp-mi)}\quad\tactkind{rule}\index{comp-mi@\texttt{comp-mi} (tactic)}\par}\nobreak
\noindent Comprehension membership intro.\par
\noindent\textit{Uses:} \texttt{comp-mem-intro}.\par
\noindent\textit{When useful:} the goal is comprehension-set membership\par

Prove that something belongs to a comprehension set by establishing the set's defining condition for it.  (Technically: comprehension-membership introduction.)

\par\smallskip{\raggedright\noindent\texttt{(comp-me hyp)}\quad\tactkind{rule}\index{comp-me@\texttt{comp-me} (tactic)}\par}\nobreak
\noindent Comprehension membership elim.\par
\noindent\textit{Uses:} \texttt{comp-mem-elim}.\par
\noindent\textit{When useful:} a hypothesis is comprehension-set membership\par

From a comprehension-membership hypothesis, extract its defining condition.  (Technically: comprehension-membership elimination.)

\par\smallskip{\raggedright\noindent\texttt{(iota-d term)}\quad\tactkind{rule}\index{iota-d@\texttt{iota-d} (tactic)}\par}\nobreak
\noindent Definite-description: discharge the IOTA uniqueness obligation for `term'.\par
\noindent\textit{Uses:} \texttt{iota-def}.\par
\noindent\textit{When useful:} the goal involves IOTA and you must discharge its uniqueness obligation\par

Justify `the unique x such that p(x)' (a definite description) by proving that exactly one such x exists.  (Technically: IOTA -- discharge the uniqueness obligation for the described term.)

\par\smallskip{\raggedright\noindent\texttt{(bu-set)}\quad\tactkind{rule}\index{bu-set@\texttt{bu-set} (tactic)}\par}\nobreak
\noindent Big-union sethood.\par
\noindent\textit{Uses:} \texttt{big-union-sethood}.\par
\noindent\textit{When useful:} the goal asserts a BIG-UNION is a set\par

Show that a big union (the union of an indexed family of sets) is itself a set.  (Technically: big-union sethood.)

\par\smallskip{\raggedright\noindent\texttt{(bu-mi w)}\quad\tactkind{rule}\index{bu-mi@\texttt{bu-mi} (tactic)}\par}\nobreak
\noindent Big-union membership intro via the index witness w.\par
\noindent\textit{Uses:} \texttt{big-union-mem-intro}.\par
\noindent\textit{When useful:} the goal is BIG-UNION membership -- you have the index witness\par

Prove that an element lies in a big union by exhibiting one index w whose set already contains it.  (Technically: big-union-membership introduction via the index witness.)

\par\smallskip{\raggedright\noindent\texttt{(bu-me hyp)}\quad\tactkind{rule}\index{bu-me@\texttt{bu-me} (tactic)}\par}\nobreak
\noindent Big-union membership elim.\par
\noindent\textit{Uses:} \texttt{big-union-mem-elim}.\par
\noindent\textit{When useful:} a hypothesis is BIG-UNION membership -- obtain its index\par

From a hypothesis that an element lies in a big union, obtain an index whose set contains it (a fresh name for that index).  (Technically: big-union-membership elimination.)


\subsection*{Conditional terms}

\par\smallskip{\raggedright\noindent\texttt{(if-true t)}\quad\tactkind{rule}\index{if-true@\texttt{if-true} (tactic)}\par}\nobreak
\noindent Reduce an (IF p a b) term on the p branch: spawns p as a subgoal; the continuation gains (= (IF p a b) a) when a is certified defined, else (== (IF p a b) a) (since 2026-10-08).\par
\noindent\textit{Uses:} \texttt{if-true}.\par
\noindent\textit{When useful:} the goal has an (IF p a b) term to evaluate on the p branch\par

Evaluate a conditional term `if p then a else b' on the assumption that p holds: you take on p as a side goal, and may then use that the conditional equals a.  (Technically: if-true reduction, spawning p.)

\par\smallskip{\raggedright\noindent\texttt{(if-false t)}\quad\tactkind{rule}\index{if-false@\texttt{if-false} (tactic)}\par}\nobreak
\noindent Reduce an (IF p a b) term on the not-p branch: spawns (NOT p); the continuation gains (= (IF p a b) b).\par
\noindent\textit{Uses:} \texttt{if-false}.\par
\noindent\textit{When useful:} the goal has an (IF p a b) term to evaluate on the not-p branch\par

Evaluate a conditional term `if p then a else b' on the assumption that p fails: you take on `not p' as a side goal, and may then use that the conditional equals b.  (Technically: if-false reduction, spawning not p.)


\subsection*{Navigation, display \& tacticals}

\par\smallskip{\raggedright\noindent\texttt{(focus n)}\index{focus@\texttt{focus} (tactic)}\par}\nobreak
\noindent Switch the focus to the n-th open goal (1-based, the POSITION in the list (show) prints: the focus first, then the others).  A position shifts as leaves close; the node number in brackets does not -- see focus-id.\par
\noindent\textit{Uses:} no kernel operation: it records no inference.\par
\par\smallskip{\raggedright\noindent\texttt{(focus-id k)}\index{focus-id@\texttt{focus-id} (tactic)}\par}\nobreak
\noindent Switch the focus to the ungrounded node whose NODE NUMBER is k, the number (show) and the Focus Workspace print in brackets before each sequent.  Any ungrounded node, not only a leaf: a driver records a leaf by position as (focus n) and an ungrounded node a rule has already fired on as (focus-id k), so a page may name either.  Recorded as a step, undone by backup-one.  The move to make when reading a goal off the panel, e.g. the side leaf of a cut: (cut L) opens  G |- L  (the side leaf) and  G, L |- goal  (the main branch), and a composite such as zero-it leaves the focus on whichever it did not close.\par
\noindent\textit{Uses:} no kernel operation: it records no inference.\par
\par\smallskip{\raggedright\noindent\texttt{(backup-one)}\quad\tactkind{meta}\index{backup-one@\texttt{backup-one} (tactic)}\par}\nobreak
\noindent Take back the last recorded command: restore the goal, the deduction graph and the proof script to the state before it.  Repeat to walk further back; (sp) clears the stack.  `undo' is an alias.\par
\noindent\textit{Uses:} no kernel operation: it records no inference.\par
\noindent\textit{When useful:} the last command was a mistake -- a greedy `di' that ate the induction, a `cut' you did not mean, a branch you want back\par

The one undo in VNB.  It restores the deduction graph -- not merely the focus -- by rolling back the journalled arrow and grounding writes and DROPPING every sequent and inference node posted since, so no node from the abandoned branch survives to be counted as an open leaf and `qed' cannot be handed a phantom obligation.  *proof-script* and the live trace are rewound with it, so the saved script and the printed proof are the proof you actually kept.  It backs up over the LAST RECORDED command: a command that changed nothing was never recorded and is not a step to back up over.  One thing is deliberately not restored -- *fresh-counter*, the eigenvariable source -- so backing up over an `ai' or `ew' that minted `u\_4' and re-running it mints `u\_5'; the proof is the same, the witness has a different name.  Depth is capped at *vnb-undo-depth* (64).

\par\smallskip{\raggedright\noindent\texttt{(undo)}\quad\tactkind{meta}\index{undo@\texttt{undo} (tactic)}\par}\nobreak
\noindent Alias for (backup-one).\par
\noindent\textit{Uses:} no kernel operation: it records no inference.\par
\noindent\textit{When useful:} the last command was a mistake -- alias for backup-one\par
\par\smallskip{\raggedright\noindent\texttt{(show)}\index{show@\texttt{show} (tactic)}\par}\nobreak
\noindent Redisplay the current proof state.\par
\noindent\textit{Uses:} no kernel operation: it records no inference.\par
\par\smallskip{\raggedright\noindent\texttt{(show-graph)}\index{show-graph@\texttt{show-graph} (tactic)}\par}\nobreak
\noindent Print EVERY node of the current proof, closed ones included: one line per sequent node in creation order (the bracketed number the panel prints, the sequent, [GROUNDED] when justified), and under each node one line per inference into it -- the rule and the numbers of the hypothesis nodes it rests on.  The focus and the open leaves are marked.  This is where a branch a composite closed is found: after (cut L), the side leaf G |- L and the main branch G, L |- goal are both nodes, and zero-it on an existential closes the main one.  Records nothing; the [Graph] button of the Focus Workspace shows it.\par
\par\smallskip{\raggedright\noindent\texttt{(dial notch)}\index{dial@\texttt{dial} (tactic)}\par}\nobreak
\noindent The PRESENTATION RESOLUTION DIAL: how much of the sequent to show.  r0 kernel S-expression, r1 surface syntax (the default, and the last INVERTIBLE notch -- what is printed re-parses to what was printed), r2 guards and typing hypotheses elided, r3 also structures destructured to their declared slot names and deep subterms elided.  (dial) reports the notch, (dial n) sets it, (dial 'auto) reads the current goal and suggests one, (dial 'why) shows the score behind the suggestion.  For ONE FORMULA rather than the proof state, (dial-wff x [n]) -- x a wff, a ``surface string'' or an S-expression, n defaulting to the dial's current notch: (dial-wff (lookup-theorem 'matmul-assoc) 3).\par
\noindent\textit{Uses:} no kernel operation: it records no inference.\par

It is a FILTER, not a set of modes, and two constraints make it one.  The rungs NEST: turning it up only ever SUBTRACTS, so everything readable at a higher notch is readable below it -- checked in the suite, atom by atom.  And the APERTURE IS ON SCREEN: every notch above r1 names itself on the line and counts what it hid, because a lossy render that does not announce its own lossiness is the defect class of `(f)' printing as `f'.  The auto-notch SUGGESTS a starting notch and then STAYS PUT: it is consulted when you ask, never per sequent, so the lens cannot move while you are reading a proof.

\par\smallskip{\raggedright\noindent\texttt{(dial-wff x notch)}\index{dial-wff@\texttt{dial-wff} (tactic)}\par}\nobreak
\noindent Present ONE FORMULA at a resolution notch, instead of the current proof state: x is a wff, a ``surface string'' or an S-expression, and notch defaults to the dial's current setting.  (dial-wff (lookup-theorem 'matmul-assoc) 3).  Same rungs, same aperture line, no proof required.\par
\noindent\textit{Uses:} no kernel operation: it records no inference.\par

`sequent->string' is the display path for a proof STATE, so the dial reaches a formula only while a proof is holding it.  A theorem looked up, a statement being drafted, a hypothesis pasted from somewhere wants the same dial, and this is it.

\par\smallskip{\raggedright\noindent\texttt{(goal-status)}\index{goal-status@\texttt{goal-status} (tactic)}\par}\nobreak
\noindent One-line summary: done / N open goals.\par
\noindent\textit{Uses:} no kernel operation: it records no inference.\par
\par\smallskip{\raggedright\noindent\texttt{(repeat thunk [cap])}\index{repeat@\texttt{repeat} (tactic)}\par}\nobreak
\noindent Run thunk until it stops changing the proof state (LCF REPEAT).\par
\noindent\textit{Uses:} no kernel operation: it records no inference.\par
\par\smallskip{\raggedright\noindent\texttt{(orelse t1 t2 ...)}\index{orelse@\texttt{orelse} (tactic)}\par}\nobreak
\noindent Run thunks in order, stop at the first that makes progress (LCF ORELSE).\par
\noindent\textit{Uses:} no kernel operation: it records no inference.\par
\par\smallskip{\raggedright\noindent\texttt{(preamble)  |  (preamble 'NAME)  |  (preamble ``file.pre'')}\quad\tactkind{composite}\index{preamble@\texttt{preamble} (tactic)}\par}\nobreak
\noindent Apply an editable RULE FILE to the focus goal, one step at a time: each rule reads `if the sequent looks like this, do that', and the first rule that matches AND makes progress is kept; the loop repeats on the new focus until the goal is done, no rule applies, or 40 steps.  Rules come from \textasciitilde{}/.vnb-preamble.pre when it exists, else the shipped preambles/default.pre (peel, split, close by context, reflexivity, arithmetic, type-term, zero-it, ineq, beta, unfold, and a tentative cut).  Prints one report: what each rule did, which rules matched but were rejected, and where it stopped and why.\par
\noindent\textit{Uses:} \texttt{cut}.\par
\noindent\textit{When useful:} you want the routine steps done for you by rules you can read and edit, with a report of what each rule did and where it stopped\par

The notes-45 `giant COND' as data you can edit.  A rule is (rule NAME (goal PATTERN) (with PATTERN) (guard ...) (do FORM) (probe ...)): patterns are raw S-expressions with schema variables ?x, matched against the goal and the assumptions (bindings shared, a ?x in binder position binds the bound variable's name); guards are typed, head, closed, not-in-context, one-of, occurs, unfolds; the action is an ordinary surface command with the variables filled in.  A rule is TRIED on a scratch copy first and KEPT only when its probe holds there and again on the proof (grounded: the goal closed; progress, the default: something changed and the open goals did not multiply; changed; (lands F)); otherwise nothing it did is kept.  A rule whose action is (cut F) is a TENTATIVE CUT: the rules (without their cuts) are run on F, and F is proven in place or left open and reported as OWED -- a support is never installed.  Each kept step is recorded AS ITSELF, so the page of the proof is the ordinary steps and replays without any rule file.  A malformed rule is reported by name and left out; a rule that raises is reported and skipped.  The value is a verdict: status (done, owed, stalled, cap, error), the steps, the rules fired and rejected, the owed claims, the report text.  (preamble '(induct) ...) with CLAUSE LISTS is the older strategy pipeline, also (preamble-clauses ...).  (Technically: composite -- no rule of its own; every action is kernel-checked; transactional per firing.)

\par\smallskip{\raggedright\noindent\texttt{(quietly thunk)}\index{quietly@\texttt{quietly} (tactic)}\par}\nobreak
\noindent Run thunk with state-dump output suppressed; returns its value.\par
\noindent\textit{Uses:} no kernel operation: it records no inference.\par

\subsection*{Choosing and naming witnesses}

\par\smallskip{\raggedright\noindent\texttt{(minimize! '(v1 ... vk) GUARD MEASURE)}\quad\tactkind{composite}\index{minimize"!@\texttt{minimize"!} (tactic)}\par}\nobreak
\noindent Choose v1..vk satisfying GUARD with the NN-valued MEASURE as small as possible: lands the witnesses and their minimality, and opens the two obligations any minimisation owes (MEASURE lands in NN; GUARD is satisfiable).\par
\noindent\textit{Uses:} \texttt{and-elim} \texttt{and-intro} \texttt{arith-forsome} \texttt{arith-ground} \texttt{arith-simplify} \texttt{assumption} \texttt{cartesian-decompose} \texttt{cut} \texttt{detach} \texttt{eq-subst} \texttt{forall-elim} \texttt{forall-intro} \texttt{forsome-elim} \texttt{forsome-intro} \texttt{iff-elim} \texttt{iff-intro} \texttt{implies-intro} \texttt{macete} \texttt{not-elim} \texttt{not-intro} \texttt{or-elim} \texttt{reflexivity} \texttt{sep-mem-elim} \texttt{sep-mem-intro} \texttt{theorem-assumption} \texttt{tuple-equality-decompose}.\par
\noindent\textit{When useful:} the goal falls to a `least such' / minimal-counterexample argument -- descent proofs (sqrt 2, sqrt 3 irrational), least-degree or least-pivot witnesses\par

The `least such' / minimal-counterexample step, mechanised.  You give three things: the variables to choose, a GUARD formula they must satisfy, and a natural-number-valued MEASURE term to make small.  On the main branch it hands you fresh eigenconstants w1..wk with two hypotheses -- GUARD holds of them, and nothing satisfying GUARD has a strictly smaller MEASURE -- and it opens exactly the two side goals a minimisation genuinely owes: that MEASURE lands in NN wherever GUARD holds (the TYPE goal), and that GUARD holds somewhere (the NONEMPTY goal).  Everything in between -- forming the value set, well-ordering it, unpacking the least witness, and restating minimality in terms of your variables rather than the set -- is done for you.  Every vi must occur in MEASURE (a variable the measure ignores is not one you are minimising over).  Because it quantifies over the vk at the META level, no lambda and no FUN(U,NN) typing obligation ever enters the logic -- it takes a formula and a term, which is what a driver has in hand, and never forms the set U at all.  Returns (list (w1 ... wk) TYPE-node NONEMPTY-node); either node is \#f when that obligation was already in context up to alpha, so test before refocusing.  (Technically: composite -- it drives cut / di / ai / ew / sep-mi / sep-me / fact / inst / detach! / subst / ass / rfl and adds no kernel rule; its one mathematical appeal is `nn-least-element', the well-ordering of NN, proven from ord-well-ordered.  NO choice principle is used: well-ordering returns a MEMBER of a separation set, and sep-me recovers the witness.)

\par\smallskip{\raggedright\noindent\texttt{(obtain LANE)}\quad\tactkind{composite}\index{obtain@\texttt{obtain} (tactic)}\par}\nobreak
\noindent Run LANE -- a thunk whose forward step lands a `there exists' into context -- then eliminate that existential and RETURN the fresh witness's name, read off the proof state rather than guessed.\par
\noindent\textit{Uses:} every one of the 64 kernel operations.\par
\noindent\textit{When useful:} a forward step yields `there exists ...' and you want to name the witness for later use, without guessing the engine's eigenvariable\par

The `let w be such an x' step, with the eigenvariable named for you.  LANE is a zero-argument procedure whose effect is to land some `there exists v. P(v)' among your hypotheses -- e.g. (lambda () (fact 'nn-3-div-square w)).  obtain runs it, spots the existential that newly appeared, eliminates it (introducing a fresh eigenvariable and landing P at that witness), unpacks any conjunction in the body, and hands back the eigenvariable's NAME -- identified by free-variable set-difference (the symbol now in context that was not there before), so the driver never guesses what the engine called it.  This is the robust way to name a descended witness: (define k (obtain (lambda () ...))), then use k.  Returns \#f (with a note) if the lane landed no existential.  (Technically: composite -- a guarded forward discharge followed by forsome-elim (ai) and AND-splitting.  Plain existential elimination, so it owes nothing and uses no choice.)

\par\smallskip{\raggedright\noindent\texttt{(vlet (n1 ...) FORMER)}\quad\tactkind{composite}\index{vlet@\texttt{vlet} (tactic)}\par}\nobreak
\noindent Bind proof-object names from the current state.  FORMER is (match PATTERN) -- bind each name to its slot in a context formula matching PATTERN -- or (choice [v body]) -- eliminate an in-context existential and bind its witness.\par
\noindent\textit{Uses:} \texttt{and-elim} \texttt{cut} \texttt{forsome-elim} \texttt{iff-elim} \texttt{not-elim} \texttt{or-elim}.\par
\noindent\textit{When useful:} you need to NAME a witness or a matched subterm from the proof state, rather than navigate to it by shape\par

A binding form for the pieces of a proof, so a driver NAMES what it needs instead of navigating to it by shape.  Two FORMERs.  (match PATTERN): the listed names ARE the holes (wildcards) in PATTERN; everything else is literal.  vlet searches the context for a formula or subterm matching PATTERN and binds each name to what filled its slot -- pure selection, no proof step, no obligation; a hard error if nothing matches (you named a piece of something absent).  (choice): eliminate the sole existential in context and bind its witness.  (choice v body): present-else-debt -- if `there exists v. body' is already in context up to alpha, eliminate it; otherwise cut it, LEAVE the existence side-goal open as a debt leaf, and eliminate on the main branch -- either way the witness is bound and body[witness] is landed.  vlet's (choice) and `obtain' are the same underlying mechanism -- eliminate an existential, name the witness by free-variable difference -- differing only in what they take: obtain runs a lane and names what it just produced, vlet names an existential already (or, with a body, about to be) in context.  No choice AXIOM is used; a single witness is plain existential-elimination.  (Technically: a define-syntax expanding to (define n ...) over vlet--match / vlet--choice!; composite, no kernel rule beyond forsome-elim and, for present-else-debt, cut.)




%% =========================================================
\section{Focus management}

After each command the focus advances automatically to the first
ungrounded subgoal just introduced.  If there are no new subgoals, it
moves to any remaining open goal.

\begin{remark}
Every proof command operates exclusively on the focus node; the other
open goals are untouched.  There can be several open goals simultaneously
(for example, after \texttt{(ci)} on a conjunction goal, or after
\texttt{(cut f)} which always creates two subgoals).  \texttt{(show)}
displays them all, numbering the non-focus goals so you can use those
numbers with \texttt{(focus $n$)}.  Until you call \texttt{(focus $n$)},
the system always works on whichever goal it chose automatically.
\end{remark}

In the interactive session use:
\begin{VNBcode}
(show)       ; list all open goals (numbered) and the current focus
(focus n)    ; switch focus to the n-th open goal (1-based)
\end{VNBcode}

The underlying API (for use outside the interactive session) is:
\begin{VNBcode}
(focus-on ps sqn)            ; focus on a specific sequent node
(focus-on-first-open ps)     ; focus on the first open goal
(proof-open-goals ps)        ; list all open goals
\end{VNBcode}

%% =========================================================
\section{Macetes}
\label{sec:macetes}

A \emph{macete} is a named, composable rewrite strategy.  When a
theorem is installed, its macete is automatically derived by extracting a
source pattern and replacement pattern from the theorem's logical form:

\begin{center}
\begin{tabular}{ll}
\toprule
Theorem form & Rewrite \\
\midrule
$\forall\bar{x}.\;(s = r)$ & $s \mapsto r$ \\
$\forall\bar{x}.\;(s \Leftrightarrow r)$ & $s \mapsto r$ \\
$\forall\bar{x}.\;(\text{conds} \Rightarrow (s = r))$ & $s \mapsto r$ when conds hold \\
$\forall\bar{x}.\;p$ & $p \mapsto \top$ \\
\bottomrule
\end{tabular}
\end{center}

The quantifiers need not all be leading.  A \texttt{forall} occurring in
positive position --- under an \texttt{implies} consequent or inside an
\texttt{and} --- is pulled to the front before the source and replacement
patterns are extracted, so a theorem such as
\texttt{forall([a,b], $H$ implies forall(x, $s$ iff $r$))} still yields the
rewrite $s \mapsto r$ (conditional on $H$).

Macetes can be composed:
\begin{VNBcode}
(macete-series m1 m2 m3)   ; apply in sequence
(macete-repeat m)           ; apply until no change
(macete-try m)              ; apply m, succeed even if m fails
(macete-choice m1 m2 m3)   ; first one that applies
\end{VNBcode}

To install a custom macete:
\begin{VNBcode}
(install-theorem! 'eq-symm
  (make-wff-from-string "forall([x, y], x = y iff y = x)"))

; Use it:
(cmd-apply-macete ps 'eq-symm)
\end{VNBcode}

\subsection{Applying macetes to wffs outside a proof}

\begin{sloppypar}The function \texttt{apply-macete} applies a named macete directly to a
\texttt{wff} without creating a proof state:\end{sloppypar}

\begin{itemize}
\item \texttt{(apply-macete \textit{name} \textit{w})} --- rewrites \textit{w};
  returns a new \texttt{wff} if any rewrite fires, or \texttt{\#f} if the
  pattern matches nowhere in \textit{w}.
\end{itemize}

This is purely syntactic: no deduction graph is involved and no
theorem is proved.  It is useful for inspecting rewrite behavior,
preprocessing formulas, or building higher-level tools outside of proofs.

\paragraph{Unconditional equation macete.}
\begin{VNBcode}
(install-theorem! 'nn-zero-succ
  (make-wff-from-string "forall([n in NN], succ(n) = n + 1)"))

(apply-macete 'nn-zero-succ (make-wff-from-string "succ(3) = x"))
; => wff for "3 + 1 = x"

(apply-macete 'nn-zero-succ (make-wff-from-string "3 in nn"))
; => #f               -- macete doesn't fire
\end{VNBcode}

\paragraph{Biconditional (\texttt{iff}) macete.}
\begin{VNBcode}
(install-theorem! 'eq-symm
  (make-wff-from-string "forall([x, y], x = y iff y = x)"))

(apply-macete 'eq-symm (make-wff-from-string "n + m = 0"))
; => wff for "0 = n + m"

(apply-macete 'eq-symm (make-wff-from-string "n = m implies truth"))
; => wff for "m = n implies truth"

(apply-macete 'eq-symm (make-wff-from-string "n in nn"))
; => #f
\end{VNBcode}

\paragraph{Closure axiom macete (conditional).}
\begin{VNBcode}
;; nn-add-closed: forall([n in NN, m in NN], n + m in NN)
;; Condition (n in NN and m in NN) must be present in local context.

(apply-macete 'nn-add-closed
  (make-wff-from-string "n in nn and m in nn implies n + m in nn"))
; => wff for "n in nn and m in nn implies truth"

(apply-macete 'nn-add-closed (make-wff-from-string "n + m in nn"))
; => #f  -- conditions not in context; rewrite correctly suppressed
\end{VNBcode}

\begin{remark}
\texttt{apply-macete} applies the rewrite depth-first throughout the
formula (outermost first within each subexpression).  Conditions of a
conditional macete are checked against the \emph{local context}
accumulated as the rewriter descends: when descending into the
consequent of \texttt{implies}, the antecedent is added to the local
context; when descending into the right conjunct of \texttt{and}, the
left conjunct is added; and so on (Monk~\cite{Monk1988}).  A condition that
matches a formula already present in the local context is considered
discharged; if any condition is unmet at a candidate site, the rewrite
is suppressed at that site.
\end{remark}

%% =========================================================
\section{Example proofs}

\subsection{Identity}

\begin{VNBcode}
(sp (make-wff-from-string "x in NN implies x in NN"))
(di)    ; assume (x in nn), prove (x in nn)
(ass)   ; it is in context
; => Proof complete.
\end{VNBcode}

\subsection{\texttt{and} elimination}

\begin{VNBcode}
(sp (make-wff-from-string "n in NN and m in NN implies n in NN"))
(di)                          ; assume (n in nn) and (m in nn)
(ai "n in NN and m in NN")   ; split into (n in nn) and (m in nn)
(ass)                         ; prove (n in nn)
\end{VNBcode}

\subsection{VNB axiom instance}

Proving \texttt{"a in b implies a in set"} from the base axiom:

\begin{VNBcode}
(sp (make-wff-from-string "a in b implies a in SET"))
(di)
(ta 'membership-implies-sethood)
(inst "forall([a, b], a in b implies a in set)" 'a)
(inst "forall([b], a in b implies a in set)" 'b)
(bc "a in b implies a in set")
(ass)
\end{VNBcode}

\subsection{Double negation (classical)}

\begin{VNBcode}
(sp (make-wff-from-string "not(not(x in NN)) implies x in NN"))
(di)                              ; assume not(not((x in nn))), prove (x in nn)
(pbc)                             ; assume not((x in nn)), prove falsity
(ai "not(not(x in NN))")         ; not(not(P)) and not(P) yield falsity
\end{VNBcode}

\subsection{Ordered tuples and Cartesian products}

\begin{VNBcode}
; From x in A, y in B, z in C prove [x,y,z] in cartesian(A,B,C)
(sp (make-wff-from-string
      "x in A and y in B and z in C implies [x, y, z] in cartesian(A, B, C)"))
(di)
(ai "x in A and y in B and z in C")
(ai "y in B and z in C")
(ci)              ; cartesian-intro generates 3 subgoals
(ass) (ass) (ass)

; NTH projection: nth(2, [x, y, z]) = y
(sp (make-wff-from-string "nth(2, [x, y, z]) = y"))
(nth-r)           ; reduces goal to (y = y)
(rfl)             ; closes by reflexivity
\end{VNBcode}

\subsection{Mathematical structures}

Algebraic structures are declared with \texttt{declare-structure}.
No quoting is needed: carrier names and operation signatures are written
directly.

\begin{VNBcode}
(declare-structure GROUP
  (carriers CARR)
  (op OPR (CARTESIAN CARR CARR) CARR)  ; operation: CARR x CARR -> CARR
  (constant IDEN CARR)                 ; identity: IDEN in CARR
  (op INV CARR CARR))                  ; inverse: CARR -> CARR
; Accessor indices: CARR -> 1, OPR -> 2, IDEN -> 3, INV -> 4
; (the group laws are folded in by additional (property ...) clauses)
\end{VNBcode}

\begin{sloppypar}The \texttt{(op \textit{name} \textit{domain} \textit{range})} clause declares
$\textit{name}(s) \in \texttt{fun}(\textit{domain}, \textit{range})$.
Higher-arity operations write the Cartesian product explicitly:
\texttt{(op opr (cartesian carr carr) carr)} for a binary operation.
A \texttt{(constant \textit{name} S)} clause declares an element of set $S$.
Accessor indices are assigned sequentially: carriers first (starting
at~1), then operations in declaration order.\end{sloppypar}

A \texttt{declare-structure} call installs two things automatically:
\begin{itemize}
\item \emph{Accessor macetes}: \texttt{carr}$(s)$ rewrites to
  \texttt{nth(1,$s$)}, \texttt{opr}$(s)$ to \texttt{nth(2,$s$)}, etc.
\item \emph{Definitional biconditional} \texttt{is-\textit{name}}: the axiom
  $\forall s.\;\texttt{is-group}(s) \Leftrightarrow
   (\texttt{carr}(s)\in\mathrm{set} \wedge \texttt{opr}(s)\in\cdots)$.
\end{itemize}

Given \texttt{is-group(g)} in context, the \texttt{is-group} macete
unfolds to:
\begin{VNBsnip}
CARR(g) in SET
and OPR(g) in fun(cartesian(CARR(g), CARR(g)), CARR(g))
and IDEN(g) in CARR(g)
and INV(g) in fun(CARR(g), CARR(g))
\end{VNBsnip}

\paragraph{Pre-loaded structures.}
The following structures are defined in \texttt{structure-library/} (one
file per structure) and loaded automatically.

\begin{center}\scriptsize
\begin{tabular}{@{}p{2.3cm}p{4.3cm}p{5.3cm}@{}}
\toprule
\textbf{Structure} & \textbf{Accessor slots (index)} & \textbf{Characteristic axioms} \\
\midrule
\texttt{semigroup} &
  \texttt{carr}\,(1), \texttt{opr}\,(2) &
  \texttt{semigroup-assoc} \\[4pt]
\texttt{monoid} &
  \texttt{carr}\,(1), \texttt{opr}\,(2), \texttt{iden}\,(3) &
  \texttt{monoid-assoc},
  \texttt{monoid-left-id},
  \texttt{monoid-right-id} \\[4pt]
\texttt{group} &
  \texttt{carr}\,(1), \texttt{opr}\,(2), \texttt{iden}\,(3), \texttt{inv}\,(4) &
  \texttt{group-assoc},
  \texttt{group-left-id},
  \texttt{group-left-inv} \\[4pt]
\texttt{ring} &
  \texttt{carr}\,(1), \texttt{add}\,(2), \texttt{mul}\,(3),
  \texttt{neg}\,(4), \texttt{zero}\,(5), \texttt{one}\,(6) &
  \texttt{ring-add-assoc}, \texttt{ring-add-comm},
  \texttt{ring-add-left-id}, \texttt{ring-add-left-inv},
  \texttt{ring-mul-assoc}, \texttt{ring-mul-left-id},
  \texttt{ring-mul-right-id},
  \texttt{ring-left-dist}, \texttt{ring-right-dist} \\[4pt]
\texttt{metric-space} &
  \texttt{pts}\,(1), \texttt{dist}\,(2) &
  \texttt{metric-pos},
  \texttt{metric-self-zero},
  \texttt{metric-zero-eq},
  \texttt{metric-sym},
  \texttt{metric-triangle} \\
\bottomrule
\end{tabular}
\end{center}

\begin{remark}
Each entry above is citable by name with \texttt{(ta 'name)}, which transports
it into the current proof.  They do \emph{not} all have the same standing,
however, and the column heading is a convenience rather than a claim.

A structure declaration generates an \texttt{is-\textit{x}} defining biconditional, and
each operation-property conjunct of it, unfolded, \emph{is} one of the
structure's laws.  Stating such a law separately therefore adds no mathematical
content, and the question is only how the system records that fact.  Three
answers are in use, and they cost a citing \texttt{qed} the same --- nothing ---
by three different routes:

\begin{itemize}
\item \emph{Proved.}  The \textsc{Metric-Space} laws (\texttt{metric-pos},
  \texttt{metric-self-zero}, \texttt{metric-zero-eq}, \texttt{metric-sym},
  \texttt{metric-triangle}) are established by \texttt{prove-metric-law!} in
  \texttt{structure-library/metric-laws.scm}; the \textsc{Group} laws
  (\texttt{group-assoc}, \texttt{group-left-id}, \texttt{group-left-inv},
  \texttt{group-identity-in}) and \texttt{abelian-group-opr-comm} by
  \texttt{stl--project!} in \texttt{structure-library/subtype-laws.scm}.  The
  unfold is \emph{checked}, not asserted to exist.
\item \emph{Installed and stamped \texttt{definitional}.}  The \textsc{Ring}
  laws are installed by \texttt{add-axiom!} in
  \texttt{structure-library/ring.scm} and then swept with
  \texttt{register-provenance!\ \ldots\ 'definitional} in the same file.  A
  \texttt{definitional} fact contributes the empty set to every bill, exactly as
  a \texttt{primitive} one does.
\item \emph{Wrapped at installation.}  \texttt{structure-library/module.scm}
  reaches the same place by the other door, installing its projections inside
  \texttt{(fluid-let ((*current-provenance* 'definitional)) \ldots)}.
\end{itemize}

Two laws in the table are genuinely \emph{derived} rather than projected, and
they do carry a warrant: \texttt{ring-mul-zero-left} and
\texttt{ring-mul-zero-right} ($0\cdot a = 0$, which needs distributivity and
cancellation, not merely the definition of a ring).
\texttt{reference/PROOF-DEBT.md} is the authority on which is which; this table
is not.  For example, to apply
\texttt{group-left-inv} (which asserts
$\forall s,a.\;\texttt{is-group}(s)\Rightarrow a\in\texttt{carr}(s)
 \Rightarrow \texttt{opr}(s)(\texttt{inv}(s)(a),a)=\texttt{iden}(s)$),
supply the group and the element as instantiation witnesses.
\end{remark}

If you prefer to work with raw tuples directly, the n-ary constructors
are still available:
\begin{VNBcode}
; Prove a tuple is in a Cartesian product
(sp (make-wff-from-string
  "G in power(U) and op in fun(cartesian(G,G),G) and e in G and inv in fun(G,G)
   implies [G, op, e, inv]
   in cartesian(power(U), fun(cartesian(G,G),G), G, fun(G,G))"))
(di)
(ai "G in power(U) and op in fun(cartesian(G,G),G) and e in G and inv in fun(G,G)")
(ai "op in fun(cartesian(G,G),G) and e in G and inv in fun(G,G)")
(ai "e in G and inv in fun(G,G)")
(ci)                              ; cartesian-intro: 4 subgoals
(ass) (ass) (ass) (ass)
\end{VNBcode}

\subsection{Applying a macete inside a proof}

The short-form command \texttt{(mac 'name)} applies the named macete to
the \emph{current goal}.  It unifies the macete's source pattern against
the goal formula (outermost-first, depth-first) and replaces the first
matched subterm with the replacement pattern.

\paragraph{Biconditional macete.}
\texttt{pairing-membership} states
$\forall a\,b\,x.\;x \in \mathrm{pair}(a,b) \Leftrightarrow (x=a \vee x=b)$.
Its source pattern is \verb|x in pair(a, b)| and its replacement is
\verb|x = a or x = b|.

\begin{VNBcode}
(sp (make-wff-from-string "2 in pair(2, 3)"))
; proof state:  []  |-  2 in pair(2, 3)
(mac 'pairing-membership)
; source: x in pair(a, b)   replacement: x = a or x = b
; unifies: x=2, a=2, b=3
; goal rewritten:  2 in pair(2, 3)  ->  2 = 2 or 2 = 3
; proof state:  []  |-  2 = 2 or 2 = 3
(oi-l)      ; or-intro-left: new goal  2 = 2
(rfl)       ; reflexivity
; => Proof complete.
\end{VNBcode}

\paragraph{Rewrite fires inside a larger formula.}
The rewrite finds its source pattern at the first matching position,
which may be nested inside a connective.

\begin{VNBcode}
(sp (make-wff-from-string "x in pair(a, b) implies x = a or x = b"))
; proof state:  []  |-  (x in pair(a,b)) implies (x = a or x = b)
(mac 'pairing-membership)
; pattern "x in pair(a, b)" matched at the left of implies
; goal rewritten:  (x = a or x = b) implies (x = a or x = b)
(di)        ; discharge implies: push antecedent to context
; context: ["x = a or x = b"]   goal: "x = a or x = b"
(ass)
; => Proof complete.
\end{VNBcode}

\begin{remark}
Macetes rewrite the goal; they do not rewrite the context.  If the
matching pattern appears only in an assumption, \texttt{(mac 'name)}
will not fire.  To use an assumption as a rewrite source, retrieve it
with \texttt{(ai ...)} so that it appears in the goal.
\end{remark}

\paragraph{Conditional macete.}
For a conditional theorem
$\forall\bar{x}.\;(\mathit{cond} \Rightarrow s = r)$, the macete still
fires unconditionally: it rewrites $s \mapsto r$ without generating
side-goals for $\mathit{cond}$.  The goal $s$ is replaced by~$r$
(or by $\top$ when the conclusion is not an equation or biconditional),
and the resulting $\top$ goal is closed with \texttt{(ass)}.

\begin{VNBcode}
; nn-add-closed:
;   forall([a in NN, b in NN], a + b in NN)
; As a macete: source = "a + b in nn"
;              replacement = truth
(sp (make-wff-from-string "2 + 3 in NN"))
; proof state:  []  |-  2 + 3 in NN
(mac 'nn-add-closed)
; unifies a=2, b=3
; conditions (a in NN, b in NN) not discharged
; goal rewritten:  2 + 3 in NN  ->  truth
(ass)       ; truth is trivially grounded
; => Proof complete.
\end{VNBcode}

\begin{remark}
The conditions of a conditional macete are silently dropped.  This
makes \texttt{(mac 'name)} fast for goals whose conditions are obviously
satisfied, but it is the user's responsibility to verify them; the
proof sketch above is technically unsound without separate proofs that
$2 \in \mathtt{nn}$ and $3 \in \mathtt{nn}$.  For fully verified
condition discharge, use \texttt{(ta 'name)} followed by explicit
instantiation and back-chaining.
\end{remark}

\subsection{Ground arithmetic with \texttt{(arith)}}

The \texttt{(arith)} command (primitive \texttt{cmd-arith}) closes any
goal whose formula reduces to true under exact arithmetic evaluation:

\begin{VNBcode}
(sp (make-wff-from-string "1/3 + 1/6 = 1/2"))
(arith)                              ; closes by exact rational arithmetic

(sp (make-wff-from-string "2 ^ 10 = 1024"))
(arith)                              ; 2^10 evaluates to 1024

(sp (make-wff-from-string "3 + 4 <= 2 * 5"))
(arith)                              ; 7 <= 10
\end{VNBcode}

\subsection{The calculator pattern}

\texttt{(arith)} also recognises a goal of the form
\begin{VNBsnip}
forsome([n in CLASS], EXPR = n)
\end{VNBsnip}
where \texttt{EXPR} is a ground arithmetic expression and \texttt{CLASS}
is one of \texttt{nn}, \texttt{zz}, \texttt{qq}, \texttt{rr},
\texttt{cc}.  It evaluates \texttt{EXPR}, checks that the result belongs
to \texttt{CLASS}, and closes the goal with that value as the witness.
This lets the prover double as a calculator:

\begin{VNBcode}
(sp (make-wff-from-string "forsome([n in NN], 2^10 = n)"))
(arith)   ; closes with witness 1024

(sp (make-wff-from-string "forsome([q in QQ], 1/3 + 1/4 = q)"))
(arith)   ; closes with witness 7/12
\end{VNBcode}

\begin{remark}
The symmetric form \texttt{n = EXPR} is also accepted.
The pattern is intentionally narrow: \texttt{EXPR} must be fully ground
(no free variables), and the existential variable must appear alone on
one side of the equality.
\end{remark}

\subsection{\texorpdfstring{$\beta$}{beta}-reduction with \texttt{(beta)}}

Applied \texttt{lambda} (or \texttt{lambdoid}) terms are reduced
in-place, then the underlying numeric or logical content is closed
by another rule:

\begin{VNBcode}
(sp (make-wff-from-string "lambda([x in nn], 2 * x)(5) = 10"))
(beta)   ; goal becomes (2 * 5 = 10)
(arith)  ; closes
\end{VNBcode}

\subsection{Auto-installing arithmetic theorems}

\texttt{vnb-create-num-theorem} verifies a closed arithmetic wff by
ground evaluation and installs it as a named theorem.  The theorem is
then available via \texttt{(ta \textit{name})}:

\begin{VNBcode}
(vnb-create-num-theorem 'pythag-3-4-5 "3 ^ 2 + 4 ^ 2 = 5 ^ 2")
; => pythag-3-4-5

(sp (make-wff-from-string "truth implies 3 ^ 2 + 4 ^ 2 = 5 ^ 2"))
(di) (ta 'pythag-3-4-5) (ass)
; => Proof complete.
\end{VNBcode}

False or undecidable formulas are rejected:

\begin{VNBcode}
(vnb-create-num-theorem 'bogus "1 + 1 = 3")
; ERROR: formula evaluates to FALSE

(vnb-create-num-theorem 'bogus2 "x = 5")
; ERROR: not a decidable closed arithmetic sentence
\end{VNBcode}

\subsection{Installing a proven theorem}
\label{sec:qed-and-subst}

A completed proof does not by itself add anything to the library.  Use
\texttt{(qed \textit{name})} to install the root assertion as a named
theorem:

\begin{VNBcode}
(sp (make-wff-from-string "n in NN and m in NN implies m in NN and n in NN"))
(di) (ai "n in NN and m in NN") (di) (ass) (ass)
(qed 'and-comm)
; => and-comm
; (lookup-theorem 'and-comm) returns the formula S-expression.
\end{VNBcode}

\paragraph{What \texttt{subst-wff} is and is not.}
\texttt{subst-wff} applies capture-avoiding substitution to a wff's
underlying S-expression and returns a new \texttt{wff} (re-validated
through \texttt{make-wff}).  It is \emph{not} a logical inference;
after calling it you must reprove the result.  It is useful for
generating related goals from a known formula:

\begin{VNBcode}
(define w
  (subst-wff '((n . (in x nn)) (m . (in y rr)))
             (lookup-theorem 'and-comm)))
; w is the wff for (x in nn) and (y in rr) implies (y in rr) and (x in nn)
\end{VNBcode}

\subsection{Saving and replaying proof scripts}
\label{sec:replay}

Every interactive command you issue between \texttt{(sp\ \ldots)} and
\texttt{(qed\ \ldots)} is appended to a global accumulator
\texttt{*proof-script*}.  \texttt{(qed \textit{name})} installs the
proven theorem and saves the script under the same name; you can also
call \texttt{(save-proof \textit{name})} on demand.

\begin{VNBcode}
(sp (make-wff-from-string "n in NN and m in NN implies m in NN and n in NN"))
(di) (ai "n in NN and m in NN") (di) (ass) (ass)
(qed 'and-comm)
; theorem AND script installed under the same name

(lookup-proof 'and-comm)
; => ((di) (ai (and (in n nn) (in m nn))) (di) (ass) (ass))
\end{VNBcode}

\paragraph{Replay verbatim.}
\texttt{(replay-proof \textit{name})} re-runs the saved commands
sequentially against the current proof state.  Useful for redoing a
proof in a fresh session, or applying an identical script to a sequent
of the same shape.

%% =========================================================
\subsection{\texttt{nn} induction for type properties}
\label{sec:ex-ni-type}

The \texttt{(ni)} command reduces a goal
$\forall n \in \mathtt{nn}.\; P(n)$ to two subgoals:
a base case $P(0)$ (focused first) and a step case
$\forall n \in \mathtt{nn}.\; P(n) \Rightarrow P(\mathrm{succ}(n))$.
The following proof shows that $\mathit{prod\text{-}ord}(m, f, n)$
stays in the carrier $\mathtt{carr}(m)$ for every $n \in \mathtt{nn}$,
deriving \texttt{prod-ord-type} from first principles.

\begin{VNBcode}
(sp (make-wff '(IMPLIES (IS-MONOID m)
               (IMPLIES (IN f (FUN NN (CARR m)))
                 (FORALL n (IMPLIES (IN n NN)
                   (IN (PROD-ORD m f n) (CARR m))))))))
(di) (di)   ; IS-MONOID(m) and f in FUN(NN,CARR(m)) enter context
(ni)        ; NN-induction; BASE gets focus first

;; --- BASE: PROD-ORD(m,f,0) in CARR(m) ---
;; Strategy: PROD-ORD(m,f,0) = IDEN(m) (prod-ord-zero),
;;           IDEN(m) in CARR(m)            (monoid-identity-in),
;;           so by eq-subst-membership the goal closes.
(ta 'eq-subst-membership)
(inst ... '(PROD-ORD m f 0)) (inst ... '(IDEN m)) (inst ... '(CARR m))
(bc '(IMPLIES (AND (= (PROD-ORD m f 0) (IDEN m)) (IN (IDEN m) (CARR m)))
              (IN (PROD-ORD m f 0) (CARR m))))
(di)    ; AND-split: focus = equality goal
(ta 'prod-ord-zero) (inst ... 'm) (inst ... 'f) (ass)
        ; focus jumps: manually refocus on membership goal
(ta 'monoid-identity-in) (inst ... 'm)
(bc '(IMPLIES (IS-MONOID m) (IN (IDEN m) (CARR m)))) (ass)
        ; BASE complete; focus moves to STEP

;; --- STEP: forall([n in NN], IH => PROD-ORD(m,f,succ n) in CARR(m)) ---
(di)    ; peel FORALL n -> n_k; (IN n_k NN) enters context
;; n_k is the bound name introduced by (di); fetch it via:
;;   (cadr (wff-formula
;;           (car (sequent-node-assumptions
;;                  (proof-state-focus *ps*)))))
(di)    ; peel IH: PROD-ORD(m,f,n_k) in CARR(m) enters context
;; Strategy:
;;   PROD-ORD(m,f,succ n_k)
;;     = OPR(m)(PROD-ORD(m,f,n_k), f(n_k))   (prod-ord-succ)
;;   then close via monoid-carrier-closed-opr.
(ta 'eq-subst-membership) (inst ...) (inst ...) (inst ...)
(bc ...)
(di)    ; AND-split: focus = equality goal
(ta 'prod-ord-succ)
(inst ...) (inst ...) (inst ...)
(bc ...) (ass)   ; n_k in NN closes it
;; focus moves to membership goal: OPR(m)(...) in CARR(m)
(ta 'monoid-carrier-closed-opr)
(inst ...) (inst ...) (inst ...)
(bc ...)
(di) (ass)   ; IS-MONOID(m)
(di) (ass)   ; PROD-ORD(m,f,n_k) in CARR(m) [IH]
;; f(n_k) in CARR(m) via fun-apply-type:
(ta 'fun-apply-type)
(inst ...) (inst ...) (inst ...) (inst ...)
(bc ...) (di) (ass) (ass)
;; closes f in FUN(NN,CARR(m)) and n_k in NN
\end{VNBcode}

\begin{remark}
The \texttt{inst} chains are elided above; each peels one leading
\texttt{forall} from a theorem schema.  Two subtleties arise in practice.
First, \texttt{(di)} on an AND-goal always focuses on the left (equality)
conjunct first.  After \texttt{(ass)} closes that conjunct, automatic
focus advancement jumps to the earliest ungrounded goal in the deduction
graph --- which may be the STEP node, not the right (membership) conjunct.
The fix is to capture the membership node
immediately after the AND-split:
\begin{VNBcode}
(di)   ; AND-split
(let* ((dg       (proof-state-dg *ps*))
       (mem-node (car (reverse (dg-sequent-nodes dg)))))
  ... prove equality ... (ass)   ; focus jumps
  (set-proof-state-focus! *ps* mem-node)
  ... prove membership ...)
\end{VNBcode}

Second, \texttt{alpha-equiv?} --- used by \texttt{(inst)} and
\texttt{(bc)} when searching assumptions --- relies on structural
equality of formula heads.  Compound heads such as
\texttt{((opr m) a b)} (a curried two-argument application; see
Section~\ref{sec:apply-tupling}) are compared correctly only when
the head pair is the same \emph{object}.  If \texttt{alpha-equiv?}
fails to locate a formula you know is in context, the likely cause
is a compound-head mismatch; the fix is in
\texttt{expressions.scm}'s \texttt{alpha-equiv-under?}.
\end{remark}

\begin{remark}[The curried/tupled apply convention]
\label{sec:apply-tupling}
VNB's S-expression form \texttt{(f a$_1$ \ldots a$_n$)} is what the
string-syntax parser emits from \texttt{f(a$_1$,\,\ldots,\,a$_n$)} ---
verbatim, no implicit tupling.  An operation typed
\texttt{f $\in$ fun(cartesian(A$_1$,\,\ldots,\,A$_n$),~B)}, however,
takes one argument from the Cartesian product.  The two views are
reconciled by the library axioms
\begin{equation*}
\texttt{apply-tupling-$n$:} \qquad
(f\ a_1\ \ldots\ a_n)\ \mathrel{==}\ (f\ (\texttt{list}\ a_1\ \ldots\ a_n))
\quad\text{for each arity }n
\end{equation*}
(installed for $n=1,2,3$ in \texttt{theorem-library/axioms.scm};
add more arities by declaration as needed).  Quasi-equality
$\mathrel{==}$, not partial $=$: the convention holds vacuously when
either side is undefined.  Combined with \texttt{fun-apply-type} and
\texttt{quasi-eq-def}, these give the $n$-ary apply closure
\((f\ a_1\ \ldots\ a_n)\in B\) from the operation's signature and
the membership of each $a_i$ in its declared component.
\end{remark}

\subsection{Infinite descent: \texorpdfstring{$\sqrt3$}{sqrt3} is irrational}
\label{sec:ex-sqrt3}

Stated without reals or division --- the form the checker proves, entirely
inside $\mathtt{nn}$ --- $\sqrt3$ irrational is
\[
  \forall p,q \in \mathtt{nn}.\quad
  q \neq 0 \;\Longrightarrow\; p\cdot p \neq 3\cdot(q\cdot q) ,
\]
that is, the goal
\begin{VNBsnip}
(FORALL p (IMPLIES (IN p NN) (FORALL q (IMPLIES (IN q NN)
  (IMPLIES (NOT (= q 0)) (NOT (= (* p p) (* 3 (* q q)))))))))
\end{VNBsnip}

\paragraph{The argument.}
Infinite descent.  Suppose the equation has a solution with $q \neq 0$, and
among all such solutions choose one whose numerator $w$ is \emph{least}.  Since
$3 \mid w\cdot w$ we have $3 \mid w$, say $w = 3k$; substituting and cancelling
the factor $3$ gives $q_0\cdot q_0 = 3\cdot(k\cdot k)$, whence $3 \mid q_0$, say
$q_0 = 3m_0$; cancelling once more gives $k\cdot k = 3\cdot(m_0\cdot m_0)$.  Thus
$(k,m_0)$ is again a solution, with $k \neq 0$ and $k < 3k = w$ --- a numerator
strictly below the least one.  Contradiction.

The entire mathematical content is the single implication
\begin{equation}\label{eq:3divsq}
  3 \mid p\cdot p \;\Longrightarrow\; 3 \mid p ;
\end{equation}
everything else --- clearing denominators, the two cancellations, the order
step $k < 3k$ --- is bookkeeping.

\paragraph{Why parity is not special.}
The $\sqrt2$ proof is the same skeleton with \emph{even} for \emph{divisible by
$3$}: there \eqref{eq:3divsq} reads ``$p\cdot p$ even $\Rightarrow p$ even'' and
follows from parity, every natural being $2k$ or $2k+1$.  For $\sqrt3$ it follows
from the trichotomy --- every natural is $3k$, $3k+1$, or $3k+2$ --- together
with the fact that the three residues square to $0,1,1$ modulo $3$, so a square
is divisible by $3$ only when its root is.  Parity is the $r=2$ instance of a
uniform statement about residues modulo $r$; the $\sqrt3$ development is the
$\sqrt2$ development with the residue system widened by one, and nothing changes
in kind.

\paragraph{The lemma ladder.}
The mod-$3$ arithmetic (\texttt{theorem-library/nn-mod3-proof.scm}) is built by
induction from the two Peano recursion equations for $+$ and $\cdot$, in exact
parallel with the parity ladder (\texttt{nn-parity-proof.scm}).  Each lemma is
\emph{proven}; none is asserted.

\begin{center}
\begin{tabular}{lll}
\toprule
parity ($r=2$) & this proof ($r=3$) & statement \\
\midrule
\texttt{nn-two-mul-succ}     & \texttt{nn-three-mul-succ} & $r\cdot\mathrm{succ}(k) = \mathrm{succ}^{\,r}(r\cdot k)$ \\
\texttt{nn-parity}           & \texttt{nn-trichotomy-3}   & every $n$ is $r\cdot k + i$, residue $i<r$ \\
\texttt{nn-parity-exclusive} & \texttt{nn-3-not-succ}     & $r\cdot x$ is never $r\cdot y + 1$ \\
\texttt{nn-even-square}      & \texttt{nn-3-div-square}   & $r \mid p\cdot p \Rightarrow r \mid p$ \\
\bottomrule
\end{tabular}
\end{center}

\noindent For $r=3$ a single residue-inequality suffices: both nonzero residues
square to $r\cdot j + 1$, so only ``not $\equiv 1$'' is ever needed.

\paragraph{The descent, mechanized.}
The driver (\texttt{theorem-library/sqrt3-proof.scm}) states the goal and hands
the descent to three tools; the intermediate steps are elided.

\begin{VNBcode}
(sp (make-wff '(FORALL p (IMPLIES (IN p NN) (FORALL q (IMPLIES (IN q NN)
     (IMPLIES (NOT (= q 0)) (NOT (= (* p p) (* 3 (* q q)))))))))))
(di)(di)(di)(di)(di)(di)                 ; intro p,q, hyps; assume p*p = 3*(q*q)

(define R (minimize! '(pv) (guard 'pv) 'pv))    ; choose the LEAST numerator w
;; ... discharge minimize!'s two obligations: measure in NN; guard satisfiable

(define k (obtain (lambda () (fact 'nn-3-div-square w))))   ; 3|w  =>  w = 3k
;; ... q0*q0 = 3*(k*k) by have! ; 3|q0 => q0 = 3*m0 by obtain ; k*k = 3*(m0*m0)

(inst+ minf k)              ; minimality gives  w <= k
(fact 'nn-lt-triple k)      ; but k < 3k = w -- contradiction
\end{VNBcode}

\begin{itemize}
\item \texttt{minimize!} supplies the least-numerator witness together with its
  two honest obligations --- the measure lands in $\mathtt{nn}$, and the guard
  is satisfiable --- and appeals to no choice principle.
\item \texttt{obtain} names each descended witness --- the $k$ with $w = 3k$,
  the $m_0$ with $q_0 = 3m_0$ --- by \emph{reading it off the proof state} (the
  variable left free once the existential is eliminated), never by guessing a
  machine-generated name.  This is what makes the descent survive a re-run.
\item \texttt{have!} discharges each intermediate equality and typing from
  context, so the script reads as the argument does: ``we have
  $q_0\cdot q_0 = 3\cdot(k\cdot k)$; we have $k \neq 0$; \ldots''.
\end{itemize}

\begin{remark}[The \texttt{minimize!} step]
\label{remark:minimize}
The enigmatic-looking \texttt{(minimize! '(pv) (guard 'pv) 'pv)} decodes as:
choose the single number \texttt{pv} satisfying the descent guard, with
\texttt{pv} itself as the quantity to minimise --- that is, take the
\emph{least} numerator meeting the guard.

The second argument to \texttt{minimize!} is the \textbf{guard formula} itself:
an ordinary VNB formula, free in the variables being chosen (here \texttt{pv}),
stating the property those values must satisfy.  \texttt{guard} is \emph{not} a
built-in and is nothing special to \texttt{minimize!} --- it is a helper the
proof author writes (named \texttt{s3-guard} in \texttt{sqrt3-proof.scm}), a
one-argument function that takes a term and returns the guard formula about it.
So \texttt{(guard 'pv)} yields the formula with \texttt{pv} free --- exactly what
\texttt{minimize!} wants --- and the \emph{same} helper, called later at the
chosen witness $w$ as \texttt{(guard $w$)}, restates that property of $w$.  Here
the property is ``\texttt{pv} is a natural number that is the numerator of a
fraction whose square is~$3$''; written out, the helper is just a formula
constructor:

\begin{VNBcode}
(define (s3-guard pv)
  (list 'AND (list 'IN pv 'NN)
        (list 'FORSOME 'q_
          (list 'AND (list 'IN 'q_ 'NN)
                (list 'AND (list 'NOT (list '= 'q_ 0))
                      (list '= (list '* pv pv) (list '* 3 (list '* 'q_ 'q_))))))))
\end{VNBcode}

Passing the formula literally would work just as well; the helper only spares the
author from writing it twice.  In general,
\[
  \texttt{(minimize! '($v_1$ \ldots\ $v_k$) GUARD MEASURE)}
\]
means ``choose $v_1,\ldots,v_k$ satisfying the formula \texttt{GUARD} with the
$\mathtt{nn}$-valued term \texttt{MEASURE} as small as possible.''  On the main
branch it lands two hypotheses, for fresh eigenconstants $w_1,\ldots,w_k$: that
\texttt{GUARD} holds of them, and that nothing satisfying \texttt{GUARD} has a
strictly smaller \texttt{MEASURE}.  It leaves exactly the two obligations any
minimisation owes --- that \texttt{MEASURE} lands in $\mathtt{nn}$ wherever
\texttt{GUARD} holds, and that \texttt{GUARD} holds somewhere.  Everything
between --- forming the value set, well-ordering it, unpacking the least
witness, restating minimality over the $v_i$ rather than the set --- it does for
you.  Every $v_i$ must occur in \texttt{MEASURE} (a variable the measure ignores
is not one you are minimising over).

It adds no kernel rule --- it chains \texttt{cut}, \texttt{di}, \texttt{ai},
\texttt{ew}, \ldots --- and uses \emph{no choice principle}: its one appeal is
\texttt{nn-least-element}, the well-ordering of $\mathtt{nn}$, and well-ordering
returns a genuine \emph{member} of the value set, from which the witness is read
back.  Stated instead as a first-order lemma, the same fact would force every
call site to tuple its variables into a Cartesian product, build a lambda, and
prove that lambda inhabits $\mathtt{fun}(U,\mathtt{nn})$ --- a typing obligation
larger than the minimisation it buys.  As a tactic it quantifies over the $v_i$
at the meta level and never forms $U$ at all.
\end{remark}

\begin{remark}[\texttt{obtain} versus \texttt{vlet}]
\texttt{obtain} and the \texttt{choice} former of \texttt{vlet} are the same
mechanism --- eliminate an existential and read its witness off by
free-variable difference --- but they differ in what they take.  \texttt{obtain}
takes a \emph{lane}: it runs a forward step and names the witness that step
\emph{just produced}.  That is the right fit inside a descent, where the relevant
existential is freshly landed and typically sits beside others: at the
``$3 \mid w$'' step the conclusion $\exists k.\,w = 3k$ coexists in context with
the evidence $\exists m.\,w\cdot w = 3m$ cut in to derive it, and \texttt{obtain}
selects the new one automatically.  Concretely, that one step reads
\begin{VNBsnip}
;; as written in the proof:
(define k (obtain (lambda () (fact 'nn-3-div-square w))))
;; the vlet alternative -- land the existential, then name it by its body:
(fact 'nn-3-div-square w)
(define k (vlet--choice! 'k '(AND (IN k NN) (= w (* 3 k)))))
\end{VNBsnip}
The \texttt{vlet} body names the eigenconstant $w$, a value produced at run
time, so this is the procedure \texttt{vlet--choice!}, not the
\texttt{(vlet \ldots)} macro (whose body is a fixed literal pattern).
\texttt{vlet} instead destructures a
\emph{named, stable} formula: its \texttt{match} former binds the holes of a
pattern (``the $a$ and $b$ in $x\cdot b = a$''), and its \texttt{choice} former
eliminates a \emph{designated} existential.  The rule of thumb: \texttt{obtain}
to name whatever a forward step produced, \texttt{vlet} to take apart a formula
you can already name.
\end{remark}

\begin{remark}
The proof adds no asserted step of its own: the mod-$3$ ladder is proven and the
descent cites it, exactly as \texttt{sqrt2-proof.scm} cites the parity ladder.
A new irrationality --- $\sqrt r$ for any non-square $r$ --- is then a matter of
widening the residue ladder, not of inventing an argument.  That division of
labour is the lesson: the author states the intermediate claims, and the machine
supplies the least witness (\texttt{minimize!}), the fresh names
(\texttt{obtain}), and the bookkeeping (\texttt{have!}, \texttt{from-context!}).
\end{remark}

\paragraph{The same proof at the interactive surface.}
By the transcript principle (Section~\ref{sec:two-surfaces}), most of this
script is \emph{already} a sequence of interactive short forms: \texttt{di},
\texttt{ai}, \texttt{ew}, \texttt{fact}, \texttt{cut}, \texttt{subst},
\texttt{crs}, \texttt{ass}, and \texttt{minimize!} are exactly what a user
types.  What the driver adds is the orchestration a human otherwise does by eye,
and it is of two kinds.  \emph{Focus}: after a branching command one selects the
next open leaf from the frontier, where the script calls \texttt{dk-focus!}.
\emph{Names}: each \texttt{(ai\ (forsome \ldots))} skolemizes to a machine-chosen
eigenconstant, printed in the resulting sequent, which the user reads and types
into the following commands.  The $3\mid w$ step, done by hand, is
\begin{VNBsnip}
(fact 'nn-3-div-square w)                          ; lands  exists k. w = 3k
(ai '(FORSOME k (AND (IN k NN) (= w (* 3 k)))))    ; skolemize; sequent now shows k_899
;; ... read k_899 off the sequent; type it into the subsequent (subst ...) / (fact ...)
\end{VNBsnip}
\texttt{obtain} is exactly this pair automated --- run the forward step, read the
fresh name back by free-variable difference --- so that a \emph{script}, unlike a
transcript with \texttt{k\_899} hard-coded, replays unchanged when the
eigenvariable counter lands on a different number.

%% =========================================================


%% =========================================================
\chapter{Provenance}
\label{ch:provenance}
\index{provenance}

A theorem in a running VNB was either proved in that session or
installed from a record of an earlier proof.  Either way it may rest on
supports that were asserted rather than proved.  And the copy of VNB
that holds it was obtained from somewhere, and may have been changed
since.  This chapter describes how each of these is recorded and how it
is checked:
\begin{itemize}
\item the proof certificate, which lets a load install a theorem without
  running its proof (Section~\ref{sec:proof-certificate});
\item the bill, which names the asserted facts a theorem rests on, and
  the stamp of the VNB test (Sections~\ref{sec:bill}
  and~\ref{sec:test-stamp});
\item signed trees, by which a user records that a copy was checked and
  detects any later change to it (Section~\ref{sec:signed-trees});
\item signed releases, by which the author records the files of a
  published version (Section~\ref{sec:signed-releases}).
\end{itemize}
The last two concern the copy of VNB rather than a single theorem: a
certificate is a plain file of the copy, so what a load accepts as
proven is only as sound as the copy it is read from.  This chapter says
what the two signatures are and what they protect.  The commands that
make and check them are in Chapter~\ref{ch:checking}
(Sections~\ref{sec:signed-trees-use} and~\ref{sec:signed-releases-use}).
The word is used in the sense it has for a single theorem, whose
provenance (\texttt{proven}, \texttt{certified}, \texttt{asserted},
\ldots) says where it came from: here the question is asked of the
library and of the copy as well.

\section{The proof certificate}\label{sec:proof-certificate}
\index{certificate!proof}\index{VNB\_CERTIFIED@\texttt{VNB\_CERTIFIED}}
\index{exam}

\begin{description}
\item[Certifies] that the proof of a named theorem ran to \texttt{qed} and
  passed the inference check, on a stated kernel, and what that proof
  depended on.  It is the mechanism by which the
  everyday load installs a theorem without running its proof.
\item[Issued by] a load that ran the proof: the \emph{exam},
  \texttt{VNB\_CERTIFIED=off ./prover --build-band}, which runs every proof
  and rewrites the store; or an ordinary load meeting a theorem with no
  valid record, which proves it and writes one.
\item[Form] one file per proof file, \texttt{certificates/}$\langle$load
  key$\rangle$\texttt{.cert}, one record per \texttt{qed} in installation
  order.  A record holds the theorem's name, the hash of its statement and
  the statement itself; \texttt{cites}, the names the ledger read off the
  script, each with the hash of its installed statement; \texttt{defs},
  the closure of the statement's defined heads through their definitions
  (the bodies of \texttt{def-functoid}s, the axioms of \texttt{def-predicate}s
  and \texttt{def-constant}s; proven theorems are not followed), with their hashes;
  \texttt{uses}, every macete the kernel saw fire in the proof (an accessor
  read by \texttt{slot} among them), and every \texttt{def-functoid} body and
  predicate defining iff a checker read while accepting an inference, with
  the hash of the rewrite rule or definition, the kernel's own list kept
  independent of the ledger's; the oracles the script called and the bill and oracles
  \texttt{qed} announced; and the kernel hash.  A record carries no date and no host.  A hash is
  the MD5 of the canonical text of a formula, with counter-minted names
  relabelled in order of first occurrence so that it does not depend on
  how much was proved before.  The \emph{kernel hash} is the MD5 of the
  source of the files in \texttt{*cert-kernel-files*}
  (\texttt{certificates.scm}): the files that may call
  \texttt{dg-apply-rule!}, the deduction-graph file, \texttt{expressions},
  \texttt{sequents} and the four checker files; tactics, the kit and
  the drivers are not in it, because every inference they produce is
  checked.
\item[Read by] \texttt{sp}, at load, under \texttt{VNB\_CERTIFIED=on}
  (the default) or \texttt{strict}.  A record is valid when the statement
  being started hashes to its statement hash, every cited, defined and
  used name hashes the same in this tree, and the kernel hash is this
  tree's.  On a valid record the theorem is installed with provenance
  \texttt{certified}, its bill recomputed from the recorded citations, and
  the driver's steps are skipped.  Under \texttt{on} a theorem without a
  valid record is proved and its record written; under \texttt{strict}
  nothing is proved, the uncovered theorems are listed and the load exits
  non-zero.  The bill line of such a theorem reads \texttt{certified modulo
  \{\ldots\}} (a record carries no date and no host, so an exam that changes
  nothing rewrites nothing); the load summary counts \texttt{proven} and
  \texttt{certified} separately, and the word \texttt{proven} alone means
  the proof ran in this image.  The page audit covers the proofs that ran.
\item[Invalidated by] a change to the statement, to the statement of a
  cited name, to a definition the statement mentions or a macete the proof
  fired, or to any kernel file, which invalidates every record at once.  A
  record without a \texttt{uses} field is not valid.
\item[Does not] certify what the record cannot see: a functoid reduced by
  \texttt{beta} needs no entry (a \texttt{beta} contracts only a lambda record
  that reached the goal through the statement, a cited theorem or a named
  unfold, each of which is keyed), and the exam is the backstop for what
  the kernel hash does not cover: the tactics and the kit.  A certificate is not a proof:
  it says a proof ran and was checked.  The design is
  \texttt{docs/certificates-2026-09-24.md}.
\end{description}

\section{The bill, and \texttt{proof-rules}}\label{sec:bill}
\index{bill}\index{proof-rules@\texttt{proof-rules}}

Two further things are certified about a completed proof, neither called
a certificate in the code.  The \emph{bill} that \texttt{qed} prints,
\texttt{proven modulo \{\ldots\} [trust: \ldots]}, is the kernel's
certification of the theorem relative to the asserted facts it names
(Section~\ref{sec:proof-debt}): the theorem holds if those do.  The
command \texttt{(proof-rules)} walks the grounded derivation of the
current proof from its root and tallies every kernel inference by its
rule tag; because every tag is one of the fixed set of kernel operations,
this shows that a proof driven by composite tactics rests on nothing but
those operations, independently of the command trace, which is lossy.

\section{The certification stamp of the VNB test}\label{sec:test-stamp}
\index{certificate!VNB test stamp}

The classification pages of Chapter~\ref{ch:classification} report a
result as ``certified by the VNB test on \textit{date}''.  The stamp is
written to the file \texttt{reference/certification.scm} by
\texttt{vnb-test}, one
\texttt{(certify!\ '\textit{name} "\textit{date}")} per theorem the test
re-proved to \texttt{qed}, and read back at the next load.  It certifies a
run of the test, not a proof.  The everyday mechanism is the proof
certificate above: the exam that writes the store is the run that
re-proves the library.

%% =========================================================
\section{Signed trees}
\label{sec:signed-trees}
\index{signed tree}\index{vnb-trust@\texttt{vnb-trust}}\index{vnb-start@\texttt{vnb-start}}

This section says what a signed tree is and what it protects.  The
commands, \texttt{vnb-trust} and \texttt{vnb-start}, the contents of the
signed list, and the procedure for work that must withstand an
adversary are in Section~\ref{sec:signed-trees-use}.

A user obtains VNB as a copy of its repository: a clone from GitHub,
or an archive of it.  That copy is one directory,
\texttt{\textasciitilde/prover}, called the \emph{tree}.  It holds
\begin{itemize}
\item the Scheme source of the kernel, the tactics and the interface,
  and the Emacs files of the workspace;
\item the launchers \texttt{VNB} and \texttt{prover}, which are short
  scripts;
\item the libraries of definitions (\texttt{structure-library/}) and of
  proofs (\texttt{theorem-library/}), and one proof certificate per
  proof file (\texttt{certificates/});
\item the documentation (\texttt{docs/}, \texttt{reference/});
\item the compiled code (\texttt{.com} files) and the \emph{band}, a
  saved image of the loaded library from which VNB starts.  However,
  the repository does not contain them: a user who clones it makes them on
  the first run.  A copy taken from a machine that has built them
  brings them along.
\end{itemize}
There is no stand-alone executable in the tree.  The program that runs
it is MIT Scheme, with GNU Emacs for the workspace; both are installed
on the computer separately and are not part of the tree.  But the
compiled files and the band are machine code that MIT Scheme runs, so
they are as much the program as the source is, and the signature
covers them.

VNB does not rerun every proof each time it starts.  It installs most
theorems from their proof certificates
(Section~\ref{sec:proof-certificate}), which are plain files in the
tree.  So a program, a person, or an accident that changes files in
the tree can change what VNB accepts as proven, and nothing at
start-up would show it.

A \emph{signed tree} makes such a change visible.  Once the user is
satisfied that a copy of VNB is a good one, the user \emph{accepts}
it: the hash of every file of the tree, the documents aside, is
recorded in a list, and the list is signed with the user's own key.  From then on, a copy is
compared with the lists the user has signed before it is used.  The
picture is
\[
  \begin{array}{c}
  \text{you obtain a copy}\\
  \downarrow\\
  \text{you check the copy}\\
  \downarrow\\
  \text{you sign its hashes}\\
  \downarrow\\
  \text{every later use is compared with them.}
  \end{array}
\]

\paragraph{How a user checks that a copy is good.}  Checking is up to
the user.  Signing is not a check: it only records that the user made
one and was satisfied.  Three ways of checking, from the lightest to
the most thorough:
\begin{itemize}
\item checking the signature of a published release
  (Section~\ref{sec:signed-releases}): the copy is then the one its
  author published;
\item reading what changed since the last accepted copy, which the
  comparison lists file by file;
\item rerunning proofs.  A build with \texttt{VNB\_CERTIFIED=off}
  ignores the certificates and runs every proof through the kernel; it
  is called the \emph{exam}, and takes hours for the whole library.  A
  user can also rerun only some files, for instance those that
  changed.
\end{itemize}
Only the third establishes anything about the proofs themselves.  The
first two establish where the files came from.

\paragraph{What it protects, and what it does not.}  The private key
exists only on the user's computer.  Nothing elsewhere can sign a tree
for the user.  So
\begin{itemize}
\item a copy that differs from every tree the user accepted is refused
  at the next \texttt{vnb-start}, with the differing files named;
\item a proof certificate (Section~\ref{sec:proof-certificate}) is not
  signed, and everything in it can be computed from the tree.  Anyone
  who can write the files of a tree can therefore add a statement to a
  library file and write a certificate for it that the certified load
  accepts.  The statement is then installed without its proof, with the
  provenance \texttt{certified} and an empty bill, and a theorem proved
  from it looks proved.  The signature detects the changed files at the
  next \texttt{vnb-start}; the exam, which ignores the certificates and
  runs every proof, fails at the statement, whose proof does not
  prove it;
\item VNB started without \texttt{vnb-start}, for instance from Emacs
  directly or by \texttt{\textasciitilde/prover/VNB}, is not checked,
  and neither is a tree accepted without a check: the certificate above
  passes both;
\item the check is made when VNB starts.  A file that VNB reads later
  in a session, for instance a library file loaded again, is not
  checked again;
\item MIT Scheme and Emacs are outside the tree and are not checked;
  they are as trustworthy as the way they were installed;
\item the web browser is outside the tree too.  VNB opens its home page
  and the reference pages in it.  The pages are documents and the
  signature does not cover them, so a changed page could show a wrong
  statement; the library that VNB loads is checked, the pages that
  describe it are not;
\item while VNB runs, Emacs answers requests from the home page at a
  local address (\texttt{127.0.0.1}, port 8973).  It carries out only a
  fixed list of nine commands, such as starting a proof or opening the
  calculator, and never runs anything it is sent;
\item nothing here defends a computer that an intruder controls: such
  an intruder can change \texttt{vnb-start}, \texttt{vnb-trust} or
  \texttt{ssh-keygen} itself.
\end{itemize}
The signature says ``I accepted this tree.''  It is not a proof, and it
shows no more about the proofs than the check that preceded it.


%% =========================================================
\section{Signed releases}
\label{sec:signed-releases}
\index{signed release}

This section says what a signed release is and what it shows.  How the
author makes one and how anyone checks one are in
Section~\ref{sec:signed-releases-use}.

A signed tree (Section~\ref{sec:signed-trees}) records a check made by
the user.  A \emph{signed release} records one made by the author: it
lets anyone who obtains a copy of VNB check that the files are the
ones the author published.

\paragraph{What a release is.}  A release has a version,
\texttt{v}\textit{YYYY.MM.DD}, and two files committed under
\texttt{releases/}:
\begin{itemize}
\item the \emph{release manifest}: the commit, then the SHA-256 hash of
  every file git tracks at that commit, one line per file, in the order
  \texttt{git ls-files} gives;
\item the author's SSH signature of the manifest, in the namespace
  \texttt{vnb-release}.
\end{itemize}

\paragraph{The author's key.}  The author's public key is published in two places, under two
different accounts, so that taking over one account is not enough to
replace it:
\begin{itemize}
\item on GitHub, in the author's list of signing keys
  (\url{https://api.github.com/users/JavierThayer/ssh_signing_keys});
\item away from GitHub, in the author's Bluesky profile
  (\url{https://bsky.app/profile/javiertf.bsky.social}): its
  fingerprint, as text in the profile description and in the banner
  image.
\end{itemize}
The key is an ed25519 key, added 2026-11-17, with fingerprint
\begin{center}
\texttt{SHA256:pvv9BBY6uAVZrlz0H7G/xub84NareBQos+3Af8jHTvk}
\end{center}
The README and this manual print it too, for convenience.  They are
not further places to check it, since they come from the same
repository as the release.

\paragraph{What it shows, and what it does not.}  A valid signature
protects a user who obtains a copy of the release from
\begin{itemize}
\item a doctored copy: a fake fork, a tampered mirror, a file changed in
  transit;
\item a forged release, since the private key exists only on the
  author's computer;
\item a planted proof certificate, since certificates are files of the
  release.
\end{itemize}
It does not protect from
\begin{itemize}
\item a change after the download, unless the check is made again;
  \texttt{vnb-start}, used with the downloader's own key, makes it at
  every start;
\item a computer that an intruder already controls;
\item a stolen signing key, until the author withdraws it;
\item a takeover of both of the author's accounts, GitHub and Bluesky,
  together.  A takeover of the GitHub account alone can replace the key
  on GitHub, in the README and in this manual at once; the Bluesky
  copy is what exposes it, which is why it must be compared too.
\end{itemize}

Nor does a valid signature show that the code is harmless: VNB is a
program, and running it means trusting its author or reading it.  And
it shows nothing about the proofs.  For those, no trust in the author
is needed: the exam, run by the user, reruns every proof through the
kernel (Section~\ref{sec:signed-trees-use}, ``For work that must
withstand an adversary'').


%% =========================================================
\chapter{Definitions}
\label{ch:defs}

The forms described in this chapter are \emph{Scheme procedures} called
at the meta-level --- in a \texttt{.scm} source file or at the Scheme
REPL --- not expressions in the VNB object language.  Each call installs
axioms or macetes as a side effect; nothing is
proved.  The distinction between meta-level definition and object-level
proof bears repeating: \texttt{declare-structure}, \texttt{def-constant},
and their companions manipulate the library record directly, while
proofs are constructed separately using the proof commands of
Chapter~\ref{ch:proofs}.

%% =========================================================
\section{Constants and predicates}
\label{sec:def-constant}

A \emph{defined constant} is a new constant symbol introduced by a
collection of characterizing axioms.  It differs from a plain
\texttt{add-axiom!} call only in bookkeeping: the axioms are
recorded under a separate \texttt{definitions} slot in the library record,
so the system can distinguish what was \emph{postulated} from what was
\emph{introduced by definition}.

\begin{VNBcode}
(def-constant 'NAME
  '(axiom-name-1 formula-1)
  '(axiom-name-2 formula-2)
  ...)
\end{VNBcode}

\begin{sloppypar}Each extra argument is a two-element list \texttt{(axiom-name formula)}.
The formulas are installed in the global theorem and macete tables exactly
as \texttt{add-axiom!} would install them, so they are immediately
available to proof commands and macetes.\end{sloppypar}

\textbf{Existence and uniqueness are the user's responsibility.}  The
system accepts the definition without checking that the described object
exists or is unique.  For constants where this matters (e.g., a least
element), you should prove existence and uniqueness before invoking
\texttt{def-constant}, or at least note the obligation explicitly.

Example---introducing the first uncountable ordinal $\omega_1$:
\begin{VNBcode}
(def-constant 'OMEGA1
  (list 'omega1-is-ordinal
        (parse-string "OMEGA1 in ORD"))
  (list 'omega1-is-uncountable
        (parse-string "not(countable(OMEGA1))"))
  (list 'omega1-is-least-unc
        (parse-string "forall([alpha in ORD],
          not(countable(alpha)) implies OMEGA1 <= alpha)")))
\end{VNBcode}

To display all defined constants and their characterizing axioms:
\begin{VNBcode}
(display-definitions)
\end{VNBcode}

%% =========================================================
\section{Predicates}
\label{sec:def-predicate}

An $n$-ary \emph{defined predicate} is introduced by \texttt{def-predicate}.
It installs a single characterising axiom --- an \texttt{iff} --- that
unfolds the predicate in terms of its arguments.

\begin{VNBcode}
(def-predicate 'PRED-NAME '(p1 p2 ...)  body)
\end{VNBcode}

Installs one axiom named \texttt{PRED-NAME}:
\begin{VNBsnip}
forall([p1, p2, ...], PRED-NAME(p1, p2, ...) iff body)
\end{VNBsnip}

The axiom is immediately available via \texttt{(ta 'PRED-NAME)} in a proof.

\paragraph{Example: Cauchy sequence in a metric space.}
\begin{VNBcode}
(def-predicate 'IS-CAUCHY-SEQ '(s f)
  '(AND (IS-METRIC-SPACE s)
   (AND (IN f (FUN NN (PTS s)))
        (FORALL eps (IMPLIES (POS-RR eps)
          (FORSOME N (AND (IN N NN)
            (FORALL m (IMPLIES (IN m NN)
              (FORALL n_ (IMPLIES (IN n_ NN)
                (IMPLIES (AND (<= N m) (<= N n_))
                  (<= ((DIST s) (f m) (f n_)) eps)))))))))))))
\end{VNBcode}

\noindent
This is the predicate shipped in
\texttt{structure-library/metric-completeness.scm}; \texttt{pos-rr}
(``$r \in \mathtt{rr}$ and $r > 0$'') is the convenience predicate introduced
next.  It installs:
\begin{VNBsnip}
forall([s, f],
  IS-CAUCHY-SEQ(s, f) iff
    IS-METRIC-SPACE(s)
    and f in fun(NN, PTS(s))
    and forall([eps], POS-RR(eps) implies
          forsome([N in NN],
            forall([m in NN, n_ in NN], N <= m and N <= n_ implies
              DIST(s)(f(m), f(n_)) <= eps)))
\end{VNBsnip}

\paragraph{Order predicates: \texttt{<} and \texttt{pos-rr}.}  Two
pre-installed convenience predicates relieve the two most pervasive
analytic idioms.  The kernel has \texttt{<=} as a primitive predicate
but no strict \texttt{<}; ``$r > 0$ for $r \in \mathtt{rr}$'' similarly
expands to a three-conjunct mouthful.

\begin{VNBcode}
(def-predicate '< '(x y)
  '(AND (<= x y) (NOT (= x y))))

(def-predicate 'POS-RR '(r)
  '(AND (IN r RR) (<= 0 r) (NOT (= 0 r))))
\end{VNBcode}

Installed in \texttt{structure-library/order-predicates.scm} (loads
right after \texttt{number-systems}).  Each is just a defined-by-iff
predicate; ground reasoning still unfolds through the iff (arith-eval
has not been extended to dispatch on \texttt{<} directly).

\begin{remark}
Opt-in.  Existing statements (the metric-topology
vocabulary) retain the verbose forms for backward consistency; new
axioms going forward should prefer \texttt{<} and \texttt{pos-rr}.
\end{remark}

%% =========================================================
\section{Recursive definitions}
\label{sec:recursive-defs}

Two Scheme helpers introduce function constants by recursion and install
the defining equations as axioms.  Both rely on the same
\texttt{add-definition!} mechanism as \texttt{def-constant}
(Section~\ref{sec:def-constant}), so every axiom is immediately
available in proofs.

\begin{prop}[Primitive recursion on $\mathtt{nn}$]\label{prop:nn-rec}
Let $A \in \mathtt{set}$ and let
\[
  \mathcal{F}(A) \;=\; \bigcup_{n \in \mathtt{nn}}\,\mathtt{fun}(\mathit{ord\text{-}segment}(n),\, A)
\]
be the set of all functions defined on a finite initial segment of $\mathtt{nn}$
with values in $A$ (including the empty function at $n = 0$, since
$\mathit{ord\text{-}segment}(0) = \emptyset$).
For any $\mathcal{E} \in \mathtt{fun}(\mathcal{F}(A),\, A)$ there is a unique
$f \in \mathtt{fun}(\mathtt{nn},\, A)$ such that
\[
  f(n) \;=\; \mathcal{E}\!\bigl(f \!\restriction\! \mathit{ord\text{-}segment}(n)\bigr)
  \qquad \text{for every } n \in \mathtt{nn}.
\]
\end{prop}

\begin{remark}
This is the special case of Proposition~\ref{prop:tfrec-set} with
$U = \mathtt{nn}$; it follows from the well-ordering of $\mathtt{nn}$ by the
standard least-counterexample argument.
Since $\mathit{ord\text{-}segment}(n)$ is the integer interval $\{0, 1, \ldots, n-1\}$,
the extension function $\mathcal{E}$ sees the entire history
$f(0), f(1), \ldots, f(n-1)$ when computing $f(n)$.
The \emph{predecessor form} --- $f(0) = b$ and
$f(\mathrm{succ}(n)) = s(n,\, f(n))$ for $n \in \mathtt{nn}$ --- is the special case
where $\mathcal{E}$ only inspects the last value; this covers sums, products,
and most concrete recursive definitions.
\end{remark}

\subsection{\texttt{def-by-nn-recursion}: primitive recursion on \texttt{nn}}

Defines a function $F(p_1, \ldots, p_k, n)$ by two-clause recursion on the
last argument $n \in \mathtt{nn}$, with arbitrary extra parameters
$p_1, \ldots, p_k$:
\[
  F(p_1,\ldots,p_k,0) = b, \qquad
  F(p_1,\ldots,p_k,\mathrm{succ}(n)) = s(p_1,\ldots,p_k,n,F(p_1,\ldots,p_k,n)).
\]

\begin{VNBcode}
(def-by-nn-recursion 'F
  '(p1 p2 ...)      ; extra parameters ('() for a plain NN->A function)
  base-val          ; F(p1,p2,..., 0) = base-val
  '(n val)          ; step variables
  succ-expr)        ; F(p1,p2,..., succ(n)) = succ-expr
                    ;   (val stands for F(p1,p2,...,n))
\end{VNBcode}

Two definitional axioms are installed, as quasi-equations (\texttt{==}:
both sides undefined, or both defined and equal), with the step guarded by
$n \in \mathtt{nn}$:
\begin{VNBsnip}
F-zero: forall([p1,p2,...], F(p1,p2,...,0) == base-val)
F-succ: forall([p1,p2,..., n in NN],
          F(p1,p2,...,succ(n)) == succ-expr[val:=F(p1,p2,...,n)])
\end{VNBsnip}

\paragraph{Example: factorial (no parameters).}
\begin{VNBcode}
(def-by-nn-recursion 'FACT '()
  1                           ; FACT(0) = 1
  '(n val) '(* (succ n) val)) ; FACT(succ(n)) = succ(n) * FACT(n)
\end{VNBcode}

\paragraph{Example: power of 2 (no parameters).}
\begin{VNBcode}
(def-by-nn-recursion 'POW2 '()
  1                       ; POW2(0) = 1
  '(n val) '(* 2 val))   ; POW2(succ(n)) = 2 * POW2(n)
\end{VNBcode}

\paragraph{Example: partial sum (one parameter).}
\begin{sloppypar}For $f \in \mathtt{fun}(\mathtt{nn}, A)$, define
$\mathrm{sum}(f, n) = \sum_{k=0}^{n} f(k)$:\end{sloppypar}
\begin{VNBcode}
(def-by-nn-recursion 'SUM '(f)
  '(f 0)                          ; SUM(f, 0) = f(0)
  '(n val) '(+ val (f (succ n)))) ; SUM(f, succ(n)) = SUM(f,n) + f(succ(n))

; Installs SUM-zero: forall([f], SUM(f,0) = f(0))
; Installs SUM-succ: forall([f, n in NN], SUM(f,succ(n)) = SUM(f,n) + f(succ(n)))
\end{VNBcode}

\subsection{\texttt{def-by-ord-recursion}: transfinite recursion on \texttt{ord}}

For functions on all ordinals, use the three-clause schema
(Proposition~\ref{prop:tfrec-class}, Section~\ref{sec:ord-recursion}):
\begin{itemize}
\item \textbf{Base} ($\xi = 0$):\quad $F(0) = b$
\item \textbf{Successor} ($\xi = \mathrm{succ\_ord}(\alpha)$):\quad
      $F(\mathrm{succ\_ord}(\alpha)) = s(\alpha,\, F(\alpha))$
\item \textbf{Limit} ($\mathrm{limit\text{-}ord}(\xi)$):\quad
      $F(\xi) = \ell\!\left(\xi,\, F\!\restriction\!\mathrm{ord\text{-}segment}(\xi)\right)$
\end{itemize}

\begin{VNBcode}
(def-by-ord-recursion 'F
  '(p1 p2 ...)        ; extra parameters ('() for none)
  base-val            ; F(p1,p2,..., 0) = base-val
  '(alpha val)        ; successor: F(p..., succ_ORD(alpha)) = succ-expr
  succ-expr           ;   where val stands for F(p..., alpha)
  '(lam)              ; limit: F(p..., lam) = lim-expr
  lim-expr)           ;   where lam is the limit ordinal
\end{VNBcode}

Three definitional axioms are installed, as quasi-equations (\texttt{==}),
guarded by \texttt{ord} and \texttt{limit-ord}, with the parameters
universally quantified outside (\texttt{structure-library/ordinals.scm}):
\begin{VNBsnip}
F-zero:  F(0) == base-val
F-succ:  forall([alpha in ORD],
           F(succ_ORD(alpha)) == succ-expr[val:=F(alpha)])
F-limit: forall([lambda], LIMIT-ORD(lambda) implies F(lambda) == lim-expr)
\end{VNBsnip}

\paragraph{Example: identity function on ordinals.}
\begin{VNBcode}
(def-by-ord-recursion 'COPY-ORD '()
  0                                         ; COPY-ORD(0) = 0
  '(alpha val) '(succ_ORD val)              ; COPY-ORD(succ(a)) = succ(COPY-ORD(a))
  '(lam) '(SUP-ORD (ORD-SEGMENT lam)))      ; COPY-ORD(lam) = sup of predecessors
\end{VNBcode}

\paragraph{Example: ordinal addition $\alpha + \beta$ (fixing $\alpha$).}

Defining $A_\alpha(\beta) = \alpha + \beta$ as a function of $\beta$:
\begin{VNBcode}
(def-by-ord-recursion 'ORD-ADD '(a_fix)     ; the fixed addend is a parameter
  'a_fix                                    ; ORD-ADD(a_fix, 0) = a_fix
  '(beta val) '(succ_ORD val)              ; ORD-ADD(succ(b)) = succ(ORD-ADD(b))
  '(lam) '(SUP-ORD (ORD-SEGMENT lam)))     ; ORD-ADD(lam) = sup of ORD-ADD on lam
\end{VNBcode}

(Here \texttt{alpha-fixed} is the constant playing the role of $\alpha$.
A full definition would use nested ordinal recursion.)

%% =========================================================
\section{Mathematical structures}
\label{sec:declare-structure}

\subsection{The \texttt{declare-structure} form}

\texttt{declare-structure} declares a named mathematical structure at the
meta-level --- in the Scheme REPL or a \texttt{.scm} source file, as the
chapter opening explains, not inside the object language.  It registers
the structure, installs \texttt{nth}-based accessor
macetes for every carrier set and every operation, and posts a
definitional axiom for the associated \texttt{is-\textit{name}} predicate.
No quoting is needed.

\begin{VNBcode}
(declare-structure NAME
  (carrier C)
  (carriers C1 C2 ...)
  (op OPNAME DOMAIN RANGE)
  (constant CNAME SET)
  (substructure SNAME TYPE)
  (derived DNAME BASE EXPR)
  (family FNAME CARRIER)
  (family-fun MNAME FAMILY RANGE)
  (relation RNAME CARRIER)
  (property PRED ACC ...)
  (law "FORMULA"))
\end{VNBcode}

The clauses other than \texttt{property} and \texttt{law} describe the
structure's \emph{shape} --- its components and the role each plays.  Every
shape clause may be repeated, once per component of that kind (two operations
are two \texttt{op} clauses); each occurrence declares one component (a
\texttt{carriers} clause declares one per name), and contributes one typing conjunct to the
characteristic law (two for \texttt{derived}: a typing and the equation that
fixes the component).  The \texttt{property} and \texttt{law} clauses are
different in kind: each names a \emph{characteristic law} the components must
satisfy.  The distinction is developed in the remark on shape versus properties
below.

\begin{itemize}

\item \texttt{(carrier C)}
  Declares one carrier set: $\mathtt{C}(s) \in \textsc{set}$.  The clause
  takes exactly one name; \texttt{(carrier C D)} is an error.  Like every
  other shape clause it may occur any number of times and at any position.

\item \texttt{(carriers C1 C2 ...)}
  Shorthand for \texttt{(carrier C1) (carrier C2) \ldots}: it declares
  exactly the same components, in the same slots.  It is the only clause that
  declares several components at once.

\item \textbf{Slot numbering.}  An instance \texttt{s} of the structure is a
  list with one entry per component.  The components are numbered
  $1, 2, \ldots, n$ in the order in which their clauses are \emph{written} in
  the declaration, read left to right, whatever their kinds.  The order in which
  the clause kinds are listed above plays no role, and the declaration may
  place the clauses in any order.  A \texttt{carrier}, \texttt{op},
  \texttt{constant}, \texttt{substructure}, \texttt{derived}, \texttt{family},
  \texttt{family-fun} or \texttt{relation} clause takes the next number;
  \texttt{carriers} takes one number per name, left to right;
  \texttt{property} and \texttt{law} take none.  For example, in
\begin{VNBcode}
(declare-structure POINTED-SET
  (constant BASEPT PSET)
  (carrier PSET))
\end{VNBcode}
  there are two components: the carrier \texttt{PSET}, declared by the second
  clause, and the constant \texttt{BASEPT}, declared by the first, whose clause
  says that it is an element of the carrier, $\mathtt{BASEPT}(s) \in
  \mathtt{PSET}(s)$.  The name \texttt{PSET} occurs in both clauses but denotes
  one component.  \texttt{BASEPT} is slot~1 and \texttt{PSET} is slot~2,
  because the constant's clause is written first.  The accessor of the
  $k$-th component is the macete $\mathtt{ACC}(s) \to \mathtt{nth}(k, s)$.
  Hence $\mathtt{nth}(i, s)$ is the $i$-th carrier only when the carriers
  are declared before every other component.  The slot numbers matter only to
  the accessor macetes and to an instance written out as a literal list.  A
  statement or proof refers to a component by its accessor
  (\texttt{PSET(s)}), never by $\mathtt{nth}(k, s)$.  In \texttt{FIELD}
  (\texttt{structure-library/field.scm}) the carrier \texttt{CARR} is slot~1,
  the operations and constants occupy slots~2--6, and the derived carrier
  \texttt{NON-ZERO} is slot~7.  Two structures may share an accessor name only
  when it occupies the same slot number in both, since the accessor macete is
  global.  A declaration that gives an existing accessor name another slot is
  refused with an error naming the structure and the slot that already hold it,
  and nothing is declared; a declaration that names one component twice is
  refused the same way.

\item \texttt{(op NAME DOMAIN RANGE)}
  Declares a function-valued slot: $\mathtt{NAME}(s) \in
  \mathtt{fun}(\mathtt{DOMAIN}, \mathtt{RANGE})$.  Every VNB function is
  unary on its domain, so binary and higher-arity operations use an explicit
  Cartesian product: write \texttt{(op mul (cartesian A A) A)}, not
  \texttt{(op mul (carr A) A)}.  All carrier and op names in domain and range
  positions are automatically expanded to accessor applications.

  The unary typing is reconciled with the curried application form
  $((\mathtt{mul}\,s)\,a\,b)$ used throughout the library by the
  \texttt{apply-tupling-$n$} axioms
  (Section~\ref{sec:curried-tupled-apply}): the $n$-arg apply equals the
  unary apply to the $n$-tuple $(\texttt{list}\,a_1\,\ldots\,a_n)$.
  Together with \texttt{fun-apply-type} this closes
  $((\mathtt{mul}\,s)\,a\,b) \in \mathtt{carr}(s)$ from the typing of $\mathtt{mul}(s)$
  and $a,b \in \mathtt{carr}(s)$.

\item \texttt{(constant \textit{name} \textit{set})}
  Declares a constant-valued slot: $\textit{name}(s) \in \textit{set}(s)$.

\item \texttt{(substructure NAME TYPE)}
  Declares a slot that holds a whole structure of kind \texttt{TYPE}
  (for example a vector space's base \texttt{field}).  It is typed by
  \texttt{is-\textit{type}}$(\mathtt{NAME}(s))$ rather than by the bare
  $\in \textsc{set}$ of a carrier.  An op signature may reach into it
  with foreign accessors --- e.g.\
  \texttt{(op smul (cartesian (carr K) V) V)}, where \texttt{K} is a
  \texttt{(substructure K field)} slot and \texttt{(carr K)} is the base
  field's carrier.

\item \texttt{(derived NAME BASE EXPR)}
  Declares a carrier \emph{carved out} of another rather than an independent
  one.  \texttt{NAME(s)} still occupies its own slot, but \texttt{is-\textit{name}}
  pins its value with an equation, \texttt{NAME(s) = EXPR}, where \texttt{EXPR}
  is a term over the other accessors --- not the bare $\in\textsc{set}$ a plain
  carrier gets.  \texttt{field} declares
  \texttt{(derived non-zero carr (difference carr (singleton zero)))}: the unit
  group's carrier is the field carrier minus zero.  Because its value is fixed
  by the other slots, a derived carrier owes a homomorphism no map of its own;
  it rides the base carrier's (Section~\ref{sec:morphisms}).

\item \texttt{(family NAME CARRIER)}
  Declares a component that is a family of subsets of a carrier --- a
  $\sigma$-algebra, a topology.  Its typing conjunct is
  $\mathtt{NAME}(s) \in \mathtt{power}(\mathtt{power}(\mathtt{CARRIER}(s)))$.
  What distinguishes it from a \texttt{constant} of that type is the generated
  morphism: a homomorphism must pull members of the family back to members of the
  family (Section~\ref{sec:morphisms}).  \texttt{top-space} declares
  \texttt{(family opens pts)}.

\item \texttt{(family-fun NAME FAMILY RANGE)}
  Declares a function on a family component --- a measure.  Its typing conjunct
  is $\mathtt{NAME}(s) \in \mathtt{fun}(\mathtt{FAMILY}(s), \mathtt{RANGE})$; its
  morphism clause is transport along preimages.

\item \texttt{(relation NAME CARRIER)}
  Declares a component that is a relation on a carrier, that is, a set of pairs of
  its points.  Its typing conjunct is
  $\mathtt{NAME}(s) \in \mathtt{power}(\mathtt{CARRIER}(s) \times \mathtt{CARRIER}(s))$;
  a homomorphism must send related points to related points.  \texttt{setoid}
  declares \texttt{(relation rel pts)}.

\item \texttt{(property PRED ACC ...)}
  Names a characteristic law --- \emph{not} a component.  \texttt{PRED}
  is a property predicate (Section~\ref{sec:shape-vs-property}), applied
  to the named accessors, and is conjoined into the
  \texttt{is-\textit{name}} definition.  For example
  \texttt{(property is-metric D X)} folds the metric axioms on
  $\mathtt{dist}(s)$ over $\mathtt{pts}(s)$ into \texttt{is-metric-space}.

\item \texttt{(law "\textit{formula}")}
  Adds a raw conjunct to \texttt{is-\textit{name}}, written in the surface
  string syntax and parsed like any user formula (its free variable is the
  instance \texttt{s}).  It is the escape hatch for a characteristic law that no
  named \texttt{property} predicate captures --- \texttt{module}'s scalar-action
  axioms are laws of this kind.  A \texttt{property} clause is just the special
  case where the law already has a name in \texttt{operation-properties.scm}.

\end{itemize}

An accessor is a projection and nothing more (its index is the slot
number defined above): $\mathtt{dist}(s)$ is $\mathtt{nth}(2, s)$ for \emph{every}
object $s$, a ring included, and no contradiction can come of writing it,
since what the second component of a metric space \emph{is} lives in the
characteristic law $\mathtt{is\text{-}metric\text{-}space}(s)$ and not in
the accessor; the load-time \texttt{accessor-type-audit} then refuses any
library statement that applies an accessor to a variable guarded by a
structure without that slot.

Every slot name --- carrier, operation, constant, or substructure --- is
registered as a global \emph{accessor constant} of that name.  The
particular names are the declarer's to choose (\texttt{carr} for a ring's
carrier, \texttt{pts} for a metric space's points, \texttt{opr} for a
group's operation); nothing forces one spelling over another.  But once
declared, a name is \emph{reserved}: because the reader case-folds
(Remark~\ref{rem:case-fold}), using it as a bound variable anywhere ---
\texttt{forall([carr in \ldots], \ldots)} --- makes the binder read as the
accessor rather than the variable, and the load-time
\texttt{constant-binder-audit} rejects it outright.  This is why the
property-predicate parameters below (\texttt{op}, \texttt{crr}, \ldots)
are spelled to avoid the accessor constants they stand in for.

\begin{remark}[What the form permits, and what the idiom favours]
\texttt{declare-structure} is general.  The \texttt{carrier} clause may be
repeated, so a structure admits any number of carriers, and operations of arbitrary signature over them,
so an arbitrary finitary algebraic system --- many-sorted, in the sense of
Mal'cev's \emph{algebraic systems} --- can be declared.  A module, for instance,
has a ring of scalars (a \texttt{substructure}) beside its carrier of vectors.

What the library's \emph{idiom} favours is narrower than what the form
\emph{permits}.  Most examples --- groups, rings, fields --- have a single
distinguished carrier, and their homomorphisms are functions on it; the generated
morphism predicate (Section~\ref{sec:morphisms}) takes that as its default.  But
neither is imposed by the mechanism.  A structure may carry several independent
carriers, and its generated homomorphism is then one map per carrier; and
\texttt{declare-hom!} may replace the notion of morphism entirely with one that is
not a carrier-function at all --- continuity, for a topological space.  The
single-carrier reading, with homomorphisms as carrier maps, is a convention of the
material at hand, not a constraint the framework enforces.
\end{remark}

\begin{remark}
A \texttt{declare-structure} declaration involves three distinct things,
which it is worth keeping clearly separated.
\begin{enumerate}
\item \textbf{The name} --- a symbol such as \texttt{ring} or
  \texttt{metric-space}, used as a label in declarations and prose.
\item \textbf{An instance} --- a specific VNB list, which is an
  \emph{element of \textsc{set}} (a concrete set-theoretic object).
  For example,
  \[
    \mathtt{cc\text{-}ms} \;=\;
      [\,\mathtt{cc},\;\mathtt{vnb\text{-}lambda}([\mathtt{x},\mathtt{y}],\,
      \mathtt{magnitude}(\mathtt{x}-\mathtt{y}))\,]
  \]
  is a 2-element VNB list.  Internally it is \texttt{(list cc (vnb-lambda
  (list x y) (cartesian cc cc) (magnitude (- x y))))}.  The accessors \texttt{pts(s)} and
  \texttt{dist(s)} are \texttt{nth(1,~s)} and \texttt{nth(2,~s)}.
\item \textbf{The associated class} --- the collection
  $\{\,s : \mathit{is\text{-}name}(s)\,\}$ of \emph{all} VNB lists
  satisfying the structure's typing conditions.  \emph{This class
  shares the structure's name}: \texttt{declare-structure \textit{name}}
  installs both the predicate \texttt{is-\textit{name}} and the class
  \texttt{\textit{name}}, related by the axiom \texttt{\textit{name}-class}:
  \[
    \forall s.\; s \in \mathit{name} \;\Leftrightarrow\; \mathit{is\text{-}name}(s).
  \]
  The class is a proper class (defined by a predicate), not an
  element of \textsc{set}.  Naming the class is what makes the
  natural quantification idiom
  $\texttt{forall}([s \in \mathit{name}], \ldots)$
  work --- without it, one would have to write the unbounded form
  $\texttt{forall}([s], \mathit{is\text{-}name}(s) \Rightarrow \ldots)$.
\end{enumerate}
This matches the standard textbook presentation in which a ring, say,
is defined as a 6-tuple $(R,+,\times,-,0,1)$ satisfying certain axioms:
the tuple is the instance (point~2) and the collection of all such tuples
is the class (point~3), now an addressable name.
\end{remark}

\paragraph{How a declared structure is printed.}
The book of theorems and definitions (\texttt{docs/library-book.pdf}) and the
definitions block of \texttt{display-proof} print a declared structure from its
characteristic law, in three parts.  First, the components, numbered by their
\texttt{nth} index; each line gives the accessor as a projection
($\mathtt{add}(s) = \mathtt{nth}(2, s)$), the kind of clause that declared it
(\texttt{op}), and its typing conjunct.  Second, the laws: the conjuncts
contributed by the \texttt{property} and \texttt{law} clauses.  Third, one
sentence stating that \texttt{is-\textit{name}}$(s)$ holds exactly when $s$ has the
declared length, every component has its typing, and every law holds.  Every
conjunct of the law is printed once, in the first or the second part.  A
refinement (\texttt{same-shape-as}) is printed as its parent together with the
laws it adds.

\subsection{What \texttt{declare-structure} installs}

\paragraph{Accessor macetes.}
For each carrier $C_i$ (at \texttt{nth} index $i$) and each operation
$\mathit{op}_j$ (at \texttt{nth} index $m+j$), a macete is installed that
rewrites the accessor term to the corresponding \texttt{nth} projection:
\[
  C_i(s) \;\longmapsto\; \texttt{nth}(i,\; s)
  \qquad
  \mathit{op}_j(s) \;\longmapsto\; \texttt{nth}(m{+}j,\; s)
\]
The schema variable matches any expression, so the macete fires anywhere
in a goal.

\paragraph{\texttt{is-\textit{name}} axiom.}
The predicate \texttt{is-\textit{name}} is axiomatized as a conjunction
of two groups of conjuncts: the \emph{shape} (a length condition plus a
typing condition per slot) and the \emph{properties} (one conjunct per
\texttt{property} clause):
\begin{align*}
  \forall s.\;\mathit{is\text{-}name}(s) \;\Leftrightarrow\; &
  \mathtt{length}(s) = n
  && \text{\small(shape: length)}\\
  &\land\; C_1(s) \in \mathtt{set} \land \cdots \land C_m(s) \in \mathtt{set}
  && \text{\small(shape: carriers)}\\
  &\land\; \mathit{op}_1(s) \in \mathtt{fun}(\mathtt{dom}_1, \mathit{rng}_1)
  \land \cdots
  && \text{\small(shape: ops/constants)}\\
  &\land\; \mathtt{PRED}_1(\mathit{acc}(s), \ldots) \land \cdots
  && \text{\small(properties)}
\end{align*}
where $n = m + k$ is the total number of components ($m$ carriers and $k$
operations/constants).  The length condition ties the predicate to the
instance side of the three-part picture: only a VNB list of exactly $n$
elements can satisfy \texttt{is-\textit{name}}.  A declaration with no
\texttt{property} clause has the property group empty --- \texttt{is-\textit{name}}
is then a pure shape predicate.

\begin{sloppypar}The axiom is named \texttt{is-\textit{name}} and is immediately available
via \texttt{cmd-theorem-assumption} and \texttt{cmd-apply-macete}.\end{sloppypar}

\paragraph{\textit{name}-class axiom.}
A second axiom named \texttt{\textit{name}-class} ties the predicate to
the class:
\[
  \forall s.\; s \in \mathit{name} \;\Leftrightarrow\; \mathit{is\text{-}name}(s).
\]
The symbol $\mathit{name}$ thereby denotes the proper class
$\{s \mid \mathit{is\text{-}name}(s)\}$, ready to use as the
quantification range:
\begin{VNBsnip}
forall([r in RING],        some-claim-about(r))   ; bounded over the ring class
forall([E in NORMEDSPACE], some-claim-about(E))   ; bounded over normed spaces
\end{VNBsnip}
which the parser expands to
$\forall s.\; s \in \mathit{name} \Rightarrow \ldots$; the
\texttt{\textit{name}-class} axiom then converts the antecedent to
the predicate form $\mathit{is\text{-}name}(s)$ exactly when needed.

\subsection{Shape versus properties}
\label{sec:shape-vs-property}

A structure declaration says two separate things about an instance $s$.
Its \textbf{shape} is the list of components and the role each plays: how long
$s$ is, which slots are carriers, which are operations of which
signature, which are constants or substructures.  Its \textbf{properties}
are the characteristic laws those components must obey: that the
operation is associative, that the constant is a two-sided identity, that
the distance function is a metric.  Shape is what \texttt{carriers},
\texttt{op}, \texttt{constant}, and \texttt{substructure} contribute;
properties are what \texttt{property} clauses contribute.  Both are
conjoined into \texttt{is-\textit{name}}, so the predicate means
\emph{``$s$ is a $\mathit{name}$,''} not merely \emph{``$s$ is
$\mathit{name}$-shaped.''}

\paragraph{Property predicates.}
Each \texttt{property} clause names a predicate defined once, on its own,
in \texttt{operation-properties.scm}, and applies it to the structure's
accessors:
\begin{description}[leftmargin=2.4cm, style=nextline]
\item[\texttt{is-associative}] $(op, crr)$ --- $op$ is associative on $crr$.
\item[\texttt{is-commutative}] $(op, crr)$ --- $op$ is commutative on $crr$.
\item[\texttt{is-identity}] $(op, unit, crr)$ --- $unit$ is a two-sided identity.
\item[\texttt{has-inverses}] $(op, unit, invop, crr)$ --- $invop$ gives two-sided inverses.
\item[\texttt{is-distributive}] $(addop, mulop, crr)$ --- $mulop$ distributes over $addop$.
\item[\texttt{is-metric}] $(dist, crr)$ --- $dist$ is a metric (nonneg, zero on the diagonal, symmetric, triangle).
\item[\texttt{is-norm}, \texttt{is-group-norm}] --- norms for \texttt{normed-field} / \texttt{normed-ag}.
\end{description}
The parameter names (\texttt{op}, \texttt{crr}, \texttt{unit}, \ldots)
deliberately avoid the registered accessor constants
(\texttt{carr}, \texttt{mul}, \texttt{iden}, \texttt{pts}, \texttt{dist}, \ldots):
the reader case-folds, so a bound \texttt{carr} would collide with the
carrier accessor \texttt{carr}.

\begin{remark}
Folding the laws into \texttt{is-\textit{name}} is not cosmetic.  Without
the property group, two structures of the same shape would have
\emph{equivalent} predicates: \texttt{is-abelian-group} would be
$\Leftrightarrow$ \texttt{is-group}, since both are 4-slot
$(\mathtt{carr}, \mathtt{mul}, \mathtt{id}, \mathtt{inv})$ structures.  A
commutativity axiom stated separately would then force \emph{every}
group commutative.  With the property clause, \texttt{is-abelian-group}
is \texttt{is-group}'s body \emph{plus} \texttt{is-commutative} --- a
genuine sub-predicate, and non-commutative groups remain expressible.
\end{remark}

\begin{remark}
The equational laws are additionally available under their own names
(\texttt{monoid-assoc}, \texttt{metric-triangle}, \ldots) as separate
axioms, for citing a single law directly in a proof.  Once a structure
carries the law through a \texttt{property} clause, these named
restatements are derivable by projection rather than independent
assumptions: the metric laws, for instance, are proven modulo~0 by
projecting \texttt{is-metric} out of \texttt{is-metric-space}.  The
ring-tower predicates
(\texttt{commutative-ring}, \texttt{field}, \ldots) still spell their
\texttt{is-\textit{name}} out by a hand-written
$\texttt{is-}\mathit{name} \Leftrightarrow \texttt{is-}\mathit{parent}
\land \text{properties}$ axiom rather than through \texttt{property}
clauses --- equivalent in effect, a candidate for later unification.
\end{remark}

\subsection{Pre-loaded algebraic structures}
\label{sec:algebraic}

The core algebraic species are declared in \texttt{structure-library/} (one file
per structure) and loaded at startup.  The declarations below are \emph{generated}
from the live catalogue (\texttt{reference/STRUCTURE-INDEX.md}), so their slots and
laws match the running system.  Names are shown as \emph{stored}: the reader and
MIT~Scheme fold case, so \texttt{monoid} and \texttt{monoid}, \texttt{carr} and
\texttt{carr}, denote the same symbol.  Each \texttt{(property \ldots)} clause folds
one characteristic law into the structure's \texttt{is-X} predicate.

%% AUTO-GENERATED by docs/gen-structures.py from reference/STRUCTURE-INDEX.md -- do not edit.
%% Regenerate:  python3 docs/gen-structures.py   (after a prover load refreshes STRUCTURE-INDEX.md)

\subsubsection{Semigroup}
Declared in \texttt{structure-library/semigroup.scm}.
\begin{VNBcode}
(declare-structure semigroup
  (carriers carr)
  (op opr (cartesian carr carr) carr)
  (property is-associative opr carr))
\end{VNBcode}
The characteristic laws are the \texttt{(property \ldots)} clauses; the theorems quantifying over \texttt{is-semigroup}, with statements, are in \texttt{reference/STRUCTURE-INDEX.md} and \texttt{reference/BY-TOPIC.md}.

\subsubsection{Monoid}
Declared in \texttt{structure-library/monoid.scm}.
\begin{VNBcode}
(declare-structure monoid
  (carriers carr)
  (op opr (cartesian carr carr) carr)
  (constant iden carr)
  (property is-associative opr carr)
  (property is-identity opr iden carr))
\end{VNBcode}
The characteristic laws are the \texttt{(property \ldots)} clauses; the theorems quantifying over \texttt{is-monoid}, with statements, are in \texttt{reference/STRUCTURE-INDEX.md} and \texttt{reference/BY-TOPIC.md}.

\subsubsection{Commutative monoid}
Declared in \texttt{structure-library/monoid.scm}.
\begin{VNBcode}
(declare-structure comm-monoid
  (carriers carr)
  (op opr (cartesian carr carr) carr)
  (constant iden carr)
  (property is-associative opr carr)
  (property is-identity opr iden carr)
  (property is-commutative opr carr))
\end{VNBcode}
The characteristic laws are the \texttt{(property \ldots)} clauses; the theorems quantifying over \texttt{is-comm-monoid}, with statements, are in \texttt{reference/STRUCTURE-INDEX.md} and \texttt{reference/BY-TOPIC.md}.

\subsubsection{Group}
Declared in \texttt{structure-library/group.scm}.
\begin{VNBcode}
(declare-structure group
  (carriers carr)
  (op opr (cartesian carr carr) carr)
  (constant iden carr)
  (op inv carr carr)
  (property is-associative opr carr)
  (property is-identity opr iden carr)
  (property has-inverses opr iden inv carr))
\end{VNBcode}
The characteristic laws are the \texttt{(property \ldots)} clauses; the theorems quantifying over \texttt{is-group}, with statements, are in \texttt{reference/STRUCTURE-INDEX.md} and \texttt{reference/BY-TOPIC.md}.

\subsubsection{Abelian group}
Declared in \texttt{structure-library/abelian-group.scm}.
\begin{VNBcode}
(declare-structure abelian-group
  (carriers carr)
  (op opr (cartesian carr carr) carr)
  (constant iden carr)
  (op inv carr carr)
  (property is-associative opr carr)
  (property is-identity opr iden carr)
  (property has-inverses opr iden inv carr)
  (property is-commutative opr carr))
\end{VNBcode}
The characteristic laws are the \texttt{(property \ldots)} clauses; the theorems quantifying over \texttt{is-abelian-group}, with statements, are in \texttt{reference/STRUCTURE-INDEX.md} and \texttt{reference/BY-TOPIC.md}.

\subsubsection{Ring}
Declared in \texttt{structure-library/ring.scm}.
\begin{VNBcode}
(declare-structure ring
  (carriers carr)
  (op add (cartesian carr carr) carr)
  (op mul (cartesian carr carr) carr)
  (op neg carr carr)
  (constant zero carr)
  (constant one carr)
  (property is-associative add carr)
  (property is-commutative add carr)
  (property is-identity add zero carr)
  (property has-inverses add zero neg carr)
  (property is-associative mul carr)
  (property is-identity mul one carr)
  (property is-distributive add mul carr))
\end{VNBcode}
The characteristic laws are the \texttt{(property \ldots)} clauses; the theorems quantifying over \texttt{is-ring}, with statements, are in \texttt{reference/STRUCTURE-INDEX.md} and \texttt{reference/BY-TOPIC.md}.

\subsubsection{Metric space}
Declared in \texttt{structure-library/metric-space.scm}.
\begin{VNBcode}
(declare-structure metric-space
  (carriers pts)
  (op dist (cartesian pts pts) rr)
  (property is-metric dist pts))
\end{VNBcode}
The characteristic laws are the \texttt{(property \ldots)} clauses; the theorems quantifying over \texttt{is-metric-space}, with statements, are in \texttt{reference/STRUCTURE-INDEX.md} and \texttt{reference/BY-TOPIC.md}.



\subsection{Further structure definitions}

The same mechanism handles any arity of carriers and operations.  For
instance, a Banach space over $\mathbb{C}$:

\begin{VNBcode}
(declare-structure NORMEDSPACE
  (carriers VECTORS)
  (op SCALAR-TIMES (CARTESIAN CC VECTORS) VECTORS)
  (op PLUS         (CARTESIAN VECTORS VECTORS) VECTORS)
  (op MINUS        (CARTESIAN VECTORS VECTORS) VECTORS)
  (op NORM         VECTORS RR))
; VECTORS->1, SCALAR-TIMES->2, PLUS->3, MINUS->4, NORM->5
\end{VNBcode}

This installs five accessor macetes and adds the following axiom:
\begin{VNBsnip}
forall([s],
  IS-NORMEDSPACE(s) iff
    VECTORS(s) in SET
    and SCALAR-TIMES(s) in fun(cartesian(CC, VECTORS(s)), VECTORS(s))
    and PLUS(s) in fun(cartesian(VECTORS(s), VECTORS(s)), VECTORS(s))
    and MINUS(s) in fun(cartesian(VECTORS(s), VECTORS(s)), VECTORS(s))
    and NORM(s) in fun(VECTORS(s), RR))
\end{VNBsnip}

A characteristic axiom relating \texttt{minus} to \texttt{plus} and
\texttt{scalar-times} is installed separately.  Here the axiom is written
in S-expression form rather than string notation, because the string
parser reads \texttt{-1} as unary minus applied to \texttt{1} rather than
as a negative integer literal; the S-expression \texttt{-1} is a genuine
Scheme integer and avoids that ambiguity.

\begin{VNBcode}
(add-axiom! 'normedspace-minus
  '(FORALL s
     (IMPLIES (IS-NORMEDSPACE s)
       (FORALL x (IMPLIES (IN x (VECTORS s))
         (FORALL y (IMPLIES (IN y (VECTORS s))
           (= ((MINUS s) x y)
              ((PLUS s) x ((SCALAR-TIMES s) -1 y))))))))))
\end{VNBcode}

A vector space over an arbitrary field:
\begin{VNBcode}
(declare-structure VECTORSPACE
  (carriers FIELD VECTORS)
  (op SCALAR-TIMES (CARTESIAN FIELD VECTORS) VECTORS)
  (op PLUS         (CARTESIAN VECTORS VECTORS) VECTORS)
  (constant ZERO VECTORS)
  (op NEG          VECTORS VECTORS))
; FIELD->1, VECTORS->2, SCALAR-TIMES->3, PLUS->4, ZERO->5, NEG->6
\end{VNBcode}

\begin{remark}
\texttt{normedspace} and \texttt{vectorspace} above are illustrations of
multi-carrier declaration, not pre-loaded structures.  The library's actual
normed and vector spaces are \texttt{normed-vector-space} (carrier \texttt{vec},
with \texttt{vadd}/\texttt{vzero}/\texttt{vneg}/\texttt{act}/\texttt{vnrm} and a
\texttt{(substructure scal ring)} for the scalars) and \texttt{vector-space}
(declared \texttt{(same-shape-as module)} with \texttt{(law "is-field(scal(s))")});
see \texttt{reference/STRUCTURE-INDEX.md}.
\end{remark}

%%----------------------------------------------------------
\subsection{Refinement: \texttt{same-shape-as}}
\label{sec:same-shape-as}

A \emph{refinement} is a species with the same components as another but more
laws: a commutative ring is a ring whose multiplication commutes; an integral
domain, a commutative ring with no zero divisors.  Re-declaring the parent's
every slot would be verbose and would risk drifting from it.  A
\texttt{(same-shape-as PARENT)} clause instead inherits the parent's shape
wholesale --- its slots, its accessor indices, its shape conjuncts --- and the
declaration adds only the new \texttt{(law "\ldots")} clauses:

\begin{VNBcode}
(declare-structure COMMUTATIVE-RING
  (same-shape-as RING)
  (law "forall([a in carr(s), b in carr(s)], mul(s)(a, b) = mul(s)(b, a))"))
\end{VNBcode}

The installed predicate is the parent's, strengthened, with
\texttt{is-\textit{parent}} as its \emph{first, literal} conjunct:
\[
  \forall s.\;\mathit{is\text{-}commutative\text{-}ring}(s)
    \;\Leftrightarrow\;
    \mathit{is\text{-}ring}(s)\ \wedge\ (\text{multiplication commutes}).
\]
Two things follow.  The refinement is transparently a sub-predicate of its
parent --- the subtype machinery (\texttt{subtype-laws.scm}) reads
\texttt{is-ring} straight off the first conjunct --- and \texttt{carr},
\texttt{add}, \texttt{mul}, \ldots keep exactly the indices they carry under
\texttt{ring}, so every accessor macete is shared.

Because the shape is inherited, a \texttt{same-shape-as} declaration \emph{may
not} declare slots of its own: a \texttt{(carriers \ldots)}, \texttt{(op
\ldots)}, \texttt{(constant \ldots)}, or \texttt{(substructure \ldots)} clause
beside it is a load error.  A species of genuinely different shape is a
different structure, related to this one not by refinement but by a functor
(Section~\ref{sec:view-as}).

%%----------------------------------------------------------
\subsection{Declaring an instance: \texttt{declare-instance!}}
\label{sec:declare-instance}

A structure \emph{species} is a class; an \emph{instance} is one concrete tuple
belonging to it --- $\mathbb{Z}$ as a ring, $\mathbb{C}$ as a metric space.
\texttt{declare-instance!} names such a tuple, defines it, and wires it into the
catalogue in one call:

\begin{VNBcode}
(declare-instance! 'ZZ-RING 'RING 'zz-ring-def
  '(ZZ binplus bintimes binneg 0 1))
\end{VNBcode}

The arguments are the instance name, the species it instantiates, the name to
give its defining equation, and the tuple of slot values in declared slot order.
The call:
\begin{itemize}
\item \textbf{checks the tuple against the shape} --- a tuple of the wrong length
  for \texttt{ring} fails the load rather than silently asserting a false
  \texttt{is-ring};
\item \textbf{installs the defining equation} under the given name
  (\texttt{zz-ring-def}: $\mathit{zz\text{-}ring} = \mathtt{list}(\ldots)$);
\item \textbf{installs one per-slot value macete}, \texttt{\textit{name}@\textit{acc}},
  so \texttt{(mul zz-ring)} reduces to \texttt{bintimes} in a single
  \texttt{slot} step instead of exposing \texttt{(nth 3 zz-ring)}.  These use
  \texttt{==}, not \texttt{=}: the slot value is unconditionally that term, so no
  definedness obligation is owed;
\item \textbf{registers the instance} as a definitional structure with the
  species as parent, so it earns a card in \texttt{describe-structure}, a node in
  the structure graph, and a place in \texttt{known-structures}.
\end{itemize}
The instance still owes an \texttt{is-\textit{species}} theorem to be
\emph{certified}; once that is proved, \texttt{specialize-structure} (below)
transports the species' theorems to it.

%%----------------------------------------------------------
\subsection{Structure specialization: \texttt{specialize-structure}}
\label{sec:specialize-structure}

Every characteristic axiom and every theorem proved about a structure class
has the generic shape
\[
  \forall s.\;\mathit{is\text{-}name}(s) \;\Rightarrow\; P[s].
\]
Once a specific instance $h$ has been \emph{certified}---meaning a theorem
of the form $\mathit{is\text{-}ring}(h)$ has been proved and installed---every
such generic statement can be specialized to $h$ by universal instantiation
followed by modus ponens, yielding the concrete statement $P[h]$.

\texttt{specialize-structure} automates this transport in bulk:

\begin{VNBcode}
(specialize-structure 'HARP 'RING 'harp-is-ring)
; => N theorems specialized for HARP.
\end{VNBcode}

The interactive short form is \texttt{(spec instance struct is-thm)}.

\paragraph{Arguments.}
\begin{itemize}
\item \textit{instance-name} --- symbol naming the specific structure, e.g.\
  \texttt{harp}.
\item \textit{struct-name} --- the structure class, e.g.\ \texttt{ring}.
  The predicate \texttt{is-ring} is derived automatically.
\item \textit{is-thm-name} --- name of an installed theorem
  whose formula is exactly $\mathit{is\text{-}ring}(\mathtt{harp})$.
\end{itemize}

\paragraph{What gets specialized.}
Every entry in the global theorem table (axioms and proved theorems alike)
whose formula matches
\[
  \forall s.\;\mathit{is\text{-}ring}(s) \;\Rightarrow\; P[s]
\]
--- with $s$ the bound variable appearing as the sole argument to
$\mathit{is\text{-}ring}$ --- is specialized.  The result $P[\mathtt{harp}]$
is installed under the name
$\langle\textit{original\text{-}name}\rangle\texttt{-harp}$.
Axioms such as \texttt{ring-add-assoc} are included, so
\texttt{ring-add-assoc-harp} is immediately available after the call.

\paragraph{Soundness.}
No new axioms are introduced.  Each specialized theorem $P[h]$ follows
from the certified $\mathit{is\text{-}ring}(h)$ and the generic
$\forall s.\;\mathit{is\text{-}ring}(s)\Rightarrow P[s]$ by two standard
inferences: universal instantiation and modus ponens.
\texttt{specialize-structure} performs these in bulk without requiring a
separate interactive proof for each one.  The certification check at the
start of the call verifies that \textit{is-thm-name} has the expected form
before any specialization occurs; an incorrect theorem name raises an error.

\begin{remark}
Theorems about sub-structures are \emph{not} automatically transported.
For instance, theorems of the form
$\forall s.\;\mathit{is\text{-}group}(s)\Rightarrow\ldots$
are not touched by \texttt{(spec 'harp 'ring ...)}.  To obtain them you
would need to certify $\mathit{is\text{-}group}(\mathtt{harp})$ separately
and call \texttt{spec} again for the \texttt{group} class.
\end{remark}

\paragraph{Example: the integers as a ring.}
The constant \texttt{zz-ring} is defined in \texttt{numeric-instances.scm} as
the 6-tuple $[\mathtt{zz},\,+,\,\times,\,-,\,0,\,1]$.  The axiom
\texttt{zz-is-ring} asserts $\mathit{is\text{-}ring}(\mathit{zz\text{-}ring})$.
(We use $\mathbb{Z}$, not $\mathbb{C}$: \texttt{cc-normed-field} is a 7-tuple normed
field, so it cannot witness the length-6 \texttt{is-ring} on the same tuple;
it reaches the ring world through the \texttt{normed-field-as-*} views.)

\begin{VNBcode}
(spec 'ZZ-RING 'RING 'zz-is-ring)
; => n theorems specialized for ZZ-RING.
; Installs: ring-add-assoc-zz-ring, ring-add-comm-zz-ring,
;           ring-mul-assoc-zz-ring, ring-left-dist-zz-ring, ...
\end{VNBcode}

A proof about $\mathbb{Z}$ can then retrieve any of these directly:
\begin{VNBcode}
(ta 'ring-add-comm-zz-ring)
; adds  ADD(ZZ-RING)(x, y) = ADD(ZZ-RING)(y, x)  to the context
\end{VNBcode}

%% =========================================================
\section{\texttt{def-functoid}}
\label{sec:def-functoid}

\texttt{def-functoid} introduces a named functoid at the proof-checker
level.  It is a definitional device, not an object of the library: the
name it installs can be applied and reasoned about, but there is no
formal quantifier ranging over functoids.

The general form is:
\begin{VNBcode}
(def-functoid name (x) body)
\end{VNBcode}
where \texttt{x} is a formal parameter ranging over classes and
\texttt{body} is a class-valued expression, typically built from
accessor macetes installed by \texttt{declare-structure}.

\paragraph{Example: forgetful functor NORMEDSPACE $\to$ METRIC-SPACE.}
A normed space carries an induced metric $d(x,y) = \|x - y\|$.
Since the domain is the proper class of all normed spaces, \texttt{fun}
cannot express this mapping and a functoid is required:

\begin{VNBcode}
(def-functoid 'NormedToMetric '(S)
  '(LIST (VECTORS S)
         (VNB-LAMBDA (LIST x y) (CARTESIAN (CARR S) (CARR S)) ((NORM S) ((MINUS S) x y)))))
\end{VNBcode}

The result is a tuple \texttt{[vectors(S), D]} matching the
\texttt{metric-space} signature (carrier \texttt{pts}, distance \texttt{dist}).
That it satisfies \texttt{is-metric-space} is a theorem, not automatic:
the metric axioms follow from the norm axioms and the definition of
\texttt{minus}.

\paragraph{Example: forgetful functor NORMEDSPACE $\to$ VECTORSPACE.}
Dropping the norm recovers the underlying vector space:

\begin{VNBcode}
(def-functoid 'NormedToVector '(S)
  '(LIST (VECTORS S) (SCALAR-TIMES S) (PLUS S) (MINUS S)))
\end{VNBcode}

\paragraph{Example: the carrier functor.}
\begin{VNBcode}
(def-functoid 'GroupCarrier '(G) '(CARR G))
\end{VNBcode}
This maps every group to its underlying carrier set.  The result is
always a set (\texttt{is-group} entails \texttt{carr(G) in set}),
even though the domain is a proper class.

\paragraph{Transfinite recursion and functoids.}
\begin{sloppypar}The transfinite recursion scheme
(Section~\ref{sec:ord-recursion}) produces a functoid, not a
set-theoretic function: its domain is the proper class $\mathtt{ord}$
of all ordinals, so the result cannot live in $\mathtt{fun}(\mathtt{ord}, B)$
for any set-valued $B$.  A typical pattern is to define the step
functor $\mathcal{E}$ as a functoid via \texttt{def-functoid} and
then invoke \texttt{def-by-ord-recursion}:\end{sloppypar}
The signature has seven arguments: the new functoid name, the list of
parameters (\texttt{'()} for none), the base value, plus a binding spec
and body for both the successor and limit cases.  The successor body is evaluated with the bound name
referring to the recursive value at the predecessor; the limit body
is evaluated with the bound name referring to the partial functoid
restricted to all earlier ordinals.
\begin{VNBcode}
(def-by-ord-recursion 'MySeq '()
  'a0                                 ; base value at 0
  '(prev) '(succ_ORD prev)            ; successor: MySeq(succ alpha) = succ_ORD(prev)
  '(part) '(SUP-ORD (image part)))    ; limit:    MySeq(lambda) = sup of partial seq
; => installs definitional axioms
;      MySeq(0) == a0
;      IN alpha ORD => MySeq(succ_ORD alpha) == succ_ORD(MySeq alpha)
;      LIMIT-ORD lambda => MySeq(lambda) == SUP-ORD(image-of MySeq below lambda)
\end{VNBcode}

%%% -----------------------------------------------------------------------
\section{Notation: \texttt{notation!}}
\label{sec:notation}

Every head the system knows --- a predicate, a functoid, a structure predicate,
even a primitive like \texttt{=} or \texttt{in} --- has an entry in a single
operator table.  \texttt{def-predicate}, \texttt{def-functoid} and
\texttt{declare-structure} fill in the \emph{structural} part of the entry (its
kind, its arity) at the definition point.  \texttt{notation!} fills in the
\emph{reading} --- how the head renders in \LaTeX, in English, and to the
parser.  The two are orthogonal: re-declaring a head updates the structural slots
and keeps whatever notation was declared, so a \texttt{notation!} may sit beside
the definition or well after it.

It is a keyword form; a call names only the slots it has to set:
\begin{VNBcode}
(notation! 'IS-RING 'noun "ring" 'article "a")
(notation! 'IS-EUCLIDEAN-RING 'kind 'predicate 'arity 1
           'noun "Euclidean ring" 'article "a")
(notation! 'IS-COUNTABLE-PSEUDOMETRIC-FAMILY 'kind 'predicate 'arity 2
           'english "$1 is a countable family of pseudometrics on $2")
\end{VNBcode}

The recognised keys are
\texttt{tex}, \texttt{english}, \texttt{noun}, \texttt{article},
\texttt{prec}, \texttt{infix?}, \texttt{arity}, and \texttt{kind}.

\begin{itemize}
\item \texttt{tex} and \texttt{english} are \emph{templates} --- strings with
  placeholders \texttt{\$1}\ldots\texttt{\$9} for the rendered arguments.  A
  binary operator combining both with \texttt{prec} (a precedence for
  parenthesisation) and \texttt{infix?} reads, schematically,
  \texttt{(notation! '+ 'tex "\$1 + \$2" 'english "\$1 plus \$2" 'prec 60 'infix?
  \#t)}.  When a template cannot express the rendering, \texttt{english} may
  instead be a \emph{procedure} from the list of argument strings to a string.
\item \texttt{noun} and \texttt{article} are the special case of a unary
  structure predicate: given them, \texttt{is-ring($s$)} reads ``$s$ is a
  ring'' throughout the prose renderers, and the article joins correctly
  (``a ring'', ``an $R$-module'').
\item \texttt{kind} and \texttt{arity} let a \texttt{notation!} register a head
  the def-forms never introduced --- a primitive such as \texttt{=} or
  \texttt{in} --- so it, too, gets a reading.
\end{itemize}

One table, one declaration form, read by every renderer: \texttt{expr->tex} and
\texttt{wff-english} for formulas, \texttt{describe-structure} for the structure
cards, and the generated \texttt{OPERATORS.md}.  A head with no
\texttt{notation!} still works --- it falls back to prefix application,
\texttt{is-ring($s$)} --- so notation is an addition, never a prerequisite.

%%% -----------------------------------------------------------------------
\section{The optional infix layer}
\label{sec:infix-layer}

The TeX printer writes an applied structure operation in functional
form: \texttt{((add K) a b)} prints as $\mathtt{add}(K)(a, b)$.  The
\emph{infix layer} (\texttt{tex-infix.scm}) writes it instead as
\[
  a +_K b .
\]
It is off by default.  While it is off the printer prints exactly what
it printed before.  It changes only TeX output: the panels, the
proof reader and \texttt{display-proof}.  The parser, \texttt{pp} and
the kernel never see it.

It differs from \texttt{notation!} (Section~\ref{sec:notation}) in
two ways.  A \texttt{notation!} declares the reading of one head and is
always in force.  The layer reads an \emph{accessor applied to a
structure}, so one row serves \texttt{add} in every ring, module and
field, and it is switched on and off as a whole.

\subsection{Turning it on}

In the workspace, the menu \textsf{VNB > Printer Macros} has the
toggle \textsf{Printer Macros On}; the panels repaint at once.  At the
REPL:
\begin{VNBcode}
(tex-infix! #t)                          ; on for every printout
(tex-infix! #f)                          ; off
(display-proof 'NAME "file.pdf" 'infix)  ; on for this printout only
(infix-tex 'NAME)                        ; the TeX of a theorem, layer on
(with-tex-infix CONV thunk)              ; one printout, conventions CONV
\end{VNBcode}

\subsection{What it prints}

An accessor applied to a structure term $S$, then to its arguments,
is written with the subscript $S$:
\[
\begin{array}{ll@{\qquad}ll}
\texttt{((add S) a b)}  & a +_S b        & \texttt{((mul S) a b)}  & a \times_S b \\
\texttt{((opr S) a b)}  & a \cdot_S b    & \texttt{((act S) c v)}  & c \cdot_S v \\
\texttt{((neg S) a)}    & -_S\, a        & \texttt{((inv S) a)}    & a^{-1}_S \\
\texttt{(zero S)}       & 0_S            & \texttt{(iden S)}       & e_S \\
\texttt{((fnrm S) a)}   & \lvert a \rvert_S & \texttt{((vnrm S) v)} & \lVert v \rVert_S \\
\texttt{((ip S) u v)}   & \langle u, v \rangle_S & \texttt{((vstar S) x)} & x^{*}_S
\end{array}
\]
The full list is the base table (below).  $S$ is read through views:
a \texttt{def-functor} view such as
\texttt{(normed-field-as-commutative-ring K)} is written $K$, since the
view's slots are $K$'s slots.

A defined head can also have a template.  The shipped templates print
the sums and products over an index:
\[
  \texttt{series-partial-sum(c, n)} \quad\text{as}\quad \sum_{i < n} c(i).
\]
When the summand is a \texttt{vnb-lambda} over $\mathtt{nn}$, the
lambda's own variable is used and its body printed, as in
$\sum_{k < n} c(k)\, x^k$.  Otherwise the printer invents an index
variable that clashes with no symbol of the formula.

\subsection{Faithfulness}

The layer changes notation, never the term.  Nothing is simplified or
rewritten: $a + (-b)$ prints as
\[
  a +_S -_S\, b ,
\]
never as $a - b$.  Every subterm the plain printer would show is still
there.

Parentheses follow the binding strengths: $+$ is weakest, then the
products, then unary minus, then the postfix operations.  A chain of
one associative operation of one structure nested to the \emph{left}
drops its parentheses.  Nested to the right it keeps them, so an
associativity law still reads as one:
\[
  a + b + c \;=\; a + (b + c).
\]

\subsection{Conventions: when the subscript is dropped}

The conventions are a list, like a typing block:
\begin{itemize}
\item \texttt{(auto)}, the default.  Drop the subscript when one
  structure is in play in the formula being printed and it is a
  variable, not a named constant such as \texttt{zz-ring}.  Its
  constants $0_S$, $1_S$, $e_S$ keep their subscript when the formula
  also holds a numeral, so $0_K$ never reads as the number $0$.
\item \texttt{(keep)}.  Always write the subscript.
\item \texttt{(drop S \ldots)}.  Drop the subscript of these structure
  terms.
\item \texttt{(symbol HEAD "tex")}.  Write \texttt{HEAD}'s operation
  with this symbol.
\end{itemize}
``The formula being printed'' is the term given to the outermost
printer call.  A display broken over rows is scanned row by row.

\subsection{The tables}

The operations, their TeX and the default conventions are data, not
code.  The base table is \texttt{notation/base.tbl} in the tree; its
head describes the format.  A row reads
\begin{VNBcode}
; KIND     HEAD      PREC  ASSOC  TEX
infix      add       10    assoc  +
postfix    inv       40    -      ^{-1}
outfix     vnrm      100   -      \norm{#1}
apply      series-partial-sum  5  -  \sum_{#i < #2} #1(#i)
\end{VNBcode}
The kinds are \texttt{infix}, \texttt{prefix}, \texttt{postfix},
\texttt{const}, \texttt{outfix} (a template in \texttt{\#1}),
\texttt{bracket} (in \texttt{\#1} and \texttt{\#2}) and \texttt{apply}
(a template for a defined head).  An \texttt{apply} template must use
every argument: it may not hide one.  There are also \texttt{macro}
rows, \texttt{off HEAD}, and \texttt{convention} rows.

Further tables are read after the base table, and a later row replaces
an earlier row of the same head.  The order is:
\begin{enumerate}
\item the tables named, one per line, in
  \texttt{\textasciitilde/vnb-notation/active}.  A name is
  \texttt{\textasciitilde/vnb-notation/NAME.tbl} when that file exists,
  else the tree's \texttt{notation/NAME.tbl};
\item the personal table \texttt{\textasciitilde/vnb-notation/personal.tbl},
  when it exists.
\end{enumerate}
The menu \textsf{VNB > Printer Macros} turns each table on or off, and
\textsf{Edit or New Table\ldots} and \textsf{Edit Personal Table} open
one for editing.

A changed table is re-read at the next printout with the layer on, or
at once by \textsf{Reload} or \texttt{(tex-infix-reload!)}.  A reload
with any row in error changes nothing.  The errors are listed by file
and line, and the table in force stays.

%%% -----------------------------------------------------------------------
\section{Cardinality: \texttt{card}}
\label{sec:cardinality}

\texttt{card} is a functoid (\texttt{structure-library/cardinality.scm}) defined by
a description: $\mathtt{card}(A)$ is the least ordinal $\alpha$ such that
$A$ bijects onto $\mathit{ord\text{-}segment}(\alpha)$, with $A$ first in
\texttt{bijection}.  No axiom is asserted about it.  The eight facts that
once stood on the primitive shelf are theorems under their old names and
in their old shapes: \texttt{card-segment} (\texttt{card-defined.scm}),
\texttt{card-empty} (\texttt{card-finite.scm}), \texttt{finite-set-induction}
(\texttt{rake-card-star-laws.scm}), and \texttt{card-in-ord},
\texttt{card-insert}, \texttt{card-finite-bij}, \texttt{card-union-disjoint},
\texttt{card-image-injection} (\texttt{card-laws.scm}).

\begin{remark}
Because the domain of $\mathtt{card}$ is the proper class $\mathtt{set}$,
it cannot be expressed as an element of $\mathtt{fun}(\mathtt{set}, \mathtt{ord})$
(that class is not a set).  It is therefore a functoid, on the same
footing as $\mathit{ord\text{-}segment}$.
\end{remark}

\begin{remark}
The definition is a description, and the library proves nothing about
$\mathtt{card}(A)$ for a proper class $A$: the term $\mathtt{card}(\mathtt{ord})$
is syntactically well-formed, but no ordinal is exhibited that satisfies the
description, and a description is never certified as denoting without a
proof (Section~\ref{sec:kernel-ops-detail}).  Any equational use of
$\mathtt{card}(A)$ outside the sethood-conditioned theorems below will
not discharge.  Convention: only invoke $\mathtt{card}$ on a class you
have proven to be a set.
\end{remark}

\subsection*{Theorems}

\begin{description}
\item[\texttt{card-in-ord}]
  $A \in \mathtt{set} \Rightarrow \mathtt{card}(A) \in \mathtt{ord}$.

\item[\texttt{card-empty}]
  $\mathtt{card}(\emptyset) = 0$.

\item[\texttt{card-insert}]
  $A \in \mathtt{set} \wedge \mathtt{card}(A) \in \mathtt{nn} \wedge x \notin A \;\Rightarrow\;
   \mathtt{card}(A \cup \{x\}) = \mathit{succ_{ord}}(\mathtt{card}(A))$.\\
  Here $\{x\} = \mathtt{pair}(x,x)$.  The finiteness guard is needed:
  for $A = x = \omega$ the set $A \cup \{x\}$ is $\omega + 1$, which
  bijects with $\omega$, so its cardinality is $\omega$ and not $\omega + 1$.

\item[\texttt{card-segment}]
  $n \in \mathtt{nn} \;\Rightarrow\; \mathtt{card}(\mathit{ord\text{-}segment}(n)) = n$.\\
  Bridge between cardinality and the $\mathtt{nn}$-indexed model of finite sets.

\item[\texttt{card-finite-bij}]
  $A \in \mathtt{set} \wedge \mathtt{card}(A) \in \mathtt{nn} \;\Rightarrow\;$
  there exists a bijection
  $\varphi : \mathit{ord\text{-}segment}(\mathtt{card}(A)) \to A$.

\item[\texttt{card-union-disjoint}]
  $A, B \in \mathtt{set},\;
   \mathtt{card}(A), \mathtt{card}(B) \in \mathtt{nn},\;
   A \cap B = \emptyset \;\Rightarrow\;
   \mathtt{card}(A \cup B) = \mathtt{card}(A) + \mathtt{card}(B)$.\\
  Finite additivity, with the same finiteness guard; the $+$ on the right is
  $\mathtt{nn}$ addition.
\end{description}

\subsection*{Finiteness criterion}

A set $A$ is \emph{finite} iff $\mathtt{card}(A) \in \mathtt{nn}$.
For finite sets, \texttt{card-segment} says
$\mathtt{card}(\mathit{ord\text{-}segment}(n)) = n$, so the initial
segment $\{0, 1, \ldots, n-1\}$ has cardinality exactly $n$.  The
theorem \texttt{card-finite-bij} supplies an enumeration bijection for
any finite set, which is the key ingredient for defining unordered
finite sums in Section~\ref{sec:sequences}.

%%% -----------------------------------------------------------------------
\section{Morphisms of structures}
\label{sec:morphisms}

Declaring a structure declares its morphisms too.  From the slot list alone,
\texttt{declare-structure} generates a homomorphism predicate
\texttt{is-hom-\textit{name}}$(a, b, f_1 \ldots f_k)$ --- one map $f_i$ per
\emph{independent} carrier --- whose meaning is the obvious one: each map is
typed between the corresponding carriers, and every operation and constant is
preserved.

\begin{itemize}
\item \textbf{carrier} $c$: a map \texttt{$f$ in fun($c(a)$, $c(b)$)}.
\item \textbf{operation} \texttt{OP}: preservation ---
  \texttt{$f$(OP($a$)($x_1,\ldots$)) = OP($b$)($f(x_1),\ldots$)} for all
  arguments in their sorts (each argument mapped by its own carrier's map).
\item \textbf{constant} \texttt{C}: \texttt{$f$(C($a$)) = C($b$)}.
\item \textbf{substructure} \texttt{K}: agreement, \texttt{K($a$) = K($b$)} ---
  morphisms are taken between structures over the \emph{same} base.
\item \textbf{derived} carrier: nothing.  Its value is pinned by
  \texttt{is-\textit{name}} (Section~\ref{sec:declare-structure}), so it rides
  the base carrier's map and owes no map of its own.  This is why a field
  homomorphism is one map, not two.
\end{itemize}

\noindent A group's generated hom, for instance, is the single map \texttt{$f$
in fun(carr($a$), carr($b$))} with \texttt{$f$(opr($a$)($x$,$y$)) =
opr($b$)($f(x)$,$f(y)$)} and \texttt{$f$(iden($a$)) = iden($b$)} --- exactly a
group homomorphism.  The definition is installed with \emph{definitional}
provenance, so it costs no warrant.

\subsection{Overriding the default: \texttt{declare-hom!}}
\label{sec:declare-hom}

The generated clause for each slot kind is the standard one: preservation for
an operation or a constant, the pullback condition for a \texttt{family} slot
(the arrows of \texttt{top-space} are the continuous maps, those of
\texttt{measurable-space} the measurable maps), transport for a
\texttt{family-fun} slot, ``related points go to related points'' for a
\texttt{relation} slot.  For a \texttt{metric-space} the generated arrows
preserve the distance: they are the isometries, and a \texttt{notation!}
beside the structure says so, since ``homomorphism'' is then the wrong word
(Section~\ref{sec:notation}).  Where no generated clause is the right one, a
species declares its own morphism notion with \texttt{declare-hom!}, beside the
structure like \texttt{notation!}:

\begin{VNBcode}
(declare-hom! 'METRIZABLE-TOP-SPACE '(s t f)
  '(FORALL u (IMPLIES (IN u (OPENS t))
      (IN (PREIMAGE s f u) (OPENS s)))))
\end{VNBcode}

The argument list is the two objects and one map per carrier.  The body states
\emph{only} what is characteristic: the object typings
\texttt{is-metrizable-top-space($s$)}, \texttt{is-metrizable-top-space($t$)}
and the map typing \texttt{$f$ in fun(pts($s$), pts($t$))} are supplied
automatically and need not be repeated.  The declaration replaces the generated
\texttt{is-hom-metrizable-top-space} entirely.  Two species override today:
\texttt{metrizable-top-space} (a refinement inherits its parent's shape, not
its arrows, so continuity is restated) and \texttt{ringoid} (a ring
homomorphism of the underlying rings that preserves the ideal).  The table of
Section~\ref{sec:declare-category} says, for every structure, which of the
three cases its arrows are.

\subsection{Several categories on one species: \texttt{declare-category!}}
\label{sec:declare-category}

A structure is a category: its objects are its models, its arrows the maps
satisfying \texttt{is-hom-\textit{name}}, composed as functions.  That is the
species' \emph{default} category, and for an algebraic species it is the right
one.  For a metric space the default is the category of isometries, which is
not the one analysis works in: the Lipschitz maps, the uniformly continuous
maps and the continuous maps each form a category on the same objects; on the
countably-metrised spaces (\texttt{c-metric-space}) the continuous maps for
the canonical metric form one, which rests on the canonical metric being a
metric (\texttt{c-metric-is-metric-space}).  A species may therefore carry
several categories.  A second one is declared
beside the structure:

\begin{VNBcode}
(declare-category! 'LIPSCHITZ 'METRIC-SPACE '(a b f)
  '(FORSOME lpk_ (AND (IN lpk_ RR) (AND (<= 0 lpk_)
     (FORALL lpx_ (IMPLIES (IN lpx_ (PTS a))
       (FORALL lpy_ (IMPLIES (IN lpy_ (PTS a))
         (<= ((DIST b) (f lpx_) (f lpy_))
             (* lpk_ ((DIST a) lpx_ lpy_)))))))))))
\end{VNBcode}

The body states what is characteristic of an arrow $f \colon a \to b$; the
object typings and the carrier typing of each map are supplied, as for
\texttt{declare-hom!}.  The declaration installs, definitionally, the arrow
predicate \texttt{is-lipschitz-arrow}$(a, b, f)$ and the hom-set
\texttt{hom-lipschitz}$(a, b)$, and registers three \emph{obligations} as
statements, not as theorems: \texttt{hom-lipschitz-id} (the identity is an
arrow), \texttt{hom-lipschitz-compose} (arrows compose) and
\texttt{hom-lipschitz-in-set} (the hom-set is a set).  That the arrows of a
declared category contain the identities and compose is where the
mathematics is (Lipschitz constants multiply; an $\varepsilon$--$\delta$
argument for the continuous maps), so the generator does not stamp it: the
obligations are proven in \texttt{theorem-library/categories.scm}, and
\texttt{category-obligations-audit}, run at the end of every load, lists
every obligation that is not a proven theorem of exactly its statement.
Until that list is empty the declared category is not yet a category.  The
default category is untouched by the declaration.

The categories the library knows are listed below.  The listing is generated
at every load (\texttt{reference/CATEGORIES.md}, also a page of the browser
reference): the default category of each structure with the status of its
three laws, the declared categories with their defining equivalences, and the
inclusions between categories on the same objects --- an isometry is
$1$-Lipschitz, a Lipschitz map is uniformly continuous, a uniformly
continuous map is continuous, a linear isometry is bounded linear --- which
the page finds by the shape of the installed theorems.

%% AUTO-GENERATED by docs/gen-categories.py from reference/CATEGORIES.md -- do not edit.
%% Regenerate:  python3 docs/gen-categories.py   (after a prover load refreshes CATEGORIES.md)

% the long identifiers break only with loose spacing; \sloppy, scoped to this file
\begingroup\sloppy
\subsubsection{The default categories}
One per structure with a morphism predicate.  The laws column reads ``proven'' when the identity law, the composition law and the sethood of the hom-set are all proven theorems; otherwise it names the law and its status.

\begin{longtable}{@{}>{\raggedright\arraybackslash}p{0.22\linewidth}>{\raggedright\arraybackslash}p{0.30\linewidth}>{\raggedright\arraybackslash}p{0.20\linewidth}>{\raggedright\arraybackslash}p{0.18\linewidth}@{}}
\toprule
structure & arrows & defined & laws \\
\midrule
\endhead
\texttt{abelian-\allowbreak{}group} & f is a homomorphism from a to b & generated & proven \\
\texttt{bundle} & f is a homomorphism from a to b & generated & id: not installed; compose: not installed; in-set: not installed \\
\texttt{c-\allowbreak{}metric-\allowbreak{}space} & f is a homomorphism from a to b & generated & proven \\
\texttt{comm-\allowbreak{}monoid} & f is a homomorphism from a to b & generated & proven \\
\texttt{commutative-\allowbreak{}banach-\allowbreak{}algebra} & f is a homomorphism from a to b & refinement of \texttt{unital-\allowbreak{}banach-\allowbreak{}algebra} & id: not installed; compose: not installed; in-set: not installed \\
\texttt{commutative-\allowbreak{}c-\allowbreak{}star-\allowbreak{}algebra} & f is a homomorphism from a to b & refinement of \texttt{unital-\allowbreak{}c-\allowbreak{}star-\allowbreak{}algebra} & id: not installed; compose: not installed; in-set: not installed \\
\texttt{commutative-\allowbreak{}ring} & f is a homomorphism from a to b & refinement of \texttt{ring} & proven \\
\texttt{complex-\allowbreak{}banach-\allowbreak{}algebra} & f is a homomorphism from a to b & refinement of \texttt{complex-\allowbreak{}normed-\allowbreak{}algebra} & id: not installed; compose: not installed; in-set: not installed \\
\texttt{complex-\allowbreak{}banach-\allowbreak{}space} & f is a homomorphism from a to b & refinement of \texttt{complex-\allowbreak{}normed-\allowbreak{}vector-\allowbreak{}space} & id: not installed; compose: not installed; in-set: not installed \\
\texttt{complex-\allowbreak{}inner-\allowbreak{}product-\allowbreak{}space} & f is a homomorphism from a to b & generated & proven \\
\texttt{complex-\allowbreak{}normed-\allowbreak{}algebra} & f is a homomorphism from a to b & generated & id: not installed; compose: not installed; in-set: not installed \\
\texttt{complex-\allowbreak{}normed-\allowbreak{}vector-\allowbreak{}space} & f is a homomorphism from a to b & generated & id: not installed; compose: not installed; in-set: not installed \\
\texttt{complex-\allowbreak{}pre-\allowbreak{}hilbert-\allowbreak{}space} & f is a homomorphism from a to b & refinement of \texttt{complex-\allowbreak{}semi-\allowbreak{}inner-\allowbreak{}product-\allowbreak{}space} & id: not installed; compose: not installed; in-set: not installed \\
\texttt{complex-\allowbreak{}semi-\allowbreak{}inner-\allowbreak{}product-\allowbreak{}space} & f is a homomorphism from a to b & generated & id: not installed; compose: not installed; in-set: not installed \\
\texttt{directed-\allowbreak{}set} & f is a homomorphism from a to b & generated & id: not installed; compose: not installed; in-set: not installed \\
\texttt{euclidean-\allowbreak{}ring} & f is a homomorphism from a to b & refinement of \texttt{integral-\allowbreak{}domain} & proven \\
\texttt{field} & f is a homomorphism from a to b & generated & proven \\
\texttt{field-\allowbreak{}ring} & f is a homomorphism from a to b & refinement of \texttt{commutative-\allowbreak{}ring} & proven \\
\texttt{group} & f is a homomorphism from a to b & generated & proven \\
\texttt{infiltrated-\allowbreak{}space} & f1 is a morphism of infiltrated strand spaces from a to b & generated & id: not installed; compose: not installed; in-set: not installed \\
\texttt{integral-\allowbreak{}domain} & f is a homomorphism from a to b & refinement of \texttt{commutative-\allowbreak{}ring} & proven \\
\texttt{measurable-\allowbreak{}space} & f is a measurable map from a to b & generated & proven \\
\texttt{measure-\allowbreak{}space} & f is a measure-preserving measurable map from a to b & generated & proven \\
\texttt{metric-\allowbreak{}group} & f is an isometric homomorphism from a to b & generated & id: not installed; compose: not installed; in-set: not installed \\
\texttt{metric-\allowbreak{}space} & f is an isometry from a to b & generated & proven \\
\texttt{metrizable-\allowbreak{}top-\allowbreak{}space} & f is continuous from a to b & by \texttt{declare-\allowbreak{}hom!} & proven \\
\texttt{module} & f is a homomorphism from a to b & generated & proven \\
\texttt{monoid} & f is a homomorphism from a to b & generated & proven \\
\texttt{normed-\allowbreak{}ag} & f is a homomorphism from a to b & generated & proven \\
\texttt{normed-\allowbreak{}field} & f is a homomorphism from a to b & generated & proven \\
\texttt{normed-\allowbreak{}vector-\allowbreak{}space} & f is a homomorphism from a to b & generated & proven \\
\texttt{nu-\allowbreak{}c-\allowbreak{}star-\allowbreak{}algebra} & f is a homomorphism from a to b & refinement of \texttt{nu-\allowbreak{}involutive-\allowbreak{}banach-\allowbreak{}algebra} & id: not installed; compose: not installed; in-set: not installed \\
\texttt{nu-\allowbreak{}involutive-\allowbreak{}banach-\allowbreak{}algebra} & f is a homomorphism from a to b & generated & id: not installed; compose: not installed; in-set: not installed \\
\texttt{pid} & f is a homomorphism from a to b & refinement of \texttt{integral-\allowbreak{}domain} & proven \\
\texttt{pk-\allowbreak{}space} & f is a morphism of public-key strand spaces from a to b & generated & id: not installed; compose: not installed; in-set: not installed \\
\texttt{pre-\allowbreak{}hilbert-\allowbreak{}space} & f is a homomorphism from a to b & refinement of \texttt{semi-\allowbreak{}inner-\allowbreak{}product-\allowbreak{}space} & id: not installed; compose: not installed; in-set: not installed \\
\texttt{pseudometric-\allowbreak{}space} & f is a homomorphism from a to b & generated & proven \\
\texttt{ring} & f is a homomorphism from a to b & generated & proven \\
\texttt{ringoid} & f is a ringoid homomorphism from a to b & by \texttt{declare-\allowbreak{}hom!} & proven \\
\texttt{semi-\allowbreak{}inner-\allowbreak{}product-\allowbreak{}space} & f is a homomorphism from a to b & generated & id: not installed; compose: not installed; in-set: not installed \\
\texttt{semigroup} & f is a homomorphism from a to b & generated & proven \\
\texttt{setoid} & f is a setoid morphism from a to b & generated & proven \\
\texttt{sk-\allowbreak{}space} & f is a morphism of shared-key strand spaces from a to b & generated & id: not installed; compose: not installed; in-set: not installed \\
\texttt{strand-\allowbreak{}space} & f1 is a morphism of strand spaces from a to b & generated & id: not installed; compose: not installed; in-set: not installed \\
\texttt{top-\allowbreak{}group} & f is a continuous homomorphism from a to b & generated & id: not installed; compose: not installed; in-set: not installed \\
\texttt{top-\allowbreak{}space} & f is continuous from a to b & generated & proven \\
\texttt{unital-\allowbreak{}banach-\allowbreak{}algebra} & f is a homomorphism from a to b & generated & id: not installed; compose: not installed; in-set: not installed \\
\texttt{unital-\allowbreak{}c-\allowbreak{}star-\allowbreak{}algebra} & f is a homomorphism from a to b & refinement of \texttt{unital-\allowbreak{}involutive-\allowbreak{}banach-\allowbreak{}algebra} & id: not installed; compose: not installed; in-set: not installed \\
\texttt{unital-\allowbreak{}involutive-\allowbreak{}banach-\allowbreak{}algebra} & f is a homomorphism from a to b & generated & id: not installed; compose: not installed; in-set: not installed \\
\texttt{vector-\allowbreak{}space} & f is a homomorphism from a to b & refinement of \texttt{module} & proven \\
\bottomrule
\end{longtable}

\subsubsection{The declared categories}

\paragraph{\texttt{bounded-\allowbreak{}linear,\allowbreak{} on normed-\allowbreak{}vector-\allowbreak{}space}}
\emph{Declared in} \texttt{theorem-\allowbreak{}library/\allowbreak{}categories.scm}.

Objects: the models of \texttt{is-\allowbreak{}normed-\allowbreak{}vector-\allowbreak{}space}.  Arrows: \texttt{is-\allowbreak{}bounded-\allowbreak{}linear-\allowbreak{}arrow(a,\allowbreak{} b,\allowbreak{} f)} --- f is a bounded linear map from a to b.  Hom-set: \texttt{hom-\allowbreak{}bounded-\allowbreak{}linear(a,\allowbreak{} b)} --- the bounded linear maps from a to b.

Definition (\texttt{is-\allowbreak{}bounded-\allowbreak{}linear-\allowbreak{}arrow-\allowbreak{}def}, definitional):
\begin{VNBsnip}
forall([a, b, f],
  is-bounded-linear-arrow(a, b, f)
  iff is-normed-vector-space(a)
      and is-normed-vector-space(b)
      and f in fun(vec(a), vec(b))
      and scal(a) = scal(b)
      and forall([x1_ in vec(a), x2_ in vec(a)],
            f((vadd(a))(x1_, x2_)) = (vadd(b))(f(x1_), f(x2_)))
      and forall([x1_ in carr(scal(a)), x2_ in vec(a)],
            f((act(a))(x1_, x2_)) = (act(b))(x1_, f(x2_)))
      and forsome([blk_ in rr],
            0 <= blk_
            and forall([x1_ in vec(a)],
                  (vnrm(b))(f(x1_)) <= blk_ * (vnrm(a))(x1_))))
\end{VNBsnip}

Laws: \texttt{hom-\allowbreak{}bounded-\allowbreak{}linear-\allowbreak{}id} --- proven; \texttt{hom-\allowbreak{}bounded-\allowbreak{}linear-\allowbreak{}compose} --- proven; \texttt{hom-\allowbreak{}bounded-\allowbreak{}linear-\allowbreak{}in-\allowbreak{}set} --- proven; membership \texttt{hom-\allowbreak{}bounded-\allowbreak{}linear-\allowbreak{}member-\allowbreak{}iff} --- proven; Hom functors \texttt{hom-\allowbreak{}bounded-\allowbreak{}linear-\allowbreak{}post-\allowbreak{}type} --- proven, \texttt{hom-\allowbreak{}bounded-\allowbreak{}linear-\allowbreak{}pre-\allowbreak{}type} --- proven.

\paragraph{\texttt{c-\allowbreak{}continuous,\allowbreak{} on c-\allowbreak{}metric-\allowbreak{}space}}
\emph{Declared in} \texttt{theorem-\allowbreak{}library/\allowbreak{}c-\allowbreak{}continuous-\allowbreak{}category.scm}.

Objects: the models of \texttt{is-\allowbreak{}c-\allowbreak{}metric-\allowbreak{}space}.  Arrows: \texttt{is-\allowbreak{}c-\allowbreak{}continuous-\allowbreak{}arrow(a,\allowbreak{} b,\allowbreak{} f)} --- f is a continuous map of countably-metrised spaces from a to b.  Hom-set: \texttt{hom-\allowbreak{}c-\allowbreak{}continuous(a,\allowbreak{} b)} --- the continuous maps of countably-metrised spaces from a to b.

Definition (\texttt{is-\allowbreak{}c-\allowbreak{}continuous-\allowbreak{}arrow-\allowbreak{}def}, definitional):
\begin{VNBsnip}
forall([a, b, f],
  is-c-continuous-arrow(a, b, f)
  iff is-c-metric-space(a)
      and is-c-metric-space(b)
      and f in fun(pts(a), pts(b))
      and is-continuous(c-metric(a), c-metric(b), f))
\end{VNBsnip}

Laws: \texttt{hom-\allowbreak{}c-\allowbreak{}continuous-\allowbreak{}id} --- proven; \texttt{hom-\allowbreak{}c-\allowbreak{}continuous-\allowbreak{}compose} --- proven; \texttt{hom-\allowbreak{}c-\allowbreak{}continuous-\allowbreak{}in-\allowbreak{}set} --- proven; membership \texttt{hom-\allowbreak{}c-\allowbreak{}continuous-\allowbreak{}member-\allowbreak{}iff} --- proven; Hom functors \texttt{hom-\allowbreak{}c-\allowbreak{}continuous-\allowbreak{}post-\allowbreak{}type} --- proven, \texttt{hom-\allowbreak{}c-\allowbreak{}continuous-\allowbreak{}pre-\allowbreak{}type} --- proven.

\paragraph{\texttt{continuous,\allowbreak{} on metric-\allowbreak{}space}}
\emph{Declared in} \texttt{theorem-\allowbreak{}library/\allowbreak{}categories.scm}.

Objects: the models of \texttt{is-\allowbreak{}metric-\allowbreak{}space}.  Arrows: \texttt{is-\allowbreak{}continuous-\allowbreak{}arrow(a,\allowbreak{} b,\allowbreak{} f)} --- f is a continuous map from a to b.  Hom-set: \texttt{hom-\allowbreak{}continuous(a,\allowbreak{} b)} --- the continuous maps from a to b.

Definition (\texttt{is-\allowbreak{}continuous-\allowbreak{}arrow-\allowbreak{}def}, definitional):
\begin{VNBsnip}
forall([a, b, f],
  is-continuous-arrow(a, b, f)
  iff is-metric-space(a)
      and is-metric-space(b)
      and f in fun(pts(a), pts(b))
      and is-continuous(a, b, f))
\end{VNBsnip}

Laws: \texttt{hom-\allowbreak{}continuous-\allowbreak{}id} --- proven; \texttt{hom-\allowbreak{}continuous-\allowbreak{}compose} --- proven; \texttt{hom-\allowbreak{}continuous-\allowbreak{}in-\allowbreak{}set} --- proven; membership \texttt{hom-\allowbreak{}continuous-\allowbreak{}member-\allowbreak{}iff} --- proven; Hom functors \texttt{hom-\allowbreak{}continuous-\allowbreak{}post-\allowbreak{}type} --- proven, \texttt{hom-\allowbreak{}continuous-\allowbreak{}pre-\allowbreak{}type} --- proven.

\paragraph{\texttt{lipschitz,\allowbreak{} on metric-\allowbreak{}space}}
\emph{Declared in} \texttt{theorem-\allowbreak{}library/\allowbreak{}categories.scm}.

Objects: the models of \texttt{is-\allowbreak{}metric-\allowbreak{}space}.  Arrows: \texttt{is-\allowbreak{}lipschitz-\allowbreak{}arrow(a,\allowbreak{} b,\allowbreak{} f)} --- f is a Lipschitz map from a to b.  Hom-set: \texttt{hom-\allowbreak{}lipschitz(a,\allowbreak{} b)} --- the Lipschitz maps from a to b.

Definition (\texttt{is-\allowbreak{}lipschitz-\allowbreak{}arrow-\allowbreak{}def}, definitional):
\begin{VNBsnip}
forall([a, b, f],
  is-lipschitz-arrow(a, b, f)
  iff is-metric-space(a)
      and is-metric-space(b)
      and f in fun(pts(a), pts(b))
      and forsome([lpk_ in rr],
            0 <= lpk_
            and forall([lpx_ in pts(a), lpy_ in pts(a)],
                  (dist(b))(f(lpx_), f(lpy_)) <= lpk_ * (dist(a))(lpx_, lpy_))))
\end{VNBsnip}

Laws: \texttt{hom-\allowbreak{}lipschitz-\allowbreak{}id} --- proven; \texttt{hom-\allowbreak{}lipschitz-\allowbreak{}compose} --- proven; \texttt{hom-\allowbreak{}lipschitz-\allowbreak{}in-\allowbreak{}set} --- proven; membership \texttt{hom-\allowbreak{}lipschitz-\allowbreak{}member-\allowbreak{}iff} --- proven; Hom functors \texttt{hom-\allowbreak{}lipschitz-\allowbreak{}post-\allowbreak{}type} --- proven, \texttt{hom-\allowbreak{}lipschitz-\allowbreak{}pre-\allowbreak{}type} --- proven.

\paragraph{\texttt{uniformly-\allowbreak{}continuous,\allowbreak{} on metric-\allowbreak{}space}}
\emph{Declared in} \texttt{theorem-\allowbreak{}library/\allowbreak{}categories.scm}.

Objects: the models of \texttt{is-\allowbreak{}metric-\allowbreak{}space}.  Arrows: \texttt{is-\allowbreak{}uniformly-\allowbreak{}continuous-\allowbreak{}arrow(a,\allowbreak{} b,\allowbreak{} f)} --- f is a uniformly continuous map from a to b.  Hom-set: \texttt{hom-\allowbreak{}uniformly-\allowbreak{}continuous(a,\allowbreak{} b)} --- the uniformly continuous maps from a to b.

Definition (\texttt{is-\allowbreak{}uniformly-\allowbreak{}continuous-\allowbreak{}arrow-\allowbreak{}def}, definitional):
\begin{VNBsnip}
forall([a, b, f],
  is-uniformly-continuous-arrow(a, b, f)
  iff is-metric-space(a)
      and is-metric-space(b)
      and f in fun(pts(a), pts(b))
      and is-uniformly-continuous(a, b, f))
\end{VNBsnip}

Laws: \texttt{hom-\allowbreak{}uniformly-\allowbreak{}continuous-\allowbreak{}id} --- proven; \texttt{hom-\allowbreak{}uniformly-\allowbreak{}continuous-\allowbreak{}compose} --- proven; \texttt{hom-\allowbreak{}uniformly-\allowbreak{}continuous-\allowbreak{}in-\allowbreak{}set} --- proven; membership \texttt{hom-\allowbreak{}uniformly-\allowbreak{}continuous-\allowbreak{}member-\allowbreak{}iff} --- proven; Hom functors \texttt{hom-\allowbreak{}uniformly-\allowbreak{}continuous-\allowbreak{}post-\allowbreak{}type} --- proven, \texttt{hom-\allowbreak{}uniformly-\allowbreak{}continuous-\allowbreak{}pre-\allowbreak{}type} --- proven.

\subsubsection{Inclusions}
Theorems of the form $P(a, b, f) \Rightarrow Q(a, b, f)$ for two arrow predicates, or $H_1(a, b) \subseteq H_2(a, b)$ for two hom-sets, found by shape.

\begin{itemize}
\item \texttt{hom-\allowbreak{}metric-\allowbreak{}space-\allowbreak{}is-\allowbreak{}lipschitz-\allowbreak{}arrow} (proven): \texttt{is-\allowbreak{}hom-\allowbreak{}metric-\allowbreak{}space(a,\allowbreak{} b,\allowbreak{} f) implies is-\allowbreak{}lipschitz-\allowbreak{}arrow(a,\allowbreak{} b,\allowbreak{} f)}
\item \texttt{hom-\allowbreak{}normed-\allowbreak{}vector-\allowbreak{}space-\allowbreak{}is-\allowbreak{}bounded-\allowbreak{}linear-\allowbreak{}arrow} (proven): \texttt{is-\allowbreak{}hom-\allowbreak{}normed-\allowbreak{}vector-\allowbreak{}space(a,\allowbreak{} b,\allowbreak{} f) implies is-\allowbreak{}bounded-\allowbreak{}linear-\allowbreak{}arrow(a,\allowbreak{} b,\allowbreak{} f)}
\item \texttt{lipschitz-\allowbreak{}arrow-\allowbreak{}is-\allowbreak{}uniformly-\allowbreak{}continuous-\allowbreak{}arrow} (proven): \texttt{is-\allowbreak{}lipschitz-\allowbreak{}arrow(a,\allowbreak{} b,\allowbreak{} f) implies is-\allowbreak{}uniformly-\allowbreak{}continuous-\allowbreak{}arrow(a,\allowbreak{} b,\allowbreak{} f)}
\item \texttt{uniformly-\allowbreak{}continuous-\allowbreak{}arrow-\allowbreak{}is-\allowbreak{}continuous-\allowbreak{}arrow} (proven): \texttt{is-\allowbreak{}uniformly-\allowbreak{}continuous-\allowbreak{}arrow(a,\allowbreak{} b,\allowbreak{} f) implies is-\allowbreak{}continuous-\allowbreak{}arrow(a,\allowbreak{} b,\allowbreak{} f)}
\end{itemize}
\par\endgroup


\section{Viewing one structure as another: \texttt{view-as}}
\label{sec:view-as}

A ring's additive part $(\mathtt{carr}, \mathtt{add}, \mathtt{zero},
\mathtt{neg})$ is an abelian group; its multiplicative part
$(\mathtt{carr}, \mathtt{mul}, \mathtt{one})$ is a monoid.  A
\emph{view-as} records such a relationship: a way of \emph{seeing} an
instance of one structure as an instance of another.  The relationship
is asymmetric (the ring's additive group forgets the multiplication;
there is no route back) and not unique (the same ring has both an
additive-group view and a multiplicative-monoid view).  We call these
\emph{views}, not ``bridges'': a bridge sounds symmetric, and the
asymmetry and the multiplicity are exactly the point.

A view is a \emph{functoid} (Remark~\ref{rem:functoids-quantified}):
\texttt{def-functor} installs the \texttt{def-functoid} that builds the
target instance as the list of the source's slots, plus the definitional
axiom that the result satisfies the target predicate.  Its domain, the
class of all instances of the source structure, is a proper class; a
view is a member of nothing and, like every functoid, outside the range
of the quantifiers.

\subsection{The \texttt{def-functor} form}

\begin{VNBcode}
(def-functor NAME
  SOURCE-STRUCT  (src-comp-1 src-comp-2 ... src-comp-n)
  TARGET-STRUCT  (tgt-slot-1 tgt-slot-2 ... tgt-slot-n))
\end{VNBcode}

The two component lists are flat lists of accessor symbols of equal
length, matched position by position: the value at the target's $k$-th
slot is the source's $k$-th named component.  The target list must be
the target structure's slots in their declared order (the form checks
this); the source list names, in that same order, which source accessor
fills each target slot.  So a view is a \textbf{slot map written in the
target's coordinate system}.

A view-as is a \emph{forgetful functor} between structure categories, and a
forgetful functor does exactly one or both of two things (recall
\S\ref{sec:shape-vs-property}: a structure is a \emph{shape} plus
\emph{properties}):

\begin{itemize}
\item \textbf{forget shape} --- the source list keeps only a sub-list of the
  source's slots; the rest drop.  (A ring's additive group forgets
  \texttt{mul}/\texttt{one}; a field's multiplicative group keeps the
  \texttt{non-zero} carrier slot and forgets the additive ones.)
\item \textbf{forget properties} --- \texttt{is-\textit{target}} folds fewer laws than
  \texttt{is-\textit{source}}, so the image satisfies less (a commutative ring viewed
  as a ring forgets commutativity).
\end{itemize}

\noindent These are the \emph{only} two forgetful operations.  Writing the kept
slots in the target's coordinate order (a ring's \texttt{add} filling a group's
\texttt{mul} slot, or a non-first slot serving as the carrier) is just
notation, not a third kind --- it is what lets shape-forgetting be expressed at
all.  Genuinely \emph{non}-forgetful functors --- the quotient, the completion
--- are constructions, not views; they are written as plain functoids.

\begin{VNBcode}
;;; A ring's additive structure is an abelian group.
;;; ABELIAN-GROUP's slots, in order, are (CARR OPR IDEN INV); the view fills
;;; them from the ring's (CARR ADD ZERO NEG):
(def-functor 'RING-ADDITIVE-AG
  'RING          '(CARR ADD ZERO NEG)
  'ABELIAN-GROUP '(CARR OPR IDEN INV))
\end{VNBcode}

\subsection{What \texttt{def-functor} installs}

\begin{enumerate}
\item \textbf{A constructor functoid} \texttt{(\textit{name} s)} whose defining
  equation gathers the named source components into a target-shaped list:
  \[
    \mathtt{(ring\text{-}additive\text{-}ag}\;r) \;=\;
      \mathtt{list}(\mathtt{carr}(r),\, \mathtt{add}(r),\,
                    \mathtt{zero}(r),\, \mathtt{neg}(r)).
  \]
  Inside a proof, \texttt{(mul (ring-additive-ag r))} reduces to
  \texttt{(add r)} by unfolding the functoid and the accessor.

\item \textbf{A typing axiom}, named
  \texttt{\textit{name}-is-\textit{target}} (here
  \texttt{ring-additive-ag-is-abelian-group}):
  \[
    \forall r.\; \mathit{is\text{-}source}(r) \;\Rightarrow\;
      \mathit{is\text{-}target}(\mathit{name}(r)).
  \]
  During the library-build phase this is installed as an \emph{axiom}
  (a warranted assertion, Section~\ref{sec:warrants}), not proved;
  promoting it to a theorem is a separate task.

\item \textbf{Auto-specializations.}  Every theorem of the form
  $\forall s.\; \mathit{is\text{-}target}(s) \Rightarrow P[s]$ already in
  the library is re-installed as
  $\forall r.\; \mathit{is\text{-}source}(r) \Rightarrow P[\mathit{name}(r)]$,
  with the target-slot accessors reduced to source components (so
  $\mathtt{mul}(s)$ becomes $\mathtt{add}(r)$, etc.), under the name
  \texttt{\textit{original}-\textit{name}}.  Thus declaring the additive
  view immediately carries every abelian-group theorem over to the
  additive structure of every ring.
\end{enumerate}

\begin{remark}
A view-as is a covariant functor in waiting: the constructor is the
object map (an instance of the source becomes an instance of the
target), and the auto-specialization transports the target's theorems
along it.  Homomorphisms are now first-class
(Section~\ref{sec:morphisms}); the morphism half --- transporting a
homomorphism across a view --- remains deferred.  The
auto-specializer currently matches only theorems guarded by exactly
\texttt{is-\textit{target}}; it does not chain through structure-inclusion axioms,
so same-shape refinements (a commutative ring's additive view, say) are
declared directly rather than derived.
\end{remark}

\begin{remark}
Same-shape inclusions --- every abelian group is a group, every
commutative ring is a ring --- are \emph{not} views.  They are recorded
as plain $\texttt{is-}X \Rightarrow \texttt{is-}Y$ typing axioms, because
the constructor would be a redundant identity list.  A \texttt{view-as}
earns its keep when source and target have different shapes, so the
constructor genuinely projects.
\end{remark}

\subsection{Constructed functors: \texttt{def-constructed-functor}}
\label{sec:constructed-functor}

A \texttt{view-as} selects the target's slots from among the source's --- its
object map is a re-shuffling of accessors.  Some functors are not like that.
The metric topology $\tau_d$ is \emph{computed} from the distance, not read off a
slot; the distance of a normed field is \emph{computed} from the norm.  For these
the object map is an arbitrary term, and the form is
\texttt{def-constructed-functor}:

\begin{VNBcode}
;;; A normed field is a metric space, with d(x, y) = norm(x - y).
(def-constructed-functor 'NF-METRIC-SPACE 'NORMED-FIELD 'METRIC-SPACE '(nf)
  '(LIST (CARR nf)
         (VNB-LAMBDA (LIST x y) (CARTESIAN (CARR nf) (CARR nf)) ((FNRM nf) ((ADD nf) x ((NEG nf) y))))))
\end{VNBcode}

The arguments are the functor name, its source and target species, the
parameter list, and the defining term --- installed as an ordinary functoid
(Section~\ref{sec:def-functoid}), so \texttt{(nf-metric-space nf)} rewrites to
that tuple.

The decisive difference from \texttt{view-as} is what it \emph{claims}.  A view's
typing axiom is installed as a warranted assertion; a constructed functor
\textbf{asserts nothing}.  ``An isometry is continuous'' and ``the norm-difference
is a metric'' are theorems with content, not free consequences of a slot
correspondence.  So the form only \emph{records} the two obligations any functor
owes, and \texttt{functor-obligation-audit} reports at load any still
undischarged:
\[
  \begin{aligned}
    \texttt{\textit{name}-is-\textit{tgt}} &:\quad
      \mathit{is\text{-}src}(r) \Rightarrow \mathit{is\text{-}tgt}(\mathit{name}(r)), \\
    \texttt{\textit{name}-functorial} &:\quad
      \mathit{is\text{-}hom\text{-}src}(a, b, f) \Rightarrow
      \mathit{is\text{-}hom\text{-}tgt}(\mathit{name}(a), \mathit{name}(b), f).
  \end{aligned}
\]
Both are available as goals, ready for \texttt{sp}.  A functor whose obligations
are unproved is a functor you do not yet have.

%%% -----------------------------------------------------------------------
\section{Monoid products and ring sums}
\label{sec:sequences}

\texttt{sequences.scm} defines two constants by primitive recursion on
$\mathtt{nn}$:

\begin{itemize}
\item $\mathit{prod\text{-}ord}(m, f, n)$ --- product of $f(0) * f(1) * \cdots * f(n{-}1)$
  in a monoid $m$.
\item $\mathtt{sum}(r, f, n)$ --- sum $f(0) + f(1) + \cdots + f(n{-}1)$
  in a ring $r$.
\end{itemize}

Both are $\mathtt{nn}$-indexed: the third argument $n$ counts how many
terms are included.  The empty product ($n = 0$) is the identity;
$\mathtt{sum}(r, f, 0) = \mathtt{zero}(r)$.

\subsection{\texttt{prod-ord}: monoid product}

\paragraph{Signature.}
$\mathit{prod\text{-}ord}(m, f, n)$ where $m$ is a \textsc{Monoid},
$f \in \mathtt{fun}(\mathtt{nn}, \mathtt{carr}(m))$, $n \in \mathtt{nn}$.

\paragraph{Defining recursion.}
Installed by \texttt{def-by-nn-recursion} as two named axioms:

\begin{align*}
\mathtt{prod\text{-}ord\text{-}zero}: &\quad
  \forall m\, f.\;
  \mathit{prod\text{-}ord}(m, f, 0) = \mathtt{id}(m) \\
\mathtt{prod\text{-}ord\text{-}succ}: &\quad
  \forall m\, f\, n \in \mathtt{nn}.\;
  \mathit{prod\text{-}ord}(m, f, \mathtt{succ}(n))
  = \mathit{prod\text{-}ord}(m, f, n) * f(n)
\end{align*}

where $*$ denotes $\mathtt{mul}(m)$.

\paragraph{Additional theorems.}
\begin{description}
\item[\texttt{prod-ord-type}]
  $\mathit{is\text{-}monoid}(m) \wedge f \in \mathtt{fun}(\mathtt{nn}, \mathtt{carr}(m))
  \wedge n \in \mathtt{nn} \;\Rightarrow\;
  \mathit{prod\text{-}ord}(m, f, n) \in \mathtt{carr}(m)$.
  A theorem (\texttt{theorem-library/rake-asm-2.scm}), proved from
  \texttt{prod-ord-zero}, \texttt{prod-ord-succ}, \texttt{monoid-identity-in},
  and \texttt{monoid-carrier-closed-opr} by $\mathtt{nn}$-induction
  (see Section~\ref{sec:ex-ni-type}).

\item[\texttt{prod-ord-singleton}]
  $\mathit{is\text{-}monoid}(m) \wedge f \in \mathtt{fun}(\mathtt{nn}, \mathtt{carr}(m))
  \;\Rightarrow\; \mathit{prod\text{-}ord}(m, f, 1) = f(0)$.\\
  (Derivable from \texttt{prod-ord-succ} at $n=0$, \texttt{prod-ord-zero},
  and \texttt{monoid-left-id}.)
\end{description}

\subsection{\texttt{sum}: ring summation}

\texttt{sum} is the special case of \texttt{prod-ord} for the additive
monoid of a ring---\textit{i.e.}, $\mathtt{mul} \to \mathtt{add}$,
$\mathtt{id} \to \mathtt{zero}$---defined independently for direct use.

\paragraph{Signature.}
$\mathtt{sum}(r, f, n)$ where $r$ is a \textsc{Ring},
$f \in \mathtt{fun}(\mathtt{nn}, \mathtt{carr}(r))$, $n \in \mathtt{nn}$.

\paragraph{Defining recursion.}

\begin{align*}
\mathtt{sum\text{-}zero}: &\quad
  \forall r\, f.\; \mathtt{sum}(r, f, 0) = \mathtt{zero}(r) \\
\mathtt{sum\text{-}succ}: &\quad
  \forall r\, f\, n \in \mathtt{nn}.\;
  \mathtt{sum}(r, f, \mathtt{succ}(n))
  = \mathtt{sum}(r, f, n) + f(n)
\end{align*}

\paragraph{Additional theorems.}
\begin{description}
\item[\texttt{sum-type}]
  $\mathit{is\text{-}ring}(r) \wedge f \in \mathtt{fun}(\mathtt{nn}, \mathtt{carr}(r))
  \wedge n \in \mathtt{nn} \;\Rightarrow\; \mathtt{sum}(r, f, n) \in \mathtt{carr}(r)$.
  A theorem (\texttt{theorem-library/rake-asm-2.scm}), proved from
  \texttt{sum-zero}, \texttt{sum-succ}, \texttt{ring-zero-in},
  and \texttt{ring-carrier-closed-add} by $\mathtt{nn}$-induction
  (Section~\ref{sec:ex-ni-type}).

\item[\texttt{sum-singleton}]
  $\mathit{is\text{-}ring}(r) \wedge f \in \mathtt{fun}(\mathtt{nn}, \mathtt{carr}(r))
  \;\Rightarrow\; \mathtt{sum}(r, f, 1) = f(0)$.

\item[\texttt{sum-left-scalar}]
  $a \in \mathtt{carr}(r) \wedge f \in \mathtt{fun}(\mathtt{nn}, \mathtt{carr}(r))
  \wedge n \in \mathtt{nn} \;\Rightarrow\;$\\
  $\mathtt{sum}(r,\, \lambda i.\, a \cdot f(i),\, n)
  = a \cdot \mathtt{sum}(r, f, n)$.\\
  (Derivable by $\mathtt{nn}$-induction using \texttt{ring-left-dist}.)
\end{description}

\begin{remark}
The unordered finite sum $\sum_{x \in S} f(x)$ over an arbitrary finite
set $S$ (not indexed by $\mathtt{nn}$) requires commutativity of
addition and a well-definedness proof (independence of the enumeration
bijection from \texttt{card-finite-bij}).  This more general
$\mathit{sum\text{-}set}$ operation is implemented (\texttt{sequences.scm}); see
Section~\ref{sec:sum-set} below.
\end{remark}

\begin{remark}
The connection $\mathtt{sum}(r, f, n) = \mathit{prod\text{-}ord}(m, f, n)$
where $m = [\mathtt{carr}(r),\, \mathtt{add}(r),\, \mathtt{zero}(r)]$ is the
additive monoid of $r$ is a theorem (by structural induction on $n$),
not an axiom.
\end{remark}

\paragraph{Typical proof pattern.}
To use \texttt{sum} in a proof, unfold the recursion manually with
\texttt{(ta 'sum-succ)} and \texttt{(ta 'sum-zero)}, then close
numeric identities with \texttt{(rs)} or \texttt{(arith)}.

\begin{VNBcode}
;; Prove SUM(r, f, 2) = f(0) + f(1)
(sp (make-wff-from-string
  "forall([r, f], IS-RING(r) and f in FUN(NN, CARR(r)) implies
     SUM(r, f, 2) = ADD(r)(SUM(r, f, 1), f(1))"))
(ta 'sum-succ)    ; unfold outer step
(ta 'sum-singleton)
(rfl)
\end{VNBcode}

%% =========================================================
\section{Set-indexed sums and products}
\label{sec:sum-set-prod-set}

\texttt{sum} and \texttt{prod-ord} index by an initial segment
$\{0, \ldots, n-1\}$ of $\mathtt{nn}$, so the order in which the
operation is applied is built into the definition.  Once the operation
is commutative, the value depends only on the \emph{set} of indices.
\texttt{sum-set} and \texttt{prod-set} package this directly.

\subsection{\texttt{sum-set}: ring sum over a finite set}
\label{sec:sum-set}

\[
   \mathit{sum\text{-}set}(r,\, S,\, f) \;=\; \sum_{x \in S} f(x)
\]
The additive monoid of a ring is commutative
(\texttt{ring-add-comm}), so the value depends only on $S$, not on any
enumeration.  \texttt{sum-set} is a \texttt{def-functoid}:
$\mathit{sum\text{-}set}(r, S, f) = \mathit{finsum}(\mathit{ring\text{-}additive\text{-}ag}(r), f, S)$
(\texttt{structure-library/finsum.scm}).  Its four laws are theorems, not axioms
(\texttt{sum-set-empty}, \texttt{sum-set-singleton-defined},
\texttt{sum-set-disjoint-union-defined}, \texttt{sum-set-type-defined}, in
\texttt{theorem-library/rake-sum-set-defined.scm}), each in its strongest true
form: the statements below omit the typing and finiteness hypotheses on the
index sets that the theorems carry.
\begin{description}
\item[\texttt{sum-set-empty}]
  $\mathit{sum\text{-}set}(r,\, \emptyset,\, f) = \mathtt{zero}(r)$.
\item[\texttt{sum-set-singleton}]
  $\mathit{is\text{-}ring}(r) \Rightarrow
   \mathit{sum\text{-}set}(r,\, \{x\},\, f) = f(x)$.
\item[\texttt{sum-set-disjoint-union}]
  $\mathit{is\text{-}ring}(r) \land S_1, S_2 \in \mathtt{set}
   \land S_1 \cap S_2 = \emptyset \land f \in \mathtt{fun}(S_1 \cup S_2,\, \mathtt{carr}(r))
   \Rightarrow$
  \[
    \mathit{sum\text{-}set}(r,\, S_1 \cup S_2,\, f)
    \;=\; \mathtt{add}(r)(\mathit{sum\text{-}set}(r, S_1, f),\,
                          \mathit{sum\text{-}set}(r, S_2, f))
  \]
\item[\texttt{sum-set-type}]
  $\mathit{is\text{-}ring}(r) \land X \in \mathtt{set}
   \land S \subseteq X \land f \in \mathtt{fun}(X,\, \mathtt{carr}(r))
   \Rightarrow \mathit{sum\text{-}set}(r, S, f) \in \mathtt{carr}(r)$.
  \\
  The explicit superset $X$ keeps the typing applicable in the
  disjoint-union axiom: $f$ is typed on $S_1 \cup S_2$, and the RHS
  uses sums over $S_1, S_2$ separately --- each $S_i \subseteq X$
  with $X = S_1 \cup S_2$.
\end{description}

\subsection{\texttt{prod-set}: commutative-monoid product over a finite set}

\[
   \mathit{prod\text{-}set}(\mathit{cm},\, S,\, f) \;=\; \prod_{x \in S} f(x)
\]
Mirrors \texttt{sum-set} but in a commutative monoid: it is the
\texttt{def-functoid} $\mathit{prod\text{-}set}(\mathit{cm}, S, f) =
\mathit{finsum}(\mathit{cm}, f, S)$ (\texttt{structure-library/finprod.scm}), and
its four laws are theorems (\texttt{prod-set-empty},
\texttt{prod-set-singleton-defined}, \texttt{prod-set-disjoint-union-defined},
\texttt{prod-set-type-defined}, in \texttt{theorem-library/rake-prod-set-defined.scm}),
with the same hypotheses omitted below.
\begin{description}
\item[\texttt{prod-set-empty}]
  $\mathit{prod\text{-}set}(\mathit{cm},\, \emptyset,\, f) = \mathtt{id}(\mathit{cm})$.
\item[\texttt{prod-set-singleton}]
  $\mathit{is\text{-}comm\text{-}monoid}(\mathit{cm}) \Rightarrow
   \mathit{prod\text{-}set}(\mathit{cm},\, \{x\},\, f) = f(x)$.
\item[\texttt{prod-set-disjoint-union}]
  $\mathit{is\text{-}comm\text{-}monoid}(\mathit{cm}) \land
   S_1, S_2 \in \mathtt{set} \land S_1 \cap S_2 = \emptyset \land
   f \in \mathtt{fun}(S_1 \cup S_2,\, \mathtt{carr}(\mathit{cm})) \Rightarrow$
  \[
    \mathit{prod\text{-}set}(\mathit{cm},\, S_1 \cup S_2,\, f)
    \;=\; \mathtt{mul}(\mathit{cm})(\mathit{prod\text{-}set}(\mathit{cm}, S_1, f),\,
                                    \mathit{prod\text{-}set}(\mathit{cm}, S_2, f))
  \]
\item[\texttt{prod-set-type}]
  $\mathit{is\text{-}comm\text{-}monoid}(\mathit{cm}) \land X \in \mathtt{set}
   \land S \subseteq X \land f \in \mathtt{fun}(X,\, \mathtt{carr}(\mathit{cm}))
   \Rightarrow \mathit{prod\text{-}set}(\mathit{cm}, S, f) \in \mathtt{carr}(\mathit{cm}).$
\end{description}

\begin{remark}
The connection $\mathtt{sum}(r, f, n) =
\mathit{sum\text{-}set}(r,\, \mathit{ord\text{-}segment}(n),\, f)$
relates the two presentations: \texttt{sum} indexes by the canonical
initial segment, \texttt{sum-set} by an arbitrary finite set.
The equation is a theorem (by $\mathtt{nn}$-induction), not an axiom.
A general unordered $\mathit{sum\text{-}set}(r, S, f)$ can be
\emph{computed} via any enumeration $\mathit{bij}:\{0,\ldots,|S|-1\}\to S$
(existence by \texttt{card-finite-bij}) as
$\mathtt{sum}(r,\, \lambda i.\, f(\mathit{bij}(i)),\, |S|)$;
independence of the enumeration is the permutation lemma, a future
derived theorem.
\end{remark}

\begin{remark}
The three recursive laws (\texttt{*-empty},
\texttt{*-singleton}, \texttt{*-disjoint-union}) are stated separately so that direct
unfolding is available in proofs.  The disjointness condition
$S_1 \cap S_2 = \emptyset$ is essential: without it the right-hand side
double-counts $S_1 \cap S_2$.
\end{remark}

%% =========================================================
\section{Ring products}
\label{sec:ring-prod}

\texttt{ring.scm} defines the binary ring product, and
\texttt{sequences.scm} extends it to finite indexed products.
Both follow Path~2: the operations are \emph{total} functoids, defined
for any arguments; typing theorems fire conditionally when the inputs
are rings.

\subsection{Zero ring: \texttt{zero-ring}}

\texttt{zero-ring} is the terminal ring: carrier $\{0\}$, all operations
returning $0$, with $\mathtt{zero} = \mathtt{one} = 0$.

\begin{VNBsnip}
ZERO-RING = [{0},
             lambda([p, q], 0),   ; ADD
             lambda([p, q], 0),   ; MUL
             lambda([p], 0),      ; NEG
             0,                   ; ZERO
             0]                   ; ONE
\end{VNBsnip}

\begin{description}
\item[\texttt{zero-ring-def}]
  $\mathit{zero\text{-}ring} = [\{0\},\,
   \lambda(p,q).\,0,\; \lambda(p,q).\,0,\; \lambda(p).\,0,\; 0,\; 0]$.
\item[\texttt{zero-ring-is-ring}]
  $\mathit{is\text{-}ring}(\mathit{zero\text{-}ring})$.
  (A theorem, proved in \texttt{theorem-library/rake-asm-1.scm}.)
\end{description}

\begin{remark}
The zero ring is the terminal object in the category of rings: there is
exactly one ring homomorphism from \emph{any} ring $R$ to
$\mathit{zero\text{-}ring}$ (mapping everything to $0$).  It is also
the identity for \texttt{ring-prod} up to isomorphism:
$\mathit{ring\text{-}prod}(\mathit{zero\text{-}ring}, R) \cong R$.
\end{remark}

\subsection{Binary product: \texttt{ring-prod}}

\texttt{ring-prod(X,~Y)} is the product ring with carrier
$\mathtt{cartesian}(\mathtt{carr}(X),\,\mathtt{carr}(Y))$ and
componentwise operations.

\paragraph{Definition.}
Installed as a \texttt{def-functoid} rewrite:

\begin{VNBsnip}
RING-PROD(X, Y) =
  [CARTESIAN(CARR(X), CARR(Y)),
   lambda([p, q], [ADD(X)(nth(1,p), nth(1,q)),
                   ADD(Y)(nth(2,p), nth(2,q))]),
   lambda([p, q], [MUL(X)(nth(1,p), nth(1,q)),
                   MUL(Y)(nth(2,p), nth(2,q))]),
   lambda([p],    [NEG(X)(nth(1,p)),
                   NEG(Y)(nth(2,p))]),
   [ZERO(X), ZERO(Y)],
   [ONE(X),  ONE(Y)]]
\end{VNBsnip}

Elements of $\mathtt{carr}(\mathit{ring\text{-}prod}(X,Y))$ are pairs $[a,b]$
with $a \in \mathtt{carr}(X)$ and $b \in \mathtt{carr}(Y)$;
all six operations act componentwise.

\paragraph{Accessor expansion.}
The \texttt{ring-prod} macete chain is:
\[
  \mathtt{carr}(\mathit{ring\text{-}prod}(X,Y))
  \;\longrightarrow\; \mathtt{cartesian}(\mathtt{carr}(X),\,\mathtt{carr}(Y))
\]
and similarly for \texttt{add}, \texttt{mul}, \texttt{neg}, \texttt{zero},
\texttt{one}: each expands by first applying the \texttt{ring-prod}
rewrite and then the accessor macete.

\paragraph{Typing theorem.}
\begin{description}
\item[\texttt{ring-prod-is-ring}]
  $\mathit{is\text{-}ring}(X) \wedge \mathit{is\text{-}ring}(Y)
  \;\Rightarrow\; \mathit{is\text{-}ring}(\mathit{ring\text{-}prod}(X,Y))$.
  (A theorem, proved in \texttt{theorem-library/rake-asm-1.scm}.)
\end{description}

\begin{remark}[Totality]
\texttt{ring-prod} is total: $\mathit{ring\text{-}prod}(X,Y)$ is a
well-defined tuple for \emph{any} $X$ and $Y$.  When neither is a ring
--- when bongos are passed in --- the result is some tuple that satisfies
no ring properties.  \texttt{ring-prod-is-ring} is silent on this case (its
antecedent is false), so no contradiction arises.  This is the VNB
tradition: totality and typing are separate concerns.
\end{remark}

\subsection{Finite indexed product: \texttt{ring-prod-n}}

$\mathit{ring\text{-}prod\text{-}n}(f, n)$ is the $n$-fold product ring
$f(0) \times f(1) \times \cdots \times f(n-1)$, defined by primitive
recursion on~$n$.

\paragraph{Signature.}
$f$ is any function on $\mathtt{nn}$ (not required to map to rings;
the typing axiom states when the result is a ring).
$n \in \mathtt{nn}$.

\paragraph{Defining recursion.}
Installed by \texttt{def-by-nn-recursion}:

\begin{align*}
\mathtt{ring\text{-}prod\text{-}n\text{-}zero}: &\quad
  \forall f.\; \mathit{ring\text{-}prod\text{-}n}(f, 0) = \mathit{zero\text{-}ring} \\
\mathtt{ring\text{-}prod\text{-}n\text{-}succ}: &\quad
  \forall f,\, n \in \mathtt{nn}.\;
  \mathit{ring\text{-}prod\text{-}n}(f, \mathtt{succ}(n))
  = \mathit{ring\text{-}prod}(\mathit{ring\text{-}prod\text{-}n}(f, n),\; f(n))
\end{align*}

\paragraph{Typing theorem.}
\begin{description}
\item[\texttt{ring-prod-n-is-ring}]
  $(\forall i \in \mathtt{nn}.\;\mathit{is\text{-}ring}(f(i)))
  \;\wedge\; n \in \mathtt{nn}
  \;\Rightarrow\; \mathit{is\text{-}ring}(\mathit{ring\text{-}prod\text{-}n}(f,n))$.\\
  (A theorem, proved by $\mathtt{nn}$-induction from \texttt{zero-ring-is-ring}
  and \texttt{ring-prod-is-ring}, in \texttt{theorem-library/rake-asm-1.scm}.)
\end{description}

\begin{remark}
The indexed-family product $\prod_{i \in I} f(i)$ for an arbitrary
index set $I$ (not necessarily $\mathtt{nn}$) will use
\texttt{big-union($i$, $I$, carr($f(i)$))} as the function codomain.
The big-union binder is now in place
(Section~\ref{sec:class-set-discipline}); the family product is the next
step.
\end{remark}

%% =========================================================
\section{Basic rings: officialising $\mathtt{nn}, \mathtt{zz}, \mathtt{qq}, \mathtt{rr}, \mathtt{cc}$}
\label{sec:basic-rings}

The basic numeric domains $\mathtt{nn}, \mathtt{zz}, \mathtt{qq},
\mathtt{rr}, \mathtt{cc}$ ship with $n$-ary arithmetic operators
$(+\; x\; y\; z)$, $(*\; a\; b\; c\; d)$, decided by the
\texttt{arith-eval} oracle.  This is the form arithmetic is taught and
computed in: a flat sum or product of many arguments.

The algebraic structures \texttt{ring} and \texttt{comm-monoid}, by
contrast, require their operations to be \emph{binary} ---
elements of $\mathtt{fun}(\mathtt{cartesian}(A, A),\, A)$.  The file
\texttt{numeric-instances.scm} (loaded after \texttt{complex.scm}) bridges
the two presentations.

\begin{remark}
The duality is intentional.  We learned arithmetic as $n$-ary
operations; we learned abstract algebra with binary ones.  Both views
have their place: keep $(+\; x\; y\; z)$ for everyday calculation, but
when the goal is a generic ring/monoid theorem, route it through the
binary form.
\end{remark}

\subsection{Polymorphic binary operators}

Three named constants wrap the $n$-ary operators at arity 2 (and 1 for
negation).  They are polymorphic in the sense that the \emph{same}
\texttt{binplus} is typed for every numeric domain.

\begin{description}
\item[\texttt{binplus-apply}]
  $\forall x,\, y.\;\mathtt{binplus}(x, y) = x + y$.
\item[\texttt{bintimes-apply}]
  $\forall x,\, y.\;\mathtt{bintimes}(x, y) = x \cdot y$.
\item[\texttt{binneg-apply}]
  $\forall x.\;\mathtt{binneg}(x) = -x$.
\end{description}

There are no per-domain typing axioms for them.  The axioms
\texttt{binplus-in-fun-}$d$ that once gave them a \texttt{fun} typing
were deleted on 2026-08-29 (Section~\ref{sec:nums}, where the reason is given);
the \texttt{fun} typing that a structure instance needs is obtained inside
the instance proof by \texttt{lam-t}.

\subsection{Ring instances: $\mathtt{zz}, \mathtt{qq}, \mathtt{rr}, \mathtt{cc}$ as rings}

Each domain pairs with the binary operators above to form a
\texttt{ring} instance --- a 6-element VNB list:
\[
   \mathit{zz\text{-}ring} \;=\;
   [\,\mathtt{zz},\; \mathtt{binplus},\; \mathtt{bintimes},\;
     \mathtt{binneg},\; 0,\; 1\,]
\]
with definitional axiom \texttt{zz-ring-def}, and analogously
\texttt{qq-ring-def} (also a 6-tuple).  \texttt{rr-normed-field} and \texttt{cc-normed-field}
are 7-tuples in the \texttt{normed-field} shape (norm at slot~7).

A structure predicate pins \texttt{length(s) = n}, so a tuple witnesses only
the predicates of its own shape.  The \texttt{is-ring} witnesses are thus
installed only for the genuine 6-tuples: \texttt{zz-is-ring} and
\texttt{qq-is-ring}.  The 7-tuples carry \texttt{rr-is-normed-field} /
\texttt{cc-is-normed-field} and reach the ring world through the
\texttt{normed-field-as-*} views, not an \texttt{is-ring} assertion (that
would assert \texttt{length}~$=6$ of a length-7 list).  \texttt{qq} as a field
is the 8-tuple \texttt{qq-field} (\texttt{qq-field-is-field}), since
\texttt{field} is an 8-slot shape.

\paragraph{Use via \texttt{specialize-structure}.}
Once the \texttt{is-ring} witness is in place, every generic ring theorem can
be transported to the instance:
\begin{VNBcode}
(specialize-structure 'ZZ-RING 'RING 'zz-is-ring)
;; Installs ring-add-comm-zz-ring, ring-mul-assoc-zz-ring, ...
\end{VNBcode}
This is what turns abstract ring theorems into usable facts about
$\mathtt{zz}$.  The same call works for \texttt{qq-ring} (the other genuine
6-tuple ring); for $\mathtt{rr}$/$\mathtt{cc}$ the analogue runs through the
\texttt{normed-field-as-*} views instead.

\subsection{$\mathtt{nn}$ as a commutative monoid under addition}

$\mathtt{nn}$ is not a ring (no negation), but it is a commutative
monoid under $+$ with identity $0$:
\[
   \mathit{nn\text{-}add\text{-}monoid} \;=\; [\,\mathtt{nn},\;
     \mathtt{binplus},\; 0\,]
\]
with \texttt{nn-add-monoid-def} and witness
\texttt{nn-add-monoid-is-comm-monoid}.  A multiplicative
comm-monoid for $\mathtt{nn}$ under $\cdot$ is left for later.

\begin{remark}
The $n$-ary arithmetic operators do not disappear.  Any proof can
freely mix the two presentations: use $(+\; a\; b\; c)$ for ground
calculations (closed by \texttt{(arith)}), and route through the
binary form whenever the goal is a generic algebraic claim about a
ring or monoid.  The apply axioms (\texttt{binplus-apply}, etc.)
mediate translation in either direction.
\end{remark}

%% =========================================================
\section{$\mathtt{cc}$ as a metric space: \texttt{cc-ms}}
\label{sec:cc-ms}

The \texttt{metric-space} structure (Section~\ref{sec:declare-structure})
is one of the basic algebraic structures.  Its sole instance shipped
with the prover is \texttt{cc-ms}, the complex plane with the standard
metric:
\[
   \mathit{cc\text{-}ms} \;=\; \bigl[\,\mathtt{cc},\;
     \lambda([x, y],\; \mathtt{magnitude}(x - y))\,\bigr]
\]
i.e.\ the 2-element VNB list whose carrier is $\mathtt{cc}$ and whose
distance function is $d(x, y) = |x - y|$ (the user's
``$\mathtt{abs}(y - x)$'' --- the standard Euclidean metric on the
complex plane, expressed via the \texttt{magnitude} primitive).

\paragraph{Defining axiom.}
\begin{description}
\item[\texttt{cc-ms-def}]
  $\mathit{cc\text{-}ms} = [\mathtt{cc},\; \lambda([x, y],\;
   \mathtt{magnitude}(x - y))]$
\item[\texttt{cc-is-metric-space}]
  $\mathit{is\text{-}metric\text{-}space}(\mathit{cc\text{-}ms})$.
  A theorem (\texttt{theorem-library/cc-metric-space-proof.scm}): the
  proof unfolds the defining equivalence of the structure predicate and
  discharges each metric-space axiom (positivity, self-zero,
  separation, symmetry, triangle inequality) from the corresponding
  property of \texttt{magnitude}.
\end{description}

\paragraph{Use via \texttt{specialize-structure}.}
Once \texttt{cc-is-metric-space} is in place, every generic
metric-space theorem can be transported to \texttt{cc-ms}:
\begin{VNBcode}
(specialize-structure 'CC-MS 'METRIC-SPACE 'cc-is-metric-space)
;; Installs metric-pos-cc-ms, metric-triangle-cc-ms, ...
\end{VNBcode}

\paragraph{Completeness.}
The predicates \texttt{is-cauchy-seq}, \texttt{converges-to} and
\texttt{is-complete} are defined at the \texttt{metric-space} level
(\texttt{structure-library/metric-completeness.scm}).  Completeness of
$\mathtt{cc}$ is the theorem \texttt{cc-complete},
$\mathit{is\text{-}complete}(\mathit{cc\text{-}ms})$
(\texttt{theorem-library/cc-complete-proof.scm}), proved from the
completeness of $\mathtt{rr}$ through the real and imaginary parts.

%% =========================================================
\section{Extended reals: $\mathtt{rr}^{*}$}
\label{sec:extended-reals}

The extended real line
\[
   \mathtt{rr}^{*} \;=\; \mathtt{rr} \;\cup\; \{-\infty,\, +\infty\}
\]
adds two new constants $\mathtt{pos\text{-}inf}$ and
$\mathtt{neg\text{-}inf}$ at the top and bottom of the real line.  The
file \texttt{extended-reals.scm} (loaded after \texttt{numeric-instances.scm})
installs the carrier class and the ordering.

\begin{remark}
The class is spelled \texttt{rr-star}.  It was written \texttt{RR*}
until 2026-08-24, and that spelling could not be tokenized: the
\texttt{*} lexes as an operator, so \texttt{x in rr*} parsed as a
product rather than as a membership.  \texttt{rr-star} is the spelling
the axioms below already use, and it parses from a string like any
other name.
\end{remark}

\subsection{Membership}

\begin{description}
\item[\texttt{rr-star-membership}]
  $\forall x.\; x \in \mathtt{rr}^{*} \;\Leftrightarrow\;
   x \in \mathtt{rr} \lor x = \mathtt{pos\text{-}inf}
                     \lor x = \mathtt{neg\text{-}inf}.$
\item[\texttt{rr-subset-rr-star}]
  $\mathtt{rr} \subseteq \mathtt{rr}^{*}$ (immediate from the iff).
\item[\texttt{pos-inf-in-rr-star}]
  $\mathtt{pos\text{-}inf} \in \mathtt{rr}^{*}$.
\item[\texttt{neg-inf-in-rr-star}]
  $\mathtt{neg\text{-}inf} \in \mathtt{rr}^{*}$.
\end{description}

\subsection{Distinctness}

The two infinities are distinct from each other and from every real
number.  Without these, the membership iff could be satisfied
vacuously --- e.g.\ \texttt{pos-inf} could coincide with $0$.
\begin{description}
\item[\texttt{pos-inf-neq-neg-inf}]
  $\mathtt{pos\text{-}inf} \neq \mathtt{neg\text{-}inf}$.
\item[\texttt{pos-inf-not-in-rr}]
  $\mathtt{pos\text{-}inf} \notin \mathtt{rr}$.
\item[\texttt{neg-inf-not-in-rr}]
  $\mathtt{neg\text{-}inf} \notin \mathtt{rr}$.
\end{description}

\subsection{Ordering}

The existing $\le$ on $\mathtt{rr}$ is extended to $\mathtt{rr}^{*}$ by
two bound axioms.  Reflexivity at the infinities ($\mathtt{pos\text{-}inf}
\le \mathtt{pos\text{-}inf}$, etc.) follows by substituting
$x := \mathtt{pos\text{-}inf}$ (or $\mathtt{neg\text{-}inf}$) in their
respective bounds.
\begin{description}
\item[\texttt{pos-inf-upper-bound}]
  $\forall x \in \mathtt{rr}^{*}.\; x \le \mathtt{pos\text{-}inf}.$
\item[\texttt{neg-inf-lower-bound}]
  $\forall x \in \mathtt{rr}^{*}.\; \mathtt{neg\text{-}inf} \le x.$
\end{description}

\begin{remark}
Arithmetic on the extended reals is intentionally \emph{not} provided.
The standard pitfalls --- $\mathtt{pos\text{-}inf} + \mathtt{neg\text{-}inf}$,
$\mathtt{pos\text{-}inf} - \mathtt{pos\text{-}inf}$,
$0 \cdot \mathtt{pos\text{-}inf}$, \ldots --- have no canonical
resolution, and the convergence machinery that motivates
$\mathtt{rr}^{*}$ does not need them.  $\mathtt{rr}^{*}$ here is purely
an \emph{ordered carrier}: a place to put limits and bounds that may
fail to be finite.
\end{remark}

%% =========================================================
\section{Matrices over a ring}
\label{sec:matrices}

Matrices are built in three layers, and the middle one is the point of
the design.  A matrix is first a \emph{set-theoretic} object: a list of
rows, of a pinned size.  Nothing is proved at that level.  All matrix
algebra then goes through a \emph{tabulation bridge}---construct a
matrix by giving its entry function, prove two matrices equal by
comparing entries---so that no proof performs list surgery.  The ring
operations are defined on top of the bridge.

\subsection{What a matrix is: \texttt{mat} and \texttt{entry}}

\paragraph{Signature.}
$\mathtt{mat}(m, n, X)$ is the class of $m \times n$ matrices with
entries in $X$; $\mathtt{entry}(M, i, j)$ is the $(i,j)$ entry.

\begin{align*}
\mathtt{size}(P) &\;=\; [\,\mathtt{length}(P),\;
   \mathbf{if}\; \mathtt{length}(P) = 0 \;\mathbf{then}\; 0
   \;\mathbf{else}\; \mathtt{length}(\mathtt{nth}(1, P))\,] \\
\mathtt{mat}(m, n, X) &\;=\; \mathtt{sep}\, P \in \mathtt{matrix}(X).\;
   n \in \mathtt{nn} \;\wedge\;
   (\mathtt{length}(P) = 0 \Rightarrow m = 0) \\
 &\qquad\qquad \wedge\;
   (\mathtt{length}(P) \neq 0 \Rightarrow \mathtt{size}(P) = [m, n]) \\
\mathtt{entry}(M, i, j) &\;=\; \mathtt{nth}(j,\, \mathtt{nth}(i, M))
\end{align*}

\noindent
$\mathtt{size}$ is total.  A matrix with no rows is $[\,]$, and it is
the unique member of $\mathtt{mat}(0, n, X)$ for every natural $n$: a
matrix with no rows does not determine its column count, and membership
in $\mathtt{mat}$ does not pretend that it does.  Both clauses are
decisions of 2026-09-16 (\texttt{docs/size-mat-surgery-2026-09-16.md}).
Until then $\mathtt{mat}(m,n,X)$ was $\mathtt{sep}\, P \in
\mathtt{matrix}(X).\; \mathtt{size}(P) = [m,n]$ with $\mathtt{size}$
reading the column count off the first row, so no zero-row matrix had a
size and no $\mathtt{mat}(0,n,X)$ had a member.  Proofs read a
$\mathtt{mat}$ membership only through the lemmas of
\texttt{theorem-library/mat-basics.scm} (\texttt{mat-in-matrix},
\texttt{mat-length}, \texttt{mat-cols}, \texttt{mat-rows-in-nn},
\texttt{mat-intro}, \texttt{nil-in-mat}, \ldots), never by unfolding the
definition.

Indices are \emph{1-based}, matching the library's $\mathtt{nth}$
convention: row $i$ is $\mathtt{nth}(i, M)$, and the index range of a
row or column is $\mathtt{interval}(1, n)$.  Because
$\mathtt{mat}(m,n,X)$ is a $\mathtt{sep}$ over $\mathtt{matrix}(X)$ it
inherits sethood and the $\mathtt{sep}$ membership characterisation for
free (\texttt{mat-is-set}, \texttt{matrix-membership},
\texttt{matrix-sethood}), and $\mathtt{mat}(m,n,X)$ really is a subclass
of $\mathtt{matrix}(X)$.

\subsection{The tabulation bridge: \texttt{matof}}

$\mathtt{matof}(m, n, g)$ is the $m \times n$ matrix whose $(i,j)$ entry
is $g(i,j)$, defined as the $\mathtt{iota}$ of the obvious description.
Two facts make it usable, and they are of different kinds:

\begin{description}
\item[\texttt{matrix-entry-extensionality}]
  two matrices of the same size with equal entries throughout are equal.
  This is \emph{uniqueness} for the description.

\item[\texttt{matof-exists}]
  for natural $m$ and $n$, and $g$ defined throughout the index box, the
  tabulation of $g$ over the box exists.  This is \emph{existence} for
  the description, and it is a theorem
  (\texttt{theorem-library/tuple-tabulation.scm}): the box is finite, so
  the tabulation is built row by row by list recursion.
\end{description}

\noindent
Together the two pin $\mathtt{matof}$ down, and \texttt{entry-of-matof}
and \texttt{matof-in-mat} are derived from them by \texttt{iota-def}
(\texttt{tuple-tabulation.scm}, \texttt{matof-in-mat.scm}).  All three
are theorems guarded on $m, n \in \mathtt{nn}$, and a \texttt{qed} that
cites them pays nothing.

The guard is not decoration.  Until 2026-09-16 \texttt{matof-exists}
was an asserted support with $m$ and $n$ unquantified, and the library
proved \texttt{falsity} from it: instantiated at $m := \mathtt{ord}$ it
gave a matrix with $\mathtt{ord}$ rows, \texttt{mat-rows-in-nn} read
$\mathtt{ord} \in \mathtt{nn}$ off that matrix, and Burali-Forti refuted
it (the derivation is \texttt{scratchpad/mx/mx-probe3.scm}, kept as a
control that must \emph{fail}).  The assertion had been harmless for
weeks; the lemma that met it was proved much later for an unrelated
reason.  That is the case for the bill ledger of
Section~\ref{sec:proof-debt}: an asserted step is a liability whose
size is not known when it is taken.

Every construction below is a $\mathtt{matof}$, so each carries its
typing lemma and its entry lemma, and proofs about it reduce to
arithmetic on entries.

\subsection{Arithmetic}

Writing $P_{ij}$ for $\mathtt{entry}(P,i,j)$---an abbreviation of the
argument, the operators stay prefix:

\begin{align*}
\mathtt{matadd}(A, P, Q) &\;=\;
  \mathtt{matof}(m, n,\; \lambda i\, j.\;
     \mathtt{add}(A)(P_{ij},\, Q_{ij})) \\
\mathtt{matmul}(A, P, Q) &\;=\;
  \mathtt{matof}(m, k,\; \lambda i\, k.\; \mu(i,k)),
  \quad\text{where} \\
 &\qquad \mu(i,k) \;=\;
     \textstyle\mathtt{finsum}_{j \in \mathtt{interval}(1,n)}\;
     \mathtt{mul}(A)(P_{ij},\, Q_{jk}) \\
\mathtt{zeromat}(A, m, n) &\;=\;
  \mathtt{matof}(m, n,\; \lambda i\, j.\; \mathtt{zero}(A)) \\
\mathtt{identmat}(A, n) &\;=\;
  \mathtt{matof}(n, n,\; \lambda i\, j.\; \delta(i,j)),
  \quad\text{where} \\
 &\qquad \delta(i,j) \;=\;
     \mathbf{if}\; i = j \;\mathbf{then}\; \mathtt{one}(A)
     \;\mathbf{else}\; \mathtt{zero}(A)
\end{align*}

\noindent
with $\mathtt{matneg}(A,P)$ and $\mathtt{matscale}(A,r,P)$ the same
shape.  The dimensions are read off $\mathtt{size}(P)$ and
$\mathtt{size}(Q)$ rather than passed: multiplication is
\emph{total}, garbage-in/garbage-out off its typed domain, and all the
content lives in \texttt{matmul-type} and \texttt{matmul-entry}, both
theorems.

Reading the dimensions off the factors has one consequence that the
statements must carry.  The column count of $\mathtt{matmul}(A,P,Q)$ is
read off $\mathtt{size}(Q)$; when the middle dimension $n$ is $0$, $Q$
is $[\,]$, $\mathtt{size}([\,]) = [0,0]$, and the product of an
$m \times 0$ matrix with it is $m \times 0$, not $m \times k$.  So every
law about a product whose middle dimension may be $0$ while both outer
dimensions are at least $1$ is false there, and each carries a guard
stating when it holds: an entry law, whose user holds an index in the
interval, carries $1 \le n$; a typing law carries
$n = 0 \Rightarrow (m = 0 \vee k = 0)$; associativity, with two inner
dimensions, carries a pair of such conditions (G1, G2 in the design
note).  The guards are premises placed after the binders, so a citation
of the law reads exactly as it did before the guard.

The whole-matrix laws (addition: associativity, commutativity, the zero and
negation laws; multiplication: associativity, the guarded distributive laws
\texttt{matmul-left-dist-guarded} and \texttt{matmul-right-dist-guarded}; the
identity laws \texttt{identmat-left-identity} and \texttt{identmat-right-identity})
are theorems of the library, proved in \texttt{theorem-library/}
(\texttt{rake-matrix-laws.scm}, \texttt{rake-identmat.scm},
\texttt{matmul-assoc-proof.scm}); \texttt{structure-library/matrix.scm} no
longer installs any support.  Whether a particular law is proved or still
rests on an asserted fact is reported, with its bill, by
\texttt{reference/PROOF-DEBT.md} and the browser reference; the generated
pages are the authority.

\subsection{The matrix ring: \texttt{mat-ring}}

$\mathit{mat\text{-}ring}(A, n)$ packages the $n \times n$ matrices as a
\textsc{Ring} 6-tuple, currying each operation into a map on the carrier
so it fits the ring shape:
\begin{align*}
  [\;&\mathtt{mat}(n,n,\mathtt{carr}(A)),\;
     \lambda P\,Q.\,\mathtt{matadd}(A,P,Q),\;
     \lambda P\,Q.\,\mathtt{matmul}(A,P,Q),\\
    &\lambda P.\,\mathtt{matneg}(A,P),\;
     \mathtt{zeromat}(A,n,n),\;
     \mathtt{identmat}(A,n)\;]
\end{align*}
\texttt{mat-ring-is-ring} proves it satisfies $\mathit{is\text{-}ring}$,
and the read-offs \texttt{mat-ring-carr}, \texttt{-add}, \texttt{-mul},
\texttt{-zero}, \texttt{-one}, \texttt{-neg} let a proof move between
the packaged ring and the bare operations.

\subsection{Cutting and extending}

\begin{align*}
\mathtt{submat}(S, p, q) &\;=\;
  \mathtt{matof}(p, q,\; \lambda i\, j.\;
    \mathtt{entry}(S, \mathtt{succ}(i), \mathtt{succ}(j)))
\end{align*}

\noindent
deletes the first row and the first column (the shape the Smith-normal-form
recursion wants), while $\mathtt{block}$, $\mathit{snoc\text{-}col}$ and
$\mathit{snoc\text{-}row}$ extend a matrix by a column or a row.  Each has a
typing lemma and an entry lemma (\texttt{submat-type},
\texttt{entry-of-submat}, \texttt{snoc-col-type}, \texttt{snoc-col-last},
\dots).  Deleting an \emph{arbitrary} row and column is a different
operation, $\mathtt{minor}$, defined with the determinant in
Section~\ref{sec:determinant}.

\subsection{Elementary matrices}

$\mathtt{matunit}(A, n, k, l)$ is the matrix unit---$\mathtt{one}$ at
$(k,l)$, $\mathtt{zero}$ elsewhere---and its column-shift lemma
\texttt{matunit-col-shift} is the engine of the elementary-operation
arc.  On top of it sit the three elementary matrices
$\mathit{elem\text{-}f}(A,n,k,l)$ (exchange), $\mathit{elem\text{-}g}(A,n,r,k,l)$
(add a multiple), and $\mathit{elem\text{-}h}(A,n,r,k)$ (scale), each with
its column action (\texttt{elem-*-action}, Prop.~3.5 of
\texttt{algebraic-numbers.pdf}), its row action, and its inverse
(\texttt{elem-*-inverse}, Cor.~3.6).  $\mathit{mat\text{-}equiv}$ is the
equivalence relation they generate; it is reflexive and transitive
(\texttt{mat-equiv-refl}, \texttt{mat-equiv-trans}) and stable under
multiplication on either side.

\begin{remark}
The browser reference groups this whole layer under \textsc{Linear
algebra} in \texttt{reference/BY-TOPIC.md}, in the same reading order,
each entry linking to the file that proves it.
\end{remark}

%% =========================================================
\section{Determinants}
\label{sec:determinant}

The determinant is defined by \emph{recursive cofactor expansion along
the first row}---deliberately the pedestrian route.  There is no
exterior algebra here, no symmetric group and no sign of a permutation:
the Leibniz sum would have required $\mathtt{sgn}$ and a well-definedness
proof that the library does not have, whereas cofactor expansion needs
only the recursion on size, minors, and $\mathtt{finsum}$.

\subsection{\texttt{minor}: deleting a row and a column}

\begin{align*}
  \mathtt{minor}(S, r, c, n) &\;=\;
  \mathtt{matof}(n, n,\; \lambda i\, j.\;
    \mathtt{entry}(S,\, \mathit{skip}(i,r),\, \mathit{skip}(j,c)))
  \quad\text{where}\\
  \mathit{skip}(i, r) &\;=\;
    \mathbf{if}\; i < r \;\mathbf{then}\; i \;\mathbf{else}\; \mathtt{succ}(i)
\end{align*}

\noindent
($\mathit{skip}$ is not an operator of the library; the source writes the
conditional inline with \texttt{if}.)  This is
the $n \times n$ matrix obtained from a
$\mathtt{succ}(n) \times \mathtt{succ}(n)$ matrix $S$ by deleting row $r$
and column $c$.  \texttt{minor-type} states that it lands in
$\mathtt{mat}(n, n, \mathtt{carr}(R))$.

\subsection{\texttt{det}: the defining recursion}

Two axioms, installed as \emph{definitional}
(\texttt{add-definition!} under provenance \texttt{definitional}),
so they carry no warrant and owe no proof debt:

\begin{align*}
\mathtt{det\text{-}zero}: &\quad \mathtt{det}(R, 0, A) \simeq \mathtt{one}(R) \\
\mathtt{det\text{-}cofactor}: &\quad
  \mathtt{det}(R, \mathtt{succ}(n), A) \;\simeq\; {} \\
 &\qquad \textstyle\mathtt{finsum}_{j \in \mathtt{interval}(1,\,\mathtt{succ}(n))}\;
   \sigma_j \\
 &\qquad\text{where }\;
   \sigma_j = (-1)^{1+j}_{R} \cdot \mathtt{entry}(A, 1, j) \cdot
   \mathtt{det}(R, n, M_j) \\
 &\qquad\text{and }\; M_j = \mathtt{minor}(A, 1, j, n)
\end{align*}

\noindent
where the sum is $\mathtt{finsum}$ over $\mathit{ring\text{-}additive\text{-}ag}(R)$
and the sign $(-1)^{1+j}$ is taken \emph{in the ring}, as
$\mathtt{mpow}$ of $R$'s multiplicative monoid applied to
$\mathtt{neg}(R)(\mathtt{one}(R))$.  The $0 \times 0$ case is the empty
product, $\mathtt{one}(R)$, and the $1 \times 1$ case falls out of the
recursion: $\mathtt{det}(R,1,A) = \mathtt{entry}(A,1,1)$.

The axioms are hand-stated rather than emitted by
\texttt{def-by-nn-recursion} because the matrix \emph{shrinks} in the
recursive call---to its own minor---so the size is not a fixed recursion
parameter.  Their consistency rests on $\mathtt{det}$ being well defined
by this recursion, which is the projection of the legitimate
size-indexed recursion $\mathtt{mat}(n,n) \to \mathtt{carr}$.

\subsection{Properties}

\begin{remark}
The statements in this subsection are theorems of the library
(\texttt{det-in-carrier} in \texttt{rake-mat-typing.scm}; \texttt{det-1x1},
\texttt{det-2x2}, \texttt{det-identity} in \texttt{rake-det-small.scm};
\texttt{det-alternating-rows} and \texttt{det-expand-row} in \texttt{det-rows.scm};
\texttt{det-multiplicative} in \texttt{det-multiplicative.scm};
\texttt{det-transpose} and \texttt{det-expand-col} in \texttt{det-transpose.scm}).
The bill of a proof that cites one is reported by \texttt{reference/PROOF-DEBT.md}.
\end{remark}

\begin{description}
\item[\texttt{det-in-carrier}]
  for $A \in \mathtt{mat}(n,n,\mathtt{carr}(r))$ over a ring $r$,
  $\mathtt{det}(r,n,A) \in \mathtt{carr}(r)$.

\item[\texttt{det-1x1}]
  $\mathtt{det}(r,1,A) = \mathtt{entry}(A,1,1)$.

\item[\texttt{det-2x2}]
  $\mathtt{det}(r,2,A) = A_{11}A_{22} - A_{12}A_{21}$.

\item[\texttt{det-identity}]
  $\mathtt{det}(r, n, \mathtt{one}(\mathit{mat\text{-}ring}(r,n)))
   = \mathtt{one}(r)$.

\item[\texttt{det-alternating-rows}]
  two equal rows force $\mathtt{det}(r,n,A) = \mathtt{zero}(r)$.

\item[\texttt{det-multiplicative}]
  over a \emph{commutative} ring,
  \[
    \mathtt{det}(r,n,\mathtt{matmul}(r,P,Q)) \;=\;
    \mathtt{mul}(r)(\mathtt{det}(r,n,P),\, \mathtt{det}(r,n,Q)).
  \]
\end{description}

\begin{remark}
The following facts are deferred: Cramer's rule, the adjugate, and
$\det(A) \ne 0 \Leftrightarrow A$ invertible.  Transposition
(\texttt{structure-library/transpose.scm}), $\det(A^{T}) = \det(A)$ and
expansion along an arbitrary row or column are proved.
\end{remark}

\paragraph{Computing with the recursion.}
The axiom is keyed on the pattern
$\mathtt{det}(R, \mathtt{succ}(n), A)$, while a concrete goal names a
literal size, $\mathtt{det}(r, 1, A)$.  The rewriter bridges the two in
both directions.  Matching is handled by \texttt{match-expr}, which
matches $\mathtt{succ}(n)$ against the numeral $1$ with $n := 0$.  The
macete's conditions are handled by \texttt{condition-holds?}, which
compares them to the context with numerals and $\mathtt{succ}$-towers
collapsed to one form, and discharges a ground arithmetic side condition
such as $0 \in \mathtt{nn}$ outright.  So the plain \texttt{mac} fires
on a literal size:

\begin{VNBcode}
(sp (make-wff '(FORALL (r a)
     (IMPLIES (IS-RING r)
       (IMPLIES (IN a (MAT 1 1 (CARR r)))
         (= (DET r 1 a) (ENTRY a 1 1)))))))
(di)(di)(di)
(mac 'det-cofactor)   ; goal becomes the FINSUM over INTERVAL(1, succ 0)
\end{VNBcode}


%% =========================================================
\chapter{Structures as categories}
\label{ch:categories}

A structure declared with \texttt{declare-structure}
(Chapter~\ref{ch:defs}) comes with its morphisms: the maps between two
of its models that preserve the structure.  The models and their
morphisms form a category.  This chapter describes how VNB sets those
categories up, how a species can carry more than one, and how one
structure is viewed as another.  Such a view is a functor.

%%% -----------------------------------------------------------------------
\section{Morphisms of structures}
\label{sec:morphisms}

Declaring a structure declares its morphisms too.  From the slot list alone,
\texttt{declare-structure} generates a homomorphism predicate
\texttt{is-hom-\textit{name}}$(a, b, f_1 \ldots f_k)$ --- one map $f_i$ per
\emph{independent} carrier --- whose meaning is the obvious one: each map is
typed between the corresponding carriers, and every operation and constant is
preserved.

\begin{itemize}
\item \textbf{carrier} $c$: a map \texttt{$f$ in fun($c(a)$, $c(b)$)}.
\item \textbf{operation} \texttt{OP}: preservation ---
  \texttt{$f$(OP($a$)($x_1,\ldots$)) = OP($b$)($f(x_1),\ldots$)} for all
  arguments in their sorts (each argument mapped by its own carrier's map).
\item \textbf{constant} \texttt{C}: \texttt{$f$(C($a$)) = C($b$)}.
\item \textbf{substructure} \texttt{K}: agreement, \texttt{K($a$) = K($b$)} ---
  morphisms are taken between structures over the \emph{same} base.
\item \textbf{derived} carrier: nothing.  Its value is pinned by
  \texttt{is-\textit{name}} (Section~\ref{sec:declare-structure}), so it rides
  the base carrier's map and owes no map of its own.  This is why a field
  homomorphism is one map, not two.
\end{itemize}

\noindent A group's generated hom, for instance, is the single map \texttt{$f$
in fun(carr($a$), carr($b$))} with \texttt{$f$(opr($a$)($x$,$y$)) =
opr($b$)($f(x)$,$f(y)$)} and \texttt{$f$(iden($a$)) = iden($b$)} --- exactly a
group homomorphism.  The definition is installed with \emph{definitional}
provenance, so it costs no warrant.

\subsection{Overriding the default: \texttt{declare-hom!}}
\label{sec:declare-hom}

The generated clause for each slot kind is the standard one: preservation for
an operation or a constant, the pullback condition for a \texttt{family} slot
(the arrows of \texttt{top-space} are the continuous maps, those of
\texttt{measurable-space} the measurable maps), transport for a
\texttt{family-fun} slot, ``related points go to related points'' for a
\texttt{relation} slot.  For a \texttt{metric-space} the generated arrows
preserve the distance: they are the isometries, and a \texttt{notation!}
beside the structure says so, since ``homomorphism'' is then the wrong word
(Section~\ref{sec:notation}).  Where no generated clause is the right one, a
species declares its own morphism notion with \texttt{declare-hom!}, beside the
structure like \texttt{notation!}:

\begin{VNBcode}
(declare-hom! 'METRIZABLE-TOP-SPACE '(s t f)
  '(FORALL u (IMPLIES (IN u (OPENS t))
      (IN (PREIMAGE s f u) (OPENS s)))))
\end{VNBcode}

The argument list is the two objects and one map per carrier.  The body states
\emph{only} what is characteristic: the object typings
\texttt{is-metrizable-top-space($s$)}, \texttt{is-metrizable-top-space($t$)}
and the map typing \texttt{$f$ in fun(pts($s$), pts($t$))} are supplied
automatically and need not be repeated.  The declaration replaces the generated
\texttt{is-hom-metrizable-top-space} entirely.  Two species override today:
\texttt{metrizable-top-space} (a refinement inherits its parent's shape, not
its arrows, so continuity is restated) and \texttt{ringoid} (a ring
homomorphism of the underlying rings that preserves the ideal).  The table of
Section~\ref{sec:declare-category} says, for every structure, which of the
three cases its arrows are.

\subsection{Several categories on one species: \texttt{declare-category!}}
\label{sec:declare-category}

A structure is a category: its objects are its models, its arrows the maps
satisfying \texttt{is-hom-\textit{name}}, composed as functions.  That is the
species' \emph{default} category, and for an algebraic species it is the right
one.  For a metric space the default is the category of isometries, which is
not the one analysis works in: the Lipschitz maps, the uniformly continuous
maps and the continuous maps each form a category on the same objects; on the
countably-metrised spaces (\texttt{c-metric-space}) the continuous maps for
the canonical metric form one, which rests on the canonical metric being a
metric (\texttt{c-metric-is-metric-space}).  A species may therefore carry
several categories.  A second one is declared
beside the structure:

\begin{VNBcode}
(declare-category! 'LIPSCHITZ 'METRIC-SPACE '(a b f)
  '(FORSOME lpk_ (AND (IN lpk_ RR) (AND (<= 0 lpk_)
     (FORALL lpx_ (IMPLIES (IN lpx_ (PTS a))
       (FORALL lpy_ (IMPLIES (IN lpy_ (PTS a))
         (<= ((DIST b) (f lpx_) (f lpy_))
             (* lpk_ ((DIST a) lpx_ lpy_)))))))))))
\end{VNBcode}

The body states what is characteristic of an arrow $f \colon a \to b$; the
object typings and the carrier typing of each map are supplied, as for
\texttt{declare-hom!}.  The declaration installs, definitionally, the arrow
predicate \texttt{is-lipschitz-arrow}$(a, b, f)$ and the hom-set
\texttt{hom-lipschitz}$(a, b)$, and registers three \emph{obligations} as
statements, not as theorems: \texttt{hom-lipschitz-id} (the identity is an
arrow), \texttt{hom-lipschitz-compose} (arrows compose) and
\texttt{hom-lipschitz-in-set} (the hom-set is a set).  That the arrows of a
declared category contain the identities and compose is where the
mathematics is (Lipschitz constants multiply; an $\varepsilon$--$\delta$
argument for the continuous maps), so the generator does not stamp it: the
obligations are proven in \texttt{theorem-library/categories.scm}, and
\texttt{category-obligations-audit}, run at the end of every load, lists
every obligation that is not a proven theorem of exactly its statement.
Until that list is empty the declared category is not yet a category.  The
default category is untouched by the declaration.

The categories the library knows are listed below.  The listing is generated
at every load (\texttt{reference/CATEGORIES.md}, also a page of the browser
reference): the default category of each structure with the status of its
three laws, the declared categories with their defining equivalences, and the
inclusions between categories on the same objects --- an isometry is
$1$-Lipschitz, a Lipschitz map is uniformly continuous, a uniformly
continuous map is continuous, a linear isometry is bounded linear --- which
the page finds by the shape of the installed theorems.

%% AUTO-GENERATED by docs/gen-categories.py from reference/CATEGORIES.md -- do not edit.
%% Regenerate:  python3 docs/gen-categories.py   (after a prover load refreshes CATEGORIES.md)

% the long identifiers break only with loose spacing; \sloppy, scoped to this file
\begingroup\sloppy
\subsubsection{The default categories}
One per structure with a morphism predicate.  The laws column reads ``proven'' when the identity law, the composition law and the sethood of the hom-set are all proven theorems; otherwise it names the law and its status.

\begin{longtable}{@{}>{\raggedright\arraybackslash}p{0.22\linewidth}>{\raggedright\arraybackslash}p{0.30\linewidth}>{\raggedright\arraybackslash}p{0.20\linewidth}>{\raggedright\arraybackslash}p{0.18\linewidth}@{}}
\toprule
structure & arrows & defined & laws \\
\midrule
\endhead
\texttt{abelian-\allowbreak{}group} & f is a homomorphism from a to b & generated & proven \\
\texttt{bundle} & f is a homomorphism from a to b & generated & id: not installed; compose: not installed; in-set: not installed \\
\texttt{c-\allowbreak{}metric-\allowbreak{}space} & f is a homomorphism from a to b & generated & proven \\
\texttt{comm-\allowbreak{}monoid} & f is a homomorphism from a to b & generated & proven \\
\texttt{commutative-\allowbreak{}banach-\allowbreak{}algebra} & f is a homomorphism from a to b & refinement of \texttt{unital-\allowbreak{}banach-\allowbreak{}algebra} & id: not installed; compose: not installed; in-set: not installed \\
\texttt{commutative-\allowbreak{}c-\allowbreak{}star-\allowbreak{}algebra} & f is a homomorphism from a to b & refinement of \texttt{unital-\allowbreak{}c-\allowbreak{}star-\allowbreak{}algebra} & id: not installed; compose: not installed; in-set: not installed \\
\texttt{commutative-\allowbreak{}ring} & f is a homomorphism from a to b & refinement of \texttt{ring} & proven \\
\texttt{complex-\allowbreak{}banach-\allowbreak{}algebra} & f is a homomorphism from a to b & refinement of \texttt{complex-\allowbreak{}normed-\allowbreak{}algebra} & id: not installed; compose: not installed; in-set: not installed \\
\texttt{complex-\allowbreak{}banach-\allowbreak{}space} & f is a homomorphism from a to b & refinement of \texttt{complex-\allowbreak{}normed-\allowbreak{}vector-\allowbreak{}space} & id: not installed; compose: not installed; in-set: not installed \\
\texttt{complex-\allowbreak{}inner-\allowbreak{}product-\allowbreak{}space} & f is a homomorphism from a to b & generated & proven \\
\texttt{complex-\allowbreak{}normed-\allowbreak{}algebra} & f is a homomorphism from a to b & generated & id: not installed; compose: not installed; in-set: not installed \\
\texttt{complex-\allowbreak{}normed-\allowbreak{}vector-\allowbreak{}space} & f is a homomorphism from a to b & generated & id: not installed; compose: not installed; in-set: not installed \\
\texttt{complex-\allowbreak{}pre-\allowbreak{}hilbert-\allowbreak{}space} & f is a homomorphism from a to b & refinement of \texttt{complex-\allowbreak{}semi-\allowbreak{}inner-\allowbreak{}product-\allowbreak{}space} & id: not installed; compose: not installed; in-set: not installed \\
\texttt{complex-\allowbreak{}semi-\allowbreak{}inner-\allowbreak{}product-\allowbreak{}space} & f is a homomorphism from a to b & generated & id: not installed; compose: not installed; in-set: not installed \\
\texttt{directed-\allowbreak{}set} & f is a homomorphism from a to b & generated & id: not installed; compose: not installed; in-set: not installed \\
\texttt{euclidean-\allowbreak{}ring} & f is a homomorphism from a to b & refinement of \texttt{integral-\allowbreak{}domain} & proven \\
\texttt{field} & f is a homomorphism from a to b & generated & proven \\
\texttt{field-\allowbreak{}ring} & f is a homomorphism from a to b & refinement of \texttt{commutative-\allowbreak{}ring} & proven \\
\texttt{group} & f is a homomorphism from a to b & generated & proven \\
\texttt{infiltrated-\allowbreak{}space} & f1 is a morphism of infiltrated strand spaces from a to b & generated & id: not installed; compose: not installed; in-set: not installed \\
\texttt{integral-\allowbreak{}domain} & f is a homomorphism from a to b & refinement of \texttt{commutative-\allowbreak{}ring} & proven \\
\texttt{measurable-\allowbreak{}space} & f is a measurable map from a to b & generated & proven \\
\texttt{measure-\allowbreak{}space} & f is a measure-preserving measurable map from a to b & generated & proven \\
\texttt{metric-\allowbreak{}group} & f is an isometric homomorphism from a to b & generated & id: not installed; compose: not installed; in-set: not installed \\
\texttt{metric-\allowbreak{}space} & f is an isometry from a to b & generated & proven \\
\texttt{metrizable-\allowbreak{}top-\allowbreak{}space} & f is continuous from a to b & by \texttt{declare-\allowbreak{}hom!} & proven \\
\texttt{module} & f is a homomorphism from a to b & generated & proven \\
\texttt{monoid} & f is a homomorphism from a to b & generated & proven \\
\texttt{normed-\allowbreak{}ag} & f is a homomorphism from a to b & generated & proven \\
\texttt{normed-\allowbreak{}field} & f is a homomorphism from a to b & generated & proven \\
\texttt{normed-\allowbreak{}vector-\allowbreak{}space} & f is a homomorphism from a to b & generated & proven \\
\texttt{nu-\allowbreak{}c-\allowbreak{}star-\allowbreak{}algebra} & f is a homomorphism from a to b & refinement of \texttt{nu-\allowbreak{}involutive-\allowbreak{}banach-\allowbreak{}algebra} & id: not installed; compose: not installed; in-set: not installed \\
\texttt{nu-\allowbreak{}involutive-\allowbreak{}banach-\allowbreak{}algebra} & f is a homomorphism from a to b & generated & id: not installed; compose: not installed; in-set: not installed \\
\texttt{pid} & f is a homomorphism from a to b & refinement of \texttt{integral-\allowbreak{}domain} & proven \\
\texttt{pk-\allowbreak{}space} & f is a morphism of public-key strand spaces from a to b & generated & id: not installed; compose: not installed; in-set: not installed \\
\texttt{pre-\allowbreak{}hilbert-\allowbreak{}space} & f is a homomorphism from a to b & refinement of \texttt{semi-\allowbreak{}inner-\allowbreak{}product-\allowbreak{}space} & id: not installed; compose: not installed; in-set: not installed \\
\texttt{pseudometric-\allowbreak{}space} & f is a homomorphism from a to b & generated & proven \\
\texttt{ring} & f is a homomorphism from a to b & generated & proven \\
\texttt{ringoid} & f is a ringoid homomorphism from a to b & by \texttt{declare-\allowbreak{}hom!} & proven \\
\texttt{semi-\allowbreak{}inner-\allowbreak{}product-\allowbreak{}space} & f is a homomorphism from a to b & generated & id: not installed; compose: not installed; in-set: not installed \\
\texttt{semigroup} & f is a homomorphism from a to b & generated & proven \\
\texttt{setoid} & f is a setoid morphism from a to b & generated & proven \\
\texttt{sk-\allowbreak{}space} & f is a morphism of shared-key strand spaces from a to b & generated & id: not installed; compose: not installed; in-set: not installed \\
\texttt{strand-\allowbreak{}space} & f1 is a morphism of strand spaces from a to b & generated & id: not installed; compose: not installed; in-set: not installed \\
\texttt{top-\allowbreak{}group} & f is a continuous homomorphism from a to b & generated & id: not installed; compose: not installed; in-set: not installed \\
\texttt{top-\allowbreak{}space} & f is continuous from a to b & generated & proven \\
\texttt{unital-\allowbreak{}banach-\allowbreak{}algebra} & f is a homomorphism from a to b & generated & id: not installed; compose: not installed; in-set: not installed \\
\texttt{unital-\allowbreak{}c-\allowbreak{}star-\allowbreak{}algebra} & f is a homomorphism from a to b & refinement of \texttt{unital-\allowbreak{}involutive-\allowbreak{}banach-\allowbreak{}algebra} & id: not installed; compose: not installed; in-set: not installed \\
\texttt{unital-\allowbreak{}involutive-\allowbreak{}banach-\allowbreak{}algebra} & f is a homomorphism from a to b & generated & id: not installed; compose: not installed; in-set: not installed \\
\texttt{vector-\allowbreak{}space} & f is a homomorphism from a to b & refinement of \texttt{module} & proven \\
\bottomrule
\end{longtable}

\subsubsection{The declared categories}

\paragraph{\texttt{bounded-\allowbreak{}linear,\allowbreak{} on normed-\allowbreak{}vector-\allowbreak{}space}}
\emph{Declared in} \texttt{theorem-\allowbreak{}library/\allowbreak{}categories.scm}.

Objects: the models of \texttt{is-\allowbreak{}normed-\allowbreak{}vector-\allowbreak{}space}.  Arrows: \texttt{is-\allowbreak{}bounded-\allowbreak{}linear-\allowbreak{}arrow(a,\allowbreak{} b,\allowbreak{} f)} --- f is a bounded linear map from a to b.  Hom-set: \texttt{hom-\allowbreak{}bounded-\allowbreak{}linear(a,\allowbreak{} b)} --- the bounded linear maps from a to b.

Definition (\texttt{is-\allowbreak{}bounded-\allowbreak{}linear-\allowbreak{}arrow-\allowbreak{}def}, definitional):
\begin{VNBsnip}
forall([a, b, f],
  is-bounded-linear-arrow(a, b, f)
  iff is-normed-vector-space(a)
      and is-normed-vector-space(b)
      and f in fun(vec(a), vec(b))
      and scal(a) = scal(b)
      and forall([x1_ in vec(a), x2_ in vec(a)],
            f((vadd(a))(x1_, x2_)) = (vadd(b))(f(x1_), f(x2_)))
      and forall([x1_ in carr(scal(a)), x2_ in vec(a)],
            f((act(a))(x1_, x2_)) = (act(b))(x1_, f(x2_)))
      and forsome([blk_ in rr],
            0 <= blk_
            and forall([x1_ in vec(a)],
                  (vnrm(b))(f(x1_)) <= blk_ * (vnrm(a))(x1_))))
\end{VNBsnip}

Laws: \texttt{hom-\allowbreak{}bounded-\allowbreak{}linear-\allowbreak{}id} --- proven; \texttt{hom-\allowbreak{}bounded-\allowbreak{}linear-\allowbreak{}compose} --- proven; \texttt{hom-\allowbreak{}bounded-\allowbreak{}linear-\allowbreak{}in-\allowbreak{}set} --- proven; membership \texttt{hom-\allowbreak{}bounded-\allowbreak{}linear-\allowbreak{}member-\allowbreak{}iff} --- proven; Hom functors \texttt{hom-\allowbreak{}bounded-\allowbreak{}linear-\allowbreak{}post-\allowbreak{}type} --- proven, \texttt{hom-\allowbreak{}bounded-\allowbreak{}linear-\allowbreak{}pre-\allowbreak{}type} --- proven.

\paragraph{\texttt{c-\allowbreak{}continuous,\allowbreak{} on c-\allowbreak{}metric-\allowbreak{}space}}
\emph{Declared in} \texttt{theorem-\allowbreak{}library/\allowbreak{}c-\allowbreak{}continuous-\allowbreak{}category.scm}.

Objects: the models of \texttt{is-\allowbreak{}c-\allowbreak{}metric-\allowbreak{}space}.  Arrows: \texttt{is-\allowbreak{}c-\allowbreak{}continuous-\allowbreak{}arrow(a,\allowbreak{} b,\allowbreak{} f)} --- f is a continuous map of countably-metrised spaces from a to b.  Hom-set: \texttt{hom-\allowbreak{}c-\allowbreak{}continuous(a,\allowbreak{} b)} --- the continuous maps of countably-metrised spaces from a to b.

Definition (\texttt{is-\allowbreak{}c-\allowbreak{}continuous-\allowbreak{}arrow-\allowbreak{}def}, definitional):
\begin{VNBsnip}
forall([a, b, f],
  is-c-continuous-arrow(a, b, f)
  iff is-c-metric-space(a)
      and is-c-metric-space(b)
      and f in fun(pts(a), pts(b))
      and is-continuous(c-metric(a), c-metric(b), f))
\end{VNBsnip}

Laws: \texttt{hom-\allowbreak{}c-\allowbreak{}continuous-\allowbreak{}id} --- proven; \texttt{hom-\allowbreak{}c-\allowbreak{}continuous-\allowbreak{}compose} --- proven; \texttt{hom-\allowbreak{}c-\allowbreak{}continuous-\allowbreak{}in-\allowbreak{}set} --- proven; membership \texttt{hom-\allowbreak{}c-\allowbreak{}continuous-\allowbreak{}member-\allowbreak{}iff} --- proven; Hom functors \texttt{hom-\allowbreak{}c-\allowbreak{}continuous-\allowbreak{}post-\allowbreak{}type} --- proven, \texttt{hom-\allowbreak{}c-\allowbreak{}continuous-\allowbreak{}pre-\allowbreak{}type} --- proven.

\paragraph{\texttt{continuous,\allowbreak{} on metric-\allowbreak{}space}}
\emph{Declared in} \texttt{theorem-\allowbreak{}library/\allowbreak{}categories.scm}.

Objects: the models of \texttt{is-\allowbreak{}metric-\allowbreak{}space}.  Arrows: \texttt{is-\allowbreak{}continuous-\allowbreak{}arrow(a,\allowbreak{} b,\allowbreak{} f)} --- f is a continuous map from a to b.  Hom-set: \texttt{hom-\allowbreak{}continuous(a,\allowbreak{} b)} --- the continuous maps from a to b.

Definition (\texttt{is-\allowbreak{}continuous-\allowbreak{}arrow-\allowbreak{}def}, definitional):
\begin{VNBsnip}
forall([a, b, f],
  is-continuous-arrow(a, b, f)
  iff is-metric-space(a)
      and is-metric-space(b)
      and f in fun(pts(a), pts(b))
      and is-continuous(a, b, f))
\end{VNBsnip}

Laws: \texttt{hom-\allowbreak{}continuous-\allowbreak{}id} --- proven; \texttt{hom-\allowbreak{}continuous-\allowbreak{}compose} --- proven; \texttt{hom-\allowbreak{}continuous-\allowbreak{}in-\allowbreak{}set} --- proven; membership \texttt{hom-\allowbreak{}continuous-\allowbreak{}member-\allowbreak{}iff} --- proven; Hom functors \texttt{hom-\allowbreak{}continuous-\allowbreak{}post-\allowbreak{}type} --- proven, \texttt{hom-\allowbreak{}continuous-\allowbreak{}pre-\allowbreak{}type} --- proven.

\paragraph{\texttt{lipschitz,\allowbreak{} on metric-\allowbreak{}space}}
\emph{Declared in} \texttt{theorem-\allowbreak{}library/\allowbreak{}categories.scm}.

Objects: the models of \texttt{is-\allowbreak{}metric-\allowbreak{}space}.  Arrows: \texttt{is-\allowbreak{}lipschitz-\allowbreak{}arrow(a,\allowbreak{} b,\allowbreak{} f)} --- f is a Lipschitz map from a to b.  Hom-set: \texttt{hom-\allowbreak{}lipschitz(a,\allowbreak{} b)} --- the Lipschitz maps from a to b.

Definition (\texttt{is-\allowbreak{}lipschitz-\allowbreak{}arrow-\allowbreak{}def}, definitional):
\begin{VNBsnip}
forall([a, b, f],
  is-lipschitz-arrow(a, b, f)
  iff is-metric-space(a)
      and is-metric-space(b)
      and f in fun(pts(a), pts(b))
      and forsome([lpk_ in rr],
            0 <= lpk_
            and forall([lpx_ in pts(a), lpy_ in pts(a)],
                  (dist(b))(f(lpx_), f(lpy_)) <= lpk_ * (dist(a))(lpx_, lpy_))))
\end{VNBsnip}

Laws: \texttt{hom-\allowbreak{}lipschitz-\allowbreak{}id} --- proven; \texttt{hom-\allowbreak{}lipschitz-\allowbreak{}compose} --- proven; \texttt{hom-\allowbreak{}lipschitz-\allowbreak{}in-\allowbreak{}set} --- proven; membership \texttt{hom-\allowbreak{}lipschitz-\allowbreak{}member-\allowbreak{}iff} --- proven; Hom functors \texttt{hom-\allowbreak{}lipschitz-\allowbreak{}post-\allowbreak{}type} --- proven, \texttt{hom-\allowbreak{}lipschitz-\allowbreak{}pre-\allowbreak{}type} --- proven.

\paragraph{\texttt{uniformly-\allowbreak{}continuous,\allowbreak{} on metric-\allowbreak{}space}}
\emph{Declared in} \texttt{theorem-\allowbreak{}library/\allowbreak{}categories.scm}.

Objects: the models of \texttt{is-\allowbreak{}metric-\allowbreak{}space}.  Arrows: \texttt{is-\allowbreak{}uniformly-\allowbreak{}continuous-\allowbreak{}arrow(a,\allowbreak{} b,\allowbreak{} f)} --- f is a uniformly continuous map from a to b.  Hom-set: \texttt{hom-\allowbreak{}uniformly-\allowbreak{}continuous(a,\allowbreak{} b)} --- the uniformly continuous maps from a to b.

Definition (\texttt{is-\allowbreak{}uniformly-\allowbreak{}continuous-\allowbreak{}arrow-\allowbreak{}def}, definitional):
\begin{VNBsnip}
forall([a, b, f],
  is-uniformly-continuous-arrow(a, b, f)
  iff is-metric-space(a)
      and is-metric-space(b)
      and f in fun(pts(a), pts(b))
      and is-uniformly-continuous(a, b, f))
\end{VNBsnip}

Laws: \texttt{hom-\allowbreak{}uniformly-\allowbreak{}continuous-\allowbreak{}id} --- proven; \texttt{hom-\allowbreak{}uniformly-\allowbreak{}continuous-\allowbreak{}compose} --- proven; \texttt{hom-\allowbreak{}uniformly-\allowbreak{}continuous-\allowbreak{}in-\allowbreak{}set} --- proven; membership \texttt{hom-\allowbreak{}uniformly-\allowbreak{}continuous-\allowbreak{}member-\allowbreak{}iff} --- proven; Hom functors \texttt{hom-\allowbreak{}uniformly-\allowbreak{}continuous-\allowbreak{}post-\allowbreak{}type} --- proven, \texttt{hom-\allowbreak{}uniformly-\allowbreak{}continuous-\allowbreak{}pre-\allowbreak{}type} --- proven.

\subsubsection{Inclusions}
Theorems of the form $P(a, b, f) \Rightarrow Q(a, b, f)$ for two arrow predicates, or $H_1(a, b) \subseteq H_2(a, b)$ for two hom-sets, found by shape.

\begin{itemize}
\item \texttt{hom-\allowbreak{}metric-\allowbreak{}space-\allowbreak{}is-\allowbreak{}lipschitz-\allowbreak{}arrow} (proven): \texttt{is-\allowbreak{}hom-\allowbreak{}metric-\allowbreak{}space(a,\allowbreak{} b,\allowbreak{} f) implies is-\allowbreak{}lipschitz-\allowbreak{}arrow(a,\allowbreak{} b,\allowbreak{} f)}
\item \texttt{hom-\allowbreak{}normed-\allowbreak{}vector-\allowbreak{}space-\allowbreak{}is-\allowbreak{}bounded-\allowbreak{}linear-\allowbreak{}arrow} (proven): \texttt{is-\allowbreak{}hom-\allowbreak{}normed-\allowbreak{}vector-\allowbreak{}space(a,\allowbreak{} b,\allowbreak{} f) implies is-\allowbreak{}bounded-\allowbreak{}linear-\allowbreak{}arrow(a,\allowbreak{} b,\allowbreak{} f)}
\item \texttt{lipschitz-\allowbreak{}arrow-\allowbreak{}is-\allowbreak{}uniformly-\allowbreak{}continuous-\allowbreak{}arrow} (proven): \texttt{is-\allowbreak{}lipschitz-\allowbreak{}arrow(a,\allowbreak{} b,\allowbreak{} f) implies is-\allowbreak{}uniformly-\allowbreak{}continuous-\allowbreak{}arrow(a,\allowbreak{} b,\allowbreak{} f)}
\item \texttt{uniformly-\allowbreak{}continuous-\allowbreak{}arrow-\allowbreak{}is-\allowbreak{}continuous-\allowbreak{}arrow} (proven): \texttt{is-\allowbreak{}uniformly-\allowbreak{}continuous-\allowbreak{}arrow(a,\allowbreak{} b,\allowbreak{} f) implies is-\allowbreak{}continuous-\allowbreak{}arrow(a,\allowbreak{} b,\allowbreak{} f)}
\end{itemize}
\par\endgroup


\section{Viewing one structure as another: \texttt{view-as}}
\label{sec:view-as}

A ring's additive part $(\mathtt{carr}, \mathtt{add}, \mathtt{zero},
\mathtt{neg})$ is an abelian group; its multiplicative part
$(\mathtt{carr}, \mathtt{mul}, \mathtt{one})$ is a monoid.  A
\emph{view-as} records such a relationship: a way of \emph{seeing} an
instance of one structure as an instance of another.  The relationship
is asymmetric (the ring's additive group forgets the multiplication;
there is no route back) and not unique (the same ring has both an
additive-group view and a multiplicative-monoid view).  We call these
\emph{views}, not ``bridges'': a bridge sounds symmetric, and the
asymmetry and the multiplicity are exactly the point.

A view is a \emph{functoid} (Remark~\ref{rem:functoids-quantified}):
\texttt{def-functor} installs the \texttt{def-functoid} that builds the
target instance as the list of the source's slots, plus the definitional
axiom that the result satisfies the target predicate.  Its domain, the
class of all instances of the source structure, is a proper class; a
view is a member of nothing and, like every functoid, outside the range
of the quantifiers.

\subsection{The \texttt{def-functor} form}

\begin{VNBcode}
(def-functor NAME
  SOURCE-STRUCT  (src-comp-1 src-comp-2 ... src-comp-n)
  TARGET-STRUCT  (tgt-slot-1 tgt-slot-2 ... tgt-slot-n))
\end{VNBcode}

The two component lists are flat lists of accessor symbols of equal
length, matched position by position: the value at the target's $k$-th
slot is the source's $k$-th named component.  The target list must be
the target structure's slots in their declared order (the form checks
this); the source list names, in that same order, which source accessor
fills each target slot.  So a view is a \textbf{slot map written in the
target's coordinate system}.

A view-as is a \emph{forgetful functor} between structure categories, and a
forgetful functor does exactly one or both of two things (recall
\S\ref{sec:shape-vs-property}: a structure is a \emph{shape} plus
\emph{properties}):

\begin{itemize}
\item \textbf{forget shape} --- the source list keeps only a sub-list of the
  source's slots; the rest drop.  (A ring's additive group forgets
  \texttt{mul}/\texttt{one}; a field's multiplicative group keeps the
  \texttt{non-zero} carrier slot and forgets the additive ones.)
\item \textbf{forget properties} --- \texttt{is-\textit{target}} folds fewer laws than
  \texttt{is-\textit{source}}, so the image satisfies less (a commutative ring viewed
  as a ring forgets commutativity).
\end{itemize}

\noindent These are the \emph{only} two forgetful operations.  Writing the kept
slots in the target's coordinate order (a ring's \texttt{add} filling a group's
\texttt{mul} slot, or a non-first slot serving as the carrier) is just
notation, not a third kind --- it is what lets shape-forgetting be expressed at
all.  Genuinely \emph{non}-forgetful functors --- the quotient, the completion
--- are constructions, not views; they are written as plain functoids.

\begin{VNBcode}
;;; A ring's additive structure is an abelian group.
;;; ABELIAN-GROUP's slots, in order, are (CARR OPR IDEN INV); the view fills
;;; them from the ring's (CARR ADD ZERO NEG):
(def-functor 'RING-ADDITIVE-AG
  'RING          '(CARR ADD ZERO NEG)
  'ABELIAN-GROUP '(CARR OPR IDEN INV))
\end{VNBcode}

\subsection{What \texttt{def-functor} installs}

\begin{enumerate}
\item \textbf{A constructor functoid} \texttt{(\textit{name} s)} whose defining
  equation gathers the named source components into a target-shaped list:
  \[
    \mathtt{(ring\text{-}additive\text{-}ag}\;r) \;=\;
      \mathtt{list}(\mathtt{carr}(r),\, \mathtt{add}(r),\,
                    \mathtt{zero}(r),\, \mathtt{neg}(r)).
  \]
  Inside a proof, \texttt{(mul (ring-additive-ag r))} reduces to
  \texttt{(add r)} by unfolding the functoid and the accessor.

\item \textbf{A typing axiom}, named
  \texttt{\textit{name}-is-\textit{target}} (here
  \texttt{ring-additive-ag-is-abelian-group}):
  \[
    \forall r.\; \mathit{is\text{-}source}(r) \;\Rightarrow\;
      \mathit{is\text{-}target}(\mathit{name}(r)).
  \]
  During the library-build phase this is installed as an \emph{axiom}
  (a warranted assertion, Section~\ref{sec:warrants}), not proved;
  promoting it to a theorem is a separate task.

\item \textbf{Auto-specializations.}  Every theorem of the form
  $\forall s.\; \mathit{is\text{-}target}(s) \Rightarrow P[s]$ already in
  the library is re-installed as
  $\forall r.\; \mathit{is\text{-}source}(r) \Rightarrow P[\mathit{name}(r)]$,
  with the target-slot accessors reduced to source components (so
  $\mathtt{mul}(s)$ becomes $\mathtt{add}(r)$, etc.), under the name
  \texttt{\textit{original}-\textit{name}}.  Thus declaring the additive
  view immediately carries every abelian-group theorem over to the
  additive structure of every ring.
\end{enumerate}

\begin{remark}
A view-as is a covariant functor in waiting: the constructor is the
object map (an instance of the source becomes an instance of the
target), and the auto-specialization transports the target's theorems
along it.  Homomorphisms are now first-class
(Section~\ref{sec:morphisms}); the morphism half --- transporting a
homomorphism across a view --- remains deferred.  The
auto-specializer currently matches only theorems guarded by exactly
\texttt{is-\textit{target}}; it does not chain through structure-inclusion axioms,
so same-shape refinements (a commutative ring's additive view, say) are
declared directly rather than derived.
\end{remark}

\begin{remark}
Same-shape inclusions --- every abelian group is a group, every
commutative ring is a ring --- are \emph{not} views.  They are recorded
as plain $\texttt{is-}X \Rightarrow \texttt{is-}Y$ typing axioms, because
the constructor would be a redundant identity list.  A \texttt{view-as}
earns its keep when source and target have different shapes, so the
constructor genuinely projects.
\end{remark}

\subsection{Constructed functors: \texttt{def-constructed-functor}}
\label{sec:constructed-functor}

A \texttt{view-as} selects the target's slots from among the source's --- its
object map is a re-shuffling of accessors.  Some functors are not like that.
The metric topology $\tau_d$ is \emph{computed} from the distance, not read off a
slot; the distance of a normed field is \emph{computed} from the norm.  For these
the object map is an arbitrary term, and the form is
\texttt{def-constructed-functor}:

\begin{VNBcode}
;;; A normed field is a metric space, with d(x, y) = norm(x - y).
(def-constructed-functor 'NF-METRIC-SPACE 'NORMED-FIELD 'METRIC-SPACE '(nf)
  '(LIST (CARR nf)
         (VNB-LAMBDA (LIST x y) (CARTESIAN (CARR nf) (CARR nf)) ((FNRM nf) ((ADD nf) x ((NEG nf) y))))))
\end{VNBcode}

The arguments are the functor name, its source and target species, the
parameter list, and the defining term --- installed as an ordinary functoid
(Section~\ref{sec:def-functoid}), so \texttt{(nf-metric-space nf)} rewrites to
that tuple.

The decisive difference from \texttt{view-as} is what it \emph{claims}.  A view's
typing axiom is installed as a warranted assertion; a constructed functor
\textbf{asserts nothing}.  ``An isometry is continuous'' and ``the norm-difference
is a metric'' are theorems with content, not free consequences of a slot
correspondence.  So the form only \emph{records} the two obligations any functor
owes, and \texttt{functor-obligation-audit} reports at load any still
undischarged:
\[
  \begin{aligned}
    \texttt{\textit{name}-is-\textit{tgt}} &:\quad
      \mathit{is\text{-}src}(r) \Rightarrow \mathit{is\text{-}tgt}(\mathit{name}(r)), \\
    \texttt{\textit{name}-functorial} &:\quad
      \mathit{is\text{-}hom\text{-}src}(a, b, f) \Rightarrow
      \mathit{is\text{-}hom\text{-}tgt}(\mathit{name}(a), \mathit{name}(b), f).
  \end{aligned}
\]
Both are available as goals, ready for \texttt{sp}.  A functor whose obligations
are unproved is a functor you do not yet have.


\part{Self-verification}\label{part:meta}

%% =========================================================
%% ch-meta-claim.tex -- Part 2 (\ref{part:meta}), chapter 1: the claim and the trusted base.
%% Cites the design specification as \DS{n} (app-design-spec.tex).
%% Uses only the manual's own preamble; the two macros below are provided here in case this
%% chapter is built before ch-proofs.tex defines its figure macros.
\providecommand{\DS}[1]{\hyperref[ds:#1]{DS-#1}}
\providecommand{\VNBfig}[3]{\[\frac{\displaystyle #2}{\displaystyle #3}\;\;\mbox{\texttt{#1}}\]}
\providecommand{\VNBnohyp}{\phantom{\Gamma}}

\chapter{Self-verification: the claim and the trusted base}
\label{ch:meta-claim}
\index{META}\index{self-verification}

This part describes what VNB does to check its own kernel, what it
proves about it, and what it leaves to hand-written argument or to
assumption.  The work has two purposes (\DS{1}): to make VNB credible as
a proof checker, and to show that code can be modelled formally and its
correctness proven relative to a stated base.

The operations and their checkers are described in
Sections~\ref{sec:kernel-boundary} and~\ref{sec:inference-checked}; the
closed list of files that may record an inference in
Section~\ref{sec:kernel-checked}; the proof certificate in
Section~\ref{sec:proof-certificate}.  They are not repeated here.

%% =========================================================
\section{The claim}\label{sec:meta-the-claim}

\subsection{Consistency is relative}\label{sec:meta-relative}

VNB cannot prove its own consistency, and no check in this part can
either.  This is G\"odel's second incompleteness theorem~\cite{Godel1931}:
a consistent theory with effectively listed axioms, strong enough to
express arithmetic, does not prove its own consistency.  VNB is such a
theory.

So no statement of the form ``VNB cannot prove \textsc{falsity}'' is made
(\DS{2}).  Every consistency statement is relative.  A derivation of
\textsc{falsity}, once found, is conclusive.  Finding none is evidence,
never proof.

\subsection{Version-X}

Every claim is relative to a fixed class theory, called
\emph{version-X}.  Version-X is NBG with global choice (\DS{3}).  Global
choice is needed because \texttt{choice-axiom} gives
$\mathrm{CHOICE}(A) \in A$ for every inhabited class $A$.
$\mathrm{CHOICE}(A)$ denotes exactly when $A$ is inhabited (\DS{52}).  The
definedness certificate never accepts a $\mathrm{CHOICE}$ term by its
shape: it is not a total head.

What VNB calls a class is read as a class of version-X (\DS{4}).  The
oracles are excluded from the formal reasoning and assumed correct
(\DS{5}).

NBG proves the same statements about sets as ZF, and is consistent
exactly when ZF is.  The claim therefore says that VNB is as safe as
ordinary set theory, under the assumptions listed below.

\subsection{The statement}

The claim, in the form this part establishes it, is:
\begin{center}
\emph{Relative to version-X, and assuming the oracles correct, every
inference the kernel accepts is admissible over the core.}
\end{center}
``Admissible over the core'' is defined in
Chapter~\ref{ch:meta-k2}: whenever the hypotheses of the inference are
derivable by the rules of a small, fixed system (the \emph{core}), so is
its conclusion.

\subsection{From the claim to relative consistency}\label{sec:meta-chain}

The step from the claim to relative consistency has five links.  Only
the first is a VNB theorem.
\begin{enumerate}
\item \emph{Admissibility} (formal).  Each non-oracle kernel rule is
  admissible over the core.  This is K2 (\DS{8}), stated rule by rule in
  Chapter~\ref{ch:meta-k2}; its coverage is in
  Section~\ref{sec:meta-coverage}.
\item \emph{Faithfulness} (hand review).  Each transcribed relation is the
  relation its Scheme checker enforces.  This is K3 (\DS{9}).
\item \emph{Soundness of the core} (hand-written).  Each core rule is
  sound for version-X under the reading of classes of \DS{4} and the
  reading of functoids of Chapter~\ref{ch:meta-functoids}; each axiom of
  the catalog holds under that reading (\DS{6}).  The argument is
  Chapter~\ref{ch:meta-soundness}.
\item \emph{Oracles} (assumption, \DS{5}).
\item \emph{Infrastructure} (assumption).  The expression layer, the write
  point and grounding, the install, the certificate store and MIT Scheme
  behave as described.
\end{enumerate}

\begin{prop}[informal; conditional]\label{prop:meta-relcons}
If version-X is consistent, links 1--5 hold, and every installed theorem
was installed by \texttt{qed} over a grounded deduction graph whose every
inference passed its checker, then the installed theorems do not include
\textup{\textsc{falsity}}.
\end{prop}

The proposition is stated to make the dependency explicit.  It is not a
VNB theorem.  What is not yet established for each link is listed in
Chapter~\ref{ch:meta-limits}.

The META theorems are themselves VNB theorems.  They presuppose the
consistency of the system they are proven in, and many of them use the
oracles \texttt{arith} and \texttt{ineq}.

%% =========================================================
\section{The trusted base}\label{sec:meta-base}

\subsection{The write point and the check}

\texttt{dg-apply-rule!} (\texttt{deduction-graphs.scm}) is the only
procedure that writes an inference into a deduction graph.  The
\emph{kernel} consists of all the procedures that call it (\DS{11}).
Before anything is written, \texttt{dg-apply-rule!} hands the proposed inference to the checker
registered for the rule's name (\DS{12}).  A missing checker or a
refusal is an error, and nothing reaches the graph.

Each rule is a relation between the premises of an inference, its
conclusion and its side data.  A checker receives an inference after it
is built and tests whether these stand in the relation of its rule; it
builds nothing itself.  It is written
independently of the procedure that builds the inference, and calls
neither that procedure nor a copy of it.  So the builders and the
tactics carry no soundness: a wrong builder can only fail to prove.

The four checker files register 64 checkers: 40 in
\texttt{rule-checkers-logic.scm}, 13 in \texttt{-schema.scm}, 4 in
\texttt{-rewrite.scm} and 7 in \texttt{-oracle.scm}.  The 57 that are
not oracles are the subject of Chapter~\ref{ch:meta-k2}.

No rule is accepted on trust: the checker registered for each rule
tests the inference, and none accepts every inference.  Two procedures
could register such an always-accepting checker: \texttt{dg-check-arm-with!}, for a documented
tag without a checker, and \texttt{rco-trust-oracle!}.  The first
registers nothing in a complete load, because every documented tag has
its checker; the second has no caller.  The suite checks both.

\paragraph{One exception: the definedness test is shared.}  Several
rules are sound only when some term denotes.  For instance,
\texttt{forall-elim} passes from $\forall x\,A$ to $A[t/x]$ only for a
defined $t$.  When no hypothesis of the inference says that $t$ is
defined, the checker must decide it, and for this it does not have a
test of its own.  It calls \texttt{chk-elim-certified?}
(\texttt{rule-checkers-logic.scm}), which tries a few clauses and then
falls back on \texttt{pi--defined?}, the same definedness test that the
builders use.

The checkers of \texttt{forall-elim}, \texttt{forsome-intro},
\texttt{if-true}, \texttt{if-false}, \texttt{tuples-intro},
\texttt{union-intro} and \texttt{lambda-type} take this route.  The
checker of \texttt{reflexivity} and the value test of the macete door
call \texttt{pi--defined?} directly.

The clauses of \texttt{pi--defined?} are themselves proven: each is
transcribed as a VNB relation and derived over the core, so a term that
a clause accepts does denote (Section~\ref{sec:meta-lemma-star}).  What
the checkers share with the builders is only the Scheme code that
applies those clauses.  A mistake in that code, one that reports a term
as defined when no clause covers it, would reach the builder and its
checker alike, where elsewhere the checker would catch the builder's
mistake.  Whether Scheme code computes its transcribed relation is the
question K3 asks of every checker (Chapter~\ref{ch:meta-limits}).  A
second definedness
test, written separately for the checkers, would remove the overlap; if
one is adopted, it accepts no more than \texttt{pi--defined?} does
(\DS{32}).

\subsection{How ``every inference is checked'' is enforced}

\begin{itemize}
\item The switch \texttt{*dg-check-inferences?*} is armed at load, after
  the four checker files.  \texttt{VNB\_NO\_RULE\_CHECK=1} is the control
  that turns it off.
\item A refusal prints a line beginning \texttt{;VNB RULE CHECKER
  REFUSED}, also inside \texttt{quietly}, and is recorded.  The load ends
  with \texttt{inference checking: $N$ verified, $M$ refusal(s)}; an
  ordinary load (Section~\ref{sec:meta-gates}) is an error when $M>0$.
\item The suite checks that the switch is on after the load, that more
  than a thousand inferences were verified, that no rule is accepted on
  trust, and that an unregistered tag is refused.
\item Each checker has one accepted and one refused control, and a new
  kernel operation comes with its checker and both controls (\DS{13}).
\end{itemize}

\subsection{The code that may write}

Only the files listed in \texttt{*cert-kernel-caller-files*}
(\texttt{certificates.scm}) may call \texttt{dg-apply-rule!}:
\texttt{primitive-inferences}, \texttt{macetes}, \texttt{library},
\texttt{arith-eval} and the four oracle files.  The fatal
\texttt{kernel-callers-audit} refuses a call, or a use as a value, of
\texttt{dg-apply-rule!} anywhere else (\DS{11};
Section~\ref{sec:kernel-checked}).

\subsection{The axioms and the primitive shelf}

The axioms are assumed.  The set and class formation schemas are kernel
rules, not entries of the axiom table.  The shelf of primitives grows
only by an explicit decision, one fact at a time, with its claim
recorded in the file (\DS{15}).

The stamps \texttt{primitive} and \texttt{definitional} bill nothing.  A
stamp therefore records what it assumes (\DS{15}), and a guarded
definitional equation says nothing off its guards (\DS{31}).

\subsection{The gates that stop a load}\label{sec:meta-gates}

An \emph{ordinary load} is any load of VNB without the variable
\texttt{VNB\_KEEP\_GOING=1}: the prover, the band build and the exam.
That variable selects a development mode, the \emph{keep-going load}, in
which an error in a proof file is recorded and the load goes on, so that
every broken file is listed in one load.  A band built that way has holes
in its library, and the end of the load says so.  (This is unrelated to
\texttt{VNB\_CERTIFIED=strict}, which concerns certificates,
Chapter~\ref{ch:provenance}.)

Among the end-of-load gates, the following stop an ordinary load and are
only reported in a keep-going one:
\begin{itemize}
\item inference checking: any refusal;
\item the \texttt{qed}-failure gate: a \texttt{qed} that failed;
\item the after-\texttt{qed} gate: a command run on a proof that a
  \texttt{qed} had just closed;
\item the guarded-definition gate: a guarded defined head in a
  certificate table or typed in $\mathrm{FUN}(\ldots)$ (\DS{31});
\item the view gate: a view typing installed without its proof (\DS{36});
\item \texttt{connective-arity-audit} and \texttt{constant-binder-audit};
\item the category obligations of \texttt{declare-category!};
\item the predicativity gates for \texttt{comp} and for the binder terms
  (\DS{38});
\item the functoid-hypotheses gate: a violation of H1--H4 or H6
  (\DS{35}; Chapter~\ref{ch:meta-functoids});
\item the library gate;
\item \texttt{kernel-callers-audit}.
\end{itemize}
Under keep-going these gates report and the load continues.  The
free-variable audit, the head-registry sweep and the duplicate audits
warn only.

\subsection{The certified load}\label{sec:meta-certified}

The everyday load installs a theorem from its certificate without
running its proof (Section~\ref{sec:proof-certificate}).  The inferences
of a certified theorem are therefore not re-checked in that load: they
were checked by the run that wrote the record.

The certificate key covers the statement, the definitions of its defined
heads, followed transitively, and every definition or macete a checker
read while accepting an inference (\DS{17}).  The kernel hash covers the
files that may write, \texttt{deduction-graphs}, \texttt{expressions},
\texttt{sequents} and the four checker files; a change to any of them
invalidates every certificate.

The run in which every proof is re-executed is the nightly exam
(Chapter~\ref{ch:checking}).  A kernel change lands only after an exam of
a tree that contains it (\DS{16}).

%% =========================================================
\section{Predicativity: staying within NBG}\label{sec:meta-predicative}

NBG allows class comprehension only over conditions whose quantifiers
range over sets.  The kernel keeps to that (\DS{38}).
\begin{itemize}
\item The membership rules of \texttt{comp}, \texttt{sep},
  \texttt{big-union} and \texttt{vnb-lambda}, and the three induction
  rules, refuse a condition or body with a quantifier not confined to
  sets.  The test reads through defined heads.
\item The definedness certificate treats a \texttt{comp}, \texttt{sep},
  \texttt{big-union} or \texttt{vnb-lambda} term as defined only when its
  condition or body is predicative.
\item \texttt{comp-predicativity-audit} and
  \texttt{binder-predicativity-audit} (\texttt{proof-debt.scm}) list every
  impredicative term named by a proof, an installed formula or a functoid
  body.  Both gates stop an ordinary load.
\item The induction census prints
  \texttt{;; NBG induction census: $n$ impredicative induction(s)}; with
  \texttt{*nbg-induction-fatal?*} on, $n=0$ in any load that completes.
\end{itemize}
The encoded core is NBG-strength in the same way (\DS{39}): each core
schema that forms a class carries the premise $\mathrm{NBG\text{-}OK}$ on
its condition or body (Section~\ref{sec:meta-core}).


%% =========================================================
%% ch-meta-k2.tex -- Part 2 (\ref{part:meta}), chapter 2: admissibility of the kernel rules over the core (K2).
%% Sources: structure-library/meta-*.scm, theorem-library/meta-*-laws.scm; the per-rule ledger
%% reference/K2-LEDGER.md.  Cites the design specification as \DS{n}.
\providecommand{\DS}[1]{\hyperref[ds:#1]{DS-#1}}
\providecommand{\VNBfig}[3]{\[\frac{\displaystyle #2}{\displaystyle #3}\;\;\mbox{\texttt{#1}}\]}
\providecommand{\VNBfigS}[3]{\[\frac{\displaystyle\begin{array}{c}#2\end{array}}{\displaystyle #3}\;\;\mbox{\texttt{#1}}\]}
\providecommand{\VNBnohyp}{\phantom{\Gamma}}
\providecommand{\FV}{\mathrm{FV}}
%% A core figure in a table cell: \Rl{hypotheses}{conclusion}.
\providecommand{\Rl}[2]{$\frac{#1}{#2}$}
\providecommand{\Sq}{\Rightarrow}

\chapter{Admissibility over the core}
\label{ch:meta-k2}

This chapter states what is proven, inside VNB, about the kernel's
rules.  The route is K2 (\DS{8}): formulas and sequents are encoded as
sets, a small system of rules called the \emph{core} is defined, and each
non-oracle kernel rule is proven \emph{admissible} over it.  What has to
be trusted is thereby reduced from the 57 checked relations to the core.

The definitions are in \texttt{structure-library/meta-*.scm}, the
theorems in the files of the same names in \texttt{theorem-library/},
all in the library \texttt{meta}.  Every theorem bills
\texttt{modulo 0}.

%% =========================================================
\section{The encoding}\label{sec:meta-encoding}

The rules are schemas, and VNB has no formula variables, so formulas
are VNB objects.  The encoding is locally nameless (\DS{10}): a bound
variable is a de~Bruijn index, a free variable a name.
Alpha-equivalence is then equality, and substitution for a free name is
structural.

Every node is a literal list whose first entry is a numeral tag:
\[
\begin{gathered}
\mathrm{EFVAR}(n)=[0,n],\quad \mathrm{EBVAR}(k)=[1,k],\quad
\mathrm{ECONST}(c)=[2,c],\\
\mathrm{EAPP}(h,l)=[3,h,l],\quad \mathrm{EBIND}(h,n,a,b)=[4,h,n,a,b].
\end{gathered}
\]
$h$ is a numeral head code, $l$ a list of nodes, $a$ the list of
arguments outside a binder (the domain of \texttt{sep},
\texttt{big-union}, \texttt{vnb-lambda}) and $b$ its body.

$\mathrm{EXPR}$ is a union over $\mathtt{nn}$ of stages.  The free
names $\FV$, the opening of an index $\mathrm{OPEN}$, the substitution
$\mathrm{SUBST}$ and local closure $\mathrm{LC}$ are defined by
recursion on the stage, with the usual substitution lemmas.  A
well-formedness predicate separates terms from formulas by head code.
$\mathrm{CLF}$ is the set of sentences (closed formulas), $\mathrm{CLT}$
the set of closed terms.

Predicate constants have codes ($\mathrm{PRED\text{-}CODE}$).  A
functoid application is $\mathrm{EAPP}(h,l)$ with $h\geq 32$, read
through a table $\mathrm{FUN\text{-}DEF}$ of rows $(h, n, b)$
(\DS{44}).  The tables hold the rows the META theorems need; filling
them for the whole library is the translator's task (K3).

In what follows formulas are written in ordinary notation for their
encodings: $A \wedge B$ for $\mathrm{ENC\text{-}AND}(A,B)$, $B[x:=t]$ for
$\mathrm{OPEN}(B,0,t)$, and so on.

%% =========================================================
\section{Theory, sequents and derivability}\label{sec:meta-derivable}

Fix a \emph{theory} $T$: a set of encoded sentences, standing for the
axioms of the catalog and the theorems installed at the time.  A cited
theorem is a member of $T$ and nothing more.

A \emph{sequent} $\Gamma \Longrightarrow A$ is a pair $[\Gamma, A]$ with
$\Gamma$ a finite set of sentences ($\mathrm{CTX\text{-}SET}$) and $A$ a
sentence; $\mathrm{SEQ\text{-}SET}$ is the set of sequents.  Contexts are
sets, so ``$\Gamma, A$'' is $\Gamma \cup \{A\}$.

Each core rule is a relation $\mathrm{CR\text{-}}x(T, \mathit{hs}, c)$
between a list $\mathit{hs}$ of hypothesis sequents and a conclusion
sequent $c$.  $\mathrm{CORE\text{-}RULE}(T, \mathit{hs}, c)$ is their
disjunction.  $\mathrm{DERIVABLE}(T)$ is the least subset of
$\mathrm{SEQ\text{-}SET}$ closed under
$\mathrm{CORE\text{-}RULE}(T,\cdot,\cdot)$.

We write $\vdash_T \Gamma \Longrightarrow A$ for
$[\Gamma, A] \in \mathrm{DERIVABLE}(T)$.  Each core rule has a closure
law, generated in \texttt{meta-core-laws.scm} and named
\texttt{derivable-}$x$: if the hypotheses of an instance are derivable,
so is its conclusion.

%% =========================================================
\section{The core}\label{sec:meta-core}

The core is chosen so that each member is a rule of natural deduction
with its side condition, or a one-line schema that a reader checks
against the intended meaning of a constructor (\DS{40}).  A relation
enters the core only by a design decision, with its claim written beside
it in \texttt{meta-core.scm} (\DS{14}).

$\mathrm{CORE\text{-}RULE}$ has 56 disjuncts.  They are listed below as
figures, hypotheses above the line.  A side condition follows in
brackets.  In the set-formation and definedness schemas, $p$ or $b$
carries the premise $\mathrm{NBG\text{-}OK}$ (every quantifier confined
to sets, read through the tables), written [NBG] (\DS{39}).

{\small
\begin{longtable}{@{}p{0.22\linewidth}p{\dimexpr0.78\linewidth-2\tabcolsep\relax}@{}}
\toprule
core rule & figure \\
\midrule
\endhead
\multicolumn{2}{@{}l}{\emph{Logic} (21)} \\
and-intro & \Rl{\Gamma\Sq A \quad \Gamma\Sq B}{\Gamma\Sq A\wedge B} \\
and-elim & \Rl{\Gamma, A, B\Sq G}{\Gamma\Sq G}\ \ [$A\wedge B\in\Gamma$] \\
or-intro-left, -right & \Rl{\Gamma\Sq A}{\Gamma\Sq A\vee B}\ \ \Rl{\Gamma\Sq B}{\Gamma\Sq A\vee B} \\
or-elim & \Rl{\Gamma,A\Sq G \quad \Gamma,B\Sq G}{\Gamma\Sq G}\ \ [$A\vee B\in\Gamma$] \\
implies-intro & \Rl{\Gamma,A\Sq B}{\Gamma\Sq A\to B} \\
backchain & \Rl{\Gamma\Sq A}{\Gamma\Sq G}\ \ [$A\to G\in\Gamma$] \\
not-intro & \Rl{\Gamma,A\Sq\textsc{falsity}}{\Gamma\Sq\neg A} \\
not-elim & \Rl{\VNBnohyp}{\Gamma\Sq G}\ \ [$\neg A\in\Gamma$, $A\in\Gamma$] \\
iff-intro & \Rl{\Gamma,A\Sq B \quad \Gamma,B\Sq A}{\Gamma\Sq A\leftrightarrow B} \\
iff-elim & \Rl{\Gamma, A\to B, B\to A\Sq G}{\Gamma\Sq G}\ \ [$A\leftrightarrow B\in\Gamma$] \\
forall-intro & \Rl{\Gamma\Sq B[x:=y]}{\Gamma\Sq\forall x.\,B}\ \ [$y\notin\FV(\Gamma)\cup\FV(B)$] \\
forall-elim & \Rl{\Gamma, B[x:=t]\Sq G \quad \Gamma\Sq t=t}{\Gamma\Sq G}\ \ [$\forall x.\,B\in\Gamma$, $t\in\mathrm{CLT}$] \\
forsome-intro & \Rl{\Gamma\Sq B[x:=t] \quad \Gamma\Sq t=t}{\Gamma\Sq\exists x.\,B}\ \ [$t\in\mathrm{CLT}$] \\
forsome-elim & \Rl{\Gamma, B[x:=y]\Sq G}{\Gamma\Sq G}\ \ [$\exists x.\,B\in\Gamma$, $y\notin\FV(\Gamma)\cup\FV(G)$] \\
assumption & \Rl{\VNBnohyp}{\Gamma\Sq A}\ \ [$A\in\Gamma$] \\
theorem-assumption & \Rl{\VNBnohyp}{\Gamma\Sq A}\ \ [$A\in T$] \\
cut & \Rl{\Gamma\Sq A \quad \Gamma,A\Sq G}{\Gamma\Sq G} \\
weakening & \Rl{\Delta\Sq G}{\Gamma\Sq G}\ \ [$\Delta\subseteq\Gamma$] \\
proof-by-contradiction & \Rl{\Gamma,\neg G\Sq\textsc{falsity}}{\Gamma\Sq G} \\
truth-intro & \Rl{\VNBnohyp}{\Gamma\Sq\textsc{truth}} \\
\midrule
\multicolumn{2}{@{}l}{\emph{Equality} (4)} \\
eq-subst, eq-subst-sym & \Rl{\Gamma\Sq P[z:=t]}{\Gamma\Sq P[z:=s]}\ \ [one of $s=t$, $s\equiv t$, $t=s$, $t\equiv s$ in $\Gamma$] \\
reflexivity & \Rl{\VNBnohyp}{\Gamma\Sq t=t}\ \ [$t$ a variable or an atomic constant] \\
quasi-reflexivity & \Rl{\VNBnohyp}{\Gamma\Sq t\equiv t} \\
\midrule
\multicolumn{2}{@{}l}{\emph{Structural membership and tuples} (7)} \\
union-mem & \Rl{\Gamma\Sq S_1=S_1 \ \cdots\ \Gamma\Sq S_n=S_n}{\Gamma\Sq x\in\bigcup(S_1..S_n)\leftrightarrow x\in S_1\vee\cdots\vee x\in S_n} \\
intersection-mem & \Rl{\VNBnohyp}{\Gamma\Sq x\in\bigcap(S_1..S_n)\leftrightarrow x\in S_1\wedge\cdots\wedge x\in S_n}\ \ [$n\geq 1$] \\
cartesian-mem & \Rl{\VNBnohyp}{\begin{array}{@{}l@{}}\Gamma\Sq x\in\mathrm{CART}(A_1..A_n)\leftrightarrow{}\\ \quad x\in\mathrm{TUPLES}(\mathrm{SET})\wedge\mathrm{LENGTH}(x)=n\wedge{}\\ \quad\bigwedge_i\mathrm{NTH}(i,x)\in A_i\end{array}}\ \ [$n\geq1$] \\
tuples-mem & \Rl{\Gamma\Sq A=A}{\Gamma\Sq[e_1..e_n]\in\mathrm{TUPLES}(A)\leftrightarrow\bigwedge_i e_i\in A} \\
nth-eq & \Rl{\Gamma\Sq L=L}{\Gamma\Sq\mathrm{NTH}(k,L)\equiv e_k}\ \ [$L=[e_1..e_n]$, $1\leq k\leq n$] \\
length-eq & \Rl{\Gamma\Sq L=L}{\Gamma\Sq\mathrm{LENGTH}(L)\equiv n} \\
tuple-eq & \Rl{\Gamma\Sq L=L \quad \Gamma\Sq M=M}{\Gamma\Sq L=M\leftrightarrow\bigwedge_i a_i=b_i}\ \ [$L=[a_1..a_n]$, $M=[b_1..b_n]$] \\
\midrule
\multicolumn{2}{@{}l}{\emph{Conditional terms} (2)} \\
if-true & \Rl{\Gamma\Sq p \quad \Gamma, \mathrm{IF}(p,a,b)\equiv a\Sq G}{\Gamma\Sq G} \\
if-false & \Rl{\Gamma\Sq \neg p \quad \Gamma, \mathrm{IF}(p,a,b)\equiv b\Sq G}{\Gamma\Sq G} \\
\midrule
\multicolumn{2}{@{}l}{\emph{Set formation, descriptions, functions} (9)} \\
sep-sethood & \Rl{\Gamma\Sq A\in\mathrm{SET}}{\Gamma\Sq\mathrm{SEP}\,x{\in}A.\,p\in\mathrm{SET}}\ \ [NBG] \\
sep-mem & \Rl{\VNBnohyp}{\Gamma\Sq t\in\mathrm{SEP}\,x{\in}A.\,p\leftrightarrow t\in A\wedge p[x:=t]}\ \ [NBG] \\
comp-mem & \Rl{\VNBnohyp}{\Gamma\Sq t\in\mathrm{COMP}\,x.\,p\leftrightarrow t\in\mathrm{SET}\wedge p[x:=t]}\ \ [NBG] \\
big-union-sethood & \Rl{\Gamma\Sq A\in\mathrm{SET} \quad \Gamma\Sq\forall z.\,(z\in A\to b\in\mathrm{SET})}{\Gamma\Sq\bigcup z{\in}A.\,b\in\mathrm{SET}}\ \ [NBG] \\
big-union-mem & \Rl{\VNBnohyp}{\Gamma\Sq x\in\bigcup z{\in}A.\,b\leftrightarrow\exists z.\,(z\in A\wedge x\in b)}\ \ [NBG] \\
iota-def & \Rl{\Gamma\Sq\exists x.\,(B\wedge\forall y.\,(B[x:=y]\to x=y)) \quad \Gamma, B[x:=\iota x.B]\Sq G}{\Gamma\Sq G} \\
iota-in-elim & \Rl{\Gamma\Sq\iota x.B=\iota x.B \quad \Gamma, B[x:=\iota x.B]\Sq G}{\Gamma\Sq G} \\
lambda-type & \Rl{\Gamma\Sq\forall x.\,(x\in A\to b\in D) \quad \Gamma\Sq A\in\mathrm{SET} \quad \Gamma\Sq D=D}{\Gamma\Sq\lambda x{\in}A.\,b\in\mathrm{FUN}(A,D)}\ \ [NBG; one binder] \\
lambda-beta-eq & \Rl{\Gamma\Sq t\in A \quad \Gamma\Sq b[x:=t]\in\mathrm{SET}}{\Gamma\Sq(\lambda x{\in}A.\,b)(t)\equiv b[x:=t]}\ \ [NBG] \\
\midrule
\multicolumn{2}{@{}l}{\emph{Definedness} (9; \DS{42})} \\
def-strict-in-l, -r & \Rl{\Gamma\Sq w\in A}{\Gamma\Sq w=w}\ \ \Rl{\Gamma\Sq w\in A}{\Gamma\Sq A=A} \\
def-total & \Rl{\Gamma\Sq t_1=t_1 \ \cdots\ \Gamma\Sq t_n=t_n}{\Gamma\Sq h(t_1..t_n)=h(t_1..t_n)}\ \ [$h$ a total constructor] \\
def-binder & \Rl{\Gamma\Sq A=A}{\Gamma\Sq Qx{\in}A.\,b=Qx{\in}A.\,b}\ \ [$Q$ one of SEP, BIG-UNION, VNB-LAMBDA; NBG] \\
def-comp & \Rl{\VNBnohyp}{\Gamma\Sq\mathrm{COMP}\,x.\,p=\mathrm{COMP}\,x.\,p}\ \ [NBG] \\
def-list & \Rl{\Gamma\Sq e_1\in\mathrm{SET}\ \cdots\ \Gamma\Sq e_n\in\mathrm{SET}}{\Gamma\Sq[e_1..e_n]=[e_1..e_n]} \\
def-strict-arg & \Rl{\Gamma\Sq h(t_1..t_n)=h(t_1..t_n)}{\Gamma\Sq t_i=t_i}\ \ [$h$ a strict constructor] \\
def-strict-dom & \Rl{\Gamma\Sq Qx{\in}A.\,b=Qx{\in}A.\,b}{\Gamma\Sq A=A} \\
def-strict-pred & \Rl{\Gamma\Sq P(t_1..t_n)}{\Gamma\Sq t_i=t_i}\ \ [$P$ a predicate constant] \\
\midrule
\multicolumn{2}{@{}l}{\emph{Congruence, functoids, tuples} (4; \DS{43}--\DS{45})} \\
iota-cr & \Rl{\Gamma\Sq B[x:=y]\leftrightarrow C[x:=y]}{\Gamma\Sq\iota x.B\equiv\iota x.C}\ \ [$y$ fresh] \\
tbind-cr & \Rl{\Gamma\Sq B[x:=y]\leftrightarrow C[x:=y]}{\Gamma\Sq Qx{\in}A.\,B\equiv Qx{\in}A.\,C}\ \ [$Q$ SEP or COMP; for BIG-UNION, VNB-LAMBDA the premise is $b[y]\equiv c[y]$; $y$ fresh] \\
functoid-beta-eq & \Rl{\VNBnohyp}{\Gamma\Sq F(l)\equiv b[l]}\ \ [$(F,|l|,b)$ a row of FUN-DEF] \\
tuple-eta & \Rl{\Gamma\Sq x\in\mathrm{CART}(A_1..A_n)}{\Gamma\Sq x=[\mathrm{NTH}(1,x)..\mathrm{NTH}(n,x)]}\ \ [$n\geq1$] \\
\bottomrule
\end{longtable}
}

Where a strict conclusion would assert that a term denotes, the core
concludes a quasi-equation instead: the conditional rules, lambda-beta,
the two congruence rules and functoid beta.  Where the kernel reads the
definedness of a term off its certificate, the core states it as a
premise: forall-elim, forsome-intro, iota-in-elim and the list schemas.
The differences between such a core rule and the kernel's figure are the
bridges of Section~\ref{sec:meta-bridges}.

A schema whose conclusion makes a class defined that its other premises
do not carries that class's definedness as a premise (\DS{53}):
tuples-mem the class $A$, union-mem every entry $S_i$, lambda-type the
codomain $D$.  $\mathrm{TUPLES}$, $\mathrm{UNION}$ and $\mathrm{FUN}$ are
strict in these arguments, and without the premise the schema is false
at the empty tuple, at a non-denoting entry, or over an empty domain.
lambda-beta-eq carries the value's sethood, $b[x:=t]\in\mathrm{SET}$: a
pair with a proper class as its second component does not denote.

Each schema beyond the logic, equality and structural groups carries its
claim beside it in \texttt{meta-core.scm}; where a schema is sound only
under a precondition, the precondition is written there and pinned by a
suite control.  The strictness of a predicate application, for example,
is sound only because every instantiation of a definition owes the
definedness of the instance (\DS{19}, \DS{42}).  The suite checks that
the four schemas above carry their premises.

%% =========================================================
\section{The statement of K2}\label{sec:meta-k2-statement}

For each non-oracle kernel rule $R$, $\mathrm{REL\text{-}}R(T,
\mathit{hs}, c)$ is the relation its checker enforces, transcribed by
hand into VNB.  Each transcription is kept beside the checker's line
numbers, to be reviewed as a pair (\DS{9}).

An instance of $R$ is a figure
\[
  \frac{\Gamma_1 \Longrightarrow A_1 \qquad \cdots \qquad \Gamma_n \Longrightarrow A_n}
       {\Gamma \Longrightarrow A}\;\;\mbox{\textit{R}}
\]
with $\mathrm{REL\text{-}}R(T, [\Gamma_1\Sq A_1, \ldots,
\Gamma_n\Sq A_n], \Gamma\Sq A)$.

\begin{dfn}
$R$ is \emph{admissible over the core} of $T$ when every instance with
derivable hypotheses has a derivable conclusion:
\[
  \frac{\vdash_T \Gamma_1 \Longrightarrow A_1 \qquad \cdots \qquad \vdash_T \Gamma_n \Longrightarrow A_n}
       {\vdash_T \Gamma \Longrightarrow A}
\]
\end{dfn}

\begin{claim}[K2, per rule]\label{claim:meta-k2}
Every non-oracle kernel rule is admissible over the core.  As a VNB
theorem, for a rule $R$:
\[
\begin{aligned}
&\forall T, \mathit{hs}, c.\;
\mathrm{REL\text{-}}R(T,\mathit{hs},c) \Rightarrow{}\\
&\qquad\bigl(\mathrm{HYPS\text{-}IN}(\mathit{hs}, \mathrm{DERIVABLE}(T)) \Rightarrow
c \in \mathrm{DERIVABLE}(T)\bigr).
\end{aligned}
\]
\end{claim}

$\mathrm{HYPS\text{-}IN}(\mathit{hs}, S)$ says that $\mathit{hs}$ is a
list, a set, every entry of which is in $S$ (\DS{47}).

For a rule that is itself a core rule, the claim is its closure law.
The content lies in the rules that are not: there the proof is a
derivation template, built by recursion on the syntax.  What the
theorems establish is that trusting the kernel's relations is no more
than trusting the core.  They do not establish that the core is sound
(link~3 of Section~\ref{sec:meta-chain}), nor that the Scheme of a
checker computes its relation (link~2).

\paragraph{Four kinds of proof.}
\begin{enumerate}
\item \emph{Derivation templates.}  A rule outside the core is derived
  by a template (\texttt{adm-detach}, \texttt{adm-union-intro},
  \texttt{adm-tuples-intro}, \ldots).
\item \emph{Contextual rewriting.}  The rules that rewrite at an
  arbitrary position (\texttt{macete}, \texttt{macete-hyp},
  \texttt{lambda-beta}, \texttt{lambda-beta-hyp}, \texttt{nth-reduce},
  \texttt{length-reduce}, \texttt{functoid-beta}, and the two
  decompositions) are reduced to a theorem CR: a rewrite valid in the
  local context of a position is valid at that position.  CR is proven by
  induction on positions, by stages: stages 1--5 cover the connectives,
  atoms, quantifiers, \texttt{if}, the arguments of a term and the domain
  of a binder; stage~6 the bodies of \texttt{sep}, \texttt{big-union} and
  \texttt{vnb-lambda}, through tbind-cr; stage~7 the bodies of
  \texttt{comp} and \texttt{iota}.  A rewrite at several positions is a
  chain of one-position rewrites, and its admissibility follows from CR
  by induction on the length of the chain
  (Section~\ref{sec:meta-positions}).
\item \emph{Definedness rows.}  Each clause of the certificate
  \texttt{pi--defined?} is transcribed as a relation
  $\mathrm{DEF\text{-}}x$ and derived over the core or admitted as a
  schema (Section~\ref{sec:meta-lemma-star}).
\item \emph{Inductions.}  The three induction rules are derived from the
  class-form induction axioms of the catalog and comp-mem, for a
  predicative body.
\end{enumerate}

%% =========================================================
\section{Coverage}\label{sec:meta-coverage}

Each of the 57 non-oracle rules has one of the following forms.

\begin{center}
\small
\begin{tabular}{@{}lp{\dimexpr0.88\linewidth-2\tabcolsep\relax}@{}}
\toprule
form & meaning \\
\midrule
C & The kernel figure is a core figure.  The theorem is the closure law. \\
B & The kernel figure differs from a core figure by a small step, the
  \emph{bridge}: the kernel replaces a formula the core keeps, or omits a
  premise the core states.  The kernel relation is transcribed and its
  admissibility proven (\texttt{adm-}$R$\texttt{-k}). \\
B* & As B, when the premise the kernel omits is the definedness
  $\Gamma \Longrightarrow t = t$ of a term the kernel certifies.  The
  admissibility is proven for a term certified by a member of the context
  (\texttt{adm-}$R$\texttt{-k}) and for a term certified by any
  composition of the definedness rows of lemma~(*)
  (\texttt{adm-}$R$\texttt{-kd}; Section~\ref{sec:meta-lemma-star}).
  A certificate that uses an open row is not covered. \\
D & Not a core rule.  Transcribed and derived in full over the core
  (\texttt{adm-}$R$). \\
PS & A rewriting rule, admissible for a rewrite at any finite set of
  positions, at any stage of CR (\texttt{adm-}$R$\texttt{-s};
  Section~\ref{sec:meta-positions}). \\
\bottomrule
\end{tabular}
\end{center}

\begin{center}
\small
\begin{tabular}{@{}lrp{\dimexpr0.82\linewidth-4\tabcolsep\relax}@{}}
\toprule
form & count & rules \\
\midrule
C & 19 & and-intro, or-intro-left, or-intro-right, implies-intro,
  not-intro, iff-intro, forall-intro, not-elim, assumption,
  theorem-assumption, cut, weakening, backchain, proof-by-contradiction,
  eq-subst, quasi-reflexivity, sep-sethood, big-union-sethood, iota-def \\
B & 5 & and-elim, or-elim, iff-elim, forsome-elim, iota-in-elim \\
B* & 5 & forsome-intro, forall-elim, reflexivity, if-true, if-false \\
D & 19 & detach, cartesian-intro, cartesian-elim, tuples-intro,
  tuples-elim, union-intro, union-elim, intersection-intro,
  intersection-elim, nn-induction, transfinite-induction,
  transfinite-induction-3cases, sep-mem-intro, sep-mem-elim,
  comp-mem-intro, comp-mem-elim, big-union-mem-intro, big-union-mem-elim,
  lambda-type \\
PS & 9 & nth-reduce, length-reduce, functoid-beta, lambda-beta,
  lambda-beta-hyp, macete, macete-hyp, cartesian-decompose,
  tuple-equality-decompose \\
\midrule
all & 57 & \\
\bottomrule
\end{tabular}
\end{center}

So Claim~\ref{claim:meta-k2} is proven in full for the C, B, D and PS
rows, and for the B* rows for every certificate built from the rows of
lemma~(*).  One qualification applies to the D rows tuples-intro,
union-intro and lambda-type, which owe the definedness of a class
(Section~\ref{sec:meta-owed}).  The D theorem is for the figure with
the owed premise.  Where the kernel omits the premise because its
certificate accepts the class, the figure is covered at the empty tuple
and for lambda-type at one binder, for a class certified by a member of
the context; it is open for union-intro and for lambda-type at a binder
list.

The ledger \texttt{reference/K2-LEDGER.md} lists, for each of the 57
rules, its form, the relation, the theorems that discharge it and what
is open.

%% =========================================================
\section{The bridged rules}\label{sec:meta-bridges}

For each B and B* rule: the kernel figure, the core figure, the bridge,
and the theorem.  The kernel relations are transcribed in
\texttt{structure-library/meta-nbg.scm}, the theorems proven in section~B
of \texttt{theorem-library/meta-nbg-laws.scm}.  $\Gamma'$ is the context
with the eliminated formula removed.

\subsection{The kernel replaces, the core keeps}

The kernel's elimination rules remove the formula they eliminate from
the context; the core's keep it.  The bridge is one weakening.

\VNBfig{and-elim (kernel)}{\Gamma', A, B \Longrightarrow G}{\Gamma \Longrightarrow G}
\VNBfig{and-elim (core)}{\Gamma, A, B \Longrightarrow G}{\Gamma \Longrightarrow G}
with $A \wedge B \in \Gamma$ and $\Gamma' = \Gamma \setminus \{A\wedge B\}$.
Bridge: $\Gamma', A, B \subseteq \Gamma, A, B$, so weakening takes the
kernel's hypothesis to the core's.  Theorem \texttt{adm-and-elim-k}.

\VNBfig{or-elim (kernel)}{\Gamma', A \Longrightarrow G \qquad \Gamma', B \Longrightarrow G}{\Gamma \Longrightarrow G}
with $A \vee B \in \Gamma$.  Core: the same with $\Gamma$ for $\Gamma'$.
Bridge: weakening in each branch.  Theorem \texttt{adm-or-elim-k}.

\VNBfig{iff-elim (kernel)}{\Gamma', A \to B, B \to A \Longrightarrow G}{\Gamma \Longrightarrow G}
with $A \leftrightarrow B \in \Gamma$.  Bridge: weakening.  Theorem
\texttt{adm-iff-elim-k}.

\VNBfig{forsome-elim (kernel)}{\Gamma', B[x:=y] \Longrightarrow G}{\Gamma \Longrightarrow G}
with $\exists x.\,B \in \Gamma$ and $y \notin \FV(\Gamma') \cup \FV(G)
\cup \FV(\exists x.\,B)$.  The core asks $y \notin \FV(\Gamma) \cup
\FV(G)$.  Bridge: a free name of a member of $\Gamma$ is free in a member
of $\Gamma'$ or in $\exists x.\,B$ (\texttt{ctx-fv-minus}), so the
kernel's eigenvariable condition gives the core's; then weakening.
Theorem \texttt{adm-forsome-elim-k}.

\subsection{The kernel reads a definedness off the context}

\VNBfig{iota-in-elim (kernel)}{\Gamma, B[x:=\iota x.B] \Longrightarrow G}{\Gamma \Longrightarrow G}
with $\iota x.B \in A$, $\iota x.B = u$ or $u = \iota x.B$ in $\Gamma$.
\VNBfig{iota-in-elim (core)}{\Gamma \Longrightarrow \iota x.B = \iota x.B \qquad \Gamma, B[x:=\iota x.B] \Longrightarrow G}{\Gamma \Longrightarrow G}
Bridge: the core's first hypothesis is derived from the context member,
by assumption and the strictness of $\in$ (def-strict-in-l) or of $=$
(\texttt{adm-def-strict-eq-l}, \texttt{-r}).  Theorem
\texttt{adm-iota-in-elim-k}.

\subsection{The kernel omits a definedness premise it certifies}

For the five B* rules the kernel omits the premise
$\Gamma \Longrightarrow t = t$ when its certificate accepts $t$ in
$\Gamma$.  When the certificate declines, the kernel posts that premise
and its figure is the core's: that case is C.

Each B* rule has two transcribed kernel forms.  In the form
$\mathrm{REL\text{-}}R\mathrm{\text{-}K}$ the term is certified by one
member of the context; in the form $\mathrm{REL\text{-}}R\mathrm{\text{-}KD}$
the sequent $\Gamma \Longrightarrow t = t$ is a member of
$\mathrm{DEF\text{-}DERIVABLE}(T)$, the sequents the definedness rows
reach (Section~\ref{sec:meta-lemma-star}).  The theorems are
\texttt{adm-}$R$\texttt{-k} and \texttt{adm-}$R$\texttt{-kd}.

\VNBfig{forall-elim (kernel, $t$ certified)}{\Gamma, B[x:=t] \Longrightarrow G}{\Gamma \Longrightarrow G}
\VNBfig{forall-elim (core)}{\Gamma, B[x:=t] \Longrightarrow G \qquad \Gamma \Longrightarrow t = t}{\Gamma \Longrightarrow G}
with $\forall x.\,B \in \Gamma$.  Bridge: $\vdash_T \Gamma \Longrightarrow
t=t$ whenever $\Gamma$ certifies $t$; this is lemma~(*).  Theorems
\texttt{adm-forall-elim-k} and \texttt{adm-forall-elim-kd}.

\VNBfig{forsome-intro (kernel, $t$ certified)}{\Gamma \Longrightarrow B[x:=t]}{\Gamma \Longrightarrow \exists x.\,B}
\VNBfig{forsome-intro (core)}{\Gamma \Longrightarrow B[x:=t] \qquad \Gamma \Longrightarrow t = t}{\Gamma \Longrightarrow \exists x.\,B}
The kernel accepts $t$ certified by the instance itself (a strict atom
among the sentences $B[x:=t]$ makes true) or by the context.  Bridge:
from the instance, by cut with the walks \texttt{adm-true-walk} and
\texttt{adm-strict-atom}; from the context or the rows, lemma~(*).  Theorems
\texttt{adm-forsome-intro-k} and \texttt{adm-forsome-intro-kd}.

\VNBfig{reflexivity (kernel)}{\VNBnohyp}{\Gamma \Longrightarrow t = t}
with $\Gamma$ certifying $t$.  The core's reflexivity fires only for a
variable or an atomic constant.  Bridge: lemma~(*).  Theorems
\texttt{adm-reflexivity-k} and \texttt{adm-reflexivity-kd}.  The
reflexivity checker also accepts $\Gamma \Longrightarrow \textsc{truth}$
with no hypothesis; that figure is the core's truth-intro (theorem
\texttt{adm-reflexivity-kt}).  The kernel has no other operation that
concludes \textsc{truth}.

\VNBfig{if-true (kernel, $a$ certified)}{\Gamma \Longrightarrow p \qquad \Gamma, \mathrm{IF}(p,a,b) = a \Longrightarrow G}{\Gamma \Longrightarrow G}
\VNBfig{if-true (core)}{\Gamma \Longrightarrow p \qquad \Gamma, \mathrm{IF}(p,a,b) \equiv a \Longrightarrow G}{\Gamma \Longrightarrow G}
Bridge: from $\Gamma \Longrightarrow a = a$ (lemma~(*)) and the
quasi-equation the core hands over, the strict equation (a
strictness law of $=$), then a cut.  The kernel hands over the strict
equation only when $a$ is certified, otherwise the quasi-equation
(\DS{20}).  Theorems \texttt{adm-if-true-k} and
\texttt{adm-if-true-kd}; if-false likewise with $\neg p$ and $b$,
\texttt{adm-if-false-k} and \texttt{adm-if-false-kd}.

\subsection{A hypothesis that adds nothing}

The checkers of forall-elim, forsome-elim, if-true, if-false and
iota-in-elim also accept a hypothesis $\Delta \Longrightarrow G$ with the
conclusion's goal and $\Delta \subseteq \Gamma$: the formula the rule
would add is already in the context.  That figure is one weakening
(theorem \texttt{adm-adds-nothing-k}).

When the bound variable does not occur in the body, forsome-elim imposes
no eigenvariable condition.  That instance is an instance of the kernel
form for any fresh $y$, since opening a vacuous body does not change it;
the vacuous forall-elim is the kernel form at a fresh variable $y$, which
the rows certify by reflexivity.

%% =========================================================
\section{Rules that owe the definedness of a class}\label{sec:meta-owed}

Three kernel rules conclude a membership, or a $\mathrm{FUN}$ statement,
that makes a class defined which their other hypotheses do not
(\DS{53}).  Each owes the definedness of that class unless its
certificate accepts the class in $\Gamma$.
\begin{itemize}
\item \texttt{tuples-intro} at the empty tuple owes
  $\Gamma \Longrightarrow A = A$.  With $n \geq 1$ components the
  hypothesis $\Gamma \Longrightarrow a_1 \in A$ makes $A$ defined, and
  nothing is owed.
\item \texttt{union-intro} $k$ owes $\Gamma \Longrightarrow S_j = S_j$ for
  each entry $S_j$ other than $S_k$; the hypothesis
  $\Gamma \Longrightarrow x \in S_k$ makes $S_k$ defined.
\item \texttt{lambda-type} owes $\Gamma \Longrightarrow B = B$ for the
  codomain $B$.  Over an empty domain the typing hypothesis holds for
  every $B$.
\end{itemize}
The owed sequents are hypotheses of the inference, and the checker
accepts each omitted one only when \texttt{chk-elim-certified?} accepts
the class.  The core schemas tuples-mem, union-mem and lambda-type carry
the same premises (Section~\ref{sec:meta-core}).

\paragraph{The relations.}
$\mathrm{REL\text{-}TUPLES\text{-}INTRO}$ is the figure with $n \geq 1$
components, and $\mathrm{REL\text{-}TUPLES\text{-}INTRO\text{-}0}$ the
empty tuple with the owed premise (\texttt{meta-core.scm}).
$\mathrm{REL\text{-}UNION\text{-}INTRO}$ is the owed form: the
membership, then one definedness sequent for each entry.
$\mathrm{REL\text{-}LAMBDA\text{-}TYPE\text{-}1}$ to
$\text{-}4$ and $\mathrm{REL\text{-}LAMBDA\text{-}TYPE\text{-}N}$
(\texttt{meta-ind.scm}) carry $\Gamma \Longrightarrow B = B$ as their last
hypothesis.  The certified forms
$\mathrm{REL\text{-}TUPLES\text{-}INTRO\text{-}0K}$ and
$\mathrm{REL\text{-}LAMBDA\text{-}TYPE\text{-}K}$ (one binder;
\texttt{meta-nbg.scm}) drop the premise for a class certified by a
member of the context, as the -K forms of the B* rules do.

\paragraph{The theorems.}  \texttt{adm-tuples-intro},
\texttt{adm-tuples-intro-0} and \texttt{adm-tuples-intro-0k};
\texttt{adm-union-intro}; \texttt{adm-lambda-type-1} to \texttt{-4},
\texttt{adm-lambda-type-n} and \texttt{adm-lambda-type-k}.  The closure
law used by \texttt{adm-union-intro}, \texttt{derivable-union-mem-c},
assumes the definedness of every entry derivable.
\texttt{derivable-union-mem-u} obtains those premises from the
definedness of the union itself, by the strictness of $\mathrm{UNION}$
in its arguments.  The decomposition lemmas \texttt{der-union-decompose}
and \texttt{der-union-binary} carry the same assumption.

\paragraph{The rewriting rules.}  The local instances of the union
decomposition, $\mathrm{J\text{-}UDEC}$, require each entry's definedness
derivable in the local context.  Those of lambda-beta,
$\mathrm{J\text{-}LB}$, require the value's sethood derivable there: the
kernel's \texttt{lambda-beta} owes $\Gamma \Longrightarrow b[x:=u] \in
\mathrm{SET}$ when the value may be a proper class, and the core's
lambda-beta-eq has it as a premise.

\paragraph{What is open.}  The kernel omits an owed sequent that its
certificate accepts.  For \texttt{union-intro} at any entry, and for
\texttt{lambda-type} at a binder list, that omission is not
transcribed.  The -K forms cover a class certified by one member of the
context; no form covers a class certified through a composition of the
rows of lemma~(*).

\subsection{lambda-type at a binder list}

The core has one-binder lambdas only.  A kernel term
$\mathrm{vnb\text{-}lambda}([x_1..x_n], A_1 \times \cdots \times A_n, b)$
is encoded as a one-binder lambda over the product,
$\mathrm{ENC\text{-}LAM}(\mathrm{ENC\text{-}CART}(a), c)$, where $c$ is $b$
with each $x_i$ replaced by the projection $\mathrm{NTH}(i, \cdot)$ of the
bound variable.  The relation fixes $c$ by the one property the rule
needs: opened at a name $y$, $c$ is $b$ opened at the projections of $y$,
in the order in which instantiating the $n$ guarded universals opens
them.

\texttt{adm-lambda-type-n} proves the kernel figure admissible at every
length $n \geq 1$; the fixed lengths 1 to 4 are proven separately in
literal form.  For a fresh $z$ in the product, the chain of guarded
universals is instantiated at the projections of $z$ by induction on the
number of guards left (\texttt{ltch-inst}).  This gives $c$ opened at $z$
in $B$.  Then implies-intro, forall-intro at $z$, and the core
lambda-type.

The theorem is relative to the encoding of a destructuring lambda,
which is the translator's (K3).  Two of its properties are argued in
the file, not proven: every member of a product is the list of its
projections (the core's tuple-eta says it), and $c$ is predicative
exactly when $b$ is.

%% =========================================================
\section{Rewriting at several positions}\label{sec:meta-positions}

The kernel's rewriting rules rewrite at every position their walk
explains, in one inference.  The META relations read such an inference
as a chain of one-position rewrites.
$\mathrm{REWRS}(J)$ is the least set of rewrite tuples that contains
each one-position rewrite justified by $J$ (at any stage of CR) and the
empty rewrite, and is closed under composition.
$\mathrm{REL\text{-}}R\mathrm{\text{-}S}$ is the rule $R$ over the $J$
set of its one-position form, rewritten by a member of
$\mathrm{REWRS}(J)$ (\texttt{meta-nbg.scm}).

The chain is the kernel's relation when the positions are taken in the
order of the kernel's walk.  The macete walk stops at the first
explained position, so the positions are disjoint.  The local contexts
are read off different sides at a connective: at $A \wedge B$ the right
conjunct is rewritten under the original left, the left under the
rewritten right; at $A \vee B$ the left disjunct under the negation of
the original right, the right under the negation of the rewritten left;
at $A \to B$ the consequent under the original antecedent.  Taking the
right conjunct, the left disjunct and the consequent first, the context
the $k$-th step computes is the walk's.  The reduction rules read no
local context: their reachability, created redexes and contractions
inside a contractum included, is a chain of single contractions.

\texttt{adm-rewrite-s} and \texttt{adm-rewrite-hyp-s} prove, by the
least-set induction of $\mathrm{REWRS}$ from the one-position law, that a
rewrite at any finite set of positions is admissible for a goal and for
an assumption.  The rows are \texttt{adm-}$R$\texttt{-s} for nth-reduce,
length-reduce, functoid-beta, lambda-beta and its hypothesis form,
macete and macete-hyp (with their several-variable forms), the union and
intersection decompositions, cartesian-decompose and
tuple-equality-decompose.

Composition adds no restriction, because each step's justification is
judged on the formula the earlier steps produced.  What stays outside
is what the one-position forms exclude: the licence of lambda-beta
through its binder's own scope, the binder-domain reads of the list
definedness checks, and any justification outside the rule's $J$ set.

%% =========================================================
\section{The definedness lemma (*)}\label{sec:meta-lemma-star}

\begin{lem}[(*)]
If the certificate \texttt{pi--defined?} accepts the term $t$ in the
context $\Gamma$, and every clause it uses is a row of
$\mathrm{DEF\text{-}RULE}$, then $\vdash_T \Gamma \Longrightarrow t = t$.
\end{lem}

The lemma is the bridge of every B* rule.  Each clause of the
certificate is transcribed as a relation $\mathrm{DEF\text{-}}x$
(\texttt{structure-library/meta-def.scm}), and each relation is either a
core schema (its closure law) or derived over the core.  A clause that
certifies a compound term from certified parts is the corresponding
relation applied to the parts.

$\mathrm{DEF\text{-}RULE}(T, \mathit{hs}, c)$ is the disjunction of the
derived and admitted rows of the table below, together with the three
core rules that start a reading: assumption (a member of the context),
theorem-assumption (a member of $T$, such as a definition a predicate
row reads) and reflexivity (an atom).  The rows $\mathrm{DEF\text{-}APP}$
and $\mathrm{DEF\text{-}LEN\text{-}TUPLES}$ enter only when $T$ is a set of sentences
holding the encoded axiom the row rests on.  $\mathrm{DEF\text{-}DERIVABLE}(T)$ is the
least set of sequents closed under $\mathrm{DEF\text{-}RULE}(T,\cdot,\cdot)$,
defined as $\mathrm{DERIVABLE}(T)$ is.

\begin{thm}[\texttt{def-derivable-sub}]
For every $T$,
\[
\mathrm{DEF\text{-}DERIVABLE}(T) \subseteq \mathrm{DERIVABLE}(T).
\]
\end{thm}

The proof is \texttt{def-closed-derivable}: $\mathrm{DERIVABLE}(T)$ is
closed under every row of $\mathrm{DEF\text{-}RULE}$, each row citing its
closure law.  The theorem is lemma~(*) for every composition of the
rows.  That the kernel's acceptance of $t$, through clauses that are
rows, puts $\Gamma \Longrightarrow t = t$ in
$\mathrm{DEF\text{-}DERIVABLE}(T)$ is a faithfulness claim, part of K3
(\DS{9}).

{\small
\begin{longtable}{@{}p{\dimexpr0.4\linewidth-2\tabcolsep\relax}p{\dimexpr0.28\linewidth-2\tabcolsep\relax}p{0.32\linewidth}@{}}
\toprule
clause & relation & status \\
\midrule
\endhead
atom (variable, constant) & reflexivity & core \\
$t \in \_$ or $\_ \in t$ in context & def-strict-in-l, -r & core \\
$t = \_$, $\_ = t$ in context & DEF-STRICT-EQ-L, -R & derived \\
$t$ outside binders in a true $\in$, $=$, $\leq$, $<$, or in a
  separation condition & DEF-STRICT-HYP & derived \\
structure hypothesis $\mathrm{IS\text{-}X}(t)$, $t \in X$ & DEF-STRUCT-HYP & derived \\
predicate hypothesis (any arity, leading existentials, nested, parent
  chain); accessor under $\mathrm{IS\text{-}X}(s)$ & DEF-PRED-HYP-N, DEF-CERT-HYP & derived \\
literal list of sets & def-list & core \\
total class constructor on certified arguments & def-total & core \\
\texttt{sep}, \texttt{big-union}, \texttt{vnb-lambda} over a certified
  domain, predicative body & def-binder & core \\
\texttt{comp}, predicative condition & def-comp & core \\
$\mathrm{NTH}(k, L)$, $L$ certified & DEF-NTH-LIST & derived \\
functoid with a certified instance; NTH or SEP through a functoid
  & DEF-FUN-INST, -NTH, -SEP, -STRICT & derived \\
$f(a)$ with $f \in \mathrm{FUN}(D,\_)$, $a \in D$ & DEF-APP & derived, $T$ holding \texttt{fun-domain-apply-def} \\
$\mathrm{LENGTH}(t)$ with $t \in \mathrm{TUPLES}(A)$ in context & DEF-LEN-TUPLES & derived, $T$ holding \texttt{length-in-nn} \\
$\mathrm{LENGTH}(t)$, $t$ typed in a class of tuples through a
  \texttt{sep} domain or a functoid body & -- & open \\
ground arithmetic; arithmetic on certified numbers & -- & open \\
applied structure operation on typed arguments & -- & open \\
view clause $\mathit{ACC}(V(r))$, the whole list certified (\DS{28}) & -- & no row \\
\bottomrule
\end{longtable}
}

The derived rows are proven in the laws files of
\texttt{meta-def.scm} and \texttt{meta-ind.scm}; the theorem for each is
in the ledger.  The open rows have no theorem.  Ground arithmetic needs
head codes for the arithmetic heads and their closure axioms in $T$; the
applied structure operation needs an encoding of an application to
several arguments and the slot's domain term; \texttt{LENGTH} through a
\texttt{sep} domain or a functoid body needs the certificate's walk to a
class of tuples transcribed.  All three are part of K3.  The view clause
is not admissible as a row.  The retired clause for $\mathrm{NTH}$ of a
literal list is not a row.  A certificate that uses an open clause is not
covered by the B* theorems (Chapter~\ref{ch:meta-limits}).

%% =========================================================
\section{Regression}\label{sec:meta-regression}

The META theorems are checked as every theorem is: by the certified load
by day and by the exam at night.  Their certificates follow the
definitions of the statement transitively (\DS{17}), so a change to a
core relation re-proves every META theorem that rests on it.

Each kernel repair comes with a reproducer and suite controls, at least
one of which fails on the code before the repair.  Each core schema with
a precondition has a control for it; each congruence and beta rule has a
control that it concludes only the quasi-equation.


%% =========================================================
%% ch-meta-soundness.tex -- Part 2 (\ref{part:meta}), the soundness of the core for version-X (link 3).
%% Hand-written (DS-6).  Cites the design specification as \DS{n}.  Placed after ch-meta-k2.
\providecommand{\DS}[1]{\hyperref[ds:#1]{DS-#1}}
\providecommand{\Sq}{\Rightarrow}
\providecommand{\FV}{\mathrm{FV}}

\chapter{Soundness of the core}
\label{ch:meta-soundness}

This chapter is link~3 of Section~\ref{sec:meta-chain}.  It argues
that each of the 56 core rules of Section~\ref{sec:meta-core} is sound
for version-X, and that each axiom of the catalog holds, under one
reading of VNB's terms and formulas.

The argument is written by hand (\DS{6}).  It is not a VNB theorem
(Section~\ref{sec:meta-not-claimed}).  The functoid part is
Chapter~\ref{ch:meta-functoids}; here every formula is read after the
elimination of its defined heads.

%% =========================================================
\section{The reading}\label{sec:meta-reading}

\subsection{Models}

Let $M$ be a model of version-X, NBG with global choice
(\DS{3})~\cite{vonNeumann1925,Bernays1937,Godel1940}.  The objects of
$M$ are its classes; a set is a class that is a member of a class.  Fix
in $M$ a well-ordering $W$ of the universe.  NBG with global choice
provides one.

A VNB class is read as a class of $M$ (\DS{4}).  Free names
(variables) range over all classes of $M$.  An \emph{assignment}
$\sigma$ gives each free name a class.

\subsection{Terms}

A term denotes a class of $M$ or nothing.  We write $[\![t]\!]\sigma$
for the denotation, and say that $t$ \emph{denotes} when there is one.
The clauses:
\begin{itemize}
\item A free name denotes $\sigma$ of it.  A primitive constant denotes
  a fixed class: $\mathrm{SET}$ the class of all sets, $\mathrm{ORD}$
  the class of von Neumann ordinals, $\mathrm{NN}$ the von Neumann
  naturals, $\mathrm{EMPTY\text{-}SET}$ the empty set, the numerals and
  the number classes as in Section~\ref{sec:meta-numbers}.
\item \emph{Lists.}  The list $[e_1..e_n]$ denotes when every $e_i$
  denotes a set (\DS{24}).  Its value is the Kuratowski pair
  $(n, s)$, where $s$ is the set of Kuratowski pairs $(i, e_i)$,
  $1\leq i\leq n$.  These pairs are NBG's own, not VNB's; a VNB pair
  $[x, y]$ is the list of length 2.  Every list is a set.
\item $\mathrm{LENGTH}(c)$ is $n$ when $c$ is a list of length $n$, and
  $0$ for every other class.  $\mathrm{NTH}(k, c)$ is the $k$-th entry
  when $c$ is a list and $1\leq k\leq n$, and nothing otherwise.
  $\mathrm{CONS}(x, c)$, for a set $x$, is the list with first entry $x$
  followed by the entries of $c$ when $c$ is a list, and $[x]$
  otherwise; for a proper class $x$ it is nothing.
\item The class constructors $\mathrm{UNION}$, $\mathrm{INTERSECTION}$,
  $\mathrm{CARTESIAN}$, $\mathrm{TUPLES}$ and $\mathrm{FUN}$ denote
  when every argument denotes.  They are then the classes of sets that
  their membership schemas describe.  $\mathrm{CARTESIAN}(A_1..A_n)$ is the class of lists of
  length $n$ whose $i$-th entry is in $A_i$; $\mathrm{TUPLES}(A)$ the
  class of lists all of whose entries are in $A$; $\mathrm{FUN}(A)$ the
  class of sets $f$ of VNB pairs, single-valued, with domain exactly
  $A$; $\mathrm{FUN}(A, B)$ those members of $\mathrm{FUN}(A)$ whose
  values lie in $B$.  $\mathrm{UNION}()$ is the empty set; the nullary
  intersection and product are fixed arbitrarily, since no rule or
  axiom mentions them.
\item $\mathrm{succ}$ and $\mathrm{succ\_ORD}$ are total: at a set $x$
  the value is $x\cup\{x\}$, at a proper class the class itself.
\item \emph{Application.}  $\mathrm{APPLY}(f, x)$ is the unique $v$ with
  $[x, v]\in f$, when $f$ and $x$ denote and exactly one such $v$
  exists, and nothing otherwise.  Such a $v$ is a set.  An application
  of arity $n\geq 2$, $f(a_1,\ldots,a_n)$, is
  $\mathrm{APPLY}(f, [a_1..a_n])$.
\item $\mathrm{IF}(p, a, b)$ is the denotation of $a$ when $p$ is true,
  of $b$ otherwise, and nothing when that branch denotes nothing.
\item $\iota x.\,B$ denotes the unique class $C$ such that $B$ holds at
  $x := C$, and nothing when there is no such class or more than one.
\item $\mathrm{CHOICE}(A)$ denotes the $W$-least member of $A$ when $A$
  denotes an inhabited class, and nothing otherwise (\DS{52}).
\item \emph{Binder terms.}  Over a denoting domain $A$:
  $\mathrm{SEP}\,x{\in}A.\,p$ is the class of the members of $A$ at which
  $p$ holds; $\bigcup x{\in}A.\,b$ the class of the members of the
  denoting values $b[x]$, $x\in A$; $\lambda x{\in}A.\,b$ the class of
  the pairs $[x, b[x]]$, $x\in A$, $b[x]$ a set.
  $\mathrm{COMP}\,x.\,p$ is the class of the sets at which $p$ holds.
  Each denotes when that class exists in $M$, and nothing otherwise.
  Over a non-denoting domain the binder term denotes nothing.
\end{itemize}
Each other constructor is read by its defining axioms in the catalog
(Section~\ref{sec:meta-axioms}).  A guarded definition is read as a
description: off every guard its head denotes nothing (\DS{31}).

\subsection{Formulas and sequents}

The atoms $\in$, $=$, $\leq$ and $<$ are strict: an atom with a
non-denoting argument is false (\DS{18}).  So $t = t$ says that $t$
denotes.  $s\equiv t$ holds when neither side denotes, or both denote the
same class.  A predicate application $P(t_1..t_n)$ is false when some
$t_i$ does not denote, and otherwise holds when the values satisfy $P$.
Primitive predicates are classes of $M$; a defined predicate is read
through its definition.

The connectives are classical.  The quantifiers range over all classes
of $M$.  A formula containing defined heads means its elimination
(\DS{7}; Chapter~\ref{ch:meta-functoids}).

A sequent $\Gamma\Sq A$ is \emph{valid in $M$} when every assignment
that makes each member of $\Gamma$ true makes $A$ true.  A theory $T$ is
\emph{true in $M$} when each of its members is true at every
assignment.

\begin{dfn}\label{dfn:meta-sound}
A core rule $\mathrm{CR}\text{-}x$ is \emph{sound} when, for every model
$M$ of version-X and every theory $T$ true in $M$, each instance of
$\mathrm{CR}\text{-}x(T, \mathit{hs}, c)$ whose hypotheses are valid in
$M$ has a conclusion valid in $M$.
\end{dfn}

A definedness side sequent $\Gamma\Sq t=t$ is an ordinary hypothesis
under this definition.  Nothing else is needed for it.

\subsection{Two lemmas}

\paragraph{Substitution.}  Let $t$ be a closed term that denotes at
$\sigma$.  Then $B[x:=t]$ is true at $\sigma$ exactly when $B$ is true
at $\sigma$ with $x$ given the value of $t$.  The same holds for the
denotation of a term.  If $s\equiv t$ holds at $\sigma$, then $P[z:=s]$
and $P[z:=t]$ have the same truth value at $\sigma$.  Both statements
are proven by induction on the expression.  They hold because each clause
above reads a subterm only through its denotation or the absence of one.
The encoding is locally nameless (\DS{10}), so no capture arises.

\paragraph{Elimination.}  Let $p$ satisfy $\mathrm{NBG\text{-}OK}$: every
quantifier and description in it is confined to sets, read through the
tables $\mathrm{FUN\text{-}DEF}$ and $\mathrm{PRED\text{-}DEF}$ (\DS{39}).
Then $p$ is equivalent in $M$ to a formula of NBG all of whose
quantifiers range over sets, with class parameters: the values of its
free names, the primitive classes and $W$.

The proof replaces each atom about a compound term by its defining
condition.  $y\in\mathrm{SEP}\,x{\in}A.\,q$ becomes
$y\in A\wedge q[x:=y]$.  $y=\mathrm{APPLY}(f,x)$ becomes
$[x,y]\in f\wedge\forall y'\,([x,y']\in f\to y'=y)$, where $y'$ is
confined by the membership.  $y=\mathrm{CHOICE}(A)$ becomes
$y\in A\wedge\forall z\,(z\in A\to\neg\,zWy)$.  A list becomes its
Kuratowski code.  A confined description becomes the confined unique
existence it abbreviates.  Hence the class $\{x\in\mathrm{SET} : p\}$
exists in $M$, by the class comprehension of NBG.

\begin{prop}[hand-written, conditional]\label{prop:meta-sound}
Assume the hypotheses of Section~\ref{sec:meta-sound-hyp}.  Every core
rule is sound.  Consequently, if $T$ is true in a model of version-X,
every sequent in $\mathrm{DERIVABLE}(T)$ is valid in it, and
$\{\}\Sq\textsc{falsity}$ is not derivable.
\end{prop}

The consequence is an induction on $\mathrm{DERIVABLE}(T)$, the least
set closed under the rules.  $\{\}\Sq\textsc{falsity}$ is valid in no
model.  If version-X is consistent it has a model, so $T$, when true in
one, derives no \textsc{falsity}.  Section~\ref{sec:meta-axioms} argues
that the catalog is true in every model of version-X, with the primitive
heads read as stated there.  A theorem enters $T$ only after it is derived
(links~1 and~2), so by induction on the order of installation it is true
as well.

%% =========================================================
\section{The rules}\label{sec:meta-rules}

The groups follow the table of Section~\ref{sec:meta-core}.  In each
argument, $\sigma$ is an assignment that makes $\Gamma$ true, and a
hypothesis is used through its validity at $\sigma$.

\subsection{Logic (21 rules)}

and-intro, and-elim, or-intro-left, or-intro-right, or-elim,
implies-intro, backchain, not-intro, not-elim, iff-intro, iff-elim,
assumption, cut, weakening and proof-by-contradiction are the
propositional rules of classical natural deduction in the sequent
form~\cite{Gentzen1935}.  Each is sound by the truth tables.  The
elimination rules read a formula of the context; that changes nothing in
the argument.  truth-intro concludes \textsc{truth},
true everywhere.  theorem-assumption concludes a member of $T$, true by
the assumption on $T$.

The four quantifier rules are those of a partial first-order logic in
which variables always denote and a term may not, as in the logic of
IMPS~\cite{IMPS} and in~\cite{Farmer2023}.
\begin{itemize}
\item forall-elim and forsome-intro carry the hypothesis $\Gamma\Sq t=t$
  (\DS{19}, \DS{40}).  At $\sigma$, $t$ denotes a class $C$.  By the
  substitution lemma, $B[x:=t]$ is $B$ at $x:=C$.  The universal gives
  it; the existential is witnessed by it.  Without the hypothesis
  neither rule is sound: $\forall x.\,x=x$ is true, and
  $\iota x.\,\textsc{falsity} = \iota x.\,\textsc{falsity}$ is false.
\item forall-intro: $y\notin\FV(\Gamma)\cup\FV(B)$.  For a class $C$,
  the assignment $\sigma[y:=C]$ still makes $\Gamma$ true, so $B[x:=y]$
  holds there, which is $B$ at $x:=C$.  forsome-elim is the dual
  argument, with $y\notin\FV(G)$.
\end{itemize}

\subsection{Equality (4 rules)}

reflexivity concludes $t=t$ for $t$ a free name or an atomic constant.
A free name denotes under every assignment.  An atomic constant denotes
by hypothesis~(S2).  quasi-reflexivity holds by the definition of
$\equiv$.

eq-subst and eq-subst-sym need one of $s=t$, $s\equiv t$, $t=s$,
$t\equiv s$ in $\Gamma$ (\DS{41}).  Each makes $s\equiv t$ true at
$\sigma$.  The substitution lemma then gives $P[z:=s]$ from $P[z:=t]$.
The rule is sound for $\equiv$, not only for $=$, because the reading
is compositional in the denotation or its absence.

\subsection{Structural membership and tuples (7 rules)}

Each schema is the membership condition of the reading.  The work is at
the arguments that may not denote.
\begin{itemize}
\item union-mem, with one hypothesis $\Gamma\Sq S_i=S_i$ per entry
  (\DS{53}).  All entries denote, so $\bigcup(S_1..S_n)$ is the union of
  their values, and the equivalence is its definition.  If $x$ does not
  denote, both sides are false.  At $n=0$ the union is empty and the
  right side is \textsc{falsity}.  Without the hypotheses the schema is
  false: with $S_2$ non-denoting and $x\in S_1$ true, the right side holds
  and the left does not, since $\mathrm{UNION}$ is strict.
\item intersection-mem, $n\geq 1$.  If some $S_i$ does not denote, both
  sides are false.  Otherwise the equivalence is the definition.  No
  hypothesis is needed, since every entry occurs in a strict atom on the
  right.
\item cartesian-mem, $n\geq1$.  The same argument.  The right side's
  conjuncts $\mathrm{NTH}(i,x)\in A_i$ are false at a non-denoting
  $A_i$, and the product does not denote.
\item tuples-mem, with the hypothesis $\Gamma\Sq A=A$ (\DS{53}).  If some
  $e_i$ is a proper class or does not denote, the list does not denote
  and $e_i\in A$ is false: both sides are false.  Otherwise the list is
  in $\mathrm{TUPLES}(A)$ exactly when each entry is in $A$.  At $n=0$
  the list $[\,]$ denotes and lies in $\mathrm{TUPLES}(A)$ for every
  denoting $A$; the right side is \textsc{truth}.  This is where the
  hypothesis is needed.
\item nth-eq and length-eq carry $\Gamma\Sq L=L$ for the literal
  $L=[e_1..e_n]$ (\DS{25}).  $L$ denotes, so each $e_k$ is a set,
  $\mathrm{NTH}(k,L)$ is $e_k$ for $1\leq k\leq n$, and
  $\mathrm{LENGTH}(L)$ is $n$.  The conclusion is the quasi-equation.
\item tuple-eq carries $\Gamma\Sq L=L$ and $\Gamma\Sq M=M$.  Two
  denoting lists of the same length are equal exactly when their entries
  are, by the uniqueness of Kuratowski pairs.  This is also the claim of
  the primitive \texttt{tuple-equality-decompose} (\DS{50}).
\end{itemize}

\subsection{Conditional terms (2 rules)}

The rule if-true has the hypotheses $\Gamma\Sq p$ and
$\Gamma,\ \mathrm{IF}(p,a,b)\equiv a\Sq G$, and the conclusion
$\Gamma\Sq G$.  At $\sigma$, $p$ is true, so $\mathrm{IF}(p,a,b)$
has the denotation of $a$, or none when $a$ has none.  The added
quasi-equation is therefore true at $\sigma$, and the second hypothesis
gives $G$.  if-false is the same with $\neg p$ and $b$.  The core
concludes the quasi-equation, not the strict one: the strict equation
would assert that the branch denotes (\DS{20}).

\subsection{Set formation, descriptions, functions (9 rules)}

Every schema here that forms a class carries $\mathrm{NBG\text{-}OK}$ on
its condition or body.  By the elimination lemma the class it names
exists in $M$.  This is the only place where the strength of NBG, and
not of Morse--Kelley, is at stake (\DS{38}, \DS{39}).
\begin{itemize}
\item sep-mem and comp-mem are the definitions of the two classes.  At a
  non-denoting $t$ both sides are false: $t\in A$, $t\in\mathrm{SET}$
  are strict.
\item sep-sethood: a subclass of a set is a set (separation in NBG).
\item big-union-mem is the definition of the union of the denoting values.
  $\exists z\,(z\in A\wedge x\in b)$ is confined by $z\in A$.
\item big-union-sethood: $A$ is a set and each value $b[z]$, $z\in A$, is
  a set.  By replacement the values form a set; its union is a set.
\item iota-def: the first hypothesis says that exactly one class
  satisfies $B$.  So $\iota x.\,B$ denotes it, and $B[x:=\iota x.\,B]$ is
  true by the substitution lemma.  The second hypothesis gives $G$.
\item iota-in-elim: $\iota x.\,B$ denotes, so it denotes the unique
  class satisfying $B$, and $B[x:=\iota x.\,B]$ holds.
\item lambda-type carries $\Gamma\Sq\forall x.\,(x\in A\to b\in D)$,
  $\Gamma\Sq A\in\mathrm{SET}$ and $\Gamma\Sq D=D$ (\DS{53}).  Let $L$ be
  $\lambda x{\in}A.\,b$.  $L$ is a class, by the elimination lemma.  Each
  value $b[x]$, $x\in A$, is in $D$, so it is a set, and $L$ contains a
  pair for every $x\in A$.  $L$ is single-valued with domain $A$.  It is a
  set by replacement, since $A$ is a set.  $D$ denotes, so
  $\mathrm{FUN}(A,D)$ denotes, and $L$ is a member.  Without $D=D$ the
  instance $A=\mathrm{EMPTY\text{-}SET}$, $D$ non-denoting has a true
  first hypothesis and a false conclusion.
\item lambda-beta-eq carries $\Gamma\Sq t\in A$ and
  $\Gamma\Sq b[x:=t]\in\mathrm{SET}$.  $t$ and $b[t]$ are sets, so
  $[t, b[t]]$ denotes and lies in $L$.  A pair $[x, v]\in L$ with $x=t$
  has $v=b[t]$.  So $\mathrm{APPLY}(L,t)$ is $b[t]$, and the
  quasi-equation holds.  Without the second hypothesis the rule is
  unsound: a pair whose second entry is a proper class does not denote,
  so a lambda term with a class-valued body has no pair at that point,
  and the application denotes nothing while $b[t]$ denotes.
\end{itemize}

\subsection{Definedness (9 rules; \DS{42})}

\begin{itemize}
\item def-strict-in-l, def-strict-in-r: a true membership has denoting
  sides.
\item def-total: $\mathrm{UNION}$, $\mathrm{INTERSECTION}$,
  $\mathrm{CARTESIAN}$, $\mathrm{TUPLES}$, $\mathrm{FUN}$, $\mathrm{succ}$
  and $\mathrm{succ\_ORD}$ denote at denoting arguments.  For the class
  constructors this is comprehension over a predicative condition of the
  arguments.  For the two successors it is their totality.
\item def-binder: $\mathrm{SEP}$, $\bigcup$ and $\lambda$ over a
  denoting domain, with an $\mathrm{NBG\text{-}OK}$ body, denote by the
  elimination lemma.  def-comp: likewise for $\mathrm{COMP}\,x.\,p$.
\item def-list: a list of sets denotes (\DS{24}).
\item def-strict-arg: the heads $\mathrm{UNION}$ through $\mathrm{FUN}$
  (with $\mathrm{NTH}$, $\mathrm{LENGTH}$, $\mathrm{LIST}$ and
  $\mathrm{APPLY}$), $\mathrm{succ}$ and $\mathrm{succ\_ORD}$ denote only
  at denoting arguments.  For $\mathrm{LIST}$ this is the other half of
  \DS{24}.  def-strict-dom: a binder term denotes only over a denoting
  domain.
\item def-strict-pred: a true predicate application has denoting
  arguments, by the reading.  The reading agrees with the definition of
  a defined predicate $P$ only at denoting arguments.  This is enough,
  because every instantiation of a definition owes the definedness of the
  instance (\DS{19}): a definition is never applied at a term that may
  not denote.
\end{itemize}

\subsection{Congruence, functoids, tuples (4 rules)}

\begin{itemize}
\item iota-cr: from $\Gamma\Sq B[x:=y]\leftrightarrow C[x:=y]$, $y$ fresh,
  infer $\iota x.\,B\equiv\iota x.\,C$.  By forall-intro's argument $B$ and
  $C$ hold at the same classes.  So the two descriptions denote the same
  class, or both denote nothing.
\item tbind-cr: for $\mathrm{SEP}$ and $\mathrm{COMP}$ the premise is
  the equivalence of the conditions, for $\bigcup$ and $\lambda$ the
  quasi-equation of the bodies, at a fresh $y$.  The two classes have the
  same members, or neither exists (\DS{43}).  No predicativity premise is
  needed: the conclusion is a quasi-equation, true when neither side
  denotes.
\item functoid-beta-eq: $F(l)\equiv b[l]$ for a row $(F,|l|,b)$ of
  $\mathrm{FUN\text{-}DEF}$.  The meaning of $F(l)$ is that of its
  instantiated body (\DS{7}, \DS{34}), at every argument list, so the
  quasi-equation holds.  It rests on hypothesis~(S3).
\item tuple-eta, $n\geq1$: a member $x$ of $\mathrm{CART}(A_1..A_n)$ is a
  list of length $n$ of sets, and a list is the list of its own entries.
  The right side denotes, being equal to $x$ (\DS{45}).
\end{itemize}

%% =========================================================
\section{The axioms}\label{sec:meta-axioms}

The catalog is the set of axioms in \texttt{reference/THEOREMS.md},
section ``Axioms -- the fixed VNB base''.  Its sources are the base
theory of \texttt{library.scm} and the primitive shelf:
\texttt{theorem-library/axioms.scm}, \texttt{number-systems.scm},
and, in \texttt{structure-library/}, \texttt{ordinals.scm},
\texttt{list-recursion.scm}, \texttt{injection.scm},
\texttt{extended-reals.scm}, \texttt{order-lemmas.scm} and
\texttt{ca-lft-vocab.scm}.  Each
axiom falls into one of the classes below.

\subsection{NBG axioms and their direct consequences}

\begin{itemize}
\item Extensionality: \texttt{class-extensionality},
  \texttt{extensionality}.
\item Sets: \texttt{membership-implies-sethood},
  \texttt{empty-set-is-set}, \texttt{empty-set-has-no-members},
  \texttt{subset-set}.
\item Pairing, power and union: \texttt{pairing},
  \texttt{pairing-membership}, \texttt{power-set},
  \texttt{power-set-membership}, \texttt{union-set-closure}.
\item Separation: \texttt{intersection-set-closure} (one factor a set),
  \texttt{complement-in-set-closure}.
\item Replacement: \texttt{image-set}.  Here $\mathrm{IMAGE}(\varphi, S)$
  is the class of sets $w$ with $\varphi(x)=w$ for some $x\in S$
  (\texttt{image-membership-iff}).  This is the range of the
  single-valued part of $\varphi$ over $S$, a set when $S$ is one.
  $\varphi$ ranges over classes: a bare functoid is never an instance
  (\DS{33}).
\item Infinity: \texttt{nn-is-set}.
\item Global choice: \texttt{choice-axiom}, by the reading of
  $\mathrm{CHOICE}$ (\DS{52}).
\item Theorems of NBG: \texttt{burali-forti}; \texttt{nn-induction} and
  \texttt{transfinite-induction}, at every class $c$.  Every non-empty
  class of ordinals has a least member.
\item Equality.  The axioms \texttt{equality-symmetry},
  \texttt{equality-transitivity} and \texttt{eq-subst-membership} are
  true for strict equality.
\end{itemize}

\subsection{Definitional axioms}

Each of these fixes the meaning of a head under the reading.  It is
therefore true by choice of that meaning.  A membership condition is
predicative, read through defined heads, as the binder-predicativity
gate checks for every definitional or primitive membership equivalence
(\DS{38}).  So the class it describes exists.
\begin{itemize}
\item Application and domain: \texttt{app-graph} (the reading of
  $\mathrm{APPLY}$), \texttt{dom-membership}, \texttt{is-fun-def}, and
  the definitional \texttt{apply-tupling-2} and \texttt{-3}, the
  application of arity $n\geq2$.
\item Functions: \texttt{fun-no-junk}, \texttt{fun-codomain-iff},
  \texttt{fun-domain-apply-def}.  These are the reading of
  $\mathrm{FUN}$.  At $x\notin A$ no pair $[x, v]$ is in $f$, so $f(x)$
  denotes nothing.
\item Boolean operations.  The axioms \texttt{union-membership},
  \texttt{intersection-membership},
  \texttt{complement-in-membership} and \texttt{subset-def}; and the
  schemas \texttt{union-decompose},
  \texttt{intersection-decompose}.  The two schemas are true at every
  instance whose entries denote.  The macete door owes each entry's
  definedness (\DS{19}), and the nullary intersection is refused.
\item Constructions: \texttt{make-set-membership},
  \texttt{make-set-empty}, \texttt{partial-fun-membership},
  \texttt{partial-fun-binary-membership}, \texttt{res-apply},
  \texttt{res-typing}, \texttt{res-codomain}.  $\mathrm{RES}(f,b)$ is the
  set of the pairs of $f$ whose first entry is in $b$.
\item Quasi-equality: \texttt{quasi-eq-reflexivity}, \texttt{-symmetry},
  \texttt{-transitivity}, \texttt{quasi-eq-def},
  \texttt{quasi-eq-subst-membership}.  The last holds because $b\in S$
  makes $b$ denote, and then $a\equiv b$ is $a=b$.
\item Ordinals: $\mathrm{ORD\text{-}LE}$ is inclusion on $\mathrm{ORD}$
  and false elsewhere; $\mathrm{succ\_ORD}(\alpha)=\alpha\cup\{\alpha\}$;
  $\mathrm{ORD\text{-}SEGMENT}(\alpha)=\alpha$;
  $\mathrm{SUP\text{-}ORD}(a)=\bigcup a$.  Under these, \texttt{ord-le-*},
  \texttt{ord-succ-*}, \texttt{ord-segment-*}, \texttt{sup-ord-*},
  \texttt{ord-zero-least}, \texttt{limit-ord-is-sup},
  \texttt{nn-subset-ord}, \texttt{ord-succ-nn} and
  \texttt{ord-le-nn-compat} are facts about von Neumann ordinals.
  \texttt{ord-lt-iff} and \texttt{limit-ord-iff} define two predicates.
\item Structures.  For each declared structure $X$, $\mathrm{IS\text{-}}X$
  is a predicate defined by its shape and laws.  Its last conjunct is
  $s\in\mathrm{SET}$, and a refinement inherits it through its parent.
  The class axiom $\forall s\,(s\in X\leftrightarrow\mathrm{IS\text{-}}X(s))$
  is then comprehension over that condition.  It holds because the
  condition implies sethood and is predicative (the gate above).  The
  morphism and hom-set definitions are definitions of new heads.
\end{itemize}

\subsection{The two sethood axioms}

\texttt{fun-set-iff} and \texttt{cartesian-set-iff} are implications.
Their names are kept:
\begin{itemize}
\item \texttt{fun-set-iff}:
  \[\forall A\,\forall B\,(A\in\mathrm{SET}\wedge B\in\mathrm{SET}\to\mathrm{FUN}(A,B)\in\mathrm{SET});\]
\item \texttt{cartesian-set-iff}:
  \[\forall A\,\forall B\,(A\in\mathrm{SET}\wedge B\in\mathrm{SET}\to\mathrm{CARTESIAN}(A,B)\in\mathrm{SET}).\]
\end{itemize}
\emph{cartesian-set-iff.}  The map $(a,b)\mapsto[a,b]$ is a class
function on the NBG product $A\times B$, a set.  By replacement its
range is a set, and that range is $\mathrm{CARTESIAN}(A,B)$.

\emph{fun-set-iff.}  A member of $\mathrm{FUN}(A,B)$ is a set of pairs
$[x,y]$ with $x\in A$ and $y\in B$, so a subset of
$\mathrm{CARTESIAN}(A,B)$.  Hence $\mathrm{FUN}(A,B)$ is a subclass of
$\mathrm{POWER}(\mathrm{CARTESIAN}(A,B))$, a set.

The converses do not hold in general.  $\mathrm{CARTESIAN}(\mathrm{EMPTY\text{-}SET},\mathrm{ORD})$
is empty, hence a set.  $\mathrm{FUN}(\mathrm{EMPTY\text{-}SET},\mathrm{ORD})$
is $\{\mathrm{EMPTY\text{-}SET}\}$.  $\mathrm{FUN}(\mathrm{ORD},B)$ is
empty for every $B$, since a set function has a set domain.  The
converses that hold are theorems, proven in
\texttt{theorem-library/set-iff-converses.scm}:
\begin{itemize}
\item \texttt{fun-domain-in-set}: an inhabited $\mathrm{FUN}(A,B)$ has a
  set domain $A$;
\item \texttt{fun-set-converse}: if $A$ is an inhabited set and
  $\mathrm{FUN}(A,B)$ is a set, then $B$ is a set;
\item \texttt{cartesian-set-converse-1} and \texttt{-2}: if
  $\mathrm{CARTESIAN}(A,B)$ is a set and one factor is inhabited, the
  other factor is a set.
\end{itemize}

\subsection{Consequences under the reading}

These follow from the reading by a short argument.
\begin{itemize}
\item Functions.  The axioms \texttt{dom-of-fun},
  \texttt{dom-fun-membership}, \texttt{fun-apply-type},
  \texttt{fun-elements-are-sets},
  \texttt{fun-domain-elements-are-sets},
  \texttt{fun-domain-extensionality} and
  \texttt{partial-fun-binary-sethood}.  The last is the union, over the subsets $b$ of $A$, of the sets
  $\mathrm{FUN}(b,C)$.
\item Lists: \texttt{list-sethood}, \texttt{tuples-sethood},
  \texttt{length-in-nn}, \texttt{length-of-empty}, \texttt{nth-in-range},
  \texttt{tuple-members-are-sets}, \texttt{make-set-sethood}.  A list's
  entries are sets.  The lists over a set $A$ form a set: they are coded
  by finite sequences from $A$.
\item List generation, primitive: \texttt{tuple-length-zero} and
  \texttt{tuple-cons-decompose}.  Under the reading every list of length
  $n+1$ is $\mathrm{CONS}$ of its first entry onto the list of the
  others.
\item List construction, definitional: \texttt{empty-in-tuples},
  \texttt{cons-in-tuples}, \texttt{length-cons}, \texttt{nth-cons-1},
  \texttt{nth-cons-succ}, \texttt{makeset-cons}.  The equations are
  unguarded in $L$.  They hold at a non-list $L$ because
  $\mathrm{LENGTH}(L)$ is $0$ and $\mathrm{CONS}(x,L)$ is $[x]$.  At that
  $L$, $\mathrm{NTH}$ has no index in range and $\mathrm{MAKE\text{-}SET}(L)$
  is empty.  A reading in which $\mathrm{LENGTH}$ or $\mathrm{CONS}$
  denotes nothing off the lists would make \texttt{length-cons} and
  \texttt{nth-cons-1} false.
\item \texttt{nn-not-a-list} (\DS{49}).  Every list is a Kuratowski pair
  $\{\{n\},\{n,s\}\}$, and no member of a Kuratowski pair is empty.
  $0$ is empty, and every natural $n\geq1$ has $0$ as a member.
  So no natural is a list.
\item The primitive rule \texttt{cartesian-decompose} is tuple eta with
  the membership condition, and \texttt{tuple-equality-decompose} is
  tuple-eq (\DS{50}).
\end{itemize}

\subsection{The number systems}\label{sec:meta-numbers}

The axioms of \texttt{number-systems.scm}, the order and infinity
primitives and \texttt{recip-defined-nonzero} hold in a standard
construction, recoded so that the inclusions are literal.

Build $\mathbb{C}$ over the Dedekind reals in $M$.  Recode it by an
injection into the sets that is the identity from the standard copy of
$\omega$ onto the von Neumann naturals.  Then
$\mathrm{NN}\subseteq\mathrm{ZZ}\subseteq\mathrm{QQ}\subseteq\mathrm{RR}\subseteq\mathrm{CC}$
holds literally, $0$ is the empty set, and $\mathrm{succ}$ on
$\mathrm{NN}$ is $n+1$.  The arithmetic heads, the order, $|\cdot|$,
$\max$, $\min$, conjugation and the powers are transported.  Then:
\begin{itemize}
\item the field, order and completeness axioms (\texttt{nn-*},
  \texttt{zz-*}, \texttt{qq-*}, \texttt{rr-*}, \texttt{cc-*},
  \texttt{power-*}) hold, \texttt{rr-sup-*} included;
\item $\mathrm{recip}$ is the reciprocal on the non-zero numbers and
  denotes nothing at $0$, which gives \texttt{recip-defined-nonzero};
\item $\mathrm{RR\text{-}STAR}$ is $\mathrm{RR}$ with two fresh sets
  $\mathrm{NEG\text{-}INF}$ and $\mathrm{POS\text{-}INF}$.  $\leq$ is the
  order of the extended line on $\mathrm{RR\text{-}STAR}$ and false
  elsewhere.  This gives \texttt{pos-inf-upper-bound},
  \texttt{neg-inf-lower-bound}, \texttt{rr-star-leq-antisymmetric} and
  \texttt{rr-star-leq-domain}.  The two points must be sets for the
  definitional \texttt{rr-star-membership} to hold;
\item $\mathrm{CC\text{-}INF}$ is a fresh set outside $\mathrm{CC}$
  (\texttt{cc-inf-in-set}, \texttt{cc-inf-not-in-cc});
\item the definitional \texttt{binary-minus-def},
  \texttt{binary-divide-def}, \texttt{rr-abs-def}, \texttt{rr-max-def}
  and \texttt{rr-min-def} define new heads.
\end{itemize}
The recoding does not disturb the reading of lists: no axiom says
whether a number other than a natural is a list.

%% =========================================================
\section{The hypotheses}\label{sec:meta-sound-hyp}

Proposition~\ref{prop:meta-sound} assumes the following.
\begin{enumerate}
\item[(S1)] Version-X is consistent.  The oracles are correct
  (\DS{5}).
\item[(S2)] Every atomic constant code of the encoding names a
  denoting class.  A defined constant whose definition may fail is not
  encoded as an atomic constant.
\item[(S3)] The rows of $\mathrm{FUN\text{-}DEF}$ are the eliminations
  of the kernel's defined heads.  The hypotheses H1--H6 of
  Chapter~\ref{ch:meta-functoids} hold.
\item[(S4)] Every predicate of $T$ without a row of
  $\mathrm{PRED\text{-}DEF}$ has a predicative definition, so the
  elimination lemma reads through it.
\item[(S5)] $T$ consists of the catalog and of installed theorems.  The
  argument covers the axioms of the catalog and the definitions read
  as above.  A guarded definition is read as a description
  (Chapter~\ref{ch:meta-functoids}, \DS{31}).
\end{enumerate}
(S2)--(S4) are obligations of the translator, K3
(Section~\ref{sec:meta-open}).  The definitions of the library are
covered by their general form, not one by one.


%% =========================================================
%% ch-meta-functoids.tex -- Part 2 (\ref{part:meta}), chapter 3: the hand-written argument for functoids.
\providecommand{\DS}[1]{\hyperref[ds:#1]{DS-#1}}

\chapter{Functoids}
\label{ch:meta-functoids}

The encoding of Chapter~\ref{ch:meta-k2} has no functoids of its own
beyond the table $\mathrm{FUN\text{-}DEF}$.  The soundness of the kernel's
functoid apparatus is therefore covered by a hand-written argument
(\DS{6}, \DS{7}).  This chapter states the argument and the hypotheses it
needs, and says how the kernel enforces each.

%% =========================================================
\section{What is eliminated}

Three kinds of head are \emph{defined heads}:
\begin{itemize}
\item a \texttt{def-functoid} head $F$ with parameters $\bar p$ and body
  $b_F$: its unfold rewrites $F(\bar p)$ to $b_F$ with no condition
  (Chapter~\ref{ch:defs});
\item a structure accessor $\mathit{ACC}$ of slot $k$: it rewrites
  $\mathit{ACC}(s)$ to $\mathrm{NTH}(k, s)$;
\item an anonymous record $R$ (a \texttt{lambdoid}) with bindings $\bar x$
  and body $b_R$, applied as \texttt{(apply-functoid $R$ $\bar t$)}.
\end{itemize}
$L_0$ is the part of the language in which none of them occurs.

\begin{dfn}[the elimination map $E$]
On a formula in which every defined head is applied to as many arguments
as it has parameters, $E$ is defined by recursion, arguments first:
\[
\begin{aligned}
E(F(t_1,\ldots,t_n)) &= E(b_F)[\bar p := E(\bar t)],\\
E(\mathit{ACC}(s)) &= \mathrm{NTH}(k, E(s)),\\
E(\texttt{apply-functoid}(R, t_1,\ldots,t_n)) &= E(b_R)[\bar x := E(\bar t)],
\end{aligned}
\]
with simultaneous capture-avoiding substitution; $E$ commutes with every
other constructor, connective and binder.
\end{dfn}

The meaning of a formula with defined heads is the meaning of its
E-image (\DS{7}).  ``$F(\bar t)$ denotes'' means ``$E(F(\bar t))$
denotes''; nothing more is asserted about $F$.

A record's domains are dropped by $E$: \texttt{functoid-beta} does not
read them, so they have no logical content.

\paragraph{Two lemmas.}  \emph{Substitution:}
$E(e[x:=t]) =_\alpha E(e)[x:=E(t)]$.  \emph{Definedness through the
body:} $F(\bar t)$ denotes exactly when $E(b_F)[\bar p := E(\bar t)]$
does.  In particular a defined head is not strict: $F(\bar t)$ may denote
while some $t_i$ does not.  This is the schematic reading of an unfolding
(\DS{34}).

%% =========================================================
\section{The proposition}

\begin{prop}[hand-written, conditional]\label{prop:meta-fun}
Assume version-X, the soundness of the functoid-free rules for $L_0$
(link~3 of Section~\ref{sec:meta-chain}, restricted to $L_0$), the
oracles and the infrastructure, and the hypotheses H1--H6 below.  For
every inference accepted by its checker, either its E-image is the
identity (one hypothesis with the E-image of the conclusion), or its
E-image is an inference between $L_0$ sequents accepted by some kernel
rule.

Consequently every theorem installed by \texttt{qed} over a grounded
graph with an empty bill has a valid E-image; since
$E(\textup{\textsc{falsity}}) = \textup{\textsc{falsity}}$, no such theorem is
\textup{\textsc{falsity}}, provided version-X is consistent.
\end{prop}

The proof goes rule by rule.
\begin{itemize}
\item \emph{An unfold} (\texttt{mac} at $F$, goal or hypothesis) and an
  accessor rewrite owe nothing; under $E$ they are the identity.  Owing
  nothing is correct because the substitution lemma holds at every
  argument, denoting or not.
\item \emph{A functor projection} $\mathit{NAME}@\mathit{ACC}$ owes the
  definedness of the instantiated list (\DS{29}); its E-image is an
  \texttt{nth-reduce} of a literal list, which owes the same (\DS{25}).
\item \emph{The view clause} certifies $\mathit{ACC}(V(r))$ only when the
  whole instantiated list is certified (\DS{28}); its E-image is then
  certified by the NTH clause.
\item \emph{The strictness walks} read a \texttt{def-functoid} application
  through its instance and an applied record through its contraction
  (\DS{21}); every other head they enter is, under $E$, a strict head of
  $L_0$.  So a term the walk reaches is at a strict position of the
  E-image.
\item \emph{The definedness certificate and the value tests} read a
  defined head through its instance (\DS{30}); by the two lemmas each
  verdict is the verdict on the E-image.
\item \emph{Every other rule} treats $F(\bar a)$ as an application with a
  constant head.  Such a head is never substituted, and the matchers match
  a constant head only by identity; so the E-image of an accepted
  inference is an accepted inference of the same rule.
\item \emph{Inside a record body} a record's variable is a placeholder for
  an arbitrary term, denoting or not.  Along the local-context walk of the
  rewriting rules, an assumption mentioning a crossed bound variable is
  dropped, and a term mentioning a record variable is not certified
  (\DS{27}).  A step inside the body is therefore valid after any term is
  substituted for the variable.
\end{itemize}

\paragraph{Guarded definitions.}  Heads defined by \texttt{def-constant},
the two recursion commands and \texttt{def-operator} are not eliminated:
their equations are axioms of $T$.  They are read as descriptions: off
every guard the head denotes nothing (\DS{31}).  This reading satisfies
every defining equation provided no two cases hold at once with
different values.  Under it such a head is strict, and the two readings
-- non-strict for \texttt{def-functoid} heads, strict for defined
functions -- are kept apart by H6.

%% =========================================================
\section{The hypotheses, and how each is enforced}

\begin{center}
\small
\begin{tabular}{@{}lp{\dimexpr0.44\linewidth-2\tabcolsep\relax}p{\dimexpr0.48\linewidth-2\tabcolsep\relax}@{}}
\toprule
& hypothesis & enforcement \\
\midrule
H1 & the parameters of a \texttt{def-functoid} are distinct symbols
  & functoid-hypotheses gate, fatal (\DS{35}) \\
H2 & a body names only heads defined before it (no recursion, no forward
  reference) & functoid-hypotheses gate, fatal \\
H3 & one definition per head; no theorem proven under an earlier
  definition survives a redefinition & functoid-hypotheses gate, fatal; the
  certificate key records the body's hash (\DS{17}) \\
H4 & the free variables of a body are among its parameters
  & functoid-hypotheses gate, fatal \\
H5 & every formula written into the graph is eliminable: a defined head
  occurs only applied, never as a term or as a quantifier instance or
  witness & refused at formula construction, by \texttt{forall-elim}, by
  \texttt{forsome-intro} and by \texttt{reflexivity} (\DS{33}) \\
H6 & no head is both a \texttt{def-functoid} and a library definition
  & functoid-hypotheses gate, fatal \\
\bottomrule
\end{tabular}
\end{center}

H4 carries weight.  With $G(x) := y$, $y$ free, the unfold rewrites
$G(0)$ to $y$ inside $\forall y$, so $\forall y.\,G(0)=y$ would be
provable; instantiated at $0$ and at $1$ it gives $0=1$.  The gate
refuses such a body.

The functoid-hypotheses gate (\texttt{functoid-gate.scm}) runs at the end
of every load, a cold load, a certified build and an extension of the
band alike.  It reads the registry of functoids and the log of
definitions, which records for each \texttt{def-functoid} the heads and
constants defined at that moment.  In a strict load a violation of any of
H1--H4 or H6 is an error, with one line naming the hypotheses violated;
under keep-going it is reported and the load continues.  The library has
no violation.

Two consequences for the library.  A body may name a
constant or a head only if it is defined earlier in the load order
(H2); an instance made by \texttt{declare-instance!} counts from the file
that makes it.  A structure species used as a class in a body, such as
$\mathrm{DIRECTED\text{-}SET}$ in $\mathrm{PTWISE\text{-}NET\text{-}LIMITS}$,
is a class constant through the class axiom of its
\texttt{declare-structure}.  H4 reads every declared structure species
as a constant, not as a free variable of the body.  A species is not a
registered constant: that would make it head-only, and a membership in
it would be refused as a formula.

%% =========================================================
\section{The J-MACN bound}

$\mathrm{J\text{-}MACN}(T)$, the local instances of the macetes of $T$
with any number of schema variables, bounds its argument list by
$\mathrm{SET}$.  The bound holds on the intended domain but is not
derivable from the body alone (\DS{47}).

It excludes no kernel step.  For an actual macete step the encoder takes
the list of encoded substitution values, each a member of
$\mathrm{EXPR\text{-}SET}$, which is a set; a literal list of sets is a
set.  The bound thus holds of every encoded kernel step, provided the
encoder produces a literal list of encoded terms, which is the
translator's obligation (K3).  $\mathrm{HYPS\text{-}IN}$ makes the same
kind of bound on a list of hypothesis sequents.


%% =========================================================
%% ch-meta-limits.tex -- Part 2 (\ref{part:meta}), chapter 4: what is not claimed, and what is open.
\providecommand{\DS}[1]{\hyperref[ds:#1]{DS-#1}}

\chapter{Limits of the self-verification}
\label{ch:meta-limits}

%% =========================================================
\section{What is not claimed}\label{sec:meta-not-claimed}

\begin{itemize}
\item No absolute consistency statement of any kind (\DS{2}).
\item No formal soundness theorem for the core (K1).  Within VNB, K1 could
  be a theorem only relative to set models: VNB cannot define truth over
  the class universe, and cannot prove that a model exists.  The
  soundness of the core for version-X is link~3 of
  Section~\ref{sec:meta-chain}, a hand-written argument (\DS{6}).
\item No formal statement that a checker computes its relation (K3).  The
  transcription is by hand (\DS{9}).
\item Nothing about the seven oracle relations, \texttt{ineq},
  \texttt{arith-ground}, \texttt{arith-forsome}, \texttt{arith-simplify},
  \texttt{comm-ring-simplify}, \texttt{ring-simplify} and \texttt{sos}
  (\DS{5}).  They are checked -- \texttt{ineq} and \texttt{sos} by a
  certificate the checker verifies by arithmetic, the others by an
  independent second decision (Section~\ref{sec:oracles}) -- but the
  checks add no proving power and are not a proof of the oracle.
\item Nothing about the certificate store and its hash, \texttt{qed} and
  the install, the undo journal, the Scheme of the expression layer
  (\texttt{alpha-equiv?}, \texttt{subst-free}, the matchers), or MIT
  Scheme.  The encoded operations restate the expression layer; they do
  not verify it.
\item The META theorems are VNB theorems with empty bills, but most of
  their certificates list the oracles \texttt{arith} or \texttt{ineq}.
  They rest on the oracle assumption as every theorem does.
\end{itemize}

%% =========================================================
\section{What is open}\label{sec:meta-open}

\subsection{Coverage of K2}

\begin{itemize}
\item The open rows of lemma~(*): ground arithmetic and arithmetic on
  certified numbers, $\mathrm{LENGTH}$ of a term typed in a class of
  tuples through a \texttt{sep} domain or a functoid body, an applied
  structure operation, and the view clause
  (Section~\ref{sec:meta-lemma-star}).  The B* rows are proven for every
  certificate built from the other rows; a certificate that uses an open
  clause is not covered.
\item The rewriting rules (PS) are proven at any finite set of positions
  (Section~\ref{sec:meta-positions}).  The \texttt{lambda-beta} licence
  through its binder's own scope, the binder-domain reads of the list
  definedness checks, a macete with a chain of conditions, and the
  numeral and successor bridge of the macete matcher are not covered.
\item The certified omissions of the rules that owe the definedness of a
  class (Section~\ref{sec:meta-owed}): the figure in which the kernel
  omits the owed sequent because its certificate accepts the class is
  open for \texttt{union-intro} at any entry and for
  \texttt{lambda-type} at a binder list.  For the empty tuple and for
  \texttt{lambda-type} at one binder it is covered when one member of
  the context certifies the class, not for a composition of the rows.
\end{itemize}

\subsection{Faithfulness (K3)}

K3 is deferred until K2 is complete (\DS{9}).  Its obligations:
\begin{itemize}
\item that each actual kernel step is a member of its transcribed
  relation, and that each C row is faithful (named syntax against locally
  nameless; eigenvariable and capture conditions);
\item the rows of $\mathrm{FUN\text{-}DEF}$, with the lambda-lifting of
  records, and of the predicate table $\mathrm{PRED\text{-}DEF}$, which has
  no row yet; until it has, every definition in $T$ of a predicate without
  a row must have a predicative body;
\item the head codes still missing (\texttt{pair}, \texttt{singleton},
  \texttt{power}, \texttt{choice}, the arithmetic heads, the structure
  operations);
\item that $T$ contains the encoded catalog axioms the induction and
  class-extensionality theorems assume;
\item that when the kernel certifies a term through clauses that are
  rows of lemma~(*), its definedness sequent is among the sequents the
  rows reach (Section~\ref{sec:meta-lemma-star});
\item that the kernel's predicativity test of a \texttt{comp} condition
  decides exactly $\mathrm{NBG\text{-}OK}$, the premise of def-comp;
\item the encoding of a destructuring lambda over a product, with the
  two properties argued for it: every member of a product is the list of
  its projections, and the encoded body is predicative exactly when the
  body is (Section~\ref{sec:meta-owed});
\item that the chain of one-position rewrites, in the order of
  Section~\ref{sec:meta-positions}, is the relation the rewriting
  checkers enforce;
\item that a value the kernel's value test accepts as a set has its
  sethood derivable, as $\mathrm{J\text{-}LB}$ requires;
\item the independence from the builder's certificate of the checkers
  that re-run it (Section~\ref{sec:meta-base}; \DS{32}).
\end{itemize}

\subsection{Hand-written arguments}

\begin{itemize}
\item The soundness of each core rule for version-X, and the truth of
  each axiom of the catalog, under the readings of \DS{4} and \DS{7}
  (link~3), are argued in Chapter~\ref{ch:meta-soundness} and, for
  functoids, in Chapter~\ref{ch:meta-functoids}.  The argument is
  conditional on the hypotheses (S1)--(S5) of
  Section~\ref{sec:meta-sound-hyp}.
\item The definitions of the library are covered by their general form
  (definitional equations of new heads, guarded definitions read as
  descriptions), not one by one.
\end{itemize}

\subsection{Enforcement}

\begin{itemize}
\item The macete builder offers a rewrite at a record variable, which the
  checker refuses.  The disagreement is in the safe direction.
\end{itemize}


\part{Using VNB}\label{part:using}

%% =========================================================
\chapter{Mathematical content}

\section{Number systems} \label{sec:nums}

The file \texttt{number-systems.scm} declares five number systems as
class constants and installs their arithmetic axioms in the
library.  All declarations take effect at load time.

\subsection{The standard number domains}

\begin{itemize}
\item \texttt{nn} ($\mathbb{N}$) Non-negative integers $\{0, 1, 2, \ldots\}$
\item \texttt{zz} ($\mathbb{Z}$) Integers
\item \texttt{qq} ($\mathbb{Q}$) Rationals
\item \texttt{rr} ($\mathbb{R}$) Reals
\item \texttt{cc} ($\mathbb{C}$) Complex numbers
\end{itemize}

Sethood axioms (\texttt{nn-is-set} through \texttt{cc-is-set}) and the inclusion chain
\[
  \mathbb{N} \subseteq \mathbb{Z} \subseteq \mathbb{Q} \subseteq \mathbb{R} \subseteq \mathbb{C}
\]
\begin{sloppypar}are provided as axioms \texttt{nn-subset-zz}, \texttt{zz-subset-qq},
\texttt{qq-subset-rr}, \texttt{rr-subset-cc}.\end{sloppypar}

\subsection{Arithmetic operations}

Arithmetic operations are shared across all number systems: there is one
\texttt{+}, one \texttt{*}, one \texttt{-}, and so on.  Membership
determines which system an argument belongs to; closure axioms guarantee
the result stays in the appropriate system.  For example,
\texttt{forall([x in nn, y in nn], x + y in nn)} is the
closure axiom for \texttt{+} on \texttt{nn}.

\begin{center}
\begin{tabular}{lll}
\toprule
System & Operations & Distinguished constants \\
\midrule
\texttt{nn} &
  \texttt{+}, \texttt{*}, \texttt{succ}, \texttt{<=} &
  \texttt{0}, \texttt{1} \\[3pt]
\texttt{zz} &
  \texttt{+}, \texttt{*}, \texttt{-} (unary), \texttt{<=} &
  \texttt{0}, \texttt{1} \\[3pt]
\texttt{qq} &
  \texttt{+}, \texttt{*}, \texttt{-}, \texttt{recip}, \texttt{<=} &
  \texttt{0}, \texttt{1} \\[3pt]
\texttt{rr} &
  \texttt{+}, \texttt{*}, \texttt{-}, \texttt{recip}, \texttt{abs}, \texttt{<=} &
  \texttt{0}, \texttt{1} \\[3pt]
\texttt{cc} &
  \texttt{+}, \texttt{*}, \texttt{-}, \texttt{recip}, \texttt{conjugate} &
  \texttt{0}, \texttt{1} \\
\bottomrule
\end{tabular}
\end{center}

\begin{remark}\label{remark:27135-34942-842417}
%
The infix notation \texttt{x \^{} n} is the preferred way to write
exponentiation; it parses to the same internal form as
\texttt{power(x, n)}.  The typing axiom is the natural closure form,
matching the two-argument application syntax:
\[
  \forall x\,\forall n.\;
    (x \in \mathtt{cc} \land n \in \mathtt{nn})
    \;\Rightarrow\; \texttt{power}(x,n) \in \mathtt{cc}
  \qquad\text{(axiom \texttt{power-typing-nonneg})}
\]
The defining axioms are \texttt{power-zero} ($x^0 = 1$) and
\texttt{power-succ} ($x^{\mathrm{succ}(n)} = x \cdot x^n$ for
$n \in \mathtt{nn}$).  A conditional equation \texttt{power-neg}
covers negative integer exponents: $x^{-n} = \mathrm{recip}(x^n)$
for $n \in \mathtt{nn}$, $x \neq 0$; but this does not extend the
\texttt{fun} typing to $\mathtt{zz}$, since $0^{-1}$ is undefined.
Complex exponentiation $x^z$ for $z \in \mathtt{cc}$ will be defined
later via the power series for \texttt{exp}.
%
\end{remark}
%
\begin{remark}\label{remark:27135-34995-146430}
%
In particular, $0^0 = 1$.  This is the mathematically correct value
for combinatorial and power-series purposes, as Knuth argues, and it is
what MIT~Scheme's exact arithmetic returns.  GNU/Maxima raises an error
here.  IMPS treated $0^0$ as undefined, which is logically sound but
forces theorems such as the binomial theorem to carry an explicit
$x \neq 0$ or $y \neq 0$ hypothesis to exclude the special case;
this system avoids that overhead by defining $0^0 = 1$.
%
\end{remark}

The arithmetic decision procedure \texttt{vnb-do-arith} (and the
\texttt{(arith)} command) evaluates \texttt{x \^{} n} on ground
numeric arguments.

\noindent
Integer literals are constants directly: \texttt{0}, \texttt{1},
\texttt{7}, \texttt{-3}, etc.\ are elements of \texttt{zz} (and of
\texttt{qq}, \texttt{rr}, \texttt{cc} by inclusion); non-negative integer
literals are additionally elements of \texttt{nn}.  No subscripted
forms (\texttt{0\_ZZ}, \texttt{1\_RR}, etc.) are used.  In MIT~Scheme
numeric literals are first-class objects, so this representation is
natural.

Integer literals must be in canonical decimal form: no leading zeros
(other than \texttt{0} itself), so \texttt{007} and \texttt{000075} are
rejected by the parser.  \texttt{-0} normalises to \texttt{0}.

The ordinal operators \texttt{ord-le}, \texttt{ord-lt}, and
\texttt{succ\_ord} are named apart from the numeric \texttt{<=} and
\texttt{succ}.  The first two were written \texttt{<=\_ORD} and
\texttt{<\_ORD} until 2026-08-24; those spellings could not be tokenized
(\texttt{read-op} consumes the \texttt{<} and stops, and \texttt{\_}
cannot begin a token), so no formula mentioning them could be retyped
from its own printed form.  \texttt{succ\_ord} keeps its subscript: it
begins with a letter and reads back.

\subsection{The arithmetic symbols: what they are, and what the axioms say}
\label{sec:arith-symbols}

This subsection states, in full, the logical status of \texttt{+},
\texttt{*}, unary \texttt{-}, \texttt{recip}, \texttt{abs},
\texttt{conjugate} and \texttt{succ}.  The point is small but it is easy to
misread, and a reader who assumes that \texttt{+} is ``the function
$\mathbb{C} \times \mathbb{C} \to \mathbb{C}$'' will be surprised by one
consequence below.

\paragraph{What kind of thing \texttt{+} is.}
\texttt{+} is a \emph{primitive operator symbol} of the signature.  The
operator census (\texttt{reference/OPERATORS.md}) files such a head under
``syntax: the head denotes nothing'': it forms terms and denotes nothing itself.  It is
registered in the constant-head registry
(Section~\ref{sec:constant-registry}) with kind \texttt{operator} by
\texttt{wff.scm}; it is applied syntactically, \texttt{(+ a b)}; it is never
substituted for; and it is not an object of the theory.  In particular it is
\emph{not} a function in the sense of Section~\ref{sec:equality} --- not an
element of any \texttt{fun($A$, $B$)}, not a set of ordered pairs, not built
by \texttt{vnb-lambda}, not typed by \texttt{lam-t}, not reduced by
\texttt{lam-b} --- and it is \emph{not} a functoid (Section~\ref{sec:def-functoid}):
no \texttt{def-functoid} or \texttt{lambdoid} introduces it and \texttt{beta} does
nothing to it.  \texttt{+} standing alone is not a term; \texttt{+ in set} is
not a formula anyone can establish or refute.  The same holds of every symbol
in the list above.

\paragraph{Arity.}
The parser reads a chain \texttt{a + b + c} as the \emph{flat} node
\texttt{(+ a b c)} (\texttt{parser.scm}, \texttt{p-parse-add}), while a
parenthesised right operand \texttt{a + (b + c)} keeps its nesting; the same
for \texttt{*}.  The supports \texttt{nary-plus-3}, \texttt{nary-plus-4},
\texttt{nary-plus-5} and their \texttt{nary-times} siblings
(\texttt{structure-library/numeric-instances.scm}) read the flat node as the
left-nested binary sum.  They are deliberately asserted: they fix a
\emph{reading} of the parser's node and carry no mathematical content.  So
\texttt{+} is $n$-ary in the syntax and binary in the axioms, and nothing is
lost in saying ``binary'' from here on.

\paragraph{What the axioms say about it --- and what they do not.}
For each number class $C$ among \texttt{nn}, \texttt{zz}, \texttt{qq},
\texttt{rr}, \texttt{cc} the base (\texttt{number-systems.scm}) installs the
\emph{closure} axiom
\[
  \forall a\,\forall b.\; (a \in C \land b \in C) \Rightarrow a + b \in C
  \qquad (\texttt{nn-add-closed}, \ldots, \texttt{cc-add-closed})
\]
and the algebraic identities of $C$ (commutativity, associativity, units,
inverses where they exist, distributivity), \emph{each guarded} by the
membership of its arguments in $C$.  That is the whole of what the theory
says about \texttt{+}.  It does \emph{not} say any of the following:
\begin{enumerate}
\item a domain: there is no axiom
  \texttt{+ in fun(cartesian(cc, cc), cc)} or anything of that shape;
\item the converse of closure: neither
  $a + b \in \mathtt{cc} \Rightarrow a \in \mathtt{cc} \land b \in \mathtt{cc}$
  nor $(a + b = a + b) \Rightarrow a \in \mathtt{cc}$;
\item that $a + b$ is undefined when an argument lies outside \texttt{cc}.
\end{enumerate}

\paragraph{The consequence, stated plainly.}
For two objects $a$ and $b$ about which the context knows nothing --- two
sets, two ordinals, a metric space and a matrix --- the term $a + b$ is
\emph{underdetermined}: the theory does not say whether it denotes, and if it
denotes it does not say what.  Underdetermined is not the same as undefined.
Undefinedness propagates (Section~\ref{sec:equality}): if \texttt{george(s)}
is undefined then \texttt{george(s) + 1} is undefined.  But definedness does
not propagate \emph{backwards} to a typing: from $a + b = a + b$ nothing
follows about $a$.  This is the ordinary situation of a function symbol in a
logic of partial terms.  The symbol carries no domain; its domain is whatever
the axioms let one derive; and here they let one derive membership of the
sum \emph{only} from membership of the summands.

\paragraph{The intended reading, and the others.}
The intended model takes \texttt{+} to be addition on $\mathbb{C}$, with its
restrictions to $\mathbb{R}, \mathbb{Q}, \mathbb{Z}, \mathbb{N}$.  Every
theorem of the system holds in that model, because (Chapter~\ref{ch:intro})
a theorem holds under \emph{every} reading of the primitives that satisfies
their axioms.  A reading in which \texttt{+} is total on all objects, and a
reading in which it is defined on exactly $\mathbb{C} \times \mathbb{C}$,
both satisfy the axioms; nothing provable distinguishes them, exactly as
nothing provable distinguishes Dedekind cuts from Cauchy classes.

\paragraph{Why it is harmless in a proof.}
Every mechanism that draws a conclusion from a term headed by \texttt{+}
demands the arguments typed first.
\begin{itemize}
\item The \emph{definedness certificate} consulted by \texttt{forall-elim}
  and \texttt{rfl} (\texttt{primitive-inferences.scm}, the heads in
  \texttt{*pi-arith-total-heads*}: \texttt{+ - * min max abs succ \^{}})
  certifies $a + b$ only when \emph{every} argument is a ground numeral or is
  typed in \texttt{nn}, \texttt{zz}, \texttt{qq}, \texttt{rr} or \texttt{cc}
  (or under \texttt{pos-rr}) by an assumption in the context
  (\texttt{pi--number-typed?}).  Instantiating a universal at $a + b$ with
  $a$, $b$ untyped therefore posts $a + b = a + b$ as an open obligation, and
  \texttt{rfl} refuses to close it (``not certified DEFINED'').
\item \texttt{crs} (\texttt{structure-library/comm-ring-simplify.scm},
  \texttt{ring-vars-ok?} / \texttt{cring-vars-ok?}) declines an identity
  unless every \emph{source} generator --- including those that cancel in the
  normal form --- is a variable of the typed universal prefix or is typed in
  the number domain, respectively in \texttt{carr($R$)}, by the context.
\item \texttt{ineq} certifies an atom only from a standalone assumption
  \texttt{$t$ in rr}; \texttt{arith} evaluates ground numerals only.
\item The algebraic laws are guarded, so a \texttt{fact} or macete use of one
  carries the typings as antecedents or side conditions.
\end{itemize}
Hence no derivation in the library has drawn a conclusion about $a + b$
without first establishing $a \in C$ and $b \in C$.  The underdetermination
is unobservable from inside a proof; it is visible only in a
\emph{statement}, see the last paragraph.

\paragraph{Why this design, and not a function.}
\begin{enumerate}
\item One symbol serves five systems under the inclusion chain, so no
  compatibility axioms between adjacent systems are needed (above).  As a
  function on $\mathbb{C} \times \mathbb{C}$ the five closure facts would
  still be required --- that $n + m \in \mathtt{nn}$ for $n, m \in \mathtt{nn}$
  does not follow from a codomain of $\mathbb{C}$ --- and the lower systems
  would have to be its restrictions.
\item Numerals are literals of the signature, and ground arithmetic is
  evaluated by Scheme on exact rationals (\texttt{arith-eval-term}):
  \texttt{2 + 3 = 5} needs no proof that an application is defined.
\item The oracles \texttt{crs}, \texttt{rs}, \texttt{ineq}, \texttt{sos},
  \texttt{arith} and the macetes \texttt{binary-minus-def},
  \texttt{binary-divide-def} all operate on the syntactic tree
  \texttt{(+ a b)}.  An application form would have to be unfolded before
  any of them could see it.
\item \emph{The function typing was tried and was unsound.}  Until
  2026-08-29 \texttt{numeric-instances.scm} carried axioms of the
  form \texttt{binplus in fun(cartesian(zz, zz), zz)}, one per operation and
  domain, described in their comment as closure facts.  They were not.
  Membership in \texttt{fun($A$, $B$)} pins the domain \emph{exactly}
  (\texttt{dom-of-fun} draws the equation out), so one object asserted into
  function classes over \texttt{zz} and over \texttt{rr} proved
  $\mathtt{zz} = \mathtt{rr}$; with all fourteen, $\mathtt{nn} = \mathtt{zz} =
  \mathtt{qq} = \mathtt{rr} = \mathtt{cc}$, and with \texttt{cc-i-squared},
  $0 \le i \cdot i = -1$: FALSITY with zero open leaves.  No proof ever took
  the route, but the axiom set had no model.  The fourteen were deleted and
  nothing replaced them, because the typing conjunct they existed to
  discharge is proved by \texttt{lam-t} from the closure axiom and sethood.
  A single symbol shared by five domains \emph{cannot} be a member of a
  \texttt{fun} class; a symbol with closure axioms can.
\end{enumerate}
The cost of changing the design would be the whole tree --- every statement
containing a sum, the parser and printer, the oracles, the certificate
clause, the macete surface --- and the one fact gained, ``$a + b$ defined
implies $a, b \in \mathtt{cc}$'', has never been needed by a proof.

\paragraph{Where the function form lives.}
When addition must be \emph{named as an object} --- as the \texttt{add} slot
of a ring instance, as something to quantify over, as a function to apply
through \texttt{app-graph} --- the object is \texttt{binplus}
(\texttt{structure-library/numeric-instances.scm}), tied to the syntax by the
unconditional \texttt{binplus(x, y) == x + y} (\texttt{binplus-apply}); and
inside a structure the slot holds a tupled \texttt{vnb-lambda} over the
instance's own carrier, so \texttt{add(cc-normed-field)} \emph{is} an element
of \texttt{fun(cartesian(cc, cc), cc)}, by the structure typing, with a
domain of exactly $\mathbb{C} \times \mathbb{C}$.  The two surfaces meet in
theorems such as \texttt{cc-additive-ag-opr} (\texttt{theorem-library/cc-series.scm}),
and \texttt{crs} reaches \texttt{add($R$)}, \texttt{mul($R$)}, \texttt{neg($R$)}
once \texttt{is-commutative-ring($R$)} is in the context.

\paragraph{The rule for authors of statements.}
The one place the underdetermination shows is a \emph{statement} that applies
\texttt{+} to the members of a class it has not placed inside a number system.
\texttt{forall([x in u, y in u], x + y in u)} with \texttt{u} unconstrained
says what the author means only for a \texttt{u} that lies in \texttt{cc}; for
any other \texttt{u} its truth value is underdetermined, and a definition
built on it is then a definition of something nobody intends.  The library
met exactly this case: the convexity predicate of the plane was first written
with the combination clause alone, and no theorem could read \texttt{u subset
cc} out of it (every theorem that needed the inclusion carried it beside the
predicate, from \texttt{holomorphic-on} or \texttt{is-open}).  Since
2026-10-05 it is \texttt{is-convex-cc(u)}, with the conjunct \texttt{u subset
cc} first and the name saying the ambient plane.  State the class inside the
number system --- a conjunct or antecedent \texttt{u subset cc} --- or type
the parameters in the statement itself, and the problem does not arise.  The
same applies to every symbol of the list at the head of this subsection.

\subsection{Axioms installed}

Each number system comes with: a sethood axiom, inclusion from the
previous system, closure of each operation, and algebraic identities
appropriate to its structure.

\begin{itemize}
\item \texttt{nn}: closed under \texttt{+} and \texttt{*}; constants
  \texttt{0} (additive unit) and \texttt{1} (multiplicative unit);
  commutativity, associativity, and distributivity for \texttt{+} and
  \texttt{*}.  Unary \texttt{-} and \texttt{recip} are not total here.
\item \texttt{zz}: $\mathtt{nn}$ axioms plus closure under unary
  \texttt{-} (additive inverse); $x + (-x) = 0$.
\item \texttt{qq}: $\mathtt{zz}$ axioms plus \texttt{recip} on nonzero
  elements; ordering \texttt{<=} compatible with \texttt{+} and \texttt{*}.
\item \texttt{rr}: $\mathtt{qq}$ axioms plus \texttt{abs} satisfying the
  triangle inequality, with \texttt{abs($x$) = 0} iff \texttt{$x$ = 0}.
\item \texttt{cc}: $\mathtt{rr}$ axioms plus \texttt{conjugate}, an
  involution compatible with $+$ and $\cdot$, with
  $a \cdot \overline{a} \in \mathbb{R}$ and $a \cdot \overline{a} \geq 0$.
  No ordering on $\mathtt{cc}$.
\end{itemize}

\noindent The abstract structural names---commutative monoid,
commutative ring, ordered field---and the corresponding officialised
ring instances are deferred to Section~\ref{sec:basic-rings}.  At this
point we are simply listing the operations and identities each
domain carries.

Because there is a single \texttt{+}, \texttt{*}, etc.\ shared across all
systems, no compatibility axioms between adjacent systems are needed:
there is nothing to be compatible with.

\begin{remark}
Each of \texttt{nn}, \texttt{zz}, and \texttt{qq} carries an infinite
family of membership axioms, one per numeral appearing in the library.
When a positive-integer numeral $n$ is introduced, it is implicitly
axiomatized as a iterated sum $1 + 1 + \cdots + 1$ ($n$ times), so
$n \in \mathtt{nn}$.  For \texttt{zz}, negative numerals are similarly
axiomatized as the additive inverse of such sums, giving $-n \in
\mathtt{zz}$.  These axioms are generated on demand: entering a numeral
literal in a formula automatically extends the library with the
corresponding membership fact.
\end{remark}

\subsection{Example use in a proof}

\begin{VNBcode}
; Prove: 0 is in RR
(sp (make-wff-from-string "0 in RR"))
(ta 'rr-zero-in)                ; introduce theorem (0 in rr)
(ass)                           ; goal is now in assumptions

; Prove: for reals a, b, (a + b) is in RR
(sp (make-wff-from-string "a in RR and b in RR implies a + b in RR"))
(di)                              ; assume (a in rr) and (b in rr)
(ai "a in RR and b in RR")       ; split into (a in rr) and (b in rr)
(ta 'rr-add-closed)               ; closure axiom: forall([a, b], ...)
(inst "forall([a, b], a in rr and b in rr implies a + b in rr)" 'a)
(inst "forall([b], a in rr and b in rr implies a + b in rr)" 'b)
(bc "a in rr and b in rr implies a + b in rr")
(ass)
; => Proof complete.
\end{VNBcode}

%% =========================================================
\section{Ordinals}
\label{sec:ordinals}

\subsection{The class \texttt{ord}}

\texttt{ord} is a primitive class constant, declared in the same style as
\texttt{nn}, \texttt{rr}, and \texttt{cc}: no explicit construction of
ordinals is needed.  The class is characterised entirely by its axioms.

\begin{itemize}
\item \texttt{ord} the class of all ordinals (a proper class)
\item \texttt{$\alpha$ in ord} $\alpha$ is an ordinal
\end{itemize}

\subsection{The ordering predicate}

Ordinals are equipped with a primitive ordering predicate \texttt{ord-le}
(and its strict companion \texttt{ord-lt}) that is independent of set
membership.  The ordering has no required connection to $\in$.

\begin{itemize}
\item \texttt{ord-le($\alpha$, $\beta$)} \quad $\alpha \leq \beta$ in the ordinal ordering
\item \texttt{ord-lt($\alpha$, $\beta$)} \quad $\alpha < \beta$ (strict)
\item \texttt{succ\_ord($\alpha$)} \quad successor ordinal of $\alpha$
\end{itemize}

\subsection{Axioms}

The following axioms are installed in \texttt{ordinals.scm}.

\begin{itemize}
\item \texttt{burali-forti}
  $\mathtt{ord} \notin \mathtt{set}$: \texttt{ord} is a proper class.
  Every \emph{set} of ordinals has a supremum in \texttt{ord}
  (axiom \texttt{sup-ord-in}); the supremum need not lie in the set itself
  (e.g.\ $\sup \mathtt{nn} = \omega \notin \mathtt{nn}$).
  If \texttt{ord} were a set, its supremum $\alpha$ would be an ordinal,
  hence $\mathit{succ\_ord}(\alpha) \in \mathtt{ord}$ would strictly
  exceed $\alpha$---contradicting that $\alpha$ bounds \texttt{ord}.
\item \texttt{ord-le-closure}
  $\alpha \leq_O \beta \;\Rightarrow\; \alpha \in \mathtt{ord} \land \beta \in \mathtt{ord}$
\item \texttt{ord-le-refl}
  $\alpha \in \mathtt{ord} \;\Rightarrow\; \alpha \leq_O \alpha$
\item \texttt{ord-le-antisymm}
  $\alpha \leq_O \beta \land \beta \leq_O \alpha \;\Rightarrow\; \alpha = \beta$
\item \texttt{ord-le-trans}
  $\alpha \leq_O \beta \land \beta \leq_O \gamma \;\Rightarrow\; \alpha \leq_O \gamma$
\item \texttt{ord-le-total}
  $\alpha,\beta \in \mathtt{ord} \;\Rightarrow\; \alpha \leq_O \beta \lor \beta \leq_O \alpha$
\item \texttt{ord-zero-least}
  $\alpha \in \mathtt{ord} \;\Rightarrow\; 0 \leq_O \alpha$
\item \texttt{ord-lt-iff}
  $\alpha <_O \beta \;\Leftrightarrow\; \alpha \leq_O \beta \land \alpha \neq \beta$
\item \texttt{ord-succ-in}
  $\alpha \in \mathtt{ord} \;\Rightarrow\; \mathit{succ\_ord}(\alpha) \in \mathtt{ord}$
\item \texttt{ord-succ-above}
  $\alpha \in \mathtt{ord} \;\Rightarrow\; \alpha <_O \mathit{succ\_ord}(\alpha)$
\item \texttt{ord-succ-immediate}
  $\alpha,\beta \in \mathtt{ord},\; \alpha <_O \beta \;\Rightarrow\; \mathit{succ\_ord}(\alpha) \leq_O \beta$
\item \texttt{ord-succ-injective}
  $\mathit{succ\_ord}(\alpha) = \mathit{succ\_ord}(\beta) \;\Rightarrow\; \alpha = \beta$
\item \begin{sloppypar}\texttt{limit-ord-iff}
  $\mathit{limit\text{-}ord}(\lambda) \;\Leftrightarrow\; \lambda \in \mathtt{ord} \land \lambda \neq 0 \land \lnot\exists\alpha.\,\lambda = \mathit{succ\_ord}(\alpha)$\end{sloppypar}
\item \texttt{ord-segment-is-set}
  $\alpha \in \mathtt{ord} \;\Rightarrow\; \mathit{ord\text{-}segment}(\alpha) \in \mathtt{set}$
\item \texttt{ord-segment-membership}
  $x \in \mathit{ord\text{-}segment}(\alpha) \;\Leftrightarrow\; x <_O \alpha$
\item \texttt{ord-segment-zero}
  $\mathit{ord\text{-}segment}(0) = \emptyset$
\item \texttt{ord-segment-succ}
  \[x \in \mathit{ord\text{-}segment}(\mathit{succ\_ord}(\alpha))
  \;\Leftrightarrow\; x \in \mathit{ord\text{-}segment}(\alpha) \lor x = \alpha\]
\item \texttt{sup-ord-in}
  $A \in \mathtt{set},\; A \subseteq \mathtt{ord} \;\Rightarrow\; \mathit{sup\text{-}ord}(A) \in \mathtt{ord}$
\item \texttt{sup-ord-upper}
  $\mathit{sup\text{-}ord}(A)$ is an upper bound of $A$
\item \texttt{sup-ord-least}
  $\mathit{sup\text{-}ord}(A)$ is the least upper bound
\item \texttt{sup-ord-empty}
  $\mathit{sup\text{-}ord}(\emptyset) = 0$
\item \texttt{transfinite-induction}
  see \S\ref{sec:ord-induction}
\end{itemize}

Here $\leq_O$ and $<_O$ abbreviate \texttt{ord-le} and \texttt{ord-lt}.

\subsection{Relationship to number systems}

\texttt{nn} is a subset of \texttt{ord} (axiom \texttt{nn-subset-ord}):
\[
  \forall n.\; n \in \mathtt{nn} \;\Rightarrow\; n \in \mathtt{ord}
\]
Consequently $0 \in \mathtt{ord}$ (from \texttt{nn-zero-in} and
\texttt{nn-subset-ord}).  The ordinal successor agrees with the numeric
successor on $\mathtt{nn}$ (axiom \texttt{ord-succ-nn}), and
\texttt{ord-le} restricted to $\mathtt{nn}$ coincides with the numeric
\texttt{<=} (axiom \texttt{ord-le-nn-compat}).  The number systems
\texttt{zz}, \texttt{qq}, \texttt{rr}, \texttt{cc} are not subsets of
\texttt{ord}.

\subsection{Transfinite induction}
\label{sec:ord-induction}

\paragraph{The induction axiom.}

\begin{sloppypar}The axiom \texttt{transfinite-induction} is the \emph{complete} (strong)
form: a class $C$ that contains every ordinal whose strict predecessors are
all in $C$ must contain every ordinal.\end{sloppypar}
\[
  \forall C.\;
  \Bigl(\forall\alpha \in \mathtt{ord}.\;
        (\forall\beta <_O \alpha.\; \beta \in C) \Rightarrow \alpha \in C\Bigr)
  \;\Rightarrow\;
  \forall\alpha \in \mathtt{ord}.\; \alpha \in C
\]

\paragraph{The \texttt{(tfi)} proof command.}

\texttt{(tfi)} applies transfinite induction to the current focus goal.
The goal must have the form
\[
  \forall \alpha.\; \alpha \in \mathtt{ord} \;\Rightarrow\; P(\alpha)
\]
and \texttt{(tfi)} replaces it with the single inductive-step subgoal
\[
  \forall \alpha.\;
    \bigl(\alpha \in \mathtt{ord} \;\land\;
          \forall\beta <_O \alpha.\; P(\beta)\bigr)
    \;\Rightarrow\; P(\alpha)
\]

\paragraph{The \texttt{(tfi3)} proof command.}

\texttt{(tfi3)} splits the same goal into three subgoals matching the
standard Peano-style cases:
\begin{enumerate}
\item \textbf{Base:} $P(0)$.
\item \textbf{Successor:} $\forall\alpha \in \mathtt{ord}.\; P(\alpha) \Rightarrow P(\mathit{succ\_ord}(\alpha))$.
\item \textbf{Limit:} $\forall\lambda.\;
  \mathit{limit\text{-}ord}(\lambda) \land
  (\forall\beta <_O \lambda.\; P(\beta)) \;\Rightarrow\; P(\lambda)$.
\end{enumerate}

\subsection{Transfinite recursion}
\label{sec:ord-recursion}

The theoretical basis for transfinite recursive definitions is the
following pair of propositions.  Both are proved by the standard
least-counterexample argument using well-foundedness of $\leq_O$.

\begin{sloppypar}
\begin{prop}[Transfinite recursion, sets]\label{prop:tfrec-set}
Let $U = \mathit{ord\text{-}segment}(\alpha)$ for some $\alpha \in \mathtt{ord}$,
let $A \in \mathtt{set}$, and let
\[
  \mathcal{F}(U, A) \;=\; \bigcup_{x \in U}\,\mathtt{fun}(\mathit{ord\text{-}segment}(x),\, A)
\]
be the set of all functions defined on an initial segment within $U$ with
values in $A$.  For any $\mathcal{E} \in \mathtt{fun}(\mathcal{F}(U,A),\, A)$ there
is a unique $f \in \mathtt{fun}(U, A)$ such that
\[
  f(x) \;=\; \mathcal{E}\bigl(f \!\restriction\! \mathit{ord\text{-}segment}(x)\bigr)
  \quad \text{for every } x \in U.
\]
\end{prop}
\end{sloppypar}

\begin{prop}[Transfinite recursion, class version]\label{prop:tfrec-class}
Let $A$ be a class, and let $\mathcal{E}$ be a functoid from the class of
all functions defined on initial ordinal segments with values in $A$ to $A$
(a proper-class domain, so $\mathcal{E}$ cannot be an element of
$\mathtt{fun}$).
Then there is a unique functoid $f$ on $\mathtt{ord}$ with values in $A$ such that
\[
  f(\xi) \;=\; \mathcal{E}\bigl(f \!\restriction\! \mathit{ord\text{-}segment}(\xi)\bigr)
  \quad \text{for every } \xi \in \mathtt{ord}.
\]
\end{prop}

\begin{remark}
Both $\mathcal{E}$ and $f$ in Proposition~\ref{prop:tfrec-class} are functoids, not
elements of $\mathtt{fun}(A,B)$: their domains are proper classes ($\mathtt{ord}$
for $f$; the class of all initial-segment functions for $\mathcal{E}$), and
$\mathtt{fun}$ requires set-valued domains.  Each is nonetheless a well-defined
class of ordered pairs with the function property.
\end{remark}

The Scheme helper \texttt{def-by-ord-recursion} implements this pattern;
see Section~\ref{sec:recursive-defs}.

%% =========================================================
\section{The base library}
\label{sec:axioms}

The base library \texttt{VNB-LIBRARY} is loaded automatically and
provides the following named axioms (all usable via
\texttt{cmd-theorem-assumption}):

\begin{itemize}
\item \texttt{membership-implies-sethood}
  $\forall a\,\forall b.\;(a \in b) \Rightarrow a \in \mathtt{set}$
\item \texttt{extensionality}
  $\forall a\,\forall b.\;(a \in \mathtt{set} \land b \in \mathtt{set}) \Rightarrow
  (a = b \Leftrightarrow \forall x.\,(x \in a \Leftrightarrow x \in b))$
\item \texttt{empty-set-is-set}
  $\texttt{empty-set} \in \mathtt{set}$
\item \texttt{empty-set-has-no-members}
  $\forall x.\;\lnot(x \in \texttt{empty-set})$
\item \texttt{pairing}
  $\forall a\,\forall b.\;(a \in \mathtt{set} \land b \in \mathtt{set}) \Rightarrow \{a,b\} \in \mathtt{set}$
\item \texttt{pairing-membership}
  $\forall a\,\forall b.\;(a \in \mathtt{set} \land b \in \mathtt{set}) \Rightarrow
  \forall x.\;(x \in \{a,b\}) \Leftrightarrow (x=a \lor x=b)$
\item \texttt{power-set}
  $\forall a.\;a \in \mathtt{set} \Rightarrow \texttt{power}(a) \in \mathtt{set}$
\item \texttt{power-set-membership}
  $\forall a\,\forall x.\;(x \in \texttt{power}(a)) \Leftrightarrow
  (x \in \mathtt{set} \land \forall z.\,(z \in x \Rightarrow z \in a))$
\end{itemize}

A theorem is added by proving it.  State it with \texttt{sp}, prove it,
and install it with \texttt{qed}, which also installs its macete:
\begin{VNBcode}
(sp "forall([x in nn], x in nn)")
(di) (ass)
(qed 'my-thm)
\end{VNBcode}
\texttt{qed} refuses unless the proof is complete, and it is the only
command that records a result as proven.  A statement that is to be
assumed for the time being is added with \texttt{support} and given a
stated justification with \texttt{warrant!}
(Chapter~\ref{ch:classification}).  It is then recorded as asserted, and
every theorem that depends on it says so when it is installed.


%% Definitions moved to ch-defs.tex


%% =========================================================
\section{Pending axioms}
\label{sec:pending-axioms}

The axioms in this section are needed for full set-theoretic reasoning
but are not yet installed.  They fall into two groups.

\subsection*{Installed in \texttt{axioms.scm}}

The following were previously pending but are now installed.
They are straightforward first-order statements requiring no new
expression primitives.

\begin{itemize}

\item \texttt{equality-symmetry}
  $a = b \Rightarrow b = a$.

\item \texttt{equality-transitivity}
  $a = b \land b = c \Rightarrow a = c$.

\item \texttt{fun-apply-type}
  $(f \in \mathtt{fun}(A,B)) \land (x \in A) \Rightarrow f(x) \in B$.
  Derivable from \texttt{fun-codomain-iff} but named for direct use.

\item \texttt{list-sethood}
  $A \in \mathtt{set} \land L \in \mathtt{tuples}(A) \Rightarrow L \in \mathtt{set}$.
  Derivable from \texttt{subset-set} (elements of any class are sets)
  but named for direct use in structure instance proofs.

\item \texttt{eq-subst-membership}
  $a = b \land b \in S \Rightarrow a \in S$.
  A first-order instance of \texttt{equality-substitution} restricted to
  membership; derivable from \texttt{equality-transitivity} and
  extensionality, but named for direct use in type proofs.

\item \texttt{quasi-eq-reflexivity}\ $(a \mathbin{==} a)$.

\item \texttt{quasi-eq-symmetry}\ $(a \mathbin{==} b) \Rightarrow (b \mathbin{==} a)$.

\item \texttt{quasi-eq-transitivity}\
  $(a \mathbin{==} b) \land (b \mathbin{==} c) \Rightarrow (a \mathbin{==} c)$.

\item \texttt{quasi-eq-def}\
  $(a = a) \Rightarrow ((a \mathbin{==} b) \Leftrightarrow (a = b))$;
  when $a$ is defined, quasi-equality coincides with strict equality.

\end{itemize}

The following were already installed (noted here for completeness):
\texttt{fun-domain-apply-def}, \texttt{fun-domain-extensionality},
\texttt{fun-codomain-iff} and \texttt{choice-axiom} in \texttt{library.scm};
\texttt{nn-subset-ord} and all \texttt{ord-*} axioms in
\texttt{structure-library/ordinals.scm}, not in \texttt{library.scm}
(see Section~\ref{sec:ordinals}).  The ordinal axioms carry
\texttt{primitive} provenance and so contribute nothing to any proof's bill.

\subsection*{List recursion and list induction}
\label{sec:list-recursion}

Both were listed as pending in earlier editions of this manual, and both are
now installed.  \texttt{structure-library/list-recursion.scm} adds the
constructor \texttt{cons}$(x, L)$ --- the name is \texttt{cons}, not
\texttt{prepend} --- with four defining equations, stamped
\texttt{definitional} because they constrain a new head and so add no
consequence about the old vocabulary:
\[
  \texttt{length}(\texttt{cons}(x,L)) = \texttt{succ}(\texttt{length}(L)),
  \qquad
  \texttt{nth}(1, \texttt{cons}(x,L)) = x,
\]
with \texttt{nth-cons-succ} shifting the remaining indices and
\texttt{makeset-cons} relating \texttt{make-set} to a union with a singleton.
Alongside them are two \emph{generation} axioms, which are claims about
\texttt{tuples} rather than about \texttt{cons} and are therefore asserted and
warranted: \texttt{tuple-length-zero} (a tuple of length $0$ is the empty
tuple) and \texttt{tuple-cons-decompose} (a tuple of length $\texttt{succ}(n)$
is $x$ prepended to a tuple of length $n$).

The induction principle is then a \emph{theorem}, not a schema-level primitive
inference: \texttt{theorem-library/tuples-induction.scm} proves
\[
  \forall A, C.\;
  [\,]\in C \;\wedge\;
  (\forall x, M.\; x \in A \wedge M \in \mathtt{tuples}(A) \wedge M \in C
     \Rightarrow \texttt{cons}(x,M) \in C)
  \;\Rightarrow\;
  \forall L \in \mathtt{tuples}(A).\; L \in C
\]
by induction on the \emph{length}, using \texttt{nn-induction} and the two
generation axioms.  It is stated in class form, quantifying over a class $C$
with membership in it as the conclusion, exactly as \texttt{nn-induction} and
\texttt{finite-set-induction} are; VNB quantifies over classes, so one theorem
does the work a schema would do in a first-order presentation.  Its bill reads
\texttt{modulo \{tuple-length-zero, tuple-cons-decompose\}}, which is the honest
report: the induction is derived, and what it rests on is the statement that
tuples are finite.

\subsection*{Still pending}

Four items once listed here are installed kernel rules and not pending:
\texttt{equality-substitution} is the rule \texttt{eq-subst}
(\texttt{pi-eq-subst!}), \texttt{lambda-type} is \texttt{pi-lambda-type!},
\texttt{lambda-beta} is \texttt{pi-lambda-beta!} (with
\texttt{lambda-beta-hyp} for an assumption), and \texttt{iota-def} is
\texttt{pi-iota-def!}, all in \texttt{primitive-inferences.scm}
(Section~\ref{sec:kernel-boundary}).

The remaining items require schema-level infrastructure --- a primitive
inference rather than a first-order axiom.  (An earlier edition also listed
\texttt{separation}, \texttt{union-set}, \texttt{comp} and \texttt{prepend}
as missing expression primitives.  They are not missing: \texttt{sep},
\texttt{comp} and \texttt{big-union} are registered term-forming heads
(\texttt{wff.scm}) with intro/elim rules in \texttt{primitive-inferences.scm},
and they are documented in \texttt{reference/KERNEL-RULES.md}.)

\begin{itemize}

\item \texttt{separation}
  $A \in \mathtt{set} \Rightarrow \{x \in A \mid p\} \in \mathtt{set}$.
  Requires the \texttt{separation} expression primitive.

\item \texttt{union-set} and \texttt{union-set-membership}
  $A \in \mathtt{set} \Rightarrow \bigcup A \in \mathtt{set}$, and
  $x \in \bigcup A \Leftrightarrow \exists a \in A.\; x \in a$.
  Requires the \texttt{union-set} (big union) expression primitive,
  distinct from the binary \texttt{union} constructor.

\item \texttt{cartesian-n-ary-sethood}
  All $A_i \in \mathtt{set} \Rightarrow A_1 \times \cdots \times A_n \in \mathtt{set}$.
  The binary case is installed (\texttt{cartesian-set-iff}); no fact of this
  name exists in the code, and the $n$-ary case is not installed under any name.
  %% TODO(docrot): item 12b -- whether the n-ary case follows from
  %% cartesian-set-iff by the tupling convention was not settled; no theorem found.

\item \texttt{comp-membership}
  $x \in \{y \mid p\} \Leftrightarrow p[y \mapsto x]$.
  Requires the \texttt{comp} (class comprehension) expression primitive.

\item \texttt{infinity}
  There exists an inductive set:
  $\exists A \in \mathtt{set}.\;\emptyset \in A \land
   \forall x \in A.\; x \cup \{x\} \in A$.

\end{itemize}

%% =========================================================


%% =========================================================
\chapter{Getting started}
\label{ch:quickstart}

This chapter walks through installing the system, the interactive workspace, and two complete worked examples.  It uses the vocabulary of Part~\ref{part:system}; a reader who wants to run something first may start here.

%% (moved from ch-intro.tex)
\label{sec:quickstart-chap}

\section{Installation and startup}

\subsection{Requirements}

MIT/GNU~Scheme~12.1, and GNU~Emacs~\cite{Emacs} (for the interactive workspace).
A \texttt{Makefile} at the root of the tree wraps the steps below
(\texttt{make help} lists them).
No other dependencies.

\subsection{Directory layout}

The prover lives in \texttt{\textasciitilde/prover/}: the executable
script \texttt{prover} and the loader \texttt{load.scm} are the only
two files you need to know about to get started.  See
Appendix~\ref{app:source} for a per-file overview.

\subsection{Compile the tree first}
\label{sec:compile-first}

Compiled binaries (\texttt{.com}, \texttt{.bin}) are excluded from both the
repository and the distributed archive, so a freshly cloned or unpacked tree is
\emph{interpreted}.  It will run correctly, and nothing in its output will say
that it is doing so; only the clock will.  On a modest machine the difference is
between a library load of seconds and one of minutes, and likewise for the
full test suite.

Compile once, after obtaining the tree.  This includes one load from source:

\begin{VNBsnip}
~/prover/VNB-with-compile    # compile every source file, then launch
\end{VNBsnip}

Thereafter recompile only what has changed.  From a session that has already
loaded \texttt{load.scm}:

\begin{VNBsnip}
(compile-vnb!)               # incremental: skips files whose .com is current
\end{VNBsnip}

After editing a source file, recompile it before running anything.  The
loader compares modification times: a \texttt{.com} older than its
\texttt{.scm} is ignored and the file loads from source.  That is harmless
for a proof script and costly for a core file, and it is the usual
explanation for a session that has become slow.  The question to ask is
whether \emph{that file} is compiled.  A small number of files are compiled
by nobody, deliberately: those that use a top-level macro, which
\texttt{compile-vnb!} lists in \texttt{*vnb-no-compile-files*} and skips.  They
load from source at no measurable cost, their work being done in the compiled
procedures they call.

\section{Starting the checker}

\begin{VNBsnip}
~/prover/prover              # interactive REPL
~/prover/prover proof.scm    # run a proof script and exit
~/prover/prover -i proof.scm # run a script, then stay interactive
\end{VNBsnip}

The interactive session is a standard MIT~Scheme REPL with the proof
checker pre-loaded.
\section{The VNB workspace}
\label{sec:emacs}

The VNB workspace is a styled Emacs session built for one task:
writing proofs.  Start it from the shell:

\begin{VNBsnip}
~/prover/VNB
\end{VNBsnip}

\noindent
A graphical Emacs frame opens with the VNB menu and toolbar already
in place, the prover ready to launch on first use, and a Home
Workspace inviting you to begin.

\begin{remark}
The launcher uses \texttt{emacs -Q}, so it ignores your personal
Emacs init file.  The session is dedicated to the prover; nothing
you do inside it affects your usual Emacs setup.
\end{remark}

\subsection{The Home Workspace}

The first workspace you see.  It is a list of buttons, each with a line of
explanation beside it (\texttt{emacs/vnb-launch.el}, \texttt{vnb-launch--paint-workspace}).
They fall into groups; the list is not exhaustive and the workspace itself is
the authority.

\begin{itemize}
\item \textbf{Starting work.}  \textbf{Start Proof} enters a formula in string
  syntax, switches to the Focus Workspace and begins the proof;
  \textbf{Load Proof Drive} runs a prepared script and stops on the leaf it
  leaves open; \textbf{Calculator} evaluates an arithmetic expression;
  \textbf{Build Structure} defines a new structure and saves it to a file in
  \texttt{structure-library/}.
\item \textbf{Looking things up.}  \textbf{What Is\ldots?}, \textbf{Show Theorems},
  \textbf{Find Theorem}, \textbf{Show PSS}, \textbf{Definitions},
  \textbf{Describe Structure}, \textbf{Structure Manual} and
  \textbf{Notes} show a structure, a number or constant, the installed
  theorems, the Proof Support Set, definitions, a structure card, all
  structures in one document, or a hand-written note.  \textbf{Proof as PDF}
  typesets a stored proof.
\item \textbf{Browsing the library.}  \textbf{Browse Library},
  \textbf{Fingerprint Index}, \textbf{Structure Graph}, \textbf{Library
  (HTML)} and \textbf{Browser Home} open the generated reference pages, some
  of them in a browser.
\item \textbf{Learning and scratch work.}  \textbf{Examples} steps through
  worked example proofs; \textbf{Scratch Workspace} is a Lisp-interaction
  sheet whose \texttt{C-j} sends an expression or region to the prover.
\item \textbf{Session.}  \textbf{Save Session} writes every proof completed
  in the session to one re-loadable script; \textbf{Quit} shuts down the
  prover and exits.
\end{itemize}

\noindent
Most buttons have a single-letter key (for example \texttt{s} starts a proof,
\texttt{t} shows theorems, \texttt{q} quits); the key legend is printed at the
foot of the workspace.  Press \texttt{g} to refresh.

\subsection{The Focus Workspace}

The working surface.  After \textbf{Start Proof} the workspace shows
the goal you are currently focused on, with its assumptions, and, above it, a
short row of buttons (\texttt{vnb-launch--paint-proof-1}):

\begin{itemize}
\item \textbf{What now?} --- asks the copilot what to try on the focused
  goal; each suggested move can be run from there.
\item \textbf{What is\ldots?} --- looks up a structure, number or constant.
\item \textbf{Preamble} --- runs the editable rule file on the focused goal:
  routine steps, then a report of what fired and where it stopped.
\item \textbf{Graph} --- walks the deduction graph one sequent node at a
  time, closed branches included (shown only while a proof is open).
\item \textbf{Undo} --- rolls back the last step (shown only while a proof
  is open).
\end{itemize}

\noindent
The tactics themselves are not buttons.  They run from single-letter keys,
printed at the foot of the workspace in groups: \emph{goal} (\texttt{d} direct
inference, \texttt{u} undo), \emph{hypotheses} (\texttt{a} assume,
\texttt{A} assume all, \texttt{D} decompose a hypothesis, \texttt{e}
separation elimination), \emph{rewriting} (\texttt{m} on the goal,
\texttt{M} on a hypothesis), \emph{citing} (\texttt{t} theorem, \texttt{F}
fact, \texttt{i} instantiate, \texttt{w} witness, \texttt{b} backchain,
\texttt{B} cite a lemma), \emph{closers} (\texttt{p} propositional, \texttt{y}
and \texttt{Y} syllogism), \emph{moving} (\texttt{f} focus, \texttt{o} overview,
\texttt{q} qed) and \emph{workspace} (\texttt{h} home, \texttt{r} scratch pad,
\texttt{S} scratch workspace, \texttt{T} toggle TeX rendering, \texttt{W} save
the proof script, \texttt{g} refresh).  The same moves are in the VNB menu and
can be typed on the scratch sheet.

\begin{remark}
Many more tactics are available in the prover than there are
keys.  The keys cover the most common everyday moves; the
full set is documented in Chapter~\ref{ch:writing-proofs}, and any
tactic can be typed directly in the REPL described below.
\end{remark}

\subsection{The Proof Overview Workspace}

A bird's-eye view of every open goal in the current proof.  Each
goal is rendered as assumptions above a sequent arrow above the
goal formula, with a \textbf{Focus on this} button beside it.
The current focus is marked with a small triangle.  Use this
workspace to navigate large proofs or to remember which goals are
still open.

Press \texttt{f} to return to the Focus Workspace, \texttt{h} for
Home, \texttt{n}/\texttt{p} to jump between goals, and
\texttt{RET} to focus on the goal at point.

\subsection{Error reporting}

When the prover rejects something --- a malformed formula, a tactic
that does not apply --- the error appears in two places at once:

\begin{itemize}
\item the echo area at the bottom of the frame flashes the message
  immediately;
\item a red \texttt{! Error:} panel appears at the top of the
  current workspace and stays there until the next successful
  action clears it.
\end{itemize}

\noindent
No detail is hidden in a side window: whatever the prover said,
you see.

\subsection{Display preferences}

The \textbf{Display} submenu of the VNB menu offers the usual
adjustments: frame size, position, font family and size,
foreground/background/cursor colours.  Any change is saved
automatically to \texttt{\textasciitilde/.vnb-display.el}, so the
next \texttt{VNB} launch reopens the frame exactly the way you
left it.  \textbf{Reset to Defaults} restores the built-in
appearance and removes the saved file.

\subsection{Advanced: the bare REPL}

Sitting behind the styled workspaces is an ordinary MIT~Scheme
REPL.  Open it with \textbf{Show Scratch Pad (advanced)} from the VNB
menu, or skip the workspaces entirely:

\begin{VNBsnip}
~/prover/VNB --no-workspace       # plain comint layout
~/prover/prover                   # no Emacs at all
\end{VNBsnip}

\noindent
The plain Emacs layout is a two-window arrangement: the
\texttt{*VNB*} REPL on the left, the live proof-state display
\texttt{*VNB State*} on the right.  Any Scheme expression you type
at the REPL prompt is evaluated; \texttt{(show)} refreshes the
state display, and short forms like \texttt{(di)} or \texttt{(ass)}
work exactly as they do from the workspace keys.

\paragraph{Scratch sheet of proof commands.}
\texttt{M-x vnb-command-buffer} opens a \texttt{*VNB Commands*}
window similar to Emacs's own \texttt{*scratch*}: type an
S-expression, press \texttt{C-j}, and the result is sent to the
running prover and inserted below.  \texttt{TAB} completes proof
command names; \texttt{C-c C-d} describes the command at point.

\paragraph{Programmatic evaluation.}
\texttt{(vnb-eval-string STR \&optional TIMEOUT)} sends \texttt{STR}
to the running prover, waits up to \texttt{TIMEOUT} seconds
(default~30) for the next REPL prompt, and returns the Scheme
return value as an Emacs string.  Use it from your own Elisp when
integrating the prover into a larger setup.

%% =========================================================
\section{A complete proof}
\label{sec:quickstart}

\begin{sloppypar}
Launch the workspace as in Section~\ref{sec:emacs} (\texttt{\textasciitilde/prover/VNB}).
The Home Workspace appears; the proof state below is what you would
see after stepping through the example in the Focus Workspace.  The
same commands typed at the bare REPL produce identical output.
\end{sloppypar}

We prove the classical syllogism: all men are mortal; Socrates is a
man; therefore Socrates is mortal.  There is no ``theory'' to extend
(Section~\ref{sec:structures-not-theories}): the premises are either
hypotheses of one statement, or the axioms of a structure.  Both
readings follow; run them in this order in one session, because the
second reserves the names it declares.

\paragraph{First reading: one statement, free variables.}
The classes \texttt{man} and \texttt{mortal} and the individual
\texttt{socrates} are free variables of the theorem; the premises are
its antecedent.

\begin{VNBsnip}
(sp (make-wff-from-string
  "(forall([x in man], x in mortal) and socrates in man)
     implies socrates in mortal"))
(di)
(ai 1)
(inst "forall([x in man], x in mortal)" "socrates")
(bc   "socrates in man implies socrates in mortal")
(ass)
(qed  'socrates-is-mortal)
\end{VNBsnip}

Final \texttt{*VNB State*}: \texttt{Proof complete.}  After \texttt{(di)}
the antecedent is assumption~1; \texttt{(ai 1)} splits the conjunction
into two assumptions; \texttt{inst} lands the implication for
\texttt{socrates}; \texttt{bc} turns the goal into its antecedent, which
\texttt{ass} finds in the context.

\paragraph{Second reading: a structure.}
Declare the species whose models are the worlds in which the syllogism
holds: a carrier of people, a carrier of mortals, a distinguished person,
and one law.  The theorem then quantifies over every model.

\begin{VNBsnip}
(declare-structure HUMANITY
  (instance-var s)
  (carriers PEOPLE MORTALS)
  (constant SOCRATES PEOPLE)
  (law "forall([x in PEOPLE(s)], x in MORTALS(s))"))

(sp (make-wff-from-string
  "forall([s in HUMANITY], SOCRATES(s) in MORTALS(s))"))
(di)
(mac-h 'humanity-class 1)
(mac-h 'is-humanity 1)
(dk-split-all!)
(inst "forall([x in PEOPLE(s)], x in MORTALS(s))" "SOCRATES(s)")
(bc   "SOCRATES(s) in PEOPLE(s) implies SOCRATES(s) in MORTALS(s)")
(ass)
(qed  'socrates-is-mortal-in-humanity)
\end{VNBsnip}

\texttt{(di)} lands \texttt{s in humanity}; the class axiom
\texttt{humanity-class} rewrites it to \texttt{is-humanity(s)}, and the
definition \texttt{is-humanity} unfolds that to the structure's typing
conditions and its law, one conjunction, which \texttt{dk-split-all!}
splits into separate assumptions; the rest is the first reading with
\texttt{socrates(s)}, \texttt{people(s)}, \texttt{mortals(s)} in place of
the free variables.

\begin{remark}
The same proofs in the styled workspace: \textbf{Start Proof} with the
statement, then the keys of the Focus Workspace (\texttt{d} for
\texttt{di}, \texttt{a} for \texttt{ass}, \ldots); a command with an
argument asks for it in the minibuffer.  The structure declaration is typed once in the
\texttt{*VNB Commands*} scratch sheet, or put in a file loaded at start.
The error panel (Section~\ref{sec:emacs}) catches typos as you go.
\end{remark}

\paragraph{What each command does.}
\begin{sloppypar}
\begin{itemize}
\item \texttt{(sp w)} starts a new proof with goal \texttt{w}.
\item \texttt{(di)} direct inference: introduces the goal's leading
  universals and the antecedent of an implication as assumptions.
\item \texttt{(ai f)} antecedent inference on assumption \texttt{f}
  (a formula or its 1-based index): a conjunction is split, an
  implication whose antecedent is in context is detached.
\item \texttt{(mac-h 'name k)} rewrites assumption \texttt{k} by the
  named definition or theorem; \texttt{(dk-split-all!)} splits every
  conjunction in the context.
\item \texttt{(inst f t)} instantiates the universally quantified
  assumption \texttt{f} at term \texttt{t}.
\item \texttt{(bc f)} backchains: assumption \texttt{f} is $P
  \Rightarrow Q$, current goal is $Q$, new goal is $P$.
\item \texttt{(ass)} closes a goal that appears in the assumption
  list.
\item \texttt{(qed 'name)} installs the completed proof as a named
  theorem.
\end{itemize}
\end{sloppypar}

\begin{remark}
The ``open goals'' count in \texttt{*VNB State*} climbs as the proof
progresses: each inference adds a node to the deduction graph, and
the parent stays open until its child grounds.  Completion is the
root grounding, not the count returning to zero in entry order.
Keep going until \texttt{Proof complete.} appears.
\end{remark}

\paragraph{The same proof, scripted.}
For batch use, the first reading in the long forms:
\begin{VNBsnip}
(let* ((wic (make-wff-from-string
              "(forall([x in man], x in mortal) and socrates in man)
                 implies socrates in mortal"))
       (ps  (start-proof wic))
       (ps  (cmd-direct-inference ps))
       (ps  (cmd-antecedent-inference ps
              (parse-string
                "forall([x in man], x in mortal) and socrates in man")))
       (ps  (cmd-instantiate ps
              (parse-string "forall([x in man], x in mortal)")
              'socrates))
       (ps  (cmd-backchain ps
              (parse-string
                "socrates in man implies socrates in mortal")))
       (ps  (cmd-assumption ps)))
  (cmd-qed ps 'socrates-is-mortal))
\end{VNBsnip}
All examples in this manual use the interactive short forms; the long
form is available for scripting.

\paragraph{What \texttt{make-wff} checks.}
\texttt{make-wff} (and \texttt{make-wff-from-string}, which calls it)
performs a position-typed walk of its argument and signals an error
rather than wrap a malformed expression.  The checks are:
\begin{sloppypar}
\begin{itemize}
\item \textbf{Numbers in wff positions} are rejected:
  \texttt{(make-wff-from-string "1 and 35")} $\Rightarrow$ ``number in wff
  position 1''.
\item \textbf{Term-forming heads in wff positions} are rejected:
  \texttt{(make-wff-from-string "choice(A)")} as a top-level wff fails because
  \texttt{choice} produces a term, not a wff.
\item \textbf{Wff-forming heads in term positions} are rejected:
  \texttt{(make-wff-from-string "(A and B) in C")} fails because \texttt{A and B}
  is a wff and the first argument of \texttt{in} must be a term.
\item \textbf{Bare symbols in wff positions} are rejected.  Only
  \texttt{truth} and \texttt{falsity} are valid bare-symbol wffs;
  anything else (e.g.\ \texttt{A implies B}) is an error.
\item \textbf{Free-symbol role conflicts}: a symbol that appears free
  in both wff and term positions in the same formula is rejected.
\item \textbf{Wff-position binders}: a binder whose body uses the
  bound variable in wff position is rejected within its scope.  This
  catches \texttt{forall([p], p implies p)}: \texttt{p} is bound as a
  class variable but the body uses it as a proposition.
\item \textbf{Bound vs.\ free same-name use} is \emph{accepted with a
  warning}.  In \texttt{forall([a], a in nn) implies a in zz} the
  inner \texttt{a} is bound and the outer \texttt{a} is free.  The
  validator emits a \texttt{;; warning} but does not reject.
\item \textbf{Arities and bound-variable shapes} are checked: every
  connective and quantifier gets the right number of arguments;
  \texttt{forall}, \texttt{forsome}, \texttt{iota}, \texttt{comp},
  \texttt{sep}, \texttt{big-union} and \texttt{vnb-lambda} bound variables must be
  symbols (\texttt{validate-wff!}, \texttt{wff.scm}); the parser requires each
  binding of a \texttt{lambdoid} to have the form \texttt{x in A}, and
  rejects \texttt{lambda}, which is not a binder.
\end{itemize}
\end{sloppypar}
%
\begin{remark}\label{remark:27135-34761-138906}
%
The set-theoretic \texttt{power(A)} (power set) and arithmetic
\texttt{x \^{} n} (exponentiation) share the same underlying symbol because
MIT~Scheme reads identifiers case-insensitively; \texttt{make-wff}
distinguishes them by arity (1 vs.\ 2 arguments).  All downstream
operations (\texttt{free-vars}, \texttt{subst-free},
\texttt{alpha-equiv?}, macete pattern matching) honour the same
arity-based dispatch.

%
\end{remark}

%% =========================================================
\section{A macete example}
\label{sec:quickstart-macete}

Macetes are named rewrite rules derived from installed theorems.
To apply a macete outside a proof, install a theorem and call
\texttt{apply-macete}:

\begin{VNBcode}
(install-theorem! 'eq-symm
  (make-wff-from-string "forall([x, y], x = y iff y = x)"))

(apply-macete 'eq-symm (make-wff-from-string "3 + 5 = 8"))
; => wff for "8 = 3 + 5"

(apply-macete 'eq-symm (make-wff-from-string "n = m implies truth"))
; => wff for "m = n implies truth"

(apply-macete 'eq-symm (make-wff-from-string "3 in nn"))
; => #f          -- pattern x = y does not match; no rewrite fires
\end{VNBcode}

\texttt{apply-macete} accepts and returns a \texttt{wff} object; it
returns \texttt{\#f} when the pattern matches nowhere.
To apply a macete as a proof step, use
\texttt{(mac \textit{name})} inside a proof (see
Section~\ref{sec:macetes}).

\begin{remark}
To try the four lines above in the styled workspace, open the
\texttt{*VNB Commands*} scratch sheet (\texttt{M-x vnb-command-buffer})
and press \texttt{C-j} after each S-expression; the prover's reply
appears on the line below.  The Home, Focus, and Proof Overview
workspaces are for \emph{proofs}; ad-hoc expression evaluation lives
on the scratch sheet or in the bare REPL.
\end{remark}

%% =========================================================


%% =========================================================


%% =========================================================
\chapter{Writing proofs}
\label{ch:writing-proofs}

Chapter~\ref{ch:proofs} describes the deductive system: sequents,
deduction graphs, the kernel operations that change them, and how
each inference is checked.  This chapter is about using it: starting a
proof, the commands and tactics that build one, managing the focus,
rewriting with macetes, and complete examples.  Every command here
works by calling kernel operations; none adds an inference of its own.

\section{Starting and inspecting a proof}

\begin{VNBcode}
; Start a new proof of a formula
(define ps (start-proof (make-wff-from-string
  "n in NN and m in NN implies n in NN")))

; Display the current state
(print-proof-state ps)

; Check if done
(proof-done? ps)   ; => #f or #t

; List open goals
(proof-open-goals ps)
\end{VNBcode}

The variable \texttt{ps} holds the \emph{proof state}: a deduction graph
plus the current \emph{focus} sequent node (the subgoal currently being
worked on).  Every proof command returns a new proof state with an
updated focus.

\texttt{start-proof} takes a \texttt{wff} (produced by \texttt{make-wff} or \texttt{make-wff-from-string}).
The root sequent has no assumptions: a proof begins from the bare goal,
and every hypothesis arrives through an inference.  All sequent assumptions and assertions are
\texttt{wff} objects throughout the proof; the library name propagates
automatically to every subformula produced by inference rules.

%% =========================================================
\section{Proof commands}
\label{sec:commands}

Commands are Scheme procedures of the form
\[
  \texttt{(cmd-}\mathit{name}\;\; \mathit{ps}\;\; \mathit{arg}_1 \ldots\texttt{)}
\]
Each command applies one primitive inference to the focus node, then
advances the focus to the first new ungrounded subgoal.

\subsection{Logical connectives}

\begin{itemize}

\item \texttt{(cmd-direct-inference ps)}
  Decompose the goal by its top-level connective:
  \begin{itemize}
  \item \texttt{p and q} $\to$ two subgoals $p$, $q$.
  \item \texttt{p implies q} $\to$ assume $p$, prove $q$.
  \item \texttt{p iff q} $\to$ two subgoals (each direction).
  \item \texttt{not(p)} $\to$ assume $p$, prove \texttt{falsity}.
  \item \texttt{forall([\ldots], p)} $\to$ peel the entire leading
    \texttt{forall} chain in one step, introducing a fresh variable for
    each bound variable.  If the substituted body is
    \texttt{$y$ in $A$ implies \ldots}, absorb \texttt{$y$ in $A$} into
    the assumptions and continue peeling.  Stop at the first
    non-\texttt{forall} head.
    Example: goal \texttt{forall([x in A, y], body)} yields one subgoal
    $x_0 \in A \;\Longrightarrow\; \texttt{body}[x/x_0,\,y/y_0]$.
  \end{itemize}

\item \texttt{(cmd-antecedent-inference ps $f$)}
  Decompose assumption $f$ by its top-level connective:
  \begin{itemize}
  \item \texttt{p and q} $\to$ replace with $p$ and $q$ separately.
  \item \texttt{p or q} $\to$ case split: branch on $p$, branch on $q$.
  \item \texttt{not(p)} $\to$ derive \texttt{falsity} if $p$ also in context.
  \item \texttt{forsome([\ldots], p)} $\to$ introduce a fresh witness for $x$.
  \end{itemize}

\item \texttt{(cmd-assumption ps)}
  Close goal immediately if it matches an assumption or is \texttt{truth}.

\item \texttt{(cmd-arith ps)}
  Close the current goal by ground arithmetic evaluation.  The goal must
  be a closed formula (no free variables) built from numeric literals,
  $+$, $*$, $-$, \texttt{recip}, \texttt{abs}, \texttt{succ},
  \texttt{conjugate}, $\leq$, $=$,
  $\in \mathtt{nn}/\mathtt{zz}/\mathtt{qq}/\mathtt{rr}/\mathtt{cc}$,
  and the propositional connectives.  Exact (rational) arithmetic only;
  no axiom chain is constructed.
  Related: \texttt{(vnb-do-arith \textit{f})} returns \texttt{\#t}/\texttt{\#f}
  for use outside a proof.

\item \texttt{(cmd-ring-simplify ps)}
  Close a goal by reducing both sides to their canonical form in the free
  associative $\mathbb{Z}$-algebra.  Addition is commutative (both sides
  are brought to the same sorted canonical order); multiplication is
  non-commutative.  Generators are detected automatically and no generator
  list need be supplied: \texttt{(rs)} treats any \emph{term} that
  \texttt{arith-eval} cannot reduce to a number as a generator of the free
  ring --- a bare symbol such as \texttt{x}, or a compound term such as
  \texttt{f(x)} whose head is not one of $+\,-\,*$.  A compound generator is
  a single opaque indeterminate: \texttt{f(x)} is atomic and its argument is
  not examined.

  The goal may be in either of two forms:
  \begin{itemize}
  \item \textbf{Bare equality} \texttt{$e_1$ = $e_2$}: generators are
    certified by \texttt{(in $v$ $D$)} assumptions already in the sequent
    context.  Typically reached after \texttt{(di)} peels the leading
    \texttt{forall} chain.
  \item \textbf{Universally quantified identity}
    \texttt{forall([$x_1$ in $D_1$, \ldots], $e_1$ = $e_2$)}: the
    quantifier typings themselves certify the generators.  \texttt{(di)}
    is not needed; \texttt{(rs)} closes the goal in one step.
  \end{itemize}

  In both cases every generator must be certified as an element of one of
  the five ring domains $\mathtt{nn}$, $\mathtt{zz}$, $\mathtt{qq}$,
  $\mathtt{rr}$, $\mathtt{cc}$ --- a symbol \texttt{x} by an
  \texttt{(in x $D$)} typing, a compound generator \texttt{f(x)} by an
  \texttt{(in f(x) $D$)} assumption.  This certification is also what keeps
  the rule sound on partial terms.  Since \texttt{=} is strict (it requires
  both sides defined), \texttt{3 + f(x) = 1 + f(x) + 2}
  holds only when \texttt{f(x)} is defined; membership entails definedness,
  so demanding \texttt{(in f(x) $D$)} supplies exactly that guarantee.  A
  generator that is not certified leaves the goal untouched --- the command
  never makes an unsound step.

  \begin{VNBcode}
;;; Directly on a universally quantified goal -- no (di) needed:
(sp (make-wff-from-string
  "forall([a in ZZ, b in ZZ, c in ZZ], a + b + c = c + b + a)"))
(rs)    ; closes in one step

;;; After (di), a compound term certified in context becomes a generator:
(sp (make-wff-from-string
  "forall([x in ZZ], (f(x) in ZZ) implies 3 + f(x) = 1 + f(x) + 2)"))
(di) (di) (rs)    ; di peels the forall then the hypothesis; rs then closes

;;; Coefficient arithmetic:
(sp (make-wff-from-string "forall([x in ZZ], 2*x + 3*x = 5*x)"))
(rs)
  \end{VNBcode}

  The primitive inference is \texttt{pi-ring-simplify!}, in
  \texttt{structure-library/ring-simplify.scm}; it normalizes each side to a
  canonical sum-of-monomials polynomial (\texttt{vnb->poly}) and compares.  The interactive short form is \texttt{(rs)}.

\item \texttt{(cmd-comm-ring-simplify ps)}
  The commutative counterpart of \texttt{(cmd-ring-simplify ps)}.  Where
  \texttt{(rs)} fixes the order of factors within a monomial,
  \texttt{(crs)} treats multiplication as commutative: a monomial is a
  \emph{multiset} of generators, canonicalised as a sorted generator
  list, so \texttt{x*y} and \texttt{y*x} share one normal form and
  commutativity lives in the representation rather than in any rewrite
  step.  Each side reduces to a sum of such monomials over \texttt{zz},
  and the goal closes iff the two normal forms coincide.  Use it for any
  identity whose proof needs reordering of factors, which \texttt{(rs)}
  cannot perform.

  The goal may take the same two shapes as for \texttt{(rs)} (a bare
  equality with generators certified by context typings, or a typed
  universally quantified identity), and additionally a \textbf{generic}
  form over an arbitrary ring: \texttt{forall($R$, is-commutative-ring($R$)
  implies \ldots)} built from \texttt{(add $R$)}, \texttt{(mul $R$)},
  \texttt{(neg $R$)}, \texttt{(zero $R$)}, \texttt{(one $R$)}, with the
  generators ranging over the carrier \texttt{(carr $R$)}.

  Generators are detected and certified exactly as for \texttt{(rs)}: a
  symbol or an opaque compound term, certified by an \texttt{(in $v$ $D$)}
  typing in the concrete form or an \texttt{(in $v$ (carr $R$))} assumption in
  the generic form.  The same partiality caveat applies --- a compound
  generator fires only once certified, which in VNB guarantees it is
  defined.

  \begin{VNBcode}
;;; Reordering of factors -- (rs) cannot close this, (crs) does:
(sp (make-wff-from-string
  "forall([x in ZZ, y in ZZ], x*y*x = x*x*y)"))
(crs)   ; closes in one step

;;; Commutative expansion:
(sp (make-wff-from-string
  "forall([x in ZZ, y in ZZ], (x+y)*(x+y) = x*x + 2*x*y + y*y)"))
(crs)
  \end{VNBcode}

  The procedure is a trusted oracle (it compares canonical normal forms
  rather than mechanising the ring properties), carrying the warrant
  \texttt{comm-ring-simplify} of kind \texttt{well-known}.  The primitive
  inference is \texttt{pi-comm-ring-simplify!} in
  \texttt{comm-ring-simplify.scm}.  The interactive short form is
  \texttt{(crs)}.

\item \texttt{(cmd-reflexivity ps)}
  Close goal \texttt{a = a} when \texttt{a} is demonstrably defined.
  Fails if \texttt{a} may be undefined.

\item \texttt{(cmd-quasi-reflexivity ps)}
  Close goal \texttt{a == a} unconditionally, for any term \texttt{a}.

\item \texttt{(cmd-cut ps $L$)}
  Introduce lemma $L$: subgoal~1 proves $L$; subgoal~2 has $L$ as assumption.

\item \texttt{(cmd-or-intro-left ps)}
  To prove \texttt{p or q}, prove $p$ instead.

\item \texttt{(cmd-or-intro-right ps)}
  To prove \texttt{p or q}, prove $q$ instead.

\item \texttt{(cmd-proof-by-contradiction ps)}
  To prove $p$, assume \texttt{not(p)} and derive \texttt{falsity}.

\item \texttt{(cmd-instantiate ps $F$\ $t$)}
  From \texttt{forall([$x$], $p$)} $= F$ in context, add $p[x/t]$.

\item \texttt{(cmd-exists-witness ps $t$)}
  To prove \texttt{forsome([$x$], $p$)}, prove $p[x/t]$.

\item \texttt{(cmd-weaken ps $f$)}
  Remove formula $f$ from the assumptions.

\item \texttt{(cmd-backchain ps $F$)}
  From \texttt{p implies goal} $= F$ in context, generate subgoal~$p$.

\item \texttt{(cmd-theorem-assumption ps $\mathit{name}$)}
  Add the named theorem (or axiom) to the assumptions.

\item \texttt{(cmd-apply-macete ps $\mathit{name}$)}
  Apply named macete: rewrite the goal using the associated theorem.

\item \begin{sloppypar}\texttt{(cmd-eq-subst ps $F$)}
  Equality substitution (Leibniz).  From an equality \texttt{$s$ = $t$} $= F$
  in the assumptions (either orientation), rewrite every occurrence of $s$
  in the goal to $t$.  Short form \texttt{(subst $F$)}.  This realises the
  schema \texttt{$s$ = $t$} $\wedge$ \texttt{$P[t]$} $\Rightarrow$
  \texttt{$P[s]$} for an arbitrary context $P$; sound for partial equality
  because \texttt{$s$ = $t$} already certifies both sides
  defined.\end{sloppypar}

\item \texttt{(cmd-qed ps $\mathit{name}$)}
  After a proof completes, install its root assertion as a named theorem.
  Errors if the proof is not complete.  This is the only command that
  records a result as \emph{proven}.

\end{itemize}

\subsection{N-ary constructors}
\label{sec:nary-proof-commands}

\begin{itemize}

\item \begin{sloppypar}\texttt{(cmd-cartesian-intro ps)}
  To prove \texttt{[$a_1$,\ldots,$a_n$] in cartesian($A_1$,\ldots,$A_n$)},
  generate $n$ subgoals \texttt{$a_i$ in $A_i$}.\end{sloppypar}

\item \begin{sloppypar}\texttt{(cmd-cartesian-elim ps $F$\ $k$)}
  From \texttt{[$a_1$,\ldots,$a_n$] in cartesian($A_1$,\ldots,$A_n$)} $= F$
  in context, add \texttt{$a_k$ in $A_k$}.\end{sloppypar}

\item \texttt{(cmd-tuples-intro ps)}
  To prove \texttt{[$a_1$,\ldots,$a_n$] in tuples($A$)},
  generate $n$ subgoals \texttt{$a_i$ in $A$}.
  For the empty sequence \texttt{[]} the goal closes immediately (0 subgoals).

\item \begin{sloppypar}\texttt{(cmd-tuples-elim ps $F$\ $k$)}
  From \texttt{[$a_1$,\ldots,$a_n$] in tuples($A$)} $= F$
  in context, add \texttt{$a_k$ in $A$}.\end{sloppypar}

\item \texttt{(cmd-nth-reduce ps)}
  Rewrite every occurrence of \texttt{nth($k$, [$\ldots$])} in the goal to
  the $k$-th element (1-indexed).

\item \begin{sloppypar}\texttt{(cmd-functoid-beta ps)}
  Rewrite every occurrence of
  \texttt{lambda([$x_1$ in $A_1$,\ldots], $\mathit{body}$)($v_1$,\ldots)}
  (or \texttt{lambdoid}) in the goal to
  $\mathit{body}[x_1 := v_1]\cdots$ (capture-avoiding).\end{sloppypar}

\item \texttt{(cmd-union-intro ps $k$)}
  To prove \texttt{$x$ in union($A_1$,\ldots,$A_n$)},
  prove \texttt{$x$ in $A_k$} instead.

\item \texttt{(cmd-union-elim ps $F$)}
  From \texttt{$x$ in union($A_1$,\ldots,$A_n$)} $= F$ in context,
  case-split: one branch per $A_i$ with \texttt{$x$ in $A_i$} assumed.

\item \begin{sloppypar}\texttt{(cmd-intersection-intro ps)}
  To prove \texttt{$x$ in intersection($A_1$,\ldots,$A_n$)},
  generate $n$ subgoals \texttt{$x$ in $A_i$}.\end{sloppypar}

\item \begin{sloppypar}\texttt{(cmd-intersection-elim ps $F$\ $k$)}
  From \texttt{$x$ in intersection($A_1$,\ldots,$A_n$)} $= F$ in context,
  add \texttt{$x$ in $A_k$}.\end{sloppypar}

\end{itemize}

%% =========================================================
\section{Two surfaces, one script}
\label{sec:two-surfaces}

Every proof step exists in two forms.

The \emph{builder} \texttt{(cmd-\textit{name}\ \textit{ps}\ \textit{arg}
\ldots)} (Section~\ref{sec:commands}) takes a proof state explicitly and
\emph{returns} a new one; it has no side effects and never prints.  It is
the form to call from Scheme code that manipulates proof states as values.

The \emph{interactive short form} \texttt{(\textit{name}\ \textit{arg}
\ldots)} (Section~\ref{sec:interactive}) is a thin wrapper around the
builder: it reads the global proof state \texttt{*ps*}, applies the
builder, stores the result back in \texttt{*ps*}, and calls \texttt{(show)}
so the new state is displayed.  This is the form typed in a live session.

\begin{center}
\begin{tabular}{@{}lll@{}}
\toprule
Builder & Short form & \\
\midrule
\texttt{(cmd-direct-inference ps)}        & \texttt{(di)}    & no argument \\
\texttt{(cmd-antecedent-inference ps $f$)}& \texttt{(ai $f$)}& formula argument \\
\texttt{(cmd-backchain ps $F$)}           & \texttt{(bc $F$)}& formula argument \\
\texttt{(cmd-instantiate ps $F$ $t$)}     & \texttt{(inst $F$ $t$)} & formula + term \\
\texttt{(cmd-apply-macete ps $n$)}        & \texttt{(mac $n$)} & theorem name \\
\bottomrule
\end{tabular}
\end{center}

\noindent
The wrapper does one more thing: it \emph{records itself}.  Each short form
appends an entry \texttt{(\textit{name}\ .\ \textit{args})} to the global
accumulator \texttt{*proof-script*}, which \texttt{(sp \ldots)} clears at
the start of a proof.  So there is no separate scripting language to learn:
\textbf{the script is the transcript} --- the very sequence of short forms
you typed.  \texttt{(qed \textit{name})} and \texttt{(save-proof
\textit{name})} store that sequence; \texttt{(replay-proof \textit{name})}
re-runs it (Section~\ref{sec:replay}).

\subsection{Supplying arguments}
\label{sec:arg-collection}

A short form takes one of a few kinds of argument, and each kind is
collected the same way wherever it appears.

\begin{itemize}
\item \textbf{Nothing.}  The structural steps --- \texttt{di}, \texttt{ass},
  \texttt{pbc}, \texttt{rs}, \texttt{rfl}, \texttt{ci}, \texttt{ni},
  \ldots\ --- read the focus goal and need no argument.

\item \textbf{A formula} (\texttt{ai}, \texttt{bc}, \texttt{wk},
  \texttt{cut}, \texttt{ue}, the assumption cited by \texttt{inst} and
  \texttt{mac-h}).  It may be supplied three interchangeable ways:
  \begin{itemize}
  \item a \textbf{surface string}, parsed for you ---
    \texttt{"n in nn and m in nn"};
  \item a \textbf{raw S-expression} ---
    \texttt{'(and (in n nn) (in m nn))};
  \item for a command that cites an existing \emph{assumption}, a
    \textbf{1-based assumption index}: a positive integer $k$ selects the
    $k$-th assumption of the focus sequent, exactly as numbered in the
    display.  \texttt{(ai 1)} stands in for retyping the first assumption.
  \end{itemize}

\item \textbf{A term} (\texttt{ew}, the witness slot of \texttt{inst},
  \texttt{if-true}/\texttt{if-false}).  A string or an S-expression; there
  is no index form, since a term is not an assumption to point at.

\item \textbf{A name} (\texttt{ta}, \texttt{mac}, \texttt{bc*},
  \texttt{qed}).  A symbol naming an installed theorem or macete.

\item \textbf{A small integer} (\texttt{ui $k$}, \texttt{ce $f$ $k$},
  \texttt{focus $n$}).  Selects a branch, a tuple component, or an open
  goal --- not an assumption index.
\end{itemize}

\begin{VNBcode}
(sp (make-wff-from-string "n in NN and m in NN implies n in NN"))
(di)        ; assume the conjunction (now assumption #1)
(ai 1)      ; crack assumption #1 -- the index, not the retyped formula
(ass)
(qed 'and-elim-demo)

(lookup-proof 'and-elim-demo)
; => ((di) (ai 1) (ass))     -- the index is recorded verbatim
\end{VNBcode}

\begin{remark}
The assumption-index form is accepted by the commands that \emph{cite} an
assumption --- \texttt{ai}, \texttt{bc}, \texttt{wk}, \texttt{inst},
\texttt{ce}, \texttt{te}, \texttt{ie}, \texttt{mac-h}.  It is recorded
verbatim in the script, so a replay re-resolves ``assumption~1'' against
the focus sequent it actually reaches --- robust to the formula's exact
text, but dependent on the assumption \emph{order} being the same on
replay.  When that is in doubt, give the formula itself.
\end{remark}

\subsection{Argument collection in a front end: \texttt{bc*}, \texttt{B}, \texttt{B+}}
\label{sec:bc-star-collection}

\texttt{bc*} (Section~\ref{sec:interactive}) matches a theorem's conclusion
against the goal and instantiates whatever that match determines.  A schema
variable the match leaves \emph{undetermined} must be supplied --- as a
binding list \texttt{((\textit{v}\ \textit{val})\ \ldots)} of Scheme
expressions.  The pure query \texttt{(bc*-undetermined '\textit{name})}
reports exactly which variables those are, without touching the proof.

A front end uses that query to collect the values \emph{by name} instead of
making the user write the quoted binding list: it asks for each undetermined
variable's value as an ordinary \emph{term string} (\texttt{"ran(f)"},
\texttt{"nn"}) and parses it.  In the Emacs Focus workspace this is the
\texttt{B} key --- the argument \emph{builder}.  Its companion key
\texttt{B+} is a saturating closer: it only commits a citation it
can first \emph{prove} will close the focus goal
(\texttt{bc*-can-close?} matches the conclusion \emph{and} checks every
hypothesis is already an assumption), so it never leaves the proof worse
than it found it.

%% =========================================================
\section{Interactive short-form commands}\label{sec:tactics-ref}
\label{sec:interactive}

\texttt{interactive.scm} provides the short-form commands --- one- or two-letter
wrappers around the \texttt{cmd-*} proof procedures --- that you type in a live
proof session, at the bare REPL or through the Emacs workspace.  Each mutates the
global proof state \texttt{*ps*} and calls \texttt{(show)} automatically; the
exceptions, which return a value rather than mutating state, are \texttt{sp},
\texttt{show}, \texttt{pp}, \texttt{fa}, \texttt{fs}, \texttt{vnb-unwind},
\texttt{vnb-wind}, \texttt{parse-string}, and \texttt{make-wff-from-string}.

The reference below is \emph{generated} from the live \texttt{*tactic-help*}
registry (through \texttt{reference/TACTICS.md}), so it lists the current commands
with their signatures and trust \emph{kind} and cannot drift from the prover.  At
the REPL, \texttt{(tactics)} prints the same menu and \texttt{(tactics 'name)} the
full per-command help.

%% AUTO-GENERATED by docs/gen-tactics-ref.py from reference/TACTICS.md -- do not edit.
%% Regenerate:  python3 docs/gen-tactics-ref.py   (after a prover load refreshes TACTICS.md)

\DeclareRobustCommand{\tactkind}[1]{{\footnotesize\textsc{#1}}}

Every tactic is tagged with a \emph{kind} grounded in the \texttt{dg-apply-rule!} tag it emits, so the trusted base is legible: \texttt{rule} is a single primitive kernel inference (the fixed trusted base); \texttt{oracle} is a trusted decision procedure run as a black box; \texttt{composite} only chains kernel rules, adding no new inference; \texttt{meta} performs no deduction (session, search, navigation).  A proof's trust surface is exactly its \texttt{rule} steps plus whichever \texttt{oracle}s and asserted premises it cites.

The same reference is available at the REPL with \texttt{(tactics 'name)} and in the browser file \texttt{reference/TACTICS.md}; all three are generated from one registry, so they never disagree.

An entry's \emph{Uses} line names every kernel operation the command's own code can cause to be recorded, those recorded by the procedures it calls included. It is the bound on what the command can do, whether or not any proof in the library exercises it, and it is measured by reading the source rather than declared alongside the command. One further set is common to every command and is stated once instead of in every entry: the hook that runs after any command which changed the proof, to close the definedness sequents an instantiation left owed. Section~\ref{sec:command-uses} tabulates the same figures for every command at once.

\subsection*{Starting \& finishing a proof}

\par\smallskip{\raggedright\noindent\texttt{(sp goal)}\quad\tactkind{meta}\index{sp@\texttt{sp} (tactic)}\par}\nobreak
\noindent Start a proof of `goal' (a ``string'', a raw S-expr, or a wff); clears the script.\par
\noindent\textit{Uses:} no kernel operation: it records no inference.\par
\noindent\textit{When useful:} starting a new proof\par

Begin a new proof: state what you want to prove.  The goal may be a ``string'' in surface syntax, a raw S-expr, or an already-built wff (e.g. (wff ``...'') or (make-wff '(...)));  sp coerces a string/S-expr for you.  This becomes your one open goal and clears any previous script.  (Technically: set the proof goal.)

\par\smallskip{\raggedright\noindent\texttt{(qed 'name)}\quad\tactkind{meta}\index{qed@\texttt{qed} (tactic)}\par}\nobreak
\noindent Install the finished proof as theorem `name' and save its replayable script.\par
\noindent\textit{Uses:} no kernel operation: it records no inference.\par
\noindent\textit{When useful:} no open goals remain -- record the result\par

Finish: once there are no open goals, this records the result as a named theorem you can cite later, and saves the replayable script.  (Technically: install-theorem! plus script capture; reports the asserted facts the proof still rests on, `proven modulo \{...\}'.)

\par\smallskip{\raggedright\noindent\texttt{(save-proof 'name)}\quad\tactkind{meta}\index{save-proof@\texttt{save-proof} (tactic)}\par}\nobreak
\noindent Save the current script under a name without finishing.\par
\noindent\textit{Uses:} no kernel operation: it records no inference.\par

Save the current (possibly unfinished) proof script under a name, to resume or replay later.  (Technically: snapshot the script without installing a theorem.)

\par\smallskip{\raggedright\noindent\texttt{(replay-proof 'name [subst])}\quad\tactkind{meta}\index{replay-proof@\texttt{replay-proof} (tactic)}\par}\nobreak
\noindent Re-run a saved script on the current goal, optionally renaming free vars.\par
\noindent\textit{Uses:} every one of the 64 kernel operations.\par

The optional subst is an alist ((old . new) ...) applied to every command argument before replay -- this is how one proof is reused at fresh eigenvariables.


\subsection*{Decomposing the goal}

\par\smallskip{\raggedright\noindent\texttt{(di)}\quad\tactkind{rule}\index{di@\texttt{di} (tactic)}\par}\nobreak
\noindent Direct inference: split an AND goal, move an IMPLIES antecedent into the assumptions, or introduce a fresh eigenvariable for a leading FORALL.\par
\noindent\textit{Uses:} \texttt{and-intro} \texttt{forall-intro} \texttt{iff-intro} \texttt{implies-intro} \texttt{not-intro}.\par
\noindent\textit{When useful:} the goal head is AND / IMPLIES / FORALL -- decompose before deciding\par

Break the goal into its pieces.  A conjunction `A and B' splits into two goals, A and B.  An implication `P implies Q' assumes P (P becomes a hypothesis you may use) and leaves you to prove Q.  A universal `for all x, ...' fixes an arbitrary x (a new constant standing for `any x') and asks for the statement at that x.  Apply it repeatedly until the goal is a single atomic statement.  (Technically: the goal-side introduction rules for AND / IMPLIES / FORALL; the constant introduced for a FORALL is an eigenvariable.  Pairs with mac, which unfolds a definition in the goal.)

\par\smallskip{\raggedright\noindent\texttt{(pbc)}\quad\tactkind{rule}\index{pbc@\texttt{pbc} (tactic)}\par}\nobreak
\noindent Proof by contradiction: assume the goal's negation, prove FALSITY.\par
\noindent\textit{Uses:} \texttt{proof-by-contradiction}.\par
\noindent\textit{When useful:} a direct argument stalls and assuming the negation gives something concrete to work with\par

Argue by contradiction: assume the goal is false and derive an absurdity.  (Technically: classical reductio -- adds the negation of the goal as a hypothesis and changes the goal to FALSITY.)

\par\smallskip{\raggedright\noindent\texttt{(oi-l)}\quad\tactkind{rule}\index{oi-l@\texttt{oi-l} (tactic)}\par}\nobreak
\noindent OR-intro left: reduce an (OR a b) goal to a.\par
\noindent\textit{Uses:} \texttt{or-intro-left}.\par
\noindent\textit{When useful:} the goal is (OR a b) and the LEFT alternative is provable\par

Prove a disjunction `A or B' by proving the LEFT alternative, A.  (Technically: or-introduction, left.)

\par\smallskip{\raggedright\noindent\texttt{(oi-r)}\quad\tactkind{rule}\index{oi-r@\texttt{oi-r} (tactic)}\par}\nobreak
\noindent OR-intro right: reduce an (OR a b) goal to b.\par
\noindent\textit{Uses:} \texttt{or-intro-right}.\par
\noindent\textit{When useful:} the goal is (OR a b) and the RIGHT alternative is provable\par

Prove a disjunction `A or B' by proving the RIGHT alternative, B.  (Technically: or-introduction, right.)

\par\smallskip{\raggedright\noindent\texttt{(ew term)}\quad\tactkind{rule}\index{ew@\texttt{ew} (tactic)}\par}\nobreak
\noindent Existential witness: discharge a FORSOME goal by supplying the witness term.\par
\noindent\textit{Uses:} \texttt{forsome-intro}.\par
\noindent\textit{When useful:} the goal is FORSOME and you have a witness term in mind\par

Prove `there exists an x with property P' by exhibiting a specific witness: you supply the term, and the goal becomes `P holds of that term'.  The everyday `take x = ...' step.  (Technically: existential introduction for a FORSOME goal.)

\par\smallskip{\raggedright\noindent\texttt{(ci)}\quad\tactkind{rule}\index{ci@\texttt{ci} (tactic)}\par}\nobreak
\noindent Cartesian intro: prove a CARTESIAN-product membership component-wise.\par
\noindent\textit{Uses:} \texttt{cartesian-intro}.\par
\noindent\textit{When useful:} the goal is membership in a CARTESIAN product\par

Prove that a pair (or tuple) belongs to a Cartesian product A x B by proving each coordinate lies in its factor.  (Technically: Cartesian-membership introduction, component-wise.)

\par\smallskip{\raggedright\noindent\texttt{(ti)}\quad\tactkind{rule}\index{ti@\texttt{ti} (tactic)}\par}\nobreak
\noindent Tuple intro: prove a tuple/LIST membership component-wise.\par
\noindent\textit{Uses:} \texttt{tuples-intro}.\par
\noindent\textit{When useful:} the goal is membership in a set of tuples / LIST\par

Prove that a list belongs to the set of tuples over A by proving each entry lies in A.  (Technically: tuple-membership introduction, component-wise.)

\par\smallskip{\raggedright\noindent\texttt{(ii)}\quad\tactkind{rule}\index{ii@\texttt{ii} (tactic)}\par}\nobreak
\noindent Intersection intro: prove (IN x (INTERSECTION ...)) for each branch.\par
\noindent\textit{Uses:} \texttt{intersection-intro}.\par
\noindent\textit{When useful:} the goal is membership in an INTERSECTION\par

Prove that x lies in an intersection of sets by proving it lies in each of them.  (Technically: intersection-membership introduction.)

\par\smallskip{\raggedright\noindent\texttt{(ui k)}\quad\tactkind{rule}\index{ui@\texttt{ui} (tactic)}\par}\nobreak
\noindent Union intro: reduce a goal (IN x (UNION ...)) to membership of x in the k-th set (1-based).\par
\noindent\textit{Uses:} \texttt{union-intro}.\par
\noindent\textit{When useful:} the goal is membership in a UNION and you know which branch\par

Prove that x lies in a union of sets by proving it lies in the k-th one (you choose which).  (Technically: union-membership introduction at branch k.)

\par\smallskip{\raggedright\noindent\texttt{(ni)}\quad\tactkind{rule}\index{ni@\texttt{ni} (tactic)}\par}\nobreak
\noindent Natural-number induction on the goal's leading FORALL over NN.\par
\noindent\textit{Uses:} \texttt{nn-induction}.\par
\noindent\textit{When useful:} the goal is `forall n in NN, P(n)'\par

Prove `for all natural numbers n, P(n)' by induction: show P(0), and that P(n) implies P(n+1).  (Technically: natural-number induction on the goal's leading FORALL over NN.)


\subsection*{Using a hypothesis}

\par\smallskip{\raggedright\noindent\texttt{(ai hyp)}\quad\tactkind{rule}\index{ai@\texttt{ai} (tactic)}\par}\nobreak
\noindent Antecedent inference: decompose a cited assumption -- AND-split, OR-into-cases, or FORSOME-elimination to a fresh eigenvariable.\par
\noindent\textit{Uses:} \texttt{and-elim} \texttt{forsome-elim} \texttt{iff-elim} \texttt{not-elim} \texttt{or-elim}.\par
\noindent\textit{When useful:} a hypothesis is an AND / OR / FORSOME to break apart\par

Break a hypothesis into its pieces -- the mirror image of di, but on the assumptions instead of the goal.  From a hypothesis `A and B' you get both A and B; from `A or B' you split into two cases (prove the goal in each); from `there exists x, P(x)' you get a fresh name for such an x together with P(x).  Cite the hypothesis by its number in the display, its formula, or a ``string''.  (Technically: the hypothesis-side elimination rules for AND / OR / FORSOME; the fresh name is an eigenvariable.)

\par\smallskip{\raggedright\noindent\texttt{(ass)}\quad\tactkind{rule}\index{ass@\texttt{ass} (tactic)}\par}\nobreak
\noindent Close the goal by an assumption alpha-equivalent to it.\par
\noindent\textit{Uses:} \texttt{assumption}.\par
\noindent\textit{When useful:} the goal already appears (up to bound-var renaming) among the hypotheses\par

Finish the goal because it is already one of your hypotheses -- the `that is exactly what we assumed' step.  Renaming of dummy/bound variables doesn't matter (`there exists x, P(x)' closes `there exists y, P(y)').  (Technically: the assumption rule, matching up to alpha-equivalence.)

\par\smallskip{\raggedright\noindent\texttt{(inst forall-hyp term)}\quad\tactkind{rule}\index{inst@\texttt{inst} (tactic)}\par}\nobreak
\noindent Instantiate a universally-quantified assumption at `term', adding the instance to context.\par
\noindent\textit{Uses:} \texttt{forall-elim}.\par
\noindent\textit{When useful:} you have a forall-hypothesis and a specific value to use it at\par

Use a `for all x, ...' hypothesis at a particular value: you supply the term, and the statement with x replaced by that term is added to your hypotheses.  (Technically: universal instantiation of an assumption.)

\par\smallskip{\raggedright\noindent\texttt{(inst+ forall-hyp term)}\quad\tactkind{composite}\index{inst+@\texttt{inst+} (tactic)}\par}\nobreak
\noindent Instantiate an in-context universal at `term', then forward-detach any guards whose antecedents are in context.\par
\noindent\textit{Uses:} \texttt{arith-forsome} \texttt{arith-ground} \texttt{arith-simplify} \texttt{cut} \texttt{detach} \texttt{forall-elim}.\par
\noindent\textit{When useful:} a GUARDED forall-hyp whose guards are dischargeable from context (e.g. the metric laws post-grind)\par

inst followed by detach: from `forall x. (x in S) => P(x)' and `t in S' already known, this lands `P(t)' directly (peeling nested guards level by level), instead of leaving the guarded implication for you to detach by hand.  The hypothesis-side analogue of `fact' (which assembles a THEOREM); this assembles an in-context UNIVERSAL.  scout's inst lane emits these to close witness-needing goals like the metric laws.  (Technically: pi-instantiate! then pi-detach! while the consequent stays a guard with an in-context antecedent.)

\par\smallskip{\raggedright\noindent\texttt{(mp)}\quad\tactkind{composite}\index{mp@\texttt{mp} (tactic)}\par}\nobreak
\noindent Close a goal that is an INSTANCE of a universal already in your context.  No arguments: the term is DERIVED by matching, not supplied.\par
\noindent\textit{Uses:} \texttt{arith-forsome} \texttt{arith-ground} \texttt{arith-simplify} \texttt{assumption} \texttt{cut} \texttt{detach} \texttt{forall-elim}.\par
\noindent\textit{When useful:} the goal IS an instance of a universal you already have -- no term to supply, it is derived by matching\par

The syllogism, with nothing to fill in.  From `forall([thing in human], thing in mortal)' together with `socrates in human', it closes `socrates in mortal'.  Where inst+ makes you name the term, mp works it out: it matches the universal's conclusion against the goal, which either determines the bound variable or fails outright, and then checks that the universal's guard at that term is something you already have.  Nothing is searched for and nothing is ranked -- if the goal is an instance there is exactly one term it can be.  It DECLINES, naming what it found, when no universal in the context has the goal as an instance, and when more than one does: a closer that silently picks between two universals is one you cannot predict.  One bound variable in this version; for `forall([a in A, b in B], ...)' use inst+.  (Technically: composite -- inst+ at the derived term, then ass.  It adds no kernel rule and no debt; the bill is what typing the two steps by hand would bill, which is nothing.)

\par\smallskip{\raggedright\noindent\texttt{(grind-and-mp)}\quad\tactkind{composite}\index{grind-and-mp@\texttt{grind-and-mp} (tactic)}\par}\nobreak
\noindent Tidy the goal with grind, then close it by modus ponens -- the one-button form of the syllogism, for a sentence typed exactly as it reads.\par
\noindent\textit{Uses:} \texttt{and-elim} \texttt{and-intro} \texttt{arith-forsome} \texttt{arith-ground} \texttt{arith-simplify} \texttt{assumption} \texttt{cut} \texttt{detach} \texttt{forall-elim} \texttt{forall-intro} \texttt{forsome-elim} \texttt{iff-elim} \texttt{iff-intro} \texttt{implies-intro} \texttt{macete-hyp} \texttt{not-elim} \texttt{not-intro} \texttt{or-elim}.\par
\noindent\textit{When useful:} a sentence typed as it reads -- `all men are mortal and socrates is a man implies socrates is mortal' -- closing in one move\par

Runs the two steps a syllogism actually takes on a sentence you have just typed.  `forall([thing in human], thing in mortal) and socrates in human implies socrates in mortal' is not yet a syllogism: it is an implication whose antecedent is a conjunction.  grind does the BOOKKEEPING -- split the AND, move the antecedents into the assumptions -- which leaves the goal `socrates in mortal' with the universal and `socrates in human' as hypotheses; mp then does the LOGIC (universal instantiation, then modus ponens).  They stay separate tactics because they are different kinds of step, and mp on its own is what to reach for once the goal is already tidy.  The recorded script keeps both steps, so the page reads as the two moves rather than as this wrapper.  (Technically: composite -- grind, quietly since on a tidy goal it legitimately declines, then mp.  Adds no kernel rule and no debt.)

\par\smallskip{\raggedright\noindent\texttt{(detach! impl)}\quad\tactkind{rule}\index{detach"!@\texttt{detach"!} (tactic)}\par}\nobreak
\noindent Forward modus ponens: from an in-context (IMPLIES A B) whose A is also in context, leave B in context.\par
\noindent\textit{Uses:} \texttt{detach}.\par
\noindent\textit{When useful:} you have both P and (P => Q) in context and want Q\par

If you have both `P' and `P implies Q' among your hypotheses, this adds `Q' to them.  It grows what you KNOW (the hypotheses) rather than changing the goal -- the forward counterpart of bc.  (Technically: forward modus ponens, the kernel rule pi-detach!; sound since P and P=>Q give Q.)

\par\smallskip{\raggedright\noindent\texttt{(fact 'thm term ...)}\quad\tactkind{composite}\index{fact@\texttt{fact} (tactic)}\par}\nobreak
\noindent Forward APPLICATION of a theorem: bring it in, instantiate its leading universals with the terms, and auto-detach every antecedent already in context, landing the consequent as a hypothesis.\par
\noindent\textit{Uses:} \texttt{arith-forsome} \texttt{arith-ground} \texttt{arith-simplify} \texttt{cut} \texttt{detach} \texttt{forall-elim} \texttt{theorem-assumption}.\par
\noindent\textit{When useful:} a library law `forall x. H(x) => P(x)' whose P you want landed as a hypothesis\par

Handles interleaved forall/implies (e.g. forall s. IS-X(s) => forall a. a in CARR(s) => P): consumes one term per FORALL, detaches each IMPLIES whose antecedent is in context.  The forward-assembly workhorse -- a law `forall x. H(x) => P(x)' becomes the usable fact P in one call, instead of ta + inst* + cut/backchain.  See theorem-library/module-zero-act.scm.

\par\smallskip{\raggedright\noindent\texttt{(ce hyp k)}\quad\tactkind{rule}\index{ce@\texttt{ce} (tactic)}\par}\nobreak
\noindent Cartesian elim: project the k-th component out of a CARTESIAN-membership assumption.\par
\noindent\textit{Uses:} \texttt{cartesian-elim}.\par
\noindent\textit{When useful:} a hypothesis is CARTESIAN-product membership\par

From a hypothesis that a tuple lies in a Cartesian product, extract that its k-th coordinate lies in the k-th factor.  (Technically: Cartesian-membership elimination, k-th projection.)

\par\smallskip{\raggedright\noindent\texttt{(te hyp k)}\quad\tactkind{rule}\index{te@\texttt{te} (tactic)}\par}\nobreak
\noindent Tuple elim: project the k-th component out of a tuple-membership assumption.\par
\noindent\textit{Uses:} \texttt{tuples-elim}.\par
\noindent\textit{When useful:} a hypothesis is tuple membership\par

From a hypothesis that something is a tuple over A, extract that its k-th entry lies in A.  (Technically: tuple-membership elimination, k-th projection.)

\par\smallskip{\raggedright\noindent\texttt{(ie hyp k)}\quad\tactkind{rule}\index{ie@\texttt{ie} (tactic)}\par}\nobreak
\noindent Intersection elim: extract the k-th branch of an INTERSECTION-membership assumption.\par
\noindent\textit{Uses:} \texttt{intersection-elim}.\par
\noindent\textit{When useful:} a hypothesis is INTERSECTION membership\par

From a hypothesis that x lies in an intersection, extract that x lies in the k-th set.  (Technically: intersection-membership elimination.)

\par\smallskip{\raggedright\noindent\texttt{(ue hyp)}\quad\tactkind{rule}\index{ue@\texttt{ue} (tactic)}\par}\nobreak
\noindent Union elim: split a cited (IN x (UNION ...)) membership assumption into one subgoal per set -- union-side case analysis, the dual of `ui'.\par
\noindent\textit{Uses:} \texttt{union-elim}.\par
\noindent\textit{When useful:} a hypothesis is UNION membership -- case-split on it\par

From a hypothesis that x lies in a union, split into cases -- one for each set x might belong to -- and prove the goal in each.  (Technically: union-membership elimination, case analysis; the dual of ui.)

\par\smallskip{\raggedright\noindent\texttt{(cut formula)}\quad\tactkind{rule}\index{cut@\texttt{cut} (tactic)}\par}\nobreak
\noindent Cut: prove `formula' as a side subgoal, then continue the main goal with `formula' added as an assumption (Gentzen cut).\par
\noindent\textit{Uses:} \texttt{cut}.\par
\noindent\textit{When useful:} you want to prove a lemma on the spot, then use it\par

Introduce a lemma you prove on the spot.  You state a formula; it becomes a side goal to prove, and on the main line you may then use it as a hypothesis.  The standard `we first show that ..., and now using it ...' move.  (Technically: Gentzen's cut -- nothing is left assumed-but-unproved, the side goal discharges it.)

\par\smallskip{\raggedright\noindent\texttt{(wk hyp)}\quad\tactkind{rule}\index{wk@\texttt{wk} (tactic)}\par}\nobreak
\noindent Weaken: drop a cited assumption from the context to tidy the hypothesis list.  `hyp' may be a formula, a ``string'', or a 1-based assumption index.\par
\noindent\textit{Uses:} \texttt{weakening}.\par
\noindent\textit{When useful:} the hypothesis list is cluttered with something no longer needed\par

Discard a hypothesis you no longer need, to keep the assumption list readable.  (Technically: weakening -- removing a hypothesis is always sound.)

\par\smallskip{\raggedright\noindent\texttt{(keep hyp ...)}\index{keep@\texttt{keep} (tactic)}\par}\nobreak
\noindent Keep only the cited assumptions and drop every other one, as ONE recorded step.  Each `hyp' may be a formula, a ``string'', or a 1-based assumption index; they are recorded as formulas.\par
\noindent\textit{When useful:} the context is deep and only a few hypotheses matter for the next step (prop, ineq)\par

Prune the context down to what the next step needs -- what a driver does before `prop' or `ineq' in a deep context.  (Technically: one weakening per dropped hypothesis, on the same leaf, recorded once; the driver kit's dk-only! is this command.)


\subsection*{Rewriting}

\par\smallskip{\raggedright\noindent\texttt{(mac 'name)}\quad\tactkind{rule}\index{mac@\texttt{mac} (tactic)}\par}\nobreak
\noindent Rewrite the GOAL with an equivalence macete (unfold a definition, apply an iff/=/== law).  Fires only where the macete's side-conditions already hold in context.\par
\noindent\textit{Uses:} \texttt{cartesian-decompose} \texttt{macete} \texttt{tuple-equality-decompose}.\par
\noindent\textit{When useful:} the GOAL has a defined predicate to unfold / an identity to apply, side-conditions already met\par

Rewrite the goal using a definition or a known equivalence/equality (a `macete').  For instance, replace a defined predicate by what it stands for, or apply an identity.  It fires only where the law's side-conditions already hold in your hypotheses; where they don't, it leaves that spot untouched and looks deeper inside.  (Technically: goal-side rewriting by an equivalence/equality macete; all-or-nothing -- it does NOT spawn unmet side-conditions as goals.  Its hypothesis-side cousin is mac-h; the minor-premise-spawning variant is macm.)

\par\smallskip{\raggedright\noindent\texttt{(slot 'acc)}\index{slot@\texttt{slot} (tactic)}\par}\nobreak
\noindent Reduce a structure ACCESSOR to its projection: (CARR s) becomes (NTH 1 s).  The one door for accessor reductions -- use it, never mac, on an accessor name.\par
\noindent\textit{Uses:} \texttt{cartesian-decompose} \texttt{macete} \texttt{tuple-equality-decompose}.\par

Replace an accessor by the tuple position it stands for: (CARR s) is slot 1, so it becomes (NTH 1 s).  Pair it with nth-r to compute on a concrete structure -- (MUL ZZ-RING) reduces to bintimes.  It refuses anything that is not an accessor.  (Technically: fires the accessor's macete, which def-structure installs at declaration.  It exists so that ALL accessor reductions go through ONE procedure: the reduction is currently global and unconditional -- index k for every argument -- and making it structure-relative later, guarded by IS-X(s), should change this door and not its callers.  accessor-callsite-audit fails the suite if any file fires an accessor macete by name.)

\par\smallskip{\raggedright\noindent\texttt{(slot-h 'acc hyp)}\index{slot-h@\texttt{slot-h} (tactic)}\par}\nobreak
\noindent Reduce a structure ACCESSOR to its projection inside a cited HYPOTHESIS -- `slot', hypothesis-side.\par
\noindent\textit{Uses:} \texttt{macete-hyp}.\par

The hypothesis-side door of slot: replace (PTS RR-MS) by the tuple position it stands for inside an assumption you cite, instead of in the goal.  Before it existed, a proof needing that reduction in a hypothesis reached for mac-h with the accessor macete's own name (rr-ms@pts) -- firing an accessor macete BY NAME, which is what the accessor pin (accessor-callsite-audit) exists to prevent, and which kept the suite red.  It refuses anything that is not an accessor, and it ERRORS rather than guess when the hypothesis mentions two different instances of the same accessor: after the first rewrite the assumption is a different formula, so a silent half-rewrite would be the failure mode.  (Technically: slot's projection-macete lookup run over the cited assumption, then cmd-apply-macete-to-assumption.)

\par\smallskip{\raggedright\noindent\texttt{(macm 'name)}\quad\tactkind{rule}\index{macm@\texttt{macm} (tactic)}\par}\nobreak
\noindent Like mac, but a CONDITIONAL macete fires even when its side-conditions are not yet in context: each unmet condition is left as a new subgoal.\par
\noindent\textit{Uses:} \texttt{cartesian-decompose} \texttt{macete} \texttt{tuple-equality-decompose}.\par
\noindent\textit{When useful:} the GOAL has a CONDITIONAL identity/definition to apply whose side-conditions you are willing to leave as subgoals\par

The minor-premises variant of mac.  mac only rewrites where the law's side-conditions ALREADY hold; macm rewrites regardless, spawning each unmet condition as a fresh goal to discharge -- the goal-side mirror of what mac-h already does on hypotheses.  This is what makes conditional finite-sum lemmas (finsum-*, sum-*, ...) usable as goal rewrites instead of only forward via fact.  (Technically: goal-side rewriting by a conditional macete; it emits the SAME single `macete' kernel rule as mac -- the unmet conditions become that one step's minor-premise subgoals.  It introduces NO new inference rule.)

\par\smallskip{\raggedright\noindent\texttt{(mac-h 'name hyp)}\quad\tactkind{rule}\index{mac-h@\texttt{mac-h} (tactic)}\par}\nobreak
\noindent Rewrite a cited ASSUMPTION in place with an equivalence macete; any side-condition not already in context is spawned as a new subgoal.\par
\noindent\textit{Uses:} \texttt{macete-hyp}.\par
\noindent\textit{When useful:} a HYPOTHESIS has a definition to unfold / an equivalence to apply\par

Like mac, but rewrites inside one of your HYPOTHESES instead of the goal -- replacing it by an equivalent statement (e.g. unfolding a definition in a hypothesis).  If the rewrite law has a side-condition you haven't established, that side-condition becomes a new goal to prove.  (Technically: hypothesis-side rewriting by Leibniz substitution of equivalents; sound because the macete is a genuine equivalence under its side-conditions, which are spawned as goals.)

\par\smallskip{\raggedright\noindent\texttt{(mac-h*)}\quad\tactkind{composite}\index{mac-h*@\texttt{mac-h*} (tactic)}\par}\nobreak
\noindent Saturating mac-h: repeatedly unfold every defined predicate in the hypotheses and split the conjunctions they expose, until nothing is left folded.  One step in the trace instead of a dozen.\par
\noindent\textit{Uses:} \texttt{and-elim} \texttt{forsome-elim} \texttt{iff-elim} \texttt{macete-hyp} \texttt{not-elim} \texttt{or-elim}.\par
\noindent\textit{When useful:} several hypotheses are folded definitions / conjunctions to open at once\par

The hands-free version of mac-h.  Instead of naming each definition and each conjunction by hand -- (mac-h 'IS-METRIC-SPACE A1), split, (mac-h 'is-metric A2), split, ... -- (mac-h*) keeps unfolding any assumption whose head is a defined predicate and splitting any AND assumption, restarting until a full pass changes nothing.  It records as a SINGLE step, so the proof (and its PDF) is far shorter.  It adds no kernel rule: every unfold is the same kernel-checked mac-h, every split the same ai -- just driven to a fixpoint.  Ask (what-now) to see which hypotheses it would touch first.

\par\smallskip{\raggedright\noindent\texttt{(grind)}\quad\tactkind{composite}\index{grind@\texttt{grind} (tactic)}\par}\nobreak
\noindent Saturate the no-choice moves at the focus: decompose the goal connective (di) and break open the hypotheses (mac-h*), repeatedly, until neither fires.  The whole introduce-everything prefix in one step.\par
\noindent\textit{Uses:} \texttt{and-elim} \texttt{and-intro} \texttt{forall-intro} \texttt{forsome-elim} \texttt{iff-elim} \texttt{iff-intro} \texttt{implies-intro} \texttt{macete-hyp} \texttt{not-elim} \texttt{not-intro} \texttt{or-elim}.\par
\noindent\textit{When useful:} right after sp, or any time the focus has connective/definitional structure -- normalize first\par

The deterministic normalizer.  It keeps applying di -- which strips a leading FORALL/IMPLIES/AND, introducing the bound variables and moving antecedents into your hypotheses -- and mac-h*, which unfolds defined predicates and splits conjunctions in the hypotheses, until the goal head is no longer a connective and no hypothesis is foldable.  None of these moves involves CHOOSING a lemma, so there is never anything to reconsider: grind just puts the focus into normal form (e.g. a metric-law goal becomes `(d s)(x,y) = (d s)(y,x)' with the metric's defining properties sitting unfolded in context).  It collapses the di...di mac-h* prefix that bloats proofs into a single step, and adds no kernel rule.

\par\smallskip{\raggedright\noindent\texttt{(scout [depth [branch [nodes]]])}\quad\tactkind{meta}\index{scout@\texttt{scout} (tactic)}\par}\nobreak
\noindent Speculatively search to a bounded depth across INDEPENDENT scratch branches; RETURNS a nested list (number-of-branches (d b) goal best-partials closing-branches) rather than printing.  Never touches the live proof.\par
\noindent\textit{Uses:} every one of the 64 kernel operations.\par
\noindent\textit{When useful:} you are stuck and want the machine to TRY move-sequences (returns data)\par

The copilot's deep lane.  Where (what-now)/(tt) suggest one move, scout actually TRIES sequences: it clones the focus into a fresh deduction graph per branch (so backtracking is free -- a dead branch is just discarded), and does a BEST-FIRST search (frontier ordered by open subgoals left, ties broken deeper-first so it dives toward a closure) whose alphabet is (grind), the closers (ass/rfl/crs/arith), the top rewrite-index (mac) rules, the parameterless backchain (bc*) lemmas, and the INST LANE -- (inst+ <universal hyp> <context-typed term>), instantiating a `forall' hypothesis at a term the context already types and detaching the guard.  It RETURNS the result as data -- (examined-count (depth branch) goal best-partials closing-branches), where best-partials is ((open-goals-left (form ...)) ...) most-reduced first and closing-branches is ((form ...) ...) shortest first -- so you can pick it apart programmatically; use (scout-show ...) for the readable REPL report.  Every move is the real, kernel-checked tactic, so a closing branch is a genuine proof when you adopt it with (scout-run k).  A CITATION GUARD suppresses circular closures: a branch that discharges the goal by citing a library theorem that is the goal -- DIRECTLY (statement alpha-equal to the goal, `P proved by P') or TRANSITIVELY (the cited theorem was itself proven using the goal, i.e. a name alpha-equal to the goal is in its proof-debt closure -- the `compact=>complete by compact=>bongo + bongo=>complete' trap) -- is dropped from closing-branches, the theorem name recorded in *scout-citations*, and scout-show notes `the goal is already theorem X'.  (Loops built purely from independent ASSERTED facts have no live dependency edge to detect; that is warrant/debt hygiene, not scout's to check.)  Bounds: depth (default 6), per-node fan-out branch (default 3), total nodes (default 600).  El cheapo: it does not reason about which move is wise (only the witness terms are guessed, from what the context already types -- a bad guess just dies on the clone); it brute-forces a small tree and you eyeball the survivors.  The inst lane closes the metric laws past grind; goals that need a witness scout can't type (a fresh existential, a constructed term) still surface under best-partials.

\par\smallskip{\raggedright\noindent\texttt{(scout-show [depth [branch [nodes]]])}\quad\tactkind{meta}\index{scout-show@\texttt{scout-show} (tactic)}\par}\nobreak
\noindent Like (scout), but PRINTS the human-readable report (examined count, goal, closing branches or best partials) to the REPL.  Returns the same nested list (scout) does.\par
\noindent\textit{Uses:} every one of the 64 kernel operations.\par
\noindent\textit{When useful:} same as scout but you want the readable report at the REPL\par

The eyeball version of scout: same search, same return value, but it also prints the numbered CLOSING branches (or, when none close, the best partials with how many goals each leaves open).  Use scout-show interactively, (scout) when you want to consume the result as data.

\par\smallskip{\raggedright\noindent\texttt{(scout-run k)}\quad\tactkind{composite}\index{scout-run@\texttt{scout-run} (tactic)}\par}\nobreak
\noindent Adopt closing branch k from the last (scout)/(scout-show) onto the live proof, running its tactics for real (they record and display normally).\par
\noindent\textit{Uses:} no kernel operation: it records no inference.\par
\noindent\textit{When useful:} scout found a closing branch you want to adopt onto the live proof\par

After scout finds CLOSING branches [1], [2], ..., (scout-run k) replays branch k's tactic forms through the real tactics on your live *ps*, from the same focus scout cloned -- so the proof advances and the steps are recorded for the script / PDF exactly as if you had typed them.  Works after either (scout) or (scout-show); both stash the closing branches.

\par\smallskip{\raggedright\noindent\texttt{(prep 'ineq)}\index{prep@\texttt{prep} (tactic)}\par}\nobreak
\noindent Why will a tactic not fire here, and what must be done first?  Runs the tactic's OWN preconditions one at a time, marks each ok / PREP / STOP, names the library lemma that repairs each unmet one, and prints a PLAN it has checked by running it on a scratch clone.  READ-ONLY: the live proof is never touched.  (prep 'ineq) is the only method so far.\par
\noindent\textit{Uses:} no kernel operation: it records no inference.\par

A tactic reports failure as a boolean, so every unmet precondition comes out as the same \#f and the same warning -- (ineq) says `goal not a linear-RR consequence of the named assumptions' whether your goal merely needs a di, or an atom needs a typing fact the library already proves, or the goal is simply false.  prep runs the same predicates separately and tells you WHICH failed.  For ineq the obligations are: the goal must BE an order relation (repair: di, counted by measuring, since one di consumes a typed FORALL and its guard together); every atom must be certified in RR by a LITERAL (IN t RR) scan -- (IN k NN) does not count, the oracle does no subtype reasoning (repair: a coercion lemma, e.g. nn-in-rr); some assumption must be order-shaped, since a fact like not(k=0) is invisible to Farkas (repair: a bridge lemma, e.g. nn-pos-of-nonzero); and the goal must actually follow, which only Fourier-Motzkin decides -- a STOP there means no prepping will help, which is the useful answer.  The repair search is by SHAPE over the whole library and is verified by SPECULATION: a candidate counts only if ineq then fires, not merely because it landed something order-shaped.  Every repair that works is reported, not just the first -- they come out alphabetical, which is no order of merit, and if the goal is itself a library theorem then citing IT is one of the closers.  Worked case: |- forall k in NN. \textasciitilde{}(k=0) => k < 2k.  prep returns (di) (di) (fact 'nn-in-rr 'k) (fact 'nn-pos-of-nonzero 'k) (ineq 4) -- six steps, where the hand-built cut/crs route through k = k+0 < k+k = 2k took thirteen and billed two extra transitivity lemmas.  A tactic joins the table with (prep-method! 'FUBA proc); crs and ass are the obvious next two.

\par\smallskip{\raggedright\noindent\texttt{(subst '(= s t))  |  (subst '(== s t))}\quad\tactkind{rule}\index{subst@\texttt{subst} (tactic)}\par}\nobreak
\noindent Rewrite s -> t throughout the goal, using an equation s = t -- or a quasi-equation s == t -- that is in context, in EITHER orientation (Leibniz substitution).  Reaches every position, operator (head) position included; a binder whose variable the equation mentions is not entered.\par
\noindent\textit{Uses:} \texttt{eq-subst}.\par
\noindent\textit{When useful:} you have an equation s = t (or a quasi-equation s == t) in context and want to rewrite the goal by it -- wherever the term occurs, operator position included\par

Use an equation among your hypotheses to replace s by t everywhere in the goal.  BOTH equalities license it: `=' is VNB's partial equality, and `s = t' already entails s = s and t = t, so t is defined wherever s was and the rewrite never replaces a defined term by an undefined one; `==' is quasi-equality (same definedness, equal where defined), which is a congruence and so substitutes exactly as `=' does -- which you need, since the partial-op recursion and bridge facts are stated with `=='.  The head you pass need not match the head in context, and neither need the orientation: all four of (= s t), (= t s), (== s t), (== t s) are searched for, and the goal is always rewritten s -> t.  LIMIT: the rewrite walk reaches argument positions only, so it is a silent no-op on a term in OPERATOR position -- (subst '(= (VADD md) ...)) will not touch the goal ((VADD md) x y).  Structure accessors are almost always in operator position; use the equation as a macete there (mac on the goal, mac-h on a hypothesis).  (Technically: Leibniz substitution from an in-context equality or quasi-equality.)

\par\smallskip{\raggedright\noindent\texttt{(beta)}\quad\tactkind{rule}\index{beta@\texttt{beta} (tactic)}\par}\nobreak
\noindent Beta-reduce a functoid application in the goal.\par
\noindent\textit{Uses:} \texttt{functoid-beta}.\par
\noindent\textit{When useful:} the goal has a functoid applied to an argument\par

Simplify a function-expression applied to an argument by substituting the argument into its body.  (Technically: beta-reduction of a functoid application in the goal.)

\par\smallskip{\raggedright\noindent\texttt{(nth-r)}\quad\tactkind{rule}\index{nth-r@\texttt{nth-r} (tactic)}\par}\nobreak
\noindent Reduce an NTH applied to a literal LIST in the goal.\par
\noindent\textit{Uses:} \texttt{nth-reduce}.\par
\noindent\textit{When useful:} the goal has NTH of a literal LIST\par

Simplify `the k-th entry of an explicit list [a1, a2, ...]' to that entry.  (Technically: reduce NTH applied to a literal LIST.)

\par\smallskip{\raggedright\noindent\texttt{(len-r)}\quad\tactkind{rule}\index{len-r@\texttt{len-r} (tactic)}\par}\nobreak
\noindent Reduce a LENGTH applied to a literal LIST in the goal.\par
\noindent\textit{Uses:} \texttt{length-reduce}.\par
\noindent\textit{When useful:} the goal has LENGTH of a literal LIST\par

Simplify `the length of an explicit list [a1, ..., an]' to the count n.  Structural (the dual of nth-r): the spine of a list literal is total, so its length is n regardless of whether the entries are defined.

\par\smallskip{\raggedright\noindent\texttt{(rfl)}\quad\tactkind{rule}\index{rfl@\texttt{rfl} (tactic)}\par}\nobreak
\noindent Close a reflexive equality goal (t = t).\par
\noindent\textit{Uses:} \texttt{reflexivity}.\par
\noindent\textit{When useful:} the goal is `t = t' with t manifestly defined\par

Close a goal `t = t' -- a thing equals itself.  (Technically: reflexivity of equality.)

\par\smallskip{\raggedright\noindent\texttt{(qrfl)}\quad\tactkind{rule}\index{qrfl@\texttt{qrfl} (tactic)}\par}\nobreak
\noindent Quasi-reflexivity: close t = t under the partial-equality definedness reading.\par
\noindent\textit{Uses:} \texttt{quasi-reflexivity}.\par
\noindent\textit{When useful:} the goal is `t = t' but t might be UNDEFINED (partial =)\par

Close `t = t' under the partial-equality reading, where asserting t = t also asserts that t is DEFINED.  Use this rather than rfl when t might be undefined.  (Technically: quasi-reflexivity; see the partial-equality convention, where `t = t' is the definedness predicate.)


\subsection*{Propositional logic}

\par\smallskip{\raggedright\noindent\texttt{(prop)}\index{prop@\texttt{prop} (tactic)}\par}\nobreak
\noindent Decide the goal by PROPOSITIONAL logic from the context, and close it if it follows.  Every non-connective formula -- `x in a', an equation, a whole `forall(...)' -- is one opaque atom; AND/OR/NOT/IMPLIES/IFF are read.  On a goal that does NOT follow it prints the countermodel (which atom must be true, which false) and leaves the proof untouched.  Adds no trust: it decides semantically, then discharges through di/ai/oi/ass/use-em/have!/detach!, so the bill is unchanged and the recorded script is the ordinary step-by-step proof.\par
\noindent\textit{Uses:} \texttt{and-elim} \texttt{and-intro} \texttt{arith-forsome} \texttt{arith-ground} \texttt{arith-simplify} \texttt{assumption} \texttt{cartesian-decompose} \texttt{cut} \texttt{detach} \texttt{eq-subst} \texttt{forall-elim} \texttt{forall-intro} \texttt{forsome-elim} \texttt{iff-elim} \texttt{iff-intro} \texttt{implies-intro} \texttt{macete} \texttt{not-elim} \texttt{not-intro} \texttt{or-elim} \texttt{or-intro-left} \texttt{or-intro-right} \texttt{proof-by-contradiction} \texttt{theorem-assumption} \texttt{tuple-equality-decompose}.\par
\noindent\textit{When useful:} the goal follows from the hypotheses by AND/OR/NOT/IMPLIES/IFF alone -- no quantifier or equality reasoning needed\par

Close a goal that follows from the hypotheses by pure propositional reasoning -- the leaves where you can SEE the answer and still have to pick between (oi-l)(ass), (oi-r) plus a conjunction split, and (ai) on a negation.  It treats each atomic statement as a black box, so it will not instantiate a quantifier or reason about equality: if the goal needs an instance, land the instance first (fact / inst+) and run it again.  When it declines it names a falsifying assignment, which usually tells you exactly which hypothesis is missing.  (Technically: three-valued evaluation over the atoms decides the entailment; the proof is then replayed through the kernel rules, case-splitting with excluded middle where a disjunctive goal needs it.)


\subsection*{Arithmetic \& ring oracles}

\par\smallskip{\raggedright\noindent\texttt{(arith)}\quad\tactkind{oracle}\index{arith@\texttt{arith} (tactic)}\par}\nobreak
\noindent Discharge a ground arithmetic goal by evaluation.\par
\noindent\textit{Uses:} \texttt{arith-forsome} \texttt{arith-ground} \texttt{arith-simplify} \texttt{cartesian-decompose} \texttt{eq-subst} \texttt{macete} \texttt{tuple-equality-decompose}.\par
\noindent\textit{When useful:} the goal is ground arithmetic -- no variables\par

Close a goal that is a concrete numerical fact with no variables -- e.g. 2 + 3 = 5, or 7 in NN -- by just computing it.  (Technically: decision by ground arithmetic evaluation.)

\par\smallskip{\raggedright\noindent\texttt{(rs)}\quad\tactkind{oracle}\index{rs@\texttt{rs} (tactic)}\par}\nobreak
\noindent Ring-simplify the goal (normal form over the ambient ring).\par
\noindent\textit{Uses:} \texttt{ring-simplify}.\par
\noindent\textit{When useful:} the goal is a (possibly non-commutative) ring-word identity\par

Simplify the goal to a normal form in the ambient ring (which need not be commutative).  The non-commutative companion of crs.  (Technically: ring-simplify to a canonical word form.)

\par\smallskip{\raggedright\noindent\texttt{(crs)}\quad\tactkind{oracle}\index{crs@\texttt{crs} (tactic)}\par}\nobreak
\noindent Commutative-ring decision procedure: prove a polynomial identity over ZZ[generators].  Expands literal powers, so (x+y)\textasciicircum{}2 = ... closes directly.\par
\noindent\textit{Uses:} \texttt{comm-ring-simplify}.\par
\noindent\textit{When useful:} the goal is a commutative-ring identity, e.g. (x+y)\textasciicircum{}2 = x\textasciicircum{}2+2xy+y\textasciicircum{}2\par

Prove a polynomial identity that holds in EVERY commutative ring -- e.g. (x+y)\textasciicircum{}2 = x\textasciicircum{}2 + 2xy + y\textasciicircum{}2 -- by reducing both sides to a canonical sum-of-monomials form and checking they agree.  A genuine decision procedure: a true commutative-ring identity closes, a non-identity is refused.  Literal powers are expanded for you.  (Technically: normal form over the free commutative ring ZZ[generators].)

\par\smallskip{\raggedright\noindent\texttt{(simp [target])}\quad\tactkind{oracle}\index{simp@\texttt{simp} (tactic)}\par}\nobreak
\noindent Rewrite a commutative-ring SUBTERM of the goal to canonical form, IN PLACE (e.g. (x+y)\textasciicircum{}2 inside a larger goal becomes x\textasciicircum{}2 + 2*x*y + y\textasciicircum{}2).  Works on BOTH surfaces: concrete number domains (+ * - \textasciicircum{} over NN/ZZ/QQ/RR/CC) and a generic ring s ((ADD s)/(MUL s)/(NEG s), carrier (CARR s)).  No arg = outermost ring subterm; ``term'' targets a specific one.  Sound by cut + crs + eq-subst (no new kernel rule); needs the subterm's generators typed in context (true post-di), else refuses and names them.\par
\noindent\textit{Uses:} \texttt{comm-ring-simplify} \texttt{cut} \texttt{eq-subst}.\par
\noindent\textit{When useful:} a ring SUBTERM of a larger goal should be put in normal form in place\par
\par\smallskip{\raggedright\noindent\texttt{(type-term)}\quad\tactkind{composite}\index{type-term@\texttt{type-term} (tactic)}\par}\nobreak
\noindent Close a typing goal (IN t C).  In a number class (NN ZZ QQ RR CC) t may be built from context-typed atoms (or atoms typed in a subclass: NN in ZZ, ZZ in RR, ...), numerals, + * -, and powers, typed by the class's closure laws, bottom-up.  In ANY class C, t may be an atom typed in a subclass, or an application f(a) with f in fun(A, C) in context and a typable in A, recursively: f(3 * x + 1) in B from f in fun(zz, B) and x in zz.  Declines -- with a warning, and nothing changed -- when there is no plan.\par
\noindent\textit{Uses:} \texttt{assumption} \texttt{comm-ring-simplify} \texttt{cut} \texttt{detach} \texttt{eq-subst} \texttt{forall-elim} \texttt{theorem-assumption}.\par
\noindent\textit{When useful:} the goal types a term: a polynomial in context-typed atoms in NN/ZZ/QQ/RR/CC, e.g. 3*x\textasciicircum{}2 + 3*x*y + y\textasciicircum{}2 in ZZ, or f(a) in any class C with f in fun(A, C) in context\par

The typing leaf that arithmetic leaves behind: `3 * x\textasciicircum{}2 + 3 * (x * y) + y\textasciicircum{}2 in ZZ' with x and y integers.  It plans the whole typing first -- which closure law types each subterm, or which subclass inclusion types an atom -- and only then runs it, so a term it cannot type leaves the proof untouched.  Each subterm's typing is proved on its own lane (cut) and lands in the context once.  A power x\textasciicircum{}k is typed by the power law in RR and CC; in ZZ, QQ and NN, where the library has no power law, by the identity x\textasciicircum{}k = x * ... * x (proved by crs) and the multiplication law, so the exponent must be a numeral there.  Subtraction in QQ goes through a - b = a + (-b) the same way; NN has no minus and declines.  Division and recip are not handled.  The kit twin (dk-type! t C) lands the typing in the context and returns it; in-rr hands its power goals to the same planner.  (Technically: composite -- cut, theorem-assumption, forall-elim, detach, assumption, eq-subst, and comm-ring-simplify for the power identities; no rule of its own.  Transactional: a run that does not close the goal is rolled back, script and undo stack included.)

\par\smallskip{\raggedright\noindent\texttt{(ew-poly)}\quad\tactkind{composite}\index{ew-poly@\texttt{ew-poly} (tactic)}\par}\nobreak
\noindent Close an existential goal forsome([a in C], L = R) (or unguarded) in which a occurs once, linearly: compute the witness by EXACT POLYNOMIAL DIVISION, supply it (ew), type it (type-term) and close the identity (crs).  One command for `(x+y)\textasciicircum{}3 = x\textasciicircum{}3 + y*a'.  Declines, with the reason and nothing changed, when the division leaves a remainder.\par
\noindent\textit{Uses:} \texttt{and-intro} \texttt{comm-ring-simplify} \texttt{forsome-intro}.\par
\noindent\textit{When useful:} the goal is forsome(a, L = R) with a occurring once, linearly, and the witness is a quotient of polynomials, e.g. (x+y)\textasciicircum{}3 = x\textasciicircum{}3 + y*a\par

The `find a such that ...' step when the answer is a quotient: for (x + y)\textasciicircum{}3 = x\textasciicircum{}3 + y * a the witness is ((x+y)\textasciicircum{}3 - x\textasciicircum{}3) / y = 3x\textasciicircum{}2 + 3xy + y\textasciicircum{}2.  Both sides are put in the sum-of-monomials normal form crs uses (integer coefficients over the generators, the witness variable among them); the equation must be of degree one in the witness variable, L - R = A*a + B, and the witness is -B/A by exact division of polynomials.  It declines, saying why, when a occurs other than exactly once, a side is not a polynomial, the equation is not linear in a, the coefficient A is 0, the division leaves a remainder, or the quotient has a fractional coefficient in ZZ or NN.  Then the recorded steps are ordinary ones: ew at the witness (written with + - * \textasciicircum{}, highest degree first), di to split the typing from the identity, type-term on the typing, crs on the identity.  what-now proposes it on such a goal, and scout's ew lane tries the same witness.  (Technically: composite -- forsome-intro, and-intro, the type-term steps, comm-ring-simplify; no rule of its own.  Everything is checked before anything runs, and the run is transactional besides.)

\par\smallskip{\raggedright\noindent\texttt{(expand)}\quad\tactkind{composite}\index{expand@\texttt{expand} (tactic)}\par}\nobreak
\noindent SHOW the expansion as a step: rewrite each polynomial side of an equation in the goal (also inside an existential or a conjunction) to its sum-of-monomials normal form, highest degree first.  (x + y)\textasciicircum{}3 = x\textasciicircum{}3 + y*a becomes x\textasciicircum{}3 + 3*x\textasciicircum{}2*y + 3*x*y\textasciicircum{}2 + y\textasciicircum{}3 = x\textasciicircum{}3 + y*a.  Declines, with the reason and nothing changed, when no side changes.\par
\noindent\textit{Uses:} \texttt{comm-ring-simplify} \texttt{cut} \texttt{eq-subst}.\par
\noindent\textit{When useful:} a side of an equation in the goal is an unexpanded polynomial and you want to SEE its expansion, e.g. (x+y)\textasciicircum{}3\par

The step a hand proof writes as `expanding, ...': the goal after it says what the equation says, term by term.  Each side t is proved equal to its normal form nf(t) by crs on a lane, and the goal is rewritten by subst, so the page shows the equation and the rewrite and both are checked.  A side is left alone when it mentions the variable an existential binds (that is not a term of the context -- the witness is what the goal asks for), when an atom of it is not typed in a number class (crs could not check it), or when it is already in normal form.  what-now offers it first on such a goal, and names the binomial theorem (binomial-theorem, n = k) when a power of a sum is expanded -- cited as what the expansion instantiates, not fired: the library states it over an abstract commutative ring as a finite sum.  (Technically: composite -- cut, comm-ring-simplify, eq-subst; no rule of its own; transactional.)

\par\smallskip{\raggedright\noindent\texttt{(zero-it)}\quad\tactkind{composite}\index{zero-it@\texttt{zero-it} (tactic)}\par}\nobreak
\noindent Move an equation goal P = Q (or P == Q, bare or under typed universals) to R = 0, where R is the normal form of P - Q (crs's calculator, highest degree first); prints `P - Q normalises to R' and returns R.  An inequality goal P <= Q or P < Q (the order of RR; atoms typed in NN, ZZ or QQ are lifted) is moved to 0 <= R or 0 < R, R the normal form of Q - P, the branch closed through rr-le-from-diff-nonneg / rr-lt-from-diff-pos; a numeral R decides it (0 <= -1: FALSE, nothing changed; 0 < 2: closed).  CC and an abstract ring have no order and are declined.  R = 0: closed by crs.  R a non-zero numeral: says the goal is FALSE and changes nothing.  Otherwise ONE recorded step leaves the goal R = 0, the atoms' typings landed or listed as owed side leaves.\par
\noindent\textit{Uses:} \texttt{comm-ring-simplify} \texttt{cut} \texttt{eq-subst}.\par
\noindent\textit{When useful:} an equation P = Q over a ring, or an inequality P <= Q / P < Q over the reals, and you want to see (and prove) what it reduces to: R = 0 (0 <= R), or a FALSE numeral that marks a dead path\par

The user's notes-42: `if in the end it means I have to prove 1 = 2, I can see I've been barking up the wrong tree'.  The ring is read off the typings of the atoms: the number surface (+ - * \textasciicircum{}, numerals; atoms typed in NN ZZ QQ RR CC, the largest class wins, NN lifting) or an abstract ring a with IS-COMMUTATIVE-RING(a) in context ((ADD a) (MUL a) (NEG a) (ZERO a) (ONE a)).  NN alone is refused for the rewrite (no subtraction).  Each atom of P and Q must be typed in the ring: an untyped one is typed by type-term's plan or, for recip(t) with not(t = 0) in context, by the recip-closure law; what cannot be typed is CUT as an owed side leaf `g in D', listed in the printout, and the new goal carries it as a hypothesis -- the notes' `[the eliminated terms all defined]'.  Then R = 0 is cut, and the branch R = 0 |- P = Q is closed by crs (P = Q + R), subst (P -> Q + R, then R -> 0) and crs; no sub-zero law and no typing of P or Q is needed.  Universals typed in the ring are peeled first (di, inside the same step).  The step records as (zero-it) with no argument: replay recomputes the normal form from the goal.  what-now shows the normal form on an equation goal and offers (zero-it).  (Technically: composite -- cut, comm-ring-simplify, eq-subst, plus di / theorem-assumption / forall-elim for the peeling and typings; the inner steps run with recording suppressed on the bc*-handler path, their citations kept for the bill; transactional.)

\par\smallskip{\raggedright\noindent\texttt{(supply)}\index{supply@\texttt{supply} (tactic)}\par}\nobreak
\noindent Close an inequality goal the way a hand proof would: beta-reduce any applied lambda, land the typing certificates and the standard bounds the linear oracle cannot see -- including the TRIANGLE inequality at the summands of an abs of a sum -- then call ineq over every premise.  Rehearses the whole sequence on a throwaway copy and does NOTHING unless it closes; on a miss it prints the sequence it tried.  Adds no trust: it drives lam-b, fact and ineq, each of which records itself.\par
\noindent\textit{Uses:} every one of the 64 kernel operations.\par

The committing form of what-now's SUPPLY THE ORACLE lane.  `ineq' reads abs(...), max(...) and (dist(s))(x,y) as opaque atoms and refuses any premise whose atoms are not certified real, so an inequality that is obviously true can fail for want of a typing nobody mentioned.  This lands them.  The one real idea is subadditivity: where the goal bounds abs(u + v), the useful intermediate term is abs(u) + abs(v), and the decomposition is in the term itself -- nothing is searched for.  Because a lambda reduction cannot be undone and can owe an unprovable leaf, the whole sequence is rehearsed before any of it is run.  (Technically: certificate closure to depth 2 over the forward-citation lane, plus the curated bound table, then Fourier-Motzkin.)

\par\smallskip{\raggedright\noindent\texttt{(ineq i1 i2 ...)}\quad\tactkind{oracle}\index{ineq@\texttt{ineq} (tactic)}\par}\nobreak
\noindent Close a linear-inequality goal over RR (<= < > >= = between RR terms) as a consequence of the named assumptions (1-based indices), via the Fourier-Motzkin/Farkas oracle.  Linearizes over + - * and the binplus/binneg/bintimes aliases; every MAXIMAL non-arithmetic subterm is an atom that must be certified in RR.  (Does NOT see through a generic ring's (ADD s)/(MUL s) -- those become opaque atoms.)\par
\noindent\textit{Uses:} \texttt{ineq}.\par
\noindent\textit{When useful:} a LINEAR <= / < over RR that follows from inequalities you can cite\par

Close a LINEAR inequality over the reals that follows from inequalities you cite (by their hypothesis numbers) -- by chaining them, adding them, and scaling by positive constants.  Anything that is not built from + - and multiplication-by-constants is treated as an opaque quantity, so it handles e.g. the triangle inequality where the distances are unknowns.  For NONLINEAR (polynomial) inequalities use sos instead.  (Technically: Fourier-Motzkin / Farkas over the ordered field RR; every atom must be certified real.)

\par\smallskip{\raggedright\noindent\texttt{(sos ``c1'' ``c2'' ...)}\quad\tactkind{oracle}\index{sos@\texttt{sos} (tactic)}\par}\nobreak
\noindent Sum-of-squares closer for a nonstrict polynomial inequality a <= b over RR (the nonlinear companion of (ineq)).  You supply the terms to be SQUARED; it finds the nonnegative coefficients.  (tactics 'sos) for the worked example.\par
\noindent\textit{Uses:} \texttt{sos}.\par
\noindent\textit{When useful:} a NONSTRICT POLYNOMIAL <= over RR (the nonlinear cousin of ineq)\par

A sum-of-squares closer.  To prove  a <= b  it is enough to exhibit b - a  as a sum of squares (each obviously >= 0), reducing the inequality to an algebraic identity.  You hand sos the terms to be SQUARED -- the c\_i, NOT the squares -- and it solves for nonnegative rationals lambda\_i with

\begin{VNBsnip}
      b - a  =  lambda_1 c_1^2 + ... + lambda_n c_n^2
\end{VNBsnip}

(matched coefficient-by-coefficient in crs's commutative-ring normal form), then PRINTS the lambda\_i it found.  You do not supply them.

Worked example.  Goal  forall([x in rr, y in rr], x*y <= x\textasciicircum{}2 + y\textasciicircum{}2).  Here b - a = x\textasciicircum{}2 - x*y + y\textasciicircum{}2 = 1/2(x-y)\textasciicircum{}2 + 1/2 x\textasciicircum{}2 + 1/2 y\textasciicircum{}2,  so the things squared are  x-y, x, y:

\begin{VNBsnip}
      (sos "x - y" "x" "y")
          ;; closes, printing  1/2(x-y)^2 + 1/2 x^2 + 1/2 y^2
\end{VNBsnip}

2*x*y <= x\textasciicircum{}2 + y\textasciicircum{}2 needs only one square:  (sos ``x - y'').  The three-variable a*b+b*c+c*a <= a\textasciicircum{}2+b\textasciicircum{}2+c\textasciicircum{}2  wants  (sos ``a - b'' ``b - c'' ``c - a'').

Notes. - Write the c\_i with the goal's own variable names.  di preserves them, so sos works before OR after di; skip di and sos peels the typed forall([... in rr]) wrapper itself. - Every generator (variable) must be certified in RR -- a goal binder x in rr  or an  (IN g RR)  assumption. - A wrong or insufficient certificate refuses cleanly: nothing is closed, you simply try other squares. - Strict (<) goals are refused -- a square can be 0, so squares alone never force a strict inequality.


\subsection*{Chained reasoning}

\par\smallskip{\raggedright\noindent\texttt{(calc L0 (rel1 L1 [just1]) (rel2 L2 [just2]) ...)}\quad\tactkind{composite}\index{calc@\texttt{calc} (tactic)}\par}\nobreak
\noindent Ground the focus goal (REL L0 Ln) by a CHAIN of intermediaries L0 rel1 L1 rel2 L2 ... reln Ln: proves each link and composes them into the endpoint.  A link no lane can close is LEFT OPEN as a leaf -- the refinement point, and the PSS candidate.  Relations: = == (folded by transitivity), < <= (folded through the co-*-trans lemmas, with = steps riding along), IFF (chained per direction).  RETURNS the links, the open ones, and their PSS candidates as an alist.\par
\noindent\textit{Uses:} every one of the 64 kernel operations.\par
\noindent\textit{When useful:} the goal is (REL L0 Ln) and you can WRITE the chain of intermediaries that gets there -- or you want the one link that will not close isolated as a PSS candidate\par

Write the argument the way you would on paper -- as a chain through intermediate terms -- and let the machine prove each link:

\begin{VNBsnip}
      Goal  k < 2*k   (with k in NN, ~(k=0))
      (calc 'k '(= (+ k 0)) '(< (+ k k)) '(= (* 2 k)))
          ;; i.e.  k = k+0 < k+k = 2*k
\end{VNBsnip}

Each step names a RELATION and the next LINE; calc cuts the link (rel L(i-1) Li), dispatches a lane at it, and on success composes the whole chain into the goal.  It is pure bookkeeping over the trusted tactics -- no kernel rule, no axiom of its own -- so a calc proof bills exactly the lemmas its lanes and composers cite.

The COMPOSER is forced by the relations you used, and there are three families: `cong' for = and ==, folded by eq-trans; `order' for < and <=, folded through the co-lt-trans / co-le-lt-trans / co-eq-lt-trans ... lemmas (a = step rewrites an endpoint of the running relation and rides along, which is why the chain above needs no separate arithmetic); and `iff', which di's the goal into its two directions and chains each with ai/detach!, because VNB cannot quantify over propositions and so has NO first-order iff-trans lemma to fold with.

The JUSTIFICATION of a link is optional and defaults to 'auto -- try the relation's lanes: crs then arith for = / ==; for < / <= the NN->RR bridge then ineq over the order-shaped premises; grind for IFF.  Otherwise pass 'scout (a small scout search), a macete/theorem NAME (mac it, then close), an explicit tactic form like '(fact 'nn-pos-of-nonzero 'k) eval'd at the link's focus, or 'open / \#f to leave the link open ON PURPOSE.

Sanity first: calc ERRORS before touching the proof if the goal's head is not the composed relation, or its LHS is not L0, or its RHS is not Ln.  A link already in the main context is taken as a GIVEN and skipped -- cutting it would be an alpha self-loop that opens no leaf.

What it is FOR is the open links.  A chain whose links are all closed is a proof; a chain with one open link has isolated your obstacle to a single formula, printed with its free variables and their context typing -- which is the PSS candidate to state as a lemma.  Refine by inserting more intermediaries until each link is discoverable.  (See also (prep 'ineq), which answers the other question: why a lane will not fire.)

\par\smallskip{\raggedright\noindent\texttt{(have! CLAIM [THUNK])}\quad\tactkind{composite}\index{have"!@\texttt{have"!} (tactic)}\par}\nobreak
\noindent Assert CLAIM as an intermediate step and carry on with it as a hypothesis: cut CLAIM, discharge the side goal from context (or by THUNK), and stay on the main branch.\par
\noindent\textit{Uses:} \texttt{and-intro} \texttt{arith-forsome} \texttt{arith-ground} \texttt{arith-simplify} \texttt{assumption} \texttt{cartesian-decompose} \texttt{cut} \texttt{detach} \texttt{eq-subst} \texttt{forall-elim} \texttt{forall-intro} \texttt{iff-intro} \texttt{implies-intro} \texttt{macete} \texttt{not-intro} \texttt{theorem-assumption} \texttt{tuple-equality-decompose}.\par
\noindent\textit{When useful:} you want to state an intermediate fact that follows immediately from context and continue with it -- the `we have X' step\par

The `we have X' step.  You state an intermediate fact CLAIM that follows immediately from what is already known; have! cuts it, proves the resulting side goal automatically, and returns you to the main branch with CLAIM now available as a hypothesis.  The automatic discharge (`from-context!') closes the everyday cases -- a conjunction, splitwise; an (IN (a*b) NN) by nn-mul-closed on the factors; a numeral membership by arith; otherwise the goal is already a context assumption up to alpha.  When the side goal needs more than that, pass a THUNK -- any tactic sequence -- as the second argument to discharge it your way.  It ERRORS, never silently no-ops, if CLAIM is already in context up to alpha: the cut would self-loop, opening one child and no main branch.  So a script reads as the argument does -- `we have q0*q0 = 3*(k*k); we have k =/= 0; ...'.  (Technically: composite -- cut then from-context! or THUNK; adds no kernel rule.)


\subsection*{Backchaining with a theorem}

\par\smallskip{\raggedright\noindent\texttt{(ta 'name)}\quad\tactkind{rule}\index{ta@\texttt{ta} (tactic)}\par}\nobreak
\noindent Theorem-assumption: bring the named installed theorem into context as an assumption.\par
\noindent\textit{Uses:} \texttt{theorem-assumption}.\par
\noindent\textit{When useful:} you want a library theorem available as a hypothesis\par

Bring an already-proved theorem into your current hypotheses so you can use it (instantiate it, detach from it, ...).  (Technically: adds the named installed theorem as an assumption.)

\par\smallskip{\raggedright\noindent\texttt{(wbc ['name])}\quad\tactkind{composite}\index{wbc@\texttt{wbc} (tactic)}\par}\nobreak
\noindent Witness-shape backchain: on an `exists v. ...' goal, cite a PSS lemma that MANUFACTURES a witness of that shape, so you can finish with inst+/grind/ew.\par
\noindent\textit{Uses:} \texttt{and-elim} \texttt{arith-forsome} \texttt{arith-ground} \texttt{arith-simplify} \texttt{cut} \texttt{detach} \texttt{forall-elim} \texttt{forsome-elim} \texttt{iff-elim} \texttt{not-elim} \texttt{or-elim} \texttt{theorem-assumption}.\par
\noindent\textit{When useful:} an EXISTENCE goal whose witness a PSS lemma manufactures (diagonalization / block-family)\par

The discovery move for existence proofs.  Faced with `there exists phi with ...', it looks up the produces-witness index for lemmas whose conclusion builds the RIGHT KIND of object -- the same existential-witness shape -- even when the rest of the statement differs (which is why plain bc* misses them: bc* demands the WHOLE conclusion unify).  E.g. on `exists phi. STRICTLY-MONO-NN(phi) and IS-CAUCHY-SEQ(SUBSEQ f phi)' it reaches for `diagonalization' (which makes a STRICTLY-MONO-NN phi from a nested family); on a nested-family goal it reaches for `block-family'.  With no argument it cites the top retrieved producer; (wbc 'name) cites the one you name.  It brings the lemma in (via ta); you then discharge its hypotheses (inst+), skolemize its existential (grind), and supply its witness to your goal (ew).  Circular `headline' lemmas that already conclude your exact goal are filtered out.  (See (witness-producers goal) for the raw retrieval; what-now lists the (wbc 'name) candidates on any existential goal.)

\par\smallskip{\raggedright\noindent\texttt{(bc impl)}\quad\tactkind{rule}\index{bc@\texttt{bc} (tactic)}\par}\nobreak
\noindent Backchain the goal through an (IMPLIES A B) already in context: if the goal matches B, the new goal is A.\par
\noindent\textit{Uses:} \texttt{backchain}.\par
\noindent\textit{When useful:} you have (P => Q) in context and the goal is Q\par

Work backwards through an implication you already have.  If `P implies Q' is among your hypotheses and your goal is Q, this reduces the goal to proving P.  The everyday `to get Q it's enough to show P'.  (Technically: backchaining the goal through an in-context implication whose conclusion matches.)

\par\smallskip{\raggedright\noindent\texttt{(bc* 'thm [((v val)...)] h1 h2 ...)}\quad\tactkind{composite}\index{bc*@\texttt{bc*} (tactic)}\par}\nobreak
\noindent Matching backchain: unify the theorem's conclusion against the goal, then leave its (instantiated) antecedents as subgoals -- optional handlers hk run on the k-th subgoal.\par
\noindent\textit{Uses:} \texttt{assumption} \texttt{backchain} \texttt{cut} \texttt{forall-elim} \texttt{theorem-assumption}.\par
\noindent\textit{When useful:} a library theorem's CONCLUSION matches your goal\par

Apply a known theorem to your goal.  If you have a theorem `if A and B then C' and your goal is its conclusion C -- matching the theorem's variables to yours, and the names of any dummy/bound variables don't matter (`there exists a net N' applies to a goal `there exists a net F') -- this replaces `prove the goal' with `prove A' and `prove B', the theorem's hypotheses.  The everyday `by Theorem X it suffices to show A and B'.  You may attach a handler to each hypothesis to dispatch it; values the match can't determine you supply as ((v val) ...).  (Technically: peels the theorem's leading universals and implications, matches the conclusion -- alpha-aware on bound variables -- and replays ta/inst/cut/bc automatically.)


\subsection*{Lambda, comprehension \& description}

\par\smallskip{\raggedright\noindent\texttt{(lam-t)}\quad\tactkind{rule}\index{lam-t@\texttt{lam-t} (tactic)}\par}\nobreak
\noindent VNB-LAMBDA typing: reduce (IN (VNB-LAMBDA ...) (FUN A B)) to TWO obligations -- the body, and that the domain is a set.\par
\noindent\textit{Uses:} \texttt{lambda-type}.\par
\noindent\textit{When useful:} the goal types a VNB-LAMBDA into FUN(A,B)\par

Show that a function defined by a formula (`x |-> ...') maps A into B -- reduces to showing that, for an arbitrary input in A, the value lies in B, AND that A is a set.  The second is not a formality: carrying the domain in the term stops one lambda from typing into FUN(A,B) for every A, but says nothing about A being a set, and a lambda over a proper class is not a function.  The lambda's declared domain must also be the FUN's, or the rule does not apply at all.  (Technically: VNB-LAMBDA typing into FUN A B.)

\par\smallskip{\raggedright\noindent\texttt{(lam-b)}\quad\tactkind{rule}\index{lam-b@\texttt{lam-b} (tactic)}\par}\nobreak
\noindent VNB-LAMBDA beta: reduce an applied lambda to its substituted body.\par
\noindent\textit{Uses:} \texttt{lambda-beta}.\par
\noindent\textit{When useful:} the goal has a VNB-LAMBDA applied to an argument\par

Simplify a function `x |-> e(x)' applied to an argument a to e(a) -- the body with the argument substituted in.  A lambda carries its DOMAIN, and it is only defined on that domain, so the reduction is licensed only where a is known to be in it: if the membership is neither in the context nor supplied by an enclosing binder, the step still fires but leaves (IN a A) as an extra subgoal.  Land the typing fact BEFORE the lam-b and no subgoal appears.  (Technically: VNB-LAMBDA beta-reduction.)

\par\smallskip{\raggedright\noindent\texttt{(lam-b-h hyp)}\quad\tactkind{rule}\index{lam-b-h@\texttt{lam-b-h} (tactic)}\par}\nobreak
\noindent VNB-LAMBDA beta in a cited ASSUMPTION -- what mac-h is to mac.\par
\noindent\textit{Uses:} \texttt{lambda-beta-hyp}.\par

Beta-reduce an applied lambda inside a hypothesis, in place.  Needed because a `fact' that instantiates a theorem's function variable at a lambda lands the APPLIED lambda in the CONTEXT, where the goal-side `lam-b' cannot reach it: union-of-opens-open at the identity family g := x |-> x lands is-open(md, big-union(i, fam, (x |-> x)(i))).  Without this the proof must detour through a cut beta-equation and a subst.  Cites nothing, so it adds no debt.

\par\smallskip{\raggedright\noindent\texttt{(sep-set)}\quad\tactkind{rule}\index{sep-set@\texttt{sep-set} (tactic)}\par}\nobreak
\noindent Separation sethood: the separation set \{x in A | p\} is a set.\par
\noindent\textit{Uses:} \texttt{sep-sethood}.\par
\noindent\textit{When useful:} the goal asserts a separation set \{ x in A | p \} is a set\par

Show that a set-builder set \{x in A | p(x)\} is genuinely a set.  (Technically: separation sethood -- a subclass of a set is a set.)

\par\smallskip{\raggedright\noindent\texttt{(sep-mi)}\quad\tactkind{rule}\index{sep-mi@\texttt{sep-mi} (tactic)}\par}\nobreak
\noindent Separation membership intro: prove (IN t \{x in A | p\}).\par
\noindent\textit{Uses:} \texttt{sep-mem-intro}.\par
\noindent\textit{When useful:} the goal is membership in a separation set \{ x in A | p \}\par

Prove that a term t belongs to \{x in A | p(x)\} by showing t lies in A and satisfies the condition p.  (Technically: separation-membership introduction.)

\par\smallskip{\raggedright\noindent\texttt{(sep-me hyp)}\quad\tactkind{rule}\index{sep-me@\texttt{sep-me} (tactic)}\par}\nobreak
\noindent Separation membership elim: split a separation-membership assumption into A-membership and the predicate.\par
\noindent\textit{Uses:} \texttt{sep-mem-elim}.\par
\noindent\textit{When useful:} a hypothesis is separation-set membership -- unpack it\par

From a hypothesis that t belongs to \{x in A | p(x)\}, extract the two facts it packs: t lies in A, and p(t) holds.  (Technically: separation-membership elimination.)

\par\smallskip{\raggedright\noindent\texttt{(comp-mi)}\quad\tactkind{rule}\index{comp-mi@\texttt{comp-mi} (tactic)}\par}\nobreak
\noindent Comprehension membership intro.\par
\noindent\textit{Uses:} \texttt{comp-mem-intro}.\par
\noindent\textit{When useful:} the goal is comprehension-set membership\par

Prove that something belongs to a comprehension set by establishing the set's defining condition for it.  (Technically: comprehension-membership introduction.)

\par\smallskip{\raggedright\noindent\texttt{(comp-me hyp)}\quad\tactkind{rule}\index{comp-me@\texttt{comp-me} (tactic)}\par}\nobreak
\noindent Comprehension membership elim.\par
\noindent\textit{Uses:} \texttt{comp-mem-elim}.\par
\noindent\textit{When useful:} a hypothesis is comprehension-set membership\par

From a comprehension-membership hypothesis, extract its defining condition.  (Technically: comprehension-membership elimination.)

\par\smallskip{\raggedright\noindent\texttt{(iota-d term)}\quad\tactkind{rule}\index{iota-d@\texttt{iota-d} (tactic)}\par}\nobreak
\noindent Definite-description: discharge the IOTA uniqueness obligation for `term'.\par
\noindent\textit{Uses:} \texttt{iota-def}.\par
\noindent\textit{When useful:} the goal involves IOTA and you must discharge its uniqueness obligation\par

Justify `the unique x such that p(x)' (a definite description) by proving that exactly one such x exists.  (Technically: IOTA -- discharge the uniqueness obligation for the described term.)

\par\smallskip{\raggedright\noindent\texttt{(bu-set)}\quad\tactkind{rule}\index{bu-set@\texttt{bu-set} (tactic)}\par}\nobreak
\noindent Big-union sethood.\par
\noindent\textit{Uses:} \texttt{big-union-sethood}.\par
\noindent\textit{When useful:} the goal asserts a BIG-UNION is a set\par

Show that a big union (the union of an indexed family of sets) is itself a set.  (Technically: big-union sethood.)

\par\smallskip{\raggedright\noindent\texttt{(bu-mi w)}\quad\tactkind{rule}\index{bu-mi@\texttt{bu-mi} (tactic)}\par}\nobreak
\noindent Big-union membership intro via the index witness w.\par
\noindent\textit{Uses:} \texttt{big-union-mem-intro}.\par
\noindent\textit{When useful:} the goal is BIG-UNION membership -- you have the index witness\par

Prove that an element lies in a big union by exhibiting one index w whose set already contains it.  (Technically: big-union-membership introduction via the index witness.)

\par\smallskip{\raggedright\noindent\texttt{(bu-me hyp)}\quad\tactkind{rule}\index{bu-me@\texttt{bu-me} (tactic)}\par}\nobreak
\noindent Big-union membership elim.\par
\noindent\textit{Uses:} \texttt{big-union-mem-elim}.\par
\noindent\textit{When useful:} a hypothesis is BIG-UNION membership -- obtain its index\par

From a hypothesis that an element lies in a big union, obtain an index whose set contains it (a fresh name for that index).  (Technically: big-union-membership elimination.)


\subsection*{Conditional terms}

\par\smallskip{\raggedright\noindent\texttt{(if-true t)}\quad\tactkind{rule}\index{if-true@\texttt{if-true} (tactic)}\par}\nobreak
\noindent Reduce an (IF p a b) term on the p branch: spawns p as a subgoal; the continuation gains (= (IF p a b) a) when a is certified defined, else (== (IF p a b) a) (since 2026-10-08).\par
\noindent\textit{Uses:} \texttt{if-true}.\par
\noindent\textit{When useful:} the goal has an (IF p a b) term to evaluate on the p branch\par

Evaluate a conditional term `if p then a else b' on the assumption that p holds: you take on p as a side goal, and may then use that the conditional equals a.  (Technically: if-true reduction, spawning p.)

\par\smallskip{\raggedright\noindent\texttt{(if-false t)}\quad\tactkind{rule}\index{if-false@\texttt{if-false} (tactic)}\par}\nobreak
\noindent Reduce an (IF p a b) term on the not-p branch: spawns (NOT p); the continuation gains (= (IF p a b) b).\par
\noindent\textit{Uses:} \texttt{if-false}.\par
\noindent\textit{When useful:} the goal has an (IF p a b) term to evaluate on the not-p branch\par

Evaluate a conditional term `if p then a else b' on the assumption that p fails: you take on `not p' as a side goal, and may then use that the conditional equals b.  (Technically: if-false reduction, spawning not p.)


\subsection*{Navigation, display \& tacticals}

\par\smallskip{\raggedright\noindent\texttt{(focus n)}\index{focus@\texttt{focus} (tactic)}\par}\nobreak
\noindent Switch the focus to the n-th open goal (1-based, the POSITION in the list (show) prints: the focus first, then the others).  A position shifts as leaves close; the node number in brackets does not -- see focus-id.\par
\noindent\textit{Uses:} no kernel operation: it records no inference.\par
\par\smallskip{\raggedright\noindent\texttt{(focus-id k)}\index{focus-id@\texttt{focus-id} (tactic)}\par}\nobreak
\noindent Switch the focus to the ungrounded node whose NODE NUMBER is k, the number (show) and the Focus Workspace print in brackets before each sequent.  Any ungrounded node, not only a leaf: a driver records a leaf by position as (focus n) and an ungrounded node a rule has already fired on as (focus-id k), so a page may name either.  Recorded as a step, undone by backup-one.  The move to make when reading a goal off the panel, e.g. the side leaf of a cut: (cut L) opens  G |- L  (the side leaf) and  G, L |- goal  (the main branch), and a composite such as zero-it leaves the focus on whichever it did not close.\par
\noindent\textit{Uses:} no kernel operation: it records no inference.\par
\par\smallskip{\raggedright\noindent\texttt{(backup-one)}\quad\tactkind{meta}\index{backup-one@\texttt{backup-one} (tactic)}\par}\nobreak
\noindent Take back the last recorded command: restore the goal, the deduction graph and the proof script to the state before it.  Repeat to walk further back; (sp) clears the stack.  `undo' is an alias.\par
\noindent\textit{Uses:} no kernel operation: it records no inference.\par
\noindent\textit{When useful:} the last command was a mistake -- a greedy `di' that ate the induction, a `cut' you did not mean, a branch you want back\par

The one undo in VNB.  It restores the deduction graph -- not merely the focus -- by rolling back the journalled arrow and grounding writes and DROPPING every sequent and inference node posted since, so no node from the abandoned branch survives to be counted as an open leaf and `qed' cannot be handed a phantom obligation.  *proof-script* and the live trace are rewound with it, so the saved script and the printed proof are the proof you actually kept.  It backs up over the LAST RECORDED command: a command that changed nothing was never recorded and is not a step to back up over.  One thing is deliberately not restored -- *fresh-counter*, the eigenvariable source -- so backing up over an `ai' or `ew' that minted `u\_4' and re-running it mints `u\_5'; the proof is the same, the witness has a different name.  Depth is capped at *vnb-undo-depth* (64).

\par\smallskip{\raggedright\noindent\texttt{(undo)}\quad\tactkind{meta}\index{undo@\texttt{undo} (tactic)}\par}\nobreak
\noindent Alias for (backup-one).\par
\noindent\textit{Uses:} no kernel operation: it records no inference.\par
\noindent\textit{When useful:} the last command was a mistake -- alias for backup-one\par
\par\smallskip{\raggedright\noindent\texttt{(show)}\index{show@\texttt{show} (tactic)}\par}\nobreak
\noindent Redisplay the current proof state.\par
\noindent\textit{Uses:} no kernel operation: it records no inference.\par
\par\smallskip{\raggedright\noindent\texttt{(show-graph)}\index{show-graph@\texttt{show-graph} (tactic)}\par}\nobreak
\noindent Print EVERY node of the current proof, closed ones included: one line per sequent node in creation order (the bracketed number the panel prints, the sequent, [GROUNDED] when justified), and under each node one line per inference into it -- the rule and the numbers of the hypothesis nodes it rests on.  The focus and the open leaves are marked.  This is where a branch a composite closed is found: after (cut L), the side leaf G |- L and the main branch G, L |- goal are both nodes, and zero-it on an existential closes the main one.  Records nothing; the [Graph] button of the Focus Workspace shows it.\par
\par\smallskip{\raggedright\noindent\texttt{(dial notch)}\index{dial@\texttt{dial} (tactic)}\par}\nobreak
\noindent The PRESENTATION RESOLUTION DIAL: how much of the sequent to show.  r0 kernel S-expression, r1 surface syntax (the default, and the last INVERTIBLE notch -- what is printed re-parses to what was printed), r2 guards and typing hypotheses elided, r3 also structures destructured to their declared slot names and deep subterms elided.  (dial) reports the notch, (dial n) sets it, (dial 'auto) reads the current goal and suggests one, (dial 'why) shows the score behind the suggestion.  For ONE FORMULA rather than the proof state, (dial-wff x [n]) -- x a wff, a ``surface string'' or an S-expression, n defaulting to the dial's current notch: (dial-wff (lookup-theorem 'matmul-assoc) 3).\par
\noindent\textit{Uses:} no kernel operation: it records no inference.\par

It is a FILTER, not a set of modes, and two constraints make it one.  The rungs NEST: turning it up only ever SUBTRACTS, so everything readable at a higher notch is readable below it -- checked in the suite, atom by atom.  And the APERTURE IS ON SCREEN: every notch above r1 names itself on the line and counts what it hid, because a lossy render that does not announce its own lossiness is the defect class of `(f)' printing as `f'.  The auto-notch SUGGESTS a starting notch and then STAYS PUT: it is consulted when you ask, never per sequent, so the lens cannot move while you are reading a proof.

\par\smallskip{\raggedright\noindent\texttt{(dial-wff x notch)}\index{dial-wff@\texttt{dial-wff} (tactic)}\par}\nobreak
\noindent Present ONE FORMULA at a resolution notch, instead of the current proof state: x is a wff, a ``surface string'' or an S-expression, and notch defaults to the dial's current setting.  (dial-wff (lookup-theorem 'matmul-assoc) 3).  Same rungs, same aperture line, no proof required.\par
\noindent\textit{Uses:} no kernel operation: it records no inference.\par

`sequent->string' is the display path for a proof STATE, so the dial reaches a formula only while a proof is holding it.  A theorem looked up, a statement being drafted, a hypothesis pasted from somewhere wants the same dial, and this is it.

\par\smallskip{\raggedright\noindent\texttt{(goal-status)}\index{goal-status@\texttt{goal-status} (tactic)}\par}\nobreak
\noindent One-line summary: done / N open goals.\par
\noindent\textit{Uses:} no kernel operation: it records no inference.\par
\par\smallskip{\raggedright\noindent\texttt{(repeat thunk [cap])}\index{repeat@\texttt{repeat} (tactic)}\par}\nobreak
\noindent Run thunk until it stops changing the proof state (LCF REPEAT).\par
\noindent\textit{Uses:} no kernel operation: it records no inference.\par
\par\smallskip{\raggedright\noindent\texttt{(orelse t1 t2 ...)}\index{orelse@\texttt{orelse} (tactic)}\par}\nobreak
\noindent Run thunks in order, stop at the first that makes progress (LCF ORELSE).\par
\noindent\textit{Uses:} no kernel operation: it records no inference.\par
\par\smallskip{\raggedright\noindent\texttt{(preamble)  |  (preamble 'NAME)  |  (preamble ``file.pre'')}\quad\tactkind{composite}\index{preamble@\texttt{preamble} (tactic)}\par}\nobreak
\noindent Apply an editable RULE FILE to the focus goal, one step at a time: each rule reads `if the sequent looks like this, do that', and the first rule that matches AND makes progress is kept; the loop repeats on the new focus until the goal is done, no rule applies, or 40 steps.  Rules come from \textasciitilde{}/.vnb-preamble.pre when it exists, else the shipped preambles/default.pre (peel, split, close by context, reflexivity, arithmetic, type-term, zero-it, ineq, beta, unfold, and a tentative cut).  Prints one report: what each rule did, which rules matched but were rejected, and where it stopped and why.\par
\noindent\textit{Uses:} \texttt{cut}.\par
\noindent\textit{When useful:} you want the routine steps done for you by rules you can read and edit, with a report of what each rule did and where it stopped\par

The notes-45 `giant COND' as data you can edit.  A rule is (rule NAME (goal PATTERN) (with PATTERN) (guard ...) (do FORM) (probe ...)): patterns are raw S-expressions with schema variables ?x, matched against the goal and the assumptions (bindings shared, a ?x in binder position binds the bound variable's name); guards are typed, head, closed, not-in-context, one-of, occurs, unfolds; the action is an ordinary surface command with the variables filled in.  A rule is TRIED on a scratch copy first and KEPT only when its probe holds there and again on the proof (grounded: the goal closed; progress, the default: something changed and the open goals did not multiply; changed; (lands F)); otherwise nothing it did is kept.  A rule whose action is (cut F) is a TENTATIVE CUT: the rules (without their cuts) are run on F, and F is proven in place or left open and reported as OWED -- a support is never installed.  Each kept step is recorded AS ITSELF, so the page of the proof is the ordinary steps and replays without any rule file.  A malformed rule is reported by name and left out; a rule that raises is reported and skipped.  The value is a verdict: status (done, owed, stalled, cap, error), the steps, the rules fired and rejected, the owed claims, the report text.  (preamble '(induct) ...) with CLAUSE LISTS is the older strategy pipeline, also (preamble-clauses ...).  (Technically: composite -- no rule of its own; every action is kernel-checked; transactional per firing.)

\par\smallskip{\raggedright\noindent\texttt{(quietly thunk)}\index{quietly@\texttt{quietly} (tactic)}\par}\nobreak
\noindent Run thunk with state-dump output suppressed; returns its value.\par
\noindent\textit{Uses:} no kernel operation: it records no inference.\par

\subsection*{Choosing and naming witnesses}

\par\smallskip{\raggedright\noindent\texttt{(minimize! '(v1 ... vk) GUARD MEASURE)}\quad\tactkind{composite}\index{minimize"!@\texttt{minimize"!} (tactic)}\par}\nobreak
\noindent Choose v1..vk satisfying GUARD with the NN-valued MEASURE as small as possible: lands the witnesses and their minimality, and opens the two obligations any minimisation owes (MEASURE lands in NN; GUARD is satisfiable).\par
\noindent\textit{Uses:} \texttt{and-elim} \texttt{and-intro} \texttt{arith-forsome} \texttt{arith-ground} \texttt{arith-simplify} \texttt{assumption} \texttt{cartesian-decompose} \texttt{cut} \texttt{detach} \texttt{eq-subst} \texttt{forall-elim} \texttt{forall-intro} \texttt{forsome-elim} \texttt{forsome-intro} \texttt{iff-elim} \texttt{iff-intro} \texttt{implies-intro} \texttt{macete} \texttt{not-elim} \texttt{not-intro} \texttt{or-elim} \texttt{reflexivity} \texttt{sep-mem-elim} \texttt{sep-mem-intro} \texttt{theorem-assumption} \texttt{tuple-equality-decompose}.\par
\noindent\textit{When useful:} the goal falls to a `least such' / minimal-counterexample argument -- descent proofs (sqrt 2, sqrt 3 irrational), least-degree or least-pivot witnesses\par

The `least such' / minimal-counterexample step, mechanised.  You give three things: the variables to choose, a GUARD formula they must satisfy, and a natural-number-valued MEASURE term to make small.  On the main branch it hands you fresh eigenconstants w1..wk with two hypotheses -- GUARD holds of them, and nothing satisfying GUARD has a strictly smaller MEASURE -- and it opens exactly the two side goals a minimisation genuinely owes: that MEASURE lands in NN wherever GUARD holds (the TYPE goal), and that GUARD holds somewhere (the NONEMPTY goal).  Everything in between -- forming the value set, well-ordering it, unpacking the least witness, and restating minimality in terms of your variables rather than the set -- is done for you.  Every vi must occur in MEASURE (a variable the measure ignores is not one you are minimising over).  Because it quantifies over the vk at the META level, no lambda and no FUN(U,NN) typing obligation ever enters the logic -- it takes a formula and a term, which is what a driver has in hand, and never forms the set U at all.  Returns (list (w1 ... wk) TYPE-node NONEMPTY-node); either node is \#f when that obligation was already in context up to alpha, so test before refocusing.  (Technically: composite -- it drives cut / di / ai / ew / sep-mi / sep-me / fact / inst / detach! / subst / ass / rfl and adds no kernel rule; its one mathematical appeal is `nn-least-element', the well-ordering of NN, proven from ord-well-ordered.  NO choice principle is used: well-ordering returns a MEMBER of a separation set, and sep-me recovers the witness.)

\par\smallskip{\raggedright\noindent\texttt{(obtain LANE)}\quad\tactkind{composite}\index{obtain@\texttt{obtain} (tactic)}\par}\nobreak
\noindent Run LANE -- a thunk whose forward step lands a `there exists' into context -- then eliminate that existential and RETURN the fresh witness's name, read off the proof state rather than guessed.\par
\noindent\textit{Uses:} every one of the 64 kernel operations.\par
\noindent\textit{When useful:} a forward step yields `there exists ...' and you want to name the witness for later use, without guessing the engine's eigenvariable\par

The `let w be such an x' step, with the eigenvariable named for you.  LANE is a zero-argument procedure whose effect is to land some `there exists v. P(v)' among your hypotheses -- e.g. (lambda () (fact 'nn-3-div-square w)).  obtain runs it, spots the existential that newly appeared, eliminates it (introducing a fresh eigenvariable and landing P at that witness), unpacks any conjunction in the body, and hands back the eigenvariable's NAME -- identified by free-variable set-difference (the symbol now in context that was not there before), so the driver never guesses what the engine called it.  This is the robust way to name a descended witness: (define k (obtain (lambda () ...))), then use k.  Returns \#f (with a note) if the lane landed no existential.  (Technically: composite -- a guarded forward discharge followed by forsome-elim (ai) and AND-splitting.  Plain existential elimination, so it owes nothing and uses no choice.)

\par\smallskip{\raggedright\noindent\texttt{(vlet (n1 ...) FORMER)}\quad\tactkind{composite}\index{vlet@\texttt{vlet} (tactic)}\par}\nobreak
\noindent Bind proof-object names from the current state.  FORMER is (match PATTERN) -- bind each name to its slot in a context formula matching PATTERN -- or (choice [v body]) -- eliminate an in-context existential and bind its witness.\par
\noindent\textit{Uses:} \texttt{and-elim} \texttt{cut} \texttt{forsome-elim} \texttt{iff-elim} \texttt{not-elim} \texttt{or-elim}.\par
\noindent\textit{When useful:} you need to NAME a witness or a matched subterm from the proof state, rather than navigate to it by shape\par

A binding form for the pieces of a proof, so a driver NAMES what it needs instead of navigating to it by shape.  Two FORMERs.  (match PATTERN): the listed names ARE the holes (wildcards) in PATTERN; everything else is literal.  vlet searches the context for a formula or subterm matching PATTERN and binds each name to what filled its slot -- pure selection, no proof step, no obligation; a hard error if nothing matches (you named a piece of something absent).  (choice): eliminate the sole existential in context and bind its witness.  (choice v body): present-else-debt -- if `there exists v. body' is already in context up to alpha, eliminate it; otherwise cut it, LEAVE the existence side-goal open as a debt leaf, and eliminate on the main branch -- either way the witness is bound and body[witness] is landed.  vlet's (choice) and `obtain' are the same underlying mechanism -- eliminate an existential, name the witness by free-variable difference -- differing only in what they take: obtain runs a lane and names what it just produced, vlet names an existential already (or, with a body, about to be) in context.  No choice AXIOM is used; a single witness is plain existential-elimination.  (Technically: a define-syntax expanding to (define n ...) over vlet--match / vlet--choice!; composite, no kernel rule beyond forsome-elim and, for present-else-debt, cut.)




%% =========================================================
\section{Focus management}

After each command the focus advances automatically to the first
ungrounded subgoal just introduced.  If there are no new subgoals, it
moves to any remaining open goal.

\begin{remark}
Every proof command operates exclusively on the focus node; the other
open goals are untouched.  There can be several open goals simultaneously
(for example, after \texttt{(ci)} on a conjunction goal, or after
\texttt{(cut f)} which always creates two subgoals).  \texttt{(show)}
displays them all, numbering the non-focus goals so you can use those
numbers with \texttt{(focus $n$)}.  Until you call \texttt{(focus $n$)},
the system always works on whichever goal it chose automatically.
\end{remark}

In the interactive session use:
\begin{VNBcode}
(show)       ; list all open goals (numbered) and the current focus
(focus n)    ; switch focus to the n-th open goal (1-based)
\end{VNBcode}

The underlying API (for use outside the interactive session) is:
\begin{VNBcode}
(focus-on ps sqn)            ; focus on a specific sequent node
(focus-on-first-open ps)     ; focus on the first open goal
(proof-open-goals ps)        ; list all open goals
\end{VNBcode}

%% =========================================================
\section{Macetes}
\label{sec:macetes}

A \emph{macete} is a named, composable rewrite strategy.  When a
theorem is installed, its macete is automatically derived by extracting a
source pattern and replacement pattern from the theorem's logical form:

\begin{center}
\begin{tabular}{ll}
\toprule
Theorem form & Rewrite \\
\midrule
$\forall\bar{x}.\;(s = r)$ & $s \mapsto r$ \\
$\forall\bar{x}.\;(s \Leftrightarrow r)$ & $s \mapsto r$ \\
$\forall\bar{x}.\;(\text{conds} \Rightarrow (s = r))$ & $s \mapsto r$ when conds hold \\
$\forall\bar{x}.\;p$ & $p \mapsto \top$ \\
\bottomrule
\end{tabular}
\end{center}

The quantifiers need not all be leading.  A \texttt{forall} occurring in
positive position --- under an \texttt{implies} consequent or inside an
\texttt{and} --- is pulled to the front before the source and replacement
patterns are extracted, so a theorem such as
\texttt{forall([a,b], $H$ implies forall(x, $s$ iff $r$))} still yields the
rewrite $s \mapsto r$ (conditional on $H$).

Macetes can be composed:
\begin{VNBcode}
(macete-series m1 m2 m3)   ; apply in sequence
(macete-repeat m)           ; apply until no change
(macete-try m)              ; apply m, succeed even if m fails
(macete-choice m1 m2 m3)   ; first one that applies
\end{VNBcode}

To install a custom macete:
\begin{VNBcode}
(install-theorem! 'eq-symm
  (make-wff-from-string "forall([x, y], x = y iff y = x)"))

; Use it:
(cmd-apply-macete ps 'eq-symm)
\end{VNBcode}

\subsection{Applying macetes to wffs outside a proof}

\begin{sloppypar}The function \texttt{apply-macete} applies a named macete directly to a
\texttt{wff} without creating a proof state:\end{sloppypar}

\begin{itemize}
\item \texttt{(apply-macete \textit{name} \textit{w})} --- rewrites \textit{w};
  returns a new \texttt{wff} if any rewrite fires, or \texttt{\#f} if the
  pattern matches nowhere in \textit{w}.
\end{itemize}

This is purely syntactic: no deduction graph is involved and no
theorem is proved.  It is useful for inspecting rewrite behavior,
preprocessing formulas, or building higher-level tools outside of proofs.

\paragraph{Unconditional equation macete.}
\begin{VNBcode}
(install-theorem! 'nn-zero-succ
  (make-wff-from-string "forall([n in NN], succ(n) = n + 1)"))

(apply-macete 'nn-zero-succ (make-wff-from-string "succ(3) = x"))
; => wff for "3 + 1 = x"

(apply-macete 'nn-zero-succ (make-wff-from-string "3 in nn"))
; => #f               -- macete doesn't fire
\end{VNBcode}

\paragraph{Biconditional (\texttt{iff}) macete.}
\begin{VNBcode}
(install-theorem! 'eq-symm
  (make-wff-from-string "forall([x, y], x = y iff y = x)"))

(apply-macete 'eq-symm (make-wff-from-string "n + m = 0"))
; => wff for "0 = n + m"

(apply-macete 'eq-symm (make-wff-from-string "n = m implies truth"))
; => wff for "m = n implies truth"

(apply-macete 'eq-symm (make-wff-from-string "n in nn"))
; => #f
\end{VNBcode}

\paragraph{Closure axiom macete (conditional).}
\begin{VNBcode}
;; nn-add-closed: forall([n in NN, m in NN], n + m in NN)
;; Condition (n in NN and m in NN) must be present in local context.

(apply-macete 'nn-add-closed
  (make-wff-from-string "n in nn and m in nn implies n + m in nn"))
; => wff for "n in nn and m in nn implies truth"

(apply-macete 'nn-add-closed (make-wff-from-string "n + m in nn"))
; => #f  -- conditions not in context; rewrite correctly suppressed
\end{VNBcode}

\begin{remark}
\texttt{apply-macete} applies the rewrite depth-first throughout the
formula (outermost first within each subexpression).  Conditions of a
conditional macete are checked against the \emph{local context}
accumulated as the rewriter descends: when descending into the
consequent of \texttt{implies}, the antecedent is added to the local
context; when descending into the right conjunct of \texttt{and}, the
left conjunct is added; and so on (Monk~\cite{Monk1988}).  A condition that
matches a formula already present in the local context is considered
discharged; if any condition is unmet at a candidate site, the rewrite
is suppressed at that site.
\end{remark}

%% =========================================================
\section{Example proofs}

\subsection{Identity}

\begin{VNBcode}
(sp (make-wff-from-string "x in NN implies x in NN"))
(di)    ; assume (x in nn), prove (x in nn)
(ass)   ; it is in context
; => Proof complete.
\end{VNBcode}

\subsection{\texttt{and} elimination}

\begin{VNBcode}
(sp (make-wff-from-string "n in NN and m in NN implies n in NN"))
(di)                          ; assume (n in nn) and (m in nn)
(ai "n in NN and m in NN")   ; split into (n in nn) and (m in nn)
(ass)                         ; prove (n in nn)
\end{VNBcode}

\subsection{VNB axiom instance}

Proving \texttt{"a in b implies a in set"} from the base axiom:

\begin{VNBcode}
(sp (make-wff-from-string "a in b implies a in SET"))
(di)
(ta 'membership-implies-sethood)
(inst "forall([a, b], a in b implies a in set)" 'a)
(inst "forall([b], a in b implies a in set)" 'b)
(bc "a in b implies a in set")
(ass)
\end{VNBcode}

\subsection{Double negation (classical)}

\begin{VNBcode}
(sp (make-wff-from-string "not(not(x in NN)) implies x in NN"))
(di)                              ; assume not(not((x in nn))), prove (x in nn)
(pbc)                             ; assume not((x in nn)), prove falsity
(ai "not(not(x in NN))")         ; not(not(P)) and not(P) yield falsity
\end{VNBcode}

\subsection{Ordered tuples and Cartesian products}

\begin{VNBcode}
; From x in A, y in B, z in C prove [x,y,z] in cartesian(A,B,C)
(sp (make-wff-from-string
      "x in A and y in B and z in C implies [x, y, z] in cartesian(A, B, C)"))
(di)
(ai "x in A and y in B and z in C")
(ai "y in B and z in C")
(ci)              ; cartesian-intro generates 3 subgoals
(ass) (ass) (ass)

; NTH projection: nth(2, [x, y, z]) = y
(sp (make-wff-from-string "nth(2, [x, y, z]) = y"))
(nth-r)           ; reduces goal to (y = y)
(rfl)             ; closes by reflexivity
\end{VNBcode}

\subsection{Mathematical structures}

Algebraic structures are declared with \texttt{declare-structure}.
No quoting is needed: carrier names and operation signatures are written
directly.

\begin{VNBcode}
(declare-structure GROUP
  (carriers CARR)
  (op OPR (CARTESIAN CARR CARR) CARR)  ; operation: CARR x CARR -> CARR
  (constant IDEN CARR)                 ; identity: IDEN in CARR
  (op INV CARR CARR))                  ; inverse: CARR -> CARR
; Accessor indices: CARR -> 1, OPR -> 2, IDEN -> 3, INV -> 4
; (the group laws are folded in by additional (property ...) clauses)
\end{VNBcode}

\begin{sloppypar}The \texttt{(op \textit{name} \textit{domain} \textit{range})} clause declares
$\textit{name}(s) \in \texttt{fun}(\textit{domain}, \textit{range})$.
Higher-arity operations write the Cartesian product explicitly:
\texttt{(op opr (cartesian carr carr) carr)} for a binary operation.
A \texttt{(constant \textit{name} S)} clause declares an element of set $S$.
Accessor indices are assigned sequentially: carriers first (starting
at~1), then operations in declaration order.\end{sloppypar}

A \texttt{declare-structure} call installs two things automatically:
\begin{itemize}
\item \emph{Accessor macetes}: \texttt{carr}$(s)$ rewrites to
  \texttt{nth(1,$s$)}, \texttt{opr}$(s)$ to \texttt{nth(2,$s$)}, etc.
\item \emph{Definitional biconditional} \texttt{is-\textit{name}}: the axiom
  $\forall s.\;\texttt{is-group}(s) \Leftrightarrow
   (\texttt{carr}(s)\in\mathrm{set} \wedge \texttt{opr}(s)\in\cdots)$.
\end{itemize}

Given \texttt{is-group(g)} in context, the \texttt{is-group} macete
unfolds to:
\begin{VNBsnip}
CARR(g) in SET
and OPR(g) in fun(cartesian(CARR(g), CARR(g)), CARR(g))
and IDEN(g) in CARR(g)
and INV(g) in fun(CARR(g), CARR(g))
\end{VNBsnip}

\paragraph{Pre-loaded structures.}
The following structures are defined in \texttt{structure-library/} (one
file per structure) and loaded automatically.

\begin{center}\scriptsize
\begin{tabular}{@{}p{2.3cm}p{4.3cm}p{5.3cm}@{}}
\toprule
\textbf{Structure} & \textbf{Accessor slots (index)} & \textbf{Characteristic axioms} \\
\midrule
\texttt{semigroup} &
  \texttt{carr}\,(1), \texttt{opr}\,(2) &
  \texttt{semigroup-assoc} \\[4pt]
\texttt{monoid} &
  \texttt{carr}\,(1), \texttt{opr}\,(2), \texttt{iden}\,(3) &
  \texttt{monoid-assoc},
  \texttt{monoid-left-id},
  \texttt{monoid-right-id} \\[4pt]
\texttt{group} &
  \texttt{carr}\,(1), \texttt{opr}\,(2), \texttt{iden}\,(3), \texttt{inv}\,(4) &
  \texttt{group-assoc},
  \texttt{group-left-id},
  \texttt{group-left-inv} \\[4pt]
\texttt{ring} &
  \texttt{carr}\,(1), \texttt{add}\,(2), \texttt{mul}\,(3),
  \texttt{neg}\,(4), \texttt{zero}\,(5), \texttt{one}\,(6) &
  \texttt{ring-add-assoc}, \texttt{ring-add-comm},
  \texttt{ring-add-left-id}, \texttt{ring-add-left-inv},
  \texttt{ring-mul-assoc}, \texttt{ring-mul-left-id},
  \texttt{ring-mul-right-id},
  \texttt{ring-left-dist}, \texttt{ring-right-dist} \\[4pt]
\texttt{metric-space} &
  \texttt{pts}\,(1), \texttt{dist}\,(2) &
  \texttt{metric-pos},
  \texttt{metric-self-zero},
  \texttt{metric-zero-eq},
  \texttt{metric-sym},
  \texttt{metric-triangle} \\
\bottomrule
\end{tabular}
\end{center}

\begin{remark}
Each entry above is citable by name with \texttt{(ta 'name)}, which transports
it into the current proof.  They do \emph{not} all have the same standing,
however, and the column heading is a convenience rather than a claim.

A structure declaration generates an \texttt{is-\textit{x}} defining biconditional, and
each operation-property conjunct of it, unfolded, \emph{is} one of the
structure's laws.  Stating such a law separately therefore adds no mathematical
content, and the question is only how the system records that fact.  Three
answers are in use, and they cost a citing \texttt{qed} the same --- nothing ---
by three different routes:

\begin{itemize}
\item \emph{Proved.}  The \textsc{Metric-Space} laws (\texttt{metric-pos},
  \texttt{metric-self-zero}, \texttt{metric-zero-eq}, \texttt{metric-sym},
  \texttt{metric-triangle}) are established by \texttt{prove-metric-law!} in
  \texttt{structure-library/metric-laws.scm}; the \textsc{Group} laws
  (\texttt{group-assoc}, \texttt{group-left-id}, \texttt{group-left-inv},
  \texttt{group-identity-in}) and \texttt{abelian-group-opr-comm} by
  \texttt{stl--project!} in \texttt{structure-library/subtype-laws.scm}.  The
  unfold is \emph{checked}, not asserted to exist.
\item \emph{Installed and stamped \texttt{definitional}.}  The \textsc{Ring}
  laws are installed by \texttt{add-axiom!} in
  \texttt{structure-library/ring.scm} and then swept with
  \texttt{register-provenance!\ \ldots\ 'definitional} in the same file.  A
  \texttt{definitional} fact contributes the empty set to every bill, exactly as
  a \texttt{primitive} one does.
\item \emph{Wrapped at installation.}  \texttt{structure-library/module.scm}
  reaches the same place by the other door, installing its projections inside
  \texttt{(fluid-let ((*current-provenance* 'definitional)) \ldots)}.
\end{itemize}

Two laws in the table are genuinely \emph{derived} rather than projected, and
they do carry a warrant: \texttt{ring-mul-zero-left} and
\texttt{ring-mul-zero-right} ($0\cdot a = 0$, which needs distributivity and
cancellation, not merely the definition of a ring).
\texttt{reference/PROOF-DEBT.md} is the authority on which is which; this table
is not.  For example, to apply
\texttt{group-left-inv} (which asserts
$\forall s,a.\;\texttt{is-group}(s)\Rightarrow a\in\texttt{carr}(s)
 \Rightarrow \texttt{opr}(s)(\texttt{inv}(s)(a),a)=\texttt{iden}(s)$),
supply the group and the element as instantiation witnesses.
\end{remark}

If you prefer to work with raw tuples directly, the n-ary constructors
are still available:
\begin{VNBcode}
; Prove a tuple is in a Cartesian product
(sp (make-wff-from-string
  "G in power(U) and op in fun(cartesian(G,G),G) and e in G and inv in fun(G,G)
   implies [G, op, e, inv]
   in cartesian(power(U), fun(cartesian(G,G),G), G, fun(G,G))"))
(di)
(ai "G in power(U) and op in fun(cartesian(G,G),G) and e in G and inv in fun(G,G)")
(ai "op in fun(cartesian(G,G),G) and e in G and inv in fun(G,G)")
(ai "e in G and inv in fun(G,G)")
(ci)                              ; cartesian-intro: 4 subgoals
(ass) (ass) (ass) (ass)
\end{VNBcode}

\subsection{Applying a macete inside a proof}

The short-form command \texttt{(mac 'name)} applies the named macete to
the \emph{current goal}.  It unifies the macete's source pattern against
the goal formula (outermost-first, depth-first) and replaces the first
matched subterm with the replacement pattern.

\paragraph{Biconditional macete.}
\texttt{pairing-membership} states
$\forall a\,b\,x.\;x \in \mathrm{pair}(a,b) \Leftrightarrow (x=a \vee x=b)$.
Its source pattern is \verb|x in pair(a, b)| and its replacement is
\verb|x = a or x = b|.

\begin{VNBcode}
(sp (make-wff-from-string "2 in pair(2, 3)"))
; proof state:  []  |-  2 in pair(2, 3)
(mac 'pairing-membership)
; source: x in pair(a, b)   replacement: x = a or x = b
; unifies: x=2, a=2, b=3
; goal rewritten:  2 in pair(2, 3)  ->  2 = 2 or 2 = 3
; proof state:  []  |-  2 = 2 or 2 = 3
(oi-l)      ; or-intro-left: new goal  2 = 2
(rfl)       ; reflexivity
; => Proof complete.
\end{VNBcode}

\paragraph{Rewrite fires inside a larger formula.}
The rewrite finds its source pattern at the first matching position,
which may be nested inside a connective.

\begin{VNBcode}
(sp (make-wff-from-string "x in pair(a, b) implies x = a or x = b"))
; proof state:  []  |-  (x in pair(a,b)) implies (x = a or x = b)
(mac 'pairing-membership)
; pattern "x in pair(a, b)" matched at the left of implies
; goal rewritten:  (x = a or x = b) implies (x = a or x = b)
(di)        ; discharge implies: push antecedent to context
; context: ["x = a or x = b"]   goal: "x = a or x = b"
(ass)
; => Proof complete.
\end{VNBcode}

\begin{remark}
Macetes rewrite the goal; they do not rewrite the context.  If the
matching pattern appears only in an assumption, \texttt{(mac 'name)}
will not fire.  To use an assumption as a rewrite source, retrieve it
with \texttt{(ai ...)} so that it appears in the goal.
\end{remark}

\paragraph{Conditional macete.}
For a conditional theorem
$\forall\bar{x}.\;(\mathit{cond} \Rightarrow s = r)$, the macete still
fires unconditionally: it rewrites $s \mapsto r$ without generating
side-goals for $\mathit{cond}$.  The goal $s$ is replaced by~$r$
(or by $\top$ when the conclusion is not an equation or biconditional),
and the resulting $\top$ goal is closed with \texttt{(ass)}.

\begin{VNBcode}
; nn-add-closed:
;   forall([a in NN, b in NN], a + b in NN)
; As a macete: source = "a + b in nn"
;              replacement = truth
(sp (make-wff-from-string "2 + 3 in NN"))
; proof state:  []  |-  2 + 3 in NN
(mac 'nn-add-closed)
; unifies a=2, b=3
; conditions (a in NN, b in NN) not discharged
; goal rewritten:  2 + 3 in NN  ->  truth
(ass)       ; truth is trivially grounded
; => Proof complete.
\end{VNBcode}

\begin{remark}
The conditions of a conditional macete are silently dropped.  This
makes \texttt{(mac 'name)} fast for goals whose conditions are obviously
satisfied, but it is the user's responsibility to verify them; the
proof sketch above is technically unsound without separate proofs that
$2 \in \mathtt{nn}$ and $3 \in \mathtt{nn}$.  For fully verified
condition discharge, use \texttt{(ta 'name)} followed by explicit
instantiation and back-chaining.
\end{remark}

\subsection{Ground arithmetic with \texttt{(arith)}}

The \texttt{(arith)} command (primitive \texttt{cmd-arith}) closes any
goal whose formula reduces to true under exact arithmetic evaluation:

\begin{VNBcode}
(sp (make-wff-from-string "1/3 + 1/6 = 1/2"))
(arith)                              ; closes by exact rational arithmetic

(sp (make-wff-from-string "2 ^ 10 = 1024"))
(arith)                              ; 2^10 evaluates to 1024

(sp (make-wff-from-string "3 + 4 <= 2 * 5"))
(arith)                              ; 7 <= 10
\end{VNBcode}

\subsection{The calculator pattern}

\texttt{(arith)} also recognises a goal of the form
\begin{VNBsnip}
forsome([n in CLASS], EXPR = n)
\end{VNBsnip}
where \texttt{EXPR} is a ground arithmetic expression and \texttt{CLASS}
is one of \texttt{nn}, \texttt{zz}, \texttt{qq}, \texttt{rr},
\texttt{cc}.  It evaluates \texttt{EXPR}, checks that the result belongs
to \texttt{CLASS}, and closes the goal with that value as the witness.
This lets the prover double as a calculator:

\begin{VNBcode}
(sp (make-wff-from-string "forsome([n in NN], 2^10 = n)"))
(arith)   ; closes with witness 1024

(sp (make-wff-from-string "forsome([q in QQ], 1/3 + 1/4 = q)"))
(arith)   ; closes with witness 7/12
\end{VNBcode}

\begin{remark}
The symmetric form \texttt{n = EXPR} is also accepted.
The pattern is intentionally narrow: \texttt{EXPR} must be fully ground
(no free variables), and the existential variable must appear alone on
one side of the equality.
\end{remark}

\subsection{\texorpdfstring{$\beta$}{beta}-reduction with \texttt{(beta)}}

Applied \texttt{lambda} (or \texttt{lambdoid}) terms are reduced
in-place, then the underlying numeric or logical content is closed
by another rule:

\begin{VNBcode}
(sp (make-wff-from-string "lambda([x in nn], 2 * x)(5) = 10"))
(beta)   ; goal becomes (2 * 5 = 10)
(arith)  ; closes
\end{VNBcode}

\subsection{Auto-installing arithmetic theorems}

\texttt{vnb-create-num-theorem} verifies a closed arithmetic wff by
ground evaluation and installs it as a named theorem.  The theorem is
then available via \texttt{(ta \textit{name})}:

\begin{VNBcode}
(vnb-create-num-theorem 'pythag-3-4-5 "3 ^ 2 + 4 ^ 2 = 5 ^ 2")
; => pythag-3-4-5

(sp (make-wff-from-string "truth implies 3 ^ 2 + 4 ^ 2 = 5 ^ 2"))
(di) (ta 'pythag-3-4-5) (ass)
; => Proof complete.
\end{VNBcode}

False or undecidable formulas are rejected:

\begin{VNBcode}
(vnb-create-num-theorem 'bogus "1 + 1 = 3")
; ERROR: formula evaluates to FALSE

(vnb-create-num-theorem 'bogus2 "x = 5")
; ERROR: not a decidable closed arithmetic sentence
\end{VNBcode}

\subsection{Installing a proven theorem}
\label{sec:qed-and-subst}

A completed proof does not by itself add anything to the library.  Use
\texttt{(qed \textit{name})} to install the root assertion as a named
theorem:

\begin{VNBcode}
(sp (make-wff-from-string "n in NN and m in NN implies m in NN and n in NN"))
(di) (ai "n in NN and m in NN") (di) (ass) (ass)
(qed 'and-comm)
; => and-comm
; (lookup-theorem 'and-comm) returns the formula S-expression.
\end{VNBcode}

\paragraph{What \texttt{subst-wff} is and is not.}
\texttt{subst-wff} applies capture-avoiding substitution to a wff's
underlying S-expression and returns a new \texttt{wff} (re-validated
through \texttt{make-wff}).  It is \emph{not} a logical inference;
after calling it you must reprove the result.  It is useful for
generating related goals from a known formula:

\begin{VNBcode}
(define w
  (subst-wff '((n . (in x nn)) (m . (in y rr)))
             (lookup-theorem 'and-comm)))
; w is the wff for (x in nn) and (y in rr) implies (y in rr) and (x in nn)
\end{VNBcode}

\subsection{Saving and replaying proof scripts}
\label{sec:replay}

Every interactive command you issue between \texttt{(sp\ \ldots)} and
\texttt{(qed\ \ldots)} is appended to a global accumulator
\texttt{*proof-script*}.  \texttt{(qed \textit{name})} installs the
proven theorem and saves the script under the same name; you can also
call \texttt{(save-proof \textit{name})} on demand.

\begin{VNBcode}
(sp (make-wff-from-string "n in NN and m in NN implies m in NN and n in NN"))
(di) (ai "n in NN and m in NN") (di) (ass) (ass)
(qed 'and-comm)
; theorem AND script installed under the same name

(lookup-proof 'and-comm)
; => ((di) (ai (and (in n nn) (in m nn))) (di) (ass) (ass))
\end{VNBcode}

\paragraph{Replay verbatim.}
\texttt{(replay-proof \textit{name})} re-runs the saved commands
sequentially against the current proof state.  Useful for redoing a
proof in a fresh session, or applying an identical script to a sequent
of the same shape.

%% =========================================================
\subsection{\texttt{nn} induction for type properties}
\label{sec:ex-ni-type}

The \texttt{(ni)} command reduces a goal
$\forall n \in \mathtt{nn}.\; P(n)$ to two subgoals:
a base case $P(0)$ (focused first) and a step case
$\forall n \in \mathtt{nn}.\; P(n) \Rightarrow P(\mathrm{succ}(n))$.
The following proof shows that $\mathit{prod\text{-}ord}(m, f, n)$
stays in the carrier $\mathtt{carr}(m)$ for every $n \in \mathtt{nn}$,
deriving \texttt{prod-ord-type} from first principles.

\begin{VNBcode}
(sp (make-wff '(IMPLIES (IS-MONOID m)
               (IMPLIES (IN f (FUN NN (CARR m)))
                 (FORALL n (IMPLIES (IN n NN)
                   (IN (PROD-ORD m f n) (CARR m))))))))
(di) (di)   ; IS-MONOID(m) and f in FUN(NN,CARR(m)) enter context
(ni)        ; NN-induction; BASE gets focus first

;; --- BASE: PROD-ORD(m,f,0) in CARR(m) ---
;; Strategy: PROD-ORD(m,f,0) = IDEN(m) (prod-ord-zero),
;;           IDEN(m) in CARR(m)            (monoid-identity-in),
;;           so by eq-subst-membership the goal closes.
(ta 'eq-subst-membership)
(inst ... '(PROD-ORD m f 0)) (inst ... '(IDEN m)) (inst ... '(CARR m))
(bc '(IMPLIES (AND (= (PROD-ORD m f 0) (IDEN m)) (IN (IDEN m) (CARR m)))
              (IN (PROD-ORD m f 0) (CARR m))))
(di)    ; AND-split: focus = equality goal
(ta 'prod-ord-zero) (inst ... 'm) (inst ... 'f) (ass)
        ; focus jumps: manually refocus on membership goal
(ta 'monoid-identity-in) (inst ... 'm)
(bc '(IMPLIES (IS-MONOID m) (IN (IDEN m) (CARR m)))) (ass)
        ; BASE complete; focus moves to STEP

;; --- STEP: forall([n in NN], IH => PROD-ORD(m,f,succ n) in CARR(m)) ---
(di)    ; peel FORALL n -> n_k; (IN n_k NN) enters context
;; n_k is the bound name introduced by (di); fetch it via:
;;   (cadr (wff-formula
;;           (car (sequent-node-assumptions
;;                  (proof-state-focus *ps*)))))
(di)    ; peel IH: PROD-ORD(m,f,n_k) in CARR(m) enters context
;; Strategy:
;;   PROD-ORD(m,f,succ n_k)
;;     = OPR(m)(PROD-ORD(m,f,n_k), f(n_k))   (prod-ord-succ)
;;   then close via monoid-carrier-closed-opr.
(ta 'eq-subst-membership) (inst ...) (inst ...) (inst ...)
(bc ...)
(di)    ; AND-split: focus = equality goal
(ta 'prod-ord-succ)
(inst ...) (inst ...) (inst ...)
(bc ...) (ass)   ; n_k in NN closes it
;; focus moves to membership goal: OPR(m)(...) in CARR(m)
(ta 'monoid-carrier-closed-opr)
(inst ...) (inst ...) (inst ...)
(bc ...)
(di) (ass)   ; IS-MONOID(m)
(di) (ass)   ; PROD-ORD(m,f,n_k) in CARR(m) [IH]
;; f(n_k) in CARR(m) via fun-apply-type:
(ta 'fun-apply-type)
(inst ...) (inst ...) (inst ...) (inst ...)
(bc ...) (di) (ass) (ass)
;; closes f in FUN(NN,CARR(m)) and n_k in NN
\end{VNBcode}

\begin{remark}
The \texttt{inst} chains are elided above; each peels one leading
\texttt{forall} from a theorem schema.  Two subtleties arise in practice.
First, \texttt{(di)} on an AND-goal always focuses on the left (equality)
conjunct first.  After \texttt{(ass)} closes that conjunct, automatic
focus advancement jumps to the earliest ungrounded goal in the deduction
graph --- which may be the STEP node, not the right (membership) conjunct.
The fix is to capture the membership node
immediately after the AND-split:
\begin{VNBcode}
(di)   ; AND-split
(let* ((dg       (proof-state-dg *ps*))
       (mem-node (car (reverse (dg-sequent-nodes dg)))))
  ... prove equality ... (ass)   ; focus jumps
  (set-proof-state-focus! *ps* mem-node)
  ... prove membership ...)
\end{VNBcode}

Second, \texttt{alpha-equiv?} --- used by \texttt{(inst)} and
\texttt{(bc)} when searching assumptions --- relies on structural
equality of formula heads.  Compound heads such as
\texttt{((opr m) a b)} (a curried two-argument application; see
Section~\ref{sec:apply-tupling}) are compared correctly only when
the head pair is the same \emph{object}.  If \texttt{alpha-equiv?}
fails to locate a formula you know is in context, the likely cause
is a compound-head mismatch; the fix is in
\texttt{expressions.scm}'s \texttt{alpha-equiv-under?}.
\end{remark}

\begin{remark}[The curried/tupled apply convention]
\label{sec:apply-tupling}
VNB's S-expression form \texttt{(f a$_1$ \ldots a$_n$)} is what the
string-syntax parser emits from \texttt{f(a$_1$,\,\ldots,\,a$_n$)} ---
verbatim, no implicit tupling.  An operation typed
\texttt{f $\in$ fun(cartesian(A$_1$,\,\ldots,\,A$_n$),~B)}, however,
takes one argument from the Cartesian product.  The two views are
reconciled by the library definitions
\begin{equation*}
\texttt{apply-tupling-$n$:} \qquad
(f\ a_1\ \ldots\ a_n)\ \mathrel{==}\ (f\ (\texttt{list}\ a_1\ \ldots\ a_n))
\quad\text{for each arity }n
\end{equation*}
(definitions for $n=2,3$ in \texttt{theorem-library/axioms.scm};
add more arities by declaration as needed).  There is none for $n=1$:
$(f\ a) \mathrel{==} (f\ (\texttt{list}\ a))$ would give $a=[a]$ for every set $a$
(Section~\ref{sec:curried-tupled-apply}), so $[\,]$ would be a list of length one.  Quasi-equality
$\mathrel{==}$, not partial $=$: the convention holds vacuously when
either side is undefined.  Combined with \texttt{fun-apply-type} and
\texttt{quasi-eq-def}, these give the $n$-ary apply closure
\((f\ a_1\ \ldots\ a_n)\in B\) from the operation's signature and
the membership of each $a_i$ in its declared component.
\end{remark}

\subsection{Infinite descent: \texorpdfstring{$\sqrt3$}{sqrt3} is irrational}
\label{sec:ex-sqrt3}

Stated without reals or division --- the form the checker proves, entirely
inside $\mathtt{nn}$ --- $\sqrt3$ irrational is
\[
  \forall p,q \in \mathtt{nn}.\quad
  q \neq 0 \;\Longrightarrow\; p\cdot p \neq 3\cdot(q\cdot q) ,
\]
that is, the goal
\begin{VNBsnip}
(FORALL p (IMPLIES (IN p NN) (FORALL q (IMPLIES (IN q NN)
  (IMPLIES (NOT (= q 0)) (NOT (= (* p p) (* 3 (* q q)))))))))
\end{VNBsnip}

\paragraph{The argument.}
Infinite descent.  Suppose the equation has a solution with $q \neq 0$, and
among all such solutions choose one whose numerator $w$ is \emph{least}.  Since
$3 \mid w\cdot w$ we have $3 \mid w$, say $w = 3k$; substituting and cancelling
the factor $3$ gives $q_0\cdot q_0 = 3\cdot(k\cdot k)$, whence $3 \mid q_0$, say
$q_0 = 3m_0$; cancelling once more gives $k\cdot k = 3\cdot(m_0\cdot m_0)$.  Thus
$(k,m_0)$ is again a solution, with $k \neq 0$ and $k < 3k = w$ --- a numerator
strictly below the least one.  Contradiction.

The entire mathematical content is the single implication
\begin{equation}\label{eq:3divsq}
  3 \mid p\cdot p \;\Longrightarrow\; 3 \mid p ;
\end{equation}
everything else --- clearing denominators, the two cancellations, the order
step $k < 3k$ --- is bookkeeping.

\paragraph{Why parity is not special.}
The $\sqrt2$ proof is the same skeleton with \emph{even} for \emph{divisible by
$3$}: there \eqref{eq:3divsq} reads ``$p\cdot p$ even $\Rightarrow p$ even'' and
follows from parity, every natural being $2k$ or $2k+1$.  For $\sqrt3$ it follows
from the trichotomy --- every natural is $3k$, $3k+1$, or $3k+2$ --- together
with the fact that the three residues square to $0,1,1$ modulo $3$, so a square
is divisible by $3$ only when its root is.  Parity is the $r=2$ instance of a
uniform statement about residues modulo $r$; the $\sqrt3$ development is the
$\sqrt2$ development with the residue system widened by one, and nothing changes
in kind.

\paragraph{The lemma ladder.}
The mod-$3$ arithmetic (\texttt{theorem-library/nn-mod3-proof.scm}) is built by
induction from the two Peano recursion equations for $+$ and $\cdot$, in exact
parallel with the parity ladder (\texttt{nn-parity-proof.scm}).  Each lemma is
\emph{proven}; none is asserted.

\begin{center}
\begin{tabular}{lll}
\toprule
parity ($r=2$) & this proof ($r=3$) & statement \\
\midrule
\texttt{nn-two-mul-succ}     & \texttt{nn-three-mul-succ} & $r\cdot\mathrm{succ}(k) = \mathrm{succ}^{\,r}(r\cdot k)$ \\
\texttt{nn-parity}           & \texttt{nn-trichotomy-3}   & every $n$ is $r\cdot k + i$, residue $i<r$ \\
\texttt{nn-parity-exclusive} & \texttt{nn-3-not-succ}     & $r\cdot x$ is never $r\cdot y + 1$ \\
\texttt{nn-even-square}      & \texttt{nn-3-div-square}   & $r \mid p\cdot p \Rightarrow r \mid p$ \\
\bottomrule
\end{tabular}
\end{center}

\noindent For $r=3$ a single residue-inequality suffices: both nonzero residues
square to $r\cdot j + 1$, so only ``not $\equiv 1$'' is ever needed.

\paragraph{The descent, mechanized.}
The driver (\texttt{theorem-library/sqrt3-proof.scm}) states the goal and hands
the descent to three tools; the intermediate steps are elided.

\begin{VNBcode}
(sp (make-wff '(FORALL p (IMPLIES (IN p NN) (FORALL q (IMPLIES (IN q NN)
     (IMPLIES (NOT (= q 0)) (NOT (= (* p p) (* 3 (* q q)))))))))))
(di)(di)(di)(di)(di)(di)                 ; intro p,q, hyps; assume p*p = 3*(q*q)

(define R (minimize! '(pv) (guard 'pv) 'pv))    ; choose the LEAST numerator w
;; ... discharge minimize!'s two obligations: measure in NN; guard satisfiable

(define k (obtain (lambda () (fact 'nn-3-div-square w))))   ; 3|w  =>  w = 3k
;; ... q0*q0 = 3*(k*k) by have! ; 3|q0 => q0 = 3*m0 by obtain ; k*k = 3*(m0*m0)

(inst+ minf k)              ; minimality gives  w <= k
(fact 'nn-lt-triple k)      ; but k < 3k = w -- contradiction
\end{VNBcode}

\begin{itemize}
\item \texttt{minimize!} supplies the least-numerator witness together with its
  two honest obligations --- the measure lands in $\mathtt{nn}$, and the guard
  is satisfiable --- and appeals to no choice principle.
\item \texttt{obtain} names each descended witness --- the $k$ with $w = 3k$,
  the $m_0$ with $q_0 = 3m_0$ --- by \emph{reading it off the proof state} (the
  variable left free once the existential is eliminated), never by guessing a
  machine-generated name.  This is what makes the descent survive a re-run.
\item \texttt{have!} discharges each intermediate equality and typing from
  context, so the script reads as the argument does: ``we have
  $q_0\cdot q_0 = 3\cdot(k\cdot k)$; we have $k \neq 0$; \ldots''.
\end{itemize}

\begin{remark}[The \texttt{minimize!} step]
\label{remark:minimize}
The enigmatic-looking \texttt{(minimize! '(pv) (guard 'pv) 'pv)} decodes as:
choose the single number \texttt{pv} satisfying the descent guard, with
\texttt{pv} itself as the quantity to minimise --- that is, take the
\emph{least} numerator meeting the guard.

The second argument to \texttt{minimize!} is the \textbf{guard formula} itself:
an ordinary VNB formula, free in the variables being chosen (here \texttt{pv}),
stating the property those values must satisfy.  \texttt{guard} is \emph{not} a
built-in and is nothing special to \texttt{minimize!} --- it is a helper the
proof author writes (named \texttt{s3-guard} in \texttt{sqrt3-proof.scm}), a
one-argument function that takes a term and returns the guard formula about it.
So \texttt{(guard 'pv)} yields the formula with \texttt{pv} free --- exactly what
\texttt{minimize!} wants --- and the \emph{same} helper, called later at the
chosen witness $w$ as \texttt{(guard $w$)}, restates that property of $w$.  Here
the property is ``\texttt{pv} is a natural number that is the numerator of a
fraction whose square is~$3$''; written out, the helper is just a formula
constructor:

\begin{VNBcode}
(define (s3-guard pv)
  (list 'AND (list 'IN pv 'NN)
        (list 'FORSOME 'q_
          (list 'AND (list 'IN 'q_ 'NN)
                (list 'AND (list 'NOT (list '= 'q_ 0))
                      (list '= (list '* pv pv) (list '* 3 (list '* 'q_ 'q_))))))))
\end{VNBcode}

Passing the formula literally would work just as well; the helper only spares the
author from writing it twice.  In general,
\[
  \texttt{(minimize! '($v_1$ \ldots\ $v_k$) GUARD MEASURE)}
\]
means ``choose $v_1,\ldots,v_k$ satisfying the formula \texttt{GUARD} with the
$\mathtt{nn}$-valued term \texttt{MEASURE} as small as possible.''  On the main
branch it lands two hypotheses, for fresh eigenconstants $w_1,\ldots,w_k$: that
\texttt{GUARD} holds of them, and that nothing satisfying \texttt{GUARD} has a
strictly smaller \texttt{MEASURE}.  It leaves exactly the two obligations any
minimisation owes --- that \texttt{MEASURE} lands in $\mathtt{nn}$ wherever
\texttt{GUARD} holds, and that \texttt{GUARD} holds somewhere.  Everything
between --- forming the value set, well-ordering it, unpacking the least
witness, restating minimality over the $v_i$ rather than the set --- it does for
you.  Every $v_i$ must occur in \texttt{MEASURE} (a variable the measure ignores
is not one you are minimising over).

It adds no kernel rule --- it chains \texttt{cut}, \texttt{di}, \texttt{ai},
\texttt{ew}, \ldots --- and uses \emph{no choice principle}: its one appeal is
\texttt{nn-least-element}, the well-ordering of $\mathtt{nn}$, and well-ordering
returns a genuine \emph{member} of the value set, from which the witness is read
back.  Stated instead as a first-order lemma, the same fact would force every
call site to tuple its variables into a Cartesian product, build a lambda, and
prove that lambda inhabits $\mathtt{fun}(U,\mathtt{nn})$ --- a typing obligation
larger than the minimisation it buys.  As a tactic it quantifies over the $v_i$
at the meta level and never forms $U$ at all.
\end{remark}

\begin{remark}[\texttt{obtain} versus \texttt{vlet}]
\texttt{obtain} and the \texttt{choice} former of \texttt{vlet} are the same
mechanism --- eliminate an existential and read its witness off by
free-variable difference --- but they differ in what they take.  \texttt{obtain}
takes a \emph{lane}: it runs a forward step and names the witness that step
\emph{just produced}.  That is the right fit inside a descent, where the relevant
existential is freshly landed and typically sits beside others: at the
``$3 \mid w$'' step the conclusion $\exists k.\,w = 3k$ coexists in context with
the evidence $\exists m.\,w\cdot w = 3m$ cut in to derive it, and \texttt{obtain}
selects the new one automatically.  Concretely, that one step reads
\begin{VNBsnip}
;; as written in the proof:
(define k (obtain (lambda () (fact 'nn-3-div-square w))))
;; the vlet alternative -- land the existential, then name it by its body:
(fact 'nn-3-div-square w)
(define k (vlet--choice! 'k '(AND (IN k NN) (= w (* 3 k)))))
\end{VNBsnip}
The \texttt{vlet} body names the eigenconstant $w$, a value produced at run
time, so this is the procedure \texttt{vlet--choice!}, not the
\texttt{(vlet \ldots)} macro (whose body is a fixed literal pattern).
\texttt{vlet} instead destructures a
\emph{named, stable} formula: its \texttt{match} former binds the holes of a
pattern (``the $a$ and $b$ in $x\cdot b = a$''), and its \texttt{choice} former
eliminates a \emph{designated} existential.  The rule of thumb: \texttt{obtain}
to name whatever a forward step produced, \texttt{vlet} to take apart a formula
you can already name.
\end{remark}

\begin{remark}
The proof adds no asserted step of its own: the mod-$3$ ladder is proven and the
descent cites it, exactly as \texttt{sqrt2-proof.scm} cites the parity ladder.
A new irrationality --- $\sqrt r$ for any non-square $r$ --- is then a matter of
widening the residue ladder, not of inventing an argument.  That division of
labour is the lesson: the author states the intermediate claims, and the machine
supplies the least witness (\texttt{minimize!}), the fresh names
(\texttt{obtain}), and the bookkeeping (\texttt{have!}, \texttt{from-context!}).
\end{remark}

\paragraph{The same proof at the interactive surface.}
By the transcript principle (Section~\ref{sec:two-surfaces}), most of this
script is \emph{already} a sequence of interactive short forms: \texttt{di},
\texttt{ai}, \texttt{ew}, \texttt{fact}, \texttt{cut}, \texttt{subst},
\texttt{crs}, \texttt{ass}, and \texttt{minimize!} are exactly what a user
types.  What the driver adds is the orchestration a human otherwise does by eye,
and it is of two kinds.  \emph{Focus}: after a branching command one selects the
next open leaf from the frontier, where the script calls \texttt{dk-focus!}.
\emph{Names}: each \texttt{(ai\ (forsome \ldots))} skolemizes to a machine-chosen
eigenconstant, printed in the resulting sequent, which the user reads and types
into the following commands.  The $3\mid w$ step, done by hand, is
\begin{VNBsnip}
(fact 'nn-3-div-square w)                          ; lands  exists k. w = 3k
(ai '(FORSOME k (AND (IN k NN) (= w (* 3 k)))))    ; skolemize; sequent now shows k_899
;; ... read k_899 off the sequent; type it into the subsequent (subst ...) / (fact ...)
\end{VNBsnip}
\texttt{obtain} is exactly this pair automated --- run the forward step, read the
fresh name back by free-variable difference --- so that a \emph{script}, unlike a
transcript with \texttt{k\_899} hard-coded, replays unchanged when the
eigenvariable counter lands on a different number.

%% =========================================================


%% =========================================================
\chapter{Working with formulas}
\label{ch:formulas}

Chapter~\ref{sec:language} describes the expression language.  This
chapter gives the Scheme procedures that work with its formulas:
\begin{itemize}
\item the parser and the printer, which turn a string into a formula and
  back (Section~\ref{sec:parser-printer});
\item the accessors that name the parts of a formula
  (Section~\ref{sec:parts});
\item the \texttt{wff} record that holds a checked formula
  (Section~\ref{chap:syntax-contexts});
\item the two ways of changing how a formula is printed:
  \texttt{notation!} (Section~\ref{sec:notation}) and the optional infix
  layer (Section~\ref{sec:infix-layer}).
\end{itemize}

%% =========================================================

\section{Parser and printer}
\label{sec:parser-printer}

Four entry points provide string-level access to VNB formulas:

\begin{VNBsnip}
(make-wff-from-string str)  ; string -> <wff>  (validates via make-wff)
(parse-string str)          ; string -> raw S-expression (no validation)
(wff->string w)             ; <wff>  -> string
(pp w)                      ; display wff->string then newline
\end{VNBsnip}

\noindent
\texttt{parse-string} converts a string to a raw S-expression without
validation; it is used internally by the interactive commands and is
available for low-level \texttt{cmd-*} calls.  For normal use, pass
formula strings directly to \texttt{ai}, \texttt{bc}, \texttt{inst},
\texttt{cut}, \texttt{wk}, and \texttt{vnb-create-num-theorem}.
\texttt{make-wff-from-string} is the right choice when constructing a
formula to prove (for \texttt{sp}).

\paragraph{Operator precedence (high to low).}
\begin{center}
\small
\begin{tabular}{lll}
\toprule
Level & Operators & Associativity \\
\midrule
primary & atoms, \texttt{f(a,b)}, \texttt{[a,b]}, \texttt{(expr)} & --- \\
power & \texttt{\textasciicircum} & right \\
unary & \texttt{-x} & right \\
mul & \texttt{*}, \texttt{/} & left (n-ary \texttt{*}) \\
add & \texttt{+}, \texttt{-} & left (n-ary \texttt{+}) \\
cmp & \texttt{in}, \texttt{=}, \texttt{==}, \texttt{<=} & binary \\
and & \texttt{and} & left (n-ary) \\
or & \texttt{or} & left (n-ary) \\
iff & \texttt{iff} & left (n-ary) \\
implies & \texttt{implies} & right \\
\bottomrule
\end{tabular}
\end{center}

\paragraph{Division and negation sugar.}
\begin{VNBsnip}
x / y    =>  (* x (recip y))   ; division is syntactic sugar
-x       =>  (- x)             ; unary minus; lower prec than ^
-x ^ 2   =>  (- (power x 2))  ; i.e. -(x^2), not (-x)^2
\end{VNBsnip}

\paragraph{Quantifier syntax.}
Quantifiers use the binding-list form:
\begin{VNBsnip}
forall([x], body)              ; bare variable (all classes)
forall([x in A], body)         ; guarded variable
forall([[a, b, c] in A], body) ; tuple destructuring
forall([x in A, y, z in B], body)  ; mixed bindings
forsome([x in NN], body)       ; existential; same forms
\end{VNBsnip}

\noindent
The bracket syntax after the function name distinguishes quantifiers
from ordinary function calls.

\paragraph{Display conventions (\texttt{wff->string} / \texttt{pp}).}
\begin{itemize}
\item \begin{sloppypar}\textbf{Quantifiers} print in binding-list form:
  the primitive expanded form is reverse-engineered.
  \texttt{forall([x in nn, y], p(x, y))}\end{sloppypar}
\item \textbf{Logical connectives} \texttt{and}, \texttt{or}, \texttt{iff}
  print n-ary infix with surrounding parentheses:
  \texttt{((p and q) and r)}.
  \texttt{implies} prints binary infix: \texttt{(p implies q)}.
\item \textbf{Comparisons} \texttt{in}, \texttt{=}, \texttt{==},
  \texttt{<=} print binary infix: \texttt{(x in nn)}, \texttt{(x = y)}.
\item \textbf{Arithmetic} \texttt{+}, \texttt{*} print n-ary infix:
  \texttt{(x + y + z)}.  Binary \texttt{-} prints as \texttt{(x - y)};
  unary as \texttt{-x}.  Exponentiation as \texttt{(x \textasciicircum{} y)}.
\item \textbf{Lists} \texttt{[a, b, c]} in both input and output: the
  list syntax is the same in string form (internally stored as \texttt{list}).
\item \textbf{All other compound terms} use functional notation
  \texttt{f(x, y, z)}: \texttt{nth}, \texttt{cartesian}, \texttt{fun},
  \texttt{power} (unary set power), \texttt{not}, and user-defined heads.
  Nullary symbols print without parentheses.
\end{itemize}

\paragraph{Example round-trips.}
\begin{VNBsnip}
(pp (make-wff-from-string "forall([x in NN, y in NN], x + y = y + x)"))
; => forall([x in nn, y in nn], ((x + y) = (y + x)))

(pp (make-wff-from-string "not(P) implies falsity"))
; => (not(p) implies falsity)

(pp (make-wff-from-string "2 ^ 10 = 1024"))
; => ((2 ^ 10) = 1024)

(expression->string '(forall x (implies (in x nn)
                       (forall y (implies (in y nn) (= (+ x y) (+ y x)))))))
; => "forall([x in nn, y in nn], ((x + y) = (y + x)))"
\end{VNBsnip}

\noindent
All identifiers are case-folded to lowercase by MIT~Scheme, so
\texttt{NN} and \texttt{nn} are the same symbol; \texttt{FORALL} and
\texttt{forall} are also identical.  The printer always emits lowercase.

%% =========================================================

\section{Naming the parts of a formula}
\label{sec:parts}

The parser and printer of the previous section are the surface-conforming
way to \emph{build} and \emph{display} a formula.  This section gives the
surface-conforming way to \emph{take one apart}.

A script or a piece of scratch-workspace code must never reach into the
internal S-expression by position --- \texttt{(caddr (cadr f))} and the
like.  Such code is welded to a representation accident and breaks the
moment the internal form shifts (the \texttt{and}-nesting, the
\texttt{+}/\texttt{binplus} split, a new sugar): it is the
representation-independence sin one level down.  The accessors in
\texttt{parts.scm} name the parts the way the \emph{external} syntax names
them, so the code that uses them survives changes to the internal form.

There are three entry points: \texttt{formula-kind} classifies the top
construct, \texttt{part} navigates to a named sub-part, and \texttt{match}
fills in the blanks of a surface pattern.  All three accept a wff, a
surface string, or a raw S-expression interchangeably.

\subsection{\texttt{formula-kind} --- what is this, at the top?}

\texttt{(formula-kind \textit{x})} returns a single symbol naming the
outermost surface construct, without committing to any internal detail
below it:

\begin{VNBsnip}
(formula-kind "p(x) implies q(x)")       ; => implication
(formula-kind "x = y")                   ; => equation
(formula-kind "p(x) and q(x) and r(x)")  ; => conjunction
(formula-kind "forall([x in nn], p(x))") ; => universal
(formula-kind "x in nn")                 ; => membership
(formula-kind "f(a, b)")                 ; => application
(formula-kind '((MUL R) a b))            ; => op-application
\end{VNBsnip}

\noindent
The complete set of tags follows.

\begin{itemize}
\item \texttt{truth}, \texttt{falsity} --- the logical constants.
\item \texttt{atom} --- a bare symbol; \texttt{number} --- a numeric literal.
\item \texttt{implication} (\texttt{implies}), \texttt{biconditional}
  (\texttt{iff}), \texttt{conjunction} (\texttt{and}), \texttt{disjunction}
  (\texttt{or}), \texttt{negation} (\texttt{not}) --- the connectives.
\item \texttt{equation} (\texttt{=}), \texttt{equivalence} (\texttt{==}),
  \texttt{membership} (\texttt{in}), \texttt{inequality} (\texttt{<=}),
  \texttt{subset} (\texttt{subset}) --- the binary relations.
\item \texttt{universal} (\texttt{forall}), \texttt{existential}
  (\texttt{forsome}), \texttt{description} (\texttt{iota}) --- the binders.
\item \texttt{arithmetic} (\texttt{+}, \texttt{*}).
\item \texttt{set-construction} (\texttt{union}, \texttt{intersection},
  \texttt{cartesian}).
\item \texttt{op-application} --- a structure operator such as
  \texttt{(mul R)} applied to arguments.
\item \texttt{application} --- a named predicate or function applied.
\item \texttt{unknown} --- anything that matches none of the above.
\end{itemize}

\begin{remark}
This vocabulary is meant to be \emph{closed}.  Each tag names a construct
the surface grammar actually has, so the list grows only if the grammar
itself does --- a rare event.  The pull to add a tag for some derived or
compound form should be resisted: such forms already classify as
\texttt{application}, \texttt{op-application}, or \texttt{unknown}, and a
proliferation of fine-grained tags would couple the classifier to passing
notational fashion rather than to the language.
\end{remark}

\subsection{\texttt{part} --- navigate to a sub-part}

\texttt{(part \textit{x} \textit{step} \ldots)} walks a path of selector
steps and returns the part reached.  A single result is itself a wff (it
prints through the normal verbalizer and round-trips); a \emph{plural}
selector returns a \emph{list} of wffs and must be the last step.

\begin{center}
\footnotesize
\begin{tabular}{@{}llp{6.0cm}@{}}
\toprule
Selector & Applies to & Yields \\
\midrule
\texttt{antecedent}, \texttt{consequent} & implication & the two sides \\[2pt]
\texttt{lhs}/\texttt{left}, \texttt{rhs}/\texttt{right} & any binary relation (\texttt{= == iff in <= subset}) & the two sides \\[2pt]
\texttt{negand} & negation & the negated formula \\[2pt]
\texttt{conjuncts}, \texttt{disjuncts} & conjunction / disjunction & \emph{(plural)} the flattened list of conjuncts / disjuncts \\[2pt]
\texttt{(conjunct \textit{N})}, \texttt{(disjunct \textit{N})} & conjunction / disjunction & the \textit{N}-th, 1-indexed \\[2pt]
\texttt{binder} & quantifier (\texttt{forall}/\texttt{forsome}/\texttt{iota}) & the bound variable \\[2pt]
\texttt{matrix}/\texttt{scope}/\texttt{body} & quantifier & the scope \\[2pt]
\texttt{operator}/\texttt{op}/\texttt{head} & application & the operator (e.g.\ \texttt{(mul r)}) \\[2pt]
\texttt{operands}/\texttt{args} & application & \emph{(plural)} the argument list \\[2pt]
\texttt{(operand \textit{N})} & application & the \textit{N}-th argument, 1-indexed \\
\bottomrule
\end{tabular}
\end{center}

\begin{VNBsnip}
(part "p(x) implies q(x)" 'antecedent)         ; => p(x)
(part "x = y" 'lhs)                            ; => x
(part "p(x) and q(x) and r(x)" '(conjunct 2))  ; => q(x)
(part "p(x) and q(x) and r(x)" 'conjuncts)     ; => (p(x) q(x) r(x))
(part "f(a, b, c)" '(operand 2))               ; => b
(part "f(a, b, c)" 'operands)                  ; => (a b c)
(part "forall([x in nn], p(x))" 'binder)       ; => x
(part "forall([x in nn], p(x))" 'matrix)       ; => (x in nn implies p(x))

;; steps chain left-to-right: descend, then descend again
(part "p(x) implies y = z" 'consequent 'lhs)   ; => y
\end{VNBsnip}

\noindent
The selectors are \emph{total over their declared shape and loud
otherwise}: \texttt{(part "x = y" 'antecedent)} is an \emph{error}, not a
silent \texttt{\#f} --- the same partial-\texttt{=} definedness discipline
that runs through the rest of the system (Section~\ref{sec:funapp}).

\begin{remark}
The accessors honour the \emph{surface structure} of a formula --- the
string-syntax view laid out in Section~\ref{sec:parser-printer} --- not the
internal S-expression that stores it.  \texttt{conjuncts} flattens the
associative spine, so \texttt{p and q and r} yields three conjuncts however
the binary \texttt{and} happens to be nested.  \texttt{operands} de-tuples the apply-tupling convention
(Section~\ref{sec:curried-tupled-apply}), so the single internal
\texttt{(list a b)} argument comes back as the operand list
\texttt{(a~b)}.  And the binary structure-op application \texttt{(mul
R)(a, b)} is first-class: \texttt{operator} gives \texttt{(mul r)},
\texttt{operands} gives \texttt{(a b)}.
\end{remark}

\subsection{\texttt{match} --- fill the blanks of a pattern}

\texttt{(match \textit{pattern-string} \textit{x})} matches a surface
pattern against \textit{x}.  The pattern is an ordinary formula in which
some sub-parts are left \emph{blank}, each blank written as a name with a
leading question mark --- \texttt{?P}, \texttt{?Q}, \texttt{?a} --- standing
for an arbitrary sub-formula or term.  The call returns an association list
\texttt{((?v\ .\ \textit{wff})\ \ldots)} pairing each blank with what it
matched, or \texttt{\#f} if the pattern does not fit.

\begin{VNBsnip}
(match "?P implies ?Q" "x in nn implies x in zz")
; => ((?q . x in zz) (?p . x in nn))

(match "?P and ?Q" "x in nn implies x in zz")
; => #f
\end{VNBsnip}

\noindent
Matching reuses the macete one-way matcher
(Section~\ref{sec:macetes}) rather than introducing a second pattern
engine, and binder desugaring is applied to both sides so a quantified
pattern matches a quantified target.  Where \texttt{part} is the right tool
when the shape is known in advance, \texttt{match} is the right tool when
the shape is the thing being tested.

\begin{remark}
Input is coerced uniformly: \texttt{part} and friends accept a wff, a
surface string, or a raw S-expression.  Given a wff or string they hand
back wffs (re-wrapped in the same library); given a bare S-expression there
is no wff parent to re-wrap against, so the parts come back as raw
S-expressions.  One asymmetry follows from \texttt{make-wff} validation: a
standalone \emph{term} string such as \texttt{"x + y"} is not a wff, so
\texttt{(formula-kind "x + y")} signals an error --- classify a term by
passing the S-expression, \texttt{(formula-kind '(+ x y))}, or reach it as
a \texttt{part} of the formula it sits in.
\end{remark}

%% =========================================================


%% =========================================================
\section{Formulas as typed objects: \texttt{wff}}
\label{chap:syntax-contexts}

A \emph{well-formed formula} in this system is not a bare S-expression
but a typed object: a formula that has passed the well-formedness check.  In the REPL it prints as:
\begin{VNBsnip}
#[wff N hD "formula string"]
\end{VNBsnip}
where \texttt{N} is MIT Scheme's object hash, \texttt{hD} is the
formula's digest in hexadecimal, and the quoted string is the
pretty-printed formula.  Every formula in a sequent is a \texttt{wff}.

\texttt{(make-wff expr)} constructs such an object, from a string in
surface syntax or from an S-expression.  \texttt{(free-vars (wff->sexp p))}
returns the variables free in the formula: \texttt{free-vars} works on
the expression, not on the \texttt{wff} record.

Two procedures expose the internals of a \texttt{wff}:
\begin{VNBsnip}
(wff->sexp p)        ; => raw S-expression (no output)
(display-contents p) ; prints kind, library name, formula
\end{VNBsnip}
\texttt{wff->sexp} is defined in \texttt{wff.scm};
\texttt{display-contents} is defined in \texttt{contexts.scm}.
Both are available in interactive sessions.

Example (the base library is \texttt{VNB-LIBRARY}):
\begin{VNBsnip}
(define p (make-wff "forall([x in A], 0 <= funny-name)"))
; => #[wff 12 "forall([x in a], 0 <= funny-name)"]

(free-vars (wff->sexp p))
; => (a funny-name)   ; x is bound by the forall
\end{VNBsnip}


%% =========================================================
\section{Notation: \texttt{notation!}}
\label{sec:notation}

Every head the system knows --- a predicate, a functoid, a structure predicate,
even a primitive like \texttt{=} or \texttt{in} --- has an entry in a single
operator table.  \texttt{def-predicate}, \texttt{def-functoid} and
\texttt{declare-structure} fill in the \emph{structural} part of the entry (its
kind, its arity) at the definition point.  \texttt{notation!} fills in the
\emph{reading} --- how the head renders in \LaTeX, in English, and to the
parser.  The two are orthogonal: re-declaring a head updates the structural slots
and keeps whatever notation was declared, so a \texttt{notation!} may sit beside
the definition or well after it.

It is a keyword form; a call names only the slots it has to set:
\begin{VNBcode}
(notation! 'IS-RING 'noun "ring" 'article "a")
(notation! 'IS-EUCLIDEAN-RING 'kind 'predicate 'arity 1
           'noun "Euclidean ring" 'article "a")
(notation! 'IS-COUNTABLE-PSEUDOMETRIC-FAMILY 'kind 'predicate 'arity 2
           'english "$1 is a countable family of pseudometrics on $2")
\end{VNBcode}

The recognised keys are
\texttt{tex}, \texttt{english}, \texttt{noun}, \texttt{article},
\texttt{prec}, \texttt{infix?}, \texttt{arity}, and \texttt{kind}.

\begin{itemize}
\item \texttt{tex} and \texttt{english} are \emph{templates} --- strings with
  placeholders \texttt{\$1}\ldots\texttt{\$9} for the rendered arguments.  A
  binary operator combining both with \texttt{prec} (a precedence for
  parenthesisation) and \texttt{infix?} reads, schematically,
  \texttt{(notation! '+ 'tex "\$1 + \$2" 'english "\$1 plus \$2" 'prec 60 'infix?
  \#t)}.  When a template cannot express the rendering, \texttt{english} may
  instead be a \emph{procedure} from the list of argument strings to a string.
\item \texttt{noun} and \texttt{article} are the special case of a unary
  structure predicate: given them, \texttt{is-ring($s$)} reads ``$s$ is a
  ring'' throughout the prose renderers, and the article joins correctly
  (``a ring'', ``an $R$-module'').
\item \texttt{kind} and \texttt{arity} let a \texttt{notation!} register a head
  the def-forms never introduced --- a primitive such as \texttt{=} or
  \texttt{in} --- so it, too, gets a reading.
\end{itemize}

One table, one declaration form, read by every renderer: \texttt{expr->tex} and
\texttt{wff-english} for formulas, \texttt{describe-structure} for the structure
cards, and the generated \texttt{OPERATORS.md}.  A head with no
\texttt{notation!} still works --- it falls back to prefix application,
\texttt{is-ring($s$)} --- so notation is an addition, never a prerequisite.

%%% -----------------------------------------------------------------------
\section{The optional infix layer}
\label{sec:infix-layer}

The TeX printer writes an applied structure operation in functional
form: \texttt{((add K) a b)} prints as $\mathtt{add}(K)(a, b)$.  The
\emph{infix layer} (\texttt{tex-infix.scm}) writes it instead as
\[
  a +_K b .
\]
It is off by default, and while it is off the printer writes the
functional form.  It changes only TeX output: the panels, the
proof reader and \texttt{display-proof}.  The parser, \texttt{pp} and
the kernel never see it.

It differs from \texttt{notation!} (Section~\ref{sec:notation}) in
two ways.  A \texttt{notation!} declares the reading of one head and is
always in force.  The layer reads an \emph{accessor applied to a
structure}, so one row serves \texttt{add} in every ring, module and
field, and it is switched on and off as a whole.

\subsection{Turning it on}

In the workspace, the menu \textsf{VNB > Printer Macros} has the
toggle \textsf{Printer Macros On}; the panels repaint at once.  At the
REPL:
\begin{VNBcode}
(tex-infix! #t)                          ; on for every printout
(tex-infix! #f)                          ; off
(display-proof 'NAME "file.pdf" 'infix)  ; on for this printout only
(infix-tex 'NAME)                        ; the TeX of a theorem, layer on
(with-tex-infix CONV thunk)              ; one printout, conventions CONV
\end{VNBcode}

\subsection{What it prints}

An accessor applied to a structure term $S$, then to its arguments,
is written with the subscript $S$:
\[
\begin{array}{ll@{\qquad}ll}
\texttt{((add S) a b)}  & a +_S b        & \texttt{((mul S) a b)}  & a \times_S b \\
\texttt{((opr S) a b)}  & a \cdot_S b    & \texttt{((act S) c v)}  & c \cdot_S v \\
\texttt{((neg S) a)}    & -_S\, a        & \texttt{((inv S) a)}    & a^{-1}_S \\
\texttt{(zero S)}       & 0_S            & \texttt{(iden S)}       & e_S \\
\texttt{((fnrm S) a)}   & \lvert a \rvert_S & \texttt{((vnrm S) v)} & \lVert v \rVert_S \\
\texttt{((ip S) u v)}   & \langle u, v \rangle_S & \texttt{((vstar S) x)} & x^{*}_S
\end{array}
\]
The full list is the base table (below).  $S$ is read through views:
a \texttt{def-functor} view such as
\texttt{(normed-field-as-commutative-ring K)} is written $K$, since the
view's slots are $K$'s slots.

A defined head can also have a template.  The shipped templates print
the sums and products over an index:
\[
  \texttt{series-partial-sum(c, n)} \quad\text{as}\quad \sum_{i < n} c(i).
\]
When the summand is a \texttt{vnb-lambda} over $\mathtt{nn}$, the
lambda's own variable is used and its body printed, as in
$\sum_{k < n} c(k)\, x^k$.  Otherwise the printer invents an index
variable that clashes with no symbol of the formula.

\subsection{Faithfulness}

The layer changes notation, never the term.  Nothing is simplified or
rewritten: $a + (-b)$ prints as
\[
  a +_S -_S\, b ,
\]
never as $a - b$.  Every subterm the plain printer would show is still
there.

Parentheses follow the binding strengths: $+$ is weakest, then the
products, then unary minus, then the postfix operations.  A chain of
one associative operation of one structure nested to the \emph{left}
drops its parentheses.  Nested to the right it keeps them, so an
associativity law still reads as one:
\[
  a + b + c \;=\; a + (b + c).
\]

\subsection{Conventions: when the subscript is dropped}

The conventions are a list, like a typing block:
\begin{itemize}
\item \texttt{(auto)}, the default.  Drop the subscript when one
  structure is in play in the formula being printed and it is a
  variable, not a named constant such as \texttt{zz-ring}.  Its
  constants $0_S$, $1_S$, $e_S$ keep their subscript when the formula
  also holds a numeral, so $0_K$ never reads as the number $0$.
\item \texttt{(keep)}.  Always write the subscript.
\item \texttt{(drop S \ldots)}.  Drop the subscript of these structure
  terms.
\item \texttt{(symbol HEAD "tex")}.  Write \texttt{HEAD}'s operation
  with this symbol.
\end{itemize}
``The formula being printed'' is the term given to the outermost
printer call.  A display broken over rows is scanned row by row.

\subsection{The tables}

The operations, their TeX and the default conventions are data, not
code.  The base table is \texttt{notation/base.tbl} in the tree; its
head describes the format.  A row reads
\begin{VNBcode}
; KIND     HEAD      PREC  ASSOC  TEX
infix      add       10    assoc  +
postfix    inv       40    -      ^{-1}
outfix     vnrm      100   -      \norm{#1}
apply      series-partial-sum  5  -  \sum_{#i < #2} #1(#i)
\end{VNBcode}
The kinds are \texttt{infix}, \texttt{prefix}, \texttt{postfix},
\texttt{const}, \texttt{outfix} (a template in \texttt{\#1}),
\texttt{bracket} (in \texttt{\#1} and \texttt{\#2}) and \texttt{apply}
(a template for a defined head).  An \texttt{apply} template must use
every argument: it may not hide one.  There are also \texttt{macro}
rows, \texttt{off HEAD}, and \texttt{convention} rows.

Further tables are read after the base table, and a later row replaces
an earlier row of the same head.  The order is:
\begin{enumerate}
\item the tables named, one per line, in
  \texttt{\textasciitilde/vnb-notation/active}.  A name is
  \texttt{\textasciitilde/vnb-notation/NAME.tbl} when that file exists,
  else the tree's \texttt{notation/NAME.tbl};
\item the personal table \texttt{\textasciitilde/vnb-notation/personal.tbl},
  when it exists.
\end{enumerate}
The menu \textsf{VNB > Printer Macros} turns each table on or off, and
\textsf{Edit or New Table\ldots} and \textsf{Edit Personal Table} open
one for editing.

A changed table is re-read at the next printout with the layer on, or
at once by \textsf{Reload} or \texttt{(tex-infix-reload!)}.  A reload
with any row in error changes nothing.  The errors are listed by file
and line, and the table in force stays.


%% =========================================================
\chapter{Checking results and copies}
\label{ch:checking}

Chapter~\ref{ch:classification} sorts results into kinds, and
Chapter~\ref{ch:provenance} describes the records behind them.  This
chapter gives the commands that check them:
\begin{itemize}
\item the exam, which reruns every proof of the library
  (Section~\ref{sec:vnb-test});
\item the REPL commands that show the classification of one result
  (Section~\ref{sec:status-query});
\item the commands for signed trees (Section~\ref{sec:signed-trees-use})
  and for signed releases (Section~\ref{sec:signed-releases-use}).  What
  the two signatures are and what they protect is in
  Sections~\ref{sec:signed-trees} and~\ref{sec:signed-releases}.
\end{itemize}

%% =========================================================

\section{The exam}
\label{sec:vnb-test}

The \emph{exam} is the standing examination of the \textbf{theorem}
bucket.  It is a full library build in which every proof script runs,
every proof page is typed back in and must reach a grounded graph, and
every certificate in \texttt{certificates/} is rewritten from the run.
Definitions and PSS members are excused by construction: the former have
nothing to prove, the latter are accepted on a warrant.

\begin{VNBsnip}
vnb-band WORKER --build --exam
\end{VNBsnip}

\noindent
runs it on a worker machine; \texttt{vnb-nightly-exam} runs it every
night, skips a tree that the last exam already passed, and records the
outcome in \texttt{last-exam.json}.  The everyday build is the
\emph{certified} build: a theorem whose certificate is still valid is
installed without running its proof, and only the theorems whose
statement, citations or definitions changed since the exam are proved
again.  A change to a kernel file invalidates every certificate, so the
everyday build after such a change is itself an exam.

\begin{remark}
The exam is \emph{not} the regression suite.  \texttt{test-suite.scm}
checks that the machinery behaves; the exam checks that the library's
theorems still derive.  The two answer different questions and both are
run.
\end{remark}

\begin{remark}
The script \texttt{./vnb-test} and its stamp
\nolinkurl{reference/certification.scm} are not maintained; the exam is
run by the command above.
\end{remark}

%% =========================================================
\section{Inspecting a result}
\label{sec:status-query}

Three commands expose the classification at the REPL.

\begin{itemize}
\item \texttt{(status 'name)} prints, on one screen, a result's
  provenance, whether it is a PSS member, any warrant it carries, and the
  bill.

\item \texttt{(status-audit)} writes \texttt{reference/STATUS-AUDIT.md}: a
  census of the whole library by kind, flagging hygiene issues --- asserted
  results with no warrant, theorem-shaped names that have drifted out of
  the theorem bucket, and the like.

\item \texttt{(pss-letter 'name)} drafts a plain-English ``recommendation
  letter'' for a candidate PSS member: an English rendering of what it
  says, its current standing, and an assessment against the admission
  gates (true; a sound, terminating rewrite shape; general rather than
  bespoke; non-redundant).  It is an aid to deciding whether a result
  belongs in the curated support set, not an automatic admission.
\end{itemize}


%% =========================================================
\section{Using signed trees}
\label{sec:signed-trees-use}

This section gives the commands for the signed trees of
Section~\ref{sec:signed-trees}.

\paragraph{The pieces.}  Two scripts, \texttt{vnb-trust} and
\texttt{vnb-start}, whose master copies are in the tree under
\texttt{trust/}.  The user copies them once, after checking the
tree, into \texttt{\textasciitilde/bin}, outside the tree, so that a
later change to the tree cannot change them too:
\begin{verbatim}
  cp trust/vnb-trust trust/vnb-start ~/bin/
\end{verbatim}
\begin{itemize}
\item \texttt{vnb-trust} has the commands \texttt{setup KEY} (once),
  \texttt{accept DIR}, \texttt{verify DIR}, \texttt{list} and
  \texttt{release-sign} (the last is for an author signing a
  release).
\item \texttt{vnb-start} runs \texttt{vnb-trust verify} on the tree and
  starts VNB only if it succeeds.  On a mismatch it lists the changed
  files and offers to stop, to fetch a fresh copy, or to start anyway
  once the word \texttt{unverified} is typed in full.
\end{itemize}
The signature is an SSH signature (\texttt{ssh-keygen -Y sign}) in its
own namespace, \texttt{vnb-tree}, so it cannot be mistaken for any
other signature made with the same key.  The hashes are SHA-256.

\paragraph{The manifest.}  To \emph{accept} a tree is to run
\texttt{vnb-trust accept DIR} on it.  The command hashes the files,
shows a summary and asks \texttt{Accept and sign this tree? [y/N]}.  On
\texttt{y}, and the key's passphrase, it writes a \emph{manifest} to
\texttt{\textasciitilde/vnb-signed/}, with its signature beside it: a
few header lines, then one line per file, its SHA-256 hash and its
path.  Every file of the tree is
listed except
\begin{itemize}
\item \texttt{.git/}, \texttt{scratch/}, \texttt{scratchpad/}, and the
  files \texttt{*.bci}, \texttt{*.ext}, \texttt{*.elc};
\item the documents under \texttt{docs/} and \texttt{reference/}, which
  VNB never loads.  The \texttt{.sexp} and \texttt{.scm} files there are
  listed, because the load reads some of them.
\end{itemize}
So the source, the compiled files, the band, the certificates and the
Emacs files are all covered.  A session of VNB does not change any
listed file, so an accepted tree stays accepted while it is used.  The
manifest and its signature stay on the user's computer.

\paragraph{Updating.}  A new copy, a new release or an update fetched
from elsewhere, is compared in the same way and, until the user
accepts it, is refused with the list of files that differ.  The user
then checks it by one of the three ways of Section~\ref{sec:signed-trees}
and accepts it, or keeps
the old copy.

Two rules keep these signatures meaningful.  The copies of \texttt{vnb-trust} and
\texttt{vnb-start} in \texttt{\textasciitilde/bin} must not be replaced
without reading the change.  And the signing key should have a
passphrase.  The setup takes a few commands; they are in the
repository's \texttt{README.md}, under ``Keeping your copy checked''.

\paragraph{For work that must withstand an adversary.}  The steps are
those above, with every choice made the strict way.  Three things
differ from ordinary use.
\begin{itemize}
\item The copy is a copy of a published release, and its signature is
  checked (Section~\ref{sec:signed-releases}) with the author's key
  compared in both places it is published, on GitHub and in the
  author's Bluesky profile.  Ordinary use may
  accept a copy on lighter grounds.
\item Before accepting, the user runs the exam on the user's own
  computer:
\begin{verbatim}
  VNB_CERTIFIED=off ./prover --build-band
\end{verbatim}
  It reruns every proof through the kernel and rebuilds the
  certificates and the band.  This is the step that matters most.  In
  ordinary use the certificates and the band come from the author;
  after the exam, the user relies on nothing the author or the
  author's server did, only on the kernel, whose source is in the
  tree.
\item VNB is started only with \texttt{vnb-start}, never directly.
\end{itemize}
The limits listed under ``What it protects, and what it does not''
(Section~\ref{sec:signed-trees}) still apply.


%% =========================================================
\section{Making and checking a release}
\label{sec:signed-releases-use}

This section gives the commands for the signed releases of
Section~\ref{sec:signed-releases}.

\paragraph{How the author makes one.}  The order is:
\begin{enumerate}
\item the author checks the tree, by the exam, and accepts it
  (Section~\ref{sec:signed-trees-use});
\item on the machine that holds the repository, \texttt{vnb-release-manifest} \textit{VERSION}
  writes the manifest straight from git's object store;
\item on the author's computer, \texttt{vnb-trust release-sign}
  compares every file of the manifest with the same file in the
  accepted tree, and signs only if all agree.  Files under
  \texttt{scratch/} and \texttt{scratchpad/}, which VNB does not load,
  are counted and shown, not compared;
\item the signature is copied back, both files are committed, and the
  release commit is tagged and pushed.
\end{enumerate}
So a release is exactly a tree the author accepted.

\paragraph{How anyone checks one.}  With standard tools only:
\texttt{ssh-keygen -Y verify} for the signature, then
\texttt{git ls-files -z | xargs -0 sha256sum} at the release commit,
compared with the manifest.  The commands are in the repository's
\texttt{README.md}, under ``Verifying a release''.



%% =========================================================
\appendix
\chapter{Source file overview}
\label{app:source}

The prover sources live in \texttt{\textasciitilde/prover/}.  This appendix is a
map of the tree, not a design document: it says what each file \emph{is}, so a
reader following a cross-reference or extending the system can find the right
place.  The layout is:

\begin{itemize}
\item \textbf{Root} (\texttt{\textasciitilde/prover/*.scm}) --- the proof
  kernel, the reader and library record, the interactive command layer, and the
  documentation, proof-printing, and proof-driving subsystems.
\item \textbf{\texttt{structure-library/}} --- definitions, structures,
  vocabulary, and warranted supports: the mathematical objects the system knows
  at startup.
\item \textbf{\texttt{theorem-library/}} --- proofs that reach \texttt{(qed
  \ldots)}, loaded, gated, and counted.
\item \textbf{\texttt{calculus/}} --- probes and stress tests.  \emph{Mostly}
  outside the load path, but not wholly: \texttt{load.scm} names
  \texttt{calculus/compact-tb-proof},
  \texttt{calculus/compact-complete-proof} and
  \texttt{calculus/finite-ball-subcover-proof}, and it gives every
  \texttt{calculus/} file the same fresh-environment containment it gives
  \texttt{theorem-library/} precisely because such a file can be loaded.  An
  edit here is therefore not automatically inert.
\item \textbf{\texttt{reference/}} --- for the most part generated at every
  load (\texttt{PSS.md}, \texttt{THEOREMS.md}, \texttt{DEFINITIONS.md},
  \texttt{PROOF-DEBT.md}, \texttt{TACTICS.md}, \ldots), and those files must
  not be edited by hand; \texttt{LIBRARY.md}, \texttt{KERNEL.md},
  \texttt{KERNEL-RULES.md}, \texttt{REVIEW.md}, \texttt{USABILITY-REVIEW.md}
  and \texttt{VNB-TEST.md} are written by hand.
\item \textbf{\texttt{certificates/}} --- the proof-certificate store, one
  file per proof file.
\item \textbf{\texttt{libraries.sexp}} --- the assignment of files to
  libraries.
\item \textbf{\texttt{emacs/}} --- the workspace.
\item \textbf{\texttt{preambles/}} --- the preamble rule files.
\item \textbf{\texttt{notation/}} --- the printer tables (\texttt{base.tbl} and others).
\item \textbf{\texttt{docs/}} --- this manual and the design notes.
\item \textbf{\texttt{archive/}} --- retired and spliced originals.
\end{itemize}

\begin{remark}
The file listings below are \emph{generated} from the tree by
\texttt{docs/gen-source-appendix.py}, which harvests each file's
\texttt{;;;\ \textit{name} -- \textit{description}} header.  Regenerate after
adding or renaming files so this appendix cannot drift from the sources it
describes:
\begin{VNBsnip}
python3 docs/gen-source-appendix.py
\end{VNBsnip}
\end{remark}

%% AUTO-GENERATED by docs/gen-source-appendix.py -- do not edit by hand.
%% Regenerate:  python3 docs/gen-source-appendix.py
%% Descriptions are harvested from each file's `;;; name -- desc' header.

\section{Kernel and subsystems (\texttt{\textasciitilde/prover/})}
\label{sec:src-root}
The 82 Scheme files at the root of the tree: the proof kernel, the reader, the theory record, the interactive command layer, and the documentation / proof-printing / driver subsystems.

\begin{longtable}{@{}p{0.30\textwidth}p{0.63\textwidth}@{}}
\toprule
\textbf{File} & \textbf{Description} \\
\midrule
\endhead
\texttt{all-theorems.scm} & THE BOOK OF ALL PROVEN THEOREMS (the user's roadmap, \\
\texttt{arith-eval.scm} & ground arithmetic decision procedure \\
\texttt{audit.scm} & library hygiene diagnostics \\
\texttt{band-dump.scm} & dump what a saved band holds, one deterministic line per name \\
\texttt{calc.scm} & the ``directive chain checker'' (notes-27): ground a goal by a \\
\texttt{certificates.scm} & PROOF CERTIFICATES: a load that runs no proof \\
\texttt{cite-index.scm} & MEASURED citation affinity: P(theorem cited | head in goal) \\
\texttt{clobber-guard.scm} & the gate CLAUDE.md says does not exist \\
\texttt{contexts.scm} & make-wff, the interactive entry point for building a <wff> \\
\texttt{contra.scm} & (contra): close a branch whose CONTEXT is arithmetically \\
\texttt{core-capture.scm} & CORE REPLAY, steps 1-3: capture, the kernel bill, the lowering \\
\texttt{core-replay-checks.scm} & the suite checks of the CORE REPLAYER (core-replay.scm), loaded by \\
\texttt{core-replay.scm} & the CORE REPLAYER (core replay, design step 4) and the \\
\texttt{counterexample.scm} & probes that try to REFUTE the focus goal \\
\texttt{deduction-graphs.scm} & deduction graph: sequent nodes and inference nodes \\
\texttt{driver-kit.scm} & the proof-driving helpers that more than one file needs \\
\texttt{equal-patch.scm} & MIT/GNU Scheme 12.1's equal? without its per-call hash table \\
\texttt{errors.scm} & VNB error value and guard wrapper \\
\texttt{expressions.scm} & FOL expressions over VNB-style set theory \\
\texttt{extend-band.scm} & extend a saved band instead of rebuilding it (2026-09-23) \\
\texttt{functoid-gate.scm} & the def-functoid DEFINITION LOG and the load-end gate on \\
\texttt{glossary.scm} & THE ALPHABETICAL GLOSSARY: every name in VNB, in one list \\
\texttt{harvest.scm} & MEASURE the What Now panel against the library's own proofs \\
\texttt{ineq-supply.scm} & the What Now lane that SUPPLIES the inequalities the \\
\texttt{input-context.scm} & generalized context-discharge + operator-remap engine \\
\texttt{interactive.scm} & short-form commands for interactive proof sessions \\
\texttt{kernel-map-static.scm} & INSTRUMENT, not library \\
\texttt{kernel-map-trace.scm} & INSTRUMENT, not library \\
\texttt{libraries.scm} & THE LIBRARIES (the user, notes-61, 2026-10-06: "start thinking of \\
\texttt{library-book.scm} & THE BOOK OF THEOREMS AND DEFINITIONS, NO PROOFS (the user, \\
\texttt{library.scm} & theory management \\
\texttt{load.scm} & load the proof checker in dependency order \\
\texttt{macetes.scm} & macete mechanism \\
\texttt{major-theorems.scm} & the LaTeX list of MAJOR THEOREMS \\
\texttt{minimize.scm} & `minimize!', the least-measure tactic \\
\texttt{mint-attribution.scm} & \textit{(no header)} \\
\texttt{mp.scm} & `(mp)': close a goal that is an INSTANCE of a universal you already \\
\texttt{mutation-check.scm} & IS THE MUST-NOT-PROVE CORPUS ACTUALLY DETECTING ANYTHING? \\
\texttt{number-systems.scm} & number system constants and axioms \\
\texttt{operators.scm} & ONE table, keyed by head symbol \\
\texttt{page-audit-support.scm} & \textit{(no header)} \\
\texttt{page-audit.scm} & THE GATE ON THE PRINTED PAGE \\
\texttt{parser.scm} & VNB string expression parser \\
\texttt{parts.scm} & accessing the PARTS of a formula by SURFACE syntax \\
\texttt{preamble.scm} & rule-based drivers a human can edit (notes-45, batch 41, 2026-09-28), \\
\texttt{prep.scm} & (prep 'FUBA): why will FUBA not fire here, and what must be done \\
\texttt{presentation.scm} & THE RESOLUTION DIAL \\
\texttt{primitive-inferences.scm} & the primitive inference kernel \\
\texttt{proof-commands.scm} & higher-level proof commands (tactics) \\
\texttt{proof-debt.scm} & the warrant / proof-debt ledger \\
\texttt{proof-map.scm} & (proof-map): see the WHOLE proof tree, closed branches \\
\texttt{proof-reader.scm} & the HUMAN-READABLE proof sketch (``reader mode''), the \\
\texttt{proof-tex.scm} & render a completed proof as a LaTeX step-trace \\
\texttt{prop.scm} & `prop': decide a propositional entailment, then PROVE it \\
\texttt{proven-theorems.scm} & TOMBSTONE (2026-05-27) \\
\texttt{push-not.scm} & (push-not-h k): push a NOT inward through one quantifier or \\
\texttt{ra-diverge.scm} & \textit{(no header)} \\
\texttt{reference-html.scm} & the browser pages of the VNB reference, built by \\
\texttt{replay-audit.scm} & does a recorded proof SCRIPT replay to a grounded graph? \\
\texttt{rule-checkers-logic.scm} & the 42 `logic' kernel operations, CHECKED \\
\texttt{rule-checkers-oracle.scm} & inference checkers for the seven ORACLE rules \\
\texttt{rule-checkers-rewrite.scm} & checkers for the four `rewrite' kernel operations \\
\texttt{rule-checkers-schema.scm} & checkers for the thirteen `schema' kernel operations \\
\texttt{sequents.scm} & sequents over <wff> objects \\
\texttt{sketch.scm} & the structured-proof surface (sketch/step/obtain/qed-sketch), \\
\texttt{statement-types.scm} & THE IMPLIED TYPES OF A STATEMENT (the user's request, 2026-10-17) \\
\texttt{strapp-fast.scm} & \textit{(no header)} \\
\texttt{strapp-patch.scm} & MIT/GNU Scheme 12.1's string-append without its two 1024-slot \\
\texttt{structures.scm} & mathematical structure definitions \\
\texttt{suggest.scm} & assumption-pattern scanner for forward-move discovery \\
\texttt{tactic-uses-data.scm} & GENERATED by docs/gen-tactic-uses.py from \\
\texttt{tactics-help.scm} & a self-describing registry of the interactive tactics \\
\texttt{test-suite-negative.scm} & the MUST-NOT-PROVE corpus \\
\texttt{test-suite.scm} & VNB system regression and demonstration tests \\
\texttt{tex-infix.scm} & THE OPTIONAL INFIX LAYER for the TeX printer (printer side) \\
\texttt{tex-output.scm} & render VNB s-expression formulas as LaTeX math \\
\texttt{transport.scm} & TRANSPORTABLE FACTS \\
\texttt{vlet.scm} & (vlet (names...) FORMER): bind names to proof objects \\
\texttt{wff-english.scm} & render a wff as an English sentence \\
\texttt{wff.scm} & the <wff> record type (well-formed formula in a theory) \\
\texttt{witness-tactics.scm} & getting a witness out of a quantified hypothesis \\
\texttt{zero-it.scm} & (zero-it): move an equation goal P = Q to its normal form R = 0 \\
\bottomrule
\end{longtable}

\section{\texttt{structure-library/}}
\label{sec:src-structure-library}
The 322 files that install definitions, structures, vocabulary, and warranted supports.  One file per structure or cluster; the always-current per-structure listing (declaration, slots, parent, theorems) is the generated \texttt{reference/STRUCTURE-INDEX.md}.

\begin{longtable}{@{}p{0.30\textwidth}p{0.63\textwidth}@{}}
\toprule
\textbf{File} & \textbf{Description} \\
\midrule
\endhead
\texttt{abelian-group.scm} & ABELIAN-GROUP structure \\
\texttt{algnum-vocab.scm} & ALGEBRAIC ELEMENTS OF C (ALGNUM, algebraic-numbers bundle B7, 2026-10-26) \\
\texttt{analytic-log.scm} & ANALYTICITY AT A POINT (the notes' Def 2.11) AND THE \\
\texttt{antiderivative.scm} & Definition 4.6 (calculus.pdf), as VOCABULARY \\
\texttt{ascoli-arzela-statement.scm} & \textit{(no header)} \\
\texttt{ba-amp-vocab.scm} & Thayer, Banach Algebras and Applications, ch. 6 (the amplification in the proof of \\
\texttt{ba-bh2-vocab.scm} & wave of 2026-11-09/10 (agent BA-BH2), wired 2026-11-10 from scratchpad/babh2-1109/ba-bh2-vocab.scm \\
\texttt{ba-cx-vocab.scm} & wave of 2026-11-09/10 (agent BA-1), wired 2026-11-10 from scratchpad/ba1-1109/ba-cx-vocab.scm \\
\texttt{ba-gns-vocab.scm} & THE GNS CONSTRUCTION, the pre-Hilbert stage: BanachAlgebras.pdf Theorem 4.13 and \\
\texttt{ba-gns2-vocab.scm} & THE GNS CONSTRUCTION ON THE HILBERT COMPLETION and THE RADICAL: \\
\texttt{ba-neumann-vocab.scm} & wave of 2026-11-09/10 (agent BA-1), wired 2026-11-10 from scratchpad/ba1-1109/ba-neumann-vocab.scm \\
\texttt{ba-ops-vocab.scm} & wave of 2026-11-09/10 (agent BA-BH), wired 2026-11-10 from scratchpad/babh-1109/ba-ops-vocab.scm \\
\texttt{ba-optop-vocab.scm} & wave of 2026-11-09/10 (agent BA-OT1), wired 2026-11-10 from scratchpad/baot1-1109/ba-optop-vocab.scm \\
\texttt{ba-quotient-na-vocab.scm} & wave of 2026-11-09/10 (agent BA-1E), wired 2026-11-10 from scratchpad/ba1e-1109/ba-quotient-na-vocab.scm \\
\texttt{ba-quotient-vocab.scm} & wave of 2026-11-09/10 (agent BA-1C), wired 2026-11-10 from scratchpad/ba1c-1109/ba-quotient-vocab.scm \\
\texttt{ba-rep-vocab.scm} & POSITIVE FUNCTIONALS, STATES AND REPRESENTATIONS of a unital involutive \\
\texttt{ba-trace-dual-vocab.scm} & Thayer, Banach Algebras and Applications, chapter 5 s.6-s.8 (spaces of \\
\texttt{ba-unbdd-vocab.scm} & wave of 2026-11-09/10 (agent BA-G1), wired 2026-11-10 from scratchpad/bag1-1109/ba-unbdd-vocab.scm \\
\texttt{ba-vn-vocab.scm} & Thayer, Banach Algebras and Applications, chapter 6 (density theorems): the \\
\texttt{ba-vocab.scm} & wave of 2026-11-09/10 (agent BA-1), wired 2026-11-10 from scratchpad/ba1-1109/ba-vocab.scm \\
\texttt{ba2-vocab.scm} & wave of 2026-11-09/10 (agent BA-2), wired 2026-11-10 from scratchpad/ba2-1109/ba2-vocab.scm \\
\texttt{ba2b-gen-vocab.scm} & wave of 2026-11-09/10 (agent BA-2B), wired 2026-11-10 from scratchpad/ba2b-1109/ba2b-gen-vocab.scm \\
\texttt{baire-category.scm} & interior/closure, nowhere-dense, meager/nonmeager, and \\
\texttt{basw-vocab.scm} & wave of 2026-11-09/10 (agent BA-SW), wired 2026-11-10 from scratchpad/basw-1109/basw-vocab.scm \\
\texttt{basw2-vocab.scm} & wave of 2026-11-09/10 (agent BA-SW2), wired 2026-11-10 from scratchpad/basw2-1109/basw2-vocab.scm \\
\texttt{basw3-vocab.scm} & wave of 2026-11-09/10 (agent BA-SW3), wired 2026-11-10 from scratchpad/basw3-1109/basw3-vocab.scm \\
\texttt{bidual.scm} & \textit{(no header)} \\
\texttt{bijection-derived.scm} & the three BIJECTION projection lemmas, PROVEN \\
\texttt{bijection.scm} & bijections between classes, and INVERSE-BIJ via CHOICE \\
\texttt{binomial.scm} & the Binomial Theorem for commutative rings, over the \\
\texttt{block-family-combinatorial.scm} & \textit{(no header)} \\
\texttt{bounded-metric.scm} & the bounded metric  rho = d/(1+d)  of a metric space, \\
\texttt{c-metric-space.scm} & C-METRIC-SPACE, a set carrying a COUNTABLE FAMILY OF \\
\texttt{ca-ac-vocab.scm} & UNIFORM CONVERGENCE ON COMPACT SUBSETS of an open set of CC \\
\texttt{ca-angle-vocab.scm} & WAVE 2 (2026-10-27..30), wired 2026-10-30 from scratchpad/ch3e/ch3e-angle.scm (vocab) \\
\texttt{ca-homotopy-vocab.scm} & \textit{(no header)} \\
\texttt{ca-hopf-vocab.scm} & THE HOPF FIBRATION, complex-analysis.pdf s. 1.6.3 ((13)-(16), Prop. 1.22, \\
\texttt{ca-lft-vocab.scm} & THE EXTENDED PLANE AND LINEAR FRACTIONAL TRANSFORMATIONS, \\
\texttt{ca-line-vocab.scm} & WAVE 2 (2026-10-27..30), wired 2026-10-30 from scratchpad/ch3e/ch3f-line.scm (vocab), scratchpad/ch3e/ch3f-refl.scm (vocab) \\
\texttt{ca1-vocab.scm} & THE COMPLEX NUMBERS AS 2-BY-2 REAL MATRICES (complex-analysis.pdf, \\
\texttt{ca5-function-spaces.scm} & THE VOCABULARY OF CHAPTER 5 (FUNCTION SPACES) of the \\
\texttt{ca51-vocab.scm} & local compactness and exhaustions by compact sets (CA-51, 2026-10-26) \\
\texttt{calc3-vocab.scm} & vocabulary for calculus.pdf chapter 3 (metric topology): \\
\texttt{calint-vocab.scm} & WAVE 2 (2026-10-27..30), wired 2026-10-30 from scratchpad/calint/calint-vocab.scm \\
\texttt{cardinality.scm} & CARD, the cardinality of a set, DEFINED \\
\texttt{catools-gal-vocab.scm} & divisibility in an integral domain \\
\texttt{cauchy-subsequence.scm} & \textit{(no header)} \\
\texttt{cc-coords.scm} & CC and cartesian(RR, RR), as VOCABULARY \\
\texttt{cc-elementary.scm} & THE ELEMENTARY FUNCTIONS exp, cos, sin ON CC, DEFINED BY \\
\texttt{cc-pi.scm} & THE CONSTANT pi, DEFINED FROM THE KERNEL OF exp (the notes 2.3.3) \\
\texttt{cc-power-series.scm} & THE VOCABULARY OF COMPLEX POWER SERIES (the notes 2.2) \\
\texttt{ccr-ch3-vocab.scm} & wave of 2026-11-09/10 (agent CCR-1109-B), wired 2026-11-10 from scratchpad/ccr1109b/ccr-ch3-vocab.scm \\
\texttt{ccr-ch3b-vocab.scm} & wave of 2026-11-09/10 (agent CCR-1109-B2), wired 2026-11-10 from scratchpad/ccr1109b2/ccr-ch3b-vocab.scm \\
\texttt{ccr-ch3c-vocab.scm} & wave of 2026-11-09/10 (agent CCR-1109-B3), wired 2026-11-10 from scratchpad/ccr1109b3/ccr-ch3c-vocab.scm \\
\texttt{ccr1-vocab.scm} & ULTRAMETRIC SPACES (CCR.pdf, chapter 1, Definition 1.1 and the objects of \\
\texttt{ccr1109a2-vocab.scm} & wave of 2026-11-09/10 (agent CCR-1109-A2), wired 2026-11-10 from scratchpad/ccr1109a2/ccr1109a2-vocab.scm \\
\texttt{ccr1109a3-laurent-vocab.scm} & wave of 2026-11-09/10 (agent CCR-1109-A3), wired 2026-11-10 from scratchpad/ccr1109a3/ccr1109a3-laurent-vocab.scm \\
\texttt{ccr1109a4-laurent-lim-vocab.scm} & wave of 2026-11-09/10 (agent CCR-1109-A4), wired 2026-11-10 from scratchpad/ccr1109a4/ccr1109a4-laurent-lim-vocab.scm \\
\texttt{ccr1109a5-vocab.scm} & wave of 2026-11-09/10 (agent CCR-1109-A5), wired 2026-11-10 from scratchpad/ccr1109a5/ccr1109a5-vocab.scm \\
\texttt{ccr1109a6-vocab.scm} & wave of 2026-11-09/10 (agent CCR-1109-A6), wired 2026-11-10 from scratchpad/ccr1109a6/ccr1109a6-vocab.scm \\
\texttt{ccr2-vocab.scm} & NON-ARCHIMEDEAN VALUED FIELDS AND THEIR DOMAINS (CCR.pdf chapters 1-2 and 4: \\
\texttt{ccr3-vocab.scm} & MODULES IN A LOCAL NV FIELD (CCR.pdf Definition 4.11) \\
\texttt{ccr4-vocab.scm} & ADDITIVE CHARACTERS OF A NORMED FIELD AND THE DUALITY MAP (CCR.pdf chapter 5, Definitions \\
\texttt{ccr4b-vocab.scm} & wave of 2026-11-09/10 (agent CCR-1109-C3), wired 2026-11-10 from scratchpad/ccr1109c3/ccr4b-vocab.scm \\
\texttt{ccr4c-subnf-vocab.scm} & wave of 2026-11-09/10 (agent CCR-1109-C4), wired 2026-11-10 from scratchpad/ccr1109c4/ccr4c-subnf-vocab.scm \\
\texttt{ccr4d-ext-vocab.scm} & wave of 2026-11-09/10 (agent CCR-1109-C5), wired 2026-11-10 from scratchpad/ccr1109c5/ccr4d-ext-vocab.scm \\
\texttt{ccr4f-vocab.scm} & THE DIGIT MAP AND THE COORDINATE MAP AS FUNCTIONS (CCR-1111 finding F1) \\
\texttt{ccr4x-ext-vocab.scm} & wave of 2026-11-09/10 (agent CCR-1109-C2), wired 2026-11-10 from scratchpad/ccr1109c2/ccr4x-ext-vocab.scm \\
\texttt{ccr4x-vocab.scm} & wave of 2026-11-09/10 (agent CCR-1109-C), wired 2026-11-10 from scratchpad/ccr1109c/ccr4x-vocab.scm \\
\texttt{ccr5-vocab.scm} & wave of 2026-11-09/10 (agent CCR-1109-D), wired 2026-11-10 from scratchpad/ccr1109d/ccr5-vocab.scm \\
\texttt{ccr6e-vocab.scm} & wave of 2026-11-09/10 (agent CCR-1109-D3), wired 2026-11-10 from scratchpad/ccr1109d3/ccr6e-vocab.scm \\
\texttt{cg1-vocab.scm} & vocabulary for the compact-groups notes, chapter 5 section 2 (CG-1) \\
\texttt{cg2-vocab.scm} & set products and set powers in a group (compact-groups notes, Prop 5.8: \\
\texttt{cg4-vocab.scm} & compact-groups notes, sections 1 and 5 (CG-4 session) \\
\texttt{ch3a-vocab.scm} & automorphisms of an open subset of CC (CH3-A, 2026-10-26) \\
\texttt{ch4a-vocab.scm} & WAVE 2 (2026-10-27..30), wired 2026-10-30 from scratchpad/ch4a/ch4a-vocab.scm \\
\texttt{ch4b-vocab.scm} & WAVE 2 (2026-10-27..30), wired 2026-10-30 from scratchpad/ch4b/ch4b-vocab.scm \\
\texttt{ch4d-vocab.scm} & WAVE 2 (2026-10-27..30), wired 2026-10-30 from scratchpad/ch4d/ch4d-vocab.scm \\
\texttt{cnvs-ba-vocab.scm} & wave of 2026-11-09/10 (agent CNVS), wired 2026-11-10 from scratchpad/cnvs-1109/cnvs-ba-vocab.scm \\
\texttt{cnvs-comp-vocab.scm} & wave of 2026-11-09/10 (agent CNVS2), wired 2026-11-10 from scratchpad/cnvs2-1109/cnvs2-comp-vocab.scm \\
\texttt{cnvs-vocab.scm} & wave of 2026-11-09/10 (agent CNVS), wired 2026-11-10 from scratchpad/cnvs-1109/cnvs-vocab.scm \\
\texttt{comm-ring-simplify.scm} & commutative-ring identity decision procedure \\
\texttt{commutative-ring.scm} & COMMUTATIVE-RING: a ring whose MUL is commutative \\
\texttt{compactness.scm} & compactness of a metric space and the four-way \\
\texttt{complex-inner-product.scm} & COMPLEX-INNER-PRODUCT-SPACE: a vector space over \\
\texttt{complex-measure.scm} & WAVE 2 (2026-10-27..30), wired 2026-10-30 from scratchpad/mf/mf-cvocab.scm \\
\texttt{complex.scm} & CC as a metric space \\
\texttt{compose-typing.scm} & RAN (range) + mindless type-closure lemmas for \\
\texttt{compose.scm} & COMPOSE: composition of functions, f o g \\
\texttt{convex-functional.scm} & \textit{(no header)} \\
\texttt{cvx-vocab.scm} & convex combinations and the convex hull in a real normed vector space \\
\texttt{definitional-reclass.scm} & reclassify hand-written DEFINING iffs \\
\texttt{derivative.scm} & the Caratheodory derivative, as VOCABULARY \\
\texttt{determinant.scm} & the determinant of a square matrix over a ring, by \\
\texttt{diff-on.scm} & THE DERIVATIVE ON AN OPEN SET, OVER A NORMED FIELD \\
\texttt{doubling-functions.scm} & WAVE 2 (2026-10-27..30), wired 2026-10-30 from scratchpad/mf/mf-mivocab.scm \\
\texttt{dual-space.scm} & THE DUAL of a real normed vector space, as an OBJECT \\
\texttt{elementary-matrix.scm} & matrix units and elementary column operations \\
\texttt{euclidean-ring.scm} & EUCLIDEAN-RING: an integral domain admitting a \\
\texttt{extended-arith.scm} & the rest of the arithmetic and order vocabulary of \\
\texttt{extended-reals-pos.scm} & RR-POS-STAR = [0, +inf], the nonnegative extended reals \\
\texttt{extended-reals-updown.scm} & \textit{(no header)} \\
\texttt{extended-reals.scm} & the extended real line RR-STAR = [-inf, +inf] \\
\texttt{extended-sum.scm} & \textit{(no header)} \\
\texttt{extreme-value.scm} & the closed interval [a,b] and the Extreme Value Theorem \\
\texttt{field.scm} & FIELD: a commutative ring in which every nonzero element \\
\texttt{finite-dimensional.scm} & finite-dimensional vector spaces, Zorn-free \\
\texttt{finprod.scm} & FINPROD / PROD-RING: finite PRODUCT of a function over a \\
\texttt{finsum.scm} & FINSUM: summation of an abelian-group-valued function \\
\texttt{flat-regular.scm} & PUBLIC-KEY SPACES WHOSE REGULAR CONTENTS ARE FLAT: the one property of the \\
\texttt{frechet-banach-steinhaus.scm} & BOUNDED SETS and EQUICONTINUOUS FAMILIES of linear \\
\texttt{frechet-metric.scm} & the METRIC of a seminormed space whose seminorm family is \\
\texttt{frechet-open-mapping.scm} & Frechet spaces, continuous linear maps, and the \\
\texttt{functor-invariance.scm} & a construction that cannot see the functor \\
\texttt{functoriality.scm} & every view-as is a FUNCTOR, and here is the proof \\
\texttt{gal-ch-vocab.scm} & vocabulary for the general Cayley-Hamilton theorem (GALOIS B6, 2026-10-26) \\
\texttt{gal-char-vocab.scm} & prime natural numbers and the characteristic of a field (galois-theory notes, \\
\texttt{gal-charpoly-vocab.scm} & the companion matrix and the characteristic polynomial (galois-theory notes, \\
\texttt{gal-gcd-vocab.scm} & a greatest common divisor of a set of ring elements (galois-theory notes, \\
\texttt{gal-op-vocab.scm} & a linear operator T on a module E over F as an F[x]-module: the iterates T\textasciicircum{}k, the \\
\texttt{gal-poly-vocab.scm} & evaluation of a formal polynomial and the polynomial x - c \\
\texttt{gal-ufd-vocab.scm} & associates and unique factorization domains (galois-theory notes, section 1) \\
\texttt{gamma-continuation-vocab.scm} & WAVE 2 (2026-10-27..30), wired 2026-10-30 from scratchpad/ch3g/ch3g-gamma.scm (vocab) \\
\texttt{gamma-function.scm} & THE IMPROPER INTEGRAL ON (0, +inf) AND THE GAMMA FUNCTION \\
\texttt{gn-vocab.scm} & BANACH ALGEBRAS: the natural order of an involutive algebra (Def. 1.34, (15)) and the \\
\texttt{group-cosets-vocab.scm} & \textit{(no header)} \\
\texttt{group-zpow-vocab.scm} & integer powers in a group (ALGNUM, 2026-10-27) \\
\texttt{group.scm} & GROUP structure \\
\texttt{hahn-banach-sublinear.scm} & \textit{(no header)} \\
\texttt{hausdorff-vocab.scm} & \textit{(no header)} \\
\texttt{hc-vocab.scm} & wave of 2026-11-09/10 (agent HAAR-C), wired 2026-11-10 from scratchpad/haarc-1109/hc-vocab.scm \\
\texttt{hc2-vocab.scm} & wave of 2026-11-09/10 (agent HAAR-C2), wired 2026-11-10 from scratchpad/haarc2-1109/hc2-vocab.scm \\
\texttt{hc3-vocab.scm} & wave of 2026-11-09/10 (agent HAAR-C3), wired 2026-11-10 from scratchpad/haarc3-1109/hc3-vocab.scm \\
\texttt{hilbert-completion-vocab.scm} & \textit{(no header)} \\
\texttt{hilbert-onb-vocab.scm} & wave of 2026-11-09/10 (agent BA-ONB), wired 2026-11-10 from scratchpad/baonb-1109/hilbert-onb-vocab.scm \\
\texttt{hilbert-space.scm} & HILBERT SPACES: the metric space of a complex inner \\
\texttt{id-fun.scm} & ID-FUN: the identity map of a set \\
\texttt{ideal.scm} & ideals, principal ideals, and principal-ideal domains over a \\
\texttt{ineq-oracle.scm} & the (ineq) decision procedure: close a linear-inequality \\
\texttt{inf-subsets.scm} & the class of infinite subsets of a given set \\
\texttt{infinite-products-vocab.scm} & WAVE 2 (2026-10-27..30), wired 2026-10-30 from scratchpad/ch6/ch6-prod.scm (vocab) \\
\texttt{injection.scm} & INJECTION class, IMAGE operator, and the \\
\texttt{integral-domain-laws.scm} & integral-domain-cancel-zero, PROVEN from the \\
\texttt{integral-domain.scm} & INTEGRAL-DOMAIN: a commutative ring with ONE /= ZERO \\
\texttt{integral.scm} & \textit{(no header)} \\
\texttt{intersection-of.scm} & INTERSECTION-OF(c), the intersection of the members \\
\texttt{interval-calculus.scm} & THE VOCABULARY OF THE CALCULUS ON AN INTERVAL \\
\texttt{interval-extremum.scm} & LOCAL MAXIMUM AND LOCAL MINIMUM OF A FUNCTION ON \\
\texttt{interval-taylor.scm} & THE VOCABULARY OF TAYLOR'S FORMULA ON AN INTERVAL \\
\texttt{ipf-vocab.scm} & INNER PRODUCT FORMS ON A COMPLEX MODULE, the module carrying such a form, and the \\
\texttt{is-set-laws.scm} & \textit{(no header)} \\
\texttt{isom-vocab.scm} & vocabulary for calculus.pdf ch. 3, sections 6 and 8: uniform and \\
\texttt{isom2-rrvs-vocab.scm} & the real numbers as a FIELD instance, the vector space of \\
\texttt{isom2-vocab.scm} & the group structure of the isometry group (calculus.pdf ch. 3 \\
\texttt{km-vocab.scm} & wave of 2026-11-09/10 (agent BA-KM), wired 2026-11-10 from scratchpad/bakm-1109/km-vocab.scm \\
\texttt{l1-dual.scm} & WAVE 2 (2026-10-27..30), wired 2026-10-30 from scratchpad/md/me-du-vocab.scm \\
\texttt{l2-families-vocab.scm} & wave of 2026-11-09/10 (agent BA-ONB2), wired 2026-11-10 from scratchpad/baonb2-1109/bo2-l2-vocab.scm \\
\texttt{l2-integral-operator.scm} & wave of 2026-11-09/10 (agent TH3-L2C), wired 2026-11-10 from scratchpad/th3-1109/l2c/kc-vocab.scm \\
\texttt{l2c-vocab.scm} & wave of 2026-11-09/10 (agent TH3-L2), wired 2026-11-10 from scratchpad/th3-1109/l2/l2c-vocab.scm \\
\texttt{la-det-vocab.scm} & triangular matrices and the formal adjoint (LA-DET, 2026-10-18) \\
\texttt{la-elem-vocab.scm} & elementary column matrices as a class, the subgroup they \\
\texttt{la-z-vocab.scm} & subdeterminants, equivalence by elementary operations and the \\
\texttt{la-z2-vocab.scm} & the principal sequence of a subgroup (linear algebra, the \\
\texttt{la-z3-vocab.scm} & quotient modules, the ring as a module over itself, the cyclic \\
\texttt{la-z4-vocab.scm} & the free module Z\textasciicircum{}n, the image of a sequence under a map, the \\
\texttt{la-z5-vocab.scm} & the norm N\_E(F) of a submodule (linear algebra, the definition \\
\texttt{la-z6-vocab.scm} & polynomials in n indeterminates and the permutation of the indeterminates \\
\texttt{lebesgue-translate.scm} & the TRANSLATE of a set of reals \\
\texttt{line-coefficient.scm} & \textit{(no header)} \\
\texttt{linear-arith.scm} & standalone linear-arithmetic decision core \\
\texttt{linear-functional.scm} & bounded linear functionals and the dual norm on a \\
\texttt{list-recursion.scm} & CONS, and the recursive characterization of LENGTH \\
\texttt{little-o.scm} & calculus.pdf Section 2.2, the calculus of o/O notation, in a \\
\texttt{local-measure.scm} & WAVE 2 (2026-10-27..30), wired 2026-10-30 from scratchpad/mf/mf-mgvocab.scm \\
\texttt{lp-inequalities.scm} & POWERS IN [0,+inf] AND THE p-TH POWER INTEGRAL I\_p \\
\texttt{lp-seminorm.scm} & THE L\textasciicircum{}p SEMINORM OF A REAL-VALUED FUNCTION \\
\texttt{m-riesz-vocab.scm} & ORDERED VECTOR SPACES: BanachAlgebras.pdf section 4.6, before Proposition 4.32 \\
\texttt{mat-equiv.scm} & the matrix EQUIVALENCE relation \textasciitilde{} (algebraic-numbers.pdf \\
\texttt{matrix.scm} & \textit{(no header)} \\
\texttt{mb-vocab.scm} & vocabulary for the M-B wave (construction-of-measures.pdf 2.9, 3.2-3.4, \\
\texttt{mb12-vocab.scm} & vocabulary for construction-of-measures.pdf Appendix A.3: the \\
\texttt{mb35-souslin-vocab.scm} & vocabulary for construction of measures, chapter 5 section 1.3-1.4: sieves and \\
\texttt{mb35-vocab.scm} & vocabulary for construction of measures, chapter 5 section 2 (M-B35, 2026-10-26) \\
\texttt{mb4-vocab.scm} & vocabulary for construction of measures, chapter 5 (M-B4, 2026-10-26) \\
\texttt{mb67-vocab.scm} & vocabulary for construction of measures, chapter 5 sections 1.4-1.6: the Baire space N\textasciicircum{}N \\
\texttt{measurable-space.scm} & MEASURABLE-SPACE and MEASURE-SPACE as declared \\
\texttt{measure-ch1-vocab.scm} & ATOMS OF AN ALGEBRA, FINITE UNIONS, RECTANGLES OF R\textasciicircum{}n, THE PULLBACK SIGMA-ALGEBRA AND \\
\texttt{measure-step.scm} & WAVE 2 (2026-10-27..30), wired 2026-10-30 from scratchpad/md/md-vocab.scm \\
\texttt{measure.scm} & OUTER MEASURES, CARATHEODORY MEASURABILITY, MEASURES \\
\texttt{meta-core.scm} & THE CORE OF K2: well-formedness, sequents, the core rule relations, \\
\texttt{meta-cr.scm} & CONTEXTUAL REWRITING (CR) over the encoded core: the one-step rewrite relation \\
\texttt{meta-def.scm} & K2 PHASE 5: the clauses of the kernel's definedness certificate pi--defined? \\
\texttt{meta-ind.scm} & K2 PHASE 4: the kernel relations of the induction rules, the set-formation rules \\
\texttt{meta-nbg.scm} & K2, agent METACORE (2026-10-29): the kernel forms of the ten C+ rows of the K2 ledger \\
\texttt{meta-syntax.scm} & THE SYNTAX OF THE KERNEL AS A VNB OBJECT: locally nameless expressions, \\
\texttt{metric-completeness.scm} & Cauchy sequences, convergence, completeness \\
\texttt{metric-completion.scm} & \textit{(no header)} \\
\texttt{metric-connected.scm} & CONNECTEDNESS OF A SUBSET OF A METRIC SPACE \\
\texttt{metric-continuity.scm} & continuous maps between metric spaces \\
\texttt{metric-laws.scm} & the five metric laws, PROVEN from the definition \\
\texttt{metric-open-sets.scm} & \textit{(no header)} \\
\texttt{metric-outer-measure.scm} & the distance between sets, METRIC OUTER MEASURES and \\
\texttt{metric-space.scm} & METRIC-SPACE structure \\
\texttt{metric-subspace.scm} & the METRIC SUBSPACE of a metric space, and the \\
\texttt{metric-topology.scm} & balls, r-nets, total-boundedness \\
\texttt{minkowski-functional.scm} & \textit{(no header)} \\
\texttt{mod-seq.scm} & matrices acting on sequences of module elements \\
\texttt{module.scm} & MODULE over a ring, with the scalar action ACT \\
\texttt{monoid-power.scm} & MPOW: the natural-number power x\textasciicircum{}n in a monoid \\
\texttt{monoid.scm} & MONOID and COMM-MONOID structures \\
\texttt{monotone-class.scm} & MONOTONE CLASSES \\
\texttt{mvp-vocab.scm} & WAVE 2 (2026-10-27..30), wired 2026-10-30 from scratchpad/calint2/mvp-vocab.scm \\
\texttt{nets.scm} & \textit{(no header)} \\
\texttt{nn-arith.scm} & the recursion equations for + and * on NN \\
\texttt{nn-minus.scm} & truncated natural subtraction (monus), as VOCABULARY \\
\texttt{normed-ag-metric.scm} & the metric space underlying a normed abelian group \\
\texttt{normed-ag.scm} & NORMED-AG structure (normed abelian group) \\
\texttt{normed-field-metric.scm} & the metric space underlying a normed field \\
\texttt{normed-field.scm} & NORMED-FIELD: a commutative ring with multiplicative \\
\texttt{normed-vector-space.scm} & NORMED-VECTOR-SPACE: a real normed vector space \\
\texttt{nsl.scm} & THE NEEDHAM-SCHROEDER-LOWE PROTOCOL ON THE USER'S STRAND MODEL: names, public keys, \\
\texttt{nso.scm} & THE ORIGINAL NEEDHAM-SCHROEDER PROTOCOL (Thayer, Herzog, Guttman, section 5.3), ON THE \\
\texttt{nth-deriv.scm} & THE ITERATED DERIVATIVE ON AN OPEN SET, OVER A NORMED FIELD \\
\texttt{numeric-instances.scm} & \textit{(no header)} \\
\texttt{nvlocal-vocab.scm} & WAVE 2 (2026-10-27..30), wired 2026-10-30 from scratchpad/ccrqp/nvlocal-vocab.scm \\
\texttt{nvs-metric.scm} & the metric induced by the norm of a real normed vector \\
\texttt{one-sided-derivative.scm} & THE ONE-SIDED DERIVATIVES OF THE USER'S NOTES, \\
\texttt{operation-properties.scm} & named properties of operations \\
\texttt{or-missing-guarantee-instance.scm} & \textit{(no header)} \\
\texttt{order-lemmas.scm} & the finite order calculus of RR \\
\texttt{order-predicates.scm} & derived order predicates on the kernel <= \\
\texttt{order-zorn.scm} & abstract partial orders and ZORN'S LEMMA (statement only) \\
\texttt{ordinals.scm} & ordinal numbers: axioms, ordering, induction, recursion \\
\texttt{otway-rees.scm} & THE OTWAY-REES PROTOCOL ON THE USER'S STRAND MODEL (Thayer, Herzog, Guttman, \\
\texttt{outer-measure-construction.scm} & METHOD I OUTER MEASURES \\
\texttt{padic-vocab.scm} & THE p-ADIC ABSOLUTE VALUE ON QQ (CCR.pdf section 1.2) \\
\texttt{path-integral.scm} & THE VOCABULARY OF THE INTEGRAL ALONG A PATH \\
\texttt{poly-degree.scm} & DEG, LEADCOEF and MONOMIAL for the monoid algebra's \\
\texttt{poly-fn-vocab.scm} & POLYNOMIAL FUNCTIONS (CA-MONTEL ``POLY'', 2026-10-26) \\
\texttt{polynomial.scm} & the monoid-algebra construction A[M], and polynomials as \\
\texttt{positive-functional.scm} & \textit{(no header)} \\
\texttt{power-series.scm} & real power series  Sum\_\{n>=0\} coef(n) x\textasciicircum{}n  in RR \\
\texttt{prod-of-sums.scm} & \textit{(no header)} \\
\texttt{product-measure.scm} & PRODUCT SIGMA-ALGEBRAS, SECTIONS, DYNKIN SYSTEMS \\
\texttt{product-metric.scm} & the countable product of metric spaces, as a metric \\
\texttt{product2-vocab.scm} & THE BINARY PRODUCT OF TWO METRIC SPACES (definitions only) \\
\texttt{pruefer-vocab.scm} & WAVE 2 (2026-10-27..30), wired 2026-10-30 from scratchpad/ccrpr/pruefer-vocab.scm \\
\texttt{pseudometric.scm} & PSEUDOMETRIC-SPACE, the topology GENERATED by a countable \\
\texttt{push-forward.scm} & \textit{(no header)} \\
\texttt{qag-vocab.scm} & WAVE 2 (2026-10-27..30), wired 2026-10-30 from scratchpad/ccrpr/qag-vocab.scm \\
\texttt{qp-vocab.scm} & WAVE 2 (2026-10-27..30), wired 2026-10-30 from scratchpad/ccrqp/qp-vocab.scm \\
\texttt{qq-fractions.scm} & QQ is the field of fractions of the integers \\
\texttt{rabasics-vocab.scm} & vocabulary for the RA-BASICS wave (calculus.pdf 1.13, 3.7 \\
\texttt{rb2-vocab.scm} & vocabulary for calculus.pdf 2.27 and for path-connectedness \\
\texttt{real-integral.scm} & THE INTEGRAL OF A REAL-VALUED FUNCTION: positive and \\
\texttt{real-powers.scm} & the real power RPOW and the square root SQRT \\
\texttt{reduce.scm} & REDUCE: left-fold of a binary op over a finite indexed \\
\texttt{references.scm} & \textit{(no header)} \\
\texttt{regulated-primitive.scm} & COUNTABLE SETS, ONE-SIDED LIMITS WITHIN A SET, \\
\texttt{regulated.scm} & docs/calculus.pdf CHAPTER 4 SECTION 1: the vocabulary of \\
\texttt{riesz-markov-model.scm} & \textit{(no header)} \\
\texttt{ring-pow-vocab.scm} & wave of 2026-11-09/10 (agent CCR-1109-E), wired 2026-11-10 from scratchpad/ccr1109e/ring-pow-vocab.scm \\
\texttt{ring-power.scm} & RING-POWER: the natural-number power x\textasciicircum{}n in a commutative ring \\
\texttt{ring-simplify.scm} & polynomial-normal-form ring identity decision procedure \\
\texttt{ring.scm} & RING, RING-PROD, ZERO-RING \\
\texttt{ringoid.scm} & a ring with a distinguished two-sided ideal \\
\texttt{rn-vocab.scm} & E = R\textasciicircum{}n WITH THE SUP NORM, calculus.pdf ch. 2 s. 8 (32), (36)-(37), (44) \\
\texttt{road-juxtaposition.scm} & THE JUXTAPOSITION OF ROADS, THE TRIANGULAR ROAD AND THE LENGTH \\
\texttt{rr-ineq.scm} & PROTOTYPE of the rr-ineq applier (plan item 2, the "dumb \\
\texttt{scalar-inequalities.scm} & a library of important elementary inequalities \\
\texttt{segments-triangles.scm} & THE VOCABULARY OF CAUCHY-GOURSAT (complex-analysis.pdf \\
\texttt{semigroup.scm} & SEMIGROUP structure \\
\texttt{seminorm-extension.scm} & \textit{(no header)} \\
\texttt{seminorm-hahn-banach.scm} & seminorms on a K-vector space and the \\
\texttt{separable.scm} & separability of a metric space (a dense sequence), the \\
\texttt{seq-compact-product.scm} & \textit{(no header)} \\
\texttt{sequences.scm} & PROD-ORD, SUM, SUM-AG: products and sums over initial NN-segments \\
\texttt{set-basics.scm} & the NBG facts about SET that the kernel rules do not give \\
\texttt{set-vocabulary.scm} & two set constructors that the rest of the library is \\
\texttt{setoid.scm} & SETOID = (X, REL): a set with an equivalence relation, \\
\texttt{sigma-algebra.scm} & ALGEBRAS and SIGMA-ALGEBRAS of sets \\
\texttt{signed-measure.scm} & WAVE 2 (2026-10-27..30), wired 2026-10-30 from scratchpad/mf/mf-vocab.scm \\
\texttt{simple-integral.scm} & THE INTEGRAL OF A [0,+inf]-VALUED FUNCTION, DEFINED \\
\texttt{sip-complex-vocab.scm} & THE COMPLEX STRUCTURE J on a semi-inner-product space (SIP-1, \\
\texttt{sip-hilbert-vocab.scm} & REAL HILBERT SPACES: the norm and the metric of a \\
\texttt{sip-l2-vocab.scm} & EQUALITY ALMOST EVERYWHERE as a setoid (SIP-1, 2026-10-18) \\
\texttt{sip-vocab.scm} & SEMI-INNER-PRODUCT SPACES (SIP-1, 2026-10-18) \\
\texttt{sos-arith.scm} & exact-rational LP feasibility with witness extraction \\
\texttt{sos-oracle.scm} & the (sos) sum-of-squares oracle: close a nonstrict \\
\texttt{sqn.scm} & SQN(a): the sequences in a, i.e. FUN(NN, a) \\
\texttt{step-maps.scm} & WAVE 2 (2026-10-27..30), wired 2026-10-30 from scratchpad/mf/mf-mcvocab.scm \\
\texttt{strand-adversary.scm} & THE ADVERSARY: its traces, the infiltrated strand space \\
\texttt{strand-honesty.scm} & HONEST IDEALS (Thayer, Herzog, Guttman, \textasciitilde{}/docs/jcs\_strand\_spaces.pdf, Definition 6.10), \\
\texttt{strand-ideals.scm} & IDEALS OF THE MESSAGE ALGEBRA (Thayer, Herzog, Guttman, "Strand spaces: \\
\texttt{strand-messages.scm} & THE MESSAGES OF STRAND SPACES, ON THE USER'S MODEL, WITH ENCRYPTION \\
\texttt{strand-spaces.scm} & STRAND SPACES, NODES, BUNDLES ON THE USER'S MODEL \\
\texttt{subsequence-capture.scm} & \textit{(no header)} \\
\texttt{subtype-laws.scm} & the trivial subtype-subsumption laws, PROVEN \\
\texttt{summability.scm} & unconditional summability of a normed-AG-valued function \\
\texttt{tb11-vocab.scm} & vocabulary for topological-vector-spaces Appendix A.3 (TB-11): \\
\texttt{tg-vocab.scm} & topological groups (topological-vector-space.pdf Appendix A.3, TB-11) \\
\texttt{th3-vocab.scm} & wave of 2026-11-09/10 (agent TH3-A), wired 2026-11-10 from scratchpad/th3-1109/th3-vocab.scm \\
\texttt{th3a2-vocab.scm} & wave of 2026-11-09/10 (agent TH3-A2), wired 2026-11-10 from scratchpad/th3-1109/a2/th3a2-vocab.scm \\
\texttt{top-closed.scm} & \textit{(no header)} \\
\texttt{top-continuity.scm} & COMPACTNESS of a topological space, and the SUBSPACE \\
\texttt{top-product.scm} & \textit{(no header)} \\
\texttt{top-space.scm} & TOPOLOGICAL SPACES, and the functor Met -> Top \\
\texttt{transpose.scm} & the TRANSPOSE of a matrix (batch 32, 2026-09-25 \\
\texttt{trunc-metric.scm} & the TRUNCATED METRIC  d'(x,y) = min(1, d(x,y))  of a \\
\texttt{tvs-product-line.scm} & \textit{(no header)} \\
\texttt{tvs-vocab.scm} & \textit{(no header)} \\
\texttt{user-additions.scm} & structures saved via Build Structure \\
\texttt{views.scm} & view-as declarations between structures \\
\texttt{weak-star-separable.scm} & \textit{(no header)} \\
\texttt{weak-star.scm} & \textit{(no header)} \\
\texttt{winding-number.scm} & THE WINDING NUMBER AND THE CIRCLE AS A ROAD \\
\texttt{wot-compact-vocab.scm} & agent BA-1117-TRACE (round 2), 2026-11-17 \\
\texttt{zero-order-vocab.scm} & WAVE 2 (2026-10-27..30), wired 2026-10-30 from scratchpad/ch6/zero-order-vocab.scm \\
\texttt{zz-action.scm} & ZZ-ACT: the integer action on an abelian group \\
\texttt{zz-arith.scm} & ZZ is GENERATED by NN \\
\texttt{zz-divisibility.scm} & divisibility, Bezout combinations, gcd and coprimality \\
\texttt{zz-nt-vocab.scm} & \textit{(no header)} \\
\bottomrule
\end{longtable}

\section{\texttt{theorem-library/} and \texttt{calculus/}}
\label{sec:src-libraries}
\texttt{theorem-library/} holds 1112 proof scripts, each reaching \texttt{(qed \ldots)} and gated into the library load; by convention \texttt{X-proof.scm} proves \texttt{X}.  \texttt{calculus/} holds 21 probes and stress tests that are \emph{not} in \texttt{load.scm}.  Rather than enumerate them here, the theorems they prove are indexed, with statements and per-file links, in the generated browser references \texttt{reference/THEOREMS.md} (flat) and \texttt{reference/BY-TOPIC.md} (grouped like a textbook table of contents).

\section{Emacs interface (\texttt{emacs/}) and generated docs}
\label{sec:src-emacs}
The workspace front end lives in \texttt{emacs/}:
\begin{itemize}
\item \texttt{core-replay-stub.scm}
\item \texttt{vnb-cmd-gen.el}
\item \texttt{vnb-commands.lisp}
\item \texttt{vnb-complete.el}
\item \texttt{vnb-launch.el}
\item \texttt{vnb-panel-check.el}
\item \texttt{vnb.el}
\end{itemize}
The \texttt{reference/} directory is entirely \emph{generated} on prover load (\texttt{STRUCTURE-INDEX}, \texttt{TACTICS}, \texttt{OPERATORS}, \texttt{THEOREMS}, \texttt{BY-TOPIC}, \texttt{PROOF-DEBT}, \ldots) as \texttt{.md}, converted to the browser HTML by \texttt{build-reference-html} (reference-html.scm); it is never hand-edited.


%% =========================================================


%% =========================================================
%% app-design-spec.tex -- GENERATED from design-spec.md
%% (rules only; the Source lines stay in the .md for the integrator).
\chapter{Design specification}\label{app:design-spec}

The rules below are the design decisions the kernel and its self-verification follow.
Part~\ref{part:meta} cites them as DS-$n$.

\section*{Purpose and claim}
\begin{description}[style=nextline,leftmargin=1.4cm]
\item[DS-1]\phantomsection\label{ds:1} The self-verification work has two purposes: to make VNB credible as a proof checker, and to show that code can be modelled formally and its correctness proven relative to a stated base.
\item[DS-2]\phantomsection\label{ds:2} No statement of the form ``VNB cannot prove \textsc{falsity}'' is made. Every consistency statement is relative: if version-X is consistent (and the stated assumptions hold), the installed theorems do not include \textsc{falsity}. A derivation of \textsc{falsity}, once found, is conclusive; finding none is evidence only.
\item[DS-3]\phantomsection\label{ds:3} Version-X, the class theory every claim is relative to, is NBG with global choice. (Morse-Kelley is not used: the NBG claim is the stronger one, and the MK claim follows from it.)
\item[DS-4]\phantomsection\label{ds:4} A VNB class is read as a class of version-X.
\item[DS-5]\phantomsection\label{ds:5} The oracles are excluded from the formal reasoning and assumed correct.
\item[DS-6]\phantomsection\label{ds:6} The arguments that are not VNB theorems (the soundness of the core for version-X, the functoid argument, the faithfulness of the transcriptions) are written by hand and are labelled as such.
\item[DS-7]\phantomsection\label{ds:7} Functoids are covered by a hand-written argument, not by the formal encoding: a def-functoid head and an anonymous record are \emph{eliminated} (replaced by their instantiated bodies), and the meaning of a formula with functoids is the meaning of its elimination.
\end{description}
\section*{Route}
\begin{description}[style=nextline,leftmargin=1.4cm]
\item[DS-8]\phantomsection\label{ds:8} The formal route is K2: each non-oracle kernel rule is proven \emph{admissible} over a small core of rules, inside VNB, with formulas encoded as sets. Semantic soundness of the core (K1) is not attempted formally.
\item[DS-9]\phantomsection\label{ds:9} The faithfulness of each Scheme checker to its transcribed relation (K3) is deferred until K2 is complete. Its depth (a second independent transcription, or checkers rewritten as VNB-definable functions run through a small translator) is decided after a prototype.
\item[DS-10]\phantomsection\label{ds:10} The encoding is locally nameless (a bound variable is a de Bruijn index, a free variable a name); every node is a literal list whose first entry is a numeral tag; terms and formulas form one category separated by a well-formedness predicate; a context is a finite set of sentences.
\end{description}
\section*{The trusted base}
\begin{description}[style=nextline,leftmargin=1.4cm]
\item[DS-11]\phantomsection\label{ds:11} The kernel consists of all the procedures that call \texttt{dg-apply-rule!}. Only the files listed in \texttt{*cert-kernel-caller-files*} may call it; a fatal gate refuses any other call.
\item[DS-12]\phantomsection\label{ds:12} Every inference is checked before it is recorded, by a checker registered for its rule. A checker receives the inference after it is built and tests whether its premises, conclusion and side data stand in the relation of its rule; it builds nothing itself, and it is written independently of the procedure that builds the inference; it calls neither that procedure nor a copy. A missing checker or a refusal records nothing.
\item[DS-13]\phantomsection\label{ds:13} A new kernel operation comes with its checker and with one accepted and one refused suite control.
\item[DS-14]\phantomsection\label{ds:14} A kernel change, and a new core rule, enter only by an explicit design decision, one at a time, with the claim written beside the code. Proof repairs (statement guards) are preferred to kernel additions.
\item[DS-15]\phantomsection\label{ds:15} The primitive shelf grows only by an explicit foundational decision, one fact at a time, with the claim recorded in the file. A \texttt{primitive} or \texttt{definitional} stamp records what it assumes.
\item[DS-16]\phantomsection\label{ds:16} A kernel change lands only after a full exam (every proof run) of a tree that contains it. Several candidate trees may be examined in one night, one per worker.
\item[DS-17]\phantomsection\label{ds:17} The certificate key of a theorem follows the definitions of the statement's defined heads transitively, and records every definition and macete a checker read while accepting an inference of its proof.
\end{description}
\section*{Definedness}
\begin{description}[style=nextline,leftmargin=1.4cm]
\item[DS-18]\phantomsection\label{ds:18} The atoms IN, =, $\leq$ and < are strict: an atom with a non-denoting argument is false, so t = t states that t denotes. Quasi-equality s == t holds when neither side denotes or both denote the same object.
\item[DS-19]\phantomsection\label{ds:19} Every instantiation owes the definedness of the instance (the LUTINS rule): forall-elim, forsome-intro and the instantiation of a theorem macete post  G $\Rightarrow$ t = t  unless the certificate accepts t in G.
\item[DS-20]\phantomsection\label{ds:20} if-true and if-false hand over the strict equation IF(p, a, b) = a (= b) only when the branch is certified; otherwise the quasi-equation.
\item[DS-21]\phantomsection\label{ds:21} The strictness walks of the certificate stop at IF and at the binders, enter a SEP, BIG-UNION or VNB-LAMBDA through its domain only, read a def-functoid application through its instantiated body, and read an applied record through its contraction.
\item[DS-22]\phantomsection\label{ds:22} A true membership y in SEP x D. P certifies the strict terms of P at y.
\item[DS-23]\phantomsection\label{ds:23} The certificate accepts arithmetic over arguments that are themselves certified numbers, recursively. It has no clause for a structure term built from untyped variables; such definedness is landed in the proof.
\item[DS-24]\phantomsection\label{ds:24} A literal list denotes exactly when every component is a set.
\item[DS-25]\phantomsection\label{ds:25} nth-reduce, length-reduce and tuple-equality-decompose require the list's definedness (certified, or owed as L = L); NTH(k, t) is certified only when t is.
\item[DS-26]\phantomsection\label{ds:26} A one-entry list binder destructures a 1-tuple: [a] is not a. A LIST binder needs a literal CARTESIAN domain of the same length.
\item[DS-27]\phantomsection\label{ds:27} A term that mentions a record's bound variable is not certified under the record's binder.
\item[DS-28]\phantomsection\label{ds:28} The view clause certifies ACC(V(r)) only when the whole instantiated list V(r) is certified.
\item[DS-29]\phantomsection\label{ds:29} A functor-projection rewrite owes the list's definedness; the generated invariance equations are guarded by IS-src(r). The stamp stays \texttt{definitional}, its claim conditional.
\item[DS-30]\phantomsection\label{ds:30} The value certificate reads an unknown head as ``may be a class''; only justified rows certify a value as a set or nothing. A head is judged through its instantiated body, never by a shortcut on the structure its argument belongs to.
\item[DS-31]\phantomsection\label{ds:31} A guarded equation of a \emph{definition} (def-constant, def-by-nn-recursion, def-by-ord-recursion, def-operator) is read as a description: off every guard the head denotes nothing. A theorem of the same shape says nothing off its guards. def-by-ord-recursion writes its equations as quasi-equations. A guarded head is kept out of every table that would certify it total and out of every FUN typing (fatal gate).
\item[DS-32]\phantomsection\label{ds:32} The independent definedness certificate, when adopted, mirrors the builder's relation and does not widen it (its accessor-through-functoid reading stays narrow).
\end{description}
\section*{Functoids and operators}
\begin{description}[style=nextline,leftmargin=1.4cm]
\item[DS-33]\phantomsection\label{ds:33} Functoids stand outside quantification: a quantifier instance or witness is an object, never a def-functoid name, an anonymous record, or a bare operator or predicate name. A bare operator, functoid or predicate name is refused as a term.
\item[DS-34]\phantomsection\label{ds:34} A functoid unfolding is schematic: F(a) == body[a] holds at every argument, denoting or not, and owes nothing.
\item[DS-35]\phantomsection\label{ds:35} The hypotheses of the functoid argument (H1 distinct parameters, H2 no forward or recursive reference, H3 one definition per head, H4 no free variable in a body outside the parameters, H6 no head both a def-functoid and a library definition) are checked by a load-end gate, fatal in a strict load and reported under keep-going.
\item[DS-36]\phantomsection\label{ds:36} def-functor proves each view typing it would install; a view whose obligation fails is refused.
\item[DS-37]\phantomsection\label{ds:37} The install grader accepts the symbol its own defining formula defines.
\end{description}
\section*{Predicativity (staying within NBG)}
\begin{description}[style=nextline,leftmargin=1.4cm]
\item[DS-38]\phantomsection\label{ds:38} Comprehension is predicative: COMP, SEP, BIG-UNION and VNB-LAMBDA, and the three induction rules, accept a condition or body only when every quantifier in it is confined to sets, read through defined heads. The gates are fatal in a strict load.
\item[DS-39]\phantomsection\label{ds:39} The encoded core is NBG-strength: each schema that forms a class carries the premise NBG-OK on its condition or body.
\end{description}
\section*{The core}
\begin{description}[style=nextline,leftmargin=1.4cm]
\item[DS-40]\phantomsection\label{ds:40} The core consists of the natural-deduction rules with their side conditions and of one-line schemas a reader checks against the intended meaning of a constructor. In the core, forall-elim and forsome-intro always owe t = t, reflexivity fires only on a variable or an atomic constant, and the list schemas carry the list's definedness.
\item[DS-41]\phantomsection\label{ds:41} The core's eq-subst is two-way (both orientations of the context equation).
\item[DS-42]\phantomsection\label{ds:42} The core admits, each with its claim: strictness of IN on both sides, total constructors on denoting arguments, binder terms over a denoting domain, comprehension, lists of sets; strictness of the constructors in their arguments and of the binder terms in their domain; strictness of a predicate application, with predicate codes, sound only because every instantiation owes definedness.
\item[DS-43]\phantomsection\label{ds:43} The core admits the congruence rule for descriptions (IOTA-CR) and for the term binders SEP, COMP, BIG-UNION, VNB-LAMBDA (TBIND-CR); both conclude a quasi-equation only.
\item[DS-44]\phantomsection\label{ds:44} The core admits a functoid table FUN-DEF (rows filled by the K3 translator; anonymous records lambda-lifted into rows) and functoid-beta-eq, a quasi-equation at every argument list.
\item[DS-45]\phantomsection\label{ds:45} The core admits tuple eta: a member of an n-ary product is the list of its own components.
\item[DS-46]\phantomsection\label{ds:46} The core keeps truth-intro although the kernel operation is retired; contraposition is retired from the kernel and from the encoding.
\item[DS-47]\phantomsection\label{ds:47} A META relation that bounds a witness by SET only on its intended domain says so in the hand-written text: the argument list of J-MACN, and the hypothesis list of HYPS-IN.
\item[DS-48]\phantomsection\label{ds:48} A relation transcribing a kernel rewrite requires its two formulas to be sentences (J-CART, as J-TEQ, J-UDEC, J-IDEC).
\end{description}
\section*{The odd primitives}
\begin{description}[style=nextline,leftmargin=1.4cm]
\item[DS-49]\phantomsection\label{ds:49} No natural number is a list (\texttt{nn-not-a-list}, primitive). It is included because the strand-space token law needs a set of tokens free of lists; its consistency is shown by an explicit interpretation.
\item[DS-50]\phantomsection\label{ds:50} tuple-equality-decompose is a primitive with its claim recorded: two literal lists that both denote are equal iff they have the same length and equal components.
\end{description}
\section*{Documentation}
\begin{description}[style=nextline,leftmargin=1.4cm]
\item[DS-51]\phantomsection\label{ds:51} The manual states present-tense rules and cites this specification; dates, decision history and the findings that led to a rule stay in the META report.
\end{description}
\section*{Partial terms}
\begin{description}[style=nextline,leftmargin=1.4cm]
\item[DS-52]\phantomsection\label{ds:52} $\mathrm{CHOICE}(A)$ denotes iff $A$ is inhabited; then it is a member of $A$. No rule infers anything from $\mathrm{CHOICE}(A) = \mathrm{CHOICE}(A)$ beyond that.
\item[DS-53]\phantomsection\label{ds:53} A rule that makes a membership (or a $\mathrm{FUN}$ statement) true owes the definedness of every class argument it does not itself check, unless the certificate already accepts it: \texttt{tuples-intro} at the empty tuple, every entry of \texttt{union-intro}, the codomain of \texttt{lambda-type}. The matching core rules carry the same premise.
\end{description}


%% =========================================================
\chapter{Regression and demonstration test suite}
\label{app:tests}

The file \texttt{test-suite.scm} is the regression suite of the VNB
proof checker: some two and a half thousand checks in sixty-odd
sections.  This appendix reproduces the first of them (sections 1 to 7
of the file), which serve also as worked examples: ground arithmetic,
axiom installation, the parser, wff validation, macetes with local
context, proof commands, the number-system chain, and arithmetic
over~$\mathbb{Z}$, $\mathbb{Q}$, $\mathbb{R}$, $\mathbb{C}$.  The
remainder covers the kernel checkers, the gates, the tactics and the
driver kit.  Running it requires only the command below; each test
prints \texttt{PASS} or \texttt{FAIL} and the final line gives the
aggregate count.

\begin{VNBcode}
;;; Run the suite (it loads the system itself):
;;;   mit-scheme --heap 200000 --quiet --load test-suite.scm < /dev/null
\end{VNBcode}

\section{Test harness}

\begin{sloppypar}The harness uses three predicates (\texttt{check}, \texttt{check-true},
\texttt{check-false}), one error-expectation helper
(\texttt{check-error}), and one proof-completion helper
(\texttt{check-proof}).  Two convenience wrappers, \texttt{wff-str} and
\texttt{macete-str}, keep individual test lines short.\end{sloppypar}

\begin{VNBcode}
(define *pass-count* 0)
(define *fail-count* 0)

(define (check label thunk expected)
  (let ((got (thunk)))
    (if (equal? got expected)
        (begin (set! *pass-count* (+ *pass-count* 1))
               (display "  PASS  ") (display label) (newline))
        (begin (set! *fail-count* (+ *fail-count* 1))
               (display "  FAIL  ") (display label)
               (display " -- expected ") (write expected)
               (display " got ") (write got) (newline)))))

(define (check-true  label thunk) (check label thunk #t))
(define (check-false label thunk) (check label thunk #f))

(define (check-error label thunk)
  ;; Pass when thunk returns a <vnb-error> (i.e. input was rejected).
  (let ((result (thunk)))
    (if (vnb-error? result)
        (begin (set! *pass-count* (+ *pass-count* 1))
               (display "  PASS  ") (display label) (newline))
        (begin (set! *fail-count* (+ *fail-count* 1))
               (display "  FAIL  ") (display label)
               (display " -- expected error, got ") (write result) (newline)))))

(define (check-proof label thunk)
  ;; Pass when thunk completes without raising.
  (call-with-current-continuation
    (lambda (k)
      (with-exception-handler
        (lambda (e)
          (set! *fail-count* (+ *fail-count* 1))
          (display "  FAIL  ") (display label)
          (display " -- exception: ")
          (display (condition/report-string e)) (newline)
          (k #f))
        thunk)
      (set! *pass-count* (+ *pass-count* 1))
      (display "  PASS  ") (display label) (newline))))

;;; wff-str: extract display string from a wff (strips surrounding quotes).
(define (wff-str w)
  (let ((s (wff->string w)))
    (substring s 1 (- (string-length s) 1))))

;;; macete-str: apply a named macete to a formula given as a string.
;;; Returns the result display string, or #f if the macete doesn't fire.
(define (macete-str name str)
  (let ((r (apply-macete name (make-wff-from-string str))))
    (and r (wff-str r))))

;;; roundtrip: parse a formula string and print it back.
(define (roundtrip str)
  (let ((w (make-wff-from-string str)))
    (if (vnb-error? w) #f (wff-str w))))
\end{VNBcode}

\section{Ground arithmetic evaluation}

The built-in evaluator decides membership and equalities over exact
numbers without invoking the proof kernel.

\begin{VNBcode}
(check-true  "0 in NN"        (lambda () (arith-membership-check 0   'NN)))
(check-true  "100 in NN"      (lambda () (arith-membership-check 100 'NN)))
(check-false "-1 not in NN"   (lambda () (arith-membership-check -1  'NN)))
(check-false "1/2 not in NN"  (lambda () (arith-membership-check 1/2 'NN)))

(check-true  "0 in ZZ"        (lambda () (arith-membership-check 0   'ZZ)))
(check-true  "-3 in ZZ"       (lambda () (arith-membership-check -3  'ZZ)))
(check-false "1/2 not in ZZ"  (lambda () (arith-membership-check 1/2 'ZZ)))

(check-true  "1/3 in QQ"      (lambda () (arith-membership-check 1/3 'QQ)))
(check-false "1/3 not in ZZ"  (lambda () (arith-membership-check 1/3 'ZZ)))

(check "2 + 3 = 5"       (lambda () (arith-eval-term '(+ 2 3)))     5)
(check "3 * 4 = 12"      (lambda () (arith-eval-term '(* 3 4)))     12)
(check "10 + (-3) = 7"   (lambda () (arith-eval-term '(+ 10 (- 3)))) 7)
(check "succ(0) = 1"     (lambda () (arith-eval-term '(succ 0)))    1)
(check "power(2,3) = 8"  (lambda () (arith-eval-term '(power 2 3))) 8)

(check-true  "= 2+3 5"   (lambda () (arith-eval-formula '(= (+ 2 3) 5))))
(check-false "= 2+3 6"   (lambda () (arith-eval-formula '(= (+ 2 3) 6))))
(check-true  "<= 0 5"    (lambda () (arith-eval-formula '(<= 0 5))))
(check-false "<= 5 0"    (lambda () (arith-eval-formula '(<= 5 0))))
\end{VNBcode}

\section{Axioms installed}

Every axiom loaded by \texttt{load.scm} is simultaneously a theorem
(available for \texttt{ta} and backchaining) and a macete entry.
The \texttt{axiom-present?} helper checks the theorem table.

\begin{VNBcode}
(define (axiom-present? name)
  (if (hash-table-ref/default *theorem-table* name #f) #t #f))

(for-each (lambda (name)
            (check-true (symbol->string name)
                        (lambda () (axiom-present? name))))
  '(membership-implies-sethood extensionality
    empty-set-is-set empty-set-has-no-members
    pairing pairing-membership
    power-set power-set-membership
    union-set-closure intersection-set-closure
    fun-set-iff fun-elements-are-sets
    fun-domain-apply-def fun-codomain-iff
    cartesian-set-iff tuples-sethood
    make-set-membership make-set-sethood make-set-empty
    length-of-empty length-in-nn nth-in-range
    nn-is-set nn-zero-in nn-succ-closed
    nn-add-closed nn-mul-closed nn-add-comm nn-mul-comm
    nn-add-assoc nn-mul-assoc nn-distributive nn-one-mul nn-add-zero
    power-typing-nonneg power-zero power-succ power-neg
    zz-is-set zz-zero-in zz-one-in
    qq-is-set rr-is-set cc-is-set
    nn-subset-zz zz-subset-qq qq-subset-rr rr-subset-cc
    burali-forti nn-subset-ord ord-succ-in transfinite-induction))

;;; ZZplus must be completely absent (superseded by ZZ axioms)
(check-false "zzplus-is-set absent"
             (lambda () (axiom-present? 'zzplus-is-set)))
\end{VNBcode}

\section{Parser round-trips}

The parser produces an internal S-expression; \texttt{wff->string} prints
it back.  Round-tripping through \texttt{make-wff-from-string} and
\texttt{wff-str} is the primary correctness check for the surface syntax.

\begin{VNBcode}
(check "n in nn"         (lambda () (roundtrip "n in nn"))         "n in nn")
(check "n = m"           (lambda () (roundtrip "n = m"))           "n = m")
(check "truth"           (lambda () (roundtrip "truth"))           "truth")
(check "falsity"         (lambda () (roundtrip "falsity"))         "falsity")
(check "n + m = m + n"   (lambda () (roundtrip "n + m = m + n"))   "n + m = m + n")

(check "A and B"         (lambda () (roundtrip "n in nn and m in nn"))
                         "n in nn and m in nn")
(check "A implies B"     (lambda () (roundtrip "n in nn implies n in zz"))
                         "n in nn implies n in zz")
(check "not A"           (lambda () (roundtrip "not(n in nn)"))    "not(n in nn)")

(check-true "forall bare"
  (lambda () (wff? (make-wff-from-string "forall([x], x in nn implies x in zz)"))))
(check-true "forall restricted"
  (lambda () (wff? (make-wff-from-string "forall([x in nn], x in zz)"))))
(check-true "forsome"
  (lambda () (wff? (make-wff-from-string "forsome([x in nn], x = 0)"))))
(check-true "sep {x in A | p}"
  (lambda () (wff? (make-wff-from-string "z in {x in a | x in nn}"))))
\end{VNBcode}

\section{Wff validation}

\begin{sloppypar}\texttt{make-wff} and \texttt{make-wff-from-string} validate the formula
and return a \texttt{<vnb-error>} (rather than raising) on
ill-formed input.\end{sloppypar}

\begin{VNBcode}
(check-true  "IN accepted"      (lambda () (wff? (make-wff-from-string
                                                   "n in nn"))))
(check-true  "IMPLIES accepted" (lambda () (wff? (make-wff-from-string
                                                   "n in nn implies n in zz"))))
(check-true  "NOT accepted"     (lambda () (wff? (make-wff-from-string
                                                   "not(n in nn)"))))
(check-true  "FORALL accepted"  (lambda () (wff? (make-wff-from-string
                                                   "forall([x in nn], x in zz)"))))

;;; Invalid: bare number in wff position
(check-error "bare number rejected"
             (lambda () (make-wff-from-string "1")))

;;; Invalid: term-forming head (UNION) in wff position
(check-error "union in wff position rejected"
             (lambda () (make-wff '(UNION x y))))
\end{VNBcode}

\section[\texttt{apply-macete} and local context (Monk 1988)]{\texttt{apply-macete} and local context (Monk~\cite{Monk1988})}

Conditions of a conditional macete are checked against the local
context accumulated as the rewriter descends.  Entering the
consequent of \texttt{implies} adds the antecedent; entering the
right conjunct of \texttt{and} adds the left conjunct; entering the
right disjunct of \texttt{or} adds the negation of the left disjunct.

\subsection*{Unconditional macetes}

\begin{VNBcode}
(install-theorem! 'eq-symm
  (make-wff-from-string "forall([x, y], x = y iff y = x)"))

(check "eq-symm: 3+5=8 -> 8=3+5"
       (lambda () (macete-str 'eq-symm "3 + 5 = 8"))
       "8 = 3 + 5")
(check "eq-symm: rewrites inside implies"
       (lambda () (macete-str 'eq-symm "n = m implies truth"))
       "m = n implies truth")
(check "eq-symm: no match on n in nn"
       (lambda () (macete-str 'eq-symm "n in nn"))
       #f)
(check "eq-symm: both sides rewritten"
       (lambda () (macete-str 'eq-symm "n = m and m = k"))
       "m = n and k = m")

;;; A propositional schema cannot be stated through make-wff, which
;;; refuses a quantifier over a formula position; it is installed as a
;;; macete from its S-expression.
(install-theorem! 'and-comm-iff
  '(FORALL p (FORALL q (IFF (AND p q) (AND q p)))))

(check "and-comm: (A and B) -> (B and A)"
       (lambda () (macete-str 'and-comm-iff "n in nn and m in nn"))
       "m in nn and n in nn")
(check "and-comm: fires inside implies antecedent"
       (lambda () (macete-str 'and-comm-iff "n in nn and m in nn implies truth"))
       "m in nn and n in nn implies truth")
\end{VNBcode}

\subsection*{Conditional macetes}

\begin{VNBcode}
;;; nn-add-comm: forall([a,b in NN], a+b = b+a)
;;; Without context the rewrite is suppressed.
(check "nn-add-comm: no fire on bare sum"
       (lambda () (macete-str 'nn-add-comm "n + m = 0"))   #f)

;;; Antecedent (AND (IN n NN) (IN m NN)) is flattened into local
;;; context when entering the consequent.
(check "nn-add-comm: fires inside implies consequent"
       (lambda () (macete-str 'nn-add-comm
                              "n in nn and m in nn implies n + m = m + n"))
       "n in nn and m in nn implies m + n = n + m")

;;; nn-add-zero: forall([n in NN], n + 0 = n)
(check "nn-add-zero: no fire without context"
       (lambda () (macete-str 'nn-add-zero "n + 0 = 5"))   #f)
(check "nn-add-zero: fires inside implies consequent"
       (lambda () (macete-str 'nn-add-zero "n in nn implies n + 0 = 5"))
       "n in nn implies n = 5")
;;; Entering right conjunct of AND adds left conjunct to context.
(check "nn-add-zero: fires in right conjunct of AND"
       (lambda () (macete-str 'nn-add-zero "k in nn and k + 0 = 5"))
       "k in nn and k = 5")
;;; Wrong order: condition is right conjunct — no context available yet.
(check "nn-add-zero: no fire when condition is right conjunct"
       (lambda () (macete-str 'nn-add-zero "k + 0 = 5 and k in nn"))  #f)

;;; nn-add-closed: rewrites n+m in NN to truth when conditions hold.
(check "nn-add-closed: fires inside implies, rewrites to truth"
       (lambda () (macete-str 'nn-add-closed
                              "n in nn and m in nn implies n + m in nn"))
       "n in nn and m in nn implies truth")
(check "nn-add-closed: no fire without context"
       (lambda () (macete-str 'nn-add-closed "n + m in nn"))  #f)

;;; power-succ: forall([x in CC, n in NN], power(x,succ(n)) = x*power(x,n))
(check "power-succ: fires with two conditions in context"
       (lambda () (macete-str 'power-succ
                              "z in cc and k in nn implies power(z, succ(k)) = w"))
       "z in cc and k in nn implies z * z ^ k = w")
\end{VNBcode}

\section{Proof commands}

Each sub-test proves a simple theorem using one or two rules.
\texttt{check-proof} verifies that the proof completes without exception.

\begin{VNBcode}
;;; Arithmetic (arith) — ground facts
(check-proof "arith: 0 in NN"
  (lambda () (sp (make-wff-from-string "0 in nn")) (arith)))
(check-proof "arith: 2 + 3 = 5"
  (lambda () (sp (make-wff-from-string "2 + 3 = 5")) (arith)))
(check-proof "arith: -5 in ZZ"
  (lambda () (sp (make-wff-from-string "-5 in zz")) (arith)))
(check-proof "arith: 1/3 in QQ"
  (lambda () (sp (make-wff-from-string "1/3 in qq")) (arith)))

;;; Reflexivity (rfl) and quasi-reflexivity (qrfl)
(check-proof "rfl: n = n"
  (lambda () (sp (make-wff-from-string "n = n")) (rfl)))
(check-proof "qrfl: n == n"
  (lambda () (sp (make-wff-from-string "n == n")) (qrfl)))

;;; Direct inference (di) — introduce implication, split conjunction/AND
(check-proof "di + ass: A implies A"
  (lambda ()
    (sp (make-wff-from-string "n in nn implies n in nn"))
    (di) (ass)))

;;; Antecedent inference (ai) — decompose an assumption
(check-proof "di + ai: prove A and B from A, B"
  (lambda ()
    (sp (make-wff-from-string "n in nn and m in nn implies m in nn and n in nn"))
    (di)
    (ai (make-wff-from-string "n in nn and m in nn"))
    (di) (ass) (ass)))

;;; OR introduction (oi-l, oi-r)
(check-proof "oi-l: A implies (A or B)"
  (lambda ()
    (sp (make-wff-from-string "n in nn implies n in nn or m in nn"))
    (di) (oi-l) (ass)))
(check-proof "oi-r: A implies (B or A)"
  (lambda ()
    (sp (make-wff-from-string "n in nn implies m in nn or n in nn"))
    (di) (oi-r) (ass)))

;;; Theorem assumption (ta) + universal instantiation (ui)
(check-proof "ta + ui: n+0 = n"
  (lambda ()
    (sp (make-wff-from-string "n in nn implies n + 0 = n"))
    (di) (ta 'nn-add-zero) (ui "n") (ass)))
(check-proof "ta + ui: n+m = m+n"
  (lambda ()
    (sp (make-wff-from-string "n in nn and m in nn implies n + m = m + n"))
    (di) (ta 'nn-add-comm) (ui "n") (ui "m") (ass)))
(check-proof "ta + ui: n in NN => n in ZZ"
  (lambda ()
    (sp (make-wff-from-string "n in nn implies n in zz"))
    (di) (ta 'nn-subset-zz) (ui "n") (ass)))
(check-proof "ta + ui: power(x,succ n)"
  (lambda ()
    (sp (make-wff-from-string
          "x in cc and n in nn implies power(x, succ(n)) = x * power(x, n)"))
    (di) (ta 'power-succ) (ui "x") (ui "n") (ass)))

;;; Backchain (bc) -- bc takes an implication in context, not a theorem
;;; name: the theorem is brought in with ta and instantiated first.
(check-proof "bc: n in NN and m in NN"
  (lambda ()
    (sp (make-wff-from-string "n in nn and m in nn implies n + m in nn"))
    (di) (ai (make-wff-from-string "n in nn and m in nn"))
    (ta 'nn-add-closed)
    (inst '(FORALL n (FORALL m (IMPLIES (AND (IN n NN) (IN m NN)) (IN (+ n m) NN)))) 'n)
    (inst '(FORALL m (IMPLIES (AND (IN n NN) (IN m NN)) (IN (+ n m) NN))) 'm)
    (bc '(IMPLIES (AND (IN n NN) (IN m NN)) (IN (+ n m) NN)))
    (ass) (ass) (ass)
    (unless (proof-done? *ps*) (error "bc test: proof did not close"))))

;;; Proof by contradiction (pbc)
(check-proof "pbc: double negation"
  (lambda ()
    (sp (make-wff-from-string "not(not(n in nn)) implies n in nn"))
    (di) (pbc) (ass)))

;;; Macete inside a proof (mac)
(check-proof "mac: eq-symm rewrites goal"
  (lambda ()
    (sp (make-wff-from-string "n = m implies m = n"))
    (di) (mac 'eq-symm) (ass)))

;;; Exists-witness (ew)
(check-proof "ew: exists x in NN. x = 0"
  (lambda ()
    (sp (make-wff-from-string "forsome([x in nn], x = 0)"))
    (ew "0") (arith) (arith)))

;;; Cut — introduce an intermediate lemma
(check-proof "cut: introduce intermediate fact"
  (lambda ()
    (sp (make-wff-from-string "n in nn implies n + 0 = n"))
    (di)
    (cut (make-wff-from-string "n in nn implies n + 0 = n"))
    (di) (ta 'nn-add-zero) (ui "n") (ass)
    (di) (ass)))

;;; Functoid (lambdoid) beta reduction (beta)
(check-proof "beta: lambdoid([x in nn], x+1)(3) = 3+1"
  (lambda ()
    (sp (make-wff-from-string "lambdoid([x in nn], x + 1)(3) = 3 + 1"))
    (beta) (rfl)))
\end{VNBcode}

\section{Number-system inclusion chain}

\begin{VNBcode}
(for-each
  (lambda (triple)
    (let ((val (car triple)) (sys (cadr triple)) (expected (caddr triple)))
      (check (string-append (number->string val) " in " (symbol->string sys))
             (lambda () (arith-membership-check val sys))
             expected)))
  '((0    NN #t)  (1    NN #t)  (-1   NN #f)  (1/2  NN #f)
    (0    ZZ #t)  (-5   ZZ #t)  (1/2  ZZ #f)
    (1/3  QQ #t)  (-7   QQ #t)  (1/3  ZZ #f)
    (0    RR #t)  (-7   RR #t)
    (0    CC #t)  (-5   CC #t)))
\end{VNBcode}

\section{Arithmetic in \texorpdfstring{$\mathbb{Z}$, $\mathbb{Q}$, $\mathbb{R}$, $\mathbb{C}$}{Z, Q, R, C}}

\begin{VNBcode}
(check-proof "arith: 0 in ZZ"
  (lambda () (sp (make-wff-from-string "0 in zz")) (arith)))
(check-proof "arith: -3 in ZZ"
  (lambda () (sp (make-wff-from-string "-3 in zz")) (arith)))
(check-proof "arith: 1/3 in QQ"
  (lambda () (sp (make-wff-from-string "1/3 in qq")) (arith)))
(check-proof "arith: 1/3 not in ZZ"
  (lambda () (sp (make-wff-from-string "not(1/3 in zz)")) (arith)))

(check "1/3 in qq" (lambda () (arith-eval-formula '(IN 1/3 QQ)))  #t)
(check "1/3 in zz" (lambda () (arith-eval-formula '(IN 1/3 ZZ)))  #f)
(check "-3 in zz"  (lambda () (arith-eval-formula '(IN -3 ZZ)))   #t)
(check "-3 in qq"  (lambda () (arith-eval-formula '(IN -3 QQ)))   #t)
\end{VNBcode}

\section{Summary line}

\begin{VNBcode}
(newline)
(display "=== SUMMARY: ")
(display *pass-count*) (display " passed, ")
(display *fail-count*) (display " failed ===\n")
\end{VNBcode}

\noindent A clean run ends with a line of the form \verb|=== SUMMARY: N passed, 0 failed ===|.

\section{Other checks}

The workspace has its own check, \texttt{emacs --batch -l
emacs/vnb-panel-check.el}, ending in \texttt{=== PANEL CHECK: N passed, M
failed ===}.  \texttt{test-suite-negative.scm} is a corpus of statements
that must \emph{not} be provable.  The suite can start from a saved band
instead of loading the library.  In every case it is the summary line,
not the exit status, that shows the run completed.

%% =========================================================


%% =========================================================
% app-glossary-generated.tex -- GENERATED by docs/gen-glossary.py from
% reference/GLOSSARY.md.  Do not edit; re-run `make gen' in docs/.

\chapter{Glossary: every name, A to Z}
\label{app:glossary}

Every other listing in this manual and in the browser reference is grouped by
something: operators by class, structures by structure, tactics by task.  This
one has the ordering nobody has to learn.  It carries the whole
\emph{vocabulary} -- 2241 names: 35 structure, 15 refinement, 20 instance, 33 view, 843 predicate, 1005 functoid, 54 accessor, 25 defined-fn, 57 operator, 10 primitive, 94 tactic -- and no results: the theorems, axioms and
supports are catalogued in \texttt{THEOREMS.md}, \texttt{BY-OPERATOR.md} and
\texttt{PROOF-DEBT.md}, and each entry here says how many of them mention it.

At the REPL, \texttt{(glossary 'NAME)} prints one entry and
\texttt{(glossary "pat")} every name containing \texttt{pat}.

\begingroup
\small
\setlength{\parindent}{0pt}

\section*{Symbols}
\addcontentsline{toc}{section}{Symbols}

\noindent\textbf{\texttt{*}}\quad\textit{operator}\index{*@\texttt{*}}\\
a term-forming head \\
mentioned by 3226 result(s)

\smallskip

\noindent\textbf{\texttt{+}}\quad\textit{operator}\index{+@\texttt{+}}\\
a term-forming head \\
mentioned by 2546 result(s)

\smallskip

\noindent\textbf{\texttt{-}}\quad\textit{operator}\index{-@\texttt{-}}\\
a term-forming head \\
characterized by \texttt{is-pole}  \textit{(definitional)}:
\begin{quote}\ttfamily\footnotesize\raggedright\noindent
forall(\allowbreak{}[mrf\_, mra\_, mrn\_], is-pole(\allowbreak{}mrf\_, mra\_, mrn\_) iff mra\_ in cc and mrn\_ in nn and 1 <= mrn\_ and forsome(\allowbreak{}[mrr\_ in rr], 0 < mrr\_ and forsome(\allowbreak{}[mrg\_], holomorphic-on(\allowbreak{}ball(\allowbreak{}cc-ms, mra\_, mrr\_), mrg\_) and not(\allowbreak{}mrg\_(\allowbreak{}mra\_) = 0) and forall(\allowbreak{}[mrz\_ in difference(\allowbreak{}ball(\allowbreak{}cc-ms, mra\_, mrr\_), singleton(\allowbreak{}mra\_))], mrf\_(\allowbreak{}mrz\_) = recip(\allowbreak{}mrz\_ - mra\_) \textasciicircum{} mrn\_ * mrg\_(\allowbreak{}mrz\_)))))
\end{quote}
also: \texttt{opg-rep} \texttt{lfts-def} \texttt{zz-prime} \texttt{cc-ms-def} \texttt{cchs-succ} \texttt{rr-ms-def} \texttt{cayley-def} \texttt{cc-ms@dist} ... (1503 in all) \\
mentioned by 2368 result(s)

\smallskip

\noindent\textbf{\texttt{/}}\quad\textit{operator}\index{/@\texttt{/}}\\
a term-forming head \\
mentioned by 73 result(s)

\smallskip

\noindent\textbf{\texttt{<(x, y)}}\quad\textit{predicate}\index{<(x, y)@\texttt{<(x, y)}}\\
reads: \$1 is less than \$2 \\
defined by \texttt{<} \\
characterized by \texttt{<}  \textit{(definitional)}:
\begin{quote}\ttfamily\footnotesize\raggedright\noindent
forall(\allowbreak{}[x, y], x < y iff x <= y and not(\allowbreak{}x = y))
\end{quote}
mentioned by 2094 result(s)

\smallskip

\noindent\textbf{\texttt{<=}}\quad\textit{primitive}\index{<=@\texttt{<=}}\\
reads: \$1 is at most \$2 \\
a kernel relation \\
mentioned by 3760 result(s)

\smallskip

\noindent\textbf{\texttt{=}}\quad\textit{primitive}\index{=@\texttt{=}}\\
reads: \$1 equals \$2 \\
a kernel relation

\smallskip

\noindent\textbf{\texttt{==}}\quad\textit{primitive}\index{==@\texttt{==}}\\
reads: \$1 is identical to \$2 \\
a kernel relation

\smallskip

\noindent\textbf{\texttt{>}}\quad\textit{primitive}\index{>@\texttt{>}}\\
reads: \$1 is greater than \$2 \\
a kernel relation

\smallskip

\noindent\textbf{\texttt{>=}}\quad\textit{primitive}\index{>=@\texttt{>=}}\\
reads: \$1 is at least \$2 \\
a kernel relation

\smallskip


\section*{A}
\addcontentsline{toc}{section}{A}

\noindent\textbf{\texttt{abelian-group}}\quad\textit{structure}\index{abelian-group@\texttt{abelian-group}}\\
predicate \texttt{is-abelian-group}; slots: carr opr iden inv

\smallskip

\noindent\textbf{\texttt{abelian-group-as-monoid(s)}}\quad\textit{view}\index{abelian-group-as-monoid(s)@\texttt{abelian-group-as-monoid(s)}}\\
an abelian-group viewed as a monoid \\
characterized by \texttt{mpow-add-abelian-group-as-monoid}  \textit{(definitional)}:
\begin{quote}\ttfamily\footnotesize\raggedright\noindent
forall(\allowbreak{}[r], is-abelian-group(\allowbreak{}r) implies forall(\allowbreak{}[x in carr(\allowbreak{}r), j in nn, k in nn], mpow(\allowbreak{}abelian-group-as-monoid(\allowbreak{}r), x, j + k) = (\allowbreak{}opr(\allowbreak{}r))(\allowbreak{}mpow(\allowbreak{}abelian-group-as-monoid(\allowbreak{}r), x, j), mpow(\allowbreak{}abelian-group-as-monoid(\allowbreak{}r), x, k))))
\end{quote}
also: \texttt{mpow-one-abelian-group-as-monoid} \texttt{abelian-group-as-monoid-is-monoid} \texttt{abelian-group-as-monoid-is-monoid-module-vector-ag} \texttt{abelian-group-as-monoid-is-monoid-ring-additive-ag} \texttt{abelian-group-as-monoid-is-monoid-field-additive-ag} \texttt{abelian-group-as-monoid-is-monoid-normed-field-additive-ag} \texttt{abelian-group-as-monoid-is-monoid-normed-ag-as-abelian-group} \texttt{abelian-group-as-monoid-is-monoid-commutative-ring-additive-ag} ... (28 in all) \\
mentioned by 46 result(s)

\smallskip

\noindent\textbf{\texttt{abs}}\quad\textit{operator}\index{abs@\texttt{abs}}\\
a term-forming head \\
characterized by \texttt{rr-abs-def}  \textit{(definitional)}:
\begin{quote}\ttfamily\footnotesize\raggedright\noindent
forall(\allowbreak{}[x in rr], (\allowbreak{}0 <= x implies abs(\allowbreak{}x) = x) and (\allowbreak{}not(\allowbreak{}0 <= x) implies abs(\allowbreak{}x) = -x))
\end{quote}
also: \texttt{rn-maxabs-succ} \texttt{ps-absolutely-converges-at} \texttt{rr-le-abs} \texttt{rr-abs-neg} \texttt{bsw-ucl-abs} \texttt{ps-abs-term} \texttt{rr-abs-mult} \texttt{rr-abs-zero} ... (78 in all) \\
mentioned by 559 result(s)

\smallskip

\noindent\textbf{\texttt{abs-fn(omega, f)}}\quad\textit{functoid}\index{abs-fn(omega, f)@\texttt{abs-fn(omega, f)}}\\
reads: the absolute value of \$2 on \$1 \\
definition:
\begin{quote}\ttfamily\footnotesize\raggedright\noindent
vnb-lambda(\allowbreak{}afx\_, omega, abs(\allowbreak{}f(\allowbreak{}afx\_)))
\end{quote}
characterized by \texttt{abs-fn-fun}  \textit{(proven)}:
\begin{quote}\ttfamily\footnotesize\raggedright\noindent
forall(\allowbreak{}[omega, f in fun(\allowbreak{}omega, rr)], abs-fn(\allowbreak{}omega, f) in fun(\allowbreak{}omega, rr-pos-star))
\end{quote}
also: \texttt{abs-fn-value} \texttt{abs-fn-unfold} \texttt{lp-norm-abs-fn} \texttt{abs-fn-eq-parts} \texttt{integral-abs-fn} \texttt{abs-fn-real-measurable} \\
mentioned by 27 result(s)

\smallskip

\noindent\textbf{\texttt{abs-op(kqv\_, kqt\_)}}\quad\textit{functoid}\index{abs-op(kqv_, kqt_)@\texttt{abs-op(kqv\_, kqt\_)}}\\
reads: the absolute value of \$2 \\
definition:
\begin{quote}\ttfamily\footnotesize\raggedright\noindent
op-sqrt(\allowbreak{}kqv\_, bop-comp(\allowbreak{}kqv\_, op-adjoint(\allowbreak{}kqv\_, kqv\_, kqt\_), kqt\_))
\end{quote}
characterized by \texttt{abs-op-unfold}  \textit{(proven)}:
\begin{quote}\ttfamily\footnotesize\raggedright\noindent
forall(\allowbreak{}[kqv\_, kqt\_], abs-op(\allowbreak{}kqv\_, kqt\_) == op-sqrt(\allowbreak{}kqv\_, bop-comp(\allowbreak{}kqv\_, op-adjoint(\allowbreak{}kqv\_, kqv\_, kqt\_), kqt\_)))
\end{quote}
mentioned by 2 result(s)

\smallskip

\noindent\textbf{\texttt{act(s)}}\quad\textit{accessor}\index{act(s)@\texttt{act(s)}}\\
reads: the scalar action of \$1 \\
a structure slot \\
characterized by \texttt{rr-nvs@act}  \textit{(definitional)}:
\begin{quote}\ttfamily\footnotesize\raggedright\noindent
act(\allowbreak{}rr-nvs) == vnb-lambda(\allowbreak{}[r\_, x\_], cartesian(\allowbreak{}rr, rr), r\_ * x\_)
\end{quote}
also: \texttt{qq-line@act} \texttt{cips-act-type} \texttt{is-e-character} \texttt{module-act-type} \texttt{is-e-bicharacter} \texttt{module-act-unital} \texttt{sipq-act-csip-as-sip} \texttt{submodule-act-closed} ... (671 in all) \\
mentioned by 3363 result(s)

\smallskip

\noindent\textbf{\texttt{add(s)}}\quad\textit{accessor}\index{add(s)@\texttt{add(s)}}\\
reads: the addition of \$1 \\
a structure slot \\
characterized by \texttt{rq-add}  \textit{(definitional)}:
\begin{quote}\ttfamily\footnotesize\raggedright\noindent
forall(\allowbreak{}[r], add(\allowbreak{}ringoid-quotient-ring(\allowbreak{}r)) == descend2(\allowbreak{}ringoid-setoid(\allowbreak{}r), vnb-lambda(\allowbreak{}[a, b], cartesian(\allowbreak{}carr(\allowbreak{}r), carr(\allowbreak{}r)), class(\allowbreak{}ringoid-setoid(\allowbreak{}r), (\allowbreak{}add(\allowbreak{}r))(\allowbreak{}a, b)))))
\end{quote}
also: \texttt{qq-ring@add} \texttt{zz-ring@add} \texttt{qq-field@add} \texttt{rr-field@add} \texttt{ring-add-comm} \texttt{ring-add-assoc} \texttt{is-add-subgroup} \texttt{is-add-character} ... (145 in all) \\
mentioned by 882 result(s)

\smallskip

\noindent\textbf{\texttt{adj(r, n, a)}}\quad\textit{functoid}\index{adj(r, n, a)@\texttt{adj(r, n, a)}}\\
reads: adj(\$3) \\
definition:
\begin{quote}\ttfamily\footnotesize\raggedright\noindent
matof(\allowbreak{}succ(\allowbreak{}n), succ(\allowbreak{}n), vnb-lambda(\allowbreak{}[adu\_, adv\_], cartesian(\allowbreak{}interval(\allowbreak{}1, succ(\allowbreak{}n)), interval(\allowbreak{}1, succ(\allowbreak{}n))), (\allowbreak{}mul(\allowbreak{}r))(\allowbreak{}mpow(\allowbreak{}ring-multiplicative-monoid(\allowbreak{}r), (\allowbreak{}neg(\allowbreak{}r))(\allowbreak{}one(\allowbreak{}r)), adu\_ + adv\_), det(\allowbreak{}r, n, minor(\allowbreak{}a, adv\_, adu\_, n)))))
\end{quote}
characterized by \texttt{adj-mul}  \textit{(proven)}:
\begin{quote}\ttfamily\footnotesize\raggedright\noindent
forall(\allowbreak{}[r, n, a], is-commutative-ring(\allowbreak{}r) implies n in nn implies a in mat(\allowbreak{}succ(\allowbreak{}n), succ(\allowbreak{}n), carr(\allowbreak{}r)) implies matmul(\allowbreak{}r, adj(\allowbreak{}r, n, a), a) = matscale(\allowbreak{}r, det(\allowbreak{}r, succ(\allowbreak{}n), a), identmat(\allowbreak{}r, succ(\allowbreak{}n))))
\end{quote}
also: \texttt{mul-adj} \texttt{adj-type} \texttt{adj-entry} \texttt{adj-entries-integral} \\
mentioned by 9 result(s)

\smallskip

\noindent\textbf{\texttt{adjoin-char-value(a6vg\_, a6vh\_, a6vc\_, a6vb\_, a6vt\_, a6vy\_)}}\quad\textit{functoid}\index{adjoin-char-value(a6vg_, a6vh_, a6vc_, a6vb_, a6vt_, a6vy_)@\texttt{adjoin-char-value(a6vg\_, a6vh\_, a6vc\_, a6vb\_, a6vt\_, a6vy\_)}}\\
reads: the value at \$6 of the extension of the character \$3 of \$2 to \$2 + <\$4> with \$4 |-> exp(\$5) \\
definition:
\begin{quote}\ttfamily\footnotesize\raggedright\noindent
iota(\allowbreak{}a6vw\_, a6vw\_ in cc and forsome(\allowbreak{}[a6vq\_ in a6vh\_, a6vl\_ in zz], a6vy\_ = (\allowbreak{}opr(\allowbreak{}a6vg\_))(\allowbreak{}a6vq\_, zz-act(\allowbreak{}a6vg\_, a6vl\_, a6vb\_)) and a6vw\_ = a6vc\_(\allowbreak{}a6vq\_) * cc-exp(\allowbreak{}a6vl\_ * a6vt\_)))
\end{quote}
characterized by \texttt{adjoin-char-value-spec}  \textit{(proven)}:
\begin{quote}\ttfamily\footnotesize\raggedright\noindent
forall(\allowbreak{}[a6g], is-abelian-group(\allowbreak{}a6g) implies forall(\allowbreak{}[a6h], is-subgroup(\allowbreak{}a6g, a6h) implies forall(\allowbreak{}[a6c in fun(\allowbreak{}a6h, unit-circle)], forall(\allowbreak{}[a6u in a6h, a6v in a6h], a6c(\allowbreak{}(\allowbreak{}opr(\allowbreak{}a6g))(\allowbreak{}a6u, a6v)) = a6c(\allowbreak{}a6u) * a6c(\allowbreak{}a6v)) implies forall(\allowbreak{}[a6b in carr(\allowbreak{}a6g), a6m in nn], 1 <= a6m implies zz-act(\allowbreak{}a6g, a6m, a6b) in a6h implies forall(\allowbreak{}[a6j in nn], 1 <= a6j implies a6j < a6m implies not(\allowbreak{}zz-act(\allowbreak{}a6g, a6j, a6b) in a6h)) implies forall(\allowbreak{}[a6t in cc], a6c(\allowbreak{}zz-act(\allowbreak{}a6g, a6m, a6b)) = cc-exp(\allowbreak{}a6m * a6t) implies forall(\allowbreak{}[a6pp in a6h, a6kk in zz], adjoin-char-value(\allowbreak{}a6g, a6h, a6c, a6b, a6t, (\allowbreak{}opr(\allowbreak{}a6g))(\allowbreak{}a6pp, zz-act(\allowbreak{}a6g, a6kk, a6b))) = a6c(\allowbreak{}a6pp) * cc-exp(\allowbreak{}a6kk * a6t)))))))
\end{quote}
mentioned by 3 result(s)

\smallskip

\noindent\textbf{\texttt{adm-def-cert-hyp(mdt\_)}}\quad\textit{predicate}\index{adm-def-cert-hyp(mdt_)@\texttt{adm-def-cert-hyp(mdt\_)}}\\
reads: the core over \$1 is closed under def-cert-hyp \\
defined by \texttt{adm-def-cert-hyp} \\
characterized by \texttt{adm-def-cert-hyp}  \textit{(definitional)}:
\begin{quote}\ttfamily\footnotesize\raggedright\noindent
forall(\allowbreak{}[mdt\_], adm-def-cert-hyp(\allowbreak{}mdt\_) iff forall(\allowbreak{}[mdhs\_, mdc\_], def-cert-hyp(\allowbreak{}mdt\_, mdhs\_, mdc\_) implies hyps-in(\allowbreak{}mdhs\_, derivable(\allowbreak{}mdt\_)) implies mdc\_ in derivable(\allowbreak{}mdt\_)))
\end{quote}
also: \texttt{adm-def-cert-hyp-holds} \\
mentioned by 3 result(s)

\smallskip

\noindent\textbf{\texttt{adm-def-comp(mdt\_)}}\quad\textit{predicate}\index{adm-def-comp(mdt_)@\texttt{adm-def-comp(mdt\_)}}\\
reads: the core over \$1 is closed under def-comp \\
defined by \texttt{adm-def-comp} \\
characterized by \texttt{adm-def-comp}  \textit{(definitional)}:
\begin{quote}\ttfamily\footnotesize\raggedright\noindent
forall(\allowbreak{}[mdt\_], adm-def-comp(\allowbreak{}mdt\_) iff forall(\allowbreak{}[mdhs\_, mdc\_], def-comp(\allowbreak{}mdt\_, mdhs\_, mdc\_) implies hyps-in(\allowbreak{}mdhs\_, derivable(\allowbreak{}mdt\_)) implies mdc\_ in derivable(\allowbreak{}mdt\_)))
\end{quote}
also: \texttt{adm-def-comp-holds} \\
mentioned by 4 result(s)

\smallskip

\noindent\textbf{\texttt{adm-def-fun-inst(mdt\_)}}\quad\textit{predicate}\index{adm-def-fun-inst(mdt_)@\texttt{adm-def-fun-inst(mdt\_)}}\\
reads: the core over \$1 is closed under def-fun-inst \\
defined by \texttt{adm-def-fun-inst} \\
characterized by \texttt{adm-def-fun-inst}  \textit{(definitional)}:
\begin{quote}\ttfamily\footnotesize\raggedright\noindent
forall(\allowbreak{}[mdt\_], adm-def-fun-inst(\allowbreak{}mdt\_) iff forall(\allowbreak{}[mdhs\_, mdc\_], def-fun-inst(\allowbreak{}mdt\_, mdhs\_, mdc\_) implies hyps-in(\allowbreak{}mdhs\_, derivable(\allowbreak{}mdt\_)) implies mdc\_ in derivable(\allowbreak{}mdt\_)))
\end{quote}
mentioned by 2 result(s)

\smallskip

\noindent\textbf{\texttt{adm-def-fun-nth(mdt\_)}}\quad\textit{predicate}\index{adm-def-fun-nth(mdt_)@\texttt{adm-def-fun-nth(mdt\_)}}\\
reads: the core over \$1 is closed under def-fun-nth \\
defined by \texttt{adm-def-fun-nth} \\
characterized by \texttt{adm-def-fun-nth}  \textit{(definitional)}:
\begin{quote}\ttfamily\footnotesize\raggedright\noindent
forall(\allowbreak{}[mdt\_], adm-def-fun-nth(\allowbreak{}mdt\_) iff forall(\allowbreak{}[mdhs\_, mdc\_], def-fun-nth(\allowbreak{}mdt\_, mdhs\_, mdc\_) implies hyps-in(\allowbreak{}mdhs\_, derivable(\allowbreak{}mdt\_)) implies mdc\_ in derivable(\allowbreak{}mdt\_)))
\end{quote}
mentioned by 2 result(s)

\smallskip

\noindent\textbf{\texttt{adm-def-fun-sep(mdt\_)}}\quad\textit{predicate}\index{adm-def-fun-sep(mdt_)@\texttt{adm-def-fun-sep(mdt\_)}}\\
reads: the core over \$1 is closed under def-fun-sep \\
defined by \texttt{adm-def-fun-sep} \\
characterized by \texttt{adm-def-fun-sep}  \textit{(definitional)}:
\begin{quote}\ttfamily\footnotesize\raggedright\noindent
forall(\allowbreak{}[mdt\_], adm-def-fun-sep(\allowbreak{}mdt\_) iff forall(\allowbreak{}[mdhs\_, mdc\_], def-fun-sep(\allowbreak{}mdt\_, mdhs\_, mdc\_) implies hyps-in(\allowbreak{}mdhs\_, derivable(\allowbreak{}mdt\_)) implies mdc\_ in derivable(\allowbreak{}mdt\_)))
\end{quote}
mentioned by 2 result(s)

\smallskip

\noindent\textbf{\texttt{adm-def-fun-strict(mdt\_)}}\quad\textit{predicate}\index{adm-def-fun-strict(mdt_)@\texttt{adm-def-fun-strict(mdt\_)}}\\
reads: the core over \$1 is closed under def-fun-strict \\
defined by \texttt{adm-def-fun-strict} \\
characterized by \texttt{adm-def-fun-strict}  \textit{(definitional)}:
\begin{quote}\ttfamily\footnotesize\raggedright\noindent
forall(\allowbreak{}[mdt\_], adm-def-fun-strict(\allowbreak{}mdt\_) iff forall(\allowbreak{}[mdhs\_, mdc\_], def-fun-strict(\allowbreak{}mdt\_, mdhs\_, mdc\_) implies hyps-in(\allowbreak{}mdhs\_, derivable(\allowbreak{}mdt\_)) implies mdc\_ in derivable(\allowbreak{}mdt\_)))
\end{quote}
mentioned by 2 result(s)

\smallskip

\noindent\textbf{\texttt{adm-def-nth-list(mdt\_)}}\quad\textit{predicate}\index{adm-def-nth-list(mdt_)@\texttt{adm-def-nth-list(mdt\_)}}\\
reads: the core over \$1 is closed under def-nth-list \\
defined by \texttt{adm-def-nth-list} \\
characterized by \texttt{adm-def-nth-list}  \textit{(definitional)}:
\begin{quote}\ttfamily\footnotesize\raggedright\noindent
forall(\allowbreak{}[mdt\_], adm-def-nth-list(\allowbreak{}mdt\_) iff forall(\allowbreak{}[mdhs\_, mdc\_], def-nth-list(\allowbreak{}mdt\_, mdhs\_, mdc\_) implies hyps-in(\allowbreak{}mdhs\_, derivable(\allowbreak{}mdt\_)) implies mdc\_ in derivable(\allowbreak{}mdt\_)))
\end{quote}
also: \texttt{adm-def-nth-list-holds} \\
mentioned by 3 result(s)

\smallskip

\noindent\textbf{\texttt{adm-def-pred-hyp(mdt\_)}}\quad\textit{predicate}\index{adm-def-pred-hyp(mdt_)@\texttt{adm-def-pred-hyp(mdt\_)}}\\
reads: the core over \$1 is closed under def-pred-hyp \\
defined by \texttt{adm-def-pred-hyp} \\
characterized by \texttt{adm-def-pred-hyp}  \textit{(definitional)}:
\begin{quote}\ttfamily\footnotesize\raggedright\noindent
forall(\allowbreak{}[mdt\_], adm-def-pred-hyp(\allowbreak{}mdt\_) iff forall(\allowbreak{}[mdhs\_, mdc\_], def-pred-hyp(\allowbreak{}mdt\_, mdhs\_, mdc\_) implies hyps-in(\allowbreak{}mdhs\_, derivable(\allowbreak{}mdt\_)) implies mdc\_ in derivable(\allowbreak{}mdt\_)))
\end{quote}
also: \texttt{adm-def-pred-hyp-holds} \\
mentioned by 3 result(s)

\smallskip

\noindent\textbf{\texttt{adm-def-pred-hyp-n(mdt\_)}}\quad\textit{predicate}\index{adm-def-pred-hyp-n(mdt_)@\texttt{adm-def-pred-hyp-n(mdt\_)}}\\
reads: the core over \$1 is closed under def-pred-hyp-n \\
defined by \texttt{adm-def-pred-hyp-n} \\
characterized by \texttt{adm-def-pred-hyp-n}  \textit{(definitional)}:
\begin{quote}\ttfamily\footnotesize\raggedright\noindent
forall(\allowbreak{}[mdt\_], adm-def-pred-hyp-n(\allowbreak{}mdt\_) iff forall(\allowbreak{}[mdhs\_, mdc\_], def-pred-hyp-n(\allowbreak{}mdt\_, mdhs\_, mdc\_) implies hyps-in(\allowbreak{}mdhs\_, derivable(\allowbreak{}mdt\_)) implies mdc\_ in derivable(\allowbreak{}mdt\_)))
\end{quote}
also: \texttt{adm-def-pred-hyp-n-holds} \\
mentioned by 3 result(s)

\smallskip

\noindent\textbf{\texttt{adm-def-pred-unfold(mdt\_)}}\quad\textit{predicate}\index{adm-def-pred-unfold(mdt_)@\texttt{adm-def-pred-unfold(mdt\_)}}\\
reads: the core over \$1 is closed under def-pred-unfold \\
defined by \texttt{adm-def-pred-unfold} \\
characterized by \texttt{adm-def-pred-unfold}  \textit{(definitional)}:
\begin{quote}\ttfamily\footnotesize\raggedright\noindent
forall(\allowbreak{}[mdt\_], adm-def-pred-unfold(\allowbreak{}mdt\_) iff forall(\allowbreak{}[mdhs\_, mdc\_], def-pred-unfold(\allowbreak{}mdt\_, mdhs\_, mdc\_) implies hyps-in(\allowbreak{}mdhs\_, derivable(\allowbreak{}mdt\_)) implies mdc\_ in derivable(\allowbreak{}mdt\_)))
\end{quote}
also: \texttt{adm-def-pred-unfold-holds} \\
mentioned by 3 result(s)

\smallskip

\noindent\textbf{\texttt{adm-def-pred-unfold-n(mdt\_)}}\quad\textit{predicate}\index{adm-def-pred-unfold-n(mdt_)@\texttt{adm-def-pred-unfold-n(mdt\_)}}\\
reads: the core over \$1 is closed under def-pred-unfold-n \\
defined by \texttt{adm-def-pred-unfold-n} \\
characterized by \texttt{adm-def-pred-unfold-n}  \textit{(definitional)}:
\begin{quote}\ttfamily\footnotesize\raggedright\noindent
forall(\allowbreak{}[mdt\_], adm-def-pred-unfold-n(\allowbreak{}mdt\_) iff forall(\allowbreak{}[mdhs\_, mdc\_], def-pred-unfold-n(\allowbreak{}mdt\_, mdhs\_, mdc\_) implies hyps-in(\allowbreak{}mdhs\_, derivable(\allowbreak{}mdt\_)) implies mdc\_ in derivable(\allowbreak{}mdt\_)))
\end{quote}
also: \texttt{adm-def-pred-unfold-n-holds} \\
mentioned by 3 result(s)

\smallskip

\noindent\textbf{\texttt{adm-def-strict-arg(mdt\_)}}\quad\textit{predicate}\index{adm-def-strict-arg(mdt_)@\texttt{adm-def-strict-arg(mdt\_)}}\\
reads: the core over \$1 is closed under def-strict-arg \\
defined by \texttt{adm-def-strict-arg} \\
characterized by \texttt{adm-def-strict-arg}  \textit{(definitional)}:
\begin{quote}\ttfamily\footnotesize\raggedright\noindent
forall(\allowbreak{}[mdt\_], adm-def-strict-arg(\allowbreak{}mdt\_) iff forall(\allowbreak{}[mdhs\_, mdc\_], def-strict-arg(\allowbreak{}mdt\_, mdhs\_, mdc\_) implies hyps-in(\allowbreak{}mdhs\_, derivable(\allowbreak{}mdt\_)) implies mdc\_ in derivable(\allowbreak{}mdt\_)))
\end{quote}
also: \texttt{adm-def-strict-arg-holds} \\
mentioned by 3 result(s)

\smallskip

\noindent\textbf{\texttt{adm-def-strict-dom(mdt\_)}}\quad\textit{predicate}\index{adm-def-strict-dom(mdt_)@\texttt{adm-def-strict-dom(mdt\_)}}\\
reads: the core over \$1 is closed under def-strict-dom \\
defined by \texttt{adm-def-strict-dom} \\
characterized by \texttt{adm-def-strict-dom}  \textit{(definitional)}:
\begin{quote}\ttfamily\footnotesize\raggedright\noindent
forall(\allowbreak{}[mdt\_], adm-def-strict-dom(\allowbreak{}mdt\_) iff forall(\allowbreak{}[mdhs\_, mdc\_], def-strict-dom(\allowbreak{}mdt\_, mdhs\_, mdc\_) implies hyps-in(\allowbreak{}mdhs\_, derivable(\allowbreak{}mdt\_)) implies mdc\_ in derivable(\allowbreak{}mdt\_)))
\end{quote}
also: \texttt{adm-def-strict-dom-holds} \\
mentioned by 3 result(s)

\smallskip

\noindent\textbf{\texttt{adm-def-strict-hyp(mdt\_)}}\quad\textit{predicate}\index{adm-def-strict-hyp(mdt_)@\texttt{adm-def-strict-hyp(mdt\_)}}\\
reads: the core over \$1 is closed under def-strict-hyp \\
defined by \texttt{adm-def-strict-hyp} \\
characterized by \texttt{adm-def-strict-hyp}  \textit{(definitional)}:
\begin{quote}\ttfamily\footnotesize\raggedright\noindent
forall(\allowbreak{}[mdt\_], adm-def-strict-hyp(\allowbreak{}mdt\_) iff forall(\allowbreak{}[mdhs\_, mdc\_], def-strict-hyp(\allowbreak{}mdt\_, mdhs\_, mdc\_) implies hyps-in(\allowbreak{}mdhs\_, derivable(\allowbreak{}mdt\_)) implies mdc\_ in derivable(\allowbreak{}mdt\_)))
\end{quote}
also: \texttt{adm-def-strict-hyp-holds} \\
mentioned by 3 result(s)

\smallskip

\noindent\textbf{\texttt{adm-def-strict-in-l(mdt\_)}}\quad\textit{predicate}\index{adm-def-strict-in-l(mdt_)@\texttt{adm-def-strict-in-l(mdt\_)}}\\
reads: the core over \$1 is closed under def-strict-in-l \\
defined by \texttt{adm-def-strict-in-l} \\
characterized by \texttt{adm-def-strict-in-l}  \textit{(definitional)}:
\begin{quote}\ttfamily\footnotesize\raggedright\noindent
forall(\allowbreak{}[mdt\_], adm-def-strict-in-l(\allowbreak{}mdt\_) iff forall(\allowbreak{}[mdhs\_, mdc\_], def-strict-in-l(\allowbreak{}mdt\_, mdhs\_, mdc\_) implies hyps-in(\allowbreak{}mdhs\_, derivable(\allowbreak{}mdt\_)) implies mdc\_ in derivable(\allowbreak{}mdt\_)))
\end{quote}
also: \texttt{adm-def-strict-in-l-holds} \\
mentioned by 5 result(s)

\smallskip

\noindent\textbf{\texttt{adm-def-strict-in-r(mdt\_)}}\quad\textit{predicate}\index{adm-def-strict-in-r(mdt_)@\texttt{adm-def-strict-in-r(mdt\_)}}\\
reads: the core over \$1 is closed under def-strict-in-r \\
defined by \texttt{adm-def-strict-in-r} \\
characterized by \texttt{adm-def-strict-in-r}  \textit{(definitional)}:
\begin{quote}\ttfamily\footnotesize\raggedright\noindent
forall(\allowbreak{}[mdt\_], adm-def-strict-in-r(\allowbreak{}mdt\_) iff forall(\allowbreak{}[mdhs\_, mdc\_], def-strict-in-r(\allowbreak{}mdt\_, mdhs\_, mdc\_) implies hyps-in(\allowbreak{}mdhs\_, derivable(\allowbreak{}mdt\_)) implies mdc\_ in derivable(\allowbreak{}mdt\_)))
\end{quote}
also: \texttt{adm-def-strict-in-r-holds} \\
mentioned by 4 result(s)

\smallskip

\noindent\textbf{\texttt{adm-def-strict-pred(mdt\_)}}\quad\textit{predicate}\index{adm-def-strict-pred(mdt_)@\texttt{adm-def-strict-pred(mdt\_)}}\\
reads: the core over \$1 is closed under def-strict-pred \\
defined by \texttt{adm-def-strict-pred} \\
characterized by \texttt{adm-def-strict-pred}  \textit{(definitional)}:
\begin{quote}\ttfamily\footnotesize\raggedright\noindent
forall(\allowbreak{}[mdt\_], adm-def-strict-pred(\allowbreak{}mdt\_) iff forall(\allowbreak{}[mdhs\_, mdc\_], def-strict-pred(\allowbreak{}mdt\_, mdhs\_, mdc\_) implies hyps-in(\allowbreak{}mdhs\_, derivable(\allowbreak{}mdt\_)) implies mdc\_ in derivable(\allowbreak{}mdt\_)))
\end{quote}
also: \texttt{adm-def-strict-pred-holds} \\
mentioned by 3 result(s)

\smallskip

\noindent\textbf{\texttt{adm-def-struct-hyp(mdt\_)}}\quad\textit{predicate}\index{adm-def-struct-hyp(mdt_)@\texttt{adm-def-struct-hyp(mdt\_)}}\\
reads: the core over \$1 is closed under def-struct-hyp \\
defined by \texttt{adm-def-struct-hyp} \\
characterized by \texttt{adm-def-struct-hyp}  \textit{(definitional)}:
\begin{quote}\ttfamily\footnotesize\raggedright\noindent
forall(\allowbreak{}[mdt\_], adm-def-struct-hyp(\allowbreak{}mdt\_) iff forall(\allowbreak{}[mdhs\_, mdc\_], def-struct-hyp(\allowbreak{}mdt\_, mdhs\_, mdc\_) implies hyps-in(\allowbreak{}mdhs\_, derivable(\allowbreak{}mdt\_)) implies mdc\_ in derivable(\allowbreak{}mdt\_)))
\end{quote}
also: \texttt{adm-def-struct-hyp-holds} \\
mentioned by 3 result(s)

\smallskip

\noindent\textbf{\texttt{adv-ids(s)}}\quad\textit{accessor}\index{adv-ids(s)@\texttt{adv-ids(s)}}\\
reads: the adv-ids slot of \$1 \\
a structure slot \\
characterized by \texttt{cx-isp@adv-ids}  \textit{(definitional)}:
\begin{quote}\ttfamily\footnotesize\raggedright\noindent
adv-ids(\allowbreak{}cx-isp) == interval(\allowbreak{}5, 8)
\end{quote}
mentioned by 20 result(s)

\smallskip

\noindent\textbf{\texttt{adv-tokens(s)}}\quad\textit{accessor}\index{adv-tokens(s)@\texttt{adv-tokens(s)}}\\
reads: the adv-tokens slot of \$1 \\
a structure slot \\
characterized by \texttt{cx-isp@adv-tokens}  \textit{(definitional)}:
\begin{quote}\ttfamily\footnotesize\raggedright\noindent
adv-tokens(\allowbreak{}cx-isp) == empty-set
\end{quote}
mentioned by 47 result(s)

\smallskip

\noindent\textbf{\texttt{ae-rel(omega, ca, mu, sf)}}\quad\textit{functoid}\index{ae-rel(omega, ca, mu, sf)@\texttt{ae-rel(omega, ca, mu, sf)}}\\
reads: equality almost everywhere on \$4 with respect to \$3 \\
definition:
\begin{quote}\ttfamily\footnotesize\raggedright\noindent
\{aep\_ in cartesian(\allowbreak{}sf, sf): forsome(\allowbreak{}[aen\_ in set], is-null-set(\allowbreak{}omega, ca, mu, aen\_) and forall(\allowbreak{}[aex\_ in omega], not(\allowbreak{}aex\_ in aen\_) implies (\allowbreak{}nth(\allowbreak{}1, aep\_))(\allowbreak{}aex\_) = (\allowbreak{}nth(\allowbreak{}2, aep\_))(\allowbreak{}aex\_)))\}
\end{quote}
characterized by \texttt{ae-rel-in}  \textit{(proven)}:
\begin{quote}\ttfamily\footnotesize\raggedright\noindent
forall(\allowbreak{}[omega, ca, mu, sf, af\_ in sf, ag\_ in sf, ns\_], is-null-set(\allowbreak{}omega, ca, mu, ns\_) implies forall(\allowbreak{}[ax\_ in omega], not(\allowbreak{}ax\_ in ns\_) implies af\_(\allowbreak{}ax\_) = ag\_(\allowbreak{}ax\_)) implies [af\_, ag\_] in ae-rel(\allowbreak{}omega, ca, mu, sf))
\end{quote}
also: \texttt{ae-rel-out} \texttt{ae-rel-unfold} \texttt{ae-rel-is-equivalence} \texttt{l2-null-rel-is-ae-rel} \\
mentioned by 7 result(s)

\smallskip

\noindent\textbf{\texttt{ae-setoid(omega, ca, mu, sf)}}\quad\textit{functoid}\index{ae-setoid(omega, ca, mu, sf)@\texttt{ae-setoid(omega, ca, mu, sf)}}\\
reads: \$4 modulo equality almost everywhere with respect to \$3 \\
definition:
\begin{quote}\ttfamily\footnotesize\raggedright\noindent
[sf, ae-rel(\allowbreak{}omega, ca, mu, sf)]
\end{quote}
characterized by \texttt{ae-setoid-is-setoid}  \textit{(proven)}:
\begin{quote}\ttfamily\footnotesize\raggedright\noindent
forall(\allowbreak{}[omega, ca, mu, sf], is-measure(\allowbreak{}omega, ca, mu) implies sf in set implies sf subset fun(\allowbreak{}omega, rr) implies is-setoid(\allowbreak{}ae-setoid(\allowbreak{}omega, ca, mu, sf)))
\end{quote}
also: \texttt{l2-setoid-is-ae-setoid} \\
mentioned by 5 result(s)

\smallskip

\noindent\textbf{\texttt{aff-der(hmp\_, hmq\_, hmc\_, hmd\_)}}\quad\textit{functoid}\index{aff-der(hmp_, hmq_, hmc_, hmd_)@\texttt{aff-der(hmp\_, hmq\_, hmc\_, hmd\_)}}\\
reads: the derivative of the affine segment from \$1 to \$2 on [\$3, \$4] \\
definition:
\begin{quote}\ttfamily\footnotesize\raggedright\noindent
vnb-lambda(\allowbreak{}hmz\_, ccint(\allowbreak{}hmc\_, hmd\_), recip(\allowbreak{}hmd\_ - hmc\_) * (\allowbreak{}hmq\_ - hmp\_))
\end{quote}
mentioned by 5 result(s)

\smallskip

\noindent\textbf{\texttt{aff-seg(hmp\_, hmq\_, hmc\_, hmd\_)}}\quad\textit{functoid}\index{aff-seg(hmp_, hmq_, hmc_, hmd_)@\texttt{aff-seg(hmp\_, hmq\_, hmc\_, hmd\_)}}\\
reads: the affine segment from \$1 to \$2 on [\$3, \$4] \\
definition:
\begin{quote}\ttfamily\footnotesize\raggedright\noindent
vnb-lambda(\allowbreak{}hmz\_, ccint(\allowbreak{}hmc\_, hmd\_), hmp\_ + (\allowbreak{}(\allowbreak{}hmz\_ - hmc\_) * recip(\allowbreak{}hmd\_ - hmc\_)) * (\allowbreak{}hmq\_ - hmp\_))
\end{quote}
characterized by \texttt{aff-seg-ends}  \textit{(proven)}:
\begin{quote}\ttfamily\footnotesize\raggedright\noindent
forall(\allowbreak{}[p, q, c, d], p in cc implies q in cc implies c in rr implies d in rr implies c < d implies (\allowbreak{}aff-seg(\allowbreak{}p, q, c, d))(\allowbreak{}c) = p and (\allowbreak{}aff-seg(\allowbreak{}p, q, c, d))(\allowbreak{}d) = q)
\end{quote}
also: \texttt{aff-seg-is-road} \texttt{aff-seg-in-convex} \\
mentioned by 8 result(s)

\smallskip

\noindent\textbf{\texttt{ag-zz-module(qag\_)}}\quad\textit{functoid}\index{ag-zz-module(qag_)@\texttt{ag-zz-module(qag\_)}}\\
reads: the abelian group \$1 as a Z-module \\
definition:
\begin{quote}\ttfamily\footnotesize\raggedright\noindent
[zz-ring, carr(\allowbreak{}qag\_), opr(\allowbreak{}qag\_), iden(\allowbreak{}qag\_), inv(\allowbreak{}qag\_), vnb-lambda(\allowbreak{}[azk\_, aza\_], cartesian(\allowbreak{}zz, carr(\allowbreak{}qag\_)), zz-act(\allowbreak{}qag\_, azk\_, aza\_))]
\end{quote}
characterized by \texttt{ag-zz-module-is-module}  \textit{(proven)}:
\begin{quote}\ttfamily\footnotesize\raggedright\noindent
forall(\allowbreak{}[g], is-abelian-group(\allowbreak{}g) implies is-module(\allowbreak{}ag-zz-module(\allowbreak{}g)))
\end{quote}
mentioned by 48 result(s)

\smallskip

\noindent\textbf{\texttt{agree-length(alx\_, aly\_)}}\quad\textit{functoid}\index{agree-length(alx_, aly_)@\texttt{agree-length(alx\_, aly\_)}}\\
reads: the length of the common prefix of \$1 and \$2 \\
definition:
\begin{quote}\ttfamily\footnotesize\raggedright\noindent
iota(\allowbreak{}alk\_, alk\_ in nn and not(\allowbreak{}alx\_(\allowbreak{}alk\_) = aly\_(\allowbreak{}alk\_)) and forall(\allowbreak{}[alj\_ in nn], alj\_ < alk\_ implies alx\_(\allowbreak{}alj\_) = aly\_(\allowbreak{}alj\_)))
\end{quote}
characterized by \texttt{agree-length-sym}  \textit{(proven)}:
\begin{quote}\ttfamily\footnotesize\raggedright\noindent
forall(\allowbreak{}[sx\_, xs, ys], xs in fun(\allowbreak{}nn, sx\_) implies ys in fun(\allowbreak{}nn, sx\_) implies not(\allowbreak{}xs = ys) implies agree-length(\allowbreak{}xs, ys) = agree-length(\allowbreak{}ys, xs))
\end{quote}
also: \texttt{agree-length-prop} \\
mentioned by 4 result(s)

\smallskip

\noindent\textbf{\texttt{ai}}\quad\textit{tactic}\index{ai@\texttt{ai}}\\
\texttt{(ai hyp)} -- Antecedent inference: decompose a cited assumption -- AND-split, OR-into-cases, or FORSOME-elimination to a fresh eigenvariable. \\
declared in \texttt{tactics-help.scm}

\smallskip

\noindent\textbf{\texttt{algebraic-numbers-over(kk)}}\quad\textit{functoid}\index{algebraic-numbers-over(kk)@\texttt{algebraic-numbers-over(kk)}}\\
reads: the algebraic numbers over \$1 \\
definition:
\begin{quote}\ttfamily\footnotesize\raggedright\noindent
\{anth\_ in cc: is-algebraic-over-field(\allowbreak{}kk, anth\_)\}
\end{quote}
characterized by \texttt{algebraic-numbers-over-unfold}  \textit{(proven)}:
\begin{quote}\ttfamily\footnotesize\raggedright\noindent
forall(\allowbreak{}[kk], algebraic-numbers-over(\allowbreak{}kk) == \{anth\_ in cc: is-algebraic-over-field(\allowbreak{}kk, anth\_)\})
\end{quote}
also: \texttt{algebraic-numbers-over-qq-countable} \texttt{algebraic-numbers-over-subfield-field} \\
mentioned by 8 result(s)

\smallskip

\noindent\textbf{\texttt{all-confined(mnbb\_)}}\quad\textit{predicate}\index{all-confined(mnbb_)@\texttt{all-confined(mnbb\_)}}\\
reads: the universal with body \$1 is confined to sets \\
defined by \texttt{all-confined} \\
characterized by \texttt{all-confined}  \textit{(definitional)}:
\begin{quote}\ttfamily\footnotesize\raggedright\noindent
forall(\allowbreak{}[mnbb\_], all-confined(\allowbreak{}mnbb\_) iff forsome(\allowbreak{}[mnbqx\_, mnbqy\_, mnbqd\_, mnbqc\_], mnbqx\_ in expr-set and mnbqy\_ in expr-set and mnbqd\_ in nn and mnbqc\_ in expr-set and [enc-imp(\allowbreak{}mnbqx\_, mnbqy\_), mnbqd\_] in all-spine(\allowbreak{}mnbb\_) and enc-in(\allowbreak{}ebvar(\allowbreak{}mnbqd\_), mnbqc\_) in conj-reach(\allowbreak{}mnbqx\_)))
\end{quote}
also: \texttt{all-confined-guarded} \\
mentioned by 5 result(s)

\smallskip

\noindent\textbf{\texttt{all-spine(mnbb\_)}}\quad\textit{functoid}\index{all-spine(mnbb_)@\texttt{all-spine(mnbb\_)}}\\
reads: the implicational spine of the body \$1, with depths \\
definition:
\begin{quote}\ttfamily\footnotesize\raggedright\noindent
\{mnbz\_ in cartesian(\allowbreak{}expr-set, nn): forall(\allowbreak{}[mnbv\_ in power(\allowbreak{}cartesian(\allowbreak{}expr-set, nn))], [mnbb\_, 0] in mnbv\_ implies all-spine-closed(\allowbreak{}mnbv\_) implies mnbz\_ in mnbv\_)\}
\end{quote}
characterized by \texttt{all-spine-least}  \textit{(proven)}:
\begin{quote}\ttfamily\footnotesize\raggedright\noindent
forall(\allowbreak{}[b, mnvxs\_ in power(\allowbreak{}cartesian(\allowbreak{}expr-set, nn))], [b, 0] in mnvxs\_ implies all-spine-closed(\allowbreak{}mnvxs\_) implies forall(\allowbreak{}[x in all-spine(\allowbreak{}b)], x in mnvxs\_))
\end{quote}
also: \texttt{all-spine-unfold} \\
mentioned by 5 result(s)

\smallskip

\noindent\textbf{\texttt{all-spine-closed(mnbw\_)}}\quad\textit{predicate}\index{all-spine-closed(mnbw_)@\texttt{all-spine-closed(mnbw\_)}}\\
reads: \$1 is closed under steps down implicational spines \\
defined by \texttt{all-spine-closed} \\
characterized by \texttt{all-spine-closed}  \textit{(definitional)}:
\begin{quote}\ttfamily\footnotesize\raggedright\noindent
forall(\allowbreak{}[mnbw\_], all-spine-closed(\allowbreak{}mnbw\_) iff forall(\allowbreak{}[mnbce\_ in mnbw\_, mnbcf\_ in cartesian(\allowbreak{}expr-set, nn)], all-spine-step(\allowbreak{}mnbce\_, mnbcf\_) implies mnbcf\_ in mnbw\_))
\end{quote}
mentioned by 5 result(s)

\smallskip

\noindent\textbf{\texttt{all-spine-step(mnbu\_, mnbt\_)}}\quad\textit{predicate}\index{all-spine-step(mnbu_, mnbt_)@\texttt{all-spine-step(mnbu\_, mnbt\_)}}\\
reads: \$2 is one step down the implicational spine from \$1 \\
defined by \texttt{all-spine-step} \\
characterized by \texttt{all-spine-step}  \textit{(definitional)}:
\begin{quote}\ttfamily\footnotesize\raggedright\noindent
forall(\allowbreak{}[mnbu\_, mnbt\_], all-spine-step(\allowbreak{}mnbu\_, mnbt\_) iff forsome(\allowbreak{}[mnbsc\_, mnbsd\_], mnbsc\_ in expr-set and mnbsd\_ in nn and mnbu\_ = [enc-all(\allowbreak{}mnbsc\_), mnbsd\_] and mnbt\_ = [mnbsc\_, mnbsd\_ + 1]) or forsome(\allowbreak{}[mnbsx\_, mnbsy\_, mnbse\_], mnbsx\_ in expr-set and mnbsy\_ in expr-set and mnbse\_ in nn and mnbu\_ = [enc-imp(\allowbreak{}mnbsx\_, mnbsy\_), mnbse\_] and mnbt\_ = [mnbsy\_, mnbse\_]))
\end{quote}
mentioned by 4 result(s)

\smallskip

\noindent\textbf{\texttt{amp-step(amv\_, amw\_, amf\_)}}\quad\textit{functoid}\index{amp-step(amv_, amw_, amf_)@\texttt{amp-step(amv\_, amw\_, amf\_)}}\\
reads: the representation T |-> \$3(T) (+) T \\
definition:
\begin{quote}\ttfamily\footnotesize\raggedright\noindent
vnb-lambda(\allowbreak{}amt\_, bounded-ops(\allowbreak{}amv\_, amv\_), hsum-diag(\allowbreak{}amw\_, amv\_, amf\_(\allowbreak{}amt\_), amt\_))
\end{quote}
characterized by \texttt{amp-step-type}  \textit{(proven)}:
\begin{quote}\ttfamily\footnotesize\raggedright\noindent
forall(\allowbreak{}[w, v, f\_], is-complex-inner-product-space(\allowbreak{}w) implies is-complex-inner-product-space(\allowbreak{}v) implies f\_ in fun(\allowbreak{}bounded-ops(\allowbreak{}v, v), bounded-ops(\allowbreak{}w, w)) implies amp-step(\allowbreak{}v, w, f\_) in fun(\allowbreak{}bounded-ops(\allowbreak{}v, v), bounded-ops(\allowbreak{}hsum(\allowbreak{}w, v), hsum(\allowbreak{}w, v))))
\end{quote}
also: \texttt{amp-step-apply} \texttt{amp-step-star-rep} \texttt{amp-step-norm-ge-1} \texttt{amp-step-norm-ge-2} \\
mentioned by 9 result(s)

\smallskip

\noindent\textbf{\texttt{angle(ang1\_, ang2\_)}}\quad\textit{functoid}\index{angle(ang1_, ang2_)@\texttt{angle(ang1\_, ang2\_)}}\\
reads: the angle from \$2 to \$1 \\
definition:
\begin{quote}\ttfamily\footnotesize\raggedright\noindent
pi + imag-part(\allowbreak{}choice(\allowbreak{}\{angz\_ in cc: -pi <= imag-part(\allowbreak{}angz\_) and imag-part(\allowbreak{}angz\_) < pi and cc-exp(\allowbreak{}angz\_) = -(\allowbreak{}ang1\_ * conjugate(\allowbreak{}ang2\_))\}))
\end{quote}
characterized by \texttt{angle-lt-pi-iff}  \textit{(proven)}:
\begin{quote}\ttfamily\footnotesize\raggedright\noindent
forall(\allowbreak{}[anz\_, anv\_], anz\_ in cc implies anv\_ in cc implies not(\allowbreak{}anz\_ = 0) implies not(\allowbreak{}anv\_ = 0) implies angle(\allowbreak{}anz\_, anv\_) < pi iff 0 < imag-part(\allowbreak{}anz\_ * conjugate(\allowbreak{}anv\_)) or imag-part(\allowbreak{}anz\_ * conjugate(\allowbreak{}anv\_)) = 0 and 0 < real-part(\allowbreak{}anz\_ * conjugate(\allowbreak{}anv\_)))
\end{quote}
also: \texttt{interior-angle-sign} \\
mentioned by 8 result(s)

\smallskip

\noindent\textbf{\texttt{annulus(anla\_, anlr\_, anls\_)}}\quad\textit{functoid}\index{annulus(anla_, anlr_, anls_)@\texttt{annulus(anla\_, anlr\_, anls\_)}}\\
reads: the open annulus about \$1 between the radii \$2 and \$3 \\
definition:
\begin{quote}\ttfamily\footnotesize\raggedright\noindent
\{anlz\_ in cc: anlr\_ < magnitude(\allowbreak{}anlz\_ - anla\_) and magnitude(\allowbreak{}anlz\_ - anla\_) < anls\_\}
\end{quote}
characterized by \texttt{annulus-open}  \textit{(proven)}:
\begin{quote}\ttfamily\footnotesize\raggedright\noindent
forall(\allowbreak{}[a in cc, rm in rr, rp in rr], is-open(\allowbreak{}cc-ms, annulus(\allowbreak{}a, rm, rp)))
\end{quote}
also: \texttt{annulus-unfold} \texttt{annulus-shift-facts} \\
mentioned by 11 result(s)

\smallskip

\noindent\textbf{\texttt{antiderivable-fns(ada\_, adb\_)}}\quad\textit{functoid}\index{antiderivable-fns(ada_, adb_)@\texttt{antiderivable-fns(ada\_, adb\_)}}\\
reads: the set of functions antiderivable on [\$1, \$2] \\
definition:
\begin{quote}\ttfamily\footnotesize\raggedright\noindent
\{adp\_ in fun(\allowbreak{}rr, rr): is-antiderivable(\allowbreak{}adp\_, ada\_, adb\_)\}
\end{quote}
characterized by \texttt{antiderivable-fns-unfold}  \textit{(proven)}:
\begin{quote}\ttfamily\footnotesize\raggedright\noindent
forall(\allowbreak{}[a, b], antiderivable-fns(\allowbreak{}a, b) == \{adp\_ in fun(\allowbreak{}rr, rr): is-antiderivable(\allowbreak{}adp\_, a, b)\})
\end{quote}
also: \texttt{antiderivable-fns-membership} \texttt{antiderivable-fns-vector-space} \\
mentioned by 5 result(s)

\smallskip

\noindent\textbf{\texttt{app-ar-ok(mch\_, mcn\_)}}\quad\textit{predicate}\index{app-ar-ok(mch_, mcn_)@\texttt{app-ar-ok(mch\_, mcn\_)}}\\
reads: \$2 is an arity allowed for the head \$1 \\
defined by \texttt{app-ar-ok} \\
characterized by \texttt{app-ar-ok}  \textit{(definitional)}:
\begin{quote}\ttfamily\footnotesize\raggedright\noindent
forall(\allowbreak{}[mch\_, mcn\_], app-ar-ok(\allowbreak{}mch\_, mcn\_) iff 10 <= mch\_ or mch\_ <= 1 and mcn\_ = 0 or mch\_ = 2 and mcn\_ = 1 or 3 <= mch\_ and mcn\_ = 2)
\end{quote}
mentioned by 8 result(s)

\smallskip

\noindent\textbf{\texttt{apply-functoid}}\quad\textit{operator}\index{apply-functoid@\texttt{apply-functoid}}\\
a term-forming head

\smallskip

\noindent\textbf{\texttt{arith}}\quad\textit{tactic}\index{arith@\texttt{arith}}\\
\texttt{(arith)} -- Discharge a ground arithmetic goal by evaluation. \\
declared in \texttt{tactics-help.scm}

\smallskip

\noindent\textbf{\texttt{ass}}\quad\textit{tactic}\index{ass@\texttt{ass}}\\
\texttt{(ass)} -- Close the goal by an assumption alpha-equivalent to it. \\
declared in \texttt{tactics-help.scm}

\smallskip

\noindent\textbf{\texttt{associated(gss\_, gsa\_, gsb\_)}}\quad\textit{predicate}\index{associated(gss_, gsa_, gsb_)@\texttt{associated(gss\_, gsa\_, gsb\_)}}\\
reads: \$2 and \$3 are associated in \$1 \\
defined by \texttt{associated} \\
characterized by \texttt{associated}  \textit{(definitional)}:
\begin{quote}\ttfamily\footnotesize\raggedright\noindent
forall(\allowbreak{}[gss\_, gsa\_, gsb\_], associated(\allowbreak{}gss\_, gsa\_, gsb\_) iff gsa\_ in carr(\allowbreak{}gss\_) and gsb\_ in carr(\allowbreak{}gss\_) and forsome(\allowbreak{}[gsw\_], is-unit(\allowbreak{}gss\_, gsw\_) and gsb\_ = (\allowbreak{}mul(\allowbreak{}gss\_))(\allowbreak{}gsw\_, gsa\_)))
\end{quote}
also: \texttt{sdet-gcd-associated} \texttt{mutual-divides-associated} \texttt{prime-nat-associated-equal} \\
mentioned by 11 result(s)

\smallskip

\noindent\textbf{\texttt{atoms(ca)}}\quad\textit{functoid}\index{atoms(ca)@\texttt{atoms(ca)}}\\
reads: the atoms of \$1 \\
definition:
\begin{quote}\ttfamily\footnotesize\raggedright\noindent
\{mat\_ in ca: is-atom(\allowbreak{}ca, mat\_)\}
\end{quote}
characterized by \texttt{atoms-unfold}  \textit{(proven)}:
\begin{quote}\ttfamily\footnotesize\raggedright\noindent
forall(\allowbreak{}[ca], atoms(\allowbreak{}ca) == \{mat\_ in ca: is-atom(\allowbreak{}ca, mat\_)\})
\end{quote}
also: \texttt{atoms-membership} \texttt{finite-algebra-atoms} \texttt{atoms-pairwise-disjoint} \\
mentioned by 9 result(s)

\smallskip

\noindent\textbf{\texttt{ax-class-ext}}\quad\textit{constant}\index{ax-class-ext@\texttt{ax-class-ext}}\\
reads: the encoded axiom class-extensionality \\
defined by \texttt{ax-class-ext-def}

\smallskip

\noindent\textbf{\texttt{ax-fun-app}}\quad\textit{constant}\index{ax-fun-app@\texttt{ax-fun-app}}\\
reads: the encoded axiom fun-domain-apply-def \\
defined by \texttt{ax-fun-app-def}

\smallskip

\noindent\textbf{\texttt{ax-limit-iff}}\quad\textit{constant}\index{ax-limit-iff@\texttt{ax-limit-iff}}\\
reads: the encoded axiom limit-ord-iff \\
defined by \texttt{ax-limit-iff-def}

\smallskip

\noindent\textbf{\texttt{ax-mem-set}}\quad\textit{constant}\index{ax-mem-set@\texttt{ax-mem-set}}\\
reads: the encoded axiom membership-implies-sethood \\
defined by \texttt{ax-mem-set-def}

\smallskip

\noindent\textbf{\texttt{ax-nn-ind}}\quad\textit{constant}\index{ax-nn-ind@\texttt{ax-nn-ind}}\\
reads: the encoded axiom nn-induction \\
defined by \texttt{ax-nn-ind-def}

\smallskip

\noindent\textbf{\texttt{ax-ord-succ-above}}\quad\textit{constant}\index{ax-ord-succ-above@\texttt{ax-ord-succ-above}}\\
reads: the encoded axiom ord-succ-above \\
defined by \texttt{ax-ord-succ-above-def}

\smallskip

\noindent\textbf{\texttt{ax-succ-nn}}\quad\textit{constant}\index{ax-succ-nn@\texttt{ax-succ-nn}}\\
reads: the encoded axiom nn-succ-closed \\
defined by \texttt{ax-succ-nn-def}

\smallskip

\noindent\textbf{\texttt{ax-tfi}}\quad\textit{constant}\index{ax-tfi@\texttt{ax-tfi}}\\
reads: the encoded axiom transfinite-induction \\
defined by \texttt{ax-tfi-def}

\smallskip

\noindent\textbf{\texttt{ax-zero-nn}}\quad\textit{constant}\index{ax-zero-nn@\texttt{ax-zero-nn}}\\
reads: the encoded axiom nn-zero-in \\
defined by \texttt{ax-zero-nn-def}

\smallskip


\section*{B}
\addcontentsline{toc}{section}{B}

\noindent\textbf{\texttt{ba-metric-space(baa\_)}}\quad\textit{functoid}\index{ba-metric-space(baa_)@\texttt{ba-metric-space(baa\_)}}\\
reads: the metric space of the Banach algebra \$1 \\
definition:
\begin{quote}\ttfamily\footnotesize\raggedright\noindent
nag-metric-space(\allowbreak{}uba-as-normed-ag(\allowbreak{}baa\_))
\end{quote}
mentioned by 31 result(s)

\smallskip

\noindent\textbf{\texttt{backup-one}}\quad\textit{tactic}\index{backup-one@\texttt{backup-one}}\\
\texttt{(backup-one)} -- Take back the last recorded command: restore the goal, the deduction graph and the proof script to the state before it.  Repeat to walk further back; (sp) clears the stack.  `undo' is an alias. \\
declared in \texttt{tactics-help.scm}

\smallskip

\noindent\textbf{\texttt{baire-code(bcm\_)}}\quad\textit{functoid}\index{baire-code(bcm_)@\texttt{baire-code(bcm\_)}}\\
reads: the sequence of naturals coded by \$1 \\
definition:
\begin{quote}\ttfamily\footnotesize\raggedright\noindent
vnb-lambda(\allowbreak{}bck\_, nn, nnfst(\allowbreak{}snd-iter(\allowbreak{}bcm\_, bck\_)))
\end{quote}
characterized by \texttt{baire-code-facts}  \textit{(proven)}:
\begin{quote}\ttfamily\footnotesize\raggedright\noindent
forall(\allowbreak{}[m in nn], baire-code(\allowbreak{}m) in fun(\allowbreak{}nn, nn) and forall(\allowbreak{}[k in nn], (\allowbreak{}baire-code(\allowbreak{}m))(\allowbreak{}k) = nnfst(\allowbreak{}snd-iter(\allowbreak{}m, k))))
\end{quote}
mentioned by 1 result(s)

\smallskip

\noindent\textbf{\texttt{baire-ms}}\quad\textit{constant}\index{baire-ms@\texttt{baire-ms}}\\
reads: the Baire space of sequences of naturals \\
defined by \texttt{baire-ms-def}

\smallskip

\noindent\textbf{\texttt{ball(s, c, r)}}\quad\textit{functoid}\index{ball(s, c, r)@\texttt{ball(s, c, r)}}\\
definition:
\begin{quote}\ttfamily\footnotesize\raggedright\noindent
\{y in pts(\allowbreak{}s): (\allowbreak{}dist(\allowbreak{}s))(\allowbreak{}c, y) <= r and not(\allowbreak{}(\allowbreak{}dist(\allowbreak{}s))(\allowbreak{}c, y) = r)\}
\end{quote}
characterized by \texttt{ball-is-set}  \textit{(proven)}:
\begin{quote}\ttfamily\footnotesize\raggedright\noindent
forall(\allowbreak{}[s], is-metric-space(\allowbreak{}s) implies forall(\allowbreak{}[x, r], ball(\allowbreak{}s, x, r) in set))
\end{quote}
also: \texttt{ball-is-open} \texttt{bk-ball-open} \texttt{hurwitz-ball} \texttt{ball-point-le} \texttt{ball-point-ne} \texttt{ms-agree-ball} \texttt{subspace-ball} \texttt{ball-center-in} ... (69 in all) \\
mentioned by 279 result(s)

\smallskip

\noindent\textbf{\texttt{ball-cover(s, r)}}\quad\textit{functoid}\index{ball-cover(s, r)@\texttt{ball-cover(s, r)}}\\
definition:
\begin{quote}\ttfamily\footnotesize\raggedright\noindent
image(\allowbreak{}vnb-lambda(\allowbreak{}c, pts(\allowbreak{}s), ball(\allowbreak{}s, c, r)), pts(\allowbreak{}s))
\end{quote}
characterized by \texttt{ball-cover-unfold}  \textit{(proven)}:
\begin{quote}\ttfamily\footnotesize\raggedright\noindent
forall(\allowbreak{}[s, r], ball-cover(\allowbreak{}s, r) == image(\allowbreak{}vnb-lambda(\allowbreak{}c, pts(\allowbreak{}s), ball(\allowbreak{}s, c, r)), pts(\allowbreak{}s)))
\end{quote}
also: \texttt{ball-cover-mem-fwd} \texttt{ball-cover-is-open-cover} \\
mentioned by 6 result(s)

\smallskip

\noindent\textbf{\texttt{bc}}\quad\textit{tactic}\index{bc@\texttt{bc}}\\
\texttt{(bc impl)} -- Backchain the goal through an (IMPLIES A B) already in context: if the goal matches B, the new goal is A. \\
declared in \texttt{tactics-help.scm}

\smallskip

\noindent\textbf{\texttt{bc*}}\quad\textit{tactic}\index{bc*@\texttt{bc*}}\\
\texttt{(bc* 'thm [((v val)...)] h1 h2 ...)} -- Matching backchain: unify the theorem's conclusion against the goal, then leave its (instantiated) antecedents as subgoals -- optional handlers hk run on the k-th subgoal. \\
declared in \texttt{tactics-help.scm}

\smallskip

\noindent\textbf{\texttt{bdd-metric(s)}}\quad\textit{functoid}\index{bdd-metric(s)@\texttt{bdd-metric(s)}}\\
definition:
\begin{quote}\ttfamily\footnotesize\raggedright\noindent
[pts(\allowbreak{}s), vnb-lambda(\allowbreak{}[u, v], cartesian(\allowbreak{}pts(\allowbreak{}s), pts(\allowbreak{}s)), (\allowbreak{}/)(\allowbreak{}(\allowbreak{}dist(\allowbreak{}s))(\allowbreak{}u, v), 1 + (\allowbreak{}dist(\allowbreak{}s))(\allowbreak{}u, v)))]
\end{quote}
characterized by \texttt{bdd-metric-bounded}  \textit{(proven)}:
\begin{quote}\ttfamily\footnotesize\raggedright\noindent
forall(\allowbreak{}[s], is-metric-space(\allowbreak{}s) implies forall(\allowbreak{}[x in pts(\allowbreak{}s), y in pts(\allowbreak{}s)], (\allowbreak{}dist(\allowbreak{}bdd-metric(\allowbreak{}s)))(\allowbreak{}x, y) < 1))
\end{quote}
also: \texttt{bdd-metric-carrier} \texttt{bdd-metric-dist-le} \texttt{bdd-metric-distance} \texttt{bdd-metric-open-bwd} \texttt{bdd-metric-open-fwd} \texttt{bdd-metric-dist-reflect} \texttt{bdd-metric-weight-bound} \texttt{bdd-metric-converges-bwd} ... (14 in all) \\
mentioned by 48 result(s)

\smallskip

\noindent\textbf{\texttt{bernstein-basis(n, x)}}\quad\textit{functoid}\index{bernstein-basis(n, x)@\texttt{bernstein-basis(n, x)}}\\
reads: the Bernstein basis family of degree \$1 at \$2 \\
definition:
\begin{quote}\ttfamily\footnotesize\raggedright\noindent
comb-kk(\allowbreak{}normed-field-as-commutative-ring(\allowbreak{}rr-normed-field), x, 1 - x, n)
\end{quote}
characterized by \texttt{bernstein-basis-null}  \textit{(proven)}:
\begin{quote}\ttfamily\footnotesize\raggedright\noindent
forall(\allowbreak{}[n\_, x\_, k\_], n\_ in nn implies x\_ in rr implies k\_ in zz implies k\_ < 0 implies (\allowbreak{}bernstein-basis(\allowbreak{}n\_, x\_))(\allowbreak{}k\_) = 0)
\end{quote}
also: \texttt{bernstein-basis-succ} \texttt{bernstein-basis-above} \texttt{bernstein-basis-in-fun} \texttt{bernstein-basis-nonneg} \texttt{bernstein-basis-unfold} \texttt{bernstein-basis-poly-fn} \texttt{bernstein-basis-ptwise-in-rr} \texttt{bernstein-basis-antiderivable} ... (11 in all) \\
mentioned by 32 result(s)

\smallskip

\noindent\textbf{\texttt{bernstein-poly(f, n, x)}}\quad\textit{functoid}\index{bernstein-poly(f, n, x)@\texttt{bernstein-poly(f, n, x)}}\\
reads: the degree-\$2 Bernstein polynomial of \$1 at \$3 \\
definition:
\begin{quote}\ttfamily\footnotesize\raggedright\noindent
series-partial-sum(\allowbreak{}vnb-lambda(\allowbreak{}l\_, nn, f(\allowbreak{}(\allowbreak{}/)(\allowbreak{}l\_, n)) * (\allowbreak{}bernstein-basis(\allowbreak{}n, x))(\allowbreak{}l\_)), succ(\allowbreak{}n))
\end{quote}
characterized by \texttt{bernstein-poly-antiderivable}  \textit{(proven)}:
\begin{quote}\ttfamily\footnotesize\raggedright\noindent
forall(\allowbreak{}[f, n\_, a, b], f in fun(\allowbreak{}rr, rr) implies n\_ in nn implies not(\allowbreak{}n\_ = 0) implies a in rr and b in rr and a < b implies is-antiderivable(\allowbreak{}vnb-lambda(\allowbreak{}x, rr, bernstein-poly(\allowbreak{}f, n\_, x)), a, b))
\end{quote}
also: \texttt{bernstein-poly-unit-coordinate-antiderivable} \\
mentioned by 6 result(s)

\smallskip

\noindent\textbf{\texttt{beta}}\quad\textit{tactic}\index{beta@\texttt{beta}}\\
\texttt{(beta)} -- Beta-reduce a functoid application in the goal. \\
declared in \texttt{tactics-help.scm}

\smallskip

\noindent\textbf{\texttt{bh-act(bov\_)}}\quad\textit{functoid}\index{bh-act(bov_)@\texttt{bh-act(bov\_)}}\\
definition:
\begin{quote}\ttfamily\footnotesize\raggedright\noindent
vnb-lambda(\allowbreak{}[boc\_, bot\_], cartesian(\allowbreak{}cc, bounded-ops(\allowbreak{}bov\_, bov\_)), bop-scale(\allowbreak{}bov\_, bov\_, boc\_, bot\_))
\end{quote}
characterized by \texttt{bh-act-add}  \textit{(proven)}:
\begin{quote}\ttfamily\footnotesize\raggedright\noindent
forall(\allowbreak{}[v], is-complex-inner-product-space(\allowbreak{}v) implies forall(\allowbreak{}[r\_ in carr(\allowbreak{}normed-field-as-commutative-ring(\allowbreak{}cc-normed-field)), s\_ in carr(\allowbreak{}normed-field-as-commutative-ring(\allowbreak{}cc-normed-field)), x\_ in bounded-ops(\allowbreak{}v, v)], (\allowbreak{}bh-act(\allowbreak{}v))(\allowbreak{}(\allowbreak{}add(\allowbreak{}normed-field-as-commutative-ring(\allowbreak{}cc-normed-field)))(\allowbreak{}r\_, s\_), x\_) = (\allowbreak{}bh-vadd(\allowbreak{}v))(\allowbreak{}(\allowbreak{}bh-act(\allowbreak{}v))(\allowbreak{}r\_, x\_), (\allowbreak{}bh-act(\allowbreak{}v))(\allowbreak{}s\_, x\_))))
\end{quote}
also: \texttt{bh-act-mul} \texttt{bh-act-one} \texttt{bh-act-type} \texttt{bh-act-vadd} \texttt{bh-act-value} \\
mentioned by 31 result(s)

\smallskip

\noindent\textbf{\texttt{bh-algebra(bov\_)}}\quad\textit{functoid}\index{bh-algebra(bov_)@\texttt{bh-algebra(bov\_)}}\\
reads: the algebra of bounded operators on \$1 \\
definition:
\begin{quote}\ttfamily\footnotesize\raggedright\noindent
[normed-field-as-commutative-ring(\allowbreak{}cc-normed-field), bounded-ops(\allowbreak{}bov\_, bov\_), bh-vadd(\allowbreak{}bov\_), bop-zero(\allowbreak{}bov\_, bov\_), bh-vneg(\allowbreak{}bov\_), bh-act(\allowbreak{}bov\_), bh-vnrm(\allowbreak{}bov\_), bh-vmul(\allowbreak{}bov\_), bop-id(\allowbreak{}bov\_), bh-vstar(\allowbreak{}bov\_)]
\end{quote}
characterized by \texttt{bh-algebra-unital-c-star}  \textit{(proven)}:
\begin{quote}\ttfamily\footnotesize\raggedright\noindent
forall(\allowbreak{}[v], is-hilbert-space(\allowbreak{}v) implies forsome(\allowbreak{}[bnx\_ in vec(\allowbreak{}v)], not(\allowbreak{}bnx\_ = vzero(\allowbreak{}v))) implies is-unital-c-star-algebra(\allowbreak{}bh-algebra(\allowbreak{}v)))
\end{quote}
mentioned by 20 result(s)

\smallskip

\noindent\textbf{\texttt{bh-module(otv\_)}}\quad\textit{functoid}\index{bh-module(otv_)@\texttt{bh-module(otv\_)}}\\
reads: the vector space of bounded operators on \$1 \\
definition:
\begin{quote}\ttfamily\footnotesize\raggedright\noindent
[normed-field-as-commutative-ring(\allowbreak{}cc-normed-field), bounded-ops(\allowbreak{}otv\_, otv\_), bh-vadd(\allowbreak{}otv\_), bop-zero(\allowbreak{}otv\_, otv\_), bh-vneg(\allowbreak{}otv\_), bh-act(\allowbreak{}otv\_)]
\end{quote}
characterized by \texttt{bh-module-act}  \textit{(proven)}:
\begin{quote}\ttfamily\footnotesize\raggedright\noindent
forall(\allowbreak{}[v], is-complex-inner-product-space(\allowbreak{}v) implies act(\allowbreak{}bh-module(\allowbreak{}v)) == bh-act(\allowbreak{}v))
\end{quote}
also: \texttt{bh-module-vec} \texttt{bh-module-scal} \texttt{bh-module-vadd} \texttt{bh-module-vneg} \texttt{bh-module-vzero} \texttt{bh-module-is-module} \\
mentioned by 39 result(s)

\smallskip

\noindent\textbf{\texttt{bh-vadd(bov\_)}}\quad\textit{functoid}\index{bh-vadd(bov_)@\texttt{bh-vadd(bov\_)}}\\
definition:
\begin{quote}\ttfamily\footnotesize\raggedright\noindent
vnb-lambda(\allowbreak{}[bos\_, bot\_], cartesian(\allowbreak{}bounded-ops(\allowbreak{}bov\_, bov\_), bounded-ops(\allowbreak{}bov\_, bov\_)), bop-sum(\allowbreak{}bov\_, bov\_, bos\_, bot\_))
\end{quote}
characterized by \texttt{bh-vadd-inv}  \textit{(proven)}:
\begin{quote}\ttfamily\footnotesize\raggedright\noindent
forall(\allowbreak{}[v], is-complex-inner-product-space(\allowbreak{}v) implies has-inverses(\allowbreak{}bh-vadd(\allowbreak{}v), bop-zero(\allowbreak{}v, v), bh-vneg(\allowbreak{}v), bounded-ops(\allowbreak{}v, v)))
\end{quote}
also: \texttt{bh-vadd-comm} \texttt{bh-vadd-type} \texttt{bh-vadd-assoc} \texttt{bh-vadd-ident} \texttt{bh-vadd-value} \\
mentioned by 30 result(s)

\smallskip

\noindent\textbf{\texttt{bh-vmul(bov\_)}}\quad\textit{functoid}\index{bh-vmul(bov_)@\texttt{bh-vmul(bov\_)}}\\
definition:
\begin{quote}\ttfamily\footnotesize\raggedright\noindent
vnb-lambda(\allowbreak{}[bos\_, bot\_], cartesian(\allowbreak{}bounded-ops(\allowbreak{}bov\_, bov\_), bounded-ops(\allowbreak{}bov\_, bov\_)), bop-comp(\allowbreak{}bov\_, bos\_, bot\_))
\end{quote}
characterized by \texttt{bh-vmul-type}  \textit{(proven)}:
\begin{quote}\ttfamily\footnotesize\raggedright\noindent
forall(\allowbreak{}[v], is-complex-inner-product-space(\allowbreak{}v) implies bh-vmul(\allowbreak{}v) in fun(\allowbreak{}cartesian(\allowbreak{}bounded-ops(\allowbreak{}v, v), bounded-ops(\allowbreak{}v, v)), bounded-ops(\allowbreak{}v, v)))
\end{quote}
also: \texttt{bh-vmul-assoc} \texttt{bh-vmul-value} \texttt{bh-vmul-act-left} \texttt{bh-vmul-act-right} \texttt{bh-vmul-unit-left} \texttt{bh-vmul-vadd-left} \texttt{bh-vmul-unit-right} \texttt{bh-vmul-vadd-right} \\
mentioned by 31 result(s)

\smallskip

\noindent\textbf{\texttt{bh-vneg(bov\_)}}\quad\textit{functoid}\index{bh-vneg(bov_)@\texttt{bh-vneg(bov\_)}}\\
definition:
\begin{quote}\ttfamily\footnotesize\raggedright\noindent
vnb-lambda(\allowbreak{}bot\_, bounded-ops(\allowbreak{}bov\_, bov\_), bop-neg(\allowbreak{}bov\_, bov\_, bot\_))
\end{quote}
characterized by \texttt{bh-vneg-type}  \textit{(proven)}:
\begin{quote}\ttfamily\footnotesize\raggedright\noindent
forall(\allowbreak{}[v], is-complex-inner-product-space(\allowbreak{}v) implies bh-vneg(\allowbreak{}v) in fun(\allowbreak{}bounded-ops(\allowbreak{}v, v), bounded-ops(\allowbreak{}v, v)))
\end{quote}
also: \texttt{bh-vneg-value} \\
mentioned by 14 result(s)

\smallskip

\noindent\textbf{\texttt{bh-vnrm(bov\_)}}\quad\textit{functoid}\index{bh-vnrm(bov_)@\texttt{bh-vnrm(bov\_)}}\\
definition:
\begin{quote}\ttfamily\footnotesize\raggedright\noindent
vnb-lambda(\allowbreak{}bot\_, bounded-ops(\allowbreak{}bov\_, bov\_), opnorm(\allowbreak{}bov\_, bov\_, bot\_))
\end{quote}
characterized by \texttt{bh-vnrm-act}  \textit{(proven)}:
\begin{quote}\ttfamily\footnotesize\raggedright\noindent
forall(\allowbreak{}[v], is-complex-inner-product-space(\allowbreak{}v) implies forall(\allowbreak{}[r\_ in carr(\allowbreak{}normed-field-as-commutative-ring(\allowbreak{}cc-normed-field)), x\_ in bounded-ops(\allowbreak{}v, v)], (\allowbreak{}bh-vnrm(\allowbreak{}v))(\allowbreak{}(\allowbreak{}bh-act(\allowbreak{}v))(\allowbreak{}r\_, x\_)) = magnitude(\allowbreak{}r\_) * (\allowbreak{}bh-vnrm(\allowbreak{}v))(\allowbreak{}x\_)))
\end{quote}
also: \texttt{bh-vnrm-type} \texttt{bh-vnrm-vone} \texttt{bh-vnrm-cstar} \texttt{bh-vnrm-value} \texttt{bh-vnrm-vstar} \texttt{bh-vnrm-nonneg} \texttt{bh-vnrm-submult} \texttt{bh-vnrm-triangle} ... (9 in all) \\
mentioned by 26 result(s)

\smallskip

\noindent\textbf{\texttt{bh-vstar(bov\_)}}\quad\textit{functoid}\index{bh-vstar(bov_)@\texttt{bh-vstar(bov\_)}}\\
definition:
\begin{quote}\ttfamily\footnotesize\raggedright\noindent
vnb-lambda(\allowbreak{}bot\_, bounded-ops(\allowbreak{}bov\_, bov\_), op-adjoint(\allowbreak{}bov\_, bov\_, bot\_))
\end{quote}
characterized by \texttt{bh-vstar-act}  \textit{(proven)}:
\begin{quote}\ttfamily\footnotesize\raggedright\noindent
forall(\allowbreak{}[v], is-hilbert-space(\allowbreak{}v) implies forall(\allowbreak{}[r\_ in cc, x\_ in bounded-ops(\allowbreak{}v, v)], (\allowbreak{}bh-vstar(\allowbreak{}v))(\allowbreak{}(\allowbreak{}bh-act(\allowbreak{}v))(\allowbreak{}r\_, x\_)) = (\allowbreak{}bh-act(\allowbreak{}v))(\allowbreak{}conjugate(\allowbreak{}r\_), (\allowbreak{}bh-vstar(\allowbreak{}v))(\allowbreak{}x\_))))
\end{quote}
also: \texttt{bh-vstar-type} \texttt{bh-vstar-vadd} \texttt{bh-vstar-vmul} \texttt{bh-vstar-value} \texttt{bh-vstar-vstar} \\
mentioned by 19 result(s)

\smallskip

\noindent\textbf{\texttt{bichar-of(chk\_, chc\_, cha\_)}}\quad\textit{functoid}\index{bichar-of(chk_, chc_, cha_)@\texttt{bichar-of(chk\_, chc\_, cha\_)}}\\
reads: the bicharacter (x, y) |-> \$2(\$3 x y) of \$1 \\
definition:
\begin{quote}\ttfamily\footnotesize\raggedright\noindent
vnb-lambda(\allowbreak{}[chu\_, chv\_], cartesian(\allowbreak{}carr(\allowbreak{}chk\_), carr(\allowbreak{}chk\_)), chc\_(\allowbreak{}(\allowbreak{}mul(\allowbreak{}chk\_))(\allowbreak{}cha\_, (\allowbreak{}mul(\allowbreak{}chk\_))(\allowbreak{}chu\_, chv\_))))
\end{quote}
mentioned by 13 result(s)

\smallskip

\noindent\textbf{\texttt{bichar-polar(c5k\_, c5w\_, c5v\_)}}\quad\textit{functoid}\index{bichar-polar(c5k_, c5w_, c5v_)@\texttt{bichar-polar(c5k\_, c5w\_, c5v\_)}}\\
reads: the polar of \$3 with respect to \$2 \\
definition:
\begin{quote}\ttfamily\footnotesize\raggedright\noindent
\{c5h\_ in carr(\allowbreak{}c5k\_): forall(\allowbreak{}[c5x\_ in c5v\_], c5w\_(\allowbreak{}c5h\_, c5x\_) = 1)\}
\end{quote}
mentioned by 6 result(s)

\smallskip

\noindent\textbf{\texttt{bicommutant(vnv\_, vns\_)}}\quad\textit{functoid}\index{bicommutant(vnv_, vns_)@\texttt{bicommutant(vnv\_, vns\_)}}\\
reads: the bicommutant of \$2 in the bounded operators on \$1 \\
definition:
\begin{quote}\ttfamily\footnotesize\raggedright\noindent
commutant(\allowbreak{}vnv\_, commutant(\allowbreak{}vnv\_, vns\_))
\end{quote}
characterized by \texttt{bicommutant-unfold}  \textit{(proven)}:
\begin{quote}\ttfamily\footnotesize\raggedright\noindent
forall(\allowbreak{}[v, s], bicommutant(\allowbreak{}v, s) == commutant(\allowbreak{}v, commutant(\allowbreak{}v, s)))
\end{quote}
also: \texttt{subset-bicommutant} \texttt{bicommutant-approx-at} \texttt{rep-image-bicommutant} \texttt{bicommutant-von-neumann} \texttt{bicommutant-finite-approx} \texttt{bicommutant-in-sot-closure} \texttt{bicommutant-is-sot-closure} \texttt{bicommutant-vector-closure} ... (11 in all) \\
mentioned by 19 result(s)

\smallskip

\noindent\textbf{\texttt{bidual-eval(m, x)}}\quad\textit{functoid}\index{bidual-eval(m, x)@\texttt{bidual-eval(m, x)}}\\
reads: the evaluation at \$2 of the bounded linear functionals on \$1 \\
definition:
\begin{quote}\ttfamily\footnotesize\raggedright\noindent
vnb-lambda(\allowbreak{}bdef\_, dual-vec(\allowbreak{}m), bdef\_(\allowbreak{}x))
\end{quote}
characterized by \texttt{bidual-eval-fun}  \textit{(proven)}:
\begin{quote}\ttfamily\footnotesize\raggedright\noindent
forall(\allowbreak{}[m, x], is-normed-vector-space(\allowbreak{}m) implies x in vec(\allowbreak{}m) implies bidual-eval(\allowbreak{}m, x) in fun(\allowbreak{}dual-vec(\allowbreak{}m), rr))
\end{quote}
also: \texttt{bidual-eval-apply} \texttt{bidual-eval-bound} \texttt{bidual-eval-linear} \texttt{bidual-eval-additive} \texttt{bidual-eval-isometry} \texttt{bidual-eval-in-bidual} \texttt{bidual-eval-homogeneous} \\
mentioned by 12 result(s)

\smallskip

\noindent\textbf{\texttt{big-union}}\quad\textit{operator}\index{big-union@\texttt{big-union}}\\
a term-forming head \\
characterized by \texttt{big-union-mono}  \textit{(proven)}:
\begin{quote}\ttfamily\footnotesize\raggedright\noindent
forall(\allowbreak{}[a\_, b\_], a\_ subset b\_ implies big-union(\allowbreak{}u\_, a\_, u\_) subset big-union(\allowbreak{}u\_, b\_, u\_))
\end{quote}
also: \texttt{power-big-union-closed} \texttt{complement-in-big-union} \texttt{big-union-of-empty-family} \texttt{big-union-of-union-family} \texttt{big-union-of-singleton-family} \\
mentioned by 221 result(s)

\smallskip

\noindent\textbf{\texttt{bijection}}\quad\textit{operator}\index{bijection@\texttt{bijection}}\\
a term-forming head \\
characterized by \texttt{bijection-membership-iff}  \textit{(definitional)}:
\begin{quote}\ttfamily\footnotesize\raggedright\noindent
forall(\allowbreak{}[x, y, phi], phi in bijection(\allowbreak{}x, y) iff phi in fun(\allowbreak{}x, y) and forall(\allowbreak{}[a in x, b in x], phi(\allowbreak{}a) = phi(\allowbreak{}b) implies a = b) and forall(\allowbreak{}[w in y], forsome(\allowbreak{}[z in x], phi(\allowbreak{}z) = w)))
\end{quote}
also: \texttt{lft-bijection} \texttt{swap-bijection} \texttt{bijection-in-fun} \texttt{lzt-lc-bijection} \texttt{bijection-compose} \texttt{bijection-set-iff} \texttt{card-bijection-eq} \texttt{zermelo-bijection} ... (35 in all) \\
mentioned by 117 result(s)

\smallskip

\noindent\textbf{\texttt{bind-sig(mch\_, mcnb\_, mcn\_)}}\quad\textit{predicate}\index{bind-sig(mch_, mcnb_, mcn_)@\texttt{bind-sig(mch\_, mcnb\_, mcn\_)}}\\
reads: the head \$1 binds \$2 variables over \$3 outside arguments \\
defined by \texttt{bind-sig} \\
characterized by \texttt{bind-sig}  \textit{(definitional)}:
\begin{quote}\ttfamily\footnotesize\raggedright\noindent
forall(\allowbreak{}[mch\_, mcnb\_, mcn\_], bind-sig(\allowbreak{}mch\_, mcnb\_, mcn\_) iff 20 <= mch\_ and mch\_ <= 21 and mcnb\_ = 1 and mcn\_ = 0 or (\allowbreak{}mch\_ = 22 or mch\_ = 24 or mch\_ = 26) and mcnb\_ = 1 and mcn\_ = 1 or (\allowbreak{}mch\_ = 23 or mch\_ = 25) and mcnb\_ = 1 and mcn\_ = 0 or mch\_ = 27 and mcnb\_ = 0 and mcn\_ = 2)
\end{quote}
mentioned by 6 result(s)

\smallskip

\noindent\textbf{\texttt{binneg}}\quad\textit{operator}\index{binneg@\texttt{binneg}}\\
a term-forming head \\
characterized by \texttt{binneg-apply}  \textit{(definitional)}:
\begin{quote}\ttfamily\footnotesize\raggedright\noindent
forall(\allowbreak{}[x], binneg(\allowbreak{}x) == -x)
\end{quote}
mentioned by 6 result(s)

\smallskip

\noindent\textbf{\texttt{binplus}}\quad\textit{operator}\index{binplus@\texttt{binplus}}\\
a term-forming head \\
characterized by \texttt{binplus-apply}  \textit{(definitional)}:
\begin{quote}\ttfamily\footnotesize\raggedright\noindent
forall(\allowbreak{}[x, y], binplus(\allowbreak{}x, y) == x + y)
\end{quote}
mentioned by 22 result(s)

\smallskip

\noindent\textbf{\texttt{bintimes}}\quad\textit{operator}\index{bintimes@\texttt{bintimes}}\\
a term-forming head \\
characterized by \texttt{bintimes-apply}  \textit{(definitional)}:
\begin{quote}\ttfamily\footnotesize\raggedright\noindent
forall(\allowbreak{}[x, y], bintimes(\allowbreak{}x, y) == x * y)
\end{quote}
mentioned by 20 result(s)

\smallskip

\noindent\textbf{\texttt{bk-f(bfg\_, bfv\_, bfx\_)}}\quad\textit{functoid}\index{bk-f(bfg_, bfv_, bfx_)@\texttt{bk-f(bfg\_, bfv\_, bfx\_)}}\\
reads: the regularised gauge of \$3 for \$1 and \$2 \\
definition:
\begin{quote}\ttfamily\footnotesize\raggedright\noindent
subadd-reg(\allowbreak{}grp-monoid(\allowbreak{}bfg\_), vnb-lambda(\allowbreak{}bfz\_, carr(\allowbreak{}bfg\_), bk-gauge(\allowbreak{}bfv\_, bfz\_)), bfx\_)
\end{quote}
characterized by \texttt{bk-f-iden}  \textit{(proven)}:
\begin{quote}\ttfamily\footnotesize\raggedright\noindent
forall(\allowbreak{}[s, g, vs], is-group-topology(\allowbreak{}s, g) implies is-bk-sequence(\allowbreak{}s, g, vs) implies bk-f(\allowbreak{}g, vs, iden(\allowbreak{}g)) = 0)
\end{quote}
also: \texttt{bk-f-bounds} \texttt{bk-f-unfold} \texttt{bk-f-subadditive} \\
mentioned by 11 result(s)

\smallskip

\noindent\textbf{\texttt{bk-gauge(bgv\_, bgx\_)}}\quad\textit{functoid}\index{bk-gauge(bgv_, bgx_)@\texttt{bk-gauge(bgv\_, bgx\_)}}\\
reads: the gauge of \$2 relative to the neighbourhoods \$1 \\
definition:
\begin{quote}\ttfamily\footnotesize\raggedright\noindent
einf(\allowbreak{}\{bgr\_ in rr: forsome(\allowbreak{}[bgk\_ in nn], bgx\_ in bgv\_(\allowbreak{}bgk\_) and bgr\_ = num(\allowbreak{}1/2) \textasciicircum{} bgk\_)\})
\end{quote}
characterized by \texttt{bk-gauge-le}  \textit{(proven)}:
\begin{quote}\ttfamily\footnotesize\raggedright\noindent
forall(\allowbreak{}[vs, x, k in nn], x in vs(\allowbreak{}k) implies bk-gauge(\allowbreak{}vs, x) in rr and 0 <= bk-gauge(\allowbreak{}vs, x) and bk-gauge(\allowbreak{}vs, x) <= num(\allowbreak{}1/2) \textasciicircum{} k)
\end{quote}
also: \texttt{bk-gauge-char} \texttt{bk-gauge-iden} \texttt{bk-gauge-exact} \texttt{bk-gauge-power} \texttt{bk-gauge-triple} \texttt{bk-gauge-unfold} \\
mentioned by 13 result(s)

\smallskip

\noindent\textbf{\texttt{bk-metric(bmg\_, bmv\_)}}\quad\textit{functoid}\index{bk-metric(bmg_, bmv_)@\texttt{bk-metric(bmg\_, bmv\_)}}\\
reads: the Birkhoff-Kakutani metric of \$1 for \$2 \\
definition:
\begin{quote}\ttfamily\footnotesize\raggedright\noindent
[carr(\allowbreak{}bmg\_), vnb-lambda(\allowbreak{}[bmx\_, bmy\_], cartesian(\allowbreak{}carr(\allowbreak{}bmg\_), carr(\allowbreak{}bmg\_)), bk-f(\allowbreak{}bmg\_, bmv\_, (\allowbreak{}opr(\allowbreak{}bmg\_))(\allowbreak{}(\allowbreak{}inv(\allowbreak{}bmg\_))(\allowbreak{}bmy\_), bmx\_)) + bk-f(\allowbreak{}bmg\_, bmv\_, (\allowbreak{}opr(\allowbreak{}bmg\_))(\allowbreak{}(\allowbreak{}inv(\allowbreak{}bmg\_))(\allowbreak{}bmx\_), bmy\_)))]
\end{quote}
characterized by \texttt{bk-metric-pts}  \textit{(proven)}:
\begin{quote}\ttfamily\footnotesize\raggedright\noindent
forall(\allowbreak{}[s, g, vs], is-group-topology(\allowbreak{}s, g) implies is-bk-sequence(\allowbreak{}s, g, vs) implies pts(\allowbreak{}bk-metric(\allowbreak{}g, vs)) == carr(\allowbreak{}g))
\end{quote}
also: \texttt{bk-metric-dist} \texttt{bk-metric-self} \texttt{bk-metric-dist-fun} \texttt{bk-metric-topology} \texttt{bk-metric-pseudometric} \texttt{bk-metric-left-invariant} \texttt{bk-metric-is-metric-space} \\
mentioned by 21 result(s)

\smallskip

\noindent\textbf{\texttt{block(p, k, l)}}\quad\textit{functoid}\index{block(p, k, l)@\texttt{block(p, k, l)}}\\
definition:
\begin{quote}\ttfamily\footnotesize\raggedright\noindent
matof(\allowbreak{}k, l, vnb-lambda(\allowbreak{}[i, j], cartesian(\allowbreak{}interval(\allowbreak{}1, k), interval(\allowbreak{}1, l)), entry(\allowbreak{}p, i, j)))
\end{quote}
characterized by \texttt{block-type}  \textit{(proven)}:
\begin{quote}\ttfamily\footnotesize\raggedright\noindent
forall(\allowbreak{}[m, n, x, p, k, l], k in nn implies l in nn implies p in mat(\allowbreak{}m, n, x) implies k <= m implies l <= n implies block(\allowbreak{}p, k, l) in mat(\allowbreak{}k, l, x))
\end{quote}
also: \texttt{map-seq-block} \texttt{entry-of-block} \\
mentioned by 60 result(s)

\smallskip

\noindent\textbf{\texttt{bnodes(s)}}\quad\textit{accessor}\index{bnodes(s)@\texttt{bnodes(s)}}\\
reads: the bnodes slot of \$1 \\
a structure slot \\
characterized by \texttt{cx-c@bnodes}  \textit{(definitional)}:
\begin{quote}\ttfamily\footnotesize\raggedright\noindent
bnodes(\allowbreak{}cx-c) == nodes(\allowbreak{}cx-sp)
\end{quote}
mentioned by 70 result(s)

\smallskip

\noindent\textbf{\texttt{bop-comp(bov\_, bos\_, bot\_)}}\quad\textit{functoid}\index{bop-comp(bov_, bos_, bot_)@\texttt{bop-comp(bov\_, bos\_, bot\_)}}\\
reads: \$2 \$3 \\
definition:
\begin{quote}\ttfamily\footnotesize\raggedright\noindent
vnb-lambda(\allowbreak{}box\_, vec(\allowbreak{}bov\_), bos\_(\allowbreak{}bot\_(\allowbreak{}box\_)))
\end{quote}
characterized by \texttt{bop-comp-apply}  \textit{(proven)}:
\begin{quote}\ttfamily\footnotesize\raggedright\noindent
forall(\allowbreak{}[v, s\_, t\_, x\_], s\_ in fun(\allowbreak{}vec(\allowbreak{}v), vec(\allowbreak{}v)) implies t\_ in fun(\allowbreak{}vec(\allowbreak{}v), vec(\allowbreak{}v)) implies x\_ in vec(\allowbreak{}v) implies (\allowbreak{}bop-comp(\allowbreak{}v, s\_, t\_))(\allowbreak{}x\_) = s\_(\allowbreak{}t\_(\allowbreak{}x\_)))
\end{quote}
also: \texttt{bop-comp-assoc} \texttt{bop-comp-bounded} \texttt{bop-comp-id-left} \texttt{bop-comp-id-right} \texttt{bop-comp-sum-left} \texttt{bop-comp-apply-gen} \texttt{bop-comp-sum-right} \texttt{bop-comp-scale-left} ... (10 in all) \\
mentioned by 94 result(s)

\smallskip

\noindent\textbf{\texttt{bop-id(bov\_)}}\quad\textit{functoid}\index{bop-id(bov_)@\texttt{bop-id(bov\_)}}\\
reads: the identity operator of \$1 \\
definition:
\begin{quote}\ttfamily\footnotesize\raggedright\noindent
vnb-lambda(\allowbreak{}box\_, vec(\allowbreak{}bov\_), box\_)
\end{quote}
characterized by \texttt{bop-id-apply}  \textit{(proven)}:
\begin{quote}\ttfamily\footnotesize\raggedright\noindent
forall(\allowbreak{}[v, x\_ in vec(\allowbreak{}v)], (\allowbreak{}bop-id(\allowbreak{}v))(\allowbreak{}x\_) = x\_)
\end{quote}
also: \texttt{bop-id-bounded} \\
mentioned by 28 result(s)

\smallskip

\noindent\textbf{\texttt{bop-neg(bov\_, bow\_, bot\_)}}\quad\textit{functoid}\index{bop-neg(bov_, bow_, bot_)@\texttt{bop-neg(bov\_, bow\_, bot\_)}}\\
reads: - \$3 \\
definition:
\begin{quote}\ttfamily\footnotesize\raggedright\noindent
vnb-lambda(\allowbreak{}box\_, vec(\allowbreak{}bov\_), (\allowbreak{}vneg(\allowbreak{}bow\_))(\allowbreak{}bot\_(\allowbreak{}box\_)))
\end{quote}
characterized by \texttt{bop-neg-apply}  \textit{(proven)}:
\begin{quote}\ttfamily\footnotesize\raggedright\noindent
forall(\allowbreak{}[v, w, t\_, x\_], is-complex-inner-product-space(\allowbreak{}w) implies t\_ in fun(\allowbreak{}vec(\allowbreak{}v), vec(\allowbreak{}w)) implies x\_ in vec(\allowbreak{}v) implies (\allowbreak{}bop-neg(\allowbreak{}v, w, t\_))(\allowbreak{}x\_) = (\allowbreak{}vneg(\allowbreak{}w))(\allowbreak{}t\_(\allowbreak{}x\_)))
\end{quote}
also: \texttt{bop-neg-unfold} \texttt{bop-neg-bounded} \\
mentioned by 21 result(s)

\smallskip

\noindent\textbf{\texttt{bop-normed-ag(bov\_, bow\_)}}\quad\textit{functoid}\index{bop-normed-ag(bov_, bow_)@\texttt{bop-normed-ag(bov\_, bow\_)}}\\
reads: the normed group of bounded operators from \$1 to \$2 \\
definition:
\begin{quote}\ttfamily\footnotesize\raggedright\noindent
[bounded-ops(\allowbreak{}bov\_, bow\_), bop-vadd(\allowbreak{}bov\_, bow\_), bop-zero(\allowbreak{}bov\_, bow\_), bop-vinv(\allowbreak{}bov\_, bow\_), bop-vnrm(\allowbreak{}bov\_, bow\_)]
\end{quote}
characterized by \texttt{bop-normed-ag-cnvs}  \textit{(proven)}:
\begin{quote}\ttfamily\footnotesize\raggedright\noindent
forall(\allowbreak{}[v, w], is-complex-inner-product-space(\allowbreak{}v) implies is-complex-inner-product-space(\allowbreak{}w) implies bop-normed-ag(\allowbreak{}v, w) == cbop-normed-ag(\allowbreak{}cips-cnvs(\allowbreak{}v), cips-cnvs(\allowbreak{}w)))
\end{quote}
also: \texttt{bop-normed-ag-is-normed-ag} \\
mentioned by 37 result(s)

\smallskip

\noindent\textbf{\texttt{bop-scale(bov\_, bow\_, boc\_, bot\_)}}\quad\textit{functoid}\index{bop-scale(bov_, bow_, boc_, bot_)@\texttt{bop-scale(bov\_, bow\_, boc\_, bot\_)}}\\
reads: \$3 \$4 \\
definition:
\begin{quote}\ttfamily\footnotesize\raggedright\noindent
vnb-lambda(\allowbreak{}box\_, vec(\allowbreak{}bov\_), (\allowbreak{}act(\allowbreak{}bow\_))(\allowbreak{}boc\_, bot\_(\allowbreak{}box\_)))
\end{quote}
characterized by \texttt{bop-scale-add}  \textit{(proven)}:
\begin{quote}\ttfamily\footnotesize\raggedright\noindent
forall(\allowbreak{}[v, w, c\_, d\_, t\_], is-complex-inner-product-space(\allowbreak{}v) implies is-complex-inner-product-space(\allowbreak{}w) implies c\_ in cc implies d\_ in cc implies t\_ in bounded-ops(\allowbreak{}v, w) implies bop-scale(\allowbreak{}v, w, c\_ + d\_, t\_) = bop-sum(\allowbreak{}v, w, bop-scale(\allowbreak{}v, w, c\_, t\_), bop-scale(\allowbreak{}v, w, d\_, t\_)))
\end{quote}
also: \texttt{bop-scale-mul} \texttt{bop-scale-one} \texttt{bop-scale-sum} \texttt{bop-scale-apply} \texttt{bop-scale-bounded} \\
mentioned by 35 result(s)

\smallskip

\noindent\textbf{\texttt{bop-sum(bov\_, bow\_, bos\_, bot\_)}}\quad\textit{functoid}\index{bop-sum(bov_, bow_, bos_, bot_)@\texttt{bop-sum(bov\_, bow\_, bos\_, bot\_)}}\\
reads: \$3 + \$4 \\
definition:
\begin{quote}\ttfamily\footnotesize\raggedright\noindent
vnb-lambda(\allowbreak{}box\_, vec(\allowbreak{}bov\_), (\allowbreak{}vadd(\allowbreak{}bow\_))(\allowbreak{}bos\_(\allowbreak{}box\_), bot\_(\allowbreak{}box\_)))
\end{quote}
characterized by \texttt{bop-sum-comm}  \textit{(proven)}:
\begin{quote}\ttfamily\footnotesize\raggedright\noindent
forall(\allowbreak{}[v, w, s\_, t\_], is-complex-inner-product-space(\allowbreak{}v) implies is-complex-inner-product-space(\allowbreak{}w) implies s\_ in bounded-ops(\allowbreak{}v, w) implies t\_ in bounded-ops(\allowbreak{}v, w) implies bop-sum(\allowbreak{}v, w, s\_, t\_) = bop-sum(\allowbreak{}v, w, t\_, s\_))
\end{quote}
also: \texttt{bop-sum-apply} \texttt{bop-sum-assoc} \texttt{bop-sum-unfold} \texttt{bop-sum-bounded} \texttt{bop-sum-neg-left} \texttt{bop-sum-neg-right} \texttt{bop-sum-zero-left} \texttt{bop-sum-zero-right} \\
mentioned by 51 result(s)

\smallskip

\noindent\textbf{\texttt{bop-vadd(bov\_, bow\_)}}\quad\textit{functoid}\index{bop-vadd(bov_, bow_)@\texttt{bop-vadd(bov\_, bow\_)}}\\
definition:
\begin{quote}\ttfamily\footnotesize\raggedright\noindent
vnb-lambda(\allowbreak{}[bos\_, bot\_], cartesian(\allowbreak{}bounded-ops(\allowbreak{}bov\_, bow\_), bounded-ops(\allowbreak{}bov\_, bow\_)), bop-sum(\allowbreak{}bov\_, bow\_, bos\_, bot\_))
\end{quote}
characterized by \texttt{bop-vadd-cnvs}  \textit{(proven)}:
\begin{quote}\ttfamily\footnotesize\raggedright\noindent
forall(\allowbreak{}[v, w], is-complex-inner-product-space(\allowbreak{}v) implies is-complex-inner-product-space(\allowbreak{}w) implies bop-vadd(\allowbreak{}v, w) == cbop-vadd(\allowbreak{}cips-cnvs(\allowbreak{}v), cips-cnvs(\allowbreak{}w)))
\end{quote}
also: \texttt{bop-vadd-value} \\
mentioned by 6 result(s)

\smallskip

\noindent\textbf{\texttt{bop-vinv(bov\_, bow\_)}}\quad\textit{functoid}\index{bop-vinv(bov_, bow_)@\texttt{bop-vinv(bov\_, bow\_)}}\\
definition:
\begin{quote}\ttfamily\footnotesize\raggedright\noindent
vnb-lambda(\allowbreak{}bot\_, bounded-ops(\allowbreak{}bov\_, bow\_), bop-neg(\allowbreak{}bov\_, bow\_, bot\_))
\end{quote}
characterized by \texttt{bop-vinv-cnvs}  \textit{(proven)}:
\begin{quote}\ttfamily\footnotesize\raggedright\noindent
forall(\allowbreak{}[v, w], is-complex-inner-product-space(\allowbreak{}v) implies is-complex-inner-product-space(\allowbreak{}w) implies bop-vinv(\allowbreak{}v, w) == cbop-vinv(\allowbreak{}cips-cnvs(\allowbreak{}v), cips-cnvs(\allowbreak{}w)))
\end{quote}
also: \texttt{bop-vinv-value} \\
mentioned by 6 result(s)

\smallskip

\noindent\textbf{\texttt{bop-vnrm(bov\_, bow\_)}}\quad\textit{functoid}\index{bop-vnrm(bov_, bow_)@\texttt{bop-vnrm(bov\_, bow\_)}}\\
definition:
\begin{quote}\ttfamily\footnotesize\raggedright\noindent
vnb-lambda(\allowbreak{}bot\_, bounded-ops(\allowbreak{}bov\_, bow\_), opnorm(\allowbreak{}bov\_, bow\_, bot\_))
\end{quote}
characterized by \texttt{bop-vnrm-cnvs}  \textit{(proven)}:
\begin{quote}\ttfamily\footnotesize\raggedright\noindent
forall(\allowbreak{}[v, w], is-complex-inner-product-space(\allowbreak{}v) implies is-complex-inner-product-space(\allowbreak{}w) implies bop-vnrm(\allowbreak{}v, w) = cbop-vnrm(\allowbreak{}cips-cnvs(\allowbreak{}v), cips-cnvs(\allowbreak{}w)))
\end{quote}
also: \texttt{bop-vnrm-value} \\
mentioned by 6 result(s)

\smallskip

\noindent\textbf{\texttt{bop-zero(bov\_, bow\_)}}\quad\textit{functoid}\index{bop-zero(bov_, bow_)@\texttt{bop-zero(bov\_, bow\_)}}\\
reads: the zero operator from \$1 to \$2 \\
definition:
\begin{quote}\ttfamily\footnotesize\raggedright\noindent
vnb-lambda(\allowbreak{}box\_, vec(\allowbreak{}bov\_), vzero(\allowbreak{}bow\_))
\end{quote}
characterized by \texttt{bop-zero-cnvs}  \textit{(proven)}:
\begin{quote}\ttfamily\footnotesize\raggedright\noindent
forall(\allowbreak{}[v, w], is-complex-inner-product-space(\allowbreak{}v) implies is-complex-inner-product-space(\allowbreak{}w) implies bop-zero(\allowbreak{}v, w) == cbop-zero(\allowbreak{}cips-cnvs(\allowbreak{}v), cips-cnvs(\allowbreak{}w)))
\end{quote}
also: \texttt{bop-zero-apply} \texttt{bop-zero-bounded} \texttt{bop-zero-compact} \\
mentioned by 45 result(s)

\smallskip

\noindent\textbf{\texttt{border(a, b, m, p, q)}}\quad\textit{functoid}\index{border(a, b, m, p, q)@\texttt{border(a, b, m, p, q)}}\\
definition:
\begin{quote}\ttfamily\footnotesize\raggedright\noindent
matof(\allowbreak{}succ(\allowbreak{}p), succ(\allowbreak{}q), vnb-lambda(\allowbreak{}[i, j], cartesian(\allowbreak{}interval(\allowbreak{}1, succ(\allowbreak{}p)), interval(\allowbreak{}1, succ(\allowbreak{}q))), if(\allowbreak{}i = 1, if(\allowbreak{}j = 1, b, zero(\allowbreak{}a)), if(\allowbreak{}j = 1, zero(\allowbreak{}a), entry(\allowbreak{}m, nn-minus(\allowbreak{}i, 1), nn-minus(\allowbreak{}j, 1))))))
\end{quote}
characterized by \texttt{bordering-uba-as-ring}  \textit{(definitional)}:
\begin{quote}\ttfamily\footnotesize\raggedright\noindent
forall(\allowbreak{}[r], is-unital-banach-algebra(\allowbreak{}r) implies forall(\allowbreak{}[p, q, b, w, g, e], p in nn implies q in nn implies b in vec(\allowbreak{}r) implies w in mat(\allowbreak{}p, q, vec(\allowbreak{}r)) implies g in mat(\allowbreak{}succ(\allowbreak{}p), succ(\allowbreak{}q), vec(\allowbreak{}r)) implies g = border(\allowbreak{}uba-as-ring(\allowbreak{}r), b, w, p, q) implies mat-equiv(\allowbreak{}uba-as-ring(\allowbreak{}r), p, q, w, e) implies mat-equiv(\allowbreak{}uba-as-ring(\allowbreak{}r), succ(\allowbreak{}p), succ(\allowbreak{}q), g, border(\allowbreak{}uba-as-ring(\allowbreak{}r), b, e, p, q))))
\end{quote}
also: \texttt{border-mult-uba-as-ring} \texttt{border-identity-uba-as-ring} \texttt{border-scale-row-uba-as-ring} \texttt{border-staircase-uba-as-ring} \texttt{border-invertible-uba-as-ring} \texttt{border-is-diagonal-uba-as-ring} \texttt{bordered-eq-border-uba-as-ring} \texttt{bordering-uba-as-ring-uiba-as-uba} ... (31 in all) \\
mentioned by 90 result(s)

\smallskip

\noindent\textbf{\texttt{borel-separated(s, a, b)}}\quad\textit{predicate}\index{borel-separated(s, a, b)@\texttt{borel-separated(s, a, b)}}\\
reads: \$2 and \$3 are separated by a Borel set of \$1 \\
defined by \texttt{borel-separated} \\
characterized by \texttt{borel-separated}  \textit{(definitional)}:
\begin{quote}\ttfamily\footnotesize\raggedright\noindent
forall(\allowbreak{}[s, a, b], borel-separated(\allowbreak{}s, a, b) iff forsome(\allowbreak{}[bsc\_ in borel-sets(\allowbreak{}s)], a subset bsc\_ and forall(\allowbreak{}[bsx\_ in b], not(\allowbreak{}bsx\_ in bsc\_))))
\end{quote}
also: \texttt{borel-separated-mono} \texttt{borel-separated-union} \\
mentioned by 9 result(s)

\smallskip

\noindent\textbf{\texttt{borel-sets(m)}}\quad\textit{functoid}\index{borel-sets(m)@\texttt{borel-sets(m)}}\\
reads: the Borel subsets of the metric space \$1 \\
definition:
\begin{quote}\ttfamily\footnotesize\raggedright\noindent
sigma-generated(\allowbreak{}pts(\allowbreak{}m), \{mou\_ in power(\allowbreak{}pts(\allowbreak{}m)): is-open(\allowbreak{}m, mou\_)\})
\end{quote}
characterized by \texttt{borel-sets-unfold}  \textit{(proven)}:
\begin{quote}\ttfamily\footnotesize\raggedright\noindent
forall(\allowbreak{}[m], borel-sets(\allowbreak{}m) == sigma-generated(\allowbreak{}pts(\allowbreak{}m), \{mou\_ in power(\allowbreak{}pts(\allowbreak{}m)): is-open(\allowbreak{}m, mou\_)\}))
\end{quote}
also: \texttt{borel-sets-product} \texttt{borel-sets-subspace} \texttt{borel-sets-caratheodory-measurable} \\
mentioned by 112 result(s)

\smallskip

\noindent\textbf{\texttt{bounded-ops(bov\_, bow\_)}}\quad\textit{functoid}\index{bounded-ops(bov_, bow_)@\texttt{bounded-ops(bov\_, bow\_)}}\\
reads: the bounded operators from \$1 to \$2 \\
definition:
\begin{quote}\ttfamily\footnotesize\raggedright\noindent
\{bot\_ in fun(\allowbreak{}vec(\allowbreak{}bov\_), vec(\allowbreak{}bow\_)): is-bounded-op(\allowbreak{}bov\_, bow\_, bot\_)\}
\end{quote}
characterized by \texttt{bounded-ops-is-set}  \textit{(proven)}:
\begin{quote}\ttfamily\footnotesize\raggedright\noindent
forall(\allowbreak{}[v, w], is-complex-inner-product-space(\allowbreak{}v) implies is-complex-inner-product-space(\allowbreak{}w) implies bounded-ops(\allowbreak{}v, w) in set)
\end{quote}
also: \texttt{bounded-ops-unfold} \texttt{bounded-ops-membership} \\
mentioned by 399 result(s)

\smallskip

\noindent\textbf{\texttt{bsend(s)}}\quad\textit{accessor}\index{bsend(s)@\texttt{bsend(s)}}\\
reads: the bsend slot of \$1 \\
a structure slot \\
characterized by \texttt{cx-c@bsend}  \textit{(definitional)}:
\begin{quote}\ttfamily\footnotesize\raggedright\noindent
bsend(\allowbreak{}cx-c) == \{cxe\_ in cartesian(\allowbreak{}nodes(\allowbreak{}cx-sp), nodes(\allowbreak{}cx-sp)): cxe\_ = [[1, 1], [1, 2]] or cxe\_ = [[2, 2], [1, 5]] or cxe\_ = [[2, 5], [1, 3]] or cxe\_ = [[3, 5], [1, 4]] or cxe\_ = [[2, 3], [3, 2]] or cxe\_ = [[4, 2], [1, 6]] or cxe\_ = [[2, 4], [1, 7]] or cxe\_ = [[2, 7], [1, 8]] or cxe\_ = [[3, 7], [2, 8]] or cxe\_ = [[3, 8], [2, 1]]\}
\end{quote}
mentioned by 16 result(s)

\smallskip

\noindent\textbf{\texttt{bspace(s)}}\quad\textit{accessor}\index{bspace(s)@\texttt{bspace(s)}}\\
reads: the bspace slot of \$1 \\
a structure slot \\
characterized by \texttt{cx-c@bspace}  \textit{(definitional)}:
\begin{quote}\ttfamily\footnotesize\raggedright\noindent
bspace(\allowbreak{}cx-c) == cx-sp
\end{quote}
mentioned by 59 result(s)

\smallskip

\noindent\textbf{\texttt{bu-me}}\quad\textit{tactic}\index{bu-me@\texttt{bu-me}}\\
\texttt{(bu-me hyp)} -- Big-union membership elim. \\
declared in \texttt{tactics-help.scm}

\smallskip

\noindent\textbf{\texttt{bu-mi}}\quad\textit{tactic}\index{bu-mi@\texttt{bu-mi}}\\
\texttt{(bu-mi w)} -- Big-union membership intro via the index witness w. \\
declared in \texttt{tactics-help.scm}

\smallskip

\noindent\textbf{\texttt{bu-set}}\quad\textit{tactic}\index{bu-set@\texttt{bu-set}}\\
\texttt{(bu-set)} -- Big-union sethood. \\
declared in \texttt{tactics-help.scm}

\smallskip

\noindent\textbf{\texttt{bundle}}\quad\textit{structure}\index{bundle@\texttt{bundle}}\\
predicate \texttt{is-bundle}; slots: bspace bnodes bsend

\smallskip


\section*{C}
\addcontentsline{toc}{section}{C}

\noindent\textbf{\texttt{c-int(phi, a, b)}}\quad\textit{functoid}\index{c-int(phi, a, b)@\texttt{c-int(phi, a, b)}}\\
reads: the integral of \$1 from \$2 to \$3 \\
definition:
\begin{quote}\ttfamily\footnotesize\raggedright\noindent
iota(\allowbreak{}v\_, v\_ in rr and forsome(\allowbreak{}[f\_], is-antiderivative(\allowbreak{}f\_, phi, a, b) and v\_ = f\_(\allowbreak{}b) - f\_(\allowbreak{}a)))
\end{quote}
characterized by \texttt{c-int-add}  \textit{(proven)}:
\begin{quote}\ttfamily\footnotesize\raggedright\noindent
forall(\allowbreak{}[phi, psi, a, b], is-antiderivable(\allowbreak{}phi, a, b) implies is-antiderivable(\allowbreak{}psi, a, b) implies c-int(\allowbreak{}vnb-lambda(\allowbreak{}x, rr, phi(\allowbreak{}x) + psi(\allowbreak{}x)), a, b) = c-int(\allowbreak{}phi, a, b) + c-int(\allowbreak{}psi, a, b))
\end{quote}
also: \texttt{c-int-in-rr} \texttt{c-int-scale} \texttt{c-int-value} \texttt{c-int-or-anti} \texttt{c-int-or-flip} \texttt{c-int-add-bounds} \texttt{c-int-add-continuous} \texttt{c-int-scale-continuous} ... (12 in all) \\
mentioned by 23 result(s)

\smallskip

\noindent\textbf{\texttt{c-int-or(phi, a, b)}}\quad\textit{functoid}\index{c-int-or(phi, a, b)@\texttt{c-int-or(phi, a, b)}}\\
reads: the oriented integral of \$1 from \$2 to \$3 \\
definition:
\begin{quote}\ttfamily\footnotesize\raggedright\noindent
if(\allowbreak{}is-antiderivable(\allowbreak{}phi, a, b), c-int(\allowbreak{}phi, a, b), if(\allowbreak{}is-antiderivable(\allowbreak{}phi, b, a), -c-int(\allowbreak{}phi, b, a), 0))
\end{quote}
characterized by \texttt{c-int-or-anti}  \textit{(proven)}:
\begin{quote}\ttfamily\footnotesize\raggedright\noindent
forall(\allowbreak{}[phi, a, b], is-antiderivable(\allowbreak{}phi, a, b) implies c-int-or(\allowbreak{}phi, a, b) = c-int(\allowbreak{}phi, a, b))
\end{quote}
also: \texttt{c-int-or-flip} \texttt{c-int-or-none} \texttt{c-int-or-in-rr} \texttt{c-int-or-scale} \texttt{c-int-or-value} \texttt{c-int-or-additive} \texttt{c-int-or-degenerate} \texttt{c-int-or-scale-ordered} \\
mentioned by 25 result(s)

\smallskip

\noindent\textbf{\texttt{c-metric(s)}}\quad\textit{functoid}\index{c-metric(s)@\texttt{c-metric(s)}}\\
reads: the canonical metric space of \$1 \\
definition:
\begin{quote}\ttfamily\footnotesize\raggedright\noindent
c-metric-w(\allowbreak{}s, vnb-lambda(\allowbreak{}n, nn, (\allowbreak{}/)(\allowbreak{}1, 2 \textasciicircum{} (\allowbreak{}n + 1))))
\end{quote}
characterized by \texttt{c-metric-is-metric-space}  \textit{(proven)}:
\begin{quote}\ttfamily\footnotesize\raggedright\noindent
forall(\allowbreak{}[s], is-c-metric-space(\allowbreak{}s) implies is-metric-space(\allowbreak{}c-metric(\allowbreak{}s)))
\end{quote}
also: \texttt{c-metric-canonical-carrier} \\
mentioned by 10 result(s)

\smallskip

\noindent\textbf{\texttt{c-metric-space}}\quad\textit{structure}\index{c-metric-space@\texttt{c-metric-space}}\\
predicate \texttt{is-c-metric-space}; slots: pts dists

\smallskip

\noindent\textbf{\texttt{c-metric-w(s, w)}}\quad\textit{functoid}\index{c-metric-w(s, w)@\texttt{c-metric-w(s, w)}}\\
reads: the canonical metric of \$1 with weights \$2 \\
definition:
\begin{quote}\ttfamily\footnotesize\raggedright\noindent
[pts(\allowbreak{}s), vnb-lambda(\allowbreak{}[u, v], cartesian(\allowbreak{}pts(\allowbreak{}s), pts(\allowbreak{}s)), iota(\allowbreak{}l, l in rr and series-converges-to(\allowbreak{}vnb-lambda(\allowbreak{}k, nn, w(\allowbreak{}k) * min(\allowbreak{}1, (\allowbreak{}(\allowbreak{}dists(\allowbreak{}s))(\allowbreak{}k))(\allowbreak{}u, v))), l)))]
\end{quote}
characterized by \texttt{c-metric-w-is-metric-space}  \textit{(proven)}:
\begin{quote}\ttfamily\footnotesize\raggedright\noindent
forall(\allowbreak{}[s], is-c-metric-space(\allowbreak{}s) implies forall(\allowbreak{}[w], summable-weight(\allowbreak{}w) implies is-metric-space(\allowbreak{}c-metric-w(\allowbreak{}s, w))))
\end{quote}
mentioned by 19 result(s)

\smallskip

\noindent\textbf{\texttt{c-pseudo-family(s)}}\quad\textit{functoid}\index{c-pseudo-family(s)@\texttt{c-pseudo-family(s)}}\\
reads: the family of pseudometric spaces of \$1 \\
definition:
\begin{quote}\ttfamily\footnotesize\raggedright\noindent
vnb-lambda(\allowbreak{}k, nn, [pts(\allowbreak{}s), (\allowbreak{}dists(\allowbreak{}s))(\allowbreak{}k)])
\end{quote}

\smallskip

\noindent\textbf{\texttt{cact(v)}}\quad\textit{functoid}\index{cact(v)@\texttt{cact(v)}}\\
reads: the complex scalar multiplication of \$1 \\
definition:
\begin{quote}\ttfamily\footnotesize\raggedright\noindent
vnb-lambda(\allowbreak{}[cxz\_, cxx\_], cartesian(\allowbreak{}cc, vec(\allowbreak{}v)), (\allowbreak{}vadd(\allowbreak{}v))(\allowbreak{}(\allowbreak{}act(\allowbreak{}v))(\allowbreak{}real-part(\allowbreak{}cxz\_), cxx\_), (\allowbreak{}act(\allowbreak{}v))(\allowbreak{}imag-part(\allowbreak{}cxz\_), (\allowbreak{}jop(\allowbreak{}v))(\allowbreak{}cxx\_))))
\end{quote}
characterized by \texttt{cnvs-cact-i}  \textit{(proven)}:
\begin{quote}\ttfamily\footnotesize\raggedright\noindent
forall(\allowbreak{}[v], is-complex-normed-vector-space(\allowbreak{}v) implies forall(\allowbreak{}[x\_ in vec(\allowbreak{}v)], (\allowbreak{}cact(\allowbreak{}v))(\allowbreak{}1i, x\_) = (\allowbreak{}jop(\allowbreak{}v))(\allowbreak{}x\_)))
\end{quote}
also: \texttt{cna-cnvs-cact} \texttt{cnvs-cact-add} \texttt{cnvs-cact-fun} \texttt{cnvs-cact-mul} \texttt{cnvs-cact-one} \texttt{csip-cact-add} \texttt{csip-cact-fun} \texttt{csip-cact-mul} ... (26 in all) \\
mentioned by 73 result(s)

\smallskip

\noindent\textbf{\texttt{calc}}\quad\textit{tactic}\index{calc@\texttt{calc}}\\
\texttt{(calc L0 (rel1 L1 [just1]) (rel2 L2 [just2]) ...)} -- Ground the focus goal (REL L0 Ln) by a CHAIN of intermediaries L0 rel1 L1 rel2 L2 ... reln Ln: proves each link and composes them into the endpoint.  A link no lane can close is LEFT OPEN as a leaf -- the refinement point, and the PSS candidate.  Relations: = == (folded by transitivity), < <= (folded through the co-*-trans lemmas, with = steps riding along), IFF (chained per direction).  RETURNS the links, the open ones, and their PSS candidates as an alist. \\
declared in \texttt{tactics-help.scm}

\smallskip

\noindent\textbf{\texttt{canon-onb(kqv\_)}}\quad\textit{functoid}\index{canon-onb(kqv_)@\texttt{canon-onb(kqv\_)}}\\
reads: the chosen orthonormal basis of \$1 \\
definition:
\begin{quote}\ttfamily\footnotesize\raggedright\noindent
choice(\allowbreak{}onb-sets(\allowbreak{}kqv\_))
\end{quote}
characterized by \texttt{eigen-canon-onb}  \textit{(proven)}:
\begin{quote}\ttfamily\footnotesize\raggedright\noindent
forall(\allowbreak{}[v, t\_, c\_], is-complex-inner-product-space(\allowbreak{}v) implies t\_ in bounded-ops(\allowbreak{}v, v) implies c\_ in cc implies canon-onb(\allowbreak{}sub-cips(\allowbreak{}v, eigenspace(\allowbreak{}v, t\_, c\_))) subset eigenspace(\allowbreak{}v, t\_, c\_) and is-orthonormal-set(\allowbreak{}v, canon-onb(\allowbreak{}sub-cips(\allowbreak{}v, eigenspace(\allowbreak{}v, t\_, c\_)))) and forall(\allowbreak{}[tfy\_ in eigenspace(\allowbreak{}v, t\_, c\_)], forall(\allowbreak{}[tfu\_ in canon-onb(\allowbreak{}sub-cips(\allowbreak{}v, eigenspace(\allowbreak{}v, t\_, c\_)))], (\allowbreak{}ip(\allowbreak{}v))(\allowbreak{}tfy\_, tfu\_) = 0) implies tfy\_ = vzero(\allowbreak{}v)))
\end{quote}
also: \texttt{canon-onb-is-onb} \texttt{canon-onb-unfold} \\
mentioned by 13 result(s)

\smallskip

\noindent\textbf{\texttt{cantor-ms}}\quad\textit{constant}\index{cantor-ms@\texttt{cantor-ms}}\\
reads: the Cantor space \{0, 1\}\textasciicircum{}NN \\
defined by \texttt{cantor-ms-def}

\smallskip

\noindent\textbf{\texttt{caratheodory-sets(omega, mu)}}\quad\textit{functoid}\index{caratheodory-sets(omega, mu)@\texttt{caratheodory-sets(omega, mu)}}\\
reads: the \$2-measurable subsets of \$1 \\
definition:
\begin{quote}\ttfamily\footnotesize\raggedright\noindent
\{es\_ in power(\allowbreak{}omega): is-caratheodory-measurable(\allowbreak{}omega, mu, es\_)\}
\end{quote}
characterized by \texttt{caratheodory-sets-membership}  \textit{(definitional)}:
\begin{quote}\ttfamily\footnotesize\raggedright\noindent
forall(\allowbreak{}[omega, mu, es\_], es\_ in caratheodory-sets(\allowbreak{}omega, mu) iff es\_ in power(\allowbreak{}omega) and is-caratheodory-measurable(\allowbreak{}omega, mu, es\_))
\end{quote}
mentioned by 40 result(s)

\smallskip

\noindent\textbf{\texttt{card(a\_)}}\quad\textit{functoid}\index{card(a_)@\texttt{card(a\_)}}\\
reads: the cardinal of \$1 \\
definition:
\begin{quote}\ttfamily\footnotesize\raggedright\noindent
iota(\allowbreak{}alpha, alpha in ord and forsome(\allowbreak{}[phi], phi in bijection(\allowbreak{}a\_, ord-segment(\allowbreak{}alpha))) and forall(\allowbreak{}[beta], ord-lt(\allowbreak{}beta, alpha) implies not(\allowbreak{}forsome(\allowbreak{}[psi], psi in bijection(\allowbreak{}a\_, ord-segment(\allowbreak{}beta))))))
\end{quote}
characterized by \texttt{quotient-card-norm-cna-as-module}  \textit{(definitional)}:
\begin{quote}\ttfamily\footnotesize\raggedright\noindent
forall(\allowbreak{}[r], is-complex-normed-algebra(\allowbreak{}r) implies scal(\allowbreak{}r) = zz-ring implies forall(\allowbreak{}[n, u, sm], n in nn implies u in mat(\allowbreak{}n, 1, vec(\allowbreak{}r)) implies generates(\allowbreak{}cna-as-module(\allowbreak{}r), n, u) implies rel-free(\allowbreak{}cna-as-module(\allowbreak{}r), n, u) implies is-submodule(\allowbreak{}cna-as-module(\allowbreak{}r), sm) implies not(\allowbreak{}submodule-norm(\allowbreak{}cna-as-module(\allowbreak{}r), sm) = 0) implies card(\allowbreak{}vec(\allowbreak{}quotient-module(\allowbreak{}cna-as-module(\allowbreak{}r), sm))) = submodule-norm(\allowbreak{}cna-as-module(\allowbreak{}r), sm)))
\end{quote}
also: \texttt{quotient-card-norm-sip-as-module} \texttt{quotient-card-norm-uba-as-module} \texttt{quotient-card-norm-cnvs-as-module} \texttt{quotient-card-norm-sip-as-module-csip-as-sip} \texttt{quotient-card-norm-uba-as-module-uiba-as-uba} \texttt{card-bij} \texttt{card-pair} \texttt{dsum-card} ... (65 in all) \\
mentioned by 798 result(s)

\smallskip

\noindent\textbf{\texttt{carr(s)}}\quad\textit{accessor}\index{carr(s)@\texttt{carr(s)}}\\
reads: the carrier of \$1 \\
a structure slot \\
characterized by \texttt{rq-carr}  \textit{(definitional)}:
\begin{quote}\ttfamily\footnotesize\raggedright\noindent
forall(\allowbreak{}[r], carr(\allowbreak{}ringoid-quotient-ring(\allowbreak{}r)) == quotient(\allowbreak{}ringoid-setoid(\allowbreak{}r)))
\end{quote}
also: \texttt{cc-nag@carr} \texttt{qq-ring@carr} \texttt{zz-ring@carr} \texttt{qq-field@carr} \texttt{rr-field@carr} \texttt{lft-group@carr} \texttt{cips-scal-carrier} \texttt{nf-metric-carrier} ... (160 in all) \\
mentioned by 4783 result(s)

\smallskip

\noindent\textbf{\texttt{cartesian}}\quad\textit{operator}\index{cartesian@\texttt{cartesian}}\\
a term-forming head \\
characterized by \texttt{cartesian-set-iff}  \textit{(primitive)}:
\begin{quote}\ttfamily\footnotesize\raggedright\noindent
forall(\allowbreak{}[a, b], cartesian(\allowbreak{}a, b) in set iff a in set and b in set)
\end{quote}
also: \texttt{cartesian-nth} \texttt{card-cartesian} \texttt{cartesian-mono} \texttt{cartesian3-set} \texttt{cartesian4-set} \texttt{cartesian5-set} \texttt{esum-cartesian} \texttt{card-cartesian-nn} ... (19 in all) \\
mentioned by 1221 result(s)

\smallskip

\noindent\textbf{\texttt{cauchy-setoid(m)}}\quad\textit{functoid}\index{cauchy-setoid(m)@\texttt{cauchy-setoid(m)}}\\
definition:
\begin{quote}\ttfamily\footnotesize\raggedright\noindent
[cseq(\allowbreak{}m), crel(\allowbreak{}m)]
\end{quote}
characterized by \texttt{rkt-cauchy-setoid-pts}  \textit{(proven)}:
\begin{quote}\ttfamily\footnotesize\raggedright\noindent
forall(\allowbreak{}[m], is-metric-space(\allowbreak{}m) implies pts(\allowbreak{}cauchy-setoid(\allowbreak{}m)) == cseq(\allowbreak{}m))
\end{quote}
also: \texttt{rko2-cauchy-setoid-rel} \texttt{cauchy-setoid-is-setoid} \texttt{rko2-cauchy-setoid-length} \\
mentioned by 81 result(s)

\smallskip

\noindent\textbf{\texttt{cayley}}\quad\textit{constant}\index{cayley@\texttt{cayley}}\\
reads: the Cayley transform z |-> (z - i)/(z + i) \\
defined by \texttt{cayley-def} \\
characterized by \texttt{cayley-at}  \textit{(proven)}:
\begin{quote}\ttfamily\footnotesize\raggedright\noindent
forall(\allowbreak{}[w in cc], not(\allowbreak{}w + 1i = 0) implies cayley(\allowbreak{}w) = (\allowbreak{}w - 1i) * recip(\allowbreak{}w + 1i))
\end{quote}
also: \texttt{cayley-mag} \texttt{cayley-pole} \texttt{cayley-disc-iff} \texttt{cayley-maps-real} \texttt{cayley-circle-iff} \texttt{cayley-maps-upper} \texttt{cayley-upper-disc} \texttt{cayley-real-circle} ... (10 in all) \\
mentioned by 18 result(s)

\smallskip

\noindent\textbf{\texttt{cayley-graph(ubv\_, ubg\_)}}\quad\textit{functoid}\index{cayley-graph(ubv_, ubg_)@\texttt{cayley-graph(ubv\_, ubg\_)}}\\
reads: the Cayley transform of \$2 \\
definition:
\begin{quote}\ttfamily\footnotesize\raggedright\noindent
\{ubp\_ in cartesian(\allowbreak{}vec(\allowbreak{}ubv\_), vec(\allowbreak{}ubv\_)): forsome(\allowbreak{}[ubq\_ in ubg\_], ubp\_ = [(\allowbreak{}vadd(\allowbreak{}ubv\_))(\allowbreak{}(\allowbreak{}act(\allowbreak{}ubv\_))(\allowbreak{}1i, nth(\allowbreak{}1, ubq\_)), nth(\allowbreak{}2, ubq\_)), (\allowbreak{}vadd(\allowbreak{}ubv\_))(\allowbreak{}(\allowbreak{}act(\allowbreak{}ubv\_))(\allowbreak{}-1i, nth(\allowbreak{}1, ubq\_)), nth(\allowbreak{}2, ubq\_))])\}
\end{quote}
characterized by \texttt{cayley-graph-is-set}  \textit{(proven)}:
\begin{quote}\ttfamily\footnotesize\raggedright\noindent
forall(\allowbreak{}[v, gg], is-complex-inner-product-space(\allowbreak{}v) implies cayley-graph(\allowbreak{}v, gg) in set)
\end{quote}
also: \texttt{cayley-graph-linear} \texttt{cayley-graph-subset} \texttt{cayley-graph-injective} \\
mentioned by 26 result(s)

\smallskip

\noindent\textbf{\texttt{cbop-comp(cqu\_, cqs\_, cqt\_)}}\quad\textit{functoid}\index{cbop-comp(cqu_, cqs_, cqt_)@\texttt{cbop-comp(cqu\_, cqs\_, cqt\_)}}\\
reads: \$2 \$3 \\
definition:
\begin{quote}\ttfamily\footnotesize\raggedright\noindent
vnb-lambda(\allowbreak{}cqx\_, vec(\allowbreak{}cqu\_), cqs\_(\allowbreak{}cqt\_(\allowbreak{}cqx\_)))
\end{quote}
characterized by \texttt{cbop-comp-fun}  \textit{(proven)}:
\begin{quote}\ttfamily\footnotesize\raggedright\noindent
forall(\allowbreak{}[u, v, w, s\_, t\_], is-complex-normed-vector-space(\allowbreak{}u) implies t\_ in fun(\allowbreak{}vec(\allowbreak{}u), vec(\allowbreak{}v)) implies s\_ in fun(\allowbreak{}vec(\allowbreak{}v), vec(\allowbreak{}w)) implies cbop-comp(\allowbreak{}u, s\_, t\_) in fun(\allowbreak{}vec(\allowbreak{}u), vec(\allowbreak{}w)))
\end{quote}
also: \texttt{cbop-comp-apply} \texttt{cbop-comp-assoc} \texttt{cbop-comp-bounded} \texttt{cbop-comp-id-left} \texttt{cbop-comp-id-right} \texttt{cbop-comp-sum-left} \texttt{cbop-comp-sum-right} \\
mentioned by 15 result(s)

\smallskip

\noindent\textbf{\texttt{cbop-id(cqu\_)}}\quad\textit{functoid}\index{cbop-id(cqu_)@\texttt{cbop-id(cqu\_)}}\\
reads: the identity operator of \$1 \\
definition:
\begin{quote}\ttfamily\footnotesize\raggedright\noindent
vnb-lambda(\allowbreak{}cqx\_, vec(\allowbreak{}cqu\_), cqx\_)
\end{quote}
characterized by \texttt{cbop-id-fun}  \textit{(proven)}:
\begin{quote}\ttfamily\footnotesize\raggedright\noindent
forall(\allowbreak{}[u], is-complex-normed-vector-space(\allowbreak{}u) implies cbop-id(\allowbreak{}u) in fun(\allowbreak{}vec(\allowbreak{}u), vec(\allowbreak{}u)))
\end{quote}
also: \texttt{cbop-id-apply} \texttt{cbop-id-bounded} \\
mentioned by 8 result(s)

\smallskip

\noindent\textbf{\texttt{cbop-neg(cbv\_, cbw\_, cbt\_)}}\quad\textit{functoid}\index{cbop-neg(cbv_, cbw_, cbt_)@\texttt{cbop-neg(cbv\_, cbw\_, cbt\_)}}\\
definition:
\begin{quote}\ttfamily\footnotesize\raggedright\noindent
vnb-lambda(\allowbreak{}cbx\_, vec(\allowbreak{}cbv\_), (\allowbreak{}vneg(\allowbreak{}cbw\_))(\allowbreak{}cbt\_(\allowbreak{}cbx\_)))
\end{quote}
characterized by \texttt{cbop-neg-apply}  \textit{(proven)}:
\begin{quote}\ttfamily\footnotesize\raggedright\noindent
forall(\allowbreak{}[v, w, t\_, x\_], is-complex-normed-vector-space(\allowbreak{}w) implies t\_ in fun(\allowbreak{}vec(\allowbreak{}v), vec(\allowbreak{}w)) implies x\_ in vec(\allowbreak{}v) implies (\allowbreak{}cbop-neg(\allowbreak{}v, w, t\_))(\allowbreak{}x\_) = (\allowbreak{}vneg(\allowbreak{}w))(\allowbreak{}t\_(\allowbreak{}x\_)))
\end{quote}
also: \texttt{cbop-neg-unfold} \texttt{cbop-neg-bounded} \\
mentioned by 15 result(s)

\smallskip

\noindent\textbf{\texttt{cbop-normed-ag(cbv\_, cbw\_)}}\quad\textit{functoid}\index{cbop-normed-ag(cbv_, cbw_)@\texttt{cbop-normed-ag(cbv\_, cbw\_)}}\\
reads: the normed group of bounded operators from \$1 to \$2 \\
definition:
\begin{quote}\ttfamily\footnotesize\raggedright\noindent
[cnvs-bounded-ops(\allowbreak{}cbv\_, cbw\_), cbop-vadd(\allowbreak{}cbv\_, cbw\_), cbop-zero(\allowbreak{}cbv\_, cbw\_), cbop-vinv(\allowbreak{}cbv\_, cbw\_), cbop-vnrm(\allowbreak{}cbv\_, cbw\_)]
\end{quote}
characterized by \texttt{cbop-normed-ag-is-normed-ag}  \textit{(proven)}:
\begin{quote}\ttfamily\footnotesize\raggedright\noindent
forall(\allowbreak{}[v, w], is-complex-normed-vector-space(\allowbreak{}v) implies is-complex-normed-vector-space(\allowbreak{}w) implies is-normed-ag(\allowbreak{}cbop-normed-ag(\allowbreak{}v, w)))
\end{quote}
mentioned by 22 result(s)

\smallskip

\noindent\textbf{\texttt{cbop-scale(cbv\_, cbw\_, cbc\_, cbt\_)}}\quad\textit{functoid}\index{cbop-scale(cbv_, cbw_, cbc_, cbt_)@\texttt{cbop-scale(cbv\_, cbw\_, cbc\_, cbt\_)}}\\
definition:
\begin{quote}\ttfamily\footnotesize\raggedright\noindent
vnb-lambda(\allowbreak{}cbx\_, vec(\allowbreak{}cbv\_), (\allowbreak{}cact(\allowbreak{}cbw\_))(\allowbreak{}cbc\_, cbt\_(\allowbreak{}cbx\_)))
\end{quote}
characterized by \texttt{cbop-scale-add}  \textit{(proven)}:
\begin{quote}\ttfamily\footnotesize\raggedright\noindent
forall(\allowbreak{}[v, w, c\_, d\_, t\_], is-complex-normed-vector-space(\allowbreak{}v) implies is-complex-normed-vector-space(\allowbreak{}w) implies c\_ in cc implies d\_ in cc implies t\_ in cnvs-bounded-ops(\allowbreak{}v, w) implies cbop-scale(\allowbreak{}v, w, c\_ + d\_, t\_) = cbop-sum(\allowbreak{}v, w, cbop-scale(\allowbreak{}v, w, c\_, t\_), cbop-scale(\allowbreak{}v, w, d\_, t\_)))
\end{quote}
also: \texttt{cbop-scale-mul} \texttt{cbop-scale-one} \texttt{cbop-scale-sum} \texttt{cbop-scale-apply} \texttt{cbop-scale-bounded} \\
mentioned by 13 result(s)

\smallskip

\noindent\textbf{\texttt{cbop-sum(cbv\_, cbw\_, cbs\_, cbt\_)}}\quad\textit{functoid}\index{cbop-sum(cbv_, cbw_, cbs_, cbt_)@\texttt{cbop-sum(cbv\_, cbw\_, cbs\_, cbt\_)}}\\
definition:
\begin{quote}\ttfamily\footnotesize\raggedright\noindent
vnb-lambda(\allowbreak{}cbx\_, vec(\allowbreak{}cbv\_), (\allowbreak{}vadd(\allowbreak{}cbw\_))(\allowbreak{}cbs\_(\allowbreak{}cbx\_), cbt\_(\allowbreak{}cbx\_)))
\end{quote}
characterized by \texttt{cbop-sum-comm}  \textit{(proven)}:
\begin{quote}\ttfamily\footnotesize\raggedright\noindent
forall(\allowbreak{}[v, w, s\_, t\_], is-complex-normed-vector-space(\allowbreak{}v) implies is-complex-normed-vector-space(\allowbreak{}w) implies s\_ in cnvs-bounded-ops(\allowbreak{}v, w) implies t\_ in cnvs-bounded-ops(\allowbreak{}v, w) implies cbop-sum(\allowbreak{}v, w, s\_, t\_) = cbop-sum(\allowbreak{}v, w, t\_, s\_))
\end{quote}
also: \texttt{cbop-sum-apply} \texttt{cbop-sum-assoc} \texttt{cbop-sum-unfold} \texttt{cbop-sum-bounded} \texttt{cbop-sum-neg-left} \texttt{cbop-sum-neg-right} \texttt{cbop-sum-zero-left} \texttt{cbop-sum-zero-right} \\
mentioned by 30 result(s)

\smallskip

\noindent\textbf{\texttt{cbop-vadd(cbv\_, cbw\_)}}\quad\textit{functoid}\index{cbop-vadd(cbv_, cbw_)@\texttt{cbop-vadd(cbv\_, cbw\_)}}\\
definition:
\begin{quote}\ttfamily\footnotesize\raggedright\noindent
vnb-lambda(\allowbreak{}[cbs\_, cbt\_], cartesian(\allowbreak{}cnvs-bounded-ops(\allowbreak{}cbv\_, cbw\_), cnvs-bounded-ops(\allowbreak{}cbv\_, cbw\_)), cbop-sum(\allowbreak{}cbv\_, cbw\_, cbs\_, cbt\_))
\end{quote}
characterized by \texttt{cbop-vadd-value}  \textit{(proven)}:
\begin{quote}\ttfamily\footnotesize\raggedright\noindent
forall(\allowbreak{}[v, w, s\_, t\_], is-complex-normed-vector-space(\allowbreak{}v) implies is-complex-normed-vector-space(\allowbreak{}w) implies s\_ in cnvs-bounded-ops(\allowbreak{}v, w) implies t\_ in cnvs-bounded-ops(\allowbreak{}v, w) implies (\allowbreak{}cbop-vadd(\allowbreak{}v, w))(\allowbreak{}s\_, t\_) = cbop-sum(\allowbreak{}v, w, s\_, t\_))
\end{quote}
mentioned by 6 result(s)

\smallskip

\noindent\textbf{\texttt{cbop-vinv(cbv\_, cbw\_)}}\quad\textit{functoid}\index{cbop-vinv(cbv_, cbw_)@\texttt{cbop-vinv(cbv\_, cbw\_)}}\\
definition:
\begin{quote}\ttfamily\footnotesize\raggedright\noindent
vnb-lambda(\allowbreak{}cbt\_, cnvs-bounded-ops(\allowbreak{}cbv\_, cbw\_), cbop-neg(\allowbreak{}cbv\_, cbw\_, cbt\_))
\end{quote}
characterized by \texttt{cbop-vinv-value}  \textit{(proven)}:
\begin{quote}\ttfamily\footnotesize\raggedright\noindent
forall(\allowbreak{}[v, w, t\_], is-complex-normed-vector-space(\allowbreak{}v) implies is-complex-normed-vector-space(\allowbreak{}w) implies t\_ in cnvs-bounded-ops(\allowbreak{}v, w) implies (\allowbreak{}cbop-vinv(\allowbreak{}v, w))(\allowbreak{}t\_) = cbop-neg(\allowbreak{}v, w, t\_))
\end{quote}
mentioned by 6 result(s)

\smallskip

\noindent\textbf{\texttt{cbop-vnrm(cbv\_, cbw\_)}}\quad\textit{functoid}\index{cbop-vnrm(cbv_, cbw_)@\texttt{cbop-vnrm(cbv\_, cbw\_)}}\\
definition:
\begin{quote}\ttfamily\footnotesize\raggedright\noindent
vnb-lambda(\allowbreak{}cbt\_, cnvs-bounded-ops(\allowbreak{}cbv\_, cbw\_), cnvs-opnorm(\allowbreak{}cbv\_, cbw\_, cbt\_))
\end{quote}
characterized by \texttt{cbop-vnrm-value}  \textit{(proven)}:
\begin{quote}\ttfamily\footnotesize\raggedright\noindent
forall(\allowbreak{}[v, w, t\_], is-complex-normed-vector-space(\allowbreak{}v) implies is-complex-normed-vector-space(\allowbreak{}w) implies t\_ in cnvs-bounded-ops(\allowbreak{}v, w) implies (\allowbreak{}cbop-vnrm(\allowbreak{}v, w))(\allowbreak{}t\_) = cnvs-opnorm(\allowbreak{}v, w, t\_))
\end{quote}
mentioned by 6 result(s)

\smallskip

\noindent\textbf{\texttt{cbop-zero(cbv\_, cbw\_)}}\quad\textit{functoid}\index{cbop-zero(cbv_, cbw_)@\texttt{cbop-zero(cbv\_, cbw\_)}}\\
definition:
\begin{quote}\ttfamily\footnotesize\raggedright\noindent
vnb-lambda(\allowbreak{}cbx\_, vec(\allowbreak{}cbv\_), vzero(\allowbreak{}cbw\_))
\end{quote}
characterized by \texttt{cbop-zero-apply}  \textit{(proven)}:
\begin{quote}\ttfamily\footnotesize\raggedright\noindent
forall(\allowbreak{}[v, w, x\_], is-complex-normed-vector-space(\allowbreak{}w) implies x\_ in vec(\allowbreak{}v) implies (\allowbreak{}cbop-zero(\allowbreak{}v, w))(\allowbreak{}x\_) = vzero(\allowbreak{}w))
\end{quote}
also: \texttt{cbop-zero-bounded} \\
mentioned by 21 result(s)

\smallskip

\noindent\textbf{\texttt{cc-cont-fns(u)}}\quad\textit{functoid}\index{cc-cont-fns(u)@\texttt{cc-cont-fns(u)}}\\
reads: the continuous complex functions on \$1 \\
definition:
\begin{quote}\ttfamily\footnotesize\raggedright\noindent
\{ccf\_ in fun(\allowbreak{}u, cc): forall(\allowbreak{}[ccy\_ in u], is-continuous-at(\allowbreak{}subspace-ms(\allowbreak{}nf-metric-space(\allowbreak{}cc-normed-field), u), nf-metric-space(\allowbreak{}cc-normed-field), ccf\_, ccy\_))\}
\end{quote}
characterized by \texttt{cc-cont-fns-add}  \textit{(proven)}:
\begin{quote}\ttfamily\footnotesize\raggedright\noindent
forall(\allowbreak{}[u], u subset cc implies forall(\allowbreak{}[a in cc-cont-fns(\allowbreak{}u), b in cc-cont-fns(\allowbreak{}u)], vnb-lambda(\allowbreak{}cfz\_, u, a(\allowbreak{}cfz\_) + b(\allowbreak{}cfz\_)) in cc-cont-fns(\allowbreak{}u)))
\end{quote}
also: \texttt{cc-cont-fns-fun} \texttt{cc-cont-fns-neg} \texttt{cc-cont-fns-diff} \texttt{cc-cont-fns-zero} \texttt{cc-cont-fns-scale} \texttt{cc-cont-fns-is-set} \texttt{cc-cont-fns-unfold} \texttt{cc-cont-fns-add-neg} ... (10 in all) \\
mentioned by 61 result(s)

\smallskip

\noindent\textbf{\texttt{cc-coords}}\quad\textit{defined-fn}\index{cc-coords@\texttt{cc-coords}}\\
defined by \texttt{cc-coords-def} \\
characterized by \texttt{cc-coords-add}  \textit{(proven)}:
\begin{quote}\ttfamily\footnotesize\raggedright\noindent
forall(\allowbreak{}[z in cc, w in cc], cc-coords(\allowbreak{}z + w) = [real-part(\allowbreak{}z) + real-part(\allowbreak{}w), imag-part(\allowbreak{}z) + imag-part(\allowbreak{}w)])
\end{quote}
also: \texttt{cc-coords-apply} \texttt{cc-coords-dist-ge} \texttt{cc-coords-dist-le} \texttt{cc-coords-of-pair} \texttt{cc-coords-real-mul} \\
mentioned by 12 result(s)

\smallskip

\noindent\textbf{\texttt{cc-cos(z)}}\quad\textit{functoid}\index{cc-cos(z)@\texttt{cc-cos(z)}}\\
reads: cos \$1 \\
definition:
\begin{quote}\ttfamily\footnotesize\raggedright\noindent
cc-series-limit(\allowbreak{}vnb-lambda(\allowbreak{}cek\_, nn, (\allowbreak{}num(\allowbreak{}-1) \textasciicircum{} cek\_ * recip(\allowbreak{}factorial(\allowbreak{}2 * cek\_))) * z \textasciicircum{} (\allowbreak{}2 * cek\_)))
\end{quote}
characterized by \texttt{cc-cos-pi}  \textit{(proven)}:
\begin{quote}\ttfamily\footnotesize\raggedright\noindent
cc-cos(\allowbreak{}pi) = num(\allowbreak{}-1)
\end{quote}
also: \texttt{cc-cos-add} \texttt{cc-cos-exp} \texttt{cc-cos-even} \texttt{cc-cos-zero} \texttt{cc-cos-deriv} \texttt{cc-cos-sin-pi} \texttt{cc-cos-two-pi} \texttt{cc-cos-unfold} ... (27 in all) \\
mentioned by 46 result(s)

\smallskip

\noindent\textbf{\texttt{cc-exhaust(u, n)}}\quad\textit{functoid}\index{cc-exhaust(u, n)@\texttt{cc-exhaust(u, n)}}\\
reads: the \$2-th compact set of the exhaustion of \$1 \\
definition:
\begin{quote}\ttfamily\footnotesize\raggedright\noindent
\{cxz\_ in cc: magnitude(\allowbreak{}cxz\_) <= n and ball(\allowbreak{}cc-ms, cxz\_, recip(\allowbreak{}n + 1)) subset u\}
\end{quote}
characterized by \texttt{cc-exhaust-mono}  \textit{(proven)}:
\begin{quote}\ttfamily\footnotesize\raggedright\noindent
forall(\allowbreak{}[u, a in nn, b in nn], a <= b implies cc-exhaust(\allowbreak{}u, a) subset cc-exhaust(\allowbreak{}u, b))
\end{quote}
also: \texttt{cc-exhaust-covers} \texttt{cc-exhaust-subset} \texttt{cc-exhaust-unfold} \texttt{cc-exhaust-compact} \texttt{cc-exhaust-interior} \texttt{cc-exhaust-limit-in} \texttt{cc-exhaust-dense-in-u} \texttt{cc-exhaust-dense-seqs} ... (18 in all) \\
mentioned by 48 result(s)

\smallskip

\noindent\textbf{\texttt{cc-exhaust-seminorms(u, v)}}\quad\textit{functoid}\index{cc-exhaust-seminorms(u, v)@\texttt{cc-exhaust-seminorms(u, v)}}\\
reads: the sup seminorms on the exhaustion of \$1, on \$2 \\
definition:
\begin{quote}\ttfamily\footnotesize\raggedright\noindent
vnb-lambda(\allowbreak{}cfk\_, nn, vnb-lambda(\allowbreak{}cfa\_, v, cc-sup-norm(\allowbreak{}cfa\_, cc-exhaust(\allowbreak{}u, cfk\_))))
\end{quote}
characterized by \texttt{cc-exhaust-seminorms-family}  \textit{(proven)}:
\begin{quote}\ttfamily\footnotesize\raggedright\noindent
forall(\allowbreak{}[u], is-open(\allowbreak{}nf-metric-space(\allowbreak{}cc-normed-field), u) implies is-seminorm-family(\allowbreak{}cc-normed-field, cc-fn-module(\allowbreak{}u, cc-cont-fns(\allowbreak{}u)), cc-exhaust-seminorms(\allowbreak{}u, cc-cont-fns(\allowbreak{}u))))
\end{quote}
also: \texttt{cc-exhaust-seminorms-unfold} \\
mentioned by 22 result(s)

\smallskip

\noindent\textbf{\texttt{cc-exp(z)}}\quad\textit{functoid}\index{cc-exp(z)@\texttt{cc-exp(z)}}\\
reads: exp \$1 \\
definition:
\begin{quote}\ttfamily\footnotesize\raggedright\noindent
cc-series-limit(\allowbreak{}vnb-lambda(\allowbreak{}cek\_, nn, recip(\allowbreak{}factorial(\allowbreak{}cek\_)) * z \textasciicircum{} cek\_))
\end{quote}
characterized by \texttt{cc-exp-i}  \textit{(proven)}:
\begin{quote}\ttfamily\footnotesize\raggedright\noindent
forall(\allowbreak{}[cexz\_ in cc], cc-exp(\allowbreak{}1i * cexz\_) = cc-cos(\allowbreak{}cexz\_) + 1i * cc-sin(\allowbreak{}cexz\_))
\end{quote}
also: \texttt{cc-exp-add} \texttt{cc-exp-log} \texttt{cc-exp-neg} \texttt{cc-exp-conj} \texttt{cc-exp-i-pi} \texttt{cc-exp-real} \texttt{cc-exp-tail} \texttt{cc-exp-zero} ... (52 in all) \\
mentioned by 234 result(s)

\smallskip

\noindent\textbf{\texttt{cc-family-sum(f)}}\quad\textit{functoid}\index{cc-family-sum(f)@\texttt{cc-family-sum(f)}}\\
reads: the sum in CC of the family \$1 \\
definition:
\begin{quote}\ttfamily\footnotesize\raggedright\noindent
choice(\allowbreak{}\{csw\_ in cc: sums-to(\allowbreak{}cc-nag, f, csw\_)\})
\end{quote}
characterized by \texttt{cc-family-sum-value}  \textit{(proven)}:
\begin{quote}\ttfamily\footnotesize\raggedright\noindent
forall(\allowbreak{}[f in fun(\allowbreak{}dom(\allowbreak{}f), cc), r in cc], sums-to(\allowbreak{}cc-nag, f, r) implies cc-family-sum(\allowbreak{}f) = r)
\end{quote}
mentioned by 23 result(s)

\smallskip

\noindent\textbf{\texttt{cc-fn-module(u, v)}}\quad\textit{functoid}\index{cc-fn-module(u, v)@\texttt{cc-fn-module(u, v)}}\\
reads: the module of the functions \$2 on \$1, with pointwise operations \\
definition:
\begin{quote}\ttfamily\footnotesize\raggedright\noindent
[normed-field-as-commutative-ring(\allowbreak{}cc-normed-field), v, vnb-lambda(\allowbreak{}[cfa\_, cfb\_], cartesian(\allowbreak{}v, v), vnb-lambda(\allowbreak{}cfz\_, u, cfa\_(\allowbreak{}cfz\_) + cfb\_(\allowbreak{}cfz\_))), vnb-lambda(\allowbreak{}cfz\_, u, 0), vnb-lambda(\allowbreak{}cfa\_, v, vnb-lambda(\allowbreak{}cfz\_, u, -cfa\_(\allowbreak{}cfz\_))), vnb-lambda(\allowbreak{}[cfc\_, cfa\_], cartesian(\allowbreak{}cc, v), vnb-lambda(\allowbreak{}cfz\_, u, cfc\_ * cfa\_(\allowbreak{}cfz\_)))]
\end{quote}
characterized by \texttt{cc-fn-module-act}  \textit{(proven)}:
\begin{quote}\ttfamily\footnotesize\raggedright\noindent
forall(\allowbreak{}[u in set, v in set], act(\allowbreak{}cc-fn-module(\allowbreak{}u, v)) == vnb-lambda(\allowbreak{}[cfc\_, cfa\_], cartesian(\allowbreak{}cc, v), vnb-lambda(\allowbreak{}cfz\_, u, cfc\_ * cfa\_(\allowbreak{}cfz\_))))
\end{quote}
also: \texttt{cc-fn-module-vec} \texttt{cc-fn-module-scal} \texttt{cc-fn-module-vadd} \texttt{cc-fn-module-vneg} \texttt{cc-fn-module-vzero} \texttt{cc-fn-module-is-module} \texttt{cc-fn-module-is-module-of} \texttt{cc-fn-module-cauchy-uniform} ... (9 in all) \\
mentioned by 35 result(s)

\smallskip

\noindent\textbf{\texttt{cc-fraction-field(aa)}}\quad\textit{functoid}\index{cc-fraction-field(aa)@\texttt{cc-fraction-field(aa)}}\\
reads: the field of fractions of \$1 inside the complex numbers \\
definition:
\begin{quote}\ttfamily\footnotesize\raggedright\noindent
\{anfq\_ in cc: forsome(\allowbreak{}[anfa\_ in aa, anfb\_ in aa], not(\allowbreak{}anfb\_ = 0) and anfq\_ * anfb\_ = anfa\_)\}
\end{quote}
characterized by \texttt{cc-fraction-field-sub}  \textit{(proven)}:
\begin{quote}\ttfamily\footnotesize\raggedright\noindent
forall(\allowbreak{}[aa], cc-fraction-field(\allowbreak{}aa) subset \{anfq\_ in cc: forsome(\allowbreak{}[anfa\_ in aa, anfb\_ in aa], not(\allowbreak{}anfb\_ = 0) and anfq\_ * anfb\_ = anfa\_)\})
\end{quote}
also: \texttt{cc-fraction-field-sup} \\
mentioned by 5 result(s)

\smallskip

\noindent\textbf{\texttt{cc-hat}}\quad\textit{constant}\index{cc-hat@\texttt{cc-hat}}\\
reads: the extended complex plane C u \{oo\} \\
defined by \texttt{cc-hat-def}

\smallskip

\noindent\textbf{\texttt{cc-hat-opens}}\quad\textit{constant}\index{cc-hat-opens@\texttt{cc-hat-opens}}\\
reads: the open sets of the extended complex plane \\
defined by \texttt{cc-hat-opens-def}

\smallskip

\noindent\textbf{\texttt{cc-hat-top}}\quad\textit{instance}\index{cc-hat-top@\texttt{cc-hat-top}}\\
reads: the extended complex plane with its topology \\
an instance of top-space; predicate \texttt{is-cc-hat-top}; same shape as top-space

\smallskip

\noindent\textbf{\texttt{cc-hol-bounded(u, a)}}\quad\textit{predicate}\index{cc-hol-bounded(u, a)@\texttt{cc-hol-bounded(u, a)}}\\
reads: \$2 is a bounded set of holomorphic functions on \$1 \\
defined by \texttt{cc-hol-bounded} \\
characterized by \texttt{cc-hol-bounded}  \textit{(definitional)}:
\begin{quote}\ttfamily\footnotesize\raggedright\noindent
forall(\allowbreak{}[u, a], cc-hol-bounded(\allowbreak{}u, a) iff a subset cc-hol-fns(\allowbreak{}u) and forall(\allowbreak{}[chn\_ in nn], forsome(\allowbreak{}[chm\_ in rr], forall(\allowbreak{}[chf\_ in a, chz\_ in cc-exhaust(\allowbreak{}u, chn\_)], magnitude(\allowbreak{}chf\_(\allowbreak{}chz\_)) <= chm\_))))
\end{quote}
also: \texttt{cc-hol-bounded-mono} \texttt{cc-hol-bounded-subseq} \texttt{cc-hol-bounded-equicontinuous} \\
mentioned by 10 result(s)

\smallskip

\noindent\textbf{\texttt{cc-hol-fns(u)}}\quad\textit{functoid}\index{cc-hol-fns(u)@\texttt{cc-hol-fns(u)}}\\
reads: the holomorphic functions on \$1 \\
definition:
\begin{quote}\ttfamily\footnotesize\raggedright\noindent
\{ccf\_ in fun(\allowbreak{}u, cc): holomorphic-on(\allowbreak{}u, ccf\_)\}
\end{quote}
characterized by \texttt{cc-hol-fns-add}  \textit{(proven)}:
\begin{quote}\ttfamily\footnotesize\raggedright\noindent
forall(\allowbreak{}[u], is-open(\allowbreak{}nf-metric-space(\allowbreak{}cc-normed-field), u) implies forall(\allowbreak{}[a in cc-hol-fns(\allowbreak{}u), b in cc-hol-fns(\allowbreak{}u)], vnb-lambda(\allowbreak{}cfz\_, u, a(\allowbreak{}cfz\_) + b(\allowbreak{}cfz\_)) in cc-hol-fns(\allowbreak{}u)))
\end{quote}
also: \texttt{cc-hol-fns-neg} \texttt{cc-hol-fns-zero} \texttt{cc-hol-fns-scale} \texttt{cc-hol-fns-closed} \texttt{cc-hol-fns-unfold} \texttt{cc-hol-fns-add-neg} \texttt{cc-hol-fns-frechet} \texttt{cc-hol-fns-membership} ... (9 in all) \\
mentioned by 36 result(s)

\smallskip

\noindent\textbf{\texttt{cc-inf}}\quad\textit{constant}\index{cc-inf@\texttt{cc-inf}}\\
reads: the point at infinity of the extended complex plane

\smallskip

\noindent\textbf{\texttt{cc-int(pphi\_, a, b)}}\quad\textit{functoid}\index{cc-int(pphi_, a, b)@\texttt{cc-int(pphi\_, a, b)}}\\
reads: the complex integral of \$1 from \$2 to \$3 \\
definition:
\begin{quote}\ttfamily\footnotesize\raggedright\noindent
pw-int(\allowbreak{}vnb-lambda(\allowbreak{}pat\_, ccint(\allowbreak{}a, b), real-part(\allowbreak{}pphi\_(\allowbreak{}pat\_))), a, b) + pw-int(\allowbreak{}vnb-lambda(\allowbreak{}pat\_, ccint(\allowbreak{}a, b), imag-part(\allowbreak{}pphi\_(\allowbreak{}pat\_))), a, b) * 1i
\end{quote}
characterized by \texttt{cc-int-sum}  \textit{(proven)}:
\begin{quote}\ttfamily\footnotesize\raggedright\noindent
forall(\allowbreak{}[a, b], a in rr implies b in rr implies forall(\allowbreak{}[pphi\_, ppsi\_, pchi\_], pphi\_ in fun(\allowbreak{}ccint(\allowbreak{}a, b), cc) implies ppsi\_ in fun(\allowbreak{}ccint(\allowbreak{}a, b), cc) implies pchi\_ in fun(\allowbreak{}ccint(\allowbreak{}a, b), cc) implies forall(\allowbreak{}[pay\_ in ccint(\allowbreak{}a, b)], pchi\_(\allowbreak{}pay\_) == pphi\_(\allowbreak{}pay\_) + ppsi\_(\allowbreak{}pay\_)) implies forall(\allowbreak{}[pwf\_, paw\_, pauh\_, pbw\_], is-primitive(\allowbreak{}pwf\_, vnb-lambda(\allowbreak{}pat\_, ccint(\allowbreak{}a, b), real-part(\allowbreak{}pphi\_(\allowbreak{}pat\_))), a, b) implies is-primitive(\allowbreak{}paw\_, vnb-lambda(\allowbreak{}pat\_, ccint(\allowbreak{}a, b), imag-part(\allowbreak{}pphi\_(\allowbreak{}pat\_))), a, b) implies is-primitive(\allowbreak{}pauh\_, vnb-lambda(\allowbreak{}pat\_, ccint(\allowbreak{}a, b), real-part(\allowbreak{}ppsi\_(\allowbreak{}pat\_))), a, b) implies is-primitive(\allowbreak{}pbw\_, vnb-lambda(\allowbreak{}pat\_, ccint(\allowbreak{}a, b), imag-part(\allowbreak{}ppsi\_(\allowbreak{}pat\_))), a, b) implies cc-int(\allowbreak{}pchi\_, a, b) = cc-int(\allowbreak{}pphi\_, a, b) + cc-int(\allowbreak{}ppsi\_, a, b))))
\end{quote}
also: \texttt{cc-int-in-cc} \texttt{cc-int-split} \texttt{cc-int-value} \texttt{cc-int-unfold} \texttt{cc-int-reflect} \texttt{cc-int-adjacent} \texttt{cc-int-real-mul} \texttt{cc-int-abs-bound} ... (16 in all) \\
mentioned by 57 result(s)

\smallskip

\noindent\textbf{\texttt{cc-lincomb(n, b, g)}}\quad\textit{functoid}\index{cc-lincomb(n, b, g)@\texttt{cc-lincomb(n, b, g)}}\\
reads: the sum of \$2(j) \$3(j) for j = 1, ..., \$1 \\
definition:
\begin{quote}\ttfamily\footnotesize\raggedright\noindent
finsum(\allowbreak{}ring-additive-ag(\allowbreak{}normed-field-as-commutative-ring(\allowbreak{}cc-normed-field)), vnb-lambda(\allowbreak{}almj\_, interval(\allowbreak{}1, n), b(\allowbreak{}almj\_) * g(\allowbreak{}almj\_)), interval(\allowbreak{}1, n))
\end{quote}
characterized by \texttt{cc-lincomb-add}  \textit{(proven)}:
\begin{quote}\ttfamily\footnotesize\raggedright\noindent
forall(\allowbreak{}[n, b1, b2, g], n in nn implies b1 in fun(\allowbreak{}interval(\allowbreak{}1, n), cc) implies b2 in fun(\allowbreak{}interval(\allowbreak{}1, n), cc) implies g in fun(\allowbreak{}interval(\allowbreak{}1, n), cc) implies cc-lincomb(\allowbreak{}n, b1, g) + cc-lincomb(\allowbreak{}n, b2, g) = cc-lincomb(\allowbreak{}n, vnb-lambda(\allowbreak{}amj\_, interval(\allowbreak{}1, n), b1(\allowbreak{}amj\_) + b2(\allowbreak{}amj\_)), g))
\end{quote}
also: \texttt{cc-lincomb-unit} \texttt{cc-lincomb-scale} \texttt{cc-lincomb-unfold} \texttt{cc-lincomb-zero-generators} \texttt{cc-lincomb-zero-coefficients} \\
mentioned by 23 result(s)

\smallskip

\noindent\textbf{\texttt{cc-line(clnp\_, clnd\_)}}\quad\textit{functoid}\index{cc-line(clnp_, clnd_)@\texttt{cc-line(clnp\_, clnd\_)}}\\
reads: the straight line through \$1 in the direction \$2 \\
definition:
\begin{quote}\ttfamily\footnotesize\raggedright\noindent
\{clnz\_ in cc: imag-part(\allowbreak{}(\allowbreak{}clnz\_ - clnp\_) * conjugate(\allowbreak{}clnd\_)) = 0\}
\end{quote}
characterized by \texttt{cc-line-unfold}  \textit{(proven)}:
\begin{quote}\ttfamily\footnotesize\raggedright\noindent
forall(\allowbreak{}[clnp\_, clnd\_], cc-line(\allowbreak{}clnp\_, clnd\_) == \{clnz\_ in cc: imag-part(\allowbreak{}(\allowbreak{}clnz\_ - clnp\_) * conjugate(\allowbreak{}clnd\_)) = 0\})
\end{quote}
mentioned by 11 result(s)

\smallskip

\noindent\textbf{\texttt{cc-log(clgw\_)}}\quad\textit{functoid}\index{cc-log(clgw_)@\texttt{cc-log(clgw\_)}}\\
reads: the principal logarithm of \$1 \\
definition:
\begin{quote}\ttfamily\footnotesize\raggedright\noindent
iota(\allowbreak{}clgz\_, clgz\_ in cc and -pi < imag-part(\allowbreak{}clgz\_) and imag-part(\allowbreak{}clgz\_) < pi and cc-exp(\allowbreak{}clgz\_) = clgw\_)
\end{quote}
characterized by \texttt{cc-log-exp}  \textit{(proven)}:
\begin{quote}\ttfamily\footnotesize\raggedright\noindent
forall(\allowbreak{}[clgz\_ in cc], -pi < imag-part(\allowbreak{}clgz\_) implies imag-part(\allowbreak{}clgz\_) < pi implies cc-log(\allowbreak{}cc-exp(\allowbreak{}clgz\_)) = clgz\_)
\end{quote}
also: \texttt{cc-log-mul} \texttt{cc-log-one} \texttt{cc-log-char} \texttt{cc-log-real} \texttt{cc-log-spec} \texttt{cc-log-deriv} \texttt{cc-log-polar} \texttt{cc-log-scale} ... (26 in all) \\
mentioned by 74 result(s)

\smallskip

\noindent\textbf{\texttt{cc-matrices}}\quad\textit{constant}\index{cc-matrices@\texttt{cc-matrices}}\\
reads: the set of 2-by-2 real matrices of the form [[a, b], [-b, a]] \\
defined by \texttt{cc-matrices-def}

\smallskip

\noindent\textbf{\texttt{cc-matrix-of(a, b)}}\quad\textit{functoid}\index{cc-matrix-of(a, b)@\texttt{cc-matrix-of(a, b)}}\\
reads: the 2-by-2 matrix with rows [\$1, \$2] and [-\$2, \$1] \\
definition:
\begin{quote}\ttfamily\footnotesize\raggedright\noindent
[[a, b], [-b, a]]
\end{quote}
characterized by \texttt{cc-matrix-of-add}  \textit{(proven)}:
\begin{quote}\ttfamily\footnotesize\raggedright\noindent
forall(\allowbreak{}[a in rr, b in rr, c in rr, d in rr], matadd(\allowbreak{}normed-field-as-commutative-ring(\allowbreak{}rr-normed-field), cc-matrix-of(\allowbreak{}a, b), cc-matrix-of(\allowbreak{}c, d)) = cc-matrix-of(\allowbreak{}a + c, b + d))
\end{quote}
also: \texttt{cc-matrix-of-mul} \texttt{cc-matrix-of-scale} \texttt{cc-matrix-of-in-mat} \texttt{cc-matrix-of-entries} \texttt{cc-matrix-of-injective} \\
mentioned by 15 result(s)

\smallskip

\noindent\textbf{\texttt{cc-matrix-rep}}\quad\textit{constant}\index{cc-matrix-rep@\texttt{cc-matrix-rep}}\\
reads: the representation of a complex number a + b i as the matrix [[a, b], [-b, a]] \\
defined by \texttt{cc-matrix-rep-def} \\
characterized by \texttt{cc-matrix-rep-add}  \textit{(proven)}:
\begin{quote}\ttfamily\footnotesize\raggedright\noindent
forall(\allowbreak{}[z in cc, w in cc], cc-matrix-rep(\allowbreak{}z + w) = matadd(\allowbreak{}normed-field-as-commutative-ring(\allowbreak{}rr-normed-field), cc-matrix-rep(\allowbreak{}z), cc-matrix-rep(\allowbreak{}w)))
\end{quote}
also: \texttt{cc-matrix-rep-mul} \texttt{cc-matrix-rep-one} \texttt{cc-matrix-rep-apply} \texttt{cc-matrix-rep-real-mul} \texttt{cc-matrix-representation} \\
mentioned by 11 result(s)

\smallskip

\noindent\textbf{\texttt{cc-mid(cmx\_, cmy\_)}}\quad\textit{functoid}\index{cc-mid(cmx_, cmy_)@\texttt{cc-mid(cmx\_, cmy\_)}}\\
reads: the midpoint of \$1 and \$2 \\
definition:
\begin{quote}\ttfamily\footnotesize\raggedright\noindent
(\allowbreak{}cmx\_ + cmy\_) * recip(\allowbreak{}2)
\end{quote}
characterized by \texttt{cc-mid-affine}  \textit{(proven)}:
\begin{quote}\ttfamily\footnotesize\raggedright\noindent
forall(\allowbreak{}[goux\_ in cc, gouy\_ in cc], cc-mid(\allowbreak{}goux\_, gouy\_) = goux\_ + num(\allowbreak{}1/2) * (\allowbreak{}gouy\_ - goux\_) and cc-mid(\allowbreak{}goux\_, gouy\_) = gouy\_ + num(\allowbreak{}1/2) * (\allowbreak{}goux\_ - gouy\_))
\end{quote}
also: \texttt{cc-mid-coords} \texttt{cc-mid-unfold} \texttt{cc-mid-in-seg-trace} \\
mentioned by 16 result(s)

\smallskip

\noindent\textbf{\texttt{cc-ms}}\quad\textit{instance}\index{cc-ms@\texttt{cc-ms}}\\
an instance of metric-space; predicate \texttt{is-cc-ms}; same shape as metric-space

\smallskip

\noindent\textbf{\texttt{cc-nag}}\quad\textit{instance}\index{cc-nag@\texttt{cc-nag}}\\
an instance of normed-ag; predicate \texttt{is-cc-nag}; same shape as normed-ag

\smallskip

\noindent\textbf{\texttt{cc-normed-field}}\quad\textit{instance}\index{cc-normed-field@\texttt{cc-normed-field}}\\
an instance of normed-field; predicate \texttt{is-cc-normed-field}; same shape as normed-field

\smallskip

\noindent\textbf{\texttt{cc-of-pair}}\quad\textit{defined-fn}\index{cc-of-pair@\texttt{cc-of-pair}}\\
defined by \texttt{cc-of-pair-def} \\
characterized by \texttt{cc-of-pair-add}  \textit{(proven)}:
\begin{quote}\ttfamily\footnotesize\raggedright\noindent
forall(\allowbreak{}[a1\_ in rr, b1\_ in rr, a2\_ in rr, b2\_ in rr], cc-of-pair(\allowbreak{}[a1\_ + a2\_, b1\_ + b2\_]) = cc-of-pair(\allowbreak{}[a1\_, b1\_]) + cc-of-pair(\allowbreak{}[a2\_, b2\_]))
\end{quote}
also: \texttt{cc-of-pair-apply} \texttt{cc-of-pair-real-mul} \texttt{cc-of-pair-of-coords} \\
mentioned by 10 result(s)

\smallskip

\noindent\textbf{\texttt{cc-partial-prod}}\quad\textit{functoid}\index{cc-partial-prod@\texttt{cc-partial-prod}}\\
reads: the product of the first \$2 values of \$1 \\
defined by \texttt{cc-partial-prod-zero} \texttt{cc-partial-prod-succ} \\
characterized by \texttt{cc-partial-prod-succ}  \textit{(definitional)}:
\begin{quote}\ttfamily\footnotesize\raggedright\noindent
forall(\allowbreak{}[cpf\_, cpn\_ in nn], cc-partial-prod(\allowbreak{}cpf\_, succ(\allowbreak{}cpn\_)) == cc-partial-prod(\allowbreak{}cpf\_, cpn\_) * cpf\_(\allowbreak{}cpn\_))
\end{quote}
also: \texttt{cc-partial-prod-zero} \texttt{cc-partial-prod-congr} \texttt{cc-partial-prod-in-cc} \texttt{cc-partial-prod-zero-factor} \\
mentioned by 11 result(s)

\smallskip

\noindent\textbf{\texttt{cc-series-converges(f)}}\quad\textit{predicate}\index{cc-series-converges(f)@\texttt{cc-series-converges(f)}}\\
reads: the complex series \$1 converges \\
defined by \texttt{cc-series-converges} \\
characterized by \texttt{cc-series-converges}  \textit{(definitional)}:
\begin{quote}\ttfamily\footnotesize\raggedright\noindent
forall(\allowbreak{}[f], cc-series-converges(\allowbreak{}f) iff converges(\allowbreak{}cc-ms, vnb-lambda(\allowbreak{}k, nn, cc-series-partial-sum(\allowbreak{}f, k))))
\end{quote}
mentioned by 24 result(s)

\smallskip

\noindent\textbf{\texttt{cc-series-converges-to(f, l)}}\quad\textit{predicate}\index{cc-series-converges-to(f, l)@\texttt{cc-series-converges-to(f, l)}}\\
reads: the complex series \$1 converges to \$2 \\
defined by \texttt{cc-series-converges-to} \\
characterized by \texttt{cc-series-converges-to}  \textit{(definitional)}:
\begin{quote}\ttfamily\footnotesize\raggedright\noindent
forall(\allowbreak{}[f, l], cc-series-converges-to(\allowbreak{}f, l) iff converges-to(\allowbreak{}cc-ms, vnb-lambda(\allowbreak{}k, nn, cc-series-partial-sum(\allowbreak{}f, k)), l))
\end{quote}
mentioned by 16 result(s)

\smallskip

\noindent\textbf{\texttt{cc-series-limit(f)}}\quad\textit{functoid}\index{cc-series-limit(f)@\texttt{cc-series-limit(f)}}\\
reads: the sum of the complex series \$1 \\
definition:
\begin{quote}\ttfamily\footnotesize\raggedright\noindent
iota(\allowbreak{}cpsl\_, cc-series-converges-to(\allowbreak{}f, cpsl\_))
\end{quote}
characterized by \texttt{cc-series-limit-in-cc}  \textit{(proven)}:
\begin{quote}\ttfamily\footnotesize\raggedright\noindent
forall(\allowbreak{}[f in fun(\allowbreak{}nn, cc)], cc-series-converges(\allowbreak{}f) implies cc-series-limit(\allowbreak{}f) in cc)
\end{quote}
also: \texttt{cc-series-limit-linear} \texttt{cc-series-limit-conjugate} \texttt{cc-series-limit-converges-to} \texttt{cc-series-limit-of-converges-to} \\
mentioned by 72 result(s)

\smallskip

\noindent\textbf{\texttt{cc-series-partial-sum(f, k)}}\quad\textit{functoid}\index{cc-series-partial-sum(f, k)@\texttt{cc-series-partial-sum(f, k)}}\\
reads: the \$2-th partial sum of the complex series \$1 \\
definition:
\begin{quote}\ttfamily\footnotesize\raggedright\noindent
sum-ag(\allowbreak{}normed-field-additive-ag(\allowbreak{}cc-normed-field), f, k)
\end{quote}
characterized by \texttt{cc-series-partial-sum-succ}  \textit{(proven)}:
\begin{quote}\ttfamily\footnotesize\raggedright\noindent
forall(\allowbreak{}[f, k in nn], cc-series-partial-sum(\allowbreak{}f, k) in cc implies f(\allowbreak{}k) in cc implies cc-series-partial-sum(\allowbreak{}f, succ(\allowbreak{}k)) == cc-series-partial-sum(\allowbreak{}f, k) + f(\allowbreak{}k))
\end{quote}
also: \texttt{cc-series-partial-sum-zero} \texttt{cc-series-partial-sum-in-cc} \texttt{cc-series-partial-sum-unfold} \texttt{cc-series-partial-sum-lincomb} \texttt{cc-series-partial-sum-magnitude-le} \\
mentioned by 136 result(s)

\smallskip

\noindent\textbf{\texttt{cc-sin(z)}}\quad\textit{functoid}\index{cc-sin(z)@\texttt{cc-sin(z)}}\\
reads: sin \$1 \\
definition:
\begin{quote}\ttfamily\footnotesize\raggedright\noindent
cc-series-limit(\allowbreak{}vnb-lambda(\allowbreak{}cek\_, nn, (\allowbreak{}num(\allowbreak{}-1) \textasciicircum{} cek\_ * recip(\allowbreak{}factorial(\allowbreak{}succ(\allowbreak{}2 * cek\_)))) * z \textasciicircum{} succ(\allowbreak{}2 * cek\_)))
\end{quote}
characterized by \texttt{cc-sin-pi}  \textit{(proven)}:
\begin{quote}\ttfamily\footnotesize\raggedright\noindent
cc-sin(\allowbreak{}pi) = 0
\end{quote}
also: \texttt{cc-sin-add} \texttt{cc-sin-exp} \texttt{cc-sin-odd} \texttt{cc-sin-deriv} \texttt{cc-sin-two-pi} \texttt{cc-sin-unfold} \texttt{cc-sin-half-pi} \texttt{cc-sin-pos-half} ... (12 in all) \\
mentioned by 34 result(s)

\smallskip

\noindent\textbf{\texttt{cc-sn-module}}\quad\textit{constant}\index{cc-sn-module@\texttt{cc-sn-module}}\\
reads: the complex numbers as a vector space over themselves \\
defined by \texttt{cc-sn-module-def}

\smallskip

\noindent\textbf{\texttt{cc-sn-pf}}\quad\textit{constant}\index{cc-sn-pf@\texttt{cc-sn-pf}}\\
reads: the absolute value as a family of seminorms on the complex numbers \\
defined by \texttt{cc-sn-pf-def}

\smallskip

\noindent\textbf{\texttt{cc-sn-top}}\quad\textit{constant}\index{cc-sn-top@\texttt{cc-sn-top}}\\
reads: the complex numbers with the topology of the absolute value \\
defined by \texttt{cc-sn-top-def}

\smallskip

\noindent\textbf{\texttt{cc-sup-norm(u, cpsxs\_)}}\quad\textit{functoid}\index{cc-sup-norm(u, cpsxs_)@\texttt{cc-sup-norm(u, cpsxs\_)}}\\
reads: the supremum norm of \$1 on \$2 \\
definition:
\begin{quote}\ttfamily\footnotesize\raggedright\noindent
esup(\allowbreak{}image(\allowbreak{}vnb-lambda(\allowbreak{}cpsx\_, cpsxs\_, magnitude(\allowbreak{}u(\allowbreak{}cpsx\_))), cpsxs\_))
\end{quote}
characterized by \texttt{cc-sup-norm-bound}  \textit{(proven)}:
\begin{quote}\ttfamily\footnotesize\raggedright\noindent
forall(\allowbreak{}[u, cpsxs\_, cpsm\_ in rr], 0 <= cpsm\_ implies forall(\allowbreak{}[cpsx\_ in cpsxs\_], u(\allowbreak{}cpsx\_) in cc) implies forall(\allowbreak{}[cpsx\_ in cpsxs\_], magnitude(\allowbreak{}u(\allowbreak{}cpsx\_)) <= cpsm\_) implies cc-sup-norm(\allowbreak{}u, cpsxs\_) in rr and 0 <= cc-sup-norm(\allowbreak{}u, cpsxs\_) and cc-sup-norm(\allowbreak{}u, cpsxs\_) <= cpsm\_)
\end{quote}
also: \texttt{cc-sup-norm-upper} \\
mentioned by 17 result(s)

\smallskip

\noindent\textbf{\texttt{cc2-sphere}}\quad\textit{constant}\index{cc2-sphere@\texttt{cc2-sphere}}\\
reads: the unit sphere S\textasciicircum{}3 of C\textasciicircum{}2 \\
defined by \texttt{cc2-sphere-def}

\smallskip

\noindent\textbf{\texttt{cchs}}\quad\textit{functoid}\index{cchs@\texttt{cchs}}\\
reads: the existential chain of \$1 in the last \$5 of the \$3 entries of \$2, with the innermost list \$4 \\
defined by \texttt{cchs-zero} \texttt{cchs-succ} \\
characterized by \texttt{cchs-succ}  \textit{(definitional)}:
\begin{quote}\ttfamily\footnotesize\raggedright\noindent
forall(\allowbreak{}[m57x\_, m57a\_, m57n\_, m57l\_, m57k\_ in nn], cchs(\allowbreak{}m57x\_, m57a\_, m57n\_, m57l\_, succ(\allowbreak{}m57k\_)) == enc-ex(\allowbreak{}enc-and(\allowbreak{}enc-in(\allowbreak{}ebvar(\allowbreak{}0), nth(\allowbreak{}m57n\_ - m57k\_, m57a\_)), cchs(\allowbreak{}m57x\_, m57a\_, m57n\_, m57l\_, m57k\_))))
\end{quote}
also: \texttt{cchs-zero} \texttt{cchs-bwd} \texttt{cchs-fwd} \texttt{cchs-idx} \texttt{cchs-expr} \texttt{cchs-open} \\
mentioned by 11 result(s)

\smallskip

\noindent\textbf{\texttt{ccint(a, b)}}\quad\textit{functoid}\index{ccint(a, b)@\texttt{ccint(a, b)}}\\
reads: the closed real interval [\$1, \$2] \\
definition:
\begin{quote}\ttfamily\footnotesize\raggedright\noindent
\{x in rr: a <= x and x <= b\}
\end{quote}
characterized by \texttt{has-deriv-on-ccint}  \textit{(definitional)}:
\begin{quote}\ttfamily\footnotesize\raggedright\noindent
forall(\allowbreak{}[hdcg\_, hdcd\_, hdca\_, hdcb\_], has-deriv-on-ccint(\allowbreak{}hdcg\_, hdcd\_, hdca\_, hdcb\_) iff hdca\_ in rr and hdcb\_ in rr and hdca\_ < hdcb\_ and hdcg\_ in fun(\allowbreak{}ccint(\allowbreak{}hdca\_, hdcb\_), rr) and hdcd\_ in fun(\allowbreak{}ccint(\allowbreak{}hdca\_, hdcb\_), rr) and forall(\allowbreak{}[hdct\_ in ooint(\allowbreak{}hdca\_, hdcb\_)], has-deriv-at(\allowbreak{}hdcg\_, hdct\_, hdcd\_(\allowbreak{}hdct\_))) and has-right-deriv-at(\allowbreak{}hdcg\_, hdca\_, hdcd\_(\allowbreak{}hdca\_)) and has-left-deriv-at(\allowbreak{}hdcg\_, hdcb\_, hdcd\_(\allowbreak{}hdcb\_)))
\end{quote}
also: \texttt{ccint-creep} \texttt{ccint-parts} \texttt{clamp-in-ccint} \texttt{ccint-affine-in} \texttt{ccint-elt-in-rr} \texttt{ccint-subset-rr} \texttt{translate-ccint} \texttt{ccint-membership} ... (65 in all) \\
mentioned by 563 result(s)

\smallskip

\noindent\textbf{\texttt{ccont(kxt\_)}}\quad\textit{functoid}\index{ccont(kxt_)@\texttt{ccont(kxt\_)}}\\
reads: the continuous complex functions on \$1 \\
definition:
\begin{quote}\ttfamily\footnotesize\raggedright\noindent
\{kxf\_ in fun(\allowbreak{}pts(\allowbreak{}kxt\_), cc): is-hom-top-space(\allowbreak{}kxt\_, metric-top(\allowbreak{}cc-ms), kxf\_)\}
\end{quote}
characterized by \texttt{ccont-intro}  \textit{(proven)}:
\begin{quote}\ttfamily\footnotesize\raggedright\noindent
forall(\allowbreak{}[t, f], is-top-space(\allowbreak{}t) implies f in fun(\allowbreak{}pts(\allowbreak{}t), cc) implies forall(\allowbreak{}[p in pts(\allowbreak{}t), eps], pos-rr(\allowbreak{}eps) implies forsome(\allowbreak{}[u in opens(\allowbreak{}t)], p in u and forall(\allowbreak{}[q in u], magnitude(\allowbreak{}f(\allowbreak{}p) - f(\allowbreak{}q)) < eps))) implies f in ccont(\allowbreak{}t))
\end{quote}
also: \texttt{ccont-local} \texttt{ccont-parts} \texttt{ccont-act-in} \texttt{ccont-is-set} \texttt{ccont-unfold} \texttt{ccont-act-add} \texttt{ccont-act-mul} \texttt{ccont-act-one} ... (107 in all) \\
mentioned by 163 result(s)

\smallskip

\noindent\textbf{\texttt{ccont-act(kxt\_)}}\quad\textit{functoid}\index{ccont-act(kxt_)@\texttt{ccont-act(kxt\_)}}\\
reads: pointwise scalar multiplication on C(\$1) \\
definition:
\begin{quote}\ttfamily\footnotesize\raggedright\noindent
vnb-lambda(\allowbreak{}[kxc\_, kxf\_], cartesian(\allowbreak{}cc, ccont(\allowbreak{}kxt\_)), vnb-lambda(\allowbreak{}kxp\_, pts(\allowbreak{}kxt\_), kxc\_ * kxf\_(\allowbreak{}kxp\_)))
\end{quote}
characterized by \texttt{ccont-act-in}  \textit{(proven)}:
\begin{quote}\ttfamily\footnotesize\raggedright\noindent
forall(\allowbreak{}[t], is-top-space(\allowbreak{}t) implies forall(\allowbreak{}[c in cc, f in ccont(\allowbreak{}t)], (\allowbreak{}ccont-act(\allowbreak{}t))(\allowbreak{}c, f) in ccont(\allowbreak{}t)))
\end{quote}
also: \texttt{ccont-act-add} \texttt{ccont-act-mul} \texttt{ccont-act-one} \texttt{ccont-act-val} \texttt{ccont-act-type} \texttt{ccont-act-vadd} \\
mentioned by 55 result(s)

\smallskip

\noindent\textbf{\texttt{ccont-algebra(kxt\_)}}\quad\textit{functoid}\index{ccont-algebra(kxt_)@\texttt{ccont-algebra(kxt\_)}}\\
reads: the C*-algebra of continuous complex functions on \$1 \\
definition:
\begin{quote}\ttfamily\footnotesize\raggedright\noindent
[normed-field-as-commutative-ring(\allowbreak{}cc-normed-field), ccont(\allowbreak{}kxt\_), ccont-vadd(\allowbreak{}kxt\_), ccont-vzero(\allowbreak{}kxt\_), ccont-vneg(\allowbreak{}kxt\_), ccont-act(\allowbreak{}kxt\_), ccont-vnrm(\allowbreak{}kxt\_), ccont-vmul(\allowbreak{}kxt\_), ccont-vone(\allowbreak{}kxt\_), ccont-vstar(\allowbreak{}kxt\_)]
\end{quote}
characterized by \texttt{ccont-algebra-act}  \textit{(proven)}:
\begin{quote}\ttfamily\footnotesize\raggedright\noindent
forall(\allowbreak{}[t], is-compact-t(\allowbreak{}t) implies act(\allowbreak{}ccont-algebra(\allowbreak{}t)) == ccont-act(\allowbreak{}t))
\end{quote}
also: \texttt{ccont-algebra-uba} \texttt{ccont-algebra-vec} \texttt{ccont-algebra-scal} \texttt{ccont-algebra-vadd} \texttt{ccont-algebra-vmul} \texttt{ccont-algebra-vneg} \texttt{ccont-algebra-vnrm} \texttt{ccont-algebra-vone} ... (12 in all) \\
mentioned by 51 result(s)

\smallskip

\noindent\textbf{\texttt{ccont-vadd(kxt\_)}}\quad\textit{functoid}\index{ccont-vadd(kxt_)@\texttt{ccont-vadd(kxt\_)}}\\
reads: pointwise addition on C(\$1) \\
definition:
\begin{quote}\ttfamily\footnotesize\raggedright\noindent
vnb-lambda(\allowbreak{}[kxf\_, kxg\_], cartesian(\allowbreak{}ccont(\allowbreak{}kxt\_), ccont(\allowbreak{}kxt\_)), vnb-lambda(\allowbreak{}kxp\_, pts(\allowbreak{}kxt\_), kxf\_(\allowbreak{}kxp\_) + kxg\_(\allowbreak{}kxp\_)))
\end{quote}
characterized by \texttt{ccont-vadd-in}  \textit{(proven)}:
\begin{quote}\ttfamily\footnotesize\raggedright\noindent
forall(\allowbreak{}[t], is-top-space(\allowbreak{}t) implies forall(\allowbreak{}[f in ccont(\allowbreak{}t), g in ccont(\allowbreak{}t)], (\allowbreak{}ccont-vadd(\allowbreak{}t))(\allowbreak{}f, g) in ccont(\allowbreak{}t)))
\end{quote}
also: \texttt{ccont-vadd-inv} \texttt{ccont-vadd-val} \texttt{ccont-vadd-comm} \texttt{ccont-vadd-type} \texttt{ccont-vadd-assoc} \texttt{ccont-vadd-ident} \\
mentioned by 54 result(s)

\smallskip

\noindent\textbf{\texttt{ccont-vmul(kxt\_)}}\quad\textit{functoid}\index{ccont-vmul(kxt_)@\texttt{ccont-vmul(kxt\_)}}\\
reads: pointwise multiplication on C(\$1) \\
definition:
\begin{quote}\ttfamily\footnotesize\raggedright\noindent
vnb-lambda(\allowbreak{}[kxf\_, kxg\_], cartesian(\allowbreak{}ccont(\allowbreak{}kxt\_), ccont(\allowbreak{}kxt\_)), vnb-lambda(\allowbreak{}kxp\_, pts(\allowbreak{}kxt\_), kxf\_(\allowbreak{}kxp\_) * kxg\_(\allowbreak{}kxp\_)))
\end{quote}
characterized by \texttt{ccont-vmul-in}  \textit{(proven)}:
\begin{quote}\ttfamily\footnotesize\raggedright\noindent
forall(\allowbreak{}[t], is-top-space(\allowbreak{}t) implies forall(\allowbreak{}[f in ccont(\allowbreak{}t), g in ccont(\allowbreak{}t)], (\allowbreak{}ccont-vmul(\allowbreak{}t))(\allowbreak{}f, g) in ccont(\allowbreak{}t)))
\end{quote}
also: \texttt{ccont-vmul-val} \texttt{ccont-vmul-comm} \texttt{ccont-vmul-type} \texttt{ccont-vmul-assoc} \texttt{ccont-vmul-act-left} \texttt{ccont-vmul-act-right} \texttt{ccont-vmul-unit-left} \texttt{ccont-vmul-vadd-left} ... (10 in all) \\
mentioned by 59 result(s)

\smallskip

\noindent\textbf{\texttt{ccont-vneg(kxt\_)}}\quad\textit{functoid}\index{ccont-vneg(kxt_)@\texttt{ccont-vneg(kxt\_)}}\\
reads: pointwise negation on C(\$1) \\
definition:
\begin{quote}\ttfamily\footnotesize\raggedright\noindent
vnb-lambda(\allowbreak{}kxf\_, ccont(\allowbreak{}kxt\_), vnb-lambda(\allowbreak{}kxp\_, pts(\allowbreak{}kxt\_), -kxf\_(\allowbreak{}kxp\_)))
\end{quote}
characterized by \texttt{ccont-vneg-in}  \textit{(proven)}:
\begin{quote}\ttfamily\footnotesize\raggedright\noindent
forall(\allowbreak{}[t], is-top-space(\allowbreak{}t) implies forall(\allowbreak{}[f in ccont(\allowbreak{}t)], (\allowbreak{}ccont-vneg(\allowbreak{}t))(\allowbreak{}f) in ccont(\allowbreak{}t)))
\end{quote}
also: \texttt{ccont-vneg-val} \texttt{ccont-vneg-type} \\
mentioned by 32 result(s)

\smallskip

\noindent\textbf{\texttt{ccont-vnrm(kxt\_)}}\quad\textit{functoid}\index{ccont-vnrm(kxt_)@\texttt{ccont-vnrm(kxt\_)}}\\
reads: the supremum norm on C(\$1) \\
definition:
\begin{quote}\ttfamily\footnotesize\raggedright\noindent
vnb-lambda(\allowbreak{}kxf\_, ccont(\allowbreak{}kxt\_), sup-norm-t(\allowbreak{}kxt\_, kxf\_))
\end{quote}
characterized by \texttt{ccont-vnrm-act}  \textit{(proven)}:
\begin{quote}\ttfamily\footnotesize\raggedright\noindent
forall(\allowbreak{}[t], is-compact-t(\allowbreak{}t) implies forall(\allowbreak{}[r\_ in carr(\allowbreak{}normed-field-as-commutative-ring(\allowbreak{}cc-normed-field)), x\_ in ccont(\allowbreak{}t)], (\allowbreak{}ccont-vnrm(\allowbreak{}t))(\allowbreak{}(\allowbreak{}ccont-act(\allowbreak{}t))(\allowbreak{}r\_, x\_)) = magnitude(\allowbreak{}r\_) * (\allowbreak{}ccont-vnrm(\allowbreak{}t))(\allowbreak{}x\_)))
\end{quote}
also: \texttt{ccont-vnrm-val} \texttt{ccont-vnrm-type} \texttt{ccont-vnrm-vone} \texttt{ccont-vnrm-cstar} \texttt{ccont-vnrm-vstar} \texttt{ccont-vnrm-nonneg} \texttt{ccont-vnrm-submult} \texttt{ccont-vnrm-triangle} ... (9 in all) \\
mentioned by 45 result(s)

\smallskip

\noindent\textbf{\texttt{ccont-vone(kxt\_)}}\quad\textit{functoid}\index{ccont-vone(kxt_)@\texttt{ccont-vone(kxt\_)}}\\
reads: the constant function 1 on \$1 \\
definition:
\begin{quote}\ttfamily\footnotesize\raggedright\noindent
vnb-lambda(\allowbreak{}kxp\_, pts(\allowbreak{}kxt\_), 1)
\end{quote}
characterized by \texttt{ccont-vone-val}  \textit{(proven)}:
\begin{quote}\ttfamily\footnotesize\raggedright\noindent
forall(\allowbreak{}[t], is-top-space(\allowbreak{}t) implies forall(\allowbreak{}[p in pts(\allowbreak{}t)], (\allowbreak{}ccont-vone(\allowbreak{}t))(\allowbreak{}p) = 1))
\end{quote}
also: \texttt{ccont-vone-type} \\
mentioned by 20 result(s)

\smallskip

\noindent\textbf{\texttt{ccont-vstar(kxt\_)}}\quad\textit{functoid}\index{ccont-vstar(kxt_)@\texttt{ccont-vstar(kxt\_)}}\\
reads: pointwise conjugation on C(\$1) \\
definition:
\begin{quote}\ttfamily\footnotesize\raggedright\noindent
vnb-lambda(\allowbreak{}kxf\_, ccont(\allowbreak{}kxt\_), vnb-lambda(\allowbreak{}kxp\_, pts(\allowbreak{}kxt\_), conjugate(\allowbreak{}kxf\_(\allowbreak{}kxp\_))))
\end{quote}
characterized by \texttt{ccont-vstar-in}  \textit{(proven)}:
\begin{quote}\ttfamily\footnotesize\raggedright\noindent
forall(\allowbreak{}[t], is-top-space(\allowbreak{}t) implies forall(\allowbreak{}[f in ccont(\allowbreak{}t)], (\allowbreak{}ccont-vstar(\allowbreak{}t))(\allowbreak{}f) in ccont(\allowbreak{}t)))
\end{quote}
also: \texttt{ccont-vstar-act} \texttt{ccont-vstar-val} \texttt{ccont-vstar-type} \texttt{ccont-vstar-vadd} \texttt{ccont-vstar-vmul} \texttt{ccont-vstar-vstar} \\
mentioned by 21 result(s)

\smallskip

\noindent\textbf{\texttt{ccont-vzero(kxt\_)}}\quad\textit{functoid}\index{ccont-vzero(kxt_)@\texttt{ccont-vzero(kxt\_)}}\\
reads: the zero function on \$1 \\
definition:
\begin{quote}\ttfamily\footnotesize\raggedright\noindent
vnb-lambda(\allowbreak{}kxp\_, pts(\allowbreak{}kxt\_), 0)
\end{quote}
characterized by \texttt{ccont-vzero-val}  \textit{(proven)}:
\begin{quote}\ttfamily\footnotesize\raggedright\noindent
forall(\allowbreak{}[t], is-top-space(\allowbreak{}t) implies forall(\allowbreak{}[p in pts(\allowbreak{}t)], (\allowbreak{}ccont-vzero(\allowbreak{}t))(\allowbreak{}p) = 0))
\end{quote}
also: \texttt{ccont-vzero-type} \\
mentioned by 41 result(s)

\smallskip

\noindent\textbf{\texttt{cdual-norm(baa\_, baf\_)}}\quad\textit{functoid}\index{cdual-norm(baa_, baf_)@\texttt{cdual-norm(baa\_, baf\_)}}\\
reads: the norm of the functional \$2 on \$1 \\
definition:
\begin{quote}\ttfamily\footnotesize\raggedright\noindent
esup(\allowbreak{}image(\allowbreak{}vnb-lambda(\allowbreak{}bdx\_, \{bdy\_ in vec(\allowbreak{}baa\_): (\allowbreak{}vnrm(\allowbreak{}baa\_))(\allowbreak{}bdy\_) <= 1\}, magnitude(\allowbreak{}baf\_(\allowbreak{}bdx\_))), \{bdy\_ in vec(\allowbreak{}baa\_): (\allowbreak{}vnrm(\allowbreak{}baa\_))(\allowbreak{}bdy\_) <= 1\}))
\end{quote}
mentioned by 5 result(s)

\smallskip

\noindent\textbf{\texttt{cdual-vec(baa\_)}}\quad\textit{functoid}\index{cdual-vec(baa_)@\texttt{cdual-vec(baa\_)}}\\
reads: the bounded complex-linear functionals on \$1 \\
definition:
\begin{quote}\ttfamily\footnotesize\raggedright\noindent
\{baf\_ in fun(\allowbreak{}vec(\allowbreak{}baa\_), cc): is-cc-linear-on(\allowbreak{}baa\_, baf\_) and forsome(\allowbreak{}[bam\_ in rr], forall(\allowbreak{}[bax\_ in vec(\allowbreak{}baa\_)], magnitude(\allowbreak{}baf\_(\allowbreak{}bax\_)) <= bam\_ * (\allowbreak{}vnrm(\allowbreak{}baa\_))(\allowbreak{}bax\_)))\}
\end{quote}
mentioned by 3 result(s)

\smallskip

\noindent\textbf{\texttt{ce}}\quad\textit{tactic}\index{ce@\texttt{ce}}\\
\texttt{(ce hyp k)} -- Cartesian elim: project the k-th component out of a CARTESIAN-membership assumption. \\
declared in \texttt{tactics-help.scm}

\smallskip

\noindent\textbf{\texttt{centre-set(s, r, f)}}\quad\textit{functoid}\index{centre-set(s, r, f)@\texttt{centre-set(s, r, f)}}\\
definition:
\begin{quote}\ttfamily\footnotesize\raggedright\noindent
image(\allowbreak{}vnb-lambda(\allowbreak{}b, f, choice(\allowbreak{}centres(\allowbreak{}s, b, r))), f)
\end{quote}
characterized by \texttt{centre-set-unfold}  \textit{(proven)}:
\begin{quote}\ttfamily\footnotesize\raggedright\noindent
forall(\allowbreak{}[s, r, f], centre-set(\allowbreak{}s, r, f) == image(\allowbreak{}vnb-lambda(\allowbreak{}b, f, choice(\allowbreak{}centres(\allowbreak{}s, b, r))), f))
\end{quote}
also: \texttt{centre-set-finite-guarded} \texttt{centre-set-contains-choice} \\
mentioned by 4 result(s)

\smallskip

\noindent\textbf{\texttt{centres(s, b, r)}}\quad\textit{functoid}\index{centres(s, b, r)@\texttt{centres(s, b, r)}}\\
definition:
\begin{quote}\ttfamily\footnotesize\raggedright\noindent
\{c in pts(\allowbreak{}s): ball(\allowbreak{}s, c, r) = b\}
\end{quote}
characterized by \texttt{centres-unfold}  \textit{(proven)}:
\begin{quote}\ttfamily\footnotesize\raggedright\noindent
forall(\allowbreak{}[s, b, r], centres(\allowbreak{}s, b, r) == \{c in pts(\allowbreak{}s): ball(\allowbreak{}s, c, r) = b\})
\end{quote}
also: \texttt{centres-ball-eq} \texttt{centres-mem-build} \texttt{centres-in-carrier} \\
mentioned by 10 result(s)

\smallskip

\noindent\textbf{\texttt{cert-closed(md12t\_, md12cg\_, md12cw\_, md12x\_)}}\quad\textit{predicate}\index{cert-closed(md12t_, md12cg_, md12cw_, md12x_)@\texttt{cert-closed(md12t\_, md12cg\_, md12cw\_, md12x\_)}}\\
reads: \$4 is closed under the steps certifying \$3 in \$2 over \$1 \\
defined by \texttt{cert-closed} \\
characterized by \texttt{cert-closed}  \textit{(definitional)}:
\begin{quote}\ttfamily\footnotesize\raggedright\noindent
forall(\allowbreak{}[md12t\_, md12cg\_, md12cw\_, md12x\_], cert-closed(\allowbreak{}md12t\_, md12cg\_, md12cw\_, md12x\_) iff forall(\allowbreak{}[md12e\_ in expr-set], cert-step(\allowbreak{}md12t\_, md12cg\_, md12cw\_, md12x\_, md12e\_) implies md12e\_ in md12x\_))
\end{quote}
mentioned by 5 result(s)

\smallskip

\noindent\textbf{\texttt{cert-set(md12t\_, md12cg\_, md12cw\_)}}\quad\textit{functoid}\index{cert-set(md12t_, md12cg_, md12cw_)@\texttt{cert-set(md12t\_, md12cg\_, md12cw\_)}}\\
reads: the sentences certifying \$3 in \$2 over \$1 \\
definition:
\begin{quote}\ttfamily\footnotesize\raggedright\noindent
\{md12zz\_ in expr-set: forall(\allowbreak{}[md12v\_ in power(\allowbreak{}expr-set)], cert-closed(\allowbreak{}md12t\_, md12cg\_, md12cw\_, md12v\_) implies md12zz\_ in md12v\_)\}
\end{quote}
characterized by \texttt{cert-set-least}  \textit{(proven)}:
\begin{quote}\ttfamily\footnotesize\raggedright\noindent
forall(\allowbreak{}[tt, g, w, md12xs\_ in power(\allowbreak{}expr-set)], cert-closed(\allowbreak{}tt, g, w, md12xs\_) implies forall(\allowbreak{}[x in cert-set(\allowbreak{}tt, g, w)], x in md12xs\_))
\end{quote}
also: \texttt{cert-set-unfold} \\
mentioned by 18 result(s)

\smallskip

\noindent\textbf{\texttt{cert-step(md12t\_, md12cg\_, md12cw\_, md12x\_, md12f\_)}}\quad\textit{predicate}\index{cert-step(md12t_, md12cg_, md12cw_, md12x_, md12f_)@\texttt{cert-step(md12t\_, md12cg\_, md12cw\_, md12x\_, md12f\_)}}\\
reads: \$5 certifies \$3 in \$2 over \$1 by one step into \$4 \\
defined by \texttt{cert-step} \\
characterized by \texttt{cert-step}  \textit{(definitional)}:
\begin{quote}\ttfamily\footnotesize\raggedright\noindent
forall(\allowbreak{}[md12t\_, md12cg\_, md12cw\_, md12x\_, md12f\_], cert-step(\allowbreak{}md12t\_, md12cg\_, md12cw\_, md12x\_, md12f\_) iff strict-atom(\allowbreak{}md12f\_, md12cw\_) or forsome(\allowbreak{}[md12f2\_ in true-reach(\allowbreak{}md12f\_)], md12f2\_ in md12x\_) or forsome(\allowbreak{}[md12b\_, md12y\_], md12f\_ = enc-ex(\allowbreak{}md12b\_) and md12y\_ in nn and not(\allowbreak{}md12y\_ in union(\allowbreak{}ctx-fv(\allowbreak{}ctx-add(\allowbreak{}md12cg\_, md12f\_)), expr-fv(\allowbreak{}enc-eq(\allowbreak{}md12cw\_, md12cw\_)))) and expr-open(\allowbreak{}md12b\_, 0, efvar(\allowbreak{}md12y\_)) in md12x\_) or forsome(\allowbreak{}[md12l\_, md12k\_, md12d\_, md12b2\_], md12f\_ = eapp(\allowbreak{}10, md12l\_) and nth(\allowbreak{}1, md12l\_) = econst(\allowbreak{}md12k\_) and pred-code(\allowbreak{}md12k\_) and [md12cg\_, md12d\_] in derivable(\allowbreak{}md12t\_) and enc-iff(\allowbreak{}md12f\_, md12b2\_) in inst-reach(\allowbreak{}md12d\_, md12l\_) and md12b2\_ in md12x\_) or forsome(\allowbreak{}[md12a\_, md12s\_, md12d3\_, md12b3\_], md12f\_ = enc-in(\allowbreak{}md12s\_, md12a\_) and [md12cg\_, md12d3\_] in derivable(\allowbreak{}md12t\_) and enc-iff(\allowbreak{}md12f\_, md12b3\_) in inst-reach(\allowbreak{}md12d3\_, [md12a\_, md12s\_]) and md12b3\_ in md12x\_))
\end{quote}
mentioned by 4 result(s)

\smallskip

\noindent\textbf{\texttt{cf-cons(b, c)}}\quad\textit{functoid}\index{cf-cons(b, c)@\texttt{cf-cons(b, c)}}\\
reads: the sequence \$1 followed by \$2 \\
definition:
\begin{quote}\ttfamily\footnotesize\raggedright\noindent
vnb-lambda(\allowbreak{}pfk\_, nn, b * 0 \textasciicircum{} pfk\_ + (\allowbreak{}1 - 0 \textasciicircum{} pfk\_) * c(\allowbreak{}nn-minus(\allowbreak{}pfk\_, 1)))
\end{quote}
characterized by \texttt{cf-cons-succ}  \textit{(proven)}:
\begin{quote}\ttfamily\footnotesize\raggedright\noindent
forall(\allowbreak{}[pfb\_ in cc, pfc\_ in fun(\allowbreak{}nn, cc), pfj\_ in nn], (\allowbreak{}cf-cons(\allowbreak{}pfb\_, pfc\_))(\allowbreak{}succ(\allowbreak{}pfj\_)) = pfc\_(\allowbreak{}pfj\_))
\end{quote}
also: \texttt{cf-cons-zero} \texttt{cf-cons-in-fun} \texttt{cf-cons-unfold} \\
mentioned by 10 result(s)

\smallskip

\noindent\textbf{\texttt{char-converges(c5k\_, c5s\_, c5c\_)}}\quad\textit{predicate}\index{char-converges(c5k_, c5s_, c5c_)@\texttt{char-converges(c5k\_, c5s\_, c5c\_)}}\\
reads: \$2 converges to \$3 uniformly on the compact subsets of \$1 \\
defined by \texttt{char-converges} \\
characterized by \texttt{char-converges}  \textit{(definitional)}:
\begin{quote}\ttfamily\footnotesize\raggedright\noindent
forall(\allowbreak{}[c5k\_, c5s\_, c5c\_], char-converges(\allowbreak{}c5k\_, c5s\_, c5c\_) iff forall(\allowbreak{}[c5e\_], c5e\_ subset carr(\allowbreak{}c5k\_) implies is-compact(\allowbreak{}subspace-ms(\allowbreak{}nf-metric-space(\allowbreak{}c5k\_), c5e\_)) implies forall(\allowbreak{}[c5p\_], pos-rr(\allowbreak{}c5p\_) implies forsome(\allowbreak{}[c5n\_ in nn], forall(\allowbreak{}[c5m\_ in nn], c5n\_ <= c5m\_ implies forall(\allowbreak{}[c5x\_ in c5e\_], magnitude(\allowbreak{}(\allowbreak{}c5s\_(\allowbreak{}c5m\_))(\allowbreak{}c5x\_) - c5c\_(\allowbreak{}c5x\_)) < c5p\_))))))
\end{quote}
mentioned by 5 result(s)

\smallskip

\noindent\textbf{\texttt{char-mat(f, a, n)}}\quad\textit{functoid}\index{char-mat(f, a, n)@\texttt{char-mat(f, a, n)}}\\
reads: the matrix x I - \$2 over \$1[x] \\
definition:
\begin{quote}\ttfamily\footnotesize\raggedright\noindent
matof(\allowbreak{}n, n, vnb-lambda(\allowbreak{}[gmi\_, gmj\_], cartesian(\allowbreak{}interval(\allowbreak{}1, n), interval(\allowbreak{}1, n)), (\allowbreak{}add(\allowbreak{}poly(\allowbreak{}f)))(\allowbreak{}if(\allowbreak{}gmi\_ = gmj\_, monomial(\allowbreak{}f, one(\allowbreak{}f), 1), zero(\allowbreak{}poly(\allowbreak{}f))), (\allowbreak{}neg(\allowbreak{}poly(\allowbreak{}f)))(\allowbreak{}monomial(\allowbreak{}f, entry(\allowbreak{}a, gmi\_, gmj\_), 0)))))
\end{quote}
characterized by \texttt{char-mat-entry-uba-as-ring}  \textit{(definitional)}:
\begin{quote}\ttfamily\footnotesize\raggedright\noindent
forall(\allowbreak{}[r], is-unital-banach-algebra(\allowbreak{}r) implies forall(\allowbreak{}[nv\_ in nn, am\_ in mat(\allowbreak{}nv\_, nv\_, vec(\allowbreak{}r))], char-mat(\allowbreak{}uba-as-ring(\allowbreak{}r), am\_, nv\_) in mat(\allowbreak{}nv\_, nv\_, carr(\allowbreak{}poly(\allowbreak{}uba-as-ring(\allowbreak{}r)))) and forall(\allowbreak{}[i\_ in interval(\allowbreak{}1, nv\_), j\_ in interval(\allowbreak{}1, nv\_)], entry(\allowbreak{}char-mat(\allowbreak{}uba-as-ring(\allowbreak{}r), am\_, nv\_), i\_, j\_) = (\allowbreak{}add(\allowbreak{}poly(\allowbreak{}uba-as-ring(\allowbreak{}r))))(\allowbreak{}if(\allowbreak{}i\_ = j\_, monomial(\allowbreak{}uba-as-ring(\allowbreak{}r), vone(\allowbreak{}r), 1), zero(\allowbreak{}poly(\allowbreak{}uba-as-ring(\allowbreak{}r)))), (\allowbreak{}neg(\allowbreak{}poly(\allowbreak{}uba-as-ring(\allowbreak{}r))))(\allowbreak{}monomial(\allowbreak{}uba-as-ring(\allowbreak{}r), entry(\allowbreak{}am\_, i\_, j\_), 0))))))
\end{quote}
also: \texttt{char-mat-entry-uba-as-ring-uiba-as-uba} \texttt{char-mat-entry} \texttt{char-mat-kills-basis} \texttt{char-mat-companion-row1} \texttt{char-mat-companion-minor11} \texttt{char-mat-companion-corner-det} \\
mentioned by 10 result(s)

\smallskip

\noindent\textbf{\texttt{char-poly(f, a, n)}}\quad\textit{functoid}\index{char-poly(f, a, n)@\texttt{char-poly(f, a, n)}}\\
reads: the characteristic polynomial of \$2 \\
definition:
\begin{quote}\ttfamily\footnotesize\raggedright\noindent
det(\allowbreak{}poly(\allowbreak{}f), n, char-mat(\allowbreak{}f, a, n))
\end{quote}
mentioned by 7 result(s)

\smallskip

\noindent\textbf{\texttt{characters(baa\_)}}\quad\textit{functoid}\index{characters(baa_)@\texttt{characters(baa\_)}}\\
reads: the characters of \$1 \\
definition:
\begin{quote}\ttfamily\footnotesize\raggedright\noindent
\{bac\_ in fun(\allowbreak{}vec(\allowbreak{}baa\_), cc): is-cc-linear-on(\allowbreak{}baa\_, bac\_) and forall(\allowbreak{}[bax\_ in vec(\allowbreak{}baa\_), bay\_ in vec(\allowbreak{}baa\_)], bac\_(\allowbreak{}(\allowbreak{}vmul(\allowbreak{}baa\_))(\allowbreak{}bax\_, bay\_)) = bac\_(\allowbreak{}bax\_) * bac\_(\allowbreak{}bay\_)) and bac\_(\allowbreak{}vone(\allowbreak{}baa\_)) = 1\}
\end{quote}
characterized by \texttt{characters-in-set}  \textit{(proven)}:
\begin{quote}\ttfamily\footnotesize\raggedright\noindent
forall(\allowbreak{}[a], is-unital-banach-algebra(\allowbreak{}a) implies characters(\allowbreak{}a) in set)
\end{quote}
also: \texttt{bcs-characters-view} \texttt{uiba-characters-view} \texttt{bas-characters-unfold} \texttt{ccont-characters-view} \texttt{ccx-characters-unfold} \texttt{characters-closed-in-product} \\
mentioned by 58 result(s)

\smallskip

\noindent\textbf{\texttt{choice}}\quad\textit{operator}\index{choice@\texttt{choice}}\\
a term-forming head \\
characterized by \texttt{choice-axiom}  \textit{(primitive)}:
\begin{quote}\ttfamily\footnotesize\raggedright\noindent
forall(\allowbreak{}[a], forsome(\allowbreak{}[x], x in a) implies choice(\allowbreak{}a) in a)
\end{quote}
also: \texttt{vitali-choice} \texttt{vitali-q-choice} \texttt{quotient-member-choice} \texttt{vitali-choice-set-finite} \texttt{centre-set-contains-choice} \texttt{vitali-q-choice-set-finite} \\
mentioned by 49 result(s)

\smallskip

\noindent\textbf{\texttt{choose(n, m)}}\quad\textit{functoid}\index{choose(n, m)@\texttt{choose(n, m)}}\\
definition:
\begin{quote}\ttfamily\footnotesize\raggedright\noindent
card(\allowbreak{}choose-set(\allowbreak{}n, m))
\end{quote}
characterized by \texttt{choose-n-0}  \textit{(proven)}:
\begin{quote}\ttfamily\footnotesize\raggedright\noindent
forall(\allowbreak{}[n in nn], choose(\allowbreak{}n, 0) = succ(\allowbreak{}0))
\end{quote}
also: \texttt{choose-succ} \texttt{choose-in-nn} \texttt{choose-0-succ} \texttt{choose-diagonal} \texttt{choose-above-zero} \texttt{choose-card-unfold} \texttt{choose-times-factorial} \texttt{choose-times-factorial-all} \\
mentioned by 20 result(s)

\smallskip

\noindent\textbf{\texttt{choose-set(n, m)}}\quad\textit{functoid}\index{choose-set(n, m)@\texttt{choose-set(n, m)}}\\
definition:
\begin{quote}\ttfamily\footnotesize\raggedright\noindent
\{a in power(\allowbreak{}ord-segment(\allowbreak{}n)): card(\allowbreak{}a) = m\}
\end{quote}
characterized by \texttt{choose-set-split}  \textit{(proven)}:
\begin{quote}\ttfamily\footnotesize\raggedright\noindent
forall(\allowbreak{}[n\_ in nn, k\_ in nn], choose-set(\allowbreak{}succ(\allowbreak{}n\_), succ(\allowbreak{}k\_)) = union(\allowbreak{}choose-set(\allowbreak{}n\_, succ(\allowbreak{}k\_)), image(\allowbreak{}vnb-lambda(\allowbreak{}s\_, choose-set(\allowbreak{}n\_, k\_), union(\allowbreak{}s\_, pair(\allowbreak{}n\_, n\_))), choose-set(\allowbreak{}n\_, k\_))))
\end{quote}
also: \texttt{choose-set-widen} \texttt{choose-set-unfold} \texttt{choose-set-mem-iff} \texttt{choose-set-avoid-in} \texttt{choose-set-insert-in} \texttt{choose-set-remove-in} \texttt{choose-set-avoids-point} \texttt{choose-set-split-disjoint} ... (9 in all) \\
mentioned by 17 result(s)

\smallskip

\noindent\textbf{\texttt{ci}}\quad\textit{tactic}\index{ci@\texttt{ci}}\\
\texttt{(ci)} -- Cartesian intro: prove a CARTESIAN-product membership component-wise. \\
declared in \texttt{tactics-help.scm}

\smallskip

\noindent\textbf{\texttt{cips-cnvs(w)}}\quad\textit{functoid}\index{cips-cnvs(w)@\texttt{cips-cnvs(w)}}\\
reads: \$1 as a complex normed vector space \\
definition:
\begin{quote}\ttfamily\footnotesize\raggedright\noindent
cmod-cnvs(\allowbreak{}complex-inner-product-space-as-module(\allowbreak{}w), cips-vnrm(\allowbreak{}w))
\end{quote}
characterized by \texttt{cips-cnvs-vec}  \textit{(proven)}:
\begin{quote}\ttfamily\footnotesize\raggedright\noindent
forall(\allowbreak{}[w], is-complex-inner-product-space(\allowbreak{}w) implies vec(\allowbreak{}cips-cnvs(\allowbreak{}w)) == vec(\allowbreak{}w))
\end{quote}
also: \texttt{cips-cnvs-cact} \texttt{cips-cnvs-scal} \texttt{cips-cnvs-vadd} \texttt{cips-cnvs-vneg} \texttt{cips-cnvs-vnrm} \texttt{cips-cnvs-vzero} \texttt{cips-cnvs-is-cnvs} \texttt{cips-cnvs-jop-value} ... (11 in all) \\
mentioned by 42 result(s)

\smallskip

\noindent\textbf{\texttt{cips-j(w)}}\quad\textit{functoid}\index{cips-j(w)@\texttt{cips-j(w)}}\\
reads: multiplication by i on \$1 \\
definition:
\begin{quote}\ttfamily\footnotesize\raggedright\noindent
vnb-lambda(\allowbreak{}cxx\_, vec(\allowbreak{}w), (\allowbreak{}act(\allowbreak{}w))(\allowbreak{}1i, cxx\_))
\end{quote}
characterized by \texttt{cips-j-fun}  \textit{(proven)}:
\begin{quote}\ttfamily\footnotesize\raggedright\noindent
forall(\allowbreak{}[w], is-complex-inner-product-space(\allowbreak{}w) implies cips-j(\allowbreak{}w) in fun(\allowbreak{}vec(\allowbreak{}w), vec(\allowbreak{}w)))
\end{quote}
mentioned by 3 result(s)

\smallskip

\noindent\textbf{\texttt{cips-ract(w)}}\quad\textit{functoid}\index{cips-ract(w)@\texttt{cips-ract(w)}}\\
definition:
\begin{quote}\ttfamily\footnotesize\raggedright\noindent
vnb-lambda(\allowbreak{}[cxr\_, cxx\_], cartesian(\allowbreak{}rr, vec(\allowbreak{}w)), (\allowbreak{}act(\allowbreak{}w))(\allowbreak{}cxr\_, cxx\_))
\end{quote}
characterized by \texttt{cips-ract-fun}  \textit{(proven)}:
\begin{quote}\ttfamily\footnotesize\raggedright\noindent
forall(\allowbreak{}[w], is-complex-inner-product-space(\allowbreak{}w) implies cips-ract(\allowbreak{}w) in fun(\allowbreak{}cartesian(\allowbreak{}rr, vec(\allowbreak{}w)), vec(\allowbreak{}w)))
\end{quote}
mentioned by 7 result(s)

\smallskip

\noindent\textbf{\texttt{cips-real(w)}}\quad\textit{functoid}\index{cips-real(w)@\texttt{cips-real(w)}}\\
reads: \$1 as a real space with complex structure \\
definition:
\begin{quote}\ttfamily\footnotesize\raggedright\noindent
[normed-field-as-commutative-ring(\allowbreak{}rr-normed-field), vec(\allowbreak{}w), vadd(\allowbreak{}w), vzero(\allowbreak{}w), vneg(\allowbreak{}w), cips-ract(\allowbreak{}w), cips-rip(\allowbreak{}w), cips-j(\allowbreak{}w)]
\end{quote}
characterized by \texttt{cips-real-act}  \textit{(proven)}:
\begin{quote}\ttfamily\footnotesize\raggedright\noindent
forall(\allowbreak{}[w], is-complex-inner-product-space(\allowbreak{}w) implies act(\allowbreak{}cips-real(\allowbreak{}w)) == cips-ract(\allowbreak{}w))
\end{quote}
also: \texttt{cips-real-jop} \texttt{cips-real-rip} \texttt{cips-real-sip} \texttt{cips-real-vec} \texttt{cips-real-scal} \texttt{cips-real-vadd} \texttt{cips-real-vneg} \texttt{cips-real-vzero} ... (22 in all) \\
mentioned by 33 result(s)

\smallskip

\noindent\textbf{\texttt{cips-rip(w)}}\quad\textit{functoid}\index{cips-rip(w)@\texttt{cips-rip(w)}}\\
reads: the real part of the inner product of \$1 \\
definition:
\begin{quote}\ttfamily\footnotesize\raggedright\noindent
vnb-lambda(\allowbreak{}[cxx\_, cxy\_], cartesian(\allowbreak{}vec(\allowbreak{}w), vec(\allowbreak{}w)), real-part(\allowbreak{}(\allowbreak{}ip(\allowbreak{}w))(\allowbreak{}cxx\_, cxy\_)))
\end{quote}
characterized by \texttt{cips-rip-fun}  \textit{(proven)}:
\begin{quote}\ttfamily\footnotesize\raggedright\noindent
forall(\allowbreak{}[w], is-complex-inner-product-space(\allowbreak{}w) implies cips-rip(\allowbreak{}w) in fun(\allowbreak{}cartesian(\allowbreak{}vec(\allowbreak{}w), vec(\allowbreak{}w)), rr))
\end{quote}
also: \texttt{cips-rip-unfold} \\
mentioned by 9 result(s)

\smallskip

\noindent\textbf{\texttt{cips-rsip(w)}}\quad\textit{functoid}\index{cips-rsip(w)@\texttt{cips-rsip(w)}}\\
reads: the underlying real semi-inner-product space of \$1 \\
definition:
\begin{quote}\ttfamily\footnotesize\raggedright\noindent
[normed-field-as-commutative-ring(\allowbreak{}rr-normed-field), vec(\allowbreak{}w), vadd(\allowbreak{}w), vzero(\allowbreak{}w), vneg(\allowbreak{}w), cips-ract(\allowbreak{}w), cips-rip(\allowbreak{}w)]
\end{quote}
characterized by \texttt{cips-rsip-act}  \textit{(proven)}:
\begin{quote}\ttfamily\footnotesize\raggedright\noindent
forall(\allowbreak{}[w], is-complex-inner-product-space(\allowbreak{}w) implies act(\allowbreak{}cips-rsip(\allowbreak{}w)) == cips-ract(\allowbreak{}w))
\end{quote}
also: \texttt{cips-rsip-rip} \texttt{cips-rsip-vec} \texttt{cips-rsip-comm} \texttt{cips-rsip-scal} \texttt{cips-rsip-vadd} \texttt{cips-rsip-vneg} \texttt{cips-rsip-assoc} \texttt{cips-rsip-vzero} ... (30 in all) \\
mentioned by 49 result(s)

\smallskip

\noindent\textbf{\texttt{cips-vnrm(w)}}\quad\textit{functoid}\index{cips-vnrm(w)@\texttt{cips-vnrm(w)}}\\
reads: the norm of the inner product space \$1 \\
definition:
\begin{quote}\ttfamily\footnotesize\raggedright\noindent
vnb-lambda(\allowbreak{}cnx\_, vec(\allowbreak{}w), ip-norm(\allowbreak{}w, cnx\_))
\end{quote}
characterized by \texttt{cips-vnrm-fun}  \textit{(proven)}:
\begin{quote}\ttfamily\footnotesize\raggedright\noindent
forall(\allowbreak{}[w], is-complex-inner-product-space(\allowbreak{}w) implies cips-vnrm(\allowbreak{}w) in fun(\allowbreak{}vec(\allowbreak{}w), rr))
\end{quote}
also: \texttt{cips-vnrm-value} \\
mentioned by 6 result(s)

\smallskip

\noindent\textbf{\texttt{circle-deriv(wnc\_, wnr\_)}}\quad\textit{functoid}\index{circle-deriv(wnc_, wnr_)@\texttt{circle-deriv(wnc\_, wnr\_)}}\\
reads: the derivative of the circle of centre \$1 and radius \$2 \\
definition:
\begin{quote}\ttfamily\footnotesize\raggedright\noindent
vnb-lambda(\allowbreak{}wnt\_, ccint(\allowbreak{}0, 2 * pi), 1i * (\allowbreak{}wnr\_ * cc-exp(\allowbreak{}1i * wnt\_)))
\end{quote}
mentioned by 52 result(s)

\smallskip

\noindent\textbf{\texttt{circle-path(wnc\_, wnr\_)}}\quad\textit{functoid}\index{circle-path(wnc_, wnr_)@\texttt{circle-path(wnc\_, wnr\_)}}\\
reads: the circle of centre \$1 and radius \$2 \\
definition:
\begin{quote}\ttfamily\footnotesize\raggedright\noindent
vnb-lambda(\allowbreak{}wnt\_, ccint(\allowbreak{}0, 2 * pi), wnc\_ + wnr\_ * cc-exp(\allowbreak{}1i * wnt\_))
\end{quote}
characterized by \texttt{circle-path-closed}  \textit{(proven)}:
\begin{quote}\ttfamily\footnotesize\raggedright\noindent
forall(\allowbreak{}[c in cc, r in rr], (\allowbreak{}circle-path(\allowbreak{}c, r))(\allowbreak{}2 * pi) = (\allowbreak{}circle-path(\allowbreak{}c, r))(\allowbreak{}0))
\end{quote}
mentioned by 70 result(s)

\smallskip

\noindent\textbf{\texttt{ck-norm(cks\_, ckf\_)}}\quad\textit{functoid}\index{ck-norm(cks_, ckf_)@\texttt{ck-norm(cks\_, ckf\_)}}\\
reads: the supremum norm of \$2 on \$1 \\
definition:
\begin{quote}\ttfamily\footnotesize\raggedright\noindent
esup(\allowbreak{}image(\allowbreak{}vnb-lambda(\allowbreak{}ckp\_, pts(\allowbreak{}cks\_), abs(\allowbreak{}ckf\_(\allowbreak{}ckp\_))), pts(\allowbreak{}cks\_)))
\end{quote}
characterized by \texttt{ck-norm-real}  \textit{(proven)}:
\begin{quote}\ttfamily\footnotesize\raggedright\noindent
forall(\allowbreak{}[s, h], is-compact(\allowbreak{}s) implies h in cont-fns(\allowbreak{}s) implies ck-norm(\allowbreak{}s, h) in rr and 0 <= ck-norm(\allowbreak{}s, h))
\end{quote}
also: \texttt{ck-norm-bound} \texttt{ck-norm-upper} \texttt{ck-norm-act-ge} \texttt{ck-norm-act-le} \texttt{ck-norm-add-le} \texttt{ck-norm-zero-elim} \texttt{ck-norm-zero-intro} \\
mentioned by 12 result(s)

\smallskip

\noindent\textbf{\texttt{ck-nvs(cks\_)}}\quad\textit{functoid}\index{ck-nvs(cks_)@\texttt{ck-nvs(cks\_)}}\\
reads: the normed space of the continuous real functions on \$1 with the supremum norm \\
definition:
\begin{quote}\ttfamily\footnotesize\raggedright\noindent
[normed-field-as-commutative-ring(\allowbreak{}rr-normed-field), cont-fns(\allowbreak{}cks\_), vnb-lambda(\allowbreak{}[ckaf\_, ckag\_], cartesian(\allowbreak{}cont-fns(\allowbreak{}cks\_), cont-fns(\allowbreak{}cks\_)), vnb-lambda(\allowbreak{}sx\_, pts(\allowbreak{}cks\_), ckaf\_(\allowbreak{}sx\_) + ckag\_(\allowbreak{}sx\_))), vnb-lambda(\allowbreak{}sx\_, pts(\allowbreak{}cks\_), 0), vnb-lambda(\allowbreak{}cknf\_, cont-fns(\allowbreak{}cks\_), vnb-lambda(\allowbreak{}sx\_, pts(\allowbreak{}cks\_), -cknf\_(\allowbreak{}sx\_))), vnb-lambda(\allowbreak{}[cksr\_, cksf\_], cartesian(\allowbreak{}rr, cont-fns(\allowbreak{}cks\_)), vnb-lambda(\allowbreak{}sx\_, pts(\allowbreak{}cks\_), cksr\_ * cksf\_(\allowbreak{}sx\_))), vnb-lambda(\allowbreak{}ckvf\_, cont-fns(\allowbreak{}cks\_), ck-norm(\allowbreak{}cks\_, ckvf\_))]
\end{quote}
characterized by \texttt{ck-nvs-vec}  \textit{(proven)}:
\begin{quote}\ttfamily\footnotesize\raggedright\noindent
forall(\allowbreak{}[s], is-metric-space(\allowbreak{}s) implies vec(\allowbreak{}ck-nvs(\allowbreak{}s)) == cont-fns(\allowbreak{}s))
\end{quote}
also: \texttt{ck-nvs-scal} \texttt{ck-nvs-complete} \texttt{ck-nvs-is-normed-vector-space} \\
mentioned by 116 result(s)

\smallskip

\noindent\textbf{\texttt{cl2-hilbert(omega, ca, mu)}}\quad\textit{functoid}\index{cl2-hilbert(omega, ca, mu)@\texttt{cl2-hilbert(omega, ca, mu)}}\\
reads: the complex Hilbert space L\textasciicircum{}2 of \$3 on \$1 \\
definition:
\begin{quote}\ttfamily\footnotesize\raggedright\noindent
csip-cips(\allowbreak{}complexify(\allowbreak{}sip-quotient(\allowbreak{}l2-space(\allowbreak{}omega, ca, mu))))
\end{quote}
characterized by \texttt{cl2-hilbert-rep}  \textit{(proven)}:
\begin{quote}\ttfamily\footnotesize\raggedright\noindent
forall(\allowbreak{}[omega, ca, mu], is-measure(\allowbreak{}omega, ca, mu) implies forall(\allowbreak{}[x\_ in vec(\allowbreak{}cl2-hilbert(\allowbreak{}omega, ca, mu))], forsome(\allowbreak{}[f1 in l2-set(\allowbreak{}omega, ca, mu), f2 in l2-set(\allowbreak{}omega, ca, mu)], x\_ = [class(\allowbreak{}sip-setoid(\allowbreak{}l2-space(\allowbreak{}omega, ca, mu)), f1), class(\allowbreak{}sip-setoid(\allowbreak{}l2-space(\allowbreak{}omega, ca, mu)), f2)])))
\end{quote}
also: \texttt{cl2-hilbert-vec} \texttt{cl2-hilbert-unfold} \texttt{cl2-hilbert-ip-self} \texttt{cl2-hilbert-pair-in} \texttt{cl2-hilbert-ip-classes} \texttt{cl2-hilbert-act-classes} \texttt{cl2-hilbert-vadd-classes} \texttt{cl2-hilbert-is-hilbert-space} \\
mentioned by 23 result(s)

\smallskip

\noindent\textbf{\texttt{cl2-int-op(kcxs\_, kcca\_, kcmu\_, kcys\_, kccb\_, kcnu\_, kck1\_, kck2\_)}}\quad\textit{functoid}\index{cl2-int-op(kcxs_, kcca_, kcmu_, kcys_, kccb_, kcnu_, kck1_, kck2_)@\texttt{cl2-int-op(kcxs\_, kcca\_, kcmu\_, kcys\_, kccb\_, kcnu\_, kck1\_, kck2\_)}}\\
reads: the integral operator with kernel \$7 + i \$8 \\
definition:
\begin{quote}\ttfamily\footnotesize\raggedright\noindent
vnb-lambda(\allowbreak{}kcz\_, vec(\allowbreak{}cl2-hilbert(\allowbreak{}kcys\_, kccb\_, kcnu\_)), [class(\allowbreak{}sip-setoid(\allowbreak{}l2-space(\allowbreak{}kcxs\_, kcca\_, kcmu\_)), (\allowbreak{}vadd(\allowbreak{}l2-space(\allowbreak{}kcxs\_, kcca\_, kcmu\_)))(\allowbreak{}kernel-apply(\allowbreak{}kcxs\_, kcys\_, kccb\_, kcnu\_, kck1\_, choice(\allowbreak{}nth(\allowbreak{}1, kcz\_))), (\allowbreak{}vneg(\allowbreak{}l2-space(\allowbreak{}kcxs\_, kcca\_, kcmu\_)))(\allowbreak{}kernel-apply(\allowbreak{}kcxs\_, kcys\_, kccb\_, kcnu\_, kck2\_, choice(\allowbreak{}nth(\allowbreak{}2, kcz\_)))))), class(\allowbreak{}sip-setoid(\allowbreak{}l2-space(\allowbreak{}kcxs\_, kcca\_, kcmu\_)), (\allowbreak{}vadd(\allowbreak{}l2-space(\allowbreak{}kcxs\_, kcca\_, kcmu\_)))(\allowbreak{}kernel-apply(\allowbreak{}kcxs\_, kcys\_, kccb\_, kcnu\_, kck1\_, choice(\allowbreak{}nth(\allowbreak{}2, kcz\_))), kernel-apply(\allowbreak{}kcxs\_, kcys\_, kccb\_, kcnu\_, kck2\_, choice(\allowbreak{}nth(\allowbreak{}1, kcz\_)))))])
\end{quote}
characterized by \texttt{cl2-int-op-act}  \textit{(proven)}:
\begin{quote}\ttfamily\footnotesize\raggedright\noindent
forall(\allowbreak{}[xs, ca, mu, ys, cb, nu, k1, k2, c\_, z\_], is-measure(\allowbreak{}xs, ca, mu) implies is-measure(\allowbreak{}ys, cb, nu) implies is-sigma-finite(\allowbreak{}xs, ca, mu) implies is-sigma-finite(\allowbreak{}ys, cb, nu) implies k1 in l2-set(\allowbreak{}cartesian(\allowbreak{}xs, ys), product-sigma(\allowbreak{}xs, ca, ys, cb), product-measure(\allowbreak{}xs, ca, mu, ys, cb, nu)) implies k2 in l2-set(\allowbreak{}cartesian(\allowbreak{}xs, ys), product-sigma(\allowbreak{}xs, ca, ys, cb), product-measure(\allowbreak{}xs, ca, mu, ys, cb, nu)) implies c\_ in cc implies z\_ in vec(\allowbreak{}cl2-hilbert(\allowbreak{}ys, cb, nu)) implies (\allowbreak{}cl2-int-op(\allowbreak{}xs, ca, mu, ys, cb, nu, k1, k2))(\allowbreak{}(\allowbreak{}act(\allowbreak{}cl2-hilbert(\allowbreak{}ys, cb, nu)))(\allowbreak{}c\_, z\_)) = (\allowbreak{}act(\allowbreak{}cl2-hilbert(\allowbreak{}xs, ca, mu)))(\allowbreak{}c\_, (\allowbreak{}cl2-int-op(\allowbreak{}xs, ca, mu, ys, cb, nu, k1, k2))(\allowbreak{}z\_)))
\end{quote}
also: \texttt{cl2-int-op-add} \texttt{cl2-int-op-fun} \texttt{cl2-int-op-linear} \texttt{cl2-int-op-classes} \texttt{cl2-int-op-ip-bound} \texttt{cl2-int-op-norm-bound} \\
mentioned by 10 result(s)

\smallskip

\noindent\textbf{\texttt{cl2-space(omega, ca, mu)}}\quad\textit{functoid}\index{cl2-space(omega, ca, mu)@\texttt{cl2-space(omega, ca, mu)}}\\
reads: the complex L\textasciicircum{}2 space of \$3 on \$1 \\
definition:
\begin{quote}\ttfamily\footnotesize\raggedright\noindent
complexify(\allowbreak{}l2-space(\allowbreak{}omega, ca, mu))
\end{quote}
characterized by \texttt{cl2-space-is-csip}  \textit{(proven)}:
\begin{quote}\ttfamily\footnotesize\raggedright\noindent
forall(\allowbreak{}[omega, ca, mu], is-measure(\allowbreak{}omega, ca, mu) implies is-complex-semi-inner-product-space(\allowbreak{}cl2-space(\allowbreak{}omega, ca, mu)))
\end{quote}
mentioned by 3 result(s)

\smallskip

\noindent\textbf{\texttt{clamp(a, b, x)}}\quad\textit{functoid}\index{clamp(a, b, x)@\texttt{clamp(a, b, x)}}\\
reads: \$3 clamped to [\$1, \$2] \\
definition:
\begin{quote}\ttfamily\footnotesize\raggedright\noindent
min(\allowbreak{}max(\allowbreak{}x, a), b)
\end{quote}
characterized by \texttt{clamp-fixes}  \textit{(proven)}:
\begin{quote}\ttfamily\footnotesize\raggedright\noindent
forall(\allowbreak{}[a in rr, b in rr, x in rr], a <= x implies x <= b implies clamp(\allowbreak{}a, b, x) = x)
\end{quote}
also: \texttt{clamp-in-rr} \texttt{clamp-in-ccint} \texttt{clamp-lipschitz} \texttt{clamp-left-const} \texttt{clamp-right-const} \texttt{clamp-compose-fixes} \texttt{clamp-compose-lam-in-fun} \texttt{clamp-compose-continuous-at} \\
mentioned by 15 result(s)

\smallskip

\noindent\textbf{\texttt{class(s, a)}}\quad\textit{functoid}\index{class(s, a)@\texttt{class(s, a)}}\\
definition:
\begin{quote}\ttfamily\footnotesize\raggedright\noindent
\{b in pts(\allowbreak{}s): related(\allowbreak{}s, a, b)\}
\end{quote}
characterized by \texttt{sip-class-eq-iff-csip-as-sip}  \textit{(definitional)}:
\begin{quote}\ttfamily\footnotesize\raggedright\noindent
forall(\allowbreak{}[r], is-complex-semi-inner-product-space(\allowbreak{}r) implies forall(\allowbreak{}[a\_ in vec(\allowbreak{}r), b\_ in vec(\allowbreak{}r)], class(\allowbreak{}sip-setoid(\allowbreak{}csip-as-sip(\allowbreak{}r)), a\_) = class(\allowbreak{}sip-setoid(\allowbreak{}csip-as-sip(\allowbreak{}r)), b\_) iff (\allowbreak{}rip(\allowbreak{}r))(\allowbreak{}(\allowbreak{}vadd(\allowbreak{}r))(\allowbreak{}a\_, (\allowbreak{}vneg(\allowbreak{}r))(\allowbreak{}b\_)), (\allowbreak{}vadd(\allowbreak{}r))(\allowbreak{}a\_, (\allowbreak{}vneg(\allowbreak{}r))(\allowbreak{}b\_))) = 0))
\end{quote}
also: \texttt{qmod-class-eq-iff-cna-as-module} \texttt{qmod-class-eq-iff-sip-as-module} \texttt{qmod-class-eq-iff-uba-as-module} \texttt{qmod-class-eq-iff-cnvs-as-module} \texttt{qmod-class-eq-iff-sip-as-module-csip-as-sip} \texttt{qmod-class-eq-iff-uba-as-module-uiba-as-uba} \texttt{class-self} \texttt{class-eq-iff} ... (47 in all) \\
mentioned by 534 result(s)

\smallskip

\noindent\textbf{\texttt{clf}}\quad\textit{constant}\index{clf@\texttt{clf}}\\
reads: the sentences (locally closed well-formed formulas) \\
defined by \texttt{clf-def}

\smallskip

\noindent\textbf{\texttt{clopen-sets(s)}}\quad\textit{functoid}\index{clopen-sets(s)@\texttt{clopen-sets(s)}}\\
reads: the family of the subsets of \$1 that are open and closed \\
definition:
\begin{quote}\ttfamily\footnotesize\raggedright\noindent
\{cou\_ in power(\allowbreak{}pts(\allowbreak{}s)): is-open(\allowbreak{}s, cou\_) and is-closed(\allowbreak{}s, cou\_)\}
\end{quote}
mentioned by 40 result(s)

\smallskip

\noindent\textbf{\texttt{closed-annulus(anla\_, anlr\_, anls\_)}}\quad\textit{functoid}\index{closed-annulus(anla_, anlr_, anls_)@\texttt{closed-annulus(anla\_, anlr\_, anls\_)}}\\
reads: the closed annulus about \$1 between the radii \$2 and \$3 \\
definition:
\begin{quote}\ttfamily\footnotesize\raggedright\noindent
\{anlz\_ in cc: anlr\_ <= magnitude(\allowbreak{}anlz\_ - anla\_) and magnitude(\allowbreak{}anlz\_ - anla\_) <= anls\_\}
\end{quote}

\smallskip

\noindent\textbf{\texttt{closed-ball(s, c, r)}}\quad\textit{functoid}\index{closed-ball(s, c, r)@\texttt{closed-ball(s, c, r)}}\\
definition:
\begin{quote}\ttfamily\footnotesize\raggedright\noindent
\{y in pts(\allowbreak{}s): (\allowbreak{}dist(\allowbreak{}s))(\allowbreak{}c, y) <= r\}
\end{quote}
characterized by \texttt{qp-zp-closed-ball}  \textit{(proven)}:
\begin{quote}\ttfamily\footnotesize\raggedright\noindent
forall(\allowbreak{}[p], zz-prime(\allowbreak{}p) implies zp(\allowbreak{}p) = closed-ball(\allowbreak{}qp(\allowbreak{}p), class(\allowbreak{}cauchy-setoid(\allowbreak{}nf-metric-space(\allowbreak{}qq-padic(\allowbreak{}p))), vnb-lambda(\allowbreak{}qsn\_, nn, 0)), 1))
\end{quote}
also: \texttt{cc-closed-ball-mem} \texttt{closed-ball-is-set} \texttt{ball-in-closed-ball} \texttt{closed-ball-in-ball} \texttt{closed-ball-in-power} \texttt{cc-closed-ball-origin} \texttt{closed-ball-center-in} \texttt{closed-ball-is-closed} ... (26 in all) \\
mentioned by 76 result(s)

\smallskip

\noindent\textbf{\texttt{closed-conv-hull(m, t, e)}}\quad\textit{functoid}\index{closed-conv-hull(m, t, e)@\texttt{closed-conv-hull(m, t, e)}}\\
reads: the closed convex hull of \$3 in \$1 for the topology \$2 \\
definition:
\begin{quote}\ttfamily\footnotesize\raggedright\noindent
\{khx\_ in vec(\allowbreak{}m): forall(\allowbreak{}[khc\_ in set], is-convex-subset(\allowbreak{}m, khc\_) implies is-top-closed(\allowbreak{}t, khc\_) implies e subset khc\_ implies khx\_ in khc\_)\}
\end{quote}
characterized by \texttt{closed-conv-hull-mem}  \textit{(proven)}:
\begin{quote}\ttfamily\footnotesize\raggedright\noindent
forall(\allowbreak{}[m, t, e, z], z in closed-conv-hull(\allowbreak{}m, t, e) iff z in vec(\allowbreak{}m) and forall(\allowbreak{}[khc\_ in set], is-convex-subset(\allowbreak{}m, khc\_) implies is-top-closed(\allowbreak{}t, khc\_) implies e subset khc\_ implies z in khc\_))
\end{quote}
also: \texttt{closed-conv-hull-least} \texttt{closed-conv-hull-closed} \texttt{closed-conv-hull-convex} \texttt{closed-conv-hull-unfold} \texttt{closed-conv-hull-contains} \\
mentioned by 10 result(s)

\smallskip

\noindent\textbf{\texttt{closed-intervals}}\quad\textit{primitive}\index{closed-intervals@\texttt{closed-intervals}}\\
reads: the family of the empty set and the closed intervals of the real line \\
defined by \texttt{closed-intervals-def}

\smallskip

\noindent\textbf{\texttt{closure(s, a)}}\quad\textit{functoid}\index{closure(s, a)@\texttt{closure(s, a)}}\\
definition:
\begin{quote}\ttfamily\footnotesize\raggedright\noindent
\{x in pts(\allowbreak{}s): forall(\allowbreak{}[u], u in power(\allowbreak{}pts(\allowbreak{}s)) and is-open(\allowbreak{}s, u) and x in u implies forsome(\allowbreak{}[y in u], y in a))\}
\end{quote}
characterized by \texttt{ms-closure-tb}  \textit{(proven)}:
\begin{quote}\ttfamily\footnotesize\raggedright\noindent
forall(\allowbreak{}[s], is-metric-space(\allowbreak{}s) implies forall(\allowbreak{}[a], a subset pts(\allowbreak{}s) implies is-totally-bounded-in(\allowbreak{}s, a) implies is-totally-bounded-in(\allowbreak{}s, closure(\allowbreak{}s, a))))
\end{quote}
also: \texttt{closure-is-set} \texttt{nf-closure-add} \texttt{nf-closure-inv} \texttt{nf-closure-mul} \texttt{nf-closure-neg} \texttt{bic-closure-act} \texttt{bic-closure-add} \texttt{bic-closure-mul} ... (46 in all) \\
mentioned by 98 result(s)

\smallskip

\noindent\textbf{\texttt{clt}}\quad\textit{constant}\index{clt@\texttt{clt}}\\
reads: the locally closed well-formed terms \\
defined by \texttt{clt-def}

\smallskip

\noindent\textbf{\texttt{cluster-point(s, f, x)}}\quad\textit{predicate}\index{cluster-point(s, f, x)@\texttt{cluster-point(s, f, x)}}\\
reads: \$3 is a cluster point of the sequence \$2 in \$1 \\
defined by \texttt{cluster-point} \\
characterized by \texttt{cluster-point}  \textit{(definitional)}:
\begin{quote}\ttfamily\footnotesize\raggedright\noindent
forall(\allowbreak{}[s, f, x], cluster-point(\allowbreak{}s, f, x) iff is-metric-space(\allowbreak{}s) and f in fun(\allowbreak{}nn, pts(\allowbreak{}s)) and x in pts(\allowbreak{}s) and forall(\allowbreak{}[eps], pos-rr(\allowbreak{}eps) implies forall(\allowbreak{}[m in nn], forsome(\allowbreak{}[n in nn], m <= n and (\allowbreak{}dist(\allowbreak{}s))(\allowbreak{}f(\allowbreak{}n), x) < eps))))
\end{quote}
also: \texttt{compact-iff-cluster-point} \texttt{cluster-point-has-convergent-subseq} \\
mentioned by 12 result(s)

\smallskip

\noindent\textbf{\texttt{cmod-cnvs(m, nf)}}\quad\textit{functoid}\index{cmod-cnvs(m, nf)@\texttt{cmod-cnvs(m, nf)}}\\
reads: the complex module \$1 with norm \$2, as a real space with complex structure \\
definition:
\begin{quote}\ttfamily\footnotesize\raggedright\noindent
[normed-field-as-commutative-ring(\allowbreak{}rr-normed-field), vec(\allowbreak{}m), vadd(\allowbreak{}m), vzero(\allowbreak{}m), vneg(\allowbreak{}m), cmod-ract(\allowbreak{}m), nf, cmod-j(\allowbreak{}m)]
\end{quote}
characterized by \texttt{cmod-cnvs-act}  \textit{(proven)}:
\begin{quote}\ttfamily\footnotesize\raggedright\noindent
forall(\allowbreak{}[m], is-module(\allowbreak{}m) implies scal(\allowbreak{}m) = normed-field-as-commutative-ring(\allowbreak{}cc-normed-field) implies forall(\allowbreak{}[nfn], is-cmod-norm(\allowbreak{}m, nfn) implies act(\allowbreak{}cmod-cnvs(\allowbreak{}m, nfn)) == cmod-ract(\allowbreak{}m)))
\end{quote}
also: \texttt{cmod-cnvs-jop} \texttt{cmod-cnvs-nvs} \texttt{cmod-cnvs-vec} \texttt{cmod-cnvs-scal} \texttt{cmod-cnvs-vadd} \texttt{cmod-cnvs-vneg} \texttt{cmod-cnvs-vnrm} \texttt{cmod-cnvs-vzero} ... (21 in all) \\
mentioned by 34 result(s)

\smallskip

\noindent\textbf{\texttt{cmod-j(m)}}\quad\textit{functoid}\index{cmod-j(m)@\texttt{cmod-j(m)}}\\
reads: multiplication by i on \$1 \\
definition:
\begin{quote}\ttfamily\footnotesize\raggedright\noindent
vnb-lambda(\allowbreak{}cnx\_, vec(\allowbreak{}m), (\allowbreak{}act(\allowbreak{}m))(\allowbreak{}1i, cnx\_))
\end{quote}
characterized by \texttt{cmod-j-fun}  \textit{(proven)}:
\begin{quote}\ttfamily\footnotesize\raggedright\noindent
forall(\allowbreak{}[m], is-module(\allowbreak{}m) implies scal(\allowbreak{}m) = normed-field-as-commutative-ring(\allowbreak{}cc-normed-field) implies cmod-j(\allowbreak{}m) in fun(\allowbreak{}vec(\allowbreak{}m), vec(\allowbreak{}m)))
\end{quote}
mentioned by 3 result(s)

\smallskip

\noindent\textbf{\texttt{cmod-nvs(m, nf)}}\quad\textit{functoid}\index{cmod-nvs(m, nf)@\texttt{cmod-nvs(m, nf)}}\\
reads: the real normed space underlying the complex module \$1 with norm \$2 \\
definition:
\begin{quote}\ttfamily\footnotesize\raggedright\noindent
[normed-field-as-commutative-ring(\allowbreak{}rr-normed-field), vec(\allowbreak{}m), vadd(\allowbreak{}m), vzero(\allowbreak{}m), vneg(\allowbreak{}m), cmod-ract(\allowbreak{}m), nf]
\end{quote}
characterized by \texttt{cmod-nvs-act}  \textit{(proven)}:
\begin{quote}\ttfamily\footnotesize\raggedright\noindent
forall(\allowbreak{}[m], is-module(\allowbreak{}m) implies scal(\allowbreak{}m) = normed-field-as-commutative-ring(\allowbreak{}cc-normed-field) implies forall(\allowbreak{}[nfn], is-cmod-norm(\allowbreak{}m, nfn) implies act(\allowbreak{}cmod-nvs(\allowbreak{}m, nfn)) == cmod-ract(\allowbreak{}m)))
\end{quote}
also: \texttt{cmod-nvs-vec} \texttt{cmod-nvs-comm} \texttt{cmod-nvs-scal} \texttt{cmod-nvs-vadd} \texttt{cmod-nvs-vneg} \texttt{cmod-nvs-vnrm} \texttt{cmod-nvs-assoc} \texttt{cmod-nvs-homog} ... (29 in all) \\
mentioned by 46 result(s)

\smallskip

\noindent\textbf{\texttt{cmod-ract(m)}}\quad\textit{functoid}\index{cmod-ract(m)@\texttt{cmod-ract(m)}}\\
reads: the real scalar multiplication of the complex module \$1 \\
definition:
\begin{quote}\ttfamily\footnotesize\raggedright\noindent
vnb-lambda(\allowbreak{}[cnr\_, cnx\_], cartesian(\allowbreak{}rr, vec(\allowbreak{}m)), (\allowbreak{}act(\allowbreak{}m))(\allowbreak{}cnr\_, cnx\_))
\end{quote}
characterized by \texttt{cmod-ract-fun}  \textit{(proven)}:
\begin{quote}\ttfamily\footnotesize\raggedright\noindent
forall(\allowbreak{}[m], is-module(\allowbreak{}m) implies scal(\allowbreak{}m) = normed-field-as-commutative-ring(\allowbreak{}cc-normed-field) implies cmod-ract(\allowbreak{}m) in fun(\allowbreak{}cartesian(\allowbreak{}rr, vec(\allowbreak{}m)), vec(\allowbreak{}m)))
\end{quote}
mentioned by 7 result(s)

\smallskip

\noindent\textbf{\texttt{cna-as-module(s)}}\quad\textit{view}\index{cna-as-module(s)@\texttt{cna-as-module(s)}}\\
a complex-normed-algebra viewed as a module \\
characterized by \texttt{mvag-id-cna-as-module}  \textit{(definitional)}:
\begin{quote}\ttfamily\footnotesize\raggedright\noindent
forall(\allowbreak{}[r], is-complex-normed-algebra(\allowbreak{}r) implies iden(\allowbreak{}module-vector-ag(\allowbreak{}cna-as-module(\allowbreak{}r))) = vzero(\allowbreak{}r))
\end{quote}
also: \texttt{mvag-op-cna-as-module} \texttt{cna-as-module-is-module} \texttt{mvag-carr-cna-as-module} \texttt{qmod-comm-cna-as-module} \texttt{qmod-assoc-cna-as-module} \texttt{spans-span-cna-as-module} \texttt{hbs-lf-zero-cna-as-module} \texttt{hbs-sl-zero-cna-as-module} ... (155 in all) \\
mentioned by 334 result(s)

\smallskip

\noindent\textbf{\texttt{cna-as-normed-ag(s)}}\quad\textit{view}\index{cna-as-normed-ag(s)@\texttt{cna-as-normed-ag(s)}}\\
a complex-normed-algebra viewed as a normed-ag \\
characterized by \texttt{cna-as-normed-ag-is-normed-ag}  \textit{(definitional)}:
\begin{quote}\ttfamily\footnotesize\raggedright\noindent
forall(\allowbreak{}[r], is-complex-normed-algebra(\allowbreak{}r) implies is-normed-ag(\allowbreak{}cna-as-normed-ag(\allowbreak{}r)))
\end{quote}
also: \texttt{rb2-dist-union-cna-as-normed-ag} \texttt{sums-to-unique-cna-as-normed-ag} \texttt{rb2-dist-subset-cna-as-normed-ag} \texttt{hom-normed-ag-id-cna-as-normed-ag} \texttt{rb2-nrm-finsum-le-cna-as-normed-ag} \texttt{nag-metric-carrier-cna-as-normed-ag} \texttt{is-normed-ag-is-set-cna-as-normed-ag} \texttt{nag-metric-distance-cna-as-normed-ag} ... (14 in all) \\
mentioned by 45 result(s)

\smallskip

\noindent\textbf{\texttt{cna-cnvs(cna\_)}}\quad\textit{functoid}\index{cna-cnvs(cna_)@\texttt{cna-cnvs(cna\_)}}\\
reads: the normed algebra \$1 as a complex normed vector space \\
definition:
\begin{quote}\ttfamily\footnotesize\raggedright\noindent
cmod-cnvs(\allowbreak{}cna-as-module(\allowbreak{}cna\_), vnrm(\allowbreak{}cna\_))
\end{quote}
characterized by \texttt{cna-cnvs-vec}  \textit{(proven)}:
\begin{quote}\ttfamily\footnotesize\raggedright\noindent
forall(\allowbreak{}[a], is-complex-normed-algebra(\allowbreak{}a) implies vec(\allowbreak{}cna-cnvs(\allowbreak{}a)) == vec(\allowbreak{}a))
\end{quote}
also: \texttt{cna-cnvs-cact} \texttt{cna-cnvs-vadd} \texttt{cna-cnvs-vneg} \texttt{cna-cnvs-vnrm} \texttt{cna-cnvs-vzero} \texttt{cna-cnvs-is-cnvs} \texttt{cna-cnvs-normed-ag} \texttt{cna-cnvs-metric-space} \\
mentioned by 18 result(s)

\smallskip

\noindent\textbf{\texttt{cna-metric-space(baa\_)}}\quad\textit{functoid}\index{cna-metric-space(baa_)@\texttt{cna-metric-space(baa\_)}}\\
reads: the metric space of the normed algebra \$1 \\
definition:
\begin{quote}\ttfamily\footnotesize\raggedright\noindent
nag-metric-space(\allowbreak{}cna-as-normed-ag(\allowbreak{}baa\_))
\end{quote}
mentioned by 2 result(s)

\smallskip

\noindent\textbf{\texttt{cnvs-as-module(s)}}\quad\textit{view}\index{cnvs-as-module(s)@\texttt{cnvs-as-module(s)}}\\
a complex-normed-vector-space viewed as a module \\
characterized by \texttt{mvag-id-cnvs-as-module}  \textit{(definitional)}:
\begin{quote}\ttfamily\footnotesize\raggedright\noindent
forall(\allowbreak{}[r], is-complex-normed-vector-space(\allowbreak{}r) implies iden(\allowbreak{}module-vector-ag(\allowbreak{}cnvs-as-module(\allowbreak{}r))) = vzero(\allowbreak{}r))
\end{quote}
also: \texttt{mvag-op-cnvs-as-module} \texttt{cnvs-as-module-is-module} \texttt{mvag-carr-cnvs-as-module} \texttt{qmod-comm-cnvs-as-module} \texttt{pv-act-fun-cnvs-as-module} \texttt{qmod-assoc-cnvs-as-module} \texttt{spans-span-cnvs-as-module} \texttt{hbs-lf-zero-cnvs-as-module} ... (176 in all) \\
mentioned by 381 result(s)

\smallskip

\noindent\textbf{\texttt{cnvs-as-normed-ag(s)}}\quad\textit{view}\index{cnvs-as-normed-ag(s)@\texttt{cnvs-as-normed-ag(s)}}\\
a complex-normed-vector-space viewed as a normed-ag \\
characterized by \texttt{cnvs-as-normed-ag-is-normed-ag}  \textit{(definitional)}:
\begin{quote}\ttfamily\footnotesize\raggedright\noindent
forall(\allowbreak{}[r], is-complex-normed-vector-space(\allowbreak{}r) implies is-normed-ag(\allowbreak{}cnvs-as-normed-ag(\allowbreak{}r)))
\end{quote}
also: \texttt{rb2-dist-union-cnvs-as-normed-ag} \texttt{sums-to-unique-cnvs-as-normed-ag} \texttt{rb2-dist-subset-cnvs-as-normed-ag} \texttt{hom-normed-ag-id-cnvs-as-normed-ag} \texttt{rb2-nrm-finsum-le-cnvs-as-normed-ag} \texttt{nag-metric-carrier-cnvs-as-normed-ag} \texttt{is-normed-ag-is-set-cnvs-as-normed-ag} \texttt{nag-metric-distance-cnvs-as-normed-ag} ... (15 in all) \\
mentioned by 29 result(s)

\smallskip

\noindent\textbf{\texttt{cnvs-as-nvs(s)}}\quad\textit{view}\index{cnvs-as-nvs(s)@\texttt{cnvs-as-nvs(s)}}\\
a complex-normed-vector-space viewed as a normed-vector-space \\
characterized by \texttt{nvs-dist-shift-cnvs-as-nvs}  \textit{(definitional)}:
\begin{quote}\ttfamily\footnotesize\raggedright\noindent
forall(\allowbreak{}[r], is-complex-normed-vector-space(\allowbreak{}r) implies forall(\allowbreak{}[u\_ in vec(\allowbreak{}r), w\_ in vec(\allowbreak{}r)], (\allowbreak{}dist(\allowbreak{}nvs-metric-space(\allowbreak{}cnvs-as-nvs(\allowbreak{}r))))(\allowbreak{}u\_, (\allowbreak{}vadd(\allowbreak{}r))(\allowbreak{}u\_, w\_)) = (\allowbreak{}vnrm(\allowbreak{}r))(\allowbreak{}w\_)))
\end{quote}
also: \texttt{dual-act-typing-cnvs-as-nvs} \texttt{dual-vec-in-set-cnvs-as-nvs} \texttt{lfe-line-subset-cnvs-as-nvs} \texttt{nvs-norm-family-cnvs-as-nvs} \texttt{bidual-view-scal-cnvs-as-nvs} \texttt{dual-vadd-typing-cnvs-as-nvs} \texttt{dual-vneg-typing-cnvs-as-nvs} \texttt{dual-vnrm-typing-cnvs-as-nvs} ... (50 in all) \\
mentioned by 94 result(s)

\smallskip

\noindent\textbf{\texttt{cnvs-bounded-ops(cbv\_, cbw\_)}}\quad\textit{functoid}\index{cnvs-bounded-ops(cbv_, cbw_)@\texttt{cnvs-bounded-ops(cbv\_, cbw\_)}}\\
reads: the bounded operators from \$1 to \$2 \\
definition:
\begin{quote}\ttfamily\footnotesize\raggedright\noindent
\{cbt\_ in fun(\allowbreak{}vec(\allowbreak{}cbv\_), vec(\allowbreak{}cbw\_)): is-cnvs-bounded-op(\allowbreak{}cbv\_, cbw\_, cbt\_)\}
\end{quote}
mentioned by 72 result(s)

\smallskip

\noindent\textbf{\texttt{cnvs-metric-space(cnv\_)}}\quad\textit{functoid}\index{cnvs-metric-space(cnv_)@\texttt{cnvs-metric-space(cnv\_)}}\\
reads: the metric space of the normed space \$1 \\
definition:
\begin{quote}\ttfamily\footnotesize\raggedright\noindent
nag-metric-space(\allowbreak{}cnvs-as-normed-ag(\allowbreak{}cnv\_))
\end{quote}
characterized by \texttt{cna-cnvs-metric-space}  \textit{(proven)}:
\begin{quote}\ttfamily\footnotesize\raggedright\noindent
forall(\allowbreak{}[a], is-complex-normed-algebra(\allowbreak{}a) implies cnvs-metric-space(\allowbreak{}cna-cnvs(\allowbreak{}a)) == cna-metric-space(\allowbreak{}a))
\end{quote}
also: \texttt{cips-cnvs-metric-space} \\
mentioned by 18 result(s)

\smallskip

\noindent\textbf{\texttt{cnvs-opnorm(cbv\_, cbw\_, cbt\_)}}\quad\textit{functoid}\index{cnvs-opnorm(cbv_, cbw_, cbt_)@\texttt{cnvs-opnorm(cbv\_, cbw\_, cbt\_)}}\\
reads: the operator norm of \$3 \\
definition:
\begin{quote}\ttfamily\footnotesize\raggedright\noindent
sup(\allowbreak{}\{cba\_ in rr: cba\_ = 0 or forsome(\allowbreak{}[cbx\_ in vec(\allowbreak{}cbv\_)], not(\allowbreak{}(\allowbreak{}vnrm(\allowbreak{}cbv\_))(\allowbreak{}cbx\_) = 0) and cba\_ = (\allowbreak{}vnrm(\allowbreak{}cbw\_))(\allowbreak{}cbt\_(\allowbreak{}cbx\_)) * recip(\allowbreak{}(\allowbreak{}vnrm(\allowbreak{}cbv\_))(\allowbreak{}cbx\_)))\})
\end{quote}
characterized by \texttt{cnvs-opnorm-neg}  \textit{(proven)}:
\begin{quote}\ttfamily\footnotesize\raggedright\noindent
forall(\allowbreak{}[v, w, t\_], is-complex-normed-vector-space(\allowbreak{}v) implies is-complex-normed-vector-space(\allowbreak{}w) implies t\_ in cnvs-bounded-ops(\allowbreak{}v, w) implies cnvs-opnorm(\allowbreak{}v, w, cbop-neg(\allowbreak{}v, w, t\_)) = cnvs-opnorm(\allowbreak{}v, w, t\_))
\end{quote}
also: \texttt{cnvs-opnorm-spec} \texttt{cnvs-opnorm-zero} \texttt{cnvs-opnorm-bound} \texttt{cnvs-opnorm-in-rr} \texttt{cnvs-opnorm-scale} \texttt{cnvs-opnorm-nonneg} \texttt{opnorm-cnvs-opnorm} \texttt{cnvs-opnorm-comp-le} ... (11 in all) \\
mentioned by 22 result(s)

\smallskip

\noindent\textbf{\texttt{coint(a, b)}}\quad\textit{functoid}\index{coint(a, b)@\texttt{coint(a, b)}}\\
reads: the interval from \$1 included to \$2 excluded \\
definition:
\begin{quote}\ttfamily\footnotesize\raggedright\noindent
\{osq\_ in rr: a <= osq\_ and osq\_ < b\}
\end{quote}
characterized by \texttt{coint-shrink}  \textit{(proven)}:
\begin{quote}\ttfamily\footnotesize\raggedright\noindent
forall(\allowbreak{}[osx\_, osr\_, osv\_], osx\_ in rr implies pos-rr(\allowbreak{}osr\_) implies pos-rr(\allowbreak{}osv\_) implies osv\_ <= osr\_ implies coint(\allowbreak{}osx\_, osx\_ + osv\_) subset coint(\allowbreak{}osx\_, osx\_ + osr\_))
\end{quote}
also: \texttt{coint-of-ooint} \texttt{coint-elt-in-rr} \texttt{coint-subset-rr} \texttt{coint-membership} \texttt{coint-subset-ooint} \texttt{coint-left-endpoint} \texttt{lebesgue-outer-coint} \texttt{coint-left-limit-point} ... (9 in all) \\
mentioned by 19 result(s)

\smallskip

\noindent\textbf{\texttt{col-as-row(gcr-a\_, gcr-n\_, gcr-j\_)}}\quad\textit{functoid}\index{col-as-row(gcr-a_, gcr-n_, gcr-j_)@\texttt{col-as-row(gcr-a\_, gcr-n\_, gcr-j\_)}}\\
reads: column \$3 of \$1 as a row \\
definition:
\begin{quote}\ttfamily\footnotesize\raggedright\noindent
matof(\allowbreak{}1, gcr-n\_, vnb-lambda(\allowbreak{}[gcri\_, gcrk\_], cartesian(\allowbreak{}interval(\allowbreak{}1, 1), interval(\allowbreak{}1, gcr-n\_)), entry(\allowbreak{}gcr-a\_, gcrk\_, gcr-j\_)))
\end{quote}
characterized by \texttt{col-as-row-type}  \textit{(proven)}:
\begin{quote}\ttfamily\footnotesize\raggedright\noindent
forall(\allowbreak{}[gcra\_, n in nn, x, a in mat(\allowbreak{}n, n, x), j in interval(\allowbreak{}1, n)], col-as-row(\allowbreak{}a, n, j) in mat(\allowbreak{}1, n, x))
\end{quote}
also: \texttt{col-as-row-entry} \\
mentioned by 5 result(s)

\smallskip

\noindent\textbf{\texttt{col-equiv(r, n, a, b)}}\quad\textit{predicate}\index{col-equiv(r, n, a, b)@\texttt{col-equiv(r, n, a, b)}}\\
reads: \$3 and \$4 are column equivalent \\
defined by \texttt{col-equiv} \\
characterized by \texttt{col-equiv}  \textit{(definitional)}:
\begin{quote}\ttfamily\footnotesize\raggedright\noindent
forall(\allowbreak{}[r, n, a, b], col-equiv(\allowbreak{}r, n, a, b) iff forsome(\allowbreak{}[cep\_ in elem-gen(\allowbreak{}r, n)], b = matmul(\allowbreak{}r, a, cep\_)))
\end{quote}
also: \texttt{col-equiv-mat} \texttt{col-equiv-elem} \texttt{col-equiv-refl} \texttt{col-equiv-trans} \texttt{col-equiv-identity-in-elem-gen} \texttt{invertible-iff-col-equiv-identity} \\
mentioned by 17 result(s)

\smallskip

\noindent\textbf{\texttt{col-of(q, n, j)}}\quad\textit{functoid}\index{col-of(q, n, j)@\texttt{col-of(q, n, j)}}\\
reads: the column \$3 of \$1 \\
definition:
\begin{quote}\ttfamily\footnotesize\raggedright\noindent
matof(\allowbreak{}n, 1, vnb-lambda(\allowbreak{}[zkvi\_, zkvu\_], cartesian(\allowbreak{}interval(\allowbreak{}1, n), interval(\allowbreak{}1, 1)), entry(\allowbreak{}q, zkvi\_, j)))
\end{quote}
characterized by \texttt{cq4-col-of-zero}  \textit{(proven)}:
\begin{quote}\ttfamily\footnotesize\raggedright\noindent
forall(\allowbreak{}[r], is-ring(\allowbreak{}r) implies forall(\allowbreak{}[n in nn, j in interval(\allowbreak{}1, n)], col-of(\allowbreak{}zeromat(\allowbreak{}r, n, n), n, j) = zeromat(\allowbreak{}r, n, 1)))
\end{quote}
mentioned by 9 result(s)

\smallskip

\noindent\textbf{\texttt{col-step-row(r, n, b, i)}}\quad\textit{predicate}\index{col-step-row(r, n, b, i)@\texttt{col-step-row(r, n, b, i)}}\\
reads: row \$4 of \$3 is a row of a reduced column step matrix \\
defined by \texttt{col-step-row} \\
characterized by \texttt{col-step-row}  \textit{(definitional)}:
\begin{quote}\ttfamily\footnotesize\raggedright\noindent
forall(\allowbreak{}[r, n, b, i], col-step-row(\allowbreak{}r, n, b, i) iff forall(\allowbreak{}[csj\_ in interval(\allowbreak{}1, n)], i < csj\_ implies entry(\allowbreak{}b, i, csj\_) = zero(\allowbreak{}r)) and (\allowbreak{}i in interval(\allowbreak{}1, n) implies (\allowbreak{}entry(\allowbreak{}b, i, i) = zero(\allowbreak{}r) or entry(\allowbreak{}b, i, i) = one(\allowbreak{}r)) and (\allowbreak{}entry(\allowbreak{}b, i, i) = one(\allowbreak{}r) implies forall(\allowbreak{}[csj\_ in interval(\allowbreak{}1, n)], not(\allowbreak{}csj\_ = i) implies entry(\allowbreak{}b, i, csj\_) = zero(\allowbreak{}r)))))
\end{quote}
mentioned by 15 result(s)

\smallskip

\noindent\textbf{\texttt{collapse-at(k\_)}}\quad\textit{functoid}\index{collapse-at(k_)@\texttt{collapse-at(k\_)}}\\
reads: the collapse of NN at \$1 \\
definition:
\begin{quote}\ttfamily\footnotesize\raggedright\noindent
vnb-lambda(\allowbreak{}z\_, nn, if(\allowbreak{}k\_ < z\_, pred(\allowbreak{}z\_), z\_))
\end{quote}
characterized by \texttt{collapse-at-hi}  \textit{(proven)}:
\begin{quote}\ttfamily\footnotesize\raggedright\noindent
forall(\allowbreak{}[k\_ in nn, j\_ in nn], k\_ < j\_ implies (\allowbreak{}collapse-at(\allowbreak{}k\_))(\allowbreak{}j\_) = pred(\allowbreak{}j\_))
\end{quote}
also: \texttt{collapse-at-lo} \texttt{collapse-at-lt} \texttt{collapse-at-inj} \texttt{collapse-at-in-nn} \texttt{collapse-at-in-seg} \\
mentioned by 9 result(s)

\smallskip

\noindent\textbf{\texttt{comb-kk}}\quad\textit{defined-fn}\index{comb-kk@\texttt{comb-kk}}\\
defined by \texttt{comb-kk-zero} \texttt{comb-kk-succ} \\
characterized by \texttt{comb-kk-succ}  \textit{(definitional)}:
\begin{quote}\ttfamily\footnotesize\raggedright\noindent
forall(\allowbreak{}[r, x, y, m in nn], comb-kk(\allowbreak{}r, x, y, succ(\allowbreak{}m)) == vnb-lambda(\allowbreak{}k, zz, (\allowbreak{}add(\allowbreak{}r))(\allowbreak{}(\allowbreak{}mul(\allowbreak{}r))(\allowbreak{}x, (\allowbreak{}comb-kk(\allowbreak{}r, x, y, m))(\allowbreak{}k - 1)), (\allowbreak{}mul(\allowbreak{}r))(\allowbreak{}y, (\allowbreak{}comb-kk(\allowbreak{}r, x, y, m))(\allowbreak{}k)))))
\end{quote}
also: \texttt{comb-kk-zero} \texttt{comb-kk-0-0} \texttt{comb-kk-null} \texttt{comb-kk-above} \texttt{comb-kk-in-fun} \texttt{comb-kk-null-ind} \texttt{comb-kk-above-ind} \texttt{comb-kk-in-fun-ind} ... (9 in all) \\
mentioned by 24 result(s)

\smallskip

\noindent\textbf{\texttt{comm-monoid}}\quad\textit{structure}\index{comm-monoid@\texttt{comm-monoid}}\\
predicate \texttt{is-comm-monoid}; slots: carr opr iden

\smallskip

\noindent\textbf{\texttt{commutant(vnv\_, vns\_)}}\quad\textit{functoid}\index{commutant(vnv_, vns_)@\texttt{commutant(vnv\_, vns\_)}}\\
reads: the commutant of \$2 in the bounded operators on \$1 \\
definition:
\begin{quote}\ttfamily\footnotesize\raggedright\noindent
\{vnt\_ in bounded-ops(\allowbreak{}vnv\_, vnv\_): forall(\allowbreak{}[vna\_ in vns\_], bop-comp(\allowbreak{}vnv\_, vnt\_, vna\_) = bop-comp(\allowbreak{}vnv\_, vna\_, vnt\_))\}
\end{quote}
characterized by \texttt{commutant-unfold}  \textit{(proven)}:
\begin{quote}\ttfamily\footnotesize\raggedright\noindent
forall(\allowbreak{}[v, s], commutant(\allowbreak{}v, s) == \{vnt\_ in bounded-ops(\allowbreak{}v, v): forall(\allowbreak{}[vna\_ in s], bop-comp(\allowbreak{}v, vnt\_, vna\_) = bop-comp(\allowbreak{}v, vna\_, vnt\_))\})
\end{quote}
also: \texttt{bicommutant-unfold} \texttt{commutant-subset-bo} \texttt{commutant-membership} \texttt{commutant-sot-closed} \texttt{commutant-von-neumann} \texttt{commutant-star-algebra} \texttt{orbit-projection-commutant} \texttt{cyclic-commutant-separating} ... (11 in all) \\
mentioned by 18 result(s)

\smallskip

\noindent\textbf{\texttt{commutative-banach-algebra}}\quad\textit{refinement}\index{commutative-banach-algebra@\texttt{commutative-banach-algebra}}\\
predicate \texttt{is-commutative-banach-algebra}; same shape as unital-banach-algebra

\smallskip

\noindent\textbf{\texttt{commutative-c-star-algebra}}\quad\textit{refinement}\index{commutative-c-star-algebra@\texttt{commutative-c-star-algebra}}\\
predicate \texttt{is-commutative-c-star-algebra}; same shape as unital-c-star-algebra

\smallskip

\noindent\textbf{\texttt{commutative-ring}}\quad\textit{refinement}\index{commutative-ring@\texttt{commutative-ring}}\\
predicate \texttt{is-commutative-ring}; same shape as ring

\smallskip

\noindent\textbf{\texttt{commutative-ring-additive-ag(s)}}\quad\textit{view}\index{commutative-ring-additive-ag(s)@\texttt{commutative-ring-additive-ag(s)}}\\
a commutative-ring viewed as an abelian-group \\
characterized by \texttt{commutative-ring-additive-ag-is-abelian-group}  \textit{(definitional)}:
\begin{quote}\ttfamily\footnotesize\raggedright\noindent
forall(\allowbreak{}[r], is-commutative-ring(\allowbreak{}r) implies is-abelian-group(\allowbreak{}commutative-ring-additive-ag(\allowbreak{}r)))
\end{quote}
also: \texttt{abelian-group-as-monoid-is-monoid-commutative-ring-additive-ag} \texttt{commutative-ring-additive-ag-is-abelian-group-normed-field-as-commutative-ring} \texttt{abelian-group-as-monoid-is-monoid-commutative-ring-additive-ag-normed-field-as-commutative-ring} \texttt{commutative-ring-additive-ag-functorial} \\
mentioned by 68 result(s)

\smallskip

\noindent\textbf{\texttt{commutative-ring-multiplicative-cm(s)}}\quad\textit{view}\index{commutative-ring-multiplicative-cm(s)@\texttt{commutative-ring-multiplicative-cm(s)}}\\
a commutative-ring viewed as a comm-monoid \\
characterized by \texttt{commutative-ring-multiplicative-cm-is-comm-monoid}  \textit{(definitional)}:
\begin{quote}\ttfamily\footnotesize\raggedright\noindent
forall(\allowbreak{}[r], is-commutative-ring(\allowbreak{}r) implies is-comm-monoid(\allowbreak{}commutative-ring-multiplicative-cm(\allowbreak{}r)))
\end{quote}
also: \texttt{commutative-ring-multiplicative-cm-is-comm-monoid-normed-field-as-commutative-ring} \texttt{commutative-ring-multiplicative-cm-functorial} \\
mentioned by 36 result(s)

\smallskip

\noindent\textbf{\texttt{comp}}\quad\textit{operator}\index{comp@\texttt{comp}}\\
a term-forming head \\
mentioned by 4 result(s)

\smallskip

\noindent\textbf{\texttt{comp-me}}\quad\textit{tactic}\index{comp-me@\texttt{comp-me}}\\
\texttt{(comp-me hyp)} -- Comprehension membership elim. \\
declared in \texttt{tactics-help.scm}

\smallskip

\noindent\textbf{\texttt{comp-mi}}\quad\textit{tactic}\index{comp-mi@\texttt{comp-mi}}\\
\texttt{(comp-mi)} -- Comprehension membership intro. \\
declared in \texttt{tactics-help.scm}

\smallskip

\noindent\textbf{\texttt{compact-ops(kqv\_, kqw\_)}}\quad\textit{functoid}\index{compact-ops(kqv_, kqw_)@\texttt{compact-ops(kqv\_, kqw\_)}}\\
reads: the compact operators from \$1 to \$2 \\
definition:
\begin{quote}\ttfamily\footnotesize\raggedright\noindent
\{kqt\_ in bounded-ops(\allowbreak{}kqv\_, kqw\_): is-compact-op(\allowbreak{}kqv\_, kqw\_, kqt\_)\}
\end{quote}
characterized by \texttt{compact-ops-unfold}  \textit{(proven)}:
\begin{quote}\ttfamily\footnotesize\raggedright\noindent
forall(\allowbreak{}[kqv\_, kqw\_], compact-ops(\allowbreak{}kqv\_, kqw\_) == \{kqt\_ in bounded-ops(\allowbreak{}kqv\_, kqw\_): is-compact-op(\allowbreak{}kqv\_, kqw\_, kqt\_)\})
\end{quote}
also: \texttt{compact-ops-membership} \\
mentioned by 5 result(s)

\smallskip

\noindent\textbf{\texttt{companion(f, a, n)}}\quad\textit{functoid}\index{companion(f, a, n)@\texttt{companion(f, a, n)}}\\
reads: the companion matrix of the monic polynomial with coefficients \$2 \\
definition:
\begin{quote}\ttfamily\footnotesize\raggedright\noindent
matof(\allowbreak{}n, n, vnb-lambda(\allowbreak{}[gci\_, gcj\_], cartesian(\allowbreak{}interval(\allowbreak{}1, n), interval(\allowbreak{}1, n)), if(\allowbreak{}gcj\_ = n, (\allowbreak{}neg(\allowbreak{}f))(\allowbreak{}a(\allowbreak{}gci\_)), if(\allowbreak{}gci\_ = succ(\allowbreak{}gcj\_), one(\allowbreak{}f), zero(\allowbreak{}f)))))
\end{quote}
characterized by \texttt{companion-entry-uba-as-ring}  \textit{(definitional)}:
\begin{quote}\ttfamily\footnotesize\raggedright\noindent
forall(\allowbreak{}[r], is-unital-banach-algebra(\allowbreak{}r) implies forall(\allowbreak{}[nv\_ in nn, av\_], forall(\allowbreak{}[i\_ in interval(\allowbreak{}1, nv\_)], av\_(\allowbreak{}i\_) in vec(\allowbreak{}r)) implies forall(\allowbreak{}[i\_ in interval(\allowbreak{}1, nv\_), j\_ in interval(\allowbreak{}1, nv\_)], entry(\allowbreak{}companion(\allowbreak{}uba-as-ring(\allowbreak{}r), av\_, nv\_), i\_, j\_) = if(\allowbreak{}j\_ = nv\_, (\allowbreak{}vneg(\allowbreak{}r))(\allowbreak{}av\_(\allowbreak{}i\_)), if(\allowbreak{}i\_ = succ(\allowbreak{}j\_), vone(\allowbreak{}r), vzero(\allowbreak{}r))))))
\end{quote}
also: \texttt{companion-in-mat-uba-as-ring} \texttt{companion-congruence-uba-as-ring} \texttt{companion-entry-uba-as-ring-uiba-as-uba} \texttt{companion-in-mat-uba-as-ring-uiba-as-uba} \texttt{companion-congruence-uba-as-ring-uiba-as-uba} \texttt{companion-entry} \texttt{companion-in-mat} \texttt{krylov-companion} ... (14 in all) \\
mentioned by 41 result(s)

\smallskip

\noindent\textbf{\texttt{complement-in}}\quad\textit{operator}\index{complement-in@\texttt{complement-in}}\\
a term-forming head \\
characterized by \texttt{complement-in-membership}  \textit{(primitive)}:
\begin{quote}\ttfamily\footnotesize\raggedright\noindent
forall(\allowbreak{}[a, b, x], x in complement-in(\allowbreak{}a, b) iff x in a and not(\allowbreak{}x in b))
\end{quote}
also: \texttt{complement-in-set-closure} \texttt{union-complement-in} \texttt{complement-in-double} \texttt{complement-in-subset} \texttt{complement-in-big-union} \texttt{complement-in-empty-set} \texttt{inter-complement-in-shift} \texttt{power-complement-in-member} ... (15 in all) \\
mentioned by 139 result(s)

\smallskip

\noindent\textbf{\texttt{completion(m)}}\quad\textit{functoid}\index{completion(m)@\texttt{completion(m)}}\\
definition:
\begin{quote}\ttfamily\footnotesize\raggedright\noindent
[quotient(\allowbreak{}cauchy-setoid(\allowbreak{}m)), completion-dist(\allowbreak{}m)]
\end{quote}
characterized by \texttt{rkt-completion-pts}  \textit{(proven)}:
\begin{quote}\ttfamily\footnotesize\raggedright\noindent
forall(\allowbreak{}[m], is-metric-space(\allowbreak{}m) implies pts(\allowbreak{}completion(\allowbreak{}m)) == quotient(\allowbreak{}cauchy-setoid(\allowbreak{}m)))
\end{quote}
also: \texttt{rko2-completion-dist} \texttt{completion-is-complete} \texttt{completion-is-metric-space} \\
mentioned by 36 result(s)

\smallskip

\noindent\textbf{\texttt{completion-dist(m)}}\quad\textit{functoid}\index{completion-dist(m)@\texttt{completion-dist(m)}}\\
definition:
\begin{quote}\ttfamily\footnotesize\raggedright\noindent
vnb-lambda(\allowbreak{}p, cartesian(\allowbreak{}quotient(\allowbreak{}cauchy-setoid(\allowbreak{}m)), quotient(\allowbreak{}cauchy-setoid(\allowbreak{}m))), iota(\allowbreak{}dval, dval in rr and forsome(\allowbreak{}[f, g], f in cseq(\allowbreak{}m) and g in cseq(\allowbreak{}m) and nth(\allowbreak{}1, p) = class(\allowbreak{}cauchy-setoid(\allowbreak{}m), f) and nth(\allowbreak{}2, p) = class(\allowbreak{}cauchy-setoid(\allowbreak{}m), g) and converges-to(\allowbreak{}rr-ms, dist-seq(\allowbreak{}m, f, g), dval))))
\end{quote}
characterized by \texttt{rko2-completion-dist}  \textit{(proven)}:
\begin{quote}\ttfamily\footnotesize\raggedright\noindent
forall(\allowbreak{}[m], is-metric-space(\allowbreak{}m) implies dist(\allowbreak{}completion(\allowbreak{}m)) == completion-dist(\allowbreak{}m))
\end{quote}
mentioned by 2 result(s)

\smallskip

\noindent\textbf{\texttt{completion-measure(omega, ca, mu)}}\quad\textit{functoid}\index{completion-measure(omega, ca, mu)@\texttt{completion-measure(omega, ca, mu)}}\\
reads: the completion of the measure \$3 on \$2 \\
definition:
\begin{quote}\ttfamily\footnotesize\raggedright\noindent
res(\allowbreak{}method-i-outer(\allowbreak{}omega, ca, mu), completion-sigma(\allowbreak{}omega, ca, mu))
\end{quote}
characterized by \texttt{completion-measure-value}  \textit{(proven)}:
\begin{quote}\ttfamily\footnotesize\raggedright\noindent
forall(\allowbreak{}[omega, ca, mu, aa, ns], is-measure(\allowbreak{}omega, ca, mu) implies aa in ca implies is-null-set(\allowbreak{}omega, ca, mu, ns) implies (\allowbreak{}completion-measure(\allowbreak{}omega, ca, mu))(\allowbreak{}union(\allowbreak{}aa, ns)) = mu(\allowbreak{}aa))
\end{quote}
also: \texttt{completion-measure-unfold} \texttt{completion-measure-unique} \texttt{completion-measure-extends} \texttt{completion-measure-complete} \texttt{completion-measure-is-measure} \texttt{completion-measure-is-cara-measure} \\
mentioned by 14 result(s)

\smallskip

\noindent\textbf{\texttt{completion-sigma(omega, ca, mu)}}\quad\textit{functoid}\index{completion-sigma(omega, ca, mu)@\texttt{completion-sigma(omega, ca, mu)}}\\
reads: the completion of the sigma-algebra \$2 with respect to \$3 \\
definition:
\begin{quote}\ttfamily\footnotesize\raggedright\noindent
\{cps\_ in power(\allowbreak{}omega): forsome(\allowbreak{}[cpa\_ in ca, cpn\_ in power(\allowbreak{}omega)], is-null-set(\allowbreak{}omega, ca, mu, cpn\_) and cps\_ = union(\allowbreak{}cpa\_, cpn\_))\}
\end{quote}
characterized by \texttt{completion-sigma-rep-in}  \textit{(proven)}:
\begin{quote}\ttfamily\footnotesize\raggedright\noindent
forall(\allowbreak{}[omega, ca, mu, aa, ns], is-measure(\allowbreak{}omega, ca, mu) implies aa in ca implies is-null-set(\allowbreak{}omega, ca, mu, ns) implies union(\allowbreak{}aa, ns) in completion-sigma(\allowbreak{}omega, ca, mu))
\end{quote}
also: \texttt{completion-sigma-unfold} \texttt{completion-sigma-algebra} \texttt{completion-sigma-contains} \texttt{completion-sigma-membership} \texttt{completion-sigma-subset-cara} \texttt{completion-sigma-subset-power} \\
mentioned by 20 result(s)

\smallskip

\noindent\textbf{\texttt{complex-banach-algebra}}\quad\textit{refinement}\index{complex-banach-algebra@\texttt{complex-banach-algebra}}\\
predicate \texttt{is-complex-banach-algebra}; same shape as complex-normed-algebra

\smallskip

\noindent\textbf{\texttt{complex-banach-space}}\quad\textit{refinement}\index{complex-banach-space@\texttt{complex-banach-space}}\\
predicate \texttt{is-complex-banach-space}; same shape as complex-normed-vector-space

\smallskip

\noindent\textbf{\texttt{complex-inner-product-space}}\quad\textit{structure}\index{complex-inner-product-space@\texttt{complex-inner-product-space}}\\
predicate \texttt{is-complex-inner-product-space}; slots: scal vec vadd vzero vneg act ip

\smallskip

\noindent\textbf{\texttt{complex-inner-product-space-as-module(s)}}\quad\textit{view}\index{complex-inner-product-space-as-module(s)@\texttt{complex-inner-product-space-as-module(s)}}\\
a complex-inner-product-space viewed as a module \\
characterized by \texttt{complex-inner-product-space-as-module-is-module}  \textit{(definitional)}:
\begin{quote}\ttfamily\footnotesize\raggedright\noindent
forall(\allowbreak{}[r], is-complex-inner-product-space(\allowbreak{}r) implies is-module(\allowbreak{}complex-inner-product-space-as-module(\allowbreak{}r)))
\end{quote}
also: \texttt{module-vector-ag-is-abelian-group-complex-inner-product-space-as-module} \texttt{abelian-group-as-monoid-is-monoid-module-vector-ag-complex-inner-product-space-as-module} \texttt{complex-inner-product-space-as-module-functorial} \\
mentioned by 32 result(s)

\smallskip

\noindent\textbf{\texttt{complex-integral(omega, ca, mu, u)}}\quad\textit{functoid}\index{complex-integral(omega, ca, mu, u)@\texttt{complex-integral(omega, ca, mu, u)}}\\
reads: the integral of the complex function \$4 with respect to \$3 \\
definition:
\begin{quote}\ttfamily\footnotesize\raggedright\noindent
rintegral(\allowbreak{}omega, ca, mu, vnb-lambda(\allowbreak{}cix\_, omega, real-part(\allowbreak{}u(\allowbreak{}cix\_)))) + 1i * rintegral(\allowbreak{}omega, ca, mu, vnb-lambda(\allowbreak{}cix\_, omega, imag-part(\allowbreak{}u(\allowbreak{}cix\_))))
\end{quote}
characterized by \texttt{complex-integral-unfold}  \textit{(proven)}:
\begin{quote}\ttfamily\footnotesize\raggedright\noindent
forall(\allowbreak{}[omega, ca, mu, u], complex-integral(\allowbreak{}omega, ca, mu, u) == rintegral(\allowbreak{}omega, ca, mu, vnb-lambda(\allowbreak{}cix\_, omega, real-part(\allowbreak{}u(\allowbreak{}cix\_)))) + 1i * rintegral(\allowbreak{}omega, ca, mu, vnb-lambda(\allowbreak{}cix\_, omega, imag-part(\allowbreak{}u(\allowbreak{}cix\_)))))
\end{quote}
mentioned by 4 result(s)

\smallskip

\noindent\textbf{\texttt{complex-normed-algebra}}\quad\textit{structure}\index{complex-normed-algebra@\texttt{complex-normed-algebra}}\\
predicate \texttt{is-complex-normed-algebra}; slots: scal vec vadd vzero vneg act vnrm vmul

\smallskip

\noindent\textbf{\texttt{complex-normed-vector-space}}\quad\textit{structure}\index{complex-normed-vector-space@\texttt{complex-normed-vector-space}}\\
predicate \texttt{is-complex-normed-vector-space}; slots: scal vec vadd vzero vneg act vnrm jop

\smallskip

\noindent\textbf{\texttt{complex-pre-hilbert-space}}\quad\textit{refinement}\index{complex-pre-hilbert-space@\texttt{complex-pre-hilbert-space}}\\
predicate \texttt{is-complex-pre-hilbert-space}; same shape as complex-semi-inner-product-space

\smallskip

\noindent\textbf{\texttt{complex-semi-inner-product-space}}\quad\textit{structure}\index{complex-semi-inner-product-space@\texttt{complex-semi-inner-product-space}}\\
predicate \texttt{is-complex-semi-inner-product-space}; slots: scal vec vadd vzero vneg act rip jop

\smallskip

\noindent\textbf{\texttt{complexify(v)}}\quad\textit{functoid}\index{complexify(v)@\texttt{complexify(v)}}\\
reads: the complexification of \$1 \\
definition:
\begin{quote}\ttfamily\footnotesize\raggedright\noindent
[scal(\allowbreak{}v), cartesian(\allowbreak{}vec(\allowbreak{}v), vec(\allowbreak{}v)), cx-vadd(\allowbreak{}v), [vzero(\allowbreak{}v), vzero(\allowbreak{}v)], cx-vneg(\allowbreak{}v), cx-act(\allowbreak{}v), cx-rip(\allowbreak{}v), cx-j(\allowbreak{}v)]
\end{quote}
characterized by \texttt{complexify-act}  \textit{(proven)}:
\begin{quote}\ttfamily\footnotesize\raggedright\noindent
forall(\allowbreak{}[v], is-semi-inner-product-space(\allowbreak{}v) implies act(\allowbreak{}complexify(\allowbreak{}v)) == cx-act(\allowbreak{}v))
\end{quote}
also: \texttt{complexify-jop} \texttt{complexify-rip} \texttt{complexify-sip} \texttt{complexify-vec} \texttt{complexify-scal} \texttt{complexify-vadd} \texttt{complexify-vneg} \texttt{complexify-vzero} ... (28 in all) \\
mentioned by 75 result(s)

\smallskip

\noindent\textbf{\texttt{compose(f, g)}}\quad\textit{functoid}\index{compose(f, g)@\texttt{compose(f, g)}}\\
definition:
\begin{quote}\ttfamily\footnotesize\raggedright\noindent
vnb-lambda(\allowbreak{}z\_, dom(\allowbreak{}g), f(\allowbreak{}g(\allowbreak{}z\_)))
\end{quote}
characterized by \texttt{lft-compose}  \textit{(proven)}:
\begin{quote}\ttfamily\footnotesize\raggedright\noindent
forall(\allowbreak{}[a in cc, b in cc, c in cc, d in cc, a2 in cc, b2 in cc, c2 in cc, d2 in cc], not(\allowbreak{}a * d - b * c = 0) implies not(\allowbreak{}a2 * d2 - b2 * c2 = 0) implies compose(\allowbreak{}lft(\allowbreak{}a, b, c, d), lft(\allowbreak{}a2, b2, c2, d2)) = lft(\allowbreak{}a * a2 + b * c2, a * b2 + b * d2, c * a2 + d * c2, c * b2 + d * d2))
\end{quote}
also: \texttt{pre-compose} \texttt{compose-type} \texttt{post-compose} \texttt{compose-apply} \texttt{compose-assoc} \texttt{compose-id-left} \texttt{hom-pid-compose} \texttt{compose-id-right} ... (45 in all) \\
mentioned by 170 result(s)

\smallskip

\noindent\textbf{\texttt{conj-closed(mnbw\_)}}\quad\textit{predicate}\index{conj-closed(mnbw_)@\texttt{conj-closed(mnbw\_)}}\\
reads: \$1 is closed under the sides of conjunctions \\
defined by \texttt{conj-closed} \\
characterized by \texttt{conj-closed}  \textit{(definitional)}:
\begin{quote}\ttfamily\footnotesize\raggedright\noindent
forall(\allowbreak{}[mnbw\_], conj-closed(\allowbreak{}mnbw\_) iff forall(\allowbreak{}[mnbce\_ in mnbw\_, mnbcf\_ in expr-set], conj-step(\allowbreak{}mnbce\_, mnbcf\_) implies mnbcf\_ in mnbw\_))
\end{quote}
mentioned by 4 result(s)

\smallskip

\noindent\textbf{\texttt{conj-reach(mnbx\_)}}\quad\textit{functoid}\index{conj-reach(mnbx_)@\texttt{conj-reach(mnbx\_)}}\\
reads: the conjuncts of \$1 \\
definition:
\begin{quote}\ttfamily\footnotesize\raggedright\noindent
\{mnbz\_ in expr-set: forall(\allowbreak{}[mnbv\_ in power(\allowbreak{}expr-set)], mnbx\_ in mnbv\_ implies conj-closed(\allowbreak{}mnbv\_) implies mnbz\_ in mnbv\_)\}
\end{quote}
characterized by \texttt{conj-reach-unfold}  \textit{(proven)}:
\begin{quote}\ttfamily\footnotesize\raggedright\noindent
forall(\allowbreak{}[mnbx\_], conj-reach(\allowbreak{}mnbx\_) == \{mnbz\_ in expr-set: forall(\allowbreak{}[mnbv\_ in power(\allowbreak{}expr-set)], mnbx\_ in mnbv\_ implies conj-closed(\allowbreak{}mnbv\_) implies mnbz\_ in mnbv\_)\})
\end{quote}
mentioned by 6 result(s)

\smallskip

\noindent\textbf{\texttt{conj-set(cjsa\_)}}\quad\textit{functoid}\index{conj-set(cjsa_)@\texttt{conj-set(cjsa\_)}}\\
reads: the set of complex conjugates of the elements of \$1 \\
definition:
\begin{quote}\ttfamily\footnotesize\raggedright\noindent
\{cjsz\_ in cc: conjugate(\allowbreak{}cjsz\_) in cjsa\_\}
\end{quote}
characterized by \texttt{conj-set-open}  \textit{(proven)}:
\begin{quote}\ttfamily\footnotesize\raggedright\noindent
forall(\allowbreak{}[rxa\_], is-open(\allowbreak{}nf-metric-space(\allowbreak{}cc-normed-field), rxa\_) implies is-open(\allowbreak{}nf-metric-space(\allowbreak{}cc-normed-field), conj-set(\allowbreak{}rxa\_)))
\end{quote}
also: \texttt{conj-set-unfold} \\
mentioned by 5 result(s)

\smallskip

\noindent\textbf{\texttt{conj-step(mnbx\_, mnby\_)}}\quad\textit{predicate}\index{conj-step(mnbx_, mnby_)@\texttt{conj-step(mnbx\_, mnby\_)}}\\
reads: \$2 is a side of the conjunction \$1 \\
defined by \texttt{conj-step} \\
characterized by \texttt{conj-step}  \textit{(definitional)}:
\begin{quote}\ttfamily\footnotesize\raggedright\noindent
forall(\allowbreak{}[mnbx\_, mnby\_], conj-step(\allowbreak{}mnbx\_, mnby\_) iff forsome(\allowbreak{}[mnbcp\_, mnbcq\_], mnbcp\_ in expr-set and mnbcq\_ in expr-set and mnbx\_ = enc-and(\allowbreak{}mnbcp\_, mnbcq\_) and (\allowbreak{}mnby\_ = mnbcp\_ or mnby\_ = mnbcq\_)))
\end{quote}
mentioned by 4 result(s)

\smallskip

\noindent\textbf{\texttt{conjugate}}\quad\textit{operator}\index{conjugate@\texttt{conjugate}}\\
a term-forming head \\
characterized by \texttt{cc-conjugate-i}  \textit{(primitive)}:
\begin{quote}\ttfamily\footnotesize\raggedright\noindent
conjugate(\allowbreak{}1i) = -1i
\end{quote}
also: \texttt{cc-conjugate-add} \texttt{cc-conjugate-mul} \texttt{cc-conjugate-closed} \texttt{cc-conjugate-fixes-rr} \texttt{cc-conjugate-involution} \texttt{cc-conjugate-neg} \texttt{cc-conjugate-sub} \texttt{cc-conjugate-zero} ... (15 in all) \\
mentioned by 216 result(s)

\smallskip

\noindent\textbf{\texttt{connected-component(ccs\_, ccx\_)}}\quad\textit{functoid}\index{connected-component(ccs_, ccx_)@\texttt{connected-component(ccs\_, ccx\_)}}\\
reads: the connected component of \$2 in \$1 \\
definition:
\begin{quote}\ttfamily\footnotesize\raggedright\noindent
big-union(\allowbreak{}ccc\_, \{ccc\_ in power(\allowbreak{}pts(\allowbreak{}ccs\_)): is-connected-t(\allowbreak{}ccs\_, ccc\_) and ccx\_ in ccc\_\}, ccc\_)
\end{quote}
characterized by \texttt{connected-component-elim}  \textit{(proven)}:
\begin{quote}\ttfamily\footnotesize\raggedright\noindent
forall(\allowbreak{}[s, x, y in connected-component(\allowbreak{}s, x)], forsome(\allowbreak{}[c], is-connected-t(\allowbreak{}s, c) and x in c and y in c))
\end{quote}
also: \texttt{connected-component-self} \texttt{connected-component-intro} \texttt{connected-component-unfold} \texttt{connected-component-self-map} \texttt{connected-component-connected} \\
mentioned by 21 result(s)

\smallskip

\noindent\textbf{\texttt{cons}}\quad\textit{operator}\index{cons@\texttt{cons}}\\
a term-forming head \\
characterized by \texttt{nth-cons-1}  \textit{(definitional)}:
\begin{quote}\ttfamily\footnotesize\raggedright\noindent
forall(\allowbreak{}[x, l], x in set implies nth(\allowbreak{}1, cons(\allowbreak{}x, l)) = x)
\end{quote}
also: \texttt{length-cons} \texttt{makeset-cons} \texttt{nth-cons-succ} \texttt{cons-in-tuples} \texttt{tuple-cons-decompose} \\
mentioned by 11 result(s)

\smallskip

\noindent\textbf{\texttt{cont-fns(s)}}\quad\textit{functoid}\index{cont-fns(s)@\texttt{cont-fns(s)}}\\
reads: the set of continuous real functions on \$1 \\
definition:
\begin{quote}\ttfamily\footnotesize\raggedright\noindent
\{cfh\_ in fun(\allowbreak{}pts(\allowbreak{}s), rr): is-continuous(\allowbreak{}s, rr-ms, cfh\_)\}
\end{quote}
characterized by \texttt{rmc-cont-fns-add}  \textit{(proven)}:
\begin{quote}\ttfamily\footnotesize\raggedright\noindent
forall(\allowbreak{}[s, f, g], f in cont-fns(\allowbreak{}s) implies g in cont-fns(\allowbreak{}s) implies vnb-lambda(\allowbreak{}sx\_, pts(\allowbreak{}s), f(\allowbreak{}sx\_) + g(\allowbreak{}sx\_)) in cont-fns(\allowbreak{}s))
\end{quote}
also: \texttt{rmc-cont-fns-mem} \texttt{ck-cont-fns-is-set} \texttt{rmc-cont-fns-intro} \texttt{rmc-cont-fns-parts} \texttt{rmc-cont-fns-scale} \texttt{ck-cont-fns-subset-fun} \\
mentioned by 201 result(s)

\smallskip

\noindent\textbf{\texttt{content-of(ssd\_)}}\quad\textit{functoid}\index{content-of(ssd_)@\texttt{content-of(ssd\_)}}\\
reads: the content of the directed message \$1 \\
definition:
\begin{quote}\ttfamily\footnotesize\raggedright\noindent
nth(\allowbreak{}2, ssd\_)
\end{quote}

\smallskip

\noindent\textbf{\texttt{conv-hull(cvm\_, cve\_)}}\quad\textit{functoid}\index{conv-hull(cvm_, cve_)@\texttt{conv-hull(cvm\_, cve\_)}}\\
reads: the convex hull of \$2 in \$1 \\
definition:
\begin{quote}\ttfamily\footnotesize\raggedright\noindent
big-union(\allowbreak{}cvn\_, nn, conv-level(\allowbreak{}cvm\_, cve\_, cvn\_))
\end{quote}
characterized by \texttt{conv-hull-elim}  \textit{(proven)}:
\begin{quote}\ttfamily\footnotesize\raggedright\noindent
forall(\allowbreak{}[m, e, z in conv-hull(\allowbreak{}m, e)], forsome(\allowbreak{}[k in nn], z in conv-level(\allowbreak{}m, e, k)))
\end{quote}
also: \texttt{conv-hull-intro} \texttt{conv-hull-least} \texttt{conv-hull-convex} \texttt{conv-hull-unfold} \texttt{conv-hull-sub-vec} \texttt{conv-hull-contains} \texttt{conv-hull-convex-step} \texttt{conv-hull-totally-bounded} ... (9 in all) \\
mentioned by 46 result(s)

\smallskip

\noindent\textbf{\texttt{conv-join(cjm\_, cjc\_, cja\_)}}\quad\textit{functoid}\index{conv-join(cjm_, cjc_, cja_)@\texttt{conv-join(cjm\_, cjc\_, cja\_)}}\\
reads: the join of \$2 to the point \$3 in \$1 \\
definition:
\begin{quote}\ttfamily\footnotesize\raggedright\noindent
\{cjz\_ in vec(\allowbreak{}cjm\_): forsome(\allowbreak{}[cjy\_ in cjc\_, cjt\_ in rr], 0 <= cjt\_ and cjt\_ <= 1 and cjz\_ = (\allowbreak{}vadd(\allowbreak{}cjm\_))(\allowbreak{}(\allowbreak{}act(\allowbreak{}cjm\_))(\allowbreak{}1 - cjt\_, cjy\_), (\allowbreak{}act(\allowbreak{}cjm\_))(\allowbreak{}cjt\_, cja\_)))\}
\end{quote}
characterized by \texttt{conv-join-elim}  \textit{(proven)}:
\begin{quote}\ttfamily\footnotesize\raggedright\noindent
forall(\allowbreak{}[m, c, a, z in conv-join(\allowbreak{}m, c, a)], z in vec(\allowbreak{}m) and forsome(\allowbreak{}[y in c, t in rr], 0 <= t and t <= 1 and z = (\allowbreak{}vadd(\allowbreak{}m))(\allowbreak{}(\allowbreak{}act(\allowbreak{}m))(\allowbreak{}1 - t, y), (\allowbreak{}act(\allowbreak{}m))(\allowbreak{}t, a))))
\end{quote}
also: \texttt{conv-join-intro} \texttt{conv-join-convex} \texttt{conv-join-unfold} \texttt{conv-join-totally-bounded} \texttt{conv-join-with-point-convex} \\
mentioned by 7 result(s)

\smallskip

\noindent\textbf{\texttt{conv-level}}\quad\textit{functoid}\index{conv-level@\texttt{conv-level}}\\
reads: the convex combinations of \$3 + 1 points of \$2 in \$1 \\
defined by \texttt{conv-level-zero} \texttt{conv-level-succ} \\
characterized by \texttt{conv-level-succ}  \textit{(definitional)}:
\begin{quote}\ttfamily\footnotesize\raggedright\noindent
forall(\allowbreak{}[cvm\_, cve\_, cvn\_ in nn], conv-level(\allowbreak{}cvm\_, cve\_, succ(\allowbreak{}cvn\_)) == \{cvz\_ in vec(\allowbreak{}cvm\_): forsome(\allowbreak{}[cvx\_ in conv-level(\allowbreak{}cvm\_, cve\_, cvn\_), cvy\_ in cve\_, cvt\_ in rr], 0 <= cvt\_ and cvt\_ <= 1 and cvz\_ = (\allowbreak{}vadd(\allowbreak{}cvm\_))(\allowbreak{}(\allowbreak{}act(\allowbreak{}cvm\_))(\allowbreak{}1 - cvt\_, cvx\_), (\allowbreak{}act(\allowbreak{}cvm\_))(\allowbreak{}cvt\_, cvy\_)))\})
\end{quote}
also: \texttt{conv-level-zero} \texttt{conv-level-is-set} \texttt{conv-level-sub-vec} \texttt{conv-level-in-convex} \texttt{conv-level-succ-elim} \texttt{conv-level-succ-intro} \\
mentioned by 15 result(s)

\smallskip

\noindent\textbf{\texttt{conv3(cva\_, cvb\_, cvc\_)}}\quad\textit{functoid}\index{conv3(cva_, cvb_, cvc_)@\texttt{conv3(cva\_, cvb\_, cvc\_)}}\\
reads: the convex hull of \$1, \$2, \$3 \\
definition:
\begin{quote}\ttfamily\footnotesize\raggedright\noindent
\{cvz\_ in cc: forsome(\allowbreak{}[cvl\_ in ccint(\allowbreak{}0, 1), cvm\_ in ccint(\allowbreak{}0, 1), cvn\_ in ccint(\allowbreak{}0, 1)], (\allowbreak{}cvl\_ + cvm\_) + cvn\_ = 1 and cvz\_ = (\allowbreak{}cvl\_ * cva\_ + cvm\_ * cvb\_) + cvn\_ * cvc\_)\}
\end{quote}
characterized by \texttt{conv3-intro}  \textit{(proven)}:
\begin{quote}\ttfamily\footnotesize\raggedright\noindent
forall(\allowbreak{}[goua\_ in cc, goub\_ in cc, gouc\_ in cc, goul\_ in ccint(\allowbreak{}0, 1), goum\_ in ccint(\allowbreak{}0, 1), goun\_ in ccint(\allowbreak{}0, 1), gouz\_], (\allowbreak{}goul\_ + goum\_) + goun\_ = 1 implies gouz\_ = (\allowbreak{}goul\_ * goua\_ + goum\_ * goub\_) + goun\_ * gouc\_ implies gouz\_ in conv3(\allowbreak{}goua\_, goub\_, gouc\_))
\end{quote}
also: \texttt{conv3-mid-in} \texttt{bary-in-conv3} \texttt{conv3-sub-hull} \texttt{conv3-in-convex} \texttt{conv3-is-convex} \texttt{conv3-seg-trace} \texttt{conv3-member-iff} \texttt{conv3-near-vertex} ... (15 in all) \\
mentioned by 65 result(s)

\smallskip

\noindent\textbf{\texttt{converges(s, f)}}\quad\textit{predicate}\index{converges(s, f)@\texttt{converges(s, f)}}\\
reads: \$2 converges in \$1 \\
defined by \texttt{converges} \\
characterized by \texttt{converges}  \textit{(definitional)}:
\begin{quote}\ttfamily\footnotesize\raggedright\noindent
forall(\allowbreak{}[s, f], converges(\allowbreak{}s, f) iff forsome(\allowbreak{}[l], converges-to(\allowbreak{}s, f, l)))
\end{quote}
also: \texttt{converges-on} \texttt{ps-converges-at} \texttt{series-converges} \texttt{cc-series-converges} \texttt{complete-cauchy-converges} \texttt{ban-nsum-converges} \texttt{rr-cauchy-converges} \texttt{baire-cauchy-converges} ... (13 in all) \\
mentioned by 30 result(s)

\smallskip

\noindent\textbf{\texttt{converges-along(s, g, b, p)}}\quad\textit{predicate}\index{converges-along(s, g, b, p)@\texttt{converges-along(s, g, b, p)}}\\
reads: \$2 converges to \$4 along \$3 in \$1 \\
defined by \texttt{converges-along} \\
characterized by \texttt{converges-along}  \textit{(definitional)}:
\begin{quote}\ttfamily\footnotesize\raggedright\noindent
forall(\allowbreak{}[s, g, b, p], converges-along(\allowbreak{}s, g, b, p) iff is-metric-space(\allowbreak{}s) and g in fun(\allowbreak{}nn, pts(\allowbreak{}s)) and p in pts(\allowbreak{}s) and forall(\allowbreak{}[eps], pos-rr(\allowbreak{}eps) implies forsome(\allowbreak{}[n in nn], forall(\allowbreak{}[i in nn], i in b implies n <= i implies (\allowbreak{}dist(\allowbreak{}s))(\allowbreak{}g(\allowbreak{}i), p) <= eps))))
\end{quote}
mentioned by 7 result(s)

\smallskip

\noindent\textbf{\texttt{converges-on(s, fam, dseq)}}\quad\textit{predicate}\index{converges-on(s, fam, dseq)@\texttt{converges-on(s, fam, dseq)}}\\
reads: \$2 converges at every point of the sequence \$3 in \$1 \\
defined by \texttt{converges-on} \\
characterized by \texttt{converges-on}  \textit{(definitional)}:
\begin{quote}\ttfamily\footnotesize\raggedright\noindent
forall(\allowbreak{}[s, fam, dseq], converges-on(\allowbreak{}s, fam, dseq) iff is-metric-space(\allowbreak{}s) and fam in fun(\allowbreak{}nn, fun(\allowbreak{}pts(\allowbreak{}s), rr)) and dseq in fun(\allowbreak{}nn, pts(\allowbreak{}s)) and forall(\allowbreak{}[m in nn], converges(\allowbreak{}rr-ms, vnb-lambda(\allowbreak{}k, nn, (\allowbreak{}fam(\allowbreak{}k))(\allowbreak{}dseq(\allowbreak{}m))))))
\end{quote}
mentioned by 7 result(s)

\smallskip

\noindent\textbf{\texttt{converges-ptwise-to(cvps\_, cvpt\_, cvpf\_, cvpg\_)}}\quad\textit{predicate}\index{converges-ptwise-to(cvps_, cvpt_, cvpf_, cvpg_)@\texttt{converges-ptwise-to(cvps\_, cvpt\_, cvpf\_, cvpg\_)}}\\
reads: the sequence \$3 converges pointwise to \$4 as functions from \$1 to \$2 \\
defined by \texttt{converges-ptwise-to} \\
characterized by \texttt{converges-ptwise-to}  \textit{(definitional)}:
\begin{quote}\ttfamily\footnotesize\raggedright\noindent
forall(\allowbreak{}[cvps\_, cvpt\_, cvpf\_, cvpg\_], converges-ptwise-to(\allowbreak{}cvps\_, cvpt\_, cvpf\_, cvpg\_) iff is-metric-space(\allowbreak{}cvps\_) and is-metric-space(\allowbreak{}cvpt\_) and cvpf\_ in fun(\allowbreak{}nn, fun(\allowbreak{}pts(\allowbreak{}cvps\_), pts(\allowbreak{}cvpt\_))) and cvpg\_ in fun(\allowbreak{}pts(\allowbreak{}cvps\_), pts(\allowbreak{}cvpt\_)) and forall(\allowbreak{}[cvpx\_ in pts(\allowbreak{}cvps\_), cvpe\_], pos-rr(\allowbreak{}cvpe\_) implies forsome(\allowbreak{}[cvpn\_ in nn], forall(\allowbreak{}[cvpk\_ in nn], cvpn\_ <= cvpk\_ implies (\allowbreak{}dist(\allowbreak{}cvpt\_))(\allowbreak{}(\allowbreak{}cvpf\_(\allowbreak{}cvpk\_))(\allowbreak{}cvpx\_), cvpg\_(\allowbreak{}cvpx\_)) <= cvpe\_))))
\end{quote}
mentioned by 15 result(s)

\smallskip

\noindent\textbf{\texttt{converges-to(s, f, l)}}\quad\textit{predicate}\index{converges-to(s, f, l)@\texttt{converges-to(s, f, l)}}\\
reads: \$2 converges to \$3 in \$1 \\
defined by \texttt{converges-to} \\
characterized by \texttt{converges-to}  \textit{(definitional)}:
\begin{quote}\ttfamily\footnotesize\raggedright\noindent
forall(\allowbreak{}[s, f, l], converges-to(\allowbreak{}s, f, l) iff is-metric-space(\allowbreak{}s) and f in fun(\allowbreak{}nn, pts(\allowbreak{}s)) and l in pts(\allowbreak{}s) and forall(\allowbreak{}[eps], pos-rr(\allowbreak{}eps) implies forsome(\allowbreak{}[n in nn], forall(\allowbreak{}[n\_ in nn], n <= n\_ implies (\allowbreak{}dist(\allowbreak{}s))(\allowbreak{}f(\allowbreak{}n\_), l) <= eps))))
\end{quote}
also: \texttt{prod-converges-to} \texttt{ps-converges-to-at} \texttt{series-converges-to} \texttt{cc-series-converges-to} \texttt{rr-converges-to-abs} \texttt{converges-to-transfer} \texttt{pullback-converges-to} \texttt{rr-const-converges-to} ... (18 in all) \\
mentioned by 296 result(s)

\smallskip

\noindent\textbf{\texttt{converges-uniformly(s, seq, g)}}\quad\textit{predicate}\index{converges-uniformly(s, seq, g)@\texttt{converges-uniformly(s, seq, g)}}\\
reads: \$2 converges uniformly to \$3 on \$1 \\
defined by \texttt{converges-uniformly} \\
characterized by \texttt{converges-uniformly}  \textit{(definitional)}:
\begin{quote}\ttfamily\footnotesize\raggedright\noindent
forall(\allowbreak{}[s, seq, g], converges-uniformly(\allowbreak{}s, seq, g) iff is-metric-space(\allowbreak{}s) and seq in fun(\allowbreak{}nn, fun(\allowbreak{}pts(\allowbreak{}s), rr)) and g in fun(\allowbreak{}pts(\allowbreak{}s), rr) and forall(\allowbreak{}[eps], pos-rr(\allowbreak{}eps) implies forsome(\allowbreak{}[cap in nn], forall(\allowbreak{}[k], k in nn and cap <= k implies forall(\allowbreak{}[x in pts(\allowbreak{}s)], abs(\allowbreak{}(\allowbreak{}seq(\allowbreak{}k))(\allowbreak{}x) - g(\allowbreak{}x)) < eps)))))
\end{quote}
mentioned by 8 result(s)

\smallskip

\noindent\textbf{\texttt{converges-uniformly-to(cvus\_, cvut\_, cvuf\_, cvug\_)}}\quad\textit{predicate}\index{converges-uniformly-to(cvus_, cvut_, cvuf_, cvug_)@\texttt{converges-uniformly-to(cvus\_, cvut\_, cvuf\_, cvug\_)}}\\
reads: the sequence \$3 converges uniformly to \$4 as functions from \$1 to \$2 \\
defined by \texttt{converges-uniformly-to} \\
characterized by \texttt{converges-uniformly-to}  \textit{(definitional)}:
\begin{quote}\ttfamily\footnotesize\raggedright\noindent
forall(\allowbreak{}[cvus\_, cvut\_, cvuf\_, cvug\_], converges-uniformly-to(\allowbreak{}cvus\_, cvut\_, cvuf\_, cvug\_) iff is-metric-space(\allowbreak{}cvus\_) and is-metric-space(\allowbreak{}cvut\_) and cvuf\_ in fun(\allowbreak{}nn, fun(\allowbreak{}pts(\allowbreak{}cvus\_), pts(\allowbreak{}cvut\_))) and cvug\_ in fun(\allowbreak{}pts(\allowbreak{}cvus\_), pts(\allowbreak{}cvut\_)) and forall(\allowbreak{}[cvue\_], pos-rr(\allowbreak{}cvue\_) implies forsome(\allowbreak{}[cvun\_ in nn], forall(\allowbreak{}[cvuk\_ in nn], cvun\_ <= cvuk\_ implies forall(\allowbreak{}[cvux\_ in pts(\allowbreak{}cvus\_)], (\allowbreak{}dist(\allowbreak{}cvut\_))(\allowbreak{}(\allowbreak{}cvuf\_(\allowbreak{}cvuk\_))(\allowbreak{}cvux\_), cvug\_(\allowbreak{}cvux\_)) <= cvue\_)))))
\end{quote}
mentioned by 10 result(s)

\smallskip

\noindent\textbf{\texttt{coord-permute(cpx\_, cpg\_)}}\quad\textit{functoid}\index{coord-permute(cpx_, cpg_)@\texttt{coord-permute(cpx\_, cpg\_)}}\\
reads: the coordinatewise permutation of the sequences in \$1 by \$2 \\
definition:
\begin{quote}\ttfamily\footnotesize\raggedright\noindent
vnb-lambda(\allowbreak{}cpw\_, fun(\allowbreak{}nn, cpx\_), vnb-lambda(\allowbreak{}cpk\_, nn, (\allowbreak{}cpg\_(\allowbreak{}cpk\_))(\allowbreak{}cpw\_(\allowbreak{}cpk\_))))
\end{quote}
characterized by \texttt{coord-permute-transitive}  \textit{(proven)}:
\begin{quote}\ttfamily\footnotesize\raggedright\noindent
forall(\allowbreak{}[sx\_ in set, xs in fun(\allowbreak{}nn, sx\_), ys in fun(\allowbreak{}nn, sx\_)], forsome(\allowbreak{}[sg in fun(\allowbreak{}nn, bijection(\allowbreak{}sx\_, sx\_))], (\allowbreak{}coord-permute(\allowbreak{}sx\_, sg))(\allowbreak{}xs) = ys))
\end{quote}
mentioned by 1 result(s)

\smallskip

\noindent\textbf{\texttt{core-closed(mct\_, mcs\_)}}\quad\textit{predicate}\index{core-closed(mct_, mcs_)@\texttt{core-closed(mct\_, mcs\_)}}\\
reads: \$2 is closed under the core rules over \$1 \\
defined by \texttt{core-closed} \\
characterized by \texttt{core-closed}  \textit{(definitional)}:
\begin{quote}\ttfamily\footnotesize\raggedright\noindent
forall(\allowbreak{}[mct\_, mcs\_], core-closed(\allowbreak{}mct\_, mcs\_) iff forall(\allowbreak{}[mchs\_, mcc\_], core-rule(\allowbreak{}mct\_, mchs\_, mcc\_) implies hyps-in(\allowbreak{}mchs\_, mcs\_) implies mcc\_ in seq-set implies mcc\_ in mcs\_))
\end{quote}
also: \texttt{core-closed-apply} \\
mentioned by 6 result(s)

\smallskip

\noindent\textbf{\texttt{core-rule(mct\_, mchs\_, mcc\_)}}\quad\textit{predicate}\index{core-rule(mct_, mchs_, mcc_)@\texttt{core-rule(mct\_, mchs\_, mcc\_)}}\\
reads: some core rule over \$1 relates \$2 to \$3 \\
defined by \texttt{core-rule} \\
characterized by \texttt{core-rule}  \textit{(definitional)}:
\begin{quote}\ttfamily\footnotesize\raggedright\noindent
forall(\allowbreak{}[mct\_, mchs\_, mcc\_], core-rule(\allowbreak{}mct\_, mchs\_, mcc\_) iff cr-and-intro(\allowbreak{}mct\_, mchs\_, mcc\_) or cr-and-elim(\allowbreak{}mct\_, mchs\_, mcc\_) or cr-or-intro-left(\allowbreak{}mct\_, mchs\_, mcc\_) or cr-or-intro-right(\allowbreak{}mct\_, mchs\_, mcc\_) or cr-or-elim(\allowbreak{}mct\_, mchs\_, mcc\_) or cr-implies-intro(\allowbreak{}mct\_, mchs\_, mcc\_) or cr-backchain(\allowbreak{}mct\_, mchs\_, mcc\_) or cr-not-intro(\allowbreak{}mct\_, mchs\_, mcc\_) or cr-not-elim(\allowbreak{}mct\_, mchs\_, mcc\_) or cr-iff-intro(\allowbreak{}mct\_, mchs\_, mcc\_) or cr-iff-elim(\allowbreak{}mct\_, mchs\_, mcc\_) or cr-forall-intro(\allowbreak{}mct\_, mchs\_, mcc\_) or cr-forall-elim(\allowbreak{}mct\_, mchs\_, mcc\_) or cr-forsome-intro(\allowbreak{}mct\_, mchs\_, mcc\_) or cr-forsome-elim(\allowbreak{}mct\_, mchs\_, mcc\_) or cr-assumption(\allowbreak{}mct\_, mchs\_, mcc\_) or cr-theorem-assumption(\allowbreak{}mct\_, mchs\_, mcc\_) or cr-cut(\allowbreak{}mct\_, mchs\_, mcc\_) or cr-weakening(\allowbreak{}mct\_, mchs\_, mcc\_) or cr-contradiction(\allowbreak{}mct\_, mchs\_, mcc\_) or cr-truth-intro(\allowbreak{}mct\_, mchs\_, mcc\_) or cr-eq-subst(\allowbreak{}mct\_, mchs\_, mcc\_) or cr-eq-subst-sym(\allowbreak{}mct\_, mchs\_, mcc\_) or cr-reflexivity(\allowbreak{}mct\_, mchs\_, mcc\_) or cr-quasi-reflexivity(\allowbreak{}mct\_, mchs\_, mcc\_) or cr-union-mem(\allowbreak{}mct\_, mchs\_, mcc\_) or cr-intersection-mem(\allowbreak{}mct\_, mchs\_, mcc\_) or cr-cartesian-mem(\allowbreak{}mct\_, mchs\_, mcc\_) or cr-tuples-mem(\allowbreak{}mct\_, mchs\_, mcc\_) or cr-nth-eq(\allowbreak{}mct\_, mchs\_, mcc\_) or cr-length-eq(\allowbreak{}mct\_, mchs\_, mcc\_) or cr-tuple-eq(\allowbreak{}mct\_, mchs\_, mcc\_) or cr-if-true(\allowbreak{}mct\_, mchs\_, mcc\_) or cr-if-false(\allowbreak{}mct\_, mchs\_, mcc\_) or cr-sep-sethood(\allowbreak{}mct\_, mchs\_, mcc\_) or cr-sep-mem(\allowbreak{}mct\_, mchs\_, mcc\_) or cr-comp-mem(\allowbreak{}mct\_, mchs\_, mcc\_) or cr-big-union-sethood(\allowbreak{}mct\_, mchs\_, mcc\_) or cr-big-union-mem(\allowbreak{}mct\_, mchs\_, mcc\_) or cr-iota-def(\allowbreak{}mct\_, mchs\_, mcc\_) or cr-iota-in-elim(\allowbreak{}mct\_, mchs\_, mcc\_) or cr-lambda-type(\allowbreak{}mct\_, mchs\_, mcc\_) or cr-lambda-beta-eq(\allowbreak{}mct\_, mchs\_, mcc\_) or cr-def-strict-in-l(\allowbreak{}mct\_, mchs\_, mcc\_) or cr-def-strict-in-r(\allowbreak{}mct\_, mchs\_, mcc\_) or cr-def-total(\allowbreak{}mct\_, mchs\_, mcc\_) or cr-def-binder(\allowbreak{}mct\_, mchs\_, mcc\_) or cr-def-comp(\allowbreak{}mct\_, mchs\_, mcc\_) or cr-def-list(\allowbreak{}mct\_, mchs\_, mcc\_) or cr-def-strict-arg(\allowbreak{}mct\_, mchs\_, mcc\_) or cr-def-strict-dom(\allowbreak{}mct\_, mchs\_, mcc\_) or cr-def-strict-pred(\allowbreak{}mct\_, mchs\_, mcc\_) or cr-iota-cr(\allowbreak{}mct\_, mchs\_, mcc\_) or cr-tbind-cr(\allowbreak{}mct\_, mchs\_, mcc\_) or cr-functoid-beta-eq(\allowbreak{}mct\_, mchs\_, mcc\_) or cr-tuple-eta(\allowbreak{}mct\_, mchs\_, mcc\_))
\end{quote}
mentioned by 6 result(s)

\smallskip

\noindent\textbf{\texttt{cos}}\quad\textit{operator}\index{cos@\texttt{cos}}\\
a term-forming head

\smallskip

\noindent\textbf{\texttt{coset-norm(bqna\_, bqni\_, bqnc\_)}}\quad\textit{functoid}\index{coset-norm(bqna_, bqni_, bqnc_)@\texttt{coset-norm(bqna\_, bqni\_, bqnc\_)}}\\
reads: the norm of the coset \$3 of \$2 in \$1 \\
definition:
\begin{quote}\ttfamily\footnotesize\raggedright\noindent
einf(\allowbreak{}\{bqnw\_ in rr: forsome(\allowbreak{}[bqnu\_ in bqnc\_], bqnw\_ = (\allowbreak{}vnrm(\allowbreak{}bqna\_))(\allowbreak{}bqnu\_))\})
\end{quote}
characterized by \texttt{coset-norm-le}  \textit{(proven)}:
\begin{quote}\ttfamily\footnotesize\raggedright\noindent
forall(\allowbreak{}[a], is-unital-banach-algebra(\allowbreak{}a) implies forall(\allowbreak{}[ii], is-two-sided-ideal(\allowbreak{}a, ii) implies forall(\allowbreak{}[c in quotient(\allowbreak{}qmod-setoid(\allowbreak{}uba-as-module(\allowbreak{}a), ii)), u in c], coset-norm(\allowbreak{}a, ii, c) <= (\allowbreak{}vnrm(\allowbreak{}a))(\allowbreak{}u))))
\end{quote}
also: \texttt{coset-norm-real} \texttt{coset-norm-homog} \texttt{coset-norm-act-le} \texttt{coset-norm-approx} \texttt{coset-norm-submult} \texttt{coset-norm-definite} \texttt{coset-norm-triangle} \texttt{coset-norm-scale-greatest} ... (9 in all) \\
mentioned by 15 result(s)

\smallskip

\noindent\textbf{\texttt{count-hits}}\quad\textit{functoid}\index{count-hits@\texttt{count-hits}}\\
reads: the number of k < \$3 with \$1(k) = \$2 \\
defined by \texttt{count-hits-zero} \texttt{count-hits-succ} \\
characterized by \texttt{count-hits-succ}  \textit{(definitional)}:
\begin{quote}\ttfamily\footnotesize\raggedright\noindent
forall(\allowbreak{}[chs\_, chz\_, chn\_ in nn], count-hits(\allowbreak{}chs\_, chz\_, succ(\allowbreak{}chn\_)) == if(\allowbreak{}chs\_(\allowbreak{}chn\_) = chz\_, succ(\allowbreak{}count-hits(\allowbreak{}chs\_, chz\_, chn\_)), count-hits(\allowbreak{}chs\_, chz\_, chn\_)))
\end{quote}
also: \texttt{count-hits-zero} \texttt{count-hits-none} \texttt{count-hits-in-nn} \\
mentioned by 11 result(s)

\smallskip

\noindent\textbf{\texttt{cover-sums(omega, cls, pm, es)}}\quad\textit{functoid}\index{cover-sums(omega, cls, pm, es)@\texttt{cover-sums(omega, cls, pm, es)}}\\
reads: the set of sums of \$3 over the sequences in \$2 covering \$4 \\
definition:
\begin{quote}\ttfamily\footnotesize\raggedright\noindent
\{cvv\_ in rr-pos-star: forsome(\allowbreak{}[cvs\_ in fun(\allowbreak{}nn, cls)], es subset big-union(\allowbreak{}cvn\_, nn, cvs\_(\allowbreak{}cvn\_)) and cvv\_ = esum(\allowbreak{}vnb-lambda(\allowbreak{}cvk\_, nn, pm(\allowbreak{}cvs\_(\allowbreak{}cvk\_)))))\}
\end{quote}
characterized by \texttt{cover-sums-elim}  \textit{(proven)}:
\begin{quote}\ttfamily\footnotesize\raggedright\noindent
forall(\allowbreak{}[omega, cls, pm, es, v\_ in cover-sums(\allowbreak{}omega, cls, pm, es)], forsome(\allowbreak{}[cvs\_ in fun(\allowbreak{}nn, cls)], es subset big-union(\allowbreak{}cvn\_, nn, cvs\_(\allowbreak{}cvn\_)) and v\_ = esum(\allowbreak{}vnb-lambda(\allowbreak{}cvk\_, nn, pm(\allowbreak{}cvs\_(\allowbreak{}cvk\_))))))
\end{quote}
also: \texttt{cover-sums-intro} \texttt{cover-sums-subset} \texttt{cover-sums-unfold} \\
mentioned by 15 result(s)

\smallskip

\noindent\textbf{\texttt{cps-radius(cf)}}\quad\textit{functoid}\index{cps-radius(cf)@\texttt{cps-radius(cf)}}\\
reads: the radius of convergence of the power series with coefficients \$1 \\
definition:
\begin{quote}\ttfamily\footnotesize\raggedright\noindent
esup(\allowbreak{}\{cpsr\_ in rr: 0 <= cpsr\_ and series-converges(\allowbreak{}vnb-lambda(\allowbreak{}cpsk\_, nn, magnitude(\allowbreak{}cf(\allowbreak{}cpsk\_)) * cpsr\_ \textasciicircum{} cpsk\_))\})
\end{quote}
characterized by \texttt{cps-radius-in}  \textit{(proven)}:
\begin{quote}\ttfamily\footnotesize\raggedright\noindent
forall(\allowbreak{}[cf], cps-radius(\allowbreak{}cf) in rr-pos-star)
\end{quote}
also: \texttt{cps-radius-mono} \texttt{cps-radius-root} \texttt{cps-radius-zero} \texttt{cps-radius-recip} \texttt{cps-radius-approx} \texttt{cps-radius-factor} \texttt{cps-radius-unfold} \texttt{cps-radius-derived} ... (11 in all) \\
mentioned by 67 result(s)

\smallskip

\noindent\textbf{\texttt{cr-and-elim(mct\_, mchs\_, mcc\_)}}\quad\textit{predicate}\index{cr-and-elim(mct_, mchs_, mcc_)@\texttt{cr-and-elim(mct\_, mchs\_, mcc\_)}}\\
reads: the core and-elim relates \$2 to \$3 \\
defined by \texttt{cr-and-elim} \\
characterized by \texttt{cr-and-elim}  \textit{(definitional)}:
\begin{quote}\ttfamily\footnotesize\raggedright\noindent
forall(\allowbreak{}[mct\_, mchs\_, mcc\_], cr-and-elim(\allowbreak{}mct\_, mchs\_, mcc\_) iff forsome(\allowbreak{}[mcg\_, mcf\_, mca\_, mcb\_], mcc\_ = [mcg\_, mcf\_] and enc-and(\allowbreak{}mca\_, mcb\_) in mcg\_ and mchs\_ = [[ctx-add(\allowbreak{}ctx-add(\allowbreak{}mcg\_, mca\_), mcb\_), mcf\_]]))
\end{quote}
mentioned by 5 result(s)

\smallskip

\noindent\textbf{\texttt{cr-and-intro(mct\_, mchs\_, mcc\_)}}\quad\textit{predicate}\index{cr-and-intro(mct_, mchs_, mcc_)@\texttt{cr-and-intro(mct\_, mchs\_, mcc\_)}}\\
reads: the core and-intro relates \$2 to \$3 \\
defined by \texttt{cr-and-intro} \\
characterized by \texttt{cr-and-intro}  \textit{(definitional)}:
\begin{quote}\ttfamily\footnotesize\raggedright\noindent
forall(\allowbreak{}[mct\_, mchs\_, mcc\_], cr-and-intro(\allowbreak{}mct\_, mchs\_, mcc\_) iff forsome(\allowbreak{}[mcg\_, mca\_, mcb\_], mcc\_ = [mcg\_, enc-and(\allowbreak{}mca\_, mcb\_)] and mchs\_ = [[mcg\_, mca\_], [mcg\_, mcb\_]]))
\end{quote}
mentioned by 5 result(s)

\smallskip

\noindent\textbf{\texttt{cr-assumption(mct\_, mchs\_, mcc\_)}}\quad\textit{predicate}\index{cr-assumption(mct_, mchs_, mcc_)@\texttt{cr-assumption(mct\_, mchs\_, mcc\_)}}\\
reads: the core assumption relates \$2 to \$3 \\
defined by \texttt{cr-assumption} \\
characterized by \texttt{cr-assumption}  \textit{(definitional)}:
\begin{quote}\ttfamily\footnotesize\raggedright\noindent
forall(\allowbreak{}[mct\_, mchs\_, mcc\_], cr-assumption(\allowbreak{}mct\_, mchs\_, mcc\_) iff forsome(\allowbreak{}[mcg\_, mca\_], mcc\_ = [mcg\_, mca\_] and mca\_ in mcg\_ and mchs\_ = []))
\end{quote}
mentioned by 5 result(s)

\smallskip

\noindent\textbf{\texttt{cr-backchain(mct\_, mchs\_, mcc\_)}}\quad\textit{predicate}\index{cr-backchain(mct_, mchs_, mcc_)@\texttt{cr-backchain(mct\_, mchs\_, mcc\_)}}\\
reads: the core backchain relates \$2 to \$3 \\
defined by \texttt{cr-backchain} \\
characterized by \texttt{cr-backchain}  \textit{(definitional)}:
\begin{quote}\ttfamily\footnotesize\raggedright\noindent
forall(\allowbreak{}[mct\_, mchs\_, mcc\_], cr-backchain(\allowbreak{}mct\_, mchs\_, mcc\_) iff forsome(\allowbreak{}[mcg\_, mcf\_, mca\_], mcc\_ = [mcg\_, mcf\_] and enc-imp(\allowbreak{}mca\_, mcf\_) in mcg\_ and mchs\_ = [[mcg\_, mca\_]]))
\end{quote}
mentioned by 5 result(s)

\smallskip

\noindent\textbf{\texttt{cr-big-union-mem(mct\_, mchs\_, mcc\_)}}\quad\textit{predicate}\index{cr-big-union-mem(mct_, mchs_, mcc_)@\texttt{cr-big-union-mem(mct\_, mchs\_, mcc\_)}}\\
reads: the core big-union-mem relates \$2 to \$3 \\
defined by \texttt{cr-big-union-mem} \\
characterized by \texttt{cr-big-union-mem}  \textit{(definitional)}:
\begin{quote}\ttfamily\footnotesize\raggedright\noindent
forall(\allowbreak{}[mct\_, mchs\_, mcc\_], cr-big-union-mem(\allowbreak{}mct\_, mchs\_, mcc\_) iff forsome(\allowbreak{}[mcg\_, mca\_, mcb\_, mcx\_], mcc\_ = [mcg\_, enc-iff(\allowbreak{}enc-in(\allowbreak{}mcx\_, enc-bu(\allowbreak{}mca\_, mcb\_)), enc-ex(\allowbreak{}enc-and(\allowbreak{}enc-in(\allowbreak{}ebvar(\allowbreak{}0), mca\_), enc-in(\allowbreak{}mcx\_, mcb\_))))] and nbg-ok(\allowbreak{}mcb\_) and mchs\_ = []))
\end{quote}
mentioned by 5 result(s)

\smallskip

\noindent\textbf{\texttt{cr-big-union-sethood(mct\_, mchs\_, mcc\_)}}\quad\textit{predicate}\index{cr-big-union-sethood(mct_, mchs_, mcc_)@\texttt{cr-big-union-sethood(mct\_, mchs\_, mcc\_)}}\\
reads: the core big-union-sethood relates \$2 to \$3 \\
defined by \texttt{cr-big-union-sethood} \\
characterized by \texttt{cr-big-union-sethood}  \textit{(definitional)}:
\begin{quote}\ttfamily\footnotesize\raggedright\noindent
forall(\allowbreak{}[mct\_, mchs\_, mcc\_], cr-big-union-sethood(\allowbreak{}mct\_, mchs\_, mcc\_) iff forsome(\allowbreak{}[mcg\_, mca\_, mcb\_], mcc\_ = [mcg\_, enc-in(\allowbreak{}enc-bu(\allowbreak{}mca\_, mcb\_), enc-setc)] and nbg-ok(\allowbreak{}mcb\_) and mchs\_ = [[mcg\_, enc-in(\allowbreak{}mca\_, enc-setc)], [mcg\_, enc-all(\allowbreak{}enc-imp(\allowbreak{}enc-in(\allowbreak{}ebvar(\allowbreak{}0), mca\_), enc-in(\allowbreak{}mcb\_, enc-setc)))]]))
\end{quote}
mentioned by 5 result(s)

\smallskip

\noindent\textbf{\texttt{cr-cartesian-mem(mct\_, mchs\_, mcc\_)}}\quad\textit{predicate}\index{cr-cartesian-mem(mct_, mchs_, mcc_)@\texttt{cr-cartesian-mem(mct\_, mchs\_, mcc\_)}}\\
reads: the core cartesian-mem relates \$2 to \$3 \\
defined by \texttt{cr-cartesian-mem} \\
characterized by \texttt{cr-cartesian-mem}  \textit{(definitional)}:
\begin{quote}\ttfamily\footnotesize\raggedright\noindent
forall(\allowbreak{}[mct\_, mchs\_, mcc\_], cr-cartesian-mem(\allowbreak{}mct\_, mchs\_, mcc\_) iff forsome(\allowbreak{}[mcg\_, mcx\_, mcl\_], mcc\_ = [mcg\_, enc-iff(\allowbreak{}enc-in(\allowbreak{}mcx\_, enc-cart(\allowbreak{}mcl\_)), enc-and(\allowbreak{}enc-in(\allowbreak{}mcx\_, enc-tuples(\allowbreak{}enc-setc)), enc-and(\allowbreak{}enc-eq(\allowbreak{}enc-length(\allowbreak{}mcx\_), enc-num(\allowbreak{}length(\allowbreak{}mcl\_))), enc-cartl(\allowbreak{}mcx\_, mcl\_, length(\allowbreak{}mcl\_)))))] and 1 <= length(\allowbreak{}mcl\_) and mchs\_ = []))
\end{quote}
mentioned by 5 result(s)

\smallskip

\noindent\textbf{\texttt{cr-comp-mem(mct\_, mchs\_, mcc\_)}}\quad\textit{predicate}\index{cr-comp-mem(mct_, mchs_, mcc_)@\texttt{cr-comp-mem(mct\_, mchs\_, mcc\_)}}\\
reads: the core comp-mem relates \$2 to \$3 \\
defined by \texttt{cr-comp-mem} \\
characterized by \texttt{cr-comp-mem}  \textit{(definitional)}:
\begin{quote}\ttfamily\footnotesize\raggedright\noindent
forall(\allowbreak{}[mct\_, mchs\_, mcc\_], cr-comp-mem(\allowbreak{}mct\_, mchs\_, mcc\_) iff forsome(\allowbreak{}[mcg\_, mcp\_, mctm\_], mcc\_ = [mcg\_, enc-iff(\allowbreak{}enc-in(\allowbreak{}mctm\_, enc-comp(\allowbreak{}mcp\_)), enc-and(\allowbreak{}enc-in(\allowbreak{}mctm\_, enc-setc), expr-open(\allowbreak{}mcp\_, 0, mctm\_)))] and nbg-ok(\allowbreak{}mcp\_) and mchs\_ = []))
\end{quote}
mentioned by 5 result(s)

\smallskip

\noindent\textbf{\texttt{cr-contradiction(mct\_, mchs\_, mcc\_)}}\quad\textit{predicate}\index{cr-contradiction(mct_, mchs_, mcc_)@\texttt{cr-contradiction(mct\_, mchs\_, mcc\_)}}\\
reads: the core proof-by-contradiction relates \$2 to \$3 \\
defined by \texttt{cr-contradiction} \\
characterized by \texttt{cr-contradiction}  \textit{(definitional)}:
\begin{quote}\ttfamily\footnotesize\raggedright\noindent
forall(\allowbreak{}[mct\_, mchs\_, mcc\_], cr-contradiction(\allowbreak{}mct\_, mchs\_, mcc\_) iff forsome(\allowbreak{}[mcg\_, mcf\_], mcc\_ = [mcg\_, mcf\_] and mchs\_ = [[ctx-add(\allowbreak{}mcg\_, enc-not(\allowbreak{}mcf\_)), enc-falsity]]))
\end{quote}
mentioned by 5 result(s)

\smallskip

\noindent\textbf{\texttt{cr-cut(mct\_, mchs\_, mcc\_)}}\quad\textit{predicate}\index{cr-cut(mct_, mchs_, mcc_)@\texttt{cr-cut(mct\_, mchs\_, mcc\_)}}\\
reads: the core cut relates \$2 to \$3 \\
defined by \texttt{cr-cut} \\
characterized by \texttt{cr-cut}  \textit{(definitional)}:
\begin{quote}\ttfamily\footnotesize\raggedright\noindent
forall(\allowbreak{}[mct\_, mchs\_, mcc\_], cr-cut(\allowbreak{}mct\_, mchs\_, mcc\_) iff forsome(\allowbreak{}[mcg\_, mcf\_, mca\_], mcc\_ = [mcg\_, mcf\_] and mchs\_ = [[mcg\_, mca\_], [ctx-add(\allowbreak{}mcg\_, mca\_), mcf\_]]))
\end{quote}
mentioned by 5 result(s)

\smallskip

\noindent\textbf{\texttt{cr-def-binder(mct\_, mchs\_, mcc\_)}}\quad\textit{predicate}\index{cr-def-binder(mct_, mchs_, mcc_)@\texttt{cr-def-binder(mct\_, mchs\_, mcc\_)}}\\
reads: the core definedness rule binder-domain relates \$2 to \$3 \\
defined by \texttt{cr-def-binder} \\
characterized by \texttt{cr-def-binder}  \textit{(definitional)}:
\begin{quote}\ttfamily\footnotesize\raggedright\noindent
forall(\allowbreak{}[mct\_, mchs\_, mcc\_], cr-def-binder(\allowbreak{}mct\_, mchs\_, mcc\_) iff forsome(\allowbreak{}[mcg\_, mchd\_, mca\_, mcb\_], mcc\_ = [mcg\_, enc-eq(\allowbreak{}ebind(\allowbreak{}mchd\_, 1, [mca\_], mcb\_), ebind(\allowbreak{}mchd\_, 1, [mca\_], mcb\_))] and (\allowbreak{}mchd\_ = 22 or mchd\_ = 24 or mchd\_ = 26) and nbg-ok(\allowbreak{}mcb\_) and mchs\_ = [[mcg\_, enc-eq(\allowbreak{}mca\_, mca\_)]]))
\end{quote}
mentioned by 5 result(s)

\smallskip

\noindent\textbf{\texttt{cr-def-comp(mct\_, mchs\_, mcc\_)}}\quad\textit{predicate}\index{cr-def-comp(mct_, mchs_, mcc_)@\texttt{cr-def-comp(mct\_, mchs\_, mcc\_)}}\\
reads: the core definedness rule comprehension relates \$2 to \$3 \\
defined by \texttt{cr-def-comp} \\
characterized by \texttt{cr-def-comp}  \textit{(definitional)}:
\begin{quote}\ttfamily\footnotesize\raggedright\noindent
forall(\allowbreak{}[mct\_, mchs\_, mcc\_], cr-def-comp(\allowbreak{}mct\_, mchs\_, mcc\_) iff forsome(\allowbreak{}[mcg\_, mcp\_], mcc\_ = [mcg\_, enc-eq(\allowbreak{}enc-comp(\allowbreak{}mcp\_), enc-comp(\allowbreak{}mcp\_))] and nbg-ok(\allowbreak{}mcp\_) and mchs\_ = []))
\end{quote}
mentioned by 5 result(s)

\smallskip

\noindent\textbf{\texttt{cr-def-list(mct\_, mchs\_, mcc\_)}}\quad\textit{predicate}\index{cr-def-list(mct_, mchs_, mcc_)@\texttt{cr-def-list(mct\_, mchs\_, mcc\_)}}\\
reads: the core definedness rule list relates \$2 to \$3 \\
defined by \texttt{cr-def-list} \\
characterized by \texttt{cr-def-list}  \textit{(definitional)}:
\begin{quote}\ttfamily\footnotesize\raggedright\noindent
forall(\allowbreak{}[mct\_, mchs\_, mcc\_], cr-def-list(\allowbreak{}mct\_, mchs\_, mcc\_) iff forsome(\allowbreak{}[mcg\_, mcl\_], mcc\_ = [mcg\_, enc-eq(\allowbreak{}enc-list(\allowbreak{}mcl\_), enc-list(\allowbreak{}mcl\_))] and length(\allowbreak{}mchs\_) = length(\allowbreak{}mcl\_) and forall(\allowbreak{}[mci\_ in nn], 1 <= mci\_ implies mci\_ <= length(\allowbreak{}mcl\_) implies nth(\allowbreak{}mci\_, mchs\_) = [mcg\_, enc-in(\allowbreak{}nth(\allowbreak{}mci\_, mcl\_), enc-setc)])))
\end{quote}
mentioned by 5 result(s)

\smallskip

\noindent\textbf{\texttt{cr-def-strict-arg(mct\_, mchs\_, mcc\_)}}\quad\textit{predicate}\index{cr-def-strict-arg(mct_, mchs_, mcc_)@\texttt{cr-def-strict-arg(mct\_, mchs\_, mcc\_)}}\\
reads: the core definedness rule strict-argument relates \$2 to \$3 \\
defined by \texttt{cr-def-strict-arg} \\
characterized by \texttt{cr-def-strict-arg}  \textit{(definitional)}:
\begin{quote}\ttfamily\footnotesize\raggedright\noindent
forall(\allowbreak{}[mct\_, mchs\_, mcc\_], cr-def-strict-arg(\allowbreak{}mct\_, mchs\_, mcc\_) iff forsome(\allowbreak{}[mcg\_, mchd\_, mcl\_, mci\_], mcc\_ = [mcg\_, enc-eq(\allowbreak{}nth(\allowbreak{}mci\_, mcl\_), nth(\allowbreak{}mci\_, mcl\_))] and strict-head(\allowbreak{}mchd\_) and mci\_ in nn and 1 <= mci\_ and mci\_ <= length(\allowbreak{}mcl\_) and mchs\_ = [[mcg\_, enc-eq(\allowbreak{}eapp(\allowbreak{}mchd\_, mcl\_), eapp(\allowbreak{}mchd\_, mcl\_))]]))
\end{quote}
mentioned by 5 result(s)

\smallskip

\noindent\textbf{\texttt{cr-def-strict-dom(mct\_, mchs\_, mcc\_)}}\quad\textit{predicate}\index{cr-def-strict-dom(mct_, mchs_, mcc_)@\texttt{cr-def-strict-dom(mct\_, mchs\_, mcc\_)}}\\
reads: the core definedness rule strict-domain relates \$2 to \$3 \\
defined by \texttt{cr-def-strict-dom} \\
characterized by \texttt{cr-def-strict-dom}  \textit{(definitional)}:
\begin{quote}\ttfamily\footnotesize\raggedright\noindent
forall(\allowbreak{}[mct\_, mchs\_, mcc\_], cr-def-strict-dom(\allowbreak{}mct\_, mchs\_, mcc\_) iff forsome(\allowbreak{}[mcg\_, mchd\_, mca\_, mcb\_], mcc\_ = [mcg\_, enc-eq(\allowbreak{}mca\_, mca\_)] and (\allowbreak{}mchd\_ = 22 or mchd\_ = 24 or mchd\_ = 26) and mchs\_ = [[mcg\_, enc-eq(\allowbreak{}ebind(\allowbreak{}mchd\_, 1, [mca\_], mcb\_), ebind(\allowbreak{}mchd\_, 1, [mca\_], mcb\_))]]))
\end{quote}
mentioned by 5 result(s)

\smallskip

\noindent\textbf{\texttt{cr-def-strict-in-l(mct\_, mchs\_, mcc\_)}}\quad\textit{predicate}\index{cr-def-strict-in-l(mct_, mchs_, mcc_)@\texttt{cr-def-strict-in-l(mct\_, mchs\_, mcc\_)}}\\
reads: the core definedness rule strict-in-left relates \$2 to \$3 \\
defined by \texttt{cr-def-strict-in-l} \\
characterized by \texttt{cr-def-strict-in-l}  \textit{(definitional)}:
\begin{quote}\ttfamily\footnotesize\raggedright\noindent
forall(\allowbreak{}[mct\_, mchs\_, mcc\_], cr-def-strict-in-l(\allowbreak{}mct\_, mchs\_, mcc\_) iff forsome(\allowbreak{}[mcg\_, mcw\_, mca\_], mcc\_ = [mcg\_, enc-eq(\allowbreak{}mcw\_, mcw\_)] and mchs\_ = [[mcg\_, enc-in(\allowbreak{}mcw\_, mca\_)]]))
\end{quote}
mentioned by 5 result(s)

\smallskip

\noindent\textbf{\texttt{cr-def-strict-in-r(mct\_, mchs\_, mcc\_)}}\quad\textit{predicate}\index{cr-def-strict-in-r(mct_, mchs_, mcc_)@\texttt{cr-def-strict-in-r(mct\_, mchs\_, mcc\_)}}\\
reads: the core definedness rule strict-in-right relates \$2 to \$3 \\
defined by \texttt{cr-def-strict-in-r} \\
characterized by \texttt{cr-def-strict-in-r}  \textit{(definitional)}:
\begin{quote}\ttfamily\footnotesize\raggedright\noindent
forall(\allowbreak{}[mct\_, mchs\_, mcc\_], cr-def-strict-in-r(\allowbreak{}mct\_, mchs\_, mcc\_) iff forsome(\allowbreak{}[mcg\_, mcw\_, mca\_], mcc\_ = [mcg\_, enc-eq(\allowbreak{}mca\_, mca\_)] and mchs\_ = [[mcg\_, enc-in(\allowbreak{}mcw\_, mca\_)]]))
\end{quote}
mentioned by 5 result(s)

\smallskip

\noindent\textbf{\texttt{cr-def-strict-pred(mct\_, mchs\_, mcc\_)}}\quad\textit{predicate}\index{cr-def-strict-pred(mct_, mchs_, mcc_)@\texttt{cr-def-strict-pred(mct\_, mchs\_, mcc\_)}}\\
reads: the core definedness rule strict-predicate relates \$2 to \$3 \\
defined by \texttt{cr-def-strict-pred} \\
characterized by \texttt{cr-def-strict-pred}  \textit{(definitional)}:
\begin{quote}\ttfamily\footnotesize\raggedright\noindent
forall(\allowbreak{}[mct\_, mchs\_, mcc\_], cr-def-strict-pred(\allowbreak{}mct\_, mchs\_, mcc\_) iff forsome(\allowbreak{}[mcg\_, mcl\_, mck\_, mci\_], mcc\_ = [mcg\_, enc-eq(\allowbreak{}nth(\allowbreak{}mci\_, mcl\_), nth(\allowbreak{}mci\_, mcl\_))] and nth(\allowbreak{}1, mcl\_) = econst(\allowbreak{}mck\_) and pred-code(\allowbreak{}mck\_) and mci\_ in nn and 2 <= mci\_ and mci\_ <= length(\allowbreak{}mcl\_) and mchs\_ = [[mcg\_, eapp(\allowbreak{}10, mcl\_)]]))
\end{quote}
mentioned by 7 result(s)

\smallskip

\noindent\textbf{\texttt{cr-def-total(mct\_, mchs\_, mcc\_)}}\quad\textit{predicate}\index{cr-def-total(mct_, mchs_, mcc_)@\texttt{cr-def-total(mct\_, mchs\_, mcc\_)}}\\
reads: the core definedness rule total-constructor relates \$2 to \$3 \\
defined by \texttt{cr-def-total} \\
characterized by \texttt{cr-def-total}  \textit{(definitional)}:
\begin{quote}\ttfamily\footnotesize\raggedright\noindent
forall(\allowbreak{}[mct\_, mchs\_, mcc\_], cr-def-total(\allowbreak{}mct\_, mchs\_, mcc\_) iff forsome(\allowbreak{}[mcg\_, mchd\_, mcl\_], mcc\_ = [mcg\_, enc-eq(\allowbreak{}eapp(\allowbreak{}mchd\_, mcl\_), eapp(\allowbreak{}mchd\_, mcl\_))] and total-head(\allowbreak{}mchd\_) and length(\allowbreak{}mchs\_) = length(\allowbreak{}mcl\_) and forall(\allowbreak{}[mci\_ in nn], 1 <= mci\_ implies mci\_ <= length(\allowbreak{}mcl\_) implies nth(\allowbreak{}mci\_, mchs\_) = [mcg\_, enc-eq(\allowbreak{}nth(\allowbreak{}mci\_, mcl\_), nth(\allowbreak{}mci\_, mcl\_))])))
\end{quote}
mentioned by 5 result(s)

\smallskip

\noindent\textbf{\texttt{cr-eq-subst(mct\_, mchs\_, mcc\_)}}\quad\textit{predicate}\index{cr-eq-subst(mct_, mchs_, mcc_)@\texttt{cr-eq-subst(mct\_, mchs\_, mcc\_)}}\\
reads: the core eq-subst relates \$2 to \$3 \\
defined by \texttt{cr-eq-subst} \\
characterized by \texttt{cr-eq-subst}  \textit{(definitional)}:
\begin{quote}\ttfamily\footnotesize\raggedright\noindent
forall(\allowbreak{}[mct\_, mchs\_, mcc\_], cr-eq-subst(\allowbreak{}mct\_, mchs\_, mcc\_) iff forsome(\allowbreak{}[mcg\_, mcp\_, mcz\_, mcs\_, mctm\_], mcc\_ = [mcg\_, expr-subst(\allowbreak{}mcp\_, mcz\_, mcs\_)] and mcz\_ in nn and mcp\_ in expr-set and (\allowbreak{}enc-eq(\allowbreak{}mcs\_, mctm\_) in mcg\_ or enc-qeq(\allowbreak{}mcs\_, mctm\_) in mcg\_) and mchs\_ = [[mcg\_, expr-subst(\allowbreak{}mcp\_, mcz\_, mctm\_)]]))
\end{quote}
mentioned by 5 result(s)

\smallskip

\noindent\textbf{\texttt{cr-eq-subst-sym(mct\_, mchs\_, mcc\_)}}\quad\textit{predicate}\index{cr-eq-subst-sym(mct_, mchs_, mcc_)@\texttt{cr-eq-subst-sym(mct\_, mchs\_, mcc\_)}}\\
reads: the core eq-subst (other orientation) relates \$2 to \$3 \\
defined by \texttt{cr-eq-subst-sym} \\
characterized by \texttt{cr-eq-subst-sym}  \textit{(definitional)}:
\begin{quote}\ttfamily\footnotesize\raggedright\noindent
forall(\allowbreak{}[mct\_, mchs\_, mcc\_], cr-eq-subst-sym(\allowbreak{}mct\_, mchs\_, mcc\_) iff forsome(\allowbreak{}[mcg\_, mcp\_, mcz\_, mcs\_, mctm\_], mcc\_ = [mcg\_, expr-subst(\allowbreak{}mcp\_, mcz\_, mcs\_)] and mcz\_ in nn and mcp\_ in expr-set and (\allowbreak{}enc-eq(\allowbreak{}mctm\_, mcs\_) in mcg\_ or enc-qeq(\allowbreak{}mctm\_, mcs\_) in mcg\_) and mchs\_ = [[mcg\_, expr-subst(\allowbreak{}mcp\_, mcz\_, mctm\_)]]))
\end{quote}
mentioned by 5 result(s)

\smallskip

\noindent\textbf{\texttt{cr-forall-elim(mct\_, mchs\_, mcc\_)}}\quad\textit{predicate}\index{cr-forall-elim(mct_, mchs_, mcc_)@\texttt{cr-forall-elim(mct\_, mchs\_, mcc\_)}}\\
reads: the core forall-elim relates \$2 to \$3 \\
defined by \texttt{cr-forall-elim} \\
characterized by \texttt{cr-forall-elim}  \textit{(definitional)}:
\begin{quote}\ttfamily\footnotesize\raggedright\noindent
forall(\allowbreak{}[mct\_, mchs\_, mcc\_], cr-forall-elim(\allowbreak{}mct\_, mchs\_, mcc\_) iff forsome(\allowbreak{}[mcg\_, mcf\_, mcb\_, mctm\_], mcc\_ = [mcg\_, mcf\_] and enc-all(\allowbreak{}mcb\_) in mcg\_ and mctm\_ in clt and mchs\_ = [[ctx-add(\allowbreak{}mcg\_, expr-open(\allowbreak{}mcb\_, 0, mctm\_)), mcf\_], [mcg\_, enc-eq(\allowbreak{}mctm\_, mctm\_)]]))
\end{quote}
mentioned by 5 result(s)

\smallskip

\noindent\textbf{\texttt{cr-forall-intro(mct\_, mchs\_, mcc\_)}}\quad\textit{predicate}\index{cr-forall-intro(mct_, mchs_, mcc_)@\texttt{cr-forall-intro(mct\_, mchs\_, mcc\_)}}\\
reads: the core forall-intro relates \$2 to \$3 \\
defined by \texttt{cr-forall-intro} \\
characterized by \texttt{cr-forall-intro}  \textit{(definitional)}:
\begin{quote}\ttfamily\footnotesize\raggedright\noindent
forall(\allowbreak{}[mct\_, mchs\_, mcc\_], cr-forall-intro(\allowbreak{}mct\_, mchs\_, mcc\_) iff forsome(\allowbreak{}[mcg\_, mcb\_, mcy\_], mcc\_ = [mcg\_, enc-all(\allowbreak{}mcb\_)] and mcy\_ in nn and not(\allowbreak{}mcy\_ in union(\allowbreak{}ctx-fv(\allowbreak{}mcg\_), expr-fv(\allowbreak{}mcb\_))) and mchs\_ = [[mcg\_, expr-open(\allowbreak{}mcb\_, 0, efvar(\allowbreak{}mcy\_))]]))
\end{quote}
mentioned by 5 result(s)

\smallskip

\noindent\textbf{\texttt{cr-forsome-elim(mct\_, mchs\_, mcc\_)}}\quad\textit{predicate}\index{cr-forsome-elim(mct_, mchs_, mcc_)@\texttt{cr-forsome-elim(mct\_, mchs\_, mcc\_)}}\\
reads: the core forsome-elim relates \$2 to \$3 \\
defined by \texttt{cr-forsome-elim} \\
characterized by \texttt{cr-forsome-elim}  \textit{(definitional)}:
\begin{quote}\ttfamily\footnotesize\raggedright\noindent
forall(\allowbreak{}[mct\_, mchs\_, mcc\_], cr-forsome-elim(\allowbreak{}mct\_, mchs\_, mcc\_) iff forsome(\allowbreak{}[mcg\_, mcf\_, mcb\_, mcy\_], mcc\_ = [mcg\_, mcf\_] and enc-ex(\allowbreak{}mcb\_) in mcg\_ and mcy\_ in nn and not(\allowbreak{}mcy\_ in union(\allowbreak{}ctx-fv(\allowbreak{}mcg\_), expr-fv(\allowbreak{}mcf\_))) and mchs\_ = [[ctx-add(\allowbreak{}mcg\_, expr-open(\allowbreak{}mcb\_, 0, efvar(\allowbreak{}mcy\_))), mcf\_]]))
\end{quote}
mentioned by 5 result(s)

\smallskip

\noindent\textbf{\texttt{cr-forsome-intro(mct\_, mchs\_, mcc\_)}}\quad\textit{predicate}\index{cr-forsome-intro(mct_, mchs_, mcc_)@\texttt{cr-forsome-intro(mct\_, mchs\_, mcc\_)}}\\
reads: the core forsome-intro relates \$2 to \$3 \\
defined by \texttt{cr-forsome-intro} \\
characterized by \texttt{cr-forsome-intro}  \textit{(definitional)}:
\begin{quote}\ttfamily\footnotesize\raggedright\noindent
forall(\allowbreak{}[mct\_, mchs\_, mcc\_], cr-forsome-intro(\allowbreak{}mct\_, mchs\_, mcc\_) iff forsome(\allowbreak{}[mcg\_, mcb\_, mctm\_], mcc\_ = [mcg\_, enc-ex(\allowbreak{}mcb\_)] and mctm\_ in clt and mchs\_ = [[mcg\_, expr-open(\allowbreak{}mcb\_, 0, mctm\_)], [mcg\_, enc-eq(\allowbreak{}mctm\_, mctm\_)]]))
\end{quote}
mentioned by 5 result(s)

\smallskip

\noindent\textbf{\texttt{cr-functoid-beta-eq(mct\_, mchs\_, mcc\_)}}\quad\textit{predicate}\index{cr-functoid-beta-eq(mct_, mchs_, mcc_)@\texttt{cr-functoid-beta-eq(mct\_, mchs\_, mcc\_)}}\\
reads: the core functoid-beta-eq relates \$2 to \$3 \\
defined by \texttt{cr-functoid-beta-eq} \\
characterized by \texttt{cr-functoid-beta-eq}  \textit{(definitional)}:
\begin{quote}\ttfamily\footnotesize\raggedright\noindent
forall(\allowbreak{}[mct\_, mchs\_, mcc\_], cr-functoid-beta-eq(\allowbreak{}mct\_, mchs\_, mcc\_) iff forsome(\allowbreak{}[mcg\_, mchd\_, mcl\_, mcb\_], mcc\_ = [mcg\_, enc-qeq(\allowbreak{}eapp(\allowbreak{}mchd\_, mcl\_), openl(\allowbreak{}mcb\_, mcl\_, length(\allowbreak{}mcl\_)))] and fun-def(\allowbreak{}mchd\_, length(\allowbreak{}mcl\_), mcb\_) and mchs\_ = []))
\end{quote}
mentioned by 6 result(s)

\smallskip

\noindent\textbf{\texttt{cr-if-false(mct\_, mchs\_, mcc\_)}}\quad\textit{predicate}\index{cr-if-false(mct_, mchs_, mcc_)@\texttt{cr-if-false(mct\_, mchs\_, mcc\_)}}\\
reads: the core if-false relates \$2 to \$3 \\
defined by \texttt{cr-if-false} \\
characterized by \texttt{cr-if-false}  \textit{(definitional)}:
\begin{quote}\ttfamily\footnotesize\raggedright\noindent
forall(\allowbreak{}[mct\_, mchs\_, mcc\_], cr-if-false(\allowbreak{}mct\_, mchs\_, mcc\_) iff forsome(\allowbreak{}[mcg\_, mcf\_, mcp\_, mca\_, mcb\_], mcc\_ = [mcg\_, mcf\_] and mchs\_ = [[mcg\_, enc-not(\allowbreak{}mcp\_)], [ctx-add(\allowbreak{}mcg\_, enc-qeq(\allowbreak{}enc-if(\allowbreak{}mcp\_, mca\_, mcb\_), mcb\_)), mcf\_]] and mca\_ in clt and mcb\_ in clt))
\end{quote}
mentioned by 5 result(s)

\smallskip

\noindent\textbf{\texttt{cr-if-true(mct\_, mchs\_, mcc\_)}}\quad\textit{predicate}\index{cr-if-true(mct_, mchs_, mcc_)@\texttt{cr-if-true(mct\_, mchs\_, mcc\_)}}\\
reads: the core if-true relates \$2 to \$3 \\
defined by \texttt{cr-if-true} \\
characterized by \texttt{cr-if-true}  \textit{(definitional)}:
\begin{quote}\ttfamily\footnotesize\raggedright\noindent
forall(\allowbreak{}[mct\_, mchs\_, mcc\_], cr-if-true(\allowbreak{}mct\_, mchs\_, mcc\_) iff forsome(\allowbreak{}[mcg\_, mcf\_, mcp\_, mca\_, mcb\_], mcc\_ = [mcg\_, mcf\_] and mchs\_ = [[mcg\_, mcp\_], [ctx-add(\allowbreak{}mcg\_, enc-qeq(\allowbreak{}enc-if(\allowbreak{}mcp\_, mca\_, mcb\_), mca\_)), mcf\_]] and mca\_ in clt and mcb\_ in clt))
\end{quote}
mentioned by 5 result(s)

\smallskip

\noindent\textbf{\texttt{cr-iff-elim(mct\_, mchs\_, mcc\_)}}\quad\textit{predicate}\index{cr-iff-elim(mct_, mchs_, mcc_)@\texttt{cr-iff-elim(mct\_, mchs\_, mcc\_)}}\\
reads: the core iff-elim relates \$2 to \$3 \\
defined by \texttt{cr-iff-elim} \\
characterized by \texttt{cr-iff-elim}  \textit{(definitional)}:
\begin{quote}\ttfamily\footnotesize\raggedright\noindent
forall(\allowbreak{}[mct\_, mchs\_, mcc\_], cr-iff-elim(\allowbreak{}mct\_, mchs\_, mcc\_) iff forsome(\allowbreak{}[mcg\_, mcf\_, mca\_, mcb\_], mcc\_ = [mcg\_, mcf\_] and enc-iff(\allowbreak{}mca\_, mcb\_) in mcg\_ and mchs\_ = [[ctx-add(\allowbreak{}ctx-add(\allowbreak{}mcg\_, enc-imp(\allowbreak{}mca\_, mcb\_)), enc-imp(\allowbreak{}mcb\_, mca\_)), mcf\_]]))
\end{quote}
mentioned by 5 result(s)

\smallskip

\noindent\textbf{\texttt{cr-iff-intro(mct\_, mchs\_, mcc\_)}}\quad\textit{predicate}\index{cr-iff-intro(mct_, mchs_, mcc_)@\texttt{cr-iff-intro(mct\_, mchs\_, mcc\_)}}\\
reads: the core iff-intro relates \$2 to \$3 \\
defined by \texttt{cr-iff-intro} \\
characterized by \texttt{cr-iff-intro}  \textit{(definitional)}:
\begin{quote}\ttfamily\footnotesize\raggedright\noindent
forall(\allowbreak{}[mct\_, mchs\_, mcc\_], cr-iff-intro(\allowbreak{}mct\_, mchs\_, mcc\_) iff forsome(\allowbreak{}[mcg\_, mca\_, mcb\_], mcc\_ = [mcg\_, enc-iff(\allowbreak{}mca\_, mcb\_)] and mchs\_ = [[ctx-add(\allowbreak{}mcg\_, mca\_), mcb\_], [ctx-add(\allowbreak{}mcg\_, mcb\_), mca\_]]))
\end{quote}
mentioned by 5 result(s)

\smallskip

\noindent\textbf{\texttt{cr-implies-intro(mct\_, mchs\_, mcc\_)}}\quad\textit{predicate}\index{cr-implies-intro(mct_, mchs_, mcc_)@\texttt{cr-implies-intro(mct\_, mchs\_, mcc\_)}}\\
reads: the core implies-intro relates \$2 to \$3 \\
defined by \texttt{cr-implies-intro} \\
characterized by \texttt{cr-implies-intro}  \textit{(definitional)}:
\begin{quote}\ttfamily\footnotesize\raggedright\noindent
forall(\allowbreak{}[mct\_, mchs\_, mcc\_], cr-implies-intro(\allowbreak{}mct\_, mchs\_, mcc\_) iff forsome(\allowbreak{}[mcg\_, mca\_, mcb\_], mcc\_ = [mcg\_, enc-imp(\allowbreak{}mca\_, mcb\_)] and mchs\_ = [[ctx-add(\allowbreak{}mcg\_, mca\_), mcb\_]]))
\end{quote}
mentioned by 5 result(s)

\smallskip

\noindent\textbf{\texttt{cr-intersection-mem(mct\_, mchs\_, mcc\_)}}\quad\textit{predicate}\index{cr-intersection-mem(mct_, mchs_, mcc_)@\texttt{cr-intersection-mem(mct\_, mchs\_, mcc\_)}}\\
reads: the core intersection-mem relates \$2 to \$3 \\
defined by \texttt{cr-intersection-mem} \\
characterized by \texttt{cr-intersection-mem}  \textit{(definitional)}:
\begin{quote}\ttfamily\footnotesize\raggedright\noindent
forall(\allowbreak{}[mct\_, mchs\_, mcc\_], cr-intersection-mem(\allowbreak{}mct\_, mchs\_, mcc\_) iff forsome(\allowbreak{}[mcg\_, mcx\_, mcl\_], mcc\_ = [mcg\_, enc-iff(\allowbreak{}enc-in(\allowbreak{}mcx\_, enc-inter(\allowbreak{}mcl\_)), enc-andl(\allowbreak{}mcx\_, mcl\_, length(\allowbreak{}mcl\_)))] and 1 <= length(\allowbreak{}mcl\_) and mchs\_ = []))
\end{quote}
mentioned by 5 result(s)

\smallskip

\noindent\textbf{\texttt{cr-iota-cr(mct\_, mchs\_, mcc\_)}}\quad\textit{predicate}\index{cr-iota-cr(mct_, mchs_, mcc_)@\texttt{cr-iota-cr(mct\_, mchs\_, mcc\_)}}\\
reads: the core congruence rule for descriptions relates \$2 to \$3 \\
defined by \texttt{cr-iota-cr} \\
characterized by \texttt{cr-iota-cr}  \textit{(definitional)}:
\begin{quote}\ttfamily\footnotesize\raggedright\noindent
forall(\allowbreak{}[mct\_, mchs\_, mcc\_], cr-iota-cr(\allowbreak{}mct\_, mchs\_, mcc\_) iff forsome(\allowbreak{}[mcg\_, mcp\_, mcq\_, mcy\_], mcc\_ = [mcg\_, enc-qeq(\allowbreak{}enc-iota(\allowbreak{}mcp\_), enc-iota(\allowbreak{}mcq\_))] and mcy\_ in nn and not(\allowbreak{}mcy\_ in union(\allowbreak{}ctx-fv(\allowbreak{}mcg\_), expr-fv(\allowbreak{}enc-qeq(\allowbreak{}enc-iota(\allowbreak{}mcp\_), enc-iota(\allowbreak{}mcq\_))))) and mchs\_ = [[mcg\_, enc-iff(\allowbreak{}expr-open(\allowbreak{}mcp\_, 0, efvar(\allowbreak{}mcy\_)), expr-open(\allowbreak{}mcq\_, 0, efvar(\allowbreak{}mcy\_)))]]))
\end{quote}
mentioned by 6 result(s)

\smallskip

\noindent\textbf{\texttt{cr-iota-def(mct\_, mchs\_, mcc\_)}}\quad\textit{predicate}\index{cr-iota-def(mct_, mchs_, mcc_)@\texttt{cr-iota-def(mct\_, mchs\_, mcc\_)}}\\
reads: the core iota-def relates \$2 to \$3 \\
defined by \texttt{cr-iota-def} \\
characterized by \texttt{cr-iota-def}  \textit{(definitional)}:
\begin{quote}\ttfamily\footnotesize\raggedright\noindent
forall(\allowbreak{}[mct\_, mchs\_, mcc\_], cr-iota-def(\allowbreak{}mct\_, mchs\_, mcc\_) iff forsome(\allowbreak{}[mcg\_, mcf\_, mcp\_], mcc\_ = [mcg\_, mcf\_] and mchs\_ = [[mcg\_, enc-ex(\allowbreak{}enc-and(\allowbreak{}mcp\_, enc-all(\allowbreak{}enc-imp(\allowbreak{}mcp\_, enc-eq(\allowbreak{}ebvar(\allowbreak{}1), ebvar(\allowbreak{}0))))))], [ctx-add(\allowbreak{}mcg\_, expr-open(\allowbreak{}mcp\_, 0, enc-iota(\allowbreak{}mcp\_))), mcf\_]]))
\end{quote}
mentioned by 5 result(s)

\smallskip

\noindent\textbf{\texttt{cr-iota-in-elim(mct\_, mchs\_, mcc\_)}}\quad\textit{predicate}\index{cr-iota-in-elim(mct_, mchs_, mcc_)@\texttt{cr-iota-in-elim(mct\_, mchs\_, mcc\_)}}\\
reads: the core iota-in-elim relates \$2 to \$3 \\
defined by \texttt{cr-iota-in-elim} \\
characterized by \texttt{cr-iota-in-elim}  \textit{(definitional)}:
\begin{quote}\ttfamily\footnotesize\raggedright\noindent
forall(\allowbreak{}[mct\_, mchs\_, mcc\_], cr-iota-in-elim(\allowbreak{}mct\_, mchs\_, mcc\_) iff forsome(\allowbreak{}[mcg\_, mcf\_, mcp\_], mcc\_ = [mcg\_, mcf\_] and mchs\_ = [[mcg\_, enc-eq(\allowbreak{}enc-iota(\allowbreak{}mcp\_), enc-iota(\allowbreak{}mcp\_))], [ctx-add(\allowbreak{}mcg\_, expr-open(\allowbreak{}mcp\_, 0, enc-iota(\allowbreak{}mcp\_))), mcf\_]]))
\end{quote}
mentioned by 5 result(s)

\smallskip

\noindent\textbf{\texttt{cr-lambda-beta-eq(mct\_, mchs\_, mcc\_)}}\quad\textit{predicate}\index{cr-lambda-beta-eq(mct_, mchs_, mcc_)@\texttt{cr-lambda-beta-eq(mct\_, mchs\_, mcc\_)}}\\
reads: the core lambda-beta-eq relates \$2 to \$3 \\
defined by \texttt{cr-lambda-beta-eq} \\
characterized by \texttt{cr-lambda-beta-eq}  \textit{(definitional)}:
\begin{quote}\ttfamily\footnotesize\raggedright\noindent
forall(\allowbreak{}[mct\_, mchs\_, mcc\_], cr-lambda-beta-eq(\allowbreak{}mct\_, mchs\_, mcc\_) iff forsome(\allowbreak{}[mcg\_, mca\_, mcb\_, mctm\_], mcc\_ = [mcg\_, enc-qeq(\allowbreak{}enc-apply(\allowbreak{}enc-lam(\allowbreak{}mca\_, mcb\_), mctm\_), expr-open(\allowbreak{}mcb\_, 0, mctm\_))] and nbg-ok(\allowbreak{}mcb\_) and mchs\_ = [[mcg\_, enc-in(\allowbreak{}mctm\_, mca\_)]]))
\end{quote}
mentioned by 5 result(s)

\smallskip

\noindent\textbf{\texttt{cr-lambda-type(mct\_, mchs\_, mcc\_)}}\quad\textit{predicate}\index{cr-lambda-type(mct_, mchs_, mcc_)@\texttt{cr-lambda-type(mct\_, mchs\_, mcc\_)}}\\
reads: the core lambda-type relates \$2 to \$3 \\
defined by \texttt{cr-lambda-type} \\
characterized by \texttt{cr-lambda-type}  \textit{(definitional)}:
\begin{quote}\ttfamily\footnotesize\raggedright\noindent
forall(\allowbreak{}[mct\_, mchs\_, mcc\_], cr-lambda-type(\allowbreak{}mct\_, mchs\_, mcc\_) iff forsome(\allowbreak{}[mcg\_, mca\_, mcb\_, mcd\_], mcc\_ = [mcg\_, enc-in(\allowbreak{}enc-lam(\allowbreak{}mca\_, mcb\_), enc-fun(\allowbreak{}mca\_, mcd\_))] and nbg-ok(\allowbreak{}mcb\_) and mchs\_ = [[mcg\_, enc-all(\allowbreak{}enc-imp(\allowbreak{}enc-in(\allowbreak{}ebvar(\allowbreak{}0), mca\_), enc-in(\allowbreak{}mcb\_, mcd\_)))], [mcg\_, enc-in(\allowbreak{}mca\_, enc-setc)]]))
\end{quote}
mentioned by 5 result(s)

\smallskip

\noindent\textbf{\texttt{cr-length-eq(mct\_, mchs\_, mcc\_)}}\quad\textit{predicate}\index{cr-length-eq(mct_, mchs_, mcc_)@\texttt{cr-length-eq(mct\_, mchs\_, mcc\_)}}\\
reads: the core length-eq relates \$2 to \$3 \\
defined by \texttt{cr-length-eq} \\
characterized by \texttt{cr-length-eq}  \textit{(definitional)}:
\begin{quote}\ttfamily\footnotesize\raggedright\noindent
forall(\allowbreak{}[mct\_, mchs\_, mcc\_], cr-length-eq(\allowbreak{}mct\_, mchs\_, mcc\_) iff forsome(\allowbreak{}[mcg\_, mcl\_], mcc\_ = [mcg\_, enc-qeq(\allowbreak{}enc-length(\allowbreak{}enc-list(\allowbreak{}mcl\_)), enc-num(\allowbreak{}length(\allowbreak{}mcl\_)))] and mchs\_ = [[mcg\_, enc-eq(\allowbreak{}enc-list(\allowbreak{}mcl\_), enc-list(\allowbreak{}mcl\_))]]))
\end{quote}
mentioned by 5 result(s)

\smallskip

\noindent\textbf{\texttt{cr-not-elim(mct\_, mchs\_, mcc\_)}}\quad\textit{predicate}\index{cr-not-elim(mct_, mchs_, mcc_)@\texttt{cr-not-elim(mct\_, mchs\_, mcc\_)}}\\
reads: the core not-elim relates \$2 to \$3 \\
defined by \texttt{cr-not-elim} \\
characterized by \texttt{cr-not-elim}  \textit{(definitional)}:
\begin{quote}\ttfamily\footnotesize\raggedright\noindent
forall(\allowbreak{}[mct\_, mchs\_, mcc\_], cr-not-elim(\allowbreak{}mct\_, mchs\_, mcc\_) iff forsome(\allowbreak{}[mcg\_, mcf\_, mca\_], mcc\_ = [mcg\_, mcf\_] and enc-not(\allowbreak{}mca\_) in mcg\_ and mca\_ in mcg\_ and mchs\_ = []))
\end{quote}
mentioned by 5 result(s)

\smallskip

\noindent\textbf{\texttt{cr-not-intro(mct\_, mchs\_, mcc\_)}}\quad\textit{predicate}\index{cr-not-intro(mct_, mchs_, mcc_)@\texttt{cr-not-intro(mct\_, mchs\_, mcc\_)}}\\
reads: the core not-intro relates \$2 to \$3 \\
defined by \texttt{cr-not-intro} \\
characterized by \texttt{cr-not-intro}  \textit{(definitional)}:
\begin{quote}\ttfamily\footnotesize\raggedright\noindent
forall(\allowbreak{}[mct\_, mchs\_, mcc\_], cr-not-intro(\allowbreak{}mct\_, mchs\_, mcc\_) iff forsome(\allowbreak{}[mcg\_, mca\_], mcc\_ = [mcg\_, enc-not(\allowbreak{}mca\_)] and mchs\_ = [[ctx-add(\allowbreak{}mcg\_, mca\_), enc-falsity]]))
\end{quote}
mentioned by 5 result(s)

\smallskip

\noindent\textbf{\texttt{cr-nth-eq(mct\_, mchs\_, mcc\_)}}\quad\textit{predicate}\index{cr-nth-eq(mct_, mchs_, mcc_)@\texttt{cr-nth-eq(mct\_, mchs\_, mcc\_)}}\\
reads: the core nth-eq relates \$2 to \$3 \\
defined by \texttt{cr-nth-eq} \\
characterized by \texttt{cr-nth-eq}  \textit{(definitional)}:
\begin{quote}\ttfamily\footnotesize\raggedright\noindent
forall(\allowbreak{}[mct\_, mchs\_, mcc\_], cr-nth-eq(\allowbreak{}mct\_, mchs\_, mcc\_) iff forsome(\allowbreak{}[mcg\_, mcl\_, mck\_], mcc\_ = [mcg\_, enc-qeq(\allowbreak{}enc-nth(\allowbreak{}enc-num(\allowbreak{}mck\_), enc-list(\allowbreak{}mcl\_)), nth(\allowbreak{}mck\_, mcl\_))] and mck\_ in nn and 1 <= mck\_ and mck\_ <= length(\allowbreak{}mcl\_) and mchs\_ = [[mcg\_, enc-eq(\allowbreak{}enc-list(\allowbreak{}mcl\_), enc-list(\allowbreak{}mcl\_))]]))
\end{quote}
mentioned by 5 result(s)

\smallskip

\noindent\textbf{\texttt{cr-or-elim(mct\_, mchs\_, mcc\_)}}\quad\textit{predicate}\index{cr-or-elim(mct_, mchs_, mcc_)@\texttt{cr-or-elim(mct\_, mchs\_, mcc\_)}}\\
reads: the core or-elim relates \$2 to \$3 \\
defined by \texttt{cr-or-elim} \\
characterized by \texttt{cr-or-elim}  \textit{(definitional)}:
\begin{quote}\ttfamily\footnotesize\raggedright\noindent
forall(\allowbreak{}[mct\_, mchs\_, mcc\_], cr-or-elim(\allowbreak{}mct\_, mchs\_, mcc\_) iff forsome(\allowbreak{}[mcg\_, mcf\_, mca\_, mcb\_], mcc\_ = [mcg\_, mcf\_] and enc-or(\allowbreak{}mca\_, mcb\_) in mcg\_ and mchs\_ = [[ctx-add(\allowbreak{}mcg\_, mca\_), mcf\_], [ctx-add(\allowbreak{}mcg\_, mcb\_), mcf\_]]))
\end{quote}
mentioned by 5 result(s)

\smallskip

\noindent\textbf{\texttt{cr-or-intro-left(mct\_, mchs\_, mcc\_)}}\quad\textit{predicate}\index{cr-or-intro-left(mct_, mchs_, mcc_)@\texttt{cr-or-intro-left(mct\_, mchs\_, mcc\_)}}\\
reads: the core or-intro-left relates \$2 to \$3 \\
defined by \texttt{cr-or-intro-left} \\
characterized by \texttt{cr-or-intro-left}  \textit{(definitional)}:
\begin{quote}\ttfamily\footnotesize\raggedright\noindent
forall(\allowbreak{}[mct\_, mchs\_, mcc\_], cr-or-intro-left(\allowbreak{}mct\_, mchs\_, mcc\_) iff forsome(\allowbreak{}[mcg\_, mca\_, mcb\_], mcc\_ = [mcg\_, enc-or(\allowbreak{}mca\_, mcb\_)] and mchs\_ = [[mcg\_, mca\_]]))
\end{quote}
mentioned by 5 result(s)

\smallskip

\noindent\textbf{\texttt{cr-or-intro-right(mct\_, mchs\_, mcc\_)}}\quad\textit{predicate}\index{cr-or-intro-right(mct_, mchs_, mcc_)@\texttt{cr-or-intro-right(mct\_, mchs\_, mcc\_)}}\\
reads: the core or-intro-right relates \$2 to \$3 \\
defined by \texttt{cr-or-intro-right} \\
characterized by \texttt{cr-or-intro-right}  \textit{(definitional)}:
\begin{quote}\ttfamily\footnotesize\raggedright\noindent
forall(\allowbreak{}[mct\_, mchs\_, mcc\_], cr-or-intro-right(\allowbreak{}mct\_, mchs\_, mcc\_) iff forsome(\allowbreak{}[mcg\_, mca\_, mcb\_], mcc\_ = [mcg\_, enc-or(\allowbreak{}mca\_, mcb\_)] and mchs\_ = [[mcg\_, mcb\_]]))
\end{quote}
mentioned by 5 result(s)

\smallskip

\noindent\textbf{\texttt{cr-quasi-reflexivity(mct\_, mchs\_, mcc\_)}}\quad\textit{predicate}\index{cr-quasi-reflexivity(mct_, mchs_, mcc_)@\texttt{cr-quasi-reflexivity(mct\_, mchs\_, mcc\_)}}\\
reads: the core quasi-reflexivity relates \$2 to \$3 \\
defined by \texttt{cr-quasi-reflexivity} \\
characterized by \texttt{cr-quasi-reflexivity}  \textit{(definitional)}:
\begin{quote}\ttfamily\footnotesize\raggedright\noindent
forall(\allowbreak{}[mct\_, mchs\_, mcc\_], cr-quasi-reflexivity(\allowbreak{}mct\_, mchs\_, mcc\_) iff forsome(\allowbreak{}[mcg\_, mctm\_], mcc\_ = [mcg\_, enc-qeq(\allowbreak{}mctm\_, mctm\_)] and mchs\_ = []))
\end{quote}
mentioned by 5 result(s)

\smallskip

\noindent\textbf{\texttt{cr-reflexivity(mct\_, mchs\_, mcc\_)}}\quad\textit{predicate}\index{cr-reflexivity(mct_, mchs_, mcc_)@\texttt{cr-reflexivity(mct\_, mchs\_, mcc\_)}}\\
reads: the core reflexivity relates \$2 to \$3 \\
defined by \texttt{cr-reflexivity} \\
characterized by \texttt{cr-reflexivity}  \textit{(definitional)}:
\begin{quote}\ttfamily\footnotesize\raggedright\noindent
forall(\allowbreak{}[mct\_, mchs\_, mcc\_], cr-reflexivity(\allowbreak{}mct\_, mchs\_, mcc\_) iff forsome(\allowbreak{}[mcg\_, mctm\_, mcn\_], mcc\_ = [mcg\_, enc-eq(\allowbreak{}mctm\_, mctm\_)] and mcn\_ in nn and (\allowbreak{}mctm\_ = efvar(\allowbreak{}mcn\_) or mctm\_ = econst(\allowbreak{}mcn\_)) and mchs\_ = []))
\end{quote}
mentioned by 5 result(s)

\smallskip

\noindent\textbf{\texttt{cr-sep-mem(mct\_, mchs\_, mcc\_)}}\quad\textit{predicate}\index{cr-sep-mem(mct_, mchs_, mcc_)@\texttt{cr-sep-mem(mct\_, mchs\_, mcc\_)}}\\
reads: the core sep-mem relates \$2 to \$3 \\
defined by \texttt{cr-sep-mem} \\
characterized by \texttt{cr-sep-mem}  \textit{(definitional)}:
\begin{quote}\ttfamily\footnotesize\raggedright\noindent
forall(\allowbreak{}[mct\_, mchs\_, mcc\_], cr-sep-mem(\allowbreak{}mct\_, mchs\_, mcc\_) iff forsome(\allowbreak{}[mcg\_, mca\_, mcp\_, mctm\_], mcc\_ = [mcg\_, enc-iff(\allowbreak{}enc-in(\allowbreak{}mctm\_, enc-sep(\allowbreak{}mca\_, mcp\_)), enc-and(\allowbreak{}enc-in(\allowbreak{}mctm\_, mca\_), expr-open(\allowbreak{}mcp\_, 0, mctm\_)))] and nbg-ok(\allowbreak{}mcp\_) and mchs\_ = []))
\end{quote}
mentioned by 5 result(s)

\smallskip

\noindent\textbf{\texttt{cr-sep-sethood(mct\_, mchs\_, mcc\_)}}\quad\textit{predicate}\index{cr-sep-sethood(mct_, mchs_, mcc_)@\texttt{cr-sep-sethood(mct\_, mchs\_, mcc\_)}}\\
reads: the core sep-sethood relates \$2 to \$3 \\
defined by \texttt{cr-sep-sethood} \\
characterized by \texttt{cr-sep-sethood}  \textit{(definitional)}:
\begin{quote}\ttfamily\footnotesize\raggedright\noindent
forall(\allowbreak{}[mct\_, mchs\_, mcc\_], cr-sep-sethood(\allowbreak{}mct\_, mchs\_, mcc\_) iff forsome(\allowbreak{}[mcg\_, mca\_, mcp\_], mcc\_ = [mcg\_, enc-in(\allowbreak{}enc-sep(\allowbreak{}mca\_, mcp\_), enc-setc)] and nbg-ok(\allowbreak{}mcp\_) and mchs\_ = [[mcg\_, enc-in(\allowbreak{}mca\_, enc-setc)]]))
\end{quote}
mentioned by 5 result(s)

\smallskip

\noindent\textbf{\texttt{cr-tbind-cr(mct\_, mchs\_, mcc\_)}}\quad\textit{predicate}\index{cr-tbind-cr(mct_, mchs_, mcc_)@\texttt{cr-tbind-cr(mct\_, mchs\_, mcc\_)}}\\
reads: the core congruence rule for the term binders relates \$2 to \$3 \\
defined by \texttt{cr-tbind-cr} \\
characterized by \texttt{cr-tbind-cr}  \textit{(definitional)}:
\begin{quote}\ttfamily\footnotesize\raggedright\noindent
forall(\allowbreak{}[mct\_, mchs\_, mcc\_], cr-tbind-cr(\allowbreak{}mct\_, mchs\_, mcc\_) iff forsome(\allowbreak{}[mcg\_, mchd\_, mcol\_, mcp\_, mcq\_, mcy\_, mcf\_], mcc\_ = [mcg\_, enc-qeq(\allowbreak{}ebind(\allowbreak{}mchd\_, 1, mcol\_, mcp\_), ebind(\allowbreak{}mchd\_, 1, mcol\_, mcq\_))] and mcy\_ in nn and not(\allowbreak{}mcy\_ in union(\allowbreak{}ctx-fv(\allowbreak{}mcg\_), expr-fv(\allowbreak{}enc-qeq(\allowbreak{}ebind(\allowbreak{}mchd\_, 1, mcol\_, mcp\_), ebind(\allowbreak{}mchd\_, 1, mcol\_, mcq\_))))) and (\allowbreak{}(\allowbreak{}mchd\_ = 22 or mchd\_ = 23) and mcf\_ = enc-iff(\allowbreak{}expr-open(\allowbreak{}mcp\_, 0, efvar(\allowbreak{}mcy\_)), expr-open(\allowbreak{}mcq\_, 0, efvar(\allowbreak{}mcy\_))) or (\allowbreak{}mchd\_ = 24 or mchd\_ = 26) and mcf\_ = enc-qeq(\allowbreak{}expr-open(\allowbreak{}mcp\_, 0, efvar(\allowbreak{}mcy\_)), expr-open(\allowbreak{}mcq\_, 0, efvar(\allowbreak{}mcy\_)))) and mchs\_ = [[mcg\_, mcf\_]]))
\end{quote}
mentioned by 7 result(s)

\smallskip

\noindent\textbf{\texttt{cr-theorem-assumption(mct\_, mchs\_, mcc\_)}}\quad\textit{predicate}\index{cr-theorem-assumption(mct_, mchs_, mcc_)@\texttt{cr-theorem-assumption(mct\_, mchs\_, mcc\_)}}\\
reads: the core theorem-assumption relates \$2 to \$3 \\
defined by \texttt{cr-theorem-assumption} \\
characterized by \texttt{cr-theorem-assumption}  \textit{(definitional)}:
\begin{quote}\ttfamily\footnotesize\raggedright\noindent
forall(\allowbreak{}[mct\_, mchs\_, mcc\_], cr-theorem-assumption(\allowbreak{}mct\_, mchs\_, mcc\_) iff forsome(\allowbreak{}[mcg\_, mca\_], mcc\_ = [mcg\_, mca\_] and mca\_ in mct\_ and mchs\_ = []))
\end{quote}
mentioned by 5 result(s)

\smallskip

\noindent\textbf{\texttt{cr-truth-intro(mct\_, mchs\_, mcc\_)}}\quad\textit{predicate}\index{cr-truth-intro(mct_, mchs_, mcc_)@\texttt{cr-truth-intro(mct\_, mchs\_, mcc\_)}}\\
reads: the core truth-intro relates \$2 to \$3 \\
defined by \texttt{cr-truth-intro} \\
characterized by \texttt{cr-truth-intro}  \textit{(definitional)}:
\begin{quote}\ttfamily\footnotesize\raggedright\noindent
forall(\allowbreak{}[mct\_, mchs\_, mcc\_], cr-truth-intro(\allowbreak{}mct\_, mchs\_, mcc\_) iff forsome(\allowbreak{}[mcg\_], mcc\_ = [mcg\_, enc-truth] and mchs\_ = []))
\end{quote}
mentioned by 5 result(s)

\smallskip

\noindent\textbf{\texttt{cr-tuple-eq(mct\_, mchs\_, mcc\_)}}\quad\textit{predicate}\index{cr-tuple-eq(mct_, mchs_, mcc_)@\texttt{cr-tuple-eq(mct\_, mchs\_, mcc\_)}}\\
reads: the core tuple-eq relates \$2 to \$3 \\
defined by \texttt{cr-tuple-eq} \\
characterized by \texttt{cr-tuple-eq}  \textit{(definitional)}:
\begin{quote}\ttfamily\footnotesize\raggedright\noindent
forall(\allowbreak{}[mct\_, mchs\_, mcc\_], cr-tuple-eq(\allowbreak{}mct\_, mchs\_, mcc\_) iff forsome(\allowbreak{}[mcg\_, mcl\_, mca\_], mcc\_ = [mcg\_, enc-iff(\allowbreak{}enc-eq(\allowbreak{}enc-list(\allowbreak{}mcl\_), enc-list(\allowbreak{}mca\_)), enc-eql(\allowbreak{}mcl\_, mca\_, length(\allowbreak{}mcl\_)))] and length(\allowbreak{}mcl\_) = length(\allowbreak{}mca\_) and mchs\_ = [[mcg\_, enc-eq(\allowbreak{}enc-list(\allowbreak{}mcl\_), enc-list(\allowbreak{}mcl\_))], [mcg\_, enc-eq(\allowbreak{}enc-list(\allowbreak{}mca\_), enc-list(\allowbreak{}mca\_))]]))
\end{quote}
mentioned by 5 result(s)

\smallskip

\noindent\textbf{\texttt{cr-tuple-eta(mct\_, mchs\_, mcc\_)}}\quad\textit{predicate}\index{cr-tuple-eta(mct_, mchs_, mcc_)@\texttt{cr-tuple-eta(mct\_, mchs\_, mcc\_)}}\\
reads: the core tuple-eta relates \$2 to \$3 \\
defined by \texttt{cr-tuple-eta} \\
characterized by \texttt{cr-tuple-eta}  \textit{(definitional)}:
\begin{quote}\ttfamily\footnotesize\raggedright\noindent
forall(\allowbreak{}[mct\_, mchs\_, mcc\_], cr-tuple-eta(\allowbreak{}mct\_, mchs\_, mcc\_) iff forsome(\allowbreak{}[mcg\_, mcx\_, mca\_, mcl\_], mcc\_ = [mcg\_, enc-eq(\allowbreak{}mcx\_, enc-list(\allowbreak{}mcl\_))] and 1 <= length(\allowbreak{}mca\_) and length(\allowbreak{}mcl\_) = length(\allowbreak{}mca\_) and forall(\allowbreak{}[mci\_ in nn], 1 <= mci\_ implies mci\_ <= length(\allowbreak{}mca\_) implies nth(\allowbreak{}mci\_, mcl\_) = enc-nth(\allowbreak{}enc-num(\allowbreak{}mci\_), mcx\_)) and mchs\_ = [[mcg\_, enc-in(\allowbreak{}mcx\_, enc-cart(\allowbreak{}mca\_))]]))
\end{quote}
mentioned by 5 result(s)

\smallskip

\noindent\textbf{\texttt{cr-tuples-mem(mct\_, mchs\_, mcc\_)}}\quad\textit{predicate}\index{cr-tuples-mem(mct_, mchs_, mcc_)@\texttt{cr-tuples-mem(mct\_, mchs\_, mcc\_)}}\\
reads: the core tuples-mem relates \$2 to \$3 \\
defined by \texttt{cr-tuples-mem} \\
characterized by \texttt{cr-tuples-mem}  \textit{(definitional)}:
\begin{quote}\ttfamily\footnotesize\raggedright\noindent
forall(\allowbreak{}[mct\_, mchs\_, mcc\_], cr-tuples-mem(\allowbreak{}mct\_, mchs\_, mcc\_) iff forsome(\allowbreak{}[mcg\_, mcl\_, mca\_], mcc\_ = [mcg\_, enc-iff(\allowbreak{}enc-in(\allowbreak{}enc-list(\allowbreak{}mcl\_), enc-tuples(\allowbreak{}mca\_)), enc-ande(\allowbreak{}mcl\_, mca\_, length(\allowbreak{}mcl\_)))] and mchs\_ = []))
\end{quote}
mentioned by 5 result(s)

\smallskip

\noindent\textbf{\texttt{cr-union-mem(mct\_, mchs\_, mcc\_)}}\quad\textit{predicate}\index{cr-union-mem(mct_, mchs_, mcc_)@\texttt{cr-union-mem(mct\_, mchs\_, mcc\_)}}\\
reads: the core union-mem relates \$2 to \$3 \\
defined by \texttt{cr-union-mem} \\
characterized by \texttt{cr-union-mem}  \textit{(definitional)}:
\begin{quote}\ttfamily\footnotesize\raggedright\noindent
forall(\allowbreak{}[mct\_, mchs\_, mcc\_], cr-union-mem(\allowbreak{}mct\_, mchs\_, mcc\_) iff forsome(\allowbreak{}[mcg\_, mcx\_, mcl\_], mcc\_ = [mcg\_, enc-iff(\allowbreak{}enc-in(\allowbreak{}mcx\_, enc-union(\allowbreak{}mcl\_)), enc-orl(\allowbreak{}mcx\_, mcl\_, length(\allowbreak{}mcl\_)))] and mchs\_ = []))
\end{quote}
mentioned by 5 result(s)

\smallskip

\noindent\textbf{\texttt{cr-weakening(mct\_, mchs\_, mcc\_)}}\quad\textit{predicate}\index{cr-weakening(mct_, mchs_, mcc_)@\texttt{cr-weakening(mct\_, mchs\_, mcc\_)}}\\
reads: the core weakening relates \$2 to \$3 \\
defined by \texttt{cr-weakening} \\
characterized by \texttt{cr-weakening}  \textit{(definitional)}:
\begin{quote}\ttfamily\footnotesize\raggedright\noindent
forall(\allowbreak{}[mct\_, mchs\_, mcc\_], cr-weakening(\allowbreak{}mct\_, mchs\_, mcc\_) iff forsome(\allowbreak{}[mcg\_, mcf\_, mcd\_], mcc\_ = [mcg\_, mcf\_] and mcd\_ subset mcg\_ and mchs\_ = [[mcd\_, mcf\_]]))
\end{quote}
mentioned by 5 result(s)

\smallskip

\noindent\textbf{\texttt{crel(m)}}\quad\textit{functoid}\index{crel(m)@\texttt{crel(m)}}\\
definition:
\begin{quote}\ttfamily\footnotesize\raggedright\noindent
\{p in cartesian(\allowbreak{}cseq(\allowbreak{}m), cseq(\allowbreak{}m)): cseq-equiv(\allowbreak{}m, nth(\allowbreak{}1, p), nth(\allowbreak{}2, p))\}
\end{quote}
characterized by \texttt{rko2-crel-elim}  \textit{(proven)}:
\begin{quote}\ttfamily\footnotesize\raggedright\noindent
forall(\allowbreak{}[m, p\_, q\_], [p\_, q\_] in crel(\allowbreak{}m) implies cseq-equiv(\allowbreak{}m, p\_, q\_))
\end{quote}
also: \texttt{rko2-crel-intro} \texttt{rko2-crel-is-set} \texttt{rko2-crel-unfold} \\
mentioned by 11 result(s)

\smallskip

\noindent\textbf{\texttt{crip(v)}}\quad\textit{functoid}\index{crip(v)@\texttt{crip(v)}}\\
reads: the complex form of \$1 \\
definition:
\begin{quote}\ttfamily\footnotesize\raggedright\noindent
vnb-lambda(\allowbreak{}[cxx\_, cxy\_], cartesian(\allowbreak{}vec(\allowbreak{}v), vec(\allowbreak{}v)), (\allowbreak{}rip(\allowbreak{}v))(\allowbreak{}cxx\_, cxy\_) - 1i * (\allowbreak{}rip(\allowbreak{}v))(\allowbreak{}(\allowbreak{}jop(\allowbreak{}v))(\allowbreak{}cxx\_), cxy\_))
\end{quote}
characterized by \texttt{csip-crip-fun}  \textit{(proven)}:
\begin{quote}\ttfamily\footnotesize\raggedright\noindent
forall(\allowbreak{}[v], is-complex-semi-inner-product-space(\allowbreak{}v) implies crip(\allowbreak{}v) in fun(\allowbreak{}cartesian(\allowbreak{}vec(\allowbreak{}v), vec(\allowbreak{}v)), cc))
\end{quote}
also: \texttt{csip-crip-self} \texttt{csip-crip-value} \texttt{csip-crip-unique} \texttt{csip-crip-type-pt} \texttt{csip-crip-add-left} \texttt{csip-crip-conj-sym} \texttt{csip-crip-homog-left} \texttt{csip-crip-sesquilinear} \\
mentioned by 17 result(s)

\smallskip

\noindent\textbf{\texttt{crs}}\quad\textit{tactic}\index{crs@\texttt{crs}}\\
\texttt{(crs)} -- Commutative-ring decision procedure: prove a polynomial identity over ZZ[generators].  Expands literal powers, so (x+y)\textasciicircum{}2 = ... closes directly. \\
declared in \texttt{tactics-help.scm}

\smallskip

\noindent\textbf{\texttt{cseq(m)}}\quad\textit{functoid}\index{cseq(m)@\texttt{cseq(m)}}\\
definition:
\begin{quote}\ttfamily\footnotesize\raggedright\noindent
\{f in fun(\allowbreak{}nn, pts(\allowbreak{}m)): is-cauchy-seq(\allowbreak{}m, f)\}
\end{quote}
characterized by \texttt{qp-cs-cseq}  \textit{(proven)}:
\begin{quote}\ttfamily\footnotesize\raggedright\noindent
forall(\allowbreak{}[p], qp-cs(\allowbreak{}p) == cseq(\allowbreak{}nf-metric-space(\allowbreak{}qq-padic(\allowbreak{}p))))
\end{quote}
also: \texttt{r7q-cseq-cauchy} \texttt{rkt-cseq-is-set} \texttt{form-cseq-unfold} \texttt{rko2-cseq-in-fun} \texttt{rko2-cseq-unfold} \texttt{r12d-diag-in-cseq} \texttt{rkt-embed-seq-in-cseq} \\
mentioned by 33 result(s)

\smallskip

\noindent\textbf{\texttt{cseq-equiv(m, f, g)}}\quad\textit{predicate}\index{cseq-equiv(m, f, g)@\texttt{cseq-equiv(m, f, g)}}\\
reads: \$2 and \$3 are equivalent Cauchy sequences in \$1 \\
defined by \texttt{cseq-equiv} \\
characterized by \texttt{cseq-equiv}  \textit{(definitional)}:
\begin{quote}\ttfamily\footnotesize\raggedright\noindent
forall(\allowbreak{}[m, f, g], cseq-equiv(\allowbreak{}m, f, g) iff converges-to(\allowbreak{}rr-ms, dist-seq(\allowbreak{}m, f, g), 0))
\end{quote}
mentioned by 14 result(s)

\smallskip

\noindent\textbf{\texttt{cseq-module(md, fq)}}\quad\textit{functoid}\index{cseq-module(md, fq)@\texttt{cseq-module(md, fq)}}\\
reads: the module of the Cauchy sequences of (\$1, \$2) \\
definition:
\begin{quote}\ttfamily\footnotesize\raggedright\noindent
[scal(\allowbreak{}md), form-cseq(\allowbreak{}md, fq), vnb-lambda(\allowbreak{}[hcf\_, hcg\_], cartesian(\allowbreak{}form-cseq(\allowbreak{}md, fq), form-cseq(\allowbreak{}md, fq)), seq-vadd(\allowbreak{}md, hcf\_, hcg\_)), seq-zero(\allowbreak{}md), vnb-lambda(\allowbreak{}hcf\_, form-cseq(\allowbreak{}md, fq), seq-vneg(\allowbreak{}md, hcf\_)), vnb-lambda(\allowbreak{}[hca\_, hcf\_], cartesian(\allowbreak{}cc, form-cseq(\allowbreak{}md, fq)), seq-act(\allowbreak{}md, hca\_, hcf\_))]
\end{quote}
characterized by \texttt{cseq-module-act}  \textit{(proven)}:
\begin{quote}\ttfamily\footnotesize\raggedright\noindent
forall(\allowbreak{}[md, fq], is-module(\allowbreak{}md) implies scal(\allowbreak{}md) = normed-field-as-commutative-ring(\allowbreak{}cc-normed-field) implies is-inner-product-form(\allowbreak{}md, fq) implies act(\allowbreak{}cseq-module(\allowbreak{}md, fq)) == vnb-lambda(\allowbreak{}[hca\_, hcf\_], cartesian(\allowbreak{}cc, form-cseq(\allowbreak{}md, fq)), seq-act(\allowbreak{}md, hca\_, hcf\_)))
\end{quote}
also: \texttt{cseq-module-vec} \texttt{cseq-module-scal} \texttt{cseq-module-vadd} \texttt{cseq-module-vneg} \texttt{cseq-module-vzero} \texttt{cseq-module-is-module} \texttt{cseq-module-act-computes} \texttt{cseq-module-vadd-computes} ... (9 in all) \\
mentioned by 34 result(s)

\smallskip

\noindent\textbf{\texttt{csip-as-sip(s)}}\quad\textit{view}\index{csip-as-sip(s)@\texttt{csip-as-sip(s)}}\\
a complex-semi-inner-product-space viewed as a semi-inner-product-space \\
characterized by \texttt{sipq-act-csip-as-sip}  \textit{(definitional)}:
\begin{quote}\ttfamily\footnotesize\raggedright\noindent
forall(\allowbreak{}[r], is-complex-semi-inner-product-space(\allowbreak{}r) implies act(\allowbreak{}sip-quotient(\allowbreak{}csip-as-sip(\allowbreak{}r))) == sip-qact(\allowbreak{}csip-as-sip(\allowbreak{}r)))
\end{quote}
also: \texttt{sipq-rip-csip-as-sip} \texttt{sipq-vec-csip-as-sip} \texttt{sipq-comm-csip-as-sip} \texttt{sipq-scal-csip-as-sip} \texttt{sipq-vadd-csip-as-sip} \texttt{sipq-vneg-csip-as-sip} \texttt{sip-rel-in-csip-as-sip} \texttt{sipq-assoc-csip-as-sip} ... (215 in all) \\
mentioned by 715 result(s)

\smallskip

\noindent\textbf{\texttt{csip-cips(v)}}\quad\textit{functoid}\index{csip-cips(v)@\texttt{csip-cips(v)}}\\
reads: \$1 as a complex inner product space \\
definition:
\begin{quote}\ttfamily\footnotesize\raggedright\noindent
[normed-field-as-commutative-ring(\allowbreak{}cc-normed-field), vec(\allowbreak{}v), vadd(\allowbreak{}v), vzero(\allowbreak{}v), vneg(\allowbreak{}v), cact(\allowbreak{}v), crip(\allowbreak{}v)]
\end{quote}
characterized by \texttt{csip-cips-ip}  \textit{(proven)}:
\begin{quote}\ttfamily\footnotesize\raggedright\noindent
forall(\allowbreak{}[v], is-complex-semi-inner-product-space(\allowbreak{}v) implies ip(\allowbreak{}csip-cips(\allowbreak{}v)) == crip(\allowbreak{}v))
\end{quote}
also: \texttt{csip-cips-act} \texttt{csip-cips-vec} \texttt{csip-cips-comm} \texttt{csip-cips-scal} \texttt{csip-cips-vadd} \texttt{csip-cips-vneg} \texttt{csip-cips-assoc} \texttt{csip-cips-vzero} ... (30 in all) \\
mentioned by 98 result(s)

\smallskip

\noindent\textbf{\texttt{csip-cnvs(v)}}\quad\textit{functoid}\index{csip-cnvs(v)@\texttt{csip-cnvs(v)}}\\
reads: \$1 as a complex normed vector space \\
definition:
\begin{quote}\ttfamily\footnotesize\raggedright\noindent
cips-cnvs(\allowbreak{}csip-cips(\allowbreak{}v))
\end{quote}
characterized by \texttt{csip-cnvs-vec}  \textit{(proven)}:
\begin{quote}\ttfamily\footnotesize\raggedright\noindent
forall(\allowbreak{}[v], is-complex-pre-hilbert-space(\allowbreak{}v) implies vec(\allowbreak{}csip-cnvs(\allowbreak{}v)) == vec(\allowbreak{}v))
\end{quote}
also: \texttt{csip-cnvs-cact} \texttt{csip-cnvs-vadd} \texttt{csip-cnvs-vneg} \texttt{csip-cnvs-vzero} \texttt{csip-cnvs-unfold} \texttt{csip-cnvs-vnrm-value} \\
mentioned by 15 result(s)

\smallskip

\noindent\textbf{\texttt{csip-qj(v)}}\quad\textit{functoid}\index{csip-qj(v)@\texttt{csip-qj(v)}}\\
definition:
\begin{quote}\ttfamily\footnotesize\raggedright\noindent
descend(\allowbreak{}sip-setoid(\allowbreak{}csip-as-sip(\allowbreak{}v)), vnb-lambda(\allowbreak{}sqa\_, vec(\allowbreak{}v), class(\allowbreak{}sip-setoid(\allowbreak{}csip-as-sip(\allowbreak{}v)), (\allowbreak{}jop(\allowbreak{}v))(\allowbreak{}sqa\_))))
\end{quote}
characterized by \texttt{csip-qj-fun}  \textit{(proven)}:
\begin{quote}\ttfamily\footnotesize\raggedright\noindent
forall(\allowbreak{}[v], is-complex-semi-inner-product-space(\allowbreak{}v) implies csip-qj(\allowbreak{}v) in fun(\allowbreak{}quotient(\allowbreak{}sip-setoid(\allowbreak{}csip-as-sip(\allowbreak{}v))), quotient(\allowbreak{}sip-setoid(\allowbreak{}csip-as-sip(\allowbreak{}v)))))
\end{quote}
also: \texttt{csip-qj-computes} \\
mentioned by 5 result(s)

\smallskip

\noindent\textbf{\texttt{csip-quotient(v)}}\quad\textit{functoid}\index{csip-quotient(v)@\texttt{csip-quotient(v)}}\\
reads: the quotient of \$1 by its null vectors \\
definition:
\begin{quote}\ttfamily\footnotesize\raggedright\noindent
[scal(\allowbreak{}v), quotient(\allowbreak{}sip-setoid(\allowbreak{}csip-as-sip(\allowbreak{}v))), sip-qvadd(\allowbreak{}csip-as-sip(\allowbreak{}v)), class(\allowbreak{}sip-setoid(\allowbreak{}csip-as-sip(\allowbreak{}v)), vzero(\allowbreak{}v)), sip-qvneg(\allowbreak{}csip-as-sip(\allowbreak{}v)), sip-qact(\allowbreak{}csip-as-sip(\allowbreak{}v)), sip-qrip(\allowbreak{}csip-as-sip(\allowbreak{}v)), csip-qj(\allowbreak{}v)]
\end{quote}
characterized by \texttt{csip-quotient-act}  \textit{(proven)}:
\begin{quote}\ttfamily\footnotesize\raggedright\noindent
forall(\allowbreak{}[v], is-complex-semi-inner-product-space(\allowbreak{}v) implies act(\allowbreak{}csip-quotient(\allowbreak{}v)) == sip-qact(\allowbreak{}csip-as-sip(\allowbreak{}v)))
\end{quote}
also: \texttt{csip-quotient-jop} \texttt{csip-quotient-rip} \texttt{csip-quotient-sip} \texttt{csip-quotient-vec} \texttt{csip-quotient-scal} \texttt{csip-quotient-vadd} \texttt{csip-quotient-vneg} \texttt{csip-quotient-vzero} ... (27 in all) \\
mentioned by 45 result(s)

\smallskip

\noindent\textbf{\texttt{cstar-gen(b2va\_, b2vx\_)}}\quad\textit{functoid}\index{cstar-gen(b2va_, b2vx_)@\texttt{cstar-gen(b2va\_, b2vx\_)}}\\
reads: the C*-subalgebra of \$1 generated by \$2 \\
definition:
\begin{quote}\ttfamily\footnotesize\raggedright\noindent
intersection-of(\allowbreak{}\{b2vs\_ in power(\allowbreak{}vec(\allowbreak{}b2va\_)): is-closed-star-subalgebra-set(\allowbreak{}b2va\_, b2vs\_) and b2vx\_ in b2vs\_\})
\end{quote}
characterized by \texttt{cstar-gen-least}  \textit{(proven)}:
\begin{quote}\ttfamily\footnotesize\raggedright\noindent
forall(\allowbreak{}[a], is-unital-involutive-banach-algebra(\allowbreak{}a) implies forall(\allowbreak{}[x, s], is-closed-star-subalgebra-set(\allowbreak{}a, s) implies x in s implies cstar-gen(\allowbreak{}a, x) subset s))
\end{quote}
also: \texttt{cstar-gen-subset} \texttt{cstar-gen-contains} \\
mentioned by 3 result(s)

\smallskip

\noindent\textbf{\texttt{ctx-add(mcg\_, mca\_)}}\quad\textit{functoid}\index{ctx-add(mcg_, mca_)@\texttt{ctx-add(mcg\_, mca\_)}}\\
reads: the context \$1 extended by \$2 \\
definition:
\begin{quote}\ttfamily\footnotesize\raggedright\noindent
union(\allowbreak{}mcg\_, pair(\allowbreak{}mca\_, mca\_))
\end{quote}
characterized by \texttt{ctx-add-ctx}  \textit{(proven)}:
\begin{quote}\ttfamily\footnotesize\raggedright\noindent
forall(\allowbreak{}[g, a], g in ctx-set implies a in clf implies ctx-add(\allowbreak{}g, a) in ctx-set)
\end{quote}
also: \texttt{ctx-add-has} \texttt{ctx-add-sub} \texttt{ctx-add-keeps} \texttt{ctx-add-unfold} \\
mentioned by 89 result(s)

\smallskip

\noindent\textbf{\texttt{ctx-fv(mcg\_)}}\quad\textit{functoid}\index{ctx-fv(mcg_)@\texttt{ctx-fv(mcg\_)}}\\
reads: the free names of the context \$1 \\
definition:
\begin{quote}\ttfamily\footnotesize\raggedright\noindent
big-union(\allowbreak{}mcw\_, mcg\_, expr-fv(\allowbreak{}mcw\_))
\end{quote}
characterized by \texttt{ctx-fv-add}  \textit{(proven)}:
\begin{quote}\ttfamily\footnotesize\raggedright\noindent
forall(\allowbreak{}[g, h, y], h in set implies y in ctx-fv(\allowbreak{}ctx-add(\allowbreak{}g, h)) implies y in ctx-fv(\allowbreak{}g) or y in expr-fv(\allowbreak{}h))
\end{quote}
also: \texttt{ctx-fv-minus} \texttt{ctx-fv-finite} \texttt{ctx-fv-member} \texttt{ctx-fv-unfold} \\
mentioned by 23 result(s)

\smallskip

\noindent\textbf{\texttt{ctx-set}}\quad\textit{constant}\index{ctx-set@\texttt{ctx-set}}\\
reads: the contexts (finite sets of sentences) \\
defined by \texttt{ctx-def}

\smallskip

\noindent\textbf{\texttt{cu-pairs(uat\_)}}\quad\textit{functoid}\index{cu-pairs(uat_)@\texttt{cu-pairs(uat\_)}}\\
reads: the pairs [C, U] of a closed set inside an open set of \$1 \\
definition:
\begin{quote}\ttfamily\footnotesize\raggedright\noindent
\{ucw\_ in cartesian(\allowbreak{}power(\allowbreak{}pts(\allowbreak{}uat\_)), opens(\allowbreak{}uat\_)): is-top-closed(\allowbreak{}uat\_, nth(\allowbreak{}1, ucw\_)) and nth(\allowbreak{}1, ucw\_) subset nth(\allowbreak{}2, ucw\_)\}
\end{quote}
mentioned by 41 result(s)

\smallskip

\noindent\textbf{\texttt{cut}}\quad\textit{tactic}\index{cut@\texttt{cut}}\\
\texttt{(cut formula)} -- Cut: prove \texttt{formula' as a side subgoal, then continue the main goal with }formula' added as an assumption (Gentzen cut). \\
declared in \texttt{tactics-help.scm}

\smallskip

\noindent\textbf{\texttt{cx-act(v)}}\quad\textit{functoid}\index{cx-act(v)@\texttt{cx-act(v)}}\\
definition:
\begin{quote}\ttfamily\footnotesize\raggedright\noindent
vnb-lambda(\allowbreak{}[cr\_, cp\_], cartesian(\allowbreak{}carr(\allowbreak{}scal(\allowbreak{}v)), cartesian(\allowbreak{}vec(\allowbreak{}v), vec(\allowbreak{}v))), [(\allowbreak{}act(\allowbreak{}v))(\allowbreak{}cr\_, nth(\allowbreak{}1, cp\_)), (\allowbreak{}act(\allowbreak{}v))(\allowbreak{}cr\_, nth(\allowbreak{}2, cp\_))])
\end{quote}
characterized by \texttt{cx-act-fun}  \textit{(proven)}:
\begin{quote}\ttfamily\footnotesize\raggedright\noindent
forall(\allowbreak{}[v], is-semi-inner-product-space(\allowbreak{}v) implies cx-act(\allowbreak{}v) in fun(\allowbreak{}cartesian(\allowbreak{}rr, cartesian(\allowbreak{}vec(\allowbreak{}v), vec(\allowbreak{}v))), cartesian(\allowbreak{}vec(\allowbreak{}v), vec(\allowbreak{}v))))
\end{quote}
also: \texttt{cx-act-value} \\
mentioned by 9 result(s)

\smallskip

\noindent\textbf{\texttt{cx-c}}\quad\textit{instance}\index{cx-c@\texttt{cx-c}}\\
an instance of bundle; predicate \texttt{is-cx-c}; same shape as bundle

\smallskip

\noindent\textbf{\texttt{cx-isp}}\quad\textit{instance}\index{cx-isp@\texttt{cx-isp}}\\
an instance of infiltrated-space; predicate \texttt{is-cx-isp}; same shape as infiltrated-space

\smallskip

\noindent\textbf{\texttt{cx-j(v)}}\quad\textit{functoid}\index{cx-j(v)@\texttt{cx-j(v)}}\\
definition:
\begin{quote}\ttfamily\footnotesize\raggedright\noindent
vnb-lambda(\allowbreak{}cp\_, cartesian(\allowbreak{}vec(\allowbreak{}v), vec(\allowbreak{}v)), [(\allowbreak{}vneg(\allowbreak{}v))(\allowbreak{}nth(\allowbreak{}2, cp\_)), nth(\allowbreak{}1, cp\_)])
\end{quote}
characterized by \texttt{cx-j-fun}  \textit{(proven)}:
\begin{quote}\ttfamily\footnotesize\raggedright\noindent
forall(\allowbreak{}[v], is-semi-inner-product-space(\allowbreak{}v) implies cx-j(\allowbreak{}v) in fun(\allowbreak{}cartesian(\allowbreak{}vec(\allowbreak{}v), vec(\allowbreak{}v)), cartesian(\allowbreak{}vec(\allowbreak{}v), vec(\allowbreak{}v))))
\end{quote}
also: \texttt{cx-j-value} \\
mentioned by 5 result(s)

\smallskip

\noindent\textbf{\texttt{cx-orp}}\quad\textit{instance}\index{cx-orp@\texttt{cx-orp}}\\
an instance of sk-space; predicate \texttt{is-cx-orp}; same shape as sk-space

\smallskip

\noindent\textbf{\texttt{cx-rip(v)}}\quad\textit{functoid}\index{cx-rip(v)@\texttt{cx-rip(v)}}\\
definition:
\begin{quote}\ttfamily\footnotesize\raggedright\noindent
vnb-lambda(\allowbreak{}[cp\_, cq\_], cartesian(\allowbreak{}cartesian(\allowbreak{}vec(\allowbreak{}v), vec(\allowbreak{}v)), cartesian(\allowbreak{}vec(\allowbreak{}v), vec(\allowbreak{}v))), (\allowbreak{}rip(\allowbreak{}v))(\allowbreak{}nth(\allowbreak{}1, cp\_), nth(\allowbreak{}1, cq\_)) + (\allowbreak{}rip(\allowbreak{}v))(\allowbreak{}nth(\allowbreak{}2, cp\_), nth(\allowbreak{}2, cq\_)))
\end{quote}
characterized by \texttt{cx-rip-fun}  \textit{(proven)}:
\begin{quote}\ttfamily\footnotesize\raggedright\noindent
forall(\allowbreak{}[v], is-semi-inner-product-space(\allowbreak{}v) implies cx-rip(\allowbreak{}v) in fun(\allowbreak{}cartesian(\allowbreak{}cartesian(\allowbreak{}vec(\allowbreak{}v), vec(\allowbreak{}v)), cartesian(\allowbreak{}vec(\allowbreak{}v), vec(\allowbreak{}v))), rr))
\end{quote}
also: \texttt{cx-rip-value} \\
mentioned by 9 result(s)

\smallskip

\noindent\textbf{\texttt{cx-sip(v)}}\quad\textit{functoid}\index{cx-sip(v)@\texttt{cx-sip(v)}}\\
reads: the real space underlying the complexification of \$1 \\
definition:
\begin{quote}\ttfamily\footnotesize\raggedright\noindent
[scal(\allowbreak{}v), cartesian(\allowbreak{}vec(\allowbreak{}v), vec(\allowbreak{}v)), cx-vadd(\allowbreak{}v), [vzero(\allowbreak{}v), vzero(\allowbreak{}v)], cx-vneg(\allowbreak{}v), cx-act(\allowbreak{}v), cx-rip(\allowbreak{}v)]
\end{quote}
characterized by \texttt{cx-sip-act}  \textit{(proven)}:
\begin{quote}\ttfamily\footnotesize\raggedright\noindent
forall(\allowbreak{}[v], is-semi-inner-product-space(\allowbreak{}v) implies act(\allowbreak{}cx-sip(\allowbreak{}v)) == cx-act(\allowbreak{}v))
\end{quote}
also: \texttt{cx-sip-rip} \texttt{cx-sip-vec} \texttt{cx-sip-comm} \texttt{cx-sip-scal} \texttt{cx-sip-vadd} \texttt{cx-sip-vneg} \texttt{cx-sip-assoc} \texttt{cx-sip-vzero} ... (30 in all) \\
mentioned by 49 result(s)

\smallskip

\noindent\textbf{\texttt{cx-sp}}\quad\textit{instance}\index{cx-sp@\texttt{cx-sp}}\\
an instance of strand-space; predicate \texttt{is-cx-sp}; same shape as strand-space

\smallskip

\noindent\textbf{\texttt{cx-vadd(v)}}\quad\textit{functoid}\index{cx-vadd(v)@\texttt{cx-vadd(v)}}\\
definition:
\begin{quote}\ttfamily\footnotesize\raggedright\noindent
vnb-lambda(\allowbreak{}[cp\_, cq\_], cartesian(\allowbreak{}cartesian(\allowbreak{}vec(\allowbreak{}v), vec(\allowbreak{}v)), cartesian(\allowbreak{}vec(\allowbreak{}v), vec(\allowbreak{}v))), [(\allowbreak{}vadd(\allowbreak{}v))(\allowbreak{}nth(\allowbreak{}1, cp\_), nth(\allowbreak{}1, cq\_)), (\allowbreak{}vadd(\allowbreak{}v))(\allowbreak{}nth(\allowbreak{}2, cp\_), nth(\allowbreak{}2, cq\_))])
\end{quote}
characterized by \texttt{cx-vadd-fun}  \textit{(proven)}:
\begin{quote}\ttfamily\footnotesize\raggedright\noindent
forall(\allowbreak{}[v], is-semi-inner-product-space(\allowbreak{}v) implies cx-vadd(\allowbreak{}v) in fun(\allowbreak{}cartesian(\allowbreak{}cartesian(\allowbreak{}vec(\allowbreak{}v), vec(\allowbreak{}v)), cartesian(\allowbreak{}vec(\allowbreak{}v), vec(\allowbreak{}v))), cartesian(\allowbreak{}vec(\allowbreak{}v), vec(\allowbreak{}v))))
\end{quote}
also: \texttt{cx-vadd-value} \\
mentioned by 9 result(s)

\smallskip

\noindent\textbf{\texttt{cx-vneg(v)}}\quad\textit{functoid}\index{cx-vneg(v)@\texttt{cx-vneg(v)}}\\
definition:
\begin{quote}\ttfamily\footnotesize\raggedright\noindent
vnb-lambda(\allowbreak{}cp\_, cartesian(\allowbreak{}vec(\allowbreak{}v), vec(\allowbreak{}v)), [(\allowbreak{}vneg(\allowbreak{}v))(\allowbreak{}nth(\allowbreak{}1, cp\_)), (\allowbreak{}vneg(\allowbreak{}v))(\allowbreak{}nth(\allowbreak{}2, cp\_))])
\end{quote}
characterized by \texttt{cx-vneg-fun}  \textit{(proven)}:
\begin{quote}\ttfamily\footnotesize\raggedright\noindent
forall(\allowbreak{}[v], is-semi-inner-product-space(\allowbreak{}v) implies cx-vneg(\allowbreak{}v) in fun(\allowbreak{}cartesian(\allowbreak{}vec(\allowbreak{}v), vec(\allowbreak{}v)), cartesian(\allowbreak{}vec(\allowbreak{}v), vec(\allowbreak{}v))))
\end{quote}
also: \texttt{cx-vneg-value} \\
mentioned by 9 result(s)

\smallskip

\noindent\textbf{\texttt{cyclic-subspace(gcs-md\_, gcs-t\_, gcs-v\_)}}\quad\textit{functoid}\index{cyclic-subspace(gcs-md_, gcs-t_, gcs-v_)@\texttt{cyclic-subspace(gcs-md\_, gcs-t\_, gcs-v\_)}}\\
reads: the \$2-cyclic subspace generated by \$3 \\
definition:
\begin{quote}\ttfamily\footnotesize\raggedright\noindent
\{gcsz\_ in vec(\allowbreak{}gcs-md\_): forsome(\allowbreak{}[gcsq\_ in carr(\allowbreak{}poly(\allowbreak{}scal(\allowbreak{}gcs-md\_)))], gcsz\_ = poly-op(\allowbreak{}gcs-md\_, gcs-t\_, gcsq\_, gcs-v\_))\}
\end{quote}

\smallskip

\noindent\textbf{\texttt{cyclic-sum(l, a)}}\quad\textit{functoid}\index{cyclic-sum(l, a)@\texttt{cyclic-sum(l, a)}}\\
reads: the direct sum of the cyclic groups Z/(\$2\_i), i = 1, ..., \$1 \\
definition:
\begin{quote}\ttfamily\footnotesize\raggedright\noindent
direct-sum(\allowbreak{}zz-ring, l, vnb-lambda(\allowbreak{}csi\_, interval(\allowbreak{}1, l), zz-mod(\allowbreak{}a(\allowbreak{}csi\_))))
\end{quote}
characterized by \texttt{cyclic-sum-iso-core}  \textit{(proven)}:
\begin{quote}\ttfamily\footnotesize\raggedright\noindent
forall(\allowbreak{}[n, al, be, h], n in nn implies al in fun(\allowbreak{}interval(\allowbreak{}1, n), zz) implies be in fun(\allowbreak{}interval(\allowbreak{}1, n), zz) implies forall(\allowbreak{}[ifi\_ in interval(\allowbreak{}1, n)], 0 <= al(\allowbreak{}ifi\_)) implies forall(\allowbreak{}[ifi\_ in interval(\allowbreak{}1, n)], 0 <= be(\allowbreak{}ifi\_)) implies forall(\allowbreak{}[ifi\_, ifj\_], ifi\_ in interval(\allowbreak{}1, n) implies ifj\_ in interval(\allowbreak{}1, n) implies ifi\_ <= ifj\_ implies zz-divides(\allowbreak{}al(\allowbreak{}ifi\_), al(\allowbreak{}ifj\_))) implies forall(\allowbreak{}[ifi\_, ifj\_], ifi\_ in interval(\allowbreak{}1, n) implies ifj\_ in interval(\allowbreak{}1, n) implies ifi\_ <= ifj\_ implies zz-divides(\allowbreak{}be(\allowbreak{}ifi\_), be(\allowbreak{}ifj\_))) implies is-iso-module(\allowbreak{}cyclic-sum(\allowbreak{}n, al), cyclic-sum(\allowbreak{}n, be), h) implies al = be)
\end{quote}
also: \texttt{cyclic-sum-iso-unique} \\
mentioned by 18 result(s)

\smallskip

\noindent\textbf{\texttt{cyl(cyw\_, cyn\_)}}\quad\textit{functoid}\index{cyl(cyw_, cyn_)@\texttt{cyl(cyw\_, cyn\_)}}\\
reads: the sequences of naturals agreeing with \$1 below \$2 \\
definition:
\begin{quote}\ttfamily\footnotesize\raggedright\noindent
\{cyz\_ in fun(\allowbreak{}nn, nn): forall(\allowbreak{}[cyj\_ in nn], cyj\_ < cyn\_ implies cyz\_(\allowbreak{}cyj\_) = cyw\_(\allowbreak{}cyj\_))\}
\end{quote}
characterized by \texttt{cyl-sub}  \textit{(proven)}:
\begin{quote}\ttfamily\footnotesize\raggedright\noindent
forall(\allowbreak{}[w in fun(\allowbreak{}nn, nn), n in nn], cyl(\allowbreak{}w, n) subset fun(\allowbreak{}nn, nn))
\end{quote}
also: \texttt{cyl-zero} \texttt{cyl-agree} \texttt{cyl-limit} \texttt{cyl-split} \texttt{cyl-unfold} \texttt{image-cyl-power} \texttt{image-cyl-split} \\
mentioned by 14 result(s)

\smallskip


\section*{D}
\addcontentsline{toc}{section}{D}

\noindent\textbf{\texttt{dbi(m, xs)}}\quad\textit{functoid}\index{dbi(m, xs)@\texttt{dbi(m, xs)}}\\
reads: the intervals [-||x\_n||, ||x\_n||] of the sequence \$2 in \$1 \\
definition:
\begin{quote}\ttfamily\footnotesize\raggedright\noindent
vnb-lambda(\allowbreak{}dbn\_, nn, subspace-ms(\allowbreak{}rr-ms, ccint(\allowbreak{}-(\allowbreak{}vnrm(\allowbreak{}m))(\allowbreak{}xs(\allowbreak{}dbn\_)), (\allowbreak{}vnrm(\allowbreak{}m))(\allowbreak{}xs(\allowbreak{}dbn\_)))))
\end{quote}
characterized by \texttt{dbi-at}  \textit{(proven)}:
\begin{quote}\ttfamily\footnotesize\raggedright\noindent
forall(\allowbreak{}[m, xs, n], is-normed-vector-space(\allowbreak{}m) implies xs in fun(\allowbreak{}nn, vec(\allowbreak{}m)) implies n in nn implies (\allowbreak{}dbi(\allowbreak{}m, xs))(\allowbreak{}n) == subspace-ms(\allowbreak{}rr-ms, ccint(\allowbreak{}-(\allowbreak{}vnrm(\allowbreak{}m))(\allowbreak{}xs(\allowbreak{}n)), (\allowbreak{}vnrm(\allowbreak{}m))(\allowbreak{}xs(\allowbreak{}n)))))
\end{quote}
also: \texttt{dbi-seq-compact} \texttt{dbi-product-facts} \\
mentioned by 8 result(s)

\smallskip

\noindent\textbf{\texttt{dbj(m, xs)}}\quad\textit{functoid}\index{dbj(m, xs)@\texttt{dbj(m, xs)}}\\
reads: the evaluation of the dual unit ball of \$1 along the sequence \$2 \\
definition:
\begin{quote}\ttfamily\footnotesize\raggedright\noindent
vnb-lambda(\allowbreak{}dbf\_, dual-unit-ball(\allowbreak{}m), vnb-lambda(\allowbreak{}dbn\_, nn, dbf\_(\allowbreak{}xs(\allowbreak{}dbn\_))))
\end{quote}
characterized by \texttt{dbj-type}  \textit{(proven)}:
\begin{quote}\ttfamily\footnotesize\raggedright\noindent
forall(\allowbreak{}[m, xs], is-normed-vector-space(\allowbreak{}m) implies xs in fun(\allowbreak{}nn, vec(\allowbreak{}m)) implies dbj(\allowbreak{}m, xs) in fun(\allowbreak{}dual-unit-ball(\allowbreak{}m), product-carrier(\allowbreak{}dbi(\allowbreak{}m, xs))))
\end{quote}
also: \texttt{dbj-value} \texttt{dbj-coord-seq} \texttt{dbj-injective} \texttt{dbj-compose-value} \texttt{dbj-diagonal-value} \texttt{dbj-coord-converges} \\
mentioned by 14 result(s)

\smallskip

\noindent\textbf{\texttt{dc-iter}}\quad\textit{defined-fn}\index{dc-iter@\texttt{dc-iter}}\\
defined by \texttt{dc-iter-zero} \texttt{dc-iter-succ} \\
characterized by \texttt{dc-iter-succ}  \textit{(definitional)}:
\begin{quote}\ttfamily\footnotesize\raggedright\noindent
forall(\allowbreak{}[s, a, nxt, n in nn], dc-iter(\allowbreak{}s, a, nxt, succ(\allowbreak{}n)) == choice(\allowbreak{}\{y\_ in s: y\_ in nxt(\allowbreak{}n, dc-iter(\allowbreak{}s, a, nxt, n))\}))
\end{quote}
also: \texttt{dc-iter-zero} \texttt{dc-iter-in} \texttt{dc-iter-step} \\
mentioned by 13 result(s)

\smallskip

\noindent\textbf{\texttt{def-app(mdt\_, mdhs\_, mdc\_)}}\quad\textit{predicate}\index{def-app(mdt_, mdhs_, mdc_)@\texttt{def-app(mdt\_, mdhs\_, mdc\_)}}\\
reads: the definedness rule application relates \$2 to \$3 \\
defined by \texttt{def-app} \\
characterized by \texttt{def-app}  \textit{(definitional)}:
\begin{quote}\ttfamily\footnotesize\raggedright\noindent
forall(\allowbreak{}[mdt\_, mdhs\_, mdc\_], def-app(\allowbreak{}mdt\_, mdhs\_, mdc\_) iff mdc\_ in seq-set and forsome(\allowbreak{}[mdg\_, mdf\_, mdd\_, mda\_], mdc\_ = [mdg\_, enc-eq(\allowbreak{}enc-apply(\allowbreak{}mdf\_, mda\_), enc-apply(\allowbreak{}mdf\_, mda\_))] and mdhs\_ = [[mdg\_, enc-in(\allowbreak{}mdf\_, enc-fun1(\allowbreak{}mdd\_))], [mdg\_, enc-in(\allowbreak{}mda\_, mdd\_)]]))
\end{quote}
also: \texttt{adm-def-app} \\
mentioned by 3 result(s)

\smallskip

\noindent\textbf{\texttt{def-binder(mdt\_, mdhs\_, mdc\_)}}\quad\textit{predicate}\index{def-binder(mdt_, mdhs_, mdc_)@\texttt{def-binder(mdt\_, mdhs\_, mdc\_)}}\\
reads: the definedness rule binder-domain relates \$2 to \$3 \\
defined by \texttt{def-binder} \\
characterized by \texttt{def-binder}  \textit{(definitional)}:
\begin{quote}\ttfamily\footnotesize\raggedright\noindent
forall(\allowbreak{}[mdt\_, mdhs\_, mdc\_], def-binder(\allowbreak{}mdt\_, mdhs\_, mdc\_) iff mdc\_ in seq-set and forsome(\allowbreak{}[mdg\_, mdh\_, mda\_, mdb\_], mdc\_ = [mdg\_, enc-eq(\allowbreak{}ebind(\allowbreak{}mdh\_, 1, [mda\_], mdb\_), ebind(\allowbreak{}mdh\_, 1, [mda\_], mdb\_))] and (\allowbreak{}mdh\_ = 22 or mdh\_ = 24 or mdh\_ = 26) and nbg-ok(\allowbreak{}mdb\_) and mdhs\_ = [[mdg\_, enc-eq(\allowbreak{}mda\_, mda\_)]]))
\end{quote}
also: \texttt{def-binder-admitted} \\
mentioned by 3 result(s)

\smallskip

\noindent\textbf{\texttt{def-cert-hyp(mdt\_, mdhs\_, mdc\_)}}\quad\textit{predicate}\index{def-cert-hyp(mdt_, mdhs_, mdc_)@\texttt{def-cert-hyp(mdt\_, mdhs\_, mdc\_)}}\\
reads: the definedness rule certifying-hypothesis relates \$2 to \$3 \\
defined by \texttt{def-cert-hyp} \\
characterized by \texttt{def-cert-hyp}  \textit{(definitional)}:
\begin{quote}\ttfamily\footnotesize\raggedright\noindent
forall(\allowbreak{}[mdt\_, mdhs\_, mdc\_], def-cert-hyp(\allowbreak{}mdt\_, mdhs\_, mdc\_) iff mdc\_ in seq-set and forsome(\allowbreak{}[mdg\_, mdf\_, mdw\_], mdc\_ = [mdg\_, enc-eq(\allowbreak{}mdw\_, mdw\_)] and mdf\_ in cert-set(\allowbreak{}mdt\_, mdg\_, mdw\_) and mdhs\_ = [[mdg\_, mdf\_]]))
\end{quote}
also: \texttt{adm-def-cert-hyp} \texttt{def-cert-hyp-derived} \\
mentioned by 5 result(s)

\smallskip

\noindent\textbf{\texttt{def-comp(mdt\_, mdhs\_, mdc\_)}}\quad\textit{predicate}\index{def-comp(mdt_, mdhs_, mdc_)@\texttt{def-comp(mdt\_, mdhs\_, mdc\_)}}\\
reads: the definedness rule comprehension relates \$2 to \$3 \\
defined by \texttt{def-comp} \\
characterized by \texttt{def-comp}  \textit{(definitional)}:
\begin{quote}\ttfamily\footnotesize\raggedright\noindent
forall(\allowbreak{}[mdt\_, mdhs\_, mdc\_], def-comp(\allowbreak{}mdt\_, mdhs\_, mdc\_) iff mdc\_ in seq-set and forsome(\allowbreak{}[mdg\_, mdp\_], mdc\_ = [mdg\_, enc-eq(\allowbreak{}enc-comp(\allowbreak{}mdp\_), enc-comp(\allowbreak{}mdp\_))] and nbg-ok(\allowbreak{}mdp\_) and mdhs\_ = []))
\end{quote}
also: \texttt{adm-def-comp} \texttt{def-comp-admitted} \\
mentioned by 5 result(s)

\smallskip

\noindent\textbf{\texttt{def-comp-closed(mit\_)}}\quad\textit{predicate}\index{def-comp-closed(mit_)@\texttt{def-comp-closed(mit\_)}}\\
reads: every closed predicative comprehension is derivably defined over \$1 \\
defined by \texttt{def-comp-closed} \\
characterized by \texttt{def-comp-closed}  \textit{(definitional)}:
\begin{quote}\ttfamily\footnotesize\raggedright\noindent
forall(\allowbreak{}[mit\_], def-comp-closed(\allowbreak{}mit\_) iff forall(\allowbreak{}[mig\_, miq\_], mig\_ in ctx-set implies enc-comp(\allowbreak{}miq\_) in clt implies nbg-ok(\allowbreak{}miq\_) implies [mig\_, enc-eq(\allowbreak{}enc-comp(\allowbreak{}miq\_), enc-comp(\allowbreak{}miq\_))] in derivable(\allowbreak{}mit\_)))
\end{quote}
also: \texttt{def-comp-closed-holds} \texttt{def-comp-closed-from-schema} \\
mentioned by 7 result(s)

\smallskip

\noindent\textbf{\texttt{def-fun-inst(mdt\_, mdhs\_, mdc\_)}}\quad\textit{predicate}\index{def-fun-inst(mdt_, mdhs_, mdc_)@\texttt{def-fun-inst(mdt\_, mdhs\_, mdc\_)}}\\
reads: the definedness rule functoid-instance relates \$2 to \$3 \\
defined by \texttt{def-fun-inst} \\
characterized by \texttt{def-fun-inst}  \textit{(definitional)}:
\begin{quote}\ttfamily\footnotesize\raggedright\noindent
forall(\allowbreak{}[mdt\_, mdhs\_, mdc\_], def-fun-inst(\allowbreak{}mdt\_, mdhs\_, mdc\_) iff mdc\_ in seq-set and forsome(\allowbreak{}[mdg\_, mdh\_, mdl\_, mdb\_], mdc\_ = [mdg\_, enc-eq(\allowbreak{}eapp(\allowbreak{}mdh\_, mdl\_), eapp(\allowbreak{}mdh\_, mdl\_))] and fun-def(\allowbreak{}mdh\_, length(\allowbreak{}mdl\_), mdb\_) and mdhs\_ = [[mdg\_, enc-eq(\allowbreak{}openl(\allowbreak{}mdb\_, mdl\_, length(\allowbreak{}mdl\_)), openl(\allowbreak{}mdb\_, mdl\_, length(\allowbreak{}mdl\_)))]]))
\end{quote}
also: \texttt{adm-def-fun-inst} \texttt{def-fun-inst-derived} \\
mentioned by 5 result(s)

\smallskip

\noindent\textbf{\texttt{def-fun-nth(mdt\_, mdhs\_, mdc\_)}}\quad\textit{predicate}\index{def-fun-nth(mdt_, mdhs_, mdc_)@\texttt{def-fun-nth(mdt\_, mdhs\_, mdc\_)}}\\
reads: the definedness rule nth-of-functoid relates \$2 to \$3 \\
defined by \texttt{def-fun-nth} \\
characterized by \texttt{def-fun-nth}  \textit{(definitional)}:
\begin{quote}\ttfamily\footnotesize\raggedright\noindent
forall(\allowbreak{}[mdt\_, mdhs\_, mdc\_], def-fun-nth(\allowbreak{}mdt\_, mdhs\_, mdc\_) iff mdc\_ in seq-set and forsome(\allowbreak{}[mdg\_, mdh\_, mdl\_, mdb\_, mdk\_], mdc\_ = [mdg\_, enc-eq(\allowbreak{}enc-nth(\allowbreak{}enc-num(\allowbreak{}mdk\_), eapp(\allowbreak{}mdh\_, mdl\_)), enc-nth(\allowbreak{}enc-num(\allowbreak{}mdk\_), eapp(\allowbreak{}mdh\_, mdl\_)))] and fun-def(\allowbreak{}mdh\_, length(\allowbreak{}mdl\_), mdb\_) and mdk\_ in nn and mdhs\_ = [[mdg\_, enc-eq(\allowbreak{}enc-nth(\allowbreak{}enc-num(\allowbreak{}mdk\_), openl(\allowbreak{}mdb\_, mdl\_, length(\allowbreak{}mdl\_))), enc-nth(\allowbreak{}enc-num(\allowbreak{}mdk\_), openl(\allowbreak{}mdb\_, mdl\_, length(\allowbreak{}mdl\_))))]]))
\end{quote}
also: \texttt{adm-def-fun-nth} \texttt{def-fun-nth-derived} \\
mentioned by 5 result(s)

\smallskip

\noindent\textbf{\texttt{def-fun-sep(mdt\_, mdhs\_, mdc\_)}}\quad\textit{predicate}\index{def-fun-sep(mdt_, mdhs_, mdc_)@\texttt{def-fun-sep(mdt\_, mdhs\_, mdc\_)}}\\
reads: the definedness rule functoid-membership relates \$2 to \$3 \\
defined by \texttt{def-fun-sep} \\
characterized by \texttt{def-fun-sep}  \textit{(definitional)}:
\begin{quote}\ttfamily\footnotesize\raggedright\noindent
forall(\allowbreak{}[mdt\_, mdhs\_, mdc\_], def-fun-sep(\allowbreak{}mdt\_, mdhs\_, mdc\_) iff mdc\_ in seq-set and forsome(\allowbreak{}[mdg\_, mdh\_, mdl\_, mdb\_, md11y\_], mdc\_ = [mdg\_, enc-in(\allowbreak{}md11y\_, openl(\allowbreak{}mdb\_, mdl\_, length(\allowbreak{}mdl\_)))] and fun-def(\allowbreak{}mdh\_, length(\allowbreak{}mdl\_), mdb\_) and mdhs\_ = [[mdg\_, enc-in(\allowbreak{}md11y\_, eapp(\allowbreak{}mdh\_, mdl\_))]]))
\end{quote}
also: \texttt{adm-def-fun-sep} \texttt{def-fun-sep-derived} \\
mentioned by 5 result(s)

\smallskip

\noindent\textbf{\texttt{def-fun-strict(mdt\_, mdhs\_, mdc\_)}}\quad\textit{predicate}\index{def-fun-strict(mdt_, mdhs_, mdc_)@\texttt{def-fun-strict(mdt\_, mdhs\_, mdc\_)}}\\
reads: the definedness rule functoid-strict-step relates \$2 to \$3 \\
defined by \texttt{def-fun-strict} \\
characterized by \texttt{def-fun-strict}  \textit{(definitional)}:
\begin{quote}\ttfamily\footnotesize\raggedright\noindent
forall(\allowbreak{}[mdt\_, mdhs\_, mdc\_], def-fun-strict(\allowbreak{}mdt\_, mdhs\_, mdc\_) iff mdc\_ in seq-set and forsome(\allowbreak{}[mdg\_, mdh\_, mdl\_, mdb\_], mdc\_ = [mdg\_, enc-eq(\allowbreak{}openl(\allowbreak{}mdb\_, mdl\_, length(\allowbreak{}mdl\_)), openl(\allowbreak{}mdb\_, mdl\_, length(\allowbreak{}mdl\_)))] and fun-def(\allowbreak{}mdh\_, length(\allowbreak{}mdl\_), mdb\_) and mdhs\_ = [[mdg\_, enc-eq(\allowbreak{}eapp(\allowbreak{}mdh\_, mdl\_), eapp(\allowbreak{}mdh\_, mdl\_))]]))
\end{quote}
also: \texttt{adm-def-fun-strict} \texttt{def-fun-strict-derived} \\
mentioned by 5 result(s)

\smallskip

\noindent\textbf{\texttt{def-in-closed(mit\_)}}\quad\textit{predicate}\index{def-in-closed(mit_)@\texttt{def-in-closed(mit\_)}}\\
reads: a derivable member is derivably defined over \$1 \\
defined by \texttt{def-in-closed} \\
characterized by \texttt{def-in-closed}  \textit{(definitional)}:
\begin{quote}\ttfamily\footnotesize\raggedright\noindent
forall(\allowbreak{}[mit\_], def-in-closed(\allowbreak{}mit\_) iff forall(\allowbreak{}[mig\_, miw\_, mia\_], [mig\_, enc-in(\allowbreak{}miw\_, mia\_)] in derivable(\allowbreak{}mit\_) implies [mig\_, enc-eq(\allowbreak{}miw\_, miw\_)] in derivable(\allowbreak{}mit\_)))
\end{quote}
also: \texttt{def-in-closed-holds} \texttt{def-in-closed-from-strict} \\
mentioned by 6 result(s)

\smallskip

\noindent\textbf{\texttt{def-list(mdt\_, mdhs\_, mdc\_)}}\quad\textit{predicate}\index{def-list(mdt_, mdhs_, mdc_)@\texttt{def-list(mdt\_, mdhs\_, mdc\_)}}\\
reads: the definedness rule list relates \$2 to \$3 \\
defined by \texttt{def-list} \\
characterized by \texttt{def-list}  \textit{(definitional)}:
\begin{quote}\ttfamily\footnotesize\raggedright\noindent
forall(\allowbreak{}[mdt\_, mdhs\_, mdc\_], def-list(\allowbreak{}mdt\_, mdhs\_, mdc\_) iff mdc\_ in seq-set and forsome(\allowbreak{}[mdg\_, mdl\_], mdc\_ = [mdg\_, enc-eq(\allowbreak{}enc-list(\allowbreak{}mdl\_), enc-list(\allowbreak{}mdl\_))] and length(\allowbreak{}mdhs\_) = length(\allowbreak{}mdl\_) and forall(\allowbreak{}[mdi\_ in nn], 1 <= mdi\_ implies mdi\_ <= length(\allowbreak{}mdl\_) implies nth(\allowbreak{}mdi\_, mdhs\_) = [mdg\_, enc-in(\allowbreak{}nth(\allowbreak{}mdi\_, mdl\_), enc-setc)])))
\end{quote}
also: \texttt{def-list-admitted} \\
mentioned by 3 result(s)

\smallskip

\noindent\textbf{\texttt{def-nth-list(mdt\_, mdhs\_, mdc\_)}}\quad\textit{predicate}\index{def-nth-list(mdt_, mdhs_, mdc_)@\texttt{def-nth-list(mdt\_, mdhs\_, mdc\_)}}\\
reads: the definedness rule nth-of-list relates \$2 to \$3 \\
defined by \texttt{def-nth-list} \\
characterized by \texttt{def-nth-list}  \textit{(definitional)}:
\begin{quote}\ttfamily\footnotesize\raggedright\noindent
forall(\allowbreak{}[mdt\_, mdhs\_, mdc\_], def-nth-list(\allowbreak{}mdt\_, mdhs\_, mdc\_) iff mdc\_ in seq-set and forsome(\allowbreak{}[mdg\_, mdl\_, mdk\_], mdc\_ = [mdg\_, enc-eq(\allowbreak{}enc-nth(\allowbreak{}enc-num(\allowbreak{}mdk\_), enc-list(\allowbreak{}mdl\_)), enc-nth(\allowbreak{}enc-num(\allowbreak{}mdk\_), enc-list(\allowbreak{}mdl\_)))] and mdk\_ in nn and 1 <= mdk\_ and mdk\_ <= length(\allowbreak{}mdl\_) and length(\allowbreak{}mdhs\_) = length(\allowbreak{}mdl\_) and forall(\allowbreak{}[mdi\_ in nn], 1 <= mdi\_ implies mdi\_ <= length(\allowbreak{}mdl\_) implies nth(\allowbreak{}mdi\_, mdhs\_) = [mdg\_, enc-in(\allowbreak{}nth(\allowbreak{}mdi\_, mdl\_), enc-setc)])))
\end{quote}
also: \texttt{adm-def-nth-list} \texttt{def-nth-list-derived} \\
mentioned by 5 result(s)

\smallskip

\noindent\textbf{\texttt{def-nth-lit(mdt\_, mdhs\_, mdc\_)}}\quad\textit{predicate}\index{def-nth-lit(mdt_, mdhs_, mdc_)@\texttt{def-nth-lit(mdt\_, mdhs\_, mdc\_)}}\\
reads: the kernel's definedness clause nth-of-literal relates \$2 to \$3 \\
defined by \texttt{def-nth-lit} \\
characterized by \texttt{def-nth-lit}  \textit{(definitional)}:
\begin{quote}\ttfamily\footnotesize\raggedright\noindent
forall(\allowbreak{}[mdt\_, mdhs\_, mdc\_], def-nth-lit(\allowbreak{}mdt\_, mdhs\_, mdc\_) iff mdc\_ in seq-set and forsome(\allowbreak{}[mdg\_, mdl\_, mdk\_], mdc\_ = [mdg\_, enc-eq(\allowbreak{}enc-nth(\allowbreak{}enc-num(\allowbreak{}mdk\_), enc-list(\allowbreak{}mdl\_)), enc-nth(\allowbreak{}enc-num(\allowbreak{}mdk\_), enc-list(\allowbreak{}mdl\_)))] and mdk\_ in nn and 1 <= mdk\_ and mdk\_ <= length(\allowbreak{}mdl\_) and mdhs\_ = [[mdg\_, enc-eq(\allowbreak{}nth(\allowbreak{}mdk\_, mdl\_), nth(\allowbreak{}mdk\_, mdl\_))]]))
\end{quote}
mentioned by 2 result(s)

\smallskip

\noindent\textbf{\texttt{def-pred-hyp(mdt\_, mdhs\_, mdc\_)}}\quad\textit{predicate}\index{def-pred-hyp(mdt_, mdhs_, mdc_)@\texttt{def-pred-hyp(mdt\_, mdhs\_, mdc\_)}}\\
reads: the definedness rule predicate-hypothesis relates \$2 to \$3 \\
defined by \texttt{def-pred-hyp} \\
characterized by \texttt{def-pred-hyp}  \textit{(definitional)}:
\begin{quote}\ttfamily\footnotesize\raggedright\noindent
forall(\allowbreak{}[mdt\_, mdhs\_, mdc\_], def-pred-hyp(\allowbreak{}mdt\_, mdhs\_, mdc\_) iff mdc\_ in seq-set and forsome(\allowbreak{}[mdg\_, mdk\_, mds\_, mdb\_, md10f\_, mdw\_], mdc\_ = [mdg\_, enc-eq(\allowbreak{}mdw\_, mdw\_)] and pred-code(\allowbreak{}mdk\_) and mdhs\_ = [[mdg\_, enc-pred1(\allowbreak{}econst(\allowbreak{}mdk\_), mds\_)], [mdg\_, enc-all(\allowbreak{}enc-iff(\allowbreak{}enc-pred1(\allowbreak{}econst(\allowbreak{}mdk\_), ebvar(\allowbreak{}0)), mdb\_))]] and md10f\_ in true-reach(\allowbreak{}expr-open(\allowbreak{}mdb\_, 0, mds\_)) and strict-atom(\allowbreak{}md10f\_, mdw\_)))
\end{quote}
also: \texttt{adm-def-pred-hyp} \texttt{def-pred-hyp-derived} \\
mentioned by 5 result(s)

\smallskip

\noindent\textbf{\texttt{def-pred-hyp-n(mdt\_, mdhs\_, mdc\_)}}\quad\textit{predicate}\index{def-pred-hyp-n(mdt_, mdhs_, mdc_)@\texttt{def-pred-hyp-n(mdt\_, mdhs\_, mdc\_)}}\\
reads: the definedness rule predicate-hypothesis (any arity) relates \$2 to \$3 \\
defined by \texttt{def-pred-hyp-n} \\
characterized by \texttt{def-pred-hyp-n}  \textit{(definitional)}:
\begin{quote}\ttfamily\footnotesize\raggedright\noindent
forall(\allowbreak{}[mdt\_, mdhs\_, mdc\_], def-pred-hyp-n(\allowbreak{}mdt\_, mdhs\_, mdc\_) iff mdc\_ in seq-set and forsome(\allowbreak{}[mdg\_, mdl\_, mdk\_, mdd\_, mdb\_, mdw\_], mdc\_ = [mdg\_, enc-eq(\allowbreak{}mdw\_, mdw\_)] and nth(\allowbreak{}1, mdl\_) = econst(\allowbreak{}mdk\_) and pred-code(\allowbreak{}mdk\_) and enc-iff(\allowbreak{}eapp(\allowbreak{}10, mdl\_), mdb\_) in inst-reach(\allowbreak{}mdd\_, mdl\_) and mdb\_ in cert-set(\allowbreak{}mdt\_, mdg\_, mdw\_) and mdhs\_ = [[mdg\_, eapp(\allowbreak{}10, mdl\_)], [mdg\_, mdd\_]]))
\end{quote}
also: \texttt{adm-def-pred-hyp-n} \texttt{def-pred-hyp-n-derived} \\
mentioned by 5 result(s)

\smallskip

\noindent\textbf{\texttt{def-pred-unfold(mdt\_, mdhs\_, mdc\_)}}\quad\textit{predicate}\index{def-pred-unfold(mdt_, mdhs_, mdc_)@\texttt{def-pred-unfold(mdt\_, mdhs\_, mdc\_)}}\\
reads: the definedness rule predicate-unfolding relates \$2 to \$3 \\
defined by \texttt{def-pred-unfold} \\
characterized by \texttt{def-pred-unfold}  \textit{(definitional)}:
\begin{quote}\ttfamily\footnotesize\raggedright\noindent
forall(\allowbreak{}[mdt\_, mdhs\_, mdc\_], def-pred-unfold(\allowbreak{}mdt\_, mdhs\_, mdc\_) iff mdc\_ in seq-set and forsome(\allowbreak{}[mdg\_, mdk\_, mds\_, mdb\_], mdc\_ = [mdg\_, expr-open(\allowbreak{}mdb\_, 0, mds\_)] and pred-code(\allowbreak{}mdk\_) and mdhs\_ = [[mdg\_, enc-pred1(\allowbreak{}econst(\allowbreak{}mdk\_), mds\_)], [mdg\_, enc-all(\allowbreak{}enc-iff(\allowbreak{}enc-pred1(\allowbreak{}econst(\allowbreak{}mdk\_), ebvar(\allowbreak{}0)), mdb\_))]]))
\end{quote}
also: \texttt{adm-def-pred-unfold} \texttt{def-pred-unfold-derived} \\
mentioned by 5 result(s)

\smallskip

\noindent\textbf{\texttt{def-pred-unfold-n(mdt\_, mdhs\_, mdc\_)}}\quad\textit{predicate}\index{def-pred-unfold-n(mdt_, mdhs_, mdc_)@\texttt{def-pred-unfold-n(mdt\_, mdhs\_, mdc\_)}}\\
reads: the definedness rule predicate-unfolding (any arity) relates \$2 to \$3 \\
defined by \texttt{def-pred-unfold-n} \\
characterized by \texttt{def-pred-unfold-n}  \textit{(definitional)}:
\begin{quote}\ttfamily\footnotesize\raggedright\noindent
forall(\allowbreak{}[mdt\_, mdhs\_, mdc\_], def-pred-unfold-n(\allowbreak{}mdt\_, mdhs\_, mdc\_) iff mdc\_ in seq-set and forsome(\allowbreak{}[mdg\_, mdl\_, mdk\_, mdd\_, mdb\_], mdc\_ = [mdg\_, mdb\_] and nth(\allowbreak{}1, mdl\_) = econst(\allowbreak{}mdk\_) and pred-code(\allowbreak{}mdk\_) and enc-iff(\allowbreak{}eapp(\allowbreak{}10, mdl\_), mdb\_) in inst-reach(\allowbreak{}mdd\_, mdl\_) and mdhs\_ = [[mdg\_, eapp(\allowbreak{}10, mdl\_)], [mdg\_, mdd\_]]))
\end{quote}
also: \texttt{adm-def-pred-unfold-n} \texttt{def-pred-unfold-n-derived} \\
mentioned by 5 result(s)

\smallskip

\noindent\textbf{\texttt{def-strict-arg(mdt\_, mdhs\_, mdc\_)}}\quad\textit{predicate}\index{def-strict-arg(mdt_, mdhs_, mdc_)@\texttt{def-strict-arg(mdt\_, mdhs\_, mdc\_)}}\\
reads: the definedness rule strict-argument relates \$2 to \$3 \\
defined by \texttt{def-strict-arg} \\
characterized by \texttt{def-strict-arg}  \textit{(definitional)}:
\begin{quote}\ttfamily\footnotesize\raggedright\noindent
forall(\allowbreak{}[mdt\_, mdhs\_, mdc\_], def-strict-arg(\allowbreak{}mdt\_, mdhs\_, mdc\_) iff mdc\_ in seq-set and forsome(\allowbreak{}[mdg\_, mdh\_, mdl\_, mdi\_], mdc\_ = [mdg\_, enc-eq(\allowbreak{}nth(\allowbreak{}mdi\_, mdl\_), nth(\allowbreak{}mdi\_, mdl\_))] and strict-head(\allowbreak{}mdh\_) and mdi\_ in nn and 1 <= mdi\_ and mdi\_ <= length(\allowbreak{}mdl\_) and mdhs\_ = [[mdg\_, enc-eq(\allowbreak{}eapp(\allowbreak{}mdh\_, mdl\_), eapp(\allowbreak{}mdh\_, mdl\_))]]))
\end{quote}
also: \texttt{adm-def-strict-arg} \texttt{def-strict-arg-admitted} \\
mentioned by 5 result(s)

\smallskip

\noindent\textbf{\texttt{def-strict-dom(mdt\_, mdhs\_, mdc\_)}}\quad\textit{predicate}\index{def-strict-dom(mdt_, mdhs_, mdc_)@\texttt{def-strict-dom(mdt\_, mdhs\_, mdc\_)}}\\
reads: the definedness rule strict-domain relates \$2 to \$3 \\
defined by \texttt{def-strict-dom} \\
characterized by \texttt{def-strict-dom}  \textit{(definitional)}:
\begin{quote}\ttfamily\footnotesize\raggedright\noindent
forall(\allowbreak{}[mdt\_, mdhs\_, mdc\_], def-strict-dom(\allowbreak{}mdt\_, mdhs\_, mdc\_) iff mdc\_ in seq-set and forsome(\allowbreak{}[mdg\_, mdh\_, mda\_, mdb\_], mdc\_ = [mdg\_, enc-eq(\allowbreak{}mda\_, mda\_)] and (\allowbreak{}mdh\_ = 22 or mdh\_ = 24 or mdh\_ = 26) and mdhs\_ = [[mdg\_, enc-eq(\allowbreak{}ebind(\allowbreak{}mdh\_, 1, [mda\_], mdb\_), ebind(\allowbreak{}mdh\_, 1, [mda\_], mdb\_))]]))
\end{quote}
also: \texttt{adm-def-strict-dom} \texttt{def-strict-dom-admitted} \\
mentioned by 5 result(s)

\smallskip

\noindent\textbf{\texttt{def-strict-eq-l(mdt\_, mdhs\_, mdc\_)}}\quad\textit{predicate}\index{def-strict-eq-l(mdt_, mdhs_, mdc_)@\texttt{def-strict-eq-l(mdt\_, mdhs\_, mdc\_)}}\\
reads: the definedness rule strict-eq-left relates \$2 to \$3 \\
defined by \texttt{def-strict-eq-l} \\
characterized by \texttt{def-strict-eq-l}  \textit{(definitional)}:
\begin{quote}\ttfamily\footnotesize\raggedright\noindent
forall(\allowbreak{}[mdt\_, mdhs\_, mdc\_], def-strict-eq-l(\allowbreak{}mdt\_, mdhs\_, mdc\_) iff mdc\_ in seq-set and forsome(\allowbreak{}[mdg\_, mds\_, mdu\_], mdc\_ = [mdg\_, enc-eq(\allowbreak{}mds\_, mds\_)] and mdhs\_ = [[mdg\_, enc-eq(\allowbreak{}mds\_, mdu\_)]]))
\end{quote}
also: \texttt{adm-def-strict-eq-l} \\
mentioned by 3 result(s)

\smallskip

\noindent\textbf{\texttt{def-strict-eq-r(mdt\_, mdhs\_, mdc\_)}}\quad\textit{predicate}\index{def-strict-eq-r(mdt_, mdhs_, mdc_)@\texttt{def-strict-eq-r(mdt\_, mdhs\_, mdc\_)}}\\
reads: the definedness rule strict-eq-right relates \$2 to \$3 \\
defined by \texttt{def-strict-eq-r} \\
characterized by \texttt{def-strict-eq-r}  \textit{(definitional)}:
\begin{quote}\ttfamily\footnotesize\raggedright\noindent
forall(\allowbreak{}[mdt\_, mdhs\_, mdc\_], def-strict-eq-r(\allowbreak{}mdt\_, mdhs\_, mdc\_) iff mdc\_ in seq-set and forsome(\allowbreak{}[mdg\_, mds\_, mdu\_], mdc\_ = [mdg\_, enc-eq(\allowbreak{}mdu\_, mdu\_)] and mdhs\_ = [[mdg\_, enc-eq(\allowbreak{}mds\_, mdu\_)]]))
\end{quote}
also: \texttt{adm-def-strict-eq-r} \\
mentioned by 3 result(s)

\smallskip

\noindent\textbf{\texttt{def-strict-hyp(mdt\_, mdhs\_, mdc\_)}}\quad\textit{predicate}\index{def-strict-hyp(mdt_, mdhs_, mdc_)@\texttt{def-strict-hyp(mdt\_, mdhs\_, mdc\_)}}\\
reads: the definedness rule strict-hypothesis relates \$2 to \$3 \\
defined by \texttt{def-strict-hyp} \\
characterized by \texttt{def-strict-hyp}  \textit{(definitional)}:
\begin{quote}\ttfamily\footnotesize\raggedright\noindent
forall(\allowbreak{}[mdt\_, mdhs\_, mdc\_], def-strict-hyp(\allowbreak{}mdt\_, mdhs\_, mdc\_) iff mdc\_ in seq-set and forsome(\allowbreak{}[mdg\_, mdf\_, md7h2\_, mdw\_], mdc\_ = [mdg\_, enc-eq(\allowbreak{}mdw\_, mdw\_)] and mdhs\_ = [[mdg\_, mdf\_]] and md7h2\_ in true-reach(\allowbreak{}mdf\_) and strict-atom(\allowbreak{}md7h2\_, mdw\_)))
\end{quote}
also: \texttt{adm-def-strict-hyp} \texttt{def-strict-hyp-derived} \\
mentioned by 5 result(s)

\smallskip

\noindent\textbf{\texttt{def-strict-in-l(mdt\_, mdhs\_, mdc\_)}}\quad\textit{predicate}\index{def-strict-in-l(mdt_, mdhs_, mdc_)@\texttt{def-strict-in-l(mdt\_, mdhs\_, mdc\_)}}\\
reads: the definedness rule strict-in-left relates \$2 to \$3 \\
defined by \texttt{def-strict-in-l} \\
characterized by \texttt{def-strict-in-l}  \textit{(definitional)}:
\begin{quote}\ttfamily\footnotesize\raggedright\noindent
forall(\allowbreak{}[mdt\_, mdhs\_, mdc\_], def-strict-in-l(\allowbreak{}mdt\_, mdhs\_, mdc\_) iff mdc\_ in seq-set and forsome(\allowbreak{}[mdg\_, mdw\_, mda\_], mdc\_ = [mdg\_, enc-eq(\allowbreak{}mdw\_, mdw\_)] and mdhs\_ = [[mdg\_, enc-in(\allowbreak{}mdw\_, mda\_)]]))
\end{quote}
also: \texttt{adm-def-strict-in-l} \texttt{def-strict-in-l-admitted} \\
mentioned by 5 result(s)

\smallskip

\noindent\textbf{\texttt{def-strict-in-r(mdt\_, mdhs\_, mdc\_)}}\quad\textit{predicate}\index{def-strict-in-r(mdt_, mdhs_, mdc_)@\texttt{def-strict-in-r(mdt\_, mdhs\_, mdc\_)}}\\
reads: the definedness rule strict-in-right relates \$2 to \$3 \\
defined by \texttt{def-strict-in-r} \\
characterized by \texttt{def-strict-in-r}  \textit{(definitional)}:
\begin{quote}\ttfamily\footnotesize\raggedright\noindent
forall(\allowbreak{}[mdt\_, mdhs\_, mdc\_], def-strict-in-r(\allowbreak{}mdt\_, mdhs\_, mdc\_) iff mdc\_ in seq-set and forsome(\allowbreak{}[mdg\_, mdw\_, mda\_], mdc\_ = [mdg\_, enc-eq(\allowbreak{}mda\_, mda\_)] and mdhs\_ = [[mdg\_, enc-in(\allowbreak{}mdw\_, mda\_)]]))
\end{quote}
also: \texttt{adm-def-strict-in-r} \texttt{def-strict-in-r-admitted} \\
mentioned by 5 result(s)

\smallskip

\noindent\textbf{\texttt{def-strict-pred(mdt\_, mdhs\_, mdc\_)}}\quad\textit{predicate}\index{def-strict-pred(mdt_, mdhs_, mdc_)@\texttt{def-strict-pred(mdt\_, mdhs\_, mdc\_)}}\\
reads: the definedness rule strict-predicate relates \$2 to \$3 \\
defined by \texttt{def-strict-pred} \\
characterized by \texttt{def-strict-pred}  \textit{(definitional)}:
\begin{quote}\ttfamily\footnotesize\raggedright\noindent
forall(\allowbreak{}[mdt\_, mdhs\_, mdc\_], def-strict-pred(\allowbreak{}mdt\_, mdhs\_, mdc\_) iff mdc\_ in seq-set and forsome(\allowbreak{}[mdg\_, mdl\_, mdk\_, mdi\_], mdc\_ = [mdg\_, enc-eq(\allowbreak{}nth(\allowbreak{}mdi\_, mdl\_), nth(\allowbreak{}mdi\_, mdl\_))] and nth(\allowbreak{}1, mdl\_) = econst(\allowbreak{}mdk\_) and pred-code(\allowbreak{}mdk\_) and mdi\_ in nn and 2 <= mdi\_ and mdi\_ <= length(\allowbreak{}mdl\_) and mdhs\_ = [[mdg\_, eapp(\allowbreak{}10, mdl\_)]]))
\end{quote}
also: \texttt{adm-def-strict-pred} \texttt{def-strict-pred-admitted} \\
mentioned by 5 result(s)

\smallskip

\noindent\textbf{\texttt{def-struct-hyp(mdt\_, mdhs\_, mdc\_)}}\quad\textit{predicate}\index{def-struct-hyp(mdt_, mdhs_, mdc_)@\texttt{def-struct-hyp(mdt\_, mdhs\_, mdc\_)}}\\
reads: the definedness rule structure-hypothesis relates \$2 to \$3 \\
defined by \texttt{def-struct-hyp} \\
characterized by \texttt{def-struct-hyp}  \textit{(definitional)}:
\begin{quote}\ttfamily\footnotesize\raggedright\noindent
forall(\allowbreak{}[mdt\_, mdhs\_, mdc\_], def-struct-hyp(\allowbreak{}mdt\_, mdhs\_, mdc\_) iff mdc\_ in seq-set and forsome(\allowbreak{}[mdg\_, mdk\_, mds\_], mdc\_ = [mdg\_, enc-eq(\allowbreak{}mds\_, mds\_)] and pred-code(\allowbreak{}mdk\_) and mdhs\_ = [[mdg\_, enc-pred1(\allowbreak{}econst(\allowbreak{}mdk\_), mds\_)]]))
\end{quote}
also: \texttt{adm-def-struct-hyp} \texttt{def-struct-hyp-derived} \\
mentioned by 5 result(s)

\smallskip

\noindent\textbf{\texttt{def-total(mdt\_, mdhs\_, mdc\_)}}\quad\textit{predicate}\index{def-total(mdt_, mdhs_, mdc_)@\texttt{def-total(mdt\_, mdhs\_, mdc\_)}}\\
reads: the definedness rule total-constructor relates \$2 to \$3 \\
defined by \texttt{def-total} \\
characterized by \texttt{def-total}  \textit{(definitional)}:
\begin{quote}\ttfamily\footnotesize\raggedright\noindent
forall(\allowbreak{}[mdt\_, mdhs\_, mdc\_], def-total(\allowbreak{}mdt\_, mdhs\_, mdc\_) iff mdc\_ in seq-set and forsome(\allowbreak{}[mdg\_, mdh\_, mdl\_], mdc\_ = [mdg\_, enc-eq(\allowbreak{}eapp(\allowbreak{}mdh\_, mdl\_), eapp(\allowbreak{}mdh\_, mdl\_))] and total-head(\allowbreak{}mdh\_) and length(\allowbreak{}mdhs\_) = length(\allowbreak{}mdl\_) and forall(\allowbreak{}[mdi\_ in nn], 1 <= mdi\_ implies mdi\_ <= length(\allowbreak{}mdl\_) implies nth(\allowbreak{}mdi\_, mdhs\_) = [mdg\_, enc-eq(\allowbreak{}nth(\allowbreak{}mdi\_, mdl\_), nth(\allowbreak{}mdi\_, mdl\_))])))
\end{quote}
also: \texttt{def-total-admitted} \\
mentioned by 3 result(s)

\smallskip

\noindent\textbf{\texttt{deg(a, p)}}\quad\textit{functoid}\index{deg(a, p)@\texttt{deg(a, p)}}\\
reads: the degree of \$2 \\
definition:
\begin{quote}\ttfamily\footnotesize\raggedright\noindent
iota(\allowbreak{}n\_, n\_ in nn and deg-bound(\allowbreak{}a, p, n\_) and forall(\allowbreak{}[j\_ in nn], deg-bound(\allowbreak{}a, p, j\_) implies n\_ <= j\_))
\end{quote}
characterized by \texttt{deg-drop-uba-as-ring}  \textit{(definitional)}:
\begin{quote}\ttfamily\footnotesize\raggedright\noindent
forall(\allowbreak{}[r], is-unital-banach-algebra(\allowbreak{}r) implies forall(\allowbreak{}[h\_ in carr(\allowbreak{}poly(\allowbreak{}uba-as-ring(\allowbreak{}r))), n\_ in nn], deg-bound(\allowbreak{}uba-as-ring(\allowbreak{}r), h\_, n\_) implies h\_(\allowbreak{}n\_) = vzero(\allowbreak{}r) implies h\_ = zero(\allowbreak{}poly(\allowbreak{}uba-as-ring(\allowbreak{}r))) or deg(\allowbreak{}uba-as-ring(\allowbreak{}r), h\_) < n\_))
\end{quote}
also: \texttt{monic-poly-deg-uba-as-ring} \texttt{deg-drop-uba-as-ring-uiba-as-uba} \texttt{monic-poly-deg-uba-as-ring-uiba-as-uba} \texttt{deg-is} \texttt{deg-drop} \texttt{deg-add-le} \texttt{deg-unfold} \texttt{monomial-deg} ... (14 in all) \\
mentioned by 48 result(s)

\smallskip

\noindent\textbf{\texttt{deg-bound(a, p, n)}}\quad\textit{predicate}\index{deg-bound(a, p, n)@\texttt{deg-bound(a, p, n)}}\\
reads: \$3 bounds the support of \$2 \\
defined by \texttt{deg-bound} \\
characterized by \texttt{deg-bound}  \textit{(definitional)}:
\begin{quote}\ttfamily\footnotesize\raggedright\noindent
forall(\allowbreak{}[a, p, n], deg-bound(\allowbreak{}a, p, n) iff forall(\allowbreak{}[k\_ in nn], n < k\_ implies p(\allowbreak{}k\_) = zero(\allowbreak{}a)))
\end{quote}
also: \texttt{monomial-deg-bound-uba-as-ring} \texttt{poly-mul-deg-bound-uba-as-ring} \texttt{monomial-deg-bound-uba-as-ring-uiba-as-uba} \texttt{poly-mul-deg-bound-uba-as-ring-uiba-as-uba} \texttt{deg-bound-apply} \texttt{monomial-deg-bound} \texttt{poly-add-deg-bound} \texttt{poly-mul-deg-bound} ... (9 in all) \\
mentioned by 52 result(s)

\smallskip

\noindent\textbf{\texttt{delete-at}}\quad\textit{operator}\index{delete-at@\texttt{delete-at}}\\
a term-forming head

\smallskip

\noindent\textbf{\texttt{deriv(f, a)}}\quad\textit{functoid}\index{deriv(f, a)@\texttt{deriv(f, a)}}\\
definition:
\begin{quote}\ttfamily\footnotesize\raggedright\noindent
iota(\allowbreak{}l, is-diff-at(\allowbreak{}f, a, l))
\end{quote}
characterized by \texttt{nth-deriv-succ}  \textit{(definitional)}:
\begin{quote}\ttfamily\footnotesize\raggedright\noindent
forall(\allowbreak{}[f, n in nn], nth-deriv(\allowbreak{}f, succ(\allowbreak{}n)) == vnb-lambda(\allowbreak{}x, rr, deriv(\allowbreak{}nth-deriv(\allowbreak{}f, n), x)))
\end{quote}
also: \texttt{nth-deriv-one} \texttt{poly-deriv-fn} \texttt{deriv-chain-value} \texttt{poly-deriv-fn-zero} \texttt{deriv-of-is-diff-at} \texttt{nth-deriv-succ-inner} \\
mentioned by 14 result(s)

\smallskip

\noindent\textbf{\texttt{deriv-on(k, u, f, a)}}\quad\textit{functoid}\index{deriv-on(k, u, f, a)@\texttt{deriv-on(k, u, f, a)}}\\
reads: the derivative of \$3 on \$2 at \$4 over \$1 \\
definition:
\begin{quote}\ttfamily\footnotesize\raggedright\noindent
iota(\allowbreak{}l, is-diff-on(\allowbreak{}k, u, f, a, l))
\end{quote}
characterized by \texttt{nth-deriv-on-succ}  \textit{(definitional)}:
\begin{quote}\ttfamily\footnotesize\raggedright\noindent
forall(\allowbreak{}[ndvk\_, ndvu\_, ndvf\_, ndvn\_ in nn], nth-deriv-on(\allowbreak{}ndvk\_, ndvu\_, ndvf\_, succ(\allowbreak{}ndvn\_)) == vnb-lambda(\allowbreak{}ndvz\_, ndvu\_, deriv-on(\allowbreak{}ndvk\_, ndvu\_, nth-deriv-on(\allowbreak{}ndvk\_, ndvu\_, ndvf\_, ndvn\_), ndvz\_)))
\end{quote}
also: \texttt{deriv-on-unfold} \texttt{deriv-on-of-is-diff-on} \texttt{nth-deriv-on-succ-apply} \texttt{cc-deriv-on-of-is-diff-on} \\
mentioned by 52 result(s)

\smallskip

\noindent\textbf{\texttt{deriv-v(m, f, a)}}\quad\textit{functoid}\index{deriv-v(m, f, a)@\texttt{deriv-v(m, f, a)}}\\
definition:
\begin{quote}\ttfamily\footnotesize\raggedright\noindent
iota(\allowbreak{}l, is-diff-at-v(\allowbreak{}m, f, a, l))
\end{quote}
characterized by \texttt{nth-deriv-v-succ}  \textit{(definitional)}:
\begin{quote}\ttfamily\footnotesize\raggedright\noindent
forall(\allowbreak{}[m, f, n in nn], nth-deriv-v(\allowbreak{}m, f, succ(\allowbreak{}n)) == vnb-lambda(\allowbreak{}x, rr, deriv-v(\allowbreak{}m, nth-deriv-v(\allowbreak{}m, f, n), x)))
\end{quote}
mentioned by 2 result(s)

\smallskip

\noindent\textbf{\texttt{derivable(mct\_)}}\quad\textit{functoid}\index{derivable(mct_)@\texttt{derivable(mct\_)}}\\
reads: the sequents derivable in the core over \$1 \\
definition:
\begin{quote}\ttfamily\footnotesize\raggedright\noindent
\{mcq\_ in seq-set: forall(\allowbreak{}[mcs\_ in power(\allowbreak{}seq-set)], core-closed(\allowbreak{}mct\_, mcs\_) implies mcq\_ in mcs\_)\}
\end{quote}
characterized by \texttt{derivable-cut}  \textit{(proven)}:
\begin{quote}\ttfamily\footnotesize\raggedright\noindent
forall(\allowbreak{}[tt, hs, c], cr-cut(\allowbreak{}tt, hs, c) implies hyps-in(\allowbreak{}hs, derivable(\allowbreak{}tt)) implies c in seq-set implies c in derivable(\allowbreak{}tt))
\end{quote}
also: \texttt{derivable-cut-c} \texttt{derivable-closed} \texttt{derivable-nth-eq} \texttt{derivable-unfold} \texttt{derivable-if-true} \texttt{derivable-iota-cr} \texttt{derivable-or-elim} \texttt{derivable-sep-mem} ... (68 in all) \\
mentioned by 272 result(s)

\smallskip

\noindent\textbf{\texttt{descend(s, f)}}\quad\textit{functoid}\index{descend(s, f)@\texttt{descend(s, f)}}\\
definition:
\begin{quote}\ttfamily\footnotesize\raggedright\noindent
vnb-lambda(\allowbreak{}c, quotient(\allowbreak{}s), iota(\allowbreak{}z, z in set and forsome(\allowbreak{}[a in c], z = f(\allowbreak{}a))))
\end{quote}
characterized by \texttt{descend-in-fun}  \textit{(proven)}:
\begin{quote}\ttfamily\footnotesize\raggedright\noindent
forall(\allowbreak{}[s], is-setoid(\allowbreak{}s) implies forall(\allowbreak{}[z, f], f in fun(\allowbreak{}pts(\allowbreak{}s), z) and respects(\allowbreak{}s, f) implies descend(\allowbreak{}s, f) in fun(\allowbreak{}quotient(\allowbreak{}s), z)))
\end{quote}
also: \texttt{descend-computes} \\
mentioned by 6 result(s)

\smallskip

\noindent\textbf{\texttt{descend2(s, f)}}\quad\textit{functoid}\index{descend2(s, f)@\texttt{descend2(s, f)}}\\
definition:
\begin{quote}\ttfamily\footnotesize\raggedright\noindent
vnb-lambda(\allowbreak{}[c, d], cartesian(\allowbreak{}quotient(\allowbreak{}s), quotient(\allowbreak{}s)), iota(\allowbreak{}z, z in set and forsome(\allowbreak{}[a in c, b in d], z = f(\allowbreak{}a, b))))
\end{quote}
characterized by \texttt{descend2-in-fun}  \textit{(proven)}:
\begin{quote}\ttfamily\footnotesize\raggedright\noindent
forall(\allowbreak{}[s], is-setoid(\allowbreak{}s) implies forall(\allowbreak{}[z, f in fun(\allowbreak{}cartesian(\allowbreak{}pts(\allowbreak{}s), pts(\allowbreak{}s)), z)], respects2(\allowbreak{}s, f) implies descend2(\allowbreak{}s, f) in fun(\allowbreak{}cartesian(\allowbreak{}quotient(\allowbreak{}s), quotient(\allowbreak{}s)), z)))
\end{quote}
also: \texttt{descend2-computes} \\
mentioned by 7 result(s)

\smallskip

\noindent\textbf{\texttt{det}}\quad\textit{primitive}\index{det@\texttt{det}}\\
reads: det(\$3) \\
defined by \texttt{det-zero} \texttt{det-cofactor} \\
characterized by \texttt{det-zero}  \textit{(definitional)}:
\begin{quote}\ttfamily\footnotesize\raggedright\noindent
forall(\allowbreak{}[r, a], det(\allowbreak{}r, 0, a) == one(\allowbreak{}r))
\end{quote}
also: \texttt{det-cofactor} \texttt{det-identity-ind-uba-as-ring} \texttt{det-in-carrier-ind-uba-as-ring} \texttt{norm-data-det-unique-cna-as-module} \texttt{norm-data-det-unique-sip-as-module} \texttt{norm-data-det-unique-uba-as-module} \texttt{norm-data-det-unique-cnvs-as-module} \texttt{change-of-basis-det-unit-cna-as-module} ... (95 in all) \\
mentioned by 341 result(s)

\smallskip

\noindent\textbf{\texttt{detach!}}\quad\textit{tactic}\index{detach"!@\texttt{detach"!}}\\
\texttt{(detach! impl)} -- Forward modus ponens: from an in-context (IMPLIES A B) whose A is also in context, leave B in context. \\
declared in \texttt{tactics-help.scm}

\smallskip

\noindent\textbf{\texttt{di}}\quad\textit{tactic}\index{di@\texttt{di}}\\
\texttt{(di)} -- Direct inference: split an AND goal, move an IMPLIES antecedent into the assumptions, or introduce a fresh eigenvariable for a leading FORALL. \\
declared in \texttt{tactics-help.scm}

\smallskip

\noindent\textbf{\texttt{diag-mat(n, f)}}\quad\textit{functoid}\index{diag-mat(n, f)@\texttt{diag-mat(n, f)}}\\
reads: the diagonal matrix with diagonal \$2(1), ..., \$2(\$1) \\
definition:
\begin{quote}\ttfamily\footnotesize\raggedright\noindent
matof(\allowbreak{}n, n, vnb-lambda(\allowbreak{}[dga\_, dgb\_], cartesian(\allowbreak{}interval(\allowbreak{}1, n), interval(\allowbreak{}1, n)), (\allowbreak{}mul(\allowbreak{}zz-ring))(\allowbreak{}f(\allowbreak{}dga\_), entry(\allowbreak{}identmat(\allowbreak{}zz-ring, n), dga\_, dgb\_))))
\end{quote}
characterized by \texttt{diag-mat-diag}  \textit{(proven)}:
\begin{quote}\ttfamily\footnotesize\raggedright\noindent
forall(\allowbreak{}[n, f, i], n in nn implies f in fun(\allowbreak{}interval(\allowbreak{}1, n), zz) implies i in interval(\allowbreak{}1, n) implies entry(\allowbreak{}diag-mat(\allowbreak{}n, f), i, i) = f(\allowbreak{}i))
\end{quote}
also: \texttt{diag-mat-type} \texttt{diag-mat-entry} \texttt{diag-mat-diagonal} \\
mentioned by 11 result(s)

\smallskip

\noindent\textbf{\texttt{dial}}\quad\textit{tactic}\index{dial@\texttt{dial}}\\
\texttt{(dial notch)} -- The PRESENTATION RESOLUTION DIAL: how much of the sequent to show.  r0 kernel S-expression, r1 surface syntax (the default, and the last INVERTIBLE notch -- what is printed re-parses to what was printed), r2 guards and typing hypotheses elided, r3 also structures destructured to their declared slot names and deep subterms elided.  (dial) reports the notch, (dial n) sets it, (dial 'auto) reads the current goal and suggests one, (dial 'why) shows the score behind the suggestion.  For ONE FORMULA rather than the proof state, (dial-wff x [n]) -- x a wff, a "surface string" or an S-expression, n defaulting to the dial's current notch: (dial-wff (lookup-theorem 'matmul-assoc) 3). \\
declared in \texttt{tactics-help.scm}

\smallskip

\noindent\textbf{\texttt{dial-wff}}\quad\textit{tactic}\index{dial-wff@\texttt{dial-wff}}\\
\texttt{(dial-wff x notch)} -- Present ONE FORMULA at a resolution notch, instead of the current proof state: x is a wff, a "surface string" or an S-expression, and notch defaults to the dial's current setting.  (dial-wff (lookup-theorem 'matmul-assoc) 3).  Same rungs, same aperture line, no proof required. \\
declared in \texttt{tactics-help.scm}

\smallskip

\noindent\textbf{\texttt{diam(mhm\_, mhs\_)}}\quad\textit{functoid}\index{diam(mhm_, mhs_)@\texttt{diam(mhm\_, mhs\_)}}\\
reads: the diameter of \$2 in \$1 \\
definition:
\begin{quote}\ttfamily\footnotesize\raggedright\noindent
esup(\allowbreak{}\{mhw\_ in rr-pos-star: forsome(\allowbreak{}[mhx\_ in mhs\_, mhy\_ in mhs\_], mhw\_ = (\allowbreak{}dist(\allowbreak{}mhm\_))(\allowbreak{}mhx\_, mhy\_))\})
\end{quote}
characterized by \texttt{diam-in}  \textit{(proven)}:
\begin{quote}\ttfamily\footnotesize\raggedright\noindent
forall(\allowbreak{}[m, a], diam(\allowbreak{}m, a) in rr-pos-star)
\end{quote}
also: \texttt{diam-mono} \texttt{diam-empty} \texttt{diam-unfold} \texttt{diam-compare} \texttt{diam-dist-le} \texttt{subspace-diam} \texttt{diam-le-of-dist} \\
mentioned by 25 result(s)

\smallskip

\noindent\textbf{\texttt{difference(u\_, b\_)}}\quad\textit{functoid}\index{difference(u_, b_)@\texttt{difference(u\_, b\_)}}\\
reads: the set difference of \$1 and \$2 \\
definition:
\begin{quote}\ttfamily\footnotesize\raggedright\noindent
complement-in(\allowbreak{}u\_, b\_)
\end{quote}
characterized by \texttt{difference-set}  \textit{(proven)}:
\begin{quote}\ttfamily\footnotesize\raggedright\noindent
forall(\allowbreak{}[x in set, b], difference(\allowbreak{}x, b) in set)
\end{quote}
also: \texttt{difference-subset} \texttt{difference-unfold} \texttt{difference-in-power} \texttt{translate-difference} \texttt{difference-empty-left} \texttt{difference-membership} \texttt{insert-difference-cancel} \texttt{insert-difference-outside} ... (9 in all) \\
mentioned by 221 result(s)

\smallskip

\noindent\textbf{\texttt{dir-deriv(m, f, a, eta)}}\quad\textit{functoid}\index{dir-deriv(m, f, a, eta)@\texttt{dir-deriv(m, f, a, eta)}}\\
definition:
\begin{quote}\ttfamily\footnotesize\raggedright\noindent
deriv(\allowbreak{}seg-curve(\allowbreak{}m, f, a, eta), 0)
\end{quote}
characterized by \texttt{dir-deriv-increment-cnvs-as-nvs}  \textit{(definitional)}:
\begin{quote}\ttfamily\footnotesize\raggedright\noindent
forall(\allowbreak{}[r], is-complex-normed-vector-space(\allowbreak{}r) implies forall(\allowbreak{}[f in fun(\allowbreak{}vec(\allowbreak{}r), rr), b\_ in vec(\allowbreak{}r), zeta in vec(\allowbreak{}r)], forall(\allowbreak{}[x\_ in vec(\allowbreak{}r)], forsome(\allowbreak{}[k\_], is-dir-diff-at(\allowbreak{}cnvs-as-nvs(\allowbreak{}r), f, x\_, zeta, k\_))) implies forsome(\allowbreak{}[th in rr], 0 < th and th < 1 and f(\allowbreak{}(\allowbreak{}vadd(\allowbreak{}r))(\allowbreak{}b\_, zeta)) - f(\allowbreak{}b\_) = dir-deriv(\allowbreak{}cnvs-as-nvs(\allowbreak{}r), f, (\allowbreak{}vadd(\allowbreak{}r))(\allowbreak{}b\_, (\allowbreak{}act(\allowbreak{}r))(\allowbreak{}th, zeta)), zeta))))
\end{quote}
also: \texttt{dir-deriv-corner-split-cnvs-as-nvs} \texttt{dir-deriv-increment-scaled-cnvs-as-nvs} \texttt{dir-deriv-value} \texttt{dir-deriv-linear} \texttt{dir-deriv-additive} \texttt{dir-deriv-increment} \texttt{dir-deriv-corner-split} \texttt{dir-deriv-increment-scaled} \\
mentioned by 12 result(s)

\smallskip

\noindent\textbf{\texttt{dir-of(ssd\_)}}\quad\textit{functoid}\index{dir-of(ssd_)@\texttt{dir-of(ssd\_)}}\\
reads: the direction of the directed message \$1 \\
definition:
\begin{quote}\ttfamily\footnotesize\raggedright\noindent
nth(\allowbreak{}1, ssd\_)
\end{quote}

\smallskip

\noindent\textbf{\texttt{direct-sum(s, n, f)}}\quad\textit{functoid}\index{direct-sum(s, n, f)@\texttt{direct-sum(s, n, f)}}\\
reads: the direct sum of the \$2 modules \$3 over \$1 \\
definition:
\begin{quote}\ttfamily\footnotesize\raggedright\noindent
[s, dsum-vec(\allowbreak{}n, f), dsum-vadd(\allowbreak{}n, f), dsum-vzero(\allowbreak{}n, f), dsum-vneg(\allowbreak{}n, f), dsum-act(\allowbreak{}s, n, f)]
\end{quote}
characterized by \texttt{direct-sum-is-module}  \textit{(proven)}:
\begin{quote}\ttfamily\footnotesize\raggedright\noindent
forall(\allowbreak{}[s, n, f], is-ring(\allowbreak{}s) implies n in nn implies forall(\allowbreak{}[i\_ in interval(\allowbreak{}1, n)], is-module(\allowbreak{}f(\allowbreak{}i\_))) implies forall(\allowbreak{}[i\_ in interval(\allowbreak{}1, n)], scal(\allowbreak{}f(\allowbreak{}i\_)) = s) implies is-module(\allowbreak{}direct-sum(\allowbreak{}s, n, f)))
\end{quote}
mentioned by 76 result(s)

\smallskip

\noindent\textbf{\texttt{directed(sst\_, ssk\_)}}\quad\textit{functoid}\index{directed(sst_, ssk_)@\texttt{directed(sst\_, ssk\_)}}\\
reads: the directed messages over the tokens \$1 and the keys \$2 \\
definition:
\begin{quote}\ttfamily\footnotesize\raggedright\noindent
cartesian(\allowbreak{}pair(\allowbreak{}sent, received), msg-set(\allowbreak{}sst\_, ssk\_))
\end{quote}
characterized by \texttt{directed-intro}  \textit{(proven)}:
\begin{quote}\ttfamily\footnotesize\raggedright\noindent
forall(\allowbreak{}[t, k, g in msg-set(\allowbreak{}t, k)], [sent, g] in directed(\allowbreak{}t, k) and [received, g] in directed(\allowbreak{}t, k))
\end{quote}
also: \texttt{directed-parts} \texttt{directed-unfold} \texttt{node-step-in-directed} \\
mentioned by 8 result(s)

\smallskip

\noindent\textbf{\texttt{directed-set}}\quad\textit{structure}\index{directed-set@\texttt{directed-set}}\\
predicate \texttt{is-directed-set}; slots: pts dleq

\smallskip

\noindent\textbf{\texttt{discrete-ms(dsx\_)}}\quad\textit{functoid}\index{discrete-ms(dsx_)@\texttt{discrete-ms(dsx\_)}}\\
reads: \$1 with the discrete metric \\
definition:
\begin{quote}\ttfamily\footnotesize\raggedright\noindent
[dsx\_, vnb-lambda(\allowbreak{}[dsu\_, dsv\_], cartesian(\allowbreak{}dsx\_, dsx\_), if(\allowbreak{}dsu\_ = dsv\_, 0, 1))]
\end{quote}
characterized by \texttt{discrete-ms-pts}  \textit{(proven)}:
\begin{quote}\ttfamily\footnotesize\raggedright\noindent
forall(\allowbreak{}[x in set], pts(\allowbreak{}discrete-ms(\allowbreak{}x)) == x)
\end{quote}
also: \texttt{discrete-ms-dist-fun} \texttt{discrete-ms-dist-bounds} \texttt{discrete-ms-dist-values} \texttt{discrete-ms-is-metric-space} \\
mentioned by 11 result(s)

\smallskip

\noindent\textbf{\texttt{discrete-space(a\_)}}\quad\textit{functoid}\index{discrete-space(a_)@\texttt{discrete-space(a\_)}}\\
reads: the discrete space on \$1 \\
definition:
\begin{quote}\ttfamily\footnotesize\raggedright\noindent
[a\_, power(\allowbreak{}a\_)]
\end{quote}
characterized by \texttt{discrete-space-is-top-space}  \textit{(proven)}:
\begin{quote}\ttfamily\footnotesize\raggedright\noindent
forall(\allowbreak{}[a\_ in set], is-top-space(\allowbreak{}discrete-space(\allowbreak{}a\_)))
\end{quote}
also: \texttt{discrete-space-is-measurable-space} \\
mentioned by 8 result(s)

\smallskip

\noindent\textbf{\texttt{disjoint-sum-set(ssx\_)}}\quad\textit{functoid}\index{disjoint-sum-set(ssx_)@\texttt{disjoint-sum-set(ssx\_)}}\\
reads: the disjoint union of the sets \$1 \\
definition:
\begin{quote}\ttfamily\footnotesize\raggedright\noindent
\{ssp\_ in cartesian(\allowbreak{}nn, big-union(\allowbreak{}ssk\_, nn, ssx\_(\allowbreak{}ssk\_))): nth(\allowbreak{}2, ssp\_) in ssx\_(\allowbreak{}nth(\allowbreak{}1, ssp\_))\}
\end{quote}
mentioned by 9 result(s)

\smallskip

\noindent\textbf{\texttt{dist(s)}}\quad\textit{accessor}\index{dist(s)@\texttt{dist(s)}}\\
reads: the distance function of \$1 \\
a structure slot \\
characterized by \texttt{cc-ms@dist}  \textit{(definitional)}:
\begin{quote}\ttfamily\footnotesize\raggedright\noindent
dist(\allowbreak{}cc-ms) == vnb-lambda(\allowbreak{}[x, y], cartesian(\allowbreak{}cc, cc), magnitude(\allowbreak{}x - y))
\end{quote}
also: \texttt{rr-ms@dist} \texttt{nf-metric-distance} \texttt{nvs-dist-shift-cnvs-as-nvs} \texttt{nvs-metric-distance-cnvs-as-nvs} \texttt{rb2-dist-union-cna-as-normed-ag} \texttt{rb2-dist-union-uba-as-normed-ag} \texttt{rb2-dist-subset-cna-as-normed-ag} \texttt{rb2-dist-subset-uba-as-normed-ag} ... (170 in all) \\
mentioned by 587 result(s)

\smallskip

\noindent\textbf{\texttt{dist-seq(m, f, g)}}\quad\textit{functoid}\index{dist-seq(m, f, g)@\texttt{dist-seq(m, f, g)}}\\
definition:
\begin{quote}\ttfamily\footnotesize\raggedright\noindent
vnb-lambda(\allowbreak{}n, nn, (\allowbreak{}dist(\allowbreak{}m))(\allowbreak{}f(\allowbreak{}n), g(\allowbreak{}n)))
\end{quote}
characterized by \texttt{dist-seq-limit}  \textit{(proven)}:
\begin{quote}\ttfamily\footnotesize\raggedright\noindent
forall(\allowbreak{}[md, fq, f, g], is-module(\allowbreak{}md) implies scal(\allowbreak{}md) = normed-field-as-commutative-ring(\allowbreak{}cc-normed-field) implies is-inner-product-form(\allowbreak{}md, fq) implies f in form-cseq(\allowbreak{}md, fq) implies g in form-cseq(\allowbreak{}md, fq) implies converges-to(\allowbreak{}rr-ms, dist-seq(\allowbreak{}form-ms(\allowbreak{}md, fq), f, g), sqrt(\allowbreak{}(\allowbreak{}limit-form(\allowbreak{}md, fq))(\allowbreak{}seq-vadd(\allowbreak{}md, f, seq-vneg(\allowbreak{}md, g)), seq-vadd(\allowbreak{}md, f, seq-vneg(\allowbreak{}md, g))))))
\end{quote}
also: \texttt{r7q-dist-seq-sym} \texttt{rko2-dist-seq-at} \texttt{rko2-dist-seq-type} \texttt{r7q-dist-seq-cauchy} \texttt{rko2-dist-seq-unfold} \texttt{r7q-dist-seq-converges} \\
mentioned by 26 result(s)

\smallskip

\noindent\textbf{\texttt{dists(s)}}\quad\textit{functoid}\index{dists(s)@\texttt{dists(s)}}\\
reads: the pseudometric family of \$1 \\
characterized by \texttt{frm-c-space-dists}  \textit{(proven)}:
\begin{quote}\ttfamily\footnotesize\raggedright\noindent
forall(\allowbreak{}[m, fam], vec(\allowbreak{}m) in set implies dists(\allowbreak{}fr-c-space(\allowbreak{}m, fam)) == vnb-lambda(\allowbreak{}frk\_, nn, vnb-lambda(\allowbreak{}[fru\_, frv\_], cartesian(\allowbreak{}vec(\allowbreak{}m), vec(\allowbreak{}m)), (\allowbreak{}fam(\allowbreak{}frk\_))(\allowbreak{}(\allowbreak{}vadd(\allowbreak{}m))(\allowbreak{}fru\_, (\allowbreak{}vneg(\allowbreak{}m))(\allowbreak{}frv\_))))))
\end{quote}
also: \texttt{frm-c-space-dists-apply} \\
mentioned by 38 result(s)

\smallskip

\noindent\textbf{\texttt{divides(gds\_, gda\_, gdb\_)}}\quad\textit{predicate}\index{divides(gds_, gda_, gdb_)@\texttt{divides(gds\_, gda\_, gdb\_)}}\\
reads: \$2 divides \$3 in \$1 \\
defined by \texttt{divides} \\
characterized by \texttt{divides}  \textit{(definitional)}:
\begin{quote}\ttfamily\footnotesize\raggedright\noindent
forall(\allowbreak{}[gds\_, gda\_, gdb\_], divides(\allowbreak{}gds\_, gda\_, gdb\_) iff gda\_ in carr(\allowbreak{}gds\_) and gdb\_ in carr(\allowbreak{}gds\_) and forsome(\allowbreak{}[gdk\_ in carr(\allowbreak{}gds\_)], gdb\_ = (\allowbreak{}mul(\allowbreak{}gds\_))(\allowbreak{}gdk\_, gda\_)))
\end{quote}
also: \texttt{zz-ring-divides-iff} \texttt{mutual-divides-associated} \texttt{prime-divides-prime-product} \texttt{poly-root-iff-x-minus-divides} \\
mentioned by 16 result(s)

\smallskip

\noindent\textbf{\texttt{dleq(s)}}\quad\textit{accessor}\index{dleq(s)@\texttt{dleq(s)}}\\
reads: the dleq slot of \$1 \\
a structure slot \\
characterized by \texttt{nn-directed-dleq}  \textit{(proven)}:
\begin{quote}\ttfamily\footnotesize\raggedright\noindent
dleq(\allowbreak{}nn-directed) == \{ndp\_ in cartesian(\allowbreak{}nn, nn): nth(\allowbreak{}1, ndp\_) <= nth(\allowbreak{}2, ndp\_)\}
\end{quote}
mentioned by 20 result(s)

\smallskip

\noindent\textbf{\texttt{dom}}\quad\textit{operator}\index{dom@\texttt{dom}}\\
a term-forming head \\
characterized by \texttt{dom-of-fun}  \textit{(primitive)}:
\begin{quote}\ttfamily\footnotesize\raggedright\noindent
forall(\allowbreak{}[a, f], f in fun(\allowbreak{}a) and a in set implies dom(\allowbreak{}f) in set and dom(\allowbreak{}f) = a)
\end{quote}
also: \texttt{dom-membership} \texttt{dom-fun-membership} \\
mentioned by 58 result(s)

\smallskip

\noindent\textbf{\texttt{dsum-act(s, n, f)}}\quad\textit{functoid}\index{dsum-act(s, n, f)@\texttt{dsum-act(s, n, f)}}\\
definition:
\begin{quote}\ttfamily\footnotesize\raggedright\noindent
vnb-lambda(\allowbreak{}[dqr\_, dsf\_], cartesian(\allowbreak{}carr(\allowbreak{}s), dsum-vec(\allowbreak{}n, f)), vnb-lambda(\allowbreak{}dsi\_, interval(\allowbreak{}1, n), (\allowbreak{}act(\allowbreak{}f(\allowbreak{}dsi\_)))(\allowbreak{}dqr\_, dsf\_(\allowbreak{}dsi\_))))
\end{quote}
characterized by \texttt{dsum-act-fun0}  \textit{(proven)}:
\begin{quote}\ttfamily\footnotesize\raggedright\noindent
forall(\allowbreak{}[s, n, f], is-ring(\allowbreak{}s) implies n in nn implies forall(\allowbreak{}[i\_ in interval(\allowbreak{}1, n)], is-module(\allowbreak{}f(\allowbreak{}i\_))) implies forall(\allowbreak{}[i\_ in interval(\allowbreak{}1, n)], scal(\allowbreak{}f(\allowbreak{}i\_)) = s) implies dsum-act(\allowbreak{}s, n, f) in fun(\allowbreak{}cartesian(\allowbreak{}carr(\allowbreak{}s), dsum-vec(\allowbreak{}n, f)), dsum-vec(\allowbreak{}n, f)))
\end{quote}
also: \texttt{metric-group-group@dsum-act} \\
mentioned by 5 result(s)

\smallskip

\noindent\textbf{\texttt{dsum-vadd(n, f)}}\quad\textit{functoid}\index{dsum-vadd(n, f)@\texttt{dsum-vadd(n, f)}}\\
definition:
\begin{quote}\ttfamily\footnotesize\raggedright\noindent
vnb-lambda(\allowbreak{}[dsf\_, dsg\_], cartesian(\allowbreak{}dsum-vec(\allowbreak{}n, f), dsum-vec(\allowbreak{}n, f)), vnb-lambda(\allowbreak{}dsi\_, interval(\allowbreak{}1, n), (\allowbreak{}vadd(\allowbreak{}f(\allowbreak{}dsi\_)))(\allowbreak{}dsf\_(\allowbreak{}dsi\_), dsg\_(\allowbreak{}dsi\_))))
\end{quote}
characterized by \texttt{dsum-vadd-fun0}  \textit{(proven)}:
\begin{quote}\ttfamily\footnotesize\raggedright\noindent
forall(\allowbreak{}[s, n, f], is-ring(\allowbreak{}s) implies n in nn implies forall(\allowbreak{}[i\_ in interval(\allowbreak{}1, n)], is-module(\allowbreak{}f(\allowbreak{}i\_))) implies forall(\allowbreak{}[i\_ in interval(\allowbreak{}1, n)], scal(\allowbreak{}f(\allowbreak{}i\_)) = s) implies dsum-vadd(\allowbreak{}n, f) in fun(\allowbreak{}cartesian(\allowbreak{}dsum-vec(\allowbreak{}n, f), dsum-vec(\allowbreak{}n, f)), dsum-vec(\allowbreak{}n, f)))
\end{quote}
mentioned by 3 result(s)

\smallskip

\noindent\textbf{\texttt{dsum-vec(n, f)}}\quad\textit{functoid}\index{dsum-vec(n, f)@\texttt{dsum-vec(n, f)}}\\
definition:
\begin{quote}\ttfamily\footnotesize\raggedright\noindent
\{dsv\_ in fun(\allowbreak{}interval(\allowbreak{}1, n), big-union(\allowbreak{}dsk\_, interval(\allowbreak{}1, n), vec(\allowbreak{}f(\allowbreak{}dsk\_)))): forall(\allowbreak{}[dsj\_ in interval(\allowbreak{}1, n)], dsv\_(\allowbreak{}dsj\_) in vec(\allowbreak{}f(\allowbreak{}dsj\_)))\}
\end{quote}
characterized by \texttt{dsum-vec-in}  \textit{(proven)}:
\begin{quote}\ttfamily\footnotesize\raggedright\noindent
forall(\allowbreak{}[n, f, g\_ in fun(\allowbreak{}interval(\allowbreak{}1, n), big-union(\allowbreak{}dsk\_, interval(\allowbreak{}1, n), vec(\allowbreak{}f(\allowbreak{}dsk\_))))], forall(\allowbreak{}[dsj\_ in interval(\allowbreak{}1, n)], g\_(\allowbreak{}dsj\_) in vec(\allowbreak{}f(\allowbreak{}dsj\_))) implies g\_ in dsum-vec(\allowbreak{}n, f))
\end{quote}
also: \texttt{dsum-vec-out} \texttt{dsum-vec-set} \texttt{dsum-vec-unfold} \texttt{dsum-vec-set-gen} \\
mentioned by 27 result(s)

\smallskip

\noindent\textbf{\texttt{dsum-vneg(n, f)}}\quad\textit{functoid}\index{dsum-vneg(n, f)@\texttt{dsum-vneg(n, f)}}\\
definition:
\begin{quote}\ttfamily\footnotesize\raggedright\noindent
vnb-lambda(\allowbreak{}dsf\_, dsum-vec(\allowbreak{}n, f), vnb-lambda(\allowbreak{}dsi\_, interval(\allowbreak{}1, n), (\allowbreak{}vneg(\allowbreak{}f(\allowbreak{}dsi\_)))(\allowbreak{}dsf\_(\allowbreak{}dsi\_))))
\end{quote}
characterized by \texttt{dsum-vneg-fun0}  \textit{(proven)}:
\begin{quote}\ttfamily\footnotesize\raggedright\noindent
forall(\allowbreak{}[s, n, f], is-ring(\allowbreak{}s) implies n in nn implies forall(\allowbreak{}[i\_ in interval(\allowbreak{}1, n)], is-module(\allowbreak{}f(\allowbreak{}i\_))) implies forall(\allowbreak{}[i\_ in interval(\allowbreak{}1, n)], scal(\allowbreak{}f(\allowbreak{}i\_)) = s) implies dsum-vneg(\allowbreak{}n, f) in fun(\allowbreak{}dsum-vec(\allowbreak{}n, f), dsum-vec(\allowbreak{}n, f)))
\end{quote}
mentioned by 3 result(s)

\smallskip

\noindent\textbf{\texttt{dsum-vzero(n, f)}}\quad\textit{functoid}\index{dsum-vzero(n, f)@\texttt{dsum-vzero(n, f)}}\\
definition:
\begin{quote}\ttfamily\footnotesize\raggedright\noindent
vnb-lambda(\allowbreak{}dsi\_, interval(\allowbreak{}1, n), vzero(\allowbreak{}f(\allowbreak{}dsi\_)))
\end{quote}
characterized by \texttt{dsum-vzero-in0}  \textit{(proven)}:
\begin{quote}\ttfamily\footnotesize\raggedright\noindent
forall(\allowbreak{}[s, n, f], is-ring(\allowbreak{}s) implies n in nn implies forall(\allowbreak{}[i\_ in interval(\allowbreak{}1, n)], is-module(\allowbreak{}f(\allowbreak{}i\_))) implies forall(\allowbreak{}[i\_ in interval(\allowbreak{}1, n)], scal(\allowbreak{}f(\allowbreak{}i\_)) = s) implies dsum-vzero(\allowbreak{}n, f) in dsum-vec(\allowbreak{}n, f))
\end{quote}
mentioned by 3 result(s)

\smallskip

\noindent\textbf{\texttt{dual(m)}}\quad\textit{functoid}\index{dual(m)@\texttt{dual(m)}}\\
reads: the dual space of \$1 \\
definition:
\begin{quote}\ttfamily\footnotesize\raggedright\noindent
[normed-field-as-commutative-ring(\allowbreak{}rr-normed-field), dual-vec(\allowbreak{}m), vnb-lambda(\allowbreak{}[dlaf\_, dlag\_], cartesian(\allowbreak{}dual-vec(\allowbreak{}m), dual-vec(\allowbreak{}m)), vnb-lambda(\allowbreak{}dlax\_, vec(\allowbreak{}m), dlaf\_(\allowbreak{}dlax\_) + dlag\_(\allowbreak{}dlax\_))), vnb-lambda(\allowbreak{}dlzx\_, vec(\allowbreak{}m), 0), vnb-lambda(\allowbreak{}dlnf\_, dual-vec(\allowbreak{}m), vnb-lambda(\allowbreak{}dlnx\_, vec(\allowbreak{}m), -dlnf\_(\allowbreak{}dlnx\_))), vnb-lambda(\allowbreak{}[dlsr\_, dlsf\_], cartesian(\allowbreak{}rr, dual-vec(\allowbreak{}m)), vnb-lambda(\allowbreak{}dlsx\_, vec(\allowbreak{}m), dlsr\_ * dlsf\_(\allowbreak{}dlsx\_))), vnb-lambda(\allowbreak{}dlvf\_, dual-vec(\allowbreak{}m), dual-norm(\allowbreak{}m, dlvf\_))]
\end{quote}
characterized by \texttt{dual-act-typing-cnvs-as-nvs}  \textit{(definitional)}:
\begin{quote}\ttfamily\footnotesize\raggedright\noindent
forall(\allowbreak{}[r], is-complex-normed-vector-space(\allowbreak{}r) implies act(\allowbreak{}dual(\allowbreak{}cnvs-as-nvs(\allowbreak{}r))) in fun(\allowbreak{}cartesian(\allowbreak{}rr, dual-vec(\allowbreak{}cnvs-as-nvs(\allowbreak{}r))), dual-vec(\allowbreak{}cnvs-as-nvs(\allowbreak{}r))))
\end{quote}
also: \texttt{dual-vadd-typing-cnvs-as-nvs} \texttt{dual-vneg-typing-cnvs-as-nvs} \texttt{dual-vnrm-typing-cnvs-as-nvs} \texttt{dual-vzero-in-dual-cnvs-as-nvs} \texttt{dual-vzero-readoff-cnvs-as-nvs} \texttt{dual-is-normed-vector-space-cnvs-as-nvs} \texttt{dual-act-apply} \texttt{dual-act-typing} ... (38 in all) \\
mentioned by 60 result(s)

\smallskip

\noindent\textbf{\texttt{dual-ball-ms(m, xs)}}\quad\textit{functoid}\index{dual-ball-ms(m, xs)@\texttt{dual-ball-ms(m, xs)}}\\
reads: the dual unit ball of \$1 with the metric of the dense sequence \$2 \\
definition:
\begin{quote}\ttfamily\footnotesize\raggedright\noindent
pullback-ms(\allowbreak{}product-metric(\allowbreak{}dbi(\allowbreak{}m, xs)), dual-unit-ball(\allowbreak{}m), dbj(\allowbreak{}m, xs))
\end{quote}
characterized by \texttt{dual-ball-ms-facts}  \textit{(proven)}:
\begin{quote}\ttfamily\footnotesize\raggedright\noindent
forall(\allowbreak{}[m, xs], is-normed-vector-space(\allowbreak{}m) implies xs in fun(\allowbreak{}nn, vec(\allowbreak{}m)) implies forall(\allowbreak{}[dz\_ in vec(\allowbreak{}m), deps\_], pos-rr(\allowbreak{}deps\_) implies forsome(\allowbreak{}[dk\_ in nn], (\allowbreak{}vnrm(\allowbreak{}m))(\allowbreak{}(\allowbreak{}vadd(\allowbreak{}m))(\allowbreak{}dz\_, (\allowbreak{}vneg(\allowbreak{}m))(\allowbreak{}xs(\allowbreak{}dk\_)))) < deps\_)) implies is-metric-space(\allowbreak{}dual-ball-ms(\allowbreak{}m, xs)) and pts(\allowbreak{}dual-ball-ms(\allowbreak{}m, xs)) == dual-unit-ball(\allowbreak{}m) and dual-ball-ms(\allowbreak{}m, xs) == pullback-ms(\allowbreak{}product-metric(\allowbreak{}dbi(\allowbreak{}m, xs)), dual-unit-ball(\allowbreak{}m), dbj(\allowbreak{}m, xs)))
\end{quote}
mentioned by 13 result(s)

\smallskip

\noindent\textbf{\texttt{dual-norm(m, f)}}\quad\textit{functoid}\index{dual-norm(m, f)@\texttt{dual-norm(m, f)}}\\
definition:
\begin{quote}\ttfamily\footnotesize\raggedright\noindent
iota(\allowbreak{}c\_, c\_ in rr and 0 <= c\_ and forall(\allowbreak{}[x\_ in vec(\allowbreak{}m)], abs(\allowbreak{}f(\allowbreak{}x\_)) <= c\_ * (\allowbreak{}vnrm(\allowbreak{}m))(\allowbreak{}x\_)) and forall(\allowbreak{}[d\_], d\_ in rr and 0 <= d\_ and forall(\allowbreak{}[x\_ in vec(\allowbreak{}m)], abs(\allowbreak{}f(\allowbreak{}x\_)) <= d\_ * (\allowbreak{}vnrm(\allowbreak{}m))(\allowbreak{}x\_)) implies c\_ <= d\_))
\end{quote}
characterized by \texttt{dual-norm-spec}  \textit{(proven)}:
\begin{quote}\ttfamily\footnotesize\raggedright\noindent
forall(\allowbreak{}[m, f], is-normed-vector-space(\allowbreak{}m) implies is-bounded-linear-functional(\allowbreak{}m, f) implies dual-norm(\allowbreak{}m, f) in rr and 0 <= dual-norm(\allowbreak{}m, f) and forall(\allowbreak{}[x\_ in vec(\allowbreak{}m)], abs(\allowbreak{}f(\allowbreak{}x\_)) <= dual-norm(\allowbreak{}m, f) * (\allowbreak{}vnrm(\allowbreak{}m))(\allowbreak{}x\_)) and forall(\allowbreak{}[d\_], d\_ in rr and 0 <= d\_ and forall(\allowbreak{}[x\_ in vec(\allowbreak{}m)], abs(\allowbreak{}f(\allowbreak{}x\_)) <= d\_ * (\allowbreak{}vnrm(\allowbreak{}m))(\allowbreak{}x\_)) implies dual-norm(\allowbreak{}m, f) <= d\_))
\end{quote}
also: \texttt{dual-norm-nonneg-nvs} \texttt{dual-norm-is-bound-nvs} \texttt{dual-norm-le-bound-nvs} \texttt{dual-norm-value-proven} \texttt{dual-norm-zero-is-vzero} \\
mentioned by 22 result(s)

\smallskip

\noindent\textbf{\texttt{dual-norm-on(m, s, f)}}\quad\textit{functoid}\index{dual-norm-on(m, s, f)@\texttt{dual-norm-on(m, s, f)}}\\
definition:
\begin{quote}\ttfamily\footnotesize\raggedright\noindent
iota(\allowbreak{}c\_, c\_ in rr and 0 <= c\_ and forall(\allowbreak{}[x\_ in s], abs(\allowbreak{}f(\allowbreak{}x\_)) <= c\_ * (\allowbreak{}vnrm(\allowbreak{}m))(\allowbreak{}x\_)) and forall(\allowbreak{}[d\_], d\_ in rr and 0 <= d\_ and forall(\allowbreak{}[x\_ in s], abs(\allowbreak{}f(\allowbreak{}x\_)) <= d\_ * (\allowbreak{}vnrm(\allowbreak{}m))(\allowbreak{}x\_)) implies c\_ <= d\_))
\end{quote}
characterized by \texttt{dual-norm-on-spec}  \textit{(proven)}:
\begin{quote}\ttfamily\footnotesize\raggedright\noindent
forall(\allowbreak{}[m, s, f], is-normed-vector-space(\allowbreak{}m) implies s subset vec(\allowbreak{}m) implies is-bounded-linear-functional-on(\allowbreak{}m, s, f) implies dual-norm-on(\allowbreak{}m, s, f) in rr and 0 <= dual-norm-on(\allowbreak{}m, s, f) and forall(\allowbreak{}[x\_ in s], abs(\allowbreak{}f(\allowbreak{}x\_)) <= dual-norm-on(\allowbreak{}m, s, f) * (\allowbreak{}vnrm(\allowbreak{}m))(\allowbreak{}x\_)) and forall(\allowbreak{}[d\_], d\_ in rr and 0 <= d\_ and forall(\allowbreak{}[x\_ in s], abs(\allowbreak{}f(\allowbreak{}x\_)) <= d\_ * (\allowbreak{}vnrm(\allowbreak{}m))(\allowbreak{}x\_)) implies dual-norm-on(\allowbreak{}m, s, f) <= d\_))
\end{quote}
also: \texttt{dual-norm-on-nonneg-nvs} \texttt{dual-norm-on-is-bound-nvs} \texttt{dual-norm-on-le-bound-nvs} \\
mentioned by 11 result(s)

\smallskip

\noindent\textbf{\texttt{dual-unit-ball(m)}}\quad\textit{functoid}\index{dual-unit-ball(m)@\texttt{dual-unit-ball(m)}}\\
reads: the closed unit ball of the dual of the normed space \$1 \\
definition:
\begin{quote}\ttfamily\footnotesize\raggedright\noindent
\{dbf\_ in dual-vec(\allowbreak{}m): dual-norm(\allowbreak{}m, dbf\_) <= 1\}
\end{quote}
characterized by \texttt{dual-unit-ball-in-set-cnvs-as-nvs}  \textit{(definitional)}:
\begin{quote}\ttfamily\footnotesize\raggedright\noindent
forall(\allowbreak{}[r], is-complex-normed-vector-space(\allowbreak{}r) implies dual-unit-ball(\allowbreak{}cnvs-as-nvs(\allowbreak{}r)) in set)
\end{quote}
also: \texttt{dual-unit-ball-weak-star-compact-cnvs-as-nvs} \texttt{dual-unit-ball-closed-in-line-product-cnvs-as-nvs} \texttt{dual-unit-ball-mem} \texttt{dual-unit-ball-in-set} \texttt{dual-unit-ball-unfold} \texttt{dual-unit-ball-of-no-failure} \texttt{dual-unit-ball-weak-star-compact} \texttt{dual-unit-ball-closed-in-line-product} \\
mentioned by 46 result(s)

\smallskip

\noindent\textbf{\texttt{dual-vec(m)}}\quad\textit{functoid}\index{dual-vec(m)@\texttt{dual-vec(m)}}\\
reads: the bounded linear functionals on \$1 \\
definition:
\begin{quote}\ttfamily\footnotesize\raggedright\noindent
\{f\_ in fun(\allowbreak{}vec(\allowbreak{}m), rr): is-bounded-linear-functional(\allowbreak{}m, f\_)\}
\end{quote}
characterized by \texttt{dual-vec-in-set-cnvs-as-nvs}  \textit{(definitional)}:
\begin{quote}\ttfamily\footnotesize\raggedright\noindent
forall(\allowbreak{}[r], is-complex-normed-vector-space(\allowbreak{}r) implies dual-vec(\allowbreak{}cnvs-as-nvs(\allowbreak{}r)) in set)
\end{quote}
also: \texttt{dual-vec-blf} \texttt{dual-vec-fun} \texttt{dual-vec-intro} \texttt{dual-vec-in-set} \texttt{dual-vec-additive} \texttt{dual-vec-homogeneous} \texttt{dual-vec-membership-proven} \\
mentioned by 74 result(s)

\smallskip

\noindent\textbf{\texttt{dual-weak-star-top(m)}}\quad\textit{functoid}\index{dual-weak-star-top(m)@\texttt{dual-weak-star-top(m)}}\\
reads: the weak* topology on the dual of the normed space \$1 \\
definition:
\begin{quote}\ttfamily\footnotesize\raggedright\noindent
subspace-top(\allowbreak{}tproduct-top(\allowbreak{}vec(\allowbreak{}m), vnb-lambda(\allowbreak{}b6n\_, vec(\allowbreak{}m), metric-top(\allowbreak{}rr-ms))), dual-vec(\allowbreak{}m))
\end{quote}
characterized by \texttt{dual-weak-star-top-unfold}  \textit{(proven)}:
\begin{quote}\ttfamily\footnotesize\raggedright\noindent
forall(\allowbreak{}[m], dual-weak-star-top(\allowbreak{}m) == subspace-top(\allowbreak{}tproduct-top(\allowbreak{}vec(\allowbreak{}m), vnb-lambda(\allowbreak{}b6n\_, vec(\allowbreak{}m), metric-top(\allowbreak{}rr-ms))), dual-vec(\allowbreak{}m)))
\end{quote}
mentioned by 9 result(s)

\smallskip

\noindent\textbf{\texttt{duality-map(chk\_, chc\_, cha\_)}}\quad\textit{functoid}\index{duality-map(chk_, chc_, cha_)@\texttt{duality-map(chk\_, chc\_, cha\_)}}\\
reads: the character x |-> \$2(\$3 x) of \$1 \\
definition:
\begin{quote}\ttfamily\footnotesize\raggedright\noindent
vnb-lambda(\allowbreak{}chx\_, carr(\allowbreak{}chk\_), chc\_(\allowbreak{}(\allowbreak{}mul(\allowbreak{}chk\_))(\allowbreak{}cha\_, chx\_)))
\end{quote}
mentioned by 15 result(s)

\smallskip

\noindent\textbf{\texttt{dynkin-generated(omega, ce)}}\quad\textit{functoid}\index{dynkin-generated(omega, ce)@\texttt{dynkin-generated(omega, ce)}}\\
reads: the Dynkin system of subsets of \$1 generated by \$2 \\
definition:
\begin{quote}\ttfamily\footnotesize\raggedright\noindent
\{dga\_ in power(\allowbreak{}omega): forall(\allowbreak{}[dgb\_], dgb\_ in power(\allowbreak{}power(\allowbreak{}omega)) and is-dynkin-system(\allowbreak{}omega, dgb\_) and ce subset dgb\_ implies dga\_ in dgb\_)\}
\end{quote}
characterized by \texttt{dynkin-generated-unfold}  \textit{(proven)}:
\begin{quote}\ttfamily\footnotesize\raggedright\noindent
forall(\allowbreak{}[omega, ce], dynkin-generated(\allowbreak{}omega, ce) == \{dga\_ in power(\allowbreak{}omega): forall(\allowbreak{}[dgb\_], dgb\_ in power(\allowbreak{}power(\allowbreak{}omega)) and is-dynkin-system(\allowbreak{}omega, dgb\_) and ce subset dgb\_ implies dga\_ in dgb\_)\})
\end{quote}
also: \texttt{dynkin-generated-is-dynkin} \texttt{dynkin-generated-membership} \\
mentioned by 5 result(s)

\smallskip


\section*{E}
\addcontentsline{toc}{section}{E}

\noindent\textbf{\texttt{e-polar(c6vm\_, c6vw\_, c6vv\_)}}\quad\textit{functoid}\index{e-polar(c6vm_, c6vw_, c6vv_)@\texttt{e-polar(c6vm\_, c6vw\_, c6vv\_)}}\\
reads: the polar of \$3 in \$1 with respect to \$2 \\
definition:
\begin{quote}\ttfamily\footnotesize\raggedright\noindent
\{c6vh\_ in vec(\allowbreak{}c6vm\_): forall(\allowbreak{}[c6vx\_ in c6vv\_], c6vw\_(\allowbreak{}c6vh\_, c6vx\_) = 1)\}
\end{quote}
mentioned by 6 result(s)

\smallskip

\noindent\textbf{\texttt{eapp(mxh\_, mxl\_)}}\quad\textit{functoid}\index{eapp(mxh_, mxl_)@\texttt{eapp(mxh\_, mxl\_)}}\\
reads: the application of the head \$1 to the list \$2 \\
definition:
\begin{quote}\ttfamily\footnotesize\raggedright\noindent
[3, mxh\_, mxl\_]
\end{quote}
characterized by \texttt{fv-eapp}  \textit{(proven)}:
\begin{quote}\ttfamily\footnotesize\raggedright\noindent
forall(\allowbreak{}[h, l], h in nn implies l in tuples(\allowbreak{}expr-set) implies expr-fv(\allowbreak{}eapp(\allowbreak{}h, l)) == big-union(\allowbreak{}mxw\_, make-set(\allowbreak{}l), expr-fv(\allowbreak{}mxw\_)))
\end{quote}
also: \texttt{lcd-eapp} \texttt{wks-eapp} \texttt{open-eapp} \texttt{subst-eapp} \texttt{eapp-unfold} \texttt{lc-eapp-iff} \texttt{clt-eapp-arg} \texttt{eapp-in-expr} ... (13 in all) \\
mentioned by 171 result(s)

\smallskip

\noindent\textbf{\texttt{ebind(mxh\_, mxn\_, mxa\_, mxb\_)}}\quad\textit{functoid}\index{ebind(mxh_, mxn_, mxa_, mxb_)@\texttt{ebind(mxh\_, mxn\_, mxa\_, mxb\_)}}\\
reads: the binder \$1 of \$2 variables over the list \$3 with body \$4 \\
definition:
\begin{quote}\ttfamily\footnotesize\raggedright\noindent
[4, mxh\_, mxn\_, mxa\_, mxb\_]
\end{quote}
characterized by \texttt{fv-ebind}  \textit{(proven)}:
\begin{quote}\ttfamily\footnotesize\raggedright\noindent
forall(\allowbreak{}[h, nb, a, b], h in nn implies nb in nn implies a in tuples(\allowbreak{}expr-set) implies b in expr-set implies expr-fv(\allowbreak{}ebind(\allowbreak{}h, nb, a, b)) == union(\allowbreak{}big-union(\allowbreak{}mxw\_, make-set(\allowbreak{}a), expr-fv(\allowbreak{}mxw\_)), expr-fv(\allowbreak{}b)))
\end{quote}
also: \texttt{lcd-ebind} \texttt{wks-ebind} \texttt{open-ebind} \texttt{subst-ebind} \texttt{ebind-unfold} \texttt{clt-ebind-dom} \texttt{ebind-in-expr} \texttt{ebind-args-expr} ... (13 in all) \\
mentioned by 66 result(s)

\smallskip

\noindent\textbf{\texttt{ebvar(mxk\_)}}\quad\textit{functoid}\index{ebvar(mxk_)@\texttt{ebvar(mxk\_)}}\\
reads: the bound variable of index \$1 \\
definition:
\begin{quote}\ttfamily\footnotesize\raggedright\noindent
[1, mxk\_]
\end{quote}
characterized by \texttt{fv-ebvar}  \textit{(proven)}:
\begin{quote}\ttfamily\footnotesize\raggedright\noindent
forall(\allowbreak{}[k in nn], expr-fv(\allowbreak{}ebvar(\allowbreak{}k)) == empty-set)
\end{quote}
also: \texttt{lcd-ebvar} \texttt{wks-ebvar} \texttt{subst-ebvar} \texttt{ebvar-unfold} \texttt{ebvar-in-expr} \texttt{open-ebvar-hit} \texttt{open-ebvar-miss} \\
mentioned by 85 result(s)

\smallskip

\noindent\textbf{\texttt{econst(mxc\_)}}\quad\textit{functoid}\index{econst(mxc_)@\texttt{econst(mxc\_)}}\\
reads: the constant of code \$1 \\
definition:
\begin{quote}\ttfamily\footnotesize\raggedright\noindent
[2, mxc\_]
\end{quote}
characterized by \texttt{fv-econst}  \textit{(proven)}:
\begin{quote}\ttfamily\footnotesize\raggedright\noindent
forall(\allowbreak{}[c in nn], expr-fv(\allowbreak{}econst(\allowbreak{}c)) == empty-set)
\end{quote}
also: \texttt{lcd-econst} \texttt{wks-econst} \texttt{open-econst} \texttt{subst-econst} \texttt{econst-unfold} \texttt{econst-in-expr} \\
mentioned by 63 result(s)

\smallskip

\noindent\textbf{\texttt{econverges-to(f, l)}}\quad\textit{predicate}\index{econverges-to(f, l)@\texttt{econverges-to(f, l)}}\\
reads: the sequence \$1 converges to \$2 in the extended reals \\
defined by \texttt{econverges-to} \\
characterized by \texttt{econverges-to}  \textit{(definitional)}:
\begin{quote}\ttfamily\footnotesize\raggedright\noindent
forall(\allowbreak{}[f, l], econverges-to(\allowbreak{}f, l) iff f in fun(\allowbreak{}nn, rr-pos-star) and l in rr-pos-star and eliminf(\allowbreak{}f) = l and elimsup(\allowbreak{}f) = l)
\end{quote}
mentioned by 10 result(s)

\smallskip

\noindent\textbf{\texttt{efvar(mxn\_)}}\quad\textit{functoid}\index{efvar(mxn_)@\texttt{efvar(mxn\_)}}\\
reads: the free variable named \$1 \\
definition:
\begin{quote}\ttfamily\footnotesize\raggedright\noindent
[0, mxn\_]
\end{quote}
characterized by \texttt{fv-efvar}  \textit{(proven)}:
\begin{quote}\ttfamily\footnotesize\raggedright\noindent
forall(\allowbreak{}[n in nn], expr-fv(\allowbreak{}efvar(\allowbreak{}n)) == pair(\allowbreak{}n, n))
\end{quote}
also: \texttt{clt-efvar} \texttt{lcd-efvar} \texttt{wks-efvar} \texttt{open-efvar} \texttt{efvar-unfold} \texttt{efvar-in-expr} \texttt{subst-efvar-hit} \texttt{subst-efvar-miss} \\
mentioned by 66 result(s)

\smallskip

\noindent\textbf{\texttt{eigenspace(kqv\_, kqt\_, kqc\_)}}\quad\textit{functoid}\index{eigenspace(kqv_, kqt_, kqc_)@\texttt{eigenspace(kqv\_, kqt\_, kqc\_)}}\\
reads: the eigenspace of \$2 for \$3 \\
definition:
\begin{quote}\ttfamily\footnotesize\raggedright\noindent
\{kqx\_ in vec(\allowbreak{}kqv\_): kqt\_(\allowbreak{}kqx\_) = (\allowbreak{}act(\allowbreak{}kqv\_))(\allowbreak{}kqc\_, kqx\_)\}
\end{quote}
characterized by \texttt{eigenspace-unfold}  \textit{(proven)}:
\begin{quote}\ttfamily\footnotesize\raggedright\noindent
forall(\allowbreak{}[kqv\_, kqt\_, kqc\_], eigenspace(\allowbreak{}kqv\_, kqt\_, kqc\_) == \{kqx\_ in vec(\allowbreak{}kqv\_): kqt\_(\allowbreak{}kqx\_) = (\allowbreak{}act(\allowbreak{}kqv\_))(\allowbreak{}kqc\_, kqx\_)\})
\end{quote}
also: \texttt{eigenspace-submodule} \texttt{eigenspace-membership} \texttt{eigenspace-onb-finite} \texttt{sa-eigenspaces-orthogonal} \\
mentioned by 19 result(s)

\smallskip

\noindent\textbf{\texttt{einf(s)}}\quad\textit{functoid}\index{einf(s)@\texttt{einf(s)}}\\
reads: the infimum of \$1 in the extended reals \\
definition:
\begin{quote}\ttfamily\footnotesize\raggedright\noindent
esup(\allowbreak{}elower-bounds(\allowbreak{}s))
\end{quote}
characterized by \texttt{einf-in}  \textit{(proven)}:
\begin{quote}\ttfamily\footnotesize\raggedright\noindent
forall(\allowbreak{}[cpss\_], cpss\_ subset rr-pos-star implies einf(\allowbreak{}cpss\_) in rr-pos-star)
\end{quote}
also: \texttt{einf-lower} \texttt{einf-greatest} \texttt{einf-lt-witness} \texttt{einf-etimes-bound} \texttt{einf-reflection-tail} \\
mentioned by 34 result(s)

\smallskip

\noindent\textbf{\texttt{el-pow}}\quad\textit{functoid}\index{el-pow@\texttt{el-pow}}\\
reads: the \$3-th power of \$2 in \$1 \\
defined by \texttt{el-pow-zero} \texttt{el-pow-succ} \\
characterized by \texttt{el-pow-succ}  \textit{(definitional)}:
\begin{quote}\ttfamily\footnotesize\raggedright\noindent
forall(\allowbreak{}[bpa\_, bpx\_, bpn\_ in nn], el-pow(\allowbreak{}bpa\_, bpx\_, succ(\allowbreak{}bpn\_)) == (\allowbreak{}vmul(\allowbreak{}bpa\_))(\allowbreak{}bpx\_, el-pow(\allowbreak{}bpa\_, bpx\_, bpn\_)))
\end{quote}
also: \texttt{el-pow-zero} \texttt{b2r-el-pow-sa} \texttt{bae-el-pow-one} \texttt{baf-el-pow-add} \texttt{bsr-el-pow-act} \texttt{b2r-el-pow-view} \texttt{bae-el-pow-comm} \texttt{bae-el-pow-norm} ... (11 in all) \\
mentioned by 45 result(s)

\smallskip

\noindent\textbf{\texttt{elem-equiv(r, m, n, a, b)}}\quad\textit{predicate}\index{elem-equiv(r, m, n, a, b)@\texttt{elem-equiv(r, m, n, a, b)}}\\
reads: \$5 is obtained from \$4 by elementary row and column operations \\
defined by \texttt{elem-equiv} \\
characterized by \texttt{elem-equiv}  \textit{(definitional)}:
\begin{quote}\ttfamily\footnotesize\raggedright\noindent
forall(\allowbreak{}[r, m, n, a, b], elem-equiv(\allowbreak{}r, m, n, a, b) iff forsome(\allowbreak{}[eqp\_, eqq\_], eqp\_ in elem-gen(\allowbreak{}r, m) and eqq\_ in elem-gen(\allowbreak{}r, n) and b = matmul(\allowbreak{}r, matmul(\allowbreak{}r, eqp\_, a), eqq\_)))
\end{quote}
also: \texttt{elem-equiv-trans} \texttt{elem-equiv-symmetric} \\
mentioned by 12 result(s)

\smallskip

\noindent\textbf{\texttt{elem-f(a, n, k, l)}}\quad\textit{functoid}\index{elem-f(a, n, k, l)@\texttt{elem-f(a, n, k, l)}}\\
definition:
\begin{quote}\ttfamily\footnotesize\raggedright\noindent
matof(\allowbreak{}n, n, vnb-lambda(\allowbreak{}[i, j], cartesian(\allowbreak{}interval(\allowbreak{}1, n), interval(\allowbreak{}1, n)), if(\allowbreak{}i = k and j = l or i = l and j = k or i = j and not(\allowbreak{}i = k) and not(\allowbreak{}i = l), one(\allowbreak{}a), zero(\allowbreak{}a))))
\end{quote}
characterized by \texttt{elem-f-type-uba-as-ring}  \textit{(definitional)}:
\begin{quote}\ttfamily\footnotesize\raggedright\noindent
forall(\allowbreak{}[r], is-unital-banach-algebra(\allowbreak{}r) implies forall(\allowbreak{}[n, k, l], n in nn implies elem-f(\allowbreak{}uba-as-ring(\allowbreak{}r), n, k, l) in mat(\allowbreak{}n, n, vec(\allowbreak{}r))))
\end{quote}
also: \texttt{elem-f-action-uba-as-ring} \texttt{elem-f-inverse-uba-as-ring} \texttt{elem-f-symmetric-uba-as-ring} \texttt{elem-f-invertible-uba-as-ring} \texttt{elem-f-row-action-uba-as-ring} \texttt{elem-f-type-uba-as-ring-uiba-as-uba} \texttt{elem-f-action-uba-as-ring-uiba-as-uba} \texttt{elem-f-inverse-uba-as-ring-uiba-as-uba} ... (40 in all) \\
mentioned by 119 result(s)

\smallskip

\noindent\textbf{\texttt{elem-g(a, n, r, k, l)}}\quad\textit{functoid}\index{elem-g(a, n, r, k, l)@\texttt{elem-g(a, n, r, k, l)}}\\
definition:
\begin{quote}\ttfamily\footnotesize\raggedright\noindent
matof(\allowbreak{}n, n, vnb-lambda(\allowbreak{}[i, j], cartesian(\allowbreak{}interval(\allowbreak{}1, n), interval(\allowbreak{}1, n)), if(\allowbreak{}i = j, one(\allowbreak{}a), if(\allowbreak{}i = k and j = l, r, zero(\allowbreak{}a)))))
\end{quote}
characterized by \texttt{elem-g-type-uba-as-ring}  \textit{(definitional)}:
\begin{quote}\ttfamily\footnotesize\raggedright\noindent
forall(\allowbreak{}[r], is-unital-banach-algebra(\allowbreak{}r) implies forall(\allowbreak{}[n, r\_35051, k, l], n in nn implies r\_35051 in vec(\allowbreak{}r) implies elem-g(\allowbreak{}uba-as-ring(\allowbreak{}r), n, r\_35051, k, l) in mat(\allowbreak{}n, n, vec(\allowbreak{}r))))
\end{quote}
also: \texttt{elem-g-action-uba-as-ring} \texttt{elem-g-inverse-uba-as-ring} \texttt{elem-g-invertible-uba-as-ring} \texttt{elem-g-param-cong-uba-as-ring} \texttt{elem-g-row-action-uba-as-ring} \texttt{elem-g-type-uba-as-ring-uiba-as-uba} \texttt{elem-g-action-uba-as-ring-uiba-as-uba} \texttt{elem-g-inverse-uba-as-ring-uiba-as-uba} ... (38 in all) \\
mentioned by 113 result(s)

\smallskip

\noindent\textbf{\texttt{elem-gen(r, n)}}\quad\textit{functoid}\index{elem-gen(r, n)@\texttt{elem-gen(r, n)}}\\
reads: the set of products of elementary column \$2-by-\$2 matrices over \$1 \\
definition:
\begin{quote}\ttfamily\footnotesize\raggedright\noindent
intersection-of(\allowbreak{}\{egs\_ in power(\allowbreak{}mat(\allowbreak{}n, n, carr(\allowbreak{}r))): identmat(\allowbreak{}r, n) in egs\_ and forall(\allowbreak{}[egx\_ in egs\_, ege\_ in set], is-elementary(\allowbreak{}r, n, ege\_) implies matmul(\allowbreak{}r, egx\_, ege\_) in egs\_)\})
\end{quote}
characterized by \texttt{elem-gen-mat}  \textit{(proven)}:
\begin{quote}\ttfamily\footnotesize\raggedright\noindent
forall(\allowbreak{}[r, n, p], is-ring(\allowbreak{}r) implies n in nn implies p in elem-gen(\allowbreak{}r, n) implies p in mat(\allowbreak{}n, n, carr(\allowbreak{}r)))
\end{quote}
also: \texttt{elem-gen-mul} \texttt{elem-gen-step} \texttt{elem-gen-unfold} \texttt{elem-gen-inverse} \texttt{elem-gen-det-unit} \texttt{elem-gen-identity} \texttt{elem-gen-induction} \texttt{elem-gen-invertible} ... (13 in all) \\
mentioned by 31 result(s)

\smallskip

\noindent\textbf{\texttt{elem-h(a, n, r, k)}}\quad\textit{functoid}\index{elem-h(a, n, r, k)@\texttt{elem-h(a, n, r, k)}}\\
definition:
\begin{quote}\ttfamily\footnotesize\raggedright\noindent
matof(\allowbreak{}n, n, vnb-lambda(\allowbreak{}[i, j], cartesian(\allowbreak{}interval(\allowbreak{}1, n), interval(\allowbreak{}1, n)), if(\allowbreak{}i = j, if(\allowbreak{}i = k, r, one(\allowbreak{}a)), zero(\allowbreak{}a))))
\end{quote}
characterized by \texttt{elem-h-type-uba-as-ring}  \textit{(definitional)}:
\begin{quote}\ttfamily\footnotesize\raggedright\noindent
forall(\allowbreak{}[r], is-unital-banach-algebra(\allowbreak{}r) implies forall(\allowbreak{}[n, r\_35052, k], n in nn implies r\_35052 in vec(\allowbreak{}r) implies elem-h(\allowbreak{}uba-as-ring(\allowbreak{}r), n, r\_35052, k) in mat(\allowbreak{}n, n, vec(\allowbreak{}r))))
\end{quote}
also: \texttt{elem-h-action-uba-as-ring} \texttt{elem-h-inverse-uba-as-ring} \texttt{elem-h-row-action-uba-as-ring} \texttt{elem-h-type-uba-as-ring-uiba-as-uba} \texttt{elem-h-action-uba-as-ring-uiba-as-uba} \texttt{elem-h-inverse-uba-as-ring-uiba-as-uba} \texttt{elem-h-row-action-uba-as-ring-uiba-as-uba} \texttt{det-elem-h} ... (22 in all) \\
mentioned by 88 result(s)

\smallskip

\noindent\textbf{\texttt{eliminf(f)}}\quad\textit{functoid}\index{eliminf(f)@\texttt{eliminf(f)}}\\
reads: the lower limit of the sequence \$1 \\
definition:
\begin{quote}\ttfamily\footnotesize\raggedright\noindent
esup(\allowbreak{}image(\allowbreak{}vnb-lambda(\allowbreak{}n\_, nn, einf(\allowbreak{}etail(\allowbreak{}f, n\_))), nn))
\end{quote}
characterized by \texttt{eliminf-le-elimsup}  \textit{(proven)}:
\begin{quote}\ttfamily\footnotesize\raggedright\noindent
forall(\allowbreak{}[a in fun(\allowbreak{}nn, rr-pos-star)], eliminf(\allowbreak{}a) <= elimsup(\allowbreak{}a))
\end{quote}
also: \texttt{eliminf-reflection} \texttt{eliminf-in-rr-pos-star} \texttt{eliminf-eventually-above} \texttt{rr-converges-nonneg-eliminf} \\
mentioned by 12 result(s)

\smallskip

\noindent\textbf{\texttt{elimsup(f)}}\quad\textit{functoid}\index{elimsup(f)@\texttt{elimsup(f)}}\\
reads: the upper limit of the sequence \$1 \\
definition:
\begin{quote}\ttfamily\footnotesize\raggedright\noindent
einf(\allowbreak{}image(\allowbreak{}vnb-lambda(\allowbreak{}n\_, nn, esup(\allowbreak{}etail(\allowbreak{}f, n\_))), nn))
\end{quote}
characterized by \texttt{elimsup-le-bound}  \textit{(proven)}:
\begin{quote}\ttfamily\footnotesize\raggedright\noindent
forall(\allowbreak{}[a in fun(\allowbreak{}nn, rr-pos-star), b in rr-pos-star], forall(\allowbreak{}[n in nn], a(\allowbreak{}n) <= b) implies elimsup(\allowbreak{}a) <= b)
\end{quote}
also: \texttt{eliminf-le-elimsup} \texttt{elimsup-eventually-le} \texttt{elimsup-in-rr-pos-star} \texttt{elimsup-eventually-below} \texttt{elimsup-frequently-above} \\
mentioned by 23 result(s)

\smallskip

\noindent\textbf{\texttt{ell-one}}\quad\textit{constant}\index{ell-one@\texttt{ell-one}}\\
reads: the absolutely summable complex sequences \\
defined by \texttt{ell-one-def}

\smallskip

\noindent\textbf{\texttt{ell-two}}\quad\textit{constant}\index{ell-two@\texttt{ell-two}}\\
reads: the square-summable complex sequences \\
defined by \texttt{ell-two-def}

\smallskip

\noindent\textbf{\texttt{elower-bounds(s)}}\quad\textit{functoid}\index{elower-bounds(s)@\texttt{elower-bounds(s)}}\\
reads: the set of lower bounds of \$1 in the extended reals \\
definition:
\begin{quote}\ttfamily\footnotesize\raggedright\noindent
\{b\_ in rr-pos-star: forall(\allowbreak{}[x\_ in s], b\_ <= x\_)\}
\end{quote}

\smallskip

\noindent\textbf{\texttt{embed(m)}}\quad\textit{functoid}\index{embed(m)@\texttt{embed(m)}}\\
definition:
\begin{quote}\ttfamily\footnotesize\raggedright\noindent
vnb-lambda(\allowbreak{}u, pts(\allowbreak{}m), class(\allowbreak{}cauchy-setoid(\allowbreak{}m), embed-seq(\allowbreak{}m, u)))
\end{quote}
characterized by \texttt{embed-in-fun}  \textit{(proven)}:
\begin{quote}\ttfamily\footnotesize\raggedright\noindent
forall(\allowbreak{}[m], is-metric-space(\allowbreak{}m) implies embed(\allowbreak{}m) in fun(\allowbreak{}pts(\allowbreak{}m), pts(\allowbreak{}completion(\allowbreak{}m))))
\end{quote}
also: \texttt{qp-embed-abs} \texttt{qp-embed-add} \texttt{qp-embed-mul} \texttt{qp-embed-neg} \texttt{qp-embed-one} \texttt{qe-embed-diff} \texttt{qp-embed-zero} \texttt{qp-gpow-embed} ... (11 in all) \\
mentioned by 57 result(s)

\smallskip

\noindent\textbf{\texttt{embed-seq(m, u)}}\quad\textit{functoid}\index{embed-seq(m, u)@\texttt{embed-seq(m, u)}}\\
definition:
\begin{quote}\ttfamily\footnotesize\raggedright\noindent
vnb-lambda(\allowbreak{}n, nn, u)
\end{quote}
characterized by \texttt{rkt-embed-seq-cauchy}  \textit{(proven)}:
\begin{quote}\ttfamily\footnotesize\raggedright\noindent
forall(\allowbreak{}[m, u], is-metric-space(\allowbreak{}m) implies u in pts(\allowbreak{}m) implies is-cauchy-seq(\allowbreak{}m, embed-seq(\allowbreak{}m, u)))
\end{quote}
also: \texttt{rkt-embed-seq-in-fun} \texttt{rkt-embed-seq-in-cseq} \\
mentioned by 9 result(s)

\smallskip

\noindent\textbf{\texttt{enc-all(mcb\_)}}\quad\textit{functoid}\index{enc-all(mcb_)@\texttt{enc-all(mcb\_)}}\\
reads: the encoded universal with body \$1 \\
definition:
\begin{quote}\ttfamily\footnotesize\raggedright\noindent
ebind(\allowbreak{}20, 1, [], mcb\_)
\end{quote}
characterized by \texttt{enc-all-unfold}  \textit{(proven)}:
\begin{quote}\ttfamily\footnotesize\raggedright\noindent
forall(\allowbreak{}[mrp\_], enc-all(\allowbreak{}mrp\_) == ebind(\allowbreak{}20, 1, [], mrp\_))
\end{quote}
mentioned by 61 result(s)

\smallskip

\noindent\textbf{\texttt{enc-and(mcp\_, mcq\_)}}\quad\textit{functoid}\index{enc-and(mcp_, mcq_)@\texttt{enc-and(mcp\_, mcq\_)}}\\
reads: the encoded conjunction of \$1 and \$2 \\
definition:
\begin{quote}\ttfamily\footnotesize\raggedright\noindent
eapp(\allowbreak{}3, [mcp\_, mcq\_])
\end{quote}
characterized by \texttt{enc-ande-succ}  \textit{(definitional)}:
\begin{quote}\ttfamily\footnotesize\raggedright\noindent
forall(\allowbreak{}[mcl\_, mca\_, mcn\_ in nn], enc-ande(\allowbreak{}mcl\_, mca\_, succ(\allowbreak{}mcn\_)) == enc-and(\allowbreak{}enc-ande(\allowbreak{}mcl\_, mca\_, mcn\_), enc-in(\allowbreak{}nth(\allowbreak{}succ(\allowbreak{}mcn\_), mcl\_), mca\_)))
\end{quote}
also: \texttt{enc-andl-succ} \texttt{enc-and-unfold} \\
mentioned by 68 result(s)

\smallskip

\noindent\textbf{\texttt{enc-ande}}\quad\textit{functoid}\index{enc-ande@\texttt{enc-ande}}\\
reads: the conjunction of the first \$3 entries of \$1 in \$2 \\
defined by \texttt{enc-ande-zero} \texttt{enc-ande-succ} \\
characterized by \texttt{enc-ande-succ}  \textit{(definitional)}:
\begin{quote}\ttfamily\footnotesize\raggedright\noindent
forall(\allowbreak{}[mcl\_, mca\_, mcn\_ in nn], enc-ande(\allowbreak{}mcl\_, mca\_, succ(\allowbreak{}mcn\_)) == enc-and(\allowbreak{}enc-ande(\allowbreak{}mcl\_, mca\_, mcn\_), enc-in(\allowbreak{}nth(\allowbreak{}succ(\allowbreak{}mcn\_), mcl\_), mca\_)))
\end{quote}
also: \texttt{enc-ande-zero} \texttt{enc-ande-clf} \texttt{enc-ande-der} \texttt{enc-ande-elim-der} \\
mentioned by 9 result(s)

\smallskip

\noindent\textbf{\texttt{enc-andl}}\quad\textit{functoid}\index{enc-andl@\texttt{enc-andl}}\\
reads: the conjunction of \$1 in the first \$3 entries of \$2 \\
defined by \texttt{enc-andl-zero} \texttt{enc-andl-succ} \\
characterized by \texttt{enc-andl-succ}  \textit{(definitional)}:
\begin{quote}\ttfamily\footnotesize\raggedright\noindent
forall(\allowbreak{}[mcx\_, mcl\_, mcn\_ in nn], enc-andl(\allowbreak{}mcx\_, mcl\_, succ(\allowbreak{}mcn\_)) == enc-and(\allowbreak{}enc-andl(\allowbreak{}mcx\_, mcl\_, mcn\_), enc-in(\allowbreak{}mcx\_, nth(\allowbreak{}succ(\allowbreak{}mcn\_), mcl\_))))
\end{quote}
also: \texttt{enc-andl-zero} \texttt{enc-andl-clf} \texttt{enc-andl-der} \texttt{enc-andl-elim-der} \\
mentioned by 10 result(s)

\smallskip

\noindent\textbf{\texttt{enc-apply(mcp\_, mcq\_)}}\quad\textit{functoid}\index{enc-apply(mcp_, mcq_)@\texttt{enc-apply(mcp\_, mcq\_)}}\\
reads: the encoded application of \$1 to \$2 \\
definition:
\begin{quote}\ttfamily\footnotesize\raggedright\noindent
eapp(\allowbreak{}18, [mcp\_, mcq\_])
\end{quote}
characterized by \texttt{enc-apply-unfold}  \textit{(proven)}:
\begin{quote}\ttfamily\footnotesize\raggedright\noindent
forall(\allowbreak{}[mip\_, miq\_], enc-apply(\allowbreak{}mip\_, miq\_) == eapp(\allowbreak{}18, [mip\_, miq\_]))
\end{quote}
mentioned by 8 result(s)

\smallskip

\noindent\textbf{\texttt{enc-bu(mca\_, mcb\_)}}\quad\textit{functoid}\index{enc-bu(mca_, mcb_)@\texttt{enc-bu(mca\_, mcb\_)}}\\
reads: the encoded union over \$1 of \$2 \\
definition:
\begin{quote}\ttfamily\footnotesize\raggedright\noindent
ebind(\allowbreak{}24, 1, [mca\_], mcb\_)
\end{quote}
characterized by \texttt{enc-bu-unfold}  \textit{(proven)}:
\begin{quote}\ttfamily\footnotesize\raggedright\noindent
forall(\allowbreak{}[mia\_, mip\_], enc-bu(\allowbreak{}mia\_, mip\_) == ebind(\allowbreak{}24, 1, [mia\_], mip\_))
\end{quote}
mentioned by 13 result(s)

\smallskip

\noindent\textbf{\texttt{enc-cart(mcl\_)}}\quad\textit{functoid}\index{enc-cart(mcl_)@\texttt{enc-cart(mcl\_)}}\\
reads: the encoded product of the list \$1 \\
definition:
\begin{quote}\ttfamily\footnotesize\raggedright\noindent
eapp(\allowbreak{}13, mcl\_)
\end{quote}
characterized by \texttt{enc-cart-unfold}  \textit{(proven)}:
\begin{quote}\ttfamily\footnotesize\raggedright\noindent
forall(\allowbreak{}[l], enc-cart(\allowbreak{}l) == eapp(\allowbreak{}13, l))
\end{quote}
mentioned by 14 result(s)

\smallskip

\noindent\textbf{\texttt{enc-cartl}}\quad\textit{functoid}\index{enc-cartl@\texttt{enc-cartl}}\\
reads: the componentwise memberships of \$1 in the first \$3 entries of \$2 \\
defined by \texttt{enc-cartl-zero} \texttt{enc-cartl-succ} \\
characterized by \texttt{enc-cartl-succ}  \textit{(definitional)}:
\begin{quote}\ttfamily\footnotesize\raggedright\noindent
forall(\allowbreak{}[mcx\_, mcl\_, mcn\_ in nn], enc-cartl(\allowbreak{}mcx\_, mcl\_, succ(\allowbreak{}mcn\_)) == enc-and(\allowbreak{}enc-cartl(\allowbreak{}mcx\_, mcl\_, mcn\_), enc-in(\allowbreak{}enc-nth(\allowbreak{}enc-num(\allowbreak{}succ(\allowbreak{}mcn\_)), mcx\_), nth(\allowbreak{}succ(\allowbreak{}mcn\_), mcl\_))))
\end{quote}
also: \texttt{enc-cartl-zero} \texttt{enc-cartl-clf} \texttt{enc-cartl-der} \texttt{enc-cartl-elim-der} \\
mentioned by 9 result(s)

\smallskip

\noindent\textbf{\texttt{enc-cch}}\quad\textit{functoid}\index{enc-cch@\texttt{enc-cch}}\\
reads: the existential chain of \$1 in the last \$4 of the \$3 entries of \$2 \\
defined by \texttt{enc-cch-zero} \texttt{enc-cch-succ} \\
characterized by \texttt{enc-cch-succ}  \textit{(definitional)}:
\begin{quote}\ttfamily\footnotesize\raggedright\noindent
forall(\allowbreak{}[mix\_, mia\_, min\_, mik\_ in nn], enc-cch(\allowbreak{}mix\_, mia\_, min\_, succ(\allowbreak{}mik\_)) == enc-ex(\allowbreak{}enc-and(\allowbreak{}enc-in(\allowbreak{}ebvar(\allowbreak{}0), nth(\allowbreak{}min\_ - mik\_, mia\_)), enc-cch(\allowbreak{}mix\_, mia\_, min\_, mik\_))))
\end{quote}
also: \texttt{enc-cch-zero} \\
mentioned by 7 result(s)

\smallskip

\noindent\textbf{\texttt{enc-comp(mcb\_)}}\quad\textit{functoid}\index{enc-comp(mcb_)@\texttt{enc-comp(mcb\_)}}\\
reads: the encoded comprehension by \$1 \\
definition:
\begin{quote}\ttfamily\footnotesize\raggedright\noindent
ebind(\allowbreak{}23, 1, [], mcb\_)
\end{quote}
characterized by \texttt{enc-comp-unfold}  \textit{(proven)}:
\begin{quote}\ttfamily\footnotesize\raggedright\noindent
forall(\allowbreak{}[mip\_], enc-comp(\allowbreak{}mip\_) == ebind(\allowbreak{}23, 1, [], mip\_))
\end{quote}
mentioned by 18 result(s)

\smallskip

\noindent\textbf{\texttt{enc-eq(mcp\_, mcq\_)}}\quad\textit{functoid}\index{enc-eq(mcp_, mcq_)@\texttt{enc-eq(mcp\_, mcq\_)}}\\
reads: the encoded equation \$1 = \$2 \\
definition:
\begin{quote}\ttfamily\footnotesize\raggedright\noindent
eapp(\allowbreak{}8, [mcp\_, mcq\_])
\end{quote}
characterized by \texttt{enc-eql-succ}  \textit{(definitional)}:
\begin{quote}\ttfamily\footnotesize\raggedright\noindent
forall(\allowbreak{}[mcl\_, mca\_, mcn\_ in nn], enc-eql(\allowbreak{}mcl\_, mca\_, succ(\allowbreak{}mcn\_)) == enc-and(\allowbreak{}enc-eql(\allowbreak{}mcl\_, mca\_, mcn\_), enc-eq(\allowbreak{}nth(\allowbreak{}succ(\allowbreak{}mcn\_), mcl\_), nth(\allowbreak{}succ(\allowbreak{}mcn\_), mca\_))))
\end{quote}
also: \texttt{enc-eq-unfold} \\
mentioned by 136 result(s)

\smallskip

\noindent\textbf{\texttt{enc-eql}}\quad\textit{functoid}\index{enc-eql@\texttt{enc-eql}}\\
reads: the componentwise equations of the first \$3 entries of \$1 and \$2 \\
defined by \texttt{enc-eql-zero} \texttt{enc-eql-succ} \\
characterized by \texttt{enc-eql-succ}  \textit{(definitional)}:
\begin{quote}\ttfamily\footnotesize\raggedright\noindent
forall(\allowbreak{}[mcl\_, mca\_, mcn\_ in nn], enc-eql(\allowbreak{}mcl\_, mca\_, succ(\allowbreak{}mcn\_)) == enc-and(\allowbreak{}enc-eql(\allowbreak{}mcl\_, mca\_, mcn\_), enc-eq(\allowbreak{}nth(\allowbreak{}succ(\allowbreak{}mcn\_), mcl\_), nth(\allowbreak{}succ(\allowbreak{}mcn\_), mca\_))))
\end{quote}
also: \texttt{enc-eql-zero} \texttt{enc-eql-clf} \\
mentioned by 8 result(s)

\smallskip

\noindent\textbf{\texttt{enc-ex(mcb\_)}}\quad\textit{functoid}\index{enc-ex(mcb_)@\texttt{enc-ex(mcb\_)}}\\
reads: the encoded existential with body \$1 \\
definition:
\begin{quote}\ttfamily\footnotesize\raggedright\noindent
ebind(\allowbreak{}21, 1, [], mcb\_)
\end{quote}
characterized by \texttt{enc-ex-unfold}  \textit{(proven)}:
\begin{quote}\ttfamily\footnotesize\raggedright\noindent
forall(\allowbreak{}[mrp\_], enc-ex(\allowbreak{}mrp\_) == ebind(\allowbreak{}21, 1, [], mrp\_))
\end{quote}
mentioned by 30 result(s)

\smallskip

\noindent\textbf{\texttt{enc-falsity}}\quad\textit{constant}\index{enc-falsity@\texttt{enc-falsity}}\\
reads: the encoded FALSITY \\
defined by \texttt{enc-falsity-def}

\smallskip

\noindent\textbf{\texttt{enc-fun(mcp\_, mcq\_)}}\quad\textit{functoid}\index{enc-fun(mcp_, mcq_)@\texttt{enc-fun(mcp\_, mcq\_)}}\\
reads: the encoded functions from \$1 to \$2 \\
definition:
\begin{quote}\ttfamily\footnotesize\raggedright\noindent
eapp(\allowbreak{}19, [mcp\_, mcq\_])
\end{quote}
mentioned by 2 result(s)

\smallskip

\noindent\textbf{\texttt{enc-fun1(mdp\_)}}\quad\textit{functoid}\index{enc-fun1(mdp_)@\texttt{enc-fun1(mdp\_)}}\\
reads: the encoded functions with domain \$1 \\
definition:
\begin{quote}\ttfamily\footnotesize\raggedright\noindent
eapp(\allowbreak{}19, [mdp\_])
\end{quote}
characterized by \texttt{enc-fun1-unfold}  \textit{(proven)}:
\begin{quote}\ttfamily\footnotesize\raggedright\noindent
forall(\allowbreak{}[mip\_], enc-fun1(\allowbreak{}mip\_) == eapp(\allowbreak{}19, [mip\_]))
\end{quote}
mentioned by 6 result(s)

\smallskip

\noindent\textbf{\texttt{enc-if(mcp\_, mca\_, mcb\_)}}\quad\textit{functoid}\index{enc-if(mcp_, mca_, mcb_)@\texttt{enc-if(mcp\_, mca\_, mcb\_)}}\\
reads: the encoded conditional term on \$1 \\
definition:
\begin{quote}\ttfamily\footnotesize\raggedright\noindent
ebind(\allowbreak{}27, 0, [mca\_, mcb\_], mcp\_)
\end{quote}
mentioned by 14 result(s)

\smallskip

\noindent\textbf{\texttt{enc-iff(mcp\_, mcq\_)}}\quad\textit{functoid}\index{enc-iff(mcp_, mcq_)@\texttt{enc-iff(mcp\_, mcq\_)}}\\
reads: the encoded equivalence of \$1 and \$2 \\
definition:
\begin{quote}\ttfamily\footnotesize\raggedright\noindent
eapp(\allowbreak{}6, [mcp\_, mcq\_])
\end{quote}
characterized by \texttt{enc-iff-unfold}  \textit{(proven)}:
\begin{quote}\ttfamily\footnotesize\raggedright\noindent
forall(\allowbreak{}[p, q], enc-iff(\allowbreak{}p, q) == eapp(\allowbreak{}6, [p, q]))
\end{quote}
mentioned by 76 result(s)

\smallskip

\noindent\textbf{\texttt{enc-imp(mcp\_, mcq\_)}}\quad\textit{functoid}\index{enc-imp(mcp_, mcq_)@\texttt{enc-imp(mcp\_, mcq\_)}}\\
reads: the encoded implication from \$1 to \$2 \\
definition:
\begin{quote}\ttfamily\footnotesize\raggedright\noindent
eapp(\allowbreak{}5, [mcp\_, mcq\_])
\end{quote}
characterized by \texttt{enc-imp-unfold}  \textit{(proven)}:
\begin{quote}\ttfamily\footnotesize\raggedright\noindent
forall(\allowbreak{}[p, q], enc-imp(\allowbreak{}p, q) == eapp(\allowbreak{}5, [p, q]))
\end{quote}
mentioned by 60 result(s)

\smallskip

\noindent\textbf{\texttt{enc-in(mcp\_, mcq\_)}}\quad\textit{functoid}\index{enc-in(mcp_, mcq_)@\texttt{enc-in(mcp\_, mcq\_)}}\\
reads: the encoded membership of \$1 in \$2 \\
definition:
\begin{quote}\ttfamily\footnotesize\raggedright\noindent
eapp(\allowbreak{}7, [mcp\_, mcq\_])
\end{quote}
characterized by \texttt{enc-in-unfold}  \textit{(proven)}:
\begin{quote}\ttfamily\footnotesize\raggedright\noindent
forall(\allowbreak{}[p, q], enc-in(\allowbreak{}p, q) == eapp(\allowbreak{}7, [p, q]))
\end{quote}
mentioned by 168 result(s)

\smallskip

\noindent\textbf{\texttt{enc-inter(mcl\_)}}\quad\textit{functoid}\index{enc-inter(mcl_)@\texttt{enc-inter(mcl\_)}}\\
reads: the encoded intersection of the list \$1 \\
definition:
\begin{quote}\ttfamily\footnotesize\raggedright\noindent
eapp(\allowbreak{}12, mcl\_)
\end{quote}
characterized by \texttt{enc-inter-unfold}  \textit{(proven)}:
\begin{quote}\ttfamily\footnotesize\raggedright\noindent
forall(\allowbreak{}[l], enc-inter(\allowbreak{}l) == eapp(\allowbreak{}12, l))
\end{quote}
mentioned by 11 result(s)

\smallskip

\noindent\textbf{\texttt{enc-iota(mcb\_)}}\quad\textit{functoid}\index{enc-iota(mcb_)@\texttt{enc-iota(mcb\_)}}\\
reads: the encoded description by \$1 \\
definition:
\begin{quote}\ttfamily\footnotesize\raggedright\noindent
ebind(\allowbreak{}25, 1, [], mcb\_)
\end{quote}
characterized by \texttt{enc-iota-unfold}  \textit{(proven)}:
\begin{quote}\ttfamily\footnotesize\raggedright\noindent
forall(\allowbreak{}[mrp\_], enc-iota(\allowbreak{}mrp\_) == ebind(\allowbreak{}25, 1, [], mrp\_))
\end{quote}
mentioned by 14 result(s)

\smallskip

\noindent\textbf{\texttt{enc-lam(mca\_, mcb\_)}}\quad\textit{functoid}\index{enc-lam(mca_, mcb_)@\texttt{enc-lam(mca\_, mcb\_)}}\\
reads: the encoded lambda over \$1 with body \$2 \\
definition:
\begin{quote}\ttfamily\footnotesize\raggedright\noindent
ebind(\allowbreak{}26, 1, [mca\_], mcb\_)
\end{quote}
mentioned by 4 result(s)

\smallskip

\noindent\textbf{\texttt{enc-le(md10x\_, md10y\_)}}\quad\textit{functoid}\index{enc-le(md10x_, md10y_)@\texttt{enc-le(md10x\_, md10y\_)}}\\
reads: the encoded atom \$1 <= \$2 \\
definition:
\begin{quote}\ttfamily\footnotesize\raggedright\noindent
eapp(\allowbreak{}10, [econst(\allowbreak{}11), md10x\_, md10y\_])
\end{quote}
characterized by \texttt{enc-le-unfold}  \textit{(proven)}:
\begin{quote}\ttfamily\footnotesize\raggedright\noindent
forall(\allowbreak{}[mip\_, miq\_], enc-le(\allowbreak{}mip\_, miq\_) == eapp(\allowbreak{}10, [econst(\allowbreak{}11), mip\_, miq\_]))
\end{quote}
mentioned by 3 result(s)

\smallskip

\noindent\textbf{\texttt{enc-length(mcp\_)}}\quad\textit{functoid}\index{enc-length(mcp_)@\texttt{enc-length(mcp\_)}}\\
reads: the encoded length of \$1 \\
definition:
\begin{quote}\ttfamily\footnotesize\raggedright\noindent
eapp(\allowbreak{}16, [mcp\_])
\end{quote}
characterized by \texttt{enc-length-unfold}  \textit{(proven)}:
\begin{quote}\ttfamily\footnotesize\raggedright\noindent
forall(\allowbreak{}[mip\_], enc-length(\allowbreak{}mip\_) == eapp(\allowbreak{}16, [mip\_]))
\end{quote}
mentioned by 6 result(s)

\smallskip

\noindent\textbf{\texttt{enc-limit(mip\_)}}\quad\textit{functoid}\index{enc-limit(mip_)@\texttt{enc-limit(mip\_)}}\\
reads: the encoded LIMIT-ORD(\$1) \\
definition:
\begin{quote}\ttfamily\footnotesize\raggedright\noindent
eapp(\allowbreak{}10, [econst(\allowbreak{}9), mip\_])
\end{quote}
characterized by \texttt{enc-limit-unfold}  \textit{(proven)}:
\begin{quote}\ttfamily\footnotesize\raggedright\noindent
forall(\allowbreak{}[mip\_], enc-limit(\allowbreak{}mip\_) == eapp(\allowbreak{}10, [econst(\allowbreak{}9), mip\_]))
\end{quote}
mentioned by 8 result(s)

\smallskip

\noindent\textbf{\texttt{enc-list(mcl\_)}}\quad\textit{functoid}\index{enc-list(mcl_)@\texttt{enc-list(mcl\_)}}\\
reads: the encoded list of the entries of \$1 \\
definition:
\begin{quote}\ttfamily\footnotesize\raggedright\noindent
eapp(\allowbreak{}17, mcl\_)
\end{quote}
characterized by \texttt{enc-list-unfold}  \textit{(proven)}:
\begin{quote}\ttfamily\footnotesize\raggedright\noindent
forall(\allowbreak{}[l], enc-list(\allowbreak{}l) == eapp(\allowbreak{}17, l))
\end{quote}
mentioned by 34 result(s)

\smallskip

\noindent\textbf{\texttt{enc-lt(md10x\_, md10y\_)}}\quad\textit{functoid}\index{enc-lt(md10x_, md10y_)@\texttt{enc-lt(md10x\_, md10y\_)}}\\
reads: the encoded atom \$1 < \$2 \\
definition:
\begin{quote}\ttfamily\footnotesize\raggedright\noindent
eapp(\allowbreak{}10, [econst(\allowbreak{}13), md10x\_, md10y\_])
\end{quote}
characterized by \texttt{enc-lt-unfold}  \textit{(proven)}:
\begin{quote}\ttfamily\footnotesize\raggedright\noindent
forall(\allowbreak{}[mip\_, miq\_], enc-lt(\allowbreak{}mip\_, miq\_) == eapp(\allowbreak{}10, [econst(\allowbreak{}13), mip\_, miq\_]))
\end{quote}
mentioned by 3 result(s)

\smallskip

\noindent\textbf{\texttt{enc-nnc}}\quad\textit{constant}\index{enc-nnc@\texttt{enc-nnc}}\\
reads: the encoded class NN \\
defined by \texttt{enc-nnc-def}

\smallskip

\noindent\textbf{\texttt{enc-not(mcp\_)}}\quad\textit{functoid}\index{enc-not(mcp_)@\texttt{enc-not(mcp\_)}}\\
reads: the encoded negation of \$1 \\
definition:
\begin{quote}\ttfamily\footnotesize\raggedright\noindent
eapp(\allowbreak{}2, [mcp\_])
\end{quote}
characterized by \texttt{enc-not-unfold}  \textit{(proven)}:
\begin{quote}\ttfamily\footnotesize\raggedright\noindent
forall(\allowbreak{}[p], enc-not(\allowbreak{}p) == eapp(\allowbreak{}2, [p]))
\end{quote}
mentioned by 26 result(s)

\smallskip

\noindent\textbf{\texttt{enc-nth(mcp\_, mcq\_)}}\quad\textit{functoid}\index{enc-nth(mcp_, mcq_)@\texttt{enc-nth(mcp\_, mcq\_)}}\\
reads: the encoded \$1-th entry of \$2 \\
definition:
\begin{quote}\ttfamily\footnotesize\raggedright\noindent
eapp(\allowbreak{}15, [mcp\_, mcq\_])
\end{quote}
characterized by \texttt{enc-nth-unfold}  \textit{(proven)}:
\begin{quote}\ttfamily\footnotesize\raggedright\noindent
forall(\allowbreak{}[mdp7\_, mdq7\_], enc-nth(\allowbreak{}mdp7\_, mdq7\_) == eapp(\allowbreak{}15, [mdp7\_, mdq7\_]))
\end{quote}
mentioned by 17 result(s)

\smallskip

\noindent\textbf{\texttt{enc-num(mck\_)}}\quad\textit{functoid}\index{enc-num(mck_)@\texttt{enc-num(mck\_)}}\\
reads: the encoded numeral \$1 \\
definition:
\begin{quote}\ttfamily\footnotesize\raggedright\noindent
econst(\allowbreak{}2 * mck\_)
\end{quote}
characterized by \texttt{enc-num-unfold}  \textit{(proven)}:
\begin{quote}\ttfamily\footnotesize\raggedright\noindent
forall(\allowbreak{}[mik\_ in nn], enc-num(\allowbreak{}mik\_) = econst(\allowbreak{}2 * mik\_))
\end{quote}
mentioned by 35 result(s)

\smallskip

\noindent\textbf{\texttt{enc-or(mcp\_, mcq\_)}}\quad\textit{functoid}\index{enc-or(mcp_, mcq_)@\texttt{enc-or(mcp\_, mcq\_)}}\\
reads: the encoded disjunction of \$1 and \$2 \\
definition:
\begin{quote}\ttfamily\footnotesize\raggedright\noindent
eapp(\allowbreak{}4, [mcp\_, mcq\_])
\end{quote}
characterized by \texttt{enc-orl-succ}  \textit{(definitional)}:
\begin{quote}\ttfamily\footnotesize\raggedright\noindent
forall(\allowbreak{}[mcx\_, mcl\_, mcn\_ in nn], enc-orl(\allowbreak{}mcx\_, mcl\_, succ(\allowbreak{}mcn\_)) == enc-or(\allowbreak{}enc-orl(\allowbreak{}mcx\_, mcl\_, mcn\_), enc-in(\allowbreak{}mcx\_, nth(\allowbreak{}succ(\allowbreak{}mcn\_), mcl\_))))
\end{quote}
also: \texttt{enc-or-unfold} \\
mentioned by 25 result(s)

\smallskip

\noindent\textbf{\texttt{enc-ordc}}\quad\textit{constant}\index{enc-ordc@\texttt{enc-ordc}}\\
reads: the encoded class ORD \\
defined by \texttt{enc-ordc-def}

\smallskip

\noindent\textbf{\texttt{enc-ordlt(mip\_, miq\_)}}\quad\textit{functoid}\index{enc-ordlt(mip_, miq_)@\texttt{enc-ordlt(mip\_, miq\_)}}\\
reads: the encoded ordinal order \$1 < \$2 \\
definition:
\begin{quote}\ttfamily\footnotesize\raggedright\noindent
eapp(\allowbreak{}10, [econst(\allowbreak{}7), mip\_, miq\_])
\end{quote}
characterized by \texttt{enc-ordlt-unfold}  \textit{(proven)}:
\begin{quote}\ttfamily\footnotesize\raggedright\noindent
forall(\allowbreak{}[mip\_, miq\_], enc-ordlt(\allowbreak{}mip\_, miq\_) == eapp(\allowbreak{}10, [econst(\allowbreak{}7), mip\_, miq\_]))
\end{quote}
mentioned by 14 result(s)

\smallskip

\noindent\textbf{\texttt{enc-orl}}\quad\textit{functoid}\index{enc-orl@\texttt{enc-orl}}\\
reads: the disjunction of \$1 in the first \$3 entries of \$2 \\
defined by \texttt{enc-orl-zero} \texttt{enc-orl-succ} \\
characterized by \texttt{enc-orl-succ}  \textit{(definitional)}:
\begin{quote}\ttfamily\footnotesize\raggedright\noindent
forall(\allowbreak{}[mcx\_, mcl\_, mcn\_ in nn], enc-orl(\allowbreak{}mcx\_, mcl\_, succ(\allowbreak{}mcn\_)) == enc-or(\allowbreak{}enc-orl(\allowbreak{}mcx\_, mcl\_, mcn\_), enc-in(\allowbreak{}mcx\_, nth(\allowbreak{}succ(\allowbreak{}mcn\_), mcl\_))))
\end{quote}
also: \texttt{enc-orl-zero} \\
mentioned by 11 result(s)

\smallskip

\noindent\textbf{\texttt{enc-pred1(md10n\_, md10p\_)}}\quad\textit{functoid}\index{enc-pred1(md10n_, md10p_)@\texttt{enc-pred1(md10n\_, md10p\_)}}\\
reads: the encoded atom of the one-place predicate named \$1 at \$2 \\
definition:
\begin{quote}\ttfamily\footnotesize\raggedright\noindent
eapp(\allowbreak{}10, [md10n\_, md10p\_])
\end{quote}
characterized by \texttt{enc-pred1-unfold}  \textit{(proven)}:
\begin{quote}\ttfamily\footnotesize\raggedright\noindent
forall(\allowbreak{}[mip\_, miq\_], enc-pred1(\allowbreak{}mip\_, miq\_) == eapp(\allowbreak{}10, [mip\_, miq\_]))
\end{quote}
mentioned by 10 result(s)

\smallskip

\noindent\textbf{\texttt{enc-qeq(mcp\_, mcq\_)}}\quad\textit{functoid}\index{enc-qeq(mcp_, mcq_)@\texttt{enc-qeq(mcp\_, mcq\_)}}\\
reads: the encoded quasi-equation \$1 == \$2 \\
definition:
\begin{quote}\ttfamily\footnotesize\raggedright\noindent
eapp(\allowbreak{}9, [mcp\_, mcq\_])
\end{quote}
characterized by \texttt{enc-qeq-unfold}  \textit{(proven)}:
\begin{quote}\ttfamily\footnotesize\raggedright\noindent
forall(\allowbreak{}[p, q], enc-qeq(\allowbreak{}p, q) == eapp(\allowbreak{}9, [p, q]))
\end{quote}
mentioned by 40 result(s)

\smallskip

\noindent\textbf{\texttt{enc-sep(mca\_, mcb\_)}}\quad\textit{functoid}\index{enc-sep(mca_, mcb_)@\texttt{enc-sep(mca\_, mcb\_)}}\\
reads: the encoded separation of \$1 by \$2 \\
definition:
\begin{quote}\ttfamily\footnotesize\raggedright\noindent
ebind(\allowbreak{}22, 1, [mca\_], mcb\_)
\end{quote}
characterized by \texttt{enc-sep-unfold}  \textit{(proven)}:
\begin{quote}\ttfamily\footnotesize\raggedright\noindent
forall(\allowbreak{}[mia\_, mip\_], enc-sep(\allowbreak{}mia\_, mip\_) == ebind(\allowbreak{}22, 1, [mia\_], mip\_))
\end{quote}
mentioned by 19 result(s)

\smallskip

\noindent\textbf{\texttt{enc-setc}}\quad\textit{constant}\index{enc-setc@\texttt{enc-setc}}\\
reads: the encoded class SET \\
defined by \texttt{enc-setc-def}

\smallskip

\noindent\textbf{\texttt{enc-succ(mip\_)}}\quad\textit{functoid}\index{enc-succ(mip_)@\texttt{enc-succ(mip\_)}}\\
reads: the encoded successor of \$1 \\
definition:
\begin{quote}\ttfamily\footnotesize\raggedright\noindent
eapp(\allowbreak{}30, [mip\_])
\end{quote}
characterized by \texttt{enc-succ-unfold}  \textit{(proven)}:
\begin{quote}\ttfamily\footnotesize\raggedright\noindent
forall(\allowbreak{}[mip\_], enc-succ(\allowbreak{}mip\_) == eapp(\allowbreak{}30, [mip\_]))
\end{quote}
mentioned by 10 result(s)

\smallskip

\noindent\textbf{\texttt{enc-succo(mip\_)}}\quad\textit{functoid}\index{enc-succo(mip_)@\texttt{enc-succo(mip\_)}}\\
reads: the encoded ordinal successor of \$1 \\
definition:
\begin{quote}\ttfamily\footnotesize\raggedright\noindent
eapp(\allowbreak{}31, [mip\_])
\end{quote}
characterized by \texttt{enc-succo-unfold}  \textit{(proven)}:
\begin{quote}\ttfamily\footnotesize\raggedright\noindent
forall(\allowbreak{}[mip\_], enc-succo(\allowbreak{}mip\_) == eapp(\allowbreak{}31, [mip\_]))
\end{quote}
mentioned by 10 result(s)

\smallskip

\noindent\textbf{\texttt{enc-truth}}\quad\textit{constant}\index{enc-truth@\texttt{enc-truth}}\\
reads: the encoded TRUTH \\
defined by \texttt{enc-truth-def}

\smallskip

\noindent\textbf{\texttt{enc-tuples(mcp\_)}}\quad\textit{functoid}\index{enc-tuples(mcp_)@\texttt{enc-tuples(mcp\_)}}\\
reads: the encoded tuples over \$1 \\
definition:
\begin{quote}\ttfamily\footnotesize\raggedright\noindent
eapp(\allowbreak{}14, [mcp\_])
\end{quote}
characterized by \texttt{enc-tuples-unfold}  \textit{(proven)}:
\begin{quote}\ttfamily\footnotesize\raggedright\noindent
forall(\allowbreak{}[p], enc-tuples(\allowbreak{}p) == eapp(\allowbreak{}14, [p]))
\end{quote}
mentioned by 11 result(s)

\smallskip

\noindent\textbf{\texttt{enc-union(mcl\_)}}\quad\textit{functoid}\index{enc-union(mcl_)@\texttt{enc-union(mcl\_)}}\\
reads: the encoded union of the list \$1 \\
definition:
\begin{quote}\ttfamily\footnotesize\raggedright\noindent
eapp(\allowbreak{}11, mcl\_)
\end{quote}
characterized by \texttt{enc-union-unfold}  \textit{(proven)}:
\begin{quote}\ttfamily\footnotesize\raggedright\noindent
forall(\allowbreak{}[l], enc-union(\allowbreak{}l) == eapp(\allowbreak{}11, l))
\end{quote}
mentioned by 12 result(s)

\smallskip

\noindent\textbf{\texttt{encr(smg\_, smk\_)}}\quad\textit{functoid}\index{encr(smg_, smk_)@\texttt{encr(smg\_, smk\_)}}\\
reads: the encryption of \$1 under the key \$2 \\
definition:
\begin{quote}\ttfamily\footnotesize\raggedright\noindent
[2, smg\_, smk\_]
\end{quote}
characterized by \texttt{is-encrypt-trace}  \textit{(definitional)}:
\begin{quote}\ttfamily\footnotesize\raggedright\noindent
forall(\allowbreak{}[sasp\_, sas\_], is-encrypt-trace(\allowbreak{}sasp\_, sas\_) iff forsome(\allowbreak{}[sak\_ in keys(\allowbreak{}sasp\_), sag\_ in msg-set(\allowbreak{}tokens(\allowbreak{}sasp\_), keys(\allowbreak{}sasp\_))], sas\_ = [[received, sak\_], [received, sag\_], [sent, encr(\allowbreak{}sag\_, sak\_)]]))
\end{quote}
also: \texttt{encr-unfold} \texttt{lst-ne-encr} \texttt{encr-in-msgs} \texttt{encr-ne-list} \texttt{subterm-encr} \texttt{encr-ne-token} \texttt{msg-rank-encr} \texttt{encr-injective} ... (25 in all) \\
mentioned by 117 result(s)

\smallskip

\noindent\textbf{\texttt{entry(m, i, j)}}\quad\textit{functoid}\index{entry(m, i, j)@\texttt{entry(m, i, j)}}\\
reads: the entry of \$1 at row \$2, column \$3 \\
definition:
\begin{quote}\ttfamily\footnotesize\raggedright\noindent
nth(\allowbreak{}j, nth(\allowbreak{}i, m))
\end{quote}
characterized by \texttt{matmul-entry-uba-as-ring}  \textit{(definitional)}:
\begin{quote}\ttfamily\footnotesize\raggedright\noindent
forall(\allowbreak{}[r], is-unital-banach-algebra(\allowbreak{}r) implies forall(\allowbreak{}[m, n, k, p, q], p in mat(\allowbreak{}m, n, vec(\allowbreak{}r)) implies q in mat(\allowbreak{}n, k, vec(\allowbreak{}r)) implies 1 <= n implies forall(\allowbreak{}[i in interval(\allowbreak{}1, m), c in interval(\allowbreak{}1, k)], entry(\allowbreak{}matmul(\allowbreak{}uba-as-ring(\allowbreak{}r), p, q), i, c) = finsum(\allowbreak{}ring-additive-ag(\allowbreak{}uba-as-ring(\allowbreak{}r)), vnb-lambda(\allowbreak{}j, interval(\allowbreak{}1, n), (\allowbreak{}vmul(\allowbreak{}r))(\allowbreak{}entry(\allowbreak{}p, i, j), entry(\allowbreak{}q, j, c))), interval(\allowbreak{}1, n)))))
\end{quote}
also: \texttt{char-mat-entry-uba-as-ring} \texttt{matact-entry-cna-as-module} \texttt{matact-entry-sip-as-module} \texttt{matact-entry-uba-as-module} \texttt{companion-entry-uba-as-ring} \texttt{matact-entry-cnvs-as-module} \texttt{triple-entry-left-uba-as-ring} \texttt{triple-entry-right-uba-as-ring} ... (107 in all) \\
mentioned by 1285 result(s)

\smallskip

\noindent\textbf{\texttt{enum-fam(ag, f, phi, n)}}\quad\textit{functoid}\index{enum-fam(ag, f, phi, n)@\texttt{enum-fam(ag, f, phi, n)}}\\
definition:
\begin{quote}\ttfamily\footnotesize\raggedright\noindent
vnb-lambda(\allowbreak{}i, nn, if(\allowbreak{}i in ord-segment(\allowbreak{}n), f(\allowbreak{}phi(\allowbreak{}i)), iden(\allowbreak{}ag)))
\end{quote}
characterized by \texttt{enum-fam-value}  \textit{(proven)}:
\begin{quote}\ttfamily\footnotesize\raggedright\noindent
forall(\allowbreak{}[n in nn, ag, u, phi, i in ord-segment(\allowbreak{}n)], (\allowbreak{}enum-fam(\allowbreak{}ag, u, phi, n))(\allowbreak{}i) == u(\allowbreak{}phi(\allowbreak{}i)))
\end{quote}
also: \texttt{enum-fam-in-fun} \texttt{enum-fam-comm-monoid-in-fun} \texttt{metric-group-group@enum-fam} \\
mentioned by 12 result(s)

\smallskip

\noindent\textbf{\texttt{ep-norm(omega, ca, mu, epp\_, epf\_)}}\quad\textit{functoid}\index{ep-norm(omega, ca, mu, epp_, epf_)@\texttt{ep-norm(omega, ca, mu, epp\_, epf\_)}}\\
reads: the \$4-th power integral norm of \$5 over \$1 with respect to \$3 \\
definition:
\begin{quote}\ttfamily\footnotesize\raggedright\noindent
epow(\allowbreak{}integral(\allowbreak{}omega, ca, mu, ptwise-epow(\allowbreak{}omega, epf\_, epp\_)), recip(\allowbreak{}epp\_))
\end{quote}
characterized by \texttt{ep-norm-one}  \textit{(proven)}:
\begin{quote}\ttfamily\footnotesize\raggedright\noindent
forall(\allowbreak{}[omega, ca, mu, h], is-measure(\allowbreak{}omega, ca, mu) implies is-measurable-fn(\allowbreak{}omega, ca, h) implies ep-norm(\allowbreak{}omega, ca, mu, 1, h) = integral(\allowbreak{}omega, ca, mu, h))
\end{quote}
also: \texttt{ep-norm-mono} \texttt{ep-norm-scale} \texttt{ep-norm-unfold} \texttt{ep-norm-zero-null} \\
mentioned by 12 result(s)

\smallskip

\noindent\textbf{\texttt{eplus}}\quad\textit{operator}\index{eplus@\texttt{eplus}}\\
a term-forming head \\
characterized by \texttt{xsi-eplus}  \textit{(proven)}:
\begin{quote}\ttfamily\footnotesize\raggedright\noindent
forall(\allowbreak{}[xs, ca, mu, ys, cb, nu, f, g, x], is-measure(\allowbreak{}xs, ca, mu) implies is-measure(\allowbreak{}ys, cb, nu) implies is-measurable-fn(\allowbreak{}cartesian(\allowbreak{}xs, ys), product-sigma(\allowbreak{}xs, ca, ys, cb), f) implies is-measurable-fn(\allowbreak{}cartesian(\allowbreak{}xs, ys), product-sigma(\allowbreak{}xs, ca, ys, cb), g) implies x in xs implies (\allowbreak{}x-section-integral(\allowbreak{}xs, ys, cb, nu, ptwise-eplus(\allowbreak{}cartesian(\allowbreak{}xs, ys), f, g)))(\allowbreak{}x) = eplus(\allowbreak{}(\allowbreak{}x-section-integral(\allowbreak{}xs, ys, cb, nu, f))(\allowbreak{}x), (\allowbreak{}x-section-integral(\allowbreak{}xs, ys, cb, nu, g))(\allowbreak{}x)))
\end{quote}
also: \texttt{eplus-comm} \texttt{eplus-mono} \texttt{eplus-assoc} \texttt{eplus-swap4} \texttt{eplus-closed} \texttt{eplus-le-left} \texttt{epow-eplus-le} \texttt{esup-eplus-eq} ... (19 in all) \\
mentioned by 99 result(s)

\smallskip

\noindent\textbf{\texttt{eplus-down(edx\_, edy\_)}}\quad\textit{functoid}\index{eplus-down(edx_, edy_)@\texttt{eplus-down(edx\_, edy\_)}}\\
reads: the extended sum of \$1 and \$2 in [-inf, +inf) \\
definition:
\begin{quote}\ttfamily\footnotesize\raggedright\noindent
if(\allowbreak{}edx\_ = neg-inf or edy\_ = neg-inf, neg-inf, edx\_ + edy\_)
\end{quote}
characterized by \texttt{eplus-down-comm}  \textit{(proven)}:
\begin{quote}\ttfamily\footnotesize\raggedright\noindent
forall(\allowbreak{}[x in rr-ext-down, y in rr-ext-down], eplus-down(\allowbreak{}x, y) = eplus-down(\allowbreak{}y, x))
\end{quote}
also: \texttt{eplus-down-mono} \texttt{eplus-down-real} \texttt{eplus-down-assoc} \texttt{eplus-down-closed} \texttt{eplus-down-zero-left} \texttt{eplus-down-zero-right} \texttt{eplus-down-neg-inf-left} \texttt{eplus-down-neg-inf-right} \\
mentioned by 16 result(s)

\smallskip

\noindent\textbf{\texttt{eplus-up(eux\_, euy\_)}}\quad\textit{functoid}\index{eplus-up(eux_, euy_)@\texttt{eplus-up(eux\_, euy\_)}}\\
reads: the extended sum of \$1 and \$2 in (-inf, +inf] \\
definition:
\begin{quote}\ttfamily\footnotesize\raggedright\noindent
if(\allowbreak{}eux\_ = pos-inf or euy\_ = pos-inf, pos-inf, eux\_ + euy\_)
\end{quote}
characterized by \texttt{eplus-up-comm}  \textit{(proven)}:
\begin{quote}\ttfamily\footnotesize\raggedright\noindent
forall(\allowbreak{}[x in rr-ext-up, y in rr-ext-up], eplus-up(\allowbreak{}x, y) = eplus-up(\allowbreak{}y, x))
\end{quote}
also: \texttt{eplus-up-mono} \texttt{eplus-up-real} \texttt{eplus-up-assoc} \texttt{eplus-up-closed} \texttt{eplus-up-zero-left} \texttt{eplus-up-zero-right} \texttt{eplus-up-agrees-eplus} \texttt{eplus-up-pos-inf-left} ... (9 in all) \\
mentioned by 18 result(s)

\smallskip

\noindent\textbf{\texttt{epow(ept\_, epp\_)}}\quad\textit{functoid}\index{epow(ept_, epp_)@\texttt{epow(ept\_, epp\_)}}\\
reads: \$1 to the power \$2 in [0,+inf] \\
definition:
\begin{quote}\ttfamily\footnotesize\raggedright\noindent
if(\allowbreak{}ept\_ = pos-inf, pos-inf, rpow-star(\allowbreak{}ept\_, epp\_))
\end{quote}
characterized by \texttt{epow-in}  \textit{(proven)}:
\begin{quote}\ttfamily\footnotesize\raggedright\noindent
forall(\allowbreak{}[t\_ in rr-pos-star, p\_ in rr], 0 < p\_ implies epow(\allowbreak{}t\_, p\_) in rr-pos-star)
\end{quote}
also: \texttt{epow-one} \texttt{epow-epow} \texttt{epow-mono} \texttt{epow-real} \texttt{epow-zero} \texttt{epow-split} \texttt{epow-etimes} \texttt{epow-unfold} ... (24 in all) \\
mentioned by 47 result(s)

\smallskip

\noindent\textbf{\texttt{eps-rel(cuvs\_, cuve\_)}}\quad\textit{functoid}\index{eps-rel(cuvs_, cuve_)@\texttt{eps-rel(cuvs\_, cuve\_)}}\\
reads: the relation d(x, y) < \$2 on \$1 \\
definition:
\begin{quote}\ttfamily\footnotesize\raggedright\noindent
\{cuvp\_ in cartesian(\allowbreak{}pts(\allowbreak{}cuvs\_), pts(\allowbreak{}cuvs\_)): (\allowbreak{}dist(\allowbreak{}cuvs\_))(\allowbreak{}nth(\allowbreak{}1, cuvp\_), nth(\allowbreak{}2, cuvp\_)) < cuve\_\}
\end{quote}
characterized by \texttt{eps-rel-mem}  \textit{(proven)}:
\begin{quote}\ttfamily\footnotesize\raggedright\noindent
forall(\allowbreak{}[s, e, u, v], u in pts(\allowbreak{}s) implies v in pts(\allowbreak{}s) implies [u, v] in eps-rel(\allowbreak{}s, e) iff (\allowbreak{}dist(\allowbreak{}s))(\allowbreak{}u, v) < e)
\end{quote}
also: \texttt{eps-rel-set} \texttt{eps-rel-unfold} \texttt{eps-rel-equivalence} \\
mentioned by 8 result(s)

\smallskip

\noindent\textbf{\texttt{eps-setoid(cuvs\_, cuve\_)}}\quad\textit{functoid}\index{eps-setoid(cuvs_, cuve_)@\texttt{eps-setoid(cuvs\_, cuve\_)}}\\
reads: the points of \$1 under the relation d(x, y) < \$2 \\
definition:
\begin{quote}\ttfamily\footnotesize\raggedright\noindent
[pts(\allowbreak{}cuvs\_), eps-rel(\allowbreak{}cuvs\_, cuve\_)]
\end{quote}
characterized by \texttt{eps-setoid-pts}  \textit{(proven)}:
\begin{quote}\ttfamily\footnotesize\raggedright\noindent
forall(\allowbreak{}[s, e], is-metric-space(\allowbreak{}s) implies pts(\allowbreak{}eps-setoid(\allowbreak{}s, e)) == pts(\allowbreak{}s))
\end{quote}
also: \texttt{eps-setoid-rel} \texttt{eps-setoid-is-setoid} \\
mentioned by 9 result(s)

\smallskip

\noindent\textbf{\texttt{eqv(mrt\_)}}\quad\textit{functoid}\index{eqv(mrt_)@\texttt{eqv(mrt\_)}}\\
reads: the derivable equivalences over \$1 \\
definition:
\begin{quote}\ttfamily\footnotesize\raggedright\noindent
\{mrx\_ in rw-univ: eqv-ok(\allowbreak{}mrt\_, mrx\_)\}
\end{quote}
characterized by \texttt{eqv-ex}  \textit{(proven)}:
\begin{quote}\ttfamily\footnotesize\raggedright\noindent
forall(\allowbreak{}[jj, tt, x], rw-ex(\allowbreak{}jj, eqv(\allowbreak{}tt), x) implies x in rw-univ implies x in eqv(\allowbreak{}tt))
\end{quote}
also: \texttt{eqv-all} \texttt{eqv-not} \texttt{eqv-top} \texttt{eqv-or-l} \texttt{eqv-or-r} \texttt{eqv-pred} \texttt{eqv-tapp} \texttt{eqv-and-l} ... (29 in all) \\
mentioned by 66 result(s)

\smallskip

\noindent\textbf{\texttt{eqv-ok(mrt\_, mrx\_)}}\quad\textit{predicate}\index{eqv-ok(mrt_, mrx_)@\texttt{eqv-ok(mrt\_, mrx\_)}}\\
reads: \$2 relates two sentences equivalent, or two closed terms quasi-equal, over \$1 in its context \\
defined by \texttt{eqv-ok} \\
characterized by \texttt{eqv-ok}  \textit{(definitional)}:
\begin{quote}\ttfamily\footnotesize\raggedright\noindent
forall(\allowbreak{}[mrt\_, mrx\_], eqv-ok(\allowbreak{}mrt\_, mrx\_) iff forsome(\allowbreak{}[mrg\_, mre\_, mrf\_], mrx\_ = [2, mrg\_, mre\_, mrf\_] and mre\_ in clf and mrf\_ in clf and [mrg\_, enc-iff(\allowbreak{}mre\_, mrf\_)] in derivable(\allowbreak{}mrt\_) or mrx\_ = [1, mrg\_, mre\_, mrf\_] and mre\_ in clt and mrf\_ in clt and [mrg\_, enc-qeq(\allowbreak{}mre\_, mrf\_)] in derivable(\allowbreak{}mrt\_)))
\end{quote}
mentioned by 4 result(s)

\smallskip

\noindent\textbf{\texttt{esum(f\_)}}\quad\textit{functoid}\index{esum(f_)@\texttt{esum(f\_)}}\\
reads: the sum in [0,+inf] of the family \$1 \\
definition:
\begin{quote}\ttfamily\footnotesize\raggedright\noindent
esup(\allowbreak{}\{v\_ in rr-pos-star: forsome(\allowbreak{}[t\_], (\allowbreak{}t\_ in set and card(\allowbreak{}t\_) in nn and t\_ subset dom(\allowbreak{}f\_)) and v\_ = finsum(\allowbreak{}rr-pos-star-add-monoid, f\_, t\_))\})
\end{quote}
characterized by \texttt{esum-in}  \textit{(proven)}:
\begin{quote}\ttfamily\footnotesize\raggedright\noindent
forall(\allowbreak{}[f in fun(\allowbreak{}dom(\allowbreak{}f), rr-pos-star)], esum(\allowbreak{}f) in rr-pos-star)
\end{quote}
also: \texttt{esum-swap} \texttt{esum-least} \texttt{esum-scale} \texttt{esum-upper} \texttt{esum-nn-type} \texttt{esum-nn-zero} \texttt{esum-swap-le} \texttt{l2-dist-esum} ... (40 in all) \\
mentioned by 98 result(s)

\smallskip

\noindent\textbf{\texttt{esup(s\_)}}\quad\textit{functoid}\index{esup(s_)@\texttt{esup(s\_)}}\\
reads: the least upper bound in [0,+inf] of \$1 \\
definition:
\begin{quote}\ttfamily\footnotesize\raggedright\noindent
iota(\allowbreak{}bb\_, bb\_ in rr-pos-star and forall(\allowbreak{}[x in s\_], x <= bb\_) and forall(\allowbreak{}[b], b in rr-pos-star and forall(\allowbreak{}[x in s\_], x <= b) implies bb\_ <= b))
\end{quote}
characterized by \texttt{esup-in}  \textit{(proven)}:
\begin{quote}\ttfamily\footnotesize\raggedright\noindent
forall(\allowbreak{}[s], s subset rr-pos-star implies esup(\allowbreak{}s) in rr-pos-star)
\end{quote}
also: \texttt{esup-prop} \texttt{esup-empty} \texttt{esup-least} \texttt{esup-upper} \texttt{esup-eplus-eq} \texttt{esup-eplus-le} \texttt{esup-grid-max} \texttt{esup-etimes-eq} ... (18 in all) \\
mentioned by 69 result(s)

\smallskip

\noindent\textbf{\texttt{etail(f, n)}}\quad\textit{functoid}\index{etail(f, n)@\texttt{etail(f, n)}}\\
reads: the set of values of the sequence \$1 from index \$2 on \\
definition:
\begin{quote}\ttfamily\footnotesize\raggedright\noindent
image(\allowbreak{}f, \{p\_ in nn: n <= p\_\})
\end{quote}
characterized by \texttt{etail-unfold}  \textit{(proven)}:
\begin{quote}\ttfamily\footnotesize\raggedright\noindent
forall(\allowbreak{}[f, n], etail(\allowbreak{}f, n) == image(\allowbreak{}f, \{p\_ in nn: n <= p\_\}))
\end{quote}
also: \texttt{etail-subset-rr-pos-star} \\
mentioned by 4 result(s)

\smallskip

\noindent\textbf{\texttt{etimes(x, y)}}\quad\textit{operator}\index{etimes(x, y)@\texttt{etimes(x, y)}}\\
reads: \$1 times \$2 in the extended reals \\
a term-forming head \\
characterized by \texttt{epow-etimes}  \textit{(proven)}:
\begin{quote}\ttfamily\footnotesize\raggedright\noindent
forall(\allowbreak{}[c\_ in rr], 0 < c\_ implies forall(\allowbreak{}[u\_ in rr-pos-star, p\_ in rr], 0 < p\_ implies epow(\allowbreak{}etimes(\allowbreak{}c\_, u\_), p\_) = etimes(\allowbreak{}epow(\allowbreak{}c\_, p\_), epow(\allowbreak{}u\_, p\_))))
\end{quote}
also: \texttt{etimes-comm} \texttt{etimes-closed} \texttt{esup-etimes-eq} \texttt{esup-etimes-ge} \texttt{etimes-real-in} \texttt{etimes-one-left} \texttt{etimes-mono-left} \texttt{etimes-mono-star} ... (30 in all) \\
mentioned by 124 result(s)

\smallskip

\noindent\textbf{\texttt{euclidean-gauges(s)}}\quad\textit{functoid}\index{euclidean-gauges(s)@\texttt{euclidean-gauges(s)}}\\
definition:
\begin{quote}\ttfamily\footnotesize\raggedright\noindent
\{dg in fun(\allowbreak{}carr(\allowbreak{}s), nn): has-div-remainder(\allowbreak{}s, dg)\}
\end{quote}
characterized by \texttt{euclidean-gauges-unfold}  \textit{(proven)}:
\begin{quote}\ttfamily\footnotesize\raggedright\noindent
forall(\allowbreak{}[s], euclidean-gauges(\allowbreak{}s) == \{dg in fun(\allowbreak{}carr(\allowbreak{}s), nn): has-div-remainder(\allowbreak{}s, dg)\})
\end{quote}
mentioned by 5 result(s)

\smallskip

\noindent\textbf{\texttt{euclidean-ring}}\quad\textit{refinement}\index{euclidean-ring@\texttt{euclidean-ring}}\\
predicate \texttt{is-euclidean-ring}; same shape as integral-domain

\smallskip

\noindent\textbf{\texttt{euler-gamma}}\quad\textit{constant}\index{euler-gamma@\texttt{euler-gamma}}\\
reads: Euler's constant \\
defined by \texttt{euler-gamma-def}

\smallskip

\noindent\textbf{\texttt{euler-phi(n)}}\quad\textit{functoid}\index{euler-phi(n)@\texttt{euler-phi(n)}}\\
reads: Euler's totient of \$1 \\
definition:
\begin{quote}\ttfamily\footnotesize\raggedright\noindent
card(\allowbreak{}\{ephk\_ in interval(\allowbreak{}1, n): zz-coprime(\allowbreak{}ephk\_, n)\})
\end{quote}
characterized by \texttt{euler-phi-pos}  \textit{(proven)}:
\begin{quote}\ttfamily\footnotesize\raggedright\noindent
forall(\allowbreak{}[n in nn], 1 <= n implies euler-phi(\allowbreak{}n) in nn and 1 <= euler-phi(\allowbreak{}n))
\end{quote}
also: \texttt{euler-phi-product} \texttt{euler-phi-lower-bound} \texttt{euler-phi-prime-power} \texttt{euler-phi-ratio-limit} \texttt{euler-phi-ratio-bounded} \texttt{euler-phi-multiplicative} \texttt{euler-phi-ratio-prime-power} \texttt{euler-phi-ratio-multiplicative} \\
mentioned by 21 result(s)

\smallskip

\noindent\textbf{\texttt{ev-op(gev-md\_, gev-t\_)}}\quad\textit{functoid}\index{ev-op(gev-md_, gev-t_)@\texttt{ev-op(gev-md\_, gev-t\_)}}\\
reads: the evaluation p |-> p(\$2) \\
definition:
\begin{quote}\ttfamily\footnotesize\raggedright\noindent
vnb-lambda(\allowbreak{}gevp\_, carr(\allowbreak{}poly(\allowbreak{}scal(\allowbreak{}gev-md\_))), op-of(\allowbreak{}gev-md\_, gev-t\_, gevp\_))
\end{quote}
characterized by \texttt{ev-op-apply}  \textit{(proven)}:
\begin{quote}\ttfamily\footnotesize\raggedright\noindent
forall(\allowbreak{}[md, tt], is-hom-module(\allowbreak{}md, md, tt) implies is-commutative-ring(\allowbreak{}scal(\allowbreak{}md)) implies forall(\allowbreak{}[p in carr(\allowbreak{}poly(\allowbreak{}scal(\allowbreak{}md)))], (\allowbreak{}ev-op(\allowbreak{}md, tt))(\allowbreak{}p) = op-of(\allowbreak{}md, tt, p)))
\end{quote}
also: \texttt{ev-op-ring-hom} \\
mentioned by 5 result(s)

\smallskip

\noindent\textbf{\texttt{eval-at-points(a, th, m)}}\quad\textit{functoid}\index{eval-at-points(a, th, m)@\texttt{eval-at-points(a, th, m)}}\\
reads: the evaluation at the points \$2(1), ..., \$2(\$3 + 1) \\
definition:
\begin{quote}\ttfamily\footnotesize\raggedright\noindent
vnb-lambda(\allowbreak{}gep\_, poly-le(\allowbreak{}a, m), vnb-lambda(\allowbreak{}gei\_, interval(\allowbreak{}1, succ(\allowbreak{}m)), poly-eval(\allowbreak{}a, gep\_, th(\allowbreak{}gei\_))))
\end{quote}
characterized by \texttt{eval-at-points-bijective-iff-distinct}  \textit{(proven)}:
\begin{quote}\ttfamily\footnotesize\raggedright\noindent
forall(\allowbreak{}[fk], is-field-ring(\allowbreak{}fk) implies forall(\allowbreak{}[m\_ in nn, th\_], forall(\allowbreak{}[i\_ in interval(\allowbreak{}1, succ(\allowbreak{}m\_))], th\_(\allowbreak{}i\_) in carr(\allowbreak{}fk)) implies eval-at-points(\allowbreak{}fk, th\_, m\_) in bijection(\allowbreak{}poly-le(\allowbreak{}fk, m\_), fun(\allowbreak{}interval(\allowbreak{}1, succ(\allowbreak{}m\_)), carr(\allowbreak{}fk))) iff forall(\allowbreak{}[i\_ in interval(\allowbreak{}1, succ(\allowbreak{}m\_)), j\_ in interval(\allowbreak{}1, succ(\allowbreak{}m\_))], th\_(\allowbreak{}i\_) = th\_(\allowbreak{}j\_) implies i\_ = j\_)))
\end{quote}
mentioned by 2 result(s)

\smallskip

\noindent\textbf{\texttt{even(x)}}\quad\textit{predicate}\index{even(x)@\texttt{even(x)}}\\
reads: \$1 is even \\
defined by \texttt{even} \\
characterized by \texttt{even}  \textit{(definitional)}:
\begin{quote}\ttfamily\footnotesize\raggedright\noindent
forall(\allowbreak{}[x], even(\allowbreak{}x) iff forsome(\allowbreak{}[k in zz], x = 2 * k))
\end{quote}
also: \texttt{even-def} \texttt{zz-even-mul} \texttt{zz-even-iff-not-odd} \texttt{zz-not-even-and-odd} \\
mentioned by 9 result(s)

\smallskip

\noindent\textbf{\texttt{ew}}\quad\textit{tactic}\index{ew@\texttt{ew}}\\
\texttt{(ew term)} -- Existential witness: discharge a FORSOME goal by supplying the witness term. \\
declared in \texttt{tactics-help.scm}

\smallskip

\noindent\textbf{\texttt{ew-poly}}\quad\textit{tactic}\index{ew-poly@\texttt{ew-poly}}\\
\texttt{(ew-poly)} -- Close an existential goal forsome([a in C], L = R) (or unguarded) in which a occurs once, linearly: compute the witness by EXACT POLYNOMIAL DIVISION, supply it (ew), type it (type-term) and close the identity (crs).  One command for `(x+y)\textasciicircum{}3 = x\textasciicircum{}3 + y*a'.  Declines, with the reason and nothing changed, when the division leaves a remainder. \\
declared in \texttt{tactics-help.scm}

\smallskip

\noindent\textbf{\texttt{ex-confined(mnbb\_)}}\quad\textit{predicate}\index{ex-confined(mnbb_)@\texttt{ex-confined(mnbb\_)}}\\
reads: the existential with body \$1 is confined to sets \\
defined by \texttt{ex-confined} \\
characterized by \texttt{ex-confined}  \textit{(definitional)}:
\begin{quote}\ttfamily\footnotesize\raggedright\noindent
forall(\allowbreak{}[mnbb\_], ex-confined(\allowbreak{}mnbb\_) iff forsome(\allowbreak{}[mnbqd\_, mnbqc\_], mnbqd\_ in nn and mnbqc\_ in expr-set and [enc-in(\allowbreak{}ebvar(\allowbreak{}mnbqd\_), mnbqc\_), mnbqd\_] in ex-spine(\allowbreak{}mnbb\_)))
\end{quote}
also: \texttt{ex-confined-guarded} \\
mentioned by 5 result(s)

\smallskip

\noindent\textbf{\texttt{ex-spine(mnbb\_)}}\quad\textit{functoid}\index{ex-spine(mnbb_)@\texttt{ex-spine(mnbb\_)}}\\
reads: the conjunctive spine of the body \$1, with depths \\
definition:
\begin{quote}\ttfamily\footnotesize\raggedright\noindent
\{mnbz\_ in cartesian(\allowbreak{}expr-set, nn): forall(\allowbreak{}[mnbv\_ in power(\allowbreak{}cartesian(\allowbreak{}expr-set, nn))], [mnbb\_, 0] in mnbv\_ implies ex-spine-closed(\allowbreak{}mnbv\_) implies mnbz\_ in mnbv\_)\}
\end{quote}
characterized by \texttt{ex-spine-unfold}  \textit{(proven)}:
\begin{quote}\ttfamily\footnotesize\raggedright\noindent
forall(\allowbreak{}[mnbb\_], ex-spine(\allowbreak{}mnbb\_) == \{mnbz\_ in cartesian(\allowbreak{}expr-set, nn): forall(\allowbreak{}[mnbv\_ in power(\allowbreak{}cartesian(\allowbreak{}expr-set, nn))], [mnbb\_, 0] in mnbv\_ implies ex-spine-closed(\allowbreak{}mnbv\_) implies mnbz\_ in mnbv\_)\})
\end{quote}
mentioned by 4 result(s)

\smallskip

\noindent\textbf{\texttt{ex-spine-closed(mnbw\_)}}\quad\textit{predicate}\index{ex-spine-closed(mnbw_)@\texttt{ex-spine-closed(mnbw\_)}}\\
reads: \$1 is closed under steps down conjunctive spines \\
defined by \texttt{ex-spine-closed} \\
characterized by \texttt{ex-spine-closed}  \textit{(definitional)}:
\begin{quote}\ttfamily\footnotesize\raggedright\noindent
forall(\allowbreak{}[mnbw\_], ex-spine-closed(\allowbreak{}mnbw\_) iff forall(\allowbreak{}[mnbce\_ in mnbw\_, mnbcf\_ in cartesian(\allowbreak{}expr-set, nn)], ex-spine-step(\allowbreak{}mnbce\_, mnbcf\_) implies mnbcf\_ in mnbw\_))
\end{quote}
mentioned by 4 result(s)

\smallskip

\noindent\textbf{\texttt{ex-spine-step(mnbu\_, mnbt\_)}}\quad\textit{predicate}\index{ex-spine-step(mnbu_, mnbt_)@\texttt{ex-spine-step(mnbu\_, mnbt\_)}}\\
reads: \$2 is one step down the conjunctive spine from \$1 \\
defined by \texttt{ex-spine-step} \\
characterized by \texttt{ex-spine-step}  \textit{(definitional)}:
\begin{quote}\ttfamily\footnotesize\raggedright\noindent
forall(\allowbreak{}[mnbu\_, mnbt\_], ex-spine-step(\allowbreak{}mnbu\_, mnbt\_) iff forsome(\allowbreak{}[mnbsc\_, mnbsd\_], mnbsc\_ in expr-set and mnbsd\_ in nn and mnbu\_ = [enc-ex(\allowbreak{}mnbsc\_), mnbsd\_] and mnbt\_ = [mnbsc\_, mnbsd\_ + 1]) or forsome(\allowbreak{}[mnbsx\_, mnbsy\_, mnbse\_], mnbsx\_ in expr-set and mnbsy\_ in expr-set and mnbse\_ in nn and mnbu\_ = [enc-and(\allowbreak{}mnbsx\_, mnbsy\_), mnbse\_] and (\allowbreak{}mnbt\_ = [mnbsx\_, mnbse\_] or mnbt\_ = [mnbsy\_, mnbse\_])))
\end{quote}
mentioned by 4 result(s)

\smallskip

\noindent\textbf{\texttt{exp}}\quad\textit{operator}\index{exp@\texttt{exp}}\\
a term-forming head

\smallskip

\noindent\textbf{\texttt{expand}}\quad\textit{tactic}\index{expand@\texttt{expand}}\\
\texttt{(expand)} -- SHOW the expansion as a step: rewrite each polynomial side of an equation in the goal (also inside an existential or a conjunction) to its sum-of-monomials normal form, highest degree first.  (x + y)\textasciicircum{}3 = x\textasciicircum{}3 + y*a becomes x\textasciicircum{}3 + 3*x\textasciicircum{}2*y + 3*x*y\textasciicircum{}2 + y\textasciicircum{}3 = x\textasciicircum{}3 + y*a.  Declines, with the reason and nothing changed, when no side changes. \\
declared in \texttt{tactics-help.scm}

\smallskip

\noindent\textbf{\texttt{expr-fv(mxe\_)}}\quad\textit{functoid}\index{expr-fv(mxe_)@\texttt{expr-fv(mxe\_)}}\\
reads: the free names of \$1 \\
definition:
\begin{quote}\ttfamily\footnotesize\raggedright\noindent
(\allowbreak{}fv-stage(\allowbreak{}expr-rank(\allowbreak{}mxe\_)))(\allowbreak{}mxe\_)
\end{quote}
characterized by \texttt{expr-fv-at}  \textit{(proven)}:
\begin{quote}\ttfamily\footnotesize\raggedright\noindent
forall(\allowbreak{}[j, e], j in nn implies e in expr-stage(\allowbreak{}j) implies expr-fv(\allowbreak{}e) == (\allowbreak{}fv-stage(\allowbreak{}j))(\allowbreak{}e))
\end{quote}
mentioned by 58 result(s)

\smallskip

\noindent\textbf{\texttt{expr-lc(mxe\_)}}\quad\textit{predicate}\index{expr-lc(mxe_)@\texttt{expr-lc(mxe\_)}}\\
reads: \$1 is locally closed \\
defined by \texttt{expr-lc} \\
characterized by \texttt{expr-lc}  \textit{(definitional)}:
\begin{quote}\ttfamily\footnotesize\raggedright\noindent
forall(\allowbreak{}[mxe\_], expr-lc(\allowbreak{}mxe\_) iff 0 in expr-lcd(\allowbreak{}mxe\_))
\end{quote}
mentioned by 18 result(s)

\smallskip

\noindent\textbf{\texttt{expr-lcd(mxe\_)}}\quad\textit{functoid}\index{expr-lcd(mxe_)@\texttt{expr-lcd(mxe\_)}}\\
reads: the depths at which \$1 is closed \\
definition:
\begin{quote}\ttfamily\footnotesize\raggedright\noindent
(\allowbreak{}lcd-stage(\allowbreak{}expr-rank(\allowbreak{}mxe\_)))(\allowbreak{}mxe\_)
\end{quote}
characterized by \texttt{expr-lcd-at}  \textit{(proven)}:
\begin{quote}\ttfamily\footnotesize\raggedright\noindent
forall(\allowbreak{}[j, e], j in nn implies e in expr-stage(\allowbreak{}j) implies expr-lcd(\allowbreak{}e) == (\allowbreak{}lcd-stage(\allowbreak{}j))(\allowbreak{}e))
\end{quote}
mentioned by 18 result(s)

\smallskip

\noindent\textbf{\texttt{expr-open(mxe\_, mxk\_, mxu\_)}}\quad\textit{functoid}\index{expr-open(mxe_, mxk_, mxu_)@\texttt{expr-open(mxe\_, mxk\_, mxu\_)}}\\
reads: \$1 with the index \$2 opened by \$3 \\
definition:
\begin{quote}\ttfamily\footnotesize\raggedright\noindent
(\allowbreak{}open-stage(\allowbreak{}mxu\_, expr-rank(\allowbreak{}mxe\_)))(\allowbreak{}mxe\_, mxk\_)
\end{quote}
characterized by \texttt{expr-open-at}  \textit{(proven)}:
\begin{quote}\ttfamily\footnotesize\raggedright\noindent
forall(\allowbreak{}[u, j, e, ix], j in nn implies e in expr-stage(\allowbreak{}j) implies ix in nn implies expr-open(\allowbreak{}e, ix, u) == (\allowbreak{}open-stage(\allowbreak{}u, j))(\allowbreak{}e, ix))
\end{quote}
mentioned by 119 result(s)

\smallskip

\noindent\textbf{\texttt{expr-rank(mxe\_)}}\quad\textit{functoid}\index{expr-rank(mxe_)@\texttt{expr-rank(mxe\_)}}\\
reads: the rank of the expression \$1 \\
definition:
\begin{quote}\ttfamily\footnotesize\raggedright\noindent
iota(\allowbreak{}mxr\_, mxr\_ in nn and mxe\_ in expr-stage(\allowbreak{}mxr\_) and forall(\allowbreak{}[mxj\_ in nn], mxe\_ in expr-stage(\allowbreak{}mxj\_) implies mxr\_ <= mxj\_))
\end{quote}
characterized by \texttt{expr-rank-app}  \textit{(proven)}:
\begin{quote}\ttfamily\footnotesize\raggedright\noindent
forall(\allowbreak{}[h, l], h in nn implies l in tuples(\allowbreak{}expr-set) implies forsome(\allowbreak{}[k in nn], expr-rank(\allowbreak{}eapp(\allowbreak{}h, l)) = succ(\allowbreak{}k) and l in tuples(\allowbreak{}expr-stage(\allowbreak{}k)) and not(\allowbreak{}eapp(\allowbreak{}h, l) in expr-stage(\allowbreak{}k))))
\end{quote}
also: \texttt{expr-rank-bind} \texttt{expr-rank-bvar} \texttt{expr-rank-fvar} \texttt{expr-rank-prop} \texttt{expr-rank-const} \\
mentioned by 9 result(s)

\smallskip

\noindent\textbf{\texttt{expr-set}}\quad\textit{constant}\index{expr-set@\texttt{expr-set}}\\
reads: the expressions \\
defined by \texttt{expr-set-def}

\smallskip

\noindent\textbf{\texttt{expr-stage}}\quad\textit{functoid}\index{expr-stage@\texttt{expr-stage}}\\
reads: the stage \$1 of the expressions \\
defined by \texttt{expr-stage-zero} \texttt{expr-stage-succ} \\
characterized by \texttt{expr-stage-succ}  \textit{(definitional)}:
\begin{quote}\ttfamily\footnotesize\raggedright\noindent
forall(\allowbreak{}[mxn\_ in nn], expr-stage(\allowbreak{}succ(\allowbreak{}mxn\_)) == union(\allowbreak{}expr-stage(\allowbreak{}mxn\_), union(\allowbreak{}cartesian(\allowbreak{}pair(\allowbreak{}3, 3), nn, tuples(\allowbreak{}expr-stage(\allowbreak{}mxn\_))), cartesian(\allowbreak{}pair(\allowbreak{}4, 4), nn, nn, tuples(\allowbreak{}expr-stage(\allowbreak{}mxn\_)), expr-stage(\allowbreak{}mxn\_)))))
\end{quote}
also: \texttt{expr-stage-zero} \texttt{expr-stage-add} \texttt{expr-stage-app} \texttt{expr-stage-set} \texttt{expr-stage-bind} \texttt{expr-stage-bvar} \texttt{expr-stage-fvar} \texttt{expr-stage-keep} ... (11 in all) \\
mentioned by 77 result(s)

\smallskip

\noindent\textbf{\texttt{expr-subst(mxe\_, mxm\_, mxu\_)}}\quad\textit{functoid}\index{expr-subst(mxe_, mxm_, mxu_)@\texttt{expr-subst(mxe\_, mxm\_, mxu\_)}}\\
reads: \$1 with the free variable \$2 replaced by \$3 \\
definition:
\begin{quote}\ttfamily\footnotesize\raggedright\noindent
(\allowbreak{}subst-stage(\allowbreak{}mxm\_, mxu\_, expr-rank(\allowbreak{}mxe\_)))(\allowbreak{}mxe\_)
\end{quote}
characterized by \texttt{expr-subst-at}  \textit{(proven)}:
\begin{quote}\ttfamily\footnotesize\raggedright\noindent
forall(\allowbreak{}[m, u, j, e], j in nn implies e in expr-stage(\allowbreak{}j) implies expr-subst(\allowbreak{}e, m, u) == (\allowbreak{}subst-stage(\allowbreak{}m, u, j))(\allowbreak{}e))
\end{quote}
mentioned by 41 result(s)

\smallskip

\noindent\textbf{\texttt{expr-wks(mce\_)}}\quad\textit{functoid}\index{expr-wks(mce_)@\texttt{expr-wks(mce\_)}}\\
reads: the kinds of the expression \$1 \\
definition:
\begin{quote}\ttfamily\footnotesize\raggedright\noindent
(\allowbreak{}wks-stage(\allowbreak{}expr-rank(\allowbreak{}mce\_)))(\allowbreak{}mce\_)
\end{quote}
characterized by \texttt{expr-wks-at}  \textit{(proven)}:
\begin{quote}\ttfamily\footnotesize\raggedright\noindent
forall(\allowbreak{}[j, e], j in nn implies e in expr-stage(\allowbreak{}j) implies expr-wks(\allowbreak{}e) == (\allowbreak{}wks-stage(\allowbreak{}j))(\allowbreak{}e))
\end{quote}
mentioned by 28 result(s)

\smallskip

\noindent\textbf{\texttt{ext-abs(x4k\_, x4f\_, x4md\_, x4n\_, x4u\_, x4xi\_)}}\quad\textit{functoid}\index{ext-abs(x4k_, x4f_, x4md_, x4n_, x4u_, x4xi_)@\texttt{ext-abs(x4k\_, x4f\_, x4md\_, x4n\_, x4u\_, x4xi\_)}}\\
reads: |\$6|\_\$2 = |det M(\$6)| \\
definition:
\begin{quote}\ttfamily\footnotesize\raggedright\noindent
(\allowbreak{}fnrm(\allowbreak{}x4k\_))(\allowbreak{}ext-norm(\allowbreak{}x4k\_, x4f\_, x4md\_, x4n\_, x4u\_, x4xi\_))
\end{quote}
characterized by \texttt{ext-abs-mul}  \textit{(proven)}:
\begin{quote}\ttfamily\footnotesize\raggedright\noindent
forall(\allowbreak{}[k, f, md, n, u], is-fin-ext(\allowbreak{}k, f, md, n, u) implies forall(\allowbreak{}[xi in carr(\allowbreak{}f), eta in carr(\allowbreak{}f)], ext-abs(\allowbreak{}k, f, md, n, u, (\allowbreak{}mul(\allowbreak{}f))(\allowbreak{}xi, eta)) = ext-abs(\allowbreak{}k, f, md, n, u, xi) * ext-abs(\allowbreak{}k, f, md, n, u, eta)))
\end{quote}
also: \texttt{ext-abs-one} \texttt{ext-abs-type} \texttt{ext-abs-zero} \texttt{ext-abs-scalar} \texttt{ext-abs-nonzero} \\
mentioned by 13 result(s)

\smallskip

\noindent\textbf{\texttt{ext-abs-root(x4k\_, x4f\_, x4md\_, x4n\_, x4u\_, x4xi\_)}}\quad\textit{functoid}\index{ext-abs-root(x4k_, x4f_, x4md_, x4n_, x4u_, x4xi_)@\texttt{ext-abs-root(x4k\_, x4f\_, x4md\_, x4n\_, x4u\_, x4xi\_)}}\\
reads: |\$6|\_\$2\textasciicircum{}(1/\$4) \\
definition:
\begin{quote}\ttfamily\footnotesize\raggedright\noindent
rpow-star(\allowbreak{}ext-abs(\allowbreak{}x4k\_, x4f\_, x4md\_, x4n\_, x4u\_, x4xi\_), recip-star(\allowbreak{}x4n\_))
\end{quote}
characterized by \texttt{ext-abs-root-scalar}  \textit{(proven)}:
\begin{quote}\ttfamily\footnotesize\raggedright\noindent
forall(\allowbreak{}[k, f, md, n, u], is-fin-ext(\allowbreak{}k, f, md, n, u) implies forall(\allowbreak{}[x4l in carr(\allowbreak{}k)], ext-abs-root(\allowbreak{}k, f, md, n, u, (\allowbreak{}act(\allowbreak{}md))(\allowbreak{}x4l, one(\allowbreak{}f))) = (\allowbreak{}fnrm(\allowbreak{}k))(\allowbreak{}x4l)))
\end{quote}
mentioned by 2 result(s)

\smallskip

\noindent\textbf{\texttt{ext-add(eax\_, eay\_)}}\quad\textit{functoid}\index{ext-add(eax_, eay_)@\texttt{ext-add(eax\_, eay\_)}}\\
reads: the extended sum of \$1 and \$2 \\
definition:
\begin{quote}\ttfamily\footnotesize\raggedright\noindent
if(\allowbreak{}eax\_ = pos-inf or eay\_ = pos-inf, pos-inf, if(\allowbreak{}eax\_ = neg-inf or eay\_ = neg-inf, neg-inf, eax\_ + eay\_))
\end{quote}
characterized by \texttt{ext-add-in}  \textit{(proven)}:
\begin{quote}\ttfamily\footnotesize\raggedright\noindent
forall(\allowbreak{}[x, y], x in rr-star implies y in rr-star implies ext-add(\allowbreak{}x, y) in rr-star)
\end{quote}
also: \texttt{ext-add-real} \texttt{ext-add-ge-left} \texttt{ext-add-ge-right} \texttt{ext-add-le-eplus} \texttt{ext-add-zero-left} \texttt{ext-add-zero-right} \texttt{ext-add-neg-inf-left} \texttt{ext-add-pos-inf-left} ... (10 in all) \\
mentioned by 28 result(s)

\smallskip

\noindent\textbf{\texttt{ext-degree(x4k\_, x4f\_, x4md\_)}}\quad\textit{functoid}\index{ext-degree(x4k_, x4f_, x4md_)@\texttt{ext-degree(x4k\_, x4f\_, x4md\_)}}\\
reads: the degree [\$2 : \$1] \\
definition:
\begin{quote}\ttfamily\footnotesize\raggedright\noindent
iota(\allowbreak{}x4d\_, forsome(\allowbreak{}[x4b\_], is-fin-ext(\allowbreak{}x4k\_, x4f\_, x4md\_, x4d\_, x4b\_)))
\end{quote}

\smallskip

\noindent\textbf{\texttt{ext-neg(enx\_)}}\quad\textit{functoid}\index{ext-neg(enx_)@\texttt{ext-neg(enx\_)}}\\
reads: the negative of \$1 in the extended reals \\
definition:
\begin{quote}\ttfamily\footnotesize\raggedright\noindent
if(\allowbreak{}enx\_ = pos-inf, neg-inf, if(\allowbreak{}enx\_ = neg-inf, pos-inf, -enx\_))
\end{quote}
characterized by \texttt{ext-neg-in}  \textit{(proven)}:
\begin{quote}\ttfamily\footnotesize\raggedright\noindent
forall(\allowbreak{}[x in rr-star], ext-neg(\allowbreak{}x) in rr-star)
\end{quote}
also: \texttt{ext-neg-real} \texttt{ext-neg-zero} \texttt{ext-neg-add-up} \texttt{ext-neg-gt-flip} \texttt{ext-neg-lt-flip} \texttt{ext-neg-neg-inf} \texttt{ext-neg-pos-inf} \texttt{ext-neg-add-down} ... (11 in all) \\
mentioned by 27 result(s)

\smallskip

\noindent\textbf{\texttt{ext-norm(x4k\_, x4f\_, x4md\_, x4n\_, x4u\_, x4xi\_)}}\quad\textit{functoid}\index{ext-norm(x4k_, x4f_, x4md_, x4n_, x4u_, x4xi_)@\texttt{ext-norm(x4k\_, x4f\_, x4md\_, x4n\_, x4u\_, x4xi\_)}}\\
reads: the norm det M(\$6) of \$6 \\
definition:
\begin{quote}\ttfamily\footnotesize\raggedright\noindent
det(\allowbreak{}normed-field-as-commutative-ring(\allowbreak{}x4k\_), x4n\_, mult-mat(\allowbreak{}x4md\_, x4n\_, x4u\_, x4f\_, x4xi\_))
\end{quote}

\smallskip

\noindent\textbf{\texttt{ext-partial-sum(x, n)}}\quad\textit{functoid}\index{ext-partial-sum(x, n)@\texttt{ext-partial-sum(x, n)}}\\
reads: the \$2-th partial sum of the extended series \$1 \\
definition:
\begin{quote}\ttfamily\footnotesize\raggedright\noindent
if(\allowbreak{}forsome(\allowbreak{}[epk\_ in nn], epk\_ < n and x(\allowbreak{}epk\_) = pos-inf), pos-inf, if(\allowbreak{}forsome(\allowbreak{}[epk\_ in nn], epk\_ < n and x(\allowbreak{}epk\_) = neg-inf), neg-inf, series-partial-sum(\allowbreak{}x, n)))
\end{quote}
mentioned by 35 result(s)

\smallskip

\noindent\textbf{\texttt{ext-series-converges-to(x, l)}}\quad\textit{predicate}\index{ext-series-converges-to(x, l)@\texttt{ext-series-converges-to(x, l)}}\\
reads: the extended series \$1 converges to \$2 \\
defined by \texttt{ext-series-converges-to} \\
characterized by \texttt{ext-series-converges-to}  \textit{(definitional)}:
\begin{quote}\ttfamily\footnotesize\raggedright\noindent
forall(\allowbreak{}[x, l], ext-series-converges-to(\allowbreak{}x, l) iff l in rr-star and (\allowbreak{}l in rr implies forall(\allowbreak{}[ece\_], pos-rr(\allowbreak{}ece\_) implies forsome(\allowbreak{}[ecm\_ in nn], forall(\allowbreak{}[ecn\_ in nn], ecm\_ <= ecn\_ implies ext-partial-sum(\allowbreak{}x, ecn\_) in rr and abs(\allowbreak{}ext-partial-sum(\allowbreak{}x, ecn\_) - l) < ece\_)))) and (\allowbreak{}l = pos-inf implies forall(\allowbreak{}[ecb\_ in rr], forsome(\allowbreak{}[ecm\_ in nn], forall(\allowbreak{}[ecn\_ in nn], ecm\_ <= ecn\_ implies ecb\_ < ext-partial-sum(\allowbreak{}x, ecn\_))))) and (\allowbreak{}l = neg-inf implies forall(\allowbreak{}[ecb\_ in rr], forsome(\allowbreak{}[ecm\_ in nn], forall(\allowbreak{}[ecn\_ in nn], ecm\_ <= ecn\_ implies ext-partial-sum(\allowbreak{}x, ecn\_) < ecb\_)))))
\end{quote}
mentioned by 16 result(s)

\smallskip

\noindent\textbf{\texttt{ext-sub(esx\_, esy\_)}}\quad\textit{functoid}\index{ext-sub(esx_, esy_)@\texttt{ext-sub(esx\_, esy\_)}}\\
reads: the extended difference of \$1 and \$2 \\
definition:
\begin{quote}\ttfamily\footnotesize\raggedright\noindent
ext-add(\allowbreak{}esx\_, ext-neg(\allowbreak{}esy\_))
\end{quote}
mentioned by 3 result(s)

\smallskip

\noindent\textbf{\texttt{extend-const(f, a, b)}}\quad\textit{functoid}\index{extend-const(f, a, b)@\texttt{extend-const(f, a, b)}}\\
reads: \$1 extended to the whole line by constants outside \$2 to \$3 \\
definition:
\begin{quote}\ttfamily\footnotesize\raggedright\noindent
vnb-lambda(\allowbreak{}ecx\_, rr, f(\allowbreak{}clamp(\allowbreak{}a, b, ecx\_)))
\end{quote}
characterized by \texttt{extend-const-fixes}  \textit{(proven)}:
\begin{quote}\ttfamily\footnotesize\raggedright\noindent
forall(\allowbreak{}[a in rr, b in rr, icz\_ in rr], a <= icz\_ implies icz\_ <= b implies forall(\allowbreak{}[f], (\allowbreak{}extend-const(\allowbreak{}f, a, b))(\allowbreak{}icz\_) == f(\allowbreak{}icz\_)))
\end{quote}
also: \texttt{extend-const-value} \texttt{extend-const-in-fun} \texttt{extend-const-deriv-bwd} \texttt{extend-const-deriv-fwd} \texttt{extend-const-deriv-iff} \texttt{extend-const-ccint-fixes} \texttt{extend-const-continuous-at} \texttt{regulated-extend-const-is-regulated-on} \\
mentioned by 20 result(s)

\smallskip

\noindent\textbf{\texttt{extends-on(s, g, f)}}\quad\textit{predicate}\index{extends-on(s, g, f)@\texttt{extends-on(s, g, f)}}\\
reads: \$2 agrees with \$3 on \$1 \\
defined by \texttt{extends-on} \\
characterized by \texttt{extends-on}  \textit{(definitional)}:
\begin{quote}\ttfamily\footnotesize\raggedright\noindent
forall(\allowbreak{}[s, g, f], extends-on(\allowbreak{}s, g, f) iff forall(\allowbreak{}[x\_ in s], g(\allowbreak{}x\_) = f(\allowbreak{}x\_)))
\end{quote}
also: \texttt{extends-on-trans} \\
mentioned by 30 result(s)

\smallskip

\noindent\textbf{\texttt{extreme-points(m, ks)}}\quad\textit{functoid}\index{extreme-points(m, ks)@\texttt{extreme-points(m, ks)}}\\
reads: the set of extreme points of \$2 in \$1 \\
definition:
\begin{quote}\ttfamily\footnotesize\raggedright\noindent
\{kez\_ in vec(\allowbreak{}m): is-extreme-point(\allowbreak{}m, ks, kez\_)\}
\end{quote}
characterized by \texttt{extreme-points-mem}  \textit{(proven)}:
\begin{quote}\ttfamily\footnotesize\raggedright\noindent
forall(\allowbreak{}[m, ks, z], z in extreme-points(\allowbreak{}m, ks) iff z in vec(\allowbreak{}m) and is-extreme-point(\allowbreak{}m, ks, z))
\end{quote}
also: \texttt{extreme-points-unfold} \texttt{extreme-points-nonempty} \\
mentioned by 7 result(s)

\smallskip


\section*{F}
\addcontentsline{toc}{section}{F}

\noindent\textbf{\texttt{fact}}\quad\textit{tactic}\index{fact@\texttt{fact}}\\
\texttt{(fact 'thm term ...)} -- Forward APPLICATION of a theorem: bring it in, instantiate its leading universals with the terms, and auto-detach every antecedent already in context, landing the consequent as a hypothesis. \\
declared in \texttt{tactics-help.scm}

\smallskip

\noindent\textbf{\texttt{fact-sums(fsm\_, fsg\_, fsx\_)}}\quad\textit{functoid}\index{fact-sums(fsm_, fsg_, fsx_)@\texttt{fact-sums(fsm\_, fsg\_, fsx\_)}}\\
reads: the sums of \$2 over the factorisations of \$3 in \$1 \\
definition:
\begin{quote}\ttfamily\footnotesize\raggedright\noindent
\{fsr\_ in rr: forsome(\allowbreak{}[fsn\_ in nn, fsu\_ in fun(\allowbreak{}nn, carr(\allowbreak{}fsm\_))], mprod(\allowbreak{}fsm\_, fsu\_, fsn\_) = fsx\_ and fsr\_ = series-partial-sum(\allowbreak{}vnb-lambda(\allowbreak{}fsi\_, nn, fsg\_(\allowbreak{}fsu\_(\allowbreak{}fsi\_))), fsn\_))\}
\end{quote}
characterized by \texttt{fact-sums-intro}  \textit{(proven)}:
\begin{quote}\ttfamily\footnotesize\raggedright\noindent
forall(\allowbreak{}[m], is-monoid(\allowbreak{}m) implies forall(\allowbreak{}[g in fun(\allowbreak{}carr(\allowbreak{}m), rr)], forall(\allowbreak{}[x in carr(\allowbreak{}m)], 0 <= g(\allowbreak{}x)) implies g(\allowbreak{}iden(\allowbreak{}m)) = 0 implies forall(\allowbreak{}[x in carr(\allowbreak{}m), y in carr(\allowbreak{}m), z in carr(\allowbreak{}m)], g(\allowbreak{}(\allowbreak{}opr(\allowbreak{}m))(\allowbreak{}(\allowbreak{}opr(\allowbreak{}m))(\allowbreak{}x, y), z)) <= 2 * max(\allowbreak{}g(\allowbreak{}x), max(\allowbreak{}g(\allowbreak{}y), g(\allowbreak{}z)))) implies forall(\allowbreak{}[u in fun(\allowbreak{}nn, carr(\allowbreak{}m)), n in nn], series-partial-sum(\allowbreak{}vnb-lambda(\allowbreak{}msi\_, nn, g(\allowbreak{}u(\allowbreak{}msi\_))), n) in fact-sums(\allowbreak{}m, g, mprod(\allowbreak{}m, u, n)))))
\end{quote}
also: \texttt{fact-sums-concat} \texttt{fact-sums-nonneg} \texttt{fact-sums-unfold} \\
mentioned by 7 result(s)

\smallskip

\noindent\textbf{\texttt{factorial}}\quad\textit{defined-fn}\index{factorial@\texttt{factorial}}\\
defined by \texttt{factorial-zero} \texttt{factorial-succ} \\
characterized by \texttt{factorial-succ}  \textit{(definitional)}:
\begin{quote}\ttfamily\footnotesize\raggedright\noindent
forall(\allowbreak{}[k in nn], factorial(\allowbreak{}succ(\allowbreak{}k)) = succ(\allowbreak{}k) * factorial(\allowbreak{}k))
\end{quote}
also: \texttt{factorial-zero} \texttt{factorial-mono} \texttt{factorial-in-nn} \texttt{gamma-factorial} \texttt{factorial-real-pos} \texttt{rr-recip-factorial} \texttt{recip-factorial-in-rr} \texttt{choose-times-factorial} ... (11 in all) \\
mentioned by 120 result(s)

\smallskip

\noindent\textbf{\texttt{falling}}\quad\textit{defined-fn}\index{falling@\texttt{falling}}\\
defined by \texttt{falling-zero} \texttt{falling-succ} \\
characterized by \texttt{falling-succ}  \textit{(definitional)}:
\begin{quote}\ttfamily\footnotesize\raggedright\noindent
forall(\allowbreak{}[n, k in nn], falling(\allowbreak{}n, succ(\allowbreak{}k)) == (\allowbreak{}n - k) * falling(\allowbreak{}n, k))
\end{quote}
also: \texttt{falling-zero} \texttt{falling-in-nn} \texttt{falling-in-nn-ind} \texttt{falling-succ-succ} \texttt{falling-zero-succ} \texttt{rb5-count-falling-aux} \texttt{injection-count-falling} \\
mentioned by 18 result(s)

\smallskip

\noindent\textbf{\texttt{fam-of-list}}\quad\textit{defined-fn}\index{fam-of-list@\texttt{fam-of-list}}\\
defined by \texttt{fam-of-list-apply} \\
characterized by \texttt{fam-of-list-apply}  \textit{(definitional)}:
\begin{quote}\ttfamily\footnotesize\raggedright\noindent
forall(\allowbreak{}[l, i], i in nn and succ(\allowbreak{}i) <= length(\allowbreak{}l) implies (\allowbreak{}fam-of-list(\allowbreak{}l))(\allowbreak{}i) = nth(\allowbreak{}succ(\allowbreak{}i), l))
\end{quote}
mentioned by 18 result(s)

\smallskip

\noindent\textbf{\texttt{field}}\quad\textit{structure}\index{field@\texttt{field}}\\
predicate \texttt{is-field}; slots: carr add mul neg zero one non-zero mul-inv

\smallskip

\noindent\textbf{\texttt{field-additive-ag(s)}}\quad\textit{view}\index{field-additive-ag(s)@\texttt{field-additive-ag(s)}}\\
a field viewed as an abelian-group \\
characterized by \texttt{field-additive-ag-is-abelian-group}  \textit{(definitional)}:
\begin{quote}\ttfamily\footnotesize\raggedright\noindent
forall(\allowbreak{}[r], is-field(\allowbreak{}r) implies is-abelian-group(\allowbreak{}field-additive-ag(\allowbreak{}r)))
\end{quote}
also: \texttt{abelian-group-as-monoid-is-monoid-field-additive-ag} \texttt{field-additive-ag-functorial} \\
mentioned by 3 result(s)

\smallskip

\noindent\textbf{\texttt{field-as-euclidean-ring(s)}}\quad\textit{view}\index{field-as-euclidean-ring(s)@\texttt{field-as-euclidean-ring(s)}}\\
a field viewed as an euclidean-ring \\
characterized by \texttt{field-as-euclidean-ring-is-euclidean-ring}  \textit{(definitional)}:
\begin{quote}\ttfamily\footnotesize\raggedright\noindent
forall(\allowbreak{}[r], is-field(\allowbreak{}r) implies is-euclidean-ring(\allowbreak{}field-as-euclidean-ring(\allowbreak{}r)))
\end{quote}
also: \texttt{field-as-euclidean-ring-functorial} \\
mentioned by 2 result(s)

\smallskip

\noindent\textbf{\texttt{field-as-integral-domain(s)}}\quad\textit{view}\index{field-as-integral-domain(s)@\texttt{field-as-integral-domain(s)}}\\
a field viewed as an integral-domain \\
characterized by \texttt{field-as-integral-domain-is-integral-domain}  \textit{(definitional)}:
\begin{quote}\ttfamily\footnotesize\raggedright\noindent
forall(\allowbreak{}[r], is-field(\allowbreak{}r) implies is-integral-domain(\allowbreak{}field-as-integral-domain(\allowbreak{}r)))
\end{quote}
also: \texttt{field-as-integral-domain-functorial} \\
mentioned by 32 result(s)

\smallskip

\noindent\textbf{\texttt{field-multiplicative-group(s)}}\quad\textit{view}\index{field-multiplicative-group(s)@\texttt{field-multiplicative-group(s)}}\\
a field viewed as a group \\
characterized by \texttt{field-multiplicative-group-is-group}  \textit{(definitional)}:
\begin{quote}\ttfamily\footnotesize\raggedright\noindent
forall(\allowbreak{}[r], is-field(\allowbreak{}r) implies is-group(\allowbreak{}field-multiplicative-group(\allowbreak{}r)))
\end{quote}
mentioned by 1 result(s)

\smallskip

\noindent\textbf{\texttt{field-ring}}\quad\textit{refinement}\index{field-ring@\texttt{field-ring}}\\
predicate \texttt{is-field-ring}; same shape as commutative-ring

\smallskip

\noindent\textbf{\texttt{fin-enum(s)}}\quad\textit{functoid}\index{fin-enum(s)@\texttt{fin-enum(s)}}\\
definition:
\begin{quote}\ttfamily\footnotesize\raggedright\noindent
choice(\allowbreak{}bijection(\allowbreak{}ord-segment(\allowbreak{}card(\allowbreak{}s)), s))
\end{quote}
characterized by \texttt{fin-enum-is-bijection}  \textit{(proven)}:
\begin{quote}\ttfamily\footnotesize\raggedright\noindent
forall(\allowbreak{}[s in set], card(\allowbreak{}s) in nn implies fin-enum(\allowbreak{}s) in bijection(\allowbreak{}ord-segment(\allowbreak{}card(\allowbreak{}s)), s))
\end{quote}
mentioned by 3 result(s)

\smallskip

\noindent\textbf{\texttt{fin-subset-monoid(a\_)}}\quad\textit{functoid}\index{fin-subset-monoid(a_)@\texttt{fin-subset-monoid(a\_)}}\\
reads: the monoid of finite subsets of \$1 under union \\
definition:
\begin{quote}\ttfamily\footnotesize\raggedright\noindent
[fin-subsets(\allowbreak{}a\_), vnb-lambda(\allowbreak{}[x\_, y\_], cartesian(\allowbreak{}fin-subsets(\allowbreak{}a\_), fin-subsets(\allowbreak{}a\_)), union(\allowbreak{}x\_, y\_)), empty-set]
\end{quote}
characterized by \texttt{fin-subset-monoid-is-comm-monoid}  \textit{(proven)}:
\begin{quote}\ttfamily\footnotesize\raggedright\noindent
forall(\allowbreak{}[a\_ in set], is-comm-monoid(\allowbreak{}fin-subset-monoid(\allowbreak{}a\_)))
\end{quote}
mentioned by 7 result(s)

\smallskip

\noindent\textbf{\texttt{fin-subsets(a\_)}}\quad\textit{functoid}\index{fin-subsets(a_)@\texttt{fin-subsets(a\_)}}\\
reads: the finite subsets of \$1 \\
definition:
\begin{quote}\ttfamily\footnotesize\raggedright\noindent
\{t\_ in power(\allowbreak{}a\_): card(\allowbreak{}t\_) in nn\}
\end{quote}
characterized by \texttt{fin-subsets-is-set}  \textit{(proven)}:
\begin{quote}\ttfamily\footnotesize\raggedright\noindent
forall(\allowbreak{}[a\_ in set], fin-subsets(\allowbreak{}a\_) in set)
\end{quote}
also: \texttt{countable-fin-subsets} \texttt{fin-subsets-has-empty} \texttt{fin-subsets-membership} \texttt{fin-subsets-union-closed} \\
mentioned by 11 result(s)

\smallskip

\noindent\textbf{\texttt{finite-rank-ops(kqv\_, kqw\_)}}\quad\textit{functoid}\index{finite-rank-ops(kqv_, kqw_)@\texttt{finite-rank-ops(kqv\_, kqw\_)}}\\
reads: the finite-rank operators from \$1 to \$2 \\
definition:
\begin{quote}\ttfamily\footnotesize\raggedright\noindent
\{kqt\_ in bounded-ops(\allowbreak{}kqv\_, kqw\_): is-finite-rank-op(\allowbreak{}kqv\_, kqw\_, kqt\_)\}
\end{quote}
characterized by \texttt{finite-rank-ops-unfold}  \textit{(proven)}:
\begin{quote}\ttfamily\footnotesize\raggedright\noindent
forall(\allowbreak{}[kqv\_, kqw\_], finite-rank-ops(\allowbreak{}kqv\_, kqw\_) == \{kqt\_ in bounded-ops(\allowbreak{}kqv\_, kqw\_): is-finite-rank-op(\allowbreak{}kqv\_, kqw\_, kqt\_)\})
\end{quote}
also: \texttt{finite-rank-ops-membership} \\
mentioned by 4 result(s)

\smallskip

\noindent\textbf{\texttt{finite-unions(omega, ci)}}\quad\textit{functoid}\index{finite-unions(omega, ci)@\texttt{finite-unions(omega, ci)}}\\
reads: the finite unions of members of \$2 (subsets of \$1) \\
definition:
\begin{quote}\ttfamily\footnotesize\raggedright\noindent
\{mfu\_ in power(\allowbreak{}omega): forsome(\allowbreak{}[mff\_ in set], card(\allowbreak{}mff\_) in nn and mff\_ subset ci and mfu\_ = big-union(\allowbreak{}mfw\_, mff\_, mfw\_))\}
\end{quote}
characterized by \texttt{finite-unions-elim}  \textit{(proven)}:
\begin{quote}\ttfamily\footnotesize\raggedright\noindent
forall(\allowbreak{}[omega, ci, s in finite-unions(\allowbreak{}omega, ci)], s in power(\allowbreak{}omega) and forsome(\allowbreak{}[mff\_ in set], card(\allowbreak{}mff\_) in nn and mff\_ subset ci and s = big-union(\allowbreak{}mfw\_, mff\_, mfw\_)))
\end{quote}
also: \texttt{finite-unions-intro} \texttt{finite-unions-unfold} \texttt{finite-unions-algebra} \texttt{finite-unions-empty-in} \texttt{finite-unions-member-in} \texttt{finite-unions-inter-closed} \texttt{finite-unions-inter-member} \texttt{finite-unions-union-closed} ... (10 in all) \\
mentioned by 13 result(s)

\smallskip

\noindent\textbf{\texttt{finprod(m, f, s)}}\quad\textit{functoid}\index{finprod(m, f, s)@\texttt{finprod(m, f, s)}}\\
definition:
\begin{quote}\ttfamily\footnotesize\raggedright\noindent
finsum(\allowbreak{}m, f, s)
\end{quote}

\smallskip

\noindent\textbf{\texttt{finsum(ag, f, s)}}\quad\textit{functoid}\index{finsum(ag, f, s)@\texttt{finsum(ag, f, s)}}\\
reads: the sum of \$1 over the first \$2 indices \\
definition:
\begin{quote}\ttfamily\footnotesize\raggedright\noindent
sum-ag(\allowbreak{}ag, enum-fam(\allowbreak{}ag, f, fin-enum(\allowbreak{}s), card(\allowbreak{}s)), card(\allowbreak{}s))
\end{quote}
characterized by \texttt{finsum-type-ring-additive-ag}  \textit{(definitional)}:
\begin{quote}\ttfamily\footnotesize\raggedright\noindent
forall(\allowbreak{}[r], is-ring(\allowbreak{}r) implies forall(\allowbreak{}[s in set], card(\allowbreak{}s) in nn implies forall(\allowbreak{}[f in fun(\allowbreak{}s, carr(\allowbreak{}r))], finsum(\allowbreak{}ring-additive-ag(\allowbreak{}r), f, s) in carr(\allowbreak{}r))))
\end{quote}
also: \texttt{finsum-all-id-ring-additive-ag} \texttt{rb2-nrm-finsum-le-cna-as-normed-ag} \texttt{rb2-nrm-finsum-le-uba-as-normed-ag} \texttt{rb2-nrm-finsum-le-cnvs-as-normed-ag} \texttt{finsum-act-collect-gen-cna-as-module} \texttt{finsum-act-collect-gen-sip-as-module} \texttt{finsum-act-collect-gen-uba-as-module} \texttt{finsum-act-distrib-gen-cna-as-module} ... (129 in all) \\
mentioned by 690 result(s)

\smallskip

\noindent\textbf{\texttt{finsupp(a, m)}}\quad\textit{functoid}\index{finsupp(a, m)@\texttt{finsupp(a, m)}}\\
reads: the finitely-supported functions \$2 -> \$1 \\
definition:
\begin{quote}\ttfamily\footnotesize\raggedright\noindent
\{f in fun(\allowbreak{}carr(\allowbreak{}m), carr(\allowbreak{}a)): card(\allowbreak{}supp(\allowbreak{}a, m, f)) in nn\}
\end{quote}
characterized by \texttt{finsupp-in-set}  \textit{(proven)}:
\begin{quote}\ttfamily\footnotesize\raggedright\noindent
forall(\allowbreak{}[a\_, m\_], is-ring(\allowbreak{}a\_) implies is-monoid(\allowbreak{}m\_) implies finsupp(\allowbreak{}a\_, m\_) in set)
\end{quote}
also: \texttt{finsupp-unfold} \texttt{finsupp-membership} \texttt{laurent-trunc-finsupp} \texttt{monalg-pullback-finsupp} \\
mentioned by 115 result(s)

\smallskip

\noindent\textbf{\texttt{fker(m, f)}}\quad\textit{functoid}\index{fker(m, f)@\texttt{fker(m, f)}}\\
reads: the kernel of \$2 on \$1 \\
definition:
\begin{quote}\ttfamily\footnotesize\raggedright\noindent
\{fkz\_ in vec(\allowbreak{}m): f(\allowbreak{}fkz\_) = 0\}
\end{quote}
characterized by \texttt{fker-submodule}  \textit{(proven)}:
\begin{quote}\ttfamily\footnotesize\raggedright\noindent
forall(\allowbreak{}[m, f], is-module(\allowbreak{}m) implies scal(\allowbreak{}m) = normed-field-as-commutative-ring(\allowbreak{}rr-normed-field) implies is-linear-functional(\allowbreak{}m, f) implies is-submodule(\allowbreak{}m, fker(\allowbreak{}m, f)))
\end{quote}
also: \texttt{fker-membership} \\
mentioned by 10 result(s)

\smallskip

\noindent\textbf{\texttt{fn-x-section(ys, f, x)}}\quad\textit{functoid}\index{fn-x-section(ys, f, x)@\texttt{fn-x-section(ys, f, x)}}\\
reads: the function \$2 with first argument fixed at \$3, on \$1 \\
definition:
\begin{quote}\ttfamily\footnotesize\raggedright\noindent
vnb-lambda(\allowbreak{}fsy\_, ys, f(\allowbreak{}[x, fsy\_]))
\end{quote}
characterized by \texttt{fn-x-section-apply}  \textit{(proven)}:
\begin{quote}\ttfamily\footnotesize\raggedright\noindent
forall(\allowbreak{}[ys, f, x, z in ys], (\allowbreak{}fn-x-section(\allowbreak{}ys, f, x))(\allowbreak{}z) == f(\allowbreak{}[x, z]))
\end{quote}
also: \texttt{fn-x-section-measurable} \\
mentioned by 17 result(s)

\smallskip

\noindent\textbf{\texttt{fn-y-section(xs, f, y)}}\quad\textit{functoid}\index{fn-y-section(xs, f, y)@\texttt{fn-y-section(xs, f, y)}}\\
reads: the function \$2 with second argument fixed at \$3, on \$1 \\
definition:
\begin{quote}\ttfamily\footnotesize\raggedright\noindent
vnb-lambda(\allowbreak{}fsx\_, xs, f(\allowbreak{}[fsx\_, y]))
\end{quote}
characterized by \texttt{fn-y-section-apply}  \textit{(proven)}:
\begin{quote}\ttfamily\footnotesize\raggedright\noindent
forall(\allowbreak{}[xs, f, y, z in xs], (\allowbreak{}fn-y-section(\allowbreak{}xs, f, y))(\allowbreak{}z) == f(\allowbreak{}[z, y]))
\end{quote}
also: \texttt{fn-y-section-measurable} \\
mentioned by 7 result(s)

\smallskip

\noindent\textbf{\texttt{fnrm(s)}}\quad\textit{accessor}\index{fnrm(s)@\texttt{fnrm(s)}}\\
reads: the norm on the scalar field of \$1 \\
a structure slot \\
characterized by \texttt{cc-normed-field@fnrm}  \textit{(definitional)}:
\begin{quote}\ttfamily\footnotesize\raggedright\noindent
fnrm(\allowbreak{}cc-normed-field) == vnb-lambda(\allowbreak{}x\_, cc, magnitude(\allowbreak{}x\_))
\end{quote}
also: \texttt{rr-normed-field@fnrm} \texttt{pq-fnrm} \texttt{qn-fnrm} \texttt{cc-nf-fnrm} \texttt{rr-nf-fnrm} \texttt{hbs-rr-fnrm} \texttt{sub-nf-fnrm} \texttt{triv-nf-fnrm} ... (15 in all) \\
mentioned by 356 result(s)

\smallskip

\noindent\textbf{\texttt{focus}}\quad\textit{tactic}\index{focus@\texttt{focus}}\\
\texttt{(focus n)} -- Switch the focus to the n-th open goal (1-based, the POSITION in the list (show) prints: the focus first, then the others).  A position shifts as leaves close; the node number in brackets does not -- see focus-id. \\
declared in \texttt{tactics-help.scm}

\smallskip

\noindent\textbf{\texttt{focus-id}}\quad\textit{tactic}\index{focus-id@\texttt{focus-id}}\\
\texttt{(focus-id k)} -- Switch the focus to the ungrounded node whose NODE NUMBER is k, the number (show) and the Focus Workspace print in brackets before each sequent.  Any ungrounded node, not only a leaf: a driver records a leaf by position as (focus n) and an ungrounded node a rule has already fired on as (focus-id k), so a page may name either.  Recorded as a step, undone by backup-one.  The move to make when reading a goal off the panel, e.g. the side leaf of a cut: (cut L) opens  G |- L  (the side leaf) and  G, L |- goal  (the main branch), and a composite such as zero-it leaves the focus on whichever it did not close. \\
declared in \texttt{tactics-help.scm}

\smallskip

\noindent\textbf{\texttt{form-cseq(md, fq)}}\quad\textit{functoid}\index{form-cseq(md, fq)@\texttt{form-cseq(md, fq)}}\\
reads: the Cauchy sequences of the inner product space (\$1, \$2) \\
definition:
\begin{quote}\ttfamily\footnotesize\raggedright\noindent
cseq(\allowbreak{}form-ms(\allowbreak{}md, fq))
\end{quote}
characterized by \texttt{form-cseq-in}  \textit{(proven)}:
\begin{quote}\ttfamily\footnotesize\raggedright\noindent
forall(\allowbreak{}[md, fq, f], is-module(\allowbreak{}md) implies scal(\allowbreak{}md) = normed-field-as-commutative-ring(\allowbreak{}cc-normed-field) implies is-inner-product-form(\allowbreak{}md, fq) implies is-cauchy-seq(\allowbreak{}form-ms(\allowbreak{}md, fq), f) implies f in form-cseq(\allowbreak{}md, fq))
\end{quote}
also: \texttt{form-cseq-fun} \texttt{form-cseq-out} \texttt{form-cseq-unfold} \\
mentioned by 40 result(s)

\smallskip

\noindent\textbf{\texttt{form-ms(md, fq)}}\quad\textit{functoid}\index{form-ms(md, fq)@\texttt{form-ms(md, fq)}}\\
reads: the metric space of the inner product \$2 on \$1 \\
definition:
\begin{quote}\ttfamily\footnotesize\raggedright\noindent
ip-metric-space(\allowbreak{}module-with-form(\allowbreak{}md, fq))
\end{quote}
characterized by \texttt{form-ms-pts}  \textit{(proven)}:
\begin{quote}\ttfamily\footnotesize\raggedright\noindent
forall(\allowbreak{}[md, fq], is-module(\allowbreak{}md) implies scal(\allowbreak{}md) = normed-field-as-commutative-ring(\allowbreak{}cc-normed-field) implies is-inner-product-form(\allowbreak{}md, fq) implies pts(\allowbreak{}form-ms(\allowbreak{}md, fq)) == vec(\allowbreak{}md))
\end{quote}
also: \texttt{form-ms-dist} \texttt{form-ms-is-metric-space} \\
mentioned by 45 result(s)

\smallskip

\noindent\textbf{\texttt{form-null(ipfm\_, ipff\_)}}\quad\textit{functoid}\index{form-null(ipfm_, ipff_)@\texttt{form-null(ipfm\_, ipff\_)}}\\
reads: the null space of the form \$2 \\
definition:
\begin{quote}\ttfamily\footnotesize\raggedright\noindent
\{ipfx\_ in vec(\allowbreak{}ipfm\_): ipff\_(\allowbreak{}ipfx\_, ipfx\_) = 0\}
\end{quote}
characterized by \texttt{form-null-in}  \textit{(proven)}:
\begin{quote}\ttfamily\footnotesize\raggedright\noindent
forall(\allowbreak{}[md, fq, x\_ in vec(\allowbreak{}md)], fq(\allowbreak{}x\_, x\_) = 0 implies x\_ in form-null(\allowbreak{}md, fq))
\end{quote}
also: \texttt{form-null-out} \texttt{form-null-unfold} \texttt{form-null-submodule} \\
mentioned by 38 result(s)

\smallskip

\noindent\textbf{\texttt{form-operator(ubv\_, ubd\_, ubf\_)}}\quad\textit{functoid}\index{form-operator(ubv_, ubd_, ubf_)@\texttt{form-operator(ubv\_, ubd\_, ubf\_)}}\\
reads: the operator of the form \$3 \\
definition:
\begin{quote}\ttfamily\footnotesize\raggedright\noindent
\{ubp\_ in cartesian(\allowbreak{}vec(\allowbreak{}ubv\_), vec(\allowbreak{}ubv\_)): nth(\allowbreak{}1, ubp\_) in ubd\_ and forall(\allowbreak{}[ubx\_ in ubd\_], (\allowbreak{}ip(\allowbreak{}ubv\_))(\allowbreak{}ubx\_, nth(\allowbreak{}2, ubp\_)) = ubf\_(\allowbreak{}ubx\_, nth(\allowbreak{}1, ubp\_)))\}
\end{quote}
characterized by \texttt{form-operator-subset}  \textit{(proven)}:
\begin{quote}\ttfamily\footnotesize\raggedright\noindent
forall(\allowbreak{}[v, dd, bgq\_], is-complex-inner-product-space(\allowbreak{}v) implies form-operator(\allowbreak{}v, dd, bgq\_) subset cartesian(\allowbreak{}vec(\allowbreak{}v), vec(\allowbreak{}v)))
\end{quote}
also: \texttt{form-operator-pair-iff} \texttt{adjoint-compose-form-operator} \texttt{adjoint-compose-form-operator-pair-iff} \texttt{adjoint-compose-form-operator-pointwise} \\
mentioned by 9 result(s)

\smallskip

\noindent\textbf{\texttt{form-plus(ubv\_, ubd\_, ubf\_)}}\quad\textit{functoid}\index{form-plus(ubv_, ubd_, ubf_)@\texttt{form-plus(ubv\_, ubd\_, ubf\_)}}\\
reads: the form \$3 plus the inner product \\
definition:
\begin{quote}\ttfamily\footnotesize\raggedright\noindent
vnb-lambda(\allowbreak{}[ubx\_, uby\_], cartesian(\allowbreak{}ubd\_, ubd\_), ubf\_(\allowbreak{}ubx\_, uby\_) + (\allowbreak{}ip(\allowbreak{}ubv\_))(\allowbreak{}ubx\_, uby\_))
\end{quote}
characterized by \texttt{form-plus-inner-product-form}  \textit{(proven)}:
\begin{quote}\ttfamily\footnotesize\raggedright\noindent
forall(\allowbreak{}[v, dd, q], is-complex-inner-product-space(\allowbreak{}v) implies is-submodule(\allowbreak{}v, dd) implies is-positive-form(\allowbreak{}v, dd, q) implies is-inner-product-form(\allowbreak{}complex-inner-product-space-as-module(\allowbreak{}sub-cips(\allowbreak{}v, dd)), form-plus(\allowbreak{}v, dd, q)))
\end{quote}
also: \texttt{form-plus-space-as-module-with-form} \\
mentioned by 3 result(s)

\smallskip

\noindent\textbf{\texttt{form-plus-space(ubv\_, ubd\_, ubf\_)}}\quad\textit{functoid}\index{form-plus-space(ubv_, ubd_, ubf_)@\texttt{form-plus-space(ubv\_, ubd\_, ubf\_)}}\\
reads: the space \$2 with the inner product of \$3 plus that of \$1 \\
definition:
\begin{quote}\ttfamily\footnotesize\raggedright\noindent
[normed-field-as-commutative-ring(\allowbreak{}cc-normed-field), ubd\_, vnb-lambda(\allowbreak{}[ubx\_, uby\_], cartesian(\allowbreak{}ubd\_, ubd\_), (\allowbreak{}vadd(\allowbreak{}ubv\_))(\allowbreak{}ubx\_, uby\_)), vzero(\allowbreak{}ubv\_), vnb-lambda(\allowbreak{}ubx\_, ubd\_, (\allowbreak{}vneg(\allowbreak{}ubv\_))(\allowbreak{}ubx\_)), vnb-lambda(\allowbreak{}[ubk\_, ubx\_], cartesian(\allowbreak{}cc, ubd\_), (\allowbreak{}act(\allowbreak{}ubv\_))(\allowbreak{}ubk\_, ubx\_)), form-plus(\allowbreak{}ubv\_, ubd\_, ubf\_)]
\end{quote}
characterized by \texttt{form-plus-space-is-cips}  \textit{(proven)}:
\begin{quote}\ttfamily\footnotesize\raggedright\noindent
forall(\allowbreak{}[v, dd, q], is-complex-inner-product-space(\allowbreak{}v) implies is-submodule(\allowbreak{}v, dd) implies is-positive-form(\allowbreak{}v, dd, q) implies is-complex-inner-product-space(\allowbreak{}form-plus-space(\allowbreak{}v, dd, q)))
\end{quote}
also: \texttt{form-plus-space-as-module-with-form} \\
mentioned by 3 result(s)

\smallskip

\noindent\textbf{\texttt{form-qip(ipfm\_, ipff\_)}}\quad\textit{functoid}\index{form-qip(ipfm_, ipff_)@\texttt{form-qip(ipfm\_, ipff\_)}}\\
reads: the form induced by \$2 on the quotient by its null space \\
definition:
\begin{quote}\ttfamily\footnotesize\raggedright\noindent
descend2(\allowbreak{}qmod-setoid(\allowbreak{}ipfm\_, form-null(\allowbreak{}ipfm\_, ipff\_)), ipff\_)
\end{quote}
characterized by \texttt{form-qip-type}  \textit{(proven)}:
\begin{quote}\ttfamily\footnotesize\raggedright\noindent
forall(\allowbreak{}[md, fq], is-module(\allowbreak{}md) implies scal(\allowbreak{}md) = normed-field-as-commutative-ring(\allowbreak{}cc-normed-field) implies is-psd-form(\allowbreak{}md, fq) implies form-qip(\allowbreak{}md, fq) in fun(\allowbreak{}cartesian(\allowbreak{}quotient(\allowbreak{}qmod-setoid(\allowbreak{}md, form-null(\allowbreak{}md, fq))), quotient(\allowbreak{}qmod-setoid(\allowbreak{}md, form-null(\allowbreak{}md, fq)))), cc))
\end{quote}
also: \texttt{form-qip-computes} \\
mentioned by 10 result(s)

\smallskip

\noindent\textbf{\texttt{form-quotient(ipfm\_, ipff\_)}}\quad\textit{functoid}\index{form-quotient(ipfm_, ipff_)@\texttt{form-quotient(ipfm\_, ipff\_)}}\\
reads: the quotient of \$1 by the null space of \$2, with the induced inner product \\
definition:
\begin{quote}\ttfamily\footnotesize\raggedright\noindent
module-with-form(\allowbreak{}quotient-module(\allowbreak{}ipfm\_, form-null(\allowbreak{}ipfm\_, ipff\_)), form-qip(\allowbreak{}ipfm\_, ipff\_))
\end{quote}
characterized by \texttt{form-quotient-ip}  \textit{(proven)}:
\begin{quote}\ttfamily\footnotesize\raggedright\noindent
forall(\allowbreak{}[md, fq], is-module(\allowbreak{}md) implies scal(\allowbreak{}md) = normed-field-as-commutative-ring(\allowbreak{}cc-normed-field) implies is-psd-form(\allowbreak{}md, fq) implies ip(\allowbreak{}form-quotient(\allowbreak{}md, fq)) == form-qip(\allowbreak{}md, fq))
\end{quote}
also: \texttt{form-quotient-act} \texttt{form-quotient-vec} \texttt{form-quotient-vadd} \texttt{form-quotient-vneg} \texttt{form-quotient-is-cips} \\
mentioned by 11 result(s)

\smallskip

\noindent\textbf{\texttt{fp(qsp\_)}}\quad\textit{functoid}\index{fp(qsp_)@\texttt{fp(qsp\_)}}\\
reads: the field F\_\$1 = Z / \$1 Z \\
definition:
\begin{quote}\ttfamily\footnotesize\raggedright\noindent
ringoid-quotient-ring(\allowbreak{}fp-ringoid(\allowbreak{}qsp\_))
\end{quote}
characterized by \texttt{fp-is-field}  \textit{(proven)}:
\begin{quote}\ttfamily\footnotesize\raggedright\noindent
forall(\allowbreak{}[p], zz-prime(\allowbreak{}p) implies is-field-ring(\allowbreak{}fp(\allowbreak{}p)))
\end{quote}
also: \texttt{nv-residue-fp-iso} \\
mentioned by 4 result(s)

\smallskip

\noindent\textbf{\texttt{fp-ringoid(qsp\_)}}\quad\textit{functoid}\index{fp-ringoid(qsp_)@\texttt{fp-ringoid(qsp\_)}}\\
reads: the integers with the ideal \$1 Z \\
definition:
\begin{quote}\ttfamily\footnotesize\raggedright\noindent
[zz, vnb-lambda(\allowbreak{}[x\_, y\_], cartesian(\allowbreak{}zz, zz), x\_ + y\_), vnb-lambda(\allowbreak{}[x\_, y\_], cartesian(\allowbreak{}zz, zz), x\_ * y\_), vnb-lambda(\allowbreak{}x\_, zz, -x\_), 0, 1, \{fpn\_ in zz: zz-divides(\allowbreak{}qsp\_, fpn\_)\}]
\end{quote}
mentioned by 17 result(s)

\smallskip

\noindent\textbf{\texttt{fr-ball(m, fam, n, eps)}}\quad\textit{functoid}\index{fr-ball(m, fam, n, eps)@\texttt{fr-ball(m, fam, n, eps)}}\\
definition:
\begin{quote}\ttfamily\footnotesize\raggedright\noindent
\{x in vec(\allowbreak{}m): forall(\allowbreak{}[k], k in nn and k <= n implies (\allowbreak{}fam(\allowbreak{}k))(\allowbreak{}x) < eps)\}
\end{quote}
characterized by \texttt{fom-fr-ball-unfold}  \textit{(proven)}:
\begin{quote}\ttfamily\footnotesize\raggedright\noindent
forall(\allowbreak{}[m, fam, n, eps], fr-ball(\allowbreak{}m, fam, n, eps) == \{x in vec(\allowbreak{}m): forall(\allowbreak{}[k], k in nn and k <= n implies (\allowbreak{}fam(\allowbreak{}k))(\allowbreak{}x) < eps)\})
\end{quote}
mentioned by 4 result(s)

\smallskip

\noindent\textbf{\texttt{fr-c-space(m, fam)}}\quad\textit{functoid}\index{fr-c-space(m, fam)@\texttt{fr-c-space(m, fam)}}\\
reads: the pseudometric family of the seminorms \$2 on \$1 \\
definition:
\begin{quote}\ttfamily\footnotesize\raggedright\noindent
[vec(\allowbreak{}m), vnb-lambda(\allowbreak{}frk\_, nn, vnb-lambda(\allowbreak{}[fru\_, frv\_], cartesian(\allowbreak{}vec(\allowbreak{}m), vec(\allowbreak{}m)), (\allowbreak{}fam(\allowbreak{}frk\_))(\allowbreak{}(\allowbreak{}vadd(\allowbreak{}m))(\allowbreak{}fru\_, (\allowbreak{}vneg(\allowbreak{}m))(\allowbreak{}frv\_)))))]
\end{quote}
mentioned by 7 result(s)

\smallskip

\noindent\textbf{\texttt{fr-cauchy(m, fam, seq)}}\quad\textit{predicate}\index{fr-cauchy(m, fam, seq)@\texttt{fr-cauchy(m, fam, seq)}}\\
reads: \$3 is Cauchy in (\$1, \$2) \\
defined by \texttt{fr-cauchy} \\
characterized by \texttt{fr-cauchy}  \textit{(definitional)}:
\begin{quote}\ttfamily\footnotesize\raggedright\noindent
forall(\allowbreak{}[m, fam, seq], fr-cauchy(\allowbreak{}m, fam, seq) iff forall(\allowbreak{}[k in nn, eps], pos-rr(\allowbreak{}eps) implies forsome(\allowbreak{}[cap in nn], forall(\allowbreak{}[i], i in nn and cap <= i implies forall(\allowbreak{}[j], j in nn and cap <= j implies (\allowbreak{}fam(\allowbreak{}k))(\allowbreak{}(\allowbreak{}vadd(\allowbreak{}m))(\allowbreak{}seq(\allowbreak{}i), (\allowbreak{}vneg(\allowbreak{}m))(\allowbreak{}seq(\allowbreak{}j)))) < eps)))))
\end{quote}
mentioned by 10 result(s)

\smallskip

\noindent\textbf{\texttt{fr-conv(m, fam, seq, lim)}}\quad\textit{predicate}\index{fr-conv(m, fam, seq, lim)@\texttt{fr-conv(m, fam, seq, lim)}}\\
reads: \$3 converges to \$4 in (\$1, \$2) \\
defined by \texttt{fr-conv} \\
characterized by \texttt{fr-conv}  \textit{(definitional)}:
\begin{quote}\ttfamily\footnotesize\raggedright\noindent
forall(\allowbreak{}[m, fam, seq, lim], fr-conv(\allowbreak{}m, fam, seq, lim) iff forall(\allowbreak{}[k in nn, eps], pos-rr(\allowbreak{}eps) implies forsome(\allowbreak{}[cap in nn], forall(\allowbreak{}[i], i in nn and cap <= i implies (\allowbreak{}fam(\allowbreak{}k))(\allowbreak{}(\allowbreak{}vadd(\allowbreak{}m))(\allowbreak{}seq(\allowbreak{}i), (\allowbreak{}vneg(\allowbreak{}m))(\allowbreak{}lim))) < eps))))
\end{quote}
also: \texttt{cc-fr-conv-transfer} \texttt{cc-fr-conv-iff-unif-on-compacta} \\
mentioned by 21 result(s)

\smallskip

\noindent\textbf{\texttt{fr-metric(m, fam)}}\quad\textit{functoid}\index{fr-metric(m, fam)@\texttt{fr-metric(m, fam)}}\\
reads: the metric of the seminorms \$2 on \$1 \\
definition:
\begin{quote}\ttfamily\footnotesize\raggedright\noindent
dist(\allowbreak{}fr-metric-space(\allowbreak{}m, fam))
\end{quote}
mentioned by 2 result(s)

\smallskip

\noindent\textbf{\texttt{fr-metric-space(m, fam)}}\quad\textit{functoid}\index{fr-metric-space(m, fam)@\texttt{fr-metric-space(m, fam)}}\\
reads: the metric space of the seminorms \$2 on \$1 \\
definition:
\begin{quote}\ttfamily\footnotesize\raggedright\noindent
c-metric(\allowbreak{}fr-c-space(\allowbreak{}m, fam))
\end{quote}
mentioned by 26 result(s)

\smallskip

\noindent\textbf{\texttt{fun}}\quad\textit{operator}\index{fun@\texttt{fun}}\\
a term-forming head \\
characterized by \texttt{is-degree-function}  \textit{(definitional)}:
\begin{quote}\ttfamily\footnotesize\raggedright\noindent
forall(\allowbreak{}[gfs\_, gfn\_], is-degree-function(\allowbreak{}gfs\_, gfn\_) iff is-integral-domain(\allowbreak{}gfs\_) and gfn\_ in fun(\allowbreak{}carr(\allowbreak{}gfs\_), nn) and forall(\allowbreak{}[gfa\_ in carr(\allowbreak{}gfs\_), gfb\_ in carr(\allowbreak{}gfs\_)], not(\allowbreak{}gfa\_ = zero(\allowbreak{}gfs\_)) implies not(\allowbreak{}gfb\_ = zero(\allowbreak{}gfs\_)) implies gfn\_(\allowbreak{}gfa\_) <= gfn\_(\allowbreak{}(\allowbreak{}mul(\allowbreak{}gfs\_))(\allowbreak{}gfa\_, gfb\_))) and forall(\allowbreak{}[gfa\_ in carr(\allowbreak{}gfs\_), gfb\_ in carr(\allowbreak{}gfs\_)], not(\allowbreak{}gfa\_ = zero(\allowbreak{}gfs\_)) implies forsome(\allowbreak{}[gfq\_ in carr(\allowbreak{}gfs\_)], gfb\_ = (\allowbreak{}mul(\allowbreak{}gfs\_))(\allowbreak{}gfa\_, gfq\_) or gfn\_(\allowbreak{}(\allowbreak{}add(\allowbreak{}gfs\_))(\allowbreak{}gfb\_, (\allowbreak{}neg(\allowbreak{}gfs\_))(\allowbreak{}(\allowbreak{}mul(\allowbreak{}gfs\_))(\allowbreak{}gfa\_, gfq\_)))) < gfn\_(\allowbreak{}gfa\_))))
\end{quote}
also: \texttt{is-convex-functional} \texttt{is-linear-functional} \texttt{is-positive-functional} \texttt{is-cc-linear-functional} \texttt{is-linear-functional-on} \texttt{sip-act-fun-csip-as-sip} \texttt{sip-rip-fun-csip-as-sip} \texttt{is-bounded-l1-functional} ... (470 in all) \\
mentioned by 4544 result(s)

\smallskip

\noindent\textbf{\texttt{fun-def(mcfh\_, mcfn\_, mcfb\_)}}\quad\textit{predicate}\index{fun-def(mcfh_, mcfn_, mcfb_)@\texttt{fun-def(mcfh\_, mcfn\_, mcfb\_)}}\\
reads: the functoid code \$1 has arity \$2 and body \$3 \\
defined by \texttt{fun-def} \\
characterized by \texttt{fun-def}  \textit{(definitional)}:
\begin{quote}\ttfamily\footnotesize\raggedright\noindent
forall(\allowbreak{}[mcfh\_, mcfn\_, mcfb\_], fun-def(\allowbreak{}mcfh\_, mcfn\_, mcfb\_) iff mcfh\_ = 32 and mcfn\_ = 2 and mcfb\_ = enc-sep(\allowbreak{}ebvar(\allowbreak{}0), enc-in(\allowbreak{}ebvar(\allowbreak{}0), ebvar(\allowbreak{}2))))
\end{quote}
also: \texttt{fun-def-heads} \\
mentioned by 15 result(s)

\smallskip

\noindent\textbf{\texttt{functional-le(g4a\_, g4g\_, g4f\_)}}\quad\textit{predicate}\index{functional-le(g4a_, g4g_, g4f_)@\texttt{functional-le(g4a\_, g4g\_, g4f\_)}}\\
reads: \$2 <= \$3 as functionals on \$1 \\
defined by \texttt{functional-le} \\
characterized by \texttt{functional-le}  \textit{(definitional)}:
\begin{quote}\ttfamily\footnotesize\raggedright\noindent
forall(\allowbreak{}[g4a\_, g4g\_, g4f\_], functional-le(\allowbreak{}g4a\_, g4g\_, g4f\_) iff forall(\allowbreak{}[g4x\_ in vec(\allowbreak{}g4a\_)], g4g\_(\allowbreak{}(\allowbreak{}vmul(\allowbreak{}g4a\_))(\allowbreak{}(\allowbreak{}vstar(\allowbreak{}g4a\_))(\allowbreak{}g4x\_), g4x\_)) <= g4f\_(\allowbreak{}(\allowbreak{}vmul(\allowbreak{}g4a\_))(\allowbreak{}(\allowbreak{}vstar(\allowbreak{}g4a\_))(\allowbreak{}g4x\_), g4x\_))))
\end{quote}
mentioned by 5 result(s)

\smallskip

\noindent\textbf{\texttt{fv-stage}}\quad\textit{functoid}\index{fv-stage@\texttt{fv-stage}}\\
reads: the free-name map on the stage \$1 \\
defined by \texttt{fv-stage-zero} \texttt{fv-stage-succ} \\
characterized by \texttt{fv-stage-succ}  \textit{(definitional)}:
\begin{quote}\ttfamily\footnotesize\raggedright\noindent
forall(\allowbreak{}[mxn\_ in nn], fv-stage(\allowbreak{}succ(\allowbreak{}mxn\_)) == vnb-lambda(\allowbreak{}mxe\_, expr-stage(\allowbreak{}succ(\allowbreak{}mxn\_)), if(\allowbreak{}mxe\_ in expr-stage(\allowbreak{}mxn\_), (\allowbreak{}fv-stage(\allowbreak{}mxn\_))(\allowbreak{}mxe\_), if(\allowbreak{}nth(\allowbreak{}1, mxe\_) = 3, big-union(\allowbreak{}mxw\_, make-set(\allowbreak{}nth(\allowbreak{}3, mxe\_)), (\allowbreak{}fv-stage(\allowbreak{}mxn\_))(\allowbreak{}mxw\_)), union(\allowbreak{}big-union(\allowbreak{}mxw\_, make-set(\allowbreak{}nth(\allowbreak{}4, mxe\_)), (\allowbreak{}fv-stage(\allowbreak{}mxn\_))(\allowbreak{}mxw\_)), (\allowbreak{}fv-stage(\allowbreak{}mxn\_))(\allowbreak{}nth(\allowbreak{}5, mxe\_)))))))
\end{quote}
also: \texttt{fv-stage-zero} \texttt{fv-stage-coh} \texttt{fv-stage-step} \\
mentioned by 10 result(s)

\smallskip


\section*{G}
\addcontentsline{toc}{section}{G}

\noindent\textbf{\texttt{gamma-ext(gez\_)}}\quad\textit{functoid}\index{gamma-ext(gez_)@\texttt{gamma-ext(gez\_)}}\\
reads: the analytic continuation of the Gamma function at \$1 \\
definition:
\begin{quote}\ttfamily\footnotesize\raggedright\noindent
gamma-function(\allowbreak{}gez\_ + choice(\allowbreak{}\{gek\_ in nn: 0 < real-part(\allowbreak{}gez\_) + gek\_\})) * recip(\allowbreak{}poch(\allowbreak{}gez\_, choice(\allowbreak{}\{gek\_ in nn: 0 < real-part(\allowbreak{}gez\_) + gek\_\})))
\end{quote}
characterized by \texttt{gamma-ext-value}  \textit{(proven)}:
\begin{quote}\ttfamily\footnotesize\raggedright\noindent
forall(\allowbreak{}[gxn\_ in nn, gxz\_ in \{gdz\_ in cc: forall(\allowbreak{}[gdn\_ in nn], not(\allowbreak{}gdz\_ + gdn\_ = 0))\}], 0 < real-part(\allowbreak{}gxz\_) + gxn\_ implies gamma-ext(\allowbreak{}gxz\_) = gamma-function(\allowbreak{}gxz\_ + gxn\_) * recip(\allowbreak{}poch(\allowbreak{}gxz\_, gxn\_)))
\end{quote}
also: \texttt{gamma-ext-extends} \texttt{gamma-ext-recursion} \texttt{gamma-ext-holomorphic} \\
mentioned by 8 result(s)

\smallskip

\noindent\textbf{\texttt{gamma-function(gmz\_)}}\quad\textit{functoid}\index{gamma-function(gmz_)@\texttt{gamma-function(gmz\_)}}\\
reads: Gamma(\$1) \\
definition:
\begin{quote}\ttfamily\footnotesize\raggedright\noindent
improper-int(\allowbreak{}gamma-integrand(\allowbreak{}gmz\_))
\end{quote}
mentioned by 20 result(s)

\smallskip

\noindent\textbf{\texttt{gamma-integrand(gmz\_)}}\quad\textit{functoid}\index{gamma-integrand(gmz_)@\texttt{gamma-integrand(gmz\_)}}\\
reads: the Gamma integrand t |-> t\textasciicircum{}(\$1 - 1) e\textasciicircum{}(-t) \\
definition:
\begin{quote}\ttfamily\footnotesize\raggedright\noindent
vnb-lambda(\allowbreak{}gmw\_, \{gmu\_ in cc: 0 < real-part(\allowbreak{}gmu\_)\}, cc-exp(\allowbreak{}(\allowbreak{}gmz\_ - 1) * cc-log(\allowbreak{}gmw\_)) * cc-exp(\allowbreak{}-gmw\_))
\end{quote}
characterized by \texttt{gamma-integrand-locint}  \textit{(proven)}:
\begin{quote}\ttfamily\footnotesize\raggedright\noindent
forall(\allowbreak{}[z in cc, ipr\_ in rr, ipq\_ in rr], 0 < ipr\_ implies ipr\_ < ipq\_ implies cc-int(\allowbreak{}gamma-integrand(\allowbreak{}z), ipr\_, ipq\_) in cc)
\end{quote}
also: \texttt{gamma-integrand-expansion} \texttt{gamma-integrand-magnitude} \texttt{gamma-integrand-quadratic} \texttt{gamma-integrand-bound-large} \texttt{gamma-integrand-bound-small} \texttt{gamma-integrand-holomorphic} \texttt{gamma-integrand-bound-interval} \\
mentioned by 21 result(s)

\smallskip

\noindent\textbf{\texttt{gamma-product}}\quad\textit{constant}\index{gamma-product@\texttt{gamma-product}}\\
reads: the Gamma function of Definition 6.27 (the reciprocal of the canonical product) \\
defined by \texttt{gamma-product-def} \\
mentioned by 1 result(s)

\smallskip

\noindent\textbf{\texttt{gauge(s)}}\quad\textit{functoid}\index{gauge(s)@\texttt{gauge(s)}}\\
definition:
\begin{quote}\ttfamily\footnotesize\raggedright\noindent
choice(\allowbreak{}euclidean-gauges(\allowbreak{}s))
\end{quote}
characterized by \texttt{gauge-is-degree}  \textit{(proven)}:
\begin{quote}\ttfamily\footnotesize\raggedright\noindent
forall(\allowbreak{}[s], is-euclidean-ring(\allowbreak{}s) implies gauge(\allowbreak{}s) in fun(\allowbreak{}carr(\allowbreak{}s), nn) and has-div-remainder(\allowbreak{}s, gauge(\allowbreak{}s)))
\end{quote}
mentioned by 19 result(s)

\smallskip

\noindent\textbf{\texttt{gdist(s)}}\quad\textit{accessor}\index{gdist(s)@\texttt{gdist(s)}}\\
reads: the distance of \$1 \\
a structure slot \\
characterized by \texttt{opp-mg-gdist}  \textit{(proven)}:
\begin{quote}\ttfamily\footnotesize\raggedright\noindent
forall(\allowbreak{}[g], is-metric-group(\allowbreak{}g) implies gdist(\allowbreak{}opp-metric-group(\allowbreak{}g)) == gdist(\allowbreak{}g))
\end{quote}
also: \texttt{make-metric-group-gdist} \\
mentioned by 10 result(s)

\smallskip

\noindent\textbf{\texttt{gelfand-dual-map(baa\_, bab\_, baf\_)}}\quad\textit{functoid}\index{gelfand-dual-map(baa_, bab_, baf_)@\texttt{gelfand-dual-map(baa\_, bab\_, baf\_)}}\\
reads: the dual map of \$3 on the Gelfand spaces \\
definition:
\begin{quote}\ttfamily\footnotesize\raggedright\noindent
vnb-lambda(\allowbreak{}bgw\_, characters(\allowbreak{}bab\_), vnb-lambda(\allowbreak{}bgx\_, vec(\allowbreak{}baa\_), bgw\_(\allowbreak{}baf\_(\allowbreak{}bgx\_))))
\end{quote}

\smallskip

\noindent\textbf{\texttt{gelfand-space(baa\_)}}\quad\textit{functoid}\index{gelfand-space(baa_)@\texttt{gelfand-space(baa\_)}}\\
reads: the Gelfand space of \$1 \\
definition:
\begin{quote}\ttfamily\footnotesize\raggedright\noindent
subspace-top(\allowbreak{}tproduct-top(\allowbreak{}vec(\allowbreak{}baa\_), vnb-lambda(\allowbreak{}bgx\_, vec(\allowbreak{}baa\_), metric-top(\allowbreak{}cc-ms))), characters(\allowbreak{}baa\_))
\end{quote}
characterized by \texttt{gelfand-space-pts}  \textit{(proven)}:
\begin{quote}\ttfamily\footnotesize\raggedright\noindent
forall(\allowbreak{}[a], is-unital-banach-algebra(\allowbreak{}a) implies pts(\allowbreak{}gelfand-space(\allowbreak{}a)) == characters(\allowbreak{}a))
\end{quote}
also: \texttt{gelfand-space-compact} \texttt{b2g-gelfand-space-view} \\
mentioned by 12 result(s)

\smallskip

\noindent\textbf{\texttt{gelfand-transform(baa\_)}}\quad\textit{functoid}\index{gelfand-transform(baa_)@\texttt{gelfand-transform(baa\_)}}\\
reads: the Gelfand transform of \$1 \\
definition:
\begin{quote}\ttfamily\footnotesize\raggedright\noindent
vnb-lambda(\allowbreak{}bgx\_, vec(\allowbreak{}baa\_), vnb-lambda(\allowbreak{}bgc\_, characters(\allowbreak{}baa\_), bgc\_(\allowbreak{}bgx\_)))
\end{quote}
characterized by \texttt{gelfand-transform-act}  \textit{(proven)}:
\begin{quote}\ttfamily\footnotesize\raggedright\noindent
forall(\allowbreak{}[a], is-unital-banach-algebra(\allowbreak{}a) implies forall(\allowbreak{}[z in cc, x in vec(\allowbreak{}a), c in characters(\allowbreak{}a)], (\allowbreak{}(\allowbreak{}gelfand-transform(\allowbreak{}a))(\allowbreak{}(\allowbreak{}act(\allowbreak{}a))(\allowbreak{}z, x)))(\allowbreak{}c) = z * (\allowbreak{}(\allowbreak{}gelfand-transform(\allowbreak{}a))(\allowbreak{}x))(\allowbreak{}c)))
\end{quote}
also: \texttt{gelfand-transform-add} \texttt{gelfand-transform-mul} \texttt{gelfand-transform-one} \texttt{gelfand-transform-value} \texttt{gelfand-transform-in-fun} \texttt{gelfand-transform-bounded} \texttt{b2g-gelfand-transform-view} \texttt{gelfand-transform-in-ccont} ... (9 in all) \\
mentioned by 23 result(s)

\smallskip

\noindent\textbf{\texttt{generated-subgroup(gcg\_, gcx\_)}}\quad\textit{functoid}\index{generated-subgroup(gcg_, gcx_)@\texttt{generated-subgroup(gcg\_, gcx\_)}}\\
reads: the subgroup of \$1 generated by \$2 \\
definition:
\begin{quote}\ttfamily\footnotesize\raggedright\noindent
\{gcy\_ in carr(\allowbreak{}gcg\_): forall(\allowbreak{}[gcs\_ in power(\allowbreak{}carr(\allowbreak{}gcg\_))], is-subgroup(\allowbreak{}gcg\_, gcs\_) implies gcx\_ in gcs\_ implies gcy\_ in gcs\_)\}
\end{quote}
characterized by \texttt{generated-subgroup-unfold}  \textit{(proven)}:
\begin{quote}\ttfamily\footnotesize\raggedright\noindent
forall(\allowbreak{}[g, x], generated-subgroup(\allowbreak{}g, x) == \{gcy\_ in carr(\allowbreak{}g): forall(\allowbreak{}[gcs\_ in power(\allowbreak{}carr(\allowbreak{}g))], is-subgroup(\allowbreak{}g, gcs\_) implies x in gcs\_ implies gcy\_ in gcs\_)\})
\end{quote}
also: \texttt{generated-subgroup-is-subgroup} \\
mentioned by 4 result(s)

\smallskip

\noindent\textbf{\texttt{generates(md, n, u)}}\quad\textit{predicate}\index{generates(md, n, u)@\texttt{generates(md, n, u)}}\\
reads: the \$2 vectors \$3 generate \$1 \\
defined by \texttt{generates} \\
characterized by \texttt{generates}  \textit{(definitional)}:
\begin{quote}\ttfamily\footnotesize\raggedright\noindent
forall(\allowbreak{}[md, n, u], generates(\allowbreak{}md, n, u) iff forall(\allowbreak{}[x\_ in vec(\allowbreak{}md)], forsome(\allowbreak{}[c\_ in mat(\allowbreak{}1, n, carr(\allowbreak{}scal(\allowbreak{}md)))], x\_ = lincomb(\allowbreak{}md, n, c\_, u))))
\end{quote}
also: \texttt{generates-transport-cna-as-module} \texttt{generates-transport-sip-as-module} \texttt{generates-transport-uba-as-module} \texttt{generates-transport-cnvs-as-module} \texttt{generates-coeff-matrix-cna-as-module} \texttt{generates-coeff-matrix-sip-as-module} \texttt{generates-coeff-matrix-uba-as-module} \texttt{generates-coeff-matrix-cnvs-as-module} ... (22 in all) \\
mentioned by 315 result(s)

\smallskip

\noindent\textbf{\texttt{gns-form(g4ga\_, g4gf\_)}}\quad\textit{functoid}\index{gns-form(g4ga_, g4gf_)@\texttt{gns-form(g4ga\_, g4gf\_)}}\\
reads: the form (x, y) |-> \$2(y* x) on \$1 \\
definition:
\begin{quote}\ttfamily\footnotesize\raggedright\noindent
vnb-lambda(\allowbreak{}[g4gx\_, g4gy\_], cartesian(\allowbreak{}vec(\allowbreak{}g4ga\_), vec(\allowbreak{}g4ga\_)), g4gf\_(\allowbreak{}(\allowbreak{}vmul(\allowbreak{}g4ga\_))(\allowbreak{}(\allowbreak{}vstar(\allowbreak{}g4ga\_))(\allowbreak{}g4gy\_), g4gx\_)))
\end{quote}
characterized by \texttt{gns-form-psd}  \textit{(proven)}:
\begin{quote}\ttfamily\footnotesize\raggedright\noindent
forall(\allowbreak{}[a], is-unital-involutive-banach-algebra(\allowbreak{}a) implies forall(\allowbreak{}[f], is-positive-lf(\allowbreak{}a, f) implies is-psd-form(\allowbreak{}gns-module(\allowbreak{}a), gns-form(\allowbreak{}a, f))))
\end{quote}
also: \texttt{gns-form-value} \\
mentioned by 9 result(s)

\smallskip

\noindent\textbf{\texttt{gns-module(g4ga\_)}}\quad\textit{functoid}\index{gns-module(g4ga_)@\texttt{gns-module(g4ga\_)}}\\
reads: \$1 as a complex module \\
definition:
\begin{quote}\ttfamily\footnotesize\raggedright\noindent
uba-as-module(\allowbreak{}uiba-as-uba(\allowbreak{}g4ga\_))
\end{quote}
characterized by \texttt{gns-module-is-module}  \textit{(proven)}:
\begin{quote}\ttfamily\footnotesize\raggedright\noindent
forall(\allowbreak{}[a], is-unital-involutive-banach-algebra(\allowbreak{}a) implies is-module(\allowbreak{}gns-module(\allowbreak{}a)))
\end{quote}
mentioned by 20 result(s)

\smallskip

\noindent\textbf{\texttt{gns-op(g4ga\_, g4gf\_, g4gx\_)}}\quad\textit{functoid}\index{gns-op(g4ga_, g4gf_, g4gx_)@\texttt{gns-op(g4ga\_, g4gf\_, g4gx\_)}}\\
reads: the GNS operator of \$3 \\
definition:
\begin{quote}\ttfamily\footnotesize\raggedright\noindent
descend(\allowbreak{}gns-setoid(\allowbreak{}g4ga\_, g4gf\_), vnb-lambda(\allowbreak{}g4gb\_, vec(\allowbreak{}g4ga\_), class(\allowbreak{}gns-setoid(\allowbreak{}g4ga\_, g4gf\_), (\allowbreak{}vmul(\allowbreak{}g4ga\_))(\allowbreak{}g4gx\_, g4gb\_))))
\end{quote}
characterized by \texttt{gns-op-act}  \textit{(proven)}:
\begin{quote}\ttfamily\footnotesize\raggedright\noindent
forall(\allowbreak{}[a], is-unital-involutive-banach-algebra(\allowbreak{}a) implies forall(\allowbreak{}[f], is-positive-lf(\allowbreak{}a, f) implies forall(\allowbreak{}[x in vec(\allowbreak{}a), c in cc, u in vec(\allowbreak{}gns-pre(\allowbreak{}a, f))], (\allowbreak{}gns-op(\allowbreak{}a, f, x))(\allowbreak{}(\allowbreak{}act(\allowbreak{}gns-pre(\allowbreak{}a, f)))(\allowbreak{}c, u)) = (\allowbreak{}act(\allowbreak{}gns-pre(\allowbreak{}a, f)))(\allowbreak{}c, (\allowbreak{}gns-op(\allowbreak{}a, f, x))(\allowbreak{}u)))))
\end{quote}
also: \texttt{gns-op-mul} \texttt{gns-op-sum} \texttt{gns-op-type} \texttt{gns-op-unit} \texttt{gns-op-vadd} \texttt{gns-op-scale} \texttt{gns-op-adjoint} \texttt{gns-op-bounded} ... (9 in all) \\
mentioned by 23 result(s)

\smallskip

\noindent\textbf{\texttt{gns-pre(g4ga\_, g4gf\_)}}\quad\textit{functoid}\index{gns-pre(g4ga_, g4gf_)@\texttt{gns-pre(g4ga\_, g4gf\_)}}\\
reads: the GNS pre-Hilbert space of \$2 \\
definition:
\begin{quote}\ttfamily\footnotesize\raggedright\noindent
form-quotient(\allowbreak{}gns-module(\allowbreak{}g4ga\_), gns-form(\allowbreak{}g4ga\_, g4gf\_))
\end{quote}
characterized by \texttt{gns-pre-ip}  \textit{(proven)}:
\begin{quote}\ttfamily\footnotesize\raggedright\noindent
forall(\allowbreak{}[a], is-unital-involutive-banach-algebra(\allowbreak{}a) implies forall(\allowbreak{}[f], is-positive-lf(\allowbreak{}a, f) implies ip(\allowbreak{}gns-pre(\allowbreak{}a, f)) == form-qip(\allowbreak{}gns-module(\allowbreak{}a), gns-form(\allowbreak{}a, f))))
\end{quote}
also: \texttt{gns-pre-act} \texttt{gns-pre-vec} \texttt{gns-pre-vadd} \texttt{gns-pre-is-cips} \texttt{gns-pre-act-computes} \texttt{gns-pre-vadd-computes} \\
mentioned by 37 result(s)

\smallskip

\noindent\textbf{\texttt{gns-setoid(g4ga\_, g4gf\_)}}\quad\textit{functoid}\index{gns-setoid(g4ga_, g4gf_)@\texttt{gns-setoid(g4ga\_, g4gf\_)}}\\
reads: congruence modulo the null space of \$2 \\
definition:
\begin{quote}\ttfamily\footnotesize\raggedright\noindent
qmod-setoid(\allowbreak{}gns-module(\allowbreak{}g4ga\_), form-null(\allowbreak{}gns-module(\allowbreak{}g4ga\_), gns-form(\allowbreak{}g4ga\_, g4gf\_)))
\end{quote}
characterized by \texttt{gns-setoid-pts}  \textit{(proven)}:
\begin{quote}\ttfamily\footnotesize\raggedright\noindent
forall(\allowbreak{}[a], is-unital-involutive-banach-algebra(\allowbreak{}a) implies forall(\allowbreak{}[f], is-positive-lf(\allowbreak{}a, f) implies pts(\allowbreak{}gns-setoid(\allowbreak{}a, f)) == vec(\allowbreak{}a)))
\end{quote}
also: \texttt{gns-setoid-is-setoid} \\
mentioned by 16 result(s)

\smallskip

\noindent\textbf{\texttt{gns-vector(g4ga\_, g4gf\_)}}\quad\textit{functoid}\index{gns-vector(g4ga_, g4gf_)@\texttt{gns-vector(g4ga\_, g4gf\_)}}\\
reads: the GNS cyclic vector of \$2 \\
definition:
\begin{quote}\ttfamily\footnotesize\raggedright\noindent
class(\allowbreak{}gns-setoid(\allowbreak{}g4ga\_, g4gf\_), vone(\allowbreak{}g4ga\_))
\end{quote}
characterized by \texttt{gns-vector-in}  \textit{(proven)}:
\begin{quote}\ttfamily\footnotesize\raggedright\noindent
forall(\allowbreak{}[a], is-unital-involutive-banach-algebra(\allowbreak{}a) implies forall(\allowbreak{}[f], is-positive-lf(\allowbreak{}a, f) implies gns-vector(\allowbreak{}a, f) in vec(\allowbreak{}gns-pre(\allowbreak{}a, f))))
\end{quote}
mentioned by 6 result(s)

\smallskip

\noindent\textbf{\texttt{goal-status}}\quad\textit{tactic}\index{goal-status@\texttt{goal-status}}\\
\texttt{(goal-status)} -- One-line summary: done / N open goals. \\
declared in \texttt{tactics-help.scm}

\smallskip

\noindent\textbf{\texttt{good-approx(hmu\_, hmg\_, hmr\_, hma\_, hmb\_)}}\quad\textit{predicate}\index{good-approx(hmu_, hmg_, hmr_, hma_, hmb_)@\texttt{good-approx(hmu\_, hmg\_, hmr\_, hma\_, hmb\_)}}\\
reads: \$3 is a good approximation to \$2 in \$1 \\
defined by \texttt{good-approx} \\
characterized by \texttt{good-approx}  \textit{(definitional)}:
\begin{quote}\ttfamily\footnotesize\raggedright\noindent
forall(\allowbreak{}[hmu\_, hmg\_, hmr\_, hma\_, hmb\_], good-approx(\allowbreak{}hmu\_, hmg\_, hmr\_, hma\_, hmb\_) iff is-path(\allowbreak{}hmg\_, hma\_, hmb\_) and is-path(\allowbreak{}hmr\_, hma\_, hmb\_) and trace(\allowbreak{}hmg\_, hma\_, hmb\_) subset hmu\_ and trace(\allowbreak{}hmr\_, hma\_, hmb\_) subset hmu\_ and hmr\_(\allowbreak{}hma\_) = hmg\_(\allowbreak{}hma\_) and hmr\_(\allowbreak{}hmb\_) = hmg\_(\allowbreak{}hmb\_) and forsome(\allowbreak{}[hmn\_ in nn], 1 <= hmn\_ and forsome(\allowbreak{}[hmt\_ in fun(\allowbreak{}nn, rr)], hmt\_(\allowbreak{}0) = hma\_ and hmt\_(\allowbreak{}hmn\_) = hmb\_ and forsome(\allowbreak{}[hmd\_ in fun(\allowbreak{}nn, power(\allowbreak{}cc))], forall(\allowbreak{}[hmk\_ in nn], hmk\_ < hmn\_ implies hmt\_(\allowbreak{}hmk\_) < hmt\_(\allowbreak{}succ(\allowbreak{}hmk\_)) and is-open(\allowbreak{}nf-metric-space(\allowbreak{}cc-normed-field), hmd\_(\allowbreak{}hmk\_)) and is-convex-cc(\allowbreak{}hmd\_(\allowbreak{}hmk\_)) and hmd\_(\allowbreak{}hmk\_) subset hmu\_ and forall(\allowbreak{}[hmw\_ in ccint(\allowbreak{}hmt\_(\allowbreak{}hmk\_), hmt\_(\allowbreak{}succ(\allowbreak{}hmk\_)))], hmg\_(\allowbreak{}hmw\_) in hmd\_(\allowbreak{}hmk\_) and hmr\_(\allowbreak{}hmw\_) in hmd\_(\allowbreak{}hmk\_)))))))
\end{quote}
also: \texttt{good-approx-balls} \texttt{good-approx-polygon} \texttt{good-approx-integral} \texttt{good-approx-homotopic} \texttt{polygonal-good-approx} \\
mentioned by 11 result(s)

\smallskip

\noindent\textbf{\texttt{good-sub(m, s, f, t)}}\quad\textit{predicate}\index{good-sub(m, s, f, t)@\texttt{good-sub(m, s, f, t)}}\\
reads: \$3 has a norm-preserving extension from \$2 to \$4, in \$1 \\
defined by \texttt{good-sub} \\
characterized by \texttt{good-sub}  \textit{(definitional)}:
\begin{quote}\ttfamily\footnotesize\raggedright\noindent
forall(\allowbreak{}[m, s, f, t], good-sub(\allowbreak{}m, s, f, t) iff forsome(\allowbreak{}[g\_], npe(\allowbreak{}m, s, f, t, g\_)))
\end{quote}
also: \texttt{good-sub-in-power-cnvs-as-nvs} \texttt{good-sub-in-power} \texttt{good-sub-self-nvs} \texttt{good-sub-submodule} \\
mentioned by 8 result(s)

\smallskip

\noindent\textbf{\texttt{gpow}}\quad\textit{functoid}\index{gpow@\texttt{gpow}}\\
reads: the \$3-th power of \$2 in \$1 \\
defined by \texttt{gpow-zero} \texttt{gpow-succ} \\
characterized by \texttt{gpow-succ}  \textit{(definitional)}:
\begin{quote}\ttfamily\footnotesize\raggedright\noindent
forall(\allowbreak{}[gpg\_, gpx\_, gpn\_ in nn], gpow(\allowbreak{}gpg\_, gpx\_, succ(\allowbreak{}gpn\_)) == (\allowbreak{}opr(\allowbreak{}gpg\_))(\allowbreak{}gpx\_, gpow(\allowbreak{}gpg\_, gpx\_, gpn\_)))
\end{quote}
also: \texttt{gpow-zero} \texttt{gpow-add} \texttt{gpow-one} \texttt{gpow-prod} \texttt{gpow-type} \texttt{hom-cc-gpow} \texttt{gpow-commute} \texttt{qp-gpow-embed} ... (13 in all) \\
mentioned by 46 result(s)

\smallskip

\noindent\textbf{\texttt{graph-adjoint(ubv\_, ubg\_)}}\quad\textit{functoid}\index{graph-adjoint(ubv_, ubg_)@\texttt{graph-adjoint(ubv\_, ubg\_)}}\\
reads: the adjoint of the graph \$2 \\
definition:
\begin{quote}\ttfamily\footnotesize\raggedright\noindent
\{ubp\_ in cartesian(\allowbreak{}vec(\allowbreak{}ubv\_), vec(\allowbreak{}ubv\_)): forall(\allowbreak{}[ubq\_ in ubg\_], (\allowbreak{}ip(\allowbreak{}ubv\_))(\allowbreak{}nth(\allowbreak{}1, ubq\_), nth(\allowbreak{}2, ubp\_)) = (\allowbreak{}ip(\allowbreak{}ubv\_))(\allowbreak{}nth(\allowbreak{}2, ubq\_), nth(\allowbreak{}1, ubp\_)))\}
\end{quote}
characterized by \texttt{graph-adjoint-via-j}  \textit{(proven)}:
\begin{quote}\ttfamily\footnotesize\raggedright\noindent
forall(\allowbreak{}[v, gg], is-complex-inner-product-space(\allowbreak{}v) implies gg subset cartesian(\allowbreak{}vec(\allowbreak{}v), vec(\allowbreak{}v)) implies graph-adjoint(\allowbreak{}v, gg) = perp(\allowbreak{}hsum(\allowbreak{}v, v), graph-j(\allowbreak{}v, gg)) and graph-adjoint(\allowbreak{}v, gg) = graph-j(\allowbreak{}v, perp(\allowbreak{}hsum(\allowbreak{}v, v), gg)))
\end{quote}
also: \texttt{graph-adjoint-closed} \texttt{graph-adjoint-is-set} \texttt{graph-adjoint-j-perp} \texttt{graph-adjoint-linear} \texttt{graph-adjoint-perp-j} \texttt{graph-adjoint-subset} \texttt{graph-adjoint-pair-iff} \texttt{graph-adjoint-of-inverse} ... (13 in all) \\
mentioned by 50 result(s)

\smallskip

\noindent\textbf{\texttt{graph-closure(ubv\_, ubg\_)}}\quad\textit{functoid}\index{graph-closure(ubv_, ubg_)@\texttt{graph-closure(ubv\_, ubg\_)}}\\
reads: the closure of the graph \$2 \\
definition:
\begin{quote}\ttfamily\footnotesize\raggedright\noindent
closure(\allowbreak{}ip-metric-space(\allowbreak{}hsum(\allowbreak{}ubv\_, ubv\_)), ubg\_)
\end{quote}
characterized by \texttt{graph-closure-double-adjoint}  \textit{(proven)}:
\begin{quote}\ttfamily\footnotesize\raggedright\noindent
forall(\allowbreak{}[v, gg], is-hilbert-space(\allowbreak{}v) implies is-linear-graph(\allowbreak{}v, gg) implies graph-closure(\allowbreak{}v, gg) = graph-adjoint(\allowbreak{}v, graph-adjoint(\allowbreak{}v, gg)))
\end{quote}
also: \texttt{graph-closure-is-double-adjoint} \\
mentioned by 8 result(s)

\smallskip

\noindent\textbf{\texttt{graph-compose(ubv\_, ubg\_, ubh\_)}}\quad\textit{functoid}\index{graph-compose(ubv_, ubg_, ubh_)@\texttt{graph-compose(ubv\_, ubg\_, ubh\_)}}\\
reads: the composite graph \$2 \$3 \\
definition:
\begin{quote}\ttfamily\footnotesize\raggedright\noindent
\{ubp\_ in cartesian(\allowbreak{}vec(\allowbreak{}ubv\_), vec(\allowbreak{}ubv\_)): forsome(\allowbreak{}[uby\_ in vec(\allowbreak{}ubv\_)], [nth(\allowbreak{}1, ubp\_), uby\_] in ubh\_ and [uby\_, nth(\allowbreak{}2, ubp\_)] in ubg\_)\}
\end{quote}
characterized by \texttt{graph-compose-subset}  \textit{(proven)}:
\begin{quote}\ttfamily\footnotesize\raggedright\noindent
forall(\allowbreak{}[v, gg, hh], is-complex-inner-product-space(\allowbreak{}v) implies graph-compose(\allowbreak{}v, gg, hh) subset cartesian(\allowbreak{}vec(\allowbreak{}v), vec(\allowbreak{}v)))
\end{quote}
also: \texttt{graph-compose-pair-iff} \\
mentioned by 14 result(s)

\smallskip

\noindent\textbf{\texttt{graph-dom(ubv\_, ubg\_)}}\quad\textit{functoid}\index{graph-dom(ubv_, ubg_)@\texttt{graph-dom(ubv\_, ubg\_)}}\\
reads: the domain of the graph \$2 \\
definition:
\begin{quote}\ttfamily\footnotesize\raggedright\noindent
\{ubx\_ in vec(\allowbreak{}ubv\_): forsome(\allowbreak{}[uby\_ in vec(\allowbreak{}ubv\_)], [ubx\_, uby\_] in ubg\_)\}
\end{quote}
characterized by \texttt{graph-dom-unfold}  \textit{(proven)}:
\begin{quote}\ttfamily\footnotesize\raggedright\noindent
forall(\allowbreak{}[v, uu], graph-dom(\allowbreak{}v, uu) == \{ubx\_ in vec(\allowbreak{}v): forsome(\allowbreak{}[uby\_ in vec(\allowbreak{}v)], [ubx\_, uby\_] in uu)\})
\end{quote}
also: \texttt{graph-dom-submodule} \texttt{graph-dom-of-inverse} \\
mentioned by 16 result(s)

\smallskip

\noindent\textbf{\texttt{graph-inverse(ubv\_, ubg\_)}}\quad\textit{functoid}\index{graph-inverse(ubv_, ubg_)@\texttt{graph-inverse(ubv\_, ubg\_)}}\\
reads: the inverse graph of \$2 \\
definition:
\begin{quote}\ttfamily\footnotesize\raggedright\noindent
\{ubp\_ in cartesian(\allowbreak{}vec(\allowbreak{}ubv\_), vec(\allowbreak{}ubv\_)): [nth(\allowbreak{}2, ubp\_), nth(\allowbreak{}1, ubp\_)] in ubg\_\}
\end{quote}
characterized by \texttt{graph-inverse-is-set}  \textit{(proven)}:
\begin{quote}\ttfamily\footnotesize\raggedright\noindent
forall(\allowbreak{}[v, gg], is-complex-inner-product-space(\allowbreak{}v) implies graph-inverse(\allowbreak{}v, gg) in set)
\end{quote}
also: \texttt{graph-inverse-linear} \texttt{graph-inverse-subset} \texttt{graph-inverse-pair-iff} \texttt{graph-inverse-involutive} \\
mentioned by 21 result(s)

\smallskip

\noindent\textbf{\texttt{graph-j(ubv\_, ubg\_)}}\quad\textit{functoid}\index{graph-j(ubv_, ubg_)@\texttt{graph-j(ubv\_, ubg\_)}}\\
reads: the image of \$2 under J \\
definition:
\begin{quote}\ttfamily\footnotesize\raggedright\noindent
\{ubp\_ in cartesian(\allowbreak{}vec(\allowbreak{}ubv\_), vec(\allowbreak{}ubv\_)): [nth(\allowbreak{}2, ubp\_), (\allowbreak{}vneg(\allowbreak{}ubv\_))(\allowbreak{}nth(\allowbreak{}1, ubp\_))] in ubg\_\}
\end{quote}
characterized by \texttt{graph-j-is-set}  \textit{(proven)}:
\begin{quote}\ttfamily\footnotesize\raggedright\noindent
forall(\allowbreak{}[v, gg], is-complex-inner-product-space(\allowbreak{}v) implies graph-j(\allowbreak{}v, gg) in set)
\end{quote}
also: \texttt{graph-j-subset} \texttt{graph-j-of-perp} \texttt{graph-j-j-linear} \texttt{graph-j-pair-iff} \\
mentioned by 13 result(s)

\smallskip

\noindent\textbf{\texttt{graph-ms(s, a, g)}}\quad\textit{functoid}\index{graph-ms(s, a, g)@\texttt{graph-ms(s, a, g)}}\\
reads: the graph space of \$3 over \$2 in \$1 \\
definition:
\begin{quote}\ttfamily\footnotesize\raggedright\noindent
[a, vnb-lambda(\allowbreak{}[gmu\_, gmv\_], cartesian(\allowbreak{}a, a), (\allowbreak{}dist(\allowbreak{}s))(\allowbreak{}gmu\_, gmv\_) + abs(\allowbreak{}g(\allowbreak{}gmu\_) - g(\allowbreak{}gmv\_)))]
\end{quote}
characterized by \texttt{graph-ms-pts}  \textit{(proven)}:
\begin{quote}\ttfamily\footnotesize\raggedright\noindent
forall(\allowbreak{}[s, a, g], a in set implies pts(\allowbreak{}graph-ms(\allowbreak{}s, a, g)) == a)
\end{quote}
also: \texttt{graph-ms-dist} \texttt{graph-ms-near} \texttt{graph-ms-polish} \texttt{graph-ms-complete} \texttt{graph-ms-dist-slot} \texttt{graph-ms-homeomorphism} \texttt{graph-ms-is-metric-space} \\
mentioned by 11 result(s)

\smallskip

\noindent\textbf{\texttt{graph-of(gfx\_, gfy\_, gff\_)}}\quad\textit{functoid}\index{graph-of(gfx_, gfy_, gff_)@\texttt{graph-of(gfx\_, gfy\_, gff\_)}}\\
reads: the graph of \$3 from \$1 to \$2 \\
definition:
\begin{quote}\ttfamily\footnotesize\raggedright\noindent
\{grp\_ in cartesian(\allowbreak{}gfx\_, gfy\_): gff\_(\allowbreak{}nth(\allowbreak{}1, grp\_)) = nth(\allowbreak{}2, grp\_)\}
\end{quote}
characterized by \texttt{graph-of-unfold}  \textit{(proven)}:
\begin{quote}\ttfamily\footnotesize\raggedright\noindent
forall(\allowbreak{}[a, b, f], graph-of(\allowbreak{}a, b, f) == \{grp\_ in cartesian(\allowbreak{}a, b): f(\allowbreak{}nth(\allowbreak{}1, grp\_)) = nth(\allowbreak{}2, grp\_)\})
\end{quote}
mentioned by 3 result(s)

\smallskip

\noindent\textbf{\texttt{graph-ran(ubv\_, ubg\_)}}\quad\textit{functoid}\index{graph-ran(ubv_, ubg_)@\texttt{graph-ran(ubv\_, ubg\_)}}\\
reads: the range of the graph \$2 \\
definition:
\begin{quote}\ttfamily\footnotesize\raggedright\noindent
\{uby\_ in vec(\allowbreak{}ubv\_): forsome(\allowbreak{}[ubx\_ in vec(\allowbreak{}ubv\_)], [ubx\_, uby\_] in ubg\_)\}
\end{quote}
characterized by \texttt{graph-ran-unfold}  \textit{(proven)}:
\begin{quote}\ttfamily\footnotesize\raggedright\noindent
forall(\allowbreak{}[v, uu], graph-ran(\allowbreak{}v, uu) == \{uby\_ in vec(\allowbreak{}v): forsome(\allowbreak{}[ubx\_ in vec(\allowbreak{}v)], [ubx\_, uby\_] in uu)\})
\end{quote}
also: \texttt{graph-ran-submodule} \\
mentioned by 16 result(s)

\smallskip

\noindent\textbf{\texttt{greedy-chain(phi, grd, porel)}}\quad\textit{functoid}\index{greedy-chain(phi, grd, porel)@\texttt{greedy-chain(phi, grd, porel)}}\\
definition:
\begin{quote}\ttfamily\footnotesize\raggedright\noindent
zkept(\allowbreak{}phi, grd, porel, card(\allowbreak{}grd))
\end{quote}
characterized by \texttt{greedy-chain-is-chain}  \textit{(proven)}:
\begin{quote}\ttfamily\footnotesize\raggedright\noindent
forall(\allowbreak{}[phi, grd, porel], is-partial-order(\allowbreak{}grd, porel) implies is-chain(\allowbreak{}grd, porel, greedy-chain(\allowbreak{}phi, grd, porel)))
\end{quote}
mentioned by 1 result(s)

\smallskip

\noindent\textbf{\texttt{grind}}\quad\textit{tactic}\index{grind@\texttt{grind}}\\
\texttt{(grind)} -- Saturate the no-choice moves at the focus: decompose the goal connective (di) and break open the hypotheses (mac-h*), repeatedly, until neither fires.  The whole introduce-everything prefix in one step. \\
declared in \texttt{tactics-help.scm}

\smallskip

\noindent\textbf{\texttt{grind-and-mp}}\quad\textit{tactic}\index{grind-and-mp@\texttt{grind-and-mp}}\\
\texttt{(grind-and-mp)} -- Tidy the goal with grind, then close it by modus ponens -- the one-button form of the syllogism, for a sentence typed exactly as it reads. \\
declared in \texttt{tactics-help.scm}

\smallskip

\noindent\textbf{\texttt{group}}\quad\textit{structure}\index{group@\texttt{group}}\\
predicate \texttt{is-group}; slots: carr opr iden inv

\smallskip

\noindent\textbf{\texttt{group-famprod}}\quad\textit{functoid}\index{group-famprod@\texttt{group-famprod}}\\
reads: the product of the groups \$1(1), ..., \$1(\$2) \\
defined by \texttt{group-famprod-zero} \texttt{group-famprod-succ} \\
characterized by \texttt{group-famprod-succ}  \textit{(definitional)}:
\begin{quote}\ttfamily\footnotesize\raggedright\noindent
forall(\allowbreak{}[gfam\_, gfn\_ in nn], group-famprod(\allowbreak{}gfam\_, succ(\allowbreak{}gfn\_)) == group-prod(\allowbreak{}group-famprod(\allowbreak{}gfam\_, gfn\_), gfam\_(\allowbreak{}succ(\allowbreak{}gfn\_))))
\end{quote}
also: \texttt{group-famprod-zero} \texttt{group-famprod-facts} \texttt{group-famprod-card-coprime} \texttt{group-famprod-card-divisible} \\
mentioned by 9 result(s)

\smallskip

\noindent\textbf{\texttt{group-famtuple}}\quad\textit{functoid}\index{group-famtuple@\texttt{group-famtuple}}\\
reads: the tuple (\$1(1), ..., \$1(\$2)) \\
defined by \texttt{group-famtuple-zero} \texttt{group-famtuple-succ} \\
characterized by \texttt{group-famtuple-succ}  \textit{(definitional)}:
\begin{quote}\ttfamily\footnotesize\raggedright\noindent
forall(\allowbreak{}[xfam\_, xfn\_ in nn], group-famtuple(\allowbreak{}xfam\_, succ(\allowbreak{}xfn\_)) == [group-famtuple(\allowbreak{}xfam\_, xfn\_), xfam\_(\allowbreak{}succ(\allowbreak{}xfn\_))])
\end{quote}
also: \texttt{group-famtuple-zero} \\
mentioned by 7 result(s)

\smallskip

\noindent\textbf{\texttt{group-prod(x, y)}}\quad\textit{functoid}\index{group-prod(x, y)@\texttt{group-prod(x, y)}}\\
reads: the direct product of the groups \$1 and \$2 \\
definition:
\begin{quote}\ttfamily\footnotesize\raggedright\noindent
[cartesian(\allowbreak{}carr(\allowbreak{}x), carr(\allowbreak{}y)), vnb-lambda(\allowbreak{}[gpp\_, gpq\_], cartesian(\allowbreak{}cartesian(\allowbreak{}carr(\allowbreak{}x), carr(\allowbreak{}y)), cartesian(\allowbreak{}carr(\allowbreak{}x), carr(\allowbreak{}y))), [(\allowbreak{}opr(\allowbreak{}x))(\allowbreak{}nth(\allowbreak{}1, gpp\_), nth(\allowbreak{}1, gpq\_)), (\allowbreak{}opr(\allowbreak{}y))(\allowbreak{}nth(\allowbreak{}2, gpp\_), nth(\allowbreak{}2, gpq\_))]), [iden(\allowbreak{}x), iden(\allowbreak{}y)], vnb-lambda(\allowbreak{}gpp\_, cartesian(\allowbreak{}carr(\allowbreak{}x), carr(\allowbreak{}y)), [(\allowbreak{}inv(\allowbreak{}x))(\allowbreak{}nth(\allowbreak{}1, gpp\_)), (\allowbreak{}inv(\allowbreak{}y))(\allowbreak{}nth(\allowbreak{}2, gpp\_))])]
\end{quote}
characterized by \texttt{group-prod-inv}  \textit{(proven)}:
\begin{quote}\ttfamily\footnotesize\raggedright\noindent
forall(\allowbreak{}[x], is-group(\allowbreak{}x) implies forall(\allowbreak{}[y], is-group(\allowbreak{}y) implies inv(\allowbreak{}group-prod(\allowbreak{}x, y)) == vnb-lambda(\allowbreak{}gpp\_, cartesian(\allowbreak{}carr(\allowbreak{}x), carr(\allowbreak{}y)), [(\allowbreak{}inv(\allowbreak{}x))(\allowbreak{}nth(\allowbreak{}1, gpp\_)), (\allowbreak{}inv(\allowbreak{}y))(\allowbreak{}nth(\allowbreak{}2, gpp\_))])))
\end{quote}
also: \texttt{group-prod-opr} \texttt{group-prod-carr} \texttt{group-prod-iden} \texttt{group-prod-inv-pair} \texttt{group-prod-is-group} \texttt{group-prod-opr-pair} \texttt{group-prod-generator-left} \\
mentioned by 22 result(s)

\smallskip

\noindent\textbf{\texttt{group-zpow(g, x, k)}}\quad\textit{functoid}\index{group-zpow(g, x, k)@\texttt{group-zpow(g, x, k)}}\\
reads: the \$3-th power of \$2 in \$1 (an integer exponent) \\
definition:
\begin{quote}\ttfamily\footnotesize\raggedright\noindent
if(\allowbreak{}0 <= k, gpow(\allowbreak{}g, x, k), (\allowbreak{}inv(\allowbreak{}g))(\allowbreak{}gpow(\allowbreak{}g, x, 0 - k)))
\end{quote}
mentioned by 40 result(s)

\smallskip

\noindent\textbf{\texttt{grp-monoid(gmg\_)}}\quad\textit{functoid}\index{grp-monoid(gmg_)@\texttt{grp-monoid(gmg\_)}}\\
reads: the monoid underlying the group \$1 \\
definition:
\begin{quote}\ttfamily\footnotesize\raggedright\noindent
[carr(\allowbreak{}gmg\_), opr(\allowbreak{}gmg\_), iden(\allowbreak{}gmg\_)]
\end{quote}
characterized by \texttt{grp-monoid-opr}  \textit{(proven)}:
\begin{quote}\ttfamily\footnotesize\raggedright\noindent
forall(\allowbreak{}[s], is-group(\allowbreak{}s) implies opr(\allowbreak{}grp-monoid(\allowbreak{}s)) == opr(\allowbreak{}s))
\end{quote}
also: \texttt{grp-monoid-carr} \texttt{grp-monoid-iden} \texttt{grp-monoid-is-monoid} \\
mentioned by 9 result(s)

\smallskip


\section*{H}
\addcontentsline{toc}{section}{H}

\noindent\textbf{\texttt{haar-averages(hag\_)}}\quad\textit{functoid}\index{haar-averages(hag_)@\texttt{haar-averages(hag\_)}}\\
reads: the averaging operators of the compact metric group \$1 \\
definition:
\begin{quote}\ttfamily\footnotesize\raggedright\noindent
\{haq\_ in fun(\allowbreak{}cont-fns(\allowbreak{}metric-group-ms(\allowbreak{}hag\_)), cont-fns(\allowbreak{}metric-group-ms(\allowbreak{}hag\_))): forall(\allowbreak{}[hau\_ in cont-fns(\allowbreak{}metric-group-ms(\allowbreak{}hag\_)), hav\_ in cont-fns(\allowbreak{}metric-group-ms(\allowbreak{}hag\_))], haq\_(\allowbreak{}(\allowbreak{}vadd(\allowbreak{}ck-nvs(\allowbreak{}metric-group-ms(\allowbreak{}hag\_))))(\allowbreak{}hau\_, hav\_)) = (\allowbreak{}vadd(\allowbreak{}ck-nvs(\allowbreak{}metric-group-ms(\allowbreak{}hag\_))))(\allowbreak{}haq\_(\allowbreak{}hau\_), haq\_(\allowbreak{}hav\_))) and forall(\allowbreak{}[har\_ in rr, hau\_ in cont-fns(\allowbreak{}metric-group-ms(\allowbreak{}hag\_))], haq\_(\allowbreak{}(\allowbreak{}act(\allowbreak{}ck-nvs(\allowbreak{}metric-group-ms(\allowbreak{}hag\_))))(\allowbreak{}har\_, hau\_)) = (\allowbreak{}act(\allowbreak{}ck-nvs(\allowbreak{}metric-group-ms(\allowbreak{}hag\_))))(\allowbreak{}har\_, haq\_(\allowbreak{}hau\_))) and forall(\allowbreak{}[hau\_ in cont-fns(\allowbreak{}metric-group-ms(\allowbreak{}hag\_)), hav\_ in cont-fns(\allowbreak{}metric-group-ms(\allowbreak{}hag\_))], (\allowbreak{}dist(\allowbreak{}nvs-metric-space(\allowbreak{}ck-nvs(\allowbreak{}metric-group-ms(\allowbreak{}hag\_)))))(\allowbreak{}haq\_(\allowbreak{}hau\_), haq\_(\allowbreak{}hav\_)) <= (\allowbreak{}dist(\allowbreak{}nvs-metric-space(\allowbreak{}ck-nvs(\allowbreak{}metric-group-ms(\allowbreak{}hag\_)))))(\allowbreak{}hau\_, hav\_)) and forall(\allowbreak{}[hac\_ in rr], haq\_(\allowbreak{}vnb-lambda(\allowbreak{}hix\_, pts(\allowbreak{}metric-group-ms(\allowbreak{}hag\_)), hac\_)) = vnb-lambda(\allowbreak{}hix\_, pts(\allowbreak{}metric-group-ms(\allowbreak{}hag\_)), hac\_)) and forall(\allowbreak{}[hau\_ in cont-fns(\allowbreak{}metric-group-ms(\allowbreak{}hag\_))], haq\_(\allowbreak{}hau\_) in conv-hull(\allowbreak{}ck-nvs(\allowbreak{}metric-group-ms(\allowbreak{}hag\_)), image(\allowbreak{}vnb-lambda(\allowbreak{}hqa\_, pts(\allowbreak{}metric-group-ms(\allowbreak{}hag\_)), vnb-lambda(\allowbreak{}hix\_, pts(\allowbreak{}metric-group-ms(\allowbreak{}hag\_)), hau\_(\allowbreak{}(\allowbreak{}opr(\allowbreak{}hag\_))(\allowbreak{}hqa\_, hix\_)))), pts(\allowbreak{}metric-group-ms(\allowbreak{}hag\_))))) and forall(\allowbreak{}[hab\_ in pts(\allowbreak{}metric-group-ms(\allowbreak{}hag\_)), hau\_ in cont-fns(\allowbreak{}metric-group-ms(\allowbreak{}hag\_))], haq\_(\allowbreak{}vnb-lambda(\allowbreak{}hix\_, pts(\allowbreak{}metric-group-ms(\allowbreak{}hag\_)), hau\_(\allowbreak{}(\allowbreak{}opr(\allowbreak{}hag\_))(\allowbreak{}hix\_, hab\_)))) = vnb-lambda(\allowbreak{}hix\_, pts(\allowbreak{}metric-group-ms(\allowbreak{}hag\_)), (\allowbreak{}haq\_(\allowbreak{}hau\_))(\allowbreak{}(\allowbreak{}opr(\allowbreak{}hag\_))(\allowbreak{}hix\_, hab\_))))\}
\end{quote}
mentioned by 11 result(s)

\smallskip

\noindent\textbf{\texttt{haar-integral(hng\_)}}\quad\textit{functoid}\index{haar-integral(hng_)@\texttt{haar-integral(hng\_)}}\\
reads: the Haar integral of the compact metric group \$1 \\
definition:
\begin{quote}\ttfamily\footnotesize\raggedright\noindent
iota(\allowbreak{}hnu\_, is-haar-integral(\allowbreak{}hng\_, hnu\_))
\end{quote}
characterized by \texttt{haar-integral-value}  \textit{(proven)}:
\begin{quote}\ttfamily\footnotesize\raggedright\noindent
forall(\allowbreak{}[g, f], is-metric-group(\allowbreak{}g) implies is-compact(\allowbreak{}metric-group-ms(\allowbreak{}g)) implies f in cont-fns(\allowbreak{}metric-group-ms(\allowbreak{}g)) implies (\allowbreak{}haar-integral(\allowbreak{}g))(\allowbreak{}f) = haar-mean(\allowbreak{}g, f))
\end{quote}
also: \texttt{haar-integral-unfold} \texttt{haar-integral-is-haar} \\
mentioned by 5 result(s)

\smallskip

\noindent\textbf{\texttt{haar-mean(hmg\_, hmf\_)}}\quad\textit{functoid}\index{haar-mean(hmg_, hmf_)@\texttt{haar-mean(hmg\_, hmf\_)}}\\
reads: the Haar mean of \$2 on the compact metric group \$1 \\
definition:
\begin{quote}\ttfamily\footnotesize\raggedright\noindent
iota(\allowbreak{}hmc\_, hmc\_ in rr and vnb-lambda(\allowbreak{}hix\_, pts(\allowbreak{}metric-group-ms(\allowbreak{}hmg\_)), hmc\_) in closure(\allowbreak{}nvs-metric-space(\allowbreak{}ck-nvs(\allowbreak{}metric-group-ms(\allowbreak{}hmg\_))), conv-hull(\allowbreak{}ck-nvs(\allowbreak{}metric-group-ms(\allowbreak{}hmg\_)), image(\allowbreak{}vnb-lambda(\allowbreak{}hqa\_, pts(\allowbreak{}metric-group-ms(\allowbreak{}hmg\_)), vnb-lambda(\allowbreak{}hix\_, pts(\allowbreak{}metric-group-ms(\allowbreak{}hmg\_)), hmf\_(\allowbreak{}(\allowbreak{}opr(\allowbreak{}hmg\_))(\allowbreak{}hqa\_, hix\_)))), pts(\allowbreak{}metric-group-ms(\allowbreak{}hmg\_))))))
\end{quote}
characterized by \texttt{haar-mean-eq}  \textit{(proven)}:
\begin{quote}\ttfamily\footnotesize\raggedright\noindent
forall(\allowbreak{}[g, f, c], is-metric-group(\allowbreak{}g) implies is-compact(\allowbreak{}metric-group-ms(\allowbreak{}g)) implies f in cont-fns(\allowbreak{}metric-group-ms(\allowbreak{}g)) implies c in rr implies vnb-lambda(\allowbreak{}hix\_, pts(\allowbreak{}metric-group-ms(\allowbreak{}g)), c) in closure(\allowbreak{}nvs-metric-space(\allowbreak{}ck-nvs(\allowbreak{}metric-group-ms(\allowbreak{}g))), conv-hull(\allowbreak{}ck-nvs(\allowbreak{}metric-group-ms(\allowbreak{}g)), image(\allowbreak{}vnb-lambda(\allowbreak{}hqa\_, pts(\allowbreak{}metric-group-ms(\allowbreak{}g)), vnb-lambda(\allowbreak{}hix\_, pts(\allowbreak{}metric-group-ms(\allowbreak{}g)), f(\allowbreak{}(\allowbreak{}opr(\allowbreak{}g))(\allowbreak{}hqa\_, hix\_)))), pts(\allowbreak{}metric-group-ms(\allowbreak{}g))))) implies haar-mean(\allowbreak{}g, f) = c)
\end{quote}
also: \texttt{haar-mean-act} \texttt{haar-mean-add} \texttt{haar-mean-avg} \texttt{haar-mean-opp} \texttt{haar-mean-spec} \texttt{haar-mean-const} \texttt{haar-mean-nonneg} \texttt{haar-mean-unfold} ... (12 in all) \\
mentioned by 27 result(s)

\smallskip

\noindent\textbf{\texttt{half-sq-iter}}\quad\textit{functoid}\index{half-sq-iter@\texttt{half-sq-iter}}\\
reads: the n-th square-root iterate of \$2 \\
defined by \texttt{half-sq-iter-zero} \texttt{half-sq-iter-succ} \\
characterized by \texttt{half-sq-iter-succ}  \textit{(definitional)}:
\begin{quote}\ttfamily\footnotesize\raggedright\noindent
forall(\allowbreak{}[g4a\_, g4y\_, g4n\_ in nn], half-sq-iter(\allowbreak{}g4a\_, g4y\_, succ(\allowbreak{}g4n\_)) == (\allowbreak{}act(\allowbreak{}g4a\_))(\allowbreak{}num(\allowbreak{}1/2), (\allowbreak{}vadd(\allowbreak{}g4a\_))(\allowbreak{}g4y\_, (\allowbreak{}vmul(\allowbreak{}g4a\_))(\allowbreak{}half-sq-iter(\allowbreak{}g4a\_, g4y\_, g4n\_), half-sq-iter(\allowbreak{}g4a\_, g4y\_, g4n\_)))))
\end{quote}
also: \texttt{half-sq-iter-zero} \\
mentioned by 7 result(s)

\smallskip

\noindent\textbf{\texttt{has-characteristic(ghc\_, gcp\_)}}\quad\textit{predicate}\index{has-characteristic(ghc_, gcp_)@\texttt{has-characteristic(ghc\_, gcp\_)}}\\
reads: \$1 has characteristic \$2 \\
defined by \texttt{has-characteristic} \\
characterized by \texttt{has-characteristic}  \textit{(definitional)}:
\begin{quote}\ttfamily\footnotesize\raggedright\noindent
forall(\allowbreak{}[ghc\_, gcp\_], has-characteristic(\allowbreak{}ghc\_, gcp\_) iff is-prime-nat(\allowbreak{}gcp\_) and zz-act(\allowbreak{}ring-additive-ag(\allowbreak{}ghc\_), gcp\_, one(\allowbreak{}ghc\_)) = zero(\allowbreak{}ghc\_))
\end{quote}
mentioned by 24 result(s)

\smallskip

\noindent\textbf{\texttt{has-closed-graph(m1, fam1, m2, fam2, tt)}}\quad\textit{predicate}\index{has-closed-graph(m1, fam1, m2, fam2, tt)@\texttt{has-closed-graph(m1, fam1, m2, fam2, tt)}}\\
reads: \$5 has a closed graph from (\$1, \$2) to (\$3, \$4) \\
defined by \texttt{has-closed-graph} \\
characterized by \texttt{has-closed-graph}  \textit{(definitional)}:
\begin{quote}\ttfamily\footnotesize\raggedright\noindent
forall(\allowbreak{}[m1, fam1, m2, fam2, tt], has-closed-graph(\allowbreak{}m1, fam1, m2, fam2, tt) iff is-linear-map(\allowbreak{}m1, m2, tt) and forall(\allowbreak{}[seq in fun(\allowbreak{}nn, vec(\allowbreak{}m1)), p1 in vec(\allowbreak{}m1), p2 in vec(\allowbreak{}m2)], fr-conv(\allowbreak{}m1, fam1, seq, p1) and fr-conv(\allowbreak{}m2, fam2, vnb-lambda(\allowbreak{}i, nn, tt(\allowbreak{}seq(\allowbreak{}i))), p2) implies p2 = tt(\allowbreak{}p1)))
\end{quote}
mentioned by 4 result(s)

\smallskip

\noindent\textbf{\texttt{has-countable-nbhd-base(cbs\_, cbx\_)}}\quad\textit{predicate}\index{has-countable-nbhd-base(cbs_, cbx_)@\texttt{has-countable-nbhd-base(cbs\_, cbx\_)}}\\
reads: \$2 has a countable fundamental system of neighbourhoods in \$1 \\
defined by \texttt{has-countable-nbhd-base} \\
characterized by \texttt{has-countable-nbhd-base}  \textit{(definitional)}:
\begin{quote}\ttfamily\footnotesize\raggedright\noindent
forall(\allowbreak{}[cbs\_, cbx\_], has-countable-nbhd-base(\allowbreak{}cbs\_, cbx\_) iff forsome(\allowbreak{}[cbn\_ in fun(\allowbreak{}nn, opens(\allowbreak{}cbs\_))], forall(\allowbreak{}[cbk\_ in nn], cbx\_ in cbn\_(\allowbreak{}cbk\_)) and forall(\allowbreak{}[cbw\_ in opens(\allowbreak{}cbs\_)], cbx\_ in cbw\_ implies forsome(\allowbreak{}[cbk\_ in nn], cbn\_(\allowbreak{}cbk\_) subset cbw\_))))
\end{quote}
mentioned by 4 result(s)

\smallskip

\noindent\textbf{\texttt{has-deriv-at(f, a, l)}}\quad\textit{predicate}\index{has-deriv-at(f, a, l)@\texttt{has-deriv-at(f, a, l)}}\\
reads: \$1 has derivative \$3 at \$2 \\
defined by \texttt{has-deriv-at} \\
characterized by \texttt{has-deriv-at}  \textit{(definitional)}:
\begin{quote}\ttfamily\footnotesize\raggedright\noindent
forall(\allowbreak{}[f, a, l], has-deriv-at(\allowbreak{}f, a, l) iff forsome(\allowbreak{}[ice\_], pos-rr(\allowbreak{}ice\_) and is-diff-on(\allowbreak{}rr-normed-field, ooint(\allowbreak{}a - ice\_, a + ice\_), restrict(\allowbreak{}f, ooint(\allowbreak{}a - ice\_, a + ice\_)), a, l)))
\end{quote}
also: \texttt{has-deriv-at-sum} \texttt{has-deriv-at-chain} \texttt{has-deriv-at-const} \texttt{has-deriv-at-in-rr} \texttt{has-deriv-at-intro} \texttt{has-deriv-at-local} \texttt{has-deriv-at-affine} \texttt{has-deriv-at-shrink} ... (26 in all) \\
mentioned by 79 result(s)

\smallskip

\noindent\textbf{\texttt{has-deriv-at-within(f, w, x, l)}}\quad\textit{predicate}\index{has-deriv-at-within(f, w, x, l)@\texttt{has-deriv-at-within(f, w, x, l)}}\\
reads: \$1 has derivative \$4 at \$3 within \$2 \\
defined by \texttt{has-deriv-at-within} \\
characterized by \texttt{has-deriv-at-within}  \textit{(definitional)}:
\begin{quote}\ttfamily\footnotesize\raggedright\noindent
forall(\allowbreak{}[f, w, x, l], has-deriv-at-within(\allowbreak{}f, w, x, l) iff restrict(\allowbreak{}f, w) in fun(\allowbreak{}w, rr) and x in w and l in rr and forsome(\allowbreak{}[osp\_ in fun(\allowbreak{}w, rr)], is-continuous-at(\allowbreak{}subspace-ms(\allowbreak{}rr-ms, w), rr-ms, osp\_, x) and osp\_(\allowbreak{}x) = l and forall(\allowbreak{}[ost\_ in w], f(\allowbreak{}ost\_) - f(\allowbreak{}x) = osp\_(\allowbreak{}ost\_) * (\allowbreak{}ost\_ - x))))
\end{quote}
mentioned by 25 result(s)

\smallskip

\noindent\textbf{\texttt{has-deriv-on-ccint(hdcg\_, hdcd\_, hdca\_, hdcb\_)}}\quad\textit{predicate}\index{has-deriv-on-ccint(hdcg_, hdcd_, hdca_, hdcb_)@\texttt{has-deriv-on-ccint(hdcg\_, hdcd\_, hdca\_, hdcb\_)}}\\
reads: \$2 is the derivative of \$1 on [\$3, \$4] (one-sided at the end points) \\
defined by \texttt{has-deriv-on-ccint} \\
characterized by \texttt{has-deriv-on-ccint}  \textit{(definitional)}:
\begin{quote}\ttfamily\footnotesize\raggedright\noindent
forall(\allowbreak{}[hdcg\_, hdcd\_, hdca\_, hdcb\_], has-deriv-on-ccint(\allowbreak{}hdcg\_, hdcd\_, hdca\_, hdcb\_) iff hdca\_ in rr and hdcb\_ in rr and hdca\_ < hdcb\_ and hdcg\_ in fun(\allowbreak{}ccint(\allowbreak{}hdca\_, hdcb\_), rr) and hdcd\_ in fun(\allowbreak{}ccint(\allowbreak{}hdca\_, hdcb\_), rr) and forall(\allowbreak{}[hdct\_ in ooint(\allowbreak{}hdca\_, hdcb\_)], has-deriv-at(\allowbreak{}hdcg\_, hdct\_, hdcd\_(\allowbreak{}hdct\_))) and has-right-deriv-at(\allowbreak{}hdcg\_, hdca\_, hdcd\_(\allowbreak{}hdca\_)) and has-left-deriv-at(\allowbreak{}hdcg\_, hdcb\_, hdcd\_(\allowbreak{}hdcb\_)))
\end{quote}
mentioned by 7 result(s)

\smallskip

\noindent\textbf{\texttt{has-dir-deriv(m, f, a, eta, l)}}\quad\textit{predicate}\index{has-dir-deriv(m, f, a, eta, l)@\texttt{has-dir-deriv(m, f, a, eta, l)}}\\
reads: \$2 has directional derivative \$5 at \$3 in the direction \$4 \\
defined by \texttt{has-dir-deriv} \\
characterized by \texttt{has-dir-deriv}  \textit{(definitional)}:
\begin{quote}\ttfamily\footnotesize\raggedright\noindent
forall(\allowbreak{}[m, f, a, eta, l], has-dir-deriv(\allowbreak{}m, f, a, eta, l) iff has-deriv-at(\allowbreak{}seg-curve(\allowbreak{}m, f, a, eta), 0, l))
\end{quote}
also: \texttt{has-dir-deriv-add} \texttt{has-dir-deriv-step} \texttt{has-dir-deriv-scale} \texttt{has-dir-deriv-shift} \\
mentioned by 14 result(s)

\smallskip

\noindent\textbf{\texttt{has-div-remainder(s, dg)}}\quad\textit{predicate}\index{has-div-remainder(s, dg)@\texttt{has-div-remainder(s, dg)}}\\
reads: \$1 has division with remainder for the degree function \$2 \\
defined by \texttt{has-div-remainder} \\
characterized by \texttt{has-div-remainder}  \textit{(definitional)}:
\begin{quote}\ttfamily\footnotesize\raggedright\noindent
forall(\allowbreak{}[s, dg], has-div-remainder(\allowbreak{}s, dg) iff forall(\allowbreak{}[a\_ in carr(\allowbreak{}s), b in carr(\allowbreak{}s)], not(\allowbreak{}b = zero(\allowbreak{}s)) implies forsome(\allowbreak{}[q in carr(\allowbreak{}s), r in carr(\allowbreak{}s)], a\_ = (\allowbreak{}add(\allowbreak{}s))(\allowbreak{}(\allowbreak{}mul(\allowbreak{}s))(\allowbreak{}q, b), r) and (\allowbreak{}r = zero(\allowbreak{}s) or succ(\allowbreak{}dg(\allowbreak{}r)) <= dg(\allowbreak{}b)))))
\end{quote}
mentioned by 8 result(s)

\smallskip

\noindent\textbf{\texttt{has-finite-subcover(s, c)}}\quad\textit{predicate}\index{has-finite-subcover(s, c)@\texttt{has-finite-subcover(s, c)}}\\
reads: a finite subfamily of \$2 covers \$1 \\
defined by \texttt{has-finite-subcover} \\
characterized by \texttt{has-finite-subcover}  \textit{(definitional)}:
\begin{quote}\ttfamily\footnotesize\raggedright\noindent
forall(\allowbreak{}[s, c], has-finite-subcover(\allowbreak{}s, c) iff forsome(\allowbreak{}[tpg\_ in set], tpg\_ subset c and card(\allowbreak{}tpg\_) in nn and big-union(\allowbreak{}tpv\_, tpg\_, tpv\_) == pts(\allowbreak{}s)))
\end{quote}
also: \texttt{has-finite-subcover-classes} \\
mentioned by 7 result(s)

\smallskip

\noindent\textbf{\texttt{has-fip(s, c)}}\quad\textit{predicate}\index{has-fip(s, c)@\texttt{has-fip(s, c)}}\\
reads: \$2 has the finite intersection property in \$1 \\
defined by \texttt{has-fip} \\
characterized by \texttt{has-fip}  \textit{(definitional)}:
\begin{quote}\ttfamily\footnotesize\raggedright\noindent
forall(\allowbreak{}[s, c], has-fip(\allowbreak{}s, c) iff is-metric-space(\allowbreak{}s) and forsome(\allowbreak{}[a0], a0 in c) and forall(\allowbreak{}[a in c], is-closed(\allowbreak{}s, a)) and forall(\allowbreak{}[f], f subset c implies card(\allowbreak{}f) in nn implies forsome(\allowbreak{}[a0], a0 in f) implies forsome(\allowbreak{}[p], p in intersection-of(\allowbreak{}f))))
\end{quote}
mentioned by 6 result(s)

\smallskip

\noindent\textbf{\texttt{has-inverses(op, unit, invop, crr)}}\quad\textit{predicate}\index{has-inverses(op, unit, invop, crr)@\texttt{has-inverses(op, unit, invop, crr)}}\\
reads: every element of \$4 has an inverse under \$1, given by \$3, with unit \$2 \\
defined by \texttt{has-inverses} \\
characterized by \texttt{has-inverses}  \textit{(definitional)}:
\begin{quote}\ttfamily\footnotesize\raggedright\noindent
forall(\allowbreak{}[op, unit, invop, crr], has-inverses(\allowbreak{}op, unit, invop, crr) iff forall(\allowbreak{}[u in crr], op(\allowbreak{}invop(\allowbreak{}u), u) = unit and op(\allowbreak{}u, invop(\allowbreak{}u)) = unit))
\end{quote}
mentioned by 62 result(s)

\smallskip

\noindent\textbf{\texttt{has-left-deriv-at(f, x, l)}}\quad\textit{predicate}\index{has-left-deriv-at(f, x, l)@\texttt{has-left-deriv-at(f, x, l)}}\\
reads: \$1 has left derivative \$3 at \$2 \\
defined by \texttt{has-left-deriv-at} \\
characterized by \texttt{has-left-deriv-at}  \textit{(definitional)}:
\begin{quote}\ttfamily\footnotesize\raggedright\noindent
forall(\allowbreak{}[f, x, l], has-left-deriv-at(\allowbreak{}f, x, l) iff forsome(\allowbreak{}[ose\_], pos-rr(\allowbreak{}ose\_) and has-deriv-at-within(\allowbreak{}f, ocint(\allowbreak{}x - ose\_, x), x, l)))
\end{quote}
also: \texttt{has-left-deriv-at-sum} \texttt{has-left-deriv-at-in-rr} \texttt{has-left-deriv-at-intro} \texttt{has-left-deriv-at-local} \texttt{has-left-deriv-at-shrink} \texttt{has-left-deriv-at-unique} \texttt{has-left-deriv-at-pt-in-rr} \\
mentioned by 20 result(s)

\smallskip

\noindent\textbf{\texttt{has-mean-value-property(mvu\_, mvf\_)}}\quad\textit{predicate}\index{has-mean-value-property(mvu_, mvf_)@\texttt{has-mean-value-property(mvu\_, mvf\_)}}\\
reads: \$2 has the mean value property on \$1 \\
defined by \texttt{has-mean-value-property} \\
characterized by \texttt{has-mean-value-property}  \textit{(definitional)}:
\begin{quote}\ttfamily\footnotesize\raggedright\noindent
forall(\allowbreak{}[mvu\_, mvf\_], has-mean-value-property(\allowbreak{}mvu\_, mvf\_) iff is-open(\allowbreak{}nf-metric-space(\allowbreak{}cc-normed-field), mvu\_) and mvf\_ in fun(\allowbreak{}mvu\_, cc) and forall(\allowbreak{}[rgy\_ in mvu\_], is-continuous-at(\allowbreak{}subspace-ms(\allowbreak{}nf-metric-space(\allowbreak{}cc-normed-field), mvu\_), nf-metric-space(\allowbreak{}cc-normed-field), mvf\_, rgy\_)) and forall(\allowbreak{}[mvc\_ in cc, mvr\_ in rr], 0 < mvr\_ implies closure(\allowbreak{}cc-ms, ball(\allowbreak{}cc-ms, mvc\_, mvr\_)) subset mvu\_ implies mvf\_(\allowbreak{}mvc\_) = recip(\allowbreak{}2 * pi) * cc-int(\allowbreak{}vnb-lambda(\allowbreak{}mvt\_, ccint(\allowbreak{}0, 2 * pi), mvf\_(\allowbreak{}mvc\_ + mvr\_ * cc-exp(\allowbreak{}1i * mvt\_))), 0, 2 * pi)))
\end{quote}
also: \texttt{holomorphic-has-mean-value-property} \\
mentioned by 10 result(s)

\smallskip

\noindent\textbf{\texttt{has-right-deriv-at(f, x, l)}}\quad\textit{predicate}\index{has-right-deriv-at(f, x, l)@\texttt{has-right-deriv-at(f, x, l)}}\\
reads: \$1 has right derivative \$3 at \$2 \\
defined by \texttt{has-right-deriv-at} \\
characterized by \texttt{has-right-deriv-at}  \textit{(definitional)}:
\begin{quote}\ttfamily\footnotesize\raggedright\noindent
forall(\allowbreak{}[f, x, l], has-right-deriv-at(\allowbreak{}f, x, l) iff forsome(\allowbreak{}[ose\_], pos-rr(\allowbreak{}ose\_) and has-deriv-at-within(\allowbreak{}f, coint(\allowbreak{}x, x + ose\_), x, l)))
\end{quote}
also: \texttt{has-right-deriv-at-sum} \texttt{has-right-deriv-at-in-rr} \texttt{has-right-deriv-at-intro} \texttt{has-right-deriv-at-local} \texttt{has-right-deriv-at-shrink} \texttt{has-right-deriv-at-unique} \texttt{has-right-deriv-at-pt-in-rr} \\
mentioned by 20 result(s)

\smallskip

\noindent\textbf{\texttt{has-zero-of-order(zou\_, zof\_, zoa\_, zom\_)}}\quad\textit{predicate}\index{has-zero-of-order(zou_, zof_, zoa_, zom_)@\texttt{has-zero-of-order(zou\_, zof\_, zoa\_, zom\_)}}\\
reads: \$2 has a zero of order \$4 at \$3 (on \$1) \\
defined by \texttt{has-zero-of-order} \\
characterized by \texttt{has-zero-of-order}  \textit{(definitional)}:
\begin{quote}\ttfamily\footnotesize\raggedright\noindent
forall(\allowbreak{}[zou\_, zof\_, zoa\_, zom\_], has-zero-of-order(\allowbreak{}zou\_, zof\_, zoa\_, zom\_) iff holomorphic-on(\allowbreak{}zou\_, zof\_) and zoa\_ in zou\_ and zom\_ in nn and forsome(\allowbreak{}[zoh\_], holomorphic-on(\allowbreak{}zou\_, zoh\_) and not(\allowbreak{}zoh\_(\allowbreak{}zoa\_) = 0) and forall(\allowbreak{}[zox\_ in zou\_], zof\_(\allowbreak{}zox\_) = (\allowbreak{}zox\_ - zoa\_) \textasciicircum{} zom\_ * zoh\_(\allowbreak{}zox\_))))
\end{quote}
mentioned by 36 result(s)

\smallskip

\noindent\textbf{\texttt{hausdorff(mhm\_, mhl\_)}}\quad\textit{functoid}\index{hausdorff(mhm_, mhl_)@\texttt{hausdorff(mhm\_, mhl\_)}}\\
reads: the Hausdorff outer measure of dimension \$2 on \$1 \\
definition:
\begin{quote}\ttfamily\footnotesize\raggedright\noindent
method-ii(\allowbreak{}mhm\_, power(\allowbreak{}pts(\allowbreak{}mhm\_)), hausdorff-pm(\allowbreak{}mhm\_, mhl\_))
\end{quote}
characterized by \texttt{is-hausdorff-dim}  \textit{(definitional)}:
\begin{quote}\ttfamily\footnotesize\raggedright\noindent
forall(\allowbreak{}[mhm\_, mhs\_, mhl\_], is-hausdorff-dim(\allowbreak{}mhm\_, mhs\_, mhl\_) iff mhl\_ in rr and 0 <= mhl\_ and forall(\allowbreak{}[mhr\_ in rr], mhl\_ < mhr\_ implies (\allowbreak{}hausdorff(\allowbreak{}mhm\_, mhr\_))(\allowbreak{}mhs\_) = 0) and forall(\allowbreak{}[mhr\_ in rr], 0 < mhr\_ implies mhr\_ < mhl\_ implies (\allowbreak{}hausdorff(\allowbreak{}mhm\_, mhr\_))(\allowbreak{}mhs\_) = pos-inf))
\end{quote}
also: \texttt{hausdorff-unfold} \texttt{hausdorff-compare} \texttt{hausdorff-scale-pos} \texttt{hausdorff-positive-infinite} \texttt{hausdorff-metric-outer-measure} \\
mentioned by 9 result(s)

\smallskip

\noindent\textbf{\texttt{hausdorff-phi(mhm\_, mhf\_)}}\quad\textit{functoid}\index{hausdorff-phi(mhm_, mhf_)@\texttt{hausdorff-phi(mhm\_, mhf\_)}}\\
reads: the generalized Hausdorff outer measure H\textasciicircum{}\$2 on \$1 \\
definition:
\begin{quote}\ttfamily\footnotesize\raggedright\noindent
method-ii(\allowbreak{}mhm\_, power(\allowbreak{}pts(\allowbreak{}mhm\_)), phi-pm(\allowbreak{}mhm\_, mhf\_))
\end{quote}
characterized by \texttt{hausdorff-phi-compare}  \textit{(proven)}:
\begin{quote}\ttfamily\footnotesize\raggedright\noindent
forall(\allowbreak{}[m, m1, c, phi], is-metric-space(\allowbreak{}m) implies is-metric-space(\allowbreak{}m1) implies pts(\allowbreak{}m1) = pts(\allowbreak{}m) implies c in rr implies 0 < c implies forall(\allowbreak{}[x in pts(\allowbreak{}m), y in pts(\allowbreak{}m)], (\allowbreak{}dist(\allowbreak{}m))(\allowbreak{}x, y) <= c * (\allowbreak{}dist(\allowbreak{}m1))(\allowbreak{}x, y)) implies is-doubling(\allowbreak{}phi) implies forsome(\allowbreak{}[mhk\_], pos-rr(\allowbreak{}mhk\_) and forall(\allowbreak{}[es in power(\allowbreak{}pts(\allowbreak{}m))], (\allowbreak{}hausdorff-phi(\allowbreak{}m, phi))(\allowbreak{}es) <= etimes(\allowbreak{}mhk\_, (\allowbreak{}hausdorff-phi(\allowbreak{}m1, phi))(\allowbreak{}es)))))
\end{quote}
also: \texttt{hausdorff-phi-subspace} \\
mentioned by 3 result(s)

\smallskip

\noindent\textbf{\texttt{hausdorff-pm(mhm\_, mhl\_)}}\quad\textit{functoid}\index{hausdorff-pm(mhm_, mhl_)@\texttt{hausdorff-pm(mhm\_, mhl\_)}}\\
reads: A |-> diam(A)\textasciicircum{}\$2 (in \$1) \\
definition:
\begin{quote}\ttfamily\footnotesize\raggedright\noindent
vnb-lambda(\allowbreak{}mhh\_, power(\allowbreak{}pts(\allowbreak{}mhm\_)), epow(\allowbreak{}diam(\allowbreak{}mhm\_, mhh\_), mhl\_))
\end{quote}
characterized by \texttt{hausdorff-pm-value}  \textit{(proven)}:
\begin{quote}\ttfamily\footnotesize\raggedright\noindent
forall(\allowbreak{}[m, lam, a in power(\allowbreak{}pts(\allowbreak{}m))], (\allowbreak{}hausdorff-pm(\allowbreak{}m, lam))(\allowbreak{}a) == epow(\allowbreak{}diam(\allowbreak{}m, a), lam))
\end{quote}
also: \texttt{hausdorff-pm-premeasure} \\
mentioned by 7 result(s)

\smallskip

\noindent\textbf{\texttt{have!}}\quad\textit{tactic}\index{have"!@\texttt{have"!}}\\
\texttt{(have! CLAIM [THUNK])} -- Assert CLAIM as an intermediate step and carry on with it as a hypothesis: cut CLAIM, discharge the side goal from context (or by THUNK), and stay on the main branch. \\
declared in \texttt{tactics-help.scm}

\smallskip

\noindent\textbf{\texttt{hc-embed(md, fq)}}\quad\textit{functoid}\index{hc-embed(md, fq)@\texttt{hc-embed(md, fq)}}\\
reads: the embedding of (\$1, \$2) in its Hilbert completion \\
definition:
\begin{quote}\ttfamily\footnotesize\raggedright\noindent
vnb-lambda(\allowbreak{}hcx\_, vec(\allowbreak{}md), class(\allowbreak{}qmod-setoid(\allowbreak{}cseq-module(\allowbreak{}md, fq), form-null(\allowbreak{}cseq-module(\allowbreak{}md, fq), limit-form(\allowbreak{}md, fq))), embed-seq(\allowbreak{}form-ms(\allowbreak{}md, fq), hcx\_)))
\end{quote}
characterized by \texttt{hc-embed-ip}  \textit{(proven)}:
\begin{quote}\ttfamily\footnotesize\raggedright\noindent
forall(\allowbreak{}[md, fq, x\_, y\_], is-module(\allowbreak{}md) implies scal(\allowbreak{}md) = normed-field-as-commutative-ring(\allowbreak{}cc-normed-field) implies is-inner-product-form(\allowbreak{}md, fq) implies x\_ in vec(\allowbreak{}md) implies y\_ in vec(\allowbreak{}md) implies (\allowbreak{}ip(\allowbreak{}hilbert-completion(\allowbreak{}md, fq)))(\allowbreak{}(\allowbreak{}hc-embed(\allowbreak{}md, fq))(\allowbreak{}x\_), (\allowbreak{}hc-embed(\allowbreak{}md, fq))(\allowbreak{}y\_)) = fq(\allowbreak{}x\_, y\_))
\end{quote}
also: \texttt{hc-embed-act} \texttt{hc-embed-near} \texttt{hc-embed-type} \texttt{hc-embed-vadd} \texttt{hc-embed-dense} \texttt{hc-embed-inverse} \texttt{hc-embed-isometry} \texttt{hc-embed-injective} \\
mentioned by 20 result(s)

\smallskip

\noindent\textbf{\texttt{hd-argk(mch\_)}}\quad\textit{functoid}\index{hd-argk(mch_)@\texttt{hd-argk(mch\_)}}\\
reads: the kind of the arguments of the head \$1 \\
definition:
\begin{quote}\ttfamily\footnotesize\raggedright\noindent
if(\allowbreak{}mch\_ <= 6, 2, 1)
\end{quote}
mentioned by 6 result(s)

\smallskip

\noindent\textbf{\texttt{hd-body(mch\_)}}\quad\textit{functoid}\index{hd-body(mch_)@\texttt{hd-body(mch\_)}}\\
reads: the kind of the body of a binder of head \$1 \\
definition:
\begin{quote}\ttfamily\footnotesize\raggedright\noindent
if(\allowbreak{}mch\_ = 24 or mch\_ = 26, 1, 2)
\end{quote}
mentioned by 4 result(s)

\smallskip

\noindent\textbf{\texttt{hd-bres(mch\_)}}\quad\textit{functoid}\index{hd-bres(mch_)@\texttt{hd-bres(mch\_)}}\\
reads: the kind of a binder of head \$1 \\
definition:
\begin{quote}\ttfamily\footnotesize\raggedright\noindent
if(\allowbreak{}mch\_ <= 21, 2, 1)
\end{quote}
mentioned by 4 result(s)

\smallskip

\noindent\textbf{\texttt{hd-res(mch\_)}}\quad\textit{functoid}\index{hd-res(mch_)@\texttt{hd-res(mch\_)}}\\
reads: the kind of an application of the head \$1 \\
definition:
\begin{quote}\ttfamily\footnotesize\raggedright\noindent
if(\allowbreak{}mch\_ <= 10, 2, 1)
\end{quote}
mentioned by 6 result(s)

\smallskip

\noindent\textbf{\texttt{hilbert-completion(md, fq)}}\quad\textit{functoid}\index{hilbert-completion(md, fq)@\texttt{hilbert-completion(md, fq)}}\\
reads: the Hilbert completion of the inner product space (\$1, \$2) \\
definition:
\begin{quote}\ttfamily\footnotesize\raggedright\noindent
form-quotient(\allowbreak{}cseq-module(\allowbreak{}md, fq), limit-form(\allowbreak{}md, fq))
\end{quote}
characterized by \texttt{hilbert-completion-is-cips}  \textit{(proven)}:
\begin{quote}\ttfamily\footnotesize\raggedright\noindent
forall(\allowbreak{}[md, fq], is-module(\allowbreak{}md) implies scal(\allowbreak{}md) = normed-field-as-commutative-ring(\allowbreak{}cc-normed-field) implies is-inner-product-form(\allowbreak{}md, fq) implies is-complex-inner-product-space(\allowbreak{}hilbert-completion(\allowbreak{}md, fq)))
\end{quote}
also: \texttt{hilbert-completion-complete} \texttt{hilbert-completion-universal} \texttt{hilbert-completion-is-hilbert} \\
mentioned by 38 result(s)

\smallskip

\noindent\textbf{\texttt{holomorphic-on(u, f)}}\quad\textit{predicate}\index{holomorphic-on(u, f)@\texttt{holomorphic-on(u, f)}}\\
reads: \$2 is holomorphic on \$1 \\
defined by \texttt{holomorphic-on} \\
characterized by \texttt{holomorphic-on}  \textit{(definitional)}:
\begin{quote}\ttfamily\footnotesize\raggedright\noindent
forall(\allowbreak{}[u, f], holomorphic-on(\allowbreak{}u, f) iff is-open(\allowbreak{}nf-metric-space(\allowbreak{}cc-normed-field), u) and f in fun(\allowbreak{}u, cc) and forall(\allowbreak{}[a in u], forsome(\allowbreak{}[l], is-diff-on(\allowbreak{}cc-normed-field, u, f, a, l))))
\end{quote}
also: \texttt{holomorphic-on-diff} \texttt{holomorphic-on-open} \texttt{holomorphic-on-in-fun} \texttt{holomorphic-on-restrict} \texttt{holomorphic-on-iff-cc-ms} \texttt{holomorphic-on-continuous} \texttt{holomorphic-on-open-cc-ms} \\
mentioned by 344 result(s)

\smallskip

\noindent\textbf{\texttt{hom-abelian-group(a, b)}}\quad\textit{functoid}\index{hom-abelian-group(a, b)@\texttt{hom-abelian-group(a, b)}}\\
reads: the abelian group morphisms from \$1 to \$2 \\
definition:
\begin{quote}\ttfamily\footnotesize\raggedright\noindent
\{homf\_ in fun(\allowbreak{}carr(\allowbreak{}a), carr(\allowbreak{}b)): is-hom-abelian-group(\allowbreak{}a, b, homf\_)\}
\end{quote}
characterized by \texttt{hom-abelian-group-in-set}  \textit{(proven)}:
\begin{quote}\ttfamily\footnotesize\raggedright\noindent
forall(\allowbreak{}[a, b], is-abelian-group(\allowbreak{}a) implies is-abelian-group(\allowbreak{}b) implies hom-abelian-group(\allowbreak{}a, b) in set)
\end{quote}
also: \texttt{hom-abelian-group-pre-type} \texttt{hom-abelian-group-post-type} \texttt{hom-abelian-group-member-iff} \\
mentioned by 5 result(s)

\smallskip

\noindent\textbf{\texttt{hom-bounded-linear(a, b)}}\quad\textit{functoid}\index{hom-bounded-linear(a, b)@\texttt{hom-bounded-linear(a, b)}}\\
reads: the bounded linear maps from \$1 to \$2 \\
definition:
\begin{quote}\ttfamily\footnotesize\raggedright\noindent
\{homf\_ in fun(\allowbreak{}vec(\allowbreak{}a), vec(\allowbreak{}b)): is-bounded-linear-arrow(\allowbreak{}a, b, homf\_)\}
\end{quote}
characterized by \texttt{hom-bounded-linear-in-set}  \textit{(proven)}:
\begin{quote}\ttfamily\footnotesize\raggedright\noindent
forall(\allowbreak{}[a, b], is-normed-vector-space(\allowbreak{}a) implies is-normed-vector-space(\allowbreak{}b) implies hom-bounded-linear(\allowbreak{}a, b) in set)
\end{quote}
also: \texttt{hom-bounded-linear-pre-type} \texttt{hom-bounded-linear-post-type} \texttt{hom-bounded-linear-member-iff} \texttt{hom-normed-vector-space-subset-hom-bounded-linear} \\
mentioned by 6 result(s)

\smallskip

\noindent\textbf{\texttt{hom-bundle(a, b)}}\quad\textit{functoid}\index{hom-bundle(a, b)@\texttt{hom-bundle(a, b)}}\\
reads: the bundle morphisms from \$1 to \$2 \\
definition:
\begin{quote}\ttfamily\footnotesize\raggedright\noindent
\{homf\_ in fun(\allowbreak{}bnodes(\allowbreak{}a), bnodes(\allowbreak{}b)): is-hom-bundle(\allowbreak{}a, b, homf\_)\}
\end{quote}

\smallskip

\noindent\textbf{\texttt{hom-c-continuous(a, b)}}\quad\textit{functoid}\index{hom-c-continuous(a, b)@\texttt{hom-c-continuous(a, b)}}\\
reads: the continuous maps of countably-metrised spaces from \$1 to \$2 \\
definition:
\begin{quote}\ttfamily\footnotesize\raggedright\noindent
\{homf\_ in fun(\allowbreak{}pts(\allowbreak{}a), pts(\allowbreak{}b)): is-c-continuous-arrow(\allowbreak{}a, b, homf\_)\}
\end{quote}
characterized by \texttt{hom-c-continuous-in-set}  \textit{(proven)}:
\begin{quote}\ttfamily\footnotesize\raggedright\noindent
forall(\allowbreak{}[a, b], is-c-metric-space(\allowbreak{}a) implies is-c-metric-space(\allowbreak{}b) implies hom-c-continuous(\allowbreak{}a, b) in set)
\end{quote}
also: \texttt{hom-c-continuous-pre-type} \texttt{hom-c-continuous-post-type} \texttt{hom-c-continuous-member-iff} \\
mentioned by 5 result(s)

\smallskip

\noindent\textbf{\texttt{hom-c-metric-space(a, b)}}\quad\textit{functoid}\index{hom-c-metric-space(a, b)@\texttt{hom-c-metric-space(a, b)}}\\
reads: the c metric space morphisms from \$1 to \$2 \\
definition:
\begin{quote}\ttfamily\footnotesize\raggedright\noindent
\{homf\_ in fun(\allowbreak{}pts(\allowbreak{}a), pts(\allowbreak{}b)): is-hom-c-metric-space(\allowbreak{}a, b, homf\_)\}
\end{quote}
characterized by \texttt{hom-c-metric-space-in-set}  \textit{(proven)}:
\begin{quote}\ttfamily\footnotesize\raggedright\noindent
forall(\allowbreak{}[a, b], is-c-metric-space(\allowbreak{}a) implies is-c-metric-space(\allowbreak{}b) implies hom-c-metric-space(\allowbreak{}a, b) in set)
\end{quote}
also: \texttt{hom-c-metric-space-pre-type} \texttt{hom-c-metric-space-post-type} \texttt{hom-c-metric-space-member-iff} \\
mentioned by 5 result(s)

\smallskip

\noindent\textbf{\texttt{hom-comm-monoid(a, b)}}\quad\textit{functoid}\index{hom-comm-monoid(a, b)@\texttt{hom-comm-monoid(a, b)}}\\
reads: the comm monoid morphisms from \$1 to \$2 \\
definition:
\begin{quote}\ttfamily\footnotesize\raggedright\noindent
\{homf\_ in fun(\allowbreak{}carr(\allowbreak{}a), carr(\allowbreak{}b)): is-hom-comm-monoid(\allowbreak{}a, b, homf\_)\}
\end{quote}
characterized by \texttt{hom-comm-monoid-in-set}  \textit{(proven)}:
\begin{quote}\ttfamily\footnotesize\raggedright\noindent
forall(\allowbreak{}[a, b], is-comm-monoid(\allowbreak{}a) implies is-comm-monoid(\allowbreak{}b) implies hom-comm-monoid(\allowbreak{}a, b) in set)
\end{quote}
also: \texttt{hom-comm-monoid-pre-type} \texttt{hom-comm-monoid-post-type} \texttt{hom-comm-monoid-member-iff} \\
mentioned by 5 result(s)

\smallskip

\noindent\textbf{\texttt{hom-commutative-banach-algebra(a, b)}}\quad\textit{functoid}\index{hom-commutative-banach-algebra(a, b)@\texttt{hom-commutative-banach-algebra(a, b)}}\\
reads: the commutative banach algebra morphisms from \$1 to \$2 \\
definition:
\begin{quote}\ttfamily\footnotesize\raggedright\noindent
\{homf\_ in fun(\allowbreak{}vec(\allowbreak{}a), vec(\allowbreak{}b)): is-hom-commutative-banach-algebra(\allowbreak{}a, b, homf\_)\}
\end{quote}

\smallskip

\noindent\textbf{\texttt{hom-commutative-c-star-algebra(a, b)}}\quad\textit{functoid}\index{hom-commutative-c-star-algebra(a, b)@\texttt{hom-commutative-c-star-algebra(a, b)}}\\
reads: the commutative c star algebra morphisms from \$1 to \$2 \\
definition:
\begin{quote}\ttfamily\footnotesize\raggedright\noindent
\{homf\_ in fun(\allowbreak{}vec(\allowbreak{}a), vec(\allowbreak{}b)): is-hom-commutative-c-star-algebra(\allowbreak{}a, b, homf\_)\}
\end{quote}

\smallskip

\noindent\textbf{\texttt{hom-commutative-ring(a, b)}}\quad\textit{functoid}\index{hom-commutative-ring(a, b)@\texttt{hom-commutative-ring(a, b)}}\\
reads: the commutative ring morphisms from \$1 to \$2 \\
definition:
\begin{quote}\ttfamily\footnotesize\raggedright\noindent
\{homf\_ in fun(\allowbreak{}carr(\allowbreak{}a), carr(\allowbreak{}b)): is-hom-commutative-ring(\allowbreak{}a, b, homf\_)\}
\end{quote}
characterized by \texttt{hom-commutative-ring-in-set}  \textit{(proven)}:
\begin{quote}\ttfamily\footnotesize\raggedright\noindent
forall(\allowbreak{}[a, b], is-commutative-ring(\allowbreak{}a) implies is-commutative-ring(\allowbreak{}b) implies hom-commutative-ring(\allowbreak{}a, b) in set)
\end{quote}
also: \texttt{hom-commutative-ring-pre-type} \texttt{hom-commutative-ring-post-type} \texttt{hom-commutative-ring-member-iff} \\
mentioned by 5 result(s)

\smallskip

\noindent\textbf{\texttt{hom-complex-banach-algebra(a, b)}}\quad\textit{functoid}\index{hom-complex-banach-algebra(a, b)@\texttt{hom-complex-banach-algebra(a, b)}}\\
reads: the complex banach algebra morphisms from \$1 to \$2 \\
definition:
\begin{quote}\ttfamily\footnotesize\raggedright\noindent
\{homf\_ in fun(\allowbreak{}vec(\allowbreak{}a), vec(\allowbreak{}b)): is-hom-complex-banach-algebra(\allowbreak{}a, b, homf\_)\}
\end{quote}

\smallskip

\noindent\textbf{\texttt{hom-complex-banach-space(a, b)}}\quad\textit{functoid}\index{hom-complex-banach-space(a, b)@\texttt{hom-complex-banach-space(a, b)}}\\
reads: the complex banach space morphisms from \$1 to \$2 \\
definition:
\begin{quote}\ttfamily\footnotesize\raggedright\noindent
\{homf\_ in fun(\allowbreak{}vec(\allowbreak{}a), vec(\allowbreak{}b)): is-hom-complex-banach-space(\allowbreak{}a, b, homf\_)\}
\end{quote}

\smallskip

\noindent\textbf{\texttt{hom-complex-inner-product-space(a, b)}}\quad\textit{functoid}\index{hom-complex-inner-product-space(a, b)@\texttt{hom-complex-inner-product-space(a, b)}}\\
reads: the complex inner product space morphisms from \$1 to \$2 \\
definition:
\begin{quote}\ttfamily\footnotesize\raggedright\noindent
\{homf\_ in fun(\allowbreak{}vec(\allowbreak{}a), vec(\allowbreak{}b)): is-hom-complex-inner-product-space(\allowbreak{}a, b, homf\_)\}
\end{quote}
characterized by \texttt{hom-complex-inner-product-space-in-set}  \textit{(proven)}:
\begin{quote}\ttfamily\footnotesize\raggedright\noindent
forall(\allowbreak{}[a, b], is-complex-inner-product-space(\allowbreak{}a) implies is-complex-inner-product-space(\allowbreak{}b) implies hom-complex-inner-product-space(\allowbreak{}a, b) in set)
\end{quote}
also: \texttt{hom-complex-inner-product-space-pre-type} \texttt{hom-complex-inner-product-space-post-type} \texttt{hom-complex-inner-product-space-member-iff} \\
mentioned by 5 result(s)

\smallskip

\noindent\textbf{\texttt{hom-complex-normed-algebra(a, b)}}\quad\textit{functoid}\index{hom-complex-normed-algebra(a, b)@\texttt{hom-complex-normed-algebra(a, b)}}\\
reads: the complex normed algebra morphisms from \$1 to \$2 \\
definition:
\begin{quote}\ttfamily\footnotesize\raggedright\noindent
\{homf\_ in fun(\allowbreak{}vec(\allowbreak{}a), vec(\allowbreak{}b)): is-hom-complex-normed-algebra(\allowbreak{}a, b, homf\_)\}
\end{quote}

\smallskip

\noindent\textbf{\texttt{hom-complex-normed-vector-space(a, b)}}\quad\textit{functoid}\index{hom-complex-normed-vector-space(a, b)@\texttt{hom-complex-normed-vector-space(a, b)}}\\
reads: the complex normed vector space morphisms from \$1 to \$2 \\
definition:
\begin{quote}\ttfamily\footnotesize\raggedright\noindent
\{homf\_ in fun(\allowbreak{}vec(\allowbreak{}a), vec(\allowbreak{}b)): is-hom-complex-normed-vector-space(\allowbreak{}a, b, homf\_)\}
\end{quote}

\smallskip

\noindent\textbf{\texttt{hom-complex-pre-hilbert-space(a, b)}}\quad\textit{functoid}\index{hom-complex-pre-hilbert-space(a, b)@\texttt{hom-complex-pre-hilbert-space(a, b)}}\\
reads: the complex pre hilbert space morphisms from \$1 to \$2 \\
definition:
\begin{quote}\ttfamily\footnotesize\raggedright\noindent
\{homf\_ in fun(\allowbreak{}vec(\allowbreak{}a), vec(\allowbreak{}b)): is-hom-complex-pre-hilbert-space(\allowbreak{}a, b, homf\_)\}
\end{quote}

\smallskip

\noindent\textbf{\texttt{hom-complex-semi-inner-product-space(a, b)}}\quad\textit{functoid}\index{hom-complex-semi-inner-product-space(a, b)@\texttt{hom-complex-semi-inner-product-space(a, b)}}\\
reads: the complex semi inner product space morphisms from \$1 to \$2 \\
definition:
\begin{quote}\ttfamily\footnotesize\raggedright\noindent
\{homf\_ in fun(\allowbreak{}vec(\allowbreak{}a), vec(\allowbreak{}b)): is-hom-complex-semi-inner-product-space(\allowbreak{}a, b, homf\_)\}
\end{quote}

\smallskip

\noindent\textbf{\texttt{hom-continuous(a, b)}}\quad\textit{functoid}\index{hom-continuous(a, b)@\texttt{hom-continuous(a, b)}}\\
reads: the continuous maps from \$1 to \$2 \\
definition:
\begin{quote}\ttfamily\footnotesize\raggedright\noindent
\{homf\_ in fun(\allowbreak{}pts(\allowbreak{}a), pts(\allowbreak{}b)): is-continuous-arrow(\allowbreak{}a, b, homf\_)\}
\end{quote}
characterized by \texttt{hom-continuous-in-set}  \textit{(proven)}:
\begin{quote}\ttfamily\footnotesize\raggedright\noindent
forall(\allowbreak{}[a, b], is-metric-space(\allowbreak{}a) implies is-metric-space(\allowbreak{}b) implies hom-continuous(\allowbreak{}a, b) in set)
\end{quote}
also: \texttt{hom-continuous-pre-type} \texttt{hom-continuous-post-type} \texttt{hom-continuous-member-iff} \texttt{hom-uniformly-continuous-subset-hom-continuous} \\
mentioned by 6 result(s)

\smallskip

\noindent\textbf{\texttt{hom-directed-set(a, b)}}\quad\textit{functoid}\index{hom-directed-set(a, b)@\texttt{hom-directed-set(a, b)}}\\
reads: the directed set morphisms from \$1 to \$2 \\
definition:
\begin{quote}\ttfamily\footnotesize\raggedright\noindent
\{homf\_ in fun(\allowbreak{}pts(\allowbreak{}a), pts(\allowbreak{}b)): is-hom-directed-set(\allowbreak{}a, b, homf\_)\}
\end{quote}

\smallskip

\noindent\textbf{\texttt{hom-euclidean-ring(a, b)}}\quad\textit{functoid}\index{hom-euclidean-ring(a, b)@\texttt{hom-euclidean-ring(a, b)}}\\
reads: the euclidean ring morphisms from \$1 to \$2 \\
definition:
\begin{quote}\ttfamily\footnotesize\raggedright\noindent
\{homf\_ in fun(\allowbreak{}carr(\allowbreak{}a), carr(\allowbreak{}b)): is-hom-euclidean-ring(\allowbreak{}a, b, homf\_)\}
\end{quote}
characterized by \texttt{hom-euclidean-ring-in-set}  \textit{(proven)}:
\begin{quote}\ttfamily\footnotesize\raggedright\noindent
forall(\allowbreak{}[a, b], is-euclidean-ring(\allowbreak{}a) implies is-euclidean-ring(\allowbreak{}b) implies hom-euclidean-ring(\allowbreak{}a, b) in set)
\end{quote}
also: \texttt{hom-euclidean-ring-pre-type} \texttt{hom-euclidean-ring-post-type} \texttt{hom-euclidean-ring-member-iff} \\
mentioned by 5 result(s)

\smallskip

\noindent\textbf{\texttt{hom-field(a, b)}}\quad\textit{functoid}\index{hom-field(a, b)@\texttt{hom-field(a, b)}}\\
reads: the field morphisms from \$1 to \$2 \\
definition:
\begin{quote}\ttfamily\footnotesize\raggedright\noindent
\{homf\_ in fun(\allowbreak{}carr(\allowbreak{}a), carr(\allowbreak{}b)): is-hom-field(\allowbreak{}a, b, homf\_)\}
\end{quote}
characterized by \texttt{hom-field-in-set}  \textit{(proven)}:
\begin{quote}\ttfamily\footnotesize\raggedright\noindent
forall(\allowbreak{}[a, b], is-field(\allowbreak{}a) implies is-field(\allowbreak{}b) implies hom-field(\allowbreak{}a, b) in set)
\end{quote}
also: \texttt{hom-field-pre-type} \texttt{hom-field-post-type} \texttt{hom-field-member-iff} \\
mentioned by 5 result(s)

\smallskip

\noindent\textbf{\texttt{hom-field-ring(a, b)}}\quad\textit{functoid}\index{hom-field-ring(a, b)@\texttt{hom-field-ring(a, b)}}\\
reads: the field ring morphisms from \$1 to \$2 \\
definition:
\begin{quote}\ttfamily\footnotesize\raggedright\noindent
\{homf\_ in fun(\allowbreak{}carr(\allowbreak{}a), carr(\allowbreak{}b)): is-hom-field-ring(\allowbreak{}a, b, homf\_)\}
\end{quote}
characterized by \texttt{hom-field-ring-in-set}  \textit{(proven)}:
\begin{quote}\ttfamily\footnotesize\raggedright\noindent
forall(\allowbreak{}[a, b], is-field-ring(\allowbreak{}a) implies is-field-ring(\allowbreak{}b) implies hom-field-ring(\allowbreak{}a, b) in set)
\end{quote}
also: \texttt{hom-field-ring-pre-type} \texttt{hom-field-ring-post-type} \texttt{hom-field-ring-member-iff} \\
mentioned by 5 result(s)

\smallskip

\noindent\textbf{\texttt{hom-group(a, b)}}\quad\textit{functoid}\index{hom-group(a, b)@\texttt{hom-group(a, b)}}\\
reads: the group morphisms from \$1 to \$2 \\
definition:
\begin{quote}\ttfamily\footnotesize\raggedright\noindent
\{homf\_ in fun(\allowbreak{}carr(\allowbreak{}a), carr(\allowbreak{}b)): is-hom-group(\allowbreak{}a, b, homf\_)\}
\end{quote}
characterized by \texttt{hom-group-in-set}  \textit{(proven)}:
\begin{quote}\ttfamily\footnotesize\raggedright\noindent
forall(\allowbreak{}[a, b], is-group(\allowbreak{}a) implies is-group(\allowbreak{}b) implies hom-group(\allowbreak{}a, b) in set)
\end{quote}
also: \texttt{hom-group-pre-type} \texttt{hom-group-post-type} \texttt{hom-group-member-iff} \\
mentioned by 5 result(s)

\smallskip

\noindent\textbf{\texttt{hom-infiltrated-space(a, b)}}\quad\textit{functoid}\index{hom-infiltrated-space(a, b)@\texttt{hom-infiltrated-space(a, b)}}\\
reads: the infiltrated space morphisms from \$1 to \$2 \\
definition:
\begin{quote}\ttfamily\footnotesize\raggedright\noindent
\{homp\_ in cartesian(\allowbreak{}fun(\allowbreak{}adv-ids(\allowbreak{}a), adv-ids(\allowbreak{}b)), fun(\allowbreak{}adv-tokens(\allowbreak{}a), adv-tokens(\allowbreak{}b))): is-hom-infiltrated-space(\allowbreak{}a, b, nth(\allowbreak{}1, homp\_), nth(\allowbreak{}2, homp\_))\}
\end{quote}

\smallskip

\noindent\textbf{\texttt{hom-integral-domain(a, b)}}\quad\textit{functoid}\index{hom-integral-domain(a, b)@\texttt{hom-integral-domain(a, b)}}\\
reads: the integral domain morphisms from \$1 to \$2 \\
definition:
\begin{quote}\ttfamily\footnotesize\raggedright\noindent
\{homf\_ in fun(\allowbreak{}carr(\allowbreak{}a), carr(\allowbreak{}b)): is-hom-integral-domain(\allowbreak{}a, b, homf\_)\}
\end{quote}
characterized by \texttt{hom-integral-domain-in-set}  \textit{(proven)}:
\begin{quote}\ttfamily\footnotesize\raggedright\noindent
forall(\allowbreak{}[a, b], is-integral-domain(\allowbreak{}a) implies is-integral-domain(\allowbreak{}b) implies hom-integral-domain(\allowbreak{}a, b) in set)
\end{quote}
also: \texttt{hom-integral-domain-pre-type} \texttt{hom-integral-domain-post-type} \texttt{hom-integral-domain-member-iff} \\
mentioned by 5 result(s)

\smallskip

\noindent\textbf{\texttt{hom-lipschitz(a, b)}}\quad\textit{functoid}\index{hom-lipschitz(a, b)@\texttt{hom-lipschitz(a, b)}}\\
reads: the Lipschitz maps from \$1 to \$2 \\
definition:
\begin{quote}\ttfamily\footnotesize\raggedright\noindent
\{homf\_ in fun(\allowbreak{}pts(\allowbreak{}a), pts(\allowbreak{}b)): is-lipschitz-arrow(\allowbreak{}a, b, homf\_)\}
\end{quote}
characterized by \texttt{hom-lipschitz-in-set}  \textit{(proven)}:
\begin{quote}\ttfamily\footnotesize\raggedright\noindent
forall(\allowbreak{}[a, b], is-metric-space(\allowbreak{}a) implies is-metric-space(\allowbreak{}b) implies hom-lipschitz(\allowbreak{}a, b) in set)
\end{quote}
also: \texttt{hom-lipschitz-pre-type} \texttt{hom-lipschitz-post-type} \texttt{hom-lipschitz-member-iff} \texttt{hom-metric-space-subset-hom-lipschitz} \texttt{hom-lipschitz-subset-hom-uniformly-continuous} \\
mentioned by 7 result(s)

\smallskip

\noindent\textbf{\texttt{hom-measurable-space(a, b)}}\quad\textit{functoid}\index{hom-measurable-space(a, b)@\texttt{hom-measurable-space(a, b)}}\\
reads: the measurable space morphisms from \$1 to \$2 \\
definition:
\begin{quote}\ttfamily\footnotesize\raggedright\noindent
\{homf\_ in fun(\allowbreak{}pts(\allowbreak{}a), pts(\allowbreak{}b)): is-hom-measurable-space(\allowbreak{}a, b, homf\_)\}
\end{quote}
characterized by \texttt{hom-measurable-space-in-set}  \textit{(proven)}:
\begin{quote}\ttfamily\footnotesize\raggedright\noindent
forall(\allowbreak{}[a, b], is-measurable-space(\allowbreak{}a) implies is-measurable-space(\allowbreak{}b) implies hom-measurable-space(\allowbreak{}a, b) in set)
\end{quote}
also: \texttt{hom-measurable-space-pre-type} \texttt{hom-measurable-space-post-type} \texttt{hom-measurable-space-member-iff} \\
mentioned by 5 result(s)

\smallskip

\noindent\textbf{\texttt{hom-measure-space(a, b)}}\quad\textit{functoid}\index{hom-measure-space(a, b)@\texttt{hom-measure-space(a, b)}}\\
reads: the measure space morphisms from \$1 to \$2 \\
definition:
\begin{quote}\ttfamily\footnotesize\raggedright\noindent
\{homf\_ in fun(\allowbreak{}pts(\allowbreak{}a), pts(\allowbreak{}b)): is-hom-measure-space(\allowbreak{}a, b, homf\_)\}
\end{quote}
characterized by \texttt{hom-measure-space-in-set}  \textit{(proven)}:
\begin{quote}\ttfamily\footnotesize\raggedright\noindent
forall(\allowbreak{}[a, b], is-measure-space(\allowbreak{}a) implies is-measure-space(\allowbreak{}b) implies hom-measure-space(\allowbreak{}a, b) in set)
\end{quote}
also: \texttt{hom-measure-space-pre-type} \texttt{hom-measure-space-post-type} \texttt{hom-measure-space-member-iff} \\
mentioned by 5 result(s)

\smallskip

\noindent\textbf{\texttt{hom-metric-group(a, b)}}\quad\textit{functoid}\index{hom-metric-group(a, b)@\texttt{hom-metric-group(a, b)}}\\
reads: the metric group morphisms from \$1 to \$2 \\
definition:
\begin{quote}\ttfamily\footnotesize\raggedright\noindent
\{homf\_ in fun(\allowbreak{}carr(\allowbreak{}a), carr(\allowbreak{}b)): is-hom-metric-group(\allowbreak{}a, b, homf\_)\}
\end{quote}

\smallskip

\noindent\textbf{\texttt{hom-metric-space(a, b)}}\quad\textit{functoid}\index{hom-metric-space(a, b)@\texttt{hom-metric-space(a, b)}}\\
reads: the metric space morphisms from \$1 to \$2 \\
definition:
\begin{quote}\ttfamily\footnotesize\raggedright\noindent
\{homf\_ in fun(\allowbreak{}pts(\allowbreak{}a), pts(\allowbreak{}b)): is-hom-metric-space(\allowbreak{}a, b, homf\_)\}
\end{quote}
characterized by \texttt{hom-metric-space-in-set}  \textit{(proven)}:
\begin{quote}\ttfamily\footnotesize\raggedright\noindent
forall(\allowbreak{}[a, b], is-metric-space(\allowbreak{}a) implies is-metric-space(\allowbreak{}b) implies hom-metric-space(\allowbreak{}a, b) in set)
\end{quote}
also: \texttt{hom-metric-space-pre-type} \texttt{hom-metric-space-post-type} \texttt{hom-metric-space-member-iff} \texttt{hom-metric-space-subset-hom-lipschitz} \\
mentioned by 6 result(s)

\smallskip

\noindent\textbf{\texttt{hom-metrizable-top-space(a, b)}}\quad\textit{functoid}\index{hom-metrizable-top-space(a, b)@\texttt{hom-metrizable-top-space(a, b)}}\\
reads: the metrizable top space morphisms from \$1 to \$2 \\
definition:
\begin{quote}\ttfamily\footnotesize\raggedright\noindent
\{homf\_ in fun(\allowbreak{}pts(\allowbreak{}a), pts(\allowbreak{}b)): is-hom-metrizable-top-space(\allowbreak{}a, b, homf\_)\}
\end{quote}
characterized by \texttt{hom-metrizable-top-space-in-set}  \textit{(proven)}:
\begin{quote}\ttfamily\footnotesize\raggedright\noindent
forall(\allowbreak{}[a, b], is-metrizable-top-space(\allowbreak{}a) implies is-metrizable-top-space(\allowbreak{}b) implies hom-metrizable-top-space(\allowbreak{}a, b) in set)
\end{quote}
also: \texttt{hom-metrizable-top-space-pre-type} \texttt{hom-metrizable-top-space-post-type} \texttt{hom-metrizable-top-space-member-iff} \\
mentioned by 5 result(s)

\smallskip

\noindent\textbf{\texttt{hom-module(a, b)}}\quad\textit{functoid}\index{hom-module(a, b)@\texttt{hom-module(a, b)}}\\
reads: the module morphisms from \$1 to \$2 \\
definition:
\begin{quote}\ttfamily\footnotesize\raggedright\noindent
\{homf\_ in fun(\allowbreak{}vec(\allowbreak{}a), vec(\allowbreak{}b)): is-hom-module(\allowbreak{}a, b, homf\_)\}
\end{quote}
characterized by \texttt{hom-module-in-set}  \textit{(proven)}:
\begin{quote}\ttfamily\footnotesize\raggedright\noindent
forall(\allowbreak{}[a, b], is-module(\allowbreak{}a) implies is-module(\allowbreak{}b) implies hom-module(\allowbreak{}a, b) in set)
\end{quote}
also: \texttt{hom-module-pre-type} \texttt{hom-module-post-type} \texttt{hom-module-member-iff} \\
mentioned by 5 result(s)

\smallskip

\noindent\textbf{\texttt{hom-monoid(a, b)}}\quad\textit{functoid}\index{hom-monoid(a, b)@\texttt{hom-monoid(a, b)}}\\
reads: the monoid morphisms from \$1 to \$2 \\
definition:
\begin{quote}\ttfamily\footnotesize\raggedright\noindent
\{homf\_ in fun(\allowbreak{}carr(\allowbreak{}a), carr(\allowbreak{}b)): is-hom-monoid(\allowbreak{}a, b, homf\_)\}
\end{quote}
characterized by \texttt{hom-monoid-in-set}  \textit{(proven)}:
\begin{quote}\ttfamily\footnotesize\raggedright\noindent
forall(\allowbreak{}[a, b], is-monoid(\allowbreak{}a) implies is-monoid(\allowbreak{}b) implies hom-monoid(\allowbreak{}a, b) in set)
\end{quote}
also: \texttt{hom-monoid-pre-type} \texttt{hom-monoid-post-type} \texttt{hom-monoid-member-iff} \\
mentioned by 5 result(s)

\smallskip

\noindent\textbf{\texttt{hom-normed-ag(a, b)}}\quad\textit{functoid}\index{hom-normed-ag(a, b)@\texttt{hom-normed-ag(a, b)}}\\
reads: the normed ag morphisms from \$1 to \$2 \\
definition:
\begin{quote}\ttfamily\footnotesize\raggedright\noindent
\{homf\_ in fun(\allowbreak{}carr(\allowbreak{}a), carr(\allowbreak{}b)): is-hom-normed-ag(\allowbreak{}a, b, homf\_)\}
\end{quote}
characterized by \texttt{hom-normed-ag-in-set}  \textit{(proven)}:
\begin{quote}\ttfamily\footnotesize\raggedright\noindent
forall(\allowbreak{}[a, b], is-normed-ag(\allowbreak{}a) implies is-normed-ag(\allowbreak{}b) implies hom-normed-ag(\allowbreak{}a, b) in set)
\end{quote}
also: \texttt{hom-normed-ag-pre-type} \texttt{hom-normed-ag-post-type} \texttt{hom-normed-ag-member-iff} \\
mentioned by 5 result(s)

\smallskip

\noindent\textbf{\texttt{hom-normed-field(a, b)}}\quad\textit{functoid}\index{hom-normed-field(a, b)@\texttt{hom-normed-field(a, b)}}\\
reads: the normed field morphisms from \$1 to \$2 \\
definition:
\begin{quote}\ttfamily\footnotesize\raggedright\noindent
\{homf\_ in fun(\allowbreak{}carr(\allowbreak{}a), carr(\allowbreak{}b)): is-hom-normed-field(\allowbreak{}a, b, homf\_)\}
\end{quote}
characterized by \texttt{hom-normed-field-in-set}  \textit{(proven)}:
\begin{quote}\ttfamily\footnotesize\raggedright\noindent
forall(\allowbreak{}[a, b], is-normed-field(\allowbreak{}a) implies is-normed-field(\allowbreak{}b) implies hom-normed-field(\allowbreak{}a, b) in set)
\end{quote}
also: \texttt{hom-normed-field-pre-type} \texttt{hom-normed-field-post-type} \texttt{hom-normed-field-member-iff} \\
mentioned by 5 result(s)

\smallskip

\noindent\textbf{\texttt{hom-normed-vector-space(a, b)}}\quad\textit{functoid}\index{hom-normed-vector-space(a, b)@\texttt{hom-normed-vector-space(a, b)}}\\
reads: the normed vector space morphisms from \$1 to \$2 \\
definition:
\begin{quote}\ttfamily\footnotesize\raggedright\noindent
\{homf\_ in fun(\allowbreak{}vec(\allowbreak{}a), vec(\allowbreak{}b)): is-hom-normed-vector-space(\allowbreak{}a, b, homf\_)\}
\end{quote}
characterized by \texttt{hom-normed-vector-space-in-set}  \textit{(proven)}:
\begin{quote}\ttfamily\footnotesize\raggedright\noindent
forall(\allowbreak{}[a, b], is-normed-vector-space(\allowbreak{}a) implies is-normed-vector-space(\allowbreak{}b) implies hom-normed-vector-space(\allowbreak{}a, b) in set)
\end{quote}
also: \texttt{hom-normed-vector-space-pre-type} \texttt{hom-normed-vector-space-post-type} \texttt{hom-normed-vector-space-member-iff} \texttt{hom-normed-vector-space-subset-hom-bounded-linear} \\
mentioned by 6 result(s)

\smallskip

\noindent\textbf{\texttt{hom-nu-c-star-algebra(a, b)}}\quad\textit{functoid}\index{hom-nu-c-star-algebra(a, b)@\texttt{hom-nu-c-star-algebra(a, b)}}\\
reads: the nu c star algebra morphisms from \$1 to \$2 \\
definition:
\begin{quote}\ttfamily\footnotesize\raggedright\noindent
\{homf\_ in fun(\allowbreak{}vec(\allowbreak{}a), vec(\allowbreak{}b)): is-hom-nu-c-star-algebra(\allowbreak{}a, b, homf\_)\}
\end{quote}

\smallskip

\noindent\textbf{\texttt{hom-nu-involutive-banach-algebra(a, b)}}\quad\textit{functoid}\index{hom-nu-involutive-banach-algebra(a, b)@\texttt{hom-nu-involutive-banach-algebra(a, b)}}\\
reads: the nu involutive banach algebra morphisms from \$1 to \$2 \\
definition:
\begin{quote}\ttfamily\footnotesize\raggedright\noindent
\{homf\_ in fun(\allowbreak{}vec(\allowbreak{}a), vec(\allowbreak{}b)): is-hom-nu-involutive-banach-algebra(\allowbreak{}a, b, homf\_)\}
\end{quote}

\smallskip

\noindent\textbf{\texttt{hom-pid(a, b)}}\quad\textit{functoid}\index{hom-pid(a, b)@\texttt{hom-pid(a, b)}}\\
reads: the pid morphisms from \$1 to \$2 \\
definition:
\begin{quote}\ttfamily\footnotesize\raggedright\noindent
\{homf\_ in fun(\allowbreak{}carr(\allowbreak{}a), carr(\allowbreak{}b)): is-hom-pid(\allowbreak{}a, b, homf\_)\}
\end{quote}
characterized by \texttt{hom-pid-in-set}  \textit{(proven)}:
\begin{quote}\ttfamily\footnotesize\raggedright\noindent
forall(\allowbreak{}[a, b], is-pid(\allowbreak{}a) implies is-pid(\allowbreak{}b) implies hom-pid(\allowbreak{}a, b) in set)
\end{quote}
also: \texttt{hom-pid-pre-type} \texttt{hom-pid-post-type} \texttt{hom-pid-member-iff} \\
mentioned by 5 result(s)

\smallskip

\noindent\textbf{\texttt{hom-pk-space(a, b)}}\quad\textit{functoid}\index{hom-pk-space(a, b)@\texttt{hom-pk-space(a, b)}}\\
reads: the pk space morphisms from \$1 to \$2 \\
definition:
\begin{quote}\ttfamily\footnotesize\raggedright\noindent
\{homf\_ in fun(\allowbreak{}names(\allowbreak{}a), names(\allowbreak{}b)): is-hom-pk-space(\allowbreak{}a, b, homf\_)\}
\end{quote}

\smallskip

\noindent\textbf{\texttt{hom-pre-hilbert-space(a, b)}}\quad\textit{functoid}\index{hom-pre-hilbert-space(a, b)@\texttt{hom-pre-hilbert-space(a, b)}}\\
reads: the pre hilbert space morphisms from \$1 to \$2 \\
definition:
\begin{quote}\ttfamily\footnotesize\raggedright\noindent
\{homf\_ in fun(\allowbreak{}vec(\allowbreak{}a), vec(\allowbreak{}b)): is-hom-pre-hilbert-space(\allowbreak{}a, b, homf\_)\}
\end{quote}

\smallskip

\noindent\textbf{\texttt{hom-pseudometric-space(a, b)}}\quad\textit{functoid}\index{hom-pseudometric-space(a, b)@\texttt{hom-pseudometric-space(a, b)}}\\
reads: the pseudometric space morphisms from \$1 to \$2 \\
definition:
\begin{quote}\ttfamily\footnotesize\raggedright\noindent
\{homf\_ in fun(\allowbreak{}pts(\allowbreak{}a), pts(\allowbreak{}b)): is-hom-pseudometric-space(\allowbreak{}a, b, homf\_)\}
\end{quote}
characterized by \texttt{hom-pseudometric-space-in-set}  \textit{(proven)}:
\begin{quote}\ttfamily\footnotesize\raggedright\noindent
forall(\allowbreak{}[a, b], is-pseudometric-space(\allowbreak{}a) implies is-pseudometric-space(\allowbreak{}b) implies hom-pseudometric-space(\allowbreak{}a, b) in set)
\end{quote}
also: \texttt{hom-pseudometric-space-pre-type} \texttt{hom-pseudometric-space-post-type} \texttt{hom-pseudometric-space-member-iff} \\
mentioned by 5 result(s)

\smallskip

\noindent\textbf{\texttt{hom-ring(a, b)}}\quad\textit{functoid}\index{hom-ring(a, b)@\texttt{hom-ring(a, b)}}\\
reads: the ring morphisms from \$1 to \$2 \\
definition:
\begin{quote}\ttfamily\footnotesize\raggedright\noindent
\{homf\_ in fun(\allowbreak{}carr(\allowbreak{}a), carr(\allowbreak{}b)): is-hom-ring(\allowbreak{}a, b, homf\_)\}
\end{quote}
characterized by \texttt{hom-ring-in-set}  \textit{(proven)}:
\begin{quote}\ttfamily\footnotesize\raggedright\noindent
forall(\allowbreak{}[a, b], is-ring(\allowbreak{}a) implies is-ring(\allowbreak{}b) implies hom-ring(\allowbreak{}a, b) in set)
\end{quote}
also: \texttt{hom-ring-pre-type} \texttt{hom-ring-post-type} \texttt{hom-ring-member-iff} \\
mentioned by 5 result(s)

\smallskip

\noindent\textbf{\texttt{hom-ringoid(a, b)}}\quad\textit{functoid}\index{hom-ringoid(a, b)@\texttt{hom-ringoid(a, b)}}\\
reads: the ringoid morphisms from \$1 to \$2 \\
definition:
\begin{quote}\ttfamily\footnotesize\raggedright\noindent
\{homf\_ in fun(\allowbreak{}carr(\allowbreak{}a), carr(\allowbreak{}b)): is-hom-ringoid(\allowbreak{}a, b, homf\_)\}
\end{quote}
characterized by \texttt{hom-ringoid-in-set}  \textit{(proven)}:
\begin{quote}\ttfamily\footnotesize\raggedright\noindent
forall(\allowbreak{}[a, b], is-ringoid(\allowbreak{}a) implies is-ringoid(\allowbreak{}b) implies hom-ringoid(\allowbreak{}a, b) in set)
\end{quote}
also: \texttt{hom-ringoid-pre-type} \texttt{hom-ringoid-post-type} \texttt{hom-ringoid-member-iff} \\
mentioned by 5 result(s)

\smallskip

\noindent\textbf{\texttt{hom-semi-inner-product-space(a, b)}}\quad\textit{functoid}\index{hom-semi-inner-product-space(a, b)@\texttt{hom-semi-inner-product-space(a, b)}}\\
reads: the semi inner product space morphisms from \$1 to \$2 \\
definition:
\begin{quote}\ttfamily\footnotesize\raggedright\noindent
\{homf\_ in fun(\allowbreak{}vec(\allowbreak{}a), vec(\allowbreak{}b)): is-hom-semi-inner-product-space(\allowbreak{}a, b, homf\_)\}
\end{quote}

\smallskip

\noindent\textbf{\texttt{hom-semigroup(a, b)}}\quad\textit{functoid}\index{hom-semigroup(a, b)@\texttt{hom-semigroup(a, b)}}\\
reads: the semigroup morphisms from \$1 to \$2 \\
definition:
\begin{quote}\ttfamily\footnotesize\raggedright\noindent
\{homf\_ in fun(\allowbreak{}carr(\allowbreak{}a), carr(\allowbreak{}b)): is-hom-semigroup(\allowbreak{}a, b, homf\_)\}
\end{quote}
characterized by \texttt{hom-semigroup-in-set}  \textit{(proven)}:
\begin{quote}\ttfamily\footnotesize\raggedright\noindent
forall(\allowbreak{}[a, b], is-semigroup(\allowbreak{}a) implies is-semigroup(\allowbreak{}b) implies hom-semigroup(\allowbreak{}a, b) in set)
\end{quote}
also: \texttt{hom-semigroup-pre-type} \texttt{hom-semigroup-post-type} \texttt{hom-semigroup-member-iff} \\
mentioned by 5 result(s)

\smallskip

\noindent\textbf{\texttt{hom-setoid(a, b)}}\quad\textit{functoid}\index{hom-setoid(a, b)@\texttt{hom-setoid(a, b)}}\\
reads: the setoid morphisms from \$1 to \$2 \\
definition:
\begin{quote}\ttfamily\footnotesize\raggedright\noindent
\{homf\_ in fun(\allowbreak{}pts(\allowbreak{}a), pts(\allowbreak{}b)): is-hom-setoid(\allowbreak{}a, b, homf\_)\}
\end{quote}
characterized by \texttt{hom-setoid-in-set}  \textit{(proven)}:
\begin{quote}\ttfamily\footnotesize\raggedright\noindent
forall(\allowbreak{}[a, b], is-setoid(\allowbreak{}a) implies is-setoid(\allowbreak{}b) implies hom-setoid(\allowbreak{}a, b) in set)
\end{quote}
also: \texttt{hom-setoid-pre-type} \texttt{hom-setoid-post-type} \texttt{hom-setoid-member-iff} \\
mentioned by 5 result(s)

\smallskip

\noindent\textbf{\texttt{hom-sk-space(a, b)}}\quad\textit{functoid}\index{hom-sk-space(a, b)@\texttt{hom-sk-space(a, b)}}\\
reads: the sk space morphisms from \$1 to \$2 \\
definition:
\begin{quote}\ttfamily\footnotesize\raggedright\noindent
\{homf\_ in fun(\allowbreak{}sk-names(\allowbreak{}a), sk-names(\allowbreak{}b)): is-hom-sk-space(\allowbreak{}a, b, homf\_)\}
\end{quote}

\smallskip

\noindent\textbf{\texttt{hom-top-group(a, b)}}\quad\textit{functoid}\index{hom-top-group(a, b)@\texttt{hom-top-group(a, b)}}\\
reads: the top group morphisms from \$1 to \$2 \\
definition:
\begin{quote}\ttfamily\footnotesize\raggedright\noindent
\{homf\_ in fun(\allowbreak{}carr(\allowbreak{}a), carr(\allowbreak{}b)): is-hom-top-group(\allowbreak{}a, b, homf\_)\}
\end{quote}

\smallskip

\noindent\textbf{\texttt{hom-top-space(a, b)}}\quad\textit{functoid}\index{hom-top-space(a, b)@\texttt{hom-top-space(a, b)}}\\
reads: the top space morphisms from \$1 to \$2 \\
definition:
\begin{quote}\ttfamily\footnotesize\raggedright\noindent
\{homf\_ in fun(\allowbreak{}pts(\allowbreak{}a), pts(\allowbreak{}b)): is-hom-top-space(\allowbreak{}a, b, homf\_)\}
\end{quote}
characterized by \texttt{hom-top-space-in-set}  \textit{(proven)}:
\begin{quote}\ttfamily\footnotesize\raggedright\noindent
forall(\allowbreak{}[a, b], is-top-space(\allowbreak{}a) implies is-top-space(\allowbreak{}b) implies hom-top-space(\allowbreak{}a, b) in set)
\end{quote}
also: \texttt{hom-top-space-pre-type} \texttt{hom-top-space-post-type} \texttt{hom-top-space-member-iff} \\
mentioned by 5 result(s)

\smallskip

\noindent\textbf{\texttt{hom-uniformly-continuous(a, b)}}\quad\textit{functoid}\index{hom-uniformly-continuous(a, b)@\texttt{hom-uniformly-continuous(a, b)}}\\
reads: the uniformly continuous maps from \$1 to \$2 \\
definition:
\begin{quote}\ttfamily\footnotesize\raggedright\noindent
\{homf\_ in fun(\allowbreak{}pts(\allowbreak{}a), pts(\allowbreak{}b)): is-uniformly-continuous-arrow(\allowbreak{}a, b, homf\_)\}
\end{quote}
characterized by \texttt{hom-uniformly-continuous-in-set}  \textit{(proven)}:
\begin{quote}\ttfamily\footnotesize\raggedright\noindent
forall(\allowbreak{}[a, b], is-metric-space(\allowbreak{}a) implies is-metric-space(\allowbreak{}b) implies hom-uniformly-continuous(\allowbreak{}a, b) in set)
\end{quote}
also: \texttt{hom-uniformly-continuous-pre-type} \texttt{hom-uniformly-continuous-post-type} \texttt{hom-uniformly-continuous-member-iff} \texttt{hom-lipschitz-subset-hom-uniformly-continuous} \texttt{hom-uniformly-continuous-subset-hom-continuous} \\
mentioned by 7 result(s)

\smallskip

\noindent\textbf{\texttt{hom-unital-banach-algebra(a, b)}}\quad\textit{functoid}\index{hom-unital-banach-algebra(a, b)@\texttt{hom-unital-banach-algebra(a, b)}}\\
reads: the unital banach algebra morphisms from \$1 to \$2 \\
definition:
\begin{quote}\ttfamily\footnotesize\raggedright\noindent
\{homf\_ in fun(\allowbreak{}vec(\allowbreak{}a), vec(\allowbreak{}b)): is-hom-unital-banach-algebra(\allowbreak{}a, b, homf\_)\}
\end{quote}

\smallskip

\noindent\textbf{\texttt{hom-unital-c-star-algebra(a, b)}}\quad\textit{functoid}\index{hom-unital-c-star-algebra(a, b)@\texttt{hom-unital-c-star-algebra(a, b)}}\\
reads: the unital c star algebra morphisms from \$1 to \$2 \\
definition:
\begin{quote}\ttfamily\footnotesize\raggedright\noindent
\{homf\_ in fun(\allowbreak{}vec(\allowbreak{}a), vec(\allowbreak{}b)): is-hom-unital-c-star-algebra(\allowbreak{}a, b, homf\_)\}
\end{quote}

\smallskip

\noindent\textbf{\texttt{hom-unital-involutive-banach-algebra(a, b)}}\quad\textit{functoid}\index{hom-unital-involutive-banach-algebra(a, b)@\texttt{hom-unital-involutive-banach-algebra(a, b)}}\\
reads: the unital involutive banach algebra morphisms from \$1 to \$2 \\
definition:
\begin{quote}\ttfamily\footnotesize\raggedright\noindent
\{homf\_ in fun(\allowbreak{}vec(\allowbreak{}a), vec(\allowbreak{}b)): is-hom-unital-involutive-banach-algebra(\allowbreak{}a, b, homf\_)\}
\end{quote}

\smallskip

\noindent\textbf{\texttt{hom-vector-space(a, b)}}\quad\textit{functoid}\index{hom-vector-space(a, b)@\texttt{hom-vector-space(a, b)}}\\
reads: the vector space morphisms from \$1 to \$2 \\
definition:
\begin{quote}\ttfamily\footnotesize\raggedright\noindent
\{homf\_ in fun(\allowbreak{}vec(\allowbreak{}a), vec(\allowbreak{}b)): is-hom-vector-space(\allowbreak{}a, b, homf\_)\}
\end{quote}
characterized by \texttt{hom-vector-space-in-set}  \textit{(proven)}:
\begin{quote}\ttfamily\footnotesize\raggedright\noindent
forall(\allowbreak{}[a, b], is-vector-space(\allowbreak{}a) implies is-vector-space(\allowbreak{}b) implies hom-vector-space(\allowbreak{}a, b) in set)
\end{quote}
also: \texttt{hom-vector-space-pre-type} \texttt{hom-vector-space-post-type} \texttt{hom-vector-space-member-iff} \\
mentioned by 5 result(s)

\smallskip

\noindent\textbf{\texttt{homog-equiv(u, w)}}\quad\textit{predicate}\index{homog-equiv(u, w)@\texttt{homog-equiv(u, w)}}\\
reads: \$1 and \$2 are proportional homogeneous coordinates \\
defined by \texttt{homog-equiv} \\
characterized by \texttt{homog-equiv}  \textit{(definitional)}:
\begin{quote}\ttfamily\footnotesize\raggedright\noindent
forall(\allowbreak{}[u, w], homog-equiv(\allowbreak{}u, w) iff forsome(\allowbreak{}[hez\_ in cc], not(\allowbreak{}hez\_ = 0) and u = [hez\_ * nth(\allowbreak{}1, w), hez\_ * nth(\allowbreak{}2, w)]))
\end{quote}
mentioned by 6 result(s)

\smallskip

\noindent\textbf{\texttt{homog-minorant(m, p)}}\quad\textit{functoid}\index{homog-minorant(m, p)@\texttt{homog-minorant(m, p)}}\\
reads: the positively homogeneous minorant of \$2 \\
definition:
\begin{quote}\ttfamily\footnotesize\raggedright\noindent
vnb-lambda(\allowbreak{}hmx\_, vec(\allowbreak{}m), inf(\allowbreak{}homog-set(\allowbreak{}m, p, hmx\_)))
\end{quote}
characterized by \texttt{homog-minorant-sublinear}  \textit{(proven)}:
\begin{quote}\ttfamily\footnotesize\raggedright\noindent
forall(\allowbreak{}[m, p], is-module(\allowbreak{}m) implies scal(\allowbreak{}m) = normed-field-as-commutative-ring(\allowbreak{}rr-normed-field) implies is-convex-functional(\allowbreak{}m, p) implies is-sublinear(\allowbreak{}m, homog-minorant(\allowbreak{}m, p)))
\end{quote}
mentioned by 3 result(s)

\smallskip

\noindent\textbf{\texttt{homog-point(u, v)}}\quad\textit{functoid}\index{homog-point(u, v)@\texttt{homog-point(u, v)}}\\
reads: the point of the extended plane with homogeneous coordinates (\$1, \$2) \\
definition:
\begin{quote}\ttfamily\footnotesize\raggedright\noindent
if(\allowbreak{}v = 0, cc-inf, u * recip(\allowbreak{}v))
\end{quote}
characterized by \texttt{homog-point-cover}  \textit{(proven)}:
\begin{quote}\ttfamily\footnotesize\raggedright\noindent
forall(\allowbreak{}[w in cc-hat], forsome(\allowbreak{}[u in cc, v in cc], not(\allowbreak{}u = 0 and v = 0) and w = homog-point(\allowbreak{}u, v)))
\end{quote}
also: \texttt{homog-point-scale} \texttt{homog-point-at-inf} \texttt{homog-point-finite} \texttt{homog-point-in-cc-hat} \texttt{homog-point-sphere-eq} \\
mentioned by 18 result(s)

\smallskip

\noindent\textbf{\texttt{homog-set(m, p, x)}}\quad\textit{functoid}\index{homog-set(m, p, x)@\texttt{homog-set(m, p, x)}}\\
reads: the reals r with \$2(t \$3) <= t r for some t > 0 \\
definition:
\begin{quote}\ttfamily\footnotesize\raggedright\noindent
\{hmr\_ in rr: forsome(\allowbreak{}[hmt\_ in rr], 0 < hmt\_ and p(\allowbreak{}(\allowbreak{}act(\allowbreak{}m))(\allowbreak{}hmt\_, x)) <= hmt\_ * hmr\_)\}
\end{quote}
characterized by \texttt{homog-set-glb}  \textit{(proven)}:
\begin{quote}\ttfamily\footnotesize\raggedright\noindent
forall(\allowbreak{}[m, p, z], is-module(\allowbreak{}m) implies scal(\allowbreak{}m) = normed-field-as-commutative-ring(\allowbreak{}rr-normed-field) implies is-convex-functional(\allowbreak{}m, p) implies z in vec(\allowbreak{}m) implies inf(\allowbreak{}homog-set(\allowbreak{}m, p, z)) in rr and forall(\allowbreak{}[lbw\_ in homog-set(\allowbreak{}m, p, z)], inf(\allowbreak{}homog-set(\allowbreak{}m, p, z)) <= lbw\_) and forall(\allowbreak{}[lbb\_ in rr], forall(\allowbreak{}[lbw\_ in homog-set(\allowbreak{}m, p, z)], lbb\_ <= lbw\_) implies lbb\_ <= inf(\allowbreak{}homog-set(\allowbreak{}m, p, z))))
\end{quote}
also: \texttt{homog-set-sum} \texttt{homog-set-elim} \texttt{homog-set-zero} \texttt{homog-set-intro} \texttt{homog-set-scale} \texttt{homog-set-unfold} \texttt{homog-set-unscale} \\
mentioned by 21 result(s)

\smallskip

\noindent\textbf{\texttt{homotopic(hmu\_, hmg\_, hmr\_, hma\_, hmb\_)}}\quad\textit{predicate}\index{homotopic(hmu_, hmg_, hmr_, hma_, hmb_)@\texttt{homotopic(hmu\_, hmg\_, hmr\_, hma\_, hmb\_)}}\\
reads: \$2 and \$3 are homotopic in \$1 with endpoints fixed \\
defined by \texttt{homotopic} \\
characterized by \texttt{homotopic}  \textit{(definitional)}:
\begin{quote}\ttfamily\footnotesize\raggedright\noindent
forall(\allowbreak{}[hmu\_, hmg\_, hmr\_, hma\_, hmb\_], homotopic(\allowbreak{}hmu\_, hmg\_, hmr\_, hma\_, hmb\_) iff is-path(\allowbreak{}hmg\_, hma\_, hmb\_) and is-path(\allowbreak{}hmr\_, hma\_, hmb\_) and hmr\_(\allowbreak{}hma\_) = hmg\_(\allowbreak{}hma\_) and hmr\_(\allowbreak{}hmb\_) = hmg\_(\allowbreak{}hmb\_) and forsome(\allowbreak{}[hmf\_ in fun(\allowbreak{}cartesian(\allowbreak{}ccint(\allowbreak{}0, 1), ccint(\allowbreak{}hma\_, hmb\_)), hmu\_)], forall(\allowbreak{}[hmz\_ in cartesian(\allowbreak{}ccint(\allowbreak{}0, 1), ccint(\allowbreak{}hma\_, hmb\_))], is-continuous-at(\allowbreak{}product2-ms(\allowbreak{}subspace-ms(\allowbreak{}rr-ms, ccint(\allowbreak{}0, 1)), subspace-ms(\allowbreak{}rr-ms, ccint(\allowbreak{}hma\_, hmb\_))), cc-ms, hmf\_, hmz\_)) and forall(\allowbreak{}[hmt\_ in ccint(\allowbreak{}hma\_, hmb\_)], hmf\_(\allowbreak{}[0, hmt\_]) = hmg\_(\allowbreak{}hmt\_) and hmf\_(\allowbreak{}[1, hmt\_]) = hmr\_(\allowbreak{}hmt\_)) and forall(\allowbreak{}[hms\_ in ccint(\allowbreak{}0, 1)], hmf\_(\allowbreak{}[hms\_, hma\_]) = hmg\_(\allowbreak{}hma\_) and hmf\_(\allowbreak{}[hms\_, hmb\_]) = hmg\_(\allowbreak{}hmb\_))))
\end{quote}
also: \texttt{null-homotopic} \texttt{homotopic-sym} \texttt{homotopic-refl} \texttt{homotopic-trans} \texttt{homotopic-traces} \texttt{good-approx-homotopic} \texttt{straight-line-homotopic} \\
mentioned by 15 result(s)

\smallskip

\noindent\textbf{\texttt{hopf-map(w)}}\quad\textit{functoid}\index{hopf-map(w)@\texttt{hopf-map(w)}}\\
reads: the Hopf map pi(\$1) = (2 u1 conj(u2), |u1|\textasciicircum{}2 - |u2|\textasciicircum{}2) \\
definition:
\begin{quote}\ttfamily\footnotesize\raggedright\noindent
[2 * (\allowbreak{}nth(\allowbreak{}1, w) * conjugate(\allowbreak{}nth(\allowbreak{}2, w))), magnitude(\allowbreak{}nth(\allowbreak{}1, w)) * magnitude(\allowbreak{}nth(\allowbreak{}1, w)) - magnitude(\allowbreak{}nth(\allowbreak{}2, w)) * magnitude(\allowbreak{}nth(\allowbreak{}2, w))]
\end{quote}
characterized by \texttt{hopf-map-pair}  \textit{(proven)}:
\begin{quote}\ttfamily\footnotesize\raggedright\noindent
forall(\allowbreak{}[a in cc, b in cc], hopf-map(\allowbreak{}[a, b]) = [2 * (\allowbreak{}a * conjugate(\allowbreak{}b)), magnitude(\allowbreak{}a) * magnitude(\allowbreak{}a) - magnitude(\allowbreak{}b) * magnitude(\allowbreak{}b)])
\end{quote}
also: \texttt{hopf-map-in-cr} \texttt{hopf-maps-sphere} \texttt{hopf-map-continuous} \texttt{hopf-map-unimodular-scale} \texttt{hopf-map-continuous-product} \texttt{hopf-map-continuous-at-product} \\
mentioned by 19 result(s)

\smallskip

\noindent\textbf{\texttt{hopf-quotient-map}}\quad\textit{constant}\index{hopf-quotient-map@\texttt{hopf-quotient-map}}\\
reads: the bijection S\textasciicircum{}2 -> C u \{oo\} through which the Hopf map factors \\
defined by \texttt{hopf-quotient-map-def} \\
characterized by \texttt{hopf-quotient-map-onto}  \textit{(proven)}:
\begin{quote}\ttfamily\footnotesize\raggedright\noindent
forall(\allowbreak{}[w in cc-hat], forsome(\allowbreak{}[z in unit-sphere-cr], hopf-quotient-map(\allowbreak{}z) = w))
\end{quote}
also: \texttt{hopf-quotient-map-apply} \texttt{hopf-quotient-map-injective} \\
mentioned by 7 result(s)

\smallskip

\noindent\textbf{\texttt{hs-ip(kqv\_, kqw\_, kqu\_, kqt\_)}}\quad\textit{functoid}\index{hs-ip(kqv_, kqw_, kqu_, kqt_)@\texttt{hs-ip(kqv\_, kqw\_, kqu\_, kqt\_)}}\\
reads: the Hilbert-Schmidt inner product of \$3 and \$4 \\
definition:
\begin{quote}\ttfamily\footnotesize\raggedright\noindent
cc-family-sum(\allowbreak{}vnb-lambda(\allowbreak{}kqe\_, canon-onb(\allowbreak{}kqv\_), (\allowbreak{}ip(\allowbreak{}kqw\_))(\allowbreak{}kqu\_(\allowbreak{}kqe\_), kqt\_(\allowbreak{}kqe\_))))
\end{quote}
characterized by \texttt{hs-ip-self}  \textit{(proven)}:
\begin{quote}\ttfamily\footnotesize\raggedright\noindent
forall(\allowbreak{}[v, w], is-hilbert-space(\allowbreak{}v) implies is-hilbert-space(\allowbreak{}w) implies forall(\allowbreak{}[t\_], is-hilbert-schmidt(\allowbreak{}v, w, t\_) implies forall(\allowbreak{}[ix, e], is-orthonormal-basis(\allowbreak{}v, ix, e) implies hs-ip(\allowbreak{}v, w, t\_, t\_) = hs-sum(\allowbreak{}v, w, ix, e, t\_))))
\end{quote}
also: \texttt{hs-ip-unfold} \texttt{hs-ip-adjoint} \texttt{hs-ip-add-left} \texttt{hs-ip-conj-sym} \texttt{hs-ip-definite} \texttt{hs-ip-any-basis} \texttt{hs-ip-homog-left} \texttt{hs-ip-self-nonneg} \\
mentioned by 19 result(s)

\smallskip

\noindent\textbf{\texttt{hs-norm(kqv\_, kqw\_, kqt\_)}}\quad\textit{functoid}\index{hs-norm(kqv_, kqw_, kqt_)@\texttt{hs-norm(kqv\_, kqw\_, kqt\_)}}\\
reads: the Hilbert-Schmidt norm of \$3 \\
definition:
\begin{quote}\ttfamily\footnotesize\raggedright\noindent
sqrt(\allowbreak{}hs-sum(\allowbreak{}kqv\_, kqw\_, canon-onb(\allowbreak{}kqv\_), vnb-lambda(\allowbreak{}kqe\_, canon-onb(\allowbreak{}kqv\_), kqe\_), kqt\_))
\end{quote}
characterized by \texttt{hs-norm-unfold}  \textit{(proven)}:
\begin{quote}\ttfamily\footnotesize\raggedright\noindent
forall(\allowbreak{}[kqv\_, kqw\_, kqt\_], hs-norm(\allowbreak{}kqv\_, kqw\_, kqt\_) == sqrt(\allowbreak{}hs-sum(\allowbreak{}kqv\_, kqw\_, canon-onb(\allowbreak{}kqv\_), vnb-lambda(\allowbreak{}kqe\_, canon-onb(\allowbreak{}kqv\_), kqe\_), kqt\_)))
\end{quote}
also: \texttt{hs-norm-any-basis} \\
mentioned by 5 result(s)

\smallskip

\noindent\textbf{\texttt{hs-ops(kqv\_, kqw\_)}}\quad\textit{functoid}\index{hs-ops(kqv_, kqw_)@\texttt{hs-ops(kqv\_, kqw\_)}}\\
reads: the Hilbert-Schmidt operators from \$1 to \$2 \\
definition:
\begin{quote}\ttfamily\footnotesize\raggedright\noindent
\{kqt\_ in bounded-ops(\allowbreak{}kqv\_, kqw\_): is-hilbert-schmidt(\allowbreak{}kqv\_, kqw\_, kqt\_)\}
\end{quote}
characterized by \texttt{hs-ops-unfold}  \textit{(proven)}:
\begin{quote}\ttfamily\footnotesize\raggedright\noindent
forall(\allowbreak{}[kqv\_, kqw\_], hs-ops(\allowbreak{}kqv\_, kqw\_) == \{kqt\_ in bounded-ops(\allowbreak{}kqv\_, kqw\_): is-hilbert-schmidt(\allowbreak{}kqv\_, kqw\_, kqt\_)\})
\end{quote}
also: \texttt{hs-ops-membership} \\
mentioned by 4 result(s)

\smallskip

\noindent\textbf{\texttt{hs-sum(kqv\_, kqw\_, kqi\_, kqe\_, kqt\_)}}\quad\textit{functoid}\index{hs-sum(kqv_, kqw_, kqi_, kqe_, kqt_)@\texttt{hs-sum(kqv\_, kqw\_, kqi\_, kqe\_, kqt\_)}}\\
reads: the sum of ||\$5 \$4\_i||\textasciicircum{}2 over \$3 \\
definition:
\begin{quote}\ttfamily\footnotesize\raggedright\noindent
esum(\allowbreak{}vnb-lambda(\allowbreak{}kqk\_, kqi\_, ip-norm(\allowbreak{}kqw\_, kqt\_(\allowbreak{}kqe\_(\allowbreak{}kqk\_))) * ip-norm(\allowbreak{}kqw\_, kqt\_(\allowbreak{}kqe\_(\allowbreak{}kqk\_)))))
\end{quote}
characterized by \texttt{hs-sum-polar}  \textit{(proven)}:
\begin{quote}\ttfamily\footnotesize\raggedright\noindent
forall(\allowbreak{}[v, w], is-hilbert-space(\allowbreak{}v) implies is-hilbert-space(\allowbreak{}w) implies forall(\allowbreak{}[s\_, t\_], is-hilbert-schmidt(\allowbreak{}v, w, s\_) implies is-hilbert-schmidt(\allowbreak{}v, w, t\_) implies forall(\allowbreak{}[c\_ in cc, ix, e], is-orthonormal-basis(\allowbreak{}v, ix, e) implies hs-sum(\allowbreak{}v, w, ix, e, bop-sum(\allowbreak{}v, w, s\_, bop-scale(\allowbreak{}v, w, c\_, t\_))) = (\allowbreak{}hs-sum(\allowbreak{}v, w, ix, e, s\_) + conjugate(\allowbreak{}c\_) * cc-family-sum(\allowbreak{}vnb-lambda(\allowbreak{}kqk\_, ix, (\allowbreak{}ip(\allowbreak{}w))(\allowbreak{}s\_(\allowbreak{}e(\allowbreak{}kqk\_)), t\_(\allowbreak{}e(\allowbreak{}kqk\_)))))) + (\allowbreak{}c\_ * conjugate(\allowbreak{}cc-family-sum(\allowbreak{}vnb-lambda(\allowbreak{}kqk\_, ix, (\allowbreak{}ip(\allowbreak{}w))(\allowbreak{}s\_(\allowbreak{}e(\allowbreak{}kqk\_)), t\_(\allowbreak{}e(\allowbreak{}kqk\_)))))) + (\allowbreak{}c\_ * conjugate(\allowbreak{}c\_)) * hs-sum(\allowbreak{}v, w, ix, e, t\_)))))
\end{quote}
also: \texttt{hs-sum-add-le} \texttt{hs-sum-in-rps} \texttt{hs-sum-unfold} \texttt{hs-sum-product} \texttt{hs-sum-comp-left} \texttt{hs-sum-adjoint-eq} \texttt{hs-sum-compact-sa} \texttt{hs-sum-eigenbasis} ... (12 in all) \\
mentioned by 31 result(s)

\smallskip

\noindent\textbf{\texttt{hsum(bov\_, bow\_)}}\quad\textit{functoid}\index{hsum(bov_, bow_)@\texttt{hsum(bov\_, bow\_)}}\\
reads: the Hilbert direct sum of \$1 and \$2 \\
definition:
\begin{quote}\ttfamily\footnotesize\raggedright\noindent
[normed-field-as-commutative-ring(\allowbreak{}cc-normed-field), cartesian(\allowbreak{}vec(\allowbreak{}bov\_), vec(\allowbreak{}bow\_)), hsum-vadd(\allowbreak{}bov\_, bow\_), [vzero(\allowbreak{}bov\_), vzero(\allowbreak{}bow\_)], hsum-vneg(\allowbreak{}bov\_, bow\_), hsum-act(\allowbreak{}bov\_, bow\_), hsum-ip(\allowbreak{}bov\_, bow\_)]
\end{quote}
characterized by \texttt{hsum-ip-acc}  \textit{(proven)}:
\begin{quote}\ttfamily\footnotesize\raggedright\noindent
forall(\allowbreak{}[v, w], is-complex-inner-product-space(\allowbreak{}v) implies is-complex-inner-product-space(\allowbreak{}w) implies ip(\allowbreak{}hsum(\allowbreak{}v, w)) == hsum-ip(\allowbreak{}v, w))
\end{quote}
also: \texttt{hsum-is-set} \texttt{hsum-vec-eq} \texttt{hsum-act-acc} \texttt{hsum-diag-id} \texttt{hsum-ip-pair} \texttt{hsum-is-cips} \texttt{hsum-vec-acc} \texttt{hsum-act-pair} ... (35 in all) \\
mentioned by 76 result(s)

\smallskip

\noindent\textbf{\texttt{hsum-act(bov\_, bow\_)}}\quad\textit{functoid}\index{hsum-act(bov_, bow_)@\texttt{hsum-act(bov\_, bow\_)}}\\
definition:
\begin{quote}\ttfamily\footnotesize\raggedright\noindent
vnb-lambda(\allowbreak{}[bhc\_, bhp\_], cartesian(\allowbreak{}cc, cartesian(\allowbreak{}vec(\allowbreak{}bov\_), vec(\allowbreak{}bow\_))), [(\allowbreak{}act(\allowbreak{}bov\_))(\allowbreak{}bhc\_, nth(\allowbreak{}1, bhp\_)), (\allowbreak{}act(\allowbreak{}bow\_))(\allowbreak{}bhc\_, nth(\allowbreak{}2, bhp\_))])
\end{quote}
characterized by \texttt{hsum-act-acc}  \textit{(proven)}:
\begin{quote}\ttfamily\footnotesize\raggedright\noindent
forall(\allowbreak{}[v, w], is-complex-inner-product-space(\allowbreak{}v) implies is-complex-inner-product-space(\allowbreak{}w) implies act(\allowbreak{}hsum(\allowbreak{}v, w)) == hsum-act(\allowbreak{}v, w))
\end{quote}
also: \texttt{hsum-act-type} \\
mentioned by 3 result(s)

\smallskip

\noindent\textbf{\texttt{hsum-diag(amw\_, amv\_, ama\_, amb\_)}}\quad\textit{functoid}\index{hsum-diag(amw_, amv_, ama_, amb_)@\texttt{hsum-diag(amw\_, amv\_, ama\_, amb\_)}}\\
reads: the diagonal operator \$3 (+) \$4 \\
definition:
\begin{quote}\ttfamily\footnotesize\raggedright\noindent
vnb-lambda(\allowbreak{}amp\_, cartesian(\allowbreak{}vec(\allowbreak{}amw\_), vec(\allowbreak{}amv\_)), [ama\_(\allowbreak{}nth(\allowbreak{}1, amp\_)), amb\_(\allowbreak{}nth(\allowbreak{}2, amp\_))])
\end{quote}
characterized by \texttt{hsum-diag-id}  \textit{(proven)}:
\begin{quote}\ttfamily\footnotesize\raggedright\noindent
forall(\allowbreak{}[w, v], is-complex-inner-product-space(\allowbreak{}w) implies is-complex-inner-product-space(\allowbreak{}v) implies hsum-diag(\allowbreak{}w, v, bop-id(\allowbreak{}w), bop-id(\allowbreak{}v)) = bop-id(\allowbreak{}hsum(\allowbreak{}w, v)))
\end{quote}
also: \texttt{hsum-diag-sum} \texttt{hsum-diag-comp} \texttt{hsum-diag-apply} \texttt{hsum-diag-scale} \texttt{hsum-diag-value} \texttt{hsum-diag-adjoint} \texttt{hsum-diag-bounded} \\
mentioned by 17 result(s)

\smallskip

\noindent\textbf{\texttt{hsum-in1(amw\_, amv\_)}}\quad\textit{functoid}\index{hsum-in1(amw_, amv_)@\texttt{hsum-in1(amw\_, amv\_)}}\\
reads: the injection of \$1 into the Hilbert direct sum of \$1 and \$2 \\
definition:
\begin{quote}\ttfamily\footnotesize\raggedright\noindent
vnb-lambda(\allowbreak{}amx\_, vec(\allowbreak{}amw\_), [amx\_, vzero(\allowbreak{}amv\_)])
\end{quote}
characterized by \texttt{hsum-in1-apply}  \textit{(proven)}:
\begin{quote}\ttfamily\footnotesize\raggedright\noindent
forall(\allowbreak{}[w, v, x\_], is-complex-inner-product-space(\allowbreak{}w) implies is-complex-inner-product-space(\allowbreak{}v) implies x\_ in vec(\allowbreak{}w) implies (\allowbreak{}hsum-in1(\allowbreak{}w, v))(\allowbreak{}x\_) = [x\_, vzero(\allowbreak{}v)])
\end{quote}
also: \texttt{hsum-in1-bounded} \\
mentioned by 3 result(s)

\smallskip

\noindent\textbf{\texttt{hsum-in2(amw\_, amv\_)}}\quad\textit{functoid}\index{hsum-in2(amw_, amv_)@\texttt{hsum-in2(amw\_, amv\_)}}\\
reads: the injection of \$2 into the Hilbert direct sum of \$1 and \$2 \\
definition:
\begin{quote}\ttfamily\footnotesize\raggedright\noindent
vnb-lambda(\allowbreak{}amy\_, vec(\allowbreak{}amv\_), [vzero(\allowbreak{}amw\_), amy\_])
\end{quote}
characterized by \texttt{hsum-in2-apply}  \textit{(proven)}:
\begin{quote}\ttfamily\footnotesize\raggedright\noindent
forall(\allowbreak{}[w, v, y\_], is-complex-inner-product-space(\allowbreak{}w) implies is-complex-inner-product-space(\allowbreak{}v) implies y\_ in vec(\allowbreak{}v) implies (\allowbreak{}hsum-in2(\allowbreak{}w, v))(\allowbreak{}y\_) = [vzero(\allowbreak{}w), y\_])
\end{quote}
also: \texttt{hsum-in2-bounded} \\
mentioned by 3 result(s)

\smallskip

\noindent\textbf{\texttt{hsum-ip(bov\_, bow\_)}}\quad\textit{functoid}\index{hsum-ip(bov_, bow_)@\texttt{hsum-ip(bov\_, bow\_)}}\\
definition:
\begin{quote}\ttfamily\footnotesize\raggedright\noindent
vnb-lambda(\allowbreak{}[bhp\_, bhq\_], cartesian(\allowbreak{}cartesian(\allowbreak{}vec(\allowbreak{}bov\_), vec(\allowbreak{}bow\_)), cartesian(\allowbreak{}vec(\allowbreak{}bov\_), vec(\allowbreak{}bow\_))), (\allowbreak{}ip(\allowbreak{}bov\_))(\allowbreak{}nth(\allowbreak{}1, bhp\_), nth(\allowbreak{}1, bhq\_)) + (\allowbreak{}ip(\allowbreak{}bow\_))(\allowbreak{}nth(\allowbreak{}2, bhp\_), nth(\allowbreak{}2, bhq\_)))
\end{quote}
characterized by \texttt{hsum-ip-acc}  \textit{(proven)}:
\begin{quote}\ttfamily\footnotesize\raggedright\noindent
forall(\allowbreak{}[v, w], is-complex-inner-product-space(\allowbreak{}v) implies is-complex-inner-product-space(\allowbreak{}w) implies ip(\allowbreak{}hsum(\allowbreak{}v, w)) == hsum-ip(\allowbreak{}v, w))
\end{quote}
also: \texttt{hsum-ip-type} \\
mentioned by 3 result(s)

\smallskip

\noindent\textbf{\texttt{hsum-pr1(amw\_, amv\_)}}\quad\textit{functoid}\index{hsum-pr1(amw_, amv_)@\texttt{hsum-pr1(amw\_, amv\_)}}\\
reads: the projection of the Hilbert direct sum of \$1 and \$2 onto \$1 \\
definition:
\begin{quote}\ttfamily\footnotesize\raggedright\noindent
vnb-lambda(\allowbreak{}amp\_, cartesian(\allowbreak{}vec(\allowbreak{}amw\_), vec(\allowbreak{}amv\_)), nth(\allowbreak{}1, amp\_))
\end{quote}
characterized by \texttt{hsum-pr1-apply}  \textit{(proven)}:
\begin{quote}\ttfamily\footnotesize\raggedright\noindent
forall(\allowbreak{}[w, v, x\_, y\_], is-complex-inner-product-space(\allowbreak{}w) implies is-complex-inner-product-space(\allowbreak{}v) implies x\_ in vec(\allowbreak{}w) implies y\_ in vec(\allowbreak{}v) implies (\allowbreak{}hsum-pr1(\allowbreak{}w, v))(\allowbreak{}[x\_, y\_]) = x\_)
\end{quote}
also: \texttt{hsum-pr1-value} \texttt{hsum-pr1-bounded} \\
mentioned by 5 result(s)

\smallskip

\noindent\textbf{\texttt{hsum-pr2(amw\_, amv\_)}}\quad\textit{functoid}\index{hsum-pr2(amw_, amv_)@\texttt{hsum-pr2(amw\_, amv\_)}}\\
reads: the projection of the Hilbert direct sum of \$1 and \$2 onto \$2 \\
definition:
\begin{quote}\ttfamily\footnotesize\raggedright\noindent
vnb-lambda(\allowbreak{}amp\_, cartesian(\allowbreak{}vec(\allowbreak{}amw\_), vec(\allowbreak{}amv\_)), nth(\allowbreak{}2, amp\_))
\end{quote}
characterized by \texttt{hsum-pr2-apply}  \textit{(proven)}:
\begin{quote}\ttfamily\footnotesize\raggedright\noindent
forall(\allowbreak{}[w, v, x\_, y\_], is-complex-inner-product-space(\allowbreak{}w) implies is-complex-inner-product-space(\allowbreak{}v) implies x\_ in vec(\allowbreak{}w) implies y\_ in vec(\allowbreak{}v) implies (\allowbreak{}hsum-pr2(\allowbreak{}w, v))(\allowbreak{}[x\_, y\_]) = y\_)
\end{quote}
also: \texttt{hsum-pr2-value} \texttt{hsum-pr2-bounded} \\
mentioned by 5 result(s)

\smallskip

\noindent\textbf{\texttt{hsum-vadd(bov\_, bow\_)}}\quad\textit{functoid}\index{hsum-vadd(bov_, bow_)@\texttt{hsum-vadd(bov\_, bow\_)}}\\
definition:
\begin{quote}\ttfamily\footnotesize\raggedright\noindent
vnb-lambda(\allowbreak{}[bhp\_, bhq\_], cartesian(\allowbreak{}cartesian(\allowbreak{}vec(\allowbreak{}bov\_), vec(\allowbreak{}bow\_)), cartesian(\allowbreak{}vec(\allowbreak{}bov\_), vec(\allowbreak{}bow\_))), [(\allowbreak{}vadd(\allowbreak{}bov\_))(\allowbreak{}nth(\allowbreak{}1, bhp\_), nth(\allowbreak{}1, bhq\_)), (\allowbreak{}vadd(\allowbreak{}bow\_))(\allowbreak{}nth(\allowbreak{}2, bhp\_), nth(\allowbreak{}2, bhq\_))])
\end{quote}
characterized by \texttt{hsum-vadd-acc}  \textit{(proven)}:
\begin{quote}\ttfamily\footnotesize\raggedright\noindent
forall(\allowbreak{}[v, w], is-complex-inner-product-space(\allowbreak{}v) implies is-complex-inner-product-space(\allowbreak{}w) implies vadd(\allowbreak{}hsum(\allowbreak{}v, w)) == hsum-vadd(\allowbreak{}v, w))
\end{quote}
also: \texttt{hsum-vadd-type} \\
mentioned by 3 result(s)

\smallskip

\noindent\textbf{\texttt{hsum-vec(bov\_, bow\_)}}\quad\textit{functoid}\index{hsum-vec(bov_, bow_)@\texttt{hsum-vec(bov\_, bow\_)}}\\
definition:
\begin{quote}\ttfamily\footnotesize\raggedright\noindent
cartesian(\allowbreak{}vec(\allowbreak{}bov\_), vec(\allowbreak{}bow\_))
\end{quote}

\smallskip

\noindent\textbf{\texttt{hsum-vneg(bov\_, bow\_)}}\quad\textit{functoid}\index{hsum-vneg(bov_, bow_)@\texttt{hsum-vneg(bov\_, bow\_)}}\\
definition:
\begin{quote}\ttfamily\footnotesize\raggedright\noindent
vnb-lambda(\allowbreak{}bhp\_, cartesian(\allowbreak{}vec(\allowbreak{}bov\_), vec(\allowbreak{}bow\_)), [(\allowbreak{}vneg(\allowbreak{}bov\_))(\allowbreak{}nth(\allowbreak{}1, bhp\_)), (\allowbreak{}vneg(\allowbreak{}bow\_))(\allowbreak{}nth(\allowbreak{}2, bhp\_))])
\end{quote}
characterized by \texttt{hsum-vneg-acc}  \textit{(proven)}:
\begin{quote}\ttfamily\footnotesize\raggedright\noindent
forall(\allowbreak{}[v, w], is-complex-inner-product-space(\allowbreak{}v) implies is-complex-inner-product-space(\allowbreak{}w) implies vneg(\allowbreak{}hsum(\allowbreak{}v, w)) == hsum-vneg(\allowbreak{}v, w))
\end{quote}
also: \texttt{hsum-vneg-type} \\
mentioned by 3 result(s)

\smallskip

\noindent\textbf{\texttt{hyps-in(mchs\_, mcs\_)}}\quad\textit{predicate}\index{hyps-in(mchs_, mcs_)@\texttt{hyps-in(mchs\_, mcs\_)}}\\
reads: every entry of the list \$1 is in \$2 \\
defined by \texttt{hyps-in} \\
characterized by \texttt{hyps-in}  \textit{(definitional)}:
\begin{quote}\ttfamily\footnotesize\raggedright\noindent
forall(\allowbreak{}[mchs\_, mcs\_], hyps-in(\allowbreak{}mchs\_, mcs\_) iff mchs\_ in set and forall(\allowbreak{}[mci\_ in nn], 1 <= mci\_ implies mci\_ <= length(\allowbreak{}mchs\_) implies nth(\allowbreak{}mci\_, mchs\_) in mcs\_))
\end{quote}
also: \texttt{hyps-in-0} \texttt{hyps-in-1} \texttt{hyps-in-2} \texttt{hyps-in-all} \texttt{hyps-in-set} \texttt{hyps-in-1-elim} \texttt{hyps-in-2-elim} \texttt{hyps-in-3-elim} \\
mentioned by 198 result(s)

\smallskip


\section*{I}
\addcontentsline{toc}{section}{I}

\noindent\textbf{\texttt{id-fun(a\_)}}\quad\textit{functoid}\index{id-fun(a_)@\texttt{id-fun(a\_)}}\\
reads: the identity map of \$1 \\
definition:
\begin{quote}\ttfamily\footnotesize\raggedright\noindent
vnb-lambda(\allowbreak{}idfx\_, a\_, idfx\_)
\end{quote}
characterized by \texttt{id-fun-type}  \textit{(proven)}:
\begin{quote}\ttfamily\footnotesize\raggedright\noindent
forall(\allowbreak{}[a in set], id-fun(\allowbreak{}a) in fun(\allowbreak{}a, a))
\end{quote}
also: \texttt{id-fun-apply} \\
mentioned by 88 result(s)

\smallskip

\noindent\textbf{\texttt{id-of(ssn\_)}}\quad\textit{functoid}\index{id-of(ssn_)@\texttt{id-of(ssn\_)}}\\
reads: the strand identifier of the node \$1 \\
definition:
\begin{quote}\ttfamily\footnotesize\raggedright\noindent
nth(\allowbreak{}2, ssn\_)
\end{quote}
mentioned by 1 result(s)

\smallskip

\noindent\textbf{\texttt{id-rep(amv\_)}}\quad\textit{functoid}\index{id-rep(amv_)@\texttt{id-rep(amv\_)}}\\
reads: the identity representation of the operators on \$1 \\
definition:
\begin{quote}\ttfamily\footnotesize\raggedright\noindent
vnb-lambda(\allowbreak{}amt\_, bounded-ops(\allowbreak{}amv\_, amv\_), amt\_)
\end{quote}
characterized by \texttt{id-rep-type}  \textit{(proven)}:
\begin{quote}\ttfamily\footnotesize\raggedright\noindent
forall(\allowbreak{}[v], is-complex-inner-product-space(\allowbreak{}v) implies id-rep(\allowbreak{}v) in fun(\allowbreak{}bounded-ops(\allowbreak{}v, v), bounded-ops(\allowbreak{}v, v)))
\end{quote}
also: \texttt{id-rep-apply} \texttt{id-rep-star-rep} \\
mentioned by 10 result(s)

\smallskip

\noindent\textbf{\texttt{ideal-stage}}\quad\textit{functoid}\index{ideal-stage@\texttt{ideal-stage}}\\
reads: the stage \$5 of the \$3-encryptions of \$4 over the tokens \$1 and the keys \$2 \\
defined by \texttt{ideal-stage-zero} \texttt{ideal-stage-succ} \\
characterized by \texttt{ideal-stage-succ}  \textit{(definitional)}:
\begin{quote}\ttfamily\footnotesize\raggedright\noindent
forall(\allowbreak{}[sbrt\_, sbrk\_, sbrks\_, sbrs\_, sbrn\_ in nn], ideal-stage(\allowbreak{}sbrt\_, sbrk\_, sbrks\_, sbrs\_, succ(\allowbreak{}sbrn\_)) == \{sbrm\_ in msg-set(\allowbreak{}sbrt\_, sbrk\_): forsome(\allowbreak{}[sbrg\_ in key-ideal-gen(\allowbreak{}sbrt\_, sbrk\_, empty-set, ideal-stage(\allowbreak{}sbrt\_, sbrk\_, sbrks\_, sbrs\_, sbrn\_)), sbry\_ in sbrks\_], sbrm\_ = encr(\allowbreak{}sbrg\_, sbry\_))\})
\end{quote}
also: \texttt{ideal-stage-zero} \texttt{ideal-stage-in-gen} \texttt{ideal-stage-in-power} \\
mentioned by 8 result(s)

\smallskip

\noindent\textbf{\texttt{iden(s)}}\quad\textit{accessor}\index{iden(s)@\texttt{iden(s)}}\\
reads: the identity element of \$1 \\
a structure slot \\
characterized by \texttt{cc-nag@iden}  \textit{(definitional)}:
\begin{quote}\ttfamily\footnotesize\raggedright\noindent
iden(\allowbreak{}cc-nag) == 0
\end{quote}
also: \texttt{lft-group@iden} \texttt{nn-add-monoid@iden} \texttt{zz-add-monoid@iden} \texttt{rmm-iden-uba-as-ring} \texttt{raag-iden-uba-as-ring} \texttt{rr-pos-star-add-monoid@iden} \texttt{rmm-iden-uba-as-ring-uiba-as-uba} \texttt{raag-iden-uba-as-ring-uiba-as-uba} ... (50 in all) \\
mentioned by 351 result(s)

\smallskip

\noindent\textbf{\texttt{identmat(a, n)}}\quad\textit{functoid}\index{identmat(a, n)@\texttt{identmat(a, n)}}\\
definition:
\begin{quote}\ttfamily\footnotesize\raggedright\noindent
matof(\allowbreak{}n, n, vnb-lambda(\allowbreak{}[i, j], cartesian(\allowbreak{}interval(\allowbreak{}1, n), interval(\allowbreak{}1, n)), if(\allowbreak{}i = j, one(\allowbreak{}a), zero(\allowbreak{}a))))
\end{quote}
characterized by \texttt{identmat-type-uba-as-ring}  \textit{(definitional)}:
\begin{quote}\ttfamily\footnotesize\raggedright\noindent
forall(\allowbreak{}[r], is-unital-banach-algebra(\allowbreak{}r) implies forall(\allowbreak{}[n in nn], identmat(\allowbreak{}uba-as-ring(\allowbreak{}r), n) in mat(\allowbreak{}n, n, vec(\allowbreak{}r))))
\end{quote}
also: \texttt{matact-identmat-cna-as-module} \texttt{matact-identmat-sip-as-module} \texttt{matact-identmat-uba-as-module} \texttt{matact-identmat-cnvs-as-module} \texttt{identmat-invertible-uba-as-ring} \texttt{identmat-left-identity-uba-as-ring} \texttt{identmat-right-identity-uba-as-ring} \texttt{identmat-type-uba-as-ring-uiba-as-uba} ... (24 in all) \\
mentioned by 270 result(s)

\smallskip

\noindent\textbf{\texttt{idl(s)}}\quad\textit{accessor}\index{idl(s)@\texttt{idl(s)}}\\
reads: the distinguished ideal of \$1 \\
a structure slot \\
characterized by \texttt{fp-idl}  \textit{(proven)}:
\begin{quote}\ttfamily\footnotesize\raggedright\noindent
forall(\allowbreak{}[p], zz-prime(\allowbreak{}p) implies idl(\allowbreak{}fp-ringoid(\allowbreak{}p)) == \{fpn\_ in zz: zz-divides(\allowbreak{}p, fpn\_)\})
\end{quote}
also: \texttt{rs-idl} \texttt{qp-csr-idl} \\
mentioned by 20 result(s)

\smallskip

\noindent\textbf{\texttt{ids(s)}}\quad\textit{accessor}\index{ids(s)@\texttt{ids(s)}}\\
reads: the ids slot of \$1 \\
a structure slot \\
characterized by \texttt{cx-sp@ids}  \textit{(definitional)}:
\begin{quote}\ttfamily\footnotesize\raggedright\noindent
ids(\allowbreak{}cx-sp) == interval(\allowbreak{}1, 8)
\end{quote}
mentioned by 56 result(s)

\smallskip

\noindent\textbf{\texttt{idx-list(min\_)}}\quad\textit{functoid}\index{idx-list(min_)@\texttt{idx-list(min\_)}}\\
reads: the list of the indices \$1 - 1, .., 0 \\
definition:
\begin{quote}\ttfamily\footnotesize\raggedright\noindent
iota(\allowbreak{}mir\_, mir\_ in tuples(\allowbreak{}set) and length(\allowbreak{}mir\_) = min\_ and forall(\allowbreak{}[mii\_ in nn], 1 <= mii\_ implies mii\_ <= min\_ implies nth(\allowbreak{}mii\_, mir\_) = ebvar(\allowbreak{}min\_ - mii\_)))
\end{quote}
characterized by \texttt{idx-list-prop}  \textit{(proven)}:
\begin{quote}\ttfamily\footnotesize\raggedright\noindent
forall(\allowbreak{}[n in nn], idx-list(\allowbreak{}n) in tuples(\allowbreak{}expr-set) and length(\allowbreak{}idx-list(\allowbreak{}n)) = n and forall(\allowbreak{}[i in nn], 1 <= i implies i <= n implies nth(\allowbreak{}i, idx-list(\allowbreak{}n)) = ebvar(\allowbreak{}n - i)))
\end{quote}
mentioned by 5 result(s)

\smallskip

\noindent\textbf{\texttt{ie}}\quad\textit{tactic}\index{ie@\texttt{ie}}\\
\texttt{(ie hyp k)} -- Intersection elim: extract the k-th branch of an INTERSECTION-membership assumption. \\
declared in \texttt{tactics-help.scm}

\smallskip

\noindent\textbf{\texttt{if}}\quad\textit{operator}\index{if@\texttt{if}}\\
a term-forming head \\
characterized by \texttt{matunit-col-shift-uba-as-ring}  \textit{(definitional)}:
\begin{quote}\ttfamily\footnotesize\raggedright\noindent
forall(\allowbreak{}[r], is-unital-banach-algebra(\allowbreak{}r) implies forall(\allowbreak{}[m, n, p, k, l], p in mat(\allowbreak{}m, n, vec(\allowbreak{}r)) implies k in interval(\allowbreak{}1, n) implies l in interval(\allowbreak{}1, n) implies forall(\allowbreak{}[i in interval(\allowbreak{}1, m), c in interval(\allowbreak{}1, n)], entry(\allowbreak{}matmul(\allowbreak{}uba-as-ring(\allowbreak{}r), p, matunit(\allowbreak{}uba-as-ring(\allowbreak{}r), n, k, l)), i, c) = if(\allowbreak{}c = l, entry(\allowbreak{}p, i, k), vzero(\allowbreak{}r)))))
\end{quote}
also: \texttt{matunit-col-shift-uba-as-ring-uiba-as-uba} \texttt{cd7-unif-ind} \texttt{difference-value} \texttt{matunit-col-shift} \texttt{ccint-rr-star-iso-le-iff} \\
mentioned by 356 result(s)

\smallskip

\noindent\textbf{\texttt{if-false}}\quad\textit{tactic}\index{if-false@\texttt{if-false}}\\
\texttt{(if-false t)} -- Reduce an (IF p a b) term on the not-p branch: spawns (NOT p); the continuation gains (= (IF p a b) b). \\
declared in \texttt{tactics-help.scm}

\smallskip

\noindent\textbf{\texttt{if-true}}\quad\textit{tactic}\index{if-true@\texttt{if-true}}\\
\texttt{(if-true t)} -- Reduce an (IF p a b) term on the p branch: spawns p as a subgoal; the continuation gains (= (IF p a b) a) when a is certified defined, else (== (IF p a b) a) (since 2026-10-08). \\
declared in \texttt{tactics-help.scm}

\smallskip

\noindent\textbf{\texttt{ii}}\quad\textit{tactic}\index{ii@\texttt{ii}}\\
\texttt{(ii)} -- Intersection intro: prove (IN x (INTERSECTION ...)) for each branch. \\
declared in \texttt{tactics-help.scm}

\smallskip

\noindent\textbf{\texttt{imag-part}}\quad\textit{operator}\index{imag-part@\texttt{imag-part}}\\
a term-forming head \\
characterized by \texttt{imag-part-def}  \textit{(definitional)}:
\begin{quote}\ttfamily\footnotesize\raggedright\noindent
forall(\allowbreak{}[z], imag-part(\allowbreak{}z) == (\allowbreak{}z - conjugate(\allowbreak{}z)) * recip(\allowbreak{}2 * 1i))
\end{quote}
also: \texttt{imag-part-in-rr} \texttt{rr-imag-part-zero} \\
mentioned by 237 result(s)

\smallskip

\noindent\textbf{\texttt{image}}\quad\textit{operator}\index{image@\texttt{image}}\\
a term-forming head \\
characterized by \texttt{image-membership-iff}  \textit{(definitional)}:
\begin{quote}\ttfamily\footnotesize\raggedright\noindent
forall(\allowbreak{}[phi, s, w], w in image(\allowbreak{}phi, s) iff w in set and forsome(\allowbreak{}[x in s], phi(\allowbreak{}x) = w))
\end{quote}
also: \texttt{image-set} \texttt{compact-image} \texttt{fun-image-set} \texttt{card-seg-image} \texttt{image-restrict} \texttt{countable-image} \texttt{image-cyl-power} \texttt{image-cyl-split} ... (50 in all) \\
mentioned by 231 result(s)

\smallskip

\noindent\textbf{\texttt{improper-int(ipf\_)}}\quad\textit{functoid}\index{improper-int(ipf_)@\texttt{improper-int(ipf\_)}}\\
reads: the improper integral of \$1 over (0, +inf) \\
definition:
\begin{quote}\ttfamily\footnotesize\raggedright\noindent
iota(\allowbreak{}ipv\_, improper-int-converges-to(\allowbreak{}ipf\_, ipv\_))
\end{quote}
characterized by \texttt{improper-int-value}  \textit{(proven)}:
\begin{quote}\ttfamily\footnotesize\raggedright\noindent
forall(\allowbreak{}[f], forall(\allowbreak{}[ipr\_ in rr, ipq\_ in rr], 0 < ipr\_ implies ipr\_ < ipq\_ implies cc-int(\allowbreak{}f, ipr\_, ipq\_) in cc) implies forall(\allowbreak{}[gml\_], improper-int-converges-to(\allowbreak{}f, gml\_) implies improper-int(\allowbreak{}f) = gml\_))
\end{quote}
mentioned by 2 result(s)

\smallskip

\noindent\textbf{\texttt{improper-int-converges-to(ipf\_, ipl\_)}}\quad\textit{predicate}\index{improper-int-converges-to(ipf_, ipl_)@\texttt{improper-int-converges-to(ipf\_, ipl\_)}}\\
reads: the improper integral of \$1 over (0, +inf) converges to \$2 \\
defined by \texttt{improper-int-converges-to} \\
characterized by \texttt{improper-int-converges-to}  \textit{(definitional)}:
\begin{quote}\ttfamily\footnotesize\raggedright\noindent
forall(\allowbreak{}[ipf\_, ipl\_], improper-int-converges-to(\allowbreak{}ipf\_, ipl\_) iff ipl\_ in cc and forall(\allowbreak{}[ipe\_], pos-rr(\allowbreak{}ipe\_) implies forsome(\allowbreak{}[ipd\_], pos-rr(\allowbreak{}ipd\_) and forsome(\allowbreak{}[ipm\_ in rr], forall(\allowbreak{}[ipr\_ in rr, ipq\_ in rr], 0 < ipr\_ implies ipr\_ < ipd\_ implies ipm\_ < ipq\_ implies ipr\_ < ipq\_ implies magnitude(\allowbreak{}cc-int(\allowbreak{}ipf\_, ipr\_, ipq\_) - ipl\_) <= ipe\_)))))
\end{quote}
mentioned by 9 result(s)

\smallskip

\noindent\textbf{\texttt{in}}\quad\textit{primitive}\index{in@\texttt{in}}\\
reads: \$1 is in \$2 \\
a kernel relation \\
characterized by \texttt{hyps-in}  \textit{(definitional)}:
\begin{quote}\ttfamily\footnotesize\raggedright\noindent
forall(\allowbreak{}[mchs\_, mcs\_], hyps-in(\allowbreak{}mchs\_, mcs\_) iff mchs\_ in set and forall(\allowbreak{}[mci\_ in nn], 1 <= mci\_ implies mci\_ <= length(\allowbreak{}mchs\_) implies nth(\allowbreak{}mci\_, mchs\_) in mcs\_))
\end{quote}
also: \texttt{is-ring} \texttt{is-chain} \texttt{is-lusin} \texttt{inst-step} \texttt{def-binder} \texttt{is-fin-ext} \texttt{is-or-init} \texttt{is-ringoid} ... (4243 in all) \\
mentioned by 22565 result(s)

\smallskip

\noindent\textbf{\texttt{index-of(ssn\_)}}\quad\textit{functoid}\index{index-of(ssn_)@\texttt{index-of(ssn\_)}}\\
reads: the index of the node \$1 \\
definition:
\begin{quote}\ttfamily\footnotesize\raggedright\noindent
nth(\allowbreak{}1, ssn\_)
\end{quote}
mentioned by 1 result(s)

\smallskip

\noindent\textbf{\texttt{indicator(omega, a)}}\quad\textit{functoid}\index{indicator(omega, a)@\texttt{indicator(omega, a)}}\\
reads: the indicator function of \$2 in \$1 \\
definition:
\begin{quote}\ttfamily\footnotesize\raggedright\noindent
vnb-lambda(\allowbreak{}x\_, omega, if(\allowbreak{}x\_ in a, 1, 0))
\end{quote}
characterized by \texttt{indicator-l1}  \textit{(proven)}:
\begin{quote}\ttfamily\footnotesize\raggedright\noindent
forall(\allowbreak{}[omega, ca, mu, p], is-measure(\allowbreak{}omega, ca, mu) implies p in ca implies mu(\allowbreak{}p) in rr implies indicator(\allowbreak{}omega, p) in l1-set(\allowbreak{}omega, ca, mu) and integral(\allowbreak{}omega, ca, mu, abs-fn(\allowbreak{}omega, indicator(\allowbreak{}omega, p))) = mu(\allowbreak{}p) and rintegral(\allowbreak{}omega, ca, mu, indicator(\allowbreak{}omega, p)) = mu(\allowbreak{}p))
\end{quote}
also: \texttt{xsi-indicator} \texttt{indicator-simple} \texttt{rmc-indicator-fun} \texttt{indicator-value-in} \texttt{integral-indicator} \texttt{rmc-indicator-cont} \texttt{indicator-value-out} \texttt{rmc-indicator-values} ... (14 in all) \\
mentioned by 83 result(s)

\smallskip

\noindent\textbf{\texttt{ineq}}\quad\textit{tactic}\index{ineq@\texttt{ineq}}\\
\texttt{(ineq i1 i2 ...)} -- Close a linear-inequality goal over RR (<= < > >= = between RR terms) as a consequence of the named assumptions (1-based indices), via the Fourier-Motzkin/Farkas oracle.  Linearizes over + - * and the binplus/binneg/bintimes aliases; every MAXIMAL non-arithmetic subterm is an atom that must be certified in RR.  (Does NOT see through a generic ring's (ADD s)/(MUL s) -- those become opaque atoms.) \\
declared in \texttt{tactics-help.scm}

\smallskip

\noindent\textbf{\texttt{inf(infs\_)}}\quad\textit{functoid}\index{inf(infs_)@\texttt{inf(infs\_)}}\\
reads: the infimum of \$1 \\
definition:
\begin{quote}\ttfamily\footnotesize\raggedright\noindent
-sup(\allowbreak{}\{infv\_ in rr: forsome(\allowbreak{}[infx\_ in infs\_], infv\_ = -infx\_)\})
\end{quote}
characterized by \texttt{mink-inf-neg}  \textit{(proven)}:
\begin{quote}\ttfamily\footnotesize\raggedright\noindent
forall(\allowbreak{}[m, cs, x], is-module(\allowbreak{}m) implies scal(\allowbreak{}m) = normed-field-as-commutative-ring(\allowbreak{}rr-normed-field) implies is-convex-subset(\allowbreak{}m, cs) implies is-symmetric-subset(\allowbreak{}m, cs) implies is-absorbing-subset(\allowbreak{}m, cs) implies x in vec(\allowbreak{}m) implies inf(\allowbreak{}mink-set(\allowbreak{}m, cs, (\allowbreak{}act(\allowbreak{}m))(\allowbreak{}-1, x))) = inf(\allowbreak{}mink-set(\allowbreak{}m, cs, x)))
\end{quote}
also: \texttt{sxt-inf-le-q} \texttt{mink-inf-zero} \texttt{rr-inf-is-glb} \texttt{homog-inf-le-p} \texttt{homog-inf-zero} \texttt{mink-inf-scale} \texttt{sxt-inf-nonneg} \texttt{homog-inf-scale} ... (19 in all) \\
mentioned by 42 result(s)

\smallskip

\noindent\textbf{\texttt{inf-subsets(a)}}\quad\textit{functoid}\index{inf-subsets(a)@\texttt{inf-subsets(a)}}\\
definition:
\begin{quote}\ttfamily\footnotesize\raggedright\noindent
\{s in power(\allowbreak{}a): not(\allowbreak{}card(\allowbreak{}s) in nn)\}
\end{quote}
characterized by \texttt{inf-subsets-membership}  \textit{(definitional)}:
\begin{quote}\ttfamily\footnotesize\raggedright\noindent
forall(\allowbreak{}[a, s], s in inf-subsets(\allowbreak{}a) iff s subset a and not(\allowbreak{}card(\allowbreak{}s) in nn))
\end{quote}
also: \texttt{nn-in-inf-subsets} \texttt{inf-subsets-is-set} \texttt{inf-subsets-unfold} \\
mentioned by 19 result(s)

\smallskip

\noindent\textbf{\texttt{infiltrated-space}}\quad\textit{structure}\index{infiltrated-space@\texttt{infiltrated-space}}\\
predicate \texttt{is-infiltrated-space}; slots: ispace adv-ids adv-tokens

\smallskip

\noindent\textbf{\texttt{injection}}\quad\textit{operator}\index{injection@\texttt{injection}}\\
a term-forming head \\
characterized by \texttt{injection-membership-iff}  \textit{(definitional)}:
\begin{quote}\ttfamily\footnotesize\raggedright\noindent
forall(\allowbreak{}[x, y, phi], phi in injection(\allowbreak{}x, y) iff phi in fun(\allowbreak{}x, y) and forall(\allowbreak{}[a in x, b in x], phi(\allowbreak{}a) = phi(\allowbreak{}b) implies a = b))
\end{quote}
also: \texttt{injection-in-fun} \texttt{lzd-id-injection} \texttt{injection-set-iff} \texttt{rb5-ext-injection} \texttt{lzs-swap-injection} \texttt{injection-injective} \texttt{lz3-const-injection} \texttt{card-image-injection} ... (21 in all) \\
mentioned by 57 result(s)

\smallskip

\noindent\textbf{\texttt{injective-star(f)}}\quad\textit{predicate}\index{injective-star(f)@\texttt{injective-star(f)}}\\
reads: \$1 is injective \\
defined by \texttt{injective-star} \\
characterized by \texttt{injective-star}  \textit{(definitional)}:
\begin{quote}\ttfamily\footnotesize\raggedright\noindent
forall(\allowbreak{}[f], injective-star(\allowbreak{}f) iff forall(\allowbreak{}[u\_, v\_], f(\allowbreak{}u\_) = f(\allowbreak{}v\_) implies u\_ = v\_))
\end{quote}
also: \texttt{injection-is-injective-star} \\
mentioned by 3 result(s)

\smallskip

\noindent\textbf{\texttt{insert-last(phi, x, n)}}\quad\textit{functoid}\index{insert-last(phi, x, n)@\texttt{insert-last(phi, x, n)}}\\
definition:
\begin{quote}\ttfamily\footnotesize\raggedright\noindent
vnb-lambda(\allowbreak{}i, nn, if(\allowbreak{}i in ord-segment(\allowbreak{}n), phi(\allowbreak{}i), x))
\end{quote}

\smallskip

\noindent\textbf{\texttt{inst}}\quad\textit{tactic}\index{inst@\texttt{inst}}\\
\texttt{(inst forall-hyp term)} -- Instantiate a universally-quantified assumption at `term', adding the instance to context. \\
declared in \texttt{tactics-help.scm}

\smallskip

\noindent\textbf{\texttt{inst+}}\quad\textit{tactic}\index{inst+@\texttt{inst+}}\\
\texttt{(inst+ forall-hyp term)} -- Instantiate an in-context universal at `term', then forward-detach any guards whose antecedents are in context. \\
declared in \texttt{tactics-help.scm}

\smallskip

\noindent\textbf{\texttt{inst-closed(md12l\_, md12x\_)}}\quad\textit{predicate}\index{inst-closed(md12l_, md12x_)@\texttt{inst-closed(md12l\_, md12x\_)}}\\
reads: \$2 is closed under instantiations at the arguments of \$1 \\
defined by \texttt{inst-closed} \\
characterized by \texttt{inst-closed}  \textit{(definitional)}:
\begin{quote}\ttfamily\footnotesize\raggedright\noindent
forall(\allowbreak{}[md12l\_, md12x\_], inst-closed(\allowbreak{}md12l\_, md12x\_) iff forall(\allowbreak{}[md12p\_, md12q\_], md12p\_ in md12x\_ implies md12q\_ in set implies inst-step(\allowbreak{}md12l\_, md12p\_, md12q\_) implies md12q\_ in md12x\_))
\end{quote}
mentioned by 5 result(s)

\smallskip

\noindent\textbf{\texttt{inst-reach(md12d\_, md12l\_)}}\quad\textit{functoid}\index{inst-reach(md12d_, md12l_)@\texttt{inst-reach(md12d\_, md12l\_)}}\\
reads: the instances of \$1 at the arguments of \$2 \\
definition:
\begin{quote}\ttfamily\footnotesize\raggedright\noindent
\{md12z\_ in expr-set: forall(\allowbreak{}[md12w\_ in power(\allowbreak{}expr-set)], md12d\_ in md12w\_ implies inst-closed(\allowbreak{}md12l\_, md12w\_) implies md12z\_ in md12w\_)\}
\end{quote}
characterized by \texttt{adm-inst-reach}  \textit{(proven)}:
\begin{quote}\ttfamily\footnotesize\raggedright\noindent
forall(\allowbreak{}[tt, g, l, d, x], [g, d] in derivable(\allowbreak{}tt) implies forall(\allowbreak{}[i in nn], 2 <= i implies i <= length(\allowbreak{}l) implies [g, enc-eq(\allowbreak{}nth(\allowbreak{}i, l), nth(\allowbreak{}i, l))] in derivable(\allowbreak{}tt)) implies x in inst-reach(\allowbreak{}d, l) implies [g, x] in derivable(\allowbreak{}tt))
\end{quote}
also: \texttt{inst-reach-least} \texttt{inst-reach-unfold} \\
mentioned by 11 result(s)

\smallskip

\noindent\textbf{\texttt{inst-step(md12l\_, md12f\_, md12g\_)}}\quad\textit{predicate}\index{inst-step(md12l_, md12f_, md12g_)@\texttt{inst-step(md12l\_, md12f\_, md12g\_)}}\\
reads: \$3 is \$2 instantiated at an argument of \$1 \\
defined by \texttt{inst-step} \\
characterized by \texttt{inst-step}  \textit{(definitional)}:
\begin{quote}\ttfamily\footnotesize\raggedright\noindent
forall(\allowbreak{}[md12l\_, md12f\_, md12g\_], inst-step(\allowbreak{}md12l\_, md12f\_, md12g\_) iff forsome(\allowbreak{}[md12b\_, md12i\_], md12f\_ = enc-all(\allowbreak{}md12b\_) and md12i\_ in nn and 2 <= md12i\_ and md12i\_ <= length(\allowbreak{}md12l\_) and md12g\_ = expr-open(\allowbreak{}md12b\_, 0, nth(\allowbreak{}md12i\_, md12l\_))))
\end{quote}
mentioned by 4 result(s)

\smallskip

\noindent\textbf{\texttt{integral(omega, ca, mu, f)}}\quad\textit{functoid}\index{integral(omega, ca, mu, f)@\texttt{integral(omega, ca, mu, f)}}\\
reads: the integral of \$4 over \$1 with respect to the measure \$3 \\
definition:
\begin{quote}\ttfamily\footnotesize\raggedright\noindent
esup(\allowbreak{}\{iny\_ in rr-pos-star: forsome(\allowbreak{}[ins\_ in fun(\allowbreak{}omega, rr-pos-star)], is-simple-fn(\allowbreak{}omega, ca, ins\_) and ptwise-le(\allowbreak{}omega, ins\_, f) and iny\_ = simple-integral(\allowbreak{}omega, mu, ins\_))\})
\end{quote}
characterized by \texttt{integral-in}  \textit{(proven)}:
\begin{quote}\ttfamily\footnotesize\raggedright\noindent
forall(\allowbreak{}[omega, ca, mu, f], is-measure(\allowbreak{}omega, ca, mu) implies is-measurable-fn(\allowbreak{}omega, ca, f) implies integral(\allowbreak{}omega, ca, mu, f) in rr-pos-star)
\end{quote}
also: \texttt{integral-sub} \texttt{integral-esum} \texttt{integral-ae-eq} \texttt{integral-abs-fn} \texttt{integral-finsum} \texttt{integral-markov} \texttt{integral-simple} \texttt{integral-unfold} ... (44 in all) \\
mentioned by 148 result(s)

\smallskip

\noindent\textbf{\texttt{integral-domain}}\quad\textit{refinement}\index{integral-domain@\texttt{integral-domain}}\\
predicate \texttt{is-integral-domain}; same shape as commutative-ring

\smallskip

\noindent\textbf{\texttt{interior(s, a)}}\quad\textit{functoid}\index{interior(s, a)@\texttt{interior(s, a)}}\\
definition:
\begin{quote}\ttfamily\footnotesize\raggedright\noindent
\{x in pts(\allowbreak{}s): forsome(\allowbreak{}[u in power(\allowbreak{}pts(\allowbreak{}s))], is-open(\allowbreak{}s, u) and x in u and u subset a)\}
\end{quote}
characterized by \texttt{interior-in-set}  \textit{(proven)}:
\begin{quote}\ttfamily\footnotesize\raggedright\noindent
forall(\allowbreak{}[s], is-metric-space(\allowbreak{}s) implies forall(\allowbreak{}[a], interior(\allowbreak{}s, a) subset a))
\end{quote}
also: \texttt{interior-is-set} \texttt{interior-is-open} \texttt{interior-mem-intro} \texttt{cc-exhaust-interior} \texttt{interior-angle-sign} \texttt{interior-cross-sign} \texttt{interior-in-carrier} \texttt{interior-membership} ... (12 in all) \\
mentioned by 34 result(s)

\smallskip

\noindent\textbf{\texttt{intersection}}\quad\textit{operator}\index{intersection@\texttt{intersection}}\\
a term-forming head \\
characterized by \texttt{submodule-intersection-cna-as-module}  \textit{(definitional)}:
\begin{quote}\ttfamily\footnotesize\raggedright\noindent
forall(\allowbreak{}[r], is-complex-normed-algebra(\allowbreak{}r) implies forall(\allowbreak{}[s1], is-submodule(\allowbreak{}cna-as-module(\allowbreak{}r), s1) implies forall(\allowbreak{}[s2], is-submodule(\allowbreak{}cna-as-module(\allowbreak{}r), s2) implies is-submodule(\allowbreak{}cna-as-module(\allowbreak{}r), intersection(\allowbreak{}s1, s2)))))
\end{quote}
also: \texttt{submodule-intersection-sip-as-module} \texttt{submodule-intersection-uba-as-module} \texttt{submodule-intersection-cnvs-as-module} \texttt{submodule-intersection-sip-as-module-csip-as-sip} \texttt{submodule-intersection-uba-as-module-uiba-as-uba} \texttt{intersection-decompose} \texttt{intersection-membership} \texttt{intersection-set-closure} ... (15 in all) \\
mentioned by 196 result(s)

\smallskip

\noindent\textbf{\texttt{intersection-of(c)}}\quad\textit{functoid}\index{intersection-of(c)@\texttt{intersection-of(c)}}\\
reads: the intersection of the members of \$1 \\
definition:
\begin{quote}\ttfamily\footnotesize\raggedright\noindent
\{iov\_ in big-union(\allowbreak{}iow\_, c, iow\_): forall(\allowbreak{}[iou\_ in c], iov\_ in iou\_)\}
\end{quote}
characterized by \texttt{intersection-of-insert}  \textit{(proven)}:
\begin{quote}\ttfamily\footnotesize\raggedright\noindent
forall(\allowbreak{}[tf, x0 in set], forsome(\allowbreak{}[w], w in tf) implies intersection-of(\allowbreak{}union(\allowbreak{}tf, pair(\allowbreak{}x0, x0))) = intersection(\allowbreak{}intersection-of(\allowbreak{}tf), x0))
\end{quote}
also: \texttt{intersection-of-is-set} \texttt{intersection-of-unfold} \texttt{intersection-of-antitone} \texttt{intersection-of-membership} \texttt{intersection-of-empty-family} \texttt{intersection-of-insert-empty} \texttt{complement-in-intersection-of} \texttt{intersection-of-subset-member} \\
mentioned by 26 result(s)

\smallskip

\noindent\textbf{\texttt{interval(a, b)}}\quad\textit{functoid}\index{interval(a, b)@\texttt{interval(a, b)}}\\
reads: the natural numbers from \$1 to \$2 inclusive \\
definition:
\begin{quote}\ttfamily\footnotesize\raggedright\noindent
\{i in nn: a <= i and i <= b\}
\end{quote}
characterized by \texttt{interval-hi}  \textit{(proven)}:
\begin{quote}\ttfamily\footnotesize\raggedright\noindent
forall(\allowbreak{}[a, b, i in interval(\allowbreak{}a, b)], i <= b)
\end{quote}
also: \texttt{interval-lo} \texttt{interval-1-1} \texttt{interval-card} \texttt{interval-widen} \texttt{interval-in-set} \texttt{interval-unfold} \texttt{one-in-interval} \texttt{pred-in-interval} ... (23 in all) \\
mentioned by 1805 result(s)

\smallskip

\noindent\textbf{\texttt{interval-length(s)}}\quad\textit{functoid}\index{interval-length(s)@\texttt{interval-length(s)}}\\
reads: the length (diameter) of the set of reals \$1 \\
definition:
\begin{quote}\ttfamily\footnotesize\raggedright\noindent
esup(\allowbreak{}\{odv\_ in rr-pos-star: odv\_ = 0 or forsome(\allowbreak{}[odx\_ in s, ody\_ in s], odx\_ <= ody\_ and odv\_ = ody\_ - odx\_)\})
\end{quote}
characterized by \texttt{interval-length-mono}  \textit{(proven)}:
\begin{quote}\ttfamily\footnotesize\raggedright\noindent
forall(\allowbreak{}[aa, bb], aa subset bb implies interval-length(\allowbreak{}aa) <= interval-length(\allowbreak{}bb))
\end{quote}
also: \texttt{interval-length-type} \texttt{interval-length-ccint} \texttt{interval-length-widen} \texttt{interval-length-extend} \texttt{interval-length-unfold} \texttt{interval-length-translate-le} \\
mentioned by 58 result(s)

\smallskip

\noindent\textbf{\texttt{inv(s)}}\quad\textit{accessor}\index{inv(s)@\texttt{inv(s)}}\\
reads: the inverse operation of \$1 \\
a structure slot \\
characterized by \texttt{cc-nag@inv}  \textit{(definitional)}:
\begin{quote}\ttfamily\footnotesize\raggedright\noindent
inv(\allowbreak{}cc-nag) == vnb-lambda(\allowbreak{}x\_, cc, -x\_)
\end{quote}
also: \texttt{lft-group@inv} \texttt{ras-inv-uba-as-ring} \texttt{frm-mvag-inv-cna-as-module} \texttt{frm-mvag-inv-sip-as-module} \texttt{frm-mvag-inv-uba-as-module} \texttt{frm-mvag-inv-cnvs-as-module} \texttt{ras-inv-uba-as-ring-uiba-as-uba} \texttt{frm-mvag-inv-sip-as-module-csip-as-sip} ... (54 in all) \\
mentioned by 224 result(s)

\smallskip

\noindent\textbf{\texttt{inverse-bij(phi, x, y)}}\quad\textit{functoid}\index{inverse-bij(phi, x, y)@\texttt{inverse-bij(phi, x, y)}}\\
definition:
\begin{quote}\ttfamily\footnotesize\raggedright\noindent
vnb-lambda(\allowbreak{}y\_, y, choice(\allowbreak{}\{x\_ in x: phi(\allowbreak{}x\_) = y\_\}))
\end{quote}
characterized by \texttt{lft-inverse-bij}  \textit{(proven)}:
\begin{quote}\ttfamily\footnotesize\raggedright\noindent
forall(\allowbreak{}[a in cc, b in cc, c in cc, d in cc], not(\allowbreak{}a * d - b * c = 0) implies inverse-bij(\allowbreak{}lft(\allowbreak{}a, b, c, d), cc-hat, cc-hat) = lft(\allowbreak{}d, -b, -c, a))
\end{quote}
also: \texttt{inverse-bij-left} \texttt{inverse-bij-right} \texttt{inverse-bij-in-fun} \texttt{inverse-bij-is-bijection} \\
mentioned by 57 result(s)

\smallskip

\noindent\textbf{\texttt{inverse-el(baa\_, bax\_)}}\quad\textit{functoid}\index{inverse-el(baa_, bax_)@\texttt{inverse-el(baa\_, bax\_)}}\\
reads: the inverse of \$2 in \$1 \\
definition:
\begin{quote}\ttfamily\footnotesize\raggedright\noindent
iota(\allowbreak{}bau\_, bau\_ in vec(\allowbreak{}baa\_) and (\allowbreak{}vmul(\allowbreak{}baa\_))(\allowbreak{}bau\_, bax\_) = vone(\allowbreak{}baa\_) and (\allowbreak{}vmul(\allowbreak{}baa\_))(\allowbreak{}bax\_, bau\_) = vone(\allowbreak{}baa\_))
\end{quote}
mentioned by 54 result(s)

\smallskip

\noindent\textbf{\texttt{invertibles(baa\_)}}\quad\textit{functoid}\index{invertibles(baa_)@\texttt{invertibles(baa\_)}}\\
reads: the invertible elements of \$1 \\
definition:
\begin{quote}\ttfamily\footnotesize\raggedright\noindent
\{bav\_ in vec(\allowbreak{}baa\_): is-invertible-el(\allowbreak{}baa\_, bav\_)\}
\end{quote}
characterized by \texttt{invertibles-open}  \textit{(proven)}:
\begin{quote}\ttfamily\footnotesize\raggedright\noindent
forall(\allowbreak{}[a], is-unital-banach-algebra(\allowbreak{}a) implies is-open(\allowbreak{}ba-metric-space(\allowbreak{}a), invertibles(\allowbreak{}a)))
\end{quote}
mentioned by 1 result(s)

\smallskip

\noindent\textbf{\texttt{iota}}\quad\textit{operator}\index{iota@\texttt{iota}}\\
a term-forming head \\
characterized by \texttt{pz-rel-iota}  \textit{(proven)}:
\begin{quote}\ttfamily\footnotesize\raggedright\noindent
forall(\allowbreak{}[p], zz-prime(\allowbreak{}p) implies forall(\allowbreak{}[x in qq], padic-rel(\allowbreak{}p, x, iota(\allowbreak{}cpv\_, cpv\_ in rr and padic-rel(\allowbreak{}p, x, cpv\_)))))
\end{quote}
also: \texttt{pc-char-iota} \texttt{r7q-dhat-iota} \texttt{r7q-iota-dhat} \texttt{limit-form-iota} \texttt{series-limit-iota} \\
mentioned by 43 result(s)

\smallskip

\noindent\textbf{\texttt{iota-confined(mnbb\_)}}\quad\textit{predicate}\index{iota-confined(mnbb_)@\texttt{iota-confined(mnbb\_)}}\\
reads: the description with body \$1 is confined to sets \\
defined by \texttt{iota-confined} \\
characterized by \texttt{iota-confined}  \textit{(definitional)}:
\begin{quote}\ttfamily\footnotesize\raggedright\noindent
forall(\allowbreak{}[mnbb\_], iota-confined(\allowbreak{}mnbb\_) iff forsome(\allowbreak{}[mnbqc\_ in expr-set], enc-in(\allowbreak{}ebvar(\allowbreak{}0), mnbqc\_) in conj-reach(\allowbreak{}mnbb\_)))
\end{quote}
mentioned by 4 result(s)

\smallskip

\noindent\textbf{\texttt{iota-d}}\quad\textit{tactic}\index{iota-d@\texttt{iota-d}}\\
\texttt{(iota-d term)} -- Definite-description: discharge the IOTA uniqueness obligation for `term'. \\
declared in \texttt{tactics-help.scm}

\smallskip

\noindent\textbf{\texttt{ip(s)}}\quad\textit{accessor}\index{ip(s)@\texttt{ip(s)}}\\
reads: the inner product of \$1 \\
a structure slot \\
characterized by \texttt{cips-ip-type}  \textit{(definitional)}:
\begin{quote}\ttfamily\footnotesize\raggedright\noindent
forall(\allowbreak{}[v], is-complex-inner-product-space(\allowbreak{}v) implies forall(\allowbreak{}[x\_ in vec(\allowbreak{}v), y\_ in vec(\allowbreak{}v)], (\allowbreak{}ip(\allowbreak{}v))(\allowbreak{}x\_, y\_) in cc))
\end{quote}
also: \texttt{cips-ip-add-left} \texttt{cips-ip-conj-sym} \texttt{cips-ip-definite} \texttt{cips-ip-self-real} \texttt{cips-ip-homog-left} \texttt{cips-ip-self-nonneg} \texttt{cips-ip-zero-vector} \texttt{gns-pre-ip} ... (85 in all) \\
mentioned by 284 result(s)

\smallskip

\noindent\textbf{\texttt{ip-metric-space(v)}}\quad\textit{functoid}\index{ip-metric-space(v)@\texttt{ip-metric-space(v)}}\\
reads: the metric space of the inner product space \$1 \\
definition:
\begin{quote}\ttfamily\footnotesize\raggedright\noindent
nag-metric-space(\allowbreak{}ip-normed-ag(\allowbreak{}v))
\end{quote}
characterized by \texttt{ip-metric-space-pts}  \textit{(proven)}:
\begin{quote}\ttfamily\footnotesize\raggedright\noindent
forall(\allowbreak{}[v], is-complex-inner-product-space(\allowbreak{}v) implies pts(\allowbreak{}ip-metric-space(\allowbreak{}v)) == vec(\allowbreak{}v))
\end{quote}
also: \texttt{ip-metric-space-dist} \texttt{ip-metric-space-unfold} \texttt{ip-metric-space-is-metric-space} \\
mentioned by 139 result(s)

\smallskip

\noindent\textbf{\texttt{ip-norm(v, x)}}\quad\textit{functoid}\index{ip-norm(v, x)@\texttt{ip-norm(v, x)}}\\
reads: the induced norm of \$2 in \$1 \\
definition:
\begin{quote}\ttfamily\footnotesize\raggedright\noindent
sqrt(\allowbreak{}(\allowbreak{}ip(\allowbreak{}v))(\allowbreak{}x, x))
\end{quote}
characterized by \texttt{ip-normed-ag-norm-value}  \textit{(definitional)}:
\begin{quote}\ttfamily\footnotesize\raggedright\noindent
forall(\allowbreak{}[v, x in vec(\allowbreak{}v)], (\allowbreak{}nrm(\allowbreak{}ip-normed-ag(\allowbreak{}v)))(\allowbreak{}x) == ip-norm(\allowbreak{}v, x))
\end{quote}
also: \texttt{ip-norm-sq} \texttt{ip-norm-act} \texttt{ip-norm-vneg} \texttt{ip-norm-zero} \texttt{ip-norm-in-rr} \texttt{ip-norm-nonneg} \texttt{ip-norm-unfold} \texttt{ip-norm-diff-sym} ... (9 in all) \\
mentioned by 113 result(s)

\smallskip

\noindent\textbf{\texttt{ip-normed-ag(v)}}\quad\textit{functoid}\index{ip-normed-ag(v)@\texttt{ip-normed-ag(v)}}\\
reads: the additive group of \$1 under the induced norm \\
definition:
\begin{quote}\ttfamily\footnotesize\raggedright\noindent
[vec(\allowbreak{}v), vadd(\allowbreak{}v), vzero(\allowbreak{}v), vneg(\allowbreak{}v), vnb-lambda(\allowbreak{}x\_, vec(\allowbreak{}v), ip-norm(\allowbreak{}v, x\_))]
\end{quote}
characterized by \texttt{ip-normed-ag-carrier}  \textit{(definitional)}:
\begin{quote}\ttfamily\footnotesize\raggedright\noindent
forall(\allowbreak{}[v], is-complex-inner-product-space(\allowbreak{}v) implies carr(\allowbreak{}ip-normed-ag(\allowbreak{}v)) == vec(\allowbreak{}v))
\end{quote}
also: \texttt{ip-normed-ag-norm-value} \texttt{ip-normed-ag-inv} \texttt{ip-normed-ag-opr} \texttt{onb-ip-normed-ag-iden} \texttt{ip-normed-ag-is-normed-ag} \\
mentioned by 45 result(s)

\smallskip

\noindent\textbf{\texttt{is-abelian-group(s)}}\quad\textit{predicate}\index{is-abelian-group(s)@\texttt{is-abelian-group(s)}}\\
reads: an abelian group \\
a predicate \\
characterized by \texttt{is-abelian-group}  \textit{(definitional)}:
\begin{quote}\ttfamily\footnotesize\raggedright\noindent
forall(\allowbreak{}[s], is-abelian-group(\allowbreak{}s) iff length(\allowbreak{}s) = 4 and carr(\allowbreak{}s) in set and opr(\allowbreak{}s) in fun(\allowbreak{}cartesian(\allowbreak{}carr(\allowbreak{}s), carr(\allowbreak{}s)), carr(\allowbreak{}s)) and iden(\allowbreak{}s) in carr(\allowbreak{}s) and inv(\allowbreak{}s) in fun(\allowbreak{}carr(\allowbreak{}s), carr(\allowbreak{}s)) and is-associative(\allowbreak{}opr(\allowbreak{}s), carr(\allowbreak{}s)) and is-identity(\allowbreak{}opr(\allowbreak{}s), iden(\allowbreak{}s), carr(\allowbreak{}s)) and has-inverses(\allowbreak{}opr(\allowbreak{}s), iden(\allowbreak{}s), inv(\allowbreak{}s), carr(\allowbreak{}s)) and is-commutative(\allowbreak{}opr(\allowbreak{}s), carr(\allowbreak{}s)) and s in set)
\end{quote}
also: \texttt{module-vector-ag-is-abelian-group} \texttt{ring-additive-ag-is-abelian-group} \texttt{field-additive-ag-is-abelian-group} \texttt{normed-field-additive-ag-is-abelian-group} \texttt{normed-ag-as-abelian-group-is-abelian-group} \texttt{commutative-ring-additive-ag-is-abelian-group} \texttt{ring-additive-ag-is-abelian-group-uba-as-ring} \texttt{module-vector-ag-is-abelian-group-cna-as-module} ... (28 in all) \\
mentioned by 250 result(s)

\smallskip

\noindent\textbf{\texttt{is-absolutely-continuous(omega, ca, mu, nu)}}\quad\textit{predicate}\index{is-absolutely-continuous(omega, ca, mu, nu)@\texttt{is-absolutely-continuous(omega, ca, mu, nu)}}\\
reads: \$3 is absolutely continuous with respect to \$4 on \$2 \\
defined by \texttt{is-absolutely-continuous} \\
characterized by \texttt{is-absolutely-continuous}  \textit{(definitional)}:
\begin{quote}\ttfamily\footnotesize\raggedright\noindent
forall(\allowbreak{}[omega, ca, mu, nu], is-absolutely-continuous(\allowbreak{}omega, ca, mu, nu) iff forall(\allowbreak{}[ace\_ in ca], nu(\allowbreak{}ace\_) = 0 implies mu(\allowbreak{}ace\_) = 0))
\end{quote}
mentioned by 8 result(s)

\smallskip

\noindent\textbf{\texttt{is-absolutely-summable(grp, f)}}\quad\textit{predicate}\index{is-absolutely-summable(grp, f)@\texttt{is-absolutely-summable(grp, f)}}\\
reads: \$2 is absolutely summable in \$1 \\
defined by \texttt{is-absolutely-summable} \\
characterized by \texttt{is-absolutely-summable}  \textit{(definitional)}:
\begin{quote}\ttfamily\footnotesize\raggedright\noindent
forall(\allowbreak{}[grp, f], is-absolutely-summable(\allowbreak{}grp, f) iff esum(\allowbreak{}vnb-lambda(\allowbreak{}i, dom(\allowbreak{}f), (\allowbreak{}nrm(\allowbreak{}grp))(\allowbreak{}f(\allowbreak{}i)))) in rr)
\end{quote}
mentioned by 4 result(s)

\smallskip

\noindent\textbf{\texttt{is-absorbing-subset(m, cs)}}\quad\textit{predicate}\index{is-absorbing-subset(m, cs)@\texttt{is-absorbing-subset(m, cs)}}\\
reads: \$2 is an absorbing subset of \$1 \\
defined by \texttt{is-absorbing-subset} \\
characterized by \texttt{is-absorbing-subset}  \textit{(definitional)}:
\begin{quote}\ttfamily\footnotesize\raggedright\noindent
forall(\allowbreak{}[m, cs], is-absorbing-subset(\allowbreak{}m, cs) iff cs subset vec(\allowbreak{}m) and forall(\allowbreak{}[abx\_ in vec(\allowbreak{}m)], forsome(\allowbreak{}[abr\_ in rr], 0 < abr\_ and forsome(\allowbreak{}[aby\_ in cs], abx\_ = (\allowbreak{}act(\allowbreak{}m))(\allowbreak{}abr\_, aby\_)))))
\end{quote}
mentioned by 22 result(s)

\smallskip

\noindent\textbf{\texttt{is-accumulation-point(apts\_, apta\_, aptx\_)}}\quad\textit{predicate}\index{is-accumulation-point(apts_, apta_, aptx_)@\texttt{is-accumulation-point(apts\_, apta\_, aptx\_)}}\\
reads: \$3 is an accumulation point of \$2 in the metric space \$1 \\
defined by \texttt{is-accumulation-point} \\
characterized by \texttt{is-accumulation-point}  \textit{(definitional)}:
\begin{quote}\ttfamily\footnotesize\raggedright\noindent
forall(\allowbreak{}[apts\_, apta\_, aptx\_], is-accumulation-point(\allowbreak{}apts\_, apta\_, aptx\_) iff is-metric-space(\allowbreak{}apts\_) and aptx\_ in pts(\allowbreak{}apts\_) and forall(\allowbreak{}[aptr\_], pos-rr(\allowbreak{}aptr\_) implies forsome(\allowbreak{}[apty\_ in apta\_], apty\_ in ball(\allowbreak{}apts\_, aptx\_, aptr\_) and not(\allowbreak{}apty\_ = aptx\_))))
\end{quote}
mentioned by 6 result(s)

\smallskip

\noindent\textbf{\texttt{is-add-character(chk\_, chc\_)}}\quad\textit{predicate}\index{is-add-character(chk_, chc_)@\texttt{is-add-character(chk\_, chc\_)}}\\
reads: \$2 is a character of the additive group of \$1 \\
defined by \texttt{is-add-character} \\
characterized by \texttt{is-add-character}  \textit{(definitional)}:
\begin{quote}\ttfamily\footnotesize\raggedright\noindent
forall(\allowbreak{}[chk\_, chc\_], is-add-character(\allowbreak{}chk\_, chc\_) iff chc\_ in fun(\allowbreak{}carr(\allowbreak{}chk\_), unit-circle) and forall(\allowbreak{}[chx\_ in carr(\allowbreak{}chk\_), chy\_ in carr(\allowbreak{}chk\_)], chc\_(\allowbreak{}(\allowbreak{}add(\allowbreak{}chk\_))(\allowbreak{}chx\_, chy\_)) = chc\_(\allowbreak{}chx\_) * chc\_(\allowbreak{}chy\_)) and is-continuous(\allowbreak{}nf-metric-space(\allowbreak{}chk\_), cc-ms, chc\_))
\end{quote}
mentioned by 59 result(s)

\smallskip

\noindent\textbf{\texttt{is-add-subgroup(asgk\_, asgu\_)}}\quad\textit{predicate}\index{is-add-subgroup(asgk_, asgu_)@\texttt{is-add-subgroup(asgk\_, asgu\_)}}\\
reads: \$2 is an additive subgroup of \$1 \\
defined by \texttt{is-add-subgroup} \\
characterized by \texttt{is-add-subgroup}  \textit{(definitional)}:
\begin{quote}\ttfamily\footnotesize\raggedright\noindent
forall(\allowbreak{}[asgk\_, asgu\_], is-add-subgroup(\allowbreak{}asgk\_, asgu\_) iff forall(\allowbreak{}[asx\_ in asgu\_], asx\_ in carr(\allowbreak{}asgk\_)) and zero(\allowbreak{}asgk\_) in asgu\_ and forall(\allowbreak{}[asx\_ in asgu\_, asy\_ in asgu\_], (\allowbreak{}add(\allowbreak{}asgk\_))(\allowbreak{}asx\_, asy\_) in asgu\_) and forall(\allowbreak{}[asx\_ in asgu\_], (\allowbreak{}neg(\allowbreak{}asgk\_))(\allowbreak{}asx\_) in asgu\_))
\end{quote}
mentioned by 7 result(s)

\smallskip

\noindent\textbf{\texttt{is-adversary-node(sai\_, san\_)}}\quad\textit{predicate}\index{is-adversary-node(sai_, san_)@\texttt{is-adversary-node(sai\_, san\_)}}\\
reads: \$2 is an adversary node of \$1 \\
defined by \texttt{is-adversary-node} \\
characterized by \texttt{is-adversary-node}  \textit{(definitional)}:
\begin{quote}\ttfamily\footnotesize\raggedright\noindent
forall(\allowbreak{}[sai\_, san\_], is-adversary-node(\allowbreak{}sai\_, san\_) iff nth(\allowbreak{}2, san\_) in adv-ids(\allowbreak{}sai\_))
\end{quote}
mentioned by 16 result(s)

\smallskip

\noindent\textbf{\texttt{is-adversary-trace(sasp\_, sap\_, sas\_)}}\quad\textit{predicate}\index{is-adversary-trace(sasp_, sap_, sas_)@\texttt{is-adversary-trace(sasp\_, sap\_, sas\_)}}\\
reads: \$3 is an adversary trace for the adversary knowing \$2 in the strand space \$1 \\
defined by \texttt{is-adversary-trace} \\
characterized by \texttt{is-adversary-trace}  \textit{(definitional)}:
\begin{quote}\ttfamily\footnotesize\raggedright\noindent
forall(\allowbreak{}[sasp\_, sap\_, sas\_], is-adversary-trace(\allowbreak{}sasp\_, sap\_, sas\_) iff is-text-trace(\allowbreak{}sasp\_, sap\_, sas\_) or is-flush-trace(\allowbreak{}sasp\_, sas\_) or is-tee-trace(\allowbreak{}sasp\_, sas\_) or is-concat-trace(\allowbreak{}sasp\_, sas\_) or is-separate-trace(\allowbreak{}sasp\_, sas\_) or is-encrypt-trace(\allowbreak{}sasp\_, sas\_) or is-decrypt-trace(\allowbreak{}sasp\_, sas\_))
\end{quote}
mentioned by 4 result(s)

\smallskip

\noindent\textbf{\texttt{is-affine-isometry-group(aim\_, aia\_, aig\_)}}\quad\textit{predicate}\index{is-affine-isometry-group(aim_, aia_, aig_)@\texttt{is-affine-isometry-group(aim\_, aia\_, aig\_)}}\\
reads: \$3 is a group of affine isometries of \$2 in the normed space \$1 \\
defined by \texttt{is-affine-isometry-group} \\
characterized by \texttt{is-affine-isometry-group}  \textit{(definitional)}:
\begin{quote}\ttfamily\footnotesize\raggedright\noindent
forall(\allowbreak{}[aim\_, aia\_, aig\_], is-affine-isometry-group(\allowbreak{}aim\_, aia\_, aig\_) iff aig\_ subset fun(\allowbreak{}aia\_, aia\_) and forall(\allowbreak{}[aif\_ in aig\_, aiu\_ in aia\_, aiv\_ in aia\_, ait\_ in rr], 0 <= ait\_ implies ait\_ <= 1 implies aif\_(\allowbreak{}(\allowbreak{}vadd(\allowbreak{}aim\_))(\allowbreak{}(\allowbreak{}act(\allowbreak{}aim\_))(\allowbreak{}1 - ait\_, aiu\_), (\allowbreak{}act(\allowbreak{}aim\_))(\allowbreak{}ait\_, aiv\_))) = (\allowbreak{}vadd(\allowbreak{}aim\_))(\allowbreak{}(\allowbreak{}act(\allowbreak{}aim\_))(\allowbreak{}1 - ait\_, aif\_(\allowbreak{}aiu\_)), (\allowbreak{}act(\allowbreak{}aim\_))(\allowbreak{}ait\_, aif\_(\allowbreak{}aiv\_)))) and forall(\allowbreak{}[aif\_ in aig\_, aiu\_ in aia\_, aiv\_ in aia\_], (\allowbreak{}dist(\allowbreak{}nvs-metric-space(\allowbreak{}aim\_)))(\allowbreak{}aif\_(\allowbreak{}aiu\_), aif\_(\allowbreak{}aiv\_)) = (\allowbreak{}dist(\allowbreak{}nvs-metric-space(\allowbreak{}aim\_)))(\allowbreak{}aiu\_, aiv\_)) and forall(\allowbreak{}[aif\_ in aig\_], forsome(\allowbreak{}[aih\_ in aig\_], forall(\allowbreak{}[aix\_ in aia\_], aif\_(\allowbreak{}aih\_(\allowbreak{}aix\_)) = aix\_))))
\end{quote}
mentioned by 5 result(s)

\smallskip

\noindent\textbf{\texttt{is-alg-hom(baa\_, bab\_, baf\_)}}\quad\textit{predicate}\index{is-alg-hom(baa_, bab_, baf_)@\texttt{is-alg-hom(baa\_, bab\_, baf\_)}}\\
reads: \$3 is a unital algebra homomorphism from \$1 to \$2 \\
defined by \texttt{is-alg-hom} \\
characterized by \texttt{is-alg-hom}  \textit{(definitional)}:
\begin{quote}\ttfamily\footnotesize\raggedright\noindent
forall(\allowbreak{}[baa\_, bab\_, baf\_], is-alg-hom(\allowbreak{}baa\_, bab\_, baf\_) iff baf\_ in fun(\allowbreak{}vec(\allowbreak{}baa\_), vec(\allowbreak{}bab\_)) and forall(\allowbreak{}[bax\_ in vec(\allowbreak{}baa\_), bay\_ in vec(\allowbreak{}baa\_)], baf\_(\allowbreak{}(\allowbreak{}vadd(\allowbreak{}baa\_))(\allowbreak{}bax\_, bay\_)) = (\allowbreak{}vadd(\allowbreak{}bab\_))(\allowbreak{}baf\_(\allowbreak{}bax\_), baf\_(\allowbreak{}bay\_))) and forall(\allowbreak{}[baz\_ in cc, bax\_ in vec(\allowbreak{}baa\_)], baf\_(\allowbreak{}(\allowbreak{}act(\allowbreak{}baa\_))(\allowbreak{}baz\_, bax\_)) = (\allowbreak{}act(\allowbreak{}bab\_))(\allowbreak{}baz\_, baf\_(\allowbreak{}bax\_))) and forall(\allowbreak{}[bax\_ in vec(\allowbreak{}baa\_), bay\_ in vec(\allowbreak{}baa\_)], baf\_(\allowbreak{}(\allowbreak{}vmul(\allowbreak{}baa\_))(\allowbreak{}bax\_, bay\_)) = (\allowbreak{}vmul(\allowbreak{}bab\_))(\allowbreak{}baf\_(\allowbreak{}bax\_), baf\_(\allowbreak{}bay\_))) and baf\_(\allowbreak{}vone(\allowbreak{}baa\_)) = vone(\allowbreak{}bab\_))
\end{quote}
mentioned by 13 result(s)

\smallskip

\noindent\textbf{\texttt{is-alg-subspace(baa\_, bai\_)}}\quad\textit{predicate}\index{is-alg-subspace(baa_, bai_)@\texttt{is-alg-subspace(baa\_, bai\_)}}\\
reads: \$2 is a linear subspace of \$1 \\
defined by \texttt{is-alg-subspace} \\
characterized by \texttt{is-alg-subspace}  \textit{(definitional)}:
\begin{quote}\ttfamily\footnotesize\raggedright\noindent
forall(\allowbreak{}[baa\_, bai\_], is-alg-subspace(\allowbreak{}baa\_, bai\_) iff bai\_ subset vec(\allowbreak{}baa\_) and vzero(\allowbreak{}baa\_) in bai\_ and forall(\allowbreak{}[bax\_ in bai\_, bay\_ in bai\_], (\allowbreak{}vadd(\allowbreak{}baa\_))(\allowbreak{}bax\_, bay\_) in bai\_) and forall(\allowbreak{}[baz\_ in cc, bax\_ in bai\_], (\allowbreak{}act(\allowbreak{}baa\_))(\allowbreak{}baz\_, bax\_) in bai\_))
\end{quote}
mentioned by 6 result(s)

\smallskip

\noindent\textbf{\texttt{is-algebra-of-sets(omega, ca)}}\quad\textit{predicate}\index{is-algebra-of-sets(omega, ca)@\texttt{is-algebra-of-sets(omega, ca)}}\\
reads: \$2 is an algebra of subsets of \$1 \\
defined by \texttt{is-algebra-of-sets} \\
characterized by \texttt{is-algebra-of-sets}  \textit{(definitional)}:
\begin{quote}\ttfamily\footnotesize\raggedright\noindent
forall(\allowbreak{}[omega, ca], is-algebra-of-sets(\allowbreak{}omega, ca) iff omega in set and ca subset power(\allowbreak{}omega) and forsome(\allowbreak{}[a\_], a\_ in ca) and forall(\allowbreak{}[a\_ in ca], complement-in(\allowbreak{}omega, a\_) in ca) and forall(\allowbreak{}[a\_ in ca, b\_ in ca], union(\allowbreak{}a\_, b\_) in ca) and forall(\allowbreak{}[a\_ in ca, b\_ in ca], intersection(\allowbreak{}a\_, b\_) in ca))
\end{quote}
mentioned by 18 result(s)

\smallskip

\noindent\textbf{\texttt{is-algebraic-of-degree-le(kk, th, n)}}\quad\textit{predicate}\index{is-algebraic-of-degree-le(kk, th, n)@\texttt{is-algebraic-of-degree-le(kk, th, n)}}\\
reads: \$2 is algebraic of degree at most \$3 over \$1 \\
defined by \texttt{is-algebraic-of-degree-le} \\
characterized by \texttt{is-algebraic-of-degree-le}  \textit{(definitional)}:
\begin{quote}\ttfamily\footnotesize\raggedright\noindent
forall(\allowbreak{}[kk, th, n], is-algebraic-of-degree-le(\allowbreak{}kk, th, n) iff th in cc and forsome(\allowbreak{}[anan\_ in nn], anan\_ <= n and forsome(\allowbreak{}[anac\_ in fun(\allowbreak{}nn, kk)], forall(\allowbreak{}[anak\_ in nn], succ(\allowbreak{}anan\_) <= anak\_ implies anac\_(\allowbreak{}anak\_) = 0) and not(\allowbreak{}anac\_(\allowbreak{}anan\_) = 0) and cc-series-partial-sum(\allowbreak{}vnb-lambda(\allowbreak{}pfj\_, nn, anac\_(\allowbreak{}pfj\_) * th \textasciicircum{} pfj\_), succ(\allowbreak{}anan\_)) = 0)))
\end{quote}
mentioned by 5 result(s)

\smallskip

\noindent\textbf{\texttt{is-algebraic-over(aa, r)}}\quad\textit{predicate}\index{is-algebraic-over(aa, r)@\texttt{is-algebraic-over(aa, r)}}\\
reads: \$2 is algebraic over the ring \$1 \\
defined by \texttt{is-algebraic-over} \\
characterized by \texttt{is-algebraic-over}  \textit{(definitional)}:
\begin{quote}\ttfamily\footnotesize\raggedright\noindent
forall(\allowbreak{}[aa, r], is-algebraic-over(\allowbreak{}aa, r) iff r in cc and forsome(\allowbreak{}[anan\_ in nn, anac\_ in fun(\allowbreak{}nn, aa)], forall(\allowbreak{}[anak\_ in nn], succ(\allowbreak{}anan\_) <= anak\_ implies anac\_(\allowbreak{}anak\_) = 0) and anac\_(\allowbreak{}anan\_) = 1 and cc-series-partial-sum(\allowbreak{}vnb-lambda(\allowbreak{}pfj\_, nn, anac\_(\allowbreak{}pfj\_) * r \textasciicircum{} pfj\_), succ(\allowbreak{}anan\_)) = 0))
\end{quote}
mentioned by 20 result(s)

\smallskip

\noindent\textbf{\texttt{is-algebraic-over-field(kk, th)}}\quad\textit{predicate}\index{is-algebraic-over-field(kk, th)@\texttt{is-algebraic-over-field(kk, th)}}\\
reads: \$2 is algebraic over the field \$1 \\
defined by \texttt{is-algebraic-over-field} \\
characterized by \texttt{is-algebraic-over-field}  \textit{(definitional)}:
\begin{quote}\ttfamily\footnotesize\raggedright\noindent
forall(\allowbreak{}[kk, th], is-algebraic-over-field(\allowbreak{}kk, th) iff th in cc and forsome(\allowbreak{}[anan\_ in nn, anac\_ in fun(\allowbreak{}nn, kk)], forall(\allowbreak{}[anak\_ in nn], succ(\allowbreak{}anan\_) <= anak\_ implies anac\_(\allowbreak{}anak\_) = 0) and not(\allowbreak{}anac\_(\allowbreak{}anan\_) = 0) and cc-series-partial-sum(\allowbreak{}vnb-lambda(\allowbreak{}pfj\_, nn, anac\_(\allowbreak{}pfj\_) * th \textasciicircum{} pfj\_), succ(\allowbreak{}anan\_)) = 0))
\end{quote}
mentioned by 10 result(s)

\smallskip

\noindent\textbf{\texttt{is-amplification(amv\_, ams\_, amw\_, amf\_)}}\quad\textit{predicate}\index{is-amplification(amv_, ams_, amw_, amf_)@\texttt{is-amplification(amv\_, ams\_, amw\_, amf\_)}}\\
reads: \$4 is an amplification of the operators on \$1 on \$3 adapted to \$2 \\
defined by \texttt{is-amplification} \\
characterized by \texttt{is-amplification}  \textit{(definitional)}:
\begin{quote}\ttfamily\footnotesize\raggedright\noindent
forall(\allowbreak{}[amv\_, ams\_, amw\_, amf\_], is-amplification(\allowbreak{}amv\_, ams\_, amw\_, amf\_) iff is-hilbert-space(\allowbreak{}amw\_) and is-star-rep(\allowbreak{}amv\_, amw\_, amf\_) and is-intertwined(\allowbreak{}amv\_, amw\_, id-rep(\allowbreak{}amv\_), amf\_, amv\_, ams\_) and is-intertwined(\allowbreak{}amw\_, amv\_, amf\_, id-rep(\allowbreak{}amv\_), amv\_, ams\_) and is-intertwined(\allowbreak{}amw\_, amw\_, amf\_, amf\_, amv\_, ams\_))
\end{quote}
mentioned by 8 result(s)

\smallskip

\noindent\textbf{\texttt{is-analytic-at(anlf\_, anlu\_, anla\_)}}\quad\textit{predicate}\index{is-analytic-at(anlf_, anlu_, anla_)@\texttt{is-analytic-at(anlf\_, anlu\_, anla\_)}}\\
reads: \$1 is analytic at \$3 on \$2 \\
defined by \texttt{is-analytic-at} \\
characterized by \texttt{is-analytic-at}  \textit{(definitional)}:
\begin{quote}\ttfamily\footnotesize\raggedright\noindent
forall(\allowbreak{}[anlf\_, anlu\_, anla\_], is-analytic-at(\allowbreak{}anlf\_, anlu\_, anla\_) iff is-open(\allowbreak{}cc-ms, anlu\_) and anlf\_ in fun(\allowbreak{}anlu\_, cc) and anla\_ in anlu\_ and forsome(\allowbreak{}[anlcf\_ in fun(\allowbreak{}nn, cc), anlr\_ in rr], 0 < anlr\_ and anlr\_ < cps-radius(\allowbreak{}anlcf\_) and ball(\allowbreak{}cc-ms, anla\_, anlr\_) subset anlu\_ and forall(\allowbreak{}[anlz\_ in ball(\allowbreak{}cc-ms, anla\_, anlr\_)], anlf\_(\allowbreak{}anlz\_) = cc-series-limit(\allowbreak{}vnb-lambda(\allowbreak{}anlk\_, nn, anlcf\_(\allowbreak{}anlk\_) * (\allowbreak{}anlz\_ - anla\_) \textasciicircum{} anlk\_)))))
\end{quote}
mentioned by 10 result(s)

\smallskip

\noindent\textbf{\texttt{is-annihilating-poly(gan-md\_, gan-t\_, gan-v\_, gan-k\_)}}\quad\textit{predicate}\index{is-annihilating-poly(gan-md_, gan-t_, gan-v_, gan-k_)@\texttt{is-annihilating-poly(gan-md\_, gan-t\_, gan-v\_, gan-k\_)}}\\
reads: \$4 is the annihilating polynomial of \$2 at \$3 \\
defined by \texttt{is-annihilating-poly} \\
characterized by \texttt{is-annihilating-poly}  \textit{(definitional)}:
\begin{quote}\ttfamily\footnotesize\raggedright\noindent
forall(\allowbreak{}[gan-md\_, gan-t\_, gan-v\_, gan-k\_], is-annihilating-poly(\allowbreak{}gan-md\_, gan-t\_, gan-v\_, gan-k\_) iff is-monic(\allowbreak{}scal(\allowbreak{}gan-md\_), gan-k\_) and poly-op(\allowbreak{}gan-md\_, gan-t\_, gan-k\_, gan-v\_) = vzero(\allowbreak{}gan-md\_) and forall(\allowbreak{}[ganq\_ in carr(\allowbreak{}poly(\allowbreak{}scal(\allowbreak{}gan-md\_)))], not(\allowbreak{}ganq\_ = zero(\allowbreak{}poly(\allowbreak{}scal(\allowbreak{}gan-md\_)))) implies poly-op(\allowbreak{}gan-md\_, gan-t\_, ganq\_, gan-v\_) = vzero(\allowbreak{}gan-md\_) implies deg(\allowbreak{}scal(\allowbreak{}gan-md\_), gan-k\_) <= deg(\allowbreak{}scal(\allowbreak{}gan-md\_), ganq\_)))
\end{quote}
also: \texttt{is-annihilating-poly-eq} \\
mentioned by 13 result(s)

\smallskip

\noindent\textbf{\texttt{is-antiderivable(phi, a, b)}}\quad\textit{predicate}\index{is-antiderivable(phi, a, b)@\texttt{is-antiderivable(phi, a, b)}}\\
reads: \$1 is antiderivable on the interval [\$2, \$3] \\
defined by \texttt{is-antiderivable} \\
characterized by \texttt{is-antiderivable}  \textit{(definitional)}:
\begin{quote}\ttfamily\footnotesize\raggedright\noindent
forall(\allowbreak{}[phi, a, b], is-antiderivable(\allowbreak{}phi, a, b) iff forsome(\allowbreak{}[f\_], is-antiderivative(\allowbreak{}f\_, phi, a, b)))
\end{quote}
also: \texttt{poly-is-antiderivable} \texttt{zero-is-antiderivable} \texttt{continuous-is-antiderivable} \texttt{bernstein-basis-antiderivable} \\
mentioned by 55 result(s)

\smallskip

\noindent\textbf{\texttt{is-antiderivative(f, phi, a, b)}}\quad\textit{predicate}\index{is-antiderivative(f, phi, a, b)@\texttt{is-antiderivative(f, phi, a, b)}}\\
reads: \$1 is an antiderivative of \$2 on the interval [\$3, \$4] \\
defined by \texttt{is-antiderivative} \\
characterized by \texttt{is-antiderivative}  \textit{(definitional)}:
\begin{quote}\ttfamily\footnotesize\raggedright\noindent
forall(\allowbreak{}[f, phi, a, b], is-antiderivative(\allowbreak{}f, phi, a, b) iff f in fun(\allowbreak{}rr, rr) and phi in fun(\allowbreak{}rr, rr) and a in rr and b in rr and a < b and forall(\allowbreak{}[x\_ in ccint(\allowbreak{}a, b)], is-continuous-at(\allowbreak{}rr-ms, rr-ms, f, x\_)) and forall(\allowbreak{}[th\_], th\_ in rr and a < th\_ and th\_ < b implies is-diff-at(\allowbreak{}f, th\_, phi(\allowbreak{}th\_))))
\end{quote}
mentioned by 31 result(s)

\smallskip

\noindent\textbf{\texttt{is-antisymmetric-op(v, jo)}}\quad\textit{predicate}\index{is-antisymmetric-op(v, jo)@\texttt{is-antisymmetric-op(v, jo)}}\\
reads: \$2 is antisymmetric on \$1 \\
defined by \texttt{is-antisymmetric-op} \\
characterized by \texttt{is-antisymmetric-op}  \textit{(definitional)}:
\begin{quote}\ttfamily\footnotesize\raggedright\noindent
forall(\allowbreak{}[v, jo], is-antisymmetric-op(\allowbreak{}v, jo) iff forall(\allowbreak{}[x\_ in vec(\allowbreak{}v), y\_ in vec(\allowbreak{}v)], (\allowbreak{}rip(\allowbreak{}v))(\allowbreak{}jo(\allowbreak{}x\_), y\_) = -(\allowbreak{}rip(\allowbreak{}v))(\allowbreak{}x\_, jo(\allowbreak{}y\_))))
\end{quote}
mentioned by 14 result(s)

\smallskip

\noindent\textbf{\texttt{is-associative(op, crr)}}\quad\textit{predicate}\index{is-associative(op, crr)@\texttt{is-associative(op, crr)}}\\
reads: \$1 is associative on \$2 \\
defined by \texttt{is-associative} \\
characterized by \texttt{is-associative}  \textit{(definitional)}:
\begin{quote}\ttfamily\footnotesize\raggedright\noindent
forall(\allowbreak{}[op, crr], is-associative(\allowbreak{}op, crr) iff forall(\allowbreak{}[u in crr, v in crr, w in crr], op(\allowbreak{}op(\allowbreak{}u, v), w) = op(\allowbreak{}u, op(\allowbreak{}v, w))))
\end{quote}
mentioned by 66 result(s)

\smallskip

\noindent\textbf{\texttt{is-atom(ca, a)}}\quad\textit{predicate}\index{is-atom(ca, a)@\texttt{is-atom(ca, a)}}\\
reads: \$2 is an atom (a minimal non-empty element) of \$1 \\
defined by \texttt{is-atom} \\
characterized by \texttt{is-atom}  \textit{(definitional)}:
\begin{quote}\ttfamily\footnotesize\raggedright\noindent
forall(\allowbreak{}[ca, a], is-atom(\allowbreak{}ca, a) iff a in ca and not(\allowbreak{}a = empty-set) and forall(\allowbreak{}[mab\_ in ca], mab\_ subset a implies not(\allowbreak{}mab\_ = empty-set) implies mab\_ = a))
\end{quote}
mentioned by 6 result(s)

\smallskip

\noindent\textbf{\texttt{is-baire-continuous(s, f)}}\quad\textit{predicate}\index{is-baire-continuous(s, f)@\texttt{is-baire-continuous(s, f)}}\\
reads: \$2 is a continuous map from N\textasciicircum{}N to \$1 \\
defined by \texttt{is-baire-continuous} \\
characterized by \texttt{is-baire-continuous}  \textit{(definitional)}:
\begin{quote}\ttfamily\footnotesize\raggedright\noindent
forall(\allowbreak{}[s, f], is-baire-continuous(\allowbreak{}s, f) iff f in fun(\allowbreak{}fun(\allowbreak{}nn, nn), pts(\allowbreak{}s)) and forall(\allowbreak{}[bcz\_ in fun(\allowbreak{}nn, nn), bce\_], pos-rr(\allowbreak{}bce\_) implies forsome(\allowbreak{}[bcn\_ in nn], forall(\allowbreak{}[bcw\_ in fun(\allowbreak{}nn, nn)], restrict(\allowbreak{}bcw\_, ord-segment(\allowbreak{}bcn\_)) = restrict(\allowbreak{}bcz\_, ord-segment(\allowbreak{}bcn\_)) implies (\allowbreak{}dist(\allowbreak{}s))(\allowbreak{}f(\allowbreak{}bcz\_), f(\allowbreak{}bcw\_)) <= bce\_))))
\end{quote}
mentioned by 11 result(s)

\smallskip

\noindent\textbf{\texttt{is-bicharacter(chk\_, chw\_)}}\quad\textit{predicate}\index{is-bicharacter(chk_, chw_)@\texttt{is-bicharacter(chk\_, chw\_)}}\\
reads: \$2 is a bicharacter on \$1 \\
defined by \texttt{is-bicharacter} \\
characterized by \texttt{is-bicharacter}  \textit{(definitional)}:
\begin{quote}\ttfamily\footnotesize\raggedright\noindent
forall(\allowbreak{}[chk\_, chw\_], is-bicharacter(\allowbreak{}chk\_, chw\_) iff chw\_ in fun(\allowbreak{}cartesian(\allowbreak{}carr(\allowbreak{}chk\_), carr(\allowbreak{}chk\_)), unit-circle) and forall(\allowbreak{}[chy\_ in carr(\allowbreak{}chk\_)], is-add-character(\allowbreak{}chk\_, vnb-lambda(\allowbreak{}chx\_, carr(\allowbreak{}chk\_), chw\_(\allowbreak{}chx\_, chy\_)))) and forall(\allowbreak{}[chy\_ in carr(\allowbreak{}chk\_)], is-add-character(\allowbreak{}chk\_, vnb-lambda(\allowbreak{}chx\_, carr(\allowbreak{}chk\_), chw\_(\allowbreak{}chy\_, chx\_)))))
\end{quote}
mentioned by 12 result(s)

\smallskip

\noindent\textbf{\texttt{is-bilinear-bicharacter(chk\_, chw\_)}}\quad\textit{predicate}\index{is-bilinear-bicharacter(chk_, chw_)@\texttt{is-bilinear-bicharacter(chk\_, chw\_)}}\\
reads: \$2 is a bilinear bicharacter on \$1 \\
defined by \texttt{is-bilinear-bicharacter} \\
characterized by \texttt{is-bilinear-bicharacter}  \textit{(definitional)}:
\begin{quote}\ttfamily\footnotesize\raggedright\noindent
forall(\allowbreak{}[chk\_, chw\_], is-bilinear-bicharacter(\allowbreak{}chk\_, chw\_) iff is-bicharacter(\allowbreak{}chk\_, chw\_) and forall(\allowbreak{}[cha\_ in carr(\allowbreak{}chk\_), chx\_ in carr(\allowbreak{}chk\_), chy\_ in carr(\allowbreak{}chk\_)], chw\_(\allowbreak{}(\allowbreak{}mul(\allowbreak{}chk\_))(\allowbreak{}cha\_, chx\_), chy\_) = chw\_(\allowbreak{}chx\_, (\allowbreak{}mul(\allowbreak{}chk\_))(\allowbreak{}cha\_, chy\_))))
\end{quote}
mentioned by 41 result(s)

\smallskip

\noindent\textbf{\texttt{is-bk-sequence(bss\_, bsg\_, bsv\_)}}\quad\textit{predicate}\index{is-bk-sequence(bss_, bsg_, bsv_)@\texttt{is-bk-sequence(bss\_, bsg\_, bsv\_)}}\\
reads: \$3 is a sequence of symmetric open neighbourhoods of the identity of \$2 in \$1 with cubes nested \\
defined by \texttt{is-bk-sequence} \\
characterized by \texttt{is-bk-sequence}  \textit{(definitional)}:
\begin{quote}\ttfamily\footnotesize\raggedright\noindent
forall(\allowbreak{}[bss\_, bsg\_, bsv\_], is-bk-sequence(\allowbreak{}bss\_, bsg\_, bsv\_) iff bsv\_ in fun(\allowbreak{}nn, opens(\allowbreak{}bss\_)) and bsv\_(\allowbreak{}0) = pts(\allowbreak{}bss\_) and forall(\allowbreak{}[bsk\_ in nn], iden(\allowbreak{}bsg\_) in bsv\_(\allowbreak{}bsk\_) and forall(\allowbreak{}[bsp\_ in bsv\_(\allowbreak{}bsk\_)], (\allowbreak{}inv(\allowbreak{}bsg\_))(\allowbreak{}bsp\_) in bsv\_(\allowbreak{}bsk\_)) and forall(\allowbreak{}[bsp\_ in bsv\_(\allowbreak{}succ(\allowbreak{}bsk\_)), bsq\_ in bsv\_(\allowbreak{}succ(\allowbreak{}bsk\_)), bsr\_ in bsv\_(\allowbreak{}succ(\allowbreak{}bsk\_))], (\allowbreak{}opr(\allowbreak{}bsg\_))(\allowbreak{}(\allowbreak{}opr(\allowbreak{}bsg\_))(\allowbreak{}bsp\_, bsq\_), bsr\_) in bsv\_(\allowbreak{}bsk\_))))
\end{quote}
mentioned by 39 result(s)

\smallskip

\noindent\textbf{\texttt{is-bounded-l1-functional(omega, ca, mu, f, cb)}}\quad\textit{predicate}\index{is-bounded-l1-functional(omega, ca, mu, f, cb)@\texttt{is-bounded-l1-functional(omega, ca, mu, f, cb)}}\\
reads: \$4 is a linear functional on L\textasciicircum{}1(\$1, \$3) of norm at most \$5 \\
defined by \texttt{is-bounded-l1-functional} \\
characterized by \texttt{is-bounded-l1-functional}  \textit{(definitional)}:
\begin{quote}\ttfamily\footnotesize\raggedright\noindent
forall(\allowbreak{}[omega, ca, mu, f, cb], is-bounded-l1-functional(\allowbreak{}omega, ca, mu, f, cb) iff f in fun(\allowbreak{}l1-set(\allowbreak{}omega, ca, mu), rr) and cb in rr and 0 <= cb and forall(\allowbreak{}[bfx\_ in l1-set(\allowbreak{}omega, ca, mu), bfy\_ in l1-set(\allowbreak{}omega, ca, mu)], f(\allowbreak{}vnb-lambda(\allowbreak{}sx\_, omega, bfx\_(\allowbreak{}sx\_) + bfy\_(\allowbreak{}sx\_))) = f(\allowbreak{}bfx\_) + f(\allowbreak{}bfy\_)) and forall(\allowbreak{}[bfr\_ in rr, bfx\_ in l1-set(\allowbreak{}omega, ca, mu)], f(\allowbreak{}vnb-lambda(\allowbreak{}sx\_, omega, bfr\_ * bfx\_(\allowbreak{}sx\_))) = bfr\_ * f(\allowbreak{}bfx\_)) and forall(\allowbreak{}[bfx\_ in l1-set(\allowbreak{}omega, ca, mu)], abs(\allowbreak{}f(\allowbreak{}bfx\_)) <= cb * integral(\allowbreak{}omega, ca, mu, abs-fn(\allowbreak{}omega, bfx\_))))
\end{quote}
mentioned by 9 result(s)

\smallskip

\noindent\textbf{\texttt{is-bounded-linear-arrow(a, b, f)}}\quad\textit{predicate}\index{is-bounded-linear-arrow(a, b, f)@\texttt{is-bounded-linear-arrow(a, b, f)}}\\
reads: \$3 is a bounded linear map from \$1 to \$2 \\
a predicate \\
characterized by \texttt{is-bounded-linear-arrow-def}  \textit{(definitional)}:
\begin{quote}\ttfamily\footnotesize\raggedright\noindent
forall(\allowbreak{}[a, b, f], is-bounded-linear-arrow(\allowbreak{}a, b, f) iff is-normed-vector-space(\allowbreak{}a) and is-normed-vector-space(\allowbreak{}b) and f in fun(\allowbreak{}vec(\allowbreak{}a), vec(\allowbreak{}b)) and scal(\allowbreak{}a) = scal(\allowbreak{}b) and forall(\allowbreak{}[x1\_ in vec(\allowbreak{}a), x2\_ in vec(\allowbreak{}a)], f(\allowbreak{}(\allowbreak{}vadd(\allowbreak{}a))(\allowbreak{}x1\_, x2\_)) = (\allowbreak{}vadd(\allowbreak{}b))(\allowbreak{}f(\allowbreak{}x1\_), f(\allowbreak{}x2\_))) and forall(\allowbreak{}[x1\_ in carr(\allowbreak{}scal(\allowbreak{}a)), x2\_ in vec(\allowbreak{}a)], f(\allowbreak{}(\allowbreak{}act(\allowbreak{}a))(\allowbreak{}x1\_, x2\_)) = (\allowbreak{}act(\allowbreak{}b))(\allowbreak{}x1\_, f(\allowbreak{}x2\_))) and forsome(\allowbreak{}[blk\_ in rr], 0 <= blk\_ and forall(\allowbreak{}[x1\_ in vec(\allowbreak{}a)], (\allowbreak{}vnrm(\allowbreak{}b))(\allowbreak{}f(\allowbreak{}x1\_)) <= blk\_ * (\allowbreak{}vnrm(\allowbreak{}a))(\allowbreak{}x1\_))))
\end{quote}
also: \texttt{hom-normed-vector-space-is-bounded-linear-arrow} \\
mentioned by 8 result(s)

\smallskip

\noindent\textbf{\texttt{is-bounded-linear-functional(m, f)}}\quad\textit{predicate}\index{is-bounded-linear-functional(m, f)@\texttt{is-bounded-linear-functional(m, f)}}\\
reads: \$2 is a bounded linear functional on \$1 \\
defined by \texttt{is-bounded-linear-functional} \\
characterized by \texttt{is-bounded-linear-functional}  \textit{(definitional)}:
\begin{quote}\ttfamily\footnotesize\raggedright\noindent
forall(\allowbreak{}[m, f], is-bounded-linear-functional(\allowbreak{}m, f) iff is-linear-functional(\allowbreak{}m, f) and forsome(\allowbreak{}[c\_ in rr], 0 <= c\_ and forall(\allowbreak{}[x\_ in vec(\allowbreak{}m)], abs(\allowbreak{}f(\allowbreak{}x\_)) <= c\_ * (\allowbreak{}vnrm(\allowbreak{}m))(\allowbreak{}x\_))))
\end{quote}
mentioned by 34 result(s)

\smallskip

\noindent\textbf{\texttt{is-bounded-linear-functional-on(m, s, f)}}\quad\textit{predicate}\index{is-bounded-linear-functional-on(m, s, f)@\texttt{is-bounded-linear-functional-on(m, s, f)}}\\
reads: \$3 is a bounded linear functional on the subspace \$2 of \$1 \\
defined by \texttt{is-bounded-linear-functional-on} \\
characterized by \texttt{is-bounded-linear-functional-on}  \textit{(definitional)}:
\begin{quote}\ttfamily\footnotesize\raggedright\noindent
forall(\allowbreak{}[m, s, f], is-bounded-linear-functional-on(\allowbreak{}m, s, f) iff is-linear-functional-on(\allowbreak{}m, s, f) and forsome(\allowbreak{}[c\_ in rr], 0 <= c\_ and forall(\allowbreak{}[x\_ in s], abs(\allowbreak{}f(\allowbreak{}x\_)) <= c\_ * (\allowbreak{}vnrm(\allowbreak{}m))(\allowbreak{}x\_))))
\end{quote}
mentioned by 15 result(s)

\smallskip

\noindent\textbf{\texttt{is-bounded-metric-space(s)}}\quad\textit{predicate}\index{is-bounded-metric-space(s)@\texttt{is-bounded-metric-space(s)}}\\
reads: a bounded metric space \\
defined by \texttt{is-bounded-metric-space} \\
characterized by \texttt{is-bounded-metric-space}  \textit{(definitional)}:
\begin{quote}\ttfamily\footnotesize\raggedright\noindent
forall(\allowbreak{}[s], is-bounded-metric-space(\allowbreak{}s) iff is-metric-space(\allowbreak{}s) and forsome(\allowbreak{}[b in rr], forall(\allowbreak{}[x, y], x in pts(\allowbreak{}s) implies y in pts(\allowbreak{}s) implies (\allowbreak{}dist(\allowbreak{}s))(\allowbreak{}x, y) <= b)))
\end{quote}
also: \texttt{bdd-metric-is-bounded-metric-space} \\
mentioned by 8 result(s)

\smallskip

\noindent\textbf{\texttt{is-bounded-op(bov\_, bow\_, bot\_)}}\quad\textit{predicate}\index{is-bounded-op(bov_, bow_, bot_)@\texttt{is-bounded-op(bov\_, bow\_, bot\_)}}\\
reads: \$3 is a bounded linear operator from \$1 to \$2 \\
defined by \texttt{is-bounded-op} \\
characterized by \texttt{is-bounded-op}  \textit{(definitional)}:
\begin{quote}\ttfamily\footnotesize\raggedright\noindent
forall(\allowbreak{}[bov\_, bow\_, bot\_], is-bounded-op(\allowbreak{}bov\_, bow\_, bot\_) iff is-linear-op(\allowbreak{}bov\_, bow\_, bot\_) and forsome(\allowbreak{}[bom\_ in rr], 0 <= bom\_ and forall(\allowbreak{}[box\_ in vec(\allowbreak{}bov\_)], ip-norm(\allowbreak{}bow\_, bot\_(\allowbreak{}box\_)) <= bom\_ * ip-norm(\allowbreak{}bov\_, box\_))))
\end{quote}
mentioned by 6 result(s)

\smallskip

\noindent\textbf{\texttt{is-boundedly-equivalent(phi, psi)}}\quad\textit{predicate}\index{is-boundedly-equivalent(phi, psi)@\texttt{is-boundedly-equivalent(phi, psi)}}\\
reads: \$1 and \$2 are boundedly equivalent on [0, +oo) \\
defined by \texttt{is-boundedly-equivalent} \\
characterized by \texttt{is-boundedly-equivalent}  \textit{(definitional)}:
\begin{quote}\ttfamily\footnotesize\raggedright\noindent
forall(\allowbreak{}[phi, psi], is-boundedly-equivalent(\allowbreak{}phi, psi) iff phi in fun(\allowbreak{}nnreal, nnreal) and psi in fun(\allowbreak{}nnreal, nnreal) and forsome(\allowbreak{}[bel\_, beu\_], pos-rr(\allowbreak{}bel\_) and pos-rr(\allowbreak{}beu\_) and forall(\allowbreak{}[ber\_ in rr], 0 < ber\_ implies bel\_ * phi(\allowbreak{}ber\_) <= psi(\allowbreak{}ber\_) and psi(\allowbreak{}ber\_) <= beu\_ * phi(\allowbreak{}ber\_))))
\end{quote}
mentioned by 6 result(s)

\smallskip

\noindent\textbf{\texttt{is-bundle(s)}}\quad\textit{predicate}\index{is-bundle(s)@\texttt{is-bundle(s)}}\\
reads: a bundle \\
a predicate \\
characterized by \texttt{is-bundle}  \textit{(definitional)}:
\begin{quote}\ttfamily\footnotesize\raggedright\noindent
forall(\allowbreak{}[s], is-bundle(\allowbreak{}s) iff length(\allowbreak{}s) = 3 and is-strand-space(\allowbreak{}bspace(\allowbreak{}s)) and bnodes(\allowbreak{}s) in set and bsend(\allowbreak{}s) in power(\allowbreak{}cartesian(\allowbreak{}bnodes(\allowbreak{}s), bnodes(\allowbreak{}s))) and forall(\allowbreak{}[ssx\_ in bnodes(\allowbreak{}s)], ssx\_ in nodes(\allowbreak{}bspace(\allowbreak{}s))) and card(\allowbreak{}bnodes(\allowbreak{}s)) in nn and forall(\allowbreak{}[ssx\_, ssy\_], [ssx\_, ssy\_] in bsend(\allowbreak{}s) implies sends-to(\allowbreak{}bspace(\allowbreak{}s), ssx\_, ssy\_)) and forall(\allowbreak{}[ssy\_ in bnodes(\allowbreak{}s)], is-received(\allowbreak{}bspace(\allowbreak{}s), ssy\_) implies forsome(\allowbreak{}[ssx\_], [ssx\_, ssy\_] in bsend(\allowbreak{}s) and forall(\allowbreak{}[ssz\_], [ssz\_, ssy\_] in bsend(\allowbreak{}s) implies ssz\_ = ssx\_))) and forall(\allowbreak{}[ssy\_ in bnodes(\allowbreak{}s), ssx\_], same-strand-before(\allowbreak{}bspace(\allowbreak{}s), ssx\_, ssy\_) implies ssx\_ in bnodes(\allowbreak{}s)) and forall(\allowbreak{}[ssx\_], not(\allowbreak{}[ssx\_, ssx\_] in prec-strict(\allowbreak{}bspace(\allowbreak{}s), bnodes(\allowbreak{}s), bsend(\allowbreak{}s)))) and s in set)
\end{quote}
also: \texttt{cx-c-is-bundle} \texttt{is-bundle-is-set} \\
mentioned by 74 result(s)

\smallskip

\noindent\textbf{\texttt{is-c-continuous-arrow(a, b, f)}}\quad\textit{predicate}\index{is-c-continuous-arrow(a, b, f)@\texttt{is-c-continuous-arrow(a, b, f)}}\\
reads: \$3 is a continuous map of countably-metrised spaces from \$1 to \$2 \\
a predicate \\
characterized by \texttt{is-c-continuous-arrow-def}  \textit{(definitional)}:
\begin{quote}\ttfamily\footnotesize\raggedright\noindent
forall(\allowbreak{}[a, b, f], is-c-continuous-arrow(\allowbreak{}a, b, f) iff is-c-metric-space(\allowbreak{}a) and is-c-metric-space(\allowbreak{}b) and f in fun(\allowbreak{}pts(\allowbreak{}a), pts(\allowbreak{}b)) and is-continuous(\allowbreak{}c-metric(\allowbreak{}a), c-metric(\allowbreak{}b), f))
\end{quote}
mentioned by 7 result(s)

\smallskip

\noindent\textbf{\texttt{is-c-metric-space(s)}}\quad\textit{predicate}\index{is-c-metric-space(s)@\texttt{is-c-metric-space(s)}}\\
reads: a countably-metrised space \\
a predicate \\
characterized by \texttt{is-c-metric-space}  \textit{(definitional)}:
\begin{quote}\ttfamily\footnotesize\raggedright\noindent
forall(\allowbreak{}[s], is-c-metric-space(\allowbreak{}s) iff length(\allowbreak{}s) = 2 and pts(\allowbreak{}s) in set and dists(\allowbreak{}s) in fun(\allowbreak{}nn, fun(\allowbreak{}cartesian(\allowbreak{}pts(\allowbreak{}s), pts(\allowbreak{}s)), rr)) and forall(\allowbreak{}[k\_ in nn], is-pseudometric(\allowbreak{}(\allowbreak{}dists(\allowbreak{}s))(\allowbreak{}k\_), pts(\allowbreak{}s))) and forall(\allowbreak{}[u in pts(\allowbreak{}s), v in pts(\allowbreak{}s)], forall(\allowbreak{}[k\_ in nn], (\allowbreak{}(\allowbreak{}dists(\allowbreak{}s))(\allowbreak{}k\_))(\allowbreak{}u, v) = 0) implies u = v) and s in set)
\end{quote}
also: \texttt{is-c-metric-space-is-set} \texttt{frm-c-space-is-c-metric-space} \\
mentioned by 62 result(s)

\smallskip

\noindent\textbf{\texttt{is-caratheodory-measurable(omega, mu, es)}}\quad\textit{predicate}\index{is-caratheodory-measurable(omega, mu, es)@\texttt{is-caratheodory-measurable(omega, mu, es)}}\\
reads: \$3 is Caratheodory measurable for the outer measure \$2 on \$1 \\
defined by \texttt{is-caratheodory-measurable} \\
characterized by \texttt{is-caratheodory-measurable}  \textit{(definitional)}:
\begin{quote}\ttfamily\footnotesize\raggedright\noindent
forall(\allowbreak{}[omega, mu, es], is-caratheodory-measurable(\allowbreak{}omega, mu, es) iff es subset omega and forall(\allowbreak{}[a\_ in power(\allowbreak{}omega)], mu(\allowbreak{}a\_) = eplus(\allowbreak{}mu(\allowbreak{}intersection(\allowbreak{}a\_, es)), mu(\allowbreak{}difference(\allowbreak{}a\_, es)))))
\end{quote}
mentioned by 21 result(s)

\smallskip

\noindent\textbf{\texttt{is-cauchy-net(m, t, d, x)}}\quad\textit{predicate}\index{is-cauchy-net(m, t, d, x)@\texttt{is-cauchy-net(m, t, d, x)}}\\
reads: the net \$4 on \$3 is Cauchy in the topological vector space \$1, \$2 \\
defined by \texttt{is-cauchy-net} \\
characterized by \texttt{is-cauchy-net}  \textit{(definitional)}:
\begin{quote}\ttfamily\footnotesize\raggedright\noindent
forall(\allowbreak{}[m, t, d, x], is-cauchy-net(\allowbreak{}m, t, d, x) iff is-directed-set(\allowbreak{}d) and x in fun(\allowbreak{}pts(\allowbreak{}d), vec(\allowbreak{}m)) and forall(\allowbreak{}[cnu\_ in opens(\allowbreak{}t)], vzero(\allowbreak{}m) in cnu\_ implies forsome(\allowbreak{}[cni\_ in pts(\allowbreak{}d)], forall(\allowbreak{}[cnj\_ in pts(\allowbreak{}d)], [cni\_, cnj\_] in dleq(\allowbreak{}d) implies forall(\allowbreak{}[cnk\_ in pts(\allowbreak{}d)], [cni\_, cnk\_] in dleq(\allowbreak{}d) implies (\allowbreak{}vadd(\allowbreak{}m))(\allowbreak{}x(\allowbreak{}cnj\_), (\allowbreak{}vneg(\allowbreak{}m))(\allowbreak{}x(\allowbreak{}cnk\_))) in cnu\_)))))
\end{quote}
mentioned by 6 result(s)

\smallskip

\noindent\textbf{\texttt{is-cauchy-seq(s, f)}}\quad\textit{predicate}\index{is-cauchy-seq(s, f)@\texttt{is-cauchy-seq(s, f)}}\\
reads: Cauchy \\
defined by \texttt{is-cauchy-seq} \\
characterized by \texttt{is-cauchy-seq}  \textit{(definitional)}:
\begin{quote}\ttfamily\footnotesize\raggedright\noindent
forall(\allowbreak{}[s, f], is-cauchy-seq(\allowbreak{}s, f) iff is-metric-space(\allowbreak{}s) and f in fun(\allowbreak{}nn, pts(\allowbreak{}s)) and forall(\allowbreak{}[eps], pos-rr(\allowbreak{}eps) implies forsome(\allowbreak{}[n in nn], forall(\allowbreak{}[m in nn, n\_ in nn], n <= m and n <= n\_ implies (\allowbreak{}dist(\allowbreak{}s))(\allowbreak{}f(\allowbreak{}m), f(\allowbreak{}n\_)) <= eps))))
\end{quote}
also: \texttt{subspace-is-cauchy-seq} \\
mentioned by 61 result(s)

\smallskip

\noindent\textbf{\texttt{is-cc-diff-at(g, a, l)}}\quad\textit{predicate}\index{is-cc-diff-at(g, a, l)@\texttt{is-cc-diff-at(g, a, l)}}\\
reads: \$1 is differentiable at \$2 as a complex-valued function of a real variable, with derivative \$3 \\
defined by \texttt{is-cc-diff-at} \\
characterized by \texttt{is-cc-diff-at}  \textit{(definitional)}:
\begin{quote}\ttfamily\footnotesize\raggedright\noindent
forall(\allowbreak{}[g, a, l], is-cc-diff-at(\allowbreak{}g, a, l) iff g in fun(\allowbreak{}rr, cc) and a in rr and l in cc and is-diff-at(\allowbreak{}vnb-lambda(\allowbreak{}sv\_, rr, real-part(\allowbreak{}g(\allowbreak{}sv\_))), a, real-part(\allowbreak{}l)) and is-diff-at(\allowbreak{}vnb-lambda(\allowbreak{}sv\_, rr, imag-part(\allowbreak{}g(\allowbreak{}sv\_))), a, imag-part(\allowbreak{}l)))
\end{quote}
mentioned by 14 result(s)

\smallskip

\noindent\textbf{\texttt{is-cc-linear-functional(v, f)}}\quad\textit{predicate}\index{is-cc-linear-functional(v, f)@\texttt{is-cc-linear-functional(v, f)}}\\
reads: \$2 is a linear functional on \$1 \\
defined by \texttt{is-cc-linear-functional} \\
characterized by \texttt{is-cc-linear-functional}  \textit{(definitional)}:
\begin{quote}\ttfamily\footnotesize\raggedright\noindent
forall(\allowbreak{}[v, f], is-cc-linear-functional(\allowbreak{}v, f) iff f in fun(\allowbreak{}vec(\allowbreak{}v), cc) and forall(\allowbreak{}[hlx\_ in vec(\allowbreak{}v), hly\_ in vec(\allowbreak{}v)], f(\allowbreak{}(\allowbreak{}vadd(\allowbreak{}v))(\allowbreak{}hlx\_, hly\_)) = f(\allowbreak{}hlx\_) + f(\allowbreak{}hly\_)) and forall(\allowbreak{}[hlc\_ in cc, hlx\_ in vec(\allowbreak{}v)], f(\allowbreak{}(\allowbreak{}act(\allowbreak{}v))(\allowbreak{}hlc\_, hlx\_)) = hlc\_ * f(\allowbreak{}hlx\_)))
\end{quote}
mentioned by 12 result(s)

\smallskip

\noindent\textbf{\texttt{is-cc-linear-on(baa\_, baf\_)}}\quad\textit{predicate}\index{is-cc-linear-on(baa_, baf_)@\texttt{is-cc-linear-on(baa\_, baf\_)}}\\
reads: \$2 is a complex-linear functional on \$1 \\
defined by \texttt{is-cc-linear-on} \\
characterized by \texttt{is-cc-linear-on}  \textit{(definitional)}:
\begin{quote}\ttfamily\footnotesize\raggedright\noindent
forall(\allowbreak{}[baa\_, baf\_], is-cc-linear-on(\allowbreak{}baa\_, baf\_) iff baf\_ in fun(\allowbreak{}vec(\allowbreak{}baa\_), cc) and forall(\allowbreak{}[bax\_ in vec(\allowbreak{}baa\_), bay\_ in vec(\allowbreak{}baa\_)], baf\_(\allowbreak{}(\allowbreak{}vadd(\allowbreak{}baa\_))(\allowbreak{}bax\_, bay\_)) = baf\_(\allowbreak{}bax\_) + baf\_(\allowbreak{}bay\_)) and forall(\allowbreak{}[baz\_ in cc, bax\_ in vec(\allowbreak{}baa\_)], baf\_(\allowbreak{}(\allowbreak{}act(\allowbreak{}baa\_))(\allowbreak{}baz\_, bax\_)) = baz\_ * baf\_(\allowbreak{}bax\_)))
\end{quote}
mentioned by 27 result(s)

\smallskip

\noindent\textbf{\texttt{is-cc-module-over(aa, vv)}}\quad\textit{predicate}\index{is-cc-module-over(aa, vv)@\texttt{is-cc-module-over(aa, vv)}}\\
reads: \$2 is an \$1-module inside the complex numbers \\
defined by \texttt{is-cc-module-over} \\
characterized by \texttt{is-cc-module-over}  \textit{(definitional)}:
\begin{quote}\ttfamily\footnotesize\raggedright\noindent
forall(\allowbreak{}[aa, vv], is-cc-module-over(\allowbreak{}aa, vv) iff vv subset cc and 0 in vv and forall(\allowbreak{}[amx\_ in vv, amy\_ in vv], amx\_ + amy\_ in vv) and forall(\allowbreak{}[amx\_ in vv], -amx\_ in vv) and forall(\allowbreak{}[ama\_ in aa, amx\_ in vv], ama\_ * amx\_ in vv))
\end{quote}
mentioned by 10 result(s)

\smallskip

\noindent\textbf{\texttt{is-cc-space-dim-le(kk, vv, n)}}\quad\textit{predicate}\index{is-cc-space-dim-le(kk, vv, n)@\texttt{is-cc-space-dim-le(kk, vv, n)}}\\
reads: \$2 is a \$1-vector space inside the complex numbers of dimension at most \$3 \\
defined by \texttt{is-cc-space-dim-le} \\
characterized by \texttt{is-cc-space-dim-le}  \textit{(definitional)}:
\begin{quote}\ttfamily\footnotesize\raggedright\noindent
forall(\allowbreak{}[kk, vv, n], is-cc-space-dim-le(\allowbreak{}kk, vv, n) iff is-cc-module-over(\allowbreak{}kk, vv) and forsome(\allowbreak{}[amn\_ in nn], amn\_ <= n and forsome(\allowbreak{}[amg\_ in fun(\allowbreak{}interval(\allowbreak{}1, amn\_), vv)], forall(\allowbreak{}[amv\_ in vv], forsome(\allowbreak{}[amb\_ in fun(\allowbreak{}interval(\allowbreak{}1, amn\_), kk)], amv\_ = cc-lincomb(\allowbreak{}amn\_, amb\_, amg\_))))))
\end{quote}
mentioned by 10 result(s)

\smallskip

\noindent\textbf{\texttt{is-cc-subring(kk)}}\quad\textit{predicate}\index{is-cc-subring(kk)@\texttt{is-cc-subring(kk)}}\\
reads: \$1 is a subring of the complex numbers \\
defined by \texttt{is-cc-subring} \\
characterized by \texttt{is-cc-subring}  \textit{(definitional)}:
\begin{quote}\ttfamily\footnotesize\raggedright\noindent
forall(\allowbreak{}[kk], is-cc-subring(\allowbreak{}kk) iff kk subset cc and 0 in kk and 1 in kk and forall(\allowbreak{}[psu\_ in kk, psv\_ in kk], psu\_ + psv\_ in kk) and forall(\allowbreak{}[psu\_ in kk, psv\_ in kk], psu\_ * psv\_ in kk))
\end{quote}
also: \texttt{qq-is-cc-subring} \\
mentioned by 26 result(s)

\smallskip

\noindent\textbf{\texttt{is-ccont-star-subalgebra(rst\_, rsa\_)}}\quad\textit{predicate}\index{is-ccont-star-subalgebra(rst_, rsa_)@\texttt{is-ccont-star-subalgebra(rst\_, rsa\_)}}\\
reads: \$2 is a unital self-adjoint subalgebra of the continuous complex functions on \$1 \\
defined by \texttt{is-ccont-star-subalgebra} \\
characterized by \texttt{is-ccont-star-subalgebra}  \textit{(definitional)}:
\begin{quote}\ttfamily\footnotesize\raggedright\noindent
forall(\allowbreak{}[rst\_, rsa\_], is-ccont-star-subalgebra(\allowbreak{}rst\_, rsa\_) iff forall(\allowbreak{}[csf\_ in rsa\_], csf\_ in fun(\allowbreak{}pts(\allowbreak{}rst\_), cc) and is-hom-top-space(\allowbreak{}rst\_, metric-top(\allowbreak{}cc-ms), csf\_)) and forall(\allowbreak{}[csc\_ in cc], vnb-lambda(\allowbreak{}rsz\_, pts(\allowbreak{}rst\_), csc\_) in rsa\_) and forall(\allowbreak{}[csf\_ in rsa\_, csg\_ in rsa\_], vnb-lambda(\allowbreak{}rsz\_, pts(\allowbreak{}rst\_), csf\_(\allowbreak{}rsz\_) + csg\_(\allowbreak{}rsz\_)) in rsa\_) and forall(\allowbreak{}[csf\_ in rsa\_, csg\_ in rsa\_], vnb-lambda(\allowbreak{}rsz\_, pts(\allowbreak{}rst\_), csf\_(\allowbreak{}rsz\_) * csg\_(\allowbreak{}rsz\_)) in rsa\_) and forall(\allowbreak{}[csf\_ in rsa\_], vnb-lambda(\allowbreak{}rsz\_, pts(\allowbreak{}rst\_), conjugate(\allowbreak{}csf\_(\allowbreak{}rsz\_))) in rsa\_))
\end{quote}
mentioned by 9 result(s)

\smallskip

\noindent\textbf{\texttt{is-chain(grd, porel, ch)}}\quad\textit{predicate}\index{is-chain(grd, porel, ch)@\texttt{is-chain(grd, porel, ch)}}\\
reads: \$3 is a chain in \$1 under \$2 \\
defined by \texttt{is-chain} \\
characterized by \texttt{is-chain}  \textit{(definitional)}:
\begin{quote}\ttfamily\footnotesize\raggedright\noindent
forall(\allowbreak{}[grd, porel, ch], is-chain(\allowbreak{}grd, porel, ch) iff ch subset grd and forall(\allowbreak{}[x in ch, y in ch], [x, y] in porel or [y, x] in porel))
\end{quote}
also: \texttt{zorn-empty-is-chain} \texttt{zorn-zkept-is-chain} \texttt{greedy-chain-is-chain} \texttt{zorn-zkept-limit-is-chain} \\
mentioned by 11 result(s)

\smallskip

\noindent\textbf{\texttt{is-closable-graph(ubv\_, ubg\_)}}\quad\textit{predicate}\index{is-closable-graph(ubv_, ubg_)@\texttt{is-closable-graph(ubv\_, ubg\_)}}\\
reads: \$2 is a closable graph \\
defined by \texttt{is-closable-graph} \\
characterized by \texttt{is-closable-graph}  \textit{(definitional)}:
\begin{quote}\ttfamily\footnotesize\raggedright\noindent
forall(\allowbreak{}[ubv\_, ubg\_], is-closable-graph(\allowbreak{}ubv\_, ubg\_) iff is-functional-graph(\allowbreak{}ubv\_, graph-closure(\allowbreak{}ubv\_, ubg\_)))
\end{quote}
mentioned by 6 result(s)

\smallskip

\noindent\textbf{\texttt{is-closed(s, a)}}\quad\textit{predicate}\index{is-closed(s, a)@\texttt{is-closed(s, a)}}\\
reads: \$2 is closed in \$1 \\
defined by \texttt{is-closed} \\
characterized by \texttt{is-closed}  \textit{(definitional)}:
\begin{quote}\ttfamily\footnotesize\raggedright\noindent
forall(\allowbreak{}[s, a], is-closed(\allowbreak{}s, a) iff is-metric-space(\allowbreak{}s) and a subset pts(\allowbreak{}s) and is-open(\allowbreak{}s, complement-in(\allowbreak{}pts(\allowbreak{}s), a)))
\end{quote}
also: \texttt{is-closed-graph} \texttt{perp-is-closed} \texttt{carrier-is-closed} \texttt{closure-is-closed} \texttt{closed-ball-is-closed} \texttt{compact-subspace-is-closed} \\
mentioned by 125 result(s)

\smallskip

\noindent\textbf{\texttt{is-closed-form(ubv\_, ubd\_, ubf\_)}}\quad\textit{predicate}\index{is-closed-form(ubv_, ubd_, ubf_)@\texttt{is-closed-form(ubv\_, ubd\_, ubf\_)}}\\
reads: \$3 is a closed positive quadratic form on \$1 with domain \$2 \\
defined by \texttt{is-closed-form} \\
characterized by \texttt{is-closed-form}  \textit{(definitional)}:
\begin{quote}\ttfamily\footnotesize\raggedright\noindent
forall(\allowbreak{}[ubv\_, ubd\_, ubf\_], is-closed-form(\allowbreak{}ubv\_, ubd\_, ubf\_) iff is-positive-form(\allowbreak{}ubv\_, ubd\_, ubf\_) and forall(\allowbreak{}[ubq\_ in fun(\allowbreak{}nn, ubd\_), ubx\_], converges-to(\allowbreak{}ip-metric-space(\allowbreak{}ubv\_), ubq\_, ubx\_) implies forall(\allowbreak{}[ube\_], pos-rr(\allowbreak{}ube\_) implies forsome(\allowbreak{}[ubn\_ in nn], forall(\allowbreak{}[ubm\_ in nn, ubk\_ in nn], ubn\_ <= ubm\_ implies ubn\_ <= ubk\_ implies magnitude(\allowbreak{}ubf\_(\allowbreak{}(\allowbreak{}vadd(\allowbreak{}ubv\_))(\allowbreak{}ubq\_(\allowbreak{}ubm\_), (\allowbreak{}vneg(\allowbreak{}ubv\_))(\allowbreak{}ubq\_(\allowbreak{}ubk\_))), (\allowbreak{}vadd(\allowbreak{}ubv\_))(\allowbreak{}ubq\_(\allowbreak{}ubm\_), (\allowbreak{}vneg(\allowbreak{}ubv\_))(\allowbreak{}ubq\_(\allowbreak{}ubk\_))))) <= ube\_))) implies ubx\_ in ubd\_ and forall(\allowbreak{}[ube\_], pos-rr(\allowbreak{}ube\_) implies forsome(\allowbreak{}[ubn\_ in nn], forall(\allowbreak{}[ubm\_ in nn], ubn\_ <= ubm\_ implies magnitude(\allowbreak{}ubf\_(\allowbreak{}(\allowbreak{}vadd(\allowbreak{}ubv\_))(\allowbreak{}ubq\_(\allowbreak{}ubm\_), (\allowbreak{}vneg(\allowbreak{}ubv\_))(\allowbreak{}ubx\_)), (\allowbreak{}vadd(\allowbreak{}ubv\_))(\allowbreak{}ubq\_(\allowbreak{}ubm\_), (\allowbreak{}vneg(\allowbreak{}ubv\_))(\allowbreak{}ubx\_)))) <= ube\_)))))
\end{quote}
mentioned by 2 result(s)

\smallskip

\noindent\textbf{\texttt{is-closed-graph(ubv\_, ubg\_)}}\quad\textit{predicate}\index{is-closed-graph(ubv_, ubg_)@\texttt{is-closed-graph(ubv\_, ubg\_)}}\\
reads: \$2 is a closed graph \\
defined by \texttt{is-closed-graph} \\
characterized by \texttt{is-closed-graph}  \textit{(definitional)}:
\begin{quote}\ttfamily\footnotesize\raggedright\noindent
forall(\allowbreak{}[ubv\_, ubg\_], is-closed-graph(\allowbreak{}ubv\_, ubg\_) iff is-closed(\allowbreak{}ip-metric-space(\allowbreak{}hsum(\allowbreak{}ubv\_, ubv\_)), ubg\_))
\end{quote}
mentioned by 3 result(s)

\smallskip

\noindent\textbf{\texttt{is-closed-star-subalgebra-set(b2va\_, b2vs\_)}}\quad\textit{predicate}\index{is-closed-star-subalgebra-set(b2va_, b2vs_)@\texttt{is-closed-star-subalgebra-set(b2va\_, b2vs\_)}}\\
reads: \$2 is a norm-closed unital *-subalgebra of \$1 \\
defined by \texttt{is-closed-star-subalgebra-set} \\
characterized by \texttt{is-closed-star-subalgebra-set}  \textit{(definitional)}:
\begin{quote}\ttfamily\footnotesize\raggedright\noindent
forall(\allowbreak{}[b2va\_, b2vs\_], is-closed-star-subalgebra-set(\allowbreak{}b2va\_, b2vs\_) iff b2vs\_ subset vec(\allowbreak{}b2va\_) and vone(\allowbreak{}b2va\_) in b2vs\_ and forall(\allowbreak{}[b2vx\_ in b2vs\_, b2vy\_ in b2vs\_], (\allowbreak{}vadd(\allowbreak{}b2va\_))(\allowbreak{}b2vx\_, b2vy\_) in b2vs\_) and forall(\allowbreak{}[b2vz\_ in cc, b2vx\_ in b2vs\_], (\allowbreak{}act(\allowbreak{}b2va\_))(\allowbreak{}b2vz\_, b2vx\_) in b2vs\_) and forall(\allowbreak{}[b2vx\_ in b2vs\_, b2vy\_ in b2vs\_], (\allowbreak{}vmul(\allowbreak{}b2va\_))(\allowbreak{}b2vx\_, b2vy\_) in b2vs\_) and forall(\allowbreak{}[b2vx\_ in b2vs\_], (\allowbreak{}vstar(\allowbreak{}b2va\_))(\allowbreak{}b2vx\_) in b2vs\_) and forall(\allowbreak{}[b2vy\_ in vec(\allowbreak{}b2va\_)], forall(\allowbreak{}[b2ve\_], pos-rr(\allowbreak{}b2ve\_) implies forsome(\allowbreak{}[b2vx\_ in b2vs\_], (\allowbreak{}vnrm(\allowbreak{}b2va\_))(\allowbreak{}(\allowbreak{}vadd(\allowbreak{}b2va\_))(\allowbreak{}b2vy\_, (\allowbreak{}vneg(\allowbreak{}b2va\_))(\allowbreak{}b2vx\_))) < b2ve\_)) implies b2vy\_ in b2vs\_))
\end{quote}
mentioned by 4 result(s)

\smallskip

\noindent\textbf{\texttt{is-cmod-norm(cnm\_, cnf\_)}}\quad\textit{predicate}\index{is-cmod-norm(cnm_, cnf_)@\texttt{is-cmod-norm(cnm\_, cnf\_)}}\\
reads: \$2 is a norm on the complex module \$1 \\
defined by \texttt{is-cmod-norm} \\
characterized by \texttt{is-cmod-norm}  \textit{(definitional)}:
\begin{quote}\ttfamily\footnotesize\raggedright\noindent
forall(\allowbreak{}[cnm\_, cnf\_], is-cmod-norm(\allowbreak{}cnm\_, cnf\_) iff cnf\_ in fun(\allowbreak{}vec(\allowbreak{}cnm\_), rr) and forall(\allowbreak{}[cnx\_ in vec(\allowbreak{}cnm\_)], 0 <= cnf\_(\allowbreak{}cnx\_)) and forall(\allowbreak{}[cnx\_ in vec(\allowbreak{}cnm\_)], cnf\_(\allowbreak{}cnx\_) = 0 iff cnx\_ = vzero(\allowbreak{}cnm\_)) and forall(\allowbreak{}[cnc\_ in cc, cnx\_ in vec(\allowbreak{}cnm\_)], cnf\_(\allowbreak{}(\allowbreak{}act(\allowbreak{}cnm\_))(\allowbreak{}cnc\_, cnx\_)) = magnitude(\allowbreak{}cnc\_) * cnf\_(\allowbreak{}cnx\_)) and forall(\allowbreak{}[cnx\_ in vec(\allowbreak{}cnm\_), cny\_ in vec(\allowbreak{}cnm\_)], cnf\_(\allowbreak{}(\allowbreak{}vadd(\allowbreak{}cnm\_))(\allowbreak{}cnx\_, cny\_)) <= cnf\_(\allowbreak{}cnx\_) + cnf\_(\allowbreak{}cny\_)))
\end{quote}
mentioned by 91 result(s)

\smallskip

\noindent\textbf{\texttt{is-cnvs-bounded-op(cbv\_, cbw\_, cbt\_)}}\quad\textit{predicate}\index{is-cnvs-bounded-op(cbv_, cbw_, cbt_)@\texttt{is-cnvs-bounded-op(cbv\_, cbw\_, cbt\_)}}\\
reads: \$3 is a bounded operator from \$1 to \$2 \\
defined by \texttt{is-cnvs-bounded-op} \\
characterized by \texttt{is-cnvs-bounded-op}  \textit{(definitional)}:
\begin{quote}\ttfamily\footnotesize\raggedright\noindent
forall(\allowbreak{}[cbv\_, cbw\_, cbt\_], is-cnvs-bounded-op(\allowbreak{}cbv\_, cbw\_, cbt\_) iff is-cnvs-linear-op(\allowbreak{}cbv\_, cbw\_, cbt\_) and forsome(\allowbreak{}[cbm\_ in rr], 0 <= cbm\_ and forall(\allowbreak{}[cbx\_ in vec(\allowbreak{}cbv\_)], (\allowbreak{}vnrm(\allowbreak{}cbw\_))(\allowbreak{}cbt\_(\allowbreak{}cbx\_)) <= cbm\_ * (\allowbreak{}vnrm(\allowbreak{}cbv\_))(\allowbreak{}cbx\_))))
\end{quote}
mentioned by 6 result(s)

\smallskip

\noindent\textbf{\texttt{is-cnvs-linear-op(cbv\_, cbw\_, cbt\_)}}\quad\textit{predicate}\index{is-cnvs-linear-op(cbv_, cbw_, cbt_)@\texttt{is-cnvs-linear-op(cbv\_, cbw\_, cbt\_)}}\\
reads: \$3 is a complex-linear map from \$1 to \$2 \\
defined by \texttt{is-cnvs-linear-op} \\
characterized by \texttt{is-cnvs-linear-op}  \textit{(definitional)}:
\begin{quote}\ttfamily\footnotesize\raggedright\noindent
forall(\allowbreak{}[cbv\_, cbw\_, cbt\_], is-cnvs-linear-op(\allowbreak{}cbv\_, cbw\_, cbt\_) iff cbt\_ in fun(\allowbreak{}vec(\allowbreak{}cbv\_), vec(\allowbreak{}cbw\_)) and forall(\allowbreak{}[cbx\_ in vec(\allowbreak{}cbv\_), cby\_ in vec(\allowbreak{}cbv\_)], cbt\_(\allowbreak{}(\allowbreak{}vadd(\allowbreak{}cbv\_))(\allowbreak{}cbx\_, cby\_)) = (\allowbreak{}vadd(\allowbreak{}cbw\_))(\allowbreak{}cbt\_(\allowbreak{}cbx\_), cbt\_(\allowbreak{}cby\_))) and forall(\allowbreak{}[cbc\_ in cc, cbx\_ in vec(\allowbreak{}cbv\_)], cbt\_(\allowbreak{}(\allowbreak{}cact(\allowbreak{}cbv\_))(\allowbreak{}cbc\_, cbx\_)) = (\allowbreak{}cact(\allowbreak{}cbw\_))(\allowbreak{}cbc\_, cbt\_(\allowbreak{}cbx\_))))
\end{quote}
mentioned by 6 result(s)

\smallskip

\noindent\textbf{\texttt{is-cofinal-monotone(e, d, h)}}\quad\textit{predicate}\index{is-cofinal-monotone(e, d, h)@\texttt{is-cofinal-monotone(e, d, h)}}\\
reads: \$3 is a cofinal monotone map from \$1 to \$2 \\
defined by \texttt{is-cofinal-monotone} \\
characterized by \texttt{is-cofinal-monotone}  \textit{(definitional)}:
\begin{quote}\ttfamily\footnotesize\raggedright\noindent
forall(\allowbreak{}[e, d, h], is-cofinal-monotone(\allowbreak{}e, d, h) iff h in fun(\allowbreak{}pts(\allowbreak{}e), pts(\allowbreak{}d)) and forall(\allowbreak{}[cmi\_ in pts(\allowbreak{}e), cmj\_ in pts(\allowbreak{}e)], [cmi\_, cmj\_] in dleq(\allowbreak{}e) implies [h(\allowbreak{}cmi\_), h(\allowbreak{}cmj\_)] in dleq(\allowbreak{}d)) and forall(\allowbreak{}[cmk\_ in pts(\allowbreak{}d)], forsome(\allowbreak{}[cml\_ in pts(\allowbreak{}e)], [cmk\_, h(\allowbreak{}cml\_)] in dleq(\allowbreak{}d))))
\end{quote}
mentioned by 4 result(s)

\smallskip

\noindent\textbf{\texttt{is-column-step(r, m, n, b)}}\quad\textit{predicate}\index{is-column-step(r, m, n, b)@\texttt{is-column-step(r, m, n, b)}}\\
reads: \$4 is an \$2-by-\$3 column step matrix over \$1 \\
defined by \texttt{is-column-step} \\
characterized by \texttt{is-column-step}  \textit{(definitional)}:
\begin{quote}\ttfamily\footnotesize\raggedright\noindent
forall(\allowbreak{}[r, m, n, b], is-column-step(\allowbreak{}r, m, n, b) iff forall(\allowbreak{}[csi\_ in interval(\allowbreak{}1, m), csj\_ in interval(\allowbreak{}1, n)], csi\_ < csj\_ implies entry(\allowbreak{}b, csi\_, csj\_) = zero(\allowbreak{}r)) and forall(\allowbreak{}[csi\_ in interval(\allowbreak{}1, m)], csi\_ in interval(\allowbreak{}1, n) implies entry(\allowbreak{}b, csi\_, csi\_) = zero(\allowbreak{}r) or entry(\allowbreak{}b, csi\_, csi\_) = one(\allowbreak{}r)))
\end{quote}
mentioned by 4 result(s)

\smallskip

\noindent\textbf{\texttt{is-comm-monoid(s)}}\quad\textit{predicate}\index{is-comm-monoid(s)@\texttt{is-comm-monoid(s)}}\\
reads: a comm monoid \\
a predicate \\
characterized by \texttt{is-comm-monoid}  \textit{(definitional)}:
\begin{quote}\ttfamily\footnotesize\raggedright\noindent
forall(\allowbreak{}[s], is-comm-monoid(\allowbreak{}s) iff length(\allowbreak{}s) = 3 and carr(\allowbreak{}s) in set and opr(\allowbreak{}s) in fun(\allowbreak{}cartesian(\allowbreak{}carr(\allowbreak{}s), carr(\allowbreak{}s)), carr(\allowbreak{}s)) and iden(\allowbreak{}s) in carr(\allowbreak{}s) and is-associative(\allowbreak{}opr(\allowbreak{}s), carr(\allowbreak{}s)) and is-identity(\allowbreak{}opr(\allowbreak{}s), iden(\allowbreak{}s), carr(\allowbreak{}s)) and is-commutative(\allowbreak{}opr(\allowbreak{}s), carr(\allowbreak{}s)) and s in set)
\end{quote}
also: \texttt{commutative-ring-multiplicative-cm-is-comm-monoid} \texttt{commutative-ring-multiplicative-cm-is-comm-monoid-normed-field-as-commutative-ring} \texttt{is-comm-monoid-is-set} \texttt{rr-pos-star-is-comm-monoid} \texttt{nn-add-monoid-is-comm-monoid} \texttt{zz-add-monoid-is-comm-monoid} \texttt{fin-subset-monoid-is-comm-monoid} \texttt{multi-index-monoid-is-comm-monoid} \\
mentioned by 87 result(s)

\smallskip

\noindent\textbf{\texttt{is-commutative(op, crr)}}\quad\textit{predicate}\index{is-commutative(op, crr)@\texttt{is-commutative(op, crr)}}\\
reads: \$1 is commutative on \$2 \\
defined by \texttt{is-commutative} \\
characterized by \texttt{is-commutative}  \textit{(definitional)}:
\begin{quote}\ttfamily\footnotesize\raggedright\noindent
forall(\allowbreak{}[op, crr], is-commutative(\allowbreak{}op, crr) iff forall(\allowbreak{}[u in crr, v in crr], op(\allowbreak{}u, v) = op(\allowbreak{}v, u)))
\end{quote}
mentioned by 58 result(s)

\smallskip

\noindent\textbf{\texttt{is-commutative-banach-algebra(s)}}\quad\textit{predicate}\index{is-commutative-banach-algebra(s)@\texttt{is-commutative-banach-algebra(s)}}\\
reads: a commutative Banach algebra with unit \\
a predicate \\
characterized by \texttt{is-commutative-banach-algebra-def}  \textit{(definitional)}:
\begin{quote}\ttfamily\footnotesize\raggedright\noindent
forall(\allowbreak{}[s], is-commutative-banach-algebra(\allowbreak{}s) iff is-unital-banach-algebra(\allowbreak{}s) and forall(\allowbreak{}[x\_ in vec(\allowbreak{}s), y\_ in vec(\allowbreak{}s)], (\allowbreak{}vmul(\allowbreak{}s))(\allowbreak{}x\_, y\_) = (\allowbreak{}vmul(\allowbreak{}s))(\allowbreak{}y\_, x\_)))
\end{quote}
also: \texttt{is-commutative-banach-algebra-is-set} \\
mentioned by 37 result(s)

\smallskip

\noindent\textbf{\texttt{is-commutative-c-star-algebra(s)}}\quad\textit{predicate}\index{is-commutative-c-star-algebra(s)@\texttt{is-commutative-c-star-algebra(s)}}\\
reads: a commutative C*-algebra with unit \\
a predicate \\
characterized by \texttt{is-commutative-c-star-algebra-def}  \textit{(definitional)}:
\begin{quote}\ttfamily\footnotesize\raggedright\noindent
forall(\allowbreak{}[s], is-commutative-c-star-algebra(\allowbreak{}s) iff is-unital-c-star-algebra(\allowbreak{}s) and forall(\allowbreak{}[x\_ in vec(\allowbreak{}s), y\_ in vec(\allowbreak{}s)], (\allowbreak{}vmul(\allowbreak{}s))(\allowbreak{}x\_, y\_) = (\allowbreak{}vmul(\allowbreak{}s))(\allowbreak{}y\_, x\_)))
\end{quote}
also: \texttt{is-commutative-c-star-algebra-is-set} \\
mentioned by 17 result(s)

\smallskip

\noindent\textbf{\texttt{is-commutative-ring(s)}}\quad\textit{predicate}\index{is-commutative-ring(s)@\texttt{is-commutative-ring(s)}}\\
reads: a commutative ring \\
a predicate \\
characterized by \texttt{is-commutative-ring-def}  \textit{(definitional)}:
\begin{quote}\ttfamily\footnotesize\raggedright\noindent
forall(\allowbreak{}[s], is-commutative-ring(\allowbreak{}s) iff is-ring(\allowbreak{}s) and forall(\allowbreak{}[a in carr(\allowbreak{}s), b in carr(\allowbreak{}s)], (\allowbreak{}mul(\allowbreak{}s))(\allowbreak{}a, b) = (\allowbreak{}mul(\allowbreak{}s))(\allowbreak{}b, a)))
\end{quote}
also: \texttt{normed-field-as-commutative-ring-is-commutative-ring} \texttt{qq-is-commutative-ring} \texttt{zz-is-commutative-ring} \texttt{mpoly-is-commutative-ring} \texttt{is-commutative-ring-is-set} \texttt{op-alg-is-commutative-ring} \texttt{laurent-ring-is-commutative-ring} \texttt{integral-domain-is-commutative-ring} \\
mentioned by 612 result(s)

\smallskip

\noindent\textbf{\texttt{is-compact(s)}}\quad\textit{predicate}\index{is-compact(s)@\texttt{is-compact(s)}}\\
reads: \$1 is compact \\
defined by \texttt{is-compact} \\
characterized by \texttt{is-compact}  \textit{(definitional)}:
\begin{quote}\ttfamily\footnotesize\raggedright\noindent
forall(\allowbreak{}[s], is-compact(\allowbreak{}s) iff is-metric-space(\allowbreak{}s) and forall(\allowbreak{}[c in set], is-open-cover(\allowbreak{}s, c) implies forsome(\allowbreak{}[f in set], f subset c and card(\allowbreak{}f) in nn and is-open-cover(\allowbreak{}s, f))))
\end{quote}
also: \texttt{is-compact-nv-domain} \texttt{is-compact-exhaustion} \texttt{is-compact-classes} \texttt{closed-subset-of-compact-is-compact} \\
mentioned by 282 result(s)

\smallskip

\noindent\textbf{\texttt{is-compact-exhaustion(ces\_, cek\_)}}\quad\textit{predicate}\index{is-compact-exhaustion(ces_, cek_)@\texttt{is-compact-exhaustion(ces\_, cek\_)}}\\
reads: \$2 is an exhaustive sequence of compact subsets of \$1 \\
defined by \texttt{is-compact-exhaustion} \\
characterized by \texttt{is-compact-exhaustion}  \textit{(definitional)}:
\begin{quote}\ttfamily\footnotesize\raggedright\noindent
forall(\allowbreak{}[ces\_, cek\_], is-compact-exhaustion(\allowbreak{}ces\_, cek\_) iff is-metric-space(\allowbreak{}ces\_) and cek\_ in fun(\allowbreak{}nn, power(\allowbreak{}pts(\allowbreak{}ces\_))) and forall(\allowbreak{}[cen\_ in nn], is-compact(\allowbreak{}subspace-ms(\allowbreak{}ces\_, cek\_(\allowbreak{}cen\_)))) and forall(\allowbreak{}[cen\_ in nn], forsome(\allowbreak{}[cer\_], pos-rr(\allowbreak{}cer\_) and forall(\allowbreak{}[cey\_ in cek\_(\allowbreak{}cen\_)], ball(\allowbreak{}ces\_, cey\_, cer\_) subset cek\_(\allowbreak{}cen\_ + 1)))) and forall(\allowbreak{}[cex\_ in pts(\allowbreak{}ces\_)], forsome(\allowbreak{}[cen\_ in nn], cex\_ in cek\_(\allowbreak{}cen\_))) and forall(\allowbreak{}[cea\_], cea\_ subset pts(\allowbreak{}ces\_) implies is-compact(\allowbreak{}subspace-ms(\allowbreak{}ces\_, cea\_)) implies forsome(\allowbreak{}[cen\_ in nn], cea\_ subset cek\_(\allowbreak{}cen\_))))
\end{quote}
also: \texttt{cc-exhaust-is-compact-exhaustion} \\
mentioned by 6 result(s)

\smallskip

\noindent\textbf{\texttt{is-compact-nv-domain(cwk\_, cwa\_)}}\quad\textit{predicate}\index{is-compact-nv-domain(cwk_, cwa_)@\texttt{is-compact-nv-domain(cwk\_, cwa\_)}}\\
reads: \$2 is a compact non-archimedean valued domain inside the field \$1 \\
defined by \texttt{is-compact-nv-domain} \\
characterized by \texttt{is-compact-nv-domain}  \textit{(definitional)}:
\begin{quote}\ttfamily\footnotesize\raggedright\noindent
forall(\allowbreak{}[cwk\_, cwa\_], is-compact-nv-domain(\allowbreak{}cwk\_, cwa\_) iff is-nv-domain(\allowbreak{}cwk\_, cwa\_) and is-compact(\allowbreak{}subspace-ms(\allowbreak{}nf-metric-space(\allowbreak{}cwk\_), cwa\_)))
\end{quote}
mentioned by 85 result(s)

\smallskip

\noindent\textbf{\texttt{is-compact-op(kqv\_, kqw\_, kqt\_)}}\quad\textit{predicate}\index{is-compact-op(kqv_, kqw_, kqt_)@\texttt{is-compact-op(kqv\_, kqw\_, kqt\_)}}\\
reads: \$3 is a compact operator from \$1 to \$2 \\
defined by \texttt{is-compact-op} \\
characterized by \texttt{is-compact-op}  \textit{(definitional)}:
\begin{quote}\ttfamily\footnotesize\raggedright\noindent
forall(\allowbreak{}[kqv\_, kqw\_, kqt\_], is-compact-op(\allowbreak{}kqv\_, kqw\_, kqt\_) iff kqt\_ in bounded-ops(\allowbreak{}kqv\_, kqw\_) and forall(\allowbreak{}[kqf\_ in fun(\allowbreak{}nn, vec(\allowbreak{}kqv\_))], forall(\allowbreak{}[kqn\_ in nn], ip-norm(\allowbreak{}kqv\_, kqf\_(\allowbreak{}kqn\_)) <= 1) implies forsome(\allowbreak{}[kqp\_], strictly-mono-nn(\allowbreak{}kqp\_) and forsome(\allowbreak{}[kql\_], converges-to(\allowbreak{}ip-metric-space(\allowbreak{}kqw\_), vnb-lambda(\allowbreak{}kqn\_, nn, kqt\_(\allowbreak{}kqf\_(\allowbreak{}kqp\_(\allowbreak{}kqn\_)))), kql\_)))))
\end{quote}
mentioned by 41 result(s)

\smallskip

\noindent\textbf{\texttt{is-compact-t(s)}}\quad\textit{predicate}\index{is-compact-t(s)@\texttt{is-compact-t(s)}}\\
reads: the topological space \$1 is compact \\
defined by \texttt{is-compact-t} \\
characterized by \texttt{is-compact-t}  \textit{(definitional)}:
\begin{quote}\ttfamily\footnotesize\raggedright\noindent
forall(\allowbreak{}[s], is-compact-t(\allowbreak{}s) iff is-top-space(\allowbreak{}s) and forall(\allowbreak{}[c], forall(\allowbreak{}[u in c], u in opens(\allowbreak{}s)) implies big-union(\allowbreak{}u, c, u) == pts(\allowbreak{}s) implies forsome(\allowbreak{}[f], f subset c and card(\allowbreak{}f) in nn and big-union(\allowbreak{}u, f, u) == pts(\allowbreak{}s))))
\end{quote}
mentioned by 204 result(s)

\smallskip

\noindent\textbf{\texttt{is-complete(s)}}\quad\textit{predicate}\index{is-complete(s)@\texttt{is-complete(s)}}\\
reads: complete \\
defined by \texttt{is-complete} \\
characterized by \texttt{is-complete}  \textit{(definitional)}:
\begin{quote}\ttfamily\footnotesize\raggedright\noindent
forall(\allowbreak{}[s], is-complete(\allowbreak{}s) iff is-metric-space(\allowbreak{}s) and forall(\allowbreak{}[f], is-cauchy-seq(\allowbreak{}s, f) implies converges(\allowbreak{}s, f)))
\end{quote}
also: \texttt{is-complete-nv-domain} \texttt{completion-is-complete} \\
mentioned by 89 result(s)

\smallskip

\noindent\textbf{\texttt{is-complete-measure(omega, ca, mu)}}\quad\textit{predicate}\index{is-complete-measure(omega, ca, mu)@\texttt{is-complete-measure(omega, ca, mu)}}\\
reads: \$3 is a complete measure on the sigma-algebra \$2 of subsets of \$1 \\
defined by \texttt{is-complete-measure} \\
characterized by \texttt{is-complete-measure}  \textit{(definitional)}:
\begin{quote}\ttfamily\footnotesize\raggedright\noindent
forall(\allowbreak{}[omega, ca, mu], is-complete-measure(\allowbreak{}omega, ca, mu) iff is-measure(\allowbreak{}omega, ca, mu) and forall(\allowbreak{}[nls\_], is-null-set(\allowbreak{}omega, ca, mu, nls\_) implies nls\_ in ca))
\end{quote}
mentioned by 11 result(s)

\smallskip

\noindent\textbf{\texttt{is-complete-nv-domain(c2k\_, c2a\_)}}\quad\textit{predicate}\index{is-complete-nv-domain(c2k_, c2a_)@\texttt{is-complete-nv-domain(c2k\_, c2a\_)}}\\
reads: \$2 is a complete non-archimedean valued domain inside the field \$1 \\
defined by \texttt{is-complete-nv-domain} \\
characterized by \texttt{is-complete-nv-domain}  \textit{(definitional)}:
\begin{quote}\ttfamily\footnotesize\raggedright\noindent
forall(\allowbreak{}[c2k\_, c2a\_], is-complete-nv-domain(\allowbreak{}c2k\_, c2a\_) iff is-nv-domain(\allowbreak{}c2k\_, c2a\_) and is-complete(\allowbreak{}subspace-ms(\allowbreak{}nf-metric-space(\allowbreak{}c2k\_), c2a\_)))
\end{quote}
mentioned by 7 result(s)

\smallskip

\noindent\textbf{\texttt{is-complete-subset(m, t, a)}}\quad\textit{predicate}\index{is-complete-subset(m, t, a)@\texttt{is-complete-subset(m, t, a)}}\\
reads: \$3 is complete in the topological vector space \$1, \$2 \\
defined by \texttt{is-complete-subset} \\
characterized by \texttt{is-complete-subset}  \textit{(definitional)}:
\begin{quote}\ttfamily\footnotesize\raggedright\noindent
forall(\allowbreak{}[m, t, a], is-complete-subset(\allowbreak{}m, t, a) iff a subset vec(\allowbreak{}m) and forall(\allowbreak{}[csd\_, csx\_], is-cauchy-net(\allowbreak{}m, t, csd\_, csx\_) implies forall(\allowbreak{}[csi\_ in pts(\allowbreak{}csd\_)], csx\_(\allowbreak{}csi\_) in a) implies forsome(\allowbreak{}[csl\_ in a], net-converges-to(\allowbreak{}t, csd\_, csx\_, csl\_))))
\end{quote}
mentioned by 4 result(s)

\smallskip

\noindent\textbf{\texttt{is-complex-banach-algebra(s)}}\quad\textit{predicate}\index{is-complex-banach-algebra(s)@\texttt{is-complex-banach-algebra(s)}}\\
reads: a complex Banach algebra \\
a predicate \\
characterized by \texttt{is-complex-banach-algebra-def}  \textit{(definitional)}:
\begin{quote}\ttfamily\footnotesize\raggedright\noindent
forall(\allowbreak{}[s], is-complex-banach-algebra(\allowbreak{}s) iff is-complex-normed-algebra(\allowbreak{}s) and is-complete(\allowbreak{}nag-metric-space(\allowbreak{}cna-as-normed-ag(\allowbreak{}s))))
\end{quote}
also: \texttt{uba-as-cba-is-complex-banach-algebra} \texttt{nuiba-as-cba-is-complex-banach-algebra} \texttt{uba-as-cba-is-complex-banach-algebra-uiba-as-uba} \texttt{nuiba-as-cba-is-complex-banach-algebra-uiba-as-nuiba} \texttt{is-complex-banach-algebra-is-set} \texttt{unitization-is-complex-banach-algebra} \\
mentioned by 25 result(s)

\smallskip

\noindent\textbf{\texttt{is-complex-banach-space(s)}}\quad\textit{predicate}\index{is-complex-banach-space(s)@\texttt{is-complex-banach-space(s)}}\\
reads: a complex Banach space \\
a predicate \\
characterized by \texttt{is-complex-banach-space-def}  \textit{(definitional)}:
\begin{quote}\ttfamily\footnotesize\raggedright\noindent
forall(\allowbreak{}[s], is-complex-banach-space(\allowbreak{}s) iff is-complex-normed-vector-space(\allowbreak{}s) and is-complete(\allowbreak{}cnvs-metric-space(\allowbreak{}s)))
\end{quote}
also: \texttt{is-complex-banach-space-is-set} \\
mentioned by 13 result(s)

\smallskip

\noindent\textbf{\texttt{is-complex-inner-product-space(s)}}\quad\textit{predicate}\index{is-complex-inner-product-space(s)@\texttt{is-complex-inner-product-space(s)}}\\
reads: \$1 is a complex inner product space \\
a predicate \\
characterized by \texttt{is-complex-inner-product-space}  \textit{(definitional)}:
\begin{quote}\ttfamily\footnotesize\raggedright\noindent
forall(\allowbreak{}[s], is-complex-inner-product-space(\allowbreak{}s) iff length(\allowbreak{}s) = 7 and is-ring(\allowbreak{}scal(\allowbreak{}s)) and vec(\allowbreak{}s) in set and vadd(\allowbreak{}s) in fun(\allowbreak{}cartesian(\allowbreak{}vec(\allowbreak{}s), vec(\allowbreak{}s)), vec(\allowbreak{}s)) and vzero(\allowbreak{}s) in vec(\allowbreak{}s) and vneg(\allowbreak{}s) in fun(\allowbreak{}vec(\allowbreak{}s), vec(\allowbreak{}s)) and act(\allowbreak{}s) in fun(\allowbreak{}cartesian(\allowbreak{}carr(\allowbreak{}scal(\allowbreak{}s)), vec(\allowbreak{}s)), vec(\allowbreak{}s)) and ip(\allowbreak{}s) in fun(\allowbreak{}cartesian(\allowbreak{}vec(\allowbreak{}s), vec(\allowbreak{}s)), cc) and is-associative(\allowbreak{}vadd(\allowbreak{}s), vec(\allowbreak{}s)) and is-commutative(\allowbreak{}vadd(\allowbreak{}s), vec(\allowbreak{}s)) and is-identity(\allowbreak{}vadd(\allowbreak{}s), vzero(\allowbreak{}s), vec(\allowbreak{}s)) and has-inverses(\allowbreak{}vadd(\allowbreak{}s), vzero(\allowbreak{}s), vneg(\allowbreak{}s), vec(\allowbreak{}s)) and scal(\allowbreak{}s) = normed-field-as-commutative-ring(\allowbreak{}cc-normed-field) and carr(\allowbreak{}scal(\allowbreak{}s)) = cc and forall(\allowbreak{}[r\_ in carr(\allowbreak{}scal(\allowbreak{}s)), x\_ in vec(\allowbreak{}s), y\_ in vec(\allowbreak{}s)], (\allowbreak{}act(\allowbreak{}s))(\allowbreak{}r\_, (\allowbreak{}vadd(\allowbreak{}s))(\allowbreak{}x\_, y\_)) = (\allowbreak{}vadd(\allowbreak{}s))(\allowbreak{}(\allowbreak{}act(\allowbreak{}s))(\allowbreak{}r\_, x\_), (\allowbreak{}act(\allowbreak{}s))(\allowbreak{}r\_, y\_))) and forall(\allowbreak{}[r\_ in carr(\allowbreak{}scal(\allowbreak{}s)), s\_ in carr(\allowbreak{}scal(\allowbreak{}s)), x\_ in vec(\allowbreak{}s)], (\allowbreak{}act(\allowbreak{}s))(\allowbreak{}(\allowbreak{}add(\allowbreak{}scal(\allowbreak{}s)))(\allowbreak{}r\_, s\_), x\_) = (\allowbreak{}vadd(\allowbreak{}s))(\allowbreak{}(\allowbreak{}act(\allowbreak{}s))(\allowbreak{}r\_, x\_), (\allowbreak{}act(\allowbreak{}s))(\allowbreak{}s\_, x\_))) and forall(\allowbreak{}[r\_ in carr(\allowbreak{}scal(\allowbreak{}s)), s\_ in carr(\allowbreak{}scal(\allowbreak{}s)), x\_ in vec(\allowbreak{}s)], (\allowbreak{}act(\allowbreak{}s))(\allowbreak{}(\allowbreak{}mul(\allowbreak{}scal(\allowbreak{}s)))(\allowbreak{}r\_, s\_), x\_) = (\allowbreak{}act(\allowbreak{}s))(\allowbreak{}r\_, (\allowbreak{}act(\allowbreak{}s))(\allowbreak{}s\_, x\_))) and forall(\allowbreak{}[x\_ in vec(\allowbreak{}s)], (\allowbreak{}act(\allowbreak{}s))(\allowbreak{}one(\allowbreak{}scal(\allowbreak{}s)), x\_) = x\_) and forall(\allowbreak{}[x\_ in vec(\allowbreak{}s), y\_ in vec(\allowbreak{}s), z\_ in vec(\allowbreak{}s)], (\allowbreak{}ip(\allowbreak{}s))(\allowbreak{}(\allowbreak{}vadd(\allowbreak{}s))(\allowbreak{}x\_, y\_), z\_) = (\allowbreak{}ip(\allowbreak{}s))(\allowbreak{}x\_, z\_) + (\allowbreak{}ip(\allowbreak{}s))(\allowbreak{}y\_, z\_)) and forall(\allowbreak{}[a\_ in cc, x\_ in vec(\allowbreak{}s), y\_ in vec(\allowbreak{}s)], (\allowbreak{}ip(\allowbreak{}s))(\allowbreak{}(\allowbreak{}act(\allowbreak{}s))(\allowbreak{}a\_, x\_), y\_) = a\_ * (\allowbreak{}ip(\allowbreak{}s))(\allowbreak{}x\_, y\_)) and forall(\allowbreak{}[x\_ in vec(\allowbreak{}s), y\_ in vec(\allowbreak{}s)], (\allowbreak{}ip(\allowbreak{}s))(\allowbreak{}y\_, x\_) = conjugate(\allowbreak{}(\allowbreak{}ip(\allowbreak{}s))(\allowbreak{}x\_, y\_))) and forall(\allowbreak{}[x\_ in vec(\allowbreak{}s)], (\allowbreak{}ip(\allowbreak{}s))(\allowbreak{}x\_, x\_) in rr) and forall(\allowbreak{}[x\_ in vec(\allowbreak{}s)], 0 <= (\allowbreak{}ip(\allowbreak{}s))(\allowbreak{}x\_, x\_)) and forall(\allowbreak{}[x\_ in vec(\allowbreak{}s)], (\allowbreak{}ip(\allowbreak{}s))(\allowbreak{}x\_, x\_) = 0 iff x\_ = vzero(\allowbreak{}s)) and s in set)
\end{quote}
also: \texttt{cphs-is-complex-inner-product-space} \texttt{is-complex-inner-product-space-is-set} \\
mentioned by 905 result(s)

\smallskip

\noindent\textbf{\texttt{is-complex-measurable-fn(omega, ca, u)}}\quad\textit{predicate}\index{is-complex-measurable-fn(omega, ca, u)@\texttt{is-complex-measurable-fn(omega, ca, u)}}\\
reads: \$3 is a measurable complex-valued function on \$1 with \$2 \\
defined by \texttt{is-complex-measurable-fn} \\
characterized by \texttt{is-complex-measurable-fn}  \textit{(definitional)}:
\begin{quote}\ttfamily\footnotesize\raggedright\noindent
forall(\allowbreak{}[omega, ca, u], is-complex-measurable-fn(\allowbreak{}omega, ca, u) iff u in fun(\allowbreak{}omega, cc) and is-real-measurable-fn(\allowbreak{}omega, ca, vnb-lambda(\allowbreak{}crx\_, omega, real-part(\allowbreak{}u(\allowbreak{}crx\_)))) and is-real-measurable-fn(\allowbreak{}omega, ca, vnb-lambda(\allowbreak{}crx\_, omega, imag-part(\allowbreak{}u(\allowbreak{}crx\_)))))
\end{quote}
mentioned by 6 result(s)

\smallskip

\noindent\textbf{\texttt{is-complex-measure(omega, ca, nu)}}\quad\textit{predicate}\index{is-complex-measure(omega, ca, nu)@\texttt{is-complex-measure(omega, ca, nu)}}\\
reads: \$3 is a complex measure on the sigma-algebra \$2 of subsets of \$1 \\
defined by \texttt{is-complex-measure} \\
characterized by \texttt{is-complex-measure}  \textit{(definitional)}:
\begin{quote}\ttfamily\footnotesize\raggedright\noindent
forall(\allowbreak{}[omega, ca, nu], is-complex-measure(\allowbreak{}omega, ca, nu) iff is-sigma-algebra(\allowbreak{}omega, ca) and nu in fun(\allowbreak{}ca, cc) and forsome(\allowbreak{}[cm1\_, cm2\_, cm3\_, cm4\_], is-finite-measure(\allowbreak{}omega, ca, cm1\_) and is-finite-measure(\allowbreak{}omega, ca, cm2\_) and is-finite-measure(\allowbreak{}omega, ca, cm3\_) and is-finite-measure(\allowbreak{}omega, ca, cm4\_) and forall(\allowbreak{}[cma\_ in ca], nu(\allowbreak{}cma\_) = cm1\_(\allowbreak{}cma\_) - cm2\_(\allowbreak{}cma\_) + 1i * (\allowbreak{}cm3\_(\allowbreak{}cma\_) - cm4\_(\allowbreak{}cma\_)))))
\end{quote}
mentioned by 3 result(s)

\smallskip

\noindent\textbf{\texttt{is-complex-normed-algebra(s)}}\quad\textit{predicate}\index{is-complex-normed-algebra(s)@\texttt{is-complex-normed-algebra(s)}}\\
reads: a complex normed algebra \\
a predicate \\
characterized by \texttt{is-complex-normed-algebra}  \textit{(definitional)}:
\begin{quote}\ttfamily\footnotesize\raggedright\noindent
forall(\allowbreak{}[s], is-complex-normed-algebra(\allowbreak{}s) iff length(\allowbreak{}s) = 8 and is-ring(\allowbreak{}scal(\allowbreak{}s)) and vec(\allowbreak{}s) in set and vadd(\allowbreak{}s) in fun(\allowbreak{}cartesian(\allowbreak{}vec(\allowbreak{}s), vec(\allowbreak{}s)), vec(\allowbreak{}s)) and vzero(\allowbreak{}s) in vec(\allowbreak{}s) and vneg(\allowbreak{}s) in fun(\allowbreak{}vec(\allowbreak{}s), vec(\allowbreak{}s)) and act(\allowbreak{}s) in fun(\allowbreak{}cartesian(\allowbreak{}carr(\allowbreak{}scal(\allowbreak{}s)), vec(\allowbreak{}s)), vec(\allowbreak{}s)) and vnrm(\allowbreak{}s) in fun(\allowbreak{}vec(\allowbreak{}s), rr) and vmul(\allowbreak{}s) in fun(\allowbreak{}cartesian(\allowbreak{}vec(\allowbreak{}s), vec(\allowbreak{}s)), vec(\allowbreak{}s)) and is-associative(\allowbreak{}vadd(\allowbreak{}s), vec(\allowbreak{}s)) and is-commutative(\allowbreak{}vadd(\allowbreak{}s), vec(\allowbreak{}s)) and is-identity(\allowbreak{}vadd(\allowbreak{}s), vzero(\allowbreak{}s), vec(\allowbreak{}s)) and has-inverses(\allowbreak{}vadd(\allowbreak{}s), vzero(\allowbreak{}s), vneg(\allowbreak{}s), vec(\allowbreak{}s)) and is-associative(\allowbreak{}vmul(\allowbreak{}s), vec(\allowbreak{}s)) and scal(\allowbreak{}s) = normed-field-as-commutative-ring(\allowbreak{}cc-normed-field) and carr(\allowbreak{}scal(\allowbreak{}s)) = cc and forall(\allowbreak{}[r\_ in carr(\allowbreak{}scal(\allowbreak{}s)), x\_ in vec(\allowbreak{}s), y\_ in vec(\allowbreak{}s)], (\allowbreak{}act(\allowbreak{}s))(\allowbreak{}r\_, (\allowbreak{}vadd(\allowbreak{}s))(\allowbreak{}x\_, y\_)) = (\allowbreak{}vadd(\allowbreak{}s))(\allowbreak{}(\allowbreak{}act(\allowbreak{}s))(\allowbreak{}r\_, x\_), (\allowbreak{}act(\allowbreak{}s))(\allowbreak{}r\_, y\_))) and forall(\allowbreak{}[r\_ in carr(\allowbreak{}scal(\allowbreak{}s)), s\_ in carr(\allowbreak{}scal(\allowbreak{}s)), x\_ in vec(\allowbreak{}s)], (\allowbreak{}act(\allowbreak{}s))(\allowbreak{}(\allowbreak{}add(\allowbreak{}scal(\allowbreak{}s)))(\allowbreak{}r\_, s\_), x\_) = (\allowbreak{}vadd(\allowbreak{}s))(\allowbreak{}(\allowbreak{}act(\allowbreak{}s))(\allowbreak{}r\_, x\_), (\allowbreak{}act(\allowbreak{}s))(\allowbreak{}s\_, x\_))) and forall(\allowbreak{}[r\_ in carr(\allowbreak{}scal(\allowbreak{}s)), s\_ in carr(\allowbreak{}scal(\allowbreak{}s)), x\_ in vec(\allowbreak{}s)], (\allowbreak{}act(\allowbreak{}s))(\allowbreak{}(\allowbreak{}mul(\allowbreak{}scal(\allowbreak{}s)))(\allowbreak{}r\_, s\_), x\_) = (\allowbreak{}act(\allowbreak{}s))(\allowbreak{}r\_, (\allowbreak{}act(\allowbreak{}s))(\allowbreak{}s\_, x\_))) and forall(\allowbreak{}[x\_ in vec(\allowbreak{}s)], (\allowbreak{}act(\allowbreak{}s))(\allowbreak{}one(\allowbreak{}scal(\allowbreak{}s)), x\_) = x\_) and forall(\allowbreak{}[x\_ in vec(\allowbreak{}s)], 0 <= (\allowbreak{}vnrm(\allowbreak{}s))(\allowbreak{}x\_)) and forall(\allowbreak{}[x\_ in vec(\allowbreak{}s)], (\allowbreak{}vnrm(\allowbreak{}s))(\allowbreak{}x\_) = 0 iff x\_ = vzero(\allowbreak{}s)) and forall(\allowbreak{}[r\_ in carr(\allowbreak{}scal(\allowbreak{}s)), x\_ in vec(\allowbreak{}s)], (\allowbreak{}vnrm(\allowbreak{}s))(\allowbreak{}(\allowbreak{}act(\allowbreak{}s))(\allowbreak{}r\_, x\_)) = magnitude(\allowbreak{}r\_) * (\allowbreak{}vnrm(\allowbreak{}s))(\allowbreak{}x\_)) and forall(\allowbreak{}[x\_ in vec(\allowbreak{}s), y\_ in vec(\allowbreak{}s)], (\allowbreak{}vnrm(\allowbreak{}s))(\allowbreak{}(\allowbreak{}vadd(\allowbreak{}s))(\allowbreak{}x\_, y\_)) <= (\allowbreak{}vnrm(\allowbreak{}s))(\allowbreak{}x\_) + (\allowbreak{}vnrm(\allowbreak{}s))(\allowbreak{}y\_)) and forall(\allowbreak{}[x\_ in vec(\allowbreak{}s), y\_ in vec(\allowbreak{}s), z\_ in vec(\allowbreak{}s)], (\allowbreak{}vmul(\allowbreak{}s))(\allowbreak{}x\_, (\allowbreak{}vadd(\allowbreak{}s))(\allowbreak{}y\_, z\_)) = (\allowbreak{}vadd(\allowbreak{}s))(\allowbreak{}(\allowbreak{}vmul(\allowbreak{}s))(\allowbreak{}x\_, y\_), (\allowbreak{}vmul(\allowbreak{}s))(\allowbreak{}x\_, z\_))) and forall(\allowbreak{}[x\_ in vec(\allowbreak{}s), y\_ in vec(\allowbreak{}s), z\_ in vec(\allowbreak{}s)], (\allowbreak{}vmul(\allowbreak{}s))(\allowbreak{}(\allowbreak{}vadd(\allowbreak{}s))(\allowbreak{}x\_, y\_), z\_) = (\allowbreak{}vadd(\allowbreak{}s))(\allowbreak{}(\allowbreak{}vmul(\allowbreak{}s))(\allowbreak{}x\_, z\_), (\allowbreak{}vmul(\allowbreak{}s))(\allowbreak{}y\_, z\_))) and forall(\allowbreak{}[r\_ in carr(\allowbreak{}scal(\allowbreak{}s)), x\_ in vec(\allowbreak{}s), y\_ in vec(\allowbreak{}s)], (\allowbreak{}vmul(\allowbreak{}s))(\allowbreak{}(\allowbreak{}act(\allowbreak{}s))(\allowbreak{}r\_, x\_), y\_) = (\allowbreak{}act(\allowbreak{}s))(\allowbreak{}r\_, (\allowbreak{}vmul(\allowbreak{}s))(\allowbreak{}x\_, y\_))) and forall(\allowbreak{}[r\_ in carr(\allowbreak{}scal(\allowbreak{}s)), x\_ in vec(\allowbreak{}s), y\_ in vec(\allowbreak{}s)], (\allowbreak{}vmul(\allowbreak{}s))(\allowbreak{}x\_, (\allowbreak{}act(\allowbreak{}s))(\allowbreak{}r\_, y\_)) = (\allowbreak{}act(\allowbreak{}s))(\allowbreak{}r\_, (\allowbreak{}vmul(\allowbreak{}s))(\allowbreak{}x\_, y\_))) and forall(\allowbreak{}[x\_ in vec(\allowbreak{}s), y\_ in vec(\allowbreak{}s)], (\allowbreak{}vnrm(\allowbreak{}s))(\allowbreak{}(\allowbreak{}vmul(\allowbreak{}s))(\allowbreak{}x\_, y\_)) <= (\allowbreak{}vnrm(\allowbreak{}s))(\allowbreak{}x\_) * (\allowbreak{}vnrm(\allowbreak{}s))(\allowbreak{}y\_)) and s in set)
\end{quote}
also: \texttt{is-complex-normed-algebra-is-set} \texttt{unitization-is-complex-normed-algebra} \texttt{complex-banach-algebra-is-complex-normed-algebra} \\
mentioned by 950 result(s)

\smallskip

\noindent\textbf{\texttt{is-complex-normed-vector-space(s)}}\quad\textit{predicate}\index{is-complex-normed-vector-space(s)@\texttt{is-complex-normed-vector-space(s)}}\\
reads: a complex normed vector space \\
a predicate \\
characterized by \texttt{is-complex-normed-vector-space}  \textit{(definitional)}:
\begin{quote}\ttfamily\footnotesize\raggedright\noindent
forall(\allowbreak{}[s], is-complex-normed-vector-space(\allowbreak{}s) iff length(\allowbreak{}s) = 8 and is-ring(\allowbreak{}scal(\allowbreak{}s)) and vec(\allowbreak{}s) in set and vadd(\allowbreak{}s) in fun(\allowbreak{}cartesian(\allowbreak{}vec(\allowbreak{}s), vec(\allowbreak{}s)), vec(\allowbreak{}s)) and vzero(\allowbreak{}s) in vec(\allowbreak{}s) and vneg(\allowbreak{}s) in fun(\allowbreak{}vec(\allowbreak{}s), vec(\allowbreak{}s)) and act(\allowbreak{}s) in fun(\allowbreak{}cartesian(\allowbreak{}carr(\allowbreak{}scal(\allowbreak{}s)), vec(\allowbreak{}s)), vec(\allowbreak{}s)) and vnrm(\allowbreak{}s) in fun(\allowbreak{}vec(\allowbreak{}s), rr) and jop(\allowbreak{}s) in fun(\allowbreak{}vec(\allowbreak{}s), vec(\allowbreak{}s)) and is-normed-vector-space(\allowbreak{}[scal(\allowbreak{}s), vec(\allowbreak{}s), vadd(\allowbreak{}s), vzero(\allowbreak{}s), vneg(\allowbreak{}s), act(\allowbreak{}s), vnrm(\allowbreak{}s)]) and is-real-linear-op(\allowbreak{}s, jop(\allowbreak{}s)) and is-square-minus-id(\allowbreak{}s, jop(\allowbreak{}s)) and forall(\allowbreak{}[cnc\_ in cc, cnx\_ in vec(\allowbreak{}s)], (\allowbreak{}vnrm(\allowbreak{}s))(\allowbreak{}(\allowbreak{}cact(\allowbreak{}s))(\allowbreak{}cnc\_, cnx\_)) = magnitude(\allowbreak{}cnc\_) * (\allowbreak{}vnrm(\allowbreak{}s))(\allowbreak{}cnx\_)) and s in set)
\end{quote}
also: \texttt{is-complex-normed-vector-space-is-set} \\
mentioned by 1190 result(s)

\smallskip

\noindent\textbf{\texttt{is-complex-null-set(omega, ca, nu, a)}}\quad\textit{predicate}\index{is-complex-null-set(omega, ca, nu, a)@\texttt{is-complex-null-set(omega, ca, nu, a)}}\\
reads: \$4 is a null set of the complex measure \$3 \\
defined by \texttt{is-complex-null-set} \\
characterized by \texttt{is-complex-null-set}  \textit{(definitional)}:
\begin{quote}\ttfamily\footnotesize\raggedright\noindent
forall(\allowbreak{}[omega, ca, nu, a], is-complex-null-set(\allowbreak{}omega, ca, nu, a) iff a in ca and forall(\allowbreak{}[cnb\_ in ca], cnb\_ subset a implies nu(\allowbreak{}cnb\_) = 0))
\end{quote}
mentioned by 3 result(s)

\smallskip

\noindent\textbf{\texttt{is-complex-pre-hilbert-space(s)}}\quad\textit{predicate}\index{is-complex-pre-hilbert-space(s)@\texttt{is-complex-pre-hilbert-space(s)}}\\
reads: a complex pre-Hilbert space \\
a predicate \\
characterized by \texttt{is-complex-pre-hilbert-space-def}  \textit{(definitional)}:
\begin{quote}\ttfamily\footnotesize\raggedright\noindent
forall(\allowbreak{}[s], is-complex-pre-hilbert-space(\allowbreak{}s) iff is-complex-semi-inner-product-space(\allowbreak{}s) and forall(\allowbreak{}[x\_ in vec(\allowbreak{}s)], (\allowbreak{}rip(\allowbreak{}s))(\allowbreak{}x\_, x\_) = 0 implies x\_ = vzero(\allowbreak{}s)))
\end{quote}
also: \texttt{is-complex-pre-hilbert-space-is-set} \\
mentioned by 65 result(s)

\smallskip

\noindent\textbf{\texttt{is-complex-semi-inner-product-space(s)}}\quad\textit{predicate}\index{is-complex-semi-inner-product-space(s)@\texttt{is-complex-semi-inner-product-space(s)}}\\
reads: a complex semi-inner-product space \\
a predicate \\
characterized by \texttt{is-complex-semi-inner-product-space}  \textit{(definitional)}:
\begin{quote}\ttfamily\footnotesize\raggedright\noindent
forall(\allowbreak{}[s], is-complex-semi-inner-product-space(\allowbreak{}s) iff length(\allowbreak{}s) = 8 and is-ring(\allowbreak{}scal(\allowbreak{}s)) and vec(\allowbreak{}s) in set and vadd(\allowbreak{}s) in fun(\allowbreak{}cartesian(\allowbreak{}vec(\allowbreak{}s), vec(\allowbreak{}s)), vec(\allowbreak{}s)) and vzero(\allowbreak{}s) in vec(\allowbreak{}s) and vneg(\allowbreak{}s) in fun(\allowbreak{}vec(\allowbreak{}s), vec(\allowbreak{}s)) and act(\allowbreak{}s) in fun(\allowbreak{}cartesian(\allowbreak{}carr(\allowbreak{}scal(\allowbreak{}s)), vec(\allowbreak{}s)), vec(\allowbreak{}s)) and rip(\allowbreak{}s) in fun(\allowbreak{}cartesian(\allowbreak{}vec(\allowbreak{}s), vec(\allowbreak{}s)), rr) and jop(\allowbreak{}s) in fun(\allowbreak{}vec(\allowbreak{}s), vec(\allowbreak{}s)) and is-semi-inner-product-space(\allowbreak{}[scal(\allowbreak{}s), vec(\allowbreak{}s), vadd(\allowbreak{}s), vzero(\allowbreak{}s), vneg(\allowbreak{}s), act(\allowbreak{}s), rip(\allowbreak{}s)]) and is-real-linear-op(\allowbreak{}s, jop(\allowbreak{}s)) and is-square-minus-id(\allowbreak{}s, jop(\allowbreak{}s)) and is-antisymmetric-op(\allowbreak{}s, jop(\allowbreak{}s)) and s in set)
\end{quote}
also: \texttt{is-complex-semi-inner-product-space-is-set} \\
mentioned by 1514 result(s)

\smallskip

\noindent\textbf{\texttt{is-concat-trace(sasp\_, sas\_)}}\quad\textit{predicate}\index{is-concat-trace(sasp_, sas_)@\texttt{is-concat-trace(sasp\_, sas\_)}}\\
reads: \$2 is a concatenation trace in the strand space \$1 \\
defined by \texttt{is-concat-trace} \\
characterized by \texttt{is-concat-trace}  \textit{(definitional)}:
\begin{quote}\ttfamily\footnotesize\raggedright\noindent
forall(\allowbreak{}[sasp\_, sas\_], is-concat-trace(\allowbreak{}sasp\_, sas\_) iff forsome(\allowbreak{}[sal\_ in tuples(\allowbreak{}msg-set(\allowbreak{}tokens(\allowbreak{}sasp\_), keys(\allowbreak{}sasp\_)))], length(\allowbreak{}sas\_) = succ(\allowbreak{}length(\allowbreak{}sal\_)) and forall(\allowbreak{}[sai\_ in nn], 1 <= sai\_ implies sai\_ <= length(\allowbreak{}sal\_) implies nth(\allowbreak{}sai\_, sas\_) = [received, nth(\allowbreak{}sai\_, sal\_)]) and nth(\allowbreak{}succ(\allowbreak{}length(\allowbreak{}sal\_)), sas\_) = [sent, lst(\allowbreak{}sal\_)]))
\end{quote}
mentioned by 6 result(s)

\smallskip

\noindent\textbf{\texttt{is-connected(cns\_, cna\_)}}\quad\textit{predicate}\index{is-connected(cns_, cna_)@\texttt{is-connected(cns\_, cna\_)}}\\
reads: \$2 is a connected subset of the metric space \$1 \\
defined by \texttt{is-connected} \\
characterized by \texttt{is-connected}  \textit{(definitional)}:
\begin{quote}\ttfamily\footnotesize\raggedright\noindent
forall(\allowbreak{}[cns\_, cna\_], is-connected(\allowbreak{}cns\_, cna\_) iff is-metric-space(\allowbreak{}cns\_) and cna\_ subset pts(\allowbreak{}cns\_) and forall(\allowbreak{}[cnu\_], is-open(\allowbreak{}cns\_, cnu\_) implies forall(\allowbreak{}[cnv\_], is-open(\allowbreak{}cns\_, cnv\_) implies cna\_ subset union(\allowbreak{}cnu\_, cnv\_) implies forall(\allowbreak{}[cnx\_ in cna\_], not(\allowbreak{}cnx\_ in cnu\_ and cnx\_ in cnv\_)) implies forsome(\allowbreak{}[cnx\_ in cna\_], cnx\_ in cnu\_) implies not(\allowbreak{}forsome(\allowbreak{}[cnx\_ in cna\_], cnx\_ in cnv\_)))))
\end{quote}
mentioned by 51 result(s)

\smallskip

\noindent\textbf{\texttt{is-connected-t(cts\_, cta\_)}}\quad\textit{predicate}\index{is-connected-t(cts_, cta_)@\texttt{is-connected-t(cts\_, cta\_)}}\\
reads: \$2 is a connected subset of \$1 \\
defined by \texttt{is-connected-t} \\
characterized by \texttt{is-connected-t}  \textit{(definitional)}:
\begin{quote}\ttfamily\footnotesize\raggedright\noindent
forall(\allowbreak{}[cts\_, cta\_], is-connected-t(\allowbreak{}cts\_, cta\_) iff cta\_ subset pts(\allowbreak{}cts\_) and forall(\allowbreak{}[ctu\_ in opens(\allowbreak{}cts\_), ctv\_ in opens(\allowbreak{}cts\_)], cta\_ subset union(\allowbreak{}ctu\_, ctv\_) implies forall(\allowbreak{}[ctx\_ in cta\_], not(\allowbreak{}ctx\_ in ctu\_ and ctx\_ in ctv\_)) implies forsome(\allowbreak{}[ctx\_ in cta\_], ctx\_ in ctu\_) implies not(\allowbreak{}forsome(\allowbreak{}[ctx\_ in cta\_], ctx\_ in ctv\_))))
\end{quote}
mentioned by 19 result(s)

\smallskip

\noindent\textbf{\texttt{is-cont-lin(m1, fam1, m2, fam2, tt)}}\quad\textit{predicate}\index{is-cont-lin(m1, fam1, m2, fam2, tt)@\texttt{is-cont-lin(m1, fam1, m2, fam2, tt)}}\\
reads: \$5 is a continuous linear map from (\$1, \$2) to (\$3, \$4) \\
defined by \texttt{is-cont-lin} \\
characterized by \texttt{is-cont-lin}  \textit{(definitional)}:
\begin{quote}\ttfamily\footnotesize\raggedright\noindent
forall(\allowbreak{}[m1, fam1, m2, fam2, tt], is-cont-lin(\allowbreak{}m1, fam1, m2, fam2, tt) iff is-linear-map(\allowbreak{}m1, m2, tt) and forall(\allowbreak{}[seq in fun(\allowbreak{}nn, vec(\allowbreak{}m1)), lim in vec(\allowbreak{}m1)], fr-conv(\allowbreak{}m1, fam1, seq, lim) implies fr-conv(\allowbreak{}m2, fam2, vnb-lambda(\allowbreak{}i, nn, tt(\allowbreak{}seq(\allowbreak{}i))), tt(\allowbreak{}lim))))
\end{quote}
mentioned by 18 result(s)

\smallskip

\noindent\textbf{\texttt{is-continuous(s, t, f)}}\quad\textit{predicate}\index{is-continuous(s, t, f)@\texttt{is-continuous(s, t, f)}}\\
reads: \$3 is continuous from \$1 to \$2 \\
defined by \texttt{is-continuous} \\
characterized by \texttt{is-continuous}  \textit{(definitional)}:
\begin{quote}\ttfamily\footnotesize\raggedright\noindent
forall(\allowbreak{}[s, t, f], is-continuous(\allowbreak{}s, t, f) iff is-metric-space(\allowbreak{}s) and is-metric-space(\allowbreak{}t) and f in fun(\allowbreak{}pts(\allowbreak{}s), pts(\allowbreak{}t)) and forall(\allowbreak{}[a in pts(\allowbreak{}s)], is-continuous-at(\allowbreak{}s, t, f, a)))
\end{quote}
also: \texttt{is-continuous-arrow-def} \texttt{metric-hom-is-continuous} \texttt{continuous-is-continuous-at} \texttt{uniformly-continuous-is-continuous} \\
mentioned by 91 result(s)

\smallskip

\noindent\textbf{\texttt{is-continuous-arrow(a, b, f)}}\quad\textit{predicate}\index{is-continuous-arrow(a, b, f)@\texttt{is-continuous-arrow(a, b, f)}}\\
reads: \$3 is a continuous map from \$1 to \$2 \\
a predicate \\
characterized by \texttt{is-continuous-arrow-def}  \textit{(definitional)}:
\begin{quote}\ttfamily\footnotesize\raggedright\noindent
forall(\allowbreak{}[a, b, f], is-continuous-arrow(\allowbreak{}a, b, f) iff is-metric-space(\allowbreak{}a) and is-metric-space(\allowbreak{}b) and f in fun(\allowbreak{}pts(\allowbreak{}a), pts(\allowbreak{}b)) and is-continuous(\allowbreak{}a, b, f))
\end{quote}
also: \texttt{uniformly-continuous-arrow-is-continuous-arrow} \\
mentioned by 7 result(s)

\smallskip

\noindent\textbf{\texttt{is-continuous-at(s, t, f, a)}}\quad\textit{predicate}\index{is-continuous-at(s, t, f, a)@\texttt{is-continuous-at(s, t, f, a)}}\\
reads: \$3 is continuous at \$4 \\
defined by \texttt{is-continuous-at} \\
characterized by \texttt{is-continuous-at}  \textit{(definitional)}:
\begin{quote}\ttfamily\footnotesize\raggedright\noindent
forall(\allowbreak{}[s, t, f, a], is-continuous-at(\allowbreak{}s, t, f, a) iff is-metric-space(\allowbreak{}s) and is-metric-space(\allowbreak{}t) and f in fun(\allowbreak{}pts(\allowbreak{}s), pts(\allowbreak{}t)) and a in pts(\allowbreak{}s) and forall(\allowbreak{}[eps], pos-rr(\allowbreak{}eps) implies forsome(\allowbreak{}[delta], pos-rr(\allowbreak{}delta) and forall(\allowbreak{}[b in pts(\allowbreak{}s)], (\allowbreak{}dist(\allowbreak{}s))(\allowbreak{}a, b) <= delta implies (\allowbreak{}dist(\allowbreak{}t))(\allowbreak{}f(\allowbreak{}a), f(\allowbreak{}b)) <= eps))))
\end{quote}
also: \texttt{continuous-is-continuous-at} \\
mentioned by 380 result(s)

\smallskip

\noindent\textbf{\texttt{is-continuous-on(f, d)}}\quad\textit{predicate}\index{is-continuous-on(f, d)@\texttt{is-continuous-on(f, d)}}\\
reads: \$1 is continuous on \$2 \\
defined by \texttt{is-continuous-on} \\
characterized by \texttt{is-continuous-on}  \textit{(definitional)}:
\begin{quote}\ttfamily\footnotesize\raggedright\noindent
forall(\allowbreak{}[f, d], is-continuous-on(\allowbreak{}f, d) iff f in fun(\allowbreak{}d, rr) and forall(\allowbreak{}[icp\_ in d], is-continuous-at(\allowbreak{}subspace-ms(\allowbreak{}rr-ms, d), rr-ms, f, icp\_)))
\end{quote}
mentioned by 57 result(s)

\smallskip

\noindent\textbf{\texttt{is-conull(omega, ca, mu, cs)}}\quad\textit{predicate}\index{is-conull(omega, ca, mu, cs)@\texttt{is-conull(omega, ca, mu, cs)}}\\
reads: \$4 is a conull set of the measure \$3 on \$2 \\
defined by \texttt{is-conull} \\
characterized by \texttt{is-conull}  \textit{(definitional)}:
\begin{quote}\ttfamily\footnotesize\raggedright\noindent
forall(\allowbreak{}[omega, ca, mu, cs], is-conull(\allowbreak{}omega, ca, mu, cs) iff cs in ca and mu(\allowbreak{}complement-in(\allowbreak{}omega, cs)) = 0)
\end{quote}
mentioned by 4 result(s)

\smallskip

\noindent\textbf{\texttt{is-convex-cc(cvu\_)}}\quad\textit{predicate}\index{is-convex-cc(cvu_)@\texttt{is-convex-cc(cvu\_)}}\\
reads: \$1 is a convex subset of CC \\
defined by \texttt{is-convex-cc} \\
characterized by \texttt{is-convex-cc}  \textit{(definitional)}:
\begin{quote}\ttfamily\footnotesize\raggedright\noindent
forall(\allowbreak{}[cvu\_], is-convex-cc(\allowbreak{}cvu\_) iff cvu\_ subset cc and forall(\allowbreak{}[cvx\_ in cvu\_, cvy\_ in cvu\_, cvt\_ in ccint(\allowbreak{}0, 1)], cvx\_ + cvt\_ * (\allowbreak{}cvy\_ - cvx\_) in cvu\_))
\end{quote}
mentioned by 42 result(s)

\smallskip

\noindent\textbf{\texttt{is-convex-functional(m, p)}}\quad\textit{predicate}\index{is-convex-functional(m, p)@\texttt{is-convex-functional(m, p)}}\\
reads: \$2 is a convex functional on \$1 \\
defined by \texttt{is-convex-functional} \\
characterized by \texttt{is-convex-functional}  \textit{(definitional)}:
\begin{quote}\ttfamily\footnotesize\raggedright\noindent
forall(\allowbreak{}[m, p], is-convex-functional(\allowbreak{}m, p) iff p in fun(\allowbreak{}vec(\allowbreak{}m), rr) and forall(\allowbreak{}[cfx\_ in vec(\allowbreak{}m)], 0 <= p(\allowbreak{}cfx\_)) and forall(\allowbreak{}[cfx\_ in vec(\allowbreak{}m), cfy\_ in vec(\allowbreak{}m), cft\_ in rr], 0 <= cft\_ implies cft\_ <= 1 implies p(\allowbreak{}(\allowbreak{}vadd(\allowbreak{}m))(\allowbreak{}(\allowbreak{}act(\allowbreak{}m))(\allowbreak{}cft\_, cfx\_), (\allowbreak{}act(\allowbreak{}m))(\allowbreak{}1 - cft\_, cfy\_))) <= cft\_ * p(\allowbreak{}cfx\_) + (\allowbreak{}1 - cft\_) * p(\allowbreak{}cfy\_)))
\end{quote}
mentioned by 20 result(s)

\smallskip

\noindent\textbf{\texttt{is-convex-rr(cvra\_)}}\quad\textit{predicate}\index{is-convex-rr(cvra_)@\texttt{is-convex-rr(cvra\_)}}\\
reads: \$1 is a convex subset of RR \\
defined by \texttt{is-convex-rr} \\
characterized by \texttt{is-convex-rr}  \textit{(definitional)}:
\begin{quote}\ttfamily\footnotesize\raggedright\noindent
forall(\allowbreak{}[cvra\_], is-convex-rr(\allowbreak{}cvra\_) iff cvra\_ subset rr and forall(\allowbreak{}[cvrx\_ in cvra\_, cvry\_ in cvra\_, cvrt\_ in ccint(\allowbreak{}0, 1)], cvrx\_ + cvrt\_ * (\allowbreak{}cvry\_ - cvrx\_) in cvra\_))
\end{quote}
also: \texttt{is-convex-rr-unfold} \\
mentioned by 4 result(s)

\smallskip

\noindent\textbf{\texttt{is-convex-subset(v, cs)}}\quad\textit{predicate}\index{is-convex-subset(v, cs)@\texttt{is-convex-subset(v, cs)}}\\
reads: \$2 is a convex set of vectors of \$1 \\
defined by \texttt{is-convex-subset} \\
characterized by \texttt{is-convex-subset}  \textit{(definitional)}:
\begin{quote}\ttfamily\footnotesize\raggedright\noindent
forall(\allowbreak{}[v, cs], is-convex-subset(\allowbreak{}v, cs) iff cs subset vec(\allowbreak{}v) and forall(\allowbreak{}[sxa\_ in cs, sxb\_ in cs, sxt\_ in rr], 0 <= sxt\_ implies sxt\_ <= 1 implies (\allowbreak{}vadd(\allowbreak{}v))(\allowbreak{}(\allowbreak{}act(\allowbreak{}v))(\allowbreak{}1 - sxt\_, sxa\_), (\allowbreak{}act(\allowbreak{}v))(\allowbreak{}sxt\_, sxb\_)) in cs))
\end{quote}
mentioned by 74 result(s)

\smallskip

\noindent\textbf{\texttt{is-countable(rpc\_)}}\quad\textit{predicate}\index{is-countable(rpc_)@\texttt{is-countable(rpc\_)}}\\
reads: \$1 is countable \\
defined by \texttt{is-countable} \\
characterized by \texttt{is-countable}  \textit{(definitional)}:
\begin{quote}\ttfamily\footnotesize\raggedright\noindent
forall(\allowbreak{}[rpc\_], is-countable(\allowbreak{}rpc\_) iff rpc\_ in set and (\allowbreak{}rpc\_ = empty-set or forsome(\allowbreak{}[rpen\_ in fun(\allowbreak{}nn, rpc\_)], forall(\allowbreak{}[rpy\_ in rpc\_], forsome(\allowbreak{}[rpn\_ in nn], rpy\_ = rpen\_(\allowbreak{}rpn\_))))))
\end{quote}
mentioned by 46 result(s)

\smallskip

\noindent\textbf{\texttt{is-countable-pseudometric-family(fam, ground)}}\quad\textit{predicate}\index{is-countable-pseudometric-family(fam, ground)@\texttt{is-countable-pseudometric-family(fam, ground)}}\\
reads: \$1 is a countable family of pseudometrics on \$2 \\
defined by \texttt{is-countable-pseudometric-family} \\
characterized by \texttt{is-countable-pseudometric-family}  \textit{(definitional)}:
\begin{quote}\ttfamily\footnotesize\raggedright\noindent
forall(\allowbreak{}[fam, ground], is-countable-pseudometric-family(\allowbreak{}fam, ground) iff forall(\allowbreak{}[n\_ in nn], is-pseudometric-space(\allowbreak{}fam(\allowbreak{}n\_)) and pts(\allowbreak{}fam(\allowbreak{}n\_)) == ground))
\end{quote}
mentioned by 17 result(s)

\smallskip

\noindent\textbf{\texttt{is-countably-separated(ys, cb)}}\quad\textit{predicate}\index{is-countably-separated(ys, cb)@\texttt{is-countably-separated(ys, cb)}}\\
reads: the measurable space (\$1, \$2) is countably separated \\
defined by \texttt{is-countably-separated} \\
characterized by \texttt{is-countably-separated}  \textit{(definitional)}:
\begin{quote}\ttfamily\footnotesize\raggedright\noindent
forall(\allowbreak{}[ys, cb], is-countably-separated(\allowbreak{}ys, cb) iff is-sigma-algebra(\allowbreak{}ys, cb) and forsome(\allowbreak{}[csf\_ in fun(\allowbreak{}nn, cb)], forall(\allowbreak{}[csx\_ in ys, csy\_ in ys], not(\allowbreak{}csx\_ = csy\_) implies forsome(\allowbreak{}[csk\_ in nn], csx\_ in csf\_(\allowbreak{}csk\_) and not(\allowbreak{}csy\_ in csf\_(\allowbreak{}csk\_)) or csy\_ in csf\_(\allowbreak{}csk\_) and not(\allowbreak{}csx\_ in csf\_(\allowbreak{}csk\_))))))
\end{quote}
mentioned by 11 result(s)

\smallskip

\noindent\textbf{\texttt{is-cyclic-for(vnv\_, vns\_, vnx\_)}}\quad\textit{predicate}\index{is-cyclic-for(vnv_, vns_, vnx_)@\texttt{is-cyclic-for(vnv\_, vns\_, vnx\_)}}\\
reads: \$3 is a cyclic vector for \$2 \\
defined by \texttt{is-cyclic-for} \\
characterized by \texttt{is-cyclic-for}  \textit{(definitional)}:
\begin{quote}\ttfamily\footnotesize\raggedright\noindent
forall(\allowbreak{}[vnv\_, vns\_, vnx\_], is-cyclic-for(\allowbreak{}vnv\_, vns\_, vnx\_) iff vnx\_ in vec(\allowbreak{}vnv\_) and forall(\allowbreak{}[vny\_ in vec(\allowbreak{}vnv\_), vne\_], pos-rr(\allowbreak{}vne\_) implies forsome(\allowbreak{}[vna\_ in vns\_], ip-norm(\allowbreak{}vnv\_, (\allowbreak{}vadd(\allowbreak{}vnv\_))(\allowbreak{}vny\_, (\allowbreak{}vneg(\allowbreak{}vnv\_))(\allowbreak{}vna\_(\allowbreak{}vnx\_)))) < vne\_)))
\end{quote}
mentioned by 10 result(s)

\smallskip

\noindent\textbf{\texttt{is-cyclic-vector(gcv-md\_, gcv-t\_, gcv-v\_)}}\quad\textit{predicate}\index{is-cyclic-vector(gcv-md_, gcv-t_, gcv-v_)@\texttt{is-cyclic-vector(gcv-md\_, gcv-t\_, gcv-v\_)}}\\
reads: \$3 is a cyclic vector for \$2 \\
defined by \texttt{is-cyclic-vector} \\
characterized by \texttt{is-cyclic-vector}  \textit{(definitional)}:
\begin{quote}\ttfamily\footnotesize\raggedright\noindent
forall(\allowbreak{}[gcv-md\_, gcv-t\_, gcv-v\_], is-cyclic-vector(\allowbreak{}gcv-md\_, gcv-t\_, gcv-v\_) iff gcv-v\_ in vec(\allowbreak{}gcv-md\_) and forall(\allowbreak{}[gcvw\_ in vec(\allowbreak{}gcv-md\_)], forsome(\allowbreak{}[gcvq\_ in carr(\allowbreak{}poly(\allowbreak{}scal(\allowbreak{}gcv-md\_)))], gcvw\_ = poly-op(\allowbreak{}gcv-md\_, gcv-t\_, gcvq\_, gcv-v\_))))
\end{quote}
mentioned by 7 result(s)

\smallskip

\noindent\textbf{\texttt{is-cyclic-vector-rep(g4a\_, g4v\_, g4p\_, g4e\_)}}\quad\textit{predicate}\index{is-cyclic-vector-rep(g4a_, g4v_, g4p_, g4e_)@\texttt{is-cyclic-vector-rep(g4a\_, g4v\_, g4p\_, g4e\_)}}\\
reads: \$4 is a cyclic vector of \$3 \\
defined by \texttt{is-cyclic-vector-rep} \\
characterized by \texttt{is-cyclic-vector-rep}  \textit{(definitional)}:
\begin{quote}\ttfamily\footnotesize\raggedright\noindent
forall(\allowbreak{}[g4a\_, g4v\_, g4p\_, g4e\_], is-cyclic-vector-rep(\allowbreak{}g4a\_, g4v\_, g4p\_, g4e\_) iff g4e\_ in vec(\allowbreak{}g4v\_) and closure(\allowbreak{}ip-metric-space(\allowbreak{}g4v\_), rep-orbit(\allowbreak{}g4a\_, g4v\_, g4p\_, g4e\_)) = vec(\allowbreak{}g4v\_))
\end{quote}
mentioned by 2 result(s)

\smallskip

\noindent\textbf{\texttt{is-decrypt-trace(sasp\_, sas\_)}}\quad\textit{predicate}\index{is-decrypt-trace(sasp_, sas_)@\texttt{is-decrypt-trace(sasp\_, sas\_)}}\\
reads: \$2 is a decryption trace in the strand space \$1 \\
defined by \texttt{is-decrypt-trace} \\
characterized by \texttt{is-decrypt-trace}  \textit{(definitional)}:
\begin{quote}\ttfamily\footnotesize\raggedright\noindent
forall(\allowbreak{}[sasp\_, sas\_], is-decrypt-trace(\allowbreak{}sasp\_, sas\_) iff forsome(\allowbreak{}[sak\_ in keys(\allowbreak{}sasp\_), sag\_ in msg-set(\allowbreak{}tokens(\allowbreak{}sasp\_), keys(\allowbreak{}sasp\_))], sas\_ = [[received, (\allowbreak{}kinv(\allowbreak{}sasp\_))(\allowbreak{}sak\_)], [received, encr(\allowbreak{}sag\_, sak\_)], [sent, sag\_]]))
\end{quote}
mentioned by 6 result(s)

\smallskip

\noindent\textbf{\texttt{is-degree-function(gfs\_, gfn\_)}}\quad\textit{predicate}\index{is-degree-function(gfs_, gfn_)@\texttt{is-degree-function(gfs\_, gfn\_)}}\\
reads: \$2 is a degree function on \$1 \\
defined by \texttt{is-degree-function} \\
characterized by \texttt{is-degree-function}  \textit{(definitional)}:
\begin{quote}\ttfamily\footnotesize\raggedright\noindent
forall(\allowbreak{}[gfs\_, gfn\_], is-degree-function(\allowbreak{}gfs\_, gfn\_) iff is-integral-domain(\allowbreak{}gfs\_) and gfn\_ in fun(\allowbreak{}carr(\allowbreak{}gfs\_), nn) and forall(\allowbreak{}[gfa\_ in carr(\allowbreak{}gfs\_), gfb\_ in carr(\allowbreak{}gfs\_)], not(\allowbreak{}gfa\_ = zero(\allowbreak{}gfs\_)) implies not(\allowbreak{}gfb\_ = zero(\allowbreak{}gfs\_)) implies gfn\_(\allowbreak{}gfa\_) <= gfn\_(\allowbreak{}(\allowbreak{}mul(\allowbreak{}gfs\_))(\allowbreak{}gfa\_, gfb\_))) and forall(\allowbreak{}[gfa\_ in carr(\allowbreak{}gfs\_), gfb\_ in carr(\allowbreak{}gfs\_)], not(\allowbreak{}gfa\_ = zero(\allowbreak{}gfs\_)) implies forsome(\allowbreak{}[gfq\_ in carr(\allowbreak{}gfs\_)], gfb\_ = (\allowbreak{}mul(\allowbreak{}gfs\_))(\allowbreak{}gfa\_, gfq\_) or gfn\_(\allowbreak{}(\allowbreak{}add(\allowbreak{}gfs\_))(\allowbreak{}gfb\_, (\allowbreak{}neg(\allowbreak{}gfs\_))(\allowbreak{}(\allowbreak{}mul(\allowbreak{}gfs\_))(\allowbreak{}gfa\_, gfq\_)))) < gfn\_(\allowbreak{}gfa\_))))
\end{quote}
also: \texttt{poly-degree-is-degree-function} \\
mentioned by 8 result(s)

\smallskip

\noindent\textbf{\texttt{is-dense(dns\_, dna\_)}}\quad\textit{predicate}\index{is-dense(dns_, dna_)@\texttt{is-dense(dns\_, dna\_)}}\\
reads: \$2 is dense in the metric space \$1 \\
defined by \texttt{is-dense} \\
characterized by \texttt{is-dense}  \textit{(definitional)}:
\begin{quote}\ttfamily\footnotesize\raggedright\noindent
forall(\allowbreak{}[dns\_, dna\_], is-dense(\allowbreak{}dns\_, dna\_) iff is-metric-space(\allowbreak{}dns\_) and dna\_ subset pts(\allowbreak{}dns\_) and pts(\allowbreak{}dns\_) subset closure(\allowbreak{}dns\_, dna\_))
\end{quote}
mentioned by 24 result(s)

\smallskip

\noindent\textbf{\texttt{is-dense-seq(s, dseq)}}\quad\textit{predicate}\index{is-dense-seq(s, dseq)@\texttt{is-dense-seq(s, dseq)}}\\
reads: \$2 is a dense sequence in \$1 \\
defined by \texttt{is-dense-seq} \\
characterized by \texttt{is-dense-seq}  \textit{(definitional)}:
\begin{quote}\ttfamily\footnotesize\raggedright\noindent
forall(\allowbreak{}[s, dseq], is-dense-seq(\allowbreak{}s, dseq) iff is-metric-space(\allowbreak{}s) and dseq in fun(\allowbreak{}nn, pts(\allowbreak{}s)) and forall(\allowbreak{}[x in pts(\allowbreak{}s), eps], pos-rr(\allowbreak{}eps) implies forsome(\allowbreak{}[m in nn], (\allowbreak{}dist(\allowbreak{}s))(\allowbreak{}x, dseq(\allowbreak{}m)) < eps)))
\end{quote}
mentioned by 45 result(s)

\smallskip

\noindent\textbf{\texttt{is-diagonal(a, m, n, d)}}\quad\textit{predicate}\index{is-diagonal(a, m, n, d)@\texttt{is-diagonal(a, m, n, d)}}\\
reads: \$4 is a diagonal \$2-by-\$3 matrix over \$1 \\
defined by \texttt{is-diagonal} \\
characterized by \texttt{is-diagonal}  \textit{(definitional)}:
\begin{quote}\ttfamily\footnotesize\raggedright\noindent
forall(\allowbreak{}[a, m, n, d], is-diagonal(\allowbreak{}a, m, n, d) iff forall(\allowbreak{}[i, j], i in interval(\allowbreak{}1, m) implies j in interval(\allowbreak{}1, n) implies not(\allowbreak{}i = j) implies entry(\allowbreak{}d, i, j) = zero(\allowbreak{}a)))
\end{quote}
also: \texttt{border-is-diagonal-uba-as-ring} \texttt{border-is-diagonal-uba-as-ring-uiba-as-uba} \texttt{border-is-diagonal} \\
mentioned by 54 result(s)

\smallskip

\noindent\textbf{\texttt{is-diff-at(f, a, l)}}\quad\textit{predicate}\index{is-diff-at(f, a, l)@\texttt{is-diff-at(f, a, l)}}\\
reads: \$1 is differentiable at \$2, with derivative \$3 \\
defined by \texttt{is-diff-at} \\
characterized by \texttt{is-diff-at}  \textit{(definitional)}:
\begin{quote}\ttfamily\footnotesize\raggedright\noindent
forall(\allowbreak{}[f, a, l], is-diff-at(\allowbreak{}f, a, l) iff f in fun(\allowbreak{}rr, rr) and a in rr and l in rr and forsome(\allowbreak{}[phi in fun(\allowbreak{}rr, rr)], is-continuous-at(\allowbreak{}rr-ms, rr-ms, phi, a) and phi(\allowbreak{}a) = l and forall(\allowbreak{}[x in rr], f(\allowbreak{}x) - f(\allowbreak{}a) = phi(\allowbreak{}x) * (\allowbreak{}x - a))))
\end{quote}
also: \texttt{deriv-of-is-diff-at} \\
mentioned by 113 result(s)

\smallskip

\noindent\textbf{\texttt{is-diff-at-v(m, f, a, l)}}\quad\textit{predicate}\index{is-diff-at-v(m, f, a, l)@\texttt{is-diff-at-v(m, f, a, l)}}\\
reads: \$2 is differentiable at \$3 with derivative \$4, as a curve in \$1 \\
defined by \texttt{is-diff-at-v} \\
characterized by \texttt{is-diff-at-v}  \textit{(definitional)}:
\begin{quote}\ttfamily\footnotesize\raggedright\noindent
forall(\allowbreak{}[m, f, a, l], is-diff-at-v(\allowbreak{}m, f, a, l) iff is-normed-vector-space(\allowbreak{}m) and f in fun(\allowbreak{}rr, vec(\allowbreak{}m)) and a in rr and l in vec(\allowbreak{}m) and forsome(\allowbreak{}[phi in fun(\allowbreak{}rr, vec(\allowbreak{}m))], is-continuous-at(\allowbreak{}rr-ms, nvs-metric-space(\allowbreak{}m), phi, a) and phi(\allowbreak{}a) = l and forall(\allowbreak{}[x\_ in rr], (\allowbreak{}vadd(\allowbreak{}m))(\allowbreak{}f(\allowbreak{}x\_), (\allowbreak{}vneg(\allowbreak{}m))(\allowbreak{}f(\allowbreak{}a))) = (\allowbreak{}act(\allowbreak{}m))(\allowbreak{}x\_ - a, phi(\allowbreak{}x\_)))))
\end{quote}
mentioned by 10 result(s)

\smallskip

\noindent\textbf{\texttt{is-diff-on(k, u, f, a, l)}}\quad\textit{predicate}\index{is-diff-on(k, u, f, a, l)@\texttt{is-diff-on(k, u, f, a, l)}}\\
reads: \$3, on \$2, is differentiable at \$4 over \$1, with derivative \$5 \\
defined by \texttt{is-diff-on} \\
characterized by \texttt{is-diff-on}  \textit{(definitional)}:
\begin{quote}\ttfamily\footnotesize\raggedright\noindent
forall(\allowbreak{}[k, u, f, a, l], is-diff-on(\allowbreak{}k, u, f, a, l) iff is-normed-field(\allowbreak{}k) and is-open(\allowbreak{}nf-metric-space(\allowbreak{}k), u) and f in fun(\allowbreak{}u, carr(\allowbreak{}k)) and a in u and l in carr(\allowbreak{}k) and forsome(\allowbreak{}[phi in fun(\allowbreak{}u, carr(\allowbreak{}k))], is-continuous-at(\allowbreak{}subspace-ms(\allowbreak{}nf-metric-space(\allowbreak{}k), u), nf-metric-space(\allowbreak{}k), phi, a) and phi(\allowbreak{}a) = l and forall(\allowbreak{}[x in u], (\allowbreak{}add(\allowbreak{}k))(\allowbreak{}f(\allowbreak{}x), (\allowbreak{}neg(\allowbreak{}k))(\allowbreak{}f(\allowbreak{}a))) = (\allowbreak{}mul(\allowbreak{}k))(\allowbreak{}phi(\allowbreak{}x), (\allowbreak{}add(\allowbreak{}k))(\allowbreak{}x, (\allowbreak{}neg(\allowbreak{}k))(\allowbreak{}a))))))
\end{quote}
also: \texttt{deriv-on-of-is-diff-on} \texttt{cc-deriv-on-of-is-diff-on} \texttt{nth-deriv-on-of-is-diff-on} \\
mentioned by 109 result(s)

\smallskip

\noindent\textbf{\texttt{is-dir-diff-at(m, f, a, eta, dl)}}\quad\textit{predicate}\index{is-dir-diff-at(m, f, a, eta, dl)@\texttt{is-dir-diff-at(m, f, a, eta, dl)}}\\
reads: \$2 has directional derivative \$5 at \$3 in the direction \$4, in \$1 \\
defined by \texttt{is-dir-diff-at} \\
characterized by \texttt{is-dir-diff-at}  \textit{(definitional)}:
\begin{quote}\ttfamily\footnotesize\raggedright\noindent
forall(\allowbreak{}[m, f, a, eta, dl], is-dir-diff-at(\allowbreak{}m, f, a, eta, dl) iff is-diff-at(\allowbreak{}seg-curve(\allowbreak{}m, f, a, eta), 0, dl))
\end{quote}
mentioned by 19 result(s)

\smallskip

\noindent\textbf{\texttt{is-directed-set(s)}}\quad\textit{predicate}\index{is-directed-set(s)@\texttt{is-directed-set(s)}}\\
reads: a directed set \\
a predicate \\
characterized by \texttt{is-directed-set}  \textit{(definitional)}:
\begin{quote}\ttfamily\footnotesize\raggedright\noindent
forall(\allowbreak{}[s], is-directed-set(\allowbreak{}s) iff length(\allowbreak{}s) = 2 and pts(\allowbreak{}s) in set and dleq(\allowbreak{}s) in power(\allowbreak{}cartesian(\allowbreak{}pts(\allowbreak{}s), pts(\allowbreak{}s))) and forsome(\allowbreak{}[dsn\_], dsn\_ in pts(\allowbreak{}s)) and forall(\allowbreak{}[dsa\_ in pts(\allowbreak{}s)], [dsa\_, dsa\_] in dleq(\allowbreak{}s)) and forall(\allowbreak{}[dsa\_ in pts(\allowbreak{}s), dsb\_ in pts(\allowbreak{}s), dsc\_ in pts(\allowbreak{}s)], [dsa\_, dsb\_] in dleq(\allowbreak{}s) implies [dsb\_, dsc\_] in dleq(\allowbreak{}s) implies [dsa\_, dsc\_] in dleq(\allowbreak{}s)) and forall(\allowbreak{}[dsa\_ in pts(\allowbreak{}s), dsb\_ in pts(\allowbreak{}s)], forsome(\allowbreak{}[dsc\_ in pts(\allowbreak{}s)], [dsa\_, dsc\_] in dleq(\allowbreak{}s) and [dsb\_, dsc\_] in dleq(\allowbreak{}s))) and s in set)
\end{quote}
also: \texttt{is-directed-set-is-set} \texttt{nn-directed-is-directed-set} \\
mentioned by 28 result(s)

\smallskip

\noindent\textbf{\texttt{is-discrete-ms(dms\_)}}\quad\textit{predicate}\index{is-discrete-ms(dms_)@\texttt{is-discrete-ms(dms\_)}}\\
reads: \$1 is discrete \\
defined by \texttt{is-discrete-ms} \\
characterized by \texttt{is-discrete-ms}  \textit{(definitional)}:
\begin{quote}\ttfamily\footnotesize\raggedright\noindent
forall(\allowbreak{}[dms\_], is-discrete-ms(\allowbreak{}dms\_) iff forall(\allowbreak{}[dmx\_ in pts(\allowbreak{}dms\_)], forsome(\allowbreak{}[dmr\_], pos-rr(\allowbreak{}dmr\_) and forall(\allowbreak{}[dmy\_ in ball(\allowbreak{}dms\_, dmx\_, dmr\_)], dmy\_ = dmx\_))))
\end{quote}
mentioned by 5 result(s)

\smallskip

\noindent\textbf{\texttt{is-distributive(addop, mulop, crr)}}\quad\textit{predicate}\index{is-distributive(addop, mulop, crr)@\texttt{is-distributive(addop, mulop, crr)}}\\
reads: \$2 distributes over \$1 on \$3 \\
defined by \texttt{is-distributive} \\
characterized by \texttt{is-distributive}  \textit{(definitional)}:
\begin{quote}\ttfamily\footnotesize\raggedright\noindent
forall(\allowbreak{}[addop, mulop, crr], is-distributive(\allowbreak{}addop, mulop, crr) iff forall(\allowbreak{}[u in crr, v in crr, w in crr], mulop(\allowbreak{}u, addop(\allowbreak{}v, w)) = addop(\allowbreak{}mulop(\allowbreak{}u, v), mulop(\allowbreak{}u, w)) and mulop(\allowbreak{}addop(\allowbreak{}u, v), w) = addop(\allowbreak{}mulop(\allowbreak{}u, w), mulop(\allowbreak{}v, w))))
\end{quote}
mentioned by 10 result(s)

\smallskip

\noindent\textbf{\texttt{is-doubling(mhf\_)}}\quad\textit{predicate}\index{is-doubling(mhf_)@\texttt{is-doubling(mhf\_)}}\\
reads: \$1 is a doubling function \\
defined by \texttt{is-doubling} \\
characterized by \texttt{is-doubling}  \textit{(definitional)}:
\begin{quote}\ttfamily\footnotesize\raggedright\noindent
forall(\allowbreak{}[mhf\_], is-doubling(\allowbreak{}mhf\_) iff mhf\_ in fun(\allowbreak{}nnreal, nnreal) and mhf\_(\allowbreak{}0) = 0 and forall(\allowbreak{}[mhr\_ in nnreal, mhq\_ in nnreal], mhr\_ <= mhq\_ implies mhf\_(\allowbreak{}mhr\_) <= mhf\_(\allowbreak{}mhq\_)) and forsome(\allowbreak{}[mhp\_], pos-rr(\allowbreak{}mhp\_) and forall(\allowbreak{}[mhr\_ in rr], 0 < mhr\_ implies mhf\_(\allowbreak{}2 * mhr\_) <= mhp\_ * mhf\_(\allowbreak{}mhr\_))))
\end{quote}
mentioned by 13 result(s)

\smallskip

\noindent\textbf{\texttt{is-dualizing-character(c5k\_, c5c\_)}}\quad\textit{predicate}\index{is-dualizing-character(c5k_, c5c_)@\texttt{is-dualizing-character(c5k\_, c5c\_)}}\\
reads: \$2 is a non-trivial character of \$1 of which every character of \$1 is a dilate \\
defined by \texttt{is-dualizing-character} \\
characterized by \texttt{is-dualizing-character}  \textit{(definitional)}:
\begin{quote}\ttfamily\footnotesize\raggedright\noindent
forall(\allowbreak{}[c5k\_, c5c\_], is-dualizing-character(\allowbreak{}c5k\_, c5c\_) iff is-nontrivial-character(\allowbreak{}c5k\_, c5c\_) and forall(\allowbreak{}[c5t\_], is-add-character(\allowbreak{}c5k\_, c5t\_) implies forsome(\allowbreak{}[c5a\_ in carr(\allowbreak{}c5k\_)], c5t\_ = duality-map(\allowbreak{}c5k\_, c5c\_, c5a\_))))
\end{quote}
mentioned by 7 result(s)

\smallskip

\noindent\textbf{\texttt{is-dynkin-system(omega, dd)}}\quad\textit{predicate}\index{is-dynkin-system(omega, dd)@\texttt{is-dynkin-system(omega, dd)}}\\
reads: \$2 is a Dynkin system of subsets of \$1 \\
defined by \texttt{is-dynkin-system} \\
characterized by \texttt{is-dynkin-system}  \textit{(definitional)}:
\begin{quote}\ttfamily\footnotesize\raggedright\noindent
forall(\allowbreak{}[omega, dd], is-dynkin-system(\allowbreak{}omega, dd) iff omega in set and dd subset power(\allowbreak{}omega) and omega in dd and forall(\allowbreak{}[dya\_ in dd, dyb\_ in dd], dya\_ subset dyb\_ implies complement-in(\allowbreak{}dyb\_, dya\_) in dd) and forall(\allowbreak{}[dyf\_ in fun(\allowbreak{}nn, dd)], forall(\allowbreak{}[dyn\_ in nn], dyf\_(\allowbreak{}dyn\_) subset dyf\_(\allowbreak{}succ(\allowbreak{}dyn\_))) implies big-union(\allowbreak{}dyk\_, nn, dyf\_(\allowbreak{}dyk\_)) in dd))
\end{quote}
mentioned by 11 result(s)

\smallskip

\noindent\textbf{\texttt{is-e-bicharacter(c6vk\_, c6vm\_, c6vw\_)}}\quad\textit{predicate}\index{is-e-bicharacter(c6vk_, c6vm_, c6vw_)@\texttt{is-e-bicharacter(c6vk\_, c6vm\_, c6vw\_)}}\\
reads: \$3 is a bicharacter on the vector space \$2 over \$1 \\
defined by \texttt{is-e-bicharacter} \\
characterized by \texttt{is-e-bicharacter}  \textit{(definitional)}:
\begin{quote}\ttfamily\footnotesize\raggedright\noindent
forall(\allowbreak{}[c6vk\_, c6vm\_, c6vw\_], is-e-bicharacter(\allowbreak{}c6vk\_, c6vm\_, c6vw\_) iff c6vw\_ in fun(\allowbreak{}cartesian(\allowbreak{}vec(\allowbreak{}c6vm\_), vec(\allowbreak{}c6vm\_)), unit-circle) and forall(\allowbreak{}[c6vx\_ in vec(\allowbreak{}c6vm\_), c6vy\_ in vec(\allowbreak{}c6vm\_), c6vz\_ in vec(\allowbreak{}c6vm\_)], c6vw\_(\allowbreak{}(\allowbreak{}vadd(\allowbreak{}c6vm\_))(\allowbreak{}c6vx\_, c6vy\_), c6vz\_) = c6vw\_(\allowbreak{}c6vx\_, c6vz\_) * c6vw\_(\allowbreak{}c6vy\_, c6vz\_) and c6vw\_(\allowbreak{}c6vz\_, (\allowbreak{}vadd(\allowbreak{}c6vm\_))(\allowbreak{}c6vx\_, c6vy\_)) = c6vw\_(\allowbreak{}c6vz\_, c6vx\_) * c6vw\_(\allowbreak{}c6vz\_, c6vy\_)) and forall(\allowbreak{}[c6vx\_ in vec(\allowbreak{}c6vm\_), c6vy\_ in vec(\allowbreak{}c6vm\_)], is-add-character(\allowbreak{}c6vk\_, vnb-lambda(\allowbreak{}c6va\_, carr(\allowbreak{}c6vk\_), c6vw\_(\allowbreak{}(\allowbreak{}act(\allowbreak{}c6vm\_))(\allowbreak{}c6va\_, c6vx\_), c6vy\_))) and is-add-character(\allowbreak{}c6vk\_, vnb-lambda(\allowbreak{}c6va\_, carr(\allowbreak{}c6vk\_), c6vw\_(\allowbreak{}c6vy\_, (\allowbreak{}act(\allowbreak{}c6vm\_))(\allowbreak{}c6va\_, c6vx\_))))))
\end{quote}
mentioned by 16 result(s)

\smallskip

\noindent\textbf{\texttt{is-e-bilinear-bicharacter(c6vk\_, c6vm\_, c6vw\_)}}\quad\textit{predicate}\index{is-e-bilinear-bicharacter(c6vk_, c6vm_, c6vw_)@\texttt{is-e-bilinear-bicharacter(c6vk\_, c6vm\_, c6vw\_)}}\\
reads: \$3 is a bilinear bicharacter on the vector space \$2 over \$1 \\
defined by \texttt{is-e-bilinear-bicharacter} \\
characterized by \texttt{is-e-bilinear-bicharacter}  \textit{(definitional)}:
\begin{quote}\ttfamily\footnotesize\raggedright\noindent
forall(\allowbreak{}[c6vk\_, c6vm\_, c6vw\_], is-e-bilinear-bicharacter(\allowbreak{}c6vk\_, c6vm\_, c6vw\_) iff is-e-bicharacter(\allowbreak{}c6vk\_, c6vm\_, c6vw\_) and forall(\allowbreak{}[c6va\_ in carr(\allowbreak{}c6vk\_), c6vx\_ in vec(\allowbreak{}c6vm\_), c6vy\_ in vec(\allowbreak{}c6vm\_)], c6vw\_(\allowbreak{}(\allowbreak{}act(\allowbreak{}c6vm\_))(\allowbreak{}c6va\_, c6vx\_), c6vy\_) = c6vw\_(\allowbreak{}c6vx\_, (\allowbreak{}act(\allowbreak{}c6vm\_))(\allowbreak{}c6va\_, c6vy\_))))
\end{quote}
mentioned by 14 result(s)

\smallskip

\noindent\textbf{\texttt{is-e-character(c6vk\_, c6vm\_, c6vc\_)}}\quad\textit{predicate}\index{is-e-character(c6vk_, c6vm_, c6vc_)@\texttt{is-e-character(c6vk\_, c6vm\_, c6vc\_)}}\\
reads: \$3 is a character of the additive group of the vector space \$2 over \$1 \\
defined by \texttt{is-e-character} \\
characterized by \texttt{is-e-character}  \textit{(definitional)}:
\begin{quote}\ttfamily\footnotesize\raggedright\noindent
forall(\allowbreak{}[c6vk\_, c6vm\_, c6vc\_], is-e-character(\allowbreak{}c6vk\_, c6vm\_, c6vc\_) iff c6vc\_ in fun(\allowbreak{}vec(\allowbreak{}c6vm\_), unit-circle) and forall(\allowbreak{}[c6vx\_ in vec(\allowbreak{}c6vm\_), c6vy\_ in vec(\allowbreak{}c6vm\_)], c6vc\_(\allowbreak{}(\allowbreak{}vadd(\allowbreak{}c6vm\_))(\allowbreak{}c6vx\_, c6vy\_)) = c6vc\_(\allowbreak{}c6vx\_) * c6vc\_(\allowbreak{}c6vy\_)) and forall(\allowbreak{}[c6vx\_ in vec(\allowbreak{}c6vm\_)], is-add-character(\allowbreak{}c6vk\_, vnb-lambda(\allowbreak{}c6va\_, carr(\allowbreak{}c6vk\_), c6vc\_(\allowbreak{}(\allowbreak{}act(\allowbreak{}c6vm\_))(\allowbreak{}c6va\_, c6vx\_))))))
\end{quote}
mentioned by 2 result(s)

\smallskip

\noindent\textbf{\texttt{is-e-unipolar(c6vk\_, c6vm\_, c6vn\_, c6vw\_)}}\quad\textit{predicate}\index{is-e-unipolar(c6vk_, c6vm_, c6vn_, c6vw_)@\texttt{is-e-unipolar(c6vk\_, c6vm\_, c6vn\_, c6vw\_)}}\\
reads: \$4 is a unipolar bicharacter on \$2 over \$1 for the norm \$3 \\
defined by \texttt{is-e-unipolar} \\
characterized by \texttt{is-e-unipolar}  \textit{(definitional)}:
\begin{quote}\ttfamily\footnotesize\raggedright\noindent
forall(\allowbreak{}[c6vk\_, c6vm\_, c6vn\_, c6vw\_], is-e-unipolar(\allowbreak{}c6vk\_, c6vm\_, c6vn\_, c6vw\_) iff is-norm-on(\allowbreak{}c6vk\_, c6vm\_, c6vn\_) and is-e-bilinear-bicharacter(\allowbreak{}c6vk\_, c6vm\_, c6vw\_) and e-polar(\allowbreak{}c6vm\_, c6vw\_, norm-ball(\allowbreak{}c6vm\_, c6vn\_, 1)) = norm-ball(\allowbreak{}c6vm\_, c6vn\_, 1))
\end{quote}
mentioned by 4 result(s)

\smallskip

\noindent\textbf{\texttt{is-eigenvalue(kqv\_, kqt\_, kqc\_)}}\quad\textit{predicate}\index{is-eigenvalue(kqv_, kqt_, kqc_)@\texttt{is-eigenvalue(kqv\_, kqt\_, kqc\_)}}\\
reads: \$3 is an eigenvalue of \$2 \\
defined by \texttt{is-eigenvalue} \\
characterized by \texttt{is-eigenvalue}  \textit{(definitional)}:
\begin{quote}\ttfamily\footnotesize\raggedright\noindent
forall(\allowbreak{}[kqv\_, kqt\_, kqc\_], is-eigenvalue(\allowbreak{}kqv\_, kqt\_, kqc\_) iff kqc\_ in cc and forsome(\allowbreak{}[kqx\_ in eigenspace(\allowbreak{}kqv\_, kqt\_, kqc\_)], not(\allowbreak{}kqx\_ = vzero(\allowbreak{}kqv\_))))
\end{quote}
mentioned by 9 result(s)

\smallskip

\noindent\textbf{\texttt{is-element-order(g, x, d)}}\quad\textit{predicate}\index{is-element-order(g, x, d)@\texttt{is-element-order(g, x, d)}}\\
reads: \$3 is the order of \$2 in \$1 \\
defined by \texttt{is-element-order} \\
characterized by \texttt{is-element-order}  \textit{(definitional)}:
\begin{quote}\ttfamily\footnotesize\raggedright\noindent
forall(\allowbreak{}[g, x, d], is-element-order(\allowbreak{}g, x, d) iff d in nn and 1 <= d and gpow(\allowbreak{}g, x, d) = iden(\allowbreak{}g) and forall(\allowbreak{}[eod\_ in nn], 1 <= eod\_ implies gpow(\allowbreak{}g, x, eod\_) = iden(\allowbreak{}g) implies d <= eod\_))
\end{quote}
mentioned by 26 result(s)

\smallskip

\noindent\textbf{\texttt{is-elementary(r, n, e)}}\quad\textit{predicate}\index{is-elementary(r, n, e)@\texttt{is-elementary(r, n, e)}}\\
reads: \$3 is an elementary column \$2-by-\$2 matrix over \$1 \\
defined by \texttt{is-elementary} \\
characterized by \texttt{is-elementary}  \textit{(definitional)}:
\begin{quote}\ttfamily\footnotesize\raggedright\noindent
forall(\allowbreak{}[r, n, e], is-elementary(\allowbreak{}r, n, e) iff forsome(\allowbreak{}[elk\_, ell\_], elk\_ in interval(\allowbreak{}1, n) and ell\_ in interval(\allowbreak{}1, n) and not(\allowbreak{}elk\_ = ell\_) and e = elem-f(\allowbreak{}r, n, elk\_, ell\_)) or forsome(\allowbreak{}[elr\_, elk\_, ell\_], elr\_ in carr(\allowbreak{}r) and elk\_ in interval(\allowbreak{}1, n) and ell\_ in interval(\allowbreak{}1, n) and not(\allowbreak{}elk\_ = ell\_) and e = elem-g(\allowbreak{}r, n, elr\_, elk\_, ell\_)) or forsome(\allowbreak{}[elr\_, elk\_], elr\_ in carr(\allowbreak{}r) and forsome(\allowbreak{}[els\_ in carr(\allowbreak{}r)], (\allowbreak{}mul(\allowbreak{}r))(\allowbreak{}elr\_, els\_) = one(\allowbreak{}r) and (\allowbreak{}mul(\allowbreak{}r))(\allowbreak{}els\_, elr\_) = one(\allowbreak{}r)) and elk\_ in interval(\allowbreak{}1, n) and e = elem-h(\allowbreak{}r, n, elr\_, elk\_)))
\end{quote}
also: \texttt{elem-f-is-elementary} \texttt{elem-g-is-elementary} \texttt{elem-h-is-elementary} \\
mentioned by 16 result(s)

\smallskip

\noindent\textbf{\texttt{is-encrypt-trace(sasp\_, sas\_)}}\quad\textit{predicate}\index{is-encrypt-trace(sasp_, sas_)@\texttt{is-encrypt-trace(sasp\_, sas\_)}}\\
reads: \$2 is an encryption trace in the strand space \$1 \\
defined by \texttt{is-encrypt-trace} \\
characterized by \texttt{is-encrypt-trace}  \textit{(definitional)}:
\begin{quote}\ttfamily\footnotesize\raggedright\noindent
forall(\allowbreak{}[sasp\_, sas\_], is-encrypt-trace(\allowbreak{}sasp\_, sas\_) iff forsome(\allowbreak{}[sak\_ in keys(\allowbreak{}sasp\_), sag\_ in msg-set(\allowbreak{}tokens(\allowbreak{}sasp\_), keys(\allowbreak{}sasp\_))], sas\_ = [[received, sak\_], [received, sag\_], [sent, encr(\allowbreak{}sag\_, sak\_)]]))
\end{quote}
mentioned by 6 result(s)

\smallskip

\noindent\textbf{\texttt{is-entry-point(ssp\_, ssn\_, ssi\_)}}\quad\textit{predicate}\index{is-entry-point(ssp_, ssn_, ssi_)@\texttt{is-entry-point(ssp\_, ssn\_, ssi\_)}}\\
reads: the node \$2 of \$1 is an entry point into \$3 \\
defined by \texttt{is-entry-point} \\
characterized by \texttt{is-entry-point}  \textit{(definitional)}:
\begin{quote}\ttfamily\footnotesize\raggedright\noindent
forall(\allowbreak{}[ssp\_, ssn\_, ssi\_], is-entry-point(\allowbreak{}ssp\_, ssn\_, ssi\_) iff ssn\_ in nodes(\allowbreak{}ssp\_) and is-sent(\allowbreak{}ssp\_, ssn\_) and node-content(\allowbreak{}ssp\_, ssn\_) in ssi\_ and forall(\allowbreak{}[ssm\_], same-strand-before(\allowbreak{}ssp\_, ssm\_, ssn\_) implies not(\allowbreak{}node-content(\allowbreak{}ssp\_, ssm\_) in ssi\_)))
\end{quote}
mentioned by 9 result(s)

\smallskip

\noindent\textbf{\texttt{is-eps-cauchy-seq(s, eps, y)}}\quad\textit{predicate}\index{is-eps-cauchy-seq(s, eps, y)@\texttt{is-eps-cauchy-seq(s, eps, y)}}\\
reads: \$3 is \$2-Cauchy in \$1 \\
defined by \texttt{is-eps-cauchy-seq} \\
characterized by \texttt{is-eps-cauchy-seq}  \textit{(definitional)}:
\begin{quote}\ttfamily\footnotesize\raggedright\noindent
forall(\allowbreak{}[s, eps, y], is-eps-cauchy-seq(\allowbreak{}s, eps, y) iff is-metric-space(\allowbreak{}s) and y in fun(\allowbreak{}nn, pts(\allowbreak{}s)) and forall(\allowbreak{}[m in nn, n\_ in nn], (\allowbreak{}dist(\allowbreak{}s))(\allowbreak{}y(\allowbreak{}m), y(\allowbreak{}n\_)) <= eps))
\end{quote}
mentioned by 3 result(s)

\smallskip

\noindent\textbf{\texttt{is-equicontinuous(s, t, fam)}}\quad\textit{predicate}\index{is-equicontinuous(s, t, fam)@\texttt{is-equicontinuous(s, t, fam)}}\\
reads: \$3 is equicontinuous from \$1 to \$2 \\
defined by \texttt{is-equicontinuous} \\
characterized by \texttt{is-equicontinuous}  \textit{(definitional)}:
\begin{quote}\ttfamily\footnotesize\raggedright\noindent
forall(\allowbreak{}[s, t, fam], is-equicontinuous(\allowbreak{}s, t, fam) iff is-metric-space(\allowbreak{}s) and is-metric-space(\allowbreak{}t) and fam in fun(\allowbreak{}nn, fun(\allowbreak{}pts(\allowbreak{}s), pts(\allowbreak{}t))) and forall(\allowbreak{}[x in pts(\allowbreak{}s), eps], pos-rr(\allowbreak{}eps) implies forsome(\allowbreak{}[del], pos-rr(\allowbreak{}del) and forall(\allowbreak{}[k in nn, y in pts(\allowbreak{}s)], (\allowbreak{}dist(\allowbreak{}s))(\allowbreak{}x, y) < del implies (\allowbreak{}dist(\allowbreak{}t))(\allowbreak{}(\allowbreak{}fam(\allowbreak{}k))(\allowbreak{}x), (\allowbreak{}fam(\allowbreak{}k))(\allowbreak{}y)) < eps))))
\end{quote}
mentioned by 17 result(s)

\smallskip

\noindent\textbf{\texttt{is-equicontinuous-set(eqss\_, eqst\_, eqse\_)}}\quad\textit{predicate}\index{is-equicontinuous-set(eqss_, eqst_, eqse_)@\texttt{is-equicontinuous-set(eqss\_, eqst\_, eqse\_)}}\\
reads: \$3 is an equicontinuous set of functions from \$1 to \$2 \\
defined by \texttt{is-equicontinuous-set} \\
characterized by \texttt{is-equicontinuous-set}  \textit{(definitional)}:
\begin{quote}\ttfamily\footnotesize\raggedright\noindent
forall(\allowbreak{}[eqss\_, eqst\_, eqse\_], is-equicontinuous-set(\allowbreak{}eqss\_, eqst\_, eqse\_) iff is-metric-space(\allowbreak{}eqss\_) and is-metric-space(\allowbreak{}eqst\_) and eqse\_ subset fun(\allowbreak{}pts(\allowbreak{}eqss\_), pts(\allowbreak{}eqst\_)) and forall(\allowbreak{}[eqsx\_ in pts(\allowbreak{}eqss\_), eqsp\_], pos-rr(\allowbreak{}eqsp\_) implies forsome(\allowbreak{}[eqsd\_], pos-rr(\allowbreak{}eqsd\_) and forall(\allowbreak{}[eqsf\_ in eqse\_, eqsy\_ in pts(\allowbreak{}eqss\_)], (\allowbreak{}dist(\allowbreak{}eqss\_))(\allowbreak{}eqsx\_, eqsy\_) < eqsd\_ implies (\allowbreak{}dist(\allowbreak{}eqst\_))(\allowbreak{}eqsf\_(\allowbreak{}eqsx\_), eqsf\_(\allowbreak{}eqsy\_)) < eqsp\_))))
\end{quote}
mentioned by 5 result(s)

\smallskip

\noindent\textbf{\texttt{is-equicontinuous-t(t, ee)}}\quad\textit{predicate}\index{is-equicontinuous-t(t, ee)@\texttt{is-equicontinuous-t(t, ee)}}\\
reads: \$2 is an equicontinuous set of real functions on \$1 \\
defined by \texttt{is-equicontinuous-t} \\
characterized by \texttt{is-equicontinuous-t}  \textit{(definitional)}:
\begin{quote}\ttfamily\footnotesize\raggedright\noindent
forall(\allowbreak{}[t, ee], is-equicontinuous-t(\allowbreak{}t, ee) iff ee subset fun(\allowbreak{}pts(\allowbreak{}t), rr) and forall(\allowbreak{}[eqx\_ in pts(\allowbreak{}t), eqe\_], pos-rr(\allowbreak{}eqe\_) implies forsome(\allowbreak{}[eqv\_ in opens(\allowbreak{}t)], eqx\_ in eqv\_ and forall(\allowbreak{}[eqf\_ in ee, eqy\_ in eqv\_], abs(\allowbreak{}eqf\_(\allowbreak{}eqx\_) - eqf\_(\allowbreak{}eqy\_)) < eqe\_))))
\end{quote}
mentioned by 4 result(s)

\smallskip

\noindent\textbf{\texttt{is-equivalence(rho, crr)}}\quad\textit{predicate}\index{is-equivalence(rho, crr)@\texttt{is-equivalence(rho, crr)}}\\
reads: \$1 is an equivalence relation on \$2 \\
defined by \texttt{is-equivalence} \\
characterized by \texttt{is-equivalence}  \textit{(definitional)}:
\begin{quote}\ttfamily\footnotesize\raggedright\noindent
forall(\allowbreak{}[rho, crr], is-equivalence(\allowbreak{}rho, crr) iff rho in power(\allowbreak{}cartesian(\allowbreak{}crr, crr)) and forall(\allowbreak{}[u in crr], [u, u] in rho) and forall(\allowbreak{}[u in crr, v in crr], [u, v] in rho implies [v, u] in rho) and forall(\allowbreak{}[u in crr, v in crr, w in crr], [u, v] in rho and [v, w] in rho implies [u, w] in rho))
\end{quote}
also: \texttt{sip-rel-is-equivalence-csip-as-sip} \texttt{qmod-rel-is-equivalence-cna-as-module} \texttt{qmod-rel-is-equivalence-sip-as-module} \texttt{qmod-rel-is-equivalence-uba-as-module} \texttt{qmod-rel-is-equivalence-cnvs-as-module} \texttt{qmod-rel-is-equivalence-sip-as-module-csip-as-sip} \texttt{qmod-rel-is-equivalence-uba-as-module-uiba-as-uba} \texttt{ae-rel-is-equivalence} ... (11 in all) \\
mentioned by 16 result(s)

\smallskip

\noindent\textbf{\texttt{is-ess-self-adjoint-graph(ubv\_, ubg\_)}}\quad\textit{predicate}\index{is-ess-self-adjoint-graph(ubv_, ubg_)@\texttt{is-ess-self-adjoint-graph(ubv\_, ubg\_)}}\\
reads: \$2 is an essentially self-adjoint graph \\
defined by \texttt{is-ess-self-adjoint-graph} \\
characterized by \texttt{is-ess-self-adjoint-graph}  \textit{(definitional)}:
\begin{quote}\ttfamily\footnotesize\raggedright\noindent
forall(\allowbreak{}[ubv\_, ubg\_], is-ess-self-adjoint-graph(\allowbreak{}ubv\_, ubg\_) iff is-self-adjoint-graph(\allowbreak{}ubv\_, graph-closure(\allowbreak{}ubv\_, ubg\_)))
\end{quote}
mentioned by 2 result(s)

\smallskip

\noindent\textbf{\texttt{is-essential-singularity(mrf\_, mra\_)}}\quad\textit{predicate}\index{is-essential-singularity(mrf_, mra_)@\texttt{is-essential-singularity(mrf\_, mra\_)}}\\
reads: \$1 has an essential singularity at \$2 \\
defined by \texttt{is-essential-singularity} \\
characterized by \texttt{is-essential-singularity}  \textit{(definitional)}:
\begin{quote}\ttfamily\footnotesize\raggedright\noindent
forall(\allowbreak{}[mrf\_, mra\_], is-essential-singularity(\allowbreak{}mrf\_, mra\_) iff mra\_ in cc and not(\allowbreak{}is-removable-singularity(\allowbreak{}mrf\_, mra\_)) and not(\allowbreak{}forsome(\allowbreak{}[mrn\_], is-pole(\allowbreak{}mrf\_, mra\_, mrn\_))))
\end{quote}
mentioned by 3 result(s)

\smallskip

\noindent\textbf{\texttt{is-euclidean-domain(ges\_)}}\quad\textit{predicate}\index{is-euclidean-domain(ges_)@\texttt{is-euclidean-domain(ges\_)}}\\
reads: \$1 is a Euclidean domain \\
defined by \texttt{is-euclidean-domain} \\
characterized by \texttt{is-euclidean-domain}  \textit{(definitional)}:
\begin{quote}\ttfamily\footnotesize\raggedright\noindent
forall(\allowbreak{}[ges\_], is-euclidean-domain(\allowbreak{}ges\_) iff is-integral-domain(\allowbreak{}ges\_) and forsome(\allowbreak{}[gen\_], is-degree-function(\allowbreak{}ges\_, gen\_)))
\end{quote}
also: \texttt{zz-is-euclidean-domain} \texttt{poly-is-euclidean-domain} \\
mentioned by 12 result(s)

\smallskip

\noindent\textbf{\texttt{is-euclidean-ring(s)}}\quad\textit{predicate}\index{is-euclidean-ring(s)@\texttt{is-euclidean-ring(s)}}\\
reads: a Euclidean ring \\
a predicate \\
characterized by \texttt{is-euclidean-ring-def}  \textit{(definitional)}:
\begin{quote}\ttfamily\footnotesize\raggedright\noindent
forall(\allowbreak{}[s], is-euclidean-ring(\allowbreak{}s) iff is-integral-domain(\allowbreak{}s) and forsome(\allowbreak{}[dg in fun(\allowbreak{}carr(\allowbreak{}s), nn)], forall(\allowbreak{}[a in carr(\allowbreak{}s), b in carr(\allowbreak{}s)], not(\allowbreak{}b = zero(\allowbreak{}s)) implies forsome(\allowbreak{}[q in carr(\allowbreak{}s), r in carr(\allowbreak{}s)], a = (\allowbreak{}add(\allowbreak{}s))(\allowbreak{}(\allowbreak{}mul(\allowbreak{}s))(\allowbreak{}q, b), r) and (\allowbreak{}r = zero(\allowbreak{}s) or succ(\allowbreak{}dg(\allowbreak{}r)) <= dg(\allowbreak{}b))))))
\end{quote}
also: \texttt{field-as-euclidean-ring-is-euclidean-ring} \texttt{zz-is-euclidean-ring} \texttt{is-euclidean-ring-is-set} \texttt{euclidean-domain-is-euclidean-ring} \\
mentioned by 90 result(s)

\smallskip

\noindent\textbf{\texttt{is-eventually-in(d, x, u)}}\quad\textit{predicate}\index{is-eventually-in(d, x, u)@\texttt{is-eventually-in(d, x, u)}}\\
reads: the net \$2 on \$1 is eventually in \$3 \\
defined by \texttt{is-eventually-in} \\
characterized by \texttt{is-eventually-in}  \textit{(definitional)}:
\begin{quote}\ttfamily\footnotesize\raggedright\noindent
forall(\allowbreak{}[d, x, u], is-eventually-in(\allowbreak{}d, x, u) iff forsome(\allowbreak{}[evi\_ in pts(\allowbreak{}d)], forall(\allowbreak{}[evj\_ in pts(\allowbreak{}d)], [evi\_, evj\_] in dleq(\allowbreak{}d) implies x(\allowbreak{}evj\_) in u)))
\end{quote}
mentioned by 8 result(s)

\smallskip

\noindent\textbf{\texttt{is-extremal-subset(v, b, a)}}\quad\textit{predicate}\index{is-extremal-subset(v, b, a)@\texttt{is-extremal-subset(v, b, a)}}\\
reads: \$2 is an extremal subset of \$3 in the real vector space \$1 \\
defined by \texttt{is-extremal-subset} \\
characterized by \texttt{is-extremal-subset}  \textit{(definitional)}:
\begin{quote}\ttfamily\footnotesize\raggedright\noindent
forall(\allowbreak{}[v, b, a], is-extremal-subset(\allowbreak{}v, b, a) iff is-convex-subset(\allowbreak{}v, b) and forsome(\allowbreak{}[exz\_], exz\_ in b) and b subset a and forall(\allowbreak{}[exx\_, exy\_, ext\_], exx\_ in a implies exy\_ in a implies ext\_ in rr implies 0 < ext\_ implies ext\_ < 1 implies (\allowbreak{}vadd(\allowbreak{}v))(\allowbreak{}(\allowbreak{}act(\allowbreak{}v))(\allowbreak{}ext\_, exx\_), (\allowbreak{}act(\allowbreak{}v))(\allowbreak{}1 - ext\_, exy\_)) in b implies exx\_ in b and exy\_ in b))
\end{quote}
mentioned by 11 result(s)

\smallskip

\noindent\textbf{\texttt{is-extreme-point(epm\_, epk\_, epx\_)}}\quad\textit{predicate}\index{is-extreme-point(epm_, epk_, epx_)@\texttt{is-extreme-point(epm\_, epk\_, epx\_)}}\\
reads: \$3 is an extreme point of \$2 in \$1 \\
defined by \texttt{is-extreme-point} \\
characterized by \texttt{is-extreme-point}  \textit{(definitional)}:
\begin{quote}\ttfamily\footnotesize\raggedright\noindent
forall(\allowbreak{}[epm\_, epk\_, epx\_], is-extreme-point(\allowbreak{}epm\_, epk\_, epx\_) iff epx\_ in epk\_ and forall(\allowbreak{}[epu\_ in epk\_, epv\_ in epk\_, ept\_ in rr], 0 < ept\_ implies ept\_ < 1 implies epx\_ = (\allowbreak{}vadd(\allowbreak{}epm\_))(\allowbreak{}(\allowbreak{}act(\allowbreak{}epm\_))(\allowbreak{}1 - ept\_, epu\_), (\allowbreak{}act(\allowbreak{}epm\_))(\allowbreak{}ept\_, epv\_)) implies epu\_ = epx\_ and epv\_ = epx\_))
\end{quote}
mentioned by 12 result(s)

\smallskip

\noindent\textbf{\texttt{is-faithful-bicharacter(c5k\_, c5w\_)}}\quad\textit{predicate}\index{is-faithful-bicharacter(c5k_, c5w_)@\texttt{is-faithful-bicharacter(c5k\_, c5w\_)}}\\
reads: \$2 is a faithful bicharacter on \$1 \\
defined by \texttt{is-faithful-bicharacter} \\
characterized by \texttt{is-faithful-bicharacter}  \textit{(definitional)}:
\begin{quote}\ttfamily\footnotesize\raggedright\noindent
forall(\allowbreak{}[c5k\_, c5w\_], is-faithful-bicharacter(\allowbreak{}c5k\_, c5w\_) iff is-bicharacter(\allowbreak{}c5k\_, c5w\_) and forall(\allowbreak{}[c5c\_], is-add-character(\allowbreak{}c5k\_, c5c\_) implies forsome(\allowbreak{}[c5x\_ in carr(\allowbreak{}c5k\_)], c5c\_ = vnb-lambda(\allowbreak{}c5h\_, carr(\allowbreak{}c5k\_), c5w\_(\allowbreak{}c5h\_, c5x\_)) and forall(\allowbreak{}[c5y\_ in carr(\allowbreak{}c5k\_)], c5c\_ = vnb-lambda(\allowbreak{}c5h\_, carr(\allowbreak{}c5k\_), c5w\_(\allowbreak{}c5h\_, c5y\_)) implies c5y\_ = c5x\_)) and forsome(\allowbreak{}[c5x\_ in carr(\allowbreak{}c5k\_)], c5c\_ = vnb-lambda(\allowbreak{}c5h\_, carr(\allowbreak{}c5k\_), c5w\_(\allowbreak{}c5x\_, c5h\_)) and forall(\allowbreak{}[c5y\_ in carr(\allowbreak{}c5k\_)], c5c\_ = vnb-lambda(\allowbreak{}c5h\_, carr(\allowbreak{}c5k\_), c5w\_(\allowbreak{}c5y\_, c5h\_)) implies c5y\_ = c5x\_))) and forall(\allowbreak{}[c5x\_ in carr(\allowbreak{}c5k\_)], not(\allowbreak{}c5x\_ = zero(\allowbreak{}c5k\_)) implies forsome(\allowbreak{}[c5c\_], is-add-character(\allowbreak{}c5k\_, c5c\_) and not(\allowbreak{}c5c\_(\allowbreak{}c5x\_) = 1))))
\end{quote}
mentioned by 15 result(s)

\smallskip

\noindent\textbf{\texttt{is-fg-cc-module-over(aa, vv)}}\quad\textit{predicate}\index{is-fg-cc-module-over(aa, vv)@\texttt{is-fg-cc-module-over(aa, vv)}}\\
reads: \$2 is a finitely generated \$1-module inside the complex numbers \\
defined by \texttt{is-fg-cc-module-over} \\
characterized by \texttt{is-fg-cc-module-over}  \textit{(definitional)}:
\begin{quote}\ttfamily\footnotesize\raggedright\noindent
forall(\allowbreak{}[aa, vv], is-fg-cc-module-over(\allowbreak{}aa, vv) iff is-cc-module-over(\allowbreak{}aa, vv) and forsome(\allowbreak{}[amn\_ in nn, amg\_ in fun(\allowbreak{}interval(\allowbreak{}1, amn\_), vv)], forall(\allowbreak{}[amv\_ in vv], forsome(\allowbreak{}[amb\_ in fun(\allowbreak{}interval(\allowbreak{}1, amn\_), aa)], amv\_ = cc-lincomb(\allowbreak{}amn\_, amb\_, amg\_)))))
\end{quote}
mentioned by 10 result(s)

\smallskip

\noindent\textbf{\texttt{is-field(s)}}\quad\textit{predicate}\index{is-field(s)@\texttt{is-field(s)}}\\
reads: a field \\
a predicate \\
characterized by \texttt{is-field}  \textit{(definitional)}:
\begin{quote}\ttfamily\footnotesize\raggedright\noindent
forall(\allowbreak{}[s], is-field(\allowbreak{}s) iff length(\allowbreak{}s) = 8 and carr(\allowbreak{}s) in set and add(\allowbreak{}s) in fun(\allowbreak{}cartesian(\allowbreak{}carr(\allowbreak{}s), carr(\allowbreak{}s)), carr(\allowbreak{}s)) and mul(\allowbreak{}s) in fun(\allowbreak{}cartesian(\allowbreak{}carr(\allowbreak{}s), carr(\allowbreak{}s)), carr(\allowbreak{}s)) and neg(\allowbreak{}s) in fun(\allowbreak{}carr(\allowbreak{}s), carr(\allowbreak{}s)) and zero(\allowbreak{}s) in carr(\allowbreak{}s) and one(\allowbreak{}s) in carr(\allowbreak{}s) and non-zero(\allowbreak{}s) in set and non-zero(\allowbreak{}s) = difference(\allowbreak{}carr(\allowbreak{}s), singleton(\allowbreak{}zero(\allowbreak{}s))) and mul-inv(\allowbreak{}s) in fun(\allowbreak{}non-zero(\allowbreak{}s), non-zero(\allowbreak{}s)) and is-associative(\allowbreak{}add(\allowbreak{}s), carr(\allowbreak{}s)) and is-commutative(\allowbreak{}add(\allowbreak{}s), carr(\allowbreak{}s)) and is-identity(\allowbreak{}add(\allowbreak{}s), zero(\allowbreak{}s), carr(\allowbreak{}s)) and has-inverses(\allowbreak{}add(\allowbreak{}s), zero(\allowbreak{}s), neg(\allowbreak{}s), carr(\allowbreak{}s)) and is-associative(\allowbreak{}mul(\allowbreak{}s), carr(\allowbreak{}s)) and is-commutative(\allowbreak{}mul(\allowbreak{}s), carr(\allowbreak{}s)) and is-identity(\allowbreak{}mul(\allowbreak{}s), one(\allowbreak{}s), carr(\allowbreak{}s)) and is-distributive(\allowbreak{}add(\allowbreak{}s), mul(\allowbreak{}s), carr(\allowbreak{}s)) and not(\allowbreak{}zero(\allowbreak{}s) = one(\allowbreak{}s)) and forall(\allowbreak{}[a\_ in non-zero(\allowbreak{}s)], (\allowbreak{}mul(\allowbreak{}s))(\allowbreak{}a\_, (\allowbreak{}mul-inv(\allowbreak{}s))(\allowbreak{}a\_)) = one(\allowbreak{}s)) and s in set)
\end{quote}
also: \texttt{is-field-is-set} \texttt{qq-field-is-field} \texttt{rr-field-is-field} \texttt{field-is-field-ring} \\
mentioned by 33 result(s)

\smallskip

\noindent\textbf{\texttt{is-field-ring(s)}}\quad\textit{predicate}\index{is-field-ring(s)@\texttt{is-field-ring(s)}}\\
reads: a field \\
a predicate \\
characterized by \texttt{is-field-ring-def}  \textit{(definitional)}:
\begin{quote}\ttfamily\footnotesize\raggedright\noindent
forall(\allowbreak{}[s], is-field-ring(\allowbreak{}s) iff is-commutative-ring(\allowbreak{}s) and not(\allowbreak{}one(\allowbreak{}s) = zero(\allowbreak{}s)) and forall(\allowbreak{}[a\_ in carr(\allowbreak{}s)], not(\allowbreak{}a\_ = zero(\allowbreak{}s)) implies forsome(\allowbreak{}[b\_ in carr(\allowbreak{}s)], (\allowbreak{}mul(\allowbreak{}s))(\allowbreak{}a\_, b\_) = one(\allowbreak{}s))))
\end{quote}
also: \texttt{field-is-field-ring} \texttt{is-field-ring-is-set} \texttt{rr-field-is-field-ring} \texttt{laurent-ring-is-field-ring} \\
mentioned by 153 result(s)

\smallskip

\noindent\textbf{\texttt{is-fin-ext(x4k\_, x4f\_, x4md\_, x4n\_, x4u\_)}}\quad\textit{predicate}\index{is-fin-ext(x4k_, x4f_, x4md_, x4n_, x4u_)@\texttt{is-fin-ext(x4k\_, x4f\_, x4md\_, x4n\_, x4u\_)}}\\
reads: \$2 is an extension of \$1 of degree \$4 with basis \$5 (the scalar action \$3) \\
defined by \texttt{is-fin-ext} \\
characterized by \texttt{is-fin-ext}  \textit{(definitional)}:
\begin{quote}\ttfamily\footnotesize\raggedright\noindent
forall(\allowbreak{}[x4k\_, x4f\_, x4md\_, x4n\_, x4u\_], is-fin-ext(\allowbreak{}x4k\_, x4f\_, x4md\_, x4n\_, x4u\_) iff is-normed-field(\allowbreak{}x4k\_) and is-field-ring(\allowbreak{}x4f\_) and is-module(\allowbreak{}x4md\_) and scal(\allowbreak{}x4md\_) = normed-field-as-commutative-ring(\allowbreak{}x4k\_) and vec(\allowbreak{}x4md\_) = carr(\allowbreak{}x4f\_) and vadd(\allowbreak{}x4md\_) = add(\allowbreak{}x4f\_) and vzero(\allowbreak{}x4md\_) = zero(\allowbreak{}x4f\_) and vneg(\allowbreak{}x4md\_) = neg(\allowbreak{}x4f\_) and forall(\allowbreak{}[x4r\_ in carr(\allowbreak{}x4k\_), x4x\_ in carr(\allowbreak{}x4f\_), x4y\_ in carr(\allowbreak{}x4f\_)], (\allowbreak{}mul(\allowbreak{}x4f\_))(\allowbreak{}x4x\_, (\allowbreak{}act(\allowbreak{}x4md\_))(\allowbreak{}x4r\_, x4y\_)) = (\allowbreak{}act(\allowbreak{}x4md\_))(\allowbreak{}x4r\_, (\allowbreak{}mul(\allowbreak{}x4f\_))(\allowbreak{}x4x\_, x4y\_))) and x4n\_ in nn and 1 <= x4n\_ and x4u\_ in mat(\allowbreak{}x4n\_, 1, carr(\allowbreak{}x4f\_)) and generates(\allowbreak{}x4md\_, x4n\_, x4u\_) and rel-free(\allowbreak{}x4md\_, x4n\_, x4u\_))
\end{quote}
mentioned by 32 result(s)

\smallskip

\noindent\textbf{\texttt{is-finite-cover(c, a)}}\quad\textit{predicate}\index{is-finite-cover(c, a)@\texttt{is-finite-cover(c, a)}}\\
reads: \$1 is a finite cover of \$2 \\
defined by \texttt{is-finite-cover} \\
characterized by \texttt{is-finite-cover}  \textit{(definitional)}:
\begin{quote}\ttfamily\footnotesize\raggedright\noindent
forall(\allowbreak{}[c, a], is-finite-cover(\allowbreak{}c, a) iff c in set and card(\allowbreak{}c) in nn and a subset big-union(\allowbreak{}u, c, u))
\end{quote}
mentioned by 8 result(s)

\smallskip

\noindent\textbf{\texttt{is-finite-dimensional(m)}}\quad\textit{predicate}\index{is-finite-dimensional(m)@\texttt{is-finite-dimensional(m)}}\\
reads: \$1 is finite dimensional \\
defined by \texttt{is-finite-dimensional} \\
characterized by \texttt{is-finite-dimensional}  \textit{(definitional)}:
\begin{quote}\ttfamily\footnotesize\raggedright\noindent
forall(\allowbreak{}[m], is-finite-dimensional(\allowbreak{}m) iff is-vector-space(\allowbreak{}m) and is-noetherian(\allowbreak{}m))
\end{quote}
mentioned by 8 result(s)

\smallskip

\noindent\textbf{\texttt{is-finite-measure(omega, ca, mu)}}\quad\textit{predicate}\index{is-finite-measure(omega, ca, mu)@\texttt{is-finite-measure(omega, ca, mu)}}\\
reads: \$3 is a finite countably additive measure on the sigma-algebra \$2 of subsets of \$1 \\
defined by \texttt{is-finite-measure} \\
characterized by \texttt{is-finite-measure}  \textit{(definitional)}:
\begin{quote}\ttfamily\footnotesize\raggedright\noindent
forall(\allowbreak{}[omega, ca, mu], is-finite-measure(\allowbreak{}omega, ca, mu) iff is-measure(\allowbreak{}omega, ca, mu) and mu(\allowbreak{}omega) in rr)
\end{quote}
mentioned by 11 result(s)

\smallskip

\noindent\textbf{\texttt{is-finite-rank-op(kqv\_, kqw\_, kqt\_)}}\quad\textit{predicate}\index{is-finite-rank-op(kqv_, kqw_, kqt_)@\texttt{is-finite-rank-op(kqv\_, kqw\_, kqt\_)}}\\
reads: \$3 is a finite-rank operator from \$1 to \$2 \\
defined by \texttt{is-finite-rank-op} \\
characterized by \texttt{is-finite-rank-op}  \textit{(definitional)}:
\begin{quote}\ttfamily\footnotesize\raggedright\noindent
forall(\allowbreak{}[kqv\_, kqw\_, kqt\_], is-finite-rank-op(\allowbreak{}kqv\_, kqw\_, kqt\_) iff kqt\_ in bounded-ops(\allowbreak{}kqv\_, kqw\_) and forsome(\allowbreak{}[kqi\_ in set], card(\allowbreak{}kqi\_) in nn and forsome(\allowbreak{}[kqa\_ in fun(\allowbreak{}kqi\_, vec(\allowbreak{}kqv\_)), kqb\_ in fun(\allowbreak{}kqi\_, vec(\allowbreak{}kqw\_))], forall(\allowbreak{}[kqx\_ in vec(\allowbreak{}kqv\_)], kqt\_(\allowbreak{}kqx\_) = finsum(\allowbreak{}normed-ag-as-abelian-group(\allowbreak{}ip-normed-ag(\allowbreak{}kqw\_)), vnb-lambda(\allowbreak{}kqk\_, kqi\_, (\allowbreak{}act(\allowbreak{}kqw\_))(\allowbreak{}(\allowbreak{}ip(\allowbreak{}kqv\_))(\allowbreak{}kqx\_, kqa\_(\allowbreak{}kqk\_)), kqb\_(\allowbreak{}kqk\_))), kqi\_)))))
\end{quote}
mentioned by 7 result(s)

\smallskip

\noindent\textbf{\texttt{is-flat-regular(frp\_)}}\quad\textit{predicate}\index{is-flat-regular(frp_)@\texttt{is-flat-regular(frp\_)}}\\
reads: \$1 is a public-key strand space whose regular contents are flat \\
defined by \texttt{is-flat-regular} \\
characterized by \texttt{is-flat-regular}  \textit{(definitional)}:
\begin{quote}\ttfamily\footnotesize\raggedright\noindent
forall(\allowbreak{}[frp\_], is-flat-regular(\allowbreak{}frp\_) iff is-pk-space(\allowbreak{}frp\_) and forall(\allowbreak{}[frn\_ in nodes(\allowbreak{}ispace(\allowbreak{}pk-inf(\allowbreak{}frp\_)))], is-regular-node(\allowbreak{}pk-inf(\allowbreak{}frp\_), frn\_) implies forsome(\allowbreak{}[frk\_ in keys(\allowbreak{}ispace(\allowbreak{}pk-inf(\allowbreak{}frp\_)))], forsome(\allowbreak{}[frx\_ in \{frd\_ in tokens(\allowbreak{}ispace(\allowbreak{}pk-inf(\allowbreak{}frp\_))): not(\allowbreak{}frd\_ in keys(\allowbreak{}ispace(\allowbreak{}pk-inf(\allowbreak{}frp\_))))\}], node-content(\allowbreak{}ispace(\allowbreak{}pk-inf(\allowbreak{}frp\_)), frn\_) = encr(\allowbreak{}frx\_, frk\_)) or forsome(\allowbreak{}[frl\_ in tuples(\allowbreak{}\{frd\_ in tokens(\allowbreak{}ispace(\allowbreak{}pk-inf(\allowbreak{}frp\_))): not(\allowbreak{}frd\_ in keys(\allowbreak{}ispace(\allowbreak{}pk-inf(\allowbreak{}frp\_))))\})], node-content(\allowbreak{}ispace(\allowbreak{}pk-inf(\allowbreak{}frp\_)), frn\_) = encr(\allowbreak{}lst(\allowbreak{}frl\_), frk\_)))))
\end{quote}
mentioned by 13 result(s)

\smallskip

\noindent\textbf{\texttt{is-flush-trace(sasp\_, sas\_)}}\quad\textit{predicate}\index{is-flush-trace(sasp_, sas_)@\texttt{is-flush-trace(sasp\_, sas\_)}}\\
reads: \$2 is a flushing trace in the strand space \$1 \\
defined by \texttt{is-flush-trace} \\
characterized by \texttt{is-flush-trace}  \textit{(definitional)}:
\begin{quote}\ttfamily\footnotesize\raggedright\noindent
forall(\allowbreak{}[sasp\_, sas\_], is-flush-trace(\allowbreak{}sasp\_, sas\_) iff forsome(\allowbreak{}[sag\_ in msg-set(\allowbreak{}tokens(\allowbreak{}sasp\_), keys(\allowbreak{}sasp\_))], sas\_ = [[received, sag\_]]))
\end{quote}
mentioned by 5 result(s)

\smallskip

\noindent\textbf{\texttt{is-fr-bounded(m, fam, aa)}}\quad\textit{predicate}\index{is-fr-bounded(m, fam, aa)@\texttt{is-fr-bounded(m, fam, aa)}}\\
reads: \$3 is a bounded set in (\$1, \$2) \\
defined by \texttt{is-fr-bounded} \\
characterized by \texttt{is-fr-bounded}  \textit{(definitional)}:
\begin{quote}\ttfamily\footnotesize\raggedright\noindent
forall(\allowbreak{}[m, fam, aa], is-fr-bounded(\allowbreak{}m, fam, aa) iff aa subset vec(\allowbreak{}m) and forall(\allowbreak{}[k in nn], forsome(\allowbreak{}[bd in rr], forall(\allowbreak{}[y in aa], (\allowbreak{}fam(\allowbreak{}k))(\allowbreak{}y) <= bd))))
\end{quote}
mentioned by 5 result(s)

\smallskip

\noindent\textbf{\texttt{is-fr-equicont(m1, fam1, m2, fam2, ix, tf)}}\quad\textit{predicate}\index{is-fr-equicont(m1, fam1, m2, fam2, ix, tf)@\texttt{is-fr-equicont(m1, fam1, m2, fam2, ix, tf)}}\\
reads: the family \$6 indexed by \$5 is equicontinuous from (\$1, \$2) to (\$3, \$4) \\
defined by \texttt{is-fr-equicont} \\
characterized by \texttt{is-fr-equicont}  \textit{(definitional)}:
\begin{quote}\ttfamily\footnotesize\raggedright\noindent
forall(\allowbreak{}[m1, fam1, m2, fam2, ix, tf], is-fr-equicont(\allowbreak{}m1, fam1, m2, fam2, ix, tf) iff tf in fun(\allowbreak{}ix, fun(\allowbreak{}vec(\allowbreak{}m1), vec(\allowbreak{}m2))) and forall(\allowbreak{}[k in nn], forsome(\allowbreak{}[nb in nn, ck], pos-rr(\allowbreak{}ck) and forall(\allowbreak{}[i in ix, x in vec(\allowbreak{}m1), r in rr], forall(\allowbreak{}[j in nn], j <= nb implies (\allowbreak{}fam1(\allowbreak{}j))(\allowbreak{}x) <= r) implies (\allowbreak{}fam2(\allowbreak{}k))(\allowbreak{}(\allowbreak{}tf(\allowbreak{}i))(\allowbreak{}x)) <= ck * r))))
\end{quote}
mentioned by 3 result(s)

\smallskip

\noindent\textbf{\texttt{is-fractions-of(cwk\_, cwa\_)}}\quad\textit{predicate}\index{is-fractions-of(cwk_, cwa_)@\texttt{is-fractions-of(cwk\_, cwa\_)}}\\
reads: \$1 is the field of fractions of \$2 \\
defined by \texttt{is-fractions-of} \\
characterized by \texttt{is-fractions-of}  \textit{(definitional)}:
\begin{quote}\ttfamily\footnotesize\raggedright\noindent
forall(\allowbreak{}[cwk\_, cwa\_], is-fractions-of(\allowbreak{}cwk\_, cwa\_) iff forall(\allowbreak{}[cwx\_ in carr(\allowbreak{}cwk\_)], forsome(\allowbreak{}[cwu\_ in cwa\_, cwv\_ in cwa\_], not(\allowbreak{}cwv\_ = zero(\allowbreak{}cwk\_)) and (\allowbreak{}mul(\allowbreak{}cwk\_))(\allowbreak{}cwx\_, cwv\_) = cwu\_)))
\end{quote}
also: \texttt{is-fractions-of-mono} \\
mentioned by 14 result(s)

\smallskip

\noindent\textbf{\texttt{is-frechet-structure(fld, m, fam)}}\quad\textit{predicate}\index{is-frechet-structure(fld, m, fam)@\texttt{is-frechet-structure(fld, m, fam)}}\\
reads: (\$2, \$3) is a Frechet space over \$1 \\
defined by \texttt{is-frechet-structure} \\
characterized by \texttt{is-frechet-structure}  \textit{(definitional)}:
\begin{quote}\ttfamily\footnotesize\raggedright\noindent
forall(\allowbreak{}[fld, m, fam], is-frechet-structure(\allowbreak{}fld, m, fam) iff is-seminorm-family(\allowbreak{}fld, m, fam) and forall(\allowbreak{}[seq in fun(\allowbreak{}nn, vec(\allowbreak{}m))], fr-cauchy(\allowbreak{}m, fam, seq) implies forsome(\allowbreak{}[lim in vec(\allowbreak{}m)], fr-conv(\allowbreak{}m, fam, seq, lim))))
\end{quote}
mentioned by 15 result(s)

\smallskip

\noindent\textbf{\texttt{is-fubini-set(xs, ca, mu, ys, cb, nu, cs)}}\quad\textit{predicate}\index{is-fubini-set(xs, ca, mu, ys, cb, nu, cs)@\texttt{is-fubini-set(xs, ca, mu, ys, cb, nu, cs)}}\\
reads: the iterated integrals of the indicator of \$7 over \$1 x \$4 with \$3 and \$6 exist and agree \\
defined by \texttt{is-fubini-set} \\
characterized by \texttt{is-fubini-set}  \textit{(definitional)}:
\begin{quote}\ttfamily\footnotesize\raggedright\noindent
forall(\allowbreak{}[xs, ca, mu, ys, cb, nu, cs], is-fubini-set(\allowbreak{}xs, ca, mu, ys, cb, nu, cs) iff is-measurable-fn(\allowbreak{}xs, ca, x-section-measure(\allowbreak{}xs, ys, nu, cs)) and is-measurable-fn(\allowbreak{}ys, cb, y-section-measure(\allowbreak{}xs, ys, mu, cs)) and integral(\allowbreak{}xs, ca, mu, x-section-measure(\allowbreak{}xs, ys, nu, cs)) = integral(\allowbreak{}ys, cb, nu, y-section-measure(\allowbreak{}xs, ys, mu, cs)))
\end{quote}
mentioned by 9 result(s)

\smallskip

\noindent\textbf{\texttt{is-fun}}\quad\textit{predicate}\index{is-fun@\texttt{is-fun}}\\
a predicate \\
characterized by \texttt{is-fun-def}  \textit{(primitive)}:
\begin{quote}\ttfamily\footnotesize\raggedright\noindent
forall(\allowbreak{}[f], is-fun(\allowbreak{}f) iff forsome(\allowbreak{}[a in set], f in fun(\allowbreak{}a)))
\end{quote}
mentioned by 2 result(s)

\smallskip

\noindent\textbf{\texttt{is-functional-graph(ubv\_, ubg\_)}}\quad\textit{predicate}\index{is-functional-graph(ubv_, ubg_)@\texttt{is-functional-graph(ubv\_, ubg\_)}}\\
reads: \$2 is a functional graph \\
defined by \texttt{is-functional-graph} \\
characterized by \texttt{is-functional-graph}  \textit{(definitional)}:
\begin{quote}\ttfamily\footnotesize\raggedright\noindent
forall(\allowbreak{}[ubv\_, ubg\_], is-functional-graph(\allowbreak{}ubv\_, ubg\_) iff forall(\allowbreak{}[ubx\_ in vec(\allowbreak{}ubv\_), uby\_ in vec(\allowbreak{}ubv\_), ubz\_ in vec(\allowbreak{}ubv\_)], [ubx\_, uby\_] in ubg\_ implies [ubx\_, ubz\_] in ubg\_ implies uby\_ = ubz\_))
\end{quote}
mentioned by 10 result(s)

\smallskip

\noindent\textbf{\texttt{is-gauge-countable(s)}}\quad\textit{predicate}\index{is-gauge-countable(s)@\texttt{is-gauge-countable(s)}}\\
reads: \$1 has a countable pseudometric gauge \\
defined by \texttt{is-gauge-countable} \\
characterized by \texttt{is-gauge-countable}  \textit{(definitional)}:
\begin{quote}\ttfamily\footnotesize\raggedright\noindent
forall(\allowbreak{}[s], is-gauge-countable(\allowbreak{}s) iff is-top-space(\allowbreak{}s) and forsome(\allowbreak{}[fam], is-countable-pseudometric-family(\allowbreak{}fam, pts(\allowbreak{}s)) and pseudo-gauge-top(\allowbreak{}fam, pts(\allowbreak{}s)) == s))
\end{quote}
mentioned by 9 result(s)

\smallskip

\noindent\textbf{\texttt{is-gcd-set(ggs\_, ggt\_, ggd\_)}}\quad\textit{predicate}\index{is-gcd-set(ggs_, ggt_, ggd_)@\texttt{is-gcd-set(ggs\_, ggt\_, ggd\_)}}\\
reads: \$3 is a greatest common divisor of \$2 in \$1 \\
defined by \texttt{is-gcd-set} \\
characterized by \texttt{is-gcd-set}  \textit{(definitional)}:
\begin{quote}\ttfamily\footnotesize\raggedright\noindent
forall(\allowbreak{}[ggs\_, ggt\_, ggd\_], is-gcd-set(\allowbreak{}ggs\_, ggt\_, ggd\_) iff ggd\_ in carr(\allowbreak{}ggs\_) and forall(\allowbreak{}[ggx\_ in ggt\_], divides(\allowbreak{}ggs\_, ggd\_, ggx\_)) and forall(\allowbreak{}[gge\_ in carr(\allowbreak{}ggs\_)], forall(\allowbreak{}[ggx\_ in ggt\_], divides(\allowbreak{}ggs\_, gge\_, ggx\_)) implies divides(\allowbreak{}ggs\_, gge\_, ggd\_)))
\end{quote}
mentioned by 4 result(s)

\smallskip

\noindent\textbf{\texttt{is-generator(g, x)}}\quad\textit{predicate}\index{is-generator(g, x)@\texttt{is-generator(g, x)}}\\
reads: \$2 generates the group \$1 \\
defined by \texttt{is-generator} \\
characterized by \texttt{is-generator}  \textit{(definitional)}:
\begin{quote}\ttfamily\footnotesize\raggedright\noindent
forall(\allowbreak{}[g, x], is-generator(\allowbreak{}g, x) iff x in carr(\allowbreak{}g) and forall(\allowbreak{}[geny\_ in carr(\allowbreak{}g)], forsome(\allowbreak{}[genk\_ in zz], geny\_ = group-zpow(\allowbreak{}g, x, genk\_))))
\end{quote}
mentioned by 19 result(s)

\smallskip

\noindent\textbf{\texttt{is-group(s)}}\quad\textit{predicate}\index{is-group(s)@\texttt{is-group(s)}}\\
reads: a group \\
a predicate \\
characterized by \texttt{is-group}  \textit{(definitional)}:
\begin{quote}\ttfamily\footnotesize\raggedright\noindent
forall(\allowbreak{}[s], is-group(\allowbreak{}s) iff length(\allowbreak{}s) = 4 and carr(\allowbreak{}s) in set and opr(\allowbreak{}s) in fun(\allowbreak{}cartesian(\allowbreak{}carr(\allowbreak{}s), carr(\allowbreak{}s)), carr(\allowbreak{}s)) and iden(\allowbreak{}s) in carr(\allowbreak{}s) and inv(\allowbreak{}s) in fun(\allowbreak{}carr(\allowbreak{}s), carr(\allowbreak{}s)) and is-associative(\allowbreak{}opr(\allowbreak{}s), carr(\allowbreak{}s)) and is-identity(\allowbreak{}opr(\allowbreak{}s), iden(\allowbreak{}s), carr(\allowbreak{}s)) and has-inverses(\allowbreak{}opr(\allowbreak{}s), iden(\allowbreak{}s), inv(\allowbreak{}s), carr(\allowbreak{}s)) and s in set)
\end{quote}
also: \texttt{is-group-metric} \texttt{is-group-topology} \texttt{field-multiplicative-group-is-group} \texttt{abelian-group-is-group-ring-additive-ag} \texttt{abelian-group-is-group-ring-additive-ag-uba-as-ring} \texttt{abelian-group-is-group-ring-additive-ag-uba-as-ring-uiba-as-uba} \texttt{is-group-is-set} \texttt{lft-group-is-group} ... (15 in all) \\
mentioned by 251 result(s)

\smallskip

\noindent\textbf{\texttt{is-group-metric(mgm\_, mgg\_)}}\quad\textit{predicate}\index{is-group-metric(mgm_, mgg_)@\texttt{is-group-metric(mgm\_, mgg\_)}}\\
reads: \$1 is a metric for the group \$2 making its operations continuous \\
defined by \texttt{is-group-metric} \\
characterized by \texttt{is-group-metric}  \textit{(definitional)}:
\begin{quote}\ttfamily\footnotesize\raggedright\noindent
forall(\allowbreak{}[mgm\_, mgg\_], is-group-metric(\allowbreak{}mgm\_, mgg\_) iff is-metric-space(\allowbreak{}mgm\_) and is-group(\allowbreak{}mgg\_) and pts(\allowbreak{}mgm\_) = carr(\allowbreak{}mgg\_) and forall(\allowbreak{}[mgu\_ in carr(\allowbreak{}mgg\_), mgv\_ in carr(\allowbreak{}mgg\_), mge\_], pos-rr(\allowbreak{}mge\_) implies forsome(\allowbreak{}[mgd\_], pos-rr(\allowbreak{}mgd\_) and forall(\allowbreak{}[mgp\_ in carr(\allowbreak{}mgg\_), mgq\_ in carr(\allowbreak{}mgg\_)], (\allowbreak{}dist(\allowbreak{}mgm\_))(\allowbreak{}mgu\_, mgp\_) <= mgd\_ implies (\allowbreak{}dist(\allowbreak{}mgm\_))(\allowbreak{}mgv\_, mgq\_) <= mgd\_ implies (\allowbreak{}dist(\allowbreak{}mgm\_))(\allowbreak{}(\allowbreak{}opr(\allowbreak{}mgg\_))(\allowbreak{}mgu\_, mgv\_), (\allowbreak{}opr(\allowbreak{}mgg\_))(\allowbreak{}mgp\_, mgq\_)) <= mge\_))) and is-continuous(\allowbreak{}mgm\_, mgm\_, inv(\allowbreak{}mgg\_)))
\end{quote}
mentioned by 8 result(s)

\smallskip

\noindent\textbf{\texttt{is-group-norm(nm, op, invop, unit, crr)}}\quad\textit{predicate}\index{is-group-norm(nm, op, invop, unit, crr)@\texttt{is-group-norm(nm, op, invop, unit, crr)}}\\
reads: \$1 is a group norm for \$2 on \$5 \\
defined by \texttt{is-group-norm} \\
characterized by \texttt{is-group-norm}  \textit{(definitional)}:
\begin{quote}\ttfamily\footnotesize\raggedright\noindent
forall(\allowbreak{}[nm, op, invop, unit, crr], is-group-norm(\allowbreak{}nm, op, invop, unit, crr) iff nm in fun(\allowbreak{}crr, rr) and forall(\allowbreak{}[u in crr], 0 <= nm(\allowbreak{}u) and (\allowbreak{}nm(\allowbreak{}u) = 0 iff u = unit) and nm(\allowbreak{}invop(\allowbreak{}u)) = nm(\allowbreak{}u) and forall(\allowbreak{}[v in crr], nm(\allowbreak{}op(\allowbreak{}u, v)) <= nm(\allowbreak{}u) + nm(\allowbreak{}v))))
\end{quote}
mentioned by 5 result(s)

\smallskip

\noindent\textbf{\texttt{is-group-segment(gss\_, gsg\_, gsa\_, gsp\_)}}\quad\textit{predicate}\index{is-group-segment(gss_, gsg_, gsa_, gsp_)@\texttt{is-group-segment(gss\_, gsg\_, gsa\_, gsp\_)}}\\
reads: \$4 is a one-parameter group segment of length \$3 in \$2 on \$1 \\
defined by \texttt{is-group-segment} \\
characterized by \texttt{is-group-segment}  \textit{(definitional)}:
\begin{quote}\ttfamily\footnotesize\raggedright\noindent
forall(\allowbreak{}[gss\_, gsg\_, gsa\_, gsp\_], is-group-segment(\allowbreak{}gss\_, gsg\_, gsa\_, gsp\_) iff pos-rr(\allowbreak{}gsa\_) and gsp\_ in fun(\allowbreak{}ccint(\allowbreak{}0, gsa\_), carr(\allowbreak{}gsg\_)) and forall(\allowbreak{}[gsr\_ in ccint(\allowbreak{}0, gsa\_), gsu\_ in ccint(\allowbreak{}0, gsa\_)], gsr\_ + gsu\_ <= gsa\_ implies gsp\_(\allowbreak{}gsr\_ + gsu\_) = (\allowbreak{}opr(\allowbreak{}gsg\_))(\allowbreak{}gsp\_(\allowbreak{}gsr\_), gsp\_(\allowbreak{}gsu\_))) and forall(\allowbreak{}[gst\_ in ccint(\allowbreak{}0, gsa\_), gsw\_ in opens(\allowbreak{}gss\_)], gsp\_(\allowbreak{}gst\_) in gsw\_ implies forsome(\allowbreak{}[gsd\_], pos-rr(\allowbreak{}gsd\_) and forall(\allowbreak{}[gsv\_ in ccint(\allowbreak{}0, gsa\_)], abs(\allowbreak{}gsv\_ - gst\_) < gsd\_ implies gsp\_(\allowbreak{}gsv\_) in gsw\_))))
\end{quote}
mentioned by 10 result(s)

\smallskip

\noindent\textbf{\texttt{is-group-topology(tgs\_, tgg\_)}}\quad\textit{predicate}\index{is-group-topology(tgs_, tgg_)@\texttt{is-group-topology(tgs\_, tgg\_)}}\\
reads: \$1 is a group topology for the group \$2 \\
defined by \texttt{is-group-topology} \\
characterized by \texttt{is-group-topology}  \textit{(definitional)}:
\begin{quote}\ttfamily\footnotesize\raggedright\noindent
forall(\allowbreak{}[tgs\_, tgg\_], is-group-topology(\allowbreak{}tgs\_, tgg\_) iff is-top-space(\allowbreak{}tgs\_) and is-group(\allowbreak{}tgg\_) and pts(\allowbreak{}tgs\_) = carr(\allowbreak{}tgg\_) and forall(\allowbreak{}[tga\_ in carr(\allowbreak{}tgg\_), tgb\_ in carr(\allowbreak{}tgg\_), tgw\_ in opens(\allowbreak{}tgs\_)], (\allowbreak{}opr(\allowbreak{}tgg\_))(\allowbreak{}tga\_, (\allowbreak{}inv(\allowbreak{}tgg\_))(\allowbreak{}tgb\_)) in tgw\_ implies forsome(\allowbreak{}[tgu\_ in opens(\allowbreak{}tgs\_)], tga\_ in tgu\_ and forsome(\allowbreak{}[tgv\_ in opens(\allowbreak{}tgs\_)], tgb\_ in tgv\_ and forall(\allowbreak{}[tgp\_ in tgu\_, tgq\_ in tgv\_], (\allowbreak{}opr(\allowbreak{}tgg\_))(\allowbreak{}tgp\_, (\allowbreak{}inv(\allowbreak{}tgg\_))(\allowbreak{}tgq\_)) in tgw\_)))))
\end{quote}
mentioned by 70 result(s)

\smallskip

\noindent\textbf{\texttt{is-haar-integral(hig\_, hiu\_)}}\quad\textit{predicate}\index{is-haar-integral(hig_, hiu_)@\texttt{is-haar-integral(hig\_, hiu\_)}}\\
reads: \$2 is a Haar integral of the compact metric group \$1 \\
defined by \texttt{is-haar-integral} \\
characterized by \texttt{is-haar-integral}  \textit{(definitional)}:
\begin{quote}\ttfamily\footnotesize\raggedright\noindent
forall(\allowbreak{}[hig\_, hiu\_], is-haar-integral(\allowbreak{}hig\_, hiu\_) iff is-metric-group(\allowbreak{}hig\_) and is-compact(\allowbreak{}metric-group-ms(\allowbreak{}hig\_)) and is-positive-functional(\allowbreak{}metric-group-ms(\allowbreak{}hig\_), hiu\_) and hiu\_(\allowbreak{}vnb-lambda(\allowbreak{}hix\_, pts(\allowbreak{}metric-group-ms(\allowbreak{}hig\_)), 1)) = 1 and forall(\allowbreak{}[hia\_ in pts(\allowbreak{}metric-group-ms(\allowbreak{}hig\_)), hif\_ in cont-fns(\allowbreak{}metric-group-ms(\allowbreak{}hig\_))], hiu\_(\allowbreak{}vnb-lambda(\allowbreak{}hix\_, pts(\allowbreak{}metric-group-ms(\allowbreak{}hig\_)), hif\_(\allowbreak{}(\allowbreak{}opr(\allowbreak{}hig\_))(\allowbreak{}hia\_, hix\_)))) = hiu\_(\allowbreak{}hif\_)))
\end{quote}
also: \texttt{haar-mean-is-haar-integral} \\
mentioned by 21 result(s)

\smallskip

\noindent\textbf{\texttt{is-hausdorff(s)}}\quad\textit{predicate}\index{is-hausdorff(s)@\texttt{is-hausdorff(s)}}\\
reads: \$1 is Hausdorff \\
defined by \texttt{is-hausdorff} \\
characterized by \texttt{is-hausdorff}  \textit{(definitional)}:
\begin{quote}\ttfamily\footnotesize\raggedright\noindent
forall(\allowbreak{}[s], is-hausdorff(\allowbreak{}s) iff is-top-space(\allowbreak{}s) and forall(\allowbreak{}[a, b], a in pts(\allowbreak{}s) implies b in pts(\allowbreak{}s) implies not(\allowbreak{}a == b) implies forsome(\allowbreak{}[u, v], u in opens(\allowbreak{}s) and v in opens(\allowbreak{}s) and a in u and b in v and intersection(\allowbreak{}u, v) == empty-set)))
\end{quote}
mentioned by 60 result(s)

\smallskip

\noindent\textbf{\texttt{is-hausdorff-dim(mhm\_, mhs\_, mhl\_)}}\quad\textit{predicate}\index{is-hausdorff-dim(mhm_, mhs_, mhl_)@\texttt{is-hausdorff-dim(mhm\_, mhs\_, mhl\_)}}\\
reads: \$2 has Hausdorff dimension \$3 in \$1 \\
defined by \texttt{is-hausdorff-dim} \\
characterized by \texttt{is-hausdorff-dim}  \textit{(definitional)}:
\begin{quote}\ttfamily\footnotesize\raggedright\noindent
forall(\allowbreak{}[mhm\_, mhs\_, mhl\_], is-hausdorff-dim(\allowbreak{}mhm\_, mhs\_, mhl\_) iff mhl\_ in rr and 0 <= mhl\_ and forall(\allowbreak{}[mhr\_ in rr], mhl\_ < mhr\_ implies (\allowbreak{}hausdorff(\allowbreak{}mhm\_, mhr\_))(\allowbreak{}mhs\_) = 0) and forall(\allowbreak{}[mhr\_ in rr], 0 < mhr\_ implies mhr\_ < mhl\_ implies (\allowbreak{}hausdorff(\allowbreak{}mhm\_, mhr\_))(\allowbreak{}mhs\_) = pos-inf))
\end{quote}
mentioned by 6 result(s)

\smallskip

\noindent\textbf{\texttt{is-hereditary(ca, cb)}}\quad\textit{predicate}\index{is-hereditary(ca, cb)@\texttt{is-hereditary(ca, cb)}}\\
reads: \$2 is a hereditary subset of \$1 \\
defined by \texttt{is-hereditary} \\
characterized by \texttt{is-hereditary}  \textit{(definitional)}:
\begin{quote}\ttfamily\footnotesize\raggedright\noindent
forall(\allowbreak{}[ca, cb], is-hereditary(\allowbreak{}ca, cb) iff cb subset ca and forall(\allowbreak{}[hde\_ in cb, hdf\_ in ca], hdf\_ subset hde\_ implies hdf\_ in cb))
\end{quote}
mentioned by 4 result(s)

\smallskip

\noindent\textbf{\texttt{is-hermitian-form(ipfm\_, ipff\_)}}\quad\textit{predicate}\index{is-hermitian-form(ipfm_, ipff_)@\texttt{is-hermitian-form(ipfm\_, ipff\_)}}\\
reads: \$2 is a hermitian sesquilinear form on \$1 \\
defined by \texttt{is-hermitian-form} \\
characterized by \texttt{is-hermitian-form}  \textit{(definitional)}:
\begin{quote}\ttfamily\footnotesize\raggedright\noindent
forall(\allowbreak{}[ipfm\_, ipff\_], is-hermitian-form(\allowbreak{}ipfm\_, ipff\_) iff ipff\_ in fun(\allowbreak{}cartesian(\allowbreak{}vec(\allowbreak{}ipfm\_), vec(\allowbreak{}ipfm\_)), cc) and forall(\allowbreak{}[ipfx\_ in vec(\allowbreak{}ipfm\_), ipfy\_ in vec(\allowbreak{}ipfm\_), ipfz\_ in vec(\allowbreak{}ipfm\_)], ipff\_(\allowbreak{}(\allowbreak{}vadd(\allowbreak{}ipfm\_))(\allowbreak{}ipfx\_, ipfy\_), ipfz\_) = ipff\_(\allowbreak{}ipfx\_, ipfz\_) + ipff\_(\allowbreak{}ipfy\_, ipfz\_)) and forall(\allowbreak{}[ipfa\_ in cc, ipfx\_ in vec(\allowbreak{}ipfm\_), ipfy\_ in vec(\allowbreak{}ipfm\_)], ipff\_(\allowbreak{}(\allowbreak{}act(\allowbreak{}ipfm\_))(\allowbreak{}ipfa\_, ipfx\_), ipfy\_) = ipfa\_ * ipff\_(\allowbreak{}ipfx\_, ipfy\_)) and forall(\allowbreak{}[ipfx\_ in vec(\allowbreak{}ipfm\_), ipfy\_ in vec(\allowbreak{}ipfm\_)], ipff\_(\allowbreak{}ipfy\_, ipfx\_) = conjugate(\allowbreak{}ipff\_(\allowbreak{}ipfx\_, ipfy\_))))
\end{quote}
mentioned by 17 result(s)

\smallskip

\noindent\textbf{\texttt{is-hilbert-schmidt(kqv\_, kqw\_, kqt\_)}}\quad\textit{predicate}\index{is-hilbert-schmidt(kqv_, kqw_, kqt_)@\texttt{is-hilbert-schmidt(kqv\_, kqw\_, kqt\_)}}\\
reads: \$3 is a Hilbert-Schmidt operator from \$1 to \$2 \\
defined by \texttt{is-hilbert-schmidt} \\
characterized by \texttt{is-hilbert-schmidt}  \textit{(definitional)}:
\begin{quote}\ttfamily\footnotesize\raggedright\noindent
forall(\allowbreak{}[kqv\_, kqw\_, kqt\_], is-hilbert-schmidt(\allowbreak{}kqv\_, kqw\_, kqt\_) iff kqt\_ in bounded-ops(\allowbreak{}kqv\_, kqw\_) and forsome(\allowbreak{}[kqi\_ in set, kqe\_ in fun(\allowbreak{}kqi\_, vec(\allowbreak{}kqv\_))], is-orthonormal-basis(\allowbreak{}kqv\_, kqi\_, kqe\_) and hs-sum(\allowbreak{}kqv\_, kqw\_, kqi\_, kqe\_, kqt\_) in rr))
\end{quote}
mentioned by 48 result(s)

\smallskip

\noindent\textbf{\texttt{is-hilbert-space(v)}}\quad\textit{predicate}\index{is-hilbert-space(v)@\texttt{is-hilbert-space(v)}}\\
reads: a Hilbert space \\
defined by \texttt{is-hilbert-space} \\
characterized by \texttt{is-hilbert-space}  \textit{(definitional)}:
\begin{quote}\ttfamily\footnotesize\raggedright\noindent
forall(\allowbreak{}[v], is-hilbert-space(\allowbreak{}v) iff is-complex-inner-product-space(\allowbreak{}v) and is-complete(\allowbreak{}ip-metric-space(\allowbreak{}v)))
\end{quote}
also: \texttt{l2-is-hilbert-space} \texttt{cl2-hilbert-is-hilbert-space} \\
mentioned by 309 result(s)

\smallskip

\noindent\textbf{\texttt{is-holomorphic-automorphism(hau\_, hap\_)}}\quad\textit{predicate}\index{is-holomorphic-automorphism(hau_, hap_)@\texttt{is-holomorphic-automorphism(hau\_, hap\_)}}\\
reads: \$2 is an automorphism of \$1 \\
defined by \texttt{is-holomorphic-automorphism} \\
characterized by \texttt{is-holomorphic-automorphism}  \textit{(definitional)}:
\begin{quote}\ttfamily\footnotesize\raggedright\noindent
forall(\allowbreak{}[hau\_, hap\_], is-holomorphic-automorphism(\allowbreak{}hau\_, hap\_) iff hap\_ in bijection(\allowbreak{}hau\_, hau\_) and holomorphic-on(\allowbreak{}hau\_, hap\_))
\end{quote}
mentioned by 5 result(s)

\smallskip

\noindent\textbf{\texttt{is-hom-abelian-group(a, b, f)}}\quad\textit{predicate}\index{is-hom-abelian-group(a, b, f)@\texttt{is-hom-abelian-group(a, b, f)}}\\
reads: \$3 is a homomorphism from \$1 to \$2 \\
a predicate \\
characterized by \texttt{is-hom-abelian-group-def}  \textit{(definitional)}:
\begin{quote}\ttfamily\footnotesize\raggedright\noindent
forall(\allowbreak{}[a, b, f], is-hom-abelian-group(\allowbreak{}a, b, f) iff is-abelian-group(\allowbreak{}a) and is-abelian-group(\allowbreak{}b) and f in fun(\allowbreak{}carr(\allowbreak{}a), carr(\allowbreak{}b)) and forall(\allowbreak{}[x1\_ in carr(\allowbreak{}a), x2\_ in carr(\allowbreak{}a)], f(\allowbreak{}(\allowbreak{}opr(\allowbreak{}a))(\allowbreak{}x1\_, x2\_)) = (\allowbreak{}opr(\allowbreak{}b))(\allowbreak{}f(\allowbreak{}x1\_), f(\allowbreak{}x2\_))) and f(\allowbreak{}iden(\allowbreak{}a)) = iden(\allowbreak{}b) and forall(\allowbreak{}[x1\_ in carr(\allowbreak{}a)], f(\allowbreak{}(\allowbreak{}inv(\allowbreak{}a))(\allowbreak{}x1\_)) = (\allowbreak{}inv(\allowbreak{}b))(\allowbreak{}f(\allowbreak{}x1\_))))
\end{quote}
mentioned by 27 result(s)

\smallskip

\noindent\textbf{\texttt{is-hom-bundle(a, b, f)}}\quad\textit{predicate}\index{is-hom-bundle(a, b, f)@\texttt{is-hom-bundle(a, b, f)}}\\
reads: \$3 is a homomorphism from \$1 to \$2 \\
a predicate \\
characterized by \texttt{is-hom-bundle-def}  \textit{(definitional)}:
\begin{quote}\ttfamily\footnotesize\raggedright\noindent
forall(\allowbreak{}[a, b, f], is-hom-bundle(\allowbreak{}a, b, f) iff is-bundle(\allowbreak{}a) and is-bundle(\allowbreak{}b) and f in fun(\allowbreak{}bnodes(\allowbreak{}a), bnodes(\allowbreak{}b)) and bspace(\allowbreak{}a) = bspace(\allowbreak{}b) and forall(\allowbreak{}[x1\_ in bnodes(\allowbreak{}a), x2\_ in bnodes(\allowbreak{}a)], [x1\_, x2\_] in bsend(\allowbreak{}a) implies [f(\allowbreak{}x1\_), f(\allowbreak{}x2\_)] in bsend(\allowbreak{}b)))
\end{quote}
mentioned by 2 result(s)

\smallskip

\noindent\textbf{\texttt{is-hom-c-metric-space(a, b, f)}}\quad\textit{predicate}\index{is-hom-c-metric-space(a, b, f)@\texttt{is-hom-c-metric-space(a, b, f)}}\\
reads: \$3 is a homomorphism from \$1 to \$2 \\
a predicate \\
characterized by \texttt{is-hom-c-metric-space-def}  \textit{(definitional)}:
\begin{quote}\ttfamily\footnotesize\raggedright\noindent
forall(\allowbreak{}[a, b, f], is-hom-c-metric-space(\allowbreak{}a, b, f) iff is-c-metric-space(\allowbreak{}a) and is-c-metric-space(\allowbreak{}b) and f in fun(\allowbreak{}pts(\allowbreak{}a), pts(\allowbreak{}b)) and forall(\allowbreak{}[x1\_ in nn], (\allowbreak{}dists(\allowbreak{}a))(\allowbreak{}x1\_) = (\allowbreak{}dists(\allowbreak{}b))(\allowbreak{}x1\_)))
\end{quote}
mentioned by 6 result(s)

\smallskip

\noindent\textbf{\texttt{is-hom-comm-monoid(a, b, f)}}\quad\textit{predicate}\index{is-hom-comm-monoid(a, b, f)@\texttt{is-hom-comm-monoid(a, b, f)}}\\
reads: \$3 is a homomorphism from \$1 to \$2 \\
a predicate \\
characterized by \texttt{is-hom-comm-monoid-def}  \textit{(definitional)}:
\begin{quote}\ttfamily\footnotesize\raggedright\noindent
forall(\allowbreak{}[a, b, f], is-hom-comm-monoid(\allowbreak{}a, b, f) iff is-comm-monoid(\allowbreak{}a) and is-comm-monoid(\allowbreak{}b) and f in fun(\allowbreak{}carr(\allowbreak{}a), carr(\allowbreak{}b)) and forall(\allowbreak{}[x1\_ in carr(\allowbreak{}a), x2\_ in carr(\allowbreak{}a)], f(\allowbreak{}(\allowbreak{}opr(\allowbreak{}a))(\allowbreak{}x1\_, x2\_)) = (\allowbreak{}opr(\allowbreak{}b))(\allowbreak{}f(\allowbreak{}x1\_), f(\allowbreak{}x2\_))) and f(\allowbreak{}iden(\allowbreak{}a)) = iden(\allowbreak{}b))
\end{quote}
mentioned by 7 result(s)

\smallskip

\noindent\textbf{\texttt{is-hom-commutative-banach-algebra(a, b, f)}}\quad\textit{predicate}\index{is-hom-commutative-banach-algebra(a, b, f)@\texttt{is-hom-commutative-banach-algebra(a, b, f)}}\\
reads: \$3 is a homomorphism from \$1 to \$2 \\
a predicate \\
characterized by \texttt{is-hom-commutative-banach-algebra-def}  \textit{(definitional)}:
\begin{quote}\ttfamily\footnotesize\raggedright\noindent
forall(\allowbreak{}[a, b, f], is-hom-commutative-banach-algebra(\allowbreak{}a, b, f) iff is-commutative-banach-algebra(\allowbreak{}a) and is-commutative-banach-algebra(\allowbreak{}b) and is-hom-unital-banach-algebra(\allowbreak{}a, b, f))
\end{quote}
mentioned by 2 result(s)

\smallskip

\noindent\textbf{\texttt{is-hom-commutative-c-star-algebra(a, b, f)}}\quad\textit{predicate}\index{is-hom-commutative-c-star-algebra(a, b, f)@\texttt{is-hom-commutative-c-star-algebra(a, b, f)}}\\
reads: \$3 is a homomorphism from \$1 to \$2 \\
a predicate \\
characterized by \texttt{is-hom-commutative-c-star-algebra-def}  \textit{(definitional)}:
\begin{quote}\ttfamily\footnotesize\raggedright\noindent
forall(\allowbreak{}[a, b, f], is-hom-commutative-c-star-algebra(\allowbreak{}a, b, f) iff is-commutative-c-star-algebra(\allowbreak{}a) and is-commutative-c-star-algebra(\allowbreak{}b) and is-hom-unital-c-star-algebra(\allowbreak{}a, b, f))
\end{quote}
mentioned by 2 result(s)

\smallskip

\noindent\textbf{\texttt{is-hom-commutative-ring(a, b, f)}}\quad\textit{predicate}\index{is-hom-commutative-ring(a, b, f)@\texttt{is-hom-commutative-ring(a, b, f)}}\\
reads: \$3 is a homomorphism from \$1 to \$2 \\
a predicate \\
characterized by \texttt{is-hom-commutative-ring-def}  \textit{(definitional)}:
\begin{quote}\ttfamily\footnotesize\raggedright\noindent
forall(\allowbreak{}[a, b, f], is-hom-commutative-ring(\allowbreak{}a, b, f) iff is-commutative-ring(\allowbreak{}a) and is-commutative-ring(\allowbreak{}b) and is-hom-ring(\allowbreak{}a, b, f))
\end{quote}
mentioned by 14 result(s)

\smallskip

\noindent\textbf{\texttt{is-hom-complex-banach-algebra(a, b, f)}}\quad\textit{predicate}\index{is-hom-complex-banach-algebra(a, b, f)@\texttt{is-hom-complex-banach-algebra(a, b, f)}}\\
reads: \$3 is a homomorphism from \$1 to \$2 \\
a predicate \\
characterized by \texttt{is-hom-complex-banach-algebra-def}  \textit{(definitional)}:
\begin{quote}\ttfamily\footnotesize\raggedright\noindent
forall(\allowbreak{}[a, b, f], is-hom-complex-banach-algebra(\allowbreak{}a, b, f) iff is-complex-banach-algebra(\allowbreak{}a) and is-complex-banach-algebra(\allowbreak{}b) and is-hom-complex-normed-algebra(\allowbreak{}a, b, f))
\end{quote}
mentioned by 4 result(s)

\smallskip

\noindent\textbf{\texttt{is-hom-complex-banach-space(a, b, f)}}\quad\textit{predicate}\index{is-hom-complex-banach-space(a, b, f)@\texttt{is-hom-complex-banach-space(a, b, f)}}\\
reads: \$3 is a homomorphism from \$1 to \$2 \\
a predicate \\
characterized by \texttt{is-hom-complex-banach-space-def}  \textit{(definitional)}:
\begin{quote}\ttfamily\footnotesize\raggedright\noindent
forall(\allowbreak{}[a, b, f], is-hom-complex-banach-space(\allowbreak{}a, b, f) iff is-complex-banach-space(\allowbreak{}a) and is-complex-banach-space(\allowbreak{}b) and is-hom-complex-normed-vector-space(\allowbreak{}a, b, f))
\end{quote}
mentioned by 2 result(s)

\smallskip

\noindent\textbf{\texttt{is-hom-complex-inner-product-space(a, b, f)}}\quad\textit{predicate}\index{is-hom-complex-inner-product-space(a, b, f)@\texttt{is-hom-complex-inner-product-space(a, b, f)}}\\
reads: \$3 is a homomorphism from \$1 to \$2 \\
a predicate \\
characterized by \texttt{is-hom-complex-inner-product-space-def}  \textit{(definitional)}:
\begin{quote}\ttfamily\footnotesize\raggedright\noindent
forall(\allowbreak{}[a, b, f], is-hom-complex-inner-product-space(\allowbreak{}a, b, f) iff is-complex-inner-product-space(\allowbreak{}a) and is-complex-inner-product-space(\allowbreak{}b) and f in fun(\allowbreak{}vec(\allowbreak{}a), vec(\allowbreak{}b)) and scal(\allowbreak{}a) = scal(\allowbreak{}b) and forall(\allowbreak{}[x1\_ in vec(\allowbreak{}a), x2\_ in vec(\allowbreak{}a)], f(\allowbreak{}(\allowbreak{}vadd(\allowbreak{}a))(\allowbreak{}x1\_, x2\_)) = (\allowbreak{}vadd(\allowbreak{}b))(\allowbreak{}f(\allowbreak{}x1\_), f(\allowbreak{}x2\_))) and f(\allowbreak{}vzero(\allowbreak{}a)) = vzero(\allowbreak{}b) and forall(\allowbreak{}[x1\_ in vec(\allowbreak{}a)], f(\allowbreak{}(\allowbreak{}vneg(\allowbreak{}a))(\allowbreak{}x1\_)) = (\allowbreak{}vneg(\allowbreak{}b))(\allowbreak{}f(\allowbreak{}x1\_))) and forall(\allowbreak{}[x1\_ in carr(\allowbreak{}scal(\allowbreak{}a)), x2\_ in vec(\allowbreak{}a)], f(\allowbreak{}(\allowbreak{}act(\allowbreak{}a))(\allowbreak{}x1\_, x2\_)) = (\allowbreak{}act(\allowbreak{}b))(\allowbreak{}x1\_, f(\allowbreak{}x2\_))) and forall(\allowbreak{}[x1\_ in vec(\allowbreak{}a), x2\_ in vec(\allowbreak{}a)], (\allowbreak{}ip(\allowbreak{}a))(\allowbreak{}x1\_, x2\_) = (\allowbreak{}ip(\allowbreak{}b))(\allowbreak{}f(\allowbreak{}x1\_), f(\allowbreak{}x2\_))))
\end{quote}
mentioned by 7 result(s)

\smallskip

\noindent\textbf{\texttt{is-hom-complex-normed-algebra(a, b, f)}}\quad\textit{predicate}\index{is-hom-complex-normed-algebra(a, b, f)@\texttt{is-hom-complex-normed-algebra(a, b, f)}}\\
reads: \$3 is a homomorphism from \$1 to \$2 \\
a predicate \\
characterized by \texttt{is-hom-complex-normed-algebra-def}  \textit{(definitional)}:
\begin{quote}\ttfamily\footnotesize\raggedright\noindent
forall(\allowbreak{}[a, b, f], is-hom-complex-normed-algebra(\allowbreak{}a, b, f) iff is-complex-normed-algebra(\allowbreak{}a) and is-complex-normed-algebra(\allowbreak{}b) and f in fun(\allowbreak{}vec(\allowbreak{}a), vec(\allowbreak{}b)) and scal(\allowbreak{}a) = scal(\allowbreak{}b) and forall(\allowbreak{}[x1\_ in vec(\allowbreak{}a), x2\_ in vec(\allowbreak{}a)], f(\allowbreak{}(\allowbreak{}vadd(\allowbreak{}a))(\allowbreak{}x1\_, x2\_)) = (\allowbreak{}vadd(\allowbreak{}b))(\allowbreak{}f(\allowbreak{}x1\_), f(\allowbreak{}x2\_))) and f(\allowbreak{}vzero(\allowbreak{}a)) = vzero(\allowbreak{}b) and forall(\allowbreak{}[x1\_ in vec(\allowbreak{}a)], f(\allowbreak{}(\allowbreak{}vneg(\allowbreak{}a))(\allowbreak{}x1\_)) = (\allowbreak{}vneg(\allowbreak{}b))(\allowbreak{}f(\allowbreak{}x1\_))) and forall(\allowbreak{}[x1\_ in carr(\allowbreak{}scal(\allowbreak{}a)), x2\_ in vec(\allowbreak{}a)], f(\allowbreak{}(\allowbreak{}act(\allowbreak{}a))(\allowbreak{}x1\_, x2\_)) = (\allowbreak{}act(\allowbreak{}b))(\allowbreak{}x1\_, f(\allowbreak{}x2\_))) and forall(\allowbreak{}[x1\_ in vec(\allowbreak{}a)], (\allowbreak{}vnrm(\allowbreak{}a))(\allowbreak{}x1\_) = (\allowbreak{}vnrm(\allowbreak{}b))(\allowbreak{}f(\allowbreak{}x1\_))) and forall(\allowbreak{}[x1\_ in vec(\allowbreak{}a), x2\_ in vec(\allowbreak{}a)], f(\allowbreak{}(\allowbreak{}vmul(\allowbreak{}a))(\allowbreak{}x1\_, x2\_)) = (\allowbreak{}vmul(\allowbreak{}b))(\allowbreak{}f(\allowbreak{}x1\_), f(\allowbreak{}x2\_))))
\end{quote}
mentioned by 6 result(s)

\smallskip

\noindent\textbf{\texttt{is-hom-complex-normed-vector-space(a, b, f)}}\quad\textit{predicate}\index{is-hom-complex-normed-vector-space(a, b, f)@\texttt{is-hom-complex-normed-vector-space(a, b, f)}}\\
reads: \$3 is a homomorphism from \$1 to \$2 \\
a predicate \\
characterized by \texttt{is-hom-complex-normed-vector-space-def}  \textit{(definitional)}:
\begin{quote}\ttfamily\footnotesize\raggedright\noindent
forall(\allowbreak{}[a, b, f], is-hom-complex-normed-vector-space(\allowbreak{}a, b, f) iff is-complex-normed-vector-space(\allowbreak{}a) and is-complex-normed-vector-space(\allowbreak{}b) and f in fun(\allowbreak{}vec(\allowbreak{}a), vec(\allowbreak{}b)) and scal(\allowbreak{}a) = scal(\allowbreak{}b) and forall(\allowbreak{}[x1\_ in vec(\allowbreak{}a), x2\_ in vec(\allowbreak{}a)], f(\allowbreak{}(\allowbreak{}vadd(\allowbreak{}a))(\allowbreak{}x1\_, x2\_)) = (\allowbreak{}vadd(\allowbreak{}b))(\allowbreak{}f(\allowbreak{}x1\_), f(\allowbreak{}x2\_))) and f(\allowbreak{}vzero(\allowbreak{}a)) = vzero(\allowbreak{}b) and forall(\allowbreak{}[x1\_ in vec(\allowbreak{}a)], f(\allowbreak{}(\allowbreak{}vneg(\allowbreak{}a))(\allowbreak{}x1\_)) = (\allowbreak{}vneg(\allowbreak{}b))(\allowbreak{}f(\allowbreak{}x1\_))) and forall(\allowbreak{}[x1\_ in carr(\allowbreak{}scal(\allowbreak{}a)), x2\_ in vec(\allowbreak{}a)], f(\allowbreak{}(\allowbreak{}act(\allowbreak{}a))(\allowbreak{}x1\_, x2\_)) = (\allowbreak{}act(\allowbreak{}b))(\allowbreak{}x1\_, f(\allowbreak{}x2\_))) and forall(\allowbreak{}[x1\_ in vec(\allowbreak{}a)], (\allowbreak{}vnrm(\allowbreak{}a))(\allowbreak{}x1\_) = (\allowbreak{}vnrm(\allowbreak{}b))(\allowbreak{}f(\allowbreak{}x1\_))) and forall(\allowbreak{}[x1\_ in vec(\allowbreak{}a)], f(\allowbreak{}(\allowbreak{}jop(\allowbreak{}a))(\allowbreak{}x1\_)) = (\allowbreak{}jop(\allowbreak{}b))(\allowbreak{}f(\allowbreak{}x1\_))))
\end{quote}
mentioned by 7 result(s)

\smallskip

\noindent\textbf{\texttt{is-hom-complex-pre-hilbert-space(a, b, f)}}\quad\textit{predicate}\index{is-hom-complex-pre-hilbert-space(a, b, f)@\texttt{is-hom-complex-pre-hilbert-space(a, b, f)}}\\
reads: \$3 is a homomorphism from \$1 to \$2 \\
a predicate \\
characterized by \texttt{is-hom-complex-pre-hilbert-space-def}  \textit{(definitional)}:
\begin{quote}\ttfamily\footnotesize\raggedright\noindent
forall(\allowbreak{}[a, b, f], is-hom-complex-pre-hilbert-space(\allowbreak{}a, b, f) iff is-complex-pre-hilbert-space(\allowbreak{}a) and is-complex-pre-hilbert-space(\allowbreak{}b) and is-hom-complex-semi-inner-product-space(\allowbreak{}a, b, f))
\end{quote}
mentioned by 2 result(s)

\smallskip

\noindent\textbf{\texttt{is-hom-complex-semi-inner-product-space(a, b, f)}}\quad\textit{predicate}\index{is-hom-complex-semi-inner-product-space(a, b, f)@\texttt{is-hom-complex-semi-inner-product-space(a, b, f)}}\\
reads: \$3 is a homomorphism from \$1 to \$2 \\
a predicate \\
characterized by \texttt{is-hom-complex-semi-inner-product-space-def}  \textit{(definitional)}:
\begin{quote}\ttfamily\footnotesize\raggedright\noindent
forall(\allowbreak{}[a, b, f], is-hom-complex-semi-inner-product-space(\allowbreak{}a, b, f) iff is-complex-semi-inner-product-space(\allowbreak{}a) and is-complex-semi-inner-product-space(\allowbreak{}b) and f in fun(\allowbreak{}vec(\allowbreak{}a), vec(\allowbreak{}b)) and scal(\allowbreak{}a) = scal(\allowbreak{}b) and forall(\allowbreak{}[x1\_ in vec(\allowbreak{}a), x2\_ in vec(\allowbreak{}a)], f(\allowbreak{}(\allowbreak{}vadd(\allowbreak{}a))(\allowbreak{}x1\_, x2\_)) = (\allowbreak{}vadd(\allowbreak{}b))(\allowbreak{}f(\allowbreak{}x1\_), f(\allowbreak{}x2\_))) and f(\allowbreak{}vzero(\allowbreak{}a)) = vzero(\allowbreak{}b) and forall(\allowbreak{}[x1\_ in vec(\allowbreak{}a)], f(\allowbreak{}(\allowbreak{}vneg(\allowbreak{}a))(\allowbreak{}x1\_)) = (\allowbreak{}vneg(\allowbreak{}b))(\allowbreak{}f(\allowbreak{}x1\_))) and forall(\allowbreak{}[x1\_ in carr(\allowbreak{}scal(\allowbreak{}a)), x2\_ in vec(\allowbreak{}a)], f(\allowbreak{}(\allowbreak{}act(\allowbreak{}a))(\allowbreak{}x1\_, x2\_)) = (\allowbreak{}act(\allowbreak{}b))(\allowbreak{}x1\_, f(\allowbreak{}x2\_))) and forall(\allowbreak{}[x1\_ in vec(\allowbreak{}a), x2\_ in vec(\allowbreak{}a)], (\allowbreak{}rip(\allowbreak{}a))(\allowbreak{}x1\_, x2\_) = (\allowbreak{}rip(\allowbreak{}b))(\allowbreak{}f(\allowbreak{}x1\_), f(\allowbreak{}x2\_))) and forall(\allowbreak{}[x1\_ in vec(\allowbreak{}a)], f(\allowbreak{}(\allowbreak{}jop(\allowbreak{}a))(\allowbreak{}x1\_)) = (\allowbreak{}jop(\allowbreak{}b))(\allowbreak{}f(\allowbreak{}x1\_))))
\end{quote}
mentioned by 5 result(s)

\smallskip

\noindent\textbf{\texttt{is-hom-directed-set(a, b, f)}}\quad\textit{predicate}\index{is-hom-directed-set(a, b, f)@\texttt{is-hom-directed-set(a, b, f)}}\\
reads: \$3 is a homomorphism from \$1 to \$2 \\
a predicate \\
characterized by \texttt{is-hom-directed-set-def}  \textit{(definitional)}:
\begin{quote}\ttfamily\footnotesize\raggedright\noindent
forall(\allowbreak{}[a, b, f], is-hom-directed-set(\allowbreak{}a, b, f) iff is-directed-set(\allowbreak{}a) and is-directed-set(\allowbreak{}b) and f in fun(\allowbreak{}pts(\allowbreak{}a), pts(\allowbreak{}b)) and forall(\allowbreak{}[x1\_ in pts(\allowbreak{}a), x2\_ in pts(\allowbreak{}a)], [x1\_, x2\_] in dleq(\allowbreak{}a) implies [f(\allowbreak{}x1\_), f(\allowbreak{}x2\_)] in dleq(\allowbreak{}b)))
\end{quote}
mentioned by 2 result(s)

\smallskip

\noindent\textbf{\texttt{is-hom-euclidean-ring(a, b, f)}}\quad\textit{predicate}\index{is-hom-euclidean-ring(a, b, f)@\texttt{is-hom-euclidean-ring(a, b, f)}}\\
reads: \$3 is a homomorphism from \$1 to \$2 \\
a predicate \\
characterized by \texttt{is-hom-euclidean-ring-def}  \textit{(definitional)}:
\begin{quote}\ttfamily\footnotesize\raggedright\noindent
forall(\allowbreak{}[a, b, f], is-hom-euclidean-ring(\allowbreak{}a, b, f) iff is-euclidean-ring(\allowbreak{}a) and is-euclidean-ring(\allowbreak{}b) and is-hom-integral-domain(\allowbreak{}a, b, f))
\end{quote}
mentioned by 7 result(s)

\smallskip

\noindent\textbf{\texttt{is-hom-field(a, b, f)}}\quad\textit{predicate}\index{is-hom-field(a, b, f)@\texttt{is-hom-field(a, b, f)}}\\
reads: \$3 is a homomorphism from \$1 to \$2 \\
a predicate \\
characterized by \texttt{is-hom-field-def}  \textit{(definitional)}:
\begin{quote}\ttfamily\footnotesize\raggedright\noindent
forall(\allowbreak{}[a, b, f], is-hom-field(\allowbreak{}a, b, f) iff is-field(\allowbreak{}a) and is-field(\allowbreak{}b) and f in fun(\allowbreak{}carr(\allowbreak{}a), carr(\allowbreak{}b)) and forall(\allowbreak{}[x1\_ in carr(\allowbreak{}a), x2\_ in carr(\allowbreak{}a)], f(\allowbreak{}(\allowbreak{}add(\allowbreak{}a))(\allowbreak{}x1\_, x2\_)) = (\allowbreak{}add(\allowbreak{}b))(\allowbreak{}f(\allowbreak{}x1\_), f(\allowbreak{}x2\_))) and forall(\allowbreak{}[x1\_ in carr(\allowbreak{}a), x2\_ in carr(\allowbreak{}a)], f(\allowbreak{}(\allowbreak{}mul(\allowbreak{}a))(\allowbreak{}x1\_, x2\_)) = (\allowbreak{}mul(\allowbreak{}b))(\allowbreak{}f(\allowbreak{}x1\_), f(\allowbreak{}x2\_))) and forall(\allowbreak{}[x1\_ in carr(\allowbreak{}a)], f(\allowbreak{}(\allowbreak{}neg(\allowbreak{}a))(\allowbreak{}x1\_)) = (\allowbreak{}neg(\allowbreak{}b))(\allowbreak{}f(\allowbreak{}x1\_))) and f(\allowbreak{}zero(\allowbreak{}a)) = zero(\allowbreak{}b) and f(\allowbreak{}one(\allowbreak{}a)) = one(\allowbreak{}b) and forall(\allowbreak{}[x1\_ in non-zero(\allowbreak{}a)], f(\allowbreak{}(\allowbreak{}mul-inv(\allowbreak{}a))(\allowbreak{}x1\_)) = (\allowbreak{}mul-inv(\allowbreak{}b))(\allowbreak{}f(\allowbreak{}x1\_))))
\end{quote}
mentioned by 10 result(s)

\smallskip

\noindent\textbf{\texttt{is-hom-field-ring(a, b, f)}}\quad\textit{predicate}\index{is-hom-field-ring(a, b, f)@\texttt{is-hom-field-ring(a, b, f)}}\\
reads: \$3 is a homomorphism from \$1 to \$2 \\
a predicate \\
characterized by \texttt{is-hom-field-ring-def}  \textit{(definitional)}:
\begin{quote}\ttfamily\footnotesize\raggedright\noindent
forall(\allowbreak{}[a, b, f], is-hom-field-ring(\allowbreak{}a, b, f) iff is-field-ring(\allowbreak{}a) and is-field-ring(\allowbreak{}b) and is-hom-commutative-ring(\allowbreak{}a, b, f))
\end{quote}
mentioned by 6 result(s)

\smallskip

\noindent\textbf{\texttt{is-hom-group(a, b, f)}}\quad\textit{predicate}\index{is-hom-group(a, b, f)@\texttt{is-hom-group(a, b, f)}}\\
reads: \$3 is a homomorphism from \$1 to \$2 \\
a predicate \\
characterized by \texttt{is-hom-group-def}  \textit{(definitional)}:
\begin{quote}\ttfamily\footnotesize\raggedright\noindent
forall(\allowbreak{}[a, b, f], is-hom-group(\allowbreak{}a, b, f) iff is-group(\allowbreak{}a) and is-group(\allowbreak{}b) and f in fun(\allowbreak{}carr(\allowbreak{}a), carr(\allowbreak{}b)) and forall(\allowbreak{}[x1\_ in carr(\allowbreak{}a), x2\_ in carr(\allowbreak{}a)], f(\allowbreak{}(\allowbreak{}opr(\allowbreak{}a))(\allowbreak{}x1\_, x2\_)) = (\allowbreak{}opr(\allowbreak{}b))(\allowbreak{}f(\allowbreak{}x1\_), f(\allowbreak{}x2\_))) and f(\allowbreak{}iden(\allowbreak{}a)) = iden(\allowbreak{}b) and forall(\allowbreak{}[x1\_ in carr(\allowbreak{}a)], f(\allowbreak{}(\allowbreak{}inv(\allowbreak{}a))(\allowbreak{}x1\_)) = (\allowbreak{}inv(\allowbreak{}b))(\allowbreak{}f(\allowbreak{}x1\_))))
\end{quote}
mentioned by 13 result(s)

\smallskip

\noindent\textbf{\texttt{is-hom-infiltrated-space(a, b, f1, f2)}}\quad\textit{predicate}\index{is-hom-infiltrated-space(a, b, f1, f2)@\texttt{is-hom-infiltrated-space(a, b, f1, f2)}}\\
reads: \$3 is a morphism of infiltrated strand spaces from \$1 to \$2 \\
a predicate \\
characterized by \texttt{is-hom-infiltrated-space-def}  \textit{(definitional)}:
\begin{quote}\ttfamily\footnotesize\raggedright\noindent
forall(\allowbreak{}[a, b, f1, f2], is-hom-infiltrated-space(\allowbreak{}a, b, f1, f2) iff is-infiltrated-space(\allowbreak{}a) and is-infiltrated-space(\allowbreak{}b) and f1 in fun(\allowbreak{}adv-ids(\allowbreak{}a), adv-ids(\allowbreak{}b)) and f2 in fun(\allowbreak{}adv-tokens(\allowbreak{}a), adv-tokens(\allowbreak{}b)) and ispace(\allowbreak{}a) = ispace(\allowbreak{}b))
\end{quote}
mentioned by 2 result(s)

\smallskip

\noindent\textbf{\texttt{is-hom-integral-domain(a, b, f)}}\quad\textit{predicate}\index{is-hom-integral-domain(a, b, f)@\texttt{is-hom-integral-domain(a, b, f)}}\\
reads: \$3 is a homomorphism from \$1 to \$2 \\
a predicate \\
characterized by \texttt{is-hom-integral-domain-def}  \textit{(definitional)}:
\begin{quote}\ttfamily\footnotesize\raggedright\noindent
forall(\allowbreak{}[a, b, f], is-hom-integral-domain(\allowbreak{}a, b, f) iff is-integral-domain(\allowbreak{}a) and is-integral-domain(\allowbreak{}b) and is-hom-commutative-ring(\allowbreak{}a, b, f))
\end{quote}
mentioned by 12 result(s)

\smallskip

\noindent\textbf{\texttt{is-hom-measurable-space(s, t, f)}}\quad\textit{predicate}\index{is-hom-measurable-space(s, t, f)@\texttt{is-hom-measurable-space(s, t, f)}}\\
reads: \$3 is a measurable map from \$1 to \$2 \\
a predicate \\
characterized by \texttt{is-hom-measurable-space-def}  \textit{(definitional)}:
\begin{quote}\ttfamily\footnotesize\raggedright\noindent
forall(\allowbreak{}[s, t, f], is-hom-measurable-space(\allowbreak{}s, t, f) iff is-measurable-space(\allowbreak{}s) and is-measurable-space(\allowbreak{}t) and f in fun(\allowbreak{}pts(\allowbreak{}s), pts(\allowbreak{}t)) and forall(\allowbreak{}[u in sigma(\allowbreak{}t)], preimage(\allowbreak{}s, f, u) in sigma(\allowbreak{}s)))
\end{quote}
mentioned by 7 result(s)

\smallskip

\noindent\textbf{\texttt{is-hom-measure-space(s, t, f)}}\quad\textit{predicate}\index{is-hom-measure-space(s, t, f)@\texttt{is-hom-measure-space(s, t, f)}}\\
reads: \$3 is a measure-preserving measurable map from \$1 to \$2 \\
a predicate \\
characterized by \texttt{is-hom-measure-space-def}  \textit{(definitional)}:
\begin{quote}\ttfamily\footnotesize\raggedright\noindent
forall(\allowbreak{}[s, t, f], is-hom-measure-space(\allowbreak{}s, t, f) iff is-measure-space(\allowbreak{}s) and is-measure-space(\allowbreak{}t) and f in fun(\allowbreak{}pts(\allowbreak{}s), pts(\allowbreak{}t)) and forall(\allowbreak{}[u in sigma(\allowbreak{}t)], preimage(\allowbreak{}s, f, u) in sigma(\allowbreak{}s)) and forall(\allowbreak{}[u in sigma(\allowbreak{}t)], (\allowbreak{}meas(\allowbreak{}t))(\allowbreak{}u) = (\allowbreak{}meas(\allowbreak{}s))(\allowbreak{}preimage(\allowbreak{}s, f, u))))
\end{quote}
mentioned by 7 result(s)

\smallskip

\noindent\textbf{\texttt{is-hom-metric-group(a, b, f)}}\quad\textit{predicate}\index{is-hom-metric-group(a, b, f)@\texttt{is-hom-metric-group(a, b, f)}}\\
reads: \$3 is an isometric homomorphism from \$1 to \$2 \\
a predicate \\
characterized by \texttt{is-hom-metric-group-def}  \textit{(definitional)}:
\begin{quote}\ttfamily\footnotesize\raggedright\noindent
forall(\allowbreak{}[a, b, f], is-hom-metric-group(\allowbreak{}a, b, f) iff is-metric-group(\allowbreak{}a) and is-metric-group(\allowbreak{}b) and f in fun(\allowbreak{}carr(\allowbreak{}a), carr(\allowbreak{}b)) and forall(\allowbreak{}[x1\_ in carr(\allowbreak{}a), x2\_ in carr(\allowbreak{}a)], f(\allowbreak{}(\allowbreak{}opr(\allowbreak{}a))(\allowbreak{}x1\_, x2\_)) = (\allowbreak{}opr(\allowbreak{}b))(\allowbreak{}f(\allowbreak{}x1\_), f(\allowbreak{}x2\_))) and f(\allowbreak{}iden(\allowbreak{}a)) = iden(\allowbreak{}b) and forall(\allowbreak{}[x1\_ in carr(\allowbreak{}a)], f(\allowbreak{}(\allowbreak{}inv(\allowbreak{}a))(\allowbreak{}x1\_)) = (\allowbreak{}inv(\allowbreak{}b))(\allowbreak{}f(\allowbreak{}x1\_))) and forall(\allowbreak{}[x1\_ in carr(\allowbreak{}a), x2\_ in carr(\allowbreak{}a)], (\allowbreak{}gdist(\allowbreak{}a))(\allowbreak{}x1\_, x2\_) = (\allowbreak{}gdist(\allowbreak{}b))(\allowbreak{}f(\allowbreak{}x1\_), f(\allowbreak{}x2\_))))
\end{quote}
mentioned by 4 result(s)

\smallskip

\noindent\textbf{\texttt{is-hom-metric-space(a, b, f)}}\quad\textit{predicate}\index{is-hom-metric-space(a, b, f)@\texttt{is-hom-metric-space(a, b, f)}}\\
reads: \$3 is an isometry from \$1 to \$2 \\
a predicate \\
characterized by \texttt{is-hom-metric-space-def}  \textit{(definitional)}:
\begin{quote}\ttfamily\footnotesize\raggedright\noindent
forall(\allowbreak{}[a, b, f], is-hom-metric-space(\allowbreak{}a, b, f) iff is-metric-space(\allowbreak{}a) and is-metric-space(\allowbreak{}b) and f in fun(\allowbreak{}pts(\allowbreak{}a), pts(\allowbreak{}b)) and forall(\allowbreak{}[x1\_ in pts(\allowbreak{}a), x2\_ in pts(\allowbreak{}a)], (\allowbreak{}dist(\allowbreak{}a))(\allowbreak{}x1\_, x2\_) = (\allowbreak{}dist(\allowbreak{}b))(\allowbreak{}f(\allowbreak{}x1\_), f(\allowbreak{}x2\_))))
\end{quote}
mentioned by 10 result(s)

\smallskip

\noindent\textbf{\texttt{is-hom-metrizable-top-space(s, t, f)}}\quad\textit{predicate}\index{is-hom-metrizable-top-space(s, t, f)@\texttt{is-hom-metrizable-top-space(s, t, f)}}\\
reads: \$3 is continuous from \$1 to \$2 \\
a predicate \\
characterized by \texttt{is-hom-metrizable-top-space-def}  \textit{(definitional)}:
\begin{quote}\ttfamily\footnotesize\raggedright\noindent
forall(\allowbreak{}[s, t, f], is-hom-metrizable-top-space(\allowbreak{}s, t, f) iff is-metrizable-top-space(\allowbreak{}s) and is-metrizable-top-space(\allowbreak{}t) and f in fun(\allowbreak{}pts(\allowbreak{}s), pts(\allowbreak{}t)) and forall(\allowbreak{}[u in opens(\allowbreak{}t)], preimage(\allowbreak{}s, f, u) in opens(\allowbreak{}s)))
\end{quote}
mentioned by 7 result(s)

\smallskip

\noindent\textbf{\texttt{is-hom-module(a, b, f)}}\quad\textit{predicate}\index{is-hom-module(a, b, f)@\texttt{is-hom-module(a, b, f)}}\\
reads: \$3 is a homomorphism from \$1 to \$2 \\
a predicate \\
characterized by \texttt{is-hom-module-def}  \textit{(definitional)}:
\begin{quote}\ttfamily\footnotesize\raggedright\noindent
forall(\allowbreak{}[a, b, f], is-hom-module(\allowbreak{}a, b, f) iff is-module(\allowbreak{}a) and is-module(\allowbreak{}b) and f in fun(\allowbreak{}vec(\allowbreak{}a), vec(\allowbreak{}b)) and scal(\allowbreak{}a) = scal(\allowbreak{}b) and forall(\allowbreak{}[x1\_ in vec(\allowbreak{}a), x2\_ in vec(\allowbreak{}a)], f(\allowbreak{}(\allowbreak{}vadd(\allowbreak{}a))(\allowbreak{}x1\_, x2\_)) = (\allowbreak{}vadd(\allowbreak{}b))(\allowbreak{}f(\allowbreak{}x1\_), f(\allowbreak{}x2\_))) and f(\allowbreak{}vzero(\allowbreak{}a)) = vzero(\allowbreak{}b) and forall(\allowbreak{}[x1\_ in vec(\allowbreak{}a)], f(\allowbreak{}(\allowbreak{}vneg(\allowbreak{}a))(\allowbreak{}x1\_)) = (\allowbreak{}vneg(\allowbreak{}b))(\allowbreak{}f(\allowbreak{}x1\_))) and forall(\allowbreak{}[x1\_ in carr(\allowbreak{}scal(\allowbreak{}a)), x2\_ in vec(\allowbreak{}a)], f(\allowbreak{}(\allowbreak{}act(\allowbreak{}a))(\allowbreak{}x1\_, x2\_)) = (\allowbreak{}act(\allowbreak{}b))(\allowbreak{}x1\_, f(\allowbreak{}x2\_))))
\end{quote}
mentioned by 178 result(s)

\smallskip

\noindent\textbf{\texttt{is-hom-monoid(a, b, f)}}\quad\textit{predicate}\index{is-hom-monoid(a, b, f)@\texttt{is-hom-monoid(a, b, f)}}\\
reads: \$3 is a homomorphism from \$1 to \$2 \\
a predicate \\
characterized by \texttt{is-hom-monoid-def}  \textit{(definitional)}:
\begin{quote}\ttfamily\footnotesize\raggedright\noindent
forall(\allowbreak{}[a, b, f], is-hom-monoid(\allowbreak{}a, b, f) iff is-monoid(\allowbreak{}a) and is-monoid(\allowbreak{}b) and f in fun(\allowbreak{}carr(\allowbreak{}a), carr(\allowbreak{}b)) and forall(\allowbreak{}[x1\_ in carr(\allowbreak{}a), x2\_ in carr(\allowbreak{}a)], f(\allowbreak{}(\allowbreak{}opr(\allowbreak{}a))(\allowbreak{}x1\_, x2\_)) = (\allowbreak{}opr(\allowbreak{}b))(\allowbreak{}f(\allowbreak{}x1\_), f(\allowbreak{}x2\_))) and f(\allowbreak{}iden(\allowbreak{}a)) = iden(\allowbreak{}b))
\end{quote}
mentioned by 10 result(s)

\smallskip

\noindent\textbf{\texttt{is-hom-normed-ag(a, b, f)}}\quad\textit{predicate}\index{is-hom-normed-ag(a, b, f)@\texttt{is-hom-normed-ag(a, b, f)}}\\
reads: \$3 is a homomorphism from \$1 to \$2 \\
a predicate \\
characterized by \texttt{is-hom-normed-ag-def}  \textit{(definitional)}:
\begin{quote}\ttfamily\footnotesize\raggedright\noindent
forall(\allowbreak{}[a, b, f], is-hom-normed-ag(\allowbreak{}a, b, f) iff is-normed-ag(\allowbreak{}a) and is-normed-ag(\allowbreak{}b) and f in fun(\allowbreak{}carr(\allowbreak{}a), carr(\allowbreak{}b)) and forall(\allowbreak{}[x1\_ in carr(\allowbreak{}a), x2\_ in carr(\allowbreak{}a)], f(\allowbreak{}(\allowbreak{}opr(\allowbreak{}a))(\allowbreak{}x1\_, x2\_)) = (\allowbreak{}opr(\allowbreak{}b))(\allowbreak{}f(\allowbreak{}x1\_), f(\allowbreak{}x2\_))) and f(\allowbreak{}iden(\allowbreak{}a)) = iden(\allowbreak{}b) and forall(\allowbreak{}[x1\_ in carr(\allowbreak{}a)], f(\allowbreak{}(\allowbreak{}inv(\allowbreak{}a))(\allowbreak{}x1\_)) = (\allowbreak{}inv(\allowbreak{}b))(\allowbreak{}f(\allowbreak{}x1\_))) and forall(\allowbreak{}[x1\_ in carr(\allowbreak{}a)], (\allowbreak{}nrm(\allowbreak{}a))(\allowbreak{}x1\_) = (\allowbreak{}nrm(\allowbreak{}b))(\allowbreak{}f(\allowbreak{}x1\_))))
\end{quote}
mentioned by 15 result(s)

\smallskip

\noindent\textbf{\texttt{is-hom-normed-field(a, b, f)}}\quad\textit{predicate}\index{is-hom-normed-field(a, b, f)@\texttt{is-hom-normed-field(a, b, f)}}\\
reads: \$3 is a homomorphism from \$1 to \$2 \\
a predicate \\
characterized by \texttt{is-hom-normed-field-def}  \textit{(definitional)}:
\begin{quote}\ttfamily\footnotesize\raggedright\noindent
forall(\allowbreak{}[a, b, f], is-hom-normed-field(\allowbreak{}a, b, f) iff is-normed-field(\allowbreak{}a) and is-normed-field(\allowbreak{}b) and f in fun(\allowbreak{}carr(\allowbreak{}a), carr(\allowbreak{}b)) and forall(\allowbreak{}[x1\_ in carr(\allowbreak{}a), x2\_ in carr(\allowbreak{}a)], f(\allowbreak{}(\allowbreak{}add(\allowbreak{}a))(\allowbreak{}x1\_, x2\_)) = (\allowbreak{}add(\allowbreak{}b))(\allowbreak{}f(\allowbreak{}x1\_), f(\allowbreak{}x2\_))) and forall(\allowbreak{}[x1\_ in carr(\allowbreak{}a), x2\_ in carr(\allowbreak{}a)], f(\allowbreak{}(\allowbreak{}mul(\allowbreak{}a))(\allowbreak{}x1\_, x2\_)) = (\allowbreak{}mul(\allowbreak{}b))(\allowbreak{}f(\allowbreak{}x1\_), f(\allowbreak{}x2\_))) and forall(\allowbreak{}[x1\_ in carr(\allowbreak{}a)], f(\allowbreak{}(\allowbreak{}neg(\allowbreak{}a))(\allowbreak{}x1\_)) = (\allowbreak{}neg(\allowbreak{}b))(\allowbreak{}f(\allowbreak{}x1\_))) and f(\allowbreak{}zero(\allowbreak{}a)) = zero(\allowbreak{}b) and f(\allowbreak{}one(\allowbreak{}a)) = one(\allowbreak{}b) and forall(\allowbreak{}[x1\_ in carr(\allowbreak{}a)], (\allowbreak{}fnrm(\allowbreak{}a))(\allowbreak{}x1\_) = (\allowbreak{}fnrm(\allowbreak{}b))(\allowbreak{}f(\allowbreak{}x1\_))))
\end{quote}
mentioned by 9 result(s)

\smallskip

\noindent\textbf{\texttt{is-hom-normed-vector-space(a, b, f)}}\quad\textit{predicate}\index{is-hom-normed-vector-space(a, b, f)@\texttt{is-hom-normed-vector-space(a, b, f)}}\\
reads: \$3 is a homomorphism from \$1 to \$2 \\
a predicate \\
characterized by \texttt{is-hom-normed-vector-space-def}  \textit{(definitional)}:
\begin{quote}\ttfamily\footnotesize\raggedright\noindent
forall(\allowbreak{}[a, b, f], is-hom-normed-vector-space(\allowbreak{}a, b, f) iff is-normed-vector-space(\allowbreak{}a) and is-normed-vector-space(\allowbreak{}b) and f in fun(\allowbreak{}vec(\allowbreak{}a), vec(\allowbreak{}b)) and scal(\allowbreak{}a) = scal(\allowbreak{}b) and forall(\allowbreak{}[x1\_ in vec(\allowbreak{}a), x2\_ in vec(\allowbreak{}a)], f(\allowbreak{}(\allowbreak{}vadd(\allowbreak{}a))(\allowbreak{}x1\_, x2\_)) = (\allowbreak{}vadd(\allowbreak{}b))(\allowbreak{}f(\allowbreak{}x1\_), f(\allowbreak{}x2\_))) and f(\allowbreak{}vzero(\allowbreak{}a)) = vzero(\allowbreak{}b) and forall(\allowbreak{}[x1\_ in vec(\allowbreak{}a)], f(\allowbreak{}(\allowbreak{}vneg(\allowbreak{}a))(\allowbreak{}x1\_)) = (\allowbreak{}vneg(\allowbreak{}b))(\allowbreak{}f(\allowbreak{}x1\_))) and forall(\allowbreak{}[x1\_ in carr(\allowbreak{}scal(\allowbreak{}a)), x2\_ in vec(\allowbreak{}a)], f(\allowbreak{}(\allowbreak{}act(\allowbreak{}a))(\allowbreak{}x1\_, x2\_)) = (\allowbreak{}act(\allowbreak{}b))(\allowbreak{}x1\_, f(\allowbreak{}x2\_))) and forall(\allowbreak{}[x1\_ in vec(\allowbreak{}a)], (\allowbreak{}vnrm(\allowbreak{}a))(\allowbreak{}x1\_) = (\allowbreak{}vnrm(\allowbreak{}b))(\allowbreak{}f(\allowbreak{}x1\_))))
\end{quote}
mentioned by 11 result(s)

\smallskip

\noindent\textbf{\texttt{is-hom-nu-c-star-algebra(a, b, f)}}\quad\textit{predicate}\index{is-hom-nu-c-star-algebra(a, b, f)@\texttt{is-hom-nu-c-star-algebra(a, b, f)}}\\
reads: \$3 is a homomorphism from \$1 to \$2 \\
a predicate \\
characterized by \texttt{is-hom-nu-c-star-algebra-def}  \textit{(definitional)}:
\begin{quote}\ttfamily\footnotesize\raggedright\noindent
forall(\allowbreak{}[a, b, f], is-hom-nu-c-star-algebra(\allowbreak{}a, b, f) iff is-nu-c-star-algebra(\allowbreak{}a) and is-nu-c-star-algebra(\allowbreak{}b) and is-hom-nu-involutive-banach-algebra(\allowbreak{}a, b, f))
\end{quote}
mentioned by 2 result(s)

\smallskip

\noindent\textbf{\texttt{is-hom-nu-involutive-banach-algebra(a, b, f)}}\quad\textit{predicate}\index{is-hom-nu-involutive-banach-algebra(a, b, f)@\texttt{is-hom-nu-involutive-banach-algebra(a, b, f)}}\\
reads: \$3 is a homomorphism from \$1 to \$2 \\
a predicate \\
characterized by \texttt{is-hom-nu-involutive-banach-algebra-def}  \textit{(definitional)}:
\begin{quote}\ttfamily\footnotesize\raggedright\noindent
forall(\allowbreak{}[a, b, f], is-hom-nu-involutive-banach-algebra(\allowbreak{}a, b, f) iff is-nu-involutive-banach-algebra(\allowbreak{}a) and is-nu-involutive-banach-algebra(\allowbreak{}b) and f in fun(\allowbreak{}vec(\allowbreak{}a), vec(\allowbreak{}b)) and scal(\allowbreak{}a) = scal(\allowbreak{}b) and forall(\allowbreak{}[x1\_ in vec(\allowbreak{}a), x2\_ in vec(\allowbreak{}a)], f(\allowbreak{}(\allowbreak{}vadd(\allowbreak{}a))(\allowbreak{}x1\_, x2\_)) = (\allowbreak{}vadd(\allowbreak{}b))(\allowbreak{}f(\allowbreak{}x1\_), f(\allowbreak{}x2\_))) and f(\allowbreak{}vzero(\allowbreak{}a)) = vzero(\allowbreak{}b) and forall(\allowbreak{}[x1\_ in vec(\allowbreak{}a)], f(\allowbreak{}(\allowbreak{}vneg(\allowbreak{}a))(\allowbreak{}x1\_)) = (\allowbreak{}vneg(\allowbreak{}b))(\allowbreak{}f(\allowbreak{}x1\_))) and forall(\allowbreak{}[x1\_ in carr(\allowbreak{}scal(\allowbreak{}a)), x2\_ in vec(\allowbreak{}a)], f(\allowbreak{}(\allowbreak{}act(\allowbreak{}a))(\allowbreak{}x1\_, x2\_)) = (\allowbreak{}act(\allowbreak{}b))(\allowbreak{}x1\_, f(\allowbreak{}x2\_))) and forall(\allowbreak{}[x1\_ in vec(\allowbreak{}a)], (\allowbreak{}vnrm(\allowbreak{}a))(\allowbreak{}x1\_) = (\allowbreak{}vnrm(\allowbreak{}b))(\allowbreak{}f(\allowbreak{}x1\_))) and forall(\allowbreak{}[x1\_ in vec(\allowbreak{}a), x2\_ in vec(\allowbreak{}a)], f(\allowbreak{}(\allowbreak{}vmul(\allowbreak{}a))(\allowbreak{}x1\_, x2\_)) = (\allowbreak{}vmul(\allowbreak{}b))(\allowbreak{}f(\allowbreak{}x1\_), f(\allowbreak{}x2\_))) and forall(\allowbreak{}[x1\_ in vec(\allowbreak{}a)], f(\allowbreak{}(\allowbreak{}istar(\allowbreak{}a))(\allowbreak{}x1\_)) = (\allowbreak{}istar(\allowbreak{}b))(\allowbreak{}f(\allowbreak{}x1\_))))
\end{quote}
mentioned by 6 result(s)

\smallskip

\noindent\textbf{\texttt{is-hom-pid(a, b, f)}}\quad\textit{predicate}\index{is-hom-pid(a, b, f)@\texttt{is-hom-pid(a, b, f)}}\\
reads: \$3 is a homomorphism from \$1 to \$2 \\
a predicate \\
characterized by \texttt{is-hom-pid-def}  \textit{(definitional)}:
\begin{quote}\ttfamily\footnotesize\raggedright\noindent
forall(\allowbreak{}[a, b, f], is-hom-pid(\allowbreak{}a, b, f) iff is-pid(\allowbreak{}a) and is-pid(\allowbreak{}b) and is-hom-integral-domain(\allowbreak{}a, b, f))
\end{quote}
mentioned by 6 result(s)

\smallskip

\noindent\textbf{\texttt{is-hom-pk-space(a, b, f)}}\quad\textit{predicate}\index{is-hom-pk-space(a, b, f)@\texttt{is-hom-pk-space(a, b, f)}}\\
reads: \$3 is a morphism of public-key strand spaces from \$1 to \$2 \\
a predicate \\
characterized by \texttt{is-hom-pk-space-def}  \textit{(definitional)}:
\begin{quote}\ttfamily\footnotesize\raggedright\noindent
forall(\allowbreak{}[a, b, f], is-hom-pk-space(\allowbreak{}a, b, f) iff is-pk-space(\allowbreak{}a) and is-pk-space(\allowbreak{}b) and f in fun(\allowbreak{}names(\allowbreak{}a), names(\allowbreak{}b)) and pk-inf(\allowbreak{}a) = pk-inf(\allowbreak{}b) and forall(\allowbreak{}[x1\_ in names(\allowbreak{}a)], (\allowbreak{}pubkey(\allowbreak{}a))(\allowbreak{}x1\_) = (\allowbreak{}pubkey(\allowbreak{}b))(\allowbreak{}f(\allowbreak{}x1\_))))
\end{quote}
mentioned by 2 result(s)

\smallskip

\noindent\textbf{\texttt{is-hom-pre-hilbert-space(a, b, f)}}\quad\textit{predicate}\index{is-hom-pre-hilbert-space(a, b, f)@\texttt{is-hom-pre-hilbert-space(a, b, f)}}\\
reads: \$3 is a homomorphism from \$1 to \$2 \\
a predicate \\
characterized by \texttt{is-hom-pre-hilbert-space-def}  \textit{(definitional)}:
\begin{quote}\ttfamily\footnotesize\raggedright\noindent
forall(\allowbreak{}[a, b, f], is-hom-pre-hilbert-space(\allowbreak{}a, b, f) iff is-pre-hilbert-space(\allowbreak{}a) and is-pre-hilbert-space(\allowbreak{}b) and is-hom-semi-inner-product-space(\allowbreak{}a, b, f))
\end{quote}
mentioned by 2 result(s)

\smallskip

\noindent\textbf{\texttt{is-hom-pseudometric-space(a, b, f)}}\quad\textit{predicate}\index{is-hom-pseudometric-space(a, b, f)@\texttt{is-hom-pseudometric-space(a, b, f)}}\\
reads: \$3 is a homomorphism from \$1 to \$2 \\
a predicate \\
characterized by \texttt{is-hom-pseudometric-space-def}  \textit{(definitional)}:
\begin{quote}\ttfamily\footnotesize\raggedright\noindent
forall(\allowbreak{}[a, b, f], is-hom-pseudometric-space(\allowbreak{}a, b, f) iff is-pseudometric-space(\allowbreak{}a) and is-pseudometric-space(\allowbreak{}b) and f in fun(\allowbreak{}pts(\allowbreak{}a), pts(\allowbreak{}b)) and forall(\allowbreak{}[x1\_ in pts(\allowbreak{}a), x2\_ in pts(\allowbreak{}a)], (\allowbreak{}dist(\allowbreak{}a))(\allowbreak{}x1\_, x2\_) = (\allowbreak{}dist(\allowbreak{}b))(\allowbreak{}f(\allowbreak{}x1\_), f(\allowbreak{}x2\_))))
\end{quote}
mentioned by 6 result(s)

\smallskip

\noindent\textbf{\texttt{is-hom-ring(a, b, f)}}\quad\textit{predicate}\index{is-hom-ring(a, b, f)@\texttt{is-hom-ring(a, b, f)}}\\
reads: \$3 is a homomorphism from \$1 to \$2 \\
a predicate \\
characterized by \texttt{is-hom-ring-def}  \textit{(definitional)}:
\begin{quote}\ttfamily\footnotesize\raggedright\noindent
forall(\allowbreak{}[a, b, f], is-hom-ring(\allowbreak{}a, b, f) iff is-ring(\allowbreak{}a) and is-ring(\allowbreak{}b) and f in fun(\allowbreak{}carr(\allowbreak{}a), carr(\allowbreak{}b)) and forall(\allowbreak{}[x1\_ in carr(\allowbreak{}a), x2\_ in carr(\allowbreak{}a)], f(\allowbreak{}(\allowbreak{}add(\allowbreak{}a))(\allowbreak{}x1\_, x2\_)) = (\allowbreak{}add(\allowbreak{}b))(\allowbreak{}f(\allowbreak{}x1\_), f(\allowbreak{}x2\_))) and forall(\allowbreak{}[x1\_ in carr(\allowbreak{}a), x2\_ in carr(\allowbreak{}a)], f(\allowbreak{}(\allowbreak{}mul(\allowbreak{}a))(\allowbreak{}x1\_, x2\_)) = (\allowbreak{}mul(\allowbreak{}b))(\allowbreak{}f(\allowbreak{}x1\_), f(\allowbreak{}x2\_))) and forall(\allowbreak{}[x1\_ in carr(\allowbreak{}a)], f(\allowbreak{}(\allowbreak{}neg(\allowbreak{}a))(\allowbreak{}x1\_)) = (\allowbreak{}neg(\allowbreak{}b))(\allowbreak{}f(\allowbreak{}x1\_))) and f(\allowbreak{}zero(\allowbreak{}a)) = zero(\allowbreak{}b) and f(\allowbreak{}one(\allowbreak{}a)) = one(\allowbreak{}b))
\end{quote}
also: \texttt{is-hom-ringoid-def} \\
mentioned by 22 result(s)

\smallskip

\noindent\textbf{\texttt{is-hom-ringoid(a, b, f)}}\quad\textit{predicate}\index{is-hom-ringoid(a, b, f)@\texttt{is-hom-ringoid(a, b, f)}}\\
reads: \$3 is a ringoid homomorphism from \$1 to \$2 \\
a predicate \\
characterized by \texttt{is-hom-ringoid-def}  \textit{(definitional)}:
\begin{quote}\ttfamily\footnotesize\raggedright\noindent
forall(\allowbreak{}[a, b, f], is-hom-ringoid(\allowbreak{}a, b, f) iff is-ringoid(\allowbreak{}a) and is-ringoid(\allowbreak{}b) and f in fun(\allowbreak{}carr(\allowbreak{}a), carr(\allowbreak{}b)) and is-hom-ring(\allowbreak{}ringoid-as-ring(\allowbreak{}a), ringoid-as-ring(\allowbreak{}b), f) and forall(\allowbreak{}[x\_ in idl(\allowbreak{}a)], f(\allowbreak{}x\_) in idl(\allowbreak{}b)))
\end{quote}
mentioned by 7 result(s)

\smallskip

\noindent\textbf{\texttt{is-hom-semi-inner-product-space(a, b, f)}}\quad\textit{predicate}\index{is-hom-semi-inner-product-space(a, b, f)@\texttt{is-hom-semi-inner-product-space(a, b, f)}}\\
reads: \$3 is a homomorphism from \$1 to \$2 \\
a predicate \\
characterized by \texttt{is-hom-semi-inner-product-space-def}  \textit{(definitional)}:
\begin{quote}\ttfamily\footnotesize\raggedright\noindent
forall(\allowbreak{}[a, b, f], is-hom-semi-inner-product-space(\allowbreak{}a, b, f) iff is-semi-inner-product-space(\allowbreak{}a) and is-semi-inner-product-space(\allowbreak{}b) and f in fun(\allowbreak{}vec(\allowbreak{}a), vec(\allowbreak{}b)) and scal(\allowbreak{}a) = scal(\allowbreak{}b) and forall(\allowbreak{}[x1\_ in vec(\allowbreak{}a), x2\_ in vec(\allowbreak{}a)], f(\allowbreak{}(\allowbreak{}vadd(\allowbreak{}a))(\allowbreak{}x1\_, x2\_)) = (\allowbreak{}vadd(\allowbreak{}b))(\allowbreak{}f(\allowbreak{}x1\_), f(\allowbreak{}x2\_))) and f(\allowbreak{}vzero(\allowbreak{}a)) = vzero(\allowbreak{}b) and forall(\allowbreak{}[x1\_ in vec(\allowbreak{}a)], f(\allowbreak{}(\allowbreak{}vneg(\allowbreak{}a))(\allowbreak{}x1\_)) = (\allowbreak{}vneg(\allowbreak{}b))(\allowbreak{}f(\allowbreak{}x1\_))) and forall(\allowbreak{}[x1\_ in carr(\allowbreak{}scal(\allowbreak{}a)), x2\_ in vec(\allowbreak{}a)], f(\allowbreak{}(\allowbreak{}act(\allowbreak{}a))(\allowbreak{}x1\_, x2\_)) = (\allowbreak{}act(\allowbreak{}b))(\allowbreak{}x1\_, f(\allowbreak{}x2\_))) and forall(\allowbreak{}[x1\_ in vec(\allowbreak{}a), x2\_ in vec(\allowbreak{}a)], (\allowbreak{}rip(\allowbreak{}a))(\allowbreak{}x1\_, x2\_) = (\allowbreak{}rip(\allowbreak{}b))(\allowbreak{}f(\allowbreak{}x1\_), f(\allowbreak{}x2\_))))
\end{quote}
mentioned by 6 result(s)

\smallskip

\noindent\textbf{\texttt{is-hom-semigroup(a, b, f)}}\quad\textit{predicate}\index{is-hom-semigroup(a, b, f)@\texttt{is-hom-semigroup(a, b, f)}}\\
reads: \$3 is a homomorphism from \$1 to \$2 \\
a predicate \\
characterized by \texttt{is-hom-semigroup-def}  \textit{(definitional)}:
\begin{quote}\ttfamily\footnotesize\raggedright\noindent
forall(\allowbreak{}[a, b, f], is-hom-semigroup(\allowbreak{}a, b, f) iff is-semigroup(\allowbreak{}a) and is-semigroup(\allowbreak{}b) and f in fun(\allowbreak{}carr(\allowbreak{}a), carr(\allowbreak{}b)) and forall(\allowbreak{}[x1\_ in carr(\allowbreak{}a), x2\_ in carr(\allowbreak{}a)], f(\allowbreak{}(\allowbreak{}opr(\allowbreak{}a))(\allowbreak{}x1\_, x2\_)) = (\allowbreak{}opr(\allowbreak{}b))(\allowbreak{}f(\allowbreak{}x1\_), f(\allowbreak{}x2\_))))
\end{quote}
mentioned by 6 result(s)

\smallskip

\noindent\textbf{\texttt{is-hom-setoid(a, b, f)}}\quad\textit{predicate}\index{is-hom-setoid(a, b, f)@\texttt{is-hom-setoid(a, b, f)}}\\
reads: \$3 is a setoid morphism from \$1 to \$2 \\
a predicate \\
characterized by \texttt{is-hom-setoid-def}  \textit{(definitional)}:
\begin{quote}\ttfamily\footnotesize\raggedright\noindent
forall(\allowbreak{}[a, b, f], is-hom-setoid(\allowbreak{}a, b, f) iff is-setoid(\allowbreak{}a) and is-setoid(\allowbreak{}b) and f in fun(\allowbreak{}pts(\allowbreak{}a), pts(\allowbreak{}b)) and forall(\allowbreak{}[x1\_ in pts(\allowbreak{}a), x2\_ in pts(\allowbreak{}a)], [x1\_, x2\_] in rel(\allowbreak{}a) implies [f(\allowbreak{}x1\_), f(\allowbreak{}x2\_)] in rel(\allowbreak{}b)))
\end{quote}
mentioned by 6 result(s)

\smallskip

\noindent\textbf{\texttt{is-hom-sk-space(a, b, f)}}\quad\textit{predicate}\index{is-hom-sk-space(a, b, f)@\texttt{is-hom-sk-space(a, b, f)}}\\
reads: \$3 is a morphism of shared-key strand spaces from \$1 to \$2 \\
a predicate \\
characterized by \texttt{is-hom-sk-space-def}  \textit{(definitional)}:
\begin{quote}\ttfamily\footnotesize\raggedright\noindent
forall(\allowbreak{}[a, b, f], is-hom-sk-space(\allowbreak{}a, b, f) iff is-sk-space(\allowbreak{}a) and is-sk-space(\allowbreak{}b) and f in fun(\allowbreak{}sk-names(\allowbreak{}a), sk-names(\allowbreak{}b)) and sk-inf(\allowbreak{}a) = sk-inf(\allowbreak{}b) and forall(\allowbreak{}[x1\_ in sk-names(\allowbreak{}a)], (\allowbreak{}ltkey(\allowbreak{}a))(\allowbreak{}x1\_) = (\allowbreak{}ltkey(\allowbreak{}b))(\allowbreak{}f(\allowbreak{}x1\_))))
\end{quote}
mentioned by 2 result(s)

\smallskip

\noindent\textbf{\texttt{is-hom-strand-space(a, b, f1, f2, f3)}}\quad\textit{predicate}\index{is-hom-strand-space(a, b, f1, f2, f3)@\texttt{is-hom-strand-space(a, b, f1, f2, f3)}}\\
reads: \$3 is a morphism of strand spaces from \$1 to \$2 \\
a predicate \\
characterized by \texttt{is-hom-strand-space-def}  \textit{(definitional)}:
\begin{quote}\ttfamily\footnotesize\raggedright\noindent
forall(\allowbreak{}[a, b, f1, f2, f3], is-hom-strand-space(\allowbreak{}a, b, f1, f2, f3) iff is-strand-space(\allowbreak{}a) and is-strand-space(\allowbreak{}b) and f1 in fun(\allowbreak{}tokens(\allowbreak{}a), tokens(\allowbreak{}b)) and f2 in fun(\allowbreak{}ids(\allowbreak{}a), ids(\allowbreak{}b)) and f3 in fun(\allowbreak{}keys(\allowbreak{}a), keys(\allowbreak{}b)) and forall(\allowbreak{}[x1\_ in ids(\allowbreak{}a)], (\allowbreak{}strand-map(\allowbreak{}a))(\allowbreak{}x1\_) = (\allowbreak{}strand-map(\allowbreak{}b))(\allowbreak{}f2(\allowbreak{}x1\_))) and forall(\allowbreak{}[x1\_ in keys(\allowbreak{}a)], f3(\allowbreak{}(\allowbreak{}kinv(\allowbreak{}a))(\allowbreak{}x1\_)) = (\allowbreak{}kinv(\allowbreak{}b))(\allowbreak{}f3(\allowbreak{}x1\_))))
\end{quote}
mentioned by 2 result(s)

\smallskip

\noindent\textbf{\texttt{is-hom-top-group(s, t, f)}}\quad\textit{predicate}\index{is-hom-top-group(s, t, f)@\texttt{is-hom-top-group(s, t, f)}}\\
reads: \$3 is a continuous homomorphism from \$1 to \$2 \\
a predicate \\
characterized by \texttt{is-hom-top-group-def}  \textit{(definitional)}:
\begin{quote}\ttfamily\footnotesize\raggedright\noindent
forall(\allowbreak{}[s, t, f], is-hom-top-group(\allowbreak{}s, t, f) iff is-top-group(\allowbreak{}s) and is-top-group(\allowbreak{}t) and f in fun(\allowbreak{}carr(\allowbreak{}s), carr(\allowbreak{}t)) and forall(\allowbreak{}[x1\_ in carr(\allowbreak{}s), x2\_ in carr(\allowbreak{}s)], f(\allowbreak{}(\allowbreak{}opr(\allowbreak{}s))(\allowbreak{}x1\_, x2\_)) = (\allowbreak{}opr(\allowbreak{}t))(\allowbreak{}f(\allowbreak{}x1\_), f(\allowbreak{}x2\_))) and f(\allowbreak{}iden(\allowbreak{}s)) = iden(\allowbreak{}t) and forall(\allowbreak{}[x1\_ in carr(\allowbreak{}s)], f(\allowbreak{}(\allowbreak{}inv(\allowbreak{}s))(\allowbreak{}x1\_)) = (\allowbreak{}inv(\allowbreak{}t))(\allowbreak{}f(\allowbreak{}x1\_))) and forall(\allowbreak{}[u in topology(\allowbreak{}t)], \{prmx\_ in carr(\allowbreak{}s): f(\allowbreak{}prmx\_) in u\} in topology(\allowbreak{}s)))
\end{quote}
mentioned by 4 result(s)

\smallskip

\noindent\textbf{\texttt{is-hom-top-space(s, t, f)}}\quad\textit{predicate}\index{is-hom-top-space(s, t, f)@\texttt{is-hom-top-space(s, t, f)}}\\
reads: \$3 is continuous from \$1 to \$2 \\
a predicate \\
characterized by \texttt{is-hom-top-space-def}  \textit{(definitional)}:
\begin{quote}\ttfamily\footnotesize\raggedright\noindent
forall(\allowbreak{}[s, t, f], is-hom-top-space(\allowbreak{}s, t, f) iff is-top-space(\allowbreak{}s) and is-top-space(\allowbreak{}t) and f in fun(\allowbreak{}pts(\allowbreak{}s), pts(\allowbreak{}t)) and forall(\allowbreak{}[u in opens(\allowbreak{}t)], preimage(\allowbreak{}s, f, u) in opens(\allowbreak{}s)))
\end{quote}
mentioned by 98 result(s)

\smallskip

\noindent\textbf{\texttt{is-hom-unital-banach-algebra(a, b, f)}}\quad\textit{predicate}\index{is-hom-unital-banach-algebra(a, b, f)@\texttt{is-hom-unital-banach-algebra(a, b, f)}}\\
reads: \$3 is a homomorphism from \$1 to \$2 \\
a predicate \\
characterized by \texttt{is-hom-unital-banach-algebra-def}  \textit{(definitional)}:
\begin{quote}\ttfamily\footnotesize\raggedright\noindent
forall(\allowbreak{}[a, b, f], is-hom-unital-banach-algebra(\allowbreak{}a, b, f) iff is-unital-banach-algebra(\allowbreak{}a) and is-unital-banach-algebra(\allowbreak{}b) and f in fun(\allowbreak{}vec(\allowbreak{}a), vec(\allowbreak{}b)) and scal(\allowbreak{}a) = scal(\allowbreak{}b) and forall(\allowbreak{}[x1\_ in vec(\allowbreak{}a), x2\_ in vec(\allowbreak{}a)], f(\allowbreak{}(\allowbreak{}vadd(\allowbreak{}a))(\allowbreak{}x1\_, x2\_)) = (\allowbreak{}vadd(\allowbreak{}b))(\allowbreak{}f(\allowbreak{}x1\_), f(\allowbreak{}x2\_))) and f(\allowbreak{}vzero(\allowbreak{}a)) = vzero(\allowbreak{}b) and forall(\allowbreak{}[x1\_ in vec(\allowbreak{}a)], f(\allowbreak{}(\allowbreak{}vneg(\allowbreak{}a))(\allowbreak{}x1\_)) = (\allowbreak{}vneg(\allowbreak{}b))(\allowbreak{}f(\allowbreak{}x1\_))) and forall(\allowbreak{}[x1\_ in carr(\allowbreak{}scal(\allowbreak{}a)), x2\_ in vec(\allowbreak{}a)], f(\allowbreak{}(\allowbreak{}act(\allowbreak{}a))(\allowbreak{}x1\_, x2\_)) = (\allowbreak{}act(\allowbreak{}b))(\allowbreak{}x1\_, f(\allowbreak{}x2\_))) and forall(\allowbreak{}[x1\_ in vec(\allowbreak{}a)], (\allowbreak{}vnrm(\allowbreak{}a))(\allowbreak{}x1\_) = (\allowbreak{}vnrm(\allowbreak{}b))(\allowbreak{}f(\allowbreak{}x1\_))) and forall(\allowbreak{}[x1\_ in vec(\allowbreak{}a), x2\_ in vec(\allowbreak{}a)], f(\allowbreak{}(\allowbreak{}vmul(\allowbreak{}a))(\allowbreak{}x1\_, x2\_)) = (\allowbreak{}vmul(\allowbreak{}b))(\allowbreak{}f(\allowbreak{}x1\_), f(\allowbreak{}x2\_))) and f(\allowbreak{}vone(\allowbreak{}a)) = vone(\allowbreak{}b))
\end{quote}
mentioned by 9 result(s)

\smallskip

\noindent\textbf{\texttt{is-hom-unital-c-star-algebra(a, b, f)}}\quad\textit{predicate}\index{is-hom-unital-c-star-algebra(a, b, f)@\texttt{is-hom-unital-c-star-algebra(a, b, f)}}\\
reads: \$3 is a homomorphism from \$1 to \$2 \\
a predicate \\
characterized by \texttt{is-hom-unital-c-star-algebra-def}  \textit{(definitional)}:
\begin{quote}\ttfamily\footnotesize\raggedright\noindent
forall(\allowbreak{}[a, b, f], is-hom-unital-c-star-algebra(\allowbreak{}a, b, f) iff is-unital-c-star-algebra(\allowbreak{}a) and is-unital-c-star-algebra(\allowbreak{}b) and is-hom-unital-involutive-banach-algebra(\allowbreak{}a, b, f))
\end{quote}
mentioned by 4 result(s)

\smallskip

\noindent\textbf{\texttt{is-hom-unital-involutive-banach-algebra(a, b, f)}}\quad\textit{predicate}\index{is-hom-unital-involutive-banach-algebra(a, b, f)@\texttt{is-hom-unital-involutive-banach-algebra(a, b, f)}}\\
reads: \$3 is a homomorphism from \$1 to \$2 \\
a predicate \\
characterized by \texttt{is-hom-unital-involutive-banach-algebra-def}  \textit{(definitional)}:
\begin{quote}\ttfamily\footnotesize\raggedright\noindent
forall(\allowbreak{}[a, b, f], is-hom-unital-involutive-banach-algebra(\allowbreak{}a, b, f) iff is-unital-involutive-banach-algebra(\allowbreak{}a) and is-unital-involutive-banach-algebra(\allowbreak{}b) and f in fun(\allowbreak{}vec(\allowbreak{}a), vec(\allowbreak{}b)) and scal(\allowbreak{}a) = scal(\allowbreak{}b) and forall(\allowbreak{}[x1\_ in vec(\allowbreak{}a), x2\_ in vec(\allowbreak{}a)], f(\allowbreak{}(\allowbreak{}vadd(\allowbreak{}a))(\allowbreak{}x1\_, x2\_)) = (\allowbreak{}vadd(\allowbreak{}b))(\allowbreak{}f(\allowbreak{}x1\_), f(\allowbreak{}x2\_))) and f(\allowbreak{}vzero(\allowbreak{}a)) = vzero(\allowbreak{}b) and forall(\allowbreak{}[x1\_ in vec(\allowbreak{}a)], f(\allowbreak{}(\allowbreak{}vneg(\allowbreak{}a))(\allowbreak{}x1\_)) = (\allowbreak{}vneg(\allowbreak{}b))(\allowbreak{}f(\allowbreak{}x1\_))) and forall(\allowbreak{}[x1\_ in carr(\allowbreak{}scal(\allowbreak{}a)), x2\_ in vec(\allowbreak{}a)], f(\allowbreak{}(\allowbreak{}act(\allowbreak{}a))(\allowbreak{}x1\_, x2\_)) = (\allowbreak{}act(\allowbreak{}b))(\allowbreak{}x1\_, f(\allowbreak{}x2\_))) and forall(\allowbreak{}[x1\_ in vec(\allowbreak{}a)], (\allowbreak{}vnrm(\allowbreak{}a))(\allowbreak{}x1\_) = (\allowbreak{}vnrm(\allowbreak{}b))(\allowbreak{}f(\allowbreak{}x1\_))) and forall(\allowbreak{}[x1\_ in vec(\allowbreak{}a), x2\_ in vec(\allowbreak{}a)], f(\allowbreak{}(\allowbreak{}vmul(\allowbreak{}a))(\allowbreak{}x1\_, x2\_)) = (\allowbreak{}vmul(\allowbreak{}b))(\allowbreak{}f(\allowbreak{}x1\_), f(\allowbreak{}x2\_))) and f(\allowbreak{}vone(\allowbreak{}a)) = vone(\allowbreak{}b) and forall(\allowbreak{}[x1\_ in vec(\allowbreak{}a)], f(\allowbreak{}(\allowbreak{}vstar(\allowbreak{}a))(\allowbreak{}x1\_)) = (\allowbreak{}vstar(\allowbreak{}b))(\allowbreak{}f(\allowbreak{}x1\_))))
\end{quote}
mentioned by 6 result(s)

\smallskip

\noindent\textbf{\texttt{is-hom-vector-space(a, b, f)}}\quad\textit{predicate}\index{is-hom-vector-space(a, b, f)@\texttt{is-hom-vector-space(a, b, f)}}\\
reads: \$3 is a homomorphism from \$1 to \$2 \\
a predicate \\
characterized by \texttt{is-hom-vector-space-def}  \textit{(definitional)}:
\begin{quote}\ttfamily\footnotesize\raggedright\noindent
forall(\allowbreak{}[a, b, f], is-hom-vector-space(\allowbreak{}a, b, f) iff is-vector-space(\allowbreak{}a) and is-vector-space(\allowbreak{}b) and is-hom-module(\allowbreak{}a, b, f))
\end{quote}
mentioned by 6 result(s)

\smallskip

\noindent\textbf{\texttt{is-homeomorphism(hmes\_, hmet\_, hmef\_)}}\quad\textit{predicate}\index{is-homeomorphism(hmes_, hmet_, hmef_)@\texttt{is-homeomorphism(hmes\_, hmet\_, hmef\_)}}\\
reads: \$3 is a homeomorphism from \$1 onto \$2 \\
defined by \texttt{is-homeomorphism} \\
characterized by \texttt{is-homeomorphism}  \textit{(definitional)}:
\begin{quote}\ttfamily\footnotesize\raggedright\noindent
forall(\allowbreak{}[hmes\_, hmet\_, hmef\_], is-homeomorphism(\allowbreak{}hmes\_, hmet\_, hmef\_) iff hmef\_ in bijection(\allowbreak{}pts(\allowbreak{}hmes\_), pts(\allowbreak{}hmet\_)) and is-continuous(\allowbreak{}hmes\_, hmet\_, hmef\_) and is-continuous(\allowbreak{}hmet\_, hmes\_, inverse-bij(\allowbreak{}hmef\_, pts(\allowbreak{}hmes\_), pts(\allowbreak{}hmet\_))))
\end{quote}
also: \texttt{compact-continuous-bijection-is-homeomorphism} \\
mentioned by 6 result(s)

\smallskip

\noindent\textbf{\texttt{is-honest(shp\_, shc\_, shi\_)}}\quad\textit{predicate}\index{is-honest(shp_, shc_, shi_)@\texttt{is-honest(shp\_, shc\_, shi\_)}}\\
reads: \$3 is honest relative to the bundle \$2 of the infiltrated space \$1 \\
defined by \texttt{is-honest} \\
characterized by \texttt{is-honest}  \textit{(definitional)}:
\begin{quote}\ttfamily\footnotesize\raggedright\noindent
forall(\allowbreak{}[shp\_, shc\_, shi\_], is-honest(\allowbreak{}shp\_, shc\_, shi\_) iff forall(\allowbreak{}[shn\_ in bnodes(\allowbreak{}shc\_)], is-adversary-node(\allowbreak{}shp\_, shn\_) implies is-entry-point(\allowbreak{}ispace(\allowbreak{}shp\_), shn\_, shi\_) implies is-text-trace(\allowbreak{}ispace(\allowbreak{}shp\_), adv-tokens(\allowbreak{}shp\_), (\allowbreak{}strand-map(\allowbreak{}ispace(\allowbreak{}shp\_)))(\allowbreak{}nth(\allowbreak{}2, shn\_)))))
\end{quote}
mentioned by 3 result(s)

\smallskip

\noindent\textbf{\texttt{is-ideal(s, i)}}\quad\textit{predicate}\index{is-ideal(s, i)@\texttt{is-ideal(s, i)}}\\
reads: \$2 is an ideal of \$1 \\
defined by \texttt{is-ideal} \\
characterized by \texttt{is-ideal}  \textit{(definitional)}:
\begin{quote}\ttfamily\footnotesize\raggedright\noindent
forall(\allowbreak{}[s, i], is-ideal(\allowbreak{}s, i) iff is-commutative-ring(\allowbreak{}s) and i subset carr(\allowbreak{}s) and zero(\allowbreak{}s) in i and forall(\allowbreak{}[a in i, b in i], (\allowbreak{}add(\allowbreak{}s))(\allowbreak{}a, b) in i) and forall(\allowbreak{}[a in i], (\allowbreak{}neg(\allowbreak{}s))(\allowbreak{}a) in i) and forall(\allowbreak{}[r in carr(\allowbreak{}s), a in i], (\allowbreak{}mul(\allowbreak{}s))(\allowbreak{}r, a) in i))
\end{quote}
also: \texttt{lastcoeff-set-is-ideal-cna-as-module} \texttt{lastcoeff-set-is-ideal-sip-as-module} \texttt{lastcoeff-set-is-ideal-uba-as-module} \texttt{lastcoeff-set-is-ideal-cnvs-as-module} \texttt{lastcoeff-set-is-ideal-sip-as-module-csip-as-sip} \texttt{lastcoeff-set-is-ideal-uba-as-module-uiba-as-uba} \texttt{lastcoeff-set-is-ideal} \texttt{zz-bezout-set-is-ideal} \\
mentioned by 16 result(s)

\smallskip

\noindent\textbf{\texttt{is-ideal-in(i\_, carr\_, add\_, mul\_, neg\_, zero\_)}}\quad\textit{predicate}\index{is-ideal-in(i_, carr_, add_, mul_, neg_, zero_)@\texttt{is-ideal-in(i\_, carr\_, add\_, mul\_, neg\_, zero\_)}}\\
reads: \$1 is a two-sided ideal of the ring with carrier \$2 \\
defined by \texttt{is-ideal-in} \\
characterized by \texttt{is-ideal-in}  \textit{(definitional)}:
\begin{quote}\ttfamily\footnotesize\raggedright\noindent
forall(\allowbreak{}[i\_, carr\_, add\_, mul\_, neg\_, zero\_], is-ideal-in(\allowbreak{}i\_, carr\_, add\_, mul\_, neg\_, zero\_) iff i\_ subset carr\_ and zero\_ in i\_ and forall(\allowbreak{}[a, b], a in i\_ implies b in i\_ implies add\_(\allowbreak{}a, b) in i\_) and forall(\allowbreak{}[a in i\_], neg\_(\allowbreak{}a) in i\_) and forall(\allowbreak{}[x\_, a], x\_ in carr\_ implies a in i\_ implies mul\_(\allowbreak{}x\_, a) in i\_) and forall(\allowbreak{}[x\_, a], x\_ in carr\_ implies a in i\_ implies mul\_(\allowbreak{}a, x\_) in i\_))
\end{quote}
mentioned by 4 result(s)

\smallskip

\noindent\textbf{\texttt{is-identity(op, unit, crr)}}\quad\textit{predicate}\index{is-identity(op, unit, crr)@\texttt{is-identity(op, unit, crr)}}\\
reads: \$2 is an identity for \$1 on \$3 \\
defined by \texttt{is-identity} \\
characterized by \texttt{is-identity}  \textit{(definitional)}:
\begin{quote}\ttfamily\footnotesize\raggedright\noindent
forall(\allowbreak{}[op, unit, crr], is-identity(\allowbreak{}op, unit, crr) iff forall(\allowbreak{}[u in crr], op(\allowbreak{}unit, u) = u and op(\allowbreak{}u, unit) = u))
\end{quote}
mentioned by 66 result(s)

\smallskip

\noindent\textbf{\texttt{is-infiltrated-space(s)}}\quad\textit{predicate}\index{is-infiltrated-space(s)@\texttt{is-infiltrated-space(s)}}\\
reads: an infiltrated strand space \\
a predicate \\
characterized by \texttt{is-infiltrated-space}  \textit{(definitional)}:
\begin{quote}\ttfamily\footnotesize\raggedright\noindent
forall(\allowbreak{}[s], is-infiltrated-space(\allowbreak{}s) iff length(\allowbreak{}s) = 3 and is-strand-space(\allowbreak{}ispace(\allowbreak{}s)) and adv-ids(\allowbreak{}s) in set and adv-tokens(\allowbreak{}s) in set and forall(\allowbreak{}[sap\_ in adv-ids(\allowbreak{}s)], sap\_ in ids(\allowbreak{}ispace(\allowbreak{}s))) and forall(\allowbreak{}[sat\_ in adv-tokens(\allowbreak{}s)], sat\_ in tokens(\allowbreak{}ispace(\allowbreak{}s))) and forall(\allowbreak{}[sap\_ in adv-ids(\allowbreak{}s)], is-adversary-trace(\allowbreak{}ispace(\allowbreak{}s), adv-tokens(\allowbreak{}s), (\allowbreak{}strand-map(\allowbreak{}ispace(\allowbreak{}s)))(\allowbreak{}sap\_))) and s in set)
\end{quote}
also: \texttt{cx-isp-is-infiltrated-space} \texttt{is-infiltrated-space-is-set} \\
mentioned by 25 result(s)

\smallskip

\noindent\textbf{\texttt{is-init-strand(nlp\_, nls\_, nla\_, nlb\_, nlm\_, nln\_)}}\quad\textit{predicate}\index{is-init-strand(nlp_, nls_, nla_, nlb_, nlm_, nln_)@\texttt{is-init-strand(nlp\_, nls\_, nla\_, nlb\_, nlm\_, nln\_)}}\\
reads: \$2 is an initiator's strand of \$1 for the names \$3, \$4 and the values \$5, \$6 \\
defined by \texttt{is-init-strand} \\
characterized by \texttt{is-init-strand}  \textit{(definitional)}:
\begin{quote}\ttfamily\footnotesize\raggedright\noindent
forall(\allowbreak{}[nlp\_, nls\_, nla\_, nlb\_, nlm\_, nln\_], is-init-strand(\allowbreak{}nlp\_, nls\_, nla\_, nlb\_, nlm\_, nln\_) iff nla\_ in names(\allowbreak{}nlp\_) and nlb\_ in names(\allowbreak{}nlp\_) and nlm\_ in tokens(\allowbreak{}ispace(\allowbreak{}pk-inf(\allowbreak{}nlp\_))) and nln\_ in tokens(\allowbreak{}ispace(\allowbreak{}pk-inf(\allowbreak{}nlp\_))) and not(\allowbreak{}nlm\_ in names(\allowbreak{}nlp\_)) and not(\allowbreak{}nlm\_ in keys(\allowbreak{}ispace(\allowbreak{}pk-inf(\allowbreak{}nlp\_)))) and not(\allowbreak{}nln\_ in keys(\allowbreak{}ispace(\allowbreak{}pk-inf(\allowbreak{}nlp\_)))) and nls\_ = [[sent, encr(\allowbreak{}lst(\allowbreak{}[nlm\_, nla\_]), (\allowbreak{}pubkey(\allowbreak{}nlp\_))(\allowbreak{}nlb\_))], [received, encr(\allowbreak{}lst(\allowbreak{}[nlm\_, nln\_, nlb\_]), (\allowbreak{}pubkey(\allowbreak{}nlp\_))(\allowbreak{}nla\_))], [sent, encr(\allowbreak{}nln\_, (\allowbreak{}pubkey(\allowbreak{}nlp\_))(\allowbreak{}nlb\_))]])
\end{quote}
mentioned by 16 result(s)

\smallskip

\noindent\textbf{\texttt{is-inner-product-form(ipfm\_, ipff\_)}}\quad\textit{predicate}\index{is-inner-product-form(ipfm_, ipff_)@\texttt{is-inner-product-form(ipfm\_, ipff\_)}}\\
reads: \$2 is an inner product on \$1 \\
defined by \texttt{is-inner-product-form} \\
characterized by \texttt{is-inner-product-form}  \textit{(definitional)}:
\begin{quote}\ttfamily\footnotesize\raggedright\noindent
forall(\allowbreak{}[ipfm\_, ipff\_], is-inner-product-form(\allowbreak{}ipfm\_, ipff\_) iff is-psd-form(\allowbreak{}ipfm\_, ipff\_) and forall(\allowbreak{}[ipfx\_ in vec(\allowbreak{}ipfm\_)], ipff\_(\allowbreak{}ipfx\_, ipfx\_) = 0 implies ipfx\_ = vzero(\allowbreak{}ipfm\_)))
\end{quote}
mentioned by 140 result(s)

\smallskip

\noindent\textbf{\texttt{is-integrable(omega, ca, mu, f)}}\quad\textit{predicate}\index{is-integrable(omega, ca, mu, f)@\texttt{is-integrable(omega, ca, mu, f)}}\\
reads: \$4 is integrable on \$1 with respect to the measure \$3 \\
defined by \texttt{is-integrable} \\
characterized by \texttt{is-integrable}  \textit{(definitional)}:
\begin{quote}\ttfamily\footnotesize\raggedright\noindent
forall(\allowbreak{}[omega, ca, mu, f], is-integrable(\allowbreak{}omega, ca, mu, f) iff is-real-measurable-fn(\allowbreak{}omega, ca, f) and integral(\allowbreak{}omega, ca, mu, pos-part(\allowbreak{}omega, f)) in rr and integral(\allowbreak{}omega, ca, mu, neg-part(\allowbreak{}omega, f)) in rr)
\end{quote}
mentioned by 58 result(s)

\smallskip

\noindent\textbf{\texttt{is-integral-domain(s)}}\quad\textit{predicate}\index{is-integral-domain(s)@\texttt{is-integral-domain(s)}}\\
reads: an integral domain \\
a predicate \\
characterized by \texttt{is-integral-domain-def}  \textit{(definitional)}:
\begin{quote}\ttfamily\footnotesize\raggedright\noindent
forall(\allowbreak{}[s], is-integral-domain(\allowbreak{}s) iff is-commutative-ring(\allowbreak{}s) and not(\allowbreak{}one(\allowbreak{}s) = zero(\allowbreak{}s)) and forall(\allowbreak{}[a in carr(\allowbreak{}s), b in carr(\allowbreak{}s)], (\allowbreak{}mul(\allowbreak{}s))(\allowbreak{}a, b) = zero(\allowbreak{}s) implies a = zero(\allowbreak{}s) or b = zero(\allowbreak{}s)))
\end{quote}
also: \texttt{field-as-integral-domain-is-integral-domain} \texttt{normed-field-as-integral-domain-is-integral-domain} \texttt{qq-is-integral-domain} \texttt{zz-is-integral-domain} \texttt{poly-is-integral-domain} \texttt{is-integral-domain-is-set} \texttt{field-ring-is-integral-domain} \texttt{euclidean-ring-is-integral-domain} \\
mentioned by 56 result(s)

\smallskip

\noindent\textbf{\texttt{is-intertwined(amu\_, amz\_, amr\_, amq\_, amv\_, ams\_)}}\quad\textit{predicate}\index{is-intertwined(amu_, amz_, amr_, amq_, amv_, ams_)@\texttt{is-intertwined(amu\_, amz\_, amr\_, amq\_, amv\_, ams\_)}}\\
reads: every operator from \$1 to \$2 intertwining \$3 and \$4 on \$6 intertwines them on the bicommutant of \$6 \\
defined by \texttt{is-intertwined} \\
characterized by \texttt{is-intertwined}  \textit{(definitional)}:
\begin{quote}\ttfamily\footnotesize\raggedright\noindent
forall(\allowbreak{}[amu\_, amz\_, amr\_, amq\_, amv\_, ams\_], is-intertwined(\allowbreak{}amu\_, amz\_, amr\_, amq\_, amv\_, ams\_) iff forall(\allowbreak{}[amk\_ in bounded-ops(\allowbreak{}amu\_, amz\_)], forall(\allowbreak{}[ama\_ in ams\_, amx\_ in vec(\allowbreak{}amu\_)], amk\_(\allowbreak{}(\allowbreak{}amr\_(\allowbreak{}ama\_))(\allowbreak{}amx\_)) = (\allowbreak{}amq\_(\allowbreak{}ama\_))(\allowbreak{}amk\_(\allowbreak{}amx\_))) implies forall(\allowbreak{}[amt\_ in bicommutant(\allowbreak{}amv\_, ams\_), amx\_ in vec(\allowbreak{}amu\_)], amk\_(\allowbreak{}(\allowbreak{}amr\_(\allowbreak{}amt\_))(\allowbreak{}amx\_)) = (\allowbreak{}amq\_(\allowbreak{}amt\_))(\allowbreak{}amk\_(\allowbreak{}amx\_)))))
\end{quote}
mentioned by 7 result(s)

\smallskip

\noindent\textbf{\texttt{is-interval(es)}}\quad\textit{predicate}\index{is-interval(es)@\texttt{is-interval(es)}}\\
reads: \$1 is an interval of real numbers \\
defined by \texttt{is-interval} \\
characterized by \texttt{is-interval}  \textit{(definitional)}:
\begin{quote}\ttfamily\footnotesize\raggedright\noindent
forall(\allowbreak{}[es], is-interval(\allowbreak{}es) iff es subset rr and forall(\allowbreak{}[mia\_ in es, mib\_ in es, mix\_ in rr], mia\_ <= mix\_ implies mix\_ <= mib\_ implies mix\_ in es))
\end{quote}
also: \texttt{is-interval-rr} \texttt{is-interval-inter} \\
mentioned by 14 result(s)

\smallskip

\noindent\textbf{\texttt{is-invariant-ccs(icm\_, ica\_, icg\_, icb\_)}}\quad\textit{predicate}\index{is-invariant-ccs(icm_, ica_, icg_, icb_)@\texttt{is-invariant-ccs(icm\_, ica\_, icg\_, icb\_)}}\\
reads: \$4 is a nonempty closed convex subset of \$2 in \$1 invariant under \$3 \\
defined by \texttt{is-invariant-ccs} \\
characterized by \texttt{is-invariant-ccs}  \textit{(definitional)}:
\begin{quote}\ttfamily\footnotesize\raggedright\noindent
forall(\allowbreak{}[icm\_, ica\_, icg\_, icb\_], is-invariant-ccs(\allowbreak{}icm\_, ica\_, icg\_, icb\_) iff icb\_ subset ica\_ and forsome(\allowbreak{}[icz\_], icz\_ in icb\_) and is-convex-subset(\allowbreak{}icm\_, icb\_) and is-closed(\allowbreak{}subspace-ms(\allowbreak{}nvs-metric-space(\allowbreak{}icm\_), ica\_), icb\_) and forall(\allowbreak{}[icf\_ in icg\_, icx\_ in icb\_], icf\_(\allowbreak{}icx\_) in icb\_))
\end{quote}
mentioned by 4 result(s)

\smallskip

\noindent\textbf{\texttt{is-invariant-factors(l, a)}}\quad\textit{predicate}\index{is-invariant-factors(l, a)@\texttt{is-invariant-factors(l, a)}}\\
reads: \$2\_1, ..., \$2\_\$1 are invariant factors (non-negative, none equal to 1, each dividing the next) \\
defined by \texttt{is-invariant-factors} \\
characterized by \texttt{is-invariant-factors}  \textit{(definitional)}:
\begin{quote}\ttfamily\footnotesize\raggedright\noindent
forall(\allowbreak{}[l, a], is-invariant-factors(\allowbreak{}l, a) iff a in fun(\allowbreak{}interval(\allowbreak{}1, l), zz) and forall(\allowbreak{}[ifi\_ in interval(\allowbreak{}1, l)], 0 <= a(\allowbreak{}ifi\_)) and forall(\allowbreak{}[ifi\_ in interval(\allowbreak{}1, l)], not(\allowbreak{}a(\allowbreak{}ifi\_) = 1)) and forall(\allowbreak{}[ifi\_, ifj\_], ifi\_ in interval(\allowbreak{}1, l) implies ifj\_ in interval(\allowbreak{}1, l) implies ifi\_ <= ifj\_ implies zz-divides(\allowbreak{}a(\allowbreak{}ifi\_), a(\allowbreak{}ifj\_))))
\end{quote}
mentioned by 18 result(s)

\smallskip

\noindent\textbf{\texttt{is-invariant-subspace-rep(g4a\_, g4v\_, g4p\_, g4k\_)}}\quad\textit{predicate}\index{is-invariant-subspace-rep(g4a_, g4v_, g4p_, g4k_)@\texttt{is-invariant-subspace-rep(g4a\_, g4v\_, g4p\_, g4k\_)}}\\
reads: \$4 is a closed subspace invariant under \$3 \\
defined by \texttt{is-invariant-subspace-rep} \\
characterized by \texttt{is-invariant-subspace-rep}  \textit{(definitional)}:
\begin{quote}\ttfamily\footnotesize\raggedright\noindent
forall(\allowbreak{}[g4a\_, g4v\_, g4p\_, g4k\_], is-invariant-subspace-rep(\allowbreak{}g4a\_, g4v\_, g4p\_, g4k\_) iff is-submodule(\allowbreak{}g4v\_, g4k\_) and is-closed(\allowbreak{}ip-metric-space(\allowbreak{}g4v\_), g4k\_) and forall(\allowbreak{}[g4x\_ in vec(\allowbreak{}g4a\_), g4y\_ in g4k\_], (\allowbreak{}g4p\_(\allowbreak{}g4x\_))(\allowbreak{}g4y\_) in g4k\_))
\end{quote}
mentioned by 4 result(s)

\smallskip

\noindent\textbf{\texttt{is-invertible-el(baa\_, bax\_)}}\quad\textit{predicate}\index{is-invertible-el(baa_, bax_)@\texttt{is-invertible-el(baa\_, bax\_)}}\\
reads: \$2 is invertible in \$1 \\
defined by \texttt{is-invertible-el} \\
characterized by \texttt{is-invertible-el}  \textit{(definitional)}:
\begin{quote}\ttfamily\footnotesize\raggedright\noindent
forall(\allowbreak{}[baa\_, bax\_], is-invertible-el(\allowbreak{}baa\_, bax\_) iff bax\_ in vec(\allowbreak{}baa\_) and forsome(\allowbreak{}[bau\_ in vec(\allowbreak{}baa\_)], (\allowbreak{}vmul(\allowbreak{}baa\_))(\allowbreak{}bau\_, bax\_) = vone(\allowbreak{}baa\_) and (\allowbreak{}vmul(\allowbreak{}baa\_))(\allowbreak{}bax\_, bau\_) = vone(\allowbreak{}baa\_)))
\end{quote}
mentioned by 76 result(s)

\smallskip

\noindent\textbf{\texttt{is-invertible-mat(a, n, u)}}\quad\textit{predicate}\index{is-invertible-mat(a, n, u)@\texttt{is-invertible-mat(a, n, u)}}\\
reads: \$3 is an invertible \$2-by-\$2 matrix over \$1 \\
defined by \texttt{is-invertible-mat} \\
characterized by \texttt{is-invertible-mat}  \textit{(definitional)}:
\begin{quote}\ttfamily\footnotesize\raggedright\noindent
forall(\allowbreak{}[a, n, u], is-invertible-mat(\allowbreak{}a, n, u) iff u in mat(\allowbreak{}n, n, carr(\allowbreak{}a)) and forsome(\allowbreak{}[v in mat(\allowbreak{}n, n, carr(\allowbreak{}a))], matmul(\allowbreak{}a, u, v) = identmat(\allowbreak{}a, n) and matmul(\allowbreak{}a, v, u) = identmat(\allowbreak{}a, n)))
\end{quote}
mentioned by 177 result(s)

\smallskip

\noindent\textbf{\texttt{is-irreducible-mat(k, n, tm)}}\quad\textit{predicate}\index{is-irreducible-mat(k, n, tm)@\texttt{is-irreducible-mat(k, n, tm)}}\\
reads: the \$2-by-\$2 matrix \$3 over \$1 is irreducible \\
defined by \texttt{is-irreducible-mat} \\
characterized by \texttt{is-irreducible-mat}  \textit{(definitional)}:
\begin{quote}\ttfamily\footnotesize\raggedright\noindent
forall(\allowbreak{}[k, n, tm], is-irreducible-mat(\allowbreak{}k, n, tm) iff tm in mat(\allowbreak{}n, n, carr(\allowbreak{}k)) and forall(\allowbreak{}[zkvp\_ in mat(\allowbreak{}n, 1, carr(\allowbreak{}k))], not(\allowbreak{}zkvp\_ = zeromat(\allowbreak{}normed-field-as-commutative-ring(\allowbreak{}k), n, 1)) implies forall(\allowbreak{}[zkvx\_ in mat(\allowbreak{}n, 1, carr(\allowbreak{}k))], forsome(\allowbreak{}[zkvc\_ in mat(\allowbreak{}n, 1, carr(\allowbreak{}k))], zkvx\_ = matmul(\allowbreak{}normed-field-as-commutative-ring(\allowbreak{}k), krylov-mat(\allowbreak{}k, n, tm, zkvp\_), zkvc\_)))))
\end{quote}
mentioned by 12 result(s)

\smallskip

\noindent\textbf{\texttt{is-irreducible-rep(g4a\_, g4v\_, g4p\_)}}\quad\textit{predicate}\index{is-irreducible-rep(g4a_, g4v_, g4p_)@\texttt{is-irreducible-rep(g4a\_, g4v\_, g4p\_)}}\\
reads: \$3 is an irreducible representation of \$1 on \$2 \\
defined by \texttt{is-irreducible-rep} \\
characterized by \texttt{is-irreducible-rep}  \textit{(definitional)}:
\begin{quote}\ttfamily\footnotesize\raggedright\noindent
forall(\allowbreak{}[g4a\_, g4v\_, g4p\_], is-irreducible-rep(\allowbreak{}g4a\_, g4v\_, g4p\_) iff is-representation(\allowbreak{}g4a\_, g4v\_, g4p\_) and forall(\allowbreak{}[g4k\_], is-invariant-subspace-rep(\allowbreak{}g4a\_, g4v\_, g4p\_, g4k\_) implies g4k\_ = vec(\allowbreak{}g4v\_) or g4k\_ = \{vzero(\allowbreak{}g4v\_)\}))
\end{quote}
mentioned by 2 result(s)

\smallskip

\noindent\textbf{\texttt{is-iso-group(qag\_, qah\_, azk\_)}}\quad\textit{predicate}\index{is-iso-group(qag_, qah_, azk_)@\texttt{is-iso-group(qag\_, qah\_, azk\_)}}\\
reads: \$3 is an isomorphism of the group \$1 onto the group \$2 \\
defined by \texttt{is-iso-group} \\
characterized by \texttt{is-iso-group}  \textit{(definitional)}:
\begin{quote}\ttfamily\footnotesize\raggedright\noindent
forall(\allowbreak{}[qag\_, qah\_, azk\_], is-iso-group(\allowbreak{}qag\_, qah\_, azk\_) iff is-hom-group(\allowbreak{}qag\_, qah\_, azk\_) and azk\_ in bijection(\allowbreak{}carr(\allowbreak{}qag\_), carr(\allowbreak{}qah\_)))
\end{quote}
mentioned by 2 result(s)

\smallskip

\noindent\textbf{\texttt{is-iso-module(a, b, f)}}\quad\textit{predicate}\index{is-iso-module(a, b, f)@\texttt{is-iso-module(a, b, f)}}\\
reads: \$3 is an isomorphism of the module \$1 onto the module \$2 \\
defined by \texttt{is-iso-module} \\
characterized by \texttt{is-iso-module}  \textit{(definitional)}:
\begin{quote}\ttfamily\footnotesize\raggedright\noindent
forall(\allowbreak{}[a, b, f], is-iso-module(\allowbreak{}a, b, f) iff is-hom-module(\allowbreak{}a, b, f) and f in bijection(\allowbreak{}vec(\allowbreak{}a), vec(\allowbreak{}b)))
\end{quote}
mentioned by 27 result(s)

\smallskip

\noindent\textbf{\texttt{is-iso-ring(a, b, f)}}\quad\textit{predicate}\index{is-iso-ring(a, b, f)@\texttt{is-iso-ring(a, b, f)}}\\
reads: \$3 is a ring isomorphism from \$1 to \$2 \\
defined by \texttt{is-iso-ring} \\
characterized by \texttt{is-iso-ring}  \textit{(definitional)}:
\begin{quote}\ttfamily\footnotesize\raggedright\noindent
forall(\allowbreak{}[a, b, f], is-iso-ring(\allowbreak{}a, b, f) iff is-hom-ring(\allowbreak{}a, b, f) and f in bijection(\allowbreak{}carr(\allowbreak{}a), carr(\allowbreak{}b)))
\end{quote}
mentioned by 5 result(s)

\smallskip

\noindent\textbf{\texttt{is-isometric-graph(ubv\_, ubg\_)}}\quad\textit{predicate}\index{is-isometric-graph(ubv_, ubg_)@\texttt{is-isometric-graph(ubv\_, ubg\_)}}\\
reads: \$2 is an isometric graph \\
defined by \texttt{is-isometric-graph} \\
characterized by \texttt{is-isometric-graph}  \textit{(definitional)}:
\begin{quote}\ttfamily\footnotesize\raggedright\noindent
forall(\allowbreak{}[ubv\_, ubg\_], is-isometric-graph(\allowbreak{}ubv\_, ubg\_) iff forall(\allowbreak{}[ubp\_ in ubg\_], ip-norm(\allowbreak{}ubv\_, nth(\allowbreak{}1, ubp\_)) = ip-norm(\allowbreak{}ubv\_, nth(\allowbreak{}2, ubp\_))))
\end{quote}
mentioned by 8 result(s)

\smallskip

\noindent\textbf{\texttt{is-isometric-on(bov\_, bow\_, bot\_, bok\_)}}\quad\textit{predicate}\index{is-isometric-on(bov_, bow_, bot_, bok_)@\texttt{is-isometric-on(bov\_, bow\_, bot\_, bok\_)}}\\
reads: \$3 is isometric on \$4 \\
defined by \texttt{is-isometric-on} \\
characterized by \texttt{is-isometric-on}  \textit{(definitional)}:
\begin{quote}\ttfamily\footnotesize\raggedright\noindent
forall(\allowbreak{}[bov\_, bow\_, bot\_, bok\_], is-isometric-on(\allowbreak{}bov\_, bow\_, bot\_, bok\_) iff forall(\allowbreak{}[box\_ in bok\_], ip-norm(\allowbreak{}bow\_, bot\_(\allowbreak{}box\_)) = ip-norm(\allowbreak{}bov\_, box\_)))
\end{quote}
mentioned by 6 result(s)

\smallskip

\noindent\textbf{\texttt{is-isometric-op(v, jo)}}\quad\textit{predicate}\index{is-isometric-op(v, jo)@\texttt{is-isometric-op(v, jo)}}\\
reads: \$2 is isometric on \$1 \\
defined by \texttt{is-isometric-op} \\
characterized by \texttt{is-isometric-op}  \textit{(definitional)}:
\begin{quote}\ttfamily\footnotesize\raggedright\noindent
forall(\allowbreak{}[v, jo], is-isometric-op(\allowbreak{}v, jo) iff forall(\allowbreak{}[x\_ in vec(\allowbreak{}v), y\_ in vec(\allowbreak{}v)], (\allowbreak{}rip(\allowbreak{}v))(\allowbreak{}jo(\allowbreak{}x\_), jo(\allowbreak{}y\_)) = (\allowbreak{}rip(\allowbreak{}v))(\allowbreak{}x\_, y\_)))
\end{quote}
mentioned by 5 result(s)

\smallskip

\noindent\textbf{\texttt{is-isometry(s, t, f)}}\quad\textit{predicate}\index{is-isometry(s, t, f)@\texttt{is-isometry(s, t, f)}}\\
reads: \$3 is an isometry from \$1 to \$2 \\
defined by \texttt{is-isometry} \\
characterized by \texttt{is-isometry}  \textit{(definitional)}:
\begin{quote}\ttfamily\footnotesize\raggedright\noindent
forall(\allowbreak{}[s, t, f], is-isometry(\allowbreak{}s, t, f) iff is-metric-space(\allowbreak{}s) and is-metric-space(\allowbreak{}t) and f in fun(\allowbreak{}pts(\allowbreak{}s), pts(\allowbreak{}t)) and forall(\allowbreak{}[x in pts(\allowbreak{}s), y in pts(\allowbreak{}s)], (\allowbreak{}dist(\allowbreak{}t))(\allowbreak{}f(\allowbreak{}x), f(\allowbreak{}y)) = (\allowbreak{}dist(\allowbreak{}s))(\allowbreak{}x, y)))
\end{quote}
mentioned by 6 result(s)

\smallskip

\noindent\textbf{\texttt{is-k-bilinear-form(c6vk\_, c6vm\_, c6vl\_)}}\quad\textit{predicate}\index{is-k-bilinear-form(c6vk_, c6vm_, c6vl_)@\texttt{is-k-bilinear-form(c6vk\_, c6vm\_, c6vl\_)}}\\
reads: \$3 is a \$1-bilinear form on \$2 \\
defined by \texttt{is-k-bilinear-form} \\
characterized by \texttt{is-k-bilinear-form}  \textit{(definitional)}:
\begin{quote}\ttfamily\footnotesize\raggedright\noindent
forall(\allowbreak{}[c6vk\_, c6vm\_, c6vl\_], is-k-bilinear-form(\allowbreak{}c6vk\_, c6vm\_, c6vl\_) iff c6vl\_ in fun(\allowbreak{}cartesian(\allowbreak{}vec(\allowbreak{}c6vm\_), vec(\allowbreak{}c6vm\_)), carr(\allowbreak{}c6vk\_)) and forall(\allowbreak{}[c6vx\_ in vec(\allowbreak{}c6vm\_), c6vy\_ in vec(\allowbreak{}c6vm\_), c6vz\_ in vec(\allowbreak{}c6vm\_)], c6vl\_(\allowbreak{}(\allowbreak{}vadd(\allowbreak{}c6vm\_))(\allowbreak{}c6vx\_, c6vy\_), c6vz\_) = (\allowbreak{}add(\allowbreak{}c6vk\_))(\allowbreak{}c6vl\_(\allowbreak{}c6vx\_, c6vz\_), c6vl\_(\allowbreak{}c6vy\_, c6vz\_)) and c6vl\_(\allowbreak{}c6vz\_, (\allowbreak{}vadd(\allowbreak{}c6vm\_))(\allowbreak{}c6vx\_, c6vy\_)) = (\allowbreak{}add(\allowbreak{}c6vk\_))(\allowbreak{}c6vl\_(\allowbreak{}c6vz\_, c6vx\_), c6vl\_(\allowbreak{}c6vz\_, c6vy\_))) and forall(\allowbreak{}[c6va\_ in carr(\allowbreak{}c6vk\_), c6vx\_ in vec(\allowbreak{}c6vm\_), c6vy\_ in vec(\allowbreak{}c6vm\_)], c6vl\_(\allowbreak{}(\allowbreak{}act(\allowbreak{}c6vm\_))(\allowbreak{}c6va\_, c6vx\_), c6vy\_) = (\allowbreak{}mul(\allowbreak{}c6vk\_))(\allowbreak{}c6va\_, c6vl\_(\allowbreak{}c6vx\_, c6vy\_)) and c6vl\_(\allowbreak{}c6vx\_, (\allowbreak{}act(\allowbreak{}c6vm\_))(\allowbreak{}c6va\_, c6vy\_)) = (\allowbreak{}mul(\allowbreak{}c6vk\_))(\allowbreak{}c6va\_, c6vl\_(\allowbreak{}c6vx\_, c6vy\_))))
\end{quote}
mentioned by 15 result(s)

\smallskip

\noindent\textbf{\texttt{is-k-linear(m, f)}}\quad\textit{predicate}\index{is-k-linear(m, f)@\texttt{is-k-linear(m, f)}}\\
reads: \$2 is a K-linear functional on \$1 \\
defined by \texttt{is-k-linear} \\
characterized by \texttt{is-k-linear}  \textit{(definitional)}:
\begin{quote}\ttfamily\footnotesize\raggedright\noindent
forall(\allowbreak{}[m, f], is-k-linear(\allowbreak{}m, f) iff is-module(\allowbreak{}m) and f in fun(\allowbreak{}vec(\allowbreak{}m), carr(\allowbreak{}scal(\allowbreak{}m))) and forall(\allowbreak{}[x in vec(\allowbreak{}m), y in vec(\allowbreak{}m)], f(\allowbreak{}(\allowbreak{}vadd(\allowbreak{}m))(\allowbreak{}x, y)) = (\allowbreak{}add(\allowbreak{}scal(\allowbreak{}m)))(\allowbreak{}f(\allowbreak{}x), f(\allowbreak{}y))) and forall(\allowbreak{}[lam in carr(\allowbreak{}scal(\allowbreak{}m)), x in vec(\allowbreak{}m)], f(\allowbreak{}(\allowbreak{}act(\allowbreak{}m))(\allowbreak{}lam, x)) = (\allowbreak{}mul(\allowbreak{}scal(\allowbreak{}m)))(\allowbreak{}lam, f(\allowbreak{}x))))
\end{quote}
mentioned by 13 result(s)

\smallskip

\noindent\textbf{\texttt{is-k-linear-on(m, s, f)}}\quad\textit{predicate}\index{is-k-linear-on(m, s, f)@\texttt{is-k-linear-on(m, s, f)}}\\
reads: \$3 is a K-linear functional on the submodule \$2 of \$1 \\
defined by \texttt{is-k-linear-on} \\
characterized by \texttt{is-k-linear-on}  \textit{(definitional)}:
\begin{quote}\ttfamily\footnotesize\raggedright\noindent
forall(\allowbreak{}[m, s, f], is-k-linear-on(\allowbreak{}m, s, f) iff is-module(\allowbreak{}m) and is-submodule(\allowbreak{}m, s) and f in fun(\allowbreak{}s, carr(\allowbreak{}scal(\allowbreak{}m))) and forall(\allowbreak{}[x in s, y in s], f(\allowbreak{}(\allowbreak{}vadd(\allowbreak{}m))(\allowbreak{}x, y)) = (\allowbreak{}add(\allowbreak{}scal(\allowbreak{}m)))(\allowbreak{}f(\allowbreak{}x), f(\allowbreak{}y))) and forall(\allowbreak{}[lam in carr(\allowbreak{}scal(\allowbreak{}m)), x in s], f(\allowbreak{}(\allowbreak{}act(\allowbreak{}m))(\allowbreak{}lam, x)) = (\allowbreak{}mul(\allowbreak{}scal(\allowbreak{}m)))(\allowbreak{}lam, f(\allowbreak{}x))))
\end{quote}
mentioned by 12 result(s)

\smallskip

\noindent\textbf{\texttt{is-k-vector-space(c6vk\_, c6vm\_)}}\quad\textit{predicate}\index{is-k-vector-space(c6vk_, c6vm_)@\texttt{is-k-vector-space(c6vk\_, c6vm\_)}}\\
reads: \$2 is a vector space over \$1 \\
defined by \texttt{is-k-vector-space} \\
characterized by \texttt{is-k-vector-space}  \textit{(definitional)}:
\begin{quote}\ttfamily\footnotesize\raggedright\noindent
forall(\allowbreak{}[c6vk\_, c6vm\_], is-k-vector-space(\allowbreak{}c6vk\_, c6vm\_) iff is-module(\allowbreak{}c6vm\_) and scal(\allowbreak{}c6vm\_) = normed-field-as-commutative-ring(\allowbreak{}c6vk\_))
\end{quote}
mentioned by 53 result(s)

\smallskip

\noindent\textbf{\texttt{is-key-ideal(sbit\_, sbik\_, sbiks\_, sbii\_)}}\quad\textit{predicate}\index{is-key-ideal(sbit_, sbik_, sbiks_, sbii_)@\texttt{is-key-ideal(sbit\_, sbik\_, sbiks\_, sbii\_)}}\\
reads: \$4 is a \$3-ideal of the messages over the tokens \$1 and the keys \$2 \\
defined by \texttt{is-key-ideal} \\
characterized by \texttt{is-key-ideal}  \textit{(definitional)}:
\begin{quote}\ttfamily\footnotesize\raggedright\noindent
forall(\allowbreak{}[sbit\_, sbik\_, sbiks\_, sbii\_], is-key-ideal(\allowbreak{}sbit\_, sbik\_, sbiks\_, sbii\_) iff forall(\allowbreak{}[sbix\_ in sbii\_], sbix\_ in msg-set(\allowbreak{}sbit\_, sbik\_)) and forall(\allowbreak{}[sbil\_, sbij\_], sbil\_ in tuples(\allowbreak{}msg-set(\allowbreak{}sbit\_, sbik\_)) implies sbij\_ in nn implies 1 <= sbij\_ implies sbij\_ <= length(\allowbreak{}sbil\_) implies nth(\allowbreak{}sbij\_, sbil\_) in sbii\_ implies lst(\allowbreak{}sbil\_) in sbii\_) and forall(\allowbreak{}[sbih\_, sbiy\_], sbih\_ in sbii\_ implies sbiy\_ in sbiks\_ implies encr(\allowbreak{}sbih\_, sbiy\_) in sbii\_))
\end{quote}
mentioned by 9 result(s)

\smallskip

\noindent\textbf{\texttt{is-laurent(lrk\_, lrf\_)}}\quad\textit{predicate}\index{is-laurent(lrk_, lrf_)@\texttt{is-laurent(lrk\_, lrf\_)}}\\
reads: \$2 is a Laurent series over \$1 \\
defined by \texttt{is-laurent} \\
characterized by \texttt{is-laurent}  \textit{(definitional)}:
\begin{quote}\ttfamily\footnotesize\raggedright\noindent
forall(\allowbreak{}[lrk\_, lrf\_], is-laurent(\allowbreak{}lrk\_, lrf\_) iff lrf\_ in fun(\allowbreak{}zz, carr(\allowbreak{}lrk\_)) and forsome(\allowbreak{}[lrb\_ in zz], forall(\allowbreak{}[lrj\_ in zz], lrj\_ < lrb\_ implies lrf\_(\allowbreak{}lrj\_) = zero(\allowbreak{}lrk\_))))
\end{quote}
mentioned by 4 result(s)

\smallskip

\noindent\textbf{\texttt{is-left-ideal(baa\_, bai\_)}}\quad\textit{predicate}\index{is-left-ideal(baa_, bai_)@\texttt{is-left-ideal(baa\_, bai\_)}}\\
reads: \$2 is a left ideal of \$1 \\
defined by \texttt{is-left-ideal} \\
characterized by \texttt{is-left-ideal}  \textit{(definitional)}:
\begin{quote}\ttfamily\footnotesize\raggedright\noindent
forall(\allowbreak{}[baa\_, bai\_], is-left-ideal(\allowbreak{}baa\_, bai\_) iff is-alg-subspace(\allowbreak{}baa\_, bai\_) and forall(\allowbreak{}[bay\_ in vec(\allowbreak{}baa\_), bax\_ in bai\_], (\allowbreak{}vmul(\allowbreak{}baa\_))(\allowbreak{}bay\_, bax\_) in bai\_))
\end{quote}
mentioned by 23 result(s)

\smallskip

\noindent\textbf{\texttt{is-left-invariant(lim\_, lig\_)}}\quad\textit{predicate}\index{is-left-invariant(lim_, lig_)@\texttt{is-left-invariant(lim\_, lig\_)}}\\
reads: the distance of \$1 is left invariant under \$2 \\
defined by \texttt{is-left-invariant} \\
characterized by \texttt{is-left-invariant}  \textit{(definitional)}:
\begin{quote}\ttfamily\footnotesize\raggedright\noindent
forall(\allowbreak{}[lim\_, lig\_], is-left-invariant(\allowbreak{}lim\_, lig\_) iff forall(\allowbreak{}[lia\_ in carr(\allowbreak{}lig\_), lix\_ in carr(\allowbreak{}lig\_), liy\_ in carr(\allowbreak{}lig\_)], (\allowbreak{}dist(\allowbreak{}lim\_))(\allowbreak{}(\allowbreak{}opr(\allowbreak{}lig\_))(\allowbreak{}lia\_, lix\_), (\allowbreak{}opr(\allowbreak{}lig\_))(\allowbreak{}lia\_, liy\_)) = (\allowbreak{}dist(\allowbreak{}lim\_))(\allowbreak{}lix\_, liy\_)))
\end{quote}
mentioned by 5 result(s)

\smallskip

\noindent\textbf{\texttt{is-left-limit(f, x, l)}}\quad\textit{predicate}\index{is-left-limit(f, x, l)@\texttt{is-left-limit(f, x, l)}}\\
reads: \$3 is the limit of \$1 as the argument approaches \$2 from the left \\
defined by \texttt{is-left-limit} \\
characterized by \texttt{is-left-limit}  \textit{(definitional)}:
\begin{quote}\ttfamily\footnotesize\raggedright\noindent
forall(\allowbreak{}[f, x, l], is-left-limit(\allowbreak{}f, x, l) iff f in fun(\allowbreak{}rr, rr) and x in rr and l in rr and forall(\allowbreak{}[eps\_], pos-rr(\allowbreak{}eps\_) implies forsome(\allowbreak{}[delta\_], pos-rr(\allowbreak{}delta\_) and forall(\allowbreak{}[t\_], t\_ in rr and x - delta\_ <= t\_ and t\_ < x implies abs(\allowbreak{}f(\allowbreak{}t\_) - l) <= eps\_))))
\end{quote}
mentioned by 6 result(s)

\smallskip

\noindent\textbf{\texttt{is-left-limit-within(rpf\_, rpw\_, rpx\_, rpl\_)}}\quad\textit{predicate}\index{is-left-limit-within(rpf_, rpw_, rpx_, rpl_)@\texttt{is-left-limit-within(rpf\_, rpw\_, rpx\_, rpl\_)}}\\
reads: \$1 has left limit \$4 at \$3 within \$2 \\
defined by \texttt{is-left-limit-within} \\
characterized by \texttt{is-left-limit-within}  \textit{(definitional)}:
\begin{quote}\ttfamily\footnotesize\raggedright\noindent
forall(\allowbreak{}[rpf\_, rpw\_, rpx\_, rpl\_], is-left-limit-within(\allowbreak{}rpf\_, rpw\_, rpx\_, rpl\_) iff rpw\_ subset rr and rpf\_ in fun(\allowbreak{}rpw\_, rr) and rpx\_ in rr and rpl\_ in rr and forall(\allowbreak{}[rple\_], pos-rr(\allowbreak{}rple\_) implies forsome(\allowbreak{}[rpld\_], pos-rr(\allowbreak{}rpld\_) and forall(\allowbreak{}[rplt\_ in rpw\_], rpx\_ - rpld\_ < rplt\_ implies rplt\_ < rpx\_ implies abs(\allowbreak{}rpf\_(\allowbreak{}rplt\_) - rpl\_) < rple\_))))
\end{quote}
mentioned by 17 result(s)

\smallskip

\noindent\textbf{\texttt{is-limit-along(lmas\_, lmat\_, lmaa\_, lmaf\_, lmax\_, lmay\_)}}\quad\textit{predicate}\index{is-limit-along(lmas_, lmat_, lmaa_, lmaf_, lmax_, lmay_)@\texttt{is-limit-along(lmas\_, lmat\_, lmaa\_, lmaf\_, lmax\_, lmay\_)}}\\
reads: \$6 is the limit of \$4(a) as a tends to \$5 through \$3 \\
defined by \texttt{is-limit-along} \\
characterized by \texttt{is-limit-along}  \textit{(definitional)}:
\begin{quote}\ttfamily\footnotesize\raggedright\noindent
forall(\allowbreak{}[lmas\_, lmat\_, lmaa\_, lmaf\_, lmax\_, lmay\_], is-limit-along(\allowbreak{}lmas\_, lmat\_, lmaa\_, lmaf\_, lmax\_, lmay\_) iff is-metric-space(\allowbreak{}lmas\_) and is-metric-space(\allowbreak{}lmat\_) and lmay\_ in pts(\allowbreak{}lmat\_) and forall(\allowbreak{}[lmae\_], pos-rr(\allowbreak{}lmae\_) implies forsome(\allowbreak{}[lmad\_], pos-rr(\allowbreak{}lmad\_) and forall(\allowbreak{}[lmaz\_ in lmaa\_], (\allowbreak{}dist(\allowbreak{}lmas\_))(\allowbreak{}lmax\_, lmaz\_) <= lmad\_ implies (\allowbreak{}dist(\allowbreak{}lmat\_))(\allowbreak{}lmaf\_(\allowbreak{}lmaz\_), lmay\_) <= lmae\_))))
\end{quote}
mentioned by 7 result(s)

\smallskip

\noindent\textbf{\texttt{is-linear-functional(m, f)}}\quad\textit{predicate}\index{is-linear-functional(m, f)@\texttt{is-linear-functional(m, f)}}\\
reads: \$2 is a linear functional on \$1 \\
defined by \texttt{is-linear-functional} \\
characterized by \texttt{is-linear-functional}  \textit{(definitional)}:
\begin{quote}\ttfamily\footnotesize\raggedright\noindent
forall(\allowbreak{}[m, f], is-linear-functional(\allowbreak{}m, f) iff f in fun(\allowbreak{}vec(\allowbreak{}m), rr) and forall(\allowbreak{}[x\_ in vec(\allowbreak{}m), y\_ in vec(\allowbreak{}m)], f(\allowbreak{}(\allowbreak{}vadd(\allowbreak{}m))(\allowbreak{}x\_, y\_)) = f(\allowbreak{}x\_) + f(\allowbreak{}y\_)) and forall(\allowbreak{}[r\_ in rr, x\_ in vec(\allowbreak{}m)], f(\allowbreak{}(\allowbreak{}act(\allowbreak{}m))(\allowbreak{}r\_, x\_)) = r\_ * f(\allowbreak{}x\_)))
\end{quote}
mentioned by 39 result(s)

\smallskip

\noindent\textbf{\texttt{is-linear-functional-on(m, s, f)}}\quad\textit{predicate}\index{is-linear-functional-on(m, s, f)@\texttt{is-linear-functional-on(m, s, f)}}\\
reads: \$3 is a linear functional on the subspace \$2 of \$1 \\
defined by \texttt{is-linear-functional-on} \\
characterized by \texttt{is-linear-functional-on}  \textit{(definitional)}:
\begin{quote}\ttfamily\footnotesize\raggedright\noindent
forall(\allowbreak{}[m, s, f], is-linear-functional-on(\allowbreak{}m, s, f) iff f in fun(\allowbreak{}s, rr) and forall(\allowbreak{}[x\_ in s, y\_ in s], f(\allowbreak{}(\allowbreak{}vadd(\allowbreak{}m))(\allowbreak{}x\_, y\_)) = f(\allowbreak{}x\_) + f(\allowbreak{}y\_)) and forall(\allowbreak{}[r\_ in rr, x\_ in s], f(\allowbreak{}(\allowbreak{}act(\allowbreak{}m))(\allowbreak{}r\_, x\_)) = r\_ * f(\allowbreak{}x\_)))
\end{quote}
mentioned by 116 result(s)

\smallskip

\noindent\textbf{\texttt{is-linear-graph(ubv\_, ubg\_)}}\quad\textit{predicate}\index{is-linear-graph(ubv_, ubg_)@\texttt{is-linear-graph(ubv\_, ubg\_)}}\\
reads: \$2 is a linear graph in \$1 x \$1 \\
defined by \texttt{is-linear-graph} \\
characterized by \texttt{is-linear-graph}  \textit{(definitional)}:
\begin{quote}\ttfamily\footnotesize\raggedright\noindent
forall(\allowbreak{}[ubv\_, ubg\_], is-linear-graph(\allowbreak{}ubv\_, ubg\_) iff is-submodule(\allowbreak{}hsum(\allowbreak{}ubv\_, ubv\_), ubg\_))
\end{quote}
mentioned by 28 result(s)

\smallskip

\noindent\textbf{\texttt{is-linear-map(m1, m2, tt)}}\quad\textit{predicate}\index{is-linear-map(m1, m2, tt)@\texttt{is-linear-map(m1, m2, tt)}}\\
reads: \$3 is a linear map from \$1 to \$2 \\
defined by \texttt{is-linear-map} \\
characterized by \texttt{is-linear-map}  \textit{(definitional)}:
\begin{quote}\ttfamily\footnotesize\raggedright\noindent
forall(\allowbreak{}[m1, m2, tt], is-linear-map(\allowbreak{}m1, m2, tt) iff is-module(\allowbreak{}m1) and is-module(\allowbreak{}m2) and scal(\allowbreak{}m1) = scal(\allowbreak{}m2) and tt in fun(\allowbreak{}vec(\allowbreak{}m1), vec(\allowbreak{}m2)) and forall(\allowbreak{}[x in vec(\allowbreak{}m1), y in vec(\allowbreak{}m1)], tt(\allowbreak{}(\allowbreak{}vadd(\allowbreak{}m1))(\allowbreak{}x, y)) = (\allowbreak{}vadd(\allowbreak{}m2))(\allowbreak{}tt(\allowbreak{}x), tt(\allowbreak{}y))) and forall(\allowbreak{}[lam in carr(\allowbreak{}scal(\allowbreak{}m1)), x in vec(\allowbreak{}m1)], tt(\allowbreak{}(\allowbreak{}act(\allowbreak{}m1))(\allowbreak{}lam, x)) = (\allowbreak{}act(\allowbreak{}m2))(\allowbreak{}lam, tt(\allowbreak{}x))))
\end{quote}
mentioned by 60 result(s)

\smallskip

\noindent\textbf{\texttt{is-linear-op(bov\_, bow\_, bot\_)}}\quad\textit{predicate}\index{is-linear-op(bov_, bow_, bot_)@\texttt{is-linear-op(bov\_, bow\_, bot\_)}}\\
reads: \$3 is a linear operator from \$1 to \$2 \\
defined by \texttt{is-linear-op} \\
characterized by \texttt{is-linear-op}  \textit{(definitional)}:
\begin{quote}\ttfamily\footnotesize\raggedright\noindent
forall(\allowbreak{}[bov\_, bow\_, bot\_], is-linear-op(\allowbreak{}bov\_, bow\_, bot\_) iff bot\_ in fun(\allowbreak{}vec(\allowbreak{}bov\_), vec(\allowbreak{}bow\_)) and forall(\allowbreak{}[box\_ in vec(\allowbreak{}bov\_), boy\_ in vec(\allowbreak{}bov\_)], bot\_(\allowbreak{}(\allowbreak{}vadd(\allowbreak{}bov\_))(\allowbreak{}box\_, boy\_)) = (\allowbreak{}vadd(\allowbreak{}bow\_))(\allowbreak{}bot\_(\allowbreak{}box\_), bot\_(\allowbreak{}boy\_))) and forall(\allowbreak{}[boc\_ in cc, box\_ in vec(\allowbreak{}bov\_)], bot\_(\allowbreak{}(\allowbreak{}act(\allowbreak{}bov\_))(\allowbreak{}boc\_, box\_)) = (\allowbreak{}act(\allowbreak{}bow\_))(\allowbreak{}boc\_, bot\_(\allowbreak{}box\_))))
\end{quote}
mentioned by 9 result(s)

\smallskip

\noindent\textbf{\texttt{is-lipschitz-arrow(a, b, f)}}\quad\textit{predicate}\index{is-lipschitz-arrow(a, b, f)@\texttt{is-lipschitz-arrow(a, b, f)}}\\
reads: \$3 is a Lipschitz map from \$1 to \$2 \\
a predicate \\
characterized by \texttt{is-lipschitz-arrow-def}  \textit{(definitional)}:
\begin{quote}\ttfamily\footnotesize\raggedright\noindent
forall(\allowbreak{}[a, b, f], is-lipschitz-arrow(\allowbreak{}a, b, f) iff is-metric-space(\allowbreak{}a) and is-metric-space(\allowbreak{}b) and f in fun(\allowbreak{}pts(\allowbreak{}a), pts(\allowbreak{}b)) and forsome(\allowbreak{}[lpk\_ in rr], 0 <= lpk\_ and forall(\allowbreak{}[lpx\_ in pts(\allowbreak{}a), lpy\_ in pts(\allowbreak{}a)], (\allowbreak{}dist(\allowbreak{}b))(\allowbreak{}f(\allowbreak{}lpx\_), f(\allowbreak{}lpy\_)) <= lpk\_ * (\allowbreak{}dist(\allowbreak{}a))(\allowbreak{}lpx\_, lpy\_))))
\end{quote}
also: \texttt{hom-metric-space-is-lipschitz-arrow} \\
mentioned by 8 result(s)

\smallskip

\noindent\textbf{\texttt{is-local-hom(lhs1\_, lhg1\_, lhs2\_, lhg2\_, lhv\_, lhf\_)}}\quad\textit{predicate}\index{is-local-hom(lhs1_, lhg1_, lhs2_, lhg2_, lhv_, lhf_)@\texttt{is-local-hom(lhs1\_, lhg1\_, lhs2\_, lhg2\_, lhv\_, lhf\_)}}\\
reads: \$6 is a local homomorphism on \$5 from the group \$2 on \$1 to the group \$4 on \$3 \\
defined by \texttt{is-local-hom} \\
characterized by \texttt{is-local-hom}  \textit{(definitional)}:
\begin{quote}\ttfamily\footnotesize\raggedright\noindent
forall(\allowbreak{}[lhs1\_, lhg1\_, lhs2\_, lhg2\_, lhv\_, lhf\_], is-local-hom(\allowbreak{}lhs1\_, lhg1\_, lhs2\_, lhg2\_, lhv\_, lhf\_) iff lhv\_ in opens(\allowbreak{}lhs1\_) and iden(\allowbreak{}lhg1\_) in lhv\_ and lhf\_ in fun(\allowbreak{}lhv\_, carr(\allowbreak{}lhg2\_)) and forall(\allowbreak{}[lhx\_ in lhv\_, lho\_ in opens(\allowbreak{}lhs2\_)], lhf\_(\allowbreak{}lhx\_) in lho\_ implies forsome(\allowbreak{}[lhu\_ in opens(\allowbreak{}lhs1\_)], lhx\_ in lhu\_ and forall(\allowbreak{}[lhy\_ in lhu\_], lhy\_ in lhv\_ implies lhf\_(\allowbreak{}lhy\_) in lho\_))) and forsome(\allowbreak{}[lhw\_ in opens(\allowbreak{}lhs1\_)], iden(\allowbreak{}lhg1\_) in lhw\_ and forall(\allowbreak{}[lhx\_ in lhw\_], (\allowbreak{}inv(\allowbreak{}lhg1\_))(\allowbreak{}lhx\_) in lhw\_) and forall(\allowbreak{}[lhx\_ in lhw\_, lhy\_ in lhw\_], (\allowbreak{}opr(\allowbreak{}lhg1\_))(\allowbreak{}lhx\_, lhy\_) in lhv\_ and lhf\_(\allowbreak{}(\allowbreak{}opr(\allowbreak{}lhg1\_))(\allowbreak{}lhx\_, lhy\_)) = (\allowbreak{}opr(\allowbreak{}lhg2\_))(\allowbreak{}lhf\_(\allowbreak{}lhx\_), lhf\_(\allowbreak{}lhy\_)))))
\end{quote}
mentioned by 4 result(s)

\smallskip

\noindent\textbf{\texttt{is-local-max-at(f, d, lmt\_)}}\quad\textit{predicate}\index{is-local-max-at(f, d, lmt_)@\texttt{is-local-max-at(f, d, lmt\_)}}\\
reads: \$1 has a local maximum on \$2 at \$3 \\
defined by \texttt{is-local-max-at} \\
characterized by \texttt{is-local-max-at}  \textit{(definitional)}:
\begin{quote}\ttfamily\footnotesize\raggedright\noindent
forall(\allowbreak{}[f, d, lmt\_], is-local-max-at(\allowbreak{}f, d, lmt\_) iff f in fun(\allowbreak{}d, rr) and lmt\_ in d and forsome(\allowbreak{}[lma\_], pos-rr(\allowbreak{}lma\_) and forall(\allowbreak{}[lmx\_ in d], abs(\allowbreak{}lmx\_ - lmt\_) <= lma\_ implies f(\allowbreak{}lmx\_) <= f(\allowbreak{}lmt\_))))
\end{quote}
mentioned by 8 result(s)

\smallskip

\noindent\textbf{\texttt{is-local-measure(omega, ca, cb, rho)}}\quad\textit{predicate}\index{is-local-measure(omega, ca, cb, rho)@\texttt{is-local-measure(omega, ca, cb, rho)}}\\
reads: \$4 is a local countably additive measure on the hereditary subset \$3 of the sigma-algebra \$2 on \$1 \\
defined by \texttt{is-local-measure} \\
characterized by \texttt{is-local-measure}  \textit{(definitional)}:
\begin{quote}\ttfamily\footnotesize\raggedright\noindent
forall(\allowbreak{}[omega, ca, cb, rho], is-local-measure(\allowbreak{}omega, ca, cb, rho) iff is-sigma-algebra(\allowbreak{}omega, ca) and is-hereditary(\allowbreak{}ca, cb) and rho in fun(\allowbreak{}cb, rr-pos-star) and forall(\allowbreak{}[lme\_ in cb], is-measure(\allowbreak{}lme\_, restrict-sigma(\allowbreak{}ca, lme\_), restrict(\allowbreak{}rho, restrict-sigma(\allowbreak{}ca, lme\_)))))
\end{quote}
mentioned by 27 result(s)

\smallskip

\noindent\textbf{\texttt{is-local-min-at(f, d, lmt\_)}}\quad\textit{predicate}\index{is-local-min-at(f, d, lmt_)@\texttt{is-local-min-at(f, d, lmt\_)}}\\
reads: \$1 has a local minimum on \$2 at \$3 \\
defined by \texttt{is-local-min-at} \\
characterized by \texttt{is-local-min-at}  \textit{(definitional)}:
\begin{quote}\ttfamily\footnotesize\raggedright\noindent
forall(\allowbreak{}[f, d, lmt\_], is-local-min-at(\allowbreak{}f, d, lmt\_) iff f in fun(\allowbreak{}d, rr) and lmt\_ in d and forsome(\allowbreak{}[lma\_], pos-rr(\allowbreak{}lma\_) and forall(\allowbreak{}[lmx\_ in d], abs(\allowbreak{}lmx\_ - lmt\_) <= lma\_ implies f(\allowbreak{}lmt\_) <= f(\allowbreak{}lmx\_))))
\end{quote}
mentioned by 4 result(s)

\smallskip

\noindent\textbf{\texttt{is-local-nv-field(lck\_)}}\quad\textit{predicate}\index{is-local-nv-field(lck_)@\texttt{is-local-nv-field(lck\_)}}\\
reads: \$1 is a local NV field \\
defined by \texttt{is-local-nv-field} \\
characterized by \texttt{is-local-nv-field}  \textit{(definitional)}:
\begin{quote}\ttfamily\footnotesize\raggedright\noindent
forall(\allowbreak{}[lck\_], is-local-nv-field(\allowbreak{}lck\_) iff is-nv-field(\allowbreak{}lck\_) and not(\allowbreak{}is-discrete-ms(\allowbreak{}nf-metric-space(\allowbreak{}lck\_))) and is-locally-compact(\allowbreak{}nf-metric-space(\allowbreak{}lck\_)))
\end{quote}
mentioned by 25 result(s)

\smallskip

\noindent\textbf{\texttt{is-local-uniformizer(c5k\_, c5u\_)}}\quad\textit{predicate}\index{is-local-uniformizer(c5k_, c5u_)@\texttt{is-local-uniformizer(c5k\_, c5u\_)}}\\
reads: \$1 has compact integers and \$2 is a uniformizing parameter of \$1 \\
defined by \texttt{is-local-uniformizer} \\
characterized by \texttt{is-local-uniformizer}  \textit{(definitional)}:
\begin{quote}\ttfamily\footnotesize\raggedright\noindent
forall(\allowbreak{}[c5k\_, c5u\_], is-local-uniformizer(\allowbreak{}c5k\_, c5u\_) iff is-compact-nv-domain(\allowbreak{}c5k\_, nv-integers(\allowbreak{}c5k\_)) and c5u\_ in nv-integers(\allowbreak{}c5k\_) and nv-maximal(\allowbreak{}c5k\_, nv-integers(\allowbreak{}c5k\_)) = nv-principal(\allowbreak{}c5k\_, nv-integers(\allowbreak{}c5k\_), c5u\_) and not(\allowbreak{}c5u\_ = zero(\allowbreak{}c5k\_)))
\end{quote}
mentioned by 38 result(s)

\smallskip

\noindent\textbf{\texttt{is-locally-compact(lcs\_)}}\quad\textit{predicate}\index{is-locally-compact(lcs_)@\texttt{is-locally-compact(lcs\_)}}\\
reads: \$1 is locally compact \\
defined by \texttt{is-locally-compact} \\
characterized by \texttt{is-locally-compact}  \textit{(definitional)}:
\begin{quote}\ttfamily\footnotesize\raggedright\noindent
forall(\allowbreak{}[lcs\_], is-locally-compact(\allowbreak{}lcs\_) iff is-metric-space(\allowbreak{}lcs\_) and forall(\allowbreak{}[lca\_ in pts(\allowbreak{}lcs\_)], forsome(\allowbreak{}[lcr\_], pos-rr(\allowbreak{}lcr\_) and is-compact(\allowbreak{}subspace-ms(\allowbreak{}lcs\_, closed-ball(\allowbreak{}lcs\_, lca\_, lcr\_))))))
\end{quote}
mentioned by 9 result(s)

\smallskip

\noindent\textbf{\texttt{is-locally-constant-fn(c5k\_, c5f\_)}}\quad\textit{predicate}\index{is-locally-constant-fn(c5k_, c5f_)@\texttt{is-locally-constant-fn(c5k\_, c5f\_)}}\\
reads: \$2 is locally constant on \$1 \\
defined by \texttt{is-locally-constant-fn} \\
characterized by \texttt{is-locally-constant-fn}  \textit{(definitional)}:
\begin{quote}\ttfamily\footnotesize\raggedright\noindent
forall(\allowbreak{}[c5k\_, c5f\_], is-locally-constant-fn(\allowbreak{}c5k\_, c5f\_) iff forall(\allowbreak{}[c5x\_ in carr(\allowbreak{}c5k\_)], forsome(\allowbreak{}[c5r\_], pos-rr(\allowbreak{}c5r\_) and forall(\allowbreak{}[c5y\_ in carr(\allowbreak{}c5k\_)], (\allowbreak{}fnrm(\allowbreak{}c5k\_))(\allowbreak{}(\allowbreak{}add(\allowbreak{}c5k\_))(\allowbreak{}c5y\_, (\allowbreak{}neg(\allowbreak{}c5k\_))(\allowbreak{}c5x\_))) <= c5r\_ implies c5f\_(\allowbreak{}c5y\_) = c5f\_(\allowbreak{}c5x\_)))))
\end{quote}
mentioned by 4 result(s)

\smallskip

\noindent\textbf{\texttt{is-locally-convex-tvs(m, t)}}\quad\textit{predicate}\index{is-locally-convex-tvs(m, t)@\texttt{is-locally-convex-tvs(m, t)}}\\
reads: \$2 is a locally convex vector space topology on the real vector space \$1 \\
defined by \texttt{is-locally-convex-tvs} \\
characterized by \texttt{is-locally-convex-tvs}  \textit{(definitional)}:
\begin{quote}\ttfamily\footnotesize\raggedright\noindent
forall(\allowbreak{}[m, t], is-locally-convex-tvs(\allowbreak{}m, t) iff is-tvs(\allowbreak{}rr-normed-field, m, t) and forall(\allowbreak{}[lcx\_ in vec(\allowbreak{}m), lcu\_ in opens(\allowbreak{}t)], lcx\_ in lcu\_ implies forsome(\allowbreak{}[lcv\_ in opens(\allowbreak{}t)], lcx\_ in lcv\_ and lcv\_ subset lcu\_ and is-convex-subset(\allowbreak{}m, lcv\_))))
\end{quote}
mentioned by 9 result(s)

\smallskip

\noindent\textbf{\texttt{is-locally-null(ca, cb, rho, es)}}\quad\textit{predicate}\index{is-locally-null(ca, cb, rho, es)@\texttt{is-locally-null(ca, cb, rho, es)}}\\
reads: \$4 is locally null relative to the local measure \$3 on \$2 in \$1 \\
defined by \texttt{is-locally-null} \\
characterized by \texttt{is-locally-null}  \textit{(definitional)}:
\begin{quote}\ttfamily\footnotesize\raggedright\noindent
forall(\allowbreak{}[ca, cb, rho, es], is-locally-null(\allowbreak{}ca, cb, rho, es) iff es in ca and forall(\allowbreak{}[lna\_ in cb], rho(\allowbreak{}intersection(\allowbreak{}lna\_, es)) = 0))
\end{quote}
mentioned by 12 result(s)

\smallskip

\noindent\textbf{\texttt{is-lower-triangular(a, n, d)}}\quad\textit{predicate}\index{is-lower-triangular(a, n, d)@\texttt{is-lower-triangular(a, n, d)}}\\
reads: \$3 is a lower triangular \$2-by-\$2 matrix over \$1 \\
defined by \texttt{is-lower-triangular} \\
characterized by \texttt{is-lower-triangular}  \textit{(definitional)}:
\begin{quote}\ttfamily\footnotesize\raggedright\noindent
forall(\allowbreak{}[a, n, d], is-lower-triangular(\allowbreak{}a, n, d) iff forall(\allowbreak{}[ltu\_, ltv\_], ltu\_ in interval(\allowbreak{}1, n) implies ltv\_ in interval(\allowbreak{}1, n) implies ltu\_ < ltv\_ implies entry(\allowbreak{}d, ltu\_, ltv\_) = zero(\allowbreak{}a)))
\end{quote}
mentioned by 9 result(s)

\smallskip

\noindent\textbf{\texttt{is-lusin(s, e)}}\quad\textit{predicate}\index{is-lusin(s, e)@\texttt{is-lusin(s, e)}}\\
reads: \$2 is a Lusin subset of \$1 \\
defined by \texttt{is-lusin} \\
characterized by \texttt{is-lusin}  \textit{(definitional)}:
\begin{quote}\ttfamily\footnotesize\raggedright\noindent
forall(\allowbreak{}[s, e], is-lusin(\allowbreak{}s, e) iff forsome(\allowbreak{}[lzt\_], is-complete(\allowbreak{}lzt\_) and is-separable(\allowbreak{}lzt\_) and forsome(\allowbreak{}[lzh\_], is-continuous(\allowbreak{}lzt\_, subspace-ms(\allowbreak{}s, e), lzh\_) and image(\allowbreak{}lzh\_, pts(\allowbreak{}lzt\_)) = e and forall(\allowbreak{}[lzx\_ in pts(\allowbreak{}lzt\_), lzy\_ in pts(\allowbreak{}lzt\_)], lzh\_(\allowbreak{}lzx\_) = lzh\_(\allowbreak{}lzy\_) implies lzx\_ = lzy\_))))
\end{quote}
mentioned by 11 result(s)

\smallskip

\noindent\textbf{\texttt{is-mat-bound(k, m, n, a, c)}}\quad\textit{predicate}\index{is-mat-bound(k, m, n, a, c)@\texttt{is-mat-bound(k, m, n, a, c)}}\\
reads: \$4 is a \$2-by-\$3 matrix over \$1 with max norm at most \$5 \\
defined by \texttt{is-mat-bound} \\
characterized by \texttt{is-mat-bound}  \textit{(definitional)}:
\begin{quote}\ttfamily\footnotesize\raggedright\noindent
forall(\allowbreak{}[k, m, n, a, c], is-mat-bound(\allowbreak{}k, m, n, a, c) iff a in mat(\allowbreak{}m, n, carr(\allowbreak{}k)) and forall(\allowbreak{}[cbvi\_ in interval(\allowbreak{}1, m), cbvj\_ in interval(\allowbreak{}1, n)], (\allowbreak{}fnrm(\allowbreak{}k))(\allowbreak{}entry(\allowbreak{}a, cbvi\_, cbvj\_)) <= c))
\end{quote}
mentioned by 33 result(s)

\smallskip

\noindent\textbf{\texttt{is-mat-close(k, m, n, a, b, e)}}\quad\textit{predicate}\index{is-mat-close(k, m, n, a, b, e)@\texttt{is-mat-close(k, m, n, a, b, e)}}\\
reads: \$4 and \$5 are \$2-by-\$3 matrices over \$1 with max norm of the difference at most \$6 \\
defined by \texttt{is-mat-close} \\
characterized by \texttt{is-mat-close}  \textit{(definitional)}:
\begin{quote}\ttfamily\footnotesize\raggedright\noindent
forall(\allowbreak{}[k, m, n, a, b, e], is-mat-close(\allowbreak{}k, m, n, a, b, e) iff a in mat(\allowbreak{}m, n, carr(\allowbreak{}k)) and b in mat(\allowbreak{}m, n, carr(\allowbreak{}k)) and forall(\allowbreak{}[cdvi\_ in interval(\allowbreak{}1, m), cdvj\_ in interval(\allowbreak{}1, n)], (\allowbreak{}fnrm(\allowbreak{}k))(\allowbreak{}(\allowbreak{}add(\allowbreak{}k))(\allowbreak{}entry(\allowbreak{}a, cdvi\_, cdvj\_), (\allowbreak{}neg(\allowbreak{}k))(\allowbreak{}entry(\allowbreak{}b, cdvi\_, cdvj\_)))) <= e))
\end{quote}
mentioned by 8 result(s)

\smallskip

\noindent\textbf{\texttt{is-matrix-of(gmx-md\_, gmx-n\_, gmx-u\_, gmx-t\_, gmx-a\_)}}\quad\textit{predicate}\index{is-matrix-of(gmx-md_, gmx-n_, gmx-u_, gmx-t_, gmx-a_)@\texttt{is-matrix-of(gmx-md\_, gmx-n\_, gmx-u\_, gmx-t\_, gmx-a\_)}}\\
reads: \$5 is the matrix of \$4 in the basis \$3 \\
defined by \texttt{is-matrix-of} \\
characterized by \texttt{is-matrix-of}  \textit{(definitional)}:
\begin{quote}\ttfamily\footnotesize\raggedright\noindent
forall(\allowbreak{}[gmx-md\_, gmx-n\_, gmx-u\_, gmx-t\_, gmx-a\_], is-matrix-of(\allowbreak{}gmx-md\_, gmx-n\_, gmx-u\_, gmx-t\_, gmx-a\_) iff gmx-a\_ in mat(\allowbreak{}gmx-n\_, gmx-n\_, carr(\allowbreak{}scal(\allowbreak{}gmx-md\_))) and forall(\allowbreak{}[gmxj\_ in interval(\allowbreak{}1, gmx-n\_)], gmx-t\_(\allowbreak{}entry(\allowbreak{}gmx-u\_, gmxj\_, 1)) = lincomb(\allowbreak{}gmx-md\_, gmx-n\_, col-as-row(\allowbreak{}gmx-a\_, gmx-n\_, gmxj\_), gmx-u\_)))
\end{quote}
mentioned by 46 result(s)

\smallskip

\noindent\textbf{\texttt{is-maximal(grd, porel, mx)}}\quad\textit{predicate}\index{is-maximal(grd, porel, mx)@\texttt{is-maximal(grd, porel, mx)}}\\
reads: \$3 is maximal in \$1 under \$2 \\
defined by \texttt{is-maximal} \\
characterized by \texttt{is-maximal}  \textit{(definitional)}:
\begin{quote}\ttfamily\footnotesize\raggedright\noindent
forall(\allowbreak{}[grd, porel, mx], is-maximal(\allowbreak{}grd, porel, mx) iff mx in grd and forall(\allowbreak{}[y in grd], [mx, y] in porel implies [y, mx] in porel))
\end{quote}
mentioned by 4 result(s)

\smallskip

\noindent\textbf{\texttt{is-maximal-ideal(baa\_, bai\_)}}\quad\textit{predicate}\index{is-maximal-ideal(baa_, bai_)@\texttt{is-maximal-ideal(baa\_, bai\_)}}\\
reads: \$2 is a maximal proper ideal of \$1 \\
defined by \texttt{is-maximal-ideal} \\
characterized by \texttt{is-maximal-ideal}  \textit{(definitional)}:
\begin{quote}\ttfamily\footnotesize\raggedright\noindent
forall(\allowbreak{}[baa\_, bai\_], is-maximal-ideal(\allowbreak{}baa\_, bai\_) iff is-two-sided-ideal(\allowbreak{}baa\_, bai\_) and not(\allowbreak{}bai\_ = vec(\allowbreak{}baa\_)) and forall(\allowbreak{}[baj\_], is-two-sided-ideal(\allowbreak{}baa\_, baj\_) implies bai\_ subset baj\_ implies not(\allowbreak{}baj\_ = vec(\allowbreak{}baa\_)) implies baj\_ = bai\_))
\end{quote}
mentioned by 22 result(s)

\smallskip

\noindent\textbf{\texttt{is-maximal-left-ideal(baa\_, bai\_)}}\quad\textit{predicate}\index{is-maximal-left-ideal(baa_, bai_)@\texttt{is-maximal-left-ideal(baa\_, bai\_)}}\\
reads: \$2 is a maximal proper left ideal of \$1 \\
defined by \texttt{is-maximal-left-ideal} \\
characterized by \texttt{is-maximal-left-ideal}  \textit{(definitional)}:
\begin{quote}\ttfamily\footnotesize\raggedright\noindent
forall(\allowbreak{}[baa\_, bai\_], is-maximal-left-ideal(\allowbreak{}baa\_, bai\_) iff is-left-ideal(\allowbreak{}baa\_, bai\_) and not(\allowbreak{}bai\_ = vec(\allowbreak{}baa\_)) and forall(\allowbreak{}[baj\_], is-left-ideal(\allowbreak{}baa\_, baj\_) implies bai\_ subset baj\_ implies not(\allowbreak{}baj\_ = vec(\allowbreak{}baa\_)) implies baj\_ = bai\_))
\end{quote}
mentioned by 4 result(s)

\smallskip

\noindent\textbf{\texttt{is-maximal-right-ideal(baa\_, bai\_)}}\quad\textit{predicate}\index{is-maximal-right-ideal(baa_, bai_)@\texttt{is-maximal-right-ideal(baa\_, bai\_)}}\\
reads: \$2 is a maximal proper right ideal of \$1 \\
defined by \texttt{is-maximal-right-ideal} \\
characterized by \texttt{is-maximal-right-ideal}  \textit{(definitional)}:
\begin{quote}\ttfamily\footnotesize\raggedright\noindent
forall(\allowbreak{}[baa\_, bai\_], is-maximal-right-ideal(\allowbreak{}baa\_, bai\_) iff is-right-ideal(\allowbreak{}baa\_, bai\_) and not(\allowbreak{}bai\_ = vec(\allowbreak{}baa\_)) and forall(\allowbreak{}[baj\_], is-right-ideal(\allowbreak{}baa\_, baj\_) implies bai\_ subset baj\_ implies not(\allowbreak{}baj\_ = vec(\allowbreak{}baa\_)) implies baj\_ = bai\_))
\end{quote}
mentioned by 2 result(s)

\smallskip

\noindent\textbf{\texttt{is-meager(s, a)}}\quad\textit{predicate}\index{is-meager(s, a)@\texttt{is-meager(s, a)}}\\
reads: \$2 is meager in \$1 \\
defined by \texttt{is-meager} \\
characterized by \texttt{is-meager}  \textit{(definitional)}:
\begin{quote}\ttfamily\footnotesize\raggedright\noindent
forall(\allowbreak{}[s, a], is-meager(\allowbreak{}s, a) iff is-metric-space(\allowbreak{}s) and a subset pts(\allowbreak{}s) and forsome(\allowbreak{}[ee in fun(\allowbreak{}nn, power(\allowbreak{}pts(\allowbreak{}s)))], forall(\allowbreak{}[n in nn], is-nowhere-dense(\allowbreak{}s, ee(\allowbreak{}n))) and a subset big-union(\allowbreak{}n, nn, ee(\allowbreak{}n))))
\end{quote}
mentioned by 4 result(s)

\smallskip

\noindent\textbf{\texttt{is-measurable-fn(omega, ca, f)}}\quad\textit{predicate}\index{is-measurable-fn(omega, ca, f)@\texttt{is-measurable-fn(omega, ca, f)}}\\
reads: \$3 is a measurable [0,+inf]-valued function on \$1 with \$2 \\
defined by \texttt{is-measurable-fn} \\
characterized by \texttt{is-measurable-fn}  \textit{(definitional)}:
\begin{quote}\ttfamily\footnotesize\raggedright\noindent
forall(\allowbreak{}[omega, ca, f], is-measurable-fn(\allowbreak{}omega, ca, f) iff is-sigma-algebra(\allowbreak{}omega, ca) and f in fun(\allowbreak{}omega, rr-pos-star) and forall(\allowbreak{}[alpha\_ in rr], \{x\_ in omega: alpha\_ < f(\allowbreak{}x\_)\} in ca))
\end{quote}
mentioned by 157 result(s)

\smallskip

\noindent\textbf{\texttt{is-measurable-map(omega, ca, omega2, cb, f)}}\quad\textit{predicate}\index{is-measurable-map(omega, ca, omega2, cb, f)@\texttt{is-measurable-map(omega, ca, omega2, cb, f)}}\\
reads: \$5 is a measurable map from \$1 with \$2 to \$3 with \$4 \\
defined by \texttt{is-measurable-map} \\
characterized by \texttt{is-measurable-map}  \textit{(definitional)}:
\begin{quote}\ttfamily\footnotesize\raggedright\noindent
forall(\allowbreak{}[omega, ca, omega2, cb, f], is-measurable-map(\allowbreak{}omega, ca, omega2, cb, f) iff is-sigma-algebra(\allowbreak{}omega, ca) and is-sigma-algebra(\allowbreak{}omega2, cb) and f in fun(\allowbreak{}omega, omega2) and forall(\allowbreak{}[b\_ in cb], \{x\_ in omega: f(\allowbreak{}x\_) in b\_\} in ca))
\end{quote}
mentioned by 53 result(s)

\smallskip

\noindent\textbf{\texttt{is-measurable-space(s)}}\quad\textit{predicate}\index{is-measurable-space(s)@\texttt{is-measurable-space(s)}}\\
reads: a measurable space \\
a predicate \\
characterized by \texttt{is-measurable-space}  \textit{(definitional)}:
\begin{quote}\ttfamily\footnotesize\raggedright\noindent
forall(\allowbreak{}[s], is-measurable-space(\allowbreak{}s) iff length(\allowbreak{}s) = 2 and pts(\allowbreak{}s) in set and sigma(\allowbreak{}s) in power(\allowbreak{}power(\allowbreak{}pts(\allowbreak{}s))) and is-sigma-algebra(\allowbreak{}pts(\allowbreak{}s), sigma(\allowbreak{}s)) and s in set)
\end{quote}
also: \texttt{measure-space-as-measurable-space-is-measurable-space} \texttt{is-measurable-space-is-set} \texttt{discrete-space-is-measurable-space} \\
mentioned by 13 result(s)

\smallskip

\noindent\textbf{\texttt{is-measurable-step-fn(omega, ca, s)}}\quad\textit{predicate}\index{is-measurable-step-fn(omega, ca, s)@\texttt{is-measurable-step-fn(omega, ca, s)}}\\
reads: \$3 is a measurable real step function on \$1 with \$2 \\
defined by \texttt{is-measurable-step-fn} \\
characterized by \texttt{is-measurable-step-fn}  \textit{(definitional)}:
\begin{quote}\ttfamily\footnotesize\raggedright\noindent
forall(\allowbreak{}[omega, ca, s], is-measurable-step-fn(\allowbreak{}omega, ca, s) iff is-real-measurable-fn(\allowbreak{}omega, ca, s) and card(\allowbreak{}image(\allowbreak{}s, omega)) in nn)
\end{quote}
mentioned by 4 result(s)

\smallskip

\noindent\textbf{\texttt{is-measure(omega, ca, mu)}}\quad\textit{predicate}\index{is-measure(omega, ca, mu)@\texttt{is-measure(omega, ca, mu)}}\\
reads: \$3 is a countably additive measure on the sigma-algebra \$2 of subsets of \$1 \\
defined by \texttt{is-measure} \\
characterized by \texttt{is-measure}  \textit{(definitional)}:
\begin{quote}\ttfamily\footnotesize\raggedright\noindent
forall(\allowbreak{}[omega, ca, mu], is-measure(\allowbreak{}omega, ca, mu) iff is-sigma-algebra(\allowbreak{}omega, ca) and mu in fun(\allowbreak{}ca, rr-pos-star) and mu(\allowbreak{}empty-set) = 0 and forall(\allowbreak{}[f\_ in fun(\allowbreak{}nn, ca)], forall(\allowbreak{}[i\_, j\_], i\_ in nn implies j\_ in nn implies not(\allowbreak{}i\_ = j\_) implies intersection(\allowbreak{}f\_(\allowbreak{}i\_), f\_(\allowbreak{}j\_)) = empty-set) implies mu(\allowbreak{}big-union(\allowbreak{}n\_, nn, f\_(\allowbreak{}n\_))) = esum(\allowbreak{}vnb-lambda(\allowbreak{}n\_, nn, mu(\allowbreak{}f\_(\allowbreak{}n\_))))))
\end{quote}
also: \texttt{is-measure-space} \texttt{neg-variation-is-measure} \texttt{pos-variation-is-measure} \texttt{product-measure-is-measure} \texttt{total-variation-is-measure} \texttt{pullback-measure-is-measure} \texttt{completion-measure-is-measure} \\
mentioned by 490 result(s)

\smallskip

\noindent\textbf{\texttt{is-measure-isomorphism(omega, ca, mu, omega2, cb, nu, tt)}}\quad\textit{predicate}\index{is-measure-isomorphism(omega, ca, mu, omega2, cb, nu, tt)@\texttt{is-measure-isomorphism(omega, ca, mu, omega2, cb, nu, tt)}}\\
reads: \$7 is a measure isomorphism from (\$1, \$3) to (\$4, \$6) \\
defined by \texttt{is-measure-isomorphism} \\
characterized by \texttt{is-measure-isomorphism}  \textit{(definitional)}:
\begin{quote}\ttfamily\footnotesize\raggedright\noindent
forall(\allowbreak{}[omega, ca, mu, omega2, cb, nu, tt], is-measure-isomorphism(\allowbreak{}omega, ca, mu, omega2, cb, nu, tt) iff is-measurable-map(\allowbreak{}omega, ca, omega2, cb, tt) and forall(\allowbreak{}[mib\_ in cb], nu(\allowbreak{}mib\_) = mu(\allowbreak{}\{mix\_ in omega: tt(\allowbreak{}mix\_) in mib\_\})) and forsome(\allowbreak{}[mis\_], is-measurable-map(\allowbreak{}omega2, cb, omega, ca, mis\_) and forall(\allowbreak{}[mix\_ in omega], mis\_(\allowbreak{}tt(\allowbreak{}mix\_)) = mix\_) and forall(\allowbreak{}[miy\_ in omega2], tt(\allowbreak{}mis\_(\allowbreak{}miy\_)) = miy\_)))
\end{quote}
mentioned by 3 result(s)

\smallskip

\noindent\textbf{\texttt{is-measure-space(s)}}\quad\textit{predicate}\index{is-measure-space(s)@\texttt{is-measure-space(s)}}\\
reads: a measure space \\
a predicate \\
characterized by \texttt{is-measure-space}  \textit{(definitional)}:
\begin{quote}\ttfamily\footnotesize\raggedright\noindent
forall(\allowbreak{}[s], is-measure-space(\allowbreak{}s) iff length(\allowbreak{}s) = 3 and pts(\allowbreak{}s) in set and sigma(\allowbreak{}s) in power(\allowbreak{}power(\allowbreak{}pts(\allowbreak{}s))) and meas(\allowbreak{}s) in fun(\allowbreak{}sigma(\allowbreak{}s), rr-pos-star) and is-measure(\allowbreak{}pts(\allowbreak{}s), sigma(\allowbreak{}s), meas(\allowbreak{}s)) and s in set)
\end{quote}
also: \texttt{is-measure-space-is-set} \\
mentioned by 12 result(s)

\smallskip

\noindent\textbf{\texttt{is-meromorphic-at(mrf\_, mra\_)}}\quad\textit{predicate}\index{is-meromorphic-at(mrf_, mra_)@\texttt{is-meromorphic-at(mrf\_, mra\_)}}\\
reads: \$1 is meromorphic at \$2 \\
defined by \texttt{is-meromorphic-at} \\
characterized by \texttt{is-meromorphic-at}  \textit{(definitional)}:
\begin{quote}\ttfamily\footnotesize\raggedright\noindent
forall(\allowbreak{}[mrf\_, mra\_], is-meromorphic-at(\allowbreak{}mrf\_, mra\_) iff mra\_ in cc and forsome(\allowbreak{}[mrn\_ in nn, mrr\_ in rr], 0 < mrr\_ and forsome(\allowbreak{}[mrg\_], holomorphic-on(\allowbreak{}ball(\allowbreak{}cc-ms, mra\_, mrr\_), mrg\_) and forall(\allowbreak{}[mrz\_ in difference(\allowbreak{}ball(\allowbreak{}cc-ms, mra\_, mrr\_), singleton(\allowbreak{}mra\_))], mrf\_(\allowbreak{}mrz\_) = recip(\allowbreak{}mrz\_ - mra\_) \textasciicircum{} mrn\_ * mrg\_(\allowbreak{}mrz\_)))))
\end{quote}
also: \texttt{pole-is-meromorphic-at} \texttt{removable-is-meromorphic-at} \texttt{holomorphic-is-meromorphic-at} \\
mentioned by 13 result(s)

\smallskip

\noindent\textbf{\texttt{is-metric(dst, crr)}}\quad\textit{predicate}\index{is-metric(dst, crr)@\texttt{is-metric(dst, crr)}}\\
reads: \$1 is a metric on \$2 \\
defined by \texttt{is-metric} \\
characterized by \texttt{is-metric}  \textit{(definitional)}:
\begin{quote}\ttfamily\footnotesize\raggedright\noindent
forall(\allowbreak{}[dst, crr], is-metric(\allowbreak{}dst, crr) iff forall(\allowbreak{}[u in crr], dst(\allowbreak{}u, u) = 0 and forall(\allowbreak{}[v in crr], 0 <= dst(\allowbreak{}u, v) and (\allowbreak{}dst(\allowbreak{}u, v) = 0 implies u = v) and dst(\allowbreak{}u, v) = dst(\allowbreak{}v, u) and forall(\allowbreak{}[w in crr], dst(\allowbreak{}u, w) <= dst(\allowbreak{}u, v) + dst(\allowbreak{}v, w)))))
\end{quote}
also: \texttt{is-metric-group} \texttt{is-metric-space} \\
mentioned by 6 result(s)

\smallskip

\noindent\textbf{\texttt{is-metric-group(s)}}\quad\textit{predicate}\index{is-metric-group(s)@\texttt{is-metric-group(s)}}\\
reads: a metric group \\
a predicate \\
characterized by \texttt{is-metric-group}  \textit{(definitional)}:
\begin{quote}\ttfamily\footnotesize\raggedright\noindent
forall(\allowbreak{}[s], is-metric-group(\allowbreak{}s) iff length(\allowbreak{}s) = 5 and carr(\allowbreak{}s) in set and opr(\allowbreak{}s) in fun(\allowbreak{}cartesian(\allowbreak{}carr(\allowbreak{}s), carr(\allowbreak{}s)), carr(\allowbreak{}s)) and iden(\allowbreak{}s) in carr(\allowbreak{}s) and inv(\allowbreak{}s) in fun(\allowbreak{}carr(\allowbreak{}s), carr(\allowbreak{}s)) and gdist(\allowbreak{}s) in fun(\allowbreak{}cartesian(\allowbreak{}carr(\allowbreak{}s), carr(\allowbreak{}s)), rr) and is-associative(\allowbreak{}opr(\allowbreak{}s), carr(\allowbreak{}s)) and is-identity(\allowbreak{}opr(\allowbreak{}s), iden(\allowbreak{}s), carr(\allowbreak{}s)) and has-inverses(\allowbreak{}opr(\allowbreak{}s), iden(\allowbreak{}s), inv(\allowbreak{}s), carr(\allowbreak{}s)) and is-metric(\allowbreak{}gdist(\allowbreak{}s), carr(\allowbreak{}s)) and forall(\allowbreak{}[mgu\_ in carr(\allowbreak{}s), mgv\_ in carr(\allowbreak{}s), mge\_], pos-rr(\allowbreak{}mge\_) implies forsome(\allowbreak{}[mgd\_], pos-rr(\allowbreak{}mgd\_) and forall(\allowbreak{}[mgp\_ in carr(\allowbreak{}s), mgq\_ in carr(\allowbreak{}s)], (\allowbreak{}gdist(\allowbreak{}s))(\allowbreak{}mgu\_, mgp\_) <= mgd\_ implies (\allowbreak{}gdist(\allowbreak{}s))(\allowbreak{}mgv\_, mgq\_) <= mgd\_ implies (\allowbreak{}gdist(\allowbreak{}s))(\allowbreak{}(\allowbreak{}opr(\allowbreak{}s))(\allowbreak{}mgu\_, mgv\_), (\allowbreak{}opr(\allowbreak{}s))(\allowbreak{}mgp\_, mgq\_)) <= mge\_))) and forall(\allowbreak{}[mga\_ in carr(\allowbreak{}s), mge\_], pos-rr(\allowbreak{}mge\_) implies forsome(\allowbreak{}[mgd\_], pos-rr(\allowbreak{}mgd\_) and forall(\allowbreak{}[mgb\_ in carr(\allowbreak{}s)], (\allowbreak{}gdist(\allowbreak{}s))(\allowbreak{}mga\_, mgb\_) <= mgd\_ implies (\allowbreak{}gdist(\allowbreak{}s))(\allowbreak{}(\allowbreak{}inv(\allowbreak{}s))(\allowbreak{}mga\_), (\allowbreak{}inv(\allowbreak{}s))(\allowbreak{}mgb\_)) <= mge\_))) and s in set)
\end{quote}
also: \texttt{is-metric-group-is-set} \texttt{opp-mg-is-metric-group} \texttt{isometry-group-is-metric-group} \texttt{make-metric-group-is-metric-group} \\
mentioned by 179 result(s)

\smallskip

\noindent\textbf{\texttt{is-metric-outer-measure(m, mu)}}\quad\textit{predicate}\index{is-metric-outer-measure(m, mu)@\texttt{is-metric-outer-measure(m, mu)}}\\
reads: \$2 is a metric outer measure on the metric space \$1 \\
defined by \texttt{is-metric-outer-measure} \\
characterized by \texttt{is-metric-outer-measure}  \textit{(definitional)}:
\begin{quote}\ttfamily\footnotesize\raggedright\noindent
forall(\allowbreak{}[m, mu], is-metric-outer-measure(\allowbreak{}m, mu) iff is-metric-space(\allowbreak{}m) and is-outer-measure(\allowbreak{}pts(\allowbreak{}m), mu) and forall(\allowbreak{}[moa\_ in power(\allowbreak{}pts(\allowbreak{}m)), mob\_ in power(\allowbreak{}pts(\allowbreak{}m))], not(\allowbreak{}set-dist(\allowbreak{}m, moa\_, mob\_) = 0) implies mu(\allowbreak{}union(\allowbreak{}moa\_, mob\_)) = eplus(\allowbreak{}mu(\allowbreak{}moa\_), mu(\allowbreak{}mob\_))))
\end{quote}
mentioned by 17 result(s)

\smallskip

\noindent\textbf{\texttt{is-metric-space(s)}}\quad\textit{predicate}\index{is-metric-space(s)@\texttt{is-metric-space(s)}}\\
reads: \$1 is a metric space \\
a predicate \\
characterized by \texttt{is-metric-space}  \textit{(definitional)}:
\begin{quote}\ttfamily\footnotesize\raggedright\noindent
forall(\allowbreak{}[s], is-metric-space(\allowbreak{}s) iff length(\allowbreak{}s) = 2 and pts(\allowbreak{}s) in set and dist(\allowbreak{}s) in fun(\allowbreak{}cartesian(\allowbreak{}pts(\allowbreak{}s), pts(\allowbreak{}s)), rr) and is-metric(\allowbreak{}dist(\allowbreak{}s), pts(\allowbreak{}s)) and s in set)
\end{quote}
also: \texttt{nag-metric-space-is-metric-space-cna-as-normed-ag} \texttt{nag-metric-space-is-metric-space-uba-as-normed-ag} \texttt{nag-metric-space-is-metric-space-cnvs-as-normed-ag} \texttt{nag-metric-space-is-metric-space-uba-as-normed-ag-uiba-as-uba} \texttt{cc-is-metric-space} \texttt{rr-is-metric-space} \texttt{nf-cc-is-metric-space} \texttt{nf-rr-is-metric-space} ... (37 in all) \\
mentioned by 747 result(s)

\smallskip

\noindent\textbf{\texttt{is-metrizable-top-space(s)}}\quad\textit{predicate}\index{is-metrizable-top-space(s)@\texttt{is-metrizable-top-space(s)}}\\
reads: \$1 is a metrizable topological space \\
a predicate \\
characterized by \texttt{is-metrizable-top-space-def}  \textit{(definitional)}:
\begin{quote}\ttfamily\footnotesize\raggedright\noindent
forall(\allowbreak{}[s], is-metrizable-top-space(\allowbreak{}s) iff is-top-space(\allowbreak{}s) and forsome(\allowbreak{}[md in set], is-metric-space(\allowbreak{}md) and s == [pts(\allowbreak{}md), \{u in power(\allowbreak{}pts(\allowbreak{}md)): is-open(\allowbreak{}md, u)\}]))
\end{quote}
also: \texttt{is-metrizable-top-space-is-set} \texttt{metric-top-is-metrizable-top-space} \\
mentioned by 26 result(s)

\smallskip

\noindent\textbf{\texttt{is-minimal-poly(gmp-md\_, gmp-t\_, gmp-k\_)}}\quad\textit{predicate}\index{is-minimal-poly(gmp-md_, gmp-t_, gmp-k_)@\texttt{is-minimal-poly(gmp-md\_, gmp-t\_, gmp-k\_)}}\\
reads: \$3 is the minimum polynomial of \$2 \\
defined by \texttt{is-minimal-poly} \\
characterized by \texttt{is-minimal-poly}  \textit{(definitional)}:
\begin{quote}\ttfamily\footnotesize\raggedright\noindent
forall(\allowbreak{}[gmp-md\_, gmp-t\_, gmp-k\_], is-minimal-poly(\allowbreak{}gmp-md\_, gmp-t\_, gmp-k\_) iff is-monic(\allowbreak{}scal(\allowbreak{}gmp-md\_), gmp-k\_) and forall(\allowbreak{}[gmpw\_ in vec(\allowbreak{}gmp-md\_)], poly-op(\allowbreak{}gmp-md\_, gmp-t\_, gmp-k\_, gmpw\_) = vzero(\allowbreak{}gmp-md\_)) and forall(\allowbreak{}[gmpq\_ in carr(\allowbreak{}poly(\allowbreak{}scal(\allowbreak{}gmp-md\_)))], not(\allowbreak{}gmpq\_ = zero(\allowbreak{}poly(\allowbreak{}scal(\allowbreak{}gmp-md\_)))) implies forall(\allowbreak{}[gmpw\_ in vec(\allowbreak{}gmp-md\_)], poly-op(\allowbreak{}gmp-md\_, gmp-t\_, gmpq\_, gmpw\_) = vzero(\allowbreak{}gmp-md\_)) implies deg(\allowbreak{}scal(\allowbreak{}gmp-md\_), gmp-k\_) <= deg(\allowbreak{}scal(\allowbreak{}gmp-md\_), gmpq\_)))
\end{quote}
mentioned by 5 result(s)

\smallskip

\noindent\textbf{\texttt{is-module(s)}}\quad\textit{predicate}\index{is-module(s)@\texttt{is-module(s)}}\\
reads: \$1 is a module \\
a predicate \\
characterized by \texttt{is-module}  \textit{(definitional)}:
\begin{quote}\ttfamily\footnotesize\raggedright\noindent
forall(\allowbreak{}[s], is-module(\allowbreak{}s) iff length(\allowbreak{}s) = 6 and is-ring(\allowbreak{}scal(\allowbreak{}s)) and vec(\allowbreak{}s) in set and vadd(\allowbreak{}s) in fun(\allowbreak{}cartesian(\allowbreak{}vec(\allowbreak{}s), vec(\allowbreak{}s)), vec(\allowbreak{}s)) and vzero(\allowbreak{}s) in vec(\allowbreak{}s) and vneg(\allowbreak{}s) in fun(\allowbreak{}vec(\allowbreak{}s), vec(\allowbreak{}s)) and act(\allowbreak{}s) in fun(\allowbreak{}cartesian(\allowbreak{}carr(\allowbreak{}scal(\allowbreak{}s)), vec(\allowbreak{}s)), vec(\allowbreak{}s)) and is-associative(\allowbreak{}vadd(\allowbreak{}s), vec(\allowbreak{}s)) and is-commutative(\allowbreak{}vadd(\allowbreak{}s), vec(\allowbreak{}s)) and is-identity(\allowbreak{}vadd(\allowbreak{}s), vzero(\allowbreak{}s), vec(\allowbreak{}s)) and has-inverses(\allowbreak{}vadd(\allowbreak{}s), vzero(\allowbreak{}s), vneg(\allowbreak{}s), vec(\allowbreak{}s)) and forall(\allowbreak{}[r\_ in carr(\allowbreak{}scal(\allowbreak{}s)), x\_ in vec(\allowbreak{}s), y\_ in vec(\allowbreak{}s)], (\allowbreak{}act(\allowbreak{}s))(\allowbreak{}r\_, (\allowbreak{}vadd(\allowbreak{}s))(\allowbreak{}x\_, y\_)) = (\allowbreak{}vadd(\allowbreak{}s))(\allowbreak{}(\allowbreak{}act(\allowbreak{}s))(\allowbreak{}r\_, x\_), (\allowbreak{}act(\allowbreak{}s))(\allowbreak{}r\_, y\_))) and forall(\allowbreak{}[r\_ in carr(\allowbreak{}scal(\allowbreak{}s)), s\_ in carr(\allowbreak{}scal(\allowbreak{}s)), x\_ in vec(\allowbreak{}s)], (\allowbreak{}act(\allowbreak{}s))(\allowbreak{}(\allowbreak{}add(\allowbreak{}scal(\allowbreak{}s)))(\allowbreak{}r\_, s\_), x\_) = (\allowbreak{}vadd(\allowbreak{}s))(\allowbreak{}(\allowbreak{}act(\allowbreak{}s))(\allowbreak{}r\_, x\_), (\allowbreak{}act(\allowbreak{}s))(\allowbreak{}s\_, x\_))) and forall(\allowbreak{}[r\_ in carr(\allowbreak{}scal(\allowbreak{}s)), s\_ in carr(\allowbreak{}scal(\allowbreak{}s)), x\_ in vec(\allowbreak{}s)], (\allowbreak{}act(\allowbreak{}s))(\allowbreak{}(\allowbreak{}mul(\allowbreak{}scal(\allowbreak{}s)))(\allowbreak{}r\_, s\_), x\_) = (\allowbreak{}act(\allowbreak{}s))(\allowbreak{}r\_, (\allowbreak{}act(\allowbreak{}s))(\allowbreak{}s\_, x\_))) and forall(\allowbreak{}[x\_ in vec(\allowbreak{}s)], (\allowbreak{}act(\allowbreak{}s))(\allowbreak{}one(\allowbreak{}scal(\allowbreak{}s)), x\_) = x\_) and s in set)
\end{quote}
also: \texttt{cna-as-module-is-module} \texttt{sip-as-module-is-module} \texttt{uba-as-module-is-module} \texttt{cnvs-as-module-is-module} \texttt{prod-rmod-is-module-cnvs-as-module} \texttt{sip-as-module-is-module-csip-as-sip} \texttt{uba-as-module-is-module-uiba-as-uba} \texttt{ring-as-module-is-module-uba-as-ring} ... (37 in all) \\
mentioned by 1079 result(s)

\smallskip

\noindent\textbf{\texttt{is-monic(gmo-f\_, gmo-p\_)}}\quad\textit{predicate}\index{is-monic(gmo-f_, gmo-p_)@\texttt{is-monic(gmo-f\_, gmo-p\_)}}\\
reads: \$2 is a monic polynomial over \$1 \\
defined by \texttt{is-monic} \\
characterized by \texttt{is-monic}  \textit{(definitional)}:
\begin{quote}\ttfamily\footnotesize\raggedright\noindent
forall(\allowbreak{}[gmo-f\_, gmo-p\_], is-monic(\allowbreak{}gmo-f\_, gmo-p\_) iff gmo-p\_ in carr(\allowbreak{}poly(\allowbreak{}gmo-f\_)) and gmo-p\_(\allowbreak{}deg(\allowbreak{}gmo-f\_, gmo-p\_)) = one(\allowbreak{}gmo-f\_))
\end{quote}
mentioned by 9 result(s)

\smallskip

\noindent\textbf{\texttt{is-monoid(s)}}\quad\textit{predicate}\index{is-monoid(s)@\texttt{is-monoid(s)}}\\
reads: a monoid \\
a predicate \\
characterized by \texttt{is-monoid}  \textit{(definitional)}:
\begin{quote}\ttfamily\footnotesize\raggedright\noindent
forall(\allowbreak{}[s], is-monoid(\allowbreak{}s) iff length(\allowbreak{}s) = 3 and carr(\allowbreak{}s) in set and opr(\allowbreak{}s) in fun(\allowbreak{}cartesian(\allowbreak{}carr(\allowbreak{}s), carr(\allowbreak{}s)), carr(\allowbreak{}s)) and iden(\allowbreak{}s) in carr(\allowbreak{}s) and is-associative(\allowbreak{}opr(\allowbreak{}s), carr(\allowbreak{}s)) and is-identity(\allowbreak{}opr(\allowbreak{}s), iden(\allowbreak{}s), carr(\allowbreak{}s)) and s in set)
\end{quote}
also: \texttt{abelian-group-as-monoid-is-monoid} \texttt{ring-multiplicative-monoid-is-monoid} \texttt{ring-multiplicative-monoid-is-monoid-uba-as-ring} \texttt{abelian-group-as-monoid-is-monoid-module-vector-ag} \texttt{abelian-group-as-monoid-is-monoid-ring-additive-ag} \texttt{abelian-group-as-monoid-is-monoid-field-additive-ag} \texttt{abelian-group-as-monoid-is-monoid-normed-field-additive-ag} \texttt{abelian-group-as-monoid-is-monoid-normed-ag-as-abelian-group} ... (32 in all) \\
mentioned by 164 result(s)

\smallskip

\noindent\textbf{\texttt{is-monotone-class(omega, mc)}}\quad\textit{predicate}\index{is-monotone-class(omega, mc)@\texttt{is-monotone-class(omega, mc)}}\\
reads: \$2 is a monotone class of subsets of \$1 \\
defined by \texttt{is-monotone-class} \\
characterized by \texttt{is-monotone-class}  \textit{(definitional)}:
\begin{quote}\ttfamily\footnotesize\raggedright\noindent
forall(\allowbreak{}[omega, mc], is-monotone-class(\allowbreak{}omega, mc) iff omega in set and mc subset power(\allowbreak{}omega) and forall(\allowbreak{}[mcf\_ in fun(\allowbreak{}nn, mc)], forall(\allowbreak{}[mcn\_ in nn], mcf\_(\allowbreak{}mcn\_) subset mcf\_(\allowbreak{}succ(\allowbreak{}mcn\_))) implies big-union(\allowbreak{}mck\_, nn, mcf\_(\allowbreak{}mck\_)) in mc) and forall(\allowbreak{}[mcg\_ in fun(\allowbreak{}nn, mc)], forall(\allowbreak{}[mcm\_ in nn], mcg\_(\allowbreak{}succ(\allowbreak{}mcm\_)) subset mcg\_(\allowbreak{}mcm\_)) implies complement-in(\allowbreak{}omega, big-union(\allowbreak{}mcj\_, nn, complement-in(\allowbreak{}omega, mcg\_(\allowbreak{}mcj\_)))) in mc))
\end{quote}
also: \texttt{monotone-class-generated-is-monotone-class} \\
mentioned by 11 result(s)

\smallskip

\noindent\textbf{\texttt{is-ms-sequence(ms)}}\quad\textit{predicate}\index{is-ms-sequence(ms)@\texttt{is-ms-sequence(ms)}}\\
reads: \$1 is a sequence of metric spaces \\
defined by \texttt{is-ms-sequence} \\
characterized by \texttt{is-ms-sequence}  \textit{(definitional)}:
\begin{quote}\ttfamily\footnotesize\raggedright\noindent
forall(\allowbreak{}[ms], is-ms-sequence(\allowbreak{}ms) iff forall(\allowbreak{}[n in nn], is-metric-space(\allowbreak{}ms(\allowbreak{}n))))
\end{quote}
mentioned by 48 result(s)

\smallskip

\noindent\textbf{\texttt{is-multiplicative(f)}}\quad\textit{predicate}\index{is-multiplicative(f)@\texttt{is-multiplicative(f)}}\\
reads: \$1 is multiplicative \\
defined by \texttt{is-multiplicative} \\
characterized by \texttt{is-multiplicative}  \textit{(definitional)}:
\begin{quote}\ttfamily\footnotesize\raggedright\noindent
forall(\allowbreak{}[f], is-multiplicative(\allowbreak{}f) iff f in fun(\allowbreak{}nn, rr) and forall(\allowbreak{}[mlm\_ in nn, mln\_ in nn], 1 <= mlm\_ implies 1 <= mln\_ implies zz-coprime(\allowbreak{}mlm\_, mln\_) implies f(\allowbreak{}mlm\_ * mln\_) = f(\allowbreak{}mlm\_) * f(\allowbreak{}mln\_)))
\end{quote}
mentioned by 5 result(s)

\smallskip

\noindent\textbf{\texttt{is-nf-subfield(bzk\_, bzs\_)}}\quad\textit{predicate}\index{is-nf-subfield(bzk_, bzs_)@\texttt{is-nf-subfield(bzk\_, bzs\_)}}\\
reads: \$2 is a subfield of \$1 \\
defined by \texttt{is-nf-subfield} \\
characterized by \texttt{is-nf-subfield}  \textit{(definitional)}:
\begin{quote}\ttfamily\footnotesize\raggedright\noindent
forall(\allowbreak{}[bzk\_, bzs\_], is-nf-subfield(\allowbreak{}bzk\_, bzs\_) iff forall(\allowbreak{}[bzx\_ in bzs\_], bzx\_ in carr(\allowbreak{}bzk\_)) and zero(\allowbreak{}bzk\_) in bzs\_ and one(\allowbreak{}bzk\_) in bzs\_ and forall(\allowbreak{}[bzx\_ in bzs\_, bzy\_ in bzs\_], (\allowbreak{}add(\allowbreak{}bzk\_))(\allowbreak{}bzx\_, bzy\_) in bzs\_) and forall(\allowbreak{}[bzx\_ in bzs\_, bzy\_ in bzs\_], (\allowbreak{}mul(\allowbreak{}bzk\_))(\allowbreak{}bzx\_, bzy\_) in bzs\_) and forall(\allowbreak{}[bzx\_ in bzs\_], (\allowbreak{}neg(\allowbreak{}bzk\_))(\allowbreak{}bzx\_) in bzs\_) and forall(\allowbreak{}[bzx\_ in bzs\_], not(\allowbreak{}bzx\_ = zero(\allowbreak{}bzk\_)) implies forsome(\allowbreak{}[bzy\_ in bzs\_], (\allowbreak{}mul(\allowbreak{}bzk\_))(\allowbreak{}bzx\_, bzy\_) = one(\allowbreak{}bzk\_))))
\end{quote}
mentioned by 76 result(s)

\smallskip

\noindent\textbf{\texttt{is-nilpotent-op(gnp-md\_, gnp-t\_)}}\quad\textit{predicate}\index{is-nilpotent-op(gnp-md_, gnp-t_)@\texttt{is-nilpotent-op(gnp-md\_, gnp-t\_)}}\\
reads: \$2 is nilpotent \\
defined by \texttt{is-nilpotent-op} \\
characterized by \texttt{is-nilpotent-op}  \textit{(definitional)}:
\begin{quote}\ttfamily\footnotesize\raggedright\noindent
forall(\allowbreak{}[gnp-md\_, gnp-t\_], is-nilpotent-op(\allowbreak{}gnp-md\_, gnp-t\_) iff forsome(\allowbreak{}[gnpk\_ in nn], forall(\allowbreak{}[gnpw\_ in vec(\allowbreak{}gnp-md\_)], op-iter(\allowbreak{}gnp-md\_, gnp-t\_, gnpw\_, gnpk\_) = vzero(\allowbreak{}gnp-md\_))))
\end{quote}
mentioned by 7 result(s)

\smallskip

\noindent\textbf{\texttt{is-noetherian(m)}}\quad\textit{predicate}\index{is-noetherian(m)@\texttt{is-noetherian(m)}}\\
reads: \$1 is Noetherian \\
defined by \texttt{is-noetherian} \\
characterized by \texttt{is-noetherian}  \textit{(definitional)}:
\begin{quote}\ttfamily\footnotesize\raggedright\noindent
forall(\allowbreak{}[m], is-noetherian(\allowbreak{}m) iff is-module(\allowbreak{}m) and forall(\allowbreak{}[f\_ in fun(\allowbreak{}nn, power(\allowbreak{}vec(\allowbreak{}m)))], forall(\allowbreak{}[n\_ in nn], is-submodule(\allowbreak{}m, f\_(\allowbreak{}n\_))) and forall(\allowbreak{}[n\_ in nn], f\_(\allowbreak{}n\_) subset f\_(\allowbreak{}succ(\allowbreak{}n\_))) implies forsome(\allowbreak{}[k\_ in nn], forall(\allowbreak{}[n\_], n\_ in nn and k\_ <= n\_ implies f\_(\allowbreak{}n\_) = f\_(\allowbreak{}k\_)))))
\end{quote}
mentioned by 5 result(s)

\smallskip

\noindent\textbf{\texttt{is-nonmeager(s, a)}}\quad\textit{predicate}\index{is-nonmeager(s, a)@\texttt{is-nonmeager(s, a)}}\\
reads: \$2 is nonmeager in \$1 \\
defined by \texttt{is-nonmeager} \\
characterized by \texttt{is-nonmeager}  \textit{(definitional)}:
\begin{quote}\ttfamily\footnotesize\raggedright\noindent
forall(\allowbreak{}[s, a], is-nonmeager(\allowbreak{}s, a) iff is-metric-space(\allowbreak{}s) and a subset pts(\allowbreak{}s) and not(\allowbreak{}is-meager(\allowbreak{}s, a)))
\end{quote}
mentioned by 3 result(s)

\smallskip

\noindent\textbf{\texttt{is-nontrivial-character(chk\_, chc\_)}}\quad\textit{predicate}\index{is-nontrivial-character(chk_, chc_)@\texttt{is-nontrivial-character(chk\_, chc\_)}}\\
reads: \$2 is a non-trivial character of the additive group of \$1 \\
defined by \texttt{is-nontrivial-character} \\
characterized by \texttt{is-nontrivial-character}  \textit{(definitional)}:
\begin{quote}\ttfamily\footnotesize\raggedright\noindent
forall(\allowbreak{}[chk\_, chc\_], is-nontrivial-character(\allowbreak{}chk\_, chc\_) iff is-add-character(\allowbreak{}chk\_, chc\_) and forsome(\allowbreak{}[chx\_ in carr(\allowbreak{}chk\_)], not(\allowbreak{}chc\_(\allowbreak{}chx\_) = 1)))
\end{quote}
mentioned by 30 result(s)

\smallskip

\noindent\textbf{\texttt{is-nontrivial-nv(cpk\_)}}\quad\textit{predicate}\index{is-nontrivial-nv(cpk_)@\texttt{is-nontrivial-nv(cpk\_)}}\\
reads: the valuation of \$1 is non-trivial \\
defined by \texttt{is-nontrivial-nv} \\
characterized by \texttt{is-nontrivial-nv}  \textit{(definitional)}:
\begin{quote}\ttfamily\footnotesize\raggedright\noindent
forall(\allowbreak{}[cpk\_], is-nontrivial-nv(\allowbreak{}cpk\_) iff forsome(\allowbreak{}[cpa\_ in carr(\allowbreak{}cpk\_)], not(\allowbreak{}(\allowbreak{}fnrm(\allowbreak{}cpk\_))(\allowbreak{}cpa\_) = 0) and not(\allowbreak{}(\allowbreak{}fnrm(\allowbreak{}cpk\_))(\allowbreak{}cpa\_) = 1)))
\end{quote}
mentioned by 4 result(s)

\smallskip

\noindent\textbf{\texttt{is-nontrivially-valued(fld)}}\quad\textit{predicate}\index{is-nontrivially-valued(fld)@\texttt{is-nontrivially-valued(fld)}}\\
reads: \$1 is a non-trivially valued normed field \\
defined by \texttt{is-nontrivially-valued} \\
characterized by \texttt{is-nontrivially-valued}  \textit{(definitional)}:
\begin{quote}\ttfamily\footnotesize\raggedright\noindent
forall(\allowbreak{}[fld], is-nontrivially-valued(\allowbreak{}fld) iff is-normed-field(\allowbreak{}fld) and forsome(\allowbreak{}[c in carr(\allowbreak{}fld)], 1 < (\allowbreak{}fnrm(\allowbreak{}fld))(\allowbreak{}c) and forsome(\allowbreak{}[b in carr(\allowbreak{}fld)], (\allowbreak{}mul(\allowbreak{}fld))(\allowbreak{}c, b) = one(\allowbreak{}fld))))
\end{quote}
mentioned by 12 result(s)

\smallskip

\noindent\textbf{\texttt{is-norm(nm, addop, mulop, zr, crr)}}\quad\textit{predicate}\index{is-norm(nm, addop, mulop, zr, crr)@\texttt{is-norm(nm, addop, mulop, zr, crr)}}\\
reads: \$1 is a norm on \$5 \\
defined by \texttt{is-norm} \\
characterized by \texttt{is-norm}  \textit{(definitional)}:
\begin{quote}\ttfamily\footnotesize\raggedright\noindent
forall(\allowbreak{}[nm, addop, mulop, zr, crr], is-norm(\allowbreak{}nm, addop, mulop, zr, crr) iff nm in fun(\allowbreak{}crr, rr) and forall(\allowbreak{}[a in crr], 0 <= nm(\allowbreak{}a) and (\allowbreak{}nm(\allowbreak{}a) = 0 iff a = zr) and forall(\allowbreak{}[b in crr], nm(\allowbreak{}mulop(\allowbreak{}a, b)) = nm(\allowbreak{}a) * nm(\allowbreak{}b) and nm(\allowbreak{}addop(\allowbreak{}a, b)) <= nm(\allowbreak{}a) + nm(\allowbreak{}b))))
\end{quote}
also: \texttt{is-normed-field} \\
mentioned by 4 result(s)

\smallskip

\noindent\textbf{\texttt{is-norm-data(md, sm, n, xs, ys, am)}}\quad\textit{predicate}\index{is-norm-data(md, sm, n, xs, ys, am)@\texttt{is-norm-data(md, sm, n, xs, ys, am)}}\\
reads: \$4 freely generates \$1, \$5 is relation-free and spans \$2, both of length \$3, and \$5 = \$6 . \$4 \\
defined by \texttt{is-norm-data} \\
characterized by \texttt{is-norm-data}  \textit{(definitional)}:
\begin{quote}\ttfamily\footnotesize\raggedright\noindent
forall(\allowbreak{}[md, sm, n, xs, ys, am], is-norm-data(\allowbreak{}md, sm, n, xs, ys, am) iff n in nn and xs in mat(\allowbreak{}n, 1, vec(\allowbreak{}md)) and generates(\allowbreak{}md, n, xs) and rel-free(\allowbreak{}md, n, xs) and ys in mat(\allowbreak{}n, 1, vec(\allowbreak{}md)) and spans(\allowbreak{}md, n, ys, sm) and rel-free(\allowbreak{}md, n, ys) and am in mat(\allowbreak{}n, n, zz) and ys = matact(\allowbreak{}md, am, xs))
\end{quote}
mentioned by 79 result(s)

\smallskip

\noindent\textbf{\texttt{is-norm-on(k, md, nu)}}\quad\textit{predicate}\index{is-norm-on(k, md, nu)@\texttt{is-norm-on(k, md, nu)}}\\
reads: \$3 is a norm on the vector space \$2 over \$1 \\
defined by \texttt{is-norm-on} \\
characterized by \texttt{is-norm-on}  \textit{(definitional)}:
\begin{quote}\ttfamily\footnotesize\raggedright\noindent
forall(\allowbreak{}[k, md, nu], is-norm-on(\allowbreak{}k, md, nu) iff is-module(\allowbreak{}md) and scal(\allowbreak{}md) = normed-field-as-commutative-ring(\allowbreak{}k) and nu in fun(\allowbreak{}vec(\allowbreak{}md), rr) and forall(\allowbreak{}[zkvx\_ in vec(\allowbreak{}md)], 0 <= nu(\allowbreak{}zkvx\_)) and forall(\allowbreak{}[zkvx\_ in vec(\allowbreak{}md)], nu(\allowbreak{}zkvx\_) = 0 iff zkvx\_ = vzero(\allowbreak{}md)) and forall(\allowbreak{}[zkvx\_ in vec(\allowbreak{}md), zkvy\_ in vec(\allowbreak{}md)], nu(\allowbreak{}(\allowbreak{}vadd(\allowbreak{}md))(\allowbreak{}zkvx\_, zkvy\_)) <= nu(\allowbreak{}zkvx\_) + nu(\allowbreak{}zkvy\_)) and forall(\allowbreak{}[zkvr\_ in carr(\allowbreak{}k), zkvx\_ in vec(\allowbreak{}md)], nu(\allowbreak{}(\allowbreak{}act(\allowbreak{}md))(\allowbreak{}zkvr\_, zkvx\_)) = (\allowbreak{}fnrm(\allowbreak{}k))(\allowbreak{}zkvr\_) * nu(\allowbreak{}zkvx\_)))
\end{quote}
mentioned by 11 result(s)

\smallskip

\noindent\textbf{\texttt{is-normal-el(baa\_, bax\_)}}\quad\textit{predicate}\index{is-normal-el(baa_, bax_)@\texttt{is-normal-el(baa\_, bax\_)}}\\
reads: \$2 is normal in \$1 \\
defined by \texttt{is-normal-el} \\
characterized by \texttt{is-normal-el}  \textit{(definitional)}:
\begin{quote}\ttfamily\footnotesize\raggedright\noindent
forall(\allowbreak{}[baa\_, bax\_], is-normal-el(\allowbreak{}baa\_, bax\_) iff bax\_ in vec(\allowbreak{}baa\_) and (\allowbreak{}vmul(\allowbreak{}baa\_))(\allowbreak{}(\allowbreak{}vstar(\allowbreak{}baa\_))(\allowbreak{}bax\_), bax\_) = (\allowbreak{}vmul(\allowbreak{}baa\_))(\allowbreak{}bax\_, (\allowbreak{}vstar(\allowbreak{}baa\_))(\allowbreak{}bax\_)))
\end{quote}
mentioned by 2 result(s)

\smallskip

\noindent\textbf{\texttt{is-normal-subgroup(sgg\_, sgh\_)}}\quad\textit{predicate}\index{is-normal-subgroup(sgg_, sgh_)@\texttt{is-normal-subgroup(sgg\_, sgh\_)}}\\
reads: \$2 is a normal subgroup of \$1 \\
defined by \texttt{is-normal-subgroup} \\
characterized by \texttt{is-normal-subgroup}  \textit{(definitional)}:
\begin{quote}\ttfamily\footnotesize\raggedright\noindent
forall(\allowbreak{}[sgg\_, sgh\_], is-normal-subgroup(\allowbreak{}sgg\_, sgh\_) iff is-subgroup(\allowbreak{}sgg\_, sgh\_) and forall(\allowbreak{}[sgu\_ in carr(\allowbreak{}sgg\_), sga\_ in sgh\_], (\allowbreak{}opr(\allowbreak{}sgg\_))(\allowbreak{}(\allowbreak{}opr(\allowbreak{}sgg\_))(\allowbreak{}sgu\_, sga\_), (\allowbreak{}inv(\allowbreak{}sgg\_))(\allowbreak{}sgu\_)) in sgh\_))
\end{quote}
mentioned by 4 result(s)

\smallskip

\noindent\textbf{\texttt{is-normally-convergent(u, cpsxs\_)}}\quad\textit{predicate}\index{is-normally-convergent(u, cpsxs_)@\texttt{is-normally-convergent(u, cpsxs\_)}}\\
reads: the series of the functions \$1 is normally convergent on \$2 \\
defined by \texttt{is-normally-convergent} \\
characterized by \texttt{is-normally-convergent}  \textit{(definitional)}:
\begin{quote}\ttfamily\footnotesize\raggedright\noindent
forall(\allowbreak{}[u, cpsxs\_], is-normally-convergent(\allowbreak{}u, cpsxs\_) iff forall(\allowbreak{}[cpsi\_ in nn], cc-sup-norm(\allowbreak{}u(\allowbreak{}cpsi\_), cpsxs\_) in rr) and series-converges(\allowbreak{}vnb-lambda(\allowbreak{}cpsi\_, nn, cc-sup-norm(\allowbreak{}u(\allowbreak{}cpsi\_), cpsxs\_))))
\end{quote}
mentioned by 14 result(s)

\smallskip

\noindent\textbf{\texttt{is-normed-ag(s)}}\quad\textit{predicate}\index{is-normed-ag(s)@\texttt{is-normed-ag(s)}}\\
reads: a normed ag \\
a predicate \\
characterized by \texttt{is-normed-ag}  \textit{(definitional)}:
\begin{quote}\ttfamily\footnotesize\raggedright\noindent
forall(\allowbreak{}[s], is-normed-ag(\allowbreak{}s) iff length(\allowbreak{}s) = 5 and carr(\allowbreak{}s) in set and opr(\allowbreak{}s) in fun(\allowbreak{}cartesian(\allowbreak{}carr(\allowbreak{}s), carr(\allowbreak{}s)), carr(\allowbreak{}s)) and iden(\allowbreak{}s) in carr(\allowbreak{}s) and inv(\allowbreak{}s) in fun(\allowbreak{}carr(\allowbreak{}s), carr(\allowbreak{}s)) and nrm(\allowbreak{}s) in fun(\allowbreak{}carr(\allowbreak{}s), rr) and is-associative(\allowbreak{}opr(\allowbreak{}s), carr(\allowbreak{}s)) and is-identity(\allowbreak{}opr(\allowbreak{}s), iden(\allowbreak{}s), carr(\allowbreak{}s)) and has-inverses(\allowbreak{}opr(\allowbreak{}s), iden(\allowbreak{}s), inv(\allowbreak{}s), carr(\allowbreak{}s)) and is-commutative(\allowbreak{}opr(\allowbreak{}s), carr(\allowbreak{}s)) and is-group-norm(\allowbreak{}nrm(\allowbreak{}s), opr(\allowbreak{}s), inv(\allowbreak{}s), iden(\allowbreak{}s), carr(\allowbreak{}s)) and s in set)
\end{quote}
also: \texttt{cna-as-normed-ag-is-normed-ag} \texttt{uba-as-normed-ag-is-normed-ag} \texttt{cnvs-as-normed-ag-is-normed-ag} \texttt{uba-as-normed-ag-is-normed-ag-uiba-as-uba} \texttt{normed-vector-space-as-normed-ag-is-normed-ag} \texttt{normed-vector-space-as-normed-ag-is-normed-ag-cnvs-as-nvs} \texttt{cc-nag-is-normed-ag} \texttt{is-normed-ag-is-set} ... (12 in all) \\
mentioned by 46 result(s)

\smallskip

\noindent\textbf{\texttt{is-normed-field(s)}}\quad\textit{predicate}\index{is-normed-field(s)@\texttt{is-normed-field(s)}}\\
reads: \$1 is a normed field \\
a predicate \\
characterized by \texttt{is-normed-field}  \textit{(definitional)}:
\begin{quote}\ttfamily\footnotesize\raggedright\noindent
forall(\allowbreak{}[s], is-normed-field(\allowbreak{}s) iff length(\allowbreak{}s) = 7 and carr(\allowbreak{}s) in set and add(\allowbreak{}s) in fun(\allowbreak{}cartesian(\allowbreak{}carr(\allowbreak{}s), carr(\allowbreak{}s)), carr(\allowbreak{}s)) and mul(\allowbreak{}s) in fun(\allowbreak{}cartesian(\allowbreak{}carr(\allowbreak{}s), carr(\allowbreak{}s)), carr(\allowbreak{}s)) and neg(\allowbreak{}s) in fun(\allowbreak{}carr(\allowbreak{}s), carr(\allowbreak{}s)) and zero(\allowbreak{}s) in carr(\allowbreak{}s) and one(\allowbreak{}s) in carr(\allowbreak{}s) and fnrm(\allowbreak{}s) in fun(\allowbreak{}carr(\allowbreak{}s), rr) and is-associative(\allowbreak{}add(\allowbreak{}s), carr(\allowbreak{}s)) and is-commutative(\allowbreak{}add(\allowbreak{}s), carr(\allowbreak{}s)) and is-identity(\allowbreak{}add(\allowbreak{}s), zero(\allowbreak{}s), carr(\allowbreak{}s)) and has-inverses(\allowbreak{}add(\allowbreak{}s), zero(\allowbreak{}s), neg(\allowbreak{}s), carr(\allowbreak{}s)) and is-associative(\allowbreak{}mul(\allowbreak{}s), carr(\allowbreak{}s)) and is-commutative(\allowbreak{}mul(\allowbreak{}s), carr(\allowbreak{}s)) and is-identity(\allowbreak{}mul(\allowbreak{}s), one(\allowbreak{}s), carr(\allowbreak{}s)) and is-distributive(\allowbreak{}add(\allowbreak{}s), mul(\allowbreak{}s), carr(\allowbreak{}s)) and is-norm(\allowbreak{}fnrm(\allowbreak{}s), add(\allowbreak{}s), mul(\allowbreak{}s), zero(\allowbreak{}s), carr(\allowbreak{}s)) and not(\allowbreak{}zero(\allowbreak{}s) = one(\allowbreak{}s)) and forall(\allowbreak{}[a\_ in carr(\allowbreak{}s)], not(\allowbreak{}a\_ = zero(\allowbreak{}s)) implies forsome(\allowbreak{}[b\_ in carr(\allowbreak{}s)], not(\allowbreak{}b\_ = zero(\allowbreak{}s)) and (\allowbreak{}mul(\allowbreak{}s))(\allowbreak{}a\_, b\_) = one(\allowbreak{}s))) and s in set)
\end{quote}
also: \texttt{cc-is-normed-field} \texttt{qp-is-normed-field} \texttt{rr-is-normed-field} \texttt{is-normed-field-is-set} \texttt{sub-nf-is-normed-field} \texttt{triv-nf-is-normed-field} \texttt{qq-padic-is-normed-field} \texttt{laurent-nf-is-normed-field} \\
mentioned by 471 result(s)

\smallskip

\noindent\textbf{\texttt{is-normed-vector-space(s)}}\quad\textit{predicate}\index{is-normed-vector-space(s)@\texttt{is-normed-vector-space(s)}}\\
reads: \$1 is a normed vector space \\
a predicate \\
characterized by \texttt{is-normed-vector-space}  \textit{(definitional)}:
\begin{quote}\ttfamily\footnotesize\raggedright\noindent
forall(\allowbreak{}[s], is-normed-vector-space(\allowbreak{}s) iff length(\allowbreak{}s) = 7 and is-ring(\allowbreak{}scal(\allowbreak{}s)) and vec(\allowbreak{}s) in set and vadd(\allowbreak{}s) in fun(\allowbreak{}cartesian(\allowbreak{}vec(\allowbreak{}s), vec(\allowbreak{}s)), vec(\allowbreak{}s)) and vzero(\allowbreak{}s) in vec(\allowbreak{}s) and vneg(\allowbreak{}s) in fun(\allowbreak{}vec(\allowbreak{}s), vec(\allowbreak{}s)) and act(\allowbreak{}s) in fun(\allowbreak{}cartesian(\allowbreak{}carr(\allowbreak{}scal(\allowbreak{}s)), vec(\allowbreak{}s)), vec(\allowbreak{}s)) and vnrm(\allowbreak{}s) in fun(\allowbreak{}vec(\allowbreak{}s), rr) and is-associative(\allowbreak{}vadd(\allowbreak{}s), vec(\allowbreak{}s)) and is-commutative(\allowbreak{}vadd(\allowbreak{}s), vec(\allowbreak{}s)) and is-identity(\allowbreak{}vadd(\allowbreak{}s), vzero(\allowbreak{}s), vec(\allowbreak{}s)) and has-inverses(\allowbreak{}vadd(\allowbreak{}s), vzero(\allowbreak{}s), vneg(\allowbreak{}s), vec(\allowbreak{}s)) and scal(\allowbreak{}s) = normed-field-as-commutative-ring(\allowbreak{}rr-normed-field) and forall(\allowbreak{}[r\_ in carr(\allowbreak{}scal(\allowbreak{}s)), x\_ in vec(\allowbreak{}s), y\_ in vec(\allowbreak{}s)], (\allowbreak{}act(\allowbreak{}s))(\allowbreak{}r\_, (\allowbreak{}vadd(\allowbreak{}s))(\allowbreak{}x\_, y\_)) = (\allowbreak{}vadd(\allowbreak{}s))(\allowbreak{}(\allowbreak{}act(\allowbreak{}s))(\allowbreak{}r\_, x\_), (\allowbreak{}act(\allowbreak{}s))(\allowbreak{}r\_, y\_))) and forall(\allowbreak{}[r\_ in carr(\allowbreak{}scal(\allowbreak{}s)), s\_ in carr(\allowbreak{}scal(\allowbreak{}s)), x\_ in vec(\allowbreak{}s)], (\allowbreak{}act(\allowbreak{}s))(\allowbreak{}(\allowbreak{}add(\allowbreak{}scal(\allowbreak{}s)))(\allowbreak{}r\_, s\_), x\_) = (\allowbreak{}vadd(\allowbreak{}s))(\allowbreak{}(\allowbreak{}act(\allowbreak{}s))(\allowbreak{}r\_, x\_), (\allowbreak{}act(\allowbreak{}s))(\allowbreak{}s\_, x\_))) and forall(\allowbreak{}[r\_ in carr(\allowbreak{}scal(\allowbreak{}s)), s\_ in carr(\allowbreak{}scal(\allowbreak{}s)), x\_ in vec(\allowbreak{}s)], (\allowbreak{}act(\allowbreak{}s))(\allowbreak{}(\allowbreak{}mul(\allowbreak{}scal(\allowbreak{}s)))(\allowbreak{}r\_, s\_), x\_) = (\allowbreak{}act(\allowbreak{}s))(\allowbreak{}r\_, (\allowbreak{}act(\allowbreak{}s))(\allowbreak{}s\_, x\_))) and forall(\allowbreak{}[x\_ in vec(\allowbreak{}s)], (\allowbreak{}act(\allowbreak{}s))(\allowbreak{}one(\allowbreak{}scal(\allowbreak{}s)), x\_) = x\_) and forall(\allowbreak{}[x\_ in vec(\allowbreak{}s)], 0 <= (\allowbreak{}vnrm(\allowbreak{}s))(\allowbreak{}x\_)) and forall(\allowbreak{}[x\_ in vec(\allowbreak{}s)], (\allowbreak{}vnrm(\allowbreak{}s))(\allowbreak{}x\_) = 0 iff x\_ = vzero(\allowbreak{}s)) and forall(\allowbreak{}[r\_ in carr(\allowbreak{}scal(\allowbreak{}s)), x\_ in vec(\allowbreak{}s)], (\allowbreak{}vnrm(\allowbreak{}s))(\allowbreak{}(\allowbreak{}act(\allowbreak{}s))(\allowbreak{}r\_, x\_)) = abs(\allowbreak{}r\_) * (\allowbreak{}vnrm(\allowbreak{}s))(\allowbreak{}x\_)) and forall(\allowbreak{}[x\_ in vec(\allowbreak{}s), y\_ in vec(\allowbreak{}s)], (\allowbreak{}vnrm(\allowbreak{}s))(\allowbreak{}(\allowbreak{}vadd(\allowbreak{}s))(\allowbreak{}x\_, y\_)) <= (\allowbreak{}vnrm(\allowbreak{}s))(\allowbreak{}x\_) + (\allowbreak{}vnrm(\allowbreak{}s))(\allowbreak{}y\_)) and s in set)
\end{quote}
also: \texttt{cnvs-as-nvs-is-normed-vector-space} \texttt{dual-is-normed-vector-space-cnvs-as-nvs} \texttt{rn-is-normed-vector-space} \texttt{dual-is-normed-vector-space} \texttt{ck-nvs-is-normed-vector-space} \texttt{is-normed-vector-space-is-set} \texttt{rr-nvs-is-normed-vector-space} \\
mentioned by 434 result(s)

\smallskip

\noindent\textbf{\texttt{is-nowhere-dense(s, a)}}\quad\textit{predicate}\index{is-nowhere-dense(s, a)@\texttt{is-nowhere-dense(s, a)}}\\
reads: \$2 is nowhere dense in \$1 \\
defined by \texttt{is-nowhere-dense} \\
characterized by \texttt{is-nowhere-dense}  \textit{(definitional)}:
\begin{quote}\ttfamily\footnotesize\raggedright\noindent
forall(\allowbreak{}[s, a], is-nowhere-dense(\allowbreak{}s, a) iff is-metric-space(\allowbreak{}s) and a subset pts(\allowbreak{}s) and interior(\allowbreak{}s, closure(\allowbreak{}s, a)) = empty-set)
\end{quote}
mentioned by 6 result(s)

\smallskip

\noindent\textbf{\texttt{is-nsl-space(nlp\_)}}\quad\textit{predicate}\index{is-nsl-space(nlp_)@\texttt{is-nsl-space(nlp\_)}}\\
reads: \$1 is a Needham-Schroeder-Lowe space \\
defined by \texttt{is-nsl-space} \\
characterized by \texttt{is-nsl-space}  \textit{(definitional)}:
\begin{quote}\ttfamily\footnotesize\raggedright\noindent
forall(\allowbreak{}[nlp\_], is-nsl-space(\allowbreak{}nlp\_) iff is-pk-space(\allowbreak{}nlp\_) and forall(\allowbreak{}[nli\_ in ids(\allowbreak{}ispace(\allowbreak{}pk-inf(\allowbreak{}nlp\_)))], nli\_ in adv-ids(\allowbreak{}pk-inf(\allowbreak{}nlp\_)) or forsome(\allowbreak{}[nla\_, nlb\_, nlm\_, nln\_], is-init-strand(\allowbreak{}nlp\_, (\allowbreak{}strand-map(\allowbreak{}ispace(\allowbreak{}pk-inf(\allowbreak{}nlp\_))))(\allowbreak{}nli\_), nla\_, nlb\_, nlm\_, nln\_)) or forsome(\allowbreak{}[nla\_, nlb\_, nlm\_, nln\_], is-resp-strand(\allowbreak{}nlp\_, (\allowbreak{}strand-map(\allowbreak{}ispace(\allowbreak{}pk-inf(\allowbreak{}nlp\_))))(\allowbreak{}nli\_), nla\_, nlb\_, nlm\_, nln\_))))
\end{quote}
mentioned by 23 result(s)

\smallskip

\noindent\textbf{\texttt{is-nso-init-strand(nlp\_, nls\_, nla\_, nlb\_, nlm\_, nln\_)}}\quad\textit{predicate}\index{is-nso-init-strand(nlp_, nls_, nla_, nlb_, nlm_, nln_)@\texttt{is-nso-init-strand(nlp\_, nls\_, nla\_, nlb\_, nlm\_, nln\_)}}\\
reads: \$2 is an original Needham-Schroeder initiator's strand of \$1 for the names \$3, \$4 and the values \$5, \$6 \\
defined by \texttt{is-nso-init-strand} \\
characterized by \texttt{is-nso-init-strand}  \textit{(definitional)}:
\begin{quote}\ttfamily\footnotesize\raggedright\noindent
forall(\allowbreak{}[nlp\_, nls\_, nla\_, nlb\_, nlm\_, nln\_], is-nso-init-strand(\allowbreak{}nlp\_, nls\_, nla\_, nlb\_, nlm\_, nln\_) iff nla\_ in names(\allowbreak{}nlp\_) and nlb\_ in names(\allowbreak{}nlp\_) and nlm\_ in tokens(\allowbreak{}ispace(\allowbreak{}pk-inf(\allowbreak{}nlp\_))) and nln\_ in tokens(\allowbreak{}ispace(\allowbreak{}pk-inf(\allowbreak{}nlp\_))) and not(\allowbreak{}nlm\_ in names(\allowbreak{}nlp\_)) and not(\allowbreak{}nlm\_ in keys(\allowbreak{}ispace(\allowbreak{}pk-inf(\allowbreak{}nlp\_)))) and not(\allowbreak{}nln\_ in keys(\allowbreak{}ispace(\allowbreak{}pk-inf(\allowbreak{}nlp\_)))) and nls\_ = [[sent, encr(\allowbreak{}lst(\allowbreak{}[nlm\_, nla\_]), (\allowbreak{}pubkey(\allowbreak{}nlp\_))(\allowbreak{}nlb\_))], [received, encr(\allowbreak{}lst(\allowbreak{}[nlm\_, nln\_]), (\allowbreak{}pubkey(\allowbreak{}nlp\_))(\allowbreak{}nla\_))], [sent, encr(\allowbreak{}nln\_, (\allowbreak{}pubkey(\allowbreak{}nlp\_))(\allowbreak{}nlb\_))]])
\end{quote}
mentioned by 4 result(s)

\smallskip

\noindent\textbf{\texttt{is-nso-resp-strand(nlp\_, nls\_, nla\_, nlb\_, nlm\_, nln\_)}}\quad\textit{predicate}\index{is-nso-resp-strand(nlp_, nls_, nla_, nlb_, nlm_, nln_)@\texttt{is-nso-resp-strand(nlp\_, nls\_, nla\_, nlb\_, nlm\_, nln\_)}}\\
reads: \$2 is an original Needham-Schroeder responder's strand of \$1 for the names \$3, \$4 and the values \$5, \$6 \\
defined by \texttt{is-nso-resp-strand} \\
characterized by \texttt{is-nso-resp-strand}  \textit{(definitional)}:
\begin{quote}\ttfamily\footnotesize\raggedright\noindent
forall(\allowbreak{}[nlp\_, nls\_, nla\_, nlb\_, nlm\_, nln\_], is-nso-resp-strand(\allowbreak{}nlp\_, nls\_, nla\_, nlb\_, nlm\_, nln\_) iff nla\_ in names(\allowbreak{}nlp\_) and nlb\_ in names(\allowbreak{}nlp\_) and nlm\_ in tokens(\allowbreak{}ispace(\allowbreak{}pk-inf(\allowbreak{}nlp\_))) and nln\_ in tokens(\allowbreak{}ispace(\allowbreak{}pk-inf(\allowbreak{}nlp\_))) and not(\allowbreak{}nln\_ in names(\allowbreak{}nlp\_)) and not(\allowbreak{}nlm\_ in keys(\allowbreak{}ispace(\allowbreak{}pk-inf(\allowbreak{}nlp\_)))) and not(\allowbreak{}nln\_ in keys(\allowbreak{}ispace(\allowbreak{}pk-inf(\allowbreak{}nlp\_)))) and nls\_ = [[received, encr(\allowbreak{}lst(\allowbreak{}[nlm\_, nla\_]), (\allowbreak{}pubkey(\allowbreak{}nlp\_))(\allowbreak{}nlb\_))], [sent, encr(\allowbreak{}lst(\allowbreak{}[nlm\_, nln\_]), (\allowbreak{}pubkey(\allowbreak{}nlp\_))(\allowbreak{}nla\_))], [received, encr(\allowbreak{}nln\_, (\allowbreak{}pubkey(\allowbreak{}nlp\_))(\allowbreak{}nlb\_))]])
\end{quote}
mentioned by 7 result(s)

\smallskip

\noindent\textbf{\texttt{is-nso-space(nlp\_)}}\quad\textit{predicate}\index{is-nso-space(nlp_)@\texttt{is-nso-space(nlp\_)}}\\
reads: \$1 is an original Needham-Schroeder space \\
defined by \texttt{is-nso-space} \\
characterized by \texttt{is-nso-space}  \textit{(definitional)}:
\begin{quote}\ttfamily\footnotesize\raggedright\noindent
forall(\allowbreak{}[nlp\_], is-nso-space(\allowbreak{}nlp\_) iff is-pk-space(\allowbreak{}nlp\_) and forall(\allowbreak{}[nli\_ in ids(\allowbreak{}ispace(\allowbreak{}pk-inf(\allowbreak{}nlp\_)))], nli\_ in adv-ids(\allowbreak{}pk-inf(\allowbreak{}nlp\_)) or forsome(\allowbreak{}[nla\_, nlb\_, nlm\_, nln\_], is-nso-init-strand(\allowbreak{}nlp\_, (\allowbreak{}strand-map(\allowbreak{}ispace(\allowbreak{}pk-inf(\allowbreak{}nlp\_))))(\allowbreak{}nli\_), nla\_, nlb\_, nlm\_, nln\_)) or forsome(\allowbreak{}[nla\_, nlb\_, nlm\_, nln\_], is-nso-resp-strand(\allowbreak{}nlp\_, (\allowbreak{}strand-map(\allowbreak{}ispace(\allowbreak{}pk-inf(\allowbreak{}nlp\_))))(\allowbreak{}nli\_), nla\_, nlb\_, nlm\_, nln\_))))
\end{quote}
mentioned by 13 result(s)

\smallskip

\noindent\textbf{\texttt{is-nu-c-star-algebra(s)}}\quad\textit{predicate}\index{is-nu-c-star-algebra(s)@\texttt{is-nu-c-star-algebra(s)}}\\
reads: a C*-algebra \\
a predicate \\
characterized by \texttt{is-nu-c-star-algebra-def}  \textit{(definitional)}:
\begin{quote}\ttfamily\footnotesize\raggedright\noindent
forall(\allowbreak{}[s], is-nu-c-star-algebra(\allowbreak{}s) iff is-nu-involutive-banach-algebra(\allowbreak{}s) and forall(\allowbreak{}[x\_ in vec(\allowbreak{}s)], (\allowbreak{}vnrm(\allowbreak{}s))(\allowbreak{}(\allowbreak{}vmul(\allowbreak{}s))(\allowbreak{}(\allowbreak{}istar(\allowbreak{}s))(\allowbreak{}x\_), x\_)) = (\allowbreak{}vnrm(\allowbreak{}s))(\allowbreak{}x\_) * (\allowbreak{}vnrm(\allowbreak{}s))(\allowbreak{}x\_)))
\end{quote}
also: \texttt{is-nu-c-star-algebra-is-set} \\
mentioned by 7 result(s)

\smallskip

\noindent\textbf{\texttt{is-nu-involutive-banach-algebra(s)}}\quad\textit{predicate}\index{is-nu-involutive-banach-algebra(s)@\texttt{is-nu-involutive-banach-algebra(s)}}\\
reads: an involutive Banach algebra \\
a predicate \\
characterized by \texttt{is-nu-involutive-banach-algebra}  \textit{(definitional)}:
\begin{quote}\ttfamily\footnotesize\raggedright\noindent
forall(\allowbreak{}[s], is-nu-involutive-banach-algebra(\allowbreak{}s) iff length(\allowbreak{}s) = 9 and is-ring(\allowbreak{}scal(\allowbreak{}s)) and vec(\allowbreak{}s) in set and vadd(\allowbreak{}s) in fun(\allowbreak{}cartesian(\allowbreak{}vec(\allowbreak{}s), vec(\allowbreak{}s)), vec(\allowbreak{}s)) and vzero(\allowbreak{}s) in vec(\allowbreak{}s) and vneg(\allowbreak{}s) in fun(\allowbreak{}vec(\allowbreak{}s), vec(\allowbreak{}s)) and act(\allowbreak{}s) in fun(\allowbreak{}cartesian(\allowbreak{}carr(\allowbreak{}scal(\allowbreak{}s)), vec(\allowbreak{}s)), vec(\allowbreak{}s)) and vnrm(\allowbreak{}s) in fun(\allowbreak{}vec(\allowbreak{}s), rr) and vmul(\allowbreak{}s) in fun(\allowbreak{}cartesian(\allowbreak{}vec(\allowbreak{}s), vec(\allowbreak{}s)), vec(\allowbreak{}s)) and istar(\allowbreak{}s) in fun(\allowbreak{}vec(\allowbreak{}s), vec(\allowbreak{}s)) and is-complex-banach-algebra(\allowbreak{}[scal(\allowbreak{}s), vec(\allowbreak{}s), vadd(\allowbreak{}s), vzero(\allowbreak{}s), vneg(\allowbreak{}s), act(\allowbreak{}s), vnrm(\allowbreak{}s), vmul(\allowbreak{}s)]) and carr(\allowbreak{}scal(\allowbreak{}s)) = cc and forall(\allowbreak{}[x\_ in vec(\allowbreak{}s), y\_ in vec(\allowbreak{}s)], (\allowbreak{}istar(\allowbreak{}s))(\allowbreak{}(\allowbreak{}vadd(\allowbreak{}s))(\allowbreak{}x\_, y\_)) = (\allowbreak{}vadd(\allowbreak{}s))(\allowbreak{}(\allowbreak{}istar(\allowbreak{}s))(\allowbreak{}x\_), (\allowbreak{}istar(\allowbreak{}s))(\allowbreak{}y\_))) and forall(\allowbreak{}[r\_ in cc, x\_ in vec(\allowbreak{}s)], (\allowbreak{}istar(\allowbreak{}s))(\allowbreak{}(\allowbreak{}act(\allowbreak{}s))(\allowbreak{}r\_, x\_)) = (\allowbreak{}act(\allowbreak{}s))(\allowbreak{}conjugate(\allowbreak{}r\_), (\allowbreak{}istar(\allowbreak{}s))(\allowbreak{}x\_))) and forall(\allowbreak{}[x\_ in vec(\allowbreak{}s)], (\allowbreak{}istar(\allowbreak{}s))(\allowbreak{}(\allowbreak{}istar(\allowbreak{}s))(\allowbreak{}x\_)) = x\_) and forall(\allowbreak{}[x\_ in vec(\allowbreak{}s), y\_ in vec(\allowbreak{}s)], (\allowbreak{}istar(\allowbreak{}s))(\allowbreak{}(\allowbreak{}vmul(\allowbreak{}s))(\allowbreak{}x\_, y\_)) = (\allowbreak{}vmul(\allowbreak{}s))(\allowbreak{}(\allowbreak{}istar(\allowbreak{}s))(\allowbreak{}y\_), (\allowbreak{}istar(\allowbreak{}s))(\allowbreak{}x\_))) and forall(\allowbreak{}[x\_ in vec(\allowbreak{}s)], (\allowbreak{}vnrm(\allowbreak{}s))(\allowbreak{}(\allowbreak{}istar(\allowbreak{}s))(\allowbreak{}x\_)) = (\allowbreak{}vnrm(\allowbreak{}s))(\allowbreak{}x\_)) and s in set)
\end{quote}
also: \texttt{uiba-as-nuiba-is-nu-involutive-banach-algebra} \texttt{is-nu-involutive-banach-algebra-is-set} \\
mentioned by 139 result(s)

\smallskip

\noindent\textbf{\texttt{is-null-set(omega, ca, mu, ns)}}\quad\textit{predicate}\index{is-null-set(omega, ca, mu, ns)@\texttt{is-null-set(omega, ca, mu, ns)}}\\
reads: \$4 is a null set of the measure \$3 on \$2 \\
defined by \texttt{is-null-set} \\
characterized by \texttt{is-null-set}  \textit{(definitional)}:
\begin{quote}\ttfamily\footnotesize\raggedright\noindent
forall(\allowbreak{}[omega, ca, mu, ns], is-null-set(\allowbreak{}omega, ca, mu, ns) iff forsome(\allowbreak{}[nla\_ in ca], ns subset nla\_ and mu(\allowbreak{}nla\_) = 0))
\end{quote}
mentioned by 53 result(s)

\smallskip

\noindent\textbf{\texttt{is-nv-content(cw3k\_, cw3p\_, cw3c\_)}}\quad\textit{predicate}\index{is-nv-content(cw3k_, cw3p_, cw3c_)@\texttt{is-nv-content(cw3k\_, cw3p\_, cw3c\_)}}\\
reads: \$3 is the content of the polynomial \$2 over \$1 \\
defined by \texttt{is-nv-content} \\
characterized by \texttt{is-nv-content}  \textit{(definitional)}:
\begin{quote}\ttfamily\footnotesize\raggedright\noindent
forall(\allowbreak{}[cw3k\_, cw3p\_, cw3c\_], is-nv-content(\allowbreak{}cw3k\_, cw3p\_, cw3c\_) iff forall(\allowbreak{}[cw3i\_ in nn], (\allowbreak{}fnrm(\allowbreak{}cw3k\_))(\allowbreak{}cw3p\_(\allowbreak{}cw3i\_)) <= cw3c\_) and forsome(\allowbreak{}[cw3i\_ in nn], (\allowbreak{}fnrm(\allowbreak{}cw3k\_))(\allowbreak{}cw3p\_(\allowbreak{}cw3i\_)) = cw3c\_))
\end{quote}
mentioned by 7 result(s)

\smallskip

\noindent\textbf{\texttt{is-nv-content-one(cw3k\_, cw3p\_)}}\quad\textit{predicate}\index{is-nv-content-one(cw3k_, cw3p_)@\texttt{is-nv-content-one(cw3k\_, cw3p\_)}}\\
reads: the polynomial \$2 over \$1 has content 1 \\
defined by \texttt{is-nv-content-one} \\
characterized by \texttt{is-nv-content-one}  \textit{(definitional)}:
\begin{quote}\ttfamily\footnotesize\raggedright\noindent
forall(\allowbreak{}[cw3k\_, cw3p\_], is-nv-content-one(\allowbreak{}cw3k\_, cw3p\_) iff forall(\allowbreak{}[cw3i\_ in nn], (\allowbreak{}fnrm(\allowbreak{}cw3k\_))(\allowbreak{}cw3p\_(\allowbreak{}cw3i\_)) <= 1) and forsome(\allowbreak{}[cw3i\_ in nn], (\allowbreak{}fnrm(\allowbreak{}cw3k\_))(\allowbreak{}cw3p\_(\allowbreak{}cw3i\_)) = 1))
\end{quote}
mentioned by 3 result(s)

\smallskip

\noindent\textbf{\texttt{is-nv-domain(cwk\_, cwa\_)}}\quad\textit{predicate}\index{is-nv-domain(cwk_, cwa_)@\texttt{is-nv-domain(cwk\_, cwa\_)}}\\
reads: \$2 is a non-archimedean valued domain inside the field \$1 \\
defined by \texttt{is-nv-domain} \\
characterized by \texttt{is-nv-domain}  \textit{(definitional)}:
\begin{quote}\ttfamily\footnotesize\raggedright\noindent
forall(\allowbreak{}[cwk\_, cwa\_], is-nv-domain(\allowbreak{}cwk\_, cwa\_) iff is-nv-field(\allowbreak{}cwk\_) and forall(\allowbreak{}[cwx\_ in cwa\_], cwx\_ in carr(\allowbreak{}cwk\_)) and zero(\allowbreak{}cwk\_) in cwa\_ and one(\allowbreak{}cwk\_) in cwa\_ and forall(\allowbreak{}[cwx\_ in cwa\_, cwy\_ in cwa\_], (\allowbreak{}add(\allowbreak{}cwk\_))(\allowbreak{}cwx\_, cwy\_) in cwa\_) and forall(\allowbreak{}[cwx\_ in cwa\_, cwy\_ in cwa\_], (\allowbreak{}mul(\allowbreak{}cwk\_))(\allowbreak{}cwx\_, cwy\_) in cwa\_) and forall(\allowbreak{}[cwx\_ in cwa\_], (\allowbreak{}neg(\allowbreak{}cwk\_))(\allowbreak{}cwx\_) in cwa\_) and forall(\allowbreak{}[cwx\_ in cwa\_], (\allowbreak{}fnrm(\allowbreak{}cwk\_))(\allowbreak{}cwx\_) <= 1))
\end{quote}
mentioned by 84 result(s)

\smallskip

\noindent\textbf{\texttt{is-nv-expansion(nek\_, new\_, ner\_, nea\_, nen\_, nex\_)}}\quad\textit{predicate}\index{is-nv-expansion(nek_, new_, ner_, nea_, nen_, nex_)@\texttt{is-nv-expansion(nek\_, new\_, ner\_, nea\_, nen\_, nex\_)}}\\
reads: \$6 = sum over k >= \$5 of \$4(k) \$2\textasciicircum{}k with digits in \$3 \\
defined by \texttt{is-nv-expansion} \\
characterized by \texttt{is-nv-expansion}  \textit{(definitional)}:
\begin{quote}\ttfamily\footnotesize\raggedright\noindent
forall(\allowbreak{}[nek\_, new\_, ner\_, nea\_, nen\_, nex\_], is-nv-expansion(\allowbreak{}nek\_, new\_, ner\_, nea\_, nen\_, nex\_) iff nea\_ in fun(\allowbreak{}zz, ner\_) and forall(\allowbreak{}[nei\_ in zz], nei\_ < nen\_ implies nea\_(\allowbreak{}nei\_) = zero(\allowbreak{}nek\_)) and converges-to(\allowbreak{}nf-metric-space(\allowbreak{}nek\_), vnb-lambda(\allowbreak{}nem\_, nn, nv-psum(\allowbreak{}nek\_, new\_, nea\_, nen\_, nem\_)), nex\_))
\end{quote}
mentioned by 17 result(s)

\smallskip

\noindent\textbf{\texttt{is-nv-field(cwk\_)}}\quad\textit{predicate}\index{is-nv-field(cwk_)@\texttt{is-nv-field(cwk\_)}}\\
reads: \$1 is a non-archimedean valued field \\
defined by \texttt{is-nv-field} \\
characterized by \texttt{is-nv-field}  \textit{(definitional)}:
\begin{quote}\ttfamily\footnotesize\raggedright\noindent
forall(\allowbreak{}[cwk\_], is-nv-field(\allowbreak{}cwk\_) iff is-normed-field(\allowbreak{}cwk\_) and forall(\allowbreak{}[cwx\_ in carr(\allowbreak{}cwk\_), cwy\_ in carr(\allowbreak{}cwk\_)], (\allowbreak{}fnrm(\allowbreak{}cwk\_))(\allowbreak{}(\allowbreak{}add(\allowbreak{}cwk\_))(\allowbreak{}cwx\_, cwy\_)) <= max(\allowbreak{}(\allowbreak{}fnrm(\allowbreak{}cwk\_))(\allowbreak{}cwx\_), (\allowbreak{}fnrm(\allowbreak{}cwk\_))(\allowbreak{}cwy\_))))
\end{quote}
also: \texttt{qp-is-nv-field} \texttt{sub-nf-is-nv-field} \texttt{qq-padic-is-nv-field} \texttt{laurent-nf-is-nv-field} \\
mentioned by 172 result(s)

\smallskip

\noindent\textbf{\texttt{is-nv-ideal(cwk\_, cwa\_, cwi\_)}}\quad\textit{predicate}\index{is-nv-ideal(cwk_, cwa_, cwi_)@\texttt{is-nv-ideal(cwk\_, cwa\_, cwi\_)}}\\
reads: \$3 is an ideal of \$2 \\
defined by \texttt{is-nv-ideal} \\
characterized by \texttt{is-nv-ideal}  \textit{(definitional)}:
\begin{quote}\ttfamily\footnotesize\raggedright\noindent
forall(\allowbreak{}[cwk\_, cwa\_, cwi\_], is-nv-ideal(\allowbreak{}cwk\_, cwa\_, cwi\_) iff forall(\allowbreak{}[cwx\_ in cwi\_], cwx\_ in cwa\_) and zero(\allowbreak{}cwk\_) in cwi\_ and forall(\allowbreak{}[cwx\_ in cwi\_, cwy\_ in cwi\_], (\allowbreak{}add(\allowbreak{}cwk\_))(\allowbreak{}cwx\_, cwy\_) in cwi\_) and forall(\allowbreak{}[cwx\_ in cwi\_], (\allowbreak{}neg(\allowbreak{}cwk\_))(\allowbreak{}cwx\_) in cwi\_) and forall(\allowbreak{}[cwr\_ in cwa\_, cwx\_ in cwi\_], (\allowbreak{}mul(\allowbreak{}cwk\_))(\allowbreak{}cwr\_, cwx\_) in cwi\_))
\end{quote}
mentioned by 16 result(s)

\smallskip

\noindent\textbf{\texttt{is-nv-module(cw3k\_, cw3a\_, cw3e\_)}}\quad\textit{predicate}\index{is-nv-module(cw3k_, cw3a_, cw3e_)@\texttt{is-nv-module(cw3k\_, cw3a\_, cw3e\_)}}\\
reads: \$3 is a module over \$2 in the field \$1 \\
defined by \texttt{is-nv-module} \\
characterized by \texttt{is-nv-module}  \textit{(definitional)}:
\begin{quote}\ttfamily\footnotesize\raggedright\noindent
forall(\allowbreak{}[cw3k\_, cw3a\_, cw3e\_], is-nv-module(\allowbreak{}cw3k\_, cw3a\_, cw3e\_) iff forall(\allowbreak{}[cw3x\_ in cw3e\_], cw3x\_ in carr(\allowbreak{}cw3k\_)) and zero(\allowbreak{}cw3k\_) in cw3e\_ and forall(\allowbreak{}[cw3x\_ in cw3e\_, cw3y\_ in cw3e\_], (\allowbreak{}add(\allowbreak{}cw3k\_))(\allowbreak{}cw3x\_, cw3y\_) in cw3e\_) and forall(\allowbreak{}[cw3x\_ in cw3e\_], (\allowbreak{}neg(\allowbreak{}cw3k\_))(\allowbreak{}cw3x\_) in cw3e\_) and forall(\allowbreak{}[cw3r\_ in cw3a\_, cw3x\_ in cw3e\_], (\allowbreak{}mul(\allowbreak{}cw3k\_))(\allowbreak{}cw3r\_, cw3x\_) in cw3e\_))
\end{quote}
mentioned by 12 result(s)

\smallskip

\noindent\textbf{\texttt{is-nv-pid(cwk\_, cwa\_)}}\quad\textit{predicate}\index{is-nv-pid(cwk_, cwa_)@\texttt{is-nv-pid(cwk\_, cwa\_)}}\\
reads: \$2 is a principal ideal domain \\
defined by \texttt{is-nv-pid} \\
characterized by \texttt{is-nv-pid}  \textit{(definitional)}:
\begin{quote}\ttfamily\footnotesize\raggedright\noindent
forall(\allowbreak{}[cwk\_, cwa\_], is-nv-pid(\allowbreak{}cwk\_, cwa\_) iff is-nv-domain(\allowbreak{}cwk\_, cwa\_) and forall(\allowbreak{}[cwi\_], is-nv-ideal(\allowbreak{}cwk\_, cwa\_, cwi\_) implies forsome(\allowbreak{}[cwg\_ in cwa\_], cwi\_ = nv-principal(\allowbreak{}cwk\_, cwa\_, cwg\_))))
\end{quote}
mentioned by 13 result(s)

\smallskip

\noindent\textbf{\texttt{is-nv-rep-set(rpk\_, rpa\_, rpr\_)}}\quad\textit{predicate}\index{is-nv-rep-set(rpk_, rpa_, rpr_)@\texttt{is-nv-rep-set(rpk\_, rpa\_, rpr\_)}}\\
reads: \$3 is a representative set for \$2 modulo its maximal ideal in \$1 \\
defined by \texttt{is-nv-rep-set} \\
characterized by \texttt{is-nv-rep-set}  \textit{(definitional)}:
\begin{quote}\ttfamily\footnotesize\raggedright\noindent
forall(\allowbreak{}[rpk\_, rpa\_, rpr\_], is-nv-rep-set(\allowbreak{}rpk\_, rpa\_, rpr\_) iff forall(\allowbreak{}[rpx\_ in rpr\_], rpx\_ in rpa\_) and zero(\allowbreak{}rpk\_) in rpr\_ and one(\allowbreak{}rpk\_) in rpr\_ and forall(\allowbreak{}[rpx\_ in rpa\_], forsome(\allowbreak{}[rpy\_ in rpr\_], (\allowbreak{}fnrm(\allowbreak{}rpk\_))(\allowbreak{}(\allowbreak{}add(\allowbreak{}rpk\_))(\allowbreak{}rpx\_, (\allowbreak{}neg(\allowbreak{}rpk\_))(\allowbreak{}rpy\_))) < 1)) and forall(\allowbreak{}[rpx\_ in rpr\_, rpy\_ in rpr\_], (\allowbreak{}fnrm(\allowbreak{}rpk\_))(\allowbreak{}(\allowbreak{}add(\allowbreak{}rpk\_))(\allowbreak{}rpx\_, (\allowbreak{}neg(\allowbreak{}rpk\_))(\allowbreak{}rpy\_))) < 1 implies rpx\_ = rpy\_))
\end{quote}
mentioned by 37 result(s)

\smallskip

\noindent\textbf{\texttt{is-nv-valuation(x4f\_, x4nu\_)}}\quad\textit{predicate}\index{is-nv-valuation(x4f_, x4nu_)@\texttt{is-nv-valuation(x4f\_, x4nu\_)}}\\
reads: \$2 is a non-archimedean valuation of the field \$1 \\
defined by \texttt{is-nv-valuation} \\
characterized by \texttt{is-nv-valuation}  \textit{(definitional)}:
\begin{quote}\ttfamily\footnotesize\raggedright\noindent
forall(\allowbreak{}[x4f\_, x4nu\_], is-nv-valuation(\allowbreak{}x4f\_, x4nu\_) iff x4nu\_ in fun(\allowbreak{}carr(\allowbreak{}x4f\_), rr) and forall(\allowbreak{}[x4x\_ in carr(\allowbreak{}x4f\_)], 0 <= x4nu\_(\allowbreak{}x4x\_)) and forall(\allowbreak{}[x4x\_ in carr(\allowbreak{}x4f\_)], x4nu\_(\allowbreak{}x4x\_) = 0 iff x4x\_ = zero(\allowbreak{}x4f\_)) and forall(\allowbreak{}[x4x\_ in carr(\allowbreak{}x4f\_), x4y\_ in carr(\allowbreak{}x4f\_)], x4nu\_(\allowbreak{}(\allowbreak{}mul(\allowbreak{}x4f\_))(\allowbreak{}x4x\_, x4y\_)) = x4nu\_(\allowbreak{}x4x\_) * x4nu\_(\allowbreak{}x4y\_)) and forall(\allowbreak{}[x4x\_ in carr(\allowbreak{}x4f\_), x4y\_ in carr(\allowbreak{}x4f\_)], x4nu\_(\allowbreak{}(\allowbreak{}add(\allowbreak{}x4f\_))(\allowbreak{}x4x\_, x4y\_)) <= max(\allowbreak{}x4nu\_(\allowbreak{}x4x\_), x4nu\_(\allowbreak{}x4y\_))))
\end{quote}
mentioned by 5 result(s)

\smallskip

\noindent\textbf{\texttt{is-one-parameter-group(ops\_, opg\_, opp\_)}}\quad\textit{predicate}\index{is-one-parameter-group(ops_, opg_, opp_)@\texttt{is-one-parameter-group(ops\_, opg\_, opp\_)}}\\
reads: \$3 is a one-parameter group of \$2 on \$1 \\
defined by \texttt{is-one-parameter-group} \\
characterized by \texttt{is-one-parameter-group}  \textit{(definitional)}:
\begin{quote}\ttfamily\footnotesize\raggedright\noindent
forall(\allowbreak{}[ops\_, opg\_, opp\_], is-one-parameter-group(\allowbreak{}ops\_, opg\_, opp\_) iff opp\_ in fun(\allowbreak{}rr, carr(\allowbreak{}opg\_)) and forall(\allowbreak{}[oqr\_ in rr, oqu\_ in rr], opp\_(\allowbreak{}oqr\_ + oqu\_) = (\allowbreak{}opr(\allowbreak{}opg\_))(\allowbreak{}opp\_(\allowbreak{}oqr\_), opp\_(\allowbreak{}oqu\_))) and forall(\allowbreak{}[oqt\_ in rr, oqw\_ in opens(\allowbreak{}ops\_)], opp\_(\allowbreak{}oqt\_) in oqw\_ implies forsome(\allowbreak{}[oqd\_], pos-rr(\allowbreak{}oqd\_) and forall(\allowbreak{}[oqv\_ in rr], abs(\allowbreak{}oqv\_ - oqt\_) < oqd\_ implies opp\_(\allowbreak{}oqv\_) in oqw\_))))
\end{quote}
mentioned by 8 result(s)

\smallskip

\noindent\textbf{\texttt{is-open(s, u)}}\quad\textit{predicate}\index{is-open(s, u)@\texttt{is-open(s, u)}}\\
reads: \$2 is open in \$1 \\
defined by \texttt{is-open} \\
characterized by \texttt{is-open}  \textit{(definitional)}:
\begin{quote}\ttfamily\footnotesize\raggedright\noindent
forall(\allowbreak{}[s, u], is-open(\allowbreak{}s, u) iff is-metric-space(\allowbreak{}s) and u subset pts(\allowbreak{}s) and forall(\allowbreak{}[y in u], forsome(\allowbreak{}[r], pos-rr(\allowbreak{}r) and ball(\allowbreak{}s, y, r) subset u)))
\end{quote}
also: \texttt{is-open-cover} \texttt{ball-is-open} \texttt{empty-is-open} \texttt{carrier-is-open} \texttt{interior-is-open} \texttt{is-open-in-power} \\
mentioned by 383 result(s)

\smallskip

\noindent\textbf{\texttt{is-open-cover(s, c)}}\quad\textit{predicate}\index{is-open-cover(s, c)@\texttt{is-open-cover(s, c)}}\\
reads: \$2 is an open cover of \$1 \\
defined by \texttt{is-open-cover} \\
characterized by \texttt{is-open-cover}  \textit{(definitional)}:
\begin{quote}\ttfamily\footnotesize\raggedright\noindent
forall(\allowbreak{}[s, c], is-open-cover(\allowbreak{}s, c) iff is-metric-space(\allowbreak{}s) and forall(\allowbreak{}[u in c], is-open(\allowbreak{}s, u)) and big-union(\allowbreak{}u, c, u) == pts(\allowbreak{}s))
\end{quote}
also: \texttt{ball-cover-is-open-cover} \\
mentioned by 10 result(s)

\smallskip

\noindent\textbf{\texttt{is-open-gauge(fam, ground, u)}}\quad\textit{predicate}\index{is-open-gauge(fam, ground, u)@\texttt{is-open-gauge(fam, ground, u)}}\\
reads: \$3 is open in the gauge topology of the family \$1 on \$2 \\
defined by \texttt{is-open-gauge} \\
characterized by \texttt{is-open-gauge}  \textit{(definitional)}:
\begin{quote}\ttfamily\footnotesize\raggedright\noindent
forall(\allowbreak{}[fam, ground, u], is-open-gauge(\allowbreak{}fam, ground, u) iff u subset ground and forall(\allowbreak{}[y in u], forsome(\allowbreak{}[k\_ in nn, r], pos-rr(\allowbreak{}r) and forall(\allowbreak{}[x in ground], forall(\allowbreak{}[n\_ in nn], n\_ <= k\_ implies (\allowbreak{}dist(\allowbreak{}fam(\allowbreak{}n\_)))(\allowbreak{}y, x) < r) implies x in u))))
\end{quote}
mentioned by 13 result(s)

\smallskip

\noindent\textbf{\texttt{is-open-lin-map(m1, fam1, m2, fam2, tt)}}\quad\textit{predicate}\index{is-open-lin-map(m1, fam1, m2, fam2, tt)@\texttt{is-open-lin-map(m1, fam1, m2, fam2, tt)}}\\
reads: \$5 is an open linear map from (\$1, \$2) to (\$3, \$4) \\
defined by \texttt{is-open-lin-map} \\
characterized by \texttt{is-open-lin-map}  \textit{(definitional)}:
\begin{quote}\ttfamily\footnotesize\raggedright\noindent
forall(\allowbreak{}[m1, fam1, m2, fam2, tt], is-open-lin-map(\allowbreak{}m1, fam1, m2, fam2, tt) iff is-linear-map(\allowbreak{}m1, m2, tt) and forall(\allowbreak{}[n in nn, eps], pos-rr(\allowbreak{}eps) implies forsome(\allowbreak{}[n\_ in nn, eps\_], pos-rr(\allowbreak{}eps\_) and fr-ball(\allowbreak{}m2, fam2, n\_, eps\_) subset image(\allowbreak{}tt, fr-ball(\allowbreak{}m1, fam1, n, eps)))))
\end{quote}
mentioned by 4 result(s)

\smallskip

\noindent\textbf{\texttt{is-or-init(orip\_, oris\_, oria\_, orib\_, orin\_, orim\_, orik\_)}}\quad\textit{predicate}\index{is-or-init(orip_, oris_, oria_, orib_, orin_, orim_, orik_)@\texttt{is-or-init(orip\_, oris\_, oria\_, orib\_, orin\_, orim\_, orik\_)}}\\
reads: \$2 is an Otway-Rees initiator's strand of \$1 for \$3, \$4, \$5, \$6, \$7 \\
defined by \texttt{is-or-init} \\
characterized by \texttt{is-or-init}  \textit{(definitional)}:
\begin{quote}\ttfamily\footnotesize\raggedright\noindent
forall(\allowbreak{}[orip\_, oris\_, oria\_, orib\_, orin\_, orim\_, orik\_], is-or-init(\allowbreak{}orip\_, oris\_, oria\_, orib\_, orin\_, orim\_, orik\_) iff oria\_ in sk-names(\allowbreak{}orip\_) and orib\_ in sk-names(\allowbreak{}orip\_) and is-or-text(\allowbreak{}orip\_, orin\_) and is-or-text(\allowbreak{}orip\_, orim\_) and orik\_ in keys(\allowbreak{}ispace(\allowbreak{}sk-inf(\allowbreak{}orip\_))) and oris\_ = [[sent, lst(\allowbreak{}[orim\_, oria\_, orib\_, encr(\allowbreak{}lst(\allowbreak{}[orin\_, orim\_, oria\_, orib\_]), (\allowbreak{}ltkey(\allowbreak{}orip\_))(\allowbreak{}oria\_))])], [received, lst(\allowbreak{}[orim\_, encr(\allowbreak{}lst(\allowbreak{}[orin\_, orik\_]), (\allowbreak{}ltkey(\allowbreak{}orip\_))(\allowbreak{}oria\_))])]])
\end{quote}
mentioned by 15 result(s)

\smallskip

\noindent\textbf{\texttt{is-or-resp(orrp\_, orrs\_, orra\_, orrb\_, orrn\_, orrm\_, orrk\_, orrh\_, orrg\_)}}\quad\textit{predicate}\index{is-or-resp(orrp_, orrs_, orra_, orrb_, orrn_, orrm_, orrk_, orrh_, orrg_)@\texttt{is-or-resp(orrp\_, orrs\_, orra\_, orrb\_, orrn\_, orrm\_, orrk\_, orrh\_, orrg\_)}}\\
reads: \$2 is an Otway-Rees responder's strand of \$1 for \$3, \$4, \$5, \$6, \$7, \$8, \$9 \\
defined by \texttt{is-or-resp} \\
characterized by \texttt{is-or-resp}  \textit{(definitional)}:
\begin{quote}\ttfamily\footnotesize\raggedright\noindent
forall(\allowbreak{}[orrp\_, orrs\_, orra\_, orrb\_, orrn\_, orrm\_, orrk\_, orrh\_, orrg\_], is-or-resp(\allowbreak{}orrp\_, orrs\_, orra\_, orrb\_, orrn\_, orrm\_, orrk\_, orrh\_, orrg\_) iff orra\_ in sk-names(\allowbreak{}orrp\_) and orrb\_ in sk-names(\allowbreak{}orrp\_) and is-or-text(\allowbreak{}orrp\_, orrn\_) and is-or-text(\allowbreak{}orrp\_, orrm\_) and orrk\_ in keys(\allowbreak{}ispace(\allowbreak{}sk-inf(\allowbreak{}orrp\_))) and orrh\_ in msg-set(\allowbreak{}tokens(\allowbreak{}ispace(\allowbreak{}sk-inf(\allowbreak{}orrp\_))), keys(\allowbreak{}ispace(\allowbreak{}sk-inf(\allowbreak{}orrp\_)))) and orrg\_ in msg-set(\allowbreak{}tokens(\allowbreak{}ispace(\allowbreak{}sk-inf(\allowbreak{}orrp\_))), keys(\allowbreak{}ispace(\allowbreak{}sk-inf(\allowbreak{}orrp\_)))) and not(\allowbreak{}orrn\_ = orrm\_) and not(\allowbreak{}is-subterm(\allowbreak{}tokens(\allowbreak{}ispace(\allowbreak{}sk-inf(\allowbreak{}orrp\_))), keys(\allowbreak{}ispace(\allowbreak{}sk-inf(\allowbreak{}orrp\_))), orrn\_, orrh\_)) and orrs\_ = [[received, lst(\allowbreak{}[orrm\_, orra\_, orrb\_, orrh\_])], [sent, lst(\allowbreak{}[orrm\_, orra\_, orrb\_, orrh\_, encr(\allowbreak{}lst(\allowbreak{}[orrn\_, orrm\_, orra\_, orrb\_]), (\allowbreak{}ltkey(\allowbreak{}orrp\_))(\allowbreak{}orrb\_))])], [received, lst(\allowbreak{}[orrm\_, orrg\_, encr(\allowbreak{}lst(\allowbreak{}[orrn\_, orrk\_]), (\allowbreak{}ltkey(\allowbreak{}orrp\_))(\allowbreak{}orrb\_))])], [sent, lst(\allowbreak{}[orrm\_, orrg\_])]])
\end{quote}
mentioned by 15 result(s)

\smallskip

\noindent\textbf{\texttt{is-or-serv(orsp\_, orss\_, orsa\_, orsb\_, orsn\_, orso\_, orsm\_, orsk\_)}}\quad\textit{predicate}\index{is-or-serv(orsp_, orss_, orsa_, orsb_, orsn_, orso_, orsm_, orsk_)@\texttt{is-or-serv(orsp\_, orss\_, orsa\_, orsb\_, orsn\_, orso\_, orsm\_, orsk\_)}}\\
reads: \$2 is an Otway-Rees server's strand of \$1 for \$3, \$4, \$5, \$6, \$7, \$8 \\
defined by \texttt{is-or-serv} \\
characterized by \texttt{is-or-serv}  \textit{(definitional)}:
\begin{quote}\ttfamily\footnotesize\raggedright\noindent
forall(\allowbreak{}[orsp\_, orss\_, orsa\_, orsb\_, orsn\_, orso\_, orsm\_, orsk\_], is-or-serv(\allowbreak{}orsp\_, orss\_, orsa\_, orsb\_, orsn\_, orso\_, orsm\_, orsk\_) iff orsa\_ in sk-names(\allowbreak{}orsp\_) and orsb\_ in sk-names(\allowbreak{}orsp\_) and is-or-text(\allowbreak{}orsp\_, orsn\_) and is-or-text(\allowbreak{}orsp\_, orso\_) and is-or-text(\allowbreak{}orsp\_, orsm\_) and orsk\_ in keys(\allowbreak{}ispace(\allowbreak{}sk-inf(\allowbreak{}orsp\_))) and not(\allowbreak{}orsk\_ in adv-tokens(\allowbreak{}sk-inf(\allowbreak{}orsp\_))) and forall(\allowbreak{}[orsx\_ in sk-names(\allowbreak{}orsp\_)], not(\allowbreak{}orsk\_ = (\allowbreak{}ltkey(\allowbreak{}orsp\_))(\allowbreak{}orsx\_))) and (\allowbreak{}kinv(\allowbreak{}ispace(\allowbreak{}sk-inf(\allowbreak{}orsp\_))))(\allowbreak{}orsk\_) = orsk\_ and orss\_ = [[received, lst(\allowbreak{}[orsm\_, orsa\_, orsb\_, encr(\allowbreak{}lst(\allowbreak{}[orsn\_, orsm\_, orsa\_, orsb\_]), (\allowbreak{}ltkey(\allowbreak{}orsp\_))(\allowbreak{}orsa\_)), encr(\allowbreak{}lst(\allowbreak{}[orso\_, orsm\_, orsa\_, orsb\_]), (\allowbreak{}ltkey(\allowbreak{}orsp\_))(\allowbreak{}orsb\_))])], [sent, lst(\allowbreak{}[orsm\_, encr(\allowbreak{}lst(\allowbreak{}[orsn\_, orsk\_]), (\allowbreak{}ltkey(\allowbreak{}orsp\_))(\allowbreak{}orsa\_)), encr(\allowbreak{}lst(\allowbreak{}[orso\_, orsk\_]), (\allowbreak{}ltkey(\allowbreak{}orsp\_))(\allowbreak{}orsb\_))])]])
\end{quote}
mentioned by 13 result(s)

\smallskip

\noindent\textbf{\texttt{is-or-space(orqp\_)}}\quad\textit{predicate}\index{is-or-space(orqp_)@\texttt{is-or-space(orqp\_)}}\\
reads: \$1 is an Otway-Rees strand space \\
defined by \texttt{is-or-space} \\
characterized by \texttt{is-or-space}  \textit{(definitional)}:
\begin{quote}\ttfamily\footnotesize\raggedright\noindent
forall(\allowbreak{}[orqp\_], is-or-space(\allowbreak{}orqp\_) iff is-sk-space(\allowbreak{}orqp\_) and forall(\allowbreak{}[orqi\_ in ids(\allowbreak{}ispace(\allowbreak{}sk-inf(\allowbreak{}orqp\_)))], orqi\_ in adv-ids(\allowbreak{}sk-inf(\allowbreak{}orqp\_)) or forsome(\allowbreak{}[orqa\_, orqb\_, orqn\_, orqm\_, orqk\_], is-or-init(\allowbreak{}orqp\_, (\allowbreak{}strand-map(\allowbreak{}ispace(\allowbreak{}sk-inf(\allowbreak{}orqp\_))))(\allowbreak{}orqi\_), orqa\_, orqb\_, orqn\_, orqm\_, orqk\_)) or forsome(\allowbreak{}[orqa\_, orqb\_, orqn\_, orqm\_, orqk\_, orqh\_, orqg\_], is-or-resp(\allowbreak{}orqp\_, (\allowbreak{}strand-map(\allowbreak{}ispace(\allowbreak{}sk-inf(\allowbreak{}orqp\_))))(\allowbreak{}orqi\_), orqa\_, orqb\_, orqn\_, orqm\_, orqk\_, orqh\_, orqg\_)) or forsome(\allowbreak{}[orqa\_, orqb\_, orqn\_, orqo\_, orqm\_, orqk\_], is-or-serv(\allowbreak{}orqp\_, (\allowbreak{}strand-map(\allowbreak{}ispace(\allowbreak{}sk-inf(\allowbreak{}orqp\_))))(\allowbreak{}orqi\_), orqa\_, orqb\_, orqn\_, orqo\_, orqm\_, orqk\_))))
\end{quote}
also: \texttt{cx-orp-is-or-space} \\
mentioned by 16 result(s)

\smallskip

\noindent\textbf{\texttt{is-or-text(ortp\_, ortx\_)}}\quad\textit{predicate}\index{is-or-text(ortp_, ortx_)@\texttt{is-or-text(ortp\_, ortx\_)}}\\
reads: \$2 is a text (neither a name nor a key) of \$1 \\
defined by \texttt{is-or-text} \\
characterized by \texttt{is-or-text}  \textit{(definitional)}:
\begin{quote}\ttfamily\footnotesize\raggedright\noindent
forall(\allowbreak{}[ortp\_, ortx\_], is-or-text(\allowbreak{}ortp\_, ortx\_) iff ortx\_ in tokens(\allowbreak{}ispace(\allowbreak{}sk-inf(\allowbreak{}ortp\_))) and not(\allowbreak{}ortx\_ in sk-names(\allowbreak{}ortp\_)) and not(\allowbreak{}ortx\_ in keys(\allowbreak{}ispace(\allowbreak{}sk-inf(\allowbreak{}ortp\_)))))
\end{quote}
mentioned by 9 result(s)

\smallskip

\noindent\textbf{\texttt{is-ord(x)}}\quad\textit{predicate}\index{is-ord(x)@\texttt{is-ord(x)}}\\
reads: \$1 is an ordinal \\
a predicate \\
characterized by \texttt{is-ord}  \textit{(definitional)}:
\begin{quote}\ttfamily\footnotesize\raggedright\noindent
forall(\allowbreak{}[x], is-ord(\allowbreak{}x) iff x in ord)
\end{quote}
mentioned by 2 result(s)

\smallskip

\noindent\textbf{\texttt{is-order-iso(oisf\_, oisa\_, oisb\_)}}\quad\textit{predicate}\index{is-order-iso(oisf_, oisa_, oisb_)@\texttt{is-order-iso(oisf\_, oisa\_, oisb\_)}}\\
reads: \$1 is an order isomorphism of \$2 onto \$3 \\
defined by \texttt{is-order-iso} \\
characterized by \texttt{is-order-iso}  \textit{(definitional)}:
\begin{quote}\ttfamily\footnotesize\raggedright\noindent
forall(\allowbreak{}[oisf\_, oisa\_, oisb\_], is-order-iso(\allowbreak{}oisf\_, oisa\_, oisb\_) iff oisf\_ in bijection(\allowbreak{}oisa\_, oisb\_) and forall(\allowbreak{}[oisx\_ in oisa\_, oisy\_ in oisa\_], oisx\_ < oisy\_ implies oisf\_(\allowbreak{}oisx\_) < oisf\_(\allowbreak{}oisy\_)))
\end{quote}
mentioned by 4 result(s)

\smallskip

\noindent\textbf{\texttt{is-orthogonal-projection(v, ks, pk)}}\quad\textit{predicate}\index{is-orthogonal-projection(v, ks, pk)@\texttt{is-orthogonal-projection(v, ks, pk)}}\\
reads: \$3 is the orthogonal projection of \$1 onto \$2 \\
defined by \texttt{is-orthogonal-projection} \\
characterized by \texttt{is-orthogonal-projection}  \textit{(definitional)}:
\begin{quote}\ttfamily\footnotesize\raggedright\noindent
forall(\allowbreak{}[v, ks, pk], is-orthogonal-projection(\allowbreak{}v, ks, pk) iff pk in fun(\allowbreak{}vec(\allowbreak{}v), vec(\allowbreak{}v)) and forall(\allowbreak{}[opz\_ in vec(\allowbreak{}v)], pk(\allowbreak{}opz\_) in ks and (\allowbreak{}vadd(\allowbreak{}v))(\allowbreak{}opz\_, (\allowbreak{}vneg(\allowbreak{}v))(\allowbreak{}pk(\allowbreak{}opz\_))) in perp(\allowbreak{}v, ks)))
\end{quote}
also: \texttt{proj-op-is-orthogonal-projection} \\
mentioned by 19 result(s)

\smallskip

\noindent\textbf{\texttt{is-orthonormal-basis(v, ix, e)}}\quad\textit{predicate}\index{is-orthonormal-basis(v, ix, e)@\texttt{is-orthonormal-basis(v, ix, e)}}\\
reads: \$3 is an orthonormal basis of \$1 indexed by \$2 \\
defined by \texttt{is-orthonormal-basis} \\
characterized by \texttt{is-orthonormal-basis}  \textit{(definitional)}:
\begin{quote}\ttfamily\footnotesize\raggedright\noindent
forall(\allowbreak{}[v, ix, e], is-orthonormal-basis(\allowbreak{}v, ix, e) iff is-orthonormal-family(\allowbreak{}v, ix, e) and forall(\allowbreak{}[onx\_ in vec(\allowbreak{}v)], forall(\allowbreak{}[oni\_ in ix], (\allowbreak{}ip(\allowbreak{}v))(\allowbreak{}onx\_, e(\allowbreak{}oni\_)) = 0) implies onx\_ = vzero(\allowbreak{}v)))
\end{quote}
mentioned by 65 result(s)

\smallskip

\noindent\textbf{\texttt{is-orthonormal-family(v, ix, e)}}\quad\textit{predicate}\index{is-orthonormal-family(v, ix, e)@\texttt{is-orthonormal-family(v, ix, e)}}\\
reads: \$3 is an orthonormal family in \$1 indexed by \$2 \\
defined by \texttt{is-orthonormal-family} \\
characterized by \texttt{is-orthonormal-family}  \textit{(definitional)}:
\begin{quote}\ttfamily\footnotesize\raggedright\noindent
forall(\allowbreak{}[v, ix, e], is-orthonormal-family(\allowbreak{}v, ix, e) iff e in fun(\allowbreak{}ix, vec(\allowbreak{}v)) and forall(\allowbreak{}[oni\_ in ix], (\allowbreak{}ip(\allowbreak{}v))(\allowbreak{}e(\allowbreak{}oni\_), e(\allowbreak{}oni\_)) = 1) and forall(\allowbreak{}[oni\_ in ix, onj\_ in ix], not(\allowbreak{}oni\_ = onj\_) implies (\allowbreak{}ip(\allowbreak{}v))(\allowbreak{}e(\allowbreak{}oni\_), e(\allowbreak{}onj\_)) = 0))
\end{quote}
mentioned by 34 result(s)

\smallskip

\noindent\textbf{\texttt{is-orthonormal-set(v, os)}}\quad\textit{predicate}\index{is-orthonormal-set(v, os)@\texttt{is-orthonormal-set(v, os)}}\\
reads: \$2 is an orthonormal set in \$1 \\
defined by \texttt{is-orthonormal-set} \\
characterized by \texttt{is-orthonormal-set}  \textit{(definitional)}:
\begin{quote}\ttfamily\footnotesize\raggedright\noindent
forall(\allowbreak{}[v, os], is-orthonormal-set(\allowbreak{}v, os) iff os subset vec(\allowbreak{}v) and forall(\allowbreak{}[osa\_ in os], (\allowbreak{}ip(\allowbreak{}v))(\allowbreak{}osa\_, osa\_) = 1) and forall(\allowbreak{}[osa\_ in os, osb\_ in os], not(\allowbreak{}osa\_ = osb\_) implies (\allowbreak{}ip(\allowbreak{}v))(\allowbreak{}osa\_, osb\_) = 0))
\end{quote}
mentioned by 8 result(s)

\smallskip

\noindent\textbf{\texttt{is-outer-measure(omega, mu)}}\quad\textit{predicate}\index{is-outer-measure(omega, mu)@\texttt{is-outer-measure(omega, mu)}}\\
reads: \$2 is an outer measure on \$1 \\
defined by \texttt{is-outer-measure} \\
characterized by \texttt{is-outer-measure}  \textit{(definitional)}:
\begin{quote}\ttfamily\footnotesize\raggedright\noindent
forall(\allowbreak{}[omega, mu], is-outer-measure(\allowbreak{}omega, mu) iff omega in set and mu in fun(\allowbreak{}power(\allowbreak{}omega), rr-pos-star) and mu(\allowbreak{}empty-set) = 0 and forall(\allowbreak{}[a\_, b\_], a\_ in power(\allowbreak{}omega) implies b\_ in power(\allowbreak{}omega) implies a\_ subset b\_ implies mu(\allowbreak{}a\_) <= mu(\allowbreak{}b\_)) and forall(\allowbreak{}[f\_ in fun(\allowbreak{}nn, power(\allowbreak{}omega))], mu(\allowbreak{}big-union(\allowbreak{}n\_, nn, f\_(\allowbreak{}n\_))) <= esum(\allowbreak{}vnb-lambda(\allowbreak{}n\_, nn, mu(\allowbreak{}f\_(\allowbreak{}n\_))))))
\end{quote}
also: \texttt{method-i-outer-is-outer-measure} \\
mentioned by 20 result(s)

\smallskip

\noindent\textbf{\texttt{is-partial-isometry(bov\_, bot\_)}}\quad\textit{predicate}\index{is-partial-isometry(bov_, bot_)@\texttt{is-partial-isometry(bov\_, bot\_)}}\\
reads: \$2 is a partial isometry in \$1 \\
defined by \texttt{is-partial-isometry} \\
characterized by \texttt{is-partial-isometry}  \textit{(definitional)}:
\begin{quote}\ttfamily\footnotesize\raggedright\noindent
forall(\allowbreak{}[bov\_, bot\_], is-partial-isometry(\allowbreak{}bov\_, bot\_) iff bot\_ in bounded-ops(\allowbreak{}bov\_, bov\_) and forsome(\allowbreak{}[bok\_], is-submodule(\allowbreak{}bov\_, bok\_) and is-closed(\allowbreak{}ip-metric-space(\allowbreak{}bov\_), bok\_) and is-isometric-on(\allowbreak{}bov\_, bov\_, bot\_, bok\_) and forall(\allowbreak{}[box\_ in perp(\allowbreak{}bov\_, bok\_)], bot\_(\allowbreak{}box\_) = vzero(\allowbreak{}bov\_))))
\end{quote}
mentioned by 12 result(s)

\smallskip

\noindent\textbf{\texttt{is-partial-order(grd, porel)}}\quad\textit{predicate}\index{is-partial-order(grd, porel)@\texttt{is-partial-order(grd, porel)}}\\
reads: \$2 partially orders \$1 \\
defined by \texttt{is-partial-order} \\
characterized by \texttt{is-partial-order}  \textit{(definitional)}:
\begin{quote}\ttfamily\footnotesize\raggedright\noindent
forall(\allowbreak{}[grd, porel], is-partial-order(\allowbreak{}grd, porel) iff grd in set and porel in set and forall(\allowbreak{}[x in grd], [x, x] in porel) and forall(\allowbreak{}[x in grd, y in grd], [x, y] in porel and [y, x] in porel implies x = y) and forall(\allowbreak{}[x in grd, y in grd, z in grd], [x, y] in porel and [y, z] in porel implies [x, z] in porel))
\end{quote}
mentioned by 14 result(s)

\smallskip

\noindent\textbf{\texttt{is-partition(p, n, a, b)}}\quad\textit{predicate}\index{is-partition(p, n, a, b)@\texttt{is-partition(p, n, a, b)}}\\
reads: \$1 is a partition of [\$3, \$4] into \$2 pieces \\
defined by \texttt{is-partition} \\
characterized by \texttt{is-partition}  \textit{(definitional)}:
\begin{quote}\ttfamily\footnotesize\raggedright\noindent
forall(\allowbreak{}[p, n, a, b], is-partition(\allowbreak{}p, n, a, b) iff n in nn and 1 <= n and a in rr and b in rr and a < b and p in fun(\allowbreak{}nn, rr) and p(\allowbreak{}0) = a and p(\allowbreak{}n) = b and forall(\allowbreak{}[i\_, j\_], i\_ in interval(\allowbreak{}0, n) and j\_ in interval(\allowbreak{}0, n) and i\_ < j\_ implies p(\allowbreak{}i\_) < p(\allowbreak{}j\_)))
\end{quote}
mentioned by 9 result(s)

\smallskip

\noindent\textbf{\texttt{is-path(pgam, a, b)}}\quad\textit{predicate}\index{is-path(pgam, a, b)@\texttt{is-path(pgam, a, b)}}\\
reads: \$1 is a path on the interval [\$2, \$3] \\
defined by \texttt{is-path} \\
characterized by \texttt{is-path}  \textit{(definitional)}:
\begin{quote}\ttfamily\footnotesize\raggedright\noindent
forall(\allowbreak{}[pgam, a, b], is-path(\allowbreak{}pgam, a, b) iff a in rr and b in rr and a < b and pgam in fun(\allowbreak{}ccint(\allowbreak{}a, b), cc) and forall(\allowbreak{}[pax\_ in ccint(\allowbreak{}a, b)], is-continuous-at(\allowbreak{}subspace-ms(\allowbreak{}rr-ms, ccint(\allowbreak{}a, b)), cc-ms, pgam, pax\_)))
\end{quote}
also: \texttt{is-path-in-fun} \texttt{is-road-is-path} \texttt{is-path-endpoints} \texttt{is-path-of-coords} \texttt{is-path-continuous} \\
mentioned by 26 result(s)

\smallskip

\noindent\textbf{\texttt{is-path-connected(pcs\_, pca\_)}}\quad\textit{predicate}\index{is-path-connected(pcs_, pca_)@\texttt{is-path-connected(pcs\_, pca\_)}}\\
reads: \$2 is path-connected in the metric space \$1 \\
defined by \texttt{is-path-connected} \\
characterized by \texttt{is-path-connected}  \textit{(definitional)}:
\begin{quote}\ttfamily\footnotesize\raggedright\noindent
forall(\allowbreak{}[pcs\_, pca\_], is-path-connected(\allowbreak{}pcs\_, pca\_) iff is-metric-space(\allowbreak{}pcs\_) and pca\_ subset pts(\allowbreak{}pcs\_) and forall(\allowbreak{}[pcx\_ in pca\_, pcy\_ in pca\_], forsome(\allowbreak{}[pcg\_ in fun(\allowbreak{}ccint(\allowbreak{}0, 1), pca\_)], is-continuous(\allowbreak{}subspace-ms(\allowbreak{}rr-ms, ccint(\allowbreak{}0, 1)), pcs\_, pcg\_) and pcg\_(\allowbreak{}0) = pcx\_ and pcg\_(\allowbreak{}1) = pcy\_)))
\end{quote}
mentioned by 5 result(s)

\smallskip

\noindent\textbf{\texttt{is-pi-system(pp)}}\quad\textit{predicate}\index{is-pi-system(pp)@\texttt{is-pi-system(pp)}}\\
reads: \$1 is closed under binary intersection \\
defined by \texttt{is-pi-system} \\
characterized by \texttt{is-pi-system}  \textit{(definitional)}:
\begin{quote}\ttfamily\footnotesize\raggedright\noindent
forall(\allowbreak{}[pp], is-pi-system(\allowbreak{}pp) iff forall(\allowbreak{}[pya\_ in pp, pyb\_ in pp], intersection(\allowbreak{}pya\_, pyb\_) in pp))
\end{quote}
mentioned by 5 result(s)

\smallskip

\noindent\textbf{\texttt{is-pid(s)}}\quad\textit{predicate}\index{is-pid(s)@\texttt{is-pid(s)}}\\
reads: a pid \\
a predicate \\
characterized by \texttt{is-pid-def}  \textit{(definitional)}:
\begin{quote}\ttfamily\footnotesize\raggedright\noindent
forall(\allowbreak{}[s], is-pid(\allowbreak{}s) iff is-integral-domain(\allowbreak{}s) and forall(\allowbreak{}[i in set], is-ideal(\allowbreak{}s, i) implies forsome(\allowbreak{}[a in carr(\allowbreak{}s)], i = principal-ideal(\allowbreak{}s, a))))
\end{quote}
also: \texttt{is-pid-is-set} \texttt{poly-is-pid-and-ufd} \texttt{euclidean-domain-is-pid} \\
mentioned by 13 result(s)

\smallskip

\noindent\textbf{\texttt{is-piecewise-continuous(f, a, b)}}\quad\textit{predicate}\index{is-piecewise-continuous(f, a, b)@\texttt{is-piecewise-continuous(f, a, b)}}\\
reads: \$1 is piecewise continuous on [\$2, \$3] \\
defined by \texttt{is-piecewise-continuous} \\
characterized by \texttt{is-piecewise-continuous}  \textit{(definitional)}:
\begin{quote}\ttfamily\footnotesize\raggedright\noindent
forall(\allowbreak{}[f, a, b], is-piecewise-continuous(\allowbreak{}f, a, b) iff f in fun(\allowbreak{}rr, rr) and a in rr and b in rr and a < b and forsome(\allowbreak{}[n\_, p\_], is-partition(\allowbreak{}p\_, n\_, a, b) and forall(\allowbreak{}[i\_], i\_ in interval(\allowbreak{}0, n\_) and i\_ < n\_ implies forall(\allowbreak{}[t\_], t\_ in rr and p\_(\allowbreak{}i\_) < t\_ and t\_ < p\_(\allowbreak{}succ(\allowbreak{}i\_)) implies is-continuous-at(\allowbreak{}rr-ms, rr-ms, f, t\_)))))
\end{quote}
mentioned by 2 result(s)

\smallskip

\noindent\textbf{\texttt{is-pk-space(s)}}\quad\textit{predicate}\index{is-pk-space(s)@\texttt{is-pk-space(s)}}\\
reads: a public-key strand space \\
a predicate \\
characterized by \texttt{is-pk-space}  \textit{(definitional)}:
\begin{quote}\ttfamily\footnotesize\raggedright\noindent
forall(\allowbreak{}[s], is-pk-space(\allowbreak{}s) iff length(\allowbreak{}s) = 3 and is-infiltrated-space(\allowbreak{}pk-inf(\allowbreak{}s)) and names(\allowbreak{}s) in set and pubkey(\allowbreak{}s) in fun(\allowbreak{}names(\allowbreak{}s), keys(\allowbreak{}ispace(\allowbreak{}pk-inf(\allowbreak{}s)))) and forall(\allowbreak{}[nlx\_ in names(\allowbreak{}s)], nlx\_ in tokens(\allowbreak{}ispace(\allowbreak{}pk-inf(\allowbreak{}s)))) and forall(\allowbreak{}[nlx\_ in names(\allowbreak{}s)], not(\allowbreak{}nlx\_ in keys(\allowbreak{}ispace(\allowbreak{}pk-inf(\allowbreak{}s))))) and forall(\allowbreak{}[nlx\_, nly\_], nlx\_ in names(\allowbreak{}s) implies nly\_ in names(\allowbreak{}s) implies (\allowbreak{}pubkey(\allowbreak{}s))(\allowbreak{}nlx\_) = (\allowbreak{}pubkey(\allowbreak{}s))(\allowbreak{}nly\_) implies nlx\_ = nly\_) and s in set)
\end{quote}
also: \texttt{is-pk-space-is-set} \\
mentioned by 25 result(s)

\smallskip

\noindent\textbf{\texttt{is-point-separating(rst\_, rsa\_)}}\quad\textit{predicate}\index{is-point-separating(rst_, rsa_)@\texttt{is-point-separating(rst\_, rsa\_)}}\\
reads: \$2 separates the points of \$1 \\
defined by \texttt{is-point-separating} \\
characterized by \texttt{is-point-separating}  \textit{(definitional)}:
\begin{quote}\ttfamily\footnotesize\raggedright\noindent
forall(\allowbreak{}[rst\_, rsa\_], is-point-separating(\allowbreak{}rst\_, rsa\_) iff forall(\allowbreak{}[rsp\_ in pts(\allowbreak{}rst\_), rsq\_ in pts(\allowbreak{}rst\_)], not(\allowbreak{}rsp\_ = rsq\_) implies forsome(\allowbreak{}[rsf\_ in rsa\_], not(\allowbreak{}rsf\_(\allowbreak{}rsp\_) = rsf\_(\allowbreak{}rsq\_)))))
\end{quote}
mentioned by 10 result(s)

\smallskip

\noindent\textbf{\texttt{is-pole(mrf\_, mra\_, mrn\_)}}\quad\textit{predicate}\index{is-pole(mrf_, mra_, mrn_)@\texttt{is-pole(mrf\_, mra\_, mrn\_)}}\\
reads: \$1 has a pole of order \$3 at \$2 \\
defined by \texttt{is-pole} \\
characterized by \texttt{is-pole}  \textit{(definitional)}:
\begin{quote}\ttfamily\footnotesize\raggedright\noindent
forall(\allowbreak{}[mrf\_, mra\_, mrn\_], is-pole(\allowbreak{}mrf\_, mra\_, mrn\_) iff mra\_ in cc and mrn\_ in nn and 1 <= mrn\_ and forsome(\allowbreak{}[mrr\_ in rr], 0 < mrr\_ and forsome(\allowbreak{}[mrg\_], holomorphic-on(\allowbreak{}ball(\allowbreak{}cc-ms, mra\_, mrr\_), mrg\_) and not(\allowbreak{}mrg\_(\allowbreak{}mra\_) = 0) and forall(\allowbreak{}[mrz\_ in difference(\allowbreak{}ball(\allowbreak{}cc-ms, mra\_, mrr\_), singleton(\allowbreak{}mra\_))], mrf\_(\allowbreak{}mrz\_) = recip(\allowbreak{}mrz\_ - mra\_) \textasciicircum{} mrn\_ * mrg\_(\allowbreak{}mrz\_)))))
\end{quote}
mentioned by 24 result(s)

\smallskip

\noindent\textbf{\texttt{is-polish(s)}}\quad\textit{predicate}\index{is-polish(s)@\texttt{is-polish(s)}}\\
reads: \$1 is a Polish space \\
defined by \texttt{is-polish} \\
characterized by \texttt{is-polish}  \textit{(definitional)}:
\begin{quote}\ttfamily\footnotesize\raggedright\noindent
forall(\allowbreak{}[s], is-polish(\allowbreak{}s) iff is-complete(\allowbreak{}s) and is-separable(\allowbreak{}s))
\end{quote}
mentioned by 36 result(s)

\smallskip

\noindent\textbf{\texttt{is-poly-deg(kk, a, f, n)}}\quad\textit{predicate}\index{is-poly-deg(kk, a, f, n)@\texttt{is-poly-deg(kk, a, f, n)}}\\
reads: \$3 is a polynomial function of degree \$4 on \$2 with coefficients in \$1 \\
defined by \texttt{is-poly-deg} \\
characterized by \texttt{is-poly-deg}  \textit{(definitional)}:
\begin{quote}\ttfamily\footnotesize\raggedright\noindent
forall(\allowbreak{}[kk, a, f, n], is-poly-deg(\allowbreak{}kk, a, f, n) iff a subset kk and f in fun(\allowbreak{}a, kk) and forsome(\allowbreak{}[pfc\_ in fun(\allowbreak{}nn, kk)], forall(\allowbreak{}[pfk\_ in nn], succ(\allowbreak{}n) <= pfk\_ implies pfc\_(\allowbreak{}pfk\_) = 0) and not(\allowbreak{}pfc\_(\allowbreak{}n) = 0) and forall(\allowbreak{}[pfz\_ in a], f(\allowbreak{}pfz\_) = cc-series-partial-sum(\allowbreak{}vnb-lambda(\allowbreak{}pfj\_, nn, pfc\_(\allowbreak{}pfj\_) * pfz\_ \textasciicircum{} pfj\_), succ(\allowbreak{}n)))))
\end{quote}
mentioned by 12 result(s)

\smallskip

\noindent\textbf{\texttt{is-poly-fn-on(kk, a, f)}}\quad\textit{predicate}\index{is-poly-fn-on(kk, a, f)@\texttt{is-poly-fn-on(kk, a, f)}}\\
reads: \$3 is a polynomial function on \$2 with coefficients in \$1 \\
defined by \texttt{is-poly-fn-on} \\
characterized by \texttt{is-poly-fn-on}  \textit{(definitional)}:
\begin{quote}\ttfamily\footnotesize\raggedright\noindent
forall(\allowbreak{}[kk, a, f], is-poly-fn-on(\allowbreak{}kk, a, f) iff a subset kk and f in fun(\allowbreak{}a, kk) and forsome(\allowbreak{}[pfn\_ in nn, pfc\_ in fun(\allowbreak{}nn, kk)], forall(\allowbreak{}[pfk\_ in nn], pfn\_ <= pfk\_ implies pfc\_(\allowbreak{}pfk\_) = 0) and forall(\allowbreak{}[pfz\_ in a], f(\allowbreak{}pfz\_) = cc-series-partial-sum(\allowbreak{}vnb-lambda(\allowbreak{}pfj\_, nn, pfc\_(\allowbreak{}pfj\_) * pfz\_ \textasciicircum{} pfj\_), pfn\_))))
\end{quote}
mentioned by 24 result(s)

\smallskip

\noindent\textbf{\texttt{is-positive-form(ubv\_, ubd\_, ubf\_)}}\quad\textit{predicate}\index{is-positive-form(ubv_, ubd_, ubf_)@\texttt{is-positive-form(ubv\_, ubd\_, ubf\_)}}\\
reads: \$3 is a positive quadratic form on \$1 with domain \$2 \\
defined by \texttt{is-positive-form} \\
characterized by \texttt{is-positive-form}  \textit{(definitional)}:
\begin{quote}\ttfamily\footnotesize\raggedright\noindent
forall(\allowbreak{}[ubv\_, ubd\_, ubf\_], is-positive-form(\allowbreak{}ubv\_, ubd\_, ubf\_) iff is-quadratic-form(\allowbreak{}ubv\_, ubd\_, ubf\_) and forall(\allowbreak{}[ubx\_ in ubd\_], ubf\_(\allowbreak{}ubx\_, ubx\_) in rr and 0 <= ubf\_(\allowbreak{}ubx\_, ubx\_)))
\end{quote}
mentioned by 6 result(s)

\smallskip

\noindent\textbf{\texttt{is-positive-functional(s, lam)}}\quad\textit{predicate}\index{is-positive-functional(s, lam)@\texttt{is-positive-functional(s, lam)}}\\
reads: \$2 is a positive linear functional on the continuous real functions on \$1 \\
defined by \texttt{is-positive-functional} \\
characterized by \texttt{is-positive-functional}  \textit{(definitional)}:
\begin{quote}\ttfamily\footnotesize\raggedright\noindent
forall(\allowbreak{}[s, lam], is-positive-functional(\allowbreak{}s, lam) iff lam in fun(\allowbreak{}cont-fns(\allowbreak{}s), rr) and forall(\allowbreak{}[pff\_ in cont-fns(\allowbreak{}s), pfg\_ in cont-fns(\allowbreak{}s)], lam(\allowbreak{}vnb-lambda(\allowbreak{}sx\_, pts(\allowbreak{}s), pff\_(\allowbreak{}sx\_) + pfg\_(\allowbreak{}sx\_))) = lam(\allowbreak{}pff\_) + lam(\allowbreak{}pfg\_)) and forall(\allowbreak{}[pfc\_ in rr, pff\_ in cont-fns(\allowbreak{}s)], lam(\allowbreak{}vnb-lambda(\allowbreak{}sx\_, pts(\allowbreak{}s), pfc\_ * pff\_(\allowbreak{}sx\_))) = pfc\_ * lam(\allowbreak{}pff\_)) and forall(\allowbreak{}[pff\_ in cont-fns(\allowbreak{}s)], forall(\allowbreak{}[pfx\_ in pts(\allowbreak{}s)], 0 <= pff\_(\allowbreak{}pfx\_)) implies 0 <= lam(\allowbreak{}pff\_)))
\end{quote}
mentioned by 42 result(s)

\smallskip

\noindent\textbf{\texttt{is-positive-lf(g4a\_, g4f\_)}}\quad\textit{predicate}\index{is-positive-lf(g4a_, g4f_)@\texttt{is-positive-lf(g4a\_, g4f\_)}}\\
reads: \$2 is a positive linear functional on \$1 \\
defined by \texttt{is-positive-lf} \\
characterized by \texttt{is-positive-lf}  \textit{(definitional)}:
\begin{quote}\ttfamily\footnotesize\raggedright\noindent
forall(\allowbreak{}[g4a\_, g4f\_], is-positive-lf(\allowbreak{}g4a\_, g4f\_) iff is-cc-linear-on(\allowbreak{}g4a\_, g4f\_) and forall(\allowbreak{}[g4x\_ in vec(\allowbreak{}g4a\_)], g4f\_(\allowbreak{}(\allowbreak{}vmul(\allowbreak{}g4a\_))(\allowbreak{}(\allowbreak{}vstar(\allowbreak{}g4a\_))(\allowbreak{}g4x\_), g4x\_)) in rr and 0 <= g4f\_(\allowbreak{}(\allowbreak{}vmul(\allowbreak{}g4a\_))(\allowbreak{}(\allowbreak{}vstar(\allowbreak{}g4a\_))(\allowbreak{}g4x\_), g4x\_))))
\end{quote}
mentioned by 77 result(s)

\smallskip

\noindent\textbf{\texttt{is-positive-op(kqv\_, kqt\_)}}\quad\textit{predicate}\index{is-positive-op(kqv_, kqt_)@\texttt{is-positive-op(kqv\_, kqt\_)}}\\
reads: \$2 is a positive operator on \$1 \\
defined by \texttt{is-positive-op} \\
characterized by \texttt{is-positive-op}  \textit{(definitional)}:
\begin{quote}\ttfamily\footnotesize\raggedright\noindent
forall(\allowbreak{}[kqv\_, kqt\_], is-positive-op(\allowbreak{}kqv\_, kqt\_) iff is-self-adjoint-op(\allowbreak{}kqv\_, kqt\_) and forall(\allowbreak{}[kqx\_ in vec(\allowbreak{}kqv\_)], (\allowbreak{}ip(\allowbreak{}kqv\_))(\allowbreak{}kqt\_(\allowbreak{}kqx\_), kqx\_) in rr and 0 <= (\allowbreak{}ip(\allowbreak{}kqv\_))(\allowbreak{}kqt\_(\allowbreak{}kqx\_), kqx\_)))
\end{quote}
mentioned by 4 result(s)

\smallskip

\noindent\textbf{\texttt{is-power-bounded(k, n, tm)}}\quad\textit{predicate}\index{is-power-bounded(k, n, tm)@\texttt{is-power-bounded(k, n, tm)}}\\
reads: the \$2-by-\$2 matrix \$3 over \$1 is power bounded \\
defined by \texttt{is-power-bounded} \\
characterized by \texttt{is-power-bounded}  \textit{(definitional)}:
\begin{quote}\ttfamily\footnotesize\raggedright\noindent
forall(\allowbreak{}[k, n, tm], is-power-bounded(\allowbreak{}k, n, tm) iff forsome(\allowbreak{}[cbvc\_], is-power-bounded-by(\allowbreak{}k, n, tm, cbvc\_)))
\end{quote}
mentioned by 14 result(s)

\smallskip

\noindent\textbf{\texttt{is-power-bounded-by(k, n, tm, c)}}\quad\textit{predicate}\index{is-power-bounded-by(k, n, tm, c)@\texttt{is-power-bounded-by(k, n, tm, c)}}\\
reads: the powers of the \$2-by-\$2 matrix \$3 over \$1 have max norm at most \$4 \\
defined by \texttt{is-power-bounded-by} \\
characterized by \texttt{is-power-bounded-by}  \textit{(definitional)}:
\begin{quote}\ttfamily\footnotesize\raggedright\noindent
forall(\allowbreak{}[k, n, tm, c], is-power-bounded-by(\allowbreak{}k, n, tm, c) iff tm in mat(\allowbreak{}n, n, carr(\allowbreak{}k)) and c in rr and 0 <= c and forall(\allowbreak{}[cbvk\_ in nn], is-mat-bound(\allowbreak{}k, n, n, mat-power(\allowbreak{}k, n, tm, cbvk\_), c)))
\end{quote}
mentioned by 7 result(s)

\smallskip

\noindent\textbf{\texttt{is-pre-hilbert-space(s)}}\quad\textit{predicate}\index{is-pre-hilbert-space(s)@\texttt{is-pre-hilbert-space(s)}}\\
reads: a pre-Hilbert space \\
a predicate \\
characterized by \texttt{is-pre-hilbert-space-def}  \textit{(definitional)}:
\begin{quote}\ttfamily\footnotesize\raggedright\noindent
forall(\allowbreak{}[s], is-pre-hilbert-space(\allowbreak{}s) iff is-semi-inner-product-space(\allowbreak{}s) and forall(\allowbreak{}[x\_ in vec(\allowbreak{}s)], (\allowbreak{}rip(\allowbreak{}s))(\allowbreak{}x\_, x\_) = 0 implies x\_ = vzero(\allowbreak{}s)))
\end{quote}
also: \texttt{is-pre-hilbert-space-is-set} \\
mentioned by 52 result(s)

\smallskip

\noindent\textbf{\texttt{is-prec-closed(spc\_, snd\_, ssd\_, srl\_)}}\quad\textit{predicate}\index{is-prec-closed(spc_, snd_, ssd_, srl_)@\texttt{is-prec-closed(spc\_, snd\_, ssd\_, srl\_)}}\\
reads: \$4 contains \$3 and the strand order of \$1 on \$2, and is transitive \\
defined by \texttt{is-prec-closed} \\
characterized by \texttt{is-prec-closed}  \textit{(definitional)}:
\begin{quote}\ttfamily\footnotesize\raggedright\noindent
forall(\allowbreak{}[spc\_, snd\_, ssd\_, srl\_], is-prec-closed(\allowbreak{}spc\_, snd\_, ssd\_, srl\_) iff forall(\allowbreak{}[ssx\_, ssy\_], [ssx\_, ssy\_] in ssd\_ implies [ssx\_, ssy\_] in srl\_) and forall(\allowbreak{}[ssx\_ in snd\_, ssy\_ in snd\_], same-strand-before(\allowbreak{}spc\_, ssx\_, ssy\_) implies [ssx\_, ssy\_] in srl\_) and forall(\allowbreak{}[ssx\_, ssy\_, ssz\_], [ssx\_, ssy\_] in srl\_ implies [ssy\_, ssz\_] in srl\_ implies [ssx\_, ssz\_] in srl\_))
\end{quote}
mentioned by 6 result(s)

\smallskip

\noindent\textbf{\texttt{is-premeasure(omega, cls, pm)}}\quad\textit{predicate}\index{is-premeasure(omega, cls, pm)@\texttt{is-premeasure(omega, cls, pm)}}\\
reads: \$3 is a pre-measure on the family \$2 of subsets of \$1 \\
defined by \texttt{is-premeasure} \\
characterized by \texttt{is-premeasure}  \textit{(definitional)}:
\begin{quote}\ttfamily\footnotesize\raggedright\noindent
forall(\allowbreak{}[omega, cls, pm], is-premeasure(\allowbreak{}omega, cls, pm) iff omega in set and cls subset power(\allowbreak{}omega) and empty-set in cls and pm in fun(\allowbreak{}cls, rr-pos-star) and pm(\allowbreak{}empty-set) = 0)
\end{quote}
also: \texttt{measure-is-premeasure} \\
mentioned by 24 result(s)

\smallskip

\noindent\textbf{\texttt{is-prime-element(gps\_, gpp\_)}}\quad\textit{predicate}\index{is-prime-element(gps_, gpp_)@\texttt{is-prime-element(gps\_, gpp\_)}}\\
reads: \$2 is a prime of \$1 \\
defined by \texttt{is-prime-element} \\
characterized by \texttt{is-prime-element}  \textit{(definitional)}:
\begin{quote}\ttfamily\footnotesize\raggedright\noindent
forall(\allowbreak{}[gps\_, gpp\_], is-prime-element(\allowbreak{}gps\_, gpp\_) iff gpp\_ in carr(\allowbreak{}gps\_) and not(\allowbreak{}gpp\_ = zero(\allowbreak{}gps\_)) and not(\allowbreak{}is-unit(\allowbreak{}gps\_, gpp\_)) and forall(\allowbreak{}[gpa\_ in carr(\allowbreak{}gps\_), gpb\_ in carr(\allowbreak{}gps\_)], gpp\_ = (\allowbreak{}mul(\allowbreak{}gps\_))(\allowbreak{}gpa\_, gpb\_) implies is-unit(\allowbreak{}gps\_, gpa\_) or is-unit(\allowbreak{}gps\_, gpb\_)))
\end{quote}
mentioned by 11 result(s)

\smallskip

\noindent\textbf{\texttt{is-prime-nat(gpn\_)}}\quad\textit{predicate}\index{is-prime-nat(gpn_)@\texttt{is-prime-nat(gpn\_)}}\\
reads: \$1 is a prime number \\
defined by \texttt{is-prime-nat} \\
characterized by \texttt{is-prime-nat}  \textit{(definitional)}:
\begin{quote}\ttfamily\footnotesize\raggedright\noindent
forall(\allowbreak{}[gpn\_], is-prime-nat(\allowbreak{}gpn\_) iff gpn\_ in nn and is-prime-element(\allowbreak{}zz-ring, gpn\_))
\end{quote}
mentioned by 50 result(s)

\smallskip

\noindent\textbf{\texttt{is-prime-seq(gqs\_, gqp\_, gqn\_)}}\quad\textit{predicate}\index{is-prime-seq(gqs_, gqp_, gqn_)@\texttt{is-prime-seq(gqs\_, gqp\_, gqn\_)}}\\
reads: \$2(1), ..., \$2(\$3) are primes of \$1 \\
defined by \texttt{is-prime-seq} \\
characterized by \texttt{is-prime-seq}  \textit{(definitional)}:
\begin{quote}\ttfamily\footnotesize\raggedright\noindent
forall(\allowbreak{}[gqs\_, gqp\_, gqn\_], is-prime-seq(\allowbreak{}gqs\_, gqp\_, gqn\_) iff forall(\allowbreak{}[gqi\_ in interval(\allowbreak{}1, gqn\_)], is-prime-element(\allowbreak{}gqs\_, gqp\_(\allowbreak{}gqi\_))))
\end{quote}
mentioned by 7 result(s)

\smallskip

\noindent\textbf{\texttt{is-primitive(rpg\_, rpf\_, rpa\_, rpb\_)}}\quad\textit{predicate}\index{is-primitive(rpg_, rpf_, rpa_, rpb_)@\texttt{is-primitive(rpg\_, rpf\_, rpa\_, rpb\_)}}\\
reads: \$1 is a primitive of \$2 on the interval [\$3, \$4] \\
defined by \texttt{is-primitive} \\
characterized by \texttt{is-primitive}  \textit{(definitional)}:
\begin{quote}\ttfamily\footnotesize\raggedright\noindent
forall(\allowbreak{}[rpg\_, rpf\_, rpa\_, rpb\_], is-primitive(\allowbreak{}rpg\_, rpf\_, rpa\_, rpb\_) iff rpa\_ in rr and rpb\_ in rr and rpa\_ < rpb\_ and rpg\_ in fun(\allowbreak{}ccint(\allowbreak{}rpa\_, rpb\_), rr) and rpf\_ in fun(\allowbreak{}ccint(\allowbreak{}rpa\_, rpb\_), rr) and is-continuous-on(\allowbreak{}rpg\_, ccint(\allowbreak{}rpa\_, rpb\_)) and forsome(\allowbreak{}[rps\_], is-countable(\allowbreak{}rps\_) and rps\_ subset ccint(\allowbreak{}rpa\_, rpb\_) and forall(\allowbreak{}[rpt\_ in ooint(\allowbreak{}rpa\_, rpb\_)], not(\allowbreak{}rpt\_ in rps\_) implies has-deriv-at(\allowbreak{}rpg\_, rpt\_, rpf\_(\allowbreak{}rpt\_)))))
\end{quote}
mentioned by 83 result(s)

\smallskip

\noindent\textbf{\texttt{is-principal-sequence(md, n, sm, es, dm)}}\quad\textit{predicate}\index{is-principal-sequence(md, n, sm, es, dm)@\texttt{is-principal-sequence(md, n, sm, es, dm)}}\\
reads: the diagonal of \$5 is a principal sequence for \$3 corresponding to the basis \$4 of \$1 \\
defined by \texttt{is-principal-sequence} \\
characterized by \texttt{is-principal-sequence}  \textit{(definitional)}:
\begin{quote}\ttfamily\footnotesize\raggedright\noindent
forall(\allowbreak{}[md, n, sm, es, dm], is-principal-sequence(\allowbreak{}md, n, sm, es, dm) iff es in mat(\allowbreak{}n, 1, vec(\allowbreak{}md)) and generates(\allowbreak{}md, n, es) and rel-free(\allowbreak{}md, n, es) and dm in mat(\allowbreak{}n, n, zz) and is-diagonal(\allowbreak{}zz-ring, n, n, dm) and forall(\allowbreak{}[psi\_ in interval(\allowbreak{}1, n)], 0 <= entry(\allowbreak{}dm, psi\_, psi\_)) and forall(\allowbreak{}[psi\_, psj\_], psi\_ in interval(\allowbreak{}1, n) implies psj\_ in interval(\allowbreak{}1, n) implies psi\_ <= psj\_ implies zz-divides(\allowbreak{}entry(\allowbreak{}dm, psi\_, psi\_), entry(\allowbreak{}dm, psj\_, psj\_))) and spans(\allowbreak{}md, n, matact(\allowbreak{}md, dm, es), sm))
\end{quote}
mentioned by 57 result(s)

\smallskip

\noindent\textbf{\texttt{is-projection-el(baa\_, bax\_)}}\quad\textit{predicate}\index{is-projection-el(baa_, bax_)@\texttt{is-projection-el(baa\_, bax\_)}}\\
reads: \$2 is a projection in \$1 \\
defined by \texttt{is-projection-el} \\
characterized by \texttt{is-projection-el}  \textit{(definitional)}:
\begin{quote}\ttfamily\footnotesize\raggedright\noindent
forall(\allowbreak{}[baa\_, bax\_], is-projection-el(\allowbreak{}baa\_, bax\_) iff bax\_ in vec(\allowbreak{}baa\_) and (\allowbreak{}vstar(\allowbreak{}baa\_))(\allowbreak{}bax\_) = bax\_ and (\allowbreak{}vmul(\allowbreak{}baa\_))(\allowbreak{}bax\_, bax\_) = bax\_)
\end{quote}
mentioned by 2 result(s)

\smallskip

\noindent\textbf{\texttt{is-projection-op(bov\_, bot\_)}}\quad\textit{predicate}\index{is-projection-op(bov_, bot_)@\texttt{is-projection-op(bov\_, bot\_)}}\\
reads: \$2 is a projection on \$1 \\
defined by \texttt{is-projection-op} \\
characterized by \texttt{is-projection-op}  \textit{(definitional)}:
\begin{quote}\ttfamily\footnotesize\raggedright\noindent
forall(\allowbreak{}[bov\_, bot\_], is-projection-op(\allowbreak{}bov\_, bot\_) iff bot\_ in bounded-ops(\allowbreak{}bov\_, bov\_) and op-adjoint(\allowbreak{}bov\_, bov\_, bot\_) = bot\_ and bop-comp(\allowbreak{}bov\_, bot\_, bot\_) = bot\_)
\end{quote}
mentioned by 31 result(s)

\smallskip

\noindent\textbf{\texttt{is-psd-form(ipfm\_, ipff\_)}}\quad\textit{predicate}\index{is-psd-form(ipfm_, ipff_)@\texttt{is-psd-form(ipfm\_, ipff\_)}}\\
reads: \$2 is a positive semi-definite hermitian form on \$1 \\
defined by \texttt{is-psd-form} \\
characterized by \texttt{is-psd-form}  \textit{(definitional)}:
\begin{quote}\ttfamily\footnotesize\raggedright\noindent
forall(\allowbreak{}[ipfm\_, ipff\_], is-psd-form(\allowbreak{}ipfm\_, ipff\_) iff is-hermitian-form(\allowbreak{}ipfm\_, ipff\_) and forall(\allowbreak{}[ipfx\_ in vec(\allowbreak{}ipfm\_)], ipff\_(\allowbreak{}ipfx\_, ipfx\_) in rr) and forall(\allowbreak{}[ipfx\_ in vec(\allowbreak{}ipfm\_)], 0 <= ipff\_(\allowbreak{}ipfx\_, ipfx\_)))
\end{quote}
mentioned by 41 result(s)

\smallskip

\noindent\textbf{\texttt{is-pseudometric(dst, crr)}}\quad\textit{predicate}\index{is-pseudometric(dst, crr)@\texttt{is-pseudometric(dst, crr)}}\\
reads: \$1 is a pseudometric on \$2 \\
defined by \texttt{is-pseudometric} \\
characterized by \texttt{is-pseudometric}  \textit{(definitional)}:
\begin{quote}\ttfamily\footnotesize\raggedright\noindent
forall(\allowbreak{}[dst, crr], is-pseudometric(\allowbreak{}dst, crr) iff forall(\allowbreak{}[u in crr], dst(\allowbreak{}u, u) = 0 and forall(\allowbreak{}[v in crr], 0 <= dst(\allowbreak{}u, v) and dst(\allowbreak{}u, v) = dst(\allowbreak{}v, u) and forall(\allowbreak{}[w in crr], dst(\allowbreak{}u, w) <= dst(\allowbreak{}u, v) + dst(\allowbreak{}v, w)))))
\end{quote}
also: \texttt{is-pseudometric-space} \\
mentioned by 16 result(s)

\smallskip

\noindent\textbf{\texttt{is-pseudometric-space(s)}}\quad\textit{predicate}\index{is-pseudometric-space(s)@\texttt{is-pseudometric-space(s)}}\\
reads: a pseudometric space \\
a predicate \\
characterized by \texttt{is-pseudometric-space}  \textit{(definitional)}:
\begin{quote}\ttfamily\footnotesize\raggedright\noindent
forall(\allowbreak{}[s], is-pseudometric-space(\allowbreak{}s) iff length(\allowbreak{}s) = 2 and pts(\allowbreak{}s) in set and dist(\allowbreak{}s) in fun(\allowbreak{}cartesian(\allowbreak{}pts(\allowbreak{}s), pts(\allowbreak{}s)), rr) and is-pseudometric(\allowbreak{}dist(\allowbreak{}s), pts(\allowbreak{}s)) and s in set)
\end{quote}
also: \texttt{is-pseudometric-space-is-set} \\
mentioned by 15 result(s)

\smallskip

\noindent\textbf{\texttt{is-pt-open(t, u)}}\quad\textit{predicate}\index{is-pt-open(t, u)@\texttt{is-pt-open(t, u)}}\\
reads: \$2 is open in the product of \$1 and R \\
defined by \texttt{is-pt-open} \\
characterized by \texttt{is-pt-open}  \textit{(definitional)}:
\begin{quote}\ttfamily\footnotesize\raggedright\noindent
forall(\allowbreak{}[t, u], is-pt-open(\allowbreak{}t, u) iff u subset cartesian(\allowbreak{}pts(\allowbreak{}t), rr) and forall(\allowbreak{}[ptw\_ in u], forsome(\allowbreak{}[ptv\_ in opens(\allowbreak{}t)], nth(\allowbreak{}1, ptw\_) in ptv\_ and forsome(\allowbreak{}[pte\_], pos-rr(\allowbreak{}pte\_) and pt-box(\allowbreak{}t, ptv\_, nth(\allowbreak{}2, ptw\_), pte\_) subset u))))
\end{quote}
mentioned by 8 result(s)

\smallskip

\noindent\textbf{\texttt{is-ptwise-cauchy(s, fam)}}\quad\textit{predicate}\index{is-ptwise-cauchy(s, fam)@\texttt{is-ptwise-cauchy(s, fam)}}\\
reads: \$2 is pointwise Cauchy on \$1 \\
defined by \texttt{is-ptwise-cauchy} \\
characterized by \texttt{is-ptwise-cauchy}  \textit{(definitional)}:
\begin{quote}\ttfamily\footnotesize\raggedright\noindent
forall(\allowbreak{}[s, fam], is-ptwise-cauchy(\allowbreak{}s, fam) iff is-metric-space(\allowbreak{}s) and fam in fun(\allowbreak{}nn, fun(\allowbreak{}pts(\allowbreak{}s), rr)) and forall(\allowbreak{}[x in pts(\allowbreak{}s), eps], pos-rr(\allowbreak{}eps) implies forsome(\allowbreak{}[cap in nn], forall(\allowbreak{}[k], k in nn and cap <= k implies forall(\allowbreak{}[l], l in nn and cap <= l implies abs(\allowbreak{}(\allowbreak{}fam(\allowbreak{}k))(\allowbreak{}x) - (\allowbreak{}fam(\allowbreak{}l))(\allowbreak{}x)) < eps)))))
\end{quote}
mentioned by 4 result(s)

\smallskip

\noindent\textbf{\texttt{is-pure-state(g4a\_, g4f\_)}}\quad\textit{predicate}\index{is-pure-state(g4a_, g4f_)@\texttt{is-pure-state(g4a\_, g4f\_)}}\\
reads: \$2 is a pure state of \$1 \\
defined by \texttt{is-pure-state} \\
characterized by \texttt{is-pure-state}  \textit{(definitional)}:
\begin{quote}\ttfamily\footnotesize\raggedright\noindent
forall(\allowbreak{}[g4a\_, g4f\_], is-pure-state(\allowbreak{}g4a\_, g4f\_) iff g4f\_ in states(\allowbreak{}g4a\_) and forall(\allowbreak{}[g4g\_], is-positive-lf(\allowbreak{}g4a\_, g4g\_) implies functional-le(\allowbreak{}g4a\_, g4g\_, g4f\_) implies forsome(\allowbreak{}[g4l\_ in rr], 0 <= g4l\_ and g4l\_ <= 1 and forall(\allowbreak{}[g4x\_ in vec(\allowbreak{}g4a\_)], g4g\_(\allowbreak{}g4x\_) = g4l\_ * g4f\_(\allowbreak{}g4x\_)))))
\end{quote}
mentioned by 4 result(s)

\smallskip

\noindent\textbf{\texttt{is-pw-antiderivable(pphi\_, a, b)}}\quad\textit{predicate}\index{is-pw-antiderivable(pphi_, a, b)@\texttt{is-pw-antiderivable(pphi\_, a, b)}}\\
reads: \$1 has a primitive on the interval [\$2, \$3] \\
defined by \texttt{is-pw-antiderivable} \\
characterized by \texttt{is-pw-antiderivable}  \textit{(definitional)}:
\begin{quote}\ttfamily\footnotesize\raggedright\noindent
forall(\allowbreak{}[pphi\_, a, b], is-pw-antiderivable(\allowbreak{}pphi\_, a, b) iff forsome(\allowbreak{}[paw\_], is-primitive(\allowbreak{}paw\_, pphi\_, a, b)))
\end{quote}
mentioned by 2 result(s)

\smallskip

\noindent\textbf{\texttt{is-pw-continuous-on(pphi\_, a, b)}}\quad\textit{predicate}\index{is-pw-continuous-on(pphi_, a, b)@\texttt{is-pw-continuous-on(pphi\_, a, b)}}\\
reads: \$1 is piecewise continuous on the interval [\$2, \$3] \\
defined by \texttt{is-pw-continuous-on} \\
characterized by \texttt{is-pw-continuous-on}  \textit{(definitional)}:
\begin{quote}\ttfamily\footnotesize\raggedright\noindent
forall(\allowbreak{}[pphi\_, a, b], is-pw-continuous-on(\allowbreak{}pphi\_, a, b) iff a in rr and b in rr and a < b and pphi\_ in fun(\allowbreak{}ccint(\allowbreak{}a, b), rr) and forsome(\allowbreak{}[pas\_ in set], card(\allowbreak{}pas\_) in nn and pas\_ subset ccint(\allowbreak{}a, b) and forall(\allowbreak{}[pat\_ in ooint(\allowbreak{}a, b)], not(\allowbreak{}pat\_ in pas\_) implies is-continuous-at(\allowbreak{}subspace-ms(\allowbreak{}rr-ms, ccint(\allowbreak{}a, b)), rr-ms, pphi\_, pat\_))))
\end{quote}
mentioned by 6 result(s)

\smallskip

\noindent\textbf{\texttt{is-qp-digits(qsp\_, qta\_, qtn\_)}}\quad\textit{predicate}\index{is-qp-digits(qsp_, qta_, qtn_)@\texttt{is-qp-digits(qsp\_, qta\_, qtn\_)}}\\
reads: \$2 is a \$1-adic digit sequence starting at \$3 \\
defined by \texttt{is-qp-digits} \\
characterized by \texttt{is-qp-digits}  \textit{(definitional)}:
\begin{quote}\ttfamily\footnotesize\raggedright\noindent
forall(\allowbreak{}[qsp\_, qta\_, qtn\_], is-qp-digits(\allowbreak{}qsp\_, qta\_, qtn\_) iff qta\_ in fun(\allowbreak{}zz, zz) and forall(\allowbreak{}[qtk\_ in zz], 0 <= qta\_(\allowbreak{}qtk\_) and qta\_(\allowbreak{}qtk\_) < qsp\_ and (\allowbreak{}qtk\_ < qtn\_ implies qta\_(\allowbreak{}qtk\_) = 0)))
\end{quote}
mentioned by 4 result(s)

\smallskip

\noindent\textbf{\texttt{is-qp-expansion(qsp\_, qta\_, qtn\_, qtx\_)}}\quad\textit{predicate}\index{is-qp-expansion(qsp_, qta_, qtn_, qtx_)@\texttt{is-qp-expansion(qsp\_, qta\_, qtn\_, qtx\_)}}\\
reads: \$4 = sum over k >= \$3 of \$2(k) \$1\textasciicircum{}k in Q\_\$1 \\
defined by \texttt{is-qp-expansion} \\
characterized by \texttt{is-qp-expansion}  \textit{(definitional)}:
\begin{quote}\ttfamily\footnotesize\raggedright\noindent
forall(\allowbreak{}[qsp\_, qta\_, qtn\_, qtx\_], is-qp-expansion(\allowbreak{}qsp\_, qta\_, qtn\_, qtx\_) iff is-qp-digits(\allowbreak{}qsp\_, qta\_, qtn\_) and converges-to(\allowbreak{}qp(\allowbreak{}qsp\_), vnb-lambda(\allowbreak{}qtm\_, nn, qp-psum(\allowbreak{}qsp\_, qta\_, qtn\_, qtm\_)), qtx\_))
\end{quote}
mentioned by 11 result(s)

\smallskip

\noindent\textbf{\texttt{is-quadratic-form(ubv\_, ubd\_, ubf\_)}}\quad\textit{predicate}\index{is-quadratic-form(ubv_, ubd_, ubf_)@\texttt{is-quadratic-form(ubv\_, ubd\_, ubf\_)}}\\
reads: \$3 is a quadratic form on \$1 with domain \$2 \\
defined by \texttt{is-quadratic-form} \\
characterized by \texttt{is-quadratic-form}  \textit{(definitional)}:
\begin{quote}\ttfamily\footnotesize\raggedright\noindent
forall(\allowbreak{}[ubv\_, ubd\_, ubf\_], is-quadratic-form(\allowbreak{}ubv\_, ubd\_, ubf\_) iff is-submodule(\allowbreak{}ubv\_, ubd\_) and is-dense(\allowbreak{}ip-metric-space(\allowbreak{}ubv\_), ubd\_) and ubf\_ in fun(\allowbreak{}cartesian(\allowbreak{}ubd\_, ubd\_), cc) and forall(\allowbreak{}[ubx\_ in ubd\_, uby\_ in ubd\_, ubz\_ in ubd\_], ubf\_(\allowbreak{}(\allowbreak{}vadd(\allowbreak{}ubv\_))(\allowbreak{}ubx\_, uby\_), ubz\_) = ubf\_(\allowbreak{}ubx\_, ubz\_) + ubf\_(\allowbreak{}uby\_, ubz\_)) and forall(\allowbreak{}[ubk\_ in cc, ubx\_ in ubd\_, uby\_ in ubd\_], ubf\_(\allowbreak{}(\allowbreak{}act(\allowbreak{}ubv\_))(\allowbreak{}ubk\_, ubx\_), uby\_) = ubk\_ * ubf\_(\allowbreak{}ubx\_, uby\_)) and forall(\allowbreak{}[ubx\_ in ubd\_, uby\_ in ubd\_], ubf\_(\allowbreak{}uby\_, ubx\_) = conjugate(\allowbreak{}ubf\_(\allowbreak{}ubx\_, uby\_))))
\end{quote}
mentioned by 4 result(s)

\smallskip

\noindent\textbf{\texttt{is-r-net(s, f, a, r)}}\quad\textit{predicate}\index{is-r-net(s, f, a, r)@\texttt{is-r-net(s, f, a, r)}}\\
reads: \$2 is an \$4-net for \$3 in \$1 \\
defined by \texttt{is-r-net} \\
characterized by \texttt{is-r-net}  \textit{(definitional)}:
\begin{quote}\ttfamily\footnotesize\raggedright\noindent
forall(\allowbreak{}[s, f, a, r], is-r-net(\allowbreak{}s, f, a, r) iff f subset a and forall(\allowbreak{}[p in a], forsome(\allowbreak{}[c in f], (\allowbreak{}dist(\allowbreak{}s))(\allowbreak{}c, p) <= r and not(\allowbreak{}(\allowbreak{}dist(\allowbreak{}s))(\allowbreak{}c, p) = r))))
\end{quote}
mentioned by 5 result(s)

\smallskip

\noindent\textbf{\texttt{is-rcont-subalgebra(rst\_, rsa\_)}}\quad\textit{predicate}\index{is-rcont-subalgebra(rst_, rsa_)@\texttt{is-rcont-subalgebra(rst\_, rsa\_)}}\\
reads: \$2 is a unital subalgebra of the continuous real functions on \$1 \\
defined by \texttt{is-rcont-subalgebra} \\
characterized by \texttt{is-rcont-subalgebra}  \textit{(definitional)}:
\begin{quote}\ttfamily\footnotesize\raggedright\noindent
forall(\allowbreak{}[rst\_, rsa\_], is-rcont-subalgebra(\allowbreak{}rst\_, rsa\_) iff rsa\_ subset rcont(\allowbreak{}rst\_) and forall(\allowbreak{}[rsc\_ in rr], vnb-lambda(\allowbreak{}rsz\_, pts(\allowbreak{}rst\_), rsc\_) in rsa\_) and forall(\allowbreak{}[rsf\_ in rsa\_, rsg\_ in rsa\_], vnb-lambda(\allowbreak{}rsz\_, pts(\allowbreak{}rst\_), rsf\_(\allowbreak{}rsz\_) + rsg\_(\allowbreak{}rsz\_)) in rsa\_) and forall(\allowbreak{}[rsf\_ in rsa\_, rsg\_ in rsa\_], vnb-lambda(\allowbreak{}rsz\_, pts(\allowbreak{}rst\_), rsf\_(\allowbreak{}rsz\_) * rsg\_(\allowbreak{}rsz\_)) in rsa\_))
\end{quote}
mentioned by 17 result(s)

\smallskip

\noindent\textbf{\texttt{is-real-hilbert-space(v)}}\quad\textit{predicate}\index{is-real-hilbert-space(v)@\texttt{is-real-hilbert-space(v)}}\\
reads: a real Hilbert space \\
defined by \texttt{is-real-hilbert-space} \\
characterized by \texttt{is-real-hilbert-space}  \textit{(definitional)}:
\begin{quote}\ttfamily\footnotesize\raggedright\noindent
forall(\allowbreak{}[v], is-real-hilbert-space(\allowbreak{}v) iff is-pre-hilbert-space(\allowbreak{}v) and is-complete(\allowbreak{}sip-metric-space(\allowbreak{}v)))
\end{quote}
mentioned by 17 result(s)

\smallskip

\noindent\textbf{\texttt{is-real-linear-op(v, jo)}}\quad\textit{predicate}\index{is-real-linear-op(v, jo)@\texttt{is-real-linear-op(v, jo)}}\\
reads: \$2 is a real-linear operator on \$1 \\
defined by \texttt{is-real-linear-op} \\
characterized by \texttt{is-real-linear-op}  \textit{(definitional)}:
\begin{quote}\ttfamily\footnotesize\raggedright\noindent
forall(\allowbreak{}[v, jo], is-real-linear-op(\allowbreak{}v, jo) iff jo in fun(\allowbreak{}vec(\allowbreak{}v), vec(\allowbreak{}v)) and forall(\allowbreak{}[x\_ in vec(\allowbreak{}v), y\_ in vec(\allowbreak{}v)], jo(\allowbreak{}(\allowbreak{}vadd(\allowbreak{}v))(\allowbreak{}x\_, y\_)) = (\allowbreak{}vadd(\allowbreak{}v))(\allowbreak{}jo(\allowbreak{}x\_), jo(\allowbreak{}y\_))) and forall(\allowbreak{}[a\_ in rr, x\_ in vec(\allowbreak{}v)], jo(\allowbreak{}(\allowbreak{}act(\allowbreak{}v))(\allowbreak{}a\_, x\_)) = (\allowbreak{}act(\allowbreak{}v))(\allowbreak{}a\_, jo(\allowbreak{}x\_))))
\end{quote}
mentioned by 14 result(s)

\smallskip

\noindent\textbf{\texttt{is-real-measurable-fn(omega, ca, f)}}\quad\textit{predicate}\index{is-real-measurable-fn(omega, ca, f)@\texttt{is-real-measurable-fn(omega, ca, f)}}\\
reads: \$3 is a measurable real-valued function on \$1 with \$2 \\
defined by \texttt{is-real-measurable-fn} \\
characterized by \texttt{is-real-measurable-fn}  \textit{(definitional)}:
\begin{quote}\ttfamily\footnotesize\raggedright\noindent
forall(\allowbreak{}[omega, ca, f], is-real-measurable-fn(\allowbreak{}omega, ca, f) iff f in fun(\allowbreak{}omega, rr) and is-measurable-fn(\allowbreak{}omega, ca, pos-part(\allowbreak{}omega, f)) and is-measurable-fn(\allowbreak{}omega, ca, neg-part(\allowbreak{}omega, f)))
\end{quote}
mentioned by 87 result(s)

\smallskip

\noindent\textbf{\texttt{is-real-orthogonal-projection(v, ks, pk)}}\quad\textit{predicate}\index{is-real-orthogonal-projection(v, ks, pk)@\texttt{is-real-orthogonal-projection(v, ks, pk)}}\\
reads: \$3 is the orthogonal projection of \$1 onto \$2 \\
defined by \texttt{is-real-orthogonal-projection} \\
characterized by \texttt{is-real-orthogonal-projection}  \textit{(definitional)}:
\begin{quote}\ttfamily\footnotesize\raggedright\noindent
forall(\allowbreak{}[v, ks, pk], is-real-orthogonal-projection(\allowbreak{}v, ks, pk) iff pk in fun(\allowbreak{}vec(\allowbreak{}v), vec(\allowbreak{}v)) and forall(\allowbreak{}[sop\_ in vec(\allowbreak{}v)], pk(\allowbreak{}sop\_) in ks and (\allowbreak{}vadd(\allowbreak{}v))(\allowbreak{}sop\_, (\allowbreak{}vneg(\allowbreak{}v))(\allowbreak{}pk(\allowbreak{}sop\_))) in sip-perp(\allowbreak{}v, ks)))
\end{quote}
mentioned by 17 result(s)

\smallskip

\noindent\textbf{\texttt{is-received(ssp\_, ssn\_)}}\quad\textit{predicate}\index{is-received(ssp_, ssn_)@\texttt{is-received(ssp\_, ssn\_)}}\\
reads: the node \$2 of \$1 is received \\
defined by \texttt{is-received} \\
characterized by \texttt{is-received}  \textit{(definitional)}:
\begin{quote}\ttfamily\footnotesize\raggedright\noindent
forall(\allowbreak{}[ssp\_, ssn\_], is-received(\allowbreak{}ssp\_, ssn\_) iff node-dir(\allowbreak{}ssp\_, ssn\_) = received)
\end{quote}
mentioned by 14 result(s)

\smallskip

\noindent\textbf{\texttt{is-regular-measure(s, mu)}}\quad\textit{predicate}\index{is-regular-measure(s, mu)@\texttt{is-regular-measure(s, mu)}}\\
reads: \$2 is a regular measure on \$1 \\
defined by \texttt{is-regular-measure} \\
characterized by \texttt{is-regular-measure}  \textit{(definitional)}:
\begin{quote}\ttfamily\footnotesize\raggedright\noindent
forall(\allowbreak{}[s, mu], is-regular-measure(\allowbreak{}s, mu) iff is-measure(\allowbreak{}pts(\allowbreak{}s), borel-sets(\allowbreak{}s), mu) and mu(\allowbreak{}pts(\allowbreak{}s)) in rr and forall(\allowbreak{}[rge\_ in borel-sets(\allowbreak{}s)], mu(\allowbreak{}rge\_) = esup(\allowbreak{}\{rgv\_ in rr-pos-star: forsome(\allowbreak{}[rgk\_ in power(\allowbreak{}pts(\allowbreak{}s))], rgk\_ subset rge\_ and is-compact(\allowbreak{}subspace-ms(\allowbreak{}s, rgk\_)) and rgv\_ = mu(\allowbreak{}rgk\_))\})))
\end{quote}
mentioned by 8 result(s)

\smallskip

\noindent\textbf{\texttt{is-regular-node(sai\_, san\_)}}\quad\textit{predicate}\index{is-regular-node(sai_, san_)@\texttt{is-regular-node(sai\_, san\_)}}\\
reads: \$2 is a regular node of \$1 \\
defined by \texttt{is-regular-node} \\
characterized by \texttt{is-regular-node}  \textit{(definitional)}:
\begin{quote}\ttfamily\footnotesize\raggedright\noindent
forall(\allowbreak{}[sai\_, san\_], is-regular-node(\allowbreak{}sai\_, san\_) iff not(\allowbreak{}nth(\allowbreak{}2, san\_) in adv-ids(\allowbreak{}sai\_)))
\end{quote}
mentioned by 40 result(s)

\smallskip

\noindent\textbf{\texttt{is-regulated(f, a, b)}}\quad\textit{predicate}\index{is-regulated(f, a, b)@\texttt{is-regulated(f, a, b)}}\\
reads: \$1 is regulated on the interval [\$2, \$3] -- both one-sided limits exist at every point \\
defined by \texttt{is-regulated} \\
characterized by \texttt{is-regulated}  \textit{(definitional)}:
\begin{quote}\ttfamily\footnotesize\raggedright\noindent
forall(\allowbreak{}[f, a, b], is-regulated(\allowbreak{}f, a, b) iff f in fun(\allowbreak{}rr, rr) and a in rr and b in rr and a < b and forall(\allowbreak{}[x\_ in ccint(\allowbreak{}a, b)], (\allowbreak{}x\_ < b implies forsome(\allowbreak{}[l\_], is-right-limit(\allowbreak{}f, x\_, l\_))) and (\allowbreak{}a < x\_ implies forsome(\allowbreak{}[l\_], is-left-limit(\allowbreak{}f, x\_, l\_)))))
\end{quote}
also: \texttt{step-fn-is-regulated} \texttt{continuous-is-regulated} \texttt{restrict-regulated-on-is-regulated} \texttt{regulated-extend-const-is-regulated-on} \\
mentioned by 7 result(s)

\smallskip

\noindent\textbf{\texttt{is-regulated-on(rpf\_, rpa\_, rpb\_)}}\quad\textit{predicate}\index{is-regulated-on(rpf_, rpa_, rpb_)@\texttt{is-regulated-on(rpf\_, rpa\_, rpb\_)}}\\
reads: \$1 is regulated on the interval [\$2, \$3] \\
defined by \texttt{is-regulated-on} \\
characterized by \texttt{is-regulated-on}  \textit{(definitional)}:
\begin{quote}\ttfamily\footnotesize\raggedright\noindent
forall(\allowbreak{}[rpf\_, rpa\_, rpb\_], is-regulated-on(\allowbreak{}rpf\_, rpa\_, rpb\_) iff rpa\_ in rr and rpb\_ in rr and rpa\_ < rpb\_ and rpf\_ in fun(\allowbreak{}ccint(\allowbreak{}rpa\_, rpb\_), rr) and forall(\allowbreak{}[rpx\_ in ccint(\allowbreak{}rpa\_, rpb\_)], rpx\_ < rpb\_ implies forsome(\allowbreak{}[rpl\_], is-right-limit-within(\allowbreak{}rpf\_, ccint(\allowbreak{}rpa\_, rpb\_), rpx\_, rpl\_))) and forall(\allowbreak{}[rpx\_ in ccint(\allowbreak{}rpa\_, rpb\_)], rpa\_ < rpx\_ implies forsome(\allowbreak{}[rpl\_], is-left-limit-within(\allowbreak{}rpf\_, ccint(\allowbreak{}rpa\_, rpb\_), rpx\_, rpl\_))))
\end{quote}
also: \texttt{regulated-extend-const-is-regulated-on} \\
mentioned by 33 result(s)

\smallskip

\noindent\textbf{\texttt{is-removable-singularity(mrf\_, mra\_)}}\quad\textit{predicate}\index{is-removable-singularity(mrf_, mra_)@\texttt{is-removable-singularity(mrf\_, mra\_)}}\\
reads: \$1 has a removable singularity at \$2 \\
defined by \texttt{is-removable-singularity} \\
characterized by \texttt{is-removable-singularity}  \textit{(definitional)}:
\begin{quote}\ttfamily\footnotesize\raggedright\noindent
forall(\allowbreak{}[mrf\_, mra\_], is-removable-singularity(\allowbreak{}mrf\_, mra\_) iff mra\_ in cc and forsome(\allowbreak{}[mrr\_ in rr], 0 < mrr\_ and forsome(\allowbreak{}[mrg\_], holomorphic-on(\allowbreak{}ball(\allowbreak{}cc-ms, mra\_, mrr\_), mrg\_) and forall(\allowbreak{}[mrz\_ in difference(\allowbreak{}ball(\allowbreak{}cc-ms, mra\_, mrr\_), singleton(\allowbreak{}mra\_))], mrf\_(\allowbreak{}mrz\_) = mrg\_(\allowbreak{}mrz\_)))))
\end{quote}
mentioned by 12 result(s)

\smallskip

\noindent\textbf{\texttt{is-representation(g4a\_, g4v\_, g4p\_)}}\quad\textit{predicate}\index{is-representation(g4a_, g4v_, g4p_)@\texttt{is-representation(g4a\_, g4v\_, g4p\_)}}\\
reads: \$3 is a representation of \$1 on \$2 \\
defined by \texttt{is-representation} \\
characterized by \texttt{is-representation}  \textit{(definitional)}:
\begin{quote}\ttfamily\footnotesize\raggedright\noindent
forall(\allowbreak{}[g4a\_, g4v\_, g4p\_], is-representation(\allowbreak{}g4a\_, g4v\_, g4p\_) iff is-hilbert-space(\allowbreak{}g4v\_) and is-star-hom(\allowbreak{}g4a\_, bh-algebra(\allowbreak{}g4v\_), g4p\_))
\end{quote}
mentioned by 15 result(s)

\smallskip

\noindent\textbf{\texttt{is-resp-strand(nlp\_, nls\_, nla\_, nlb\_, nlm\_, nln\_)}}\quad\textit{predicate}\index{is-resp-strand(nlp_, nls_, nla_, nlb_, nlm_, nln_)@\texttt{is-resp-strand(nlp\_, nls\_, nla\_, nlb\_, nlm\_, nln\_)}}\\
reads: \$2 is a responder's strand of \$1 for the names \$3, \$4 and the values \$5, \$6 \\
defined by \texttt{is-resp-strand} \\
characterized by \texttt{is-resp-strand}  \textit{(definitional)}:
\begin{quote}\ttfamily\footnotesize\raggedright\noindent
forall(\allowbreak{}[nlp\_, nls\_, nla\_, nlb\_, nlm\_, nln\_], is-resp-strand(\allowbreak{}nlp\_, nls\_, nla\_, nlb\_, nlm\_, nln\_) iff nla\_ in names(\allowbreak{}nlp\_) and nlb\_ in names(\allowbreak{}nlp\_) and nlm\_ in tokens(\allowbreak{}ispace(\allowbreak{}pk-inf(\allowbreak{}nlp\_))) and nln\_ in tokens(\allowbreak{}ispace(\allowbreak{}pk-inf(\allowbreak{}nlp\_))) and not(\allowbreak{}nln\_ in names(\allowbreak{}nlp\_)) and not(\allowbreak{}nlm\_ in keys(\allowbreak{}ispace(\allowbreak{}pk-inf(\allowbreak{}nlp\_)))) and not(\allowbreak{}nln\_ in keys(\allowbreak{}ispace(\allowbreak{}pk-inf(\allowbreak{}nlp\_)))) and nls\_ = [[received, encr(\allowbreak{}lst(\allowbreak{}[nlm\_, nla\_]), (\allowbreak{}pubkey(\allowbreak{}nlp\_))(\allowbreak{}nlb\_))], [sent, encr(\allowbreak{}lst(\allowbreak{}[nlm\_, nln\_, nlb\_]), (\allowbreak{}pubkey(\allowbreak{}nlp\_))(\allowbreak{}nla\_))], [received, encr(\allowbreak{}nln\_, (\allowbreak{}pubkey(\allowbreak{}nlp\_))(\allowbreak{}nlb\_))]])
\end{quote}
mentioned by 15 result(s)

\smallskip

\noindent\textbf{\texttt{is-right-ideal(baa\_, bai\_)}}\quad\textit{predicate}\index{is-right-ideal(baa_, bai_)@\texttt{is-right-ideal(baa\_, bai\_)}}\\
reads: \$2 is a right ideal of \$1 \\
defined by \texttt{is-right-ideal} \\
characterized by \texttt{is-right-ideal}  \textit{(definitional)}:
\begin{quote}\ttfamily\footnotesize\raggedright\noindent
forall(\allowbreak{}[baa\_, bai\_], is-right-ideal(\allowbreak{}baa\_, bai\_) iff is-alg-subspace(\allowbreak{}baa\_, bai\_) and forall(\allowbreak{}[bay\_ in vec(\allowbreak{}baa\_), bax\_ in bai\_], (\allowbreak{}vmul(\allowbreak{}baa\_))(\allowbreak{}bax\_, bay\_) in bai\_))
\end{quote}
mentioned by 6 result(s)

\smallskip

\noindent\textbf{\texttt{is-right-limit(f, x, l)}}\quad\textit{predicate}\index{is-right-limit(f, x, l)@\texttt{is-right-limit(f, x, l)}}\\
reads: \$3 is the limit of \$1 as the argument approaches \$2 from the right \\
defined by \texttt{is-right-limit} \\
characterized by \texttt{is-right-limit}  \textit{(definitional)}:
\begin{quote}\ttfamily\footnotesize\raggedright\noindent
forall(\allowbreak{}[f, x, l], is-right-limit(\allowbreak{}f, x, l) iff f in fun(\allowbreak{}rr, rr) and x in rr and l in rr and forall(\allowbreak{}[eps\_], pos-rr(\allowbreak{}eps\_) implies forsome(\allowbreak{}[delta\_], pos-rr(\allowbreak{}delta\_) and forall(\allowbreak{}[t\_], t\_ in rr and x < t\_ and t\_ <= x + delta\_ implies abs(\allowbreak{}f(\allowbreak{}t\_) - l) <= eps\_))))
\end{quote}
mentioned by 6 result(s)

\smallskip

\noindent\textbf{\texttt{is-right-limit-within(rpf\_, rpw\_, rpx\_, rpl\_)}}\quad\textit{predicate}\index{is-right-limit-within(rpf_, rpw_, rpx_, rpl_)@\texttt{is-right-limit-within(rpf\_, rpw\_, rpx\_, rpl\_)}}\\
reads: \$1 has right limit \$4 at \$3 within \$2 \\
defined by \texttt{is-right-limit-within} \\
characterized by \texttt{is-right-limit-within}  \textit{(definitional)}:
\begin{quote}\ttfamily\footnotesize\raggedright\noindent
forall(\allowbreak{}[rpf\_, rpw\_, rpx\_, rpl\_], is-right-limit-within(\allowbreak{}rpf\_, rpw\_, rpx\_, rpl\_) iff rpw\_ subset rr and rpf\_ in fun(\allowbreak{}rpw\_, rr) and rpx\_ in rr and rpl\_ in rr and forall(\allowbreak{}[rpe\_], pos-rr(\allowbreak{}rpe\_) implies forsome(\allowbreak{}[rpdl\_], pos-rr(\allowbreak{}rpdl\_) and forall(\allowbreak{}[rpt\_ in rpw\_], rpx\_ < rpt\_ implies rpt\_ < rpx\_ + rpdl\_ implies abs(\allowbreak{}rpf\_(\allowbreak{}rpt\_) - rpl\_) < rpe\_))))
\end{quote}
mentioned by 17 result(s)

\smallskip

\noindent\textbf{\texttt{is-ring(s)}}\quad\textit{predicate}\index{is-ring(s)@\texttt{is-ring(s)}}\\
reads: \$1 is a ring \\
a predicate \\
characterized by \texttt{is-ring}  \textit{(definitional)}:
\begin{quote}\ttfamily\footnotesize\raggedright\noindent
forall(\allowbreak{}[s], is-ring(\allowbreak{}s) iff length(\allowbreak{}s) = 6 and carr(\allowbreak{}s) in set and add(\allowbreak{}s) in fun(\allowbreak{}cartesian(\allowbreak{}carr(\allowbreak{}s), carr(\allowbreak{}s)), carr(\allowbreak{}s)) and mul(\allowbreak{}s) in fun(\allowbreak{}cartesian(\allowbreak{}carr(\allowbreak{}s), carr(\allowbreak{}s)), carr(\allowbreak{}s)) and neg(\allowbreak{}s) in fun(\allowbreak{}carr(\allowbreak{}s), carr(\allowbreak{}s)) and zero(\allowbreak{}s) in carr(\allowbreak{}s) and one(\allowbreak{}s) in carr(\allowbreak{}s) and is-associative(\allowbreak{}add(\allowbreak{}s), carr(\allowbreak{}s)) and is-commutative(\allowbreak{}add(\allowbreak{}s), carr(\allowbreak{}s)) and is-identity(\allowbreak{}add(\allowbreak{}s), zero(\allowbreak{}s), carr(\allowbreak{}s)) and has-inverses(\allowbreak{}add(\allowbreak{}s), zero(\allowbreak{}s), neg(\allowbreak{}s), carr(\allowbreak{}s)) and is-associative(\allowbreak{}mul(\allowbreak{}s), carr(\allowbreak{}s)) and is-identity(\allowbreak{}mul(\allowbreak{}s), one(\allowbreak{}s), carr(\allowbreak{}s)) and is-distributive(\allowbreak{}add(\allowbreak{}s), mul(\allowbreak{}s), carr(\allowbreak{}s)) and s in set)
\end{quote}
also: \texttt{uba-as-ring-is-ring} \texttt{ringoid-as-ring-is-ring} \texttt{poly-is-ring-uba-as-ring} \texttt{mat-ring-is-ring-uba-as-ring} \texttt{ring-prod-is-ring-uba-as-ring} \texttt{uba-as-ring-is-ring-uiba-as-uba} \texttt{poly-is-ring-uba-as-ring-uiba-as-uba} \texttt{mat-ring-is-ring-uba-as-ring-uiba-as-uba} ... (26 in all) \\
mentioned by 1020 result(s)

\smallskip

\noindent\textbf{\texttt{is-ring-iso(rib\_, ric\_, rif\_)}}\quad\textit{predicate}\index{is-ring-iso(rib_, ric_, rif_)@\texttt{is-ring-iso(rib\_, ric\_, rif\_)}}\\
reads: \$3 is a ring isomorphism from \$1 onto \$2 \\
defined by \texttt{is-ring-iso} \\
characterized by \texttt{is-ring-iso}  \textit{(definitional)}:
\begin{quote}\ttfamily\footnotesize\raggedright\noindent
forall(\allowbreak{}[rib\_, ric\_, rif\_], is-ring-iso(\allowbreak{}rib\_, ric\_, rif\_) iff is-hom-ring(\allowbreak{}rib\_, ric\_, rif\_) and forall(\allowbreak{}[riu\_ in carr(\allowbreak{}rib\_), riv\_ in carr(\allowbreak{}rib\_)], rif\_(\allowbreak{}riu\_) = rif\_(\allowbreak{}riv\_) implies riu\_ = riv\_) and forall(\allowbreak{}[riw\_ in carr(\allowbreak{}ric\_)], forsome(\allowbreak{}[riz\_ in carr(\allowbreak{}rib\_)], rif\_(\allowbreak{}riz\_) = riw\_)))
\end{quote}
mentioned by 6 result(s)

\smallskip

\noindent\textbf{\texttt{is-ringoid(s)}}\quad\textit{predicate}\index{is-ringoid(s)@\texttt{is-ringoid(s)}}\\
reads: a ringoid \\
a predicate \\
characterized by \texttt{is-ringoid}  \textit{(definitional)}:
\begin{quote}\ttfamily\footnotesize\raggedright\noindent
forall(\allowbreak{}[s], is-ringoid(\allowbreak{}s) iff length(\allowbreak{}s) = 7 and carr(\allowbreak{}s) in set and add(\allowbreak{}s) in fun(\allowbreak{}cartesian(\allowbreak{}carr(\allowbreak{}s), carr(\allowbreak{}s)), carr(\allowbreak{}s)) and mul(\allowbreak{}s) in fun(\allowbreak{}cartesian(\allowbreak{}carr(\allowbreak{}s), carr(\allowbreak{}s)), carr(\allowbreak{}s)) and neg(\allowbreak{}s) in fun(\allowbreak{}carr(\allowbreak{}s), carr(\allowbreak{}s)) and zero(\allowbreak{}s) in carr(\allowbreak{}s) and one(\allowbreak{}s) in carr(\allowbreak{}s) and idl(\allowbreak{}s) in power(\allowbreak{}carr(\allowbreak{}s)) and is-associative(\allowbreak{}add(\allowbreak{}s), carr(\allowbreak{}s)) and is-commutative(\allowbreak{}add(\allowbreak{}s), carr(\allowbreak{}s)) and is-identity(\allowbreak{}add(\allowbreak{}s), zero(\allowbreak{}s), carr(\allowbreak{}s)) and has-inverses(\allowbreak{}add(\allowbreak{}s), zero(\allowbreak{}s), neg(\allowbreak{}s), carr(\allowbreak{}s)) and is-associative(\allowbreak{}mul(\allowbreak{}s), carr(\allowbreak{}s)) and is-identity(\allowbreak{}mul(\allowbreak{}s), one(\allowbreak{}s), carr(\allowbreak{}s)) and is-distributive(\allowbreak{}add(\allowbreak{}s), mul(\allowbreak{}s), carr(\allowbreak{}s)) and is-ideal-in(\allowbreak{}idl(\allowbreak{}s), carr(\allowbreak{}s), add(\allowbreak{}s), mul(\allowbreak{}s), neg(\allowbreak{}s), zero(\allowbreak{}s)) and s in set)
\end{quote}
also: \texttt{fp-is-ringoid} \texttt{rs-is-ringoid} \texttt{is-ringoid-is-set} \texttt{qp-csr-is-ringoid} \\
mentioned by 112 result(s)

\smallskip

\noindent\textbf{\texttt{is-road(pgam, dgam, a, b)}}\quad\textit{predicate}\index{is-road(pgam, dgam, a, b)@\texttt{is-road(pgam, dgam, a, b)}}\\
reads: \$1 is a road with derivative \$2 on the interval [\$3, \$4] \\
defined by \texttt{is-road} \\
characterized by \texttt{is-road}  \textit{(definitional)}:
\begin{quote}\ttfamily\footnotesize\raggedright\noindent
forall(\allowbreak{}[pgam, dgam, a, b], is-road(\allowbreak{}pgam, dgam, a, b) iff is-path(\allowbreak{}pgam, a, b) and dgam in fun(\allowbreak{}ccint(\allowbreak{}a, b), cc) and is-regulated-on(\allowbreak{}vnb-lambda(\allowbreak{}pat\_, ccint(\allowbreak{}a, b), real-part(\allowbreak{}dgam(\allowbreak{}pat\_))), a, b) and is-regulated-on(\allowbreak{}vnb-lambda(\allowbreak{}pat\_, ccint(\allowbreak{}a, b), imag-part(\allowbreak{}dgam(\allowbreak{}pat\_))), a, b) and is-primitive(\allowbreak{}vnb-lambda(\allowbreak{}pat\_, ccint(\allowbreak{}a, b), real-part(\allowbreak{}pgam(\allowbreak{}pat\_))), vnb-lambda(\allowbreak{}pat\_, ccint(\allowbreak{}a, b), real-part(\allowbreak{}dgam(\allowbreak{}pat\_))), a, b) and is-primitive(\allowbreak{}vnb-lambda(\allowbreak{}pat\_, ccint(\allowbreak{}a, b), imag-part(\allowbreak{}pgam(\allowbreak{}pat\_))), vnb-lambda(\allowbreak{}pat\_, ccint(\allowbreak{}a, b), imag-part(\allowbreak{}dgam(\allowbreak{}pat\_))), a, b))
\end{quote}
also: \texttt{juxta-is-road} \texttt{circle-is-road} \texttt{aff-seg-is-road} \texttt{is-road-is-path} \texttt{segment-is-road} \texttt{opposite-is-road} \texttt{seg-path-is-road} \texttt{tri-road-is-road} ... (15 in all) \\
mentioned by 161 result(s)

\smallskip

\noindent\textbf{\texttt{is-self-adjoint-el(baa\_, bax\_)}}\quad\textit{predicate}\index{is-self-adjoint-el(baa_, bax_)@\texttt{is-self-adjoint-el(baa\_, bax\_)}}\\
reads: \$2 is self-adjoint in \$1 \\
defined by \texttt{is-self-adjoint-el} \\
characterized by \texttt{is-self-adjoint-el}  \textit{(definitional)}:
\begin{quote}\ttfamily\footnotesize\raggedright\noindent
forall(\allowbreak{}[baa\_, bax\_], is-self-adjoint-el(\allowbreak{}baa\_, bax\_) iff bax\_ in vec(\allowbreak{}baa\_) and (\allowbreak{}vstar(\allowbreak{}baa\_))(\allowbreak{}bax\_) = bax\_)
\end{quote}
mentioned by 16 result(s)

\smallskip

\noindent\textbf{\texttt{is-self-adjoint-graph(ubv\_, ubg\_)}}\quad\textit{predicate}\index{is-self-adjoint-graph(ubv_, ubg_)@\texttt{is-self-adjoint-graph(ubv\_, ubg\_)}}\\
reads: \$2 is a self-adjoint graph \\
defined by \texttt{is-self-adjoint-graph} \\
characterized by \texttt{is-self-adjoint-graph}  \textit{(definitional)}:
\begin{quote}\ttfamily\footnotesize\raggedright\noindent
forall(\allowbreak{}[ubv\_, ubg\_], is-self-adjoint-graph(\allowbreak{}ubv\_, ubg\_) iff graph-adjoint(\allowbreak{}ubv\_, ubg\_) = ubg\_)
\end{quote}
mentioned by 16 result(s)

\smallskip

\noindent\textbf{\texttt{is-self-adjoint-op(kqv\_, kqt\_)}}\quad\textit{predicate}\index{is-self-adjoint-op(kqv_, kqt_)@\texttt{is-self-adjoint-op(kqv\_, kqt\_)}}\\
reads: \$2 is a self-adjoint operator on \$1 \\
defined by \texttt{is-self-adjoint-op} \\
characterized by \texttt{is-self-adjoint-op}  \textit{(definitional)}:
\begin{quote}\ttfamily\footnotesize\raggedright\noindent
forall(\allowbreak{}[kqv\_, kqt\_], is-self-adjoint-op(\allowbreak{}kqv\_, kqt\_) iff kqt\_ in bounded-ops(\allowbreak{}kqv\_, kqv\_) and forall(\allowbreak{}[kqx\_ in vec(\allowbreak{}kqv\_), kqy\_ in vec(\allowbreak{}kqv\_)], (\allowbreak{}ip(\allowbreak{}kqv\_))(\allowbreak{}kqt\_(\allowbreak{}kqx\_), kqy\_) = (\allowbreak{}ip(\allowbreak{}kqv\_))(\allowbreak{}kqx\_, kqt\_(\allowbreak{}kqy\_))))
\end{quote}
mentioned by 31 result(s)

\smallskip

\noindent\textbf{\texttt{is-semi-inner-product-space(s)}}\quad\textit{predicate}\index{is-semi-inner-product-space(s)@\texttt{is-semi-inner-product-space(s)}}\\
reads: a semi-inner-product space \\
a predicate \\
characterized by \texttt{is-semi-inner-product-space}  \textit{(definitional)}:
\begin{quote}\ttfamily\footnotesize\raggedright\noindent
forall(\allowbreak{}[s], is-semi-inner-product-space(\allowbreak{}s) iff length(\allowbreak{}s) = 7 and is-ring(\allowbreak{}scal(\allowbreak{}s)) and vec(\allowbreak{}s) in set and vadd(\allowbreak{}s) in fun(\allowbreak{}cartesian(\allowbreak{}vec(\allowbreak{}s), vec(\allowbreak{}s)), vec(\allowbreak{}s)) and vzero(\allowbreak{}s) in vec(\allowbreak{}s) and vneg(\allowbreak{}s) in fun(\allowbreak{}vec(\allowbreak{}s), vec(\allowbreak{}s)) and act(\allowbreak{}s) in fun(\allowbreak{}cartesian(\allowbreak{}carr(\allowbreak{}scal(\allowbreak{}s)), vec(\allowbreak{}s)), vec(\allowbreak{}s)) and rip(\allowbreak{}s) in fun(\allowbreak{}cartesian(\allowbreak{}vec(\allowbreak{}s), vec(\allowbreak{}s)), rr) and is-associative(\allowbreak{}vadd(\allowbreak{}s), vec(\allowbreak{}s)) and is-commutative(\allowbreak{}vadd(\allowbreak{}s), vec(\allowbreak{}s)) and is-identity(\allowbreak{}vadd(\allowbreak{}s), vzero(\allowbreak{}s), vec(\allowbreak{}s)) and has-inverses(\allowbreak{}vadd(\allowbreak{}s), vzero(\allowbreak{}s), vneg(\allowbreak{}s), vec(\allowbreak{}s)) and scal(\allowbreak{}s) = normed-field-as-commutative-ring(\allowbreak{}rr-normed-field) and carr(\allowbreak{}scal(\allowbreak{}s)) = rr and forall(\allowbreak{}[r\_ in carr(\allowbreak{}scal(\allowbreak{}s)), x\_ in vec(\allowbreak{}s), y\_ in vec(\allowbreak{}s)], (\allowbreak{}act(\allowbreak{}s))(\allowbreak{}r\_, (\allowbreak{}vadd(\allowbreak{}s))(\allowbreak{}x\_, y\_)) = (\allowbreak{}vadd(\allowbreak{}s))(\allowbreak{}(\allowbreak{}act(\allowbreak{}s))(\allowbreak{}r\_, x\_), (\allowbreak{}act(\allowbreak{}s))(\allowbreak{}r\_, y\_))) and forall(\allowbreak{}[r\_ in carr(\allowbreak{}scal(\allowbreak{}s)), s\_ in carr(\allowbreak{}scal(\allowbreak{}s)), x\_ in vec(\allowbreak{}s)], (\allowbreak{}act(\allowbreak{}s))(\allowbreak{}(\allowbreak{}add(\allowbreak{}scal(\allowbreak{}s)))(\allowbreak{}r\_, s\_), x\_) = (\allowbreak{}vadd(\allowbreak{}s))(\allowbreak{}(\allowbreak{}act(\allowbreak{}s))(\allowbreak{}r\_, x\_), (\allowbreak{}act(\allowbreak{}s))(\allowbreak{}s\_, x\_))) and forall(\allowbreak{}[r\_ in carr(\allowbreak{}scal(\allowbreak{}s)), s\_ in carr(\allowbreak{}scal(\allowbreak{}s)), x\_ in vec(\allowbreak{}s)], (\allowbreak{}act(\allowbreak{}s))(\allowbreak{}(\allowbreak{}mul(\allowbreak{}scal(\allowbreak{}s)))(\allowbreak{}r\_, s\_), x\_) = (\allowbreak{}act(\allowbreak{}s))(\allowbreak{}r\_, (\allowbreak{}act(\allowbreak{}s))(\allowbreak{}s\_, x\_))) and forall(\allowbreak{}[x\_ in vec(\allowbreak{}s)], (\allowbreak{}act(\allowbreak{}s))(\allowbreak{}one(\allowbreak{}scal(\allowbreak{}s)), x\_) = x\_) and forall(\allowbreak{}[x\_ in vec(\allowbreak{}s), y\_ in vec(\allowbreak{}s)], (\allowbreak{}rip(\allowbreak{}s))(\allowbreak{}x\_, y\_) = (\allowbreak{}rip(\allowbreak{}s))(\allowbreak{}y\_, x\_)) and forall(\allowbreak{}[x\_ in vec(\allowbreak{}s), y\_ in vec(\allowbreak{}s), z\_ in vec(\allowbreak{}s)], (\allowbreak{}rip(\allowbreak{}s))(\allowbreak{}(\allowbreak{}vadd(\allowbreak{}s))(\allowbreak{}x\_, y\_), z\_) = (\allowbreak{}rip(\allowbreak{}s))(\allowbreak{}x\_, z\_) + (\allowbreak{}rip(\allowbreak{}s))(\allowbreak{}y\_, z\_)) and forall(\allowbreak{}[a\_ in rr, x\_ in vec(\allowbreak{}s), y\_ in vec(\allowbreak{}s)], (\allowbreak{}rip(\allowbreak{}s))(\allowbreak{}(\allowbreak{}act(\allowbreak{}s))(\allowbreak{}a\_, x\_), y\_) = a\_ * (\allowbreak{}rip(\allowbreak{}s))(\allowbreak{}x\_, y\_)) and forall(\allowbreak{}[x\_ in vec(\allowbreak{}s)], 0 <= (\allowbreak{}rip(\allowbreak{}s))(\allowbreak{}x\_, x\_)) and s in set)
\end{quote}
also: \texttt{csip-as-sip-is-semi-inner-product-space} \texttt{is-semi-inner-product-space-is-set} \\
mentioned by 930 result(s)

\smallskip

\noindent\textbf{\texttt{is-semigroup(s)}}\quad\textit{predicate}\index{is-semigroup(s)@\texttt{is-semigroup(s)}}\\
reads: a semigroup \\
a predicate \\
characterized by \texttt{is-semigroup}  \textit{(definitional)}:
\begin{quote}\ttfamily\footnotesize\raggedright\noindent
forall(\allowbreak{}[s], is-semigroup(\allowbreak{}s) iff length(\allowbreak{}s) = 2 and carr(\allowbreak{}s) in set and opr(\allowbreak{}s) in fun(\allowbreak{}cartesian(\allowbreak{}carr(\allowbreak{}s), carr(\allowbreak{}s)), carr(\allowbreak{}s)) and is-associative(\allowbreak{}opr(\allowbreak{}s), carr(\allowbreak{}s)) and s in set)
\end{quote}
also: \texttt{is-semigroup-is-set} \\
mentioned by 13 result(s)

\smallskip

\noindent\textbf{\texttt{is-seminorm(fld, m, p)}}\quad\textit{predicate}\index{is-seminorm(fld, m, p)@\texttt{is-seminorm(fld, m, p)}}\\
reads: \$3 is a seminorm on \$2 over the normed field \$1 \\
defined by \texttt{is-seminorm} \\
characterized by \texttt{is-seminorm}  \textit{(definitional)}:
\begin{quote}\ttfamily\footnotesize\raggedright\noindent
forall(\allowbreak{}[fld, m, p], is-seminorm(\allowbreak{}fld, m, p) iff is-normed-field(\allowbreak{}fld) and is-module(\allowbreak{}m) and scal(\allowbreak{}m) = normed-field-as-commutative-ring(\allowbreak{}fld) and p in fun(\allowbreak{}vec(\allowbreak{}m), rr) and forall(\allowbreak{}[x in vec(\allowbreak{}m)], 0 <= p(\allowbreak{}x)) and forall(\allowbreak{}[x in vec(\allowbreak{}m), y in vec(\allowbreak{}m)], p(\allowbreak{}(\allowbreak{}vadd(\allowbreak{}m))(\allowbreak{}x, y)) <= p(\allowbreak{}x) + p(\allowbreak{}y)) and forall(\allowbreak{}[lam in carr(\allowbreak{}scal(\allowbreak{}m)), x in vec(\allowbreak{}m)], p(\allowbreak{}(\allowbreak{}act(\allowbreak{}m))(\allowbreak{}lam, x)) = (\allowbreak{}fnrm(\allowbreak{}fld))(\allowbreak{}lam) * p(\allowbreak{}x)))
\end{quote}
also: \texttt{is-seminorm-family} \texttt{is-seminorm-family-on} \texttt{nvs-norm-is-seminorm-cnvs-as-nvs} \texttt{nvs-norm-is-seminorm} \texttt{minkowski-functional-is-seminorm} \\
mentioned by 71 result(s)

\smallskip

\noindent\textbf{\texttt{is-seminorm-family(fld, m, fam)}}\quad\textit{predicate}\index{is-seminorm-family(fld, m, fam)@\texttt{is-seminorm-family(fld, m, fam)}}\\
reads: \$3 is a separating countable family of seminorms on \$2 over \$1 \\
defined by \texttt{is-seminorm-family} \\
characterized by \texttt{is-seminorm-family}  \textit{(definitional)}:
\begin{quote}\ttfamily\footnotesize\raggedright\noindent
forall(\allowbreak{}[fld, m, fam], is-seminorm-family(\allowbreak{}fld, m, fam) iff is-normed-field(\allowbreak{}fld) and is-module(\allowbreak{}m) and scal(\allowbreak{}m) = normed-field-as-commutative-ring(\allowbreak{}fld) and fam in fun(\allowbreak{}nn, fun(\allowbreak{}vec(\allowbreak{}m), rr)) and forall(\allowbreak{}[k in nn], is-seminorm(\allowbreak{}fld, m, fam(\allowbreak{}k))) and forall(\allowbreak{}[x in vec(\allowbreak{}m)], not(\allowbreak{}x = vzero(\allowbreak{}m)) implies forsome(\allowbreak{}[k in nn], not(\allowbreak{}(\allowbreak{}fam(\allowbreak{}k))(\allowbreak{}x) = 0))))
\end{quote}
mentioned by 56 result(s)

\smallskip

\noindent\textbf{\texttt{is-seminorm-family-on(fld, m, ix, pf)}}\quad\textit{predicate}\index{is-seminorm-family-on(fld, m, ix, pf)@\texttt{is-seminorm-family-on(fld, m, ix, pf)}}\\
reads: \$4 is a separating family of seminorms on \$2 over \$1 indexed by \$3 \\
defined by \texttt{is-seminorm-family-on} \\
characterized by \texttt{is-seminorm-family-on}  \textit{(definitional)}:
\begin{quote}\ttfamily\footnotesize\raggedright\noindent
forall(\allowbreak{}[fld, m, ix, pf], is-seminorm-family-on(\allowbreak{}fld, m, ix, pf) iff is-normed-field(\allowbreak{}fld) and is-module(\allowbreak{}m) and scal(\allowbreak{}m) = normed-field-as-commutative-ring(\allowbreak{}fld) and ix in set and pf in fun(\allowbreak{}ix, fun(\allowbreak{}vec(\allowbreak{}m), rr)) and forall(\allowbreak{}[sni\_ in ix], is-seminorm(\allowbreak{}fld, m, pf(\allowbreak{}sni\_))) and forall(\allowbreak{}[snx\_ in vec(\allowbreak{}m)], forall(\allowbreak{}[sni\_ in ix], (\allowbreak{}pf(\allowbreak{}sni\_))(\allowbreak{}snx\_) = 0) implies snx\_ = vzero(\allowbreak{}m)))
\end{quote}
mentioned by 37 result(s)

\smallskip

\noindent\textbf{\texttt{is-seminorm-on(m, s, p)}}\quad\textit{predicate}\index{is-seminorm-on(m, s, p)@\texttt{is-seminorm-on(m, s, p)}}\\
reads: \$3 is a seminorm on the subspace \$2 of \$1 \\
defined by \texttt{is-seminorm-on} \\
characterized by \texttt{is-seminorm-on}  \textit{(definitional)}:
\begin{quote}\ttfamily\footnotesize\raggedright\noindent
forall(\allowbreak{}[m, s, p], is-seminorm-on(\allowbreak{}m, s, p) iff p in fun(\allowbreak{}s, rr) and forall(\allowbreak{}[snx\_ in s], 0 <= p(\allowbreak{}snx\_)) and forall(\allowbreak{}[snx\_ in s, sny\_ in s], p(\allowbreak{}(\allowbreak{}vadd(\allowbreak{}m))(\allowbreak{}snx\_, sny\_)) <= p(\allowbreak{}snx\_) + p(\allowbreak{}sny\_)) and forall(\allowbreak{}[snl\_ in rr, snx\_ in s], p(\allowbreak{}(\allowbreak{}act(\allowbreak{}m))(\allowbreak{}snl\_, snx\_)) = abs(\allowbreak{}snl\_) * p(\allowbreak{}snx\_)))
\end{quote}
mentioned by 20 result(s)

\smallskip

\noindent\textbf{\texttt{is-sent(ssp\_, ssn\_)}}\quad\textit{predicate}\index{is-sent(ssp_, ssn_)@\texttt{is-sent(ssp\_, ssn\_)}}\\
reads: the node \$2 of \$1 is sent \\
defined by \texttt{is-sent} \\
characterized by \texttt{is-sent}  \textit{(definitional)}:
\begin{quote}\ttfamily\footnotesize\raggedright\noindent
forall(\allowbreak{}[ssp\_, ssn\_], is-sent(\allowbreak{}ssp\_, ssn\_) iff node-dir(\allowbreak{}ssp\_, ssn\_) = sent)
\end{quote}
mentioned by 32 result(s)

\smallskip

\noindent\textbf{\texttt{is-separable(s)}}\quad\textit{predicate}\index{is-separable(s)@\texttt{is-separable(s)}}\\
reads: \$1 is separable \\
defined by \texttt{is-separable} \\
characterized by \texttt{is-separable}  \textit{(definitional)}:
\begin{quote}\ttfamily\footnotesize\raggedright\noindent
forall(\allowbreak{}[s], is-separable(\allowbreak{}s) iff is-metric-space(\allowbreak{}s) and forsome(\allowbreak{}[dseq in fun(\allowbreak{}nn, pts(\allowbreak{}s))], forall(\allowbreak{}[x in pts(\allowbreak{}s), eps], pos-rr(\allowbreak{}eps) implies forsome(\allowbreak{}[n in nn], (\allowbreak{}dist(\allowbreak{}s))(\allowbreak{}x, dseq(\allowbreak{}n)) < eps))))
\end{quote}
also: \texttt{compact-metric-is-separable} \\
mentioned by 48 result(s)

\smallskip

\noindent\textbf{\texttt{is-separate-trace(sasp\_, sas\_)}}\quad\textit{predicate}\index{is-separate-trace(sasp_, sas_)@\texttt{is-separate-trace(sasp\_, sas\_)}}\\
reads: \$2 is a separation trace in the strand space \$1 \\
defined by \texttt{is-separate-trace} \\
characterized by \texttt{is-separate-trace}  \textit{(definitional)}:
\begin{quote}\ttfamily\footnotesize\raggedright\noindent
forall(\allowbreak{}[sasp\_, sas\_], is-separate-trace(\allowbreak{}sasp\_, sas\_) iff forsome(\allowbreak{}[sal\_ in tuples(\allowbreak{}msg-set(\allowbreak{}tokens(\allowbreak{}sasp\_), keys(\allowbreak{}sasp\_)))], length(\allowbreak{}sas\_) = succ(\allowbreak{}length(\allowbreak{}sal\_)) and nth(\allowbreak{}1, sas\_) = [received, lst(\allowbreak{}sal\_)] and forall(\allowbreak{}[sai\_ in nn], 1 <= sai\_ implies sai\_ <= length(\allowbreak{}sal\_) implies nth(\allowbreak{}succ(\allowbreak{}sai\_), sas\_) = [sent, nth(\allowbreak{}sai\_, sal\_)])))
\end{quote}
mentioned by 6 result(s)

\smallskip

\noindent\textbf{\texttt{is-separating-character(sck\_, scl\_)}}\quad\textit{predicate}\index{is-separating-character(sck_, scl_)@\texttt{is-separating-character(sck\_, scl\_)}}\\
reads: \$2 is a separating character of \$1 (a character whose kernel is the integral elements) \\
defined by \texttt{is-separating-character} \\
characterized by \texttt{is-separating-character}  \textit{(definitional)}:
\begin{quote}\ttfamily\footnotesize\raggedright\noindent
forall(\allowbreak{}[sck\_, scl\_], is-separating-character(\allowbreak{}sck\_, scl\_) iff is-add-character(\allowbreak{}sck\_, scl\_) and forall(\allowbreak{}[scx\_ in carr(\allowbreak{}sck\_)], scl\_(\allowbreak{}scx\_) = 1 iff scx\_ in nv-integers(\allowbreak{}sck\_)))
\end{quote}
mentioned by 5 result(s)

\smallskip

\noindent\textbf{\texttt{is-separating-for(vnv\_, vns\_, vnx\_)}}\quad\textit{predicate}\index{is-separating-for(vnv_, vns_, vnx_)@\texttt{is-separating-for(vnv\_, vns\_, vnx\_)}}\\
reads: \$3 is a separating vector for \$2 \\
defined by \texttt{is-separating-for} \\
characterized by \texttt{is-separating-for}  \textit{(definitional)}:
\begin{quote}\ttfamily\footnotesize\raggedright\noindent
forall(\allowbreak{}[vnv\_, vns\_, vnx\_], is-separating-for(\allowbreak{}vnv\_, vns\_, vnx\_) iff vnx\_ in vec(\allowbreak{}vnv\_) and forall(\allowbreak{}[vna\_ in vns\_], vna\_(\allowbreak{}vnx\_) = vzero(\allowbreak{}vnv\_) implies vna\_ = bop-zero(\allowbreak{}vnv\_, vnv\_)))
\end{quote}
mentioned by 10 result(s)

\smallskip

\noindent\textbf{\texttt{is-sesquilinear-form(v, sqf)}}\quad\textit{predicate}\index{is-sesquilinear-form(v, sqf)@\texttt{is-sesquilinear-form(v, sqf)}}\\
reads: \$2 is a sesquilinear form on \$1 \\
defined by \texttt{is-sesquilinear-form} \\
characterized by \texttt{is-sesquilinear-form}  \textit{(definitional)}:
\begin{quote}\ttfamily\footnotesize\raggedright\noindent
forall(\allowbreak{}[v, sqf], is-sesquilinear-form(\allowbreak{}v, sqf) iff sqf in fun(\allowbreak{}cartesian(\allowbreak{}vec(\allowbreak{}v), vec(\allowbreak{}v)), cc) and forall(\allowbreak{}[x\_ in vec(\allowbreak{}v), y\_ in vec(\allowbreak{}v), z\_ in vec(\allowbreak{}v)], sqf(\allowbreak{}(\allowbreak{}vadd(\allowbreak{}v))(\allowbreak{}x\_, y\_), z\_) = sqf(\allowbreak{}x\_, z\_) + sqf(\allowbreak{}y\_, z\_)) and forall(\allowbreak{}[a\_ in cc, x\_ in vec(\allowbreak{}v), y\_ in vec(\allowbreak{}v)], sqf(\allowbreak{}(\allowbreak{}cact(\allowbreak{}v))(\allowbreak{}a\_, x\_), y\_) = a\_ * sqf(\allowbreak{}x\_, y\_)) and forall(\allowbreak{}[x\_ in vec(\allowbreak{}v), y\_ in vec(\allowbreak{}v)], sqf(\allowbreak{}y\_, x\_) = conjugate(\allowbreak{}sqf(\allowbreak{}x\_, y\_))))
\end{quote}
mentioned by 20 result(s)

\smallskip

\noindent\textbf{\texttt{is-set(x)}}\quad\textit{predicate}\index{is-set(x)@\texttt{is-set(x)}}\\
reads: \$1 is a set \\
a predicate \\
characterized by \texttt{is-set}  \textit{(definitional)}:
\begin{quote}\ttfamily\footnotesize\raggedright\noindent
forall(\allowbreak{}[x], is-set(\allowbreak{}x) iff x in set)
\end{quote}
mentioned by 2 result(s)

\smallskip

\noindent\textbf{\texttt{is-setoid(s)}}\quad\textit{predicate}\index{is-setoid(s)@\texttt{is-setoid(s)}}\\
reads: a setoid \\
a predicate \\
characterized by \texttt{is-setoid}  \textit{(definitional)}:
\begin{quote}\ttfamily\footnotesize\raggedright\noindent
forall(\allowbreak{}[s], is-setoid(\allowbreak{}s) iff length(\allowbreak{}s) = 2 and pts(\allowbreak{}s) in set and rel(\allowbreak{}s) in power(\allowbreak{}cartesian(\allowbreak{}pts(\allowbreak{}s), pts(\allowbreak{}s))) and is-equivalence(\allowbreak{}rel(\allowbreak{}s), pts(\allowbreak{}s)) and s in set)
\end{quote}
also: \texttt{sip-setoid-is-setoid-csip-as-sip} \texttt{qmod-setoid-is-setoid-cna-as-module} \texttt{qmod-setoid-is-setoid-sip-as-module} \texttt{qmod-setoid-is-setoid-uba-as-module} \texttt{qmod-setoid-is-setoid-cnvs-as-module} \texttt{qmod-setoid-is-setoid-sip-as-module-csip-as-sip} \texttt{qmod-setoid-is-setoid-uba-as-module-uiba-as-uba} \texttt{is-setoid-is-set} ... (15 in all) \\
mentioned by 48 result(s)

\smallskip

\noindent\textbf{\texttt{is-sieve(s, sv)}}\quad\textit{predicate}\index{is-sieve(s, sv)@\texttt{is-sieve(s, sv)}}\\
reads: \$2 is a sieve of subsets of \$1 \\
defined by \texttt{is-sieve} \\
characterized by \texttt{is-sieve}  \textit{(definitional)}:
\begin{quote}\ttfamily\footnotesize\raggedright\noindent
forall(\allowbreak{}[s, sv], is-sieve(\allowbreak{}s, sv) iff sv in fun(\allowbreak{}nstar, power(\allowbreak{}pts(\allowbreak{}s))) and forall(\allowbreak{}[svz\_ in fun(\allowbreak{}nn, nn)], forall(\allowbreak{}[sve\_], pos-rr(\allowbreak{}sve\_) implies forsome(\allowbreak{}[svb\_ in nn], forall(\allowbreak{}[svn\_ in nn], svb\_ <= svn\_ implies diam(\allowbreak{}s, sv(\allowbreak{}restrict(\allowbreak{}svz\_, ord-segment(\allowbreak{}svn\_)))) <= sve\_))) and forall(\allowbreak{}[svn\_ in nn], sv(\allowbreak{}restrict(\allowbreak{}svz\_, ord-segment(\allowbreak{}succ(\allowbreak{}svn\_)))) subset sv(\allowbreak{}restrict(\allowbreak{}svz\_, ord-segment(\allowbreak{}svn\_)))) and forsome(\allowbreak{}[svp\_ in set], forall(\allowbreak{}[svn\_ in nn], svp\_ in sv(\allowbreak{}restrict(\allowbreak{}svz\_, ord-segment(\allowbreak{}svn\_)))))))
\end{quote}
also: \texttt{is-sieve-classes} \\
mentioned by 6 result(s)

\smallskip

\noindent\textbf{\texttt{is-sigma-algebra(omega, ca)}}\quad\textit{predicate}\index{is-sigma-algebra(omega, ca)@\texttt{is-sigma-algebra(omega, ca)}}\\
reads: \$2 is a sigma-algebra of subsets of \$1 \\
defined by \texttt{is-sigma-algebra} \\
characterized by \texttt{is-sigma-algebra}  \textit{(definitional)}:
\begin{quote}\ttfamily\footnotesize\raggedright\noindent
forall(\allowbreak{}[omega, ca], is-sigma-algebra(\allowbreak{}omega, ca) iff omega in set and ca subset power(\allowbreak{}omega) and omega in ca and forall(\allowbreak{}[a\_ in ca], complement-in(\allowbreak{}omega, a\_) in ca) and forall(\allowbreak{}[f\_ in fun(\allowbreak{}nn, ca)], big-union(\allowbreak{}n\_, nn, f\_(\allowbreak{}n\_)) in ca))
\end{quote}
also: \texttt{dynkin-pi-is-sigma-algebra} \texttt{product-sigma-is-sigma-algebra} \texttt{pullback-sigma-is-sigma-algebra} \texttt{sigma-generated-is-sigma-algebra} \\
mentioned by 119 result(s)

\smallskip

\noindent\textbf{\texttt{is-sigma-finite(omega, ca, mu)}}\quad\textit{predicate}\index{is-sigma-finite(omega, ca, mu)@\texttt{is-sigma-finite(omega, ca, mu)}}\\
reads: \$3 is sigma-finite on the sigma-algebra \$2 of subsets of \$1 \\
defined by \texttt{is-sigma-finite} \\
characterized by \texttt{is-sigma-finite}  \textit{(definitional)}:
\begin{quote}\ttfamily\footnotesize\raggedright\noindent
forall(\allowbreak{}[omega, ca, mu], is-sigma-finite(\allowbreak{}omega, ca, mu) iff forsome(\allowbreak{}[sff\_ in fun(\allowbreak{}nn, ca)], forall(\allowbreak{}[sfn\_ in nn], mu(\allowbreak{}sff\_(\allowbreak{}sfn\_)) in rr) and big-union(\allowbreak{}sfk\_, nn, sff\_(\allowbreak{}sfk\_)) = omega))
\end{quote}
mentioned by 52 result(s)

\smallskip

\noindent\textbf{\texttt{is-sigma-union-of(cb, es)}}\quad\textit{predicate}\index{is-sigma-union-of(cb, es)@\texttt{is-sigma-union-of(cb, es)}}\\
reads: \$2 is the union of a sequence of members of \$1 \\
defined by \texttt{is-sigma-union-of} \\
characterized by \texttt{is-sigma-union-of}  \textit{(definitional)}:
\begin{quote}\ttfamily\footnotesize\raggedright\noindent
forall(\allowbreak{}[cb, es], is-sigma-union-of(\allowbreak{}cb, es) iff forsome(\allowbreak{}[sug\_ in fun(\allowbreak{}nn, cb)], es = big-union(\allowbreak{}sun\_, nn, sug\_(\allowbreak{}sun\_))))
\end{quote}
mentioned by 14 result(s)

\smallskip

\noindent\textbf{\texttt{is-signed-measure(omega, ca, nu)}}\quad\textit{predicate}\index{is-signed-measure(omega, ca, nu)@\texttt{is-signed-measure(omega, ca, nu)}}\\
reads: \$3 is a countably additive signed measure on the sigma-algebra \$2 of subsets of \$1 \\
defined by \texttt{is-signed-measure} \\
characterized by \texttt{is-signed-measure}  \textit{(definitional)}:
\begin{quote}\ttfamily\footnotesize\raggedright\noindent
forall(\allowbreak{}[omega, ca, nu], is-signed-measure(\allowbreak{}omega, ca, nu) iff is-sigma-algebra(\allowbreak{}omega, ca) and nu in fun(\allowbreak{}ca, rr-star) and nu(\allowbreak{}empty-set) = 0 and not(\allowbreak{}forsome(\allowbreak{}[sma\_ in ca], nu(\allowbreak{}sma\_) = pos-inf) and forsome(\allowbreak{}[smb\_ in ca], nu(\allowbreak{}smb\_) = neg-inf)) and forall(\allowbreak{}[smf\_ in fun(\allowbreak{}nn, ca)], forall(\allowbreak{}[smi\_, smj\_], smi\_ in nn implies smj\_ in nn implies not(\allowbreak{}smi\_ = smj\_) implies intersection(\allowbreak{}smf\_(\allowbreak{}smi\_), smf\_(\allowbreak{}smj\_)) = empty-set) implies ext-series-converges-to(\allowbreak{}vnb-lambda(\allowbreak{}smn\_, nn, nu(\allowbreak{}smf\_(\allowbreak{}smn\_))), nu(\allowbreak{}big-union(\allowbreak{}smn\_, nn, smf\_(\allowbreak{}smn\_))))))
\end{quote}
also: \texttt{signed-neg-is-signed-measure} \\
mentioned by 35 result(s)

\smallskip

\noindent\textbf{\texttt{is-simple-fn(omega, ca, f)}}\quad\textit{predicate}\index{is-simple-fn(omega, ca, f)@\texttt{is-simple-fn(omega, ca, f)}}\\
reads: \$3 is a simple measurable function on \$1 with \$2 \\
defined by \texttt{is-simple-fn} \\
characterized by \texttt{is-simple-fn}  \textit{(definitional)}:
\begin{quote}\ttfamily\footnotesize\raggedright\noindent
forall(\allowbreak{}[omega, ca, f], is-simple-fn(\allowbreak{}omega, ca, f) iff is-measurable-fn(\allowbreak{}omega, ca, f) and card(\allowbreak{}image(\allowbreak{}f, omega)) in nn and forall(\allowbreak{}[x\_ in omega], f(\allowbreak{}x\_) in rr))
\end{quote}
mentioned by 36 result(s)

\smallskip

\noindent\textbf{\texttt{is-simple-msg(sbst\_, sbsk\_, sbsm\_)}}\quad\textit{predicate}\index{is-simple-msg(sbst_, sbsk_, sbsm_)@\texttt{is-simple-msg(sbst\_, sbsk\_, sbsm\_)}}\\
reads: \$3 is a simple message over the tokens \$1 and the keys \$2 \\
defined by \texttt{is-simple-msg} \\
characterized by \texttt{is-simple-msg}  \textit{(definitional)}:
\begin{quote}\ttfamily\footnotesize\raggedright\noindent
forall(\allowbreak{}[sbst\_, sbsk\_, sbsm\_], is-simple-msg(\allowbreak{}sbst\_, sbsk\_, sbsm\_) iff not(\allowbreak{}forsome(\allowbreak{}[sbsl\_ in tuples(\allowbreak{}msg-set(\allowbreak{}sbst\_, sbsk\_))], sbsm\_ = lst(\allowbreak{}sbsl\_))))
\end{quote}
mentioned by 6 result(s)

\smallskip

\noindent\textbf{\texttt{is-sk-space(s)}}\quad\textit{predicate}\index{is-sk-space(s)@\texttt{is-sk-space(s)}}\\
reads: a shared-key strand space \\
a predicate \\
characterized by \texttt{is-sk-space}  \textit{(definitional)}:
\begin{quote}\ttfamily\footnotesize\raggedright\noindent
forall(\allowbreak{}[s], is-sk-space(\allowbreak{}s) iff length(\allowbreak{}s) = 3 and is-infiltrated-space(\allowbreak{}sk-inf(\allowbreak{}s)) and sk-names(\allowbreak{}s) in set and ltkey(\allowbreak{}s) in fun(\allowbreak{}sk-names(\allowbreak{}s), keys(\allowbreak{}ispace(\allowbreak{}sk-inf(\allowbreak{}s)))) and forall(\allowbreak{}[orx\_ in sk-names(\allowbreak{}s)], orx\_ in tokens(\allowbreak{}ispace(\allowbreak{}sk-inf(\allowbreak{}s)))) and forall(\allowbreak{}[orx\_ in sk-names(\allowbreak{}s)], not(\allowbreak{}orx\_ in keys(\allowbreak{}ispace(\allowbreak{}sk-inf(\allowbreak{}s))))) and forall(\allowbreak{}[orx\_, ory\_], orx\_ in sk-names(\allowbreak{}s) implies ory\_ in sk-names(\allowbreak{}s) implies (\allowbreak{}ltkey(\allowbreak{}s))(\allowbreak{}orx\_) = (\allowbreak{}ltkey(\allowbreak{}s))(\allowbreak{}ory\_) implies orx\_ = ory\_) and forall(\allowbreak{}[orx\_ in sk-names(\allowbreak{}s)], (\allowbreak{}kinv(\allowbreak{}ispace(\allowbreak{}sk-inf(\allowbreak{}s))))(\allowbreak{}(\allowbreak{}ltkey(\allowbreak{}s))(\allowbreak{}orx\_)) = (\allowbreak{}ltkey(\allowbreak{}s))(\allowbreak{}orx\_)) and s in set)
\end{quote}
also: \texttt{cx-orp-is-sk-space} \texttt{is-sk-space-is-set} \\
mentioned by 19 result(s)

\smallskip

\noindent\textbf{\texttt{is-sn-open(m, ix, pf, u)}}\quad\textit{predicate}\index{is-sn-open(m, ix, pf, u)@\texttt{is-sn-open(m, ix, pf, u)}}\\
reads: \$4 is open in the topology defined by the seminorms \$3 \\
defined by \texttt{is-sn-open} \\
characterized by \texttt{is-sn-open}  \textit{(definitional)}:
\begin{quote}\ttfamily\footnotesize\raggedright\noindent
forall(\allowbreak{}[m, ix, pf, u], is-sn-open(\allowbreak{}m, ix, pf, u) iff u subset vec(\allowbreak{}m) and forall(\allowbreak{}[snx\_ in u], forsome(\allowbreak{}[snf\_ in set], card(\allowbreak{}snf\_) in nn and snf\_ subset ix and forsome(\allowbreak{}[sne\_], pos-rr(\allowbreak{}sne\_) and sn-ball(\allowbreak{}m, pf, snf\_, snx\_, sne\_) subset u))))
\end{quote}
mentioned by 9 result(s)

\smallskip

\noindent\textbf{\texttt{is-souslin(s, e)}}\quad\textit{predicate}\index{is-souslin(s, e)@\texttt{is-souslin(s, e)}}\\
reads: \$2 is a Souslin subset of \$1 \\
defined by \texttt{is-souslin} \\
characterized by \texttt{is-souslin}  \textit{(definitional)}:
\begin{quote}\ttfamily\footnotesize\raggedright\noindent
forall(\allowbreak{}[s, e], is-souslin(\allowbreak{}s, e) iff forsome(\allowbreak{}[ssv\_], is-sieve(\allowbreak{}s, ssv\_) and souslin-set(\allowbreak{}s, ssv\_) = e))
\end{quote}
also: \texttt{lusin-is-souslin} \\
mentioned by 31 result(s)

\smallskip

\noindent\textbf{\texttt{is-square-minus-id(v, jo)}}\quad\textit{predicate}\index{is-square-minus-id(v, jo)@\texttt{is-square-minus-id(v, jo)}}\\
reads: \$2 squares to minus the identity on \$1 \\
defined by \texttt{is-square-minus-id} \\
characterized by \texttt{is-square-minus-id}  \textit{(definitional)}:
\begin{quote}\ttfamily\footnotesize\raggedright\noindent
forall(\allowbreak{}[v, jo], is-square-minus-id(\allowbreak{}v, jo) iff forall(\allowbreak{}[x\_ in vec(\allowbreak{}v)], jo(\allowbreak{}jo(\allowbreak{}x\_)) = (\allowbreak{}vneg(\allowbreak{}v))(\allowbreak{}x\_)))
\end{quote}
mentioned by 15 result(s)

\smallskip

\noindent\textbf{\texttt{is-standard-borel(xs, ca)}}\quad\textit{predicate}\index{is-standard-borel(xs, ca)@\texttt{is-standard-borel(xs, ca)}}\\
reads: the measurable space (\$1, \$2) is a standard Borel space \\
defined by \texttt{is-standard-borel} \\
characterized by \texttt{is-standard-borel}  \textit{(definitional)}:
\begin{quote}\ttfamily\footnotesize\raggedright\noindent
forall(\allowbreak{}[xs, ca], is-standard-borel(\allowbreak{}xs, ca) iff forsome(\allowbreak{}[sbs\_], is-polish(\allowbreak{}sbs\_) and forsome(\allowbreak{}[sbf\_ in bijection(\allowbreak{}xs, pts(\allowbreak{}sbs\_))], is-measurable-map(\allowbreak{}xs, ca, pts(\allowbreak{}sbs\_), borel-sets(\allowbreak{}sbs\_), sbf\_) and forall(\allowbreak{}[sba\_ in ca], image(\allowbreak{}sbf\_, sba\_) in borel-sets(\allowbreak{}sbs\_)))))
\end{quote}
mentioned by 7 result(s)

\smallskip

\noindent\textbf{\texttt{is-star-hom(baa\_, bab\_, baf\_)}}\quad\textit{predicate}\index{is-star-hom(baa_, bab_, baf_)@\texttt{is-star-hom(baa\_, bab\_, baf\_)}}\\
reads: \$3 is a unital *-homomorphism from \$1 to \$2 \\
defined by \texttt{is-star-hom} \\
characterized by \texttt{is-star-hom}  \textit{(definitional)}:
\begin{quote}\ttfamily\footnotesize\raggedright\noindent
forall(\allowbreak{}[baa\_, bab\_, baf\_], is-star-hom(\allowbreak{}baa\_, bab\_, baf\_) iff is-alg-hom(\allowbreak{}baa\_, bab\_, baf\_) and forall(\allowbreak{}[bax\_ in vec(\allowbreak{}baa\_)], baf\_(\allowbreak{}(\allowbreak{}vstar(\allowbreak{}baa\_))(\allowbreak{}bax\_)) = (\allowbreak{}vstar(\allowbreak{}bab\_))(\allowbreak{}baf\_(\allowbreak{}bax\_))))
\end{quote}
mentioned by 9 result(s)

\smallskip

\noindent\textbf{\texttt{is-star-op-algebra(vnv\_, vns\_)}}\quad\textit{predicate}\index{is-star-op-algebra(vnv_, vns_)@\texttt{is-star-op-algebra(vnv\_, vns\_)}}\\
reads: \$2 is an involutive algebra of operators on \$1 containing the identity \\
defined by \texttt{is-star-op-algebra} \\
characterized by \texttt{is-star-op-algebra}  \textit{(definitional)}:
\begin{quote}\ttfamily\footnotesize\raggedright\noindent
forall(\allowbreak{}[vnv\_, vns\_], is-star-op-algebra(\allowbreak{}vnv\_, vns\_) iff vns\_ subset bounded-ops(\allowbreak{}vnv\_, vnv\_) and bop-id(\allowbreak{}vnv\_) in vns\_ and forall(\allowbreak{}[vna\_ in vns\_, vnb\_ in vns\_], bop-sum(\allowbreak{}vnv\_, vnv\_, vna\_, vnb\_) in vns\_) and forall(\allowbreak{}[vnc\_ in cc, vna\_ in vns\_], bop-scale(\allowbreak{}vnv\_, vnv\_, vnc\_, vna\_) in vns\_) and forall(\allowbreak{}[vna\_ in vns\_, vnb\_ in vns\_], bop-comp(\allowbreak{}vnv\_, vna\_, vnb\_) in vns\_) and forall(\allowbreak{}[vna\_ in vns\_], op-adjoint(\allowbreak{}vnv\_, vnv\_, vna\_) in vns\_))
\end{quote}
mentioned by 20 result(s)

\smallskip

\noindent\textbf{\texttt{is-star-rep(amv\_, amw\_, amf\_)}}\quad\textit{predicate}\index{is-star-rep(amv_, amw_, amf_)@\texttt{is-star-rep(amv\_, amw\_, amf\_)}}\\
reads: \$3 is a unital *-representation of the operators on \$1 on \$2 \\
defined by \texttt{is-star-rep} \\
characterized by \texttt{is-star-rep}  \textit{(definitional)}:
\begin{quote}\ttfamily\footnotesize\raggedright\noindent
forall(\allowbreak{}[amv\_, amw\_, amf\_], is-star-rep(\allowbreak{}amv\_, amw\_, amf\_) iff amf\_ in fun(\allowbreak{}bounded-ops(\allowbreak{}amv\_, amv\_), bounded-ops(\allowbreak{}amw\_, amw\_)) and forall(\allowbreak{}[ama\_ in bounded-ops(\allowbreak{}amv\_, amv\_), amb\_ in bounded-ops(\allowbreak{}amv\_, amv\_)], amf\_(\allowbreak{}bop-sum(\allowbreak{}amv\_, amv\_, ama\_, amb\_)) = bop-sum(\allowbreak{}amw\_, amw\_, amf\_(\allowbreak{}ama\_), amf\_(\allowbreak{}amb\_))) and forall(\allowbreak{}[amc\_ in cc, ama\_ in bounded-ops(\allowbreak{}amv\_, amv\_)], amf\_(\allowbreak{}bop-scale(\allowbreak{}amv\_, amv\_, amc\_, ama\_)) = bop-scale(\allowbreak{}amw\_, amw\_, amc\_, amf\_(\allowbreak{}ama\_))) and forall(\allowbreak{}[ama\_ in bounded-ops(\allowbreak{}amv\_, amv\_), amb\_ in bounded-ops(\allowbreak{}amv\_, amv\_)], amf\_(\allowbreak{}bop-comp(\allowbreak{}amv\_, ama\_, amb\_)) = bop-comp(\allowbreak{}amw\_, amf\_(\allowbreak{}ama\_), amf\_(\allowbreak{}amb\_))) and amf\_(\allowbreak{}bop-id(\allowbreak{}amv\_)) = bop-id(\allowbreak{}amw\_) and forall(\allowbreak{}[ama\_ in bounded-ops(\allowbreak{}amv\_, amv\_)], amf\_(\allowbreak{}op-adjoint(\allowbreak{}amv\_, amv\_, ama\_)) = op-adjoint(\allowbreak{}amw\_, amw\_, amf\_(\allowbreak{}ama\_))))
\end{quote}
mentioned by 7 result(s)

\smallskip

\noindent\textbf{\texttt{is-step-fn(f, a, b)}}\quad\textit{predicate}\index{is-step-fn(f, a, b)@\texttt{is-step-fn(f, a, b)}}\\
reads: \$1 is a step function on [\$2, \$3] \\
defined by \texttt{is-step-fn} \\
characterized by \texttt{is-step-fn}  \textit{(definitional)}:
\begin{quote}\ttfamily\footnotesize\raggedright\noindent
forall(\allowbreak{}[f, a, b], is-step-fn(\allowbreak{}f, a, b) iff f in fun(\allowbreak{}rr, rr) and a in rr and b in rr and a < b and forsome(\allowbreak{}[n\_, p\_], is-partition(\allowbreak{}p\_, n\_, a, b) and forall(\allowbreak{}[i\_], i\_ in interval(\allowbreak{}0, n\_) and i\_ < n\_ implies forsome(\allowbreak{}[c\_ in rr], forall(\allowbreak{}[t\_], t\_ in rr and p\_(\allowbreak{}i\_) < t\_ and t\_ < p\_(\allowbreak{}succ(\allowbreak{}i\_)) implies f(\allowbreak{}t\_) = c\_)))))
\end{quote}
mentioned by 10 result(s)

\smallskip

\noindent\textbf{\texttt{is-step-map(omega, ca, s, g)}}\quad\textit{predicate}\index{is-step-map(omega, ca, s, g)@\texttt{is-step-map(omega, ca, s, g)}}\\
reads: \$4 is a step map on \$1 with \$2 into the metric space \$3 (measurable, with countably many values) \\
defined by \texttt{is-step-map} \\
characterized by \texttt{is-step-map}  \textit{(definitional)}:
\begin{quote}\ttfamily\footnotesize\raggedright\noindent
forall(\allowbreak{}[omega, ca, s, g], is-step-map(\allowbreak{}omega, ca, s, g) iff is-measurable-map(\allowbreak{}omega, ca, pts(\allowbreak{}s), borel-sets(\allowbreak{}s), g) and forsome(\allowbreak{}[stv\_ in fun(\allowbreak{}nn, pts(\allowbreak{}s))], forall(\allowbreak{}[stx\_ in omega], forsome(\allowbreak{}[stk\_ in nn], g(\allowbreak{}stx\_) = stv\_(\allowbreak{}stk\_)))))
\end{quote}
mentioned by 4 result(s)

\smallskip

\noindent\textbf{\texttt{is-strand-space(s)}}\quad\textit{predicate}\index{is-strand-space(s)@\texttt{is-strand-space(s)}}\\
reads: a strand space \\
a predicate \\
characterized by \texttt{is-strand-space}  \textit{(definitional)}:
\begin{quote}\ttfamily\footnotesize\raggedright\noindent
forall(\allowbreak{}[s], is-strand-space(\allowbreak{}s) iff length(\allowbreak{}s) = 5 and tokens(\allowbreak{}s) in set and ids(\allowbreak{}s) in set and keys(\allowbreak{}s) in set and strand-map(\allowbreak{}s) in fun(\allowbreak{}ids(\allowbreak{}s), strands(\allowbreak{}tokens(\allowbreak{}s), keys(\allowbreak{}s))) and kinv(\allowbreak{}s) in fun(\allowbreak{}keys(\allowbreak{}s), keys(\allowbreak{}s)) and is-token-set(\allowbreak{}tokens(\allowbreak{}s)) and forall(\allowbreak{}[ssx\_ in keys(\allowbreak{}s)], ssx\_ in tokens(\allowbreak{}s)) and forall(\allowbreak{}[ssx\_ in keys(\allowbreak{}s)], (\allowbreak{}kinv(\allowbreak{}s))(\allowbreak{}(\allowbreak{}kinv(\allowbreak{}s))(\allowbreak{}ssx\_)) = ssx\_) and s in set)
\end{quote}
also: \texttt{cx-sp-is-strand-space} \texttt{is-strand-space-is-set} \\
mentioned by 42 result(s)

\smallskip

\noindent\textbf{\texttt{is-strictly-below(porel, x, y)}}\quad\textit{predicate}\index{is-strictly-below(porel, x, y)@\texttt{is-strictly-below(porel, x, y)}}\\
reads: \$2 is strictly below \$3 under \$1 \\
defined by \texttt{is-strictly-below} \\
characterized by \texttt{is-strictly-below}  \textit{(definitional)}:
\begin{quote}\ttfamily\footnotesize\raggedright\noindent
forall(\allowbreak{}[porel, x, y], is-strictly-below(\allowbreak{}porel, x, y) iff [x, y] in porel and not(\allowbreak{}[y, x] in porel))
\end{quote}
mentioned by 11 result(s)

\smallskip

\noindent\textbf{\texttt{is-strictly-increasing-on(f, s)}}\quad\textit{predicate}\index{is-strictly-increasing-on(f, s)@\texttt{is-strictly-increasing-on(f, s)}}\\
reads: \$1 is strictly increasing on \$2 \\
defined by \texttt{is-strictly-increasing-on} \\
characterized by \texttt{is-strictly-increasing-on}  \textit{(definitional)}:
\begin{quote}\ttfamily\footnotesize\raggedright\noindent
forall(\allowbreak{}[f, s], is-strictly-increasing-on(\allowbreak{}f, s) iff f in fun(\allowbreak{}rr, rr) and s subset rr and forall(\allowbreak{}[u in s, v in s], u < v implies f(\allowbreak{}u) < f(\allowbreak{}v)))
\end{quote}
mentioned by 6 result(s)

\smallskip

\noindent\textbf{\texttt{is-subfield-of-cc(kk)}}\quad\textit{predicate}\index{is-subfield-of-cc(kk)@\texttt{is-subfield-of-cc(kk)}}\\
reads: \$1 is a subfield of the complex numbers \\
defined by \texttt{is-subfield-of-cc} \\
characterized by \texttt{is-subfield-of-cc}  \textit{(definitional)}:
\begin{quote}\ttfamily\footnotesize\raggedright\noindent
forall(\allowbreak{}[kk], is-subfield-of-cc(\allowbreak{}kk) iff is-subring-of-cc(\allowbreak{}kk) and forall(\allowbreak{}[ansf\_ in kk], not(\allowbreak{}ansf\_ = 0) implies recip(\allowbreak{}ansf\_) in kk))
\end{quote}
also: \texttt{qq-is-subfield-of-cc} \\
mentioned by 12 result(s)

\smallskip

\noindent\textbf{\texttt{is-subgroup(sgg\_, sgh\_)}}\quad\textit{predicate}\index{is-subgroup(sgg_, sgh_)@\texttt{is-subgroup(sgg\_, sgh\_)}}\\
reads: \$2 is a subgroup of \$1 \\
defined by \texttt{is-subgroup} \\
characterized by \texttt{is-subgroup}  \textit{(definitional)}:
\begin{quote}\ttfamily\footnotesize\raggedright\noindent
forall(\allowbreak{}[sgg\_, sgh\_], is-subgroup(\allowbreak{}sgg\_, sgh\_) iff sgh\_ subset carr(\allowbreak{}sgg\_) and iden(\allowbreak{}sgg\_) in sgh\_ and forall(\allowbreak{}[sga\_ in sgh\_, sgb\_ in sgh\_], (\allowbreak{}opr(\allowbreak{}sgg\_))(\allowbreak{}sga\_, sgb\_) in sgh\_) and forall(\allowbreak{}[sga\_ in sgh\_], (\allowbreak{}inv(\allowbreak{}sgg\_))(\allowbreak{}sga\_) in sgh\_))
\end{quote}
also: \texttt{generated-subgroup-is-subgroup} \\
mentioned by 75 result(s)

\smallskip

\noindent\textbf{\texttt{is-sublinear(m, p)}}\quad\textit{predicate}\index{is-sublinear(m, p)@\texttt{is-sublinear(m, p)}}\\
reads: \$2 is a sublinear functional on \$1 \\
defined by \texttt{is-sublinear} \\
characterized by \texttt{is-sublinear}  \textit{(definitional)}:
\begin{quote}\ttfamily\footnotesize\raggedright\noindent
forall(\allowbreak{}[m, p], is-sublinear(\allowbreak{}m, p) iff p in fun(\allowbreak{}vec(\allowbreak{}m), rr) and forall(\allowbreak{}[hsx\_ in vec(\allowbreak{}m), hsy\_ in vec(\allowbreak{}m)], p(\allowbreak{}(\allowbreak{}vadd(\allowbreak{}m))(\allowbreak{}hsx\_, hsy\_)) <= p(\allowbreak{}hsx\_) + p(\allowbreak{}hsy\_)) and forall(\allowbreak{}[hst\_ in rr], 0 <= hst\_ implies forall(\allowbreak{}[hsx\_ in vec(\allowbreak{}m)], p(\allowbreak{}(\allowbreak{}act(\allowbreak{}m))(\allowbreak{}hst\_, hsx\_)) = hst\_ * p(\allowbreak{}hsx\_))))
\end{quote}
mentioned by 127 result(s)

\smallskip

\noindent\textbf{\texttt{is-submodule(m, s)}}\quad\textit{predicate}\index{is-submodule(m, s)@\texttt{is-submodule(m, s)}}\\
reads: \$2 is a submodule of \$1 \\
defined by \texttt{is-submodule} \\
characterized by \texttt{is-submodule}  \textit{(definitional)}:
\begin{quote}\ttfamily\footnotesize\raggedright\noindent
forall(\allowbreak{}[m, s], is-submodule(\allowbreak{}m, s) iff s subset vec(\allowbreak{}m) and vzero(\allowbreak{}m) in s and forall(\allowbreak{}[x\_ in s, y\_ in s], (\allowbreak{}vadd(\allowbreak{}m))(\allowbreak{}x\_, y\_) in s) and forall(\allowbreak{}[x\_ in s], (\allowbreak{}vneg(\allowbreak{}m))(\allowbreak{}x\_) in s) and forall(\allowbreak{}[r\_ in carr(\allowbreak{}scal(\allowbreak{}m)), x\_ in s], (\allowbreak{}act(\allowbreak{}m))(\allowbreak{}r\_, x\_) in s))
\end{quote}
also: \texttt{span-is-submodule-cna-as-module} \texttt{span-is-submodule-sip-as-module} \texttt{span-is-submodule-uba-as-module} \texttt{span-is-submodule-cnvs-as-module} \texttt{whole-module-is-submodule-cna-as-module} \texttt{whole-module-is-submodule-sip-as-module} \texttt{whole-module-is-submodule-uba-as-module} \texttt{whole-module-is-submodule-cnvs-as-module} ... (19 in all) \\
mentioned by 1203 result(s)

\smallskip

\noindent\textbf{\texttt{is-subnet(d, x, e, y)}}\quad\textit{predicate}\index{is-subnet(d, x, e, y)@\texttt{is-subnet(d, x, e, y)}}\\
reads: the net \$4 on \$3 is a subnet of the net \$2 on \$1 \\
defined by \texttt{is-subnet} \\
characterized by \texttt{is-subnet}  \textit{(definitional)}:
\begin{quote}\ttfamily\footnotesize\raggedright\noindent
forall(\allowbreak{}[d, x, e, y], is-subnet(\allowbreak{}d, x, e, y) iff is-directed-set(\allowbreak{}d) and is-directed-set(\allowbreak{}e) and forsome(\allowbreak{}[snh\_], is-cofinal-monotone(\allowbreak{}e, d, snh\_) and forall(\allowbreak{}[snj\_ in pts(\allowbreak{}e)], y(\allowbreak{}snj\_) = x(\allowbreak{}snh\_(\allowbreak{}snj\_)))))
\end{quote}
mentioned by 4 result(s)

\smallskip

\noindent\textbf{\texttt{is-subring-of-cc(aa)}}\quad\textit{predicate}\index{is-subring-of-cc(aa)@\texttt{is-subring-of-cc(aa)}}\\
reads: \$1 is a subring of the complex numbers \\
defined by \texttt{is-subring-of-cc} \\
characterized by \texttt{is-subring-of-cc}  \textit{(definitional)}:
\begin{quote}\ttfamily\footnotesize\raggedright\noindent
forall(\allowbreak{}[aa], is-subring-of-cc(\allowbreak{}aa) iff is-cc-subring(\allowbreak{}aa) and forall(\allowbreak{}[ansn\_ in aa], -ansn\_ in aa))
\end{quote}
also: \texttt{cc-is-subring-of-cc} \texttt{qq-is-subring-of-cc} \texttt{zz-is-subring-of-cc} \\
mentioned by 36 result(s)

\smallskip

\noindent\textbf{\texttt{is-subsequence(s, y, f)}}\quad\textit{predicate}\index{is-subsequence(s, y, f)@\texttt{is-subsequence(s, y, f)}}\\
reads: \$2 is a subsequence of \$3 in \$1 \\
defined by \texttt{is-subsequence} \\
characterized by \texttt{is-subsequence}  \textit{(definitional)}:
\begin{quote}\ttfamily\footnotesize\raggedright\noindent
forall(\allowbreak{}[s, y, f], is-subsequence(\allowbreak{}s, y, f) iff f in fun(\allowbreak{}nn, pts(\allowbreak{}s)) and forsome(\allowbreak{}[phi], strictly-mono-nn(\allowbreak{}phi) and y = subseq(\allowbreak{}f, phi)))
\end{quote}
also: \texttt{is-subsequence-iff-typed-subsqn} \\
mentioned by 4 result(s)

\smallskip

\noindent\textbf{\texttt{is-subspace(m, s)}}\quad\textit{predicate}\index{is-subspace(m, s)@\texttt{is-subspace(m, s)}}\\
reads: \$2 is a subspace of \$1 \\
defined by \texttt{is-subspace} \\
characterized by \texttt{is-subspace}  \textit{(definitional)}:
\begin{quote}\ttfamily\footnotesize\raggedright\noindent
forall(\allowbreak{}[m, s], is-subspace(\allowbreak{}m, s) iff is-submodule(\allowbreak{}m, s))
\end{quote}
mentioned by 4 result(s)

\smallskip

\noindent\textbf{\texttt{is-subterm(smt\_, smk\_, sma\_, smb\_)}}\quad\textit{predicate}\index{is-subterm(smt_, smk_, sma_, smb_)@\texttt{is-subterm(smt\_, smk\_, sma\_, smb\_)}}\\
reads: \$3 is a subterm of \$4 over the tokens \$1 and the keys \$2 \\
defined by \texttt{is-subterm} \\
characterized by \texttt{is-subterm}  \textit{(definitional)}:
\begin{quote}\ttfamily\footnotesize\raggedright\noindent
forall(\allowbreak{}[smt\_, smk\_, sma\_, smb\_], is-subterm(\allowbreak{}smt\_, smk\_, sma\_, smb\_) iff [sma\_, smb\_] in subterm-rel(\allowbreak{}smt\_, smk\_))
\end{quote}
mentioned by 76 result(s)

\smallskip

\noindent\textbf{\texttt{is-subterm-closed(smt\_, smk\_, smr\_)}}\quad\textit{predicate}\index{is-subterm-closed(smt_, smk_, smr_)@\texttt{is-subterm-closed(smt\_, smk\_, smr\_)}}\\
reads: \$3 is closed under the subterm clauses over the tokens \$1 and the keys \$2 \\
defined by \texttt{is-subterm-closed} \\
characterized by \texttt{is-subterm-closed}  \textit{(definitional)}:
\begin{quote}\ttfamily\footnotesize\raggedright\noindent
forall(\allowbreak{}[smt\_, smk\_, smr\_], is-subterm-closed(\allowbreak{}smt\_, smk\_, smr\_) iff forall(\allowbreak{}[sma\_ in msg-set(\allowbreak{}smt\_, smk\_)], [sma\_, sma\_] in smr\_) and forall(\allowbreak{}[sma\_, sml\_, smi\_], sml\_ in tuples(\allowbreak{}msg-set(\allowbreak{}smt\_, smk\_)) implies smi\_ in nn implies 1 <= smi\_ implies smi\_ <= length(\allowbreak{}sml\_) implies [sma\_, nth(\allowbreak{}smi\_, sml\_)] in smr\_ implies [sma\_, lst(\allowbreak{}sml\_)] in smr\_) and forall(\allowbreak{}[sma\_, smg\_, smy\_], smg\_ in msg-set(\allowbreak{}smt\_, smk\_) implies smy\_ in smk\_ implies [sma\_, smg\_] in smr\_ implies [sma\_, encr(\allowbreak{}smg\_, smy\_)] in smr\_))
\end{quote}
mentioned by 5 result(s)

\smallskip

\noindent\textbf{\texttt{is-summable(grp, f)}}\quad\textit{predicate}\index{is-summable(grp, f)@\texttt{is-summable(grp, f)}}\\
reads: \$2 is summable in \$1 \\
defined by \texttt{is-summable} \\
characterized by \texttt{is-summable}  \textit{(definitional)}:
\begin{quote}\ttfamily\footnotesize\raggedright\noindent
forall(\allowbreak{}[grp, f], is-summable(\allowbreak{}grp, f) iff forsome(\allowbreak{}[r in carr(\allowbreak{}grp)], sums-to(\allowbreak{}grp, f, r)))
\end{quote}
mentioned by 6 result(s)

\smallskip

\noindent\textbf{\texttt{is-symmetric-graph(ubv\_, ubg\_)}}\quad\textit{predicate}\index{is-symmetric-graph(ubv_, ubg_)@\texttt{is-symmetric-graph(ubv\_, ubg\_)}}\\
reads: \$2 is a symmetric graph \\
defined by \texttt{is-symmetric-graph} \\
characterized by \texttt{is-symmetric-graph}  \textit{(definitional)}:
\begin{quote}\ttfamily\footnotesize\raggedright\noindent
forall(\allowbreak{}[ubv\_, ubg\_], is-symmetric-graph(\allowbreak{}ubv\_, ubg\_) iff ubg\_ subset graph-adjoint(\allowbreak{}ubv\_, ubg\_))
\end{quote}
mentioned by 6 result(s)

\smallskip

\noindent\textbf{\texttt{is-symmetric-subset(m, cs)}}\quad\textit{predicate}\index{is-symmetric-subset(m, cs)@\texttt{is-symmetric-subset(m, cs)}}\\
reads: \$2 is a symmetric subset of \$1 \\
defined by \texttt{is-symmetric-subset} \\
characterized by \texttt{is-symmetric-subset}  \textit{(definitional)}:
\begin{quote}\ttfamily\footnotesize\raggedright\noindent
forall(\allowbreak{}[m, cs], is-symmetric-subset(\allowbreak{}m, cs) iff cs subset vec(\allowbreak{}m) and forall(\allowbreak{}[syx\_ in cs], (\allowbreak{}act(\allowbreak{}m))(\allowbreak{}-1, syx\_) in cs))
\end{quote}
mentioned by 10 result(s)

\smallskip

\noindent\textbf{\texttt{is-taylor-family(d, n, a, b)}}\quad\textit{predicate}\index{is-taylor-family(d, n, a, b)@\texttt{is-taylor-family(d, n, a, b)}}\\
reads: \$1 is a Taylor family of order \$2 on \$3 to \$4 \\
defined by \texttt{is-taylor-family} \\
characterized by \texttt{is-taylor-family}  \textit{(definitional)}:
\begin{quote}\ttfamily\footnotesize\raggedright\noindent
forall(\allowbreak{}[d, n, a, b], is-taylor-family(\allowbreak{}d, n, a, b) iff n in nn and 1 <= n and a in rr and b in rr and a < b and d(\allowbreak{}n) in fun(\allowbreak{}ooint(\allowbreak{}a, b), rr) and forall(\allowbreak{}[tk\_ in nn], tk\_ < n implies d(\allowbreak{}tk\_) in fun(\allowbreak{}ccint(\allowbreak{}a, b), rr) and is-continuous-on(\allowbreak{}d(\allowbreak{}tk\_), ccint(\allowbreak{}a, b)) and forall(\allowbreak{}[tx\_ in ooint(\allowbreak{}a, b)], has-deriv-at(\allowbreak{}d(\allowbreak{}tk\_), tx\_, (\allowbreak{}d(\allowbreak{}succ(\allowbreak{}tk\_)))(\allowbreak{}tx\_)))))
\end{quote}
mentioned by 23 result(s)

\smallskip

\noindent\textbf{\texttt{is-tee-trace(sasp\_, sas\_)}}\quad\textit{predicate}\index{is-tee-trace(sasp_, sas_)@\texttt{is-tee-trace(sasp\_, sas\_)}}\\
reads: \$2 is a tee trace in the strand space \$1 \\
defined by \texttt{is-tee-trace} \\
characterized by \texttt{is-tee-trace}  \textit{(definitional)}:
\begin{quote}\ttfamily\footnotesize\raggedright\noindent
forall(\allowbreak{}[sasp\_, sas\_], is-tee-trace(\allowbreak{}sasp\_, sas\_) iff forsome(\allowbreak{}[sag\_ in msg-set(\allowbreak{}tokens(\allowbreak{}sasp\_), keys(\allowbreak{}sasp\_))], sas\_ = [[received, sag\_], [sent, sag\_], [sent, sag\_]]))
\end{quote}
mentioned by 5 result(s)

\smallskip

\noindent\textbf{\texttt{is-text-trace(sasp\_, sap\_, sas\_)}}\quad\textit{predicate}\index{is-text-trace(sasp_, sap_, sas_)@\texttt{is-text-trace(sasp\_, sap\_, sas\_)}}\\
reads: \$3 is a text trace of the adversary knowing \$2 in the strand space \$1 \\
defined by \texttt{is-text-trace} \\
characterized by \texttt{is-text-trace}  \textit{(definitional)}:
\begin{quote}\ttfamily\footnotesize\raggedright\noindent
forall(\allowbreak{}[sasp\_, sap\_, sas\_], is-text-trace(\allowbreak{}sasp\_, sap\_, sas\_) iff forsome(\allowbreak{}[saa\_ in sap\_], sas\_ = [[sent, saa\_]]))
\end{quote}
mentioned by 9 result(s)

\smallskip

\noindent\textbf{\texttt{is-token-set(smt\_)}}\quad\textit{predicate}\index{is-token-set(smt_)@\texttt{is-token-set(smt\_)}}\\
reads: \$1 is a set of tokens, none of which is a list \\
defined by \texttt{is-token-set} \\
characterized by \texttt{is-token-set}  \textit{(definitional)}:
\begin{quote}\ttfamily\footnotesize\raggedright\noindent
forall(\allowbreak{}[smt\_], is-token-set(\allowbreak{}smt\_) iff smt\_ in set and forall(\allowbreak{}[smx\_ in smt\_], not(\allowbreak{}smx\_ in tuples(\allowbreak{}set))))
\end{quote}
also: \texttt{nn-is-token-set} \\
mentioned by 34 result(s)

\smallskip

\noindent\textbf{\texttt{is-top-closed(t, a)}}\quad\textit{predicate}\index{is-top-closed(t, a)@\texttt{is-top-closed(t, a)}}\\
reads: \$2 is closed in \$1 \\
defined by \texttt{is-top-closed} \\
characterized by \texttt{is-top-closed}  \textit{(definitional)}:
\begin{quote}\ttfamily\footnotesize\raggedright\noindent
forall(\allowbreak{}[t, a], is-top-closed(\allowbreak{}t, a) iff a subset pts(\allowbreak{}t) and complement-in(\allowbreak{}pts(\allowbreak{}t), a) in opens(\allowbreak{}t))
\end{quote}
also: \texttt{metric-closed-is-top-closed} \\
mentioned by 47 result(s)

\smallskip

\noindent\textbf{\texttt{is-top-family(ix, tf)}}\quad\textit{predicate}\index{is-top-family(ix, tf)@\texttt{is-top-family(ix, tf)}}\\
reads: \$2 is a family of topological spaces indexed by \$1 \\
defined by \texttt{is-top-family} \\
characterized by \texttt{is-top-family}  \textit{(definitional)}:
\begin{quote}\ttfamily\footnotesize\raggedright\noindent
forall(\allowbreak{}[ix, tf], is-top-family(\allowbreak{}ix, tf) iff ix in set and forall(\allowbreak{}[tpi\_ in ix], is-top-space(\allowbreak{}tf(\allowbreak{}tpi\_))))
\end{quote}
mentioned by 26 result(s)

\smallskip

\noindent\textbf{\texttt{is-top-group(s)}}\quad\textit{predicate}\index{is-top-group(s)@\texttt{is-top-group(s)}}\\
reads: a topological group \\
a predicate \\
characterized by \texttt{is-top-group}  \textit{(definitional)}:
\begin{quote}\ttfamily\footnotesize\raggedright\noindent
forall(\allowbreak{}[s], is-top-group(\allowbreak{}s) iff length(\allowbreak{}s) = 5 and carr(\allowbreak{}s) in set and opr(\allowbreak{}s) in fun(\allowbreak{}cartesian(\allowbreak{}carr(\allowbreak{}s), carr(\allowbreak{}s)), carr(\allowbreak{}s)) and iden(\allowbreak{}s) in carr(\allowbreak{}s) and inv(\allowbreak{}s) in fun(\allowbreak{}carr(\allowbreak{}s), carr(\allowbreak{}s)) and topology(\allowbreak{}s) in power(\allowbreak{}power(\allowbreak{}carr(\allowbreak{}s))) and is-associative(\allowbreak{}opr(\allowbreak{}s), carr(\allowbreak{}s)) and is-identity(\allowbreak{}opr(\allowbreak{}s), iden(\allowbreak{}s), carr(\allowbreak{}s)) and has-inverses(\allowbreak{}opr(\allowbreak{}s), iden(\allowbreak{}s), inv(\allowbreak{}s), carr(\allowbreak{}s)) and empty-set in topology(\allowbreak{}s) and carr(\allowbreak{}s) in topology(\allowbreak{}s) and forall(\allowbreak{}[u in topology(\allowbreak{}s), v in topology(\allowbreak{}s)], intersection(\allowbreak{}u, v) in topology(\allowbreak{}s)) and forall(\allowbreak{}[fam in set], fam subset topology(\allowbreak{}s) implies big-union(\allowbreak{}u, fam, u) in topology(\allowbreak{}s)) and forall(\allowbreak{}[tga\_ in carr(\allowbreak{}s), tgb\_ in carr(\allowbreak{}s), tgw\_ in topology(\allowbreak{}s)], (\allowbreak{}opr(\allowbreak{}s))(\allowbreak{}tga\_, (\allowbreak{}inv(\allowbreak{}s))(\allowbreak{}tgb\_)) in tgw\_ implies forsome(\allowbreak{}[tgu\_ in topology(\allowbreak{}s)], tga\_ in tgu\_ and forsome(\allowbreak{}[tgv\_ in topology(\allowbreak{}s)], tgb\_ in tgv\_ and forall(\allowbreak{}[tgp\_ in tgu\_, tgq\_ in tgv\_], (\allowbreak{}opr(\allowbreak{}s))(\allowbreak{}tgp\_, (\allowbreak{}inv(\allowbreak{}s))(\allowbreak{}tgq\_)) in tgw\_)))) and s in set)
\end{quote}
also: \texttt{is-top-group-is-set} \texttt{metric-group-is-top-group} \texttt{isometry-group-is-top-group} \texttt{make-top-group-is-top-group} \\
mentioned by 40 result(s)

\smallskip

\noindent\textbf{\texttt{is-top-space(s)}}\quad\textit{predicate}\index{is-top-space(s)@\texttt{is-top-space(s)}}\\
reads: \$1 is a topological space \\
a predicate \\
characterized by \texttt{is-top-space}  \textit{(definitional)}:
\begin{quote}\ttfamily\footnotesize\raggedright\noindent
forall(\allowbreak{}[s], is-top-space(\allowbreak{}s) iff length(\allowbreak{}s) = 2 and pts(\allowbreak{}s) in set and opens(\allowbreak{}s) in power(\allowbreak{}power(\allowbreak{}pts(\allowbreak{}s))) and empty-set in opens(\allowbreak{}s) and pts(\allowbreak{}s) in opens(\allowbreak{}s) and forall(\allowbreak{}[u in opens(\allowbreak{}s), v in opens(\allowbreak{}s)], intersection(\allowbreak{}u, v) in opens(\allowbreak{}s)) and forall(\allowbreak{}[fam in set], fam subset opens(\allowbreak{}s) implies big-union(\allowbreak{}u, fam, u) in opens(\allowbreak{}s)) and s in set)
\end{quote}
also: \texttt{sntop-is-top-space} \texttt{is-top-space-is-set} \texttt{prod-top-is-top-space} \texttt{tprod-top-is-top-space} \texttt{metric-top-is-top-space} \texttt{subspace-top-is-top-space} \texttt{rb3-gauge-top-is-top-space} \texttt{discrete-space-is-top-space} ... (9 in all) \\
mentioned by 189 result(s)

\smallskip

\noindent\textbf{\texttt{is-totally-bounded-in(tbm\_, tba\_)}}\quad\textit{predicate}\index{is-totally-bounded-in(tbm_, tba_)@\texttt{is-totally-bounded-in(tbm\_, tba\_)}}\\
reads: \$2 is totally bounded in \$1 \\
defined by \texttt{is-totally-bounded-in} \\
characterized by \texttt{is-totally-bounded-in}  \textit{(definitional)}:
\begin{quote}\ttfamily\footnotesize\raggedright\noindent
forall(\allowbreak{}[tbm\_, tba\_], is-totally-bounded-in(\allowbreak{}tbm\_, tba\_) iff forall(\allowbreak{}[tbr\_], pos-rr(\allowbreak{}tbr\_) implies forsome(\allowbreak{}[tbf\_ in set], card(\allowbreak{}tbf\_) in nn and tbf\_ subset pts(\allowbreak{}tbm\_) and forall(\allowbreak{}[tbp\_ in tba\_], forsome(\allowbreak{}[tbc\_ in tbf\_], (\allowbreak{}dist(\allowbreak{}tbm\_))(\allowbreak{}tbc\_, tbp\_) < tbr\_)))))
\end{quote}
also: \texttt{is-totally-bounded-in-classes} \\
mentioned by 14 result(s)

\smallskip

\noindent\textbf{\texttt{is-tprod-open(ix, tf, u)}}\quad\textit{predicate}\index{is-tprod-open(ix, tf, u)@\texttt{is-tprod-open(ix, tf, u)}}\\
reads: \$3 is open in the product of the family \$2 \\
defined by \texttt{is-tprod-open} \\
characterized by \texttt{is-tprod-open}  \textit{(definitional)}:
\begin{quote}\ttfamily\footnotesize\raggedright\noindent
forall(\allowbreak{}[ix, tf, u], is-tprod-open(\allowbreak{}ix, tf, u) iff u subset tprod-pts(\allowbreak{}ix, tf) and forall(\allowbreak{}[tpw\_ in u], forsome(\allowbreak{}[tpf\_ in set], tpf\_ subset tprod-constr(\allowbreak{}ix, tf) and card(\allowbreak{}tpf\_) in nn and tpw\_ in tprod-box(\allowbreak{}ix, tf, tpf\_) and tprod-box(\allowbreak{}ix, tf, tpf\_) subset u)))
\end{quote}
also: \texttt{is-tprod-open-classes} \\
mentioned by 10 result(s)

\smallskip

\noindent\textbf{\texttt{is-trace-class(kqv\_, kqt\_)}}\quad\textit{predicate}\index{is-trace-class(kqv_, kqt_)@\texttt{is-trace-class(kqv\_, kqt\_)}}\\
reads: \$2 is a trace-class operator on \$1 \\
defined by \texttt{is-trace-class} \\
characterized by \texttt{is-trace-class}  \textit{(definitional)}:
\begin{quote}\ttfamily\footnotesize\raggedright\noindent
forall(\allowbreak{}[kqv\_, kqt\_], is-trace-class(\allowbreak{}kqv\_, kqt\_) iff kqt\_ in bounded-ops(\allowbreak{}kqv\_, kqv\_) and forsome(\allowbreak{}[kqa\_], is-hilbert-schmidt(\allowbreak{}kqv\_, kqv\_, kqa\_) and forsome(\allowbreak{}[kqb\_], is-hilbert-schmidt(\allowbreak{}kqv\_, kqv\_, kqb\_) and kqt\_ = bop-comp(\allowbreak{}kqv\_, kqa\_, kqb\_))))
\end{quote}
mentioned by 25 result(s)

\smallskip

\noindent\textbf{\texttt{is-tvs(fld, m, t)}}\quad\textit{predicate}\index{is-tvs(fld, m, t)@\texttt{is-tvs(fld, m, t)}}\\
reads: \$3 is a topological vector space structure on \$2 over \$1 \\
defined by \texttt{is-tvs} \\
characterized by \texttt{is-tvs}  \textit{(definitional)}:
\begin{quote}\ttfamily\footnotesize\raggedright\noindent
forall(\allowbreak{}[fld, m, t], is-tvs(\allowbreak{}fld, m, t) iff is-normed-field(\allowbreak{}fld) and is-module(\allowbreak{}m) and scal(\allowbreak{}m) = normed-field-as-commutative-ring(\allowbreak{}fld) and is-hausdorff(\allowbreak{}t) and pts(\allowbreak{}t) == vec(\allowbreak{}m) and forall(\allowbreak{}[tvx\_, tvy\_, tvu\_], tvx\_ in vec(\allowbreak{}m) implies tvy\_ in vec(\allowbreak{}m) implies tvu\_ in opens(\allowbreak{}t) implies (\allowbreak{}vadd(\allowbreak{}m))(\allowbreak{}tvx\_, tvy\_) in tvu\_ implies forsome(\allowbreak{}[tvv\_ in opens(\allowbreak{}t)], tvx\_ in tvv\_ and forsome(\allowbreak{}[tvw\_ in opens(\allowbreak{}t)], tvy\_ in tvw\_ and forall(\allowbreak{}[tva\_ in tvv\_, tvb\_ in tvw\_], (\allowbreak{}vadd(\allowbreak{}m))(\allowbreak{}tva\_, tvb\_) in tvu\_)))) and forall(\allowbreak{}[tvr\_, tvx\_, tvu\_], tvr\_ in carr(\allowbreak{}scal(\allowbreak{}m)) implies tvx\_ in vec(\allowbreak{}m) implies tvu\_ in opens(\allowbreak{}t) implies (\allowbreak{}act(\allowbreak{}m))(\allowbreak{}tvr\_, tvx\_) in tvu\_ implies forsome(\allowbreak{}[tvd\_], pos-rr(\allowbreak{}tvd\_) and forsome(\allowbreak{}[tvw\_ in opens(\allowbreak{}t)], tvx\_ in tvw\_ and forall(\allowbreak{}[tvs\_ in carr(\allowbreak{}scal(\allowbreak{}m))], (\allowbreak{}fnrm(\allowbreak{}fld))(\allowbreak{}(\allowbreak{}add(\allowbreak{}scal(\allowbreak{}m)))(\allowbreak{}tvs\_, (\allowbreak{}neg(\allowbreak{}scal(\allowbreak{}m)))(\allowbreak{}tvr\_))) < tvd\_ implies forall(\allowbreak{}[tva\_ in tvw\_], (\allowbreak{}act(\allowbreak{}m))(\allowbreak{}tvs\_, tva\_) in tvu\_))))))
\end{quote}
also: \texttt{nvs-norm-top-is-tvs-cnvs-as-nvs} \texttt{rn-is-tvs} \texttt{sot-is-tvs} \texttt{wot-is-tvs} \texttt{prod-is-tvs} \texttt{cc-sn-top-is-tvs} \texttt{norm-optop-is-tvs} \texttt{nvs-norm-top-is-tvs} ... (9 in all) \\
mentioned by 51 result(s)

\smallskip

\noindent\textbf{\texttt{is-two-sided-ideal(baa\_, bai\_)}}\quad\textit{predicate}\index{is-two-sided-ideal(baa_, bai_)@\texttt{is-two-sided-ideal(baa\_, bai\_)}}\\
reads: \$2 is a two-sided ideal of \$1 \\
defined by \texttt{is-two-sided-ideal} \\
characterized by \texttt{is-two-sided-ideal}  \textit{(definitional)}:
\begin{quote}\ttfamily\footnotesize\raggedright\noindent
forall(\allowbreak{}[baa\_, bai\_], is-two-sided-ideal(\allowbreak{}baa\_, bai\_) iff is-left-ideal(\allowbreak{}baa\_, bai\_) and is-right-ideal(\allowbreak{}baa\_, bai\_))
\end{quote}
mentioned by 120 result(s)

\smallskip

\noindent\textbf{\texttt{is-ufd(gus\_)}}\quad\textit{predicate}\index{is-ufd(gus_)@\texttt{is-ufd(gus\_)}}\\
reads: \$1 is a unique factorization domain \\
defined by \texttt{is-ufd} \\
characterized by \texttt{is-ufd}  \textit{(definitional)}:
\begin{quote}\ttfamily\footnotesize\raggedright\noindent
forall(\allowbreak{}[gus\_], is-ufd(\allowbreak{}gus\_) iff is-integral-domain(\allowbreak{}gus\_) and forall(\allowbreak{}[gua\_ in carr(\allowbreak{}gus\_)], not(\allowbreak{}gua\_ = zero(\allowbreak{}gus\_)) implies forsome(\allowbreak{}[guu\_], is-unit(\allowbreak{}gus\_, guu\_) and forsome(\allowbreak{}[gun\_ in nn, gup\_], is-prime-seq(\allowbreak{}gus\_, gup\_, gun\_) and gua\_ = (\allowbreak{}mul(\allowbreak{}gus\_))(\allowbreak{}guu\_, prod-ring(\allowbreak{}gus\_, gup\_, interval(\allowbreak{}1, gun\_)))))) and forall(\allowbreak{}[gun\_ in nn, gup\_], is-prime-seq(\allowbreak{}gus\_, gup\_, gun\_) implies forall(\allowbreak{}[gum\_ in nn, guq\_], is-prime-seq(\allowbreak{}gus\_, guq\_, gum\_) implies forall(\allowbreak{}[guu\_], is-unit(\allowbreak{}gus\_, guu\_) implies forall(\allowbreak{}[guv\_], is-unit(\allowbreak{}gus\_, guv\_) implies (\allowbreak{}mul(\allowbreak{}gus\_))(\allowbreak{}guu\_, prod-ring(\allowbreak{}gus\_, gup\_, interval(\allowbreak{}1, gun\_))) = (\allowbreak{}mul(\allowbreak{}gus\_))(\allowbreak{}guv\_, prod-ring(\allowbreak{}gus\_, guq\_, interval(\allowbreak{}1, gum\_))) implies gun\_ = gum\_ and forsome(\allowbreak{}[gus2\_ in bijection(\allowbreak{}interval(\allowbreak{}1, gun\_), interval(\allowbreak{}1, gum\_))], forall(\allowbreak{}[gui\_ in interval(\allowbreak{}1, gun\_)], associated(\allowbreak{}gus\_, gup\_(\allowbreak{}gui\_), guq\_(\allowbreak{}gus2\_(\allowbreak{}gui\_))))))))))
\end{quote}
also: \texttt{euclidean-domain-is-ufd} \\
mentioned by 4 result(s)

\smallskip

\noindent\textbf{\texttt{is-ultrametric(ums\_)}}\quad\textit{predicate}\index{is-ultrametric(ums_)@\texttt{is-ultrametric(ums\_)}}\\
reads: \$1 is an ultrametric space \\
defined by \texttt{is-ultrametric} \\
characterized by \texttt{is-ultrametric}  \textit{(definitional)}:
\begin{quote}\ttfamily\footnotesize\raggedright\noindent
forall(\allowbreak{}[ums\_], is-ultrametric(\allowbreak{}ums\_) iff is-metric-space(\allowbreak{}ums\_) and forall(\allowbreak{}[umx\_ in pts(\allowbreak{}ums\_), umy\_ in pts(\allowbreak{}ums\_), umz\_ in pts(\allowbreak{}ums\_)], (\allowbreak{}dist(\allowbreak{}ums\_))(\allowbreak{}umx\_, umy\_) <= max(\allowbreak{}(\allowbreak{}dist(\allowbreak{}ums\_))(\allowbreak{}umx\_, umz\_), (\allowbreak{}dist(\allowbreak{}ums\_))(\allowbreak{}umz\_, umy\_))))
\end{quote}
also: \texttt{ultra-seq-is-ultrametric} \\
mentioned by 25 result(s)

\smallskip

\noindent\textbf{\texttt{is-unif-cauchy(s, fam)}}\quad\textit{predicate}\index{is-unif-cauchy(s, fam)@\texttt{is-unif-cauchy(s, fam)}}\\
reads: \$2 is uniformly Cauchy on \$1 \\
defined by \texttt{is-unif-cauchy} \\
characterized by \texttt{is-unif-cauchy}  \textit{(definitional)}:
\begin{quote}\ttfamily\footnotesize\raggedright\noindent
forall(\allowbreak{}[s, fam], is-unif-cauchy(\allowbreak{}s, fam) iff is-metric-space(\allowbreak{}s) and fam in fun(\allowbreak{}nn, fun(\allowbreak{}pts(\allowbreak{}s), rr)) and forall(\allowbreak{}[eps], pos-rr(\allowbreak{}eps) implies forsome(\allowbreak{}[cap in nn], forall(\allowbreak{}[k], k in nn and cap <= k implies forall(\allowbreak{}[l], l in nn and cap <= l implies forall(\allowbreak{}[x in pts(\allowbreak{}s)], abs(\allowbreak{}(\allowbreak{}fam(\allowbreak{}k))(\allowbreak{}x) - (\allowbreak{}fam(\allowbreak{}l))(\allowbreak{}x)) < eps))))))
\end{quote}
mentioned by 6 result(s)

\smallskip

\noindent\textbf{\texttt{is-unif-limit-on(fam, g, a, b)}}\quad\textit{predicate}\index{is-unif-limit-on(fam, g, a, b)@\texttt{is-unif-limit-on(fam, g, a, b)}}\\
reads: \$2 is the uniform limit of the sequence \$1 on [\$3, \$4] \\
defined by \texttt{is-unif-limit-on} \\
characterized by \texttt{is-unif-limit-on}  \textit{(definitional)}:
\begin{quote}\ttfamily\footnotesize\raggedright\noindent
forall(\allowbreak{}[fam, g, a, b], is-unif-limit-on(\allowbreak{}fam, g, a, b) iff fam in fun(\allowbreak{}nn, fun(\allowbreak{}rr, rr)) and g in fun(\allowbreak{}rr, rr) and a in rr and b in rr and a < b and forall(\allowbreak{}[eps\_], pos-rr(\allowbreak{}eps\_) implies forsome(\allowbreak{}[n\_ in nn], forall(\allowbreak{}[k\_], k\_ in nn and n\_ <= k\_ implies forall(\allowbreak{}[x\_ in ccint(\allowbreak{}a, b)], abs(\allowbreak{}g(\allowbreak{}x\_) - (\allowbreak{}fam(\allowbreak{}k\_))(\allowbreak{}x\_)) <= eps\_)))))
\end{quote}
mentioned by 6 result(s)

\smallskip

\noindent\textbf{\texttt{is-uniformly-continuous(s, t, f)}}\quad\textit{predicate}\index{is-uniformly-continuous(s, t, f)@\texttt{is-uniformly-continuous(s, t, f)}}\\
reads: \$3 is uniformly continuous from \$1 to \$2 \\
defined by \texttt{is-uniformly-continuous} \\
characterized by \texttt{is-uniformly-continuous}  \textit{(definitional)}:
\begin{quote}\ttfamily\footnotesize\raggedright\noindent
forall(\allowbreak{}[s, t, f], is-uniformly-continuous(\allowbreak{}s, t, f) iff is-metric-space(\allowbreak{}s) and is-metric-space(\allowbreak{}t) and f in fun(\allowbreak{}pts(\allowbreak{}s), pts(\allowbreak{}t)) and forall(\allowbreak{}[eps], pos-rr(\allowbreak{}eps) implies forsome(\allowbreak{}[delta], pos-rr(\allowbreak{}delta) and forall(\allowbreak{}[a in pts(\allowbreak{}s), b in pts(\allowbreak{}s)], (\allowbreak{}dist(\allowbreak{}s))(\allowbreak{}a, b) <= delta implies (\allowbreak{}dist(\allowbreak{}t))(\allowbreak{}f(\allowbreak{}a), f(\allowbreak{}b)) <= eps))))
\end{quote}
also: \texttt{is-uniformly-continuous-arrow-def} \\
mentioned by 26 result(s)

\smallskip

\noindent\textbf{\texttt{is-uniformly-continuous-arrow(a, b, f)}}\quad\textit{predicate}\index{is-uniformly-continuous-arrow(a, b, f)@\texttt{is-uniformly-continuous-arrow(a, b, f)}}\\
reads: \$3 is a uniformly continuous map from \$1 to \$2 \\
a predicate \\
characterized by \texttt{is-uniformly-continuous-arrow-def}  \textit{(definitional)}:
\begin{quote}\ttfamily\footnotesize\raggedright\noindent
forall(\allowbreak{}[a, b, f], is-uniformly-continuous-arrow(\allowbreak{}a, b, f) iff is-metric-space(\allowbreak{}a) and is-metric-space(\allowbreak{}b) and f in fun(\allowbreak{}pts(\allowbreak{}a), pts(\allowbreak{}b)) and is-uniformly-continuous(\allowbreak{}a, b, f))
\end{quote}
also: \texttt{lipschitz-arrow-is-uniformly-continuous-arrow} \\
mentioned by 8 result(s)

\smallskip

\noindent\textbf{\texttt{is-unipolar(c5k\_, c5w\_)}}\quad\textit{predicate}\index{is-unipolar(c5k_, c5w_)@\texttt{is-unipolar(c5k\_, c5w\_)}}\\
reads: \$2 is a unipolar bicharacter on \$1 \\
defined by \texttt{is-unipolar} \\
characterized by \texttt{is-unipolar}  \textit{(definitional)}:
\begin{quote}\ttfamily\footnotesize\raggedright\noindent
forall(\allowbreak{}[c5k\_, c5w\_], is-unipolar(\allowbreak{}c5k\_, c5w\_) iff is-bilinear-bicharacter(\allowbreak{}c5k\_, c5w\_) and bichar-polar(\allowbreak{}c5k\_, c5w\_, nv-integers(\allowbreak{}c5k\_)) = nv-integers(\allowbreak{}c5k\_))
\end{quote}
mentioned by 18 result(s)

\smallskip

\noindent\textbf{\texttt{is-unit(gus\_, guu\_)}}\quad\textit{predicate}\index{is-unit(gus_, guu_)@\texttt{is-unit(gus\_, guu\_)}}\\
reads: \$2 is a unit of \$1 \\
defined by \texttt{is-unit} \\
characterized by \texttt{is-unit}  \textit{(definitional)}:
\begin{quote}\ttfamily\footnotesize\raggedright\noindent
forall(\allowbreak{}[gus\_, guu\_], is-unit(\allowbreak{}gus\_, guu\_) iff guu\_ in carr(\allowbreak{}gus\_) and forsome(\allowbreak{}[guv\_ in carr(\allowbreak{}gus\_)], (\allowbreak{}mul(\allowbreak{}gus\_))(\allowbreak{}guu\_, guv\_) = one(\allowbreak{}gus\_)))
\end{quote}
mentioned by 13 result(s)

\smallskip

\noindent\textbf{\texttt{is-unital-banach-algebra(s)}}\quad\textit{predicate}\index{is-unital-banach-algebra(s)@\texttt{is-unital-banach-algebra(s)}}\\
reads: a Banach algebra with unit \\
a predicate \\
characterized by \texttt{is-unital-banach-algebra}  \textit{(definitional)}:
\begin{quote}\ttfamily\footnotesize\raggedright\noindent
forall(\allowbreak{}[s], is-unital-banach-algebra(\allowbreak{}s) iff length(\allowbreak{}s) = 9 and is-ring(\allowbreak{}scal(\allowbreak{}s)) and vec(\allowbreak{}s) in set and vadd(\allowbreak{}s) in fun(\allowbreak{}cartesian(\allowbreak{}vec(\allowbreak{}s), vec(\allowbreak{}s)), vec(\allowbreak{}s)) and vzero(\allowbreak{}s) in vec(\allowbreak{}s) and vneg(\allowbreak{}s) in fun(\allowbreak{}vec(\allowbreak{}s), vec(\allowbreak{}s)) and act(\allowbreak{}s) in fun(\allowbreak{}cartesian(\allowbreak{}carr(\allowbreak{}scal(\allowbreak{}s)), vec(\allowbreak{}s)), vec(\allowbreak{}s)) and vnrm(\allowbreak{}s) in fun(\allowbreak{}vec(\allowbreak{}s), rr) and vmul(\allowbreak{}s) in fun(\allowbreak{}cartesian(\allowbreak{}vec(\allowbreak{}s), vec(\allowbreak{}s)), vec(\allowbreak{}s)) and vone(\allowbreak{}s) in vec(\allowbreak{}s) and is-complex-banach-algebra(\allowbreak{}[scal(\allowbreak{}s), vec(\allowbreak{}s), vadd(\allowbreak{}s), vzero(\allowbreak{}s), vneg(\allowbreak{}s), act(\allowbreak{}s), vnrm(\allowbreak{}s), vmul(\allowbreak{}s)]) and carr(\allowbreak{}scal(\allowbreak{}s)) = cc and forall(\allowbreak{}[x\_ in vec(\allowbreak{}s)], (\allowbreak{}vmul(\allowbreak{}s))(\allowbreak{}vone(\allowbreak{}s), x\_) = x\_) and forall(\allowbreak{}[x\_ in vec(\allowbreak{}s)], (\allowbreak{}vmul(\allowbreak{}s))(\allowbreak{}x\_, vone(\allowbreak{}s)) = x\_) and (\allowbreak{}vnrm(\allowbreak{}s))(\allowbreak{}vone(\allowbreak{}s)) = 1 and s in set)
\end{quote}
also: \texttt{uiba-as-uba-is-unital-banach-algebra} \texttt{is-unital-banach-algebra-is-set} \texttt{unitization-is-unital-banach-algebra} \\
mentioned by 1631 result(s)

\smallskip

\noindent\textbf{\texttt{is-unital-c-star-algebra(s)}}\quad\textit{predicate}\index{is-unital-c-star-algebra(s)@\texttt{is-unital-c-star-algebra(s)}}\\
reads: a C*-algebra with unit \\
a predicate \\
characterized by \texttt{is-unital-c-star-algebra-def}  \textit{(definitional)}:
\begin{quote}\ttfamily\footnotesize\raggedright\noindent
forall(\allowbreak{}[s], is-unital-c-star-algebra(\allowbreak{}s) iff is-unital-involutive-banach-algebra(\allowbreak{}s) and forall(\allowbreak{}[x\_ in vec(\allowbreak{}s)], (\allowbreak{}vnrm(\allowbreak{}s))(\allowbreak{}(\allowbreak{}vmul(\allowbreak{}s))(\allowbreak{}(\allowbreak{}vstar(\allowbreak{}s))(\allowbreak{}x\_), x\_)) = (\allowbreak{}vnrm(\allowbreak{}s))(\allowbreak{}x\_) * (\allowbreak{}vnrm(\allowbreak{}s))(\allowbreak{}x\_)))
\end{quote}
also: \texttt{is-unital-c-star-algebra-is-set} \\
mentioned by 43 result(s)

\smallskip

\noindent\textbf{\texttt{is-unital-c-star-subalgebra(c2a\_, c2b\_)}}\quad\textit{predicate}\index{is-unital-c-star-subalgebra(c2a_, c2b_)@\texttt{is-unital-c-star-subalgebra(c2a\_, c2b\_)}}\\
reads: \$1 is a unital C*-subalgebra of \$2 containing its unit \\
defined by \texttt{is-unital-c-star-subalgebra} \\
characterized by \texttt{is-unital-c-star-subalgebra}  \textit{(definitional)}:
\begin{quote}\ttfamily\footnotesize\raggedright\noindent
forall(\allowbreak{}[c2a\_, c2b\_], is-unital-c-star-subalgebra(\allowbreak{}c2a\_, c2b\_) iff is-unital-c-star-algebra(\allowbreak{}c2a\_) and is-unital-c-star-algebra(\allowbreak{}c2b\_) and forall(\allowbreak{}[c2x\_ in vec(\allowbreak{}c2a\_)], c2x\_ in vec(\allowbreak{}c2b\_)) and vone(\allowbreak{}c2a\_) = vone(\allowbreak{}c2b\_) and forall(\allowbreak{}[c2x\_ in vec(\allowbreak{}c2a\_), c2y\_ in vec(\allowbreak{}c2a\_)], (\allowbreak{}vadd(\allowbreak{}c2a\_))(\allowbreak{}c2x\_, c2y\_) = (\allowbreak{}vadd(\allowbreak{}c2b\_))(\allowbreak{}c2x\_, c2y\_)) and forall(\allowbreak{}[c2x\_ in vec(\allowbreak{}c2a\_), c2y\_ in vec(\allowbreak{}c2a\_)], (\allowbreak{}vmul(\allowbreak{}c2a\_))(\allowbreak{}c2x\_, c2y\_) = (\allowbreak{}vmul(\allowbreak{}c2b\_))(\allowbreak{}c2x\_, c2y\_)) and forall(\allowbreak{}[c2z\_ in cc, c2x\_ in vec(\allowbreak{}c2a\_)], (\allowbreak{}act(\allowbreak{}c2a\_))(\allowbreak{}c2z\_, c2x\_) = (\allowbreak{}act(\allowbreak{}c2b\_))(\allowbreak{}c2z\_, c2x\_)) and forall(\allowbreak{}[c2x\_ in vec(\allowbreak{}c2a\_)], (\allowbreak{}vneg(\allowbreak{}c2a\_))(\allowbreak{}c2x\_) = (\allowbreak{}vneg(\allowbreak{}c2b\_))(\allowbreak{}c2x\_)) and forall(\allowbreak{}[c2x\_ in vec(\allowbreak{}c2a\_)], (\allowbreak{}vstar(\allowbreak{}c2a\_))(\allowbreak{}c2x\_) = (\allowbreak{}vstar(\allowbreak{}c2b\_))(\allowbreak{}c2x\_)) and forall(\allowbreak{}[c2x\_ in vec(\allowbreak{}c2a\_)], (\allowbreak{}vnrm(\allowbreak{}c2a\_))(\allowbreak{}c2x\_) = (\allowbreak{}vnrm(\allowbreak{}c2b\_))(\allowbreak{}c2x\_)))
\end{quote}
mentioned by 31 result(s)

\smallskip

\noindent\textbf{\texttt{is-unital-involutive-banach-algebra(s)}}\quad\textit{predicate}\index{is-unital-involutive-banach-algebra(s)@\texttt{is-unital-involutive-banach-algebra(s)}}\\
reads: an involutive Banach algebra with unit \\
a predicate \\
characterized by \texttt{is-unital-involutive-banach-algebra}  \textit{(definitional)}:
\begin{quote}\ttfamily\footnotesize\raggedright\noindent
forall(\allowbreak{}[s], is-unital-involutive-banach-algebra(\allowbreak{}s) iff length(\allowbreak{}s) = 10 and is-ring(\allowbreak{}scal(\allowbreak{}s)) and vec(\allowbreak{}s) in set and vadd(\allowbreak{}s) in fun(\allowbreak{}cartesian(\allowbreak{}vec(\allowbreak{}s), vec(\allowbreak{}s)), vec(\allowbreak{}s)) and vzero(\allowbreak{}s) in vec(\allowbreak{}s) and vneg(\allowbreak{}s) in fun(\allowbreak{}vec(\allowbreak{}s), vec(\allowbreak{}s)) and act(\allowbreak{}s) in fun(\allowbreak{}cartesian(\allowbreak{}carr(\allowbreak{}scal(\allowbreak{}s)), vec(\allowbreak{}s)), vec(\allowbreak{}s)) and vnrm(\allowbreak{}s) in fun(\allowbreak{}vec(\allowbreak{}s), rr) and vmul(\allowbreak{}s) in fun(\allowbreak{}cartesian(\allowbreak{}vec(\allowbreak{}s), vec(\allowbreak{}s)), vec(\allowbreak{}s)) and vone(\allowbreak{}s) in vec(\allowbreak{}s) and vstar(\allowbreak{}s) in fun(\allowbreak{}vec(\allowbreak{}s), vec(\allowbreak{}s)) and is-unital-banach-algebra(\allowbreak{}[scal(\allowbreak{}s), vec(\allowbreak{}s), vadd(\allowbreak{}s), vzero(\allowbreak{}s), vneg(\allowbreak{}s), act(\allowbreak{}s), vnrm(\allowbreak{}s), vmul(\allowbreak{}s), vone(\allowbreak{}s)]) and carr(\allowbreak{}scal(\allowbreak{}s)) = cc and forall(\allowbreak{}[x\_ in vec(\allowbreak{}s), y\_ in vec(\allowbreak{}s)], (\allowbreak{}vstar(\allowbreak{}s))(\allowbreak{}(\allowbreak{}vadd(\allowbreak{}s))(\allowbreak{}x\_, y\_)) = (\allowbreak{}vadd(\allowbreak{}s))(\allowbreak{}(\allowbreak{}vstar(\allowbreak{}s))(\allowbreak{}x\_), (\allowbreak{}vstar(\allowbreak{}s))(\allowbreak{}y\_))) and forall(\allowbreak{}[r\_ in cc, x\_ in vec(\allowbreak{}s)], (\allowbreak{}vstar(\allowbreak{}s))(\allowbreak{}(\allowbreak{}act(\allowbreak{}s))(\allowbreak{}r\_, x\_)) = (\allowbreak{}act(\allowbreak{}s))(\allowbreak{}conjugate(\allowbreak{}r\_), (\allowbreak{}vstar(\allowbreak{}s))(\allowbreak{}x\_))) and forall(\allowbreak{}[x\_ in vec(\allowbreak{}s)], (\allowbreak{}vstar(\allowbreak{}s))(\allowbreak{}(\allowbreak{}vstar(\allowbreak{}s))(\allowbreak{}x\_)) = x\_) and forall(\allowbreak{}[x\_ in vec(\allowbreak{}s), y\_ in vec(\allowbreak{}s)], (\allowbreak{}vstar(\allowbreak{}s))(\allowbreak{}(\allowbreak{}vmul(\allowbreak{}s))(\allowbreak{}x\_, y\_)) = (\allowbreak{}vmul(\allowbreak{}s))(\allowbreak{}(\allowbreak{}vstar(\allowbreak{}s))(\allowbreak{}y\_), (\allowbreak{}vstar(\allowbreak{}s))(\allowbreak{}x\_))) and forall(\allowbreak{}[x\_ in vec(\allowbreak{}s)], (\allowbreak{}vnrm(\allowbreak{}s))(\allowbreak{}(\allowbreak{}vstar(\allowbreak{}s))(\allowbreak{}x\_)) = (\allowbreak{}vnrm(\allowbreak{}s))(\allowbreak{}x\_)) and s in set)
\end{quote}
also: \texttt{is-unital-involutive-banach-algebra-is-set} \\
mentioned by 2449 result(s)

\smallskip

\noindent\textbf{\texttt{is-unitary-el(baa\_, bax\_)}}\quad\textit{predicate}\index{is-unitary-el(baa_, bax_)@\texttt{is-unitary-el(baa\_, bax\_)}}\\
reads: \$2 is unitary in \$1 \\
defined by \texttt{is-unitary-el} \\
characterized by \texttt{is-unitary-el}  \textit{(definitional)}:
\begin{quote}\ttfamily\footnotesize\raggedright\noindent
forall(\allowbreak{}[baa\_, bax\_], is-unitary-el(\allowbreak{}baa\_, bax\_) iff bax\_ in vec(\allowbreak{}baa\_) and (\allowbreak{}vmul(\allowbreak{}baa\_))(\allowbreak{}(\allowbreak{}vstar(\allowbreak{}baa\_))(\allowbreak{}bax\_), bax\_) = vone(\allowbreak{}baa\_) and (\allowbreak{}vmul(\allowbreak{}baa\_))(\allowbreak{}bax\_, (\allowbreak{}vstar(\allowbreak{}baa\_))(\allowbreak{}bax\_)) = vone(\allowbreak{}baa\_))
\end{quote}
mentioned by 2 result(s)

\smallskip

\noindent\textbf{\texttt{is-unitary-graph(ubv\_, ubg\_)}}\quad\textit{predicate}\index{is-unitary-graph(ubv_, ubg_)@\texttt{is-unitary-graph(ubv\_, ubg\_)}}\\
reads: \$2 is the graph of a unitary operator \\
defined by \texttt{is-unitary-graph} \\
characterized by \texttt{is-unitary-graph}  \textit{(definitional)}:
\begin{quote}\ttfamily\footnotesize\raggedright\noindent
forall(\allowbreak{}[ubv\_, ubg\_], is-unitary-graph(\allowbreak{}ubv\_, ubg\_) iff is-linear-graph(\allowbreak{}ubv\_, ubg\_) and is-isometric-graph(\allowbreak{}ubv\_, ubg\_) and graph-dom(\allowbreak{}ubv\_, ubg\_) = vec(\allowbreak{}ubv\_) and graph-ran(\allowbreak{}ubv\_, ubg\_) = vec(\allowbreak{}ubv\_))
\end{quote}
mentioned by 10 result(s)

\smallskip

\noindent\textbf{\texttt{is-upper-bound(grd, porel, ch, b)}}\quad\textit{predicate}\index{is-upper-bound(grd, porel, ch, b)@\texttt{is-upper-bound(grd, porel, ch, b)}}\\
reads: \$4 is an upper bound of \$3 in \$1 under \$2 \\
defined by \texttt{is-upper-bound} \\
characterized by \texttt{is-upper-bound}  \textit{(definitional)}:
\begin{quote}\ttfamily\footnotesize\raggedright\noindent
forall(\allowbreak{}[grd, porel, ch, b], is-upper-bound(\allowbreak{}grd, porel, ch, b) iff b in grd and forall(\allowbreak{}[x in ch], [x, b] in porel))
\end{quote}
mentioned by 6 result(s)

\smallskip

\noindent\textbf{\texttt{is-upper-triangular(a, n, d)}}\quad\textit{predicate}\index{is-upper-triangular(a, n, d)@\texttt{is-upper-triangular(a, n, d)}}\\
reads: \$3 is an upper triangular \$2-by-\$2 matrix over \$1 \\
defined by \texttt{is-upper-triangular} \\
characterized by \texttt{is-upper-triangular}  \textit{(definitional)}:
\begin{quote}\ttfamily\footnotesize\raggedright\noindent
forall(\allowbreak{}[a, n, d], is-upper-triangular(\allowbreak{}a, n, d) iff forall(\allowbreak{}[ltu\_, ltv\_], ltu\_ in interval(\allowbreak{}1, n) implies ltv\_ in interval(\allowbreak{}1, n) implies ltv\_ < ltu\_ implies entry(\allowbreak{}d, ltu\_, ltv\_) = zero(\allowbreak{}a)))
\end{quote}
mentioned by 9 result(s)

\smallskip

\noindent\textbf{\texttt{is-vector-space(s)}}\quad\textit{predicate}\index{is-vector-space(s)@\texttt{is-vector-space(s)}}\\
reads: a vector space \\
a predicate \\
characterized by \texttt{is-vector-space-def}  \textit{(definitional)}:
\begin{quote}\ttfamily\footnotesize\raggedright\noindent
forall(\allowbreak{}[s], is-vector-space(\allowbreak{}s) iff is-module(\allowbreak{}s) and is-field-ring(\allowbreak{}scal(\allowbreak{}s)))
\end{quote}
also: \texttt{is-vector-space-is-set} \texttt{is-vector-space-unfold} \texttt{qq-line-is-vector-space} \texttt{rr-fn-space-is-vector-space} \\
mentioned by 22 result(s)

\smallskip

\noindent\textbf{\texttt{is-von-neumann-algebra(vnv\_, vns\_)}}\quad\textit{predicate}\index{is-von-neumann-algebra(vnv_, vns_)@\texttt{is-von-neumann-algebra(vnv\_, vns\_)}}\\
reads: \$2 is a von Neumann algebra on \$1 \\
defined by \texttt{is-von-neumann-algebra} \\
characterized by \texttt{is-von-neumann-algebra}  \textit{(definitional)}:
\begin{quote}\ttfamily\footnotesize\raggedright\noindent
forall(\allowbreak{}[vnv\_, vns\_], is-von-neumann-algebra(\allowbreak{}vnv\_, vns\_) iff is-star-op-algebra(\allowbreak{}vnv\_, vns\_) and is-top-closed(\allowbreak{}sot(\allowbreak{}vnv\_), vns\_))
\end{quote}
mentioned by 8 result(s)

\smallskip

\noindent\textbf{\texttt{isometry-group(igs\_)}}\quad\textit{functoid}\index{isometry-group(igs_)@\texttt{isometry-group(igs\_)}}\\
reads: the group of isometries of \$1 onto itself \\
definition:
\begin{quote}\ttfamily\footnotesize\raggedright\noindent
\{igf\_ in bijection(\allowbreak{}pts(\allowbreak{}igs\_), pts(\allowbreak{}igs\_)): is-isometry(\allowbreak{}igs\_, igs\_, igf\_)\}
\end{quote}
characterized by \texttt{isometry-group-id-in}  \textit{(proven)}:
\begin{quote}\ttfamily\footnotesize\raggedright\noindent
forall(\allowbreak{}[s], is-metric-space(\allowbreak{}s) implies vnb-lambda(\allowbreak{}igx\_, pts(\allowbreak{}s), igx\_) in isometry-group(\allowbreak{}s))
\end{quote}
also: \texttt{isometry-group-in-fun} \texttt{isometry-group-is-set} \texttt{isometry-group-ms-pts} \texttt{isometry-group-unfold} \texttt{isometry-group-grp-inv} \texttt{isometry-group-grp-opr} \texttt{isometry-group-ms-dist} \texttt{isometry-group-grp-carr} ... (21 in all) \\
mentioned by 34 result(s)

\smallskip

\noindent\textbf{\texttt{isometry-group-grp(igs\_)}}\quad\textit{functoid}\index{isometry-group-grp(igs_)@\texttt{isometry-group-grp(igs\_)}}\\
reads: the group of isometries of \$1 onto itself under composition \\
definition:
\begin{quote}\ttfamily\footnotesize\raggedright\noindent
[isometry-group(\allowbreak{}igs\_), vnb-lambda(\allowbreak{}[igu\_, igv\_], cartesian(\allowbreak{}isometry-group(\allowbreak{}igs\_), isometry-group(\allowbreak{}igs\_)), vnb-lambda(\allowbreak{}igx\_, pts(\allowbreak{}igs\_), igu\_(\allowbreak{}igv\_(\allowbreak{}igx\_)))), vnb-lambda(\allowbreak{}igx\_, pts(\allowbreak{}igs\_), igx\_), vnb-lambda(\allowbreak{}igi\_, isometry-group(\allowbreak{}igs\_), inverse-bij(\allowbreak{}igi\_, pts(\allowbreak{}igs\_), pts(\allowbreak{}igs\_)))]
\end{quote}
characterized by \texttt{isometry-group-grp-inv}  \textit{(proven)}:
\begin{quote}\ttfamily\footnotesize\raggedright\noindent
forall(\allowbreak{}[s], is-metric-space(\allowbreak{}s) implies inv(\allowbreak{}isometry-group-grp(\allowbreak{}s)) == vnb-lambda(\allowbreak{}igi\_, isometry-group(\allowbreak{}s), inverse-bij(\allowbreak{}igi\_, pts(\allowbreak{}s), pts(\allowbreak{}s))))
\end{quote}
also: \texttt{isometry-group-grp-opr} \texttt{isometry-group-grp-carr} \texttt{isometry-group-grp-iden} \texttt{isometry-group-grp-is-group} \\
mentioned by 15 result(s)

\smallskip

\noindent\textbf{\texttt{isometry-group-ms(ims\_)}}\quad\textit{functoid}\index{isometry-group-ms(ims_)@\texttt{isometry-group-ms(ims\_)}}\\
reads: the group of isometries of \$1 with the sup metric \\
definition:
\begin{quote}\ttfamily\footnotesize\raggedright\noindent
[isometry-group(\allowbreak{}ims\_), vnb-lambda(\allowbreak{}[imu\_, imv\_], cartesian(\allowbreak{}isometry-group(\allowbreak{}ims\_), isometry-group(\allowbreak{}ims\_)), sup-dist(\allowbreak{}ims\_, ims\_, imu\_, imv\_))]
\end{quote}
characterized by \texttt{isometry-group-ms-pts}  \textit{(proven)}:
\begin{quote}\ttfamily\footnotesize\raggedright\noindent
forall(\allowbreak{}[s], is-metric-space(\allowbreak{}s) implies pts(\allowbreak{}isometry-group-ms(\allowbreak{}s)) == isometry-group(\allowbreak{}s))
\end{quote}
also: \texttt{isometry-group-ms-dist} \texttt{isometry-group-ms-dist-fun} \texttt{isometry-group-ms-is-metric-space} \\
mentioned by 18 result(s)

\smallskip

\noindent\textbf{\texttt{ispace(s)}}\quad\textit{accessor}\index{ispace(s)@\texttt{ispace(s)}}\\
reads: the ispace slot of \$1 \\
a structure slot \\
characterized by \texttt{cx-isp@ispace}  \textit{(definitional)}:
\begin{quote}\ttfamily\footnotesize\raggedright\noindent
ispace(\allowbreak{}cx-isp) == cx-sp
\end{quote}
mentioned by 122 result(s)

\smallskip

\noindent\textbf{\texttt{istar(s)}}\quad\textit{accessor}\index{istar(s)@\texttt{istar(s)}}\\
reads: the involution of \$1 \\
a structure slot \\
characterized by \texttt{b1s-istar-fun}  \textit{(proven)}:
\begin{quote}\ttfamily\footnotesize\raggedright\noindent
forall(\allowbreak{}[a], is-nu-involutive-banach-algebra(\allowbreak{}a) implies istar(\allowbreak{}a) in fun(\allowbreak{}vec(\allowbreak{}a), vec(\allowbreak{}a)))
\end{quote}
also: \texttt{b1s-istar-isom} \texttt{b1s-istar-apply} \\
mentioned by 22 result(s)

\smallskip


\section*{J}
\addcontentsline{toc}{section}{J}

\noindent\textbf{\texttt{j-cart(mit\_)}}\quad\textit{functoid}\index{j-cart(mit_)@\texttt{j-cart(mit\_)}}\\
reads: the local instances of cartesian-decompose \\
definition:
\begin{quote}\ttfamily\footnotesize\raggedright\noindent
\{mix\_ in rw-univ: forsome(\allowbreak{}[mig\_, miz\_, mia\_], miz\_ in set and mia\_ in set and mix\_ = [2, mig\_, enc-in(\allowbreak{}miz\_, enc-cart(\allowbreak{}mia\_)), enc-cch(\allowbreak{}miz\_, mia\_, length(\allowbreak{}mia\_), length(\allowbreak{}mia\_))] and mig\_ in ctx-set and 1 <= length(\allowbreak{}mia\_) and enc-in(\allowbreak{}miz\_, enc-cart(\allowbreak{}mia\_)) in clf and enc-cch(\allowbreak{}miz\_, mia\_, length(\allowbreak{}mia\_), length(\allowbreak{}mia\_)) in clf)\}
\end{quote}
characterized by \texttt{j-cart-sub}  \textit{(proven)}:
\begin{quote}\ttfamily\footnotesize\raggedright\noindent
forall(\allowbreak{}[tt], tt subset clf implies ax-mem-set in tt implies j-cart(\allowbreak{}tt) subset eqv(\allowbreak{}tt))
\end{quote}
mentioned by 5 result(s)

\smallskip

\noindent\textbf{\texttt{j-fb(mrt\_)}}\quad\textit{functoid}\index{j-fb(mrt_)@\texttt{j-fb(mrt\_)}}\\
reads: the local functoid redexes over \$1, at any arguments \\
definition:
\begin{quote}\ttfamily\footnotesize\raggedright\noindent
\{mrx\_ in rw-univ: forsome(\allowbreak{}[mrg\_, mrh\_, mrl\_, mrb\_], mrg\_ in set and mrh\_ in set and mrl\_ in set and mrb\_ in set and mrx\_ = [1, mrg\_, eapp(\allowbreak{}mrh\_, mrl\_), openl(\allowbreak{}mrb\_, mrl\_, length(\allowbreak{}mrl\_))] and fun-def(\allowbreak{}mrh\_, length(\allowbreak{}mrl\_), mrb\_) and eapp(\allowbreak{}mrh\_, mrl\_) in clt and openl(\allowbreak{}mrb\_, mrl\_, length(\allowbreak{}mrl\_)) in clt)\}
\end{quote}
characterized by \texttt{j-fb-sub}  \textit{(proven)}:
\begin{quote}\ttfamily\footnotesize\raggedright\noindent
forall(\allowbreak{}[tt], j-fb(\allowbreak{}tt) subset eqv(\allowbreak{}tt))
\end{quote}
mentioned by 5 result(s)

\smallskip

\noindent\textbf{\texttt{j-idec(md14t\_)}}\quad\textit{functoid}\index{j-idec(md14t_)@\texttt{j-idec(md14t\_)}}\\
reads: the local instances over \$1 of the catalog's intersection-decompose \\
definition:
\begin{quote}\ttfamily\footnotesize\raggedright\noindent
\{md14y\_ in rw-univ: forsome(\allowbreak{}[md14g\_, md14x2\_, md14l2\_, md14f2\_], md14g\_ in set and md14x2\_ in set and md14l2\_ in set and md14y\_ = [2, md14g\_, enc-in(\allowbreak{}md14x2\_, enc-inter(\allowbreak{}md14l2\_)), md14f2\_] and [0, md14f2\_] in rf-set(\allowbreak{}3, md14x2\_, md14l2\_) and enc-in(\allowbreak{}md14x2\_, enc-inter(\allowbreak{}md14l2\_)) in clf and md14f2\_ in clf)\}
\end{quote}
characterized by \texttt{j-idec-sub}  \textit{(proven)}:
\begin{quote}\ttfamily\footnotesize\raggedright\noindent
forall(\allowbreak{}[tt], j-idec(\allowbreak{}tt) subset eqv(\allowbreak{}tt))
\end{quote}
mentioned by 9 result(s)

\smallskip

\noindent\textbf{\texttt{j-lb(mrt\_)}}\quad\textit{functoid}\index{j-lb(mrt_)@\texttt{j-lb(mrt\_)}}\\
reads: the local lambda redexes over \$1, the argument in the domain \\
definition:
\begin{quote}\ttfamily\footnotesize\raggedright\noindent
\{mrx\_ in rw-univ: forsome(\allowbreak{}[mrg\_, mra\_, mrb\_, mre\_], mrg\_ in set and mra\_ in set and mrb\_ in set and mre\_ in set and mrx\_ = [1, mrg\_, enc-apply(\allowbreak{}enc-lam(\allowbreak{}mra\_, mrb\_), mre\_), expr-open(\allowbreak{}mrb\_, 0, mre\_)] and nbg-ok(\allowbreak{}mrb\_) and enc-apply(\allowbreak{}enc-lam(\allowbreak{}mra\_, mrb\_), mre\_) in clt and expr-open(\allowbreak{}mrb\_, 0, mre\_) in clt and [mrg\_, enc-in(\allowbreak{}mre\_, mra\_)] in derivable(\allowbreak{}mrt\_))\}
\end{quote}
characterized by \texttt{j-lb-sub}  \textit{(proven)}:
\begin{quote}\ttfamily\footnotesize\raggedright\noindent
forall(\allowbreak{}[tt], j-lb(\allowbreak{}tt) subset eqv(\allowbreak{}tt))
\end{quote}
mentioned by 9 result(s)

\smallskip

\noindent\textbf{\texttt{j-len(mrt\_)}}\quad\textit{functoid}\index{j-len(mrt_)@\texttt{j-len(mrt\_)}}\\
reads: the local LENGTH-of-a-list redexes over \$1, the list denoting \\
definition:
\begin{quote}\ttfamily\footnotesize\raggedright\noindent
\{mrx\_ in rw-univ: forsome(\allowbreak{}[mrg\_, mrl\_], mrg\_ in set and mrl\_ in set and mrx\_ = [1, mrg\_, enc-length(\allowbreak{}enc-list(\allowbreak{}mrl\_)), enc-num(\allowbreak{}length(\allowbreak{}mrl\_))] and enc-length(\allowbreak{}enc-list(\allowbreak{}mrl\_)) in clt and enc-num(\allowbreak{}length(\allowbreak{}mrl\_)) in clt and [mrg\_, enc-eq(\allowbreak{}enc-list(\allowbreak{}mrl\_), enc-list(\allowbreak{}mrl\_))] in derivable(\allowbreak{}mrt\_))\}
\end{quote}
characterized by \texttt{j-len-sub}  \textit{(proven)}:
\begin{quote}\ttfamily\footnotesize\raggedright\noindent
forall(\allowbreak{}[tt], j-len(\allowbreak{}tt) subset eqv(\allowbreak{}tt))
\end{quote}
mentioned by 5 result(s)

\smallskip

\noindent\textbf{\texttt{j-mac(mrt\_)}}\quad\textit{functoid}\index{j-mac(mrt_)@\texttt{j-mac(mrt\_)}}\\
reads: the local instances of the one-variable macetes of \$1, the condition derivable and the instance denoting \\
definition:
\begin{quote}\ttfamily\footnotesize\raggedright\noindent
\{mrx\_ in rw-univ: forsome(\allowbreak{}[mrg\_, mrc\_, mra\_, mrb\_, mre\_], mrg\_ in set and mrc\_ in set and mra\_ in set and mrb\_ in set and mre\_ in clt and [mrg\_, expr-open(\allowbreak{}mrc\_, 0, mre\_)] in derivable(\allowbreak{}mrt\_) and [mrg\_, enc-eq(\allowbreak{}mre\_, mre\_)] in derivable(\allowbreak{}mrt\_) and (\allowbreak{}mrx\_ = [1, mrg\_, expr-open(\allowbreak{}mra\_, 0, mre\_), expr-open(\allowbreak{}mrb\_, 0, mre\_)] and enc-all(\allowbreak{}enc-imp(\allowbreak{}mrc\_, enc-qeq(\allowbreak{}mra\_, mrb\_))) in mrt\_ and enc-all(\allowbreak{}enc-imp(\allowbreak{}mrc\_, enc-qeq(\allowbreak{}mra\_, mrb\_))) in clf and expr-open(\allowbreak{}mra\_, 0, mre\_) in clt and expr-open(\allowbreak{}mrb\_, 0, mre\_) in clt or mrx\_ = [2, mrg\_, expr-open(\allowbreak{}mra\_, 0, mre\_), expr-open(\allowbreak{}mrb\_, 0, mre\_)] and enc-all(\allowbreak{}enc-imp(\allowbreak{}mrc\_, enc-iff(\allowbreak{}mra\_, mrb\_))) in mrt\_ and enc-all(\allowbreak{}enc-imp(\allowbreak{}mrc\_, enc-iff(\allowbreak{}mra\_, mrb\_))) in clf and expr-open(\allowbreak{}mra\_, 0, mre\_) in clf and expr-open(\allowbreak{}mrb\_, 0, mre\_) in clf))\}
\end{quote}
characterized by \texttt{j-mac-sub}  \textit{(proven)}:
\begin{quote}\ttfamily\footnotesize\raggedright\noindent
forall(\allowbreak{}[tt], j-mac(\allowbreak{}tt) subset eqv(\allowbreak{}tt))
\end{quote}
mentioned by 9 result(s)

\smallskip

\noindent\textbf{\texttt{j-macn(md13t\_)}}\quad\textit{functoid}\index{j-macn(md13t_)@\texttt{j-macn(md13t\_)}}\\
reads: the local instances of the macetes of \$1 of any number of schema variables, the condition derivable and every value denoting \\
definition:
\begin{quote}\ttfamily\footnotesize\raggedright\noindent
\{md13x\_ in rw-univ: forsome(\allowbreak{}[md13g\_, md13m\_, md13l\_, md13i\_, md13r\_, md13a\_, md13b\_], md13g\_ in set and md13l\_ in set and md13r\_ in set and md13a\_ in set and md13b\_ in set and md13m\_ in md13t\_ and md13m\_ in clf and md13i\_ in inst-reach(\allowbreak{}md13m\_, md13l\_) and forall(\allowbreak{}[md13j\_ in nn], 2 <= md13j\_ implies md13j\_ <= length(\allowbreak{}md13l\_) implies [md13g\_, enc-eq(\allowbreak{}nth(\allowbreak{}md13j\_, md13l\_), nth(\allowbreak{}md13j\_, md13l\_))] in derivable(\allowbreak{}md13t\_)) and (\allowbreak{}md13i\_ = md13r\_ or forsome(\allowbreak{}[md13c\_ in set], md13i\_ = enc-imp(\allowbreak{}md13c\_, md13r\_) and [md13g\_, md13c\_] in derivable(\allowbreak{}md13t\_))) and (\allowbreak{}md13r\_ = enc-qeq(\allowbreak{}md13a\_, md13b\_) and md13x\_ = [1, md13g\_, md13a\_, md13b\_] and md13a\_ in clt and md13b\_ in clt or md13r\_ = enc-iff(\allowbreak{}md13a\_, md13b\_) and md13x\_ = [2, md13g\_, md13a\_, md13b\_] and md13a\_ in clf and md13b\_ in clf))\}
\end{quote}
characterized by \texttt{j-macn-sub}  \textit{(proven)}:
\begin{quote}\ttfamily\footnotesize\raggedright\noindent
forall(\allowbreak{}[tt], j-macn(\allowbreak{}tt) subset eqv(\allowbreak{}tt))
\end{quote}
mentioned by 9 result(s)

\smallskip

\noindent\textbf{\texttt{j-nth(mrt\_)}}\quad\textit{functoid}\index{j-nth(mrt_)@\texttt{j-nth(mrt\_)}}\\
reads: the local NTH-of-a-list redexes over \$1, the list denoting \\
definition:
\begin{quote}\ttfamily\footnotesize\raggedright\noindent
\{mrx\_ in rw-univ: forsome(\allowbreak{}[mrg\_, mrl\_, mrk\_], mrg\_ in set and mrl\_ in set and mrx\_ = [1, mrg\_, enc-nth(\allowbreak{}enc-num(\allowbreak{}mrk\_), enc-list(\allowbreak{}mrl\_)), nth(\allowbreak{}mrk\_, mrl\_)] and mrk\_ in nn and 1 <= mrk\_ and mrk\_ <= length(\allowbreak{}mrl\_) and enc-nth(\allowbreak{}enc-num(\allowbreak{}mrk\_), enc-list(\allowbreak{}mrl\_)) in clt and nth(\allowbreak{}mrk\_, mrl\_) in clt and [mrg\_, enc-eq(\allowbreak{}enc-list(\allowbreak{}mrl\_), enc-list(\allowbreak{}mrl\_))] in derivable(\allowbreak{}mrt\_))\}
\end{quote}
characterized by \texttt{j-nth-sub}  \textit{(proven)}:
\begin{quote}\ttfamily\footnotesize\raggedright\noindent
forall(\allowbreak{}[tt], j-nth(\allowbreak{}tt) subset eqv(\allowbreak{}tt))
\end{quote}
mentioned by 5 result(s)

\smallskip

\noindent\textbf{\texttt{j-teq(mit\_)}}\quad\textit{functoid}\index{j-teq(mit_)@\texttt{j-teq(mit\_)}}\\
reads: the local instances of tuple-equality-decompose, both lists denoting \\
definition:
\begin{quote}\ttfamily\footnotesize\raggedright\noindent
\{mix\_ in rw-univ: forsome(\allowbreak{}[mig\_, mia\_, mib\_, mif\_], mia\_ in set and mib\_ in set and mif\_ in set and mix\_ = [2, mig\_, enc-eq(\allowbreak{}enc-list(\allowbreak{}mia\_), enc-list(\allowbreak{}mib\_)), mif\_] and mig\_ in ctx-set and length(\allowbreak{}mia\_) = length(\allowbreak{}mib\_) and enc-eq(\allowbreak{}enc-list(\allowbreak{}mia\_), enc-list(\allowbreak{}mib\_)) in clf and mif\_ in clf and (\allowbreak{}length(\allowbreak{}mia\_) = 0 and mif\_ = enc-truth or [0, mif\_] in tf-set(\allowbreak{}mia\_, mib\_)) and [mig\_, enc-eq(\allowbreak{}enc-list(\allowbreak{}mia\_), enc-list(\allowbreak{}mia\_))] in derivable(\allowbreak{}mit\_) and [mig\_, enc-eq(\allowbreak{}enc-list(\allowbreak{}mib\_), enc-list(\allowbreak{}mib\_))] in derivable(\allowbreak{}mit\_))\}
\end{quote}
characterized by \texttt{j-teq-sub}  \textit{(proven)}:
\begin{quote}\ttfamily\footnotesize\raggedright\noindent
forall(\allowbreak{}[tt], j-teq(\allowbreak{}tt) subset eqv(\allowbreak{}tt))
\end{quote}
mentioned by 5 result(s)

\smallskip

\noindent\textbf{\texttt{j-udec(md14t\_)}}\quad\textit{functoid}\index{j-udec(md14t_)@\texttt{j-udec(md14t\_)}}\\
reads: the local instances over \$1 of the catalog's union-decompose \\
definition:
\begin{quote}\ttfamily\footnotesize\raggedright\noindent
\{md14y\_ in rw-univ: forsome(\allowbreak{}[md14g\_, md14x2\_, md14l2\_, md14f2\_], md14g\_ in set and md14x2\_ in set and md14l2\_ in set and md14y\_ = [2, md14g\_, enc-in(\allowbreak{}md14x2\_, enc-union(\allowbreak{}md14l2\_)), md14f2\_] and [0, md14f2\_] in rf-set(\allowbreak{}4, md14x2\_, md14l2\_) and enc-in(\allowbreak{}md14x2\_, enc-union(\allowbreak{}md14l2\_)) in clf and md14f2\_ in clf)\}
\end{quote}
characterized by \texttt{j-udec-sub}  \textit{(proven)}:
\begin{quote}\ttfamily\footnotesize\raggedright\noindent
forall(\allowbreak{}[tt], j-udec(\allowbreak{}tt) subset eqv(\allowbreak{}tt))
\end{quote}
mentioned by 9 result(s)

\smallskip

\noindent\textbf{\texttt{jop(s)}}\quad\textit{accessor}\index{jop(s)@\texttt{jop(s)}}\\
reads: the complex structure of \$1 \\
a structure slot \\
characterized by \texttt{cnvs-jop-sq}  \textit{(proven)}:
\begin{quote}\ttfamily\footnotesize\raggedright\noindent
forall(\allowbreak{}[v], is-complex-normed-vector-space(\allowbreak{}v) implies forall(\allowbreak{}[x\_ in vec(\allowbreak{}v)], (\allowbreak{}jop(\allowbreak{}v))(\allowbreak{}(\allowbreak{}jop(\allowbreak{}v))(\allowbreak{}x\_)) = (\allowbreak{}vneg(\allowbreak{}v))(\allowbreak{}x\_)))
\end{quote}
also: \texttt{cnvs-nf-jop} \texttt{csip-jop-sq} \texttt{csip-nf-jop} \texttt{cnvs-jop-act} \texttt{cnvs-jop-add} \texttt{cnvs-jop-fun} \texttt{csip-jop-act} \texttt{csip-jop-add} ... (25 in all) \\
mentioned by 88 result(s)

\smallskip

\noindent\textbf{\texttt{juxta(jxp\_, jxq\_, jxa\_, jxb\_, jxc\_)}}\quad\textit{functoid}\index{juxta(jxp_, jxq_, jxa_, jxb_, jxc_)@\texttt{juxta(jxp\_, jxq\_, jxa\_, jxb\_, jxc\_)}}\\
reads: the juxtaposition of \$1 on [\$3, \$4] and \$2 on [\$4, \$5] \\
definition:
\begin{quote}\ttfamily\footnotesize\raggedright\noindent
vnb-lambda(\allowbreak{}jxt\_, ccint(\allowbreak{}jxa\_, jxc\_), if(\allowbreak{}jxt\_ <= jxb\_, jxp\_(\allowbreak{}jxt\_), jxq\_(\allowbreak{}jxt\_)))
\end{quote}
characterized by \texttt{juxta-trace}  \textit{(proven)}:
\begin{quote}\ttfamily\footnotesize\raggedright\noindent
forall(\allowbreak{}[p, q, a in rr, b in rr, c in rr], a <= b implies b <= c implies p in fun(\allowbreak{}ccint(\allowbreak{}a, b), cc) implies q in fun(\allowbreak{}ccint(\allowbreak{}b, c), cc) implies p(\allowbreak{}b) = q(\allowbreak{}b) implies trace(\allowbreak{}juxta(\allowbreak{}p, q, a, b, c), a, c) = union(\allowbreak{}trace(\allowbreak{}p, a, b), trace(\allowbreak{}q, b, c)))
\end{quote}
also: \texttt{juxta-in-fun} \texttt{juxta-unfold} \texttt{juxta-is-road} \texttt{winding-juxta} \texttt{line-int-juxta} \texttt{juxta-apply-left} \texttt{juxta-apply-right} \texttt{road-length-juxta} ... (10 in all) \\
mentioned by 29 result(s)

\smallskip


\section*{K}
\addcontentsline{toc}{section}{K}

\noindent\textbf{\texttt{keep}}\quad\textit{tactic}\index{keep@\texttt{keep}}\\
\texttt{(keep hyp ...)} -- Keep only the cited assumptions and drop every other one, as ONE recorded step.  Each `hyp' may be a formula, a "string", or a 1-based assumption index; they are recorded as formulas. \\
declared in \texttt{tactics-help.scm}

\smallskip

\noindent\textbf{\texttt{keep-set(phi, grd, porel, kset, alpha)}}\quad\textit{functoid}\index{keep-set(phi, grd, porel, kset, alpha)@\texttt{keep-set(phi, grd, porel, kset, alpha)}}\\
definition:
\begin{quote}\ttfamily\footnotesize\raggedright\noindent
\{y in grd: y = phi(\allowbreak{}alpha) and forall(\allowbreak{}[z in kset], [z, y] in porel)\}
\end{quote}
characterized by \texttt{keep-set-membership}  \textit{(definitional)}:
\begin{quote}\ttfamily\footnotesize\raggedright\noindent
forall(\allowbreak{}[phi, grd, porel, kset, alpha, y\_], y\_ in keep-set(\allowbreak{}phi, grd, porel, kset, alpha) iff y\_ in grd and y\_ = phi(\allowbreak{}alpha) and forall(\allowbreak{}[z in kset], [z, y\_] in porel))
\end{quote}
mentioned by 5 result(s)

\smallskip

\noindent\textbf{\texttt{kernel-apply(kaxs\_, kays\_, kacb\_, kanu\_, kak\_, kaf\_)}}\quad\textit{functoid}\index{kernel-apply(kaxs_, kays_, kacb_, kanu_, kak_, kaf_)@\texttt{kernel-apply(kaxs\_, kays\_, kacb\_, kanu\_, kak\_, kaf\_)}}\\
reads: x |-> the integral over \$2 of \$5(x, y) \$6(y) d\$4(y) \\
definition:
\begin{quote}\ttfamily\footnotesize\raggedright\noindent
vnb-lambda(\allowbreak{}sgx\_, kaxs\_, if(\allowbreak{}(\allowbreak{}x-section-integral(\allowbreak{}kaxs\_, kays\_, kacb\_, kanu\_, pos-part(\allowbreak{}cartesian(\allowbreak{}kaxs\_, kays\_), vnb-lambda(\allowbreak{}kpw\_, cartesian(\allowbreak{}kaxs\_, kays\_), kak\_(\allowbreak{}kpw\_) * kaf\_(\allowbreak{}nth(\allowbreak{}2, kpw\_))))))(\allowbreak{}sgx\_) in rr, (\allowbreak{}x-section-integral(\allowbreak{}kaxs\_, kays\_, kacb\_, kanu\_, pos-part(\allowbreak{}cartesian(\allowbreak{}kaxs\_, kays\_), vnb-lambda(\allowbreak{}kpw\_, cartesian(\allowbreak{}kaxs\_, kays\_), kak\_(\allowbreak{}kpw\_) * kaf\_(\allowbreak{}nth(\allowbreak{}2, kpw\_))))))(\allowbreak{}sgx\_), 0) - if(\allowbreak{}(\allowbreak{}x-section-integral(\allowbreak{}kaxs\_, kays\_, kacb\_, kanu\_, neg-part(\allowbreak{}cartesian(\allowbreak{}kaxs\_, kays\_), vnb-lambda(\allowbreak{}kpw\_, cartesian(\allowbreak{}kaxs\_, kays\_), kak\_(\allowbreak{}kpw\_) * kaf\_(\allowbreak{}nth(\allowbreak{}2, kpw\_))))))(\allowbreak{}sgx\_) in rr, (\allowbreak{}x-section-integral(\allowbreak{}kaxs\_, kays\_, kacb\_, kanu\_, neg-part(\allowbreak{}cartesian(\allowbreak{}kaxs\_, kays\_), vnb-lambda(\allowbreak{}kpw\_, cartesian(\allowbreak{}kaxs\_, kays\_), kak\_(\allowbreak{}kpw\_) * kaf\_(\allowbreak{}nth(\allowbreak{}2, kpw\_))))))(\allowbreak{}sgx\_), 0))
\end{quote}
characterized by \texttt{kernel-apply-spec}  \textit{(proven)}:
\begin{quote}\ttfamily\footnotesize\raggedright\noindent
forall(\allowbreak{}[xs, ca, mu, ys, cb, nu, k, xi], is-measure(\allowbreak{}xs, ca, mu) implies is-measure(\allowbreak{}ys, cb, nu) implies is-sigma-finite(\allowbreak{}xs, ca, mu) implies is-sigma-finite(\allowbreak{}ys, cb, nu) implies k in l2-set(\allowbreak{}cartesian(\allowbreak{}xs, ys), product-sigma(\allowbreak{}xs, ca, ys, cb), product-measure(\allowbreak{}xs, ca, mu, ys, cb, nu)) implies xi in l2-set(\allowbreak{}ys, cb, nu) implies kernel-apply(\allowbreak{}xs, ys, cb, nu, k, xi) in l2-set(\allowbreak{}xs, ca, mu) and integral(\allowbreak{}xs, ca, mu, vnb-lambda(\allowbreak{}lx\_, xs, (\allowbreak{}kernel-apply(\allowbreak{}xs, ys, cb, nu, k, xi))(\allowbreak{}lx\_) * (\allowbreak{}kernel-apply(\allowbreak{}xs, ys, cb, nu, k, xi))(\allowbreak{}lx\_))) <= integral(\allowbreak{}cartesian(\allowbreak{}xs, ys), product-sigma(\allowbreak{}xs, ca, ys, cb), product-measure(\allowbreak{}xs, ca, mu, ys, cb, nu), vnb-lambda(\allowbreak{}lx\_, cartesian(\allowbreak{}xs, ys), k(\allowbreak{}lx\_) * k(\allowbreak{}lx\_))) * integral(\allowbreak{}ys, cb, nu, vnb-lambda(\allowbreak{}lx\_, ys, xi(\allowbreak{}lx\_) * xi(\allowbreak{}lx\_))) and forsome(\allowbreak{}[en\_], is-null-set(\allowbreak{}xs, ca, mu, en\_) and forall(\allowbreak{}[x in xs], not(\allowbreak{}x in en\_) implies is-integrable(\allowbreak{}ys, cb, nu, vnb-lambda(\allowbreak{}lx\_, ys, (\allowbreak{}fn-x-section(\allowbreak{}ys, k, x))(\allowbreak{}lx\_) * xi(\allowbreak{}lx\_))) and (\allowbreak{}kernel-apply(\allowbreak{}xs, ys, cb, nu, k, xi))(\allowbreak{}x) = rintegral(\allowbreak{}ys, cb, nu, vnb-lambda(\allowbreak{}lx\_, ys, (\allowbreak{}fn-x-section(\allowbreak{}ys, k, x))(\allowbreak{}lx\_) * xi(\allowbreak{}lx\_))))))
\end{quote}
also: \texttt{kernel-apply-ae-eq} \texttt{kernel-apply-unfold} \texttt{kernel-apply-cschwarz} \texttt{kernel-apply-linear-ae} \\
mentioned by 8 result(s)

\smallskip

\noindent\textbf{\texttt{key-ideal-gen(sbgt\_, sbgk\_, sbgks\_, sbgs\_)}}\quad\textit{functoid}\index{key-ideal-gen(sbgt_, sbgk_, sbgks_, sbgs_)@\texttt{key-ideal-gen(sbgt\_, sbgk\_, sbgks\_, sbgs\_)}}\\
reads: the \$3-ideal generated by \$4 in the messages over the tokens \$1 and the keys \$2 \\
definition:
\begin{quote}\ttfamily\footnotesize\raggedright\noindent
\{sbgm\_ in msg-set(\allowbreak{}sbgt\_, sbgk\_): forall(\allowbreak{}[sbgq\_ in power(\allowbreak{}msg-set(\allowbreak{}sbgt\_, sbgk\_))], is-key-ideal(\allowbreak{}sbgt\_, sbgk\_, sbgks\_, sbgq\_) implies forall(\allowbreak{}[sbgz\_ in sbgs\_], sbgz\_ in sbgq\_) implies sbgm\_ in sbgq\_)\}
\end{quote}
characterized by \texttt{key-ideal-gen-elim}  \textit{(proven)}:
\begin{quote}\ttfamily\footnotesize\raggedright\noindent
forall(\allowbreak{}[t, k, ks, s, m, ii], m in key-ideal-gen(\allowbreak{}t, k, ks, s) implies ii in power(\allowbreak{}msg-set(\allowbreak{}t, k)) implies is-key-ideal(\allowbreak{}t, k, ks, ii) implies forall(\allowbreak{}[sbgz\_ in s], sbgz\_ in ii) implies m in ii)
\end{quote}
also: \texttt{key-ideal-gen-self} \texttt{key-ideal-gen-cases} \texttt{key-ideal-gen-least} \texttt{key-ideal-gen-token} \texttt{key-ideal-gen-union} \texttt{key-ideal-gen-honest} \texttt{key-ideal-gen-stages} \texttt{key-ideal-gen-unfold} ... (20 in all) \\
mentioned by 29 result(s)

\smallskip

\noindent\textbf{\texttt{keys(s)}}\quad\textit{accessor}\index{keys(s)@\texttt{keys(s)}}\\
reads: the keys slot of \$1 \\
a structure slot \\
characterized by \texttt{cx-sp@keys}  \textit{(definitional)}:
\begin{quote}\ttfamily\footnotesize\raggedright\noindent
keys(\allowbreak{}cx-sp) == interval(\allowbreak{}30, 33)
\end{quote}
also: \texttt{keys-in-tokens} \\
mentioned by 119 result(s)

\smallskip

\noindent\textbf{\texttt{kinv(s)}}\quad\textit{accessor}\index{kinv(s)@\texttt{kinv(s)}}\\
reads: the kinv slot of \$1 \\
a structure slot \\
characterized by \texttt{cx-sp@kinv}  \textit{(definitional)}:
\begin{quote}\ttfamily\footnotesize\raggedright\noindent
kinv(\allowbreak{}cx-sp) == vnb-lambda(\allowbreak{}cxk\_, interval(\allowbreak{}30, 33), cxk\_)
\end{quote}
also: \texttt{kinv-in-fun} \texttt{kinv-injective} \texttt{kinv-involution} \\
mentioned by 43 result(s)

\smallskip

\noindent\textbf{\texttt{krylov(gkr-md\_, gkr-t\_, gkr-v\_, gkr-n\_)}}\quad\textit{functoid}\index{krylov(gkr-md_, gkr-t_, gkr-v_, gkr-n_)@\texttt{krylov(gkr-md\_, gkr-t\_, gkr-v\_, gkr-n\_)}}\\
reads: the vectors \$3, \$2 \$3, ..., \$2\textasciicircum{}(\$4 - 1) \$3 \\
definition:
\begin{quote}\ttfamily\footnotesize\raggedright\noindent
matof(\allowbreak{}gkr-n\_, 1, vnb-lambda(\allowbreak{}[gkri\_, gkrj\_], cartesian(\allowbreak{}interval(\allowbreak{}1, gkr-n\_), interval(\allowbreak{}1, 1)), op-iter(\allowbreak{}gkr-md\_, gkr-t\_, gkr-v\_, gkri\_ - 1)))
\end{quote}
characterized by \texttt{krylov-entry}  \textit{(proven)}:
\begin{quote}\ttfamily\footnotesize\raggedright\noindent
forall(\allowbreak{}[md, tt], is-hom-module(\allowbreak{}md, md, tt) implies forall(\allowbreak{}[n in nn, v in vec(\allowbreak{}md), i in interval(\allowbreak{}1, n)], entry(\allowbreak{}krylov(\allowbreak{}md, tt, v, n), i, 1) = op-iter(\allowbreak{}md, tt, v, i - 1)))
\end{quote}
also: \texttt{krylov-in-mat} \texttt{lincomb-krylov} \texttt{krylov-rel-free} \texttt{krylov-companion} \texttt{krylov-generates} \texttt{krylov-low-in-span} \texttt{krylov-top-in-span} \texttt{krylov-span-invariant} ... (10 in all) \\
mentioned by 17 result(s)

\smallskip

\noindent\textbf{\texttt{krylov-mat(k, n, tm, phi)}}\quad\textit{functoid}\index{krylov-mat(k, n, tm, phi)@\texttt{krylov-mat(k, n, tm, phi)}}\\
reads: the matrix with columns \$4, \$3 \$4, ..., \$3\textasciicircum{}(\$2 - 1) \$4 \\
definition:
\begin{quote}\ttfamily\footnotesize\raggedright\noindent
matof(\allowbreak{}n, n, vnb-lambda(\allowbreak{}[zkvi\_, zkvj\_], cartesian(\allowbreak{}interval(\allowbreak{}1, n), interval(\allowbreak{}1, n)), entry(\allowbreak{}matmul(\allowbreak{}normed-field-as-commutative-ring(\allowbreak{}k), mat-power(\allowbreak{}k, n, tm, zkvj\_ - 1), phi), zkvi\_, 1)))
\end{quote}
mentioned by 5 result(s)

\smallskip


\section*{L}
\addcontentsline{toc}{section}{L}

\noindent\textbf{\texttt{l1-pairing(omega, ca, mu, f)}}\quad\textit{functoid}\index{l1-pairing(omega, ca, mu, f)@\texttt{l1-pairing(omega, ca, mu, f)}}\\
reads: the functional g |-> int \$4 g d\$3 on L\textasciicircum{}1(\$1, \$3) \\
definition:
\begin{quote}\ttfamily\footnotesize\raggedright\noindent
vnb-lambda(\allowbreak{}lpg\_, l1-set(\allowbreak{}omega, ca, mu), rintegral(\allowbreak{}omega, ca, mu, vnb-lambda(\allowbreak{}sx\_, omega, f(\allowbreak{}sx\_) * lpg\_(\allowbreak{}sx\_))))
\end{quote}
characterized by \texttt{l1-pairing-value}  \textit{(proven)}:
\begin{quote}\ttfamily\footnotesize\raggedright\noindent
forall(\allowbreak{}[omega, ca, mu, f, g in l1-set(\allowbreak{}omega, ca, mu)], (\allowbreak{}l1-pairing(\allowbreak{}omega, ca, mu, f))(\allowbreak{}g) == rintegral(\allowbreak{}omega, ca, mu, vnb-lambda(\allowbreak{}sx\_, omega, f(\allowbreak{}sx\_) * g(\allowbreak{}sx\_))))
\end{quote}
also: \texttt{l1-pairing-bounded} \texttt{l1-pairing-norm-bound} \\
mentioned by 7 result(s)

\smallskip

\noindent\textbf{\texttt{l1-set(omega, ca, mu)}}\quad\textit{functoid}\index{l1-set(omega, ca, mu)@\texttt{l1-set(omega, ca, mu)}}\\
reads: the set of integrable real functions on \$1 with respect to \$3 \\
definition:
\begin{quote}\ttfamily\footnotesize\raggedright\noindent
\{l1g\_ in fun(\allowbreak{}omega, rr): is-integrable(\allowbreak{}omega, ca, mu, l1g\_)\}
\end{quote}
characterized by \texttt{l1-set-is-set}  \textit{(proven)}:
\begin{quote}\ttfamily\footnotesize\raggedright\noindent
forall(\allowbreak{}[omega, ca, mu], omega in set implies l1-set(\allowbreak{}omega, ca, mu) in set)
\end{quote}
also: \texttt{l1-set-membership} \\
mentioned by 17 result(s)

\smallskip

\noindent\textbf{\texttt{l2-act(omega, ca, mu)}}\quad\textit{functoid}\index{l2-act(omega, ca, mu)@\texttt{l2-act(omega, ca, mu)}}\\
definition:
\begin{quote}\ttfamily\footnotesize\raggedright\noindent
vnb-lambda(\allowbreak{}[lr\_, lf\_], cartesian(\allowbreak{}rr, l2-set(\allowbreak{}omega, ca, mu)), vnb-lambda(\allowbreak{}lx\_, omega, lr\_ * lf\_(\allowbreak{}lx\_)))
\end{quote}
characterized by \texttt{l2-act-value}  \textit{(proven)}:
\begin{quote}\ttfamily\footnotesize\raggedright\noindent
forall(\allowbreak{}[omega, ca, mu], is-measure(\allowbreak{}omega, ca, mu) implies forall(\allowbreak{}[lr\_, f], lr\_ in rr implies f in l2-set(\allowbreak{}omega, ca, mu) implies (\allowbreak{}l2-act(\allowbreak{}omega, ca, mu))(\allowbreak{}lr\_, f) = vnb-lambda(\allowbreak{}lx\_, omega, lr\_ * f(\allowbreak{}lx\_))))
\end{quote}
mentioned by 4 result(s)

\smallskip

\noindent\textbf{\texttt{l2-cips(ix)}}\quad\textit{functoid}\index{l2-cips(ix)@\texttt{l2-cips(ix)}}\\
reads: l2(\$1) as a complex inner product space \\
definition:
\begin{quote}\ttfamily\footnotesize\raggedright\noindent
[normed-field-as-commutative-ring(\allowbreak{}cc-normed-field), l2-families(\allowbreak{}ix), l2f-vadd(\allowbreak{}ix), l2f-vzero(\allowbreak{}ix), l2f-vneg(\allowbreak{}ix), l2f-act(\allowbreak{}ix), l2-ip(\allowbreak{}ix)]
\end{quote}
characterized by \texttt{l2-cips-ip}  \textit{(proven)}:
\begin{quote}\ttfamily\footnotesize\raggedright\noindent
forall(\allowbreak{}[ix in set], ip(\allowbreak{}l2-cips(\allowbreak{}ix)) == l2-ip(\allowbreak{}ix))
\end{quote}
also: \texttt{l2-cips-act} \texttt{l2-cips-vec} \texttt{l2-cips-scal} \texttt{l2-cips-vadd} \texttt{l2-cips-vneg} \texttt{l2-cips-vzero} \texttt{l2-cips-is-cips} \\
mentioned by 19 result(s)

\smallskip

\noindent\textbf{\texttt{l2-families(ix)}}\quad\textit{functoid}\index{l2-families(ix)@\texttt{l2-families(ix)}}\\
reads: the square-summable complex families on \$1 \\
definition:
\begin{quote}\ttfamily\footnotesize\raggedright\noindent
\{l2c\_ in fun(\allowbreak{}ix, cc): esum(\allowbreak{}vnb-lambda(\allowbreak{}l2k\_, ix, magnitude(\allowbreak{}l2c\_(\allowbreak{}l2k\_)) * magnitude(\allowbreak{}l2c\_(\allowbreak{}l2k\_)))) in rr\}
\end{quote}
characterized by \texttt{l2-families-ext}  \textit{(proven)}:
\begin{quote}\ttfamily\footnotesize\raggedright\noindent
forall(\allowbreak{}[ix, a, b], a in l2-families(\allowbreak{}ix) implies b in l2-families(\allowbreak{}ix) implies forall(\allowbreak{}[l2e\_ in ix], a(\allowbreak{}l2e\_) = b(\allowbreak{}l2e\_)) implies a = b)
\end{quote}
also: \texttt{l2-families-elim} \texttt{l2-families-intro} \texttt{l2-families-is-set} \texttt{l2-families-unfold} \texttt{l2-families-act-closed} \texttt{l2-families-add-closed} \texttt{l2-families-neg-closed} \\
mentioned by 45 result(s)

\smallskip

\noindent\textbf{\texttt{l2-ip(ix)}}\quad\textit{functoid}\index{l2-ip(ix)@\texttt{l2-ip(ix)}}\\
reads: the inner product of l2(\$1) \\
definition:
\begin{quote}\ttfamily\footnotesize\raggedright\noindent
vnb-lambda(\allowbreak{}[l2a\_, l2b\_], cartesian(\allowbreak{}l2-families(\allowbreak{}ix), l2-families(\allowbreak{}ix)), cc-family-sum(\allowbreak{}vnb-lambda(\allowbreak{}l2k\_, ix, l2a\_(\allowbreak{}l2k\_) * conjugate(\allowbreak{}l2b\_(\allowbreak{}l2k\_)))))
\end{quote}
characterized by \texttt{l2-ip-fun}  \textit{(proven)}:
\begin{quote}\ttfamily\footnotesize\raggedright\noindent
forall(\allowbreak{}[ix in set], l2-ip(\allowbreak{}ix) in fun(\allowbreak{}cartesian(\allowbreak{}l2-families(\allowbreak{}ix), l2-families(\allowbreak{}ix)), cc))
\end{quote}
also: \texttt{l2-ip-self} \texttt{l2-ip-apply} \texttt{l2-ip-act-left} \texttt{l2-ip-add-left} \texttt{l2-ip-conj-sym} \texttt{l2-ip-self-zero} \\
mentioned by 16 result(s)

\smallskip

\noindent\textbf{\texttt{l2-rip(omega, ca, mu)}}\quad\textit{functoid}\index{l2-rip(omega, ca, mu)@\texttt{l2-rip(omega, ca, mu)}}\\
definition:
\begin{quote}\ttfamily\footnotesize\raggedright\noindent
vnb-lambda(\allowbreak{}[lf\_, lg\_], cartesian(\allowbreak{}l2-set(\allowbreak{}omega, ca, mu), l2-set(\allowbreak{}omega, ca, mu)), rintegral(\allowbreak{}omega, ca, mu, vnb-lambda(\allowbreak{}lx\_, omega, lf\_(\allowbreak{}lx\_) * lg\_(\allowbreak{}lx\_))))
\end{quote}
characterized by \texttt{l2-rip-value}  \textit{(proven)}:
\begin{quote}\ttfamily\footnotesize\raggedright\noindent
forall(\allowbreak{}[omega, ca, mu], is-measure(\allowbreak{}omega, ca, mu) implies forall(\allowbreak{}[f, g], f in l2-set(\allowbreak{}omega, ca, mu) implies g in l2-set(\allowbreak{}omega, ca, mu) implies (\allowbreak{}l2-rip(\allowbreak{}omega, ca, mu))(\allowbreak{}f, g) = rintegral(\allowbreak{}omega, ca, mu, vnb-lambda(\allowbreak{}lx\_, omega, f(\allowbreak{}lx\_) * g(\allowbreak{}lx\_)))))
\end{quote}
mentioned by 4 result(s)

\smallskip

\noindent\textbf{\texttt{l2-set(omega, ca, mu)}}\quad\textit{functoid}\index{l2-set(omega, ca, mu)@\texttt{l2-set(omega, ca, mu)}}\\
reads: the square-integrable real functions on \$1 with respect to \$3 \\
definition:
\begin{quote}\ttfamily\footnotesize\raggedright\noindent
\{lf\_ in fun(\allowbreak{}omega, rr): is-real-measurable-fn(\allowbreak{}omega, ca, lf\_) and is-integrable(\allowbreak{}omega, ca, mu, vnb-lambda(\allowbreak{}lx\_, omega, lf\_(\allowbreak{}lx\_) * lf\_(\allowbreak{}lx\_)))\}
\end{quote}
characterized by \texttt{l2-set-is-set}  \textit{(proven)}:
\begin{quote}\ttfamily\footnotesize\raggedright\noindent
forall(\allowbreak{}[omega, ca, mu], is-measure(\allowbreak{}omega, ca, mu) implies l2-set(\allowbreak{}omega, ca, mu) in set)
\end{quote}
also: \texttt{l2-set-subset-fun} \texttt{l2-setoid-is-ae-setoid} \\
mentioned by 74 result(s)

\smallskip

\noindent\textbf{\texttt{l2-space(omega, ca, mu)}}\quad\textit{functoid}\index{l2-space(omega, ca, mu)@\texttt{l2-space(omega, ca, mu)}}\\
reads: the real L\textasciicircum{}2 space of \$3 on \$1 \\
definition:
\begin{quote}\ttfamily\footnotesize\raggedright\noindent
[normed-field-as-commutative-ring(\allowbreak{}rr-normed-field), l2-set(\allowbreak{}omega, ca, mu), l2-vadd(\allowbreak{}omega, ca, mu), l2-vzero(\allowbreak{}omega), l2-vneg(\allowbreak{}omega, ca, mu), l2-act(\allowbreak{}omega, ca, mu), l2-rip(\allowbreak{}omega, ca, mu)]
\end{quote}
characterized by \texttt{l2-space-act}  \textit{(proven)}:
\begin{quote}\ttfamily\footnotesize\raggedright\noindent
forall(\allowbreak{}[omega, ca, mu], is-measure(\allowbreak{}omega, ca, mu) implies act(\allowbreak{}l2-space(\allowbreak{}omega, ca, mu)) == l2-act(\allowbreak{}omega, ca, mu))
\end{quote}
also: \texttt{l2-space-rip} \texttt{l2-space-vec} \texttt{l2-space-comm} \texttt{l2-space-scal} \texttt{l2-space-vadd} \texttt{l2-space-vneg} \texttt{l2-space-assoc} \texttt{l2-space-vzero} ... (29 in all) \\
mentioned by 84 result(s)

\smallskip

\noindent\textbf{\texttt{l2-vadd(omega, ca, mu)}}\quad\textit{functoid}\index{l2-vadd(omega, ca, mu)@\texttt{l2-vadd(omega, ca, mu)}}\\
definition:
\begin{quote}\ttfamily\footnotesize\raggedright\noindent
vnb-lambda(\allowbreak{}[lf\_, lg\_], cartesian(\allowbreak{}l2-set(\allowbreak{}omega, ca, mu), l2-set(\allowbreak{}omega, ca, mu)), vnb-lambda(\allowbreak{}lx\_, omega, lf\_(\allowbreak{}lx\_) + lg\_(\allowbreak{}lx\_)))
\end{quote}
characterized by \texttt{l2-vadd-value}  \textit{(proven)}:
\begin{quote}\ttfamily\footnotesize\raggedright\noindent
forall(\allowbreak{}[omega, ca, mu], is-measure(\allowbreak{}omega, ca, mu) implies forall(\allowbreak{}[f, g], f in l2-set(\allowbreak{}omega, ca, mu) implies g in l2-set(\allowbreak{}omega, ca, mu) implies (\allowbreak{}l2-vadd(\allowbreak{}omega, ca, mu))(\allowbreak{}f, g) = vnb-lambda(\allowbreak{}lx\_, omega, f(\allowbreak{}lx\_) + g(\allowbreak{}lx\_))))
\end{quote}
mentioned by 4 result(s)

\smallskip

\noindent\textbf{\texttt{l2-vneg(omega, ca, mu)}}\quad\textit{functoid}\index{l2-vneg(omega, ca, mu)@\texttt{l2-vneg(omega, ca, mu)}}\\
definition:
\begin{quote}\ttfamily\footnotesize\raggedright\noindent
vnb-lambda(\allowbreak{}lf\_, l2-set(\allowbreak{}omega, ca, mu), vnb-lambda(\allowbreak{}lx\_, omega, -lf\_(\allowbreak{}lx\_)))
\end{quote}
characterized by \texttt{l2-vneg-value}  \textit{(proven)}:
\begin{quote}\ttfamily\footnotesize\raggedright\noindent
forall(\allowbreak{}[omega, ca, mu], is-measure(\allowbreak{}omega, ca, mu) implies forall(\allowbreak{}[f in l2-set(\allowbreak{}omega, ca, mu)], (\allowbreak{}l2-vneg(\allowbreak{}omega, ca, mu))(\allowbreak{}f) = vnb-lambda(\allowbreak{}lx\_, omega, -f(\allowbreak{}lx\_))))
\end{quote}
mentioned by 4 result(s)

\smallskip

\noindent\textbf{\texttt{l2-vzero(omega)}}\quad\textit{functoid}\index{l2-vzero(omega)@\texttt{l2-vzero(omega)}}\\
definition:
\begin{quote}\ttfamily\footnotesize\raggedright\noindent
vnb-lambda(\allowbreak{}lx\_, omega, 0)
\end{quote}
mentioned by 2 result(s)

\smallskip

\noindent\textbf{\texttt{l2f-act(ix)}}\quad\textit{functoid}\index{l2f-act(ix)@\texttt{l2f-act(ix)}}\\
definition:
\begin{quote}\ttfamily\footnotesize\raggedright\noindent
vnb-lambda(\allowbreak{}[l2z\_, l2a\_], cartesian(\allowbreak{}cc, l2-families(\allowbreak{}ix)), vnb-lambda(\allowbreak{}l2k\_, ix, l2z\_ * l2a\_(\allowbreak{}l2k\_)))
\end{quote}
mentioned by 8 result(s)

\smallskip

\noindent\textbf{\texttt{l2f-vadd(ix)}}\quad\textit{functoid}\index{l2f-vadd(ix)@\texttt{l2f-vadd(ix)}}\\
definition:
\begin{quote}\ttfamily\footnotesize\raggedright\noindent
vnb-lambda(\allowbreak{}[l2a\_, l2b\_], cartesian(\allowbreak{}l2-families(\allowbreak{}ix), l2-families(\allowbreak{}ix)), vnb-lambda(\allowbreak{}l2k\_, ix, l2a\_(\allowbreak{}l2k\_) + l2b\_(\allowbreak{}l2k\_)))
\end{quote}
mentioned by 8 result(s)

\smallskip

\noindent\textbf{\texttt{l2f-vneg(ix)}}\quad\textit{functoid}\index{l2f-vneg(ix)@\texttt{l2f-vneg(ix)}}\\
definition:
\begin{quote}\ttfamily\footnotesize\raggedright\noindent
vnb-lambda(\allowbreak{}l2a\_, l2-families(\allowbreak{}ix), vnb-lambda(\allowbreak{}l2k\_, ix, -l2a\_(\allowbreak{}l2k\_)))
\end{quote}
mentioned by 6 result(s)

\smallskip

\noindent\textbf{\texttt{l2f-vzero(ix)}}\quad\textit{functoid}\index{l2f-vzero(ix)@\texttt{l2f-vzero(ix)}}\\
definition:
\begin{quote}\ttfamily\footnotesize\raggedright\noindent
vnb-lambda(\allowbreak{}l2k\_, ix, 0)
\end{quote}
mentioned by 7 result(s)

\smallskip

\noindent\textbf{\texttt{lam-b}}\quad\textit{tactic}\index{lam-b@\texttt{lam-b}}\\
\texttt{(lam-b)} -- VNB-LAMBDA beta: reduce an applied lambda to its substituted body. \\
declared in \texttt{tactics-help.scm}

\smallskip

\noindent\textbf{\texttt{lam-b-h}}\quad\textit{tactic}\index{lam-b-h@\texttt{lam-b-h}}\\
\texttt{(lam-b-h hyp)} -- VNB-LAMBDA beta in a cited ASSUMPTION -- what mac-h is to mac. \\
declared in \texttt{tactics-help.scm}

\smallskip

\noindent\textbf{\texttt{lam-t}}\quad\textit{tactic}\index{lam-t@\texttt{lam-t}}\\
\texttt{(lam-t)} -- VNB-LAMBDA typing: reduce (IN (VNB-LAMBDA ...) (FUN A B)) to TWO obligations -- the body, and that the domain is a set. \\
declared in \texttt{tactics-help.scm}

\smallskip

\noindent\textbf{\texttt{lastcoeff-set(md, p, u, sm)}}\quad\textit{functoid}\index{lastcoeff-set(md, p, u, sm)@\texttt{lastcoeff-set(md, p, u, sm)}}\\
definition:
\begin{quote}\ttfamily\footnotesize\raggedright\noindent
\{r\_ in carr(\allowbreak{}scal(\allowbreak{}md)): forsome(\allowbreak{}[c\_ in mat(\allowbreak{}1, succ(\allowbreak{}p), carr(\allowbreak{}scal(\allowbreak{}md)))], entry(\allowbreak{}c\_, 1, succ(\allowbreak{}p)) = r\_ and lincomb(\allowbreak{}md, succ(\allowbreak{}p), c\_, u) in sm)\}
\end{quote}
characterized by \texttt{lastcoeff-set-membership}  \textit{(definitional)}:
\begin{quote}\ttfamily\footnotesize\raggedright\noindent
forall(\allowbreak{}[md, p, u, sm, r\_], r\_ in lastcoeff-set(\allowbreak{}md, p, u, sm) iff r\_ in carr(\allowbreak{}scal(\allowbreak{}md)) and forsome(\allowbreak{}[c\_ in mat(\allowbreak{}1, succ(\allowbreak{}p), carr(\allowbreak{}scal(\allowbreak{}md)))], entry(\allowbreak{}c\_, 1, succ(\allowbreak{}p)) = r\_ and lincomb(\allowbreak{}md, succ(\allowbreak{}p), c\_, u) in sm))
\end{quote}
also: \texttt{lastcoeff-set-is-ideal-cna-as-module} \texttt{lastcoeff-set-is-ideal-sip-as-module} \texttt{lastcoeff-set-is-ideal-uba-as-module} \texttt{lastcoeff-set-is-ideal-cnvs-as-module} \texttt{lastcoeff-set-is-ideal-sip-as-module-csip-as-sip} \texttt{lastcoeff-set-is-ideal-uba-as-module-uiba-as-uba} \texttt{lastcoeff-set-is-ideal} \\
mentioned by 9 result(s)

\smallskip

\noindent\textbf{\texttt{laurent-abs(lrk\_, lrr\_, lrf\_)}}\quad\textit{functoid}\index{laurent-abs(lrk_, lrr_, lrf_)@\texttt{laurent-abs(lrk\_, lrr\_, lrf\_)}}\\
reads: the absolute value \$2\textasciicircum{}ord(\$3) of the Laurent series \$3 \\
definition:
\begin{quote}\ttfamily\footnotesize\raggedright\noindent
if(\allowbreak{}lrf\_ = laurent-zero(\allowbreak{}lrk\_), 0, rpow-star(\allowbreak{}lrr\_, laurent-ord(\allowbreak{}lrk\_, lrf\_)))
\end{quote}
characterized by \texttt{laurent-abs-mul}  \textit{(proven)}:
\begin{quote}\ttfamily\footnotesize\raggedright\noindent
forall(\allowbreak{}[kq\_], is-integral-domain(\allowbreak{}kq\_) implies forall(\allowbreak{}[rq\_ in rr], 0 < rq\_ implies forall(\allowbreak{}[fq\_ in laurent-set(\allowbreak{}kq\_), gq\_ in laurent-set(\allowbreak{}kq\_)], laurent-abs(\allowbreak{}kq\_, rq\_, laurent-mul(\allowbreak{}kq\_, fq\_, gq\_)) = laurent-abs(\allowbreak{}kq\_, rq\_, fq\_) * laurent-abs(\allowbreak{}kq\_, rq\_, gq\_))))
\end{quote}
also: \texttt{laurent-abs-zero} \texttt{laurent-abs-in-rr} \texttt{laurent-abs-ultra} \texttt{laurent-abs-le-iff} \texttt{laurent-abs-lt-one} \texttt{laurent-abs-nonzero} \texttt{laurent-abs-zero-iff} \\
mentioned by 22 result(s)

\smallskip

\noindent\textbf{\texttt{laurent-add(lrk\_, lrf\_, lrg\_)}}\quad\textit{functoid}\index{laurent-add(lrk_, lrf_, lrg_)@\texttt{laurent-add(lrk\_, lrf\_, lrg\_)}}\\
reads: the sum of the Laurent series \$2 and \$3 \\
definition:
\begin{quote}\ttfamily\footnotesize\raggedright\noindent
vnb-lambda(\allowbreak{}lrj\_, zz, (\allowbreak{}add(\allowbreak{}lrk\_))(\allowbreak{}lrf\_(\allowbreak{}lrj\_), lrg\_(\allowbreak{}lrj\_)))
\end{quote}
characterized by \texttt{laurent-add-in}  \textit{(proven)}:
\begin{quote}\ttfamily\footnotesize\raggedright\noindent
forall(\allowbreak{}[lzk], is-ring(\allowbreak{}lzk) implies forall(\allowbreak{}[lzf in laurent-set(\allowbreak{}lzk), lzg in laurent-set(\allowbreak{}lzk)], laurent-add(\allowbreak{}lzk, lzf, lzg) in laurent-set(\allowbreak{}lzk)))
\end{quote}
also: \texttt{laurent-add-comm} \texttt{laurent-add-apply} \texttt{laurent-add-assoc} \texttt{laurent-add-below} \texttt{laurent-add-monalg} \texttt{laurent-add-neg-left} \texttt{laurent-add-zero-left} \\
mentioned by 33 result(s)

\smallskip

\noindent\textbf{\texttt{laurent-coefs(lcf\_, lcc\_, lcr\_)}}\quad\textit{functoid}\index{laurent-coefs(lcf_, lcc_, lcr_)@\texttt{laurent-coefs(lcf\_, lcc\_, lcr\_)}}\\
reads: the Laurent coefficients of \$1 about \$2 read on the circle of radius \$3 \\
definition:
\begin{quote}\ttfamily\footnotesize\raggedright\noindent
vnb-lambda(\allowbreak{}lck\_, zz, if(\allowbreak{}0 <= lck\_, recip(\allowbreak{}(\allowbreak{}2 * pi) * 1i) * line-int(\allowbreak{}vnb-lambda(\allowbreak{}lcw\_, trace(\allowbreak{}circle-path(\allowbreak{}lcc\_, lcr\_), 0, 2 * pi), lcf\_(\allowbreak{}lcw\_) * recip(\allowbreak{}lcw\_ - lcc\_) \textasciicircum{} succ(\allowbreak{}lck\_)), circle-path(\allowbreak{}lcc\_, lcr\_), circle-deriv(\allowbreak{}lcc\_, lcr\_), 0, 2 * pi), recip(\allowbreak{}(\allowbreak{}2 * pi) * 1i) * line-int(\allowbreak{}vnb-lambda(\allowbreak{}lcw\_, trace(\allowbreak{}circle-path(\allowbreak{}lcc\_, lcr\_), 0, 2 * pi), lcf\_(\allowbreak{}lcw\_) * (\allowbreak{}lcw\_ - lcc\_) \textasciicircum{} (\allowbreak{}-(\allowbreak{}lck\_ + 1))), circle-path(\allowbreak{}lcc\_, lcr\_), circle-deriv(\allowbreak{}lcc\_, lcr\_), 0, 2 * pi)))
\end{quote}
characterized by \texttt{laurent-coefs-radii}  \textit{(proven)}:
\begin{quote}\ttfamily\footnotesize\raggedright\noindent
forall(\allowbreak{}[u, f, c in cc, r in rr, rho in rr], 0 < rho implies rho < r implies is-open(\allowbreak{}nf-metric-space(\allowbreak{}cc-normed-field), u) implies holomorphic-on(\allowbreak{}u, f) implies forall(\allowbreak{}[lxb\_ in cc], 0 < magnitude(\allowbreak{}lxb\_ - c) implies magnitude(\allowbreak{}lxb\_ - c) < r implies lxb\_ in u) implies forall(\allowbreak{}[lrp\_ in rr], lrp\_ < r implies lrp\_ < cps-radius(\allowbreak{}vnb-lambda(\allowbreak{}lak\_, nn, (\allowbreak{}laurent-coefs(\allowbreak{}f, c, rho))(\allowbreak{}lak\_)))) and forall(\allowbreak{}[lrm\_ in rr], 0 < lrm\_ implies recip(\allowbreak{}lrm\_) < cps-radius(\allowbreak{}vnb-lambda(\allowbreak{}lak\_, nn, (\allowbreak{}laurent-coefs(\allowbreak{}f, c, rho))(\allowbreak{}-(\allowbreak{}lak\_ + 1))))))
\end{quote}
also: \texttt{laurent-coefs-in-fun} \\
mentioned by 11 result(s)

\smallskip

\noindent\textbf{\texttt{laurent-ext(llk\_, llm\_)}}\quad\textit{functoid}\index{laurent-ext(llk_, llm_)@\texttt{laurent-ext(llk\_, llm\_)}}\\
reads: the map (p, c) |-> p + c x\textasciicircum{}\$2 on the polynomials of degree below \$2 \\
definition:
\begin{quote}\ttfamily\footnotesize\raggedright\noindent
vnb-lambda(\allowbreak{}llw\_, cartesian(\allowbreak{}laurent-polys(\allowbreak{}llk\_, llm\_), carr(\allowbreak{}llk\_)), laurent-add(\allowbreak{}llk\_, nth(\allowbreak{}1, llw\_), laurent-mono(\allowbreak{}llk\_, nth(\allowbreak{}2, llw\_), llm\_)))
\end{quote}

\smallskip

\noindent\textbf{\texttt{laurent-inv-pre}}\quad\textit{functoid}\index{laurent-inv-pre@\texttt{laurent-inv-pre}}\\
reads: the prefix of length \$5 of the inverse of \$2 \\
defined by \texttt{laurent-inv-pre-zero} \texttt{laurent-inv-pre-succ} \\
characterized by \texttt{laurent-inv-pre-succ}  \textit{(definitional)}:
\begin{quote}\ttfamily\footnotesize\raggedright\noindent
forall(\allowbreak{}[lrk\_, lrf\_, lro\_, lra\_, lrm\_ in nn], laurent-inv-pre(\allowbreak{}lrk\_, lrf\_, lro\_, lra\_, succ(\allowbreak{}lrm\_)) == laurent-add(\allowbreak{}lrk\_, laurent-inv-pre(\allowbreak{}lrk\_, lrf\_, lro\_, lra\_, lrm\_), laurent-mono(\allowbreak{}lrk\_, (\allowbreak{}neg(\allowbreak{}lrk\_))(\allowbreak{}(\allowbreak{}mul(\allowbreak{}lrk\_))(\allowbreak{}lra\_, (\allowbreak{}laurent-mul(\allowbreak{}lrk\_, lrf\_, laurent-inv-pre(\allowbreak{}lrk\_, lrf\_, lro\_, lra\_, lrm\_)))(\allowbreak{}lrm\_ + 1))), (\allowbreak{}lrm\_ + 1) + -lro\_)))
\end{quote}
also: \texttt{laurent-inv-pre-zero} \texttt{laurent-inv-pre-in} \texttt{laurent-inv-pre-step} \texttt{laurent-inv-pre-unit} \texttt{laurent-inv-pre-stable} \\
mentioned by 15 result(s)

\smallskip

\noindent\textbf{\texttt{laurent-inv-ser(lrk\_, lrf\_, lro\_, lra\_)}}\quad\textit{functoid}\index{laurent-inv-ser(lrk_, lrf_, lro_, lra_)@\texttt{laurent-inv-ser(lrk\_, lrf\_, lro\_, lra\_)}}\\
reads: the inverse series of \$2 \\
definition:
\begin{quote}\ttfamily\footnotesize\raggedright\noindent
vnb-lambda(\allowbreak{}lrj\_, zz, if(\allowbreak{}-lro\_ <= lrj\_, (\allowbreak{}laurent-inv-pre(\allowbreak{}lrk\_, lrf\_, lro\_, lra\_, lrj\_ + lro\_))(\allowbreak{}lrj\_), zero(\allowbreak{}lrk\_)))
\end{quote}
characterized by \texttt{laurent-inv-ser-in}  \textit{(proven)}:
\begin{quote}\ttfamily\footnotesize\raggedright\noindent
forall(\allowbreak{}[kq\_], is-ring(\allowbreak{}kq\_) implies forall(\allowbreak{}[fq\_ in laurent-set(\allowbreak{}kq\_), oq\_ in zz, aq\_ in carr(\allowbreak{}kq\_)], laurent-inv-ser(\allowbreak{}kq\_, fq\_, oq\_, aq\_) in laurent-set(\allowbreak{}kq\_)))
\end{quote}
also: \texttt{laurent-inv-ser-fun} \texttt{laurent-inv-ser-low} \texttt{laurent-inv-ser-agree} \texttt{laurent-inv-ser-apply} \texttt{laurent-inv-ser-right} \\
mentioned by 10 result(s)

\smallskip

\noindent\textbf{\texttt{laurent-lim(llk\_, lls\_)}}\quad\textit{functoid}\index{laurent-lim(llk_, lls_)@\texttt{laurent-lim(llk\_, lls\_)}}\\
reads: the coefficientwise limit of the sequence \$2 of Laurent series \\
definition:
\begin{quote}\ttfamily\footnotesize\raggedright\noindent
vnb-lambda(\allowbreak{}llj\_, zz, laurent-limc(\allowbreak{}llk\_, lls\_, llj\_))
\end{quote}
characterized by \texttt{laurent-lim-in}  \textit{(proven)}:
\begin{quote}\ttfamily\footnotesize\raggedright\noindent
forall(\allowbreak{}[kq\_, sq\_ in fun(\allowbreak{}nn, laurent-set(\allowbreak{}kq\_))], forall(\allowbreak{}[nq\_ in zz], forsome(\allowbreak{}[mq\_ in nn], forall(\allowbreak{}[iq\_ in nn], mq\_ <= iq\_ implies forall(\allowbreak{}[ir\_ in nn], mq\_ <= ir\_ implies forall(\allowbreak{}[lzj in zz], lzj < nq\_ implies (\allowbreak{}sq\_(\allowbreak{}iq\_))(\allowbreak{}lzj) = (\allowbreak{}sq\_(\allowbreak{}ir\_))(\allowbreak{}lzj)))))) implies laurent-lim(\allowbreak{}kq\_, sq\_) in laurent-set(\allowbreak{}kq\_))
\end{quote}
also: \texttt{laurent-lim-near} \texttt{laurent-lim-apply} \\
mentioned by 4 result(s)

\smallskip

\noindent\textbf{\texttt{laurent-limc(llk\_, lls\_, llj\_)}}\quad\textit{functoid}\index{laurent-limc(llk_, lls_, llj_)@\texttt{laurent-limc(llk\_, lls\_, llj\_)}}\\
reads: the eventual value of the coefficients at \$3 of the sequence \$2 \\
definition:
\begin{quote}\ttfamily\footnotesize\raggedright\noindent
iota(\allowbreak{}llc\_, llc\_ in carr(\allowbreak{}llk\_) and forsome(\allowbreak{}[llm\_ in nn], forall(\allowbreak{}[lln\_ in nn], llm\_ <= lln\_ implies (\allowbreak{}lls\_(\allowbreak{}lln\_))(\allowbreak{}llj\_) = llc\_)))
\end{quote}
characterized by \texttt{laurent-limc-spec}  \textit{(proven)}:
\begin{quote}\ttfamily\footnotesize\raggedright\noindent
forall(\allowbreak{}[kq\_, sq\_ in fun(\allowbreak{}nn, laurent-set(\allowbreak{}kq\_)), jq\_ in zz, mq\_ in nn], forall(\allowbreak{}[iq\_ in nn], mq\_ <= iq\_ implies (\allowbreak{}sq\_(\allowbreak{}iq\_))(\allowbreak{}jq\_) = (\allowbreak{}sq\_(\allowbreak{}mq\_))(\allowbreak{}jq\_)) implies laurent-limc(\allowbreak{}kq\_, sq\_, jq\_) in carr(\allowbreak{}kq\_) and forall(\allowbreak{}[iq\_ in nn], mq\_ <= iq\_ implies (\allowbreak{}sq\_(\allowbreak{}iq\_))(\allowbreak{}jq\_) = laurent-limc(\allowbreak{}kq\_, sq\_, jq\_)))
\end{quote}
also: \texttt{metric-group-group@laurent-limc} \\
mentioned by 5 result(s)

\smallskip

\noindent\textbf{\texttt{laurent-mono(lrk\_, lrc\_, lrt\_)}}\quad\textit{functoid}\index{laurent-mono(lrk_, lrc_, lrt_)@\texttt{laurent-mono(lrk\_, lrc\_, lrt\_)}}\\
reads: the monomial \$2 x\textasciicircum{}\$3 \\
definition:
\begin{quote}\ttfamily\footnotesize\raggedright\noindent
vnb-lambda(\allowbreak{}lrj\_, zz, if(\allowbreak{}lrj\_ = lrt\_, lrc\_, zero(\allowbreak{}lrk\_)))
\end{quote}
characterized by \texttt{laurent-mono-at}  \textit{(proven)}:
\begin{quote}\ttfamily\footnotesize\raggedright\noindent
forall(\allowbreak{}[kq\_], is-ring(\allowbreak{}kq\_) implies forall(\allowbreak{}[cq\_ in carr(\allowbreak{}kq\_), tq\_ in zz], (\allowbreak{}laurent-mono(\allowbreak{}kq\_, cq\_, tq\_))(\allowbreak{}tq\_) = cq\_))
\end{quote}
also: \texttt{laurent-mono-in} \texttt{laurent-mono-fun} \texttt{laurent-mono-off} \texttt{laurent-mono-apply} \texttt{laurent-mono-below} \\
mentioned by 20 result(s)

\smallskip

\noindent\textbf{\texttt{laurent-mul(lrk\_, lrf\_, lrg\_)}}\quad\textit{functoid}\index{laurent-mul(lrk_, lrf_, lrg_)@\texttt{laurent-mul(lrk\_, lrf\_, lrg\_)}}\\
reads: the product of the Laurent series \$2 and \$3 \\
definition:
\begin{quote}\ttfamily\footnotesize\raggedright\noindent
monalg-mul(\allowbreak{}lrk\_, zz-add-monoid, lrf\_, lrg\_)
\end{quote}
characterized by \texttt{laurent-mul-in}  \textit{(proven)}:
\begin{quote}\ttfamily\footnotesize\raggedright\noindent
forall(\allowbreak{}[kq\_], is-ring(\allowbreak{}kq\_) implies forall(\allowbreak{}[fq\_ in laurent-set(\allowbreak{}kq\_), gq\_ in laurent-set(\allowbreak{}kq\_)], laurent-mul(\allowbreak{}kq\_, fq\_, gq\_) in laurent-set(\allowbreak{}kq\_)))
\end{quote}
also: \texttt{laurent-mul-one} \texttt{laurent-mul-comm} \texttt{laurent-mul-assoc} \texttt{laurent-mul-unfold} \texttt{laurent-mul-mono-low} \texttt{laurent-mul-mono-top} \texttt{laurent-mul-zero-left} \texttt{laurent-mul-distrib-left} \\
mentioned by 36 result(s)

\smallskip

\noindent\textbf{\texttt{laurent-neg(lrk\_, lrf\_)}}\quad\textit{functoid}\index{laurent-neg(lrk_, lrf_)@\texttt{laurent-neg(lrk\_, lrf\_)}}\\
reads: the negative of the Laurent series \$2 \\
definition:
\begin{quote}\ttfamily\footnotesize\raggedright\noindent
vnb-lambda(\allowbreak{}lrj\_, zz, (\allowbreak{}neg(\allowbreak{}lrk\_))(\allowbreak{}lrf\_(\allowbreak{}lrj\_)))
\end{quote}
characterized by \texttt{laurent-neg-in}  \textit{(proven)}:
\begin{quote}\ttfamily\footnotesize\raggedright\noindent
forall(\allowbreak{}[lzk], is-ring(\allowbreak{}lzk) implies forall(\allowbreak{}[lzf in laurent-set(\allowbreak{}lzk)], laurent-neg(\allowbreak{}lzk, lzf) in laurent-set(\allowbreak{}lzk)))
\end{quote}
also: \texttt{laurent-neg-apply} \\
mentioned by 9 result(s)

\smallskip

\noindent\textbf{\texttt{laurent-nf(lrk\_, lrr\_)}}\quad\textit{functoid}\index{laurent-nf(lrk_, lrr_)@\texttt{laurent-nf(lrk\_, lrr\_)}}\\
reads: the normed field \$1((x)) with |f| = \$2\textasciicircum{}ord(f) \\
definition:
\begin{quote}\ttfamily\footnotesize\raggedright\noindent
[carr(\allowbreak{}laurent-ring(\allowbreak{}lrk\_)), add(\allowbreak{}laurent-ring(\allowbreak{}lrk\_)), mul(\allowbreak{}laurent-ring(\allowbreak{}lrk\_)), neg(\allowbreak{}laurent-ring(\allowbreak{}lrk\_)), zero(\allowbreak{}laurent-ring(\allowbreak{}lrk\_)), one(\allowbreak{}laurent-ring(\allowbreak{}lrk\_)), vnb-lambda(\allowbreak{}lrf\_, carr(\allowbreak{}laurent-ring(\allowbreak{}lrk\_)), laurent-abs(\allowbreak{}lrk\_, lrr\_, lrf\_))]
\end{quote}
characterized by \texttt{laurent-nf-add}  \textit{(proven)}:
\begin{quote}\ttfamily\footnotesize\raggedright\noindent
forall(\allowbreak{}[kq\_], is-commutative-ring(\allowbreak{}kq\_) implies forall(\allowbreak{}[rq\_ in rr], 0 < rq\_ implies add(\allowbreak{}laurent-nf(\allowbreak{}kq\_, rq\_)) == add(\allowbreak{}laurent-ring(\allowbreak{}kq\_))))
\end{quote}
also: \texttt{laurent-nf-mul} \texttt{laurent-nf-neg} \texttt{laurent-nf-one} \texttt{laurent-nf-pts} \texttt{laurent-nf-carr} \texttt{laurent-nf-dist} \texttt{laurent-nf-fnrm} \texttt{laurent-nf-zero} ... (12 in all) \\
mentioned by 54 result(s)

\smallskip

\noindent\textbf{\texttt{laurent-one(lrk\_)}}\quad\textit{functoid}\index{laurent-one(lrk_)@\texttt{laurent-one(lrk\_)}}\\
reads: the unit Laurent series over \$1 \\
definition:
\begin{quote}\ttfamily\footnotesize\raggedright\noindent
vnb-lambda(\allowbreak{}lrj\_, zz, if(\allowbreak{}lrj\_ = 0, one(\allowbreak{}lrk\_), zero(\allowbreak{}lrk\_)))
\end{quote}
characterized by \texttt{laurent-one-in}  \textit{(proven)}:
\begin{quote}\ttfamily\footnotesize\raggedright\noindent
forall(\allowbreak{}[lzk], is-ring(\allowbreak{}lzk) implies laurent-one(\allowbreak{}lzk) in laurent-set(\allowbreak{}lzk))
\end{quote}
also: \texttt{laurent-one-low} \texttt{laurent-one-off} \texttt{laurent-one-zero} \texttt{laurent-one-apply} \texttt{laurent-one-monalg} \\
mentioned by 21 result(s)

\smallskip

\noindent\textbf{\texttt{laurent-ord(lrk\_, lrf\_)}}\quad\textit{functoid}\index{laurent-ord(lrk_, lrf_)@\texttt{laurent-ord(lrk\_, lrf\_)}}\\
reads: the order of the Laurent series \$2 \\
definition:
\begin{quote}\ttfamily\footnotesize\raggedright\noindent
iota(\allowbreak{}lro\_, lro\_ in zz and not(\allowbreak{}lrf\_(\allowbreak{}lro\_) = zero(\allowbreak{}lrk\_)) and forall(\allowbreak{}[lrj\_ in zz], lrj\_ < lro\_ implies lrf\_(\allowbreak{}lrj\_) = zero(\allowbreak{}lrk\_)))
\end{quote}
characterized by \texttt{laurent-ord-eq}  \textit{(proven)}:
\begin{quote}\ttfamily\footnotesize\raggedright\noindent
forall(\allowbreak{}[kq\_], is-ring(\allowbreak{}kq\_) implies forall(\allowbreak{}[fq\_, oq\_], oq\_ in zz and not(\allowbreak{}fq\_(\allowbreak{}oq\_) = zero(\allowbreak{}kq\_)) and forall(\allowbreak{}[lzj in zz], lzj < oq\_ implies fq\_(\allowbreak{}lzj) = zero(\allowbreak{}kq\_)) implies laurent-ord(\allowbreak{}kq\_, fq\_) = oq\_))
\end{quote}
also: \texttt{laurent-ord-mul-eq} \\
mentioned by 4 result(s)

\smallskip

\noindent\textbf{\texttt{laurent-pairing(a5k\_, a5c\_)}}\quad\textit{functoid}\index{laurent-pairing(a5k_, a5c_)@\texttt{laurent-pairing(a5k\_, a5c\_)}}\\
reads: the pairing (g, f) |-> \$2(a\_0(f g)) on the Laurent series over \$1 \\
definition:
\begin{quote}\ttfamily\footnotesize\raggedright\noindent
vnb-lambda(\allowbreak{}[a5g\_, a5f\_], cartesian(\allowbreak{}laurent-set(\allowbreak{}a5k\_), laurent-set(\allowbreak{}a5k\_)), a5c\_(\allowbreak{}(\allowbreak{}laurent-mul(\allowbreak{}a5k\_, a5f\_, a5g\_))(\allowbreak{}0)))
\end{quote}
mentioned by 6 result(s)

\smallskip

\noindent\textbf{\texttt{laurent-polys(llk\_, llm\_)}}\quad\textit{functoid}\index{laurent-polys(llk_, llm_)@\texttt{laurent-polys(llk\_, llm\_)}}\\
reads: the polynomials over \$1 of degree below \$2 \\
definition:
\begin{quote}\ttfamily\footnotesize\raggedright\noindent
\{llg\_ in fun(\allowbreak{}zz, carr(\allowbreak{}llk\_)): forall(\allowbreak{}[llj\_ in zz], llj\_ < 0 implies llg\_(\allowbreak{}llj\_) = zero(\allowbreak{}llk\_)) and forall(\allowbreak{}[llj\_ in zz], llm\_ <= llj\_ implies llg\_(\allowbreak{}llj\_) = zero(\allowbreak{}llk\_))\}
\end{quote}
characterized by \texttt{laurent-polys-iff}  \textit{(proven)}:
\begin{quote}\ttfamily\footnotesize\raggedright\noindent
forall(\allowbreak{}[kq\_, mq\_, gq\_], gq\_ in laurent-polys(\allowbreak{}kq\_, mq\_) iff gq\_ in fun(\allowbreak{}zz, carr(\allowbreak{}kq\_)) and forall(\allowbreak{}[lzj in zz], lzj < 0 implies gq\_(\allowbreak{}lzj) = zero(\allowbreak{}kq\_)) and forall(\allowbreak{}[lzj in zz], mq\_ <= lzj implies gq\_(\allowbreak{}lzj) = zero(\allowbreak{}kq\_)))
\end{quote}
also: \texttt{laurent-polys-set} \texttt{laurent-polys-step} \texttt{laurent-polys-finite} \texttt{laurent-polys-zero-finite} \texttt{laurent-polys-sub-integers} \\
mentioned by 8 result(s)

\smallskip

\noindent\textbf{\texttt{laurent-ring(lrk\_)}}\quad\textit{functoid}\index{laurent-ring(lrk_)@\texttt{laurent-ring(lrk\_)}}\\
reads: the ring of formal Laurent series \$1((x)) \\
definition:
\begin{quote}\ttfamily\footnotesize\raggedright\noindent
[laurent-set(\allowbreak{}lrk\_), vnb-lambda(\allowbreak{}[lrf\_, lrg\_], cartesian(\allowbreak{}laurent-set(\allowbreak{}lrk\_), laurent-set(\allowbreak{}lrk\_)), laurent-add(\allowbreak{}lrk\_, lrf\_, lrg\_)), vnb-lambda(\allowbreak{}[lrf\_, lrg\_], cartesian(\allowbreak{}laurent-set(\allowbreak{}lrk\_), laurent-set(\allowbreak{}lrk\_)), laurent-mul(\allowbreak{}lrk\_, lrf\_, lrg\_)), vnb-lambda(\allowbreak{}lrf\_, laurent-set(\allowbreak{}lrk\_), laurent-neg(\allowbreak{}lrk\_, lrf\_)), laurent-zero(\allowbreak{}lrk\_), laurent-one(\allowbreak{}lrk\_)]
\end{quote}
characterized by \texttt{laurent-ring-add}  \textit{(proven)}:
\begin{quote}\ttfamily\footnotesize\raggedright\noindent
forall(\allowbreak{}[kq\_], is-ring(\allowbreak{}kq\_) implies add(\allowbreak{}laurent-ring(\allowbreak{}kq\_)) == vnb-lambda(\allowbreak{}[lrf\_, lrg\_], cartesian(\allowbreak{}laurent-set(\allowbreak{}kq\_), laurent-set(\allowbreak{}kq\_)), laurent-add(\allowbreak{}kq\_, lrf\_, lrg\_)))
\end{quote}
also: \texttt{laurent-ring-mul} \texttt{laurent-ring-neg} \texttt{laurent-ring-one} \texttt{laurent-ring-carr} \texttt{laurent-ring-zero} \texttt{laurent-ring-is-ring} \texttt{laurent-ring-is-field-ring} \texttt{laurent-ring-is-commutative-ring} \\
mentioned by 29 result(s)

\smallskip

\noindent\textbf{\texttt{laurent-set(lrk\_)}}\quad\textit{functoid}\index{laurent-set(lrk_)@\texttt{laurent-set(lrk\_)}}\\
reads: the formal Laurent series \$1((x)) \\
definition:
\begin{quote}\ttfamily\footnotesize\raggedright\noindent
\{lrf\_ in fun(\allowbreak{}zz, carr(\allowbreak{}lrk\_)): is-laurent(\allowbreak{}lrk\_, lrf\_)\}
\end{quote}
characterized by \texttt{laurent-set-in}  \textit{(proven)}:
\begin{quote}\ttfamily\footnotesize\raggedright\noindent
forall(\allowbreak{}[lzk, lzf in fun(\allowbreak{}zz, carr(\allowbreak{}lzk)), lzb in zz], forall(\allowbreak{}[lzj in zz], lzj < lzb implies lzf(\allowbreak{}lzj) = zero(\allowbreak{}lzk)) implies lzf in laurent-set(\allowbreak{}lzk))
\end{quote}
also: \texttt{laurent-set-out} \texttt{laurent-set-unfold} \\
mentioned by 109 result(s)

\smallskip

\noindent\textbf{\texttt{laurent-sum(lsc\_, lsa\_, lsz\_)}}\quad\textit{functoid}\index{laurent-sum(lsc_, lsa_, lsz_)@\texttt{laurent-sum(lsc\_, lsa\_, lsz\_)}}\\
reads: the sum of the Laurent series with coefficients \$1 about \$2 at \$3 \\
definition:
\begin{quote}\ttfamily\footnotesize\raggedright\noindent
cc-series-limit(\allowbreak{}vnb-lambda(\allowbreak{}lsk\_, nn, lsc\_(\allowbreak{}lsk\_) * (\allowbreak{}lsz\_ - lsa\_) \textasciicircum{} lsk\_)) + recip(\allowbreak{}lsz\_ - lsa\_) * cc-series-limit(\allowbreak{}vnb-lambda(\allowbreak{}lsk\_, nn, lsc\_(\allowbreak{}-(\allowbreak{}lsk\_ + 1)) * recip(\allowbreak{}lsz\_ - lsa\_) \textasciicircum{} lsk\_))
\end{quote}
mentioned by 11 result(s)

\smallskip

\noindent\textbf{\texttt{laurent-supp(lrk\_, lrf\_)}}\quad\textit{functoid}\index{laurent-supp(lrk_, lrf_)@\texttt{laurent-supp(lrk\_, lrf\_)}}\\
reads: the support of the Laurent series \$2 \\
definition:
\begin{quote}\ttfamily\footnotesize\raggedright\noindent
\{lrj\_ in zz: not(\allowbreak{}lrf\_(\allowbreak{}lrj\_) = zero(\allowbreak{}lrk\_))\}
\end{quote}

\smallskip

\noindent\textbf{\texttt{laurent-trunc(lrk\_, lrf\_, lrt\_)}}\quad\textit{functoid}\index{laurent-trunc(lrk_, lrf_, lrt_)@\texttt{laurent-trunc(lrk\_, lrf\_, lrt\_)}}\\
reads: the truncation of \$2 above \$3 \\
definition:
\begin{quote}\ttfamily\footnotesize\raggedright\noindent
vnb-lambda(\allowbreak{}lrj\_, zz, if(\allowbreak{}lrj\_ <= lrt\_, lrf\_(\allowbreak{}lrj\_), zero(\allowbreak{}lrk\_)))
\end{quote}
characterized by \texttt{laurent-trunc-gt}  \textit{(proven)}:
\begin{quote}\ttfamily\footnotesize\raggedright\noindent
forall(\allowbreak{}[kq\_], is-ring(\allowbreak{}kq\_) implies forall(\allowbreak{}[fq\_ in fun(\allowbreak{}zz, carr(\allowbreak{}kq\_)), tq\_ in zz, jq\_ in zz], tq\_ < jq\_ implies (\allowbreak{}laurent-trunc(\allowbreak{}kq\_, fq\_, tq\_))(\allowbreak{}jq\_) = zero(\allowbreak{}kq\_)))
\end{quote}
also: \texttt{laurent-trunc-le} \texttt{laurent-trunc-fun} \texttt{laurent-trunc-low} \texttt{laurent-trunc-poly} \texttt{laurent-trunc-agree} \texttt{laurent-trunc-finsupp} \\
mentioned by 13 result(s)

\smallskip

\noindent\textbf{\texttt{laurent-zero(lrk\_)}}\quad\textit{functoid}\index{laurent-zero(lrk_)@\texttt{laurent-zero(lrk\_)}}\\
reads: the zero Laurent series over \$1 \\
definition:
\begin{quote}\ttfamily\footnotesize\raggedright\noindent
vnb-lambda(\allowbreak{}lrj\_, zz, zero(\allowbreak{}lrk\_))
\end{quote}
characterized by \texttt{laurent-zero-in}  \textit{(proven)}:
\begin{quote}\ttfamily\footnotesize\raggedright\noindent
forall(\allowbreak{}[lzk], is-ring(\allowbreak{}lzk) implies laurent-zero(\allowbreak{}lzk) in laurent-set(\allowbreak{}lzk))
\end{quote}
also: \texttt{laurent-zero-apply} \texttt{laurent-zero-ne-one} \\
mentioned by 20 result(s)

\smallskip

\noindent\textbf{\texttt{lcd-stage}}\quad\textit{functoid}\index{lcd-stage@\texttt{lcd-stage}}\\
reads: the closing-depth map on the stage \$1 \\
defined by \texttt{lcd-stage-zero} \texttt{lcd-stage-succ} \\
characterized by \texttt{lcd-stage-succ}  \textit{(definitional)}:
\begin{quote}\ttfamily\footnotesize\raggedright\noindent
forall(\allowbreak{}[mxn\_ in nn], lcd-stage(\allowbreak{}succ(\allowbreak{}mxn\_)) == vnb-lambda(\allowbreak{}mxe\_, expr-stage(\allowbreak{}succ(\allowbreak{}mxn\_)), if(\allowbreak{}mxe\_ in expr-stage(\allowbreak{}mxn\_), (\allowbreak{}lcd-stage(\allowbreak{}mxn\_))(\allowbreak{}mxe\_), if(\allowbreak{}nth(\allowbreak{}1, mxe\_) = 3, \{mxd\_ in nn: forall(\allowbreak{}[mxw\_ in make-set(\allowbreak{}nth(\allowbreak{}3, mxe\_))], mxd\_ in (\allowbreak{}lcd-stage(\allowbreak{}mxn\_))(\allowbreak{}mxw\_))\}, \{mxd\_ in nn: forall(\allowbreak{}[mxw\_ in make-set(\allowbreak{}nth(\allowbreak{}4, mxe\_))], mxd\_ in (\allowbreak{}lcd-stage(\allowbreak{}mxn\_))(\allowbreak{}mxw\_)) and mxd\_ + nth(\allowbreak{}3, mxe\_) in (\allowbreak{}lcd-stage(\allowbreak{}mxn\_))(\allowbreak{}nth(\allowbreak{}5, mxe\_))\}))))
\end{quote}
also: \texttt{lcd-stage-zero} \texttt{lcd-stage-coh} \texttt{lcd-stage-step} \\
mentioned by 10 result(s)

\smallskip

\noindent\textbf{\texttt{ldir-deriv(ldm\_, ldf\_, lda\_, lde\_)}}\quad\textit{functoid}\index{ldir-deriv(ldm_, ldf_, lda_, lde_)@\texttt{ldir-deriv(ldm\_, ldf\_, lda\_, lde\_)}}\\
reads: the directional derivative of \$2 at \$3 in the direction \$4 (in \$1) \\
definition:
\begin{quote}\ttfamily\footnotesize\raggedright\noindent
iota(\allowbreak{}ldl\_, has-dir-deriv(\allowbreak{}ldm\_, ldf\_, lda\_, lde\_, ldl\_))
\end{quote}
characterized by \texttt{ldir-deriv-value}  \textit{(proven)}:
\begin{quote}\ttfamily\footnotesize\raggedright\noindent
forall(\allowbreak{}[m, f, a, eta, dl], has-dir-deriv(\allowbreak{}m, f, a, eta, dl) implies ldir-deriv(\allowbreak{}m, f, a, eta) = dl)
\end{quote}
mentioned by 3 result(s)

\smallskip

\noindent\textbf{\texttt{leadcoef(a, p)}}\quad\textit{functoid}\index{leadcoef(a, p)@\texttt{leadcoef(a, p)}}\\
reads: the leading coefficient of \$2 \\
definition:
\begin{quote}\ttfamily\footnotesize\raggedright\noindent
p(\allowbreak{}deg(\allowbreak{}a, p))
\end{quote}
characterized by \texttt{leadcoef-unfold}  \textit{(proven)}:
\begin{quote}\ttfamily\footnotesize\raggedright\noindent
forall(\allowbreak{}[a\_, p\_], leadcoef(\allowbreak{}a\_, p\_) == p\_(\allowbreak{}deg(\allowbreak{}a\_, p\_)))
\end{quote}
also: \texttt{leadcoef-nonzero} \texttt{monomial-leadcoef} \texttt{leadcoef-nonzero-ptwise} \\
mentioned by 6 result(s)

\smallskip

\noindent\textbf{\texttt{left-coset(gcg\_, gca\_, gch\_)}}\quad\textit{functoid}\index{left-coset(gcg_, gca_, gch_)@\texttt{left-coset(gcg\_, gca\_, gch\_)}}\\
reads: the left coset of \$3 by \$2 in \$1 \\
definition:
\begin{quote}\ttfamily\footnotesize\raggedright\noindent
\{gcx\_ in carr(\allowbreak{}gcg\_): forsome(\allowbreak{}[gcy\_ in gch\_], gcx\_ = (\allowbreak{}opr(\allowbreak{}gcg\_))(\allowbreak{}gca\_, gcy\_))\}
\end{quote}
characterized by \texttt{left-coset-eq}  \textit{(proven)}:
\begin{quote}\ttfamily\footnotesize\raggedright\noindent
forall(\allowbreak{}[g, h, a, b], is-group(\allowbreak{}g) implies is-subgroup(\allowbreak{}g, h) implies a in carr(\allowbreak{}g) implies b in left-coset(\allowbreak{}g, a, h) implies left-coset(\allowbreak{}g, b, h) = left-coset(\allowbreak{}g, a, h))
\end{quote}
also: \texttt{left-coset-member} \texttt{left-coset-unfold} \texttt{left-cosets-unfold} \texttt{qag-class-left-coset} \\
mentioned by 9 result(s)

\smallskip

\noindent\textbf{\texttt{left-cosets(gcg\_, gch\_)}}\quad\textit{functoid}\index{left-cosets(gcg_, gch_)@\texttt{left-cosets(gcg\_, gch\_)}}\\
reads: the set of left cosets of \$2 in \$1 \\
definition:
\begin{quote}\ttfamily\footnotesize\raggedright\noindent
\{gcc\_ in power(\allowbreak{}carr(\allowbreak{}gcg\_)): forsome(\allowbreak{}[gca\_ in carr(\allowbreak{}gcg\_)], gcc\_ = left-coset(\allowbreak{}gcg\_, gca\_, gch\_))\}
\end{quote}
characterized by \texttt{left-cosets-unfold}  \textit{(proven)}:
\begin{quote}\ttfamily\footnotesize\raggedright\noindent
forall(\allowbreak{}[g, h], left-cosets(\allowbreak{}g, h) == \{gcc\_ in power(\allowbreak{}carr(\allowbreak{}g)): forsome(\allowbreak{}[gca\_ in carr(\allowbreak{}g)], gcc\_ = left-coset(\allowbreak{}g, gca\_, h))\})
\end{quote}
also: \texttt{qag-carr-left-cosets} \\
mentioned by 6 result(s)

\smallskip

\noindent\textbf{\texttt{len-r}}\quad\textit{tactic}\index{len-r@\texttt{len-r}}\\
\texttt{(len-r)} -- Reduce a LENGTH applied to a literal LIST in the goal. \\
declared in \texttt{tactics-help.scm}

\smallskip

\noindent\textbf{\texttt{length}}\quad\textit{operator}\index{length@\texttt{length}}\\
a term-forming head \\
characterized by \texttt{length-cons}  \textit{(definitional)}:
\begin{quote}\ttfamily\footnotesize\raggedright\noindent
forall(\allowbreak{}[x, l], x in set implies length(\allowbreak{}cons(\allowbreak{}x, l)) = succ(\allowbreak{}length(\allowbreak{}l)))
\end{quote}
also: \texttt{cr-length-eq} \texttt{sipq-length-csip-as-sip} \texttt{qmod-length-cna-as-module} \texttt{qmod-length-sip-as-module} \texttt{qmod-length-uba-as-module} \texttt{qmod-length-cnvs-as-module} \texttt{qmod-length-sip-as-module-csip-as-sip} \texttt{qmod-length-uba-as-module-uiba-as-uba} ... (27 in all) \\
mentioned by 337 result(s)

\smallskip

\noindent\textbf{\texttt{lft(a, b, c, d)}}\quad\textit{functoid}\index{lft(a, b, c, d)@\texttt{lft(a, b, c, d)}}\\
reads: the linear fractional map z |-> (\$1 z + \$2)/(\$3 z + \$4) of the extended plane \\
definition:
\begin{quote}\ttfamily\footnotesize\raggedright\noindent
vnb-lambda(\allowbreak{}lfz\_, cc-hat, if(\allowbreak{}lfz\_ = cc-inf, if(\allowbreak{}c = 0, cc-inf, a * recip(\allowbreak{}c)), if(\allowbreak{}c * lfz\_ + d = 0, cc-inf, (\allowbreak{}a * lfz\_ + b) * recip(\allowbreak{}c * lfz\_ + d))))
\end{quote}
characterized by \texttt{lfts-def}  \textit{(definitional)}:
\begin{quote}\ttfamily\footnotesize\raggedright\noindent
lfts == \{lff\_ in fun(\allowbreak{}cc-hat, cc-hat): forsome(\allowbreak{}[lfa\_ in cc, lfb\_ in cc, lfc\_ in cc, lfd\_ in cc], not(\allowbreak{}lfa\_ * lfd\_ - lfb\_ * lfc\_ = 0) and lff\_ = lft(\allowbreak{}lfa\_, lfb\_, lfc\_, lfd\_))\}
\end{quote}
also: \texttt{lft-homog} \texttt{lft-at-inf} \texttt{lft-in-fun} \texttt{lfts-intro} \texttt{lft-at-pole} \texttt{lft-compose} \texttt{lft-at-finite} \texttt{lft-bijection} ... (17 in all) \\
mentioned by 48 result(s)

\smallskip

\noindent\textbf{\texttt{lft-group}}\quad\textit{instance}\index{lft-group@\texttt{lft-group}}\\
reads: the group of linear fractional transformations under composition \\
an instance of group; predicate \texttt{is-lft-group}; same shape as group

\smallskip

\noindent\textbf{\texttt{lfts}}\quad\textit{constant}\index{lfts@\texttt{lfts}}\\
reads: the set of linear fractional transformations of the extended plane \\
defined by \texttt{lfts-def}

\smallskip

\noindent\textbf{\texttt{limit-form(md, fq)}}\quad\textit{functoid}\index{limit-form(md, fq)@\texttt{limit-form(md, fq)}}\\
reads: the limit of the inner products of two Cauchy sequences of (\$1, \$2) \\
definition:
\begin{quote}\ttfamily\footnotesize\raggedright\noindent
vnb-lambda(\allowbreak{}[hcf\_, hcg\_], cartesian(\allowbreak{}form-cseq(\allowbreak{}md, fq), form-cseq(\allowbreak{}md, fq)), iota(\allowbreak{}hcc\_, hcc\_ in cc and converges-to(\allowbreak{}cc-ms, seq-form(\allowbreak{}fq, hcf\_, hcg\_), hcc\_)))
\end{quote}
characterized by \texttt{limit-form-rr}  \textit{(proven)}:
\begin{quote}\ttfamily\footnotesize\raggedright\noindent
forall(\allowbreak{}[md, fq, f], is-module(\allowbreak{}md) implies scal(\allowbreak{}md) = normed-field-as-commutative-ring(\allowbreak{}cc-normed-field) implies is-inner-product-form(\allowbreak{}md, fq) implies f in form-cseq(\allowbreak{}md, fq) implies converges-to(\allowbreak{}rr-ms, seq-form(\allowbreak{}fq, f, f), (\allowbreak{}limit-form(\allowbreak{}md, fq))(\allowbreak{}f, f)))
\end{quote}
also: \texttt{limit-form-psd} \texttt{limit-form-diag} \texttt{limit-form-prop} \texttt{limit-form-type} \texttt{limit-form-embed} \texttt{limit-form-value} \texttt{limit-form-add-left} \texttt{limit-form-conj-sym} ... (9 in all) \\
mentioned by 30 result(s)

\smallskip

\noindent\textbf{\texttt{limit-ord}}\quad\textit{operator}\index{limit-ord@\texttt{limit-ord}}\\
a term-forming head \\
characterized by \texttt{limit-ord-iff}  \textit{(definitional)}:
\begin{quote}\ttfamily\footnotesize\raggedright\noindent
forall(\allowbreak{}[lambda], limit-ord(\allowbreak{}lambda) iff lambda in ord and not(\allowbreak{}lambda = 0) and not(\allowbreak{}forsome(\allowbreak{}[alpha in ord], lambda = succ\_ord(\allowbreak{}alpha))))
\end{quote}
also: \texttt{limit-ord-is-sup} \\
mentioned by 13 result(s)

\smallskip

\noindent\textbf{\texttt{lin-coords(ndsm\_, ndsd\_, ndsu\_, ndsp\_, ndsv\_)}}\quad\textit{functoid}\index{lin-coords(ndsm_, ndsd_, ndsu_, ndsp_, ndsv_)@\texttt{lin-coords(ndsm\_, ndsd\_, ndsu\_, ndsp\_, ndsv\_)}}\\
reads: the coordinates over \$4 of \$5 in the sequence \$3 \\
definition:
\begin{quote}\ttfamily\footnotesize\raggedright\noindent
iota(\allowbreak{}ndsc\_, ndsc\_ in mat(\allowbreak{}1, ndsd\_, ndsp\_) and ndsv\_ = lincomb(\allowbreak{}ndsm\_, ndsd\_, ndsc\_, ndsu\_))
\end{quote}
characterized by \texttt{lin-coords-spec}  \textit{(proven)}:
\begin{quote}\ttfamily\footnotesize\raggedright\noindent
forall(\allowbreak{}[ndm, ndd, ndu, ndp], forall(\allowbreak{}[ndc1 in mat(\allowbreak{}1, ndd, ndp), ndc2 in mat(\allowbreak{}1, ndd, ndp)], lincomb(\allowbreak{}ndm, ndd, ndc1, ndu) = lincomb(\allowbreak{}ndm, ndd, ndc2, ndu) implies ndc1 = ndc2) implies forall(\allowbreak{}[ndc in mat(\allowbreak{}1, ndd, ndp), ndr], ndr = lincomb(\allowbreak{}ndm, ndd, ndc, ndu) implies lin-coords(\allowbreak{}ndm, ndd, ndu, ndp, ndr) = ndc))
\end{quote}
mentioned by 2 result(s)

\smallskip

\noindent\textbf{\texttt{lin-span(gls-md\_, gls-n\_, gls-u\_)}}\quad\textit{functoid}\index{lin-span(gls-md_, gls-n_, gls-u_)@\texttt{lin-span(gls-md\_, gls-n\_, gls-u\_)}}\\
reads: the span of \$3 \\
definition:
\begin{quote}\ttfamily\footnotesize\raggedright\noindent
\{glsx\_ in vec(\allowbreak{}gls-md\_): forsome(\allowbreak{}[glsc\_ in mat(\allowbreak{}1, gls-n\_, carr(\allowbreak{}scal(\allowbreak{}gls-md\_)))], glsx\_ = lincomb(\allowbreak{}gls-md\_, gls-n\_, glsc\_, gls-u\_))\}
\end{quote}
characterized by \texttt{lin-span-act-cna-as-module}  \textit{(definitional)}:
\begin{quote}\ttfamily\footnotesize\raggedright\noindent
forall(\allowbreak{}[r], is-complex-normed-algebra(\allowbreak{}r) implies forall(\allowbreak{}[n in nn, u in mat(\allowbreak{}n, 1, vec(\allowbreak{}r)), r\_35006 in carr(\allowbreak{}scal(\allowbreak{}r)), x in lin-span(\allowbreak{}cna-as-module(\allowbreak{}r), n, u)], (\allowbreak{}act(\allowbreak{}r))(\allowbreak{}r\_35006, x) in lin-span(\allowbreak{}cna-as-module(\allowbreak{}r), n, u)))
\end{quote}
also: \texttt{lin-span-act-sip-as-module} \texttt{lin-span-act-uba-as-module} \texttt{lin-span-sum-cna-as-module} \texttt{lin-span-sum-sip-as-module} \texttt{lin-span-sum-uba-as-module} \texttt{lin-span-act-cnvs-as-module} \texttt{lin-span-elim-cna-as-module} \texttt{lin-span-elim-sip-as-module} ... (48 in all) \\
mentioned by 80 result(s)

\smallskip

\noindent\textbf{\texttt{lincomb(md, n, c, u)}}\quad\textit{functoid}\index{lincomb(md, n, c, u)@\texttt{lincomb(md, n, c, u)}}\\
definition:
\begin{quote}\ttfamily\footnotesize\raggedright\noindent
finsum(\allowbreak{}module-vector-ag(\allowbreak{}md), vnb-lambda(\allowbreak{}j, interval(\allowbreak{}1, n), (\allowbreak{}act(\allowbreak{}md))(\allowbreak{}entry(\allowbreak{}c, 1, j), entry(\allowbreak{}u, j, 1))), interval(\allowbreak{}1, n))
\end{quote}
characterized by \texttt{lincomb-snoc-cna-as-module}  \textit{(definitional)}:
\begin{quote}\ttfamily\footnotesize\raggedright\noindent
forall(\allowbreak{}[r], is-complex-normed-algebra(\allowbreak{}r) implies forall(\allowbreak{}[n in nn, c in mat(\allowbreak{}1, n, carr(\allowbreak{}scal(\allowbreak{}r))), u in mat(\allowbreak{}n, 1, vec(\allowbreak{}r)), r\_35016 in carr(\allowbreak{}scal(\allowbreak{}r)), x in vec(\allowbreak{}r)], lincomb(\allowbreak{}cna-as-module(\allowbreak{}r), succ(\allowbreak{}n), snoc-row(\allowbreak{}c, n, r\_35016), snoc-col(\allowbreak{}u, n, x)) = (\allowbreak{}vadd(\allowbreak{}r))(\allowbreak{}lincomb(\allowbreak{}cna-as-module(\allowbreak{}r), n, c, u), (\allowbreak{}act(\allowbreak{}r))(\allowbreak{}r\_35016, x))))
\end{quote}
also: \texttt{lincomb-snoc-sip-as-module} \texttt{lincomb-snoc-uba-as-module} \texttt{lincomb-type-cna-as-module} \texttt{lincomb-type-sip-as-module} \texttt{lincomb-type-uba-as-module} \texttt{lincomb-empty-cna-as-module} \texttt{lincomb-empty-sip-as-module} \texttt{lincomb-empty-uba-as-module} ... (103 in all) \\
mentioned by 528 result(s)

\smallskip

\noindent\textbf{\texttt{line(m, v)}}\quad\textit{functoid}\index{line(m, v)@\texttt{line(m, v)}}\\
definition:
\begin{quote}\ttfamily\footnotesize\raggedright\noindent
\{y\_ in vec(\allowbreak{}m): forsome(\allowbreak{}[r\_ in rr], y\_ = (\allowbreak{}act(\allowbreak{}m))(\allowbreak{}r\_, v))\}
\end{quote}
characterized by \texttt{lfe-line-subset-cnvs-as-nvs}  \textit{(definitional)}:
\begin{quote}\ttfamily\footnotesize\raggedright\noindent
forall(\allowbreak{}[r], is-complex-normed-vector-space(\allowbreak{}r) implies forall(\allowbreak{}[v], line(\allowbreak{}cnvs-as-nvs(\allowbreak{}r), v) subset vec(\allowbreak{}r)))
\end{quote}
also: \texttt{line-has-v} \texttt{line-unfold} \texttt{lfe-line-coef} \texttt{line-coef-spec} \texttt{line-seed-at-v} \texttt{lfe-line-in-vec} \texttt{lfe-line-mem-eq} \texttt{lfe-line-subset} ... (15 in all) \\
mentioned by 19 result(s)

\smallskip

\noindent\textbf{\texttt{line-coef(m, v, u)}}\quad\textit{functoid}\index{line-coef(m, v, u)@\texttt{line-coef(m, v, u)}}\\
reads: the coefficient of \$3 on the line through \$2 \\
definition:
\begin{quote}\ttfamily\footnotesize\raggedright\noindent
iota(\allowbreak{}lnt\_, lnt\_ in rr and u = (\allowbreak{}act(\allowbreak{}m))(\allowbreak{}lnt\_, v))
\end{quote}
characterized by \texttt{line-coef-spec}  \textit{(proven)}:
\begin{quote}\ttfamily\footnotesize\raggedright\noindent
forall(\allowbreak{}[m, v, u], is-module(\allowbreak{}m) implies scal(\allowbreak{}m) = normed-field-as-commutative-ring(\allowbreak{}rr-normed-field) implies v in vec(\allowbreak{}m) implies not(\allowbreak{}v = vzero(\allowbreak{}m)) implies u in line(\allowbreak{}m, v) implies line-coef(\allowbreak{}m, v, u) in rr and u = (\allowbreak{}act(\allowbreak{}m))(\allowbreak{}line-coef(\allowbreak{}m, v, u), v))
\end{quote}
also: \texttt{line-coef-value} \\
mentioned by 7 result(s)

\smallskip

\noindent\textbf{\texttt{line-int(pf, pgam, dgam, a, b)}}\quad\textit{functoid}\index{line-int(pf, pgam, dgam, a, b)@\texttt{line-int(pf, pgam, dgam, a, b)}}\\
reads: the integral of \$1 along the road \$2 with derivative \$3 from \$4 to \$5 \\
definition:
\begin{quote}\ttfamily\footnotesize\raggedright\noindent
cc-int(\allowbreak{}vnb-lambda(\allowbreak{}pat\_, ccint(\allowbreak{}a, b), pf(\allowbreak{}pgam(\allowbreak{}pat\_)) * dgam(\allowbreak{}pat\_)), a, b)
\end{quote}
characterized by \texttt{line-int-sum}  \textit{(proven)}:
\begin{quote}\ttfamily\footnotesize\raggedright\noindent
forall(\allowbreak{}[u, f, g, pgam, dgam, a, b], u subset cc implies f in fun(\allowbreak{}u, cc) implies forall(\allowbreak{}[rgy\_ in u], is-continuous-at(\allowbreak{}subspace-ms(\allowbreak{}nf-metric-space(\allowbreak{}cc-normed-field), u), nf-metric-space(\allowbreak{}cc-normed-field), f, rgy\_)) implies g in fun(\allowbreak{}u, cc) implies forall(\allowbreak{}[rgy\_ in u], is-continuous-at(\allowbreak{}subspace-ms(\allowbreak{}nf-metric-space(\allowbreak{}cc-normed-field), u), nf-metric-space(\allowbreak{}cc-normed-field), g, rgy\_)) implies is-road(\allowbreak{}pgam, dgam, a, b) implies trace(\allowbreak{}pgam, a, b) subset u implies forall(\allowbreak{}[lih\_ in fun(\allowbreak{}u, cc)], forall(\allowbreak{}[liz\_ in u], lih\_(\allowbreak{}liz\_) == f(\allowbreak{}liz\_) + g(\allowbreak{}liz\_)) implies line-int(\allowbreak{}lih\_, pgam, dgam, a, b) = line-int(\allowbreak{}f, pgam, dgam, a, b) + line-int(\allowbreak{}g, pgam, dgam, a, b)))
\end{quote}
also: \texttt{line-int-in-cc} \texttt{line-int-juxta} \texttt{line-int-value} \texttt{line-int-concat} \texttt{line-int-of-one} \texttt{line-int-scalar} \texttt{line-int-unfold} \texttt{homotopy-line-int} ... (30 in all) \\
mentioned by 170 result(s)

\smallskip

\noindent\textbf{\texttt{linf-norm(omega, ca, mu, f)}}\quad\textit{functoid}\index{linf-norm(omega, ca, mu, f)@\texttt{linf-norm(omega, ca, mu, f)}}\\
reads: the essential supremum norm of \$4 over \$1 with respect to \$3 \\
definition:
\begin{quote}\ttfamily\footnotesize\raggedright\noindent
einf(\allowbreak{}\{lim\_ in rr-pos-star: mu(\allowbreak{}\{lix\_ in omega: lim\_ < abs(\allowbreak{}f(\allowbreak{}lix\_))\}) = 0\})
\end{quote}
characterized by \texttt{linf-norm-in}  \textit{(proven)}:
\begin{quote}\ttfamily\footnotesize\raggedright\noindent
forall(\allowbreak{}[omega, ca, mu, f], linf-norm(\allowbreak{}omega, ca, mu, f) in rr-pos-star)
\end{quote}
also: \texttt{linf-norm-essential-bound} \texttt{comp-measure-preserving-linf-norm} \\
mentioned by 11 result(s)

\smallskip

\noindent\textbf{\texttt{list}}\quad\textit{operator}\index{list@\texttt{list}}\\
a term-forming head \\
characterized by \texttt{def-list}  \textit{(definitional)}:
\begin{quote}\ttfamily\footnotesize\raggedright\noindent
forall(\allowbreak{}[mdt\_, mdhs\_, mdc\_], def-list(\allowbreak{}mdt\_, mdhs\_, mdc\_) iff mdc\_ in seq-set and forsome(\allowbreak{}[mdg\_, mdl\_], mdc\_ = [mdg\_, enc-eq(\allowbreak{}enc-list(\allowbreak{}mdl\_), enc-list(\allowbreak{}mdl\_))] and length(\allowbreak{}mdhs\_) = length(\allowbreak{}mdl\_) and forall(\allowbreak{}[mdi\_ in nn], 1 <= mdi\_ implies mdi\_ <= length(\allowbreak{}mdl\_) implies nth(\allowbreak{}mdi\_, mdhs\_) = [mdg\_, enc-in(\allowbreak{}nth(\allowbreak{}mdi\_, mdl\_), enc-setc)])))
\end{quote}
also: \texttt{cr-def-list} \texttt{def-nth-list} \texttt{der-list-hole} \texttt{nsl-list-source} \texttt{nso-list-source} \texttt{flat-list-source} \texttt{cx-strand-list-typed} \texttt{rb5-nary-plus-2-list-cc} ... (15 in all) \\
mentioned by 1774 result(s)

\smallskip

\noindent\textbf{\texttt{little-o-at(g, a)}}\quad\textit{predicate}\index{little-o-at(g, a)@\texttt{little-o-at(g, a)}}\\
reads: \$1 is little-o at \$2 \\
defined by \texttt{little-o-at} \\
characterized by \texttt{little-o-at}  \textit{(definitional)}:
\begin{quote}\ttfamily\footnotesize\raggedright\noindent
forall(\allowbreak{}[g, a], little-o-at(\allowbreak{}g, a) iff g in fun(\allowbreak{}rr, rr) and a in rr and forsome(\allowbreak{}[eps in fun(\allowbreak{}rr, rr)], is-continuous-at(\allowbreak{}rr-ms, rr-ms, eps, a) and eps(\allowbreak{}a) = 0 and forall(\allowbreak{}[x in rr], g(\allowbreak{}x) = eps(\allowbreak{}x) * (\allowbreak{}x - a))))
\end{quote}
mentioned by 6 result(s)

\smallskip

\noindent\textbf{\texttt{log(x\_)}}\quad\textit{functoid}\index{log(x_)@\texttt{log(x\_)}}\\
reads: the logarithm of \$1 \\
definition:
\begin{quote}\ttfamily\footnotesize\raggedright\noindent
c-int-or(\allowbreak{}vnb-lambda(\allowbreak{}z\_, rr, recip-star(\allowbreak{}z\_)), 1, x\_)
\end{quote}
characterized by \texttt{log-mul}  \textit{(proven)}:
\begin{quote}\ttfamily\footnotesize\raggedright\noindent
forall(\allowbreak{}[x\_, y\_], x\_ in rr implies y\_ in rr implies 0 < x\_ implies 0 < y\_ implies log(\allowbreak{}x\_ * y\_) = log(\allowbreak{}x\_) + log(\allowbreak{}y\_))
\end{quote}
also: \texttt{log-one} \texttt{log-diff} \texttt{log-deriv} \texttt{log-in-rr} \texttt{log-pprod} \texttt{log-r-exp} \texttt{log-recip} \texttt{r-exp-log} ... (39 in all) \\
mentioned by 104 result(s)

\smallskip

\noindent\textbf{\texttt{lp-norm(omega, ca, mu, lpp\_, lpf\_)}}\quad\textit{functoid}\index{lp-norm(omega, ca, mu, lpp_, lpf_)@\texttt{lp-norm(omega, ca, mu, lpp\_, lpf\_)}}\\
reads: the L\textasciicircum{}\$4 seminorm of \$5 over \$1 with respect to \$3 \\
definition:
\begin{quote}\ttfamily\footnotesize\raggedright\noindent
ep-norm(\allowbreak{}omega, ca, mu, lpp\_, abs-fn(\allowbreak{}omega, lpf\_))
\end{quote}
characterized by \texttt{lp-norm-abs-fn}  \textit{(proven)}:
\begin{quote}\ttfamily\footnotesize\raggedright\noindent
forall(\allowbreak{}[omega, ca, mu, p, g], is-measure(\allowbreak{}omega, ca, mu) implies is-real-measurable-fn(\allowbreak{}omega, ca, g) implies lp-norm(\allowbreak{}omega, ca, mu, p, abs-fn(\allowbreak{}omega, g)) == lp-norm(\allowbreak{}omega, ca, mu, p, g))
\end{quote}
also: \texttt{lp-norm-sub-le} \texttt{lp-norm-unfold} \texttt{l2-dist-lp-norm} \texttt{lp-norm-as-epow} \texttt{lp-norm-triangle} \texttt{lp-norm-zero-iff} \texttt{lp-norm-finite-iff} \texttt{lp-norm-homogeneous} ... (9 in all) \\
mentioned by 26 result(s)

\smallskip

\noindent\textbf{\texttt{lst(sml\_)}}\quad\textit{functoid}\index{lst(sml_)@\texttt{lst(sml\_)}}\\
reads: the list message of \$1 \\
definition:
\begin{quote}\ttfamily\footnotesize\raggedright\noindent
[1, sml\_]
\end{quote}
characterized by \texttt{lst-unfold}  \textit{(proven)}:
\begin{quote}\ttfamily\footnotesize\raggedright\noindent
forall(\allowbreak{}[l], lst(\allowbreak{}l) == [1, l])
\end{quote}
also: \texttt{lst-ne-encr} \texttt{lst-injective} \texttt{token-not-lst} \texttt{lst-in-tuples-set} \\
mentioned by 107 result(s)

\smallskip

\noindent\textbf{\texttt{ltkey(s)}}\quad\textit{accessor}\index{ltkey(s)@\texttt{ltkey(s)}}\\
reads: the ltkey slot of \$1 \\
a structure slot \\
characterized by \texttt{cx-orp@ltkey}  \textit{(definitional)}:
\begin{quote}\ttfamily\footnotesize\raggedright\noindent
ltkey(\allowbreak{}cx-orp) == vnb-lambda(\allowbreak{}cxn\_, interval(\allowbreak{}10, 11), cxn\_ + 20)
\end{quote}
also: \texttt{or-ltkey-injective-at} \texttt{or-ltkey-not-regular-origin} \texttt{or-ltkey-no-adversary-origin} \\
mentioned by 28 result(s)

\smallskip


\section*{M}
\addcontentsline{toc}{section}{M}

\noindent\textbf{\texttt{mac}}\quad\textit{tactic}\index{mac@\texttt{mac}}\\
\texttt{(mac 'name)} -- Rewrite the GOAL with an equivalence macete (unfold a definition, apply an iff/=/== law).  Fires only where the macete's side-conditions already hold in context. \\
declared in \texttt{tactics-help.scm}

\smallskip

\noindent\textbf{\texttt{mac-h}}\quad\textit{tactic}\index{mac-h@\texttt{mac-h}}\\
\texttt{(mac-h 'name hyp)} -- Rewrite a cited ASSUMPTION in place with an equivalence macete; any side-condition not already in context is spawned as a new subgoal. \\
declared in \texttt{tactics-help.scm}

\smallskip

\noindent\textbf{\texttt{mac-h*}}\quad\textit{tactic}\index{mac-h*@\texttt{mac-h*}}\\
\texttt{(mac-h*)} -- Saturating mac-h: repeatedly unfold every defined predicate in the hypotheses and split the conjunctions they expose, until nothing is left folded.  One step in the trace instead of a dozen. \\
declared in \texttt{tactics-help.scm}

\smallskip

\noindent\textbf{\texttt{macm}}\quad\textit{tactic}\index{macm@\texttt{macm}}\\
\texttt{(macm 'name)} -- Like mac, but a CONDITIONAL macete fires even when its side-conditions are not yet in context: each unmet condition is left as a new subgoal. \\
declared in \texttt{tactics-help.scm}

\smallskip

\noindent\textbf{\texttt{magnitude}}\quad\textit{operator}\index{magnitude@\texttt{magnitude}}\\
a term-forming head \\
characterized by \texttt{magnitude-def}  \textit{(definitional)}:
\begin{quote}\ttfamily\footnotesize\raggedright\noindent
forall(\allowbreak{}[z], magnitude(\allowbreak{}z) == sqrt(\allowbreak{}real-part(\allowbreak{}z) * real-part(\allowbreak{}z) + imag-part(\allowbreak{}z) * imag-part(\allowbreak{}z)))
\end{quote}
also: \texttt{cc-magnitude-i} \texttt{half-magnitude} \texttt{cc-magnitude-sq} \texttt{pprod-magnitude} \texttt{rr-magnitude-sq} \texttt{cc-exp-magnitude} \texttt{cc-magnitude-mul} \texttt{cc-magnitude-neg} ... (48 in all) \\
mentioned by 650 result(s)

\smallskip

\noindent\textbf{\texttt{make-metric-group(mmm\_, mmg\_)}}\quad\textit{functoid}\index{make-metric-group(mmm_, mmg_)@\texttt{make-metric-group(mmm\_, mmg\_)}}\\
reads: the group \$2 with the distance of \$1 \\
definition:
\begin{quote}\ttfamily\footnotesize\raggedright\noindent
[carr(\allowbreak{}mmg\_), opr(\allowbreak{}mmg\_), iden(\allowbreak{}mmg\_), inv(\allowbreak{}mmg\_), dist(\allowbreak{}mmm\_)]
\end{quote}
characterized by \texttt{make-metric-group-inv}  \textit{(proven)}:
\begin{quote}\ttfamily\footnotesize\raggedright\noindent
forall(\allowbreak{}[ms, g], is-metric-space(\allowbreak{}ms) implies is-group(\allowbreak{}g) implies inv(\allowbreak{}make-metric-group(\allowbreak{}ms, g)) == inv(\allowbreak{}g))
\end{quote}
also: \texttt{make-metric-group-opr} \texttt{make-metric-group-carr} \texttt{make-metric-group-iden} \texttt{make-metric-group-gdist} \texttt{make-metric-group-is-metric-group} \\
mentioned by 13 result(s)

\smallskip

\noindent\textbf{\texttt{make-set}}\quad\textit{operator}\index{make-set@\texttt{make-set}}\\
a term-forming head \\
characterized by \texttt{make-set-empty}  \textit{(primitive)}:
\begin{quote}\ttfamily\footnotesize\raggedright\noindent
\{\} = empty-set
\end{quote}
also: \texttt{make-set-sethood} \texttt{make-set-membership} \\
mentioned by 86 result(s)

\smallskip

\noindent\textbf{\texttt{make-top-group(mks\_, mkg\_)}}\quad\textit{functoid}\index{make-top-group(mks_, mkg_)@\texttt{make-top-group(mks\_, mkg\_)}}\\
reads: the group \$2 with the topology of \$1 \\
definition:
\begin{quote}\ttfamily\footnotesize\raggedright\noindent
[carr(\allowbreak{}mkg\_), opr(\allowbreak{}mkg\_), iden(\allowbreak{}mkg\_), inv(\allowbreak{}mkg\_), opens(\allowbreak{}mks\_)]
\end{quote}
characterized by \texttt{make-top-group-inv}  \textit{(proven)}:
\begin{quote}\ttfamily\footnotesize\raggedright\noindent
forall(\allowbreak{}[s, g], is-top-space(\allowbreak{}s) implies is-group(\allowbreak{}g) implies inv(\allowbreak{}make-top-group(\allowbreak{}s, g)) == inv(\allowbreak{}g))
\end{quote}
also: \texttt{make-top-group-opr} \texttt{make-top-group-carr} \texttt{make-top-group-iden} \texttt{make-top-group-topology} \texttt{make-top-group-is-top-group} \\
mentioned by 13 result(s)

\smallskip

\noindent\textbf{\texttt{map-seq(h, n, u)}}\quad\textit{functoid}\index{map-seq(h, n, u)@\texttt{map-seq(h, n, u)}}\\
reads: the image under \$1 of the sequence \$3 of length \$2 \\
definition:
\begin{quote}\ttfamily\footnotesize\raggedright\noindent
matof(\allowbreak{}n, 1, vnb-lambda(\allowbreak{}[msa\_, msb\_], cartesian(\allowbreak{}interval(\allowbreak{}1, n), interval(\allowbreak{}1, 1)), h(\allowbreak{}entry(\allowbreak{}u, msa\_, msb\_))))
\end{quote}
characterized by \texttt{map-seq-type}  \textit{(proven)}:
\begin{quote}\ttfamily\footnotesize\raggedright\noindent
forall(\allowbreak{}[h, v, w, n, u], h in fun(\allowbreak{}v, w) implies n in nn implies u in mat(\allowbreak{}n, 1, v) implies map-seq(\allowbreak{}h, n, u) in mat(\allowbreak{}n, 1, w))
\end{quote}
also: \texttt{map-seq-block} \texttt{map-seq-entry} \texttt{map-seq-ident} \texttt{map-seq-compose} \\
mentioned by 22 result(s)

\smallskip

\noindent\textbf{\texttt{mapl(mxf\_, mxl\_)}}\quad\textit{functoid}\index{mapl(mxf_, mxl_)@\texttt{mapl(mxf\_, mxl\_)}}\\
reads: the list of the values of \$1 at the entries of \$2 \\
definition:
\begin{quote}\ttfamily\footnotesize\raggedright\noindent
iota(\allowbreak{}mxr\_, mxr\_ in tuples(\allowbreak{}set) and length(\allowbreak{}mxr\_) = length(\allowbreak{}mxl\_) and forall(\allowbreak{}[mxi\_ in nn], 1 <= mxi\_ implies mxi\_ <= length(\allowbreak{}mxl\_) implies nth(\allowbreak{}mxi\_, mxr\_) = mxf\_(\allowbreak{}nth(\allowbreak{}mxi\_, mxl\_))))
\end{quote}
characterized by \texttt{mapl-prop}  \textit{(proven)}:
\begin{quote}\ttfamily\footnotesize\raggedright\noindent
forall(\allowbreak{}[cls, mxg\_, l], length(\allowbreak{}l) in nn implies forall(\allowbreak{}[i in nn], 1 <= i implies i <= length(\allowbreak{}l) implies mxg\_(\allowbreak{}nth(\allowbreak{}i, l)) in cls) implies mapl(\allowbreak{}mxg\_, l) in tuples(\allowbreak{}cls) and length(\allowbreak{}mapl(\allowbreak{}mxg\_, l)) = length(\allowbreak{}l) and forall(\allowbreak{}[i in nn], 1 <= i implies i <= length(\allowbreak{}l) implies nth(\allowbreak{}i, mapl(\allowbreak{}mxg\_, l)) = mxg\_(\allowbreak{}nth(\allowbreak{}i, l))))
\end{quote}
mentioned by 7 result(s)

\smallskip

\noindent\textbf{\texttt{maplk(mxf\_, mxl\_, mxk\_)}}\quad\textit{functoid}\index{maplk(mxf_, mxl_, mxk_)@\texttt{maplk(mxf\_, mxl\_, mxk\_)}}\\
reads: the list of the values of \$1 at the entries of \$2 and at \$3 \\
definition:
\begin{quote}\ttfamily\footnotesize\raggedright\noindent
iota(\allowbreak{}mxr\_, mxr\_ in tuples(\allowbreak{}set) and length(\allowbreak{}mxr\_) = length(\allowbreak{}mxl\_) and forall(\allowbreak{}[mxi\_ in nn], 1 <= mxi\_ implies mxi\_ <= length(\allowbreak{}mxl\_) implies nth(\allowbreak{}mxi\_, mxr\_) = mxf\_(\allowbreak{}nth(\allowbreak{}mxi\_, mxl\_), mxk\_)))
\end{quote}
characterized by \texttt{maplk-nil}  \textit{(proven)}:
\begin{quote}\ttfamily\footnotesize\raggedright\noindent
forall(\allowbreak{}[u, k], u in expr-set implies k in nn implies maplk(\allowbreak{}open-fun(\allowbreak{}u), [], k) = [])
\end{quote}
also: \texttt{maplk-open} \texttt{maplk-prop} \\
mentioned by 12 result(s)

\smallskip

\noindent\textbf{\texttt{mat(m, n, x)}}\quad\textit{functoid}\index{mat(m, n, x)@\texttt{mat(m, n, x)}}\\
reads: the \$1-by-\$2 matrices over \$3 \\
definition:
\begin{quote}\ttfamily\footnotesize\raggedright\noindent
\{p in matrix(\allowbreak{}x): n in nn and (\allowbreak{}length(\allowbreak{}p) = 0 implies m = 0) and (\allowbreak{}not(\allowbreak{}length(\allowbreak{}p) = 0) implies size(\allowbreak{}p) = [m, n])\}
\end{quote}
characterized by \texttt{is-mat-bound}  \textit{(definitional)}:
\begin{quote}\ttfamily\footnotesize\raggedright\noindent
forall(\allowbreak{}[k, m, n, a, c], is-mat-bound(\allowbreak{}k, m, n, a, c) iff a in mat(\allowbreak{}m, n, carr(\allowbreak{}k)) and forall(\allowbreak{}[cbvi\_ in interval(\allowbreak{}1, m), cbvj\_ in interval(\allowbreak{}1, n)], (\allowbreak{}fnrm(\allowbreak{}k))(\allowbreak{}entry(\allowbreak{}a, cbvi\_, cbvj\_)) <= c))
\end{quote}
also: \texttt{is-mat-close} \texttt{is-matrix-of} \texttt{cc-matrices-def} \texttt{is-invertible-mat} \texttt{is-irreducible-mat} \texttt{matadd-comm-uba-as-ring} \texttt{matadd-type-uba-as-ring} \texttt{matmul-type-uba-as-ring} ... (307 in all) \\
mentioned by 2371 result(s)

\smallskip

\noindent\textbf{\texttt{mat-act(gma-md\_, gma-n\_, gma-p\_, gma-u\_)}}\quad\textit{functoid}\index{mat-act(gma-md_, gma-n_, gma-p_, gma-u_)@\texttt{mat-act(gma-md\_, gma-n\_, gma-p\_, gma-u\_)}}\\
reads: the matrix \$3 of maps applied to the column \$4 \\
definition:
\begin{quote}\ttfamily\footnotesize\raggedright\noindent
matof(\allowbreak{}gma-n\_, 1, vnb-lambda(\allowbreak{}[gmai\_, gmaj\_], cartesian(\allowbreak{}interval(\allowbreak{}1, gma-n\_), interval(\allowbreak{}1, 1)), finsum(\allowbreak{}module-vector-ag(\allowbreak{}gma-md\_), vnb-lambda(\allowbreak{}gmak\_, interval(\allowbreak{}1, gma-n\_), (\allowbreak{}entry(\allowbreak{}gma-p\_, gmai\_, gmak\_))(\allowbreak{}entry(\allowbreak{}gma-u\_, gmak\_, 1))), interval(\allowbreak{}1, gma-n\_))))
\end{quote}
characterized by \texttt{mat-act-entry}  \textit{(proven)}:
\begin{quote}\ttfamily\footnotesize\raggedright\noindent
forall(\allowbreak{}[md, tt], is-hom-module(\allowbreak{}md, md, tt) implies is-commutative-ring(\allowbreak{}scal(\allowbreak{}md)) implies forall(\allowbreak{}[n in nn, p in mat(\allowbreak{}n, n, ops(\allowbreak{}md, tt)), u in mat(\allowbreak{}n, 1, vec(\allowbreak{}md))], mat-act(\allowbreak{}md, n, p, u) in mat(\allowbreak{}n, 1, vec(\allowbreak{}md)) and forall(\allowbreak{}[ri in interval(\allowbreak{}1, n)], entry(\allowbreak{}mat-act(\allowbreak{}md, n, p, u), ri, 1) = finsum(\allowbreak{}module-vector-ag(\allowbreak{}md), vnb-lambda(\allowbreak{}gmak\_, interval(\allowbreak{}1, n), (\allowbreak{}entry(\allowbreak{}p, ri, gmak\_))(\allowbreak{}entry(\allowbreak{}u, gmak\_, 1))), interval(\allowbreak{}1, n)))))
\end{quote}
also: \texttt{mat-act-assoc-entry} \\
mentioned by 5 result(s)

\smallskip

\noindent\textbf{\texttt{mat-equiv(a, m, n, c, d)}}\quad\textit{predicate}\index{mat-equiv(a, m, n, c, d)@\texttt{mat-equiv(a, m, n, c, d)}}\\
reads: \$4 and \$5 are equivalent \$2-by-\$3 matrices over \$1 \\
defined by \texttt{mat-equiv} \\
characterized by \texttt{mat-equiv}  \textit{(definitional)}:
\begin{quote}\ttfamily\footnotesize\raggedright\noindent
forall(\allowbreak{}[a, m, n, c, d], mat-equiv(\allowbreak{}a, m, n, c, d) iff forsome(\allowbreak{}[u], is-invertible-mat(\allowbreak{}a, m, u) and forsome(\allowbreak{}[v], is-invertible-mat(\allowbreak{}a, n, v) and d = matmul(\allowbreak{}a, matmul(\allowbreak{}a, u, c), v))))
\end{quote}
also: \texttt{mat-equiv-refl-uba-as-ring} \texttt{mat-equiv-trans-uba-as-ring} \texttt{mat-equiv-left-mult-uba-as-ring} \texttt{mat-equiv-right-mult-uba-as-ring} \texttt{mat-equiv-target-is-mat-uba-as-ring} \texttt{mat-equiv-refl-uba-as-ring-uiba-as-uba} \texttt{mat-equiv-trans-uba-as-ring-uiba-as-uba} \texttt{mat-equiv-left-mult-uba-as-ring-uiba-as-uba} ... (16 in all) \\
mentioned by 54 result(s)

\smallskip

\noindent\textbf{\texttt{mat-lft(s)}}\quad\textit{functoid}\index{mat-lft(s)@\texttt{mat-lft(s)}}\\
reads: the linear fractional transformation L\_\$1 of the 2-by-2 matrix \$1 \\
definition:
\begin{quote}\ttfamily\footnotesize\raggedright\noindent
lft(\allowbreak{}entry(\allowbreak{}s, 1, 1), entry(\allowbreak{}s, 1, 2), entry(\allowbreak{}s, 2, 1), entry(\allowbreak{}s, 2, 2))
\end{quote}
characterized by \texttt{mat-lft-hom}  \textit{(proven)}:
\begin{quote}\ttfamily\footnotesize\raggedright\noindent
forall(\allowbreak{}[s], is-invertible-mat(\allowbreak{}normed-field-as-commutative-ring(\allowbreak{}cc-normed-field), 2, s) implies forall(\allowbreak{}[t], is-invertible-mat(\allowbreak{}normed-field-as-commutative-ring(\allowbreak{}cc-normed-field), 2, t) implies mat-lft(\allowbreak{}matmul(\allowbreak{}normed-field-as-commutative-ring(\allowbreak{}cc-normed-field), s, t)) = compose(\allowbreak{}mat-lft(\allowbreak{}s), mat-lft(\allowbreak{}t))))
\end{quote}
also: \texttt{mat-lft-onto} \texttt{mat-lft-kernel} \texttt{mat-lft-in-lfts} \texttt{mat-lft-literal} \\
mentioned by 13 result(s)

\smallskip

\noindent\textbf{\texttt{mat-map(gmm-h\_, gmm-m\_, gmm-n\_, gmm-a\_)}}\quad\textit{functoid}\index{mat-map(gmm-h_, gmm-m_, gmm-n_, gmm-a_)@\texttt{mat-map(gmm-h\_, gmm-m\_, gmm-n\_, gmm-a\_)}}\\
reads: the matrix \$4 with \$1 applied to every entry \\
definition:
\begin{quote}\ttfamily\footnotesize\raggedright\noindent
matof(\allowbreak{}gmm-m\_, gmm-n\_, vnb-lambda(\allowbreak{}[gmmi\_, gmmj\_], cartesian(\allowbreak{}interval(\allowbreak{}1, gmm-m\_), interval(\allowbreak{}1, gmm-n\_)), gmm-h\_(\allowbreak{}entry(\allowbreak{}gmm-a\_, gmmi\_, gmmj\_))))
\end{quote}
characterized by \texttt{mat-map-type}  \textit{(proven)}:
\begin{quote}\ttfamily\footnotesize\raggedright\noindent
forall(\allowbreak{}[m in nn, n in nn, x, y, h in fun(\allowbreak{}x, y), a in mat(\allowbreak{}m, n, x)], mat-map(\allowbreak{}h, m, n, a) in mat(\allowbreak{}m, n, y))
\end{quote}
also: \texttt{mat-map-entry} \texttt{minor-mat-map} \\
mentioned by 9 result(s)

\smallskip

\noindent\textbf{\texttt{mat-poly-eval(k, n, p, tm)}}\quad\textit{functoid}\index{mat-poly-eval(k, n, p, tm)@\texttt{mat-poly-eval(k, n, p, tm)}}\\
reads: \$3(\$4) \\
definition:
\begin{quote}\ttfamily\footnotesize\raggedright\noindent
sum-ag(\allowbreak{}ring-additive-ag(\allowbreak{}mat-ring(\allowbreak{}normed-field-as-commutative-ring(\allowbreak{}k), n)), vnb-lambda(\allowbreak{}zkvk\_, nn, matscale(\allowbreak{}normed-field-as-commutative-ring(\allowbreak{}k), p(\allowbreak{}zkvk\_), mat-power(\allowbreak{}k, n, tm, zkvk\_))), succ(\allowbreak{}deg(\allowbreak{}normed-field-as-commutative-ring(\allowbreak{}k), p)))
\end{quote}
mentioned by 2 result(s)

\smallskip

\noindent\textbf{\texttt{mat-power(k, n, tm, j)}}\quad\textit{functoid}\index{mat-power(k, n, tm, j)@\texttt{mat-power(k, n, tm, j)}}\\
reads: \$3\textasciicircum{}\$4 \\
definition:
\begin{quote}\ttfamily\footnotesize\raggedright\noindent
mpow(\allowbreak{}ring-multiplicative-monoid(\allowbreak{}mat-ring(\allowbreak{}normed-field-as-commutative-ring(\allowbreak{}k), n)), tm, j)
\end{quote}
characterized by \texttt{cq4-mat-power-type}  \textit{(proven)}:
\begin{quote}\ttfamily\footnotesize\raggedright\noindent
forall(\allowbreak{}[k], is-normed-field(\allowbreak{}k) implies forall(\allowbreak{}[n in nn, tm in mat(\allowbreak{}n, n, carr(\allowbreak{}k)), j in nn], mat-power(\allowbreak{}k, n, tm, j) in mat(\allowbreak{}n, n, carr(\allowbreak{}k))))
\end{quote}
also: \texttt{cb3-mat-power-unfold} \\
mentioned by 17 result(s)

\smallskip

\noindent\textbf{\texttt{mat-ring(a, n)}}\quad\textit{functoid}\index{mat-ring(a, n)@\texttt{mat-ring(a, n)}}\\
definition:
\begin{quote}\ttfamily\footnotesize\raggedright\noindent
[mat(\allowbreak{}n, n, carr(\allowbreak{}a)), vnb-lambda(\allowbreak{}[p, q], cartesian(\allowbreak{}mat(\allowbreak{}n, n, carr(\allowbreak{}a)), mat(\allowbreak{}n, n, carr(\allowbreak{}a))), matadd(\allowbreak{}a, p, q)), vnb-lambda(\allowbreak{}[p, q], cartesian(\allowbreak{}mat(\allowbreak{}n, n, carr(\allowbreak{}a)), mat(\allowbreak{}n, n, carr(\allowbreak{}a))), matmul(\allowbreak{}a, p, q)), vnb-lambda(\allowbreak{}p, mat(\allowbreak{}n, n, carr(\allowbreak{}a)), matneg(\allowbreak{}a, p)), zeromat(\allowbreak{}a, n, n), identmat(\allowbreak{}a, n)]
\end{quote}
characterized by \texttt{mat-ring-is-ring-uba-as-ring}  \textit{(definitional)}:
\begin{quote}\ttfamily\footnotesize\raggedright\noindent
forall(\allowbreak{}[r], is-unital-banach-algebra(\allowbreak{}r) implies forall(\allowbreak{}[n in nn], is-ring(\allowbreak{}mat-ring(\allowbreak{}uba-as-ring(\allowbreak{}r), n))))
\end{quote}
also: \texttt{mat-ring-is-ring-uba-as-ring-uiba-as-uba} \texttt{mat-ring-add} \texttt{mat-ring-mul} \texttt{mat-ring-neg} \texttt{mat-ring-one} \texttt{mat-ring-carr} \texttt{mat-ring-zero} \texttt{mat-ring-is-ring} \\
mentioned by 28 result(s)

\smallskip

\noindent\textbf{\texttt{matact(md, p, u)}}\quad\textit{functoid}\index{matact(md, p, u)@\texttt{matact(md, p, u)}}\\
definition:
\begin{quote}\ttfamily\footnotesize\raggedright\noindent
matof(\allowbreak{}nth(\allowbreak{}1, size(\allowbreak{}p)), nth(\allowbreak{}2, size(\allowbreak{}u)), vnb-lambda(\allowbreak{}[i, c], cartesian(\allowbreak{}interval(\allowbreak{}1, nth(\allowbreak{}1, size(\allowbreak{}p))), interval(\allowbreak{}1, nth(\allowbreak{}2, size(\allowbreak{}u)))), finsum(\allowbreak{}module-vector-ag(\allowbreak{}md), vnb-lambda(\allowbreak{}j, interval(\allowbreak{}1, nth(\allowbreak{}2, size(\allowbreak{}p))), (\allowbreak{}act(\allowbreak{}md))(\allowbreak{}entry(\allowbreak{}p, i, j), entry(\allowbreak{}u, j, c))), interval(\allowbreak{}1, nth(\allowbreak{}2, size(\allowbreak{}p))))))
\end{quote}
characterized by \texttt{matact-type-cna-as-module}  \textit{(definitional)}:
\begin{quote}\ttfamily\footnotesize\raggedright\noindent
forall(\allowbreak{}[r], is-complex-normed-algebra(\allowbreak{}r) implies forall(\allowbreak{}[m, n, q, p, u], p in mat(\allowbreak{}m, n, carr(\allowbreak{}scal(\allowbreak{}r))) implies u in mat(\allowbreak{}n, q, vec(\allowbreak{}r)) implies (\allowbreak{}n = 0 implies m = 0 or q = 0) implies matact(\allowbreak{}cna-as-module(\allowbreak{}r), p, u) in mat(\allowbreak{}m, q, vec(\allowbreak{}r))))
\end{quote}
also: \texttt{matact-type-sip-as-module} \texttt{matact-type-uba-as-module} \texttt{matact-assoc-cna-as-module} \texttt{matact-assoc-sip-as-module} \texttt{matact-assoc-uba-as-module} \texttt{matact-entry-cna-as-module} \texttt{matact-entry-sip-as-module} \texttt{matact-entry-uba-as-module} ... (66 in all) \\
mentioned by 565 result(s)

\smallskip

\noindent\textbf{\texttt{matadd(a, p, q)}}\quad\textit{functoid}\index{matadd(a, p, q)@\texttt{matadd(a, p, q)}}\\
definition:
\begin{quote}\ttfamily\footnotesize\raggedright\noindent
matof(\allowbreak{}nth(\allowbreak{}1, size(\allowbreak{}p)), nth(\allowbreak{}2, size(\allowbreak{}p)), vnb-lambda(\allowbreak{}[i, j], cartesian(\allowbreak{}interval(\allowbreak{}1, nth(\allowbreak{}1, size(\allowbreak{}p))), interval(\allowbreak{}1, nth(\allowbreak{}2, size(\allowbreak{}p)))), (\allowbreak{}add(\allowbreak{}a))(\allowbreak{}entry(\allowbreak{}p, i, j), entry(\allowbreak{}q, i, j))))
\end{quote}
characterized by \texttt{matadd-comm-uba-as-ring}  \textit{(definitional)}:
\begin{quote}\ttfamily\footnotesize\raggedright\noindent
forall(\allowbreak{}[r], is-unital-banach-algebra(\allowbreak{}r) implies forall(\allowbreak{}[m, n, p, q], p in mat(\allowbreak{}m, n, vec(\allowbreak{}r)) implies q in mat(\allowbreak{}m, n, vec(\allowbreak{}r)) implies matadd(\allowbreak{}uba-as-ring(\allowbreak{}r), p, q) = matadd(\allowbreak{}uba-as-ring(\allowbreak{}r), q, p)))
\end{quote}
also: \texttt{matadd-type-uba-as-ring} \texttt{matadd-assoc-uba-as-ring} \texttt{matadd-neg-left-uba-as-ring} \texttt{matadd-neg-right-uba-as-ring} \texttt{matadd-zero-left-uba-as-ring} \texttt{matadd-zero-right-uba-as-ring} \texttt{matadd-comm-uba-as-ring-uiba-as-uba} \texttt{matadd-type-uba-as-ring-uiba-as-uba} ... (25 in all) \\
mentioned by 191 result(s)

\smallskip

\noindent\textbf{\texttt{matmul(a, p, q)}}\quad\textit{functoid}\index{matmul(a, p, q)@\texttt{matmul(a, p, q)}}\\
definition:
\begin{quote}\ttfamily\footnotesize\raggedright\noindent
matof(\allowbreak{}nth(\allowbreak{}1, size(\allowbreak{}p)), nth(\allowbreak{}2, size(\allowbreak{}q)), vnb-lambda(\allowbreak{}[i, k], cartesian(\allowbreak{}interval(\allowbreak{}1, nth(\allowbreak{}1, size(\allowbreak{}p))), interval(\allowbreak{}1, nth(\allowbreak{}2, size(\allowbreak{}q)))), finsum(\allowbreak{}ring-additive-ag(\allowbreak{}a), vnb-lambda(\allowbreak{}j, interval(\allowbreak{}1, nth(\allowbreak{}2, size(\allowbreak{}p))), (\allowbreak{}mul(\allowbreak{}a))(\allowbreak{}entry(\allowbreak{}p, i, j), entry(\allowbreak{}q, j, k))), interval(\allowbreak{}1, nth(\allowbreak{}2, size(\allowbreak{}p))))))
\end{quote}
characterized by \texttt{matmul-type-uba-as-ring}  \textit{(definitional)}:
\begin{quote}\ttfamily\footnotesize\raggedright\noindent
forall(\allowbreak{}[r], is-unital-banach-algebra(\allowbreak{}r) implies forall(\allowbreak{}[m, n, k, p, q], p in mat(\allowbreak{}m, n, vec(\allowbreak{}r)) implies q in mat(\allowbreak{}n, k, vec(\allowbreak{}r)) implies (\allowbreak{}n = 0 implies m = 0 or k = 0) implies matmul(\allowbreak{}uba-as-ring(\allowbreak{}r), p, q) in mat(\allowbreak{}m, k, vec(\allowbreak{}r))))
\end{quote}
also: \texttt{matmul-assoc-uba-as-ring} \texttt{matmul-entry-uba-as-ring} \texttt{matmul-type-uba-as-ring-uiba-as-uba} \texttt{rb4-matmul-entry-single-uba-as-ring} \texttt{matmul-assoc-uba-as-ring-uiba-as-uba} \texttt{matmul-entry-uba-as-ring-uiba-as-uba} \texttt{matmul-left-dist-guarded-uba-as-ring} \texttt{matmul-right-dist-guarded-uba-as-ring} ... (27 in all) \\
mentioned by 665 result(s)

\smallskip

\noindent\textbf{\texttt{matneg(a, p)}}\quad\textit{functoid}\index{matneg(a, p)@\texttt{matneg(a, p)}}\\
definition:
\begin{quote}\ttfamily\footnotesize\raggedright\noindent
matof(\allowbreak{}nth(\allowbreak{}1, size(\allowbreak{}p)), nth(\allowbreak{}2, size(\allowbreak{}p)), vnb-lambda(\allowbreak{}[i, j], cartesian(\allowbreak{}interval(\allowbreak{}1, nth(\allowbreak{}1, size(\allowbreak{}p))), interval(\allowbreak{}1, nth(\allowbreak{}2, size(\allowbreak{}p)))), (\allowbreak{}neg(\allowbreak{}a))(\allowbreak{}entry(\allowbreak{}p, i, j))))
\end{quote}
characterized by \texttt{matneg-type-uba-as-ring}  \textit{(definitional)}:
\begin{quote}\ttfamily\footnotesize\raggedright\noindent
forall(\allowbreak{}[r], is-unital-banach-algebra(\allowbreak{}r) implies forall(\allowbreak{}[m, n, p in mat(\allowbreak{}m, n, vec(\allowbreak{}r))], matneg(\allowbreak{}uba-as-ring(\allowbreak{}r), p) in mat(\allowbreak{}m, n, vec(\allowbreak{}r))))
\end{quote}
also: \texttt{matneg-type-uba-as-ring-uiba-as-uba} \texttt{matneg-type} \texttt{matneg-entry} \\
mentioned by 35 result(s)

\smallskip

\noindent\textbf{\texttt{matof(m, n, g)}}\quad\textit{functoid}\index{matof(m, n, g)@\texttt{matof(m, n, g)}}\\
definition:
\begin{quote}\ttfamily\footnotesize\raggedright\noindent
iota(\allowbreak{}p, p in mat(\allowbreak{}m, n, image(\allowbreak{}g, cartesian(\allowbreak{}interval(\allowbreak{}1, m), interval(\allowbreak{}1, n)))) and forall(\allowbreak{}[i in interval(\allowbreak{}1, m), j in interval(\allowbreak{}1, n)], entry(\allowbreak{}p, i, j) = g(\allowbreak{}i, j)))
\end{quote}
characterized by \texttt{matof-in-mat}  \textit{(proven)}:
\begin{quote}\ttfamily\footnotesize\raggedright\noindent
forall(\allowbreak{}[m, n, x, g], m in nn implies n in nn implies forall(\allowbreak{}[i in interval(\allowbreak{}1, m), j in interval(\allowbreak{}1, n)], g(\allowbreak{}i, j) in x) implies matof(\allowbreak{}m, n, g) in mat(\allowbreak{}m, n, x))
\end{quote}
also: \texttt{entry-of-matof} \texttt{entry-of-matof-guarded} \\
mentioned by 35 result(s)

\smallskip

\noindent\textbf{\texttt{matrix}}\quad\textit{operator}\index{matrix@\texttt{matrix}}\\
a term-forming head \\
characterized by \texttt{matrix-membership}  \textit{(definitional)}:
\begin{quote}\ttfamily\footnotesize\raggedright\noindent
forall(\allowbreak{}[s, m], m in matrix(\allowbreak{}s) iff m in tuples(\allowbreak{}tuples(\allowbreak{}s)) and forall(\allowbreak{}[i], i in nn and 1 <= i and i <= length(\allowbreak{}m) implies forall(\allowbreak{}[j], j in nn and 1 <= j and j <= length(\allowbreak{}m) implies length(\allowbreak{}nth(\allowbreak{}i, m)) = length(\allowbreak{}nth(\allowbreak{}j, m)))))
\end{quote}
also: \texttt{matrix-mono} \texttt{mat-in-matrix} \texttt{matrix-sethood} \\
mentioned by 9 result(s)

\smallskip

\noindent\textbf{\texttt{matscale(a, r, p)}}\quad\textit{functoid}\index{matscale(a, r, p)@\texttt{matscale(a, r, p)}}\\
definition:
\begin{quote}\ttfamily\footnotesize\raggedright\noindent
matof(\allowbreak{}nth(\allowbreak{}1, size(\allowbreak{}p)), nth(\allowbreak{}2, size(\allowbreak{}p)), vnb-lambda(\allowbreak{}[i, j], cartesian(\allowbreak{}interval(\allowbreak{}1, nth(\allowbreak{}1, size(\allowbreak{}p))), interval(\allowbreak{}1, nth(\allowbreak{}2, size(\allowbreak{}p)))), (\allowbreak{}mul(\allowbreak{}a))(\allowbreak{}r, entry(\allowbreak{}p, i, j))))
\end{quote}
characterized by \texttt{matscale-type-uba-as-ring}  \textit{(definitional)}:
\begin{quote}\ttfamily\footnotesize\raggedright\noindent
forall(\allowbreak{}[r], is-unital-banach-algebra(\allowbreak{}r) implies forall(\allowbreak{}[m, n, r\_35061, p], r\_35061 in vec(\allowbreak{}r) implies p in mat(\allowbreak{}m, n, vec(\allowbreak{}r)) implies matscale(\allowbreak{}uba-as-ring(\allowbreak{}r), r\_35061, p) in mat(\allowbreak{}m, n, vec(\allowbreak{}r))))
\end{quote}
also: \texttt{matscale-type-uba-as-ring-uiba-as-uba} \texttt{matscale-one} \texttt{matscale-type} \texttt{matscale-entry} \texttt{matscale-matscale} \texttt{transpose-matscale} \texttt{matmul-matscale-left} \texttt{cq4-matscale-mul-left} ... (12 in all) \\
mentioned by 88 result(s)

\smallskip

\noindent\textbf{\texttt{matunit(a, n, k, l)}}\quad\textit{functoid}\index{matunit(a, n, k, l)@\texttt{matunit(a, n, k, l)}}\\
definition:
\begin{quote}\ttfamily\footnotesize\raggedright\noindent
matof(\allowbreak{}n, n, vnb-lambda(\allowbreak{}[i, j], cartesian(\allowbreak{}interval(\allowbreak{}1, n), interval(\allowbreak{}1, n)), if(\allowbreak{}i = k and j = l, one(\allowbreak{}a), zero(\allowbreak{}a))))
\end{quote}
characterized by \texttt{matunit-type-uba-as-ring}  \textit{(definitional)}:
\begin{quote}\ttfamily\footnotesize\raggedright\noindent
forall(\allowbreak{}[r], is-unital-banach-algebra(\allowbreak{}r) implies forall(\allowbreak{}[n, k, l], n in nn implies matunit(\allowbreak{}uba-as-ring(\allowbreak{}r), n, k, l) in mat(\allowbreak{}n, n, vec(\allowbreak{}r))))
\end{quote}
also: \texttt{matunit-col-shift-uba-as-ring} \texttt{matunit-summand-type-uba-as-ring} \texttt{matunit-type-uba-as-ring-uiba-as-uba} \texttt{matunit-col-shift-uba-as-ring-uiba-as-uba} \texttt{matunit-summand-type-uba-as-ring-uiba-as-uba} \texttt{matunit-type} \texttt{entry-of-matunit} \texttt{matunit-col-shift} ... (11 in all) \\
mentioned by 26 result(s)

\smallskip

\noindent\textbf{\texttt{max}}\quad\textit{operator}\index{max@\texttt{max}}\\
a term-forming head \\
characterized by \texttt{rr-max-def}  \textit{(definitional)}:
\begin{quote}\ttfamily\footnotesize\raggedright\noindent
forall(\allowbreak{}[x in rr, y in rr], (\allowbreak{}y <= x implies max(\allowbreak{}x, y) = x) and (\allowbreak{}not(\allowbreak{}y <= x) implies max(\allowbreak{}x, y) = y))
\end{quote}
also: \texttt{rn-maxabs-succ} \texttt{rr-max-le} \texttt{rr-max-lt} \texttt{ultra-max} \texttt{rr-max-cases} \texttt{nn-max-closed} \texttt{rr-max-closed} \texttt{rr-le-max-left} ... (14 in all) \\
mentioned by 57 result(s)

\smallskip

\noindent\textbf{\texttt{meas(s)}}\quad\textit{accessor}\index{meas(s)@\texttt{meas(s)}}\\
reads: the measure of \$1 \\
a structure slot \\
characterized by \texttt{is-measure-space}  \textit{(definitional)}:
\begin{quote}\ttfamily\footnotesize\raggedright\noindent
forall(\allowbreak{}[s], is-measure-space(\allowbreak{}s) iff length(\allowbreak{}s) = 3 and pts(\allowbreak{}s) in set and sigma(\allowbreak{}s) in power(\allowbreak{}power(\allowbreak{}pts(\allowbreak{}s))) and meas(\allowbreak{}s) in fun(\allowbreak{}sigma(\allowbreak{}s), rr-pos-star) and is-measure(\allowbreak{}pts(\allowbreak{}s), sigma(\allowbreak{}s), meas(\allowbreak{}s)) and s in set)
\end{quote}
also: \texttt{is-hom-measure-space-def} \\
mentioned by 4 result(s)

\smallskip

\noindent\textbf{\texttt{measurable-space}}\quad\textit{structure}\index{measurable-space@\texttt{measurable-space}}\\
predicate \texttt{is-measurable-space}; slots: pts sigma

\smallskip

\noindent\textbf{\texttt{measure-space}}\quad\textit{structure}\index{measure-space@\texttt{measure-space}}\\
predicate \texttt{is-measure-space}; slots: pts sigma meas

\smallskip

\noindent\textbf{\texttt{measure-space-as-measurable-space(s)}}\quad\textit{view}\index{measure-space-as-measurable-space(s)@\texttt{measure-space-as-measurable-space(s)}}\\
a measure-space viewed as a measurable-space \\
characterized by \texttt{measure-space-as-measurable-space-is-measurable-space}  \textit{(definitional)}:
\begin{quote}\ttfamily\footnotesize\raggedright\noindent
forall(\allowbreak{}[r], is-measure-space(\allowbreak{}r) implies is-measurable-space(\allowbreak{}measure-space-as-measurable-space(\allowbreak{}r)))
\end{quote}
also: \texttt{measure-space-as-measurable-space-functorial} \\
mentioned by 2 result(s)

\smallskip

\noindent\textbf{\texttt{meromorphic-on(mru\_, mrf\_)}}\quad\textit{predicate}\index{meromorphic-on(mru_, mrf_)@\texttt{meromorphic-on(mru\_, mrf\_)}}\\
reads: \$2 is meromorphic on \$1 \\
defined by \texttt{meromorphic-on} \\
characterized by \texttt{meromorphic-on}  \textit{(definitional)}:
\begin{quote}\ttfamily\footnotesize\raggedright\noindent
forall(\allowbreak{}[mru\_, mrf\_], meromorphic-on(\allowbreak{}mru\_, mrf\_) iff is-open(\allowbreak{}nf-metric-space(\allowbreak{}cc-normed-field), mru\_) and forsome(\allowbreak{}[mrs\_], mrs\_ subset mru\_ and holomorphic-on(\allowbreak{}difference(\allowbreak{}mru\_, mrs\_), mrf\_) and forall(\allowbreak{}[mrx\_ in mrs\_], is-removable-singularity(\allowbreak{}mrf\_, mrx\_) or forsome(\allowbreak{}[mrn\_], is-pole(\allowbreak{}mrf\_, mrx\_, mrn\_)))))
\end{quote}
also: \texttt{meromorphic-on-sum} \texttt{meromorphic-on-product} \\
mentioned by 4 result(s)

\smallskip

\noindent\textbf{\texttt{method-i-outer(omega, cls, pm)}}\quad\textit{functoid}\index{method-i-outer(omega, cls, pm)@\texttt{method-i-outer(omega, cls, pm)}}\\
reads: the method I outer measure on \$1 of the pre-measure \$3 on \$2 \\
definition:
\begin{quote}\ttfamily\footnotesize\raggedright\noindent
vnb-lambda(\allowbreak{}cve\_, power(\allowbreak{}omega), einf(\allowbreak{}cover-sums(\allowbreak{}omega, cls, pm, cve\_)))
\end{quote}
characterized by \texttt{method-i-outer-type}  \textit{(proven)}:
\begin{quote}\ttfamily\footnotesize\raggedright\noindent
forall(\allowbreak{}[omega, cls, pm], omega in set implies method-i-outer(\allowbreak{}omega, cls, pm) in fun(\allowbreak{}power(\allowbreak{}omega), rr-pos-star))
\end{quote}
also: \texttt{method-i-outer-value} \texttt{method-i-outer-extends} \texttt{method-i-outer-attained} \texttt{method-i-outer-null-iff} \texttt{method-i-outer-measurable} \texttt{method-i-outer-of-measure} \texttt{method-i-outer-finite-approx} \texttt{method-i-outer-is-outer-measure} ... (13 in all) \\
mentioned by 93 result(s)

\smallskip

\noindent\textbf{\texttt{method-ii(mhm\_, mhc\_, mhp\_)}}\quad\textit{functoid}\index{method-ii(mhm_, mhc_, mhp_)@\texttt{method-ii(mhm\_, mhc\_, mhp\_)}}\\
reads: the method-II outer measure of the pre-measure \$3 on \$2 (in \$1) \\
definition:
\begin{quote}\ttfamily\footnotesize\raggedright\noindent
vnb-lambda(\allowbreak{}mhe\_, power(\allowbreak{}pts(\allowbreak{}mhm\_)), esup(\allowbreak{}\{mhv\_ in rr-pos-star: forsome(\allowbreak{}[mhd\_ in rr], pos-rr(\allowbreak{}mhd\_) and mhv\_ = (\allowbreak{}method-ii-delta(\allowbreak{}mhm\_, mhc\_, mhp\_, mhd\_))(\allowbreak{}mhe\_))\}))
\end{quote}
characterized by \texttt{method-ii-in}  \textit{(proven)}:
\begin{quote}\ttfamily\footnotesize\raggedright\noindent
forall(\allowbreak{}[m, cls, pm, es in power(\allowbreak{}pts(\allowbreak{}m))], (\allowbreak{}method-ii(\allowbreak{}m, cls, pm))(\allowbreak{}es) in rr-pos-star)
\end{quote}
also: \texttt{method-ii-seq} \texttt{method-ii-least} \texttt{method-ii-value} \texttt{method-ii-compare} \texttt{method-ii-delta-le} \texttt{method-ii-outer-measure} \texttt{method-ii-metric-outer-measure} \\
mentioned by 12 result(s)

\smallskip

\noindent\textbf{\texttt{method-ii-delta(mhm\_, mhc\_, mhp\_, mhd\_)}}\quad\textit{functoid}\index{method-ii-delta(mhm_, mhc_, mhp_, mhd_)@\texttt{method-ii-delta(mhm\_, mhc\_, mhp\_, mhd\_)}}\\
reads: the method-II outer measure at scale \$4 of the pre-measure \$3 on \$2 (in \$1) \\
definition:
\begin{quote}\ttfamily\footnotesize\raggedright\noindent
method-i-outer(\allowbreak{}pts(\allowbreak{}mhm\_), \{mha\_ in mhc\_: diam(\allowbreak{}mhm\_, mha\_) <= mhd\_\}, restrict(\allowbreak{}mhp\_, \{mha\_ in mhc\_: diam(\allowbreak{}mhm\_, mha\_) <= mhd\_\}))
\end{quote}
characterized by \texttt{method-ii-delta-le}  \textit{(proven)}:
\begin{quote}\ttfamily\footnotesize\raggedright\noindent
forall(\allowbreak{}[m, cls, pm, es, dl], is-premeasure(\allowbreak{}pts(\allowbreak{}m), cls, pm) implies es in power(\allowbreak{}pts(\allowbreak{}m)) implies pos-rr(\allowbreak{}dl) implies (\allowbreak{}method-ii-delta(\allowbreak{}m, cls, pm, dl))(\allowbreak{}es) <= (\allowbreak{}method-ii(\allowbreak{}m, cls, pm))(\allowbreak{}es))
\end{quote}
also: \texttt{method-ii-delta-outer} \texttt{method-ii-delta-value} \texttt{method-ii-delta-antitone} \texttt{method-ii-delta-separated} \texttt{method-ii-delta-scale-below} \texttt{method-ii-delta-subspace-ge} \texttt{method-ii-delta-subspace-le} \\
mentioned by 16 result(s)

\smallskip

\noindent\textbf{\texttt{metric-group}}\quad\textit{structure}\index{metric-group@\texttt{metric-group}}\\
predicate \texttt{is-metric-group}; slots: carr opr iden inv gdist

\smallskip

\noindent\textbf{\texttt{metric-group-group(mgr\_)}}\quad\textit{functoid}\index{metric-group-group(mgr_)@\texttt{metric-group-group(mgr\_)}}\\
reads: the group of \$1 \\
definition:
\begin{quote}\ttfamily\footnotesize\raggedright\noindent
[carr(\allowbreak{}mgr\_), opr(\allowbreak{}mgr\_), iden(\allowbreak{}mgr\_), inv(\allowbreak{}mgr\_)]
\end{quote}
characterized by \texttt{metric-group-group-metric}  \textit{(proven)}:
\begin{quote}\ttfamily\footnotesize\raggedright\noindent
forall(\allowbreak{}[g], is-metric-group(\allowbreak{}g) implies is-group-metric(\allowbreak{}metric-group-ms(\allowbreak{}g), metric-group-group(\allowbreak{}g)))
\end{quote}
also: \texttt{metric-group-group-is-group} \texttt{metric-group-group@dsum-act} \texttt{metric-group-group@enum-fam} \texttt{metric-group-group-functorial} \texttt{metric-group-group@laurent-limc} \texttt{metric-group-group@metric-group-group} \texttt{metric-group-group@abelian-group-as-monoid} \texttt{metric-group-group@normed-ag-as-abelian-group} \\
mentioned by 27 result(s)

\smallskip

\noindent\textbf{\texttt{metric-group-ms(mgr\_)}}\quad\textit{functoid}\index{metric-group-ms(mgr_)@\texttt{metric-group-ms(mgr\_)}}\\
reads: the metric space of \$1 \\
definition:
\begin{quote}\ttfamily\footnotesize\raggedright\noindent
[carr(\allowbreak{}mgr\_), gdist(\allowbreak{}mgr\_)]
\end{quote}
characterized by \texttt{metric-group-ms-functorial}  \textit{(proven)}:
\begin{quote}\ttfamily\footnotesize\raggedright\noindent
forall(\allowbreak{}[a, b, f], is-hom-metric-group(\allowbreak{}a, b, f) implies is-hom-metric-space(\allowbreak{}metric-group-ms(\allowbreak{}a), metric-group-ms(\allowbreak{}b), f))
\end{quote}
also: \texttt{metric-group-ms-is-metric-space} \\
mentioned by 146 result(s)

\smallskip

\noindent\textbf{\texttt{metric-space}}\quad\textit{structure}\index{metric-space@\texttt{metric-space}}\\
predicate \texttt{is-metric-space}; slots: pts dist

\smallskip

\noindent\textbf{\texttt{metric-top(md)}}\quad\textit{functoid}\index{metric-top(md)@\texttt{metric-top(md)}}\\
reads: the metric topology of \$1 \\
definition:
\begin{quote}\ttfamily\footnotesize\raggedright\noindent
[pts(\allowbreak{}md), \{u in power(\allowbreak{}pts(\allowbreak{}md)): is-open(\allowbreak{}md, u)\}]
\end{quote}
characterized by \texttt{metric-top-hom}  \textit{(proven)}:
\begin{quote}\ttfamily\footnotesize\raggedright\noindent
forall(\allowbreak{}[s, t, f], is-metric-space(\allowbreak{}s) implies is-metric-space(\allowbreak{}t) implies is-hom-top-space(\allowbreak{}metric-top(\allowbreak{}s), metric-top(\allowbreak{}t), f) iff is-continuous(\allowbreak{}s, t, f))
\end{quote}
also: \texttt{metric-top-pts} \texttt{cc-metric-top-pts} \texttt{rr-metric-top-pts} \texttt{metric-top-compact} \texttt{metric-top@preimage} \texttt{metric-top@sup-dist} \texttt{metric-top-opens-iff} \texttt{metric-top-functorial} ... (13 in all) \\
mentioned by 108 result(s)

\smallskip

\noindent\textbf{\texttt{metrizable-top-space}}\quad\textit{refinement}\index{metrizable-top-space@\texttt{metrizable-top-space}}\\
predicate \texttt{is-metrizable-top-space}; same shape as top-space

\smallskip

\noindent\textbf{\texttt{mi-add(x, u, v)}}\quad\textit{functoid}\index{mi-add(x, u, v)@\texttt{mi-add(x, u, v)}}\\
reads: the pointwise sum of the multi-indices \$2 and \$3 \\
definition:
\begin{quote}\ttfamily\footnotesize\raggedright\noindent
vnb-lambda(\allowbreak{}mia\_, x, u(\allowbreak{}mia\_) + v(\allowbreak{}mia\_))
\end{quote}
characterized by \texttt{mi-add-fun}  \textit{(proven)}:
\begin{quote}\ttfamily\footnotesize\raggedright\noindent
forall(\allowbreak{}[ixs\_, u\_, v\_], ixs\_ in set implies u\_ in fun(\allowbreak{}ixs\_, nn) implies v\_ in fun(\allowbreak{}ixs\_, nn) implies mi-add(\allowbreak{}ixs\_, u\_, v\_) in fun(\allowbreak{}ixs\_, nn))
\end{quote}
also: \texttt{mi-add-comm} \texttt{mi-add-apply} \texttt{mi-add-assoc} \texttt{mi-add-zero-left} \texttt{mi-add-zero-right} \\
mentioned by 15 result(s)

\smallskip

\noindent\textbf{\texttt{mi-reindex(x, s, b)}}\quad\textit{functoid}\index{mi-reindex(x, s, b)@\texttt{mi-reindex(x, s, b)}}\\
reads: the multi-index \$3 composed with \$2 \\
definition:
\begin{quote}\ttfamily\footnotesize\raggedright\noindent
vnb-lambda(\allowbreak{}mir\_, x, b(\allowbreak{}s(\allowbreak{}mir\_)))
\end{quote}
characterized by \texttt{mi-reindex-add}  \textit{(proven)}:
\begin{quote}\ttfamily\footnotesize\raggedright\noindent
forall(\allowbreak{}[ixs\_, s\_, u\_, v\_], ixs\_ in set implies s\_ in fun(\allowbreak{}ixs\_, ixs\_) implies u\_ in fun(\allowbreak{}ixs\_, nn) implies v\_ in fun(\allowbreak{}ixs\_, nn) implies mi-reindex(\allowbreak{}ixs\_, s\_, mi-add(\allowbreak{}ixs\_, u\_, v\_)) = mi-add(\allowbreak{}ixs\_, mi-reindex(\allowbreak{}ixs\_, s\_, u\_), mi-reindex(\allowbreak{}ixs\_, s\_, v\_)))
\end{quote}
also: \texttt{mi-reindex-fun} \texttt{mi-reindex-zero} \texttt{mi-reindex-apply} \texttt{mi-reindex-inverse} \\
mentioned by 10 result(s)

\smallskip

\noindent\textbf{\texttt{mi-zero(x)}}\quad\textit{functoid}\index{mi-zero(x)@\texttt{mi-zero(x)}}\\
reads: the zero multi-index on \$1 \\
definition:
\begin{quote}\ttfamily\footnotesize\raggedright\noindent
vnb-lambda(\allowbreak{}mia\_, x, 0)
\end{quote}
characterized by \texttt{mi-zero-fun}  \textit{(proven)}:
\begin{quote}\ttfamily\footnotesize\raggedright\noindent
forall(\allowbreak{}[ixs\_ in set], mi-zero(\allowbreak{}ixs\_) in fun(\allowbreak{}ixs\_, nn))
\end{quote}
also: \texttt{mi-zero-apply} \\
mentioned by 11 result(s)

\smallskip

\noindent\textbf{\texttt{min}}\quad\textit{operator}\index{min@\texttt{min}}\\
a term-forming head \\
characterized by \texttt{rr-min-def}  \textit{(definitional)}:
\begin{quote}\ttfamily\footnotesize\raggedright\noindent
forall(\allowbreak{}[x in rr, y in rr], (\allowbreak{}x <= y implies min(\allowbreak{}x, y) = x) and (\allowbreak{}not(\allowbreak{}x <= y) implies min(\allowbreak{}x, y) = y))
\end{quote}
also: \texttt{rr-le-min} \texttt{rr-lt-min} \texttt{rr-min-cases} \texttt{nn-min-closed} \texttt{rr-min-closed} \texttt{rr-min-le-left} \texttt{rr-min-left-le} \texttt{rr-min-le-right} ... (13 in all) \\
mentioned by 40 result(s)

\smallskip

\noindent\textbf{\texttt{minimize!}}\quad\textit{tactic}\index{minimize"!@\texttt{minimize"!}}\\
\texttt{(minimize! '(v1 ... vk) GUARD MEASURE)} -- Choose v1..vk satisfying GUARD with the NN-valued MEASURE as small as possible: lands the witnesses and their minimality, and opens the two obligations any minimisation owes (MEASURE lands in NN; GUARD is satisfiable). \\
declared in \texttt{tactics-help.scm}

\smallskip

\noindent\textbf{\texttt{mink-set(m, cs, x)}}\quad\textit{functoid}\index{mink-set(m, cs, x)@\texttt{mink-set(m, cs, x)}}\\
reads: the positive reals r with \$3 in r times \$2 \\
definition:
\begin{quote}\ttfamily\footnotesize\raggedright\noindent
\{mkr\_ in rr: 0 < mkr\_ and forsome(\allowbreak{}[mky\_ in cs], x = (\allowbreak{}act(\allowbreak{}m))(\allowbreak{}mkr\_, mky\_))\}
\end{quote}
characterized by \texttt{mink-set-glb}  \textit{(proven)}:
\begin{quote}\ttfamily\footnotesize\raggedright\noindent
forall(\allowbreak{}[m, cs, z], is-module(\allowbreak{}m) implies scal(\allowbreak{}m) = normed-field-as-commutative-ring(\allowbreak{}rr-normed-field) implies is-convex-subset(\allowbreak{}m, cs) implies is-absorbing-subset(\allowbreak{}m, cs) implies vzero(\allowbreak{}m) in cs implies z in vec(\allowbreak{}m) implies inf(\allowbreak{}mink-set(\allowbreak{}m, cs, z)) in rr and forall(\allowbreak{}[lbw\_ in mink-set(\allowbreak{}m, cs, z)], inf(\allowbreak{}mink-set(\allowbreak{}m, cs, z)) <= lbw\_) and forall(\allowbreak{}[lbb\_ in rr], forall(\allowbreak{}[lbw\_ in mink-set(\allowbreak{}m, cs, z)], lbb\_ <= lbw\_) implies lbb\_ <= inf(\allowbreak{}mink-set(\allowbreak{}m, cs, z))))
\end{quote}
also: \texttt{mink-set-neg} \texttt{mink-set-sum} \texttt{mink-set-elim} \texttt{mink-set-zero} \texttt{mink-set-intro} \texttt{mink-set-scale} \texttt{mink-set-unfold} \\
mentioned by 26 result(s)

\smallskip

\noindent\textbf{\texttt{minkowski-functional(m, cs)}}\quad\textit{functoid}\index{minkowski-functional(m, cs)@\texttt{minkowski-functional(m, cs)}}\\
reads: the Minkowski functional of \$2 in \$1 \\
definition:
\begin{quote}\ttfamily\footnotesize\raggedright\noindent
vnb-lambda(\allowbreak{}mkx\_, vec(\allowbreak{}m), inf(\allowbreak{}mink-set(\allowbreak{}m, cs, mkx\_)))
\end{quote}
characterized by \texttt{minkowski-functional-is-seminorm}  \textit{(proven)}:
\begin{quote}\ttfamily\footnotesize\raggedright\noindent
forall(\allowbreak{}[m, cs], is-module(\allowbreak{}m) implies scal(\allowbreak{}m) = normed-field-as-commutative-ring(\allowbreak{}rr-normed-field) implies is-convex-subset(\allowbreak{}m, cs) implies is-symmetric-subset(\allowbreak{}m, cs) implies is-absorbing-subset(\allowbreak{}m, cs) implies is-seminorm(\allowbreak{}rr-normed-field, m, minkowski-functional(\allowbreak{}m, cs)))
\end{quote}
mentioned by 5 result(s)

\smallskip

\noindent\textbf{\texttt{minor(s, r, c, n)}}\quad\textit{functoid}\index{minor(s, r, c, n)@\texttt{minor(s, r, c, n)}}\\
definition:
\begin{quote}\ttfamily\footnotesize\raggedright\noindent
matof(\allowbreak{}n, n, vnb-lambda(\allowbreak{}[i, j], cartesian(\allowbreak{}interval(\allowbreak{}1, n), interval(\allowbreak{}1, n)), entry(\allowbreak{}s, if(\allowbreak{}i < r, i, succ(\allowbreak{}i)), if(\allowbreak{}j < c, j, succ(\allowbreak{}j)))))
\end{quote}
characterized by \texttt{minor-type}  \textit{(proven)}:
\begin{quote}\ttfamily\footnotesize\raggedright\noindent
forall(\allowbreak{}[r, s, p, q, n], is-ring(\allowbreak{}r) implies p in nn implies q in nn implies n in nn implies s in mat(\allowbreak{}succ(\allowbreak{}n), succ(\allowbreak{}n), carr(\allowbreak{}r)) implies minor(\allowbreak{}s, p, q, n) in mat(\allowbreak{}n, n, carr(\allowbreak{}r)))
\end{quote}
also: \texttt{minor-entry} \texttt{cp-lin-minor} \texttt{minor-mat-map} \texttt{identmat-minor} \texttt{cai-minor-bound} \texttt{cb3-minor-bound} \texttt{cd6-minor-close} \texttt{dtr-minor-minor} ... (20 in all) \\
mentioned by 60 result(s)

\smallskip

\noindent\textbf{\texttt{module}}\quad\textit{structure}\index{module@\texttt{module}}\\
predicate \texttt{is-module}; slots: scal vec vadd vzero vneg act

\smallskip

\noindent\textbf{\texttt{module-vector-ag(s)}}\quad\textit{view}\index{module-vector-ag(s)@\texttt{module-vector-ag(s)}}\\
a module viewed as an abelian-group \\
characterized by \texttt{module-vector-ag-is-abelian-group}  \textit{(definitional)}:
\begin{quote}\ttfamily\footnotesize\raggedright\noindent
forall(\allowbreak{}[r], is-module(\allowbreak{}r) implies is-abelian-group(\allowbreak{}module-vector-ag(\allowbreak{}r)))
\end{quote}
also: \texttt{module-vector-ag-is-abelian-group-cna-as-module} \texttt{module-vector-ag-is-abelian-group-sip-as-module} \texttt{module-vector-ag-is-abelian-group-uba-as-module} \texttt{module-vector-ag-is-abelian-group-cnvs-as-module} \texttt{abelian-group-as-monoid-is-monoid-module-vector-ag} \texttt{module-vector-ag-is-abelian-group-sip-as-module-csip-as-sip} \texttt{module-vector-ag-is-abelian-group-uba-as-module-uiba-as-uba} \texttt{module-vector-ag-is-abelian-group-normed-vector-space-as-module} ... (21 in all) \\
mentioned by 476 result(s)

\smallskip

\noindent\textbf{\texttt{module-with-form(ipfm\_, ipff\_)}}\quad\textit{functoid}\index{module-with-form(ipfm_, ipff_)@\texttt{module-with-form(ipfm\_, ipff\_)}}\\
reads: \$1 with the inner product \$2 \\
definition:
\begin{quote}\ttfamily\footnotesize\raggedright\noindent
[scal(\allowbreak{}ipfm\_), vec(\allowbreak{}ipfm\_), vadd(\allowbreak{}ipfm\_), vzero(\allowbreak{}ipfm\_), vneg(\allowbreak{}ipfm\_), act(\allowbreak{}ipfm\_), ipff\_]
\end{quote}
characterized by \texttt{module-with-form-ip}  \textit{(proven)}:
\begin{quote}\ttfamily\footnotesize\raggedright\noindent
forall(\allowbreak{}[md], is-module(\allowbreak{}md) implies forall(\allowbreak{}[fq in fun(\allowbreak{}cartesian(\allowbreak{}vec(\allowbreak{}md), vec(\allowbreak{}md)), cc)], ip(\allowbreak{}module-with-form(\allowbreak{}md, fq)) == fq))
\end{quote}
also: \texttt{module-with-form-act} \texttt{module-with-form-vec} \texttt{module-with-form-scal} \texttt{module-with-form-vadd} \texttt{module-with-form-vneg} \texttt{module-with-form-vzero} \texttt{form-plus-space-as-module-with-form} \\
mentioned by 17 result(s)

\smallskip

\noindent\textbf{\texttt{monalg(a, m)}}\quad\textit{functoid}\index{monalg(a, m)@\texttt{monalg(a, m)}}\\
reads: the monoid algebra \$1[\$2] \\
definition:
\begin{quote}\ttfamily\footnotesize\raggedright\noindent
[finsupp(\allowbreak{}a, m), vnb-lambda(\allowbreak{}[f, g], cartesian(\allowbreak{}finsupp(\allowbreak{}a, m), finsupp(\allowbreak{}a, m)), monalg-add(\allowbreak{}a, m, f, g)), vnb-lambda(\allowbreak{}[f, g], cartesian(\allowbreak{}finsupp(\allowbreak{}a, m), finsupp(\allowbreak{}a, m)), monalg-mul(\allowbreak{}a, m, f, g)), vnb-lambda(\allowbreak{}f, finsupp(\allowbreak{}a, m), monalg-neg(\allowbreak{}a, m, f)), monalg-zero(\allowbreak{}a, m), monalg-one(\allowbreak{}a, m)]
\end{quote}
characterized by \texttt{monalg-carr}  \textit{(proven)}:
\begin{quote}\ttfamily\footnotesize\raggedright\noindent
forall(\allowbreak{}[a\_, m\_], is-ring(\allowbreak{}a\_) implies is-monoid(\allowbreak{}m\_) implies carr(\allowbreak{}monalg(\allowbreak{}a\_, m\_)) == finsupp(\allowbreak{}a\_, m\_))
\end{quote}
also: \texttt{monalg-comm} \texttt{monalg-add-op} \texttt{monalg-mul-op} \texttt{monalg-neg-op} \texttt{monalg-is-ring} \texttt{monalg-one-slot} \texttt{monalg-zero-slot} \texttt{monalg-pullback-iso} \\
mentioned by 16 result(s)

\smallskip

\noindent\textbf{\texttt{monalg-add(a, m, f, g)}}\quad\textit{functoid}\index{monalg-add(a, m, f, g)@\texttt{monalg-add(a, m, f, g)}}\\
definition:
\begin{quote}\ttfamily\footnotesize\raggedright\noindent
vnb-lambda(\allowbreak{}x\_, carr(\allowbreak{}m), (\allowbreak{}add(\allowbreak{}a))(\allowbreak{}f(\allowbreak{}x\_), g(\allowbreak{}x\_)))
\end{quote}
characterized by \texttt{monalg-add-op}  \textit{(proven)}:
\begin{quote}\ttfamily\footnotesize\raggedright\noindent
forall(\allowbreak{}[a\_, m\_], is-ring(\allowbreak{}a\_) implies is-monoid(\allowbreak{}m\_) implies add(\allowbreak{}monalg(\allowbreak{}a\_, m\_)) == vnb-lambda(\allowbreak{}[f, g], cartesian(\allowbreak{}finsupp(\allowbreak{}a\_, m\_), finsupp(\allowbreak{}a\_, m\_)), monalg-add(\allowbreak{}a\_, m\_, f, g)))
\end{quote}
also: \texttt{monalg-add-fun} \texttt{monalg-add-comm} \texttt{monalg-add-apply} \texttt{monalg-add-assoc} \texttt{monalg-add-neg-left} \texttt{monalg-add-neg-right} \texttt{monalg-add-zero-left} \texttt{monalg-add-zero-right} \\
mentioned by 32 result(s)

\smallskip

\noindent\textbf{\texttt{monalg-mul(a, m, f, g)}}\quad\textit{functoid}\index{monalg-mul(a, m, f, g)@\texttt{monalg-mul(a, m, f, g)}}\\
definition:
\begin{quote}\ttfamily\footnotesize\raggedright\noindent
vnb-lambda(\allowbreak{}x\_, carr(\allowbreak{}m), finsum(\allowbreak{}ring-additive-ag(\allowbreak{}a), vnb-lambda(\allowbreak{}p, \{p in cartesian(\allowbreak{}supp(\allowbreak{}a, m, f), supp(\allowbreak{}a, m, g)): (\allowbreak{}opr(\allowbreak{}m))(\allowbreak{}nth(\allowbreak{}1, p), nth(\allowbreak{}2, p)) = x\_\}, (\allowbreak{}mul(\allowbreak{}a))(\allowbreak{}f(\allowbreak{}nth(\allowbreak{}1, p)), g(\allowbreak{}nth(\allowbreak{}2, p)))), \{p in cartesian(\allowbreak{}supp(\allowbreak{}a, m, f), supp(\allowbreak{}a, m, g)): (\allowbreak{}opr(\allowbreak{}m))(\allowbreak{}nth(\allowbreak{}1, p), nth(\allowbreak{}2, p)) = x\_\}))
\end{quote}
characterized by \texttt{monalg-mul-op}  \textit{(proven)}:
\begin{quote}\ttfamily\footnotesize\raggedright\noindent
forall(\allowbreak{}[a\_, m\_], is-ring(\allowbreak{}a\_) implies is-monoid(\allowbreak{}m\_) implies mul(\allowbreak{}monalg(\allowbreak{}a\_, m\_)) == vnb-lambda(\allowbreak{}[f, g], cartesian(\allowbreak{}finsupp(\allowbreak{}a\_, m\_), finsupp(\allowbreak{}a\_, m\_)), monalg-mul(\allowbreak{}a\_, m\_, f, g)))
\end{quote}
also: \texttt{monalg-mul-fun} \texttt{monalg-mul-comm} \texttt{monalg-mul-apply} \texttt{monalg-mul-assoc} \texttt{monalg-mul-coef-in} \texttt{monalg-mul-comm-at} \texttt{monalg-mul-assoc-at} \texttt{monalg-mul-supp-in-image} ... (11 in all) \\
mentioned by 122 result(s)

\smallskip

\noindent\textbf{\texttt{monalg-neg(a, m, f)}}\quad\textit{functoid}\index{monalg-neg(a, m, f)@\texttt{monalg-neg(a, m, f)}}\\
definition:
\begin{quote}\ttfamily\footnotesize\raggedright\noindent
vnb-lambda(\allowbreak{}x\_, carr(\allowbreak{}m), (\allowbreak{}neg(\allowbreak{}a))(\allowbreak{}f(\allowbreak{}x\_)))
\end{quote}
characterized by \texttt{monalg-neg-op}  \textit{(proven)}:
\begin{quote}\ttfamily\footnotesize\raggedright\noindent
forall(\allowbreak{}[a\_, m\_], is-ring(\allowbreak{}a\_) implies is-monoid(\allowbreak{}m\_) implies neg(\allowbreak{}monalg(\allowbreak{}a\_, m\_)) == vnb-lambda(\allowbreak{}f, finsupp(\allowbreak{}a\_, m\_), monalg-neg(\allowbreak{}a\_, m\_, f)))
\end{quote}
also: \texttt{monalg-neg-fun} \texttt{monalg-neg-apply} \\
mentioned by 11 result(s)

\smallskip

\noindent\textbf{\texttt{monalg-one(a, m)}}\quad\textit{functoid}\index{monalg-one(a, m)@\texttt{monalg-one(a, m)}}\\
definition:
\begin{quote}\ttfamily\footnotesize\raggedright\noindent
vnb-lambda(\allowbreak{}x\_, carr(\allowbreak{}m), if(\allowbreak{}x\_ = iden(\allowbreak{}m), one(\allowbreak{}a), zero(\allowbreak{}a)))
\end{quote}
characterized by \texttt{monalg-one-in}  \textit{(proven)}:
\begin{quote}\ttfamily\footnotesize\raggedright\noindent
forall(\allowbreak{}[a\_, m\_], is-ring(\allowbreak{}a\_) implies is-monoid(\allowbreak{}m\_) implies monalg-one(\allowbreak{}a\_, m\_) in finsupp(\allowbreak{}a\_, m\_))
\end{quote}
also: \texttt{monalg-one-left} \texttt{monalg-one-slot} \texttt{monalg-one-apply} \texttt{monalg-one-left-at} \texttt{monalg-one-left-fn} \texttt{monalg-one-right-at} \texttt{monalg-one-right-fn} \texttt{monalg-one-supp-subset} \\
mentioned by 20 result(s)

\smallskip

\noindent\textbf{\texttt{monalg-pullback(a, m, ph, f)}}\quad\textit{functoid}\index{monalg-pullback(a, m, ph, f)@\texttt{monalg-pullback(a, m, ph, f)}}\\
reads: the pullback of \$4 along \$3 \\
definition:
\begin{quote}\ttfamily\footnotesize\raggedright\noindent
vnb-lambda(\allowbreak{}mpx\_, carr(\allowbreak{}m), f(\allowbreak{}ph(\allowbreak{}mpx\_)))
\end{quote}
characterized by \texttt{monalg-pullback-add}  \textit{(proven)}:
\begin{quote}\ttfamily\footnotesize\raggedright\noindent
forall(\allowbreak{}[a\_, m\_, ph\_, ps\_, f\_, g\_], is-ring(\allowbreak{}a\_) implies is-monoid(\allowbreak{}m\_) implies ph\_ in fun(\allowbreak{}carr(\allowbreak{}m\_), carr(\allowbreak{}m\_)) implies ps\_ in fun(\allowbreak{}carr(\allowbreak{}m\_), carr(\allowbreak{}m\_)) implies forall(\allowbreak{}[pbu\_ in carr(\allowbreak{}m\_)], ps\_(\allowbreak{}ph\_(\allowbreak{}pbu\_)) = pbu\_) implies f\_ in finsupp(\allowbreak{}a\_, m\_) implies g\_ in finsupp(\allowbreak{}a\_, m\_) implies monalg-pullback(\allowbreak{}a\_, m\_, ph\_, monalg-add(\allowbreak{}a\_, m\_, f\_, g\_)) = monalg-add(\allowbreak{}a\_, m\_, monalg-pullback(\allowbreak{}a\_, m\_, ph\_, f\_), monalg-pullback(\allowbreak{}a\_, m\_, ph\_, g\_)))
\end{quote}
also: \texttt{monalg-pullback-fun} \texttt{monalg-pullback-iso} \texttt{monalg-pullback-mul} \texttt{monalg-pullback-neg} \texttt{monalg-pullback-one} \texttt{monalg-pullback-zero} \texttt{monalg-pullback-apply} \texttt{monalg-pullback-mul-at} ... (19 in all) \\
mentioned by 31 result(s)

\smallskip

\noindent\textbf{\texttt{monalg-zero(a, m)}}\quad\textit{functoid}\index{monalg-zero(a, m)@\texttt{monalg-zero(a, m)}}\\
definition:
\begin{quote}\ttfamily\footnotesize\raggedright\noindent
vnb-lambda(\allowbreak{}x\_, carr(\allowbreak{}m), zero(\allowbreak{}a))
\end{quote}
characterized by \texttt{monalg-zero-in}  \textit{(proven)}:
\begin{quote}\ttfamily\footnotesize\raggedright\noindent
forall(\allowbreak{}[a\_, m\_], is-ring(\allowbreak{}a\_) implies is-monoid(\allowbreak{}m\_) implies monalg-zero(\allowbreak{}a\_, m\_) in finsupp(\allowbreak{}a\_, m\_))
\end{quote}
also: \texttt{monalg-zero-fun} \texttt{monalg-zero-slot} \texttt{monalg-zero-apply} \\
mentioned by 30 result(s)

\smallskip

\noindent\textbf{\texttt{monic-poly(f, a, n)}}\quad\textit{functoid}\index{monic-poly(f, a, n)@\texttt{monic-poly(f, a, n)}}\\
reads: the monic polynomial of degree \$3 with lower coefficients \$2 \\
definition:
\begin{quote}\ttfamily\footnotesize\raggedright\noindent
vnb-lambda(\allowbreak{}gmk\_, nn, if(\allowbreak{}gmk\_ = n, one(\allowbreak{}f), if(\allowbreak{}gmk\_ < n, a(\allowbreak{}succ(\allowbreak{}gmk\_)), zero(\allowbreak{}f))))
\end{quote}
characterized by \texttt{monic-poly-deg-uba-as-ring}  \textit{(definitional)}:
\begin{quote}\ttfamily\footnotesize\raggedright\noindent
forall(\allowbreak{}[r], is-unital-banach-algebra(\allowbreak{}r) implies not(\allowbreak{}vone(\allowbreak{}r) = vzero(\allowbreak{}r)) implies forall(\allowbreak{}[nv\_ in nn, av\_], forall(\allowbreak{}[i\_ in interval(\allowbreak{}1, nv\_)], av\_(\allowbreak{}i\_) in vec(\allowbreak{}r)) implies deg(\allowbreak{}uba-as-ring(\allowbreak{}r), monic-poly(\allowbreak{}uba-as-ring(\allowbreak{}r), av\_, nv\_)) = nv\_ and is-monic(\allowbreak{}uba-as-ring(\allowbreak{}r), monic-poly(\allowbreak{}uba-as-ring(\allowbreak{}r), av\_, nv\_))))
\end{quote}
also: \texttt{monic-poly-low-uba-as-ring} \texttt{monic-poly-top-uba-as-ring} \texttt{monic-poly-high-uba-as-ring} \texttt{monic-poly-in-poly-uba-as-ring} \texttt{monic-poly-congruence-uba-as-ring} \texttt{monic-poly-deg-uba-as-ring-uiba-as-uba} \texttt{monic-poly-low-uba-as-ring-uiba-as-uba} \texttt{monic-poly-top-uba-as-ring-uiba-as-uba} ... (19 in all) \\
mentioned by 74 result(s)

\smallskip

\noindent\textbf{\texttt{monoid}}\quad\textit{structure}\index{monoid@\texttt{monoid}}\\
predicate \texttt{is-monoid}; slots: carr opr iden

\smallskip

\noindent\textbf{\texttt{monomial(a, c, n)}}\quad\textit{functoid}\index{monomial(a, c, n)@\texttt{monomial(a, c, n)}}\\
reads: the monomial \$2 x\textasciicircum{}\$3 over \$1 \\
definition:
\begin{quote}\ttfamily\footnotesize\raggedright\noindent
vnb-lambda(\allowbreak{}x\_, nn, if(\allowbreak{}x\_ = n, c, zero(\allowbreak{}a)))
\end{quote}
characterized by \texttt{monomial-deg-bound-uba-as-ring}  \textit{(definitional)}:
\begin{quote}\ttfamily\footnotesize\raggedright\noindent
forall(\allowbreak{}[r], is-unital-banach-algebra(\allowbreak{}r) implies forall(\allowbreak{}[c\_ in vec(\allowbreak{}r), e\_ in nn], deg-bound(\allowbreak{}uba-as-ring(\allowbreak{}r), monomial(\allowbreak{}uba-as-ring(\allowbreak{}r), c\_, e\_), e\_)))
\end{quote}
also: \texttt{poly-monomial-zero-uba-as-ring} \texttt{poly-one-is-monomial-uba-as-ring} \texttt{monomial-deg-bound-uba-as-ring-uiba-as-uba} \texttt{poly-monomial-zero-uba-as-ring-uiba-as-uba} \texttt{poly-one-is-monomial-uba-as-ring-uiba-as-uba} \texttt{monomial-at} \texttt{monomial-deg} \texttt{monomial-off} ... (21 in all) \\
mentioned by 62 result(s)

\smallskip

\noindent\textbf{\texttt{monotone-class-generated(omega, ce)}}\quad\textit{functoid}\index{monotone-class-generated(omega, ce)@\texttt{monotone-class-generated(omega, ce)}}\\
reads: the monotone class of subsets of \$1 generated by \$2 \\
definition:
\begin{quote}\ttfamily\footnotesize\raggedright\noindent
\{mca\_ in power(\allowbreak{}omega): forall(\allowbreak{}[mcb\_], mcb\_ in power(\allowbreak{}power(\allowbreak{}omega)) and is-monotone-class(\allowbreak{}omega, mcb\_) and ce subset mcb\_ implies mca\_ in mcb\_)\}
\end{quote}
characterized by \texttt{monotone-class-generated-unfold}  \textit{(proven)}:
\begin{quote}\ttfamily\footnotesize\raggedright\noindent
forall(\allowbreak{}[omega, ce], monotone-class-generated(\allowbreak{}omega, ce) == \{mca\_ in power(\allowbreak{}omega): forall(\allowbreak{}[mcb\_], mcb\_ in power(\allowbreak{}power(\allowbreak{}omega)) and is-monotone-class(\allowbreak{}omega, mcb\_) and ce subset mcb\_ implies mca\_ in mcb\_)\})
\end{quote}
also: \texttt{monotone-class-generated-membership} \texttt{monotone-class-generated-is-monotone-class} \\
mentioned by 5 result(s)

\smallskip

\noindent\textbf{\texttt{mp}}\quad\textit{tactic}\index{mp@\texttt{mp}}\\
\texttt{(mp)} -- Close a goal that is an INSTANCE of a universal already in your context.  No arguments: the term is DERIVED by matching, not supplied. \\
declared in \texttt{tactics-help.scm}

\smallskip

\noindent\textbf{\texttt{mpoly(r, n)}}\quad\textit{functoid}\index{mpoly(r, n)@\texttt{mpoly(r, n)}}\\
reads: the polynomial ring \$1[x\_1, ..., x\_\$2] \\
definition:
\begin{quote}\ttfamily\footnotesize\raggedright\noindent
monalg(\allowbreak{}r, multi-index-monoid(\allowbreak{}interval(\allowbreak{}1, n)))
\end{quote}
characterized by \texttt{mpoly-is-ring}  \textit{(proven)}:
\begin{quote}\ttfamily\footnotesize\raggedright\noindent
forall(\allowbreak{}[r\_, n\_], is-ring(\allowbreak{}r\_) implies n\_ in nn implies is-ring(\allowbreak{}mpoly(\allowbreak{}r\_, n\_)))
\end{quote}
also: \texttt{mpoly-is-commutative-ring} \\
mentioned by 3 result(s)

\smallskip

\noindent\textbf{\texttt{mpow}}\quad\textit{defined-fn}\index{mpow@\texttt{mpow}}\\
defined by \texttt{mpow-zero} \texttt{mpow-succ} \\
characterized by \texttt{mpow-succ}  \textit{(definitional)}:
\begin{quote}\ttfamily\footnotesize\raggedright\noindent
forall(\allowbreak{}[m, x, n in nn], mpow(\allowbreak{}m, x, succ(\allowbreak{}n)) == (\allowbreak{}opr(\allowbreak{}m))(\allowbreak{}x, mpow(\allowbreak{}m, x, n)))
\end{quote}
also: \texttt{mpow-zero} \texttt{ring-mpow-one-uba-as-ring} \texttt{mpow-add-abelian-group-as-monoid} \texttt{mpow-one-abelian-group-as-monoid} \texttt{ring-mpow-one-uba-as-ring-uiba-as-uba} \texttt{mpow-add} \texttt{mpow-one} \texttt{mpow-mult} ... (24 in all) \\
mentioned by 103 result(s)

\smallskip

\noindent\textbf{\texttt{mprod}}\quad\textit{functoid}\index{mprod@\texttt{mprod}}\\
reads: the product in \$1 of the first \$3 terms of \$2 \\
defined by \texttt{mprod-zero} \texttt{mprod-succ} \\
characterized by \texttt{mprod-succ}  \textit{(definitional)}:
\begin{quote}\ttfamily\footnotesize\raggedright\noindent
forall(\allowbreak{}[mpm\_, mpu\_, mpn\_ in nn], mprod(\allowbreak{}mpm\_, mpu\_, succ(\allowbreak{}mpn\_)) == (\allowbreak{}opr(\allowbreak{}mpm\_))(\allowbreak{}mprod(\allowbreak{}mpm\_, mpu\_, mpn\_), mpu\_(\allowbreak{}mpn\_)))
\end{quote}
also: \texttt{mprod-zero} \texttt{mprod-in-carr} \texttt{mprod-agree-below} \texttt{mprod-in-carr-ind} \texttt{mprod-shift-split-ind} \\
mentioned by 15 result(s)

\smallskip

\noindent\textbf{\texttt{msg-rank(smt\_, smk\_, smm\_)}}\quad\textit{functoid}\index{msg-rank(smt_, smk_, smm_)@\texttt{msg-rank(smt\_, smk\_, smm\_)}}\\
reads: the rank of the message \$3 over the tokens \$1 and the keys \$2 \\
definition:
\begin{quote}\ttfamily\footnotesize\raggedright\noindent
iota(\allowbreak{}smr\_, smr\_ in nn and smm\_ in msg-stage(\allowbreak{}smt\_, smk\_, smr\_) and forall(\allowbreak{}[smj\_ in nn], smm\_ in msg-stage(\allowbreak{}smt\_, smk\_, smj\_) implies smr\_ <= smj\_))
\end{quote}
characterized by \texttt{msg-rank-encr}  \textit{(proven)}:
\begin{quote}\ttfamily\footnotesize\raggedright\noindent
forall(\allowbreak{}[t, k], is-token-set(\allowbreak{}t) implies forall(\allowbreak{}[g, y], g in msg-set(\allowbreak{}t, k) implies y in k implies succ(\allowbreak{}msg-rank(\allowbreak{}t, k, g)) <= msg-rank(\allowbreak{}t, k, encr(\allowbreak{}g, y))))
\end{quote}
also: \texttt{msg-rank-prop} \texttt{msg-rank-token} \texttt{msg-rank-component} \\
mentioned by 6 result(s)

\smallskip

\noindent\textbf{\texttt{msg-set(smt\_, smk\_)}}\quad\textit{functoid}\index{msg-set(smt_, smk_)@\texttt{msg-set(smt\_, smk\_)}}\\
reads: the messages over the tokens \$1 and the keys \$2 \\
definition:
\begin{quote}\ttfamily\footnotesize\raggedright\noindent
big-union(\allowbreak{}smn\_, nn, msg-stage(\allowbreak{}smt\_, smk\_, smn\_))
\end{quote}
characterized by \texttt{msg-set-elim}  \textit{(proven)}:
\begin{quote}\ttfamily\footnotesize\raggedright\noindent
forall(\allowbreak{}[t, k, m in msg-set(\allowbreak{}t, k)], forsome(\allowbreak{}[n in nn], m in msg-stage(\allowbreak{}t, k, n)))
\end{quote}
also: \texttt{msg-set-intro} \texttt{msg-set-is-set} \texttt{msg-set-unfold} \\
mentioned by 107 result(s)

\smallskip

\noindent\textbf{\texttt{msg-stage}}\quad\textit{functoid}\index{msg-stage@\texttt{msg-stage}}\\
reads: the stage \$3 of the messages over the tokens \$1 and the keys \$2 \\
defined by \texttt{msg-stage-zero} \texttt{msg-stage-succ} \\
characterized by \texttt{msg-stage-succ}  \textit{(definitional)}:
\begin{quote}\ttfamily\footnotesize\raggedright\noindent
forall(\allowbreak{}[smt\_, smk\_, smn\_ in nn], msg-stage(\allowbreak{}smt\_, smk\_, succ(\allowbreak{}smn\_)) == union(\allowbreak{}msg-stage(\allowbreak{}smt\_, smk\_, smn\_), union(\allowbreak{}cartesian(\allowbreak{}pair(\allowbreak{}1, 1), tuples(\allowbreak{}msg-stage(\allowbreak{}smt\_, smk\_, smn\_))), cartesian(\allowbreak{}pair(\allowbreak{}2, 2), msg-stage(\allowbreak{}smt\_, smk\_, smn\_), smk\_))))
\end{quote}
also: \texttt{msg-stage-zero} \texttt{msg-stage-add} \texttt{msg-stage-set} \texttt{msg-stage-encr} \texttt{msg-stage-keep} \texttt{msg-stage-list} \texttt{msg-stage-token} \texttt{msg-stage-encr-inv} ... (11 in all) \\
mentioned by 20 result(s)

\smallskip

\noindent\textbf{\texttt{mu-pinf(prp\_)}}\quad\textit{functoid}\index{mu-pinf(prp_)@\texttt{mu-pinf(prp\_)}}\\
reads: the group of the \$1-power roots of unity \\
definition:
\begin{quote}\ttfamily\footnotesize\raggedright\noindent
\{mpz\_ in cc: forsome(\allowbreak{}[mpk\_ in nn], mpz\_ \textasciicircum{} prp\_ \textasciicircum{} mpk\_ = 1)\}
\end{quote}
mentioned by 2 result(s)

\smallskip

\noindent\textbf{\texttt{mul(s)}}\quad\textit{accessor}\index{mul(s)@\texttt{mul(s)}}\\
reads: the multiplication of \$1 \\
a structure slot \\
characterized by \texttt{rq-mul}  \textit{(definitional)}:
\begin{quote}\ttfamily\footnotesize\raggedright\noindent
forall(\allowbreak{}[r], mul(\allowbreak{}ringoid-quotient-ring(\allowbreak{}r)) == descend2(\allowbreak{}ringoid-setoid(\allowbreak{}r), vnb-lambda(\allowbreak{}[a, b], cartesian(\allowbreak{}carr(\allowbreak{}r), carr(\allowbreak{}r)), class(\allowbreak{}ringoid-setoid(\allowbreak{}r), (\allowbreak{}mul(\allowbreak{}r))(\allowbreak{}a, b)))))
\end{quote}
also: \texttt{qq-ring@mul} \texttt{zz-ring@mul} \texttt{qq-field@mul} \texttt{rr-field@mul} \texttt{ring-mul-assoc} \texttt{ring-mul-left-id} \texttt{ring-mul-right-id} \texttt{cc-normed-field@mul} ... (178 in all) \\
mentioned by 1155 result(s)

\smallskip

\noindent\textbf{\texttt{mul-inv(s)}}\quad\textit{accessor}\index{mul-inv(s)@\texttt{mul-inv(s)}}\\
reads: the multiplicative inverse of \$1 \\
a structure slot \\
characterized by \texttt{qq-field@mul-inv}  \textit{(definitional)}:
\begin{quote}\ttfamily\footnotesize\raggedright\noindent
mul-inv(\allowbreak{}qq-field) == vnb-lambda(\allowbreak{}x\_, difference(\allowbreak{}qq, singleton(\allowbreak{}0)), recip(\allowbreak{}x\_))
\end{quote}
also: \texttt{rr-field@mul-inv} \\
mentioned by 8 result(s)

\smallskip

\noindent\textbf{\texttt{mult-mat(x4md\_, x4n\_, x4u\_, x4f\_, x4xi\_)}}\quad\textit{functoid}\index{mult-mat(x4md_, x4n_, x4u_, x4f_, x4xi_)@\texttt{mult-mat(x4md\_, x4n\_, x4u\_, x4f\_, x4xi\_)}}\\
reads: M(\$5), the matrix of multiplication by \$5 \\
definition:
\begin{quote}\ttfamily\footnotesize\raggedright\noindent
op-mat(\allowbreak{}x4md\_, x4n\_, x4u\_, mult-op(\allowbreak{}x4f\_, x4xi\_))
\end{quote}
characterized by \texttt{ext-mult-mat-add}  \textit{(proven)}:
\begin{quote}\ttfamily\footnotesize\raggedright\noindent
forall(\allowbreak{}[k, f, md, n, u], is-fin-ext(\allowbreak{}k, f, md, n, u) implies forall(\allowbreak{}[xi in carr(\allowbreak{}f), eta in carr(\allowbreak{}f)], mult-mat(\allowbreak{}md, n, u, f, (\allowbreak{}add(\allowbreak{}f))(\allowbreak{}xi, eta)) = matadd(\allowbreak{}normed-field-as-commutative-ring(\allowbreak{}k), mult-mat(\allowbreak{}md, n, u, f, xi), mult-mat(\allowbreak{}md, n, u, f, eta))))
\end{quote}
also: \texttt{ext-mult-mat-mul} \texttt{ext-mult-mat-one} \texttt{ext-mult-mat-spec} \texttt{ext-mult-mat-rigid} \texttt{ext-mult-mat-scalar} \\
mentioned by 14 result(s)

\smallskip

\noindent\textbf{\texttt{mult-op(x4f\_, x4xi\_)}}\quad\textit{functoid}\index{mult-op(x4f_, x4xi_)@\texttt{mult-op(x4f\_, x4xi\_)}}\\
reads: multiplication by \$2 in \$1 \\
definition:
\begin{quote}\ttfamily\footnotesize\raggedright\noindent
vnb-lambda(\allowbreak{}x4p\_, carr(\allowbreak{}x4f\_), (\allowbreak{}mul(\allowbreak{}x4f\_))(\allowbreak{}x4xi\_, x4p\_))
\end{quote}
characterized by \texttt{mult-op-type}  \textit{(proven)}:
\begin{quote}\ttfamily\footnotesize\raggedright\noindent
forall(\allowbreak{}[f], is-ring(\allowbreak{}f) implies forall(\allowbreak{}[xi in carr(\allowbreak{}f)], mult-op(\allowbreak{}f, xi) in fun(\allowbreak{}carr(\allowbreak{}f), carr(\allowbreak{}f))))
\end{quote}
also: \texttt{mult-op-value} \texttt{ext-mult-op-hom} \\
mentioned by 5 result(s)

\smallskip

\noindent\textbf{\texttt{multi-index-monoid(x)}}\quad\textit{functoid}\index{multi-index-monoid(x)@\texttt{multi-index-monoid(x)}}\\
reads: the additive monoid of multi-indices on \$1 \\
definition:
\begin{quote}\ttfamily\footnotesize\raggedright\noindent
[fun(\allowbreak{}x, nn), vnb-lambda(\allowbreak{}[miu\_, miv\_], cartesian(\allowbreak{}fun(\allowbreak{}x, nn), fun(\allowbreak{}x, nn)), mi-add(\allowbreak{}x, miu\_, miv\_)), mi-zero(\allowbreak{}x)]
\end{quote}
characterized by \texttt{multi-index-monoid-opr}  \textit{(proven)}:
\begin{quote}\ttfamily\footnotesize\raggedright\noindent
forall(\allowbreak{}[ixs\_ in set], opr(\allowbreak{}multi-index-monoid(\allowbreak{}ixs\_)) == vnb-lambda(\allowbreak{}[miu\_, miv\_], cartesian(\allowbreak{}fun(\allowbreak{}ixs\_, nn), fun(\allowbreak{}ixs\_, nn)), mi-add(\allowbreak{}ixs\_, miu\_, miv\_)))
\end{quote}
also: \texttt{multi-index-monoid-carr} \texttt{multi-index-monoid-iden} \texttt{multi-index-monoid-is-comm-monoid} \\
mentioned by 8 result(s)

\smallskip


\section*{N}
\addcontentsline{toc}{section}{N}

\noindent\textbf{\texttt{nag-metric-space(nag)}}\quad\textit{functoid}\index{nag-metric-space(nag)@\texttt{nag-metric-space(nag)}}\\
definition:
\begin{quote}\ttfamily\footnotesize\raggedright\noindent
[carr(\allowbreak{}nag), vnb-lambda(\allowbreak{}[u, v], cartesian(\allowbreak{}carr(\allowbreak{}nag), carr(\allowbreak{}nag)), (\allowbreak{}nrm(\allowbreak{}nag))(\allowbreak{}(\allowbreak{}opr(\allowbreak{}nag))(\allowbreak{}u, (\allowbreak{}inv(\allowbreak{}nag))(\allowbreak{}v))))]
\end{quote}
characterized by \texttt{nag-metric-space-is-metric-space-cna-as-normed-ag}  \textit{(definitional)}:
\begin{quote}\ttfamily\footnotesize\raggedright\noindent
forall(\allowbreak{}[r], is-complex-normed-algebra(\allowbreak{}r) implies is-metric-space(\allowbreak{}nag-metric-space(\allowbreak{}cna-as-normed-ag(\allowbreak{}r))))
\end{quote}
also: \texttt{nag-metric-space-is-metric-space-uba-as-normed-ag} \texttt{nag-metric-space-is-metric-space-cnvs-as-normed-ag} \texttt{nag-metric-space-is-metric-space-uba-as-normed-ag-uiba-as-uba} \texttt{nag-metric-space-is-metric-space} \\
mentioned by 111 result(s)

\smallskip

\noindent\textbf{\texttt{names(s)}}\quad\textit{accessor}\index{names(s)@\texttt{names(s)}}\\
reads: the names slot of \$1 \\
a structure slot \\
mentioned by 14 result(s)

\smallskip

\noindent\textbf{\texttt{nbg-closed(mnbw\_)}}\quad\textit{predicate}\index{nbg-closed(mnbw_)@\texttt{nbg-closed(mnbw\_)}}\\
reads: \$1 is closed under the steps of the predicativity walk \\
defined by \texttt{nbg-closed} \\
characterized by \texttt{nbg-closed}  \textit{(definitional)}:
\begin{quote}\ttfamily\footnotesize\raggedright\noindent
forall(\allowbreak{}[mnbw\_], nbg-closed(\allowbreak{}mnbw\_) iff forall(\allowbreak{}[mnbce\_ in mnbw\_, mnbcf\_ in expr-set], nbg-step(\allowbreak{}mnbce\_, mnbcf\_) implies mnbcf\_ in mnbw\_))
\end{quote}
mentioned by 5 result(s)

\smallskip

\noindent\textbf{\texttt{nbg-ok(mnbp\_)}}\quad\textit{predicate}\index{nbg-ok(mnbp_)@\texttt{nbg-ok(mnbp\_)}}\\
reads: \$1 is predicative: every quantifier node in it is confined to sets \\
defined by \texttt{nbg-ok} \\
characterized by \texttt{nbg-ok}  \textit{(definitional)}:
\begin{quote}\ttfamily\footnotesize\raggedright\noindent
forall(\allowbreak{}[mnbp\_], nbg-ok(\allowbreak{}mnbp\_) iff forall(\allowbreak{}[mnboe\_ in nbg-walk(\allowbreak{}mnbp\_)], qnode-ok(\allowbreak{}mnboe\_)))
\end{quote}
also: \texttt{nbg-ok-truth} \texttt{nbg-ok-refuses-unguarded} \\
mentioned by 55 result(s)

\smallskip

\noindent\textbf{\texttt{nbg-step(mnbe\_, mnbf\_)}}\quad\textit{predicate}\index{nbg-step(mnbe_, mnbf_)@\texttt{nbg-step(mnbe\_, mnbf\_)}}\\
reads: \$2 is one step of the predicativity walk below \$1 \\
defined by \texttt{nbg-step} \\
characterized by \texttt{nbg-step}  \textit{(definitional)}:
\begin{quote}\ttfamily\footnotesize\raggedright\noindent
forall(\allowbreak{}[mnbe\_, mnbf\_], nbg-step(\allowbreak{}mnbe\_, mnbf\_) iff forsome(\allowbreak{}[mnb1h\_, mnb1l\_, mnb1i\_], mnb1h\_ in nn and mnb1l\_ in set and mnb1i\_ in nn and mnbe\_ = eapp(\allowbreak{}mnb1h\_, mnb1l\_) and 1 <= mnb1i\_ and mnb1i\_ <= length(\allowbreak{}mnb1l\_) and mnbf\_ = nth(\allowbreak{}mnb1i\_, mnb1l\_)) or forsome(\allowbreak{}[mnb2h\_, mnb2n\_, mnb2a\_, mnb2b\_, mnb2i\_], mnb2h\_ in nn and mnb2n\_ in nn and mnb2a\_ in set and mnb2b\_ in expr-set and mnb2i\_ in nn and mnbe\_ = ebind(\allowbreak{}mnb2h\_, mnb2n\_, mnb2a\_, mnb2b\_) and 1 <= mnb2i\_ and mnb2i\_ <= length(\allowbreak{}mnb2a\_) and mnbf\_ = nth(\allowbreak{}mnb2i\_, mnb2a\_)) or forsome(\allowbreak{}[mnb3h\_, mnb3n\_, mnb3a\_], mnb3h\_ in nn and mnb3n\_ in nn and mnb3a\_ in set and mnbe\_ = ebind(\allowbreak{}mnb3h\_, mnb3n\_, mnb3a\_, mnbf\_)) or forsome(\allowbreak{}[mnb4h\_, mnb4l\_], mnb4h\_ in nn and mnb4l\_ in set and mnbe\_ = eapp(\allowbreak{}mnb4h\_, mnb4l\_) and fun-def(\allowbreak{}mnb4h\_, length(\allowbreak{}mnb4l\_), mnbf\_)) or forsome(\allowbreak{}[mnb5l\_, mnb5k\_], mnb5l\_ in set and mnb5k\_ in nn and mnbe\_ = eapp(\allowbreak{}10, mnb5l\_) and nth(\allowbreak{}1, mnb5l\_) = econst(\allowbreak{}mnb5k\_) and pred-def(\allowbreak{}mnb5k\_, length(\allowbreak{}mnb5l\_), mnbf\_)))
\end{quote}
mentioned by 4 result(s)

\smallskip

\noindent\textbf{\texttt{nbg-walk(mnbp\_)}}\quad\textit{functoid}\index{nbg-walk(mnbp_)@\texttt{nbg-walk(mnbp\_)}}\\
reads: the expressions the predicativity walk reaches from \$1 \\
definition:
\begin{quote}\ttfamily\footnotesize\raggedright\noindent
\{mnbz\_ in expr-set: forall(\allowbreak{}[mnbv\_ in power(\allowbreak{}expr-set)], mnbp\_ in mnbv\_ implies nbg-closed(\allowbreak{}mnbv\_) implies mnbz\_ in mnbv\_)\}
\end{quote}
characterized by \texttt{nbg-walk-base}  \textit{(proven)}:
\begin{quote}\ttfamily\footnotesize\raggedright\noindent
forall(\allowbreak{}[p in expr-set], p in nbg-walk(\allowbreak{}p))
\end{quote}
also: \texttt{nbg-walk-least} \texttt{nbg-walk-unfold} \\
mentioned by 6 result(s)

\smallskip

\noindent\textbf{\texttt{neg(s)}}\quad\textit{accessor}\index{neg(s)@\texttt{neg(s)}}\\
reads: the negation of \$1 \\
a structure slot \\
characterized by \texttt{rq-neg}  \textit{(definitional)}:
\begin{quote}\ttfamily\footnotesize\raggedright\noindent
forall(\allowbreak{}[r], neg(\allowbreak{}ringoid-quotient-ring(\allowbreak{}r)) == descend(\allowbreak{}ringoid-setoid(\allowbreak{}r), vnb-lambda(\allowbreak{}a, carr(\allowbreak{}r), class(\allowbreak{}ringoid-setoid(\allowbreak{}r), (\allowbreak{}neg(\allowbreak{}r))(\allowbreak{}a)))))
\end{quote}
also: \texttt{qq-ring@neg} \texttt{zz-ring@neg} \texttt{qq-field@neg} \texttt{rr-field@neg} \texttt{ringoid-ideal-neg} \texttt{cc-normed-field@neg} \texttt{rr-normed-field@neg} \texttt{poly-neg-at-uba-as-ring} ... (120 in all) \\
mentioned by 655 result(s)

\smallskip

\noindent\textbf{\texttt{neg-part(omega, f)}}\quad\textit{functoid}\index{neg-part(omega, f)@\texttt{neg-part(omega, f)}}\\
reads: the negative part of \$2 on \$1 \\
definition:
\begin{quote}\ttfamily\footnotesize\raggedright\noindent
vnb-lambda(\allowbreak{}npx\_, omega, if(\allowbreak{}0 <= f(\allowbreak{}npx\_), 0, -f(\allowbreak{}npx\_)))
\end{quote}
characterized by \texttt{neg-part-fun}  \textit{(proven)}:
\begin{quote}\ttfamily\footnotesize\raggedright\noindent
forall(\allowbreak{}[omega, f in fun(\allowbreak{}omega, rr)], neg-part(\allowbreak{}omega, f) in fun(\allowbreak{}omega, rr-pos-star))
\end{quote}
also: \texttt{neg-part-value} \texttt{neg-part-unfold} \texttt{pfw-neg-part-comp} \texttt{pos-neg-parts-ptwise} \texttt{neg-part-measurable-of-diff} \\
mentioned by 31 result(s)

\smallskip

\noindent\textbf{\texttt{neg-rr}}\quad\textit{predicate}\index{neg-rr@\texttt{neg-rr}}\\
reads: a negative real \\
a predicate

\smallskip

\noindent\textbf{\texttt{neg-variation(ca, nu)}}\quad\textit{functoid}\index{neg-variation(ca, nu)@\texttt{neg-variation(ca, nu)}}\\
reads: the negative variation of \$2 \\
definition:
\begin{quote}\ttfamily\footnotesize\raggedright\noindent
pos-variation(\allowbreak{}ca, signed-neg(\allowbreak{}ca, nu))
\end{quote}
characterized by \texttt{neg-variation-finite}  \textit{(proven)}:
\begin{quote}\ttfamily\footnotesize\raggedright\noindent
forall(\allowbreak{}[omega, ca, nu], is-signed-measure(\allowbreak{}omega, ca, nu) implies forall(\allowbreak{}[a\_ in ca], not(\allowbreak{}nu(\allowbreak{}a\_) = neg-inf)) implies forall(\allowbreak{}[a\_ in ca], (\allowbreak{}neg-variation(\allowbreak{}ca, nu))(\allowbreak{}a\_) in rr))
\end{quote}
also: \texttt{neg-variation-unfold} \texttt{neg-variation-is-measure} \\
mentioned by 14 result(s)

\smallskip

\noindent\textbf{\texttt{net-converges-to(t, d, x, a)}}\quad\textit{predicate}\index{net-converges-to(t, d, x, a)@\texttt{net-converges-to(t, d, x, a)}}\\
reads: the net \$3 on \$2 converges to \$4 in \$1 \\
defined by \texttt{net-converges-to} \\
characterized by \texttt{net-converges-to}  \textit{(definitional)}:
\begin{quote}\ttfamily\footnotesize\raggedright\noindent
forall(\allowbreak{}[t, d, x, a], net-converges-to(\allowbreak{}t, d, x, a) iff is-top-space(\allowbreak{}t) and is-directed-set(\allowbreak{}d) and x in fun(\allowbreak{}pts(\allowbreak{}d), pts(\allowbreak{}t)) and a in pts(\allowbreak{}t) and forall(\allowbreak{}[ncu\_ in opens(\allowbreak{}t)], a in ncu\_ implies is-eventually-in(\allowbreak{}d, x, ncu\_)))
\end{quote}
mentioned by 21 result(s)

\smallskip

\noindent\textbf{\texttt{net-unif-conv-on(d, p, f, k)}}\quad\textit{predicate}\index{net-unif-conv-on(d, p, f, k)@\texttt{net-unif-conv-on(d, p, f, k)}}\\
reads: the net \$2 on \$1 converges to \$3 uniformly on \$4 \\
defined by \texttt{net-unif-conv-on} \\
characterized by \texttt{net-unif-conv-on}  \textit{(definitional)}:
\begin{quote}\ttfamily\footnotesize\raggedright\noindent
forall(\allowbreak{}[d, p, f, k], net-unif-conv-on(\allowbreak{}d, p, f, k) iff forall(\allowbreak{}[uce\_], pos-rr(\allowbreak{}uce\_) implies forsome(\allowbreak{}[uci\_ in pts(\allowbreak{}d)], forall(\allowbreak{}[ucj\_ in pts(\allowbreak{}d)], [uci\_, ucj\_] in dleq(\allowbreak{}d) implies forall(\allowbreak{}[ucy\_ in k], abs(\allowbreak{}(\allowbreak{}p(\allowbreak{}ucj\_))(\allowbreak{}ucy\_) - f(\allowbreak{}ucy\_)) < uce\_)))))
\end{quote}
mentioned by 4 result(s)

\smallskip

\noindent\textbf{\texttt{neumann-sum}}\quad\textit{functoid}\index{neumann-sum@\texttt{neumann-sum}}\\
reads: the sum of the powers of \$2 below \$3 in \$1 \\
defined by \texttt{neumann-sum-zero} \texttt{neumann-sum-succ} \\
characterized by \texttt{neumann-sum-succ}  \textit{(definitional)}:
\begin{quote}\ttfamily\footnotesize\raggedright\noindent
forall(\allowbreak{}[bna\_, bnx\_, bnk\_ in nn], neumann-sum(\allowbreak{}bna\_, bnx\_, succ(\allowbreak{}bnk\_)) == (\allowbreak{}vadd(\allowbreak{}bna\_))(\allowbreak{}neumann-sum(\allowbreak{}bna\_, bnx\_, bnk\_), el-pow(\allowbreak{}bna\_, bnx\_, bnk\_)))
\end{quote}
also: \texttt{neumann-sum-zero} \\
mentioned by 15 result(s)

\smallskip

\noindent\textbf{\texttt{nf-metric-space(nf)}}\quad\textit{functoid}\index{nf-metric-space(nf)@\texttt{nf-metric-space(nf)}}\\
definition:
\begin{quote}\ttfamily\footnotesize\raggedright\noindent
[carr(\allowbreak{}nf), vnb-lambda(\allowbreak{}[x, y], cartesian(\allowbreak{}carr(\allowbreak{}nf), carr(\allowbreak{}nf)), (\allowbreak{}fnrm(\allowbreak{}nf))(\allowbreak{}(\allowbreak{}add(\allowbreak{}nf))(\allowbreak{}x, (\allowbreak{}neg(\allowbreak{}nf))(\allowbreak{}y))))]
\end{quote}
characterized by \texttt{nf-metric-space-unfold}  \textit{(proven)}:
\begin{quote}\ttfamily\footnotesize\raggedright\noindent
forall(\allowbreak{}[nf], nf-metric-space(\allowbreak{}nf) == [carr(\allowbreak{}nf), vnb-lambda(\allowbreak{}[snx\_, sny\_], cartesian(\allowbreak{}carr(\allowbreak{}nf), carr(\allowbreak{}nf)), (\allowbreak{}fnrm(\allowbreak{}nf))(\allowbreak{}(\allowbreak{}add(\allowbreak{}nf))(\allowbreak{}snx\_, (\allowbreak{}neg(\allowbreak{}nf))(\allowbreak{}sny\_))))])
\end{quote}
also: \texttt{nf-metric-space-is-metric-space} \\
mentioned by 625 result(s)

\smallskip

\noindent\textbf{\texttt{nf-poly-over(ndsk\_, ndsf\_)}}\quad\textit{functoid}\index{nf-poly-over(ndsk_, ndsf_)@\texttt{nf-poly-over(ndsk\_, ndsf\_)}}\\
reads: the polynomials over \$1 with coefficients in \$2 \\
definition:
\begin{quote}\ttfamily\footnotesize\raggedright\noindent
\{ndsp\_ in carr(\allowbreak{}poly(\allowbreak{}normed-field-as-commutative-ring(\allowbreak{}ndsk\_))): forall(\allowbreak{}[ndsi\_ in nn], ndsp\_(\allowbreak{}ndsi\_) in ndsf\_)\}
\end{quote}

\smallskip

\noindent\textbf{\texttt{nf-rat-at(ndsk\_, ndsf\_, ndsw\_)}}\quad\textit{functoid}\index{nf-rat-at(ndsk_, ndsf_, ndsw_)@\texttt{nf-rat-at(ndsk\_, ndsf\_, ndsw\_)}}\\
reads: the quotients of values at \$3 of polynomials over \$2 in \$1 \\
definition:
\begin{quote}\ttfamily\footnotesize\raggedright\noindent
\{ndsx\_ in carr(\allowbreak{}ndsk\_): forsome(\allowbreak{}[ndsp\_ in nf-poly-over(\allowbreak{}ndsk\_, ndsf\_), ndsq\_ in nf-poly-over(\allowbreak{}ndsk\_, ndsf\_)], not(\allowbreak{}poly-eval(\allowbreak{}normed-field-as-commutative-ring(\allowbreak{}ndsk\_), ndsq\_, ndsw\_) = zero(\allowbreak{}ndsk\_)) and (\allowbreak{}mul(\allowbreak{}ndsk\_))(\allowbreak{}ndsx\_, poly-eval(\allowbreak{}normed-field-as-commutative-ring(\allowbreak{}ndsk\_), ndsq\_, ndsw\_)) = poly-eval(\allowbreak{}normed-field-as-commutative-ring(\allowbreak{}ndsk\_), ndsp\_, ndsw\_))\}
\end{quote}
characterized by \texttt{nf-rat-at-least}  \textit{(proven)}:
\begin{quote}\ttfamily\footnotesize\raggedright\noindent
forall(\allowbreak{}[k], is-normed-field(\allowbreak{}k) implies forall(\allowbreak{}[fsub], is-nf-subfield(\allowbreak{}k, fsub) implies forall(\allowbreak{}[s], is-nf-subfield(\allowbreak{}k, s) implies forall(\allowbreak{}[pcz in fsub], pcz in s) implies forall(\allowbreak{}[pcw in s, pcx in nf-rat-at(\allowbreak{}k, fsub, pcw)], pcx in s))))
\end{quote}
also: \texttt{nf-rat-at-contains} \texttt{nf-rat-at-subfield} \\
mentioned by 4 result(s)

\smallskip

\noindent\textbf{\texttt{nf-subfield-gen(bzk\_, bzg\_)}}\quad\textit{functoid}\index{nf-subfield-gen(bzk_, bzg_)@\texttt{nf-subfield-gen(bzk\_, bzg\_)}}\\
reads: the subfield of \$1 generated by \$2 \\
definition:
\begin{quote}\ttfamily\footnotesize\raggedright\noindent
\{bzz\_ in carr(\allowbreak{}bzk\_): forall(\allowbreak{}[bzs\_ in power(\allowbreak{}carr(\allowbreak{}bzk\_))], is-nf-subfield(\allowbreak{}bzk\_, bzs\_) implies bzg\_ subset bzs\_ implies bzz\_ in bzs\_)\}
\end{quote}
characterized by \texttt{nf-subfield-gen-sub}  \textit{(proven)}:
\begin{quote}\ttfamily\footnotesize\raggedright\noindent
forall(\allowbreak{}[k], is-normed-field(\allowbreak{}k) implies forall(\allowbreak{}[x], nf-subfield-gen(\allowbreak{}k, x) subset carr(\allowbreak{}k)))
\end{quote}
also: \texttt{nf-subfield-gen-least} \texttt{nf-subfield-gen-unfold} \texttt{nf-subfield-gen-contains} \texttt{nf-subfield-gen-is-subfield} \\
mentioned by 9 result(s)

\smallskip

\noindent\textbf{\texttt{nf-subfield-module(fmk\_, fms\_)}}\quad\textit{functoid}\index{nf-subfield-module(fmk_, fms_)@\texttt{nf-subfield-module(fmk\_, fms\_)}}\\
reads: \$1 as a module over its subfield \$2 \\
definition:
\begin{quote}\ttfamily\footnotesize\raggedright\noindent
[normed-field-as-commutative-ring(\allowbreak{}sub-nf(\allowbreak{}fmk\_, fms\_)), carr(\allowbreak{}normed-field-as-commutative-ring(\allowbreak{}fmk\_)), add(\allowbreak{}normed-field-as-commutative-ring(\allowbreak{}fmk\_)), zero(\allowbreak{}normed-field-as-commutative-ring(\allowbreak{}fmk\_)), neg(\allowbreak{}normed-field-as-commutative-ring(\allowbreak{}fmk\_)), vnb-lambda(\allowbreak{}[fmr\_, fmx\_], cartesian(\allowbreak{}fms\_, carr(\allowbreak{}normed-field-as-commutative-ring(\allowbreak{}fmk\_))), (\allowbreak{}mul(\allowbreak{}normed-field-as-commutative-ring(\allowbreak{}fmk\_)))(\allowbreak{}fmr\_, fmx\_))]
\end{quote}
characterized by \texttt{nf-subfield-module-act}  \textit{(proven)}:
\begin{quote}\ttfamily\footnotesize\raggedright\noindent
forall(\allowbreak{}[k], is-normed-field(\allowbreak{}k) implies forall(\allowbreak{}[s], is-nf-subfield(\allowbreak{}k, s) implies act(\allowbreak{}nf-subfield-module(\allowbreak{}k, s)) == vnb-lambda(\allowbreak{}[fmr\_, fmx\_], cartesian(\allowbreak{}s, carr(\allowbreak{}normed-field-as-commutative-ring(\allowbreak{}k))), (\allowbreak{}mul(\allowbreak{}normed-field-as-commutative-ring(\allowbreak{}k)))(\allowbreak{}fmr\_, fmx\_))))
\end{quote}
also: \texttt{nf-subfield-module-vec} \texttt{nf-subfield-module-scal} \texttt{nf-subfield-module-vadd} \texttt{nf-subfield-module-vneg} \texttt{nf-subfield-module-vzero} \texttt{nf-subfield-module-is-module} \texttt{nf-subfield-module-vector-ag} \\
mentioned by 19 result(s)

\smallskip

\noindent\textbf{\texttt{ni}}\quad\textit{tactic}\index{ni@\texttt{ni}}\\
\texttt{(ni)} -- Natural-number induction on the goal's leading FORALL over NN. \\
declared in \texttt{tactics-help.scm}

\smallskip

\noindent\textbf{\texttt{nn-add-monoid}}\quad\textit{instance}\index{nn-add-monoid@\texttt{nn-add-monoid}}\\
an instance of comm-monoid; predicate \texttt{is-nn-add-monoid}; same shape as comm-monoid

\smallskip

\noindent\textbf{\texttt{nn-directed}}\quad\textit{constant}\index{nn-directed@\texttt{nn-directed}}\\
reads: the natural numbers as a directed set \\
defined by \texttt{nn-directed-def}

\smallskip

\noindent\textbf{\texttt{nn-enum(s)}}\quad\textit{functoid}\index{nn-enum(s)@\texttt{nn-enum(s)}}\\
definition:
\begin{quote}\ttfamily\footnotesize\raggedright\noindent
choice(\allowbreak{}\{f in fun(\allowbreak{}nn, s): forall(\allowbreak{}[m in nn, n in nn], m < n implies f(\allowbreak{}m) < f(\allowbreak{}n))\})
\end{quote}
characterized by \texttt{nn-enum-spec}  \textit{(proven)}:
\begin{quote}\ttfamily\footnotesize\raggedright\noindent
forall(\allowbreak{}[s in inf-subsets(\allowbreak{}nn)], nn-enum(\allowbreak{}s) in fun(\allowbreak{}nn, s) and forall(\allowbreak{}[m in nn, n in nn], m < n implies (\allowbreak{}nn-enum(\allowbreak{}s))(\allowbreak{}m) < (\allowbreak{}nn-enum(\allowbreak{}s))(\allowbreak{}n)))
\end{quote}
mentioned by 1 result(s)

\smallskip

\noindent\textbf{\texttt{nn-minus}}\quad\textit{defined-fn}\index{nn-minus@\texttt{nn-minus}}\\
defined by \texttt{nn-minus-def} \\
characterized by \texttt{nn-minus-def}  \textit{(definitional)}:
\begin{quote}\ttfamily\footnotesize\raggedright\noindent
forall(\allowbreak{}[k in nn, l in nn], nn-minus(\allowbreak{}k, l) = if(\allowbreak{}l <= k, k - l, 0))
\end{quote}
also: \texttt{nn-minus-1-inj} \texttt{nn-minus-in-nn} \texttt{nn-minus-succ-1} \texttt{succ-nn-minus-1} \\
mentioned by 30 result(s)

\smallskip

\noindent\textbf{\texttt{nnfst(n\_)}}\quad\textit{functoid}\index{nnfst(n_)@\texttt{nnfst(n\_)}}\\
reads: the first component of \$1 \\
definition:
\begin{quote}\ttfamily\footnotesize\raggedright\noindent
iota(\allowbreak{}a\_, a\_ in nn and forsome(\allowbreak{}[b\_ in nn], nnpair(\allowbreak{}a\_, b\_) = n\_))
\end{quote}
characterized by \texttt{nnfst-type}  \textit{(proven)}:
\begin{quote}\ttfamily\footnotesize\raggedright\noindent
forall(\allowbreak{}[n\_ in nn], nnfst(\allowbreak{}n\_) in nn)
\end{quote}
also: \texttt{nnfst-nnpair} \\
mentioned by 35 result(s)

\smallskip

\noindent\textbf{\texttt{nnpair(i\_, j\_)}}\quad\textit{functoid}\index{nnpair(i_, j_)@\texttt{nnpair(i\_, j\_)}}\\
reads: the Cantor code of (\$1, \$2) \\
definition:
\begin{quote}\ttfamily\footnotesize\raggedright\noindent
trinum(\allowbreak{}i\_ + j\_) + j\_
\end{quote}
characterized by \texttt{nnpair-inj}  \textit{(proven)}:
\begin{quote}\ttfamily\footnotesize\raggedright\noindent
forall(\allowbreak{}[i\_, j\_, ii\_, jj\_], i\_ in nn implies j\_ in nn implies ii\_ in nn implies jj\_ in nn implies nnpair(\allowbreak{}i\_, j\_) = nnpair(\allowbreak{}ii\_, jj\_) implies i\_ = ii\_ and j\_ = jj\_)
\end{quote}
also: \texttt{nnpair-onto} \texttt{nnpair-type} \texttt{nnfst-nnpair} \texttt{nnsnd-nnpair} \texttt{nnpair-lt-right} \\
mentioned by 16 result(s)

\smallskip

\noindent\textbf{\texttt{nnreal}}\quad\textit{constant}\index{nnreal@\texttt{nnreal}}\\
reads: [0, +oo) \\
defined by \texttt{nnreal-def}

\smallskip

\noindent\textbf{\texttt{nnsnd(n\_)}}\quad\textit{functoid}\index{nnsnd(n_)@\texttt{nnsnd(n\_)}}\\
reads: the second component of \$1 \\
definition:
\begin{quote}\ttfamily\footnotesize\raggedright\noindent
iota(\allowbreak{}b\_, b\_ in nn and forsome(\allowbreak{}[a\_ in nn], nnpair(\allowbreak{}a\_, b\_) = n\_))
\end{quote}
characterized by \texttt{nnsnd-type}  \textit{(proven)}:
\begin{quote}\ttfamily\footnotesize\raggedright\noindent
forall(\allowbreak{}[n\_ in nn], nnsnd(\allowbreak{}n\_) in nn)
\end{quote}
also: \texttt{nnsnd-nnpair} \\
mentioned by 26 result(s)

\smallskip

\noindent\textbf{\texttt{node-content(ssp\_, ssn\_)}}\quad\textit{functoid}\index{node-content(ssp_, ssn_)@\texttt{node-content(ssp\_, ssn\_)}}\\
reads: the content of the node \$2 of \$1 \\
definition:
\begin{quote}\ttfamily\footnotesize\raggedright\noindent
nth(\allowbreak{}2, nth(\allowbreak{}nth(\allowbreak{}1, ssn\_), (\allowbreak{}strand-map(\allowbreak{}ssp\_))(\allowbreak{}nth(\allowbreak{}2, ssn\_))))
\end{quote}
characterized by \texttt{node-content-eq}  \textit{(proven)}:
\begin{quote}\ttfamily\footnotesize\raggedright\noindent
forall(\allowbreak{}[spa, n, k, d], nth(\allowbreak{}1, n) = k implies nth(\allowbreak{}k, (\allowbreak{}strand-map(\allowbreak{}spa))(\allowbreak{}nth(\allowbreak{}2, n))) = d implies node-content(\allowbreak{}spa, n) == nth(\allowbreak{}2, d) and node-dir(\allowbreak{}spa, n) == nth(\allowbreak{}1, d))
\end{quote}
also: \texttt{text-node-content} \texttt{node-content-unfold} \texttt{node-content-in-msgs} \\
mentioned by 76 result(s)

\smallskip

\noindent\textbf{\texttt{node-dir(ssp\_, ssn\_)}}\quad\textit{functoid}\index{node-dir(ssp_, ssn_)@\texttt{node-dir(ssp\_, ssn\_)}}\\
reads: the direction of the node \$2 of \$1 \\
definition:
\begin{quote}\ttfamily\footnotesize\raggedright\noindent
nth(\allowbreak{}1, nth(\allowbreak{}nth(\allowbreak{}1, ssn\_), (\allowbreak{}strand-map(\allowbreak{}ssp\_))(\allowbreak{}nth(\allowbreak{}2, ssn\_))))
\end{quote}
characterized by \texttt{node-dir-cases}  \textit{(proven)}:
\begin{quote}\ttfamily\footnotesize\raggedright\noindent
forall(\allowbreak{}[spa, n], is-strand-space(\allowbreak{}spa) implies n in nodes(\allowbreak{}spa) implies node-dir(\allowbreak{}spa, n) = sent or node-dir(\allowbreak{}spa, n) = received)
\end{quote}
also: \texttt{node-dir-unfold} \\
mentioned by 10 result(s)

\smallskip

\noindent\textbf{\texttt{nodes(ssp\_)}}\quad\textit{functoid}\index{nodes(ssp_)@\texttt{nodes(ssp\_)}}\\
reads: the nodes of the strand space \$1 \\
definition:
\begin{quote}\ttfamily\footnotesize\raggedright\noindent
\{ssn\_ in cartesian(\allowbreak{}nn, ids(\allowbreak{}ssp\_)): 1 <= nth(\allowbreak{}1, ssn\_) and nth(\allowbreak{}1, ssn\_) <= length(\allowbreak{}(\allowbreak{}strand-map(\allowbreak{}ssp\_))(\allowbreak{}nth(\allowbreak{}2, ssn\_)))\}
\end{quote}
characterized by \texttt{cx-c@bnodes}  \textit{(definitional)}:
\begin{quote}\ttfamily\footnotesize\raggedright\noindent
bnodes(\allowbreak{}cx-c) == nodes(\allowbreak{}cx-sp)
\end{quote}
also: \texttt{cx-nodes-set} \texttt{nodes-unfold} \texttt{cx-nodes-finite} \texttt{cx-nodes-in-box} \texttt{resp-trace-nodes} \texttt{init-strand-nodes} \texttt{resp-strand-nodes} \texttt{three-strand-nodes} ... (9 in all) \\
mentioned by 75 result(s)

\smallskip

\noindent\textbf{\texttt{non-zero(s)}}\quad\textit{accessor}\index{non-zero(s)@\texttt{non-zero(s)}}\\
reads: the non-zero elements of \$1 \\
a structure slot \\
characterized by \texttt{qq-field@non-zero}  \textit{(definitional)}:
\begin{quote}\ttfamily\footnotesize\raggedright\noindent
non-zero(\allowbreak{}qq-field) == difference(\allowbreak{}qq, singleton(\allowbreak{}0))
\end{quote}
also: \texttt{rr-field@non-zero} \texttt{field-non-zero-carrier} \texttt{hom-field-non-zero} \\
mentioned by 11 result(s)

\smallskip

\noindent\textbf{\texttt{nonneg-rr}}\quad\textit{predicate}\index{nonneg-rr@\texttt{nonneg-rr}}\\
reads: a nonnegative real \\
a predicate

\smallskip

\noindent\textbf{\texttt{norm-ball(c6vm\_, c6vn\_, c6vr\_)}}\quad\textit{functoid}\index{norm-ball(c6vm_, c6vn_, c6vr_)@\texttt{norm-ball(c6vm\_, c6vn\_, c6vr\_)}}\\
reads: the closed ball of radius \$3 about 0 in \$1 for the norm \$2 \\
definition:
\begin{quote}\ttfamily\footnotesize\raggedright\noindent
\{c6vx\_ in vec(\allowbreak{}c6vm\_): c6vn\_(\allowbreak{}c6vx\_) <= c6vr\_\}
\end{quote}
mentioned by 6 result(s)

\smallskip

\noindent\textbf{\texttt{norm-optop(otv\_)}}\quad\textit{functoid}\index{norm-optop(otv_)@\texttt{norm-optop(otv\_)}}\\
reads: the norm topology on the bounded operators on \$1 \\
definition:
\begin{quote}\ttfamily\footnotesize\raggedright\noindent
seminorm-top(\allowbreak{}bh-module(\allowbreak{}otv\_), pair(\allowbreak{}0, 0), norm-optop-pf(\allowbreak{}otv\_))
\end{quote}
characterized by \texttt{norm-optop-is-tvs}  \textit{(proven)}:
\begin{quote}\ttfamily\footnotesize\raggedright\noindent
forall(\allowbreak{}[v], is-complex-inner-product-space(\allowbreak{}v) implies is-tvs(\allowbreak{}cc-normed-field, bh-module(\allowbreak{}v), norm-optop(\allowbreak{}v)))
\end{quote}
mentioned by 2 result(s)

\smallskip

\noindent\textbf{\texttt{norm-optop-pf(otv\_)}}\quad\textit{functoid}\index{norm-optop-pf(otv_)@\texttt{norm-optop-pf(otv\_)}}\\
reads: the operator norm as a family of seminorms on the bounded operators on \$1 \\
definition:
\begin{quote}\ttfamily\footnotesize\raggedright\noindent
vnb-lambda(\allowbreak{}otk\_, pair(\allowbreak{}0, 0), vnb-lambda(\allowbreak{}ott\_, bounded-ops(\allowbreak{}otv\_, otv\_), opnorm(\allowbreak{}otv\_, otv\_, ott\_)))
\end{quote}
characterized by \texttt{norm-optop-pf-type}  \textit{(proven)}:
\begin{quote}\ttfamily\footnotesize\raggedright\noindent
forall(\allowbreak{}[v], is-complex-inner-product-space(\allowbreak{}v) implies norm-optop-pf(\allowbreak{}v) in fun(\allowbreak{}pair(\allowbreak{}0, 0), fun(\allowbreak{}bounded-ops(\allowbreak{}v, v), rr)))
\end{quote}
also: \texttt{norm-optop-pf-apply} \\
mentioned by 7 result(s)

\smallskip

\noindent\textbf{\texttt{normed-ag}}\quad\textit{structure}\index{normed-ag@\texttt{normed-ag}}\\
predicate \texttt{is-normed-ag}; slots: carr opr iden inv nrm

\smallskip

\noindent\textbf{\texttt{normed-ag-as-abelian-group(s)}}\quad\textit{view}\index{normed-ag-as-abelian-group(s)@\texttt{normed-ag-as-abelian-group(s)}}\\
a normed-ag viewed as an abelian-group \\
characterized by \texttt{normed-ag-as-abelian-group-is-abelian-group}  \textit{(definitional)}:
\begin{quote}\ttfamily\footnotesize\raggedright\noindent
forall(\allowbreak{}[r], is-normed-ag(\allowbreak{}r) implies is-abelian-group(\allowbreak{}normed-ag-as-abelian-group(\allowbreak{}r)))
\end{quote}
also: \texttt{abelian-group-as-monoid-is-monoid-normed-ag-as-abelian-group} \texttt{normed-ag-as-abelian-group-is-abelian-group-cna-as-normed-ag} \texttt{normed-ag-as-abelian-group-is-abelian-group-uba-as-normed-ag} \texttt{normed-ag-as-abelian-group-is-abelian-group-cnvs-as-normed-ag} \texttt{normed-ag-as-abelian-group-is-abelian-group-uba-as-normed-ag-uiba-as-uba} \texttt{normed-ag-as-abelian-group-is-abelian-group-normed-vector-space-as-normed-ag} \texttt{abelian-group-as-monoid-is-monoid-normed-ag-as-abelian-group-cna-as-normed-ag} \texttt{abelian-group-as-monoid-is-monoid-normed-ag-as-abelian-group-uba-as-normed-ag} ... (15 in all) \\
mentioned by 78 result(s)

\smallskip

\noindent\textbf{\texttt{normed-field}}\quad\textit{structure}\index{normed-field@\texttt{normed-field}}\\
predicate \texttt{is-normed-field}; slots: carr add mul neg zero one fnrm

\smallskip

\noindent\textbf{\texttt{normed-field-additive-ag(s)}}\quad\textit{view}\index{normed-field-additive-ag(s)@\texttt{normed-field-additive-ag(s)}}\\
a normed-field viewed as an abelian-group \\
characterized by \texttt{normed-field-additive-ag-is-abelian-group}  \textit{(definitional)}:
\begin{quote}\ttfamily\footnotesize\raggedright\noindent
forall(\allowbreak{}[r], is-normed-field(\allowbreak{}r) implies is-abelian-group(\allowbreak{}normed-field-additive-ag(\allowbreak{}r)))
\end{quote}
also: \texttt{abelian-group-as-monoid-is-monoid-normed-field-additive-ag} \texttt{normed-field-additive-ag-functorial} \\
mentioned by 64 result(s)

\smallskip

\noindent\textbf{\texttt{normed-field-as-commutative-ring(s)}}\quad\textit{view}\index{normed-field-as-commutative-ring(s)@\texttt{normed-field-as-commutative-ring(s)}}\\
a normed-field viewed as a commutative-ring \\
characterized by \texttt{normed-field-as-commutative-ring-is-commutative-ring}  \textit{(definitional)}:
\begin{quote}\ttfamily\footnotesize\raggedright\noindent
forall(\allowbreak{}[r], is-normed-field(\allowbreak{}r) implies is-commutative-ring(\allowbreak{}normed-field-as-commutative-ring(\allowbreak{}r)))
\end{quote}
also: \texttt{commutative-ring-additive-ag-is-abelian-group-normed-field-as-commutative-ring} \texttt{commutative-ring-multiplicative-cm-is-comm-monoid-normed-field-as-commutative-ring} \texttt{abelian-group-as-monoid-is-monoid-commutative-ring-additive-ag-normed-field-as-commutative-ring} \texttt{normed-field-as-commutative-ring-functorial} \\
mentioned by 2206 result(s)

\smallskip

\noindent\textbf{\texttt{normed-field-as-integral-domain(s)}}\quad\textit{view}\index{normed-field-as-integral-domain(s)@\texttt{normed-field-as-integral-domain(s)}}\\
a normed-field viewed as an integral-domain \\
characterized by \texttt{normed-field-as-integral-domain-is-integral-domain}  \textit{(definitional)}:
\begin{quote}\ttfamily\footnotesize\raggedright\noindent
forall(\allowbreak{}[r], is-normed-field(\allowbreak{}r) implies is-integral-domain(\allowbreak{}normed-field-as-integral-domain(\allowbreak{}r)))
\end{quote}
also: \texttt{normed-field-as-integral-domain-functorial} \\
mentioned by 2 result(s)

\smallskip

\noindent\textbf{\texttt{normed-vector-space}}\quad\textit{structure}\index{normed-vector-space@\texttt{normed-vector-space}}\\
predicate \texttt{is-normed-vector-space}; slots: scal vec vadd vzero vneg act vnrm

\smallskip

\noindent\textbf{\texttt{normed-vector-space-as-module(s)}}\quad\textit{view}\index{normed-vector-space-as-module(s)@\texttt{normed-vector-space-as-module(s)}}\\
a normed-vector-space viewed as a module \\
characterized by \texttt{normed-vector-space-as-module-is-module}  \textit{(definitional)}:
\begin{quote}\ttfamily\footnotesize\raggedright\noindent
forall(\allowbreak{}[r], is-normed-vector-space(\allowbreak{}r) implies is-module(\allowbreak{}normed-vector-space-as-module(\allowbreak{}r)))
\end{quote}
also: \texttt{normed-vector-space-as-module-is-module-cnvs-as-nvs} \texttt{module-vector-ag-is-abelian-group-normed-vector-space-as-module} \texttt{module-vector-ag-is-abelian-group-normed-vector-space-as-module-cnvs-as-nvs} \texttt{abelian-group-as-monoid-is-monoid-module-vector-ag-normed-vector-space-as-module} \texttt{abelian-group-as-monoid-is-monoid-module-vector-ag-normed-vector-space-as-module-cnvs-as-nvs} \texttt{normed-vector-space-as-module-functorial} \\
mentioned by 95 result(s)

\smallskip

\noindent\textbf{\texttt{normed-vector-space-as-normed-ag(s)}}\quad\textit{view}\index{normed-vector-space-as-normed-ag(s)@\texttt{normed-vector-space-as-normed-ag(s)}}\\
a normed-vector-space viewed as a normed-ag \\
characterized by \texttt{normed-vector-space-as-normed-ag-is-normed-ag}  \textit{(definitional)}:
\begin{quote}\ttfamily\footnotesize\raggedright\noindent
forall(\allowbreak{}[r], is-normed-vector-space(\allowbreak{}r) implies is-normed-ag(\allowbreak{}normed-vector-space-as-normed-ag(\allowbreak{}r)))
\end{quote}
also: \texttt{normed-vector-space-as-normed-ag-is-normed-ag-cnvs-as-nvs} \texttt{normed-ag-as-abelian-group-is-abelian-group-normed-vector-space-as-normed-ag} \texttt{normed-ag-as-abelian-group-is-abelian-group-normed-vector-space-as-normed-ag-cnvs-as-nvs} \texttt{abelian-group-as-monoid-is-monoid-normed-ag-as-abelian-group-normed-vector-space-as-normed-ag} \texttt{abelian-group-as-monoid-is-monoid-normed-ag-as-abelian-group-normed-vector-space-as-normed-ag-cnvs-as-nvs} \texttt{normed-vector-space-as-normed-ag-functorial} \\
mentioned by 7 result(s)

\smallskip

\noindent\textbf{\texttt{npe(m, s, f, t, g)}}\quad\textit{predicate}\index{npe(m, s, f, t, g)@\texttt{npe(m, s, f, t, g)}}\\
reads: \$5 is a norm-preserving extension of \$3 from \$2 to \$4, in \$1 \\
defined by \texttt{npe} \\
characterized by \texttt{npe}  \textit{(definitional)}:
\begin{quote}\ttfamily\footnotesize\raggedright\noindent
forall(\allowbreak{}[m, s, f, t, g], npe(\allowbreak{}m, s, f, t, g) iff is-submodule(\allowbreak{}m, t) and s subset t and is-linear-functional-on(\allowbreak{}m, t, g) and extends-on(\allowbreak{}s, g, f) and forall(\allowbreak{}[w\_ in t], abs(\allowbreak{}g(\allowbreak{}w\_)) <= dual-norm-on(\allowbreak{}m, s, f) * (\allowbreak{}vnrm(\allowbreak{}m))(\allowbreak{}w\_)))
\end{quote}
mentioned by 4 result(s)

\smallskip

\noindent\textbf{\texttt{nrm(s)}}\quad\textit{accessor}\index{nrm(s)@\texttt{nrm(s)}}\\
reads: the norm of \$1 \\
a structure slot \\
characterized by \texttt{cc-nag@nrm}  \textit{(definitional)}:
\begin{quote}\ttfamily\footnotesize\raggedright\noindent
nrm(\allowbreak{}cc-nag) == vnb-lambda(\allowbreak{}x\_, cc, magnitude(\allowbreak{}x\_))
\end{quote}
also: \texttt{cc-nag-nrm} \texttt{bop-nag-nrm} \texttt{cbop-nag-nrm} \texttt{rb2-nrm-finsum-le} \texttt{normed-ag-nrm-type} \texttt{normed-ag-nrm-nonneg} \texttt{normed-ag-nrm-definite} \texttt{normed-ag-nrm-subadditive} \\
mentioned by 34 result(s)

\smallskip

\noindent\textbf{\texttt{nstar}}\quad\textit{constant}\index{nstar@\texttt{nstar}}\\
reads: the finite sequences of naturals \\
defined by \texttt{nstar-def}

\smallskip

\noindent\textbf{\texttt{nth}}\quad\textit{operator}\index{nth@\texttt{nth}}\\
a term-forming head \\
characterized by \texttt{cr-nth-eq}  \textit{(definitional)}:
\begin{quote}\ttfamily\footnotesize\raggedright\noindent
forall(\allowbreak{}[mct\_, mchs\_, mcc\_], cr-nth-eq(\allowbreak{}mct\_, mchs\_, mcc\_) iff forsome(\allowbreak{}[mcg\_, mcl\_, mck\_], mcc\_ = [mcg\_, enc-qeq(\allowbreak{}enc-nth(\allowbreak{}enc-num(\allowbreak{}mck\_), enc-list(\allowbreak{}mcl\_)), nth(\allowbreak{}mck\_, mcl\_))] and mck\_ in nn and 1 <= mck\_ and mck\_ <= length(\allowbreak{}mcl\_) and mchs\_ = [[mcg\_, enc-eq(\allowbreak{}enc-list(\allowbreak{}mcl\_), enc-list(\allowbreak{}mcl\_))]]))
\end{quote}
also: \texttt{nth-cons-1} \texttt{def-nth-lit} \texttt{def-nth-list} \texttt{nth-cons-succ} \texttt{nth-in-range} \texttt{ms-nth-2} \texttt{nth1-pair} \texttt{nth2-pair} ... (11 in all) \\
mentioned by 592 result(s)

\smallskip

\noindent\textbf{\texttt{nth-deriv}}\quad\textit{defined-fn}\index{nth-deriv@\texttt{nth-deriv}}\\
defined by \texttt{nth-deriv-zero} \texttt{nth-deriv-succ} \\
characterized by \texttt{nth-deriv-succ}  \textit{(definitional)}:
\begin{quote}\ttfamily\footnotesize\raggedright\noindent
forall(\allowbreak{}[f, n in nn], nth-deriv(\allowbreak{}f, succ(\allowbreak{}n)) == vnb-lambda(\allowbreak{}x, rr, deriv(\allowbreak{}nth-deriv(\allowbreak{}f, n), x)))
\end{quote}
also: \texttt{nth-deriv-zero} \texttt{nth-deriv-one} \texttt{gof-nth-deriv-dfun} \texttt{nth-deriv-monomial} \texttt{nth-deriv-succ-inner} \texttt{r8e-gof-nth-deriv-fn} \texttt{poly-nth-deriv-in-fun} \texttt{nth-deriv-monomial-gap} ... (12 in all) \\
mentioned by 62 result(s)

\smallskip

\noindent\textbf{\texttt{nth-deriv-on}}\quad\textit{functoid}\index{nth-deriv-on@\texttt{nth-deriv-on}}\\
reads: the \$4-th derivative of \$3 on \$2 \\
defined by \texttt{nth-deriv-on-zero} \texttt{nth-deriv-on-succ} \\
characterized by \texttt{nth-deriv-on-succ}  \textit{(definitional)}:
\begin{quote}\ttfamily\footnotesize\raggedright\noindent
forall(\allowbreak{}[ndvk\_, ndvu\_, ndvf\_, ndvn\_ in nn], nth-deriv-on(\allowbreak{}ndvk\_, ndvu\_, ndvf\_, succ(\allowbreak{}ndvn\_)) == vnb-lambda(\allowbreak{}ndvz\_, ndvu\_, deriv-on(\allowbreak{}ndvk\_, ndvu\_, nth-deriv-on(\allowbreak{}ndvk\_, ndvu\_, ndvf\_, ndvn\_), ndvz\_)))
\end{quote}
also: \texttt{nth-deriv-on-zero} \texttt{nth-deriv-on-local} \texttt{nth-deriv-on-in-fun} \texttt{nth-deriv-on-succ-eq} \texttt{nth-deriv-on-succ-apply} \texttt{nth-deriv-on-of-is-diff-on} \texttt{nth-deriv-on-holomorphic-ball} \\
mentioned by 30 result(s)

\smallskip

\noindent\textbf{\texttt{nth-deriv-v}}\quad\textit{defined-fn}\index{nth-deriv-v@\texttt{nth-deriv-v}}\\
defined by \texttt{nth-deriv-v-zero} \texttt{nth-deriv-v-succ} \\
characterized by \texttt{nth-deriv-v-succ}  \textit{(definitional)}:
\begin{quote}\ttfamily\footnotesize\raggedright\noindent
forall(\allowbreak{}[m, f, n in nn], nth-deriv-v(\allowbreak{}m, f, succ(\allowbreak{}n)) == vnb-lambda(\allowbreak{}x, rr, deriv-v(\allowbreak{}m, nth-deriv-v(\allowbreak{}m, f, n), x)))
\end{quote}
also: \texttt{nth-deriv-v-zero} \texttt{nth-deriv-v-in-vec} \texttt{r8e-nth-deriv-v-diff} \\
mentioned by 31 result(s)

\smallskip

\noindent\textbf{\texttt{nth-r}}\quad\textit{tactic}\index{nth-r@\texttt{nth-r}}\\
\texttt{(nth-r)} -- Reduce an NTH applied to a literal LIST in the goal. \\
declared in \texttt{tactics-help.scm}

\smallskip

\noindent\textbf{\texttt{nu-c-star-algebra}}\quad\textit{refinement}\index{nu-c-star-algebra@\texttt{nu-c-star-algebra}}\\
predicate \texttt{is-nu-c-star-algebra}; same shape as nu-involutive-banach-algebra

\smallskip

\noindent\textbf{\texttt{nu-involutive-banach-algebra}}\quad\textit{structure}\index{nu-involutive-banach-algebra@\texttt{nu-involutive-banach-algebra}}\\
predicate \texttt{is-nu-involutive-banach-algebra}; slots: scal vec vadd vzero vneg act vnrm vmul istar

\smallskip

\noindent\textbf{\texttt{nuiba-as-cba(s)}}\quad\textit{view}\index{nuiba-as-cba(s)@\texttt{nuiba-as-cba(s)}}\\
a nu-involutive-banach-algebra viewed as a complex-banach-algebra \\
characterized by \texttt{nuiba-as-cba-is-complex-banach-algebra}  \textit{(definitional)}:
\begin{quote}\ttfamily\footnotesize\raggedright\noindent
forall(\allowbreak{}[r], is-nu-involutive-banach-algebra(\allowbreak{}r) implies is-complex-banach-algebra(\allowbreak{}nuiba-as-cba(\allowbreak{}r)))
\end{quote}
also: \texttt{is-complex-banach-algebra-is-set-nuiba-as-cba} \texttt{nuiba-as-cba-is-complex-banach-algebra-uiba-as-nuiba} \texttt{is-complex-banach-algebra-is-set-nuiba-as-cba-uiba-as-nuiba} \texttt{nuiba-as-cba-functorial} \\
mentioned by 33 result(s)

\smallskip

\noindent\textbf{\texttt{null-homotopic(hmu\_, hmg\_, hma\_, hmb\_)}}\quad\textit{predicate}\index{null-homotopic(hmu_, hmg_, hma_, hmb_)@\texttt{null-homotopic(hmu\_, hmg\_, hma\_, hmb\_)}}\\
reads: the closed curve \$2 is null homotopic in \$1 \\
defined by \texttt{null-homotopic} \\
characterized by \texttt{null-homotopic}  \textit{(definitional)}:
\begin{quote}\ttfamily\footnotesize\raggedright\noindent
forall(\allowbreak{}[hmu\_, hmg\_, hma\_, hmb\_], null-homotopic(\allowbreak{}hmu\_, hmg\_, hma\_, hmb\_) iff hmg\_(\allowbreak{}hmb\_) = hmg\_(\allowbreak{}hma\_) and homotopic(\allowbreak{}hmu\_, hmg\_, vnb-lambda(\allowbreak{}hmx\_, ccint(\allowbreak{}hma\_, hmb\_), hmg\_(\allowbreak{}hma\_)), hma\_, hmb\_))
\end{quote}
also: \texttt{null-homotopic-line-int} \texttt{convex-closed-path-null-homotopic} \\
mentioned by 17 result(s)

\smallskip

\noindent\textbf{\texttt{null-rr-seq(rad)}}\quad\textit{predicate}\index{null-rr-seq(rad)@\texttt{null-rr-seq(rad)}}\\
reads: \$1 is a sequence of positive reals tending to zero \\
defined by \texttt{null-rr-seq} \\
characterized by \texttt{null-rr-seq}  \textit{(definitional)}:
\begin{quote}\ttfamily\footnotesize\raggedright\noindent
forall(\allowbreak{}[rad], null-rr-seq(\allowbreak{}rad) iff rad in fun(\allowbreak{}nn, rr) and forall(\allowbreak{}[k in nn], pos-rr(\allowbreak{}rad(\allowbreak{}k))) and forall(\allowbreak{}[eps], pos-rr(\allowbreak{}eps) implies forsome(\allowbreak{}[n in nn], forall(\allowbreak{}[k], k in nn and n <= k implies rad(\allowbreak{}k) <= eps))))
\end{quote}
also: \texttt{null-rr-seq-exists} \\
mentioned by 9 result(s)

\smallskip

\noindent\textbf{\texttt{nv-ball(cwk\_, cwa\_, cwr\_)}}\quad\textit{functoid}\index{nv-ball(cwk_, cwa_, cwr_)@\texttt{nv-ball(cwk\_, cwa\_, cwr\_)}}\\
reads: the elements of \$2 of absolute value at most \$3 \\
definition:
\begin{quote}\ttfamily\footnotesize\raggedright\noindent
\{cwx\_ in cwa\_: (\allowbreak{}fnrm(\allowbreak{}cwk\_))(\allowbreak{}cwx\_) <= cwr\_\}
\end{quote}
characterized by \texttt{nv-ball-unfold}  \textit{(proven)}:
\begin{quote}\ttfamily\footnotesize\raggedright\noindent
forall(\allowbreak{}[k, a, r], nv-ball(\allowbreak{}k, a, r) == \{cwx\_ in a: (\allowbreak{}fnrm(\allowbreak{}k))(\allowbreak{}cwx\_) <= r\})
\end{quote}
mentioned by 4 result(s)

\smallskip

\noindent\textbf{\texttt{nv-base(bzk\_, bzw\_)}}\quad\textit{functoid}\index{nv-base(bzk_, bzw_)@\texttt{nv-base(bzk\_, bzw\_)}}\\
reads: the closure in \$1 of its prime field, adjoined with \$2 in prime characteristic \\
definition:
\begin{quote}\ttfamily\footnotesize\raggedright\noindent
closure(\allowbreak{}nf-metric-space(\allowbreak{}bzk\_), nf-subfield-gen(\allowbreak{}bzk\_, nv-base-gen(\allowbreak{}bzk\_, bzw\_)))
\end{quote}
characterized by \texttt{nv-base-char-p}  \textit{(proven)}:
\begin{quote}\ttfamily\footnotesize\raggedright\noindent
forall(\allowbreak{}[k], is-normed-field(\allowbreak{}k) implies forall(\allowbreak{}[w in carr(\allowbreak{}k), p], has-characteristic(\allowbreak{}normed-field-as-commutative-ring(\allowbreak{}k), p) implies w in nv-base(\allowbreak{}k, w)))
\end{quote}
also: \texttt{nv-base-closed} \texttt{nv-base-unfold} \texttt{nv-base-subfield} \texttt{nv-base-char-0-unif} \texttt{nv-base-contains-gen} \texttt{nv-base-char-0-residue} \\
mentioned by 15 result(s)

\smallskip

\noindent\textbf{\texttt{nv-base-gen(bzk\_, bzw\_)}}\quad\textit{functoid}\index{nv-base-gen(bzk_, bzw_)@\texttt{nv-base-gen(bzk\_, bzw\_)}}\\
reads: the generators of B(\$1): \{\$2\} in prime characteristic, none in characteristic 0 \\
definition:
\begin{quote}\ttfamily\footnotesize\raggedright\noindent
\{bzz\_ in carr(\allowbreak{}bzk\_): bzz\_ = bzw\_ and forsome(\allowbreak{}[bzp\_], has-characteristic(\allowbreak{}normed-field-as-commutative-ring(\allowbreak{}bzk\_), bzp\_))\}
\end{quote}
characterized by \texttt{nv-base-gen-sub}  \textit{(proven)}:
\begin{quote}\ttfamily\footnotesize\raggedright\noindent
forall(\allowbreak{}[k], is-normed-field(\allowbreak{}k) implies forall(\allowbreak{}[w], nv-base-gen(\allowbreak{}k, w) subset carr(\allowbreak{}k)))
\end{quote}
also: \texttt{nv-base-gen-unfold} \\
mentioned by 7 result(s)

\smallskip

\noindent\textbf{\texttt{nv-digit(dgk\_, dgr\_, dgz\_)}}\quad\textit{functoid}\index{nv-digit(dgk_, dgr_, dgz_)@\texttt{nv-digit(dgk\_, dgr\_, dgz\_)}}\\
reads: the representative in \$2 of \$3 \\
definition:
\begin{quote}\ttfamily\footnotesize\raggedright\noindent
iota(\allowbreak{}dgy\_, dgy\_ in dgr\_ and (\allowbreak{}fnrm(\allowbreak{}dgk\_))(\allowbreak{}(\allowbreak{}add(\allowbreak{}dgk\_))(\allowbreak{}dgz\_, (\allowbreak{}neg(\allowbreak{}dgk\_))(\allowbreak{}dgy\_))) < 1)
\end{quote}
characterized by \texttt{nv-digit-def}  \textit{(proven)}:
\begin{quote}\ttfamily\footnotesize\raggedright\noindent
forall(\allowbreak{}[k, a, w, rp], is-compact-nv-domain(\allowbreak{}k, a) implies w in a implies nv-maximal(\allowbreak{}k, a) = nv-principal(\allowbreak{}k, a, w) implies not(\allowbreak{}w = zero(\allowbreak{}k)) implies a = nv-integers(\allowbreak{}k) implies is-nv-rep-set(\allowbreak{}k, a, rp) implies forall(\allowbreak{}[z in a], nv-digit(\allowbreak{}k, rp, z) in rp and (\allowbreak{}fnrm(\allowbreak{}k))(\allowbreak{}(\allowbreak{}add(\allowbreak{}k))(\allowbreak{}z, (\allowbreak{}neg(\allowbreak{}k))(\allowbreak{}nv-digit(\allowbreak{}k, rp, z)))) < 1))
\end{quote}
also: \texttt{nv-digit-step} \texttt{nv-digit-sequence} \\
mentioned by 5 result(s)

\smallskip

\noindent\textbf{\texttt{nv-digits(ndsk\_, ndsw\_, ndsr\_, ndsx\_)}}\quad\textit{functoid}\index{nv-digits(ndsk_, ndsw_, ndsr_, ndsx_)@\texttt{nv-digits(ndsk\_, ndsw\_, ndsr\_, ndsx\_)}}\\
reads: the digits in \$3 of the expansion of \$4 in powers of \$2 \\
definition:
\begin{quote}\ttfamily\footnotesize\raggedright\noindent
iota(\allowbreak{}ndsa\_, ndsa\_ in fun(\allowbreak{}zz, ndsr\_) and forsome(\allowbreak{}[ndsn\_ in zz], is-nv-expansion(\allowbreak{}ndsk\_, ndsw\_, ndsr\_, ndsa\_, ndsn\_, ndsx\_)))
\end{quote}
characterized by \texttt{nv-digits-fun}  \textit{(proven)}:
\begin{quote}\ttfamily\footnotesize\raggedright\noindent
forall(\allowbreak{}[k, a, w, rp], is-compact-nv-domain(\allowbreak{}k, a) implies w in a implies nv-maximal(\allowbreak{}k, a) = nv-principal(\allowbreak{}k, a, w) implies not(\allowbreak{}w = zero(\allowbreak{}k)) implies a = nv-integers(\allowbreak{}k) implies is-nv-rep-set(\allowbreak{}k, a, rp) implies vnb-lambda(\allowbreak{}ndx\_, carr(\allowbreak{}k), nv-digits(\allowbreak{}k, w, rp, ndx\_)) in fun(\allowbreak{}carr(\allowbreak{}k), fun(\allowbreak{}zz, rp)))
\end{quote}
also: \texttt{nv-digits-spec} \texttt{nv-digits-unique} \\
mentioned by 4 result(s)

\smallskip

\noindent\textbf{\texttt{nv-ideal2(cwk\_, cwa\_, cwg\_, cwh\_)}}\quad\textit{functoid}\index{nv-ideal2(cwk_, cwa_, cwg_, cwh_)@\texttt{nv-ideal2(cwk\_, cwa\_, cwg\_, cwh\_)}}\\
reads: the ideal of \$2 generated by \$3 and \$4 \\
definition:
\begin{quote}\ttfamily\footnotesize\raggedright\noindent
\{cwx\_ in cwa\_: forsome(\allowbreak{}[cwm\_ in cwa\_, cwn\_ in cwa\_], cwx\_ = (\allowbreak{}add(\allowbreak{}cwk\_))(\allowbreak{}(\allowbreak{}mul(\allowbreak{}cwk\_))(\allowbreak{}cwm\_, cwg\_), (\allowbreak{}mul(\allowbreak{}cwk\_))(\allowbreak{}cwn\_, cwh\_)))\}
\end{quote}
characterized by \texttt{nv-ideal2-add}  \textit{(proven)}:
\begin{quote}\ttfamily\footnotesize\raggedright\noindent
forall(\allowbreak{}[k, a], is-nv-domain(\allowbreak{}k, a) implies forall(\allowbreak{}[g in a, h in a, x in nv-ideal2(\allowbreak{}k, a, g, h), y in nv-ideal2(\allowbreak{}k, a, g, h)], (\allowbreak{}add(\allowbreak{}k))(\allowbreak{}x, y) in nv-ideal2(\allowbreak{}k, a, g, h)))
\end{quote}
also: \texttt{nv-ideal2-mul} \texttt{nv-ideal2-neg} \texttt{nv-ideal2-sub} \texttt{nv-ideal2-zero} \texttt{nv-ideal2-ideal} \texttt{nv-ideal2-unfold} \\
mentioned by 8 result(s)

\smallskip

\noindent\textbf{\texttt{nv-integers(cwk\_)}}\quad\textit{functoid}\index{nv-integers(cwk_)@\texttt{nv-integers(cwk\_)}}\\
reads: the ring of integral elements of \$1 \\
definition:
\begin{quote}\ttfamily\footnotesize\raggedright\noindent
\{cwx\_ in carr(\allowbreak{}cwk\_): (\allowbreak{}fnrm(\allowbreak{}cwk\_))(\allowbreak{}cwx\_) <= 1\}
\end{quote}
characterized by \texttt{pc-nv-integers-eq}  \textit{(proven)}:
\begin{quote}\ttfamily\footnotesize\raggedright\noindent
forall(\allowbreak{}[p], zz-prime(\allowbreak{}p) implies nv-integers(\allowbreak{}qq-padic(\allowbreak{}p)) = \{cwx\_ in qq: (\allowbreak{}padic-abs(\allowbreak{}p))(\allowbreak{}cwx\_) <= 1\})
\end{quote}
also: \texttt{nv-integers-domain} \texttt{nv-integers-unfold} \texttt{nv-integers-compact} \texttt{nv-integers-subgroup} \texttt{nv-integers-iff-divisibility} \\
mentioned by 152 result(s)

\smallskip

\noindent\textbf{\texttt{nv-maximal(cwk\_, cwa\_)}}\quad\textit{functoid}\index{nv-maximal(cwk_, cwa_)@\texttt{nv-maximal(cwk\_, cwa\_)}}\\
reads: the elements of \$2 of absolute value less than one \\
definition:
\begin{quote}\ttfamily\footnotesize\raggedright\noindent
\{cwx\_ in cwa\_: (\allowbreak{}fnrm(\allowbreak{}cwk\_))(\allowbreak{}cwx\_) < 1\}
\end{quote}
characterized by \texttt{nv-maximal-add}  \textit{(proven)}:
\begin{quote}\ttfamily\footnotesize\raggedright\noindent
forall(\allowbreak{}[k, a], is-nv-domain(\allowbreak{}k, a) implies forall(\allowbreak{}[x in nv-maximal(\allowbreak{}k, a), y in nv-maximal(\allowbreak{}k, a)], (\allowbreak{}add(\allowbreak{}k))(\allowbreak{}x, y) in nv-maximal(\allowbreak{}k, a)))
\end{quote}
also: \texttt{nv-maximal-mul} \texttt{nv-maximal-neg} \texttt{nv-maximal-sub} \texttt{nv-maximal-norm} \texttt{nv-maximal-zero} \texttt{nv-maximal-ideal} \texttt{nv-maximal-in-zf} \texttt{nv-maximal-proper} ... (10 in all) \\
mentioned by 95 result(s)

\smallskip

\noindent\textbf{\texttt{nv-multiples(cw3k\_, cw3a\_, cw3g\_)}}\quad\textit{functoid}\index{nv-multiples(cw3k_, cw3a_, cw3g_)@\texttt{nv-multiples(cw3k\_, cw3a\_, cw3g\_)}}\\
reads: the multiples in \$1 of \$3 by elements of \$2 \\
definition:
\begin{quote}\ttfamily\footnotesize\raggedright\noindent
\{cw3x\_ in carr(\allowbreak{}cw3k\_): forsome(\allowbreak{}[cw3m\_ in cw3a\_], cw3x\_ = (\allowbreak{}mul(\allowbreak{}cw3k\_))(\allowbreak{}cw3m\_, cw3g\_))\}
\end{quote}
characterized by \texttt{nv-multiples-unfold}  \textit{(proven)}:
\begin{quote}\ttfamily\footnotesize\raggedright\noindent
forall(\allowbreak{}[k, a, g], nv-multiples(\allowbreak{}k, a, g) == \{cw3x\_ in carr(\allowbreak{}k): forsome(\allowbreak{}[cw3m\_ in a], cw3x\_ = (\allowbreak{}mul(\allowbreak{}k))(\allowbreak{}cw3m\_, g))\})
\end{quote}
mentioned by 5 result(s)

\smallskip

\noindent\textbf{\texttt{nv-principal(cwk\_, cwa\_, cwg\_)}}\quad\textit{functoid}\index{nv-principal(cwk_, cwa_, cwg_)@\texttt{nv-principal(cwk\_, cwa\_, cwg\_)}}\\
reads: the principal ideal of \$2 generated by \$3 \\
definition:
\begin{quote}\ttfamily\footnotesize\raggedright\noindent
\{cwx\_ in cwa\_: forsome(\allowbreak{}[cwm\_ in cwa\_], cwx\_ = (\allowbreak{}mul(\allowbreak{}cwk\_))(\allowbreak{}cwm\_, cwg\_))\}
\end{quote}
characterized by \texttt{nv-principal-self}  \textit{(proven)}:
\begin{quote}\ttfamily\footnotesize\raggedright\noindent
forall(\allowbreak{}[k, a], is-nv-domain(\allowbreak{}k, a) implies forall(\allowbreak{}[g in a], g in nv-principal(\allowbreak{}k, a, g)))
\end{quote}
also: \texttt{nv-principal-unfold} \texttt{nv-principal-norm-le} \texttt{nv-principal-of-norm} \\
mentioned by 86 result(s)

\smallskip

\noindent\textbf{\texttt{nv-psum(ntk\_, ntw\_, nta\_, ntn\_, ntm\_)}}\quad\textit{functoid}\index{nv-psum(ntk_, ntw_, nta_, ntn_, ntm_)@\texttt{nv-psum(ntk\_, ntw\_, nta\_, ntn\_, ntm\_)}}\\
reads: the \$5-th partial sum of the series with digits \$3 from \$4 \\
definition:
\begin{quote}\ttfamily\footnotesize\raggedright\noindent
sum-ag(\allowbreak{}normed-field-additive-ag(\allowbreak{}ntk\_), nv-term(\allowbreak{}ntk\_, ntw\_, nta\_, ntn\_), ntm\_)
\end{quote}
characterized by \texttt{nv-psum-succ}  \textit{(proven)}:
\begin{quote}\ttfamily\footnotesize\raggedright\noindent
forall(\allowbreak{}[k], is-normed-field(\allowbreak{}k) implies forall(\allowbreak{}[w, a, n, m in nn], nv-psum(\allowbreak{}k, w, a, n, succ(\allowbreak{}m)) == (\allowbreak{}add(\allowbreak{}k))(\allowbreak{}nv-psum(\allowbreak{}k, w, a, n, m), (\allowbreak{}mul(\allowbreak{}k))(\allowbreak{}a(\allowbreak{}n + m), nv-zpow(\allowbreak{}k, w, n + m)))))
\end{quote}
also: \texttt{nv-psum-type} \texttt{nv-psum-zero} \\
mentioned by 13 result(s)

\smallskip

\noindent\textbf{\texttt{nv-res-ringoid(rsk\_, rsa\_)}}\quad\textit{functoid}\index{nv-res-ringoid(rsk_, rsa_)@\texttt{nv-res-ringoid(rsk\_, rsa\_)}}\\
reads: the ring \$2 with its ideal M(\$2) \\
definition:
\begin{quote}\ttfamily\footnotesize\raggedright\noindent
[rsa\_, vnb-lambda(\allowbreak{}[rsu\_, rsv\_], cartesian(\allowbreak{}rsa\_, rsa\_), (\allowbreak{}add(\allowbreak{}rsk\_))(\allowbreak{}rsu\_, rsv\_)), vnb-lambda(\allowbreak{}[rsu\_, rsv\_], cartesian(\allowbreak{}rsa\_, rsa\_), (\allowbreak{}mul(\allowbreak{}rsk\_))(\allowbreak{}rsu\_, rsv\_)), vnb-lambda(\allowbreak{}rsu\_, rsa\_, (\allowbreak{}neg(\allowbreak{}rsk\_))(\allowbreak{}rsu\_)), zero(\allowbreak{}rsk\_), one(\allowbreak{}rsk\_), nv-maximal(\allowbreak{}rsk\_, rsa\_)]
\end{quote}
mentioned by 26 result(s)

\smallskip

\noindent\textbf{\texttt{nv-residue(rsk\_, rsa\_)}}\quad\textit{functoid}\index{nv-residue(rsk_, rsa_)@\texttt{nv-residue(rsk\_, rsa\_)}}\\
reads: the residue ring \$2 / M(\$2) \\
definition:
\begin{quote}\ttfamily\footnotesize\raggedright\noindent
ringoid-quotient-ring(\allowbreak{}nv-res-ringoid(\allowbreak{}rsk\_, rsa\_))
\end{quote}
characterized by \texttt{nv-residue-fp-iso}  \textit{(proven)}:
\begin{quote}\ttfamily\footnotesize\raggedright\noindent
forall(\allowbreak{}[p, k, a, io], zz-prime(\allowbreak{}p) implies is-nv-domain(\allowbreak{}k, a) implies io in fun(\allowbreak{}zz, a) implies forall(\allowbreak{}[im in zz, jn in zz], io(\allowbreak{}im + jn) = (\allowbreak{}add(\allowbreak{}k))(\allowbreak{}io(\allowbreak{}im), io(\allowbreak{}jn)) and io(\allowbreak{}im * jn) = (\allowbreak{}mul(\allowbreak{}k))(\allowbreak{}io(\allowbreak{}im), io(\allowbreak{}jn))) implies forall(\allowbreak{}[im in zz], io(\allowbreak{}-im) = (\allowbreak{}neg(\allowbreak{}k))(\allowbreak{}io(\allowbreak{}im))) implies io(\allowbreak{}0) = zero(\allowbreak{}k) implies io(\allowbreak{}1) = one(\allowbreak{}k) implies forall(\allowbreak{}[im in zz], io(\allowbreak{}im) in nv-maximal(\allowbreak{}k, a) iff zz-divides(\allowbreak{}p, im)) implies forall(\allowbreak{}[xa in a], forsome(\allowbreak{}[im in zz], (\allowbreak{}add(\allowbreak{}k))(\allowbreak{}xa, (\allowbreak{}neg(\allowbreak{}k))(\allowbreak{}io(\allowbreak{}im))) in nv-maximal(\allowbreak{}k, a))) implies forsome(\allowbreak{}[ff], is-ring-iso(\allowbreak{}fp(\allowbreak{}p), nv-residue(\allowbreak{}k, a), ff)))
\end{quote}
also: \texttt{nv-residue-unfold} \texttt{nv-residue-zero-class} \\
mentioned by 13 result(s)

\smallskip

\noindent\textbf{\texttt{nv-term(ntk\_, ntw\_, nta\_, ntn\_)}}\quad\textit{functoid}\index{nv-term(ntk_, ntw_, nta_, ntn_)@\texttt{nv-term(ntk\_, ntw\_, nta\_, ntn\_)}}\\
reads: the terms a(\$4 + j) \$2\textasciicircum{}(\$4 + j) of the series with digits \$3 \\
definition:
\begin{quote}\ttfamily\footnotesize\raggedright\noindent
vnb-lambda(\allowbreak{}ntj\_, nn, (\allowbreak{}mul(\allowbreak{}ntk\_))(\allowbreak{}nta\_(\allowbreak{}ntn\_ + ntj\_), nv-zpow(\allowbreak{}ntk\_, ntw\_, ntn\_ + ntj\_)))
\end{quote}

\smallskip

\noindent\textbf{\texttt{nv-u1(c2k\_, c2a\_)}}\quad\textit{functoid}\index{nv-u1(c2k_, c2a_)@\texttt{nv-u1(c2k\_, c2a\_)}}\\
reads: the principal units of \$2 \\
definition:
\begin{quote}\ttfamily\footnotesize\raggedright\noindent
\{c2x\_ in c2a\_: (\allowbreak{}fnrm(\allowbreak{}c2k\_))(\allowbreak{}(\allowbreak{}add(\allowbreak{}c2k\_))(\allowbreak{}c2x\_, (\allowbreak{}neg(\allowbreak{}c2k\_))(\allowbreak{}one(\allowbreak{}c2k\_)))) < 1\}
\end{quote}
characterized by \texttt{nv-u1-mul}  \textit{(proven)}:
\begin{quote}\ttfamily\footnotesize\raggedright\noindent
forall(\allowbreak{}[k, a], is-nv-domain(\allowbreak{}k, a) implies forall(\allowbreak{}[x in nv-u1(\allowbreak{}k, a), y in nv-u1(\allowbreak{}k, a)], (\allowbreak{}mul(\allowbreak{}k))(\allowbreak{}x, y) in nv-u1(\allowbreak{}k, a)))
\end{quote}
also: \texttt{nv-u1-one} \texttt{nv-u1-sub} \texttt{nv-u1-norm} \texttt{nv-u1-intro} \texttt{nv-u1-unfold} \texttt{nv-u1-norm-one} \texttt{nv-u1-inverse-in} \\
mentioned by 13 result(s)

\smallskip

\noindent\textbf{\texttt{nv-units(cwk\_, cwa\_)}}\quad\textit{functoid}\index{nv-units(cwk_, cwa_)@\texttt{nv-units(cwk\_, cwa\_)}}\\
reads: the group of units of \$2 \\
definition:
\begin{quote}\ttfamily\footnotesize\raggedright\noindent
\{cwx\_ in cwa\_: forsome(\allowbreak{}[cwy\_ in cwa\_], (\allowbreak{}mul(\allowbreak{}cwk\_))(\allowbreak{}cwx\_, cwy\_) = one(\allowbreak{}cwk\_))\}
\end{quote}
characterized by \texttt{nv-units-unfold}  \textit{(proven)}:
\begin{quote}\ttfamily\footnotesize\raggedright\noindent
forall(\allowbreak{}[k, a], nv-units(\allowbreak{}k, a) == \{cwx\_ in a: forsome(\allowbreak{}[cwy\_ in a], (\allowbreak{}mul(\allowbreak{}k))(\allowbreak{}cwx\_, cwy\_) = one(\allowbreak{}k))\})
\end{quote}
mentioned by 33 result(s)

\smallskip

\noindent\textbf{\texttt{nv-zf(bzk\_, bzf\_)}}\quad\textit{functoid}\index{nv-zf(bzk_, bzf_)@\texttt{nv-zf(bzk\_, bzf\_)}}\\
reads: the integers of \$1 whose residue class lies in \$2 \\
definition:
\begin{quote}\ttfamily\footnotesize\raggedright\noindent
\{bzz\_ in nv-integers(\allowbreak{}bzk\_): class(\allowbreak{}ringoid-setoid(\allowbreak{}nv-res-ringoid(\allowbreak{}bzk\_, nv-integers(\allowbreak{}bzk\_))), bzz\_) in bzf\_\}
\end{quote}
mentioned by 2 result(s)

\smallskip

\noindent\textbf{\texttt{nv-zpow(zpk\_, zpw\_, zpn\_)}}\quad\textit{functoid}\index{nv-zpow(zpk_, zpw_, zpn_)@\texttt{nv-zpow(zpk\_, zpw\_, zpn\_)}}\\
reads: the power \$2\textasciicircum{}\$3 in \$1 (\$3 an integer) \\
definition:
\begin{quote}\ttfamily\footnotesize\raggedright\noindent
iota(\allowbreak{}zpy\_, zpy\_ in carr(\allowbreak{}zpk\_) and (\allowbreak{}mul(\allowbreak{}zpk\_))(\allowbreak{}zpy\_, ring-power(\allowbreak{}normed-field-as-commutative-ring(\allowbreak{}zpk\_), zpw\_, max(\allowbreak{}-zpn\_, 0))) = ring-power(\allowbreak{}normed-field-as-commutative-ring(\allowbreak{}zpk\_), zpw\_, max(\allowbreak{}zpn\_, 0)))
\end{quote}
characterized by \texttt{nv-zpow-add}  \textit{(proven)}:
\begin{quote}\ttfamily\footnotesize\raggedright\noindent
forall(\allowbreak{}[k, w], is-normed-field(\allowbreak{}k) implies w in carr(\allowbreak{}k) implies not(\allowbreak{}w = zero(\allowbreak{}k)) implies forall(\allowbreak{}[ze in zz, zf in zz], (\allowbreak{}mul(\allowbreak{}k))(\allowbreak{}nv-zpow(\allowbreak{}k, w, ze), nv-zpow(\allowbreak{}k, w, zf)) = nv-zpow(\allowbreak{}k, w, ze + zf)))
\end{quote}
also: \texttt{nv-zpow-def} \texttt{nv-zpow-inv} \texttt{nv-zpow-nat} \texttt{nv-zpow-neg} \texttt{nv-zpow-char} \texttt{nv-zpow-succ} \texttt{nv-zpow-nonzero} \texttt{nv-zpow-norm-nat} ... (10 in all) \\
mentioned by 30 result(s)

\smallskip

\noindent\textbf{\texttt{nvs-metric-space(m)}}\quad\textit{functoid}\index{nvs-metric-space(m)@\texttt{nvs-metric-space(m)}}\\
definition:
\begin{quote}\ttfamily\footnotesize\raggedright\noindent
[vec(\allowbreak{}m), vnb-lambda(\allowbreak{}[x, y], cartesian(\allowbreak{}vec(\allowbreak{}m), vec(\allowbreak{}m)), (\allowbreak{}vnrm(\allowbreak{}m))(\allowbreak{}(\allowbreak{}vadd(\allowbreak{}m))(\allowbreak{}x, (\allowbreak{}vneg(\allowbreak{}m))(\allowbreak{}y))))]
\end{quote}
mentioned by 120 result(s)

\smallskip


\section*{O}
\addcontentsline{toc}{section}{O}

\noindent\textbf{\texttt{obtain}}\quad\textit{tactic}\index{obtain@\texttt{obtain}}\\
\texttt{(obtain LANE)} -- Run LANE -- a thunk whose forward step lands a `there exists' into context -- then eliminate that existential and RETURN the fresh witness's name, read off the proof state rather than guessed. \\
declared in \texttt{tactics-help.scm}

\smallskip

\noindent\textbf{\texttt{occurs-in-space(ssp\_, ssu\_)}}\quad\textit{predicate}\index{occurs-in-space(ssp_, ssu_)@\texttt{occurs-in-space(ssp\_, ssu\_)}}\\
reads: \$2 occurs in the strand space \$1 \\
defined by \texttt{occurs-in-space} \\
characterized by \texttt{occurs-in-space}  \textit{(definitional)}:
\begin{quote}\ttfamily\footnotesize\raggedright\noindent
forall(\allowbreak{}[ssp\_, ssu\_], occurs-in-space(\allowbreak{}ssp\_, ssu\_) iff forsome(\allowbreak{}[ssn\_ in nodes(\allowbreak{}ssp\_)], is-subterm(\allowbreak{}tokens(\allowbreak{}ssp\_), keys(\allowbreak{}ssp\_), ssu\_, node-content(\allowbreak{}ssp\_, ssn\_))))
\end{quote}
mentioned by 2 result(s)

\smallskip

\noindent\textbf{\texttt{ocint(a, b)}}\quad\textit{functoid}\index{ocint(a, b)@\texttt{ocint(a, b)}}\\
reads: the interval from \$1 excluded to \$2 included \\
definition:
\begin{quote}\ttfamily\footnotesize\raggedright\noindent
\{osq\_ in rr: a < osq\_ and osq\_ <= b\}
\end{quote}
characterized by \texttt{ocint-shrink}  \textit{(proven)}:
\begin{quote}\ttfamily\footnotesize\raggedright\noindent
forall(\allowbreak{}[osx\_, osr\_, osv\_], osx\_ in rr implies pos-rr(\allowbreak{}osr\_) implies pos-rr(\allowbreak{}osv\_) implies osv\_ <= osr\_ implies ocint(\allowbreak{}osx\_ - osv\_, osx\_) subset ocint(\allowbreak{}osx\_ - osr\_, osx\_))
\end{quote}
also: \texttt{ocint-of-ooint} \texttt{ocint-elt-in-rr} \texttt{ocint-subset-rr} \texttt{ocint-membership} \texttt{ocint-subset-ooint} \texttt{lebesgue-outer-ocint} \texttt{ocint-right-endpoint} \texttt{lebesgue-measure-ocint} ... (9 in all) \\
mentioned by 19 result(s)

\smallskip

\noindent\textbf{\texttt{odd(x)}}\quad\textit{predicate}\index{odd(x)@\texttt{odd(x)}}\\
reads: \$1 is odd \\
defined by \texttt{odd} \\
characterized by \texttt{odd}  \textit{(definitional)}:
\begin{quote}\ttfamily\footnotesize\raggedright\noindent
forall(\allowbreak{}[x], odd(\allowbreak{}x) iff forsome(\allowbreak{}[k in zz], x = 2 * k + 1))
\end{quote}
also: \texttt{odd-def} \texttt{zz-odd-sq-odd} \texttt{zz-odd-mul-odd} \texttt{zz-even-iff-not-odd} \texttt{zz-not-even-and-odd} \\
mentioned by 10 result(s)

\smallskip

\noindent\textbf{\texttt{oi-l}}\quad\textit{tactic}\index{oi-l@\texttt{oi-l}}\\
\texttt{(oi-l)} -- OR-intro left: reduce an (OR a b) goal to a. \\
declared in \texttt{tactics-help.scm}

\smallskip

\noindent\textbf{\texttt{oi-r}}\quad\textit{tactic}\index{oi-r@\texttt{oi-r}}\\
\texttt{(oi-r)} -- OR-intro right: reduce an (OR a b) goal to b. \\
declared in \texttt{tactics-help.scm}

\smallskip

\noindent\textbf{\texttt{onb-sets(kqv\_)}}\quad\textit{functoid}\index{onb-sets(kqv_)@\texttt{onb-sets(kqv\_)}}\\
reads: the orthonormal basis sets of \$1 \\
definition:
\begin{quote}\ttfamily\footnotesize\raggedright\noindent
\{kqo\_ in power(\allowbreak{}vec(\allowbreak{}kqv\_)): is-orthonormal-basis(\allowbreak{}kqv\_, kqo\_, vnb-lambda(\allowbreak{}kqe\_, kqo\_, kqe\_))\}
\end{quote}
characterized by \texttt{onb-sets-unfold}  \textit{(proven)}:
\begin{quote}\ttfamily\footnotesize\raggedright\noindent
forall(\allowbreak{}[kqv\_], onb-sets(\allowbreak{}kqv\_) == \{kqo\_ in power(\allowbreak{}vec(\allowbreak{}kqv\_)): is-orthonormal-basis(\allowbreak{}kqv\_, kqo\_, vnb-lambda(\allowbreak{}kqe\_, kqo\_, kqe\_))\})
\end{quote}
also: \texttt{onb-sets-membership} \\
mentioned by 6 result(s)

\smallskip

\noindent\textbf{\texttt{one(s)}}\quad\textit{accessor}\index{one(s)@\texttt{one(s)}}\\
reads: the multiplicative identity of \$1 \\
a structure slot \\
characterized by \texttt{rq-one}  \textit{(definitional)}:
\begin{quote}\ttfamily\footnotesize\raggedright\noindent
forall(\allowbreak{}[r], one(\allowbreak{}ringoid-quotient-ring(\allowbreak{}r)) == class(\allowbreak{}ringoid-setoid(\allowbreak{}r), one(\allowbreak{}r)))
\end{quote}
also: \texttt{qq-ring@one} \texttt{zz-ring@one} \texttt{qq-field@one} \texttt{rr-field@one} \texttt{cc-normed-field@one} \texttt{rr-normed-field@one} \texttt{hbs-one-eq-cna-as-module} \texttt{hbs-one-eq-sip-as-module} ... (118 in all) \\
mentioned by 783 result(s)

\smallskip

\noindent\textbf{\texttt{ooint(a, b)}}\quad\textit{functoid}\index{ooint(a, b)@\texttt{ooint(a, b)}}\\
reads: the open interval from \$1 to \$2 \\
definition:
\begin{quote}\ttfamily\footnotesize\raggedright\noindent
\{icx\_ in rr: a < icx\_ and icx\_ < b\}
\end{quote}
characterized by \texttt{ooint-open}  \textit{(proven)}:
\begin{quote}\ttfamily\footnotesize\raggedright\noindent
forall(\allowbreak{}[a, b], a in rr implies b in rr implies is-open(\allowbreak{}rr-ms, ooint(\allowbreak{}a, b)))
\end{quote}
also: \texttt{ooint-center} \texttt{ooint-open-nf} \texttt{coint-of-ooint} \texttt{ocint-of-ooint} \texttt{ooint-elt-in-rr} \texttt{ooint-subset-rr} \texttt{ooint-membership} \texttt{ooint-reflect-in} ... (27 in all) \\
mentioned by 125 result(s)

\smallskip

\noindent\textbf{\texttt{op-add(goa-md\_, goa-t\_)}}\quad\textit{functoid}\index{op-add(goa-md_, goa-t_)@\texttt{op-add(goa-md\_, goa-t\_)}}\\
definition:
\begin{quote}\ttfamily\footnotesize\raggedright\noindent
vnb-lambda(\allowbreak{}[goaf\_, goag\_], cartesian(\allowbreak{}ops(\allowbreak{}goa-md\_, goa-t\_), ops(\allowbreak{}goa-md\_, goa-t\_)), vnb-lambda(\allowbreak{}goaw\_, vec(\allowbreak{}goa-md\_), (\allowbreak{}vadd(\allowbreak{}goa-md\_))(\allowbreak{}goaf\_(\allowbreak{}goaw\_), goag\_(\allowbreak{}goaw\_))))
\end{quote}
characterized by \texttt{op-add-of}  \textit{(proven)}:
\begin{quote}\ttfamily\footnotesize\raggedright\noindent
forall(\allowbreak{}[md, tt], is-hom-module(\allowbreak{}md, md, tt) implies is-commutative-ring(\allowbreak{}scal(\allowbreak{}md)) implies forall(\allowbreak{}[p in carr(\allowbreak{}poly(\allowbreak{}scal(\allowbreak{}md))), q in carr(\allowbreak{}poly(\allowbreak{}scal(\allowbreak{}md)))], (\allowbreak{}op-add(\allowbreak{}md, tt))(\allowbreak{}op-of(\allowbreak{}md, tt, p), op-of(\allowbreak{}md, tt, q)) = op-of(\allowbreak{}md, tt, (\allowbreak{}add(\allowbreak{}poly(\allowbreak{}scal(\allowbreak{}md))))(\allowbreak{}p, q))))
\end{quote}
mentioned by 5 result(s)

\smallskip

\noindent\textbf{\texttt{op-adjoint(bov\_, bow\_, bot\_)}}\quad\textit{functoid}\index{op-adjoint(bov_, bow_, bot_)@\texttt{op-adjoint(bov\_, bow\_, bot\_)}}\\
reads: the adjoint of \$3 \\
definition:
\begin{quote}\ttfamily\footnotesize\raggedright\noindent
vnb-lambda(\allowbreak{}boy\_, vec(\allowbreak{}bow\_), iota(\allowbreak{}boz\_, boz\_ in vec(\allowbreak{}bov\_) and forall(\allowbreak{}[box\_ in vec(\allowbreak{}bov\_)], (\allowbreak{}ip(\allowbreak{}bow\_))(\allowbreak{}bot\_(\allowbreak{}box\_), boy\_) = (\allowbreak{}ip(\allowbreak{}bov\_))(\allowbreak{}box\_, boz\_))))
\end{quote}
characterized by \texttt{op-adjoint-id}  \textit{(proven)}:
\begin{quote}\ttfamily\footnotesize\raggedright\noindent
forall(\allowbreak{}[v], is-hilbert-space(\allowbreak{}v) implies op-adjoint(\allowbreak{}v, v, bop-id(\allowbreak{}v)) = bop-id(\allowbreak{}v))
\end{quote}
also: \texttt{op-adjoint-ip} \texttt{op-adjoint-fun} \texttt{op-adjoint-neg} \texttt{op-adjoint-sum} \texttt{op-adjoint-comp} \texttt{op-adjoint-spec} \texttt{op-adjoint-type} \texttt{op-adjoint-bound} ... (18 in all) \\
mentioned by 89 result(s)

\smallskip

\noindent\textbf{\texttt{op-alg(gal-md\_, gal-t\_)}}\quad\textit{functoid}\index{op-alg(gal-md_, gal-t_)@\texttt{op-alg(gal-md\_, gal-t\_)}}\\
reads: the ring F[\$2] of the operators p(\$2) \\
definition:
\begin{quote}\ttfamily\footnotesize\raggedright\noindent
[ops(\allowbreak{}gal-md\_, gal-t\_), op-add(\allowbreak{}gal-md\_, gal-t\_), op-mul(\allowbreak{}gal-md\_, gal-t\_), op-neg(\allowbreak{}gal-md\_, gal-t\_), op-zero(\allowbreak{}gal-md\_), op-one(\allowbreak{}gal-md\_)]
\end{quote}
characterized by \texttt{op-alg-add}  \textit{(proven)}:
\begin{quote}\ttfamily\footnotesize\raggedright\noindent
forall(\allowbreak{}[md, tt], is-hom-module(\allowbreak{}md, md, tt) implies is-commutative-ring(\allowbreak{}scal(\allowbreak{}md)) implies add(\allowbreak{}op-alg(\allowbreak{}md, tt)) == op-add(\allowbreak{}md, tt))
\end{quote}
also: \texttt{op-alg-mul} \texttt{op-alg-neg} \texttt{op-alg-one} \texttt{op-alg-carr} \texttt{op-alg-zero} \texttt{op-alg-is-commutative-ring} \\
mentioned by 17 result(s)

\smallskip

\noindent\textbf{\texttt{op-dual(otv\_, otu\_)}}\quad\textit{functoid}\index{op-dual(otv_, otu_)@\texttt{op-dual(otv\_, otu\_)}}\\
reads: the linear functionals on the bounded operators on \$1 continuous for \$2 \\
definition:
\begin{quote}\ttfamily\footnotesize\raggedright\noindent
\{otf\_ in fun(\allowbreak{}bounded-ops(\allowbreak{}otv\_, otv\_), cc): is-linear-map(\allowbreak{}bh-module(\allowbreak{}otv\_), cc-sn-module, otf\_) and is-hom-top-space(\allowbreak{}otu\_, cc-sn-top, otf\_)\}
\end{quote}
characterized by \texttt{op-dual-unfold}  \textit{(proven)}:
\begin{quote}\ttfamily\footnotesize\raggedright\noindent
forall(\allowbreak{}[v, u], op-dual(\allowbreak{}v, u) == \{otf\_ in fun(\allowbreak{}bounded-ops(\allowbreak{}v, v), cc): is-linear-map(\allowbreak{}bh-module(\allowbreak{}v), cc-sn-module, otf\_) and is-hom-top-space(\allowbreak{}u, cc-sn-top, otf\_)\})
\end{quote}
mentioned by 4 result(s)

\smallskip

\noindent\textbf{\texttt{op-iter}}\quad\textit{primitive}\index{op-iter@\texttt{op-iter}}\\
reads: \$2\textasciicircum{}\$4 applied to \$3 \\
defined by \texttt{op-iter-zero} \texttt{op-iter-succ} \\
characterized by \texttt{op-iter-succ}  \textit{(definitional)}:
\begin{quote}\ttfamily\footnotesize\raggedright\noindent
forall(\allowbreak{}[goi-md\_, goi-t\_, goi-w\_, m in nn], op-iter(\allowbreak{}goi-md\_, goi-t\_, goi-w\_, succ(\allowbreak{}m)) == goi-t\_(\allowbreak{}op-iter(\allowbreak{}goi-md\_, goi-t\_, goi-w\_, m)))
\end{quote}
also: \texttt{op-iter-zero} \texttt{op-iter-act} \texttt{op-iter-add} \texttt{op-iter-one} \texttt{op-iter-type} \texttt{poly-op-iter} \texttt{op-iter-shift} \texttt{op-iter-vzero} ... (10 in all) \\
mentioned by 45 result(s)

\smallskip

\noindent\textbf{\texttt{op-kernel(bov\_, bow\_, bot\_)}}\quad\textit{functoid}\index{op-kernel(bov_, bow_, bot_)@\texttt{op-kernel(bov\_, bow\_, bot\_)}}\\
reads: the kernel of \$3 \\
definition:
\begin{quote}\ttfamily\footnotesize\raggedright\noindent
\{box\_ in vec(\allowbreak{}bov\_): bot\_(\allowbreak{}box\_) = vzero(\allowbreak{}bow\_)\}
\end{quote}
characterized by \texttt{op-kernel-subset}  \textit{(proven)}:
\begin{quote}\ttfamily\footnotesize\raggedright\noindent
forall(\allowbreak{}[v, w, t\_], is-complex-inner-product-space(\allowbreak{}v) implies op-kernel(\allowbreak{}v, w, t\_) subset vec(\allowbreak{}v))
\end{quote}
also: \texttt{op-kernel-unfold} \texttt{op-kernel-membership} \\
mentioned by 7 result(s)

\smallskip

\noindent\textbf{\texttt{op-mat(x4md\_, x4n\_, x4u\_, x4t\_)}}\quad\textit{functoid}\index{op-mat(x4md_, x4n_, x4u_, x4t_)@\texttt{op-mat(x4md\_, x4n\_, x4u\_, x4t\_)}}\\
reads: the matrix of \$4 in the basis \$3 \\
definition:
\begin{quote}\ttfamily\footnotesize\raggedright\noindent
iota(\allowbreak{}x4a\_, x4a\_ in mat(\allowbreak{}x4n\_, x4n\_, carr(\allowbreak{}scal(\allowbreak{}x4md\_))) and map-seq(\allowbreak{}x4t\_, x4n\_, x4u\_) = matact(\allowbreak{}x4md\_, x4a\_, x4u\_))
\end{quote}
characterized by \texttt{op-mat-id}  \textit{(proven)}:
\begin{quote}\ttfamily\footnotesize\raggedright\noindent
forall(\allowbreak{}[md, n, u], is-module(\allowbreak{}md) implies n in nn implies 1 <= n implies u in mat(\allowbreak{}n, 1, vec(\allowbreak{}md)) implies generates(\allowbreak{}md, n, u) implies rel-free(\allowbreak{}md, n, u) implies forall(\allowbreak{}[rr\_ in fun(\allowbreak{}vec(\allowbreak{}md), vec(\allowbreak{}md))], forall(\allowbreak{}[x4z in vec(\allowbreak{}md)], rr\_(\allowbreak{}x4z) = x4z) implies op-mat(\allowbreak{}md, n, u, rr\_) = identmat(\allowbreak{}scal(\allowbreak{}md), n)))
\end{quote}
also: \texttt{op-mat-add} \texttt{op-mat-scal} \texttt{op-mat-spec} \texttt{op-mat-unique} \texttt{op-mat-compose} \\
mentioned by 11 result(s)

\smallskip

\noindent\textbf{\texttt{op-mul(gom-md\_, gom-t\_)}}\quad\textit{functoid}\index{op-mul(gom-md_, gom-t_)@\texttt{op-mul(gom-md\_, gom-t\_)}}\\
definition:
\begin{quote}\ttfamily\footnotesize\raggedright\noindent
vnb-lambda(\allowbreak{}[gomf\_, gomg\_], cartesian(\allowbreak{}ops(\allowbreak{}gom-md\_, gom-t\_), ops(\allowbreak{}gom-md\_, gom-t\_)), vnb-lambda(\allowbreak{}gomw\_, vec(\allowbreak{}gom-md\_), gomf\_(\allowbreak{}gomg\_(\allowbreak{}gomw\_))))
\end{quote}
characterized by \texttt{op-mul-of}  \textit{(proven)}:
\begin{quote}\ttfamily\footnotesize\raggedright\noindent
forall(\allowbreak{}[md, tt], is-hom-module(\allowbreak{}md, md, tt) implies is-commutative-ring(\allowbreak{}scal(\allowbreak{}md)) implies forall(\allowbreak{}[p in carr(\allowbreak{}poly(\allowbreak{}scal(\allowbreak{}md))), q in carr(\allowbreak{}poly(\allowbreak{}scal(\allowbreak{}md)))], (\allowbreak{}op-mul(\allowbreak{}md, tt))(\allowbreak{}op-of(\allowbreak{}md, tt, p), op-of(\allowbreak{}md, tt, q)) = op-of(\allowbreak{}md, tt, (\allowbreak{}mul(\allowbreak{}poly(\allowbreak{}scal(\allowbreak{}md))))(\allowbreak{}p, q))))
\end{quote}
mentioned by 5 result(s)

\smallskip

\noindent\textbf{\texttt{op-neg(gon-md\_, gon-t\_)}}\quad\textit{functoid}\index{op-neg(gon-md_, gon-t_)@\texttt{op-neg(gon-md\_, gon-t\_)}}\\
definition:
\begin{quote}\ttfamily\footnotesize\raggedright\noindent
vnb-lambda(\allowbreak{}gonf\_, ops(\allowbreak{}gon-md\_, gon-t\_), vnb-lambda(\allowbreak{}gonw\_, vec(\allowbreak{}gon-md\_), (\allowbreak{}vneg(\allowbreak{}gon-md\_))(\allowbreak{}gonf\_(\allowbreak{}gonw\_))))
\end{quote}
characterized by \texttt{op-neg-of}  \textit{(proven)}:
\begin{quote}\ttfamily\footnotesize\raggedright\noindent
forall(\allowbreak{}[md, tt], is-hom-module(\allowbreak{}md, md, tt) implies is-commutative-ring(\allowbreak{}scal(\allowbreak{}md)) implies forall(\allowbreak{}[p in carr(\allowbreak{}poly(\allowbreak{}scal(\allowbreak{}md)))], (\allowbreak{}op-neg(\allowbreak{}md, tt))(\allowbreak{}op-of(\allowbreak{}md, tt, p)) = op-of(\allowbreak{}md, tt, (\allowbreak{}neg(\allowbreak{}poly(\allowbreak{}scal(\allowbreak{}md))))(\allowbreak{}p))))
\end{quote}
mentioned by 5 result(s)

\smallskip

\noindent\textbf{\texttt{op-of(gof-md\_, gof-t\_, gof-p\_)}}\quad\textit{functoid}\index{op-of(gof-md_, gof-t_, gof-p_)@\texttt{op-of(gof-md\_, gof-t\_, gof-p\_)}}\\
reads: the operator \$3(\$2) \\
definition:
\begin{quote}\ttfamily\footnotesize\raggedright\noindent
vnb-lambda(\allowbreak{}gofw\_, vec(\allowbreak{}gof-md\_), poly-op(\allowbreak{}gof-md\_, gof-t\_, gof-p\_, gofw\_))
\end{quote}
characterized by \texttt{op-of-apply}  \textit{(proven)}:
\begin{quote}\ttfamily\footnotesize\raggedright\noindent
forall(\allowbreak{}[md, tt], is-hom-module(\allowbreak{}md, md, tt) implies is-commutative-ring(\allowbreak{}scal(\allowbreak{}md)) implies forall(\allowbreak{}[p in carr(\allowbreak{}poly(\allowbreak{}scal(\allowbreak{}md))), w in vec(\allowbreak{}md)], (\allowbreak{}op-of(\allowbreak{}md, tt, p))(\allowbreak{}w) = poly-op(\allowbreak{}md, tt, p, w)))
\end{quote}
also: \texttt{op-of-in-ops} \\
mentioned by 18 result(s)

\smallskip

\noindent\textbf{\texttt{op-one(goo-md\_)}}\quad\textit{functoid}\index{op-one(goo-md_)@\texttt{op-one(goo-md\_)}}\\
definition:
\begin{quote}\ttfamily\footnotesize\raggedright\noindent
vnb-lambda(\allowbreak{}goow\_, vec(\allowbreak{}goo-md\_), goow\_)
\end{quote}
characterized by \texttt{op-one-of}  \textit{(proven)}:
\begin{quote}\ttfamily\footnotesize\raggedright\noindent
forall(\allowbreak{}[md, tt], is-hom-module(\allowbreak{}md, md, tt) implies is-commutative-ring(\allowbreak{}scal(\allowbreak{}md)) implies op-one(\allowbreak{}md) = op-of(\allowbreak{}md, tt, one(\allowbreak{}poly(\allowbreak{}scal(\allowbreak{}md)))))
\end{quote}
mentioned by 4 result(s)

\smallskip

\noindent\textbf{\texttt{op-range(bov\_, bow\_, bot\_)}}\quad\textit{functoid}\index{op-range(bov_, bow_, bot_)@\texttt{op-range(bov\_, bow\_, bot\_)}}\\
reads: the range of \$3 \\
definition:
\begin{quote}\ttfamily\footnotesize\raggedright\noindent
\{boy\_ in vec(\allowbreak{}bow\_): forsome(\allowbreak{}[box\_ in vec(\allowbreak{}bov\_)], boy\_ = bot\_(\allowbreak{}box\_))\}
\end{quote}
characterized by \texttt{proj-op-range}  \textit{(proven)}:
\begin{quote}\ttfamily\footnotesize\raggedright\noindent
forall(\allowbreak{}[v, kk], is-hilbert-space(\allowbreak{}v) implies is-submodule(\allowbreak{}v, kk) implies is-closed(\allowbreak{}ip-metric-space(\allowbreak{}v), kk) implies op-range(\allowbreak{}v, v, proj-op(\allowbreak{}v, kk)) = kk)
\end{quote}
also: \texttt{op-range-unfold} \texttt{op-range-membership} \texttt{projection-op-range-closed} \texttt{projection-op-range-perp-kernel} \\
mentioned by 17 result(s)

\smallskip

\noindent\textbf{\texttt{op-sqrt(kqv\_, kqt\_)}}\quad\textit{functoid}\index{op-sqrt(kqv_, kqt_)@\texttt{op-sqrt(kqv\_, kqt\_)}}\\
reads: the positive square root of \$2 \\
definition:
\begin{quote}\ttfamily\footnotesize\raggedright\noindent
iota(\allowbreak{}kqs\_, is-positive-op(\allowbreak{}kqv\_, kqs\_) and bop-comp(\allowbreak{}kqv\_, kqs\_, kqs\_) = kqt\_)
\end{quote}
characterized by \texttt{op-sqrt-unfold}  \textit{(proven)}:
\begin{quote}\ttfamily\footnotesize\raggedright\noindent
forall(\allowbreak{}[kqv\_, kqt\_], op-sqrt(\allowbreak{}kqv\_, kqt\_) == iota(\allowbreak{}kqs\_, is-positive-op(\allowbreak{}kqv\_, kqs\_) and bop-comp(\allowbreak{}kqv\_, kqs\_, kqs\_) = kqt\_))
\end{quote}
mentioned by 4 result(s)

\smallskip

\noindent\textbf{\texttt{op-support(bov\_, bow\_, bot\_)}}\quad\textit{functoid}\index{op-support(bov_, bow_, bot_)@\texttt{op-support(bov\_, bow\_, bot\_)}}\\
reads: the support of \$3 \\
definition:
\begin{quote}\ttfamily\footnotesize\raggedright\noindent
perp(\allowbreak{}bov\_, op-kernel(\allowbreak{}bov\_, bow\_, bot\_))
\end{quote}
mentioned by 4 result(s)

\smallskip

\noindent\textbf{\texttt{op-trace(kqv\_, kqt\_)}}\quad\textit{functoid}\index{op-trace(kqv_, kqt_)@\texttt{op-trace(kqv\_, kqt\_)}}\\
reads: the trace of \$2 \\
definition:
\begin{quote}\ttfamily\footnotesize\raggedright\noindent
cc-family-sum(\allowbreak{}vnb-lambda(\allowbreak{}kqe\_, canon-onb(\allowbreak{}kqv\_), (\allowbreak{}ip(\allowbreak{}kqv\_))(\allowbreak{}kqt\_(\allowbreak{}kqe\_), kqe\_)))
\end{quote}
characterized by \texttt{op-trace-add}  \textit{(proven)}:
\begin{quote}\ttfamily\footnotesize\raggedright\noindent
forall(\allowbreak{}[v], is-hilbert-space(\allowbreak{}v) implies forall(\allowbreak{}[s\_, t\_], is-trace-class(\allowbreak{}v, s\_) implies is-trace-class(\allowbreak{}v, t\_) implies op-trace(\allowbreak{}v, bop-sum(\allowbreak{}v, v, s\_, t\_)) = op-trace(\allowbreak{}v, s\_) + op-trace(\allowbreak{}v, t\_)))
\end{quote}
also: \texttt{op-trace-comm} \texttt{op-trace-sums} \texttt{op-trace-canon} \texttt{op-trace-scale} \texttt{op-trace-factor} \texttt{op-trace-unfold} \texttt{op-trace-adjoint} \texttt{op-trace-comm-bounded} \\
mentioned by 21 result(s)

\smallskip

\noindent\textbf{\texttt{op-zero(goz-md\_)}}\quad\textit{functoid}\index{op-zero(goz-md_)@\texttt{op-zero(goz-md\_)}}\\
definition:
\begin{quote}\ttfamily\footnotesize\raggedright\noindent
vnb-lambda(\allowbreak{}gozw\_, vec(\allowbreak{}goz-md\_), vzero(\allowbreak{}goz-md\_))
\end{quote}
characterized by \texttt{op-zero-of}  \textit{(proven)}:
\begin{quote}\ttfamily\footnotesize\raggedright\noindent
forall(\allowbreak{}[md, tt], is-hom-module(\allowbreak{}md, md, tt) implies is-commutative-ring(\allowbreak{}scal(\allowbreak{}md)) implies op-zero(\allowbreak{}md) = op-of(\allowbreak{}md, tt, zero(\allowbreak{}poly(\allowbreak{}scal(\allowbreak{}md)))))
\end{quote}
mentioned by 4 result(s)

\smallskip

\noindent\textbf{\texttt{open-fun(mxu\_)}}\quad\textit{functoid}\index{open-fun(mxu_)@\texttt{open-fun(mxu\_)}}\\
reads: the map opening an index with \$1 \\
definition:
\begin{quote}\ttfamily\footnotesize\raggedright\noindent
vnb-lambda(\allowbreak{}[mxw\_, mxj\_], cartesian(\allowbreak{}expr-set, nn), expr-open(\allowbreak{}mxw\_, mxj\_, mxu\_))
\end{quote}
characterized by \texttt{open-fun-unfold}  \textit{(proven)}:
\begin{quote}\ttfamily\footnotesize\raggedright\noindent
forall(\allowbreak{}[u], open-fun(\allowbreak{}u) == vnb-lambda(\allowbreak{}[mxw\_, mxj\_], cartesian(\allowbreak{}expr-set, nn), expr-open(\allowbreak{}mxw\_, mxj\_, u)))
\end{quote}
mentioned by 11 result(s)

\smallskip

\noindent\textbf{\texttt{open-stage}}\quad\textit{functoid}\index{open-stage@\texttt{open-stage}}\\
reads: the map opening an index with \$1 on the stage \$2 \\
defined by \texttt{open-stage-zero} \texttt{open-stage-succ} \\
characterized by \texttt{open-stage-succ}  \textit{(definitional)}:
\begin{quote}\ttfamily\footnotesize\raggedright\noindent
forall(\allowbreak{}[mxu\_, mxn\_ in nn], open-stage(\allowbreak{}mxu\_, succ(\allowbreak{}mxn\_)) == vnb-lambda(\allowbreak{}[mxe\_, mxk\_], cartesian(\allowbreak{}expr-stage(\allowbreak{}succ(\allowbreak{}mxn\_)), nn), if(\allowbreak{}mxe\_ in expr-stage(\allowbreak{}mxn\_), (\allowbreak{}open-stage(\allowbreak{}mxu\_, mxn\_))(\allowbreak{}mxe\_, mxk\_), if(\allowbreak{}nth(\allowbreak{}1, mxe\_) = 3, [3, nth(\allowbreak{}2, mxe\_), maplk(\allowbreak{}open-stage(\allowbreak{}mxu\_, mxn\_), nth(\allowbreak{}3, mxe\_), mxk\_)], [4, nth(\allowbreak{}2, mxe\_), nth(\allowbreak{}3, mxe\_), maplk(\allowbreak{}open-stage(\allowbreak{}mxu\_, mxn\_), nth(\allowbreak{}4, mxe\_), mxk\_), (\allowbreak{}open-stage(\allowbreak{}mxu\_, mxn\_))(\allowbreak{}nth(\allowbreak{}5, mxe\_), mxk\_ + nth(\allowbreak{}3, mxe\_))]))))
\end{quote}
also: \texttt{open-stage-zero} \texttt{open-stage-coh} \texttt{open-stage-fun} \texttt{open-stage-step} \\
mentioned by 11 result(s)

\smallskip

\noindent\textbf{\texttt{openl}}\quad\textit{functoid}\index{openl@\texttt{openl}}\\
reads: \$1 opened at the first \$3 entries of \$2 \\
defined by \texttt{openl-zero} \texttt{openl-succ} \\
characterized by \texttt{openl-succ}  \textit{(definitional)}:
\begin{quote}\ttfamily\footnotesize\raggedright\noindent
forall(\allowbreak{}[mcob\_, mcol\_, mcon\_ in nn], openl(\allowbreak{}mcob\_, mcol\_, succ(\allowbreak{}mcon\_)) == expr-open(\allowbreak{}openl(\allowbreak{}mcob\_, mcol\_, mcon\_), mcon\_, nth(\allowbreak{}succ(\allowbreak{}mcon\_), mcol\_)))
\end{quote}
also: \texttt{openl-zero} \\
mentioned by 15 result(s)

\smallskip

\noindent\textbf{\texttt{opens(s)}}\quad\textit{accessor}\index{opens(s)@\texttt{opens(s)}}\\
reads: the topology of \$1 \\
a structure slot \\
characterized by \texttt{cc-hat-top@opens}  \textit{(definitional)}:
\begin{quote}\ttfamily\footnotesize\raggedright\noindent
opens(\allowbreak{}cc-hat-top) == cc-hat-opens
\end{quote}
also: \texttt{ds-opens} \texttt{sntop-opens} \texttt{rb3-pgt-opens} \texttt{prod-top-opens} \texttt{tprod-top-opens} \texttt{subspace-top-opens} \texttt{metric-top-opens-iff} \\
mentioned by 202 result(s)

\smallskip

\noindent\textbf{\texttt{opg-ext(rpg\_, rpa\_, rpp\_)}}\quad\textit{functoid}\index{opg-ext(rpg_, rpa_, rpp_)@\texttt{opg-ext(rpg\_, rpa\_, rpp\_)}}\\
reads: the extension of the segment \$3 of length \$2 in \$1 to a one-parameter group \\
definition:
\begin{quote}\ttfamily\footnotesize\raggedright\noindent
vnb-lambda(\allowbreak{}rvt\_, rr, opg-val(\allowbreak{}rpg\_, rpa\_, rpp\_, rvt\_))
\end{quote}
characterized by \texttt{opg-ext-fun}  \textit{(proven)}:
\begin{quote}\ttfamily\footnotesize\raggedright\noindent
forall(\allowbreak{}[g, al, phi], is-group(\allowbreak{}g) implies al in rr implies 0 < al implies phi in fun(\allowbreak{}ccint(\allowbreak{}0, al), carr(\allowbreak{}g)) implies forall(\allowbreak{}[gsr\_ in ccint(\allowbreak{}0, al), gsu\_ in ccint(\allowbreak{}0, al)], gsr\_ + gsu\_ <= al implies phi(\allowbreak{}gsr\_ + gsu\_) = (\allowbreak{}opr(\allowbreak{}g))(\allowbreak{}phi(\allowbreak{}gsr\_), phi(\allowbreak{}gsu\_))) implies opg-ext(\allowbreak{}g, al, phi) in fun(\allowbreak{}rr, carr(\allowbreak{}g)))
\end{quote}
also: \texttt{opg-ext-apply} \texttt{opg-ext-extends} \\
mentioned by 5 result(s)

\smallskip

\noindent\textbf{\texttt{opg-rep(rpg\_, rpa\_, rpp\_, rpt\_, rpy\_)}}\quad\textit{predicate}\index{opg-rep(rpg_, rpa_, rpp_, rpt_, rpy_)@\texttt{opg-rep(rpg\_, rpa\_, rpp\_, rpt\_, rpy\_)}}\\
reads: \$5 is the value at \$4 of the extension of the segment \$3 of length \$2 in \$1 \\
defined by \texttt{opg-rep} \\
characterized by \texttt{opg-rep}  \textit{(definitional)}:
\begin{quote}\ttfamily\footnotesize\raggedright\noindent
forall(\allowbreak{}[rpg\_, rpa\_, rpp\_, rpt\_, rpy\_], opg-rep(\allowbreak{}rpg\_, rpa\_, rpp\_, rpt\_, rpy\_) iff 0 <= rpt\_ and forsome(\allowbreak{}[rpn\_ in nn, rpr\_ in rr], 0 <= rpr\_ and rpr\_ < rpa\_ and rpt\_ = rpn\_ * rpa\_ + rpr\_ and rpy\_ = (\allowbreak{}opr(\allowbreak{}rpg\_))(\allowbreak{}gpow(\allowbreak{}rpg\_, rpp\_(\allowbreak{}rpa\_), rpn\_), rpp\_(\allowbreak{}rpr\_))) or rpt\_ < 0 and forsome(\allowbreak{}[rpn\_ in nn, rpr\_ in rr], 0 <= rpr\_ and rpr\_ < rpa\_ and 0 - rpt\_ = rpn\_ * rpa\_ + rpr\_ and rpy\_ = (\allowbreak{}inv(\allowbreak{}rpg\_))(\allowbreak{}(\allowbreak{}opr(\allowbreak{}rpg\_))(\allowbreak{}gpow(\allowbreak{}rpg\_, rpp\_(\allowbreak{}rpa\_), rpn\_), rpp\_(\allowbreak{}rpr\_)))))
\end{quote}
also: \texttt{opg-rep-exists} \texttt{opg-rep-unique} \\
mentioned by 8 result(s)

\smallskip

\noindent\textbf{\texttt{opg-val(rpg\_, rpa\_, rpp\_, rvt\_)}}\quad\textit{functoid}\index{opg-val(rpg_, rpa_, rpp_, rvt_)@\texttt{opg-val(rpg\_, rpa\_, rpp\_, rvt\_)}}\\
reads: the value at \$4 of the extension of the segment \$3 of length \$2 in \$1 \\
definition:
\begin{quote}\ttfamily\footnotesize\raggedright\noindent
iota(\allowbreak{}rvy\_, rvy\_ in carr(\allowbreak{}rpg\_) and opg-rep(\allowbreak{}rpg\_, rpa\_, rpp\_, rvt\_, rvy\_))
\end{quote}
characterized by \texttt{opg-val-rep}  \textit{(proven)}:
\begin{quote}\ttfamily\footnotesize\raggedright\noindent
forall(\allowbreak{}[g, al, phi, t], is-group(\allowbreak{}g) implies al in rr implies 0 < al implies phi in fun(\allowbreak{}ccint(\allowbreak{}0, al), carr(\allowbreak{}g)) implies forall(\allowbreak{}[gsr\_ in ccint(\allowbreak{}0, al), gsu\_ in ccint(\allowbreak{}0, al)], gsr\_ + gsu\_ <= al implies phi(\allowbreak{}gsr\_ + gsu\_) = (\allowbreak{}opr(\allowbreak{}g))(\allowbreak{}phi(\allowbreak{}gsr\_), phi(\allowbreak{}gsu\_))) implies t in rr implies opg-val(\allowbreak{}g, al, phi, t) in carr(\allowbreak{}g) and opg-rep(\allowbreak{}g, al, phi, t, opg-val(\allowbreak{}g, al, phi, t)))
\end{quote}
also: \texttt{opg-val-value} \texttt{opg-val-nonneg} \texttt{opg-val-additive} \texttt{opg-val-negative} \texttt{opg-val-continuous} \texttt{opg-val-continuous-gt} \texttt{opg-val-additive-nonneg} \\
mentioned by 15 result(s)

\smallskip

\noindent\textbf{\texttt{opnorm(bov\_, bow\_, bot\_)}}\quad\textit{functoid}\index{opnorm(bov_, bow_, bot_)@\texttt{opnorm(bov\_, bow\_, bot\_)}}\\
reads: the operator norm of \$3 \\
definition:
\begin{quote}\ttfamily\footnotesize\raggedright\noindent
sup(\allowbreak{}\{boa\_ in rr: boa\_ = 0 or forsome(\allowbreak{}[box\_ in vec(\allowbreak{}bov\_)], not(\allowbreak{}ip-norm(\allowbreak{}bov\_, box\_) = 0) and boa\_ = ip-norm(\allowbreak{}bow\_, bot\_(\allowbreak{}box\_)) * recip(\allowbreak{}ip-norm(\allowbreak{}bov\_, box\_)))\})
\end{quote}
characterized by \texttt{opnorm-id}  \textit{(proven)}:
\begin{quote}\ttfamily\footnotesize\raggedright\noindent
forall(\allowbreak{}[v], is-complex-inner-product-space(\allowbreak{}v) implies forsome(\allowbreak{}[bnx\_ in vec(\allowbreak{}v)], not(\allowbreak{}bnx\_ = vzero(\allowbreak{}v))) implies opnorm(\allowbreak{}v, v, bop-id(\allowbreak{}v)) = 1)
\end{quote}
also: \texttt{opnorm-neg} \texttt{opnorm-spec} \texttt{opnorm-zero} \texttt{hs-opnorm-le} \texttt{opnorm-bound} \texttt{opnorm-in-rr} \texttt{opnorm-scale} \texttt{opnorm-approx} ... (17 in all) \\
mentioned by 40 result(s)

\smallskip

\noindent\textbf{\texttt{opp-group(ogg\_)}}\quad\textit{functoid}\index{opp-group(ogg_)@\texttt{opp-group(ogg\_)}}\\
reads: the opposite group of \$1 \\
definition:
\begin{quote}\ttfamily\footnotesize\raggedright\noindent
[carr(\allowbreak{}ogg\_), vnb-lambda(\allowbreak{}[ogp\_, ogq\_], cartesian(\allowbreak{}carr(\allowbreak{}ogg\_), carr(\allowbreak{}ogg\_)), (\allowbreak{}opr(\allowbreak{}ogg\_))(\allowbreak{}ogq\_, ogp\_)), iden(\allowbreak{}ogg\_), inv(\allowbreak{}ogg\_)]
\end{quote}
characterized by \texttt{opp-group-inv}  \textit{(proven)}:
\begin{quote}\ttfamily\footnotesize\raggedright\noindent
forall(\allowbreak{}[g], is-group(\allowbreak{}g) implies inv(\allowbreak{}opp-group(\allowbreak{}g)) == inv(\allowbreak{}g))
\end{quote}
also: \texttt{opp-group-opr} \texttt{opp-group-carr} \texttt{opp-group-iden} \texttt{opp-group-metric} \texttt{opp-group-is-group} \texttt{opp-group-opr-apply} \\
mentioned by 14 result(s)

\smallskip

\noindent\textbf{\texttt{opp-metric-group(omg\_)}}\quad\textit{functoid}\index{opp-metric-group(omg_)@\texttt{opp-metric-group(omg\_)}}\\
reads: the opposite metric group of \$1 \\
definition:
\begin{quote}\ttfamily\footnotesize\raggedright\noindent
make-metric-group(\allowbreak{}metric-group-ms(\allowbreak{}omg\_), opp-group(\allowbreak{}metric-group-group(\allowbreak{}omg\_)))
\end{quote}
mentioned by 19 result(s)

\smallskip

\noindent\textbf{\texttt{opr(s)}}\quad\textit{accessor}\index{opr(s)@\texttt{opr(s)}}\\
reads: the group operation of \$1 \\
a structure slot \\
characterized by \texttt{cc-nag@opr}  \textit{(definitional)}:
\begin{quote}\ttfamily\footnotesize\raggedright\noindent
opr(\allowbreak{}cc-nag) == vnb-lambda(\allowbreak{}[x\_, y\_], cartesian(\allowbreak{}cc, cc), x\_ + y\_)
\end{quote}
also: \texttt{lft-group@opr} \texttt{nn-add-monoid@opr} \texttt{zz-add-monoid@opr} \texttt{raag-opr-uba-as-ring} \texttt{rr-pos-star-add-monoid@opr} \texttt{raag-opr-uba-as-ring-uiba-as-uba} \texttt{fsm-opr} \texttt{nfag-opr} ... (47 in all) \\
mentioned by 613 result(s)

\smallskip

\noindent\textbf{\texttt{ops(gos-md\_, gos-t\_)}}\quad\textit{functoid}\index{ops(gos-md_, gos-t_)@\texttt{ops(gos-md\_, gos-t\_)}}\\
reads: the operators p(\$2), p a polynomial \\
definition:
\begin{quote}\ttfamily\footnotesize\raggedright\noindent
\{gosf\_ in fun(\allowbreak{}vec(\allowbreak{}gos-md\_), vec(\allowbreak{}gos-md\_)): forsome(\allowbreak{}[gosp\_ in carr(\allowbreak{}poly(\allowbreak{}scal(\allowbreak{}gos-md\_)))], gosf\_ = op-of(\allowbreak{}gos-md\_, gos-t\_, gosp\_))\}
\end{quote}
characterized by \texttt{ops-elim}  \textit{(proven)}:
\begin{quote}\ttfamily\footnotesize\raggedright\noindent
forall(\allowbreak{}[md, tt], is-hom-module(\allowbreak{}md, md, tt) implies is-commutative-ring(\allowbreak{}scal(\allowbreak{}md)) implies forall(\allowbreak{}[f in ops(\allowbreak{}md, tt)], f in fun(\allowbreak{}vec(\allowbreak{}md), vec(\allowbreak{}md)) and forsome(\allowbreak{}[gosp\_ in carr(\allowbreak{}poly(\allowbreak{}scal(\allowbreak{}md)))], f = op-of(\allowbreak{}md, tt, gosp\_))))
\end{quote}
also: \texttt{ops-closed} \texttt{ops-in-set} \texttt{ops-unfold} \texttt{op-of-in-ops} \\
mentioned by 13 result(s)

\smallskip

\noindent\textbf{\texttt{ord-le}}\quad\textit{predicate}\index{ord-le@\texttt{ord-le}}\\
a predicate \\
characterized by \texttt{ord-le-refl}  \textit{(primitive)}:
\begin{quote}\ttfamily\footnotesize\raggedright\noindent
forall(\allowbreak{}[alpha in ord], ord-le(\allowbreak{}alpha, alpha))
\end{quote}
also: \texttt{ord-le-total} \texttt{ord-le-trans} \texttt{sup-ord-least} \texttt{ord-le-closure} \texttt{ord-le-antisymm} \texttt{ord-le-nn-compat} \\
mentioned by 20 result(s)

\smallskip

\noindent\textbf{\texttt{ord-lt}}\quad\textit{predicate}\index{ord-lt@\texttt{ord-lt}}\\
a predicate \\
characterized by \texttt{ord-lt-iff}  \textit{(definitional)}:
\begin{quote}\ttfamily\footnotesize\raggedright\noindent
forall(\allowbreak{}[alpha, beta], ord-lt(\allowbreak{}alpha, beta) iff ord-le(\allowbreak{}alpha, beta) and not(\allowbreak{}alpha = beta))
\end{quote}
also: \texttt{ord-lt-succ-cases} \texttt{ord-lt-succ-iff-le} \\
mentioned by 33 result(s)

\smallskip

\noindent\textbf{\texttt{ord-segment}}\quad\textit{operator}\index{ord-segment@\texttt{ord-segment}}\\
a term-forming head \\
characterized by \texttt{ord-segment-succ}  \textit{(primitive)}:
\begin{quote}\ttfamily\footnotesize\raggedright\noindent
forall(\allowbreak{}[alpha in ord, x], x in ord-segment(\allowbreak{}succ\_ord(\allowbreak{}alpha)) iff x in ord-segment(\allowbreak{}alpha) or x = alpha)
\end{quote}
also: \texttt{ord-segment-zero} \texttt{ord-segment-is-set} \texttt{ord-segment-membership} \texttt{ord-segment-self} \texttt{ord-segment-trans} \texttt{ord-segment-insert} \texttt{ord-segment-nn-succ} \texttt{ord-segment-succ-sub} ... (12 in all) \\
mentioned by 236 result(s)

\smallskip

\noindent\textbf{\texttt{orelse}}\quad\textit{tactic}\index{orelse@\texttt{orelse}}\\
\texttt{(orelse t1 t2 ...)} -- Run thunks in order, stop at the first that makes progress (LCF ORELSE). \\
declared in \texttt{tactics-help.scm}

\smallskip

\noindent\textbf{\texttt{originates(ssp\_, ssu\_, ssn\_)}}\quad\textit{predicate}\index{originates(ssp_, ssu_, ssn_)@\texttt{originates(ssp\_, ssu\_, ssn\_)}}\\
reads: \$2 originates on the node \$3 of \$1 \\
defined by \texttt{originates} \\
characterized by \texttt{originates}  \textit{(definitional)}:
\begin{quote}\ttfamily\footnotesize\raggedright\noindent
forall(\allowbreak{}[ssp\_, ssu\_, ssn\_], originates(\allowbreak{}ssp\_, ssu\_, ssn\_) iff is-entry-point(\allowbreak{}ssp\_, ssn\_, \{ssw\_ in msg-set(\allowbreak{}tokens(\allowbreak{}ssp\_), keys(\allowbreak{}ssp\_)): is-subterm(\allowbreak{}tokens(\allowbreak{}ssp\_), keys(\allowbreak{}ssp\_), ssu\_, ssw\_)\}))
\end{quote}
also: \texttt{uniquely-originates} \texttt{originates-elim} \texttt{originates-sent} \texttt{originates-intro} \texttt{originates-unique} \texttt{originates-no-earlier} \texttt{or-serv-key-originates} \texttt{or-init-nonce-originates} ... (15 in all) \\
mentioned by 39 result(s)

\smallskip


\section*{P}
\addcontentsline{toc}{section}{P}

\noindent\textbf{\texttt{pad-ones(m, k, f)}}\quad\textit{functoid}\index{pad-ones(m, k, f)@\texttt{pad-ones(m, k, f)}}\\
reads: the sequence of \$2 ones followed by \$3, of length \$1 \\
definition:
\begin{quote}\ttfamily\footnotesize\raggedright\noindent
vnb-lambda(\allowbreak{}poj\_, interval(\allowbreak{}1, m), if(\allowbreak{}poj\_ <= k, 1, f(\allowbreak{}nn-minus(\allowbreak{}poj\_, k))))
\end{quote}
characterized by \texttt{pad-ones-hi}  \textit{(proven)}:
\begin{quote}\ttfamily\footnotesize\raggedright\noindent
forall(\allowbreak{}[m, k, l, f, j], m in nn implies k in nn implies l in nn implies m = k + l implies f in fun(\allowbreak{}interval(\allowbreak{}1, l), zz) implies j in interval(\allowbreak{}1, m) implies not(\allowbreak{}j <= k) implies (\allowbreak{}pad-ones(\allowbreak{}m, k, f))(\allowbreak{}j) = f(\allowbreak{}nn-minus(\allowbreak{}j, k)))
\end{quote}
also: \texttt{pad-ones-lo} \texttt{pad-ones-type} \texttt{pad-ones-chain} \texttt{pad-ones-nonneg} \\
mentioned by 10 result(s)

\smallskip

\noindent\textbf{\texttt{padic-abs(cpp\_)}}\quad\textit{functoid}\index{padic-abs(cpp_)@\texttt{padic-abs(cpp\_)}}\\
reads: the \$1-adic absolute value on the rationals \\
definition:
\begin{quote}\ttfamily\footnotesize\raggedright\noindent
vnb-lambda(\allowbreak{}cpx\_, qq, iota(\allowbreak{}cpv\_, cpv\_ in rr and padic-rel(\allowbreak{}cpp\_, cpx\_, cpv\_)))
\end{quote}
mentioned by 74 result(s)

\smallskip

\noindent\textbf{\texttt{padic-char(cpp\_)}}\quad\textit{functoid}\index{padic-char(cpp_)@\texttt{padic-char(cpp\_)}}\\
reads: the \$1-adic character x |-> exp(2 pi i \{x\}\_p) of the rationals \\
definition:
\begin{quote}\ttfamily\footnotesize\raggedright\noindent
vnb-lambda(\allowbreak{}cpx\_, qq, iota(\allowbreak{}cpw\_, cpw\_ in cc and forsome(\allowbreak{}[cpr\_ in qq], padic-frac-rep(\allowbreak{}cpp\_, cpx\_, cpr\_) and cpw\_ = cc-exp(\allowbreak{}(\allowbreak{}(\allowbreak{}2 * pi) * 1i) * cpr\_))))
\end{quote}
mentioned by 17 result(s)

\smallskip

\noindent\textbf{\texttt{padic-frac-rep(cpp\_, cpx\_, cpr\_)}}\quad\textit{predicate}\index{padic-frac-rep(cpp_, cpx_, cpr_)@\texttt{padic-frac-rep(cpp\_, cpx\_, cpr\_)}}\\
reads: \$3 is a \$1-adic fractional part of \$2 \\
defined by \texttt{padic-frac-rep} \\
characterized by \texttt{padic-frac-rep}  \textit{(definitional)}:
\begin{quote}\ttfamily\footnotesize\raggedright\noindent
forall(\allowbreak{}[cpp\_, cpx\_, cpr\_], padic-frac-rep(\allowbreak{}cpp\_, cpx\_, cpr\_) iff cpr\_ in qq and forsome(\allowbreak{}[cpz\_ in zz, cpm\_ in nn], cpr\_ * cpp\_ \textasciicircum{} cpm\_ = cpz\_) and (\allowbreak{}padic-abs(\allowbreak{}cpp\_))(\allowbreak{}cpx\_ + -cpr\_) <= 1)
\end{quote}
mentioned by 11 result(s)

\smallskip

\noindent\textbf{\texttt{padic-rel(cpp\_, cpx\_, cpv\_)}}\quad\textit{predicate}\index{padic-rel(cpp_, cpx_, cpv_)@\texttt{padic-rel(cpp\_, cpx\_, cpv\_)}}\\
reads: \$3 is the \$1-adic absolute value of \$2 \\
defined by \texttt{padic-rel} \\
characterized by \texttt{padic-rel}  \textit{(definitional)}:
\begin{quote}\ttfamily\footnotesize\raggedright\noindent
forall(\allowbreak{}[cpp\_, cpx\_, cpv\_], padic-rel(\allowbreak{}cpp\_, cpx\_, cpv\_) iff cpv\_ in rr and (\allowbreak{}cpx\_ = 0 and cpv\_ = 0 or not(\allowbreak{}cpx\_ = 0) and forsome(\allowbreak{}[cpa\_ in zz, cpb\_ in zz], not(\allowbreak{}cpb\_ = 0) and cpx\_ * cpb\_ = cpa\_ and forsome(\allowbreak{}[cpi\_ in nn, cpj\_ in nn], zz-ord-is(\allowbreak{}cpp\_, cpa\_, cpi\_) and zz-ord-is(\allowbreak{}cpp\_, cpb\_, cpj\_) and cpv\_ * 2 \textasciicircum{} cpi\_ = 2 \textasciicircum{} cpj\_))))
\end{quote}
mentioned by 16 result(s)

\smallskip

\noindent\textbf{\texttt{pair}}\quad\textit{operator}\index{pair@\texttt{pair}}\\
a term-forming head \\
characterized by \texttt{pairing}  \textit{(primitive)}:
\begin{quote}\ttfamily\footnotesize\raggedright\noindent
forall(\allowbreak{}[a, b], a in set and b in set implies pair(\allowbreak{}a, b) in set)
\end{quote}
also: \texttt{pairing-membership} \texttt{pair-self-eq} \texttt{pair-self-member} \texttt{sent-received-in-pair} \\
mentioned by 252 result(s)

\smallskip

\noindent\textbf{\texttt{partial-fun}}\quad\textit{operator}\index{partial-fun@\texttt{partial-fun}}\\
a term-forming head \\
characterized by \texttt{partial-fun-membership}  \textit{(primitive)}:
\begin{quote}\ttfamily\footnotesize\raggedright\noindent
forall(\allowbreak{}[a, f], f in partial-fun(\allowbreak{}a) iff forsome(\allowbreak{}[b in power(\allowbreak{}a)], f in fun(\allowbreak{}b)))
\end{quote}
also: \texttt{partial-fun-binary-sethood} \texttt{partial-fun-binary-membership} \\
mentioned by 5 result(s)

\smallskip

\noindent\textbf{\texttt{pbc}}\quad\textit{tactic}\index{pbc@\texttt{pbc}}\\
\texttt{(pbc)} -- Proof by contradiction: assume the goal's negation, prove FALSITY. \\
declared in \texttt{tactics-help.scm}

\smallskip

\noindent\textbf{\texttt{permutations(n)}}\quad\textit{functoid}\index{permutations(n)@\texttt{permutations(n)}}\\
definition:
\begin{quote}\ttfamily\footnotesize\raggedright\noindent
injection(\allowbreak{}ord-segment(\allowbreak{}n), ord-segment(\allowbreak{}n))
\end{quote}
characterized by \texttt{permutations-zero}  \textit{(proven)}:
\begin{quote}\ttfamily\footnotesize\raggedright\noindent
card(\allowbreak{}permutations(\allowbreak{}0)) = succ(\allowbreak{}0)
\end{quote}
mentioned by 4 result(s)

\smallskip

\noindent\textbf{\texttt{perp(v, ks)}}\quad\textit{functoid}\index{perp(v, ks)@\texttt{perp(v, ks)}}\\
reads: the orthogonal complement of \$2 in \$1 \\
definition:
\begin{quote}\ttfamily\footnotesize\raggedright\noindent
\{prp\_ in vec(\allowbreak{}v): forall(\allowbreak{}[prk\_ in ks], (\allowbreak{}ip(\allowbreak{}v))(\allowbreak{}prp\_, prk\_) = 0)\}
\end{quote}
characterized by \texttt{perp-orth}  \textit{(proven)}:
\begin{quote}\ttfamily\footnotesize\raggedright\noindent
forall(\allowbreak{}[v, ks, z\_, k\_], z\_ in perp(\allowbreak{}v, ks) implies k\_ in ks implies (\allowbreak{}ip(\allowbreak{}v))(\allowbreak{}z\_, k\_) = 0)
\end{quote}
also: \texttt{perp-subset} \texttt{perp-unfold} \texttt{perp-is-closed} \texttt{graph-j-of-perp} \texttt{perp-membership} \texttt{perp-of-closure} \texttt{hsum-perp-is-set} \texttt{hsum-perp-subset} ... (22 in all) \\
mentioned by 55 result(s)

\smallskip

\noindent\textbf{\texttt{phi-pm(mhm\_, mhf\_)}}\quad\textit{functoid}\index{phi-pm(mhm_, mhf_)@\texttt{phi-pm(mhm\_, mhf\_)}}\\
reads: A |-> \$2(diam(A)) (in \$1) \\
definition:
\begin{quote}\ttfamily\footnotesize\raggedright\noindent
vnb-lambda(\allowbreak{}mhh\_, power(\allowbreak{}pts(\allowbreak{}mhm\_)), if(\allowbreak{}diam(\allowbreak{}mhm\_, mhh\_) = pos-inf, pos-inf, mhf\_(\allowbreak{}diam(\allowbreak{}mhm\_, mhh\_))))
\end{quote}
characterized by \texttt{phi-pm-mono}  \textit{(proven)}:
\begin{quote}\ttfamily\footnotesize\raggedright\noindent
forall(\allowbreak{}[m, phi, a, b], is-metric-space(\allowbreak{}m) implies is-doubling(\allowbreak{}phi) implies a subset b implies b in power(\allowbreak{}pts(\allowbreak{}m)) implies (\allowbreak{}phi-pm(\allowbreak{}m, phi))(\allowbreak{}a) <= (\allowbreak{}phi-pm(\allowbreak{}m, phi))(\allowbreak{}b))
\end{quote}
also: \texttt{phi-pm-value} \texttt{phi-pm-subspace} \texttt{phi-pm-premeasure} \\
mentioned by 8 result(s)

\smallskip

\noindent\textbf{\texttt{pi}}\quad\textit{defined-fn}\index{pi@\texttt{pi}}\\
defined by \texttt{pi-def}

\smallskip

\noindent\textbf{\texttt{pid}}\quad\textit{refinement}\index{pid@\texttt{pid}}\\
predicate \texttt{is-pid}; same shape as integral-domain

\smallskip

\noindent\textbf{\texttt{pk-inf(s)}}\quad\textit{accessor}\index{pk-inf(s)@\texttt{pk-inf(s)}}\\
reads: the pk-inf slot of \$1 \\
a structure slot \\
mentioned by 71 result(s)

\smallskip

\noindent\textbf{\texttt{pk-space}}\quad\textit{structure}\index{pk-space@\texttt{pk-space}}\\
predicate \texttt{is-pk-space}; slots: pk-inf names pubkey

\smallskip

\noindent\textbf{\texttt{poch}}\quad\textit{functoid}\index{poch@\texttt{poch}}\\
reads: the rising factorial of \$1 of length \$2 \\
defined by \texttt{poch-zero} \texttt{poch-succ} \\
characterized by \texttt{poch-succ}  \textit{(definitional)}:
\begin{quote}\ttfamily\footnotesize\raggedright\noindent
forall(\allowbreak{}[pcz\_, pcn\_ in nn], poch(\allowbreak{}pcz\_, succ(\allowbreak{}pcn\_)) == poch(\allowbreak{}pcz\_, pcn\_) * (\allowbreak{}pcz\_ + pcn\_))
\end{quote}
also: \texttt{poch-zero} \texttt{poch-in-cc} \texttt{poch-shift} \texttt{poch-nonzero} \texttt{gamma-poch-gap} \texttt{gamma-poch-step} \texttt{poch-holomorphic} \\
mentioned by 15 result(s)

\smallskip

\noindent\textbf{\texttt{point-eval(uet\_, uex\_)}}\quad\textit{functoid}\index{point-eval(uet_, uex_)@\texttt{point-eval(uet\_, uex\_)}}\\
reads: the evaluation at \$2 on the continuous functions on \$1 \\
definition:
\begin{quote}\ttfamily\footnotesize\raggedright\noindent
vnb-lambda(\allowbreak{}uef\_, ccont(\allowbreak{}uet\_), uef\_(\allowbreak{}uex\_))
\end{quote}
characterized by \texttt{point-eval-val}  \textit{(proven)}:
\begin{quote}\ttfamily\footnotesize\raggedright\noindent
forall(\allowbreak{}[t, x in pts(\allowbreak{}t), f in ccont(\allowbreak{}t)], (\allowbreak{}point-eval(\allowbreak{}t, x))(\allowbreak{}f) = f(\allowbreak{}x))
\end{quote}
also: \texttt{point-eval-character} \texttt{point-eval-injective} \texttt{ccont-character-is-point-eval} \\
mentioned by 5 result(s)

\smallskip

\noindent\textbf{\texttt{pointwise-bounded(s, fam)}}\quad\textit{predicate}\index{pointwise-bounded(s, fam)@\texttt{pointwise-bounded(s, fam)}}\\
reads: \$2 is pointwise bounded on \$1 \\
defined by \texttt{pointwise-bounded} \\
characterized by \texttt{pointwise-bounded}  \textit{(definitional)}:
\begin{quote}\ttfamily\footnotesize\raggedright\noindent
forall(\allowbreak{}[s, fam], pointwise-bounded(\allowbreak{}s, fam) iff is-metric-space(\allowbreak{}s) and fam in fun(\allowbreak{}nn, fun(\allowbreak{}pts(\allowbreak{}s), rr)) and forall(\allowbreak{}[x in pts(\allowbreak{}s)], forsome(\allowbreak{}[bd in rr], forall(\allowbreak{}[k in nn], abs(\allowbreak{}(\allowbreak{}fam(\allowbreak{}k))(\allowbreak{}x)) <= bd))))
\end{quote}
mentioned by 5 result(s)

\smallskip

\noindent\textbf{\texttt{polar(m, t, a)}}\quad\textit{functoid}\index{polar(m, t, a)@\texttt{polar(m, t, a)}}\\
reads: the polar of \$3 in the topological dual of \$1 with topology \$2 \\
definition:
\begin{quote}\ttfamily\footnotesize\raggedright\noindent
\{plf\_ in top-dual(\allowbreak{}m, t): forall(\allowbreak{}[plx\_ in a], plf\_(\allowbreak{}plx\_) <= 1)\}
\end{quote}
characterized by \texttt{polar-mem}  \textit{(proven)}:
\begin{quote}\ttfamily\footnotesize\raggedright\noindent
forall(\allowbreak{}[m, t, a, z], z in polar(\allowbreak{}m, t, a) iff z in top-dual(\allowbreak{}m, t) and forall(\allowbreak{}[plx\_ in a], z(\allowbreak{}plx\_) <= 1))
\end{quote}
also: \texttt{polar-unfold} \texttt{polar-of-no-failure} \texttt{polar-pointwise-bounded} \texttt{polar-closed-in-line-product} \\
mentioned by 8 result(s)

\smallskip

\noindent\textbf{\texttt{pole-order(pof\_, poa\_)}}\quad\textit{functoid}\index{pole-order(pof_, poa_)@\texttt{pole-order(pof\_, poa\_)}}\\
reads: the order of the pole of \$1 at \$2 \\
definition:
\begin{quote}\ttfamily\footnotesize\raggedright\noindent
iota(\allowbreak{}pon\_, is-pole(\allowbreak{}pof\_, poa\_, pon\_))
\end{quote}
characterized by \texttt{pole-order-value}  \textit{(proven)}:
\begin{quote}\ttfamily\footnotesize\raggedright\noindent
forall(\allowbreak{}[f, a, c4on\_], is-pole(\allowbreak{}f, a, c4on\_) implies pole-order(\allowbreak{}f, a) = c4on\_)
\end{quote}
mentioned by 6 result(s)

\smallskip

\noindent\textbf{\texttt{poly(a)}}\quad\textit{functoid}\index{poly(a)@\texttt{poly(a)}}\\
reads: \$1[x] \\
definition:
\begin{quote}\ttfamily\footnotesize\raggedright\noindent
monalg(\allowbreak{}a, nn-add-monoid)
\end{quote}
characterized by \texttt{is-minimal-poly}  \textit{(definitional)}:
\begin{quote}\ttfamily\footnotesize\raggedright\noindent
forall(\allowbreak{}[gmp-md\_, gmp-t\_, gmp-k\_], is-minimal-poly(\allowbreak{}gmp-md\_, gmp-t\_, gmp-k\_) iff is-monic(\allowbreak{}scal(\allowbreak{}gmp-md\_), gmp-k\_) and forall(\allowbreak{}[gmpw\_ in vec(\allowbreak{}gmp-md\_)], poly-op(\allowbreak{}gmp-md\_, gmp-t\_, gmp-k\_, gmpw\_) = vzero(\allowbreak{}gmp-md\_)) and forall(\allowbreak{}[gmpq\_ in carr(\allowbreak{}poly(\allowbreak{}scal(\allowbreak{}gmp-md\_)))], not(\allowbreak{}gmpq\_ = zero(\allowbreak{}poly(\allowbreak{}scal(\allowbreak{}gmp-md\_)))) implies forall(\allowbreak{}[gmpw\_ in vec(\allowbreak{}gmp-md\_)], poly-op(\allowbreak{}gmp-md\_, gmp-t\_, gmpq\_, gmpw\_) = vzero(\allowbreak{}gmp-md\_)) implies deg(\allowbreak{}scal(\allowbreak{}gmp-md\_), gmp-k\_) <= deg(\allowbreak{}scal(\allowbreak{}gmp-md\_), gmpq\_)))
\end{quote}
also: \texttt{is-annihilating-poly} \texttt{poly-add-at-uba-as-ring} \texttt{poly-neg-at-uba-as-ring} \texttt{poly-carrier-uba-as-ring} \texttt{poly-coef-in-uba-as-ring} \texttt{poly-is-ring-uba-as-ring} \texttt{poly-zero-fun-uba-as-ring} \texttt{poly-mul-unfold-uba-as-ring} ... (109 in all) \\
mentioned by 310 result(s)

\smallskip

\noindent\textbf{\texttt{poly-der}}\quad\textit{functoid}\index{poly-der@\texttt{poly-der}}\\
reads: the derivative of the polygon through \$2 at the nodes \$1 \\
defined by \texttt{poly-der-zero} \texttt{poly-der-succ} \\
characterized by \texttt{poly-der-succ}  \textit{(definitional)}:
\begin{quote}\ttfamily\footnotesize\raggedright\noindent
forall(\allowbreak{}[hmt\_, hmx\_, hmm\_ in nn], poly-der(\allowbreak{}hmt\_, hmx\_, succ(\allowbreak{}hmm\_)) == juxta(\allowbreak{}poly-der(\allowbreak{}hmt\_, hmx\_, hmm\_), aff-der(\allowbreak{}hmx\_(\allowbreak{}succ(\allowbreak{}hmm\_)), hmx\_(\allowbreak{}succ(\allowbreak{}succ(\allowbreak{}hmm\_))), hmt\_(\allowbreak{}succ(\allowbreak{}hmm\_)), hmt\_(\allowbreak{}succ(\allowbreak{}succ(\allowbreak{}hmm\_)))), hmt\_(\allowbreak{}0), hmt\_(\allowbreak{}succ(\allowbreak{}hmm\_)), hmt\_(\allowbreak{}succ(\allowbreak{}succ(\allowbreak{}hmm\_)))))
\end{quote}
also: \texttt{poly-der-zero} \\
mentioned by 10 result(s)

\smallskip

\noindent\textbf{\texttt{poly-eval(a, p, t)}}\quad\textit{functoid}\index{poly-eval(a, p, t)@\texttt{poly-eval(a, p, t)}}\\
reads: the value of \$2 at \$3 \\
definition:
\begin{quote}\ttfamily\footnotesize\raggedright\noindent
sum-ag(\allowbreak{}ring-additive-ag(\allowbreak{}a), vnb-lambda(\allowbreak{}gek\_, nn, (\allowbreak{}mul(\allowbreak{}a))(\allowbreak{}p(\allowbreak{}gek\_), ring-power(\allowbreak{}a, t, gek\_))), succ(\allowbreak{}deg(\allowbreak{}a, p)))
\end{quote}
characterized by \texttt{poly-eval-add}  \textit{(proven)}:
\begin{quote}\ttfamily\footnotesize\raggedright\noindent
forall(\allowbreak{}[a\_], is-commutative-ring(\allowbreak{}a\_) implies forall(\allowbreak{}[p\_ in carr(\allowbreak{}poly(\allowbreak{}a\_)), q\_ in carr(\allowbreak{}poly(\allowbreak{}a\_)), t\_ in carr(\allowbreak{}a\_)], poly-eval(\allowbreak{}a\_, (\allowbreak{}add(\allowbreak{}poly(\allowbreak{}a\_)))(\allowbreak{}p\_, q\_), t\_) = (\allowbreak{}add(\allowbreak{}a\_))(\allowbreak{}poly-eval(\allowbreak{}a\_, p\_, t\_), poly-eval(\allowbreak{}a\_, q\_, t\_))))
\end{quote}
also: \texttt{poly-eval-mul} \texttt{poly-eval-neg} \texttt{poly-eval-type} \texttt{poly-eval-zero} \texttt{poly-eval-bound} \texttt{poly-eval-constant} \texttt{poly-eval-monomial} \texttt{poly-eval-x-minus-mul} ... (10 in all) \\
mentioned by 34 result(s)

\smallskip

\noindent\textbf{\texttt{poly-fn-deg(kk, a, f)}}\quad\textit{functoid}\index{poly-fn-deg(kk, a, f)@\texttt{poly-fn-deg(kk, a, f)}}\\
reads: the degree of the polynomial function \$3 on \$2 \\
definition:
\begin{quote}\ttfamily\footnotesize\raggedright\noindent
iota(\allowbreak{}pdgn\_, pdgn\_ in nn and is-poly-deg(\allowbreak{}kk, a, f, pdgn\_))
\end{quote}
mentioned by 4 result(s)

\smallskip

\noindent\textbf{\texttt{poly-le(a, m)}}\quad\textit{functoid}\index{poly-le(a, m)@\texttt{poly-le(a, m)}}\\
reads: the polynomials over \$1 of degree at most \$2 \\
definition:
\begin{quote}\ttfamily\footnotesize\raggedright\noindent
\{gpl\_ in carr(\allowbreak{}poly(\allowbreak{}a)): deg-bound(\allowbreak{}a, gpl\_, m)\}
\end{quote}
characterized by \texttt{poly-le-unfold}  \textit{(proven)}:
\begin{quote}\ttfamily\footnotesize\raggedright\noindent
forall(\allowbreak{}[a\_, m\_], poly-le(\allowbreak{}a\_, m\_) == \{gpl\_ in carr(\allowbreak{}poly(\allowbreak{}a\_)): deg-bound(\allowbreak{}a\_, gpl\_, m\_)\})
\end{quote}
mentioned by 5 result(s)

\smallskip

\noindent\textbf{\texttt{poly-op(gpo-md\_, gpo-t\_, gpo-p\_, gpo-w\_)}}\quad\textit{functoid}\index{poly-op(gpo-md_, gpo-t_, gpo-p_, gpo-w_)@\texttt{poly-op(gpo-md\_, gpo-t\_, gpo-p\_, gpo-w\_)}}\\
reads: \$3(\$2) applied to \$4 \\
definition:
\begin{quote}\ttfamily\footnotesize\raggedright\noindent
sum-ag(\allowbreak{}module-vector-ag(\allowbreak{}gpo-md\_), vnb-lambda(\allowbreak{}gpok\_, nn, (\allowbreak{}act(\allowbreak{}gpo-md\_))(\allowbreak{}gpo-p\_(\allowbreak{}gpok\_), op-iter(\allowbreak{}gpo-md\_, gpo-t\_, gpo-w\_, gpok\_))), succ(\allowbreak{}deg(\allowbreak{}scal(\allowbreak{}gpo-md\_), gpo-p\_)))
\end{quote}
characterized by \texttt{poly-op-act}  \textit{(proven)}:
\begin{quote}\ttfamily\footnotesize\raggedright\noindent
forall(\allowbreak{}[md, tt], is-hom-module(\allowbreak{}md, md, tt) implies is-commutative-ring(\allowbreak{}scal(\allowbreak{}md)) implies forall(\allowbreak{}[p in carr(\allowbreak{}poly(\allowbreak{}scal(\allowbreak{}md))), c in carr(\allowbreak{}scal(\allowbreak{}md)), x in vec(\allowbreak{}md)], poly-op(\allowbreak{}md, tt, p, (\allowbreak{}act(\allowbreak{}md))(\allowbreak{}c, x)) = (\allowbreak{}act(\allowbreak{}md))(\allowbreak{}c, poly-op(\allowbreak{}md, tt, p, x))))
\end{quote}
also: \texttt{poly-op-add} \texttt{poly-op-mul} \texttt{poly-op-iter} \texttt{poly-op-padd} \texttt{poly-op-pneg} \texttt{poly-op-type} \texttt{poly-op-vneg} \texttt{poly-op-bound} ... (18 in all) \\
mentioned by 60 result(s)

\smallskip

\noindent\textbf{\texttt{poly-road}}\quad\textit{functoid}\index{poly-road@\texttt{poly-road}}\\
reads: the polygon through the points \$2 at the nodes \$1, with \$3 + 1 pieces \\
defined by \texttt{poly-road-zero} \texttt{poly-road-succ} \\
characterized by \texttt{poly-road-succ}  \textit{(definitional)}:
\begin{quote}\ttfamily\footnotesize\raggedright\noindent
forall(\allowbreak{}[hmt\_, hmx\_, hmm\_ in nn], poly-road(\allowbreak{}hmt\_, hmx\_, succ(\allowbreak{}hmm\_)) == juxta(\allowbreak{}poly-road(\allowbreak{}hmt\_, hmx\_, hmm\_), aff-seg(\allowbreak{}hmx\_(\allowbreak{}succ(\allowbreak{}hmm\_)), hmx\_(\allowbreak{}succ(\allowbreak{}succ(\allowbreak{}hmm\_))), hmt\_(\allowbreak{}succ(\allowbreak{}hmm\_)), hmt\_(\allowbreak{}succ(\allowbreak{}succ(\allowbreak{}hmm\_)))), hmt\_(\allowbreak{}0), hmt\_(\allowbreak{}succ(\allowbreak{}hmm\_)), hmt\_(\allowbreak{}succ(\allowbreak{}succ(\allowbreak{}hmm\_)))))
\end{quote}
also: \texttt{poly-road-zero} \texttt{poly-road-facts} \\
mentioned by 11 result(s)

\smallskip

\noindent\textbf{\texttt{pos-cone(pca\_)}}\quad\textit{functoid}\index{pos-cone(pca_)@\texttt{pos-cone(pca\_)}}\\
reads: the positive cone of \$1 \\
definition:
\begin{quote}\ttfamily\footnotesize\raggedright\noindent
\{pcx\_ in vec(\allowbreak{}pca\_): forall(\allowbreak{}[pcc\_ in power(\allowbreak{}vec(\allowbreak{}pca\_))], forall(\allowbreak{}[pcy\_ in vec(\allowbreak{}pca\_)], (\allowbreak{}vmul(\allowbreak{}pca\_))(\allowbreak{}(\allowbreak{}vstar(\allowbreak{}pca\_))(\allowbreak{}pcy\_), pcy\_) in pcc\_) implies forall(\allowbreak{}[pcu\_ in pcc\_, pcv\_ in pcc\_], (\allowbreak{}vadd(\allowbreak{}pca\_))(\allowbreak{}pcu\_, pcv\_) in pcc\_) implies pcx\_ in pcc\_)\}
\end{quote}
characterized by \texttt{pos-cone-unfold}  \textit{(proven)}:
\begin{quote}\ttfamily\footnotesize\raggedright\noindent
forall(\allowbreak{}[a], pos-cone(\allowbreak{}a) == \{pcx\_ in vec(\allowbreak{}a): forall(\allowbreak{}[pcc\_ in power(\allowbreak{}vec(\allowbreak{}a))], forall(\allowbreak{}[pcy\_ in vec(\allowbreak{}a)], (\allowbreak{}vmul(\allowbreak{}a))(\allowbreak{}(\allowbreak{}vstar(\allowbreak{}a))(\allowbreak{}pcy\_), pcy\_) in pcc\_) implies forall(\allowbreak{}[pcu\_ in pcc\_, pcv\_ in pcc\_], (\allowbreak{}vadd(\allowbreak{}a))(\allowbreak{}pcu\_, pcv\_) in pcc\_) implies pcx\_ in pcc\_)\})
\end{quote}
also: \texttt{ccont-nonneg-pos-cone} \texttt{ccont-pos-cone-nonneg} \\
mentioned by 6 result(s)

\smallskip

\noindent\textbf{\texttt{pos-part(omega, f)}}\quad\textit{functoid}\index{pos-part(omega, f)@\texttt{pos-part(omega, f)}}\\
reads: the positive part of \$2 on \$1 \\
definition:
\begin{quote}\ttfamily\footnotesize\raggedright\noindent
vnb-lambda(\allowbreak{}ppx\_, omega, if(\allowbreak{}0 <= f(\allowbreak{}ppx\_), f(\allowbreak{}ppx\_), 0))
\end{quote}
characterized by \texttt{pos-part-fun}  \textit{(proven)}:
\begin{quote}\ttfamily\footnotesize\raggedright\noindent
forall(\allowbreak{}[omega, f in fun(\allowbreak{}omega, rr)], pos-part(\allowbreak{}omega, f) in fun(\allowbreak{}omega, rr-pos-star))
\end{quote}
also: \texttt{pos-part-value} \texttt{pos-part-unfold} \texttt{pfw-pos-part-comp} \texttt{pos-part-measurable-of-diff} \texttt{pos-part-measurable-of-tails} \\
mentioned by 32 result(s)

\smallskip

\noindent\textbf{\texttt{pos-rr(r)}}\quad\textit{predicate}\index{pos-rr(r)@\texttt{pos-rr(r)}}\\
reads: a positive real \\
defined by \texttt{pos-rr} \\
characterized by \texttt{pos-rr}  \textit{(definitional)}:
\begin{quote}\ttfamily\footnotesize\raggedright\noindent
forall(\allowbreak{}[r], pos-rr(\allowbreak{}r) iff r in rr and 0 <= r and not(\allowbreak{}0 = r))
\end{quote}
also: \texttt{rr-lt-of-pos-rr} \texttt{rr-pos-rr-in-rr} \texttt{rr-pos-rr-of-lt} \\
mentioned by 619 result(s)

\smallskip

\noindent\textbf{\texttt{pos-variation(ca, nu)}}\quad\textit{functoid}\index{pos-variation(ca, nu)@\texttt{pos-variation(ca, nu)}}\\
reads: the positive variation of \$2 \\
definition:
\begin{quote}\ttfamily\footnotesize\raggedright\noindent
vnb-lambda(\allowbreak{}pva\_, ca, esup(\allowbreak{}\{pvv\_ in rr-pos-star: forsome(\allowbreak{}[pve\_ in ca], pve\_ subset pva\_ and nu(\allowbreak{}pve\_) = pvv\_)\}))
\end{quote}
characterized by \texttt{pos-variation-ge}  \textit{(proven)}:
\begin{quote}\ttfamily\footnotesize\raggedright\noindent
forall(\allowbreak{}[ca, nu, a], nu in fun(\allowbreak{}ca, rr-star) implies a in ca implies nu(\allowbreak{}a) <= (\allowbreak{}pos-variation(\allowbreak{}ca, nu))(\allowbreak{}a))
\end{quote}
also: \texttt{pos-variation-in} \texttt{pos-variation-fun} \texttt{pos-variation-mono} \texttt{pos-variation-step} \texttt{pos-variation-chain} \texttt{pos-variation-empty} \texttt{pos-variation-least} \texttt{pos-variation-upper} ... (15 in all) \\
mentioned by 31 result(s)

\smallskip

\noindent\textbf{\texttt{positively-oriented(poa\_, pob\_, poc\_)}}\quad\textit{predicate}\index{positively-oriented(poa_, pob_, poc_)@\texttt{positively-oriented(poa\_, pob\_, poc\_)}}\\
reads: the triangle \$1 \$2 \$3 is positively oriented \\
defined by \texttt{positively-oriented} \\
characterized by \texttt{positively-oriented}  \textit{(definitional)}:
\begin{quote}\ttfamily\footnotesize\raggedright\noindent
forall(\allowbreak{}[poa\_, pob\_, poc\_], positively-oriented(\allowbreak{}poa\_, pob\_, poc\_) iff forall(\allowbreak{}[pou\_ in interior(\allowbreak{}cc-ms, conv3(\allowbreak{}poa\_, pob\_, poc\_))], angle(\allowbreak{}pob\_ - pou\_, poa\_ - pou\_) < pi and angle(\allowbreak{}poc\_ - pou\_, pob\_ - pou\_) < pi and angle(\allowbreak{}poa\_ - pou\_, poc\_ - pou\_) < pi))
\end{quote}
also: \texttt{positively-oriented-iff} \\
mentioned by 15 result(s)

\smallskip

\noindent\textbf{\texttt{power}}\quad\textit{operator}\index{power@\texttt{power}}\\
a term-forming head \\
characterized by \texttt{good-sub-in-power-cnvs-as-nvs}  \textit{(definitional)}:
\begin{quote}\ttfamily\footnotesize\raggedright\noindent
forall(\allowbreak{}[r], is-complex-normed-vector-space(\allowbreak{}r) implies forall(\allowbreak{}[s, f, t], good-sub(\allowbreak{}cnvs-as-nvs(\allowbreak{}r), s, f, t) implies t in power(\allowbreak{}vec(\allowbreak{}r))))
\end{quote}
also: \texttt{power-neg} \texttt{power-set} \texttt{power-succ} \texttt{power-zero} \texttt{power-typing-nonneg} \texttt{power-set-membership} \texttt{qo-power} \texttt{deriv-power} ... (149 in all) \\
mentioned by 1021 result(s)

\smallskip

\noindent\textbf{\texttt{powers-union(pug\_, puv\_)}}\quad\textit{functoid}\index{powers-union(pug_, puv_)@\texttt{powers-union(pug\_, puv\_)}}\\
reads: the union of the powers of \$2 in \$1 \\
definition:
\begin{quote}\ttfamily\footnotesize\raggedright\noindent
big-union(\allowbreak{}puk\_, nn, set-power(\allowbreak{}pug\_, puv\_, puk\_))
\end{quote}
characterized by \texttt{powers-union-inv}  \textit{(proven)}:
\begin{quote}\ttfamily\footnotesize\raggedright\noindent
forall(\allowbreak{}[k in nn, g, v], is-group(\allowbreak{}g) implies v subset carr(\allowbreak{}g) implies forall(\allowbreak{}[p in v], (\allowbreak{}inv(\allowbreak{}g))(\allowbreak{}p) in v) implies forall(\allowbreak{}[y in set-power(\allowbreak{}g, v, k)], (\allowbreak{}inv(\allowbreak{}g))(\allowbreak{}y) in powers-union(\allowbreak{}g, v)))
\end{quote}
also: \texttt{powers-union-elim} \texttt{powers-union-iden} \texttt{powers-union-open} \texttt{powers-union-prod} \texttt{powers-union-intro} \texttt{powers-union-is-set} \texttt{powers-union-unfold} \texttt{powers-union-contains} ... (12 in all) \\
mentioned by 20 result(s)

\smallskip

\noindent\textbf{\texttt{pprod}}\quad\textit{functoid}\index{pprod@\texttt{pprod}}\\
reads: the partial product of the first \$2 terms of \$1 \\
defined by \texttt{pprod-zero} \texttt{pprod-succ} \\
characterized by \texttt{pprod-succ}  \textit{(definitional)}:
\begin{quote}\ttfamily\footnotesize\raggedright\noindent
forall(\allowbreak{}[ppf\_, ppn\_ in nn], pprod(\allowbreak{}ppf\_, succ(\allowbreak{}ppn\_)) == pprod(\allowbreak{}ppf\_, ppn\_) * ppf\_(\allowbreak{}ppn\_))
\end{quote}
also: \texttt{pprod-zero} \texttt{log-pprod} \texttt{pprod-dec} \texttt{pprod-pos} \texttt{pprod-in-cc} \texttt{pprod-in-rr} \texttt{pprod-lower} \texttt{pprod-split} ... (23 in all) \\
mentioned by 67 result(s)

\smallskip

\noindent\textbf{\texttt{pre-hilbert-space}}\quad\textit{refinement}\index{pre-hilbert-space@\texttt{pre-hilbert-space}}\\
predicate \texttt{is-pre-hilbert-space}; same shape as semi-inner-product-space

\smallskip

\noindent\textbf{\texttt{preamble}}\quad\textit{tactic}\index{preamble@\texttt{preamble}}\\
\texttt{(preamble)  |  (preamble 'NAME)  |  (preamble "file.pre")} -- Apply an editable RULE FILE to the focus goal, one step at a time: each rule reads `if the sequent looks like this, do that', and the first rule that matches AND makes progress is kept; the loop repeats on the new focus until the goal is done, no rule applies, or 40 steps.  Rules come from \textasciitilde{}/.vnb-preamble.pre when it exists, else the shipped preambles/default.pre (peel, split, close by context, reflexivity, arithmetic, type-term, zero-it, ineq, beta, unfold, and a tentative cut).  Prints one report: what each rule did, which rules matched but were rejected, and where it stopped and why. \\
declared in \texttt{tactics-help.scm}

\smallskip

\noindent\textbf{\texttt{prec-strict(spc\_, snd\_, ssd\_)}}\quad\textit{functoid}\index{prec-strict(spc_, snd_, ssd_)@\texttt{prec-strict(spc\_, snd\_, ssd\_)}}\\
reads: the strict causal precedence on the nodes \$2 of \$1 generated by the edges \$3 \\
definition:
\begin{quote}\ttfamily\footnotesize\raggedright\noindent
\{ssp\_ in cartesian(\allowbreak{}snd\_, snd\_): forall(\allowbreak{}[ssq\_ in power(\allowbreak{}cartesian(\allowbreak{}snd\_, snd\_))], is-prec-closed(\allowbreak{}spc\_, snd\_, ssd\_, ssq\_) implies ssp\_ in ssq\_)\}
\end{quote}
characterized by \texttt{prec-strict-unfold}  \textit{(proven)}:
\begin{quote}\ttfamily\footnotesize\raggedright\noindent
forall(\allowbreak{}[spa, nd, ed], prec-strict(\allowbreak{}spa, nd, ed) == \{ssp\_ in cartesian(\allowbreak{}nd, nd): forall(\allowbreak{}[ssq\_ in power(\allowbreak{}cartesian(\allowbreak{}nd, nd))], is-prec-closed(\allowbreak{}spa, nd, ed, ssq\_) implies ssp\_ in ssq\_)\})
\end{quote}
mentioned by 14 result(s)

\smallskip

\noindent\textbf{\texttt{precedes(ssc\_)}}\quad\textit{functoid}\index{precedes(ssc_)@\texttt{precedes(ssc\_)}}\\
reads: the causal precedence of the bundle \$1 \\
definition:
\begin{quote}\ttfamily\footnotesize\raggedright\noindent
union(\allowbreak{}precedes-strict(\allowbreak{}ssc\_), \{ssp\_ in cartesian(\allowbreak{}bnodes(\allowbreak{}ssc\_), bnodes(\allowbreak{}ssc\_)): nth(\allowbreak{}1, ssp\_) = nth(\allowbreak{}2, ssp\_)\})
\end{quote}
characterized by \texttt{precedes-set}  \textit{(proven)}:
\begin{quote}\ttfamily\footnotesize\raggedright\noindent
forall(\allowbreak{}[c], is-bundle(\allowbreak{}c) implies precedes(\allowbreak{}c) in set)
\end{quote}
also: \texttt{precedes-elim} \texttt{precedes-refl} \texttt{precedes-trans} \texttt{precedes-unfold} \texttt{precedes-antisym} \texttt{precedes-intro-strict} \texttt{precedes-partial-order} \texttt{bundle-minimal-precedes} \\
mentioned by 11 result(s)

\smallskip

\noindent\textbf{\texttt{precedes-strict(ssc\_)}}\quad\textit{functoid}\index{precedes-strict(ssc_)@\texttt{precedes-strict(ssc\_)}}\\
reads: the strict causal precedence of the bundle \$1 \\
definition:
\begin{quote}\ttfamily\footnotesize\raggedright\noindent
prec-strict(\allowbreak{}bspace(\allowbreak{}ssc\_), bnodes(\allowbreak{}ssc\_), bsend(\allowbreak{}ssc\_))
\end{quote}
characterized by \texttt{precedes-strict-trans}  \textit{(proven)}:
\begin{quote}\ttfamily\footnotesize\raggedright\noindent
forall(\allowbreak{}[c, x, y, z], is-bundle(\allowbreak{}c) implies [x, y] in precedes-strict(\allowbreak{}c) implies [y, z] in precedes-strict(\allowbreak{}c) implies [x, z] in precedes-strict(\allowbreak{}c))
\end{quote}
also: \texttt{precedes-strict-unfold} \texttt{precedes-strict-irreflexive} \\
mentioned by 32 result(s)

\smallskip

\noindent\textbf{\texttt{pred(n\_)}}\quad\textit{functoid}\index{pred(n_)@\texttt{pred(n\_)}}\\
reads: the predecessor of \$1 \\
definition:
\begin{quote}\ttfamily\footnotesize\raggedright\noindent
iota(\allowbreak{}m\_, m\_ in nn and succ(\allowbreak{}m\_) = n\_)
\end{quote}
characterized by \texttt{pred-succ}  \textit{(proven)}:
\begin{quote}\ttfamily\footnotesize\raggedright\noindent
forall(\allowbreak{}[n\_ in nn], not(\allowbreak{}n\_ = 0) implies succ(\allowbreak{}pred(\allowbreak{}n\_)) = n\_)
\end{quote}
also: \texttt{pred-in-nn} \texttt{pred-of-succ} \texttt{pred-property} \\
mentioned by 8 result(s)

\smallskip

\noindent\textbf{\texttt{pred-code(mcpk\_)}}\quad\textit{predicate}\index{pred-code(mcpk_)@\texttt{pred-code(mcpk\_)}}\\
reads: the code \$1 names a predicate constant \\
defined by \texttt{pred-code} \\
characterized by \texttt{pred-code}  \textit{(definitional)}:
\begin{quote}\ttfamily\footnotesize\raggedright\noindent
forall(\allowbreak{}[mcpk\_], pred-code(\allowbreak{}mcpk\_) iff mcpk\_ = 7 or mcpk\_ = 9 or mcpk\_ = 11 or mcpk\_ = 13)
\end{quote}
also: \texttt{pred-code-nn} \\
mentioned by 26 result(s)

\smallskip

\noindent\textbf{\texttt{pred-def(mcdk\_, mcdn\_, mcdb\_)}}\quad\textit{predicate}\index{pred-def(mcdk_, mcdn_, mcdb_)@\texttt{pred-def(mcdk\_, mcdn\_, mcdb\_)}}\\
reads: the predicate code \$1 at an atom of list length \$2 has the definition body \$3 \\
defined by \texttt{pred-def} \\
characterized by \texttt{pred-def}  \textit{(definitional)}:
\begin{quote}\ttfamily\footnotesize\raggedright\noindent
forall(\allowbreak{}[mcdk\_, mcdn\_, mcdb\_], pred-def(\allowbreak{}mcdk\_, mcdn\_, mcdb\_) iff mcdk\_ in nn and mcdn\_ in nn and mcdb\_ in expr-set and 0 = 1)
\end{quote}
mentioned by 4 result(s)

\smallskip

\noindent\textbf{\texttt{prefix-dist(ppa\_, pdx\_, pdy\_)}}\quad\textit{functoid}\index{prefix-dist(ppa_, pdx_, pdy_)@\texttt{prefix-dist(ppa\_, pdx\_, pdy\_)}}\\
reads: the \$1-weighted prefix distance of \$2 and \$3 \\
definition:
\begin{quote}\ttfamily\footnotesize\raggedright\noindent
if(\allowbreak{}pdx\_ = pdy\_, 0, recip(\allowbreak{}prefix-prod(\allowbreak{}ppa\_, pdx\_, agree-length(\allowbreak{}pdx\_, pdy\_))))
\end{quote}
characterized by \texttt{prefix-dist-sym}  \textit{(proven)}:
\begin{quote}\ttfamily\footnotesize\raggedright\noindent
forall(\allowbreak{}[sx\_, a, xs, ys], a in fun(\allowbreak{}sx\_, rr) implies forall(\allowbreak{}[w in sx\_], 1 < a(\allowbreak{}w)) implies xs in fun(\allowbreak{}nn, sx\_) implies ys in fun(\allowbreak{}nn, sx\_) implies prefix-dist(\allowbreak{}a, xs, ys) = prefix-dist(\allowbreak{}a, ys, xs))
\end{quote}
also: \texttt{prefix-dist-self} \texttt{prefix-dist-large} \texttt{prefix-dist-small} \texttt{prefix-dist-ultra} \texttt{prefix-dist-value} \texttt{prefix-dist-nonneg} \\
mentioned by 15 result(s)

\smallskip

\noindent\textbf{\texttt{prefix-prod}}\quad\textit{functoid}\index{prefix-prod@\texttt{prefix-prod}}\\
reads: the product of \$1 along the first \$3 coordinates of \$2 \\
defined by \texttt{prefix-prod-zero} \texttt{prefix-prod-succ} \\
characterized by \texttt{prefix-prod-succ}  \textit{(definitional)}:
\begin{quote}\ttfamily\footnotesize\raggedright\noindent
forall(\allowbreak{}[ppa\_, ppx\_, ppn\_ in nn], prefix-prod(\allowbreak{}ppa\_, ppx\_, succ(\allowbreak{}ppn\_)) == prefix-prod(\allowbreak{}ppa\_, ppx\_, ppn\_) * ppa\_(\allowbreak{}ppx\_(\allowbreak{}ppn\_)))
\end{quote}
also: \texttt{prefix-prod-zero} \texttt{prefix-prod-mono} \texttt{prefix-prod-ge-one} \texttt{prefix-prod-agree-ind} \texttt{prefix-prod-ge-one-ind} \\
mentioned by 21 result(s)

\smallskip

\noindent\textbf{\texttt{preimage(s, f, v)}}\quad\textit{functoid}\index{preimage(s, f, v)@\texttt{preimage(s, f, v)}}\\
definition:
\begin{quote}\ttfamily\footnotesize\raggedright\noindent
\{a in pts(\allowbreak{}s): f(\allowbreak{}a) in v\}
\end{quote}
characterized by \texttt{preimage-membership}  \textit{(definitional)}:
\begin{quote}\ttfamily\footnotesize\raggedright\noindent
forall(\allowbreak{}[s, f, v, a], a in preimage(\allowbreak{}s, f, v) iff a in pts(\allowbreak{}s) and f(\allowbreak{}a) in v)
\end{quote}
also: \texttt{preimage-in-power} \texttt{metric-top@preimage} \texttt{preimage-complement} \texttt{trunc-metric@preimage} \texttt{image-of-preimage-onto} \texttt{preimage-subspace-open} \texttt{local-continuity-preimage-open} \texttt{continuous-implies-open-preimage} ... (11 in all) \\
mentioned by 32 result(s)

\smallskip

\noindent\textbf{\texttt{prep}}\quad\textit{tactic}\index{prep@\texttt{prep}}\\
\texttt{(prep 'ineq)} -- Why will a tactic not fire here, and what must be done first?  Runs the tactic's OWN preconditions one at a time, marks each ok / PREP / STOP, names the library lemma that repairs each unmet one, and prints a PLAN it has checked by running it on a scratch clone.  READ-ONLY: the live proof is never touched.  (prep 'ineq) is the only method so far. \\
declared in \texttt{tactics-help.scm}

\smallskip

\noindent\textbf{\texttt{primitive-roots(prn\_)}}\quad\textit{functoid}\index{primitive-roots(prn_)@\texttt{primitive-roots(prn\_)}}\\
reads: the primitive \$1-th roots of unity \\
definition:
\begin{quote}\ttfamily\footnotesize\raggedright\noindent
\{prz\_ in cc: prz\_ \textasciicircum{} prn\_ = 1 and forall(\allowbreak{}[prm\_ in nn], 1 <= prm\_ implies prm\_ < prn\_ implies not(\allowbreak{}prz\_ \textasciicircum{} prm\_ = 1))\}
\end{quote}
characterized by \texttt{euler-primitive-roots}  \textit{(proven)}:
\begin{quote}\ttfamily\footnotesize\raggedright\noindent
forall(\allowbreak{}[n in nn], 1 <= n implies card(\allowbreak{}primitive-roots(\allowbreak{}n)) = euler-phi(\allowbreak{}n))
\end{quote}
mentioned by 6 result(s)

\smallskip

\noindent\textbf{\texttt{principal-ideal(s, a)}}\quad\textit{functoid}\index{principal-ideal(s, a)@\texttt{principal-ideal(s, a)}}\\
definition:
\begin{quote}\ttfamily\footnotesize\raggedright\noindent
\{x in carr(\allowbreak{}s): forsome(\allowbreak{}[r in carr(\allowbreak{}s)], x = (\allowbreak{}mul(\allowbreak{}s))(\allowbreak{}r, a))\}
\end{quote}
characterized by \texttt{principal-ideal-membership}  \textit{(definitional)}:
\begin{quote}\ttfamily\footnotesize\raggedright\noindent
forall(\allowbreak{}[s, a, x], x in principal-ideal(\allowbreak{}s, a) iff x in carr(\allowbreak{}s) and forsome(\allowbreak{}[r in carr(\allowbreak{}s)], x = (\allowbreak{}mul(\allowbreak{}s))(\allowbreak{}r, a)))
\end{quote}
also: \texttt{zz-principal-ideal-in} \texttt{zz-principal-ideal-out} \texttt{principal-ideal-in-ideal} \texttt{principal-ideal-is-submodule} \\
mentioned by 34 result(s)

\smallskip

\noindent\textbf{\texttt{prod-abs-converges(paa\_)}}\quad\textit{predicate}\index{prod-abs-converges(paa_)@\texttt{prod-abs-converges(paa\_)}}\\
reads: the product of the factors 1 + \$1(k) is absolutely convergent \\
defined by \texttt{prod-abs-converges} \\
characterized by \texttt{prod-abs-converges}  \textit{(definitional)}:
\begin{quote}\ttfamily\footnotesize\raggedright\noindent
forall(\allowbreak{}[paa\_], prod-abs-converges(\allowbreak{}paa\_) iff paa\_ in fun(\allowbreak{}nn, cc) and series-converges(\allowbreak{}vnb-lambda(\allowbreak{}pak\_, nn, magnitude(\allowbreak{}paa\_(\allowbreak{}pak\_)))))
\end{quote}
mentioned by 3 result(s)

\smallskip

\noindent\textbf{\texttt{prod-converges(pdf\_)}}\quad\textit{predicate}\index{prod-converges(pdf_)@\texttt{prod-converges(pdf\_)}}\\
reads: the infinite product of \$1 converges \\
defined by \texttt{prod-converges} \\
characterized by \texttt{prod-converges}  \textit{(definitional)}:
\begin{quote}\ttfamily\footnotesize\raggedright\noindent
forall(\allowbreak{}[pdf\_], prod-converges(\allowbreak{}pdf\_) iff forsome(\allowbreak{}[pdl\_], prod-converges-to(\allowbreak{}pdf\_, pdl\_)))
\end{quote}
mentioned by 7 result(s)

\smallskip

\noindent\textbf{\texttt{prod-converges-to(pcf\_, pcl\_)}}\quad\textit{predicate}\index{prod-converges-to(pcf_, pcl_)@\texttt{prod-converges-to(pcf\_, pcl\_)}}\\
reads: the infinite product of \$1 converges to \$2 \\
defined by \texttt{prod-converges-to} \\
characterized by \texttt{prod-converges-to}  \textit{(definitional)}:
\begin{quote}\ttfamily\footnotesize\raggedright\noindent
forall(\allowbreak{}[pcf\_, pcl\_], prod-converges-to(\allowbreak{}pcf\_, pcl\_) iff pcf\_ in fun(\allowbreak{}nn, cc) and forsome(\allowbreak{}[pcm\_ in nn], forall(\allowbreak{}[pck\_ in nn], pcm\_ <= pck\_ implies not(\allowbreak{}pcf\_(\allowbreak{}pck\_) = 0)) and forsome(\allowbreak{}[pcz\_ in cc], not(\allowbreak{}pcz\_ = 0) and converges-to(\allowbreak{}cc-ms, vnb-lambda(\allowbreak{}pcn\_, nn, pprod(\allowbreak{}vnb-lambda(\allowbreak{}pcj\_, nn, pcf\_(\allowbreak{}pcm\_ + pcj\_)), pcn\_)), pcz\_) and pcl\_ = pprod(\allowbreak{}pcf\_, pcm\_) * pcz\_)))
\end{quote}
mentioned by 14 result(s)

\smallskip

\noindent\textbf{\texttt{prod-fun(pfu\_, pff\_)}}\quad\textit{functoid}\index{prod-fun(pfu_, pff_)@\texttt{prod-fun(pfu\_, pff\_)}}\\
reads: the function x |-> prod\_k (1 + \$2(k)(x)) on \$1 \\
definition:
\begin{quote}\ttfamily\footnotesize\raggedright\noindent
vnb-lambda(\allowbreak{}pfx\_, pfu\_, prod-value(\allowbreak{}vnb-lambda(\allowbreak{}pfk\_, nn, 1 + (\allowbreak{}pff\_(\allowbreak{}pfk\_))(\allowbreak{}pfx\_))))
\end{quote}
characterized by \texttt{prod-fun-holomorphic}  \textit{(proven)}:
\begin{quote}\ttfamily\footnotesize\raggedright\noindent
forall(\allowbreak{}[c6u\_], is-open(\allowbreak{}cc-ms, c6u\_) implies forall(\allowbreak{}[c6f\_ in fun(\allowbreak{}nn, fun(\allowbreak{}c6u\_, cc))], forall(\allowbreak{}[c6i\_ in nn], holomorphic-on(\allowbreak{}c6u\_, c6f\_(\allowbreak{}c6i\_))) implies prod-normal-on-compacta(\allowbreak{}c6u\_, c6f\_) implies holomorphic-on(\allowbreak{}c6u\_, prod-fun(\allowbreak{}c6u\_, c6f\_))))
\end{quote}
also: \texttt{prod-fun-partial-limit} \\
mentioned by 15 result(s)

\smallskip

\noindent\textbf{\texttt{prod-normal-on-compacta(pnu\_, pnf\_)}}\quad\textit{predicate}\index{prod-normal-on-compacta(pnu_, pnf_)@\texttt{prod-normal-on-compacta(pnu\_, pnf\_)}}\\
reads: the product of the functions 1 + \$2(k) is normally convergent on the compact subsets of \$1 \\
defined by \texttt{prod-normal-on-compacta} \\
characterized by \texttt{prod-normal-on-compacta}  \textit{(definitional)}:
\begin{quote}\ttfamily\footnotesize\raggedright\noindent
forall(\allowbreak{}[pnu\_, pnf\_], prod-normal-on-compacta(\allowbreak{}pnu\_, pnf\_) iff forall(\allowbreak{}[pnn\_ in nn], is-normally-convergent(\allowbreak{}pnf\_, cc-exhaust(\allowbreak{}pnu\_, pnn\_))))
\end{quote}
mentioned by 17 result(s)

\smallskip

\noindent\textbf{\texttt{prod-ord}}\quad\textit{defined-fn}\index{prod-ord@\texttt{prod-ord}}\\
defined by \texttt{prod-ord-zero} \texttt{prod-ord-succ} \\
characterized by \texttt{prod-ord-succ}  \textit{(definitional)}:
\begin{quote}\ttfamily\footnotesize\raggedright\noindent
forall(\allowbreak{}[m, f, n in nn], prod-ord(\allowbreak{}m, f, succ(\allowbreak{}n)) == (\allowbreak{}opr(\allowbreak{}m))(\allowbreak{}prod-ord(\allowbreak{}m, f, n), f(\allowbreak{}n)))
\end{quote}
also: \texttt{prod-ord-zero} \texttt{prod-ord-type} \texttt{prod-ord-type-ind} \texttt{prod-ord-singleton} \\
mentioned by 8 result(s)

\smallskip

\noindent\textbf{\texttt{prod-ring(r, f, s)}}\quad\textit{functoid}\index{prod-ring(r, f, s)@\texttt{prod-ring(r, f, s)}}\\
definition:
\begin{quote}\ttfamily\footnotesize\raggedright\noindent
finprod(\allowbreak{}commutative-ring-multiplicative-cm(\allowbreak{}r), f, s)
\end{quote}
characterized by \texttt{prod-ring-type}  \textit{(proven)}:
\begin{quote}\ttfamily\footnotesize\raggedright\noindent
forall(\allowbreak{}[r], is-commutative-ring(\allowbreak{}r) implies forall(\allowbreak{}[x], x in set and card(\allowbreak{}x) in nn implies forall(\allowbreak{}[f in fun(\allowbreak{}x, carr(\allowbreak{}r))], prod-ring(\allowbreak{}r, f, x) in carr(\allowbreak{}r))))
\end{quote}
also: \texttt{prod-ring-const} \texttt{prod-ring-empty} \texttt{prod-ring-insert} \texttt{prod-ring-cc-ones} \texttt{prod-ring-singleton} \texttt{prod-ring-type-ptwise} \texttt{prod-ring-type-subset} \texttt{prod-ring-zero-factor} ... (12 in all) \\
mentioned by 77 result(s)

\smallskip

\noindent\textbf{\texttt{prod-rmod(m)}}\quad\textit{functoid}\index{prod-rmod(m)@\texttt{prod-rmod(m)}}\\
reads: the module \$1 x R \\
definition:
\begin{quote}\ttfamily\footnotesize\raggedright\noindent
[scal(\allowbreak{}m), cartesian(\allowbreak{}vec(\allowbreak{}m), rr), pv-vadd(\allowbreak{}m), [vzero(\allowbreak{}m), 0], pv-vneg(\allowbreak{}m), pv-act(\allowbreak{}m)]
\end{quote}
characterized by \texttt{prod-rmod-act-cnvs-as-module}  \textit{(definitional)}:
\begin{quote}\ttfamily\footnotesize\raggedright\noindent
forall(\allowbreak{}[r], is-complex-normed-vector-space(\allowbreak{}r) implies scal(\allowbreak{}r) = normed-field-as-commutative-ring(\allowbreak{}rr-normed-field) implies act(\allowbreak{}prod-rmod(\allowbreak{}cnvs-as-module(\allowbreak{}r))) == pv-act(\allowbreak{}cnvs-as-module(\allowbreak{}r)))
\end{quote}
also: \texttt{prod-rmod-vec-cnvs-as-module} \texttt{prod-rmod-comm-cnvs-as-module} \texttt{prod-rmod-scal-cnvs-as-module} \texttt{prod-rmod-vadd-cnvs-as-module} \texttt{prod-rmod-vneg-cnvs-as-module} \texttt{prod-rmod-assoc-cnvs-as-module} \texttt{prod-rmod-vzero-cnvs-as-module} \texttt{prod-rmod-identity-cnvs-as-module} ... (29 in all) \\
mentioned by 72 result(s)

\smallskip

\noindent\textbf{\texttt{prod-set(cm, s, f)}}\quad\textit{functoid}\index{prod-set(cm, s, f)@\texttt{prod-set(cm, s, f)}}\\
definition:
\begin{quote}\ttfamily\footnotesize\raggedright\noindent
finsum(\allowbreak{}cm, f, s)
\end{quote}
characterized by \texttt{prod-set-empty}  \textit{(proven)}:
\begin{quote}\ttfamily\footnotesize\raggedright\noindent
forall(\allowbreak{}[cm, f], prod-set(\allowbreak{}cm, empty-set, f) == iden(\allowbreak{}cm))
\end{quote}
also: \texttt{prod-set-type-defined} \texttt{prod-set-singleton-defined} \texttt{prod-set-disjoint-union-defined} \\
mentioned by 7 result(s)

\smallskip

\noindent\textbf{\texttt{prod-top(t)}}\quad\textit{functoid}\index{prod-top(t)@\texttt{prod-top(t)}}\\
reads: the product topology of \$1 and R \\
definition:
\begin{quote}\ttfamily\footnotesize\raggedright\noindent
[cartesian(\allowbreak{}pts(\allowbreak{}t), rr), \{ptu\_ in power(\allowbreak{}cartesian(\allowbreak{}pts(\allowbreak{}t), rr)): is-pt-open(\allowbreak{}t, ptu\_)\}]
\end{quote}
characterized by \texttt{prod-top-pts}  \textit{(proven)}:
\begin{quote}\ttfamily\footnotesize\raggedright\noindent
forall(\allowbreak{}[t], is-top-space(\allowbreak{}t) implies pts(\allowbreak{}prod-top(\allowbreak{}t)) == cartesian(\allowbreak{}pts(\allowbreak{}t), rr))
\end{quote}
also: \texttt{prod-top-opens} \texttt{prod-top-hausdorff} \texttt{prod-top-is-top-space} \\
mentioned by 11 result(s)

\smallskip

\noindent\textbf{\texttt{prod-value(pvh\_)}}\quad\textit{functoid}\index{prod-value(pvh_)@\texttt{prod-value(pvh\_)}}\\
reads: the value of the infinite product of \$1 \\
definition:
\begin{quote}\ttfamily\footnotesize\raggedright\noindent
iota(\allowbreak{}pvl\_, pvl\_ in cc and prod-converges-to(\allowbreak{}pvh\_, pvl\_))
\end{quote}
characterized by \texttt{prod-value-spec}  \textit{(proven)}:
\begin{quote}\ttfamily\footnotesize\raggedright\noindent
forall(\allowbreak{}[c6f\_ in fun(\allowbreak{}nn, cc)], prod-converges(\allowbreak{}c6f\_) implies prod-value(\allowbreak{}c6f\_) in cc and prod-converges-to(\allowbreak{}c6f\_, prod-value(\allowbreak{}c6f\_)))
\end{quote}
also: \texttt{prod-value-split} \\
mentioned by 5 result(s)

\smallskip

\noindent\textbf{\texttt{product-carrier(ms)}}\quad\textit{functoid}\index{product-carrier(ms)@\texttt{product-carrier(ms)}}\\
definition:
\begin{quote}\ttfamily\footnotesize\raggedright\noindent
\{x in fun(\allowbreak{}nn, big-union(\allowbreak{}n, nn, pts(\allowbreak{}ms(\allowbreak{}n)))): forall(\allowbreak{}[n in nn], x(\allowbreak{}n) in pts(\allowbreak{}ms(\allowbreak{}n)))\}
\end{quote}
characterized by \texttt{product-carrier-mem}  \textit{(proven)}:
\begin{quote}\ttfamily\footnotesize\raggedright\noindent
forall(\allowbreak{}[ms, pcx], pcx in product-carrier(\allowbreak{}ms) iff pcx in fun(\allowbreak{}nn, big-union(\allowbreak{}n, nn, pts(\allowbreak{}ms(\allowbreak{}n)))) and forall(\allowbreak{}[n in nn], pcx(\allowbreak{}n) in pts(\allowbreak{}ms(\allowbreak{}n))))
\end{quote}
also: \texttt{product-carrier-coord} \texttt{product-carrier-unfold} \\
mentioned by 41 result(s)

\smallskip

\noindent\textbf{\texttt{product-measure(xs, ca, mu, ys, cb, nu)}}\quad\textit{functoid}\index{product-measure(xs, ca, mu, ys, cb, nu)@\texttt{product-measure(xs, ca, mu, ys, cb, nu)}}\\
reads: the product of the measures \$3 and \$6 \\
definition:
\begin{quote}\ttfamily\footnotesize\raggedright\noindent
vnb-lambda(\allowbreak{}pmc\_, product-sigma(\allowbreak{}xs, ca, ys, cb), integral(\allowbreak{}xs, ca, mu, x-section-measure(\allowbreak{}xs, ys, nu, pmc\_)))
\end{quote}
characterized by \texttt{product-measure-apply}  \textit{(proven)}:
\begin{quote}\ttfamily\footnotesize\raggedright\noindent
forall(\allowbreak{}[xs, ca, mu, ys, cb, nu, cs in product-sigma(\allowbreak{}xs, ca, ys, cb)], (\allowbreak{}product-measure(\allowbreak{}xs, ca, mu, ys, cb, nu))(\allowbreak{}cs) == integral(\allowbreak{}xs, ca, mu, x-section-measure(\allowbreak{}xs, ys, nu, cs)))
\end{quote}
also: \texttt{product-measure-rectangle} \texttt{product-measure-symmetric} \texttt{product-measure-is-measure} \\
mentioned by 36 result(s)

\smallskip

\noindent\textbf{\texttt{product-metric(ms)}}\quad\textit{functoid}\index{product-metric(ms)@\texttt{product-metric(ms)}}\\
definition:
\begin{quote}\ttfamily\footnotesize\raggedright\noindent
product-metric-w(\allowbreak{}ms, vnb-lambda(\allowbreak{}n, nn, (\allowbreak{}/)(\allowbreak{}1, 2 \textasciicircum{} (\allowbreak{}n + 1))))
\end{quote}
characterized by \texttt{product-metric-unfold}  \textit{(proven)}:
\begin{quote}\ttfamily\footnotesize\raggedright\noindent
forall(\allowbreak{}[ms], product-metric(\allowbreak{}ms) == product-metric-w(\allowbreak{}ms, vnb-lambda(\allowbreak{}n, nn, (\allowbreak{}/)(\allowbreak{}1, 2 \textasciicircum{} (\allowbreak{}n + 1)))))
\end{quote}
mentioned by 8 result(s)

\smallskip

\noindent\textbf{\texttt{product-metric-w(ms, w)}}\quad\textit{functoid}\index{product-metric-w(ms, w)@\texttt{product-metric-w(ms, w)}}\\
definition:
\begin{quote}\ttfamily\footnotesize\raggedright\noindent
[product-carrier(\allowbreak{}ms), vnb-lambda(\allowbreak{}[x, y], cartesian(\allowbreak{}product-carrier(\allowbreak{}ms), product-carrier(\allowbreak{}ms)), iota(\allowbreak{}l, l in rr and series-converges-to(\allowbreak{}vnb-lambda(\allowbreak{}n, nn, w(\allowbreak{}n) * (\allowbreak{}dist(\allowbreak{}bdd-metric(\allowbreak{}ms(\allowbreak{}n))))(\allowbreak{}x(\allowbreak{}n), y(\allowbreak{}n))), l)))]
\end{quote}
mentioned by 24 result(s)

\smallskip

\noindent\textbf{\texttt{product-proj(ms, n)}}\quad\textit{functoid}\index{product-proj(ms, n)@\texttt{product-proj(ms, n)}}\\
definition:
\begin{quote}\ttfamily\footnotesize\raggedright\noindent
vnb-lambda(\allowbreak{}x, product-carrier(\allowbreak{}ms), x(\allowbreak{}n))
\end{quote}
characterized by \texttt{product-projection-continuous}  \textit{(proven)}:
\begin{quote}\ttfamily\footnotesize\raggedright\noindent
forall(\allowbreak{}[ms], is-ms-sequence(\allowbreak{}ms) implies forall(\allowbreak{}[w], summable-weight(\allowbreak{}w) implies forall(\allowbreak{}[n in nn], is-continuous(\allowbreak{}product-metric-w(\allowbreak{}ms, w), ms(\allowbreak{}n), product-proj(\allowbreak{}ms, n)))))
\end{quote}
mentioned by 4 result(s)

\smallskip

\noindent\textbf{\texttt{product-sigma(xs, ca, ys, cb)}}\quad\textit{functoid}\index{product-sigma(xs, ca, ys, cb)@\texttt{product-sigma(xs, ca, ys, cb)}}\\
reads: the product sigma-algebra of \$2 and \$4 on \$1 x \$3 \\
definition:
\begin{quote}\ttfamily\footnotesize\raggedright\noindent
sigma-generated(\allowbreak{}cartesian(\allowbreak{}xs, ys), rectangles(\allowbreak{}xs, ca, ys, cb))
\end{quote}
characterized by \texttt{product-sigma-unfold}  \textit{(proven)}:
\begin{quote}\ttfamily\footnotesize\raggedright\noindent
forall(\allowbreak{}[xs, ca, ys, cb], product-sigma(\allowbreak{}xs, ca, ys, cb) == sigma-generated(\allowbreak{}cartesian(\allowbreak{}xs, ys), rectangles(\allowbreak{}xs, ca, ys, cb)))
\end{quote}
also: \texttt{rectangle-in-product-sigma} \texttt{product-sigma-is-sigma-algebra} \\
mentioned by 84 result(s)

\smallskip

\noindent\textbf{\texttt{product2-ms(s, t)}}\quad\textit{functoid}\index{product2-ms(s, t)@\texttt{product2-ms(s, t)}}\\
reads: the product metric space of \$1 and \$2 \\
definition:
\begin{quote}\ttfamily\footnotesize\raggedright\noindent
[cartesian(\allowbreak{}pts(\allowbreak{}s), pts(\allowbreak{}t)), vnb-lambda(\allowbreak{}[p2u\_, p2v\_], cartesian(\allowbreak{}cartesian(\allowbreak{}pts(\allowbreak{}s), pts(\allowbreak{}t)), cartesian(\allowbreak{}pts(\allowbreak{}s), pts(\allowbreak{}t))), max(\allowbreak{}(\allowbreak{}dist(\allowbreak{}s))(\allowbreak{}nth(\allowbreak{}1, p2u\_), nth(\allowbreak{}1, p2v\_)), (\allowbreak{}dist(\allowbreak{}t))(\allowbreak{}nth(\allowbreak{}2, p2u\_), nth(\allowbreak{}2, p2v\_))))]
\end{quote}
mentioned by 53 result(s)

\smallskip

\noindent\textbf{\texttt{proj(s)}}\quad\textit{functoid}\index{proj(s)@\texttt{proj(s)}}\\
definition:
\begin{quote}\ttfamily\footnotesize\raggedright\noindent
vnb-lambda(\allowbreak{}a, pts(\allowbreak{}s), class(\allowbreak{}s, a))
\end{quote}
characterized by \texttt{proj-in-fun}  \textit{(proven)}:
\begin{quote}\ttfamily\footnotesize\raggedright\noindent
forall(\allowbreak{}[s], is-setoid(\allowbreak{}s) implies proj(\allowbreak{}s) in fun(\allowbreak{}pts(\allowbreak{}s), quotient(\allowbreak{}s)))
\end{quote}
also: \texttt{proj-unfold} \\
mentioned by 5 result(s)

\smallskip

\noindent\textbf{\texttt{proj-op(bov\_, bok\_)}}\quad\textit{functoid}\index{proj-op(bov_, bok_)@\texttt{proj-op(bov\_, bok\_)}}\\
reads: the orthogonal projection onto \$2 \\
definition:
\begin{quote}\ttfamily\footnotesize\raggedright\noindent
iota(\allowbreak{}bop\_, is-orthogonal-projection(\allowbreak{}bov\_, bok\_, bop\_))
\end{quote}
characterized by \texttt{proj-op-range}  \textit{(proven)}:
\begin{quote}\ttfamily\footnotesize\raggedright\noindent
forall(\allowbreak{}[v, kk], is-hilbert-space(\allowbreak{}v) implies is-submodule(\allowbreak{}v, kk) implies is-closed(\allowbreak{}ip-metric-space(\allowbreak{}v), kk) implies op-range(\allowbreak{}v, v, proj-op(\allowbreak{}v, kk)) = kk)
\end{quote}
also: \texttt{proj-op-linear} \texttt{proj-op-unique} \texttt{proj-op-is-orthogonal-projection} \texttt{projection-op-is-proj-op-of-range} \\
mentioned by 17 result(s)

\smallskip

\noindent\textbf{\texttt{prop}}\quad\textit{tactic}\index{prop@\texttt{prop}}\\
\texttt{(prop)} -- Decide the goal by PROPOSITIONAL logic from the context, and close it if it follows.  Every non-connective formula -- \texttt{x in a', an equation, a whole }forall(...)' -- is one opaque atom; AND/OR/NOT/IMPLIES/IFF are read.  On a goal that does NOT follow it prints the countermodel (which atom must be true, which false) and leaves the proof untouched.  Adds no trust: it decides semantically, then discharges through di/ai/oi/ass/use-em/have!/detach!, so the bill is unchanged and the recorded script is the ordinary step-by-step proof. \\
declared in \texttt{tactics-help.scm}

\smallskip

\noindent\textbf{\texttt{pruefer(prp\_)}}\quad\textit{functoid}\index{pruefer(prp_)@\texttt{pruefer(prp\_)}}\\
reads: the Pruefer \$1-group Q\_\$1/Z\_\$1 \\
definition:
\begin{quote}\ttfamily\footnotesize\raggedright\noindent
quotient-ag(\allowbreak{}normed-field-additive-ag(\allowbreak{}qp-field(\allowbreak{}prp\_)), zp(\allowbreak{}prp\_))
\end{quote}
characterized by \texttt{qp-pruefer-order}  \textit{(proven)}:
\begin{quote}\ttfamily\footnotesize\raggedright\noindent
forall(\allowbreak{}[p], zz-prime(\allowbreak{}p) implies forall(\allowbreak{}[x in carr(\allowbreak{}qp-field(\allowbreak{}p))], not(\allowbreak{}x = zero(\allowbreak{}qp-field(\allowbreak{}p))) implies forall(\allowbreak{}[n in nn], -(\allowbreak{}qp-ord(\allowbreak{}p))(\allowbreak{}x) = n implies is-element-order(\allowbreak{}pruefer(\allowbreak{}p), class(\allowbreak{}qmod-setoid(\allowbreak{}ag-zz-module(\allowbreak{}normed-field-additive-ag(\allowbreak{}qp-field(\allowbreak{}p))), zp(\allowbreak{}p)), x), p \textasciicircum{} n))))
\end{quote}
mentioned by 6 result(s)

\smallskip

\noindent\textbf{\texttt{ps-absolutely-converges-at(coef, x)}}\quad\textit{predicate}\index{ps-absolutely-converges-at(coef, x)@\texttt{ps-absolutely-converges-at(coef, x)}}\\
reads: the power series with coefficients \$1 converges absolutely at \$2 \\
defined by \texttt{ps-absolutely-converges-at} \\
characterized by \texttt{ps-absolutely-converges-at}  \textit{(definitional)}:
\begin{quote}\ttfamily\footnotesize\raggedright\noindent
forall(\allowbreak{}[coef, x], ps-absolutely-converges-at(\allowbreak{}coef, x) iff ps-converges-at(\allowbreak{}vnb-lambda(\allowbreak{}n, nn, abs(\allowbreak{}coef(\allowbreak{}n))), abs(\allowbreak{}x)))
\end{quote}
mentioned by 3 result(s)

\smallskip

\noindent\textbf{\texttt{ps-converges-at(coef, x)}}\quad\textit{predicate}\index{ps-converges-at(coef, x)@\texttt{ps-converges-at(coef, x)}}\\
reads: the power series with coefficients \$1 converges at \$2 \\
defined by \texttt{ps-converges-at} \\
characterized by \texttt{ps-converges-at}  \textit{(definitional)}:
\begin{quote}\ttfamily\footnotesize\raggedright\noindent
forall(\allowbreak{}[coef, x], ps-converges-at(\allowbreak{}coef, x) iff converges(\allowbreak{}rr-ms, vnb-lambda(\allowbreak{}k, nn, ps-partial-sum(\allowbreak{}coef, x, k))))
\end{quote}
mentioned by 8 result(s)

\smallskip

\noindent\textbf{\texttt{ps-converges-to-at(coef, x, l)}}\quad\textit{predicate}\index{ps-converges-to-at(coef, x, l)@\texttt{ps-converges-to-at(coef, x, l)}}\\
reads: the power series with coefficients \$1 converges to \$3 at \$2 \\
defined by \texttt{ps-converges-to-at} \\
characterized by \texttt{ps-converges-to-at}  \textit{(definitional)}:
\begin{quote}\ttfamily\footnotesize\raggedright\noindent
forall(\allowbreak{}[coef, x, l], ps-converges-to-at(\allowbreak{}coef, x, l) iff converges-to(\allowbreak{}rr-ms, vnb-lambda(\allowbreak{}k, nn, ps-partial-sum(\allowbreak{}coef, x, k)), l))
\end{quote}
mentioned by 2 result(s)

\smallskip

\noindent\textbf{\texttt{ps-partial-sum(coef, x, k)}}\quad\textit{functoid}\index{ps-partial-sum(coef, x, k)@\texttt{ps-partial-sum(coef, x, k)}}\\
definition:
\begin{quote}\ttfamily\footnotesize\raggedright\noindent
sum-ag(\allowbreak{}normed-field-additive-ag(\allowbreak{}rr-normed-field), vnb-lambda(\allowbreak{}n, nn, coef(\allowbreak{}n) * x \textasciicircum{} n), k)
\end{quote}
characterized by \texttt{ps-partial-sum-as-series}  \textit{(proven)}:
\begin{quote}\ttfamily\footnotesize\raggedright\noindent
forall(\allowbreak{}[coef in fun(\allowbreak{}nn, rr), x in rr, k in nn], ps-partial-sum(\allowbreak{}coef, x, k) = series-partial-sum(\allowbreak{}vnb-lambda(\allowbreak{}n, nn, coef(\allowbreak{}n) * x \textasciicircum{} n), k))
\end{quote}
mentioned by 6 result(s)

\smallskip

\noindent\textbf{\texttt{ps-ratio-limit(coef, l)}}\quad\textit{predicate}\index{ps-ratio-limit(coef, l)@\texttt{ps-ratio-limit(coef, l)}}\\
reads: the coefficients \$1 have ratio limit \$2 \\
defined by \texttt{ps-ratio-limit} \\
characterized by \texttt{ps-ratio-limit}  \textit{(definitional)}:
\begin{quote}\ttfamily\footnotesize\raggedright\noindent
forall(\allowbreak{}[coef, l], ps-ratio-limit(\allowbreak{}coef, l) iff converges-to(\allowbreak{}rr-ms, vnb-lambda(\allowbreak{}n, nn, abs(\allowbreak{}coef(\allowbreak{}succ(\allowbreak{}n)) * recip(\allowbreak{}coef(\allowbreak{}n)))), l))
\end{quote}
mentioned by 3 result(s)

\smallskip

\noindent\textbf{\texttt{pseudo-gauge-top(fam, ground)}}\quad\textit{functoid}\index{pseudo-gauge-top(fam, ground)@\texttt{pseudo-gauge-top(fam, ground)}}\\
reads: the topology generated by the pseudometric family \$1 on \$2 \\
definition:
\begin{quote}\ttfamily\footnotesize\raggedright\noindent
[ground, \{u in power(\allowbreak{}ground): is-open-gauge(\allowbreak{}fam, ground, u)\}]
\end{quote}
mentioned by 11 result(s)

\smallskip

\noindent\textbf{\texttt{pseudometric-space}}\quad\textit{structure}\index{pseudometric-space@\texttt{pseudometric-space}}\\
predicate \texttt{is-pseudometric-space}; slots: pts dist

\smallskip

\noindent\textbf{\texttt{pt-box(t, v, c, e)}}\quad\textit{functoid}\index{pt-box(t, v, c, e)@\texttt{pt-box(t, v, c, e)}}\\
reads: the box \$2 x (\$3 - \$4, \$3 + \$4) \\
definition:
\begin{quote}\ttfamily\footnotesize\raggedright\noindent
\{ptz\_ in cartesian(\allowbreak{}pts(\allowbreak{}t), rr): nth(\allowbreak{}1, ptz\_) in v and abs(\allowbreak{}nth(\allowbreak{}2, ptz\_) - c) < e\}
\end{quote}
characterized by \texttt{pt-box-mono}  \textit{(proven)}:
\begin{quote}\ttfamily\footnotesize\raggedright\noindent
forall(\allowbreak{}[t, v1, v2, c, e1, e2], is-top-space(\allowbreak{}t) implies v1 subset v2 implies c in rr implies e1 in rr implies e2 in rr implies e1 <= e2 implies pt-box(\allowbreak{}t, v1, c, e1) subset pt-box(\allowbreak{}t, v2, c, e2))
\end{quote}
also: \texttt{pt-box-open} \texttt{pt-box-unfold} \texttt{pt-box-membership} \\
mentioned by 8 result(s)

\smallskip

\noindent\textbf{\texttt{pt-set-dist(m, x, es)}}\quad\textit{functoid}\index{pt-set-dist(m, x, es)@\texttt{pt-set-dist(m, x, es)}}\\
reads: the distance in \$1 from the point \$2 to the set \$3 \\
definition:
\begin{quote}\ttfamily\footnotesize\raggedright\noindent
einf(\allowbreak{}\{mov\_ in rr-pos-star: forsome(\allowbreak{}[moz\_ in es], mov\_ = (\allowbreak{}dist(\allowbreak{}m))(\allowbreak{}x, moz\_))\})
\end{quote}
characterized by \texttt{pt-set-dist-le}  \textit{(proven)}:
\begin{quote}\ttfamily\footnotesize\raggedright\noindent
forall(\allowbreak{}[m], is-metric-space(\allowbreak{}m) implies forall(\allowbreak{}[es], es subset pts(\allowbreak{}m) implies forall(\allowbreak{}[x in pts(\allowbreak{}m), z in es], pt-set-dist(\allowbreak{}m, x, es) <= (\allowbreak{}dist(\allowbreak{}m))(\allowbreak{}x, z))))
\end{quote}
also: \texttt{pt-set-dist-real} \texttt{pt-set-dist-unfold} \texttt{pt-set-dist-ge-elim} \texttt{pt-set-dist-ge-intro} \texttt{pt-set-dist-open-pos} \texttt{pt-set-dist-lipschitz} \\
mentioned by 20 result(s)

\smallskip

\noindent\textbf{\texttt{pts(s)}}\quad\textit{accessor}\index{pts(s)@\texttt{pts(s)}}\\
reads: the points of \$1 \\
a structure slot \\
characterized by \texttt{cc-ms@pts}  \textit{(definitional)}:
\begin{quote}\ttfamily\footnotesize\raggedright\noindent
pts(\allowbreak{}cc-ms) == cc
\end{quote}
also: \texttt{rr-ms@pts} \texttt{cc-hat-top@pts} \texttt{ringoid-setoid-pts} \texttt{sip-setoid-pts-csip-as-sip} \texttt{qmod-setoid-pts-cna-as-module} \texttt{qmod-setoid-pts-sip-as-module} \texttt{qmod-setoid-pts-uba-as-module} \texttt{qmod-setoid-pts-cnvs-as-module} ... (73 in all) \\
mentioned by 1669 result(s)

\smallskip

\noindent\textbf{\texttt{pts-left-cosets(lcs\_, lcg\_, lch\_)}}\quad\textit{functoid}\index{pts-left-cosets(lcs_, lcg_, lch_)@\texttt{pts-left-cosets(lcs\_, lcg\_, lch\_)}}\\
reads: the left cosets of \$3 in the group \$2 on \$1 \\
definition:
\begin{quote}\ttfamily\footnotesize\raggedright\noindent
\{lcu\_ in power(\allowbreak{}pts(\allowbreak{}lcs\_)): forsome(\allowbreak{}[lcx\_ in carr(\allowbreak{}lcg\_)], lcu\_ = \{lcz\_ in pts(\allowbreak{}lcs\_): (\allowbreak{}opr(\allowbreak{}lcg\_))(\allowbreak{}(\allowbreak{}inv(\allowbreak{}lcg\_))(\allowbreak{}lcx\_), lcz\_) in lch\_\})\}
\end{quote}
characterized by \texttt{pts-left-cosets-unfold}  \textit{(proven)}:
\begin{quote}\ttfamily\footnotesize\raggedright\noindent
forall(\allowbreak{}[s, g, h], pts-left-cosets(\allowbreak{}s, g, h) == \{lcu\_ in power(\allowbreak{}pts(\allowbreak{}s)): forsome(\allowbreak{}[lcx\_ in carr(\allowbreak{}g)], lcu\_ = \{lcz\_ in pts(\allowbreak{}s): (\allowbreak{}opr(\allowbreak{}g))(\allowbreak{}(\allowbreak{}inv(\allowbreak{}g))(\allowbreak{}lcx\_), lcz\_) in h\})\})
\end{quote}
mentioned by 4 result(s)

\smallskip

\noindent\textbf{\texttt{ptwise-eplus(omega, f, g)}}\quad\textit{functoid}\index{ptwise-eplus(omega, f, g)@\texttt{ptwise-eplus(omega, f, g)}}\\
reads: the pointwise sum of \$2 and \$3 on \$1 \\
definition:
\begin{quote}\ttfamily\footnotesize\raggedright\noindent
vnb-lambda(\allowbreak{}x\_, omega, eplus(\allowbreak{}f(\allowbreak{}x\_), g(\allowbreak{}x\_)))
\end{quote}
characterized by \texttt{ptwise-eplus-value}  \textit{(proven)}:
\begin{quote}\ttfamily\footnotesize\raggedright\noindent
forall(\allowbreak{}[omega, f, g, x in omega], f(\allowbreak{}x) in rr-pos-star implies g(\allowbreak{}x) in rr-pos-star implies (\allowbreak{}ptwise-eplus(\allowbreak{}omega, f, g))(\allowbreak{}x) = eplus(\allowbreak{}f(\allowbreak{}x), g(\allowbreak{}x)))
\end{quote}
also: \texttt{ptwise-eplus-value-q} \\
mentioned by 19 result(s)

\smallskip

\noindent\textbf{\texttt{ptwise-epow(epo\_, epf\_, epp\_)}}\quad\textit{functoid}\index{ptwise-epow(epo_, epf_, epp_)@\texttt{ptwise-epow(epo\_, epf\_, epp\_)}}\\
reads: the pointwise power \$3 of \$2 on \$1 \\
definition:
\begin{quote}\ttfamily\footnotesize\raggedright\noindent
vnb-lambda(\allowbreak{}epx\_, epo\_, epow(\allowbreak{}epf\_(\allowbreak{}epx\_), epp\_))
\end{quote}
characterized by \texttt{ptwise-epow-value}  \textit{(proven)}:
\begin{quote}\ttfamily\footnotesize\raggedright\noindent
forall(\allowbreak{}[omega, f in fun(\allowbreak{}omega, rr-pos-star), p\_ in rr], 0 < p\_ implies forall(\allowbreak{}[x in omega], (\allowbreak{}ptwise-epow(\allowbreak{}omega, f, p\_))(\allowbreak{}x) = epow(\allowbreak{}f(\allowbreak{}x), p\_)))
\end{quote}
mentioned by 10 result(s)

\smallskip

\noindent\textbf{\texttt{ptwise-etimes(omega, f, g)}}\quad\textit{functoid}\index{ptwise-etimes(omega, f, g)@\texttt{ptwise-etimes(omega, f, g)}}\\
reads: the pointwise product of \$2 and \$3 on \$1 \\
definition:
\begin{quote}\ttfamily\footnotesize\raggedright\noindent
vnb-lambda(\allowbreak{}x\_, omega, etimes(\allowbreak{}f(\allowbreak{}x\_), g(\allowbreak{}x\_)))
\end{quote}
characterized by \texttt{ptwise-etimes-sup}  \textit{(proven)}:
\begin{quote}\ttfamily\footnotesize\raggedright\noindent
forall(\allowbreak{}[omega, h, cf], omega in set implies h in fun(\allowbreak{}omega, rr-pos-star) implies cf in fun(\allowbreak{}nn, fun(\allowbreak{}omega, rr-pos-star)) implies forall(\allowbreak{}[n\_ in nn], ptwise-le(\allowbreak{}omega, cf(\allowbreak{}n\_), cf(\allowbreak{}succ(\allowbreak{}n\_)))) implies ptwise-etimes(\allowbreak{}omega, h, ptwise-sup(\allowbreak{}omega, cf)) = ptwise-sup(\allowbreak{}omega, vnb-lambda(\allowbreak{}mrgk\_, nn, vnb-lambda(\allowbreak{}x\_, omega, etimes(\allowbreak{}h(\allowbreak{}x\_), (\allowbreak{}cf(\allowbreak{}mrgk\_))(\allowbreak{}x\_))))))
\end{quote}
also: \texttt{ptwise-etimes-value} \texttt{ptwise-etimes-value-q} \\
mentioned by 33 result(s)

\smallskip

\noindent\textbf{\texttt{ptwise-le(omega, f, g)}}\quad\textit{predicate}\index{ptwise-le(omega, f, g)@\texttt{ptwise-le(omega, f, g)}}\\
reads: \$2 is pointwise at most \$3 on \$1 \\
defined by \texttt{ptwise-le} \\
characterized by \texttt{ptwise-le}  \textit{(definitional)}:
\begin{quote}\ttfamily\footnotesize\raggedright\noindent
forall(\allowbreak{}[omega, f, g], ptwise-le(\allowbreak{}omega, f, g) iff forall(\allowbreak{}[x\_ in omega], f(\allowbreak{}x\_) <= g(\allowbreak{}x\_)))
\end{quote}
mentioned by 19 result(s)

\smallskip

\noindent\textbf{\texttt{ptwise-liminf(omega, cf)}}\quad\textit{functoid}\index{ptwise-liminf(omega, cf)@\texttt{ptwise-liminf(omega, cf)}}\\
reads: the pointwise lower limit on \$1 of the sequence \$2 \\
definition:
\begin{quote}\ttfamily\footnotesize\raggedright\noindent
vnb-lambda(\allowbreak{}x\_, omega, eliminf(\allowbreak{}vnb-lambda(\allowbreak{}n\_, nn, (\allowbreak{}cf(\allowbreak{}n\_))(\allowbreak{}x\_))))
\end{quote}
characterized by \texttt{ptwise-liminf-apply}  \textit{(proven)}:
\begin{quote}\ttfamily\footnotesize\raggedright\noindent
forall(\allowbreak{}[omega, cf, x in omega], (\allowbreak{}ptwise-liminf(\allowbreak{}omega, cf))(\allowbreak{}x) == eliminf(\allowbreak{}vnb-lambda(\allowbreak{}n\_, nn, (\allowbreak{}cf(\allowbreak{}n\_))(\allowbreak{}x))))
\end{quote}
mentioned by 5 result(s)

\smallskip

\noindent\textbf{\texttt{ptwise-net-limits(t, ee)}}\quad\textit{functoid}\index{ptwise-net-limits(t, ee)@\texttt{ptwise-net-limits(t, ee)}}\\
reads: the pointwise limits of nets in \$2 \\
definition:
\begin{quote}\ttfamily\footnotesize\raggedright\noindent
\{pwf\_ in fun(\allowbreak{}pts(\allowbreak{}t), rr): forsome(\allowbreak{}[pwd\_, pwp\_], pwd\_ in directed-set and is-directed-set(\allowbreak{}pwd\_) and pwp\_ in fun(\allowbreak{}pts(\allowbreak{}pwd\_), ee) and forall(\allowbreak{}[pwx\_ in pts(\allowbreak{}t)], net-converges-to(\allowbreak{}metric-top(\allowbreak{}rr-ms), pwd\_, vnb-lambda(\allowbreak{}pwl\_, pts(\allowbreak{}pwd\_), (\allowbreak{}pwp\_(\allowbreak{}pwl\_))(\allowbreak{}pwx\_)), pwf\_(\allowbreak{}pwx\_))))\}
\end{quote}
characterized by \texttt{ptwise-net-limits-unfold}  \textit{(proven)}:
\begin{quote}\ttfamily\footnotesize\raggedright\noindent
forall(\allowbreak{}[t, ee], ptwise-net-limits(\allowbreak{}t, ee) == \{pwf\_ in fun(\allowbreak{}pts(\allowbreak{}t), rr): forsome(\allowbreak{}[pwd\_, pwp\_], pwd\_ in directed-set and is-directed-set(\allowbreak{}pwd\_) and pwp\_ in fun(\allowbreak{}pts(\allowbreak{}pwd\_), ee) and forall(\allowbreak{}[pwx\_ in pts(\allowbreak{}t)], net-converges-to(\allowbreak{}metric-top(\allowbreak{}rr-ms), pwd\_, vnb-lambda(\allowbreak{}pwl\_, pts(\allowbreak{}pwd\_), (\allowbreak{}pwp\_(\allowbreak{}pwl\_))(\allowbreak{}pwx\_)), pwf\_(\allowbreak{}pwx\_))))\})
\end{quote}
mentioned by 3 result(s)

\smallskip

\noindent\textbf{\texttt{ptwise-scale(omega, c, f)}}\quad\textit{functoid}\index{ptwise-scale(omega, c, f)@\texttt{ptwise-scale(omega, c, f)}}\\
reads: the function \$3 scaled by \$2 on \$1 \\
definition:
\begin{quote}\ttfamily\footnotesize\raggedright\noindent
vnb-lambda(\allowbreak{}x\_, omega, etimes(\allowbreak{}c, f(\allowbreak{}x\_)))
\end{quote}
characterized by \texttt{ptwise-scale-value}  \textit{(proven)}:
\begin{quote}\ttfamily\footnotesize\raggedright\noindent
forall(\allowbreak{}[omega, c\_, f, x in omega], c\_ in rr-pos-star implies f(\allowbreak{}x) in rr-pos-star implies (\allowbreak{}ptwise-scale(\allowbreak{}omega, c\_, f))(\allowbreak{}x) = etimes(\allowbreak{}c\_, f(\allowbreak{}x)))
\end{quote}
also: \texttt{ptwise-scale-value-q} \texttt{ptwise-scale-measurable} \\
mentioned by 23 result(s)

\smallskip

\noindent\textbf{\texttt{ptwise-sup(omega, cf)}}\quad\textit{functoid}\index{ptwise-sup(omega, cf)@\texttt{ptwise-sup(omega, cf)}}\\
reads: the pointwise supremum on \$1 of the sequence \$2 \\
definition:
\begin{quote}\ttfamily\footnotesize\raggedright\noindent
vnb-lambda(\allowbreak{}x\_, omega, esup(\allowbreak{}image(\allowbreak{}vnb-lambda(\allowbreak{}n\_, nn, (\allowbreak{}cf(\allowbreak{}n\_))(\allowbreak{}x\_)), nn)))
\end{quote}
characterized by \texttt{ptwise-sup-apply}  \textit{(proven)}:
\begin{quote}\ttfamily\footnotesize\raggedright\noindent
forall(\allowbreak{}[omega, cf, x in omega], (\allowbreak{}ptwise-sup(\allowbreak{}omega, cf))(\allowbreak{}x) == esup(\allowbreak{}image(\allowbreak{}vnb-lambda(\allowbreak{}n\_, nn, (\allowbreak{}cf(\allowbreak{}n\_))(\allowbreak{}x)), nn)))
\end{quote}
mentioned by 15 result(s)

\smallskip

\noindent\textbf{\texttt{pubkey(s)}}\quad\textit{accessor}\index{pubkey(s)@\texttt{pubkey(s)}}\\
reads: the pubkey slot of \$1 \\
a structure slot \\
mentioned by 30 result(s)

\smallskip

\noindent\textbf{\texttt{pullback-measure(omega, cb, f, nu)}}\quad\textit{functoid}\index{pullback-measure(omega, cb, f, nu)@\texttt{pullback-measure(omega, cb, f, nu)}}\\
reads: the pullback by \$3 of the measure \$4 on \$2 \\
definition:
\begin{quote}\ttfamily\footnotesize\raggedright\noindent
vnb-lambda(\allowbreak{}mpm\_, pullback-sigma(\allowbreak{}omega, cb, f), nu(\allowbreak{}image(\allowbreak{}f, mpm\_)))
\end{quote}
characterized by \texttt{pullback-measure-at}  \textit{(proven)}:
\begin{quote}\ttfamily\footnotesize\raggedright\noindent
forall(\allowbreak{}[omega, cb, omega2, f, nu], is-measure(\allowbreak{}omega2, cb, nu) implies f in fun(\allowbreak{}omega, omega2) implies forall(\allowbreak{}[mpsy\_ in omega2], forsome(\allowbreak{}[mpsx\_ in omega], f(\allowbreak{}mpsx\_) = mpsy\_)) implies forall(\allowbreak{}[a in pullback-sigma(\allowbreak{}omega, cb, f)], image(\allowbreak{}f, a) in cb and (\allowbreak{}pullback-measure(\allowbreak{}omega, cb, f, nu))(\allowbreak{}a) = nu(\allowbreak{}image(\allowbreak{}f, a))))
\end{quote}
also: \texttt{pullback-measure-type} \texttt{pullback-measure-value} \texttt{pullback-measure-unfold} \texttt{pullback-measure-unique} \texttt{pullback-measure-is-measure} \texttt{pullback-measure-exists-unique} \\
mentioned by 10 result(s)

\smallskip

\noindent\textbf{\texttt{pullback-ms(mt, ss, jj)}}\quad\textit{functoid}\index{pullback-ms(mt, ss, jj)@\texttt{pullback-ms(mt, ss, jj)}}\\
reads: the set \$2 with the distance of \$1 pulled back along \$3 \\
definition:
\begin{quote}\ttfamily\footnotesize\raggedright\noindent
[ss, vnb-lambda(\allowbreak{}[pbu\_, pbv\_], cartesian(\allowbreak{}ss, ss), (\allowbreak{}dist(\allowbreak{}mt))(\allowbreak{}jj(\allowbreak{}pbu\_), jj(\allowbreak{}pbv\_)))]
\end{quote}
mentioned by 10 result(s)

\smallskip

\noindent\textbf{\texttt{pullback-sigma(omega, cb, f)}}\quad\textit{functoid}\index{pullback-sigma(omega, cb, f)@\texttt{pullback-sigma(omega, cb, f)}}\\
reads: the pullback by \$3 of the sigma-algebra \$2 to \$1 \\
definition:
\begin{quote}\ttfamily\footnotesize\raggedright\noindent
\{mpa\_ in power(\allowbreak{}omega): forsome(\allowbreak{}[mpb\_ in cb], mpa\_ = \{mpx\_ in omega: f(\allowbreak{}mpx\_) in mpb\_\})\}
\end{quote}
characterized by \texttt{pullback-sigma-unfold}  \textit{(proven)}:
\begin{quote}\ttfamily\footnotesize\raggedright\noindent
forall(\allowbreak{}[omega, cb, f], pullback-sigma(\allowbreak{}omega, cb, f) == \{mpa\_ in power(\allowbreak{}omega): forsome(\allowbreak{}[mpb\_ in cb], mpa\_ = \{mpx\_ in omega: f(\allowbreak{}mpx\_) in mpb\_\})\})
\end{quote}
also: \texttt{pullback-sigma-algebra} \texttt{pullback-sigma-smallest} \texttt{pullback-sigma-measurable} \texttt{pullback-sigma-membership} \texttt{pullback-sigma-is-sigma-algebra} \\
mentioned by 16 result(s)

\smallskip

\noindent\textbf{\texttt{push-forward(omega, cb, mu, h)}}\quad\textit{functoid}\index{push-forward(omega, cb, mu, h)@\texttt{push-forward(omega, cb, mu, h)}}\\
reads: the push-forward along \$4 of the measure \$3, on the sets \$2 \\
definition:
\begin{quote}\ttfamily\footnotesize\raggedright\noindent
vnb-lambda(\allowbreak{}pwb\_, cb, mu(\allowbreak{}\{pwx\_ in omega: h(\allowbreak{}pwx\_) in pwb\_\}))
\end{quote}
characterized by \texttt{pfw-push-forward-value}  \textit{(proven)}:
\begin{quote}\ttfamily\footnotesize\raggedright\noindent
forall(\allowbreak{}[omega, cb, mu, h, b in cb], (\allowbreak{}push-forward(\allowbreak{}omega, cb, mu, h))(\allowbreak{}b) == mu(\allowbreak{}\{pwx\_ in omega: h(\allowbreak{}pwx\_) in b\}))
\end{quote}
also: \texttt{pfw-push-forward-unfold} \texttt{pfw-push-forward-measure} \texttt{measure-preserving-push-forward} \\
mentioned by 13 result(s)

\smallskip

\noindent\textbf{\texttt{pv-act(m)}}\quad\textit{functoid}\index{pv-act(m)@\texttt{pv-act(m)}}\\
reads: the componentwise scalar action of \$1 x R \\
definition:
\begin{quote}\ttfamily\footnotesize\raggedright\noindent
vnb-lambda(\allowbreak{}[pvr\_, pvc\_], cartesian(\allowbreak{}carr(\allowbreak{}scal(\allowbreak{}m)), cartesian(\allowbreak{}vec(\allowbreak{}m), rr)), [(\allowbreak{}act(\allowbreak{}m))(\allowbreak{}pvr\_, nth(\allowbreak{}1, pvc\_)), pvr\_ * nth(\allowbreak{}2, pvc\_)])
\end{quote}
characterized by \texttt{pv-act-fun-cnvs-as-module}  \textit{(definitional)}:
\begin{quote}\ttfamily\footnotesize\raggedright\noindent
forall(\allowbreak{}[r], is-complex-normed-vector-space(\allowbreak{}r) implies scal(\allowbreak{}r) = normed-field-as-commutative-ring(\allowbreak{}rr-normed-field) implies pv-act(\allowbreak{}cnvs-as-module(\allowbreak{}r)) in fun(\allowbreak{}cartesian(\allowbreak{}carr(\allowbreak{}scal(\allowbreak{}r)), cartesian(\allowbreak{}vec(\allowbreak{}r), rr)), cartesian(\allowbreak{}vec(\allowbreak{}r), rr)))
\end{quote}
also: \texttt{pv-act-value-cnvs-as-module} \texttt{pv-act-fun} \texttt{pv-act-value} \\
mentioned by 15 result(s)

\smallskip

\noindent\textbf{\texttt{pv-vadd(m)}}\quad\textit{functoid}\index{pv-vadd(m)@\texttt{pv-vadd(m)}}\\
reads: the componentwise addition of \$1 x R \\
definition:
\begin{quote}\ttfamily\footnotesize\raggedright\noindent
vnb-lambda(\allowbreak{}[pva\_, pvb\_], cartesian(\allowbreak{}cartesian(\allowbreak{}vec(\allowbreak{}m), rr), cartesian(\allowbreak{}vec(\allowbreak{}m), rr)), [(\allowbreak{}vadd(\allowbreak{}m))(\allowbreak{}nth(\allowbreak{}1, pva\_), nth(\allowbreak{}1, pvb\_)), nth(\allowbreak{}2, pva\_) + nth(\allowbreak{}2, pvb\_)])
\end{quote}
characterized by \texttt{pv-vadd-fun-cnvs-as-module}  \textit{(definitional)}:
\begin{quote}\ttfamily\footnotesize\raggedright\noindent
forall(\allowbreak{}[r], is-complex-normed-vector-space(\allowbreak{}r) implies scal(\allowbreak{}r) = normed-field-as-commutative-ring(\allowbreak{}rr-normed-field) implies pv-vadd(\allowbreak{}cnvs-as-module(\allowbreak{}r)) in fun(\allowbreak{}cartesian(\allowbreak{}cartesian(\allowbreak{}vec(\allowbreak{}r), rr), cartesian(\allowbreak{}vec(\allowbreak{}r), rr)), cartesian(\allowbreak{}vec(\allowbreak{}r), rr)))
\end{quote}
also: \texttt{pv-vadd-value-cnvs-as-module} \texttt{pv-vadd-fun} \texttt{pv-vadd-value} \\
mentioned by 15 result(s)

\smallskip

\noindent\textbf{\texttt{pv-vneg(m)}}\quad\textit{functoid}\index{pv-vneg(m)@\texttt{pv-vneg(m)}}\\
reads: the componentwise negation of \$1 x R \\
definition:
\begin{quote}\ttfamily\footnotesize\raggedright\noindent
vnb-lambda(\allowbreak{}pvc\_, cartesian(\allowbreak{}vec(\allowbreak{}m), rr), [(\allowbreak{}vneg(\allowbreak{}m))(\allowbreak{}nth(\allowbreak{}1, pvc\_)), -nth(\allowbreak{}2, pvc\_)])
\end{quote}
characterized by \texttt{pv-vneg-fun-cnvs-as-module}  \textit{(definitional)}:
\begin{quote}\ttfamily\footnotesize\raggedright\noindent
forall(\allowbreak{}[r], is-complex-normed-vector-space(\allowbreak{}r) implies scal(\allowbreak{}r) = normed-field-as-commutative-ring(\allowbreak{}rr-normed-field) implies pv-vneg(\allowbreak{}cnvs-as-module(\allowbreak{}r)) in fun(\allowbreak{}cartesian(\allowbreak{}vec(\allowbreak{}r), rr), cartesian(\allowbreak{}vec(\allowbreak{}r), rr)))
\end{quote}
also: \texttt{pv-vneg-value-cnvs-as-module} \texttt{pv-vneg-fun} \texttt{pv-vneg-value} \\
mentioned by 14 result(s)

\smallskip

\noindent\textbf{\texttt{pw-int(pphi\_, a, b)}}\quad\textit{functoid}\index{pw-int(pphi_, a, b)@\texttt{pw-int(pphi\_, a, b)}}\\
reads: the piecewise integral of \$1 from \$2 to \$3 \\
definition:
\begin{quote}\ttfamily\footnotesize\raggedright\noindent
iota(\allowbreak{}pav\_, pav\_ in rr and forsome(\allowbreak{}[paw\_], is-primitive(\allowbreak{}paw\_, pphi\_, a, b) and pav\_ = paw\_(\allowbreak{}b) - paw\_(\allowbreak{}a)))
\end{quote}
characterized by \texttt{pw-int-const}  \textit{(proven)}:
\begin{quote}\ttfamily\footnotesize\raggedright\noindent
forall(\allowbreak{}[a in rr, b in rr], a < b implies forall(\allowbreak{}[pac\_ in rr, pphi\_ in fun(\allowbreak{}ccint(\allowbreak{}a, b), rr)], forall(\allowbreak{}[pay\_ in ccint(\allowbreak{}a, b)], pphi\_(\allowbreak{}pay\_) == pac\_) implies pw-int(\allowbreak{}pphi\_, a, b) = pac\_ * (\allowbreak{}b - a)))
\end{quote}
also: \texttt{pw-int-in-rr} \texttt{pw-int-value} \texttt{pw-int-nonneg} \texttt{pw-int-reflect} \texttt{pw-int-adjacent} \texttt{pw-int-monotone} \texttt{pw-int-abs-bound} \texttt{pw-int-congruence} ... (11 in all) \\
mentioned by 27 result(s)

\smallskip


\section*{Q}
\addcontentsline{toc}{section}{Q}

\noindent\textbf{\texttt{qed}}\quad\textit{tactic}\index{qed@\texttt{qed}}\\
\texttt{(qed 'name)} -- Install the finished proof as theorem `name' and save its replayable script. \\
declared in \texttt{tactics-help.scm}

\smallskip

\noindent\textbf{\texttt{qmod-act(md, sm)}}\quad\textit{functoid}\index{qmod-act(md, sm)@\texttt{qmod-act(md, sm)}}\\
definition:
\begin{quote}\ttfamily\footnotesize\raggedright\noindent
vnb-lambda(\allowbreak{}[qmr\_, qmc\_], cartesian(\allowbreak{}carr(\allowbreak{}scal(\allowbreak{}md)), quotient(\allowbreak{}qmod-setoid(\allowbreak{}md, sm))), (\allowbreak{}qmod-scale(\allowbreak{}md, sm, qmr\_))(\allowbreak{}qmc\_))
\end{quote}
characterized by \texttt{qmod-act-fun0-cna-as-module}  \textit{(definitional)}:
\begin{quote}\ttfamily\footnotesize\raggedright\noindent
forall(\allowbreak{}[r], is-complex-normed-algebra(\allowbreak{}r) implies forall(\allowbreak{}[sm], is-submodule(\allowbreak{}cna-as-module(\allowbreak{}r), sm) implies qmod-act(\allowbreak{}cna-as-module(\allowbreak{}r), sm) in fun(\allowbreak{}cartesian(\allowbreak{}carr(\allowbreak{}scal(\allowbreak{}r)), quotient(\allowbreak{}qmod-setoid(\allowbreak{}cna-as-module(\allowbreak{}r), sm))), quotient(\allowbreak{}qmod-setoid(\allowbreak{}cna-as-module(\allowbreak{}r), sm)))))
\end{quote}
also: \texttt{qmod-act-fun0-sip-as-module} \texttt{qmod-act-fun0-uba-as-module} \texttt{qmod-act-fun0-cnvs-as-module} \texttt{qmod-act-fun0-sip-as-module-csip-as-sip} \texttt{qmod-act-fun0-uba-as-module-uiba-as-uba} \texttt{qmod-act-fun0} \\
mentioned by 41 result(s)

\smallskip

\noindent\textbf{\texttt{qmod-rel(md, sm)}}\quad\textit{functoid}\index{qmod-rel(md, sm)@\texttt{qmod-rel(md, sm)}}\\
reads: the congruence modulo \$2 on \$1 \\
definition:
\begin{quote}\ttfamily\footnotesize\raggedright\noindent
\{qmp\_ in cartesian(\allowbreak{}vec(\allowbreak{}md), vec(\allowbreak{}md)): (\allowbreak{}vadd(\allowbreak{}md))(\allowbreak{}nth(\allowbreak{}1, qmp\_), (\allowbreak{}vneg(\allowbreak{}md))(\allowbreak{}nth(\allowbreak{}2, qmp\_))) in sm\}
\end{quote}
characterized by \texttt{qmod-rel-in-cna-as-module}  \textit{(definitional)}:
\begin{quote}\ttfamily\footnotesize\raggedright\noindent
forall(\allowbreak{}[r], is-complex-normed-algebra(\allowbreak{}r) implies forall(\allowbreak{}[sm], is-submodule(\allowbreak{}cna-as-module(\allowbreak{}r), sm) implies forall(\allowbreak{}[a\_, b\_], a\_ in vec(\allowbreak{}r) implies b\_ in vec(\allowbreak{}r) implies (\allowbreak{}vadd(\allowbreak{}r))(\allowbreak{}a\_, (\allowbreak{}vneg(\allowbreak{}r))(\allowbreak{}b\_)) in sm implies [a\_, b\_] in qmod-rel(\allowbreak{}cna-as-module(\allowbreak{}r), sm))))
\end{quote}
also: \texttt{qmod-rel-in-sip-as-module} \texttt{qmod-rel-in-uba-as-module} \texttt{qmod-rel-in-cnvs-as-module} \texttt{qmod-rel-out-cna-as-module} \texttt{qmod-rel-out-sip-as-module} \texttt{qmod-rel-out-uba-as-module} \texttt{qmod-rel-out-cnvs-as-module} \texttt{qmod-rel-in-sip-as-module-csip-as-sip} ... (21 in all) \\
mentioned by 59 result(s)

\smallskip

\noindent\textbf{\texttt{qmod-scale(md, sm, r)}}\quad\textit{functoid}\index{qmod-scale(md, sm, r)@\texttt{qmod-scale(md, sm, r)}}\\
definition:
\begin{quote}\ttfamily\footnotesize\raggedright\noindent
descend(\allowbreak{}qmod-setoid(\allowbreak{}md, sm), vnb-lambda(\allowbreak{}qma\_, vec(\allowbreak{}md), class(\allowbreak{}qmod-setoid(\allowbreak{}md, sm), (\allowbreak{}act(\allowbreak{}md))(\allowbreak{}r, qma\_))))
\end{quote}
characterized by \texttt{qmod-scale-fun-cna-as-module}  \textit{(definitional)}:
\begin{quote}\ttfamily\footnotesize\raggedright\noindent
forall(\allowbreak{}[r], is-complex-normed-algebra(\allowbreak{}r) implies forall(\allowbreak{}[sm], is-submodule(\allowbreak{}cna-as-module(\allowbreak{}r), sm) implies forall(\allowbreak{}[r\_ in carr(\allowbreak{}scal(\allowbreak{}r))], qmod-scale(\allowbreak{}cna-as-module(\allowbreak{}r), sm, r\_) in fun(\allowbreak{}quotient(\allowbreak{}qmod-setoid(\allowbreak{}cna-as-module(\allowbreak{}r), sm)), quotient(\allowbreak{}qmod-setoid(\allowbreak{}cna-as-module(\allowbreak{}r), sm))))))
\end{quote}
also: \texttt{qmod-scale-fun-sip-as-module} \texttt{qmod-scale-fun-uba-as-module} \texttt{qmod-scale-fun-cnvs-as-module} \texttt{qmod-scale-computes-cna-as-module} \texttt{qmod-scale-computes-sip-as-module} \texttt{qmod-scale-computes-uba-as-module} \texttt{qmod-scale-computes-cnvs-as-module} \texttt{qmod-scale-fun-sip-as-module-csip-as-sip} ... (13 in all) \\
mentioned by 41 result(s)

\smallskip

\noindent\textbf{\texttt{qmod-setoid(md, sm)}}\quad\textit{functoid}\index{qmod-setoid(md, sm)@\texttt{qmod-setoid(md, sm)}}\\
reads: the vectors of \$1 under the congruence modulo \$2 \\
definition:
\begin{quote}\ttfamily\footnotesize\raggedright\noindent
[vec(\allowbreak{}md), qmod-rel(\allowbreak{}md, sm)]
\end{quote}
characterized by \texttt{qmod-setoid-pts-cna-as-module}  \textit{(definitional)}:
\begin{quote}\ttfamily\footnotesize\raggedright\noindent
forall(\allowbreak{}[r], is-complex-normed-algebra(\allowbreak{}r) implies forall(\allowbreak{}[sm], is-submodule(\allowbreak{}cna-as-module(\allowbreak{}r), sm) implies pts(\allowbreak{}qmod-setoid(\allowbreak{}cna-as-module(\allowbreak{}r), sm)) == vec(\allowbreak{}r)))
\end{quote}
also: \texttt{qmod-setoid-pts-sip-as-module} \texttt{qmod-setoid-pts-uba-as-module} \texttt{qmod-setoid-rel-cna-as-module} \texttt{qmod-setoid-rel-sip-as-module} \texttt{qmod-setoid-rel-uba-as-module} \texttt{qmod-setoid-pts-cnvs-as-module} \texttt{qmod-setoid-rel-cnvs-as-module} \texttt{qmod-setoid-is-setoid-cna-as-module} ... (20 in all) \\
mentioned by 504 result(s)

\smallskip

\noindent\textbf{\texttt{qmod-vadd(md, sm)}}\quad\textit{functoid}\index{qmod-vadd(md, sm)@\texttt{qmod-vadd(md, sm)}}\\
definition:
\begin{quote}\ttfamily\footnotesize\raggedright\noindent
descend2(\allowbreak{}qmod-setoid(\allowbreak{}md, sm), vnb-lambda(\allowbreak{}[qma\_, qmb\_], cartesian(\allowbreak{}vec(\allowbreak{}md), vec(\allowbreak{}md)), class(\allowbreak{}qmod-setoid(\allowbreak{}md, sm), (\allowbreak{}vadd(\allowbreak{}md))(\allowbreak{}qma\_, qmb\_))))
\end{quote}
characterized by \texttt{qmod-vadd-fun0-cna-as-module}  \textit{(definitional)}:
\begin{quote}\ttfamily\footnotesize\raggedright\noindent
forall(\allowbreak{}[r], is-complex-normed-algebra(\allowbreak{}r) implies forall(\allowbreak{}[sm], is-submodule(\allowbreak{}cna-as-module(\allowbreak{}r), sm) implies qmod-vadd(\allowbreak{}cna-as-module(\allowbreak{}r), sm) in fun(\allowbreak{}cartesian(\allowbreak{}quotient(\allowbreak{}qmod-setoid(\allowbreak{}cna-as-module(\allowbreak{}r), sm)), quotient(\allowbreak{}qmod-setoid(\allowbreak{}cna-as-module(\allowbreak{}r), sm))), quotient(\allowbreak{}qmod-setoid(\allowbreak{}cna-as-module(\allowbreak{}r), sm)))))
\end{quote}
also: \texttt{qmod-vadd-fun0-sip-as-module} \texttt{qmod-vadd-fun0-uba-as-module} \texttt{qmod-vadd-fun0-cnvs-as-module} \texttt{qmod-vadd-fun0-sip-as-module-csip-as-sip} \texttt{qmod-vadd-fun0-uba-as-module-uiba-as-uba} \texttt{qmod-vadd-fun0} \\
mentioned by 41 result(s)

\smallskip

\noindent\textbf{\texttt{qmod-vneg(md, sm)}}\quad\textit{functoid}\index{qmod-vneg(md, sm)@\texttt{qmod-vneg(md, sm)}}\\
definition:
\begin{quote}\ttfamily\footnotesize\raggedright\noindent
descend(\allowbreak{}qmod-setoid(\allowbreak{}md, sm), vnb-lambda(\allowbreak{}qma\_, vec(\allowbreak{}md), class(\allowbreak{}qmod-setoid(\allowbreak{}md, sm), (\allowbreak{}vneg(\allowbreak{}md))(\allowbreak{}qma\_))))
\end{quote}
characterized by \texttt{qmod-vneg-fun0-cna-as-module}  \textit{(definitional)}:
\begin{quote}\ttfamily\footnotesize\raggedright\noindent
forall(\allowbreak{}[r], is-complex-normed-algebra(\allowbreak{}r) implies forall(\allowbreak{}[sm], is-submodule(\allowbreak{}cna-as-module(\allowbreak{}r), sm) implies qmod-vneg(\allowbreak{}cna-as-module(\allowbreak{}r), sm) in fun(\allowbreak{}quotient(\allowbreak{}qmod-setoid(\allowbreak{}cna-as-module(\allowbreak{}r), sm)), quotient(\allowbreak{}qmod-setoid(\allowbreak{}cna-as-module(\allowbreak{}r), sm)))))
\end{quote}
also: \texttt{qmod-vneg-fun0-sip-as-module} \texttt{qmod-vneg-fun0-uba-as-module} \texttt{qmod-vneg-fun0-cnvs-as-module} \texttt{qmod-vneg-fun0-sip-as-module-csip-as-sip} \texttt{qmod-vneg-fun0-uba-as-module-uiba-as-uba} \texttt{qmod-vneg-fun0} \\
mentioned by 41 result(s)

\smallskip

\noindent\textbf{\texttt{qna-vmul(bqna\_, bqni\_)}}\quad\textit{functoid}\index{qna-vmul(bqna_, bqni_)@\texttt{qna-vmul(bqna\_, bqni\_)}}\\
reads: the product of \$1 modulo \$2 \\
definition:
\begin{quote}\ttfamily\footnotesize\raggedright\noindent
descend2(\allowbreak{}qmod-setoid(\allowbreak{}uba-as-module(\allowbreak{}bqna\_), bqni\_), vnb-lambda(\allowbreak{}[bqnx\_, bqny\_], cartesian(\allowbreak{}vec(\allowbreak{}bqna\_), vec(\allowbreak{}bqna\_)), class(\allowbreak{}qmod-setoid(\allowbreak{}uba-as-module(\allowbreak{}bqna\_), bqni\_), (\allowbreak{}vmul(\allowbreak{}bqna\_))(\allowbreak{}bqnx\_, bqny\_))))
\end{quote}
characterized by \texttt{qna-vmul-fun}  \textit{(proven)}:
\begin{quote}\ttfamily\footnotesize\raggedright\noindent
forall(\allowbreak{}[a], is-unital-banach-algebra(\allowbreak{}a) implies forall(\allowbreak{}[ii], is-two-sided-ideal(\allowbreak{}a, ii) implies qna-vmul(\allowbreak{}a, ii) in fun(\allowbreak{}cartesian(\allowbreak{}quotient(\allowbreak{}qmod-setoid(\allowbreak{}uba-as-module(\allowbreak{}a), ii)), quotient(\allowbreak{}qmod-setoid(\allowbreak{}uba-as-module(\allowbreak{}a), ii))), quotient(\allowbreak{}qmod-setoid(\allowbreak{}uba-as-module(\allowbreak{}a), ii)))))
\end{quote}
also: \texttt{qna-vmul-assoc} \texttt{qna-vmul-act-left} \texttt{qna-vmul-assoc-na} \texttt{qna-vmul-computes} \texttt{qna-vmul-act-right} \texttt{qna-vmul-dist-left} \texttt{qna-vmul-dist-right} \texttt{qna-vmul-act-left-na} ... (11 in all) \\
mentioned by 38 result(s)

\smallskip

\noindent\textbf{\texttt{qna-vnrm(bqna\_, bqni\_)}}\quad\textit{functoid}\index{qna-vnrm(bqna_, bqni_)@\texttt{qna-vnrm(bqna\_, bqni\_)}}\\
reads: the quotient norm of \$1 modulo \$2 \\
definition:
\begin{quote}\ttfamily\footnotesize\raggedright\noindent
vnb-lambda(\allowbreak{}bqnc\_, quotient(\allowbreak{}qmod-setoid(\allowbreak{}uba-as-module(\allowbreak{}bqna\_), bqni\_)), coset-norm(\allowbreak{}bqna\_, bqni\_, bqnc\_))
\end{quote}
characterized by \texttt{qna-vnrm-fun}  \textit{(proven)}:
\begin{quote}\ttfamily\footnotesize\raggedright\noindent
forall(\allowbreak{}[a], is-unital-banach-algebra(\allowbreak{}a) implies forall(\allowbreak{}[ii], is-two-sided-ideal(\allowbreak{}a, ii) implies qna-vnrm(\allowbreak{}a, ii) in fun(\allowbreak{}quotient(\allowbreak{}qmod-setoid(\allowbreak{}uba-as-module(\allowbreak{}a), ii)), rr)))
\end{quote}
also: \texttt{qna-vnrm-apply} \texttt{qna-vnrm-homog} \texttt{qna-vnrm-nonneg} \texttt{qna-vnrm-submult} \texttt{qna-vnrm-definite} \texttt{qna-vnrm-triangle} \\
mentioned by 18 result(s)

\smallskip

\noindent\textbf{\texttt{qnode-ok(mnbe\_)}}\quad\textit{predicate}\index{qnode-ok(mnbe_)@\texttt{qnode-ok(mnbe\_)}}\\
reads: the node \$1 is not an unconfined quantifier \\
defined by \texttt{qnode-ok} \\
characterized by \texttt{qnode-ok}  \textit{(definitional)}:
\begin{quote}\ttfamily\footnotesize\raggedright\noindent
forall(\allowbreak{}[mnbe\_], qnode-ok(\allowbreak{}mnbe\_) iff forall(\allowbreak{}[mnbqb\_ in expr-set], (\allowbreak{}mnbe\_ = enc-all(\allowbreak{}mnbqb\_) implies all-confined(\allowbreak{}mnbqb\_)) and (\allowbreak{}mnbe\_ = enc-ex(\allowbreak{}mnbqb\_) implies ex-confined(\allowbreak{}mnbqb\_)) and (\allowbreak{}mnbe\_ = enc-iota(\allowbreak{}mnbqb\_) implies iota-confined(\allowbreak{}mnbqb\_))))
\end{quote}
mentioned by 4 result(s)

\smallskip

\noindent\textbf{\texttt{qp(qsp\_)}}\quad\textit{functoid}\index{qp(qsp_)@\texttt{qp(qsp\_)}}\\
reads: the \$1-adic numbers Q\_\$1, as a metric space \\
definition:
\begin{quote}\ttfamily\footnotesize\raggedright\noindent
completion(\allowbreak{}nf-metric-space(\allowbreak{}qq-padic(\allowbreak{}qsp\_)))
\end{quote}
characterized by \texttt{is-qp-expansion}  \textit{(definitional)}:
\begin{quote}\ttfamily\footnotesize\raggedright\noindent
forall(\allowbreak{}[qsp\_, qta\_, qtn\_, qtx\_], is-qp-expansion(\allowbreak{}qsp\_, qta\_, qtn\_, qtx\_) iff is-qp-digits(\allowbreak{}qsp\_, qta\_, qtn\_) and converges-to(\allowbreak{}qp(\allowbreak{}qsp\_), vnb-lambda(\allowbreak{}qtm\_, nn, qp-psum(\allowbreak{}qsp\_, qta\_, qtn\_, qtm\_)), qtx\_))
\end{quote}
also: \texttt{qp-unfold} \texttt{qp-abs-dist} \texttt{qp-dist-norm} \texttt{qp-embed-abs} \texttt{qp-zp-compact} \texttt{qp-nf-metric-eq} \texttt{qp-zp-closed-ball} \texttt{qp-zp-totally-bounded} \\
mentioned by 26 result(s)

\smallskip

\noindent\textbf{\texttt{qp-abs(qsp\_)}}\quad\textit{functoid}\index{qp-abs(qsp_)@\texttt{qp-abs(qsp\_)}}\\
reads: the \$1-adic absolute value on Q\_\$1 \\
definition:
\begin{quote}\ttfamily\footnotesize\raggedright\noindent
vnb-lambda(\allowbreak{}qsx\_, quotient(\allowbreak{}cauchy-setoid(\allowbreak{}nf-metric-space(\allowbreak{}qq-padic(\allowbreak{}qsp\_)))), (\allowbreak{}dist(\allowbreak{}completion(\allowbreak{}nf-metric-space(\allowbreak{}qq-padic(\allowbreak{}qsp\_)))))(\allowbreak{}qsx\_, class(\allowbreak{}cauchy-setoid(\allowbreak{}nf-metric-space(\allowbreak{}qq-padic(\allowbreak{}qsp\_))), vnb-lambda(\allowbreak{}qsn\_, nn, 0))))
\end{quote}
characterized by \texttt{qp-abs-le}  \textit{(proven)}:
\begin{quote}\ttfamily\footnotesize\raggedright\noindent
forall(\allowbreak{}[p], zz-prime(\allowbreak{}p) implies forall(\allowbreak{}[x in qp-cs(\allowbreak{}p), nb in nn, bb in rr], forall(\allowbreak{}[k in nn], nb <= k implies (\allowbreak{}padic-abs(\allowbreak{}p))(\allowbreak{}x(\allowbreak{}k)) <= bb) implies (\allowbreak{}qp-abs(\allowbreak{}p))(\allowbreak{}class(\allowbreak{}cauchy-setoid(\allowbreak{}nf-metric-space(\allowbreak{}qq-padic(\allowbreak{}p))), x)) <= bb))
\end{quote}
also: \texttt{qp-abs-fun} \texttt{qp-abs-mul} \texttt{qp-abs-dist} \texttt{qp-abs-null} \texttt{qp-abs-const} \texttt{qp-abs-limit} \texttt{qp-abs-ultra} \texttt{qp-abs-nonneg} ... (10 in all) \\
mentioned by 54 result(s)

\smallskip

\noindent\textbf{\texttt{qp-cs(qsp\_)}}\quad\textit{functoid}\index{qp-cs(qsp_)@\texttt{qp-cs(qsp\_)}}\\
reads: the Cauchy sequences of rationals for the \$1-adic metric \\
definition:
\begin{quote}\ttfamily\footnotesize\raggedright\noindent
cseq(\allowbreak{}nf-metric-space(\allowbreak{}qq-padic(\allowbreak{}qsp\_)))
\end{quote}
characterized by \texttt{qp-cs-fun}  \textit{(proven)}:
\begin{quote}\ttfamily\footnotesize\raggedright\noindent
forall(\allowbreak{}[p], zz-prime(\allowbreak{}p) implies forall(\allowbreak{}[h in qp-cs(\allowbreak{}p)], h in fun(\allowbreak{}nn, qq)))
\end{quote}
also: \texttt{qp-cs-iff} \texttt{qp-cs-set} \texttt{qp-cs-cseq} \texttt{qp-csr-carr} \\
mentioned by 64 result(s)

\smallskip

\noindent\textbf{\texttt{qp-csr(qsp\_)}}\quad\textit{functoid}\index{qp-csr(qsp_)@\texttt{qp-csr(qsp\_)}}\\
reads: the ring of \$1-adic Cauchy sequences with the ideal of null sequences \\
definition:
\begin{quote}\ttfamily\footnotesize\raggedright\noindent
[qp-cs(\allowbreak{}qsp\_), qp-sadd(\allowbreak{}qsp\_), qp-smul(\allowbreak{}qsp\_), qp-sneg(\allowbreak{}qsp\_), vnb-lambda(\allowbreak{}qsn\_, nn, 0), vnb-lambda(\allowbreak{}qsn\_, nn, 1), qp-null(\allowbreak{}qsp\_)]
\end{quote}
characterized by \texttt{qp-csr-add}  \textit{(proven)}:
\begin{quote}\ttfamily\footnotesize\raggedright\noindent
forall(\allowbreak{}[p], zz-prime(\allowbreak{}p) implies add(\allowbreak{}qp-csr(\allowbreak{}p)) == qp-sadd(\allowbreak{}p))
\end{quote}
also: \texttt{qp-csr-idl} \texttt{qp-csr-mul} \texttt{qp-csr-neg} \texttt{qp-csr-one} \texttt{qp-csr-carr} \texttt{qp-csr-zero} \texttt{qp-csr-is-ringoid} \\
mentioned by 19 result(s)

\smallskip

\noindent\textbf{\texttt{qp-digit(qsp\_, qdy\_)}}\quad\textit{functoid}\index{qp-digit(qsp_, qdy_)@\texttt{qp-digit(qsp\_, qdy\_)}}\\
reads: the \$1-adic digit of \$2 \\
definition:
\begin{quote}\ttfamily\footnotesize\raggedright\noindent
iota(\allowbreak{}qdd\_, qdd\_ in zz and 0 <= qdd\_ and qdd\_ < qsp\_ and (\allowbreak{}qp-abs(\allowbreak{}qsp\_))(\allowbreak{}(\allowbreak{}add(\allowbreak{}qp-field(\allowbreak{}qsp\_)))(\allowbreak{}qdy\_, (\allowbreak{}neg(\allowbreak{}qp-field(\allowbreak{}qsp\_)))(\allowbreak{}(\allowbreak{}embed(\allowbreak{}nf-metric-space(\allowbreak{}qq-padic(\allowbreak{}qsp\_))))(\allowbreak{}qdd\_)))) < 1)
\end{quote}
mentioned by 5 result(s)

\smallskip

\noindent\textbf{\texttt{qp-field(qsp\_)}}\quad\textit{functoid}\index{qp-field(qsp_)@\texttt{qp-field(qsp\_)}}\\
reads: the field Q\_\$1 of \$1-adic numbers \\
definition:
\begin{quote}\ttfamily\footnotesize\raggedright\noindent
[carr(\allowbreak{}qp-ring(\allowbreak{}qsp\_)), add(\allowbreak{}qp-ring(\allowbreak{}qsp\_)), mul(\allowbreak{}qp-ring(\allowbreak{}qsp\_)), neg(\allowbreak{}qp-ring(\allowbreak{}qsp\_)), zero(\allowbreak{}qp-ring(\allowbreak{}qsp\_)), one(\allowbreak{}qp-ring(\allowbreak{}qsp\_)), qp-abs(\allowbreak{}qsp\_)]
\end{quote}
characterized by \texttt{qp-field-add}  \textit{(proven)}:
\begin{quote}\ttfamily\footnotesize\raggedright\noindent
forall(\allowbreak{}[p], zz-prime(\allowbreak{}p) implies add(\allowbreak{}qp-field(\allowbreak{}p)) == add(\allowbreak{}qp-ring(\allowbreak{}p)))
\end{quote}
also: \texttt{qp-field-mul} \texttt{qp-field-neg} \texttt{qp-field-one} \texttt{qp-field-carr} \texttt{qp-field-fnrm} \texttt{qp-field-zero} \\
mentioned by 111 result(s)

\smallskip

\noindent\textbf{\texttt{qp-lambda(qsp\_)}}\quad\textit{functoid}\index{qp-lambda(qsp_)@\texttt{qp-lambda(qsp\_)}}\\
reads: the character lambda of the additive group of Q\_\$1 \\
definition:
\begin{quote}\ttfamily\footnotesize\raggedright\noindent
vnb-lambda(\allowbreak{}qly\_, carr(\allowbreak{}qp-field(\allowbreak{}qsp\_)), qp-lambda-at(\allowbreak{}qsp\_, qly\_))
\end{quote}
characterized by \texttt{qp-lambda-onto-roots}  \textit{(proven)}:
\begin{quote}\ttfamily\footnotesize\raggedright\noindent
forall(\allowbreak{}[p], zz-prime(\allowbreak{}p) implies forall(\allowbreak{}[w in mu-pinf(\allowbreak{}p)], forsome(\allowbreak{}[r in qq], (\allowbreak{}qp-lambda(\allowbreak{}p))(\allowbreak{}(\allowbreak{}embed(\allowbreak{}nf-metric-space(\allowbreak{}qq-padic(\allowbreak{}p))))(\allowbreak{}r)) = w)))
\end{quote}
also: \texttt{qp-lambda-separating} \\
mentioned by 8 result(s)

\smallskip

\noindent\textbf{\texttt{qp-lambda-at(qsp\_, qlx\_)}}\quad\textit{functoid}\index{qp-lambda-at(qsp_, qlx_)@\texttt{qp-lambda-at(qsp\_, qlx\_)}}\\
reads: the value at \$2 of the \$1-adic character lambda of Q\_\$1 \\
definition:
\begin{quote}\ttfamily\footnotesize\raggedright\noindent
iota(\allowbreak{}qlw\_, qlw\_ in cc and forsome(\allowbreak{}[qlu\_ in qq], (\allowbreak{}qp-abs(\allowbreak{}qsp\_))(\allowbreak{}(\allowbreak{}add(\allowbreak{}qp-field(\allowbreak{}qsp\_)))(\allowbreak{}qlx\_, (\allowbreak{}neg(\allowbreak{}qp-field(\allowbreak{}qsp\_)))(\allowbreak{}(\allowbreak{}embed(\allowbreak{}nf-metric-space(\allowbreak{}qq-padic(\allowbreak{}qsp\_))))(\allowbreak{}qlu\_)))) <= 1 and qlw\_ = (\allowbreak{}padic-char(\allowbreak{}qsp\_))(\allowbreak{}qlu\_)))
\end{quote}
mentioned by 4 result(s)

\smallskip

\noindent\textbf{\texttt{qp-ms(qsp\_)}}\quad\textit{functoid}\index{qp-ms(qsp_)@\texttt{qp-ms(qsp\_)}}\\
reads: the rationals with the \$1-adic metric \\
definition:
\begin{quote}\ttfamily\footnotesize\raggedright\noindent
nf-metric-space(\allowbreak{}qq-padic(\allowbreak{}qsp\_))
\end{quote}

\smallskip

\noindent\textbf{\texttt{qp-null(qsp\_)}}\quad\textit{functoid}\index{qp-null(qsp_)@\texttt{qp-null(qsp\_)}}\\
reads: the null sequences of rationals for the \$1-adic metric \\
definition:
\begin{quote}\ttfamily\footnotesize\raggedright\noindent
\{qsu\_ in qp-cs(\allowbreak{}qsp\_): cseq-equiv(\allowbreak{}nf-metric-space(\allowbreak{}qq-padic(\allowbreak{}qsp\_)), qsu\_, vnb-lambda(\allowbreak{}qsn\_, nn, 0))\}
\end{quote}
characterized by \texttt{qp-null-cs}  \textit{(proven)}:
\begin{quote}\ttfamily\footnotesize\raggedright\noindent
forall(\allowbreak{}[p], zz-prime(\allowbreak{}p) implies forall(\allowbreak{}[z in qp-null(\allowbreak{}p)], z in qp-cs(\allowbreak{}p)))
\end{quote}
also: \texttt{qp-null-add} \texttt{qp-null-neg} \texttt{qp-null-zero} \texttt{qp-null-power} \texttt{qp-null-mul-left} \texttt{qp-null-mul-right} \texttt{qp-null-member-iff} \\
mentioned by 23 result(s)

\smallskip

\noindent\textbf{\texttt{qp-ord(qsp\_)}}\quad\textit{functoid}\index{qp-ord(qsp_)@\texttt{qp-ord(qsp\_)}}\\
reads: the order function ord\_\$1 on Q\_\$1 \\
definition:
\begin{quote}\ttfamily\footnotesize\raggedright\noindent
vnb-lambda(\allowbreak{}qox\_, carr(\allowbreak{}qp-field(\allowbreak{}qsp\_)), log(\allowbreak{}(\allowbreak{}qp-abs(\allowbreak{}qsp\_))(\allowbreak{}qox\_)) * recip-star(\allowbreak{}log(\allowbreak{}recip-star(\allowbreak{}2))))
\end{quote}
mentioned by 7 result(s)

\smallskip

\noindent\textbf{\texttt{qp-ppow(qsp\_, qsk\_)}}\quad\textit{functoid}\index{qp-ppow(qsp_, qsk_)@\texttt{qp-ppow(qsp\_, qsk\_)}}\\
reads: the power \$1\textasciicircum{}\$2 in Q\_\$1 (\$2 an integer) \\
definition:
\begin{quote}\ttfamily\footnotesize\raggedright\noindent
iota(\allowbreak{}qpy\_, qpy\_ in carr(\allowbreak{}qp-field(\allowbreak{}qsp\_)) and (\allowbreak{}mul(\allowbreak{}qp-field(\allowbreak{}qsp\_)))(\allowbreak{}qpy\_, (\allowbreak{}embed(\allowbreak{}nf-metric-space(\allowbreak{}qq-padic(\allowbreak{}qsp\_))))(\allowbreak{}qsp\_ \textasciicircum{} max(\allowbreak{}-qsk\_, 0))) = (\allowbreak{}embed(\allowbreak{}nf-metric-space(\allowbreak{}qq-padic(\allowbreak{}qsp\_))))(\allowbreak{}qsp\_ \textasciicircum{} max(\allowbreak{}qsk\_, 0)))
\end{quote}
characterized by \texttt{qp-ppow-nat}  \textit{(proven)}:
\begin{quote}\ttfamily\footnotesize\raggedright\noindent
forall(\allowbreak{}[p], zz-prime(\allowbreak{}p) implies forall(\allowbreak{}[n in nn], qp-ppow(\allowbreak{}p, n) = (\allowbreak{}embed(\allowbreak{}nf-metric-space(\allowbreak{}qq-padic(\allowbreak{}p))))(\allowbreak{}p \textasciicircum{} n)))
\end{quote}
mentioned by 22 result(s)

\smallskip

\noindent\textbf{\texttt{qp-psum(qsp\_, qta\_, qtn\_, qtm\_)}}\quad\textit{functoid}\index{qp-psum(qsp_, qta_, qtn_, qtm_)@\texttt{qp-psum(qsp\_, qta\_, qtn\_, qtm\_)}}\\
reads: the \$4-th partial sum of the \$1-adic series with digits \$2 from \$3 \\
definition:
\begin{quote}\ttfamily\footnotesize\raggedright\noindent
sum-ag(\allowbreak{}normed-field-additive-ag(\allowbreak{}qp-field(\allowbreak{}qsp\_)), qp-term(\allowbreak{}qsp\_, qta\_, qtn\_), qtm\_)
\end{quote}
mentioned by 13 result(s)

\smallskip

\noindent\textbf{\texttt{qp-ring(qsp\_)}}\quad\textit{functoid}\index{qp-ring(qsp_)@\texttt{qp-ring(qsp\_)}}\\
definition:
\begin{quote}\ttfamily\footnotesize\raggedright\noindent
ringoid-quotient-ring(\allowbreak{}qp-csr(\allowbreak{}qsp\_))
\end{quote}
characterized by \texttt{qp-ring-is-ring}  \textit{(proven)}:
\begin{quote}\ttfamily\footnotesize\raggedright\noindent
forall(\allowbreak{}[p], zz-prime(\allowbreak{}p) implies is-ring(\allowbreak{}qp-ring(\allowbreak{}p)))
\end{quote}
also: \texttt{qp-ring-is-commutative} \\
mentioned by 26 result(s)

\smallskip

\noindent\textbf{\texttt{qp-sadd(qsp\_)}}\quad\textit{functoid}\index{qp-sadd(qsp_)@\texttt{qp-sadd(qsp\_)}}\\
definition:
\begin{quote}\ttfamily\footnotesize\raggedright\noindent
vnb-lambda(\allowbreak{}[qsu\_, qsv\_], cartesian(\allowbreak{}qp-cs(\allowbreak{}qsp\_), qp-cs(\allowbreak{}qsp\_)), vnb-lambda(\allowbreak{}qsn\_, nn, qsu\_(\allowbreak{}qsn\_) + qsv\_(\allowbreak{}qsn\_)))
\end{quote}
characterized by \texttt{qp-sadd-at}  \textit{(proven)}:
\begin{quote}\ttfamily\footnotesize\raggedright\noindent
forall(\allowbreak{}[p], zz-prime(\allowbreak{}p) implies forall(\allowbreak{}[u in qp-cs(\allowbreak{}p), v in qp-cs(\allowbreak{}p), k in nn], (\allowbreak{}(\allowbreak{}qp-sadd(\allowbreak{}p))(\allowbreak{}u, v))(\allowbreak{}k) = u(\allowbreak{}k) + v(\allowbreak{}k)))
\end{quote}
also: \texttt{qp-sadd-fun} \texttt{qp-sadd-type} \texttt{qp-sadd-value} \\
mentioned by 18 result(s)

\smallskip

\noindent\textbf{\texttt{qp-smul(qsp\_)}}\quad\textit{functoid}\index{qp-smul(qsp_)@\texttt{qp-smul(qsp\_)}}\\
definition:
\begin{quote}\ttfamily\footnotesize\raggedright\noindent
vnb-lambda(\allowbreak{}[qsu\_, qsv\_], cartesian(\allowbreak{}qp-cs(\allowbreak{}qsp\_), qp-cs(\allowbreak{}qsp\_)), vnb-lambda(\allowbreak{}qsn\_, nn, qsu\_(\allowbreak{}qsn\_) * qsv\_(\allowbreak{}qsn\_)))
\end{quote}
characterized by \texttt{qp-smul-at}  \textit{(proven)}:
\begin{quote}\ttfamily\footnotesize\raggedright\noindent
forall(\allowbreak{}[p], zz-prime(\allowbreak{}p) implies forall(\allowbreak{}[u in qp-cs(\allowbreak{}p), v in qp-cs(\allowbreak{}p), k in nn], (\allowbreak{}(\allowbreak{}qp-smul(\allowbreak{}p))(\allowbreak{}u, v))(\allowbreak{}k) = u(\allowbreak{}k) * v(\allowbreak{}k)))
\end{quote}
also: \texttt{qp-smul-fun} \texttt{qp-smul-comm} \texttt{qp-smul-type} \texttt{qp-smul-value} \\
mentioned by 15 result(s)

\smallskip

\noindent\textbf{\texttt{qp-sneg(qsp\_)}}\quad\textit{functoid}\index{qp-sneg(qsp_)@\texttt{qp-sneg(qsp\_)}}\\
definition:
\begin{quote}\ttfamily\footnotesize\raggedright\noindent
vnb-lambda(\allowbreak{}qsu\_, qp-cs(\allowbreak{}qsp\_), vnb-lambda(\allowbreak{}qsn\_, nn, -qsu\_(\allowbreak{}qsn\_)))
\end{quote}
characterized by \texttt{qp-sneg-at}  \textit{(proven)}:
\begin{quote}\ttfamily\footnotesize\raggedright\noindent
forall(\allowbreak{}[p], zz-prime(\allowbreak{}p) implies forall(\allowbreak{}[u in qp-cs(\allowbreak{}p), k in nn], (\allowbreak{}(\allowbreak{}qp-sneg(\allowbreak{}p))(\allowbreak{}u))(\allowbreak{}k) = -u(\allowbreak{}k)))
\end{quote}
also: \texttt{qp-sneg-fun} \texttt{qp-sneg-type} \texttt{qp-sneg-value} \\
mentioned by 17 result(s)

\smallskip

\noindent\textbf{\texttt{qp-term(qsp\_, qta\_, qtn\_)}}\quad\textit{functoid}\index{qp-term(qsp_, qta_, qtn_)@\texttt{qp-term(qsp\_, qta\_, qtn\_)}}\\
reads: the terms a\_(\$3 + j) \$1\textasciicircum{}(\$3 + j) of the \$1-adic series with digits \$2 \\
definition:
\begin{quote}\ttfamily\footnotesize\raggedright\noindent
vnb-lambda(\allowbreak{}qtj\_, nn, (\allowbreak{}mul(\allowbreak{}qp-field(\allowbreak{}qsp\_)))(\allowbreak{}(\allowbreak{}embed(\allowbreak{}nf-metric-space(\allowbreak{}qq-padic(\allowbreak{}qsp\_))))(\allowbreak{}qta\_(\allowbreak{}qtn\_ + qtj\_)), qp-ppow(\allowbreak{}qsp\_, qtn\_ + qtj\_)))
\end{quote}

\smallskip

\noindent\textbf{\texttt{qq-field}}\quad\textit{instance}\index{qq-field@\texttt{qq-field}}\\
an instance of field; predicate \texttt{is-qq-field}; same shape as field

\smallskip

\noindent\textbf{\texttt{qq-line}}\quad\textit{instance}\index{qq-line@\texttt{qq-line}}\\
an instance of vector-space; predicate \texttt{is-qq-line}; same shape as vector-space

\smallskip

\noindent\textbf{\texttt{qq-nv(cpf\_)}}\quad\textit{functoid}\index{qq-nv(cpf_)@\texttt{qq-nv(cpf\_)}}\\
reads: the rationals with the norm \$1 \\
definition:
\begin{quote}\ttfamily\footnotesize\raggedright\noindent
[qq, vnb-lambda(\allowbreak{}[x\_, y\_], cartesian(\allowbreak{}qq, qq), x\_ + y\_), vnb-lambda(\allowbreak{}[x\_, y\_], cartesian(\allowbreak{}qq, qq), x\_ * y\_), vnb-lambda(\allowbreak{}x\_, qq, -x\_), 0, 1, cpf\_]
\end{quote}
mentioned by 36 result(s)

\smallskip

\noindent\textbf{\texttt{qq-padic(cpp\_)}}\quad\textit{functoid}\index{qq-padic(cpp_)@\texttt{qq-padic(cpp\_)}}\\
reads: the rationals with the \$1-adic absolute value \\
definition:
\begin{quote}\ttfamily\footnotesize\raggedright\noindent
[qq, vnb-lambda(\allowbreak{}[x\_, y\_], cartesian(\allowbreak{}qq, qq), x\_ + y\_), vnb-lambda(\allowbreak{}[x\_, y\_], cartesian(\allowbreak{}qq, qq), x\_ * y\_), vnb-lambda(\allowbreak{}x\_, qq, -x\_), 0, 1, padic-abs(\allowbreak{}cpp\_)]
\end{quote}
characterized by \texttt{qq-padic-is-nv-field}  \textit{(proven)}:
\begin{quote}\ttfamily\footnotesize\raggedright\noindent
forall(\allowbreak{}[p], zz-prime(\allowbreak{}p) implies is-nv-field(\allowbreak{}qq-padic(\allowbreak{}p)))
\end{quote}
also: \texttt{qq-padic-is-normed-field} \\
mentioned by 148 result(s)

\smallskip

\noindent\textbf{\texttt{qq-ring}}\quad\textit{instance}\index{qq-ring@\texttt{qq-ring}}\\
an instance of ring; predicate \texttt{is-qq-ring}; same shape as integral-domain

\smallskip

\noindent\textbf{\texttt{qrfl}}\quad\textit{tactic}\index{qrfl@\texttt{qrfl}}\\
\texttt{(qrfl)} -- Quasi-reflexivity: close t = t under the partial-equality definedness reading. \\
declared in \texttt{tactics-help.scm}

\smallskip

\noindent\textbf{\texttt{quietly}}\quad\textit{tactic}\index{quietly@\texttt{quietly}}\\
\texttt{(quietly thunk)} -- Run thunk with state-dump output suppressed; returns its value. \\
declared in \texttt{tactics-help.scm}

\smallskip

\noindent\textbf{\texttt{quot-norm(bqqa\_, bqqi\_, bqqx\_)}}\quad\textit{functoid}\index{quot-norm(bqqa_, bqqi_, bqqx_)@\texttt{quot-norm(bqqa\_, bqqi\_, bqqx\_)}}\\
reads: the norm of the coset of \$3 modulo \$2 in \$1 \\
definition:
\begin{quote}\ttfamily\footnotesize\raggedright\noindent
einf(\allowbreak{}\{bqqw\_ in rr: forsome(\allowbreak{}[bqqy\_ in bqqi\_], bqqw\_ = (\allowbreak{}vnrm(\allowbreak{}bqqa\_))(\allowbreak{}(\allowbreak{}vadd(\allowbreak{}bqqa\_))(\allowbreak{}bqqx\_, (\allowbreak{}vneg(\allowbreak{}bqqa\_))(\allowbreak{}bqqy\_))))\})
\end{quote}
characterized by \texttt{quot-norm-one}  \textit{(proven)}:
\begin{quote}\ttfamily\footnotesize\raggedright\noindent
forall(\allowbreak{}[a], is-unital-banach-algebra(\allowbreak{}a) implies forall(\allowbreak{}[ii], is-left-ideal(\allowbreak{}a, ii) implies not(\allowbreak{}vone(\allowbreak{}a) in ii) implies quot-norm(\allowbreak{}a, ii, vone(\allowbreak{}a)) = 1))
\end{quote}
also: \texttt{quot-norm-unfold} \texttt{quot-norm-le-norm} \texttt{coset-norm-class-quot-norm} \\
mentioned by 7 result(s)

\smallskip

\noindent\textbf{\texttt{quotient(s)}}\quad\textit{functoid}\index{quotient(s)@\texttt{quotient(s)}}\\
definition:
\begin{quote}\ttfamily\footnotesize\raggedright\noindent
image(\allowbreak{}proj(\allowbreak{}s), pts(\allowbreak{}s))
\end{quote}
characterized by \texttt{quotient-rep}  \textit{(proven)}:
\begin{quote}\ttfamily\footnotesize\raggedright\noindent
forall(\allowbreak{}[s], is-setoid(\allowbreak{}s) implies forall(\allowbreak{}[x in quotient(\allowbreak{}s)], forsome(\allowbreak{}[a in pts(\allowbreak{}s)], x = class(\allowbreak{}s, a))))
\end{quote}
also: \texttt{quotient-is-set} \texttt{quotient-unfold} \texttt{class-in-quotient} \texttt{csip-quotient-vec} \texttt{quotient-universal} \texttt{quotient-member-choice} \texttt{l2-quotient-is-ae-quotient} \\
mentioned by 198 result(s)

\smallskip

\noindent\textbf{\texttt{quotient-ag(qag\_, qah\_)}}\quad\textit{functoid}\index{quotient-ag(qag_, qah_)@\texttt{quotient-ag(qag\_, qah\_)}}\\
reads: the quotient group \$1/\$2 \\
definition:
\begin{quote}\ttfamily\footnotesize\raggedright\noindent
module-vector-ag(\allowbreak{}quotient-module(\allowbreak{}ag-zz-module(\allowbreak{}qag\_), qah\_))
\end{quote}
characterized by \texttt{quotient-ag-is-abelian-group}  \textit{(proven)}:
\begin{quote}\ttfamily\footnotesize\raggedright\noindent
forall(\allowbreak{}[g], is-abelian-group(\allowbreak{}g) implies forall(\allowbreak{}[h], is-subgroup(\allowbreak{}g, h) implies is-abelian-group(\allowbreak{}quotient-ag(\allowbreak{}g, h))))
\end{quote}
mentioned by 26 result(s)

\smallskip

\noindent\textbf{\texttt{quotient-ba(bqna\_, bqni\_)}}\quad\textit{functoid}\index{quotient-ba(bqna_, bqni_)@\texttt{quotient-ba(bqna\_, bqni\_)}}\\
reads: the quotient Banach algebra \$1/\$2 \\
definition:
\begin{quote}\ttfamily\footnotesize\raggedright\noindent
[scal(\allowbreak{}quotient-module(\allowbreak{}uba-as-module(\allowbreak{}bqna\_), bqni\_)), vec(\allowbreak{}quotient-module(\allowbreak{}uba-as-module(\allowbreak{}bqna\_), bqni\_)), vadd(\allowbreak{}quotient-module(\allowbreak{}uba-as-module(\allowbreak{}bqna\_), bqni\_)), vzero(\allowbreak{}quotient-module(\allowbreak{}uba-as-module(\allowbreak{}bqna\_), bqni\_)), vneg(\allowbreak{}quotient-module(\allowbreak{}uba-as-module(\allowbreak{}bqna\_), bqni\_)), act(\allowbreak{}quotient-module(\allowbreak{}uba-as-module(\allowbreak{}bqna\_), bqni\_)), qna-vnrm(\allowbreak{}bqna\_, bqni\_), qna-vmul(\allowbreak{}bqna\_, bqni\_), class(\allowbreak{}qmod-setoid(\allowbreak{}uba-as-module(\allowbreak{}bqna\_), bqni\_), vone(\allowbreak{}bqna\_))]
\end{quote}
characterized by \texttt{quotient-ba-is-uba}  \textit{(proven)}:
\begin{quote}\ttfamily\footnotesize\raggedright\noindent
forall(\allowbreak{}[a], is-unital-banach-algebra(\allowbreak{}a) implies forall(\allowbreak{}[ii], is-two-sided-ideal(\allowbreak{}a, ii) implies is-closed(\allowbreak{}ba-metric-space(\allowbreak{}a), ii) implies not(\allowbreak{}vone(\allowbreak{}a) in ii) implies is-unital-banach-algebra(\allowbreak{}quotient-ba(\allowbreak{}a, ii))))
\end{quote}
also: \texttt{quotient-ba-is-uba-proper} \\
mentioned by 23 result(s)

\smallskip

\noindent\textbf{\texttt{quotient-module(md, sm)}}\quad\textit{functoid}\index{quotient-module(md, sm)@\texttt{quotient-module(md, sm)}}\\
reads: the quotient module \$1/\$2 \\
definition:
\begin{quote}\ttfamily\footnotesize\raggedright\noindent
[scal(\allowbreak{}md), quotient(\allowbreak{}qmod-setoid(\allowbreak{}md, sm)), qmod-vadd(\allowbreak{}md, sm), class(\allowbreak{}qmod-setoid(\allowbreak{}md, sm), vzero(\allowbreak{}md)), qmod-vneg(\allowbreak{}md, sm), qmod-act(\allowbreak{}md, sm)]
\end{quote}
characterized by \texttt{quotient-module-is-module-cna-as-module}  \textit{(definitional)}:
\begin{quote}\ttfamily\footnotesize\raggedright\noindent
forall(\allowbreak{}[r], is-complex-normed-algebra(\allowbreak{}r) implies forall(\allowbreak{}[sm], is-submodule(\allowbreak{}cna-as-module(\allowbreak{}r), sm) implies is-module(\allowbreak{}quotient-module(\allowbreak{}cna-as-module(\allowbreak{}r), sm))))
\end{quote}
also: \texttt{quotient-module-is-module-sip-as-module} \texttt{quotient-module-is-module-uba-as-module} \texttt{quotient-module-is-module-cnvs-as-module} \texttt{quotient-module-is-module-sip-as-module-csip-as-sip} \texttt{quotient-module-is-module-uba-as-module-uiba-as-uba} \texttt{quotient-module-is-module} \\
mentioned by 680 result(s)

\smallskip

\noindent\textbf{\texttt{quotient-na(bqna\_, bqni\_)}}\quad\textit{functoid}\index{quotient-na(bqna_, bqni_)@\texttt{quotient-na(bqna\_, bqni\_)}}\\
reads: the quotient normed algebra \$1/\$2 \\
definition:
\begin{quote}\ttfamily\footnotesize\raggedright\noindent
[scal(\allowbreak{}quotient-module(\allowbreak{}uba-as-module(\allowbreak{}bqna\_), bqni\_)), vec(\allowbreak{}quotient-module(\allowbreak{}uba-as-module(\allowbreak{}bqna\_), bqni\_)), vadd(\allowbreak{}quotient-module(\allowbreak{}uba-as-module(\allowbreak{}bqna\_), bqni\_)), vzero(\allowbreak{}quotient-module(\allowbreak{}uba-as-module(\allowbreak{}bqna\_), bqni\_)), vneg(\allowbreak{}quotient-module(\allowbreak{}uba-as-module(\allowbreak{}bqna\_), bqni\_)), act(\allowbreak{}quotient-module(\allowbreak{}uba-as-module(\allowbreak{}bqna\_), bqni\_)), qna-vnrm(\allowbreak{}bqna\_, bqni\_), qna-vmul(\allowbreak{}bqna\_, bqni\_)]
\end{quote}
characterized by \texttt{quotient-na-is-cba}  \textit{(proven)}:
\begin{quote}\ttfamily\footnotesize\raggedright\noindent
forall(\allowbreak{}[a], is-unital-banach-algebra(\allowbreak{}a) implies forall(\allowbreak{}[ii], is-two-sided-ideal(\allowbreak{}a, ii) implies is-closed(\allowbreak{}ba-metric-space(\allowbreak{}a), ii) implies is-complex-banach-algebra(\allowbreak{}quotient-na(\allowbreak{}a, ii))))
\end{quote}
also: \texttt{quotient-na-is-cna} \texttt{quotient-na-unfold} \texttt{quotient-na-complete} \\
mentioned by 21 result(s)

\smallskip


\section*{R}
\addcontentsline{toc}{section}{R}

\noindent\textbf{\texttt{r-exp(y\_)}}\quad\textit{functoid}\index{r-exp(y_)@\texttt{r-exp(y\_)}}\\
reads: the exponential of \$1 \\
definition:
\begin{quote}\ttfamily\footnotesize\raggedright\noindent
iota(\allowbreak{}x\_, x\_ in rr and 0 < x\_ and log(\allowbreak{}x\_) = y\_)
\end{quote}
characterized by \texttt{log-r-exp}  \textit{(proven)}:
\begin{quote}\ttfamily\footnotesize\raggedright\noindent
forall(\allowbreak{}[y\_ in rr], log(\allowbreak{}r-exp(\allowbreak{}y\_)) = y\_)
\end{quote}
also: \texttt{r-exp-add} \texttt{r-exp-log} \texttt{r-exp-neg} \texttt{r-exp-pos} \texttt{r-exp-char} \texttt{r-exp-mono} \texttt{r-exp-prop} \texttt{r-exp-zero} ... (20 in all) \\
mentioned by 39 result(s)

\smallskip

\noindent\textbf{\texttt{ran(f)}}\quad\textit{functoid}\index{ran(f)@\texttt{ran(f)}}\\
definition:
\begin{quote}\ttfamily\footnotesize\raggedright\noindent
image(\allowbreak{}f, dom(\allowbreak{}f))
\end{quote}
characterized by \texttt{range-membership}  \textit{(proven)}:
\begin{quote}\ttfamily\footnotesize\raggedright\noindent
forall(\allowbreak{}[f, a in dom(\allowbreak{}f)], f(\allowbreak{}a) in ran(\allowbreak{}f))
\end{quote}
also: \texttt{ran-subset-codomain} \texttt{fun-range-membership} \\
mentioned by 7 result(s)

\smallskip

\noindent\textbf{\texttt{rank-one-op(kqv\_, kqw\_, kqa\_, kqb\_)}}\quad\textit{functoid}\index{rank-one-op(kqv_, kqw_, kqa_, kqb_)@\texttt{rank-one-op(kqv\_, kqw\_, kqa\_, kqb\_)}}\\
reads: the rank-one operator x |-> <x | \$3> \$4 \\
definition:
\begin{quote}\ttfamily\footnotesize\raggedright\noindent
vnb-lambda(\allowbreak{}kqx\_, vec(\allowbreak{}kqv\_), (\allowbreak{}act(\allowbreak{}kqw\_))(\allowbreak{}(\allowbreak{}ip(\allowbreak{}kqv\_))(\allowbreak{}kqx\_, kqa\_), kqb\_))
\end{quote}
characterized by \texttt{rank-one-op-hs}  \textit{(proven)}:
\begin{quote}\ttfamily\footnotesize\raggedright\noindent
forall(\allowbreak{}[v], is-hilbert-space(\allowbreak{}v) implies forall(\allowbreak{}[a\_, b\_], a\_ in vec(\allowbreak{}v) implies b\_ in vec(\allowbreak{}v) implies is-hilbert-schmidt(\allowbreak{}v, v, rank-one-op(\allowbreak{}v, v, a\_, b\_))))
\end{quote}
also: \texttt{rank-one-op-fun} \texttt{rank-one-op-apply} \texttt{rank-one-op-unfold} \texttt{rank-one-op-bounded} \texttt{rank-one-op-compact} \\
mentioned by 17 result(s)

\smallskip

\noindent\textbf{\texttt{rcont(rct\_)}}\quad\textit{functoid}\index{rcont(rct_)@\texttt{rcont(rct\_)}}\\
reads: the continuous real functions on \$1 \\
definition:
\begin{quote}\ttfamily\footnotesize\raggedright\noindent
\{rcf\_ in fun(\allowbreak{}pts(\allowbreak{}rct\_), rr): is-hom-top-space(\allowbreak{}rct\_, metric-top(\allowbreak{}rr-ms), rcf\_)\}
\end{quote}
characterized by \texttt{is-rcont-subalgebra}  \textit{(definitional)}:
\begin{quote}\ttfamily\footnotesize\raggedright\noindent
forall(\allowbreak{}[rst\_, rsa\_], is-rcont-subalgebra(\allowbreak{}rst\_, rsa\_) iff rsa\_ subset rcont(\allowbreak{}rst\_) and forall(\allowbreak{}[rsc\_ in rr], vnb-lambda(\allowbreak{}rsz\_, pts(\allowbreak{}rst\_), rsc\_) in rsa\_) and forall(\allowbreak{}[rsf\_ in rsa\_, rsg\_ in rsa\_], vnb-lambda(\allowbreak{}rsz\_, pts(\allowbreak{}rst\_), rsf\_(\allowbreak{}rsz\_) + rsg\_(\allowbreak{}rsz\_)) in rsa\_) and forall(\allowbreak{}[rsf\_ in rsa\_, rsg\_ in rsa\_], vnb-lambda(\allowbreak{}rsz\_, pts(\allowbreak{}rst\_), rsf\_(\allowbreak{}rsz\_) * rsg\_(\allowbreak{}rsz\_)) in rsa\_))
\end{quote}
also: \texttt{bsw-rcont-im} \texttt{bsw-rcont-re} \texttt{bsw-rcont-abs} \texttt{bsw-rcont-add} \texttt{bsw-rcont-lin} \texttt{bsw-rcont-set} \texttt{bsw-rcont-elim} \texttt{bsw-rcont-const} ... (17 in all) \\
mentioned by 30 result(s)

\smallskip

\noindent\textbf{\texttt{rcont-closure(rst\_, rsa\_)}}\quad\textit{functoid}\index{rcont-closure(rst_, rsa_)@\texttt{rcont-closure(rst\_, rsa\_)}}\\
reads: the uniform closure of \$2 in the continuous real functions on \$1 \\
definition:
\begin{quote}\ttfamily\footnotesize\raggedright\noindent
\{rsh\_ in rcont(\allowbreak{}rst\_): forall(\allowbreak{}[rse\_], pos-rr(\allowbreak{}rse\_) implies forsome(\allowbreak{}[rsg\_ in rsa\_], forall(\allowbreak{}[rsy\_ in pts(\allowbreak{}rst\_)], abs(\allowbreak{}rsh\_(\allowbreak{}rsy\_) - rsg\_(\allowbreak{}rsy\_)) < rse\_)))\}
\end{quote}
characterized by \texttt{bsw-rcont-closure-closed}  \textit{(proven)}:
\begin{quote}\ttfamily\footnotesize\raggedright\noindent
forall(\allowbreak{}[t, aa, h], aa subset rcont(\allowbreak{}t) implies h in rcont(\allowbreak{}t) implies forall(\allowbreak{}[swe\_], pos-rr(\allowbreak{}swe\_) implies forsome(\allowbreak{}[swg\_ in rcont-closure(\allowbreak{}t, aa)], forall(\allowbreak{}[rsy\_ in pts(\allowbreak{}t)], abs(\allowbreak{}h(\allowbreak{}rsy\_) - swg\_(\allowbreak{}rsy\_)) < swe\_))) implies h in rcont-closure(\allowbreak{}t, aa))
\end{quote}
also: \texttt{bsw-rcont-closure-unfold} \texttt{bsw-rcont-closure-contains} \\
mentioned by 11 result(s)

\smallskip

\noindent\textbf{\texttt{real-part}}\quad\textit{operator}\index{real-part@\texttt{real-part}}\\
a term-forming head \\
characterized by \texttt{real-part-def}  \textit{(definitional)}:
\begin{quote}\ttfamily\footnotesize\raggedright\noindent
forall(\allowbreak{}[z], real-part(\allowbreak{}z) == (\allowbreak{}z + conjugate(\allowbreak{}z)) * recip(\allowbreak{}2))
\end{quote}
also: \texttt{real-part-in-rr} \texttt{cc-cos-real-part} \\
mentioned by 307 result(s)

\smallskip

\noindent\textbf{\texttt{received}}\quad\textit{defined-fn}\index{received@\texttt{received}}\\
defined by \texttt{received-def}

\smallskip

\noindent\textbf{\texttt{recip}}\quad\textit{operator}\index{recip@\texttt{recip}}\\
a term-forming head \\
characterized by \texttt{recip-gamma-def}  \textit{(definitional)}:
\begin{quote}\ttfamily\footnotesize\raggedright\noindent
recip-gamma == vnb-lambda(\allowbreak{}egz\_, cc, (\allowbreak{}egz\_ * cc-exp(\allowbreak{}euler-gamma * egz\_)) * (\allowbreak{}prod-fun(\allowbreak{}cc, vnb-lambda(\allowbreak{}c6k\_, nn, vnb-lambda(\allowbreak{}c6x\_, cc, (\allowbreak{}1 - c6x\_ * recip(\allowbreak{}(\allowbreak{}vnb-lambda(\allowbreak{}c6gk\_, nn, -succ(\allowbreak{}c6gk\_)))(\allowbreak{}c6k\_))) * cc-exp(\allowbreak{}cc-series-partial-sum(\allowbreak{}vnb-lambda(\allowbreak{}clhk\_, nn, recip-star(\allowbreak{}clhk\_) * (\allowbreak{}c6x\_ * recip(\allowbreak{}(\allowbreak{}vnb-lambda(\allowbreak{}c6gk\_, nn, -succ(\allowbreak{}c6gk\_)))(\allowbreak{}c6k\_))) \textasciicircum{} clhk\_), succ(\allowbreak{}(\allowbreak{}vnb-lambda(\allowbreak{}c6gk\_, nn, 1))(\allowbreak{}c6k\_)))) - 1))))(\allowbreak{}egz\_))
\end{quote}
also: \texttt{cc-recip-closed} \texttt{qq-recip-closed} \texttt{rr-recip-closed} \texttt{cc-recip-inverse} \texttt{qq-recip-inverse} \texttt{rr-recip-inverse} \texttt{recip-defined-nonzero} \texttt{recip-near} ... (68 in all) \\
mentioned by 650 result(s)

\smallskip

\noindent\textbf{\texttt{recip-gamma}}\quad\textit{constant}\index{recip-gamma@\texttt{recip-gamma}}\\
reads: the reciprocal of the Gamma function \\
defined by \texttt{recip-gamma-def} \\
characterized by \texttt{recip-gamma-nonzero-rhp}  \textit{(proven)}:
\begin{quote}\ttfamily\footnotesize\raggedright\noindent
forall(\allowbreak{}[c6sw\_ in cc], 0 < real-part(\allowbreak{}c6sw\_) implies not(\allowbreak{}recip-gamma(\allowbreak{}c6sw\_) = 0))
\end{quote}
mentioned by 7 result(s)

\smallskip

\noindent\textbf{\texttt{recip-star(u\_)}}\quad\textit{functoid}\index{recip-star(u_)@\texttt{recip-star(u\_)}}\\
reads: the totalised reciprocal of \$1 \\
definition:
\begin{quote}\ttfamily\footnotesize\raggedright\noindent
if(\allowbreak{}u\_ = 0, 0, recip(\allowbreak{}u\_))
\end{quote}
characterized by \texttt{recip-star-zero}  \textit{(proven)}:
\begin{quote}\ttfamily\footnotesize\raggedright\noindent
recip-star(\allowbreak{}0) = 0
\end{quote}
also: \texttt{recip-star-in-rr} \texttt{recip-star-scale} \texttt{recip-star-value} \texttt{recip-star-le-one} \texttt{recip-star-nn-pos} \texttt{recip-star-inverse} \texttt{recip-star-nn-nonneg} \texttt{recip-star-lam-in-fun} ... (12 in all) \\
mentioned by 90 result(s)

\smallskip

\noindent\textbf{\texttt{rectangles(xs, ca, ys, cb)}}\quad\textit{functoid}\index{rectangles(xs, ca, ys, cb)@\texttt{rectangles(xs, ca, ys, cb)}}\\
reads: the measurable rectangles of \$1 x \$3 for \$2 and \$4 \\
definition:
\begin{quote}\ttfamily\footnotesize\raggedright\noindent
\{rcr\_ in power(\allowbreak{}cartesian(\allowbreak{}xs, ys)): forsome(\allowbreak{}[rca\_ in ca, rcb\_ in cb], rcr\_ = cartesian(\allowbreak{}rca\_, rcb\_))\}
\end{quote}
characterized by \texttt{rectangles-unfold}  \textit{(proven)}:
\begin{quote}\ttfamily\footnotesize\raggedright\noindent
forall(\allowbreak{}[xs, ca, ys, cb], rectangles(\allowbreak{}xs, ca, ys, cb) == \{rcr\_ in power(\allowbreak{}cartesian(\allowbreak{}xs, ys)): forsome(\allowbreak{}[rca\_ in ca, rcb\_ in cb], rcr\_ = cartesian(\allowbreak{}rca\_, rcb\_))\})
\end{quote}
also: \texttt{rectangles-pi-system} \texttt{rectangles-membership} \texttt{rectangle-in-rectangles} \\
mentioned by 8 result(s)

\smallskip

\noindent\textbf{\texttt{reduce}}\quad\textit{defined-fn}\index{reduce@\texttt{reduce}}\\
defined by \texttt{reduce-one} \texttt{reduce-succ} \\
characterized by \texttt{reduce-one}  \textit{(definitional)}:
\begin{quote}\ttfamily\footnotesize\raggedright\noindent
forall(\allowbreak{}[op, f], reduce(\allowbreak{}op, f, 1) == f(\allowbreak{}0))
\end{quote}
also: \texttt{reduce-succ} \texttt{sum-ag-as-reduce-ring-additive-ag} \texttt{sum-ag-as-reduce-ring-additive-ag-uba-as-ring} \texttt{sum-ag-as-reduce-ring-additive-ag-uba-as-ring-uiba-as-uba} \texttt{sum-ag-as-reduce} \texttt{sum-ag-as-reduce-ind} \\
mentioned by 38 result(s)

\smallskip

\noindent\textbf{\texttt{rel(s)}}\quad\textit{accessor}\index{rel(s)@\texttt{rel(s)}}\\
reads: the equivalence relation of \$1 \\
a structure slot \\
characterized by \texttt{sip-setoid-rel-csip-as-sip}  \textit{(definitional)}:
\begin{quote}\ttfamily\footnotesize\raggedright\noindent
forall(\allowbreak{}[r], is-complex-semi-inner-product-space(\allowbreak{}r) implies rel(\allowbreak{}sip-setoid(\allowbreak{}csip-as-sip(\allowbreak{}r))) == sip-rel(\allowbreak{}csip-as-sip(\allowbreak{}r)))
\end{quote}
also: \texttt{qmod-setoid-rel-cna-as-module} \texttt{qmod-setoid-rel-sip-as-module} \texttt{qmod-setoid-rel-uba-as-module} \texttt{qmod-setoid-rel-cnvs-as-module} \texttt{qmod-setoid-rel-sip-as-module-csip-as-sip} \texttt{qmod-setoid-rel-uba-as-module-uiba-as-uba} \texttt{eps-setoid-rel} \texttt{r6b-setoid-rel} ... (12 in all) \\
mentioned by 52 result(s)

\smallskip

\noindent\textbf{\texttt{rel-and-elim-k(mnkt\_, mnkhs\_, mnkc\_)}}\quad\textit{predicate}\index{rel-and-elim-k(mnkt_, mnkhs_, mnkc_)@\texttt{rel-and-elim-k(mnkt\_, mnkhs\_, mnkc\_)}}\\
reads: the kernel's and-elim relates \$2 to \$3 \\
defined by \texttt{rel-and-elim-k} \\
characterized by \texttt{rel-and-elim-k}  \textit{(definitional)}:
\begin{quote}\ttfamily\footnotesize\raggedright\noindent
forall(\allowbreak{}[mnkt\_, mnkhs\_, mnkc\_], rel-and-elim-k(\allowbreak{}mnkt\_, mnkhs\_, mnkc\_) iff mnkc\_ in seq-set and forsome(\allowbreak{}[mnkg\_, mnkf\_, mnka\_, mnkb\_], mnkc\_ = [mnkg\_, mnkf\_] and enc-and(\allowbreak{}mnka\_, mnkb\_) in mnkg\_ and mnkhs\_ = [[ctx-add(\allowbreak{}ctx-add(\allowbreak{}difference(\allowbreak{}mnkg\_, pair(\allowbreak{}enc-and(\allowbreak{}mnka\_, mnkb\_), enc-and(\allowbreak{}mnka\_, mnkb\_))), mnka\_), mnkb\_), mnkf\_]]))
\end{quote}
mentioned by 3 result(s)

\smallskip

\noindent\textbf{\texttt{rel-big-union-mem-elim(mit\_, mihs\_, mic\_)}}\quad\textit{predicate}\index{rel-big-union-mem-elim(mit_, mihs_, mic_)@\texttt{rel-big-union-mem-elim(mit\_, mihs\_, mic\_)}}\\
reads: the kernel's big-union-mem-elim relates \$2 to \$3 \\
defined by \texttt{rel-big-union-mem-elim} \\
characterized by \texttt{rel-big-union-mem-elim}  \textit{(definitional)}:
\begin{quote}\ttfamily\footnotesize\raggedright\noindent
forall(\allowbreak{}[mit\_, mihs\_, mic\_], rel-big-union-mem-elim(\allowbreak{}mit\_, mihs\_, mic\_) iff mic\_ in seq-set and forsome(\allowbreak{}[mig\_, mif\_, mia\_, mib\_, mix\_, mie\_], mic\_ = [mig\_, mif\_] and enc-in(\allowbreak{}mix\_, enc-bu(\allowbreak{}mia\_, mib\_)) in mig\_ and nbg-ok(\allowbreak{}mib\_) and mie\_ in nn and not(\allowbreak{}mie\_ in union(\allowbreak{}ctx-fv(\allowbreak{}mig\_), expr-fv(\allowbreak{}mif\_))) and mihs\_ = [[ctx-add(\allowbreak{}ctx-add(\allowbreak{}difference(\allowbreak{}mig\_, pair(\allowbreak{}enc-in(\allowbreak{}mix\_, enc-bu(\allowbreak{}mia\_, mib\_)), enc-in(\allowbreak{}mix\_, enc-bu(\allowbreak{}mia\_, mib\_)))), enc-in(\allowbreak{}efvar(\allowbreak{}mie\_), mia\_)), enc-in(\allowbreak{}mix\_, expr-open(\allowbreak{}mib\_, 0, efvar(\allowbreak{}mie\_)))), mif\_]]))
\end{quote}
mentioned by 3 result(s)

\smallskip

\noindent\textbf{\texttt{rel-big-union-mem-intro(mit\_, mihs\_, mic\_)}}\quad\textit{predicate}\index{rel-big-union-mem-intro(mit_, mihs_, mic_)@\texttt{rel-big-union-mem-intro(mit\_, mihs\_, mic\_)}}\\
reads: the kernel's big-union-mem-intro relates \$2 to \$3 \\
defined by \texttt{rel-big-union-mem-intro} \\
characterized by \texttt{rel-big-union-mem-intro}  \textit{(definitional)}:
\begin{quote}\ttfamily\footnotesize\raggedright\noindent
forall(\allowbreak{}[mit\_, mihs\_, mic\_], rel-big-union-mem-intro(\allowbreak{}mit\_, mihs\_, mic\_) iff mic\_ in seq-set and forsome(\allowbreak{}[mig\_, mia\_, mib\_, mix\_, miw\_], mic\_ = [mig\_, enc-in(\allowbreak{}mix\_, enc-bu(\allowbreak{}mia\_, mib\_))] and nbg-ok(\allowbreak{}mib\_) and mihs\_ = [[mig\_, enc-in(\allowbreak{}miw\_, mia\_)], [mig\_, enc-in(\allowbreak{}mix\_, expr-open(\allowbreak{}mib\_, 0, miw\_))]]))
\end{quote}
mentioned by 4 result(s)

\smallskip

\noindent\textbf{\texttt{rel-big-union-sethood(mit\_, mihs\_, mic\_)}}\quad\textit{predicate}\index{rel-big-union-sethood(mit_, mihs_, mic_)@\texttt{rel-big-union-sethood(mit\_, mihs\_, mic\_)}}\\
reads: the kernel's big-union-sethood relates \$2 to \$3 \\
defined by \texttt{rel-big-union-sethood} \\
characterized by \texttt{rel-big-union-sethood}  \textit{(definitional)}:
\begin{quote}\ttfamily\footnotesize\raggedright\noindent
forall(\allowbreak{}[mit\_, mihs\_, mic\_], rel-big-union-sethood(\allowbreak{}mit\_, mihs\_, mic\_) iff mic\_ in seq-set and forsome(\allowbreak{}[mig\_, mia\_, mib\_], mic\_ = [mig\_, enc-in(\allowbreak{}enc-bu(\allowbreak{}mia\_, mib\_), enc-setc)] and nbg-ok(\allowbreak{}mib\_) and mihs\_ = [[mig\_, enc-in(\allowbreak{}mia\_, enc-setc)], [mig\_, enc-all(\allowbreak{}enc-imp(\allowbreak{}enc-in(\allowbreak{}ebvar(\allowbreak{}0), mia\_), enc-in(\allowbreak{}mib\_, enc-setc)))]]))
\end{quote}
mentioned by 3 result(s)

\smallskip

\noindent\textbf{\texttt{rel-cartesian-decompose-1(mit\_, mihs\_, mic\_)}}\quad\textit{predicate}\index{rel-cartesian-decompose-1(mit_, mihs_, mic_)@\texttt{rel-cartesian-decompose-1(mit\_, mihs\_, mic\_)}}\\
reads: the kernel's cartesian-decompose, at one position, relates \$2 to \$3 \\
defined by \texttt{rel-cartesian-decompose-1} \\
characterized by \texttt{rel-cartesian-decompose-1}  \textit{(definitional)}:
\begin{quote}\ttfamily\footnotesize\raggedright\noindent
forall(\allowbreak{}[mit\_, mihs\_, mic\_], rel-cartesian-decompose-1(\allowbreak{}mit\_, mihs\_, mic\_) iff rel-rewrite-1(\allowbreak{}j-cart(\allowbreak{}mit\_), mihs\_, mic\_))
\end{quote}
mentioned by 3 result(s)

\smallskip

\noindent\textbf{\texttt{rel-cartesian-decompose-7(mit\_, mihs\_, mic\_)}}\quad\textit{predicate}\index{rel-cartesian-decompose-7(mit_, mihs_, mic_)@\texttt{rel-cartesian-decompose-7(mit\_, mihs\_, mic\_)}}\\
reads: the kernel's cartesian-decompose, at one position at any stage, relates \$2 to \$3 \\
defined by \texttt{rel-cartesian-decompose-7} \\
characterized by \texttt{rel-cartesian-decompose-7}  \textit{(definitional)}:
\begin{quote}\ttfamily\footnotesize\raggedright\noindent
forall(\allowbreak{}[mit\_, mihs\_, mic\_], rel-cartesian-decompose-7(\allowbreak{}mit\_, mihs\_, mic\_) iff rel-rewrite-7(\allowbreak{}j-cart(\allowbreak{}mit\_), mihs\_, mic\_))
\end{quote}
mentioned by 3 result(s)

\smallskip

\noindent\textbf{\texttt{rel-cartesian-elim(mit\_, mihs\_, mic\_)}}\quad\textit{predicate}\index{rel-cartesian-elim(mit_, mihs_, mic_)@\texttt{rel-cartesian-elim(mit\_, mihs\_, mic\_)}}\\
reads: the kernel's cartesian-elim relates \$2 to \$3 \\
defined by \texttt{rel-cartesian-elim} \\
characterized by \texttt{rel-cartesian-elim}  \textit{(definitional)}:
\begin{quote}\ttfamily\footnotesize\raggedright\noindent
forall(\allowbreak{}[mit\_, mihs\_, mic\_], rel-cartesian-elim(\allowbreak{}mit\_, mihs\_, mic\_) iff mic\_ in seq-set and forsome(\allowbreak{}[mig\_, mif\_, mie\_, mia\_, mik\_], mic\_ = [mig\_, mif\_] and enc-in(\allowbreak{}enc-list(\allowbreak{}mie\_), enc-cart(\allowbreak{}mia\_)) in mig\_ and length(\allowbreak{}mie\_) = length(\allowbreak{}mia\_) and mik\_ in nn and 1 <= mik\_ and mik\_ <= length(\allowbreak{}mie\_) and mihs\_ = [[ctx-add(\allowbreak{}mig\_, enc-in(\allowbreak{}nth(\allowbreak{}mik\_, mie\_), nth(\allowbreak{}mik\_, mia\_))), mif\_]]))
\end{quote}
mentioned by 3 result(s)

\smallskip

\noindent\textbf{\texttt{rel-cartesian-intro(mit\_, mihs\_, mic\_)}}\quad\textit{predicate}\index{rel-cartesian-intro(mit_, mihs_, mic_)@\texttt{rel-cartesian-intro(mit\_, mihs\_, mic\_)}}\\
reads: the kernel's cartesian-intro relates \$2 to \$3 \\
defined by \texttt{rel-cartesian-intro} \\
characterized by \texttt{rel-cartesian-intro}  \textit{(definitional)}:
\begin{quote}\ttfamily\footnotesize\raggedright\noindent
forall(\allowbreak{}[mit\_, mihs\_, mic\_], rel-cartesian-intro(\allowbreak{}mit\_, mihs\_, mic\_) iff mic\_ in seq-set and forsome(\allowbreak{}[mig\_, mie\_, mia\_], mic\_ = [mig\_, enc-in(\allowbreak{}enc-list(\allowbreak{}mie\_), enc-cart(\allowbreak{}mia\_))] and length(\allowbreak{}mie\_) = length(\allowbreak{}mia\_) and 1 <= length(\allowbreak{}mia\_) and length(\allowbreak{}mihs\_) = length(\allowbreak{}mie\_) and forall(\allowbreak{}[mii\_ in nn], 1 <= mii\_ implies mii\_ <= length(\allowbreak{}mie\_) implies nth(\allowbreak{}mii\_, mihs\_) = [mig\_, enc-in(\allowbreak{}nth(\allowbreak{}mii\_, mie\_), nth(\allowbreak{}mii\_, mia\_))])))
\end{quote}
mentioned by 3 result(s)

\smallskip

\noindent\textbf{\texttt{rel-comp-mem-elim(mit\_, mihs\_, mic\_)}}\quad\textit{predicate}\index{rel-comp-mem-elim(mit_, mihs_, mic_)@\texttt{rel-comp-mem-elim(mit\_, mihs\_, mic\_)}}\\
reads: the kernel's comp-mem-elim relates \$2 to \$3 \\
defined by \texttt{rel-comp-mem-elim} \\
characterized by \texttt{rel-comp-mem-elim}  \textit{(definitional)}:
\begin{quote}\ttfamily\footnotesize\raggedright\noindent
forall(\allowbreak{}[mit\_, mihs\_, mic\_], rel-comp-mem-elim(\allowbreak{}mit\_, mihs\_, mic\_) iff mic\_ in seq-set and forsome(\allowbreak{}[mig\_, mif\_, mip\_, miy\_], mic\_ = [mig\_, mif\_] and enc-in(\allowbreak{}miy\_, enc-comp(\allowbreak{}mip\_)) in mig\_ and nbg-ok(\allowbreak{}mip\_) and mihs\_ = [[ctx-add(\allowbreak{}ctx-add(\allowbreak{}difference(\allowbreak{}mig\_, pair(\allowbreak{}enc-in(\allowbreak{}miy\_, enc-comp(\allowbreak{}mip\_)), enc-in(\allowbreak{}miy\_, enc-comp(\allowbreak{}mip\_)))), enc-in(\allowbreak{}miy\_, enc-setc)), expr-open(\allowbreak{}mip\_, 0, miy\_)), mif\_]]))
\end{quote}
mentioned by 3 result(s)

\smallskip

\noindent\textbf{\texttt{rel-comp-mem-intro(mit\_, mihs\_, mic\_)}}\quad\textit{predicate}\index{rel-comp-mem-intro(mit_, mihs_, mic_)@\texttt{rel-comp-mem-intro(mit\_, mihs\_, mic\_)}}\\
reads: the kernel's comp-mem-intro relates \$2 to \$3 \\
defined by \texttt{rel-comp-mem-intro} \\
characterized by \texttt{rel-comp-mem-intro}  \textit{(definitional)}:
\begin{quote}\ttfamily\footnotesize\raggedright\noindent
forall(\allowbreak{}[mit\_, mihs\_, mic\_], rel-comp-mem-intro(\allowbreak{}mit\_, mihs\_, mic\_) iff mic\_ in seq-set and forsome(\allowbreak{}[mig\_, mip\_, miy\_], mic\_ = [mig\_, enc-in(\allowbreak{}miy\_, enc-comp(\allowbreak{}mip\_))] and nbg-ok(\allowbreak{}mip\_) and mihs\_ = [[mig\_, enc-in(\allowbreak{}miy\_, enc-setc)], [mig\_, expr-open(\allowbreak{}mip\_, 0, miy\_)]]))
\end{quote}
mentioned by 3 result(s)

\smallskip

\noindent\textbf{\texttt{rel-detach(mct\_, mchs\_, mcc\_)}}\quad\textit{predicate}\index{rel-detach(mct_, mchs_, mcc_)@\texttt{rel-detach(mct\_, mchs\_, mcc\_)}}\\
reads: the kernel's detach relates \$2 to \$3 \\
defined by \texttt{rel-detach} \\
characterized by \texttt{rel-detach}  \textit{(definitional)}:
\begin{quote}\ttfamily\footnotesize\raggedright\noindent
forall(\allowbreak{}[mct\_, mchs\_, mcc\_], rel-detach(\allowbreak{}mct\_, mchs\_, mcc\_) iff mcc\_ in seq-set and forsome(\allowbreak{}[mcg\_, mcf\_, mcp\_, mcq\_], mcc\_ = [mcg\_, mcf\_] and enc-imp(\allowbreak{}mcp\_, mcq\_) in mcg\_ and mcp\_ in mcg\_ and mchs\_ = [[ctx-add(\allowbreak{}mcg\_, mcq\_), mcf\_]]))
\end{quote}
mentioned by 3 result(s)

\smallskip

\noindent\textbf{\texttt{rel-forall-elim-k(mnkt\_, mnkhs\_, mnkc\_)}}\quad\textit{predicate}\index{rel-forall-elim-k(mnkt_, mnkhs_, mnkc_)@\texttt{rel-forall-elim-k(mnkt\_, mnkhs\_, mnkc\_)}}\\
reads: the kernel's forall-elim at a certified term relates \$2 to \$3 \\
defined by \texttt{rel-forall-elim-k} \\
characterized by \texttt{rel-forall-elim-k}  \textit{(definitional)}:
\begin{quote}\ttfamily\footnotesize\raggedright\noindent
forall(\allowbreak{}[mnkt\_, mnkhs\_, mnkc\_], rel-forall-elim-k(\allowbreak{}mnkt\_, mnkhs\_, mnkc\_) iff mnkc\_ in seq-set and forsome(\allowbreak{}[mnkg\_, mnkf\_, mnkb\_, mnktm\_, mnkh\_], mnkc\_ = [mnkg\_, mnkf\_] and enc-all(\allowbreak{}mnkb\_) in mnkg\_ and mnktm\_ in clt and mnkh\_ in mnkg\_ and mnkh\_ in cert-set(\allowbreak{}mnkt\_, mnkg\_, mnktm\_) and mnkhs\_ = [[ctx-add(\allowbreak{}mnkg\_, expr-open(\allowbreak{}mnkb\_, 0, mnktm\_)), mnkf\_]]))
\end{quote}
mentioned by 3 result(s)

\smallskip

\noindent\textbf{\texttt{rel-forsome-elim-k(mnkt\_, mnkhs\_, mnkc\_)}}\quad\textit{predicate}\index{rel-forsome-elim-k(mnkt_, mnkhs_, mnkc_)@\texttt{rel-forsome-elim-k(mnkt\_, mnkhs\_, mnkc\_)}}\\
reads: the kernel's forsome-elim relates \$2 to \$3 \\
defined by \texttt{rel-forsome-elim-k} \\
characterized by \texttt{rel-forsome-elim-k}  \textit{(definitional)}:
\begin{quote}\ttfamily\footnotesize\raggedright\noindent
forall(\allowbreak{}[mnkt\_, mnkhs\_, mnkc\_], rel-forsome-elim-k(\allowbreak{}mnkt\_, mnkhs\_, mnkc\_) iff mnkc\_ in seq-set and forsome(\allowbreak{}[mnkg\_, mnkf\_, mnkb\_, mnky\_], mnkc\_ = [mnkg\_, mnkf\_] and enc-ex(\allowbreak{}mnkb\_) in mnkg\_ and mnky\_ in nn and not(\allowbreak{}mnky\_ in union(\allowbreak{}ctx-fv(\allowbreak{}difference(\allowbreak{}mnkg\_, pair(\allowbreak{}enc-ex(\allowbreak{}mnkb\_), enc-ex(\allowbreak{}mnkb\_)))), union(\allowbreak{}expr-fv(\allowbreak{}mnkf\_), expr-fv(\allowbreak{}enc-ex(\allowbreak{}mnkb\_))))) and mnkhs\_ = [[ctx-add(\allowbreak{}difference(\allowbreak{}mnkg\_, pair(\allowbreak{}enc-ex(\allowbreak{}mnkb\_), enc-ex(\allowbreak{}mnkb\_))), expr-open(\allowbreak{}mnkb\_, 0, efvar(\allowbreak{}mnky\_))), mnkf\_]]))
\end{quote}
mentioned by 3 result(s)

\smallskip

\noindent\textbf{\texttt{rel-forsome-intro-k(mnkt\_, mnkhs\_, mnkc\_)}}\quad\textit{predicate}\index{rel-forsome-intro-k(mnkt_, mnkhs_, mnkc_)@\texttt{rel-forsome-intro-k(mnkt\_, mnkhs\_, mnkc\_)}}\\
reads: the kernel's forsome-intro with a certified witness relates \$2 to \$3 \\
defined by \texttt{rel-forsome-intro-k} \\
characterized by \texttt{rel-forsome-intro-k}  \textit{(definitional)}:
\begin{quote}\ttfamily\footnotesize\raggedright\noindent
forall(\allowbreak{}[mnkt\_, mnkhs\_, mnkc\_], rel-forsome-intro-k(\allowbreak{}mnkt\_, mnkhs\_, mnkc\_) iff mnkc\_ in seq-set and forsome(\allowbreak{}[mnkg\_, mnkb\_, mnktm\_], mnkc\_ = [mnkg\_, enc-ex(\allowbreak{}mnkb\_)] and mnktm\_ in clt and mnkhs\_ = [[mnkg\_, expr-open(\allowbreak{}mnkb\_, 0, mnktm\_)]] and (\allowbreak{}forsome(\allowbreak{}[mnkf2\_ in true-reach(\allowbreak{}expr-open(\allowbreak{}mnkb\_, 0, mnktm\_))], strict-atom(\allowbreak{}mnkf2\_, mnktm\_)) or forsome(\allowbreak{}[mnkh\_ in mnkg\_], mnkh\_ in cert-set(\allowbreak{}mnkt\_, mnkg\_, mnktm\_)))))
\end{quote}
mentioned by 3 result(s)

\smallskip

\noindent\textbf{\texttt{rel-free(md, n, u)}}\quad\textit{predicate}\index{rel-free(md, n, u)@\texttt{rel-free(md, n, u)}}\\
reads: the \$2 vectors \$3 are linearly independent in \$1 \\
defined by \texttt{rel-free} \\
characterized by \texttt{rel-free}  \textit{(definitional)}:
\begin{quote}\ttfamily\footnotesize\raggedright\noindent
forall(\allowbreak{}[md, n, u], rel-free(\allowbreak{}md, n, u) iff forall(\allowbreak{}[c\_ in mat(\allowbreak{}1, n, carr(\allowbreak{}scal(\allowbreak{}md)))], lincomb(\allowbreak{}md, n, c\_, u) = vzero(\allowbreak{}md) implies forall(\allowbreak{}[j\_ in interval(\allowbreak{}1, n)], entry(\allowbreak{}c\_, 1, j\_) = zero(\allowbreak{}scal(\allowbreak{}md)))))
\end{quote}
also: \texttt{rb4-diag-rel-free-cna-as-module} \texttt{rb4-diag-rel-free-sip-as-module} \texttt{rb4-diag-rel-free-uba-as-module} \texttt{rb4-diag-rel-free-cnvs-as-module} \texttt{rel-free-matact-cancel-cna-as-module} \texttt{rel-free-matact-cancel-sip-as-module} \texttt{rel-free-matact-cancel-uba-as-module} \texttt{rel-free-lincomb-cancel-cna-as-module} ... (32 in all) \\
mentioned by 408 result(s)

\smallskip

\noindent\textbf{\texttt{rel-functoid-beta-1(mrt\_, mrhs\_, mrc\_)}}\quad\textit{predicate}\index{rel-functoid-beta-1(mrt_, mrhs_, mrc_)@\texttt{rel-functoid-beta-1(mrt\_, mrhs\_, mrc\_)}}\\
reads: the kernel rule, at one position, relates \$2 to \$3 \\
defined by \texttt{rel-functoid-beta-1} \\
characterized by \texttt{rel-functoid-beta-1}  \textit{(definitional)}:
\begin{quote}\ttfamily\footnotesize\raggedright\noindent
forall(\allowbreak{}[mrt\_, mrhs\_, mrc\_], rel-functoid-beta-1(\allowbreak{}mrt\_, mrhs\_, mrc\_) iff rel-rewrite-1(\allowbreak{}j-fb(\allowbreak{}mrt\_), mrhs\_, mrc\_))
\end{quote}
mentioned by 3 result(s)

\smallskip

\noindent\textbf{\texttt{rel-functoid-beta-7(mrt\_, mrhs\_, mrc\_)}}\quad\textit{predicate}\index{rel-functoid-beta-7(mrt_, mrhs_, mrc_)@\texttt{rel-functoid-beta-7(mrt\_, mrhs\_, mrc\_)}}\\
reads: the kernel rule, at one position of any stage, relates \$2 to \$3 \\
defined by \texttt{rel-functoid-beta-7} \\
characterized by \texttt{rel-functoid-beta-7}  \textit{(definitional)}:
\begin{quote}\ttfamily\footnotesize\raggedright\noindent
forall(\allowbreak{}[mrt\_, mrhs\_, mrc\_], rel-functoid-beta-7(\allowbreak{}mrt\_, mrhs\_, mrc\_) iff rel-rewrite-7(\allowbreak{}j-fb(\allowbreak{}mrt\_), mrhs\_, mrc\_))
\end{quote}
mentioned by 3 result(s)

\smallskip

\noindent\textbf{\texttt{rel-if-false-k(mnkt\_, mnkhs\_, mnkc\_)}}\quad\textit{predicate}\index{rel-if-false-k(mnkt_, mnkhs_, mnkc_)@\texttt{rel-if-false-k(mnkt\_, mnkhs\_, mnkc\_)}}\\
reads: the kernel's if-false at a certified branch relates \$2 to \$3 \\
defined by \texttt{rel-if-false-k} \\
characterized by \texttt{rel-if-false-k}  \textit{(definitional)}:
\begin{quote}\ttfamily\footnotesize\raggedright\noindent
forall(\allowbreak{}[mnkt\_, mnkhs\_, mnkc\_], rel-if-false-k(\allowbreak{}mnkt\_, mnkhs\_, mnkc\_) iff mnkc\_ in seq-set and forsome(\allowbreak{}[mnkg\_, mnkf\_, mnkp\_, mnka\_, mnkb\_, mnkh\_], mnkc\_ = [mnkg\_, mnkf\_] and mnkh\_ in mnkg\_ and mnkh\_ in cert-set(\allowbreak{}mnkt\_, mnkg\_, mnkb\_) and enc-if(\allowbreak{}mnkp\_, mnka\_, mnkb\_) in clt and mnka\_ in clt and mnkb\_ in clt and mnkhs\_ = [[mnkg\_, enc-not(\allowbreak{}mnkp\_)], [ctx-add(\allowbreak{}mnkg\_, enc-eq(\allowbreak{}enc-if(\allowbreak{}mnkp\_, mnka\_, mnkb\_), mnkb\_)), mnkf\_]]))
\end{quote}
mentioned by 3 result(s)

\smallskip

\noindent\textbf{\texttt{rel-if-true-k(mnkt\_, mnkhs\_, mnkc\_)}}\quad\textit{predicate}\index{rel-if-true-k(mnkt_, mnkhs_, mnkc_)@\texttt{rel-if-true-k(mnkt\_, mnkhs\_, mnkc\_)}}\\
reads: the kernel's if-true at a certified branch relates \$2 to \$3 \\
defined by \texttt{rel-if-true-k} \\
characterized by \texttt{rel-if-true-k}  \textit{(definitional)}:
\begin{quote}\ttfamily\footnotesize\raggedright\noindent
forall(\allowbreak{}[mnkt\_, mnkhs\_, mnkc\_], rel-if-true-k(\allowbreak{}mnkt\_, mnkhs\_, mnkc\_) iff mnkc\_ in seq-set and forsome(\allowbreak{}[mnkg\_, mnkf\_, mnkp\_, mnka\_, mnkb\_, mnkh\_], mnkc\_ = [mnkg\_, mnkf\_] and mnkh\_ in mnkg\_ and mnkh\_ in cert-set(\allowbreak{}mnkt\_, mnkg\_, mnka\_) and enc-if(\allowbreak{}mnkp\_, mnka\_, mnkb\_) in clt and mnka\_ in clt and mnkb\_ in clt and mnkhs\_ = [[mnkg\_, mnkp\_], [ctx-add(\allowbreak{}mnkg\_, enc-eq(\allowbreak{}enc-if(\allowbreak{}mnkp\_, mnka\_, mnkb\_), mnka\_)), mnkf\_]]))
\end{quote}
mentioned by 3 result(s)

\smallskip

\noindent\textbf{\texttt{rel-iff-elim-k(mnkt\_, mnkhs\_, mnkc\_)}}\quad\textit{predicate}\index{rel-iff-elim-k(mnkt_, mnkhs_, mnkc_)@\texttt{rel-iff-elim-k(mnkt\_, mnkhs\_, mnkc\_)}}\\
reads: the kernel's iff-elim relates \$2 to \$3 \\
defined by \texttt{rel-iff-elim-k} \\
characterized by \texttt{rel-iff-elim-k}  \textit{(definitional)}:
\begin{quote}\ttfamily\footnotesize\raggedright\noindent
forall(\allowbreak{}[mnkt\_, mnkhs\_, mnkc\_], rel-iff-elim-k(\allowbreak{}mnkt\_, mnkhs\_, mnkc\_) iff mnkc\_ in seq-set and forsome(\allowbreak{}[mnkg\_, mnkf\_, mnka\_, mnkb\_], mnkc\_ = [mnkg\_, mnkf\_] and enc-iff(\allowbreak{}mnka\_, mnkb\_) in mnkg\_ and mnkhs\_ = [[ctx-add(\allowbreak{}ctx-add(\allowbreak{}difference(\allowbreak{}mnkg\_, pair(\allowbreak{}enc-iff(\allowbreak{}mnka\_, mnkb\_), enc-iff(\allowbreak{}mnka\_, mnkb\_))), enc-imp(\allowbreak{}mnka\_, mnkb\_)), enc-imp(\allowbreak{}mnkb\_, mnka\_)), mnkf\_]]))
\end{quote}
mentioned by 3 result(s)

\smallskip

\noindent\textbf{\texttt{rel-intersection-decompose-1(md14t\_, md14hs\_, md14c\_)}}\quad\textit{predicate}\index{rel-intersection-decompose-1(md14t_, md14hs_, md14c_)@\texttt{rel-intersection-decompose-1(md14t\_, md14hs\_, md14c\_)}}\\
reads: the kernel's rewrite by the catalog's n-ary decomposition, at one position, relates \$2 to \$3 \\
defined by \texttt{rel-intersection-decompose-1} \\
characterized by \texttt{rel-intersection-decompose-1}  \textit{(definitional)}:
\begin{quote}\ttfamily\footnotesize\raggedright\noindent
forall(\allowbreak{}[md14t\_, md14hs\_, md14c\_], rel-intersection-decompose-1(\allowbreak{}md14t\_, md14hs\_, md14c\_) iff rel-rewrite-1(\allowbreak{}j-idec(\allowbreak{}md14t\_), md14hs\_, md14c\_))
\end{quote}
mentioned by 3 result(s)

\smallskip

\noindent\textbf{\texttt{rel-intersection-decompose-7(md14t\_, md14hs\_, md14c\_)}}\quad\textit{predicate}\index{rel-intersection-decompose-7(md14t_, md14hs_, md14c_)@\texttt{rel-intersection-decompose-7(md14t\_, md14hs\_, md14c\_)}}\\
reads: the kernel's rewrite by the catalog's n-ary decomposition, at one position, relates \$2 to \$3 \\
defined by \texttt{rel-intersection-decompose-7} \\
characterized by \texttt{rel-intersection-decompose-7}  \textit{(definitional)}:
\begin{quote}\ttfamily\footnotesize\raggedright\noindent
forall(\allowbreak{}[md14t\_, md14hs\_, md14c\_], rel-intersection-decompose-7(\allowbreak{}md14t\_, md14hs\_, md14c\_) iff rel-rewrite-7(\allowbreak{}j-idec(\allowbreak{}md14t\_), md14hs\_, md14c\_))
\end{quote}
mentioned by 3 result(s)

\smallskip

\noindent\textbf{\texttt{rel-intersection-decompose-hyp-1(md14t\_, md14hs\_, md14c\_)}}\quad\textit{predicate}\index{rel-intersection-decompose-hyp-1(md14t_, md14hs_, md14c_)@\texttt{rel-intersection-decompose-hyp-1(md14t\_, md14hs\_, md14c\_)}}\\
reads: the kernel's rewrite by the catalog's n-ary decomposition, at one position, relates \$2 to \$3 \\
defined by \texttt{rel-intersection-decompose-hyp-1} \\
characterized by \texttt{rel-intersection-decompose-hyp-1}  \textit{(definitional)}:
\begin{quote}\ttfamily\footnotesize\raggedright\noindent
forall(\allowbreak{}[md14t\_, md14hs\_, md14c\_], rel-intersection-decompose-hyp-1(\allowbreak{}md14t\_, md14hs\_, md14c\_) iff rel-rewrite-hyp-1(\allowbreak{}j-idec(\allowbreak{}md14t\_), md14hs\_, md14c\_))
\end{quote}
mentioned by 3 result(s)

\smallskip

\noindent\textbf{\texttt{rel-intersection-decompose-hyp-7(md14t\_, md14hs\_, md14c\_)}}\quad\textit{predicate}\index{rel-intersection-decompose-hyp-7(md14t_, md14hs_, md14c_)@\texttt{rel-intersection-decompose-hyp-7(md14t\_, md14hs\_, md14c\_)}}\\
reads: the kernel's rewrite by the catalog's n-ary decomposition, at one position, relates \$2 to \$3 \\
defined by \texttt{rel-intersection-decompose-hyp-7} \\
characterized by \texttt{rel-intersection-decompose-hyp-7}  \textit{(definitional)}:
\begin{quote}\ttfamily\footnotesize\raggedright\noindent
forall(\allowbreak{}[md14t\_, md14hs\_, md14c\_], rel-intersection-decompose-hyp-7(\allowbreak{}md14t\_, md14hs\_, md14c\_) iff rel-rewrite-hyp-7(\allowbreak{}j-idec(\allowbreak{}md14t\_), md14hs\_, md14c\_))
\end{quote}
mentioned by 3 result(s)

\smallskip

\noindent\textbf{\texttt{rel-intersection-elim(mct\_, mchs\_, mcc\_)}}\quad\textit{predicate}\index{rel-intersection-elim(mct_, mchs_, mcc_)@\texttt{rel-intersection-elim(mct\_, mchs\_, mcc\_)}}\\
reads: the kernel's intersection-elim relates \$2 to \$3 \\
defined by \texttt{rel-intersection-elim} \\
characterized by \texttt{rel-intersection-elim}  \textit{(definitional)}:
\begin{quote}\ttfamily\footnotesize\raggedright\noindent
forall(\allowbreak{}[mct\_, mchs\_, mcc\_], rel-intersection-elim(\allowbreak{}mct\_, mchs\_, mcc\_) iff mcc\_ in seq-set and forsome(\allowbreak{}[mcg\_, mcf\_, mcx\_, mcl\_, mck\_], mcc\_ = [mcg\_, mcf\_] and enc-in(\allowbreak{}mcx\_, enc-inter(\allowbreak{}mcl\_)) in mcg\_ and mck\_ in nn and 1 <= mck\_ and mck\_ <= length(\allowbreak{}mcl\_) and mchs\_ = [[ctx-add(\allowbreak{}mcg\_, enc-in(\allowbreak{}mcx\_, nth(\allowbreak{}mck\_, mcl\_))), mcf\_]]))
\end{quote}
mentioned by 3 result(s)

\smallskip

\noindent\textbf{\texttt{rel-intersection-intro(mct\_, mchs\_, mcc\_)}}\quad\textit{predicate}\index{rel-intersection-intro(mct_, mchs_, mcc_)@\texttt{rel-intersection-intro(mct\_, mchs\_, mcc\_)}}\\
reads: the kernel's intersection-intro relates \$2 to \$3 \\
defined by \texttt{rel-intersection-intro} \\
characterized by \texttt{rel-intersection-intro}  \textit{(definitional)}:
\begin{quote}\ttfamily\footnotesize\raggedright\noindent
forall(\allowbreak{}[mct\_, mchs\_, mcc\_], rel-intersection-intro(\allowbreak{}mct\_, mchs\_, mcc\_) iff mcc\_ in seq-set and forsome(\allowbreak{}[mcg\_, mcx\_, mcl\_], mcc\_ = [mcg\_, enc-in(\allowbreak{}mcx\_, enc-inter(\allowbreak{}mcl\_))] and 1 <= length(\allowbreak{}mcl\_) and length(\allowbreak{}mchs\_) = length(\allowbreak{}mcl\_) and forall(\allowbreak{}[mci\_ in nn], 1 <= mci\_ implies mci\_ <= length(\allowbreak{}mcl\_) implies nth(\allowbreak{}mci\_, mchs\_) = [mcg\_, enc-in(\allowbreak{}mcx\_, nth(\allowbreak{}mci\_, mcl\_))])))
\end{quote}
mentioned by 3 result(s)

\smallskip

\noindent\textbf{\texttt{rel-iota-in-elim-k(mnkt\_, mnkhs\_, mnkc\_)}}\quad\textit{predicate}\index{rel-iota-in-elim-k(mnkt_, mnkhs_, mnkc_)@\texttt{rel-iota-in-elim-k(mnkt\_, mnkhs\_, mnkc\_)}}\\
reads: the kernel's iota-in-elim relates \$2 to \$3 \\
defined by \texttt{rel-iota-in-elim-k} \\
characterized by \texttt{rel-iota-in-elim-k}  \textit{(definitional)}:
\begin{quote}\ttfamily\footnotesize\raggedright\noindent
forall(\allowbreak{}[mnkt\_, mnkhs\_, mnkc\_], rel-iota-in-elim-k(\allowbreak{}mnkt\_, mnkhs\_, mnkc\_) iff mnkc\_ in seq-set and forsome(\allowbreak{}[mnkg\_, mnkf\_, mnkp\_], mnkc\_ = [mnkg\_, mnkf\_] and (\allowbreak{}forsome(\allowbreak{}[mnka\_], enc-in(\allowbreak{}enc-iota(\allowbreak{}mnkp\_), mnka\_) in mnkg\_) or forsome(\allowbreak{}[mnku\_], enc-eq(\allowbreak{}enc-iota(\allowbreak{}mnkp\_), mnku\_) in mnkg\_ or enc-eq(\allowbreak{}mnku\_, enc-iota(\allowbreak{}mnkp\_)) in mnkg\_)) and mnkhs\_ = [[ctx-add(\allowbreak{}mnkg\_, expr-open(\allowbreak{}mnkp\_, 0, enc-iota(\allowbreak{}mnkp\_))), mnkf\_]]))
\end{quote}
mentioned by 3 result(s)

\smallskip

\noindent\textbf{\texttt{rel-lambda-beta-1(mrt\_, mrhs\_, mrc\_)}}\quad\textit{predicate}\index{rel-lambda-beta-1(mrt_, mrhs_, mrc_)@\texttt{rel-lambda-beta-1(mrt\_, mrhs\_, mrc\_)}}\\
reads: the kernel rule, at one position, relates \$2 to \$3 \\
defined by \texttt{rel-lambda-beta-1} \\
characterized by \texttt{rel-lambda-beta-1}  \textit{(definitional)}:
\begin{quote}\ttfamily\footnotesize\raggedright\noindent
forall(\allowbreak{}[mrt\_, mrhs\_, mrc\_], rel-lambda-beta-1(\allowbreak{}mrt\_, mrhs\_, mrc\_) iff rel-rewrite-1(\allowbreak{}j-lb(\allowbreak{}mrt\_), mrhs\_, mrc\_))
\end{quote}
mentioned by 3 result(s)

\smallskip

\noindent\textbf{\texttt{rel-lambda-beta-7(mrt\_, mrhs\_, mrc\_)}}\quad\textit{predicate}\index{rel-lambda-beta-7(mrt_, mrhs_, mrc_)@\texttt{rel-lambda-beta-7(mrt\_, mrhs\_, mrc\_)}}\\
reads: the kernel rule, at one position of any stage, relates \$2 to \$3 \\
defined by \texttt{rel-lambda-beta-7} \\
characterized by \texttt{rel-lambda-beta-7}  \textit{(definitional)}:
\begin{quote}\ttfamily\footnotesize\raggedright\noindent
forall(\allowbreak{}[mrt\_, mrhs\_, mrc\_], rel-lambda-beta-7(\allowbreak{}mrt\_, mrhs\_, mrc\_) iff rel-rewrite-7(\allowbreak{}j-lb(\allowbreak{}mrt\_), mrhs\_, mrc\_))
\end{quote}
mentioned by 3 result(s)

\smallskip

\noindent\textbf{\texttt{rel-lambda-beta-hyp-1(mrt\_, mrhs\_, mrc\_)}}\quad\textit{predicate}\index{rel-lambda-beta-hyp-1(mrt_, mrhs_, mrc_)@\texttt{rel-lambda-beta-hyp-1(mrt\_, mrhs\_, mrc\_)}}\\
reads: the kernel rule, at one position, relates \$2 to \$3 \\
defined by \texttt{rel-lambda-beta-hyp-1} \\
characterized by \texttt{rel-lambda-beta-hyp-1}  \textit{(definitional)}:
\begin{quote}\ttfamily\footnotesize\raggedright\noindent
forall(\allowbreak{}[mrt\_, mrhs\_, mrc\_], rel-lambda-beta-hyp-1(\allowbreak{}mrt\_, mrhs\_, mrc\_) iff rel-rewrite-hyp-1(\allowbreak{}j-lb(\allowbreak{}mrt\_), mrhs\_, mrc\_))
\end{quote}
mentioned by 3 result(s)

\smallskip

\noindent\textbf{\texttt{rel-lambda-beta-hyp-7(mrt\_, mrhs\_, mrc\_)}}\quad\textit{predicate}\index{rel-lambda-beta-hyp-7(mrt_, mrhs_, mrc_)@\texttt{rel-lambda-beta-hyp-7(mrt\_, mrhs\_, mrc\_)}}\\
reads: the kernel rule, at one position of any stage, relates \$2 to \$3 \\
defined by \texttt{rel-lambda-beta-hyp-7} \\
characterized by \texttt{rel-lambda-beta-hyp-7}  \textit{(definitional)}:
\begin{quote}\ttfamily\footnotesize\raggedright\noindent
forall(\allowbreak{}[mrt\_, mrhs\_, mrc\_], rel-lambda-beta-hyp-7(\allowbreak{}mrt\_, mrhs\_, mrc\_) iff rel-rewrite-hyp-7(\allowbreak{}j-lb(\allowbreak{}mrt\_), mrhs\_, mrc\_))
\end{quote}
mentioned by 3 result(s)

\smallskip

\noindent\textbf{\texttt{rel-length-reduce-1(mrt\_, mrhs\_, mrc\_)}}\quad\textit{predicate}\index{rel-length-reduce-1(mrt_, mrhs_, mrc_)@\texttt{rel-length-reduce-1(mrt\_, mrhs\_, mrc\_)}}\\
reads: the kernel rule, at one position, relates \$2 to \$3 \\
defined by \texttt{rel-length-reduce-1} \\
characterized by \texttt{rel-length-reduce-1}  \textit{(definitional)}:
\begin{quote}\ttfamily\footnotesize\raggedright\noindent
forall(\allowbreak{}[mrt\_, mrhs\_, mrc\_], rel-length-reduce-1(\allowbreak{}mrt\_, mrhs\_, mrc\_) iff rel-rewrite-1(\allowbreak{}j-len(\allowbreak{}mrt\_), mrhs\_, mrc\_))
\end{quote}
mentioned by 3 result(s)

\smallskip

\noindent\textbf{\texttt{rel-length-reduce-7(mrt\_, mrhs\_, mrc\_)}}\quad\textit{predicate}\index{rel-length-reduce-7(mrt_, mrhs_, mrc_)@\texttt{rel-length-reduce-7(mrt\_, mrhs\_, mrc\_)}}\\
reads: the kernel rule, at one position of any stage, relates \$2 to \$3 \\
defined by \texttt{rel-length-reduce-7} \\
characterized by \texttt{rel-length-reduce-7}  \textit{(definitional)}:
\begin{quote}\ttfamily\footnotesize\raggedright\noindent
forall(\allowbreak{}[mrt\_, mrhs\_, mrc\_], rel-length-reduce-7(\allowbreak{}mrt\_, mrhs\_, mrc\_) iff rel-rewrite-7(\allowbreak{}j-len(\allowbreak{}mrt\_), mrhs\_, mrc\_))
\end{quote}
mentioned by 3 result(s)

\smallskip

\noindent\textbf{\texttt{rel-macete-1(mrt\_, mrhs\_, mrc\_)}}\quad\textit{predicate}\index{rel-macete-1(mrt_, mrhs_, mrc_)@\texttt{rel-macete-1(mrt\_, mrhs\_, mrc\_)}}\\
reads: the kernel rule, at one position, relates \$2 to \$3 \\
defined by \texttt{rel-macete-1} \\
characterized by \texttt{rel-macete-1}  \textit{(definitional)}:
\begin{quote}\ttfamily\footnotesize\raggedright\noindent
forall(\allowbreak{}[mrt\_, mrhs\_, mrc\_], rel-macete-1(\allowbreak{}mrt\_, mrhs\_, mrc\_) iff rel-rewrite-1(\allowbreak{}j-mac(\allowbreak{}mrt\_), mrhs\_, mrc\_))
\end{quote}
mentioned by 3 result(s)

\smallskip

\noindent\textbf{\texttt{rel-macete-7(mrt\_, mrhs\_, mrc\_)}}\quad\textit{predicate}\index{rel-macete-7(mrt_, mrhs_, mrc_)@\texttt{rel-macete-7(mrt\_, mrhs\_, mrc\_)}}\\
reads: the kernel rule, at one position of any stage, relates \$2 to \$3 \\
defined by \texttt{rel-macete-7} \\
characterized by \texttt{rel-macete-7}  \textit{(definitional)}:
\begin{quote}\ttfamily\footnotesize\raggedright\noindent
forall(\allowbreak{}[mrt\_, mrhs\_, mrc\_], rel-macete-7(\allowbreak{}mrt\_, mrhs\_, mrc\_) iff rel-rewrite-7(\allowbreak{}j-mac(\allowbreak{}mrt\_), mrhs\_, mrc\_))
\end{quote}
mentioned by 3 result(s)

\smallskip

\noindent\textbf{\texttt{rel-macete-hyp-1(mrt\_, mrhs\_, mrc\_)}}\quad\textit{predicate}\index{rel-macete-hyp-1(mrt_, mrhs_, mrc_)@\texttt{rel-macete-hyp-1(mrt\_, mrhs\_, mrc\_)}}\\
reads: the kernel rule, at one position, relates \$2 to \$3 \\
defined by \texttt{rel-macete-hyp-1} \\
characterized by \texttt{rel-macete-hyp-1}  \textit{(definitional)}:
\begin{quote}\ttfamily\footnotesize\raggedright\noindent
forall(\allowbreak{}[mrt\_, mrhs\_, mrc\_], rel-macete-hyp-1(\allowbreak{}mrt\_, mrhs\_, mrc\_) iff rel-rewrite-hyp-1(\allowbreak{}j-mac(\allowbreak{}mrt\_), mrhs\_, mrc\_))
\end{quote}
mentioned by 3 result(s)

\smallskip

\noindent\textbf{\texttt{rel-macete-hyp-7(mrt\_, mrhs\_, mrc\_)}}\quad\textit{predicate}\index{rel-macete-hyp-7(mrt_, mrhs_, mrc_)@\texttt{rel-macete-hyp-7(mrt\_, mrhs\_, mrc\_)}}\\
reads: the kernel rule, at one position of any stage, relates \$2 to \$3 \\
defined by \texttt{rel-macete-hyp-7} \\
characterized by \texttt{rel-macete-hyp-7}  \textit{(definitional)}:
\begin{quote}\ttfamily\footnotesize\raggedright\noindent
forall(\allowbreak{}[mrt\_, mrhs\_, mrc\_], rel-macete-hyp-7(\allowbreak{}mrt\_, mrhs\_, mrc\_) iff rel-rewrite-hyp-7(\allowbreak{}j-mac(\allowbreak{}mrt\_), mrhs\_, mrc\_))
\end{quote}
mentioned by 3 result(s)

\smallskip

\noindent\textbf{\texttt{rel-macete-n-1(md13t\_, md13hs\_, md13c\_)}}\quad\textit{predicate}\index{rel-macete-n-1(md13t_, md13hs_, md13c_)@\texttt{rel-macete-n-1(md13t\_, md13hs\_, md13c\_)}}\\
reads: the kernel's macete of any number of schema variables, at one position, relates \$2 to \$3 \\
defined by \texttt{rel-macete-n-1} \\
characterized by \texttt{rel-macete-n-1}  \textit{(definitional)}:
\begin{quote}\ttfamily\footnotesize\raggedright\noindent
forall(\allowbreak{}[md13t\_, md13hs\_, md13c\_], rel-macete-n-1(\allowbreak{}md13t\_, md13hs\_, md13c\_) iff rel-rewrite-1(\allowbreak{}j-macn(\allowbreak{}md13t\_), md13hs\_, md13c\_))
\end{quote}
mentioned by 3 result(s)

\smallskip

\noindent\textbf{\texttt{rel-macete-n-7(md13t\_, md13hs\_, md13c\_)}}\quad\textit{predicate}\index{rel-macete-n-7(md13t_, md13hs_, md13c_)@\texttt{rel-macete-n-7(md13t\_, md13hs\_, md13c\_)}}\\
reads: the kernel's macete of any number of schema variables, at one position, relates \$2 to \$3 \\
defined by \texttt{rel-macete-n-7} \\
characterized by \texttt{rel-macete-n-7}  \textit{(definitional)}:
\begin{quote}\ttfamily\footnotesize\raggedright\noindent
forall(\allowbreak{}[md13t\_, md13hs\_, md13c\_], rel-macete-n-7(\allowbreak{}md13t\_, md13hs\_, md13c\_) iff rel-rewrite-7(\allowbreak{}j-macn(\allowbreak{}md13t\_), md13hs\_, md13c\_))
\end{quote}
mentioned by 3 result(s)

\smallskip

\noindent\textbf{\texttt{rel-macete-n-hyp-1(md13t\_, md13hs\_, md13c\_)}}\quad\textit{predicate}\index{rel-macete-n-hyp-1(md13t_, md13hs_, md13c_)@\texttt{rel-macete-n-hyp-1(md13t\_, md13hs\_, md13c\_)}}\\
reads: the kernel's macete of any number of schema variables, at one position, relates \$2 to \$3 \\
defined by \texttt{rel-macete-n-hyp-1} \\
characterized by \texttt{rel-macete-n-hyp-1}  \textit{(definitional)}:
\begin{quote}\ttfamily\footnotesize\raggedright\noindent
forall(\allowbreak{}[md13t\_, md13hs\_, md13c\_], rel-macete-n-hyp-1(\allowbreak{}md13t\_, md13hs\_, md13c\_) iff rel-rewrite-hyp-1(\allowbreak{}j-macn(\allowbreak{}md13t\_), md13hs\_, md13c\_))
\end{quote}
mentioned by 3 result(s)

\smallskip

\noindent\textbf{\texttt{rel-macete-n-hyp-7(md13t\_, md13hs\_, md13c\_)}}\quad\textit{predicate}\index{rel-macete-n-hyp-7(md13t_, md13hs_, md13c_)@\texttt{rel-macete-n-hyp-7(md13t\_, md13hs\_, md13c\_)}}\\
reads: the kernel's macete of any number of schema variables, at one position, relates \$2 to \$3 \\
defined by \texttt{rel-macete-n-hyp-7} \\
characterized by \texttt{rel-macete-n-hyp-7}  \textit{(definitional)}:
\begin{quote}\ttfamily\footnotesize\raggedright\noindent
forall(\allowbreak{}[md13t\_, md13hs\_, md13c\_], rel-macete-n-hyp-7(\allowbreak{}md13t\_, md13hs\_, md13c\_) iff rel-rewrite-hyp-7(\allowbreak{}j-macn(\allowbreak{}md13t\_), md13hs\_, md13c\_))
\end{quote}
mentioned by 3 result(s)

\smallskip

\noindent\textbf{\texttt{rel-nn-induction(mit\_, mihs\_, mic\_)}}\quad\textit{predicate}\index{rel-nn-induction(mit_, mihs_, mic_)@\texttt{rel-nn-induction(mit\_, mihs\_, mic\_)}}\\
reads: the kernel's nn-induction relates \$2 to \$3 \\
defined by \texttt{rel-nn-induction} \\
characterized by \texttt{rel-nn-induction}  \textit{(definitional)}:
\begin{quote}\ttfamily\footnotesize\raggedright\noindent
forall(\allowbreak{}[mit\_, mihs\_, mic\_], rel-nn-induction(\allowbreak{}mit\_, mihs\_, mic\_) iff mic\_ in seq-set and forsome(\allowbreak{}[mig\_, mip\_, mis\_, mifs\_], mic\_ = [mig\_, enc-all(\allowbreak{}enc-imp(\allowbreak{}enc-in(\allowbreak{}ebvar(\allowbreak{}0), enc-nnc), mip\_))] and mihs\_ = [[mig\_, expr-open(\allowbreak{}mip\_, 0, enc-num(\allowbreak{}0))], [mig\_, enc-all(\allowbreak{}mis\_)]] and mifs\_ in set and mifs\_ subset nn and card(\allowbreak{}mifs\_) in nn and forall(\allowbreak{}[miy\_ in nn], not(\allowbreak{}miy\_ in mifs\_) implies expr-open(\allowbreak{}mis\_, 0, efvar(\allowbreak{}miy\_)) = enc-imp(\allowbreak{}enc-in(\allowbreak{}efvar(\allowbreak{}miy\_), enc-nnc), enc-imp(\allowbreak{}expr-open(\allowbreak{}mip\_, 0, efvar(\allowbreak{}miy\_)), expr-open(\allowbreak{}mip\_, 0, enc-succ(\allowbreak{}efvar(\allowbreak{}miy\_))))))))
\end{quote}
mentioned by 2 result(s)

\smallskip

\noindent\textbf{\texttt{rel-nn-induction-nbg(mit\_, mihs\_, mic\_)}}\quad\textit{predicate}\index{rel-nn-induction-nbg(mit_, mihs_, mic_)@\texttt{rel-nn-induction-nbg(mit\_, mihs\_, mic\_)}}\\
reads: the kernel's nn-induction at a predicative body relates \$2 to \$3 \\
defined by \texttt{rel-nn-induction-nbg} \\
characterized by \texttt{rel-nn-induction-nbg}  \textit{(definitional)}:
\begin{quote}\ttfamily\footnotesize\raggedright\noindent
forall(\allowbreak{}[mit\_, mihs\_, mic\_], rel-nn-induction-nbg(\allowbreak{}mit\_, mihs\_, mic\_) iff mic\_ in seq-set and forsome(\allowbreak{}[mig\_, mip\_, mis\_, mifs\_], mic\_ = [mig\_, enc-all(\allowbreak{}enc-imp(\allowbreak{}enc-in(\allowbreak{}ebvar(\allowbreak{}0), enc-nnc), mip\_))] and nbg-ok(\allowbreak{}mip\_) and mihs\_ = [[mig\_, expr-open(\allowbreak{}mip\_, 0, enc-num(\allowbreak{}0))], [mig\_, enc-all(\allowbreak{}mis\_)]] and mifs\_ in set and mifs\_ subset nn and card(\allowbreak{}mifs\_) in nn and forall(\allowbreak{}[miy\_ in nn], not(\allowbreak{}miy\_ in mifs\_) implies expr-open(\allowbreak{}mis\_, 0, efvar(\allowbreak{}miy\_)) = enc-imp(\allowbreak{}enc-in(\allowbreak{}efvar(\allowbreak{}miy\_), enc-nnc), enc-imp(\allowbreak{}expr-open(\allowbreak{}mip\_, 0, efvar(\allowbreak{}miy\_)), expr-open(\allowbreak{}mip\_, 0, enc-succ(\allowbreak{}efvar(\allowbreak{}miy\_))))))))
\end{quote}
mentioned by 4 result(s)

\smallskip

\noindent\textbf{\texttt{rel-nth-reduce-1(mrt\_, mrhs\_, mrc\_)}}\quad\textit{predicate}\index{rel-nth-reduce-1(mrt_, mrhs_, mrc_)@\texttt{rel-nth-reduce-1(mrt\_, mrhs\_, mrc\_)}}\\
reads: the kernel rule, at one position, relates \$2 to \$3 \\
defined by \texttt{rel-nth-reduce-1} \\
characterized by \texttt{rel-nth-reduce-1}  \textit{(definitional)}:
\begin{quote}\ttfamily\footnotesize\raggedright\noindent
forall(\allowbreak{}[mrt\_, mrhs\_, mrc\_], rel-nth-reduce-1(\allowbreak{}mrt\_, mrhs\_, mrc\_) iff rel-rewrite-1(\allowbreak{}j-nth(\allowbreak{}mrt\_), mrhs\_, mrc\_))
\end{quote}
mentioned by 3 result(s)

\smallskip

\noindent\textbf{\texttt{rel-nth-reduce-7(mrt\_, mrhs\_, mrc\_)}}\quad\textit{predicate}\index{rel-nth-reduce-7(mrt_, mrhs_, mrc_)@\texttt{rel-nth-reduce-7(mrt\_, mrhs\_, mrc\_)}}\\
reads: the kernel rule, at one position of any stage, relates \$2 to \$3 \\
defined by \texttt{rel-nth-reduce-7} \\
characterized by \texttt{rel-nth-reduce-7}  \textit{(definitional)}:
\begin{quote}\ttfamily\footnotesize\raggedright\noindent
forall(\allowbreak{}[mrt\_, mrhs\_, mrc\_], rel-nth-reduce-7(\allowbreak{}mrt\_, mrhs\_, mrc\_) iff rel-rewrite-7(\allowbreak{}j-nth(\allowbreak{}mrt\_), mrhs\_, mrc\_))
\end{quote}
mentioned by 3 result(s)

\smallskip

\noindent\textbf{\texttt{rel-or-elim-k(mnkt\_, mnkhs\_, mnkc\_)}}\quad\textit{predicate}\index{rel-or-elim-k(mnkt_, mnkhs_, mnkc_)@\texttt{rel-or-elim-k(mnkt\_, mnkhs\_, mnkc\_)}}\\
reads: the kernel's or-elim relates \$2 to \$3 \\
defined by \texttt{rel-or-elim-k} \\
characterized by \texttt{rel-or-elim-k}  \textit{(definitional)}:
\begin{quote}\ttfamily\footnotesize\raggedright\noindent
forall(\allowbreak{}[mnkt\_, mnkhs\_, mnkc\_], rel-or-elim-k(\allowbreak{}mnkt\_, mnkhs\_, mnkc\_) iff mnkc\_ in seq-set and forsome(\allowbreak{}[mnkg\_, mnkf\_, mnka\_, mnkb\_], mnkc\_ = [mnkg\_, mnkf\_] and enc-or(\allowbreak{}mnka\_, mnkb\_) in mnkg\_ and mnkhs\_ = [[ctx-add(\allowbreak{}difference(\allowbreak{}mnkg\_, pair(\allowbreak{}enc-or(\allowbreak{}mnka\_, mnkb\_), enc-or(\allowbreak{}mnka\_, mnkb\_))), mnka\_), mnkf\_], [ctx-add(\allowbreak{}difference(\allowbreak{}mnkg\_, pair(\allowbreak{}enc-or(\allowbreak{}mnka\_, mnkb\_), enc-or(\allowbreak{}mnka\_, mnkb\_))), mnkb\_), mnkf\_]]))
\end{quote}
mentioned by 3 result(s)

\smallskip

\noindent\textbf{\texttt{rel-reflexivity-k(mnkt\_, mnkhs\_, mnkc\_)}}\quad\textit{predicate}\index{rel-reflexivity-k(mnkt_, mnkhs_, mnkc_)@\texttt{rel-reflexivity-k(mnkt\_, mnkhs\_, mnkc\_)}}\\
reads: the kernel's reflexivity at a certified term relates \$2 to \$3 \\
defined by \texttt{rel-reflexivity-k} \\
characterized by \texttt{rel-reflexivity-k}  \textit{(definitional)}:
\begin{quote}\ttfamily\footnotesize\raggedright\noindent
forall(\allowbreak{}[mnkt\_, mnkhs\_, mnkc\_], rel-reflexivity-k(\allowbreak{}mnkt\_, mnkhs\_, mnkc\_) iff mnkc\_ in seq-set and forsome(\allowbreak{}[mnkg\_, mnktm\_, mnkh\_], mnkc\_ = [mnkg\_, enc-eq(\allowbreak{}mnktm\_, mnktm\_)] and mnktm\_ in clt and mnkh\_ in mnkg\_ and mnkh\_ in cert-set(\allowbreak{}mnkt\_, mnkg\_, mnktm\_) and mnkhs\_ = []))
\end{quote}
mentioned by 3 result(s)

\smallskip

\noindent\textbf{\texttt{rel-rewrite-1(mrt\_, mrhs\_, mrc\_)}}\quad\textit{predicate}\index{rel-rewrite-1(mrt_, mrhs_, mrc_)@\texttt{rel-rewrite-1(mrt\_, mrhs\_, mrc\_)}}\\
reads: \$2 is \$3 with its goal rewritten at one position justified by \$1 \\
defined by \texttt{rel-rewrite-1} \\
characterized by \texttt{rel-rewrite-1}  \textit{(definitional)}:
\begin{quote}\ttfamily\footnotesize\raggedright\noindent
forall(\allowbreak{}[mrt\_, mrhs\_, mrc\_], rel-rewrite-1(\allowbreak{}mrt\_, mrhs\_, mrc\_) iff mrc\_ in seq-set and forsome(\allowbreak{}[mrg\_, mra\_, mrb\_], mrc\_ = [mrg\_, mra\_] and mrhs\_ = [[mrg\_, mrb\_]] and [2, mrg\_, mra\_, mrb\_] in rewr(\allowbreak{}mrt\_, 5)))
\end{quote}
mentioned by 23 result(s)

\smallskip

\noindent\textbf{\texttt{rel-rewrite-7(mrt\_, mrhs\_, mrc\_)}}\quad\textit{predicate}\index{rel-rewrite-7(mrt_, mrhs_, mrc_)@\texttt{rel-rewrite-7(mrt\_, mrhs\_, mrc\_)}}\\
reads: \$2 is \$3 with its goal rewritten at one position of any stage justified by \$1 \\
defined by \texttt{rel-rewrite-7} \\
characterized by \texttt{rel-rewrite-7}  \textit{(definitional)}:
\begin{quote}\ttfamily\footnotesize\raggedright\noindent
forall(\allowbreak{}[mrt\_, mrhs\_, mrc\_], rel-rewrite-7(\allowbreak{}mrt\_, mrhs\_, mrc\_) iff mrc\_ in seq-set and forsome(\allowbreak{}[mrg\_, mra\_, mrb\_], mrc\_ = [mrg\_, mra\_] and mrhs\_ = [[mrg\_, mrb\_]] and [2, mrg\_, mra\_, mrb\_] in rewr(\allowbreak{}mrt\_, 7)))
\end{quote}
mentioned by 23 result(s)

\smallskip

\noindent\textbf{\texttt{rel-rewrite-hyp-1(mrt\_, mrhs\_, mrc\_)}}\quad\textit{predicate}\index{rel-rewrite-hyp-1(mrt_, mrhs_, mrc_)@\texttt{rel-rewrite-hyp-1(mrt\_, mrhs\_, mrc\_)}}\\
reads: \$2 is \$3 with one assumption rewritten at one position justified by \$1 \\
defined by \texttt{rel-rewrite-hyp-1} \\
characterized by \texttt{rel-rewrite-hyp-1}  \textit{(definitional)}:
\begin{quote}\ttfamily\footnotesize\raggedright\noindent
forall(\allowbreak{}[mrt\_, mrhs\_, mrc\_], rel-rewrite-hyp-1(\allowbreak{}mrt\_, mrhs\_, mrc\_) iff mrc\_ in seq-set and forsome(\allowbreak{}[mrg\_, mrf\_, mra\_, mrb\_], mrc\_ = [mrg\_, mrf\_] and mra\_ in mrg\_ and mrhs\_ = [[ctx-add(\allowbreak{}difference(\allowbreak{}mrg\_, pair(\allowbreak{}mra\_, mra\_)), mrb\_), mrf\_]] and [2, mrg\_, mra\_, mrb\_] in rewr(\allowbreak{}mrt\_, 5)))
\end{quote}
mentioned by 13 result(s)

\smallskip

\noindent\textbf{\texttt{rel-rewrite-hyp-7(mrt\_, mrhs\_, mrc\_)}}\quad\textit{predicate}\index{rel-rewrite-hyp-7(mrt_, mrhs_, mrc_)@\texttt{rel-rewrite-hyp-7(mrt\_, mrhs\_, mrc\_)}}\\
reads: \$2 is \$3 with one assumption rewritten at one position of any stage justified by \$1 \\
defined by \texttt{rel-rewrite-hyp-7} \\
characterized by \texttt{rel-rewrite-hyp-7}  \textit{(definitional)}:
\begin{quote}\ttfamily\footnotesize\raggedright\noindent
forall(\allowbreak{}[mrt\_, mrhs\_, mrc\_], rel-rewrite-hyp-7(\allowbreak{}mrt\_, mrhs\_, mrc\_) iff mrc\_ in seq-set and forsome(\allowbreak{}[mrg\_, mrf\_, mra\_, mrb\_], mrc\_ = [mrg\_, mrf\_] and mra\_ in mrg\_ and mrhs\_ = [[ctx-add(\allowbreak{}difference(\allowbreak{}mrg\_, pair(\allowbreak{}mra\_, mra\_)), mrb\_), mrf\_]] and [2, mrg\_, mra\_, mrb\_] in rewr(\allowbreak{}mrt\_, 7)))
\end{quote}
mentioned by 13 result(s)

\smallskip

\noindent\textbf{\texttt{rel-sep-mem-elim(mit\_, mihs\_, mic\_)}}\quad\textit{predicate}\index{rel-sep-mem-elim(mit_, mihs_, mic_)@\texttt{rel-sep-mem-elim(mit\_, mihs\_, mic\_)}}\\
reads: the kernel's sep-mem-elim relates \$2 to \$3 \\
defined by \texttt{rel-sep-mem-elim} \\
characterized by \texttt{rel-sep-mem-elim}  \textit{(definitional)}:
\begin{quote}\ttfamily\footnotesize\raggedright\noindent
forall(\allowbreak{}[mit\_, mihs\_, mic\_], rel-sep-mem-elim(\allowbreak{}mit\_, mihs\_, mic\_) iff mic\_ in seq-set and forsome(\allowbreak{}[mig\_, mif\_, mia\_, mip\_, miy\_], mic\_ = [mig\_, mif\_] and enc-in(\allowbreak{}miy\_, enc-sep(\allowbreak{}mia\_, mip\_)) in mig\_ and nbg-ok(\allowbreak{}mip\_) and mihs\_ = [[ctx-add(\allowbreak{}ctx-add(\allowbreak{}difference(\allowbreak{}mig\_, pair(\allowbreak{}enc-in(\allowbreak{}miy\_, enc-sep(\allowbreak{}mia\_, mip\_)), enc-in(\allowbreak{}miy\_, enc-sep(\allowbreak{}mia\_, mip\_)))), enc-in(\allowbreak{}miy\_, mia\_)), expr-open(\allowbreak{}mip\_, 0, miy\_)), mif\_]]))
\end{quote}
mentioned by 3 result(s)

\smallskip

\noindent\textbf{\texttt{rel-sep-mem-intro(mit\_, mihs\_, mic\_)}}\quad\textit{predicate}\index{rel-sep-mem-intro(mit_, mihs_, mic_)@\texttt{rel-sep-mem-intro(mit\_, mihs\_, mic\_)}}\\
reads: the kernel's sep-mem-intro relates \$2 to \$3 \\
defined by \texttt{rel-sep-mem-intro} \\
characterized by \texttt{rel-sep-mem-intro}  \textit{(definitional)}:
\begin{quote}\ttfamily\footnotesize\raggedright\noindent
forall(\allowbreak{}[mit\_, mihs\_, mic\_], rel-sep-mem-intro(\allowbreak{}mit\_, mihs\_, mic\_) iff mic\_ in seq-set and forsome(\allowbreak{}[mig\_, mia\_, mip\_, miy\_], mic\_ = [mig\_, enc-in(\allowbreak{}miy\_, enc-sep(\allowbreak{}mia\_, mip\_))] and nbg-ok(\allowbreak{}mip\_) and mihs\_ = [[mig\_, enc-in(\allowbreak{}miy\_, mia\_)], [mig\_, expr-open(\allowbreak{}mip\_, 0, miy\_)]]))
\end{quote}
mentioned by 3 result(s)

\smallskip

\noindent\textbf{\texttt{rel-sep-sethood(mit\_, mihs\_, mic\_)}}\quad\textit{predicate}\index{rel-sep-sethood(mit_, mihs_, mic_)@\texttt{rel-sep-sethood(mit\_, mihs\_, mic\_)}}\\
reads: the kernel's sep-sethood relates \$2 to \$3 \\
defined by \texttt{rel-sep-sethood} \\
characterized by \texttt{rel-sep-sethood}  \textit{(definitional)}:
\begin{quote}\ttfamily\footnotesize\raggedright\noindent
forall(\allowbreak{}[mit\_, mihs\_, mic\_], rel-sep-sethood(\allowbreak{}mit\_, mihs\_, mic\_) iff mic\_ in seq-set and forsome(\allowbreak{}[mig\_, mia\_, mip\_], mic\_ = [mig\_, enc-in(\allowbreak{}enc-sep(\allowbreak{}mia\_, mip\_), enc-setc)] and nbg-ok(\allowbreak{}mip\_) and mihs\_ = [[mig\_, enc-in(\allowbreak{}mia\_, enc-setc)]]))
\end{quote}
mentioned by 3 result(s)

\smallskip

\noindent\textbf{\texttt{rel-tfi-3cases(mit\_, mihs\_, mic\_)}}\quad\textit{predicate}\index{rel-tfi-3cases(mit_, mihs_, mic_)@\texttt{rel-tfi-3cases(mit\_, mihs\_, mic\_)}}\\
reads: the kernel's transfinite-induction-3cases relates \$2 to \$3 \\
defined by \texttt{rel-tfi-3cases} \\
characterized by \texttt{rel-tfi-3cases}  \textit{(definitional)}:
\begin{quote}\ttfamily\footnotesize\raggedright\noindent
forall(\allowbreak{}[mit\_, mihs\_, mic\_], rel-tfi-3cases(\allowbreak{}mit\_, mihs\_, mic\_) iff mic\_ in seq-set and forsome(\allowbreak{}[mig\_, mip\_, mis\_, mifs\_], mic\_ = [mig\_, enc-all(\allowbreak{}enc-imp(\allowbreak{}enc-in(\allowbreak{}ebvar(\allowbreak{}0), enc-ordc), mip\_))] and mihs\_ = [[mig\_, expr-open(\allowbreak{}mip\_, 0, enc-num(\allowbreak{}0))], [mig\_, enc-all(\allowbreak{}mis\_)], [mig\_, enc-all(\allowbreak{}enc-imp(\allowbreak{}enc-and(\allowbreak{}enc-limit(\allowbreak{}ebvar(\allowbreak{}0)), enc-all(\allowbreak{}enc-imp(\allowbreak{}enc-ordlt(\allowbreak{}ebvar(\allowbreak{}0), ebvar(\allowbreak{}1)), mip\_))), mip\_))]] and mifs\_ in set and mifs\_ subset nn and card(\allowbreak{}mifs\_) in nn and forall(\allowbreak{}[miy\_ in nn], not(\allowbreak{}miy\_ in mifs\_) implies expr-open(\allowbreak{}mis\_, 0, efvar(\allowbreak{}miy\_)) = enc-imp(\allowbreak{}enc-and(\allowbreak{}enc-in(\allowbreak{}efvar(\allowbreak{}miy\_), enc-ordc), expr-open(\allowbreak{}mip\_, 0, efvar(\allowbreak{}miy\_))), expr-open(\allowbreak{}mip\_, 0, enc-succo(\allowbreak{}efvar(\allowbreak{}miy\_)))))))
\end{quote}
mentioned by 2 result(s)

\smallskip

\noindent\textbf{\texttt{rel-tfi-3cases-nbg(mit\_, mihs\_, mic\_)}}\quad\textit{predicate}\index{rel-tfi-3cases-nbg(mit_, mihs_, mic_)@\texttt{rel-tfi-3cases-nbg(mit\_, mihs\_, mic\_)}}\\
reads: the kernel's transfinite-induction-3cases at a predicative body relates \$2 to \$3 \\
defined by \texttt{rel-tfi-3cases-nbg} \\
characterized by \texttt{rel-tfi-3cases-nbg}  \textit{(definitional)}:
\begin{quote}\ttfamily\footnotesize\raggedright\noindent
forall(\allowbreak{}[mit\_, mihs\_, mic\_], rel-tfi-3cases-nbg(\allowbreak{}mit\_, mihs\_, mic\_) iff mic\_ in seq-set and forsome(\allowbreak{}[mig\_, mip\_, mis\_, mifs\_], mic\_ = [mig\_, enc-all(\allowbreak{}enc-imp(\allowbreak{}enc-in(\allowbreak{}ebvar(\allowbreak{}0), enc-ordc), mip\_))] and nbg-ok(\allowbreak{}mip\_) and mihs\_ = [[mig\_, expr-open(\allowbreak{}mip\_, 0, enc-num(\allowbreak{}0))], [mig\_, enc-all(\allowbreak{}mis\_)], [mig\_, enc-all(\allowbreak{}enc-imp(\allowbreak{}enc-and(\allowbreak{}enc-limit(\allowbreak{}ebvar(\allowbreak{}0)), enc-all(\allowbreak{}enc-imp(\allowbreak{}enc-ordlt(\allowbreak{}ebvar(\allowbreak{}0), ebvar(\allowbreak{}1)), mip\_))), mip\_))]] and mifs\_ in set and mifs\_ subset nn and card(\allowbreak{}mifs\_) in nn and forall(\allowbreak{}[miy\_ in nn], not(\allowbreak{}miy\_ in mifs\_) implies expr-open(\allowbreak{}mis\_, 0, efvar(\allowbreak{}miy\_)) = enc-imp(\allowbreak{}enc-and(\allowbreak{}enc-in(\allowbreak{}efvar(\allowbreak{}miy\_), enc-ordc), expr-open(\allowbreak{}mip\_, 0, efvar(\allowbreak{}miy\_))), expr-open(\allowbreak{}mip\_, 0, enc-succo(\allowbreak{}efvar(\allowbreak{}miy\_)))))))
\end{quote}
mentioned by 4 result(s)

\smallskip

\noindent\textbf{\texttt{rel-transfinite-induction(mit\_, mihs\_, mic\_)}}\quad\textit{predicate}\index{rel-transfinite-induction(mit_, mihs_, mic_)@\texttt{rel-transfinite-induction(mit\_, mihs\_, mic\_)}}\\
reads: the kernel's transfinite-induction relates \$2 to \$3 \\
defined by \texttt{rel-transfinite-induction} \\
characterized by \texttt{rel-transfinite-induction}  \textit{(definitional)}:
\begin{quote}\ttfamily\footnotesize\raggedright\noindent
forall(\allowbreak{}[mit\_, mihs\_, mic\_], rel-transfinite-induction(\allowbreak{}mit\_, mihs\_, mic\_) iff mic\_ in seq-set and forsome(\allowbreak{}[mig\_, mip\_], mic\_ = [mig\_, enc-all(\allowbreak{}enc-imp(\allowbreak{}enc-in(\allowbreak{}ebvar(\allowbreak{}0), enc-ordc), mip\_))] and mihs\_ = [[mig\_, enc-all(\allowbreak{}enc-imp(\allowbreak{}enc-and(\allowbreak{}enc-in(\allowbreak{}ebvar(\allowbreak{}0), enc-ordc), enc-all(\allowbreak{}enc-imp(\allowbreak{}enc-ordlt(\allowbreak{}ebvar(\allowbreak{}0), ebvar(\allowbreak{}1)), mip\_))), mip\_))]]))
\end{quote}
mentioned by 2 result(s)

\smallskip

\noindent\textbf{\texttt{rel-transfinite-induction-nbg(mit\_, mihs\_, mic\_)}}\quad\textit{predicate}\index{rel-transfinite-induction-nbg(mit_, mihs_, mic_)@\texttt{rel-transfinite-induction-nbg(mit\_, mihs\_, mic\_)}}\\
reads: the kernel's transfinite-induction at a predicative body relates \$2 to \$3 \\
defined by \texttt{rel-transfinite-induction-nbg} \\
characterized by \texttt{rel-transfinite-induction-nbg}  \textit{(definitional)}:
\begin{quote}\ttfamily\footnotesize\raggedright\noindent
forall(\allowbreak{}[mit\_, mihs\_, mic\_], rel-transfinite-induction-nbg(\allowbreak{}mit\_, mihs\_, mic\_) iff mic\_ in seq-set and forsome(\allowbreak{}[mig\_, mip\_], mic\_ = [mig\_, enc-all(\allowbreak{}enc-imp(\allowbreak{}enc-in(\allowbreak{}ebvar(\allowbreak{}0), enc-ordc), mip\_))] and nbg-ok(\allowbreak{}mip\_) and mihs\_ = [[mig\_, enc-all(\allowbreak{}enc-imp(\allowbreak{}enc-and(\allowbreak{}enc-in(\allowbreak{}ebvar(\allowbreak{}0), enc-ordc), enc-all(\allowbreak{}enc-imp(\allowbreak{}enc-ordlt(\allowbreak{}ebvar(\allowbreak{}0), ebvar(\allowbreak{}1)), mip\_))), mip\_))]]))
\end{quote}
mentioned by 4 result(s)

\smallskip

\noindent\textbf{\texttt{rel-tuple-eq-decompose-1(mit\_, mihs\_, mic\_)}}\quad\textit{predicate}\index{rel-tuple-eq-decompose-1(mit_, mihs_, mic_)@\texttt{rel-tuple-eq-decompose-1(mit\_, mihs\_, mic\_)}}\\
reads: the kernel's tuple-equality-decompose, at one position, relates \$2 to \$3 \\
defined by \texttt{rel-tuple-eq-decompose-1} \\
characterized by \texttt{rel-tuple-eq-decompose-1}  \textit{(definitional)}:
\begin{quote}\ttfamily\footnotesize\raggedright\noindent
forall(\allowbreak{}[mit\_, mihs\_, mic\_], rel-tuple-eq-decompose-1(\allowbreak{}mit\_, mihs\_, mic\_) iff rel-rewrite-1(\allowbreak{}j-teq(\allowbreak{}mit\_), mihs\_, mic\_))
\end{quote}
mentioned by 3 result(s)

\smallskip

\noindent\textbf{\texttt{rel-tuple-eq-decompose-7(mit\_, mihs\_, mic\_)}}\quad\textit{predicate}\index{rel-tuple-eq-decompose-7(mit_, mihs_, mic_)@\texttt{rel-tuple-eq-decompose-7(mit\_, mihs\_, mic\_)}}\\
reads: the kernel's tuple-equality-decompose, at one position at any stage, relates \$2 to \$3 \\
defined by \texttt{rel-tuple-eq-decompose-7} \\
characterized by \texttt{rel-tuple-eq-decompose-7}  \textit{(definitional)}:
\begin{quote}\ttfamily\footnotesize\raggedright\noindent
forall(\allowbreak{}[mit\_, mihs\_, mic\_], rel-tuple-eq-decompose-7(\allowbreak{}mit\_, mihs\_, mic\_) iff rel-rewrite-7(\allowbreak{}j-teq(\allowbreak{}mit\_), mihs\_, mic\_))
\end{quote}
mentioned by 3 result(s)

\smallskip

\noindent\textbf{\texttt{rel-tuples-elim(mct\_, mchs\_, mcc\_)}}\quad\textit{predicate}\index{rel-tuples-elim(mct_, mchs_, mcc_)@\texttt{rel-tuples-elim(mct\_, mchs\_, mcc\_)}}\\
reads: the kernel's tuples-elim relates \$2 to \$3 \\
defined by \texttt{rel-tuples-elim} \\
characterized by \texttt{rel-tuples-elim}  \textit{(definitional)}:
\begin{quote}\ttfamily\footnotesize\raggedright\noindent
forall(\allowbreak{}[mct\_, mchs\_, mcc\_], rel-tuples-elim(\allowbreak{}mct\_, mchs\_, mcc\_) iff mcc\_ in seq-set and forsome(\allowbreak{}[mcg\_, mcf\_, mcl\_, mca\_, mck\_], mcc\_ = [mcg\_, mcf\_] and enc-in(\allowbreak{}enc-list(\allowbreak{}mcl\_), enc-tuples(\allowbreak{}mca\_)) in mcg\_ and mck\_ in nn and 1 <= mck\_ and mck\_ <= length(\allowbreak{}mcl\_) and mchs\_ = [[ctx-add(\allowbreak{}mcg\_, enc-in(\allowbreak{}nth(\allowbreak{}mck\_, mcl\_), mca\_)), mcf\_]]))
\end{quote}
mentioned by 3 result(s)

\smallskip

\noindent\textbf{\texttt{rel-tuples-intro(mct\_, mchs\_, mcc\_)}}\quad\textit{predicate}\index{rel-tuples-intro(mct_, mchs_, mcc_)@\texttt{rel-tuples-intro(mct\_, mchs\_, mcc\_)}}\\
reads: the kernel's tuples-intro relates \$2 to \$3 \\
defined by \texttt{rel-tuples-intro} \\
characterized by \texttt{rel-tuples-intro}  \textit{(definitional)}:
\begin{quote}\ttfamily\footnotesize\raggedright\noindent
forall(\allowbreak{}[mct\_, mchs\_, mcc\_], rel-tuples-intro(\allowbreak{}mct\_, mchs\_, mcc\_) iff mcc\_ in seq-set and forsome(\allowbreak{}[mcg\_, mcl\_, mca\_], mcc\_ = [mcg\_, enc-in(\allowbreak{}enc-list(\allowbreak{}mcl\_), enc-tuples(\allowbreak{}mca\_))] and length(\allowbreak{}mchs\_) = length(\allowbreak{}mcl\_) and forall(\allowbreak{}[mci\_ in nn], 1 <= mci\_ implies mci\_ <= length(\allowbreak{}mcl\_) implies nth(\allowbreak{}mci\_, mchs\_) = [mcg\_, enc-in(\allowbreak{}nth(\allowbreak{}mci\_, mcl\_), mca\_)])))
\end{quote}
mentioned by 3 result(s)

\smallskip

\noindent\textbf{\texttt{rel-union-decompose-1(md14t\_, md14hs\_, md14c\_)}}\quad\textit{predicate}\index{rel-union-decompose-1(md14t_, md14hs_, md14c_)@\texttt{rel-union-decompose-1(md14t\_, md14hs\_, md14c\_)}}\\
reads: the kernel's rewrite by the catalog's n-ary decomposition, at one position, relates \$2 to \$3 \\
defined by \texttt{rel-union-decompose-1} \\
characterized by \texttt{rel-union-decompose-1}  \textit{(definitional)}:
\begin{quote}\ttfamily\footnotesize\raggedright\noindent
forall(\allowbreak{}[md14t\_, md14hs\_, md14c\_], rel-union-decompose-1(\allowbreak{}md14t\_, md14hs\_, md14c\_) iff rel-rewrite-1(\allowbreak{}j-udec(\allowbreak{}md14t\_), md14hs\_, md14c\_))
\end{quote}
mentioned by 3 result(s)

\smallskip

\noindent\textbf{\texttt{rel-union-decompose-7(md14t\_, md14hs\_, md14c\_)}}\quad\textit{predicate}\index{rel-union-decompose-7(md14t_, md14hs_, md14c_)@\texttt{rel-union-decompose-7(md14t\_, md14hs\_, md14c\_)}}\\
reads: the kernel's rewrite by the catalog's n-ary decomposition, at one position, relates \$2 to \$3 \\
defined by \texttt{rel-union-decompose-7} \\
characterized by \texttt{rel-union-decompose-7}  \textit{(definitional)}:
\begin{quote}\ttfamily\footnotesize\raggedright\noindent
forall(\allowbreak{}[md14t\_, md14hs\_, md14c\_], rel-union-decompose-7(\allowbreak{}md14t\_, md14hs\_, md14c\_) iff rel-rewrite-7(\allowbreak{}j-udec(\allowbreak{}md14t\_), md14hs\_, md14c\_))
\end{quote}
mentioned by 3 result(s)

\smallskip

\noindent\textbf{\texttt{rel-union-decompose-hyp-1(md14t\_, md14hs\_, md14c\_)}}\quad\textit{predicate}\index{rel-union-decompose-hyp-1(md14t_, md14hs_, md14c_)@\texttt{rel-union-decompose-hyp-1(md14t\_, md14hs\_, md14c\_)}}\\
reads: the kernel's rewrite by the catalog's n-ary decomposition, at one position, relates \$2 to \$3 \\
defined by \texttt{rel-union-decompose-hyp-1} \\
characterized by \texttt{rel-union-decompose-hyp-1}  \textit{(definitional)}:
\begin{quote}\ttfamily\footnotesize\raggedright\noindent
forall(\allowbreak{}[md14t\_, md14hs\_, md14c\_], rel-union-decompose-hyp-1(\allowbreak{}md14t\_, md14hs\_, md14c\_) iff rel-rewrite-hyp-1(\allowbreak{}j-udec(\allowbreak{}md14t\_), md14hs\_, md14c\_))
\end{quote}
mentioned by 3 result(s)

\smallskip

\noindent\textbf{\texttt{rel-union-decompose-hyp-7(md14t\_, md14hs\_, md14c\_)}}\quad\textit{predicate}\index{rel-union-decompose-hyp-7(md14t_, md14hs_, md14c_)@\texttt{rel-union-decompose-hyp-7(md14t\_, md14hs\_, md14c\_)}}\\
reads: the kernel's rewrite by the catalog's n-ary decomposition, at one position, relates \$2 to \$3 \\
defined by \texttt{rel-union-decompose-hyp-7} \\
characterized by \texttt{rel-union-decompose-hyp-7}  \textit{(definitional)}:
\begin{quote}\ttfamily\footnotesize\raggedright\noindent
forall(\allowbreak{}[md14t\_, md14hs\_, md14c\_], rel-union-decompose-hyp-7(\allowbreak{}md14t\_, md14hs\_, md14c\_) iff rel-rewrite-hyp-7(\allowbreak{}j-udec(\allowbreak{}md14t\_), md14hs\_, md14c\_))
\end{quote}
mentioned by 3 result(s)

\smallskip

\noindent\textbf{\texttt{rel-union-elim(mct\_, mchs\_, mcc\_)}}\quad\textit{predicate}\index{rel-union-elim(mct_, mchs_, mcc_)@\texttt{rel-union-elim(mct\_, mchs\_, mcc\_)}}\\
reads: the kernel's union-elim relates \$2 to \$3 \\
defined by \texttt{rel-union-elim} \\
characterized by \texttt{rel-union-elim}  \textit{(definitional)}:
\begin{quote}\ttfamily\footnotesize\raggedright\noindent
forall(\allowbreak{}[mct\_, mchs\_, mcc\_], rel-union-elim(\allowbreak{}mct\_, mchs\_, mcc\_) iff mcc\_ in seq-set and forsome(\allowbreak{}[mcg\_, mcf\_, mcx\_, mcl\_], mcc\_ = [mcg\_, mcf\_] and enc-in(\allowbreak{}mcx\_, enc-union(\allowbreak{}mcl\_)) in mcg\_ and length(\allowbreak{}mchs\_) = length(\allowbreak{}mcl\_) and forall(\allowbreak{}[mci\_ in nn], 1 <= mci\_ implies mci\_ <= length(\allowbreak{}mcl\_) implies nth(\allowbreak{}mci\_, mchs\_) = [ctx-add(\allowbreak{}difference(\allowbreak{}mcg\_, pair(\allowbreak{}enc-in(\allowbreak{}mcx\_, enc-union(\allowbreak{}mcl\_)), enc-in(\allowbreak{}mcx\_, enc-union(\allowbreak{}mcl\_)))), enc-in(\allowbreak{}mcx\_, nth(\allowbreak{}mci\_, mcl\_))), mcf\_])))
\end{quote}
mentioned by 3 result(s)

\smallskip

\noindent\textbf{\texttt{rel-union-intro(mct\_, mchs\_, mcc\_)}}\quad\textit{predicate}\index{rel-union-intro(mct_, mchs_, mcc_)@\texttt{rel-union-intro(mct\_, mchs\_, mcc\_)}}\\
reads: the kernel's union-intro relates \$2 to \$3 \\
defined by \texttt{rel-union-intro} \\
characterized by \texttt{rel-union-intro}  \textit{(definitional)}:
\begin{quote}\ttfamily\footnotesize\raggedright\noindent
forall(\allowbreak{}[mct\_, mchs\_, mcc\_], rel-union-intro(\allowbreak{}mct\_, mchs\_, mcc\_) iff mcc\_ in seq-set and forsome(\allowbreak{}[mcg\_, mcx\_, mcl\_, mck\_], mcc\_ = [mcg\_, enc-in(\allowbreak{}mcx\_, enc-union(\allowbreak{}mcl\_))] and mck\_ in nn and 1 <= mck\_ and mck\_ <= length(\allowbreak{}mcl\_) and mchs\_ = [[mcg\_, enc-in(\allowbreak{}mcx\_, nth(\allowbreak{}mck\_, mcl\_))]]))
\end{quote}
mentioned by 3 result(s)

\smallskip

\noindent\textbf{\texttt{related(s, a, b)}}\quad\textit{functoid}\index{related(s, a, b)@\texttt{related(s, a, b)}}\\
definition:
\begin{quote}\ttfamily\footnotesize\raggedright\noindent
[a, b] in rel(\allowbreak{}s)
\end{quote}
characterized by \texttt{sip-related-in-csip-as-sip}  \textit{(definitional)}:
\begin{quote}\ttfamily\footnotesize\raggedright\noindent
forall(\allowbreak{}[r], is-complex-semi-inner-product-space(\allowbreak{}r) implies forall(\allowbreak{}[a\_ in vec(\allowbreak{}r), b\_ in vec(\allowbreak{}r)], (\allowbreak{}rip(\allowbreak{}r))(\allowbreak{}(\allowbreak{}vadd(\allowbreak{}r))(\allowbreak{}a\_, (\allowbreak{}vneg(\allowbreak{}r))(\allowbreak{}b\_)), (\allowbreak{}vadd(\allowbreak{}r))(\allowbreak{}a\_, (\allowbreak{}vneg(\allowbreak{}r))(\allowbreak{}b\_))) = 0 implies related(\allowbreak{}sip-setoid(\allowbreak{}csip-as-sip(\allowbreak{}r)), a\_, b\_)))
\end{quote}
also: \texttt{sip-related-out-csip-as-sip} \texttt{qmod-related-in-cna-as-module} \texttt{qmod-related-in-sip-as-module} \texttt{qmod-related-in-uba-as-module} \texttt{qmod-related-in-cnvs-as-module} \texttt{qmod-related-out-cna-as-module} \texttt{qmod-related-out-sip-as-module} \texttt{qmod-related-out-uba-as-module} ... (21 in all) \\
mentioned by 39 result(s)

\smallskip

\noindent\textbf{\texttt{rep-functional(g4a\_, g4v\_, g4p\_, g4e\_)}}\quad\textit{functoid}\index{rep-functional(g4a_, g4v_, g4p_, g4e_)@\texttt{rep-functional(g4a\_, g4v\_, g4p\_, g4e\_)}}\\
reads: the functional x |-> <\$3(x) \$4 | \$4> \\
definition:
\begin{quote}\ttfamily\footnotesize\raggedright\noindent
vnb-lambda(\allowbreak{}g4x\_, vec(\allowbreak{}g4a\_), (\allowbreak{}ip(\allowbreak{}g4v\_))(\allowbreak{}(\allowbreak{}g4p\_(\allowbreak{}g4x\_))(\allowbreak{}g4e\_), g4e\_))
\end{quote}
characterized by \texttt{rep-functional-norm}  \textit{(proven)}:
\begin{quote}\ttfamily\footnotesize\raggedright\noindent
forall(\allowbreak{}[a], is-unital-involutive-banach-algebra(\allowbreak{}a) implies forall(\allowbreak{}[v, p], is-representation(\allowbreak{}a, v, p) implies forall(\allowbreak{}[e in vec(\allowbreak{}v)], rep-functional(\allowbreak{}a, v, p, e) in cdual-vec(\allowbreak{}a) and cdual-norm(\allowbreak{}a, rep-functional(\allowbreak{}a, v, p, e)) = ip-norm(\allowbreak{}v, e) * ip-norm(\allowbreak{}v, e))))
\end{quote}
also: \texttt{rep-functional-unit} \texttt{rep-functional-value} \texttt{rep-functional-square} \texttt{rep-functional-positive} \\
mentioned by 8 result(s)

\smallskip

\noindent\textbf{\texttt{rep-image(amv\_, amw\_, amf\_, ams\_)}}\quad\textit{functoid}\index{rep-image(amv_, amw_, amf_, ams_)@\texttt{rep-image(amv\_, amw\_, amf\_, ams\_)}}\\
reads: the image of \$4 under \$3 \\
definition:
\begin{quote}\ttfamily\footnotesize\raggedright\noindent
\{amb\_ in bounded-ops(\allowbreak{}amw\_, amw\_): forsome(\allowbreak{}[amt\_ in ams\_], amb\_ = amf\_(\allowbreak{}amt\_))\}
\end{quote}
characterized by \texttt{rep-image-unfold}  \textit{(proven)}:
\begin{quote}\ttfamily\footnotesize\raggedright\noindent
forall(\allowbreak{}[v, w, f\_, s], rep-image(\allowbreak{}v, w, f\_, s) == \{amb\_ in bounded-ops(\allowbreak{}w, w): forsome(\allowbreak{}[amt\_ in s], amb\_ = f\_(\allowbreak{}amt\_))\})
\end{quote}
also: \texttt{rep-image-bicommutant} \texttt{rep-image-star-algebra} \\
mentioned by 4 result(s)

\smallskip

\noindent\textbf{\texttt{rep-orbit(g4a\_, g4v\_, g4p\_, g4e\_)}}\quad\textit{functoid}\index{rep-orbit(g4a_, g4v_, g4p_, g4e_)@\texttt{rep-orbit(g4a\_, g4v\_, g4p\_, g4e\_)}}\\
reads: the orbit \$3(\$1) \$4 \\
definition:
\begin{quote}\ttfamily\footnotesize\raggedright\noindent
\{g4y\_ in vec(\allowbreak{}g4v\_): forsome(\allowbreak{}[g4x\_ in vec(\allowbreak{}g4a\_)], g4y\_ = (\allowbreak{}g4p\_(\allowbreak{}g4x\_))(\allowbreak{}g4e\_))\}
\end{quote}
mentioned by 2 result(s)

\smallskip

\noindent\textbf{\texttt{repeat}}\quad\textit{tactic}\index{repeat@\texttt{repeat}}\\
\texttt{(repeat thunk [cap])} -- Run thunk until it stops changing the proof state (LCF REPEAT). \\
declared in \texttt{tactics-help.scm}

\smallskip

\noindent\textbf{\texttt{replay-proof}}\quad\textit{tactic}\index{replay-proof@\texttt{replay-proof}}\\
\texttt{(replay-proof 'name [subst])} -- Re-run a saved script on the current goal, optionally renaming free vars. \\
declared in \texttt{tactics-help.scm}

\smallskip

\noindent\textbf{\texttt{res}}\quad\textit{operator}\index{res@\texttt{res}}\\
a term-forming head \\
characterized by \texttt{res-apply}  \textit{(primitive)}:
\begin{quote}\ttfamily\footnotesize\raggedright\noindent
forall(\allowbreak{}[a, b, f, x], f in fun(\allowbreak{}a) and b subset a and x in b implies (\allowbreak{}res(\allowbreak{}f, b))(\allowbreak{}x) = f(\allowbreak{}x))
\end{quote}
also: \texttt{res-typing} \texttt{res-codomain} \texttt{measure-restrict} \texttt{completion-of-restrict-in} \\
mentioned by 38 result(s)

\smallskip

\noindent\textbf{\texttt{residue(mrf\_, mra\_)}}\quad\textit{functoid}\index{residue(mrf_, mra_)@\texttt{residue(mrf\_, mra\_)}}\\
reads: the residue of \$1 at \$2 \\
definition:
\begin{quote}\ttfamily\footnotesize\raggedright\noindent
iota(\allowbreak{}mrx\_, mrx\_ in cc and forsome(\allowbreak{}[mrr\_ in rr], 0 < mrr\_ and forall(\allowbreak{}[mrp\_ in rr], 0 < mrp\_ implies mrp\_ < mrr\_ implies mrx\_ = recip(\allowbreak{}(\allowbreak{}2 * pi) * 1i) * line-int(\allowbreak{}vnb-lambda(\allowbreak{}mrw\_, trace(\allowbreak{}circle-path(\allowbreak{}mra\_, mrp\_), 0, 2 * pi), mrf\_(\allowbreak{}mrw\_)), circle-path(\allowbreak{}mra\_, mrp\_), circle-deriv(\allowbreak{}mra\_, mrp\_), 0, 2 * pi))))
\end{quote}
characterized by \texttt{residue-theorem}  \textit{(proven)}:
\begin{quote}\ttfamily\footnotesize\raggedright\noindent
forall(\allowbreak{}[u], is-open(\allowbreak{}nf-metric-space(\allowbreak{}cc-normed-field), u) implies forall(\allowbreak{}[mse\_ in set], card(\allowbreak{}mse\_) in nn implies mse\_ subset u implies forall(\allowbreak{}[f], holomorphic-on(\allowbreak{}difference(\allowbreak{}u, mse\_), f) implies forall(\allowbreak{}[mqe\_ in mse\_], is-meromorphic-at(\allowbreak{}f, mqe\_)) implies forall(\allowbreak{}[mdp\_, mdq\_, mds\_, mdt\_], is-road(\allowbreak{}mdp\_, mdq\_, mds\_, mdt\_) implies trace(\allowbreak{}mdp\_, mds\_, mdt\_) subset difference(\allowbreak{}u, mse\_) implies mdp\_(\allowbreak{}mdt\_) = mdp\_(\allowbreak{}mds\_) implies null-homotopic(\allowbreak{}u, mdp\_, mds\_, mdt\_) implies recip(\allowbreak{}(\allowbreak{}2 * pi) * 1i) * line-int(\allowbreak{}f, mdp\_, mdq\_, mds\_, mdt\_) = finsum(\allowbreak{}normed-field-additive-ag(\allowbreak{}cc-normed-field), vnb-lambda(\allowbreak{}mqx\_, cc, residue(\allowbreak{}f, mqx\_) * winding(\allowbreak{}mdp\_, mdq\_, mds\_, mdt\_, mqx\_)), mse\_)))))
\end{quote}
also: \texttt{residue-of-property} \texttt{residue-theorem-step} \texttt{residue-of-decomposition} \texttt{residue-theorem-induction} \texttt{log-derivative-residue-at-pole} \texttt{log-derivative-residue-at-zero} \texttt{meromorphic-residue-decomposition} \\
mentioned by 13 result(s)

\smallskip

\noindent\textbf{\texttt{resolvent(baa\_, bax\_)}}\quad\textit{functoid}\index{resolvent(baa_, bax_)@\texttt{resolvent(baa\_, bax\_)}}\\
reads: the resolvent of \$2 in \$1 \\
definition:
\begin{quote}\ttfamily\footnotesize\raggedright\noindent
vnb-lambda(\allowbreak{}brz\_, difference(\allowbreak{}cc, spectrum(\allowbreak{}baa\_, bax\_)), inverse-el(\allowbreak{}baa\_, shift-el(\allowbreak{}baa\_, brz\_, bax\_)))
\end{quote}

\smallskip

\noindent\textbf{\texttt{respects(s, f)}}\quad\textit{functoid}\index{respects(s, f)@\texttt{respects(s, f)}}\\
definition:
\begin{quote}\ttfamily\footnotesize\raggedright\noindent
forall(\allowbreak{}[a in pts(\allowbreak{}s), b in pts(\allowbreak{}s)], related(\allowbreak{}s, a, b) implies f(\allowbreak{}a) = f(\allowbreak{}b))
\end{quote}
characterized by \texttt{sipq-lsc-respects-csip-as-sip}  \textit{(definitional)}:
\begin{quote}\ttfamily\footnotesize\raggedright\noindent
forall(\allowbreak{}[r], is-complex-semi-inner-product-space(\allowbreak{}r) implies forall(\allowbreak{}[r\_ in rr], respects(\allowbreak{}sip-setoid(\allowbreak{}csip-as-sip(\allowbreak{}r)), vnb-lambda(\allowbreak{}sqa\_, vec(\allowbreak{}r), class(\allowbreak{}sip-setoid(\allowbreak{}csip-as-sip(\allowbreak{}r)), (\allowbreak{}act(\allowbreak{}r))(\allowbreak{}r\_, sqa\_))))))
\end{quote}
also: \texttt{sipq-lvn-respects-csip-as-sip} \texttt{qmod-lsc-respects-cna-as-module} \texttt{qmod-lsc-respects-sip-as-module} \texttt{qmod-lsc-respects-uba-as-module} \texttt{qmod-lvn-respects-cna-as-module} \texttt{qmod-lvn-respects-sip-as-module} \texttt{qmod-lvn-respects-uba-as-module} \texttt{qmod-lsc-respects-cnvs-as-module} ... (20 in all) \\
mentioned by 27 result(s)

\smallskip

\noindent\textbf{\texttt{respects2(s, f)}}\quad\textit{functoid}\index{respects2(s, f)@\texttt{respects2(s, f)}}\\
definition:
\begin{quote}\ttfamily\footnotesize\raggedright\noindent
forall(\allowbreak{}[a, b, a\_, b\_], a in pts(\allowbreak{}s) implies b in pts(\allowbreak{}s) implies a\_ in pts(\allowbreak{}s) implies b\_ in pts(\allowbreak{}s) implies related(\allowbreak{}s, a, a\_) and related(\allowbreak{}s, b, b\_) implies f(\allowbreak{}a, b) = f(\allowbreak{}a\_, b\_))
\end{quote}
characterized by \texttt{form-respects2}  \textit{(proven)}:
\begin{quote}\ttfamily\footnotesize\raggedright\noindent
forall(\allowbreak{}[md, fq], is-module(\allowbreak{}md) implies scal(\allowbreak{}md) = normed-field-as-commutative-ring(\allowbreak{}cc-normed-field) implies is-psd-form(\allowbreak{}md, fq) implies respects2(\allowbreak{}qmod-setoid(\allowbreak{}md, form-null(\allowbreak{}md, fq)), fq))
\end{quote}
also: \texttt{respects2-unfold} \texttt{r6b-add-respects2} \texttt{r7n-mul-respects2} \\
mentioned by 20 result(s)

\smallskip

\noindent\textbf{\texttt{restrict(f, a)}}\quad\textit{functoid}\index{restrict(f, a)@\texttt{restrict(f, a)}}\\
reads: the restriction of \$1 to \$2 \\
definition:
\begin{quote}\ttfamily\footnotesize\raggedright\noindent
vnb-lambda(\allowbreak{}rsx\_, a, f(\allowbreak{}rsx\_))
\end{quote}
characterized by \texttt{road-restrict}  \textit{(proven)}:
\begin{quote}\ttfamily\footnotesize\raggedright\noindent
forall(\allowbreak{}[pgam, dgam, a, b, licv\_ in rr, lidv\_ in rr], is-road(\allowbreak{}pgam, dgam, a, b) implies a <= licv\_ implies licv\_ < lidv\_ implies lidv\_ <= b implies is-road(\allowbreak{}restrict(\allowbreak{}pgam, ccint(\allowbreak{}licv\_, lidv\_)), restrict(\allowbreak{}dgam, ccint(\allowbreak{}licv\_, lidv\_)), licv\_, lidv\_))
\end{quote}
also: \texttt{image-restrict} \texttt{restrict-apply} \texttt{restrict-in-fun} \texttt{tri-int-restrict} \texttt{trace-restrict-in} \texttt{zero-order-restrict} \texttt{trace-restrict-subset} \texttt{ln-restrict-continuous} ... (22 in all) \\
mentioned by 146 result(s)

\smallskip

\noindent\textbf{\texttt{restrict-sigma(rsc\_, rse\_)}}\quad\textit{functoid}\index{restrict-sigma(rsc_, rse_)@\texttt{restrict-sigma(rsc\_, rse\_)}}\\
reads: the members of \$1 contained in \$2 \\
definition:
\begin{quote}\ttfamily\footnotesize\raggedright\noindent
\{rsa\_ in rsc\_: rsa\_ subset rse\_\}
\end{quote}
characterized by \texttt{restrict-sigma-unfold}  \textit{(proven)}:
\begin{quote}\ttfamily\footnotesize\raggedright\noindent
forall(\allowbreak{}[ca, e], restrict-sigma(\allowbreak{}ca, e) == \{rsa\_ in ca: rsa\_ subset e\})
\end{quote}
also: \texttt{restrict-sigma-algebra} \\
mentioned by 7 result(s)

\smallskip

\noindent\textbf{\texttt{restvar}}\quad\textit{operator}\index{restvar@\texttt{restvar}}\\
a term-forming head \\
mentioned by 4 result(s)

\smallskip

\noindent\textbf{\texttt{rewr(mrt\_, mrk\_)}}\quad\textit{functoid}\index{rewr(mrt_, mrk_)@\texttt{rewr(mrt\_, mrk\_)}}\\
reads: the one-position rewrites justified by \$1, of stage at most \$2 \\
definition:
\begin{quote}\ttfamily\footnotesize\raggedright\noindent
\{mrx\_ in rw-univ: forall(\allowbreak{}[mrs\_ in power(\allowbreak{}rw-univ)], rw-closed(\allowbreak{}mrt\_, mrk\_, mrs\_) implies mrx\_ in mrs\_)\}
\end{quote}
characterized by \texttt{rel-rewrite-1}  \textit{(definitional)}:
\begin{quote}\ttfamily\footnotesize\raggedright\noindent
forall(\allowbreak{}[mrt\_, mrhs\_, mrc\_], rel-rewrite-1(\allowbreak{}mrt\_, mrhs\_, mrc\_) iff mrc\_ in seq-set and forsome(\allowbreak{}[mrg\_, mra\_, mrb\_], mrc\_ = [mrg\_, mra\_] and mrhs\_ = [[mrg\_, mrb\_]] and [2, mrg\_, mra\_, mrb\_] in rewr(\allowbreak{}mrt\_, 5)))
\end{quote}
also: \texttt{rel-rewrite-7} \texttt{rel-rewrite-hyp-1} \texttt{rel-rewrite-hyp-7} \texttt{rewr-least} \texttt{rewr-unfold} \\
mentioned by 25 result(s)

\smallskip

\noindent\textbf{\texttt{rf-closed(md14h\_, md14x\_, md14l\_, md14s\_)}}\quad\textit{predicate}\index{rf-closed(md14h_, md14x_, md14l_, md14s_)@\texttt{rf-closed(md14h\_, md14x\_, md14l\_, md14s\_)}}\\
reads: \$4 is closed under the right folds by \$1 of the memberships of \$2 in the entries of \$3 \\
defined by \texttt{rf-closed} \\
characterized by \texttt{rf-closed}  \textit{(definitional)}:
\begin{quote}\ttfamily\footnotesize\raggedright\noindent
forall(\allowbreak{}[md14h\_, md14x\_, md14l\_, md14s\_], rf-closed(\allowbreak{}md14h\_, md14x\_, md14l\_, md14s\_) iff forall(\allowbreak{}[md14j\_ in nn], succ(\allowbreak{}md14j\_) = length(\allowbreak{}md14l\_) implies [md14j\_, enc-in(\allowbreak{}md14x\_, nth(\allowbreak{}succ(\allowbreak{}md14j\_), md14l\_))] in md14s\_) and forall(\allowbreak{}[md14k\_, md14f\_], md14k\_ in nn implies [succ(\allowbreak{}md14k\_), md14f\_] in md14s\_ implies [md14k\_, eapp(\allowbreak{}md14h\_, [enc-in(\allowbreak{}md14x\_, nth(\allowbreak{}succ(\allowbreak{}md14k\_), md14l\_)), md14f\_])] in md14s\_))
\end{quote}
mentioned by 5 result(s)

\smallskip

\noindent\textbf{\texttt{rf-set(md14h\_, md14x\_, md14l\_)}}\quad\textit{functoid}\index{rf-set(md14h_, md14x_, md14l_)@\texttt{rf-set(md14h\_, md14x\_, md14l\_)}}\\
reads: the right folds by \$1 of the memberships of \$2 in the entries of \$3 from an index on \\
definition:
\begin{quote}\ttfamily\footnotesize\raggedright\noindent
\{md14p\_ in cartesian(\allowbreak{}nn, expr-set): forall(\allowbreak{}[md14w\_ in power(\allowbreak{}cartesian(\allowbreak{}nn, expr-set))], rf-closed(\allowbreak{}md14h\_, md14x\_, md14l\_, md14w\_) implies md14p\_ in md14w\_)\}
\end{quote}
characterized by \texttt{rf-set-least}  \textit{(proven)}:
\begin{quote}\ttfamily\footnotesize\raggedright\noindent
forall(\allowbreak{}[h, x, l, md14xs\_ in power(\allowbreak{}cartesian(\allowbreak{}nn, expr-set))], rf-closed(\allowbreak{}h, x, l, md14xs\_) implies forall(\allowbreak{}[z in rf-set(\allowbreak{}h, x, l)], z in md14xs\_))
\end{quote}
also: \texttt{rf-set-unfold} \\
mentioned by 7 result(s)

\smallskip

\noindent\textbf{\texttt{rfl}}\quad\textit{tactic}\index{rfl@\texttt{rfl}}\\
\texttt{(rfl)} -- Close a reflexive equality goal (t = t). \\
declared in \texttt{tactics-help.scm}

\smallskip

\noindent\textbf{\texttt{ring}}\quad\textit{structure}\index{ring@\texttt{ring}}\\
predicate \texttt{is-ring}; slots: carr add mul neg zero one

\smallskip

\noindent\textbf{\texttt{ring-additive-ag(s)}}\quad\textit{view}\index{ring-additive-ag(s)@\texttt{ring-additive-ag(s)}}\\
a ring viewed as an abelian-group \\
characterized by \texttt{finsum-type-ring-additive-ag}  \textit{(definitional)}:
\begin{quote}\ttfamily\footnotesize\raggedright\noindent
forall(\allowbreak{}[r], is-ring(\allowbreak{}r) implies forall(\allowbreak{}[s in set], card(\allowbreak{}s) in nn implies forall(\allowbreak{}[f in fun(\allowbreak{}s, carr(\allowbreak{}r))], finsum(\allowbreak{}ring-additive-ag(\allowbreak{}r), f, s) in carr(\allowbreak{}r))))
\end{quote}
also: \texttt{sum-ag-type-ring-additive-ag} \texttt{zz-act-assoc-ring-additive-ag} \texttt{finsum-all-id-ring-additive-ag} \texttt{zz-act-distrib-ring-additive-ag} \texttt{ring-additive-ag-is-abelian-group} \texttt{sum-ag-as-reduce-ring-additive-ag} \texttt{sum-ag-singleton-ring-additive-ag} \texttt{finsum-single-support-ring-additive-ag} ... (37 in all) \\
mentioned by 504 result(s)

\smallskip

\noindent\textbf{\texttt{ring-as-module(r)}}\quad\textit{functoid}\index{ring-as-module(r)@\texttt{ring-as-module(r)}}\\
reads: \$1 as a module over itself \\
definition:
\begin{quote}\ttfamily\footnotesize\raggedright\noindent
[r, carr(\allowbreak{}r), add(\allowbreak{}r), zero(\allowbreak{}r), neg(\allowbreak{}r), mul(\allowbreak{}r)]
\end{quote}
characterized by \texttt{ring-as-module-is-module-uba-as-ring}  \textit{(definitional)}:
\begin{quote}\ttfamily\footnotesize\raggedright\noindent
forall(\allowbreak{}[r], is-unital-banach-algebra(\allowbreak{}r) implies is-module(\allowbreak{}ring-as-module(\allowbreak{}uba-as-ring(\allowbreak{}r))))
\end{quote}
also: \texttt{ring-as-module-is-module-uba-as-ring-uiba-as-uba} \texttt{ring-as-module-is-module} \\
mentioned by 131 result(s)

\smallskip

\noindent\textbf{\texttt{ring-multiplicative-monoid(s)}}\quad\textit{view}\index{ring-multiplicative-monoid(s)@\texttt{ring-multiplicative-monoid(s)}}\\
a ring viewed as a monoid \\
characterized by \texttt{ring-multiplicative-monoid-is-monoid}  \textit{(definitional)}:
\begin{quote}\ttfamily\footnotesize\raggedright\noindent
forall(\allowbreak{}[r], is-ring(\allowbreak{}r) implies is-monoid(\allowbreak{}ring-multiplicative-monoid(\allowbreak{}r)))
\end{quote}
also: \texttt{ring-multiplicative-monoid-is-monoid-uba-as-ring} \texttt{ring-multiplicative-monoid-is-monoid-uba-as-ring-uiba-as-uba} \texttt{ring-multiplicative-monoid-functorial} \\
mentioned by 118 result(s)

\smallskip

\noindent\textbf{\texttt{ring-pow(r, n)}}\quad\textit{functoid}\index{ring-pow(r, n)@\texttt{ring-pow(r, n)}}\\
reads: the free module \$1\textasciicircum{}\$2 \\
definition:
\begin{quote}\ttfamily\footnotesize\raggedright\noindent
direct-sum(\allowbreak{}r, n, vnb-lambda(\allowbreak{}rqpi\_, interval(\allowbreak{}1, n), ring-as-module(\allowbreak{}r)))
\end{quote}
characterized by \texttt{ring-pow-ext}  \textit{(proven)}:
\begin{quote}\ttfamily\footnotesize\raggedright\noindent
forall(\allowbreak{}[r, n, x, y], is-ring(\allowbreak{}r) implies n in nn implies x in vec(\allowbreak{}ring-pow(\allowbreak{}r, n)) implies y in vec(\allowbreak{}ring-pow(\allowbreak{}r, n)) implies forall(\allowbreak{}[i\_ in interval(\allowbreak{}1, n)], x(\allowbreak{}i\_) = y(\allowbreak{}i\_)) implies x = y)
\end{quote}
also: \texttt{ring-pow-scal} \texttt{ring-pow-act-at} \texttt{ring-pow-unfold} \texttt{ring-pow-vec-in} \texttt{ring-pow-vadd-at} \texttt{ring-pow-vec-out} \texttt{ring-pow-vzero-at} \texttt{ring-pow-form-type} ... (17 in all) \\
mentioned by 50 result(s)

\smallskip

\noindent\textbf{\texttt{ring-power(r, x, n)}}\quad\textit{functoid}\index{ring-power(r, x, n)@\texttt{ring-power(r, x, n)}}\\
definition:
\begin{quote}\ttfamily\footnotesize\raggedright\noindent
mpow(\allowbreak{}commutative-ring-multiplicative-cm(\allowbreak{}r), x, n)
\end{quote}
characterized by \texttt{ring-power-zero-uba-as-ring}  \textit{(definitional)}:
\begin{quote}\ttfamily\footnotesize\raggedright\noindent
forall(\allowbreak{}[r], is-unital-banach-algebra(\allowbreak{}r) implies forall(\allowbreak{}[x], ring-power(\allowbreak{}uba-as-ring(\allowbreak{}r), x, 0) == vone(\allowbreak{}r)))
\end{quote}
also: \texttt{ring-power-zero-uba-as-ring-uiba-as-uba} \texttt{ring-power-add} \texttt{ring-power-one} \texttt{ring-power-mult} \texttt{ring-power-succ} \texttt{ring-power-type} \texttt{ring-power-zero} \texttt{ring-power-of-one} ... (13 in all) \\
mentioned by 101 result(s)

\smallskip

\noindent\textbf{\texttt{ring-prod(x, y)}}\quad\textit{functoid}\index{ring-prod(x, y)@\texttt{ring-prod(x, y)}}\\
definition:
\begin{quote}\ttfamily\footnotesize\raggedright\noindent
[cartesian(\allowbreak{}carr(\allowbreak{}x), carr(\allowbreak{}y)), vnb-lambda(\allowbreak{}[p, q], cartesian(\allowbreak{}cartesian(\allowbreak{}carr(\allowbreak{}x), carr(\allowbreak{}y)), cartesian(\allowbreak{}carr(\allowbreak{}x), carr(\allowbreak{}y))), [(\allowbreak{}add(\allowbreak{}x))(\allowbreak{}nth(\allowbreak{}1, p), nth(\allowbreak{}1, q)), (\allowbreak{}add(\allowbreak{}y))(\allowbreak{}nth(\allowbreak{}2, p), nth(\allowbreak{}2, q))]), vnb-lambda(\allowbreak{}[p, q], cartesian(\allowbreak{}cartesian(\allowbreak{}carr(\allowbreak{}x), carr(\allowbreak{}y)), cartesian(\allowbreak{}carr(\allowbreak{}x), carr(\allowbreak{}y))), [(\allowbreak{}mul(\allowbreak{}x))(\allowbreak{}nth(\allowbreak{}1, p), nth(\allowbreak{}1, q)), (\allowbreak{}mul(\allowbreak{}y))(\allowbreak{}nth(\allowbreak{}2, p), nth(\allowbreak{}2, q))]), vnb-lambda(\allowbreak{}p, cartesian(\allowbreak{}carr(\allowbreak{}x), carr(\allowbreak{}y)), [(\allowbreak{}neg(\allowbreak{}x))(\allowbreak{}nth(\allowbreak{}1, p)), (\allowbreak{}neg(\allowbreak{}y))(\allowbreak{}nth(\allowbreak{}2, p))]), [zero(\allowbreak{}x), zero(\allowbreak{}y)], [one(\allowbreak{}x), one(\allowbreak{}y)]]
\end{quote}
characterized by \texttt{ring-prod-n-succ}  \textit{(definitional)}:
\begin{quote}\ttfamily\footnotesize\raggedright\noindent
forall(\allowbreak{}[f, n in nn], ring-prod-n(\allowbreak{}f, succ(\allowbreak{}n)) == ring-prod(\allowbreak{}ring-prod-n(\allowbreak{}f, n), f(\allowbreak{}n)))
\end{quote}
also: \texttt{ring-prod-is-ring-uba-as-ring} \texttt{ring-prod-is-ring-uba-as-ring-uiba-as-uba} \texttt{ring-prod-is-ring} \\
mentioned by 5 result(s)

\smallskip

\noindent\textbf{\texttt{ring-prod-n}}\quad\textit{defined-fn}\index{ring-prod-n@\texttt{ring-prod-n}}\\
defined by \texttt{ring-prod-n-zero} \texttt{ring-prod-n-succ} \\
characterized by \texttt{ring-prod-n-succ}  \textit{(definitional)}:
\begin{quote}\ttfamily\footnotesize\raggedright\noindent
forall(\allowbreak{}[f, n in nn], ring-prod-n(\allowbreak{}f, succ(\allowbreak{}n)) == ring-prod(\allowbreak{}ring-prod-n(\allowbreak{}f, n), f(\allowbreak{}n)))
\end{quote}
also: \texttt{ring-prod-n-zero} \texttt{ring-prod-n-is-ring} \texttt{ring-prod-n-is-ring-ind} \\
mentioned by 6 result(s)

\smallskip

\noindent\textbf{\texttt{ring-row(r, n, x)}}\quad\textit{functoid}\index{ring-row(r, n, x)@\texttt{ring-row(r, n, x)}}\\
reads: the coefficient row of \$3 in \$1\textasciicircum{}\$2 \\
definition:
\begin{quote}\ttfamily\footnotesize\raggedright\noindent
matof(\allowbreak{}1, n, vnb-lambda(\allowbreak{}[rqra\_, rqrb\_], cartesian(\allowbreak{}interval(\allowbreak{}1, 1), interval(\allowbreak{}1, n)), x(\allowbreak{}rqrb\_)))
\end{quote}
characterized by \texttt{ring-row-act}  \textit{(proven)}:
\begin{quote}\ttfamily\footnotesize\raggedright\noindent
forall(\allowbreak{}[r, n, a, x], is-ring(\allowbreak{}r) implies n in nn implies a in carr(\allowbreak{}r) implies x in vec(\allowbreak{}ring-pow(\allowbreak{}r, n)) implies ring-row(\allowbreak{}r, n, (\allowbreak{}act(\allowbreak{}ring-pow(\allowbreak{}r, n)))(\allowbreak{}a, x)) = matscale(\allowbreak{}r, a, ring-row(\allowbreak{}r, n, x)))
\end{quote}
also: \texttt{ring-row-type} \texttt{ring-row-vadd} \texttt{ring-row-entry} \texttt{zz-row-is-ring-row} \\
mentioned by 26 result(s)

\smallskip

\noindent\textbf{\texttt{ring-units(r, n)}}\quad\textit{functoid}\index{ring-units(r, n)@\texttt{ring-units(r, n)}}\\
reads: the unit vectors of \$1\textasciicircum{}\$2 \\
definition:
\begin{quote}\ttfamily\footnotesize\raggedright\noindent
matof(\allowbreak{}n, 1, vnb-lambda(\allowbreak{}[rqua\_, rqub\_], cartesian(\allowbreak{}interval(\allowbreak{}1, n), interval(\allowbreak{}1, 1)), vnb-lambda(\allowbreak{}rqut\_, interval(\allowbreak{}1, n), entry(\allowbreak{}unitrow(\allowbreak{}r, n, rqua\_), 1, rqut\_))))
\end{quote}
characterized by \texttt{ring-units-type}  \textit{(proven)}:
\begin{quote}\ttfamily\footnotesize\raggedright\noindent
forall(\allowbreak{}[r, n], is-ring(\allowbreak{}r) implies n in nn implies ring-units(\allowbreak{}r, n) in mat(\allowbreak{}n, 1, vec(\allowbreak{}ring-pow(\allowbreak{}r, n))))
\end{quote}
also: \texttt{ring-units-entry} \texttt{ring-units-generate} \texttt{ring-units-rel-free} \texttt{zz-units-is-ring-units} \\
mentioned by 11 result(s)

\smallskip

\noindent\textbf{\texttt{ringoid}}\quad\textit{structure}\index{ringoid@\texttt{ringoid}}\\
predicate \texttt{is-ringoid}; slots: carr add mul neg zero one idl

\smallskip

\noindent\textbf{\texttt{ringoid-as-ring(s)}}\quad\textit{view}\index{ringoid-as-ring(s)@\texttt{ringoid-as-ring(s)}}\\
a ringoid viewed as a ring \\
characterized by \texttt{ringoid-as-ring-is-ring}  \textit{(definitional)}:
\begin{quote}\ttfamily\footnotesize\raggedright\noindent
forall(\allowbreak{}[r], is-ringoid(\allowbreak{}r) implies is-ring(\allowbreak{}ringoid-as-ring(\allowbreak{}r)))
\end{quote}
also: \texttt{ringoid-as-ring-functorial} \\
mentioned by 4 result(s)

\smallskip

\noindent\textbf{\texttt{ringoid-quotient(r)}}\quad\textit{functoid}\index{ringoid-quotient(r)@\texttt{ringoid-quotient(r)}}\\
definition:
\begin{quote}\ttfamily\footnotesize\raggedright\noindent
quotient(\allowbreak{}ringoid-setoid(\allowbreak{}r))
\end{quote}

\smallskip

\noindent\textbf{\texttt{ringoid-quotient-ring(r)}}\quad\textit{functoid}\index{ringoid-quotient-ring(r)@\texttt{ringoid-quotient-ring(r)}}\\
definition:
\begin{quote}\ttfamily\footnotesize\raggedright\noindent
[quotient(\allowbreak{}ringoid-setoid(\allowbreak{}r)), descend2(\allowbreak{}ringoid-setoid(\allowbreak{}r), vnb-lambda(\allowbreak{}[a, b], cartesian(\allowbreak{}carr(\allowbreak{}r), carr(\allowbreak{}r)), class(\allowbreak{}ringoid-setoid(\allowbreak{}r), (\allowbreak{}add(\allowbreak{}r))(\allowbreak{}a, b)))), descend2(\allowbreak{}ringoid-setoid(\allowbreak{}r), vnb-lambda(\allowbreak{}[a, b], cartesian(\allowbreak{}carr(\allowbreak{}r), carr(\allowbreak{}r)), class(\allowbreak{}ringoid-setoid(\allowbreak{}r), (\allowbreak{}mul(\allowbreak{}r))(\allowbreak{}a, b)))), descend(\allowbreak{}ringoid-setoid(\allowbreak{}r), vnb-lambda(\allowbreak{}a, carr(\allowbreak{}r), class(\allowbreak{}ringoid-setoid(\allowbreak{}r), (\allowbreak{}neg(\allowbreak{}r))(\allowbreak{}a)))), class(\allowbreak{}ringoid-setoid(\allowbreak{}r), zero(\allowbreak{}r)), class(\allowbreak{}ringoid-setoid(\allowbreak{}r), one(\allowbreak{}r))]
\end{quote}
mentioned by 22 result(s)

\smallskip

\noindent\textbf{\texttt{ringoid-rel(r)}}\quad\textit{functoid}\index{ringoid-rel(r)@\texttt{ringoid-rel(r)}}\\
definition:
\begin{quote}\ttfamily\footnotesize\raggedright\noindent
\{p in cartesian(\allowbreak{}carr(\allowbreak{}r), carr(\allowbreak{}r)): (\allowbreak{}add(\allowbreak{}r))(\allowbreak{}nth(\allowbreak{}1, p), (\allowbreak{}neg(\allowbreak{}r))(\allowbreak{}nth(\allowbreak{}2, p))) in idl(\allowbreak{}r)\}
\end{quote}
characterized by \texttt{ringoid-rel-mem}  \textit{(definitional)}:
\begin{quote}\ttfamily\footnotesize\raggedright\noindent
forall(\allowbreak{}[r, a, b], [a, b] in ringoid-rel(\allowbreak{}r) iff [a, b] in cartesian(\allowbreak{}carr(\allowbreak{}r), carr(\allowbreak{}r)) and (\allowbreak{}add(\allowbreak{}r))(\allowbreak{}a, (\allowbreak{}neg(\allowbreak{}r))(\allowbreak{}b)) in idl(\allowbreak{}r))
\end{quote}
also: \texttt{ringoid-rel-is-equivalence} \\
mentioned by 7 result(s)

\smallskip

\noindent\textbf{\texttt{ringoid-setoid(r)}}\quad\textit{functoid}\index{ringoid-setoid(r)@\texttt{ringoid-setoid(r)}}\\
definition:
\begin{quote}\ttfamily\footnotesize\raggedright\noindent
[carr(\allowbreak{}r), ringoid-rel(\allowbreak{}r)]
\end{quote}
characterized by \texttt{ringoid-setoid-pts}  \textit{(definitional)}:
\begin{quote}\ttfamily\footnotesize\raggedright\noindent
forall(\allowbreak{}[r], pts(\allowbreak{}ringoid-setoid(\allowbreak{}r)) == carr(\allowbreak{}r))
\end{quote}
also: \texttt{ringoid-setoid-is-setoid} \\
mentioned by 44 result(s)

\smallskip

\noindent\textbf{\texttt{rintegral(omega, ca, mu, f)}}\quad\textit{functoid}\index{rintegral(omega, ca, mu, f)@\texttt{rintegral(omega, ca, mu, f)}}\\
reads: the integral of the real function \$4 over \$1 with respect to \$3 \\
definition:
\begin{quote}\ttfamily\footnotesize\raggedright\noindent
integral(\allowbreak{}omega, ca, mu, pos-part(\allowbreak{}omega, f)) - integral(\allowbreak{}omega, ca, mu, neg-part(\allowbreak{}omega, f))
\end{quote}
characterized by \texttt{rintegral-add}  \textit{(proven)}:
\begin{quote}\ttfamily\footnotesize\raggedright\noindent
forall(\allowbreak{}[omega, ca, mu, f, g], is-measure(\allowbreak{}omega, ca, mu) implies is-integrable(\allowbreak{}omega, ca, mu, f) implies is-integrable(\allowbreak{}omega, ca, mu, g) implies is-integrable(\allowbreak{}omega, ca, mu, vnb-lambda(\allowbreak{}sx\_, omega, f(\allowbreak{}sx\_) + g(\allowbreak{}sx\_))) and rintegral(\allowbreak{}omega, ca, mu, vnb-lambda(\allowbreak{}sx\_, omega, f(\allowbreak{}sx\_) + g(\allowbreak{}sx\_))) = rintegral(\allowbreak{}omega, ca, mu, f) + rintegral(\allowbreak{}omega, ca, mu, g))
\end{quote}
also: \texttt{rintegral-ae-eq} \texttt{rintegral-in-rr} \texttt{rintegral-scale} \texttt{rintegral-abs-le} \texttt{rintegral-unfold} \texttt{weight-rintegral} \texttt{rintegral-density} \texttt{rintegral-zero-fn} ... (19 in all) \\
mentioned by 66 result(s)

\smallskip

\noindent\textbf{\texttt{rip(s)}}\quad\textit{accessor}\index{rip(s)@\texttt{rip(s)}}\\
reads: the form of \$1 \\
a structure slot \\
characterized by \texttt{sipq-rip-csip-as-sip}  \textit{(definitional)}:
\begin{quote}\ttfamily\footnotesize\raggedright\noindent
forall(\allowbreak{}[r], is-complex-semi-inner-product-space(\allowbreak{}r) implies rip(\allowbreak{}sip-quotient(\allowbreak{}csip-as-sip(\allowbreak{}r))) == sip-qrip(\allowbreak{}csip-as-sip(\allowbreak{}r)))
\end{quote}
also: \texttt{sip-rip-fun-csip-as-sip} \texttt{sip-rip-sym-csip-as-sip} \texttt{sip-rip-type-csip-as-sip} \texttt{sipq-rip-fun-csip-as-sip} \texttt{sipq-rip-sym-csip-as-sip} \texttt{sip-rip-add-left-csip-as-sip} \texttt{sip-rip-neg-left-csip-as-sip} \texttt{sip-rip-add-right-csip-as-sip} ... (69 in all) \\
mentioned by 297 result(s)

\smallskip

\noindent\textbf{\texttt{rm-address(s, dseq, es, y)}}\quad\textit{predicate}\index{rm-address(s, dseq, es, y)@\texttt{rm-address(s, dseq, es, y)}}\\
reads: \$4 has the address \$3 in \$1 \\
defined by \texttt{rm-address} \\
characterized by \texttt{rm-address}  \textit{(definitional)}:
\begin{quote}\ttfamily\footnotesize\raggedright\noindent
forall(\allowbreak{}[s, dseq, es, y], rm-address(\allowbreak{}s, dseq, es, y) iff y in pts(\allowbreak{}s) and forall(\allowbreak{}[ran\_ in nn], (\allowbreak{}es(\allowbreak{}ran\_) = 1 implies y in rm-ball(\allowbreak{}s, dseq, ran\_)) and (\allowbreak{}not(\allowbreak{}es(\allowbreak{}ran\_) = 1) implies y in rm-shell(\allowbreak{}s, dseq, ran\_))))
\end{quote}
mentioned by 12 result(s)

\smallskip

\noindent\textbf{\texttt{rm-address-below(s, dseq, es, nb, y)}}\quad\textit{predicate}\index{rm-address-below(s, dseq, es, nb, y)@\texttt{rm-address-below(s, dseq, es, nb, y)}}\\
reads: \$5 has the address \$3 in \$1 below \$4 \\
defined by \texttt{rm-address-below} \\
characterized by \texttt{rm-address-below}  \textit{(definitional)}:
\begin{quote}\ttfamily\footnotesize\raggedright\noindent
forall(\allowbreak{}[s, dseq, es, nb, y], rm-address-below(\allowbreak{}s, dseq, es, nb, y) iff y in pts(\allowbreak{}s) and forall(\allowbreak{}[rbn\_ in nn], rbn\_ < nb implies (\allowbreak{}es(\allowbreak{}rbn\_) = 1 implies y in rm-ball(\allowbreak{}s, dseq, rbn\_)) and (\allowbreak{}not(\allowbreak{}es(\allowbreak{}rbn\_) = 1) implies y in rm-shell(\allowbreak{}s, dseq, rbn\_))))
\end{quote}
mentioned by 5 result(s)

\smallskip

\noindent\textbf{\texttt{rm-ball(s, dseq, n)}}\quad\textit{functoid}\index{rm-ball(s, dseq, n)@\texttt{rm-ball(s, dseq, n)}}\\
reads: the \$3-th ball of the family of \$1 built on \$2 \\
definition:
\begin{quote}\ttfamily\footnotesize\raggedright\noindent
closed-ball(\allowbreak{}s, dseq(\allowbreak{}nnfst(\allowbreak{}n)), recip(\allowbreak{}nnsnd(\allowbreak{}n) + 1))
\end{quote}
mentioned by 13 result(s)

\smallskip

\noindent\textbf{\texttt{rm-code(s, dseq, x)}}\quad\textit{functoid}\index{rm-code(s, dseq, x)@\texttt{rm-code(s, dseq, x)}}\\
reads: the address of \$3 in the Cantor model of \$1 built on \$2 \\
definition:
\begin{quote}\ttfamily\footnotesize\raggedright\noindent
vnb-lambda(\allowbreak{}rox\_, nn, if(\allowbreak{}x in rm-ball(\allowbreak{}s, dseq, rox\_), 1, 0))
\end{quote}
mentioned by 8 result(s)

\smallskip

\noindent\textbf{\texttt{rm-cyl(s, dseq, a, nb)}}\quad\textit{functoid}\index{rm-cyl(s, dseq, a, nb)@\texttt{rm-cyl(s, dseq, a, nb)}}\\
reads: the cylinder of the Cantor model of \$1 on which the first \$4 coordinates are those of \$3 \\
definition:
\begin{quote}\ttfamily\footnotesize\raggedright\noindent
\{rce\_ in rm-k(\allowbreak{}s, dseq): forall(\allowbreak{}[rci\_ in nn], rci\_ < nb implies rce\_(\allowbreak{}rci\_) = a(\allowbreak{}rci\_))\}
\end{quote}
mentioned by 22 result(s)

\smallskip

\noindent\textbf{\texttt{rm-k(s, dseq)}}\quad\textit{functoid}\index{rm-k(s, dseq)@\texttt{rm-k(s, dseq)}}\\
reads: the Cantor model of \$1 built on \$2 \\
definition:
\begin{quote}\ttfamily\footnotesize\raggedright\noindent
\{rke\_ in pts(\allowbreak{}cantor-ms): forall(\allowbreak{}[rkn\_ in nn], forsome(\allowbreak{}[rky\_ in pts(\allowbreak{}s)], rm-realizes(\allowbreak{}s, dseq, rke\_, rkn\_, rky\_)))\}
\end{quote}
mentioned by 37 result(s)

\smallskip

\noindent\textbf{\texttt{rm-kms(s, dseq)}}\quad\textit{functoid}\index{rm-kms(s, dseq)@\texttt{rm-kms(s, dseq)}}\\
reads: the Cantor model of \$1 built on \$2, as a metric space \\
definition:
\begin{quote}\ttfamily\footnotesize\raggedright\noindent
subspace-ms(\allowbreak{}cantor-ms, rm-k(\allowbreak{}s, dseq))
\end{quote}
mentioned by 31 result(s)

\smallskip

\noindent\textbf{\texttt{rm-pi(s, dseq)}}\quad\textit{functoid}\index{rm-pi(s, dseq)@\texttt{rm-pi(s, dseq)}}\\
reads: the projection of the Cantor model of \$1 built on \$2 onto \$1 \\
definition:
\begin{quote}\ttfamily\footnotesize\raggedright\noindent
vnb-lambda(\allowbreak{}rpe\_, rm-k(\allowbreak{}s, dseq), iota(\allowbreak{}rpy\_, rpy\_ in pts(\allowbreak{}s) and rm-address(\allowbreak{}s, dseq, rpe\_, rpy\_)))
\end{quote}
mentioned by 12 result(s)

\smallskip

\noindent\textbf{\texttt{rm-realizes(s, dseq, es, nb, y)}}\quad\textit{predicate}\index{rm-realizes(s, dseq, es, nb, y)@\texttt{rm-realizes(s, dseq, es, nb, y)}}\\
reads: \$5 realizes the address \$3 of \$1 below \$4 \\
defined by \texttt{rm-realizes} \\
characterized by \texttt{rm-realizes}  \textit{(definitional)}:
\begin{quote}\ttfamily\footnotesize\raggedright\noindent
forall(\allowbreak{}[s, dseq, es, nb, y], rm-realizes(\allowbreak{}s, dseq, es, nb, y) iff y in pts(\allowbreak{}s) and forall(\allowbreak{}[rri\_ in nn], rri\_ < nb implies y in rm-ball(\allowbreak{}s, dseq, rri\_) iff es(\allowbreak{}rri\_) = 1))
\end{quote}
mentioned by 7 result(s)

\smallskip

\noindent\textbf{\texttt{rm-shell(s, dseq, n)}}\quad\textit{functoid}\index{rm-shell(s, dseq, n)@\texttt{rm-shell(s, dseq, n)}}\\
reads: the outer shell of the \$3-th ball of the family of \$1 built on \$2 \\
definition:
\begin{quote}\ttfamily\footnotesize\raggedright\noindent
\{rsy\_ in pts(\allowbreak{}s): recip(\allowbreak{}nnsnd(\allowbreak{}n) + 1) <= (\allowbreak{}dist(\allowbreak{}s))(\allowbreak{}dseq(\allowbreak{}nnfst(\allowbreak{}n)), rsy\_)\}
\end{quote}
mentioned by 9 result(s)

\smallskip

\noindent\textbf{\texttt{rm-two}}\quad\textit{constant}\index{rm-two@\texttt{rm-two}}\\
reads: the two-point set \{0, 1\} \\
defined by \texttt{rm-two-def}

\smallskip

\noindent\textbf{\texttt{rm-two-ms}}\quad\textit{constant}\index{rm-two-ms@\texttt{rm-two-ms}}\\
reads: the two-point metric space \{0, 1\} \\
defined by \texttt{rm-two-ms-def}

\smallskip

\noindent\textbf{\texttt{rn(n)}}\quad\textit{functoid}\index{rn(n)@\texttt{rn(n)}}\\
reads: the space of real \$1-tuples with the sup norm \\
definition:
\begin{quote}\ttfamily\footnotesize\raggedright\noindent
[normed-field-as-commutative-ring(\allowbreak{}rr-normed-field), fun(\allowbreak{}interval(\allowbreak{}1, n), rr), vnb-lambda(\allowbreak{}[rnax\_, rnay\_], cartesian(\allowbreak{}fun(\allowbreak{}interval(\allowbreak{}1, n), rr), fun(\allowbreak{}interval(\allowbreak{}1, n), rr)), vnb-lambda(\allowbreak{}rnai\_, interval(\allowbreak{}1, n), rnax\_(\allowbreak{}rnai\_) + rnay\_(\allowbreak{}rnai\_))), vnb-lambda(\allowbreak{}rnzi\_, interval(\allowbreak{}1, n), 0), vnb-lambda(\allowbreak{}rnnx\_, fun(\allowbreak{}interval(\allowbreak{}1, n), rr), vnb-lambda(\allowbreak{}rnni\_, interval(\allowbreak{}1, n), -rnnx\_(\allowbreak{}rnni\_))), vnb-lambda(\allowbreak{}[rnsr\_, rnsx\_], cartesian(\allowbreak{}rr, fun(\allowbreak{}interval(\allowbreak{}1, n), rr)), vnb-lambda(\allowbreak{}rnsi\_, interval(\allowbreak{}1, n), rnsr\_ * rnsx\_(\allowbreak{}rnsi\_))), vnb-lambda(\allowbreak{}rnvx\_, fun(\allowbreak{}interval(\allowbreak{}1, n), rr), rn-maxabs(\allowbreak{}rnvx\_, n))]
\end{quote}
characterized by \texttt{rn-vec}  \textit{(proven)}:
\begin{quote}\ttfamily\footnotesize\raggedright\noindent
forall(\allowbreak{}[n in nn], vec(\allowbreak{}rn(\allowbreak{}n)) == fun(\allowbreak{}interval(\allowbreak{}1, n), rr))
\end{quote}
also: \texttt{rn-scal} \texttt{rn-is-tvs} \texttt{rn-act-add} \texttt{rn-act-mul} \texttt{rn-act-one} \texttt{rn-norm-le} \texttt{rn-act-vadd} \texttt{rn-act-apply} ... (35 in all) \\
mentioned by 63 result(s)

\smallskip

\noindent\textbf{\texttt{rn-box(n, jj)}}\quad\textit{functoid}\index{rn-box(n, jj)@\texttt{rn-box(n, jj)}}\\
reads: the rectangle in R\textasciicircum{}\$1 with sides \$2 \\
definition:
\begin{quote}\ttfamily\footnotesize\raggedright\noindent
\{mrx\_ in fun(\allowbreak{}interval(\allowbreak{}1, n), rr): forall(\allowbreak{}[mri\_ in interval(\allowbreak{}1, n)], mrx\_(\allowbreak{}mri\_) in jj(\allowbreak{}mri\_))\}
\end{quote}
characterized by \texttt{rn-box-unfold}  \textit{(proven)}:
\begin{quote}\ttfamily\footnotesize\raggedright\noindent
forall(\allowbreak{}[n, jj], rn-box(\allowbreak{}n, jj) == \{mrx\_ in fun(\allowbreak{}interval(\allowbreak{}1, n), rr): forall(\allowbreak{}[mri\_ in interval(\allowbreak{}1, n)], mrx\_(\allowbreak{}mri\_) in jj(\allowbreak{}mri\_))\})
\end{quote}
also: \texttt{rn-box-in-power} \texttt{rn-box-membership} \\
mentioned by 12 result(s)

\smallskip

\noindent\textbf{\texttt{rn-maxabs}}\quad\textit{functoid}\index{rn-maxabs@\texttt{rn-maxabs}}\\
reads: the largest of the absolute values of the first \$2 coordinates of \$1 \\
defined by \texttt{rn-maxabs-zero} \texttt{rn-maxabs-succ} \\
characterized by \texttt{rn-maxabs-succ}  \textit{(definitional)}:
\begin{quote}\ttfamily\footnotesize\raggedright\noindent
forall(\allowbreak{}[rnmx\_, rnmk\_ in nn], rn-maxabs(\allowbreak{}rnmx\_, succ(\allowbreak{}rnmk\_)) == max(\allowbreak{}rn-maxabs(\allowbreak{}rnmx\_, rnmk\_), abs(\allowbreak{}rnmx\_(\allowbreak{}succ(\allowbreak{}rnmk\_)))))
\end{quote}
also: \texttt{rn-maxabs-zero} \texttt{rn-maxabs-ge} \texttt{rn-maxabs-le} \texttt{rn-maxabs-scale} \texttt{rn-maxabs-typing} \\
mentioned by 11 result(s)

\smallskip

\noindent\textbf{\texttt{rn-rectangles(n)}}\quad\textit{functoid}\index{rn-rectangles(n)@\texttt{rn-rectangles(n)}}\\
reads: the rectangles (cartesian products of intervals) in R\textasciicircum{}\$1 \\
definition:
\begin{quote}\ttfamily\footnotesize\raggedright\noindent
\{mrr\_ in power(\allowbreak{}fun(\allowbreak{}interval(\allowbreak{}1, n), rr)): forsome(\allowbreak{}[mrj\_ in fun(\allowbreak{}interval(\allowbreak{}1, n), power(\allowbreak{}rr))], forall(\allowbreak{}[mri\_ in interval(\allowbreak{}1, n)], is-interval(\allowbreak{}mrj\_(\allowbreak{}mri\_))) and mrr\_ = rn-box(\allowbreak{}n, mrj\_))\}
\end{quote}
characterized by \texttt{rn-rectangles-unfold}  \textit{(proven)}:
\begin{quote}\ttfamily\footnotesize\raggedright\noindent
forall(\allowbreak{}[n], rn-rectangles(\allowbreak{}n) == \{mrr\_ in power(\allowbreak{}fun(\allowbreak{}interval(\allowbreak{}1, n), rr)): forsome(\allowbreak{}[mrj\_ in fun(\allowbreak{}interval(\allowbreak{}1, n), power(\allowbreak{}rr))], forall(\allowbreak{}[mri\_ in interval(\allowbreak{}1, n)], is-interval(\allowbreak{}mrj\_(\allowbreak{}mri\_))) and mrr\_ = rn-box(\allowbreak{}n, mrj\_))\})
\end{quote}
also: \texttt{rn-rectangles-inter-closed} \texttt{rn-rectangles-subset-power} \texttt{rn-rectangles-finite-unions-algebra} \\
mentioned by 10 result(s)

\smallskip

\noindent\textbf{\texttt{rn-trunc(n, x, k)}}\quad\textit{functoid}\index{rn-trunc(n, x, k)@\texttt{rn-trunc(n, x, k)}}\\
reads: the vector of the first \$3 coordinates of \$2 in the space of real \$1-tuples \\
definition:
\begin{quote}\ttfamily\footnotesize\raggedright\noindent
vnb-lambda(\allowbreak{}rntj\_, interval(\allowbreak{}1, n), if(\allowbreak{}rntj\_ <= k, x(\allowbreak{}rntj\_), 0))
\end{quote}
characterized by \texttt{rn-trunc-full}  \textit{(proven)}:
\begin{quote}\ttfamily\footnotesize\raggedright\noindent
forall(\allowbreak{}[n, x], n in nn implies x in fun(\allowbreak{}interval(\allowbreak{}1, n), rr) implies rn-trunc(\allowbreak{}n, x, n) = x)
\end{quote}
also: \texttt{rn-trunc-succ} \texttt{rn-trunc-zero} \texttt{rn-trunc-in-vec} \texttt{rn-trunc-apply-gt} \texttt{rn-trunc-apply-le} \texttt{rn-trunc-step-norm} \\
mentioned by 13 result(s)

\smallskip

\noindent\textbf{\texttt{rn-unit(n, i)}}\quad\textit{functoid}\index{rn-unit(n, i)@\texttt{rn-unit(n, i)}}\\
reads: the unit basis vector \$2 of the space of real \$1-tuples \\
definition:
\begin{quote}\ttfamily\footnotesize\raggedright\noindent
vnb-lambda(\allowbreak{}rnuj\_, interval(\allowbreak{}1, n), if(\allowbreak{}rnuj\_ = i, 1, 0))
\end{quote}
characterized by \texttt{rn-unit-in-vec}  \textit{(proven)}:
\begin{quote}\ttfamily\footnotesize\raggedright\noindent
forall(\allowbreak{}[n, i], n in nn implies i in interval(\allowbreak{}1, n) implies rn-unit(\allowbreak{}n, i) in fun(\allowbreak{}interval(\allowbreak{}1, n), rr))
\end{quote}
also: \texttt{rn-unit-at-same} \texttt{rn-unit-at-other} \\
mentioned by 12 result(s)

\smallskip

\noindent\textbf{\texttt{road-length(rld\_, rla\_, rlb\_)}}\quad\textit{functoid}\index{road-length(rld_, rla_, rlb_)@\texttt{road-length(rld\_, rla\_, rlb\_)}}\\
reads: the length of the road with derivative \$1 on [\$2, \$3] \\
definition:
\begin{quote}\ttfamily\footnotesize\raggedright\noindent
pw-int(\allowbreak{}vnb-lambda(\allowbreak{}rlt\_, ccint(\allowbreak{}rla\_, rlb\_), magnitude(\allowbreak{}rld\_(\allowbreak{}rlt\_))), rla\_, rlb\_)
\end{quote}
characterized by \texttt{road-length-tri}  \textit{(proven)}:
\begin{quote}\ttfamily\footnotesize\raggedright\noindent
forall(\allowbreak{}[a in cc, b in cc, c in cc], road-length(\allowbreak{}tri-road-deriv(\allowbreak{}a, b, c), 0, 3) = (\allowbreak{}magnitude(\allowbreak{}b - a) + magnitude(\allowbreak{}c - b)) + magnitude(\allowbreak{}a - c))
\end{quote}
also: \texttt{road-length-in-rr} \texttt{road-length-juxta} \texttt{circle-road-length} \texttt{road-length-nonneg} \texttt{road-length-unfold} \texttt{road-length-segment} \texttt{road-length-affine-reparam} \\
mentioned by 19 result(s)

\smallskip

\noindent\textbf{\texttt{rpow}}\quad\textit{operator}\index{rpow@\texttt{rpow}}\\
a term-forming head \\
characterized by \texttt{rpow-def}  \textit{(definitional)}:
\begin{quote}\ttfamily\footnotesize\raggedright\noindent
forall(\allowbreak{}[a\_, b\_], rpow(\allowbreak{}a\_, b\_) == rpow-star(\allowbreak{}a\_, b\_))
\end{quote}
also: \texttt{rpow-add} \texttt{rpow-nat} \texttt{rpow-neg} \texttt{rpow-one} \texttt{rpow-pos} \texttt{rpow-pow} \texttt{rpow-zero} \texttt{sqrt-rpow} ... (18 in all) \\
mentioned by 41 result(s)

\smallskip

\noindent\textbf{\texttt{rpow-star(x\_, s\_)}}\quad\textit{functoid}\index{rpow-star(x_, s_)@\texttt{rpow-star(x\_, s\_)}}\\
reads: \$1 to the power \$2 \\
definition:
\begin{quote}\ttfamily\footnotesize\raggedright\noindent
if(\allowbreak{}0 < x\_, r-exp(\allowbreak{}s\_ * log(\allowbreak{}x\_)), if(\allowbreak{}0 < s\_, 0, 1))
\end{quote}
characterized by \texttt{log-rpow-star}  \textit{(proven)}:
\begin{quote}\ttfamily\footnotesize\raggedright\noindent
forall(\allowbreak{}[x\_, s\_ in rr], 0 < x\_ implies log(\allowbreak{}rpow-star(\allowbreak{}x\_, s\_)) = s\_ * log(\allowbreak{}x\_))
\end{quote}
also: \texttt{rpow-star-add} \texttt{rpow-star-nat} \texttt{rpow-star-neg} \texttt{rpow-star-one} \texttt{rpow-star-pos} \texttt{rpow-star-pow} \texttt{rpow-star-zero} \texttt{rpow-star-in-rr} ... (23 in all) \\
mentioned by 121 result(s)

\smallskip

\noindent\textbf{\texttt{rr-bounded-above(s)}}\quad\textit{predicate}\index{rr-bounded-above(s)@\texttt{rr-bounded-above(s)}}\\
reads: \$1 is bounded above \\
defined by \texttt{rr-bounded-above} \\
characterized by \texttt{rr-bounded-above}  \textit{(definitional)}:
\begin{quote}\ttfamily\footnotesize\raggedright\noindent
forall(\allowbreak{}[s], rr-bounded-above(\allowbreak{}s) iff forsome(\allowbreak{}[b], rr-upper-bound(\allowbreak{}s, b)))
\end{quote}
mentioned by 7 result(s)

\smallskip

\noindent\textbf{\texttt{rr-bounded-ms}}\quad\textit{defined-fn}\index{rr-bounded-ms@\texttt{rr-bounded-ms}}\\
defined by \texttt{rr-bounded-ms-def}

\smallskip

\noindent\textbf{\texttt{rr-ext-down}}\quad\textit{constant}\index{rr-ext-down@\texttt{rr-ext-down}}\\
reads: the real line with -infinity adjoined, [-inf, +inf) \\
defined by \texttt{rr-ext-down-def}

\smallskip

\noindent\textbf{\texttt{rr-ext-up}}\quad\textit{constant}\index{rr-ext-up@\texttt{rr-ext-up}}\\
reads: the real line with +infinity adjoined, (-inf, +inf] \\
defined by \texttt{rr-ext-up-def}

\smallskip

\noindent\textbf{\texttt{rr-field}}\quad\textit{instance}\index{rr-field@\texttt{rr-field}}\\
an instance of field; predicate \texttt{is-rr-field}; same shape as field

\smallskip

\noindent\textbf{\texttt{rr-fn-space(rfv\_)}}\quad\textit{functoid}\index{rr-fn-space(rfv_)@\texttt{rr-fn-space(rfv\_)}}\\
reads: the real vector space of the functions in \$1 under the pointwise operations \\
definition:
\begin{quote}\ttfamily\footnotesize\raggedright\noindent
[field-as-integral-domain(\allowbreak{}rr-field), rfv\_, vnb-lambda(\allowbreak{}[rff\_, rfg\_], cartesian(\allowbreak{}rfv\_, rfv\_), vnb-lambda(\allowbreak{}rfx\_, rr, rff\_(\allowbreak{}rfx\_) + rfg\_(\allowbreak{}rfx\_))), vnb-lambda(\allowbreak{}rfx\_, rr, 0), vnb-lambda(\allowbreak{}rff\_, rfv\_, vnb-lambda(\allowbreak{}rfx\_, rr, -rff\_(\allowbreak{}rfx\_))), vnb-lambda(\allowbreak{}[rfc\_, rff\_], cartesian(\allowbreak{}rr, rfv\_), vnb-lambda(\allowbreak{}rfx\_, rr, rfc\_ * rff\_(\allowbreak{}rfx\_)))]
\end{quote}
characterized by \texttt{rr-fn-space-act}  \textit{(proven)}:
\begin{quote}\ttfamily\footnotesize\raggedright\noindent
forall(\allowbreak{}[v in set], v subset fun(\allowbreak{}rr, rr) implies act(\allowbreak{}rr-fn-space(\allowbreak{}v)) == vnb-lambda(\allowbreak{}[rfc\_, rff\_], cartesian(\allowbreak{}rr, v), vnb-lambda(\allowbreak{}rfx\_, rr, rfc\_ * rff\_(\allowbreak{}rfx\_))))
\end{quote}
also: \texttt{rr-fn-space-vec} \texttt{rr-fn-space-scal} \texttt{rr-fn-space-vadd} \texttt{rr-fn-space-vneg} \texttt{rr-fn-space-vzero} \texttt{rr-fn-space-is-vector-space} \\
mentioned by 14 result(s)

\smallskip

\noindent\textbf{\texttt{rr-ms}}\quad\textit{instance}\index{rr-ms@\texttt{rr-ms}}\\
an instance of metric-space; predicate \texttt{is-rr-ms}; same shape as metric-space

\smallskip

\noindent\textbf{\texttt{rr-normed-field}}\quad\textit{instance}\index{rr-normed-field@\texttt{rr-normed-field}}\\
an instance of normed-field; predicate \texttt{is-rr-normed-field}; same shape as normed-field

\smallskip

\noindent\textbf{\texttt{rr-nvs}}\quad\textit{instance}\index{rr-nvs@\texttt{rr-nvs}}\\
an instance of normed-vector-space; predicate \texttt{is-rr-nvs}; same shape as normed-vector-space

\smallskip

\noindent\textbf{\texttt{rr-pos-star-add-monoid}}\quad\textit{instance}\index{rr-pos-star-add-monoid@\texttt{rr-pos-star-add-monoid}}\\
an instance of comm-monoid; predicate \texttt{is-rr-pos-star-add-monoid}; same shape as comm-monoid

\smallskip

\noindent\textbf{\texttt{rr-upper-bound(s, b)}}\quad\textit{predicate}\index{rr-upper-bound(s, b)@\texttt{rr-upper-bound(s, b)}}\\
reads: \$2 is an upper bound of \$1 \\
defined by \texttt{rr-upper-bound} \\
characterized by \texttt{rr-upper-bound}  \textit{(definitional)}:
\begin{quote}\ttfamily\footnotesize\raggedright\noindent
forall(\allowbreak{}[s, b], rr-upper-bound(\allowbreak{}s, b) iff b in rr and forall(\allowbreak{}[x in s], x <= b))
\end{quote}
mentioned by 9 result(s)

\smallskip

\noindent\textbf{\texttt{rs}}\quad\textit{tactic}\index{rs@\texttt{rs}}\\
\texttt{(rs)} -- Ring-simplify the goal (normal form over the ambient ring). \\
declared in \texttt{tactics-help.scm}

\smallskip

\noindent\textbf{\texttt{rw-all(mrt\_, mrs\_, mrx\_)}}\quad\textit{predicate}\index{rw-all(mrt_, mrs_, mrx_)@\texttt{rw-all(mrt\_, mrs\_, mrx\_)}}\\
reads: the rewrite clause FORALL relates \$2 to \$3 \\
defined by \texttt{rw-all} \\
characterized by \texttt{rw-all}  \textit{(definitional)}:
\begin{quote}\ttfamily\footnotesize\raggedright\noindent
forall(\allowbreak{}[mrt\_, mrs\_, mrx\_], rw-all(\allowbreak{}mrt\_, mrs\_, mrx\_) iff forsome(\allowbreak{}[mrg\_, mra\_, mrb\_], mrx\_ = [2, mrg\_, enc-all(\allowbreak{}mra\_), enc-all(\allowbreak{}mrb\_)] and enc-all(\allowbreak{}mra\_) in clf and enc-all(\allowbreak{}mrb\_) in clf and forsome(\allowbreak{}[mrfs\_ in set], mrfs\_ subset nn and card(\allowbreak{}mrfs\_) in nn and forall(\allowbreak{}[mry\_ in nn], not(\allowbreak{}mry\_ in mrfs\_) implies [2, mrg\_, expr-open(\allowbreak{}mra\_, 0, efvar(\allowbreak{}mry\_)), expr-open(\allowbreak{}mrb\_, 0, efvar(\allowbreak{}mry\_))] in mrs\_))))
\end{quote}
mentioned by 5 result(s)

\smallskip

\noindent\textbf{\texttt{rw-and-l(mrt\_, mrs\_, mrx\_)}}\quad\textit{predicate}\index{rw-and-l(mrt_, mrs_, mrx_)@\texttt{rw-and-l(mrt\_, mrs\_, mrx\_)}}\\
reads: the rewrite clause AND-left relates \$2 to \$3 \\
defined by \texttt{rw-and-l} \\
characterized by \texttt{rw-and-l}  \textit{(definitional)}:
\begin{quote}\ttfamily\footnotesize\raggedright\noindent
forall(\allowbreak{}[mrt\_, mrs\_, mrx\_], rw-and-l(\allowbreak{}mrt\_, mrs\_, mrx\_) iff forsome(\allowbreak{}[mrg\_, mra\_, mrb\_, mrc\_], mrx\_ = [2, mrg\_, enc-and(\allowbreak{}mra\_, mrc\_), enc-and(\allowbreak{}mrb\_, mrc\_)] and mrc\_ in clf and [2, ctx-add(\allowbreak{}mrg\_, mrc\_), mra\_, mrb\_] in mrs\_))
\end{quote}
mentioned by 5 result(s)

\smallskip

\noindent\textbf{\texttt{rw-and-r(mrt\_, mrs\_, mrx\_)}}\quad\textit{predicate}\index{rw-and-r(mrt_, mrs_, mrx_)@\texttt{rw-and-r(mrt\_, mrs\_, mrx\_)}}\\
reads: the rewrite clause AND-right relates \$2 to \$3 \\
defined by \texttt{rw-and-r} \\
characterized by \texttt{rw-and-r}  \textit{(definitional)}:
\begin{quote}\ttfamily\footnotesize\raggedright\noindent
forall(\allowbreak{}[mrt\_, mrs\_, mrx\_], rw-and-r(\allowbreak{}mrt\_, mrs\_, mrx\_) iff forsome(\allowbreak{}[mrg\_, mra\_, mrb\_, mrc\_], mrx\_ = [2, mrg\_, enc-and(\allowbreak{}mrc\_, mra\_), enc-and(\allowbreak{}mrc\_, mrb\_)] and mrc\_ in clf and [2, ctx-add(\allowbreak{}mrg\_, mrc\_), mra\_, mrb\_] in mrs\_))
\end{quote}
mentioned by 5 result(s)

\smallskip

\noindent\textbf{\texttt{rw-bind-dom(mrt\_, mrs\_, mrx\_)}}\quad\textit{predicate}\index{rw-bind-dom(mrt_, mrs_, mrx_)@\texttt{rw-bind-dom(mrt\_, mrs\_, mrx\_)}}\\
reads: the rewrite clause BINDER-domain relates \$2 to \$3 \\
defined by \texttt{rw-bind-dom} \\
characterized by \texttt{rw-bind-dom}  \textit{(definitional)}:
\begin{quote}\ttfamily\footnotesize\raggedright\noindent
forall(\allowbreak{}[mrt\_, mrs\_, mrx\_], rw-bind-dom(\allowbreak{}mrt\_, mrs\_, mrx\_) iff forsome(\allowbreak{}[mrg\_, mrh\_, mra\_, mrb\_, mrc\_], mrx\_ = [1, mrg\_, ebind(\allowbreak{}mrh\_, 1, [mra\_], mrc\_), ebind(\allowbreak{}mrh\_, 1, [mrb\_], mrc\_)] and (\allowbreak{}mrh\_ = 22 or mrh\_ = 24 or mrh\_ = 26) and ebind(\allowbreak{}mrh\_, 1, [mra\_], mrc\_) in clt and ebind(\allowbreak{}mrh\_, 1, [mrb\_], mrc\_) in clt and [1, mrg\_, mra\_, mrb\_] in mrs\_))
\end{quote}
mentioned by 5 result(s)

\smallskip

\noindent\textbf{\texttt{rw-closed(mrt\_, mrk\_, mrs\_)}}\quad\textit{predicate}\index{rw-closed(mrt_, mrk_, mrs_)@\texttt{rw-closed(mrt\_, mrk\_, mrs\_)}}\\
reads: \$3 is closed under the rewrite clauses of stage at most \$2 justified by \$1 \\
defined by \texttt{rw-closed} \\
characterized by \texttt{rw-closed}  \textit{(definitional)}:
\begin{quote}\ttfamily\footnotesize\raggedright\noindent
forall(\allowbreak{}[mrt\_, mrk\_, mrs\_], rw-closed(\allowbreak{}mrt\_, mrk\_, mrs\_) iff forall(\allowbreak{}[mrx\_], rw-rule(\allowbreak{}mrt\_, mrk\_, mrs\_, mrx\_) implies mrx\_ in rw-univ implies mrx\_ in mrs\_))
\end{quote}
mentioned by 12 result(s)

\smallskip

\noindent\textbf{\texttt{rw-desc-body(mrt\_, mrs\_, mrx\_)}}\quad\textit{predicate}\index{rw-desc-body(mrt_, mrs_, mrx_)@\texttt{rw-desc-body(mrt\_, mrs\_, mrx\_)}}\\
reads: the rewrite clause COMP / IOTA-body relates \$2 to \$3 \\
defined by \texttt{rw-desc-body} \\
characterized by \texttt{rw-desc-body}  \textit{(definitional)}:
\begin{quote}\ttfamily\footnotesize\raggedright\noindent
forall(\allowbreak{}[mrt\_, mrs\_, mrx\_], rw-desc-body(\allowbreak{}mrt\_, mrs\_, mrx\_) iff forsome(\allowbreak{}[mrg\_, mrh\_, mra\_, mrb\_], mrx\_ = [1, mrg\_, ebind(\allowbreak{}mrh\_, 1, [], mra\_), ebind(\allowbreak{}mrh\_, 1, [], mrb\_)] and (\allowbreak{}mrh\_ = 23 or mrh\_ = 25) and ebind(\allowbreak{}mrh\_, 1, [], mra\_) in clt and ebind(\allowbreak{}mrh\_, 1, [], mrb\_) in clt and forsome(\allowbreak{}[mrfs\_ in set], mrfs\_ subset nn and card(\allowbreak{}mrfs\_) in nn and forall(\allowbreak{}[mry\_ in nn], not(\allowbreak{}mry\_ in mrfs\_) implies [2, mrg\_, expr-open(\allowbreak{}mra\_, 0, efvar(\allowbreak{}mry\_)), expr-open(\allowbreak{}mrb\_, 0, efvar(\allowbreak{}mry\_))] in mrs\_))))
\end{quote}
mentioned by 6 result(s)

\smallskip

\noindent\textbf{\texttt{rw-ex(mrt\_, mrs\_, mrx\_)}}\quad\textit{predicate}\index{rw-ex(mrt_, mrs_, mrx_)@\texttt{rw-ex(mrt\_, mrs\_, mrx\_)}}\\
reads: the rewrite clause FORSOME relates \$2 to \$3 \\
defined by \texttt{rw-ex} \\
characterized by \texttt{rw-ex}  \textit{(definitional)}:
\begin{quote}\ttfamily\footnotesize\raggedright\noindent
forall(\allowbreak{}[mrt\_, mrs\_, mrx\_], rw-ex(\allowbreak{}mrt\_, mrs\_, mrx\_) iff forsome(\allowbreak{}[mrg\_, mra\_, mrb\_], mrx\_ = [2, mrg\_, enc-ex(\allowbreak{}mra\_), enc-ex(\allowbreak{}mrb\_)] and enc-ex(\allowbreak{}mra\_) in clf and enc-ex(\allowbreak{}mrb\_) in clf and forsome(\allowbreak{}[mrfs\_ in set], mrfs\_ subset nn and card(\allowbreak{}mrfs\_) in nn and forall(\allowbreak{}[mry\_ in nn], not(\allowbreak{}mry\_ in mrfs\_) implies [2, mrg\_, expr-open(\allowbreak{}mra\_, 0, efvar(\allowbreak{}mry\_)), expr-open(\allowbreak{}mrb\_, 0, efvar(\allowbreak{}mry\_))] in mrs\_))))
\end{quote}
mentioned by 5 result(s)

\smallskip

\noindent\textbf{\texttt{rw-if-cond(mrt\_, mrs\_, mrx\_)}}\quad\textit{predicate}\index{rw-if-cond(mrt_, mrs_, mrx_)@\texttt{rw-if-cond(mrt\_, mrs\_, mrx\_)}}\\
reads: the rewrite clause IF-condition relates \$2 to \$3 \\
defined by \texttt{rw-if-cond} \\
characterized by \texttt{rw-if-cond}  \textit{(definitional)}:
\begin{quote}\ttfamily\footnotesize\raggedright\noindent
forall(\allowbreak{}[mrt\_, mrs\_, mrx\_], rw-if-cond(\allowbreak{}mrt\_, mrs\_, mrx\_) iff forsome(\allowbreak{}[mrg\_, mra\_, mrb\_, mrc\_, mrd\_], mrx\_ = [1, mrg\_, enc-if(\allowbreak{}mra\_, mrc\_, mrd\_), enc-if(\allowbreak{}mrb\_, mrc\_, mrd\_)] and enc-if(\allowbreak{}mra\_, mrc\_, mrd\_) in clt and enc-if(\allowbreak{}mrb\_, mrc\_, mrd\_) in clt and mrc\_ in clt and mrd\_ in clt and [2, mrg\_, mra\_, mrb\_] in mrs\_))
\end{quote}
mentioned by 5 result(s)

\smallskip

\noindent\textbf{\texttt{rw-if-else(mrt\_, mrs\_, mrx\_)}}\quad\textit{predicate}\index{rw-if-else(mrt_, mrs_, mrx_)@\texttt{rw-if-else(mrt\_, mrs\_, mrx\_)}}\\
reads: the rewrite clause IF-else relates \$2 to \$3 \\
defined by \texttt{rw-if-else} \\
characterized by \texttt{rw-if-else}  \textit{(definitional)}:
\begin{quote}\ttfamily\footnotesize\raggedright\noindent
forall(\allowbreak{}[mrt\_, mrs\_, mrx\_], rw-if-else(\allowbreak{}mrt\_, mrs\_, mrx\_) iff forsome(\allowbreak{}[mrg\_, mra\_, mrb\_, mrc\_, mrd\_], mrx\_ = [1, mrg\_, enc-if(\allowbreak{}mrc\_, mrd\_, mra\_), enc-if(\allowbreak{}mrc\_, mrd\_, mrb\_)] and enc-if(\allowbreak{}mrc\_, mrd\_, mra\_) in clt and enc-if(\allowbreak{}mrc\_, mrd\_, mrb\_) in clt and mrc\_ in clf and mrd\_ in clt and [1, mrg\_, mra\_, mrb\_] in mrs\_))
\end{quote}
mentioned by 5 result(s)

\smallskip

\noindent\textbf{\texttt{rw-if-then(mrt\_, mrs\_, mrx\_)}}\quad\textit{predicate}\index{rw-if-then(mrt_, mrs_, mrx_)@\texttt{rw-if-then(mrt\_, mrs\_, mrx\_)}}\\
reads: the rewrite clause IF-then relates \$2 to \$3 \\
defined by \texttt{rw-if-then} \\
characterized by \texttt{rw-if-then}  \textit{(definitional)}:
\begin{quote}\ttfamily\footnotesize\raggedright\noindent
forall(\allowbreak{}[mrt\_, mrs\_, mrx\_], rw-if-then(\allowbreak{}mrt\_, mrs\_, mrx\_) iff forsome(\allowbreak{}[mrg\_, mra\_, mrb\_, mrc\_, mrd\_], mrx\_ = [1, mrg\_, enc-if(\allowbreak{}mrc\_, mra\_, mrd\_), enc-if(\allowbreak{}mrc\_, mrb\_, mrd\_)] and enc-if(\allowbreak{}mrc\_, mra\_, mrd\_) in clt and enc-if(\allowbreak{}mrc\_, mrb\_, mrd\_) in clt and mrc\_ in clf and mrd\_ in clt and [1, mrg\_, mra\_, mrb\_] in mrs\_))
\end{quote}
mentioned by 5 result(s)

\smallskip

\noindent\textbf{\texttt{rw-iff-l(mrt\_, mrs\_, mrx\_)}}\quad\textit{predicate}\index{rw-iff-l(mrt_, mrs_, mrx_)@\texttt{rw-iff-l(mrt\_, mrs\_, mrx\_)}}\\
reads: the rewrite clause IFF-left relates \$2 to \$3 \\
defined by \texttt{rw-iff-l} \\
characterized by \texttt{rw-iff-l}  \textit{(definitional)}:
\begin{quote}\ttfamily\footnotesize\raggedright\noindent
forall(\allowbreak{}[mrt\_, mrs\_, mrx\_], rw-iff-l(\allowbreak{}mrt\_, mrs\_, mrx\_) iff forsome(\allowbreak{}[mrg\_, mra\_, mrb\_, mrc\_], mrx\_ = [2, mrg\_, enc-iff(\allowbreak{}mra\_, mrc\_), enc-iff(\allowbreak{}mrb\_, mrc\_)] and mrc\_ in clf and [2, mrg\_, mra\_, mrb\_] in mrs\_))
\end{quote}
mentioned by 5 result(s)

\smallskip

\noindent\textbf{\texttt{rw-iff-r(mrt\_, mrs\_, mrx\_)}}\quad\textit{predicate}\index{rw-iff-r(mrt_, mrs_, mrx_)@\texttt{rw-iff-r(mrt\_, mrs\_, mrx\_)}}\\
reads: the rewrite clause IFF-right relates \$2 to \$3 \\
defined by \texttt{rw-iff-r} \\
characterized by \texttt{rw-iff-r}  \textit{(definitional)}:
\begin{quote}\ttfamily\footnotesize\raggedright\noindent
forall(\allowbreak{}[mrt\_, mrs\_, mrx\_], rw-iff-r(\allowbreak{}mrt\_, mrs\_, mrx\_) iff forsome(\allowbreak{}[mrg\_, mra\_, mrb\_, mrc\_], mrx\_ = [2, mrg\_, enc-iff(\allowbreak{}mrc\_, mra\_), enc-iff(\allowbreak{}mrc\_, mrb\_)] and mrc\_ in clf and [2, mrg\_, mra\_, mrb\_] in mrs\_))
\end{quote}
mentioned by 5 result(s)

\smallskip

\noindent\textbf{\texttt{rw-imp-l(mrt\_, mrs\_, mrx\_)}}\quad\textit{predicate}\index{rw-imp-l(mrt_, mrs_, mrx_)@\texttt{rw-imp-l(mrt\_, mrs\_, mrx\_)}}\\
reads: the rewrite clause IMPLIES-left relates \$2 to \$3 \\
defined by \texttt{rw-imp-l} \\
characterized by \texttt{rw-imp-l}  \textit{(definitional)}:
\begin{quote}\ttfamily\footnotesize\raggedright\noindent
forall(\allowbreak{}[mrt\_, mrs\_, mrx\_], rw-imp-l(\allowbreak{}mrt\_, mrs\_, mrx\_) iff forsome(\allowbreak{}[mrg\_, mra\_, mrb\_, mrc\_], mrx\_ = [2, mrg\_, enc-imp(\allowbreak{}mra\_, mrc\_), enc-imp(\allowbreak{}mrb\_, mrc\_)] and mrc\_ in clf and [2, mrg\_, mra\_, mrb\_] in mrs\_))
\end{quote}
mentioned by 5 result(s)

\smallskip

\noindent\textbf{\texttt{rw-imp-r(mrt\_, mrs\_, mrx\_)}}\quad\textit{predicate}\index{rw-imp-r(mrt_, mrs_, mrx_)@\texttt{rw-imp-r(mrt\_, mrs\_, mrx\_)}}\\
reads: the rewrite clause IMPLIES-right relates \$2 to \$3 \\
defined by \texttt{rw-imp-r} \\
characterized by \texttt{rw-imp-r}  \textit{(definitional)}:
\begin{quote}\ttfamily\footnotesize\raggedright\noindent
forall(\allowbreak{}[mrt\_, mrs\_, mrx\_], rw-imp-r(\allowbreak{}mrt\_, mrs\_, mrx\_) iff forsome(\allowbreak{}[mrg\_, mra\_, mrb\_, mrc\_], mrx\_ = [2, mrg\_, enc-imp(\allowbreak{}mrc\_, mra\_), enc-imp(\allowbreak{}mrc\_, mrb\_)] and mrc\_ in clf and [2, ctx-add(\allowbreak{}mrg\_, mrc\_), mra\_, mrb\_] in mrs\_))
\end{quote}
mentioned by 5 result(s)

\smallskip

\noindent\textbf{\texttt{rw-not(mrt\_, mrs\_, mrx\_)}}\quad\textit{predicate}\index{rw-not(mrt_, mrs_, mrx_)@\texttt{rw-not(mrt\_, mrs\_, mrx\_)}}\\
reads: the rewrite clause NOT relates \$2 to \$3 \\
defined by \texttt{rw-not} \\
characterized by \texttt{rw-not}  \textit{(definitional)}:
\begin{quote}\ttfamily\footnotesize\raggedright\noindent
forall(\allowbreak{}[mrt\_, mrs\_, mrx\_], rw-not(\allowbreak{}mrt\_, mrs\_, mrx\_) iff forsome(\allowbreak{}[mrg\_, mra\_, mrb\_], mrx\_ = [2, mrg\_, enc-not(\allowbreak{}mra\_), enc-not(\allowbreak{}mrb\_)] and [2, mrg\_, mra\_, mrb\_] in mrs\_))
\end{quote}
mentioned by 5 result(s)

\smallskip

\noindent\textbf{\texttt{rw-or-l(mrt\_, mrs\_, mrx\_)}}\quad\textit{predicate}\index{rw-or-l(mrt_, mrs_, mrx_)@\texttt{rw-or-l(mrt\_, mrs\_, mrx\_)}}\\
reads: the rewrite clause OR-left relates \$2 to \$3 \\
defined by \texttt{rw-or-l} \\
characterized by \texttt{rw-or-l}  \textit{(definitional)}:
\begin{quote}\ttfamily\footnotesize\raggedright\noindent
forall(\allowbreak{}[mrt\_, mrs\_, mrx\_], rw-or-l(\allowbreak{}mrt\_, mrs\_, mrx\_) iff forsome(\allowbreak{}[mrg\_, mra\_, mrb\_, mrc\_], mrx\_ = [2, mrg\_, enc-or(\allowbreak{}mra\_, mrc\_), enc-or(\allowbreak{}mrb\_, mrc\_)] and mrc\_ in clf and [2, ctx-add(\allowbreak{}mrg\_, enc-not(\allowbreak{}mrc\_)), mra\_, mrb\_] in mrs\_))
\end{quote}
mentioned by 5 result(s)

\smallskip

\noindent\textbf{\texttt{rw-or-r(mrt\_, mrs\_, mrx\_)}}\quad\textit{predicate}\index{rw-or-r(mrt_, mrs_, mrx_)@\texttt{rw-or-r(mrt\_, mrs\_, mrx\_)}}\\
reads: the rewrite clause OR-right relates \$2 to \$3 \\
defined by \texttt{rw-or-r} \\
characterized by \texttt{rw-or-r}  \textit{(definitional)}:
\begin{quote}\ttfamily\footnotesize\raggedright\noindent
forall(\allowbreak{}[mrt\_, mrs\_, mrx\_], rw-or-r(\allowbreak{}mrt\_, mrs\_, mrx\_) iff forsome(\allowbreak{}[mrg\_, mra\_, mrb\_, mrc\_], mrx\_ = [2, mrg\_, enc-or(\allowbreak{}mrc\_, mra\_), enc-or(\allowbreak{}mrc\_, mrb\_)] and mrc\_ in clf and [2, ctx-add(\allowbreak{}mrg\_, enc-not(\allowbreak{}mrc\_)), mra\_, mrb\_] in mrs\_))
\end{quote}
mentioned by 5 result(s)

\smallskip

\noindent\textbf{\texttt{rw-pred(mrt\_, mrs\_, mrx\_)}}\quad\textit{predicate}\index{rw-pred(mrt_, mrs_, mrx_)@\texttt{rw-pred(mrt\_, mrs\_, mrx\_)}}\\
reads: the rewrite clause PREDICATE-argument relates \$2 to \$3 \\
defined by \texttt{rw-pred} \\
characterized by \texttt{rw-pred}  \textit{(definitional)}:
\begin{quote}\ttfamily\footnotesize\raggedright\noindent
forall(\allowbreak{}[mrt\_, mrs\_, mrx\_], rw-pred(\allowbreak{}mrt\_, mrs\_, mrx\_) iff forsome(\allowbreak{}[mrg\_, mrh\_, mra\_, mrb\_, mri\_], mrx\_ = [2, mrg\_, eapp(\allowbreak{}mrh\_, mra\_), eapp(\allowbreak{}mrh\_, mrb\_)] and mrh\_ in nn and mrh\_ = 10 and eapp(\allowbreak{}mrh\_, mra\_) in clf and eapp(\allowbreak{}mrh\_, mrb\_) in clf and mra\_ in tuples(\allowbreak{}expr-set) and mrb\_ in tuples(\allowbreak{}expr-set) and length(\allowbreak{}mrb\_) = length(\allowbreak{}mra\_) and mri\_ in nn and 1 <= mri\_ and mri\_ <= length(\allowbreak{}mra\_) and forall(\allowbreak{}[mrj\_ in nn], 1 <= mrj\_ implies mrj\_ <= length(\allowbreak{}mra\_) implies not(\allowbreak{}mrj\_ = mri\_) implies nth(\allowbreak{}mrj\_, mrb\_) = nth(\allowbreak{}mrj\_, mra\_)) and [1, mrg\_, nth(\allowbreak{}mri\_, mra\_), nth(\allowbreak{}mri\_, mrb\_)] in mrs\_))
\end{quote}
mentioned by 5 result(s)

\smallskip

\noindent\textbf{\texttt{rw-rel-l(mrt\_, mrs\_, mrx\_)}}\quad\textit{predicate}\index{rw-rel-l(mrt_, mrs_, mrx_)@\texttt{rw-rel-l(mrt\_, mrs\_, mrx\_)}}\\
reads: the rewrite clause ATOM-left relates \$2 to \$3 \\
defined by \texttt{rw-rel-l} \\
characterized by \texttt{rw-rel-l}  \textit{(definitional)}:
\begin{quote}\ttfamily\footnotesize\raggedright\noindent
forall(\allowbreak{}[mrt\_, mrs\_, mrx\_], rw-rel-l(\allowbreak{}mrt\_, mrs\_, mrx\_) iff forsome(\allowbreak{}[mrg\_, mrh\_, mra\_, mrb\_, mrc\_], mrx\_ = [2, mrg\_, eapp(\allowbreak{}mrh\_, [mra\_, mrc\_]), eapp(\allowbreak{}mrh\_, [mrb\_, mrc\_])] and mrh\_ in nn and 7 <= mrh\_ and mrh\_ <= 9 and mrc\_ in clt and [1, mrg\_, mra\_, mrb\_] in mrs\_))
\end{quote}
mentioned by 5 result(s)

\smallskip

\noindent\textbf{\texttt{rw-rel-r(mrt\_, mrs\_, mrx\_)}}\quad\textit{predicate}\index{rw-rel-r(mrt_, mrs_, mrx_)@\texttt{rw-rel-r(mrt\_, mrs\_, mrx\_)}}\\
reads: the rewrite clause ATOM-right relates \$2 to \$3 \\
defined by \texttt{rw-rel-r} \\
characterized by \texttt{rw-rel-r}  \textit{(definitional)}:
\begin{quote}\ttfamily\footnotesize\raggedright\noindent
forall(\allowbreak{}[mrt\_, mrs\_, mrx\_], rw-rel-r(\allowbreak{}mrt\_, mrs\_, mrx\_) iff forsome(\allowbreak{}[mrg\_, mrh\_, mra\_, mrb\_, mrc\_], mrx\_ = [2, mrg\_, eapp(\allowbreak{}mrh\_, [mrc\_, mra\_]), eapp(\allowbreak{}mrh\_, [mrc\_, mrb\_])] and mrh\_ in nn and 7 <= mrh\_ and mrh\_ <= 9 and mrc\_ in clt and [1, mrg\_, mra\_, mrb\_] in mrs\_))
\end{quote}
mentioned by 5 result(s)

\smallskip

\noindent\textbf{\texttt{rw-rule(mrt\_, mrk\_, mrs\_, mrx\_)}}\quad\textit{predicate}\index{rw-rule(mrt_, mrk_, mrs_, mrx_)@\texttt{rw-rule(mrt\_, mrk\_, mrs\_, mrx\_)}}\\
reads: some rewrite clause of stage at most \$2, justified by \$1, concludes \$4 from \$3 \\
defined by \texttt{rw-rule} \\
characterized by \texttt{rw-rule}  \textit{(definitional)}:
\begin{quote}\ttfamily\footnotesize\raggedright\noindent
forall(\allowbreak{}[mrt\_, mrk\_, mrs\_, mrx\_], rw-rule(\allowbreak{}mrt\_, mrk\_, mrs\_, mrx\_) iff 1 <= mrk\_ and rw-top(\allowbreak{}mrt\_, mrs\_, mrx\_) or 1 <= mrk\_ and rw-not(\allowbreak{}mrt\_, mrs\_, mrx\_) or 1 <= mrk\_ and rw-and-l(\allowbreak{}mrt\_, mrs\_, mrx\_) or 1 <= mrk\_ and rw-and-r(\allowbreak{}mrt\_, mrs\_, mrx\_) or 1 <= mrk\_ and rw-or-l(\allowbreak{}mrt\_, mrs\_, mrx\_) or 1 <= mrk\_ and rw-or-r(\allowbreak{}mrt\_, mrs\_, mrx\_) or 1 <= mrk\_ and rw-imp-l(\allowbreak{}mrt\_, mrs\_, mrx\_) or 1 <= mrk\_ and rw-imp-r(\allowbreak{}mrt\_, mrs\_, mrx\_) or 1 <= mrk\_ and rw-iff-l(\allowbreak{}mrt\_, mrs\_, mrx\_) or 1 <= mrk\_ and rw-iff-r(\allowbreak{}mrt\_, mrs\_, mrx\_) or 1 <= mrk\_ and rw-rel-l(\allowbreak{}mrt\_, mrs\_, mrx\_) or 1 <= mrk\_ and rw-rel-r(\allowbreak{}mrt\_, mrs\_, mrx\_) or 2 <= mrk\_ and rw-all(\allowbreak{}mrt\_, mrs\_, mrx\_) or 2 <= mrk\_ and rw-ex(\allowbreak{}mrt\_, mrs\_, mrx\_) or 4 <= mrk\_ and rw-pred(\allowbreak{}mrt\_, mrs\_, mrx\_) or 4 <= mrk\_ and rw-tapp(\allowbreak{}mrt\_, mrs\_, mrx\_) or 5 <= mrk\_ and rw-bind-dom(\allowbreak{}mrt\_, mrs\_, mrx\_) or 6 <= mrk\_ and rw-sep-body(\allowbreak{}mrt\_, mrs\_, mrx\_) or 6 <= mrk\_ and rw-tbind-body(\allowbreak{}mrt\_, mrs\_, mrx\_) or 7 <= mrk\_ and rw-desc-body(\allowbreak{}mrt\_, mrs\_, mrx\_) or 3 <= mrk\_ and rw-if-cond(\allowbreak{}mrt\_, mrs\_, mrx\_) or 3 <= mrk\_ and rw-if-then(\allowbreak{}mrt\_, mrs\_, mrx\_) or 3 <= mrk\_ and rw-if-else(\allowbreak{}mrt\_, mrs\_, mrx\_))
\end{quote}
mentioned by 4 result(s)

\smallskip

\noindent\textbf{\texttt{rw-sep-body(mrt\_, mrs\_, mrx\_)}}\quad\textit{predicate}\index{rw-sep-body(mrt_, mrs_, mrx_)@\texttt{rw-sep-body(mrt\_, mrs\_, mrx\_)}}\\
reads: the rewrite clause SEP-body relates \$2 to \$3 \\
defined by \texttt{rw-sep-body} \\
characterized by \texttt{rw-sep-body}  \textit{(definitional)}:
\begin{quote}\ttfamily\footnotesize\raggedright\noindent
forall(\allowbreak{}[mrt\_, mrs\_, mrx\_], rw-sep-body(\allowbreak{}mrt\_, mrs\_, mrx\_) iff forsome(\allowbreak{}[mrg\_, mra\_, mrb\_, mrc\_], mrx\_ = [1, mrg\_, enc-sep(\allowbreak{}mrc\_, mra\_), enc-sep(\allowbreak{}mrc\_, mrb\_)] and enc-sep(\allowbreak{}mrc\_, mra\_) in clt and enc-sep(\allowbreak{}mrc\_, mrb\_) in clt and forsome(\allowbreak{}[mrfs\_ in set], mrfs\_ subset nn and card(\allowbreak{}mrfs\_) in nn and forall(\allowbreak{}[mry\_ in nn], not(\allowbreak{}mry\_ in mrfs\_) implies [2, mrg\_, expr-open(\allowbreak{}mra\_, 0, efvar(\allowbreak{}mry\_)), expr-open(\allowbreak{}mrb\_, 0, efvar(\allowbreak{}mry\_))] in mrs\_))))
\end{quote}
mentioned by 5 result(s)

\smallskip

\noindent\textbf{\texttt{rw-tapp(mrt\_, mrs\_, mrx\_)}}\quad\textit{predicate}\index{rw-tapp(mrt_, mrs_, mrx_)@\texttt{rw-tapp(mrt\_, mrs\_, mrx\_)}}\\
reads: the rewrite clause TERM-argument relates \$2 to \$3 \\
defined by \texttt{rw-tapp} \\
characterized by \texttt{rw-tapp}  \textit{(definitional)}:
\begin{quote}\ttfamily\footnotesize\raggedright\noindent
forall(\allowbreak{}[mrt\_, mrs\_, mrx\_], rw-tapp(\allowbreak{}mrt\_, mrs\_, mrx\_) iff forsome(\allowbreak{}[mrg\_, mrh\_, mra\_, mrb\_, mri\_], mrx\_ = [1, mrg\_, eapp(\allowbreak{}mrh\_, mra\_), eapp(\allowbreak{}mrh\_, mrb\_)] and mrh\_ in nn and 11 <= mrh\_ and eapp(\allowbreak{}mrh\_, mra\_) in clt and eapp(\allowbreak{}mrh\_, mrb\_) in clt and mra\_ in tuples(\allowbreak{}expr-set) and mrb\_ in tuples(\allowbreak{}expr-set) and length(\allowbreak{}mrb\_) = length(\allowbreak{}mra\_) and mri\_ in nn and 1 <= mri\_ and mri\_ <= length(\allowbreak{}mra\_) and forall(\allowbreak{}[mrj\_ in nn], 1 <= mrj\_ implies mrj\_ <= length(\allowbreak{}mra\_) implies not(\allowbreak{}mrj\_ = mri\_) implies nth(\allowbreak{}mrj\_, mrb\_) = nth(\allowbreak{}mrj\_, mra\_)) and [1, mrg\_, nth(\allowbreak{}mri\_, mra\_), nth(\allowbreak{}mri\_, mrb\_)] in mrs\_))
\end{quote}
mentioned by 5 result(s)

\smallskip

\noindent\textbf{\texttt{rw-tbind-body(mrt\_, mrs\_, mrx\_)}}\quad\textit{predicate}\index{rw-tbind-body(mrt_, mrs_, mrx_)@\texttt{rw-tbind-body(mrt\_, mrs\_, mrx\_)}}\\
reads: the rewrite clause BIG-UNION / LAMBDA-body relates \$2 to \$3 \\
defined by \texttt{rw-tbind-body} \\
characterized by \texttt{rw-tbind-body}  \textit{(definitional)}:
\begin{quote}\ttfamily\footnotesize\raggedright\noindent
forall(\allowbreak{}[mrt\_, mrs\_, mrx\_], rw-tbind-body(\allowbreak{}mrt\_, mrs\_, mrx\_) iff forsome(\allowbreak{}[mrg\_, mrh\_, mra\_, mrb\_, mrc\_], mrx\_ = [1, mrg\_, ebind(\allowbreak{}mrh\_, 1, [mrc\_], mra\_), ebind(\allowbreak{}mrh\_, 1, [mrc\_], mrb\_)] and (\allowbreak{}mrh\_ = 24 or mrh\_ = 26) and ebind(\allowbreak{}mrh\_, 1, [mrc\_], mra\_) in clt and ebind(\allowbreak{}mrh\_, 1, [mrc\_], mrb\_) in clt and forsome(\allowbreak{}[mrfs\_ in set], mrfs\_ subset nn and card(\allowbreak{}mrfs\_) in nn and forall(\allowbreak{}[mry\_ in nn], not(\allowbreak{}mry\_ in mrfs\_) implies [1, mrg\_, expr-open(\allowbreak{}mra\_, 0, efvar(\allowbreak{}mry\_)), expr-open(\allowbreak{}mrb\_, 0, efvar(\allowbreak{}mry\_))] in mrs\_))))
\end{quote}
mentioned by 5 result(s)

\smallskip

\noindent\textbf{\texttt{rw-top(mrt\_, mrs\_, mrx\_)}}\quad\textit{predicate}\index{rw-top(mrt_, mrs_, mrx_)@\texttt{rw-top(mrt\_, mrs\_, mrx\_)}}\\
reads: the rewrite clause TOP relates \$2 to \$3 \\
defined by \texttt{rw-top} \\
characterized by \texttt{rw-top}  \textit{(definitional)}:
\begin{quote}\ttfamily\footnotesize\raggedright\noindent
forall(\allowbreak{}[mrt\_, mrs\_, mrx\_], rw-top(\allowbreak{}mrt\_, mrs\_, mrx\_) iff mrx\_ in mrt\_)
\end{quote}
mentioned by 5 result(s)

\smallskip

\noindent\textbf{\texttt{rw-univ}}\quad\textit{constant}\index{rw-univ@\texttt{rw-univ}}\\
reads: the rewrite tuples [kind, context, expression, expression] \\
defined by \texttt{rw-univ-def}

\smallskip


\section*{S}
\addcontentsline{toc}{section}{S}

\noindent\textbf{\texttt{sa-part(baa\_)}}\quad\textit{functoid}\index{sa-part(baa_)@\texttt{sa-part(baa\_)}}\\
reads: the self-adjoint elements of \$1 \\
definition:
\begin{quote}\ttfamily\footnotesize\raggedright\noindent
\{bav\_ in vec(\allowbreak{}baa\_): (\allowbreak{}vstar(\allowbreak{}baa\_))(\allowbreak{}bav\_) = bav\_\}
\end{quote}

\smallskip

\noindent\textbf{\texttt{sa-squares(baa\_)}}\quad\textit{functoid}\index{sa-squares(baa_)@\texttt{sa-squares(baa\_)}}\\
reads: the squares of self-adjoint elements of \$1 \\
definition:
\begin{quote}\ttfamily\footnotesize\raggedright\noindent
\{bav\_ in vec(\allowbreak{}baa\_): forsome(\allowbreak{}[baw\_ in sa-part(\allowbreak{}baa\_)], bav\_ = (\allowbreak{}vmul(\allowbreak{}baa\_))(\allowbreak{}baw\_, baw\_))\}
\end{quote}

\smallskip

\noindent\textbf{\texttt{same-strand-before(ssp\_, ssn\_, ssm\_)}}\quad\textit{predicate}\index{same-strand-before(ssp_, ssn_, ssm_)@\texttt{same-strand-before(ssp\_, ssn\_, ssm\_)}}\\
reads: \$2 comes strictly before \$3 on one strand of \$1 \\
defined by \texttt{same-strand-before} \\
characterized by \texttt{same-strand-before}  \textit{(definitional)}:
\begin{quote}\ttfamily\footnotesize\raggedright\noindent
forall(\allowbreak{}[ssp\_, ssn\_, ssm\_], same-strand-before(\allowbreak{}ssp\_, ssn\_, ssm\_) iff ssn\_ in nodes(\allowbreak{}ssp\_) and ssm\_ in nodes(\allowbreak{}ssp\_) and nth(\allowbreak{}2, ssn\_) = nth(\allowbreak{}2, ssm\_) and nth(\allowbreak{}1, ssn\_) < nth(\allowbreak{}1, ssm\_))
\end{quote}
mentioned by 30 result(s)

\smallskip

\noindent\textbf{\texttt{same-strand-upto(ssp\_, ssn\_, ssm\_)}}\quad\textit{predicate}\index{same-strand-upto(ssp_, ssn_, ssm_)@\texttt{same-strand-upto(ssp\_, ssn\_, ssm\_)}}\\
reads: \$2 comes at or before \$3 on one strand of \$1 \\
defined by \texttt{same-strand-upto} \\
characterized by \texttt{same-strand-upto}  \textit{(definitional)}:
\begin{quote}\ttfamily\footnotesize\raggedright\noindent
forall(\allowbreak{}[ssp\_, ssn\_, ssm\_], same-strand-upto(\allowbreak{}ssp\_, ssn\_, ssm\_) iff ssn\_ in nodes(\allowbreak{}ssp\_) and ssm\_ in nodes(\allowbreak{}ssp\_) and nth(\allowbreak{}2, ssn\_) = nth(\allowbreak{}2, ssm\_) and nth(\allowbreak{}1, ssn\_) <= nth(\allowbreak{}1, ssm\_))
\end{quote}
mentioned by 5 result(s)

\smallskip

\noindent\textbf{\texttt{save-proof}}\quad\textit{tactic}\index{save-proof@\texttt{save-proof}}\\
\texttt{(save-proof 'name)} -- Save the current script under a name without finishing. \\
declared in \texttt{tactics-help.scm}

\smallskip

\noindent\textbf{\texttt{scal(s)}}\quad\textit{accessor}\index{scal(s)@\texttt{scal(s)}}\\
reads: the scalar ring of \$1 \\
a structure slot \\
characterized by \texttt{rr-nvs@scal}  \textit{(definitional)}:
\begin{quote}\ttfamily\footnotesize\raggedright\noindent
scal(\allowbreak{}rr-nvs) == normed-field-as-commutative-ring(\allowbreak{}rr-normed-field)
\end{quote}
also: \texttt{qq-line@scal} \texttt{cips-scal-carrier} \texttt{module-scalar-ring} \texttt{module-scalar-zero-in} \texttt{sipq-scal-csip-as-sip} \texttt{ram-ro-scal-uba-as-ring} \texttt{nvs-scal-add-cnvs-as-nvs} \texttt{nvs-scal-mul-cnvs-as-nvs} ... (241 in all) \\
mentioned by 4572 result(s)

\smallskip

\noindent\textbf{\texttt{schanuel-map(m, n, la, lb)}}\quad\textit{functoid}\index{schanuel-map(m, n, la, lb)@\texttt{schanuel-map(m, n, la, lb)}}\\
reads: the map (t, b) |-> (b - \$4(t), t + \$3(b - \$4(t))) on Z\textasciicircum{}\$1 \\
definition:
\begin{quote}\ttfamily\footnotesize\raggedright\noindent
vnb-lambda(\allowbreak{}scz\_, vec(\allowbreak{}zz-pow(\allowbreak{}m)), seq-join(\allowbreak{}m, n, (\allowbreak{}vadd(\allowbreak{}zz-pow(\allowbreak{}n)))(\allowbreak{}seq-shift(\allowbreak{}n, scz\_), (\allowbreak{}vneg(\allowbreak{}zz-pow(\allowbreak{}n)))(\allowbreak{}lb(\allowbreak{}seq-head(\allowbreak{}n, scz\_)))), (\allowbreak{}vadd(\allowbreak{}zz-pow(\allowbreak{}n)))(\allowbreak{}seq-head(\allowbreak{}n, scz\_), la(\allowbreak{}(\allowbreak{}vadd(\allowbreak{}zz-pow(\allowbreak{}n)))(\allowbreak{}seq-shift(\allowbreak{}n, scz\_), (\allowbreak{}vneg(\allowbreak{}zz-pow(\allowbreak{}n)))(\allowbreak{}lb(\allowbreak{}seq-head(\allowbreak{}n, scz\_))))))))
\end{quote}
mentioned by 13 result(s)

\smallskip

\noindent\textbf{\texttt{scout}}\quad\textit{tactic}\index{scout@\texttt{scout}}\\
\texttt{(scout [depth [branch [nodes]]])} -- Speculatively search to a bounded depth across INDEPENDENT scratch branches; RETURNS a nested list (number-of-branches (d b) goal best-partials closing-branches) rather than printing.  Never touches the live proof. \\
declared in \texttt{tactics-help.scm}

\smallskip

\noindent\textbf{\texttt{scout-run}}\quad\textit{tactic}\index{scout-run@\texttt{scout-run}}\\
\texttt{(scout-run k)} -- Adopt closing branch k from the last (scout)/(scout-show) onto the live proof, running its tactics for real (they record and display normally). \\
declared in \texttt{tactics-help.scm}

\smallskip

\noindent\textbf{\texttt{scout-show}}\quad\textit{tactic}\index{scout-show@\texttt{scout-show}}\\
\texttt{(scout-show [depth [branch [nodes]]])} -- Like (scout), but PRINTS the human-readable report (examined count, goal, closing branches or best partials) to the REPL.  Returns the same nested list (scout) does. \\
declared in \texttt{tactics-help.scm}

\smallskip

\noindent\textbf{\texttt{sdet(r, m, n, a, s)}}\quad\textit{functoid}\index{sdet(r, m, n, a, s)@\texttt{sdet(r, m, n, a, s)}}\\
reads: the set of order-\$5 subdeterminants of \$4 \\
definition:
\begin{quote}\ttfamily\footnotesize\raggedright\noindent
\{sdx\_ in carr(\allowbreak{}r): forsome(\allowbreak{}[sdp\_, sdq\_], sdp\_ in injection(\allowbreak{}interval(\allowbreak{}1, s), interval(\allowbreak{}1, m)) and sdq\_ in injection(\allowbreak{}interval(\allowbreak{}1, s), interval(\allowbreak{}1, n)) and sdx\_ = det(\allowbreak{}r, s, submatrix(\allowbreak{}a, s, sdp\_, sdq\_)))\}
\end{quote}
characterized by \texttt{sdet-elim}  \textit{(proven)}:
\begin{quote}\ttfamily\footnotesize\raggedright\noindent
forall(\allowbreak{}[r, m, n, a, s, x in sdet(\allowbreak{}r, m, n, a, s)], forsome(\allowbreak{}[sdp\_, sdq\_], sdp\_ in injection(\allowbreak{}interval(\allowbreak{}1, s), interval(\allowbreak{}1, m)) and sdq\_ in injection(\allowbreak{}interval(\allowbreak{}1, s), interval(\allowbreak{}1, n)) and x = det(\allowbreak{}r, s, submatrix(\allowbreak{}a, s, sdp\_, sdq\_))))
\end{quote}
also: \texttt{sdet-intro} \texttt{sdet-unfold} \texttt{sdet-col-add} \texttt{sdet-row-add} \texttt{sdet-col-swap} \texttt{sdet-row-swap} \texttt{sdet-col-scale} \texttt{sdet-row-scale} ... (19 in all) \\
mentioned by 24 result(s)

\smallskip

\noindent\textbf{\texttt{seg-curve(m, f, a, eta)}}\quad\textit{functoid}\index{seg-curve(m, f, a, eta)@\texttt{seg-curve(m, f, a, eta)}}\\
definition:
\begin{quote}\ttfamily\footnotesize\raggedright\noindent
vnb-lambda(\allowbreak{}u\_, rr, f(\allowbreak{}(\allowbreak{}vadd(\allowbreak{}m))(\allowbreak{}a, (\allowbreak{}act(\allowbreak{}m))(\allowbreak{}u\_, eta))))
\end{quote}
characterized by \texttt{seg-curve-diff-at-cnvs-as-nvs}  \textit{(definitional)}:
\begin{quote}\ttfamily\footnotesize\raggedright\noindent
forall(\allowbreak{}[r], is-complex-normed-vector-space(\allowbreak{}r) implies forall(\allowbreak{}[f in fun(\allowbreak{}vec(\allowbreak{}r), rr), b\_ in vec(\allowbreak{}r), eta in vec(\allowbreak{}r), s\_ in rr, k\_], is-dir-diff-at(\allowbreak{}cnvs-as-nvs(\allowbreak{}r), f, (\allowbreak{}vadd(\allowbreak{}r))(\allowbreak{}b\_, (\allowbreak{}act(\allowbreak{}r))(\allowbreak{}s\_, eta)), eta, k\_) implies is-diff-at(\allowbreak{}seg-curve(\allowbreak{}cnvs-as-nvs(\allowbreak{}r), f, b\_, eta), s\_, k\_)))
\end{quote}
also: \texttt{seg-curve-apply} \texttt{seg-curve-in-fun} \texttt{seg-curve-diff-at} \\
mentioned by 29 result(s)

\smallskip

\noindent\textbf{\texttt{seg-deriv(sgp\_, sgq\_)}}\quad\textit{functoid}\index{seg-deriv(sgp_, sgq_)@\texttt{seg-deriv(sgp\_, sgq\_)}}\\
reads: the derivative of the segment from \$1 to \$2 \\
definition:
\begin{quote}\ttfamily\footnotesize\raggedright\noindent
vnb-lambda(\allowbreak{}sgs\_, ccint(\allowbreak{}0, 1), sgq\_ - sgp\_)
\end{quote}
characterized by \texttt{seg-deriv-bridge}  \textit{(proven)}:
\begin{quote}\ttfamily\footnotesize\raggedright\noindent
forall(\allowbreak{}[goua\_, goub\_, gouw\_], goub\_ - goua\_ = gouw\_ implies seg-deriv(\allowbreak{}goua\_, goub\_) == vnb-lambda(\allowbreak{}psx\_, ccint(\allowbreak{}0, 1), gouw\_))
\end{quote}
mentioned by 14 result(s)

\smallskip

\noindent\textbf{\texttt{seg-int(sgf\_, sgp\_, sgq\_)}}\quad\textit{functoid}\index{seg-int(sgf_, sgp_, sgq_)@\texttt{seg-int(sgf\_, sgp\_, sgq\_)}}\\
reads: the integral of \$1 along the segment from \$2 to \$3 \\
definition:
\begin{quote}\ttfamily\footnotesize\raggedright\noindent
line-int(\allowbreak{}sgf\_, seg-path(\allowbreak{}sgp\_, sgq\_), seg-deriv(\allowbreak{}sgp\_, sgq\_), 0, 1)
\end{quote}
characterized by \texttt{seg-int-log}  \textit{(proven)}:
\begin{quote}\ttfamily\footnotesize\raggedright\noindent
forall(\allowbreak{}[aop\_, aoq\_, aou\_], aop\_ in cc implies aoq\_ in cc implies aou\_ in cc implies 0 < imag-part(\allowbreak{}(\allowbreak{}aoq\_ - aou\_) * conjugate(\allowbreak{}aop\_ - aou\_)) implies seg-int(\allowbreak{}vnb-lambda(\allowbreak{}wnw\_, difference(\allowbreak{}cc, singleton(\allowbreak{}aou\_)), recip(\allowbreak{}wnw\_ - aou\_)), aop\_, aoq\_) = cc-log(\allowbreak{}conjugate(\allowbreak{}1i * (\allowbreak{}aop\_ - aoq\_)) * (\allowbreak{}aoq\_ - aou\_)) - cc-log(\allowbreak{}conjugate(\allowbreak{}1i * (\allowbreak{}aop\_ - aoq\_)) * (\allowbreak{}aop\_ - aou\_)))
\end{quote}
also: \texttt{seg-int-sum} \texttt{seg-int-const} \texttt{seg-int-in-cc} \texttt{seg-int-bridge} \texttt{seg-int-unfold} \texttt{seg-int-split-at} \texttt{seg-int-abs-bound} \texttt{seg-int-congruence} ... (9 in all) \\
mentioned by 35 result(s)

\smallskip

\noindent\textbf{\texttt{seg-path(sgp\_, sgq\_)}}\quad\textit{functoid}\index{seg-path(sgp_, sgq_)@\texttt{seg-path(sgp\_, sgq\_)}}\\
reads: the segment from \$1 to \$2 \\
definition:
\begin{quote}\ttfamily\footnotesize\raggedright\noindent
vnb-lambda(\allowbreak{}sgs\_, ccint(\allowbreak{}0, 1), sgp\_ + sgs\_ * (\allowbreak{}sgq\_ - sgp\_))
\end{quote}
characterized by \texttt{seg-path-apply}  \textit{(proven)}:
\begin{quote}\ttfamily\footnotesize\raggedright\noindent
forall(\allowbreak{}[goua\_ in cc, goub\_ in cc, gout\_ in ccint(\allowbreak{}0, 1)], (\allowbreak{}seg-path(\allowbreak{}goua\_, goub\_))(\allowbreak{}gout\_) = goua\_ + gout\_ * (\allowbreak{}goub\_ - goua\_))
\end{quote}
also: \texttt{seg-path-bridge} \texttt{seg-path-is-road} \\
mentioned by 46 result(s)

\smallskip

\noindent\textbf{\texttt{semi-inner-product-space}}\quad\textit{structure}\index{semi-inner-product-space@\texttt{semi-inner-product-space}}\\
predicate \texttt{is-semi-inner-product-space}; slots: scal vec vadd vzero vneg act rip

\smallskip

\noindent\textbf{\texttt{semigroup}}\quad\textit{structure}\index{semigroup@\texttt{semigroup}}\\
predicate \texttt{is-semigroup}; slots: carr opr

\smallskip

\noindent\textbf{\texttt{seminorm-ext(m, s, p, q)}}\quad\textit{functoid}\index{seminorm-ext(m, s, p, q)@\texttt{seminorm-ext(m, s, p, q)}}\\
reads: the extension of the seminorm \$3 on \$2 dominated by \$4 \\
definition:
\begin{quote}\ttfamily\footnotesize\raggedright\noindent
vnb-lambda(\allowbreak{}sxy\_, vec(\allowbreak{}m), inf(\allowbreak{}seminorm-ext-set(\allowbreak{}m, s, p, q, sxy\_)))
\end{quote}
characterized by \texttt{seminorm-extension}  \textit{(proven)}:
\begin{quote}\ttfamily\footnotesize\raggedright\noindent
forall(\allowbreak{}[m, s, p, q], is-module(\allowbreak{}m) implies scal(\allowbreak{}m) = normed-field-as-commutative-ring(\allowbreak{}rr-normed-field) implies is-submodule(\allowbreak{}m, s) implies is-seminorm-on(\allowbreak{}m, s, p) implies is-seminorm(\allowbreak{}rr-normed-field, m, q) implies forall(\allowbreak{}[dmx\_ in s], p(\allowbreak{}dmx\_) <= q(\allowbreak{}dmx\_)) implies is-seminorm(\allowbreak{}rr-normed-field, m, seminorm-ext(\allowbreak{}m, s, p, q)) and extends-on(\allowbreak{}s, seminorm-ext(\allowbreak{}m, s, p, q), p) and forall(\allowbreak{}[y in vec(\allowbreak{}m)], (\allowbreak{}seminorm-ext(\allowbreak{}m, s, p, q))(\allowbreak{}y) <= q(\allowbreak{}y)))
\end{quote}
mentioned by 3 result(s)

\smallskip

\noindent\textbf{\texttt{seminorm-ext-set(m, s, p, q, y)}}\quad\textit{functoid}\index{seminorm-ext-set(m, s, p, q, y)@\texttt{seminorm-ext-set(m, s, p, q, y)}}\\
reads: the values \$3(v) + \$4(w) over the splittings \$5 = v + w with v in \$2 \\
definition:
\begin{quote}\ttfamily\footnotesize\raggedright\noindent
\{sxt\_ in rr: forsome(\allowbreak{}[sev\_ in s, sew\_ in vec(\allowbreak{}m)], y = (\allowbreak{}vadd(\allowbreak{}m))(\allowbreak{}sev\_, sew\_) and sxt\_ = p(\allowbreak{}sev\_) + q(\allowbreak{}sew\_))\}
\end{quote}
mentioned by 17 result(s)

\smallskip

\noindent\textbf{\texttt{seminorm-top(m, ix, pf)}}\quad\textit{functoid}\index{seminorm-top(m, ix, pf)@\texttt{seminorm-top(m, ix, pf)}}\\
reads: the vector space \$1 with the topology defined by the family of seminorms \$3 \\
definition:
\begin{quote}\ttfamily\footnotesize\raggedright\noindent
[vec(\allowbreak{}m), \{snu\_ in power(\allowbreak{}vec(\allowbreak{}m)): is-sn-open(\allowbreak{}m, ix, pf, snu\_)\}]
\end{quote}
characterized by \texttt{seminorm-top-is-tvs}  \textit{(proven)}:
\begin{quote}\ttfamily\footnotesize\raggedright\noindent
forall(\allowbreak{}[fld, m, ix, pf], is-seminorm-family-on(\allowbreak{}fld, m, ix, pf) implies is-tvs(\allowbreak{}fld, m, seminorm-top(\allowbreak{}m, ix, pf)))
\end{quote}
also: \texttt{seminorm-top-locally-convex} \\
mentioned by 58 result(s)

\smallskip

\noindent\textbf{\texttt{sends-to(ssp\_, ssn\_, ssm\_)}}\quad\textit{predicate}\index{sends-to(ssp_, ssn_, ssm_)@\texttt{sends-to(ssp\_, ssn\_, ssm\_)}}\\
reads: the node \$2 of \$1 sends to the node \$3 \\
defined by \texttt{sends-to} \\
characterized by \texttt{sends-to}  \textit{(definitional)}:
\begin{quote}\ttfamily\footnotesize\raggedright\noindent
forall(\allowbreak{}[ssp\_, ssn\_, ssm\_], sends-to(\allowbreak{}ssp\_, ssn\_, ssm\_) iff ssn\_ in nodes(\allowbreak{}ssp\_) and ssm\_ in nodes(\allowbreak{}ssp\_) and is-sent(\allowbreak{}ssp\_, ssn\_) and is-received(\allowbreak{}ssp\_, ssm\_) and node-content(\allowbreak{}ssp\_, ssn\_) = node-content(\allowbreak{}ssp\_, ssm\_))
\end{quote}
also: \texttt{sends-to-dirs} \\
mentioned by 16 result(s)

\smallskip

\noindent\textbf{\texttt{sent}}\quad\textit{defined-fn}\index{sent@\texttt{sent}}\\
defined by \texttt{sent-def}

\smallskip

\noindent\textbf{\texttt{sep}}\quad\textit{operator}\index{sep@\texttt{sep}}\\
a term-forming head \\
characterized by \texttt{ball-sep-unfold}  \textit{(proven)}:
\begin{quote}\ttfamily\footnotesize\raggedright\noindent
forall(\allowbreak{}[s, c, r], ball(\allowbreak{}s, c, r) == \{y in pts(\allowbreak{}s): (\allowbreak{}dist(\allowbreak{}s))(\allowbreak{}c, y) <= r and not(\allowbreak{}(\allowbreak{}dist(\allowbreak{}s))(\allowbreak{}c, y) = r)\})
\end{quote}
also: \texttt{bsw-cs-real-sep} \texttt{closure-sep-unfold} \texttt{dist-layer-set-sep} \texttt{interior-sep-unfold} \texttt{sep-congruence-ptwise} \texttt{sep-fiber-insert-same} \texttt{closed-ball-sep-unfold} \texttt{sep-fiber-insert-other} ... (12 in all) \\
mentioned by 967 result(s)

\smallskip

\noindent\textbf{\texttt{sep-me}}\quad\textit{tactic}\index{sep-me@\texttt{sep-me}}\\
\texttt{(sep-me hyp)} -- Separation membership elim: split a separation-membership assumption into A-membership and the predicate. \\
declared in \texttt{tactics-help.scm}

\smallskip

\noindent\textbf{\texttt{sep-mi}}\quad\textit{tactic}\index{sep-mi@\texttt{sep-mi}}\\
\texttt{(sep-mi)} -- Separation membership intro: prove (IN t \{x in A | p\}). \\
declared in \texttt{tactics-help.scm}

\smallskip

\noindent\textbf{\texttt{sep-set}}\quad\textit{tactic}\index{sep-set@\texttt{sep-set}}\\
\texttt{(sep-set)} -- Separation sethood: the separation set \{x in A | p\} is a set. \\
declared in \texttt{tactics-help.scm}

\smallskip

\noindent\textbf{\texttt{seq-act(md, a, f)}}\quad\textit{functoid}\index{seq-act(md, a, f)@\texttt{seq-act(md, a, f)}}\\
reads: the pointwise multiple of the sequence \$3 by \$2 in \$1 \\
definition:
\begin{quote}\ttfamily\footnotesize\raggedright\noindent
vnb-lambda(\allowbreak{}hcn\_, nn, (\allowbreak{}act(\allowbreak{}md))(\allowbreak{}a, f(\allowbreak{}hcn\_)))
\end{quote}
characterized by \texttt{seq-act-at}  \textit{(proven)}:
\begin{quote}\ttfamily\footnotesize\raggedright\noindent
forall(\allowbreak{}[md, fq, a\_, f, n\_], is-module(\allowbreak{}md) implies scal(\allowbreak{}md) = normed-field-as-commutative-ring(\allowbreak{}cc-normed-field) implies is-inner-product-form(\allowbreak{}md, fq) implies a\_ in cc implies f in fun(\allowbreak{}nn, vec(\allowbreak{}md)) implies n\_ in nn implies (\allowbreak{}seq-act(\allowbreak{}md, a\_, f))(\allowbreak{}n\_) = (\allowbreak{}act(\allowbreak{}md))(\allowbreak{}a\_, f(\allowbreak{}n\_)))
\end{quote}
also: \texttt{seq-act-type} \texttt{seq-act-cauchy} \texttt{cseq-act-closed} \\
mentioned by 11 result(s)

\smallskip

\noindent\textbf{\texttt{seq-compact(s)}}\quad\textit{predicate}\index{seq-compact(s)@\texttt{seq-compact(s)}}\\
reads: \$1 is sequentially compact \\
defined by \texttt{seq-compact} \\
characterized by \texttt{seq-compact}  \textit{(definitional)}:
\begin{quote}\ttfamily\footnotesize\raggedright\noindent
forall(\allowbreak{}[s], seq-compact(\allowbreak{}s) iff is-metric-space(\allowbreak{}s) and forall(\allowbreak{}[f in fun(\allowbreak{}nn, pts(\allowbreak{}s))], forsome(\allowbreak{}[phi], strictly-mono-nn(\allowbreak{}phi) and forsome(\allowbreak{}[l in pts(\allowbreak{}s)], converges-to(\allowbreak{}s, subseq(\allowbreak{}f, phi), l)))))
\end{quote}
also: \texttt{dbi-seq-compact} \texttt{seq-compact-image} \texttt{rmm-two-seq-compact} \texttt{dual-ball-seq-compact} \texttt{compact-iff-seq-compact} \texttt{compact-implies-seq-compact} \texttt{seq-compact-implies-compact} \texttt{seq-compact-countable-product} ... (10 in all) \\
mentioned by 18 result(s)

\smallskip

\noindent\textbf{\texttt{seq-cons(sck\_, scv\_)}}\quad\textit{functoid}\index{seq-cons(sck_, scv_)@\texttt{seq-cons(sck\_, scv\_)}}\\
reads: the sequence \$1 followed by \$2 \\
definition:
\begin{quote}\ttfamily\footnotesize\raggedright\noindent
vnb-lambda(\allowbreak{}scj\_, nn, if(\allowbreak{}scj\_ = 0, sck\_, scv\_(\allowbreak{}nn-minus(\allowbreak{}scj\_, 1))))
\end{quote}
characterized by \texttt{seq-cons-facts}  \textit{(proven)}:
\begin{quote}\ttfamily\footnotesize\raggedright\noindent
forall(\allowbreak{}[k in nn, v in fun(\allowbreak{}nn, nn)], seq-cons(\allowbreak{}k, v) in fun(\allowbreak{}nn, nn) and (\allowbreak{}seq-cons(\allowbreak{}k, v))(\allowbreak{}0) = k and seq-tail(\allowbreak{}seq-cons(\allowbreak{}k, v)) = v)
\end{quote}
mentioned by 1 result(s)

\smallskip

\noindent\textbf{\texttt{seq-form(fq, f, g)}}\quad\textit{functoid}\index{seq-form(fq, f, g)@\texttt{seq-form(fq, f, g)}}\\
reads: the sequence of the values of \$1 at the terms of \$2 and \$3 \\
definition:
\begin{quote}\ttfamily\footnotesize\raggedright\noindent
vnb-lambda(\allowbreak{}hcn\_, nn, fq(\allowbreak{}f(\allowbreak{}hcn\_), g(\allowbreak{}hcn\_)))
\end{quote}
characterized by \texttt{seq-form-at}  \textit{(proven)}:
\begin{quote}\ttfamily\footnotesize\raggedright\noindent
forall(\allowbreak{}[md, fq, f, g, n\_], is-module(\allowbreak{}md) implies scal(\allowbreak{}md) = normed-field-as-commutative-ring(\allowbreak{}cc-normed-field) implies is-inner-product-form(\allowbreak{}md, fq) implies f in fun(\allowbreak{}nn, vec(\allowbreak{}md)) implies g in fun(\allowbreak{}nn, vec(\allowbreak{}md)) implies n\_ in nn implies (\allowbreak{}seq-form(\allowbreak{}fq, f, g))(\allowbreak{}n\_) = fq(\allowbreak{}f(\allowbreak{}n\_), g(\allowbreak{}n\_)))
\end{quote}
also: \texttt{seq-form-type} \texttt{seq-form-cauchy} \\
mentioned by 9 result(s)

\smallskip

\noindent\textbf{\texttt{seq-head(n, z)}}\quad\textit{functoid}\index{seq-head(n, z)@\texttt{seq-head(n, z)}}\\
reads: the first \$1 terms of \$2 \\
definition:
\begin{quote}\ttfamily\footnotesize\raggedright\noindent
vnb-lambda(\allowbreak{}shi\_, interval(\allowbreak{}1, n), z(\allowbreak{}shi\_))
\end{quote}
characterized by \texttt{seq-head-at}  \textit{(proven)}:
\begin{quote}\ttfamily\footnotesize\raggedright\noindent
forall(\allowbreak{}[m, n, z, i], m in nn implies n in nn implies m = n + n implies z in vec(\allowbreak{}zz-pow(\allowbreak{}m)) implies i in interval(\allowbreak{}1, n) implies (\allowbreak{}seq-head(\allowbreak{}n, z))(\allowbreak{}i) = z(\allowbreak{}i))
\end{quote}
also: \texttt{seq-head-act} \texttt{seq-head-vec} \texttt{seq-head-join} \texttt{seq-head-vadd} \texttt{seq-head-vneg} \texttt{seq-head-vzero} \\
mentioned by 23 result(s)

\smallskip

\noindent\textbf{\texttt{seq-join(m, n, hd, tl)}}\quad\textit{functoid}\index{seq-join(m, n, hd, tl)@\texttt{seq-join(m, n, hd, tl)}}\\
reads: \$3 followed by \$4, of length \$1 \\
definition:
\begin{quote}\ttfamily\footnotesize\raggedright\noindent
vnb-lambda(\allowbreak{}jnj\_, interval(\allowbreak{}1, m), if(\allowbreak{}jnj\_ <= n, hd(\allowbreak{}jnj\_), tl(\allowbreak{}nn-minus(\allowbreak{}jnj\_, n))))
\end{quote}
characterized by \texttt{seq-join-hi}  \textit{(proven)}:
\begin{quote}\ttfamily\footnotesize\raggedright\noindent
forall(\allowbreak{}[m, n, hd, tl, j], m in nn implies n in nn implies m = n + n implies hd in vec(\allowbreak{}zz-pow(\allowbreak{}n)) implies tl in vec(\allowbreak{}zz-pow(\allowbreak{}n)) implies j in interval(\allowbreak{}1, m) implies not(\allowbreak{}j <= n) implies (\allowbreak{}seq-join(\allowbreak{}m, n, hd, tl))(\allowbreak{}j) = tl(\allowbreak{}nn-minus(\allowbreak{}j, n)))
\end{quote}
also: \texttt{seq-join-lo} \texttt{seq-join-ext} \texttt{seq-join-vec} \\
mentioned by 13 result(s)

\smallskip

\noindent\textbf{\texttt{seq-limit(f)}}\quad\textit{functoid}\index{seq-limit(f)@\texttt{seq-limit(f)}}\\
reads: the limit of the sequence \$1 \\
definition:
\begin{quote}\ttfamily\footnotesize\raggedright\noindent
if(\allowbreak{}converges(\allowbreak{}rr-ms, f), iota(\allowbreak{}lm\_, converges-to(\allowbreak{}rr-ms, f, lm\_)), 0)
\end{quote}
characterized by \texttt{seq-limit-in-rr}  \textit{(proven)}:
\begin{quote}\ttfamily\footnotesize\raggedright\noindent
forall(\allowbreak{}[f], seq-limit(\allowbreak{}f) in rr)
\end{quote}
also: \texttt{seq-limit-value} \texttt{seq-limit-converges-to} \\
mentioned by 4 result(s)

\smallskip

\noindent\textbf{\texttt{seq-proj(spk\_, spz\_)}}\quad\textit{functoid}\index{seq-proj(spk_, spz_)@\texttt{seq-proj(spk\_, spz\_)}}\\
reads: the \$1-th section of the sequence \$2 \\
definition:
\begin{quote}\ttfamily\footnotesize\raggedright\noindent
vnb-lambda(\allowbreak{}spj\_, nn, spz\_(\allowbreak{}nnpair(\allowbreak{}spk\_, spj\_)))
\end{quote}
characterized by \texttt{seq-proj-facts}  \textit{(proven)}:
\begin{quote}\ttfamily\footnotesize\raggedright\noindent
forall(\allowbreak{}[k in nn, z in fun(\allowbreak{}nn, nn)], seq-proj(\allowbreak{}k, z) in fun(\allowbreak{}nn, nn) and forall(\allowbreak{}[j in nn], (\allowbreak{}seq-proj(\allowbreak{}k, z))(\allowbreak{}j) = z(\allowbreak{}nnpair(\allowbreak{}k, j))))
\end{quote}
also: \texttt{seq-proj-prefix} \\
mentioned by 4 result(s)

\smallskip

\noindent\textbf{\texttt{seq-set}}\quad\textit{constant}\index{seq-set@\texttt{seq-set}}\\
reads: the sequents \\
defined by \texttt{seq-def}

\smallskip

\noindent\textbf{\texttt{seq-shift(n, z)}}\quad\textit{functoid}\index{seq-shift(n, z)@\texttt{seq-shift(n, z)}}\\
reads: the terms \$1 + 1, ..., \$1 + \$1 of \$2 \\
definition:
\begin{quote}\ttfamily\footnotesize\raggedright\noindent
vnb-lambda(\allowbreak{}sfi\_, interval(\allowbreak{}1, n), z(\allowbreak{}n + sfi\_))
\end{quote}
characterized by \texttt{seq-shift-at}  \textit{(proven)}:
\begin{quote}\ttfamily\footnotesize\raggedright\noindent
forall(\allowbreak{}[m, n, z, i], m in nn implies n in nn implies m = n + n implies z in vec(\allowbreak{}zz-pow(\allowbreak{}m)) implies i in interval(\allowbreak{}1, n) implies (\allowbreak{}seq-shift(\allowbreak{}n, z))(\allowbreak{}i) = z(\allowbreak{}n + i))
\end{quote}
also: \texttt{seq-shift-act} \texttt{seq-shift-vec} \texttt{seq-shift-join} \texttt{seq-shift-vadd} \texttt{seq-shift-vneg} \texttt{seq-shift-vzero} \\
mentioned by 25 result(s)

\smallskip

\noindent\textbf{\texttt{seq-space-ms(sqx\_)}}\quad\textit{functoid}\index{seq-space-ms(sqx_)@\texttt{seq-space-ms(sqx\_)}}\\
reads: the sequences in \$1 with the metric sum of 2\textasciicircum{}-k over the indices where they differ \\
definition:
\begin{quote}\ttfamily\footnotesize\raggedright\noindent
product-metric-w(\allowbreak{}vnb-lambda(\allowbreak{}sqk\_, nn, discrete-ms(\allowbreak{}sqx\_)), vnb-lambda(\allowbreak{}sqk\_, nn, (\allowbreak{}/)(\allowbreak{}1, 2 \textasciicircum{} sqk\_)))
\end{quote}
mentioned by 2 result(s)

\smallskip

\noindent\textbf{\texttt{seq-tail(stz\_)}}\quad\textit{functoid}\index{seq-tail(stz_)@\texttt{seq-tail(stz\_)}}\\
reads: the tail of the sequence \$1 \\
definition:
\begin{quote}\ttfamily\footnotesize\raggedright\noindent
vnb-lambda(\allowbreak{}stj\_, nn, stz\_(\allowbreak{}succ(\allowbreak{}stj\_)))
\end{quote}
characterized by \texttt{seq-tail-facts}  \textit{(proven)}:
\begin{quote}\ttfamily\footnotesize\raggedright\noindent
forall(\allowbreak{}[z in fun(\allowbreak{}nn, nn)], seq-tail(\allowbreak{}z) in fun(\allowbreak{}nn, nn) and forall(\allowbreak{}[j in nn], (\allowbreak{}seq-tail(\allowbreak{}z))(\allowbreak{}j) = z(\allowbreak{}succ(\allowbreak{}j))))
\end{quote}
mentioned by 2 result(s)

\smallskip

\noindent\textbf{\texttt{seq-vadd(md, f, g)}}\quad\textit{functoid}\index{seq-vadd(md, f, g)@\texttt{seq-vadd(md, f, g)}}\\
reads: the pointwise sum of the sequences \$2 and \$3 in \$1 \\
definition:
\begin{quote}\ttfamily\footnotesize\raggedright\noindent
vnb-lambda(\allowbreak{}hcn\_, nn, (\allowbreak{}vadd(\allowbreak{}md))(\allowbreak{}f(\allowbreak{}hcn\_), g(\allowbreak{}hcn\_)))
\end{quote}
characterized by \texttt{seq-vadd-at}  \textit{(proven)}:
\begin{quote}\ttfamily\footnotesize\raggedright\noindent
forall(\allowbreak{}[md, fq, f, g, n\_], is-module(\allowbreak{}md) implies scal(\allowbreak{}md) = normed-field-as-commutative-ring(\allowbreak{}cc-normed-field) implies is-inner-product-form(\allowbreak{}md, fq) implies f in fun(\allowbreak{}nn, vec(\allowbreak{}md)) implies g in fun(\allowbreak{}nn, vec(\allowbreak{}md)) implies n\_ in nn implies (\allowbreak{}seq-vadd(\allowbreak{}md, f, g))(\allowbreak{}n\_) = (\allowbreak{}vadd(\allowbreak{}md))(\allowbreak{}f(\allowbreak{}n\_), g(\allowbreak{}n\_)))
\end{quote}
also: \texttt{seq-vadd-type} \texttt{seq-vadd-cauchy} \texttt{cseq-vadd-closed} \\
mentioned by 12 result(s)

\smallskip

\noindent\textbf{\texttt{seq-vneg(md, f)}}\quad\textit{functoid}\index{seq-vneg(md, f)@\texttt{seq-vneg(md, f)}}\\
reads: the pointwise negative of the sequence \$2 in \$1 \\
definition:
\begin{quote}\ttfamily\footnotesize\raggedright\noindent
vnb-lambda(\allowbreak{}hcn\_, nn, (\allowbreak{}vneg(\allowbreak{}md))(\allowbreak{}f(\allowbreak{}hcn\_)))
\end{quote}
characterized by \texttt{seq-vneg-at}  \textit{(proven)}:
\begin{quote}\ttfamily\footnotesize\raggedright\noindent
forall(\allowbreak{}[md, fq, f, n\_], is-module(\allowbreak{}md) implies scal(\allowbreak{}md) = normed-field-as-commutative-ring(\allowbreak{}cc-normed-field) implies is-inner-product-form(\allowbreak{}md, fq) implies f in fun(\allowbreak{}nn, vec(\allowbreak{}md)) implies n\_ in nn implies (\allowbreak{}seq-vneg(\allowbreak{}md, f))(\allowbreak{}n\_) = (\allowbreak{}vneg(\allowbreak{}md))(\allowbreak{}f(\allowbreak{}n\_)))
\end{quote}
also: \texttt{seq-vneg-type} \texttt{seq-vneg-cauchy} \texttt{cseq-vneg-closed} \\
mentioned by 10 result(s)

\smallskip

\noindent\textbf{\texttt{seq-zero(md)}}\quad\textit{functoid}\index{seq-zero(md)@\texttt{seq-zero(md)}}\\
reads: the zero sequence of \$1 \\
definition:
\begin{quote}\ttfamily\footnotesize\raggedright\noindent
vnb-lambda(\allowbreak{}hcn\_, nn, vzero(\allowbreak{}md))
\end{quote}
characterized by \texttt{seq-zero-at}  \textit{(proven)}:
\begin{quote}\ttfamily\footnotesize\raggedright\noindent
forall(\allowbreak{}[md, fq, n\_], is-module(\allowbreak{}md) implies scal(\allowbreak{}md) = normed-field-as-commutative-ring(\allowbreak{}cc-normed-field) implies is-inner-product-form(\allowbreak{}md, fq) implies n\_ in nn implies (\allowbreak{}seq-zero(\allowbreak{}md))(\allowbreak{}n\_) = vzero(\allowbreak{}md))
\end{quote}
also: \texttt{seq-zero-type} \texttt{seq-zero-embed} \texttt{seq-zero-cauchy} \texttt{cseq-zero-closed} \\
mentioned by 9 result(s)

\smallskip

\noindent\textbf{\texttt{series-converges(f)}}\quad\textit{predicate}\index{series-converges(f)@\texttt{series-converges(f)}}\\
reads: the series \$1 converges \\
defined by \texttt{series-converges} \\
characterized by \texttt{series-converges}  \textit{(definitional)}:
\begin{quote}\ttfamily\footnotesize\raggedright\noindent
forall(\allowbreak{}[f], series-converges(\allowbreak{}f) iff converges(\allowbreak{}rr-ms, vnb-lambda(\allowbreak{}k, nn, series-partial-sum(\allowbreak{}f, k))))
\end{quote}
also: \texttt{series-converges-sum} \texttt{zero-series-converges} \texttt{geometric-series-converges} \texttt{series-converges-to-converges} \\
mentioned by 106 result(s)

\smallskip

\noindent\textbf{\texttt{series-converges-to(f, l)}}\quad\textit{predicate}\index{series-converges-to(f, l)@\texttt{series-converges-to(f, l)}}\\
reads: the series \$1 converges to \$2 \\
defined by \texttt{series-converges-to} \\
characterized by \texttt{series-converges-to}  \textit{(definitional)}:
\begin{quote}\ttfamily\footnotesize\raggedright\noindent
forall(\allowbreak{}[f, l], series-converges-to(\allowbreak{}f, l) iff converges-to(\allowbreak{}rr-ms, vnb-lambda(\allowbreak{}k, nn, series-partial-sum(\allowbreak{}f, k)), l))
\end{quote}
also: \texttt{series-converges-to-add} \texttt{geometric-series-converges-to} \texttt{series-converges-to-converges} \\
mentioned by 15 result(s)

\smallskip

\noindent\textbf{\texttt{series-limit(f)}}\quad\textit{functoid}\index{series-limit(f)@\texttt{series-limit(f)}}\\
reads: the sum of the series \$1 \\
definition:
\begin{quote}\ttfamily\footnotesize\raggedright\noindent
iota(\allowbreak{}lm\_, lm\_ in rr and series-converges-to(\allowbreak{}f, lm\_))
\end{quote}
characterized by \texttt{series-limit-add}  \textit{(proven)}:
\begin{quote}\ttfamily\footnotesize\raggedright\noindent
forall(\allowbreak{}[f in fun(\allowbreak{}nn, rr), g in fun(\allowbreak{}nn, rr), h in fun(\allowbreak{}nn, rr)], forall(\allowbreak{}[i\_ in nn], h(\allowbreak{}i\_) = f(\allowbreak{}i\_) + g(\allowbreak{}i\_)) implies series-converges(\allowbreak{}f) implies series-converges(\allowbreak{}g) implies series-limit(\allowbreak{}h) = series-limit(\allowbreak{}f) + series-limit(\allowbreak{}g))
\end{quote}
also: \texttt{series-limit-iota} \texttt{series-limit-in-rr} \texttt{series-limit-converges-to} \\
mentioned by 15 result(s)

\smallskip

\noindent\textbf{\texttt{series-partial-sum(f, k)}}\quad\textit{functoid}\index{series-partial-sum(f, k)@\texttt{series-partial-sum(f, k)}}\\
reads: the sum of the first \$2 terms of the series \$1 \\
definition:
\begin{quote}\ttfamily\footnotesize\raggedright\noindent
sum-ag(\allowbreak{}normed-field-additive-ag(\allowbreak{}rr-normed-field), f, k)
\end{quote}
characterized by \texttt{series-partial-sum-add}  \textit{(proven)}:
\begin{quote}\ttfamily\footnotesize\raggedright\noindent
forall(\allowbreak{}[k in nn, f in fun(\allowbreak{}nn, rr), g in fun(\allowbreak{}nn, rr), h in fun(\allowbreak{}nn, rr)], forall(\allowbreak{}[i\_ in nn], h(\allowbreak{}i\_) = f(\allowbreak{}i\_) + g(\allowbreak{}i\_)) implies series-partial-sum(\allowbreak{}h, k) = series-partial-sum(\allowbreak{}f, k) + series-partial-sum(\allowbreak{}g, k))
\end{quote}
also: \texttt{series-partial-sum-mono} \texttt{series-partial-sum-succ} \texttt{series-partial-sum-zero} \texttt{series-partial-sum-in-rr} \texttt{series-partial-sum-shift} \texttt{series-partial-sum-abs-le} \texttt{series-partial-sum-nonneg} \texttt{series-partial-sum-unfold} ... (24 in all) \\
mentioned by 168 result(s)

\smallskip

\noindent\textbf{\texttt{set-dist(m, aa, bb)}}\quad\textit{functoid}\index{set-dist(m, aa, bb)@\texttt{set-dist(m, aa, bb)}}\\
reads: the distance in \$1 between the sets \$2 and \$3 \\
definition:
\begin{quote}\ttfamily\footnotesize\raggedright\noindent
einf(\allowbreak{}\{mow\_ in rr-pos-star: forsome(\allowbreak{}[mox\_ in aa, moy\_ in bb], mow\_ = (\allowbreak{}dist(\allowbreak{}m))(\allowbreak{}mox\_, moy\_))\})
\end{quote}
characterized by \texttt{set-dist-unfold}  \textit{(proven)}:
\begin{quote}\ttfamily\footnotesize\raggedright\noindent
forall(\allowbreak{}[m, aa, bb], set-dist(\allowbreak{}m, aa, bb) == einf(\allowbreak{}\{mow\_ in rr-pos-star: forsome(\allowbreak{}[mox\_ in aa, moy\_ in bb], mow\_ = (\allowbreak{}dist(\allowbreak{}m))(\allowbreak{}mox\_, moy\_))\}))
\end{quote}
also: \texttt{set-dist-pos-elim} \texttt{set-dist-pos-intro} \texttt{set-dist-union-pos} \\
mentioned by 15 result(s)

\smallskip

\noindent\textbf{\texttt{set-power}}\quad\textit{functoid}\index{set-power@\texttt{set-power}}\\
reads: the \$3-th power of the set \$2 in \$1 \\
defined by \texttt{set-power-zero} \texttt{set-power-succ} \\
characterized by \texttt{set-power-succ}  \textit{(definitional)}:
\begin{quote}\ttfamily\footnotesize\raggedright\noindent
forall(\allowbreak{}[sqg\_, sqv\_, sqk\_ in nn], set-power(\allowbreak{}sqg\_, sqv\_, succ(\allowbreak{}sqk\_)) == set-prod(\allowbreak{}sqg\_, set-power(\allowbreak{}sqg\_, sqv\_, sqk\_), sqv\_))
\end{quote}
also: \texttt{set-power-zero} \texttt{set-power-is-set} \texttt{set-power-mono-add} \texttt{set-power-sub-carr} \texttt{set-power-mono-succ} \texttt{set-power-succ-elim} \texttt{set-power-succ-open} \texttt{set-power-zero-elim} ... (11 in all) \\
mentioned by 25 result(s)

\smallskip

\noindent\textbf{\texttt{set-prod(spg\_, spa\_, spb\_)}}\quad\textit{functoid}\index{set-prod(spg_, spa_, spb_)@\texttt{set-prod(spg\_, spa\_, spb\_)}}\\
reads: the product set \$2 \$3 in \$1 \\
definition:
\begin{quote}\ttfamily\footnotesize\raggedright\noindent
\{spz\_ in carr(\allowbreak{}spg\_): forsome(\allowbreak{}[spx\_ in spa\_, spy\_ in spb\_], spz\_ = (\allowbreak{}opr(\allowbreak{}spg\_))(\allowbreak{}spx\_, spy\_))\}
\end{quote}
characterized by \texttt{set-prod-elim}  \textit{(proven)}:
\begin{quote}\ttfamily\footnotesize\raggedright\noindent
forall(\allowbreak{}[g, a, b, z in set-prod(\allowbreak{}g, a, b)], z in carr(\allowbreak{}g) and forsome(\allowbreak{}[x in a, y in b], z = (\allowbreak{}opr(\allowbreak{}g))(\allowbreak{}x, y)))
\end{quote}
also: \texttt{set-prod-open} \texttt{set-prod-intro} \texttt{set-prod-unfold} \\
mentioned by 7 result(s)

\smallskip

\noindent\textbf{\texttt{setoid}}\quad\textit{structure}\index{setoid@\texttt{setoid}}\\
predicate \texttt{is-setoid}; slots: pts rel

\smallskip

\noindent\textbf{\texttt{shift-el(baa\_, baz\_, bax\_)}}\quad\textit{functoid}\index{shift-el(baa_, baz_, bax_)@\texttt{shift-el(baa\_, baz\_, bax\_)}}\\
reads: \$2 1 - \$3 in \$1 \\
definition:
\begin{quote}\ttfamily\footnotesize\raggedright\noindent
(\allowbreak{}vadd(\allowbreak{}baa\_))(\allowbreak{}(\allowbreak{}act(\allowbreak{}baa\_))(\allowbreak{}baz\_, vone(\allowbreak{}baa\_)), (\allowbreak{}vneg(\allowbreak{}baa\_))(\allowbreak{}bax\_))
\end{quote}
mentioned by 44 result(s)

\smallskip

\noindent\textbf{\texttt{show}}\quad\textit{tactic}\index{show@\texttt{show}}\\
\texttt{(show)} -- Redisplay the current proof state. \\
declared in \texttt{tactics-help.scm}

\smallskip

\noindent\textbf{\texttt{show-graph}}\quad\textit{tactic}\index{show-graph@\texttt{show-graph}}\\
\texttt{(show-graph)} -- Print EVERY node of the current proof, closed ones included: one line per sequent node in creation order (the bracketed number the panel prints, the sequent, [GROUNDED] when justified), and under each node one line per inference into it -- the rule and the numbers of the hypothesis nodes it rests on.  The focus and the open leaves are marked.  This is where a branch a composite closed is found: after (cut L), the side leaf G |- L and the main branch G, L |- goal are both nodes, and zero-it on an existential closes the main one.  Records nothing; the [Graph] button of the Focus Workspace shows it. \\
declared in \texttt{tactics-help.scm}

\smallskip

\noindent\textbf{\texttt{sigma(s)}}\quad\textit{accessor}\index{sigma(s)@\texttt{sigma(s)}}\\
reads: the sigma-algebra of \$1 \\
a structure slot \\
characterized by \texttt{ds-sigma}  \textit{(proven)}:
\begin{quote}\ttfamily\footnotesize\raggedright\noindent
forall(\allowbreak{}[a\_ in set], sigma(\allowbreak{}discrete-space(\allowbreak{}a\_)) == power(\allowbreak{}a\_))
\end{quote}
mentioned by 10 result(s)

\smallskip

\noindent\textbf{\texttt{sigma-generated(omega, ce)}}\quad\textit{functoid}\index{sigma-generated(omega, ce)@\texttt{sigma-generated(omega, ce)}}\\
reads: the sigma-algebra of subsets of \$1 generated by \$2 \\
definition:
\begin{quote}\ttfamily\footnotesize\raggedright\noindent
\{a\_ in power(\allowbreak{}omega): forall(\allowbreak{}[cb\_], cb\_ in power(\allowbreak{}power(\allowbreak{}omega)) and is-sigma-algebra(\allowbreak{}omega, cb\_) and ce subset cb\_ implies a\_ in cb\_)\}
\end{quote}
characterized by \texttt{sigma-generated-membership}  \textit{(definitional)}:
\begin{quote}\ttfamily\footnotesize\raggedright\noindent
forall(\allowbreak{}[omega, ce, a\_], a\_ in sigma-generated(\allowbreak{}omega, ce) iff a\_ in power(\allowbreak{}omega) and forall(\allowbreak{}[cb\_], cb\_ in power(\allowbreak{}power(\allowbreak{}omega)) and is-sigma-algebra(\allowbreak{}omega, cb\_) and ce subset cb\_ implies a\_ in cb\_))
\end{quote}
also: \texttt{sigma-generated-intro} \texttt{sigma-generated-in-every} \texttt{sigma-generated-in-power} \texttt{sigma-generated-is-sigma-algebra} \\
mentioned by 14 result(s)

\smallskip

\noindent\textbf{\texttt{signed-neg(ca, nu)}}\quad\textit{functoid}\index{signed-neg(ca, nu)@\texttt{signed-neg(ca, nu)}}\\
reads: the negative of the signed measure \$2 \\
definition:
\begin{quote}\ttfamily\footnotesize\raggedright\noindent
vnb-lambda(\allowbreak{}sna\_, ca, ext-neg(\allowbreak{}nu(\allowbreak{}sna\_)))
\end{quote}
characterized by \texttt{signed-neg-fun}  \textit{(proven)}:
\begin{quote}\ttfamily\footnotesize\raggedright\noindent
forall(\allowbreak{}[omega, ca, nu], is-sigma-algebra(\allowbreak{}omega, ca) implies nu in fun(\allowbreak{}ca, rr-star) implies signed-neg(\allowbreak{}ca, nu) in fun(\allowbreak{}ca, rr-star))
\end{quote}
also: \texttt{signed-neg-value} \texttt{signed-neg-is-signed-measure} \\
mentioned by 6 result(s)

\smallskip

\noindent\textbf{\texttt{simp}}\quad\textit{tactic}\index{simp@\texttt{simp}}\\
\texttt{(simp [target])} -- Rewrite a commutative-ring SUBTERM of the goal to canonical form, IN PLACE (e.g. (x+y)\textasciicircum{}2 inside a larger goal becomes x\textasciicircum{}2 + 2*x*y + y\textasciicircum{}2).  Works on BOTH surfaces: concrete number domains (+ * - \textasciicircum{} over NN/ZZ/QQ/RR/CC) and a generic ring s ((ADD s)/(MUL s)/(NEG s), carrier (CARR s)).  No arg = outermost ring subterm; "term" targets a specific one.  Sound by cut + crs + eq-subst (no new kernel rule); needs the subterm's generators typed in context (true post-di), else refuses and names them. \\
declared in \texttt{tactics-help.scm}

\smallskip

\noindent\textbf{\texttt{simple-integral(omega, mu, s)}}\quad\textit{functoid}\index{simple-integral(omega, mu, s)@\texttt{simple-integral(omega, mu, s)}}\\
reads: the integral of the simple function \$3 over \$1 with respect to \$2 \\
definition:
\begin{quote}\ttfamily\footnotesize\raggedright\noindent
finsum(\allowbreak{}rr-pos-star-add-monoid, vnb-lambda(\allowbreak{}siy\_, image(\allowbreak{}s, omega), etimes(\allowbreak{}siy\_, mu(\allowbreak{}\{six\_ in omega: s(\allowbreak{}six\_) = siy\_\}))), image(\allowbreak{}s, omega))
\end{quote}
characterized by \texttt{simple-integral-in}  \textit{(proven)}:
\begin{quote}\ttfamily\footnotesize\raggedright\noindent
forall(\allowbreak{}[omega, ca, mu, s], is-measure(\allowbreak{}omega, ca, mu) implies is-simple-fn(\allowbreak{}omega, ca, s) implies simple-integral(\allowbreak{}omega, mu, s) in rr-pos-star)
\end{quote}
also: \texttt{simple-integral-mono} \texttt{simple-integral-unfold} \texttt{simple-integral-additive} \texttt{simple-integral-indicator} \texttt{simple-integral-partition} \texttt{simple-integral-measure-le} \texttt{simple-integral-homogeneous} \\
mentioned by 19 result(s)

\smallskip

\noindent\textbf{\texttt{simply-connected(hmu\_)}}\quad\textit{predicate}\index{simply-connected(hmu_)@\texttt{simply-connected(hmu\_)}}\\
reads: \$1 is simply connected \\
defined by \texttt{simply-connected} \\
characterized by \texttt{simply-connected}  \textit{(definitional)}:
\begin{quote}\ttfamily\footnotesize\raggedright\noindent
forall(\allowbreak{}[hmu\_], simply-connected(\allowbreak{}hmu\_) iff is-open(\allowbreak{}nf-metric-space(\allowbreak{}cc-normed-field), hmu\_) and is-connected(\allowbreak{}cc-ms, hmu\_) and forall(\allowbreak{}[hma\_, hmb\_, hmg\_], is-path(\allowbreak{}hmg\_, hma\_, hmb\_) implies trace(\allowbreak{}hmg\_, hma\_, hmb\_) subset hmu\_ implies hmg\_(\allowbreak{}hmb\_) = hmg\_(\allowbreak{}hma\_) implies null-homotopic(\allowbreak{}hmu\_, hmg\_, hma\_, hmb\_)))
\end{quote}
mentioned by 12 result(s)

\smallskip

\noindent\textbf{\texttt{sin}}\quad\textit{operator}\index{sin@\texttt{sin}}\\
a term-forming head

\smallskip

\noindent\textbf{\texttt{singleton(y\_)}}\quad\textit{functoid}\index{singleton(y_)@\texttt{singleton(y\_)}}\\
definition:
\begin{quote}\ttfamily\footnotesize\raggedright\noindent
\{y\_\}
\end{quote}
characterized by \texttt{singleton-is-set}  \textit{(proven)}:
\begin{quote}\ttfamily\footnotesize\raggedright\noindent
forall(\allowbreak{}[y in set], singleton(\allowbreak{}y) in set)
\end{quote}
also: \texttt{singleton-unfold} \texttt{singleton-membership} \texttt{singleton-connected-t} \texttt{key-ideal-gen-singleton-least} \texttt{top-hausdorff-singleton-closed} \\
mentioned by 145 result(s)

\smallskip

\noindent\textbf{\texttt{sip-as-module(s)}}\quad\textit{view}\index{sip-as-module(s)@\texttt{sip-as-module(s)}}\\
a semi-inner-product-space viewed as a module \\
characterized by \texttt{mvag-id-sip-as-module}  \textit{(definitional)}:
\begin{quote}\ttfamily\footnotesize\raggedright\noindent
forall(\allowbreak{}[r], is-semi-inner-product-space(\allowbreak{}r) implies iden(\allowbreak{}module-vector-ag(\allowbreak{}sip-as-module(\allowbreak{}r))) = vzero(\allowbreak{}r))
\end{quote}
also: \texttt{mvag-op-sip-as-module} \texttt{mvag-carr-sip-as-module} \texttt{qmod-comm-sip-as-module} \texttt{sip-as-module-is-module} \texttt{qmod-assoc-sip-as-module} \texttt{spans-span-sip-as-module} \texttt{hbs-lf-zero-sip-as-module} \texttt{hbs-sl-zero-sip-as-module} ... (310 in all) \\
mentioned by 861 result(s)

\smallskip

\noindent\textbf{\texttt{sip-metric-space(v)}}\quad\textit{functoid}\index{sip-metric-space(v)@\texttt{sip-metric-space(v)}}\\
reads: the metric space of \$1 \\
definition:
\begin{quote}\ttfamily\footnotesize\raggedright\noindent
nag-metric-space(\allowbreak{}sip-normed-ag(\allowbreak{}v))
\end{quote}
characterized by \texttt{sip-metric-space-pts}  \textit{(proven)}:
\begin{quote}\ttfamily\footnotesize\raggedright\noindent
forall(\allowbreak{}[v], is-pre-hilbert-space(\allowbreak{}v) implies pts(\allowbreak{}sip-metric-space(\allowbreak{}v)) == vec(\allowbreak{}v))
\end{quote}
also: \texttt{sip-metric-space-dist} \texttt{sip-metric-space-unfold} \texttt{sip-metric-space-is-metric-space} \\
mentioned by 25 result(s)

\smallskip

\noindent\textbf{\texttt{sip-normed-ag(v)}}\quad\textit{functoid}\index{sip-normed-ag(v)@\texttt{sip-normed-ag(v)}}\\
reads: the normed additive group of \$1 \\
definition:
\begin{quote}\ttfamily\footnotesize\raggedright\noindent
[vec(\allowbreak{}v), vadd(\allowbreak{}v), vzero(\allowbreak{}v), vneg(\allowbreak{}v), vnb-lambda(\allowbreak{}shn\_, vec(\allowbreak{}v), sqrt(\allowbreak{}(\allowbreak{}rip(\allowbreak{}v))(\allowbreak{}shn\_, shn\_)))]
\end{quote}
characterized by \texttt{sip-normed-ag-inv}  \textit{(proven)}:
\begin{quote}\ttfamily\footnotesize\raggedright\noindent
forall(\allowbreak{}[v], is-semi-inner-product-space(\allowbreak{}v) implies inv(\allowbreak{}sip-normed-ag(\allowbreak{}v)) == vneg(\allowbreak{}v))
\end{quote}
also: \texttt{sip-normed-ag-opr} \texttt{sip-normed-ag-carrier} \texttt{sip-normed-ag-norm-value} \texttt{sip-normed-ag-is-normed-ag} \\
mentioned by 11 result(s)

\smallskip

\noindent\textbf{\texttt{sip-null(v)}}\quad\textit{functoid}\index{sip-null(v)@\texttt{sip-null(v)}}\\
reads: the null vectors of \$1 \\
definition:
\begin{quote}\ttfamily\footnotesize\raggedright\noindent
\{sqn\_ in vec(\allowbreak{}v): (\allowbreak{}rip(\allowbreak{}v))(\allowbreak{}sqn\_, sqn\_) = 0\}
\end{quote}
characterized by \texttt{sip-null-in-csip-as-sip}  \textit{(definitional)}:
\begin{quote}\ttfamily\footnotesize\raggedright\noindent
forall(\allowbreak{}[r], is-complex-semi-inner-product-space(\allowbreak{}r) implies forall(\allowbreak{}[x\_ in vec(\allowbreak{}r)], (\allowbreak{}rip(\allowbreak{}r))(\allowbreak{}x\_, x\_) = 0 implies x\_ in sip-null(\allowbreak{}csip-as-sip(\allowbreak{}r))))
\end{quote}
also: \texttt{sip-null-out-csip-as-sip} \texttt{sip-null-subspace-csip-as-sip} \texttt{sip-null-in} \texttt{sip-null-out} \texttt{sip-null-unfold} \texttt{sip-null-subspace} \\
mentioned by 8 result(s)

\smallskip

\noindent\textbf{\texttt{sip-perp(v, ks)}}\quad\textit{functoid}\index{sip-perp(v, ks)@\texttt{sip-perp(v, ks)}}\\
reads: the orthogonal complement of \$2 in \$1 \\
definition:
\begin{quote}\ttfamily\footnotesize\raggedright\noindent
\{spp\_ in vec(\allowbreak{}v): forall(\allowbreak{}[spk\_ in ks], (\allowbreak{}rip(\allowbreak{}v))(\allowbreak{}spp\_, spk\_) = 0)\}
\end{quote}
characterized by \texttt{sip-perp-orth}  \textit{(proven)}:
\begin{quote}\ttfamily\footnotesize\raggedright\noindent
forall(\allowbreak{}[v, ks, z\_, k\_], z\_ in sip-perp(\allowbreak{}v, ks) implies k\_ in ks implies (\allowbreak{}rip(\allowbreak{}v))(\allowbreak{}z\_, k\_) = 0)
\end{quote}
also: \texttt{sip-perp-subset} \texttt{sip-perp-unfold} \texttt{sip-perp-invariant} \texttt{sip-perp-membership} \texttt{sip-perp-is-submodule} \\
mentioned by 15 result(s)

\smallskip

\noindent\textbf{\texttt{sip-qact(v)}}\quad\textit{functoid}\index{sip-qact(v)@\texttt{sip-qact(v)}}\\
definition:
\begin{quote}\ttfamily\footnotesize\raggedright\noindent
vnb-lambda(\allowbreak{}[sqr\_, sqc\_], cartesian(\allowbreak{}rr, quotient(\allowbreak{}sip-setoid(\allowbreak{}v))), (\allowbreak{}sip-qscale(\allowbreak{}v, sqr\_))(\allowbreak{}sqc\_))
\end{quote}
mentioned by 14 result(s)

\smallskip

\noindent\textbf{\texttt{sip-qrip(v)}}\quad\textit{functoid}\index{sip-qrip(v)@\texttt{sip-qrip(v)}}\\
definition:
\begin{quote}\ttfamily\footnotesize\raggedright\noindent
descend2(\allowbreak{}sip-setoid(\allowbreak{}v), vnb-lambda(\allowbreak{}[sqa\_, sqb\_], cartesian(\allowbreak{}vec(\allowbreak{}v), vec(\allowbreak{}v)), (\allowbreak{}rip(\allowbreak{}v))(\allowbreak{}sqa\_, sqb\_)))
\end{quote}
mentioned by 14 result(s)

\smallskip

\noindent\textbf{\texttt{sip-qscale(v, r)}}\quad\textit{functoid}\index{sip-qscale(v, r)@\texttt{sip-qscale(v, r)}}\\
definition:
\begin{quote}\ttfamily\footnotesize\raggedright\noindent
descend(\allowbreak{}sip-setoid(\allowbreak{}v), vnb-lambda(\allowbreak{}sqa\_, vec(\allowbreak{}v), class(\allowbreak{}sip-setoid(\allowbreak{}v), (\allowbreak{}act(\allowbreak{}v))(\allowbreak{}r, sqa\_))))
\end{quote}
mentioned by 8 result(s)

\smallskip

\noindent\textbf{\texttt{sip-quotient(v)}}\quad\textit{functoid}\index{sip-quotient(v)@\texttt{sip-quotient(v)}}\\
reads: the quotient of \$1 by its null vectors \\
definition:
\begin{quote}\ttfamily\footnotesize\raggedright\noindent
[scal(\allowbreak{}v), quotient(\allowbreak{}sip-setoid(\allowbreak{}v)), sip-qvadd(\allowbreak{}v), class(\allowbreak{}sip-setoid(\allowbreak{}v), vzero(\allowbreak{}v)), sip-qvneg(\allowbreak{}v), sip-qact(\allowbreak{}v), sip-qrip(\allowbreak{}v)]
\end{quote}
characterized by \texttt{sip-quotient-is-pre-hilbert-csip-as-sip}  \textit{(definitional)}:
\begin{quote}\ttfamily\footnotesize\raggedright\noindent
forall(\allowbreak{}[r], is-complex-semi-inner-product-space(\allowbreak{}r) implies is-pre-hilbert-space(\allowbreak{}sip-quotient(\allowbreak{}csip-as-sip(\allowbreak{}r))))
\end{quote}
also: \texttt{sip-quotient-complete} \texttt{sip-quotient-is-pre-hilbert} \\
mentioned by 182 result(s)

\smallskip

\noindent\textbf{\texttt{sip-qvadd(v)}}\quad\textit{functoid}\index{sip-qvadd(v)@\texttt{sip-qvadd(v)}}\\
definition:
\begin{quote}\ttfamily\footnotesize\raggedright\noindent
descend2(\allowbreak{}sip-setoid(\allowbreak{}v), vnb-lambda(\allowbreak{}[sqa\_, sqb\_], cartesian(\allowbreak{}vec(\allowbreak{}v), vec(\allowbreak{}v)), class(\allowbreak{}sip-setoid(\allowbreak{}v), (\allowbreak{}vadd(\allowbreak{}v))(\allowbreak{}sqa\_, sqb\_))))
\end{quote}
mentioned by 14 result(s)

\smallskip

\noindent\textbf{\texttt{sip-qvneg(v)}}\quad\textit{functoid}\index{sip-qvneg(v)@\texttt{sip-qvneg(v)}}\\
definition:
\begin{quote}\ttfamily\footnotesize\raggedright\noindent
descend(\allowbreak{}sip-setoid(\allowbreak{}v), vnb-lambda(\allowbreak{}sqa\_, vec(\allowbreak{}v), class(\allowbreak{}sip-setoid(\allowbreak{}v), (\allowbreak{}vneg(\allowbreak{}v))(\allowbreak{}sqa\_))))
\end{quote}
mentioned by 14 result(s)

\smallskip

\noindent\textbf{\texttt{sip-rel(v)}}\quad\textit{functoid}\index{sip-rel(v)@\texttt{sip-rel(v)}}\\
reads: the null-difference relation of \$1 \\
definition:
\begin{quote}\ttfamily\footnotesize\raggedright\noindent
\{sqp\_ in cartesian(\allowbreak{}vec(\allowbreak{}v), vec(\allowbreak{}v)): (\allowbreak{}rip(\allowbreak{}v))(\allowbreak{}(\allowbreak{}vadd(\allowbreak{}v))(\allowbreak{}nth(\allowbreak{}1, sqp\_), (\allowbreak{}vneg(\allowbreak{}v))(\allowbreak{}nth(\allowbreak{}2, sqp\_))), (\allowbreak{}vadd(\allowbreak{}v))(\allowbreak{}nth(\allowbreak{}1, sqp\_), (\allowbreak{}vneg(\allowbreak{}v))(\allowbreak{}nth(\allowbreak{}2, sqp\_)))) = 0\}
\end{quote}
characterized by \texttt{sip-rel-in-csip-as-sip}  \textit{(definitional)}:
\begin{quote}\ttfamily\footnotesize\raggedright\noindent
forall(\allowbreak{}[r], is-complex-semi-inner-product-space(\allowbreak{}r) implies forall(\allowbreak{}[a\_ in vec(\allowbreak{}r), b\_ in vec(\allowbreak{}r)], (\allowbreak{}rip(\allowbreak{}r))(\allowbreak{}(\allowbreak{}vadd(\allowbreak{}r))(\allowbreak{}a\_, (\allowbreak{}vneg(\allowbreak{}r))(\allowbreak{}b\_)), (\allowbreak{}vadd(\allowbreak{}r))(\allowbreak{}a\_, (\allowbreak{}vneg(\allowbreak{}r))(\allowbreak{}b\_))) = 0 implies [a\_, b\_] in sip-rel(\allowbreak{}csip-as-sip(\allowbreak{}r))))
\end{quote}
also: \texttt{sip-rel-out-csip-as-sip} \texttt{sip-rel-is-equivalence-csip-as-sip} \texttt{sip-rel-in} \texttt{sip-rel-out} \texttt{sip-rel-unfold} \texttt{sip-rel-is-equivalence} \\
mentioned by 20 result(s)

\smallskip

\noindent\textbf{\texttt{sip-setoid(v)}}\quad\textit{functoid}\index{sip-setoid(v)@\texttt{sip-setoid(v)}}\\
reads: the setoid of \$1 modulo its null vectors \\
definition:
\begin{quote}\ttfamily\footnotesize\raggedright\noindent
[vec(\allowbreak{}v), sip-rel(\allowbreak{}v)]
\end{quote}
characterized by \texttt{sip-setoid-pts-csip-as-sip}  \textit{(definitional)}:
\begin{quote}\ttfamily\footnotesize\raggedright\noindent
forall(\allowbreak{}[r], is-complex-semi-inner-product-space(\allowbreak{}r) implies pts(\allowbreak{}sip-setoid(\allowbreak{}csip-as-sip(\allowbreak{}r))) == vec(\allowbreak{}r))
\end{quote}
also: \texttt{sip-setoid-rel-csip-as-sip} \texttt{sip-setoid-is-setoid-csip-as-sip} \texttt{sip-setoid-pts} \texttt{sip-setoid-rel} \texttt{sip-setoid-is-setoid} \\
mentioned by 145 result(s)

\smallskip

\noindent\textbf{\texttt{size(m)}}\quad\textit{functoid}\index{size(m)@\texttt{size(m)}}\\
definition:
\begin{quote}\ttfamily\footnotesize\raggedright\noindent
[length(\allowbreak{}m), if(\allowbreak{}length(\allowbreak{}m) = 0, 0, length(\allowbreak{}nth(\allowbreak{}1, m)))]
\end{quote}
characterized by \texttt{mat-size}  \textit{(proven)}:
\begin{quote}\ttfamily\footnotesize\raggedright\noindent
forall(\allowbreak{}[m, n, x, q in mat(\allowbreak{}m, n, x)], 1 <= m implies size(\allowbreak{}q) = [m, n])
\end{quote}
also: \texttt{size-unfold} \texttt{mat-size-cols} \texttt{mat-size-rows} \texttt{size-nonempty} \texttt{mat-size-cols-pos} \texttt{mat-size-cols-in-nn} \\
mentioned by 16 result(s)

\smallskip

\noindent\textbf{\texttt{sk-inf(s)}}\quad\textit{accessor}\index{sk-inf(s)@\texttt{sk-inf(s)}}\\
reads: the sk-inf slot of \$1 \\
a structure slot \\
characterized by \texttt{cx-orp@sk-inf}  \textit{(definitional)}:
\begin{quote}\ttfamily\footnotesize\raggedright\noindent
sk-inf(\allowbreak{}cx-orp) == cx-isp
\end{quote}
mentioned by 38 result(s)

\smallskip

\noindent\textbf{\texttt{sk-names(s)}}\quad\textit{accessor}\index{sk-names(s)@\texttt{sk-names(s)}}\\
reads: the sk-names slot of \$1 \\
a structure slot \\
characterized by \texttt{cx-orp@sk-names}  \textit{(definitional)}:
\begin{quote}\ttfamily\footnotesize\raggedright\noindent
sk-names(\allowbreak{}cx-orp) == interval(\allowbreak{}10, 11)
\end{quote}
mentioned by 25 result(s)

\smallskip

\noindent\textbf{\texttt{sk-space}}\quad\textit{structure}\index{sk-space@\texttt{sk-space}}\\
predicate \texttt{is-sk-space}; slots: sk-inf sk-names ltkey

\smallskip

\noindent\textbf{\texttt{slot}}\quad\textit{tactic}\index{slot@\texttt{slot}}\\
\texttt{(slot 'acc)} -- Reduce a structure ACCESSOR to its projection: (CARR s) becomes (NTH 1 s).  The one door for accessor reductions -- use it, never mac, on an accessor name. \\
declared in \texttt{tactics-help.scm}

\smallskip

\noindent\textbf{\texttt{slot-h}}\quad\textit{tactic}\index{slot-h@\texttt{slot-h}}\\
\texttt{(slot-h 'acc hyp)} -- Reduce a structure ACCESSOR to its projection inside a cited HYPOTHESIS -- `slot', hypothesis-side. \\
declared in \texttt{tactics-help.scm}

\smallskip

\noindent\textbf{\texttt{smith-staircase(a, m, n, d, k)}}\quad\textit{predicate}\index{smith-staircase(a, m, n, d, k)@\texttt{smith-staircase(a, m, n, d, k)}}\\
reads: \$4 is in Smith staircase form of rank \$5, as an \$2-by-\$3 matrix over \$1 \\
defined by \texttt{smith-staircase} \\
characterized by \texttt{smith-staircase}  \textit{(definitional)}:
\begin{quote}\ttfamily\footnotesize\raggedright\noindent
forall(\allowbreak{}[a, m, n, d, k], smith-staircase(\allowbreak{}a, m, n, d, k) iff is-diagonal(\allowbreak{}a, m, n, d) and k <= m and k <= n and forall(\allowbreak{}[i\_ in interval(\allowbreak{}1, k)], not(\allowbreak{}entry(\allowbreak{}d, i\_, i\_) = zero(\allowbreak{}a))) and forall(\allowbreak{}[i\_ in interval(\allowbreak{}1, m)], not(\allowbreak{}i\_ <= k) implies forall(\allowbreak{}[j\_ in interval(\allowbreak{}1, n)], entry(\allowbreak{}d, i\_, j\_) = zero(\allowbreak{}a))))
\end{quote}
mentioned by 6 result(s)

\smallskip

\noindent\textbf{\texttt{sn-ball(m, pf, ff, x, eps)}}\quad\textit{functoid}\index{sn-ball(m, pf, ff, x, eps)@\texttt{sn-ball(m, pf, ff, x, eps)}}\\
reads: the basic neighbourhood of \$4 of radius \$5 for the seminorms \$3 of the family \$2 \\
definition:
\begin{quote}\ttfamily\footnotesize\raggedright\noindent
\{snby\_ in vec(\allowbreak{}m): forall(\allowbreak{}[sni\_ in ff], (\allowbreak{}pf(\allowbreak{}sni\_))(\allowbreak{}(\allowbreak{}vadd(\allowbreak{}m))(\allowbreak{}snby\_, (\allowbreak{}vneg(\allowbreak{}m))(\allowbreak{}x))) <= eps)\}
\end{quote}
characterized by \texttt{sn-ball-mono}  \textit{(proven)}:
\begin{quote}\ttfamily\footnotesize\raggedright\noindent
forall(\allowbreak{}[fld, m, ix, pf, ff, gg, x, e1, e2], is-seminorm-family-on(\allowbreak{}fld, m, ix, pf) implies ff subset gg implies gg subset ix implies x in vec(\allowbreak{}m) implies e1 in rr implies e2 in rr implies e1 <= e2 implies sn-ball(\allowbreak{}m, pf, gg, x, e1) subset sn-ball(\allowbreak{}m, pf, ff, x, e2))
\end{quote}
also: \texttt{sn-ball-unfold} \texttt{sn-ball-membership} \\
mentioned by 9 result(s)

\smallskip

\noindent\textbf{\texttt{sn-oball(m, pf, ff, x, r)}}\quad\textit{functoid}\index{sn-oball(m, pf, ff, x, r)@\texttt{sn-oball(m, pf, ff, x, r)}}\\
reads: the open basic neighbourhood of \$4 of radius \$5 for the seminorms \$3 of the family \$2 \\
definition:
\begin{quote}\ttfamily\footnotesize\raggedright\noindent
\{snoy\_ in vec(\allowbreak{}m): forsome(\allowbreak{}[snoe\_ in rr], pos-rr(\allowbreak{}snoe\_) and forall(\allowbreak{}[sni\_ in ff], (\allowbreak{}pf(\allowbreak{}sni\_))(\allowbreak{}(\allowbreak{}vadd(\allowbreak{}m))(\allowbreak{}snoy\_, (\allowbreak{}vneg(\allowbreak{}m))(\allowbreak{}x))) + snoe\_ <= r))\}
\end{quote}
characterized by \texttt{sn-oball-open}  \textit{(proven)}:
\begin{quote}\ttfamily\footnotesize\raggedright\noindent
forall(\allowbreak{}[fld, m, ix, pf, ff, x, r], is-seminorm-family-on(\allowbreak{}fld, m, ix, pf) implies ff in set implies card(\allowbreak{}ff) in nn implies ff subset ix implies x in vec(\allowbreak{}m) implies r in rr implies is-sn-open(\allowbreak{}m, ix, pf, sn-oball(\allowbreak{}m, pf, ff, x, r)))
\end{quote}
also: \texttt{sn-oball-center} \texttt{sn-oball-convex} \texttt{sn-oball-unfold} \texttt{sn-oball-in-ball} \texttt{sn-oball-membership} \\
mentioned by 9 result(s)

\smallskip

\noindent\textbf{\texttt{snd-iter}}\quad\textit{functoid}\index{snd-iter@\texttt{snd-iter}}\\
reads: the \$2-th iterate of the second unpairing at \$1 \\
defined by \texttt{snd-iter-zero} \texttt{snd-iter-succ} \\
characterized by \texttt{snd-iter-succ}  \textit{(definitional)}:
\begin{quote}\ttfamily\footnotesize\raggedright\noindent
forall(\allowbreak{}[sim\_, sin\_ in nn], snd-iter(\allowbreak{}sim\_, succ(\allowbreak{}sin\_)) == nnsnd(\allowbreak{}snd-iter(\allowbreak{}sim\_, sin\_)))
\end{quote}
also: \texttt{snd-iter-zero} \texttt{snd-iter-type} \texttt{snd-iter-shift} \\
mentioned by 20 result(s)

\smallskip

\noindent\textbf{\texttt{snoc-col(w, n, x)}}\quad\textit{functoid}\index{snoc-col(w, n, x)@\texttt{snoc-col(w, n, x)}}\\
definition:
\begin{quote}\ttfamily\footnotesize\raggedright\noindent
matof(\allowbreak{}succ(\allowbreak{}n), 1, vnb-lambda(\allowbreak{}[i\_, j\_], cartesian(\allowbreak{}interval(\allowbreak{}1, succ(\allowbreak{}n)), interval(\allowbreak{}1, 1)), if(\allowbreak{}i\_ = succ(\allowbreak{}n), x, entry(\allowbreak{}w, i\_, 1))))
\end{quote}
characterized by \texttt{snoc-col-last}  \textit{(proven)}:
\begin{quote}\ttfamily\footnotesize\raggedright\noindent
forall(\allowbreak{}[w, n, v, x], n in nn implies w in mat(\allowbreak{}n, 1, x) implies v in x implies entry(\allowbreak{}snoc-col(\allowbreak{}w, n, v), succ(\allowbreak{}n), 1) = v)
\end{quote}
also: \texttt{snoc-col-type} \texttt{entry-of-snoc-col} \\
mentioned by 39 result(s)

\smallskip

\noindent\textbf{\texttt{snoc-row(c, n, r)}}\quad\textit{functoid}\index{snoc-row(c, n, r)@\texttt{snoc-row(c, n, r)}}\\
definition:
\begin{quote}\ttfamily\footnotesize\raggedright\noindent
matof(\allowbreak{}1, succ(\allowbreak{}n), vnb-lambda(\allowbreak{}[i\_, j\_], cartesian(\allowbreak{}interval(\allowbreak{}1, 1), interval(\allowbreak{}1, succ(\allowbreak{}n))), if(\allowbreak{}j\_ = succ(\allowbreak{}n), r, entry(\allowbreak{}c, 1, j\_))))
\end{quote}
characterized by \texttt{snoc-row-last}  \textit{(proven)}:
\begin{quote}\ttfamily\footnotesize\raggedright\noindent
forall(\allowbreak{}[c, n, r, x], n in nn implies c in mat(\allowbreak{}1, n, x) implies r in x implies entry(\allowbreak{}snoc-row(\allowbreak{}c, n, r), 1, succ(\allowbreak{}n)) = r)
\end{quote}
also: \texttt{snoc-row-type} \texttt{entry-of-snoc-row} \\
mentioned by 39 result(s)

\smallskip

\noindent\textbf{\texttt{snowflake(snfs\_, snfe\_)}}\quad\textit{functoid}\index{snowflake(snfs_, snfe_)@\texttt{snowflake(snfs\_, snfe\_)}}\\
reads: the metric space \$1 with its distance raised to the power \$2 \\
definition:
\begin{quote}\ttfamily\footnotesize\raggedright\noindent
[pts(\allowbreak{}snfs\_), vnb-lambda(\allowbreak{}[snfu\_, snfv\_], cartesian(\allowbreak{}pts(\allowbreak{}snfs\_), pts(\allowbreak{}snfs\_)), rpow-star(\allowbreak{}(\allowbreak{}dist(\allowbreak{}snfs\_))(\allowbreak{}snfu\_, snfv\_), snfe\_))]
\end{quote}
characterized by \texttt{snowflake-pts}  \textit{(proven)}:
\begin{quote}\ttfamily\footnotesize\raggedright\noindent
forall(\allowbreak{}[s, e], is-metric-space(\allowbreak{}s) implies e in rr implies 0 <= e implies pts(\allowbreak{}snowflake(\allowbreak{}s, e)) == pts(\allowbreak{}s))
\end{quote}
also: \texttt{snowflake-dist} \texttt{snowflake-dist-fun} \texttt{snowflake-is-metric-space} \\
mentioned by 7 result(s)

\smallskip

\noindent\textbf{\texttt{sos}}\quad\textit{tactic}\index{sos@\texttt{sos}}\\
\texttt{(sos "c1" "c2" ...)} -- Sum-of-squares closer for a nonstrict polynomial inequality a <= b over RR (the nonlinear companion of (ineq)).  You supply the terms to be SQUARED; it finds the nonnegative coefficients.  (tactics 'sos) for the worked example. \\
declared in \texttt{tactics-help.scm}

\smallskip

\noindent\textbf{\texttt{sot(otv\_)}}\quad\textit{functoid}\index{sot(otv_)@\texttt{sot(otv\_)}}\\
reads: the strong operator topology on the bounded operators on \$1 \\
definition:
\begin{quote}\ttfamily\footnotesize\raggedright\noindent
seminorm-top(\allowbreak{}bh-module(\allowbreak{}otv\_), vec(\allowbreak{}otv\_), sot-pf(\allowbreak{}otv\_))
\end{quote}
characterized by \texttt{sot-is-tvs}  \textit{(proven)}:
\begin{quote}\ttfamily\footnotesize\raggedright\noindent
forall(\allowbreak{}[v], is-complex-inner-product-space(\allowbreak{}v) implies is-tvs(\allowbreak{}cc-normed-field, bh-module(\allowbreak{}v), sot(\allowbreak{}v)))
\end{quote}
also: \texttt{sot-unfold} \texttt{commutant-sot-closed} \texttt{sot-to-wot-continuous} \texttt{norm-to-sot-continuous} \texttt{strong-dual-unfold-sot} \texttt{bicommutant-in-sot-closure} \texttt{bicommutant-is-sot-closure} \texttt{sot-closure-in-bicommutant} ... (11 in all) \\
mentioned by 19 result(s)

\smallskip

\noindent\textbf{\texttt{sot-pf(otv\_)}}\quad\textit{functoid}\index{sot-pf(otv_)@\texttt{sot-pf(otv\_)}}\\
reads: the seminorms T |-> ||T xi|| of the strong operator topology on \$1 \\
definition:
\begin{quote}\ttfamily\footnotesize\raggedright\noindent
vnb-lambda(\allowbreak{}otx\_, vec(\allowbreak{}otv\_), vnb-lambda(\allowbreak{}ott\_, bounded-ops(\allowbreak{}otv\_, otv\_), ip-norm(\allowbreak{}otv\_, ott\_(\allowbreak{}otx\_))))
\end{quote}
characterized by \texttt{sot-pf-type}  \textit{(proven)}:
\begin{quote}\ttfamily\footnotesize\raggedright\noindent
forall(\allowbreak{}[v], is-complex-inner-product-space(\allowbreak{}v) implies sot-pf(\allowbreak{}v) in fun(\allowbreak{}vec(\allowbreak{}v), fun(\allowbreak{}bounded-ops(\allowbreak{}v, v), rr)))
\end{quote}
also: \texttt{sot-pf-apply} \\
mentioned by 9 result(s)

\smallskip

\noindent\textbf{\texttt{souslin-set(s, sv)}}\quad\textit{functoid}\index{souslin-set(s, sv)@\texttt{souslin-set(s, sv)}}\\
reads: the Souslin set of the sieve \$2 in \$1 \\
definition:
\begin{quote}\ttfamily\footnotesize\raggedright\noindent
\{sxp\_ in pts(\allowbreak{}s): forsome(\allowbreak{}[sxz\_ in fun(\allowbreak{}nn, nn)], forall(\allowbreak{}[sxn\_ in nn], sxp\_ in sv(\allowbreak{}restrict(\allowbreak{}sxz\_, ord-segment(\allowbreak{}sxn\_)))))\}
\end{quote}
characterized by \texttt{souslin-set-unfold}  \textit{(proven)}:
\begin{quote}\ttfamily\footnotesize\raggedright\noindent
forall(\allowbreak{}[s, sv], souslin-set(\allowbreak{}s, sv) == \{sxp\_ in pts(\allowbreak{}s): forsome(\allowbreak{}[sxz\_ in fun(\allowbreak{}nn, nn)], forall(\allowbreak{}[sxn\_ in nn], sxp\_ in sv(\allowbreak{}restrict(\allowbreak{}sxz\_, ord-segment(\allowbreak{}sxn\_)))))\})
\end{quote}
mentioned by 8 result(s)

\smallskip

\noindent\textbf{\texttt{sp}}\quad\textit{tactic}\index{sp@\texttt{sp}}\\
\texttt{(sp goal)} -- Start a proof of `goal' (a "string", a raw S-expr, or a wff); clears the script. \\
declared in \texttt{tactics-help.scm}

\smallskip

\noindent\textbf{\texttt{span(md, n, u)}}\quad\textit{functoid}\index{span(md, n, u)@\texttt{span(md, n, u)}}\\
definition:
\begin{quote}\ttfamily\footnotesize\raggedright\noindent
\{x\_ in vec(\allowbreak{}md): forsome(\allowbreak{}[c\_ in mat(\allowbreak{}1, n, carr(\allowbreak{}scal(\allowbreak{}md)))], x\_ = lincomb(\allowbreak{}md, n, c\_, u))\}
\end{quote}
characterized by \texttt{span-membership}  \textit{(definitional)}:
\begin{quote}\ttfamily\footnotesize\raggedright\noindent
forall(\allowbreak{}[md, n, u, x\_], x\_ in span(\allowbreak{}md, n, u) iff x\_ in vec(\allowbreak{}md) and forsome(\allowbreak{}[c\_ in mat(\allowbreak{}1, n, carr(\allowbreak{}scal(\allowbreak{}md)))], x\_ = lincomb(\allowbreak{}md, n, c\_, u)))
\end{quote}
also: \texttt{spans-span-cna-as-module} \texttt{spans-span-sip-as-module} \texttt{spans-span-uba-as-module} \texttt{spans-span-cnvs-as-module} \texttt{span-is-submodule-cna-as-module} \texttt{span-is-submodule-sip-as-module} \texttt{span-is-submodule-uba-as-module} \texttt{span-is-submodule-cnvs-as-module} ... (23 in all) \\
mentioned by 35 result(s)

\smallskip

\noindent\textbf{\texttt{span-add-one(m, s, v)}}\quad\textit{functoid}\index{span-add-one(m, s, v)@\texttt{span-add-one(m, s, v)}}\\
definition:
\begin{quote}\ttfamily\footnotesize\raggedright\noindent
\{y\_ in vec(\allowbreak{}m): forsome(\allowbreak{}[x\_ in s, r\_ in rr], y\_ = (\allowbreak{}vadd(\allowbreak{}m))(\allowbreak{}x\_, (\allowbreak{}act(\allowbreak{}m))(\allowbreak{}r\_, v)))\}
\end{quote}
characterized by \texttt{span-add-one-unfold}  \textit{(proven)}:
\begin{quote}\ttfamily\footnotesize\raggedright\noindent
forall(\allowbreak{}[m, s, v], span-add-one(\allowbreak{}m, s, v) == \{y\_ in vec(\allowbreak{}m): forsome(\allowbreak{}[x\_ in s, r\_ in rr], y\_ = (\allowbreak{}vadd(\allowbreak{}m))(\allowbreak{}x\_, (\allowbreak{}act(\allowbreak{}m))(\allowbreak{}r\_, v)))\})
\end{quote}
also: \texttt{span-add-one-has-v-nvs} \texttt{span-add-one-submodule} \texttt{span-add-one-membership} \texttt{span-add-one-superset-nvs} \\
mentioned by 13 result(s)

\smallskip

\noindent\textbf{\texttt{span2(s2m\_, s2e\_, s2x\_)}}\quad\textit{functoid}\index{span2(s2m_, s2e_, s2x_)@\texttt{span2(s2m\_, s2e\_, s2x\_)}}\\
reads: the subspace of \$1 generated by \$2 and \$3 \\
definition:
\begin{quote}\ttfamily\footnotesize\raggedright\noindent
\{s2v\_ in vec(\allowbreak{}s2m\_): forsome(\allowbreak{}[s2r\_ in rr, s2s\_ in rr], s2v\_ = (\allowbreak{}vadd(\allowbreak{}s2m\_))(\allowbreak{}(\allowbreak{}act(\allowbreak{}s2m\_))(\allowbreak{}s2r\_, s2e\_), (\allowbreak{}act(\allowbreak{}s2m\_))(\allowbreak{}s2s\_, s2x\_)))\}
\end{quote}
characterized by \texttt{span2-unfold}  \textit{(proven)}:
\begin{quote}\ttfamily\footnotesize\raggedright\noindent
forall(\allowbreak{}[m, eta, xi], span2(\allowbreak{}m, eta, xi) == \{s2v\_ in vec(\allowbreak{}m): forsome(\allowbreak{}[s2r\_ in rr, s2s\_ in rr], s2v\_ = (\allowbreak{}vadd(\allowbreak{}m))(\allowbreak{}(\allowbreak{}act(\allowbreak{}m))(\allowbreak{}s2r\_, eta), (\allowbreak{}act(\allowbreak{}m))(\allowbreak{}s2s\_, xi)))\})
\end{quote}
mentioned by 3 result(s)

\smallskip

\noindent\textbf{\texttt{spans(md, n, u, sm)}}\quad\textit{predicate}\index{spans(md, n, u, sm)@\texttt{spans(md, n, u, sm)}}\\
reads: the \$2 vectors \$3 span \$4 in \$1 \\
defined by \texttt{spans} \\
characterized by \texttt{spans}  \textit{(definitional)}:
\begin{quote}\ttfamily\footnotesize\raggedright\noindent
forall(\allowbreak{}[md, n, u, sm], spans(\allowbreak{}md, n, u, sm) iff forall(\allowbreak{}[j\_ in interval(\allowbreak{}1, n)], entry(\allowbreak{}u, j\_, 1) in sm) and forall(\allowbreak{}[x\_ in sm], forsome(\allowbreak{}[c\_ in mat(\allowbreak{}1, n, carr(\allowbreak{}scal(\allowbreak{}md)))], x\_ = lincomb(\allowbreak{}md, n, c\_, u))))
\end{quote}
also: \texttt{spans-span-cna-as-module} \texttt{spans-span-sip-as-module} \texttt{spans-span-uba-as-module} \texttt{spans-span-cnvs-as-module} \texttt{spans-transport-cna-as-module} \texttt{spans-transport-sip-as-module} \texttt{spans-transport-uba-as-module} \texttt{spans-transport-cnvs-as-module} ... (39 in all) \\
mentioned by 104 result(s)

\smallskip

\noindent\textbf{\texttt{spectral-radius(baa\_, bax\_)}}\quad\textit{functoid}\index{spectral-radius(baa_, bax_)@\texttt{spectral-radius(baa\_, bax\_)}}\\
reads: the spectral radius of \$2 in \$1 \\
definition:
\begin{quote}\ttfamily\footnotesize\raggedright\noindent
einf(\allowbreak{}\{bsw\_ in rr: forsome(\allowbreak{}[bsn\_ in nn], 1 <= bsn\_ and bsw\_ = rpow-star(\allowbreak{}(\allowbreak{}vnrm(\allowbreak{}baa\_))(\allowbreak{}el-pow(\allowbreak{}baa\_, bax\_, bsn\_)), recip-star(\allowbreak{}bsn\_)))\})
\end{quote}
characterized by \texttt{spectral-radius-real}  \textit{(proven)}:
\begin{quote}\ttfamily\footnotesize\raggedright\noindent
forall(\allowbreak{}[a], is-unital-banach-algebra(\allowbreak{}a) implies forall(\allowbreak{}[x in vec(\allowbreak{}a)], spectral-radius(\allowbreak{}a, x) in rr and 0 <= spectral-radius(\allowbreak{}a, x)))
\end{quote}
also: \texttt{spectral-radius-limit} \texttt{hom-spectral-radius-le} \texttt{spectral-radius-unfold} \texttt{b2r-spectral-radius-view} \texttt{cstar-sa-spectral-radius} \texttt{gelfand-spectral-radius-le} \texttt{gelfand-spectral-radius-max} \texttt{gelfand-sup-spectral-radius} ... (10 in all) \\
mentioned by 17 result(s)

\smallskip

\noindent\textbf{\texttt{spectrum(baa\_, bax\_)}}\quad\textit{functoid}\index{spectrum(baa_, bax_)@\texttt{spectrum(baa\_, bax\_)}}\\
reads: the spectrum of \$2 in \$1 \\
definition:
\begin{quote}\ttfamily\footnotesize\raggedright\noindent
\{baz\_ in cc: not(\allowbreak{}is-invertible-el(\allowbreak{}baa\_, shift-el(\allowbreak{}baa\_, baz\_, bax\_)))\}
\end{quote}
characterized by \texttt{bsn-no-spectrum}  \textit{(proven)}:
\begin{quote}\ttfamily\footnotesize\raggedright\noindent
forall(\allowbreak{}[a], is-unital-banach-algebra(\allowbreak{}a) implies forall(\allowbreak{}[x in vec(\allowbreak{}a)], not(\allowbreak{}forsome(\allowbreak{}[bsq\_], bsq\_ in spectrum(\allowbreak{}a, x))) implies forall(\allowbreak{}[z in cc], is-invertible-el(\allowbreak{}a, shift-el(\allowbreak{}a, z, x)))))
\end{quote}
also: \texttt{spectrum-nonempty} \texttt{bas-spectrum-unfold} \texttt{bcp-sa-spectrum-real} \texttt{gelfand-range-spectrum} \texttt{spectrum-in-closed-ball} \texttt{spectrum-character-value} \texttt{character-value-in-spectrum} \texttt{gelfand-mazur-spectrum-point} ... (10 in all) \\
mentioned by 19 result(s)

\smallskip

\noindent\textbf{\texttt{sphere(cuvs\_, cuvc\_, cuvr\_)}}\quad\textit{functoid}\index{sphere(cuvs_, cuvc_, cuvr_)@\texttt{sphere(cuvs\_, cuvc\_, cuvr\_)}}\\
reads: the sphere of \$1 with centre \$2 and radius \$3 \\
definition:
\begin{quote}\ttfamily\footnotesize\raggedright\noindent
\{cuvy\_ in pts(\allowbreak{}cuvs\_): (\allowbreak{}dist(\allowbreak{}cuvs\_))(\allowbreak{}cuvc\_, cuvy\_) = cuvr\_\}
\end{quote}
characterized by \texttt{sphere-unfold}  \textit{(proven)}:
\begin{quote}\ttfamily\footnotesize\raggedright\noindent
forall(\allowbreak{}[s, c, r], sphere(\allowbreak{}s, c, r) == \{cuvy\_ in pts(\allowbreak{}s): (\allowbreak{}dist(\allowbreak{}s))(\allowbreak{}c, cuvy\_) = r\})
\end{quote}
also: \texttt{sphere-membership} \texttt{ultra-sphere-open} \texttt{ultra-sphere-closed} \\
mentioned by 6 result(s)

\smallskip

\noindent\textbf{\texttt{splice}}\quad\textit{operator}\index{splice@\texttt{splice}}\\
a term-forming head \\
mentioned by 4 result(s)

\smallskip

\noindent\textbf{\texttt{sq-summable-seqs(otv\_)}}\quad\textit{functoid}\index{sq-summable-seqs(otv_)@\texttt{sq-summable-seqs(otv\_)}}\\
reads: the square-summable sequences in \$1 \\
definition:
\begin{quote}\ttfamily\footnotesize\raggedright\noindent
\{ots\_ in fun(\allowbreak{}nn, vec(\allowbreak{}otv\_)): esum(\allowbreak{}vnb-lambda(\allowbreak{}otk\_, nn, ip-norm(\allowbreak{}otv\_, ots\_(\allowbreak{}otk\_)) * ip-norm(\allowbreak{}otv\_, ots\_(\allowbreak{}otk\_)))) in rr\}
\end{quote}

\smallskip

\noindent\textbf{\texttt{sqn(a)}}\quad\textit{functoid}\index{sqn(a)@\texttt{sqn(a)}}\\
reads: the sequences in \$1 \\
definition:
\begin{quote}\ttfamily\footnotesize\raggedright\noindent
fun(\allowbreak{}nn, a)
\end{quote}
characterized by \texttt{sqn-sethood}  \textit{(proven)}:
\begin{quote}\ttfamily\footnotesize\raggedright\noindent
forall(\allowbreak{}[a in set], sqn(\allowbreak{}a) in set)
\end{quote}
also: \texttt{sqn-membership} \texttt{sqn-unfold-b6s} \texttt{subsqn-preserves-sqn} \\
mentioned by 28 result(s)

\smallskip

\noindent\textbf{\texttt{sqrt(a\_)}}\quad\textit{functoid}\index{sqrt(a_)@\texttt{sqrt(a\_)}}\\
reads: the square root of \$1 \\
definition:
\begin{quote}\ttfamily\footnotesize\raggedright\noindent
iota(\allowbreak{}x\_, x\_ in rr and 0 <= x\_ and x\_ * x\_ = a\_)
\end{quote}
characterized by \texttt{sqrt-sq}  \textit{(proven)}:
\begin{quote}\ttfamily\footnotesize\raggedright\noindent
forall(\allowbreak{}[a], a in rr and 0 <= a implies sqrt(\allowbreak{}a) * sqrt(\allowbreak{}a) = a)
\end{quote}
also: \texttt{sqrt-mul} \texttt{sqrt-char} \texttt{sqrt-mono} \texttt{sqrt-prop} \texttt{sqrt-rpow} \texttt{sqrt-zero} \texttt{sqrt-of-sq} \texttt{amgm-2-sqrt} ... (19 in all) \\
mentioned by 70 result(s)

\smallskip

\noindent\textbf{\texttt{states(g4a\_)}}\quad\textit{functoid}\index{states(g4a_)@\texttt{states(g4a\_)}}\\
reads: the states of \$1 \\
definition:
\begin{quote}\ttfamily\footnotesize\raggedright\noindent
\{g4f\_ in fun(\allowbreak{}vec(\allowbreak{}g4a\_), cc): is-positive-lf(\allowbreak{}g4a\_, g4f\_) and g4f\_(\allowbreak{}vone(\allowbreak{}g4a\_)) = 1\}
\end{quote}
characterized by \texttt{states-convex}  \textit{(proven)}:
\begin{quote}\ttfamily\footnotesize\raggedright\noindent
forall(\allowbreak{}[a], is-unital-involutive-banach-algebra(\allowbreak{}a) implies forall(\allowbreak{}[g in states(\allowbreak{}a), h in states(\allowbreak{}a), r in rr], 0 <= r implies r <= 1 implies vnb-lambda(\allowbreak{}g4tw\_, vec(\allowbreak{}a), r * g(\allowbreak{}g4tw\_) + (\allowbreak{}1 - r) * h(\allowbreak{}g4tw\_)) in states(\allowbreak{}a)))
\end{quote}
also: \texttt{states-unfold} \texttt{states-norm-one} \\
mentioned by 8 result(s)

\smallskip

\noindent\textbf{\texttt{strand-map(s)}}\quad\textit{accessor}\index{strand-map(s)@\texttt{strand-map(s)}}\\
reads: the strand-map slot of \$1 \\
a structure slot \\
characterized by \texttt{cx-sp@strand-map}  \textit{(definitional)}:
\begin{quote}\ttfamily\footnotesize\raggedright\noindent
strand-map(\allowbreak{}cx-sp) == vnb-lambda(\allowbreak{}cxi\_, interval(\allowbreak{}1, 8), nth(\allowbreak{}cxi\_, [[[sent, lst(\allowbreak{}[22, 10, 11, encr(\allowbreak{}lst(\allowbreak{}[20, 22, 10, 11]), 30)])], [received, lst(\allowbreak{}[22, encr(\allowbreak{}lst(\allowbreak{}[20, 33]), 30)])]], [[received, lst(\allowbreak{}[22, 10, 11, encr(\allowbreak{}lst(\allowbreak{}[20, 22, 10, 11]), 30)])], [sent, lst(\allowbreak{}[22, 10, 11, encr(\allowbreak{}lst(\allowbreak{}[20, 22, 10, 11]), 30), encr(\allowbreak{}lst(\allowbreak{}[21, 22, 10, 11]), 31)])], [received, lst(\allowbreak{}[22, encr(\allowbreak{}lst(\allowbreak{}[20, 32]), 30), encr(\allowbreak{}lst(\allowbreak{}[21, 32]), 31)])], [sent, lst(\allowbreak{}[22, encr(\allowbreak{}lst(\allowbreak{}[20, 32]), 30)])]], [[received, lst(\allowbreak{}[22, 10, 11, encr(\allowbreak{}lst(\allowbreak{}[20, 22, 10, 11]), 30), encr(\allowbreak{}lst(\allowbreak{}[21, 22, 10, 11]), 31)])], [sent, lst(\allowbreak{}[22, encr(\allowbreak{}lst(\allowbreak{}[20, 32]), 30), encr(\allowbreak{}lst(\allowbreak{}[21, 32]), 31)])]], [[received, lst(\allowbreak{}[22, 10, 11, encr(\allowbreak{}lst(\allowbreak{}[20, 22, 10, 11]), 30), encr(\allowbreak{}lst(\allowbreak{}[21, 22, 10, 11]), 31)])], [sent, lst(\allowbreak{}[22, encr(\allowbreak{}lst(\allowbreak{}[20, 33]), 30), encr(\allowbreak{}lst(\allowbreak{}[21, 33]), 31)])]], [[received, lst(\allowbreak{}[22, 10, 11, encr(\allowbreak{}lst(\allowbreak{}[20, 22, 10, 11]), 30), encr(\allowbreak{}lst(\allowbreak{}[21, 22, 10, 11]), 31)])], [sent, lst(\allowbreak{}[22, 10, 11, encr(\allowbreak{}lst(\allowbreak{}[20, 22, 10, 11]), 30), encr(\allowbreak{}lst(\allowbreak{}[21, 22, 10, 11]), 31)])], [sent, lst(\allowbreak{}[22, 10, 11, encr(\allowbreak{}lst(\allowbreak{}[20, 22, 10, 11]), 30), encr(\allowbreak{}lst(\allowbreak{}[21, 22, 10, 11]), 31)])]], [[received, lst(\allowbreak{}[22, encr(\allowbreak{}lst(\allowbreak{}[20, 32]), 30)])]], [[received, lst(\allowbreak{}[22, encr(\allowbreak{}lst(\allowbreak{}[20, 33]), 30), encr(\allowbreak{}lst(\allowbreak{}[21, 33]), 31)])], [sent, 22], [sent, encr(\allowbreak{}lst(\allowbreak{}[20, 33]), 30)], [sent, encr(\allowbreak{}lst(\allowbreak{}[21, 33]), 31)]], [[received, 22], [received, encr(\allowbreak{}lst(\allowbreak{}[20, 33]), 30)], [sent, lst(\allowbreak{}[22, encr(\allowbreak{}lst(\allowbreak{}[20, 33]), 30)])]]]))
\end{quote}
mentioned by 99 result(s)

\smallskip

\noindent\textbf{\texttt{strand-space}}\quad\textit{structure}\index{strand-space@\texttt{strand-space}}\\
predicate \texttt{is-strand-space}; slots: tokens ids keys strand-map kinv

\smallskip

\noindent\textbf{\texttt{strands(sst\_, ssk\_)}}\quad\textit{functoid}\index{strands(sst_, ssk_)@\texttt{strands(sst\_, ssk\_)}}\\
reads: the strands over the tokens \$1 and the keys \$2 \\
definition:
\begin{quote}\ttfamily\footnotesize\raggedright\noindent
tuples(\allowbreak{}directed(\allowbreak{}sst\_, ssk\_))
\end{quote}
characterized by \texttt{strands-unfold}  \textit{(proven)}:
\begin{quote}\ttfamily\footnotesize\raggedright\noindent
forall(\allowbreak{}[t, k], strands(\allowbreak{}t, k) == tuples(\allowbreak{}directed(\allowbreak{}t, k)))
\end{quote}
mentioned by 12 result(s)

\smallskip

\noindent\textbf{\texttt{strict-atom(md7af\_, md7at\_)}}\quad\textit{predicate}\index{strict-atom(md7af_, md7at_)@\texttt{strict-atom(md7af\_, md7at\_)}}\\
reads: the atom \$1 certifies the term \$2 \\
defined by \texttt{strict-atom} \\
characterized by \texttt{strict-atom}  \textit{(definitional)}:
\begin{quote}\ttfamily\footnotesize\raggedright\noindent
forall(\allowbreak{}[md7af\_, md7at\_], strict-atom(\allowbreak{}md7af\_, md7at\_) iff forsome(\allowbreak{}[md7ax\_, md7ay\_], (\allowbreak{}md7af\_ = enc-in(\allowbreak{}md7ax\_, md7ay\_) or md7af\_ = enc-eq(\allowbreak{}md7ax\_, md7ay\_)) and (\allowbreak{}md7at\_ in strict-reach(\allowbreak{}md7ax\_) or md7at\_ in strict-reach(\allowbreak{}md7ay\_))) or forsome(\allowbreak{}[md7pl\_, md7pk\_, md7pi\_], md7af\_ = eapp(\allowbreak{}10, md7pl\_) and nth(\allowbreak{}1, md7pl\_) = econst(\allowbreak{}md7pk\_) and pred-code(\allowbreak{}md7pk\_) and md7pi\_ in nn and 2 <= md7pi\_ and md7pi\_ <= length(\allowbreak{}md7pl\_) and md7at\_ in strict-reach(\allowbreak{}nth(\allowbreak{}md7pi\_, md7pl\_))))
\end{quote}
also: \texttt{strict-atom-le} \texttt{strict-atom-lt} \texttt{adm-strict-atom} \texttt{strict-atom-pred2} \\
mentioned by 14 result(s)

\smallskip

\noindent\textbf{\texttt{strict-closed(md7x\_)}}\quad\textit{predicate}\index{strict-closed(md7x_)@\texttt{strict-closed(md7x\_)}}\\
reads: \$1 is closed under strict steps \\
defined by \texttt{strict-closed} \\
characterized by \texttt{strict-closed}  \textit{(definitional)}:
\begin{quote}\ttfamily\footnotesize\raggedright\noindent
forall(\allowbreak{}[md7x\_], strict-closed(\allowbreak{}md7x\_) iff forall(\allowbreak{}[md7p\_, md7q\_], md7p\_ in md7x\_ implies md7q\_ in set implies strict-step(\allowbreak{}md7p\_, md7q\_) implies md7q\_ in md7x\_))
\end{quote}
mentioned by 5 result(s)

\smallskip

\noindent\textbf{\texttt{strict-epigraph(m, p)}}\quad\textit{functoid}\index{strict-epigraph(m, p)@\texttt{strict-epigraph(m, p)}}\\
reads: the strict epigraph of \$2 \\
definition:
\begin{quote}\ttfamily\footnotesize\raggedright\noindent
\{sepw\_ in cartesian(\allowbreak{}vec(\allowbreak{}m), rr): p(\allowbreak{}nth(\allowbreak{}1, sepw\_)) < nth(\allowbreak{}2, sepw\_)\}
\end{quote}
characterized by \texttt{strict-epigraph-unfold}  \textit{(proven)}:
\begin{quote}\ttfamily\footnotesize\raggedright\noindent
forall(\allowbreak{}[m, p], strict-epigraph(\allowbreak{}m, p) == \{sepw\_ in cartesian(\allowbreak{}vec(\allowbreak{}m), rr): p(\allowbreak{}nth(\allowbreak{}1, sepw\_)) < nth(\allowbreak{}2, sepw\_)\})
\end{quote}
also: \texttt{strict-epigraph-open-convex} \\
mentioned by 3 result(s)

\smallskip

\noindent\textbf{\texttt{strict-head(mchd\_)}}\quad\textit{predicate}\index{strict-head(mchd_)@\texttt{strict-head(mchd\_)}}\\
reads: an application of the head \$1 is strict in its arguments \\
defined by \texttt{strict-head} \\
characterized by \texttt{strict-head}  \textit{(definitional)}:
\begin{quote}\ttfamily\footnotesize\raggedright\noindent
forall(\allowbreak{}[mchd\_], strict-head(\allowbreak{}mchd\_) iff 7 <= mchd\_ and mchd\_ <= 8 or 11 <= mchd\_ and mchd\_ <= 19 or 30 <= mchd\_ and mchd\_ <= 31)
\end{quote}
mentioned by 9 result(s)

\smallskip

\noindent\textbf{\texttt{strict-reach(md7e\_)}}\quad\textit{functoid}\index{strict-reach(md7e_)@\texttt{strict-reach(md7e\_)}}\\
reads: the expressions reached from \$1 by strict steps \\
definition:
\begin{quote}\ttfamily\footnotesize\raggedright\noindent
\{md7z\_ in expr-set: forall(\allowbreak{}[md7w\_ in power(\allowbreak{}expr-set)], md7e\_ in md7w\_ implies strict-closed(\allowbreak{}md7w\_) implies md7z\_ in md7w\_)\}
\end{quote}
characterized by \texttt{strict-reach-least}  \textit{(proven)}:
\begin{quote}\ttfamily\footnotesize\raggedright\noindent
forall(\allowbreak{}[e, md7xs\_ in power(\allowbreak{}expr-set)], e in md7xs\_ implies strict-closed(\allowbreak{}md7xs\_) implies forall(\allowbreak{}[x in strict-reach(\allowbreak{}e)], x in md7xs\_))
\end{quote}
also: \texttt{strict-reach-unfold} \\
mentioned by 9 result(s)

\smallskip

\noindent\textbf{\texttt{strict-step(md7s\_, md7t\_)}}\quad\textit{predicate}\index{strict-step(md7s_, md7t_)@\texttt{strict-step(md7s\_, md7t\_)}}\\
reads: \$2 is one strict step below \$1 \\
defined by \texttt{strict-step} \\
characterized by \texttt{strict-step}  \textit{(definitional)}:
\begin{quote}\ttfamily\footnotesize\raggedright\noindent
forall(\allowbreak{}[md7s\_, md7t\_], strict-step(\allowbreak{}md7s\_, md7t\_) iff forsome(\allowbreak{}[md7h\_, md7l\_, md7i\_], md7s\_ = eapp(\allowbreak{}md7h\_, md7l\_) and strict-head(\allowbreak{}md7h\_) and md7i\_ in nn and 1 <= md7i\_ and md7i\_ <= length(\allowbreak{}md7l\_) and md7t\_ = nth(\allowbreak{}md7i\_, md7l\_)) or forsome(\allowbreak{}[md7k\_, md7a\_, md7b\_], (\allowbreak{}md7k\_ = 22 or md7k\_ = 24 or md7k\_ = 26) and md7s\_ = ebind(\allowbreak{}md7k\_, 1, [md7a\_], md7b\_) and md7t\_ = md7a\_))
\end{quote}
mentioned by 4 result(s)

\smallskip

\noindent\textbf{\texttt{strictly-mono-nn(phi)}}\quad\textit{predicate}\index{strictly-mono-nn(phi)@\texttt{strictly-mono-nn(phi)}}\\
reads: \$1 is a strictly increasing sequence of naturals \\
defined by \texttt{strictly-mono-nn} \\
characterized by \texttt{strictly-mono-nn}  \textit{(definitional)}:
\begin{quote}\ttfamily\footnotesize\raggedright\noindent
forall(\allowbreak{}[phi], strictly-mono-nn(\allowbreak{}phi) iff phi in fun(\allowbreak{}nn, nn) and forall(\allowbreak{}[m in nn, n\_ in nn], m < n\_ implies phi(\allowbreak{}m) < phi(\allowbreak{}n\_)))
\end{quote}
mentioned by 58 result(s)

\smallskip

\noindent\textbf{\texttt{strong-dual(otv\_)}}\quad\textit{functoid}\index{strong-dual(otv_)@\texttt{strong-dual(otv\_)}}\\
reads: the strongly continuous linear functionals on the bounded operators on \$1 \\
definition:
\begin{quote}\ttfamily\footnotesize\raggedright\noindent
op-dual(\allowbreak{}otv\_, sot(\allowbreak{}otv\_))
\end{quote}
characterized by \texttt{strong-dual-bound}  \textit{(proven)}:
\begin{quote}\ttfamily\footnotesize\raggedright\noindent
forall(\allowbreak{}[v, f\_], is-complex-inner-product-space(\allowbreak{}v) implies f\_ in strong-dual(\allowbreak{}v) implies forsome(\allowbreak{}[sbf\_ in set], card(\allowbreak{}sbf\_) in nn and sbf\_ subset vec(\allowbreak{}v) and forsome(\allowbreak{}[sbc\_ in rr], pos-rr(\allowbreak{}sbc\_) and forall(\allowbreak{}[sbx\_ in bounded-ops(\allowbreak{}v, v), sbt\_ in rr], 0 <= sbt\_ implies forall(\allowbreak{}[sbi\_ in sbf\_], ip-norm(\allowbreak{}v, sbx\_(\allowbreak{}sbi\_)) <= sbt\_) implies magnitude(\allowbreak{}f\_(\allowbreak{}sbx\_)) <= sbc\_ * sbt\_))))
\end{quote}
also: \texttt{strong-dual-unfold} \texttt{strong-dual-unfold-sot} \texttt{weak-dual-subset-strong-dual} \texttt{vec-functional-in-strong-dual} \\
mentioned by 7 result(s)

\smallskip

\noindent\textbf{\texttt{sub-cips(v, kk)}}\quad\textit{functoid}\index{sub-cips(v, kk)@\texttt{sub-cips(v, kk)}}\\
reads: the subspace \$2 of \$1 as an inner product space \\
definition:
\begin{quote}\ttfamily\footnotesize\raggedright\noindent
[normed-field-as-commutative-ring(\allowbreak{}cc-normed-field), kk, vnb-lambda(\allowbreak{}[tsbx\_, tsby\_], cartesian(\allowbreak{}kk, kk), (\allowbreak{}vadd(\allowbreak{}v))(\allowbreak{}tsbx\_, tsby\_)), vzero(\allowbreak{}v), vnb-lambda(\allowbreak{}tsbx\_, kk, (\allowbreak{}vneg(\allowbreak{}v))(\allowbreak{}tsbx\_)), vnb-lambda(\allowbreak{}[tsbc\_, tsbx\_], cartesian(\allowbreak{}cc, kk), (\allowbreak{}act(\allowbreak{}v))(\allowbreak{}tsbc\_, tsbx\_)), vnb-lambda(\allowbreak{}[tsbx\_, tsby\_], cartesian(\allowbreak{}kk, kk), (\allowbreak{}ip(\allowbreak{}v))(\allowbreak{}tsbx\_, tsby\_))]
\end{quote}
characterized by \texttt{sub-cips-ip}  \textit{(proven)}:
\begin{quote}\ttfamily\footnotesize\raggedright\noindent
forall(\allowbreak{}[v, kk], is-complex-inner-product-space(\allowbreak{}v) implies is-submodule(\allowbreak{}v, kk) implies ip(\allowbreak{}sub-cips(\allowbreak{}v, kk)) == vnb-lambda(\allowbreak{}[tsbx\_, tsby\_], cartesian(\allowbreak{}kk, kk), (\allowbreak{}ip(\allowbreak{}v))(\allowbreak{}tsbx\_, tsby\_)))
\end{quote}
also: \texttt{sub-cips-act} \texttt{sub-cips-vec} \texttt{sub-cips-dist} \texttt{sub-cips-norm} \texttt{sub-cips-scal} \texttt{sub-cips-vadd} \texttt{sub-cips-vneg} \texttt{sub-cips-vzero} ... (19 in all) \\
mentioned by 50 result(s)

\smallskip

\noindent\textbf{\texttt{sub-nf(snk\_, sns\_)}}\quad\textit{functoid}\index{sub-nf(snk_, sns_)@\texttt{sub-nf(snk\_, sns\_)}}\\
reads: the normed field \$2 with the structure of \$1 restricted to it \\
definition:
\begin{quote}\ttfamily\footnotesize\raggedright\noindent
[sns\_, vnb-lambda(\allowbreak{}[snu\_, snv\_], cartesian(\allowbreak{}sns\_, sns\_), (\allowbreak{}add(\allowbreak{}snk\_))(\allowbreak{}snu\_, snv\_)), vnb-lambda(\allowbreak{}[snu\_, snv\_], cartesian(\allowbreak{}sns\_, sns\_), (\allowbreak{}mul(\allowbreak{}snk\_))(\allowbreak{}snu\_, snv\_)), vnb-lambda(\allowbreak{}snu\_, sns\_, (\allowbreak{}neg(\allowbreak{}snk\_))(\allowbreak{}snu\_)), zero(\allowbreak{}snk\_), one(\allowbreak{}snk\_), vnb-lambda(\allowbreak{}snu\_, sns\_, (\allowbreak{}fnrm(\allowbreak{}snk\_))(\allowbreak{}snu\_))]
\end{quote}
characterized by \texttt{sub-nf-add}  \textit{(proven)}:
\begin{quote}\ttfamily\footnotesize\raggedright\noindent
forall(\allowbreak{}[k], is-normed-field(\allowbreak{}k) implies forall(\allowbreak{}[s], is-nf-subfield(\allowbreak{}k, s) implies add(\allowbreak{}sub-nf(\allowbreak{}k, s)) == vnb-lambda(\allowbreak{}[snu\_, snv\_], cartesian(\allowbreak{}s, s), (\allowbreak{}add(\allowbreak{}k))(\allowbreak{}snu\_, snv\_))))
\end{quote}
also: \texttt{sub-nf-mul} \texttt{sub-nf-neg} \texttt{sub-nf-one} \texttt{sub-nf-carr} \texttt{sub-nf-fnrm} \texttt{sub-nf-zero} \texttt{sub-nf-local} \texttt{sub-nf-metric} ... (11 in all) \\
mentioned by 26 result(s)

\smallskip

\noindent\textbf{\texttt{subadd-reg(srm\_, srg\_, srx\_)}}\quad\textit{functoid}\index{subadd-reg(srm_, srg_, srx_)@\texttt{subadd-reg(srm\_, srg\_, srx\_)}}\\
reads: the infimum of the sums of \$2 over the factorisations of \$3 in \$1 \\
definition:
\begin{quote}\ttfamily\footnotesize\raggedright\noindent
einf(\allowbreak{}fact-sums(\allowbreak{}srm\_, srg\_, srx\_))
\end{quote}
characterized by \texttt{subadd-reg-bounds}  \textit{(proven)}:
\begin{quote}\ttfamily\footnotesize\raggedright\noindent
forall(\allowbreak{}[m], is-monoid(\allowbreak{}m) implies forall(\allowbreak{}[g in fun(\allowbreak{}carr(\allowbreak{}m), rr)], forall(\allowbreak{}[x in carr(\allowbreak{}m)], 0 <= g(\allowbreak{}x)) implies g(\allowbreak{}iden(\allowbreak{}m)) = 0 implies forall(\allowbreak{}[x in carr(\allowbreak{}m), y in carr(\allowbreak{}m), z in carr(\allowbreak{}m)], g(\allowbreak{}(\allowbreak{}opr(\allowbreak{}m))(\allowbreak{}(\allowbreak{}opr(\allowbreak{}m))(\allowbreak{}x, y), z)) <= 2 * max(\allowbreak{}g(\allowbreak{}x), max(\allowbreak{}g(\allowbreak{}y), g(\allowbreak{}z)))) implies forall(\allowbreak{}[x in carr(\allowbreak{}m)], subadd-reg(\allowbreak{}m, g, x) in rr and num(\allowbreak{}1/2) * g(\allowbreak{}x) <= subadd-reg(\allowbreak{}m, g, x) and subadd-reg(\allowbreak{}m, g, x) <= g(\allowbreak{}x))))
\end{quote}
also: \texttt{subadd-reg-unfold} \texttt{subadd-reg-subadditive} \\
mentioned by 6 result(s)

\smallskip

\noindent\textbf{\texttt{submat(s, p, q)}}\quad\textit{functoid}\index{submat(s, p, q)@\texttt{submat(s, p, q)}}\\
definition:
\begin{quote}\ttfamily\footnotesize\raggedright\noindent
matof(\allowbreak{}p, q, vnb-lambda(\allowbreak{}[i, j], cartesian(\allowbreak{}interval(\allowbreak{}1, p), interval(\allowbreak{}1, q)), entry(\allowbreak{}s, succ(\allowbreak{}i), succ(\allowbreak{}j))))
\end{quote}
characterized by \texttt{submat-type}  \textit{(proven)}:
\begin{quote}\ttfamily\footnotesize\raggedright\noindent
forall(\allowbreak{}[a, p, q, s], p in nn implies q in nn implies s in mat(\allowbreak{}succ(\allowbreak{}p), succ(\allowbreak{}q), carr(\allowbreak{}a)) implies submat(\allowbreak{}s, p, q) in mat(\allowbreak{}p, q, carr(\allowbreak{}a)))
\end{quote}
also: \texttt{entry-of-submat} \\
mentioned by 24 result(s)

\smallskip

\noindent\textbf{\texttt{submatrix(a, s, p, q)}}\quad\textit{functoid}\index{submatrix(a, s, p, q)@\texttt{submatrix(a, s, p, q)}}\\
reads: the \$2-by-\$2 submatrix of \$1 on the rows \$3 and the columns \$4 \\
definition:
\begin{quote}\ttfamily\footnotesize\raggedright\noindent
matof(\allowbreak{}s, s, vnb-lambda(\allowbreak{}[smi\_, smj\_], cartesian(\allowbreak{}interval(\allowbreak{}1, s), interval(\allowbreak{}1, s)), entry(\allowbreak{}a, p(\allowbreak{}smi\_), q(\allowbreak{}smj\_))))
\end{quote}
characterized by \texttt{submatrix-ext}  \textit{(proven)}:
\begin{quote}\ttfamily\footnotesize\raggedright\noindent
forall(\allowbreak{}[r, m, n, s, a, p, q, b, pb, qb], is-ring(\allowbreak{}r) implies m in nn implies n in nn implies s in nn implies a in mat(\allowbreak{}m, n, carr(\allowbreak{}r)) implies p in fun(\allowbreak{}interval(\allowbreak{}1, s), interval(\allowbreak{}1, m)) implies q in fun(\allowbreak{}interval(\allowbreak{}1, s), interval(\allowbreak{}1, n)) implies b in mat(\allowbreak{}m, n, carr(\allowbreak{}r)) implies pb in fun(\allowbreak{}interval(\allowbreak{}1, s), interval(\allowbreak{}1, m)) implies qb in fun(\allowbreak{}interval(\allowbreak{}1, s), interval(\allowbreak{}1, n)) implies forall(\allowbreak{}[sxi\_ in interval(\allowbreak{}1, s), sxj\_ in interval(\allowbreak{}1, s)], entry(\allowbreak{}a, p(\allowbreak{}sxi\_), q(\allowbreak{}sxj\_)) = entry(\allowbreak{}b, pb(\allowbreak{}sxi\_), qb(\allowbreak{}sxj\_))) implies submatrix(\allowbreak{}a, s, p, q) = submatrix(\allowbreak{}b, s, pb, qb))
\end{quote}
also: \texttt{submatrix-type} \texttt{submatrix-entry} \texttt{lzd-minor-submatrix} \texttt{submatrix-mul-diagonal} \\
mentioned by 16 result(s)

\smallskip

\noindent\textbf{\texttt{submodule-norm(md, sm)}}\quad\textit{functoid}\index{submodule-norm(md, sm)@\texttt{submodule-norm(md, sm)}}\\
reads: the norm N\_E(F) of \$2 relative to \$1 \\
definition:
\begin{quote}\ttfamily\footnotesize\raggedright\noindent
if(\allowbreak{}forsome(\allowbreak{}[nmn\_, nmx\_, nmy\_, nma\_], is-norm-data(\allowbreak{}md, sm, nmn\_, nmx\_, nmy\_, nma\_)), iota(\allowbreak{}nmd\_, forsome(\allowbreak{}[nmn\_, nmx\_, nmy\_, nma\_], is-norm-data(\allowbreak{}md, sm, nmn\_, nmx\_, nmy\_, nma\_) and nmd\_ = abs(\allowbreak{}det(\allowbreak{}zz-ring, nmn\_, nma\_)))), 0)
\end{quote}
characterized by \texttt{submodule-norm-value-cna-as-module}  \textit{(definitional)}:
\begin{quote}\ttfamily\footnotesize\raggedright\noindent
forall(\allowbreak{}[r], is-complex-normed-algebra(\allowbreak{}r) implies scal(\allowbreak{}r) = zz-ring implies forall(\allowbreak{}[sm, n, xs, ys, am], is-norm-data(\allowbreak{}cna-as-module(\allowbreak{}r), sm, n, xs, ys, am) implies submodule-norm(\allowbreak{}cna-as-module(\allowbreak{}r), sm) = abs(\allowbreak{}det(\allowbreak{}zz-ring, n, am))))
\end{quote}
also: \texttt{submodule-norm-value-sip-as-module} \texttt{submodule-norm-value-uba-as-module} \texttt{submodule-norm-value-cnvs-as-module} \texttt{submodule-norm-value-sip-as-module-csip-as-sip} \texttt{submodule-norm-value-uba-as-module-uiba-as-uba} \texttt{submodule-norm-zero} \texttt{submodule-norm-value} \\
mentioned by 70 result(s)

\smallskip

\noindent\textbf{\texttt{subseq(f, phi)}}\quad\textit{functoid}\index{subseq(f, phi)@\texttt{subseq(f, phi)}}\\
definition:
\begin{quote}\ttfamily\footnotesize\raggedright\noindent
vnb-lambda(\allowbreak{}k, nn, f(\allowbreak{}phi(\allowbreak{}k)))
\end{quote}
characterized by \texttt{is-subsequence}  \textit{(definitional)}:
\begin{quote}\ttfamily\footnotesize\raggedright\noindent
forall(\allowbreak{}[s, y, f], is-subsequence(\allowbreak{}s, y, f) iff f in fun(\allowbreak{}nn, pts(\allowbreak{}s)) and forsome(\allowbreak{}[phi], strictly-mono-nn(\allowbreak{}phi) and y = subseq(\allowbreak{}f, phi)))
\end{quote}
also: \texttt{subseq-apply} \texttt{subseq-is-fun} \texttt{diagonal-subseq} \texttt{equicont-subseq} \texttt{subseq-fun-typing} \texttt{cc-pointwise-subseq} \texttt{subseq-of-convergent} \texttt{cc-hol-bounded-subseq} ... (22 in all) \\
mentioned by 42 result(s)

\smallskip

\noindent\textbf{\texttt{subset}}\quad\textit{primitive}\index{subset@\texttt{subset}}\\
reads: \$1 is a subset of \$2 \\
a kernel relation \\
characterized by \texttt{is-convex-subset}  \textit{(definitional)}:
\begin{quote}\ttfamily\footnotesize\raggedright\noindent
forall(\allowbreak{}[v, cs], is-convex-subset(\allowbreak{}v, cs) iff cs subset vec(\allowbreak{}v) and forall(\allowbreak{}[sxa\_ in cs, sxb\_ in cs, sxt\_ in rr], 0 <= sxt\_ implies sxt\_ <= 1 implies (\allowbreak{}vadd(\allowbreak{}v))(\allowbreak{}(\allowbreak{}act(\allowbreak{}v))(\allowbreak{}1 - sxt\_, sxa\_), (\allowbreak{}act(\allowbreak{}v))(\allowbreak{}sxt\_, sxb\_)) in cs))
\end{quote}
also: \texttt{submodule-subset} \texttt{is-complete-subset} \texttt{is-extremal-subset} \texttt{is-absorbing-subset} \texttt{is-symmetric-subset} \texttt{ringoid-ideal-subset} \texttt{inf-subsets-membership} \texttt{lfe-line-subset-cnvs-as-nvs} ... (114 in all) \\
mentioned by 1600 result(s)

\smallskip

\noindent\textbf{\texttt{subspace-ms(s, a)}}\quad\textit{functoid}\index{subspace-ms(s, a)@\texttt{subspace-ms(s, a)}}\\
reads: the subspace of \$1 on \$2 \\
definition:
\begin{quote}\ttfamily\footnotesize\raggedright\noindent
[a, vnb-lambda(\allowbreak{}[sbu\_, sbv\_], cartesian(\allowbreak{}a, a), (\allowbreak{}dist(\allowbreak{}s))(\allowbreak{}sbu\_, sbv\_))]
\end{quote}
mentioned by 388 result(s)

\smallskip

\noindent\textbf{\texttt{subspace-top(t, a)}}\quad\textit{functoid}\index{subspace-top(t, a)@\texttt{subspace-top(t, a)}}\\
reads: the subspace topology of \$1 on \$2 \\
definition:
\begin{quote}\ttfamily\footnotesize\raggedright\noindent
[a, \{stw\_ in power(\allowbreak{}a): forsome(\allowbreak{}[stu\_ in opens(\allowbreak{}t)], stw\_ = intersection(\allowbreak{}stu\_, a))\}]
\end{quote}
characterized by \texttt{subspace-top-pts}  \textit{(proven)}:
\begin{quote}\ttfamily\footnotesize\raggedright\noindent
forall(\allowbreak{}[t, a in set], pts(\allowbreak{}subspace-top(\allowbreak{}t, a)) == a)
\end{quote}
also: \texttt{subspace-top-opens} \texttt{subspace-top-unfold} \texttt{subspace-top-is-top-space} \texttt{subspace-top-of-metric-top} \texttt{metric-compact-subspace-top-compact} \\
mentioned by 51 result(s)

\smallskip

\noindent\textbf{\texttt{subsqn(y, f)}}\quad\textit{predicate}\index{subsqn(y, f)@\texttt{subsqn(y, f)}}\\
reads: \$1 is a subsequence of \$2 \\
defined by \texttt{subsqn} \\
characterized by \texttt{subsqn}  \textit{(definitional)}:
\begin{quote}\ttfamily\footnotesize\raggedright\noindent
forall(\allowbreak{}[y, f], subsqn(\allowbreak{}y, f) iff forsome(\allowbreak{}[phi], strictly-mono-nn(\allowbreak{}phi) and y = subseq(\allowbreak{}f, phi)))
\end{quote}
also: \texttt{subsqn-preserves-sqn} \texttt{is-subsequence-iff-typed-subsqn} \\
mentioned by 5 result(s)

\smallskip

\noindent\textbf{\texttt{subst}}\quad\textit{tactic}\index{subst@\texttt{subst}}\\
\texttt{(subst '(= s t))  |  (subst '(== s t))} -- Rewrite s -> t throughout the goal, using an equation s = t -- or a quasi-equation s == t -- that is in context, in EITHER orientation (Leibniz substitution).  Reaches every position, operator (head) position included; a binder whose variable the equation mentions is not entered. \\
declared in \texttt{tactics-help.scm}

\smallskip

\noindent\textbf{\texttt{subst-fun(mxm\_, mxu\_)}}\quad\textit{functoid}\index{subst-fun(mxm_, mxu_)@\texttt{subst-fun(mxm\_, mxu\_)}}\\
reads: the map replacing the free variable \$1 by \$2 \\
definition:
\begin{quote}\ttfamily\footnotesize\raggedright\noindent
vnb-lambda(\allowbreak{}mxw\_, expr-set, expr-subst(\allowbreak{}mxw\_, mxm\_, mxu\_))
\end{quote}
characterized by \texttt{subst-fun-unfold}  \textit{(proven)}:
\begin{quote}\ttfamily\footnotesize\raggedright\noindent
forall(\allowbreak{}[m, u], subst-fun(\allowbreak{}m, u) == vnb-lambda(\allowbreak{}mxw\_, expr-set, expr-subst(\allowbreak{}mxw\_, m, u)))
\end{quote}
mentioned by 6 result(s)

\smallskip

\noindent\textbf{\texttt{subst-stage}}\quad\textit{functoid}\index{subst-stage@\texttt{subst-stage}}\\
reads: the map replacing the free variable \$1 by \$2 on the stage \$3 \\
defined by \texttt{subst-stage-zero} \texttt{subst-stage-succ} \\
characterized by \texttt{subst-stage-succ}  \textit{(definitional)}:
\begin{quote}\ttfamily\footnotesize\raggedright\noindent
forall(\allowbreak{}[mxm\_, mxu\_, mxn\_ in nn], subst-stage(\allowbreak{}mxm\_, mxu\_, succ(\allowbreak{}mxn\_)) == vnb-lambda(\allowbreak{}mxe\_, expr-stage(\allowbreak{}succ(\allowbreak{}mxn\_)), if(\allowbreak{}mxe\_ in expr-stage(\allowbreak{}mxn\_), (\allowbreak{}subst-stage(\allowbreak{}mxm\_, mxu\_, mxn\_))(\allowbreak{}mxe\_), if(\allowbreak{}nth(\allowbreak{}1, mxe\_) = 3, [3, nth(\allowbreak{}2, mxe\_), mapl(\allowbreak{}subst-stage(\allowbreak{}mxm\_, mxu\_, mxn\_), nth(\allowbreak{}3, mxe\_))], [4, nth(\allowbreak{}2, mxe\_), nth(\allowbreak{}3, mxe\_), mapl(\allowbreak{}subst-stage(\allowbreak{}mxm\_, mxu\_, mxn\_), nth(\allowbreak{}4, mxe\_)), (\allowbreak{}subst-stage(\allowbreak{}mxm\_, mxu\_, mxn\_))(\allowbreak{}nth(\allowbreak{}5, mxe\_))]))))
\end{quote}
also: \texttt{subst-stage-zero} \texttt{subst-stage-coh} \texttt{subst-stage-fun} \texttt{subst-stage-step} \\
mentioned by 11 result(s)

\smallskip

\noindent\textbf{\texttt{subterm-rel(smt\_, smk\_)}}\quad\textit{functoid}\index{subterm-rel(smt_, smk_)@\texttt{subterm-rel(smt\_, smk\_)}}\\
reads: the subterm relation over the tokens \$1 and the keys \$2 \\
definition:
\begin{quote}\ttfamily\footnotesize\raggedright\noindent
\{smp\_ in cartesian(\allowbreak{}msg-set(\allowbreak{}smt\_, smk\_), msg-set(\allowbreak{}smt\_, smk\_)): forall(\allowbreak{}[smq\_ in power(\allowbreak{}cartesian(\allowbreak{}msg-set(\allowbreak{}smt\_, smk\_), msg-set(\allowbreak{}smt\_, smk\_)))], is-subterm-closed(\allowbreak{}smt\_, smk\_, smq\_) implies smp\_ in smq\_)\}
\end{quote}
characterized by \texttt{subterm-rel-unfold}  \textit{(proven)}:
\begin{quote}\ttfamily\footnotesize\raggedright\noindent
forall(\allowbreak{}[t, k], subterm-rel(\allowbreak{}t, k) == \{p\_ in cartesian(\allowbreak{}msg-set(\allowbreak{}t, k), msg-set(\allowbreak{}t, k)): forall(\allowbreak{}[q\_ in power(\allowbreak{}cartesian(\allowbreak{}msg-set(\allowbreak{}t, k), msg-set(\allowbreak{}t, k)))], is-subterm-closed(\allowbreak{}t, k, q\_) implies p\_ in q\_)\})
\end{quote}
mentioned by 4 result(s)

\smallskip

\noindent\textbf{\texttt{succ}}\quad\textit{operator}\index{succ@\texttt{succ}}\\
a term-forming head \\
characterized by \texttt{sum-succ}  \textit{(definitional)}:
\begin{quote}\ttfamily\footnotesize\raggedright\noindent
forall(\allowbreak{}[r, f, n in nn], sum(\allowbreak{}r, f, succ(\allowbreak{}n)) == (\allowbreak{}add(\allowbreak{}r))(\allowbreak{}sum(\allowbreak{}r, f, n), f(\allowbreak{}n)))
\end{quote}
also: \texttt{cchs-succ} \texttt{gpow-succ} \texttt{mpow-succ} \texttt{poch-succ} \texttt{mprod-succ} \texttt{openl-succ} \texttt{pprod-succ} \texttt{el-pow-succ} ... (142 in all) \\
mentioned by 1282 result(s)

\smallskip

\noindent\textbf{\texttt{succ\_ord}}\quad\textit{operator}\index{succ_ord@\texttt{succ\_ord}}\\
a term-forming head \\
mentioned by 29 result(s)

\smallskip

\noindent\textbf{\texttt{sum}}\quad\textit{defined-fn}\index{sum@\texttt{sum}}\\
defined by \texttt{sum-zero} \texttt{sum-succ} \\
characterized by \texttt{sum-succ}  \textit{(definitional)}:
\begin{quote}\ttfamily\footnotesize\raggedright\noindent
forall(\allowbreak{}[r, f, n in nn], sum(\allowbreak{}r, f, succ(\allowbreak{}n)) == (\allowbreak{}add(\allowbreak{}r))(\allowbreak{}sum(\allowbreak{}r, f, n), f(\allowbreak{}n)))
\end{quote}
also: \texttt{sum-zero} \texttt{sum-type-uba-as-ring} \texttt{sum-singleton-uba-as-ring} \texttt{sum-left-scalar-uba-as-ring} \texttt{sum-type-uba-as-ring-uiba-as-uba} \texttt{sum-singleton-uba-as-ring-uiba-as-uba} \texttt{sum-left-scalar-uba-as-ring-uiba-as-uba} \texttt{sum-type} ... (18 in all) \\
mentioned by 50 result(s)

\smallskip

\noindent\textbf{\texttt{sum-ag}}\quad\textit{defined-fn}\index{sum-ag@\texttt{sum-ag}}\\
defined by \texttt{sum-ag-zero} \texttt{sum-ag-succ} \\
characterized by \texttt{sum-ag-succ}  \textit{(definitional)}:
\begin{quote}\ttfamily\footnotesize\raggedright\noindent
forall(\allowbreak{}[ag, f, n in nn], sum-ag(\allowbreak{}ag, f, succ(\allowbreak{}n)) == (\allowbreak{}opr(\allowbreak{}ag))(\allowbreak{}sum-ag(\allowbreak{}ag, f, n), f(\allowbreak{}n)))
\end{quote}
also: \texttt{sum-ag-zero} \texttt{sum-ag-type-ring-additive-ag} \texttt{sum-ag-as-reduce-ring-additive-ag} \texttt{sum-ag-singleton-ring-additive-ag} \texttt{sum-ag-type-ring-additive-ag-uba-as-ring} \texttt{sum-ag-as-reduce-ring-additive-ag-uba-as-ring} \texttt{sum-ag-singleton-ring-additive-ag-uba-as-ring} \texttt{sum-ag-type-ring-additive-ag-uba-as-ring-uiba-as-uba} ... (41 in all) \\
mentioned by 134 result(s)

\smallskip

\noindent\textbf{\texttt{sum-ms(smt\_)}}\quad\textit{functoid}\index{sum-ms(smt_)@\texttt{sum-ms(smt\_)}}\\
reads: the disjoint sum of the metric spaces \$1 \\
definition:
\begin{quote}\ttfamily\footnotesize\raggedright\noindent
[sum-pts(\allowbreak{}smt\_), vnb-lambda(\allowbreak{}[smu\_, smv\_], cartesian(\allowbreak{}sum-pts(\allowbreak{}smt\_), sum-pts(\allowbreak{}smt\_)), if(\allowbreak{}nth(\allowbreak{}1, smu\_) = nth(\allowbreak{}1, smv\_), (\allowbreak{}dist(\allowbreak{}bdd-metric(\allowbreak{}smt\_(\allowbreak{}nth(\allowbreak{}1, smu\_)))))(\allowbreak{}nth(\allowbreak{}2, smu\_), nth(\allowbreak{}2, smv\_)), 1))]
\end{quote}
characterized by \texttt{sum-ms-pts}  \textit{(proven)}:
\begin{quote}\ttfamily\footnotesize\raggedright\noindent
forall(\allowbreak{}[ts], pts(\allowbreak{}sum-ms(\allowbreak{}ts)) == sum-pts(\allowbreak{}ts))
\end{quote}
also: \texttt{sum-ms-polish} \texttt{sum-ms-complete} \texttt{sum-ms-dist-slot} \texttt{sum-ms-separable} \texttt{sum-ms-is-metric-space} \\
mentioned by 14 result(s)

\smallskip

\noindent\textbf{\texttt{sum-pts(smt\_)}}\quad\textit{functoid}\index{sum-pts(smt_)@\texttt{sum-pts(smt\_)}}\\
reads: the points of the disjoint sum of \$1 \\
definition:
\begin{quote}\ttfamily\footnotesize\raggedright\noindent
\{smp\_ in cartesian(\allowbreak{}nn, big-union(\allowbreak{}smk\_, nn, pts(\allowbreak{}smt\_(\allowbreak{}smk\_)))): nth(\allowbreak{}2, smp\_) in pts(\allowbreak{}smt\_(\allowbreak{}nth(\allowbreak{}1, smp\_)))\}
\end{quote}
characterized by \texttt{sum-pts-intro}  \textit{(proven)}:
\begin{quote}\ttfamily\footnotesize\raggedright\noindent
forall(\allowbreak{}[ts, k in nn, x in pts(\allowbreak{}ts(\allowbreak{}k))], [k, x] in sum-pts(\allowbreak{}ts))
\end{quote}
also: \texttt{sum-pts-parts} \texttt{sum-pts-is-set} \texttt{sum-pts-unfold} \\
mentioned by 16 result(s)

\smallskip

\noindent\textbf{\texttt{sum-set(r, s, f)}}\quad\textit{functoid}\index{sum-set(r, s, f)@\texttt{sum-set(r, s, f)}}\\
definition:
\begin{quote}\ttfamily\footnotesize\raggedright\noindent
finsum(\allowbreak{}ring-additive-ag(\allowbreak{}r), f, s)
\end{quote}
characterized by \texttt{sum-set-empty-uba-as-ring}  \textit{(definitional)}:
\begin{quote}\ttfamily\footnotesize\raggedright\noindent
forall(\allowbreak{}[r], is-unital-banach-algebra(\allowbreak{}r) implies forall(\allowbreak{}[f], sum-set(\allowbreak{}uba-as-ring(\allowbreak{}r), empty-set, f) == vzero(\allowbreak{}r)))
\end{quote}
also: \texttt{sum-set-left-scalar-uba-as-ring} \texttt{sum-set-right-scalar-uba-as-ring} \texttt{sum-set-empty-uba-as-ring-uiba-as-uba} \texttt{sum-set-singleton-defined-uba-as-ring} \texttt{sum-set-disjoint-union-defined-uba-as-ring} \texttt{sum-set-left-scalar-uba-as-ring-uiba-as-uba} \texttt{sum-set-right-scalar-uba-as-ring-uiba-as-uba} \texttt{sum-set-singleton-defined-uba-as-ring-uiba-as-uba} ... (15 in all) \\
mentioned by 71 result(s)

\smallskip

\noindent\textbf{\texttt{sum-sigma(sgx\_, sgc\_)}}\quad\textit{functoid}\index{sum-sigma(sgx_, sgc_)@\texttt{sum-sigma(sgx\_, sgc\_)}}\\
reads: the sum of the sigma-algebras \$2 on the sets \$1 \\
definition:
\begin{quote}\ttfamily\footnotesize\raggedright\noindent
\{sga\_ in power(\allowbreak{}disjoint-sum-set(\allowbreak{}sgx\_)): forall(\allowbreak{}[sgk\_ in nn], \{sgy\_ in sgx\_(\allowbreak{}sgk\_): [sgk\_, sgy\_] in sga\_\} in sgc\_(\allowbreak{}sgk\_))\}
\end{quote}
characterized by \texttt{sum-sigma-unfold}  \textit{(proven)}:
\begin{quote}\ttfamily\footnotesize\raggedright\noindent
forall(\allowbreak{}[xss, cas], sum-sigma(\allowbreak{}xss, cas) == \{sga\_ in power(\allowbreak{}disjoint-sum-set(\allowbreak{}xss)): forall(\allowbreak{}[sgk\_ in nn], \{sgy\_ in xss(\allowbreak{}sgk\_): [sgk\_, sgy\_] in sga\_\} in cas(\allowbreak{}sgk\_))\})
\end{quote}
also: \texttt{sum-sigma-algebra} \\
mentioned by 4 result(s)

\smallskip

\noindent\textbf{\texttt{summable-weight(w)}}\quad\textit{predicate}\index{summable-weight(w)@\texttt{summable-weight(w)}}\\
reads: \$1 is a summable sequence of positive weights \\
defined by \texttt{summable-weight} \\
characterized by \texttt{summable-weight}  \textit{(definitional)}:
\begin{quote}\ttfamily\footnotesize\raggedright\noindent
forall(\allowbreak{}[w], summable-weight(\allowbreak{}w) iff w in fun(\allowbreak{}nn, rr) and forall(\allowbreak{}[n in nn], 0 < w(\allowbreak{}n)) and series-converges(\allowbreak{}w))
\end{quote}
mentioned by 59 result(s)

\smallskip

\noindent\textbf{\texttt{sums-to(grp, f, r)}}\quad\textit{predicate}\index{sums-to(grp, f, r)@\texttt{sums-to(grp, f, r)}}\\
reads: \$2 sums to \$3 in \$1 \\
defined by \texttt{sums-to} \\
characterized by \texttt{sums-to}  \textit{(definitional)}:
\begin{quote}\ttfamily\footnotesize\raggedright\noindent
forall(\allowbreak{}[grp, f, r], sums-to(\allowbreak{}grp, f, r) iff forall(\allowbreak{}[eps], pos-rr(\allowbreak{}eps) implies forsome(\allowbreak{}[fin in set], card(\allowbreak{}fin) in nn and fin subset dom(\allowbreak{}f) and forall(\allowbreak{}[ext], ext in set and card(\allowbreak{}ext) in nn and fin subset ext and ext subset dom(\allowbreak{}f) implies (\allowbreak{}nrm(\allowbreak{}grp))(\allowbreak{}(\allowbreak{}opr(\allowbreak{}grp))(\allowbreak{}r, (\allowbreak{}inv(\allowbreak{}grp))(\allowbreak{}finsum(\allowbreak{}normed-ag-as-abelian-group(\allowbreak{}grp), f, ext)))) < eps))))
\end{quote}
also: \texttt{sums-to-unique-cna-as-normed-ag} \texttt{sums-to-unique-uba-as-normed-ag} \texttt{sums-to-unique-cnvs-as-normed-ag} \texttt{sums-to-unique-uba-as-normed-ag-uiba-as-uba} \texttt{bop-sums-to} \texttt{sums-to-add} \texttt{sums-to-unique} \texttt{cc-sums-to-comb} ... (16 in all) \\
mentioned by 53 result(s)

\smallskip

\noindent\textbf{\texttt{sumset(m, a, p)}}\quad\textit{functoid}\index{sumset(m, a, p)@\texttt{sumset(m, a, p)}}\\
reads: the sum set \$2 + \$3 in \$1 \\
definition:
\begin{quote}\ttfamily\footnotesize\raggedright\noindent
\{smx\_ in vec(\allowbreak{}m): forsome(\allowbreak{}[smz\_ in p], (\allowbreak{}vadd(\allowbreak{}m))(\allowbreak{}smx\_, (\allowbreak{}act(\allowbreak{}m))(\allowbreak{}-1, smz\_)) in a)\}
\end{quote}
characterized by \texttt{sumset-open}  \textit{(proven)}:
\begin{quote}\ttfamily\footnotesize\raggedright\noindent
forall(\allowbreak{}[m, t, a, p], is-tvs(\allowbreak{}rr-normed-field, m, t) implies a in opens(\allowbreak{}t) implies p subset vec(\allowbreak{}m) implies sumset(\allowbreak{}m, a, p) in opens(\allowbreak{}t))
\end{quote}
also: \texttt{sumset-convex} \texttt{sumset-unfold} \\
mentioned by 4 result(s)

\smallskip

\noindent\textbf{\texttt{sup}}\quad\textit{functoid}\index{sup@\texttt{sup}}\\
reads: the least upper bound of \$1 \\
characterized by \texttt{rr-sup-in}  \textit{(primitive)}:
\begin{quote}\ttfamily\footnotesize\raggedright\noindent
forall(\allowbreak{}[s], s subset rr implies forsome(\allowbreak{}[x], x in s) implies rr-bounded-above(\allowbreak{}s) implies sup(\allowbreak{}s) in rr)
\end{quote}
also: \texttt{rr-sup-least} \texttt{rr-sup-upper} \texttt{rr-sup-approx} \texttt{rr-neg-sup-neg-is-glb} \texttt{monotone-converges-to-sup} \texttt{taylor-remainder-sup-bound} \\
mentioned by 10 result(s)

\smallskip

\noindent\textbf{\texttt{sup-dist(sds\_, sdt\_, sdf\_, sdg\_)}}\quad\textit{functoid}\index{sup-dist(sds_, sdt_, sdf_, sdg_)@\texttt{sup-dist(sds\_, sdt\_, sdf\_, sdg\_)}}\\
reads: the supremum over the points x of \$1 of the distance in \$2 from \$3(x) to \$4(x) \\
definition:
\begin{quote}\ttfamily\footnotesize\raggedright\noindent
esup(\allowbreak{}image(\allowbreak{}vnb-lambda(\allowbreak{}sdx\_, pts(\allowbreak{}sds\_), (\allowbreak{}dist(\allowbreak{}sdt\_))(\allowbreak{}sdf\_(\allowbreak{}sdx\_), sdg\_(\allowbreak{}sdx\_))), pts(\allowbreak{}sds\_)))
\end{quote}
characterized by \texttt{sup-dist-bound}  \textit{(proven)}:
\begin{quote}\ttfamily\footnotesize\raggedright\noindent
forall(\allowbreak{}[s, t, f, g, b], is-metric-space(\allowbreak{}t) implies f in fun(\allowbreak{}pts(\allowbreak{}s), pts(\allowbreak{}t)) implies g in fun(\allowbreak{}pts(\allowbreak{}s), pts(\allowbreak{}t)) implies b in rr implies 0 <= b implies forall(\allowbreak{}[x in pts(\allowbreak{}s)], (\allowbreak{}dist(\allowbreak{}t))(\allowbreak{}f(\allowbreak{}x), g(\allowbreak{}x)) <= b) implies sup-dist(\allowbreak{}s, t, f, g) in rr and 0 <= sup-dist(\allowbreak{}s, t, f, g) and sup-dist(\allowbreak{}s, t, f, g) <= b)
\end{quote}
also: \texttt{sup-dist-upper} \texttt{sup-dist-pos-star} \texttt{metric-top@sup-dist} \texttt{trunc-metric@sup-dist} \\
mentioned by 14 result(s)

\smallskip

\noindent\textbf{\texttt{sup-norm-t(kxt\_, kxf\_)}}\quad\textit{functoid}\index{sup-norm-t(kxt_, kxf_)@\texttt{sup-norm-t(kxt\_, kxf\_)}}\\
reads: the supremum norm of \$2 on \$1 \\
definition:
\begin{quote}\ttfamily\footnotesize\raggedright\noindent
esup(\allowbreak{}image(\allowbreak{}vnb-lambda(\allowbreak{}kxp\_, pts(\allowbreak{}kxt\_), magnitude(\allowbreak{}kxf\_(\allowbreak{}kxp\_))), pts(\allowbreak{}kxt\_)))
\end{quote}
characterized by \texttt{sup-norm-t-one}  \textit{(proven)}:
\begin{quote}\ttfamily\footnotesize\raggedright\noindent
forall(\allowbreak{}[t, h], is-compact-t(\allowbreak{}t) implies forsome(\allowbreak{}[tp\_], tp\_ in pts(\allowbreak{}t)) implies h in fun(\allowbreak{}pts(\allowbreak{}t), cc) implies forall(\allowbreak{}[rz\_ in pts(\allowbreak{}t)], h(\allowbreak{}rz\_) = 1) implies sup-norm-t(\allowbreak{}t, h) = 1)
\end{quote}
also: \texttt{sup-norm-t-conj} \texttt{sup-norm-t-real} \texttt{sup-norm-t-scal} \texttt{sup-norm-t-cstar} \texttt{sup-norm-t-least} \texttt{sup-norm-t-upper} \texttt{sup-norm-t-add-le} \texttt{sup-norm-t-mul-le} ... (12 in all) \\
mentioned by 30 result(s)

\smallskip

\noindent\textbf{\texttt{sup-ord}}\quad\textit{operator}\index{sup-ord@\texttt{sup-ord}}\\
a term-forming head \\
characterized by \texttt{sup-ord-in}  \textit{(primitive)}:
\begin{quote}\ttfamily\footnotesize\raggedright\noindent
forall(\allowbreak{}[a], a in set and forall(\allowbreak{}[x in a], x in ord) implies sup-ord(\allowbreak{}a) in ord)
\end{quote}
also: \texttt{sup-ord-empty} \texttt{sup-ord-least} \texttt{sup-ord-upper} \texttt{sup-ord-succ-segment} \\
mentioned by 9 result(s)

\smallskip

\noindent\textbf{\texttt{supp(a, m, f)}}\quad\textit{functoid}\index{supp(a, m, f)@\texttt{supp(a, m, f)}}\\
reads: the support of \$3 \\
definition:
\begin{quote}\ttfamily\footnotesize\raggedright\noindent
\{x\_ in carr(\allowbreak{}m): not(\allowbreak{}f(\allowbreak{}x\_) = zero(\allowbreak{}a))\}
\end{quote}
characterized by \texttt{supp-in-set}  \textit{(proven)}:
\begin{quote}\ttfamily\footnotesize\raggedright\noindent
forall(\allowbreak{}[a\_, m\_, f\_], carr(\allowbreak{}m\_) in set implies supp(\allowbreak{}a\_, m\_, f\_) in set)
\end{quote}
also: \texttt{supp-unfold} \texttt{finsupp-unfold} \texttt{supp-add-subset} \texttt{supp-membership} \texttt{monomial-support} \texttt{finsupp-membership} \texttt{monalg-one-supp-subset} \texttt{monalg-mul-supp-in-image} ... (10 in all) \\
mentioned by 43 result(s)

\smallskip

\noindent\textbf{\texttt{supply}}\quad\textit{tactic}\index{supply@\texttt{supply}}\\
\texttt{(supply)} -- Close an inequality goal the way a hand proof would: beta-reduce any applied lambda, land the typing certificates and the standard bounds the linear oracle cannot see -- including the TRIANGLE inequality at the summands of an abs of a sum -- then call ineq over every premise.  Rehearses the whole sequence on a throwaway copy and does NOTHING unless it closes; on a miss it prints the sequence it tried.  Adds no trust: it drives lam-b, fact and ineq, each of which records itself. \\
declared in \texttt{tactics-help.scm}

\smallskip

\noindent\textbf{\texttt{sym-difference(sda\_, sdb\_)}}\quad\textit{functoid}\index{sym-difference(sda_, sdb_)@\texttt{sym-difference(sda\_, sdb\_)}}\\
reads: the symmetric difference of \$1 and \$2 \\
definition:
\begin{quote}\ttfamily\footnotesize\raggedright\noindent
union(\allowbreak{}complement-in(\allowbreak{}sda\_, sdb\_), complement-in(\allowbreak{}sdb\_, sda\_))
\end{quote}
characterized by \texttt{sym-difference-unfold}  \textit{(proven)}:
\begin{quote}\ttfamily\footnotesize\raggedright\noindent
forall(\allowbreak{}[sda\_, sdb\_], sym-difference(\allowbreak{}sda\_, sdb\_) == union(\allowbreak{}complement-in(\allowbreak{}sda\_, sdb\_), complement-in(\allowbreak{}sdb\_, sda\_)))
\end{quote}
mentioned by 4 result(s)

\smallskip


\section*{T}
\addcontentsline{toc}{section}{T}

\noindent\textbf{\texttt{ta}}\quad\textit{tactic}\index{ta@\texttt{ta}}\\
\texttt{(ta 'name)} -- Theorem-assumption: bring the named installed theorem into context as an assumption. \\
declared in \texttt{tactics-help.scm}

\smallskip

\noindent\textbf{\texttt{taylor-aux(d, n, a, b)}}\quad\textit{functoid}\index{taylor-aux(d, n, a, b)@\texttt{taylor-aux(d, n, a, b)}}\\
reads: the Taylor auxiliary function of \$1 of order \$2 on \$3 to \$4 \\
definition:
\begin{quote}\ttfamily\footnotesize\raggedright\noindent
vnb-lambda(\allowbreak{}tu\_, ccint(\allowbreak{}a, b), taylor-sum(\allowbreak{}d, n, b, tu\_))
\end{quote}
characterized by \texttt{taylor-aux-step}  \textit{(proven)}:
\begin{quote}\ttfamily\footnotesize\raggedright\noindent
forall(\allowbreak{}[itd\_, itn\_, a, b], is-taylor-family(\allowbreak{}itd\_, itn\_, a, b) implies forall(\allowbreak{}[itj\_ in nn], succ(\allowbreak{}itj\_) <= itn\_ implies forall(\allowbreak{}[itz\_ in ccint(\allowbreak{}a, b)], (\allowbreak{}taylor-aux(\allowbreak{}itd\_, succ(\allowbreak{}itj\_), a, b))(\allowbreak{}itz\_) = (\allowbreak{}taylor-aux(\allowbreak{}itd\_, itj\_, a, b))(\allowbreak{}itz\_) + (\allowbreak{}(\allowbreak{}itd\_(\allowbreak{}itj\_))(\allowbreak{}itz\_) * (\allowbreak{}b - itz\_) \textasciicircum{} itj\_) * recip(\allowbreak{}factorial(\allowbreak{}itj\_)))))
\end{quote}
also: \texttt{taylor-aux-value} \texttt{taylor-aux-in-fun} \texttt{taylor-aux-has-deriv} \texttt{taylor-aux-continuous-on} \\
mentioned by 8 result(s)

\smallskip

\noindent\textbf{\texttt{taylor-differentiable(f, a, x, n)}}\quad\textit{predicate}\index{taylor-differentiable(f, a, x, n)@\texttt{taylor-differentiable(f, a, x, n)}}\\
reads: \$1 is \$4 times differentiable from \$2 to \$3 \\
defined by \texttt{taylor-differentiable} \\
characterized by \texttt{taylor-differentiable}  \textit{(definitional)}:
\begin{quote}\ttfamily\footnotesize\raggedright\noindent
forall(\allowbreak{}[f, a, x, n], taylor-differentiable(\allowbreak{}f, a, x, n) iff forall(\allowbreak{}[k], k in nn and k <= n implies forall(\allowbreak{}[t in ccint(\allowbreak{}a, x)], is-continuous-at(\allowbreak{}rr-ms, rr-ms, nth-deriv(\allowbreak{}f, k), t))) and forall(\allowbreak{}[k], k in nn and k <= n implies forall(\allowbreak{}[t], a < t and t < x implies is-diff-at(\allowbreak{}nth-deriv(\allowbreak{}f, k), t, (\allowbreak{}nth-deriv(\allowbreak{}f, succ(\allowbreak{}k)))(\allowbreak{}t)))))
\end{quote}
mentioned by 10 result(s)

\smallskip

\noindent\textbf{\texttt{taylor-differentiable-v(m, f, a, x, n)}}\quad\textit{predicate}\index{taylor-differentiable-v(m, f, a, x, n)@\texttt{taylor-differentiable-v(m, f, a, x, n)}}\\
reads: \$2 is \$5 times differentiable from \$3 to \$4, as a curve in \$1 \\
defined by \texttt{taylor-differentiable-v} \\
characterized by \texttt{taylor-differentiable-v}  \textit{(definitional)}:
\begin{quote}\ttfamily\footnotesize\raggedright\noindent
forall(\allowbreak{}[m, f, a, x, n], taylor-differentiable-v(\allowbreak{}m, f, a, x, n) iff forall(\allowbreak{}[k], k in nn and k <= n implies forall(\allowbreak{}[t in ccint(\allowbreak{}a, x)], is-continuous-at(\allowbreak{}rr-ms, nvs-metric-space(\allowbreak{}m), nth-deriv-v(\allowbreak{}m, f, k), t))) and forall(\allowbreak{}[k], k in nn and k <= n implies forall(\allowbreak{}[t], a < t and t < x implies is-diff-at-v(\allowbreak{}m, nth-deriv-v(\allowbreak{}m, f, k), t, (\allowbreak{}nth-deriv-v(\allowbreak{}m, f, succ(\allowbreak{}k)))(\allowbreak{}t)))))
\end{quote}
mentioned by 9 result(s)

\smallskip

\noindent\textbf{\texttt{taylor-poly(f, a, n, x)}}\quad\textit{functoid}\index{taylor-poly(f, a, n, x)@\texttt{taylor-poly(f, a, n, x)}}\\
definition:
\begin{quote}\ttfamily\footnotesize\raggedright\noindent
series-partial-sum(\allowbreak{}vnb-lambda(\allowbreak{}k, nn, (\allowbreak{}(\allowbreak{}nth-deriv(\allowbreak{}f, k))(\allowbreak{}a) * (\allowbreak{}x - a) \textasciicircum{} k) * recip(\allowbreak{}factorial(\allowbreak{}k))), succ(\allowbreak{}n))
\end{quote}
characterized by \texttt{taylor-poly-in-rr}  \textit{(proven)}:
\begin{quote}\ttfamily\footnotesize\raggedright\noindent
forall(\allowbreak{}[f in fun(\allowbreak{}rr, rr), a in rr, n in nn, x in rr], forall(\allowbreak{}[k], k in nn and k <= n implies nth-deriv(\allowbreak{}f, k) in fun(\allowbreak{}rr, rr)) implies taylor-poly(\allowbreak{}f, a, n, x) in rr)
\end{quote}
also: \texttt{taylor-poly-at-center} \\
mentioned by 19 result(s)

\smallskip

\noindent\textbf{\texttt{taylor-poly-v}}\quad\textit{defined-fn}\index{taylor-poly-v@\texttt{taylor-poly-v}}\\
defined by \texttt{taylor-poly-v-zero} \texttt{taylor-poly-v-succ} \\
characterized by \texttt{taylor-poly-v-succ}  \textit{(definitional)}:
\begin{quote}\ttfamily\footnotesize\raggedright\noindent
forall(\allowbreak{}[m, f, a, x, n in nn], taylor-poly-v(\allowbreak{}m, f, a, x, succ(\allowbreak{}n)) == (\allowbreak{}vadd(\allowbreak{}m))(\allowbreak{}taylor-poly-v(\allowbreak{}m, f, a, x, n), (\allowbreak{}act(\allowbreak{}m))(\allowbreak{}(\allowbreak{}x - a) \textasciicircum{} succ(\allowbreak{}n) * recip(\allowbreak{}factorial(\allowbreak{}succ(\allowbreak{}n))), (\allowbreak{}nth-deriv-v(\allowbreak{}m, f, succ(\allowbreak{}n)))(\allowbreak{}a))))
\end{quote}
also: \texttt{taylor-poly-v-zero} \\
mentioned by 14 result(s)

\smallskip

\noindent\textbf{\texttt{taylor-rem(d, n, a, b)}}\quad\textit{functoid}\index{taylor-rem(d, n, a, b)@\texttt{taylor-rem(d, n, a, b)}}\\
reads: the Taylor remainder of \$1 of order \$2 from \$3 to \$4 \\
definition:
\begin{quote}\ttfamily\footnotesize\raggedright\noindent
(\allowbreak{}d(\allowbreak{}0))(\allowbreak{}b) - taylor-sum(\allowbreak{}d, n, b, a)
\end{quote}
characterized by \texttt{taylor-remainder-factor}  \textit{(proven)}:
\begin{quote}\ttfamily\footnotesize\raggedright\noindent
forall(\allowbreak{}[a in rr, b in rr, tcf\_, tcm\_ in nn], a < b implies forall(\allowbreak{}[tck\_ in nn], tck\_ <= succ(\allowbreak{}tcm\_) implies tcf\_(\allowbreak{}tck\_) in fun(\allowbreak{}ccint(\allowbreak{}a, b), rr)) implies forall(\allowbreak{}[tck\_ in nn], tck\_ < succ(\allowbreak{}tcm\_) implies is-continuous-on(\allowbreak{}tcf\_(\allowbreak{}tck\_), ccint(\allowbreak{}a, b))) implies forall(\allowbreak{}[tck\_ in nn], tck\_ < succ(\allowbreak{}tcm\_) implies forall(\allowbreak{}[tcx\_ in ooint(\allowbreak{}a, b)], has-deriv-at(\allowbreak{}tcf\_(\allowbreak{}tck\_), tcx\_, (\allowbreak{}tcf\_(\allowbreak{}succ(\allowbreak{}tck\_)))(\allowbreak{}tcx\_)))) implies is-continuous-on(\allowbreak{}tcf\_(\allowbreak{}succ(\allowbreak{}tcm\_)), ccint(\allowbreak{}a, b)) implies forsome(\allowbreak{}[tcs\_ in fun(\allowbreak{}ccint(\allowbreak{}a, b), rr)], is-continuous-on(\allowbreak{}tcs\_, ccint(\allowbreak{}a, b)) and tcs\_(\allowbreak{}a) = (\allowbreak{}tcf\_(\allowbreak{}succ(\allowbreak{}tcm\_)))(\allowbreak{}a) and forall(\allowbreak{}[tct\_ in ccint(\allowbreak{}a, b)], taylor-rem(\allowbreak{}tcf\_, succ(\allowbreak{}tcm\_), a, tct\_) = (\allowbreak{}tcs\_(\allowbreak{}tct\_) * (\allowbreak{}tct\_ - a) \textasciicircum{} succ(\allowbreak{}tcm\_)) * recip(\allowbreak{}factorial(\allowbreak{}succ(\allowbreak{}tcm\_)))) and forall(\allowbreak{}[tct\_ in ccint(\allowbreak{}a, b)], abs(\allowbreak{}tcs\_(\allowbreak{}tct\_)) <= sup(\allowbreak{}image(\allowbreak{}vnb-lambda(\allowbreak{}tcv\_, ccint(\allowbreak{}a, b), abs(\allowbreak{}(\allowbreak{}tcf\_(\allowbreak{}succ(\allowbreak{}tcm\_)))(\allowbreak{}tcv\_))), ccint(\allowbreak{}a, b))))))
\end{quote}
also: \texttt{taylor-remainder-sup-bound} \\
mentioned by 7 result(s)

\smallskip

\noindent\textbf{\texttt{taylor-sum(d, n, b, tv\_)}}\quad\textit{functoid}\index{taylor-sum(d, n, b, tv_)@\texttt{taylor-sum(d, n, b, tv\_)}}\\
reads: the Taylor sum of \$1 of order \$2 towards \$3 at \$4 \\
definition:
\begin{quote}\ttfamily\footnotesize\raggedright\noindent
series-partial-sum(\allowbreak{}taylor-term(\allowbreak{}d, b, tv\_), n)
\end{quote}
characterized by \texttt{taylor-sum-one}  \textit{(proven)}:
\begin{quote}\ttfamily\footnotesize\raggedright\noindent
forall(\allowbreak{}[itd\_, b in rr, itv\_ in rr], (\allowbreak{}itd\_(\allowbreak{}0))(\allowbreak{}itv\_) in rr implies taylor-sum(\allowbreak{}itd\_, succ(\allowbreak{}0), b, itv\_) = (\allowbreak{}itd\_(\allowbreak{}0))(\allowbreak{}itv\_))
\end{quote}
also: \texttt{taylor-sum-in-rr} \texttt{taylor-sum-unfold} \texttt{taylor-sum-continuous} \texttt{taylor-sum-at-endpoint} \\
mentioned by 16 result(s)

\smallskip

\noindent\textbf{\texttt{taylor-term(d, b, tv\_)}}\quad\textit{functoid}\index{taylor-term(d, b, tv_)@\texttt{taylor-term(d, b, tv\_)}}\\
reads: the Taylor summand of the family \$1 towards \$2 at \$3 \\
definition:
\begin{quote}\ttfamily\footnotesize\raggedright\noindent
vnb-lambda(\allowbreak{}tk\_, nn, (\allowbreak{}(\allowbreak{}d(\allowbreak{}tk\_))(\allowbreak{}tv\_) * (\allowbreak{}b - tv\_) \textasciicircum{} tk\_) * recip(\allowbreak{}factorial(\allowbreak{}tk\_)))
\end{quote}
characterized by \texttt{taylor-term-value}  \textit{(proven)}:
\begin{quote}\ttfamily\footnotesize\raggedright\noindent
forall(\allowbreak{}[itd\_, b, itv\_, itk\_ in nn], (\allowbreak{}taylor-term(\allowbreak{}itd\_, b, itv\_))(\allowbreak{}itk\_) == (\allowbreak{}(\allowbreak{}itd\_(\allowbreak{}itk\_))(\allowbreak{}itv\_) * (\allowbreak{}b - itv\_) \textasciicircum{} itk\_) * recip(\allowbreak{}factorial(\allowbreak{}itk\_)))
\end{quote}
mentioned by 4 result(s)

\smallskip

\noindent\textbf{\texttt{te}}\quad\textit{tactic}\index{te@\texttt{te}}\\
\texttt{(te hyp k)} -- Tuple elim: project the k-th component out of a tuple-membership assumption. \\
declared in \texttt{tactics-help.scm}

\smallskip

\noindent\textbf{\texttt{tf-closed(mitfa\_, mitfb\_, mitfs\_)}}\quad\textit{predicate}\index{tf-closed(mitfa_, mitfb_, mitfs_)@\texttt{tf-closed(mitfa\_, mitfb\_, mitfs\_)}}\\
reads: \$3 is closed under the right folds of the component equations of \$1 and \$2 \\
defined by \texttt{tf-closed} \\
characterized by \texttt{tf-closed}  \textit{(definitional)}:
\begin{quote}\ttfamily\footnotesize\raggedright\noindent
forall(\allowbreak{}[mitfa\_, mitfb\_, mitfs\_], tf-closed(\allowbreak{}mitfa\_, mitfb\_, mitfs\_) iff forall(\allowbreak{}[mitfj\_ in nn], succ(\allowbreak{}mitfj\_) = length(\allowbreak{}mitfa\_) implies [mitfj\_, enc-eq(\allowbreak{}nth(\allowbreak{}succ(\allowbreak{}mitfj\_), mitfa\_), nth(\allowbreak{}succ(\allowbreak{}mitfj\_), mitfb\_))] in mitfs\_) and forall(\allowbreak{}[mitfk\_, mitff\_], mitfk\_ in nn implies [succ(\allowbreak{}mitfk\_), mitff\_] in mitfs\_ implies [mitfk\_, enc-and(\allowbreak{}enc-eq(\allowbreak{}nth(\allowbreak{}succ(\allowbreak{}mitfk\_), mitfa\_), nth(\allowbreak{}succ(\allowbreak{}mitfk\_), mitfb\_)), mitff\_)] in mitfs\_))
\end{quote}
mentioned by 5 result(s)

\smallskip

\noindent\textbf{\texttt{tf-set(mitfa\_, mitfb\_)}}\quad\textit{functoid}\index{tf-set(mitfa_, mitfb_)@\texttt{tf-set(mitfa\_, mitfb\_)}}\\
reads: the right folds of the component equations of \$1 and \$2 from an index on \\
definition:
\begin{quote}\ttfamily\footnotesize\raggedright\noindent
\{mitfp\_ in cartesian(\allowbreak{}nn, expr-set): forall(\allowbreak{}[mitfw\_ in power(\allowbreak{}cartesian(\allowbreak{}nn, expr-set))], tf-closed(\allowbreak{}mitfa\_, mitfb\_, mitfw\_) implies mitfp\_ in mitfw\_)\}
\end{quote}
characterized by \texttt{tf-set-least}  \textit{(proven)}:
\begin{quote}\ttfamily\footnotesize\raggedright\noindent
forall(\allowbreak{}[a, b, mnvxs\_ in power(\allowbreak{}cartesian(\allowbreak{}nn, expr-set))], tf-closed(\allowbreak{}a, b, mnvxs\_) implies forall(\allowbreak{}[z in tf-set(\allowbreak{}a, b)], z in mnvxs\_))
\end{quote}
also: \texttt{tf-set-unfold} \\
mentioned by 5 result(s)

\smallskip

\noindent\textbf{\texttt{ti}}\quad\textit{tactic}\index{ti@\texttt{ti}}\\
\texttt{(ti)} -- Tuple intro: prove a tuple/LIST membership component-wise. \\
declared in \texttt{tactics-help.scm}

\smallskip

\noindent\textbf{\texttt{tokens(s)}}\quad\textit{accessor}\index{tokens(s)@\texttt{tokens(s)}}\\
reads: the tokens slot of \$1 \\
a structure slot \\
characterized by \texttt{cx-sp@tokens}  \textit{(definitional)}:
\begin{quote}\ttfamily\footnotesize\raggedright\noindent
tokens(\allowbreak{}cx-sp) == nn
\end{quote}
also: \texttt{keys-in-tokens} \\
mentioned by 121 result(s)

\smallskip

\noindent\textbf{\texttt{top-closure(otq\_, ota\_)}}\quad\textit{functoid}\index{top-closure(otq_, ota_)@\texttt{top-closure(otq\_, ota\_)}}\\
reads: the closure of \$2 in \$1 \\
definition:
\begin{quote}\ttfamily\footnotesize\raggedright\noindent
\{otx\_ in pts(\allowbreak{}otq\_): forall(\allowbreak{}[otu\_ in opens(\allowbreak{}otq\_)], otx\_ in otu\_ implies forsome(\allowbreak{}[oty\_ in ota\_], oty\_ in otu\_))\}
\end{quote}
characterized by \texttt{top-closure-unfold}  \textit{(proven)}:
\begin{quote}\ttfamily\footnotesize\raggedright\noindent
forall(\allowbreak{}[q\_, a\_], top-closure(\allowbreak{}q\_, a\_) == \{otx\_ in pts(\allowbreak{}q\_): forall(\allowbreak{}[otu\_ in opens(\allowbreak{}q\_)], otx\_ in otu\_ implies forsome(\allowbreak{}[oty\_ in a\_], oty\_ in otu\_))\})
\end{quote}
also: \texttt{top-closure-closed-mem} \\
mentioned by 7 result(s)

\smallskip

\noindent\textbf{\texttt{top-converges-to(t, f, l)}}\quad\textit{predicate}\index{top-converges-to(t, f, l)@\texttt{top-converges-to(t, f, l)}}\\
reads: \$2 converges to \$3 in the topology of \$1 \\
defined by \texttt{top-converges-to} \\
characterized by \texttt{top-converges-to}  \textit{(definitional)}:
\begin{quote}\ttfamily\footnotesize\raggedright\noindent
forall(\allowbreak{}[t, f, l], top-converges-to(\allowbreak{}t, f, l) iff f in fun(\allowbreak{}nn, pts(\allowbreak{}t)) and l in pts(\allowbreak{}t) and forall(\allowbreak{}[ut\_ in opens(\allowbreak{}t)], l in ut\_ implies forsome(\allowbreak{}[nt\_ in nn], forall(\allowbreak{}[mt\_ in nn], nt\_ <= mt\_ implies f(\allowbreak{}mt\_) in ut\_))))
\end{quote}
mentioned by 6 result(s)

\smallskip

\noindent\textbf{\texttt{top-dual(m, t)}}\quad\textit{functoid}\index{top-dual(m, t)@\texttt{top-dual(m, t)}}\\
reads: the topological dual of the topological vector space \$1 with topology \$2 \\
definition:
\begin{quote}\ttfamily\footnotesize\raggedright\noindent
\{tdf\_ in fun(\allowbreak{}vec(\allowbreak{}m), rr): is-linear-functional(\allowbreak{}m, tdf\_) and is-hom-top-space(\allowbreak{}t, seminorm-top(\allowbreak{}normed-vector-space-as-module(\allowbreak{}rr-nvs), pair(\allowbreak{}0, 0), vnb-lambda(\allowbreak{}nvk\_, pair(\allowbreak{}0, 0), vnrm(\allowbreak{}rr-nvs))), tdf\_)\}
\end{quote}
characterized by \texttt{top-dual-mem}  \textit{(proven)}:
\begin{quote}\ttfamily\footnotesize\raggedright\noindent
forall(\allowbreak{}[m, t, z], z in top-dual(\allowbreak{}m, t) iff z in fun(\allowbreak{}vec(\allowbreak{}m), rr) and is-linear-functional(\allowbreak{}m, z) and is-hom-top-space(\allowbreak{}t, seminorm-top(\allowbreak{}normed-vector-space-as-module(\allowbreak{}rr-nvs), pair(\allowbreak{}0, 0), vnb-lambda(\allowbreak{}nvk\_, pair(\allowbreak{}0, 0), vnrm(\allowbreak{}rr-nvs))), z))
\end{quote}
also: \texttt{top-dual-unfold} \\
mentioned by 10 result(s)

\smallskip

\noindent\textbf{\texttt{top-group}}\quad\textit{structure}\index{top-group@\texttt{top-group}}\\
predicate \texttt{is-top-group}; slots: carr opr iden inv topology

\smallskip

\noindent\textbf{\texttt{top-group-group(tgr\_)}}\quad\textit{functoid}\index{top-group-group(tgr_)@\texttt{top-group-group(tgr\_)}}\\
reads: the group of \$1 \\
definition:
\begin{quote}\ttfamily\footnotesize\raggedright\noindent
[carr(\allowbreak{}tgr\_), opr(\allowbreak{}tgr\_), iden(\allowbreak{}tgr\_), inv(\allowbreak{}tgr\_)]
\end{quote}
characterized by \texttt{top-group-group-is-group}  \textit{(proven)}:
\begin{quote}\ttfamily\footnotesize\raggedright\noindent
forall(\allowbreak{}[g], is-top-group(\allowbreak{}g) implies is-group(\allowbreak{}top-group-group(\allowbreak{}g)))
\end{quote}
also: \texttt{top-group-group-topology} \texttt{top-group-group-functorial} \\
mentioned by 30 result(s)

\smallskip

\noindent\textbf{\texttt{top-group-space(tgr\_)}}\quad\textit{functoid}\index{top-group-space(tgr_)@\texttt{top-group-space(tgr\_)}}\\
reads: the topological space of \$1 \\
definition:
\begin{quote}\ttfamily\footnotesize\raggedright\noindent
[carr(\allowbreak{}tgr\_), topology(\allowbreak{}tgr\_)]
\end{quote}
characterized by \texttt{top-group-space-functorial}  \textit{(proven)}:
\begin{quote}\ttfamily\footnotesize\raggedright\noindent
forall(\allowbreak{}[a, b, f], is-hom-top-group(\allowbreak{}a, b, f) implies is-hom-top-space(\allowbreak{}top-group-space(\allowbreak{}a), top-group-space(\allowbreak{}b), f))
\end{quote}
also: \texttt{top-group-space-is-top-space} \\
mentioned by 30 result(s)

\smallskip

\noindent\textbf{\texttt{top-space}}\quad\textit{structure}\index{top-space@\texttt{top-space}}\\
predicate \texttt{is-top-space}; slots: pts opens

\smallskip

\noindent\textbf{\texttt{topology(s)}}\quad\textit{accessor}\index{topology(s)@\texttt{topology(s)}}\\
reads: the topology of \$1 \\
a structure slot \\
characterized by \texttt{make-top-group-topology}  \textit{(proven)}:
\begin{quote}\ttfamily\footnotesize\raggedright\noindent
forall(\allowbreak{}[s, g], is-top-space(\allowbreak{}s) implies is-group(\allowbreak{}g) implies topology(\allowbreak{}make-top-group(\allowbreak{}s, g)) == opens(\allowbreak{}s))
\end{quote}
mentioned by 6 result(s)

\smallskip

\noindent\textbf{\texttt{total-head(mchd\_)}}\quad\textit{predicate}\index{total-head(mchd_)@\texttt{total-head(mchd\_)}}\\
reads: the head \$1 is a total class constructor \\
defined by \texttt{total-head} \\
characterized by \texttt{total-head}  \textit{(definitional)}:
\begin{quote}\ttfamily\footnotesize\raggedright\noindent
forall(\allowbreak{}[mchd\_], total-head(\allowbreak{}mchd\_) iff 11 <= mchd\_ and mchd\_ <= 14 or mchd\_ = 19 or 30 <= mchd\_ and mchd\_ <= 31)
\end{quote}
mentioned by 7 result(s)

\smallskip

\noindent\textbf{\texttt{total-variation(ca, nu)}}\quad\textit{functoid}\index{total-variation(ca, nu)@\texttt{total-variation(ca, nu)}}\\
reads: the total variation of \$2 \\
definition:
\begin{quote}\ttfamily\footnotesize\raggedright\noindent
vnb-lambda(\allowbreak{}tva\_, ca, eplus(\allowbreak{}(\allowbreak{}pos-variation(\allowbreak{}ca, nu))(\allowbreak{}tva\_), (\allowbreak{}neg-variation(\allowbreak{}ca, nu))(\allowbreak{}tva\_)))
\end{quote}
characterized by \texttt{total-variation-value}  \textit{(proven)}:
\begin{quote}\ttfamily\footnotesize\raggedright\noindent
forall(\allowbreak{}[ca, nu, a in ca], (\allowbreak{}total-variation(\allowbreak{}ca, nu))(\allowbreak{}a) == eplus(\allowbreak{}(\allowbreak{}pos-variation(\allowbreak{}ca, nu))(\allowbreak{}a), (\allowbreak{}neg-variation(\allowbreak{}ca, nu))(\allowbreak{}a)))
\end{quote}
also: \texttt{total-variation-is-measure} \\
mentioned by 3 result(s)

\smallskip

\noindent\textbf{\texttt{total-variation-norm(omega, ca, nu)}}\quad\textit{functoid}\index{total-variation-norm(omega, ca, nu)@\texttt{total-variation-norm(omega, ca, nu)}}\\
reads: the total variation norm of \$3 \\
definition:
\begin{quote}\ttfamily\footnotesize\raggedright\noindent
(\allowbreak{}total-variation(\allowbreak{}ca, nu))(\allowbreak{}omega)
\end{quote}
characterized by \texttt{total-variation-norm-finite-iff}  \textit{(proven)}:
\begin{quote}\ttfamily\footnotesize\raggedright\noindent
forall(\allowbreak{}[omega, ca, nu], is-signed-measure(\allowbreak{}omega, ca, nu) implies total-variation-norm(\allowbreak{}omega, ca, nu) in rr iff forall(\allowbreak{}[a\_ in ca], (\allowbreak{}pos-variation(\allowbreak{}ca, nu))(\allowbreak{}a\_) in rr) and forall(\allowbreak{}[a\_ in ca], (\allowbreak{}neg-variation(\allowbreak{}ca, nu))(\allowbreak{}a\_) in rr))
\end{quote}
mentioned by 2 result(s)

\smallskip

\noindent\textbf{\texttt{totally-bounded(s)}}\quad\textit{predicate}\index{totally-bounded(s)@\texttt{totally-bounded(s)}}\\
reads: \$1 is totally bounded \\
defined by \texttt{totally-bounded} \\
characterized by \texttt{totally-bounded}  \textit{(definitional)}:
\begin{quote}\ttfamily\footnotesize\raggedright\noindent
forall(\allowbreak{}[s], totally-bounded(\allowbreak{}s) iff is-metric-space(\allowbreak{}s) and forall(\allowbreak{}[r], r in rr and 0 <= r and not(\allowbreak{}0 = r) implies forsome(\allowbreak{}[f], card(\allowbreak{}f) in nn and is-r-net(\allowbreak{}s, f, pts(\allowbreak{}s), r))))
\end{quote}
also: \texttt{qp-zp-totally-bounded} \texttt{ultra-totally-bounded-iff} \texttt{compact-implies-totally-bounded} \texttt{seq-compact-implies-totally-bounded} \texttt{totally-bounded-has-cauchy-subseq-rad} \texttt{totally-bounded-has-cauchy-subsequence} \\
mentioned by 19 result(s)

\smallskip

\noindent\textbf{\texttt{tprod-box(ix, tf, bb)}}\quad\textit{functoid}\index{tprod-box(ix, tf, bb)@\texttt{tprod-box(ix, tf, bb)}}\\
reads: the box of the constraints \$3 \\
definition:
\begin{quote}\ttfamily\footnotesize\raggedright\noindent
\{tpy\_ in tprod-pts(\allowbreak{}ix, tf): forall(\allowbreak{}[tpc\_ in bb], tpy\_(\allowbreak{}nth(\allowbreak{}1, tpc\_)) in nth(\allowbreak{}2, tpc\_))\}
\end{quote}
characterized by \texttt{tprod-box-mem}  \textit{(proven)}:
\begin{quote}\ttfamily\footnotesize\raggedright\noindent
forall(\allowbreak{}[ix, tf, bb, z], z in tprod-box(\allowbreak{}ix, tf, bb) iff z in tprod-pts(\allowbreak{}ix, tf) and forall(\allowbreak{}[tpc\_ in bb], z(\allowbreak{}nth(\allowbreak{}1, tpc\_)) in nth(\allowbreak{}2, tpc\_)))
\end{quote}
also: \texttt{tprod-box-anti} \texttt{tprod-box-open} \texttt{tprod-box-unfold} \\
mentioned by 16 result(s)

\smallskip

\noindent\textbf{\texttt{tprod-constr(ix, tf)}}\quad\textit{functoid}\index{tprod-constr(ix, tf)@\texttt{tprod-constr(ix, tf)}}\\
reads: the subbasic constraints of the product of \$2 \\
definition:
\begin{quote}\ttfamily\footnotesize\raggedright\noindent
\{tpb\_ in cartesian(\allowbreak{}ix, big-union(\allowbreak{}tpj\_, ix, opens(\allowbreak{}tf(\allowbreak{}tpj\_)))): nth(\allowbreak{}2, tpb\_) in opens(\allowbreak{}tf(\allowbreak{}nth(\allowbreak{}1, tpb\_)))\}
\end{quote}
characterized by \texttt{tprod-constr-mem}  \textit{(proven)}:
\begin{quote}\ttfamily\footnotesize\raggedright\noindent
forall(\allowbreak{}[ix, tf, z], z in tprod-constr(\allowbreak{}ix, tf) iff z in cartesian(\allowbreak{}ix, big-union(\allowbreak{}tpj\_, ix, opens(\allowbreak{}tf(\allowbreak{}tpj\_)))) and nth(\allowbreak{}2, z) in opens(\allowbreak{}tf(\allowbreak{}nth(\allowbreak{}1, z))))
\end{quote}
also: \texttt{tprod-constr-unfold} \\
mentioned by 13 result(s)

\smallskip

\noindent\textbf{\texttt{tprod-pts(ix, tf)}}\quad\textit{functoid}\index{tprod-pts(ix, tf)@\texttt{tprod-pts(ix, tf)}}\\
reads: the product of the points of the family \$2 \\
definition:
\begin{quote}\ttfamily\footnotesize\raggedright\noindent
\{tpx\_ in fun(\allowbreak{}ix, big-union(\allowbreak{}tpj\_, ix, pts(\allowbreak{}tf(\allowbreak{}tpj\_)))): forall(\allowbreak{}[tpi\_ in ix], tpx\_(\allowbreak{}tpi\_) in pts(\allowbreak{}tf(\allowbreak{}tpi\_)))\}
\end{quote}
characterized by \texttt{tprod-pts-mem}  \textit{(proven)}:
\begin{quote}\ttfamily\footnotesize\raggedright\noindent
forall(\allowbreak{}[ix, tf, z], z in tprod-pts(\allowbreak{}ix, tf) iff z in fun(\allowbreak{}ix, big-union(\allowbreak{}tpj\_, ix, pts(\allowbreak{}tf(\allowbreak{}tpj\_)))) and forall(\allowbreak{}[tpi\_ in ix], z(\allowbreak{}tpi\_) in pts(\allowbreak{}tf(\allowbreak{}tpi\_))))
\end{quote}
also: \texttt{tprod-pts-unfold} \\
mentioned by 47 result(s)

\smallskip

\noindent\textbf{\texttt{tproduct-top(ix, tf)}}\quad\textit{functoid}\index{tproduct-top(ix, tf)@\texttt{tproduct-top(ix, tf)}}\\
reads: the product topology of the family \$2 \\
definition:
\begin{quote}\ttfamily\footnotesize\raggedright\noindent
[tprod-pts(\allowbreak{}ix, tf), \{tpu\_ in power(\allowbreak{}tprod-pts(\allowbreak{}ix, tf)): is-tprod-open(\allowbreak{}ix, tf, tpu\_)\}]
\end{quote}
mentioned by 36 result(s)

\smallskip

\noindent\textbf{\texttt{tproj(ix, tf, i)}}\quad\textit{functoid}\index{tproj(ix, tf, i)@\texttt{tproj(ix, tf, i)}}\\
reads: the projection onto the factor \$3 \\
definition:
\begin{quote}\ttfamily\footnotesize\raggedright\noindent
vnb-lambda(\allowbreak{}tpz\_, tprod-pts(\allowbreak{}ix, tf), tpz\_(\allowbreak{}i))
\end{quote}
characterized by \texttt{tproj-value}  \textit{(proven)}:
\begin{quote}\ttfamily\footnotesize\raggedright\noindent
forall(\allowbreak{}[ix, tf, i, y], is-top-family(\allowbreak{}ix, tf) implies i in ix implies y in tprod-pts(\allowbreak{}ix, tf) implies (\allowbreak{}tproj(\allowbreak{}ix, tf, i))(\allowbreak{}y) = y(\allowbreak{}i))
\end{quote}
also: \texttt{tproj-continuous} \\
mentioned by 9 result(s)

\smallskip

\noindent\textbf{\texttt{trace(pgam, a, b)}}\quad\textit{functoid}\index{trace(pgam, a, b)@\texttt{trace(pgam, a, b)}}\\
reads: the trace of the path \$1 on [\$2, \$3] \\
definition:
\begin{quote}\ttfamily\footnotesize\raggedright\noindent
image(\allowbreak{}pgam, ccint(\allowbreak{}a, b))
\end{quote}
characterized by \texttt{juxta-trace}  \textit{(proven)}:
\begin{quote}\ttfamily\footnotesize\raggedright\noindent
forall(\allowbreak{}[p, q, a in rr, b in rr, c in rr], a <= b implies b <= c implies p in fun(\allowbreak{}ccint(\allowbreak{}a, b), cc) implies q in fun(\allowbreak{}ccint(\allowbreak{}b, c), cc) implies p(\allowbreak{}b) = q(\allowbreak{}b) implies trace(\allowbreak{}juxta(\allowbreak{}p, q, a, b, c), a, c) = union(\allowbreak{}trace(\allowbreak{}p, a, b), trace(\allowbreak{}q, b, c)))
\end{quote}
also: \texttt{trace-is-set} \texttt{trace-unfold} \texttt{poly-trace-in} \texttt{seg-trace-elim} \texttt{trace-opposite} \texttt{trace-value-in} \texttt{tri-road-trace} \texttt{conv3-seg-trace} ... (29 in all) \\
mentioned by 232 result(s)

\smallskip

\noindent\textbf{\texttt{trace-class-ops(kqv\_)}}\quad\textit{functoid}\index{trace-class-ops(kqv_)@\texttt{trace-class-ops(kqv\_)}}\\
reads: the trace-class operators on \$1 \\
definition:
\begin{quote}\ttfamily\footnotesize\raggedright\noindent
\{kqt\_ in bounded-ops(\allowbreak{}kqv\_, kqv\_): is-trace-class(\allowbreak{}kqv\_, kqt\_)\}
\end{quote}
characterized by \texttt{trace-class-ops-unfold}  \textit{(proven)}:
\begin{quote}\ttfamily\footnotesize\raggedright\noindent
forall(\allowbreak{}[kqv\_], trace-class-ops(\allowbreak{}kqv\_) == \{kqt\_ in bounded-ops(\allowbreak{}kqv\_, kqv\_): is-trace-class(\allowbreak{}kqv\_, kqt\_)\})
\end{quote}
also: \texttt{trace-class-ops-membership} \\
mentioned by 4 result(s)

\smallskip

\noindent\textbf{\texttt{translate(es, c)}}\quad\textit{functoid}\index{translate(es, c)@\texttt{translate(es, c)}}\\
reads: the translate of the set of reals \$1 by \$2 \\
definition:
\begin{quote}\ttfamily\footnotesize\raggedright\noindent
\{ltz\_ in rr: ltz\_ - c in es\}
\end{quote}
characterized by \texttt{translate-ccint}  \textit{(proven)}:
\begin{quote}\ttfamily\footnotesize\raggedright\noindent
forall(\allowbreak{}[a\_, b\_, c\_], a\_ in rr implies b\_ in rr implies c\_ in rr implies translate(\allowbreak{}ccint(\allowbreak{}a\_, b\_), c\_) = ccint(\allowbreak{}a\_ + c\_, b\_ + c\_))
\end{quote}
also: \texttt{translate-unfold} \texttt{translate-neg-back} \texttt{translate-subset-rr} \texttt{translate-difference} \texttt{translate-membership} \texttt{translate-in-power-rr} \texttt{translate-intersection} \texttt{lebesgue-measurable-translate} ... (9 in all) \\
mentioned by 42 result(s)

\smallskip

\noindent\textbf{\texttt{transpose(a, m, n)}}\quad\textit{functoid}\index{transpose(a, m, n)@\texttt{transpose(a, m, n)}}\\
reads: \$1\textasciicircum{}t \\
definition:
\begin{quote}\ttfamily\footnotesize\raggedright\noindent
matof(\allowbreak{}n, m, vnb-lambda(\allowbreak{}[tri\_, trj\_], cartesian(\allowbreak{}interval(\allowbreak{}1, n), interval(\allowbreak{}1, m)), entry(\allowbreak{}a, trj\_, tri\_)))
\end{quote}
characterized by \texttt{det-transpose}  \textit{(proven)}:
\begin{quote}\ttfamily\footnotesize\raggedright\noindent
forall(\allowbreak{}[r, n, a], is-commutative-ring(\allowbreak{}r) implies n in nn implies a in mat(\allowbreak{}n, n, carr(\allowbreak{}r)) implies det(\allowbreak{}r, n, transpose(\allowbreak{}a, n, n)) = det(\allowbreak{}r, n, a))
\end{quote}
also: \texttt{transpose-entry} \texttt{transpose-in-mat} \texttt{transpose-matadd} \texttt{transpose-matmul} \texttt{transpose-identity} \texttt{transpose-matscale} \texttt{dtt-minor-transpose} \texttt{transpose-transpose} \\
mentioned by 40 result(s)

\smallskip

\noindent\textbf{\texttt{tri-int(sgf\_, sga\_, sgb\_, sgc\_)}}\quad\textit{functoid}\index{tri-int(sgf_, sga_, sgb_, sgc_)@\texttt{tri-int(sgf\_, sga\_, sgb\_, sgc\_)}}\\
reads: the integral of \$1 around the triangle \$2 \$3 \$4 \\
definition:
\begin{quote}\ttfamily\footnotesize\raggedright\noindent
(\allowbreak{}seg-int(\allowbreak{}sgf\_, sga\_, sgb\_) + seg-int(\allowbreak{}sgf\_, sgb\_, sgc\_)) + seg-int(\allowbreak{}sgf\_, sgc\_, sga\_)
\end{quote}
characterized by \texttt{tri-int-fan}  \textit{(proven)}:
\begin{quote}\ttfamily\footnotesize\raggedright\noindent
forall(\allowbreak{}[u, f], u subset cc implies f in fun(\allowbreak{}u, cc) implies forall(\allowbreak{}[rgy\_ in u], is-continuous-at(\allowbreak{}subspace-ms(\allowbreak{}nf-metric-space(\allowbreak{}cc-normed-field), u), nf-metric-space(\allowbreak{}cc-normed-field), f, rgy\_)) implies forall(\allowbreak{}[goxa\_ in cc, goxb\_ in cc, goxc\_ in cc], conv3(\allowbreak{}goxa\_, goxb\_, goxc\_) subset u implies forall(\allowbreak{}[goxx\_ in conv3(\allowbreak{}goxa\_, goxb\_, goxc\_)], tri-int(\allowbreak{}f, goxa\_, goxb\_, goxc\_) = (\allowbreak{}tri-int(\allowbreak{}f, goxa\_, goxb\_, goxx\_) + tri-int(\allowbreak{}f, goxb\_, goxc\_, goxx\_)) + tri-int(\allowbreak{}f, goxc\_, goxa\_, goxx\_))))
\end{quote}
also: \texttt{tri-int-sum} \texttt{tri-int-in-cc} \texttt{tri-int-unfold} \texttt{tri-int-restrict} \texttt{tri-int-abs-bound} \texttt{tri-int-collinear} \texttt{tri-int-line-ball} \texttt{tri-int-half-plane} ... (17 in all) \\
mentioned by 52 result(s)

\smallskip

\noindent\textbf{\texttt{tri-road(tra\_, trb\_, trc\_)}}\quad\textit{functoid}\index{tri-road(tra_, trb_, trc_)@\texttt{tri-road(tra\_, trb\_, trc\_)}}\\
reads: the triangular road through \$1, \$2, \$3 \\
definition:
\begin{quote}\ttfamily\footnotesize\raggedright\noindent
juxta(\allowbreak{}juxta(\allowbreak{}seg-path(\allowbreak{}tra\_, trb\_), vnb-lambda(\allowbreak{}trs\_, ccint(\allowbreak{}1, 2), (\allowbreak{}seg-path(\allowbreak{}trb\_, trc\_))(\allowbreak{}trs\_ - 1)), 0, 1, 2), vnb-lambda(\allowbreak{}trs\_, ccint(\allowbreak{}2, 3), (\allowbreak{}seg-path(\allowbreak{}trc\_, tra\_))(\allowbreak{}trs\_ - 2)), 0, 2, 3)
\end{quote}
characterized by \texttt{tri-road-trace}  \textit{(proven)}:
\begin{quote}\ttfamily\footnotesize\raggedright\noindent
forall(\allowbreak{}[a in cc, b in cc, c in cc], trace(\allowbreak{}tri-road(\allowbreak{}a, b, c), 0, 3) = union(\allowbreak{}union(\allowbreak{}trace(\allowbreak{}seg-path(\allowbreak{}a, b), 0, 1), trace(\allowbreak{}seg-path(\allowbreak{}b, c), 0, 1)), trace(\allowbreak{}seg-path(\allowbreak{}c, a), 0, 1)))
\end{quote}
also: \texttt{tri-road-unfold} \texttt{tri-road-is-road} \\
mentioned by 13 result(s)

\smallskip

\noindent\textbf{\texttt{tri-road-deriv(tra\_, trb\_, trc\_)}}\quad\textit{functoid}\index{tri-road-deriv(tra_, trb_, trc_)@\texttt{tri-road-deriv(tra\_, trb\_, trc\_)}}\\
reads: the derivative of the triangular road through \$1, \$2, \$3 \\
definition:
\begin{quote}\ttfamily\footnotesize\raggedright\noindent
juxta(\allowbreak{}juxta(\allowbreak{}seg-deriv(\allowbreak{}tra\_, trb\_), vnb-lambda(\allowbreak{}trs\_, ccint(\allowbreak{}1, 2), (\allowbreak{}seg-deriv(\allowbreak{}trb\_, trc\_))(\allowbreak{}trs\_ - 1)), 0, 1, 2), vnb-lambda(\allowbreak{}trs\_, ccint(\allowbreak{}2, 3), (\allowbreak{}seg-deriv(\allowbreak{}trc\_, tra\_))(\allowbreak{}trs\_ - 2)), 0, 2, 3)
\end{quote}
characterized by \texttt{tri-road-deriv-unfold}  \textit{(proven)}:
\begin{quote}\ttfamily\footnotesize\raggedright\noindent
forall(\allowbreak{}[tra\_, trb\_, trc\_], tri-road-deriv(\allowbreak{}tra\_, trb\_, trc\_) == juxta(\allowbreak{}juxta(\allowbreak{}seg-deriv(\allowbreak{}tra\_, trb\_), vnb-lambda(\allowbreak{}trs\_, ccint(\allowbreak{}1, 2), (\allowbreak{}seg-deriv(\allowbreak{}trb\_, trc\_))(\allowbreak{}trs\_ - 1)), 0, 1, 2), vnb-lambda(\allowbreak{}trs\_, ccint(\allowbreak{}2, 3), (\allowbreak{}seg-deriv(\allowbreak{}trc\_, tra\_))(\allowbreak{}trs\_ - 2)), 0, 2, 3))
\end{quote}
mentioned by 13 result(s)

\smallskip

\noindent\textbf{\texttt{trinum}}\quad\textit{functoid}\index{trinum@\texttt{trinum}}\\
reads: the \$1-th triangular number \\
defined by \texttt{trinum-zero} \texttt{trinum-succ} \\
characterized by \texttt{trinum-succ}  \textit{(definitional)}:
\begin{quote}\ttfamily\footnotesize\raggedright\noindent
forall(\allowbreak{}[n\_ in nn], trinum(\allowbreak{}succ(\allowbreak{}n\_)) == trinum(\allowbreak{}n\_) + succ(\allowbreak{}n\_))
\end{quote}
also: \texttt{trinum-zero} \texttt{trinum-ge} \texttt{trinum-mono} \texttt{trinum-type} \\
mentioned by 10 result(s)

\smallskip

\noindent\textbf{\texttt{triv-abs(a5k\_, a5x\_)}}\quad\textit{functoid}\index{triv-abs(a5k_, a5x_)@\texttt{triv-abs(a5k\_, a5x\_)}}\\
reads: the trivial absolute value of \$2 in \$1 \\
definition:
\begin{quote}\ttfamily\footnotesize\raggedright\noindent
if(\allowbreak{}a5x\_ = zero(\allowbreak{}a5k\_), 0, 1)
\end{quote}
characterized by \texttt{triv-abs-add}  \textit{(proven)}:
\begin{quote}\ttfamily\footnotesize\raggedright\noindent
forall(\allowbreak{}[kq\_], is-commutative-ring(\allowbreak{}kq\_) implies forall(\allowbreak{}[a5x in carr(\allowbreak{}kq\_), a5y in carr(\allowbreak{}kq\_)], triv-abs(\allowbreak{}kq\_, (\allowbreak{}add(\allowbreak{}kq\_))(\allowbreak{}a5x, a5y)) <= triv-abs(\allowbreak{}kq\_, a5x) + triv-abs(\allowbreak{}kq\_, a5y)))
\end{quote}
also: \texttt{triv-abs-mul} \texttt{triv-abs-zero} \texttt{triv-abs-range} \texttt{triv-abs-nonzero} \texttt{triv-abs-zero-iff} \\
mentioned by 12 result(s)

\smallskip

\noindent\textbf{\texttt{triv-group}}\quad\textit{constant}\index{triv-group@\texttt{triv-group}}\\
reads: the trivial group \\
defined by \texttt{triv-group-def}

\smallskip

\noindent\textbf{\texttt{triv-nf(a5k\_)}}\quad\textit{functoid}\index{triv-nf(a5k_)@\texttt{triv-nf(a5k\_)}}\\
reads: the field \$1 with the trivial absolute value \\
definition:
\begin{quote}\ttfamily\footnotesize\raggedright\noindent
[carr(\allowbreak{}a5k\_), add(\allowbreak{}a5k\_), mul(\allowbreak{}a5k\_), neg(\allowbreak{}a5k\_), zero(\allowbreak{}a5k\_), one(\allowbreak{}a5k\_), vnb-lambda(\allowbreak{}a5x\_, carr(\allowbreak{}a5k\_), triv-abs(\allowbreak{}a5k\_, a5x\_))]
\end{quote}
characterized by \texttt{triv-nf-add}  \textit{(proven)}:
\begin{quote}\ttfamily\footnotesize\raggedright\noindent
forall(\allowbreak{}[kq\_], is-ring(\allowbreak{}kq\_) implies add(\allowbreak{}triv-nf(\allowbreak{}kq\_)) == add(\allowbreak{}kq\_))
\end{quote}
also: \texttt{triv-nf-mul} \texttt{triv-nf-neg} \texttt{triv-nf-one} \texttt{triv-nf-carr} \texttt{triv-nf-fnrm} \texttt{triv-nf-zero} \texttt{triv-nf-add-character} \texttt{triv-nf-is-normed-field} \\
mentioned by 20 result(s)

\smallskip

\noindent\textbf{\texttt{true-closed(md7v\_)}}\quad\textit{predicate}\index{true-closed(md7v_)@\texttt{true-closed(md7v\_)}}\\
reads: \$1 is closed under steps of truth \\
defined by \texttt{true-closed} \\
characterized by \texttt{true-closed}  \textit{(definitional)}:
\begin{quote}\ttfamily\footnotesize\raggedright\noindent
forall(\allowbreak{}[md7v\_], true-closed(\allowbreak{}md7v\_) iff forall(\allowbreak{}[md7m\_, md7n\_], md7m\_ in md7v\_ implies md7n\_ in set implies true-step(\allowbreak{}md7m\_, md7n\_) implies md7n\_ in md7v\_))
\end{quote}
mentioned by 5 result(s)

\smallskip

\noindent\textbf{\texttt{true-reach(md7j\_)}}\quad\textit{functoid}\index{true-reach(md7j_)@\texttt{true-reach(md7j\_)}}\\
reads: the sentences reached from \$1 by steps of truth \\
definition:
\begin{quote}\ttfamily\footnotesize\raggedright\noindent
\{md7zz\_ in expr-set: forall(\allowbreak{}[md7ww\_ in power(\allowbreak{}expr-set)], md7j\_ in md7ww\_ implies true-closed(\allowbreak{}md7ww\_) implies md7zz\_ in md7ww\_)\}
\end{quote}
characterized by \texttt{true-reach-least}  \textit{(proven)}:
\begin{quote}\ttfamily\footnotesize\raggedright\noindent
forall(\allowbreak{}[e, md7xs\_ in power(\allowbreak{}expr-set)], e in md7xs\_ implies true-closed(\allowbreak{}md7xs\_) implies forall(\allowbreak{}[x in true-reach(\allowbreak{}e)], x in md7xs\_))
\end{quote}
also: \texttt{true-reach-unfold} \\
mentioned by 12 result(s)

\smallskip

\noindent\textbf{\texttt{true-step(md7f\_, md7g\_)}}\quad\textit{predicate}\index{true-step(md7f_, md7g_)@\texttt{true-step(md7f\_, md7g\_)}}\\
reads: \$2 is one step of truth below \$1 \\
defined by \texttt{true-step} \\
characterized by \texttt{true-step}  \textit{(definitional)}:
\begin{quote}\ttfamily\footnotesize\raggedright\noindent
forall(\allowbreak{}[md7f\_, md7g\_], true-step(\allowbreak{}md7f\_, md7g\_) iff forsome(\allowbreak{}[md7c\_, md7d\_], md7f\_ = enc-and(\allowbreak{}md7c\_, md7d\_) and (\allowbreak{}md7g\_ = md7c\_ or md7g\_ = md7d\_)) or forsome(\allowbreak{}[md7y\_, md7o\_, md7r\_], md7f\_ = enc-in(\allowbreak{}md7y\_, enc-sep(\allowbreak{}md7o\_, md7r\_)) and nbg-ok(\allowbreak{}md7r\_) and md7g\_ = expr-open(\allowbreak{}md7r\_, 0, md7y\_)))
\end{quote}
mentioned by 4 result(s)

\smallskip

\noindent\textbf{\texttt{trunc-metric(s)}}\quad\textit{functoid}\index{trunc-metric(s)@\texttt{trunc-metric(s)}}\\
reads: the truncated metric space of \$1 \\
definition:
\begin{quote}\ttfamily\footnotesize\raggedright\noindent
[pts(\allowbreak{}s), vnb-lambda(\allowbreak{}[u, v], cartesian(\allowbreak{}pts(\allowbreak{}s), pts(\allowbreak{}s)), min(\allowbreak{}1, (\allowbreak{}dist(\allowbreak{}s))(\allowbreak{}u, v)))]
\end{quote}
characterized by \texttt{trunc-metric-dist}  \textit{(proven)}:
\begin{quote}\ttfamily\footnotesize\raggedright\noindent
forall(\allowbreak{}[s, u\_, v\_], u\_ in pts(\allowbreak{}s) implies v\_ in pts(\allowbreak{}s) implies (\allowbreak{}dist(\allowbreak{}trunc-metric(\allowbreak{}s)))(\allowbreak{}u\_, v\_) == min(\allowbreak{}1, (\allowbreak{}dist(\allowbreak{}s))(\allowbreak{}u\_, v\_)))
\end{quote}
also: \texttt{trunc-metric-bounded} \texttt{trunc-metric@preimage} \texttt{trunc-metric@sup-dist} \texttt{trunc-metric-converges-bwd} \texttt{trunc-metric-converges-fwd} \texttt{trunc-metric-converges-iff} \texttt{trunc-metric-is-metric-space} \texttt{trunc-metric-id-uniformly-continuous} ... (9 in all) \\
mentioned by 14 result(s)

\smallskip

\noindent\textbf{\texttt{tuples}}\quad\textit{operator}\index{tuples@\texttt{tuples}}\\
a term-forming head \\
characterized by \texttt{cons-in-tuples}  \textit{(definitional)}:
\begin{quote}\ttfamily\footnotesize\raggedright\noindent
forall(\allowbreak{}[a, x, l], x in a and l in tuples(\allowbreak{}a) implies cons(\allowbreak{}x, l) in tuples(\allowbreak{}a))
\end{quote}
also: \texttt{empty-in-tuples} \texttt{tuples-sethood} \texttt{tuples-mono} \texttt{tuples-induction} \texttt{lst-in-tuples-set} \texttt{encr-in-tuples-set} \\
mentioned by 173 result(s)

\smallskip

\noindent\textbf{\texttt{type-term}}\quad\textit{tactic}\index{type-term@\texttt{type-term}}\\
\texttt{(type-term)} -- Close a typing goal (IN t C).  In a number class (NN ZZ QQ RR CC) t may be built from context-typed atoms (or atoms typed in a subclass: NN in ZZ, ZZ in RR, ...), numerals, + * -, and powers, typed by the class's closure laws, bottom-up.  In ANY class C, t may be an atom typed in a subclass, or an application f(a) with f in fun(A, C) in context and a typable in A, recursively: f(3 * x + 1) in B from f in fun(zz, B) and x in zz.  Declines -- with a warning, and nothing changed -- when there is no plan. \\
declared in \texttt{tactics-help.scm}

\smallskip


\section*{U}
\addcontentsline{toc}{section}{U}

\noindent\textbf{\texttt{uba-as-cba(s)}}\quad\textit{view}\index{uba-as-cba(s)@\texttt{uba-as-cba(s)}}\\
an unital-banach-algebra viewed as a complex-banach-algebra \\
characterized by \texttt{uba-as-cba-is-complex-banach-algebra}  \textit{(definitional)}:
\begin{quote}\ttfamily\footnotesize\raggedright\noindent
forall(\allowbreak{}[r], is-unital-banach-algebra(\allowbreak{}r) implies is-complex-banach-algebra(\allowbreak{}uba-as-cba(\allowbreak{}r)))
\end{quote}
also: \texttt{is-complex-banach-algebra-is-set-uba-as-cba} \texttt{uba-as-cba-is-complex-banach-algebra-uiba-as-uba} \texttt{is-complex-banach-algebra-is-set-uba-as-cba-uiba-as-uba} \texttt{uba-as-cba-functorial} \\
mentioned by 23 result(s)

\smallskip

\noindent\textbf{\texttt{uba-as-module(s)}}\quad\textit{view}\index{uba-as-module(s)@\texttt{uba-as-module(s)}}\\
an unital-banach-algebra viewed as a module \\
characterized by \texttt{mvag-id-uba-as-module}  \textit{(definitional)}:
\begin{quote}\ttfamily\footnotesize\raggedright\noindent
forall(\allowbreak{}[r], is-unital-banach-algebra(\allowbreak{}r) implies iden(\allowbreak{}module-vector-ag(\allowbreak{}uba-as-module(\allowbreak{}r))) = vzero(\allowbreak{}r))
\end{quote}
also: \texttt{mvag-op-uba-as-module} \texttt{mvag-carr-uba-as-module} \texttt{qmod-comm-uba-as-module} \texttt{uba-as-module-is-module} \texttt{qmod-assoc-uba-as-module} \texttt{spans-span-uba-as-module} \texttt{hbs-lf-zero-uba-as-module} \texttt{hbs-sl-zero-uba-as-module} ... (310 in all) \\
mentioned by 982 result(s)

\smallskip

\noindent\textbf{\texttt{uba-as-normed-ag(s)}}\quad\textit{view}\index{uba-as-normed-ag(s)@\texttt{uba-as-normed-ag(s)}}\\
an unital-banach-algebra viewed as a normed-ag \\
characterized by \texttt{uba-as-normed-ag-is-normed-ag}  \textit{(definitional)}:
\begin{quote}\ttfamily\footnotesize\raggedright\noindent
forall(\allowbreak{}[r], is-unital-banach-algebra(\allowbreak{}r) implies is-normed-ag(\allowbreak{}uba-as-normed-ag(\allowbreak{}r)))
\end{quote}
also: \texttt{rb2-dist-union-uba-as-normed-ag} \texttt{sums-to-unique-uba-as-normed-ag} \texttt{rb2-dist-subset-uba-as-normed-ag} \texttt{hom-normed-ag-id-uba-as-normed-ag} \texttt{rb2-nrm-finsum-le-uba-as-normed-ag} \texttt{nag-metric-carrier-uba-as-normed-ag} \texttt{is-normed-ag-is-set-uba-as-normed-ag} \texttt{nag-metric-distance-uba-as-normed-ag} ... (28 in all) \\
mentioned by 61 result(s)

\smallskip

\noindent\textbf{\texttt{uba-as-ring(s)}}\quad\textit{view}\index{uba-as-ring(s)@\texttt{uba-as-ring(s)}}\\
an unital-banach-algebra viewed as a ring \\
characterized by \texttt{ras-id-uba-as-ring}  \textit{(definitional)}:
\begin{quote}\ttfamily\footnotesize\raggedright\noindent
forall(\allowbreak{}[r], is-unital-banach-algebra(\allowbreak{}r) implies iden(\allowbreak{}ring-additive-ag(\allowbreak{}uba-as-ring(\allowbreak{}r))) = vzero(\allowbreak{}r))
\end{quote}
also: \texttt{ras-op-uba-as-ring} \texttt{rmm-op-uba-as-ring} \texttt{ras-inv-uba-as-ring} \texttt{uba-as-ring-is-ring} \texttt{deg-drop-uba-as-ring} \texttt{raag-opr-uba-as-ring} \texttt{ras-carr-uba-as-ring} \texttt{rmm-carr-uba-as-ring} ... (296 in all) \\
mentioned by 1131 result(s)

\smallskip

\noindent\textbf{\texttt{ue}}\quad\textit{tactic}\index{ue@\texttt{ue}}\\
\texttt{(ue hyp)} -- Union elim: split a cited (IN x (UNION ...)) membership assumption into one subgoal per set -- union-side case analysis, the dual of `ui'. \\
declared in \texttt{tactics-help.scm}

\smallskip

\noindent\textbf{\texttt{ui}}\quad\textit{tactic}\index{ui@\texttt{ui}}\\
\texttt{(ui k)} -- Union intro: reduce a goal (IN x (UNION ...)) to membership of x in the k-th set (1-based). \\
declared in \texttt{tactics-help.scm}

\smallskip

\noindent\textbf{\texttt{uiba-as-nuiba(s)}}\quad\textit{view}\index{uiba-as-nuiba(s)@\texttt{uiba-as-nuiba(s)}}\\
an unital-involutive-banach-algebra viewed as a nu-involutive-banach-algebra \\
characterized by \texttt{uiba-as-nuiba-is-nu-involutive-banach-algebra}  \textit{(definitional)}:
\begin{quote}\ttfamily\footnotesize\raggedright\noindent
forall(\allowbreak{}[r], is-unital-involutive-banach-algebra(\allowbreak{}r) implies is-nu-involutive-banach-algebra(\allowbreak{}uiba-as-nuiba(\allowbreak{}r)))
\end{quote}
also: \texttt{is-nu-involutive-banach-algebra-is-set-uiba-as-nuiba} \texttt{nuiba-as-cba-is-complex-banach-algebra-uiba-as-nuiba} \texttt{is-complex-banach-algebra-is-set-nuiba-as-cba-uiba-as-nuiba} \texttt{uiba-as-nuiba-functorial} \\
mentioned by 5 result(s)

\smallskip

\noindent\textbf{\texttt{uiba-as-uba(s)}}\quad\textit{view}\index{uiba-as-uba(s)@\texttt{uiba-as-uba(s)}}\\
an unital-involutive-banach-algebra viewed as an unital-banach-algebra \\
characterized by \texttt{ras-id-uba-as-ring-uiba-as-uba}  \textit{(definitional)}:
\begin{quote}\ttfamily\footnotesize\raggedright\noindent
forall(\allowbreak{}[r], is-unital-involutive-banach-algebra(\allowbreak{}r) implies iden(\allowbreak{}ring-additive-ag(\allowbreak{}uba-as-ring(\allowbreak{}uiba-as-uba(\allowbreak{}r)))) = vzero(\allowbreak{}r))
\end{quote}
also: \texttt{ras-op-uba-as-ring-uiba-as-uba} \texttt{rmm-op-uba-as-ring-uiba-as-uba} \texttt{ras-inv-uba-as-ring-uiba-as-uba} \texttt{uba-as-ring-is-ring-uiba-as-uba} \texttt{deg-drop-uba-as-ring-uiba-as-uba} \texttt{raag-opr-uba-as-ring-uiba-as-uba} \texttt{ras-carr-uba-as-ring-uiba-as-uba} \texttt{rmm-carr-uba-as-ring-uiba-as-uba} ... (321 in all) \\
mentioned by 1395 result(s)

\smallskip

\noindent\textbf{\texttt{ultra-dual(otv\_)}}\quad\textit{functoid}\index{ultra-dual(otv_)@\texttt{ultra-dual(otv\_)}}\\
reads: the ultrastrongly continuous linear functionals on the bounded operators on \$1 \\
definition:
\begin{quote}\ttfamily\footnotesize\raggedright\noindent
op-dual(\allowbreak{}otv\_, ust(\allowbreak{}otv\_))
\end{quote}

\smallskip

\noindent\textbf{\texttt{ultra-seq-ms(usx\_, usa\_)}}\quad\textit{functoid}\index{ultra-seq-ms(usx_, usa_)@\texttt{ultra-seq-ms(usx\_, usa\_)}}\\
reads: the sequences in \$1 with the \$2-weighted prefix ultrametric \\
definition:
\begin{quote}\ttfamily\footnotesize\raggedright\noindent
[fun(\allowbreak{}nn, usx\_), vnb-lambda(\allowbreak{}[usu\_, usv\_], cartesian(\allowbreak{}fun(\allowbreak{}nn, usx\_), fun(\allowbreak{}nn, usx\_)), prefix-dist(\allowbreak{}usa\_, usu\_, usv\_))]
\end{quote}
mentioned by 12 result(s)

\smallskip

\noindent\textbf{\texttt{ultraweak-dual(otv\_)}}\quad\textit{functoid}\index{ultraweak-dual(otv_)@\texttt{ultraweak-dual(otv\_)}}\\
reads: the ultraweakly continuous linear functionals on the bounded operators on \$1 \\
definition:
\begin{quote}\ttfamily\footnotesize\raggedright\noindent
op-dual(\allowbreak{}otv\_, uwt(\allowbreak{}otv\_))
\end{quote}

\smallskip

\noindent\textbf{\texttt{undo}}\quad\textit{tactic}\index{undo@\texttt{undo}}\\
\texttt{(undo)} -- Alias for (backup-one). \\
declared in \texttt{tactics-help.scm}

\smallskip

\noindent\textbf{\texttt{unif-on-compacta(u, fam, f)}}\quad\textit{predicate}\index{unif-on-compacta(u, fam, f)@\texttt{unif-on-compacta(u, fam, f)}}\\
reads: \$2 converges to \$3 uniformly on the compact subsets of \$1 \\
defined by \texttt{unif-on-compacta} \\
characterized by \texttt{unif-on-compacta}  \textit{(definitional)}:
\begin{quote}\ttfamily\footnotesize\raggedright\noindent
forall(\allowbreak{}[u, fam, f], unif-on-compacta(\allowbreak{}u, fam, f) iff forall(\allowbreak{}[wcq\_], wcq\_ subset u implies is-compact(\allowbreak{}subspace-ms(\allowbreak{}cc-ms, wcq\_)) implies forall(\allowbreak{}[wce\_], pos-rr(\allowbreak{}wce\_) implies forsome(\allowbreak{}[wcn\_ in nn], forall(\allowbreak{}[wck\_ in nn], wcn\_ <= wck\_ implies forall(\allowbreak{}[wcy\_ in wcq\_], magnitude(\allowbreak{}(\allowbreak{}fam(\allowbreak{}wck\_))(\allowbreak{}wcy\_) - f(\allowbreak{}wcy\_)) <= wce\_))))))
\end{quote}
also: \texttt{unif-on-compacta-locally} \texttt{cc-unif-on-compacta-pointwise} \texttt{cc-fr-conv-iff-unif-on-compacta} \\
mentioned by 24 result(s)

\smallskip

\noindent\textbf{\texttt{union}}\quad\textit{operator}\index{union@\texttt{union}}\\
a term-forming head \\
characterized by \texttt{rel-big-union-mem-elim}  \textit{(definitional)}:
\begin{quote}\ttfamily\footnotesize\raggedright\noindent
forall(\allowbreak{}[mit\_, mihs\_, mic\_], rel-big-union-mem-elim(\allowbreak{}mit\_, mihs\_, mic\_) iff mic\_ in seq-set and forsome(\allowbreak{}[mig\_, mif\_, mia\_, mib\_, mix\_, mie\_], mic\_ = [mig\_, mif\_] and enc-in(\allowbreak{}mix\_, enc-bu(\allowbreak{}mia\_, mib\_)) in mig\_ and nbg-ok(\allowbreak{}mib\_) and mie\_ in nn and not(\allowbreak{}mie\_ in union(\allowbreak{}ctx-fv(\allowbreak{}mig\_), expr-fv(\allowbreak{}mif\_))) and mihs\_ = [[ctx-add(\allowbreak{}ctx-add(\allowbreak{}difference(\allowbreak{}mig\_, pair(\allowbreak{}enc-in(\allowbreak{}mix\_, enc-bu(\allowbreak{}mia\_, mib\_)), enc-in(\allowbreak{}mix\_, enc-bu(\allowbreak{}mia\_, mib\_)))), enc-in(\allowbreak{}efvar(\allowbreak{}mie\_), mia\_)), enc-in(\allowbreak{}mix\_, expr-open(\allowbreak{}mib\_, 0, efvar(\allowbreak{}mie\_)))), mif\_]]))
\end{quote}
also: \texttt{rb2-dist-union-cna-as-normed-ag} \texttt{rb2-dist-union-uba-as-normed-ag} \texttt{rb2-dist-union-cnvs-as-normed-ag} \texttt{sum-set-disjoint-union-defined-uba-as-ring} \texttt{rb2-dist-union-uba-as-normed-ag-uiba-as-uba} \texttt{sum-set-disjoint-union-defined-uba-as-ring-uiba-as-uba} \texttt{union-decompose} \texttt{union-membership} ... (57 in all) \\
mentioned by 281 result(s)

\smallskip

\noindent\textbf{\texttt{uniquely-originates(ssp\_, ssu\_)}}\quad\textit{predicate}\index{uniquely-originates(ssp_, ssu_)@\texttt{uniquely-originates(ssp\_, ssu\_)}}\\
reads: \$2 originates uniquely in the strand space \$1 \\
defined by \texttt{uniquely-originates} \\
characterized by \texttt{uniquely-originates}  \textit{(definitional)}:
\begin{quote}\ttfamily\footnotesize\raggedright\noindent
forall(\allowbreak{}[ssp\_, ssu\_], uniquely-originates(\allowbreak{}ssp\_, ssu\_) iff forsome(\allowbreak{}[ssn\_], originates(\allowbreak{}ssp\_, ssu\_, ssn\_) and forall(\allowbreak{}[ssm\_], originates(\allowbreak{}ssp\_, ssu\_, ssm\_) implies ssm\_ = ssn\_)))
\end{quote}
also: \texttt{cx-na-uniquely-originates} \texttt{cx-nb-uniquely-originates} \\
mentioned by 32 result(s)

\smallskip

\noindent\textbf{\texttt{unit-act(baa\_)}}\quad\textit{functoid}\index{unit-act(baa_)@\texttt{unit-act(baa\_)}}\\
definition:
\begin{quote}\ttfamily\footnotesize\raggedright\noindent
vnb-lambda(\allowbreak{}[buc\_, bup\_], cartesian(\allowbreak{}cc, unit-vec(\allowbreak{}baa\_)), [buc\_ * nth(\allowbreak{}1, bup\_), (\allowbreak{}act(\allowbreak{}baa\_))(\allowbreak{}buc\_, nth(\allowbreak{}2, bup\_))])
\end{quote}
characterized by \texttt{b1s-uf-unit-act}  \textit{(proven)}:
\begin{quote}\ttfamily\footnotesize\raggedright\noindent
forall(\allowbreak{}[a], is-nu-involutive-banach-algebra(\allowbreak{}a) implies unit-act(\allowbreak{}nuiba-as-cba(\allowbreak{}a)) == unit-act(\allowbreak{}a))
\end{quote}
mentioned by 49 result(s)

\smallskip

\noindent\textbf{\texttt{unit-circle}}\quad\textit{constant}\index{unit-circle@\texttt{unit-circle}}\\
reads: the unit circle \{z : |z| = 1\} \\
defined by \texttt{unit-circle-def}

\smallskip

\noindent\textbf{\texttt{unit-istar(baa\_)}}\quad\textit{functoid}\index{unit-istar(baa_)@\texttt{unit-istar(baa\_)}}\\
definition:
\begin{quote}\ttfamily\footnotesize\raggedright\noindent
vnb-lambda(\allowbreak{}bup\_, unit-vec(\allowbreak{}baa\_), [conjugate(\allowbreak{}nth(\allowbreak{}1, bup\_)), (\allowbreak{}istar(\allowbreak{}baa\_))(\allowbreak{}nth(\allowbreak{}2, bup\_))])
\end{quote}
mentioned by 27 result(s)

\smallskip

\noindent\textbf{\texttt{unit-sphere-cr}}\quad\textit{constant}\index{unit-sphere-cr@\texttt{unit-sphere-cr}}\\
reads: the unit sphere S\textasciicircum{}2 of C x R \\
defined by \texttt{unit-sphere-cr-def}

\smallskip

\noindent\textbf{\texttt{unit-vadd(baa\_)}}\quad\textit{functoid}\index{unit-vadd(baa_)@\texttt{unit-vadd(baa\_)}}\\
definition:
\begin{quote}\ttfamily\footnotesize\raggedright\noindent
vnb-lambda(\allowbreak{}[bup\_, buq\_], cartesian(\allowbreak{}unit-vec(\allowbreak{}baa\_), unit-vec(\allowbreak{}baa\_)), [nth(\allowbreak{}1, bup\_) + nth(\allowbreak{}1, buq\_), (\allowbreak{}vadd(\allowbreak{}baa\_))(\allowbreak{}nth(\allowbreak{}2, bup\_), nth(\allowbreak{}2, buq\_))])
\end{quote}
characterized by \texttt{b1s-uf-unit-vadd}  \textit{(proven)}:
\begin{quote}\ttfamily\footnotesize\raggedright\noindent
forall(\allowbreak{}[a], is-nu-involutive-banach-algebra(\allowbreak{}a) implies unit-vadd(\allowbreak{}nuiba-as-cba(\allowbreak{}a)) == unit-vadd(\allowbreak{}a))
\end{quote}
mentioned by 63 result(s)

\smallskip

\noindent\textbf{\texttt{unit-vec(baa\_)}}\quad\textit{functoid}\index{unit-vec(baa_)@\texttt{unit-vec(baa\_)}}\\
definition:
\begin{quote}\ttfamily\footnotesize\raggedright\noindent
cartesian(\allowbreak{}cc, vec(\allowbreak{}baa\_))
\end{quote}
characterized by \texttt{b1s-uf-unit-vec}  \textit{(proven)}:
\begin{quote}\ttfamily\footnotesize\raggedright\noindent
forall(\allowbreak{}[a], is-nu-involutive-banach-algebra(\allowbreak{}a) implies unit-vec(\allowbreak{}nuiba-as-cba(\allowbreak{}a)) == unit-vec(\allowbreak{}a))
\end{quote}
also: \texttt{unit-vec-unfold} \\
mentioned by 76 result(s)

\smallskip

\noindent\textbf{\texttt{unit-vmul(baa\_)}}\quad\textit{functoid}\index{unit-vmul(baa_)@\texttt{unit-vmul(baa\_)}}\\
definition:
\begin{quote}\ttfamily\footnotesize\raggedright\noindent
vnb-lambda(\allowbreak{}[bup\_, buq\_], cartesian(\allowbreak{}unit-vec(\allowbreak{}baa\_), unit-vec(\allowbreak{}baa\_)), [nth(\allowbreak{}1, bup\_) * nth(\allowbreak{}1, buq\_), (\allowbreak{}vadd(\allowbreak{}baa\_))(\allowbreak{}(\allowbreak{}vadd(\allowbreak{}baa\_))(\allowbreak{}(\allowbreak{}act(\allowbreak{}baa\_))(\allowbreak{}nth(\allowbreak{}1, bup\_), nth(\allowbreak{}2, buq\_)), (\allowbreak{}act(\allowbreak{}baa\_))(\allowbreak{}nth(\allowbreak{}1, buq\_), nth(\allowbreak{}2, bup\_))), (\allowbreak{}vmul(\allowbreak{}baa\_))(\allowbreak{}nth(\allowbreak{}2, bup\_), nth(\allowbreak{}2, buq\_)))])
\end{quote}
characterized by \texttt{b1s-uf-unit-vmul}  \textit{(proven)}:
\begin{quote}\ttfamily\footnotesize\raggedright\noindent
forall(\allowbreak{}[a], is-nu-involutive-banach-algebra(\allowbreak{}a) implies unit-vmul(\allowbreak{}nuiba-as-cba(\allowbreak{}a)) == unit-vmul(\allowbreak{}a))
\end{quote}
mentioned by 51 result(s)

\smallskip

\noindent\textbf{\texttt{unit-vneg(baa\_)}}\quad\textit{functoid}\index{unit-vneg(baa_)@\texttt{unit-vneg(baa\_)}}\\
definition:
\begin{quote}\ttfamily\footnotesize\raggedright\noindent
vnb-lambda(\allowbreak{}bup\_, unit-vec(\allowbreak{}baa\_), [-nth(\allowbreak{}1, bup\_), (\allowbreak{}vneg(\allowbreak{}baa\_))(\allowbreak{}nth(\allowbreak{}2, bup\_))])
\end{quote}
characterized by \texttt{b1s-uf-unit-vneg}  \textit{(proven)}:
\begin{quote}\ttfamily\footnotesize\raggedright\noindent
forall(\allowbreak{}[a], is-nu-involutive-banach-algebra(\allowbreak{}a) implies unit-vneg(\allowbreak{}nuiba-as-cba(\allowbreak{}a)) == unit-vneg(\allowbreak{}a))
\end{quote}
mentioned by 25 result(s)

\smallskip

\noindent\textbf{\texttt{unit-vnrm(baa\_)}}\quad\textit{functoid}\index{unit-vnrm(baa_)@\texttt{unit-vnrm(baa\_)}}\\
definition:
\begin{quote}\ttfamily\footnotesize\raggedright\noindent
vnb-lambda(\allowbreak{}bup\_, unit-vec(\allowbreak{}baa\_), magnitude(\allowbreak{}nth(\allowbreak{}1, bup\_)) + (\allowbreak{}vnrm(\allowbreak{}baa\_))(\allowbreak{}nth(\allowbreak{}2, bup\_)))
\end{quote}
characterized by \texttt{b1s-uf-unit-vnrm}  \textit{(proven)}:
\begin{quote}\ttfamily\footnotesize\raggedright\noindent
forall(\allowbreak{}[a], is-nu-involutive-banach-algebra(\allowbreak{}a) implies unit-vnrm(\allowbreak{}nuiba-as-cba(\allowbreak{}a)) == unit-vnrm(\allowbreak{}a))
\end{quote}
mentioned by 35 result(s)

\smallskip

\noindent\textbf{\texttt{unital-banach-algebra}}\quad\textit{structure}\index{unital-banach-algebra@\texttt{unital-banach-algebra}}\\
predicate \texttt{is-unital-banach-algebra}; slots: scal vec vadd vzero vneg act vnrm vmul vone

\smallskip

\noindent\textbf{\texttt{unital-c-star-algebra}}\quad\textit{refinement}\index{unital-c-star-algebra@\texttt{unital-c-star-algebra}}\\
predicate \texttt{is-unital-c-star-algebra}; same shape as unital-involutive-banach-algebra

\smallskip

\noindent\textbf{\texttt{unital-involutive-banach-algebra}}\quad\textit{structure}\index{unital-involutive-banach-algebra@\texttt{unital-involutive-banach-algebra}}\\
predicate \texttt{is-unital-involutive-banach-algebra}; slots: scal vec vadd vzero vneg act vnrm vmul vone vstar

\smallskip

\noindent\textbf{\texttt{unitization(baa\_)}}\quad\textit{functoid}\index{unitization(baa_)@\texttt{unitization(baa\_)}}\\
reads: the unitization of \$1 \\
definition:
\begin{quote}\ttfamily\footnotesize\raggedright\noindent
[normed-field-as-commutative-ring(\allowbreak{}cc-normed-field), unit-vec(\allowbreak{}baa\_), unit-vadd(\allowbreak{}baa\_), [0, vzero(\allowbreak{}baa\_)], unit-vneg(\allowbreak{}baa\_), unit-act(\allowbreak{}baa\_), unit-vnrm(\allowbreak{}baa\_), unit-vmul(\allowbreak{}baa\_), [1, vzero(\allowbreak{}baa\_)]]
\end{quote}
characterized by \texttt{unitization-unfold}  \textit{(proven)}:
\begin{quote}\ttfamily\footnotesize\raggedright\noindent
forall(\allowbreak{}[a], unitization(\allowbreak{}a) == [normed-field-as-commutative-ring(\allowbreak{}cc-normed-field), unit-vec(\allowbreak{}a), unit-vadd(\allowbreak{}a), [0, vzero(\allowbreak{}a)], unit-vneg(\allowbreak{}a), unit-act(\allowbreak{}a), unit-vnrm(\allowbreak{}a), unit-vmul(\allowbreak{}a), [1, vzero(\allowbreak{}a)]])
\end{quote}
also: \texttt{unitization-is-unital-banach-algebra} \\
mentioned by 21 result(s)

\smallskip

\noindent\textbf{\texttt{unitization-na(baa\_)}}\quad\textit{functoid}\index{unitization-na(baa_)@\texttt{unitization-na(baa\_)}}\\
reads: the unitization of \$1 as a normed algebra \\
definition:
\begin{quote}\ttfamily\footnotesize\raggedright\noindent
[normed-field-as-commutative-ring(\allowbreak{}cc-normed-field), unit-vec(\allowbreak{}baa\_), unit-vadd(\allowbreak{}baa\_), [0, vzero(\allowbreak{}baa\_)], unit-vneg(\allowbreak{}baa\_), unit-act(\allowbreak{}baa\_), unit-vnrm(\allowbreak{}baa\_), unit-vmul(\allowbreak{}baa\_)]
\end{quote}
characterized by \texttt{unitization-na-unfold}  \textit{(proven)}:
\begin{quote}\ttfamily\footnotesize\raggedright\noindent
forall(\allowbreak{}[a], unitization-na(\allowbreak{}a) == [normed-field-as-commutative-ring(\allowbreak{}cc-normed-field), unit-vec(\allowbreak{}a), unit-vadd(\allowbreak{}a), [0, vzero(\allowbreak{}a)], unit-vneg(\allowbreak{}a), unit-act(\allowbreak{}a), unit-vnrm(\allowbreak{}a), unit-vmul(\allowbreak{}a)])
\end{quote}
mentioned by 27 result(s)

\smallskip

\noindent\textbf{\texttt{unitization-star(baa\_)}}\quad\textit{functoid}\index{unitization-star(baa_)@\texttt{unitization-star(baa\_)}}\\
reads: the involutive unitization of \$1 \\
definition:
\begin{quote}\ttfamily\footnotesize\raggedright\noindent
[normed-field-as-commutative-ring(\allowbreak{}cc-normed-field), unit-vec(\allowbreak{}baa\_), unit-vadd(\allowbreak{}baa\_), [0, vzero(\allowbreak{}baa\_)], unit-vneg(\allowbreak{}baa\_), unit-act(\allowbreak{}baa\_), unit-vnrm(\allowbreak{}baa\_), unit-vmul(\allowbreak{}baa\_), [1, vzero(\allowbreak{}baa\_)], unit-istar(\allowbreak{}baa\_)]
\end{quote}
mentioned by 21 result(s)

\smallskip

\noindent\textbf{\texttt{unitrow(a, n, i)}}\quad\textit{functoid}\index{unitrow(a, n, i)@\texttt{unitrow(a, n, i)}}\\
definition:
\begin{quote}\ttfamily\footnotesize\raggedright\noindent
matof(\allowbreak{}1, n, vnb-lambda(\allowbreak{}[rw, cl], cartesian(\allowbreak{}interval(\allowbreak{}1, 1), interval(\allowbreak{}1, n)), if(\allowbreak{}cl = i, one(\allowbreak{}a), zero(\allowbreak{}a))))
\end{quote}
characterized by \texttt{unitrow-type-uba-as-ring}  \textit{(definitional)}:
\begin{quote}\ttfamily\footnotesize\raggedright\noindent
forall(\allowbreak{}[r], is-unital-banach-algebra(\allowbreak{}r) implies forall(\allowbreak{}[n, i], n in nn implies unitrow(\allowbreak{}uba-as-ring(\allowbreak{}r), n, i) in mat(\allowbreak{}1, n, vec(\allowbreak{}r))))
\end{quote}
also: \texttt{lincomb-unitrow-cna-as-module} \texttt{lincomb-unitrow-sip-as-module} \texttt{lincomb-unitrow-uba-as-module} \texttt{lincomb-unitrow-cnvs-as-module} \texttt{unitrow-type-uba-as-ring-uiba-as-uba} \texttt{lincomb-unitrow-sip-as-module-csip-as-sip} \texttt{lincomb-unitrow-uba-as-module-uiba-as-uba} \texttt{unitrow-type} ... (12 in all) \\
mentioned by 92 result(s)

\smallskip

\noindent\textbf{\texttt{upd(upw\_, upn\_, upm\_)}}\quad\textit{functoid}\index{upd(upw_, upn_, upm_)@\texttt{upd(upw\_, upn\_, upm\_)}}\\
reads: \$1 below \$2, then \$3 \\
definition:
\begin{quote}\ttfamily\footnotesize\raggedright\noindent
vnb-lambda(\allowbreak{}upj\_, nn, if(\allowbreak{}upj\_ < upn\_, upw\_(\allowbreak{}upj\_), upm\_))
\end{quote}
characterized by \texttt{upd-facts}  \textit{(proven)}:
\begin{quote}\ttfamily\footnotesize\raggedright\noindent
forall(\allowbreak{}[w in fun(\allowbreak{}nn, nn), n in nn, m in nn], upd(\allowbreak{}w, n, m) in fun(\allowbreak{}nn, nn) and forall(\allowbreak{}[j in nn], j < n implies (\allowbreak{}upd(\allowbreak{}w, n, m))(\allowbreak{}j) = w(\allowbreak{}j)) and forall(\allowbreak{}[j in nn], not(\allowbreak{}j < n) implies (\allowbreak{}upd(\allowbreak{}w, n, m))(\allowbreak{}j) = m))
\end{quote}
mentioned by 5 result(s)

\smallskip

\noindent\textbf{\texttt{upper-half-plane}}\quad\textit{constant}\index{upper-half-plane@\texttt{upper-half-plane}}\\
reads: the upper half plane \{z : Im z > 0\} \\
defined by \texttt{upper-half-plane-def}

\smallskip

\noindent\textbf{\texttt{ury-apx}}\quad\textit{functoid}\index{ury-apx@\texttt{ury-apx}}\\
reads: the \$3-th Urysohn approximant at \$2 in \$1 \\
defined by \texttt{ury-apx-zero} \texttt{ury-apx-succ} \\
characterized by \texttt{ury-apx-succ}  \textit{(definitional)}:
\begin{quote}\ttfamily\footnotesize\raggedright\noindent
forall(\allowbreak{}[uat\_, uax\_, uan\_ in nn], ury-apx(\allowbreak{}uat\_, uax\_, succ(\allowbreak{}uan\_)) == vnb-lambda(\allowbreak{}uac\_, cu-pairs(\allowbreak{}uat\_), num(\allowbreak{}1/2) * (\allowbreak{}(\allowbreak{}ury-apx(\allowbreak{}uat\_, uax\_, uan\_))(\allowbreak{}ury-l(\allowbreak{}uat\_, uac\_)) + (\allowbreak{}ury-apx(\allowbreak{}uat\_, uax\_, uan\_))(\allowbreak{}ury-r(\allowbreak{}uat\_, uac\_)))))
\end{quote}
also: \texttt{ury-apx-zero} \texttt{ury-apx-mid} \texttt{ury-apx-mono} \texttt{ury-apx-on-c} \texttt{ury-apx-type} \texttt{ury-apx-off-u} \texttt{ury-apx-bounds} \texttt{ury-apx-zero-in} ... (12 in all) \\
mentioned by 25 result(s)

\smallskip

\noindent\textbf{\texttt{ury-k(uat\_, uac\_)}}\quad\textit{functoid}\index{ury-k(uat_, uac_)@\texttt{ury-k(uat\_, uac\_)}}\\
reads: the chosen closed middle of the pair \$2 in \$1 \\
definition:
\begin{quote}\ttfamily\footnotesize\raggedright\noindent
choice(\allowbreak{}\{uvk\_ in power(\allowbreak{}pts(\allowbreak{}uat\_)): is-top-closed(\allowbreak{}uat\_, uvk\_) and ury-v(\allowbreak{}uat\_, uac\_) subset uvk\_ and uvk\_ subset nth(\allowbreak{}2, uac\_)\})
\end{quote}
characterized by \texttt{ury-k-spec}  \textit{(proven)}:
\begin{quote}\ttfamily\footnotesize\raggedright\noindent
forall(\allowbreak{}[t, c], is-compact-t(\allowbreak{}t) implies is-hausdorff(\allowbreak{}t) implies c in cu-pairs(\allowbreak{}t) implies ury-k(\allowbreak{}t, c) in power(\allowbreak{}pts(\allowbreak{}t)) and is-top-closed(\allowbreak{}t, ury-k(\allowbreak{}t, c)) and ury-v(\allowbreak{}t, c) subset ury-k(\allowbreak{}t, c) and ury-k(\allowbreak{}t, c) subset nth(\allowbreak{}2, c))
\end{quote}
mentioned by 2 result(s)

\smallskip

\noindent\textbf{\texttt{ury-l(uat\_, uac\_)}}\quad\textit{functoid}\index{ury-l(uat_, uac_)@\texttt{ury-l(uat\_, uac\_)}}\\
reads: the lower half of the pair \$2 in \$1 \\
definition:
\begin{quote}\ttfamily\footnotesize\raggedright\noindent
[nth(\allowbreak{}1, uac\_), ury-v(\allowbreak{}uat\_, uac\_)]
\end{quote}
characterized by \texttt{ury-l-in}  \textit{(proven)}:
\begin{quote}\ttfamily\footnotesize\raggedright\noindent
forall(\allowbreak{}[t, c], is-compact-t(\allowbreak{}t) implies is-hausdorff(\allowbreak{}t) implies c in cu-pairs(\allowbreak{}t) implies ury-l(\allowbreak{}t, c) in cu-pairs(\allowbreak{}t))
\end{quote}
also: \texttt{ury-lr-nth} \texttt{ury-lim-mid} \\
mentioned by 9 result(s)

\smallskip

\noindent\textbf{\texttt{ury-lim(uat\_, uac\_, uax\_)}}\quad\textit{functoid}\index{ury-lim(uat_, uac_, uax_)@\texttt{ury-lim(uat\_, uac\_, uax\_)}}\\
reads: the Urysohn function of the pair \$2 in \$1 at \$3 \\
definition:
\begin{quote}\ttfamily\footnotesize\raggedright\noindent
esup(\allowbreak{}\{ulr\_ in rr: forsome(\allowbreak{}[uln\_ in nn], ulr\_ = (\allowbreak{}ury-apx(\allowbreak{}uat\_, uax\_, uln\_))(\allowbreak{}uac\_))\})
\end{quote}
characterized by \texttt{ury-lim-mid}  \textit{(proven)}:
\begin{quote}\ttfamily\footnotesize\raggedright\noindent
forall(\allowbreak{}[t, x], is-compact-t(\allowbreak{}t) implies is-hausdorff(\allowbreak{}t) implies forall(\allowbreak{}[c in cu-pairs(\allowbreak{}t)], ury-lim(\allowbreak{}t, c, x) = num(\allowbreak{}1/2) * (\allowbreak{}ury-lim(\allowbreak{}t, ury-l(\allowbreak{}t, c), x) + ury-lim(\allowbreak{}t, ury-r(\allowbreak{}t, c), x))))
\end{quote}
also: \texttt{ury-lim-osc} \texttt{ury-lim-on-c} \texttt{ury-lim-real} \texttt{ury-lim-least} \texttt{ury-lim-off-u} \texttt{ury-lim-upper} \texttt{ury-lim-unfold} \\
mentioned by 12 result(s)

\smallskip

\noindent\textbf{\texttt{ury-r(uat\_, uac\_)}}\quad\textit{functoid}\index{ury-r(uat_, uac_)@\texttt{ury-r(uat\_, uac\_)}}\\
reads: the upper half of the pair \$2 in \$1 \\
definition:
\begin{quote}\ttfamily\footnotesize\raggedright\noindent
[ury-k(\allowbreak{}uat\_, uac\_), nth(\allowbreak{}2, uac\_)]
\end{quote}
characterized by \texttt{ury-r-in}  \textit{(proven)}:
\begin{quote}\ttfamily\footnotesize\raggedright\noindent
forall(\allowbreak{}[t, c], is-compact-t(\allowbreak{}t) implies is-hausdorff(\allowbreak{}t) implies c in cu-pairs(\allowbreak{}t) implies ury-r(\allowbreak{}t, c) in cu-pairs(\allowbreak{}t))
\end{quote}
mentioned by 9 result(s)

\smallskip

\noindent\textbf{\texttt{ury-v(uat\_, uac\_)}}\quad\textit{functoid}\index{ury-v(uat_, uac_)@\texttt{ury-v(uat\_, uac\_)}}\\
reads: the chosen open middle of the pair \$2 in \$1 \\
definition:
\begin{quote}\ttfamily\footnotesize\raggedright\noindent
choice(\allowbreak{}\{uvv\_ in opens(\allowbreak{}uat\_): forsome(\allowbreak{}[uvk\_ in power(\allowbreak{}pts(\allowbreak{}uat\_))], is-top-closed(\allowbreak{}uat\_, uvk\_) and nth(\allowbreak{}1, uac\_) subset uvv\_ and uvv\_ subset uvk\_ and uvk\_ subset nth(\allowbreak{}2, uac\_))\})
\end{quote}
characterized by \texttt{ury-v-spec}  \textit{(proven)}:
\begin{quote}\ttfamily\footnotesize\raggedright\noindent
forall(\allowbreak{}[t, c], is-compact-t(\allowbreak{}t) implies is-hausdorff(\allowbreak{}t) implies c in cu-pairs(\allowbreak{}t) implies ury-v(\allowbreak{}t, c) in opens(\allowbreak{}t) and forsome(\allowbreak{}[uvk\_ in power(\allowbreak{}pts(\allowbreak{}t))], is-top-closed(\allowbreak{}t, uvk\_) and nth(\allowbreak{}1, c) subset ury-v(\allowbreak{}t, c) and ury-v(\allowbreak{}t, c) subset uvk\_ and uvk\_ subset nth(\allowbreak{}2, c)))
\end{quote}
mentioned by 3 result(s)

\smallskip

\noindent\textbf{\texttt{ust(otv\_)}}\quad\textit{functoid}\index{ust(otv_)@\texttt{ust(otv\_)}}\\
reads: the ultrastrong operator topology on the bounded operators on \$1 \\
definition:
\begin{quote}\ttfamily\footnotesize\raggedright\noindent
seminorm-top(\allowbreak{}bh-module(\allowbreak{}otv\_), sq-summable-seqs(\allowbreak{}otv\_), ust-pf(\allowbreak{}otv\_))
\end{quote}

\smallskip

\noindent\textbf{\texttt{ust-pf(otv\_)}}\quad\textit{functoid}\index{ust-pf(otv_)@\texttt{ust-pf(otv\_)}}\\
reads: the seminorms of the ultrastrong operator topology on \$1 \\
definition:
\begin{quote}\ttfamily\footnotesize\raggedright\noindent
vnb-lambda(\allowbreak{}ots\_, sq-summable-seqs(\allowbreak{}otv\_), vnb-lambda(\allowbreak{}ott\_, bounded-ops(\allowbreak{}otv\_, otv\_), sqrt(\allowbreak{}esum(\allowbreak{}vnb-lambda(\allowbreak{}otk\_, nn, ip-norm(\allowbreak{}otv\_, ott\_(\allowbreak{}ots\_(\allowbreak{}otk\_))) * ip-norm(\allowbreak{}otv\_, ott\_(\allowbreak{}ots\_(\allowbreak{}otk\_))))))))
\end{quote}

\smallskip

\noindent\textbf{\texttt{uwt(otv\_)}}\quad\textit{functoid}\index{uwt(otv_)@\texttt{uwt(otv\_)}}\\
reads: the ultraweak operator topology on the bounded operators on \$1 \\
definition:
\begin{quote}\ttfamily\footnotesize\raggedright\noindent
seminorm-top(\allowbreak{}bh-module(\allowbreak{}otv\_), cartesian(\allowbreak{}sq-summable-seqs(\allowbreak{}otv\_), sq-summable-seqs(\allowbreak{}otv\_)), uwt-pf(\allowbreak{}otv\_))
\end{quote}

\smallskip

\noindent\textbf{\texttt{uwt-pf(otv\_)}}\quad\textit{functoid}\index{uwt-pf(otv_)@\texttt{uwt-pf(otv\_)}}\\
reads: the seminorms of the ultraweak operator topology on \$1 \\
definition:
\begin{quote}\ttfamily\footnotesize\raggedright\noindent
vnb-lambda(\allowbreak{}otp\_, cartesian(\allowbreak{}sq-summable-seqs(\allowbreak{}otv\_), sq-summable-seqs(\allowbreak{}otv\_)), vnb-lambda(\allowbreak{}ott\_, bounded-ops(\allowbreak{}otv\_, otv\_), magnitude(\allowbreak{}cc-series-limit(\allowbreak{}vnb-lambda(\allowbreak{}otk\_, nn, (\allowbreak{}ip(\allowbreak{}otv\_))(\allowbreak{}ott\_(\allowbreak{}(\allowbreak{}nth(\allowbreak{}1, otp\_))(\allowbreak{}otk\_)), (\allowbreak{}nth(\allowbreak{}2, otp\_))(\allowbreak{}otk\_)))))))
\end{quote}

\smallskip


\section*{V}
\addcontentsline{toc}{section}{V}

\noindent\textbf{\texttt{vadd(s)}}\quad\textit{accessor}\index{vadd(s)@\texttt{vadd(s)}}\\
reads: the vector addition of \$1 \\
a structure slot \\
characterized by \texttt{rr-nvs@vadd}  \textit{(definitional)}:
\begin{quote}\ttfamily\footnotesize\raggedright\noindent
vadd(\allowbreak{}rr-nvs) == vnb-lambda(\allowbreak{}[x\_, y\_], cartesian(\allowbreak{}rr, rr), x\_ + y\_)
\end{quote}
also: \texttt{qq-line@vadd} \texttt{cips-vadd-type} \texttt{cips-vadd-vzero} \texttt{module-vadd-type} \texttt{sipq-vadd-csip-as-sip} \texttt{submodule-vadd-closed} \texttt{ram-ro-vadd-uba-as-ring} \texttt{lzq-d-vadd-cna-as-module} ... (274 in all) \\
mentioned by 3992 result(s)

\smallskip

\noindent\textbf{\texttt{var-permute(r, n, s, p)}}\quad\textit{functoid}\index{var-permute(r, n, s, p)@\texttt{var-permute(r, n, s, p)}}\\
reads: the polynomial \$4 with its indeterminates permuted by \$3 \\
definition:
\begin{quote}\ttfamily\footnotesize\raggedright\noindent
monalg-pullback(\allowbreak{}r, multi-index-monoid(\allowbreak{}interval(\allowbreak{}1, n)), vnb-lambda(\allowbreak{}mib\_, fun(\allowbreak{}interval(\allowbreak{}1, n), nn), mi-reindex(\allowbreak{}interval(\allowbreak{}1, n), s, mib\_)), p)
\end{quote}
characterized by \texttt{var-permute-is-ring-iso}  \textit{(proven)}:
\begin{quote}\ttfamily\footnotesize\raggedright\noindent
forall(\allowbreak{}[r\_, n\_, s\_], is-ring(\allowbreak{}r\_) implies n\_ in nn implies s\_ in bijection(\allowbreak{}interval(\allowbreak{}1, n\_), interval(\allowbreak{}1, n\_)) implies is-iso-ring(\allowbreak{}mpoly(\allowbreak{}r\_, n\_), mpoly(\allowbreak{}r\_, n\_), vnb-lambda(\allowbreak{}mpp\_, carr(\allowbreak{}mpoly(\allowbreak{}r\_, n\_)), var-permute(\allowbreak{}r\_, n\_, s\_, mpp\_))))
\end{quote}
mentioned by 1 result(s)

\smallskip

\noindent\textbf{\texttt{vec(s)}}\quad\textit{accessor}\index{vec(s)@\texttt{vec(s)}}\\
reads: the vectors of \$1 \\
a structure slot \\
characterized by \texttt{rr-nvs@vec}  \textit{(definitional)}:
\begin{quote}\ttfamily\footnotesize\raggedright\noindent
vec(\allowbreak{}rr-nvs) == rr
\end{quote}
also: \texttt{qq-line@vec} \texttt{is-cyclic-vector} \texttt{cips-ip-zero-vector} \texttt{is-cyclic-vector-rep} \texttt{sipq-vec-csip-as-sip} \texttt{is-normed-vector-space} \texttt{module-act-distrib-vec} \texttt{ram-ro-vec-uba-as-ring} ... (288 in all) \\
mentioned by 9136 result(s)

\smallskip

\noindent\textbf{\texttt{vec-functional(otv\_, otx\_, oty\_)}}\quad\textit{functoid}\index{vec-functional(otv_, otx_, oty_)@\texttt{vec-functional(otv\_, otx\_, oty\_)}}\\
reads: the functional T |-> <T \$2 | \$3> \\
definition:
\begin{quote}\ttfamily\footnotesize\raggedright\noindent
vnb-lambda(\allowbreak{}ott\_, bounded-ops(\allowbreak{}otv\_, otv\_), (\allowbreak{}ip(\allowbreak{}otv\_))(\allowbreak{}ott\_(\allowbreak{}otx\_), oty\_))
\end{quote}
characterized by \texttt{vec-functional-type}  \textit{(proven)}:
\begin{quote}\ttfamily\footnotesize\raggedright\noindent
forall(\allowbreak{}[v, xi\_, eta\_], is-complex-inner-product-space(\allowbreak{}v) implies xi\_ in vec(\allowbreak{}v) implies eta\_ in vec(\allowbreak{}v) implies vec-functional(\allowbreak{}v, xi\_, eta\_) in fun(\allowbreak{}bounded-ops(\allowbreak{}v, v), cc))
\end{quote}
also: \texttt{vec-functional-apply} \texttt{vec-functional-linear} \texttt{vec-functional-in-weak-dual} \texttt{vec-functional-in-strong-dual} \texttt{vec-functional-sot-continuous} \texttt{vec-functional-wot-continuous} \\
mentioned by 8 result(s)

\smallskip

\noindent\textbf{\texttt{vector-space}}\quad\textit{refinement}\index{vector-space@\texttt{vector-space}}\\
predicate \texttt{is-vector-space}; same shape as module

\smallskip

\noindent\textbf{\texttt{vlet}}\quad\textit{tactic}\index{vlet@\texttt{vlet}}\\
\texttt{(vlet (n1 ...) FORMER)} -- Bind proof-object names from the current state.  FORMER is (match PATTERN) -- bind each name to its slot in a context formula matching PATTERN -- or (choice [v body]) -- eliminate an in-context existential and bind its witness. \\
declared in \texttt{tactics-help.scm}

\smallskip

\noindent\textbf{\texttt{vmul(s)}}\quad\textit{accessor}\index{vmul(s)@\texttt{vmul(s)}}\\
reads: the product of \$1 \\
a structure slot \\
characterized by \texttt{bh-ro-vmul}  \textit{(proven)}:
\begin{quote}\ttfamily\footnotesize\raggedright\noindent
forall(\allowbreak{}[v], is-hilbert-space(\allowbreak{}v) implies vmul(\allowbreak{}bh-algebra(\allowbreak{}v)) == bh-vmul(\allowbreak{}v))
\end{quote}
also: \texttt{b1s-nv-vmul} \texttt{b1s-ro-vmul} \texttt{b1u-ro-vmul} \texttt{bcp-ro-vmul} \texttt{bcs-uu-vmul} \texttt{kal-uu-vmul} \texttt{kuc-ro-vmul} \texttt{qba-ro-vmul} ... (35 in all) \\
mentioned by 681 result(s)

\smallskip

\noindent\textbf{\texttt{vnb-lambda}}\quad\textit{operator}\index{vnb-lambda@\texttt{vnb-lambda}}\\
a term-forming head \\
mentioned by 3505 result(s)

\smallskip

\noindent\textbf{\texttt{vneg(s)}}\quad\textit{accessor}\index{vneg(s)@\texttt{vneg(s)}}\\
reads: the vector negation of \$1 \\
a structure slot \\
characterized by \texttt{rr-nvs@vneg}  \textit{(definitional)}:
\begin{quote}\ttfamily\footnotesize\raggedright\noindent
vneg(\allowbreak{}rr-nvs) == vnb-lambda(\allowbreak{}x\_, rr, -x\_)
\end{quote}
also: \texttt{qq-line@vneg} \texttt{module-vneg-type} \texttt{sipq-vneg-csip-as-sip} \texttt{submodule-vneg-closed} \texttt{ram-ro-vneg-uba-as-ring} \texttt{lzq-d-vneg-cna-as-module} \texttt{lzq-d-vneg-sip-as-module} \texttt{lzq-d-vneg-uba-as-module} ... (184 in all) \\
mentioned by 2093 result(s)

\smallskip

\noindent\textbf{\texttt{vnrm(s)}}\quad\textit{accessor}\index{vnrm(s)@\texttt{vnrm(s)}}\\
reads: the vector norm of \$1 \\
a structure slot \\
characterized by \texttt{rr-nvs@vnrm}  \textit{(definitional)}:
\begin{quote}\ttfamily\footnotesize\raggedright\noindent
vnrm(\allowbreak{}rr-nvs) == vnb-lambda(\allowbreak{}x\_, rr, abs(\allowbreak{}x\_))
\end{quote}
also: \texttt{nvs-vnrm-homog-cnvs-as-nvs} \texttt{nvs-vnrm-vzero-cnvs-as-nvs} \texttt{dual-vnrm-typing-cnvs-as-nvs} \texttt{nvs-vnrm-homog-c-cnvs-as-nvs} \texttt{nvs-vnrm-triangle-cnvs-as-nvs} \texttt{vnrm-real} \texttt{bh-ro-vnrm} \texttt{b1s-nv-vnrm} ... (75 in all) \\
mentioned by 521 result(s)

\smallskip

\noindent\textbf{\texttt{vone(s)}}\quad\textit{accessor}\index{vone(s)@\texttt{vone(s)}}\\
reads: the unit of \$1 \\
a structure slot \\
characterized by \texttt{bh-ro-vone}  \textit{(proven)}:
\begin{quote}\ttfamily\footnotesize\raggedright\noindent
forall(\allowbreak{}[v], is-hilbert-space(\allowbreak{}v) implies vone(\allowbreak{}bh-algebra(\allowbreak{}v)) == bop-id(\allowbreak{}v))
\end{quote}
also: \texttt{b1s-ro-vone} \texttt{b1u-ro-vone} \texttt{bcp-ro-vone} \texttt{bcs-uu-vone} \texttt{kal-uu-vone} \texttt{qba-ro-vone} \texttt{qba-vone-left} \texttt{qba-vone-norm} ... (19 in all) \\
mentioned by 318 result(s)

\smallskip

\noindent\textbf{\texttt{vstar(s)}}\quad\textit{accessor}\index{vstar(s)@\texttt{vstar(s)}}\\
reads: the involution of \$1 \\
a structure slot \\
characterized by \texttt{bh-ro-vstar}  \textit{(proven)}:
\begin{quote}\ttfamily\footnotesize\raggedright\noindent
forall(\allowbreak{}[v], is-hilbert-space(\allowbreak{}v) implies vstar(\allowbreak{}bh-algebra(\allowbreak{}v)) == bh-vstar(\allowbreak{}v))
\end{quote}
also: \texttt{b1s-ro-vstar} \texttt{bcp-sub-vstar} \texttt{kal-uiba-vstar-in} \texttt{bcs-uiba-vstar-fun} \texttt{kal-uiba-vstar-fun} \texttt{ccont-algebra-vstar} \\
mentioned by 120 result(s)

\smallskip

\noindent\textbf{\texttt{vzero(s)}}\quad\textit{accessor}\index{vzero(s)@\texttt{vzero(s)}}\\
reads: the zero vector of \$1 \\
a structure slot \\
characterized by \texttt{rr-nvs@vzero}  \textit{(definitional)}:
\begin{quote}\ttfamily\footnotesize\raggedright\noindent
vzero(\allowbreak{}rr-nvs) == 0
\end{quote}
also: \texttt{cips-vzero-in} \texttt{qq-line@vzero} \texttt{cips-vadd-vzero} \texttt{module-vzero-in} \texttt{submodule-vzero-in} \texttt{sipq-vzero-csip-as-sip} \texttt{ram-ro-vzero-uba-as-ring} \texttt{sip-vzero-in-csip-as-sip} ... (168 in all) \\
mentioned by 2022 result(s)

\smallskip


\section*{W}
\addcontentsline{toc}{section}{W}

\noindent\textbf{\texttt{wbc}}\quad\textit{tactic}\index{wbc@\texttt{wbc}}\\
\texttt{(wbc ['name])} -- Witness-shape backchain: on an `exists v. ...' goal, cite a PSS lemma that MANUFACTURES a witness of that shape, so you can finish with inst+/grind/ew. \\
declared in \texttt{tactics-help.scm}

\smallskip

\noindent\textbf{\texttt{weak-dual(otv\_)}}\quad\textit{functoid}\index{weak-dual(otv_)@\texttt{weak-dual(otv\_)}}\\
reads: the weakly continuous linear functionals on the bounded operators on \$1 \\
definition:
\begin{quote}\ttfamily\footnotesize\raggedright\noindent
op-dual(\allowbreak{}otv\_, wot(\allowbreak{}otv\_))
\end{quote}
characterized by \texttt{weak-dual-bound}  \textit{(proven)}:
\begin{quote}\ttfamily\footnotesize\raggedright\noindent
forall(\allowbreak{}[v, f\_], is-complex-inner-product-space(\allowbreak{}v) implies f\_ in weak-dual(\allowbreak{}v) implies forsome(\allowbreak{}[sbf\_ in set], card(\allowbreak{}sbf\_) in nn and sbf\_ subset cartesian(\allowbreak{}vec(\allowbreak{}v), vec(\allowbreak{}v)) and forsome(\allowbreak{}[sbc\_ in rr], pos-rr(\allowbreak{}sbc\_) and forall(\allowbreak{}[sbx\_ in bounded-ops(\allowbreak{}v, v), sbt\_ in rr], 0 <= sbt\_ implies forall(\allowbreak{}[sbi\_ in sbf\_], magnitude(\allowbreak{}(\allowbreak{}ip(\allowbreak{}v))(\allowbreak{}sbx\_(\allowbreak{}nth(\allowbreak{}1, sbi\_)), nth(\allowbreak{}2, sbi\_))) <= sbt\_) implies magnitude(\allowbreak{}f\_(\allowbreak{}sbx\_)) <= sbc\_ * sbt\_))))
\end{quote}
also: \texttt{weak-dual-unfold} \texttt{vec-functional-in-weak-dual} \texttt{weak-dual-subset-strong-dual} \\
mentioned by 5 result(s)

\smallskip

\noindent\textbf{\texttt{weak-pairing-pf(otq\_, ota\_)}}\quad\textit{functoid}\index{weak-pairing-pf(otq_, ota_)@\texttt{weak-pairing-pf(otq\_, ota\_)}}\\
reads: the seminorms x |-> |f(x)| for f in \$2 \\
definition:
\begin{quote}\ttfamily\footnotesize\raggedright\noindent
vnb-lambda(\allowbreak{}otf\_, ota\_, vnb-lambda(\allowbreak{}otx\_, vec(\allowbreak{}otq\_), magnitude(\allowbreak{}otf\_(\allowbreak{}otx\_))))
\end{quote}

\smallskip

\noindent\textbf{\texttt{weak-pairing-top(otq\_, ota\_)}}\quad\textit{functoid}\index{weak-pairing-top(otq_, ota_)@\texttt{weak-pairing-top(otq\_, ota\_)}}\\
reads: the weak topology on \$1 of the functionals \$2 \\
definition:
\begin{quote}\ttfamily\footnotesize\raggedright\noindent
seminorm-top(\allowbreak{}otq\_, ota\_, weak-pairing-pf(\allowbreak{}otq\_, ota\_))
\end{quote}

\smallskip

\noindent\textbf{\texttt{weak-star-top(m, t)}}\quad\textit{functoid}\index{weak-star-top(m, t)@\texttt{weak-star-top(m, t)}}\\
reads: the weak* topology on the topological dual of \$1 with topology \$2 \\
definition:
\begin{quote}\ttfamily\footnotesize\raggedright\noindent
subspace-top(\allowbreak{}tproduct-top(\allowbreak{}vec(\allowbreak{}m), vnb-lambda(\allowbreak{}b6n\_, vec(\allowbreak{}m), metric-top(\allowbreak{}rr-ms))), top-dual(\allowbreak{}m, t))
\end{quote}
characterized by \texttt{weak-star-top-unfold}  \textit{(proven)}:
\begin{quote}\ttfamily\footnotesize\raggedright\noindent
forall(\allowbreak{}[m, t], weak-star-top(\allowbreak{}m, t) == subspace-top(\allowbreak{}tproduct-top(\allowbreak{}vec(\allowbreak{}m), vnb-lambda(\allowbreak{}b6n\_, vec(\allowbreak{}m), metric-top(\allowbreak{}rr-ms))), top-dual(\allowbreak{}m, t)))
\end{quote}
mentioned by 3 result(s)

\smallskip

\noindent\textbf{\texttt{wf-fml(mce\_)}}\quad\textit{predicate}\index{wf-fml(mce_)@\texttt{wf-fml(mce\_)}}\\
reads: \$1 is a well-formed formula \\
defined by \texttt{wf-fml} \\
characterized by \texttt{wf-fml}  \textit{(definitional)}:
\begin{quote}\ttfamily\footnotesize\raggedright\noindent
forall(\allowbreak{}[mce\_], wf-fml(\allowbreak{}mce\_) iff 2 in expr-wks(\allowbreak{}mce\_))
\end{quote}
mentioned by 6 result(s)

\smallskip

\noindent\textbf{\texttt{wf-term(mce\_)}}\quad\textit{predicate}\index{wf-term(mce_)@\texttt{wf-term(mce\_)}}\\
reads: \$1 is a well-formed term \\
defined by \texttt{wf-term} \\
characterized by \texttt{wf-term}  \textit{(definitional)}:
\begin{quote}\ttfamily\footnotesize\raggedright\noindent
forall(\allowbreak{}[mce\_], wf-term(\allowbreak{}mce\_) iff 1 in expr-wks(\allowbreak{}mce\_))
\end{quote}
mentioned by 6 result(s)

\smallskip

\noindent\textbf{\texttt{winding(wnp\_, wnd\_, wna\_, wnb\_, wnz\_)}}\quad\textit{functoid}\index{winding(wnp_, wnd_, wna_, wnb_, wnz_)@\texttt{winding(wnp\_, wnd\_, wna\_, wnb\_, wnz\_)}}\\
reads: the winding number about \$5 of the road \$1 with derivative \$2 on [\$3, \$4] \\
definition:
\begin{quote}\ttfamily\footnotesize\raggedright\noindent
recip(\allowbreak{}(\allowbreak{}2 * pi) * 1i) * line-int(\allowbreak{}vnb-lambda(\allowbreak{}wnw\_, difference(\allowbreak{}cc, singleton(\allowbreak{}wnz\_)), recip(\allowbreak{}wnw\_ - wnz\_)), wnp\_, wnd\_, wna\_, wnb\_)
\end{quote}
characterized by \texttt{winding-in-cc}  \textit{(proven)}:
\begin{quote}\ttfamily\footnotesize\raggedright\noindent
forall(\allowbreak{}[pgam, dgam, a, b], is-road(\allowbreak{}pgam, dgam, a, b) implies forall(\allowbreak{}[wnz\_ in cc], not(\allowbreak{}wnz\_ in trace(\allowbreak{}pgam, a, b)) implies winding(\allowbreak{}pgam, dgam, a, b, wnz\_) in cc))
\end{quote}
also: \texttt{winding-juxta} \texttt{homotopy-winding} \texttt{winding-opposite} \texttt{winding-unfold-eq} \texttt{zeros-count-winding} \texttt{winding-of-composite} \texttt{winding-circle-center} \texttt{winding-circle-inside} ... (16 in all) \\
mentioned by 43 result(s)

\smallskip

\noindent\textbf{\texttt{wk}}\quad\textit{tactic}\index{wk@\texttt{wk}}\\
\texttt{(wk hyp)} -- Weaken: drop a cited assumption from the context to tidy the hypothesis list.  `hyp' may be a formula, a "string", or a 1-based assumption index. \\
declared in \texttt{tactics-help.scm}

\smallskip

\noindent\textbf{\texttt{wks-stage}}\quad\textit{functoid}\index{wks-stage@\texttt{wks-stage}}\\
reads: the kind map on the stage \$1 \\
defined by \texttt{wks-stage-zero} \texttt{wks-stage-succ} \\
characterized by \texttt{wks-stage-succ}  \textit{(definitional)}:
\begin{quote}\ttfamily\footnotesize\raggedright\noindent
forall(\allowbreak{}[mcn\_ in nn], wks-stage(\allowbreak{}succ(\allowbreak{}mcn\_)) == vnb-lambda(\allowbreak{}mce\_, expr-stage(\allowbreak{}succ(\allowbreak{}mcn\_)), if(\allowbreak{}mce\_ in expr-stage(\allowbreak{}mcn\_), (\allowbreak{}wks-stage(\allowbreak{}mcn\_))(\allowbreak{}mce\_), if(\allowbreak{}nth(\allowbreak{}1, mce\_) = 3, \{mckk\_ in nn: mckk\_ = hd-res(\allowbreak{}nth(\allowbreak{}2, mce\_)) and app-ar-ok(\allowbreak{}nth(\allowbreak{}2, mce\_), length(\allowbreak{}nth(\allowbreak{}3, mce\_))) and forall(\allowbreak{}[mcw\_ in make-set(\allowbreak{}nth(\allowbreak{}3, mce\_))], hd-argk(\allowbreak{}nth(\allowbreak{}2, mce\_)) in (\allowbreak{}wks-stage(\allowbreak{}mcn\_))(\allowbreak{}mcw\_))\}, \{mckk\_ in nn: mckk\_ = hd-bres(\allowbreak{}nth(\allowbreak{}2, mce\_)) and bind-sig(\allowbreak{}nth(\allowbreak{}2, mce\_), nth(\allowbreak{}3, mce\_), length(\allowbreak{}nth(\allowbreak{}4, mce\_))) and forall(\allowbreak{}[mcw\_ in make-set(\allowbreak{}nth(\allowbreak{}4, mce\_))], 1 in (\allowbreak{}wks-stage(\allowbreak{}mcn\_))(\allowbreak{}mcw\_)) and hd-body(\allowbreak{}nth(\allowbreak{}2, mce\_)) in (\allowbreak{}wks-stage(\allowbreak{}mcn\_))(\allowbreak{}nth(\allowbreak{}5, mce\_))\}))))
\end{quote}
also: \texttt{wks-stage-zero} \texttt{wks-stage-coh} \texttt{wks-stage-step} \\
mentioned by 10 result(s)

\smallskip

\noindent\textbf{\texttt{wot(otv\_)}}\quad\textit{functoid}\index{wot(otv_)@\texttt{wot(otv\_)}}\\
reads: the weak operator topology on the bounded operators on \$1 \\
definition:
\begin{quote}\ttfamily\footnotesize\raggedright\noindent
seminorm-top(\allowbreak{}bh-module(\allowbreak{}otv\_), cartesian(\allowbreak{}vec(\allowbreak{}otv\_), vec(\allowbreak{}otv\_)), wot-pf(\allowbreak{}otv\_))
\end{quote}
characterized by \texttt{wot-is-tvs}  \textit{(proven)}:
\begin{quote}\ttfamily\footnotesize\raggedright\noindent
forall(\allowbreak{}[v], is-complex-inner-product-space(\allowbreak{}v) implies is-tvs(\allowbreak{}cc-normed-field, bh-module(\allowbreak{}v), wot(\allowbreak{}v)))
\end{quote}
also: \texttt{wot-unfold} \texttt{sot-to-wot-continuous} \texttt{bop-adjoint-wot-continuous} \texttt{bop-left-mult-wot-continuous} \texttt{bop-right-mult-wot-continuous} \texttt{vec-functional-wot-continuous} \\
mentioned by 10 result(s)

\smallskip

\noindent\textbf{\texttt{wot-pf(otv\_)}}\quad\textit{functoid}\index{wot-pf(otv_)@\texttt{wot-pf(otv\_)}}\\
reads: the seminorms T |-> |<T xi | eta>| of the weak operator topology on \$1 \\
definition:
\begin{quote}\ttfamily\footnotesize\raggedright\noindent
vnb-lambda(\allowbreak{}otp\_, cartesian(\allowbreak{}vec(\allowbreak{}otv\_), vec(\allowbreak{}otv\_)), vnb-lambda(\allowbreak{}ott\_, bounded-ops(\allowbreak{}otv\_, otv\_), magnitude(\allowbreak{}(\allowbreak{}ip(\allowbreak{}otv\_))(\allowbreak{}ott\_(\allowbreak{}nth(\allowbreak{}1, otp\_)), nth(\allowbreak{}2, otp\_)))))
\end{quote}
characterized by \texttt{wot-pf-type}  \textit{(proven)}:
\begin{quote}\ttfamily\footnotesize\raggedright\noindent
forall(\allowbreak{}[v], is-complex-inner-product-space(\allowbreak{}v) implies wot-pf(\allowbreak{}v) in fun(\allowbreak{}cartesian(\allowbreak{}vec(\allowbreak{}v), vec(\allowbreak{}v)), fun(\allowbreak{}bounded-ops(\allowbreak{}v, v), rr)))
\end{quote}
also: \texttt{wot-pf-apply} \texttt{wot-pf-apply-pair} \\
mentioned by 11 result(s)

\smallskip


\section*{X}
\addcontentsline{toc}{section}{X}

\noindent\textbf{\texttt{x-minus(a, c)}}\quad\textit{functoid}\index{x-minus(a, c)@\texttt{x-minus(a, c)}}\\
reads: the polynomial x - \$2 over \$1 \\
definition:
\begin{quote}\ttfamily\footnotesize\raggedright\noindent
(\allowbreak{}add(\allowbreak{}poly(\allowbreak{}a)))(\allowbreak{}monomial(\allowbreak{}a, one(\allowbreak{}a), 1), monomial(\allowbreak{}a, (\allowbreak{}neg(\allowbreak{}a))(\allowbreak{}c), 0))
\end{quote}
characterized by \texttt{x-minus-degree-one}  \textit{(proven)}:
\begin{quote}\ttfamily\footnotesize\raggedright\noindent
forall(\allowbreak{}[a\_], is-commutative-ring(\allowbreak{}a\_) implies not(\allowbreak{}one(\allowbreak{}a\_) = zero(\allowbreak{}a\_)) implies forall(\allowbreak{}[al\_ in carr(\allowbreak{}a\_)], x-minus(\allowbreak{}a\_, al\_) in carr(\allowbreak{}poly(\allowbreak{}a\_)) and not(\allowbreak{}x-minus(\allowbreak{}a\_, al\_) = zero(\allowbreak{}poly(\allowbreak{}a\_))) and deg(\allowbreak{}a\_, x-minus(\allowbreak{}a\_, al\_)) = 1))
\end{quote}
also: \texttt{poly-eval-x-minus-mul} \texttt{poly-root-iff-x-minus-divides} \\
mentioned by 5 result(s)

\smallskip

\noindent\textbf{\texttt{x-section(ys, cs, x)}}\quad\textit{functoid}\index{x-section(ys, cs, x)@\texttt{x-section(ys, cs, x)}}\\
reads: the section of \$2 at the first coordinate \$3, a subset of \$1 \\
definition:
\begin{quote}\ttfamily\footnotesize\raggedright\noindent
\{xsy\_ in ys: [x, xsy\_] in cs\}
\end{quote}
characterized by \texttt{x-section-mono}  \textit{(proven)}:
\begin{quote}\ttfamily\footnotesize\raggedright\noindent
forall(\allowbreak{}[ys, aa, bb, x], aa subset bb implies x-section(\allowbreak{}ys, aa, x) subset x-section(\allowbreak{}ys, bb, x))
\end{quote}
also: \texttt{x-section-empty} \texttt{x-section-in-cut} \texttt{x-section-unfold} \texttt{x-section-difference} \texttt{x-section-measurable} \texttt{x-section-membership} \texttt{x-section-rectangle-in} \texttt{x-section-measure-apply} ... (11 in all) \\
mentioned by 25 result(s)

\smallskip

\noindent\textbf{\texttt{x-section-integral(xs, ys, cb, nu, f)}}\quad\textit{functoid}\index{x-section-integral(xs, ys, cb, nu, f)@\texttt{x-section-integral(xs, ys, cb, nu, f)}}\\
reads: the function x |-> the integral over \$2 with respect to \$4 of \$5(x, .), on \$1 \\
definition:
\begin{quote}\ttfamily\footnotesize\raggedright\noindent
vnb-lambda(\allowbreak{}six\_, xs, integral(\allowbreak{}ys, cb, nu, fn-x-section(\allowbreak{}ys, f, six\_)))
\end{quote}
characterized by \texttt{x-section-integral-type}  \textit{(proven)}:
\begin{quote}\ttfamily\footnotesize\raggedright\noindent
forall(\allowbreak{}[xs, ca, ys, cb, nu, f], is-sigma-algebra(\allowbreak{}xs, ca) implies is-sigma-algebra(\allowbreak{}ys, cb) implies is-measure(\allowbreak{}ys, cb, nu) implies is-measurable-fn(\allowbreak{}cartesian(\allowbreak{}xs, ys), product-sigma(\allowbreak{}xs, ca, ys, cb), f) implies x-section-integral(\allowbreak{}xs, ys, cb, nu, f) in fun(\allowbreak{}xs, rr-pos-star))
\end{quote}
also: \texttt{x-section-integral-apply} \texttt{x-section-integral-limit} \\
mentioned by 28 result(s)

\smallskip

\noindent\textbf{\texttt{x-section-measure(xs, ys, nu, cs)}}\quad\textit{functoid}\index{x-section-measure(xs, ys, nu, cs)@\texttt{x-section-measure(xs, ys, nu, cs)}}\\
reads: the function x |-> \$3(the section of \$4 at x) on \$1 \\
definition:
\begin{quote}\ttfamily\footnotesize\raggedright\noindent
vnb-lambda(\allowbreak{}smx\_, xs, nu(\allowbreak{}x-section(\allowbreak{}ys, cs, smx\_)))
\end{quote}
characterized by \texttt{x-section-measure-type}  \textit{(proven)}:
\begin{quote}\ttfamily\footnotesize\raggedright\noindent
forall(\allowbreak{}[xs, ca, mu, ys, cb, nu, cs], is-measure(\allowbreak{}xs, ca, mu) implies is-measure(\allowbreak{}ys, cb, nu) implies cs in product-sigma(\allowbreak{}xs, ca, ys, cb) implies x-section-measure(\allowbreak{}xs, ys, nu, cs) in fun(\allowbreak{}xs, rr-pos-star))
\end{quote}
also: \texttt{x-section-measure-apply} \texttt{x-section-measure-union} \texttt{x-section-measure-rectangle} \texttt{x-section-measure-difference} \\
mentioned by 13 result(s)

\smallskip


\section*{Y}
\addcontentsline{toc}{section}{Y}

\noindent\textbf{\texttt{y-section(xs, cs, y)}}\quad\textit{functoid}\index{y-section(xs, cs, y)@\texttt{y-section(xs, cs, y)}}\\
reads: the section of \$2 at the second coordinate \$3, a subset of \$1 \\
definition:
\begin{quote}\ttfamily\footnotesize\raggedright\noindent
\{ysx\_ in xs: [ysx\_, y] in cs\}
\end{quote}
characterized by \texttt{y-section-mono}  \textit{(proven)}:
\begin{quote}\ttfamily\footnotesize\raggedright\noindent
forall(\allowbreak{}[xs, aa, bb, y], aa subset bb implies y-section(\allowbreak{}xs, aa, y) subset y-section(\allowbreak{}xs, bb, y))
\end{quote}
also: \texttt{y-section-in-cut} \texttt{y-section-unfold} \texttt{y-section-difference} \texttt{y-section-measurable} \texttt{y-section-membership} \texttt{y-section-rectangle-in} \texttt{y-section-measure-apply} \texttt{y-section-rectangle-out} ... (10 in all) \\
mentioned by 20 result(s)

\smallskip

\noindent\textbf{\texttt{y-section-integral(xs, ca, mu, ys, f)}}\quad\textit{functoid}\index{y-section-integral(xs, ca, mu, ys, f)@\texttt{y-section-integral(xs, ca, mu, ys, f)}}\\
reads: the function y |-> the integral over \$1 with respect to \$3 of \$5(., y), on \$4 \\
definition:
\begin{quote}\ttfamily\footnotesize\raggedright\noindent
vnb-lambda(\allowbreak{}siy\_, ys, integral(\allowbreak{}xs, ca, mu, fn-y-section(\allowbreak{}xs, f, siy\_)))
\end{quote}
characterized by \texttt{y-section-integral-type}  \textit{(proven)}:
\begin{quote}\ttfamily\footnotesize\raggedright\noindent
forall(\allowbreak{}[xs, ca, mu, ys, cb, f], is-sigma-algebra(\allowbreak{}xs, ca) implies is-sigma-algebra(\allowbreak{}ys, cb) implies is-measure(\allowbreak{}xs, ca, mu) implies is-measurable-fn(\allowbreak{}cartesian(\allowbreak{}xs, ys), product-sigma(\allowbreak{}xs, ca, ys, cb), f) implies y-section-integral(\allowbreak{}xs, ca, mu, ys, f) in fun(\allowbreak{}ys, rr-pos-star))
\end{quote}
also: \texttt{y-section-integral-apply} \texttt{y-section-integral-limit} \\
mentioned by 8 result(s)

\smallskip

\noindent\textbf{\texttt{y-section-measure(xs, ys, mu, cs)}}\quad\textit{functoid}\index{y-section-measure(xs, ys, mu, cs)@\texttt{y-section-measure(xs, ys, mu, cs)}}\\
reads: the function y |-> \$3(the section of \$4 at y) on \$2 \\
definition:
\begin{quote}\ttfamily\footnotesize\raggedright\noindent
vnb-lambda(\allowbreak{}smy\_, ys, mu(\allowbreak{}y-section(\allowbreak{}xs, cs, smy\_)))
\end{quote}
characterized by \texttt{y-section-measure-type}  \textit{(proven)}:
\begin{quote}\ttfamily\footnotesize\raggedright\noindent
forall(\allowbreak{}[xs, ca, mu, ys, cb, nu, cs], is-measure(\allowbreak{}xs, ca, mu) implies is-measure(\allowbreak{}ys, cb, nu) implies cs in product-sigma(\allowbreak{}xs, ca, ys, cb) implies y-section-measure(\allowbreak{}xs, ys, mu, cs) in fun(\allowbreak{}ys, rr-pos-star))
\end{quote}
also: \texttt{y-section-measure-apply} \texttt{y-section-measure-union} \texttt{y-section-measure-rectangle} \texttt{y-section-measure-difference} \\
mentioned by 10 result(s)

\smallskip


\section*{Z}
\addcontentsline{toc}{section}{Z}

\noindent\textbf{\texttt{zen}}\quad\textit{defined-fn}\index{zen@\texttt{zen}}\\
defined by \texttt{zen-zero} \texttt{zen-succ} \texttt{zen-limit} \\
characterized by \texttt{zen-succ}  \textit{(definitional)}:
\begin{quote}\ttfamily\footnotesize\raggedright\noindent
forall(\allowbreak{}[x, alpha in ord], zen(\allowbreak{}x, succ\_ord(\allowbreak{}alpha)) == choice(\allowbreak{}\{y\_ in x: forall(\allowbreak{}[c\_ in ord], ord-le(\allowbreak{}c\_, alpha) implies not(\allowbreak{}zen(\allowbreak{}x, c\_) = y\_))\}))
\end{quote}
also: \texttt{zen-zero} \texttt{zen-limit} \texttt{zen-hits} \texttt{zen-step} \texttt{zen-exhausts} \texttt{zen-enum-in-fun} \texttt{zen-enum-injective} \texttt{zen-enum-surjective} \\
mentioned by 15 result(s)

\smallskip

\noindent\textbf{\texttt{zero(s)}}\quad\textit{accessor}\index{zero(s)@\texttt{zero(s)}}\\
reads: the additive identity of \$1 \\
a structure slot \\
characterized by \texttt{rq-zero}  \textit{(definitional)}:
\begin{quote}\ttfamily\footnotesize\raggedright\noindent
forall(\allowbreak{}[r], zero(\allowbreak{}ringoid-quotient-ring(\allowbreak{}r)) == class(\allowbreak{}ringoid-setoid(\allowbreak{}r), zero(\allowbreak{}r)))
\end{quote}
also: \texttt{sum-zero} \texttt{comb-kk-zero} \texttt{qq-ring@zero} \texttt{ring-zero-in} \texttt{zz-ring@zero} \texttt{qq-field@zero} \texttt{rr-field@zero} \texttt{ringoid-ideal-zero} ... (164 in all) \\
mentioned by 1073 result(s)

\smallskip

\noindent\textbf{\texttt{zero-it}}\quad\textit{tactic}\index{zero-it@\texttt{zero-it}}\\
\texttt{(zero-it)} -- Move an equation goal P = Q (or P == Q, bare or under typed universals) to R = 0, where R is the normal form of P - Q (crs's calculator, highest degree first); prints `P - Q normalises to R' and returns R.  An inequality goal P <= Q or P < Q (the order of RR; atoms typed in NN, ZZ or QQ are lifted) is moved to 0 <= R or 0 < R, R the normal form of Q - P, the branch closed through rr-le-from-diff-nonneg / rr-lt-from-diff-pos; a numeral R decides it (0 <= -1: FALSE, nothing changed; 0 < 2: closed).  CC and an abstract ring have no order and are declined.  R = 0: closed by crs.  R a non-zero numeral: says the goal is FALSE and changes nothing.  Otherwise ONE recorded step leaves the goal R = 0, the atoms' typings landed or listed as owed side leaves. \\
declared in \texttt{tactics-help.scm}

\smallskip

\noindent\textbf{\texttt{zero-order(zou\_, zof\_, zoa\_)}}\quad\textit{functoid}\index{zero-order(zou_, zof_, zoa_)@\texttt{zero-order(zou\_, zof\_, zoa\_)}}\\
reads: the order of the zero of \$2 at \$3 (on \$1) \\
definition:
\begin{quote}\ttfamily\footnotesize\raggedright\noindent
iota(\allowbreak{}zon\_, has-zero-of-order(\allowbreak{}zou\_, zof\_, zoa\_, zon\_))
\end{quote}
characterized by \texttt{zero-order-value}  \textit{(proven)}:
\begin{quote}\ttfamily\footnotesize\raggedright\noindent
forall(\allowbreak{}[zlu\_, zlf\_, zla\_, zlm\_], has-zero-of-order(\allowbreak{}zlu\_, zlf\_, zla\_, zlm\_) implies zero-order(\allowbreak{}zlu\_, zlf\_, zla\_) = zlm\_)
\end{quote}
mentioned by 12 result(s)

\smallskip

\noindent\textbf{\texttt{zero-ring}}\quad\textit{defined-fn}\index{zero-ring@\texttt{zero-ring}}\\
defined by \texttt{zero-ring-def}

\smallskip

\noindent\textbf{\texttt{zeromat(a, m, n)}}\quad\textit{functoid}\index{zeromat(a, m, n)@\texttt{zeromat(a, m, n)}}\\
definition:
\begin{quote}\ttfamily\footnotesize\raggedright\noindent
matof(\allowbreak{}m, n, vnb-lambda(\allowbreak{}[i, j], cartesian(\allowbreak{}interval(\allowbreak{}1, m), interval(\allowbreak{}1, n)), zero(\allowbreak{}a)))
\end{quote}
characterized by \texttt{zeromat-type-uba-as-ring}  \textit{(definitional)}:
\begin{quote}\ttfamily\footnotesize\raggedright\noindent
forall(\allowbreak{}[r], is-unital-banach-algebra(\allowbreak{}r) implies forall(\allowbreak{}[m, n], m in nn implies n in nn implies zeromat(\allowbreak{}uba-as-ring(\allowbreak{}r), m, n) in mat(\allowbreak{}m, n, vec(\allowbreak{}r))))
\end{quote}
also: \texttt{zeromat-type-uba-as-ring-uiba-as-uba} \texttt{zeromat-type} \texttt{entry-of-zeromat} \\
mentioned by 129 result(s)

\smallskip

\noindent\textbf{\texttt{zeros-accumulate-at(u, f, a)}}\quad\textit{predicate}\index{zeros-accumulate-at(u, f, a)@\texttt{zeros-accumulate-at(u, f, a)}}\\
reads: the zeros of \$2 in \$1 accumulate at \$3 \\
defined by \texttt{zeros-accumulate-at} \\
characterized by \texttt{zeros-accumulate-at}  \textit{(definitional)}:
\begin{quote}\ttfamily\footnotesize\raggedright\noindent
forall(\allowbreak{}[u, f, a], zeros-accumulate-at(\allowbreak{}u, f, a) iff forall(\allowbreak{}[acr\_], pos-rr(\allowbreak{}acr\_) implies forsome(\allowbreak{}[acx\_ in u], acx\_ in ball(\allowbreak{}cc-ms, a, acr\_) and not(\allowbreak{}acx\_ = a) and f(\allowbreak{}acx\_) = 0)))
\end{quote}
mentioned by 18 result(s)

\smallskip

\noindent\textbf{\texttt{zkept}}\quad\textit{defined-fn}\index{zkept@\texttt{zkept}}\\
defined by \texttt{zkept-zero} \texttt{zkept-succ} \texttt{zkept-limit} \\
characterized by \texttt{zkept-succ}  \textit{(definitional)}:
\begin{quote}\ttfamily\footnotesize\raggedright\noindent
forall(\allowbreak{}[phi, grd, porel, alpha in ord], zkept(\allowbreak{}phi, grd, porel, succ\_ord(\allowbreak{}alpha)) == union(\allowbreak{}zkept(\allowbreak{}phi, grd, porel, alpha), keep-set(\allowbreak{}phi, grd, porel, zkept(\allowbreak{}phi, grd, porel, alpha), alpha)))
\end{quote}
also: \texttt{zkept-zero} \texttt{zkept-limit} \texttt{zkept-defined} \texttt{zorn-zkept-is-chain} \texttt{zorn-zkept-monotone} \texttt{zorn-zkept-limit-is-chain} \\
mentioned by 11 result(s)

\smallskip

\noindent\textbf{\texttt{zp(qsp\_)}}\quad\textit{functoid}\index{zp(qsp_)@\texttt{zp(qsp\_)}}\\
reads: the ring Z\_\$1 of \$1-adic integers \\
definition:
\begin{quote}\ttfamily\footnotesize\raggedright\noindent
nv-integers(\allowbreak{}qp-field(\allowbreak{}qsp\_))
\end{quote}
characterized by \texttt{ccr-1-26-zp}  \textit{(proven)}:
\begin{quote}\ttfamily\footnotesize\raggedright\noindent
forall(\allowbreak{}[p], zz-prime(\allowbreak{}p) implies zp(\allowbreak{}p) == \{zpv\_ in pts(\allowbreak{}qp(\allowbreak{}p)): (\allowbreak{}qp-abs(\allowbreak{}p))(\allowbreak{}zpv\_) <= 1\})
\end{quote}
also: \texttt{qp-zp-compact} \texttt{qp-zp-closed-ball} \texttt{qp-zp-compact-domain} \texttt{qp-zp-totally-bounded} \\
mentioned by 29 result(s)

\smallskip

\noindent\textbf{\texttt{zpoly-alg(k, n, tm)}}\quad\textit{functoid}\index{zpoly-alg(k, n, tm)@\texttt{zpoly-alg(k, n, tm)}}\\
reads: Z(\$1)[\$3] \\
definition:
\begin{quote}\ttfamily\footnotesize\raggedright\noindent
intersection-of(\allowbreak{}\{zkvs\_ in power(\allowbreak{}mat(\allowbreak{}n, n, carr(\allowbreak{}k))): identmat(\allowbreak{}normed-field-as-commutative-ring(\allowbreak{}k), n) in zkvs\_ and zeromat(\allowbreak{}normed-field-as-commutative-ring(\allowbreak{}k), n, n) in zkvs\_ and forall(\allowbreak{}[zkvx\_ in zkvs\_], matmul(\allowbreak{}normed-field-as-commutative-ring(\allowbreak{}k), zkvx\_, tm) in zkvs\_) and forall(\allowbreak{}[zkvx\_ in zkvs\_, zkvy\_ in zkvs\_], matadd(\allowbreak{}normed-field-as-commutative-ring(\allowbreak{}k), zkvx\_, zkvy\_) in zkvs\_) and forall(\allowbreak{}[zkvx\_ in zkvs\_, zkvc\_ in nv-integers(\allowbreak{}k)], matscale(\allowbreak{}normed-field-as-commutative-ring(\allowbreak{}k), zkvc\_, zkvx\_) in zkvs\_)\})
\end{quote}
mentioned by 16 result(s)

\smallskip

\noindent\textbf{\texttt{zup}}\quad\textit{defined-fn}\index{zup@\texttt{zup}}\\
defined by \texttt{zup-zero} \texttt{zup-succ} \texttt{zup-limit} \\
characterized by \texttt{zup-succ}  \textit{(definitional)}:
\begin{quote}\ttfamily\footnotesize\raggedright\noindent
forall(\allowbreak{}[grd, porel, alpha in ord], zup(\allowbreak{}grd, porel, succ\_ord(\allowbreak{}alpha)) == choice(\allowbreak{}\{y\_ in grd: is-strictly-below(\allowbreak{}porel, zup(\allowbreak{}grd, porel, alpha), y\_)\}))
\end{quote}
also: \texttt{zup-zero} \texttt{zup-limit} \texttt{zup-tower} \texttt{zup-injective} \\
mentioned by 9 result(s)

\smallskip

\noindent\textbf{\texttt{zz-act(g\_, k\_, a\_)}}\quad\textit{functoid}\index{zz-act(g_, k_, a_)@\texttt{zz-act(g\_, k\_, a\_)}}\\
definition:
\begin{quote}\ttfamily\footnotesize\raggedright\noindent
if(\allowbreak{}k\_ in nn, mpow(\allowbreak{}abelian-group-as-monoid(\allowbreak{}g\_), a\_, k\_), (\allowbreak{}inv(\allowbreak{}g\_))(\allowbreak{}mpow(\allowbreak{}abelian-group-as-monoid(\allowbreak{}g\_), a\_, -k\_)))
\end{quote}
characterized by \texttt{zz-act-assoc-ring-additive-ag}  \textit{(definitional)}:
\begin{quote}\ttfamily\footnotesize\raggedright\noindent
forall(\allowbreak{}[r], is-ring(\allowbreak{}r) implies forall(\allowbreak{}[j in zz, k in zz, a in carr(\allowbreak{}r)], zz-act(\allowbreak{}ring-additive-ag(\allowbreak{}r), j * k, a) = zz-act(\allowbreak{}ring-additive-ag(\allowbreak{}r), j, zz-act(\allowbreak{}ring-additive-ag(\allowbreak{}r), k, a))))
\end{quote}
also: \texttt{zz-act-distrib-ring-additive-ag} \texttt{zz-act-assoc-ring-additive-ag-uba-as-ring} \texttt{zz-act-distrib-ring-additive-ag-uba-as-ring} \texttt{zz-act-assoc-ring-additive-ag-uba-as-ring-uiba-as-uba} \texttt{zz-act-distrib-ring-additive-ag-uba-as-ring-uiba-as-uba} \texttt{zz-act-add} \texttt{zz-act-neg} \texttt{zz-act-one} ... (27 in all) \\
mentioned by 112 result(s)

\smallskip

\noindent\textbf{\texttt{zz-add-monoid}}\quad\textit{instance}\index{zz-add-monoid@\texttt{zz-add-monoid}}\\
an instance of comm-monoid; predicate \texttt{is-zz-add-monoid}; same shape as comm-monoid

\smallskip

\noindent\textbf{\texttt{zz-bezout-set(a, b)}}\quad\textit{functoid}\index{zz-bezout-set(a, b)@\texttt{zz-bezout-set(a, b)}}\\
definition:
\begin{quote}\ttfamily\footnotesize\raggedright\noindent
\{z in zz: forsome(\allowbreak{}[x\_ in zz, y\_ in zz], z = x\_ * a + y\_ * b)\}
\end{quote}
characterized by \texttt{zz-bezout-set-membership}  \textit{(definitional)}:
\begin{quote}\ttfamily\footnotesize\raggedright\noindent
forall(\allowbreak{}[a, b, z\_], z\_ in zz-bezout-set(\allowbreak{}a, b) iff z\_ in zz and forsome(\allowbreak{}[x\_ in zz, y\_ in zz], z\_ = x\_ * a + y\_ * b))
\end{quote}
also: \texttt{zz-bezout-set-is-ideal} \\
mentioned by 5 result(s)

\smallskip

\noindent\textbf{\texttt{zz-coprime(a, b)}}\quad\textit{predicate}\index{zz-coprime(a, b)@\texttt{zz-coprime(a, b)}}\\
reads: \$1 and \$2 are coprime \\
defined by \texttt{zz-coprime} \\
characterized by \texttt{zz-coprime}  \textit{(definitional)}:
\begin{quote}\ttfamily\footnotesize\raggedright\noindent
forall(\allowbreak{}[a, b], zz-coprime(\allowbreak{}a, b) iff zz-is-gcd(\allowbreak{}a, b, 1))
\end{quote}
also: \texttt{zz-coprime-mul} \texttt{zz-coprime-sym} \texttt{zz-coprime-bezout} \texttt{zz-coprime-factor} \texttt{zz-coprime-congruent} \texttt{zz-coprime-divides-product} \texttt{zz-coprime-iff-no-common-prime} \\
mentioned by 36 result(s)

\smallskip

\noindent\textbf{\texttt{zz-divides(a, b)}}\quad\textit{predicate}\index{zz-divides(a, b)@\texttt{zz-divides(a, b)}}\\
reads: \$1 divides \$2 \\
defined by \texttt{zz-divides} \\
characterized by \texttt{zz-divides}  \textit{(definitional)}:
\begin{quote}\ttfamily\footnotesize\raggedright\noindent
forall(\allowbreak{}[a, b], zz-divides(\allowbreak{}a, b) iff a in zz and b in zz and forsome(\allowbreak{}[c\_ in zz], b = a * c\_))
\end{quote}
also: \texttt{zz-divides-small} \texttt{zz-divides-trans} \texttt{zz-divides-le-pos} \texttt{zz-divides-antisym} \texttt{zz-divides-lincomb} \texttt{zz-divides-mul-right} \\
mentioned by 119 result(s)

\smallskip

\noindent\textbf{\texttt{zz-is-gcd(a, b, d)}}\quad\textit{predicate}\index{zz-is-gcd(a, b, d)@\texttt{zz-is-gcd(a, b, d)}}\\
reads: \$3 is a greatest common divisor of \$1 and \$2 \\
defined by \texttt{zz-is-gcd} \\
characterized by \texttt{zz-is-gcd}  \textit{(definitional)}:
\begin{quote}\ttfamily\footnotesize\raggedright\noindent
forall(\allowbreak{}[a, b, d], zz-is-gcd(\allowbreak{}a, b, d) iff zz-divides(\allowbreak{}d, a) and zz-divides(\allowbreak{}d, b) and forall(\allowbreak{}[e\_ in zz], zz-divides(\allowbreak{}e\_, a) implies zz-divides(\allowbreak{}e\_, b) implies zz-divides(\allowbreak{}e\_, d)))
\end{quote}
mentioned by 14 result(s)

\smallskip

\noindent\textbf{\texttt{zz-is-gcd-set(s, d)}}\quad\textit{predicate}\index{zz-is-gcd-set(s, d)@\texttt{zz-is-gcd-set(s, d)}}\\
reads: \$2 is a greatest common divisor of the integers in \$1 \\
defined by \texttt{zz-is-gcd-set} \\
characterized by \texttt{zz-is-gcd-set}  \textit{(definitional)}:
\begin{quote}\ttfamily\footnotesize\raggedright\noindent
forall(\allowbreak{}[s, d], zz-is-gcd-set(\allowbreak{}s, d) iff d in zz and forall(\allowbreak{}[gsx\_ in s], zz-divides(\allowbreak{}d, gsx\_)) and forall(\allowbreak{}[gse\_ in zz], forall(\allowbreak{}[gsx\_ in s], zz-divides(\allowbreak{}gse\_, gsx\_)) implies zz-divides(\allowbreak{}gse\_, d)))
\end{quote}
mentioned by 4 result(s)

\smallskip

\noindent\textbf{\texttt{zz-mod(a)}}\quad\textit{functoid}\index{zz-mod(a)@\texttt{zz-mod(a)}}\\
reads: Z/(\$1) \\
definition:
\begin{quote}\ttfamily\footnotesize\raggedright\noindent
quotient-module(\allowbreak{}ring-as-module(\allowbreak{}zz-ring), principal-ideal(\allowbreak{}zz-ring, a))
\end{quote}
characterized by \texttt{zz-mod-vec}  \textit{(proven)}:
\begin{quote}\ttfamily\footnotesize\raggedright\noindent
forall(\allowbreak{}[a in carr(\allowbreak{}zz-ring)], vec(\allowbreak{}zz-mod(\allowbreak{}a)) == quotient(\allowbreak{}qmod-setoid(\allowbreak{}ring-as-module(\allowbreak{}zz-ring), principal-ideal(\allowbreak{}zz-ring, a))))
\end{quote}
also: \texttt{zz-mod-card} \texttt{zz-mod-scal} \texttt{zz-mod-vzero} \texttt{zz-mod-unfold} \texttt{zz-mod-vec-set} \texttt{zz-mod-class-in} \texttt{zz-mod-is-module} \texttt{zz-mod-one-trivial} ... (11 in all) \\
mentioned by 30 result(s)

\smallskip

\noindent\textbf{\texttt{zz-ord-is(cpp\_, cpn\_, cpk\_)}}\quad\textit{predicate}\index{zz-ord-is(cpp_, cpn_, cpk_)@\texttt{zz-ord-is(cpp\_, cpn\_, cpk\_)}}\\
reads: \$3 is the order of \$2 at \$1 \\
defined by \texttt{zz-ord-is} \\
characterized by \texttt{zz-ord-is}  \textit{(definitional)}:
\begin{quote}\ttfamily\footnotesize\raggedright\noindent
forall(\allowbreak{}[cpp\_, cpn\_, cpk\_], zz-ord-is(\allowbreak{}cpp\_, cpn\_, cpk\_) iff zz-divides(\allowbreak{}cpp\_ \textasciicircum{} cpk\_, cpn\_) and not(\allowbreak{}zz-divides(\allowbreak{}cpp\_ \textasciicircum{} succ(\allowbreak{}cpk\_), cpn\_)))
\end{quote}
mentioned by 18 result(s)

\smallskip

\noindent\textbf{\texttt{zz-pow(n)}}\quad\textit{functoid}\index{zz-pow(n)@\texttt{zz-pow(n)}}\\
reads: the free module Z\textasciicircum{}\$1 \\
definition:
\begin{quote}\ttfamily\footnotesize\raggedright\noindent
direct-sum(\allowbreak{}zz-ring, n, vnb-lambda(\allowbreak{}zpi\_, interval(\allowbreak{}1, n), ring-as-module(\allowbreak{}zz-ring)))
\end{quote}
characterized by \texttt{zz-pow-ext}  \textit{(proven)}:
\begin{quote}\ttfamily\footnotesize\raggedright\noindent
forall(\allowbreak{}[n, x, y], n in nn implies x in vec(\allowbreak{}zz-pow(\allowbreak{}n)) implies y in vec(\allowbreak{}zz-pow(\allowbreak{}n)) implies forall(\allowbreak{}[i\_ in interval(\allowbreak{}1, n)], x(\allowbreak{}i\_) = y(\allowbreak{}i\_)) implies x = y)
\end{quote}
also: \texttt{zz-pow-scal} \texttt{zz-pow-coord} \texttt{zz-pow-act-at} \texttt{zz-pow-unfold} \texttt{zz-pow-vec-in} \texttt{zz-pow-vadd-at} \texttt{zz-pow-vec-out} \texttt{zz-pow-vneg-at} ... (24 in all) \\
mentioned by 106 result(s)

\smallskip

\noindent\textbf{\texttt{zz-prime(cpp\_)}}\quad\textit{predicate}\index{zz-prime(cpp_)@\texttt{zz-prime(cpp\_)}}\\
reads: \$1 is a prime number \\
defined by \texttt{zz-prime} \\
characterized by \texttt{zz-prime}  \textit{(definitional)}:
\begin{quote}\ttfamily\footnotesize\raggedright\noindent
forall(\allowbreak{}[cpp\_], zz-prime(\allowbreak{}cpp\_) iff cpp\_ in zz and 1 < cpp\_ and forall(\allowbreak{}[cpd\_ in zz], zz-divides(\allowbreak{}cpd\_, cpp\_) implies cpd\_ = 1 or cpd\_ = -1 or cpd\_ = cpp\_ or cpd\_ = -cpp\_))
\end{quote}
mentioned by 381 result(s)

\smallskip

\noindent\textbf{\texttt{zz-residue(k, m)}}\quad\textit{functoid}\index{zz-residue(k, m)@\texttt{zz-residue(k, m)}}\\
reads: the residue of \$1 modulo \$2 in 1..\$2 \\
definition:
\begin{quote}\ttfamily\footnotesize\raggedright\noindent
choice(\allowbreak{}\{zrr\_ in interval(\allowbreak{}1, m): zz-divides(\allowbreak{}m, k - zrr\_)\})
\end{quote}
characterized by \texttt{zz-residue-in}  \textit{(proven)}:
\begin{quote}\ttfamily\footnotesize\raggedright\noindent
forall(\allowbreak{}[m, k], m in nn implies 1 <= m implies k in zz implies zz-residue(\allowbreak{}k, m) in interval(\allowbreak{}1, m) and zz-divides(\allowbreak{}m, k - zz-residue(\allowbreak{}k, m)))
\end{quote}
also: \texttt{zz-residue-unique} \\
mentioned by 3 result(s)

\smallskip

\noindent\textbf{\texttt{zz-ring}}\quad\textit{instance}\index{zz-ring@\texttt{zz-ring}}\\
an instance of ring; predicate \texttt{is-zz-ring}; same shape as euclidean-ring

\smallskip

\noindent\textbf{\texttt{zz-row(n, x)}}\quad\textit{functoid}\index{zz-row(n, x)@\texttt{zz-row(n, x)}}\\
reads: the coefficient row of \$2 in Z\textasciicircum{}\$1 \\
definition:
\begin{quote}\ttfamily\footnotesize\raggedright\noindent
matof(\allowbreak{}1, n, vnb-lambda(\allowbreak{}[zra\_, zrb\_], cartesian(\allowbreak{}interval(\allowbreak{}1, 1), interval(\allowbreak{}1, n)), x(\allowbreak{}zrb\_)))
\end{quote}
characterized by \texttt{zz-row-type}  \textit{(proven)}:
\begin{quote}\ttfamily\footnotesize\raggedright\noindent
forall(\allowbreak{}[n, x], n in nn implies x in vec(\allowbreak{}zz-pow(\allowbreak{}n)) implies zz-row(\allowbreak{}n, x) in mat(\allowbreak{}1, n, carr(\allowbreak{}zz-ring)))
\end{quote}
also: \texttt{zz-row-entry} \texttt{zz-row-is-ring-row} \texttt{zz-row-type-via-ring-pow} \texttt{zz-row-entry-via-ring-pow} \\
mentioned by 12 result(s)

\smallskip

\noindent\textbf{\texttt{zz-units(n)}}\quad\textit{functoid}\index{zz-units(n)@\texttt{zz-units(n)}}\\
reads: the unit vectors of Z\textasciicircum{}\$1 \\
definition:
\begin{quote}\ttfamily\footnotesize\raggedright\noindent
matof(\allowbreak{}n, 1, vnb-lambda(\allowbreak{}[zua\_, zub\_], cartesian(\allowbreak{}interval(\allowbreak{}1, n), interval(\allowbreak{}1, 1)), vnb-lambda(\allowbreak{}zut\_, interval(\allowbreak{}1, n), entry(\allowbreak{}unitrow(\allowbreak{}zz-ring, n, zua\_), 1, zut\_))))
\end{quote}
characterized by \texttt{zz-units-type}  \textit{(proven)}:
\begin{quote}\ttfamily\footnotesize\raggedright\noindent
forall(\allowbreak{}[n in nn], zz-units(\allowbreak{}n) in mat(\allowbreak{}n, 1, vec(\allowbreak{}zz-pow(\allowbreak{}n))))
\end{quote}
also: \texttt{zz-units-entry} \texttt{zz-units-generate} \texttt{zz-units-rel-free} \texttt{zz-units-is-ring-units} \texttt{zz-units-type-via-ring-pow} \texttt{zz-units-entry-via-ring-pow} \texttt{zz-units-generate-via-ring-pow} \texttt{zz-units-rel-free-via-ring-pow} \\
mentioned by 27 result(s)

\smallskip

\endgroup


%% =========================================================

\begin{thebibliography}{9}
%% Every entry below except Monk1988 is taken from the reference list of
%% F. J. Thayer, "Inteligencia artificial y modelos de lenguaje en la
%% investigacion matematica" (~/docs/article.pdf, 2026), supplied by the author.
%% Bernays1937, Godel1940 and vonNeumann1925 are the BibTeX entries the author supplied in
%% ~/docs/notes-34.txt (2026-09-20).  Farmer2023 (Springer, DOI 10.1007/978-3-031-21112-6) and LittleTheories (DOI 10.1007/3-540-55602-8_192)
%% checked 2026-11-19 against the publisher's records.  Gentzen1935 checked 2026-11-17 against EuDML
%% (eudml.org/doc/168546, /168556); Godel1931 checked 2026-11-17 (DOI 10.1007/BF01700692).  Add no entry that has not been checked against its source.

\bibitem{Bernays1937}
P.~Bernays,
A system of axiomatic set theory---Part~I,
\textit{The Journal of Symbolic Logic}, 2(2):65--77, 1937.

\bibitem{Lean}
L.~de~Moura, S.~Kong, J.~Avigad, F.~van~Doorn, J.~von~Raumer,
The Lean theorem prover (system description),
in \textit{Automated Deduction -- CADE-25}, pages 378--388, Springer, 2015.

\bibitem{IMPS}
W.~M.~Farmer, J.~D.~Guttman, F.~J.~Thayer,
IMPS: an Interactive Mathematical Proof System,
\textit{Journal of Automated Reasoning}, 11:213--248, 1993.

\bibitem{Farmer2023}
W.~M.~Farmer,
\textit{Simple Type Theory: A Practical Logic for Expressing and Reasoning
About Mathematical Ideas},
Computer Science Foundations and Applied Logic, Birkh\"auser, Cham, 2023.

\bibitem{LittleTheories}
W.~M.~Farmer, J.~D.~Guttman, F.~J.~Thayer,
Little theories,
in \textit{Automated Deduction -- CADE-11}, Lecture Notes in Computer
Science~607, pages 567--581, Springer, 1992.

\bibitem{Gentzen1935}
G.~Gentzen,
Untersuchungen \"uber das logische Schlie\ss en~I, II,
\textit{Mathematische Zeitschrift}, 39:176--210 and 405--431, 1935.

\bibitem{Godel1931}
K.~G\"odel,
\"Uber formal unentscheidbare S\"atze der Principia Mathematica und
verwandter Systeme~I,
\textit{Monatshefte f\"ur Mathematik und Physik}, 38:173--198, 1931.

\bibitem{Godel1940}
K.~G\"odel,
\textit{The Consistency of the Axiom of Choice and of the Generalized
Continuum-Hypothesis with the Axioms of Set Theory},
Annals of Mathematics Studies~3, Princeton University Press, Princeton, NJ, 1940.

\bibitem{HOL}
M.~J.~C.~Gordon, T.~F.~Melham (editors),
\textit{Introduction to HOL: A theorem proving environment for higher order logic},
Cambridge University Press, 1993.

\bibitem{MITScheme}
C.~Hanson,
\textit{MIT/GNU Scheme Reference Manual}, 2015.  VNB requires MIT/GNU Scheme 12.1.

\bibitem{Monk1988}
L.~G.~Monk,
Inference rules using local contexts,
\textit{Journal of Automated Reasoning}, 4:445--462, 1988.

\bibitem{Emacs}
R.~Stallman,
\textit{GNU Emacs}, 2024.

\bibitem{Rocq}
The Rocq Development Team,
\textit{The Rocq Prover reference manual}, Inria, 2026.  Version 9.3.

\bibitem{vonNeumann1925}
J.~von~Neumann,
Eine Axiomatisierung der Mengenlehre,
\textit{Journal f\"ur die reine und angewandte Mathematik}, 154:219--240, 1925.
\end{thebibliography}

%% =========================================================
%% The index.  Entries are emitted by the GENERATORS (gen-tactics-ref.py,
%% gen-glossary.py), one per tactic and per glossary name, so the index
%% cannot drift from the registries the way a hand-kept one would.
%% `make' runs makeindex between the two pdflatex passes.
\printindex

\end{document}
